% !TEX root = ../main.tex
The generic term ``barrier options'' refers to a class of options whose payoff depends on whether the underlying price hits a prespecified barrier during the option's lifetime.

\subsection{Knock-out Barrier}
\subsubsection{Up-and-out barrier call option}
Let us consider an up-and-out barrier call option with maturity \(T\), strike price \(K\), barrier (or call level) \(B\), and payoff

\begin{equation*}
	C = {(S_T - K)}^+ \I_{\left\{ \max\limits_{0 \le t \le T} S_t < B \right\}} =
	\begin{dcases}
		S_T - K & \text{if } \max\limits_{0 \le t \le T} S_t \le B, \\
		0       & \text{if } \max\limits_{0 \le t \le T} S_t > B,
	\end{dcases}
\end{equation*}

with \(B \ge K\).

\begin{prop}{Up-and-Out Barrier Call Option Pricing}{Up-and-Out Barrier Call Option Pricing}
	When \(K < B\), the price
	\begin{equation}
		\e^{-(T-t)r} \I_{\{M_0^t < B\}} \E^* \left[ {\left( \frac{S_{T-t}}{S_0} - K \right)}^+ \I_{\left\{ x \max\limits_{0 \le u \le T-t} \frac{S_u}{S_0} < B \right\}} \Bigg|_{x=S_t} \right]
	\end{equation}
	of the up-and-out barrier call option with maturity \(T\), strike price \(K\), and barrier level \(B\) is given by the following formula, where \(\tau = T-t\):
	\begin{align}
		C_{UO}(t) & = S_t \I_{\{M_0^t < B\}} \Bigg\{ \left[ \Phi \left( d_1(S_t, K) \right) - \Phi \left( d_1(S_t, B) \right) \right] \nonumber                                                                                   \\
		          & \quad - {\left( \frac{B}{S_t} \right)}^{1+\frac{2r}{\sigma^2}} \left[ \Phi \left( d_1{\left(\frac{B^2}{S_t}, K\right)} \right) - \Phi \left( d_1{\left(\frac{B^2}{S_t}, B\right)} \right) \right] \Bigg\} \nonumber \\
		          & \quad - \e^{-r\tau} K \I_{\{M_0^t < B\}} \Bigg\{ \left[ \Phi \left( d_2(S_t, K) \right) - \Phi \left( d_2(S_t, B) \right) \right] \nonumber                                                                   \\
		          & \quad - {\left( \frac{S_t}{B} \right)}^{1-\frac{2r}{\sigma^2}} \left[ \Phi \left( d_2{\left(\frac{B^2}{S_t}, K\right)} \right) - \Phi \left( d_2{\left(\frac{B^2}{S_t}, B\right)} \right) \right] \Bigg\},
	\end{align}
	If \(K \ge B\):
	\begin{equation}
		C_{UO}(t, S_t) = 0.
	\end{equation}
	where \(d_{1,2}(S, X)\) corresponds to the definitions in \cref{eq:D1Definition,eq:D2Definition} with spot price \(S\), strike price \(K\), and time to maturity \(\tau\). The price of the up-and-out barrier call option vanishes when \(B \le K\).
\end{prop}

\begin{figure}[H]
	\centering
	%% Creator: Matplotlib, PGF backend
%%
%% To include the figure in your LaTeX document, write
%%   \input{<filename>.pgf}
%%
%% Make sure the required packages are loaded in your preamble
%%   \usepackage{pgf}
%%
%% Also ensure that all the required font packages are loaded; for instance,
%% the lmodern package is sometimes necessary when using math font.
%%   \usepackage{lmodern}
%%
%% Figures using additional raster images can only be included by \input if
%% they are in the same directory as the main LaTeX file. For loading figures
%% from other directories you can use the `import` package
%%   \usepackage{import}
%%
%% and then include the figures with
%%   \import{<path to file>}{<filename>.pgf}
%%
%% Matplotlib used the following preamble
%%   \def\mathdefault#1{#1}
%%   \everymath=\expandafter{\the\everymath\displaystyle}
%%   \IfFileExists{scrextend.sty}{
%%     \usepackage[fontsize=10.000000pt]{scrextend}
%%   }{
%%     \renewcommand{\normalsize}{\fontsize{10.000000}{12.000000}\selectfont}
%%     \normalsize
%%   }
%%   
%%   \makeatletter\@ifpackageloaded{underscore}{}{\usepackage[strings]{underscore}}\makeatother
%%
\begingroup%
\makeatletter%
\begin{pgfpicture}%
\pgfpathrectangle{\pgfpointorigin}{\pgfqpoint{5.288298in}{5.000000in}}%
\pgfusepath{use as bounding box, clip}%
\begin{pgfscope}%
\pgfsetbuttcap%
\pgfsetmiterjoin%
\definecolor{currentfill}{rgb}{1.000000,1.000000,1.000000}%
\pgfsetfillcolor{currentfill}%
\pgfsetlinewidth{0.000000pt}%
\definecolor{currentstroke}{rgb}{1.000000,1.000000,1.000000}%
\pgfsetstrokecolor{currentstroke}%
\pgfsetdash{}{0pt}%
\pgfpathmoveto{\pgfqpoint{0.000000in}{0.000000in}}%
\pgfpathlineto{\pgfqpoint{5.288298in}{0.000000in}}%
\pgfpathlineto{\pgfqpoint{5.288298in}{5.000000in}}%
\pgfpathlineto{\pgfqpoint{0.000000in}{5.000000in}}%
\pgfpathlineto{\pgfqpoint{0.000000in}{0.000000in}}%
\pgfpathclose%
\pgfusepath{fill}%
\end{pgfscope}%
\begin{pgfscope}%
\pgfsetbuttcap%
\pgfsetmiterjoin%
\definecolor{currentfill}{rgb}{1.000000,1.000000,1.000000}%
\pgfsetfillcolor{currentfill}%
\pgfsetlinewidth{0.000000pt}%
\definecolor{currentstroke}{rgb}{0.000000,0.000000,0.000000}%
\pgfsetstrokecolor{currentstroke}%
\pgfsetstrokeopacity{0.000000}%
\pgfsetdash{}{0pt}%
\pgfpathmoveto{\pgfqpoint{0.050000in}{0.050000in}}%
\pgfpathlineto{\pgfqpoint{4.950000in}{0.050000in}}%
\pgfpathlineto{\pgfqpoint{4.950000in}{4.950000in}}%
\pgfpathlineto{\pgfqpoint{0.050000in}{4.950000in}}%
\pgfpathlineto{\pgfqpoint{0.050000in}{0.050000in}}%
\pgfpathclose%
\pgfusepath{fill}%
\end{pgfscope}%
\begin{pgfscope}%
\pgfsetbuttcap%
\pgfsetmiterjoin%
\pgfsetlinewidth{1.003750pt}%
\definecolor{currentstroke}{rgb}{0.950000,0.950000,0.950000}%
\pgfsetstrokecolor{currentstroke}%
\pgfsetstrokeopacity{0.500000}%
\pgfsetdash{}{0pt}%
\pgfpathmoveto{\pgfqpoint{0.438795in}{1.686447in}}%
\pgfpathlineto{\pgfqpoint{2.920118in}{2.434876in}}%
\pgfpathlineto{\pgfqpoint{2.932761in}{4.459308in}}%
\pgfpathlineto{\pgfqpoint{0.358066in}{3.808469in}}%
\pgfusepath{stroke}%
\end{pgfscope}%
\begin{pgfscope}%
\pgfsetbuttcap%
\pgfsetmiterjoin%
\pgfsetlinewidth{1.003750pt}%
\definecolor{currentstroke}{rgb}{0.900000,0.900000,0.900000}%
\pgfsetstrokecolor{currentstroke}%
\pgfsetstrokeopacity{0.500000}%
\pgfsetdash{}{0pt}%
\pgfpathmoveto{\pgfqpoint{2.920118in}{2.434876in}}%
\pgfpathlineto{\pgfqpoint{4.749691in}{1.337846in}}%
\pgfpathlineto{\pgfqpoint{4.834815in}{3.504368in}}%
\pgfpathlineto{\pgfqpoint{2.932761in}{4.459308in}}%
\pgfusepath{stroke}%
\end{pgfscope}%
\begin{pgfscope}%
\pgfsetbuttcap%
\pgfsetmiterjoin%
\pgfsetlinewidth{1.003750pt}%
\definecolor{currentstroke}{rgb}{0.925000,0.925000,0.925000}%
\pgfsetstrokecolor{currentstroke}%
\pgfsetstrokeopacity{0.500000}%
\pgfsetdash{}{0pt}%
\pgfpathmoveto{\pgfqpoint{0.438795in}{1.686447in}}%
\pgfpathlineto{\pgfqpoint{2.155964in}{0.447388in}}%
\pgfpathlineto{\pgfqpoint{4.749691in}{1.337846in}}%
\pgfpathlineto{\pgfqpoint{2.920118in}{2.434876in}}%
\pgfusepath{stroke}%
\end{pgfscope}%
\begin{pgfscope}%
\pgfsetbuttcap%
\pgfsetroundjoin%
\pgfsetlinewidth{0.501875pt}%
\definecolor{currentstroke}{rgb}{0.800000,0.800000,0.800000}%
\pgfsetstrokecolor{currentstroke}%
\pgfsetdash{}{0pt}%
\pgfpathmoveto{\pgfqpoint{0.459298in}{1.671653in}}%
\pgfpathlineto{\pgfqpoint{2.942077in}{2.421709in}}%
\pgfpathlineto{\pgfqpoint{2.955519in}{4.447882in}}%
\pgfusepath{stroke}%
\end{pgfscope}%
\begin{pgfscope}%
\pgfsetbuttcap%
\pgfsetroundjoin%
\pgfsetlinewidth{0.501875pt}%
\definecolor{currentstroke}{rgb}{0.800000,0.800000,0.800000}%
\pgfsetstrokecolor{currentstroke}%
\pgfsetdash{}{0pt}%
\pgfpathmoveto{\pgfqpoint{0.783380in}{1.437805in}}%
\pgfpathlineto{\pgfqpoint{3.288812in}{2.213803in}}%
\pgfpathlineto{\pgfqpoint{3.315100in}{4.267352in}}%
\pgfusepath{stroke}%
\end{pgfscope}%
\begin{pgfscope}%
\pgfsetbuttcap%
\pgfsetroundjoin%
\pgfsetlinewidth{0.501875pt}%
\definecolor{currentstroke}{rgb}{0.800000,0.800000,0.800000}%
\pgfsetstrokecolor{currentstroke}%
\pgfsetdash{}{0pt}%
\pgfpathmoveto{\pgfqpoint{1.119111in}{1.195551in}}%
\pgfpathlineto{\pgfqpoint{3.647281in}{1.998862in}}%
\pgfpathlineto{\pgfqpoint{3.687299in}{4.080487in}}%
\pgfusepath{stroke}%
\end{pgfscope}%
\begin{pgfscope}%
\pgfsetbuttcap%
\pgfsetroundjoin%
\pgfsetlinewidth{0.501875pt}%
\definecolor{currentstroke}{rgb}{0.800000,0.800000,0.800000}%
\pgfsetstrokecolor{currentstroke}%
\pgfsetdash{}{0pt}%
\pgfpathmoveto{\pgfqpoint{1.467131in}{0.944430in}}%
\pgfpathlineto{\pgfqpoint{4.018089in}{1.776522in}}%
\pgfpathlineto{\pgfqpoint{4.072794in}{3.886946in}}%
\pgfusepath{stroke}%
\end{pgfscope}%
\begin{pgfscope}%
\pgfsetbuttcap%
\pgfsetroundjoin%
\pgfsetlinewidth{0.501875pt}%
\definecolor{currentstroke}{rgb}{0.800000,0.800000,0.800000}%
\pgfsetstrokecolor{currentstroke}%
\pgfsetdash{}{0pt}%
\pgfpathmoveto{\pgfqpoint{1.828126in}{0.683946in}}%
\pgfpathlineto{\pgfqpoint{4.401884in}{1.546394in}}%
\pgfpathlineto{\pgfqpoint{4.472308in}{3.686368in}}%
\pgfusepath{stroke}%
\end{pgfscope}%
\begin{pgfscope}%
\pgfsetbuttcap%
\pgfsetroundjoin%
\pgfsetlinewidth{0.501875pt}%
\definecolor{currentstroke}{rgb}{0.800000,0.800000,0.800000}%
\pgfsetstrokecolor{currentstroke}%
\pgfsetdash{}{0pt}%
\pgfpathmoveto{\pgfqpoint{0.539260in}{1.613955in}}%
\pgfpathlineto{\pgfqpoint{3.027693in}{2.370373in}}%
\pgfpathlineto{\pgfqpoint{3.044267in}{4.403325in}}%
\pgfusepath{stroke}%
\end{pgfscope}%
\begin{pgfscope}%
\pgfsetbuttcap%
\pgfsetroundjoin%
\pgfsetlinewidth{0.501875pt}%
\definecolor{currentstroke}{rgb}{0.800000,0.800000,0.800000}%
\pgfsetstrokecolor{currentstroke}%
\pgfsetdash{}{0pt}%
\pgfpathmoveto{\pgfqpoint{0.619921in}{1.555752in}}%
\pgfpathlineto{\pgfqpoint{3.114014in}{2.318614in}}%
\pgfpathlineto{\pgfqpoint{3.133773in}{4.358388in}}%
\pgfusepath{stroke}%
\end{pgfscope}%
\begin{pgfscope}%
\pgfsetbuttcap%
\pgfsetroundjoin%
\pgfsetlinewidth{0.501875pt}%
\definecolor{currentstroke}{rgb}{0.800000,0.800000,0.800000}%
\pgfsetstrokecolor{currentstroke}%
\pgfsetdash{}{0pt}%
\pgfpathmoveto{\pgfqpoint{0.701291in}{1.497037in}}%
\pgfpathlineto{\pgfqpoint{3.201051in}{2.266426in}}%
\pgfpathlineto{\pgfqpoint{3.224047in}{4.313065in}}%
\pgfusepath{stroke}%
\end{pgfscope}%
\begin{pgfscope}%
\pgfsetbuttcap%
\pgfsetroundjoin%
\pgfsetlinewidth{0.501875pt}%
\definecolor{currentstroke}{rgb}{0.800000,0.800000,0.800000}%
\pgfsetstrokecolor{currentstroke}%
\pgfsetdash{}{0pt}%
\pgfpathmoveto{\pgfqpoint{0.866197in}{1.378047in}}%
\pgfpathlineto{\pgfqpoint{3.377306in}{2.160741in}}%
\pgfpathlineto{\pgfqpoint{3.406941in}{4.221242in}}%
\pgfusepath{stroke}%
\end{pgfscope}%
\begin{pgfscope}%
\pgfsetbuttcap%
\pgfsetroundjoin%
\pgfsetlinewidth{0.501875pt}%
\definecolor{currentstroke}{rgb}{0.800000,0.800000,0.800000}%
\pgfsetstrokecolor{currentstroke}%
\pgfsetdash{}{0pt}%
\pgfpathmoveto{\pgfqpoint{0.949751in}{1.317756in}}%
\pgfpathlineto{\pgfqpoint{3.466543in}{2.107234in}}%
\pgfpathlineto{\pgfqpoint{3.499581in}{4.174732in}}%
\pgfusepath{stroke}%
\end{pgfscope}%
\begin{pgfscope}%
\pgfsetbuttcap%
\pgfsetroundjoin%
\pgfsetlinewidth{0.501875pt}%
\definecolor{currentstroke}{rgb}{0.800000,0.800000,0.800000}%
\pgfsetstrokecolor{currentstroke}%
\pgfsetdash{}{0pt}%
\pgfpathmoveto{\pgfqpoint{1.034052in}{1.256927in}}%
\pgfpathlineto{\pgfqpoint{3.556531in}{2.053276in}}%
\pgfpathlineto{\pgfqpoint{3.593030in}{4.127815in}}%
\pgfusepath{stroke}%
\end{pgfscope}%
\begin{pgfscope}%
\pgfsetbuttcap%
\pgfsetroundjoin%
\pgfsetlinewidth{0.501875pt}%
\definecolor{currentstroke}{rgb}{0.800000,0.800000,0.800000}%
\pgfsetstrokecolor{currentstroke}%
\pgfsetdash{}{0pt}%
\pgfpathmoveto{\pgfqpoint{1.204938in}{1.133621in}}%
\pgfpathlineto{\pgfqpoint{3.738801in}{1.943985in}}%
\pgfpathlineto{\pgfqpoint{3.782399in}{4.032741in}}%
\pgfusepath{stroke}%
\end{pgfscope}%
\begin{pgfscope}%
\pgfsetbuttcap%
\pgfsetroundjoin%
\pgfsetlinewidth{0.501875pt}%
\definecolor{currentstroke}{rgb}{0.800000,0.800000,0.800000}%
\pgfsetstrokecolor{currentstroke}%
\pgfsetdash{}{0pt}%
\pgfpathmoveto{\pgfqpoint{1.291543in}{1.071129in}}%
\pgfpathlineto{\pgfqpoint{3.831103in}{1.888641in}}%
\pgfpathlineto{\pgfqpoint{3.878340in}{3.984573in}}%
\pgfusepath{stroke}%
\end{pgfscope}%
\begin{pgfscope}%
\pgfsetbuttcap%
\pgfsetroundjoin%
\pgfsetlinewidth{0.501875pt}%
\definecolor{currentstroke}{rgb}{0.800000,0.800000,0.800000}%
\pgfsetstrokecolor{currentstroke}%
\pgfsetdash{}{0pt}%
\pgfpathmoveto{\pgfqpoint{1.378937in}{1.008068in}}%
\pgfpathlineto{\pgfqpoint{3.924195in}{1.832821in}}%
\pgfpathlineto{\pgfqpoint{3.975135in}{3.935977in}}%
\pgfusepath{stroke}%
\end{pgfscope}%
\begin{pgfscope}%
\pgfsetbuttcap%
\pgfsetroundjoin%
\pgfsetlinewidth{0.501875pt}%
\definecolor{currentstroke}{rgb}{0.800000,0.800000,0.800000}%
\pgfsetstrokecolor{currentstroke}%
\pgfsetdash{}{0pt}%
\pgfpathmoveto{\pgfqpoint{1.556135in}{0.880207in}}%
\pgfpathlineto{\pgfqpoint{4.112794in}{1.719736in}}%
\pgfpathlineto{\pgfqpoint{4.171328in}{3.837476in}}%
\pgfusepath{stroke}%
\end{pgfscope}%
\begin{pgfscope}%
\pgfsetbuttcap%
\pgfsetroundjoin%
\pgfsetlinewidth{0.501875pt}%
\definecolor{currentstroke}{rgb}{0.800000,0.800000,0.800000}%
\pgfsetstrokecolor{currentstroke}%
\pgfsetdash{}{0pt}%
\pgfpathmoveto{\pgfqpoint{1.645962in}{0.815391in}}%
\pgfpathlineto{\pgfqpoint{4.208321in}{1.662457in}}%
\pgfpathlineto{\pgfqpoint{4.270751in}{3.787561in}}%
\pgfusepath{stroke}%
\end{pgfscope}%
\begin{pgfscope}%
\pgfsetbuttcap%
\pgfsetroundjoin%
\pgfsetlinewidth{0.501875pt}%
\definecolor{currentstroke}{rgb}{0.800000,0.800000,0.800000}%
\pgfsetstrokecolor{currentstroke}%
\pgfsetdash{}{0pt}%
\pgfpathmoveto{\pgfqpoint{1.736622in}{0.749974in}}%
\pgfpathlineto{\pgfqpoint{4.304681in}{1.604679in}}%
\pgfpathlineto{\pgfqpoint{4.371073in}{3.737193in}}%
\pgfusepath{stroke}%
\end{pgfscope}%
\begin{pgfscope}%
\pgfsetbuttcap%
\pgfsetroundjoin%
\pgfsetlinewidth{0.501875pt}%
\definecolor{currentstroke}{rgb}{0.800000,0.800000,0.800000}%
\pgfsetstrokecolor{currentstroke}%
\pgfsetdash{}{0pt}%
\pgfpathmoveto{\pgfqpoint{1.920488in}{0.617301in}}%
\pgfpathlineto{\pgfqpoint{4.499943in}{1.487597in}}%
\pgfpathlineto{\pgfqpoint{4.574467in}{3.635078in}}%
\pgfusepath{stroke}%
\end{pgfscope}%
\begin{pgfscope}%
\pgfsetbuttcap%
\pgfsetroundjoin%
\pgfsetlinewidth{0.501875pt}%
\definecolor{currentstroke}{rgb}{0.800000,0.800000,0.800000}%
\pgfsetstrokecolor{currentstroke}%
\pgfsetdash{}{0pt}%
\pgfpathmoveto{\pgfqpoint{2.013719in}{0.550028in}}%
\pgfpathlineto{\pgfqpoint{4.598868in}{1.428281in}}%
\pgfpathlineto{\pgfqpoint{4.677564in}{3.583317in}}%
\pgfusepath{stroke}%
\end{pgfscope}%
\begin{pgfscope}%
\pgfsetbuttcap%
\pgfsetroundjoin%
\pgfsetlinewidth{0.501875pt}%
\definecolor{currentstroke}{rgb}{0.800000,0.800000,0.800000}%
\pgfsetstrokecolor{currentstroke}%
\pgfsetdash{}{0pt}%
\pgfpathmoveto{\pgfqpoint{2.107831in}{0.482120in}}%
\pgfpathlineto{\pgfqpoint{4.698670in}{1.368439in}}%
\pgfpathlineto{\pgfqpoint{4.781611in}{3.531080in}}%
\pgfusepath{stroke}%
\end{pgfscope}%
\begin{pgfscope}%
\pgfsetbuttcap%
\pgfsetroundjoin%
\pgfsetlinewidth{0.501875pt}%
\definecolor{currentstroke}{rgb}{0.800000,0.800000,0.800000}%
\pgfsetstrokecolor{currentstroke}%
\pgfsetdash{}{0pt}%
\pgfpathmoveto{\pgfqpoint{0.361374in}{3.809305in}}%
\pgfpathlineto{\pgfqpoint{0.441976in}{1.687406in}}%
\pgfpathlineto{\pgfqpoint{2.159307in}{0.448536in}}%
\pgfusepath{stroke}%
\end{pgfscope}%
\begin{pgfscope}%
\pgfsetbuttcap%
\pgfsetroundjoin%
\pgfsetlinewidth{0.501875pt}%
\definecolor{currentstroke}{rgb}{0.800000,0.800000,0.800000}%
\pgfsetstrokecolor{currentstroke}%
\pgfsetdash{}{0pt}%
\pgfpathmoveto{\pgfqpoint{0.867596in}{3.937269in}}%
\pgfpathlineto{\pgfqpoint{0.929003in}{1.834306in}}%
\pgfpathlineto{\pgfqpoint{2.670554in}{0.624053in}}%
\pgfusepath{stroke}%
\end{pgfscope}%
\begin{pgfscope}%
\pgfsetbuttcap%
\pgfsetroundjoin%
\pgfsetlinewidth{0.501875pt}%
\definecolor{currentstroke}{rgb}{0.800000,0.800000,0.800000}%
\pgfsetstrokecolor{currentstroke}%
\pgfsetdash{}{0pt}%
\pgfpathmoveto{\pgfqpoint{1.362288in}{4.062319in}}%
\pgfpathlineto{\pgfqpoint{1.405333in}{1.977979in}}%
\pgfpathlineto{\pgfqpoint{3.169547in}{0.795364in}}%
\pgfusepath{stroke}%
\end{pgfscope}%
\begin{pgfscope}%
\pgfsetbuttcap%
\pgfsetroundjoin%
\pgfsetlinewidth{0.501875pt}%
\definecolor{currentstroke}{rgb}{0.800000,0.800000,0.800000}%
\pgfsetstrokecolor{currentstroke}%
\pgfsetdash{}{0pt}%
\pgfpathmoveto{\pgfqpoint{1.845838in}{4.184552in}}%
\pgfpathlineto{\pgfqpoint{1.871315in}{2.118531in}}%
\pgfpathlineto{\pgfqpoint{3.656720in}{0.962616in}}%
\pgfusepath{stroke}%
\end{pgfscope}%
\begin{pgfscope}%
\pgfsetbuttcap%
\pgfsetroundjoin%
\pgfsetlinewidth{0.501875pt}%
\definecolor{currentstroke}{rgb}{0.800000,0.800000,0.800000}%
\pgfsetstrokecolor{currentstroke}%
\pgfsetdash{}{0pt}%
\pgfpathmoveto{\pgfqpoint{2.318620in}{4.304063in}}%
\pgfpathlineto{\pgfqpoint{2.327282in}{2.256062in}}%
\pgfpathlineto{\pgfqpoint{4.132488in}{1.125953in}}%
\pgfusepath{stroke}%
\end{pgfscope}%
\begin{pgfscope}%
\pgfsetbuttcap%
\pgfsetroundjoin%
\pgfsetlinewidth{0.501875pt}%
\definecolor{currentstroke}{rgb}{0.800000,0.800000,0.800000}%
\pgfsetstrokecolor{currentstroke}%
\pgfsetdash{}{0pt}%
\pgfpathmoveto{\pgfqpoint{2.780988in}{4.420942in}}%
\pgfpathlineto{\pgfqpoint{2.773554in}{2.390668in}}%
\pgfpathlineto{\pgfqpoint{4.597247in}{1.285510in}}%
\pgfusepath{stroke}%
\end{pgfscope}%
\begin{pgfscope}%
\pgfsetbuttcap%
\pgfsetroundjoin%
\pgfsetlinewidth{0.501875pt}%
\definecolor{currentstroke}{rgb}{0.800000,0.800000,0.800000}%
\pgfsetstrokecolor{currentstroke}%
\pgfsetdash{}{0pt}%
\pgfpathmoveto{\pgfqpoint{0.489032in}{3.841575in}}%
\pgfpathlineto{\pgfqpoint{0.564755in}{1.724440in}}%
\pgfpathlineto{\pgfqpoint{2.288292in}{0.492818in}}%
\pgfusepath{stroke}%
\end{pgfscope}%
\begin{pgfscope}%
\pgfsetbuttcap%
\pgfsetroundjoin%
\pgfsetlinewidth{0.501875pt}%
\definecolor{currentstroke}{rgb}{0.800000,0.800000,0.800000}%
\pgfsetstrokecolor{currentstroke}%
\pgfsetdash{}{0pt}%
\pgfpathmoveto{\pgfqpoint{0.615951in}{3.873658in}}%
\pgfpathlineto{\pgfqpoint{0.686849in}{1.761266in}}%
\pgfpathlineto{\pgfqpoint{2.416491in}{0.536830in}}%
\pgfusepath{stroke}%
\end{pgfscope}%
\begin{pgfscope}%
\pgfsetbuttcap%
\pgfsetroundjoin%
\pgfsetlinewidth{0.501875pt}%
\definecolor{currentstroke}{rgb}{0.800000,0.800000,0.800000}%
\pgfsetstrokecolor{currentstroke}%
\pgfsetdash{}{0pt}%
\pgfpathmoveto{\pgfqpoint{0.742137in}{3.905555in}}%
\pgfpathlineto{\pgfqpoint{0.808263in}{1.797888in}}%
\pgfpathlineto{\pgfqpoint{2.543909in}{0.580574in}}%
\pgfusepath{stroke}%
\end{pgfscope}%
\begin{pgfscope}%
\pgfsetbuttcap%
\pgfsetroundjoin%
\pgfsetlinewidth{0.501875pt}%
\definecolor{currentstroke}{rgb}{0.800000,0.800000,0.800000}%
\pgfsetstrokecolor{currentstroke}%
\pgfsetdash{}{0pt}%
\pgfpathmoveto{\pgfqpoint{0.992334in}{3.968801in}}%
\pgfpathlineto{\pgfqpoint{1.049075in}{1.870522in}}%
\pgfpathlineto{\pgfqpoint{2.796434in}{0.667269in}}%
\pgfusepath{stroke}%
\end{pgfscope}%
\begin{pgfscope}%
\pgfsetbuttcap%
\pgfsetroundjoin%
\pgfsetlinewidth{0.501875pt}%
\definecolor{currentstroke}{rgb}{0.800000,0.800000,0.800000}%
\pgfsetstrokecolor{currentstroke}%
\pgfsetdash{}{0pt}%
\pgfpathmoveto{\pgfqpoint{1.116358in}{4.000152in}}%
\pgfpathlineto{\pgfqpoint{1.168483in}{1.906539in}}%
\pgfpathlineto{\pgfqpoint{2.921555in}{0.710225in}}%
\pgfusepath{stroke}%
\end{pgfscope}%
\begin{pgfscope}%
\pgfsetbuttcap%
\pgfsetroundjoin%
\pgfsetlinewidth{0.501875pt}%
\definecolor{currentstroke}{rgb}{0.800000,0.800000,0.800000}%
\pgfsetstrokecolor{currentstroke}%
\pgfsetdash{}{0pt}%
\pgfpathmoveto{\pgfqpoint{1.239674in}{4.031324in}}%
\pgfpathlineto{\pgfqpoint{1.287234in}{1.942357in}}%
\pgfpathlineto{\pgfqpoint{3.045924in}{0.752922in}}%
\pgfusepath{stroke}%
\end{pgfscope}%
\begin{pgfscope}%
\pgfsetbuttcap%
\pgfsetroundjoin%
\pgfsetlinewidth{0.501875pt}%
\definecolor{currentstroke}{rgb}{0.800000,0.800000,0.800000}%
\pgfsetstrokecolor{currentstroke}%
\pgfsetdash{}{0pt}%
\pgfpathmoveto{\pgfqpoint{1.484205in}{4.093137in}}%
\pgfpathlineto{\pgfqpoint{1.522785in}{2.013405in}}%
\pgfpathlineto{\pgfqpoint{3.292432in}{0.837551in}}%
\pgfusepath{stroke}%
\end{pgfscope}%
\begin{pgfscope}%
\pgfsetbuttcap%
\pgfsetroundjoin%
\pgfsetlinewidth{0.501875pt}%
\definecolor{currentstroke}{rgb}{0.800000,0.800000,0.800000}%
\pgfsetstrokecolor{currentstroke}%
\pgfsetdash{}{0pt}%
\pgfpathmoveto{\pgfqpoint{1.605432in}{4.123782in}}%
\pgfpathlineto{\pgfqpoint{1.639596in}{2.048638in}}%
\pgfpathlineto{\pgfqpoint{3.414584in}{0.879488in}}%
\pgfusepath{stroke}%
\end{pgfscope}%
\begin{pgfscope}%
\pgfsetbuttcap%
\pgfsetroundjoin%
\pgfsetlinewidth{0.501875pt}%
\definecolor{currentstroke}{rgb}{0.800000,0.800000,0.800000}%
\pgfsetstrokecolor{currentstroke}%
\pgfsetdash{}{0pt}%
\pgfpathmoveto{\pgfqpoint{1.725974in}{4.154253in}}%
\pgfpathlineto{\pgfqpoint{1.755771in}{2.083680in}}%
\pgfpathlineto{\pgfqpoint{3.536012in}{0.921175in}}%
\pgfusepath{stroke}%
\end{pgfscope}%
\begin{pgfscope}%
\pgfsetbuttcap%
\pgfsetroundjoin%
\pgfsetlinewidth{0.501875pt}%
\definecolor{currentstroke}{rgb}{0.800000,0.800000,0.800000}%
\pgfsetstrokecolor{currentstroke}%
\pgfsetdash{}{0pt}%
\pgfpathmoveto{\pgfqpoint{1.965029in}{4.214682in}}%
\pgfpathlineto{\pgfqpoint{1.986233in}{2.153193in}}%
\pgfpathlineto{\pgfqpoint{3.776715in}{1.003812in}}%
\pgfusepath{stroke}%
\end{pgfscope}%
\begin{pgfscope}%
\pgfsetbuttcap%
\pgfsetroundjoin%
\pgfsetlinewidth{0.501875pt}%
\definecolor{currentstroke}{rgb}{0.800000,0.800000,0.800000}%
\pgfsetstrokecolor{currentstroke}%
\pgfsetdash{}{0pt}%
\pgfpathmoveto{\pgfqpoint{2.083552in}{4.244642in}}%
\pgfpathlineto{\pgfqpoint{2.100530in}{2.187668in}}%
\pgfpathlineto{\pgfqpoint{3.896004in}{1.044765in}}%
\pgfusepath{stroke}%
\end{pgfscope}%
\begin{pgfscope}%
\pgfsetbuttcap%
\pgfsetroundjoin%
\pgfsetlinewidth{0.501875pt}%
\definecolor{currentstroke}{rgb}{0.800000,0.800000,0.800000}%
\pgfsetstrokecolor{currentstroke}%
\pgfsetdash{}{0pt}%
\pgfpathmoveto{\pgfqpoint{2.201414in}{4.274436in}}%
\pgfpathlineto{\pgfqpoint{2.214211in}{2.221957in}}%
\pgfpathlineto{\pgfqpoint{4.014593in}{1.085478in}}%
\pgfusepath{stroke}%
\end{pgfscope}%
\begin{pgfscope}%
\pgfsetbuttcap%
\pgfsetroundjoin%
\pgfsetlinewidth{0.501875pt}%
\definecolor{currentstroke}{rgb}{0.800000,0.800000,0.800000}%
\pgfsetstrokecolor{currentstroke}%
\pgfsetdash{}{0pt}%
\pgfpathmoveto{\pgfqpoint{2.435174in}{4.333526in}}%
\pgfpathlineto{\pgfqpoint{2.439747in}{2.289984in}}%
\pgfpathlineto{\pgfqpoint{4.249694in}{1.166191in}}%
\pgfusepath{stroke}%
\end{pgfscope}%
\begin{pgfscope}%
\pgfsetbuttcap%
\pgfsetroundjoin%
\pgfsetlinewidth{0.501875pt}%
\definecolor{currentstroke}{rgb}{0.800000,0.800000,0.800000}%
\pgfsetstrokecolor{currentstroke}%
\pgfsetdash{}{0pt}%
\pgfpathmoveto{\pgfqpoint{2.551084in}{4.362826in}}%
\pgfpathlineto{\pgfqpoint{2.551610in}{2.323725in}}%
\pgfpathlineto{\pgfqpoint{4.366219in}{1.206196in}}%
\pgfusepath{stroke}%
\end{pgfscope}%
\begin{pgfscope}%
\pgfsetbuttcap%
\pgfsetroundjoin%
\pgfsetlinewidth{0.501875pt}%
\definecolor{currentstroke}{rgb}{0.800000,0.800000,0.800000}%
\pgfsetstrokecolor{currentstroke}%
\pgfsetdash{}{0pt}%
\pgfpathmoveto{\pgfqpoint{2.666353in}{4.391964in}}%
\pgfpathlineto{\pgfqpoint{2.662878in}{2.357286in}}%
\pgfpathlineto{\pgfqpoint{4.482068in}{1.245968in}}%
\pgfusepath{stroke}%
\end{pgfscope}%
\begin{pgfscope}%
\pgfsetbuttcap%
\pgfsetroundjoin%
\pgfsetlinewidth{0.501875pt}%
\definecolor{currentstroke}{rgb}{0.800000,0.800000,0.800000}%
\pgfsetstrokecolor{currentstroke}%
\pgfsetdash{}{0pt}%
\pgfpathmoveto{\pgfqpoint{2.894993in}{4.449761in}}%
\pgfpathlineto{\pgfqpoint{2.883643in}{2.423874in}}%
\pgfpathlineto{\pgfqpoint{4.711761in}{1.324824in}}%
\pgfusepath{stroke}%
\end{pgfscope}%
\begin{pgfscope}%
\pgfsetbuttcap%
\pgfsetroundjoin%
\pgfsetlinewidth{0.501875pt}%
\definecolor{currentstroke}{rgb}{0.800000,0.800000,0.800000}%
\pgfsetstrokecolor{currentstroke}%
\pgfsetdash{}{0pt}%
\pgfpathmoveto{\pgfqpoint{4.751331in}{1.379582in}}%
\pgfpathlineto{\pgfqpoint{2.920363in}{2.473995in}}%
\pgfpathlineto{\pgfqpoint{0.437238in}{1.727366in}}%
\pgfusepath{stroke}%
\end{pgfscope}%
\begin{pgfscope}%
\pgfsetbuttcap%
\pgfsetroundjoin%
\pgfsetlinewidth{0.501875pt}%
\definecolor{currentstroke}{rgb}{0.800000,0.800000,0.800000}%
\pgfsetstrokecolor{currentstroke}%
\pgfsetdash{}{0pt}%
\pgfpathmoveto{\pgfqpoint{4.766142in}{1.756541in}}%
\pgfpathlineto{\pgfqpoint{2.922568in}{2.827104in}}%
\pgfpathlineto{\pgfqpoint{0.423181in}{2.096872in}}%
\pgfusepath{stroke}%
\end{pgfscope}%
\begin{pgfscope}%
\pgfsetbuttcap%
\pgfsetroundjoin%
\pgfsetlinewidth{0.501875pt}%
\definecolor{currentstroke}{rgb}{0.800000,0.800000,0.800000}%
\pgfsetstrokecolor{currentstroke}%
\pgfsetdash{}{0pt}%
\pgfpathmoveto{\pgfqpoint{4.781155in}{2.138644in}}%
\pgfpathlineto{\pgfqpoint{2.924801in}{3.184638in}}%
\pgfpathlineto{\pgfqpoint{0.408937in}{2.471291in}}%
\pgfusepath{stroke}%
\end{pgfscope}%
\begin{pgfscope}%
\pgfsetbuttcap%
\pgfsetroundjoin%
\pgfsetlinewidth{0.501875pt}%
\definecolor{currentstroke}{rgb}{0.800000,0.800000,0.800000}%
\pgfsetstrokecolor{currentstroke}%
\pgfsetdash{}{0pt}%
\pgfpathmoveto{\pgfqpoint{4.796374in}{2.525999in}}%
\pgfpathlineto{\pgfqpoint{2.927062in}{3.546681in}}%
\pgfpathlineto{\pgfqpoint{0.394502in}{2.850721in}}%
\pgfusepath{stroke}%
\end{pgfscope}%
\begin{pgfscope}%
\pgfsetbuttcap%
\pgfsetroundjoin%
\pgfsetlinewidth{0.501875pt}%
\definecolor{currentstroke}{rgb}{0.800000,0.800000,0.800000}%
\pgfsetstrokecolor{currentstroke}%
\pgfsetdash{}{0pt}%
\pgfpathmoveto{\pgfqpoint{4.811805in}{2.918714in}}%
\pgfpathlineto{\pgfqpoint{2.929351in}{3.913318in}}%
\pgfpathlineto{\pgfqpoint{0.379873in}{3.235262in}}%
\pgfusepath{stroke}%
\end{pgfscope}%
\begin{pgfscope}%
\pgfsetbuttcap%
\pgfsetroundjoin%
\pgfsetlinewidth{0.501875pt}%
\definecolor{currentstroke}{rgb}{0.800000,0.800000,0.800000}%
\pgfsetstrokecolor{currentstroke}%
\pgfsetdash{}{0pt}%
\pgfpathmoveto{\pgfqpoint{4.827450in}{3.316901in}}%
\pgfpathlineto{\pgfqpoint{2.931670in}{4.284638in}}%
\pgfpathlineto{\pgfqpoint{0.365045in}{3.625020in}}%
\pgfusepath{stroke}%
\end{pgfscope}%
\begin{pgfscope}%
\pgfsetrectcap%
\pgfsetroundjoin%
\pgfsetlinewidth{0.501875pt}%
\definecolor{currentstroke}{rgb}{0.000000,0.000000,0.000000}%
\pgfsetstrokecolor{currentstroke}%
\pgfsetdash{}{0pt}%
\pgfpathmoveto{\pgfqpoint{0.438795in}{1.686447in}}%
\pgfpathlineto{\pgfqpoint{2.155964in}{0.447388in}}%
\pgfusepath{stroke}%
\end{pgfscope}%
\begin{pgfscope}%
\pgfsetrectcap%
\pgfsetroundjoin%
\pgfsetlinewidth{0.501875pt}%
\definecolor{currentstroke}{rgb}{0.000000,0.000000,0.000000}%
\pgfsetstrokecolor{currentstroke}%
\pgfsetdash{}{0pt}%
\pgfpathmoveto{\pgfqpoint{0.480342in}{1.678010in}}%
\pgfpathlineto{\pgfqpoint{0.417148in}{1.658919in}}%
\pgfusepath{stroke}%
\end{pgfscope}%
\begin{pgfscope}%
\definecolor{textcolor}{rgb}{0.000000,0.000000,0.000000}%
\pgfsetstrokecolor{textcolor}%
\pgfsetfillcolor{textcolor}%
\pgftext[x=0.299177in,y=1.493571in,,top]{\color{textcolor}{\rmfamily\fontsize{10.000000}{12.000000}\selectfont\catcode`\^=\active\def^{\ifmmode\sp\else\^{}\fi}\catcode`\%=\active\def%{\%}$\mathdefault{0}$}}%
\end{pgfscope}%
\begin{pgfscope}%
\pgfsetrectcap%
\pgfsetroundjoin%
\pgfsetlinewidth{0.501875pt}%
\definecolor{currentstroke}{rgb}{0.000000,0.000000,0.000000}%
\pgfsetstrokecolor{currentstroke}%
\pgfsetdash{}{0pt}%
\pgfpathmoveto{\pgfqpoint{0.804637in}{1.444389in}}%
\pgfpathlineto{\pgfqpoint{0.740803in}{1.424617in}}%
\pgfusepath{stroke}%
\end{pgfscope}%
\begin{pgfscope}%
\definecolor{textcolor}{rgb}{0.000000,0.000000,0.000000}%
\pgfsetstrokecolor{textcolor}%
\pgfsetfillcolor{textcolor}%
\pgftext[x=0.620912in,y=1.256225in,,top]{\color{textcolor}{\rmfamily\fontsize{10.000000}{12.000000}\selectfont\catcode`\^=\active\def^{\ifmmode\sp\else\^{}\fi}\catcode`\%=\active\def%{\%}$\mathdefault{80}$}}%
\end{pgfscope}%
\begin{pgfscope}%
\pgfsetrectcap%
\pgfsetroundjoin%
\pgfsetlinewidth{0.501875pt}%
\definecolor{currentstroke}{rgb}{0.000000,0.000000,0.000000}%
\pgfsetstrokecolor{currentstroke}%
\pgfsetdash{}{0pt}%
\pgfpathmoveto{\pgfqpoint{1.140584in}{1.202374in}}%
\pgfpathlineto{\pgfqpoint{1.076102in}{1.181885in}}%
\pgfusepath{stroke}%
\end{pgfscope}%
\begin{pgfscope}%
\definecolor{textcolor}{rgb}{0.000000,0.000000,0.000000}%
\pgfsetstrokecolor{textcolor}%
\pgfsetfillcolor{textcolor}%
\pgftext[x=0.954230in,y=1.010334in,,top]{\color{textcolor}{\rmfamily\fontsize{10.000000}{12.000000}\selectfont\catcode`\^=\active\def^{\ifmmode\sp\else\^{}\fi}\catcode`\%=\active\def%{\%}$\mathdefault{160}$}}%
\end{pgfscope}%
\begin{pgfscope}%
\pgfsetrectcap%
\pgfsetroundjoin%
\pgfsetlinewidth{0.501875pt}%
\definecolor{currentstroke}{rgb}{0.000000,0.000000,0.000000}%
\pgfsetstrokecolor{currentstroke}%
\pgfsetdash{}{0pt}%
\pgfpathmoveto{\pgfqpoint{1.488820in}{0.951505in}}%
\pgfpathlineto{\pgfqpoint{1.423687in}{0.930259in}}%
\pgfusepath{stroke}%
\end{pgfscope}%
\begin{pgfscope}%
\definecolor{textcolor}{rgb}{0.000000,0.000000,0.000000}%
\pgfsetstrokecolor{textcolor}%
\pgfsetfillcolor{textcolor}%
\pgftext[x=1.299767in,y=0.755429in,,top]{\color{textcolor}{\rmfamily\fontsize{10.000000}{12.000000}\selectfont\catcode`\^=\active\def^{\ifmmode\sp\else\^{}\fi}\catcode`\%=\active\def%{\%}$\mathdefault{240}$}}%
\end{pgfscope}%
\begin{pgfscope}%
\pgfsetrectcap%
\pgfsetroundjoin%
\pgfsetlinewidth{0.501875pt}%
\definecolor{currentstroke}{rgb}{0.000000,0.000000,0.000000}%
\pgfsetstrokecolor{currentstroke}%
\pgfsetdash{}{0pt}%
\pgfpathmoveto{\pgfqpoint{1.850033in}{0.691287in}}%
\pgfpathlineto{\pgfqpoint{1.784245in}{0.669242in}}%
\pgfusepath{stroke}%
\end{pgfscope}%
\begin{pgfscope}%
\definecolor{textcolor}{rgb}{0.000000,0.000000,0.000000}%
\pgfsetstrokecolor{textcolor}%
\pgfsetfillcolor{textcolor}%
\pgftext[x=1.658207in,y=0.491006in,,top]{\color{textcolor}{\rmfamily\fontsize{10.000000}{12.000000}\selectfont\catcode`\^=\active\def^{\ifmmode\sp\else\^{}\fi}\catcode`\%=\active\def%{\%}$\mathdefault{320}$}}%
\end{pgfscope}%
\begin{pgfscope}%
\pgfsetrectcap%
\pgfsetroundjoin%
\pgfsetlinewidth{0.501875pt}%
\definecolor{currentstroke}{rgb}{0.000000,0.000000,0.000000}%
\pgfsetstrokecolor{currentstroke}%
\pgfsetdash{}{0pt}%
\pgfpathmoveto{\pgfqpoint{0.560357in}{1.620368in}}%
\pgfpathlineto{\pgfqpoint{0.497004in}{1.601110in}}%
\pgfusepath{stroke}%
\end{pgfscope}%
\begin{pgfscope}%
\pgfsetrectcap%
\pgfsetroundjoin%
\pgfsetlinewidth{0.501875pt}%
\definecolor{currentstroke}{rgb}{0.000000,0.000000,0.000000}%
\pgfsetstrokecolor{currentstroke}%
\pgfsetdash{}{0pt}%
\pgfpathmoveto{\pgfqpoint{0.641072in}{1.562221in}}%
\pgfpathlineto{\pgfqpoint{0.577558in}{1.542794in}}%
\pgfusepath{stroke}%
\end{pgfscope}%
\begin{pgfscope}%
\pgfsetrectcap%
\pgfsetroundjoin%
\pgfsetlinewidth{0.501875pt}%
\definecolor{currentstroke}{rgb}{0.000000,0.000000,0.000000}%
\pgfsetstrokecolor{currentstroke}%
\pgfsetdash{}{0pt}%
\pgfpathmoveto{\pgfqpoint{0.722495in}{1.503564in}}%
\pgfpathlineto{\pgfqpoint{0.658821in}{1.483966in}}%
\pgfusepath{stroke}%
\end{pgfscope}%
\begin{pgfscope}%
\pgfsetrectcap%
\pgfsetroundjoin%
\pgfsetlinewidth{0.501875pt}%
\definecolor{currentstroke}{rgb}{0.000000,0.000000,0.000000}%
\pgfsetstrokecolor{currentstroke}%
\pgfsetdash{}{0pt}%
\pgfpathmoveto{\pgfqpoint{0.887508in}{1.384689in}}%
\pgfpathlineto{\pgfqpoint{0.823512in}{1.364742in}}%
\pgfusepath{stroke}%
\end{pgfscope}%
\begin{pgfscope}%
\pgfsetrectcap%
\pgfsetroundjoin%
\pgfsetlinewidth{0.501875pt}%
\definecolor{currentstroke}{rgb}{0.000000,0.000000,0.000000}%
\pgfsetstrokecolor{currentstroke}%
\pgfsetdash{}{0pt}%
\pgfpathmoveto{\pgfqpoint{0.971115in}{1.324458in}}%
\pgfpathlineto{\pgfqpoint{0.906958in}{1.304333in}}%
\pgfusepath{stroke}%
\end{pgfscope}%
\begin{pgfscope}%
\pgfsetrectcap%
\pgfsetroundjoin%
\pgfsetlinewidth{0.501875pt}%
\definecolor{currentstroke}{rgb}{0.000000,0.000000,0.000000}%
\pgfsetstrokecolor{currentstroke}%
\pgfsetdash{}{0pt}%
\pgfpathmoveto{\pgfqpoint{1.055471in}{1.263689in}}%
\pgfpathlineto{\pgfqpoint{0.991151in}{1.243383in}}%
\pgfusepath{stroke}%
\end{pgfscope}%
\begin{pgfscope}%
\pgfsetrectcap%
\pgfsetroundjoin%
\pgfsetlinewidth{0.501875pt}%
\definecolor{currentstroke}{rgb}{0.000000,0.000000,0.000000}%
\pgfsetstrokecolor{currentstroke}%
\pgfsetdash{}{0pt}%
\pgfpathmoveto{\pgfqpoint{1.226464in}{1.140505in}}%
\pgfpathlineto{\pgfqpoint{1.161821in}{1.119831in}}%
\pgfusepath{stroke}%
\end{pgfscope}%
\begin{pgfscope}%
\pgfsetrectcap%
\pgfsetroundjoin%
\pgfsetlinewidth{0.501875pt}%
\definecolor{currentstroke}{rgb}{0.000000,0.000000,0.000000}%
\pgfsetstrokecolor{currentstroke}%
\pgfsetdash{}{0pt}%
\pgfpathmoveto{\pgfqpoint{1.313124in}{1.078076in}}%
\pgfpathlineto{\pgfqpoint{1.248317in}{1.057214in}}%
\pgfusepath{stroke}%
\end{pgfscope}%
\begin{pgfscope}%
\pgfsetrectcap%
\pgfsetroundjoin%
\pgfsetlinewidth{0.501875pt}%
\definecolor{currentstroke}{rgb}{0.000000,0.000000,0.000000}%
\pgfsetstrokecolor{currentstroke}%
\pgfsetdash{}{0pt}%
\pgfpathmoveto{\pgfqpoint{1.400572in}{1.015078in}}%
\pgfpathlineto{\pgfqpoint{1.335602in}{0.994026in}}%
\pgfusepath{stroke}%
\end{pgfscope}%
\begin{pgfscope}%
\pgfsetrectcap%
\pgfsetroundjoin%
\pgfsetlinewidth{0.501875pt}%
\definecolor{currentstroke}{rgb}{0.000000,0.000000,0.000000}%
\pgfsetstrokecolor{currentstroke}%
\pgfsetdash{}{0pt}%
\pgfpathmoveto{\pgfqpoint{1.577879in}{0.887347in}}%
\pgfpathlineto{\pgfqpoint{1.512582in}{0.865906in}}%
\pgfusepath{stroke}%
\end{pgfscope}%
\begin{pgfscope}%
\pgfsetrectcap%
\pgfsetroundjoin%
\pgfsetlinewidth{0.501875pt}%
\definecolor{currentstroke}{rgb}{0.000000,0.000000,0.000000}%
\pgfsetstrokecolor{currentstroke}%
\pgfsetdash{}{0pt}%
\pgfpathmoveto{\pgfqpoint{1.667760in}{0.822597in}}%
\pgfpathlineto{\pgfqpoint{1.602299in}{0.800957in}}%
\pgfusepath{stroke}%
\end{pgfscope}%
\begin{pgfscope}%
\pgfsetrectcap%
\pgfsetroundjoin%
\pgfsetlinewidth{0.501875pt}%
\definecolor{currentstroke}{rgb}{0.000000,0.000000,0.000000}%
\pgfsetstrokecolor{currentstroke}%
\pgfsetdash{}{0pt}%
\pgfpathmoveto{\pgfqpoint{1.758474in}{0.757246in}}%
\pgfpathlineto{\pgfqpoint{1.692850in}{0.735405in}}%
\pgfusepath{stroke}%
\end{pgfscope}%
\begin{pgfscope}%
\pgfsetrectcap%
\pgfsetroundjoin%
\pgfsetlinewidth{0.501875pt}%
\definecolor{currentstroke}{rgb}{0.000000,0.000000,0.000000}%
\pgfsetstrokecolor{currentstroke}%
\pgfsetdash{}{0pt}%
\pgfpathmoveto{\pgfqpoint{1.942450in}{0.624710in}}%
\pgfpathlineto{\pgfqpoint{1.876497in}{0.602458in}}%
\pgfusepath{stroke}%
\end{pgfscope}%
\begin{pgfscope}%
\pgfsetrectcap%
\pgfsetroundjoin%
\pgfsetlinewidth{0.501875pt}%
\definecolor{currentstroke}{rgb}{0.000000,0.000000,0.000000}%
\pgfsetstrokecolor{currentstroke}%
\pgfsetdash{}{0pt}%
\pgfpathmoveto{\pgfqpoint{2.035735in}{0.557508in}}%
\pgfpathlineto{\pgfqpoint{1.969618in}{0.535046in}}%
\pgfusepath{stroke}%
\end{pgfscope}%
\begin{pgfscope}%
\pgfsetrectcap%
\pgfsetroundjoin%
\pgfsetlinewidth{0.501875pt}%
\definecolor{currentstroke}{rgb}{0.000000,0.000000,0.000000}%
\pgfsetstrokecolor{currentstroke}%
\pgfsetdash{}{0pt}%
\pgfpathmoveto{\pgfqpoint{2.129902in}{0.489670in}}%
\pgfpathlineto{\pgfqpoint{2.063620in}{0.466995in}}%
\pgfusepath{stroke}%
\end{pgfscope}%
\begin{pgfscope}%
\definecolor{textcolor}{rgb}{0.000000,0.000000,0.000000}%
\pgfsetstrokecolor{textcolor}%
\pgfsetfillcolor{textcolor}%
\pgftext[x=0.883212in,y=0.671627in,,,rotate=324.186895]{\color{textcolor}{\rmfamily\fontsize{10.000000}{12.000000}\selectfont\catcode`\^=\active\def^{\ifmmode\sp\else\^{}\fi}\catcode`\%=\active\def%{\%}Time to Expiry (Days)}}%
\end{pgfscope}%
\begin{pgfscope}%
\pgfsetrectcap%
\pgfsetroundjoin%
\pgfsetlinewidth{0.501875pt}%
\definecolor{currentstroke}{rgb}{0.000000,0.000000,0.000000}%
\pgfsetstrokecolor{currentstroke}%
\pgfsetdash{}{0pt}%
\pgfpathmoveto{\pgfqpoint{4.749691in}{1.337846in}}%
\pgfpathlineto{\pgfqpoint{2.155964in}{0.447388in}}%
\pgfusepath{stroke}%
\end{pgfscope}%
\begin{pgfscope}%
\pgfsetrectcap%
\pgfsetroundjoin%
\pgfsetlinewidth{0.501875pt}%
\definecolor{currentstroke}{rgb}{0.000000,0.000000,0.000000}%
\pgfsetstrokecolor{currentstroke}%
\pgfsetdash{}{0pt}%
\pgfpathmoveto{\pgfqpoint{2.144293in}{0.459367in}}%
\pgfpathlineto{\pgfqpoint{2.189402in}{0.426825in}}%
\pgfusepath{stroke}%
\end{pgfscope}%
\begin{pgfscope}%
\definecolor{textcolor}{rgb}{0.000000,0.000000,0.000000}%
\pgfsetstrokecolor{textcolor}%
\pgfsetfillcolor{textcolor}%
\pgftext[x=2.278268in,y=0.226168in,,top]{\color{textcolor}{\rmfamily\fontsize{10.000000}{12.000000}\selectfont\catcode`\^=\active\def^{\ifmmode\sp\else\^{}\fi}\catcode`\%=\active\def%{\%}$\mathdefault{45}$}}%
\end{pgfscope}%
\begin{pgfscope}%
\pgfsetrectcap%
\pgfsetroundjoin%
\pgfsetlinewidth{0.501875pt}%
\definecolor{currentstroke}{rgb}{0.000000,0.000000,0.000000}%
\pgfsetstrokecolor{currentstroke}%
\pgfsetdash{}{0pt}%
\pgfpathmoveto{\pgfqpoint{2.655344in}{0.634623in}}%
\pgfpathlineto{\pgfqpoint{2.701042in}{0.602867in}}%
\pgfusepath{stroke}%
\end{pgfscope}%
\begin{pgfscope}%
\definecolor{textcolor}{rgb}{0.000000,0.000000,0.000000}%
\pgfsetstrokecolor{textcolor}%
\pgfsetfillcolor{textcolor}%
\pgftext[x=2.789770in,y=0.404976in,,top]{\color{textcolor}{\rmfamily\fontsize{10.000000}{12.000000}\selectfont\catcode`\^=\active\def^{\ifmmode\sp\else\^{}\fi}\catcode`\%=\active\def%{\%}$\mathdefault{60}$}}%
\end{pgfscope}%
\begin{pgfscope}%
\pgfsetrectcap%
\pgfsetroundjoin%
\pgfsetlinewidth{0.501875pt}%
\definecolor{currentstroke}{rgb}{0.000000,0.000000,0.000000}%
\pgfsetstrokecolor{currentstroke}%
\pgfsetdash{}{0pt}%
\pgfpathmoveto{\pgfqpoint{3.154154in}{0.805682in}}%
\pgfpathlineto{\pgfqpoint{3.200399in}{0.774682in}}%
\pgfusepath{stroke}%
\end{pgfscope}%
\begin{pgfscope}%
\definecolor{textcolor}{rgb}{0.000000,0.000000,0.000000}%
\pgfsetstrokecolor{textcolor}%
\pgfsetfillcolor{textcolor}%
\pgftext[x=3.288970in,y=0.579484in,,top]{\color{textcolor}{\rmfamily\fontsize{10.000000}{12.000000}\selectfont\catcode`\^=\active\def^{\ifmmode\sp\else\^{}\fi}\catcode`\%=\active\def%{\%}$\mathdefault{75}$}}%
\end{pgfscope}%
\begin{pgfscope}%
\pgfsetrectcap%
\pgfsetroundjoin%
\pgfsetlinewidth{0.501875pt}%
\definecolor{currentstroke}{rgb}{0.000000,0.000000,0.000000}%
\pgfsetstrokecolor{currentstroke}%
\pgfsetdash{}{0pt}%
\pgfpathmoveto{\pgfqpoint{3.641158in}{0.972691in}}%
\pgfpathlineto{\pgfqpoint{3.687911in}{0.942422in}}%
\pgfusepath{stroke}%
\end{pgfscope}%
\begin{pgfscope}%
\definecolor{textcolor}{rgb}{0.000000,0.000000,0.000000}%
\pgfsetstrokecolor{textcolor}%
\pgfsetfillcolor{textcolor}%
\pgftext[x=3.776307in,y=0.749845in,,top]{\color{textcolor}{\rmfamily\fontsize{10.000000}{12.000000}\selectfont\catcode`\^=\active\def^{\ifmmode\sp\else\^{}\fi}\catcode`\%=\active\def%{\%}$\mathdefault{90}$}}%
\end{pgfscope}%
\begin{pgfscope}%
\pgfsetrectcap%
\pgfsetroundjoin%
\pgfsetlinewidth{0.501875pt}%
\definecolor{currentstroke}{rgb}{0.000000,0.000000,0.000000}%
\pgfsetstrokecolor{currentstroke}%
\pgfsetdash{}{0pt}%
\pgfpathmoveto{\pgfqpoint{4.116768in}{1.135794in}}%
\pgfpathlineto{\pgfqpoint{4.163994in}{1.106229in}}%
\pgfusepath{stroke}%
\end{pgfscope}%
\begin{pgfscope}%
\definecolor{textcolor}{rgb}{0.000000,0.000000,0.000000}%
\pgfsetstrokecolor{textcolor}%
\pgfsetfillcolor{textcolor}%
\pgftext[x=4.252199in,y=0.916205in,,top]{\color{textcolor}{\rmfamily\fontsize{10.000000}{12.000000}\selectfont\catcode`\^=\active\def^{\ifmmode\sp\else\^{}\fi}\catcode`\%=\active\def%{\%}$\mathdefault{105}$}}%
\end{pgfscope}%
\begin{pgfscope}%
\pgfsetrectcap%
\pgfsetroundjoin%
\pgfsetlinewidth{0.501875pt}%
\definecolor{currentstroke}{rgb}{0.000000,0.000000,0.000000}%
\pgfsetstrokecolor{currentstroke}%
\pgfsetdash{}{0pt}%
\pgfpathmoveto{\pgfqpoint{4.581381in}{1.295125in}}%
\pgfpathlineto{\pgfqpoint{4.629045in}{1.266240in}}%
\pgfusepath{stroke}%
\end{pgfscope}%
\begin{pgfscope}%
\definecolor{textcolor}{rgb}{0.000000,0.000000,0.000000}%
\pgfsetstrokecolor{textcolor}%
\pgfsetfillcolor{textcolor}%
\pgftext[x=4.717044in,y=1.078703in,,top]{\color{textcolor}{\rmfamily\fontsize{10.000000}{12.000000}\selectfont\catcode`\^=\active\def^{\ifmmode\sp\else\^{}\fi}\catcode`\%=\active\def%{\%}$\mathdefault{120}$}}%
\end{pgfscope}%
\begin{pgfscope}%
\pgfsetrectcap%
\pgfsetroundjoin%
\pgfsetlinewidth{0.501875pt}%
\definecolor{currentstroke}{rgb}{0.000000,0.000000,0.000000}%
\pgfsetstrokecolor{currentstroke}%
\pgfsetdash{}{0pt}%
\pgfpathmoveto{\pgfqpoint{2.273228in}{0.503583in}}%
\pgfpathlineto{\pgfqpoint{2.318488in}{0.471240in}}%
\pgfusepath{stroke}%
\end{pgfscope}%
\begin{pgfscope}%
\pgfsetrectcap%
\pgfsetroundjoin%
\pgfsetlinewidth{0.501875pt}%
\definecolor{currentstroke}{rgb}{0.000000,0.000000,0.000000}%
\pgfsetstrokecolor{currentstroke}%
\pgfsetdash{}{0pt}%
\pgfpathmoveto{\pgfqpoint{2.401377in}{0.547529in}}%
\pgfpathlineto{\pgfqpoint{2.446786in}{0.515384in}}%
\pgfusepath{stroke}%
\end{pgfscope}%
\begin{pgfscope}%
\pgfsetrectcap%
\pgfsetroundjoin%
\pgfsetlinewidth{0.501875pt}%
\definecolor{currentstroke}{rgb}{0.000000,0.000000,0.000000}%
\pgfsetstrokecolor{currentstroke}%
\pgfsetdash{}{0pt}%
\pgfpathmoveto{\pgfqpoint{2.528747in}{0.591209in}}%
\pgfpathlineto{\pgfqpoint{2.574301in}{0.559259in}}%
\pgfusepath{stroke}%
\end{pgfscope}%
\begin{pgfscope}%
\pgfsetrectcap%
\pgfsetroundjoin%
\pgfsetlinewidth{0.501875pt}%
\definecolor{currentstroke}{rgb}{0.000000,0.000000,0.000000}%
\pgfsetstrokecolor{currentstroke}%
\pgfsetdash{}{0pt}%
\pgfpathmoveto{\pgfqpoint{2.781177in}{0.677775in}}%
\pgfpathlineto{\pgfqpoint{2.827015in}{0.646211in}}%
\pgfusepath{stroke}%
\end{pgfscope}%
\begin{pgfscope}%
\pgfsetrectcap%
\pgfsetroundjoin%
\pgfsetlinewidth{0.501875pt}%
\definecolor{currentstroke}{rgb}{0.000000,0.000000,0.000000}%
\pgfsetstrokecolor{currentstroke}%
\pgfsetdash{}{0pt}%
\pgfpathmoveto{\pgfqpoint{2.906252in}{0.720668in}}%
\pgfpathlineto{\pgfqpoint{2.952228in}{0.689293in}}%
\pgfusepath{stroke}%
\end{pgfscope}%
\begin{pgfscope}%
\pgfsetrectcap%
\pgfsetroundjoin%
\pgfsetlinewidth{0.501875pt}%
\definecolor{currentstroke}{rgb}{0.000000,0.000000,0.000000}%
\pgfsetstrokecolor{currentstroke}%
\pgfsetdash{}{0pt}%
\pgfpathmoveto{\pgfqpoint{3.030575in}{0.763302in}}%
\pgfpathlineto{\pgfqpoint{3.076687in}{0.732116in}}%
\pgfusepath{stroke}%
\end{pgfscope}%
\begin{pgfscope}%
\pgfsetrectcap%
\pgfsetroundjoin%
\pgfsetlinewidth{0.501875pt}%
\definecolor{currentstroke}{rgb}{0.000000,0.000000,0.000000}%
\pgfsetstrokecolor{currentstroke}%
\pgfsetdash{}{0pt}%
\pgfpathmoveto{\pgfqpoint{3.276996in}{0.847808in}}%
\pgfpathlineto{\pgfqpoint{3.323371in}{0.816993in}}%
\pgfusepath{stroke}%
\end{pgfscope}%
\begin{pgfscope}%
\pgfsetrectcap%
\pgfsetroundjoin%
\pgfsetlinewidth{0.501875pt}%
\definecolor{currentstroke}{rgb}{0.000000,0.000000,0.000000}%
\pgfsetstrokecolor{currentstroke}%
\pgfsetdash{}{0pt}%
\pgfpathmoveto{\pgfqpoint{3.399106in}{0.889683in}}%
\pgfpathlineto{\pgfqpoint{3.445610in}{0.859052in}}%
\pgfusepath{stroke}%
\end{pgfscope}%
\begin{pgfscope}%
\pgfsetrectcap%
\pgfsetroundjoin%
\pgfsetlinewidth{0.501875pt}%
\definecolor{currentstroke}{rgb}{0.000000,0.000000,0.000000}%
\pgfsetstrokecolor{currentstroke}%
\pgfsetdash{}{0pt}%
\pgfpathmoveto{\pgfqpoint{3.520491in}{0.931310in}}%
\pgfpathlineto{\pgfqpoint{3.567121in}{0.900861in}}%
\pgfusepath{stroke}%
\end{pgfscope}%
\begin{pgfscope}%
\pgfsetrectcap%
\pgfsetroundjoin%
\pgfsetlinewidth{0.501875pt}%
\definecolor{currentstroke}{rgb}{0.000000,0.000000,0.000000}%
\pgfsetstrokecolor{currentstroke}%
\pgfsetdash{}{0pt}%
\pgfpathmoveto{\pgfqpoint{3.761113in}{1.013828in}}%
\pgfpathlineto{\pgfqpoint{3.807988in}{0.983737in}}%
\pgfusepath{stroke}%
\end{pgfscope}%
\begin{pgfscope}%
\pgfsetrectcap%
\pgfsetroundjoin%
\pgfsetlinewidth{0.501875pt}%
\definecolor{currentstroke}{rgb}{0.000000,0.000000,0.000000}%
\pgfsetstrokecolor{currentstroke}%
\pgfsetdash{}{0pt}%
\pgfpathmoveto{\pgfqpoint{3.880362in}{1.054722in}}%
\pgfpathlineto{\pgfqpoint{3.927356in}{1.024808in}}%
\pgfusepath{stroke}%
\end{pgfscope}%
\begin{pgfscope}%
\pgfsetrectcap%
\pgfsetroundjoin%
\pgfsetlinewidth{0.501875pt}%
\definecolor{currentstroke}{rgb}{0.000000,0.000000,0.000000}%
\pgfsetstrokecolor{currentstroke}%
\pgfsetdash{}{0pt}%
\pgfpathmoveto{\pgfqpoint{3.998912in}{1.095377in}}%
\pgfpathlineto{\pgfqpoint{4.046023in}{1.065638in}}%
\pgfusepath{stroke}%
\end{pgfscope}%
\begin{pgfscope}%
\pgfsetrectcap%
\pgfsetroundjoin%
\pgfsetlinewidth{0.501875pt}%
\definecolor{currentstroke}{rgb}{0.000000,0.000000,0.000000}%
\pgfsetstrokecolor{currentstroke}%
\pgfsetdash{}{0pt}%
\pgfpathmoveto{\pgfqpoint{4.233937in}{1.175975in}}%
\pgfpathlineto{\pgfqpoint{4.281276in}{1.146582in}}%
\pgfusepath{stroke}%
\end{pgfscope}%
\begin{pgfscope}%
\pgfsetrectcap%
\pgfsetroundjoin%
\pgfsetlinewidth{0.501875pt}%
\definecolor{currentstroke}{rgb}{0.000000,0.000000,0.000000}%
\pgfsetstrokecolor{currentstroke}%
\pgfsetdash{}{0pt}%
\pgfpathmoveto{\pgfqpoint{4.350425in}{1.215922in}}%
\pgfpathlineto{\pgfqpoint{4.397874in}{1.186701in}}%
\pgfusepath{stroke}%
\end{pgfscope}%
\begin{pgfscope}%
\pgfsetrectcap%
\pgfsetroundjoin%
\pgfsetlinewidth{0.501875pt}%
\definecolor{currentstroke}{rgb}{0.000000,0.000000,0.000000}%
\pgfsetstrokecolor{currentstroke}%
\pgfsetdash{}{0pt}%
\pgfpathmoveto{\pgfqpoint{4.466238in}{1.255638in}}%
\pgfpathlineto{\pgfqpoint{4.513796in}{1.226586in}}%
\pgfusepath{stroke}%
\end{pgfscope}%
\begin{pgfscope}%
\pgfsetrectcap%
\pgfsetroundjoin%
\pgfsetlinewidth{0.501875pt}%
\definecolor{currentstroke}{rgb}{0.000000,0.000000,0.000000}%
\pgfsetstrokecolor{currentstroke}%
\pgfsetdash{}{0pt}%
\pgfpathmoveto{\pgfqpoint{4.695861in}{1.334383in}}%
\pgfpathlineto{\pgfqpoint{4.743629in}{1.305665in}}%
\pgfusepath{stroke}%
\end{pgfscope}%
\begin{pgfscope}%
\definecolor{textcolor}{rgb}{0.000000,0.000000,0.000000}%
\pgfsetstrokecolor{textcolor}%
\pgfsetfillcolor{textcolor}%
\pgftext[x=3.765372in,y=0.421364in,,,rotate=18.947963]{\color{textcolor}{\rmfamily\fontsize{10.000000}{12.000000}\selectfont\catcode`\^=\active\def^{\ifmmode\sp\else\^{}\fi}\catcode`\%=\active\def%{\%}Underlying (S)}}%
\end{pgfscope}%
\begin{pgfscope}%
\pgfsetrectcap%
\pgfsetroundjoin%
\pgfsetlinewidth{0.501875pt}%
\definecolor{currentstroke}{rgb}{0.000000,0.000000,0.000000}%
\pgfsetstrokecolor{currentstroke}%
\pgfsetdash{}{0pt}%
\pgfpathmoveto{\pgfqpoint{4.749691in}{1.337846in}}%
\pgfpathlineto{\pgfqpoint{4.834815in}{3.504368in}}%
\pgfusepath{stroke}%
\end{pgfscope}%
\begin{pgfscope}%
\pgfsetrectcap%
\pgfsetroundjoin%
\pgfsetlinewidth{0.501875pt}%
\definecolor{currentstroke}{rgb}{0.000000,0.000000,0.000000}%
\pgfsetstrokecolor{currentstroke}%
\pgfsetdash{}{0pt}%
\pgfpathmoveto{\pgfqpoint{4.735406in}{1.389101in}}%
\pgfpathlineto{\pgfqpoint{4.783248in}{1.360504in}}%
\pgfusepath{stroke}%
\end{pgfscope}%
\begin{pgfscope}%
\definecolor{textcolor}{rgb}{0.000000,0.000000,0.000000}%
\pgfsetstrokecolor{textcolor}%
\pgfsetfillcolor{textcolor}%
\pgftext[x=5.028906in,y=1.357205in,,top]{\color{textcolor}{\rmfamily\fontsize{10.000000}{12.000000}\selectfont\catcode`\^=\active\def^{\ifmmode\sp\else\^{}\fi}\catcode`\%=\active\def%{\%}$\mathdefault{0.0}$}}%
\end{pgfscope}%
\begin{pgfscope}%
\pgfsetrectcap%
\pgfsetroundjoin%
\pgfsetlinewidth{0.501875pt}%
\definecolor{currentstroke}{rgb}{0.000000,0.000000,0.000000}%
\pgfsetstrokecolor{currentstroke}%
\pgfsetdash{}{0pt}%
\pgfpathmoveto{\pgfqpoint{4.750099in}{1.765857in}}%
\pgfpathlineto{\pgfqpoint{4.798297in}{1.737868in}}%
\pgfusepath{stroke}%
\end{pgfscope}%
\begin{pgfscope}%
\definecolor{textcolor}{rgb}{0.000000,0.000000,0.000000}%
\pgfsetstrokecolor{textcolor}%
\pgfsetfillcolor{textcolor}%
\pgftext[x=5.045626in,y=1.734639in,,top]{\color{textcolor}{\rmfamily\fontsize{10.000000}{12.000000}\selectfont\catcode`\^=\active\def^{\ifmmode\sp\else\^{}\fi}\catcode`\%=\active\def%{\%}$\mathdefault{2.5}$}}%
\end{pgfscope}%
\begin{pgfscope}%
\pgfsetrectcap%
\pgfsetroundjoin%
\pgfsetlinewidth{0.501875pt}%
\definecolor{currentstroke}{rgb}{0.000000,0.000000,0.000000}%
\pgfsetstrokecolor{currentstroke}%
\pgfsetdash{}{0pt}%
\pgfpathmoveto{\pgfqpoint{4.764992in}{2.147752in}}%
\pgfpathlineto{\pgfqpoint{4.813551in}{2.120390in}}%
\pgfusepath{stroke}%
\end{pgfscope}%
\begin{pgfscope}%
\definecolor{textcolor}{rgb}{0.000000,0.000000,0.000000}%
\pgfsetstrokecolor{textcolor}%
\pgfsetfillcolor{textcolor}%
\pgftext[x=5.062575in,y=2.117233in,,top]{\color{textcolor}{\rmfamily\fontsize{10.000000}{12.000000}\selectfont\catcode`\^=\active\def^{\ifmmode\sp\else\^{}\fi}\catcode`\%=\active\def%{\%}$\mathdefault{5.0}$}}%
\end{pgfscope}%
\begin{pgfscope}%
\pgfsetrectcap%
\pgfsetroundjoin%
\pgfsetlinewidth{0.501875pt}%
\definecolor{currentstroke}{rgb}{0.000000,0.000000,0.000000}%
\pgfsetstrokecolor{currentstroke}%
\pgfsetdash{}{0pt}%
\pgfpathmoveto{\pgfqpoint{4.780089in}{2.534891in}}%
\pgfpathlineto{\pgfqpoint{4.829015in}{2.508177in}}%
\pgfusepath{stroke}%
\end{pgfscope}%
\begin{pgfscope}%
\definecolor{textcolor}{rgb}{0.000000,0.000000,0.000000}%
\pgfsetstrokecolor{textcolor}%
\pgfsetfillcolor{textcolor}%
\pgftext[x=5.079757in,y=2.505094in,,top]{\color{textcolor}{\rmfamily\fontsize{10.000000}{12.000000}\selectfont\catcode`\^=\active\def^{\ifmmode\sp\else\^{}\fi}\catcode`\%=\active\def%{\%}$\mathdefault{7.5}$}}%
\end{pgfscope}%
\begin{pgfscope}%
\pgfsetrectcap%
\pgfsetroundjoin%
\pgfsetlinewidth{0.501875pt}%
\definecolor{currentstroke}{rgb}{0.000000,0.000000,0.000000}%
\pgfsetstrokecolor{currentstroke}%
\pgfsetdash{}{0pt}%
\pgfpathmoveto{\pgfqpoint{4.795396in}{2.927384in}}%
\pgfpathlineto{\pgfqpoint{4.844694in}{2.901337in}}%
\pgfusepath{stroke}%
\end{pgfscope}%
\begin{pgfscope}%
\definecolor{textcolor}{rgb}{0.000000,0.000000,0.000000}%
\pgfsetstrokecolor{textcolor}%
\pgfsetfillcolor{textcolor}%
\pgftext[x=5.097177in,y=2.898331in,,top]{\color{textcolor}{\rmfamily\fontsize{10.000000}{12.000000}\selectfont\catcode`\^=\active\def^{\ifmmode\sp\else\^{}\fi}\catcode`\%=\active\def%{\%}$\mathdefault{10.0}$}}%
\end{pgfscope}%
\begin{pgfscope}%
\pgfsetrectcap%
\pgfsetroundjoin%
\pgfsetlinewidth{0.501875pt}%
\definecolor{currentstroke}{rgb}{0.000000,0.000000,0.000000}%
\pgfsetstrokecolor{currentstroke}%
\pgfsetdash{}{0pt}%
\pgfpathmoveto{\pgfqpoint{4.810915in}{3.325341in}}%
\pgfpathlineto{\pgfqpoint{4.860591in}{3.299983in}}%
\pgfusepath{stroke}%
\end{pgfscope}%
\begin{pgfscope}%
\definecolor{textcolor}{rgb}{0.000000,0.000000,0.000000}%
\pgfsetstrokecolor{textcolor}%
\pgfsetfillcolor{textcolor}%
\pgftext[x=5.114841in,y=3.297057in,,top]{\color{textcolor}{\rmfamily\fontsize{10.000000}{12.000000}\selectfont\catcode`\^=\active\def^{\ifmmode\sp\else\^{}\fi}\catcode`\%=\active\def%{\%}$\mathdefault{12.5}$}}%
\end{pgfscope}%
\begin{pgfscope}%
\definecolor{textcolor}{rgb}{0.000000,0.000000,0.000000}%
\pgfsetstrokecolor{textcolor}%
\pgfsetfillcolor{textcolor}%
\pgftext[x=5.432051in,y=2.352657in,,,rotate=87.749957]{\color{textcolor}{\rmfamily\fontsize{10.000000}{12.000000}\selectfont\catcode`\^=\active\def^{\ifmmode\sp\else\^{}\fi}\catcode`\%=\active\def%{\%}Option Price}}%
\end{pgfscope}%
\begin{pgfscope}%
\pgfpathrectangle{\pgfqpoint{0.050000in}{0.050000in}}{\pgfqpoint{4.900000in}{4.900000in}}%
\pgfusepath{clip}%
\pgfsetbuttcap%
\pgfsetroundjoin%
\definecolor{currentfill}{rgb}{0.665852,0.652072,0.800200}%
\pgfsetfillcolor{currentfill}%
\pgfsetfillopacity{0.900000}%
\pgfsetlinewidth{0.200750pt}%
\definecolor{currentstroke}{rgb}{0.200000,0.200000,0.200000}%
\pgfsetstrokecolor{currentstroke}%
\pgfsetdash{}{0pt}%
\pgfpathmoveto{\pgfqpoint{2.839953in}{3.214644in}}%
\pgfpathlineto{\pgfqpoint{2.878316in}{2.418448in}}%
\pgfpathlineto{\pgfqpoint{2.908499in}{2.389022in}}%
\pgfpathlineto{\pgfqpoint{2.938782in}{2.363234in}}%
\pgfpathlineto{\pgfqpoint{2.899726in}{2.900065in}}%
\pgfpathlineto{\pgfqpoint{2.869733in}{3.031607in}}%
\pgfpathlineto{\pgfqpoint{2.839953in}{3.214644in}}%
\pgfpathclose%
\pgfusepath{stroke,fill}%
\end{pgfscope}%
\begin{pgfscope}%
\pgfpathrectangle{\pgfqpoint{0.050000in}{0.050000in}}{\pgfqpoint{4.900000in}{4.900000in}}%
\pgfusepath{clip}%
\pgfsetbuttcap%
\pgfsetroundjoin%
\definecolor{currentfill}{rgb}{0.767290,0.757693,0.860854}%
\pgfsetfillcolor{currentfill}%
\pgfsetfillopacity{0.900000}%
\pgfsetlinewidth{0.200750pt}%
\definecolor{currentstroke}{rgb}{0.200000,0.200000,0.200000}%
\pgfsetstrokecolor{currentstroke}%
\pgfsetdash{}{0pt}%
\pgfpathmoveto{\pgfqpoint{2.899726in}{2.900065in}}%
\pgfpathlineto{\pgfqpoint{2.938782in}{2.363234in}}%
\pgfpathlineto{\pgfqpoint{2.969160in}{2.339421in}}%
\pgfpathlineto{\pgfqpoint{2.999631in}{2.316806in}}%
\pgfpathlineto{\pgfqpoint{2.960151in}{2.714228in}}%
\pgfpathlineto{\pgfqpoint{2.929875in}{2.797871in}}%
\pgfpathlineto{\pgfqpoint{2.899726in}{2.900065in}}%
\pgfpathclose%
\pgfusepath{stroke,fill}%
\end{pgfscope}%
\begin{pgfscope}%
\pgfpathrectangle{\pgfqpoint{0.050000in}{0.050000in}}{\pgfqpoint{4.900000in}{4.900000in}}%
\pgfusepath{clip}%
\pgfsetbuttcap%
\pgfsetroundjoin%
\definecolor{currentfill}{rgb}{0.820992,0.813610,0.892964}%
\pgfsetfillcolor{currentfill}%
\pgfsetfillopacity{0.900000}%
\pgfsetlinewidth{0.200750pt}%
\definecolor{currentstroke}{rgb}{0.200000,0.200000,0.200000}%
\pgfsetstrokecolor{currentstroke}%
\pgfsetdash{}{0pt}%
\pgfpathmoveto{\pgfqpoint{2.960151in}{2.714228in}}%
\pgfpathlineto{\pgfqpoint{2.999631in}{2.316806in}}%
\pgfpathlineto{\pgfqpoint{3.030194in}{2.294976in}}%
\pgfpathlineto{\pgfqpoint{3.060848in}{2.273685in}}%
\pgfpathlineto{\pgfqpoint{3.021037in}{2.581110in}}%
\pgfpathlineto{\pgfqpoint{2.990541in}{2.643173in}}%
\pgfpathlineto{\pgfqpoint{2.960151in}{2.714228in}}%
\pgfpathclose%
\pgfusepath{stroke,fill}%
\end{pgfscope}%
\begin{pgfscope}%
\pgfpathrectangle{\pgfqpoint{0.050000in}{0.050000in}}{\pgfqpoint{4.900000in}{4.900000in}}%
\pgfusepath{clip}%
\pgfsetbuttcap%
\pgfsetroundjoin%
\definecolor{currentfill}{rgb}{0.862760,0.857101,0.917939}%
\pgfsetfillcolor{currentfill}%
\pgfsetfillopacity{0.900000}%
\pgfsetlinewidth{0.200750pt}%
\definecolor{currentstroke}{rgb}{0.200000,0.200000,0.200000}%
\pgfsetstrokecolor{currentstroke}%
\pgfsetdash{}{0pt}%
\pgfpathmoveto{\pgfqpoint{3.021037in}{2.581110in}}%
\pgfpathlineto{\pgfqpoint{3.060848in}{2.273685in}}%
\pgfpathlineto{\pgfqpoint{3.091593in}{2.252777in}}%
\pgfpathlineto{\pgfqpoint{3.122429in}{2.232149in}}%
\pgfpathlineto{\pgfqpoint{3.082329in}{2.475463in}}%
\pgfpathlineto{\pgfqpoint{3.051634in}{2.525727in}}%
\pgfpathlineto{\pgfqpoint{3.021037in}{2.581110in}}%
\pgfpathclose%
\pgfusepath{stroke,fill}%
\end{pgfscope}%
\begin{pgfscope}%
\pgfpathrectangle{\pgfqpoint{0.050000in}{0.050000in}}{\pgfqpoint{4.900000in}{4.900000in}}%
\pgfusepath{clip}%
\pgfsetbuttcap%
\pgfsetroundjoin%
\definecolor{currentfill}{rgb}{0.886628,0.881953,0.932211}%
\pgfsetfillcolor{currentfill}%
\pgfsetfillopacity{0.900000}%
\pgfsetlinewidth{0.200750pt}%
\definecolor{currentstroke}{rgb}{0.200000,0.200000,0.200000}%
\pgfsetstrokecolor{currentstroke}%
\pgfsetdash{}{0pt}%
\pgfpathmoveto{\pgfqpoint{3.082329in}{2.475463in}}%
\pgfpathlineto{\pgfqpoint{3.122429in}{2.232149in}}%
\pgfpathlineto{\pgfqpoint{3.153356in}{2.211727in}}%
\pgfpathlineto{\pgfqpoint{3.184375in}{2.191460in}}%
\pgfpathlineto{\pgfqpoint{3.144004in}{2.386205in}}%
\pgfpathlineto{\pgfqpoint{3.113119in}{2.429221in}}%
\pgfpathlineto{\pgfqpoint{3.082329in}{2.475463in}}%
\pgfpathclose%
\pgfusepath{stroke,fill}%
\end{pgfscope}%
\begin{pgfscope}%
\pgfpathrectangle{\pgfqpoint{0.050000in}{0.050000in}}{\pgfqpoint{4.900000in}{4.900000in}}%
\pgfusepath{clip}%
\pgfsetbuttcap%
\pgfsetroundjoin%
\definecolor{currentfill}{rgb}{0.910496,0.906805,0.946482}%
\pgfsetfillcolor{currentfill}%
\pgfsetfillopacity{0.900000}%
\pgfsetlinewidth{0.200750pt}%
\definecolor{currentstroke}{rgb}{0.200000,0.200000,0.200000}%
\pgfsetstrokecolor{currentstroke}%
\pgfsetdash{}{0pt}%
\pgfpathmoveto{\pgfqpoint{3.144004in}{2.386205in}}%
\pgfpathlineto{\pgfqpoint{3.184375in}{2.191460in}}%
\pgfpathlineto{\pgfqpoint{3.215486in}{2.171310in}}%
\pgfpathlineto{\pgfqpoint{3.246689in}{2.151246in}}%
\pgfpathlineto{\pgfqpoint{3.206057in}{2.307624in}}%
\pgfpathlineto{\pgfqpoint{3.174984in}{2.345824in}}%
\pgfpathlineto{\pgfqpoint{3.144004in}{2.386205in}}%
\pgfpathclose%
\pgfusepath{stroke,fill}%
\end{pgfscope}%
\begin{pgfscope}%
\pgfpathrectangle{\pgfqpoint{0.050000in}{0.050000in}}{\pgfqpoint{4.900000in}{4.900000in}}%
\pgfusepath{clip}%
\pgfsetbuttcap%
\pgfsetroundjoin%
\definecolor{currentfill}{rgb}{0.970165,0.968935,0.982161}%
\pgfsetfillcolor{currentfill}%
\pgfsetfillopacity{0.900000}%
\pgfsetlinewidth{0.200750pt}%
\definecolor{currentstroke}{rgb}{0.200000,0.200000,0.200000}%
\pgfsetstrokecolor{currentstroke}%
\pgfsetdash{}{0pt}%
\pgfpathmoveto{\pgfqpoint{2.149107in}{2.147167in}}%
\pgfpathlineto{\pgfqpoint{2.192608in}{2.173405in}}%
\pgfpathlineto{\pgfqpoint{2.235898in}{2.223838in}}%
\pgfpathlineto{\pgfqpoint{2.265566in}{2.224206in}}%
\pgfpathlineto{\pgfqpoint{2.295350in}{2.223053in}}%
\pgfpathlineto{\pgfqpoint{2.252001in}{2.158941in}}%
\pgfpathlineto{\pgfqpoint{2.208479in}{2.120307in}}%
\pgfpathlineto{\pgfqpoint{2.178750in}{2.133126in}}%
\pgfpathlineto{\pgfqpoint{2.149107in}{2.147167in}}%
\pgfpathclose%
\pgfusepath{stroke,fill}%
\end{pgfscope}%
\begin{pgfscope}%
\pgfpathrectangle{\pgfqpoint{0.050000in}{0.050000in}}{\pgfqpoint{4.900000in}{4.900000in}}%
\pgfusepath{clip}%
\pgfsetbuttcap%
\pgfsetroundjoin%
\definecolor{currentfill}{rgb}{0.850827,0.844675,0.910804}%
\pgfsetfillcolor{currentfill}%
\pgfsetfillopacity{0.900000}%
\pgfsetlinewidth{0.200750pt}%
\definecolor{currentstroke}{rgb}{0.200000,0.200000,0.200000}%
\pgfsetstrokecolor{currentstroke}%
\pgfsetdash{}{0pt}%
\pgfpathmoveto{\pgfqpoint{2.235898in}{2.223838in}}%
\pgfpathlineto{\pgfqpoint{2.278946in}{2.315319in}}%
\pgfpathlineto{\pgfqpoint{2.321816in}{2.455364in}}%
\pgfpathlineto{\pgfqpoint{2.351682in}{2.457758in}}%
\pgfpathlineto{\pgfqpoint{2.381664in}{2.458313in}}%
\pgfpathlineto{\pgfqpoint{2.338548in}{2.321758in}}%
\pgfpathlineto{\pgfqpoint{2.295350in}{2.223053in}}%
\pgfpathlineto{\pgfqpoint{2.265566in}{2.224206in}}%
\pgfpathlineto{\pgfqpoint{2.235898in}{2.223838in}}%
\pgfpathclose%
\pgfusepath{stroke,fill}%
\end{pgfscope}%
\begin{pgfscope}%
\pgfpathrectangle{\pgfqpoint{0.050000in}{0.050000in}}{\pgfqpoint{4.900000in}{4.900000in}}%
\pgfusepath{clip}%
\pgfsetbuttcap%
\pgfsetroundjoin%
\definecolor{currentfill}{rgb}{1.000000,1.000000,1.000000}%
\pgfsetfillcolor{currentfill}%
\pgfsetfillopacity{0.900000}%
\pgfsetlinewidth{0.200750pt}%
\definecolor{currentstroke}{rgb}{0.200000,0.200000,0.200000}%
\pgfsetstrokecolor{currentstroke}%
\pgfsetdash{}{0pt}%
\pgfpathmoveto{\pgfqpoint{2.061694in}{2.116898in}}%
\pgfpathlineto{\pgfqpoint{2.105455in}{2.130707in}}%
\pgfpathlineto{\pgfqpoint{2.149107in}{2.147167in}}%
\pgfpathlineto{\pgfqpoint{2.178750in}{2.133126in}}%
\pgfpathlineto{\pgfqpoint{2.208479in}{2.120307in}}%
\pgfpathlineto{\pgfqpoint{2.164798in}{2.096636in}}%
\pgfpathlineto{\pgfqpoint{2.120988in}{2.079892in}}%
\pgfpathlineto{\pgfqpoint{2.091298in}{2.098216in}}%
\pgfpathlineto{\pgfqpoint{2.061694in}{2.116898in}}%
\pgfpathclose%
\pgfusepath{stroke,fill}%
\end{pgfscope}%
\begin{pgfscope}%
\pgfpathrectangle{\pgfqpoint{0.050000in}{0.050000in}}{\pgfqpoint{4.900000in}{4.900000in}}%
\pgfusepath{clip}%
\pgfsetbuttcap%
\pgfsetroundjoin%
\definecolor{currentfill}{rgb}{0.928397,0.925444,0.957186}%
\pgfsetfillcolor{currentfill}%
\pgfsetfillopacity{0.900000}%
\pgfsetlinewidth{0.200750pt}%
\definecolor{currentstroke}{rgb}{0.200000,0.200000,0.200000}%
\pgfsetstrokecolor{currentstroke}%
\pgfsetdash{}{0pt}%
\pgfpathmoveto{\pgfqpoint{3.206057in}{2.307624in}}%
\pgfpathlineto{\pgfqpoint{3.246689in}{2.151246in}}%
\pgfpathlineto{\pgfqpoint{3.277985in}{2.131246in}}%
\pgfpathlineto{\pgfqpoint{3.309374in}{2.111292in}}%
\pgfpathlineto{\pgfqpoint{3.268485in}{2.236444in}}%
\pgfpathlineto{\pgfqpoint{3.237224in}{2.271257in}}%
\pgfpathlineto{\pgfqpoint{3.206057in}{2.307624in}}%
\pgfpathclose%
\pgfusepath{stroke,fill}%
\end{pgfscope}%
\begin{pgfscope}%
\pgfpathrectangle{\pgfqpoint{0.050000in}{0.050000in}}{\pgfqpoint{4.900000in}{4.900000in}}%
\pgfusepath{clip}%
\pgfsetbuttcap%
\pgfsetroundjoin%
\definecolor{currentfill}{rgb}{0.606182,0.589942,0.764521}%
\pgfsetfillcolor{currentfill}%
\pgfsetfillopacity{0.900000}%
\pgfsetlinewidth{0.200750pt}%
\definecolor{currentstroke}{rgb}{0.200000,0.200000,0.200000}%
\pgfsetstrokecolor{currentstroke}%
\pgfsetdash{}{0pt}%
\pgfpathmoveto{\pgfqpoint{2.321816in}{2.455364in}}%
\pgfpathlineto{\pgfqpoint{2.364647in}{2.636162in}}%
\pgfpathlineto{\pgfqpoint{2.407576in}{2.841718in}}%
\pgfpathlineto{\pgfqpoint{2.437757in}{2.832179in}}%
\pgfpathlineto{\pgfqpoint{2.468035in}{2.823629in}}%
\pgfpathlineto{\pgfqpoint{2.424795in}{2.628593in}}%
\pgfpathlineto{\pgfqpoint{2.381664in}{2.458313in}}%
\pgfpathlineto{\pgfqpoint{2.351682in}{2.457758in}}%
\pgfpathlineto{\pgfqpoint{2.321816in}{2.455364in}}%
\pgfpathclose%
\pgfusepath{stroke,fill}%
\end{pgfscope}%
\begin{pgfscope}%
\pgfpathrectangle{\pgfqpoint{0.050000in}{0.050000in}}{\pgfqpoint{4.900000in}{4.900000in}}%
\pgfusepath{clip}%
\pgfsetbuttcap%
\pgfsetroundjoin%
\definecolor{currentfill}{rgb}{0.946298,0.944083,0.967889}%
\pgfsetfillcolor{currentfill}%
\pgfsetfillopacity{0.900000}%
\pgfsetlinewidth{0.200750pt}%
\definecolor{currentstroke}{rgb}{0.200000,0.200000,0.200000}%
\pgfsetstrokecolor{currentstroke}%
\pgfsetdash{}{0pt}%
\pgfpathmoveto{\pgfqpoint{2.208479in}{2.120307in}}%
\pgfpathlineto{\pgfqpoint{2.252001in}{2.158941in}}%
\pgfpathlineto{\pgfqpoint{2.295350in}{2.223053in}}%
\pgfpathlineto{\pgfqpoint{2.325246in}{2.220676in}}%
\pgfpathlineto{\pgfqpoint{2.355252in}{2.217296in}}%
\pgfpathlineto{\pgfqpoint{2.311795in}{2.144680in}}%
\pgfpathlineto{\pgfqpoint{2.268209in}{2.096690in}}%
\pgfpathlineto{\pgfqpoint{2.238297in}{2.108260in}}%
\pgfpathlineto{\pgfqpoint{2.208479in}{2.120307in}}%
\pgfpathclose%
\pgfusepath{stroke,fill}%
\end{pgfscope}%
\begin{pgfscope}%
\pgfpathrectangle{\pgfqpoint{0.050000in}{0.050000in}}{\pgfqpoint{4.900000in}{4.900000in}}%
\pgfusepath{clip}%
\pgfsetbuttcap%
\pgfsetroundjoin%
\definecolor{currentfill}{rgb}{1.000000,1.000000,1.000000}%
\pgfsetfillcolor{currentfill}%
\pgfsetfillopacity{0.900000}%
\pgfsetlinewidth{0.200750pt}%
\definecolor{currentstroke}{rgb}{0.200000,0.200000,0.200000}%
\pgfsetstrokecolor{currentstroke}%
\pgfsetdash{}{0pt}%
\pgfpathmoveto{\pgfqpoint{1.973890in}{2.090225in}}%
\pgfpathlineto{\pgfqpoint{2.017838in}{2.103547in}}%
\pgfpathlineto{\pgfqpoint{2.061694in}{2.116898in}}%
\pgfpathlineto{\pgfqpoint{2.091298in}{2.098216in}}%
\pgfpathlineto{\pgfqpoint{2.120988in}{2.079892in}}%
\pgfpathlineto{\pgfqpoint{2.077069in}{2.065640in}}%
\pgfpathlineto{\pgfqpoint{2.033054in}{2.052065in}}%
\pgfpathlineto{\pgfqpoint{2.003428in}{2.071162in}}%
\pgfpathlineto{\pgfqpoint{1.973890in}{2.090225in}}%
\pgfpathclose%
\pgfusepath{stroke,fill}%
\end{pgfscope}%
\begin{pgfscope}%
\pgfpathrectangle{\pgfqpoint{0.050000in}{0.050000in}}{\pgfqpoint{4.900000in}{4.900000in}}%
\pgfusepath{clip}%
\pgfsetbuttcap%
\pgfsetroundjoin%
\definecolor{currentfill}{rgb}{0.820992,0.813610,0.892964}%
\pgfsetfillcolor{currentfill}%
\pgfsetfillopacity{0.900000}%
\pgfsetlinewidth{0.200750pt}%
\definecolor{currentstroke}{rgb}{0.200000,0.200000,0.200000}%
\pgfsetstrokecolor{currentstroke}%
\pgfsetdash{}{0pt}%
\pgfpathmoveto{\pgfqpoint{2.295350in}{2.223053in}}%
\pgfpathlineto{\pgfqpoint{2.338548in}{2.321758in}}%
\pgfpathlineto{\pgfqpoint{2.381664in}{2.458313in}}%
\pgfpathlineto{\pgfqpoint{2.411759in}{2.457442in}}%
\pgfpathlineto{\pgfqpoint{2.441965in}{2.455416in}}%
\pgfpathlineto{\pgfqpoint{2.398620in}{2.320290in}}%
\pgfpathlineto{\pgfqpoint{2.355252in}{2.217296in}}%
\pgfpathlineto{\pgfqpoint{2.325246in}{2.220676in}}%
\pgfpathlineto{\pgfqpoint{2.295350in}{2.223053in}}%
\pgfpathclose%
\pgfusepath{stroke,fill}%
\end{pgfscope}%
\begin{pgfscope}%
\pgfpathrectangle{\pgfqpoint{0.050000in}{0.050000in}}{\pgfqpoint{4.900000in}{4.900000in}}%
\pgfusepath{clip}%
\pgfsetbuttcap%
\pgfsetroundjoin%
\definecolor{currentfill}{rgb}{0.343637,0.316571,0.607536}%
\pgfsetfillcolor{currentfill}%
\pgfsetfillopacity{0.900000}%
\pgfsetlinewidth{0.200750pt}%
\definecolor{currentstroke}{rgb}{0.200000,0.200000,0.200000}%
\pgfsetstrokecolor{currentstroke}%
\pgfsetdash{}{0pt}%
\pgfpathmoveto{\pgfqpoint{2.878892in}{3.473335in}}%
\pgfpathlineto{\pgfqpoint{2.920066in}{3.127205in}}%
\pgfpathlineto{\pgfqpoint{2.960151in}{2.714228in}}%
\pgfpathlineto{\pgfqpoint{2.990541in}{2.643173in}}%
\pgfpathlineto{\pgfqpoint{3.021037in}{2.581110in}}%
\pgfpathlineto{\pgfqpoint{2.980871in}{2.917651in}}%
\pgfpathlineto{\pgfqpoint{2.939890in}{3.223308in}}%
\pgfpathlineto{\pgfqpoint{2.909359in}{3.340994in}}%
\pgfpathlineto{\pgfqpoint{2.878892in}{3.473335in}}%
\pgfpathclose%
\pgfusepath{stroke,fill}%
\end{pgfscope}%
\begin{pgfscope}%
\pgfpathrectangle{\pgfqpoint{0.050000in}{0.050000in}}{\pgfqpoint{4.900000in}{4.900000in}}%
\pgfusepath{clip}%
\pgfsetbuttcap%
\pgfsetroundjoin%
\definecolor{currentfill}{rgb}{0.241799,0.235294,0.581592}%
\pgfsetfillcolor{currentfill}%
\pgfsetfillopacity{0.900000}%
\pgfsetlinewidth{0.200750pt}%
\definecolor{currentstroke}{rgb}{0.200000,0.200000,0.200000}%
\pgfsetstrokecolor{currentstroke}%
\pgfsetdash{}{0pt}%
\pgfpathmoveto{\pgfqpoint{2.818119in}{3.796863in}}%
\pgfpathlineto{\pgfqpoint{2.859666in}{3.418071in}}%
\pgfpathlineto{\pgfqpoint{2.899726in}{2.900065in}}%
\pgfpathlineto{\pgfqpoint{2.929875in}{2.797871in}}%
\pgfpathlineto{\pgfqpoint{2.960151in}{2.714228in}}%
\pgfpathlineto{\pgfqpoint{2.920066in}{3.127205in}}%
\pgfpathlineto{\pgfqpoint{2.878892in}{3.473335in}}%
\pgfpathlineto{\pgfqpoint{2.848482in}{3.623838in}}%
\pgfpathlineto{\pgfqpoint{2.818119in}{3.796863in}}%
\pgfpathclose%
\pgfusepath{stroke,fill}%
\end{pgfscope}%
\begin{pgfscope}%
\pgfpathrectangle{\pgfqpoint{0.050000in}{0.050000in}}{\pgfqpoint{4.900000in}{4.900000in}}%
\pgfusepath{clip}%
\pgfsetbuttcap%
\pgfsetroundjoin%
\definecolor{currentfill}{rgb}{0.994033,0.993787,0.996432}%
\pgfsetfillcolor{currentfill}%
\pgfsetfillopacity{0.900000}%
\pgfsetlinewidth{0.200750pt}%
\definecolor{currentstroke}{rgb}{0.200000,0.200000,0.200000}%
\pgfsetstrokecolor{currentstroke}%
\pgfsetdash{}{0pt}%
\pgfpathmoveto{\pgfqpoint{2.120988in}{2.079892in}}%
\pgfpathlineto{\pgfqpoint{2.164798in}{2.096636in}}%
\pgfpathlineto{\pgfqpoint{2.208479in}{2.120307in}}%
\pgfpathlineto{\pgfqpoint{2.238297in}{2.108260in}}%
\pgfpathlineto{\pgfqpoint{2.268209in}{2.096690in}}%
\pgfpathlineto{\pgfqpoint{2.224482in}{2.065683in}}%
\pgfpathlineto{\pgfqpoint{2.180624in}{2.044566in}}%
\pgfpathlineto{\pgfqpoint{2.150762in}{2.062005in}}%
\pgfpathlineto{\pgfqpoint{2.120988in}{2.079892in}}%
\pgfpathclose%
\pgfusepath{stroke,fill}%
\end{pgfscope}%
\begin{pgfscope}%
\pgfpathrectangle{\pgfqpoint{0.050000in}{0.050000in}}{\pgfqpoint{4.900000in}{4.900000in}}%
\pgfusepath{clip}%
\pgfsetbuttcap%
\pgfsetroundjoin%
\definecolor{currentfill}{rgb}{0.462976,0.440830,0.678893}%
\pgfsetfillcolor{currentfill}%
\pgfsetfillopacity{0.900000}%
\pgfsetlinewidth{0.200750pt}%
\definecolor{currentstroke}{rgb}{0.200000,0.200000,0.200000}%
\pgfsetstrokecolor{currentstroke}%
\pgfsetdash{}{0pt}%
\pgfpathmoveto{\pgfqpoint{2.939890in}{3.223308in}}%
\pgfpathlineto{\pgfqpoint{2.980871in}{2.917651in}}%
\pgfpathlineto{\pgfqpoint{3.021037in}{2.581110in}}%
\pgfpathlineto{\pgfqpoint{3.051634in}{2.525727in}}%
\pgfpathlineto{\pgfqpoint{3.082329in}{2.475463in}}%
\pgfpathlineto{\pgfqpoint{3.042049in}{2.753898in}}%
\pgfpathlineto{\pgfqpoint{3.001161in}{3.021593in}}%
\pgfpathlineto{\pgfqpoint{2.970490in}{3.117546in}}%
\pgfpathlineto{\pgfqpoint{2.939890in}{3.223308in}}%
\pgfpathclose%
\pgfusepath{stroke,fill}%
\end{pgfscope}%
\begin{pgfscope}%
\pgfpathrectangle{\pgfqpoint{0.050000in}{0.050000in}}{\pgfqpoint{4.900000in}{4.900000in}}%
\pgfusepath{clip}%
\pgfsetbuttcap%
\pgfsetroundjoin%
\definecolor{currentfill}{rgb}{0.940331,0.937870,0.964321}%
\pgfsetfillcolor{currentfill}%
\pgfsetfillopacity{0.900000}%
\pgfsetlinewidth{0.200750pt}%
\definecolor{currentstroke}{rgb}{0.200000,0.200000,0.200000}%
\pgfsetstrokecolor{currentstroke}%
\pgfsetdash{}{0pt}%
\pgfpathmoveto{\pgfqpoint{3.268485in}{2.236444in}}%
\pgfpathlineto{\pgfqpoint{3.309374in}{2.111292in}}%
\pgfpathlineto{\pgfqpoint{3.340856in}{2.091370in}}%
\pgfpathlineto{\pgfqpoint{3.372432in}{2.071467in}}%
\pgfpathlineto{\pgfqpoint{3.331288in}{2.170632in}}%
\pgfpathlineto{\pgfqpoint{3.299839in}{2.202962in}}%
\pgfpathlineto{\pgfqpoint{3.268485in}{2.236444in}}%
\pgfpathclose%
\pgfusepath{stroke,fill}%
\end{pgfscope}%
\begin{pgfscope}%
\pgfpathrectangle{\pgfqpoint{0.050000in}{0.050000in}}{\pgfqpoint{4.900000in}{4.900000in}}%
\pgfusepath{clip}%
\pgfsetbuttcap%
\pgfsetroundjoin%
\definecolor{currentfill}{rgb}{1.000000,1.000000,1.000000}%
\pgfsetfillcolor{currentfill}%
\pgfsetfillopacity{0.900000}%
\pgfsetlinewidth{0.200750pt}%
\definecolor{currentstroke}{rgb}{0.200000,0.200000,0.200000}%
\pgfsetstrokecolor{currentstroke}%
\pgfsetdash{}{0pt}%
\pgfpathmoveto{\pgfqpoint{1.885720in}{2.063507in}}%
\pgfpathlineto{\pgfqpoint{1.929851in}{2.076880in}}%
\pgfpathlineto{\pgfqpoint{1.973890in}{2.090225in}}%
\pgfpathlineto{\pgfqpoint{2.003428in}{2.071162in}}%
\pgfpathlineto{\pgfqpoint{2.033054in}{2.052065in}}%
\pgfpathlineto{\pgfqpoint{1.988946in}{2.038613in}}%
\pgfpathlineto{\pgfqpoint{1.944745in}{2.025157in}}%
\pgfpathlineto{\pgfqpoint{1.915189in}{2.044360in}}%
\pgfpathlineto{\pgfqpoint{1.885720in}{2.063507in}}%
\pgfpathclose%
\pgfusepath{stroke,fill}%
\end{pgfscope}%
\begin{pgfscope}%
\pgfpathrectangle{\pgfqpoint{0.050000in}{0.050000in}}{\pgfqpoint{4.900000in}{4.900000in}}%
\pgfusepath{clip}%
\pgfsetbuttcap%
\pgfsetroundjoin%
\definecolor{currentfill}{rgb}{0.552480,0.534025,0.732411}%
\pgfsetfillcolor{currentfill}%
\pgfsetfillopacity{0.900000}%
\pgfsetlinewidth{0.200750pt}%
\definecolor{currentstroke}{rgb}{0.200000,0.200000,0.200000}%
\pgfsetstrokecolor{currentstroke}%
\pgfsetdash{}{0pt}%
\pgfpathmoveto{\pgfqpoint{3.001161in}{3.021593in}}%
\pgfpathlineto{\pgfqpoint{3.042049in}{2.753898in}}%
\pgfpathlineto{\pgfqpoint{3.082329in}{2.475463in}}%
\pgfpathlineto{\pgfqpoint{3.113119in}{2.429221in}}%
\pgfpathlineto{\pgfqpoint{3.144004in}{2.386205in}}%
\pgfpathlineto{\pgfqpoint{3.103591in}{2.618788in}}%
\pgfpathlineto{\pgfqpoint{3.062731in}{2.852868in}}%
\pgfpathlineto{\pgfqpoint{3.031908in}{2.933800in}}%
\pgfpathlineto{\pgfqpoint{3.001161in}{3.021593in}}%
\pgfpathclose%
\pgfusepath{stroke,fill}%
\end{pgfscope}%
\begin{pgfscope}%
\pgfpathrectangle{\pgfqpoint{0.050000in}{0.050000in}}{\pgfqpoint{4.900000in}{4.900000in}}%
\pgfusepath{clip}%
\pgfsetbuttcap%
\pgfsetroundjoin%
\definecolor{currentfill}{rgb}{1.000000,1.000000,1.000000}%
\pgfsetfillcolor{currentfill}%
\pgfsetfillopacity{0.900000}%
\pgfsetlinewidth{0.200750pt}%
\definecolor{currentstroke}{rgb}{0.200000,0.200000,0.200000}%
\pgfsetstrokecolor{currentstroke}%
\pgfsetdash{}{0pt}%
\pgfpathmoveto{\pgfqpoint{2.033054in}{2.052065in}}%
\pgfpathlineto{\pgfqpoint{2.077069in}{2.065640in}}%
\pgfpathlineto{\pgfqpoint{2.120988in}{2.079892in}}%
\pgfpathlineto{\pgfqpoint{2.150762in}{2.062005in}}%
\pgfpathlineto{\pgfqpoint{2.180624in}{2.044566in}}%
\pgfpathlineto{\pgfqpoint{2.136649in}{2.028264in}}%
\pgfpathlineto{\pgfqpoint{2.092569in}{2.013901in}}%
\pgfpathlineto{\pgfqpoint{2.062768in}{2.032965in}}%
\pgfpathlineto{\pgfqpoint{2.033054in}{2.052065in}}%
\pgfpathclose%
\pgfusepath{stroke,fill}%
\end{pgfscope}%
\begin{pgfscope}%
\pgfpathrectangle{\pgfqpoint{0.050000in}{0.050000in}}{\pgfqpoint{4.900000in}{4.900000in}}%
\pgfusepath{clip}%
\pgfsetbuttcap%
\pgfsetroundjoin%
\definecolor{currentfill}{rgb}{0.252134,0.345098,0.727566}%
\pgfsetfillcolor{currentfill}%
\pgfsetfillopacity{0.900000}%
\pgfsetlinewidth{0.200750pt}%
\definecolor{currentstroke}{rgb}{0.200000,0.200000,0.200000}%
\pgfsetstrokecolor{currentstroke}%
\pgfsetdash{}{0pt}%
\pgfpathmoveto{\pgfqpoint{2.757446in}{4.226936in}}%
\pgfpathlineto{\pgfqpoint{2.799762in}{3.877558in}}%
\pgfpathlineto{\pgfqpoint{2.839953in}{3.214644in}}%
\pgfpathlineto{\pgfqpoint{2.869733in}{3.031607in}}%
\pgfpathlineto{\pgfqpoint{2.899726in}{2.900065in}}%
\pgfpathlineto{\pgfqpoint{2.859666in}{3.418071in}}%
\pgfpathlineto{\pgfqpoint{2.818119in}{3.796863in}}%
\pgfpathlineto{\pgfqpoint{2.787786in}{3.997155in}}%
\pgfpathlineto{\pgfqpoint{2.757446in}{4.226936in}}%
\pgfpathclose%
\pgfusepath{stroke,fill}%
\end{pgfscope}%
\begin{pgfscope}%
\pgfpathrectangle{\pgfqpoint{0.050000in}{0.050000in}}{\pgfqpoint{4.900000in}{4.900000in}}%
\pgfusepath{clip}%
\pgfsetbuttcap%
\pgfsetroundjoin%
\definecolor{currentfill}{rgb}{0.582314,0.565090,0.750250}%
\pgfsetfillcolor{currentfill}%
\pgfsetfillopacity{0.900000}%
\pgfsetlinewidth{0.200750pt}%
\definecolor{currentstroke}{rgb}{0.200000,0.200000,0.200000}%
\pgfsetstrokecolor{currentstroke}%
\pgfsetdash{}{0pt}%
\pgfpathmoveto{\pgfqpoint{2.381664in}{2.458313in}}%
\pgfpathlineto{\pgfqpoint{2.424795in}{2.628593in}}%
\pgfpathlineto{\pgfqpoint{2.468035in}{2.823629in}}%
\pgfpathlineto{\pgfqpoint{2.498414in}{2.815523in}}%
\pgfpathlineto{\pgfqpoint{2.528893in}{2.807408in}}%
\pgfpathlineto{\pgfqpoint{2.485366in}{2.619863in}}%
\pgfpathlineto{\pgfqpoint{2.441965in}{2.455416in}}%
\pgfpathlineto{\pgfqpoint{2.411759in}{2.457442in}}%
\pgfpathlineto{\pgfqpoint{2.381664in}{2.458313in}}%
\pgfpathclose%
\pgfusepath{stroke,fill}%
\end{pgfscope}%
\begin{pgfscope}%
\pgfpathrectangle{\pgfqpoint{0.050000in}{0.050000in}}{\pgfqpoint{4.900000in}{4.900000in}}%
\pgfusepath{clip}%
\pgfsetbuttcap%
\pgfsetroundjoin%
\definecolor{currentfill}{rgb}{0.928397,0.925444,0.957186}%
\pgfsetfillcolor{currentfill}%
\pgfsetfillopacity{0.900000}%
\pgfsetlinewidth{0.200750pt}%
\definecolor{currentstroke}{rgb}{0.200000,0.200000,0.200000}%
\pgfsetstrokecolor{currentstroke}%
\pgfsetdash{}{0pt}%
\pgfpathmoveto{\pgfqpoint{2.268209in}{2.096690in}}%
\pgfpathlineto{\pgfqpoint{2.311795in}{2.144680in}}%
\pgfpathlineto{\pgfqpoint{2.355252in}{2.217296in}}%
\pgfpathlineto{\pgfqpoint{2.385367in}{2.213074in}}%
\pgfpathlineto{\pgfqpoint{2.415589in}{2.208129in}}%
\pgfpathlineto{\pgfqpoint{2.372000in}{2.129488in}}%
\pgfpathlineto{\pgfqpoint{2.328321in}{2.074275in}}%
\pgfpathlineto{\pgfqpoint{2.298217in}{2.085403in}}%
\pgfpathlineto{\pgfqpoint{2.268209in}{2.096690in}}%
\pgfpathclose%
\pgfusepath{stroke,fill}%
\end{pgfscope}%
\begin{pgfscope}%
\pgfpathrectangle{\pgfqpoint{0.050000in}{0.050000in}}{\pgfqpoint{4.900000in}{4.900000in}}%
\pgfusepath{clip}%
\pgfsetbuttcap%
\pgfsetroundjoin%
\definecolor{currentfill}{rgb}{0.630050,0.614794,0.778793}%
\pgfsetfillcolor{currentfill}%
\pgfsetfillopacity{0.900000}%
\pgfsetlinewidth{0.200750pt}%
\definecolor{currentstroke}{rgb}{0.200000,0.200000,0.200000}%
\pgfsetstrokecolor{currentstroke}%
\pgfsetdash{}{0pt}%
\pgfpathmoveto{\pgfqpoint{3.062731in}{2.852868in}}%
\pgfpathlineto{\pgfqpoint{3.103591in}{2.618788in}}%
\pgfpathlineto{\pgfqpoint{3.144004in}{2.386205in}}%
\pgfpathlineto{\pgfqpoint{3.174984in}{2.345824in}}%
\pgfpathlineto{\pgfqpoint{3.206057in}{2.307624in}}%
\pgfpathlineto{\pgfqpoint{3.165493in}{2.502954in}}%
\pgfpathlineto{\pgfqpoint{3.124617in}{2.707654in}}%
\pgfpathlineto{\pgfqpoint{3.093634in}{2.777764in}}%
\pgfpathlineto{\pgfqpoint{3.062731in}{2.852868in}}%
\pgfpathclose%
\pgfusepath{stroke,fill}%
\end{pgfscope}%
\begin{pgfscope}%
\pgfpathrectangle{\pgfqpoint{0.050000in}{0.050000in}}{\pgfqpoint{4.900000in}{4.900000in}}%
\pgfusepath{clip}%
\pgfsetbuttcap%
\pgfsetroundjoin%
\definecolor{currentfill}{rgb}{0.791157,0.782545,0.875125}%
\pgfsetfillcolor{currentfill}%
\pgfsetfillopacity{0.900000}%
\pgfsetlinewidth{0.200750pt}%
\definecolor{currentstroke}{rgb}{0.200000,0.200000,0.200000}%
\pgfsetstrokecolor{currentstroke}%
\pgfsetdash{}{0pt}%
\pgfpathmoveto{\pgfqpoint{2.355252in}{2.217296in}}%
\pgfpathlineto{\pgfqpoint{2.398620in}{2.320290in}}%
\pgfpathlineto{\pgfqpoint{2.441965in}{2.455416in}}%
\pgfpathlineto{\pgfqpoint{2.472281in}{2.452407in}}%
\pgfpathlineto{\pgfqpoint{2.502705in}{2.448505in}}%
\pgfpathlineto{\pgfqpoint{2.459137in}{2.314151in}}%
\pgfpathlineto{\pgfqpoint{2.415589in}{2.208129in}}%
\pgfpathlineto{\pgfqpoint{2.385367in}{2.213074in}}%
\pgfpathlineto{\pgfqpoint{2.355252in}{2.217296in}}%
\pgfpathclose%
\pgfusepath{stroke,fill}%
\end{pgfscope}%
\begin{pgfscope}%
\pgfpathrectangle{\pgfqpoint{0.050000in}{0.050000in}}{\pgfqpoint{4.900000in}{4.900000in}}%
\pgfusepath{clip}%
\pgfsetbuttcap%
\pgfsetroundjoin%
\definecolor{currentfill}{rgb}{1.000000,1.000000,1.000000}%
\pgfsetfillcolor{currentfill}%
\pgfsetfillopacity{0.900000}%
\pgfsetlinewidth{0.200750pt}%
\definecolor{currentstroke}{rgb}{0.200000,0.200000,0.200000}%
\pgfsetstrokecolor{currentstroke}%
\pgfsetdash{}{0pt}%
\pgfpathmoveto{\pgfqpoint{1.797182in}{2.036678in}}%
\pgfpathlineto{\pgfqpoint{1.841498in}{2.050106in}}%
\pgfpathlineto{\pgfqpoint{1.885720in}{2.063507in}}%
\pgfpathlineto{\pgfqpoint{1.915189in}{2.044360in}}%
\pgfpathlineto{\pgfqpoint{1.944745in}{2.025157in}}%
\pgfpathlineto{\pgfqpoint{1.900452in}{2.011676in}}%
\pgfpathlineto{\pgfqpoint{1.856066in}{1.998167in}}%
\pgfpathlineto{\pgfqpoint{1.826580in}{2.017451in}}%
\pgfpathlineto{\pgfqpoint{1.797182in}{2.036678in}}%
\pgfpathclose%
\pgfusepath{stroke,fill}%
\end{pgfscope}%
\begin{pgfscope}%
\pgfpathrectangle{\pgfqpoint{0.050000in}{0.050000in}}{\pgfqpoint{4.900000in}{4.900000in}}%
\pgfusepath{clip}%
\pgfsetbuttcap%
\pgfsetroundjoin%
\definecolor{currentfill}{rgb}{0.982099,0.981361,0.989296}%
\pgfsetfillcolor{currentfill}%
\pgfsetfillopacity{0.900000}%
\pgfsetlinewidth{0.200750pt}%
\definecolor{currentstroke}{rgb}{0.200000,0.200000,0.200000}%
\pgfsetstrokecolor{currentstroke}%
\pgfsetdash{}{0pt}%
\pgfpathmoveto{\pgfqpoint{2.180624in}{2.044566in}}%
\pgfpathlineto{\pgfqpoint{2.224482in}{2.065683in}}%
\pgfpathlineto{\pgfqpoint{2.268209in}{2.096690in}}%
\pgfpathlineto{\pgfqpoint{2.298217in}{2.085403in}}%
\pgfpathlineto{\pgfqpoint{2.328321in}{2.074275in}}%
\pgfpathlineto{\pgfqpoint{2.284528in}{2.036737in}}%
\pgfpathlineto{\pgfqpoint{2.240612in}{2.010890in}}%
\pgfpathlineto{\pgfqpoint{2.210573in}{2.027542in}}%
\pgfpathlineto{\pgfqpoint{2.180624in}{2.044566in}}%
\pgfpathclose%
\pgfusepath{stroke,fill}%
\end{pgfscope}%
\begin{pgfscope}%
\pgfpathrectangle{\pgfqpoint{0.050000in}{0.050000in}}{\pgfqpoint{4.900000in}{4.900000in}}%
\pgfusepath{clip}%
\pgfsetbuttcap%
\pgfsetroundjoin%
\definecolor{currentfill}{rgb}{0.278001,0.248228,0.568289}%
\pgfsetfillcolor{currentfill}%
\pgfsetfillopacity{0.900000}%
\pgfsetlinewidth{0.200750pt}%
\definecolor{currentstroke}{rgb}{0.200000,0.200000,0.200000}%
\pgfsetstrokecolor{currentstroke}%
\pgfsetdash{}{0pt}%
\pgfpathmoveto{\pgfqpoint{2.407576in}{2.841718in}}%
\pgfpathlineto{\pgfqpoint{2.450686in}{3.058596in}}%
\pgfpathlineto{\pgfqpoint{2.494012in}{3.279828in}}%
\pgfpathlineto{\pgfqpoint{2.524494in}{3.265731in}}%
\pgfpathlineto{\pgfqpoint{2.555072in}{3.251529in}}%
\pgfpathlineto{\pgfqpoint{2.511449in}{3.033812in}}%
\pgfpathlineto{\pgfqpoint{2.468035in}{2.823629in}}%
\pgfpathlineto{\pgfqpoint{2.437757in}{2.832179in}}%
\pgfpathlineto{\pgfqpoint{2.407576in}{2.841718in}}%
\pgfpathclose%
\pgfusepath{stroke,fill}%
\end{pgfscope}%
\begin{pgfscope}%
\pgfpathrectangle{\pgfqpoint{0.050000in}{0.050000in}}{\pgfqpoint{4.900000in}{4.900000in}}%
\pgfusepath{clip}%
\pgfsetbuttcap%
\pgfsetroundjoin%
\definecolor{currentfill}{rgb}{0.952265,0.950296,0.971457}%
\pgfsetfillcolor{currentfill}%
\pgfsetfillopacity{0.900000}%
\pgfsetlinewidth{0.200750pt}%
\definecolor{currentstroke}{rgb}{0.200000,0.200000,0.200000}%
\pgfsetstrokecolor{currentstroke}%
\pgfsetdash{}{0pt}%
\pgfpathmoveto{\pgfqpoint{3.331288in}{2.170632in}}%
\pgfpathlineto{\pgfqpoint{3.372432in}{2.071467in}}%
\pgfpathlineto{\pgfqpoint{3.404102in}{2.051575in}}%
\pgfpathlineto{\pgfqpoint{3.435866in}{2.031684in}}%
\pgfpathlineto{\pgfqpoint{3.394468in}{2.108853in}}%
\pgfpathlineto{\pgfqpoint{3.362831in}{2.139304in}}%
\pgfpathlineto{\pgfqpoint{3.331288in}{2.170632in}}%
\pgfpathclose%
\pgfusepath{stroke,fill}%
\end{pgfscope}%
\begin{pgfscope}%
\pgfpathrectangle{\pgfqpoint{0.050000in}{0.050000in}}{\pgfqpoint{4.900000in}{4.900000in}}%
\pgfusepath{clip}%
\pgfsetbuttcap%
\pgfsetroundjoin%
\definecolor{currentfill}{rgb}{1.000000,1.000000,1.000000}%
\pgfsetfillcolor{currentfill}%
\pgfsetfillopacity{0.900000}%
\pgfsetlinewidth{0.200750pt}%
\definecolor{currentstroke}{rgb}{0.200000,0.200000,0.200000}%
\pgfsetstrokecolor{currentstroke}%
\pgfsetdash{}{0pt}%
\pgfpathmoveto{\pgfqpoint{1.944745in}{2.025157in}}%
\pgfpathlineto{\pgfqpoint{1.988946in}{2.038613in}}%
\pgfpathlineto{\pgfqpoint{2.033054in}{2.052065in}}%
\pgfpathlineto{\pgfqpoint{2.062768in}{2.032965in}}%
\pgfpathlineto{\pgfqpoint{2.092569in}{2.013901in}}%
\pgfpathlineto{\pgfqpoint{2.048393in}{2.000168in}}%
\pgfpathlineto{\pgfqpoint{2.004123in}{1.986586in}}%
\pgfpathlineto{\pgfqpoint{1.974389in}{2.005898in}}%
\pgfpathlineto{\pgfqpoint{1.944745in}{2.025157in}}%
\pgfpathclose%
\pgfusepath{stroke,fill}%
\end{pgfscope}%
\begin{pgfscope}%
\pgfpathrectangle{\pgfqpoint{0.050000in}{0.050000in}}{\pgfqpoint{4.900000in}{4.900000in}}%
\pgfusepath{clip}%
\pgfsetbuttcap%
\pgfsetroundjoin%
\definecolor{currentfill}{rgb}{0.689719,0.676924,0.814471}%
\pgfsetfillcolor{currentfill}%
\pgfsetfillopacity{0.900000}%
\pgfsetlinewidth{0.200750pt}%
\definecolor{currentstroke}{rgb}{0.200000,0.200000,0.200000}%
\pgfsetstrokecolor{currentstroke}%
\pgfsetdash{}{0pt}%
\pgfpathmoveto{\pgfqpoint{3.124617in}{2.707654in}}%
\pgfpathlineto{\pgfqpoint{3.165493in}{2.502954in}}%
\pgfpathlineto{\pgfqpoint{3.206057in}{2.307624in}}%
\pgfpathlineto{\pgfqpoint{3.237224in}{2.271257in}}%
\pgfpathlineto{\pgfqpoint{3.268485in}{2.236444in}}%
\pgfpathlineto{\pgfqpoint{3.227756in}{2.400813in}}%
\pgfpathlineto{\pgfqpoint{3.186830in}{2.579828in}}%
\pgfpathlineto{\pgfqpoint{3.155682in}{2.641861in}}%
\pgfpathlineto{\pgfqpoint{3.124617in}{2.707654in}}%
\pgfpathclose%
\pgfusepath{stroke,fill}%
\end{pgfscope}%
\begin{pgfscope}%
\pgfpathrectangle{\pgfqpoint{0.050000in}{0.050000in}}{\pgfqpoint{4.900000in}{4.900000in}}%
\pgfusepath{clip}%
\pgfsetbuttcap%
\pgfsetroundjoin%
\definecolor{currentfill}{rgb}{1.000000,1.000000,1.000000}%
\pgfsetfillcolor{currentfill}%
\pgfsetfillopacity{0.900000}%
\pgfsetlinewidth{0.200750pt}%
\definecolor{currentstroke}{rgb}{0.200000,0.200000,0.200000}%
\pgfsetstrokecolor{currentstroke}%
\pgfsetdash{}{0pt}%
\pgfpathmoveto{\pgfqpoint{2.092569in}{2.013901in}}%
\pgfpathlineto{\pgfqpoint{2.136649in}{2.028264in}}%
\pgfpathlineto{\pgfqpoint{2.180624in}{2.044566in}}%
\pgfpathlineto{\pgfqpoint{2.210573in}{2.027542in}}%
\pgfpathlineto{\pgfqpoint{2.240612in}{2.010890in}}%
\pgfpathlineto{\pgfqpoint{2.196579in}{1.991775in}}%
\pgfpathlineto{\pgfqpoint{2.152438in}{1.976011in}}%
\pgfpathlineto{\pgfqpoint{2.122459in}{1.994907in}}%
\pgfpathlineto{\pgfqpoint{2.092569in}{2.013901in}}%
\pgfpathclose%
\pgfusepath{stroke,fill}%
\end{pgfscope}%
\begin{pgfscope}%
\pgfpathrectangle{\pgfqpoint{0.050000in}{0.050000in}}{\pgfqpoint{4.900000in}{4.900000in}}%
\pgfusepath{clip}%
\pgfsetbuttcap%
\pgfsetroundjoin%
\definecolor{currentfill}{rgb}{1.000000,1.000000,1.000000}%
\pgfsetfillcolor{currentfill}%
\pgfsetfillopacity{0.900000}%
\pgfsetlinewidth{0.200750pt}%
\definecolor{currentstroke}{rgb}{0.200000,0.200000,0.200000}%
\pgfsetstrokecolor{currentstroke}%
\pgfsetdash{}{0pt}%
\pgfpathmoveto{\pgfqpoint{1.708274in}{2.009736in}}%
\pgfpathlineto{\pgfqpoint{1.752775in}{2.023221in}}%
\pgfpathlineto{\pgfqpoint{1.797182in}{2.036678in}}%
\pgfpathlineto{\pgfqpoint{1.826580in}{2.017451in}}%
\pgfpathlineto{\pgfqpoint{1.856066in}{1.998167in}}%
\pgfpathlineto{\pgfqpoint{1.811587in}{1.984630in}}%
\pgfpathlineto{\pgfqpoint{1.767015in}{1.971064in}}%
\pgfpathlineto{\pgfqpoint{1.737600in}{1.990429in}}%
\pgfpathlineto{\pgfqpoint{1.708274in}{2.009736in}}%
\pgfpathclose%
\pgfusepath{stroke,fill}%
\end{pgfscope}%
\begin{pgfscope}%
\pgfpathrectangle{\pgfqpoint{0.050000in}{0.050000in}}{\pgfqpoint{4.900000in}{4.900000in}}%
\pgfusepath{clip}%
\pgfsetbuttcap%
\pgfsetroundjoin%
\definecolor{currentfill}{rgb}{0.558447,0.540238,0.735978}%
\pgfsetfillcolor{currentfill}%
\pgfsetfillopacity{0.900000}%
\pgfsetlinewidth{0.200750pt}%
\definecolor{currentstroke}{rgb}{0.200000,0.200000,0.200000}%
\pgfsetstrokecolor{currentstroke}%
\pgfsetdash{}{0pt}%
\pgfpathmoveto{\pgfqpoint{2.441965in}{2.455416in}}%
\pgfpathlineto{\pgfqpoint{2.485366in}{2.619863in}}%
\pgfpathlineto{\pgfqpoint{2.528893in}{2.807408in}}%
\pgfpathlineto{\pgfqpoint{2.559475in}{2.798853in}}%
\pgfpathlineto{\pgfqpoint{2.590158in}{2.789455in}}%
\pgfpathlineto{\pgfqpoint{2.546360in}{2.608764in}}%
\pgfpathlineto{\pgfqpoint{2.502705in}{2.448505in}}%
\pgfpathlineto{\pgfqpoint{2.472281in}{2.452407in}}%
\pgfpathlineto{\pgfqpoint{2.441965in}{2.455416in}}%
\pgfpathclose%
\pgfusepath{stroke,fill}%
\end{pgfscope}%
\begin{pgfscope}%
\pgfpathrectangle{\pgfqpoint{0.050000in}{0.050000in}}{\pgfqpoint{4.900000in}{4.900000in}}%
\pgfusepath{clip}%
\pgfsetbuttcap%
\pgfsetroundjoin%
\definecolor{currentfill}{rgb}{0.910496,0.906805,0.946482}%
\pgfsetfillcolor{currentfill}%
\pgfsetfillopacity{0.900000}%
\pgfsetlinewidth{0.200750pt}%
\definecolor{currentstroke}{rgb}{0.200000,0.200000,0.200000}%
\pgfsetstrokecolor{currentstroke}%
\pgfsetdash{}{0pt}%
\pgfpathmoveto{\pgfqpoint{2.328321in}{2.074275in}}%
\pgfpathlineto{\pgfqpoint{2.372000in}{2.129488in}}%
\pgfpathlineto{\pgfqpoint{2.415589in}{2.208129in}}%
\pgfpathlineto{\pgfqpoint{2.445918in}{2.202546in}}%
\pgfpathlineto{\pgfqpoint{2.476354in}{2.196384in}}%
\pgfpathlineto{\pgfqpoint{2.432617in}{2.113169in}}%
\pgfpathlineto{\pgfqpoint{2.388826in}{2.052178in}}%
\pgfpathlineto{\pgfqpoint{2.358524in}{2.063219in}}%
\pgfpathlineto{\pgfqpoint{2.328321in}{2.074275in}}%
\pgfpathclose%
\pgfusepath{stroke,fill}%
\end{pgfscope}%
\begin{pgfscope}%
\pgfpathrectangle{\pgfqpoint{0.050000in}{0.050000in}}{\pgfqpoint{4.900000in}{4.900000in}}%
\pgfusepath{clip}%
\pgfsetbuttcap%
\pgfsetroundjoin%
\definecolor{currentfill}{rgb}{1.000000,1.000000,1.000000}%
\pgfsetfillcolor{currentfill}%
\pgfsetfillopacity{0.900000}%
\pgfsetlinewidth{0.200750pt}%
\definecolor{currentstroke}{rgb}{0.200000,0.200000,0.200000}%
\pgfsetstrokecolor{currentstroke}%
\pgfsetdash{}{0pt}%
\pgfpathmoveto{\pgfqpoint{1.856066in}{1.998167in}}%
\pgfpathlineto{\pgfqpoint{1.900452in}{2.011676in}}%
\pgfpathlineto{\pgfqpoint{1.944745in}{2.025157in}}%
\pgfpathlineto{\pgfqpoint{1.974389in}{2.005898in}}%
\pgfpathlineto{\pgfqpoint{2.004123in}{1.986586in}}%
\pgfpathlineto{\pgfqpoint{1.959759in}{1.973016in}}%
\pgfpathlineto{\pgfqpoint{1.915302in}{1.959425in}}%
\pgfpathlineto{\pgfqpoint{1.885640in}{1.978825in}}%
\pgfpathlineto{\pgfqpoint{1.856066in}{1.998167in}}%
\pgfpathclose%
\pgfusepath{stroke,fill}%
\end{pgfscope}%
\begin{pgfscope}%
\pgfpathrectangle{\pgfqpoint{0.050000in}{0.050000in}}{\pgfqpoint{4.900000in}{4.900000in}}%
\pgfusepath{clip}%
\pgfsetbuttcap%
\pgfsetroundjoin%
\definecolor{currentfill}{rgb}{0.964198,0.962722,0.978593}%
\pgfsetfillcolor{currentfill}%
\pgfsetfillopacity{0.900000}%
\pgfsetlinewidth{0.200750pt}%
\definecolor{currentstroke}{rgb}{0.200000,0.200000,0.200000}%
\pgfsetstrokecolor{currentstroke}%
\pgfsetdash{}{0pt}%
\pgfpathmoveto{\pgfqpoint{3.394468in}{2.108853in}}%
\pgfpathlineto{\pgfqpoint{3.435866in}{2.031684in}}%
\pgfpathlineto{\pgfqpoint{3.467726in}{2.011789in}}%
\pgfpathlineto{\pgfqpoint{3.499681in}{1.991882in}}%
\pgfpathlineto{\pgfqpoint{3.458029in}{2.050184in}}%
\pgfpathlineto{\pgfqpoint{3.426201in}{2.079176in}}%
\pgfpathlineto{\pgfqpoint{3.394468in}{2.108853in}}%
\pgfpathclose%
\pgfusepath{stroke,fill}%
\end{pgfscope}%
\begin{pgfscope}%
\pgfpathrectangle{\pgfqpoint{0.050000in}{0.050000in}}{\pgfqpoint{4.900000in}{4.900000in}}%
\pgfusepath{clip}%
\pgfsetbuttcap%
\pgfsetroundjoin%
\definecolor{currentfill}{rgb}{0.976132,0.975148,0.985729}%
\pgfsetfillcolor{currentfill}%
\pgfsetfillopacity{0.900000}%
\pgfsetlinewidth{0.200750pt}%
\definecolor{currentstroke}{rgb}{0.200000,0.200000,0.200000}%
\pgfsetstrokecolor{currentstroke}%
\pgfsetdash{}{0pt}%
\pgfpathmoveto{\pgfqpoint{2.240612in}{2.010890in}}%
\pgfpathlineto{\pgfqpoint{2.284528in}{2.036737in}}%
\pgfpathlineto{\pgfqpoint{2.328321in}{2.074275in}}%
\pgfpathlineto{\pgfqpoint{2.358524in}{2.063219in}}%
\pgfpathlineto{\pgfqpoint{2.388826in}{2.052178in}}%
\pgfpathlineto{\pgfqpoint{2.344948in}{2.008960in}}%
\pgfpathlineto{\pgfqpoint{2.300964in}{1.978499in}}%
\pgfpathlineto{\pgfqpoint{2.270742in}{1.994558in}}%
\pgfpathlineto{\pgfqpoint{2.240612in}{2.010890in}}%
\pgfpathclose%
\pgfusepath{stroke,fill}%
\end{pgfscope}%
\begin{pgfscope}%
\pgfpathrectangle{\pgfqpoint{0.050000in}{0.050000in}}{\pgfqpoint{4.900000in}{4.900000in}}%
\pgfusepath{clip}%
\pgfsetbuttcap%
\pgfsetroundjoin%
\definecolor{currentfill}{rgb}{0.737455,0.726628,0.843014}%
\pgfsetfillcolor{currentfill}%
\pgfsetfillopacity{0.900000}%
\pgfsetlinewidth{0.200750pt}%
\definecolor{currentstroke}{rgb}{0.200000,0.200000,0.200000}%
\pgfsetstrokecolor{currentstroke}%
\pgfsetdash{}{0pt}%
\pgfpathmoveto{\pgfqpoint{3.186830in}{2.579828in}}%
\pgfpathlineto{\pgfqpoint{3.227756in}{2.400813in}}%
\pgfpathlineto{\pgfqpoint{3.268485in}{2.236444in}}%
\pgfpathlineto{\pgfqpoint{3.299839in}{2.202962in}}%
\pgfpathlineto{\pgfqpoint{3.331288in}{2.170632in}}%
\pgfpathlineto{\pgfqpoint{3.290384in}{2.308808in}}%
\pgfpathlineto{\pgfqpoint{3.249382in}{2.465264in}}%
\pgfpathlineto{\pgfqpoint{3.218063in}{2.521091in}}%
\pgfpathlineto{\pgfqpoint{3.186830in}{2.579828in}}%
\pgfpathclose%
\pgfusepath{stroke,fill}%
\end{pgfscope}%
\begin{pgfscope}%
\pgfpathrectangle{\pgfqpoint{0.050000in}{0.050000in}}{\pgfqpoint{4.900000in}{4.900000in}}%
\pgfusepath{clip}%
\pgfsetbuttcap%
\pgfsetroundjoin%
\definecolor{currentfill}{rgb}{0.767290,0.757693,0.860854}%
\pgfsetfillcolor{currentfill}%
\pgfsetfillopacity{0.900000}%
\pgfsetlinewidth{0.200750pt}%
\definecolor{currentstroke}{rgb}{0.200000,0.200000,0.200000}%
\pgfsetstrokecolor{currentstroke}%
\pgfsetdash{}{0pt}%
\pgfpathmoveto{\pgfqpoint{2.415589in}{2.208129in}}%
\pgfpathlineto{\pgfqpoint{2.459137in}{2.314151in}}%
\pgfpathlineto{\pgfqpoint{2.502705in}{2.448505in}}%
\pgfpathlineto{\pgfqpoint{2.533236in}{2.443737in}}%
\pgfpathlineto{\pgfqpoint{2.563874in}{2.438088in}}%
\pgfpathlineto{\pgfqpoint{2.520087in}{2.304603in}}%
\pgfpathlineto{\pgfqpoint{2.476354in}{2.196384in}}%
\pgfpathlineto{\pgfqpoint{2.445918in}{2.202546in}}%
\pgfpathlineto{\pgfqpoint{2.415589in}{2.208129in}}%
\pgfpathclose%
\pgfusepath{stroke,fill}%
\end{pgfscope}%
\begin{pgfscope}%
\pgfpathrectangle{\pgfqpoint{0.050000in}{0.050000in}}{\pgfqpoint{4.900000in}{4.900000in}}%
\pgfusepath{clip}%
\pgfsetbuttcap%
\pgfsetroundjoin%
\definecolor{currentfill}{rgb}{0.266067,0.235802,0.561153}%
\pgfsetfillcolor{currentfill}%
\pgfsetfillopacity{0.900000}%
\pgfsetlinewidth{0.200750pt}%
\definecolor{currentstroke}{rgb}{0.200000,0.200000,0.200000}%
\pgfsetstrokecolor{currentstroke}%
\pgfsetdash{}{0pt}%
\pgfpathmoveto{\pgfqpoint{2.468035in}{2.823629in}}%
\pgfpathlineto{\pgfqpoint{2.511449in}{3.033812in}}%
\pgfpathlineto{\pgfqpoint{2.555072in}{3.251529in}}%
\pgfpathlineto{\pgfqpoint{2.585744in}{3.236331in}}%
\pgfpathlineto{\pgfqpoint{2.616509in}{3.219158in}}%
\pgfpathlineto{\pgfqpoint{2.572599in}{3.010150in}}%
\pgfpathlineto{\pgfqpoint{2.528893in}{2.807408in}}%
\pgfpathlineto{\pgfqpoint{2.498414in}{2.815523in}}%
\pgfpathlineto{\pgfqpoint{2.468035in}{2.823629in}}%
\pgfpathclose%
\pgfusepath{stroke,fill}%
\end{pgfscope}%
\begin{pgfscope}%
\pgfpathrectangle{\pgfqpoint{0.050000in}{0.050000in}}{\pgfqpoint{4.900000in}{4.900000in}}%
\pgfusepath{clip}%
\pgfsetbuttcap%
\pgfsetroundjoin%
\definecolor{currentfill}{rgb}{1.000000,1.000000,1.000000}%
\pgfsetfillcolor{currentfill}%
\pgfsetfillopacity{0.900000}%
\pgfsetlinewidth{0.200750pt}%
\definecolor{currentstroke}{rgb}{0.200000,0.200000,0.200000}%
\pgfsetstrokecolor{currentstroke}%
\pgfsetdash{}{0pt}%
\pgfpathmoveto{\pgfqpoint{2.004123in}{1.986586in}}%
\pgfpathlineto{\pgfqpoint{2.048393in}{2.000168in}}%
\pgfpathlineto{\pgfqpoint{2.092569in}{2.013901in}}%
\pgfpathlineto{\pgfqpoint{2.122459in}{1.994907in}}%
\pgfpathlineto{\pgfqpoint{2.152438in}{1.976011in}}%
\pgfpathlineto{\pgfqpoint{2.108195in}{1.961670in}}%
\pgfpathlineto{\pgfqpoint{2.063856in}{1.947834in}}%
\pgfpathlineto{\pgfqpoint{2.033945in}{1.967229in}}%
\pgfpathlineto{\pgfqpoint{2.004123in}{1.986586in}}%
\pgfpathclose%
\pgfusepath{stroke,fill}%
\end{pgfscope}%
\begin{pgfscope}%
\pgfpathrectangle{\pgfqpoint{0.050000in}{0.050000in}}{\pgfqpoint{4.900000in}{4.900000in}}%
\pgfusepath{clip}%
\pgfsetbuttcap%
\pgfsetroundjoin%
\definecolor{currentfill}{rgb}{1.000000,1.000000,1.000000}%
\pgfsetfillcolor{currentfill}%
\pgfsetfillopacity{0.900000}%
\pgfsetlinewidth{0.200750pt}%
\definecolor{currentstroke}{rgb}{0.200000,0.200000,0.200000}%
\pgfsetstrokecolor{currentstroke}%
\pgfsetdash{}{0pt}%
\pgfpathmoveto{\pgfqpoint{1.618993in}{1.982682in}}%
\pgfpathlineto{\pgfqpoint{1.663680in}{1.996223in}}%
\pgfpathlineto{\pgfqpoint{1.708274in}{2.009736in}}%
\pgfpathlineto{\pgfqpoint{1.737600in}{1.990429in}}%
\pgfpathlineto{\pgfqpoint{1.767015in}{1.971064in}}%
\pgfpathlineto{\pgfqpoint{1.722349in}{1.957470in}}%
\pgfpathlineto{\pgfqpoint{1.677589in}{1.943847in}}%
\pgfpathlineto{\pgfqpoint{1.648247in}{1.963294in}}%
\pgfpathlineto{\pgfqpoint{1.618993in}{1.982682in}}%
\pgfpathclose%
\pgfusepath{stroke,fill}%
\end{pgfscope}%
\begin{pgfscope}%
\pgfpathrectangle{\pgfqpoint{0.050000in}{0.050000in}}{\pgfqpoint{4.900000in}{4.900000in}}%
\pgfusepath{clip}%
\pgfsetbuttcap%
\pgfsetroundjoin%
\definecolor{currentfill}{rgb}{1.000000,1.000000,1.000000}%
\pgfsetfillcolor{currentfill}%
\pgfsetfillopacity{0.900000}%
\pgfsetlinewidth{0.200750pt}%
\definecolor{currentstroke}{rgb}{0.200000,0.200000,0.200000}%
\pgfsetstrokecolor{currentstroke}%
\pgfsetdash{}{0pt}%
\pgfpathmoveto{\pgfqpoint{2.152438in}{1.976011in}}%
\pgfpathlineto{\pgfqpoint{2.196579in}{1.991775in}}%
\pgfpathlineto{\pgfqpoint{2.240612in}{2.010890in}}%
\pgfpathlineto{\pgfqpoint{2.270742in}{1.994558in}}%
\pgfpathlineto{\pgfqpoint{2.300964in}{1.978499in}}%
\pgfpathlineto{\pgfqpoint{2.256868in}{1.956211in}}%
\pgfpathlineto{\pgfqpoint{2.212664in}{1.938573in}}%
\pgfpathlineto{\pgfqpoint{2.182506in}{1.957230in}}%
\pgfpathlineto{\pgfqpoint{2.152438in}{1.976011in}}%
\pgfpathclose%
\pgfusepath{stroke,fill}%
\end{pgfscope}%
\begin{pgfscope}%
\pgfpathrectangle{\pgfqpoint{0.050000in}{0.050000in}}{\pgfqpoint{4.900000in}{4.900000in}}%
\pgfusepath{clip}%
\pgfsetbuttcap%
\pgfsetroundjoin%
\definecolor{currentfill}{rgb}{1.000000,1.000000,1.000000}%
\pgfsetfillcolor{currentfill}%
\pgfsetfillopacity{0.900000}%
\pgfsetlinewidth{0.200750pt}%
\definecolor{currentstroke}{rgb}{0.200000,0.200000,0.200000}%
\pgfsetstrokecolor{currentstroke}%
\pgfsetdash{}{0pt}%
\pgfpathmoveto{\pgfqpoint{1.767015in}{1.971064in}}%
\pgfpathlineto{\pgfqpoint{1.811587in}{1.984630in}}%
\pgfpathlineto{\pgfqpoint{1.856066in}{1.998167in}}%
\pgfpathlineto{\pgfqpoint{1.885640in}{1.978825in}}%
\pgfpathlineto{\pgfqpoint{1.915302in}{1.959425in}}%
\pgfpathlineto{\pgfqpoint{1.870752in}{1.945807in}}%
\pgfpathlineto{\pgfqpoint{1.826108in}{1.932159in}}%
\pgfpathlineto{\pgfqpoint{1.796517in}{1.951641in}}%
\pgfpathlineto{\pgfqpoint{1.767015in}{1.971064in}}%
\pgfpathclose%
\pgfusepath{stroke,fill}%
\end{pgfscope}%
\begin{pgfscope}%
\pgfpathrectangle{\pgfqpoint{0.050000in}{0.050000in}}{\pgfqpoint{4.900000in}{4.900000in}}%
\pgfusepath{clip}%
\pgfsetbuttcap%
\pgfsetroundjoin%
\definecolor{currentfill}{rgb}{0.252380,0.347712,0.731042}%
\pgfsetfillcolor{currentfill}%
\pgfsetfillopacity{0.900000}%
\pgfsetlinewidth{0.200750pt}%
\definecolor{currentstroke}{rgb}{0.200000,0.200000,0.200000}%
\pgfsetstrokecolor{currentstroke}%
\pgfsetdash{}{0pt}%
\pgfpathmoveto{\pgfqpoint{2.494012in}{3.279828in}}%
\pgfpathlineto{\pgfqpoint{2.537563in}{3.502906in}}%
\pgfpathlineto{\pgfqpoint{2.581342in}{3.726662in}}%
\pgfpathlineto{\pgfqpoint{2.612117in}{3.708030in}}%
\pgfpathlineto{\pgfqpoint{2.642973in}{3.681861in}}%
\pgfpathlineto{\pgfqpoint{2.598916in}{3.470555in}}%
\pgfpathlineto{\pgfqpoint{2.555072in}{3.251529in}}%
\pgfpathlineto{\pgfqpoint{2.524494in}{3.265731in}}%
\pgfpathlineto{\pgfqpoint{2.494012in}{3.279828in}}%
\pgfpathclose%
\pgfusepath{stroke,fill}%
\end{pgfscope}%
\begin{pgfscope}%
\pgfpathrectangle{\pgfqpoint{0.050000in}{0.050000in}}{\pgfqpoint{4.900000in}{4.900000in}}%
\pgfusepath{clip}%
\pgfsetbuttcap%
\pgfsetroundjoin%
\definecolor{currentfill}{rgb}{0.779223,0.770119,0.867989}%
\pgfsetfillcolor{currentfill}%
\pgfsetfillopacity{0.900000}%
\pgfsetlinewidth{0.200750pt}%
\definecolor{currentstroke}{rgb}{0.200000,0.200000,0.200000}%
\pgfsetstrokecolor{currentstroke}%
\pgfsetdash{}{0pt}%
\pgfpathmoveto{\pgfqpoint{3.249382in}{2.465264in}}%
\pgfpathlineto{\pgfqpoint{3.290384in}{2.308808in}}%
\pgfpathlineto{\pgfqpoint{3.331288in}{2.170632in}}%
\pgfpathlineto{\pgfqpoint{3.362831in}{2.139304in}}%
\pgfpathlineto{\pgfqpoint{3.394468in}{2.108853in}}%
\pgfpathlineto{\pgfqpoint{3.353381in}{2.224551in}}%
\pgfpathlineto{\pgfqpoint{3.312282in}{2.361074in}}%
\pgfpathlineto{\pgfqpoint{3.280788in}{2.412018in}}%
\pgfpathlineto{\pgfqpoint{3.249382in}{2.465264in}}%
\pgfpathclose%
\pgfusepath{stroke,fill}%
\end{pgfscope}%
\begin{pgfscope}%
\pgfpathrectangle{\pgfqpoint{0.050000in}{0.050000in}}{\pgfqpoint{4.900000in}{4.900000in}}%
\pgfusepath{clip}%
\pgfsetbuttcap%
\pgfsetroundjoin%
\definecolor{currentfill}{rgb}{0.540546,0.521599,0.725275}%
\pgfsetfillcolor{currentfill}%
\pgfsetfillopacity{0.900000}%
\pgfsetlinewidth{0.200750pt}%
\definecolor{currentstroke}{rgb}{0.200000,0.200000,0.200000}%
\pgfsetstrokecolor{currentstroke}%
\pgfsetdash{}{0pt}%
\pgfpathmoveto{\pgfqpoint{2.502705in}{2.448505in}}%
\pgfpathlineto{\pgfqpoint{2.546360in}{2.608764in}}%
\pgfpathlineto{\pgfqpoint{2.590158in}{2.789455in}}%
\pgfpathlineto{\pgfqpoint{2.620942in}{2.778875in}}%
\pgfpathlineto{\pgfqpoint{2.651824in}{2.766850in}}%
\pgfpathlineto{\pgfqpoint{2.607773in}{2.594129in}}%
\pgfpathlineto{\pgfqpoint{2.563874in}{2.438088in}}%
\pgfpathlineto{\pgfqpoint{2.533236in}{2.443737in}}%
\pgfpathlineto{\pgfqpoint{2.502705in}{2.448505in}}%
\pgfpathclose%
\pgfusepath{stroke,fill}%
\end{pgfscope}%
\begin{pgfscope}%
\pgfpathrectangle{\pgfqpoint{0.050000in}{0.050000in}}{\pgfqpoint{4.900000in}{4.900000in}}%
\pgfusepath{clip}%
\pgfsetbuttcap%
\pgfsetroundjoin%
\definecolor{currentfill}{rgb}{0.892595,0.888166,0.935779}%
\pgfsetfillcolor{currentfill}%
\pgfsetfillopacity{0.900000}%
\pgfsetlinewidth{0.200750pt}%
\definecolor{currentstroke}{rgb}{0.200000,0.200000,0.200000}%
\pgfsetstrokecolor{currentstroke}%
\pgfsetdash{}{0pt}%
\pgfpathmoveto{\pgfqpoint{2.388826in}{2.052178in}}%
\pgfpathlineto{\pgfqpoint{2.432617in}{2.113169in}}%
\pgfpathlineto{\pgfqpoint{2.476354in}{2.196384in}}%
\pgfpathlineto{\pgfqpoint{2.506895in}{2.189675in}}%
\pgfpathlineto{\pgfqpoint{2.537541in}{2.182434in}}%
\pgfpathlineto{\pgfqpoint{2.493646in}{2.095694in}}%
\pgfpathlineto{\pgfqpoint{2.449730in}{2.029975in}}%
\pgfpathlineto{\pgfqpoint{2.419228in}{2.041107in}}%
\pgfpathlineto{\pgfqpoint{2.388826in}{2.052178in}}%
\pgfpathclose%
\pgfusepath{stroke,fill}%
\end{pgfscope}%
\begin{pgfscope}%
\pgfpathrectangle{\pgfqpoint{0.050000in}{0.050000in}}{\pgfqpoint{4.900000in}{4.900000in}}%
\pgfusepath{clip}%
\pgfsetbuttcap%
\pgfsetroundjoin%
\definecolor{currentfill}{rgb}{0.970165,0.968935,0.982161}%
\pgfsetfillcolor{currentfill}%
\pgfsetfillopacity{0.900000}%
\pgfsetlinewidth{0.200750pt}%
\definecolor{currentstroke}{rgb}{0.200000,0.200000,0.200000}%
\pgfsetstrokecolor{currentstroke}%
\pgfsetdash{}{0pt}%
\pgfpathmoveto{\pgfqpoint{3.458029in}{2.050184in}}%
\pgfpathlineto{\pgfqpoint{3.499681in}{1.991882in}}%
\pgfpathlineto{\pgfqpoint{3.531732in}{1.971960in}}%
\pgfpathlineto{\pgfqpoint{3.563879in}{1.952018in}}%
\pgfpathlineto{\pgfqpoint{3.521972in}{1.993961in}}%
\pgfpathlineto{\pgfqpoint{3.489952in}{2.021801in}}%
\pgfpathlineto{\pgfqpoint{3.458029in}{2.050184in}}%
\pgfpathclose%
\pgfusepath{stroke,fill}%
\end{pgfscope}%
\begin{pgfscope}%
\pgfpathrectangle{\pgfqpoint{0.050000in}{0.050000in}}{\pgfqpoint{4.900000in}{4.900000in}}%
\pgfusepath{clip}%
\pgfsetbuttcap%
\pgfsetroundjoin%
\definecolor{currentfill}{rgb}{1.000000,1.000000,1.000000}%
\pgfsetfillcolor{currentfill}%
\pgfsetfillopacity{0.900000}%
\pgfsetlinewidth{0.200750pt}%
\definecolor{currentstroke}{rgb}{0.200000,0.200000,0.200000}%
\pgfsetstrokecolor{currentstroke}%
\pgfsetdash{}{0pt}%
\pgfpathmoveto{\pgfqpoint{1.915302in}{1.959425in}}%
\pgfpathlineto{\pgfqpoint{1.959759in}{1.973016in}}%
\pgfpathlineto{\pgfqpoint{2.004123in}{1.986586in}}%
\pgfpathlineto{\pgfqpoint{2.033945in}{1.967229in}}%
\pgfpathlineto{\pgfqpoint{2.063856in}{1.947834in}}%
\pgfpathlineto{\pgfqpoint{2.019423in}{1.934136in}}%
\pgfpathlineto{\pgfqpoint{1.974895in}{1.920453in}}%
\pgfpathlineto{\pgfqpoint{1.945054in}{1.939968in}}%
\pgfpathlineto{\pgfqpoint{1.915302in}{1.959425in}}%
\pgfpathclose%
\pgfusepath{stroke,fill}%
\end{pgfscope}%
\begin{pgfscope}%
\pgfpathrectangle{\pgfqpoint{0.050000in}{0.050000in}}{\pgfqpoint{4.900000in}{4.900000in}}%
\pgfusepath{clip}%
\pgfsetbuttcap%
\pgfsetroundjoin%
\definecolor{currentfill}{rgb}{1.000000,1.000000,1.000000}%
\pgfsetfillcolor{currentfill}%
\pgfsetfillopacity{0.900000}%
\pgfsetlinewidth{0.200750pt}%
\definecolor{currentstroke}{rgb}{0.200000,0.200000,0.200000}%
\pgfsetstrokecolor{currentstroke}%
\pgfsetdash{}{0pt}%
\pgfpathmoveto{\pgfqpoint{1.529337in}{1.955514in}}%
\pgfpathlineto{\pgfqpoint{1.574212in}{1.969112in}}%
\pgfpathlineto{\pgfqpoint{1.618993in}{1.982682in}}%
\pgfpathlineto{\pgfqpoint{1.648247in}{1.963294in}}%
\pgfpathlineto{\pgfqpoint{1.677589in}{1.943847in}}%
\pgfpathlineto{\pgfqpoint{1.632735in}{1.930196in}}%
\pgfpathlineto{\pgfqpoint{1.587787in}{1.916516in}}%
\pgfpathlineto{\pgfqpoint{1.558518in}{1.936044in}}%
\pgfpathlineto{\pgfqpoint{1.529337in}{1.955514in}}%
\pgfpathclose%
\pgfusepath{stroke,fill}%
\end{pgfscope}%
\begin{pgfscope}%
\pgfpathrectangle{\pgfqpoint{0.050000in}{0.050000in}}{\pgfqpoint{4.900000in}{4.900000in}}%
\pgfusepath{clip}%
\pgfsetbuttcap%
\pgfsetroundjoin%
\definecolor{currentfill}{rgb}{0.964198,0.962722,0.978593}%
\pgfsetfillcolor{currentfill}%
\pgfsetfillopacity{0.900000}%
\pgfsetlinewidth{0.200750pt}%
\definecolor{currentstroke}{rgb}{0.200000,0.200000,0.200000}%
\pgfsetstrokecolor{currentstroke}%
\pgfsetdash{}{0pt}%
\pgfpathmoveto{\pgfqpoint{2.300964in}{1.978499in}}%
\pgfpathlineto{\pgfqpoint{2.344948in}{2.008960in}}%
\pgfpathlineto{\pgfqpoint{2.388826in}{2.052178in}}%
\pgfpathlineto{\pgfqpoint{2.419228in}{2.041107in}}%
\pgfpathlineto{\pgfqpoint{2.449730in}{2.029975in}}%
\pgfpathlineto{\pgfqpoint{2.405752in}{1.981822in}}%
\pgfpathlineto{\pgfqpoint{2.361687in}{1.947027in}}%
\pgfpathlineto{\pgfqpoint{2.331278in}{1.962668in}}%
\pgfpathlineto{\pgfqpoint{2.300964in}{1.978499in}}%
\pgfpathclose%
\pgfusepath{stroke,fill}%
\end{pgfscope}%
\begin{pgfscope}%
\pgfpathrectangle{\pgfqpoint{0.050000in}{0.050000in}}{\pgfqpoint{4.900000in}{4.900000in}}%
\pgfusepath{clip}%
\pgfsetbuttcap%
\pgfsetroundjoin%
\definecolor{currentfill}{rgb}{0.260100,0.229589,0.557586}%
\pgfsetfillcolor{currentfill}%
\pgfsetfillopacity{0.900000}%
\pgfsetlinewidth{0.200750pt}%
\definecolor{currentstroke}{rgb}{0.200000,0.200000,0.200000}%
\pgfsetstrokecolor{currentstroke}%
\pgfsetdash{}{0pt}%
\pgfpathmoveto{\pgfqpoint{2.528893in}{2.807408in}}%
\pgfpathlineto{\pgfqpoint{2.572599in}{3.010150in}}%
\pgfpathlineto{\pgfqpoint{2.616509in}{3.219158in}}%
\pgfpathlineto{\pgfqpoint{2.647363in}{3.199288in}}%
\pgfpathlineto{\pgfqpoint{2.678302in}{3.176347in}}%
\pgfpathlineto{\pgfqpoint{2.634137in}{2.982410in}}%
\pgfpathlineto{\pgfqpoint{2.590158in}{2.789455in}}%
\pgfpathlineto{\pgfqpoint{2.559475in}{2.798853in}}%
\pgfpathlineto{\pgfqpoint{2.528893in}{2.807408in}}%
\pgfpathclose%
\pgfusepath{stroke,fill}%
\end{pgfscope}%
\begin{pgfscope}%
\pgfpathrectangle{\pgfqpoint{0.050000in}{0.050000in}}{\pgfqpoint{4.900000in}{4.900000in}}%
\pgfusepath{clip}%
\pgfsetbuttcap%
\pgfsetroundjoin%
\definecolor{currentfill}{rgb}{0.743422,0.732841,0.846582}%
\pgfsetfillcolor{currentfill}%
\pgfsetfillopacity{0.900000}%
\pgfsetlinewidth{0.200750pt}%
\definecolor{currentstroke}{rgb}{0.200000,0.200000,0.200000}%
\pgfsetstrokecolor{currentstroke}%
\pgfsetdash{}{0pt}%
\pgfpathmoveto{\pgfqpoint{2.476354in}{2.196384in}}%
\pgfpathlineto{\pgfqpoint{2.520087in}{2.304603in}}%
\pgfpathlineto{\pgfqpoint{2.563874in}{2.438088in}}%
\pgfpathlineto{\pgfqpoint{2.594618in}{2.431514in}}%
\pgfpathlineto{\pgfqpoint{2.625467in}{2.423963in}}%
\pgfpathlineto{\pgfqpoint{2.581464in}{2.292029in}}%
\pgfpathlineto{\pgfqpoint{2.537541in}{2.182434in}}%
\pgfpathlineto{\pgfqpoint{2.506895in}{2.189675in}}%
\pgfpathlineto{\pgfqpoint{2.476354in}{2.196384in}}%
\pgfpathclose%
\pgfusepath{stroke,fill}%
\end{pgfscope}%
\begin{pgfscope}%
\pgfpathrectangle{\pgfqpoint{0.050000in}{0.050000in}}{\pgfqpoint{4.900000in}{4.900000in}}%
\pgfusepath{clip}%
\pgfsetbuttcap%
\pgfsetroundjoin%
\definecolor{currentfill}{rgb}{1.000000,1.000000,1.000000}%
\pgfsetfillcolor{currentfill}%
\pgfsetfillopacity{0.900000}%
\pgfsetlinewidth{0.200750pt}%
\definecolor{currentstroke}{rgb}{0.200000,0.200000,0.200000}%
\pgfsetstrokecolor{currentstroke}%
\pgfsetdash{}{0pt}%
\pgfpathmoveto{\pgfqpoint{2.063856in}{1.947834in}}%
\pgfpathlineto{\pgfqpoint{2.108195in}{1.961670in}}%
\pgfpathlineto{\pgfqpoint{2.152438in}{1.976011in}}%
\pgfpathlineto{\pgfqpoint{2.182506in}{1.957230in}}%
\pgfpathlineto{\pgfqpoint{2.212664in}{1.938573in}}%
\pgfpathlineto{\pgfqpoint{2.168356in}{1.923266in}}%
\pgfpathlineto{\pgfqpoint{2.123949in}{1.908976in}}%
\pgfpathlineto{\pgfqpoint{2.093858in}{1.928413in}}%
\pgfpathlineto{\pgfqpoint{2.063856in}{1.947834in}}%
\pgfpathclose%
\pgfusepath{stroke,fill}%
\end{pgfscope}%
\begin{pgfscope}%
\pgfpathrectangle{\pgfqpoint{0.050000in}{0.050000in}}{\pgfqpoint{4.900000in}{4.900000in}}%
\pgfusepath{clip}%
\pgfsetbuttcap%
\pgfsetroundjoin%
\definecolor{currentfill}{rgb}{1.000000,1.000000,1.000000}%
\pgfsetfillcolor{currentfill}%
\pgfsetfillopacity{0.900000}%
\pgfsetlinewidth{0.200750pt}%
\definecolor{currentstroke}{rgb}{0.200000,0.200000,0.200000}%
\pgfsetstrokecolor{currentstroke}%
\pgfsetdash{}{0pt}%
\pgfpathmoveto{\pgfqpoint{1.677589in}{1.943847in}}%
\pgfpathlineto{\pgfqpoint{1.722349in}{1.957470in}}%
\pgfpathlineto{\pgfqpoint{1.767015in}{1.971064in}}%
\pgfpathlineto{\pgfqpoint{1.796517in}{1.951641in}}%
\pgfpathlineto{\pgfqpoint{1.826108in}{1.932159in}}%
\pgfpathlineto{\pgfqpoint{1.781370in}{1.918483in}}%
\pgfpathlineto{\pgfqpoint{1.736538in}{1.904778in}}%
\pgfpathlineto{\pgfqpoint{1.707019in}{1.924342in}}%
\pgfpathlineto{\pgfqpoint{1.677589in}{1.943847in}}%
\pgfpathclose%
\pgfusepath{stroke,fill}%
\end{pgfscope}%
\begin{pgfscope}%
\pgfpathrectangle{\pgfqpoint{0.050000in}{0.050000in}}{\pgfqpoint{4.900000in}{4.900000in}}%
\pgfusepath{clip}%
\pgfsetbuttcap%
\pgfsetroundjoin%
\definecolor{currentfill}{rgb}{0.815025,0.807397,0.889396}%
\pgfsetfillcolor{currentfill}%
\pgfsetfillopacity{0.900000}%
\pgfsetlinewidth{0.200750pt}%
\definecolor{currentstroke}{rgb}{0.200000,0.200000,0.200000}%
\pgfsetstrokecolor{currentstroke}%
\pgfsetdash{}{0pt}%
\pgfpathmoveto{\pgfqpoint{3.312282in}{2.361074in}}%
\pgfpathlineto{\pgfqpoint{3.353381in}{2.224551in}}%
\pgfpathlineto{\pgfqpoint{3.394468in}{2.108853in}}%
\pgfpathlineto{\pgfqpoint{3.426201in}{2.079176in}}%
\pgfpathlineto{\pgfqpoint{3.458029in}{2.050184in}}%
\pgfpathlineto{\pgfqpoint{3.416749in}{2.146366in}}%
\pgfpathlineto{\pgfqpoint{3.375537in}{2.265171in}}%
\pgfpathlineto{\pgfqpoint{3.343864in}{2.312194in}}%
\pgfpathlineto{\pgfqpoint{3.312282in}{2.361074in}}%
\pgfpathclose%
\pgfusepath{stroke,fill}%
\end{pgfscope}%
\begin{pgfscope}%
\pgfpathrectangle{\pgfqpoint{0.050000in}{0.050000in}}{\pgfqpoint{4.900000in}{4.900000in}}%
\pgfusepath{clip}%
\pgfsetbuttcap%
\pgfsetroundjoin%
\definecolor{currentfill}{rgb}{0.247951,0.300654,0.668481}%
\pgfsetfillcolor{currentfill}%
\pgfsetfillopacity{0.900000}%
\pgfsetlinewidth{0.200750pt}%
\definecolor{currentstroke}{rgb}{0.200000,0.200000,0.200000}%
\pgfsetstrokecolor{currentstroke}%
\pgfsetdash{}{0pt}%
\pgfpathmoveto{\pgfqpoint{2.854884in}{3.567495in}}%
\pgfpathlineto{\pgfqpoint{2.897879in}{3.448839in}}%
\pgfpathlineto{\pgfqpoint{2.939890in}{3.223308in}}%
\pgfpathlineto{\pgfqpoint{2.970490in}{3.117546in}}%
\pgfpathlineto{\pgfqpoint{3.001161in}{3.021593in}}%
\pgfpathlineto{\pgfqpoint{2.959404in}{3.237200in}}%
\pgfpathlineto{\pgfqpoint{2.916723in}{3.372617in}}%
\pgfpathlineto{\pgfqpoint{2.885788in}{3.468014in}}%
\pgfpathlineto{\pgfqpoint{2.854884in}{3.567495in}}%
\pgfpathclose%
\pgfusepath{stroke,fill}%
\end{pgfscope}%
\begin{pgfscope}%
\pgfpathrectangle{\pgfqpoint{0.050000in}{0.050000in}}{\pgfqpoint{4.900000in}{4.900000in}}%
\pgfusepath{clip}%
\pgfsetbuttcap%
\pgfsetroundjoin%
\definecolor{currentfill}{rgb}{0.994033,0.993787,0.996432}%
\pgfsetfillcolor{currentfill}%
\pgfsetfillopacity{0.900000}%
\pgfsetlinewidth{0.200750pt}%
\definecolor{currentstroke}{rgb}{0.200000,0.200000,0.200000}%
\pgfsetstrokecolor{currentstroke}%
\pgfsetdash{}{0pt}%
\pgfpathmoveto{\pgfqpoint{2.212664in}{1.938573in}}%
\pgfpathlineto{\pgfqpoint{2.256868in}{1.956211in}}%
\pgfpathlineto{\pgfqpoint{2.300964in}{1.978499in}}%
\pgfpathlineto{\pgfqpoint{2.331278in}{1.962668in}}%
\pgfpathlineto{\pgfqpoint{2.361687in}{1.947027in}}%
\pgfpathlineto{\pgfqpoint{2.317522in}{1.921456in}}%
\pgfpathlineto{\pgfqpoint{2.273252in}{1.901641in}}%
\pgfpathlineto{\pgfqpoint{2.242912in}{1.920044in}}%
\pgfpathlineto{\pgfqpoint{2.212664in}{1.938573in}}%
\pgfpathclose%
\pgfusepath{stroke,fill}%
\end{pgfscope}%
\begin{pgfscope}%
\pgfpathrectangle{\pgfqpoint{0.050000in}{0.050000in}}{\pgfqpoint{4.900000in}{4.900000in}}%
\pgfusepath{clip}%
\pgfsetbuttcap%
\pgfsetroundjoin%
\definecolor{currentfill}{rgb}{0.254348,0.368627,0.758847}%
\pgfsetfillcolor{currentfill}%
\pgfsetfillopacity{0.900000}%
\pgfsetlinewidth{0.200750pt}%
\definecolor{currentstroke}{rgb}{0.200000,0.200000,0.200000}%
\pgfsetstrokecolor{currentstroke}%
\pgfsetdash{}{0pt}%
\pgfpathmoveto{\pgfqpoint{2.793130in}{3.776190in}}%
\pgfpathlineto{\pgfqpoint{2.836509in}{3.695458in}}%
\pgfpathlineto{\pgfqpoint{2.878892in}{3.473335in}}%
\pgfpathlineto{\pgfqpoint{2.909359in}{3.340994in}}%
\pgfpathlineto{\pgfqpoint{2.939890in}{3.223308in}}%
\pgfpathlineto{\pgfqpoint{2.897879in}{3.448839in}}%
\pgfpathlineto{\pgfqpoint{2.854884in}{3.567495in}}%
\pgfpathlineto{\pgfqpoint{2.824001in}{3.670584in}}%
\pgfpathlineto{\pgfqpoint{2.793130in}{3.776190in}}%
\pgfpathclose%
\pgfusepath{stroke,fill}%
\end{pgfscope}%
\begin{pgfscope}%
\pgfpathrectangle{\pgfqpoint{0.050000in}{0.050000in}}{\pgfqpoint{4.900000in}{4.900000in}}%
\pgfusepath{clip}%
\pgfsetbuttcap%
\pgfsetroundjoin%
\definecolor{currentfill}{rgb}{0.251888,0.342484,0.724091}%
\pgfsetfillcolor{currentfill}%
\pgfsetfillopacity{0.900000}%
\pgfsetlinewidth{0.200750pt}%
\definecolor{currentstroke}{rgb}{0.200000,0.200000,0.200000}%
\pgfsetstrokecolor{currentstroke}%
\pgfsetdash{}{0pt}%
\pgfpathmoveto{\pgfqpoint{2.555072in}{3.251529in}}%
\pgfpathlineto{\pgfqpoint{2.598916in}{3.470555in}}%
\pgfpathlineto{\pgfqpoint{2.642973in}{3.681861in}}%
\pgfpathlineto{\pgfqpoint{2.673895in}{3.646838in}}%
\pgfpathlineto{\pgfqpoint{2.704872in}{3.603562in}}%
\pgfpathlineto{\pgfqpoint{2.660618in}{3.422898in}}%
\pgfpathlineto{\pgfqpoint{2.616509in}{3.219158in}}%
\pgfpathlineto{\pgfqpoint{2.585744in}{3.236331in}}%
\pgfpathlineto{\pgfqpoint{2.555072in}{3.251529in}}%
\pgfpathclose%
\pgfusepath{stroke,fill}%
\end{pgfscope}%
\begin{pgfscope}%
\pgfpathrectangle{\pgfqpoint{0.050000in}{0.050000in}}{\pgfqpoint{4.900000in}{4.900000in}}%
\pgfusepath{clip}%
\pgfsetbuttcap%
\pgfsetroundjoin%
\definecolor{currentfill}{rgb}{1.000000,1.000000,1.000000}%
\pgfsetfillcolor{currentfill}%
\pgfsetfillopacity{0.900000}%
\pgfsetlinewidth{0.200750pt}%
\definecolor{currentstroke}{rgb}{0.200000,0.200000,0.200000}%
\pgfsetstrokecolor{currentstroke}%
\pgfsetdash{}{0pt}%
\pgfpathmoveto{\pgfqpoint{1.826108in}{1.932159in}}%
\pgfpathlineto{\pgfqpoint{1.870752in}{1.945807in}}%
\pgfpathlineto{\pgfqpoint{1.915302in}{1.959425in}}%
\pgfpathlineto{\pgfqpoint{1.945054in}{1.939968in}}%
\pgfpathlineto{\pgfqpoint{1.974895in}{1.920453in}}%
\pgfpathlineto{\pgfqpoint{1.930274in}{1.906750in}}%
\pgfpathlineto{\pgfqpoint{1.885558in}{1.893020in}}%
\pgfpathlineto{\pgfqpoint{1.855788in}{1.912619in}}%
\pgfpathlineto{\pgfqpoint{1.826108in}{1.932159in}}%
\pgfpathclose%
\pgfusepath{stroke,fill}%
\end{pgfscope}%
\begin{pgfscope}%
\pgfpathrectangle{\pgfqpoint{0.050000in}{0.050000in}}{\pgfqpoint{4.900000in}{4.900000in}}%
\pgfusepath{clip}%
\pgfsetbuttcap%
\pgfsetroundjoin%
\definecolor{currentfill}{rgb}{0.242537,0.243137,0.592018}%
\pgfsetfillcolor{currentfill}%
\pgfsetfillopacity{0.900000}%
\pgfsetlinewidth{0.200750pt}%
\definecolor{currentstroke}{rgb}{0.200000,0.200000,0.200000}%
\pgfsetstrokecolor{currentstroke}%
\pgfsetdash{}{0pt}%
\pgfpathmoveto{\pgfqpoint{2.916723in}{3.372617in}}%
\pgfpathlineto{\pgfqpoint{2.959404in}{3.237200in}}%
\pgfpathlineto{\pgfqpoint{3.001161in}{3.021593in}}%
\pgfpathlineto{\pgfqpoint{3.031908in}{2.933800in}}%
\pgfpathlineto{\pgfqpoint{3.062731in}{2.852868in}}%
\pgfpathlineto{\pgfqpoint{3.021144in}{3.053650in}}%
\pgfpathlineto{\pgfqpoint{2.978711in}{3.194376in}}%
\pgfpathlineto{\pgfqpoint{2.947695in}{3.281427in}}%
\pgfpathlineto{\pgfqpoint{2.916723in}{3.372617in}}%
\pgfpathclose%
\pgfusepath{stroke,fill}%
\end{pgfscope}%
\begin{pgfscope}%
\pgfpathrectangle{\pgfqpoint{0.050000in}{0.050000in}}{\pgfqpoint{4.900000in}{4.900000in}}%
\pgfusepath{clip}%
\pgfsetbuttcap%
\pgfsetroundjoin%
\definecolor{currentfill}{rgb}{1.000000,1.000000,1.000000}%
\pgfsetfillcolor{currentfill}%
\pgfsetfillopacity{0.900000}%
\pgfsetlinewidth{0.200750pt}%
\definecolor{currentstroke}{rgb}{0.200000,0.200000,0.200000}%
\pgfsetstrokecolor{currentstroke}%
\pgfsetdash{}{0pt}%
\pgfpathmoveto{\pgfqpoint{1.439303in}{1.928232in}}%
\pgfpathlineto{\pgfqpoint{1.484367in}{1.941887in}}%
\pgfpathlineto{\pgfqpoint{1.529337in}{1.955514in}}%
\pgfpathlineto{\pgfqpoint{1.558518in}{1.936044in}}%
\pgfpathlineto{\pgfqpoint{1.587787in}{1.916516in}}%
\pgfpathlineto{\pgfqpoint{1.542743in}{1.902807in}}%
\pgfpathlineto{\pgfqpoint{1.497605in}{1.889069in}}%
\pgfpathlineto{\pgfqpoint{1.468410in}{1.908680in}}%
\pgfpathlineto{\pgfqpoint{1.439303in}{1.928232in}}%
\pgfpathclose%
\pgfusepath{stroke,fill}%
\end{pgfscope}%
\begin{pgfscope}%
\pgfpathrectangle{\pgfqpoint{0.050000in}{0.050000in}}{\pgfqpoint{4.900000in}{4.900000in}}%
\pgfusepath{clip}%
\pgfsetbuttcap%
\pgfsetroundjoin%
\definecolor{currentfill}{rgb}{0.982099,0.981361,0.989296}%
\pgfsetfillcolor{currentfill}%
\pgfsetfillopacity{0.900000}%
\pgfsetlinewidth{0.200750pt}%
\definecolor{currentstroke}{rgb}{0.200000,0.200000,0.200000}%
\pgfsetstrokecolor{currentstroke}%
\pgfsetdash{}{0pt}%
\pgfpathmoveto{\pgfqpoint{3.521972in}{1.993961in}}%
\pgfpathlineto{\pgfqpoint{3.563879in}{1.952018in}}%
\pgfpathlineto{\pgfqpoint{3.596123in}{1.932052in}}%
\pgfpathlineto{\pgfqpoint{3.628463in}{1.912058in}}%
\pgfpathlineto{\pgfqpoint{3.586303in}{1.939694in}}%
\pgfpathlineto{\pgfqpoint{3.554089in}{1.966609in}}%
\pgfpathlineto{\pgfqpoint{3.521972in}{1.993961in}}%
\pgfpathclose%
\pgfusepath{stroke,fill}%
\end{pgfscope}%
\begin{pgfscope}%
\pgfpathrectangle{\pgfqpoint{0.050000in}{0.050000in}}{\pgfqpoint{4.900000in}{4.900000in}}%
\pgfusepath{clip}%
\pgfsetbuttcap%
\pgfsetroundjoin%
\definecolor{currentfill}{rgb}{0.261976,0.449673,0.866590}%
\pgfsetfillcolor{currentfill}%
\pgfsetfillopacity{0.900000}%
\pgfsetlinewidth{0.200750pt}%
\definecolor{currentstroke}{rgb}{0.200000,0.200000,0.200000}%
\pgfsetstrokecolor{currentstroke}%
\pgfsetdash{}{0pt}%
\pgfpathmoveto{\pgfqpoint{2.731374in}{3.984531in}}%
\pgfpathlineto{\pgfqpoint{2.775194in}{3.980281in}}%
\pgfpathlineto{\pgfqpoint{2.818119in}{3.796863in}}%
\pgfpathlineto{\pgfqpoint{2.848482in}{3.623838in}}%
\pgfpathlineto{\pgfqpoint{2.878892in}{3.473335in}}%
\pgfpathlineto{\pgfqpoint{2.836509in}{3.695458in}}%
\pgfpathlineto{\pgfqpoint{2.793130in}{3.776190in}}%
\pgfpathlineto{\pgfqpoint{2.762258in}{3.882156in}}%
\pgfpathlineto{\pgfqpoint{2.731374in}{3.984531in}}%
\pgfpathclose%
\pgfusepath{stroke,fill}%
\end{pgfscope}%
\begin{pgfscope}%
\pgfpathrectangle{\pgfqpoint{0.050000in}{0.050000in}}{\pgfqpoint{4.900000in}{4.900000in}}%
\pgfusepath{clip}%
\pgfsetbuttcap%
\pgfsetroundjoin%
\definecolor{currentfill}{rgb}{0.266067,0.235802,0.561153}%
\pgfsetfillcolor{currentfill}%
\pgfsetfillopacity{0.900000}%
\pgfsetlinewidth{0.200750pt}%
\definecolor{currentstroke}{rgb}{0.200000,0.200000,0.200000}%
\pgfsetstrokecolor{currentstroke}%
\pgfsetdash{}{0pt}%
\pgfpathmoveto{\pgfqpoint{2.978711in}{3.194376in}}%
\pgfpathlineto{\pgfqpoint{3.021144in}{3.053650in}}%
\pgfpathlineto{\pgfqpoint{3.062731in}{2.852868in}}%
\pgfpathlineto{\pgfqpoint{3.093634in}{2.777764in}}%
\pgfpathlineto{\pgfqpoint{3.124617in}{2.707654in}}%
\pgfpathlineto{\pgfqpoint{3.083138in}{2.892181in}}%
\pgfpathlineto{\pgfqpoint{3.040898in}{3.031945in}}%
\pgfpathlineto{\pgfqpoint{3.009777in}{3.111289in}}%
\pgfpathlineto{\pgfqpoint{2.978711in}{3.194376in}}%
\pgfpathclose%
\pgfusepath{stroke,fill}%
\end{pgfscope}%
\begin{pgfscope}%
\pgfpathrectangle{\pgfqpoint{0.050000in}{0.050000in}}{\pgfqpoint{4.900000in}{4.900000in}}%
\pgfusepath{clip}%
\pgfsetbuttcap%
\pgfsetroundjoin%
\definecolor{currentfill}{rgb}{1.000000,1.000000,1.000000}%
\pgfsetfillcolor{currentfill}%
\pgfsetfillopacity{0.900000}%
\pgfsetlinewidth{0.200750pt}%
\definecolor{currentstroke}{rgb}{0.200000,0.200000,0.200000}%
\pgfsetstrokecolor{currentstroke}%
\pgfsetdash{}{0pt}%
\pgfpathmoveto{\pgfqpoint{1.974895in}{1.920453in}}%
\pgfpathlineto{\pgfqpoint{2.019423in}{1.934136in}}%
\pgfpathlineto{\pgfqpoint{2.063856in}{1.947834in}}%
\pgfpathlineto{\pgfqpoint{2.093858in}{1.928413in}}%
\pgfpathlineto{\pgfqpoint{2.123949in}{1.908976in}}%
\pgfpathlineto{\pgfqpoint{2.079446in}{1.895062in}}%
\pgfpathlineto{\pgfqpoint{2.034848in}{1.881256in}}%
\pgfpathlineto{\pgfqpoint{2.004826in}{1.900881in}}%
\pgfpathlineto{\pgfqpoint{1.974895in}{1.920453in}}%
\pgfpathclose%
\pgfusepath{stroke,fill}%
\end{pgfscope}%
\begin{pgfscope}%
\pgfpathrectangle{\pgfqpoint{0.050000in}{0.050000in}}{\pgfqpoint{4.900000in}{4.900000in}}%
\pgfusepath{clip}%
\pgfsetbuttcap%
\pgfsetroundjoin%
\definecolor{currentfill}{rgb}{0.522645,0.502960,0.714571}%
\pgfsetfillcolor{currentfill}%
\pgfsetfillopacity{0.900000}%
\pgfsetlinewidth{0.200750pt}%
\definecolor{currentstroke}{rgb}{0.200000,0.200000,0.200000}%
\pgfsetstrokecolor{currentstroke}%
\pgfsetdash{}{0pt}%
\pgfpathmoveto{\pgfqpoint{2.563874in}{2.438088in}}%
\pgfpathlineto{\pgfqpoint{2.607773in}{2.594129in}}%
\pgfpathlineto{\pgfqpoint{2.651824in}{2.766850in}}%
\pgfpathlineto{\pgfqpoint{2.682804in}{2.753205in}}%
\pgfpathlineto{\pgfqpoint{2.713877in}{2.737843in}}%
\pgfpathlineto{\pgfqpoint{2.669596in}{2.574815in}}%
\pgfpathlineto{\pgfqpoint{2.625467in}{2.423963in}}%
\pgfpathlineto{\pgfqpoint{2.594618in}{2.431514in}}%
\pgfpathlineto{\pgfqpoint{2.563874in}{2.438088in}}%
\pgfpathclose%
\pgfusepath{stroke,fill}%
\end{pgfscope}%
\begin{pgfscope}%
\pgfpathrectangle{\pgfqpoint{0.050000in}{0.050000in}}{\pgfqpoint{4.900000in}{4.900000in}}%
\pgfusepath{clip}%
\pgfsetbuttcap%
\pgfsetroundjoin%
\definecolor{currentfill}{rgb}{0.874694,0.869527,0.925075}%
\pgfsetfillcolor{currentfill}%
\pgfsetfillopacity{0.900000}%
\pgfsetlinewidth{0.200750pt}%
\definecolor{currentstroke}{rgb}{0.200000,0.200000,0.200000}%
\pgfsetstrokecolor{currentstroke}%
\pgfsetdash{}{0pt}%
\pgfpathmoveto{\pgfqpoint{2.449730in}{2.029975in}}%
\pgfpathlineto{\pgfqpoint{2.493646in}{2.095694in}}%
\pgfpathlineto{\pgfqpoint{2.537541in}{2.182434in}}%
\pgfpathlineto{\pgfqpoint{2.568292in}{2.174660in}}%
\pgfpathlineto{\pgfqpoint{2.599149in}{2.166341in}}%
\pgfpathlineto{\pgfqpoint{2.555090in}{2.077005in}}%
\pgfpathlineto{\pgfqpoint{2.511037in}{2.007416in}}%
\pgfpathlineto{\pgfqpoint{2.480333in}{2.018753in}}%
\pgfpathlineto{\pgfqpoint{2.449730in}{2.029975in}}%
\pgfpathclose%
\pgfusepath{stroke,fill}%
\end{pgfscope}%
\begin{pgfscope}%
\pgfpathrectangle{\pgfqpoint{0.050000in}{0.050000in}}{\pgfqpoint{4.900000in}{4.900000in}}%
\pgfusepath{clip}%
\pgfsetbuttcap%
\pgfsetroundjoin%
\definecolor{currentfill}{rgb}{1.000000,1.000000,1.000000}%
\pgfsetfillcolor{currentfill}%
\pgfsetfillopacity{0.900000}%
\pgfsetlinewidth{0.200750pt}%
\definecolor{currentstroke}{rgb}{0.200000,0.200000,0.200000}%
\pgfsetstrokecolor{currentstroke}%
\pgfsetdash{}{0pt}%
\pgfpathmoveto{\pgfqpoint{1.587787in}{1.916516in}}%
\pgfpathlineto{\pgfqpoint{1.632735in}{1.930196in}}%
\pgfpathlineto{\pgfqpoint{1.677589in}{1.943847in}}%
\pgfpathlineto{\pgfqpoint{1.707019in}{1.924342in}}%
\pgfpathlineto{\pgfqpoint{1.736538in}{1.904778in}}%
\pgfpathlineto{\pgfqpoint{1.691611in}{1.891045in}}%
\pgfpathlineto{\pgfqpoint{1.646589in}{1.877282in}}%
\pgfpathlineto{\pgfqpoint{1.617144in}{1.896928in}}%
\pgfpathlineto{\pgfqpoint{1.587787in}{1.916516in}}%
\pgfpathclose%
\pgfusepath{stroke,fill}%
\end{pgfscope}%
\begin{pgfscope}%
\pgfpathrectangle{\pgfqpoint{0.050000in}{0.050000in}}{\pgfqpoint{4.900000in}{4.900000in}}%
\pgfusepath{clip}%
\pgfsetbuttcap%
\pgfsetroundjoin%
\definecolor{currentfill}{rgb}{0.254133,0.223376,0.554018}%
\pgfsetfillcolor{currentfill}%
\pgfsetfillopacity{0.900000}%
\pgfsetlinewidth{0.200750pt}%
\definecolor{currentstroke}{rgb}{0.200000,0.200000,0.200000}%
\pgfsetstrokecolor{currentstroke}%
\pgfsetdash{}{0pt}%
\pgfpathmoveto{\pgfqpoint{2.590158in}{2.789455in}}%
\pgfpathlineto{\pgfqpoint{2.634137in}{2.982410in}}%
\pgfpathlineto{\pgfqpoint{2.678302in}{3.176347in}}%
\pgfpathlineto{\pgfqpoint{2.709321in}{3.150263in}}%
\pgfpathlineto{\pgfqpoint{2.740415in}{3.121170in}}%
\pgfpathlineto{\pgfqpoint{2.696047in}{2.946945in}}%
\pgfpathlineto{\pgfqpoint{2.651824in}{2.766850in}}%
\pgfpathlineto{\pgfqpoint{2.620942in}{2.778875in}}%
\pgfpathlineto{\pgfqpoint{2.590158in}{2.789455in}}%
\pgfpathclose%
\pgfusepath{stroke,fill}%
\end{pgfscope}%
\begin{pgfscope}%
\pgfpathrectangle{\pgfqpoint{0.050000in}{0.050000in}}{\pgfqpoint{4.900000in}{4.900000in}}%
\pgfusepath{clip}%
\pgfsetbuttcap%
\pgfsetroundjoin%
\definecolor{currentfill}{rgb}{0.952265,0.950296,0.971457}%
\pgfsetfillcolor{currentfill}%
\pgfsetfillopacity{0.900000}%
\pgfsetlinewidth{0.200750pt}%
\definecolor{currentstroke}{rgb}{0.200000,0.200000,0.200000}%
\pgfsetstrokecolor{currentstroke}%
\pgfsetdash{}{0pt}%
\pgfpathmoveto{\pgfqpoint{2.361687in}{1.947027in}}%
\pgfpathlineto{\pgfqpoint{2.405752in}{1.981822in}}%
\pgfpathlineto{\pgfqpoint{2.449730in}{2.029975in}}%
\pgfpathlineto{\pgfqpoint{2.480333in}{2.018753in}}%
\pgfpathlineto{\pgfqpoint{2.511037in}{2.007416in}}%
\pgfpathlineto{\pgfqpoint{2.466948in}{1.954974in}}%
\pgfpathlineto{\pgfqpoint{2.422792in}{1.916181in}}%
\pgfpathlineto{\pgfqpoint{2.392192in}{1.931542in}}%
\pgfpathlineto{\pgfqpoint{2.361687in}{1.947027in}}%
\pgfpathclose%
\pgfusepath{stroke,fill}%
\end{pgfscope}%
\begin{pgfscope}%
\pgfpathrectangle{\pgfqpoint{0.050000in}{0.050000in}}{\pgfqpoint{4.900000in}{4.900000in}}%
\pgfusepath{clip}%
\pgfsetbuttcap%
\pgfsetroundjoin%
\definecolor{currentfill}{rgb}{0.265667,0.488889,0.918724}%
\pgfsetfillcolor{currentfill}%
\pgfsetfillopacity{0.900000}%
\pgfsetlinewidth{0.200750pt}%
\definecolor{currentstroke}{rgb}{0.200000,0.200000,0.200000}%
\pgfsetstrokecolor{currentstroke}%
\pgfsetdash{}{0pt}%
\pgfpathmoveto{\pgfqpoint{2.581342in}{3.726662in}}%
\pgfpathlineto{\pgfqpoint{2.625347in}{3.947313in}}%
\pgfpathlineto{\pgfqpoint{2.669540in}{4.148094in}}%
\pgfpathlineto{\pgfqpoint{2.700469in}{4.076559in}}%
\pgfpathlineto{\pgfqpoint{2.731374in}{3.984531in}}%
\pgfpathlineto{\pgfqpoint{2.687182in}{3.865959in}}%
\pgfpathlineto{\pgfqpoint{2.642973in}{3.681861in}}%
\pgfpathlineto{\pgfqpoint{2.612117in}{3.708030in}}%
\pgfpathlineto{\pgfqpoint{2.581342in}{3.726662in}}%
\pgfpathclose%
\pgfusepath{stroke,fill}%
\end{pgfscope}%
\begin{pgfscope}%
\pgfpathrectangle{\pgfqpoint{0.050000in}{0.050000in}}{\pgfqpoint{4.900000in}{4.900000in}}%
\pgfusepath{clip}%
\pgfsetbuttcap%
\pgfsetroundjoin%
\definecolor{currentfill}{rgb}{1.000000,1.000000,1.000000}%
\pgfsetfillcolor{currentfill}%
\pgfsetfillopacity{0.900000}%
\pgfsetlinewidth{0.200750pt}%
\definecolor{currentstroke}{rgb}{0.200000,0.200000,0.200000}%
\pgfsetstrokecolor{currentstroke}%
\pgfsetdash{}{0pt}%
\pgfpathmoveto{\pgfqpoint{2.123949in}{1.908976in}}%
\pgfpathlineto{\pgfqpoint{2.168356in}{1.923266in}}%
\pgfpathlineto{\pgfqpoint{2.212664in}{1.938573in}}%
\pgfpathlineto{\pgfqpoint{2.242912in}{1.920044in}}%
\pgfpathlineto{\pgfqpoint{2.273252in}{1.901641in}}%
\pgfpathlineto{\pgfqpoint{2.228878in}{1.885057in}}%
\pgfpathlineto{\pgfqpoint{2.184404in}{1.870091in}}%
\pgfpathlineto{\pgfqpoint{2.154131in}{1.889532in}}%
\pgfpathlineto{\pgfqpoint{2.123949in}{1.908976in}}%
\pgfpathclose%
\pgfusepath{stroke,fill}%
\end{pgfscope}%
\begin{pgfscope}%
\pgfpathrectangle{\pgfqpoint{0.050000in}{0.050000in}}{\pgfqpoint{4.900000in}{4.900000in}}%
\pgfusepath{clip}%
\pgfsetbuttcap%
\pgfsetroundjoin%
\definecolor{currentfill}{rgb}{0.361538,0.335210,0.618239}%
\pgfsetfillcolor{currentfill}%
\pgfsetfillopacity{0.900000}%
\pgfsetlinewidth{0.200750pt}%
\definecolor{currentstroke}{rgb}{0.200000,0.200000,0.200000}%
\pgfsetstrokecolor{currentstroke}%
\pgfsetdash{}{0pt}%
\pgfpathmoveto{\pgfqpoint{3.040898in}{3.031945in}}%
\pgfpathlineto{\pgfqpoint{3.083138in}{2.892181in}}%
\pgfpathlineto{\pgfqpoint{3.124617in}{2.707654in}}%
\pgfpathlineto{\pgfqpoint{3.155682in}{2.641861in}}%
\pgfpathlineto{\pgfqpoint{3.186830in}{2.579828in}}%
\pgfpathlineto{\pgfqpoint{3.145413in}{2.748155in}}%
\pgfpathlineto{\pgfqpoint{3.103320in}{2.883502in}}%
\pgfpathlineto{\pgfqpoint{3.072077in}{2.956097in}}%
\pgfpathlineto{\pgfqpoint{3.040898in}{3.031945in}}%
\pgfpathclose%
\pgfusepath{stroke,fill}%
\end{pgfscope}%
\begin{pgfscope}%
\pgfpathrectangle{\pgfqpoint{0.050000in}{0.050000in}}{\pgfqpoint{4.900000in}{4.900000in}}%
\pgfusepath{clip}%
\pgfsetbuttcap%
\pgfsetroundjoin%
\definecolor{currentfill}{rgb}{0.850827,0.844675,0.910804}%
\pgfsetfillcolor{currentfill}%
\pgfsetfillopacity{0.900000}%
\pgfsetlinewidth{0.200750pt}%
\definecolor{currentstroke}{rgb}{0.200000,0.200000,0.200000}%
\pgfsetstrokecolor{currentstroke}%
\pgfsetdash{}{0pt}%
\pgfpathmoveto{\pgfqpoint{3.375537in}{2.265171in}}%
\pgfpathlineto{\pgfqpoint{3.416749in}{2.146366in}}%
\pgfpathlineto{\pgfqpoint{3.458029in}{2.050184in}}%
\pgfpathlineto{\pgfqpoint{3.489952in}{2.021801in}}%
\pgfpathlineto{\pgfqpoint{3.521972in}{1.993961in}}%
\pgfpathlineto{\pgfqpoint{3.480495in}{2.073039in}}%
\pgfpathlineto{\pgfqpoint{3.439155in}{2.176000in}}%
\pgfpathlineto{\pgfqpoint{3.407300in}{2.219825in}}%
\pgfpathlineto{\pgfqpoint{3.375537in}{2.265171in}}%
\pgfpathclose%
\pgfusepath{stroke,fill}%
\end{pgfscope}%
\begin{pgfscope}%
\pgfpathrectangle{\pgfqpoint{0.050000in}{0.050000in}}{\pgfqpoint{4.900000in}{4.900000in}}%
\pgfusepath{clip}%
\pgfsetbuttcap%
\pgfsetroundjoin%
\definecolor{currentfill}{rgb}{0.270588,0.541176,0.988235}%
\pgfsetfillcolor{currentfill}%
\pgfsetfillopacity{0.900000}%
\pgfsetlinewidth{0.200750pt}%
\definecolor{currentstroke}{rgb}{0.200000,0.200000,0.200000}%
\pgfsetstrokecolor{currentstroke}%
\pgfsetdash{}{0pt}%
\pgfpathmoveto{\pgfqpoint{2.669540in}{4.148094in}}%
\pgfpathlineto{\pgfqpoint{2.713744in}{4.276763in}}%
\pgfpathlineto{\pgfqpoint{2.757446in}{4.226936in}}%
\pgfpathlineto{\pgfqpoint{2.787786in}{3.997155in}}%
\pgfpathlineto{\pgfqpoint{2.818119in}{3.796863in}}%
\pgfpathlineto{\pgfqpoint{2.775194in}{3.980281in}}%
\pgfpathlineto{\pgfqpoint{2.731374in}{3.984531in}}%
\pgfpathlineto{\pgfqpoint{2.700469in}{4.076559in}}%
\pgfpathlineto{\pgfqpoint{2.669540in}{4.148094in}}%
\pgfpathclose%
\pgfusepath{stroke,fill}%
\end{pgfscope}%
\begin{pgfscope}%
\pgfpathrectangle{\pgfqpoint{0.050000in}{0.050000in}}{\pgfqpoint{4.900000in}{4.900000in}}%
\pgfusepath{clip}%
\pgfsetbuttcap%
\pgfsetroundjoin%
\definecolor{currentfill}{rgb}{0.725521,0.714202,0.835879}%
\pgfsetfillcolor{currentfill}%
\pgfsetfillopacity{0.900000}%
\pgfsetlinewidth{0.200750pt}%
\definecolor{currentstroke}{rgb}{0.200000,0.200000,0.200000}%
\pgfsetstrokecolor{currentstroke}%
\pgfsetdash{}{0pt}%
\pgfpathmoveto{\pgfqpoint{2.537541in}{2.182434in}}%
\pgfpathlineto{\pgfqpoint{2.581464in}{2.292029in}}%
\pgfpathlineto{\pgfqpoint{2.625467in}{2.423963in}}%
\pgfpathlineto{\pgfqpoint{2.656419in}{2.415381in}}%
\pgfpathlineto{\pgfqpoint{2.687474in}{2.405724in}}%
\pgfpathlineto{\pgfqpoint{2.643262in}{2.276376in}}%
\pgfpathlineto{\pgfqpoint{2.599149in}{2.166341in}}%
\pgfpathlineto{\pgfqpoint{2.568292in}{2.174660in}}%
\pgfpathlineto{\pgfqpoint{2.537541in}{2.182434in}}%
\pgfpathclose%
\pgfusepath{stroke,fill}%
\end{pgfscope}%
\begin{pgfscope}%
\pgfpathrectangle{\pgfqpoint{0.050000in}{0.050000in}}{\pgfqpoint{4.900000in}{4.900000in}}%
\pgfusepath{clip}%
\pgfsetbuttcap%
\pgfsetroundjoin%
\definecolor{currentfill}{rgb}{1.000000,1.000000,1.000000}%
\pgfsetfillcolor{currentfill}%
\pgfsetfillopacity{0.900000}%
\pgfsetlinewidth{0.200750pt}%
\definecolor{currentstroke}{rgb}{0.200000,0.200000,0.200000}%
\pgfsetstrokecolor{currentstroke}%
\pgfsetdash{}{0pt}%
\pgfpathmoveto{\pgfqpoint{1.736538in}{1.904778in}}%
\pgfpathlineto{\pgfqpoint{1.781370in}{1.918483in}}%
\pgfpathlineto{\pgfqpoint{1.826108in}{1.932159in}}%
\pgfpathlineto{\pgfqpoint{1.855788in}{1.912619in}}%
\pgfpathlineto{\pgfqpoint{1.885558in}{1.893020in}}%
\pgfpathlineto{\pgfqpoint{1.840748in}{1.879261in}}%
\pgfpathlineto{\pgfqpoint{1.795843in}{1.865474in}}%
\pgfpathlineto{\pgfqpoint{1.766146in}{1.885156in}}%
\pgfpathlineto{\pgfqpoint{1.736538in}{1.904778in}}%
\pgfpathclose%
\pgfusepath{stroke,fill}%
\end{pgfscope}%
\begin{pgfscope}%
\pgfpathrectangle{\pgfqpoint{0.050000in}{0.050000in}}{\pgfqpoint{4.900000in}{4.900000in}}%
\pgfusepath{clip}%
\pgfsetbuttcap%
\pgfsetroundjoin%
\definecolor{currentfill}{rgb}{1.000000,1.000000,1.000000}%
\pgfsetfillcolor{currentfill}%
\pgfsetfillopacity{0.900000}%
\pgfsetlinewidth{0.200750pt}%
\definecolor{currentstroke}{rgb}{0.200000,0.200000,0.200000}%
\pgfsetstrokecolor{currentstroke}%
\pgfsetdash{}{0pt}%
\pgfpathmoveto{\pgfqpoint{1.348890in}{1.900834in}}%
\pgfpathlineto{\pgfqpoint{1.394144in}{1.914547in}}%
\pgfpathlineto{\pgfqpoint{1.439303in}{1.928232in}}%
\pgfpathlineto{\pgfqpoint{1.468410in}{1.908680in}}%
\pgfpathlineto{\pgfqpoint{1.497605in}{1.889069in}}%
\pgfpathlineto{\pgfqpoint{1.452371in}{1.875302in}}%
\pgfpathlineto{\pgfqpoint{1.407042in}{1.861506in}}%
\pgfpathlineto{\pgfqpoint{1.377922in}{1.881200in}}%
\pgfpathlineto{\pgfqpoint{1.348890in}{1.900834in}}%
\pgfpathclose%
\pgfusepath{stroke,fill}%
\end{pgfscope}%
\begin{pgfscope}%
\pgfpathrectangle{\pgfqpoint{0.050000in}{0.050000in}}{\pgfqpoint{4.900000in}{4.900000in}}%
\pgfusepath{clip}%
\pgfsetbuttcap%
\pgfsetroundjoin%
\definecolor{currentfill}{rgb}{0.250903,0.332026,0.710188}%
\pgfsetfillcolor{currentfill}%
\pgfsetfillopacity{0.900000}%
\pgfsetlinewidth{0.200750pt}%
\definecolor{currentstroke}{rgb}{0.200000,0.200000,0.200000}%
\pgfsetstrokecolor{currentstroke}%
\pgfsetdash{}{0pt}%
\pgfpathmoveto{\pgfqpoint{2.616509in}{3.219158in}}%
\pgfpathlineto{\pgfqpoint{2.660618in}{3.422898in}}%
\pgfpathlineto{\pgfqpoint{2.704872in}{3.603562in}}%
\pgfpathlineto{\pgfqpoint{2.735894in}{3.553499in}}%
\pgfpathlineto{\pgfqpoint{2.766957in}{3.498284in}}%
\pgfpathlineto{\pgfqpoint{2.722612in}{3.355504in}}%
\pgfpathlineto{\pgfqpoint{2.678302in}{3.176347in}}%
\pgfpathlineto{\pgfqpoint{2.647363in}{3.199288in}}%
\pgfpathlineto{\pgfqpoint{2.616509in}{3.219158in}}%
\pgfpathclose%
\pgfusepath{stroke,fill}%
\end{pgfscope}%
\begin{pgfscope}%
\pgfpathrectangle{\pgfqpoint{0.050000in}{0.050000in}}{\pgfqpoint{4.900000in}{4.900000in}}%
\pgfusepath{clip}%
\pgfsetbuttcap%
\pgfsetroundjoin%
\definecolor{currentfill}{rgb}{0.988066,0.987574,0.992864}%
\pgfsetfillcolor{currentfill}%
\pgfsetfillopacity{0.900000}%
\pgfsetlinewidth{0.200750pt}%
\definecolor{currentstroke}{rgb}{0.200000,0.200000,0.200000}%
\pgfsetstrokecolor{currentstroke}%
\pgfsetdash{}{0pt}%
\pgfpathmoveto{\pgfqpoint{2.273252in}{1.901641in}}%
\pgfpathlineto{\pgfqpoint{2.317522in}{1.921456in}}%
\pgfpathlineto{\pgfqpoint{2.361687in}{1.947027in}}%
\pgfpathlineto{\pgfqpoint{2.392192in}{1.931542in}}%
\pgfpathlineto{\pgfqpoint{2.422792in}{1.916181in}}%
\pgfpathlineto{\pgfqpoint{2.378549in}{1.887356in}}%
\pgfpathlineto{\pgfqpoint{2.334208in}{1.865195in}}%
\pgfpathlineto{\pgfqpoint{2.303683in}{1.883361in}}%
\pgfpathlineto{\pgfqpoint{2.273252in}{1.901641in}}%
\pgfpathclose%
\pgfusepath{stroke,fill}%
\end{pgfscope}%
\begin{pgfscope}%
\pgfpathrectangle{\pgfqpoint{0.050000in}{0.050000in}}{\pgfqpoint{4.900000in}{4.900000in}}%
\pgfusepath{clip}%
\pgfsetbuttcap%
\pgfsetroundjoin%
\definecolor{currentfill}{rgb}{1.000000,1.000000,1.000000}%
\pgfsetfillcolor{currentfill}%
\pgfsetfillopacity{0.900000}%
\pgfsetlinewidth{0.200750pt}%
\definecolor{currentstroke}{rgb}{0.200000,0.200000,0.200000}%
\pgfsetstrokecolor{currentstroke}%
\pgfsetdash{}{0pt}%
\pgfpathmoveto{\pgfqpoint{1.885558in}{1.893020in}}%
\pgfpathlineto{\pgfqpoint{1.930274in}{1.906750in}}%
\pgfpathlineto{\pgfqpoint{1.974895in}{1.920453in}}%
\pgfpathlineto{\pgfqpoint{2.004826in}{1.900881in}}%
\pgfpathlineto{\pgfqpoint{2.034848in}{1.881256in}}%
\pgfpathlineto{\pgfqpoint{1.990155in}{1.867460in}}%
\pgfpathlineto{\pgfqpoint{1.945367in}{1.853645in}}%
\pgfpathlineto{\pgfqpoint{1.915418in}{1.873362in}}%
\pgfpathlineto{\pgfqpoint{1.885558in}{1.893020in}}%
\pgfpathclose%
\pgfusepath{stroke,fill}%
\end{pgfscope}%
\begin{pgfscope}%
\pgfpathrectangle{\pgfqpoint{0.050000in}{0.050000in}}{\pgfqpoint{4.900000in}{4.900000in}}%
\pgfusepath{clip}%
\pgfsetbuttcap%
\pgfsetroundjoin%
\definecolor{currentfill}{rgb}{0.262714,0.457516,0.877017}%
\pgfsetfillcolor{currentfill}%
\pgfsetfillopacity{0.900000}%
\pgfsetlinewidth{0.200750pt}%
\definecolor{currentstroke}{rgb}{0.200000,0.200000,0.200000}%
\pgfsetstrokecolor{currentstroke}%
\pgfsetdash{}{0pt}%
\pgfpathmoveto{\pgfqpoint{2.642973in}{3.681861in}}%
\pgfpathlineto{\pgfqpoint{2.687182in}{3.865959in}}%
\pgfpathlineto{\pgfqpoint{2.731374in}{3.984531in}}%
\pgfpathlineto{\pgfqpoint{2.762258in}{3.882156in}}%
\pgfpathlineto{\pgfqpoint{2.793130in}{3.776190in}}%
\pgfpathlineto{\pgfqpoint{2.749128in}{3.733524in}}%
\pgfpathlineto{\pgfqpoint{2.704872in}{3.603562in}}%
\pgfpathlineto{\pgfqpoint{2.673895in}{3.646838in}}%
\pgfpathlineto{\pgfqpoint{2.642973in}{3.681861in}}%
\pgfpathclose%
\pgfusepath{stroke,fill}%
\end{pgfscope}%
\begin{pgfscope}%
\pgfpathrectangle{\pgfqpoint{0.050000in}{0.050000in}}{\pgfqpoint{4.900000in}{4.900000in}}%
\pgfusepath{clip}%
\pgfsetbuttcap%
\pgfsetroundjoin%
\definecolor{currentfill}{rgb}{0.445075,0.422191,0.668189}%
\pgfsetfillcolor{currentfill}%
\pgfsetfillopacity{0.900000}%
\pgfsetlinewidth{0.200750pt}%
\definecolor{currentstroke}{rgb}{0.200000,0.200000,0.200000}%
\pgfsetstrokecolor{currentstroke}%
\pgfsetdash{}{0pt}%
\pgfpathmoveto{\pgfqpoint{3.103320in}{2.883502in}}%
\pgfpathlineto{\pgfqpoint{3.145413in}{2.748155in}}%
\pgfpathlineto{\pgfqpoint{3.186830in}{2.579828in}}%
\pgfpathlineto{\pgfqpoint{3.218063in}{2.521091in}}%
\pgfpathlineto{\pgfqpoint{3.249382in}{2.465264in}}%
\pgfpathlineto{\pgfqpoint{3.207990in}{2.618075in}}%
\pgfpathlineto{\pgfqpoint{3.166004in}{2.747127in}}%
\pgfpathlineto{\pgfqpoint{3.134628in}{2.813921in}}%
\pgfpathlineto{\pgfqpoint{3.103320in}{2.883502in}}%
\pgfpathclose%
\pgfusepath{stroke,fill}%
\end{pgfscope}%
\begin{pgfscope}%
\pgfpathrectangle{\pgfqpoint{0.050000in}{0.050000in}}{\pgfqpoint{4.900000in}{4.900000in}}%
\pgfusepath{clip}%
\pgfsetbuttcap%
\pgfsetroundjoin%
\definecolor{currentfill}{rgb}{1.000000,1.000000,1.000000}%
\pgfsetfillcolor{currentfill}%
\pgfsetfillopacity{0.900000}%
\pgfsetlinewidth{0.200750pt}%
\definecolor{currentstroke}{rgb}{0.200000,0.200000,0.200000}%
\pgfsetstrokecolor{currentstroke}%
\pgfsetdash{}{0pt}%
\pgfpathmoveto{\pgfqpoint{1.497605in}{1.889069in}}%
\pgfpathlineto{\pgfqpoint{1.542743in}{1.902807in}}%
\pgfpathlineto{\pgfqpoint{1.587787in}{1.916516in}}%
\pgfpathlineto{\pgfqpoint{1.617144in}{1.896928in}}%
\pgfpathlineto{\pgfqpoint{1.646589in}{1.877282in}}%
\pgfpathlineto{\pgfqpoint{1.601472in}{1.863490in}}%
\pgfpathlineto{\pgfqpoint{1.556259in}{1.849669in}}%
\pgfpathlineto{\pgfqpoint{1.526888in}{1.869398in}}%
\pgfpathlineto{\pgfqpoint{1.497605in}{1.889069in}}%
\pgfpathclose%
\pgfusepath{stroke,fill}%
\end{pgfscope}%
\begin{pgfscope}%
\pgfpathrectangle{\pgfqpoint{0.050000in}{0.050000in}}{\pgfqpoint{4.900000in}{4.900000in}}%
\pgfusepath{clip}%
\pgfsetbuttcap%
\pgfsetroundjoin%
\definecolor{currentfill}{rgb}{0.988066,0.987574,0.992864}%
\pgfsetfillcolor{currentfill}%
\pgfsetfillopacity{0.900000}%
\pgfsetlinewidth{0.200750pt}%
\definecolor{currentstroke}{rgb}{0.200000,0.200000,0.200000}%
\pgfsetstrokecolor{currentstroke}%
\pgfsetdash{}{0pt}%
\pgfpathmoveto{\pgfqpoint{3.586303in}{1.939694in}}%
\pgfpathlineto{\pgfqpoint{3.628463in}{1.912058in}}%
\pgfpathlineto{\pgfqpoint{3.660902in}{1.892034in}}%
\pgfpathlineto{\pgfqpoint{3.693438in}{1.871977in}}%
\pgfpathlineto{\pgfqpoint{3.651024in}{1.887006in}}%
\pgfpathlineto{\pgfqpoint{3.618615in}{1.913173in}}%
\pgfpathlineto{\pgfqpoint{3.586303in}{1.939694in}}%
\pgfpathclose%
\pgfusepath{stroke,fill}%
\end{pgfscope}%
\begin{pgfscope}%
\pgfpathrectangle{\pgfqpoint{0.050000in}{0.050000in}}{\pgfqpoint{4.900000in}{4.900000in}}%
\pgfusepath{clip}%
\pgfsetbuttcap%
\pgfsetroundjoin%
\definecolor{currentfill}{rgb}{1.000000,1.000000,1.000000}%
\pgfsetfillcolor{currentfill}%
\pgfsetfillopacity{0.900000}%
\pgfsetlinewidth{0.200750pt}%
\definecolor{currentstroke}{rgb}{0.200000,0.200000,0.200000}%
\pgfsetstrokecolor{currentstroke}%
\pgfsetdash{}{0pt}%
\pgfpathmoveto{\pgfqpoint{2.034848in}{1.881256in}}%
\pgfpathlineto{\pgfqpoint{2.079446in}{1.895062in}}%
\pgfpathlineto{\pgfqpoint{2.123949in}{1.908976in}}%
\pgfpathlineto{\pgfqpoint{2.154131in}{1.889532in}}%
\pgfpathlineto{\pgfqpoint{2.184404in}{1.870091in}}%
\pgfpathlineto{\pgfqpoint{2.139831in}{1.855837in}}%
\pgfpathlineto{\pgfqpoint{2.095163in}{1.841851in}}%
\pgfpathlineto{\pgfqpoint{2.064960in}{1.861578in}}%
\pgfpathlineto{\pgfqpoint{2.034848in}{1.881256in}}%
\pgfpathclose%
\pgfusepath{stroke,fill}%
\end{pgfscope}%
\begin{pgfscope}%
\pgfpathrectangle{\pgfqpoint{0.050000in}{0.050000in}}{\pgfqpoint{4.900000in}{4.900000in}}%
\pgfusepath{clip}%
\pgfsetbuttcap%
\pgfsetroundjoin%
\definecolor{currentfill}{rgb}{0.258531,0.413072,0.817932}%
\pgfsetfillcolor{currentfill}%
\pgfsetfillopacity{0.900000}%
\pgfsetlinewidth{0.200750pt}%
\definecolor{currentstroke}{rgb}{0.200000,0.200000,0.200000}%
\pgfsetstrokecolor{currentstroke}%
\pgfsetdash{}{0pt}%
\pgfpathmoveto{\pgfqpoint{2.704872in}{3.603562in}}%
\pgfpathlineto{\pgfqpoint{2.749128in}{3.733524in}}%
\pgfpathlineto{\pgfqpoint{2.793130in}{3.776190in}}%
\pgfpathlineto{\pgfqpoint{2.824001in}{3.670584in}}%
\pgfpathlineto{\pgfqpoint{2.854884in}{3.567495in}}%
\pgfpathlineto{\pgfqpoint{2.811141in}{3.577953in}}%
\pgfpathlineto{\pgfqpoint{2.766957in}{3.498284in}}%
\pgfpathlineto{\pgfqpoint{2.735894in}{3.553499in}}%
\pgfpathlineto{\pgfqpoint{2.704872in}{3.603562in}}%
\pgfpathclose%
\pgfusepath{stroke,fill}%
\end{pgfscope}%
\begin{pgfscope}%
\pgfpathrectangle{\pgfqpoint{0.050000in}{0.050000in}}{\pgfqpoint{4.900000in}{4.900000in}}%
\pgfusepath{clip}%
\pgfsetbuttcap%
\pgfsetroundjoin%
\definecolor{currentfill}{rgb}{1.000000,1.000000,1.000000}%
\pgfsetfillcolor{currentfill}%
\pgfsetfillopacity{0.900000}%
\pgfsetlinewidth{0.200750pt}%
\definecolor{currentstroke}{rgb}{0.200000,0.200000,0.200000}%
\pgfsetstrokecolor{currentstroke}%
\pgfsetdash{}{0pt}%
\pgfpathmoveto{\pgfqpoint{1.646589in}{1.877282in}}%
\pgfpathlineto{\pgfqpoint{1.691611in}{1.891045in}}%
\pgfpathlineto{\pgfqpoint{1.736538in}{1.904778in}}%
\pgfpathlineto{\pgfqpoint{1.766146in}{1.885156in}}%
\pgfpathlineto{\pgfqpoint{1.795843in}{1.865474in}}%
\pgfpathlineto{\pgfqpoint{1.750843in}{1.851657in}}%
\pgfpathlineto{\pgfqpoint{1.705748in}{1.837811in}}%
\pgfpathlineto{\pgfqpoint{1.676124in}{1.857576in}}%
\pgfpathlineto{\pgfqpoint{1.646589in}{1.877282in}}%
\pgfpathclose%
\pgfusepath{stroke,fill}%
\end{pgfscope}%
\begin{pgfscope}%
\pgfpathrectangle{\pgfqpoint{0.050000in}{0.050000in}}{\pgfqpoint{4.900000in}{4.900000in}}%
\pgfusepath{clip}%
\pgfsetbuttcap%
\pgfsetroundjoin%
\definecolor{currentfill}{rgb}{0.510711,0.490534,0.707436}%
\pgfsetfillcolor{currentfill}%
\pgfsetfillopacity{0.900000}%
\pgfsetlinewidth{0.200750pt}%
\definecolor{currentstroke}{rgb}{0.200000,0.200000,0.200000}%
\pgfsetstrokecolor{currentstroke}%
\pgfsetdash{}{0pt}%
\pgfpathmoveto{\pgfqpoint{2.625467in}{2.423963in}}%
\pgfpathlineto{\pgfqpoint{2.669596in}{2.574815in}}%
\pgfpathlineto{\pgfqpoint{2.713877in}{2.737843in}}%
\pgfpathlineto{\pgfqpoint{2.745043in}{2.720732in}}%
\pgfpathlineto{\pgfqpoint{2.776298in}{2.701891in}}%
\pgfpathlineto{\pgfqpoint{2.731817in}{2.550097in}}%
\pgfpathlineto{\pgfqpoint{2.687474in}{2.405724in}}%
\pgfpathlineto{\pgfqpoint{2.656419in}{2.415381in}}%
\pgfpathlineto{\pgfqpoint{2.625467in}{2.423963in}}%
\pgfpathclose%
\pgfusepath{stroke,fill}%
\end{pgfscope}%
\begin{pgfscope}%
\pgfpathrectangle{\pgfqpoint{0.050000in}{0.050000in}}{\pgfqpoint{4.900000in}{4.900000in}}%
\pgfusepath{clip}%
\pgfsetbuttcap%
\pgfsetroundjoin%
\definecolor{currentfill}{rgb}{0.260100,0.229589,0.557586}%
\pgfsetfillcolor{currentfill}%
\pgfsetfillopacity{0.900000}%
\pgfsetlinewidth{0.200750pt}%
\definecolor{currentstroke}{rgb}{0.200000,0.200000,0.200000}%
\pgfsetstrokecolor{currentstroke}%
\pgfsetdash{}{0pt}%
\pgfpathmoveto{\pgfqpoint{2.651824in}{2.766850in}}%
\pgfpathlineto{\pgfqpoint{2.696047in}{2.946945in}}%
\pgfpathlineto{\pgfqpoint{2.740415in}{3.121170in}}%
\pgfpathlineto{\pgfqpoint{2.771582in}{3.089324in}}%
\pgfpathlineto{\pgfqpoint{2.802818in}{3.055041in}}%
\pgfpathlineto{\pgfqpoint{2.758303in}{2.902642in}}%
\pgfpathlineto{\pgfqpoint{2.713877in}{2.737843in}}%
\pgfpathlineto{\pgfqpoint{2.682804in}{2.753205in}}%
\pgfpathlineto{\pgfqpoint{2.651824in}{2.766850in}}%
\pgfpathclose%
\pgfusepath{stroke,fill}%
\end{pgfscope}%
\begin{pgfscope}%
\pgfpathrectangle{\pgfqpoint{0.050000in}{0.050000in}}{\pgfqpoint{4.900000in}{4.900000in}}%
\pgfusepath{clip}%
\pgfsetbuttcap%
\pgfsetroundjoin%
\definecolor{currentfill}{rgb}{0.874694,0.869527,0.925075}%
\pgfsetfillcolor{currentfill}%
\pgfsetfillopacity{0.900000}%
\pgfsetlinewidth{0.200750pt}%
\definecolor{currentstroke}{rgb}{0.200000,0.200000,0.200000}%
\pgfsetstrokecolor{currentstroke}%
\pgfsetdash{}{0pt}%
\pgfpathmoveto{\pgfqpoint{3.439155in}{2.176000in}}%
\pgfpathlineto{\pgfqpoint{3.480495in}{2.073039in}}%
\pgfpathlineto{\pgfqpoint{3.521972in}{1.993961in}}%
\pgfpathlineto{\pgfqpoint{3.554089in}{1.966609in}}%
\pgfpathlineto{\pgfqpoint{3.586303in}{1.939694in}}%
\pgfpathlineto{\pgfqpoint{3.544622in}{2.003659in}}%
\pgfpathlineto{\pgfqpoint{3.503143in}{2.092380in}}%
\pgfpathlineto{\pgfqpoint{3.471102in}{2.133559in}}%
\pgfpathlineto{\pgfqpoint{3.439155in}{2.176000in}}%
\pgfpathclose%
\pgfusepath{stroke,fill}%
\end{pgfscope}%
\begin{pgfscope}%
\pgfpathrectangle{\pgfqpoint{0.050000in}{0.050000in}}{\pgfqpoint{4.900000in}{4.900000in}}%
\pgfusepath{clip}%
\pgfsetbuttcap%
\pgfsetroundjoin%
\definecolor{currentfill}{rgb}{0.856794,0.850888,0.914371}%
\pgfsetfillcolor{currentfill}%
\pgfsetfillopacity{0.900000}%
\pgfsetlinewidth{0.200750pt}%
\definecolor{currentstroke}{rgb}{0.200000,0.200000,0.200000}%
\pgfsetstrokecolor{currentstroke}%
\pgfsetdash{}{0pt}%
\pgfpathmoveto{\pgfqpoint{2.511037in}{2.007416in}}%
\pgfpathlineto{\pgfqpoint{2.555090in}{2.077005in}}%
\pgfpathlineto{\pgfqpoint{2.599149in}{2.166341in}}%
\pgfpathlineto{\pgfqpoint{2.630110in}{2.157460in}}%
\pgfpathlineto{\pgfqpoint{2.661175in}{2.147995in}}%
\pgfpathlineto{\pgfqpoint{2.616947in}{2.056977in}}%
\pgfpathlineto{\pgfqpoint{2.572750in}{1.984302in}}%
\pgfpathlineto{\pgfqpoint{2.541842in}{1.995941in}}%
\pgfpathlineto{\pgfqpoint{2.511037in}{2.007416in}}%
\pgfpathclose%
\pgfusepath{stroke,fill}%
\end{pgfscope}%
\begin{pgfscope}%
\pgfpathrectangle{\pgfqpoint{0.050000in}{0.050000in}}{\pgfqpoint{4.900000in}{4.900000in}}%
\pgfusepath{clip}%
\pgfsetbuttcap%
\pgfsetroundjoin%
\definecolor{currentfill}{rgb}{1.000000,1.000000,1.000000}%
\pgfsetfillcolor{currentfill}%
\pgfsetfillopacity{0.900000}%
\pgfsetlinewidth{0.200750pt}%
\definecolor{currentstroke}{rgb}{0.200000,0.200000,0.200000}%
\pgfsetstrokecolor{currentstroke}%
\pgfsetdash{}{0pt}%
\pgfpathmoveto{\pgfqpoint{1.258094in}{1.873321in}}%
\pgfpathlineto{\pgfqpoint{1.303540in}{1.887092in}}%
\pgfpathlineto{\pgfqpoint{1.348890in}{1.900834in}}%
\pgfpathlineto{\pgfqpoint{1.377922in}{1.881200in}}%
\pgfpathlineto{\pgfqpoint{1.407042in}{1.861506in}}%
\pgfpathlineto{\pgfqpoint{1.361616in}{1.847680in}}%
\pgfpathlineto{\pgfqpoint{1.316095in}{1.833826in}}%
\pgfpathlineto{\pgfqpoint{1.287051in}{1.853603in}}%
\pgfpathlineto{\pgfqpoint{1.258094in}{1.873321in}}%
\pgfpathclose%
\pgfusepath{stroke,fill}%
\end{pgfscope}%
\begin{pgfscope}%
\pgfpathrectangle{\pgfqpoint{0.050000in}{0.050000in}}{\pgfqpoint{4.900000in}{4.900000in}}%
\pgfusepath{clip}%
\pgfsetbuttcap%
\pgfsetroundjoin%
\definecolor{currentfill}{rgb}{0.249427,0.316340,0.689335}%
\pgfsetfillcolor{currentfill}%
\pgfsetfillopacity{0.900000}%
\pgfsetlinewidth{0.200750pt}%
\definecolor{currentstroke}{rgb}{0.200000,0.200000,0.200000}%
\pgfsetstrokecolor{currentstroke}%
\pgfsetdash{}{0pt}%
\pgfpathmoveto{\pgfqpoint{2.678302in}{3.176347in}}%
\pgfpathlineto{\pgfqpoint{2.722612in}{3.355504in}}%
\pgfpathlineto{\pgfqpoint{2.766957in}{3.498284in}}%
\pgfpathlineto{\pgfqpoint{2.798059in}{3.439402in}}%
\pgfpathlineto{\pgfqpoint{2.829199in}{3.378084in}}%
\pgfpathlineto{\pgfqpoint{2.784850in}{3.271959in}}%
\pgfpathlineto{\pgfqpoint{2.740415in}{3.121170in}}%
\pgfpathlineto{\pgfqpoint{2.709321in}{3.150263in}}%
\pgfpathlineto{\pgfqpoint{2.678302in}{3.176347in}}%
\pgfpathclose%
\pgfusepath{stroke,fill}%
\end{pgfscope}%
\begin{pgfscope}%
\pgfpathrectangle{\pgfqpoint{0.050000in}{0.050000in}}{\pgfqpoint{4.900000in}{4.900000in}}%
\pgfusepath{clip}%
\pgfsetbuttcap%
\pgfsetroundjoin%
\definecolor{currentfill}{rgb}{0.516678,0.496747,0.711003}%
\pgfsetfillcolor{currentfill}%
\pgfsetfillopacity{0.900000}%
\pgfsetlinewidth{0.200750pt}%
\definecolor{currentstroke}{rgb}{0.200000,0.200000,0.200000}%
\pgfsetstrokecolor{currentstroke}%
\pgfsetdash{}{0pt}%
\pgfpathmoveto{\pgfqpoint{3.166004in}{2.747127in}}%
\pgfpathlineto{\pgfqpoint{3.207990in}{2.618075in}}%
\pgfpathlineto{\pgfqpoint{3.249382in}{2.465264in}}%
\pgfpathlineto{\pgfqpoint{3.280788in}{2.412018in}}%
\pgfpathlineto{\pgfqpoint{3.312282in}{2.361074in}}%
\pgfpathlineto{\pgfqpoint{3.270885in}{2.499296in}}%
\pgfpathlineto{\pgfqpoint{3.228975in}{2.621079in}}%
\pgfpathlineto{\pgfqpoint{3.197453in}{2.682911in}}%
\pgfpathlineto{\pgfqpoint{3.166004in}{2.747127in}}%
\pgfpathclose%
\pgfusepath{stroke,fill}%
\end{pgfscope}%
\begin{pgfscope}%
\pgfpathrectangle{\pgfqpoint{0.050000in}{0.050000in}}{\pgfqpoint{4.900000in}{4.900000in}}%
\pgfusepath{clip}%
\pgfsetbuttcap%
\pgfsetroundjoin%
\definecolor{currentfill}{rgb}{0.946298,0.944083,0.967889}%
\pgfsetfillcolor{currentfill}%
\pgfsetfillopacity{0.900000}%
\pgfsetlinewidth{0.200750pt}%
\definecolor{currentstroke}{rgb}{0.200000,0.200000,0.200000}%
\pgfsetstrokecolor{currentstroke}%
\pgfsetdash{}{0pt}%
\pgfpathmoveto{\pgfqpoint{2.422792in}{1.916181in}}%
\pgfpathlineto{\pgfqpoint{2.466948in}{1.954974in}}%
\pgfpathlineto{\pgfqpoint{2.511037in}{2.007416in}}%
\pgfpathlineto{\pgfqpoint{2.541842in}{1.995941in}}%
\pgfpathlineto{\pgfqpoint{2.572750in}{1.984302in}}%
\pgfpathlineto{\pgfqpoint{2.528540in}{1.928166in}}%
\pgfpathlineto{\pgfqpoint{2.484284in}{1.885727in}}%
\pgfpathlineto{\pgfqpoint{2.453489in}{1.900918in}}%
\pgfpathlineto{\pgfqpoint{2.422792in}{1.916181in}}%
\pgfpathclose%
\pgfusepath{stroke,fill}%
\end{pgfscope}%
\begin{pgfscope}%
\pgfpathrectangle{\pgfqpoint{0.050000in}{0.050000in}}{\pgfqpoint{4.900000in}{4.900000in}}%
\pgfusepath{clip}%
\pgfsetbuttcap%
\pgfsetroundjoin%
\definecolor{currentfill}{rgb}{1.000000,1.000000,1.000000}%
\pgfsetfillcolor{currentfill}%
\pgfsetfillopacity{0.900000}%
\pgfsetlinewidth{0.200750pt}%
\definecolor{currentstroke}{rgb}{0.200000,0.200000,0.200000}%
\pgfsetstrokecolor{currentstroke}%
\pgfsetdash{}{0pt}%
\pgfpathmoveto{\pgfqpoint{2.184404in}{1.870091in}}%
\pgfpathlineto{\pgfqpoint{2.228878in}{1.885057in}}%
\pgfpathlineto{\pgfqpoint{2.273252in}{1.901641in}}%
\pgfpathlineto{\pgfqpoint{2.303683in}{1.883361in}}%
\pgfpathlineto{\pgfqpoint{2.334208in}{1.865195in}}%
\pgfpathlineto{\pgfqpoint{2.289766in}{1.847097in}}%
\pgfpathlineto{\pgfqpoint{2.245224in}{1.831244in}}%
\pgfpathlineto{\pgfqpoint{2.214768in}{1.850660in}}%
\pgfpathlineto{\pgfqpoint{2.184404in}{1.870091in}}%
\pgfpathclose%
\pgfusepath{stroke,fill}%
\end{pgfscope}%
\begin{pgfscope}%
\pgfpathrectangle{\pgfqpoint{0.050000in}{0.050000in}}{\pgfqpoint{4.900000in}{4.900000in}}%
\pgfusepath{clip}%
\pgfsetbuttcap%
\pgfsetroundjoin%
\definecolor{currentfill}{rgb}{0.254348,0.368627,0.758847}%
\pgfsetfillcolor{currentfill}%
\pgfsetfillopacity{0.900000}%
\pgfsetlinewidth{0.200750pt}%
\definecolor{currentstroke}{rgb}{0.200000,0.200000,0.200000}%
\pgfsetstrokecolor{currentstroke}%
\pgfsetdash{}{0pt}%
\pgfpathmoveto{\pgfqpoint{2.766957in}{3.498284in}}%
\pgfpathlineto{\pgfqpoint{2.811141in}{3.577953in}}%
\pgfpathlineto{\pgfqpoint{2.854884in}{3.567495in}}%
\pgfpathlineto{\pgfqpoint{2.885788in}{3.468014in}}%
\pgfpathlineto{\pgfqpoint{2.916723in}{3.372617in}}%
\pgfpathlineto{\pgfqpoint{2.873243in}{3.417484in}}%
\pgfpathlineto{\pgfqpoint{2.829199in}{3.378084in}}%
\pgfpathlineto{\pgfqpoint{2.798059in}{3.439402in}}%
\pgfpathlineto{\pgfqpoint{2.766957in}{3.498284in}}%
\pgfpathclose%
\pgfusepath{stroke,fill}%
\end{pgfscope}%
\begin{pgfscope}%
\pgfpathrectangle{\pgfqpoint{0.050000in}{0.050000in}}{\pgfqpoint{4.900000in}{4.900000in}}%
\pgfusepath{clip}%
\pgfsetbuttcap%
\pgfsetroundjoin%
\definecolor{currentfill}{rgb}{1.000000,1.000000,1.000000}%
\pgfsetfillcolor{currentfill}%
\pgfsetfillopacity{0.900000}%
\pgfsetlinewidth{0.200750pt}%
\definecolor{currentstroke}{rgb}{0.200000,0.200000,0.200000}%
\pgfsetstrokecolor{currentstroke}%
\pgfsetdash{}{0pt}%
\pgfpathmoveto{\pgfqpoint{1.795843in}{1.865474in}}%
\pgfpathlineto{\pgfqpoint{1.840748in}{1.879261in}}%
\pgfpathlineto{\pgfqpoint{1.885558in}{1.893020in}}%
\pgfpathlineto{\pgfqpoint{1.915418in}{1.873362in}}%
\pgfpathlineto{\pgfqpoint{1.945367in}{1.853645in}}%
\pgfpathlineto{\pgfqpoint{1.900485in}{1.839802in}}%
\pgfpathlineto{\pgfqpoint{1.855507in}{1.825931in}}%
\pgfpathlineto{\pgfqpoint{1.825630in}{1.845732in}}%
\pgfpathlineto{\pgfqpoint{1.795843in}{1.865474in}}%
\pgfpathclose%
\pgfusepath{stroke,fill}%
\end{pgfscope}%
\begin{pgfscope}%
\pgfpathrectangle{\pgfqpoint{0.050000in}{0.050000in}}{\pgfqpoint{4.900000in}{4.900000in}}%
\pgfusepath{clip}%
\pgfsetbuttcap%
\pgfsetroundjoin%
\definecolor{currentfill}{rgb}{0.707620,0.695563,0.825175}%
\pgfsetfillcolor{currentfill}%
\pgfsetfillopacity{0.900000}%
\pgfsetlinewidth{0.200750pt}%
\definecolor{currentstroke}{rgb}{0.200000,0.200000,0.200000}%
\pgfsetstrokecolor{currentstroke}%
\pgfsetdash{}{0pt}%
\pgfpathmoveto{\pgfqpoint{2.599149in}{2.166341in}}%
\pgfpathlineto{\pgfqpoint{2.643262in}{2.276376in}}%
\pgfpathlineto{\pgfqpoint{2.687474in}{2.405724in}}%
\pgfpathlineto{\pgfqpoint{2.718630in}{2.394959in}}%
\pgfpathlineto{\pgfqpoint{2.749886in}{2.383067in}}%
\pgfpathlineto{\pgfqpoint{2.705476in}{2.257445in}}%
\pgfpathlineto{\pgfqpoint{2.661175in}{2.147995in}}%
\pgfpathlineto{\pgfqpoint{2.630110in}{2.157460in}}%
\pgfpathlineto{\pgfqpoint{2.599149in}{2.166341in}}%
\pgfpathclose%
\pgfusepath{stroke,fill}%
\end{pgfscope}%
\begin{pgfscope}%
\pgfpathrectangle{\pgfqpoint{0.050000in}{0.050000in}}{\pgfqpoint{4.900000in}{4.900000in}}%
\pgfusepath{clip}%
\pgfsetbuttcap%
\pgfsetroundjoin%
\definecolor{currentfill}{rgb}{1.000000,1.000000,1.000000}%
\pgfsetfillcolor{currentfill}%
\pgfsetfillopacity{0.900000}%
\pgfsetlinewidth{0.200750pt}%
\definecolor{currentstroke}{rgb}{0.200000,0.200000,0.200000}%
\pgfsetstrokecolor{currentstroke}%
\pgfsetdash{}{0pt}%
\pgfpathmoveto{\pgfqpoint{1.407042in}{1.861506in}}%
\pgfpathlineto{\pgfqpoint{1.452371in}{1.875302in}}%
\pgfpathlineto{\pgfqpoint{1.497605in}{1.889069in}}%
\pgfpathlineto{\pgfqpoint{1.526888in}{1.869398in}}%
\pgfpathlineto{\pgfqpoint{1.556259in}{1.849669in}}%
\pgfpathlineto{\pgfqpoint{1.510951in}{1.835818in}}%
\pgfpathlineto{\pgfqpoint{1.465546in}{1.821939in}}%
\pgfpathlineto{\pgfqpoint{1.436250in}{1.841752in}}%
\pgfpathlineto{\pgfqpoint{1.407042in}{1.861506in}}%
\pgfpathclose%
\pgfusepath{stroke,fill}%
\end{pgfscope}%
\begin{pgfscope}%
\pgfpathrectangle{\pgfqpoint{0.050000in}{0.050000in}}{\pgfqpoint{4.900000in}{4.900000in}}%
\pgfusepath{clip}%
\pgfsetbuttcap%
\pgfsetroundjoin%
\definecolor{currentfill}{rgb}{0.250165,0.324183,0.699762}%
\pgfsetfillcolor{currentfill}%
\pgfsetfillopacity{0.900000}%
\pgfsetlinewidth{0.200750pt}%
\definecolor{currentstroke}{rgb}{0.200000,0.200000,0.200000}%
\pgfsetstrokecolor{currentstroke}%
\pgfsetdash{}{0pt}%
\pgfpathmoveto{\pgfqpoint{2.829199in}{3.378084in}}%
\pgfpathlineto{\pgfqpoint{2.873243in}{3.417484in}}%
\pgfpathlineto{\pgfqpoint{2.916723in}{3.372617in}}%
\pgfpathlineto{\pgfqpoint{2.947695in}{3.281427in}}%
\pgfpathlineto{\pgfqpoint{2.978711in}{3.194376in}}%
\pgfpathlineto{\pgfqpoint{2.935475in}{3.260730in}}%
\pgfpathlineto{\pgfqpoint{2.891601in}{3.251799in}}%
\pgfpathlineto{\pgfqpoint{2.860379in}{3.315302in}}%
\pgfpathlineto{\pgfqpoint{2.829199in}{3.378084in}}%
\pgfpathclose%
\pgfusepath{stroke,fill}%
\end{pgfscope}%
\begin{pgfscope}%
\pgfpathrectangle{\pgfqpoint{0.050000in}{0.050000in}}{\pgfqpoint{4.900000in}{4.900000in}}%
\pgfusepath{clip}%
\pgfsetbuttcap%
\pgfsetroundjoin%
\definecolor{currentfill}{rgb}{1.000000,1.000000,1.000000}%
\pgfsetfillcolor{currentfill}%
\pgfsetfillopacity{0.900000}%
\pgfsetlinewidth{0.200750pt}%
\definecolor{currentstroke}{rgb}{0.200000,0.200000,0.200000}%
\pgfsetstrokecolor{currentstroke}%
\pgfsetdash{}{0pt}%
\pgfpathmoveto{\pgfqpoint{1.945367in}{1.853645in}}%
\pgfpathlineto{\pgfqpoint{1.990155in}{1.867460in}}%
\pgfpathlineto{\pgfqpoint{2.034848in}{1.881256in}}%
\pgfpathlineto{\pgfqpoint{2.064960in}{1.861578in}}%
\pgfpathlineto{\pgfqpoint{2.095163in}{1.841851in}}%
\pgfpathlineto{\pgfqpoint{2.050399in}{1.827940in}}%
\pgfpathlineto{\pgfqpoint{2.005540in}{1.814032in}}%
\pgfpathlineto{\pgfqpoint{1.975408in}{1.833868in}}%
\pgfpathlineto{\pgfqpoint{1.945367in}{1.853645in}}%
\pgfpathclose%
\pgfusepath{stroke,fill}%
\end{pgfscope}%
\begin{pgfscope}%
\pgfpathrectangle{\pgfqpoint{0.050000in}{0.050000in}}{\pgfqpoint{4.900000in}{4.900000in}}%
\pgfusepath{clip}%
\pgfsetbuttcap%
\pgfsetroundjoin%
\definecolor{currentfill}{rgb}{0.982099,0.981361,0.989296}%
\pgfsetfillcolor{currentfill}%
\pgfsetfillopacity{0.900000}%
\pgfsetlinewidth{0.200750pt}%
\definecolor{currentstroke}{rgb}{0.200000,0.200000,0.200000}%
\pgfsetstrokecolor{currentstroke}%
\pgfsetdash{}{0pt}%
\pgfpathmoveto{\pgfqpoint{2.334208in}{1.865195in}}%
\pgfpathlineto{\pgfqpoint{2.378549in}{1.887356in}}%
\pgfpathlineto{\pgfqpoint{2.422792in}{1.916181in}}%
\pgfpathlineto{\pgfqpoint{2.453489in}{1.900918in}}%
\pgfpathlineto{\pgfqpoint{2.484284in}{1.885727in}}%
\pgfpathlineto{\pgfqpoint{2.439955in}{1.853757in}}%
\pgfpathlineto{\pgfqpoint{2.395538in}{1.829176in}}%
\pgfpathlineto{\pgfqpoint{2.364826in}{1.847136in}}%
\pgfpathlineto{\pgfqpoint{2.334208in}{1.865195in}}%
\pgfpathclose%
\pgfusepath{stroke,fill}%
\end{pgfscope}%
\begin{pgfscope}%
\pgfpathrectangle{\pgfqpoint{0.050000in}{0.050000in}}{\pgfqpoint{4.900000in}{4.900000in}}%
\pgfusepath{clip}%
\pgfsetbuttcap%
\pgfsetroundjoin%
\definecolor{currentfill}{rgb}{1.000000,1.000000,1.000000}%
\pgfsetfillcolor{currentfill}%
\pgfsetfillopacity{0.900000}%
\pgfsetlinewidth{0.200750pt}%
\definecolor{currentstroke}{rgb}{0.200000,0.200000,0.200000}%
\pgfsetstrokecolor{currentstroke}%
\pgfsetdash{}{0pt}%
\pgfpathmoveto{\pgfqpoint{1.556259in}{1.849669in}}%
\pgfpathlineto{\pgfqpoint{1.601472in}{1.863490in}}%
\pgfpathlineto{\pgfqpoint{1.646589in}{1.877282in}}%
\pgfpathlineto{\pgfqpoint{1.676124in}{1.857576in}}%
\pgfpathlineto{\pgfqpoint{1.705748in}{1.837811in}}%
\pgfpathlineto{\pgfqpoint{1.660556in}{1.823935in}}%
\pgfpathlineto{\pgfqpoint{1.615269in}{1.810030in}}%
\pgfpathlineto{\pgfqpoint{1.585720in}{1.829879in}}%
\pgfpathlineto{\pgfqpoint{1.556259in}{1.849669in}}%
\pgfpathclose%
\pgfusepath{stroke,fill}%
\end{pgfscope}%
\begin{pgfscope}%
\pgfpathrectangle{\pgfqpoint{0.050000in}{0.050000in}}{\pgfqpoint{4.900000in}{4.900000in}}%
\pgfusepath{clip}%
\pgfsetbuttcap%
\pgfsetroundjoin%
\definecolor{currentfill}{rgb}{0.994033,0.993787,0.996432}%
\pgfsetfillcolor{currentfill}%
\pgfsetfillopacity{0.900000}%
\pgfsetlinewidth{0.200750pt}%
\definecolor{currentstroke}{rgb}{0.200000,0.200000,0.200000}%
\pgfsetstrokecolor{currentstroke}%
\pgfsetdash{}{0pt}%
\pgfpathmoveto{\pgfqpoint{3.651024in}{1.887006in}}%
\pgfpathlineto{\pgfqpoint{3.693438in}{1.871977in}}%
\pgfpathlineto{\pgfqpoint{3.726073in}{1.851884in}}%
\pgfpathlineto{\pgfqpoint{3.758812in}{1.831997in}}%
\pgfpathlineto{\pgfqpoint{3.716139in}{1.835606in}}%
\pgfpathlineto{\pgfqpoint{3.683532in}{1.861161in}}%
\pgfpathlineto{\pgfqpoint{3.651024in}{1.887006in}}%
\pgfpathclose%
\pgfusepath{stroke,fill}%
\end{pgfscope}%
\begin{pgfscope}%
\pgfpathrectangle{\pgfqpoint{0.050000in}{0.050000in}}{\pgfqpoint{4.900000in}{4.900000in}}%
\pgfusepath{clip}%
\pgfsetbuttcap%
\pgfsetroundjoin%
\definecolor{currentfill}{rgb}{0.582314,0.565090,0.750250}%
\pgfsetfillcolor{currentfill}%
\pgfsetfillopacity{0.900000}%
\pgfsetlinewidth{0.200750pt}%
\definecolor{currentstroke}{rgb}{0.200000,0.200000,0.200000}%
\pgfsetstrokecolor{currentstroke}%
\pgfsetdash{}{0pt}%
\pgfpathmoveto{\pgfqpoint{3.228975in}{2.621079in}}%
\pgfpathlineto{\pgfqpoint{3.270885in}{2.499296in}}%
\pgfpathlineto{\pgfqpoint{3.312282in}{2.361074in}}%
\pgfpathlineto{\pgfqpoint{3.343864in}{2.312194in}}%
\pgfpathlineto{\pgfqpoint{3.375537in}{2.265171in}}%
\pgfpathlineto{\pgfqpoint{3.334111in}{2.389795in}}%
\pgfpathlineto{\pgfqpoint{3.292248in}{2.503863in}}%
\pgfpathlineto{\pgfqpoint{3.260572in}{2.561451in}}%
\pgfpathlineto{\pgfqpoint{3.228975in}{2.621079in}}%
\pgfpathclose%
\pgfusepath{stroke,fill}%
\end{pgfscope}%
\begin{pgfscope}%
\pgfpathrectangle{\pgfqpoint{0.050000in}{0.050000in}}{\pgfqpoint{4.900000in}{4.900000in}}%
\pgfusepath{clip}%
\pgfsetbuttcap%
\pgfsetroundjoin%
\definecolor{currentfill}{rgb}{1.000000,1.000000,1.000000}%
\pgfsetfillcolor{currentfill}%
\pgfsetfillopacity{0.900000}%
\pgfsetlinewidth{0.200750pt}%
\definecolor{currentstroke}{rgb}{0.200000,0.200000,0.200000}%
\pgfsetstrokecolor{currentstroke}%
\pgfsetdash{}{0pt}%
\pgfpathmoveto{\pgfqpoint{1.166914in}{1.845691in}}%
\pgfpathlineto{\pgfqpoint{1.212553in}{1.859521in}}%
\pgfpathlineto{\pgfqpoint{1.258094in}{1.873321in}}%
\pgfpathlineto{\pgfqpoint{1.287051in}{1.853603in}}%
\pgfpathlineto{\pgfqpoint{1.316095in}{1.833826in}}%
\pgfpathlineto{\pgfqpoint{1.270476in}{1.819942in}}%
\pgfpathlineto{\pgfqpoint{1.224761in}{1.806028in}}%
\pgfpathlineto{\pgfqpoint{1.195794in}{1.825890in}}%
\pgfpathlineto{\pgfqpoint{1.166914in}{1.845691in}}%
\pgfpathclose%
\pgfusepath{stroke,fill}%
\end{pgfscope}%
\begin{pgfscope}%
\pgfpathrectangle{\pgfqpoint{0.050000in}{0.050000in}}{\pgfqpoint{4.900000in}{4.900000in}}%
\pgfusepath{clip}%
\pgfsetbuttcap%
\pgfsetroundjoin%
\definecolor{currentfill}{rgb}{0.247459,0.295425,0.661530}%
\pgfsetfillcolor{currentfill}%
\pgfsetfillopacity{0.900000}%
\pgfsetlinewidth{0.200750pt}%
\definecolor{currentstroke}{rgb}{0.200000,0.200000,0.200000}%
\pgfsetstrokecolor{currentstroke}%
\pgfsetdash{}{0pt}%
\pgfpathmoveto{\pgfqpoint{2.740415in}{3.121170in}}%
\pgfpathlineto{\pgfqpoint{2.784850in}{3.271959in}}%
\pgfpathlineto{\pgfqpoint{2.829199in}{3.378084in}}%
\pgfpathlineto{\pgfqpoint{2.860379in}{3.315302in}}%
\pgfpathlineto{\pgfqpoint{2.891601in}{3.251799in}}%
\pgfpathlineto{\pgfqpoint{2.847309in}{3.177671in}}%
\pgfpathlineto{\pgfqpoint{2.802818in}{3.055041in}}%
\pgfpathlineto{\pgfqpoint{2.771582in}{3.089324in}}%
\pgfpathlineto{\pgfqpoint{2.740415in}{3.121170in}}%
\pgfpathclose%
\pgfusepath{stroke,fill}%
\end{pgfscope}%
\begin{pgfscope}%
\pgfpathrectangle{\pgfqpoint{0.050000in}{0.050000in}}{\pgfqpoint{4.900000in}{4.900000in}}%
\pgfusepath{clip}%
\pgfsetbuttcap%
\pgfsetroundjoin%
\definecolor{currentfill}{rgb}{0.904529,0.900592,0.942914}%
\pgfsetfillcolor{currentfill}%
\pgfsetfillopacity{0.900000}%
\pgfsetlinewidth{0.200750pt}%
\definecolor{currentstroke}{rgb}{0.200000,0.200000,0.200000}%
\pgfsetstrokecolor{currentstroke}%
\pgfsetdash{}{0pt}%
\pgfpathmoveto{\pgfqpoint{3.503143in}{2.092380in}}%
\pgfpathlineto{\pgfqpoint{3.544622in}{2.003659in}}%
\pgfpathlineto{\pgfqpoint{3.586303in}{1.939694in}}%
\pgfpathlineto{\pgfqpoint{3.618615in}{1.913173in}}%
\pgfpathlineto{\pgfqpoint{3.651024in}{1.887006in}}%
\pgfpathlineto{\pgfqpoint{3.609134in}{1.937528in}}%
\pgfpathlineto{\pgfqpoint{3.567506in}{2.013387in}}%
\pgfpathlineto{\pgfqpoint{3.535277in}{2.052355in}}%
\pgfpathlineto{\pgfqpoint{3.503143in}{2.092380in}}%
\pgfpathclose%
\pgfusepath{stroke,fill}%
\end{pgfscope}%
\begin{pgfscope}%
\pgfpathrectangle{\pgfqpoint{0.050000in}{0.050000in}}{\pgfqpoint{4.900000in}{4.900000in}}%
\pgfusepath{clip}%
\pgfsetbuttcap%
\pgfsetroundjoin%
\definecolor{currentfill}{rgb}{1.000000,1.000000,1.000000}%
\pgfsetfillcolor{currentfill}%
\pgfsetfillopacity{0.900000}%
\pgfsetlinewidth{0.200750pt}%
\definecolor{currentstroke}{rgb}{0.200000,0.200000,0.200000}%
\pgfsetstrokecolor{currentstroke}%
\pgfsetdash{}{0pt}%
\pgfpathmoveto{\pgfqpoint{2.095163in}{1.841851in}}%
\pgfpathlineto{\pgfqpoint{2.139831in}{1.855837in}}%
\pgfpathlineto{\pgfqpoint{2.184404in}{1.870091in}}%
\pgfpathlineto{\pgfqpoint{2.214768in}{1.850660in}}%
\pgfpathlineto{\pgfqpoint{2.245224in}{1.831244in}}%
\pgfpathlineto{\pgfqpoint{2.200582in}{1.816509in}}%
\pgfpathlineto{\pgfqpoint{2.155844in}{1.802264in}}%
\pgfpathlineto{\pgfqpoint{2.125457in}{1.822078in}}%
\pgfpathlineto{\pgfqpoint{2.095163in}{1.841851in}}%
\pgfpathclose%
\pgfusepath{stroke,fill}%
\end{pgfscope}%
\begin{pgfscope}%
\pgfpathrectangle{\pgfqpoint{0.050000in}{0.050000in}}{\pgfqpoint{4.900000in}{4.900000in}}%
\pgfusepath{clip}%
\pgfsetbuttcap%
\pgfsetroundjoin%
\definecolor{currentfill}{rgb}{0.246228,0.282353,0.644152}%
\pgfsetfillcolor{currentfill}%
\pgfsetfillopacity{0.900000}%
\pgfsetlinewidth{0.200750pt}%
\definecolor{currentstroke}{rgb}{0.200000,0.200000,0.200000}%
\pgfsetstrokecolor{currentstroke}%
\pgfsetdash{}{0pt}%
\pgfpathmoveto{\pgfqpoint{2.891601in}{3.251799in}}%
\pgfpathlineto{\pgfqpoint{2.935475in}{3.260730in}}%
\pgfpathlineto{\pgfqpoint{2.978711in}{3.194376in}}%
\pgfpathlineto{\pgfqpoint{3.009777in}{3.111289in}}%
\pgfpathlineto{\pgfqpoint{3.040898in}{3.031945in}}%
\pgfpathlineto{\pgfqpoint{2.997878in}{3.111186in}}%
\pgfpathlineto{\pgfqpoint{2.954181in}{3.124722in}}%
\pgfpathlineto{\pgfqpoint{2.922868in}{3.188135in}}%
\pgfpathlineto{\pgfqpoint{2.891601in}{3.251799in}}%
\pgfpathclose%
\pgfusepath{stroke,fill}%
\end{pgfscope}%
\begin{pgfscope}%
\pgfpathrectangle{\pgfqpoint{0.050000in}{0.050000in}}{\pgfqpoint{4.900000in}{4.900000in}}%
\pgfusepath{clip}%
\pgfsetbuttcap%
\pgfsetroundjoin%
\definecolor{currentfill}{rgb}{0.272034,0.242015,0.564721}%
\pgfsetfillcolor{currentfill}%
\pgfsetfillopacity{0.900000}%
\pgfsetlinewidth{0.200750pt}%
\definecolor{currentstroke}{rgb}{0.200000,0.200000,0.200000}%
\pgfsetstrokecolor{currentstroke}%
\pgfsetdash{}{0pt}%
\pgfpathmoveto{\pgfqpoint{2.713877in}{2.737843in}}%
\pgfpathlineto{\pgfqpoint{2.758303in}{2.902642in}}%
\pgfpathlineto{\pgfqpoint{2.802818in}{3.055041in}}%
\pgfpathlineto{\pgfqpoint{2.834122in}{3.018652in}}%
\pgfpathlineto{\pgfqpoint{2.865491in}{2.980484in}}%
\pgfpathlineto{\pgfqpoint{2.820883in}{2.850058in}}%
\pgfpathlineto{\pgfqpoint{2.776298in}{2.701891in}}%
\pgfpathlineto{\pgfqpoint{2.745043in}{2.720732in}}%
\pgfpathlineto{\pgfqpoint{2.713877in}{2.737843in}}%
\pgfpathclose%
\pgfusepath{stroke,fill}%
\end{pgfscope}%
\begin{pgfscope}%
\pgfpathrectangle{\pgfqpoint{0.050000in}{0.050000in}}{\pgfqpoint{4.900000in}{4.900000in}}%
\pgfusepath{clip}%
\pgfsetbuttcap%
\pgfsetroundjoin%
\definecolor{currentfill}{rgb}{1.000000,1.000000,1.000000}%
\pgfsetfillcolor{currentfill}%
\pgfsetfillopacity{0.900000}%
\pgfsetlinewidth{0.200750pt}%
\definecolor{currentstroke}{rgb}{0.200000,0.200000,0.200000}%
\pgfsetstrokecolor{currentstroke}%
\pgfsetdash{}{0pt}%
\pgfpathmoveto{\pgfqpoint{1.705748in}{1.837811in}}%
\pgfpathlineto{\pgfqpoint{1.750843in}{1.851657in}}%
\pgfpathlineto{\pgfqpoint{1.795843in}{1.865474in}}%
\pgfpathlineto{\pgfqpoint{1.825630in}{1.845732in}}%
\pgfpathlineto{\pgfqpoint{1.855507in}{1.825931in}}%
\pgfpathlineto{\pgfqpoint{1.810434in}{1.812030in}}%
\pgfpathlineto{\pgfqpoint{1.765265in}{1.798100in}}%
\pgfpathlineto{\pgfqpoint{1.735461in}{1.817985in}}%
\pgfpathlineto{\pgfqpoint{1.705748in}{1.837811in}}%
\pgfpathclose%
\pgfusepath{stroke,fill}%
\end{pgfscope}%
\begin{pgfscope}%
\pgfpathrectangle{\pgfqpoint{0.050000in}{0.050000in}}{\pgfqpoint{4.900000in}{4.900000in}}%
\pgfusepath{clip}%
\pgfsetbuttcap%
\pgfsetroundjoin%
\definecolor{currentfill}{rgb}{0.504744,0.484321,0.703868}%
\pgfsetfillcolor{currentfill}%
\pgfsetfillopacity{0.900000}%
\pgfsetlinewidth{0.200750pt}%
\definecolor{currentstroke}{rgb}{0.200000,0.200000,0.200000}%
\pgfsetstrokecolor{currentstroke}%
\pgfsetdash{}{0pt}%
\pgfpathmoveto{\pgfqpoint{2.687474in}{2.405724in}}%
\pgfpathlineto{\pgfqpoint{2.731817in}{2.550097in}}%
\pgfpathlineto{\pgfqpoint{2.776298in}{2.701891in}}%
\pgfpathlineto{\pgfqpoint{2.807640in}{2.681377in}}%
\pgfpathlineto{\pgfqpoint{2.839068in}{2.659274in}}%
\pgfpathlineto{\pgfqpoint{2.794421in}{2.519749in}}%
\pgfpathlineto{\pgfqpoint{2.749886in}{2.383067in}}%
\pgfpathlineto{\pgfqpoint{2.718630in}{2.394959in}}%
\pgfpathlineto{\pgfqpoint{2.687474in}{2.405724in}}%
\pgfpathclose%
\pgfusepath{stroke,fill}%
\end{pgfscope}%
\begin{pgfscope}%
\pgfpathrectangle{\pgfqpoint{0.050000in}{0.050000in}}{\pgfqpoint{4.900000in}{4.900000in}}%
\pgfusepath{clip}%
\pgfsetbuttcap%
\pgfsetroundjoin%
\definecolor{currentfill}{rgb}{1.000000,1.000000,1.000000}%
\pgfsetfillcolor{currentfill}%
\pgfsetfillopacity{0.900000}%
\pgfsetlinewidth{0.200750pt}%
\definecolor{currentstroke}{rgb}{0.200000,0.200000,0.200000}%
\pgfsetstrokecolor{currentstroke}%
\pgfsetdash{}{0pt}%
\pgfpathmoveto{\pgfqpoint{1.316095in}{1.833826in}}%
\pgfpathlineto{\pgfqpoint{1.361616in}{1.847680in}}%
\pgfpathlineto{\pgfqpoint{1.407042in}{1.861506in}}%
\pgfpathlineto{\pgfqpoint{1.436250in}{1.841752in}}%
\pgfpathlineto{\pgfqpoint{1.465546in}{1.821939in}}%
\pgfpathlineto{\pgfqpoint{1.420045in}{1.808029in}}%
\pgfpathlineto{\pgfqpoint{1.374447in}{1.794090in}}%
\pgfpathlineto{\pgfqpoint{1.345227in}{1.813988in}}%
\pgfpathlineto{\pgfqpoint{1.316095in}{1.833826in}}%
\pgfpathclose%
\pgfusepath{stroke,fill}%
\end{pgfscope}%
\begin{pgfscope}%
\pgfpathrectangle{\pgfqpoint{0.050000in}{0.050000in}}{\pgfqpoint{4.900000in}{4.900000in}}%
\pgfusepath{clip}%
\pgfsetbuttcap%
\pgfsetroundjoin%
\definecolor{currentfill}{rgb}{0.844860,0.838462,0.907236}%
\pgfsetfillcolor{currentfill}%
\pgfsetfillopacity{0.900000}%
\pgfsetlinewidth{0.200750pt}%
\definecolor{currentstroke}{rgb}{0.200000,0.200000,0.200000}%
\pgfsetstrokecolor{currentstroke}%
\pgfsetdash{}{0pt}%
\pgfpathmoveto{\pgfqpoint{2.572750in}{1.984302in}}%
\pgfpathlineto{\pgfqpoint{2.616947in}{2.056977in}}%
\pgfpathlineto{\pgfqpoint{2.661175in}{2.147995in}}%
\pgfpathlineto{\pgfqpoint{2.692344in}{2.137922in}}%
\pgfpathlineto{\pgfqpoint{2.723616in}{2.127222in}}%
\pgfpathlineto{\pgfqpoint{2.679219in}{2.035436in}}%
\pgfpathlineto{\pgfqpoint{2.634872in}{1.960440in}}%
\pgfpathlineto{\pgfqpoint{2.603760in}{1.972477in}}%
\pgfpathlineto{\pgfqpoint{2.572750in}{1.984302in}}%
\pgfpathclose%
\pgfusepath{stroke,fill}%
\end{pgfscope}%
\begin{pgfscope}%
\pgfpathrectangle{\pgfqpoint{0.050000in}{0.050000in}}{\pgfqpoint{4.900000in}{4.900000in}}%
\pgfusepath{clip}%
\pgfsetbuttcap%
\pgfsetroundjoin%
\definecolor{currentfill}{rgb}{1.000000,1.000000,1.000000}%
\pgfsetfillcolor{currentfill}%
\pgfsetfillopacity{0.900000}%
\pgfsetlinewidth{0.200750pt}%
\definecolor{currentstroke}{rgb}{0.200000,0.200000,0.200000}%
\pgfsetstrokecolor{currentstroke}%
\pgfsetdash{}{0pt}%
\pgfpathmoveto{\pgfqpoint{2.245224in}{1.831244in}}%
\pgfpathlineto{\pgfqpoint{2.289766in}{1.847097in}}%
\pgfpathlineto{\pgfqpoint{2.334208in}{1.865195in}}%
\pgfpathlineto{\pgfqpoint{2.364826in}{1.847136in}}%
\pgfpathlineto{\pgfqpoint{2.395538in}{1.829176in}}%
\pgfpathlineto{\pgfqpoint{2.351025in}{1.809396in}}%
\pgfpathlineto{\pgfqpoint{2.306413in}{1.792478in}}%
\pgfpathlineto{\pgfqpoint{2.275772in}{1.811850in}}%
\pgfpathlineto{\pgfqpoint{2.245224in}{1.831244in}}%
\pgfpathclose%
\pgfusepath{stroke,fill}%
\end{pgfscope}%
\begin{pgfscope}%
\pgfpathrectangle{\pgfqpoint{0.050000in}{0.050000in}}{\pgfqpoint{4.900000in}{4.900000in}}%
\pgfusepath{clip}%
\pgfsetbuttcap%
\pgfsetroundjoin%
\definecolor{currentfill}{rgb}{0.934364,0.931657,0.960754}%
\pgfsetfillcolor{currentfill}%
\pgfsetfillopacity{0.900000}%
\pgfsetlinewidth{0.200750pt}%
\definecolor{currentstroke}{rgb}{0.200000,0.200000,0.200000}%
\pgfsetstrokecolor{currentstroke}%
\pgfsetdash{}{0pt}%
\pgfpathmoveto{\pgfqpoint{2.484284in}{1.885727in}}%
\pgfpathlineto{\pgfqpoint{2.528540in}{1.928166in}}%
\pgfpathlineto{\pgfqpoint{2.572750in}{1.984302in}}%
\pgfpathlineto{\pgfqpoint{2.603760in}{1.972477in}}%
\pgfpathlineto{\pgfqpoint{2.634872in}{1.960440in}}%
\pgfpathlineto{\pgfqpoint{2.590534in}{1.901186in}}%
\pgfpathlineto{\pgfqpoint{2.546168in}{1.855468in}}%
\pgfpathlineto{\pgfqpoint{2.515176in}{1.870584in}}%
\pgfpathlineto{\pgfqpoint{2.484284in}{1.885727in}}%
\pgfpathclose%
\pgfusepath{stroke,fill}%
\end{pgfscope}%
\begin{pgfscope}%
\pgfpathrectangle{\pgfqpoint{0.050000in}{0.050000in}}{\pgfqpoint{4.900000in}{4.900000in}}%
\pgfusepath{clip}%
\pgfsetbuttcap%
\pgfsetroundjoin%
\definecolor{currentfill}{rgb}{1.000000,1.000000,1.000000}%
\pgfsetfillcolor{currentfill}%
\pgfsetfillopacity{0.900000}%
\pgfsetlinewidth{0.200750pt}%
\definecolor{currentstroke}{rgb}{0.200000,0.200000,0.200000}%
\pgfsetstrokecolor{currentstroke}%
\pgfsetdash{}{0pt}%
\pgfpathmoveto{\pgfqpoint{1.855507in}{1.825931in}}%
\pgfpathlineto{\pgfqpoint{1.900485in}{1.839802in}}%
\pgfpathlineto{\pgfqpoint{1.945367in}{1.853645in}}%
\pgfpathlineto{\pgfqpoint{1.975408in}{1.833868in}}%
\pgfpathlineto{\pgfqpoint{2.005540in}{1.814032in}}%
\pgfpathlineto{\pgfqpoint{1.960585in}{1.800104in}}%
\pgfpathlineto{\pgfqpoint{1.915534in}{1.786148in}}%
\pgfpathlineto{\pgfqpoint{1.885475in}{1.806069in}}%
\pgfpathlineto{\pgfqpoint{1.855507in}{1.825931in}}%
\pgfpathclose%
\pgfusepath{stroke,fill}%
\end{pgfscope}%
\begin{pgfscope}%
\pgfpathrectangle{\pgfqpoint{0.050000in}{0.050000in}}{\pgfqpoint{4.900000in}{4.900000in}}%
\pgfusepath{clip}%
\pgfsetbuttcap%
\pgfsetroundjoin%
\definecolor{currentfill}{rgb}{0.242537,0.243137,0.592018}%
\pgfsetfillcolor{currentfill}%
\pgfsetfillopacity{0.900000}%
\pgfsetlinewidth{0.200750pt}%
\definecolor{currentstroke}{rgb}{0.200000,0.200000,0.200000}%
\pgfsetstrokecolor{currentstroke}%
\pgfsetdash{}{0pt}%
\pgfpathmoveto{\pgfqpoint{2.954181in}{3.124722in}}%
\pgfpathlineto{\pgfqpoint{2.997878in}{3.111186in}}%
\pgfpathlineto{\pgfqpoint{3.040898in}{3.031945in}}%
\pgfpathlineto{\pgfqpoint{3.072077in}{2.956097in}}%
\pgfpathlineto{\pgfqpoint{3.103320in}{2.883502in}}%
\pgfpathlineto{\pgfqpoint{3.060485in}{2.969919in}}%
\pgfpathlineto{\pgfqpoint{3.016961in}{2.999756in}}%
\pgfpathlineto{\pgfqpoint{2.985545in}{3.061858in}}%
\pgfpathlineto{\pgfqpoint{2.954181in}{3.124722in}}%
\pgfpathclose%
\pgfusepath{stroke,fill}%
\end{pgfscope}%
\begin{pgfscope}%
\pgfpathrectangle{\pgfqpoint{0.050000in}{0.050000in}}{\pgfqpoint{4.900000in}{4.900000in}}%
\pgfusepath{clip}%
\pgfsetbuttcap%
\pgfsetroundjoin%
\definecolor{currentfill}{rgb}{0.636017,0.621007,0.782361}%
\pgfsetfillcolor{currentfill}%
\pgfsetfillopacity{0.900000}%
\pgfsetlinewidth{0.200750pt}%
\definecolor{currentstroke}{rgb}{0.200000,0.200000,0.200000}%
\pgfsetstrokecolor{currentstroke}%
\pgfsetdash{}{0pt}%
\pgfpathmoveto{\pgfqpoint{3.292248in}{2.503863in}}%
\pgfpathlineto{\pgfqpoint{3.334111in}{2.389795in}}%
\pgfpathlineto{\pgfqpoint{3.375537in}{2.265171in}}%
\pgfpathlineto{\pgfqpoint{3.407300in}{2.219825in}}%
\pgfpathlineto{\pgfqpoint{3.439155in}{2.176000in}}%
\pgfpathlineto{\pgfqpoint{3.397679in}{2.288000in}}%
\pgfpathlineto{\pgfqpoint{3.355839in}{2.394224in}}%
\pgfpathlineto{\pgfqpoint{3.324003in}{2.448167in}}%
\pgfpathlineto{\pgfqpoint{3.292248in}{2.503863in}}%
\pgfpathclose%
\pgfusepath{stroke,fill}%
\end{pgfscope}%
\begin{pgfscope}%
\pgfpathrectangle{\pgfqpoint{0.050000in}{0.050000in}}{\pgfqpoint{4.900000in}{4.900000in}}%
\pgfusepath{clip}%
\pgfsetbuttcap%
\pgfsetroundjoin%
\definecolor{currentfill}{rgb}{1.000000,1.000000,1.000000}%
\pgfsetfillcolor{currentfill}%
\pgfsetfillopacity{0.900000}%
\pgfsetlinewidth{0.200750pt}%
\definecolor{currentstroke}{rgb}{0.200000,0.200000,0.200000}%
\pgfsetstrokecolor{currentstroke}%
\pgfsetdash{}{0pt}%
\pgfpathmoveto{\pgfqpoint{1.465546in}{1.821939in}}%
\pgfpathlineto{\pgfqpoint{1.510951in}{1.835818in}}%
\pgfpathlineto{\pgfqpoint{1.556259in}{1.849669in}}%
\pgfpathlineto{\pgfqpoint{1.585720in}{1.829879in}}%
\pgfpathlineto{\pgfqpoint{1.615269in}{1.810030in}}%
\pgfpathlineto{\pgfqpoint{1.569886in}{1.796095in}}%
\pgfpathlineto{\pgfqpoint{1.524406in}{1.782131in}}%
\pgfpathlineto{\pgfqpoint{1.494931in}{1.802065in}}%
\pgfpathlineto{\pgfqpoint{1.465546in}{1.821939in}}%
\pgfpathclose%
\pgfusepath{stroke,fill}%
\end{pgfscope}%
\begin{pgfscope}%
\pgfpathrectangle{\pgfqpoint{0.050000in}{0.050000in}}{\pgfqpoint{4.900000in}{4.900000in}}%
\pgfusepath{clip}%
\pgfsetbuttcap%
\pgfsetroundjoin%
\definecolor{currentfill}{rgb}{0.245244,0.271895,0.630250}%
\pgfsetfillcolor{currentfill}%
\pgfsetfillopacity{0.900000}%
\pgfsetlinewidth{0.200750pt}%
\definecolor{currentstroke}{rgb}{0.200000,0.200000,0.200000}%
\pgfsetstrokecolor{currentstroke}%
\pgfsetdash{}{0pt}%
\pgfpathmoveto{\pgfqpoint{2.802818in}{3.055041in}}%
\pgfpathlineto{\pgfqpoint{2.847309in}{3.177671in}}%
\pgfpathlineto{\pgfqpoint{2.891601in}{3.251799in}}%
\pgfpathlineto{\pgfqpoint{2.922868in}{3.188135in}}%
\pgfpathlineto{\pgfqpoint{2.954181in}{3.124722in}}%
\pgfpathlineto{\pgfqpoint{2.909985in}{3.077209in}}%
\pgfpathlineto{\pgfqpoint{2.865491in}{2.980484in}}%
\pgfpathlineto{\pgfqpoint{2.834122in}{3.018652in}}%
\pgfpathlineto{\pgfqpoint{2.802818in}{3.055041in}}%
\pgfpathclose%
\pgfusepath{stroke,fill}%
\end{pgfscope}%
\begin{pgfscope}%
\pgfpathrectangle{\pgfqpoint{0.050000in}{0.050000in}}{\pgfqpoint{4.900000in}{4.900000in}}%
\pgfusepath{clip}%
\pgfsetbuttcap%
\pgfsetroundjoin%
\definecolor{currentfill}{rgb}{0.695686,0.683137,0.818039}%
\pgfsetfillcolor{currentfill}%
\pgfsetfillopacity{0.900000}%
\pgfsetlinewidth{0.200750pt}%
\definecolor{currentstroke}{rgb}{0.200000,0.200000,0.200000}%
\pgfsetstrokecolor{currentstroke}%
\pgfsetdash{}{0pt}%
\pgfpathmoveto{\pgfqpoint{2.661175in}{2.147995in}}%
\pgfpathlineto{\pgfqpoint{2.705476in}{2.257445in}}%
\pgfpathlineto{\pgfqpoint{2.749886in}{2.383067in}}%
\pgfpathlineto{\pgfqpoint{2.781240in}{2.370041in}}%
\pgfpathlineto{\pgfqpoint{2.812692in}{2.355889in}}%
\pgfpathlineto{\pgfqpoint{2.768099in}{2.235061in}}%
\pgfpathlineto{\pgfqpoint{2.723616in}{2.127222in}}%
\pgfpathlineto{\pgfqpoint{2.692344in}{2.137922in}}%
\pgfpathlineto{\pgfqpoint{2.661175in}{2.147995in}}%
\pgfpathclose%
\pgfusepath{stroke,fill}%
\end{pgfscope}%
\begin{pgfscope}%
\pgfpathrectangle{\pgfqpoint{0.050000in}{0.050000in}}{\pgfqpoint{4.900000in}{4.900000in}}%
\pgfusepath{clip}%
\pgfsetbuttcap%
\pgfsetroundjoin%
\definecolor{currentfill}{rgb}{1.000000,1.000000,1.000000}%
\pgfsetfillcolor{currentfill}%
\pgfsetfillopacity{0.900000}%
\pgfsetlinewidth{0.200750pt}%
\definecolor{currentstroke}{rgb}{0.200000,0.200000,0.200000}%
\pgfsetstrokecolor{currentstroke}%
\pgfsetdash{}{0pt}%
\pgfpathmoveto{\pgfqpoint{1.075347in}{1.817944in}}%
\pgfpathlineto{\pgfqpoint{1.121179in}{1.831832in}}%
\pgfpathlineto{\pgfqpoint{1.166914in}{1.845691in}}%
\pgfpathlineto{\pgfqpoint{1.195794in}{1.825890in}}%
\pgfpathlineto{\pgfqpoint{1.224761in}{1.806028in}}%
\pgfpathlineto{\pgfqpoint{1.178948in}{1.792085in}}%
\pgfpathlineto{\pgfqpoint{1.133038in}{1.778112in}}%
\pgfpathlineto{\pgfqpoint{1.104149in}{1.798058in}}%
\pgfpathlineto{\pgfqpoint{1.075347in}{1.817944in}}%
\pgfpathclose%
\pgfusepath{stroke,fill}%
\end{pgfscope}%
\begin{pgfscope}%
\pgfpathrectangle{\pgfqpoint{0.050000in}{0.050000in}}{\pgfqpoint{4.900000in}{4.900000in}}%
\pgfusepath{clip}%
\pgfsetbuttcap%
\pgfsetroundjoin%
\definecolor{currentfill}{rgb}{1.000000,1.000000,1.000000}%
\pgfsetfillcolor{currentfill}%
\pgfsetfillopacity{0.900000}%
\pgfsetlinewidth{0.200750pt}%
\definecolor{currentstroke}{rgb}{0.200000,0.200000,0.200000}%
\pgfsetstrokecolor{currentstroke}%
\pgfsetdash{}{0pt}%
\pgfpathmoveto{\pgfqpoint{2.005540in}{1.814032in}}%
\pgfpathlineto{\pgfqpoint{2.050399in}{1.827940in}}%
\pgfpathlineto{\pgfqpoint{2.095163in}{1.841851in}}%
\pgfpathlineto{\pgfqpoint{2.125457in}{1.822078in}}%
\pgfpathlineto{\pgfqpoint{2.155844in}{1.802264in}}%
\pgfpathlineto{\pgfqpoint{2.111009in}{1.788201in}}%
\pgfpathlineto{\pgfqpoint{2.066078in}{1.774186in}}%
\pgfpathlineto{\pgfqpoint{2.035763in}{1.794138in}}%
\pgfpathlineto{\pgfqpoint{2.005540in}{1.814032in}}%
\pgfpathclose%
\pgfusepath{stroke,fill}%
\end{pgfscope}%
\begin{pgfscope}%
\pgfpathrectangle{\pgfqpoint{0.050000in}{0.050000in}}{\pgfqpoint{4.900000in}{4.900000in}}%
\pgfusepath{clip}%
\pgfsetbuttcap%
\pgfsetroundjoin%
\definecolor{currentfill}{rgb}{0.976132,0.975148,0.985729}%
\pgfsetfillcolor{currentfill}%
\pgfsetfillopacity{0.900000}%
\pgfsetlinewidth{0.200750pt}%
\definecolor{currentstroke}{rgb}{0.200000,0.200000,0.200000}%
\pgfsetstrokecolor{currentstroke}%
\pgfsetdash{}{0pt}%
\pgfpathmoveto{\pgfqpoint{2.395538in}{1.829176in}}%
\pgfpathlineto{\pgfqpoint{2.439955in}{1.853757in}}%
\pgfpathlineto{\pgfqpoint{2.484284in}{1.885727in}}%
\pgfpathlineto{\pgfqpoint{2.515176in}{1.870584in}}%
\pgfpathlineto{\pgfqpoint{2.546168in}{1.855468in}}%
\pgfpathlineto{\pgfqpoint{2.501746in}{1.820516in}}%
\pgfpathlineto{\pgfqpoint{2.457248in}{1.793509in}}%
\pgfpathlineto{\pgfqpoint{2.426345in}{1.811303in}}%
\pgfpathlineto{\pgfqpoint{2.395538in}{1.829176in}}%
\pgfpathclose%
\pgfusepath{stroke,fill}%
\end{pgfscope}%
\begin{pgfscope}%
\pgfpathrectangle{\pgfqpoint{0.050000in}{0.050000in}}{\pgfqpoint{4.900000in}{4.900000in}}%
\pgfusepath{clip}%
\pgfsetbuttcap%
\pgfsetroundjoin%
\definecolor{currentfill}{rgb}{1.000000,1.000000,1.000000}%
\pgfsetfillcolor{currentfill}%
\pgfsetfillopacity{0.900000}%
\pgfsetlinewidth{0.200750pt}%
\definecolor{currentstroke}{rgb}{0.200000,0.200000,0.200000}%
\pgfsetstrokecolor{currentstroke}%
\pgfsetdash{}{0pt}%
\pgfpathmoveto{\pgfqpoint{1.615269in}{1.810030in}}%
\pgfpathlineto{\pgfqpoint{1.660556in}{1.823935in}}%
\pgfpathlineto{\pgfqpoint{1.705748in}{1.837811in}}%
\pgfpathlineto{\pgfqpoint{1.735461in}{1.817985in}}%
\pgfpathlineto{\pgfqpoint{1.765265in}{1.798100in}}%
\pgfpathlineto{\pgfqpoint{1.720000in}{1.784139in}}%
\pgfpathlineto{\pgfqpoint{1.674639in}{1.770150in}}%
\pgfpathlineto{\pgfqpoint{1.644909in}{1.790120in}}%
\pgfpathlineto{\pgfqpoint{1.615269in}{1.810030in}}%
\pgfpathclose%
\pgfusepath{stroke,fill}%
\end{pgfscope}%
\begin{pgfscope}%
\pgfpathrectangle{\pgfqpoint{0.050000in}{0.050000in}}{\pgfqpoint{4.900000in}{4.900000in}}%
\pgfusepath{clip}%
\pgfsetbuttcap%
\pgfsetroundjoin%
\definecolor{currentfill}{rgb}{1.000000,1.000000,1.000000}%
\pgfsetfillcolor{currentfill}%
\pgfsetfillopacity{0.900000}%
\pgfsetlinewidth{0.200750pt}%
\definecolor{currentstroke}{rgb}{0.200000,0.200000,0.200000}%
\pgfsetstrokecolor{currentstroke}%
\pgfsetdash{}{0pt}%
\pgfpathmoveto{\pgfqpoint{3.716139in}{1.835606in}}%
\pgfpathlineto{\pgfqpoint{3.758812in}{1.831997in}}%
\pgfpathlineto{\pgfqpoint{3.791652in}{1.812154in}}%
\pgfpathlineto{\pgfqpoint{3.824593in}{1.792250in}}%
\pgfpathlineto{\pgfqpoint{3.781651in}{1.785260in}}%
\pgfpathlineto{\pgfqpoint{3.748845in}{1.810314in}}%
\pgfpathlineto{\pgfqpoint{3.716139in}{1.835606in}}%
\pgfpathclose%
\pgfusepath{stroke,fill}%
\end{pgfscope}%
\begin{pgfscope}%
\pgfpathrectangle{\pgfqpoint{0.050000in}{0.050000in}}{\pgfqpoint{4.900000in}{4.900000in}}%
\pgfusepath{clip}%
\pgfsetbuttcap%
\pgfsetroundjoin%
\definecolor{currentfill}{rgb}{0.922430,0.919231,0.953618}%
\pgfsetfillcolor{currentfill}%
\pgfsetfillopacity{0.900000}%
\pgfsetlinewidth{0.200750pt}%
\definecolor{currentstroke}{rgb}{0.200000,0.200000,0.200000}%
\pgfsetstrokecolor{currentstroke}%
\pgfsetdash{}{0pt}%
\pgfpathmoveto{\pgfqpoint{3.567506in}{2.013387in}}%
\pgfpathlineto{\pgfqpoint{3.609134in}{1.937528in}}%
\pgfpathlineto{\pgfqpoint{3.651024in}{1.887006in}}%
\pgfpathlineto{\pgfqpoint{3.683532in}{1.861161in}}%
\pgfpathlineto{\pgfqpoint{3.716139in}{1.835606in}}%
\pgfpathlineto{\pgfqpoint{3.674037in}{1.874100in}}%
\pgfpathlineto{\pgfqpoint{3.632251in}{1.938292in}}%
\pgfpathlineto{\pgfqpoint{3.599830in}{1.975392in}}%
\pgfpathlineto{\pgfqpoint{3.567506in}{2.013387in}}%
\pgfpathclose%
\pgfusepath{stroke,fill}%
\end{pgfscope}%
\begin{pgfscope}%
\pgfpathrectangle{\pgfqpoint{0.050000in}{0.050000in}}{\pgfqpoint{4.900000in}{4.900000in}}%
\pgfusepath{clip}%
\pgfsetbuttcap%
\pgfsetroundjoin%
\definecolor{currentfill}{rgb}{0.242199,0.210950,0.546882}%
\pgfsetfillcolor{currentfill}%
\pgfsetfillopacity{0.900000}%
\pgfsetlinewidth{0.200750pt}%
\definecolor{currentstroke}{rgb}{0.200000,0.200000,0.200000}%
\pgfsetstrokecolor{currentstroke}%
\pgfsetdash{}{0pt}%
\pgfpathmoveto{\pgfqpoint{3.016961in}{2.999756in}}%
\pgfpathlineto{\pgfqpoint{3.060485in}{2.969919in}}%
\pgfpathlineto{\pgfqpoint{3.103320in}{2.883502in}}%
\pgfpathlineto{\pgfqpoint{3.134628in}{2.813921in}}%
\pgfpathlineto{\pgfqpoint{3.166004in}{2.747127in}}%
\pgfpathlineto{\pgfqpoint{3.123326in}{2.836895in}}%
\pgfpathlineto{\pgfqpoint{3.079964in}{2.878377in}}%
\pgfpathlineto{\pgfqpoint{3.048433in}{2.938563in}}%
\pgfpathlineto{\pgfqpoint{3.016961in}{2.999756in}}%
\pgfpathclose%
\pgfusepath{stroke,fill}%
\end{pgfscope}%
\begin{pgfscope}%
\pgfpathrectangle{\pgfqpoint{0.050000in}{0.050000in}}{\pgfqpoint{4.900000in}{4.900000in}}%
\pgfusepath{clip}%
\pgfsetbuttcap%
\pgfsetroundjoin%
\definecolor{currentfill}{rgb}{1.000000,1.000000,1.000000}%
\pgfsetfillcolor{currentfill}%
\pgfsetfillopacity{0.900000}%
\pgfsetlinewidth{0.200750pt}%
\definecolor{currentstroke}{rgb}{0.200000,0.200000,0.200000}%
\pgfsetstrokecolor{currentstroke}%
\pgfsetdash{}{0pt}%
\pgfpathmoveto{\pgfqpoint{1.224761in}{1.806028in}}%
\pgfpathlineto{\pgfqpoint{1.270476in}{1.819942in}}%
\pgfpathlineto{\pgfqpoint{1.316095in}{1.833826in}}%
\pgfpathlineto{\pgfqpoint{1.345227in}{1.813988in}}%
\pgfpathlineto{\pgfqpoint{1.374447in}{1.794090in}}%
\pgfpathlineto{\pgfqpoint{1.328752in}{1.780122in}}%
\pgfpathlineto{\pgfqpoint{1.282960in}{1.766123in}}%
\pgfpathlineto{\pgfqpoint{1.253816in}{1.786106in}}%
\pgfpathlineto{\pgfqpoint{1.224761in}{1.806028in}}%
\pgfpathclose%
\pgfusepath{stroke,fill}%
\end{pgfscope}%
\begin{pgfscope}%
\pgfpathrectangle{\pgfqpoint{0.050000in}{0.050000in}}{\pgfqpoint{4.900000in}{4.900000in}}%
\pgfusepath{clip}%
\pgfsetbuttcap%
\pgfsetroundjoin%
\definecolor{currentfill}{rgb}{0.289935,0.260654,0.575425}%
\pgfsetfillcolor{currentfill}%
\pgfsetfillopacity{0.900000}%
\pgfsetlinewidth{0.200750pt}%
\definecolor{currentstroke}{rgb}{0.200000,0.200000,0.200000}%
\pgfsetstrokecolor{currentstroke}%
\pgfsetdash{}{0pt}%
\pgfpathmoveto{\pgfqpoint{2.776298in}{2.701891in}}%
\pgfpathlineto{\pgfqpoint{2.820883in}{2.850058in}}%
\pgfpathlineto{\pgfqpoint{2.865491in}{2.980484in}}%
\pgfpathlineto{\pgfqpoint{2.896926in}{2.940838in}}%
\pgfpathlineto{\pgfqpoint{2.928426in}{2.899987in}}%
\pgfpathlineto{\pgfqpoint{2.883771in}{2.790472in}}%
\pgfpathlineto{\pgfqpoint{2.839068in}{2.659274in}}%
\pgfpathlineto{\pgfqpoint{2.807640in}{2.681377in}}%
\pgfpathlineto{\pgfqpoint{2.776298in}{2.701891in}}%
\pgfpathclose%
\pgfusepath{stroke,fill}%
\end{pgfscope}%
\begin{pgfscope}%
\pgfpathrectangle{\pgfqpoint{0.050000in}{0.050000in}}{\pgfqpoint{4.900000in}{4.900000in}}%
\pgfusepath{clip}%
\pgfsetbuttcap%
\pgfsetroundjoin%
\definecolor{currentfill}{rgb}{1.000000,1.000000,1.000000}%
\pgfsetfillcolor{currentfill}%
\pgfsetfillopacity{0.900000}%
\pgfsetlinewidth{0.200750pt}%
\definecolor{currentstroke}{rgb}{0.200000,0.200000,0.200000}%
\pgfsetstrokecolor{currentstroke}%
\pgfsetdash{}{0pt}%
\pgfpathmoveto{\pgfqpoint{2.155844in}{1.802264in}}%
\pgfpathlineto{\pgfqpoint{2.200582in}{1.816509in}}%
\pgfpathlineto{\pgfqpoint{2.245224in}{1.831244in}}%
\pgfpathlineto{\pgfqpoint{2.275772in}{1.811850in}}%
\pgfpathlineto{\pgfqpoint{2.306413in}{1.792478in}}%
\pgfpathlineto{\pgfqpoint{2.261702in}{1.777118in}}%
\pgfpathlineto{\pgfqpoint{2.216894in}{1.762522in}}%
\pgfpathlineto{\pgfqpoint{2.186322in}{1.782410in}}%
\pgfpathlineto{\pgfqpoint{2.155844in}{1.802264in}}%
\pgfpathclose%
\pgfusepath{stroke,fill}%
\end{pgfscope}%
\begin{pgfscope}%
\pgfpathrectangle{\pgfqpoint{0.050000in}{0.050000in}}{\pgfqpoint{4.900000in}{4.900000in}}%
\pgfusepath{clip}%
\pgfsetbuttcap%
\pgfsetroundjoin%
\definecolor{currentfill}{rgb}{1.000000,1.000000,1.000000}%
\pgfsetfillcolor{currentfill}%
\pgfsetfillopacity{0.900000}%
\pgfsetlinewidth{0.200750pt}%
\definecolor{currentstroke}{rgb}{0.200000,0.200000,0.200000}%
\pgfsetstrokecolor{currentstroke}%
\pgfsetdash{}{0pt}%
\pgfpathmoveto{\pgfqpoint{1.765265in}{1.798100in}}%
\pgfpathlineto{\pgfqpoint{1.810434in}{1.812030in}}%
\pgfpathlineto{\pgfqpoint{1.855507in}{1.825931in}}%
\pgfpathlineto{\pgfqpoint{1.885475in}{1.806069in}}%
\pgfpathlineto{\pgfqpoint{1.915534in}{1.786148in}}%
\pgfpathlineto{\pgfqpoint{1.870388in}{1.772162in}}%
\pgfpathlineto{\pgfqpoint{1.825146in}{1.758147in}}%
\pgfpathlineto{\pgfqpoint{1.795160in}{1.778153in}}%
\pgfpathlineto{\pgfqpoint{1.765265in}{1.798100in}}%
\pgfpathclose%
\pgfusepath{stroke,fill}%
\end{pgfscope}%
\begin{pgfscope}%
\pgfpathrectangle{\pgfqpoint{0.050000in}{0.050000in}}{\pgfqpoint{4.900000in}{4.900000in}}%
\pgfusepath{clip}%
\pgfsetbuttcap%
\pgfsetroundjoin%
\definecolor{currentfill}{rgb}{0.689719,0.676924,0.814471}%
\pgfsetfillcolor{currentfill}%
\pgfsetfillopacity{0.900000}%
\pgfsetlinewidth{0.200750pt}%
\definecolor{currentstroke}{rgb}{0.200000,0.200000,0.200000}%
\pgfsetstrokecolor{currentstroke}%
\pgfsetdash{}{0pt}%
\pgfpathmoveto{\pgfqpoint{3.355839in}{2.394224in}}%
\pgfpathlineto{\pgfqpoint{3.397679in}{2.288000in}}%
\pgfpathlineto{\pgfqpoint{3.439155in}{2.176000in}}%
\pgfpathlineto{\pgfqpoint{3.471102in}{2.133559in}}%
\pgfpathlineto{\pgfqpoint{3.503143in}{2.092380in}}%
\pgfpathlineto{\pgfqpoint{3.461600in}{2.192674in}}%
\pgfpathlineto{\pgfqpoint{3.419762in}{2.291111in}}%
\pgfpathlineto{\pgfqpoint{3.387758in}{2.341910in}}%
\pgfpathlineto{\pgfqpoint{3.355839in}{2.394224in}}%
\pgfpathclose%
\pgfusepath{stroke,fill}%
\end{pgfscope}%
\begin{pgfscope}%
\pgfpathrectangle{\pgfqpoint{0.050000in}{0.050000in}}{\pgfqpoint{4.900000in}{4.900000in}}%
\pgfusepath{clip}%
\pgfsetbuttcap%
\pgfsetroundjoin%
\definecolor{currentfill}{rgb}{0.504744,0.484321,0.703868}%
\pgfsetfillcolor{currentfill}%
\pgfsetfillopacity{0.900000}%
\pgfsetlinewidth{0.200750pt}%
\definecolor{currentstroke}{rgb}{0.200000,0.200000,0.200000}%
\pgfsetstrokecolor{currentstroke}%
\pgfsetdash{}{0pt}%
\pgfpathmoveto{\pgfqpoint{2.749886in}{2.383067in}}%
\pgfpathlineto{\pgfqpoint{2.794421in}{2.519749in}}%
\pgfpathlineto{\pgfqpoint{2.839068in}{2.659274in}}%
\pgfpathlineto{\pgfqpoint{2.870581in}{2.635681in}}%
\pgfpathlineto{\pgfqpoint{2.902176in}{2.610705in}}%
\pgfpathlineto{\pgfqpoint{2.857395in}{2.483943in}}%
\pgfpathlineto{\pgfqpoint{2.812692in}{2.355889in}}%
\pgfpathlineto{\pgfqpoint{2.781240in}{2.370041in}}%
\pgfpathlineto{\pgfqpoint{2.749886in}{2.383067in}}%
\pgfpathclose%
\pgfusepath{stroke,fill}%
\end{pgfscope}%
\begin{pgfscope}%
\pgfpathrectangle{\pgfqpoint{0.050000in}{0.050000in}}{\pgfqpoint{4.900000in}{4.900000in}}%
\pgfusepath{clip}%
\pgfsetbuttcap%
\pgfsetroundjoin%
\definecolor{currentfill}{rgb}{0.243030,0.248366,0.598970}%
\pgfsetfillcolor{currentfill}%
\pgfsetfillopacity{0.900000}%
\pgfsetlinewidth{0.200750pt}%
\definecolor{currentstroke}{rgb}{0.200000,0.200000,0.200000}%
\pgfsetstrokecolor{currentstroke}%
\pgfsetdash{}{0pt}%
\pgfpathmoveto{\pgfqpoint{2.865491in}{2.980484in}}%
\pgfpathlineto{\pgfqpoint{2.909985in}{3.077209in}}%
\pgfpathlineto{\pgfqpoint{2.954181in}{3.124722in}}%
\pgfpathlineto{\pgfqpoint{2.985545in}{3.061858in}}%
\pgfpathlineto{\pgfqpoint{3.016961in}{2.999756in}}%
\pgfpathlineto{\pgfqpoint{2.972882in}{2.973841in}}%
\pgfpathlineto{\pgfqpoint{2.928426in}{2.899987in}}%
\pgfpathlineto{\pgfqpoint{2.896926in}{2.940838in}}%
\pgfpathlineto{\pgfqpoint{2.865491in}{2.980484in}}%
\pgfpathclose%
\pgfusepath{stroke,fill}%
\end{pgfscope}%
\begin{pgfscope}%
\pgfpathrectangle{\pgfqpoint{0.050000in}{0.050000in}}{\pgfqpoint{4.900000in}{4.900000in}}%
\pgfusepath{clip}%
\pgfsetbuttcap%
\pgfsetroundjoin%
\definecolor{currentfill}{rgb}{1.000000,1.000000,1.000000}%
\pgfsetfillcolor{currentfill}%
\pgfsetfillopacity{0.900000}%
\pgfsetlinewidth{0.200750pt}%
\definecolor{currentstroke}{rgb}{0.200000,0.200000,0.200000}%
\pgfsetstrokecolor{currentstroke}%
\pgfsetdash{}{0pt}%
\pgfpathmoveto{\pgfqpoint{1.374447in}{1.794090in}}%
\pgfpathlineto{\pgfqpoint{1.420045in}{1.808029in}}%
\pgfpathlineto{\pgfqpoint{1.465546in}{1.821939in}}%
\pgfpathlineto{\pgfqpoint{1.494931in}{1.802065in}}%
\pgfpathlineto{\pgfqpoint{1.524406in}{1.782131in}}%
\pgfpathlineto{\pgfqpoint{1.478829in}{1.768137in}}%
\pgfpathlineto{\pgfqpoint{1.433155in}{1.754113in}}%
\pgfpathlineto{\pgfqpoint{1.403756in}{1.774132in}}%
\pgfpathlineto{\pgfqpoint{1.374447in}{1.794090in}}%
\pgfpathclose%
\pgfusepath{stroke,fill}%
\end{pgfscope}%
\begin{pgfscope}%
\pgfpathrectangle{\pgfqpoint{0.050000in}{0.050000in}}{\pgfqpoint{4.900000in}{4.900000in}}%
\pgfusepath{clip}%
\pgfsetbuttcap%
\pgfsetroundjoin%
\definecolor{currentfill}{rgb}{1.000000,1.000000,1.000000}%
\pgfsetfillcolor{currentfill}%
\pgfsetfillopacity{0.900000}%
\pgfsetlinewidth{0.200750pt}%
\definecolor{currentstroke}{rgb}{0.200000,0.200000,0.200000}%
\pgfsetstrokecolor{currentstroke}%
\pgfsetdash{}{0pt}%
\pgfpathmoveto{\pgfqpoint{0.983390in}{1.790079in}}%
\pgfpathlineto{\pgfqpoint{1.029417in}{1.804026in}}%
\pgfpathlineto{\pgfqpoint{1.075347in}{1.817944in}}%
\pgfpathlineto{\pgfqpoint{1.104149in}{1.798058in}}%
\pgfpathlineto{\pgfqpoint{1.133038in}{1.778112in}}%
\pgfpathlineto{\pgfqpoint{1.087030in}{1.764109in}}%
\pgfpathlineto{\pgfqpoint{1.040924in}{1.750077in}}%
\pgfpathlineto{\pgfqpoint{1.012113in}{1.770108in}}%
\pgfpathlineto{\pgfqpoint{0.983390in}{1.790079in}}%
\pgfpathclose%
\pgfusepath{stroke,fill}%
\end{pgfscope}%
\begin{pgfscope}%
\pgfpathrectangle{\pgfqpoint{0.050000in}{0.050000in}}{\pgfqpoint{4.900000in}{4.900000in}}%
\pgfusepath{clip}%
\pgfsetbuttcap%
\pgfsetroundjoin%
\definecolor{currentfill}{rgb}{0.319769,0.291719,0.593264}%
\pgfsetfillcolor{currentfill}%
\pgfsetfillopacity{0.900000}%
\pgfsetlinewidth{0.200750pt}%
\definecolor{currentstroke}{rgb}{0.200000,0.200000,0.200000}%
\pgfsetstrokecolor{currentstroke}%
\pgfsetdash{}{0pt}%
\pgfpathmoveto{\pgfqpoint{3.079964in}{2.878377in}}%
\pgfpathlineto{\pgfqpoint{3.123326in}{2.836895in}}%
\pgfpathlineto{\pgfqpoint{3.166004in}{2.747127in}}%
\pgfpathlineto{\pgfqpoint{3.197453in}{2.682911in}}%
\pgfpathlineto{\pgfqpoint{3.228975in}{2.621079in}}%
\pgfpathlineto{\pgfqpoint{3.186425in}{2.711615in}}%
\pgfpathlineto{\pgfqpoint{3.143209in}{2.761237in}}%
\pgfpathlineto{\pgfqpoint{3.111555in}{2.819257in}}%
\pgfpathlineto{\pgfqpoint{3.079964in}{2.878377in}}%
\pgfpathclose%
\pgfusepath{stroke,fill}%
\end{pgfscope}%
\begin{pgfscope}%
\pgfpathrectangle{\pgfqpoint{0.050000in}{0.050000in}}{\pgfqpoint{4.900000in}{4.900000in}}%
\pgfusepath{clip}%
\pgfsetbuttcap%
\pgfsetroundjoin%
\definecolor{currentfill}{rgb}{0.994033,0.993787,0.996432}%
\pgfsetfillcolor{currentfill}%
\pgfsetfillopacity{0.900000}%
\pgfsetlinewidth{0.200750pt}%
\definecolor{currentstroke}{rgb}{0.200000,0.200000,0.200000}%
\pgfsetstrokecolor{currentstroke}%
\pgfsetdash{}{0pt}%
\pgfpathmoveto{\pgfqpoint{2.306413in}{1.792478in}}%
\pgfpathlineto{\pgfqpoint{2.351025in}{1.809396in}}%
\pgfpathlineto{\pgfqpoint{2.395538in}{1.829176in}}%
\pgfpathlineto{\pgfqpoint{2.426345in}{1.811303in}}%
\pgfpathlineto{\pgfqpoint{2.457248in}{1.793509in}}%
\pgfpathlineto{\pgfqpoint{2.412660in}{1.771940in}}%
\pgfpathlineto{\pgfqpoint{2.367976in}{1.753814in}}%
\pgfpathlineto{\pgfqpoint{2.337148in}{1.773133in}}%
\pgfpathlineto{\pgfqpoint{2.306413in}{1.792478in}}%
\pgfpathclose%
\pgfusepath{stroke,fill}%
\end{pgfscope}%
\begin{pgfscope}%
\pgfpathrectangle{\pgfqpoint{0.050000in}{0.050000in}}{\pgfqpoint{4.900000in}{4.900000in}}%
\pgfusepath{clip}%
\pgfsetbuttcap%
\pgfsetroundjoin%
\definecolor{currentfill}{rgb}{0.826959,0.819823,0.896532}%
\pgfsetfillcolor{currentfill}%
\pgfsetfillopacity{0.900000}%
\pgfsetlinewidth{0.200750pt}%
\definecolor{currentstroke}{rgb}{0.200000,0.200000,0.200000}%
\pgfsetstrokecolor{currentstroke}%
\pgfsetdash{}{0pt}%
\pgfpathmoveto{\pgfqpoint{2.634872in}{1.960440in}}%
\pgfpathlineto{\pgfqpoint{2.679219in}{2.035436in}}%
\pgfpathlineto{\pgfqpoint{2.723616in}{2.127222in}}%
\pgfpathlineto{\pgfqpoint{2.754991in}{2.115875in}}%
\pgfpathlineto{\pgfqpoint{2.786468in}{2.103867in}}%
\pgfpathlineto{\pgfqpoint{2.741904in}{2.012204in}}%
\pgfpathlineto{\pgfqpoint{2.697405in}{1.935638in}}%
\pgfpathlineto{\pgfqpoint{2.666087in}{1.948168in}}%
\pgfpathlineto{\pgfqpoint{2.634872in}{1.960440in}}%
\pgfpathclose%
\pgfusepath{stroke,fill}%
\end{pgfscope}%
\begin{pgfscope}%
\pgfpathrectangle{\pgfqpoint{0.050000in}{0.050000in}}{\pgfqpoint{4.900000in}{4.900000in}}%
\pgfusepath{clip}%
\pgfsetbuttcap%
\pgfsetroundjoin%
\definecolor{currentfill}{rgb}{1.000000,1.000000,1.000000}%
\pgfsetfillcolor{currentfill}%
\pgfsetfillopacity{0.900000}%
\pgfsetlinewidth{0.200750pt}%
\definecolor{currentstroke}{rgb}{0.200000,0.200000,0.200000}%
\pgfsetstrokecolor{currentstroke}%
\pgfsetdash{}{0pt}%
\pgfpathmoveto{\pgfqpoint{1.915534in}{1.786148in}}%
\pgfpathlineto{\pgfqpoint{1.960585in}{1.800104in}}%
\pgfpathlineto{\pgfqpoint{2.005540in}{1.814032in}}%
\pgfpathlineto{\pgfqpoint{2.035763in}{1.794138in}}%
\pgfpathlineto{\pgfqpoint{2.066078in}{1.774186in}}%
\pgfpathlineto{\pgfqpoint{2.021051in}{1.760165in}}%
\pgfpathlineto{\pgfqpoint{1.975927in}{1.746122in}}%
\pgfpathlineto{\pgfqpoint{1.945685in}{1.766165in}}%
\pgfpathlineto{\pgfqpoint{1.915534in}{1.786148in}}%
\pgfpathclose%
\pgfusepath{stroke,fill}%
\end{pgfscope}%
\begin{pgfscope}%
\pgfpathrectangle{\pgfqpoint{0.050000in}{0.050000in}}{\pgfqpoint{4.900000in}{4.900000in}}%
\pgfusepath{clip}%
\pgfsetbuttcap%
\pgfsetroundjoin%
\definecolor{currentfill}{rgb}{0.922430,0.919231,0.953618}%
\pgfsetfillcolor{currentfill}%
\pgfsetfillopacity{0.900000}%
\pgfsetlinewidth{0.200750pt}%
\definecolor{currentstroke}{rgb}{0.200000,0.200000,0.200000}%
\pgfsetstrokecolor{currentstroke}%
\pgfsetdash{}{0pt}%
\pgfpathmoveto{\pgfqpoint{2.546168in}{1.855468in}}%
\pgfpathlineto{\pgfqpoint{2.590534in}{1.901186in}}%
\pgfpathlineto{\pgfqpoint{2.634872in}{1.960440in}}%
\pgfpathlineto{\pgfqpoint{2.666087in}{1.948168in}}%
\pgfpathlineto{\pgfqpoint{2.697405in}{1.935638in}}%
\pgfpathlineto{\pgfqpoint{2.652932in}{1.873845in}}%
\pgfpathlineto{\pgfqpoint{2.608451in}{1.825234in}}%
\pgfpathlineto{\pgfqpoint{2.577259in}{1.840359in}}%
\pgfpathlineto{\pgfqpoint{2.546168in}{1.855468in}}%
\pgfpathclose%
\pgfusepath{stroke,fill}%
\end{pgfscope}%
\begin{pgfscope}%
\pgfpathrectangle{\pgfqpoint{0.050000in}{0.050000in}}{\pgfqpoint{4.900000in}{4.900000in}}%
\pgfusepath{clip}%
\pgfsetbuttcap%
\pgfsetroundjoin%
\definecolor{currentfill}{rgb}{1.000000,1.000000,1.000000}%
\pgfsetfillcolor{currentfill}%
\pgfsetfillopacity{0.900000}%
\pgfsetlinewidth{0.200750pt}%
\definecolor{currentstroke}{rgb}{0.200000,0.200000,0.200000}%
\pgfsetstrokecolor{currentstroke}%
\pgfsetdash{}{0pt}%
\pgfpathmoveto{\pgfqpoint{1.524406in}{1.782131in}}%
\pgfpathlineto{\pgfqpoint{1.569886in}{1.796095in}}%
\pgfpathlineto{\pgfqpoint{1.615269in}{1.810030in}}%
\pgfpathlineto{\pgfqpoint{1.644909in}{1.790120in}}%
\pgfpathlineto{\pgfqpoint{1.674639in}{1.770150in}}%
\pgfpathlineto{\pgfqpoint{1.629180in}{1.756130in}}%
\pgfpathlineto{\pgfqpoint{1.583625in}{1.742081in}}%
\pgfpathlineto{\pgfqpoint{1.553970in}{1.762136in}}%
\pgfpathlineto{\pgfqpoint{1.524406in}{1.782131in}}%
\pgfpathclose%
\pgfusepath{stroke,fill}%
\end{pgfscope}%
\begin{pgfscope}%
\pgfpathrectangle{\pgfqpoint{0.050000in}{0.050000in}}{\pgfqpoint{4.900000in}{4.900000in}}%
\pgfusepath{clip}%
\pgfsetbuttcap%
\pgfsetroundjoin%
\definecolor{currentfill}{rgb}{1.000000,1.000000,1.000000}%
\pgfsetfillcolor{currentfill}%
\pgfsetfillopacity{0.900000}%
\pgfsetlinewidth{0.200750pt}%
\definecolor{currentstroke}{rgb}{0.200000,0.200000,0.200000}%
\pgfsetstrokecolor{currentstroke}%
\pgfsetdash{}{0pt}%
\pgfpathmoveto{\pgfqpoint{1.133038in}{1.778112in}}%
\pgfpathlineto{\pgfqpoint{1.178948in}{1.792085in}}%
\pgfpathlineto{\pgfqpoint{1.224761in}{1.806028in}}%
\pgfpathlineto{\pgfqpoint{1.253816in}{1.786106in}}%
\pgfpathlineto{\pgfqpoint{1.282960in}{1.766123in}}%
\pgfpathlineto{\pgfqpoint{1.237070in}{1.752095in}}%
\pgfpathlineto{\pgfqpoint{1.191082in}{1.738037in}}%
\pgfpathlineto{\pgfqpoint{1.162016in}{1.758105in}}%
\pgfpathlineto{\pgfqpoint{1.133038in}{1.778112in}}%
\pgfpathclose%
\pgfusepath{stroke,fill}%
\end{pgfscope}%
\begin{pgfscope}%
\pgfpathrectangle{\pgfqpoint{0.050000in}{0.050000in}}{\pgfqpoint{4.900000in}{4.900000in}}%
\pgfusepath{clip}%
\pgfsetbuttcap%
\pgfsetroundjoin%
\definecolor{currentfill}{rgb}{0.683752,0.670711,0.810903}%
\pgfsetfillcolor{currentfill}%
\pgfsetfillopacity{0.900000}%
\pgfsetlinewidth{0.200750pt}%
\definecolor{currentstroke}{rgb}{0.200000,0.200000,0.200000}%
\pgfsetstrokecolor{currentstroke}%
\pgfsetdash{}{0pt}%
\pgfpathmoveto{\pgfqpoint{2.723616in}{2.127222in}}%
\pgfpathlineto{\pgfqpoint{2.768099in}{2.235061in}}%
\pgfpathlineto{\pgfqpoint{2.812692in}{2.355889in}}%
\pgfpathlineto{\pgfqpoint{2.844239in}{2.340627in}}%
\pgfpathlineto{\pgfqpoint{2.875882in}{2.324282in}}%
\pgfpathlineto{\pgfqpoint{2.831123in}{2.209139in}}%
\pgfpathlineto{\pgfqpoint{2.786468in}{2.103867in}}%
\pgfpathlineto{\pgfqpoint{2.754991in}{2.115875in}}%
\pgfpathlineto{\pgfqpoint{2.723616in}{2.127222in}}%
\pgfpathclose%
\pgfusepath{stroke,fill}%
\end{pgfscope}%
\begin{pgfscope}%
\pgfpathrectangle{\pgfqpoint{0.050000in}{0.050000in}}{\pgfqpoint{4.900000in}{4.900000in}}%
\pgfusepath{clip}%
\pgfsetbuttcap%
\pgfsetroundjoin%
\definecolor{currentfill}{rgb}{1.000000,1.000000,1.000000}%
\pgfsetfillcolor{currentfill}%
\pgfsetfillopacity{0.900000}%
\pgfsetlinewidth{0.200750pt}%
\definecolor{currentstroke}{rgb}{0.200000,0.200000,0.200000}%
\pgfsetstrokecolor{currentstroke}%
\pgfsetdash{}{0pt}%
\pgfpathmoveto{\pgfqpoint{2.066078in}{1.774186in}}%
\pgfpathlineto{\pgfqpoint{2.111009in}{1.788201in}}%
\pgfpathlineto{\pgfqpoint{2.155844in}{1.802264in}}%
\pgfpathlineto{\pgfqpoint{2.186322in}{1.782410in}}%
\pgfpathlineto{\pgfqpoint{2.216894in}{1.762522in}}%
\pgfpathlineto{\pgfqpoint{2.171988in}{1.748256in}}%
\pgfpathlineto{\pgfqpoint{2.126985in}{1.734109in}}%
\pgfpathlineto{\pgfqpoint{2.096485in}{1.754176in}}%
\pgfpathlineto{\pgfqpoint{2.066078in}{1.774186in}}%
\pgfpathclose%
\pgfusepath{stroke,fill}%
\end{pgfscope}%
\begin{pgfscope}%
\pgfpathrectangle{\pgfqpoint{0.050000in}{0.050000in}}{\pgfqpoint{4.900000in}{4.900000in}}%
\pgfusepath{clip}%
\pgfsetbuttcap%
\pgfsetroundjoin%
\definecolor{currentfill}{rgb}{0.946298,0.944083,0.967889}%
\pgfsetfillcolor{currentfill}%
\pgfsetfillopacity{0.900000}%
\pgfsetlinewidth{0.200750pt}%
\definecolor{currentstroke}{rgb}{0.200000,0.200000,0.200000}%
\pgfsetstrokecolor{currentstroke}%
\pgfsetdash{}{0pt}%
\pgfpathmoveto{\pgfqpoint{3.632251in}{1.938292in}}%
\pgfpathlineto{\pgfqpoint{3.674037in}{1.874100in}}%
\pgfpathlineto{\pgfqpoint{3.716139in}{1.835606in}}%
\pgfpathlineto{\pgfqpoint{3.748845in}{1.810314in}}%
\pgfpathlineto{\pgfqpoint{3.781651in}{1.785260in}}%
\pgfpathlineto{\pgfqpoint{3.739333in}{1.812939in}}%
\pgfpathlineto{\pgfqpoint{3.697382in}{1.866505in}}%
\pgfpathlineto{\pgfqpoint{3.664768in}{1.902017in}}%
\pgfpathlineto{\pgfqpoint{3.632251in}{1.938292in}}%
\pgfpathclose%
\pgfusepath{stroke,fill}%
\end{pgfscope}%
\begin{pgfscope}%
\pgfpathrectangle{\pgfqpoint{0.050000in}{0.050000in}}{\pgfqpoint{4.900000in}{4.900000in}}%
\pgfusepath{clip}%
\pgfsetbuttcap%
\pgfsetroundjoin%
\definecolor{currentfill}{rgb}{0.313802,0.285506,0.589696}%
\pgfsetfillcolor{currentfill}%
\pgfsetfillopacity{0.900000}%
\pgfsetlinewidth{0.200750pt}%
\definecolor{currentstroke}{rgb}{0.200000,0.200000,0.200000}%
\pgfsetstrokecolor{currentstroke}%
\pgfsetdash{}{0pt}%
\pgfpathmoveto{\pgfqpoint{2.839068in}{2.659274in}}%
\pgfpathlineto{\pgfqpoint{2.883771in}{2.790472in}}%
\pgfpathlineto{\pgfqpoint{2.928426in}{2.899987in}}%
\pgfpathlineto{\pgfqpoint{2.959991in}{2.858173in}}%
\pgfpathlineto{\pgfqpoint{2.991621in}{2.815605in}}%
\pgfpathlineto{\pgfqpoint{2.946957in}{2.725318in}}%
\pgfpathlineto{\pgfqpoint{2.902176in}{2.610705in}}%
\pgfpathlineto{\pgfqpoint{2.870581in}{2.635681in}}%
\pgfpathlineto{\pgfqpoint{2.839068in}{2.659274in}}%
\pgfpathclose%
\pgfusepath{stroke,fill}%
\end{pgfscope}%
\begin{pgfscope}%
\pgfpathrectangle{\pgfqpoint{0.050000in}{0.050000in}}{\pgfqpoint{4.900000in}{4.900000in}}%
\pgfusepath{clip}%
\pgfsetbuttcap%
\pgfsetroundjoin%
\definecolor{currentfill}{rgb}{1.000000,1.000000,1.000000}%
\pgfsetfillcolor{currentfill}%
\pgfsetfillopacity{0.900000}%
\pgfsetlinewidth{0.200750pt}%
\definecolor{currentstroke}{rgb}{0.200000,0.200000,0.200000}%
\pgfsetstrokecolor{currentstroke}%
\pgfsetdash{}{0pt}%
\pgfpathmoveto{\pgfqpoint{3.781651in}{1.785260in}}%
\pgfpathlineto{\pgfqpoint{3.824593in}{1.792250in}}%
\pgfpathlineto{\pgfqpoint{3.857633in}{1.772286in}}%
\pgfpathlineto{\pgfqpoint{3.890774in}{1.752262in}}%
\pgfpathlineto{\pgfqpoint{3.847620in}{1.738204in}}%
\pgfpathlineto{\pgfqpoint{3.814558in}{1.760423in}}%
\pgfpathlineto{\pgfqpoint{3.781651in}{1.785260in}}%
\pgfpathclose%
\pgfusepath{stroke,fill}%
\end{pgfscope}%
\begin{pgfscope}%
\pgfpathrectangle{\pgfqpoint{0.050000in}{0.050000in}}{\pgfqpoint{4.900000in}{4.900000in}}%
\pgfusepath{clip}%
\pgfsetbuttcap%
\pgfsetroundjoin%
\definecolor{currentfill}{rgb}{1.000000,1.000000,1.000000}%
\pgfsetfillcolor{currentfill}%
\pgfsetfillopacity{0.900000}%
\pgfsetlinewidth{0.200750pt}%
\definecolor{currentstroke}{rgb}{0.200000,0.200000,0.200000}%
\pgfsetstrokecolor{currentstroke}%
\pgfsetdash{}{0pt}%
\pgfpathmoveto{\pgfqpoint{1.674639in}{1.770150in}}%
\pgfpathlineto{\pgfqpoint{1.720000in}{1.784139in}}%
\pgfpathlineto{\pgfqpoint{1.765265in}{1.798100in}}%
\pgfpathlineto{\pgfqpoint{1.795160in}{1.778153in}}%
\pgfpathlineto{\pgfqpoint{1.825146in}{1.758147in}}%
\pgfpathlineto{\pgfqpoint{1.779806in}{1.744101in}}%
\pgfpathlineto{\pgfqpoint{1.734370in}{1.730026in}}%
\pgfpathlineto{\pgfqpoint{1.704459in}{1.750119in}}%
\pgfpathlineto{\pgfqpoint{1.674639in}{1.770150in}}%
\pgfpathclose%
\pgfusepath{stroke,fill}%
\end{pgfscope}%
\begin{pgfscope}%
\pgfpathrectangle{\pgfqpoint{0.050000in}{0.050000in}}{\pgfqpoint{4.900000in}{4.900000in}}%
\pgfusepath{clip}%
\pgfsetbuttcap%
\pgfsetroundjoin%
\definecolor{currentfill}{rgb}{0.970165,0.968935,0.982161}%
\pgfsetfillcolor{currentfill}%
\pgfsetfillopacity{0.900000}%
\pgfsetlinewidth{0.200750pt}%
\definecolor{currentstroke}{rgb}{0.200000,0.200000,0.200000}%
\pgfsetstrokecolor{currentstroke}%
\pgfsetdash{}{0pt}%
\pgfpathmoveto{\pgfqpoint{2.457248in}{1.793509in}}%
\pgfpathlineto{\pgfqpoint{2.501746in}{1.820516in}}%
\pgfpathlineto{\pgfqpoint{2.546168in}{1.855468in}}%
\pgfpathlineto{\pgfqpoint{2.577259in}{1.840359in}}%
\pgfpathlineto{\pgfqpoint{2.608451in}{1.825234in}}%
\pgfpathlineto{\pgfqpoint{2.563929in}{1.787501in}}%
\pgfpathlineto{\pgfqpoint{2.519343in}{1.758116in}}%
\pgfpathlineto{\pgfqpoint{2.488247in}{1.775784in}}%
\pgfpathlineto{\pgfqpoint{2.457248in}{1.793509in}}%
\pgfpathclose%
\pgfusepath{stroke,fill}%
\end{pgfscope}%
\begin{pgfscope}%
\pgfpathrectangle{\pgfqpoint{0.050000in}{0.050000in}}{\pgfqpoint{4.900000in}{4.900000in}}%
\pgfusepath{clip}%
\pgfsetbuttcap%
\pgfsetroundjoin%
\definecolor{currentfill}{rgb}{0.240569,0.222222,0.564214}%
\pgfsetfillcolor{currentfill}%
\pgfsetfillopacity{0.900000}%
\pgfsetlinewidth{0.200750pt}%
\definecolor{currentstroke}{rgb}{0.200000,0.200000,0.200000}%
\pgfsetstrokecolor{currentstroke}%
\pgfsetdash{}{0pt}%
\pgfpathmoveto{\pgfqpoint{2.928426in}{2.899987in}}%
\pgfpathlineto{\pgfqpoint{2.972882in}{2.973841in}}%
\pgfpathlineto{\pgfqpoint{3.016961in}{2.999756in}}%
\pgfpathlineto{\pgfqpoint{3.048433in}{2.938563in}}%
\pgfpathlineto{\pgfqpoint{3.079964in}{2.878377in}}%
\pgfpathlineto{\pgfqpoint{3.036012in}{2.869738in}}%
\pgfpathlineto{\pgfqpoint{2.991621in}{2.815605in}}%
\pgfpathlineto{\pgfqpoint{2.959991in}{2.858173in}}%
\pgfpathlineto{\pgfqpoint{2.928426in}{2.899987in}}%
\pgfpathclose%
\pgfusepath{stroke,fill}%
\end{pgfscope}%
\begin{pgfscope}%
\pgfpathrectangle{\pgfqpoint{0.050000in}{0.050000in}}{\pgfqpoint{4.900000in}{4.900000in}}%
\pgfusepath{clip}%
\pgfsetbuttcap%
\pgfsetroundjoin%
\definecolor{currentfill}{rgb}{0.731488,0.720415,0.839446}%
\pgfsetfillcolor{currentfill}%
\pgfsetfillopacity{0.900000}%
\pgfsetlinewidth{0.200750pt}%
\definecolor{currentstroke}{rgb}{0.200000,0.200000,0.200000}%
\pgfsetstrokecolor{currentstroke}%
\pgfsetdash{}{0pt}%
\pgfpathmoveto{\pgfqpoint{3.419762in}{2.291111in}}%
\pgfpathlineto{\pgfqpoint{3.461600in}{2.192674in}}%
\pgfpathlineto{\pgfqpoint{3.503143in}{2.092380in}}%
\pgfpathlineto{\pgfqpoint{3.535277in}{2.052355in}}%
\pgfpathlineto{\pgfqpoint{3.567506in}{2.013387in}}%
\pgfpathlineto{\pgfqpoint{3.525880in}{2.102828in}}%
\pgfpathlineto{\pgfqpoint{3.484026in}{2.193648in}}%
\pgfpathlineto{\pgfqpoint{3.451851in}{2.241722in}}%
\pgfpathlineto{\pgfqpoint{3.419762in}{2.291111in}}%
\pgfpathclose%
\pgfusepath{stroke,fill}%
\end{pgfscope}%
\begin{pgfscope}%
\pgfpathrectangle{\pgfqpoint{0.050000in}{0.050000in}}{\pgfqpoint{4.900000in}{4.900000in}}%
\pgfusepath{clip}%
\pgfsetbuttcap%
\pgfsetroundjoin%
\definecolor{currentfill}{rgb}{0.391373,0.366275,0.636078}%
\pgfsetfillcolor{currentfill}%
\pgfsetfillopacity{0.900000}%
\pgfsetlinewidth{0.200750pt}%
\definecolor{currentstroke}{rgb}{0.200000,0.200000,0.200000}%
\pgfsetstrokecolor{currentstroke}%
\pgfsetdash{}{0pt}%
\pgfpathmoveto{\pgfqpoint{3.143209in}{2.761237in}}%
\pgfpathlineto{\pgfqpoint{3.186425in}{2.711615in}}%
\pgfpathlineto{\pgfqpoint{3.228975in}{2.621079in}}%
\pgfpathlineto{\pgfqpoint{3.260572in}{2.561451in}}%
\pgfpathlineto{\pgfqpoint{3.292248in}{2.503863in}}%
\pgfpathlineto{\pgfqpoint{3.249802in}{2.593411in}}%
\pgfpathlineto{\pgfqpoint{3.206717in}{2.648530in}}%
\pgfpathlineto{\pgfqpoint{3.174930in}{2.704329in}}%
\pgfpathlineto{\pgfqpoint{3.143209in}{2.761237in}}%
\pgfpathclose%
\pgfusepath{stroke,fill}%
\end{pgfscope}%
\begin{pgfscope}%
\pgfpathrectangle{\pgfqpoint{0.050000in}{0.050000in}}{\pgfqpoint{4.900000in}{4.900000in}}%
\pgfusepath{clip}%
\pgfsetbuttcap%
\pgfsetroundjoin%
\definecolor{currentfill}{rgb}{1.000000,1.000000,1.000000}%
\pgfsetfillcolor{currentfill}%
\pgfsetfillopacity{0.900000}%
\pgfsetlinewidth{0.200750pt}%
\definecolor{currentstroke}{rgb}{0.200000,0.200000,0.200000}%
\pgfsetstrokecolor{currentstroke}%
\pgfsetdash{}{0pt}%
\pgfpathmoveto{\pgfqpoint{1.282960in}{1.766123in}}%
\pgfpathlineto{\pgfqpoint{1.328752in}{1.780122in}}%
\pgfpathlineto{\pgfqpoint{1.374447in}{1.794090in}}%
\pgfpathlineto{\pgfqpoint{1.403756in}{1.774132in}}%
\pgfpathlineto{\pgfqpoint{1.433155in}{1.754113in}}%
\pgfpathlineto{\pgfqpoint{1.387383in}{1.740059in}}%
\pgfpathlineto{\pgfqpoint{1.341514in}{1.725975in}}%
\pgfpathlineto{\pgfqpoint{1.312192in}{1.746080in}}%
\pgfpathlineto{\pgfqpoint{1.282960in}{1.766123in}}%
\pgfpathclose%
\pgfusepath{stroke,fill}%
\end{pgfscope}%
\begin{pgfscope}%
\pgfpathrectangle{\pgfqpoint{0.050000in}{0.050000in}}{\pgfqpoint{4.900000in}{4.900000in}}%
\pgfusepath{clip}%
\pgfsetbuttcap%
\pgfsetroundjoin%
\definecolor{currentfill}{rgb}{1.000000,1.000000,1.000000}%
\pgfsetfillcolor{currentfill}%
\pgfsetfillopacity{0.900000}%
\pgfsetlinewidth{0.200750pt}%
\definecolor{currentstroke}{rgb}{0.200000,0.200000,0.200000}%
\pgfsetstrokecolor{currentstroke}%
\pgfsetdash{}{0pt}%
\pgfpathmoveto{\pgfqpoint{0.891041in}{1.762095in}}%
\pgfpathlineto{\pgfqpoint{0.937265in}{1.776102in}}%
\pgfpathlineto{\pgfqpoint{0.983390in}{1.790079in}}%
\pgfpathlineto{\pgfqpoint{1.012113in}{1.770108in}}%
\pgfpathlineto{\pgfqpoint{1.040924in}{1.750077in}}%
\pgfpathlineto{\pgfqpoint{0.994720in}{1.736014in}}%
\pgfpathlineto{\pgfqpoint{0.948416in}{1.721922in}}%
\pgfpathlineto{\pgfqpoint{0.919685in}{1.742039in}}%
\pgfpathlineto{\pgfqpoint{0.891041in}{1.762095in}}%
\pgfpathclose%
\pgfusepath{stroke,fill}%
\end{pgfscope}%
\begin{pgfscope}%
\pgfpathrectangle{\pgfqpoint{0.050000in}{0.050000in}}{\pgfqpoint{4.900000in}{4.900000in}}%
\pgfusepath{clip}%
\pgfsetbuttcap%
\pgfsetroundjoin%
\definecolor{currentfill}{rgb}{1.000000,1.000000,1.000000}%
\pgfsetfillcolor{currentfill}%
\pgfsetfillopacity{0.900000}%
\pgfsetlinewidth{0.200750pt}%
\definecolor{currentstroke}{rgb}{0.200000,0.200000,0.200000}%
\pgfsetstrokecolor{currentstroke}%
\pgfsetdash{}{0pt}%
\pgfpathmoveto{\pgfqpoint{2.216894in}{1.762522in}}%
\pgfpathlineto{\pgfqpoint{2.261702in}{1.777118in}}%
\pgfpathlineto{\pgfqpoint{2.306413in}{1.792478in}}%
\pgfpathlineto{\pgfqpoint{2.337148in}{1.773133in}}%
\pgfpathlineto{\pgfqpoint{2.367976in}{1.753814in}}%
\pgfpathlineto{\pgfqpoint{2.323195in}{1.737697in}}%
\pgfpathlineto{\pgfqpoint{2.278316in}{1.722652in}}%
\pgfpathlineto{\pgfqpoint{2.247558in}{1.742601in}}%
\pgfpathlineto{\pgfqpoint{2.216894in}{1.762522in}}%
\pgfpathclose%
\pgfusepath{stroke,fill}%
\end{pgfscope}%
\begin{pgfscope}%
\pgfpathrectangle{\pgfqpoint{0.050000in}{0.050000in}}{\pgfqpoint{4.900000in}{4.900000in}}%
\pgfusepath{clip}%
\pgfsetbuttcap%
\pgfsetroundjoin%
\definecolor{currentfill}{rgb}{1.000000,1.000000,1.000000}%
\pgfsetfillcolor{currentfill}%
\pgfsetfillopacity{0.900000}%
\pgfsetlinewidth{0.200750pt}%
\definecolor{currentstroke}{rgb}{0.200000,0.200000,0.200000}%
\pgfsetstrokecolor{currentstroke}%
\pgfsetdash{}{0pt}%
\pgfpathmoveto{\pgfqpoint{1.825146in}{1.758147in}}%
\pgfpathlineto{\pgfqpoint{1.870388in}{1.772162in}}%
\pgfpathlineto{\pgfqpoint{1.915534in}{1.786148in}}%
\pgfpathlineto{\pgfqpoint{1.945685in}{1.766165in}}%
\pgfpathlineto{\pgfqpoint{1.975927in}{1.746122in}}%
\pgfpathlineto{\pgfqpoint{1.930708in}{1.732051in}}%
\pgfpathlineto{\pgfqpoint{1.885392in}{1.717950in}}%
\pgfpathlineto{\pgfqpoint{1.855223in}{1.738079in}}%
\pgfpathlineto{\pgfqpoint{1.825146in}{1.758147in}}%
\pgfpathclose%
\pgfusepath{stroke,fill}%
\end{pgfscope}%
\begin{pgfscope}%
\pgfpathrectangle{\pgfqpoint{0.050000in}{0.050000in}}{\pgfqpoint{4.900000in}{4.900000in}}%
\pgfusepath{clip}%
\pgfsetbuttcap%
\pgfsetroundjoin%
\definecolor{currentfill}{rgb}{0.504744,0.484321,0.703868}%
\pgfsetfillcolor{currentfill}%
\pgfsetfillopacity{0.900000}%
\pgfsetlinewidth{0.200750pt}%
\definecolor{currentstroke}{rgb}{0.200000,0.200000,0.200000}%
\pgfsetstrokecolor{currentstroke}%
\pgfsetdash{}{0pt}%
\pgfpathmoveto{\pgfqpoint{2.812692in}{2.355889in}}%
\pgfpathlineto{\pgfqpoint{2.857395in}{2.483943in}}%
\pgfpathlineto{\pgfqpoint{2.902176in}{2.610705in}}%
\pgfpathlineto{\pgfqpoint{2.933852in}{2.584459in}}%
\pgfpathlineto{\pgfqpoint{2.965610in}{2.557055in}}%
\pgfpathlineto{\pgfqpoint{2.920728in}{2.443096in}}%
\pgfpathlineto{\pgfqpoint{2.875882in}{2.324282in}}%
\pgfpathlineto{\pgfqpoint{2.844239in}{2.340627in}}%
\pgfpathlineto{\pgfqpoint{2.812692in}{2.355889in}}%
\pgfpathclose%
\pgfusepath{stroke,fill}%
\end{pgfscope}%
\begin{pgfscope}%
\pgfpathrectangle{\pgfqpoint{0.050000in}{0.050000in}}{\pgfqpoint{4.900000in}{4.900000in}}%
\pgfusepath{clip}%
\pgfsetbuttcap%
\pgfsetroundjoin%
\definecolor{currentfill}{rgb}{1.000000,1.000000,1.000000}%
\pgfsetfillcolor{currentfill}%
\pgfsetfillopacity{0.900000}%
\pgfsetlinewidth{0.200750pt}%
\definecolor{currentstroke}{rgb}{0.200000,0.200000,0.200000}%
\pgfsetstrokecolor{currentstroke}%
\pgfsetdash{}{0pt}%
\pgfpathmoveto{\pgfqpoint{1.433155in}{1.754113in}}%
\pgfpathlineto{\pgfqpoint{1.478829in}{1.768137in}}%
\pgfpathlineto{\pgfqpoint{1.524406in}{1.782131in}}%
\pgfpathlineto{\pgfqpoint{1.553970in}{1.762136in}}%
\pgfpathlineto{\pgfqpoint{1.583625in}{1.742081in}}%
\pgfpathlineto{\pgfqpoint{1.537972in}{1.728001in}}%
\pgfpathlineto{\pgfqpoint{1.492222in}{1.713891in}}%
\pgfpathlineto{\pgfqpoint{1.462643in}{1.734033in}}%
\pgfpathlineto{\pgfqpoint{1.433155in}{1.754113in}}%
\pgfpathclose%
\pgfusepath{stroke,fill}%
\end{pgfscope}%
\begin{pgfscope}%
\pgfpathrectangle{\pgfqpoint{0.050000in}{0.050000in}}{\pgfqpoint{4.900000in}{4.900000in}}%
\pgfusepath{clip}%
\pgfsetbuttcap%
\pgfsetroundjoin%
\definecolor{currentfill}{rgb}{1.000000,1.000000,1.000000}%
\pgfsetfillcolor{currentfill}%
\pgfsetfillopacity{0.900000}%
\pgfsetlinewidth{0.200750pt}%
\definecolor{currentstroke}{rgb}{0.200000,0.200000,0.200000}%
\pgfsetstrokecolor{currentstroke}%
\pgfsetdash{}{0pt}%
\pgfpathmoveto{\pgfqpoint{1.040924in}{1.750077in}}%
\pgfpathlineto{\pgfqpoint{1.087030in}{1.764109in}}%
\pgfpathlineto{\pgfqpoint{1.133038in}{1.778112in}}%
\pgfpathlineto{\pgfqpoint{1.162016in}{1.758105in}}%
\pgfpathlineto{\pgfqpoint{1.191082in}{1.738037in}}%
\pgfpathlineto{\pgfqpoint{1.144995in}{1.723949in}}%
\pgfpathlineto{\pgfqpoint{1.098810in}{1.709831in}}%
\pgfpathlineto{\pgfqpoint{1.069823in}{1.729984in}}%
\pgfpathlineto{\pgfqpoint{1.040924in}{1.750077in}}%
\pgfpathclose%
\pgfusepath{stroke,fill}%
\end{pgfscope}%
\begin{pgfscope}%
\pgfpathrectangle{\pgfqpoint{0.050000in}{0.050000in}}{\pgfqpoint{4.900000in}{4.900000in}}%
\pgfusepath{clip}%
\pgfsetbuttcap%
\pgfsetroundjoin%
\definecolor{currentfill}{rgb}{0.994033,0.993787,0.996432}%
\pgfsetfillcolor{currentfill}%
\pgfsetfillopacity{0.900000}%
\pgfsetlinewidth{0.200750pt}%
\definecolor{currentstroke}{rgb}{0.200000,0.200000,0.200000}%
\pgfsetstrokecolor{currentstroke}%
\pgfsetdash{}{0pt}%
\pgfpathmoveto{\pgfqpoint{2.367976in}{1.753814in}}%
\pgfpathlineto{\pgfqpoint{2.412660in}{1.771940in}}%
\pgfpathlineto{\pgfqpoint{2.457248in}{1.793509in}}%
\pgfpathlineto{\pgfqpoint{2.488247in}{1.775784in}}%
\pgfpathlineto{\pgfqpoint{2.519343in}{1.758116in}}%
\pgfpathlineto{\pgfqpoint{2.474676in}{1.734700in}}%
\pgfpathlineto{\pgfqpoint{2.429918in}{1.715256in}}%
\pgfpathlineto{\pgfqpoint{2.398900in}{1.734522in}}%
\pgfpathlineto{\pgfqpoint{2.367976in}{1.753814in}}%
\pgfpathclose%
\pgfusepath{stroke,fill}%
\end{pgfscope}%
\begin{pgfscope}%
\pgfpathrectangle{\pgfqpoint{0.050000in}{0.050000in}}{\pgfqpoint{4.900000in}{4.900000in}}%
\pgfusepath{clip}%
\pgfsetbuttcap%
\pgfsetroundjoin%
\definecolor{currentfill}{rgb}{1.000000,1.000000,1.000000}%
\pgfsetfillcolor{currentfill}%
\pgfsetfillopacity{0.900000}%
\pgfsetlinewidth{0.200750pt}%
\definecolor{currentstroke}{rgb}{0.200000,0.200000,0.200000}%
\pgfsetstrokecolor{currentstroke}%
\pgfsetdash{}{0pt}%
\pgfpathmoveto{\pgfqpoint{1.975927in}{1.746122in}}%
\pgfpathlineto{\pgfqpoint{2.021051in}{1.760165in}}%
\pgfpathlineto{\pgfqpoint{2.066078in}{1.774186in}}%
\pgfpathlineto{\pgfqpoint{2.096485in}{1.754176in}}%
\pgfpathlineto{\pgfqpoint{2.126985in}{1.734109in}}%
\pgfpathlineto{\pgfqpoint{2.081886in}{1.719987in}}%
\pgfpathlineto{\pgfqpoint{2.036690in}{1.705853in}}%
\pgfpathlineto{\pgfqpoint{2.006262in}{1.726018in}}%
\pgfpathlineto{\pgfqpoint{1.975927in}{1.746122in}}%
\pgfpathclose%
\pgfusepath{stroke,fill}%
\end{pgfscope}%
\begin{pgfscope}%
\pgfpathrectangle{\pgfqpoint{0.050000in}{0.050000in}}{\pgfqpoint{4.900000in}{4.900000in}}%
\pgfusepath{clip}%
\pgfsetbuttcap%
\pgfsetroundjoin%
\definecolor{currentfill}{rgb}{0.451042,0.428404,0.671757}%
\pgfsetfillcolor{currentfill}%
\pgfsetfillopacity{0.900000}%
\pgfsetlinewidth{0.200750pt}%
\definecolor{currentstroke}{rgb}{0.200000,0.200000,0.200000}%
\pgfsetstrokecolor{currentstroke}%
\pgfsetdash{}{0pt}%
\pgfpathmoveto{\pgfqpoint{3.206717in}{2.648530in}}%
\pgfpathlineto{\pgfqpoint{3.249802in}{2.593411in}}%
\pgfpathlineto{\pgfqpoint{3.292248in}{2.503863in}}%
\pgfpathlineto{\pgfqpoint{3.324003in}{2.448167in}}%
\pgfpathlineto{\pgfqpoint{3.355839in}{2.394224in}}%
\pgfpathlineto{\pgfqpoint{3.313474in}{2.481585in}}%
\pgfpathlineto{\pgfqpoint{3.270505in}{2.540199in}}%
\pgfpathlineto{\pgfqpoint{3.238575in}{2.593827in}}%
\pgfpathlineto{\pgfqpoint{3.206717in}{2.648530in}}%
\pgfpathclose%
\pgfusepath{stroke,fill}%
\end{pgfscope}%
\begin{pgfscope}%
\pgfpathrectangle{\pgfqpoint{0.050000in}{0.050000in}}{\pgfqpoint{4.900000in}{4.900000in}}%
\pgfusepath{clip}%
\pgfsetbuttcap%
\pgfsetroundjoin%
\definecolor{currentfill}{rgb}{0.260100,0.229589,0.557586}%
\pgfsetfillcolor{currentfill}%
\pgfsetfillopacity{0.900000}%
\pgfsetlinewidth{0.200750pt}%
\definecolor{currentstroke}{rgb}{0.200000,0.200000,0.200000}%
\pgfsetstrokecolor{currentstroke}%
\pgfsetdash{}{0pt}%
\pgfpathmoveto{\pgfqpoint{2.991621in}{2.815605in}}%
\pgfpathlineto{\pgfqpoint{3.036012in}{2.869738in}}%
\pgfpathlineto{\pgfqpoint{3.079964in}{2.878377in}}%
\pgfpathlineto{\pgfqpoint{3.111555in}{2.819257in}}%
\pgfpathlineto{\pgfqpoint{3.143209in}{2.761237in}}%
\pgfpathlineto{\pgfqpoint{3.099387in}{2.766287in}}%
\pgfpathlineto{\pgfqpoint{3.055079in}{2.728915in}}%
\pgfpathlineto{\pgfqpoint{3.023316in}{2.772467in}}%
\pgfpathlineto{\pgfqpoint{2.991621in}{2.815605in}}%
\pgfpathclose%
\pgfusepath{stroke,fill}%
\end{pgfscope}%
\begin{pgfscope}%
\pgfpathrectangle{\pgfqpoint{0.050000in}{0.050000in}}{\pgfqpoint{4.900000in}{4.900000in}}%
\pgfusepath{clip}%
\pgfsetbuttcap%
\pgfsetroundjoin%
\definecolor{currentfill}{rgb}{0.815025,0.807397,0.889396}%
\pgfsetfillcolor{currentfill}%
\pgfsetfillopacity{0.900000}%
\pgfsetlinewidth{0.200750pt}%
\definecolor{currentstroke}{rgb}{0.200000,0.200000,0.200000}%
\pgfsetstrokecolor{currentstroke}%
\pgfsetdash{}{0pt}%
\pgfpathmoveto{\pgfqpoint{2.697405in}{1.935638in}}%
\pgfpathlineto{\pgfqpoint{2.741904in}{2.012204in}}%
\pgfpathlineto{\pgfqpoint{2.786468in}{2.103867in}}%
\pgfpathlineto{\pgfqpoint{2.818046in}{2.091187in}}%
\pgfpathlineto{\pgfqpoint{2.849726in}{2.077828in}}%
\pgfpathlineto{\pgfqpoint{2.805001in}{1.987118in}}%
\pgfpathlineto{\pgfqpoint{2.760349in}{1.909715in}}%
\pgfpathlineto{\pgfqpoint{2.728825in}{1.922827in}}%
\pgfpathlineto{\pgfqpoint{2.697405in}{1.935638in}}%
\pgfpathclose%
\pgfusepath{stroke,fill}%
\end{pgfscope}%
\begin{pgfscope}%
\pgfpathrectangle{\pgfqpoint{0.050000in}{0.050000in}}{\pgfqpoint{4.900000in}{4.900000in}}%
\pgfusepath{clip}%
\pgfsetbuttcap%
\pgfsetroundjoin%
\definecolor{currentfill}{rgb}{1.000000,1.000000,1.000000}%
\pgfsetfillcolor{currentfill}%
\pgfsetfillopacity{0.900000}%
\pgfsetlinewidth{0.200750pt}%
\definecolor{currentstroke}{rgb}{0.200000,0.200000,0.200000}%
\pgfsetstrokecolor{currentstroke}%
\pgfsetdash{}{0pt}%
\pgfpathmoveto{\pgfqpoint{1.583625in}{1.742081in}}%
\pgfpathlineto{\pgfqpoint{1.629180in}{1.756130in}}%
\pgfpathlineto{\pgfqpoint{1.674639in}{1.770150in}}%
\pgfpathlineto{\pgfqpoint{1.704459in}{1.750119in}}%
\pgfpathlineto{\pgfqpoint{1.734370in}{1.730026in}}%
\pgfpathlineto{\pgfqpoint{1.688837in}{1.715921in}}%
\pgfpathlineto{\pgfqpoint{1.643206in}{1.701785in}}%
\pgfpathlineto{\pgfqpoint{1.613370in}{1.721964in}}%
\pgfpathlineto{\pgfqpoint{1.583625in}{1.742081in}}%
\pgfpathclose%
\pgfusepath{stroke,fill}%
\end{pgfscope}%
\begin{pgfscope}%
\pgfpathrectangle{\pgfqpoint{0.050000in}{0.050000in}}{\pgfqpoint{4.900000in}{4.900000in}}%
\pgfusepath{clip}%
\pgfsetbuttcap%
\pgfsetroundjoin%
\definecolor{currentfill}{rgb}{0.910496,0.906805,0.946482}%
\pgfsetfillcolor{currentfill}%
\pgfsetfillopacity{0.900000}%
\pgfsetlinewidth{0.200750pt}%
\definecolor{currentstroke}{rgb}{0.200000,0.200000,0.200000}%
\pgfsetstrokecolor{currentstroke}%
\pgfsetdash{}{0pt}%
\pgfpathmoveto{\pgfqpoint{2.608451in}{1.825234in}}%
\pgfpathlineto{\pgfqpoint{2.652932in}{1.873845in}}%
\pgfpathlineto{\pgfqpoint{2.697405in}{1.935638in}}%
\pgfpathlineto{\pgfqpoint{2.728825in}{1.922827in}}%
\pgfpathlineto{\pgfqpoint{2.760349in}{1.909715in}}%
\pgfpathlineto{\pgfqpoint{2.715738in}{1.845962in}}%
\pgfpathlineto{\pgfqpoint{2.671135in}{1.794864in}}%
\pgfpathlineto{\pgfqpoint{2.639742in}{1.810076in}}%
\pgfpathlineto{\pgfqpoint{2.608451in}{1.825234in}}%
\pgfpathclose%
\pgfusepath{stroke,fill}%
\end{pgfscope}%
\begin{pgfscope}%
\pgfpathrectangle{\pgfqpoint{0.050000in}{0.050000in}}{\pgfqpoint{4.900000in}{4.900000in}}%
\pgfusepath{clip}%
\pgfsetbuttcap%
\pgfsetroundjoin%
\definecolor{currentfill}{rgb}{0.773256,0.763906,0.864421}%
\pgfsetfillcolor{currentfill}%
\pgfsetfillopacity{0.900000}%
\pgfsetlinewidth{0.200750pt}%
\definecolor{currentstroke}{rgb}{0.200000,0.200000,0.200000}%
\pgfsetstrokecolor{currentstroke}%
\pgfsetdash{}{0pt}%
\pgfpathmoveto{\pgfqpoint{3.484026in}{2.193648in}}%
\pgfpathlineto{\pgfqpoint{3.525880in}{2.102828in}}%
\pgfpathlineto{\pgfqpoint{3.567506in}{2.013387in}}%
\pgfpathlineto{\pgfqpoint{3.599830in}{1.975392in}}%
\pgfpathlineto{\pgfqpoint{3.632251in}{1.938292in}}%
\pgfpathlineto{\pgfqpoint{3.590529in}{2.017661in}}%
\pgfpathlineto{\pgfqpoint{3.548643in}{2.101098in}}%
\pgfpathlineto{\pgfqpoint{3.516290in}{2.146800in}}%
\pgfpathlineto{\pgfqpoint{3.484026in}{2.193648in}}%
\pgfpathclose%
\pgfusepath{stroke,fill}%
\end{pgfscope}%
\begin{pgfscope}%
\pgfpathrectangle{\pgfqpoint{0.050000in}{0.050000in}}{\pgfqpoint{4.900000in}{4.900000in}}%
\pgfusepath{clip}%
\pgfsetbuttcap%
\pgfsetroundjoin%
\definecolor{currentfill}{rgb}{0.958231,0.956509,0.975025}%
\pgfsetfillcolor{currentfill}%
\pgfsetfillopacity{0.900000}%
\pgfsetlinewidth{0.200750pt}%
\definecolor{currentstroke}{rgb}{0.200000,0.200000,0.200000}%
\pgfsetstrokecolor{currentstroke}%
\pgfsetdash{}{0pt}%
\pgfpathmoveto{\pgfqpoint{3.697382in}{1.866505in}}%
\pgfpathlineto{\pgfqpoint{3.739333in}{1.812939in}}%
\pgfpathlineto{\pgfqpoint{3.781651in}{1.785260in}}%
\pgfpathlineto{\pgfqpoint{3.814558in}{1.760423in}}%
\pgfpathlineto{\pgfqpoint{3.847620in}{1.738204in}}%
\pgfpathlineto{\pgfqpoint{3.805029in}{1.753690in}}%
\pgfpathlineto{\pgfqpoint{3.762905in}{1.797545in}}%
\pgfpathlineto{\pgfqpoint{3.730094in}{1.831698in}}%
\pgfpathlineto{\pgfqpoint{3.697382in}{1.866505in}}%
\pgfpathclose%
\pgfusepath{stroke,fill}%
\end{pgfscope}%
\begin{pgfscope}%
\pgfpathrectangle{\pgfqpoint{0.050000in}{0.050000in}}{\pgfqpoint{4.900000in}{4.900000in}}%
\pgfusepath{clip}%
\pgfsetbuttcap%
\pgfsetroundjoin%
\definecolor{currentfill}{rgb}{0.337670,0.310358,0.603968}%
\pgfsetfillcolor{currentfill}%
\pgfsetfillopacity{0.900000}%
\pgfsetlinewidth{0.200750pt}%
\definecolor{currentstroke}{rgb}{0.200000,0.200000,0.200000}%
\pgfsetstrokecolor{currentstroke}%
\pgfsetdash{}{0pt}%
\pgfpathmoveto{\pgfqpoint{2.902176in}{2.610705in}}%
\pgfpathlineto{\pgfqpoint{2.946957in}{2.725318in}}%
\pgfpathlineto{\pgfqpoint{2.991621in}{2.815605in}}%
\pgfpathlineto{\pgfqpoint{3.023316in}{2.772467in}}%
\pgfpathlineto{\pgfqpoint{3.055079in}{2.728915in}}%
\pgfpathlineto{\pgfqpoint{3.010436in}{2.655926in}}%
\pgfpathlineto{\pgfqpoint{2.965610in}{2.557055in}}%
\pgfpathlineto{\pgfqpoint{2.933852in}{2.584459in}}%
\pgfpathlineto{\pgfqpoint{2.902176in}{2.610705in}}%
\pgfpathclose%
\pgfusepath{stroke,fill}%
\end{pgfscope}%
\begin{pgfscope}%
\pgfpathrectangle{\pgfqpoint{0.050000in}{0.050000in}}{\pgfqpoint{4.900000in}{4.900000in}}%
\pgfusepath{clip}%
\pgfsetbuttcap%
\pgfsetroundjoin%
\definecolor{currentfill}{rgb}{1.000000,1.000000,1.000000}%
\pgfsetfillcolor{currentfill}%
\pgfsetfillopacity{0.900000}%
\pgfsetlinewidth{0.200750pt}%
\definecolor{currentstroke}{rgb}{0.200000,0.200000,0.200000}%
\pgfsetstrokecolor{currentstroke}%
\pgfsetdash{}{0pt}%
\pgfpathmoveto{\pgfqpoint{1.191082in}{1.738037in}}%
\pgfpathlineto{\pgfqpoint{1.237070in}{1.752095in}}%
\pgfpathlineto{\pgfqpoint{1.282960in}{1.766123in}}%
\pgfpathlineto{\pgfqpoint{1.312192in}{1.746080in}}%
\pgfpathlineto{\pgfqpoint{1.341514in}{1.725975in}}%
\pgfpathlineto{\pgfqpoint{1.295546in}{1.711861in}}%
\pgfpathlineto{\pgfqpoint{1.249480in}{1.697717in}}%
\pgfpathlineto{\pgfqpoint{1.220236in}{1.717908in}}%
\pgfpathlineto{\pgfqpoint{1.191082in}{1.738037in}}%
\pgfpathclose%
\pgfusepath{stroke,fill}%
\end{pgfscope}%
\begin{pgfscope}%
\pgfpathrectangle{\pgfqpoint{0.050000in}{0.050000in}}{\pgfqpoint{4.900000in}{4.900000in}}%
\pgfusepath{clip}%
\pgfsetbuttcap%
\pgfsetroundjoin%
\definecolor{currentfill}{rgb}{1.000000,1.000000,1.000000}%
\pgfsetfillcolor{currentfill}%
\pgfsetfillopacity{0.900000}%
\pgfsetlinewidth{0.200750pt}%
\definecolor{currentstroke}{rgb}{0.200000,0.200000,0.200000}%
\pgfsetstrokecolor{currentstroke}%
\pgfsetdash{}{0pt}%
\pgfpathmoveto{\pgfqpoint{2.126985in}{1.734109in}}%
\pgfpathlineto{\pgfqpoint{2.171988in}{1.748256in}}%
\pgfpathlineto{\pgfqpoint{2.216894in}{1.762522in}}%
\pgfpathlineto{\pgfqpoint{2.247558in}{1.742601in}}%
\pgfpathlineto{\pgfqpoint{2.278316in}{1.722652in}}%
\pgfpathlineto{\pgfqpoint{2.233339in}{1.708123in}}%
\pgfpathlineto{\pgfqpoint{2.188265in}{1.693812in}}%
\pgfpathlineto{\pgfqpoint{2.157578in}{1.713988in}}%
\pgfpathlineto{\pgfqpoint{2.126985in}{1.734109in}}%
\pgfpathclose%
\pgfusepath{stroke,fill}%
\end{pgfscope}%
\begin{pgfscope}%
\pgfpathrectangle{\pgfqpoint{0.050000in}{0.050000in}}{\pgfqpoint{4.900000in}{4.900000in}}%
\pgfusepath{clip}%
\pgfsetbuttcap%
\pgfsetroundjoin%
\definecolor{currentfill}{rgb}{1.000000,1.000000,1.000000}%
\pgfsetfillcolor{currentfill}%
\pgfsetfillopacity{0.900000}%
\pgfsetlinewidth{0.200750pt}%
\definecolor{currentstroke}{rgb}{0.200000,0.200000,0.200000}%
\pgfsetstrokecolor{currentstroke}%
\pgfsetdash{}{0pt}%
\pgfpathmoveto{\pgfqpoint{0.798297in}{1.733991in}}%
\pgfpathlineto{\pgfqpoint{0.844719in}{1.748058in}}%
\pgfpathlineto{\pgfqpoint{0.891041in}{1.762095in}}%
\pgfpathlineto{\pgfqpoint{0.919685in}{1.742039in}}%
\pgfpathlineto{\pgfqpoint{0.948416in}{1.721922in}}%
\pgfpathlineto{\pgfqpoint{0.902014in}{1.707799in}}%
\pgfpathlineto{\pgfqpoint{0.855512in}{1.693646in}}%
\pgfpathlineto{\pgfqpoint{0.826861in}{1.713849in}}%
\pgfpathlineto{\pgfqpoint{0.798297in}{1.733991in}}%
\pgfpathclose%
\pgfusepath{stroke,fill}%
\end{pgfscope}%
\begin{pgfscope}%
\pgfpathrectangle{\pgfqpoint{0.050000in}{0.050000in}}{\pgfqpoint{4.900000in}{4.900000in}}%
\pgfusepath{clip}%
\pgfsetbuttcap%
\pgfsetroundjoin%
\definecolor{currentfill}{rgb}{0.671819,0.658285,0.803768}%
\pgfsetfillcolor{currentfill}%
\pgfsetfillopacity{0.900000}%
\pgfsetlinewidth{0.200750pt}%
\definecolor{currentstroke}{rgb}{0.200000,0.200000,0.200000}%
\pgfsetstrokecolor{currentstroke}%
\pgfsetdash{}{0pt}%
\pgfpathmoveto{\pgfqpoint{2.786468in}{2.103867in}}%
\pgfpathlineto{\pgfqpoint{2.831123in}{2.209139in}}%
\pgfpathlineto{\pgfqpoint{2.875882in}{2.324282in}}%
\pgfpathlineto{\pgfqpoint{2.907618in}{2.306888in}}%
\pgfpathlineto{\pgfqpoint{2.939447in}{2.288482in}}%
\pgfpathlineto{\pgfqpoint{2.894542in}{2.179701in}}%
\pgfpathlineto{\pgfqpoint{2.849726in}{2.077828in}}%
\pgfpathlineto{\pgfqpoint{2.818046in}{2.091187in}}%
\pgfpathlineto{\pgfqpoint{2.786468in}{2.103867in}}%
\pgfpathclose%
\pgfusepath{stroke,fill}%
\end{pgfscope}%
\begin{pgfscope}%
\pgfpathrectangle{\pgfqpoint{0.050000in}{0.050000in}}{\pgfqpoint{4.900000in}{4.900000in}}%
\pgfusepath{clip}%
\pgfsetbuttcap%
\pgfsetroundjoin%
\definecolor{currentfill}{rgb}{1.000000,1.000000,1.000000}%
\pgfsetfillcolor{currentfill}%
\pgfsetfillopacity{0.900000}%
\pgfsetlinewidth{0.200750pt}%
\definecolor{currentstroke}{rgb}{0.200000,0.200000,0.200000}%
\pgfsetstrokecolor{currentstroke}%
\pgfsetdash{}{0pt}%
\pgfpathmoveto{\pgfqpoint{3.847620in}{1.738204in}}%
\pgfpathlineto{\pgfqpoint{3.890774in}{1.752262in}}%
\pgfpathlineto{\pgfqpoint{3.924016in}{1.732176in}}%
\pgfpathlineto{\pgfqpoint{3.957360in}{1.712029in}}%
\pgfpathlineto{\pgfqpoint{3.914145in}{1.697885in}}%
\pgfpathlineto{\pgfqpoint{3.880832in}{1.718075in}}%
\pgfpathlineto{\pgfqpoint{3.847620in}{1.738204in}}%
\pgfpathclose%
\pgfusepath{stroke,fill}%
\end{pgfscope}%
\begin{pgfscope}%
\pgfpathrectangle{\pgfqpoint{0.050000in}{0.050000in}}{\pgfqpoint{4.900000in}{4.900000in}}%
\pgfusepath{clip}%
\pgfsetbuttcap%
\pgfsetroundjoin%
\definecolor{currentfill}{rgb}{1.000000,1.000000,1.000000}%
\pgfsetfillcolor{currentfill}%
\pgfsetfillopacity{0.900000}%
\pgfsetlinewidth{0.200750pt}%
\definecolor{currentstroke}{rgb}{0.200000,0.200000,0.200000}%
\pgfsetstrokecolor{currentstroke}%
\pgfsetdash{}{0pt}%
\pgfpathmoveto{\pgfqpoint{1.734370in}{1.730026in}}%
\pgfpathlineto{\pgfqpoint{1.779806in}{1.744101in}}%
\pgfpathlineto{\pgfqpoint{1.825146in}{1.758147in}}%
\pgfpathlineto{\pgfqpoint{1.855223in}{1.738079in}}%
\pgfpathlineto{\pgfqpoint{1.885392in}{1.717950in}}%
\pgfpathlineto{\pgfqpoint{1.839978in}{1.703818in}}%
\pgfpathlineto{\pgfqpoint{1.794468in}{1.689657in}}%
\pgfpathlineto{\pgfqpoint{1.764373in}{1.709872in}}%
\pgfpathlineto{\pgfqpoint{1.734370in}{1.730026in}}%
\pgfpathclose%
\pgfusepath{stroke,fill}%
\end{pgfscope}%
\begin{pgfscope}%
\pgfpathrectangle{\pgfqpoint{0.050000in}{0.050000in}}{\pgfqpoint{4.900000in}{4.900000in}}%
\pgfusepath{clip}%
\pgfsetbuttcap%
\pgfsetroundjoin%
\definecolor{currentfill}{rgb}{0.964198,0.962722,0.978593}%
\pgfsetfillcolor{currentfill}%
\pgfsetfillopacity{0.900000}%
\pgfsetlinewidth{0.200750pt}%
\definecolor{currentstroke}{rgb}{0.200000,0.200000,0.200000}%
\pgfsetstrokecolor{currentstroke}%
\pgfsetdash{}{0pt}%
\pgfpathmoveto{\pgfqpoint{2.519343in}{1.758116in}}%
\pgfpathlineto{\pgfqpoint{2.563929in}{1.787501in}}%
\pgfpathlineto{\pgfqpoint{2.608451in}{1.825234in}}%
\pgfpathlineto{\pgfqpoint{2.639742in}{1.810076in}}%
\pgfpathlineto{\pgfqpoint{2.671135in}{1.794864in}}%
\pgfpathlineto{\pgfqpoint{2.626507in}{1.754584in}}%
\pgfpathlineto{\pgfqpoint{2.581828in}{1.722911in}}%
\pgfpathlineto{\pgfqpoint{2.550536in}{1.740495in}}%
\pgfpathlineto{\pgfqpoint{2.519343in}{1.758116in}}%
\pgfpathclose%
\pgfusepath{stroke,fill}%
\end{pgfscope}%
\begin{pgfscope}%
\pgfpathrectangle{\pgfqpoint{0.050000in}{0.050000in}}{\pgfqpoint{4.900000in}{4.900000in}}%
\pgfusepath{clip}%
\pgfsetbuttcap%
\pgfsetroundjoin%
\definecolor{currentfill}{rgb}{1.000000,1.000000,1.000000}%
\pgfsetfillcolor{currentfill}%
\pgfsetfillopacity{0.900000}%
\pgfsetlinewidth{0.200750pt}%
\definecolor{currentstroke}{rgb}{0.200000,0.200000,0.200000}%
\pgfsetstrokecolor{currentstroke}%
\pgfsetdash{}{0pt}%
\pgfpathmoveto{\pgfqpoint{1.341514in}{1.725975in}}%
\pgfpathlineto{\pgfqpoint{1.387383in}{1.740059in}}%
\pgfpathlineto{\pgfqpoint{1.433155in}{1.754113in}}%
\pgfpathlineto{\pgfqpoint{1.462643in}{1.734033in}}%
\pgfpathlineto{\pgfqpoint{1.492222in}{1.713891in}}%
\pgfpathlineto{\pgfqpoint{1.446373in}{1.699751in}}%
\pgfpathlineto{\pgfqpoint{1.400427in}{1.685581in}}%
\pgfpathlineto{\pgfqpoint{1.370925in}{1.705809in}}%
\pgfpathlineto{\pgfqpoint{1.341514in}{1.725975in}}%
\pgfpathclose%
\pgfusepath{stroke,fill}%
\end{pgfscope}%
\begin{pgfscope}%
\pgfpathrectangle{\pgfqpoint{0.050000in}{0.050000in}}{\pgfqpoint{4.900000in}{4.900000in}}%
\pgfusepath{clip}%
\pgfsetbuttcap%
\pgfsetroundjoin%
\definecolor{currentfill}{rgb}{1.000000,1.000000,1.000000}%
\pgfsetfillcolor{currentfill}%
\pgfsetfillopacity{0.900000}%
\pgfsetlinewidth{0.200750pt}%
\definecolor{currentstroke}{rgb}{0.200000,0.200000,0.200000}%
\pgfsetstrokecolor{currentstroke}%
\pgfsetdash{}{0pt}%
\pgfpathmoveto{\pgfqpoint{0.948416in}{1.721922in}}%
\pgfpathlineto{\pgfqpoint{0.994720in}{1.736014in}}%
\pgfpathlineto{\pgfqpoint{1.040924in}{1.750077in}}%
\pgfpathlineto{\pgfqpoint{1.069823in}{1.729984in}}%
\pgfpathlineto{\pgfqpoint{1.098810in}{1.709831in}}%
\pgfpathlineto{\pgfqpoint{1.052526in}{1.695682in}}%
\pgfpathlineto{\pgfqpoint{1.006143in}{1.681503in}}%
\pgfpathlineto{\pgfqpoint{0.977236in}{1.701743in}}%
\pgfpathlineto{\pgfqpoint{0.948416in}{1.721922in}}%
\pgfpathclose%
\pgfusepath{stroke,fill}%
\end{pgfscope}%
\begin{pgfscope}%
\pgfpathrectangle{\pgfqpoint{0.050000in}{0.050000in}}{\pgfqpoint{4.900000in}{4.900000in}}%
\pgfusepath{clip}%
\pgfsetbuttcap%
\pgfsetroundjoin%
\definecolor{currentfill}{rgb}{1.000000,1.000000,1.000000}%
\pgfsetfillcolor{currentfill}%
\pgfsetfillopacity{0.900000}%
\pgfsetlinewidth{0.200750pt}%
\definecolor{currentstroke}{rgb}{0.200000,0.200000,0.200000}%
\pgfsetstrokecolor{currentstroke}%
\pgfsetdash{}{0pt}%
\pgfpathmoveto{\pgfqpoint{2.278316in}{1.722652in}}%
\pgfpathlineto{\pgfqpoint{2.323195in}{1.737697in}}%
\pgfpathlineto{\pgfqpoint{2.367976in}{1.753814in}}%
\pgfpathlineto{\pgfqpoint{2.398900in}{1.734522in}}%
\pgfpathlineto{\pgfqpoint{2.429918in}{1.715256in}}%
\pgfpathlineto{\pgfqpoint{2.385066in}{1.698267in}}%
\pgfpathlineto{\pgfqpoint{2.340115in}{1.682677in}}%
\pgfpathlineto{\pgfqpoint{2.309169in}{1.702676in}}%
\pgfpathlineto{\pgfqpoint{2.278316in}{1.722652in}}%
\pgfpathclose%
\pgfusepath{stroke,fill}%
\end{pgfscope}%
\begin{pgfscope}%
\pgfpathrectangle{\pgfqpoint{0.050000in}{0.050000in}}{\pgfqpoint{4.900000in}{4.900000in}}%
\pgfusepath{clip}%
\pgfsetbuttcap%
\pgfsetroundjoin%
\definecolor{currentfill}{rgb}{0.510711,0.490534,0.707436}%
\pgfsetfillcolor{currentfill}%
\pgfsetfillopacity{0.900000}%
\pgfsetlinewidth{0.200750pt}%
\definecolor{currentstroke}{rgb}{0.200000,0.200000,0.200000}%
\pgfsetstrokecolor{currentstroke}%
\pgfsetdash{}{0pt}%
\pgfpathmoveto{\pgfqpoint{3.270505in}{2.540199in}}%
\pgfpathlineto{\pgfqpoint{3.313474in}{2.481585in}}%
\pgfpathlineto{\pgfqpoint{3.355839in}{2.394224in}}%
\pgfpathlineto{\pgfqpoint{3.387758in}{2.341910in}}%
\pgfpathlineto{\pgfqpoint{3.419762in}{2.291111in}}%
\pgfpathlineto{\pgfqpoint{3.377457in}{2.375469in}}%
\pgfpathlineto{\pgfqpoint{3.334586in}{2.436054in}}%
\pgfpathlineto{\pgfqpoint{3.302507in}{2.487618in}}%
\pgfpathlineto{\pgfqpoint{3.270505in}{2.540199in}}%
\pgfpathclose%
\pgfusepath{stroke,fill}%
\end{pgfscope}%
\begin{pgfscope}%
\pgfpathrectangle{\pgfqpoint{0.050000in}{0.050000in}}{\pgfqpoint{4.900000in}{4.900000in}}%
\pgfusepath{clip}%
\pgfsetbuttcap%
\pgfsetroundjoin%
\definecolor{currentfill}{rgb}{0.313802,0.285506,0.589696}%
\pgfsetfillcolor{currentfill}%
\pgfsetfillopacity{0.900000}%
\pgfsetlinewidth{0.200750pt}%
\definecolor{currentstroke}{rgb}{0.200000,0.200000,0.200000}%
\pgfsetstrokecolor{currentstroke}%
\pgfsetdash{}{0pt}%
\pgfpathmoveto{\pgfqpoint{3.055079in}{2.728915in}}%
\pgfpathlineto{\pgfqpoint{3.099387in}{2.766287in}}%
\pgfpathlineto{\pgfqpoint{3.143209in}{2.761237in}}%
\pgfpathlineto{\pgfqpoint{3.174930in}{2.704329in}}%
\pgfpathlineto{\pgfqpoint{3.206717in}{2.648530in}}%
\pgfpathlineto{\pgfqpoint{3.163022in}{2.664334in}}%
\pgfpathlineto{\pgfqpoint{3.118808in}{2.641073in}}%
\pgfpathlineto{\pgfqpoint{3.086909in}{2.685079in}}%
\pgfpathlineto{\pgfqpoint{3.055079in}{2.728915in}}%
\pgfpathclose%
\pgfusepath{stroke,fill}%
\end{pgfscope}%
\begin{pgfscope}%
\pgfpathrectangle{\pgfqpoint{0.050000in}{0.050000in}}{\pgfqpoint{4.900000in}{4.900000in}}%
\pgfusepath{clip}%
\pgfsetbuttcap%
\pgfsetroundjoin%
\definecolor{currentfill}{rgb}{1.000000,1.000000,1.000000}%
\pgfsetfillcolor{currentfill}%
\pgfsetfillopacity{0.900000}%
\pgfsetlinewidth{0.200750pt}%
\definecolor{currentstroke}{rgb}{0.200000,0.200000,0.200000}%
\pgfsetstrokecolor{currentstroke}%
\pgfsetdash{}{0pt}%
\pgfpathmoveto{\pgfqpoint{1.885392in}{1.717950in}}%
\pgfpathlineto{\pgfqpoint{1.930708in}{1.732051in}}%
\pgfpathlineto{\pgfqpoint{1.975927in}{1.746122in}}%
\pgfpathlineto{\pgfqpoint{2.006262in}{1.726018in}}%
\pgfpathlineto{\pgfqpoint{2.036690in}{1.705853in}}%
\pgfpathlineto{\pgfqpoint{1.991397in}{1.691694in}}%
\pgfpathlineto{\pgfqpoint{1.946007in}{1.677506in}}%
\pgfpathlineto{\pgfqpoint{1.915653in}{1.697759in}}%
\pgfpathlineto{\pgfqpoint{1.885392in}{1.717950in}}%
\pgfpathclose%
\pgfusepath{stroke,fill}%
\end{pgfscope}%
\begin{pgfscope}%
\pgfpathrectangle{\pgfqpoint{0.050000in}{0.050000in}}{\pgfqpoint{4.900000in}{4.900000in}}%
\pgfusepath{clip}%
\pgfsetbuttcap%
\pgfsetroundjoin%
\definecolor{currentfill}{rgb}{0.510711,0.490534,0.707436}%
\pgfsetfillcolor{currentfill}%
\pgfsetfillopacity{0.900000}%
\pgfsetlinewidth{0.200750pt}%
\definecolor{currentstroke}{rgb}{0.200000,0.200000,0.200000}%
\pgfsetstrokecolor{currentstroke}%
\pgfsetdash{}{0pt}%
\pgfpathmoveto{\pgfqpoint{2.875882in}{2.324282in}}%
\pgfpathlineto{\pgfqpoint{2.920728in}{2.443096in}}%
\pgfpathlineto{\pgfqpoint{2.965610in}{2.557055in}}%
\pgfpathlineto{\pgfqpoint{2.997448in}{2.528600in}}%
\pgfpathlineto{\pgfqpoint{3.029365in}{2.499199in}}%
\pgfpathlineto{\pgfqpoint{2.984410in}{2.397746in}}%
\pgfpathlineto{\pgfqpoint{2.939447in}{2.288482in}}%
\pgfpathlineto{\pgfqpoint{2.907618in}{2.306888in}}%
\pgfpathlineto{\pgfqpoint{2.875882in}{2.324282in}}%
\pgfpathclose%
\pgfusepath{stroke,fill}%
\end{pgfscope}%
\begin{pgfscope}%
\pgfpathrectangle{\pgfqpoint{0.050000in}{0.050000in}}{\pgfqpoint{4.900000in}{4.900000in}}%
\pgfusepath{clip}%
\pgfsetbuttcap%
\pgfsetroundjoin%
\definecolor{currentfill}{rgb}{1.000000,1.000000,1.000000}%
\pgfsetfillcolor{currentfill}%
\pgfsetfillopacity{0.900000}%
\pgfsetlinewidth{0.200750pt}%
\definecolor{currentstroke}{rgb}{0.200000,0.200000,0.200000}%
\pgfsetstrokecolor{currentstroke}%
\pgfsetdash{}{0pt}%
\pgfpathmoveto{\pgfqpoint{1.492222in}{1.713891in}}%
\pgfpathlineto{\pgfqpoint{1.537972in}{1.728001in}}%
\pgfpathlineto{\pgfqpoint{1.583625in}{1.742081in}}%
\pgfpathlineto{\pgfqpoint{1.613370in}{1.721964in}}%
\pgfpathlineto{\pgfqpoint{1.643206in}{1.701785in}}%
\pgfpathlineto{\pgfqpoint{1.597478in}{1.687619in}}%
\pgfpathlineto{\pgfqpoint{1.551651in}{1.673423in}}%
\pgfpathlineto{\pgfqpoint{1.521891in}{1.693688in}}%
\pgfpathlineto{\pgfqpoint{1.492222in}{1.713891in}}%
\pgfpathclose%
\pgfusepath{stroke,fill}%
\end{pgfscope}%
\begin{pgfscope}%
\pgfpathrectangle{\pgfqpoint{0.050000in}{0.050000in}}{\pgfqpoint{4.900000in}{4.900000in}}%
\pgfusepath{clip}%
\pgfsetbuttcap%
\pgfsetroundjoin%
\definecolor{currentfill}{rgb}{0.809058,0.801184,0.885829}%
\pgfsetfillcolor{currentfill}%
\pgfsetfillopacity{0.900000}%
\pgfsetlinewidth{0.200750pt}%
\definecolor{currentstroke}{rgb}{0.200000,0.200000,0.200000}%
\pgfsetstrokecolor{currentstroke}%
\pgfsetdash{}{0pt}%
\pgfpathmoveto{\pgfqpoint{3.548643in}{2.101098in}}%
\pgfpathlineto{\pgfqpoint{3.590529in}{2.017661in}}%
\pgfpathlineto{\pgfqpoint{3.632251in}{1.938292in}}%
\pgfpathlineto{\pgfqpoint{3.664768in}{1.902017in}}%
\pgfpathlineto{\pgfqpoint{3.697382in}{1.866505in}}%
\pgfpathlineto{\pgfqpoint{3.655554in}{1.936517in}}%
\pgfpathlineto{\pgfqpoint{3.613620in}{2.012840in}}%
\pgfpathlineto{\pgfqpoint{3.581086in}{2.056467in}}%
\pgfpathlineto{\pgfqpoint{3.548643in}{2.101098in}}%
\pgfpathclose%
\pgfusepath{stroke,fill}%
\end{pgfscope}%
\begin{pgfscope}%
\pgfpathrectangle{\pgfqpoint{0.050000in}{0.050000in}}{\pgfqpoint{4.900000in}{4.900000in}}%
\pgfusepath{clip}%
\pgfsetbuttcap%
\pgfsetroundjoin%
\definecolor{currentfill}{rgb}{1.000000,1.000000,1.000000}%
\pgfsetfillcolor{currentfill}%
\pgfsetfillopacity{0.900000}%
\pgfsetlinewidth{0.200750pt}%
\definecolor{currentstroke}{rgb}{0.200000,0.200000,0.200000}%
\pgfsetstrokecolor{currentstroke}%
\pgfsetdash{}{0pt}%
\pgfpathmoveto{\pgfqpoint{1.098810in}{1.709831in}}%
\pgfpathlineto{\pgfqpoint{1.144995in}{1.723949in}}%
\pgfpathlineto{\pgfqpoint{1.191082in}{1.738037in}}%
\pgfpathlineto{\pgfqpoint{1.220236in}{1.717908in}}%
\pgfpathlineto{\pgfqpoint{1.249480in}{1.697717in}}%
\pgfpathlineto{\pgfqpoint{1.203315in}{1.683542in}}%
\pgfpathlineto{\pgfqpoint{1.157051in}{1.669337in}}%
\pgfpathlineto{\pgfqpoint{1.127886in}{1.689615in}}%
\pgfpathlineto{\pgfqpoint{1.098810in}{1.709831in}}%
\pgfpathclose%
\pgfusepath{stroke,fill}%
\end{pgfscope}%
\begin{pgfscope}%
\pgfpathrectangle{\pgfqpoint{0.050000in}{0.050000in}}{\pgfqpoint{4.900000in}{4.900000in}}%
\pgfusepath{clip}%
\pgfsetbuttcap%
\pgfsetroundjoin%
\definecolor{currentfill}{rgb}{0.988066,0.987574,0.992864}%
\pgfsetfillcolor{currentfill}%
\pgfsetfillopacity{0.900000}%
\pgfsetlinewidth{0.200750pt}%
\definecolor{currentstroke}{rgb}{0.200000,0.200000,0.200000}%
\pgfsetstrokecolor{currentstroke}%
\pgfsetdash{}{0pt}%
\pgfpathmoveto{\pgfqpoint{2.429918in}{1.715256in}}%
\pgfpathlineto{\pgfqpoint{2.474676in}{1.734700in}}%
\pgfpathlineto{\pgfqpoint{2.519343in}{1.758116in}}%
\pgfpathlineto{\pgfqpoint{2.550536in}{1.740495in}}%
\pgfpathlineto{\pgfqpoint{2.581828in}{1.722911in}}%
\pgfpathlineto{\pgfqpoint{2.537078in}{1.697633in}}%
\pgfpathlineto{\pgfqpoint{2.492243in}{1.676796in}}%
\pgfpathlineto{\pgfqpoint{2.461033in}{1.696015in}}%
\pgfpathlineto{\pgfqpoint{2.429918in}{1.715256in}}%
\pgfpathclose%
\pgfusepath{stroke,fill}%
\end{pgfscope}%
\begin{pgfscope}%
\pgfpathrectangle{\pgfqpoint{0.050000in}{0.050000in}}{\pgfqpoint{4.900000in}{4.900000in}}%
\pgfusepath{clip}%
\pgfsetbuttcap%
\pgfsetroundjoin%
\definecolor{currentfill}{rgb}{1.000000,1.000000,1.000000}%
\pgfsetfillcolor{currentfill}%
\pgfsetfillopacity{0.900000}%
\pgfsetlinewidth{0.200750pt}%
\definecolor{currentstroke}{rgb}{0.200000,0.200000,0.200000}%
\pgfsetstrokecolor{currentstroke}%
\pgfsetdash{}{0pt}%
\pgfpathmoveto{\pgfqpoint{2.036690in}{1.705853in}}%
\pgfpathlineto{\pgfqpoint{2.081886in}{1.719987in}}%
\pgfpathlineto{\pgfqpoint{2.126985in}{1.734109in}}%
\pgfpathlineto{\pgfqpoint{2.157578in}{1.713988in}}%
\pgfpathlineto{\pgfqpoint{2.188265in}{1.693812in}}%
\pgfpathlineto{\pgfqpoint{2.143094in}{1.679572in}}%
\pgfpathlineto{\pgfqpoint{2.097825in}{1.665340in}}%
\pgfpathlineto{\pgfqpoint{2.067210in}{1.685627in}}%
\pgfpathlineto{\pgfqpoint{2.036690in}{1.705853in}}%
\pgfpathclose%
\pgfusepath{stroke,fill}%
\end{pgfscope}%
\begin{pgfscope}%
\pgfpathrectangle{\pgfqpoint{0.050000in}{0.050000in}}{\pgfqpoint{4.900000in}{4.900000in}}%
\pgfusepath{clip}%
\pgfsetbuttcap%
\pgfsetroundjoin%
\definecolor{currentfill}{rgb}{0.367505,0.341423,0.621807}%
\pgfsetfillcolor{currentfill}%
\pgfsetfillopacity{0.900000}%
\pgfsetlinewidth{0.200750pt}%
\definecolor{currentstroke}{rgb}{0.200000,0.200000,0.200000}%
\pgfsetstrokecolor{currentstroke}%
\pgfsetdash{}{0pt}%
\pgfpathmoveto{\pgfqpoint{2.965610in}{2.557055in}}%
\pgfpathlineto{\pgfqpoint{3.010436in}{2.655926in}}%
\pgfpathlineto{\pgfqpoint{3.055079in}{2.728915in}}%
\pgfpathlineto{\pgfqpoint{3.086909in}{2.685079in}}%
\pgfpathlineto{\pgfqpoint{3.118808in}{2.641073in}}%
\pgfpathlineto{\pgfqpoint{3.074208in}{2.583427in}}%
\pgfpathlineto{\pgfqpoint{3.029365in}{2.499199in}}%
\pgfpathlineto{\pgfqpoint{2.997448in}{2.528600in}}%
\pgfpathlineto{\pgfqpoint{2.965610in}{2.557055in}}%
\pgfpathclose%
\pgfusepath{stroke,fill}%
\end{pgfscope}%
\begin{pgfscope}%
\pgfpathrectangle{\pgfqpoint{0.050000in}{0.050000in}}{\pgfqpoint{4.900000in}{4.900000in}}%
\pgfusepath{clip}%
\pgfsetbuttcap%
\pgfsetroundjoin%
\definecolor{currentfill}{rgb}{1.000000,1.000000,1.000000}%
\pgfsetfillcolor{currentfill}%
\pgfsetfillopacity{0.900000}%
\pgfsetlinewidth{0.200750pt}%
\definecolor{currentstroke}{rgb}{0.200000,0.200000,0.200000}%
\pgfsetstrokecolor{currentstroke}%
\pgfsetdash{}{0pt}%
\pgfpathmoveto{\pgfqpoint{0.705157in}{1.705767in}}%
\pgfpathlineto{\pgfqpoint{0.751777in}{1.719894in}}%
\pgfpathlineto{\pgfqpoint{0.798297in}{1.733991in}}%
\pgfpathlineto{\pgfqpoint{0.826861in}{1.713849in}}%
\pgfpathlineto{\pgfqpoint{0.855512in}{1.693646in}}%
\pgfpathlineto{\pgfqpoint{0.808910in}{1.679463in}}%
\pgfpathlineto{\pgfqpoint{0.762208in}{1.665249in}}%
\pgfpathlineto{\pgfqpoint{0.733639in}{1.685539in}}%
\pgfpathlineto{\pgfqpoint{0.705157in}{1.705767in}}%
\pgfpathclose%
\pgfusepath{stroke,fill}%
\end{pgfscope}%
\begin{pgfscope}%
\pgfpathrectangle{\pgfqpoint{0.050000in}{0.050000in}}{\pgfqpoint{4.900000in}{4.900000in}}%
\pgfusepath{clip}%
\pgfsetbuttcap%
\pgfsetroundjoin%
\definecolor{currentfill}{rgb}{0.976132,0.975148,0.985729}%
\pgfsetfillcolor{currentfill}%
\pgfsetfillopacity{0.900000}%
\pgfsetlinewidth{0.200750pt}%
\definecolor{currentstroke}{rgb}{0.200000,0.200000,0.200000}%
\pgfsetstrokecolor{currentstroke}%
\pgfsetdash{}{0pt}%
\pgfpathmoveto{\pgfqpoint{3.762905in}{1.797545in}}%
\pgfpathlineto{\pgfqpoint{3.805029in}{1.753690in}}%
\pgfpathlineto{\pgfqpoint{3.847620in}{1.738204in}}%
\pgfpathlineto{\pgfqpoint{3.880832in}{1.718075in}}%
\pgfpathlineto{\pgfqpoint{3.914145in}{1.697885in}}%
\pgfpathlineto{\pgfqpoint{3.871126in}{1.696063in}}%
\pgfpathlineto{\pgfqpoint{3.828826in}{1.731015in}}%
\pgfpathlineto{\pgfqpoint{3.795816in}{1.763999in}}%
\pgfpathlineto{\pgfqpoint{3.762905in}{1.797545in}}%
\pgfpathclose%
\pgfusepath{stroke,fill}%
\end{pgfscope}%
\begin{pgfscope}%
\pgfpathrectangle{\pgfqpoint{0.050000in}{0.050000in}}{\pgfqpoint{4.900000in}{4.900000in}}%
\pgfusepath{clip}%
\pgfsetbuttcap%
\pgfsetroundjoin%
\definecolor{currentfill}{rgb}{1.000000,1.000000,1.000000}%
\pgfsetfillcolor{currentfill}%
\pgfsetfillopacity{0.900000}%
\pgfsetlinewidth{0.200750pt}%
\definecolor{currentstroke}{rgb}{0.200000,0.200000,0.200000}%
\pgfsetstrokecolor{currentstroke}%
\pgfsetdash{}{0pt}%
\pgfpathmoveto{\pgfqpoint{1.643206in}{1.701785in}}%
\pgfpathlineto{\pgfqpoint{1.688837in}{1.715921in}}%
\pgfpathlineto{\pgfqpoint{1.734370in}{1.730026in}}%
\pgfpathlineto{\pgfqpoint{1.764373in}{1.709872in}}%
\pgfpathlineto{\pgfqpoint{1.794468in}{1.689657in}}%
\pgfpathlineto{\pgfqpoint{1.748859in}{1.675465in}}%
\pgfpathlineto{\pgfqpoint{1.703153in}{1.661243in}}%
\pgfpathlineto{\pgfqpoint{1.673134in}{1.681545in}}%
\pgfpathlineto{\pgfqpoint{1.643206in}{1.701785in}}%
\pgfpathclose%
\pgfusepath{stroke,fill}%
\end{pgfscope}%
\begin{pgfscope}%
\pgfpathrectangle{\pgfqpoint{0.050000in}{0.050000in}}{\pgfqpoint{4.900000in}{4.900000in}}%
\pgfusepath{clip}%
\pgfsetbuttcap%
\pgfsetroundjoin%
\definecolor{currentfill}{rgb}{0.803091,0.794971,0.882261}%
\pgfsetfillcolor{currentfill}%
\pgfsetfillopacity{0.900000}%
\pgfsetlinewidth{0.200750pt}%
\definecolor{currentstroke}{rgb}{0.200000,0.200000,0.200000}%
\pgfsetstrokecolor{currentstroke}%
\pgfsetdash{}{0pt}%
\pgfpathmoveto{\pgfqpoint{2.760349in}{1.909715in}}%
\pgfpathlineto{\pgfqpoint{2.805001in}{1.987118in}}%
\pgfpathlineto{\pgfqpoint{2.849726in}{2.077828in}}%
\pgfpathlineto{\pgfqpoint{2.881506in}{2.063787in}}%
\pgfpathlineto{\pgfqpoint{2.913386in}{2.049067in}}%
\pgfpathlineto{\pgfqpoint{2.868507in}{1.960058in}}%
\pgfpathlineto{\pgfqpoint{2.823705in}{1.882511in}}%
\pgfpathlineto{\pgfqpoint{2.791975in}{1.896282in}}%
\pgfpathlineto{\pgfqpoint{2.760349in}{1.909715in}}%
\pgfpathclose%
\pgfusepath{stroke,fill}%
\end{pgfscope}%
\begin{pgfscope}%
\pgfpathrectangle{\pgfqpoint{0.050000in}{0.050000in}}{\pgfqpoint{4.900000in}{4.900000in}}%
\pgfusepath{clip}%
\pgfsetbuttcap%
\pgfsetroundjoin%
\definecolor{currentfill}{rgb}{0.904529,0.900592,0.942914}%
\pgfsetfillcolor{currentfill}%
\pgfsetfillopacity{0.900000}%
\pgfsetlinewidth{0.200750pt}%
\definecolor{currentstroke}{rgb}{0.200000,0.200000,0.200000}%
\pgfsetstrokecolor{currentstroke}%
\pgfsetdash{}{0pt}%
\pgfpathmoveto{\pgfqpoint{2.671135in}{1.794864in}}%
\pgfpathlineto{\pgfqpoint{2.715738in}{1.845962in}}%
\pgfpathlineto{\pgfqpoint{2.760349in}{1.909715in}}%
\pgfpathlineto{\pgfqpoint{2.791975in}{1.896282in}}%
\pgfpathlineto{\pgfqpoint{2.823705in}{1.882511in}}%
\pgfpathlineto{\pgfqpoint{2.778954in}{1.817371in}}%
\pgfpathlineto{\pgfqpoint{2.734226in}{1.764206in}}%
\pgfpathlineto{\pgfqpoint{2.702629in}{1.779580in}}%
\pgfpathlineto{\pgfqpoint{2.671135in}{1.794864in}}%
\pgfpathclose%
\pgfusepath{stroke,fill}%
\end{pgfscope}%
\begin{pgfscope}%
\pgfpathrectangle{\pgfqpoint{0.050000in}{0.050000in}}{\pgfqpoint{4.900000in}{4.900000in}}%
\pgfusepath{clip}%
\pgfsetbuttcap%
\pgfsetroundjoin%
\definecolor{currentfill}{rgb}{1.000000,1.000000,1.000000}%
\pgfsetfillcolor{currentfill}%
\pgfsetfillopacity{0.900000}%
\pgfsetlinewidth{0.200750pt}%
\definecolor{currentstroke}{rgb}{0.200000,0.200000,0.200000}%
\pgfsetstrokecolor{currentstroke}%
\pgfsetdash{}{0pt}%
\pgfpathmoveto{\pgfqpoint{1.249480in}{1.697717in}}%
\pgfpathlineto{\pgfqpoint{1.295546in}{1.711861in}}%
\pgfpathlineto{\pgfqpoint{1.341514in}{1.725975in}}%
\pgfpathlineto{\pgfqpoint{1.370925in}{1.705809in}}%
\pgfpathlineto{\pgfqpoint{1.400427in}{1.685581in}}%
\pgfpathlineto{\pgfqpoint{1.354381in}{1.671381in}}%
\pgfpathlineto{\pgfqpoint{1.308237in}{1.657149in}}%
\pgfpathlineto{\pgfqpoint{1.278814in}{1.677464in}}%
\pgfpathlineto{\pgfqpoint{1.249480in}{1.697717in}}%
\pgfpathclose%
\pgfusepath{stroke,fill}%
\end{pgfscope}%
\begin{pgfscope}%
\pgfpathrectangle{\pgfqpoint{0.050000in}{0.050000in}}{\pgfqpoint{4.900000in}{4.900000in}}%
\pgfusepath{clip}%
\pgfsetbuttcap%
\pgfsetroundjoin%
\definecolor{currentfill}{rgb}{0.564414,0.546451,0.739546}%
\pgfsetfillcolor{currentfill}%
\pgfsetfillopacity{0.900000}%
\pgfsetlinewidth{0.200750pt}%
\definecolor{currentstroke}{rgb}{0.200000,0.200000,0.200000}%
\pgfsetstrokecolor{currentstroke}%
\pgfsetdash{}{0pt}%
\pgfpathmoveto{\pgfqpoint{3.334586in}{2.436054in}}%
\pgfpathlineto{\pgfqpoint{3.377457in}{2.375469in}}%
\pgfpathlineto{\pgfqpoint{3.419762in}{2.291111in}}%
\pgfpathlineto{\pgfqpoint{3.451851in}{2.241722in}}%
\pgfpathlineto{\pgfqpoint{3.484026in}{2.193648in}}%
\pgfpathlineto{\pgfqpoint{3.441763in}{2.274455in}}%
\pgfpathlineto{\pgfqpoint{3.398975in}{2.335843in}}%
\pgfpathlineto{\pgfqpoint{3.366741in}{2.385474in}}%
\pgfpathlineto{\pgfqpoint{3.334586in}{2.436054in}}%
\pgfpathclose%
\pgfusepath{stroke,fill}%
\end{pgfscope}%
\begin{pgfscope}%
\pgfpathrectangle{\pgfqpoint{0.050000in}{0.050000in}}{\pgfqpoint{4.900000in}{4.900000in}}%
\pgfusepath{clip}%
\pgfsetbuttcap%
\pgfsetroundjoin%
\definecolor{currentfill}{rgb}{1.000000,1.000000,1.000000}%
\pgfsetfillcolor{currentfill}%
\pgfsetfillopacity{0.900000}%
\pgfsetlinewidth{0.200750pt}%
\definecolor{currentstroke}{rgb}{0.200000,0.200000,0.200000}%
\pgfsetstrokecolor{currentstroke}%
\pgfsetdash{}{0pt}%
\pgfpathmoveto{\pgfqpoint{0.855512in}{1.693646in}}%
\pgfpathlineto{\pgfqpoint{0.902014in}{1.707799in}}%
\pgfpathlineto{\pgfqpoint{0.948416in}{1.721922in}}%
\pgfpathlineto{\pgfqpoint{0.977236in}{1.701743in}}%
\pgfpathlineto{\pgfqpoint{1.006143in}{1.681503in}}%
\pgfpathlineto{\pgfqpoint{0.959660in}{1.667294in}}%
\pgfpathlineto{\pgfqpoint{0.913078in}{1.653054in}}%
\pgfpathlineto{\pgfqpoint{0.884251in}{1.673381in}}%
\pgfpathlineto{\pgfqpoint{0.855512in}{1.693646in}}%
\pgfpathclose%
\pgfusepath{stroke,fill}%
\end{pgfscope}%
\begin{pgfscope}%
\pgfpathrectangle{\pgfqpoint{0.050000in}{0.050000in}}{\pgfqpoint{4.900000in}{4.900000in}}%
\pgfusepath{clip}%
\pgfsetbuttcap%
\pgfsetroundjoin%
\definecolor{currentfill}{rgb}{1.000000,1.000000,1.000000}%
\pgfsetfillcolor{currentfill}%
\pgfsetfillopacity{0.900000}%
\pgfsetlinewidth{0.200750pt}%
\definecolor{currentstroke}{rgb}{0.200000,0.200000,0.200000}%
\pgfsetstrokecolor{currentstroke}%
\pgfsetdash{}{0pt}%
\pgfpathmoveto{\pgfqpoint{2.188265in}{1.693812in}}%
\pgfpathlineto{\pgfqpoint{2.233339in}{1.708123in}}%
\pgfpathlineto{\pgfqpoint{2.278316in}{1.722652in}}%
\pgfpathlineto{\pgfqpoint{2.309169in}{1.702676in}}%
\pgfpathlineto{\pgfqpoint{2.340115in}{1.682677in}}%
\pgfpathlineto{\pgfqpoint{2.295067in}{1.667818in}}%
\pgfpathlineto{\pgfqpoint{2.249921in}{1.653302in}}%
\pgfpathlineto{\pgfqpoint{2.219046in}{1.673583in}}%
\pgfpathlineto{\pgfqpoint{2.188265in}{1.693812in}}%
\pgfpathclose%
\pgfusepath{stroke,fill}%
\end{pgfscope}%
\begin{pgfscope}%
\pgfpathrectangle{\pgfqpoint{0.050000in}{0.050000in}}{\pgfqpoint{4.900000in}{4.900000in}}%
\pgfusepath{clip}%
\pgfsetbuttcap%
\pgfsetroundjoin%
\definecolor{currentfill}{rgb}{0.367505,0.341423,0.621807}%
\pgfsetfillcolor{currentfill}%
\pgfsetfillopacity{0.900000}%
\pgfsetlinewidth{0.200750pt}%
\definecolor{currentstroke}{rgb}{0.200000,0.200000,0.200000}%
\pgfsetstrokecolor{currentstroke}%
\pgfsetdash{}{0pt}%
\pgfpathmoveto{\pgfqpoint{3.118808in}{2.641073in}}%
\pgfpathlineto{\pgfqpoint{3.163022in}{2.664334in}}%
\pgfpathlineto{\pgfqpoint{3.206717in}{2.648530in}}%
\pgfpathlineto{\pgfqpoint{3.238575in}{2.593827in}}%
\pgfpathlineto{\pgfqpoint{3.270505in}{2.540199in}}%
\pgfpathlineto{\pgfqpoint{3.226929in}{2.564372in}}%
\pgfpathlineto{\pgfqpoint{3.182815in}{2.552912in}}%
\pgfpathlineto{\pgfqpoint{3.150776in}{2.596991in}}%
\pgfpathlineto{\pgfqpoint{3.118808in}{2.641073in}}%
\pgfpathclose%
\pgfusepath{stroke,fill}%
\end{pgfscope}%
\begin{pgfscope}%
\pgfpathrectangle{\pgfqpoint{0.050000in}{0.050000in}}{\pgfqpoint{4.900000in}{4.900000in}}%
\pgfusepath{clip}%
\pgfsetbuttcap%
\pgfsetroundjoin%
\definecolor{currentfill}{rgb}{0.665852,0.652072,0.800200}%
\pgfsetfillcolor{currentfill}%
\pgfsetfillopacity{0.900000}%
\pgfsetlinewidth{0.200750pt}%
\definecolor{currentstroke}{rgb}{0.200000,0.200000,0.200000}%
\pgfsetstrokecolor{currentstroke}%
\pgfsetdash{}{0pt}%
\pgfpathmoveto{\pgfqpoint{2.849726in}{2.077828in}}%
\pgfpathlineto{\pgfqpoint{2.894542in}{2.179701in}}%
\pgfpathlineto{\pgfqpoint{2.939447in}{2.288482in}}%
\pgfpathlineto{\pgfqpoint{2.971368in}{2.269108in}}%
\pgfpathlineto{\pgfqpoint{3.003381in}{2.248810in}}%
\pgfpathlineto{\pgfqpoint{2.958349in}{2.146855in}}%
\pgfpathlineto{\pgfqpoint{2.913386in}{2.049067in}}%
\pgfpathlineto{\pgfqpoint{2.881506in}{2.063787in}}%
\pgfpathlineto{\pgfqpoint{2.849726in}{2.077828in}}%
\pgfpathclose%
\pgfusepath{stroke,fill}%
\end{pgfscope}%
\begin{pgfscope}%
\pgfpathrectangle{\pgfqpoint{0.050000in}{0.050000in}}{\pgfqpoint{4.900000in}{4.900000in}}%
\pgfusepath{clip}%
\pgfsetbuttcap%
\pgfsetroundjoin%
\definecolor{currentfill}{rgb}{1.000000,1.000000,1.000000}%
\pgfsetfillcolor{currentfill}%
\pgfsetfillopacity{0.900000}%
\pgfsetlinewidth{0.200750pt}%
\definecolor{currentstroke}{rgb}{0.200000,0.200000,0.200000}%
\pgfsetstrokecolor{currentstroke}%
\pgfsetdash{}{0pt}%
\pgfpathmoveto{\pgfqpoint{3.914145in}{1.697885in}}%
\pgfpathlineto{\pgfqpoint{3.957360in}{1.712029in}}%
\pgfpathlineto{\pgfqpoint{3.990806in}{1.691820in}}%
\pgfpathlineto{\pgfqpoint{4.024354in}{1.671549in}}%
\pgfpathlineto{\pgfqpoint{3.981079in}{1.657318in}}%
\pgfpathlineto{\pgfqpoint{3.947561in}{1.677633in}}%
\pgfpathlineto{\pgfqpoint{3.914145in}{1.697885in}}%
\pgfpathclose%
\pgfusepath{stroke,fill}%
\end{pgfscope}%
\begin{pgfscope}%
\pgfpathrectangle{\pgfqpoint{0.050000in}{0.050000in}}{\pgfqpoint{4.900000in}{4.900000in}}%
\pgfusepath{clip}%
\pgfsetbuttcap%
\pgfsetroundjoin%
\definecolor{currentfill}{rgb}{1.000000,1.000000,1.000000}%
\pgfsetfillcolor{currentfill}%
\pgfsetfillopacity{0.900000}%
\pgfsetlinewidth{0.200750pt}%
\definecolor{currentstroke}{rgb}{0.200000,0.200000,0.200000}%
\pgfsetstrokecolor{currentstroke}%
\pgfsetdash{}{0pt}%
\pgfpathmoveto{\pgfqpoint{1.794468in}{1.689657in}}%
\pgfpathlineto{\pgfqpoint{1.839978in}{1.703818in}}%
\pgfpathlineto{\pgfqpoint{1.885392in}{1.717950in}}%
\pgfpathlineto{\pgfqpoint{1.915653in}{1.697759in}}%
\pgfpathlineto{\pgfqpoint{1.946007in}{1.677506in}}%
\pgfpathlineto{\pgfqpoint{1.900519in}{1.663288in}}%
\pgfpathlineto{\pgfqpoint{1.854934in}{1.649040in}}%
\pgfpathlineto{\pgfqpoint{1.824654in}{1.669380in}}%
\pgfpathlineto{\pgfqpoint{1.794468in}{1.689657in}}%
\pgfpathclose%
\pgfusepath{stroke,fill}%
\end{pgfscope}%
\begin{pgfscope}%
\pgfpathrectangle{\pgfqpoint{0.050000in}{0.050000in}}{\pgfqpoint{4.900000in}{4.900000in}}%
\pgfusepath{clip}%
\pgfsetbuttcap%
\pgfsetroundjoin%
\definecolor{currentfill}{rgb}{0.958231,0.956509,0.975025}%
\pgfsetfillcolor{currentfill}%
\pgfsetfillopacity{0.900000}%
\pgfsetlinewidth{0.200750pt}%
\definecolor{currentstroke}{rgb}{0.200000,0.200000,0.200000}%
\pgfsetstrokecolor{currentstroke}%
\pgfsetdash{}{0pt}%
\pgfpathmoveto{\pgfqpoint{2.581828in}{1.722911in}}%
\pgfpathlineto{\pgfqpoint{2.626507in}{1.754584in}}%
\pgfpathlineto{\pgfqpoint{2.671135in}{1.794864in}}%
\pgfpathlineto{\pgfqpoint{2.702629in}{1.779580in}}%
\pgfpathlineto{\pgfqpoint{2.734226in}{1.764206in}}%
\pgfpathlineto{\pgfqpoint{2.689487in}{1.721644in}}%
\pgfpathlineto{\pgfqpoint{2.644709in}{1.687809in}}%
\pgfpathlineto{\pgfqpoint{2.613219in}{1.705352in}}%
\pgfpathlineto{\pgfqpoint{2.581828in}{1.722911in}}%
\pgfpathclose%
\pgfusepath{stroke,fill}%
\end{pgfscope}%
\begin{pgfscope}%
\pgfpathrectangle{\pgfqpoint{0.050000in}{0.050000in}}{\pgfqpoint{4.900000in}{4.900000in}}%
\pgfusepath{clip}%
\pgfsetbuttcap%
\pgfsetroundjoin%
\definecolor{currentfill}{rgb}{1.000000,1.000000,1.000000}%
\pgfsetfillcolor{currentfill}%
\pgfsetfillopacity{0.900000}%
\pgfsetlinewidth{0.200750pt}%
\definecolor{currentstroke}{rgb}{0.200000,0.200000,0.200000}%
\pgfsetstrokecolor{currentstroke}%
\pgfsetdash{}{0pt}%
\pgfpathmoveto{\pgfqpoint{1.400427in}{1.685581in}}%
\pgfpathlineto{\pgfqpoint{1.446373in}{1.699751in}}%
\pgfpathlineto{\pgfqpoint{1.492222in}{1.713891in}}%
\pgfpathlineto{\pgfqpoint{1.521891in}{1.693688in}}%
\pgfpathlineto{\pgfqpoint{1.551651in}{1.673423in}}%
\pgfpathlineto{\pgfqpoint{1.505726in}{1.659196in}}%
\pgfpathlineto{\pgfqpoint{1.459702in}{1.644939in}}%
\pgfpathlineto{\pgfqpoint{1.430019in}{1.665291in}}%
\pgfpathlineto{\pgfqpoint{1.400427in}{1.685581in}}%
\pgfpathclose%
\pgfusepath{stroke,fill}%
\end{pgfscope}%
\begin{pgfscope}%
\pgfpathrectangle{\pgfqpoint{0.050000in}{0.050000in}}{\pgfqpoint{4.900000in}{4.900000in}}%
\pgfusepath{clip}%
\pgfsetbuttcap%
\pgfsetroundjoin%
\definecolor{currentfill}{rgb}{1.000000,1.000000,1.000000}%
\pgfsetfillcolor{currentfill}%
\pgfsetfillopacity{0.900000}%
\pgfsetlinewidth{0.200750pt}%
\definecolor{currentstroke}{rgb}{0.200000,0.200000,0.200000}%
\pgfsetstrokecolor{currentstroke}%
\pgfsetdash{}{0pt}%
\pgfpathmoveto{\pgfqpoint{1.006143in}{1.681503in}}%
\pgfpathlineto{\pgfqpoint{1.052526in}{1.695682in}}%
\pgfpathlineto{\pgfqpoint{1.098810in}{1.709831in}}%
\pgfpathlineto{\pgfqpoint{1.127886in}{1.689615in}}%
\pgfpathlineto{\pgfqpoint{1.157051in}{1.669337in}}%
\pgfpathlineto{\pgfqpoint{1.110688in}{1.655102in}}%
\pgfpathlineto{\pgfqpoint{1.064225in}{1.640836in}}%
\pgfpathlineto{\pgfqpoint{1.035139in}{1.661200in}}%
\pgfpathlineto{\pgfqpoint{1.006143in}{1.681503in}}%
\pgfpathclose%
\pgfusepath{stroke,fill}%
\end{pgfscope}%
\begin{pgfscope}%
\pgfpathrectangle{\pgfqpoint{0.050000in}{0.050000in}}{\pgfqpoint{4.900000in}{4.900000in}}%
\pgfusepath{clip}%
\pgfsetbuttcap%
\pgfsetroundjoin%
\definecolor{currentfill}{rgb}{1.000000,1.000000,1.000000}%
\pgfsetfillcolor{currentfill}%
\pgfsetfillopacity{0.900000}%
\pgfsetlinewidth{0.200750pt}%
\definecolor{currentstroke}{rgb}{0.200000,0.200000,0.200000}%
\pgfsetstrokecolor{currentstroke}%
\pgfsetdash{}{0pt}%
\pgfpathmoveto{\pgfqpoint{2.340115in}{1.682677in}}%
\pgfpathlineto{\pgfqpoint{2.385066in}{1.698267in}}%
\pgfpathlineto{\pgfqpoint{2.429918in}{1.715256in}}%
\pgfpathlineto{\pgfqpoint{2.461033in}{1.696015in}}%
\pgfpathlineto{\pgfqpoint{2.492243in}{1.676796in}}%
\pgfpathlineto{\pgfqpoint{2.447317in}{1.658838in}}%
\pgfpathlineto{\pgfqpoint{2.402295in}{1.642614in}}%
\pgfpathlineto{\pgfqpoint{2.371157in}{1.662655in}}%
\pgfpathlineto{\pgfqpoint{2.340115in}{1.682677in}}%
\pgfpathclose%
\pgfusepath{stroke,fill}%
\end{pgfscope}%
\begin{pgfscope}%
\pgfpathrectangle{\pgfqpoint{0.050000in}{0.050000in}}{\pgfqpoint{4.900000in}{4.900000in}}%
\pgfusepath{clip}%
\pgfsetbuttcap%
\pgfsetroundjoin%
\definecolor{currentfill}{rgb}{0.844860,0.838462,0.907236}%
\pgfsetfillcolor{currentfill}%
\pgfsetfillopacity{0.900000}%
\pgfsetlinewidth{0.200750pt}%
\definecolor{currentstroke}{rgb}{0.200000,0.200000,0.200000}%
\pgfsetstrokecolor{currentstroke}%
\pgfsetdash{}{0pt}%
\pgfpathmoveto{\pgfqpoint{3.613620in}{2.012840in}}%
\pgfpathlineto{\pgfqpoint{3.655554in}{1.936517in}}%
\pgfpathlineto{\pgfqpoint{3.697382in}{1.866505in}}%
\pgfpathlineto{\pgfqpoint{3.730094in}{1.831698in}}%
\pgfpathlineto{\pgfqpoint{3.762905in}{1.797545in}}%
\pgfpathlineto{\pgfqpoint{3.720960in}{1.858852in}}%
\pgfpathlineto{\pgfqpoint{3.678967in}{1.928348in}}%
\pgfpathlineto{\pgfqpoint{3.646247in}{1.970153in}}%
\pgfpathlineto{\pgfqpoint{3.613620in}{2.012840in}}%
\pgfpathclose%
\pgfusepath{stroke,fill}%
\end{pgfscope}%
\begin{pgfscope}%
\pgfpathrectangle{\pgfqpoint{0.050000in}{0.050000in}}{\pgfqpoint{4.900000in}{4.900000in}}%
\pgfusepath{clip}%
\pgfsetbuttcap%
\pgfsetroundjoin%
\definecolor{currentfill}{rgb}{0.516678,0.496747,0.711003}%
\pgfsetfillcolor{currentfill}%
\pgfsetfillopacity{0.900000}%
\pgfsetlinewidth{0.200750pt}%
\definecolor{currentstroke}{rgb}{0.200000,0.200000,0.200000}%
\pgfsetstrokecolor{currentstroke}%
\pgfsetdash{}{0pt}%
\pgfpathmoveto{\pgfqpoint{2.939447in}{2.288482in}}%
\pgfpathlineto{\pgfqpoint{2.984410in}{2.397746in}}%
\pgfpathlineto{\pgfqpoint{3.029365in}{2.499199in}}%
\pgfpathlineto{\pgfqpoint{3.061362in}{2.468949in}}%
\pgfpathlineto{\pgfqpoint{3.093438in}{2.437941in}}%
\pgfpathlineto{\pgfqpoint{3.048437in}{2.348467in}}%
\pgfpathlineto{\pgfqpoint{3.003381in}{2.248810in}}%
\pgfpathlineto{\pgfqpoint{2.971368in}{2.269108in}}%
\pgfpathlineto{\pgfqpoint{2.939447in}{2.288482in}}%
\pgfpathclose%
\pgfusepath{stroke,fill}%
\end{pgfscope}%
\begin{pgfscope}%
\pgfpathrectangle{\pgfqpoint{0.050000in}{0.050000in}}{\pgfqpoint{4.900000in}{4.900000in}}%
\pgfusepath{clip}%
\pgfsetbuttcap%
\pgfsetroundjoin%
\definecolor{currentfill}{rgb}{1.000000,1.000000,1.000000}%
\pgfsetfillcolor{currentfill}%
\pgfsetfillopacity{0.900000}%
\pgfsetlinewidth{0.200750pt}%
\definecolor{currentstroke}{rgb}{0.200000,0.200000,0.200000}%
\pgfsetstrokecolor{currentstroke}%
\pgfsetdash{}{0pt}%
\pgfpathmoveto{\pgfqpoint{1.946007in}{1.677506in}}%
\pgfpathlineto{\pgfqpoint{1.991397in}{1.691694in}}%
\pgfpathlineto{\pgfqpoint{2.036690in}{1.705853in}}%
\pgfpathlineto{\pgfqpoint{2.067210in}{1.685627in}}%
\pgfpathlineto{\pgfqpoint{2.097825in}{1.665340in}}%
\pgfpathlineto{\pgfqpoint{2.052459in}{1.651091in}}%
\pgfpathlineto{\pgfqpoint{2.006995in}{1.636815in}}%
\pgfpathlineto{\pgfqpoint{1.976454in}{1.657192in}}%
\pgfpathlineto{\pgfqpoint{1.946007in}{1.677506in}}%
\pgfpathclose%
\pgfusepath{stroke,fill}%
\end{pgfscope}%
\begin{pgfscope}%
\pgfpathrectangle{\pgfqpoint{0.050000in}{0.050000in}}{\pgfqpoint{4.900000in}{4.900000in}}%
\pgfusepath{clip}%
\pgfsetbuttcap%
\pgfsetroundjoin%
\definecolor{currentfill}{rgb}{0.397339,0.372488,0.639646}%
\pgfsetfillcolor{currentfill}%
\pgfsetfillopacity{0.900000}%
\pgfsetlinewidth{0.200750pt}%
\definecolor{currentstroke}{rgb}{0.200000,0.200000,0.200000}%
\pgfsetstrokecolor{currentstroke}%
\pgfsetdash{}{0pt}%
\pgfpathmoveto{\pgfqpoint{3.029365in}{2.499199in}}%
\pgfpathlineto{\pgfqpoint{3.074208in}{2.583427in}}%
\pgfpathlineto{\pgfqpoint{3.118808in}{2.641073in}}%
\pgfpathlineto{\pgfqpoint{3.150776in}{2.596991in}}%
\pgfpathlineto{\pgfqpoint{3.182815in}{2.552912in}}%
\pgfpathlineto{\pgfqpoint{3.138274in}{2.508741in}}%
\pgfpathlineto{\pgfqpoint{3.093438in}{2.437941in}}%
\pgfpathlineto{\pgfqpoint{3.061362in}{2.468949in}}%
\pgfpathlineto{\pgfqpoint{3.029365in}{2.499199in}}%
\pgfpathclose%
\pgfusepath{stroke,fill}%
\end{pgfscope}%
\begin{pgfscope}%
\pgfpathrectangle{\pgfqpoint{0.050000in}{0.050000in}}{\pgfqpoint{4.900000in}{4.900000in}}%
\pgfusepath{clip}%
\pgfsetbuttcap%
\pgfsetroundjoin%
\definecolor{currentfill}{rgb}{1.000000,1.000000,1.000000}%
\pgfsetfillcolor{currentfill}%
\pgfsetfillopacity{0.900000}%
\pgfsetlinewidth{0.200750pt}%
\definecolor{currentstroke}{rgb}{0.200000,0.200000,0.200000}%
\pgfsetstrokecolor{currentstroke}%
\pgfsetdash{}{0pt}%
\pgfpathmoveto{\pgfqpoint{1.551651in}{1.673423in}}%
\pgfpathlineto{\pgfqpoint{1.597478in}{1.687619in}}%
\pgfpathlineto{\pgfqpoint{1.643206in}{1.701785in}}%
\pgfpathlineto{\pgfqpoint{1.673134in}{1.681545in}}%
\pgfpathlineto{\pgfqpoint{1.703153in}{1.661243in}}%
\pgfpathlineto{\pgfqpoint{1.657348in}{1.646990in}}%
\pgfpathlineto{\pgfqpoint{1.611445in}{1.632706in}}%
\pgfpathlineto{\pgfqpoint{1.581502in}{1.653096in}}%
\pgfpathlineto{\pgfqpoint{1.551651in}{1.673423in}}%
\pgfpathclose%
\pgfusepath{stroke,fill}%
\end{pgfscope}%
\begin{pgfscope}%
\pgfpathrectangle{\pgfqpoint{0.050000in}{0.050000in}}{\pgfqpoint{4.900000in}{4.900000in}}%
\pgfusepath{clip}%
\pgfsetbuttcap%
\pgfsetroundjoin%
\definecolor{currentfill}{rgb}{1.000000,1.000000,1.000000}%
\pgfsetfillcolor{currentfill}%
\pgfsetfillopacity{0.900000}%
\pgfsetlinewidth{0.200750pt}%
\definecolor{currentstroke}{rgb}{0.200000,0.200000,0.200000}%
\pgfsetstrokecolor{currentstroke}%
\pgfsetdash{}{0pt}%
\pgfpathmoveto{\pgfqpoint{1.157051in}{1.669337in}}%
\pgfpathlineto{\pgfqpoint{1.203315in}{1.683542in}}%
\pgfpathlineto{\pgfqpoint{1.249480in}{1.697717in}}%
\pgfpathlineto{\pgfqpoint{1.278814in}{1.677464in}}%
\pgfpathlineto{\pgfqpoint{1.308237in}{1.657149in}}%
\pgfpathlineto{\pgfqpoint{1.261994in}{1.642888in}}%
\pgfpathlineto{\pgfqpoint{1.215651in}{1.628595in}}%
\pgfpathlineto{\pgfqpoint{1.186306in}{1.648998in}}%
\pgfpathlineto{\pgfqpoint{1.157051in}{1.669337in}}%
\pgfpathclose%
\pgfusepath{stroke,fill}%
\end{pgfscope}%
\begin{pgfscope}%
\pgfpathrectangle{\pgfqpoint{0.050000in}{0.050000in}}{\pgfqpoint{4.900000in}{4.900000in}}%
\pgfusepath{clip}%
\pgfsetbuttcap%
\pgfsetroundjoin%
\definecolor{currentfill}{rgb}{0.618116,0.602368,0.771657}%
\pgfsetfillcolor{currentfill}%
\pgfsetfillopacity{0.900000}%
\pgfsetlinewidth{0.200750pt}%
\definecolor{currentstroke}{rgb}{0.200000,0.200000,0.200000}%
\pgfsetstrokecolor{currentstroke}%
\pgfsetdash{}{0pt}%
\pgfpathmoveto{\pgfqpoint{3.398975in}{2.335843in}}%
\pgfpathlineto{\pgfqpoint{3.441763in}{2.274455in}}%
\pgfpathlineto{\pgfqpoint{3.484026in}{2.193648in}}%
\pgfpathlineto{\pgfqpoint{3.516290in}{2.146800in}}%
\pgfpathlineto{\pgfqpoint{3.548643in}{2.101098in}}%
\pgfpathlineto{\pgfqpoint{3.506404in}{2.177993in}}%
\pgfpathlineto{\pgfqpoint{3.463684in}{2.239289in}}%
\pgfpathlineto{\pgfqpoint{3.431289in}{2.287126in}}%
\pgfpathlineto{\pgfqpoint{3.398975in}{2.335843in}}%
\pgfpathclose%
\pgfusepath{stroke,fill}%
\end{pgfscope}%
\begin{pgfscope}%
\pgfpathrectangle{\pgfqpoint{0.050000in}{0.050000in}}{\pgfqpoint{4.900000in}{4.900000in}}%
\pgfusepath{clip}%
\pgfsetbuttcap%
\pgfsetroundjoin%
\definecolor{currentfill}{rgb}{0.421207,0.397339,0.653918}%
\pgfsetfillcolor{currentfill}%
\pgfsetfillopacity{0.900000}%
\pgfsetlinewidth{0.200750pt}%
\definecolor{currentstroke}{rgb}{0.200000,0.200000,0.200000}%
\pgfsetstrokecolor{currentstroke}%
\pgfsetdash{}{0pt}%
\pgfpathmoveto{\pgfqpoint{3.182815in}{2.552912in}}%
\pgfpathlineto{\pgfqpoint{3.226929in}{2.564372in}}%
\pgfpathlineto{\pgfqpoint{3.270505in}{2.540199in}}%
\pgfpathlineto{\pgfqpoint{3.302507in}{2.487618in}}%
\pgfpathlineto{\pgfqpoint{3.334586in}{2.436054in}}%
\pgfpathlineto{\pgfqpoint{3.291123in}{2.466662in}}%
\pgfpathlineto{\pgfqpoint{3.247110in}{2.465012in}}%
\pgfpathlineto{\pgfqpoint{3.214926in}{2.508901in}}%
\pgfpathlineto{\pgfqpoint{3.182815in}{2.552912in}}%
\pgfpathclose%
\pgfusepath{stroke,fill}%
\end{pgfscope}%
\begin{pgfscope}%
\pgfpathrectangle{\pgfqpoint{0.050000in}{0.050000in}}{\pgfqpoint{4.900000in}{4.900000in}}%
\pgfusepath{clip}%
\pgfsetbuttcap%
\pgfsetroundjoin%
\definecolor{currentfill}{rgb}{1.000000,1.000000,1.000000}%
\pgfsetfillcolor{currentfill}%
\pgfsetfillopacity{0.900000}%
\pgfsetlinewidth{0.200750pt}%
\definecolor{currentstroke}{rgb}{0.200000,0.200000,0.200000}%
\pgfsetstrokecolor{currentstroke}%
\pgfsetdash{}{0pt}%
\pgfpathmoveto{\pgfqpoint{2.097825in}{1.665340in}}%
\pgfpathlineto{\pgfqpoint{2.143094in}{1.679572in}}%
\pgfpathlineto{\pgfqpoint{2.188265in}{1.693812in}}%
\pgfpathlineto{\pgfqpoint{2.219046in}{1.673583in}}%
\pgfpathlineto{\pgfqpoint{2.249921in}{1.653302in}}%
\pgfpathlineto{\pgfqpoint{2.204677in}{1.638924in}}%
\pgfpathlineto{\pgfqpoint{2.159336in}{1.624584in}}%
\pgfpathlineto{\pgfqpoint{2.128533in}{1.644993in}}%
\pgfpathlineto{\pgfqpoint{2.097825in}{1.665340in}}%
\pgfpathclose%
\pgfusepath{stroke,fill}%
\end{pgfscope}%
\begin{pgfscope}%
\pgfpathrectangle{\pgfqpoint{0.050000in}{0.050000in}}{\pgfqpoint{4.900000in}{4.900000in}}%
\pgfusepath{clip}%
\pgfsetbuttcap%
\pgfsetroundjoin%
\definecolor{currentfill}{rgb}{0.988066,0.987574,0.992864}%
\pgfsetfillcolor{currentfill}%
\pgfsetfillopacity{0.900000}%
\pgfsetlinewidth{0.200750pt}%
\definecolor{currentstroke}{rgb}{0.200000,0.200000,0.200000}%
\pgfsetstrokecolor{currentstroke}%
\pgfsetdash{}{0pt}%
\pgfpathmoveto{\pgfqpoint{2.492243in}{1.676796in}}%
\pgfpathlineto{\pgfqpoint{2.537078in}{1.697633in}}%
\pgfpathlineto{\pgfqpoint{2.581828in}{1.722911in}}%
\pgfpathlineto{\pgfqpoint{2.613219in}{1.705352in}}%
\pgfpathlineto{\pgfqpoint{2.644709in}{1.687809in}}%
\pgfpathlineto{\pgfqpoint{2.599871in}{1.660691in}}%
\pgfpathlineto{\pgfqpoint{2.554956in}{1.638416in}}%
\pgfpathlineto{\pgfqpoint{2.523551in}{1.657597in}}%
\pgfpathlineto{\pgfqpoint{2.492243in}{1.676796in}}%
\pgfpathclose%
\pgfusepath{stroke,fill}%
\end{pgfscope}%
\begin{pgfscope}%
\pgfpathrectangle{\pgfqpoint{0.050000in}{0.050000in}}{\pgfqpoint{4.900000in}{4.900000in}}%
\pgfusepath{clip}%
\pgfsetbuttcap%
\pgfsetroundjoin%
\definecolor{currentfill}{rgb}{1.000000,1.000000,1.000000}%
\pgfsetfillcolor{currentfill}%
\pgfsetfillopacity{0.900000}%
\pgfsetlinewidth{0.200750pt}%
\definecolor{currentstroke}{rgb}{0.200000,0.200000,0.200000}%
\pgfsetstrokecolor{currentstroke}%
\pgfsetdash{}{0pt}%
\pgfpathmoveto{\pgfqpoint{0.762208in}{1.665249in}}%
\pgfpathlineto{\pgfqpoint{0.808910in}{1.679463in}}%
\pgfpathlineto{\pgfqpoint{0.855512in}{1.693646in}}%
\pgfpathlineto{\pgfqpoint{0.884251in}{1.673381in}}%
\pgfpathlineto{\pgfqpoint{0.913078in}{1.653054in}}%
\pgfpathlineto{\pgfqpoint{0.866395in}{1.638783in}}%
\pgfpathlineto{\pgfqpoint{0.819611in}{1.624482in}}%
\pgfpathlineto{\pgfqpoint{0.790866in}{1.644897in}}%
\pgfpathlineto{\pgfqpoint{0.762208in}{1.665249in}}%
\pgfpathclose%
\pgfusepath{stroke,fill}%
\end{pgfscope}%
\begin{pgfscope}%
\pgfpathrectangle{\pgfqpoint{0.050000in}{0.050000in}}{\pgfqpoint{4.900000in}{4.900000in}}%
\pgfusepath{clip}%
\pgfsetbuttcap%
\pgfsetroundjoin%
\definecolor{currentfill}{rgb}{0.988066,0.987574,0.992864}%
\pgfsetfillcolor{currentfill}%
\pgfsetfillopacity{0.900000}%
\pgfsetlinewidth{0.200750pt}%
\definecolor{currentstroke}{rgb}{0.200000,0.200000,0.200000}%
\pgfsetstrokecolor{currentstroke}%
\pgfsetdash{}{0pt}%
\pgfpathmoveto{\pgfqpoint{3.828826in}{1.731015in}}%
\pgfpathlineto{\pgfqpoint{3.871126in}{1.696063in}}%
\pgfpathlineto{\pgfqpoint{3.914145in}{1.697885in}}%
\pgfpathlineto{\pgfqpoint{3.947561in}{1.677633in}}%
\pgfpathlineto{\pgfqpoint{3.981079in}{1.657318in}}%
\pgfpathlineto{\pgfqpoint{3.937711in}{1.643057in}}%
\pgfpathlineto{\pgfqpoint{3.895148in}{1.666583in}}%
\pgfpathlineto{\pgfqpoint{3.861936in}{1.698556in}}%
\pgfpathlineto{\pgfqpoint{3.828826in}{1.731015in}}%
\pgfpathclose%
\pgfusepath{stroke,fill}%
\end{pgfscope}%
\begin{pgfscope}%
\pgfpathrectangle{\pgfqpoint{0.050000in}{0.050000in}}{\pgfqpoint{4.900000in}{4.900000in}}%
\pgfusepath{clip}%
\pgfsetbuttcap%
\pgfsetroundjoin%
\definecolor{currentfill}{rgb}{1.000000,1.000000,1.000000}%
\pgfsetfillcolor{currentfill}%
\pgfsetfillopacity{0.900000}%
\pgfsetlinewidth{0.200750pt}%
\definecolor{currentstroke}{rgb}{0.200000,0.200000,0.200000}%
\pgfsetstrokecolor{currentstroke}%
\pgfsetdash{}{0pt}%
\pgfpathmoveto{\pgfqpoint{1.703153in}{1.661243in}}%
\pgfpathlineto{\pgfqpoint{1.748859in}{1.675465in}}%
\pgfpathlineto{\pgfqpoint{1.794468in}{1.689657in}}%
\pgfpathlineto{\pgfqpoint{1.824654in}{1.669380in}}%
\pgfpathlineto{\pgfqpoint{1.854934in}{1.649040in}}%
\pgfpathlineto{\pgfqpoint{1.809250in}{1.634761in}}%
\pgfpathlineto{\pgfqpoint{1.763469in}{1.620451in}}%
\pgfpathlineto{\pgfqpoint{1.733264in}{1.640878in}}%
\pgfpathlineto{\pgfqpoint{1.703153in}{1.661243in}}%
\pgfpathclose%
\pgfusepath{stroke,fill}%
\end{pgfscope}%
\begin{pgfscope}%
\pgfpathrectangle{\pgfqpoint{0.050000in}{0.050000in}}{\pgfqpoint{4.900000in}{4.900000in}}%
\pgfusepath{clip}%
\pgfsetbuttcap%
\pgfsetroundjoin%
\definecolor{currentfill}{rgb}{1.000000,1.000000,1.000000}%
\pgfsetfillcolor{currentfill}%
\pgfsetfillopacity{0.900000}%
\pgfsetlinewidth{0.200750pt}%
\definecolor{currentstroke}{rgb}{0.200000,0.200000,0.200000}%
\pgfsetstrokecolor{currentstroke}%
\pgfsetdash{}{0pt}%
\pgfpathmoveto{\pgfqpoint{1.308237in}{1.657149in}}%
\pgfpathlineto{\pgfqpoint{1.354381in}{1.671381in}}%
\pgfpathlineto{\pgfqpoint{1.400427in}{1.685581in}}%
\pgfpathlineto{\pgfqpoint{1.430019in}{1.665291in}}%
\pgfpathlineto{\pgfqpoint{1.459702in}{1.644939in}}%
\pgfpathlineto{\pgfqpoint{1.413578in}{1.630651in}}%
\pgfpathlineto{\pgfqpoint{1.367356in}{1.616332in}}%
\pgfpathlineto{\pgfqpoint{1.337751in}{1.636772in}}%
\pgfpathlineto{\pgfqpoint{1.308237in}{1.657149in}}%
\pgfpathclose%
\pgfusepath{stroke,fill}%
\end{pgfscope}%
\begin{pgfscope}%
\pgfpathrectangle{\pgfqpoint{0.050000in}{0.050000in}}{\pgfqpoint{4.900000in}{4.900000in}}%
\pgfusepath{clip}%
\pgfsetbuttcap%
\pgfsetroundjoin%
\definecolor{currentfill}{rgb}{0.797124,0.788758,0.878693}%
\pgfsetfillcolor{currentfill}%
\pgfsetfillopacity{0.900000}%
\pgfsetlinewidth{0.200750pt}%
\definecolor{currentstroke}{rgb}{0.200000,0.200000,0.200000}%
\pgfsetstrokecolor{currentstroke}%
\pgfsetdash{}{0pt}%
\pgfpathmoveto{\pgfqpoint{2.823705in}{1.882511in}}%
\pgfpathlineto{\pgfqpoint{2.868507in}{1.960058in}}%
\pgfpathlineto{\pgfqpoint{2.913386in}{2.049067in}}%
\pgfpathlineto{\pgfqpoint{2.945365in}{2.033671in}}%
\pgfpathlineto{\pgfqpoint{2.977443in}{2.017607in}}%
\pgfpathlineto{\pgfqpoint{2.932420in}{1.930944in}}%
\pgfpathlineto{\pgfqpoint{2.887472in}{1.853895in}}%
\pgfpathlineto{\pgfqpoint{2.855537in}{1.868386in}}%
\pgfpathlineto{\pgfqpoint{2.823705in}{1.882511in}}%
\pgfpathclose%
\pgfusepath{stroke,fill}%
\end{pgfscope}%
\begin{pgfscope}%
\pgfpathrectangle{\pgfqpoint{0.050000in}{0.050000in}}{\pgfqpoint{4.900000in}{4.900000in}}%
\pgfusepath{clip}%
\pgfsetbuttcap%
\pgfsetroundjoin%
\definecolor{currentfill}{rgb}{0.892595,0.888166,0.935779}%
\pgfsetfillcolor{currentfill}%
\pgfsetfillopacity{0.900000}%
\pgfsetlinewidth{0.200750pt}%
\definecolor{currentstroke}{rgb}{0.200000,0.200000,0.200000}%
\pgfsetstrokecolor{currentstroke}%
\pgfsetdash{}{0pt}%
\pgfpathmoveto{\pgfqpoint{2.734226in}{1.764206in}}%
\pgfpathlineto{\pgfqpoint{2.778954in}{1.817371in}}%
\pgfpathlineto{\pgfqpoint{2.823705in}{1.882511in}}%
\pgfpathlineto{\pgfqpoint{2.855537in}{1.868386in}}%
\pgfpathlineto{\pgfqpoint{2.887472in}{1.853895in}}%
\pgfpathlineto{\pgfqpoint{2.842582in}{1.787924in}}%
\pgfpathlineto{\pgfqpoint{2.797725in}{1.733119in}}%
\pgfpathlineto{\pgfqpoint{2.765924in}{1.748724in}}%
\pgfpathlineto{\pgfqpoint{2.734226in}{1.764206in}}%
\pgfpathclose%
\pgfusepath{stroke,fill}%
\end{pgfscope}%
\begin{pgfscope}%
\pgfpathrectangle{\pgfqpoint{0.050000in}{0.050000in}}{\pgfqpoint{4.900000in}{4.900000in}}%
\pgfusepath{clip}%
\pgfsetbuttcap%
\pgfsetroundjoin%
\definecolor{currentfill}{rgb}{1.000000,1.000000,1.000000}%
\pgfsetfillcolor{currentfill}%
\pgfsetfillopacity{0.900000}%
\pgfsetlinewidth{0.200750pt}%
\definecolor{currentstroke}{rgb}{0.200000,0.200000,0.200000}%
\pgfsetstrokecolor{currentstroke}%
\pgfsetdash{}{0pt}%
\pgfpathmoveto{\pgfqpoint{0.913078in}{1.653054in}}%
\pgfpathlineto{\pgfqpoint{0.959660in}{1.667294in}}%
\pgfpathlineto{\pgfqpoint{1.006143in}{1.681503in}}%
\pgfpathlineto{\pgfqpoint{1.035139in}{1.661200in}}%
\pgfpathlineto{\pgfqpoint{1.064225in}{1.640836in}}%
\pgfpathlineto{\pgfqpoint{1.017661in}{1.626539in}}%
\pgfpathlineto{\pgfqpoint{0.970998in}{1.612211in}}%
\pgfpathlineto{\pgfqpoint{0.941993in}{1.632664in}}%
\pgfpathlineto{\pgfqpoint{0.913078in}{1.653054in}}%
\pgfpathclose%
\pgfusepath{stroke,fill}%
\end{pgfscope}%
\begin{pgfscope}%
\pgfpathrectangle{\pgfqpoint{0.050000in}{0.050000in}}{\pgfqpoint{4.900000in}{4.900000in}}%
\pgfusepath{clip}%
\pgfsetbuttcap%
\pgfsetroundjoin%
\definecolor{currentfill}{rgb}{1.000000,1.000000,1.000000}%
\pgfsetfillcolor{currentfill}%
\pgfsetfillopacity{0.900000}%
\pgfsetlinewidth{0.200750pt}%
\definecolor{currentstroke}{rgb}{0.200000,0.200000,0.200000}%
\pgfsetstrokecolor{currentstroke}%
\pgfsetdash{}{0pt}%
\pgfpathmoveto{\pgfqpoint{2.249921in}{1.653302in}}%
\pgfpathlineto{\pgfqpoint{2.295067in}{1.667818in}}%
\pgfpathlineto{\pgfqpoint{2.340115in}{1.682677in}}%
\pgfpathlineto{\pgfqpoint{2.371157in}{1.662655in}}%
\pgfpathlineto{\pgfqpoint{2.402295in}{1.642614in}}%
\pgfpathlineto{\pgfqpoint{2.357175in}{1.627356in}}%
\pgfpathlineto{\pgfqpoint{2.311957in}{1.612591in}}%
\pgfpathlineto{\pgfqpoint{2.280892in}{1.632971in}}%
\pgfpathlineto{\pgfqpoint{2.249921in}{1.653302in}}%
\pgfpathclose%
\pgfusepath{stroke,fill}%
\end{pgfscope}%
\begin{pgfscope}%
\pgfpathrectangle{\pgfqpoint{0.050000in}{0.050000in}}{\pgfqpoint{4.900000in}{4.900000in}}%
\pgfusepath{clip}%
\pgfsetbuttcap%
\pgfsetroundjoin%
\definecolor{currentfill}{rgb}{1.000000,1.000000,1.000000}%
\pgfsetfillcolor{currentfill}%
\pgfsetfillopacity{0.900000}%
\pgfsetlinewidth{0.200750pt}%
\definecolor{currentstroke}{rgb}{0.200000,0.200000,0.200000}%
\pgfsetstrokecolor{currentstroke}%
\pgfsetdash{}{0pt}%
\pgfpathmoveto{\pgfqpoint{3.981079in}{1.657318in}}%
\pgfpathlineto{\pgfqpoint{4.024354in}{1.671549in}}%
\pgfpathlineto{\pgfqpoint{4.058005in}{1.651216in}}%
\pgfpathlineto{\pgfqpoint{4.091760in}{1.630821in}}%
\pgfpathlineto{\pgfqpoint{4.048424in}{1.616502in}}%
\pgfpathlineto{\pgfqpoint{4.014700in}{1.636942in}}%
\pgfpathlineto{\pgfqpoint{3.981079in}{1.657318in}}%
\pgfpathclose%
\pgfusepath{stroke,fill}%
\end{pgfscope}%
\begin{pgfscope}%
\pgfpathrectangle{\pgfqpoint{0.050000in}{0.050000in}}{\pgfqpoint{4.900000in}{4.900000in}}%
\pgfusepath{clip}%
\pgfsetbuttcap%
\pgfsetroundjoin%
\definecolor{currentfill}{rgb}{0.665852,0.652072,0.800200}%
\pgfsetfillcolor{currentfill}%
\pgfsetfillopacity{0.900000}%
\pgfsetlinewidth{0.200750pt}%
\definecolor{currentstroke}{rgb}{0.200000,0.200000,0.200000}%
\pgfsetstrokecolor{currentstroke}%
\pgfsetdash{}{0pt}%
\pgfpathmoveto{\pgfqpoint{2.913386in}{2.049067in}}%
\pgfpathlineto{\pgfqpoint{2.958349in}{2.146855in}}%
\pgfpathlineto{\pgfqpoint{3.003381in}{2.248810in}}%
\pgfpathlineto{\pgfqpoint{3.035484in}{2.227635in}}%
\pgfpathlineto{\pgfqpoint{3.067677in}{2.205630in}}%
\pgfpathlineto{\pgfqpoint{3.022538in}{2.110777in}}%
\pgfpathlineto{\pgfqpoint{2.977443in}{2.017607in}}%
\pgfpathlineto{\pgfqpoint{2.945365in}{2.033671in}}%
\pgfpathlineto{\pgfqpoint{2.913386in}{2.049067in}}%
\pgfpathclose%
\pgfusepath{stroke,fill}%
\end{pgfscope}%
\begin{pgfscope}%
\pgfpathrectangle{\pgfqpoint{0.050000in}{0.050000in}}{\pgfqpoint{4.900000in}{4.900000in}}%
\pgfusepath{clip}%
\pgfsetbuttcap%
\pgfsetroundjoin%
\definecolor{currentfill}{rgb}{1.000000,1.000000,1.000000}%
\pgfsetfillcolor{currentfill}%
\pgfsetfillopacity{0.900000}%
\pgfsetlinewidth{0.200750pt}%
\definecolor{currentstroke}{rgb}{0.200000,0.200000,0.200000}%
\pgfsetstrokecolor{currentstroke}%
\pgfsetdash{}{0pt}%
\pgfpathmoveto{\pgfqpoint{1.854934in}{1.649040in}}%
\pgfpathlineto{\pgfqpoint{1.900519in}{1.663288in}}%
\pgfpathlineto{\pgfqpoint{1.946007in}{1.677506in}}%
\pgfpathlineto{\pgfqpoint{1.976454in}{1.657192in}}%
\pgfpathlineto{\pgfqpoint{2.006995in}{1.636815in}}%
\pgfpathlineto{\pgfqpoint{1.961433in}{1.622509in}}%
\pgfpathlineto{\pgfqpoint{1.915773in}{1.608172in}}%
\pgfpathlineto{\pgfqpoint{1.885307in}{1.628638in}}%
\pgfpathlineto{\pgfqpoint{1.854934in}{1.649040in}}%
\pgfpathclose%
\pgfusepath{stroke,fill}%
\end{pgfscope}%
\begin{pgfscope}%
\pgfpathrectangle{\pgfqpoint{0.050000in}{0.050000in}}{\pgfqpoint{4.900000in}{4.900000in}}%
\pgfusepath{clip}%
\pgfsetbuttcap%
\pgfsetroundjoin%
\definecolor{currentfill}{rgb}{0.874694,0.869527,0.925075}%
\pgfsetfillcolor{currentfill}%
\pgfsetfillopacity{0.900000}%
\pgfsetlinewidth{0.200750pt}%
\definecolor{currentstroke}{rgb}{0.200000,0.200000,0.200000}%
\pgfsetstrokecolor{currentstroke}%
\pgfsetdash{}{0pt}%
\pgfpathmoveto{\pgfqpoint{3.678967in}{1.928348in}}%
\pgfpathlineto{\pgfqpoint{3.720960in}{1.858852in}}%
\pgfpathlineto{\pgfqpoint{3.762905in}{1.797545in}}%
\pgfpathlineto{\pgfqpoint{3.795816in}{1.763999in}}%
\pgfpathlineto{\pgfqpoint{3.828826in}{1.731015in}}%
\pgfpathlineto{\pgfqpoint{3.786754in}{1.784211in}}%
\pgfpathlineto{\pgfqpoint{3.744690in}{1.847172in}}%
\pgfpathlineto{\pgfqpoint{3.711781in}{1.887371in}}%
\pgfpathlineto{\pgfqpoint{3.678967in}{1.928348in}}%
\pgfpathclose%
\pgfusepath{stroke,fill}%
\end{pgfscope}%
\begin{pgfscope}%
\pgfpathrectangle{\pgfqpoint{0.050000in}{0.050000in}}{\pgfqpoint{4.900000in}{4.900000in}}%
\pgfusepath{clip}%
\pgfsetbuttcap%
\pgfsetroundjoin%
\definecolor{currentfill}{rgb}{1.000000,1.000000,1.000000}%
\pgfsetfillcolor{currentfill}%
\pgfsetfillopacity{0.900000}%
\pgfsetlinewidth{0.200750pt}%
\definecolor{currentstroke}{rgb}{0.200000,0.200000,0.200000}%
\pgfsetstrokecolor{currentstroke}%
\pgfsetdash{}{0pt}%
\pgfpathmoveto{\pgfqpoint{1.459702in}{1.644939in}}%
\pgfpathlineto{\pgfqpoint{1.505726in}{1.659196in}}%
\pgfpathlineto{\pgfqpoint{1.551651in}{1.673423in}}%
\pgfpathlineto{\pgfqpoint{1.581502in}{1.653096in}}%
\pgfpathlineto{\pgfqpoint{1.611445in}{1.632706in}}%
\pgfpathlineto{\pgfqpoint{1.565443in}{1.618392in}}%
\pgfpathlineto{\pgfqpoint{1.519342in}{1.604046in}}%
\pgfpathlineto{\pgfqpoint{1.489476in}{1.624524in}}%
\pgfpathlineto{\pgfqpoint{1.459702in}{1.644939in}}%
\pgfpathclose%
\pgfusepath{stroke,fill}%
\end{pgfscope}%
\begin{pgfscope}%
\pgfpathrectangle{\pgfqpoint{0.050000in}{0.050000in}}{\pgfqpoint{4.900000in}{4.900000in}}%
\pgfusepath{clip}%
\pgfsetbuttcap%
\pgfsetroundjoin%
\definecolor{currentfill}{rgb}{0.952265,0.950296,0.971457}%
\pgfsetfillcolor{currentfill}%
\pgfsetfillopacity{0.900000}%
\pgfsetlinewidth{0.200750pt}%
\definecolor{currentstroke}{rgb}{0.200000,0.200000,0.200000}%
\pgfsetstrokecolor{currentstroke}%
\pgfsetdash{}{0pt}%
\pgfpathmoveto{\pgfqpoint{2.644709in}{1.687809in}}%
\pgfpathlineto{\pgfqpoint{2.689487in}{1.721644in}}%
\pgfpathlineto{\pgfqpoint{2.734226in}{1.764206in}}%
\pgfpathlineto{\pgfqpoint{2.765924in}{1.748724in}}%
\pgfpathlineto{\pgfqpoint{2.797725in}{1.733119in}}%
\pgfpathlineto{\pgfqpoint{2.752871in}{1.688563in}}%
\pgfpathlineto{\pgfqpoint{2.707991in}{1.652725in}}%
\pgfpathlineto{\pgfqpoint{2.676300in}{1.670270in}}%
\pgfpathlineto{\pgfqpoint{2.644709in}{1.687809in}}%
\pgfpathclose%
\pgfusepath{stroke,fill}%
\end{pgfscope}%
\begin{pgfscope}%
\pgfpathrectangle{\pgfqpoint{0.050000in}{0.050000in}}{\pgfqpoint{4.900000in}{4.900000in}}%
\pgfusepath{clip}%
\pgfsetbuttcap%
\pgfsetroundjoin%
\definecolor{currentfill}{rgb}{0.433141,0.409765,0.661053}%
\pgfsetfillcolor{currentfill}%
\pgfsetfillopacity{0.900000}%
\pgfsetlinewidth{0.200750pt}%
\definecolor{currentstroke}{rgb}{0.200000,0.200000,0.200000}%
\pgfsetstrokecolor{currentstroke}%
\pgfsetdash{}{0pt}%
\pgfpathmoveto{\pgfqpoint{3.093438in}{2.437941in}}%
\pgfpathlineto{\pgfqpoint{3.138274in}{2.508741in}}%
\pgfpathlineto{\pgfqpoint{3.182815in}{2.552912in}}%
\pgfpathlineto{\pgfqpoint{3.214926in}{2.508901in}}%
\pgfpathlineto{\pgfqpoint{3.247110in}{2.465012in}}%
\pgfpathlineto{\pgfqpoint{3.202639in}{2.432594in}}%
\pgfpathlineto{\pgfqpoint{3.157828in}{2.373985in}}%
\pgfpathlineto{\pgfqpoint{3.125593in}{2.406261in}}%
\pgfpathlineto{\pgfqpoint{3.093438in}{2.437941in}}%
\pgfpathclose%
\pgfusepath{stroke,fill}%
\end{pgfscope}%
\begin{pgfscope}%
\pgfpathrectangle{\pgfqpoint{0.050000in}{0.050000in}}{\pgfqpoint{4.900000in}{4.900000in}}%
\pgfusepath{clip}%
\pgfsetbuttcap%
\pgfsetroundjoin%
\definecolor{currentfill}{rgb}{1.000000,1.000000,1.000000}%
\pgfsetfillcolor{currentfill}%
\pgfsetfillopacity{0.900000}%
\pgfsetlinewidth{0.200750pt}%
\definecolor{currentstroke}{rgb}{0.200000,0.200000,0.200000}%
\pgfsetstrokecolor{currentstroke}%
\pgfsetdash{}{0pt}%
\pgfpathmoveto{\pgfqpoint{1.064225in}{1.640836in}}%
\pgfpathlineto{\pgfqpoint{1.110688in}{1.655102in}}%
\pgfpathlineto{\pgfqpoint{1.157051in}{1.669337in}}%
\pgfpathlineto{\pgfqpoint{1.186306in}{1.648998in}}%
\pgfpathlineto{\pgfqpoint{1.215651in}{1.628595in}}%
\pgfpathlineto{\pgfqpoint{1.169208in}{1.614272in}}%
\pgfpathlineto{\pgfqpoint{1.122664in}{1.599918in}}%
\pgfpathlineto{\pgfqpoint{1.093400in}{1.620408in}}%
\pgfpathlineto{\pgfqpoint{1.064225in}{1.640836in}}%
\pgfpathclose%
\pgfusepath{stroke,fill}%
\end{pgfscope}%
\begin{pgfscope}%
\pgfpathrectangle{\pgfqpoint{0.050000in}{0.050000in}}{\pgfqpoint{4.900000in}{4.900000in}}%
\pgfusepath{clip}%
\pgfsetbuttcap%
\pgfsetroundjoin%
\definecolor{currentfill}{rgb}{0.659885,0.645859,0.796632}%
\pgfsetfillcolor{currentfill}%
\pgfsetfillopacity{0.900000}%
\pgfsetlinewidth{0.200750pt}%
\definecolor{currentstroke}{rgb}{0.200000,0.200000,0.200000}%
\pgfsetstrokecolor{currentstroke}%
\pgfsetdash{}{0pt}%
\pgfpathmoveto{\pgfqpoint{3.463684in}{2.239289in}}%
\pgfpathlineto{\pgfqpoint{3.506404in}{2.177993in}}%
\pgfpathlineto{\pgfqpoint{3.548643in}{2.101098in}}%
\pgfpathlineto{\pgfqpoint{3.581086in}{2.056467in}}%
\pgfpathlineto{\pgfqpoint{3.613620in}{2.012840in}}%
\pgfpathlineto{\pgfqpoint{3.571391in}{2.085597in}}%
\pgfpathlineto{\pgfqpoint{3.528724in}{2.146112in}}%
\pgfpathlineto{\pgfqpoint{3.496162in}{2.192295in}}%
\pgfpathlineto{\pgfqpoint{3.463684in}{2.239289in}}%
\pgfpathclose%
\pgfusepath{stroke,fill}%
\end{pgfscope}%
\begin{pgfscope}%
\pgfpathrectangle{\pgfqpoint{0.050000in}{0.050000in}}{\pgfqpoint{4.900000in}{4.900000in}}%
\pgfusepath{clip}%
\pgfsetbuttcap%
\pgfsetroundjoin%
\definecolor{currentfill}{rgb}{0.528612,0.509173,0.718139}%
\pgfsetfillcolor{currentfill}%
\pgfsetfillopacity{0.900000}%
\pgfsetlinewidth{0.200750pt}%
\definecolor{currentstroke}{rgb}{0.200000,0.200000,0.200000}%
\pgfsetstrokecolor{currentstroke}%
\pgfsetdash{}{0pt}%
\pgfpathmoveto{\pgfqpoint{3.003381in}{2.248810in}}%
\pgfpathlineto{\pgfqpoint{3.048437in}{2.348467in}}%
\pgfpathlineto{\pgfqpoint{3.093438in}{2.437941in}}%
\pgfpathlineto{\pgfqpoint{3.125593in}{2.406261in}}%
\pgfpathlineto{\pgfqpoint{3.157828in}{2.373985in}}%
\pgfpathlineto{\pgfqpoint{3.112803in}{2.295819in}}%
\pgfpathlineto{\pgfqpoint{3.067677in}{2.205630in}}%
\pgfpathlineto{\pgfqpoint{3.035484in}{2.227635in}}%
\pgfpathlineto{\pgfqpoint{3.003381in}{2.248810in}}%
\pgfpathclose%
\pgfusepath{stroke,fill}%
\end{pgfscope}%
\begin{pgfscope}%
\pgfpathrectangle{\pgfqpoint{0.050000in}{0.050000in}}{\pgfqpoint{4.900000in}{4.900000in}}%
\pgfusepath{clip}%
\pgfsetbuttcap%
\pgfsetroundjoin%
\definecolor{currentfill}{rgb}{1.000000,1.000000,1.000000}%
\pgfsetfillcolor{currentfill}%
\pgfsetfillopacity{0.900000}%
\pgfsetlinewidth{0.200750pt}%
\definecolor{currentstroke}{rgb}{0.200000,0.200000,0.200000}%
\pgfsetstrokecolor{currentstroke}%
\pgfsetdash{}{0pt}%
\pgfpathmoveto{\pgfqpoint{2.402295in}{1.642614in}}%
\pgfpathlineto{\pgfqpoint{2.447317in}{1.658838in}}%
\pgfpathlineto{\pgfqpoint{2.492243in}{1.676796in}}%
\pgfpathlineto{\pgfqpoint{2.523551in}{1.657597in}}%
\pgfpathlineto{\pgfqpoint{2.554956in}{1.638416in}}%
\pgfpathlineto{\pgfqpoint{2.509954in}{1.619412in}}%
\pgfpathlineto{\pgfqpoint{2.464859in}{1.602475in}}%
\pgfpathlineto{\pgfqpoint{2.433528in}{1.622553in}}%
\pgfpathlineto{\pgfqpoint{2.402295in}{1.642614in}}%
\pgfpathclose%
\pgfusepath{stroke,fill}%
\end{pgfscope}%
\begin{pgfscope}%
\pgfpathrectangle{\pgfqpoint{0.050000in}{0.050000in}}{\pgfqpoint{4.900000in}{4.900000in}}%
\pgfusepath{clip}%
\pgfsetbuttcap%
\pgfsetroundjoin%
\definecolor{currentfill}{rgb}{0.468943,0.447043,0.682461}%
\pgfsetfillcolor{currentfill}%
\pgfsetfillopacity{0.900000}%
\pgfsetlinewidth{0.200750pt}%
\definecolor{currentstroke}{rgb}{0.200000,0.200000,0.200000}%
\pgfsetstrokecolor{currentstroke}%
\pgfsetdash{}{0pt}%
\pgfpathmoveto{\pgfqpoint{3.247110in}{2.465012in}}%
\pgfpathlineto{\pgfqpoint{3.291123in}{2.466662in}}%
\pgfpathlineto{\pgfqpoint{3.334586in}{2.436054in}}%
\pgfpathlineto{\pgfqpoint{3.366741in}{2.385474in}}%
\pgfpathlineto{\pgfqpoint{3.398975in}{2.335843in}}%
\pgfpathlineto{\pgfqpoint{3.355615in}{2.371314in}}%
\pgfpathlineto{\pgfqpoint{3.311702in}{2.377776in}}%
\pgfpathlineto{\pgfqpoint{3.279368in}{2.421292in}}%
\pgfpathlineto{\pgfqpoint{3.247110in}{2.465012in}}%
\pgfpathclose%
\pgfusepath{stroke,fill}%
\end{pgfscope}%
\begin{pgfscope}%
\pgfpathrectangle{\pgfqpoint{0.050000in}{0.050000in}}{\pgfqpoint{4.900000in}{4.900000in}}%
\pgfusepath{clip}%
\pgfsetbuttcap%
\pgfsetroundjoin%
\definecolor{currentfill}{rgb}{1.000000,1.000000,1.000000}%
\pgfsetfillcolor{currentfill}%
\pgfsetfillopacity{0.900000}%
\pgfsetlinewidth{0.200750pt}%
\definecolor{currentstroke}{rgb}{0.200000,0.200000,0.200000}%
\pgfsetstrokecolor{currentstroke}%
\pgfsetdash{}{0pt}%
\pgfpathmoveto{\pgfqpoint{2.006995in}{1.636815in}}%
\pgfpathlineto{\pgfqpoint{2.052459in}{1.651091in}}%
\pgfpathlineto{\pgfqpoint{2.097825in}{1.665340in}}%
\pgfpathlineto{\pgfqpoint{2.128533in}{1.644993in}}%
\pgfpathlineto{\pgfqpoint{2.159336in}{1.624584in}}%
\pgfpathlineto{\pgfqpoint{2.113896in}{1.610240in}}%
\pgfpathlineto{\pgfqpoint{2.068359in}{1.595873in}}%
\pgfpathlineto{\pgfqpoint{2.037629in}{1.616375in}}%
\pgfpathlineto{\pgfqpoint{2.006995in}{1.636815in}}%
\pgfpathclose%
\pgfusepath{stroke,fill}%
\end{pgfscope}%
\begin{pgfscope}%
\pgfpathrectangle{\pgfqpoint{0.050000in}{0.050000in}}{\pgfqpoint{4.900000in}{4.900000in}}%
\pgfusepath{clip}%
\pgfsetbuttcap%
\pgfsetroundjoin%
\definecolor{currentfill}{rgb}{1.000000,1.000000,1.000000}%
\pgfsetfillcolor{currentfill}%
\pgfsetfillopacity{0.900000}%
\pgfsetlinewidth{0.200750pt}%
\definecolor{currentstroke}{rgb}{0.200000,0.200000,0.200000}%
\pgfsetstrokecolor{currentstroke}%
\pgfsetdash{}{0pt}%
\pgfpathmoveto{\pgfqpoint{1.611445in}{1.632706in}}%
\pgfpathlineto{\pgfqpoint{1.657348in}{1.646990in}}%
\pgfpathlineto{\pgfqpoint{1.703153in}{1.661243in}}%
\pgfpathlineto{\pgfqpoint{1.733264in}{1.640878in}}%
\pgfpathlineto{\pgfqpoint{1.763469in}{1.620451in}}%
\pgfpathlineto{\pgfqpoint{1.717588in}{1.606110in}}%
\pgfpathlineto{\pgfqpoint{1.671608in}{1.591738in}}%
\pgfpathlineto{\pgfqpoint{1.641480in}{1.612254in}}%
\pgfpathlineto{\pgfqpoint{1.611445in}{1.632706in}}%
\pgfpathclose%
\pgfusepath{stroke,fill}%
\end{pgfscope}%
\begin{pgfscope}%
\pgfpathrectangle{\pgfqpoint{0.050000in}{0.050000in}}{\pgfqpoint{4.900000in}{4.900000in}}%
\pgfusepath{clip}%
\pgfsetbuttcap%
\pgfsetroundjoin%
\definecolor{currentfill}{rgb}{1.000000,1.000000,1.000000}%
\pgfsetfillcolor{currentfill}%
\pgfsetfillopacity{0.900000}%
\pgfsetlinewidth{0.200750pt}%
\definecolor{currentstroke}{rgb}{0.200000,0.200000,0.200000}%
\pgfsetstrokecolor{currentstroke}%
\pgfsetdash{}{0pt}%
\pgfpathmoveto{\pgfqpoint{1.215651in}{1.628595in}}%
\pgfpathlineto{\pgfqpoint{1.261994in}{1.642888in}}%
\pgfpathlineto{\pgfqpoint{1.308237in}{1.657149in}}%
\pgfpathlineto{\pgfqpoint{1.337751in}{1.636772in}}%
\pgfpathlineto{\pgfqpoint{1.367356in}{1.616332in}}%
\pgfpathlineto{\pgfqpoint{1.321034in}{1.601982in}}%
\pgfpathlineto{\pgfqpoint{1.274611in}{1.587602in}}%
\pgfpathlineto{\pgfqpoint{1.245086in}{1.608130in}}%
\pgfpathlineto{\pgfqpoint{1.215651in}{1.628595in}}%
\pgfpathclose%
\pgfusepath{stroke,fill}%
\end{pgfscope}%
\begin{pgfscope}%
\pgfpathrectangle{\pgfqpoint{0.050000in}{0.050000in}}{\pgfqpoint{4.900000in}{4.900000in}}%
\pgfusepath{clip}%
\pgfsetbuttcap%
\pgfsetroundjoin%
\definecolor{currentfill}{rgb}{0.994033,0.993787,0.996432}%
\pgfsetfillcolor{currentfill}%
\pgfsetfillopacity{0.900000}%
\pgfsetlinewidth{0.200750pt}%
\definecolor{currentstroke}{rgb}{0.200000,0.200000,0.200000}%
\pgfsetstrokecolor{currentstroke}%
\pgfsetdash{}{0pt}%
\pgfpathmoveto{\pgfqpoint{3.895148in}{1.666583in}}%
\pgfpathlineto{\pgfqpoint{3.937711in}{1.643057in}}%
\pgfpathlineto{\pgfqpoint{3.981079in}{1.657318in}}%
\pgfpathlineto{\pgfqpoint{4.014700in}{1.636942in}}%
\pgfpathlineto{\pgfqpoint{4.048424in}{1.616502in}}%
\pgfpathlineto{\pgfqpoint{4.004995in}{1.602153in}}%
\pgfpathlineto{\pgfqpoint{3.961877in}{1.603969in}}%
\pgfpathlineto{\pgfqpoint{3.928461in}{1.635065in}}%
\pgfpathlineto{\pgfqpoint{3.895148in}{1.666583in}}%
\pgfpathclose%
\pgfusepath{stroke,fill}%
\end{pgfscope}%
\begin{pgfscope}%
\pgfpathrectangle{\pgfqpoint{0.050000in}{0.050000in}}{\pgfqpoint{4.900000in}{4.900000in}}%
\pgfusepath{clip}%
\pgfsetbuttcap%
\pgfsetroundjoin%
\definecolor{currentfill}{rgb}{1.000000,1.000000,1.000000}%
\pgfsetfillcolor{currentfill}%
\pgfsetfillopacity{0.900000}%
\pgfsetlinewidth{0.200750pt}%
\definecolor{currentstroke}{rgb}{0.200000,0.200000,0.200000}%
\pgfsetstrokecolor{currentstroke}%
\pgfsetdash{}{0pt}%
\pgfpathmoveto{\pgfqpoint{2.159336in}{1.624584in}}%
\pgfpathlineto{\pgfqpoint{2.204677in}{1.638924in}}%
\pgfpathlineto{\pgfqpoint{2.249921in}{1.653302in}}%
\pgfpathlineto{\pgfqpoint{2.280892in}{1.632971in}}%
\pgfpathlineto{\pgfqpoint{2.311957in}{1.612591in}}%
\pgfpathlineto{\pgfqpoint{2.266641in}{1.598047in}}%
\pgfpathlineto{\pgfqpoint{2.221227in}{1.583586in}}%
\pgfpathlineto{\pgfqpoint{2.190234in}{1.604115in}}%
\pgfpathlineto{\pgfqpoint{2.159336in}{1.624584in}}%
\pgfpathclose%
\pgfusepath{stroke,fill}%
\end{pgfscope}%
\begin{pgfscope}%
\pgfpathrectangle{\pgfqpoint{0.050000in}{0.050000in}}{\pgfqpoint{4.900000in}{4.900000in}}%
\pgfusepath{clip}%
\pgfsetbuttcap%
\pgfsetroundjoin%
\definecolor{currentfill}{rgb}{1.000000,1.000000,1.000000}%
\pgfsetfillcolor{currentfill}%
\pgfsetfillopacity{0.900000}%
\pgfsetlinewidth{0.200750pt}%
\definecolor{currentstroke}{rgb}{0.200000,0.200000,0.200000}%
\pgfsetstrokecolor{currentstroke}%
\pgfsetdash{}{0pt}%
\pgfpathmoveto{\pgfqpoint{0.819611in}{1.624482in}}%
\pgfpathlineto{\pgfqpoint{0.866395in}{1.638783in}}%
\pgfpathlineto{\pgfqpoint{0.913078in}{1.653054in}}%
\pgfpathlineto{\pgfqpoint{0.941993in}{1.632664in}}%
\pgfpathlineto{\pgfqpoint{0.970998in}{1.612211in}}%
\pgfpathlineto{\pgfqpoint{0.924233in}{1.597852in}}%
\pgfpathlineto{\pgfqpoint{0.877368in}{1.583463in}}%
\pgfpathlineto{\pgfqpoint{0.848445in}{1.604004in}}%
\pgfpathlineto{\pgfqpoint{0.819611in}{1.624482in}}%
\pgfpathclose%
\pgfusepath{stroke,fill}%
\end{pgfscope}%
\begin{pgfscope}%
\pgfpathrectangle{\pgfqpoint{0.050000in}{0.050000in}}{\pgfqpoint{4.900000in}{4.900000in}}%
\pgfusepath{clip}%
\pgfsetbuttcap%
\pgfsetroundjoin%
\definecolor{currentfill}{rgb}{0.982099,0.981361,0.989296}%
\pgfsetfillcolor{currentfill}%
\pgfsetfillopacity{0.900000}%
\pgfsetlinewidth{0.200750pt}%
\definecolor{currentstroke}{rgb}{0.200000,0.200000,0.200000}%
\pgfsetstrokecolor{currentstroke}%
\pgfsetdash{}{0pt}%
\pgfpathmoveto{\pgfqpoint{2.554956in}{1.638416in}}%
\pgfpathlineto{\pgfqpoint{2.599871in}{1.660691in}}%
\pgfpathlineto{\pgfqpoint{2.644709in}{1.687809in}}%
\pgfpathlineto{\pgfqpoint{2.676300in}{1.670270in}}%
\pgfpathlineto{\pgfqpoint{2.707991in}{1.652725in}}%
\pgfpathlineto{\pgfqpoint{2.663061in}{1.623820in}}%
\pgfpathlineto{\pgfqpoint{2.618062in}{1.600089in}}%
\pgfpathlineto{\pgfqpoint{2.586460in}{1.619247in}}%
\pgfpathlineto{\pgfqpoint{2.554956in}{1.638416in}}%
\pgfpathclose%
\pgfusepath{stroke,fill}%
\end{pgfscope}%
\begin{pgfscope}%
\pgfpathrectangle{\pgfqpoint{0.050000in}{0.050000in}}{\pgfqpoint{4.900000in}{4.900000in}}%
\pgfusepath{clip}%
\pgfsetbuttcap%
\pgfsetroundjoin%
\definecolor{currentfill}{rgb}{1.000000,1.000000,1.000000}%
\pgfsetfillcolor{currentfill}%
\pgfsetfillopacity{0.900000}%
\pgfsetlinewidth{0.200750pt}%
\definecolor{currentstroke}{rgb}{0.200000,0.200000,0.200000}%
\pgfsetstrokecolor{currentstroke}%
\pgfsetdash{}{0pt}%
\pgfpathmoveto{\pgfqpoint{1.763469in}{1.620451in}}%
\pgfpathlineto{\pgfqpoint{1.809250in}{1.634761in}}%
\pgfpathlineto{\pgfqpoint{1.854934in}{1.649040in}}%
\pgfpathlineto{\pgfqpoint{1.885307in}{1.628638in}}%
\pgfpathlineto{\pgfqpoint{1.915773in}{1.608172in}}%
\pgfpathlineto{\pgfqpoint{1.870014in}{1.593805in}}%
\pgfpathlineto{\pgfqpoint{1.824157in}{1.579407in}}%
\pgfpathlineto{\pgfqpoint{1.793766in}{1.599960in}}%
\pgfpathlineto{\pgfqpoint{1.763469in}{1.620451in}}%
\pgfpathclose%
\pgfusepath{stroke,fill}%
\end{pgfscope}%
\begin{pgfscope}%
\pgfpathrectangle{\pgfqpoint{0.050000in}{0.050000in}}{\pgfqpoint{4.900000in}{4.900000in}}%
\pgfusepath{clip}%
\pgfsetbuttcap%
\pgfsetroundjoin%
\definecolor{currentfill}{rgb}{1.000000,1.000000,1.000000}%
\pgfsetfillcolor{currentfill}%
\pgfsetfillopacity{0.900000}%
\pgfsetlinewidth{0.200750pt}%
\definecolor{currentstroke}{rgb}{0.200000,0.200000,0.200000}%
\pgfsetstrokecolor{currentstroke}%
\pgfsetdash{}{0pt}%
\pgfpathmoveto{\pgfqpoint{1.367356in}{1.616332in}}%
\pgfpathlineto{\pgfqpoint{1.413578in}{1.630651in}}%
\pgfpathlineto{\pgfqpoint{1.459702in}{1.644939in}}%
\pgfpathlineto{\pgfqpoint{1.489476in}{1.624524in}}%
\pgfpathlineto{\pgfqpoint{1.519342in}{1.604046in}}%
\pgfpathlineto{\pgfqpoint{1.473141in}{1.589670in}}%
\pgfpathlineto{\pgfqpoint{1.426840in}{1.575263in}}%
\pgfpathlineto{\pgfqpoint{1.397052in}{1.595829in}}%
\pgfpathlineto{\pgfqpoint{1.367356in}{1.616332in}}%
\pgfpathclose%
\pgfusepath{stroke,fill}%
\end{pgfscope}%
\begin{pgfscope}%
\pgfpathrectangle{\pgfqpoint{0.050000in}{0.050000in}}{\pgfqpoint{4.900000in}{4.900000in}}%
\pgfusepath{clip}%
\pgfsetbuttcap%
\pgfsetroundjoin%
\definecolor{currentfill}{rgb}{0.785190,0.776332,0.871557}%
\pgfsetfillcolor{currentfill}%
\pgfsetfillopacity{0.900000}%
\pgfsetlinewidth{0.200750pt}%
\definecolor{currentstroke}{rgb}{0.200000,0.200000,0.200000}%
\pgfsetstrokecolor{currentstroke}%
\pgfsetdash{}{0pt}%
\pgfpathmoveto{\pgfqpoint{2.887472in}{1.853895in}}%
\pgfpathlineto{\pgfqpoint{2.932420in}{1.930944in}}%
\pgfpathlineto{\pgfqpoint{2.977443in}{2.017607in}}%
\pgfpathlineto{\pgfqpoint{3.009619in}{2.000885in}}%
\pgfpathlineto{\pgfqpoint{3.041893in}{1.983520in}}%
\pgfpathlineto{\pgfqpoint{2.996737in}{1.899744in}}%
\pgfpathlineto{\pgfqpoint{2.951649in}{1.823769in}}%
\pgfpathlineto{\pgfqpoint{2.919509in}{1.839026in}}%
\pgfpathlineto{\pgfqpoint{2.887472in}{1.853895in}}%
\pgfpathclose%
\pgfusepath{stroke,fill}%
\end{pgfscope}%
\begin{pgfscope}%
\pgfpathrectangle{\pgfqpoint{0.050000in}{0.050000in}}{\pgfqpoint{4.900000in}{4.900000in}}%
\pgfusepath{clip}%
\pgfsetbuttcap%
\pgfsetroundjoin%
\definecolor{currentfill}{rgb}{0.904529,0.900592,0.942914}%
\pgfsetfillcolor{currentfill}%
\pgfsetfillopacity{0.900000}%
\pgfsetlinewidth{0.200750pt}%
\definecolor{currentstroke}{rgb}{0.200000,0.200000,0.200000}%
\pgfsetstrokecolor{currentstroke}%
\pgfsetdash{}{0pt}%
\pgfpathmoveto{\pgfqpoint{3.744690in}{1.847172in}}%
\pgfpathlineto{\pgfqpoint{3.786754in}{1.784211in}}%
\pgfpathlineto{\pgfqpoint{3.828826in}{1.731015in}}%
\pgfpathlineto{\pgfqpoint{3.861936in}{1.698556in}}%
\pgfpathlineto{\pgfqpoint{3.895148in}{1.666583in}}%
\pgfpathlineto{\pgfqpoint{3.852942in}{1.712210in}}%
\pgfpathlineto{\pgfqpoint{3.810797in}{1.768927in}}%
\pgfpathlineto{\pgfqpoint{3.777695in}{1.807705in}}%
\pgfpathlineto{\pgfqpoint{3.744690in}{1.847172in}}%
\pgfpathclose%
\pgfusepath{stroke,fill}%
\end{pgfscope}%
\begin{pgfscope}%
\pgfpathrectangle{\pgfqpoint{0.050000in}{0.050000in}}{\pgfqpoint{4.900000in}{4.900000in}}%
\pgfusepath{clip}%
\pgfsetbuttcap%
\pgfsetroundjoin%
\definecolor{currentfill}{rgb}{0.886628,0.881953,0.932211}%
\pgfsetfillcolor{currentfill}%
\pgfsetfillopacity{0.900000}%
\pgfsetlinewidth{0.200750pt}%
\definecolor{currentstroke}{rgb}{0.200000,0.200000,0.200000}%
\pgfsetstrokecolor{currentstroke}%
\pgfsetdash{}{0pt}%
\pgfpathmoveto{\pgfqpoint{2.797725in}{1.733119in}}%
\pgfpathlineto{\pgfqpoint{2.842582in}{1.787924in}}%
\pgfpathlineto{\pgfqpoint{2.887472in}{1.853895in}}%
\pgfpathlineto{\pgfqpoint{2.919509in}{1.839026in}}%
\pgfpathlineto{\pgfqpoint{2.951649in}{1.823769in}}%
\pgfpathlineto{\pgfqpoint{2.906622in}{1.757493in}}%
\pgfpathlineto{\pgfqpoint{2.861636in}{1.701475in}}%
\pgfpathlineto{\pgfqpoint{2.829629in}{1.717374in}}%
\pgfpathlineto{\pgfqpoint{2.797725in}{1.733119in}}%
\pgfpathclose%
\pgfusepath{stroke,fill}%
\end{pgfscope}%
\begin{pgfscope}%
\pgfpathrectangle{\pgfqpoint{0.050000in}{0.050000in}}{\pgfqpoint{4.900000in}{4.900000in}}%
\pgfusepath{clip}%
\pgfsetbuttcap%
\pgfsetroundjoin%
\definecolor{currentfill}{rgb}{1.000000,1.000000,1.000000}%
\pgfsetfillcolor{currentfill}%
\pgfsetfillopacity{0.900000}%
\pgfsetlinewidth{0.200750pt}%
\definecolor{currentstroke}{rgb}{0.200000,0.200000,0.200000}%
\pgfsetstrokecolor{currentstroke}%
\pgfsetdash{}{0pt}%
\pgfpathmoveto{\pgfqpoint{0.970998in}{1.612211in}}%
\pgfpathlineto{\pgfqpoint{1.017661in}{1.626539in}}%
\pgfpathlineto{\pgfqpoint{1.064225in}{1.640836in}}%
\pgfpathlineto{\pgfqpoint{1.093400in}{1.620408in}}%
\pgfpathlineto{\pgfqpoint{1.122664in}{1.599918in}}%
\pgfpathlineto{\pgfqpoint{1.076021in}{1.585532in}}%
\pgfpathlineto{\pgfqpoint{1.029276in}{1.571116in}}%
\pgfpathlineto{\pgfqpoint{1.000092in}{1.591695in}}%
\pgfpathlineto{\pgfqpoint{0.970998in}{1.612211in}}%
\pgfpathclose%
\pgfusepath{stroke,fill}%
\end{pgfscope}%
\begin{pgfscope}%
\pgfpathrectangle{\pgfqpoint{0.050000in}{0.050000in}}{\pgfqpoint{4.900000in}{4.900000in}}%
\pgfusepath{clip}%
\pgfsetbuttcap%
\pgfsetroundjoin%
\definecolor{currentfill}{rgb}{1.000000,1.000000,1.000000}%
\pgfsetfillcolor{currentfill}%
\pgfsetfillopacity{0.900000}%
\pgfsetlinewidth{0.200750pt}%
\definecolor{currentstroke}{rgb}{0.200000,0.200000,0.200000}%
\pgfsetstrokecolor{currentstroke}%
\pgfsetdash{}{0pt}%
\pgfpathmoveto{\pgfqpoint{2.311957in}{1.612591in}}%
\pgfpathlineto{\pgfqpoint{2.357175in}{1.627356in}}%
\pgfpathlineto{\pgfqpoint{2.402295in}{1.642614in}}%
\pgfpathlineto{\pgfqpoint{2.433528in}{1.622553in}}%
\pgfpathlineto{\pgfqpoint{2.464859in}{1.602475in}}%
\pgfpathlineto{\pgfqpoint{2.419666in}{1.586752in}}%
\pgfpathlineto{\pgfqpoint{2.374376in}{1.571688in}}%
\pgfpathlineto{\pgfqpoint{2.343119in}{1.592163in}}%
\pgfpathlineto{\pgfqpoint{2.311957in}{1.612591in}}%
\pgfpathclose%
\pgfusepath{stroke,fill}%
\end{pgfscope}%
\begin{pgfscope}%
\pgfpathrectangle{\pgfqpoint{0.050000in}{0.050000in}}{\pgfqpoint{4.900000in}{4.900000in}}%
\pgfusepath{clip}%
\pgfsetbuttcap%
\pgfsetroundjoin%
\definecolor{currentfill}{rgb}{0.701653,0.689350,0.821607}%
\pgfsetfillcolor{currentfill}%
\pgfsetfillopacity{0.900000}%
\pgfsetlinewidth{0.200750pt}%
\definecolor{currentstroke}{rgb}{0.200000,0.200000,0.200000}%
\pgfsetstrokecolor{currentstroke}%
\pgfsetdash{}{0pt}%
\pgfpathmoveto{\pgfqpoint{3.528724in}{2.146112in}}%
\pgfpathlineto{\pgfqpoint{3.571391in}{2.085597in}}%
\pgfpathlineto{\pgfqpoint{3.613620in}{2.012840in}}%
\pgfpathlineto{\pgfqpoint{3.646247in}{1.970153in}}%
\pgfpathlineto{\pgfqpoint{3.678967in}{1.928348in}}%
\pgfpathlineto{\pgfqpoint{3.636732in}{1.996839in}}%
\pgfpathlineto{\pgfqpoint{3.594105in}{2.056041in}}%
\pgfpathlineto{\pgfqpoint{3.561371in}{2.100704in}}%
\pgfpathlineto{\pgfqpoint{3.528724in}{2.146112in}}%
\pgfpathclose%
\pgfusepath{stroke,fill}%
\end{pgfscope}%
\begin{pgfscope}%
\pgfpathrectangle{\pgfqpoint{0.050000in}{0.050000in}}{\pgfqpoint{4.900000in}{4.900000in}}%
\pgfusepath{clip}%
\pgfsetbuttcap%
\pgfsetroundjoin%
\definecolor{currentfill}{rgb}{0.510711,0.490534,0.707436}%
\pgfsetfillcolor{currentfill}%
\pgfsetfillopacity{0.900000}%
\pgfsetlinewidth{0.200750pt}%
\definecolor{currentstroke}{rgb}{0.200000,0.200000,0.200000}%
\pgfsetstrokecolor{currentstroke}%
\pgfsetdash{}{0pt}%
\pgfpathmoveto{\pgfqpoint{3.311702in}{2.377776in}}%
\pgfpathlineto{\pgfqpoint{3.355615in}{2.371314in}}%
\pgfpathlineto{\pgfqpoint{3.398975in}{2.335843in}}%
\pgfpathlineto{\pgfqpoint{3.431289in}{2.287126in}}%
\pgfpathlineto{\pgfqpoint{3.463684in}{2.239289in}}%
\pgfpathlineto{\pgfqpoint{3.420417in}{2.278345in}}%
\pgfpathlineto{\pgfqpoint{3.376601in}{2.291471in}}%
\pgfpathlineto{\pgfqpoint{3.344112in}{2.334494in}}%
\pgfpathlineto{\pgfqpoint{3.311702in}{2.377776in}}%
\pgfpathclose%
\pgfusepath{stroke,fill}%
\end{pgfscope}%
\begin{pgfscope}%
\pgfpathrectangle{\pgfqpoint{0.050000in}{0.050000in}}{\pgfqpoint{4.900000in}{4.900000in}}%
\pgfusepath{clip}%
\pgfsetbuttcap%
\pgfsetroundjoin%
\definecolor{currentfill}{rgb}{0.462976,0.440830,0.678893}%
\pgfsetfillcolor{currentfill}%
\pgfsetfillopacity{0.900000}%
\pgfsetlinewidth{0.200750pt}%
\definecolor{currentstroke}{rgb}{0.200000,0.200000,0.200000}%
\pgfsetstrokecolor{currentstroke}%
\pgfsetdash{}{0pt}%
\pgfpathmoveto{\pgfqpoint{3.157828in}{2.373985in}}%
\pgfpathlineto{\pgfqpoint{3.202639in}{2.432594in}}%
\pgfpathlineto{\pgfqpoint{3.247110in}{2.465012in}}%
\pgfpathlineto{\pgfqpoint{3.279368in}{2.421292in}}%
\pgfpathlineto{\pgfqpoint{3.311702in}{2.377776in}}%
\pgfpathlineto{\pgfqpoint{3.267308in}{2.355553in}}%
\pgfpathlineto{\pgfqpoint{3.222537in}{2.307926in}}%
\pgfpathlineto{\pgfqpoint{3.190142in}{2.341185in}}%
\pgfpathlineto{\pgfqpoint{3.157828in}{2.373985in}}%
\pgfpathclose%
\pgfusepath{stroke,fill}%
\end{pgfscope}%
\begin{pgfscope}%
\pgfpathrectangle{\pgfqpoint{0.050000in}{0.050000in}}{\pgfqpoint{4.900000in}{4.900000in}}%
\pgfusepath{clip}%
\pgfsetbuttcap%
\pgfsetroundjoin%
\definecolor{currentfill}{rgb}{1.000000,1.000000,1.000000}%
\pgfsetfillcolor{currentfill}%
\pgfsetfillopacity{0.900000}%
\pgfsetlinewidth{0.200750pt}%
\definecolor{currentstroke}{rgb}{0.200000,0.200000,0.200000}%
\pgfsetstrokecolor{currentstroke}%
\pgfsetdash{}{0pt}%
\pgfpathmoveto{\pgfqpoint{4.048424in}{1.616502in}}%
\pgfpathlineto{\pgfqpoint{4.091760in}{1.630821in}}%
\pgfpathlineto{\pgfqpoint{4.125619in}{1.610362in}}%
\pgfpathlineto{\pgfqpoint{4.159582in}{1.589841in}}%
\pgfpathlineto{\pgfqpoint{4.116186in}{1.575434in}}%
\pgfpathlineto{\pgfqpoint{4.082253in}{1.596000in}}%
\pgfpathlineto{\pgfqpoint{4.048424in}{1.616502in}}%
\pgfpathclose%
\pgfusepath{stroke,fill}%
\end{pgfscope}%
\begin{pgfscope}%
\pgfpathrectangle{\pgfqpoint{0.050000in}{0.050000in}}{\pgfqpoint{4.900000in}{4.900000in}}%
\pgfusepath{clip}%
\pgfsetbuttcap%
\pgfsetroundjoin%
\definecolor{currentfill}{rgb}{0.665852,0.652072,0.800200}%
\pgfsetfillcolor{currentfill}%
\pgfsetfillopacity{0.900000}%
\pgfsetlinewidth{0.200750pt}%
\definecolor{currentstroke}{rgb}{0.200000,0.200000,0.200000}%
\pgfsetstrokecolor{currentstroke}%
\pgfsetdash{}{0pt}%
\pgfpathmoveto{\pgfqpoint{2.977443in}{2.017607in}}%
\pgfpathlineto{\pgfqpoint{3.022538in}{2.110777in}}%
\pgfpathlineto{\pgfqpoint{3.067677in}{2.205630in}}%
\pgfpathlineto{\pgfqpoint{3.099959in}{2.182843in}}%
\pgfpathlineto{\pgfqpoint{3.132331in}{2.159319in}}%
\pgfpathlineto{\pgfqpoint{3.087104in}{2.071680in}}%
\pgfpathlineto{\pgfqpoint{3.041893in}{1.983520in}}%
\pgfpathlineto{\pgfqpoint{3.009619in}{2.000885in}}%
\pgfpathlineto{\pgfqpoint{2.977443in}{2.017607in}}%
\pgfpathclose%
\pgfusepath{stroke,fill}%
\end{pgfscope}%
\begin{pgfscope}%
\pgfpathrectangle{\pgfqpoint{0.050000in}{0.050000in}}{\pgfqpoint{4.900000in}{4.900000in}}%
\pgfusepath{clip}%
\pgfsetbuttcap%
\pgfsetroundjoin%
\definecolor{currentfill}{rgb}{1.000000,1.000000,1.000000}%
\pgfsetfillcolor{currentfill}%
\pgfsetfillopacity{0.900000}%
\pgfsetlinewidth{0.200750pt}%
\definecolor{currentstroke}{rgb}{0.200000,0.200000,0.200000}%
\pgfsetstrokecolor{currentstroke}%
\pgfsetdash{}{0pt}%
\pgfpathmoveto{\pgfqpoint{1.915773in}{1.608172in}}%
\pgfpathlineto{\pgfqpoint{1.961433in}{1.622509in}}%
\pgfpathlineto{\pgfqpoint{2.006995in}{1.636815in}}%
\pgfpathlineto{\pgfqpoint{2.037629in}{1.616375in}}%
\pgfpathlineto{\pgfqpoint{2.068359in}{1.595873in}}%
\pgfpathlineto{\pgfqpoint{2.022723in}{1.581478in}}%
\pgfpathlineto{\pgfqpoint{1.976988in}{1.567053in}}%
\pgfpathlineto{\pgfqpoint{1.946333in}{1.587644in}}%
\pgfpathlineto{\pgfqpoint{1.915773in}{1.608172in}}%
\pgfpathclose%
\pgfusepath{stroke,fill}%
\end{pgfscope}%
\begin{pgfscope}%
\pgfpathrectangle{\pgfqpoint{0.050000in}{0.050000in}}{\pgfqpoint{4.900000in}{4.900000in}}%
\pgfusepath{clip}%
\pgfsetbuttcap%
\pgfsetroundjoin%
\definecolor{currentfill}{rgb}{1.000000,1.000000,1.000000}%
\pgfsetfillcolor{currentfill}%
\pgfsetfillopacity{0.900000}%
\pgfsetlinewidth{0.200750pt}%
\definecolor{currentstroke}{rgb}{0.200000,0.200000,0.200000}%
\pgfsetstrokecolor{currentstroke}%
\pgfsetdash{}{0pt}%
\pgfpathmoveto{\pgfqpoint{1.519342in}{1.604046in}}%
\pgfpathlineto{\pgfqpoint{1.565443in}{1.618392in}}%
\pgfpathlineto{\pgfqpoint{1.611445in}{1.632706in}}%
\pgfpathlineto{\pgfqpoint{1.641480in}{1.612254in}}%
\pgfpathlineto{\pgfqpoint{1.671608in}{1.591738in}}%
\pgfpathlineto{\pgfqpoint{1.625529in}{1.577335in}}%
\pgfpathlineto{\pgfqpoint{1.579350in}{1.562901in}}%
\pgfpathlineto{\pgfqpoint{1.549300in}{1.583505in}}%
\pgfpathlineto{\pgfqpoint{1.519342in}{1.604046in}}%
\pgfpathclose%
\pgfusepath{stroke,fill}%
\end{pgfscope}%
\begin{pgfscope}%
\pgfpathrectangle{\pgfqpoint{0.050000in}{0.050000in}}{\pgfqpoint{4.900000in}{4.900000in}}%
\pgfusepath{clip}%
\pgfsetbuttcap%
\pgfsetroundjoin%
\definecolor{currentfill}{rgb}{0.540546,0.521599,0.725275}%
\pgfsetfillcolor{currentfill}%
\pgfsetfillopacity{0.900000}%
\pgfsetlinewidth{0.200750pt}%
\definecolor{currentstroke}{rgb}{0.200000,0.200000,0.200000}%
\pgfsetstrokecolor{currentstroke}%
\pgfsetdash{}{0pt}%
\pgfpathmoveto{\pgfqpoint{3.067677in}{2.205630in}}%
\pgfpathlineto{\pgfqpoint{3.112803in}{2.295819in}}%
\pgfpathlineto{\pgfqpoint{3.157828in}{2.373985in}}%
\pgfpathlineto{\pgfqpoint{3.190142in}{2.341185in}}%
\pgfpathlineto{\pgfqpoint{3.222537in}{2.307926in}}%
\pgfpathlineto{\pgfqpoint{3.177507in}{2.240320in}}%
\pgfpathlineto{\pgfqpoint{3.132331in}{2.159319in}}%
\pgfpathlineto{\pgfqpoint{3.099959in}{2.182843in}}%
\pgfpathlineto{\pgfqpoint{3.067677in}{2.205630in}}%
\pgfpathclose%
\pgfusepath{stroke,fill}%
\end{pgfscope}%
\begin{pgfscope}%
\pgfpathrectangle{\pgfqpoint{0.050000in}{0.050000in}}{\pgfqpoint{4.900000in}{4.900000in}}%
\pgfusepath{clip}%
\pgfsetbuttcap%
\pgfsetroundjoin%
\definecolor{currentfill}{rgb}{0.946298,0.944083,0.967889}%
\pgfsetfillcolor{currentfill}%
\pgfsetfillopacity{0.900000}%
\pgfsetlinewidth{0.200750pt}%
\definecolor{currentstroke}{rgb}{0.200000,0.200000,0.200000}%
\pgfsetstrokecolor{currentstroke}%
\pgfsetdash{}{0pt}%
\pgfpathmoveto{\pgfqpoint{2.707991in}{1.652725in}}%
\pgfpathlineto{\pgfqpoint{2.752871in}{1.688563in}}%
\pgfpathlineto{\pgfqpoint{2.797725in}{1.733119in}}%
\pgfpathlineto{\pgfqpoint{2.829629in}{1.717374in}}%
\pgfpathlineto{\pgfqpoint{2.861636in}{1.701475in}}%
\pgfpathlineto{\pgfqpoint{2.816663in}{1.655228in}}%
\pgfpathlineto{\pgfqpoint{2.771677in}{1.617571in}}%
\pgfpathlineto{\pgfqpoint{2.739783in}{1.635162in}}%
\pgfpathlineto{\pgfqpoint{2.707991in}{1.652725in}}%
\pgfpathclose%
\pgfusepath{stroke,fill}%
\end{pgfscope}%
\begin{pgfscope}%
\pgfpathrectangle{\pgfqpoint{0.050000in}{0.050000in}}{\pgfqpoint{4.900000in}{4.900000in}}%
\pgfusepath{clip}%
\pgfsetbuttcap%
\pgfsetroundjoin%
\definecolor{currentfill}{rgb}{1.000000,1.000000,1.000000}%
\pgfsetfillcolor{currentfill}%
\pgfsetfillopacity{0.900000}%
\pgfsetlinewidth{0.200750pt}%
\definecolor{currentstroke}{rgb}{0.200000,0.200000,0.200000}%
\pgfsetstrokecolor{currentstroke}%
\pgfsetdash{}{0pt}%
\pgfpathmoveto{\pgfqpoint{1.122664in}{1.599918in}}%
\pgfpathlineto{\pgfqpoint{1.169208in}{1.614272in}}%
\pgfpathlineto{\pgfqpoint{1.215651in}{1.628595in}}%
\pgfpathlineto{\pgfqpoint{1.245086in}{1.608130in}}%
\pgfpathlineto{\pgfqpoint{1.274611in}{1.587602in}}%
\pgfpathlineto{\pgfqpoint{1.228089in}{1.573190in}}%
\pgfpathlineto{\pgfqpoint{1.181465in}{1.558747in}}%
\pgfpathlineto{\pgfqpoint{1.152019in}{1.579364in}}%
\pgfpathlineto{\pgfqpoint{1.122664in}{1.599918in}}%
\pgfpathclose%
\pgfusepath{stroke,fill}%
\end{pgfscope}%
\begin{pgfscope}%
\pgfpathrectangle{\pgfqpoint{0.050000in}{0.050000in}}{\pgfqpoint{4.900000in}{4.900000in}}%
\pgfusepath{clip}%
\pgfsetbuttcap%
\pgfsetroundjoin%
\definecolor{currentfill}{rgb}{1.000000,1.000000,1.000000}%
\pgfsetfillcolor{currentfill}%
\pgfsetfillopacity{0.900000}%
\pgfsetlinewidth{0.200750pt}%
\definecolor{currentstroke}{rgb}{0.200000,0.200000,0.200000}%
\pgfsetstrokecolor{currentstroke}%
\pgfsetdash{}{0pt}%
\pgfpathmoveto{\pgfqpoint{2.464859in}{1.602475in}}%
\pgfpathlineto{\pgfqpoint{2.509954in}{1.619412in}}%
\pgfpathlineto{\pgfqpoint{2.554956in}{1.638416in}}%
\pgfpathlineto{\pgfqpoint{2.586460in}{1.619247in}}%
\pgfpathlineto{\pgfqpoint{2.618062in}{1.600089in}}%
\pgfpathlineto{\pgfqpoint{2.572981in}{1.579985in}}%
\pgfpathlineto{\pgfqpoint{2.527811in}{1.562266in}}%
\pgfpathlineto{\pgfqpoint{2.496286in}{1.582379in}}%
\pgfpathlineto{\pgfqpoint{2.464859in}{1.602475in}}%
\pgfpathclose%
\pgfusepath{stroke,fill}%
\end{pgfscope}%
\begin{pgfscope}%
\pgfpathrectangle{\pgfqpoint{0.050000in}{0.050000in}}{\pgfqpoint{4.900000in}{4.900000in}}%
\pgfusepath{clip}%
\pgfsetbuttcap%
\pgfsetroundjoin%
\definecolor{currentfill}{rgb}{1.000000,1.000000,1.000000}%
\pgfsetfillcolor{currentfill}%
\pgfsetfillopacity{0.900000}%
\pgfsetlinewidth{0.200750pt}%
\definecolor{currentstroke}{rgb}{0.200000,0.200000,0.200000}%
\pgfsetstrokecolor{currentstroke}%
\pgfsetdash{}{0pt}%
\pgfpathmoveto{\pgfqpoint{2.068359in}{1.595873in}}%
\pgfpathlineto{\pgfqpoint{2.113896in}{1.610240in}}%
\pgfpathlineto{\pgfqpoint{2.159336in}{1.624584in}}%
\pgfpathlineto{\pgfqpoint{2.190234in}{1.604115in}}%
\pgfpathlineto{\pgfqpoint{2.221227in}{1.583586in}}%
\pgfpathlineto{\pgfqpoint{2.175714in}{1.569139in}}%
\pgfpathlineto{\pgfqpoint{2.130102in}{1.554679in}}%
\pgfpathlineto{\pgfqpoint{2.099183in}{1.575308in}}%
\pgfpathlineto{\pgfqpoint{2.068359in}{1.595873in}}%
\pgfpathclose%
\pgfusepath{stroke,fill}%
\end{pgfscope}%
\begin{pgfscope}%
\pgfpathrectangle{\pgfqpoint{0.050000in}{0.050000in}}{\pgfqpoint{4.900000in}{4.900000in}}%
\pgfusepath{clip}%
\pgfsetbuttcap%
\pgfsetroundjoin%
\definecolor{currentfill}{rgb}{1.000000,1.000000,1.000000}%
\pgfsetfillcolor{currentfill}%
\pgfsetfillopacity{0.900000}%
\pgfsetlinewidth{0.200750pt}%
\definecolor{currentstroke}{rgb}{0.200000,0.200000,0.200000}%
\pgfsetstrokecolor{currentstroke}%
\pgfsetdash{}{0pt}%
\pgfpathmoveto{\pgfqpoint{1.671608in}{1.591738in}}%
\pgfpathlineto{\pgfqpoint{1.717588in}{1.606110in}}%
\pgfpathlineto{\pgfqpoint{1.763469in}{1.620451in}}%
\pgfpathlineto{\pgfqpoint{1.793766in}{1.599960in}}%
\pgfpathlineto{\pgfqpoint{1.824157in}{1.579407in}}%
\pgfpathlineto{\pgfqpoint{1.778200in}{1.564977in}}%
\pgfpathlineto{\pgfqpoint{1.732144in}{1.550516in}}%
\pgfpathlineto{\pgfqpoint{1.701829in}{1.571159in}}%
\pgfpathlineto{\pgfqpoint{1.671608in}{1.591738in}}%
\pgfpathclose%
\pgfusepath{stroke,fill}%
\end{pgfscope}%
\begin{pgfscope}%
\pgfpathrectangle{\pgfqpoint{0.050000in}{0.050000in}}{\pgfqpoint{4.900000in}{4.900000in}}%
\pgfusepath{clip}%
\pgfsetbuttcap%
\pgfsetroundjoin%
\definecolor{currentfill}{rgb}{1.000000,1.000000,1.000000}%
\pgfsetfillcolor{currentfill}%
\pgfsetfillopacity{0.900000}%
\pgfsetlinewidth{0.200750pt}%
\definecolor{currentstroke}{rgb}{0.200000,0.200000,0.200000}%
\pgfsetstrokecolor{currentstroke}%
\pgfsetdash{}{0pt}%
\pgfpathmoveto{\pgfqpoint{1.274611in}{1.587602in}}%
\pgfpathlineto{\pgfqpoint{1.321034in}{1.601982in}}%
\pgfpathlineto{\pgfqpoint{1.367356in}{1.616332in}}%
\pgfpathlineto{\pgfqpoint{1.397052in}{1.595829in}}%
\pgfpathlineto{\pgfqpoint{1.426840in}{1.575263in}}%
\pgfpathlineto{\pgfqpoint{1.380439in}{1.560824in}}%
\pgfpathlineto{\pgfqpoint{1.333937in}{1.546354in}}%
\pgfpathlineto{\pgfqpoint{1.304228in}{1.567010in}}%
\pgfpathlineto{\pgfqpoint{1.274611in}{1.587602in}}%
\pgfpathclose%
\pgfusepath{stroke,fill}%
\end{pgfscope}%
\begin{pgfscope}%
\pgfpathrectangle{\pgfqpoint{0.050000in}{0.050000in}}{\pgfqpoint{4.900000in}{4.900000in}}%
\pgfusepath{clip}%
\pgfsetbuttcap%
\pgfsetroundjoin%
\definecolor{currentfill}{rgb}{1.000000,1.000000,1.000000}%
\pgfsetfillcolor{currentfill}%
\pgfsetfillopacity{0.900000}%
\pgfsetlinewidth{0.200750pt}%
\definecolor{currentstroke}{rgb}{0.200000,0.200000,0.200000}%
\pgfsetstrokecolor{currentstroke}%
\pgfsetdash{}{0pt}%
\pgfpathmoveto{\pgfqpoint{3.961877in}{1.603969in}}%
\pgfpathlineto{\pgfqpoint{4.004995in}{1.602153in}}%
\pgfpathlineto{\pgfqpoint{4.048424in}{1.616502in}}%
\pgfpathlineto{\pgfqpoint{4.082253in}{1.596000in}}%
\pgfpathlineto{\pgfqpoint{4.116186in}{1.575434in}}%
\pgfpathlineto{\pgfqpoint{4.072696in}{1.560995in}}%
\pgfpathlineto{\pgfqpoint{4.029111in}{1.546526in}}%
\pgfpathlineto{\pgfqpoint{3.995395in}{1.573267in}}%
\pgfpathlineto{\pgfqpoint{3.961877in}{1.603969in}}%
\pgfpathclose%
\pgfusepath{stroke,fill}%
\end{pgfscope}%
\begin{pgfscope}%
\pgfpathrectangle{\pgfqpoint{0.050000in}{0.050000in}}{\pgfqpoint{4.900000in}{4.900000in}}%
\pgfusepath{clip}%
\pgfsetbuttcap%
\pgfsetroundjoin%
\definecolor{currentfill}{rgb}{1.000000,1.000000,1.000000}%
\pgfsetfillcolor{currentfill}%
\pgfsetfillopacity{0.900000}%
\pgfsetlinewidth{0.200750pt}%
\definecolor{currentstroke}{rgb}{0.200000,0.200000,0.200000}%
\pgfsetstrokecolor{currentstroke}%
\pgfsetdash{}{0pt}%
\pgfpathmoveto{\pgfqpoint{2.221227in}{1.583586in}}%
\pgfpathlineto{\pgfqpoint{2.266641in}{1.598047in}}%
\pgfpathlineto{\pgfqpoint{2.311957in}{1.612591in}}%
\pgfpathlineto{\pgfqpoint{2.343119in}{1.592163in}}%
\pgfpathlineto{\pgfqpoint{2.374376in}{1.571688in}}%
\pgfpathlineto{\pgfqpoint{2.328988in}{1.556948in}}%
\pgfpathlineto{\pgfqpoint{2.283500in}{1.542347in}}%
\pgfpathlineto{\pgfqpoint{2.252315in}{1.562996in}}%
\pgfpathlineto{\pgfqpoint{2.221227in}{1.583586in}}%
\pgfpathclose%
\pgfusepath{stroke,fill}%
\end{pgfscope}%
\begin{pgfscope}%
\pgfpathrectangle{\pgfqpoint{0.050000in}{0.050000in}}{\pgfqpoint{4.900000in}{4.900000in}}%
\pgfusepath{clip}%
\pgfsetbuttcap%
\pgfsetroundjoin%
\definecolor{currentfill}{rgb}{1.000000,1.000000,1.000000}%
\pgfsetfillcolor{currentfill}%
\pgfsetfillopacity{0.900000}%
\pgfsetlinewidth{0.200750pt}%
\definecolor{currentstroke}{rgb}{0.200000,0.200000,0.200000}%
\pgfsetstrokecolor{currentstroke}%
\pgfsetdash{}{0pt}%
\pgfpathmoveto{\pgfqpoint{0.877368in}{1.583463in}}%
\pgfpathlineto{\pgfqpoint{0.924233in}{1.597852in}}%
\pgfpathlineto{\pgfqpoint{0.970998in}{1.612211in}}%
\pgfpathlineto{\pgfqpoint{1.000092in}{1.591695in}}%
\pgfpathlineto{\pgfqpoint{1.029276in}{1.571116in}}%
\pgfpathlineto{\pgfqpoint{0.982430in}{1.556669in}}%
\pgfpathlineto{\pgfqpoint{0.935483in}{1.542190in}}%
\pgfpathlineto{\pgfqpoint{0.906380in}{1.562858in}}%
\pgfpathlineto{\pgfqpoint{0.877368in}{1.583463in}}%
\pgfpathclose%
\pgfusepath{stroke,fill}%
\end{pgfscope}%
\begin{pgfscope}%
\pgfpathrectangle{\pgfqpoint{0.050000in}{0.050000in}}{\pgfqpoint{4.900000in}{4.900000in}}%
\pgfusepath{clip}%
\pgfsetbuttcap%
\pgfsetroundjoin%
\definecolor{currentfill}{rgb}{0.558447,0.540238,0.735978}%
\pgfsetfillcolor{currentfill}%
\pgfsetfillopacity{0.900000}%
\pgfsetlinewidth{0.200750pt}%
\definecolor{currentstroke}{rgb}{0.200000,0.200000,0.200000}%
\pgfsetstrokecolor{currentstroke}%
\pgfsetdash{}{0pt}%
\pgfpathmoveto{\pgfqpoint{3.376601in}{2.291471in}}%
\pgfpathlineto{\pgfqpoint{3.420417in}{2.278345in}}%
\pgfpathlineto{\pgfqpoint{3.463684in}{2.239289in}}%
\pgfpathlineto{\pgfqpoint{3.496162in}{2.192295in}}%
\pgfpathlineto{\pgfqpoint{3.528724in}{2.146112in}}%
\pgfpathlineto{\pgfqpoint{3.485539in}{2.187710in}}%
\pgfpathlineto{\pgfqpoint{3.441815in}{2.206268in}}%
\pgfpathlineto{\pgfqpoint{3.409168in}{2.248724in}}%
\pgfpathlineto{\pgfqpoint{3.376601in}{2.291471in}}%
\pgfpathclose%
\pgfusepath{stroke,fill}%
\end{pgfscope}%
\begin{pgfscope}%
\pgfpathrectangle{\pgfqpoint{0.050000in}{0.050000in}}{\pgfqpoint{4.900000in}{4.900000in}}%
\pgfusepath{clip}%
\pgfsetbuttcap%
\pgfsetroundjoin%
\definecolor{currentfill}{rgb}{0.982099,0.981361,0.989296}%
\pgfsetfillcolor{currentfill}%
\pgfsetfillopacity{0.900000}%
\pgfsetlinewidth{0.200750pt}%
\definecolor{currentstroke}{rgb}{0.200000,0.200000,0.200000}%
\pgfsetstrokecolor{currentstroke}%
\pgfsetdash{}{0pt}%
\pgfpathmoveto{\pgfqpoint{2.618062in}{1.600089in}}%
\pgfpathlineto{\pgfqpoint{2.663061in}{1.623820in}}%
\pgfpathlineto{\pgfqpoint{2.707991in}{1.652725in}}%
\pgfpathlineto{\pgfqpoint{2.739783in}{1.635162in}}%
\pgfpathlineto{\pgfqpoint{2.771677in}{1.617571in}}%
\pgfpathlineto{\pgfqpoint{2.726651in}{1.586963in}}%
\pgfpathlineto{\pgfqpoint{2.681565in}{1.561786in}}%
\pgfpathlineto{\pgfqpoint{2.649763in}{1.580937in}}%
\pgfpathlineto{\pgfqpoint{2.618062in}{1.600089in}}%
\pgfpathclose%
\pgfusepath{stroke,fill}%
\end{pgfscope}%
\begin{pgfscope}%
\pgfpathrectangle{\pgfqpoint{0.050000in}{0.050000in}}{\pgfqpoint{4.900000in}{4.900000in}}%
\pgfusepath{clip}%
\pgfsetbuttcap%
\pgfsetroundjoin%
\definecolor{currentfill}{rgb}{0.743422,0.732841,0.846582}%
\pgfsetfillcolor{currentfill}%
\pgfsetfillopacity{0.900000}%
\pgfsetlinewidth{0.200750pt}%
\definecolor{currentstroke}{rgb}{0.200000,0.200000,0.200000}%
\pgfsetstrokecolor{currentstroke}%
\pgfsetdash{}{0pt}%
\pgfpathmoveto{\pgfqpoint{3.594105in}{2.056041in}}%
\pgfpathlineto{\pgfqpoint{3.636732in}{1.996839in}}%
\pgfpathlineto{\pgfqpoint{3.678967in}{1.928348in}}%
\pgfpathlineto{\pgfqpoint{3.711781in}{1.887371in}}%
\pgfpathlineto{\pgfqpoint{3.744690in}{1.847172in}}%
\pgfpathlineto{\pgfqpoint{3.702437in}{1.911338in}}%
\pgfpathlineto{\pgfqpoint{3.659838in}{1.968820in}}%
\pgfpathlineto{\pgfqpoint{3.626927in}{2.012090in}}%
\pgfpathlineto{\pgfqpoint{3.594105in}{2.056041in}}%
\pgfpathclose%
\pgfusepath{stroke,fill}%
\end{pgfscope}%
\begin{pgfscope}%
\pgfpathrectangle{\pgfqpoint{0.050000in}{0.050000in}}{\pgfqpoint{4.900000in}{4.900000in}}%
\pgfusepath{clip}%
\pgfsetbuttcap%
\pgfsetroundjoin%
\definecolor{currentfill}{rgb}{1.000000,1.000000,1.000000}%
\pgfsetfillcolor{currentfill}%
\pgfsetfillopacity{0.900000}%
\pgfsetlinewidth{0.200750pt}%
\definecolor{currentstroke}{rgb}{0.200000,0.200000,0.200000}%
\pgfsetstrokecolor{currentstroke}%
\pgfsetdash{}{0pt}%
\pgfpathmoveto{\pgfqpoint{1.824157in}{1.579407in}}%
\pgfpathlineto{\pgfqpoint{1.870014in}{1.593805in}}%
\pgfpathlineto{\pgfqpoint{1.915773in}{1.608172in}}%
\pgfpathlineto{\pgfqpoint{1.946333in}{1.587644in}}%
\pgfpathlineto{\pgfqpoint{1.976988in}{1.567053in}}%
\pgfpathlineto{\pgfqpoint{1.931154in}{1.552596in}}%
\pgfpathlineto{\pgfqpoint{1.885221in}{1.538109in}}%
\pgfpathlineto{\pgfqpoint{1.854641in}{1.558790in}}%
\pgfpathlineto{\pgfqpoint{1.824157in}{1.579407in}}%
\pgfpathclose%
\pgfusepath{stroke,fill}%
\end{pgfscope}%
\begin{pgfscope}%
\pgfpathrectangle{\pgfqpoint{0.050000in}{0.050000in}}{\pgfqpoint{4.900000in}{4.900000in}}%
\pgfusepath{clip}%
\pgfsetbuttcap%
\pgfsetroundjoin%
\definecolor{currentfill}{rgb}{0.928397,0.925444,0.957186}%
\pgfsetfillcolor{currentfill}%
\pgfsetfillopacity{0.900000}%
\pgfsetlinewidth{0.200750pt}%
\definecolor{currentstroke}{rgb}{0.200000,0.200000,0.200000}%
\pgfsetstrokecolor{currentstroke}%
\pgfsetdash{}{0pt}%
\pgfpathmoveto{\pgfqpoint{3.810797in}{1.768927in}}%
\pgfpathlineto{\pgfqpoint{3.852942in}{1.712210in}}%
\pgfpathlineto{\pgfqpoint{3.895148in}{1.666583in}}%
\pgfpathlineto{\pgfqpoint{3.928461in}{1.635065in}}%
\pgfpathlineto{\pgfqpoint{3.961877in}{1.603969in}}%
\pgfpathlineto{\pgfqpoint{3.919530in}{1.642522in}}%
\pgfpathlineto{\pgfqpoint{3.877293in}{1.693278in}}%
\pgfpathlineto{\pgfqpoint{3.843996in}{1.730797in}}%
\pgfpathlineto{\pgfqpoint{3.810797in}{1.768927in}}%
\pgfpathclose%
\pgfusepath{stroke,fill}%
\end{pgfscope}%
\begin{pgfscope}%
\pgfpathrectangle{\pgfqpoint{0.050000in}{0.050000in}}{\pgfqpoint{4.900000in}{4.900000in}}%
\pgfusepath{clip}%
\pgfsetbuttcap%
\pgfsetroundjoin%
\definecolor{currentfill}{rgb}{0.492810,0.471895,0.696732}%
\pgfsetfillcolor{currentfill}%
\pgfsetfillopacity{0.900000}%
\pgfsetlinewidth{0.200750pt}%
\definecolor{currentstroke}{rgb}{0.200000,0.200000,0.200000}%
\pgfsetstrokecolor{currentstroke}%
\pgfsetdash{}{0pt}%
\pgfpathmoveto{\pgfqpoint{3.222537in}{2.307926in}}%
\pgfpathlineto{\pgfqpoint{3.267308in}{2.355553in}}%
\pgfpathlineto{\pgfqpoint{3.311702in}{2.377776in}}%
\pgfpathlineto{\pgfqpoint{3.344112in}{2.334494in}}%
\pgfpathlineto{\pgfqpoint{3.376601in}{2.291471in}}%
\pgfpathlineto{\pgfqpoint{3.332288in}{2.278052in}}%
\pgfpathlineto{\pgfqpoint{3.287567in}{2.240261in}}%
\pgfpathlineto{\pgfqpoint{3.255012in}{2.274267in}}%
\pgfpathlineto{\pgfqpoint{3.222537in}{2.307926in}}%
\pgfpathclose%
\pgfusepath{stroke,fill}%
\end{pgfscope}%
\begin{pgfscope}%
\pgfpathrectangle{\pgfqpoint{0.050000in}{0.050000in}}{\pgfqpoint{4.900000in}{4.900000in}}%
\pgfusepath{clip}%
\pgfsetbuttcap%
\pgfsetroundjoin%
\definecolor{currentfill}{rgb}{1.000000,1.000000,1.000000}%
\pgfsetfillcolor{currentfill}%
\pgfsetfillopacity{0.900000}%
\pgfsetlinewidth{0.200750pt}%
\definecolor{currentstroke}{rgb}{0.200000,0.200000,0.200000}%
\pgfsetstrokecolor{currentstroke}%
\pgfsetdash{}{0pt}%
\pgfpathmoveto{\pgfqpoint{1.426840in}{1.575263in}}%
\pgfpathlineto{\pgfqpoint{1.473141in}{1.589670in}}%
\pgfpathlineto{\pgfqpoint{1.519342in}{1.604046in}}%
\pgfpathlineto{\pgfqpoint{1.549300in}{1.583505in}}%
\pgfpathlineto{\pgfqpoint{1.579350in}{1.562901in}}%
\pgfpathlineto{\pgfqpoint{1.533071in}{1.548436in}}%
\pgfpathlineto{\pgfqpoint{1.486692in}{1.533939in}}%
\pgfpathlineto{\pgfqpoint{1.456720in}{1.554633in}}%
\pgfpathlineto{\pgfqpoint{1.426840in}{1.575263in}}%
\pgfpathclose%
\pgfusepath{stroke,fill}%
\end{pgfscope}%
\begin{pgfscope}%
\pgfpathrectangle{\pgfqpoint{0.050000in}{0.050000in}}{\pgfqpoint{4.900000in}{4.900000in}}%
\pgfusepath{clip}%
\pgfsetbuttcap%
\pgfsetroundjoin%
\definecolor{currentfill}{rgb}{0.779223,0.770119,0.867989}%
\pgfsetfillcolor{currentfill}%
\pgfsetfillopacity{0.900000}%
\pgfsetlinewidth{0.200750pt}%
\definecolor{currentstroke}{rgb}{0.200000,0.200000,0.200000}%
\pgfsetstrokecolor{currentstroke}%
\pgfsetdash{}{0pt}%
\pgfpathmoveto{\pgfqpoint{2.951649in}{1.823769in}}%
\pgfpathlineto{\pgfqpoint{2.996737in}{1.899744in}}%
\pgfpathlineto{\pgfqpoint{3.041893in}{1.983520in}}%
\pgfpathlineto{\pgfqpoint{3.074265in}{1.965524in}}%
\pgfpathlineto{\pgfqpoint{3.106734in}{1.946915in}}%
\pgfpathlineto{\pgfqpoint{3.061456in}{1.866464in}}%
\pgfpathlineto{\pgfqpoint{3.016236in}{1.792067in}}%
\pgfpathlineto{\pgfqpoint{2.983891in}{1.808118in}}%
\pgfpathlineto{\pgfqpoint{2.951649in}{1.823769in}}%
\pgfpathclose%
\pgfusepath{stroke,fill}%
\end{pgfscope}%
\begin{pgfscope}%
\pgfpathrectangle{\pgfqpoint{0.050000in}{0.050000in}}{\pgfqpoint{4.900000in}{4.900000in}}%
\pgfusepath{clip}%
\pgfsetbuttcap%
\pgfsetroundjoin%
\definecolor{currentfill}{rgb}{1.000000,1.000000,1.000000}%
\pgfsetfillcolor{currentfill}%
\pgfsetfillopacity{0.900000}%
\pgfsetlinewidth{0.200750pt}%
\definecolor{currentstroke}{rgb}{0.200000,0.200000,0.200000}%
\pgfsetstrokecolor{currentstroke}%
\pgfsetdash{}{0pt}%
\pgfpathmoveto{\pgfqpoint{1.029276in}{1.571116in}}%
\pgfpathlineto{\pgfqpoint{1.076021in}{1.585532in}}%
\pgfpathlineto{\pgfqpoint{1.122664in}{1.599918in}}%
\pgfpathlineto{\pgfqpoint{1.152019in}{1.579364in}}%
\pgfpathlineto{\pgfqpoint{1.181465in}{1.558747in}}%
\pgfpathlineto{\pgfqpoint{1.134741in}{1.544272in}}%
\pgfpathlineto{\pgfqpoint{1.087915in}{1.529767in}}%
\pgfpathlineto{\pgfqpoint{1.058550in}{1.550473in}}%
\pgfpathlineto{\pgfqpoint{1.029276in}{1.571116in}}%
\pgfpathclose%
\pgfusepath{stroke,fill}%
\end{pgfscope}%
\begin{pgfscope}%
\pgfpathrectangle{\pgfqpoint{0.050000in}{0.050000in}}{\pgfqpoint{4.900000in}{4.900000in}}%
\pgfusepath{clip}%
\pgfsetbuttcap%
\pgfsetroundjoin%
\definecolor{currentfill}{rgb}{1.000000,1.000000,1.000000}%
\pgfsetfillcolor{currentfill}%
\pgfsetfillopacity{0.900000}%
\pgfsetlinewidth{0.200750pt}%
\definecolor{currentstroke}{rgb}{0.200000,0.200000,0.200000}%
\pgfsetstrokecolor{currentstroke}%
\pgfsetdash{}{0pt}%
\pgfpathmoveto{\pgfqpoint{2.374376in}{1.571688in}}%
\pgfpathlineto{\pgfqpoint{2.419666in}{1.586752in}}%
\pgfpathlineto{\pgfqpoint{2.464859in}{1.602475in}}%
\pgfpathlineto{\pgfqpoint{2.496286in}{1.582379in}}%
\pgfpathlineto{\pgfqpoint{2.527811in}{1.562266in}}%
\pgfpathlineto{\pgfqpoint{2.482545in}{1.546015in}}%
\pgfpathlineto{\pgfqpoint{2.437183in}{1.530603in}}%
\pgfpathlineto{\pgfqpoint{2.405731in}{1.551168in}}%
\pgfpathlineto{\pgfqpoint{2.374376in}{1.571688in}}%
\pgfpathclose%
\pgfusepath{stroke,fill}%
\end{pgfscope}%
\begin{pgfscope}%
\pgfpathrectangle{\pgfqpoint{0.050000in}{0.050000in}}{\pgfqpoint{4.900000in}{4.900000in}}%
\pgfusepath{clip}%
\pgfsetbuttcap%
\pgfsetroundjoin%
\definecolor{currentfill}{rgb}{0.874694,0.869527,0.925075}%
\pgfsetfillcolor{currentfill}%
\pgfsetfillopacity{0.900000}%
\pgfsetlinewidth{0.200750pt}%
\definecolor{currentstroke}{rgb}{0.200000,0.200000,0.200000}%
\pgfsetstrokecolor{currentstroke}%
\pgfsetdash{}{0pt}%
\pgfpathmoveto{\pgfqpoint{2.861636in}{1.701475in}}%
\pgfpathlineto{\pgfqpoint{2.906622in}{1.757493in}}%
\pgfpathlineto{\pgfqpoint{2.951649in}{1.823769in}}%
\pgfpathlineto{\pgfqpoint{2.983891in}{1.808118in}}%
\pgfpathlineto{\pgfqpoint{3.016236in}{1.792067in}}%
\pgfpathlineto{\pgfqpoint{2.971075in}{1.725975in}}%
\pgfpathlineto{\pgfqpoint{2.925960in}{1.669158in}}%
\pgfpathlineto{\pgfqpoint{2.893746in}{1.685407in}}%
\pgfpathlineto{\pgfqpoint{2.861636in}{1.701475in}}%
\pgfpathclose%
\pgfusepath{stroke,fill}%
\end{pgfscope}%
\begin{pgfscope}%
\pgfpathrectangle{\pgfqpoint{0.050000in}{0.050000in}}{\pgfqpoint{4.900000in}{4.900000in}}%
\pgfusepath{clip}%
\pgfsetbuttcap%
\pgfsetroundjoin%
\definecolor{currentfill}{rgb}{1.000000,1.000000,1.000000}%
\pgfsetfillcolor{currentfill}%
\pgfsetfillopacity{0.900000}%
\pgfsetlinewidth{0.200750pt}%
\definecolor{currentstroke}{rgb}{0.200000,0.200000,0.200000}%
\pgfsetstrokecolor{currentstroke}%
\pgfsetdash{}{0pt}%
\pgfpathmoveto{\pgfqpoint{4.116186in}{1.575434in}}%
\pgfpathlineto{\pgfqpoint{4.159582in}{1.589841in}}%
\pgfpathlineto{\pgfqpoint{4.193650in}{1.569256in}}%
\pgfpathlineto{\pgfqpoint{4.227823in}{1.548607in}}%
\pgfpathlineto{\pgfqpoint{4.184367in}{1.534111in}}%
\pgfpathlineto{\pgfqpoint{4.150224in}{1.554804in}}%
\pgfpathlineto{\pgfqpoint{4.116186in}{1.575434in}}%
\pgfpathclose%
\pgfusepath{stroke,fill}%
\end{pgfscope}%
\begin{pgfscope}%
\pgfpathrectangle{\pgfqpoint{0.050000in}{0.050000in}}{\pgfqpoint{4.900000in}{4.900000in}}%
\pgfusepath{clip}%
\pgfsetbuttcap%
\pgfsetroundjoin%
\definecolor{currentfill}{rgb}{0.665852,0.652072,0.800200}%
\pgfsetfillcolor{currentfill}%
\pgfsetfillopacity{0.900000}%
\pgfsetlinewidth{0.200750pt}%
\definecolor{currentstroke}{rgb}{0.200000,0.200000,0.200000}%
\pgfsetstrokecolor{currentstroke}%
\pgfsetdash{}{0pt}%
\pgfpathmoveto{\pgfqpoint{3.041893in}{1.983520in}}%
\pgfpathlineto{\pgfqpoint{3.087104in}{2.071680in}}%
\pgfpathlineto{\pgfqpoint{3.132331in}{2.159319in}}%
\pgfpathlineto{\pgfqpoint{3.164792in}{2.135103in}}%
\pgfpathlineto{\pgfqpoint{3.197342in}{2.110240in}}%
\pgfpathlineto{\pgfqpoint{3.152045in}{2.029802in}}%
\pgfpathlineto{\pgfqpoint{3.106734in}{1.946915in}}%
\pgfpathlineto{\pgfqpoint{3.074265in}{1.965524in}}%
\pgfpathlineto{\pgfqpoint{3.041893in}{1.983520in}}%
\pgfpathclose%
\pgfusepath{stroke,fill}%
\end{pgfscope}%
\begin{pgfscope}%
\pgfpathrectangle{\pgfqpoint{0.050000in}{0.050000in}}{\pgfqpoint{4.900000in}{4.900000in}}%
\pgfusepath{clip}%
\pgfsetbuttcap%
\pgfsetroundjoin%
\definecolor{currentfill}{rgb}{1.000000,1.000000,1.000000}%
\pgfsetfillcolor{currentfill}%
\pgfsetfillopacity{0.900000}%
\pgfsetlinewidth{0.200750pt}%
\definecolor{currentstroke}{rgb}{0.200000,0.200000,0.200000}%
\pgfsetstrokecolor{currentstroke}%
\pgfsetdash{}{0pt}%
\pgfpathmoveto{\pgfqpoint{1.976988in}{1.567053in}}%
\pgfpathlineto{\pgfqpoint{2.022723in}{1.581478in}}%
\pgfpathlineto{\pgfqpoint{2.068359in}{1.595873in}}%
\pgfpathlineto{\pgfqpoint{2.099183in}{1.575308in}}%
\pgfpathlineto{\pgfqpoint{2.130102in}{1.554679in}}%
\pgfpathlineto{\pgfqpoint{2.084392in}{1.540193in}}%
\pgfpathlineto{\pgfqpoint{2.038582in}{1.525678in}}%
\pgfpathlineto{\pgfqpoint{2.007737in}{1.546397in}}%
\pgfpathlineto{\pgfqpoint{1.976988in}{1.567053in}}%
\pgfpathclose%
\pgfusepath{stroke,fill}%
\end{pgfscope}%
\begin{pgfscope}%
\pgfpathrectangle{\pgfqpoint{0.050000in}{0.050000in}}{\pgfqpoint{4.900000in}{4.900000in}}%
\pgfusepath{clip}%
\pgfsetbuttcap%
\pgfsetroundjoin%
\definecolor{currentfill}{rgb}{0.558447,0.540238,0.735978}%
\pgfsetfillcolor{currentfill}%
\pgfsetfillopacity{0.900000}%
\pgfsetlinewidth{0.200750pt}%
\definecolor{currentstroke}{rgb}{0.200000,0.200000,0.200000}%
\pgfsetstrokecolor{currentstroke}%
\pgfsetdash{}{0pt}%
\pgfpathmoveto{\pgfqpoint{3.132331in}{2.159319in}}%
\pgfpathlineto{\pgfqpoint{3.177507in}{2.240320in}}%
\pgfpathlineto{\pgfqpoint{3.222537in}{2.307926in}}%
\pgfpathlineto{\pgfqpoint{3.255012in}{2.274267in}}%
\pgfpathlineto{\pgfqpoint{3.287567in}{2.240261in}}%
\pgfpathlineto{\pgfqpoint{3.242549in}{2.182435in}}%
\pgfpathlineto{\pgfqpoint{3.197342in}{2.110240in}}%
\pgfpathlineto{\pgfqpoint{3.164792in}{2.135103in}}%
\pgfpathlineto{\pgfqpoint{3.132331in}{2.159319in}}%
\pgfpathclose%
\pgfusepath{stroke,fill}%
\end{pgfscope}%
\begin{pgfscope}%
\pgfpathrectangle{\pgfqpoint{0.050000in}{0.050000in}}{\pgfqpoint{4.900000in}{4.900000in}}%
\pgfusepath{clip}%
\pgfsetbuttcap%
\pgfsetroundjoin%
\definecolor{currentfill}{rgb}{1.000000,1.000000,1.000000}%
\pgfsetfillcolor{currentfill}%
\pgfsetfillopacity{0.900000}%
\pgfsetlinewidth{0.200750pt}%
\definecolor{currentstroke}{rgb}{0.200000,0.200000,0.200000}%
\pgfsetstrokecolor{currentstroke}%
\pgfsetdash{}{0pt}%
\pgfpathmoveto{\pgfqpoint{1.579350in}{1.562901in}}%
\pgfpathlineto{\pgfqpoint{1.625529in}{1.577335in}}%
\pgfpathlineto{\pgfqpoint{1.671608in}{1.591738in}}%
\pgfpathlineto{\pgfqpoint{1.701829in}{1.571159in}}%
\pgfpathlineto{\pgfqpoint{1.732144in}{1.550516in}}%
\pgfpathlineto{\pgfqpoint{1.685987in}{1.536024in}}%
\pgfpathlineto{\pgfqpoint{1.639731in}{1.521500in}}%
\pgfpathlineto{\pgfqpoint{1.609494in}{1.542233in}}%
\pgfpathlineto{\pgfqpoint{1.579350in}{1.562901in}}%
\pgfpathclose%
\pgfusepath{stroke,fill}%
\end{pgfscope}%
\begin{pgfscope}%
\pgfpathrectangle{\pgfqpoint{0.050000in}{0.050000in}}{\pgfqpoint{4.900000in}{4.900000in}}%
\pgfusepath{clip}%
\pgfsetbuttcap%
\pgfsetroundjoin%
\definecolor{currentfill}{rgb}{1.000000,1.000000,1.000000}%
\pgfsetfillcolor{currentfill}%
\pgfsetfillopacity{0.900000}%
\pgfsetlinewidth{0.200750pt}%
\definecolor{currentstroke}{rgb}{0.200000,0.200000,0.200000}%
\pgfsetstrokecolor{currentstroke}%
\pgfsetdash{}{0pt}%
\pgfpathmoveto{\pgfqpoint{1.181465in}{1.558747in}}%
\pgfpathlineto{\pgfqpoint{1.228089in}{1.573190in}}%
\pgfpathlineto{\pgfqpoint{1.274611in}{1.587602in}}%
\pgfpathlineto{\pgfqpoint{1.304228in}{1.567010in}}%
\pgfpathlineto{\pgfqpoint{1.333937in}{1.546354in}}%
\pgfpathlineto{\pgfqpoint{1.287335in}{1.531853in}}%
\pgfpathlineto{\pgfqpoint{1.240631in}{1.517320in}}%
\pgfpathlineto{\pgfqpoint{1.211002in}{1.538066in}}%
\pgfpathlineto{\pgfqpoint{1.181465in}{1.558747in}}%
\pgfpathclose%
\pgfusepath{stroke,fill}%
\end{pgfscope}%
\begin{pgfscope}%
\pgfpathrectangle{\pgfqpoint{0.050000in}{0.050000in}}{\pgfqpoint{4.900000in}{4.900000in}}%
\pgfusepath{clip}%
\pgfsetbuttcap%
\pgfsetroundjoin%
\definecolor{currentfill}{rgb}{0.940331,0.937870,0.964321}%
\pgfsetfillcolor{currentfill}%
\pgfsetfillopacity{0.900000}%
\pgfsetlinewidth{0.200750pt}%
\definecolor{currentstroke}{rgb}{0.200000,0.200000,0.200000}%
\pgfsetstrokecolor{currentstroke}%
\pgfsetdash{}{0pt}%
\pgfpathmoveto{\pgfqpoint{2.771677in}{1.617571in}}%
\pgfpathlineto{\pgfqpoint{2.816663in}{1.655228in}}%
\pgfpathlineto{\pgfqpoint{2.861636in}{1.701475in}}%
\pgfpathlineto{\pgfqpoint{2.893746in}{1.685407in}}%
\pgfpathlineto{\pgfqpoint{2.925960in}{1.669158in}}%
\pgfpathlineto{\pgfqpoint{2.880867in}{1.621533in}}%
\pgfpathlineto{\pgfqpoint{2.835771in}{1.582261in}}%
\pgfpathlineto{\pgfqpoint{2.803673in}{1.599941in}}%
\pgfpathlineto{\pgfqpoint{2.771677in}{1.617571in}}%
\pgfpathclose%
\pgfusepath{stroke,fill}%
\end{pgfscope}%
\begin{pgfscope}%
\pgfpathrectangle{\pgfqpoint{0.050000in}{0.050000in}}{\pgfqpoint{4.900000in}{4.900000in}}%
\pgfusepath{clip}%
\pgfsetbuttcap%
\pgfsetroundjoin%
\definecolor{currentfill}{rgb}{0.994033,0.993787,0.996432}%
\pgfsetfillcolor{currentfill}%
\pgfsetfillopacity{0.900000}%
\pgfsetlinewidth{0.200750pt}%
\definecolor{currentstroke}{rgb}{0.200000,0.200000,0.200000}%
\pgfsetstrokecolor{currentstroke}%
\pgfsetdash{}{0pt}%
\pgfpathmoveto{\pgfqpoint{2.527811in}{1.562266in}}%
\pgfpathlineto{\pgfqpoint{2.572981in}{1.579985in}}%
\pgfpathlineto{\pgfqpoint{2.618062in}{1.600089in}}%
\pgfpathlineto{\pgfqpoint{2.649763in}{1.580937in}}%
\pgfpathlineto{\pgfqpoint{2.681565in}{1.561786in}}%
\pgfpathlineto{\pgfqpoint{2.636403in}{1.540545in}}%
\pgfpathlineto{\pgfqpoint{2.591156in}{1.521989in}}%
\pgfpathlineto{\pgfqpoint{2.559434in}{1.542136in}}%
\pgfpathlineto{\pgfqpoint{2.527811in}{1.562266in}}%
\pgfpathclose%
\pgfusepath{stroke,fill}%
\end{pgfscope}%
\begin{pgfscope}%
\pgfpathrectangle{\pgfqpoint{0.050000in}{0.050000in}}{\pgfqpoint{4.900000in}{4.900000in}}%
\pgfusepath{clip}%
\pgfsetbuttcap%
\pgfsetroundjoin%
\definecolor{currentfill}{rgb}{1.000000,1.000000,1.000000}%
\pgfsetfillcolor{currentfill}%
\pgfsetfillopacity{0.900000}%
\pgfsetlinewidth{0.200750pt}%
\definecolor{currentstroke}{rgb}{0.200000,0.200000,0.200000}%
\pgfsetstrokecolor{currentstroke}%
\pgfsetdash{}{0pt}%
\pgfpathmoveto{\pgfqpoint{2.130102in}{1.554679in}}%
\pgfpathlineto{\pgfqpoint{2.175714in}{1.569139in}}%
\pgfpathlineto{\pgfqpoint{2.221227in}{1.583586in}}%
\pgfpathlineto{\pgfqpoint{2.252315in}{1.562996in}}%
\pgfpathlineto{\pgfqpoint{2.283500in}{1.542347in}}%
\pgfpathlineto{\pgfqpoint{2.237914in}{1.527790in}}%
\pgfpathlineto{\pgfqpoint{2.192229in}{1.513232in}}%
\pgfpathlineto{\pgfqpoint{2.161117in}{1.533987in}}%
\pgfpathlineto{\pgfqpoint{2.130102in}{1.554679in}}%
\pgfpathclose%
\pgfusepath{stroke,fill}%
\end{pgfscope}%
\begin{pgfscope}%
\pgfpathrectangle{\pgfqpoint{0.050000in}{0.050000in}}{\pgfqpoint{4.900000in}{4.900000in}}%
\pgfusepath{clip}%
\pgfsetbuttcap%
\pgfsetroundjoin%
\definecolor{currentfill}{rgb}{0.600215,0.583729,0.760953}%
\pgfsetfillcolor{currentfill}%
\pgfsetfillopacity{0.900000}%
\pgfsetlinewidth{0.200750pt}%
\definecolor{currentstroke}{rgb}{0.200000,0.200000,0.200000}%
\pgfsetstrokecolor{currentstroke}%
\pgfsetdash{}{0pt}%
\pgfpathmoveto{\pgfqpoint{3.441815in}{2.206268in}}%
\pgfpathlineto{\pgfqpoint{3.485539in}{2.187710in}}%
\pgfpathlineto{\pgfqpoint{3.528724in}{2.146112in}}%
\pgfpathlineto{\pgfqpoint{3.561371in}{2.100704in}}%
\pgfpathlineto{\pgfqpoint{3.594105in}{2.056041in}}%
\pgfpathlineto{\pgfqpoint{3.550992in}{2.099328in}}%
\pgfpathlineto{\pgfqpoint{3.507354in}{2.122272in}}%
\pgfpathlineto{\pgfqpoint{3.474543in}{2.164115in}}%
\pgfpathlineto{\pgfqpoint{3.441815in}{2.206268in}}%
\pgfpathclose%
\pgfusepath{stroke,fill}%
\end{pgfscope}%
\begin{pgfscope}%
\pgfpathrectangle{\pgfqpoint{0.050000in}{0.050000in}}{\pgfqpoint{4.900000in}{4.900000in}}%
\pgfusepath{clip}%
\pgfsetbuttcap%
\pgfsetroundjoin%
\definecolor{currentfill}{rgb}{1.000000,1.000000,1.000000}%
\pgfsetfillcolor{currentfill}%
\pgfsetfillopacity{0.900000}%
\pgfsetlinewidth{0.200750pt}%
\definecolor{currentstroke}{rgb}{0.200000,0.200000,0.200000}%
\pgfsetstrokecolor{currentstroke}%
\pgfsetdash{}{0pt}%
\pgfpathmoveto{\pgfqpoint{1.732144in}{1.550516in}}%
\pgfpathlineto{\pgfqpoint{1.778200in}{1.564977in}}%
\pgfpathlineto{\pgfqpoint{1.824157in}{1.579407in}}%
\pgfpathlineto{\pgfqpoint{1.854641in}{1.558790in}}%
\pgfpathlineto{\pgfqpoint{1.885221in}{1.538109in}}%
\pgfpathlineto{\pgfqpoint{1.839188in}{1.523589in}}%
\pgfpathlineto{\pgfqpoint{1.793055in}{1.509039in}}%
\pgfpathlineto{\pgfqpoint{1.762552in}{1.529810in}}%
\pgfpathlineto{\pgfqpoint{1.732144in}{1.550516in}}%
\pgfpathclose%
\pgfusepath{stroke,fill}%
\end{pgfscope}%
\begin{pgfscope}%
\pgfpathrectangle{\pgfqpoint{0.050000in}{0.050000in}}{\pgfqpoint{4.900000in}{4.900000in}}%
\pgfusepath{clip}%
\pgfsetbuttcap%
\pgfsetroundjoin%
\definecolor{currentfill}{rgb}{0.779223,0.770119,0.867989}%
\pgfsetfillcolor{currentfill}%
\pgfsetfillopacity{0.900000}%
\pgfsetlinewidth{0.200750pt}%
\definecolor{currentstroke}{rgb}{0.200000,0.200000,0.200000}%
\pgfsetstrokecolor{currentstroke}%
\pgfsetdash{}{0pt}%
\pgfpathmoveto{\pgfqpoint{3.659838in}{1.968820in}}%
\pgfpathlineto{\pgfqpoint{3.702437in}{1.911338in}}%
\pgfpathlineto{\pgfqpoint{3.744690in}{1.847172in}}%
\pgfpathlineto{\pgfqpoint{3.777695in}{1.807705in}}%
\pgfpathlineto{\pgfqpoint{3.810797in}{1.768927in}}%
\pgfpathlineto{\pgfqpoint{3.768513in}{1.828759in}}%
\pgfpathlineto{\pgfqpoint{3.725930in}{1.884211in}}%
\pgfpathlineto{\pgfqpoint{3.692838in}{1.926203in}}%
\pgfpathlineto{\pgfqpoint{3.659838in}{1.968820in}}%
\pgfpathclose%
\pgfusepath{stroke,fill}%
\end{pgfscope}%
\begin{pgfscope}%
\pgfpathrectangle{\pgfqpoint{0.050000in}{0.050000in}}{\pgfqpoint{4.900000in}{4.900000in}}%
\pgfusepath{clip}%
\pgfsetbuttcap%
\pgfsetroundjoin%
\definecolor{currentfill}{rgb}{1.000000,1.000000,1.000000}%
\pgfsetfillcolor{currentfill}%
\pgfsetfillopacity{0.900000}%
\pgfsetlinewidth{0.200750pt}%
\definecolor{currentstroke}{rgb}{0.200000,0.200000,0.200000}%
\pgfsetstrokecolor{currentstroke}%
\pgfsetdash{}{0pt}%
\pgfpathmoveto{\pgfqpoint{4.029111in}{1.546526in}}%
\pgfpathlineto{\pgfqpoint{4.072696in}{1.560995in}}%
\pgfpathlineto{\pgfqpoint{4.116186in}{1.575434in}}%
\pgfpathlineto{\pgfqpoint{4.150224in}{1.554804in}}%
\pgfpathlineto{\pgfqpoint{4.184367in}{1.534111in}}%
\pgfpathlineto{\pgfqpoint{4.140816in}{1.519583in}}%
\pgfpathlineto{\pgfqpoint{4.097170in}{1.505024in}}%
\pgfpathlineto{\pgfqpoint{4.063088in}{1.525807in}}%
\pgfpathlineto{\pgfqpoint{4.029111in}{1.546526in}}%
\pgfpathclose%
\pgfusepath{stroke,fill}%
\end{pgfscope}%
\begin{pgfscope}%
\pgfpathrectangle{\pgfqpoint{0.050000in}{0.050000in}}{\pgfqpoint{4.900000in}{4.900000in}}%
\pgfusepath{clip}%
\pgfsetbuttcap%
\pgfsetroundjoin%
\definecolor{currentfill}{rgb}{1.000000,1.000000,1.000000}%
\pgfsetfillcolor{currentfill}%
\pgfsetfillopacity{0.900000}%
\pgfsetlinewidth{0.200750pt}%
\definecolor{currentstroke}{rgb}{0.200000,0.200000,0.200000}%
\pgfsetstrokecolor{currentstroke}%
\pgfsetdash{}{0pt}%
\pgfpathmoveto{\pgfqpoint{1.333937in}{1.546354in}}%
\pgfpathlineto{\pgfqpoint{1.380439in}{1.560824in}}%
\pgfpathlineto{\pgfqpoint{1.426840in}{1.575263in}}%
\pgfpathlineto{\pgfqpoint{1.456720in}{1.554633in}}%
\pgfpathlineto{\pgfqpoint{1.486692in}{1.533939in}}%
\pgfpathlineto{\pgfqpoint{1.440212in}{1.519411in}}%
\pgfpathlineto{\pgfqpoint{1.393631in}{1.504851in}}%
\pgfpathlineto{\pgfqpoint{1.363738in}{1.525635in}}%
\pgfpathlineto{\pgfqpoint{1.333937in}{1.546354in}}%
\pgfpathclose%
\pgfusepath{stroke,fill}%
\end{pgfscope}%
\begin{pgfscope}%
\pgfpathrectangle{\pgfqpoint{0.050000in}{0.050000in}}{\pgfqpoint{4.900000in}{4.900000in}}%
\pgfusepath{clip}%
\pgfsetbuttcap%
\pgfsetroundjoin%
\definecolor{currentfill}{rgb}{0.528612,0.509173,0.718139}%
\pgfsetfillcolor{currentfill}%
\pgfsetfillopacity{0.900000}%
\pgfsetlinewidth{0.200750pt}%
\definecolor{currentstroke}{rgb}{0.200000,0.200000,0.200000}%
\pgfsetstrokecolor{currentstroke}%
\pgfsetdash{}{0pt}%
\pgfpathmoveto{\pgfqpoint{3.287567in}{2.240261in}}%
\pgfpathlineto{\pgfqpoint{3.332288in}{2.278052in}}%
\pgfpathlineto{\pgfqpoint{3.376601in}{2.291471in}}%
\pgfpathlineto{\pgfqpoint{3.409168in}{2.248724in}}%
\pgfpathlineto{\pgfqpoint{3.441815in}{2.206268in}}%
\pgfpathlineto{\pgfqpoint{3.397585in}{2.200421in}}%
\pgfpathlineto{\pgfqpoint{3.352924in}{2.171400in}}%
\pgfpathlineto{\pgfqpoint{3.320205in}{2.205958in}}%
\pgfpathlineto{\pgfqpoint{3.287567in}{2.240261in}}%
\pgfpathclose%
\pgfusepath{stroke,fill}%
\end{pgfscope}%
\begin{pgfscope}%
\pgfpathrectangle{\pgfqpoint{0.050000in}{0.050000in}}{\pgfqpoint{4.900000in}{4.900000in}}%
\pgfusepath{clip}%
\pgfsetbuttcap%
\pgfsetroundjoin%
\definecolor{currentfill}{rgb}{0.952265,0.950296,0.971457}%
\pgfsetfillcolor{currentfill}%
\pgfsetfillopacity{0.900000}%
\pgfsetlinewidth{0.200750pt}%
\definecolor{currentstroke}{rgb}{0.200000,0.200000,0.200000}%
\pgfsetstrokecolor{currentstroke}%
\pgfsetdash{}{0pt}%
\pgfpathmoveto{\pgfqpoint{3.877293in}{1.693278in}}%
\pgfpathlineto{\pgfqpoint{3.919530in}{1.642522in}}%
\pgfpathlineto{\pgfqpoint{3.961877in}{1.603969in}}%
\pgfpathlineto{\pgfqpoint{3.995395in}{1.573267in}}%
\pgfpathlineto{\pgfqpoint{4.029111in}{1.546526in}}%
\pgfpathlineto{\pgfqpoint{3.986523in}{1.574865in}}%
\pgfpathlineto{\pgfqpoint{3.944186in}{1.619935in}}%
\pgfpathlineto{\pgfqpoint{3.910690in}{1.656335in}}%
\pgfpathlineto{\pgfqpoint{3.877293in}{1.693278in}}%
\pgfpathclose%
\pgfusepath{stroke,fill}%
\end{pgfscope}%
\begin{pgfscope}%
\pgfpathrectangle{\pgfqpoint{0.050000in}{0.050000in}}{\pgfqpoint{4.900000in}{4.900000in}}%
\pgfusepath{clip}%
\pgfsetbuttcap%
\pgfsetroundjoin%
\definecolor{currentfill}{rgb}{1.000000,1.000000,1.000000}%
\pgfsetfillcolor{currentfill}%
\pgfsetfillopacity{0.900000}%
\pgfsetlinewidth{0.200750pt}%
\definecolor{currentstroke}{rgb}{0.200000,0.200000,0.200000}%
\pgfsetstrokecolor{currentstroke}%
\pgfsetdash{}{0pt}%
\pgfpathmoveto{\pgfqpoint{2.283500in}{1.542347in}}%
\pgfpathlineto{\pgfqpoint{2.328988in}{1.556948in}}%
\pgfpathlineto{\pgfqpoint{2.374376in}{1.571688in}}%
\pgfpathlineto{\pgfqpoint{2.405731in}{1.551168in}}%
\pgfpathlineto{\pgfqpoint{2.437183in}{1.530603in}}%
\pgfpathlineto{\pgfqpoint{2.391721in}{1.515634in}}%
\pgfpathlineto{\pgfqpoint{2.346161in}{1.500873in}}%
\pgfpathlineto{\pgfqpoint{2.314782in}{1.521639in}}%
\pgfpathlineto{\pgfqpoint{2.283500in}{1.542347in}}%
\pgfpathclose%
\pgfusepath{stroke,fill}%
\end{pgfscope}%
\begin{pgfscope}%
\pgfpathrectangle{\pgfqpoint{0.050000in}{0.050000in}}{\pgfqpoint{4.900000in}{4.900000in}}%
\pgfusepath{clip}%
\pgfsetbuttcap%
\pgfsetroundjoin%
\definecolor{currentfill}{rgb}{1.000000,1.000000,1.000000}%
\pgfsetfillcolor{currentfill}%
\pgfsetfillopacity{0.900000}%
\pgfsetlinewidth{0.200750pt}%
\definecolor{currentstroke}{rgb}{0.200000,0.200000,0.200000}%
\pgfsetstrokecolor{currentstroke}%
\pgfsetdash{}{0pt}%
\pgfpathmoveto{\pgfqpoint{0.935483in}{1.542190in}}%
\pgfpathlineto{\pgfqpoint{0.982430in}{1.556669in}}%
\pgfpathlineto{\pgfqpoint{1.029276in}{1.571116in}}%
\pgfpathlineto{\pgfqpoint{1.058550in}{1.550473in}}%
\pgfpathlineto{\pgfqpoint{1.087915in}{1.529767in}}%
\pgfpathlineto{\pgfqpoint{1.040988in}{1.515229in}}%
\pgfpathlineto{\pgfqpoint{0.993958in}{1.500660in}}%
\pgfpathlineto{\pgfqpoint{0.964675in}{1.521457in}}%
\pgfpathlineto{\pgfqpoint{0.935483in}{1.542190in}}%
\pgfpathclose%
\pgfusepath{stroke,fill}%
\end{pgfscope}%
\begin{pgfscope}%
\pgfpathrectangle{\pgfqpoint{0.050000in}{0.050000in}}{\pgfqpoint{4.900000in}{4.900000in}}%
\pgfusepath{clip}%
\pgfsetbuttcap%
\pgfsetroundjoin%
\definecolor{currentfill}{rgb}{0.976132,0.975148,0.985729}%
\pgfsetfillcolor{currentfill}%
\pgfsetfillopacity{0.900000}%
\pgfsetlinewidth{0.200750pt}%
\definecolor{currentstroke}{rgb}{0.200000,0.200000,0.200000}%
\pgfsetstrokecolor{currentstroke}%
\pgfsetdash{}{0pt}%
\pgfpathmoveto{\pgfqpoint{2.681565in}{1.561786in}}%
\pgfpathlineto{\pgfqpoint{2.726651in}{1.586963in}}%
\pgfpathlineto{\pgfqpoint{2.771677in}{1.617571in}}%
\pgfpathlineto{\pgfqpoint{2.803673in}{1.599941in}}%
\pgfpathlineto{\pgfqpoint{2.835771in}{1.582261in}}%
\pgfpathlineto{\pgfqpoint{2.790646in}{1.550059in}}%
\pgfpathlineto{\pgfqpoint{2.745469in}{1.523468in}}%
\pgfpathlineto{\pgfqpoint{2.713466in}{1.542631in}}%
\pgfpathlineto{\pgfqpoint{2.681565in}{1.561786in}}%
\pgfpathclose%
\pgfusepath{stroke,fill}%
\end{pgfscope}%
\begin{pgfscope}%
\pgfpathrectangle{\pgfqpoint{0.050000in}{0.050000in}}{\pgfqpoint{4.900000in}{4.900000in}}%
\pgfusepath{clip}%
\pgfsetbuttcap%
\pgfsetroundjoin%
\definecolor{currentfill}{rgb}{1.000000,1.000000,1.000000}%
\pgfsetfillcolor{currentfill}%
\pgfsetfillopacity{0.900000}%
\pgfsetlinewidth{0.200750pt}%
\definecolor{currentstroke}{rgb}{0.200000,0.200000,0.200000}%
\pgfsetstrokecolor{currentstroke}%
\pgfsetdash{}{0pt}%
\pgfpathmoveto{\pgfqpoint{1.885221in}{1.538109in}}%
\pgfpathlineto{\pgfqpoint{1.931154in}{1.552596in}}%
\pgfpathlineto{\pgfqpoint{1.976988in}{1.567053in}}%
\pgfpathlineto{\pgfqpoint{2.007737in}{1.546397in}}%
\pgfpathlineto{\pgfqpoint{2.038582in}{1.525678in}}%
\pgfpathlineto{\pgfqpoint{1.992673in}{1.511132in}}%
\pgfpathlineto{\pgfqpoint{1.946664in}{1.496554in}}%
\pgfpathlineto{\pgfqpoint{1.915894in}{1.517363in}}%
\pgfpathlineto{\pgfqpoint{1.885221in}{1.538109in}}%
\pgfpathclose%
\pgfusepath{stroke,fill}%
\end{pgfscope}%
\begin{pgfscope}%
\pgfpathrectangle{\pgfqpoint{0.050000in}{0.050000in}}{\pgfqpoint{4.900000in}{4.900000in}}%
\pgfusepath{clip}%
\pgfsetbuttcap%
\pgfsetroundjoin%
\definecolor{currentfill}{rgb}{1.000000,1.000000,1.000000}%
\pgfsetfillcolor{currentfill}%
\pgfsetfillopacity{0.900000}%
\pgfsetlinewidth{0.200750pt}%
\definecolor{currentstroke}{rgb}{0.200000,0.200000,0.200000}%
\pgfsetstrokecolor{currentstroke}%
\pgfsetdash{}{0pt}%
\pgfpathmoveto{\pgfqpoint{1.486692in}{1.533939in}}%
\pgfpathlineto{\pgfqpoint{1.533071in}{1.548436in}}%
\pgfpathlineto{\pgfqpoint{1.579350in}{1.562901in}}%
\pgfpathlineto{\pgfqpoint{1.609494in}{1.542233in}}%
\pgfpathlineto{\pgfqpoint{1.639731in}{1.521500in}}%
\pgfpathlineto{\pgfqpoint{1.593374in}{1.506945in}}%
\pgfpathlineto{\pgfqpoint{1.546916in}{1.492358in}}%
\pgfpathlineto{\pgfqpoint{1.516757in}{1.513181in}}%
\pgfpathlineto{\pgfqpoint{1.486692in}{1.533939in}}%
\pgfpathclose%
\pgfusepath{stroke,fill}%
\end{pgfscope}%
\begin{pgfscope}%
\pgfpathrectangle{\pgfqpoint{0.050000in}{0.050000in}}{\pgfqpoint{4.900000in}{4.900000in}}%
\pgfusepath{clip}%
\pgfsetbuttcap%
\pgfsetroundjoin%
\definecolor{currentfill}{rgb}{1.000000,1.000000,1.000000}%
\pgfsetfillcolor{currentfill}%
\pgfsetfillopacity{0.900000}%
\pgfsetlinewidth{0.200750pt}%
\definecolor{currentstroke}{rgb}{0.200000,0.200000,0.200000}%
\pgfsetstrokecolor{currentstroke}%
\pgfsetdash{}{0pt}%
\pgfpathmoveto{\pgfqpoint{1.087915in}{1.529767in}}%
\pgfpathlineto{\pgfqpoint{1.134741in}{1.544272in}}%
\pgfpathlineto{\pgfqpoint{1.181465in}{1.558747in}}%
\pgfpathlineto{\pgfqpoint{1.211002in}{1.538066in}}%
\pgfpathlineto{\pgfqpoint{1.240631in}{1.517320in}}%
\pgfpathlineto{\pgfqpoint{1.193826in}{1.502756in}}%
\pgfpathlineto{\pgfqpoint{1.146919in}{1.488160in}}%
\pgfpathlineto{\pgfqpoint{1.117371in}{1.508996in}}%
\pgfpathlineto{\pgfqpoint{1.087915in}{1.529767in}}%
\pgfpathclose%
\pgfusepath{stroke,fill}%
\end{pgfscope}%
\begin{pgfscope}%
\pgfpathrectangle{\pgfqpoint{0.050000in}{0.050000in}}{\pgfqpoint{4.900000in}{4.900000in}}%
\pgfusepath{clip}%
\pgfsetbuttcap%
\pgfsetroundjoin%
\definecolor{currentfill}{rgb}{0.576348,0.558877,0.746682}%
\pgfsetfillcolor{currentfill}%
\pgfsetfillopacity{0.900000}%
\pgfsetlinewidth{0.200750pt}%
\definecolor{currentstroke}{rgb}{0.200000,0.200000,0.200000}%
\pgfsetstrokecolor{currentstroke}%
\pgfsetdash{}{0pt}%
\pgfpathmoveto{\pgfqpoint{3.197342in}{2.110240in}}%
\pgfpathlineto{\pgfqpoint{3.242549in}{2.182435in}}%
\pgfpathlineto{\pgfqpoint{3.287567in}{2.240261in}}%
\pgfpathlineto{\pgfqpoint{3.320205in}{2.205958in}}%
\pgfpathlineto{\pgfqpoint{3.352924in}{2.171400in}}%
\pgfpathlineto{\pgfqpoint{3.307930in}{2.122570in}}%
\pgfpathlineto{\pgfqpoint{3.262708in}{2.058739in}}%
\pgfpathlineto{\pgfqpoint{3.229981in}{2.084772in}}%
\pgfpathlineto{\pgfqpoint{3.197342in}{2.110240in}}%
\pgfpathclose%
\pgfusepath{stroke,fill}%
\end{pgfscope}%
\begin{pgfscope}%
\pgfpathrectangle{\pgfqpoint{0.050000in}{0.050000in}}{\pgfqpoint{4.900000in}{4.900000in}}%
\pgfusepath{clip}%
\pgfsetbuttcap%
\pgfsetroundjoin%
\definecolor{currentfill}{rgb}{1.000000,1.000000,1.000000}%
\pgfsetfillcolor{currentfill}%
\pgfsetfillopacity{0.900000}%
\pgfsetlinewidth{0.200750pt}%
\definecolor{currentstroke}{rgb}{0.200000,0.200000,0.200000}%
\pgfsetstrokecolor{currentstroke}%
\pgfsetdash{}{0pt}%
\pgfpathmoveto{\pgfqpoint{2.437183in}{1.530603in}}%
\pgfpathlineto{\pgfqpoint{2.482545in}{1.546015in}}%
\pgfpathlineto{\pgfqpoint{2.527811in}{1.562266in}}%
\pgfpathlineto{\pgfqpoint{2.559434in}{1.542136in}}%
\pgfpathlineto{\pgfqpoint{2.591156in}{1.521989in}}%
\pgfpathlineto{\pgfqpoint{2.545816in}{1.505152in}}%
\pgfpathlineto{\pgfqpoint{2.500380in}{1.489345in}}%
\pgfpathlineto{\pgfqpoint{2.468732in}{1.509995in}}%
\pgfpathlineto{\pgfqpoint{2.437183in}{1.530603in}}%
\pgfpathclose%
\pgfusepath{stroke,fill}%
\end{pgfscope}%
\begin{pgfscope}%
\pgfpathrectangle{\pgfqpoint{0.050000in}{0.050000in}}{\pgfqpoint{4.900000in}{4.900000in}}%
\pgfusepath{clip}%
\pgfsetbuttcap%
\pgfsetroundjoin%
\definecolor{currentfill}{rgb}{0.779223,0.770119,0.867989}%
\pgfsetfillcolor{currentfill}%
\pgfsetfillopacity{0.900000}%
\pgfsetlinewidth{0.200750pt}%
\definecolor{currentstroke}{rgb}{0.200000,0.200000,0.200000}%
\pgfsetstrokecolor{currentstroke}%
\pgfsetdash{}{0pt}%
\pgfpathmoveto{\pgfqpoint{3.016236in}{1.792067in}}%
\pgfpathlineto{\pgfqpoint{3.061456in}{1.866464in}}%
\pgfpathlineto{\pgfqpoint{3.106734in}{1.946915in}}%
\pgfpathlineto{\pgfqpoint{3.139299in}{1.927710in}}%
\pgfpathlineto{\pgfqpoint{3.171962in}{1.907927in}}%
\pgfpathlineto{\pgfqpoint{3.126575in}{1.831138in}}%
\pgfpathlineto{\pgfqpoint{3.081232in}{1.758751in}}%
\pgfpathlineto{\pgfqpoint{3.048683in}{1.775612in}}%
\pgfpathlineto{\pgfqpoint{3.016236in}{1.792067in}}%
\pgfpathclose%
\pgfusepath{stroke,fill}%
\end{pgfscope}%
\begin{pgfscope}%
\pgfpathrectangle{\pgfqpoint{0.050000in}{0.050000in}}{\pgfqpoint{4.900000in}{4.900000in}}%
\pgfusepath{clip}%
\pgfsetbuttcap%
\pgfsetroundjoin%
\definecolor{currentfill}{rgb}{0.671819,0.658285,0.803768}%
\pgfsetfillcolor{currentfill}%
\pgfsetfillopacity{0.900000}%
\pgfsetlinewidth{0.200750pt}%
\definecolor{currentstroke}{rgb}{0.200000,0.200000,0.200000}%
\pgfsetstrokecolor{currentstroke}%
\pgfsetdash{}{0pt}%
\pgfpathmoveto{\pgfqpoint{3.106734in}{1.946915in}}%
\pgfpathlineto{\pgfqpoint{3.152045in}{2.029802in}}%
\pgfpathlineto{\pgfqpoint{3.197342in}{2.110240in}}%
\pgfpathlineto{\pgfqpoint{3.229981in}{2.084772in}}%
\pgfpathlineto{\pgfqpoint{3.262708in}{2.058739in}}%
\pgfpathlineto{\pgfqpoint{3.217358in}{1.985385in}}%
\pgfpathlineto{\pgfqpoint{3.171962in}{1.907927in}}%
\pgfpathlineto{\pgfqpoint{3.139299in}{1.927710in}}%
\pgfpathlineto{\pgfqpoint{3.106734in}{1.946915in}}%
\pgfpathclose%
\pgfusepath{stroke,fill}%
\end{pgfscope}%
\begin{pgfscope}%
\pgfpathrectangle{\pgfqpoint{0.050000in}{0.050000in}}{\pgfqpoint{4.900000in}{4.900000in}}%
\pgfusepath{clip}%
\pgfsetbuttcap%
\pgfsetroundjoin%
\definecolor{currentfill}{rgb}{1.000000,1.000000,1.000000}%
\pgfsetfillcolor{currentfill}%
\pgfsetfillopacity{0.900000}%
\pgfsetlinewidth{0.200750pt}%
\definecolor{currentstroke}{rgb}{0.200000,0.200000,0.200000}%
\pgfsetstrokecolor{currentstroke}%
\pgfsetdash{}{0pt}%
\pgfpathmoveto{\pgfqpoint{4.184367in}{1.534111in}}%
\pgfpathlineto{\pgfqpoint{4.227823in}{1.548607in}}%
\pgfpathlineto{\pgfqpoint{4.262103in}{1.527895in}}%
\pgfpathlineto{\pgfqpoint{4.296489in}{1.507118in}}%
\pgfpathlineto{\pgfqpoint{4.252972in}{1.492531in}}%
\pgfpathlineto{\pgfqpoint{4.218616in}{1.513353in}}%
\pgfpathlineto{\pgfqpoint{4.184367in}{1.534111in}}%
\pgfpathclose%
\pgfusepath{stroke,fill}%
\end{pgfscope}%
\begin{pgfscope}%
\pgfpathrectangle{\pgfqpoint{0.050000in}{0.050000in}}{\pgfqpoint{4.900000in}{4.900000in}}%
\pgfusepath{clip}%
\pgfsetbuttcap%
\pgfsetroundjoin%
\definecolor{currentfill}{rgb}{1.000000,1.000000,1.000000}%
\pgfsetfillcolor{currentfill}%
\pgfsetfillopacity{0.900000}%
\pgfsetlinewidth{0.200750pt}%
\definecolor{currentstroke}{rgb}{0.200000,0.200000,0.200000}%
\pgfsetstrokecolor{currentstroke}%
\pgfsetdash{}{0pt}%
\pgfpathmoveto{\pgfqpoint{2.038582in}{1.525678in}}%
\pgfpathlineto{\pgfqpoint{2.084392in}{1.540193in}}%
\pgfpathlineto{\pgfqpoint{2.130102in}{1.554679in}}%
\pgfpathlineto{\pgfqpoint{2.161117in}{1.533987in}}%
\pgfpathlineto{\pgfqpoint{2.192229in}{1.513232in}}%
\pgfpathlineto{\pgfqpoint{2.146444in}{1.498653in}}%
\pgfpathlineto{\pgfqpoint{2.100559in}{1.484046in}}%
\pgfpathlineto{\pgfqpoint{2.069522in}{1.504894in}}%
\pgfpathlineto{\pgfqpoint{2.038582in}{1.525678in}}%
\pgfpathclose%
\pgfusepath{stroke,fill}%
\end{pgfscope}%
\begin{pgfscope}%
\pgfpathrectangle{\pgfqpoint{0.050000in}{0.050000in}}{\pgfqpoint{4.900000in}{4.900000in}}%
\pgfusepath{clip}%
\pgfsetbuttcap%
\pgfsetroundjoin%
\definecolor{currentfill}{rgb}{0.868727,0.863314,0.921507}%
\pgfsetfillcolor{currentfill}%
\pgfsetfillopacity{0.900000}%
\pgfsetlinewidth{0.200750pt}%
\definecolor{currentstroke}{rgb}{0.200000,0.200000,0.200000}%
\pgfsetstrokecolor{currentstroke}%
\pgfsetdash{}{0pt}%
\pgfpathmoveto{\pgfqpoint{2.925960in}{1.669158in}}%
\pgfpathlineto{\pgfqpoint{2.971075in}{1.725975in}}%
\pgfpathlineto{\pgfqpoint{3.016236in}{1.792067in}}%
\pgfpathlineto{\pgfqpoint{3.048683in}{1.775612in}}%
\pgfpathlineto{\pgfqpoint{3.081232in}{1.758751in}}%
\pgfpathlineto{\pgfqpoint{3.035942in}{1.693287in}}%
\pgfpathlineto{\pgfqpoint{2.990699in}{1.636069in}}%
\pgfpathlineto{\pgfqpoint{2.958278in}{1.652715in}}%
\pgfpathlineto{\pgfqpoint{2.925960in}{1.669158in}}%
\pgfpathclose%
\pgfusepath{stroke,fill}%
\end{pgfscope}%
\begin{pgfscope}%
\pgfpathrectangle{\pgfqpoint{0.050000in}{0.050000in}}{\pgfqpoint{4.900000in}{4.900000in}}%
\pgfusepath{clip}%
\pgfsetbuttcap%
\pgfsetroundjoin%
\definecolor{currentfill}{rgb}{0.636017,0.621007,0.782361}%
\pgfsetfillcolor{currentfill}%
\pgfsetfillopacity{0.900000}%
\pgfsetlinewidth{0.200750pt}%
\definecolor{currentstroke}{rgb}{0.200000,0.200000,0.200000}%
\pgfsetstrokecolor{currentstroke}%
\pgfsetdash{}{0pt}%
\pgfpathmoveto{\pgfqpoint{3.507354in}{2.122272in}}%
\pgfpathlineto{\pgfqpoint{3.550992in}{2.099328in}}%
\pgfpathlineto{\pgfqpoint{3.594105in}{2.056041in}}%
\pgfpathlineto{\pgfqpoint{3.626927in}{2.012090in}}%
\pgfpathlineto{\pgfqpoint{3.659838in}{1.968820in}}%
\pgfpathlineto{\pgfqpoint{3.616785in}{2.013096in}}%
\pgfpathlineto{\pgfqpoint{3.573227in}{2.039536in}}%
\pgfpathlineto{\pgfqpoint{3.540248in}{2.080745in}}%
\pgfpathlineto{\pgfqpoint{3.507354in}{2.122272in}}%
\pgfpathclose%
\pgfusepath{stroke,fill}%
\end{pgfscope}%
\begin{pgfscope}%
\pgfpathrectangle{\pgfqpoint{0.050000in}{0.050000in}}{\pgfqpoint{4.900000in}{4.900000in}}%
\pgfusepath{clip}%
\pgfsetbuttcap%
\pgfsetroundjoin%
\definecolor{currentfill}{rgb}{1.000000,1.000000,1.000000}%
\pgfsetfillcolor{currentfill}%
\pgfsetfillopacity{0.900000}%
\pgfsetlinewidth{0.200750pt}%
\definecolor{currentstroke}{rgb}{0.200000,0.200000,0.200000}%
\pgfsetstrokecolor{currentstroke}%
\pgfsetdash{}{0pt}%
\pgfpathmoveto{\pgfqpoint{1.639731in}{1.521500in}}%
\pgfpathlineto{\pgfqpoint{1.685987in}{1.536024in}}%
\pgfpathlineto{\pgfqpoint{1.732144in}{1.550516in}}%
\pgfpathlineto{\pgfqpoint{1.762552in}{1.529810in}}%
\pgfpathlineto{\pgfqpoint{1.793055in}{1.509039in}}%
\pgfpathlineto{\pgfqpoint{1.746821in}{1.494456in}}%
\pgfpathlineto{\pgfqpoint{1.700487in}{1.479842in}}%
\pgfpathlineto{\pgfqpoint{1.670062in}{1.500704in}}%
\pgfpathlineto{\pgfqpoint{1.639731in}{1.521500in}}%
\pgfpathclose%
\pgfusepath{stroke,fill}%
\end{pgfscope}%
\begin{pgfscope}%
\pgfpathrectangle{\pgfqpoint{0.050000in}{0.050000in}}{\pgfqpoint{4.900000in}{4.900000in}}%
\pgfusepath{clip}%
\pgfsetbuttcap%
\pgfsetroundjoin%
\definecolor{currentfill}{rgb}{0.815025,0.807397,0.889396}%
\pgfsetfillcolor{currentfill}%
\pgfsetfillopacity{0.900000}%
\pgfsetlinewidth{0.200750pt}%
\definecolor{currentstroke}{rgb}{0.200000,0.200000,0.200000}%
\pgfsetstrokecolor{currentstroke}%
\pgfsetdash{}{0pt}%
\pgfpathmoveto{\pgfqpoint{3.725930in}{1.884211in}}%
\pgfpathlineto{\pgfqpoint{3.768513in}{1.828759in}}%
\pgfpathlineto{\pgfqpoint{3.810797in}{1.768927in}}%
\pgfpathlineto{\pgfqpoint{3.843996in}{1.730797in}}%
\pgfpathlineto{\pgfqpoint{3.877293in}{1.693278in}}%
\pgfpathlineto{\pgfqpoint{3.834969in}{1.748806in}}%
\pgfpathlineto{\pgfqpoint{3.792389in}{1.801992in}}%
\pgfpathlineto{\pgfqpoint{3.759113in}{1.842816in}}%
\pgfpathlineto{\pgfqpoint{3.725930in}{1.884211in}}%
\pgfpathclose%
\pgfusepath{stroke,fill}%
\end{pgfscope}%
\begin{pgfscope}%
\pgfpathrectangle{\pgfqpoint{0.050000in}{0.050000in}}{\pgfqpoint{4.900000in}{4.900000in}}%
\pgfusepath{clip}%
\pgfsetbuttcap%
\pgfsetroundjoin%
\definecolor{currentfill}{rgb}{1.000000,1.000000,1.000000}%
\pgfsetfillcolor{currentfill}%
\pgfsetfillopacity{0.900000}%
\pgfsetlinewidth{0.200750pt}%
\definecolor{currentstroke}{rgb}{0.200000,0.200000,0.200000}%
\pgfsetstrokecolor{currentstroke}%
\pgfsetdash{}{0pt}%
\pgfpathmoveto{\pgfqpoint{1.240631in}{1.517320in}}%
\pgfpathlineto{\pgfqpoint{1.287335in}{1.531853in}}%
\pgfpathlineto{\pgfqpoint{1.333937in}{1.546354in}}%
\pgfpathlineto{\pgfqpoint{1.363738in}{1.525635in}}%
\pgfpathlineto{\pgfqpoint{1.393631in}{1.504851in}}%
\pgfpathlineto{\pgfqpoint{1.346949in}{1.490259in}}%
\pgfpathlineto{\pgfqpoint{1.300165in}{1.475636in}}%
\pgfpathlineto{\pgfqpoint{1.270352in}{1.496511in}}%
\pgfpathlineto{\pgfqpoint{1.240631in}{1.517320in}}%
\pgfpathclose%
\pgfusepath{stroke,fill}%
\end{pgfscope}%
\begin{pgfscope}%
\pgfpathrectangle{\pgfqpoint{0.050000in}{0.050000in}}{\pgfqpoint{4.900000in}{4.900000in}}%
\pgfusepath{clip}%
\pgfsetbuttcap%
\pgfsetroundjoin%
\definecolor{currentfill}{rgb}{0.934364,0.931657,0.960754}%
\pgfsetfillcolor{currentfill}%
\pgfsetfillopacity{0.900000}%
\pgfsetlinewidth{0.200750pt}%
\definecolor{currentstroke}{rgb}{0.200000,0.200000,0.200000}%
\pgfsetstrokecolor{currentstroke}%
\pgfsetdash{}{0pt}%
\pgfpathmoveto{\pgfqpoint{2.835771in}{1.582261in}}%
\pgfpathlineto{\pgfqpoint{2.880867in}{1.621533in}}%
\pgfpathlineto{\pgfqpoint{2.925960in}{1.669158in}}%
\pgfpathlineto{\pgfqpoint{2.958278in}{1.652715in}}%
\pgfpathlineto{\pgfqpoint{2.990699in}{1.636069in}}%
\pgfpathlineto{\pgfqpoint{2.945486in}{1.587381in}}%
\pgfpathlineto{\pgfqpoint{2.900278in}{1.546714in}}%
\pgfpathlineto{\pgfqpoint{2.867973in}{1.564523in}}%
\pgfpathlineto{\pgfqpoint{2.835771in}{1.582261in}}%
\pgfpathclose%
\pgfusepath{stroke,fill}%
\end{pgfscope}%
\begin{pgfscope}%
\pgfpathrectangle{\pgfqpoint{0.050000in}{0.050000in}}{\pgfqpoint{4.900000in}{4.900000in}}%
\pgfusepath{clip}%
\pgfsetbuttcap%
\pgfsetroundjoin%
\definecolor{currentfill}{rgb}{0.994033,0.993787,0.996432}%
\pgfsetfillcolor{currentfill}%
\pgfsetfillopacity{0.900000}%
\pgfsetlinewidth{0.200750pt}%
\definecolor{currentstroke}{rgb}{0.200000,0.200000,0.200000}%
\pgfsetstrokecolor{currentstroke}%
\pgfsetdash{}{0pt}%
\pgfpathmoveto{\pgfqpoint{2.591156in}{1.521989in}}%
\pgfpathlineto{\pgfqpoint{2.636403in}{1.540545in}}%
\pgfpathlineto{\pgfqpoint{2.681565in}{1.561786in}}%
\pgfpathlineto{\pgfqpoint{2.713466in}{1.542631in}}%
\pgfpathlineto{\pgfqpoint{2.745469in}{1.523468in}}%
\pgfpathlineto{\pgfqpoint{2.700224in}{1.501075in}}%
\pgfpathlineto{\pgfqpoint{2.654898in}{1.481639in}}%
\pgfpathlineto{\pgfqpoint{2.622977in}{1.501823in}}%
\pgfpathlineto{\pgfqpoint{2.591156in}{1.521989in}}%
\pgfpathclose%
\pgfusepath{stroke,fill}%
\end{pgfscope}%
\begin{pgfscope}%
\pgfpathrectangle{\pgfqpoint{0.050000in}{0.050000in}}{\pgfqpoint{4.900000in}{4.900000in}}%
\pgfusepath{clip}%
\pgfsetbuttcap%
\pgfsetroundjoin%
\definecolor{currentfill}{rgb}{1.000000,1.000000,1.000000}%
\pgfsetfillcolor{currentfill}%
\pgfsetfillopacity{0.900000}%
\pgfsetlinewidth{0.200750pt}%
\definecolor{currentstroke}{rgb}{0.200000,0.200000,0.200000}%
\pgfsetstrokecolor{currentstroke}%
\pgfsetdash{}{0pt}%
\pgfpathmoveto{\pgfqpoint{2.192229in}{1.513232in}}%
\pgfpathlineto{\pgfqpoint{2.237914in}{1.527790in}}%
\pgfpathlineto{\pgfqpoint{2.283500in}{1.542347in}}%
\pgfpathlineto{\pgfqpoint{2.314782in}{1.521639in}}%
\pgfpathlineto{\pgfqpoint{2.346161in}{1.500873in}}%
\pgfpathlineto{\pgfqpoint{2.300501in}{1.486193in}}%
\pgfpathlineto{\pgfqpoint{2.254742in}{1.471530in}}%
\pgfpathlineto{\pgfqpoint{2.223437in}{1.492413in}}%
\pgfpathlineto{\pgfqpoint{2.192229in}{1.513232in}}%
\pgfpathclose%
\pgfusepath{stroke,fill}%
\end{pgfscope}%
\begin{pgfscope}%
\pgfpathrectangle{\pgfqpoint{0.050000in}{0.050000in}}{\pgfqpoint{4.900000in}{4.900000in}}%
\pgfusepath{clip}%
\pgfsetbuttcap%
\pgfsetroundjoin%
\definecolor{currentfill}{rgb}{0.558447,0.540238,0.735978}%
\pgfsetfillcolor{currentfill}%
\pgfsetfillopacity{0.900000}%
\pgfsetlinewidth{0.200750pt}%
\definecolor{currentstroke}{rgb}{0.200000,0.200000,0.200000}%
\pgfsetstrokecolor{currentstroke}%
\pgfsetdash{}{0pt}%
\pgfpathmoveto{\pgfqpoint{3.352924in}{2.171400in}}%
\pgfpathlineto{\pgfqpoint{3.397585in}{2.200421in}}%
\pgfpathlineto{\pgfqpoint{3.441815in}{2.206268in}}%
\pgfpathlineto{\pgfqpoint{3.474543in}{2.164115in}}%
\pgfpathlineto{\pgfqpoint{3.507354in}{2.122272in}}%
\pgfpathlineto{\pgfqpoint{3.463207in}{2.122908in}}%
\pgfpathlineto{\pgfqpoint{3.418611in}{2.101675in}}%
\pgfpathlineto{\pgfqpoint{3.385726in}{2.136628in}}%
\pgfpathlineto{\pgfqpoint{3.352924in}{2.171400in}}%
\pgfpathclose%
\pgfusepath{stroke,fill}%
\end{pgfscope}%
\begin{pgfscope}%
\pgfpathrectangle{\pgfqpoint{0.050000in}{0.050000in}}{\pgfqpoint{4.900000in}{4.900000in}}%
\pgfusepath{clip}%
\pgfsetbuttcap%
\pgfsetroundjoin%
\definecolor{currentfill}{rgb}{1.000000,1.000000,1.000000}%
\pgfsetfillcolor{currentfill}%
\pgfsetfillopacity{0.900000}%
\pgfsetlinewidth{0.200750pt}%
\definecolor{currentstroke}{rgb}{0.200000,0.200000,0.200000}%
\pgfsetstrokecolor{currentstroke}%
\pgfsetdash{}{0pt}%
\pgfpathmoveto{\pgfqpoint{1.793055in}{1.509039in}}%
\pgfpathlineto{\pgfqpoint{1.839188in}{1.523589in}}%
\pgfpathlineto{\pgfqpoint{1.885221in}{1.538109in}}%
\pgfpathlineto{\pgfqpoint{1.915894in}{1.517363in}}%
\pgfpathlineto{\pgfqpoint{1.946664in}{1.496554in}}%
\pgfpathlineto{\pgfqpoint{1.900554in}{1.481944in}}%
\pgfpathlineto{\pgfqpoint{1.854345in}{1.467303in}}%
\pgfpathlineto{\pgfqpoint{1.823652in}{1.488203in}}%
\pgfpathlineto{\pgfqpoint{1.793055in}{1.509039in}}%
\pgfpathclose%
\pgfusepath{stroke,fill}%
\end{pgfscope}%
\begin{pgfscope}%
\pgfpathrectangle{\pgfqpoint{0.050000in}{0.050000in}}{\pgfqpoint{4.900000in}{4.900000in}}%
\pgfusepath{clip}%
\pgfsetbuttcap%
\pgfsetroundjoin%
\definecolor{currentfill}{rgb}{0.970165,0.968935,0.982161}%
\pgfsetfillcolor{currentfill}%
\pgfsetfillopacity{0.900000}%
\pgfsetlinewidth{0.200750pt}%
\definecolor{currentstroke}{rgb}{0.200000,0.200000,0.200000}%
\pgfsetstrokecolor{currentstroke}%
\pgfsetdash{}{0pt}%
\pgfpathmoveto{\pgfqpoint{3.944186in}{1.619935in}}%
\pgfpathlineto{\pgfqpoint{3.986523in}{1.574865in}}%
\pgfpathlineto{\pgfqpoint{4.029111in}{1.546526in}}%
\pgfpathlineto{\pgfqpoint{4.063088in}{1.525807in}}%
\pgfpathlineto{\pgfqpoint{4.097170in}{1.505024in}}%
\pgfpathlineto{\pgfqpoint{4.053925in}{1.508998in}}%
\pgfpathlineto{\pgfqpoint{4.011481in}{1.548646in}}%
\pgfpathlineto{\pgfqpoint{3.977783in}{1.584048in}}%
\pgfpathlineto{\pgfqpoint{3.944186in}{1.619935in}}%
\pgfpathclose%
\pgfusepath{stroke,fill}%
\end{pgfscope}%
\begin{pgfscope}%
\pgfpathrectangle{\pgfqpoint{0.050000in}{0.050000in}}{\pgfqpoint{4.900000in}{4.900000in}}%
\pgfusepath{clip}%
\pgfsetbuttcap%
\pgfsetroundjoin%
\definecolor{currentfill}{rgb}{1.000000,1.000000,1.000000}%
\pgfsetfillcolor{currentfill}%
\pgfsetfillopacity{0.900000}%
\pgfsetlinewidth{0.200750pt}%
\definecolor{currentstroke}{rgb}{0.200000,0.200000,0.200000}%
\pgfsetstrokecolor{currentstroke}%
\pgfsetdash{}{0pt}%
\pgfpathmoveto{\pgfqpoint{4.097170in}{1.505024in}}%
\pgfpathlineto{\pgfqpoint{4.140816in}{1.519583in}}%
\pgfpathlineto{\pgfqpoint{4.184367in}{1.534111in}}%
\pgfpathlineto{\pgfqpoint{4.218616in}{1.513353in}}%
\pgfpathlineto{\pgfqpoint{4.252972in}{1.492531in}}%
\pgfpathlineto{\pgfqpoint{4.209360in}{1.477913in}}%
\pgfpathlineto{\pgfqpoint{4.165653in}{1.463263in}}%
\pgfpathlineto{\pgfqpoint{4.131359in}{1.484176in}}%
\pgfpathlineto{\pgfqpoint{4.097170in}{1.505024in}}%
\pgfpathclose%
\pgfusepath{stroke,fill}%
\end{pgfscope}%
\begin{pgfscope}%
\pgfpathrectangle{\pgfqpoint{0.050000in}{0.050000in}}{\pgfqpoint{4.900000in}{4.900000in}}%
\pgfusepath{clip}%
\pgfsetbuttcap%
\pgfsetroundjoin%
\definecolor{currentfill}{rgb}{1.000000,1.000000,1.000000}%
\pgfsetfillcolor{currentfill}%
\pgfsetfillopacity{0.900000}%
\pgfsetlinewidth{0.200750pt}%
\definecolor{currentstroke}{rgb}{0.200000,0.200000,0.200000}%
\pgfsetstrokecolor{currentstroke}%
\pgfsetdash{}{0pt}%
\pgfpathmoveto{\pgfqpoint{1.393631in}{1.504851in}}%
\pgfpathlineto{\pgfqpoint{1.440212in}{1.519411in}}%
\pgfpathlineto{\pgfqpoint{1.486692in}{1.533939in}}%
\pgfpathlineto{\pgfqpoint{1.516757in}{1.513181in}}%
\pgfpathlineto{\pgfqpoint{1.546916in}{1.492358in}}%
\pgfpathlineto{\pgfqpoint{1.500357in}{1.477740in}}%
\pgfpathlineto{\pgfqpoint{1.453697in}{1.463089in}}%
\pgfpathlineto{\pgfqpoint{1.423617in}{1.484002in}}%
\pgfpathlineto{\pgfqpoint{1.393631in}{1.504851in}}%
\pgfpathclose%
\pgfusepath{stroke,fill}%
\end{pgfscope}%
\begin{pgfscope}%
\pgfpathrectangle{\pgfqpoint{0.050000in}{0.050000in}}{\pgfqpoint{4.900000in}{4.900000in}}%
\pgfusepath{clip}%
\pgfsetbuttcap%
\pgfsetroundjoin%
\definecolor{currentfill}{rgb}{1.000000,1.000000,1.000000}%
\pgfsetfillcolor{currentfill}%
\pgfsetfillopacity{0.900000}%
\pgfsetlinewidth{0.200750pt}%
\definecolor{currentstroke}{rgb}{0.200000,0.200000,0.200000}%
\pgfsetstrokecolor{currentstroke}%
\pgfsetdash{}{0pt}%
\pgfpathmoveto{\pgfqpoint{0.993958in}{1.500660in}}%
\pgfpathlineto{\pgfqpoint{1.040988in}{1.515229in}}%
\pgfpathlineto{\pgfqpoint{1.087915in}{1.529767in}}%
\pgfpathlineto{\pgfqpoint{1.117371in}{1.508996in}}%
\pgfpathlineto{\pgfqpoint{1.146919in}{1.488160in}}%
\pgfpathlineto{\pgfqpoint{1.099910in}{1.473532in}}%
\pgfpathlineto{\pgfqpoint{1.052798in}{1.458873in}}%
\pgfpathlineto{\pgfqpoint{1.023332in}{1.479799in}}%
\pgfpathlineto{\pgfqpoint{0.993958in}{1.500660in}}%
\pgfpathclose%
\pgfusepath{stroke,fill}%
\end{pgfscope}%
\begin{pgfscope}%
\pgfpathrectangle{\pgfqpoint{0.050000in}{0.050000in}}{\pgfqpoint{4.900000in}{4.900000in}}%
\pgfusepath{clip}%
\pgfsetbuttcap%
\pgfsetroundjoin%
\definecolor{currentfill}{rgb}{1.000000,1.000000,1.000000}%
\pgfsetfillcolor{currentfill}%
\pgfsetfillopacity{0.900000}%
\pgfsetlinewidth{0.200750pt}%
\definecolor{currentstroke}{rgb}{0.200000,0.200000,0.200000}%
\pgfsetstrokecolor{currentstroke}%
\pgfsetdash{}{0pt}%
\pgfpathmoveto{\pgfqpoint{2.346161in}{1.500873in}}%
\pgfpathlineto{\pgfqpoint{2.391721in}{1.515634in}}%
\pgfpathlineto{\pgfqpoint{2.437183in}{1.530603in}}%
\pgfpathlineto{\pgfqpoint{2.468732in}{1.509995in}}%
\pgfpathlineto{\pgfqpoint{2.500380in}{1.489345in}}%
\pgfpathlineto{\pgfqpoint{2.454845in}{1.474110in}}%
\pgfpathlineto{\pgfqpoint{2.409212in}{1.459165in}}%
\pgfpathlineto{\pgfqpoint{2.377637in}{1.480048in}}%
\pgfpathlineto{\pgfqpoint{2.346161in}{1.500873in}}%
\pgfpathclose%
\pgfusepath{stroke,fill}%
\end{pgfscope}%
\begin{pgfscope}%
\pgfpathrectangle{\pgfqpoint{0.050000in}{0.050000in}}{\pgfqpoint{4.900000in}{4.900000in}}%
\pgfusepath{clip}%
\pgfsetbuttcap%
\pgfsetroundjoin%
\definecolor{currentfill}{rgb}{1.000000,1.000000,1.000000}%
\pgfsetfillcolor{currentfill}%
\pgfsetfillopacity{0.900000}%
\pgfsetlinewidth{0.200750pt}%
\definecolor{currentstroke}{rgb}{0.200000,0.200000,0.200000}%
\pgfsetstrokecolor{currentstroke}%
\pgfsetdash{}{0pt}%
\pgfpathmoveto{\pgfqpoint{1.946664in}{1.496554in}}%
\pgfpathlineto{\pgfqpoint{1.992673in}{1.511132in}}%
\pgfpathlineto{\pgfqpoint{2.038582in}{1.525678in}}%
\pgfpathlineto{\pgfqpoint{2.069522in}{1.504894in}}%
\pgfpathlineto{\pgfqpoint{2.100559in}{1.484046in}}%
\pgfpathlineto{\pgfqpoint{2.054575in}{1.469409in}}%
\pgfpathlineto{\pgfqpoint{2.008490in}{1.454740in}}%
\pgfpathlineto{\pgfqpoint{1.977529in}{1.475680in}}%
\pgfpathlineto{\pgfqpoint{1.946664in}{1.496554in}}%
\pgfpathclose%
\pgfusepath{stroke,fill}%
\end{pgfscope}%
\begin{pgfscope}%
\pgfpathrectangle{\pgfqpoint{0.050000in}{0.050000in}}{\pgfqpoint{4.900000in}{4.900000in}}%
\pgfusepath{clip}%
\pgfsetbuttcap%
\pgfsetroundjoin%
\definecolor{currentfill}{rgb}{0.976132,0.975148,0.985729}%
\pgfsetfillcolor{currentfill}%
\pgfsetfillopacity{0.900000}%
\pgfsetlinewidth{0.200750pt}%
\definecolor{currentstroke}{rgb}{0.200000,0.200000,0.200000}%
\pgfsetstrokecolor{currentstroke}%
\pgfsetdash{}{0pt}%
\pgfpathmoveto{\pgfqpoint{2.745469in}{1.523468in}}%
\pgfpathlineto{\pgfqpoint{2.790646in}{1.550059in}}%
\pgfpathlineto{\pgfqpoint{2.835771in}{1.582261in}}%
\pgfpathlineto{\pgfqpoint{2.867973in}{1.564523in}}%
\pgfpathlineto{\pgfqpoint{2.900278in}{1.546714in}}%
\pgfpathlineto{\pgfqpoint{2.855050in}{1.513045in}}%
\pgfpathlineto{\pgfqpoint{2.809780in}{1.485097in}}%
\pgfpathlineto{\pgfqpoint{2.777573in}{1.504292in}}%
\pgfpathlineto{\pgfqpoint{2.745469in}{1.523468in}}%
\pgfpathclose%
\pgfusepath{stroke,fill}%
\end{pgfscope}%
\begin{pgfscope}%
\pgfpathrectangle{\pgfqpoint{0.050000in}{0.050000in}}{\pgfqpoint{4.900000in}{4.900000in}}%
\pgfusepath{clip}%
\pgfsetbuttcap%
\pgfsetroundjoin%
\definecolor{currentfill}{rgb}{0.594248,0.577516,0.757386}%
\pgfsetfillcolor{currentfill}%
\pgfsetfillopacity{0.900000}%
\pgfsetlinewidth{0.200750pt}%
\definecolor{currentstroke}{rgb}{0.200000,0.200000,0.200000}%
\pgfsetstrokecolor{currentstroke}%
\pgfsetdash{}{0pt}%
\pgfpathmoveto{\pgfqpoint{3.262708in}{2.058739in}}%
\pgfpathlineto{\pgfqpoint{3.307930in}{2.122570in}}%
\pgfpathlineto{\pgfqpoint{3.352924in}{2.171400in}}%
\pgfpathlineto{\pgfqpoint{3.385726in}{2.136628in}}%
\pgfpathlineto{\pgfqpoint{3.418611in}{2.101675in}}%
\pgfpathlineto{\pgfqpoint{3.373652in}{2.061079in}}%
\pgfpathlineto{\pgfqpoint{3.328429in}{2.005129in}}%
\pgfpathlineto{\pgfqpoint{3.295524in}{2.032179in}}%
\pgfpathlineto{\pgfqpoint{3.262708in}{2.058739in}}%
\pgfpathclose%
\pgfusepath{stroke,fill}%
\end{pgfscope}%
\begin{pgfscope}%
\pgfpathrectangle{\pgfqpoint{0.050000in}{0.050000in}}{\pgfqpoint{4.900000in}{4.900000in}}%
\pgfusepath{clip}%
\pgfsetbuttcap%
\pgfsetroundjoin%
\definecolor{currentfill}{rgb}{1.000000,1.000000,1.000000}%
\pgfsetfillcolor{currentfill}%
\pgfsetfillopacity{0.900000}%
\pgfsetlinewidth{0.200750pt}%
\definecolor{currentstroke}{rgb}{0.200000,0.200000,0.200000}%
\pgfsetstrokecolor{currentstroke}%
\pgfsetdash{}{0pt}%
\pgfpathmoveto{\pgfqpoint{1.546916in}{1.492358in}}%
\pgfpathlineto{\pgfqpoint{1.593374in}{1.506945in}}%
\pgfpathlineto{\pgfqpoint{1.639731in}{1.521500in}}%
\pgfpathlineto{\pgfqpoint{1.670062in}{1.500704in}}%
\pgfpathlineto{\pgfqpoint{1.700487in}{1.479842in}}%
\pgfpathlineto{\pgfqpoint{1.654052in}{1.465196in}}%
\pgfpathlineto{\pgfqpoint{1.607516in}{1.450519in}}%
\pgfpathlineto{\pgfqpoint{1.577169in}{1.471471in}}%
\pgfpathlineto{\pgfqpoint{1.546916in}{1.492358in}}%
\pgfpathclose%
\pgfusepath{stroke,fill}%
\end{pgfscope}%
\begin{pgfscope}%
\pgfpathrectangle{\pgfqpoint{0.050000in}{0.050000in}}{\pgfqpoint{4.900000in}{4.900000in}}%
\pgfusepath{clip}%
\pgfsetbuttcap%
\pgfsetroundjoin%
\definecolor{currentfill}{rgb}{0.671819,0.658285,0.803768}%
\pgfsetfillcolor{currentfill}%
\pgfsetfillopacity{0.900000}%
\pgfsetlinewidth{0.200750pt}%
\definecolor{currentstroke}{rgb}{0.200000,0.200000,0.200000}%
\pgfsetstrokecolor{currentstroke}%
\pgfsetdash{}{0pt}%
\pgfpathmoveto{\pgfqpoint{3.573227in}{2.039536in}}%
\pgfpathlineto{\pgfqpoint{3.616785in}{2.013096in}}%
\pgfpathlineto{\pgfqpoint{3.659838in}{1.968820in}}%
\pgfpathlineto{\pgfqpoint{3.692838in}{1.926203in}}%
\pgfpathlineto{\pgfqpoint{3.725930in}{1.884211in}}%
\pgfpathlineto{\pgfqpoint{3.682928in}{1.928902in}}%
\pgfpathlineto{\pgfqpoint{3.639441in}{1.958078in}}%
\pgfpathlineto{\pgfqpoint{3.606291in}{1.998647in}}%
\pgfpathlineto{\pgfqpoint{3.573227in}{2.039536in}}%
\pgfpathclose%
\pgfusepath{stroke,fill}%
\end{pgfscope}%
\begin{pgfscope}%
\pgfpathrectangle{\pgfqpoint{0.050000in}{0.050000in}}{\pgfqpoint{4.900000in}{4.900000in}}%
\pgfusepath{clip}%
\pgfsetbuttcap%
\pgfsetroundjoin%
\definecolor{currentfill}{rgb}{1.000000,1.000000,1.000000}%
\pgfsetfillcolor{currentfill}%
\pgfsetfillopacity{0.900000}%
\pgfsetlinewidth{0.200750pt}%
\definecolor{currentstroke}{rgb}{0.200000,0.200000,0.200000}%
\pgfsetstrokecolor{currentstroke}%
\pgfsetdash{}{0pt}%
\pgfpathmoveto{\pgfqpoint{1.146919in}{1.488160in}}%
\pgfpathlineto{\pgfqpoint{1.193826in}{1.502756in}}%
\pgfpathlineto{\pgfqpoint{1.240631in}{1.517320in}}%
\pgfpathlineto{\pgfqpoint{1.270352in}{1.496511in}}%
\pgfpathlineto{\pgfqpoint{1.300165in}{1.475636in}}%
\pgfpathlineto{\pgfqpoint{1.253279in}{1.460981in}}%
\pgfpathlineto{\pgfqpoint{1.206290in}{1.446294in}}%
\pgfpathlineto{\pgfqpoint{1.176558in}{1.467260in}}%
\pgfpathlineto{\pgfqpoint{1.146919in}{1.488160in}}%
\pgfpathclose%
\pgfusepath{stroke,fill}%
\end{pgfscope}%
\begin{pgfscope}%
\pgfpathrectangle{\pgfqpoint{0.050000in}{0.050000in}}{\pgfqpoint{4.900000in}{4.900000in}}%
\pgfusepath{clip}%
\pgfsetbuttcap%
\pgfsetroundjoin%
\definecolor{currentfill}{rgb}{1.000000,1.000000,1.000000}%
\pgfsetfillcolor{currentfill}%
\pgfsetfillopacity{0.900000}%
\pgfsetlinewidth{0.200750pt}%
\definecolor{currentstroke}{rgb}{0.200000,0.200000,0.200000}%
\pgfsetstrokecolor{currentstroke}%
\pgfsetdash{}{0pt}%
\pgfpathmoveto{\pgfqpoint{2.500380in}{1.489345in}}%
\pgfpathlineto{\pgfqpoint{2.545816in}{1.505152in}}%
\pgfpathlineto{\pgfqpoint{2.591156in}{1.521989in}}%
\pgfpathlineto{\pgfqpoint{2.622977in}{1.501823in}}%
\pgfpathlineto{\pgfqpoint{2.654898in}{1.481639in}}%
\pgfpathlineto{\pgfqpoint{2.609482in}{1.464166in}}%
\pgfpathlineto{\pgfqpoint{2.563972in}{1.447918in}}%
\pgfpathlineto{\pgfqpoint{2.532126in}{1.468652in}}%
\pgfpathlineto{\pgfqpoint{2.500380in}{1.489345in}}%
\pgfpathclose%
\pgfusepath{stroke,fill}%
\end{pgfscope}%
\begin{pgfscope}%
\pgfpathrectangle{\pgfqpoint{0.050000in}{0.050000in}}{\pgfqpoint{4.900000in}{4.900000in}}%
\pgfusepath{clip}%
\pgfsetbuttcap%
\pgfsetroundjoin%
\definecolor{currentfill}{rgb}{0.671819,0.658285,0.803768}%
\pgfsetfillcolor{currentfill}%
\pgfsetfillopacity{0.900000}%
\pgfsetlinewidth{0.200750pt}%
\definecolor{currentstroke}{rgb}{0.200000,0.200000,0.200000}%
\pgfsetstrokecolor{currentstroke}%
\pgfsetdash{}{0pt}%
\pgfpathmoveto{\pgfqpoint{3.171962in}{1.907927in}}%
\pgfpathlineto{\pgfqpoint{3.217358in}{1.985385in}}%
\pgfpathlineto{\pgfqpoint{3.262708in}{2.058739in}}%
\pgfpathlineto{\pgfqpoint{3.295524in}{2.032179in}}%
\pgfpathlineto{\pgfqpoint{3.328429in}{2.005129in}}%
\pgfpathlineto{\pgfqpoint{3.283042in}{1.938668in}}%
\pgfpathlineto{\pgfqpoint{3.237575in}{1.866703in}}%
\pgfpathlineto{\pgfqpoint{3.204720in}{1.887585in}}%
\pgfpathlineto{\pgfqpoint{3.171962in}{1.907927in}}%
\pgfpathclose%
\pgfusepath{stroke,fill}%
\end{pgfscope}%
\begin{pgfscope}%
\pgfpathrectangle{\pgfqpoint{0.050000in}{0.050000in}}{\pgfqpoint{4.900000in}{4.900000in}}%
\pgfusepath{clip}%
\pgfsetbuttcap%
\pgfsetroundjoin%
\definecolor{currentfill}{rgb}{0.773256,0.763906,0.864421}%
\pgfsetfillcolor{currentfill}%
\pgfsetfillopacity{0.900000}%
\pgfsetlinewidth{0.200750pt}%
\definecolor{currentstroke}{rgb}{0.200000,0.200000,0.200000}%
\pgfsetstrokecolor{currentstroke}%
\pgfsetdash{}{0pt}%
\pgfpathmoveto{\pgfqpoint{3.081232in}{1.758751in}}%
\pgfpathlineto{\pgfqpoint{3.126575in}{1.831138in}}%
\pgfpathlineto{\pgfqpoint{3.171962in}{1.907927in}}%
\pgfpathlineto{\pgfqpoint{3.204720in}{1.887585in}}%
\pgfpathlineto{\pgfqpoint{3.237575in}{1.866703in}}%
\pgfpathlineto{\pgfqpoint{3.192092in}{1.793828in}}%
\pgfpathlineto{\pgfqpoint{3.146635in}{1.723812in}}%
\pgfpathlineto{\pgfqpoint{3.113883in}{1.741484in}}%
\pgfpathlineto{\pgfqpoint{3.081232in}{1.758751in}}%
\pgfpathclose%
\pgfusepath{stroke,fill}%
\end{pgfscope}%
\begin{pgfscope}%
\pgfpathrectangle{\pgfqpoint{0.050000in}{0.050000in}}{\pgfqpoint{4.900000in}{4.900000in}}%
\pgfusepath{clip}%
\pgfsetbuttcap%
\pgfsetroundjoin%
\definecolor{currentfill}{rgb}{0.844860,0.838462,0.907236}%
\pgfsetfillcolor{currentfill}%
\pgfsetfillopacity{0.900000}%
\pgfsetlinewidth{0.200750pt}%
\definecolor{currentstroke}{rgb}{0.200000,0.200000,0.200000}%
\pgfsetstrokecolor{currentstroke}%
\pgfsetdash{}{0pt}%
\pgfpathmoveto{\pgfqpoint{3.792389in}{1.801992in}}%
\pgfpathlineto{\pgfqpoint{3.834969in}{1.748806in}}%
\pgfpathlineto{\pgfqpoint{3.877293in}{1.693278in}}%
\pgfpathlineto{\pgfqpoint{3.910690in}{1.656335in}}%
\pgfpathlineto{\pgfqpoint{3.944186in}{1.619935in}}%
\pgfpathlineto{\pgfqpoint{3.901810in}{1.671215in}}%
\pgfpathlineto{\pgfqpoint{3.859224in}{1.721961in}}%
\pgfpathlineto{\pgfqpoint{3.825759in}{1.761715in}}%
\pgfpathlineto{\pgfqpoint{3.792389in}{1.801992in}}%
\pgfpathclose%
\pgfusepath{stroke,fill}%
\end{pgfscope}%
\begin{pgfscope}%
\pgfpathrectangle{\pgfqpoint{0.050000in}{0.050000in}}{\pgfqpoint{4.900000in}{4.900000in}}%
\pgfusepath{clip}%
\pgfsetbuttcap%
\pgfsetroundjoin%
\definecolor{currentfill}{rgb}{1.000000,1.000000,1.000000}%
\pgfsetfillcolor{currentfill}%
\pgfsetfillopacity{0.900000}%
\pgfsetlinewidth{0.200750pt}%
\definecolor{currentstroke}{rgb}{0.200000,0.200000,0.200000}%
\pgfsetstrokecolor{currentstroke}%
\pgfsetdash{}{0pt}%
\pgfpathmoveto{\pgfqpoint{4.252972in}{1.492531in}}%
\pgfpathlineto{\pgfqpoint{4.296489in}{1.507118in}}%
\pgfpathlineto{\pgfqpoint{4.330981in}{1.486277in}}%
\pgfpathlineto{\pgfqpoint{4.365581in}{1.465370in}}%
\pgfpathlineto{\pgfqpoint{4.322004in}{1.450693in}}%
\pgfpathlineto{\pgfqpoint{4.287434in}{1.471645in}}%
\pgfpathlineto{\pgfqpoint{4.252972in}{1.492531in}}%
\pgfpathclose%
\pgfusepath{stroke,fill}%
\end{pgfscope}%
\begin{pgfscope}%
\pgfpathrectangle{\pgfqpoint{0.050000in}{0.050000in}}{\pgfqpoint{4.900000in}{4.900000in}}%
\pgfusepath{clip}%
\pgfsetbuttcap%
\pgfsetroundjoin%
\definecolor{currentfill}{rgb}{1.000000,1.000000,1.000000}%
\pgfsetfillcolor{currentfill}%
\pgfsetfillopacity{0.900000}%
\pgfsetlinewidth{0.200750pt}%
\definecolor{currentstroke}{rgb}{0.200000,0.200000,0.200000}%
\pgfsetstrokecolor{currentstroke}%
\pgfsetdash{}{0pt}%
\pgfpathmoveto{\pgfqpoint{2.100559in}{1.484046in}}%
\pgfpathlineto{\pgfqpoint{2.146444in}{1.498653in}}%
\pgfpathlineto{\pgfqpoint{2.192229in}{1.513232in}}%
\pgfpathlineto{\pgfqpoint{2.223437in}{1.492413in}}%
\pgfpathlineto{\pgfqpoint{2.254742in}{1.471530in}}%
\pgfpathlineto{\pgfqpoint{2.208882in}{1.456855in}}%
\pgfpathlineto{\pgfqpoint{2.162923in}{1.442156in}}%
\pgfpathlineto{\pgfqpoint{2.131693in}{1.463134in}}%
\pgfpathlineto{\pgfqpoint{2.100559in}{1.484046in}}%
\pgfpathclose%
\pgfusepath{stroke,fill}%
\end{pgfscope}%
\begin{pgfscope}%
\pgfpathrectangle{\pgfqpoint{0.050000in}{0.050000in}}{\pgfqpoint{4.900000in}{4.900000in}}%
\pgfusepath{clip}%
\pgfsetbuttcap%
\pgfsetroundjoin%
\definecolor{currentfill}{rgb}{0.862760,0.857101,0.917939}%
\pgfsetfillcolor{currentfill}%
\pgfsetfillopacity{0.900000}%
\pgfsetlinewidth{0.200750pt}%
\definecolor{currentstroke}{rgb}{0.200000,0.200000,0.200000}%
\pgfsetstrokecolor{currentstroke}%
\pgfsetdash{}{0pt}%
\pgfpathmoveto{\pgfqpoint{2.990699in}{1.636069in}}%
\pgfpathlineto{\pgfqpoint{3.035942in}{1.693287in}}%
\pgfpathlineto{\pgfqpoint{3.081232in}{1.758751in}}%
\pgfpathlineto{\pgfqpoint{3.113883in}{1.741484in}}%
\pgfpathlineto{\pgfqpoint{3.146635in}{1.723812in}}%
\pgfpathlineto{\pgfqpoint{3.101223in}{1.659372in}}%
\pgfpathlineto{\pgfqpoint{3.055856in}{1.602126in}}%
\pgfpathlineto{\pgfqpoint{3.023225in}{1.619209in}}%
\pgfpathlineto{\pgfqpoint{2.990699in}{1.636069in}}%
\pgfpathclose%
\pgfusepath{stroke,fill}%
\end{pgfscope}%
\begin{pgfscope}%
\pgfpathrectangle{\pgfqpoint{0.050000in}{0.050000in}}{\pgfqpoint{4.900000in}{4.900000in}}%
\pgfusepath{clip}%
\pgfsetbuttcap%
\pgfsetroundjoin%
\definecolor{currentfill}{rgb}{0.588281,0.571303,0.753818}%
\pgfsetfillcolor{currentfill}%
\pgfsetfillopacity{0.900000}%
\pgfsetlinewidth{0.200750pt}%
\definecolor{currentstroke}{rgb}{0.200000,0.200000,0.200000}%
\pgfsetstrokecolor{currentstroke}%
\pgfsetdash{}{0pt}%
\pgfpathmoveto{\pgfqpoint{3.418611in}{2.101675in}}%
\pgfpathlineto{\pgfqpoint{3.463207in}{2.122908in}}%
\pgfpathlineto{\pgfqpoint{3.507354in}{2.122272in}}%
\pgfpathlineto{\pgfqpoint{3.540248in}{2.080745in}}%
\pgfpathlineto{\pgfqpoint{3.573227in}{2.039536in}}%
\pgfpathlineto{\pgfqpoint{3.529162in}{2.045698in}}%
\pgfpathlineto{\pgfqpoint{3.484634in}{2.031357in}}%
\pgfpathlineto{\pgfqpoint{3.451580in}{2.066576in}}%
\pgfpathlineto{\pgfqpoint{3.418611in}{2.101675in}}%
\pgfpathclose%
\pgfusepath{stroke,fill}%
\end{pgfscope}%
\begin{pgfscope}%
\pgfpathrectangle{\pgfqpoint{0.050000in}{0.050000in}}{\pgfqpoint{4.900000in}{4.900000in}}%
\pgfusepath{clip}%
\pgfsetbuttcap%
\pgfsetroundjoin%
\definecolor{currentfill}{rgb}{1.000000,1.000000,1.000000}%
\pgfsetfillcolor{currentfill}%
\pgfsetfillopacity{0.900000}%
\pgfsetlinewidth{0.200750pt}%
\definecolor{currentstroke}{rgb}{0.200000,0.200000,0.200000}%
\pgfsetstrokecolor{currentstroke}%
\pgfsetdash{}{0pt}%
\pgfpathmoveto{\pgfqpoint{1.700487in}{1.479842in}}%
\pgfpathlineto{\pgfqpoint{1.746821in}{1.494456in}}%
\pgfpathlineto{\pgfqpoint{1.793055in}{1.509039in}}%
\pgfpathlineto{\pgfqpoint{1.823652in}{1.488203in}}%
\pgfpathlineto{\pgfqpoint{1.854345in}{1.467303in}}%
\pgfpathlineto{\pgfqpoint{1.808034in}{1.452630in}}%
\pgfpathlineto{\pgfqpoint{1.761623in}{1.437924in}}%
\pgfpathlineto{\pgfqpoint{1.731007in}{1.458916in}}%
\pgfpathlineto{\pgfqpoint{1.700487in}{1.479842in}}%
\pgfpathclose%
\pgfusepath{stroke,fill}%
\end{pgfscope}%
\begin{pgfscope}%
\pgfpathrectangle{\pgfqpoint{0.050000in}{0.050000in}}{\pgfqpoint{4.900000in}{4.900000in}}%
\pgfusepath{clip}%
\pgfsetbuttcap%
\pgfsetroundjoin%
\definecolor{currentfill}{rgb}{1.000000,1.000000,1.000000}%
\pgfsetfillcolor{currentfill}%
\pgfsetfillopacity{0.900000}%
\pgfsetlinewidth{0.200750pt}%
\definecolor{currentstroke}{rgb}{0.200000,0.200000,0.200000}%
\pgfsetstrokecolor{currentstroke}%
\pgfsetdash{}{0pt}%
\pgfpathmoveto{\pgfqpoint{1.300165in}{1.475636in}}%
\pgfpathlineto{\pgfqpoint{1.346949in}{1.490259in}}%
\pgfpathlineto{\pgfqpoint{1.393631in}{1.504851in}}%
\pgfpathlineto{\pgfqpoint{1.423617in}{1.484002in}}%
\pgfpathlineto{\pgfqpoint{1.453697in}{1.463089in}}%
\pgfpathlineto{\pgfqpoint{1.406935in}{1.448407in}}%
\pgfpathlineto{\pgfqpoint{1.360070in}{1.433692in}}%
\pgfpathlineto{\pgfqpoint{1.330071in}{1.454697in}}%
\pgfpathlineto{\pgfqpoint{1.300165in}{1.475636in}}%
\pgfpathclose%
\pgfusepath{stroke,fill}%
\end{pgfscope}%
\begin{pgfscope}%
\pgfpathrectangle{\pgfqpoint{0.050000in}{0.050000in}}{\pgfqpoint{4.900000in}{4.900000in}}%
\pgfusepath{clip}%
\pgfsetbuttcap%
\pgfsetroundjoin%
\definecolor{currentfill}{rgb}{0.994033,0.993787,0.996432}%
\pgfsetfillcolor{currentfill}%
\pgfsetfillopacity{0.900000}%
\pgfsetlinewidth{0.200750pt}%
\definecolor{currentstroke}{rgb}{0.200000,0.200000,0.200000}%
\pgfsetstrokecolor{currentstroke}%
\pgfsetdash{}{0pt}%
\pgfpathmoveto{\pgfqpoint{2.654898in}{1.481639in}}%
\pgfpathlineto{\pgfqpoint{2.700224in}{1.501075in}}%
\pgfpathlineto{\pgfqpoint{2.745469in}{1.523468in}}%
\pgfpathlineto{\pgfqpoint{2.777573in}{1.504292in}}%
\pgfpathlineto{\pgfqpoint{2.809780in}{1.485097in}}%
\pgfpathlineto{\pgfqpoint{2.764448in}{1.461553in}}%
\pgfpathlineto{\pgfqpoint{2.719041in}{1.441208in}}%
\pgfpathlineto{\pgfqpoint{2.686919in}{1.461434in}}%
\pgfpathlineto{\pgfqpoint{2.654898in}{1.481639in}}%
\pgfpathclose%
\pgfusepath{stroke,fill}%
\end{pgfscope}%
\begin{pgfscope}%
\pgfpathrectangle{\pgfqpoint{0.050000in}{0.050000in}}{\pgfqpoint{4.900000in}{4.900000in}}%
\pgfusepath{clip}%
\pgfsetbuttcap%
\pgfsetroundjoin%
\definecolor{currentfill}{rgb}{0.928397,0.925444,0.957186}%
\pgfsetfillcolor{currentfill}%
\pgfsetfillopacity{0.900000}%
\pgfsetlinewidth{0.200750pt}%
\definecolor{currentstroke}{rgb}{0.200000,0.200000,0.200000}%
\pgfsetstrokecolor{currentstroke}%
\pgfsetdash{}{0pt}%
\pgfpathmoveto{\pgfqpoint{2.900278in}{1.546714in}}%
\pgfpathlineto{\pgfqpoint{2.945486in}{1.587381in}}%
\pgfpathlineto{\pgfqpoint{2.990699in}{1.636069in}}%
\pgfpathlineto{\pgfqpoint{3.023225in}{1.619209in}}%
\pgfpathlineto{\pgfqpoint{3.055856in}{1.602126in}}%
\pgfpathlineto{\pgfqpoint{3.010522in}{1.552684in}}%
\pgfpathlineto{\pgfqpoint{2.965200in}{1.510850in}}%
\pgfpathlineto{\pgfqpoint{2.932687in}{1.528826in}}%
\pgfpathlineto{\pgfqpoint{2.900278in}{1.546714in}}%
\pgfpathclose%
\pgfusepath{stroke,fill}%
\end{pgfscope}%
\begin{pgfscope}%
\pgfpathrectangle{\pgfqpoint{0.050000in}{0.050000in}}{\pgfqpoint{4.900000in}{4.900000in}}%
\pgfusepath{clip}%
\pgfsetbuttcap%
\pgfsetroundjoin%
\definecolor{currentfill}{rgb}{1.000000,1.000000,1.000000}%
\pgfsetfillcolor{currentfill}%
\pgfsetfillopacity{0.900000}%
\pgfsetlinewidth{0.200750pt}%
\definecolor{currentstroke}{rgb}{0.200000,0.200000,0.200000}%
\pgfsetstrokecolor{currentstroke}%
\pgfsetdash{}{0pt}%
\pgfpathmoveto{\pgfqpoint{2.254742in}{1.471530in}}%
\pgfpathlineto{\pgfqpoint{2.300501in}{1.486193in}}%
\pgfpathlineto{\pgfqpoint{2.346161in}{1.500873in}}%
\pgfpathlineto{\pgfqpoint{2.377637in}{1.480048in}}%
\pgfpathlineto{\pgfqpoint{2.409212in}{1.459165in}}%
\pgfpathlineto{\pgfqpoint{2.363478in}{1.444348in}}%
\pgfpathlineto{\pgfqpoint{2.317645in}{1.429574in}}%
\pgfpathlineto{\pgfqpoint{2.286144in}{1.450584in}}%
\pgfpathlineto{\pgfqpoint{2.254742in}{1.471530in}}%
\pgfpathclose%
\pgfusepath{stroke,fill}%
\end{pgfscope}%
\begin{pgfscope}%
\pgfpathrectangle{\pgfqpoint{0.050000in}{0.050000in}}{\pgfqpoint{4.900000in}{4.900000in}}%
\pgfusepath{clip}%
\pgfsetbuttcap%
\pgfsetroundjoin%
\definecolor{currentfill}{rgb}{0.982099,0.981361,0.989296}%
\pgfsetfillcolor{currentfill}%
\pgfsetfillopacity{0.900000}%
\pgfsetlinewidth{0.200750pt}%
\definecolor{currentstroke}{rgb}{0.200000,0.200000,0.200000}%
\pgfsetstrokecolor{currentstroke}%
\pgfsetdash{}{0pt}%
\pgfpathmoveto{\pgfqpoint{4.011481in}{1.548646in}}%
\pgfpathlineto{\pgfqpoint{4.053925in}{1.508998in}}%
\pgfpathlineto{\pgfqpoint{4.097170in}{1.505024in}}%
\pgfpathlineto{\pgfqpoint{4.131359in}{1.484176in}}%
\pgfpathlineto{\pgfqpoint{4.165653in}{1.463263in}}%
\pgfpathlineto{\pgfqpoint{4.121851in}{1.448581in}}%
\pgfpathlineto{\pgfqpoint{4.079184in}{1.479188in}}%
\pgfpathlineto{\pgfqpoint{4.045281in}{1.513701in}}%
\pgfpathlineto{\pgfqpoint{4.011481in}{1.548646in}}%
\pgfpathclose%
\pgfusepath{stroke,fill}%
\end{pgfscope}%
\begin{pgfscope}%
\pgfpathrectangle{\pgfqpoint{0.050000in}{0.050000in}}{\pgfqpoint{4.900000in}{4.900000in}}%
\pgfusepath{clip}%
\pgfsetbuttcap%
\pgfsetroundjoin%
\definecolor{currentfill}{rgb}{1.000000,1.000000,1.000000}%
\pgfsetfillcolor{currentfill}%
\pgfsetfillopacity{0.900000}%
\pgfsetlinewidth{0.200750pt}%
\definecolor{currentstroke}{rgb}{0.200000,0.200000,0.200000}%
\pgfsetstrokecolor{currentstroke}%
\pgfsetdash{}{0pt}%
\pgfpathmoveto{\pgfqpoint{1.854345in}{1.467303in}}%
\pgfpathlineto{\pgfqpoint{1.900554in}{1.481944in}}%
\pgfpathlineto{\pgfqpoint{1.946664in}{1.496554in}}%
\pgfpathlineto{\pgfqpoint{1.977529in}{1.475680in}}%
\pgfpathlineto{\pgfqpoint{2.008490in}{1.454740in}}%
\pgfpathlineto{\pgfqpoint{1.962304in}{1.440040in}}%
\pgfpathlineto{\pgfqpoint{1.916018in}{1.425307in}}%
\pgfpathlineto{\pgfqpoint{1.885133in}{1.446338in}}%
\pgfpathlineto{\pgfqpoint{1.854345in}{1.467303in}}%
\pgfpathclose%
\pgfusepath{stroke,fill}%
\end{pgfscope}%
\begin{pgfscope}%
\pgfpathrectangle{\pgfqpoint{0.050000in}{0.050000in}}{\pgfqpoint{4.900000in}{4.900000in}}%
\pgfusepath{clip}%
\pgfsetbuttcap%
\pgfsetroundjoin%
\definecolor{currentfill}{rgb}{1.000000,1.000000,1.000000}%
\pgfsetfillcolor{currentfill}%
\pgfsetfillopacity{0.900000}%
\pgfsetlinewidth{0.200750pt}%
\definecolor{currentstroke}{rgb}{0.200000,0.200000,0.200000}%
\pgfsetstrokecolor{currentstroke}%
\pgfsetdash{}{0pt}%
\pgfpathmoveto{\pgfqpoint{4.165653in}{1.463263in}}%
\pgfpathlineto{\pgfqpoint{4.209360in}{1.477913in}}%
\pgfpathlineto{\pgfqpoint{4.252972in}{1.492531in}}%
\pgfpathlineto{\pgfqpoint{4.287434in}{1.471645in}}%
\pgfpathlineto{\pgfqpoint{4.322004in}{1.450693in}}%
\pgfpathlineto{\pgfqpoint{4.278332in}{1.435983in}}%
\pgfpathlineto{\pgfqpoint{4.234564in}{1.421241in}}%
\pgfpathlineto{\pgfqpoint{4.200055in}{1.442285in}}%
\pgfpathlineto{\pgfqpoint{4.165653in}{1.463263in}}%
\pgfpathclose%
\pgfusepath{stroke,fill}%
\end{pgfscope}%
\begin{pgfscope}%
\pgfpathrectangle{\pgfqpoint{0.050000in}{0.050000in}}{\pgfqpoint{4.900000in}{4.900000in}}%
\pgfusepath{clip}%
\pgfsetbuttcap%
\pgfsetroundjoin%
\definecolor{currentfill}{rgb}{1.000000,1.000000,1.000000}%
\pgfsetfillcolor{currentfill}%
\pgfsetfillopacity{0.900000}%
\pgfsetlinewidth{0.200750pt}%
\definecolor{currentstroke}{rgb}{0.200000,0.200000,0.200000}%
\pgfsetstrokecolor{currentstroke}%
\pgfsetdash{}{0pt}%
\pgfpathmoveto{\pgfqpoint{1.453697in}{1.463089in}}%
\pgfpathlineto{\pgfqpoint{1.500357in}{1.477740in}}%
\pgfpathlineto{\pgfqpoint{1.546916in}{1.492358in}}%
\pgfpathlineto{\pgfqpoint{1.577169in}{1.471471in}}%
\pgfpathlineto{\pgfqpoint{1.607516in}{1.450519in}}%
\pgfpathlineto{\pgfqpoint{1.560878in}{1.435809in}}%
\pgfpathlineto{\pgfqpoint{1.514138in}{1.421066in}}%
\pgfpathlineto{\pgfqpoint{1.483870in}{1.442111in}}%
\pgfpathlineto{\pgfqpoint{1.453697in}{1.463089in}}%
\pgfpathclose%
\pgfusepath{stroke,fill}%
\end{pgfscope}%
\begin{pgfscope}%
\pgfpathrectangle{\pgfqpoint{0.050000in}{0.050000in}}{\pgfqpoint{4.900000in}{4.900000in}}%
\pgfusepath{clip}%
\pgfsetbuttcap%
\pgfsetroundjoin%
\definecolor{currentfill}{rgb}{0.707620,0.695563,0.825175}%
\pgfsetfillcolor{currentfill}%
\pgfsetfillopacity{0.900000}%
\pgfsetlinewidth{0.200750pt}%
\definecolor{currentstroke}{rgb}{0.200000,0.200000,0.200000}%
\pgfsetstrokecolor{currentstroke}%
\pgfsetdash{}{0pt}%
\pgfpathmoveto{\pgfqpoint{3.639441in}{1.958078in}}%
\pgfpathlineto{\pgfqpoint{3.682928in}{1.928902in}}%
\pgfpathlineto{\pgfqpoint{3.725930in}{1.884211in}}%
\pgfpathlineto{\pgfqpoint{3.759113in}{1.842816in}}%
\pgfpathlineto{\pgfqpoint{3.792389in}{1.801992in}}%
\pgfpathlineto{\pgfqpoint{3.749428in}{1.846627in}}%
\pgfpathlineto{\pgfqpoint{3.706006in}{1.877891in}}%
\pgfpathlineto{\pgfqpoint{3.672680in}{1.917826in}}%
\pgfpathlineto{\pgfqpoint{3.639441in}{1.958078in}}%
\pgfpathclose%
\pgfusepath{stroke,fill}%
\end{pgfscope}%
\begin{pgfscope}%
\pgfpathrectangle{\pgfqpoint{0.050000in}{0.050000in}}{\pgfqpoint{4.900000in}{4.900000in}}%
\pgfusepath{clip}%
\pgfsetbuttcap%
\pgfsetroundjoin%
\definecolor{currentfill}{rgb}{1.000000,1.000000,1.000000}%
\pgfsetfillcolor{currentfill}%
\pgfsetfillopacity{0.900000}%
\pgfsetlinewidth{0.200750pt}%
\definecolor{currentstroke}{rgb}{0.200000,0.200000,0.200000}%
\pgfsetstrokecolor{currentstroke}%
\pgfsetdash{}{0pt}%
\pgfpathmoveto{\pgfqpoint{1.052798in}{1.458873in}}%
\pgfpathlineto{\pgfqpoint{1.099910in}{1.473532in}}%
\pgfpathlineto{\pgfqpoint{1.146919in}{1.488160in}}%
\pgfpathlineto{\pgfqpoint{1.176558in}{1.467260in}}%
\pgfpathlineto{\pgfqpoint{1.206290in}{1.446294in}}%
\pgfpathlineto{\pgfqpoint{1.159199in}{1.431575in}}%
\pgfpathlineto{\pgfqpoint{1.112005in}{1.416823in}}%
\pgfpathlineto{\pgfqpoint{1.082355in}{1.437881in}}%
\pgfpathlineto{\pgfqpoint{1.052798in}{1.458873in}}%
\pgfpathclose%
\pgfusepath{stroke,fill}%
\end{pgfscope}%
\begin{pgfscope}%
\pgfpathrectangle{\pgfqpoint{0.050000in}{0.050000in}}{\pgfqpoint{4.900000in}{4.900000in}}%
\pgfusepath{clip}%
\pgfsetbuttcap%
\pgfsetroundjoin%
\definecolor{currentfill}{rgb}{1.000000,1.000000,1.000000}%
\pgfsetfillcolor{currentfill}%
\pgfsetfillopacity{0.900000}%
\pgfsetlinewidth{0.200750pt}%
\definecolor{currentstroke}{rgb}{0.200000,0.200000,0.200000}%
\pgfsetstrokecolor{currentstroke}%
\pgfsetdash{}{0pt}%
\pgfpathmoveto{\pgfqpoint{2.409212in}{1.459165in}}%
\pgfpathlineto{\pgfqpoint{2.454845in}{1.474110in}}%
\pgfpathlineto{\pgfqpoint{2.500380in}{1.489345in}}%
\pgfpathlineto{\pgfqpoint{2.532126in}{1.468652in}}%
\pgfpathlineto{\pgfqpoint{2.563972in}{1.447918in}}%
\pgfpathlineto{\pgfqpoint{2.518364in}{1.432381in}}%
\pgfpathlineto{\pgfqpoint{2.472657in}{1.417228in}}%
\pgfpathlineto{\pgfqpoint{2.440885in}{1.438225in}}%
\pgfpathlineto{\pgfqpoint{2.409212in}{1.459165in}}%
\pgfpathclose%
\pgfusepath{stroke,fill}%
\end{pgfscope}%
\begin{pgfscope}%
\pgfpathrectangle{\pgfqpoint{0.050000in}{0.050000in}}{\pgfqpoint{4.900000in}{4.900000in}}%
\pgfusepath{clip}%
\pgfsetbuttcap%
\pgfsetroundjoin%
\definecolor{currentfill}{rgb}{0.612149,0.596155,0.768089}%
\pgfsetfillcolor{currentfill}%
\pgfsetfillopacity{0.900000}%
\pgfsetlinewidth{0.200750pt}%
\definecolor{currentstroke}{rgb}{0.200000,0.200000,0.200000}%
\pgfsetstrokecolor{currentstroke}%
\pgfsetdash{}{0pt}%
\pgfpathmoveto{\pgfqpoint{3.328429in}{2.005129in}}%
\pgfpathlineto{\pgfqpoint{3.373652in}{2.061079in}}%
\pgfpathlineto{\pgfqpoint{3.418611in}{2.101675in}}%
\pgfpathlineto{\pgfqpoint{3.451580in}{2.066576in}}%
\pgfpathlineto{\pgfqpoint{3.484634in}{2.031357in}}%
\pgfpathlineto{\pgfqpoint{3.439718in}{1.998263in}}%
\pgfpathlineto{\pgfqpoint{3.394507in}{1.949695in}}%
\pgfpathlineto{\pgfqpoint{3.361423in}{1.977624in}}%
\pgfpathlineto{\pgfqpoint{3.328429in}{2.005129in}}%
\pgfpathclose%
\pgfusepath{stroke,fill}%
\end{pgfscope}%
\begin{pgfscope}%
\pgfpathrectangle{\pgfqpoint{0.050000in}{0.050000in}}{\pgfqpoint{4.900000in}{4.900000in}}%
\pgfusepath{clip}%
\pgfsetbuttcap%
\pgfsetroundjoin%
\definecolor{currentfill}{rgb}{1.000000,1.000000,1.000000}%
\pgfsetfillcolor{currentfill}%
\pgfsetfillopacity{0.900000}%
\pgfsetlinewidth{0.200750pt}%
\definecolor{currentstroke}{rgb}{0.200000,0.200000,0.200000}%
\pgfsetstrokecolor{currentstroke}%
\pgfsetdash{}{0pt}%
\pgfpathmoveto{\pgfqpoint{2.008490in}{1.454740in}}%
\pgfpathlineto{\pgfqpoint{2.054575in}{1.469409in}}%
\pgfpathlineto{\pgfqpoint{2.100559in}{1.484046in}}%
\pgfpathlineto{\pgfqpoint{2.131693in}{1.463134in}}%
\pgfpathlineto{\pgfqpoint{2.162923in}{1.442156in}}%
\pgfpathlineto{\pgfqpoint{2.116863in}{1.427426in}}%
\pgfpathlineto{\pgfqpoint{2.070702in}{1.412665in}}%
\pgfpathlineto{\pgfqpoint{2.039548in}{1.433736in}}%
\pgfpathlineto{\pgfqpoint{2.008490in}{1.454740in}}%
\pgfpathclose%
\pgfusepath{stroke,fill}%
\end{pgfscope}%
\begin{pgfscope}%
\pgfpathrectangle{\pgfqpoint{0.050000in}{0.050000in}}{\pgfqpoint{4.900000in}{4.900000in}}%
\pgfusepath{clip}%
\pgfsetbuttcap%
\pgfsetroundjoin%
\definecolor{currentfill}{rgb}{0.970165,0.968935,0.982161}%
\pgfsetfillcolor{currentfill}%
\pgfsetfillopacity{0.900000}%
\pgfsetlinewidth{0.200750pt}%
\definecolor{currentstroke}{rgb}{0.200000,0.200000,0.200000}%
\pgfsetstrokecolor{currentstroke}%
\pgfsetdash{}{0pt}%
\pgfpathmoveto{\pgfqpoint{2.809780in}{1.485097in}}%
\pgfpathlineto{\pgfqpoint{2.855050in}{1.513045in}}%
\pgfpathlineto{\pgfqpoint{2.900278in}{1.546714in}}%
\pgfpathlineto{\pgfqpoint{2.932687in}{1.528826in}}%
\pgfpathlineto{\pgfqpoint{2.965200in}{1.510850in}}%
\pgfpathlineto{\pgfqpoint{2.919867in}{1.475861in}}%
\pgfpathlineto{\pgfqpoint{2.874501in}{1.446628in}}%
\pgfpathlineto{\pgfqpoint{2.842089in}{1.465878in}}%
\pgfpathlineto{\pgfqpoint{2.809780in}{1.485097in}}%
\pgfpathclose%
\pgfusepath{stroke,fill}%
\end{pgfscope}%
\begin{pgfscope}%
\pgfpathrectangle{\pgfqpoint{0.050000in}{0.050000in}}{\pgfqpoint{4.900000in}{4.900000in}}%
\pgfusepath{clip}%
\pgfsetbuttcap%
\pgfsetroundjoin%
\definecolor{currentfill}{rgb}{0.874694,0.869527,0.925075}%
\pgfsetfillcolor{currentfill}%
\pgfsetfillopacity{0.900000}%
\pgfsetlinewidth{0.200750pt}%
\definecolor{currentstroke}{rgb}{0.200000,0.200000,0.200000}%
\pgfsetstrokecolor{currentstroke}%
\pgfsetdash{}{0pt}%
\pgfpathmoveto{\pgfqpoint{3.859224in}{1.721961in}}%
\pgfpathlineto{\pgfqpoint{3.901810in}{1.671215in}}%
\pgfpathlineto{\pgfqpoint{3.944186in}{1.619935in}}%
\pgfpathlineto{\pgfqpoint{3.977783in}{1.584048in}}%
\pgfpathlineto{\pgfqpoint{4.011481in}{1.548646in}}%
\pgfpathlineto{\pgfqpoint{3.969045in}{1.595753in}}%
\pgfpathlineto{\pgfqpoint{3.926442in}{1.643932in}}%
\pgfpathlineto{\pgfqpoint{3.892785in}{1.682707in}}%
\pgfpathlineto{\pgfqpoint{3.859224in}{1.721961in}}%
\pgfpathclose%
\pgfusepath{stroke,fill}%
\end{pgfscope}%
\begin{pgfscope}%
\pgfpathrectangle{\pgfqpoint{0.050000in}{0.050000in}}{\pgfqpoint{4.900000in}{4.900000in}}%
\pgfusepath{clip}%
\pgfsetbuttcap%
\pgfsetroundjoin%
\definecolor{currentfill}{rgb}{1.000000,1.000000,1.000000}%
\pgfsetfillcolor{currentfill}%
\pgfsetfillopacity{0.900000}%
\pgfsetlinewidth{0.200750pt}%
\definecolor{currentstroke}{rgb}{0.200000,0.200000,0.200000}%
\pgfsetstrokecolor{currentstroke}%
\pgfsetdash{}{0pt}%
\pgfpathmoveto{\pgfqpoint{1.607516in}{1.450519in}}%
\pgfpathlineto{\pgfqpoint{1.654052in}{1.465196in}}%
\pgfpathlineto{\pgfqpoint{1.700487in}{1.479842in}}%
\pgfpathlineto{\pgfqpoint{1.731007in}{1.458916in}}%
\pgfpathlineto{\pgfqpoint{1.761623in}{1.437924in}}%
\pgfpathlineto{\pgfqpoint{1.715109in}{1.423187in}}%
\pgfpathlineto{\pgfqpoint{1.668494in}{1.408417in}}%
\pgfpathlineto{\pgfqpoint{1.637958in}{1.429501in}}%
\pgfpathlineto{\pgfqpoint{1.607516in}{1.450519in}}%
\pgfpathclose%
\pgfusepath{stroke,fill}%
\end{pgfscope}%
\begin{pgfscope}%
\pgfpathrectangle{\pgfqpoint{0.050000in}{0.050000in}}{\pgfqpoint{4.900000in}{4.900000in}}%
\pgfusepath{clip}%
\pgfsetbuttcap%
\pgfsetroundjoin%
\definecolor{currentfill}{rgb}{0.683752,0.670711,0.810903}%
\pgfsetfillcolor{currentfill}%
\pgfsetfillopacity{0.900000}%
\pgfsetlinewidth{0.200750pt}%
\definecolor{currentstroke}{rgb}{0.200000,0.200000,0.200000}%
\pgfsetstrokecolor{currentstroke}%
\pgfsetdash{}{0pt}%
\pgfpathmoveto{\pgfqpoint{3.237575in}{1.866703in}}%
\pgfpathlineto{\pgfqpoint{3.283042in}{1.938668in}}%
\pgfpathlineto{\pgfqpoint{3.328429in}{2.005129in}}%
\pgfpathlineto{\pgfqpoint{3.361423in}{1.977624in}}%
\pgfpathlineto{\pgfqpoint{3.394507in}{1.949695in}}%
\pgfpathlineto{\pgfqpoint{3.349096in}{1.889883in}}%
\pgfpathlineto{\pgfqpoint{3.303573in}{1.823398in}}%
\pgfpathlineto{\pgfqpoint{3.270526in}{1.845301in}}%
\pgfpathlineto{\pgfqpoint{3.237575in}{1.866703in}}%
\pgfpathclose%
\pgfusepath{stroke,fill}%
\end{pgfscope}%
\begin{pgfscope}%
\pgfpathrectangle{\pgfqpoint{0.050000in}{0.050000in}}{\pgfqpoint{4.900000in}{4.900000in}}%
\pgfusepath{clip}%
\pgfsetbuttcap%
\pgfsetroundjoin%
\definecolor{currentfill}{rgb}{0.618116,0.602368,0.771657}%
\pgfsetfillcolor{currentfill}%
\pgfsetfillopacity{0.900000}%
\pgfsetlinewidth{0.200750pt}%
\definecolor{currentstroke}{rgb}{0.200000,0.200000,0.200000}%
\pgfsetstrokecolor{currentstroke}%
\pgfsetdash{}{0pt}%
\pgfpathmoveto{\pgfqpoint{3.484634in}{2.031357in}}%
\pgfpathlineto{\pgfqpoint{3.529162in}{2.045698in}}%
\pgfpathlineto{\pgfqpoint{3.573227in}{2.039536in}}%
\pgfpathlineto{\pgfqpoint{3.606291in}{1.998647in}}%
\pgfpathlineto{\pgfqpoint{3.639441in}{1.958078in}}%
\pgfpathlineto{\pgfqpoint{3.595455in}{1.968925in}}%
\pgfpathlineto{\pgfqpoint{3.550998in}{1.960662in}}%
\pgfpathlineto{\pgfqpoint{3.517773in}{1.996045in}}%
\pgfpathlineto{\pgfqpoint{3.484634in}{2.031357in}}%
\pgfpathclose%
\pgfusepath{stroke,fill}%
\end{pgfscope}%
\begin{pgfscope}%
\pgfpathrectangle{\pgfqpoint{0.050000in}{0.050000in}}{\pgfqpoint{4.900000in}{4.900000in}}%
\pgfusepath{clip}%
\pgfsetbuttcap%
\pgfsetroundjoin%
\definecolor{currentfill}{rgb}{1.000000,1.000000,1.000000}%
\pgfsetfillcolor{currentfill}%
\pgfsetfillopacity{0.900000}%
\pgfsetlinewidth{0.200750pt}%
\definecolor{currentstroke}{rgb}{0.200000,0.200000,0.200000}%
\pgfsetstrokecolor{currentstroke}%
\pgfsetdash{}{0pt}%
\pgfpathmoveto{\pgfqpoint{1.206290in}{1.446294in}}%
\pgfpathlineto{\pgfqpoint{1.253279in}{1.460981in}}%
\pgfpathlineto{\pgfqpoint{1.300165in}{1.475636in}}%
\pgfpathlineto{\pgfqpoint{1.330071in}{1.454697in}}%
\pgfpathlineto{\pgfqpoint{1.360070in}{1.433692in}}%
\pgfpathlineto{\pgfqpoint{1.313103in}{1.418945in}}%
\pgfpathlineto{\pgfqpoint{1.266033in}{1.404166in}}%
\pgfpathlineto{\pgfqpoint{1.236115in}{1.425263in}}%
\pgfpathlineto{\pgfqpoint{1.206290in}{1.446294in}}%
\pgfpathclose%
\pgfusepath{stroke,fill}%
\end{pgfscope}%
\begin{pgfscope}%
\pgfpathrectangle{\pgfqpoint{0.050000in}{0.050000in}}{\pgfqpoint{4.900000in}{4.900000in}}%
\pgfusepath{clip}%
\pgfsetbuttcap%
\pgfsetroundjoin%
\definecolor{currentfill}{rgb}{1.000000,1.000000,1.000000}%
\pgfsetfillcolor{currentfill}%
\pgfsetfillopacity{0.900000}%
\pgfsetlinewidth{0.200750pt}%
\definecolor{currentstroke}{rgb}{0.200000,0.200000,0.200000}%
\pgfsetstrokecolor{currentstroke}%
\pgfsetdash{}{0pt}%
\pgfpathmoveto{\pgfqpoint{2.563972in}{1.447918in}}%
\pgfpathlineto{\pgfqpoint{2.609482in}{1.464166in}}%
\pgfpathlineto{\pgfqpoint{2.654898in}{1.481639in}}%
\pgfpathlineto{\pgfqpoint{2.686919in}{1.461434in}}%
\pgfpathlineto{\pgfqpoint{2.719041in}{1.441208in}}%
\pgfpathlineto{\pgfqpoint{2.673548in}{1.423059in}}%
\pgfpathlineto{\pgfqpoint{2.627962in}{1.406326in}}%
\pgfpathlineto{\pgfqpoint{2.595917in}{1.427142in}}%
\pgfpathlineto{\pgfqpoint{2.563972in}{1.447918in}}%
\pgfpathclose%
\pgfusepath{stroke,fill}%
\end{pgfscope}%
\begin{pgfscope}%
\pgfpathrectangle{\pgfqpoint{0.050000in}{0.050000in}}{\pgfqpoint{4.900000in}{4.900000in}}%
\pgfusepath{clip}%
\pgfsetbuttcap%
\pgfsetroundjoin%
\definecolor{currentfill}{rgb}{0.773256,0.763906,0.864421}%
\pgfsetfillcolor{currentfill}%
\pgfsetfillopacity{0.900000}%
\pgfsetlinewidth{0.200750pt}%
\definecolor{currentstroke}{rgb}{0.200000,0.200000,0.200000}%
\pgfsetstrokecolor{currentstroke}%
\pgfsetdash{}{0pt}%
\pgfpathmoveto{\pgfqpoint{3.146635in}{1.723812in}}%
\pgfpathlineto{\pgfqpoint{3.192092in}{1.793828in}}%
\pgfpathlineto{\pgfqpoint{3.237575in}{1.866703in}}%
\pgfpathlineto{\pgfqpoint{3.270526in}{1.845301in}}%
\pgfpathlineto{\pgfqpoint{3.303573in}{1.823398in}}%
\pgfpathlineto{\pgfqpoint{3.258006in}{1.754611in}}%
\pgfpathlineto{\pgfqpoint{3.212445in}{1.687260in}}%
\pgfpathlineto{\pgfqpoint{3.179489in}{1.705736in}}%
\pgfpathlineto{\pgfqpoint{3.146635in}{1.723812in}}%
\pgfpathclose%
\pgfusepath{stroke,fill}%
\end{pgfscope}%
\begin{pgfscope}%
\pgfpathrectangle{\pgfqpoint{0.050000in}{0.050000in}}{\pgfqpoint{4.900000in}{4.900000in}}%
\pgfusepath{clip}%
\pgfsetbuttcap%
\pgfsetroundjoin%
\definecolor{currentfill}{rgb}{1.000000,1.000000,1.000000}%
\pgfsetfillcolor{currentfill}%
\pgfsetfillopacity{0.900000}%
\pgfsetlinewidth{0.200750pt}%
\definecolor{currentstroke}{rgb}{0.200000,0.200000,0.200000}%
\pgfsetstrokecolor{currentstroke}%
\pgfsetdash{}{0pt}%
\pgfpathmoveto{\pgfqpoint{4.322004in}{1.450693in}}%
\pgfpathlineto{\pgfqpoint{4.365581in}{1.465370in}}%
\pgfpathlineto{\pgfqpoint{4.400289in}{1.444399in}}%
\pgfpathlineto{\pgfqpoint{4.435106in}{1.423362in}}%
\pgfpathlineto{\pgfqpoint{4.391468in}{1.408592in}}%
\pgfpathlineto{\pgfqpoint{4.356682in}{1.429675in}}%
\pgfpathlineto{\pgfqpoint{4.322004in}{1.450693in}}%
\pgfpathclose%
\pgfusepath{stroke,fill}%
\end{pgfscope}%
\begin{pgfscope}%
\pgfpathrectangle{\pgfqpoint{0.050000in}{0.050000in}}{\pgfqpoint{4.900000in}{4.900000in}}%
\pgfusepath{clip}%
\pgfsetbuttcap%
\pgfsetroundjoin%
\definecolor{currentfill}{rgb}{1.000000,1.000000,1.000000}%
\pgfsetfillcolor{currentfill}%
\pgfsetfillopacity{0.900000}%
\pgfsetlinewidth{0.200750pt}%
\definecolor{currentstroke}{rgb}{0.200000,0.200000,0.200000}%
\pgfsetstrokecolor{currentstroke}%
\pgfsetdash{}{0pt}%
\pgfpathmoveto{\pgfqpoint{2.162923in}{1.442156in}}%
\pgfpathlineto{\pgfqpoint{2.208882in}{1.456855in}}%
\pgfpathlineto{\pgfqpoint{2.254742in}{1.471530in}}%
\pgfpathlineto{\pgfqpoint{2.286144in}{1.450584in}}%
\pgfpathlineto{\pgfqpoint{2.317645in}{1.429574in}}%
\pgfpathlineto{\pgfqpoint{2.271711in}{1.414798in}}%
\pgfpathlineto{\pgfqpoint{2.225677in}{1.400003in}}%
\pgfpathlineto{\pgfqpoint{2.194251in}{1.421112in}}%
\pgfpathlineto{\pgfqpoint{2.162923in}{1.442156in}}%
\pgfpathclose%
\pgfusepath{stroke,fill}%
\end{pgfscope}%
\begin{pgfscope}%
\pgfpathrectangle{\pgfqpoint{0.050000in}{0.050000in}}{\pgfqpoint{4.900000in}{4.900000in}}%
\pgfusepath{clip}%
\pgfsetbuttcap%
\pgfsetroundjoin%
\definecolor{currentfill}{rgb}{0.856794,0.850888,0.914371}%
\pgfsetfillcolor{currentfill}%
\pgfsetfillopacity{0.900000}%
\pgfsetlinewidth{0.200750pt}%
\definecolor{currentstroke}{rgb}{0.200000,0.200000,0.200000}%
\pgfsetstrokecolor{currentstroke}%
\pgfsetdash{}{0pt}%
\pgfpathmoveto{\pgfqpoint{3.055856in}{1.602126in}}%
\pgfpathlineto{\pgfqpoint{3.101223in}{1.659372in}}%
\pgfpathlineto{\pgfqpoint{3.146635in}{1.723812in}}%
\pgfpathlineto{\pgfqpoint{3.179489in}{1.705736in}}%
\pgfpathlineto{\pgfqpoint{3.212445in}{1.687260in}}%
\pgfpathlineto{\pgfqpoint{3.166917in}{1.624192in}}%
\pgfpathlineto{\pgfqpoint{3.121430in}{1.567261in}}%
\pgfpathlineto{\pgfqpoint{3.088590in}{1.584812in}}%
\pgfpathlineto{\pgfqpoint{3.055856in}{1.602126in}}%
\pgfpathclose%
\pgfusepath{stroke,fill}%
\end{pgfscope}%
\begin{pgfscope}%
\pgfpathrectangle{\pgfqpoint{0.050000in}{0.050000in}}{\pgfqpoint{4.900000in}{4.900000in}}%
\pgfusepath{clip}%
\pgfsetbuttcap%
\pgfsetroundjoin%
\definecolor{currentfill}{rgb}{1.000000,1.000000,1.000000}%
\pgfsetfillcolor{currentfill}%
\pgfsetfillopacity{0.900000}%
\pgfsetlinewidth{0.200750pt}%
\definecolor{currentstroke}{rgb}{0.200000,0.200000,0.200000}%
\pgfsetstrokecolor{currentstroke}%
\pgfsetdash{}{0pt}%
\pgfpathmoveto{\pgfqpoint{1.761623in}{1.437924in}}%
\pgfpathlineto{\pgfqpoint{1.808034in}{1.452630in}}%
\pgfpathlineto{\pgfqpoint{1.854345in}{1.467303in}}%
\pgfpathlineto{\pgfqpoint{1.885133in}{1.446338in}}%
\pgfpathlineto{\pgfqpoint{1.916018in}{1.425307in}}%
\pgfpathlineto{\pgfqpoint{1.869630in}{1.410542in}}%
\pgfpathlineto{\pgfqpoint{1.823141in}{1.395744in}}%
\pgfpathlineto{\pgfqpoint{1.792333in}{1.416867in}}%
\pgfpathlineto{\pgfqpoint{1.761623in}{1.437924in}}%
\pgfpathclose%
\pgfusepath{stroke,fill}%
\end{pgfscope}%
\begin{pgfscope}%
\pgfpathrectangle{\pgfqpoint{0.050000in}{0.050000in}}{\pgfqpoint{4.900000in}{4.900000in}}%
\pgfusepath{clip}%
\pgfsetbuttcap%
\pgfsetroundjoin%
\definecolor{currentfill}{rgb}{1.000000,1.000000,1.000000}%
\pgfsetfillcolor{currentfill}%
\pgfsetfillopacity{0.900000}%
\pgfsetlinewidth{0.200750pt}%
\definecolor{currentstroke}{rgb}{0.200000,0.200000,0.200000}%
\pgfsetstrokecolor{currentstroke}%
\pgfsetdash{}{0pt}%
\pgfpathmoveto{\pgfqpoint{1.360070in}{1.433692in}}%
\pgfpathlineto{\pgfqpoint{1.406935in}{1.448407in}}%
\pgfpathlineto{\pgfqpoint{1.453697in}{1.463089in}}%
\pgfpathlineto{\pgfqpoint{1.483870in}{1.442111in}}%
\pgfpathlineto{\pgfqpoint{1.514138in}{1.421066in}}%
\pgfpathlineto{\pgfqpoint{1.467295in}{1.406292in}}%
\pgfpathlineto{\pgfqpoint{1.420350in}{1.391485in}}%
\pgfpathlineto{\pgfqpoint{1.390163in}{1.412622in}}%
\pgfpathlineto{\pgfqpoint{1.360070in}{1.433692in}}%
\pgfpathclose%
\pgfusepath{stroke,fill}%
\end{pgfscope}%
\begin{pgfscope}%
\pgfpathrectangle{\pgfqpoint{0.050000in}{0.050000in}}{\pgfqpoint{4.900000in}{4.900000in}}%
\pgfusepath{clip}%
\pgfsetbuttcap%
\pgfsetroundjoin%
\definecolor{currentfill}{rgb}{0.994033,0.993787,0.996432}%
\pgfsetfillcolor{currentfill}%
\pgfsetfillopacity{0.900000}%
\pgfsetlinewidth{0.200750pt}%
\definecolor{currentstroke}{rgb}{0.200000,0.200000,0.200000}%
\pgfsetstrokecolor{currentstroke}%
\pgfsetdash{}{0pt}%
\pgfpathmoveto{\pgfqpoint{4.079184in}{1.479188in}}%
\pgfpathlineto{\pgfqpoint{4.121851in}{1.448581in}}%
\pgfpathlineto{\pgfqpoint{4.165653in}{1.463263in}}%
\pgfpathlineto{\pgfqpoint{4.200055in}{1.442285in}}%
\pgfpathlineto{\pgfqpoint{4.234564in}{1.421241in}}%
\pgfpathlineto{\pgfqpoint{4.190700in}{1.406467in}}%
\pgfpathlineto{\pgfqpoint{4.147300in}{1.411365in}}%
\pgfpathlineto{\pgfqpoint{4.113190in}{1.445083in}}%
\pgfpathlineto{\pgfqpoint{4.079184in}{1.479188in}}%
\pgfpathclose%
\pgfusepath{stroke,fill}%
\end{pgfscope}%
\begin{pgfscope}%
\pgfpathrectangle{\pgfqpoint{0.050000in}{0.050000in}}{\pgfqpoint{4.900000in}{4.900000in}}%
\pgfusepath{clip}%
\pgfsetbuttcap%
\pgfsetroundjoin%
\definecolor{currentfill}{rgb}{0.988066,0.987574,0.992864}%
\pgfsetfillcolor{currentfill}%
\pgfsetfillopacity{0.900000}%
\pgfsetlinewidth{0.200750pt}%
\definecolor{currentstroke}{rgb}{0.200000,0.200000,0.200000}%
\pgfsetstrokecolor{currentstroke}%
\pgfsetdash{}{0pt}%
\pgfpathmoveto{\pgfqpoint{2.719041in}{1.441208in}}%
\pgfpathlineto{\pgfqpoint{2.764448in}{1.461553in}}%
\pgfpathlineto{\pgfqpoint{2.809780in}{1.485097in}}%
\pgfpathlineto{\pgfqpoint{2.842089in}{1.465878in}}%
\pgfpathlineto{\pgfqpoint{2.874501in}{1.446628in}}%
\pgfpathlineto{\pgfqpoint{2.829080in}{1.421954in}}%
\pgfpathlineto{\pgfqpoint{2.783590in}{1.400686in}}%
\pgfpathlineto{\pgfqpoint{2.751264in}{1.420960in}}%
\pgfpathlineto{\pgfqpoint{2.719041in}{1.441208in}}%
\pgfpathclose%
\pgfusepath{stroke,fill}%
\end{pgfscope}%
\begin{pgfscope}%
\pgfpathrectangle{\pgfqpoint{0.050000in}{0.050000in}}{\pgfqpoint{4.900000in}{4.900000in}}%
\pgfusepath{clip}%
\pgfsetbuttcap%
\pgfsetroundjoin%
\definecolor{currentfill}{rgb}{1.000000,1.000000,1.000000}%
\pgfsetfillcolor{currentfill}%
\pgfsetfillopacity{0.900000}%
\pgfsetlinewidth{0.200750pt}%
\definecolor{currentstroke}{rgb}{0.200000,0.200000,0.200000}%
\pgfsetstrokecolor{currentstroke}%
\pgfsetdash{}{0pt}%
\pgfpathmoveto{\pgfqpoint{2.317645in}{1.429574in}}%
\pgfpathlineto{\pgfqpoint{2.363478in}{1.444348in}}%
\pgfpathlineto{\pgfqpoint{2.409212in}{1.459165in}}%
\pgfpathlineto{\pgfqpoint{2.440885in}{1.438225in}}%
\pgfpathlineto{\pgfqpoint{2.472657in}{1.417228in}}%
\pgfpathlineto{\pgfqpoint{2.426850in}{1.402258in}}%
\pgfpathlineto{\pgfqpoint{2.380942in}{1.387362in}}%
\pgfpathlineto{\pgfqpoint{2.349244in}{1.408500in}}%
\pgfpathlineto{\pgfqpoint{2.317645in}{1.429574in}}%
\pgfpathclose%
\pgfusepath{stroke,fill}%
\end{pgfscope}%
\begin{pgfscope}%
\pgfpathrectangle{\pgfqpoint{0.050000in}{0.050000in}}{\pgfqpoint{4.900000in}{4.900000in}}%
\pgfusepath{clip}%
\pgfsetbuttcap%
\pgfsetroundjoin%
\definecolor{currentfill}{rgb}{0.922430,0.919231,0.953618}%
\pgfsetfillcolor{currentfill}%
\pgfsetfillopacity{0.900000}%
\pgfsetlinewidth{0.200750pt}%
\definecolor{currentstroke}{rgb}{0.200000,0.200000,0.200000}%
\pgfsetstrokecolor{currentstroke}%
\pgfsetdash{}{0pt}%
\pgfpathmoveto{\pgfqpoint{2.965200in}{1.510850in}}%
\pgfpathlineto{\pgfqpoint{3.010522in}{1.552684in}}%
\pgfpathlineto{\pgfqpoint{3.055856in}{1.602126in}}%
\pgfpathlineto{\pgfqpoint{3.088590in}{1.584812in}}%
\pgfpathlineto{\pgfqpoint{3.121430in}{1.567261in}}%
\pgfpathlineto{\pgfqpoint{3.075977in}{1.517365in}}%
\pgfpathlineto{\pgfqpoint{3.030541in}{1.474597in}}%
\pgfpathlineto{\pgfqpoint{2.997818in}{1.492776in}}%
\pgfpathlineto{\pgfqpoint{2.965200in}{1.510850in}}%
\pgfpathclose%
\pgfusepath{stroke,fill}%
\end{pgfscope}%
\begin{pgfscope}%
\pgfpathrectangle{\pgfqpoint{0.050000in}{0.050000in}}{\pgfqpoint{4.900000in}{4.900000in}}%
\pgfusepath{clip}%
\pgfsetbuttcap%
\pgfsetroundjoin%
\definecolor{currentfill}{rgb}{0.743422,0.732841,0.846582}%
\pgfsetfillcolor{currentfill}%
\pgfsetfillopacity{0.900000}%
\pgfsetlinewidth{0.200750pt}%
\definecolor{currentstroke}{rgb}{0.200000,0.200000,0.200000}%
\pgfsetstrokecolor{currentstroke}%
\pgfsetdash{}{0pt}%
\pgfpathmoveto{\pgfqpoint{3.706006in}{1.877891in}}%
\pgfpathlineto{\pgfqpoint{3.749428in}{1.846627in}}%
\pgfpathlineto{\pgfqpoint{3.792389in}{1.801992in}}%
\pgfpathlineto{\pgfqpoint{3.825759in}{1.761715in}}%
\pgfpathlineto{\pgfqpoint{3.859224in}{1.721961in}}%
\pgfpathlineto{\pgfqpoint{3.816294in}{1.766152in}}%
\pgfpathlineto{\pgfqpoint{3.772929in}{1.798951in}}%
\pgfpathlineto{\pgfqpoint{3.739423in}{1.838267in}}%
\pgfpathlineto{\pgfqpoint{3.706006in}{1.877891in}}%
\pgfpathclose%
\pgfusepath{stroke,fill}%
\end{pgfscope}%
\begin{pgfscope}%
\pgfpathrectangle{\pgfqpoint{0.050000in}{0.050000in}}{\pgfqpoint{4.900000in}{4.900000in}}%
\pgfusepath{clip}%
\pgfsetbuttcap%
\pgfsetroundjoin%
\definecolor{currentfill}{rgb}{1.000000,1.000000,1.000000}%
\pgfsetfillcolor{currentfill}%
\pgfsetfillopacity{0.900000}%
\pgfsetlinewidth{0.200750pt}%
\definecolor{currentstroke}{rgb}{0.200000,0.200000,0.200000}%
\pgfsetstrokecolor{currentstroke}%
\pgfsetdash{}{0pt}%
\pgfpathmoveto{\pgfqpoint{1.916018in}{1.425307in}}%
\pgfpathlineto{\pgfqpoint{1.962304in}{1.440040in}}%
\pgfpathlineto{\pgfqpoint{2.008490in}{1.454740in}}%
\pgfpathlineto{\pgfqpoint{2.039548in}{1.433736in}}%
\pgfpathlineto{\pgfqpoint{2.070702in}{1.412665in}}%
\pgfpathlineto{\pgfqpoint{2.024440in}{1.397873in}}%
\pgfpathlineto{\pgfqpoint{1.978077in}{1.383047in}}%
\pgfpathlineto{\pgfqpoint{1.946999in}{1.404210in}}%
\pgfpathlineto{\pgfqpoint{1.916018in}{1.425307in}}%
\pgfpathclose%
\pgfusepath{stroke,fill}%
\end{pgfscope}%
\begin{pgfscope}%
\pgfpathrectangle{\pgfqpoint{0.050000in}{0.050000in}}{\pgfqpoint{4.900000in}{4.900000in}}%
\pgfusepath{clip}%
\pgfsetbuttcap%
\pgfsetroundjoin%
\definecolor{currentfill}{rgb}{1.000000,1.000000,1.000000}%
\pgfsetfillcolor{currentfill}%
\pgfsetfillopacity{0.900000}%
\pgfsetlinewidth{0.200750pt}%
\definecolor{currentstroke}{rgb}{0.200000,0.200000,0.200000}%
\pgfsetstrokecolor{currentstroke}%
\pgfsetdash{}{0pt}%
\pgfpathmoveto{\pgfqpoint{4.234564in}{1.421241in}}%
\pgfpathlineto{\pgfqpoint{4.278332in}{1.435983in}}%
\pgfpathlineto{\pgfqpoint{4.322004in}{1.450693in}}%
\pgfpathlineto{\pgfqpoint{4.356682in}{1.429675in}}%
\pgfpathlineto{\pgfqpoint{4.391468in}{1.408592in}}%
\pgfpathlineto{\pgfqpoint{4.347735in}{1.393791in}}%
\pgfpathlineto{\pgfqpoint{4.303907in}{1.378956in}}%
\pgfpathlineto{\pgfqpoint{4.269181in}{1.400132in}}%
\pgfpathlineto{\pgfqpoint{4.234564in}{1.421241in}}%
\pgfpathclose%
\pgfusepath{stroke,fill}%
\end{pgfscope}%
\begin{pgfscope}%
\pgfpathrectangle{\pgfqpoint{0.050000in}{0.050000in}}{\pgfqpoint{4.900000in}{4.900000in}}%
\pgfusepath{clip}%
\pgfsetbuttcap%
\pgfsetroundjoin%
\definecolor{currentfill}{rgb}{1.000000,1.000000,1.000000}%
\pgfsetfillcolor{currentfill}%
\pgfsetfillopacity{0.900000}%
\pgfsetlinewidth{0.200750pt}%
\definecolor{currentstroke}{rgb}{0.200000,0.200000,0.200000}%
\pgfsetstrokecolor{currentstroke}%
\pgfsetdash{}{0pt}%
\pgfpathmoveto{\pgfqpoint{1.514138in}{1.421066in}}%
\pgfpathlineto{\pgfqpoint{1.560878in}{1.435809in}}%
\pgfpathlineto{\pgfqpoint{1.607516in}{1.450519in}}%
\pgfpathlineto{\pgfqpoint{1.637958in}{1.429501in}}%
\pgfpathlineto{\pgfqpoint{1.668494in}{1.408417in}}%
\pgfpathlineto{\pgfqpoint{1.621777in}{1.393615in}}%
\pgfpathlineto{\pgfqpoint{1.574958in}{1.378780in}}%
\pgfpathlineto{\pgfqpoint{1.544500in}{1.399956in}}%
\pgfpathlineto{\pgfqpoint{1.514138in}{1.421066in}}%
\pgfpathclose%
\pgfusepath{stroke,fill}%
\end{pgfscope}%
\begin{pgfscope}%
\pgfpathrectangle{\pgfqpoint{0.050000in}{0.050000in}}{\pgfqpoint{4.900000in}{4.900000in}}%
\pgfusepath{clip}%
\pgfsetbuttcap%
\pgfsetroundjoin%
\definecolor{currentfill}{rgb}{0.630050,0.614794,0.778793}%
\pgfsetfillcolor{currentfill}%
\pgfsetfillopacity{0.900000}%
\pgfsetlinewidth{0.200750pt}%
\definecolor{currentstroke}{rgb}{0.200000,0.200000,0.200000}%
\pgfsetstrokecolor{currentstroke}%
\pgfsetdash{}{0pt}%
\pgfpathmoveto{\pgfqpoint{3.394507in}{1.949695in}}%
\pgfpathlineto{\pgfqpoint{3.439718in}{1.998263in}}%
\pgfpathlineto{\pgfqpoint{3.484634in}{2.031357in}}%
\pgfpathlineto{\pgfqpoint{3.517773in}{1.996045in}}%
\pgfpathlineto{\pgfqpoint{3.550998in}{1.960662in}}%
\pgfpathlineto{\pgfqpoint{3.506132in}{1.934379in}}%
\pgfpathlineto{\pgfqpoint{3.460943in}{1.892691in}}%
\pgfpathlineto{\pgfqpoint{3.427680in}{1.921375in}}%
\pgfpathlineto{\pgfqpoint{3.394507in}{1.949695in}}%
\pgfpathclose%
\pgfusepath{stroke,fill}%
\end{pgfscope}%
\begin{pgfscope}%
\pgfpathrectangle{\pgfqpoint{0.050000in}{0.050000in}}{\pgfqpoint{4.900000in}{4.900000in}}%
\pgfusepath{clip}%
\pgfsetbuttcap%
\pgfsetroundjoin%
\definecolor{currentfill}{rgb}{0.904529,0.900592,0.942914}%
\pgfsetfillcolor{currentfill}%
\pgfsetfillopacity{0.900000}%
\pgfsetlinewidth{0.200750pt}%
\definecolor{currentstroke}{rgb}{0.200000,0.200000,0.200000}%
\pgfsetstrokecolor{currentstroke}%
\pgfsetdash{}{0pt}%
\pgfpathmoveto{\pgfqpoint{3.926442in}{1.643932in}}%
\pgfpathlineto{\pgfqpoint{3.969045in}{1.595753in}}%
\pgfpathlineto{\pgfqpoint{4.011481in}{1.548646in}}%
\pgfpathlineto{\pgfqpoint{4.045281in}{1.513701in}}%
\pgfpathlineto{\pgfqpoint{4.079184in}{1.479188in}}%
\pgfpathlineto{\pgfqpoint{4.036678in}{1.522212in}}%
\pgfpathlineto{\pgfqpoint{3.994050in}{1.567733in}}%
\pgfpathlineto{\pgfqpoint{3.960197in}{1.605614in}}%
\pgfpathlineto{\pgfqpoint{3.926442in}{1.643932in}}%
\pgfpathclose%
\pgfusepath{stroke,fill}%
\end{pgfscope}%
\begin{pgfscope}%
\pgfpathrectangle{\pgfqpoint{0.050000in}{0.050000in}}{\pgfqpoint{4.900000in}{4.900000in}}%
\pgfusepath{clip}%
\pgfsetbuttcap%
\pgfsetroundjoin%
\definecolor{currentfill}{rgb}{1.000000,1.000000,1.000000}%
\pgfsetfillcolor{currentfill}%
\pgfsetfillopacity{0.900000}%
\pgfsetlinewidth{0.200750pt}%
\definecolor{currentstroke}{rgb}{0.200000,0.200000,0.200000}%
\pgfsetstrokecolor{currentstroke}%
\pgfsetdash{}{0pt}%
\pgfpathmoveto{\pgfqpoint{1.112005in}{1.416823in}}%
\pgfpathlineto{\pgfqpoint{1.159199in}{1.431575in}}%
\pgfpathlineto{\pgfqpoint{1.206290in}{1.446294in}}%
\pgfpathlineto{\pgfqpoint{1.236115in}{1.425263in}}%
\pgfpathlineto{\pgfqpoint{1.266033in}{1.404166in}}%
\pgfpathlineto{\pgfqpoint{1.218860in}{1.389355in}}%
\pgfpathlineto{\pgfqpoint{1.171584in}{1.374511in}}%
\pgfpathlineto{\pgfqpoint{1.141748in}{1.395700in}}%
\pgfpathlineto{\pgfqpoint{1.112005in}{1.416823in}}%
\pgfpathclose%
\pgfusepath{stroke,fill}%
\end{pgfscope}%
\begin{pgfscope}%
\pgfpathrectangle{\pgfqpoint{0.050000in}{0.050000in}}{\pgfqpoint{4.900000in}{4.900000in}}%
\pgfusepath{clip}%
\pgfsetbuttcap%
\pgfsetroundjoin%
\definecolor{currentfill}{rgb}{1.000000,1.000000,1.000000}%
\pgfsetfillcolor{currentfill}%
\pgfsetfillopacity{0.900000}%
\pgfsetlinewidth{0.200750pt}%
\definecolor{currentstroke}{rgb}{0.200000,0.200000,0.200000}%
\pgfsetstrokecolor{currentstroke}%
\pgfsetdash{}{0pt}%
\pgfpathmoveto{\pgfqpoint{2.472657in}{1.417228in}}%
\pgfpathlineto{\pgfqpoint{2.518364in}{1.432381in}}%
\pgfpathlineto{\pgfqpoint{2.563972in}{1.447918in}}%
\pgfpathlineto{\pgfqpoint{2.595917in}{1.427142in}}%
\pgfpathlineto{\pgfqpoint{2.627962in}{1.406326in}}%
\pgfpathlineto{\pgfqpoint{2.582280in}{1.390453in}}%
\pgfpathlineto{\pgfqpoint{2.536500in}{1.375064in}}%
\pgfpathlineto{\pgfqpoint{2.504528in}{1.396174in}}%
\pgfpathlineto{\pgfqpoint{2.472657in}{1.417228in}}%
\pgfpathclose%
\pgfusepath{stroke,fill}%
\end{pgfscope}%
\begin{pgfscope}%
\pgfpathrectangle{\pgfqpoint{0.050000in}{0.050000in}}{\pgfqpoint{4.900000in}{4.900000in}}%
\pgfusepath{clip}%
\pgfsetbuttcap%
\pgfsetroundjoin%
\definecolor{currentfill}{rgb}{0.647951,0.633433,0.789496}%
\pgfsetfillcolor{currentfill}%
\pgfsetfillopacity{0.900000}%
\pgfsetlinewidth{0.200750pt}%
\definecolor{currentstroke}{rgb}{0.200000,0.200000,0.200000}%
\pgfsetstrokecolor{currentstroke}%
\pgfsetdash{}{0pt}%
\pgfpathmoveto{\pgfqpoint{3.550998in}{1.960662in}}%
\pgfpathlineto{\pgfqpoint{3.595455in}{1.968925in}}%
\pgfpathlineto{\pgfqpoint{3.639441in}{1.958078in}}%
\pgfpathlineto{\pgfqpoint{3.672680in}{1.917826in}}%
\pgfpathlineto{\pgfqpoint{3.706006in}{1.877891in}}%
\pgfpathlineto{\pgfqpoint{3.662095in}{1.892685in}}%
\pgfpathlineto{\pgfqpoint{3.617709in}{1.889763in}}%
\pgfpathlineto{\pgfqpoint{3.584309in}{1.925228in}}%
\pgfpathlineto{\pgfqpoint{3.550998in}{1.960662in}}%
\pgfpathclose%
\pgfusepath{stroke,fill}%
\end{pgfscope}%
\begin{pgfscope}%
\pgfpathrectangle{\pgfqpoint{0.050000in}{0.050000in}}{\pgfqpoint{4.900000in}{4.900000in}}%
\pgfusepath{clip}%
\pgfsetbuttcap%
\pgfsetroundjoin%
\definecolor{currentfill}{rgb}{1.000000,1.000000,1.000000}%
\pgfsetfillcolor{currentfill}%
\pgfsetfillopacity{0.900000}%
\pgfsetlinewidth{0.200750pt}%
\definecolor{currentstroke}{rgb}{0.200000,0.200000,0.200000}%
\pgfsetstrokecolor{currentstroke}%
\pgfsetdash{}{0pt}%
\pgfpathmoveto{\pgfqpoint{2.070702in}{1.412665in}}%
\pgfpathlineto{\pgfqpoint{2.116863in}{1.427426in}}%
\pgfpathlineto{\pgfqpoint{2.162923in}{1.442156in}}%
\pgfpathlineto{\pgfqpoint{2.194251in}{1.421112in}}%
\pgfpathlineto{\pgfqpoint{2.225677in}{1.400003in}}%
\pgfpathlineto{\pgfqpoint{2.179542in}{1.385181in}}%
\pgfpathlineto{\pgfqpoint{2.133305in}{1.370327in}}%
\pgfpathlineto{\pgfqpoint{2.101955in}{1.391529in}}%
\pgfpathlineto{\pgfqpoint{2.070702in}{1.412665in}}%
\pgfpathclose%
\pgfusepath{stroke,fill}%
\end{pgfscope}%
\begin{pgfscope}%
\pgfpathrectangle{\pgfqpoint{0.050000in}{0.050000in}}{\pgfqpoint{4.900000in}{4.900000in}}%
\pgfusepath{clip}%
\pgfsetbuttcap%
\pgfsetroundjoin%
\definecolor{currentfill}{rgb}{1.000000,1.000000,1.000000}%
\pgfsetfillcolor{currentfill}%
\pgfsetfillopacity{0.900000}%
\pgfsetlinewidth{0.200750pt}%
\definecolor{currentstroke}{rgb}{0.200000,0.200000,0.200000}%
\pgfsetstrokecolor{currentstroke}%
\pgfsetdash{}{0pt}%
\pgfpathmoveto{\pgfqpoint{4.391468in}{1.408592in}}%
\pgfpathlineto{\pgfqpoint{4.435106in}{1.423362in}}%
\pgfpathlineto{\pgfqpoint{4.470031in}{1.402259in}}%
\pgfpathlineto{\pgfqpoint{4.426363in}{1.387443in}}%
\pgfpathlineto{\pgfqpoint{4.391468in}{1.408592in}}%
\pgfpathclose%
\pgfusepath{stroke,fill}%
\end{pgfscope}%
\begin{pgfscope}%
\pgfpathrectangle{\pgfqpoint{0.050000in}{0.050000in}}{\pgfqpoint{4.900000in}{4.900000in}}%
\pgfusepath{clip}%
\pgfsetbuttcap%
\pgfsetroundjoin%
\definecolor{currentfill}{rgb}{0.964198,0.962722,0.978593}%
\pgfsetfillcolor{currentfill}%
\pgfsetfillopacity{0.900000}%
\pgfsetlinewidth{0.200750pt}%
\definecolor{currentstroke}{rgb}{0.200000,0.200000,0.200000}%
\pgfsetstrokecolor{currentstroke}%
\pgfsetdash{}{0pt}%
\pgfpathmoveto{\pgfqpoint{2.874501in}{1.446628in}}%
\pgfpathlineto{\pgfqpoint{2.919867in}{1.475861in}}%
\pgfpathlineto{\pgfqpoint{2.965200in}{1.510850in}}%
\pgfpathlineto{\pgfqpoint{2.997818in}{1.492776in}}%
\pgfpathlineto{\pgfqpoint{3.030541in}{1.474597in}}%
\pgfpathlineto{\pgfqpoint{2.985102in}{1.438445in}}%
\pgfpathlineto{\pgfqpoint{2.939637in}{1.408018in}}%
\pgfpathlineto{\pgfqpoint{2.907017in}{1.427344in}}%
\pgfpathlineto{\pgfqpoint{2.874501in}{1.446628in}}%
\pgfpathclose%
\pgfusepath{stroke,fill}%
\end{pgfscope}%
\begin{pgfscope}%
\pgfpathrectangle{\pgfqpoint{0.050000in}{0.050000in}}{\pgfqpoint{4.900000in}{4.900000in}}%
\pgfusepath{clip}%
\pgfsetbuttcap%
\pgfsetroundjoin%
\definecolor{currentfill}{rgb}{1.000000,1.000000,1.000000}%
\pgfsetfillcolor{currentfill}%
\pgfsetfillopacity{0.900000}%
\pgfsetlinewidth{0.200750pt}%
\definecolor{currentstroke}{rgb}{0.200000,0.200000,0.200000}%
\pgfsetstrokecolor{currentstroke}%
\pgfsetdash{}{0pt}%
\pgfpathmoveto{\pgfqpoint{1.668494in}{1.408417in}}%
\pgfpathlineto{\pgfqpoint{1.715109in}{1.423187in}}%
\pgfpathlineto{\pgfqpoint{1.761623in}{1.437924in}}%
\pgfpathlineto{\pgfqpoint{1.792333in}{1.416867in}}%
\pgfpathlineto{\pgfqpoint{1.823141in}{1.395744in}}%
\pgfpathlineto{\pgfqpoint{1.776549in}{1.380914in}}%
\pgfpathlineto{\pgfqpoint{1.729855in}{1.366052in}}%
\pgfpathlineto{\pgfqpoint{1.699127in}{1.387268in}}%
\pgfpathlineto{\pgfqpoint{1.668494in}{1.408417in}}%
\pgfpathclose%
\pgfusepath{stroke,fill}%
\end{pgfscope}%
\begin{pgfscope}%
\pgfpathrectangle{\pgfqpoint{0.050000in}{0.050000in}}{\pgfqpoint{4.900000in}{4.900000in}}%
\pgfusepath{clip}%
\pgfsetbuttcap%
\pgfsetroundjoin%
\definecolor{currentfill}{rgb}{0.689719,0.676924,0.814471}%
\pgfsetfillcolor{currentfill}%
\pgfsetfillopacity{0.900000}%
\pgfsetlinewidth{0.200750pt}%
\definecolor{currentstroke}{rgb}{0.200000,0.200000,0.200000}%
\pgfsetstrokecolor{currentstroke}%
\pgfsetdash{}{0pt}%
\pgfpathmoveto{\pgfqpoint{3.303573in}{1.823398in}}%
\pgfpathlineto{\pgfqpoint{3.349096in}{1.889883in}}%
\pgfpathlineto{\pgfqpoint{3.394507in}{1.949695in}}%
\pgfpathlineto{\pgfqpoint{3.427680in}{1.921375in}}%
\pgfpathlineto{\pgfqpoint{3.460943in}{1.892691in}}%
\pgfpathlineto{\pgfqpoint{3.415521in}{1.839243in}}%
\pgfpathlineto{\pgfqpoint{3.369954in}{1.778169in}}%
\pgfpathlineto{\pgfqpoint{3.336716in}{1.801014in}}%
\pgfpathlineto{\pgfqpoint{3.303573in}{1.823398in}}%
\pgfpathclose%
\pgfusepath{stroke,fill}%
\end{pgfscope}%
\begin{pgfscope}%
\pgfpathrectangle{\pgfqpoint{0.050000in}{0.050000in}}{\pgfqpoint{4.900000in}{4.900000in}}%
\pgfusepath{clip}%
\pgfsetbuttcap%
\pgfsetroundjoin%
\definecolor{currentfill}{rgb}{1.000000,1.000000,1.000000}%
\pgfsetfillcolor{currentfill}%
\pgfsetfillopacity{0.900000}%
\pgfsetlinewidth{0.200750pt}%
\definecolor{currentstroke}{rgb}{0.200000,0.200000,0.200000}%
\pgfsetstrokecolor{currentstroke}%
\pgfsetdash{}{0pt}%
\pgfpathmoveto{\pgfqpoint{1.266033in}{1.404166in}}%
\pgfpathlineto{\pgfqpoint{1.313103in}{1.418945in}}%
\pgfpathlineto{\pgfqpoint{1.360070in}{1.433692in}}%
\pgfpathlineto{\pgfqpoint{1.390163in}{1.412622in}}%
\pgfpathlineto{\pgfqpoint{1.420350in}{1.391485in}}%
\pgfpathlineto{\pgfqpoint{1.373302in}{1.376646in}}%
\pgfpathlineto{\pgfqpoint{1.326151in}{1.361774in}}%
\pgfpathlineto{\pgfqpoint{1.296045in}{1.383003in}}%
\pgfpathlineto{\pgfqpoint{1.266033in}{1.404166in}}%
\pgfpathclose%
\pgfusepath{stroke,fill}%
\end{pgfscope}%
\begin{pgfscope}%
\pgfpathrectangle{\pgfqpoint{0.050000in}{0.050000in}}{\pgfqpoint{4.900000in}{4.900000in}}%
\pgfusepath{clip}%
\pgfsetbuttcap%
\pgfsetroundjoin%
\definecolor{currentfill}{rgb}{1.000000,1.000000,1.000000}%
\pgfsetfillcolor{currentfill}%
\pgfsetfillopacity{0.900000}%
\pgfsetlinewidth{0.200750pt}%
\definecolor{currentstroke}{rgb}{0.200000,0.200000,0.200000}%
\pgfsetstrokecolor{currentstroke}%
\pgfsetdash{}{0pt}%
\pgfpathmoveto{\pgfqpoint{2.627962in}{1.406326in}}%
\pgfpathlineto{\pgfqpoint{2.673548in}{1.423059in}}%
\pgfpathlineto{\pgfqpoint{2.719041in}{1.441208in}}%
\pgfpathlineto{\pgfqpoint{2.751264in}{1.420960in}}%
\pgfpathlineto{\pgfqpoint{2.783590in}{1.400686in}}%
\pgfpathlineto{\pgfqpoint{2.738018in}{1.381826in}}%
\pgfpathlineto{\pgfqpoint{2.692356in}{1.364573in}}%
\pgfpathlineto{\pgfqpoint{2.660109in}{1.385470in}}%
\pgfpathlineto{\pgfqpoint{2.627962in}{1.406326in}}%
\pgfpathclose%
\pgfusepath{stroke,fill}%
\end{pgfscope}%
\begin{pgfscope}%
\pgfpathrectangle{\pgfqpoint{0.050000in}{0.050000in}}{\pgfqpoint{4.900000in}{4.900000in}}%
\pgfusepath{clip}%
\pgfsetbuttcap%
\pgfsetroundjoin%
\definecolor{currentfill}{rgb}{0.773256,0.763906,0.864421}%
\pgfsetfillcolor{currentfill}%
\pgfsetfillopacity{0.900000}%
\pgfsetlinewidth{0.200750pt}%
\definecolor{currentstroke}{rgb}{0.200000,0.200000,0.200000}%
\pgfsetstrokecolor{currentstroke}%
\pgfsetdash{}{0pt}%
\pgfpathmoveto{\pgfqpoint{3.212445in}{1.687260in}}%
\pgfpathlineto{\pgfqpoint{3.258006in}{1.754611in}}%
\pgfpathlineto{\pgfqpoint{3.303573in}{1.823398in}}%
\pgfpathlineto{\pgfqpoint{3.336716in}{1.801014in}}%
\pgfpathlineto{\pgfqpoint{3.369954in}{1.778169in}}%
\pgfpathlineto{\pgfqpoint{3.324316in}{1.713577in}}%
\pgfpathlineto{\pgfqpoint{3.278662in}{1.649124in}}%
\pgfpathlineto{\pgfqpoint{3.245503in}{1.668388in}}%
\pgfpathlineto{\pgfqpoint{3.212445in}{1.687260in}}%
\pgfpathclose%
\pgfusepath{stroke,fill}%
\end{pgfscope}%
\begin{pgfscope}%
\pgfpathrectangle{\pgfqpoint{0.050000in}{0.050000in}}{\pgfqpoint{4.900000in}{4.900000in}}%
\pgfusepath{clip}%
\pgfsetbuttcap%
\pgfsetroundjoin%
\definecolor{currentfill}{rgb}{1.000000,1.000000,1.000000}%
\pgfsetfillcolor{currentfill}%
\pgfsetfillopacity{0.900000}%
\pgfsetlinewidth{0.200750pt}%
\definecolor{currentstroke}{rgb}{0.200000,0.200000,0.200000}%
\pgfsetstrokecolor{currentstroke}%
\pgfsetdash{}{0pt}%
\pgfpathmoveto{\pgfqpoint{2.225677in}{1.400003in}}%
\pgfpathlineto{\pgfqpoint{2.271711in}{1.414798in}}%
\pgfpathlineto{\pgfqpoint{2.317645in}{1.429574in}}%
\pgfpathlineto{\pgfqpoint{2.349244in}{1.408500in}}%
\pgfpathlineto{\pgfqpoint{2.380942in}{1.387362in}}%
\pgfpathlineto{\pgfqpoint{2.334934in}{1.372481in}}%
\pgfpathlineto{\pgfqpoint{2.288825in}{1.357588in}}%
\pgfpathlineto{\pgfqpoint{2.257201in}{1.378829in}}%
\pgfpathlineto{\pgfqpoint{2.225677in}{1.400003in}}%
\pgfpathclose%
\pgfusepath{stroke,fill}%
\end{pgfscope}%
\begin{pgfscope}%
\pgfpathrectangle{\pgfqpoint{0.050000in}{0.050000in}}{\pgfqpoint{4.900000in}{4.900000in}}%
\pgfusepath{clip}%
\pgfsetbuttcap%
\pgfsetroundjoin%
\definecolor{currentfill}{rgb}{1.000000,1.000000,1.000000}%
\pgfsetfillcolor{currentfill}%
\pgfsetfillopacity{0.900000}%
\pgfsetlinewidth{0.200750pt}%
\definecolor{currentstroke}{rgb}{0.200000,0.200000,0.200000}%
\pgfsetstrokecolor{currentstroke}%
\pgfsetdash{}{0pt}%
\pgfpathmoveto{\pgfqpoint{1.823141in}{1.395744in}}%
\pgfpathlineto{\pgfqpoint{1.869630in}{1.410542in}}%
\pgfpathlineto{\pgfqpoint{1.916018in}{1.425307in}}%
\pgfpathlineto{\pgfqpoint{1.946999in}{1.404210in}}%
\pgfpathlineto{\pgfqpoint{1.978077in}{1.383047in}}%
\pgfpathlineto{\pgfqpoint{1.931612in}{1.368190in}}%
\pgfpathlineto{\pgfqpoint{1.885045in}{1.353299in}}%
\pgfpathlineto{\pgfqpoint{1.854044in}{1.374555in}}%
\pgfpathlineto{\pgfqpoint{1.823141in}{1.395744in}}%
\pgfpathclose%
\pgfusepath{stroke,fill}%
\end{pgfscope}%
\begin{pgfscope}%
\pgfpathrectangle{\pgfqpoint{0.050000in}{0.050000in}}{\pgfqpoint{4.900000in}{4.900000in}}%
\pgfusepath{clip}%
\pgfsetbuttcap%
\pgfsetroundjoin%
\definecolor{currentfill}{rgb}{0.856794,0.850888,0.914371}%
\pgfsetfillcolor{currentfill}%
\pgfsetfillopacity{0.900000}%
\pgfsetlinewidth{0.200750pt}%
\definecolor{currentstroke}{rgb}{0.200000,0.200000,0.200000}%
\pgfsetstrokecolor{currentstroke}%
\pgfsetdash{}{0pt}%
\pgfpathmoveto{\pgfqpoint{3.121430in}{1.567261in}}%
\pgfpathlineto{\pgfqpoint{3.166917in}{1.624192in}}%
\pgfpathlineto{\pgfqpoint{3.212445in}{1.687260in}}%
\pgfpathlineto{\pgfqpoint{3.245503in}{1.668388in}}%
\pgfpathlineto{\pgfqpoint{3.278662in}{1.649124in}}%
\pgfpathlineto{\pgfqpoint{3.233026in}{1.587727in}}%
\pgfpathlineto{\pgfqpoint{3.187423in}{1.531423in}}%
\pgfpathlineto{\pgfqpoint{3.154374in}{1.549466in}}%
\pgfpathlineto{\pgfqpoint{3.121430in}{1.567261in}}%
\pgfpathclose%
\pgfusepath{stroke,fill}%
\end{pgfscope}%
\begin{pgfscope}%
\pgfpathrectangle{\pgfqpoint{0.050000in}{0.050000in}}{\pgfqpoint{4.900000in}{4.900000in}}%
\pgfusepath{clip}%
\pgfsetbuttcap%
\pgfsetroundjoin%
\definecolor{currentfill}{rgb}{0.773256,0.763906,0.864421}%
\pgfsetfillcolor{currentfill}%
\pgfsetfillopacity{0.900000}%
\pgfsetlinewidth{0.200750pt}%
\definecolor{currentstroke}{rgb}{0.200000,0.200000,0.200000}%
\pgfsetstrokecolor{currentstroke}%
\pgfsetdash{}{0pt}%
\pgfpathmoveto{\pgfqpoint{3.772929in}{1.798951in}}%
\pgfpathlineto{\pgfqpoint{3.816294in}{1.766152in}}%
\pgfpathlineto{\pgfqpoint{3.859224in}{1.721961in}}%
\pgfpathlineto{\pgfqpoint{3.892785in}{1.682707in}}%
\pgfpathlineto{\pgfqpoint{3.926442in}{1.643932in}}%
\pgfpathlineto{\pgfqpoint{3.883533in}{1.687359in}}%
\pgfpathlineto{\pgfqpoint{3.840218in}{1.721221in}}%
\pgfpathlineto{\pgfqpoint{3.806528in}{1.759937in}}%
\pgfpathlineto{\pgfqpoint{3.772929in}{1.798951in}}%
\pgfpathclose%
\pgfusepath{stroke,fill}%
\end{pgfscope}%
\begin{pgfscope}%
\pgfpathrectangle{\pgfqpoint{0.050000in}{0.050000in}}{\pgfqpoint{4.900000in}{4.900000in}}%
\pgfusepath{clip}%
\pgfsetbuttcap%
\pgfsetroundjoin%
\definecolor{currentfill}{rgb}{1.000000,1.000000,1.000000}%
\pgfsetfillcolor{currentfill}%
\pgfsetfillopacity{0.900000}%
\pgfsetlinewidth{0.200750pt}%
\definecolor{currentstroke}{rgb}{0.200000,0.200000,0.200000}%
\pgfsetstrokecolor{currentstroke}%
\pgfsetdash{}{0pt}%
\pgfpathmoveto{\pgfqpoint{1.420350in}{1.391485in}}%
\pgfpathlineto{\pgfqpoint{1.467295in}{1.406292in}}%
\pgfpathlineto{\pgfqpoint{1.514138in}{1.421066in}}%
\pgfpathlineto{\pgfqpoint{1.544500in}{1.399956in}}%
\pgfpathlineto{\pgfqpoint{1.574958in}{1.378780in}}%
\pgfpathlineto{\pgfqpoint{1.528035in}{1.363913in}}%
\pgfpathlineto{\pgfqpoint{1.481009in}{1.349013in}}%
\pgfpathlineto{\pgfqpoint{1.450632in}{1.370283in}}%
\pgfpathlineto{\pgfqpoint{1.420350in}{1.391485in}}%
\pgfpathclose%
\pgfusepath{stroke,fill}%
\end{pgfscope}%
\begin{pgfscope}%
\pgfpathrectangle{\pgfqpoint{0.050000in}{0.050000in}}{\pgfqpoint{4.900000in}{4.900000in}}%
\pgfusepath{clip}%
\pgfsetbuttcap%
\pgfsetroundjoin%
\definecolor{currentfill}{rgb}{1.000000,1.000000,1.000000}%
\pgfsetfillcolor{currentfill}%
\pgfsetfillopacity{0.900000}%
\pgfsetlinewidth{0.200750pt}%
\definecolor{currentstroke}{rgb}{0.200000,0.200000,0.200000}%
\pgfsetstrokecolor{currentstroke}%
\pgfsetdash{}{0pt}%
\pgfpathmoveto{\pgfqpoint{4.147300in}{1.411365in}}%
\pgfpathlineto{\pgfqpoint{4.190700in}{1.406467in}}%
\pgfpathlineto{\pgfqpoint{4.234564in}{1.421241in}}%
\pgfpathlineto{\pgfqpoint{4.269181in}{1.400132in}}%
\pgfpathlineto{\pgfqpoint{4.303907in}{1.378956in}}%
\pgfpathlineto{\pgfqpoint{4.259981in}{1.364090in}}%
\pgfpathlineto{\pgfqpoint{4.215960in}{1.349190in}}%
\pgfpathlineto{\pgfqpoint{4.181515in}{1.378012in}}%
\pgfpathlineto{\pgfqpoint{4.147300in}{1.411365in}}%
\pgfpathclose%
\pgfusepath{stroke,fill}%
\end{pgfscope}%
\begin{pgfscope}%
\pgfpathrectangle{\pgfqpoint{0.050000in}{0.050000in}}{\pgfqpoint{4.900000in}{4.900000in}}%
\pgfusepath{clip}%
\pgfsetbuttcap%
\pgfsetroundjoin%
\definecolor{currentfill}{rgb}{1.000000,1.000000,1.000000}%
\pgfsetfillcolor{currentfill}%
\pgfsetfillopacity{0.900000}%
\pgfsetlinewidth{0.200750pt}%
\definecolor{currentstroke}{rgb}{0.200000,0.200000,0.200000}%
\pgfsetstrokecolor{currentstroke}%
\pgfsetdash{}{0pt}%
\pgfpathmoveto{\pgfqpoint{4.303907in}{1.378956in}}%
\pgfpathlineto{\pgfqpoint{4.347735in}{1.393791in}}%
\pgfpathlineto{\pgfqpoint{4.391468in}{1.408592in}}%
\pgfpathlineto{\pgfqpoint{4.426363in}{1.387443in}}%
\pgfpathlineto{\pgfqpoint{4.382600in}{1.372595in}}%
\pgfpathlineto{\pgfqpoint{4.338741in}{1.357714in}}%
\pgfpathlineto{\pgfqpoint{4.303907in}{1.378956in}}%
\pgfpathclose%
\pgfusepath{stroke,fill}%
\end{pgfscope}%
\begin{pgfscope}%
\pgfpathrectangle{\pgfqpoint{0.050000in}{0.050000in}}{\pgfqpoint{4.900000in}{4.900000in}}%
\pgfusepath{clip}%
\pgfsetbuttcap%
\pgfsetroundjoin%
\definecolor{currentfill}{rgb}{1.000000,1.000000,1.000000}%
\pgfsetfillcolor{currentfill}%
\pgfsetfillopacity{0.900000}%
\pgfsetlinewidth{0.200750pt}%
\definecolor{currentstroke}{rgb}{0.200000,0.200000,0.200000}%
\pgfsetstrokecolor{currentstroke}%
\pgfsetdash{}{0pt}%
\pgfpathmoveto{\pgfqpoint{2.380942in}{1.387362in}}%
\pgfpathlineto{\pgfqpoint{2.426850in}{1.402258in}}%
\pgfpathlineto{\pgfqpoint{2.472657in}{1.417228in}}%
\pgfpathlineto{\pgfqpoint{2.504528in}{1.396174in}}%
\pgfpathlineto{\pgfqpoint{2.536500in}{1.375064in}}%
\pgfpathlineto{\pgfqpoint{2.490619in}{1.359924in}}%
\pgfpathlineto{\pgfqpoint{2.444637in}{1.344895in}}%
\pgfpathlineto{\pgfqpoint{2.412740in}{1.366160in}}%
\pgfpathlineto{\pgfqpoint{2.380942in}{1.387362in}}%
\pgfpathclose%
\pgfusepath{stroke,fill}%
\end{pgfscope}%
\begin{pgfscope}%
\pgfpathrectangle{\pgfqpoint{0.050000in}{0.050000in}}{\pgfqpoint{4.900000in}{4.900000in}}%
\pgfusepath{clip}%
\pgfsetbuttcap%
\pgfsetroundjoin%
\definecolor{currentfill}{rgb}{0.988066,0.987574,0.992864}%
\pgfsetfillcolor{currentfill}%
\pgfsetfillopacity{0.900000}%
\pgfsetlinewidth{0.200750pt}%
\definecolor{currentstroke}{rgb}{0.200000,0.200000,0.200000}%
\pgfsetstrokecolor{currentstroke}%
\pgfsetdash{}{0pt}%
\pgfpathmoveto{\pgfqpoint{2.783590in}{1.400686in}}%
\pgfpathlineto{\pgfqpoint{2.829080in}{1.421954in}}%
\pgfpathlineto{\pgfqpoint{2.874501in}{1.446628in}}%
\pgfpathlineto{\pgfqpoint{2.907017in}{1.427344in}}%
\pgfpathlineto{\pgfqpoint{2.939637in}{1.408018in}}%
\pgfpathlineto{\pgfqpoint{2.894125in}{1.382249in}}%
\pgfpathlineto{\pgfqpoint{2.848549in}{1.360056in}}%
\pgfpathlineto{\pgfqpoint{2.816018in}{1.380386in}}%
\pgfpathlineto{\pgfqpoint{2.783590in}{1.400686in}}%
\pgfpathclose%
\pgfusepath{stroke,fill}%
\end{pgfscope}%
\begin{pgfscope}%
\pgfpathrectangle{\pgfqpoint{0.050000in}{0.050000in}}{\pgfqpoint{4.900000in}{4.900000in}}%
\pgfusepath{clip}%
\pgfsetbuttcap%
\pgfsetroundjoin%
\definecolor{currentfill}{rgb}{0.916463,0.913018,0.950050}%
\pgfsetfillcolor{currentfill}%
\pgfsetfillopacity{0.900000}%
\pgfsetlinewidth{0.200750pt}%
\definecolor{currentstroke}{rgb}{0.200000,0.200000,0.200000}%
\pgfsetstrokecolor{currentstroke}%
\pgfsetdash{}{0pt}%
\pgfpathmoveto{\pgfqpoint{3.030541in}{1.474597in}}%
\pgfpathlineto{\pgfqpoint{3.075977in}{1.517365in}}%
\pgfpathlineto{\pgfqpoint{3.121430in}{1.567261in}}%
\pgfpathlineto{\pgfqpoint{3.154374in}{1.549466in}}%
\pgfpathlineto{\pgfqpoint{3.187423in}{1.531423in}}%
\pgfpathlineto{\pgfqpoint{3.141853in}{1.481358in}}%
\pgfpathlineto{\pgfqpoint{3.096304in}{1.437886in}}%
\pgfpathlineto{\pgfqpoint{3.063370in}{1.456302in}}%
\pgfpathlineto{\pgfqpoint{3.030541in}{1.474597in}}%
\pgfpathclose%
\pgfusepath{stroke,fill}%
\end{pgfscope}%
\begin{pgfscope}%
\pgfpathrectangle{\pgfqpoint{0.050000in}{0.050000in}}{\pgfqpoint{4.900000in}{4.900000in}}%
\pgfusepath{clip}%
\pgfsetbuttcap%
\pgfsetroundjoin%
\definecolor{currentfill}{rgb}{0.647951,0.633433,0.789496}%
\pgfsetfillcolor{currentfill}%
\pgfsetfillopacity{0.900000}%
\pgfsetlinewidth{0.200750pt}%
\definecolor{currentstroke}{rgb}{0.200000,0.200000,0.200000}%
\pgfsetstrokecolor{currentstroke}%
\pgfsetdash{}{0pt}%
\pgfpathmoveto{\pgfqpoint{3.460943in}{1.892691in}}%
\pgfpathlineto{\pgfqpoint{3.506132in}{1.934379in}}%
\pgfpathlineto{\pgfqpoint{3.550998in}{1.960662in}}%
\pgfpathlineto{\pgfqpoint{3.584309in}{1.925228in}}%
\pgfpathlineto{\pgfqpoint{3.617709in}{1.889763in}}%
\pgfpathlineto{\pgfqpoint{3.572899in}{1.869643in}}%
\pgfpathlineto{\pgfqpoint{3.527739in}{1.834339in}}%
\pgfpathlineto{\pgfqpoint{3.494296in}{1.863670in}}%
\pgfpathlineto{\pgfqpoint{3.460943in}{1.892691in}}%
\pgfpathclose%
\pgfusepath{stroke,fill}%
\end{pgfscope}%
\begin{pgfscope}%
\pgfpathrectangle{\pgfqpoint{0.050000in}{0.050000in}}{\pgfqpoint{4.900000in}{4.900000in}}%
\pgfusepath{clip}%
\pgfsetbuttcap%
\pgfsetroundjoin%
\definecolor{currentfill}{rgb}{1.000000,1.000000,1.000000}%
\pgfsetfillcolor{currentfill}%
\pgfsetfillopacity{0.900000}%
\pgfsetlinewidth{0.200750pt}%
\definecolor{currentstroke}{rgb}{0.200000,0.200000,0.200000}%
\pgfsetstrokecolor{currentstroke}%
\pgfsetdash{}{0pt}%
\pgfpathmoveto{\pgfqpoint{1.978077in}{1.383047in}}%
\pgfpathlineto{\pgfqpoint{2.024440in}{1.397873in}}%
\pgfpathlineto{\pgfqpoint{2.070702in}{1.412665in}}%
\pgfpathlineto{\pgfqpoint{2.101955in}{1.391529in}}%
\pgfpathlineto{\pgfqpoint{2.133305in}{1.370327in}}%
\pgfpathlineto{\pgfqpoint{2.086967in}{1.355441in}}%
\pgfpathlineto{\pgfqpoint{2.040526in}{1.340522in}}%
\pgfpathlineto{\pgfqpoint{2.009253in}{1.361818in}}%
\pgfpathlineto{\pgfqpoint{1.978077in}{1.383047in}}%
\pgfpathclose%
\pgfusepath{stroke,fill}%
\end{pgfscope}%
\begin{pgfscope}%
\pgfpathrectangle{\pgfqpoint{0.050000in}{0.050000in}}{\pgfqpoint{4.900000in}{4.900000in}}%
\pgfusepath{clip}%
\pgfsetbuttcap%
\pgfsetroundjoin%
\definecolor{currentfill}{rgb}{0.928397,0.925444,0.957186}%
\pgfsetfillcolor{currentfill}%
\pgfsetfillopacity{0.900000}%
\pgfsetlinewidth{0.200750pt}%
\definecolor{currentstroke}{rgb}{0.200000,0.200000,0.200000}%
\pgfsetstrokecolor{currentstroke}%
\pgfsetdash{}{0pt}%
\pgfpathmoveto{\pgfqpoint{3.994050in}{1.567733in}}%
\pgfpathlineto{\pgfqpoint{4.036678in}{1.522212in}}%
\pgfpathlineto{\pgfqpoint{4.079184in}{1.479188in}}%
\pgfpathlineto{\pgfqpoint{4.113190in}{1.445083in}}%
\pgfpathlineto{\pgfqpoint{4.147300in}{1.411365in}}%
\pgfpathlineto{\pgfqpoint{4.104717in}{1.450404in}}%
\pgfpathlineto{\pgfqpoint{4.062055in}{1.493210in}}%
\pgfpathlineto{\pgfqpoint{4.028002in}{1.530271in}}%
\pgfpathlineto{\pgfqpoint{3.994050in}{1.567733in}}%
\pgfpathclose%
\pgfusepath{stroke,fill}%
\end{pgfscope}%
\begin{pgfscope}%
\pgfpathrectangle{\pgfqpoint{0.050000in}{0.050000in}}{\pgfqpoint{4.900000in}{4.900000in}}%
\pgfusepath{clip}%
\pgfsetbuttcap%
\pgfsetroundjoin%
\definecolor{currentfill}{rgb}{0.677785,0.664498,0.807336}%
\pgfsetfillcolor{currentfill}%
\pgfsetfillopacity{0.900000}%
\pgfsetlinewidth{0.200750pt}%
\definecolor{currentstroke}{rgb}{0.200000,0.200000,0.200000}%
\pgfsetstrokecolor{currentstroke}%
\pgfsetdash{}{0pt}%
\pgfpathmoveto{\pgfqpoint{3.617709in}{1.889763in}}%
\pgfpathlineto{\pgfqpoint{3.662095in}{1.892685in}}%
\pgfpathlineto{\pgfqpoint{3.706006in}{1.877891in}}%
\pgfpathlineto{\pgfqpoint{3.739423in}{1.838267in}}%
\pgfpathlineto{\pgfqpoint{3.772929in}{1.798951in}}%
\pgfpathlineto{\pgfqpoint{3.729089in}{1.817044in}}%
\pgfpathlineto{\pgfqpoint{3.684773in}{1.818799in}}%
\pgfpathlineto{\pgfqpoint{3.651197in}{1.854282in}}%
\pgfpathlineto{\pgfqpoint{3.617709in}{1.889763in}}%
\pgfpathclose%
\pgfusepath{stroke,fill}%
\end{pgfscope}%
\begin{pgfscope}%
\pgfpathrectangle{\pgfqpoint{0.050000in}{0.050000in}}{\pgfqpoint{4.900000in}{4.900000in}}%
\pgfusepath{clip}%
\pgfsetbuttcap%
\pgfsetroundjoin%
\definecolor{currentfill}{rgb}{1.000000,1.000000,1.000000}%
\pgfsetfillcolor{currentfill}%
\pgfsetfillopacity{0.900000}%
\pgfsetlinewidth{0.200750pt}%
\definecolor{currentstroke}{rgb}{0.200000,0.200000,0.200000}%
\pgfsetstrokecolor{currentstroke}%
\pgfsetdash{}{0pt}%
\pgfpathmoveto{\pgfqpoint{1.574958in}{1.378780in}}%
\pgfpathlineto{\pgfqpoint{1.621777in}{1.393615in}}%
\pgfpathlineto{\pgfqpoint{1.668494in}{1.408417in}}%
\pgfpathlineto{\pgfqpoint{1.699127in}{1.387268in}}%
\pgfpathlineto{\pgfqpoint{1.729855in}{1.366052in}}%
\pgfpathlineto{\pgfqpoint{1.683059in}{1.351157in}}%
\pgfpathlineto{\pgfqpoint{1.636160in}{1.336229in}}%
\pgfpathlineto{\pgfqpoint{1.605511in}{1.357538in}}%
\pgfpathlineto{\pgfqpoint{1.574958in}{1.378780in}}%
\pgfpathclose%
\pgfusepath{stroke,fill}%
\end{pgfscope}%
\begin{pgfscope}%
\pgfpathrectangle{\pgfqpoint{0.050000in}{0.050000in}}{\pgfqpoint{4.900000in}{4.900000in}}%
\pgfusepath{clip}%
\pgfsetbuttcap%
\pgfsetroundjoin%
\definecolor{currentfill}{rgb}{1.000000,1.000000,1.000000}%
\pgfsetfillcolor{currentfill}%
\pgfsetfillopacity{0.900000}%
\pgfsetlinewidth{0.200750pt}%
\definecolor{currentstroke}{rgb}{0.200000,0.200000,0.200000}%
\pgfsetstrokecolor{currentstroke}%
\pgfsetdash{}{0pt}%
\pgfpathmoveto{\pgfqpoint{1.171584in}{1.374511in}}%
\pgfpathlineto{\pgfqpoint{1.218860in}{1.389355in}}%
\pgfpathlineto{\pgfqpoint{1.266033in}{1.404166in}}%
\pgfpathlineto{\pgfqpoint{1.296045in}{1.383003in}}%
\pgfpathlineto{\pgfqpoint{1.326151in}{1.361774in}}%
\pgfpathlineto{\pgfqpoint{1.278896in}{1.346869in}}%
\pgfpathlineto{\pgfqpoint{1.231537in}{1.331932in}}%
\pgfpathlineto{\pgfqpoint{1.201513in}{1.353255in}}%
\pgfpathlineto{\pgfqpoint{1.171584in}{1.374511in}}%
\pgfpathclose%
\pgfusepath{stroke,fill}%
\end{pgfscope}%
\begin{pgfscope}%
\pgfpathrectangle{\pgfqpoint{0.050000in}{0.050000in}}{\pgfqpoint{4.900000in}{4.900000in}}%
\pgfusepath{clip}%
\pgfsetbuttcap%
\pgfsetroundjoin%
\definecolor{currentfill}{rgb}{1.000000,1.000000,1.000000}%
\pgfsetfillcolor{currentfill}%
\pgfsetfillopacity{0.900000}%
\pgfsetlinewidth{0.200750pt}%
\definecolor{currentstroke}{rgb}{0.200000,0.200000,0.200000}%
\pgfsetstrokecolor{currentstroke}%
\pgfsetdash{}{0pt}%
\pgfpathmoveto{\pgfqpoint{2.536500in}{1.375064in}}%
\pgfpathlineto{\pgfqpoint{2.582280in}{1.390453in}}%
\pgfpathlineto{\pgfqpoint{2.627962in}{1.406326in}}%
\pgfpathlineto{\pgfqpoint{2.660109in}{1.385470in}}%
\pgfpathlineto{\pgfqpoint{2.692356in}{1.364573in}}%
\pgfpathlineto{\pgfqpoint{2.646599in}{1.348329in}}%
\pgfpathlineto{\pgfqpoint{2.600745in}{1.332677in}}%
\pgfpathlineto{\pgfqpoint{2.568572in}{1.353898in}}%
\pgfpathlineto{\pgfqpoint{2.536500in}{1.375064in}}%
\pgfpathclose%
\pgfusepath{stroke,fill}%
\end{pgfscope}%
\begin{pgfscope}%
\pgfpathrectangle{\pgfqpoint{0.050000in}{0.050000in}}{\pgfqpoint{4.900000in}{4.900000in}}%
\pgfusepath{clip}%
\pgfsetbuttcap%
\pgfsetroundjoin%
\definecolor{currentfill}{rgb}{1.000000,1.000000,1.000000}%
\pgfsetfillcolor{currentfill}%
\pgfsetfillopacity{0.900000}%
\pgfsetlinewidth{0.200750pt}%
\definecolor{currentstroke}{rgb}{0.200000,0.200000,0.200000}%
\pgfsetstrokecolor{currentstroke}%
\pgfsetdash{}{0pt}%
\pgfpathmoveto{\pgfqpoint{2.133305in}{1.370327in}}%
\pgfpathlineto{\pgfqpoint{2.179542in}{1.385181in}}%
\pgfpathlineto{\pgfqpoint{2.225677in}{1.400003in}}%
\pgfpathlineto{\pgfqpoint{2.257201in}{1.378829in}}%
\pgfpathlineto{\pgfqpoint{2.288825in}{1.357588in}}%
\pgfpathlineto{\pgfqpoint{2.242614in}{1.342670in}}%
\pgfpathlineto{\pgfqpoint{2.196301in}{1.327722in}}%
\pgfpathlineto{\pgfqpoint{2.164754in}{1.349058in}}%
\pgfpathlineto{\pgfqpoint{2.133305in}{1.370327in}}%
\pgfpathclose%
\pgfusepath{stroke,fill}%
\end{pgfscope}%
\begin{pgfscope}%
\pgfpathrectangle{\pgfqpoint{0.050000in}{0.050000in}}{\pgfqpoint{4.900000in}{4.900000in}}%
\pgfusepath{clip}%
\pgfsetbuttcap%
\pgfsetroundjoin%
\definecolor{currentfill}{rgb}{0.695686,0.683137,0.818039}%
\pgfsetfillcolor{currentfill}%
\pgfsetfillopacity{0.900000}%
\pgfsetlinewidth{0.200750pt}%
\definecolor{currentstroke}{rgb}{0.200000,0.200000,0.200000}%
\pgfsetstrokecolor{currentstroke}%
\pgfsetdash{}{0pt}%
\pgfpathmoveto{\pgfqpoint{3.369954in}{1.778169in}}%
\pgfpathlineto{\pgfqpoint{3.415521in}{1.839243in}}%
\pgfpathlineto{\pgfqpoint{3.460943in}{1.892691in}}%
\pgfpathlineto{\pgfqpoint{3.494296in}{1.863670in}}%
\pgfpathlineto{\pgfqpoint{3.527739in}{1.834339in}}%
\pgfpathlineto{\pgfqpoint{3.482318in}{1.786948in}}%
\pgfpathlineto{\pgfqpoint{3.436720in}{1.731168in}}%
\pgfpathlineto{\pgfqpoint{3.403289in}{1.754881in}}%
\pgfpathlineto{\pgfqpoint{3.369954in}{1.778169in}}%
\pgfpathclose%
\pgfusepath{stroke,fill}%
\end{pgfscope}%
\begin{pgfscope}%
\pgfpathrectangle{\pgfqpoint{0.050000in}{0.050000in}}{\pgfqpoint{4.900000in}{4.900000in}}%
\pgfusepath{clip}%
\pgfsetbuttcap%
\pgfsetroundjoin%
\definecolor{currentfill}{rgb}{0.964198,0.962722,0.978593}%
\pgfsetfillcolor{currentfill}%
\pgfsetfillopacity{0.900000}%
\pgfsetlinewidth{0.200750pt}%
\definecolor{currentstroke}{rgb}{0.200000,0.200000,0.200000}%
\pgfsetstrokecolor{currentstroke}%
\pgfsetdash{}{0pt}%
\pgfpathmoveto{\pgfqpoint{2.939637in}{1.408018in}}%
\pgfpathlineto{\pgfqpoint{2.985102in}{1.438445in}}%
\pgfpathlineto{\pgfqpoint{3.030541in}{1.474597in}}%
\pgfpathlineto{\pgfqpoint{3.063370in}{1.456302in}}%
\pgfpathlineto{\pgfqpoint{3.096304in}{1.437886in}}%
\pgfpathlineto{\pgfqpoint{3.050757in}{1.400739in}}%
\pgfpathlineto{\pgfqpoint{3.005191in}{1.369222in}}%
\pgfpathlineto{\pgfqpoint{2.972361in}{1.388646in}}%
\pgfpathlineto{\pgfqpoint{2.939637in}{1.408018in}}%
\pgfpathclose%
\pgfusepath{stroke,fill}%
\end{pgfscope}%
\begin{pgfscope}%
\pgfpathrectangle{\pgfqpoint{0.050000in}{0.050000in}}{\pgfqpoint{4.900000in}{4.900000in}}%
\pgfusepath{clip}%
\pgfsetbuttcap%
\pgfsetroundjoin%
\definecolor{currentfill}{rgb}{1.000000,1.000000,1.000000}%
\pgfsetfillcolor{currentfill}%
\pgfsetfillopacity{0.900000}%
\pgfsetlinewidth{0.200750pt}%
\definecolor{currentstroke}{rgb}{0.200000,0.200000,0.200000}%
\pgfsetstrokecolor{currentstroke}%
\pgfsetdash{}{0pt}%
\pgfpathmoveto{\pgfqpoint{1.729855in}{1.366052in}}%
\pgfpathlineto{\pgfqpoint{1.776549in}{1.380914in}}%
\pgfpathlineto{\pgfqpoint{1.823141in}{1.395744in}}%
\pgfpathlineto{\pgfqpoint{1.854044in}{1.374555in}}%
\pgfpathlineto{\pgfqpoint{1.885045in}{1.353299in}}%
\pgfpathlineto{\pgfqpoint{1.838375in}{1.338376in}}%
\pgfpathlineto{\pgfqpoint{1.791602in}{1.323420in}}%
\pgfpathlineto{\pgfqpoint{1.760680in}{1.344769in}}%
\pgfpathlineto{\pgfqpoint{1.729855in}{1.366052in}}%
\pgfpathclose%
\pgfusepath{stroke,fill}%
\end{pgfscope}%
\begin{pgfscope}%
\pgfpathrectangle{\pgfqpoint{0.050000in}{0.050000in}}{\pgfqpoint{4.900000in}{4.900000in}}%
\pgfusepath{clip}%
\pgfsetbuttcap%
\pgfsetroundjoin%
\definecolor{currentfill}{rgb}{1.000000,1.000000,1.000000}%
\pgfsetfillcolor{currentfill}%
\pgfsetfillopacity{0.900000}%
\pgfsetlinewidth{0.200750pt}%
\definecolor{currentstroke}{rgb}{0.200000,0.200000,0.200000}%
\pgfsetstrokecolor{currentstroke}%
\pgfsetdash{}{0pt}%
\pgfpathmoveto{\pgfqpoint{1.326151in}{1.361774in}}%
\pgfpathlineto{\pgfqpoint{1.373302in}{1.376646in}}%
\pgfpathlineto{\pgfqpoint{1.420350in}{1.391485in}}%
\pgfpathlineto{\pgfqpoint{1.450632in}{1.370283in}}%
\pgfpathlineto{\pgfqpoint{1.481009in}{1.349013in}}%
\pgfpathlineto{\pgfqpoint{1.433880in}{1.334081in}}%
\pgfpathlineto{\pgfqpoint{1.386647in}{1.319115in}}%
\pgfpathlineto{\pgfqpoint{1.356352in}{1.340478in}}%
\pgfpathlineto{\pgfqpoint{1.326151in}{1.361774in}}%
\pgfpathclose%
\pgfusepath{stroke,fill}%
\end{pgfscope}%
\begin{pgfscope}%
\pgfpathrectangle{\pgfqpoint{0.050000in}{0.050000in}}{\pgfqpoint{4.900000in}{4.900000in}}%
\pgfusepath{clip}%
\pgfsetbuttcap%
\pgfsetroundjoin%
\definecolor{currentfill}{rgb}{1.000000,1.000000,1.000000}%
\pgfsetfillcolor{currentfill}%
\pgfsetfillopacity{0.900000}%
\pgfsetlinewidth{0.200750pt}%
\definecolor{currentstroke}{rgb}{0.200000,0.200000,0.200000}%
\pgfsetstrokecolor{currentstroke}%
\pgfsetdash{}{0pt}%
\pgfpathmoveto{\pgfqpoint{2.692356in}{1.364573in}}%
\pgfpathlineto{\pgfqpoint{2.738018in}{1.381826in}}%
\pgfpathlineto{\pgfqpoint{2.783590in}{1.400686in}}%
\pgfpathlineto{\pgfqpoint{2.816018in}{1.380386in}}%
\pgfpathlineto{\pgfqpoint{2.848549in}{1.360056in}}%
\pgfpathlineto{\pgfqpoint{2.802896in}{1.340463in}}%
\pgfpathlineto{\pgfqpoint{2.757157in}{1.322657in}}%
\pgfpathlineto{\pgfqpoint{2.724706in}{1.343635in}}%
\pgfpathlineto{\pgfqpoint{2.692356in}{1.364573in}}%
\pgfpathclose%
\pgfusepath{stroke,fill}%
\end{pgfscope}%
\begin{pgfscope}%
\pgfpathrectangle{\pgfqpoint{0.050000in}{0.050000in}}{\pgfqpoint{4.900000in}{4.900000in}}%
\pgfusepath{clip}%
\pgfsetbuttcap%
\pgfsetroundjoin%
\definecolor{currentfill}{rgb}{0.803091,0.794971,0.882261}%
\pgfsetfillcolor{currentfill}%
\pgfsetfillopacity{0.900000}%
\pgfsetlinewidth{0.200750pt}%
\definecolor{currentstroke}{rgb}{0.200000,0.200000,0.200000}%
\pgfsetstrokecolor{currentstroke}%
\pgfsetdash{}{0pt}%
\pgfpathmoveto{\pgfqpoint{3.840218in}{1.721221in}}%
\pgfpathlineto{\pgfqpoint{3.883533in}{1.687359in}}%
\pgfpathlineto{\pgfqpoint{3.926442in}{1.643932in}}%
\pgfpathlineto{\pgfqpoint{3.960197in}{1.605614in}}%
\pgfpathlineto{\pgfqpoint{3.994050in}{1.567733in}}%
\pgfpathlineto{\pgfqpoint{3.951153in}{1.610137in}}%
\pgfpathlineto{\pgfqpoint{3.907881in}{1.644658in}}%
\pgfpathlineto{\pgfqpoint{3.874002in}{1.682797in}}%
\pgfpathlineto{\pgfqpoint{3.840218in}{1.721221in}}%
\pgfpathclose%
\pgfusepath{stroke,fill}%
\end{pgfscope}%
\begin{pgfscope}%
\pgfpathrectangle{\pgfqpoint{0.050000in}{0.050000in}}{\pgfqpoint{4.900000in}{4.900000in}}%
\pgfusepath{clip}%
\pgfsetbuttcap%
\pgfsetroundjoin%
\definecolor{currentfill}{rgb}{1.000000,1.000000,1.000000}%
\pgfsetfillcolor{currentfill}%
\pgfsetfillopacity{0.900000}%
\pgfsetlinewidth{0.200750pt}%
\definecolor{currentstroke}{rgb}{0.200000,0.200000,0.200000}%
\pgfsetstrokecolor{currentstroke}%
\pgfsetdash{}{0pt}%
\pgfpathmoveto{\pgfqpoint{4.215960in}{1.349190in}}%
\pgfpathlineto{\pgfqpoint{4.259981in}{1.364090in}}%
\pgfpathlineto{\pgfqpoint{4.303907in}{1.378956in}}%
\pgfpathlineto{\pgfqpoint{4.338741in}{1.357714in}}%
\pgfpathlineto{\pgfqpoint{4.294785in}{1.342801in}}%
\pgfpathlineto{\pgfqpoint{4.250733in}{1.327855in}}%
\pgfpathlineto{\pgfqpoint{4.215960in}{1.349190in}}%
\pgfpathclose%
\pgfusepath{stroke,fill}%
\end{pgfscope}%
\begin{pgfscope}%
\pgfpathrectangle{\pgfqpoint{0.050000in}{0.050000in}}{\pgfqpoint{4.900000in}{4.900000in}}%
\pgfusepath{clip}%
\pgfsetbuttcap%
\pgfsetroundjoin%
\definecolor{currentfill}{rgb}{0.773256,0.763906,0.864421}%
\pgfsetfillcolor{currentfill}%
\pgfsetfillopacity{0.900000}%
\pgfsetlinewidth{0.200750pt}%
\definecolor{currentstroke}{rgb}{0.200000,0.200000,0.200000}%
\pgfsetstrokecolor{currentstroke}%
\pgfsetdash{}{0pt}%
\pgfpathmoveto{\pgfqpoint{3.278662in}{1.649124in}}%
\pgfpathlineto{\pgfqpoint{3.324316in}{1.713577in}}%
\pgfpathlineto{\pgfqpoint{3.369954in}{1.778169in}}%
\pgfpathlineto{\pgfqpoint{3.403289in}{1.754881in}}%
\pgfpathlineto{\pgfqpoint{3.436720in}{1.731168in}}%
\pgfpathlineto{\pgfqpoint{3.391022in}{1.670822in}}%
\pgfpathlineto{\pgfqpoint{3.345285in}{1.609444in}}%
\pgfpathlineto{\pgfqpoint{3.311923in}{1.629474in}}%
\pgfpathlineto{\pgfqpoint{3.278662in}{1.649124in}}%
\pgfpathclose%
\pgfusepath{stroke,fill}%
\end{pgfscope}%
\begin{pgfscope}%
\pgfpathrectangle{\pgfqpoint{0.050000in}{0.050000in}}{\pgfqpoint{4.900000in}{4.900000in}}%
\pgfusepath{clip}%
\pgfsetbuttcap%
\pgfsetroundjoin%
\definecolor{currentfill}{rgb}{1.000000,1.000000,1.000000}%
\pgfsetfillcolor{currentfill}%
\pgfsetfillopacity{0.900000}%
\pgfsetlinewidth{0.200750pt}%
\definecolor{currentstroke}{rgb}{0.200000,0.200000,0.200000}%
\pgfsetstrokecolor{currentstroke}%
\pgfsetdash{}{0pt}%
\pgfpathmoveto{\pgfqpoint{2.288825in}{1.357588in}}%
\pgfpathlineto{\pgfqpoint{2.334934in}{1.372481in}}%
\pgfpathlineto{\pgfqpoint{2.380942in}{1.387362in}}%
\pgfpathlineto{\pgfqpoint{2.412740in}{1.366160in}}%
\pgfpathlineto{\pgfqpoint{2.444637in}{1.344895in}}%
\pgfpathlineto{\pgfqpoint{2.398555in}{1.329903in}}%
\pgfpathlineto{\pgfqpoint{2.352370in}{1.314908in}}%
\pgfpathlineto{\pgfqpoint{2.320548in}{1.336281in}}%
\pgfpathlineto{\pgfqpoint{2.288825in}{1.357588in}}%
\pgfpathclose%
\pgfusepath{stroke,fill}%
\end{pgfscope}%
\begin{pgfscope}%
\pgfpathrectangle{\pgfqpoint{0.050000in}{0.050000in}}{\pgfqpoint{4.900000in}{4.900000in}}%
\pgfusepath{clip}%
\pgfsetbuttcap%
\pgfsetroundjoin%
\definecolor{currentfill}{rgb}{1.000000,1.000000,1.000000}%
\pgfsetfillcolor{currentfill}%
\pgfsetfillopacity{0.900000}%
\pgfsetlinewidth{0.200750pt}%
\definecolor{currentstroke}{rgb}{0.200000,0.200000,0.200000}%
\pgfsetstrokecolor{currentstroke}%
\pgfsetdash{}{0pt}%
\pgfpathmoveto{\pgfqpoint{1.885045in}{1.353299in}}%
\pgfpathlineto{\pgfqpoint{1.931612in}{1.368190in}}%
\pgfpathlineto{\pgfqpoint{1.978077in}{1.383047in}}%
\pgfpathlineto{\pgfqpoint{2.009253in}{1.361818in}}%
\pgfpathlineto{\pgfqpoint{2.040526in}{1.340522in}}%
\pgfpathlineto{\pgfqpoint{1.993984in}{1.325571in}}%
\pgfpathlineto{\pgfqpoint{1.947339in}{1.310587in}}%
\pgfpathlineto{\pgfqpoint{1.916143in}{1.331976in}}%
\pgfpathlineto{\pgfqpoint{1.885045in}{1.353299in}}%
\pgfpathclose%
\pgfusepath{stroke,fill}%
\end{pgfscope}%
\begin{pgfscope}%
\pgfpathrectangle{\pgfqpoint{0.050000in}{0.050000in}}{\pgfqpoint{4.900000in}{4.900000in}}%
\pgfusepath{clip}%
\pgfsetbuttcap%
\pgfsetroundjoin%
\definecolor{currentfill}{rgb}{0.850827,0.844675,0.910804}%
\pgfsetfillcolor{currentfill}%
\pgfsetfillopacity{0.900000}%
\pgfsetlinewidth{0.200750pt}%
\definecolor{currentstroke}{rgb}{0.200000,0.200000,0.200000}%
\pgfsetstrokecolor{currentstroke}%
\pgfsetdash{}{0pt}%
\pgfpathmoveto{\pgfqpoint{3.187423in}{1.531423in}}%
\pgfpathlineto{\pgfqpoint{3.233026in}{1.587727in}}%
\pgfpathlineto{\pgfqpoint{3.278662in}{1.649124in}}%
\pgfpathlineto{\pgfqpoint{3.311923in}{1.629474in}}%
\pgfpathlineto{\pgfqpoint{3.345285in}{1.609444in}}%
\pgfpathlineto{\pgfqpoint{3.299549in}{1.549971in}}%
\pgfpathlineto{\pgfqpoint{3.253836in}{1.494576in}}%
\pgfpathlineto{\pgfqpoint{3.220577in}{1.513128in}}%
\pgfpathlineto{\pgfqpoint{3.187423in}{1.531423in}}%
\pgfpathclose%
\pgfusepath{stroke,fill}%
\end{pgfscope}%
\begin{pgfscope}%
\pgfpathrectangle{\pgfqpoint{0.050000in}{0.050000in}}{\pgfqpoint{4.900000in}{4.900000in}}%
\pgfusepath{clip}%
\pgfsetbuttcap%
\pgfsetroundjoin%
\definecolor{currentfill}{rgb}{1.000000,1.000000,1.000000}%
\pgfsetfillcolor{currentfill}%
\pgfsetfillopacity{0.900000}%
\pgfsetlinewidth{0.200750pt}%
\definecolor{currentstroke}{rgb}{0.200000,0.200000,0.200000}%
\pgfsetstrokecolor{currentstroke}%
\pgfsetdash{}{0pt}%
\pgfpathmoveto{\pgfqpoint{1.481009in}{1.349013in}}%
\pgfpathlineto{\pgfqpoint{1.528035in}{1.363913in}}%
\pgfpathlineto{\pgfqpoint{1.574958in}{1.378780in}}%
\pgfpathlineto{\pgfqpoint{1.605511in}{1.357538in}}%
\pgfpathlineto{\pgfqpoint{1.636160in}{1.336229in}}%
\pgfpathlineto{\pgfqpoint{1.589157in}{1.321268in}}%
\pgfpathlineto{\pgfqpoint{1.542050in}{1.306274in}}%
\pgfpathlineto{\pgfqpoint{1.511482in}{1.327677in}}%
\pgfpathlineto{\pgfqpoint{1.481009in}{1.349013in}}%
\pgfpathclose%
\pgfusepath{stroke,fill}%
\end{pgfscope}%
\begin{pgfscope}%
\pgfpathrectangle{\pgfqpoint{0.050000in}{0.050000in}}{\pgfqpoint{4.900000in}{4.900000in}}%
\pgfusepath{clip}%
\pgfsetbuttcap%
\pgfsetroundjoin%
\definecolor{currentfill}{rgb}{0.952265,0.950296,0.971457}%
\pgfsetfillcolor{currentfill}%
\pgfsetfillopacity{0.900000}%
\pgfsetlinewidth{0.200750pt}%
\definecolor{currentstroke}{rgb}{0.200000,0.200000,0.200000}%
\pgfsetstrokecolor{currentstroke}%
\pgfsetdash{}{0pt}%
\pgfpathmoveto{\pgfqpoint{4.062055in}{1.493210in}}%
\pgfpathlineto{\pgfqpoint{4.104717in}{1.450404in}}%
\pgfpathlineto{\pgfqpoint{4.147300in}{1.411365in}}%
\pgfpathlineto{\pgfqpoint{4.181515in}{1.378012in}}%
\pgfpathlineto{\pgfqpoint{4.215960in}{1.349190in}}%
\pgfpathlineto{\pgfqpoint{4.173168in}{1.380163in}}%
\pgfpathlineto{\pgfqpoint{4.130462in}{1.420218in}}%
\pgfpathlineto{\pgfqpoint{4.096208in}{1.456531in}}%
\pgfpathlineto{\pgfqpoint{4.062055in}{1.493210in}}%
\pgfpathclose%
\pgfusepath{stroke,fill}%
\end{pgfscope}%
\begin{pgfscope}%
\pgfpathrectangle{\pgfqpoint{0.050000in}{0.050000in}}{\pgfqpoint{4.900000in}{4.900000in}}%
\pgfusepath{clip}%
\pgfsetbuttcap%
\pgfsetroundjoin%
\definecolor{currentfill}{rgb}{0.671819,0.658285,0.803768}%
\pgfsetfillcolor{currentfill}%
\pgfsetfillopacity{0.900000}%
\pgfsetlinewidth{0.200750pt}%
\definecolor{currentstroke}{rgb}{0.200000,0.200000,0.200000}%
\pgfsetstrokecolor{currentstroke}%
\pgfsetdash{}{0pt}%
\pgfpathmoveto{\pgfqpoint{3.527739in}{1.834339in}}%
\pgfpathlineto{\pgfqpoint{3.572899in}{1.869643in}}%
\pgfpathlineto{\pgfqpoint{3.617709in}{1.889763in}}%
\pgfpathlineto{\pgfqpoint{3.651197in}{1.854282in}}%
\pgfpathlineto{\pgfqpoint{3.684773in}{1.818799in}}%
\pgfpathlineto{\pgfqpoint{3.640022in}{1.804238in}}%
\pgfpathlineto{\pgfqpoint{3.594900in}{1.774837in}}%
\pgfpathlineto{\pgfqpoint{3.561274in}{1.804721in}}%
\pgfpathlineto{\pgfqpoint{3.527739in}{1.834339in}}%
\pgfpathclose%
\pgfusepath{stroke,fill}%
\end{pgfscope}%
\begin{pgfscope}%
\pgfpathrectangle{\pgfqpoint{0.050000in}{0.050000in}}{\pgfqpoint{4.900000in}{4.900000in}}%
\pgfusepath{clip}%
\pgfsetbuttcap%
\pgfsetroundjoin%
\definecolor{currentfill}{rgb}{0.707620,0.695563,0.825175}%
\pgfsetfillcolor{currentfill}%
\pgfsetfillopacity{0.900000}%
\pgfsetlinewidth{0.200750pt}%
\definecolor{currentstroke}{rgb}{0.200000,0.200000,0.200000}%
\pgfsetstrokecolor{currentstroke}%
\pgfsetdash{}{0pt}%
\pgfpathmoveto{\pgfqpoint{3.684773in}{1.818799in}}%
\pgfpathlineto{\pgfqpoint{3.729089in}{1.817044in}}%
\pgfpathlineto{\pgfqpoint{3.772929in}{1.798951in}}%
\pgfpathlineto{\pgfqpoint{3.806528in}{1.759937in}}%
\pgfpathlineto{\pgfqpoint{3.840218in}{1.721221in}}%
\pgfpathlineto{\pgfqpoint{3.796444in}{1.742044in}}%
\pgfpathlineto{\pgfqpoint{3.752198in}{1.747878in}}%
\pgfpathlineto{\pgfqpoint{3.718440in}{1.783327in}}%
\pgfpathlineto{\pgfqpoint{3.684773in}{1.818799in}}%
\pgfpathclose%
\pgfusepath{stroke,fill}%
\end{pgfscope}%
\begin{pgfscope}%
\pgfpathrectangle{\pgfqpoint{0.050000in}{0.050000in}}{\pgfqpoint{4.900000in}{4.900000in}}%
\pgfusepath{clip}%
\pgfsetbuttcap%
\pgfsetroundjoin%
\definecolor{currentfill}{rgb}{1.000000,1.000000,1.000000}%
\pgfsetfillcolor{currentfill}%
\pgfsetfillopacity{0.900000}%
\pgfsetlinewidth{0.200750pt}%
\definecolor{currentstroke}{rgb}{0.200000,0.200000,0.200000}%
\pgfsetstrokecolor{currentstroke}%
\pgfsetdash{}{0pt}%
\pgfpathmoveto{\pgfqpoint{2.444637in}{1.344895in}}%
\pgfpathlineto{\pgfqpoint{2.490619in}{1.359924in}}%
\pgfpathlineto{\pgfqpoint{2.536500in}{1.375064in}}%
\pgfpathlineto{\pgfqpoint{2.568572in}{1.353898in}}%
\pgfpathlineto{\pgfqpoint{2.600745in}{1.332677in}}%
\pgfpathlineto{\pgfqpoint{2.554790in}{1.317348in}}%
\pgfpathlineto{\pgfqpoint{2.508734in}{1.302174in}}%
\pgfpathlineto{\pgfqpoint{2.476635in}{1.323566in}}%
\pgfpathlineto{\pgfqpoint{2.444637in}{1.344895in}}%
\pgfpathclose%
\pgfusepath{stroke,fill}%
\end{pgfscope}%
\begin{pgfscope}%
\pgfpathrectangle{\pgfqpoint{0.050000in}{0.050000in}}{\pgfqpoint{4.900000in}{4.900000in}}%
\pgfusepath{clip}%
\pgfsetbuttcap%
\pgfsetroundjoin%
\definecolor{currentfill}{rgb}{0.988066,0.987574,0.992864}%
\pgfsetfillcolor{currentfill}%
\pgfsetfillopacity{0.900000}%
\pgfsetlinewidth{0.200750pt}%
\definecolor{currentstroke}{rgb}{0.200000,0.200000,0.200000}%
\pgfsetstrokecolor{currentstroke}%
\pgfsetdash{}{0pt}%
\pgfpathmoveto{\pgfqpoint{2.848549in}{1.360056in}}%
\pgfpathlineto{\pgfqpoint{2.894125in}{1.382249in}}%
\pgfpathlineto{\pgfqpoint{2.939637in}{1.408018in}}%
\pgfpathlineto{\pgfqpoint{2.972361in}{1.388646in}}%
\pgfpathlineto{\pgfqpoint{3.005191in}{1.369222in}}%
\pgfpathlineto{\pgfqpoint{2.959586in}{1.342407in}}%
\pgfpathlineto{\pgfqpoint{2.913923in}{1.319301in}}%
\pgfpathlineto{\pgfqpoint{2.881184in}{1.339695in}}%
\pgfpathlineto{\pgfqpoint{2.848549in}{1.360056in}}%
\pgfpathclose%
\pgfusepath{stroke,fill}%
\end{pgfscope}%
\begin{pgfscope}%
\pgfpathrectangle{\pgfqpoint{0.050000in}{0.050000in}}{\pgfqpoint{4.900000in}{4.900000in}}%
\pgfusepath{clip}%
\pgfsetbuttcap%
\pgfsetroundjoin%
\definecolor{currentfill}{rgb}{1.000000,1.000000,1.000000}%
\pgfsetfillcolor{currentfill}%
\pgfsetfillopacity{0.900000}%
\pgfsetlinewidth{0.200750pt}%
\definecolor{currentstroke}{rgb}{0.200000,0.200000,0.200000}%
\pgfsetstrokecolor{currentstroke}%
\pgfsetdash{}{0pt}%
\pgfpathmoveto{\pgfqpoint{2.040526in}{1.340522in}}%
\pgfpathlineto{\pgfqpoint{2.086967in}{1.355441in}}%
\pgfpathlineto{\pgfqpoint{2.133305in}{1.370327in}}%
\pgfpathlineto{\pgfqpoint{2.164754in}{1.349058in}}%
\pgfpathlineto{\pgfqpoint{2.196301in}{1.327722in}}%
\pgfpathlineto{\pgfqpoint{2.149887in}{1.312742in}}%
\pgfpathlineto{\pgfqpoint{2.103369in}{1.297729in}}%
\pgfpathlineto{\pgfqpoint{2.071898in}{1.319159in}}%
\pgfpathlineto{\pgfqpoint{2.040526in}{1.340522in}}%
\pgfpathclose%
\pgfusepath{stroke,fill}%
\end{pgfscope}%
\begin{pgfscope}%
\pgfpathrectangle{\pgfqpoint{0.050000in}{0.050000in}}{\pgfqpoint{4.900000in}{4.900000in}}%
\pgfusepath{clip}%
\pgfsetbuttcap%
\pgfsetroundjoin%
\definecolor{currentfill}{rgb}{0.916463,0.913018,0.950050}%
\pgfsetfillcolor{currentfill}%
\pgfsetfillopacity{0.900000}%
\pgfsetlinewidth{0.200750pt}%
\definecolor{currentstroke}{rgb}{0.200000,0.200000,0.200000}%
\pgfsetstrokecolor{currentstroke}%
\pgfsetdash{}{0pt}%
\pgfpathmoveto{\pgfqpoint{3.096304in}{1.437886in}}%
\pgfpathlineto{\pgfqpoint{3.141853in}{1.481358in}}%
\pgfpathlineto{\pgfqpoint{3.187423in}{1.531423in}}%
\pgfpathlineto{\pgfqpoint{3.220577in}{1.513128in}}%
\pgfpathlineto{\pgfqpoint{3.253836in}{1.494576in}}%
\pgfpathlineto{\pgfqpoint{3.208153in}{1.444609in}}%
\pgfpathlineto{\pgfqpoint{3.162491in}{1.400660in}}%
\pgfpathlineto{\pgfqpoint{3.129344in}{1.419341in}}%
\pgfpathlineto{\pgfqpoint{3.096304in}{1.437886in}}%
\pgfpathclose%
\pgfusepath{stroke,fill}%
\end{pgfscope}%
\begin{pgfscope}%
\pgfpathrectangle{\pgfqpoint{0.050000in}{0.050000in}}{\pgfqpoint{4.900000in}{4.900000in}}%
\pgfusepath{clip}%
\pgfsetbuttcap%
\pgfsetroundjoin%
\definecolor{currentfill}{rgb}{1.000000,1.000000,1.000000}%
\pgfsetfillcolor{currentfill}%
\pgfsetfillopacity{0.900000}%
\pgfsetlinewidth{0.200750pt}%
\definecolor{currentstroke}{rgb}{0.200000,0.200000,0.200000}%
\pgfsetstrokecolor{currentstroke}%
\pgfsetdash{}{0pt}%
\pgfpathmoveto{\pgfqpoint{1.636160in}{1.336229in}}%
\pgfpathlineto{\pgfqpoint{1.683059in}{1.351157in}}%
\pgfpathlineto{\pgfqpoint{1.729855in}{1.366052in}}%
\pgfpathlineto{\pgfqpoint{1.760680in}{1.344769in}}%
\pgfpathlineto{\pgfqpoint{1.791602in}{1.323420in}}%
\pgfpathlineto{\pgfqpoint{1.744727in}{1.308431in}}%
\pgfpathlineto{\pgfqpoint{1.697748in}{1.293408in}}%
\pgfpathlineto{\pgfqpoint{1.666905in}{1.314852in}}%
\pgfpathlineto{\pgfqpoint{1.636160in}{1.336229in}}%
\pgfpathclose%
\pgfusepath{stroke,fill}%
\end{pgfscope}%
\begin{pgfscope}%
\pgfpathrectangle{\pgfqpoint{0.050000in}{0.050000in}}{\pgfqpoint{4.900000in}{4.900000in}}%
\pgfusepath{clip}%
\pgfsetbuttcap%
\pgfsetroundjoin%
\definecolor{currentfill}{rgb}{1.000000,1.000000,1.000000}%
\pgfsetfillcolor{currentfill}%
\pgfsetfillopacity{0.900000}%
\pgfsetlinewidth{0.200750pt}%
\definecolor{currentstroke}{rgb}{0.200000,0.200000,0.200000}%
\pgfsetstrokecolor{currentstroke}%
\pgfsetdash{}{0pt}%
\pgfpathmoveto{\pgfqpoint{1.231537in}{1.331932in}}%
\pgfpathlineto{\pgfqpoint{1.278896in}{1.346869in}}%
\pgfpathlineto{\pgfqpoint{1.326151in}{1.361774in}}%
\pgfpathlineto{\pgfqpoint{1.356352in}{1.340478in}}%
\pgfpathlineto{\pgfqpoint{1.386647in}{1.319115in}}%
\pgfpathlineto{\pgfqpoint{1.339310in}{1.304116in}}%
\pgfpathlineto{\pgfqpoint{1.291868in}{1.289085in}}%
\pgfpathlineto{\pgfqpoint{1.261655in}{1.310542in}}%
\pgfpathlineto{\pgfqpoint{1.231537in}{1.331932in}}%
\pgfpathclose%
\pgfusepath{stroke,fill}%
\end{pgfscope}%
\begin{pgfscope}%
\pgfpathrectangle{\pgfqpoint{0.050000in}{0.050000in}}{\pgfqpoint{4.900000in}{4.900000in}}%
\pgfusepath{clip}%
\pgfsetbuttcap%
\pgfsetroundjoin%
\definecolor{currentfill}{rgb}{1.000000,1.000000,1.000000}%
\pgfsetfillcolor{currentfill}%
\pgfsetfillopacity{0.900000}%
\pgfsetlinewidth{0.200750pt}%
\definecolor{currentstroke}{rgb}{0.200000,0.200000,0.200000}%
\pgfsetstrokecolor{currentstroke}%
\pgfsetdash{}{0pt}%
\pgfpathmoveto{\pgfqpoint{2.600745in}{1.332677in}}%
\pgfpathlineto{\pgfqpoint{2.646599in}{1.348329in}}%
\pgfpathlineto{\pgfqpoint{2.692356in}{1.364573in}}%
\pgfpathlineto{\pgfqpoint{2.724706in}{1.343635in}}%
\pgfpathlineto{\pgfqpoint{2.757157in}{1.322657in}}%
\pgfpathlineto{\pgfqpoint{2.711325in}{1.306010in}}%
\pgfpathlineto{\pgfqpoint{2.665395in}{1.290069in}}%
\pgfpathlineto{\pgfqpoint{2.633019in}{1.311401in}}%
\pgfpathlineto{\pgfqpoint{2.600745in}{1.332677in}}%
\pgfpathclose%
\pgfusepath{stroke,fill}%
\end{pgfscope}%
\begin{pgfscope}%
\pgfpathrectangle{\pgfqpoint{0.050000in}{0.050000in}}{\pgfqpoint{4.900000in}{4.900000in}}%
\pgfusepath{clip}%
\pgfsetbuttcap%
\pgfsetroundjoin%
\definecolor{currentfill}{rgb}{0.707620,0.695563,0.825175}%
\pgfsetfillcolor{currentfill}%
\pgfsetfillopacity{0.900000}%
\pgfsetlinewidth{0.200750pt}%
\definecolor{currentstroke}{rgb}{0.200000,0.200000,0.200000}%
\pgfsetstrokecolor{currentstroke}%
\pgfsetdash{}{0pt}%
\pgfpathmoveto{\pgfqpoint{3.436720in}{1.731168in}}%
\pgfpathlineto{\pgfqpoint{3.482318in}{1.786948in}}%
\pgfpathlineto{\pgfqpoint{3.527739in}{1.834339in}}%
\pgfpathlineto{\pgfqpoint{3.561274in}{1.804721in}}%
\pgfpathlineto{\pgfqpoint{3.594900in}{1.774837in}}%
\pgfpathlineto{\pgfqpoint{3.549488in}{1.733181in}}%
\pgfpathlineto{\pgfqpoint{3.503871in}{1.682543in}}%
\pgfpathlineto{\pgfqpoint{3.470247in}{1.707050in}}%
\pgfpathlineto{\pgfqpoint{3.436720in}{1.731168in}}%
\pgfpathclose%
\pgfusepath{stroke,fill}%
\end{pgfscope}%
\begin{pgfscope}%
\pgfpathrectangle{\pgfqpoint{0.050000in}{0.050000in}}{\pgfqpoint{4.900000in}{4.900000in}}%
\pgfusepath{clip}%
\pgfsetbuttcap%
\pgfsetroundjoin%
\definecolor{currentfill}{rgb}{1.000000,1.000000,1.000000}%
\pgfsetfillcolor{currentfill}%
\pgfsetfillopacity{0.900000}%
\pgfsetlinewidth{0.200750pt}%
\definecolor{currentstroke}{rgb}{0.200000,0.200000,0.200000}%
\pgfsetstrokecolor{currentstroke}%
\pgfsetdash{}{0pt}%
\pgfpathmoveto{\pgfqpoint{2.196301in}{1.327722in}}%
\pgfpathlineto{\pgfqpoint{2.242614in}{1.342670in}}%
\pgfpathlineto{\pgfqpoint{2.288825in}{1.357588in}}%
\pgfpathlineto{\pgfqpoint{2.320548in}{1.336281in}}%
\pgfpathlineto{\pgfqpoint{2.352370in}{1.314908in}}%
\pgfpathlineto{\pgfqpoint{2.306084in}{1.299892in}}%
\pgfpathlineto{\pgfqpoint{2.259695in}{1.284849in}}%
\pgfpathlineto{\pgfqpoint{2.227948in}{1.306319in}}%
\pgfpathlineto{\pgfqpoint{2.196301in}{1.327722in}}%
\pgfpathclose%
\pgfusepath{stroke,fill}%
\end{pgfscope}%
\begin{pgfscope}%
\pgfpathrectangle{\pgfqpoint{0.050000in}{0.050000in}}{\pgfqpoint{4.900000in}{4.900000in}}%
\pgfusepath{clip}%
\pgfsetbuttcap%
\pgfsetroundjoin%
\definecolor{currentfill}{rgb}{0.832926,0.826036,0.900100}%
\pgfsetfillcolor{currentfill}%
\pgfsetfillopacity{0.900000}%
\pgfsetlinewidth{0.200750pt}%
\definecolor{currentstroke}{rgb}{0.200000,0.200000,0.200000}%
\pgfsetstrokecolor{currentstroke}%
\pgfsetdash{}{0pt}%
\pgfpathmoveto{\pgfqpoint{3.907881in}{1.644658in}}%
\pgfpathlineto{\pgfqpoint{3.951153in}{1.610137in}}%
\pgfpathlineto{\pgfqpoint{3.994050in}{1.567733in}}%
\pgfpathlineto{\pgfqpoint{4.028002in}{1.530271in}}%
\pgfpathlineto{\pgfqpoint{4.062055in}{1.493210in}}%
\pgfpathlineto{\pgfqpoint{4.019162in}{1.534376in}}%
\pgfpathlineto{\pgfqpoint{3.975923in}{1.569210in}}%
\pgfpathlineto{\pgfqpoint{3.941854in}{1.606797in}}%
\pgfpathlineto{\pgfqpoint{3.907881in}{1.644658in}}%
\pgfpathclose%
\pgfusepath{stroke,fill}%
\end{pgfscope}%
\begin{pgfscope}%
\pgfpathrectangle{\pgfqpoint{0.050000in}{0.050000in}}{\pgfqpoint{4.900000in}{4.900000in}}%
\pgfusepath{clip}%
\pgfsetbuttcap%
\pgfsetroundjoin%
\definecolor{currentfill}{rgb}{1.000000,1.000000,1.000000}%
\pgfsetfillcolor{currentfill}%
\pgfsetfillopacity{0.900000}%
\pgfsetlinewidth{0.200750pt}%
\definecolor{currentstroke}{rgb}{0.200000,0.200000,0.200000}%
\pgfsetstrokecolor{currentstroke}%
\pgfsetdash{}{0pt}%
\pgfpathmoveto{\pgfqpoint{1.791602in}{1.323420in}}%
\pgfpathlineto{\pgfqpoint{1.838375in}{1.338376in}}%
\pgfpathlineto{\pgfqpoint{1.885045in}{1.353299in}}%
\pgfpathlineto{\pgfqpoint{1.916143in}{1.331976in}}%
\pgfpathlineto{\pgfqpoint{1.947339in}{1.310587in}}%
\pgfpathlineto{\pgfqpoint{1.900591in}{1.295569in}}%
\pgfpathlineto{\pgfqpoint{1.853739in}{1.280519in}}%
\pgfpathlineto{\pgfqpoint{1.822622in}{1.302003in}}%
\pgfpathlineto{\pgfqpoint{1.791602in}{1.323420in}}%
\pgfpathclose%
\pgfusepath{stroke,fill}%
\end{pgfscope}%
\begin{pgfscope}%
\pgfpathrectangle{\pgfqpoint{0.050000in}{0.050000in}}{\pgfqpoint{4.900000in}{4.900000in}}%
\pgfusepath{clip}%
\pgfsetbuttcap%
\pgfsetroundjoin%
\definecolor{currentfill}{rgb}{0.958231,0.956509,0.975025}%
\pgfsetfillcolor{currentfill}%
\pgfsetfillopacity{0.900000}%
\pgfsetlinewidth{0.200750pt}%
\definecolor{currentstroke}{rgb}{0.200000,0.200000,0.200000}%
\pgfsetstrokecolor{currentstroke}%
\pgfsetdash{}{0pt}%
\pgfpathmoveto{\pgfqpoint{3.005191in}{1.369222in}}%
\pgfpathlineto{\pgfqpoint{3.050757in}{1.400739in}}%
\pgfpathlineto{\pgfqpoint{3.096304in}{1.437886in}}%
\pgfpathlineto{\pgfqpoint{3.129344in}{1.419341in}}%
\pgfpathlineto{\pgfqpoint{3.162491in}{1.400660in}}%
\pgfpathlineto{\pgfqpoint{3.116836in}{1.362688in}}%
\pgfpathlineto{\pgfqpoint{3.071168in}{1.330194in}}%
\pgfpathlineto{\pgfqpoint{3.038126in}{1.349740in}}%
\pgfpathlineto{\pgfqpoint{3.005191in}{1.369222in}}%
\pgfpathclose%
\pgfusepath{stroke,fill}%
\end{pgfscope}%
\begin{pgfscope}%
\pgfpathrectangle{\pgfqpoint{0.050000in}{0.050000in}}{\pgfqpoint{4.900000in}{4.900000in}}%
\pgfusepath{clip}%
\pgfsetbuttcap%
\pgfsetroundjoin%
\definecolor{currentfill}{rgb}{0.964198,0.962722,0.978593}%
\pgfsetfillcolor{currentfill}%
\pgfsetfillopacity{0.900000}%
\pgfsetlinewidth{0.200750pt}%
\definecolor{currentstroke}{rgb}{0.200000,0.200000,0.200000}%
\pgfsetstrokecolor{currentstroke}%
\pgfsetdash{}{0pt}%
\pgfpathmoveto{\pgfqpoint{4.130462in}{1.420218in}}%
\pgfpathlineto{\pgfqpoint{4.173168in}{1.380163in}}%
\pgfpathlineto{\pgfqpoint{4.215960in}{1.349190in}}%
\pgfpathlineto{\pgfqpoint{4.250733in}{1.327855in}}%
\pgfpathlineto{\pgfqpoint{4.207549in}{1.345583in}}%
\pgfpathlineto{\pgfqpoint{4.164819in}{1.384255in}}%
\pgfpathlineto{\pgfqpoint{4.130462in}{1.420218in}}%
\pgfpathclose%
\pgfusepath{stroke,fill}%
\end{pgfscope}%
\begin{pgfscope}%
\pgfpathrectangle{\pgfqpoint{0.050000in}{0.050000in}}{\pgfqpoint{4.900000in}{4.900000in}}%
\pgfusepath{clip}%
\pgfsetbuttcap%
\pgfsetroundjoin%
\definecolor{currentfill}{rgb}{1.000000,1.000000,1.000000}%
\pgfsetfillcolor{currentfill}%
\pgfsetfillopacity{0.900000}%
\pgfsetlinewidth{0.200750pt}%
\definecolor{currentstroke}{rgb}{0.200000,0.200000,0.200000}%
\pgfsetstrokecolor{currentstroke}%
\pgfsetdash{}{0pt}%
\pgfpathmoveto{\pgfqpoint{1.386647in}{1.319115in}}%
\pgfpathlineto{\pgfqpoint{1.433880in}{1.334081in}}%
\pgfpathlineto{\pgfqpoint{1.481009in}{1.349013in}}%
\pgfpathlineto{\pgfqpoint{1.511482in}{1.327677in}}%
\pgfpathlineto{\pgfqpoint{1.542050in}{1.306274in}}%
\pgfpathlineto{\pgfqpoint{1.494840in}{1.291247in}}%
\pgfpathlineto{\pgfqpoint{1.447525in}{1.276187in}}%
\pgfpathlineto{\pgfqpoint{1.417038in}{1.297685in}}%
\pgfpathlineto{\pgfqpoint{1.386647in}{1.319115in}}%
\pgfpathclose%
\pgfusepath{stroke,fill}%
\end{pgfscope}%
\begin{pgfscope}%
\pgfpathrectangle{\pgfqpoint{0.050000in}{0.050000in}}{\pgfqpoint{4.900000in}{4.900000in}}%
\pgfusepath{clip}%
\pgfsetbuttcap%
\pgfsetroundjoin%
\definecolor{currentfill}{rgb}{0.994033,0.993787,0.996432}%
\pgfsetfillcolor{currentfill}%
\pgfsetfillopacity{0.900000}%
\pgfsetlinewidth{0.200750pt}%
\definecolor{currentstroke}{rgb}{0.200000,0.200000,0.200000}%
\pgfsetstrokecolor{currentstroke}%
\pgfsetdash{}{0pt}%
\pgfpathmoveto{\pgfqpoint{2.757157in}{1.322657in}}%
\pgfpathlineto{\pgfqpoint{2.802896in}{1.340463in}}%
\pgfpathlineto{\pgfqpoint{2.848549in}{1.360056in}}%
\pgfpathlineto{\pgfqpoint{2.881184in}{1.339695in}}%
\pgfpathlineto{\pgfqpoint{2.913923in}{1.319301in}}%
\pgfpathlineto{\pgfqpoint{2.868188in}{1.298961in}}%
\pgfpathlineto{\pgfqpoint{2.822369in}{1.280575in}}%
\pgfpathlineto{\pgfqpoint{2.789712in}{1.301637in}}%
\pgfpathlineto{\pgfqpoint{2.757157in}{1.322657in}}%
\pgfpathclose%
\pgfusepath{stroke,fill}%
\end{pgfscope}%
\begin{pgfscope}%
\pgfpathrectangle{\pgfqpoint{0.050000in}{0.050000in}}{\pgfqpoint{4.900000in}{4.900000in}}%
\pgfusepath{clip}%
\pgfsetbuttcap%
\pgfsetroundjoin%
\definecolor{currentfill}{rgb}{0.773256,0.763906,0.864421}%
\pgfsetfillcolor{currentfill}%
\pgfsetfillopacity{0.900000}%
\pgfsetlinewidth{0.200750pt}%
\definecolor{currentstroke}{rgb}{0.200000,0.200000,0.200000}%
\pgfsetstrokecolor{currentstroke}%
\pgfsetdash{}{0pt}%
\pgfpathmoveto{\pgfqpoint{3.345285in}{1.609444in}}%
\pgfpathlineto{\pgfqpoint{3.391022in}{1.670822in}}%
\pgfpathlineto{\pgfqpoint{3.436720in}{1.731168in}}%
\pgfpathlineto{\pgfqpoint{3.470247in}{1.707050in}}%
\pgfpathlineto{\pgfqpoint{3.503871in}{1.682543in}}%
\pgfpathlineto{\pgfqpoint{3.458125in}{1.626448in}}%
\pgfpathlineto{\pgfqpoint{3.412315in}{1.568271in}}%
\pgfpathlineto{\pgfqpoint{3.378749in}{1.589041in}}%
\pgfpathlineto{\pgfqpoint{3.345285in}{1.609444in}}%
\pgfpathclose%
\pgfusepath{stroke,fill}%
\end{pgfscope}%
\begin{pgfscope}%
\pgfpathrectangle{\pgfqpoint{0.050000in}{0.050000in}}{\pgfqpoint{4.900000in}{4.900000in}}%
\pgfusepath{clip}%
\pgfsetbuttcap%
\pgfsetroundjoin%
\definecolor{currentfill}{rgb}{1.000000,1.000000,1.000000}%
\pgfsetfillcolor{currentfill}%
\pgfsetfillopacity{0.900000}%
\pgfsetlinewidth{0.200750pt}%
\definecolor{currentstroke}{rgb}{0.200000,0.200000,0.200000}%
\pgfsetstrokecolor{currentstroke}%
\pgfsetdash{}{0pt}%
\pgfpathmoveto{\pgfqpoint{2.352370in}{1.314908in}}%
\pgfpathlineto{\pgfqpoint{2.398555in}{1.329903in}}%
\pgfpathlineto{\pgfqpoint{2.444637in}{1.344895in}}%
\pgfpathlineto{\pgfqpoint{2.476635in}{1.323566in}}%
\pgfpathlineto{\pgfqpoint{2.508734in}{1.302174in}}%
\pgfpathlineto{\pgfqpoint{2.462576in}{1.287062in}}%
\pgfpathlineto{\pgfqpoint{2.416317in}{1.271960in}}%
\pgfpathlineto{\pgfqpoint{2.384293in}{1.293467in}}%
\pgfpathlineto{\pgfqpoint{2.352370in}{1.314908in}}%
\pgfpathclose%
\pgfusepath{stroke,fill}%
\end{pgfscope}%
\begin{pgfscope}%
\pgfpathrectangle{\pgfqpoint{0.050000in}{0.050000in}}{\pgfqpoint{4.900000in}{4.900000in}}%
\pgfusepath{clip}%
\pgfsetbuttcap%
\pgfsetroundjoin%
\definecolor{currentfill}{rgb}{1.000000,1.000000,1.000000}%
\pgfsetfillcolor{currentfill}%
\pgfsetfillopacity{0.900000}%
\pgfsetlinewidth{0.200750pt}%
\definecolor{currentstroke}{rgb}{0.200000,0.200000,0.200000}%
\pgfsetstrokecolor{currentstroke}%
\pgfsetdash{}{0pt}%
\pgfpathmoveto{\pgfqpoint{1.947339in}{1.310587in}}%
\pgfpathlineto{\pgfqpoint{1.993984in}{1.325571in}}%
\pgfpathlineto{\pgfqpoint{2.040526in}{1.340522in}}%
\pgfpathlineto{\pgfqpoint{2.071898in}{1.319159in}}%
\pgfpathlineto{\pgfqpoint{2.103369in}{1.297729in}}%
\pgfpathlineto{\pgfqpoint{2.056749in}{1.282683in}}%
\pgfpathlineto{\pgfqpoint{2.010026in}{1.267604in}}%
\pgfpathlineto{\pgfqpoint{1.978633in}{1.289129in}}%
\pgfpathlineto{\pgfqpoint{1.947339in}{1.310587in}}%
\pgfpathclose%
\pgfusepath{stroke,fill}%
\end{pgfscope}%
\begin{pgfscope}%
\pgfpathrectangle{\pgfqpoint{0.050000in}{0.050000in}}{\pgfqpoint{4.900000in}{4.900000in}}%
\pgfusepath{clip}%
\pgfsetbuttcap%
\pgfsetroundjoin%
\definecolor{currentfill}{rgb}{0.737455,0.726628,0.843014}%
\pgfsetfillcolor{currentfill}%
\pgfsetfillopacity{0.900000}%
\pgfsetlinewidth{0.200750pt}%
\definecolor{currentstroke}{rgb}{0.200000,0.200000,0.200000}%
\pgfsetstrokecolor{currentstroke}%
\pgfsetdash{}{0pt}%
\pgfpathmoveto{\pgfqpoint{3.752198in}{1.747878in}}%
\pgfpathlineto{\pgfqpoint{3.796444in}{1.742044in}}%
\pgfpathlineto{\pgfqpoint{3.840218in}{1.721221in}}%
\pgfpathlineto{\pgfqpoint{3.874002in}{1.682797in}}%
\pgfpathlineto{\pgfqpoint{3.907881in}{1.644658in}}%
\pgfpathlineto{\pgfqpoint{3.864166in}{1.667709in}}%
\pgfpathlineto{\pgfqpoint{3.819988in}{1.677084in}}%
\pgfpathlineto{\pgfqpoint{3.786047in}{1.712461in}}%
\pgfpathlineto{\pgfqpoint{3.752198in}{1.747878in}}%
\pgfpathclose%
\pgfusepath{stroke,fill}%
\end{pgfscope}%
\begin{pgfscope}%
\pgfpathrectangle{\pgfqpoint{0.050000in}{0.050000in}}{\pgfqpoint{4.900000in}{4.900000in}}%
\pgfusepath{clip}%
\pgfsetbuttcap%
\pgfsetroundjoin%
\definecolor{currentfill}{rgb}{0.689719,0.676924,0.814471}%
\pgfsetfillcolor{currentfill}%
\pgfsetfillopacity{0.900000}%
\pgfsetlinewidth{0.200750pt}%
\definecolor{currentstroke}{rgb}{0.200000,0.200000,0.200000}%
\pgfsetstrokecolor{currentstroke}%
\pgfsetdash{}{0pt}%
\pgfpathmoveto{\pgfqpoint{3.594900in}{1.774837in}}%
\pgfpathlineto{\pgfqpoint{3.640022in}{1.804238in}}%
\pgfpathlineto{\pgfqpoint{3.684773in}{1.818799in}}%
\pgfpathlineto{\pgfqpoint{3.718440in}{1.783327in}}%
\pgfpathlineto{\pgfqpoint{3.752198in}{1.747878in}}%
\pgfpathlineto{\pgfqpoint{3.707507in}{1.738318in}}%
\pgfpathlineto{\pgfqpoint{3.662429in}{1.714357in}}%
\pgfpathlineto{\pgfqpoint{3.628618in}{1.744710in}}%
\pgfpathlineto{\pgfqpoint{3.594900in}{1.774837in}}%
\pgfpathclose%
\pgfusepath{stroke,fill}%
\end{pgfscope}%
\begin{pgfscope}%
\pgfpathrectangle{\pgfqpoint{0.050000in}{0.050000in}}{\pgfqpoint{4.900000in}{4.900000in}}%
\pgfusepath{clip}%
\pgfsetbuttcap%
\pgfsetroundjoin%
\definecolor{currentfill}{rgb}{0.844860,0.838462,0.907236}%
\pgfsetfillcolor{currentfill}%
\pgfsetfillopacity{0.900000}%
\pgfsetlinewidth{0.200750pt}%
\definecolor{currentstroke}{rgb}{0.200000,0.200000,0.200000}%
\pgfsetstrokecolor{currentstroke}%
\pgfsetdash{}{0pt}%
\pgfpathmoveto{\pgfqpoint{3.253836in}{1.494576in}}%
\pgfpathlineto{\pgfqpoint{3.299549in}{1.549971in}}%
\pgfpathlineto{\pgfqpoint{3.345285in}{1.609444in}}%
\pgfpathlineto{\pgfqpoint{3.378749in}{1.589041in}}%
\pgfpathlineto{\pgfqpoint{3.412315in}{1.568271in}}%
\pgfpathlineto{\pgfqpoint{3.366487in}{1.510933in}}%
\pgfpathlineto{\pgfqpoint{3.320671in}{1.456693in}}%
\pgfpathlineto{\pgfqpoint{3.287201in}{1.475765in}}%
\pgfpathlineto{\pgfqpoint{3.253836in}{1.494576in}}%
\pgfpathclose%
\pgfusepath{stroke,fill}%
\end{pgfscope}%
\begin{pgfscope}%
\pgfpathrectangle{\pgfqpoint{0.050000in}{0.050000in}}{\pgfqpoint{4.900000in}{4.900000in}}%
\pgfusepath{clip}%
\pgfsetbuttcap%
\pgfsetroundjoin%
\definecolor{currentfill}{rgb}{1.000000,1.000000,1.000000}%
\pgfsetfillcolor{currentfill}%
\pgfsetfillopacity{0.900000}%
\pgfsetlinewidth{0.200750pt}%
\definecolor{currentstroke}{rgb}{0.200000,0.200000,0.200000}%
\pgfsetstrokecolor{currentstroke}%
\pgfsetdash{}{0pt}%
\pgfpathmoveto{\pgfqpoint{1.542050in}{1.306274in}}%
\pgfpathlineto{\pgfqpoint{1.589157in}{1.321268in}}%
\pgfpathlineto{\pgfqpoint{1.636160in}{1.336229in}}%
\pgfpathlineto{\pgfqpoint{1.666905in}{1.314852in}}%
\pgfpathlineto{\pgfqpoint{1.697748in}{1.293408in}}%
\pgfpathlineto{\pgfqpoint{1.650665in}{1.278353in}}%
\pgfpathlineto{\pgfqpoint{1.603477in}{1.263264in}}%
\pgfpathlineto{\pgfqpoint{1.572715in}{1.284803in}}%
\pgfpathlineto{\pgfqpoint{1.542050in}{1.306274in}}%
\pgfpathclose%
\pgfusepath{stroke,fill}%
\end{pgfscope}%
\begin{pgfscope}%
\pgfpathrectangle{\pgfqpoint{0.050000in}{0.050000in}}{\pgfqpoint{4.900000in}{4.900000in}}%
\pgfusepath{clip}%
\pgfsetbuttcap%
\pgfsetroundjoin%
\definecolor{currentfill}{rgb}{1.000000,1.000000,1.000000}%
\pgfsetfillcolor{currentfill}%
\pgfsetfillopacity{0.900000}%
\pgfsetlinewidth{0.200750pt}%
\definecolor{currentstroke}{rgb}{0.200000,0.200000,0.200000}%
\pgfsetstrokecolor{currentstroke}%
\pgfsetdash{}{0pt}%
\pgfpathmoveto{\pgfqpoint{2.508734in}{1.302174in}}%
\pgfpathlineto{\pgfqpoint{2.554790in}{1.317348in}}%
\pgfpathlineto{\pgfqpoint{2.600745in}{1.332677in}}%
\pgfpathlineto{\pgfqpoint{2.633019in}{1.311401in}}%
\pgfpathlineto{\pgfqpoint{2.665395in}{1.290069in}}%
\pgfpathlineto{\pgfqpoint{2.619366in}{1.274530in}}%
\pgfpathlineto{\pgfqpoint{2.573236in}{1.259198in}}%
\pgfpathlineto{\pgfqpoint{2.540934in}{1.280718in}}%
\pgfpathlineto{\pgfqpoint{2.508734in}{1.302174in}}%
\pgfpathclose%
\pgfusepath{stroke,fill}%
\end{pgfscope}%
\begin{pgfscope}%
\pgfpathrectangle{\pgfqpoint{0.050000in}{0.050000in}}{\pgfqpoint{4.900000in}{4.900000in}}%
\pgfusepath{clip}%
\pgfsetbuttcap%
\pgfsetroundjoin%
\definecolor{currentfill}{rgb}{0.982099,0.981361,0.989296}%
\pgfsetfillcolor{currentfill}%
\pgfsetfillopacity{0.900000}%
\pgfsetlinewidth{0.200750pt}%
\definecolor{currentstroke}{rgb}{0.200000,0.200000,0.200000}%
\pgfsetstrokecolor{currentstroke}%
\pgfsetdash{}{0pt}%
\pgfpathmoveto{\pgfqpoint{2.913923in}{1.319301in}}%
\pgfpathlineto{\pgfqpoint{2.959586in}{1.342407in}}%
\pgfpathlineto{\pgfqpoint{3.005191in}{1.369222in}}%
\pgfpathlineto{\pgfqpoint{3.038126in}{1.349740in}}%
\pgfpathlineto{\pgfqpoint{3.071168in}{1.330194in}}%
\pgfpathlineto{\pgfqpoint{3.025467in}{1.302397in}}%
\pgfpathlineto{\pgfqpoint{2.979715in}{1.278399in}}%
\pgfpathlineto{\pgfqpoint{2.946766in}{1.298870in}}%
\pgfpathlineto{\pgfqpoint{2.913923in}{1.319301in}}%
\pgfpathclose%
\pgfusepath{stroke,fill}%
\end{pgfscope}%
\begin{pgfscope}%
\pgfpathrectangle{\pgfqpoint{0.050000in}{0.050000in}}{\pgfqpoint{4.900000in}{4.900000in}}%
\pgfusepath{clip}%
\pgfsetbuttcap%
\pgfsetroundjoin%
\definecolor{currentfill}{rgb}{1.000000,1.000000,1.000000}%
\pgfsetfillcolor{currentfill}%
\pgfsetfillopacity{0.900000}%
\pgfsetlinewidth{0.200750pt}%
\definecolor{currentstroke}{rgb}{0.200000,0.200000,0.200000}%
\pgfsetstrokecolor{currentstroke}%
\pgfsetdash{}{0pt}%
\pgfpathmoveto{\pgfqpoint{2.103369in}{1.297729in}}%
\pgfpathlineto{\pgfqpoint{2.149887in}{1.312742in}}%
\pgfpathlineto{\pgfqpoint{2.196301in}{1.327722in}}%
\pgfpathlineto{\pgfqpoint{2.227948in}{1.306319in}}%
\pgfpathlineto{\pgfqpoint{2.259695in}{1.284849in}}%
\pgfpathlineto{\pgfqpoint{2.213204in}{1.269774in}}%
\pgfpathlineto{\pgfqpoint{2.166610in}{1.254666in}}%
\pgfpathlineto{\pgfqpoint{2.134940in}{1.276231in}}%
\pgfpathlineto{\pgfqpoint{2.103369in}{1.297729in}}%
\pgfpathclose%
\pgfusepath{stroke,fill}%
\end{pgfscope}%
\begin{pgfscope}%
\pgfpathrectangle{\pgfqpoint{0.050000in}{0.050000in}}{\pgfqpoint{4.900000in}{4.900000in}}%
\pgfusepath{clip}%
\pgfsetbuttcap%
\pgfsetroundjoin%
\definecolor{currentfill}{rgb}{0.910496,0.906805,0.946482}%
\pgfsetfillcolor{currentfill}%
\pgfsetfillopacity{0.900000}%
\pgfsetlinewidth{0.200750pt}%
\definecolor{currentstroke}{rgb}{0.200000,0.200000,0.200000}%
\pgfsetstrokecolor{currentstroke}%
\pgfsetdash{}{0pt}%
\pgfpathmoveto{\pgfqpoint{3.162491in}{1.400660in}}%
\pgfpathlineto{\pgfqpoint{3.208153in}{1.444609in}}%
\pgfpathlineto{\pgfqpoint{3.253836in}{1.494576in}}%
\pgfpathlineto{\pgfqpoint{3.287201in}{1.475765in}}%
\pgfpathlineto{\pgfqpoint{3.320671in}{1.456693in}}%
\pgfpathlineto{\pgfqpoint{3.274877in}{1.407071in}}%
\pgfpathlineto{\pgfqpoint{3.229105in}{1.362864in}}%
\pgfpathlineto{\pgfqpoint{3.195744in}{1.381836in}}%
\pgfpathlineto{\pgfqpoint{3.162491in}{1.400660in}}%
\pgfpathclose%
\pgfusepath{stroke,fill}%
\end{pgfscope}%
\begin{pgfscope}%
\pgfpathrectangle{\pgfqpoint{0.050000in}{0.050000in}}{\pgfqpoint{4.900000in}{4.900000in}}%
\pgfusepath{clip}%
\pgfsetbuttcap%
\pgfsetroundjoin%
\definecolor{currentfill}{rgb}{1.000000,1.000000,1.000000}%
\pgfsetfillcolor{currentfill}%
\pgfsetfillopacity{0.900000}%
\pgfsetlinewidth{0.200750pt}%
\definecolor{currentstroke}{rgb}{0.200000,0.200000,0.200000}%
\pgfsetstrokecolor{currentstroke}%
\pgfsetdash{}{0pt}%
\pgfpathmoveto{\pgfqpoint{1.697748in}{1.293408in}}%
\pgfpathlineto{\pgfqpoint{1.744727in}{1.308431in}}%
\pgfpathlineto{\pgfqpoint{1.791602in}{1.323420in}}%
\pgfpathlineto{\pgfqpoint{1.822622in}{1.302003in}}%
\pgfpathlineto{\pgfqpoint{1.853739in}{1.280519in}}%
\pgfpathlineto{\pgfqpoint{1.806784in}{1.265435in}}%
\pgfpathlineto{\pgfqpoint{1.759725in}{1.250317in}}%
\pgfpathlineto{\pgfqpoint{1.728687in}{1.271897in}}%
\pgfpathlineto{\pgfqpoint{1.697748in}{1.293408in}}%
\pgfpathclose%
\pgfusepath{stroke,fill}%
\end{pgfscope}%
\begin{pgfscope}%
\pgfpathrectangle{\pgfqpoint{0.050000in}{0.050000in}}{\pgfqpoint{4.900000in}{4.900000in}}%
\pgfusepath{clip}%
\pgfsetbuttcap%
\pgfsetroundjoin%
\definecolor{currentfill}{rgb}{0.862760,0.857101,0.917939}%
\pgfsetfillcolor{currentfill}%
\pgfsetfillopacity{0.900000}%
\pgfsetlinewidth{0.200750pt}%
\definecolor{currentstroke}{rgb}{0.200000,0.200000,0.200000}%
\pgfsetstrokecolor{currentstroke}%
\pgfsetdash{}{0pt}%
\pgfpathmoveto{\pgfqpoint{3.975923in}{1.569210in}}%
\pgfpathlineto{\pgfqpoint{4.019162in}{1.534376in}}%
\pgfpathlineto{\pgfqpoint{4.062055in}{1.493210in}}%
\pgfpathlineto{\pgfqpoint{4.096208in}{1.456531in}}%
\pgfpathlineto{\pgfqpoint{4.130462in}{1.420218in}}%
\pgfpathlineto{\pgfqpoint{4.087565in}{1.459973in}}%
\pgfpathlineto{\pgfqpoint{4.044353in}{1.494824in}}%
\pgfpathlineto{\pgfqpoint{4.010089in}{1.531887in}}%
\pgfpathlineto{\pgfqpoint{3.975923in}{1.569210in}}%
\pgfpathclose%
\pgfusepath{stroke,fill}%
\end{pgfscope}%
\begin{pgfscope}%
\pgfpathrectangle{\pgfqpoint{0.050000in}{0.050000in}}{\pgfqpoint{4.900000in}{4.900000in}}%
\pgfusepath{clip}%
\pgfsetbuttcap%
\pgfsetroundjoin%
\definecolor{currentfill}{rgb}{1.000000,1.000000,1.000000}%
\pgfsetfillcolor{currentfill}%
\pgfsetfillopacity{0.900000}%
\pgfsetlinewidth{0.200750pt}%
\definecolor{currentstroke}{rgb}{0.200000,0.200000,0.200000}%
\pgfsetstrokecolor{currentstroke}%
\pgfsetdash{}{0pt}%
\pgfpathmoveto{\pgfqpoint{1.291868in}{1.289085in}}%
\pgfpathlineto{\pgfqpoint{1.339310in}{1.304116in}}%
\pgfpathlineto{\pgfqpoint{1.386647in}{1.319115in}}%
\pgfpathlineto{\pgfqpoint{1.417038in}{1.297685in}}%
\pgfpathlineto{\pgfqpoint{1.447525in}{1.276187in}}%
\pgfpathlineto{\pgfqpoint{1.400106in}{1.261093in}}%
\pgfpathlineto{\pgfqpoint{1.352581in}{1.245966in}}%
\pgfpathlineto{\pgfqpoint{1.322177in}{1.267560in}}%
\pgfpathlineto{\pgfqpoint{1.291868in}{1.289085in}}%
\pgfpathclose%
\pgfusepath{stroke,fill}%
\end{pgfscope}%
\begin{pgfscope}%
\pgfpathrectangle{\pgfqpoint{0.050000in}{0.050000in}}{\pgfqpoint{4.900000in}{4.900000in}}%
\pgfusepath{clip}%
\pgfsetbuttcap%
\pgfsetroundjoin%
\definecolor{currentfill}{rgb}{1.000000,1.000000,1.000000}%
\pgfsetfillcolor{currentfill}%
\pgfsetfillopacity{0.900000}%
\pgfsetlinewidth{0.200750pt}%
\definecolor{currentstroke}{rgb}{0.200000,0.200000,0.200000}%
\pgfsetstrokecolor{currentstroke}%
\pgfsetdash{}{0pt}%
\pgfpathmoveto{\pgfqpoint{2.665395in}{1.290069in}}%
\pgfpathlineto{\pgfqpoint{2.711325in}{1.306010in}}%
\pgfpathlineto{\pgfqpoint{2.757157in}{1.322657in}}%
\pgfpathlineto{\pgfqpoint{2.789712in}{1.301637in}}%
\pgfpathlineto{\pgfqpoint{2.822369in}{1.280575in}}%
\pgfpathlineto{\pgfqpoint{2.776460in}{1.263497in}}%
\pgfpathlineto{\pgfqpoint{2.730456in}{1.247243in}}%
\pgfpathlineto{\pgfqpoint{2.697874in}{1.268683in}}%
\pgfpathlineto{\pgfqpoint{2.665395in}{1.290069in}}%
\pgfpathclose%
\pgfusepath{stroke,fill}%
\end{pgfscope}%
\begin{pgfscope}%
\pgfpathrectangle{\pgfqpoint{0.050000in}{0.050000in}}{\pgfqpoint{4.900000in}{4.900000in}}%
\pgfusepath{clip}%
\pgfsetbuttcap%
\pgfsetroundjoin%
\definecolor{currentfill}{rgb}{0.719554,0.707989,0.832311}%
\pgfsetfillcolor{currentfill}%
\pgfsetfillopacity{0.900000}%
\pgfsetlinewidth{0.200750pt}%
\definecolor{currentstroke}{rgb}{0.200000,0.200000,0.200000}%
\pgfsetstrokecolor{currentstroke}%
\pgfsetdash{}{0pt}%
\pgfpathmoveto{\pgfqpoint{3.503871in}{1.682543in}}%
\pgfpathlineto{\pgfqpoint{3.549488in}{1.733181in}}%
\pgfpathlineto{\pgfqpoint{3.594900in}{1.774837in}}%
\pgfpathlineto{\pgfqpoint{3.628618in}{1.744710in}}%
\pgfpathlineto{\pgfqpoint{3.662429in}{1.714357in}}%
\pgfpathlineto{\pgfqpoint{3.617035in}{1.678106in}}%
\pgfpathlineto{\pgfqpoint{3.571408in}{1.632432in}}%
\pgfpathlineto{\pgfqpoint{3.537591in}{1.657665in}}%
\pgfpathlineto{\pgfqpoint{3.503871in}{1.682543in}}%
\pgfpathclose%
\pgfusepath{stroke,fill}%
\end{pgfscope}%
\begin{pgfscope}%
\pgfpathrectangle{\pgfqpoint{0.050000in}{0.050000in}}{\pgfqpoint{4.900000in}{4.900000in}}%
\pgfusepath{clip}%
\pgfsetbuttcap%
\pgfsetroundjoin%
\definecolor{currentfill}{rgb}{1.000000,1.000000,1.000000}%
\pgfsetfillcolor{currentfill}%
\pgfsetfillopacity{0.900000}%
\pgfsetlinewidth{0.200750pt}%
\definecolor{currentstroke}{rgb}{0.200000,0.200000,0.200000}%
\pgfsetstrokecolor{currentstroke}%
\pgfsetdash{}{0pt}%
\pgfpathmoveto{\pgfqpoint{2.259695in}{1.284849in}}%
\pgfpathlineto{\pgfqpoint{2.306084in}{1.299892in}}%
\pgfpathlineto{\pgfqpoint{2.352370in}{1.314908in}}%
\pgfpathlineto{\pgfqpoint{2.384293in}{1.293467in}}%
\pgfpathlineto{\pgfqpoint{2.416317in}{1.271960in}}%
\pgfpathlineto{\pgfqpoint{2.369955in}{1.256845in}}%
\pgfpathlineto{\pgfqpoint{2.323490in}{1.241704in}}%
\pgfpathlineto{\pgfqpoint{2.291542in}{1.263311in}}%
\pgfpathlineto{\pgfqpoint{2.259695in}{1.284849in}}%
\pgfpathclose%
\pgfusepath{stroke,fill}%
\end{pgfscope}%
\begin{pgfscope}%
\pgfpathrectangle{\pgfqpoint{0.050000in}{0.050000in}}{\pgfqpoint{4.900000in}{4.900000in}}%
\pgfusepath{clip}%
\pgfsetbuttcap%
\pgfsetroundjoin%
\definecolor{currentfill}{rgb}{1.000000,1.000000,1.000000}%
\pgfsetfillcolor{currentfill}%
\pgfsetfillopacity{0.900000}%
\pgfsetlinewidth{0.200750pt}%
\definecolor{currentstroke}{rgb}{0.200000,0.200000,0.200000}%
\pgfsetstrokecolor{currentstroke}%
\pgfsetdash{}{0pt}%
\pgfpathmoveto{\pgfqpoint{1.853739in}{1.280519in}}%
\pgfpathlineto{\pgfqpoint{1.900591in}{1.295569in}}%
\pgfpathlineto{\pgfqpoint{1.947339in}{1.310587in}}%
\pgfpathlineto{\pgfqpoint{1.978633in}{1.289129in}}%
\pgfpathlineto{\pgfqpoint{2.010026in}{1.267604in}}%
\pgfpathlineto{\pgfqpoint{1.963200in}{1.252492in}}%
\pgfpathlineto{\pgfqpoint{1.916269in}{1.237346in}}%
\pgfpathlineto{\pgfqpoint{1.884955in}{1.258966in}}%
\pgfpathlineto{\pgfqpoint{1.853739in}{1.280519in}}%
\pgfpathclose%
\pgfusepath{stroke,fill}%
\end{pgfscope}%
\begin{pgfscope}%
\pgfpathrectangle{\pgfqpoint{0.050000in}{0.050000in}}{\pgfqpoint{4.900000in}{4.900000in}}%
\pgfusepath{clip}%
\pgfsetbuttcap%
\pgfsetroundjoin%
\definecolor{currentfill}{rgb}{0.952265,0.950296,0.971457}%
\pgfsetfillcolor{currentfill}%
\pgfsetfillopacity{0.900000}%
\pgfsetlinewidth{0.200750pt}%
\definecolor{currentstroke}{rgb}{0.200000,0.200000,0.200000}%
\pgfsetstrokecolor{currentstroke}%
\pgfsetdash{}{0pt}%
\pgfpathmoveto{\pgfqpoint{3.071168in}{1.330194in}}%
\pgfpathlineto{\pgfqpoint{3.116836in}{1.362688in}}%
\pgfpathlineto{\pgfqpoint{3.162491in}{1.400660in}}%
\pgfpathlineto{\pgfqpoint{3.195744in}{1.381836in}}%
\pgfpathlineto{\pgfqpoint{3.229105in}{1.362864in}}%
\pgfpathlineto{\pgfqpoint{3.183342in}{1.324241in}}%
\pgfpathlineto{\pgfqpoint{3.137571in}{1.290891in}}%
\pgfpathlineto{\pgfqpoint{3.104316in}{1.310580in}}%
\pgfpathlineto{\pgfqpoint{3.071168in}{1.330194in}}%
\pgfpathclose%
\pgfusepath{stroke,fill}%
\end{pgfscope}%
\begin{pgfscope}%
\pgfpathrectangle{\pgfqpoint{0.050000in}{0.050000in}}{\pgfqpoint{4.900000in}{4.900000in}}%
\pgfusepath{clip}%
\pgfsetbuttcap%
\pgfsetroundjoin%
\definecolor{currentfill}{rgb}{1.000000,1.000000,1.000000}%
\pgfsetfillcolor{currentfill}%
\pgfsetfillopacity{0.900000}%
\pgfsetlinewidth{0.200750pt}%
\definecolor{currentstroke}{rgb}{0.200000,0.200000,0.200000}%
\pgfsetstrokecolor{currentstroke}%
\pgfsetdash{}{0pt}%
\pgfpathmoveto{\pgfqpoint{1.447525in}{1.276187in}}%
\pgfpathlineto{\pgfqpoint{1.494840in}{1.291247in}}%
\pgfpathlineto{\pgfqpoint{1.542050in}{1.306274in}}%
\pgfpathlineto{\pgfqpoint{1.572715in}{1.284803in}}%
\pgfpathlineto{\pgfqpoint{1.603477in}{1.263264in}}%
\pgfpathlineto{\pgfqpoint{1.556186in}{1.248142in}}%
\pgfpathlineto{\pgfqpoint{1.508789in}{1.232987in}}%
\pgfpathlineto{\pgfqpoint{1.478109in}{1.254621in}}%
\pgfpathlineto{\pgfqpoint{1.447525in}{1.276187in}}%
\pgfpathclose%
\pgfusepath{stroke,fill}%
\end{pgfscope}%
\begin{pgfscope}%
\pgfpathrectangle{\pgfqpoint{0.050000in}{0.050000in}}{\pgfqpoint{4.900000in}{4.900000in}}%
\pgfusepath{clip}%
\pgfsetbuttcap%
\pgfsetroundjoin%
\definecolor{currentfill}{rgb}{0.779223,0.770119,0.867989}%
\pgfsetfillcolor{currentfill}%
\pgfsetfillopacity{0.900000}%
\pgfsetlinewidth{0.200750pt}%
\definecolor{currentstroke}{rgb}{0.200000,0.200000,0.200000}%
\pgfsetstrokecolor{currentstroke}%
\pgfsetdash{}{0pt}%
\pgfpathmoveto{\pgfqpoint{3.412315in}{1.568271in}}%
\pgfpathlineto{\pgfqpoint{3.458125in}{1.626448in}}%
\pgfpathlineto{\pgfqpoint{3.503871in}{1.682543in}}%
\pgfpathlineto{\pgfqpoint{3.537591in}{1.657665in}}%
\pgfpathlineto{\pgfqpoint{3.571408in}{1.632432in}}%
\pgfpathlineto{\pgfqpoint{3.525624in}{1.580553in}}%
\pgfpathlineto{\pgfqpoint{3.479751in}{1.525665in}}%
\pgfpathlineto{\pgfqpoint{3.445982in}{1.547144in}}%
\pgfpathlineto{\pgfqpoint{3.412315in}{1.568271in}}%
\pgfpathclose%
\pgfusepath{stroke,fill}%
\end{pgfscope}%
\begin{pgfscope}%
\pgfpathrectangle{\pgfqpoint{0.050000in}{0.050000in}}{\pgfqpoint{4.900000in}{4.900000in}}%
\pgfusepath{clip}%
\pgfsetbuttcap%
\pgfsetroundjoin%
\definecolor{currentfill}{rgb}{0.761323,0.751480,0.857286}%
\pgfsetfillcolor{currentfill}%
\pgfsetfillopacity{0.900000}%
\pgfsetlinewidth{0.200750pt}%
\definecolor{currentstroke}{rgb}{0.200000,0.200000,0.200000}%
\pgfsetstrokecolor{currentstroke}%
\pgfsetdash{}{0pt}%
\pgfpathmoveto{\pgfqpoint{3.819988in}{1.677084in}}%
\pgfpathlineto{\pgfqpoint{3.864166in}{1.667709in}}%
\pgfpathlineto{\pgfqpoint{3.907881in}{1.644658in}}%
\pgfpathlineto{\pgfqpoint{3.941854in}{1.606797in}}%
\pgfpathlineto{\pgfqpoint{3.975923in}{1.569210in}}%
\pgfpathlineto{\pgfqpoint{3.932264in}{1.594051in}}%
\pgfpathlineto{\pgfqpoint{3.888150in}{1.606483in}}%
\pgfpathlineto{\pgfqpoint{3.854022in}{1.641756in}}%
\pgfpathlineto{\pgfqpoint{3.819988in}{1.677084in}}%
\pgfpathclose%
\pgfusepath{stroke,fill}%
\end{pgfscope}%
\begin{pgfscope}%
\pgfpathrectangle{\pgfqpoint{0.050000in}{0.050000in}}{\pgfqpoint{4.900000in}{4.900000in}}%
\pgfusepath{clip}%
\pgfsetbuttcap%
\pgfsetroundjoin%
\definecolor{currentfill}{rgb}{0.994033,0.993787,0.996432}%
\pgfsetfillcolor{currentfill}%
\pgfsetfillopacity{0.900000}%
\pgfsetlinewidth{0.200750pt}%
\definecolor{currentstroke}{rgb}{0.200000,0.200000,0.200000}%
\pgfsetstrokecolor{currentstroke}%
\pgfsetdash{}{0pt}%
\pgfpathmoveto{\pgfqpoint{2.822369in}{1.280575in}}%
\pgfpathlineto{\pgfqpoint{2.868188in}{1.298961in}}%
\pgfpathlineto{\pgfqpoint{2.913923in}{1.319301in}}%
\pgfpathlineto{\pgfqpoint{2.946766in}{1.298870in}}%
\pgfpathlineto{\pgfqpoint{2.979715in}{1.278399in}}%
\pgfpathlineto{\pgfqpoint{2.933896in}{1.257308in}}%
\pgfpathlineto{\pgfqpoint{2.887997in}{1.238325in}}%
\pgfpathlineto{\pgfqpoint{2.855131in}{1.259472in}}%
\pgfpathlineto{\pgfqpoint{2.822369in}{1.280575in}}%
\pgfpathclose%
\pgfusepath{stroke,fill}%
\end{pgfscope}%
\begin{pgfscope}%
\pgfpathrectangle{\pgfqpoint{0.050000in}{0.050000in}}{\pgfqpoint{4.900000in}{4.900000in}}%
\pgfusepath{clip}%
\pgfsetbuttcap%
\pgfsetroundjoin%
\definecolor{currentfill}{rgb}{0.707620,0.695563,0.825175}%
\pgfsetfillcolor{currentfill}%
\pgfsetfillopacity{0.900000}%
\pgfsetlinewidth{0.200750pt}%
\definecolor{currentstroke}{rgb}{0.200000,0.200000,0.200000}%
\pgfsetstrokecolor{currentstroke}%
\pgfsetdash{}{0pt}%
\pgfpathmoveto{\pgfqpoint{3.662429in}{1.714357in}}%
\pgfpathlineto{\pgfqpoint{3.707507in}{1.738318in}}%
\pgfpathlineto{\pgfqpoint{3.752198in}{1.747878in}}%
\pgfpathlineto{\pgfqpoint{3.786047in}{1.712461in}}%
\pgfpathlineto{\pgfqpoint{3.819988in}{1.677084in}}%
\pgfpathlineto{\pgfqpoint{3.775360in}{1.672009in}}%
\pgfpathlineto{\pgfqpoint{3.730329in}{1.653046in}}%
\pgfpathlineto{\pgfqpoint{3.696332in}{1.683797in}}%
\pgfpathlineto{\pgfqpoint{3.662429in}{1.714357in}}%
\pgfpathclose%
\pgfusepath{stroke,fill}%
\end{pgfscope}%
\begin{pgfscope}%
\pgfpathrectangle{\pgfqpoint{0.050000in}{0.050000in}}{\pgfqpoint{4.900000in}{4.900000in}}%
\pgfusepath{clip}%
\pgfsetbuttcap%
\pgfsetroundjoin%
\definecolor{currentfill}{rgb}{1.000000,1.000000,1.000000}%
\pgfsetfillcolor{currentfill}%
\pgfsetfillopacity{0.900000}%
\pgfsetlinewidth{0.200750pt}%
\definecolor{currentstroke}{rgb}{0.200000,0.200000,0.200000}%
\pgfsetstrokecolor{currentstroke}%
\pgfsetdash{}{0pt}%
\pgfpathmoveto{\pgfqpoint{2.416317in}{1.271960in}}%
\pgfpathlineto{\pgfqpoint{2.462576in}{1.287062in}}%
\pgfpathlineto{\pgfqpoint{2.508734in}{1.302174in}}%
\pgfpathlineto{\pgfqpoint{2.540934in}{1.280718in}}%
\pgfpathlineto{\pgfqpoint{2.573236in}{1.259198in}}%
\pgfpathlineto{\pgfqpoint{2.527004in}{1.243958in}}%
\pgfpathlineto{\pgfqpoint{2.480669in}{1.228745in}}%
\pgfpathlineto{\pgfqpoint{2.448442in}{1.250386in}}%
\pgfpathlineto{\pgfqpoint{2.416317in}{1.271960in}}%
\pgfpathclose%
\pgfusepath{stroke,fill}%
\end{pgfscope}%
\begin{pgfscope}%
\pgfpathrectangle{\pgfqpoint{0.050000in}{0.050000in}}{\pgfqpoint{4.900000in}{4.900000in}}%
\pgfusepath{clip}%
\pgfsetbuttcap%
\pgfsetroundjoin%
\definecolor{currentfill}{rgb}{1.000000,1.000000,1.000000}%
\pgfsetfillcolor{currentfill}%
\pgfsetfillopacity{0.900000}%
\pgfsetlinewidth{0.200750pt}%
\definecolor{currentstroke}{rgb}{0.200000,0.200000,0.200000}%
\pgfsetstrokecolor{currentstroke}%
\pgfsetdash{}{0pt}%
\pgfpathmoveto{\pgfqpoint{2.010026in}{1.267604in}}%
\pgfpathlineto{\pgfqpoint{2.056749in}{1.282683in}}%
\pgfpathlineto{\pgfqpoint{2.103369in}{1.297729in}}%
\pgfpathlineto{\pgfqpoint{2.134940in}{1.276231in}}%
\pgfpathlineto{\pgfqpoint{2.166610in}{1.254666in}}%
\pgfpathlineto{\pgfqpoint{2.119912in}{1.239524in}}%
\pgfpathlineto{\pgfqpoint{2.073111in}{1.224350in}}%
\pgfpathlineto{\pgfqpoint{2.041519in}{1.246011in}}%
\pgfpathlineto{\pgfqpoint{2.010026in}{1.267604in}}%
\pgfpathclose%
\pgfusepath{stroke,fill}%
\end{pgfscope}%
\begin{pgfscope}%
\pgfpathrectangle{\pgfqpoint{0.050000in}{0.050000in}}{\pgfqpoint{4.900000in}{4.900000in}}%
\pgfusepath{clip}%
\pgfsetbuttcap%
\pgfsetroundjoin%
\definecolor{currentfill}{rgb}{0.844860,0.838462,0.907236}%
\pgfsetfillcolor{currentfill}%
\pgfsetfillopacity{0.900000}%
\pgfsetlinewidth{0.200750pt}%
\definecolor{currentstroke}{rgb}{0.200000,0.200000,0.200000}%
\pgfsetstrokecolor{currentstroke}%
\pgfsetdash{}{0pt}%
\pgfpathmoveto{\pgfqpoint{3.320671in}{1.456693in}}%
\pgfpathlineto{\pgfqpoint{3.366487in}{1.510933in}}%
\pgfpathlineto{\pgfqpoint{3.412315in}{1.568271in}}%
\pgfpathlineto{\pgfqpoint{3.445982in}{1.547144in}}%
\pgfpathlineto{\pgfqpoint{3.479751in}{1.525665in}}%
\pgfpathlineto{\pgfqpoint{3.433841in}{1.470630in}}%
\pgfpathlineto{\pgfqpoint{3.387928in}{1.417762in}}%
\pgfpathlineto{\pgfqpoint{3.354246in}{1.437359in}}%
\pgfpathlineto{\pgfqpoint{3.320671in}{1.456693in}}%
\pgfpathclose%
\pgfusepath{stroke,fill}%
\end{pgfscope}%
\begin{pgfscope}%
\pgfpathrectangle{\pgfqpoint{0.050000in}{0.050000in}}{\pgfqpoint{4.900000in}{4.900000in}}%
\pgfusepath{clip}%
\pgfsetbuttcap%
\pgfsetroundjoin%
\definecolor{currentfill}{rgb}{0.880661,0.875740,0.928643}%
\pgfsetfillcolor{currentfill}%
\pgfsetfillopacity{0.900000}%
\pgfsetlinewidth{0.200750pt}%
\definecolor{currentstroke}{rgb}{0.200000,0.200000,0.200000}%
\pgfsetstrokecolor{currentstroke}%
\pgfsetdash{}{0pt}%
\pgfpathmoveto{\pgfqpoint{4.044353in}{1.494824in}}%
\pgfpathlineto{\pgfqpoint{4.087565in}{1.459973in}}%
\pgfpathlineto{\pgfqpoint{4.130462in}{1.420218in}}%
\pgfpathlineto{\pgfqpoint{4.164819in}{1.384255in}}%
\pgfpathlineto{\pgfqpoint{4.121916in}{1.423250in}}%
\pgfpathlineto{\pgfqpoint{4.078715in}{1.458013in}}%
\pgfpathlineto{\pgfqpoint{4.044353in}{1.494824in}}%
\pgfpathclose%
\pgfusepath{stroke,fill}%
\end{pgfscope}%
\begin{pgfscope}%
\pgfpathrectangle{\pgfqpoint{0.050000in}{0.050000in}}{\pgfqpoint{4.900000in}{4.900000in}}%
\pgfusepath{clip}%
\pgfsetbuttcap%
\pgfsetroundjoin%
\definecolor{currentfill}{rgb}{1.000000,1.000000,1.000000}%
\pgfsetfillcolor{currentfill}%
\pgfsetfillopacity{0.900000}%
\pgfsetlinewidth{0.200750pt}%
\definecolor{currentstroke}{rgb}{0.200000,0.200000,0.200000}%
\pgfsetstrokecolor{currentstroke}%
\pgfsetdash{}{0pt}%
\pgfpathmoveto{\pgfqpoint{1.603477in}{1.263264in}}%
\pgfpathlineto{\pgfqpoint{1.650665in}{1.278353in}}%
\pgfpathlineto{\pgfqpoint{1.697748in}{1.293408in}}%
\pgfpathlineto{\pgfqpoint{1.728687in}{1.271897in}}%
\pgfpathlineto{\pgfqpoint{1.759725in}{1.250317in}}%
\pgfpathlineto{\pgfqpoint{1.712562in}{1.235166in}}%
\pgfpathlineto{\pgfqpoint{1.665294in}{1.219982in}}%
\pgfpathlineto{\pgfqpoint{1.634337in}{1.241657in}}%
\pgfpathlineto{\pgfqpoint{1.603477in}{1.263264in}}%
\pgfpathclose%
\pgfusepath{stroke,fill}%
\end{pgfscope}%
\begin{pgfscope}%
\pgfpathrectangle{\pgfqpoint{0.050000in}{0.050000in}}{\pgfqpoint{4.900000in}{4.900000in}}%
\pgfusepath{clip}%
\pgfsetbuttcap%
\pgfsetroundjoin%
\definecolor{currentfill}{rgb}{1.000000,1.000000,1.000000}%
\pgfsetfillcolor{currentfill}%
\pgfsetfillopacity{0.900000}%
\pgfsetlinewidth{0.200750pt}%
\definecolor{currentstroke}{rgb}{0.200000,0.200000,0.200000}%
\pgfsetstrokecolor{currentstroke}%
\pgfsetdash{}{0pt}%
\pgfpathmoveto{\pgfqpoint{2.573236in}{1.259198in}}%
\pgfpathlineto{\pgfqpoint{2.619366in}{1.274530in}}%
\pgfpathlineto{\pgfqpoint{2.665395in}{1.290069in}}%
\pgfpathlineto{\pgfqpoint{2.697874in}{1.268683in}}%
\pgfpathlineto{\pgfqpoint{2.730456in}{1.247243in}}%
\pgfpathlineto{\pgfqpoint{2.684352in}{1.231473in}}%
\pgfpathlineto{\pgfqpoint{2.638147in}{1.215968in}}%
\pgfpathlineto{\pgfqpoint{2.605640in}{1.237615in}}%
\pgfpathlineto{\pgfqpoint{2.573236in}{1.259198in}}%
\pgfpathclose%
\pgfusepath{stroke,fill}%
\end{pgfscope}%
\begin{pgfscope}%
\pgfpathrectangle{\pgfqpoint{0.050000in}{0.050000in}}{\pgfqpoint{4.900000in}{4.900000in}}%
\pgfusepath{clip}%
\pgfsetbuttcap%
\pgfsetroundjoin%
\definecolor{currentfill}{rgb}{0.982099,0.981361,0.989296}%
\pgfsetfillcolor{currentfill}%
\pgfsetfillopacity{0.900000}%
\pgfsetlinewidth{0.200750pt}%
\definecolor{currentstroke}{rgb}{0.200000,0.200000,0.200000}%
\pgfsetstrokecolor{currentstroke}%
\pgfsetdash{}{0pt}%
\pgfpathmoveto{\pgfqpoint{2.979715in}{1.278399in}}%
\pgfpathlineto{\pgfqpoint{3.025467in}{1.302397in}}%
\pgfpathlineto{\pgfqpoint{3.071168in}{1.330194in}}%
\pgfpathlineto{\pgfqpoint{3.104316in}{1.310580in}}%
\pgfpathlineto{\pgfqpoint{3.137571in}{1.290891in}}%
\pgfpathlineto{\pgfqpoint{3.091774in}{1.262186in}}%
\pgfpathlineto{\pgfqpoint{3.045930in}{1.237330in}}%
\pgfpathlineto{\pgfqpoint{3.012770in}{1.257887in}}%
\pgfpathlineto{\pgfqpoint{2.979715in}{1.278399in}}%
\pgfpathclose%
\pgfusepath{stroke,fill}%
\end{pgfscope}%
\begin{pgfscope}%
\pgfpathrectangle{\pgfqpoint{0.050000in}{0.050000in}}{\pgfqpoint{4.900000in}{4.900000in}}%
\pgfusepath{clip}%
\pgfsetbuttcap%
\pgfsetroundjoin%
\definecolor{currentfill}{rgb}{1.000000,1.000000,1.000000}%
\pgfsetfillcolor{currentfill}%
\pgfsetfillopacity{0.900000}%
\pgfsetlinewidth{0.200750pt}%
\definecolor{currentstroke}{rgb}{0.200000,0.200000,0.200000}%
\pgfsetstrokecolor{currentstroke}%
\pgfsetdash{}{0pt}%
\pgfpathmoveto{\pgfqpoint{2.166610in}{1.254666in}}%
\pgfpathlineto{\pgfqpoint{2.213204in}{1.269774in}}%
\pgfpathlineto{\pgfqpoint{2.259695in}{1.284849in}}%
\pgfpathlineto{\pgfqpoint{2.291542in}{1.263311in}}%
\pgfpathlineto{\pgfqpoint{2.323490in}{1.241704in}}%
\pgfpathlineto{\pgfqpoint{2.276923in}{1.226533in}}%
\pgfpathlineto{\pgfqpoint{2.230251in}{1.211329in}}%
\pgfpathlineto{\pgfqpoint{2.198380in}{1.233032in}}%
\pgfpathlineto{\pgfqpoint{2.166610in}{1.254666in}}%
\pgfpathclose%
\pgfusepath{stroke,fill}%
\end{pgfscope}%
\begin{pgfscope}%
\pgfpathrectangle{\pgfqpoint{0.050000in}{0.050000in}}{\pgfqpoint{4.900000in}{4.900000in}}%
\pgfusepath{clip}%
\pgfsetbuttcap%
\pgfsetroundjoin%
\definecolor{currentfill}{rgb}{0.904529,0.900592,0.942914}%
\pgfsetfillcolor{currentfill}%
\pgfsetfillopacity{0.900000}%
\pgfsetlinewidth{0.200750pt}%
\definecolor{currentstroke}{rgb}{0.200000,0.200000,0.200000}%
\pgfsetstrokecolor{currentstroke}%
\pgfsetdash{}{0pt}%
\pgfpathmoveto{\pgfqpoint{3.229105in}{1.362864in}}%
\pgfpathlineto{\pgfqpoint{3.274877in}{1.407071in}}%
\pgfpathlineto{\pgfqpoint{3.320671in}{1.456693in}}%
\pgfpathlineto{\pgfqpoint{3.354246in}{1.437359in}}%
\pgfpathlineto{\pgfqpoint{3.387928in}{1.417762in}}%
\pgfpathlineto{\pgfqpoint{3.342029in}{1.368711in}}%
\pgfpathlineto{\pgfqpoint{3.296149in}{1.324453in}}%
\pgfpathlineto{\pgfqpoint{3.262573in}{1.343738in}}%
\pgfpathlineto{\pgfqpoint{3.229105in}{1.362864in}}%
\pgfpathclose%
\pgfusepath{stroke,fill}%
\end{pgfscope}%
\begin{pgfscope}%
\pgfpathrectangle{\pgfqpoint{0.050000in}{0.050000in}}{\pgfqpoint{4.900000in}{4.900000in}}%
\pgfusepath{clip}%
\pgfsetbuttcap%
\pgfsetroundjoin%
\definecolor{currentfill}{rgb}{1.000000,1.000000,1.000000}%
\pgfsetfillcolor{currentfill}%
\pgfsetfillopacity{0.900000}%
\pgfsetlinewidth{0.200750pt}%
\definecolor{currentstroke}{rgb}{0.200000,0.200000,0.200000}%
\pgfsetstrokecolor{currentstroke}%
\pgfsetdash{}{0pt}%
\pgfpathmoveto{\pgfqpoint{1.759725in}{1.250317in}}%
\pgfpathlineto{\pgfqpoint{1.806784in}{1.265435in}}%
\pgfpathlineto{\pgfqpoint{1.853739in}{1.280519in}}%
\pgfpathlineto{\pgfqpoint{1.884955in}{1.258966in}}%
\pgfpathlineto{\pgfqpoint{1.916269in}{1.237346in}}%
\pgfpathlineto{\pgfqpoint{1.869235in}{1.222166in}}%
\pgfpathlineto{\pgfqpoint{1.822096in}{1.206953in}}%
\pgfpathlineto{\pgfqpoint{1.790861in}{1.228669in}}%
\pgfpathlineto{\pgfqpoint{1.759725in}{1.250317in}}%
\pgfpathclose%
\pgfusepath{stroke,fill}%
\end{pgfscope}%
\begin{pgfscope}%
\pgfpathrectangle{\pgfqpoint{0.050000in}{0.050000in}}{\pgfqpoint{4.900000in}{4.900000in}}%
\pgfusepath{clip}%
\pgfsetbuttcap%
\pgfsetroundjoin%
\definecolor{currentfill}{rgb}{0.731488,0.720415,0.839446}%
\pgfsetfillcolor{currentfill}%
\pgfsetfillopacity{0.900000}%
\pgfsetlinewidth{0.200750pt}%
\definecolor{currentstroke}{rgb}{0.200000,0.200000,0.200000}%
\pgfsetstrokecolor{currentstroke}%
\pgfsetdash{}{0pt}%
\pgfpathmoveto{\pgfqpoint{3.571408in}{1.632432in}}%
\pgfpathlineto{\pgfqpoint{3.617035in}{1.678106in}}%
\pgfpathlineto{\pgfqpoint{3.662429in}{1.714357in}}%
\pgfpathlineto{\pgfqpoint{3.696332in}{1.683797in}}%
\pgfpathlineto{\pgfqpoint{3.730329in}{1.653046in}}%
\pgfpathlineto{\pgfqpoint{3.684961in}{1.621874in}}%
\pgfpathlineto{\pgfqpoint{3.639333in}{1.580967in}}%
\pgfpathlineto{\pgfqpoint{3.605322in}{1.606861in}}%
\pgfpathlineto{\pgfqpoint{3.571408in}{1.632432in}}%
\pgfpathclose%
\pgfusepath{stroke,fill}%
\end{pgfscope}%
\begin{pgfscope}%
\pgfpathrectangle{\pgfqpoint{0.050000in}{0.050000in}}{\pgfqpoint{4.900000in}{4.900000in}}%
\pgfusepath{clip}%
\pgfsetbuttcap%
\pgfsetroundjoin%
\definecolor{currentfill}{rgb}{1.000000,1.000000,1.000000}%
\pgfsetfillcolor{currentfill}%
\pgfsetfillopacity{0.900000}%
\pgfsetlinewidth{0.200750pt}%
\definecolor{currentstroke}{rgb}{0.200000,0.200000,0.200000}%
\pgfsetstrokecolor{currentstroke}%
\pgfsetdash{}{0pt}%
\pgfpathmoveto{\pgfqpoint{1.352581in}{1.245966in}}%
\pgfpathlineto{\pgfqpoint{1.400106in}{1.261093in}}%
\pgfpathlineto{\pgfqpoint{1.447525in}{1.276187in}}%
\pgfpathlineto{\pgfqpoint{1.478109in}{1.254621in}}%
\pgfpathlineto{\pgfqpoint{1.508789in}{1.232987in}}%
\pgfpathlineto{\pgfqpoint{1.461287in}{1.217797in}}%
\pgfpathlineto{\pgfqpoint{1.413680in}{1.202574in}}%
\pgfpathlineto{\pgfqpoint{1.383082in}{1.224305in}}%
\pgfpathlineto{\pgfqpoint{1.352581in}{1.245966in}}%
\pgfpathclose%
\pgfusepath{stroke,fill}%
\end{pgfscope}%
\begin{pgfscope}%
\pgfpathrectangle{\pgfqpoint{0.050000in}{0.050000in}}{\pgfqpoint{4.900000in}{4.900000in}}%
\pgfusepath{clip}%
\pgfsetbuttcap%
\pgfsetroundjoin%
\definecolor{currentfill}{rgb}{1.000000,1.000000,1.000000}%
\pgfsetfillcolor{currentfill}%
\pgfsetfillopacity{0.900000}%
\pgfsetlinewidth{0.200750pt}%
\definecolor{currentstroke}{rgb}{0.200000,0.200000,0.200000}%
\pgfsetstrokecolor{currentstroke}%
\pgfsetdash{}{0pt}%
\pgfpathmoveto{\pgfqpoint{2.730456in}{1.247243in}}%
\pgfpathlineto{\pgfqpoint{2.776460in}{1.263497in}}%
\pgfpathlineto{\pgfqpoint{2.822369in}{1.280575in}}%
\pgfpathlineto{\pgfqpoint{2.855131in}{1.259472in}}%
\pgfpathlineto{\pgfqpoint{2.887997in}{1.238325in}}%
\pgfpathlineto{\pgfqpoint{2.842011in}{1.220791in}}%
\pgfpathlineto{\pgfqpoint{2.795930in}{1.204197in}}%
\pgfpathlineto{\pgfqpoint{2.763141in}{1.225747in}}%
\pgfpathlineto{\pgfqpoint{2.730456in}{1.247243in}}%
\pgfpathclose%
\pgfusepath{stroke,fill}%
\end{pgfscope}%
\begin{pgfscope}%
\pgfpathrectangle{\pgfqpoint{0.050000in}{0.050000in}}{\pgfqpoint{4.900000in}{4.900000in}}%
\pgfusepath{clip}%
\pgfsetbuttcap%
\pgfsetroundjoin%
\definecolor{currentfill}{rgb}{1.000000,1.000000,1.000000}%
\pgfsetfillcolor{currentfill}%
\pgfsetfillopacity{0.900000}%
\pgfsetlinewidth{0.200750pt}%
\definecolor{currentstroke}{rgb}{0.200000,0.200000,0.200000}%
\pgfsetstrokecolor{currentstroke}%
\pgfsetdash{}{0pt}%
\pgfpathmoveto{\pgfqpoint{2.323490in}{1.241704in}}%
\pgfpathlineto{\pgfqpoint{2.369955in}{1.256845in}}%
\pgfpathlineto{\pgfqpoint{2.416317in}{1.271960in}}%
\pgfpathlineto{\pgfqpoint{2.448442in}{1.250386in}}%
\pgfpathlineto{\pgfqpoint{2.480669in}{1.228745in}}%
\pgfpathlineto{\pgfqpoint{2.434231in}{1.213527in}}%
\pgfpathlineto{\pgfqpoint{2.387691in}{1.198287in}}%
\pgfpathlineto{\pgfqpoint{2.355540in}{1.220030in}}%
\pgfpathlineto{\pgfqpoint{2.323490in}{1.241704in}}%
\pgfpathclose%
\pgfusepath{stroke,fill}%
\end{pgfscope}%
\begin{pgfscope}%
\pgfpathrectangle{\pgfqpoint{0.050000in}{0.050000in}}{\pgfqpoint{4.900000in}{4.900000in}}%
\pgfusepath{clip}%
\pgfsetbuttcap%
\pgfsetroundjoin%
\definecolor{currentfill}{rgb}{0.785190,0.776332,0.871557}%
\pgfsetfillcolor{currentfill}%
\pgfsetfillopacity{0.900000}%
\pgfsetlinewidth{0.200750pt}%
\definecolor{currentstroke}{rgb}{0.200000,0.200000,0.200000}%
\pgfsetstrokecolor{currentstroke}%
\pgfsetdash{}{0pt}%
\pgfpathmoveto{\pgfqpoint{3.888150in}{1.606483in}}%
\pgfpathlineto{\pgfqpoint{3.932264in}{1.594051in}}%
\pgfpathlineto{\pgfqpoint{3.975923in}{1.569210in}}%
\pgfpathlineto{\pgfqpoint{4.010089in}{1.531887in}}%
\pgfpathlineto{\pgfqpoint{4.044353in}{1.494824in}}%
\pgfpathlineto{\pgfqpoint{4.000743in}{1.521069in}}%
\pgfpathlineto{\pgfqpoint{3.956690in}{1.536124in}}%
\pgfpathlineto{\pgfqpoint{3.922372in}{1.571271in}}%
\pgfpathlineto{\pgfqpoint{3.888150in}{1.606483in}}%
\pgfpathclose%
\pgfusepath{stroke,fill}%
\end{pgfscope}%
\begin{pgfscope}%
\pgfpathrectangle{\pgfqpoint{0.050000in}{0.050000in}}{\pgfqpoint{4.900000in}{4.900000in}}%
\pgfusepath{clip}%
\pgfsetbuttcap%
\pgfsetroundjoin%
\definecolor{currentfill}{rgb}{1.000000,1.000000,1.000000}%
\pgfsetfillcolor{currentfill}%
\pgfsetfillopacity{0.900000}%
\pgfsetlinewidth{0.200750pt}%
\definecolor{currentstroke}{rgb}{0.200000,0.200000,0.200000}%
\pgfsetstrokecolor{currentstroke}%
\pgfsetdash{}{0pt}%
\pgfpathmoveto{\pgfqpoint{1.916269in}{1.237346in}}%
\pgfpathlineto{\pgfqpoint{1.963200in}{1.252492in}}%
\pgfpathlineto{\pgfqpoint{2.010026in}{1.267604in}}%
\pgfpathlineto{\pgfqpoint{2.041519in}{1.246011in}}%
\pgfpathlineto{\pgfqpoint{2.073111in}{1.224350in}}%
\pgfpathlineto{\pgfqpoint{2.026206in}{1.209141in}}%
\pgfpathlineto{\pgfqpoint{1.979197in}{1.193899in}}%
\pgfpathlineto{\pgfqpoint{1.947683in}{1.215657in}}%
\pgfpathlineto{\pgfqpoint{1.916269in}{1.237346in}}%
\pgfpathclose%
\pgfusepath{stroke,fill}%
\end{pgfscope}%
\begin{pgfscope}%
\pgfpathrectangle{\pgfqpoint{0.050000in}{0.050000in}}{\pgfqpoint{4.900000in}{4.900000in}}%
\pgfusepath{clip}%
\pgfsetbuttcap%
\pgfsetroundjoin%
\definecolor{currentfill}{rgb}{0.731488,0.720415,0.839446}%
\pgfsetfillcolor{currentfill}%
\pgfsetfillopacity{0.900000}%
\pgfsetlinewidth{0.200750pt}%
\definecolor{currentstroke}{rgb}{0.200000,0.200000,0.200000}%
\pgfsetstrokecolor{currentstroke}%
\pgfsetdash{}{0pt}%
\pgfpathmoveto{\pgfqpoint{3.730329in}{1.653046in}}%
\pgfpathlineto{\pgfqpoint{3.775360in}{1.672009in}}%
\pgfpathlineto{\pgfqpoint{3.819988in}{1.677084in}}%
\pgfpathlineto{\pgfqpoint{3.854022in}{1.641756in}}%
\pgfpathlineto{\pgfqpoint{3.888150in}{1.606483in}}%
\pgfpathlineto{\pgfqpoint{3.843585in}{1.605417in}}%
\pgfpathlineto{\pgfqpoint{3.798605in}{1.591036in}}%
\pgfpathlineto{\pgfqpoint{3.764420in}{1.622121in}}%
\pgfpathlineto{\pgfqpoint{3.730329in}{1.653046in}}%
\pgfpathclose%
\pgfusepath{stroke,fill}%
\end{pgfscope}%
\begin{pgfscope}%
\pgfpathrectangle{\pgfqpoint{0.050000in}{0.050000in}}{\pgfqpoint{4.900000in}{4.900000in}}%
\pgfusepath{clip}%
\pgfsetbuttcap%
\pgfsetroundjoin%
\definecolor{currentfill}{rgb}{0.952265,0.950296,0.971457}%
\pgfsetfillcolor{currentfill}%
\pgfsetfillopacity{0.900000}%
\pgfsetlinewidth{0.200750pt}%
\definecolor{currentstroke}{rgb}{0.200000,0.200000,0.200000}%
\pgfsetstrokecolor{currentstroke}%
\pgfsetdash{}{0pt}%
\pgfpathmoveto{\pgfqpoint{3.137571in}{1.290891in}}%
\pgfpathlineto{\pgfqpoint{3.183342in}{1.324241in}}%
\pgfpathlineto{\pgfqpoint{3.229105in}{1.362864in}}%
\pgfpathlineto{\pgfqpoint{3.262573in}{1.343738in}}%
\pgfpathlineto{\pgfqpoint{3.296149in}{1.324453in}}%
\pgfpathlineto{\pgfqpoint{3.250278in}{1.285350in}}%
\pgfpathlineto{\pgfqpoint{3.204404in}{1.251272in}}%
\pgfpathlineto{\pgfqpoint{3.170934in}{1.271124in}}%
\pgfpathlineto{\pgfqpoint{3.137571in}{1.290891in}}%
\pgfpathclose%
\pgfusepath{stroke,fill}%
\end{pgfscope}%
\begin{pgfscope}%
\pgfpathrectangle{\pgfqpoint{0.050000in}{0.050000in}}{\pgfqpoint{4.900000in}{4.900000in}}%
\pgfusepath{clip}%
\pgfsetbuttcap%
\pgfsetroundjoin%
\definecolor{currentfill}{rgb}{0.779223,0.770119,0.867989}%
\pgfsetfillcolor{currentfill}%
\pgfsetfillopacity{0.900000}%
\pgfsetlinewidth{0.200750pt}%
\definecolor{currentstroke}{rgb}{0.200000,0.200000,0.200000}%
\pgfsetstrokecolor{currentstroke}%
\pgfsetdash{}{0pt}%
\pgfpathmoveto{\pgfqpoint{3.479751in}{1.525665in}}%
\pgfpathlineto{\pgfqpoint{3.525624in}{1.580553in}}%
\pgfpathlineto{\pgfqpoint{3.571408in}{1.632432in}}%
\pgfpathlineto{\pgfqpoint{3.605322in}{1.606861in}}%
\pgfpathlineto{\pgfqpoint{3.639333in}{1.580967in}}%
\pgfpathlineto{\pgfqpoint{3.593521in}{1.533237in}}%
\pgfpathlineto{\pgfqpoint{3.547596in}{1.481685in}}%
\pgfpathlineto{\pgfqpoint{3.513622in}{1.503842in}}%
\pgfpathlineto{\pgfqpoint{3.479751in}{1.525665in}}%
\pgfpathclose%
\pgfusepath{stroke,fill}%
\end{pgfscope}%
\begin{pgfscope}%
\pgfpathrectangle{\pgfqpoint{0.050000in}{0.050000in}}{\pgfqpoint{4.900000in}{4.900000in}}%
\pgfusepath{clip}%
\pgfsetbuttcap%
\pgfsetroundjoin%
\definecolor{currentfill}{rgb}{1.000000,1.000000,1.000000}%
\pgfsetfillcolor{currentfill}%
\pgfsetfillopacity{0.900000}%
\pgfsetlinewidth{0.200750pt}%
\definecolor{currentstroke}{rgb}{0.200000,0.200000,0.200000}%
\pgfsetstrokecolor{currentstroke}%
\pgfsetdash{}{0pt}%
\pgfpathmoveto{\pgfqpoint{1.508789in}{1.232987in}}%
\pgfpathlineto{\pgfqpoint{1.556186in}{1.248142in}}%
\pgfpathlineto{\pgfqpoint{1.603477in}{1.263264in}}%
\pgfpathlineto{\pgfqpoint{1.634337in}{1.241657in}}%
\pgfpathlineto{\pgfqpoint{1.665294in}{1.219982in}}%
\pgfpathlineto{\pgfqpoint{1.617921in}{1.204764in}}%
\pgfpathlineto{\pgfqpoint{1.570442in}{1.189512in}}%
\pgfpathlineto{\pgfqpoint{1.539567in}{1.211284in}}%
\pgfpathlineto{\pgfqpoint{1.508789in}{1.232987in}}%
\pgfpathclose%
\pgfusepath{stroke,fill}%
\end{pgfscope}%
\begin{pgfscope}%
\pgfpathrectangle{\pgfqpoint{0.050000in}{0.050000in}}{\pgfqpoint{4.900000in}{4.900000in}}%
\pgfusepath{clip}%
\pgfsetbuttcap%
\pgfsetroundjoin%
\definecolor{currentfill}{rgb}{0.994033,0.993787,0.996432}%
\pgfsetfillcolor{currentfill}%
\pgfsetfillopacity{0.900000}%
\pgfsetlinewidth{0.200750pt}%
\definecolor{currentstroke}{rgb}{0.200000,0.200000,0.200000}%
\pgfsetstrokecolor{currentstroke}%
\pgfsetdash{}{0pt}%
\pgfpathmoveto{\pgfqpoint{2.887997in}{1.238325in}}%
\pgfpathlineto{\pgfqpoint{2.933896in}{1.257308in}}%
\pgfpathlineto{\pgfqpoint{2.979715in}{1.278399in}}%
\pgfpathlineto{\pgfqpoint{3.012770in}{1.257887in}}%
\pgfpathlineto{\pgfqpoint{3.045930in}{1.237330in}}%
\pgfpathlineto{\pgfqpoint{3.000025in}{1.215492in}}%
\pgfpathlineto{\pgfqpoint{2.954045in}{1.195899in}}%
\pgfpathlineto{\pgfqpoint{2.920969in}{1.217134in}}%
\pgfpathlineto{\pgfqpoint{2.887997in}{1.238325in}}%
\pgfpathclose%
\pgfusepath{stroke,fill}%
\end{pgfscope}%
\begin{pgfscope}%
\pgfpathrectangle{\pgfqpoint{0.050000in}{0.050000in}}{\pgfqpoint{4.900000in}{4.900000in}}%
\pgfusepath{clip}%
\pgfsetbuttcap%
\pgfsetroundjoin%
\definecolor{currentfill}{rgb}{1.000000,1.000000,1.000000}%
\pgfsetfillcolor{currentfill}%
\pgfsetfillopacity{0.900000}%
\pgfsetlinewidth{0.200750pt}%
\definecolor{currentstroke}{rgb}{0.200000,0.200000,0.200000}%
\pgfsetstrokecolor{currentstroke}%
\pgfsetdash{}{0pt}%
\pgfpathmoveto{\pgfqpoint{2.480669in}{1.228745in}}%
\pgfpathlineto{\pgfqpoint{2.527004in}{1.243958in}}%
\pgfpathlineto{\pgfqpoint{2.573236in}{1.259198in}}%
\pgfpathlineto{\pgfqpoint{2.605640in}{1.237615in}}%
\pgfpathlineto{\pgfqpoint{2.638147in}{1.215968in}}%
\pgfpathlineto{\pgfqpoint{2.591840in}{1.200590in}}%
\pgfpathlineto{\pgfqpoint{2.545430in}{1.185260in}}%
\pgfpathlineto{\pgfqpoint{2.512998in}{1.207036in}}%
\pgfpathlineto{\pgfqpoint{2.480669in}{1.228745in}}%
\pgfpathclose%
\pgfusepath{stroke,fill}%
\end{pgfscope}%
\begin{pgfscope}%
\pgfpathrectangle{\pgfqpoint{0.050000in}{0.050000in}}{\pgfqpoint{4.900000in}{4.900000in}}%
\pgfusepath{clip}%
\pgfsetbuttcap%
\pgfsetroundjoin%
\definecolor{currentfill}{rgb}{1.000000,1.000000,1.000000}%
\pgfsetfillcolor{currentfill}%
\pgfsetfillopacity{0.900000}%
\pgfsetlinewidth{0.200750pt}%
\definecolor{currentstroke}{rgb}{0.200000,0.200000,0.200000}%
\pgfsetstrokecolor{currentstroke}%
\pgfsetdash{}{0pt}%
\pgfpathmoveto{\pgfqpoint{2.073111in}{1.224350in}}%
\pgfpathlineto{\pgfqpoint{2.119912in}{1.239524in}}%
\pgfpathlineto{\pgfqpoint{2.166610in}{1.254666in}}%
\pgfpathlineto{\pgfqpoint{2.198380in}{1.233032in}}%
\pgfpathlineto{\pgfqpoint{2.230251in}{1.211329in}}%
\pgfpathlineto{\pgfqpoint{2.183476in}{1.196091in}}%
\pgfpathlineto{\pgfqpoint{2.136597in}{1.180820in}}%
\pgfpathlineto{\pgfqpoint{2.104804in}{1.202619in}}%
\pgfpathlineto{\pgfqpoint{2.073111in}{1.224350in}}%
\pgfpathclose%
\pgfusepath{stroke,fill}%
\end{pgfscope}%
\begin{pgfscope}%
\pgfpathrectangle{\pgfqpoint{0.050000in}{0.050000in}}{\pgfqpoint{4.900000in}{4.900000in}}%
\pgfusepath{clip}%
\pgfsetbuttcap%
\pgfsetroundjoin%
\definecolor{currentfill}{rgb}{0.844860,0.838462,0.907236}%
\pgfsetfillcolor{currentfill}%
\pgfsetfillopacity{0.900000}%
\pgfsetlinewidth{0.200750pt}%
\definecolor{currentstroke}{rgb}{0.200000,0.200000,0.200000}%
\pgfsetstrokecolor{currentstroke}%
\pgfsetdash{}{0pt}%
\pgfpathmoveto{\pgfqpoint{3.387928in}{1.417762in}}%
\pgfpathlineto{\pgfqpoint{3.433841in}{1.470630in}}%
\pgfpathlineto{\pgfqpoint{3.479751in}{1.525665in}}%
\pgfpathlineto{\pgfqpoint{3.513622in}{1.503842in}}%
\pgfpathlineto{\pgfqpoint{3.547596in}{1.481685in}}%
\pgfpathlineto{\pgfqpoint{3.501611in}{1.429088in}}%
\pgfpathlineto{\pgfqpoint{3.455608in}{1.377777in}}%
\pgfpathlineto{\pgfqpoint{3.421715in}{1.397901in}}%
\pgfpathlineto{\pgfqpoint{3.387928in}{1.417762in}}%
\pgfpathclose%
\pgfusepath{stroke,fill}%
\end{pgfscope}%
\begin{pgfscope}%
\pgfpathrectangle{\pgfqpoint{0.050000in}{0.050000in}}{\pgfqpoint{4.900000in}{4.900000in}}%
\pgfusepath{clip}%
\pgfsetbuttcap%
\pgfsetroundjoin%
\definecolor{currentfill}{rgb}{1.000000,1.000000,1.000000}%
\pgfsetfillcolor{currentfill}%
\pgfsetfillopacity{0.900000}%
\pgfsetlinewidth{0.200750pt}%
\definecolor{currentstroke}{rgb}{0.200000,0.200000,0.200000}%
\pgfsetstrokecolor{currentstroke}%
\pgfsetdash{}{0pt}%
\pgfpathmoveto{\pgfqpoint{1.665294in}{1.219982in}}%
\pgfpathlineto{\pgfqpoint{1.712562in}{1.235166in}}%
\pgfpathlineto{\pgfqpoint{1.759725in}{1.250317in}}%
\pgfpathlineto{\pgfqpoint{1.790861in}{1.228669in}}%
\pgfpathlineto{\pgfqpoint{1.822096in}{1.206953in}}%
\pgfpathlineto{\pgfqpoint{1.774852in}{1.191706in}}%
\pgfpathlineto{\pgfqpoint{1.727503in}{1.176425in}}%
\pgfpathlineto{\pgfqpoint{1.696349in}{1.198238in}}%
\pgfpathlineto{\pgfqpoint{1.665294in}{1.219982in}}%
\pgfpathclose%
\pgfusepath{stroke,fill}%
\end{pgfscope}%
\begin{pgfscope}%
\pgfpathrectangle{\pgfqpoint{0.050000in}{0.050000in}}{\pgfqpoint{4.900000in}{4.900000in}}%
\pgfusepath{clip}%
\pgfsetbuttcap%
\pgfsetroundjoin%
\definecolor{currentfill}{rgb}{1.000000,1.000000,1.000000}%
\pgfsetfillcolor{currentfill}%
\pgfsetfillopacity{0.900000}%
\pgfsetlinewidth{0.200750pt}%
\definecolor{currentstroke}{rgb}{0.200000,0.200000,0.200000}%
\pgfsetstrokecolor{currentstroke}%
\pgfsetdash{}{0pt}%
\pgfpathmoveto{\pgfqpoint{2.638147in}{1.215968in}}%
\pgfpathlineto{\pgfqpoint{2.684352in}{1.231473in}}%
\pgfpathlineto{\pgfqpoint{2.730456in}{1.247243in}}%
\pgfpathlineto{\pgfqpoint{2.763141in}{1.225747in}}%
\pgfpathlineto{\pgfqpoint{2.795930in}{1.204197in}}%
\pgfpathlineto{\pgfqpoint{2.749751in}{1.188178in}}%
\pgfpathlineto{\pgfqpoint{2.703471in}{1.172485in}}%
\pgfpathlineto{\pgfqpoint{2.670757in}{1.194258in}}%
\pgfpathlineto{\pgfqpoint{2.638147in}{1.215968in}}%
\pgfpathclose%
\pgfusepath{stroke,fill}%
\end{pgfscope}%
\begin{pgfscope}%
\pgfpathrectangle{\pgfqpoint{0.050000in}{0.050000in}}{\pgfqpoint{4.900000in}{4.900000in}}%
\pgfusepath{clip}%
\pgfsetbuttcap%
\pgfsetroundjoin%
\definecolor{currentfill}{rgb}{0.976132,0.975148,0.985729}%
\pgfsetfillcolor{currentfill}%
\pgfsetfillopacity{0.900000}%
\pgfsetlinewidth{0.200750pt}%
\definecolor{currentstroke}{rgb}{0.200000,0.200000,0.200000}%
\pgfsetstrokecolor{currentstroke}%
\pgfsetdash{}{0pt}%
\pgfpathmoveto{\pgfqpoint{3.045930in}{1.237330in}}%
\pgfpathlineto{\pgfqpoint{3.091774in}{1.262186in}}%
\pgfpathlineto{\pgfqpoint{3.137571in}{1.290891in}}%
\pgfpathlineto{\pgfqpoint{3.170934in}{1.271124in}}%
\pgfpathlineto{\pgfqpoint{3.204404in}{1.251272in}}%
\pgfpathlineto{\pgfqpoint{3.158509in}{1.221739in}}%
\pgfpathlineto{\pgfqpoint{3.112573in}{1.196068in}}%
\pgfpathlineto{\pgfqpoint{3.079198in}{1.216725in}}%
\pgfpathlineto{\pgfqpoint{3.045930in}{1.237330in}}%
\pgfpathclose%
\pgfusepath{stroke,fill}%
\end{pgfscope}%
\begin{pgfscope}%
\pgfpathrectangle{\pgfqpoint{0.050000in}{0.050000in}}{\pgfqpoint{4.900000in}{4.900000in}}%
\pgfusepath{clip}%
\pgfsetbuttcap%
\pgfsetroundjoin%
\definecolor{currentfill}{rgb}{0.803091,0.794971,0.882261}%
\pgfsetfillcolor{currentfill}%
\pgfsetfillopacity{0.900000}%
\pgfsetlinewidth{0.200750pt}%
\definecolor{currentstroke}{rgb}{0.200000,0.200000,0.200000}%
\pgfsetstrokecolor{currentstroke}%
\pgfsetdash{}{0pt}%
\pgfpathmoveto{\pgfqpoint{3.956690in}{1.536124in}}%
\pgfpathlineto{\pgfqpoint{4.000743in}{1.521069in}}%
\pgfpathlineto{\pgfqpoint{4.044353in}{1.494824in}}%
\pgfpathlineto{\pgfqpoint{4.078715in}{1.458013in}}%
\pgfpathlineto{\pgfqpoint{4.035128in}{1.484829in}}%
\pgfpathlineto{\pgfqpoint{3.991105in}{1.501047in}}%
\pgfpathlineto{\pgfqpoint{3.956690in}{1.536124in}}%
\pgfpathclose%
\pgfusepath{stroke,fill}%
\end{pgfscope}%
\begin{pgfscope}%
\pgfpathrectangle{\pgfqpoint{0.050000in}{0.050000in}}{\pgfqpoint{4.900000in}{4.900000in}}%
\pgfusepath{clip}%
\pgfsetbuttcap%
\pgfsetroundjoin%
\definecolor{currentfill}{rgb}{1.000000,1.000000,1.000000}%
\pgfsetfillcolor{currentfill}%
\pgfsetfillopacity{0.900000}%
\pgfsetlinewidth{0.200750pt}%
\definecolor{currentstroke}{rgb}{0.200000,0.200000,0.200000}%
\pgfsetstrokecolor{currentstroke}%
\pgfsetdash{}{0pt}%
\pgfpathmoveto{\pgfqpoint{2.230251in}{1.211329in}}%
\pgfpathlineto{\pgfqpoint{2.276923in}{1.226533in}}%
\pgfpathlineto{\pgfqpoint{2.323490in}{1.241704in}}%
\pgfpathlineto{\pgfqpoint{2.355540in}{1.220030in}}%
\pgfpathlineto{\pgfqpoint{2.387691in}{1.198287in}}%
\pgfpathlineto{\pgfqpoint{2.341046in}{1.183018in}}%
\pgfpathlineto{\pgfqpoint{2.294298in}{1.167716in}}%
\pgfpathlineto{\pgfqpoint{2.262224in}{1.189557in}}%
\pgfpathlineto{\pgfqpoint{2.230251in}{1.211329in}}%
\pgfpathclose%
\pgfusepath{stroke,fill}%
\end{pgfscope}%
\begin{pgfscope}%
\pgfpathrectangle{\pgfqpoint{0.050000in}{0.050000in}}{\pgfqpoint{4.900000in}{4.900000in}}%
\pgfusepath{clip}%
\pgfsetbuttcap%
\pgfsetroundjoin%
\definecolor{currentfill}{rgb}{0.743422,0.732841,0.846582}%
\pgfsetfillcolor{currentfill}%
\pgfsetfillopacity{0.900000}%
\pgfsetlinewidth{0.200750pt}%
\definecolor{currentstroke}{rgb}{0.200000,0.200000,0.200000}%
\pgfsetstrokecolor{currentstroke}%
\pgfsetdash{}{0pt}%
\pgfpathmoveto{\pgfqpoint{3.639333in}{1.580967in}}%
\pgfpathlineto{\pgfqpoint{3.684961in}{1.621874in}}%
\pgfpathlineto{\pgfqpoint{3.730329in}{1.653046in}}%
\pgfpathlineto{\pgfqpoint{3.764420in}{1.622121in}}%
\pgfpathlineto{\pgfqpoint{3.798605in}{1.591036in}}%
\pgfpathlineto{\pgfqpoint{3.753269in}{1.564617in}}%
\pgfpathlineto{\pgfqpoint{3.707648in}{1.528266in}}%
\pgfpathlineto{\pgfqpoint{3.673442in}{1.554764in}}%
\pgfpathlineto{\pgfqpoint{3.639333in}{1.580967in}}%
\pgfpathclose%
\pgfusepath{stroke,fill}%
\end{pgfscope}%
\begin{pgfscope}%
\pgfpathrectangle{\pgfqpoint{0.050000in}{0.050000in}}{\pgfqpoint{4.900000in}{4.900000in}}%
\pgfusepath{clip}%
\pgfsetbuttcap%
\pgfsetroundjoin%
\definecolor{currentfill}{rgb}{0.904529,0.900592,0.942914}%
\pgfsetfillcolor{currentfill}%
\pgfsetfillopacity{0.900000}%
\pgfsetlinewidth{0.200750pt}%
\definecolor{currentstroke}{rgb}{0.200000,0.200000,0.200000}%
\pgfsetstrokecolor{currentstroke}%
\pgfsetdash{}{0pt}%
\pgfpathmoveto{\pgfqpoint{3.296149in}{1.324453in}}%
\pgfpathlineto{\pgfqpoint{3.342029in}{1.368711in}}%
\pgfpathlineto{\pgfqpoint{3.387928in}{1.417762in}}%
\pgfpathlineto{\pgfqpoint{3.421715in}{1.397901in}}%
\pgfpathlineto{\pgfqpoint{3.455608in}{1.377777in}}%
\pgfpathlineto{\pgfqpoint{3.409610in}{1.329502in}}%
\pgfpathlineto{\pgfqpoint{3.363624in}{1.285389in}}%
\pgfpathlineto{\pgfqpoint{3.329832in}{1.305005in}}%
\pgfpathlineto{\pgfqpoint{3.296149in}{1.324453in}}%
\pgfpathclose%
\pgfusepath{stroke,fill}%
\end{pgfscope}%
\begin{pgfscope}%
\pgfpathrectangle{\pgfqpoint{0.050000in}{0.050000in}}{\pgfqpoint{4.900000in}{4.900000in}}%
\pgfusepath{clip}%
\pgfsetbuttcap%
\pgfsetroundjoin%
\definecolor{currentfill}{rgb}{1.000000,1.000000,1.000000}%
\pgfsetfillcolor{currentfill}%
\pgfsetfillopacity{0.900000}%
\pgfsetlinewidth{0.200750pt}%
\definecolor{currentstroke}{rgb}{0.200000,0.200000,0.200000}%
\pgfsetstrokecolor{currentstroke}%
\pgfsetdash{}{0pt}%
\pgfpathmoveto{\pgfqpoint{1.822096in}{1.206953in}}%
\pgfpathlineto{\pgfqpoint{1.869235in}{1.222166in}}%
\pgfpathlineto{\pgfqpoint{1.916269in}{1.237346in}}%
\pgfpathlineto{\pgfqpoint{1.947683in}{1.215657in}}%
\pgfpathlineto{\pgfqpoint{1.979197in}{1.193899in}}%
\pgfpathlineto{\pgfqpoint{1.932083in}{1.178623in}}%
\pgfpathlineto{\pgfqpoint{1.884864in}{1.163312in}}%
\pgfpathlineto{\pgfqpoint{1.853430in}{1.185167in}}%
\pgfpathlineto{\pgfqpoint{1.822096in}{1.206953in}}%
\pgfpathclose%
\pgfusepath{stroke,fill}%
\end{pgfscope}%
\begin{pgfscope}%
\pgfpathrectangle{\pgfqpoint{0.050000in}{0.050000in}}{\pgfqpoint{4.900000in}{4.900000in}}%
\pgfusepath{clip}%
\pgfsetbuttcap%
\pgfsetroundjoin%
\definecolor{currentfill}{rgb}{1.000000,1.000000,1.000000}%
\pgfsetfillcolor{currentfill}%
\pgfsetfillopacity{0.900000}%
\pgfsetlinewidth{0.200750pt}%
\definecolor{currentstroke}{rgb}{0.200000,0.200000,0.200000}%
\pgfsetstrokecolor{currentstroke}%
\pgfsetdash{}{0pt}%
\pgfpathmoveto{\pgfqpoint{1.413680in}{1.202574in}}%
\pgfpathlineto{\pgfqpoint{1.461287in}{1.217797in}}%
\pgfpathlineto{\pgfqpoint{1.508789in}{1.232987in}}%
\pgfpathlineto{\pgfqpoint{1.539567in}{1.211284in}}%
\pgfpathlineto{\pgfqpoint{1.570442in}{1.189512in}}%
\pgfpathlineto{\pgfqpoint{1.522858in}{1.174226in}}%
\pgfpathlineto{\pgfqpoint{1.475167in}{1.158906in}}%
\pgfpathlineto{\pgfqpoint{1.444375in}{1.180775in}}%
\pgfpathlineto{\pgfqpoint{1.413680in}{1.202574in}}%
\pgfpathclose%
\pgfusepath{stroke,fill}%
\end{pgfscope}%
\begin{pgfscope}%
\pgfpathrectangle{\pgfqpoint{0.050000in}{0.050000in}}{\pgfqpoint{4.900000in}{4.900000in}}%
\pgfusepath{clip}%
\pgfsetbuttcap%
\pgfsetroundjoin%
\definecolor{currentfill}{rgb}{1.000000,1.000000,1.000000}%
\pgfsetfillcolor{currentfill}%
\pgfsetfillopacity{0.900000}%
\pgfsetlinewidth{0.200750pt}%
\definecolor{currentstroke}{rgb}{0.200000,0.200000,0.200000}%
\pgfsetstrokecolor{currentstroke}%
\pgfsetdash{}{0pt}%
\pgfpathmoveto{\pgfqpoint{2.795930in}{1.204197in}}%
\pgfpathlineto{\pgfqpoint{2.842011in}{1.220791in}}%
\pgfpathlineto{\pgfqpoint{2.887997in}{1.238325in}}%
\pgfpathlineto{\pgfqpoint{2.920969in}{1.217134in}}%
\pgfpathlineto{\pgfqpoint{2.954045in}{1.195899in}}%
\pgfpathlineto{\pgfqpoint{2.907980in}{1.177888in}}%
\pgfpathlineto{\pgfqpoint{2.861822in}{1.160933in}}%
\pgfpathlineto{\pgfqpoint{2.828824in}{1.182593in}}%
\pgfpathlineto{\pgfqpoint{2.795930in}{1.204197in}}%
\pgfpathclose%
\pgfusepath{stroke,fill}%
\end{pgfscope}%
\begin{pgfscope}%
\pgfpathrectangle{\pgfqpoint{0.050000in}{0.050000in}}{\pgfqpoint{4.900000in}{4.900000in}}%
\pgfusepath{clip}%
\pgfsetbuttcap%
\pgfsetroundjoin%
\definecolor{currentfill}{rgb}{1.000000,1.000000,1.000000}%
\pgfsetfillcolor{currentfill}%
\pgfsetfillopacity{0.900000}%
\pgfsetlinewidth{0.200750pt}%
\definecolor{currentstroke}{rgb}{0.200000,0.200000,0.200000}%
\pgfsetstrokecolor{currentstroke}%
\pgfsetdash{}{0pt}%
\pgfpathmoveto{\pgfqpoint{2.387691in}{1.198287in}}%
\pgfpathlineto{\pgfqpoint{2.434231in}{1.213527in}}%
\pgfpathlineto{\pgfqpoint{2.480669in}{1.228745in}}%
\pgfpathlineto{\pgfqpoint{2.512998in}{1.207036in}}%
\pgfpathlineto{\pgfqpoint{2.545430in}{1.185260in}}%
\pgfpathlineto{\pgfqpoint{2.498917in}{1.169935in}}%
\pgfpathlineto{\pgfqpoint{2.452300in}{1.154594in}}%
\pgfpathlineto{\pgfqpoint{2.419944in}{1.176475in}}%
\pgfpathlineto{\pgfqpoint{2.387691in}{1.198287in}}%
\pgfpathclose%
\pgfusepath{stroke,fill}%
\end{pgfscope}%
\begin{pgfscope}%
\pgfpathrectangle{\pgfqpoint{0.050000in}{0.050000in}}{\pgfqpoint{4.900000in}{4.900000in}}%
\pgfusepath{clip}%
\pgfsetbuttcap%
\pgfsetroundjoin%
\definecolor{currentfill}{rgb}{0.749389,0.739054,0.850150}%
\pgfsetfillcolor{currentfill}%
\pgfsetfillopacity{0.900000}%
\pgfsetlinewidth{0.200750pt}%
\definecolor{currentstroke}{rgb}{0.200000,0.200000,0.200000}%
\pgfsetstrokecolor{currentstroke}%
\pgfsetdash{}{0pt}%
\pgfpathmoveto{\pgfqpoint{3.798605in}{1.591036in}}%
\pgfpathlineto{\pgfqpoint{3.843585in}{1.605417in}}%
\pgfpathlineto{\pgfqpoint{3.888150in}{1.606483in}}%
\pgfpathlineto{\pgfqpoint{3.922372in}{1.571271in}}%
\pgfpathlineto{\pgfqpoint{3.956690in}{1.536124in}}%
\pgfpathlineto{\pgfqpoint{3.912188in}{1.538631in}}%
\pgfpathlineto{\pgfqpoint{3.867262in}{1.528436in}}%
\pgfpathlineto{\pgfqpoint{3.832886in}{1.559803in}}%
\pgfpathlineto{\pgfqpoint{3.798605in}{1.591036in}}%
\pgfpathclose%
\pgfusepath{stroke,fill}%
\end{pgfscope}%
\begin{pgfscope}%
\pgfpathrectangle{\pgfqpoint{0.050000in}{0.050000in}}{\pgfqpoint{4.900000in}{4.900000in}}%
\pgfusepath{clip}%
\pgfsetbuttcap%
\pgfsetroundjoin%
\definecolor{currentfill}{rgb}{1.000000,1.000000,1.000000}%
\pgfsetfillcolor{currentfill}%
\pgfsetfillopacity{0.900000}%
\pgfsetlinewidth{0.200750pt}%
\definecolor{currentstroke}{rgb}{0.200000,0.200000,0.200000}%
\pgfsetstrokecolor{currentstroke}%
\pgfsetdash{}{0pt}%
\pgfpathmoveto{\pgfqpoint{1.979197in}{1.193899in}}%
\pgfpathlineto{\pgfqpoint{2.026206in}{1.209141in}}%
\pgfpathlineto{\pgfqpoint{2.073111in}{1.224350in}}%
\pgfpathlineto{\pgfqpoint{2.104804in}{1.202619in}}%
\pgfpathlineto{\pgfqpoint{2.136597in}{1.180820in}}%
\pgfpathlineto{\pgfqpoint{2.089614in}{1.165515in}}%
\pgfpathlineto{\pgfqpoint{2.042525in}{1.150175in}}%
\pgfpathlineto{\pgfqpoint{2.010811in}{1.172072in}}%
\pgfpathlineto{\pgfqpoint{1.979197in}{1.193899in}}%
\pgfpathclose%
\pgfusepath{stroke,fill}%
\end{pgfscope}%
\begin{pgfscope}%
\pgfpathrectangle{\pgfqpoint{0.050000in}{0.050000in}}{\pgfqpoint{4.900000in}{4.900000in}}%
\pgfusepath{clip}%
\pgfsetbuttcap%
\pgfsetroundjoin%
\definecolor{currentfill}{rgb}{0.785190,0.776332,0.871557}%
\pgfsetfillcolor{currentfill}%
\pgfsetfillopacity{0.900000}%
\pgfsetlinewidth{0.200750pt}%
\definecolor{currentstroke}{rgb}{0.200000,0.200000,0.200000}%
\pgfsetstrokecolor{currentstroke}%
\pgfsetdash{}{0pt}%
\pgfpathmoveto{\pgfqpoint{3.547596in}{1.481685in}}%
\pgfpathlineto{\pgfqpoint{3.593521in}{1.533237in}}%
\pgfpathlineto{\pgfqpoint{3.639333in}{1.580967in}}%
\pgfpathlineto{\pgfqpoint{3.673442in}{1.554764in}}%
\pgfpathlineto{\pgfqpoint{3.707648in}{1.528266in}}%
\pgfpathlineto{\pgfqpoint{3.661819in}{1.484597in}}%
\pgfpathlineto{\pgfqpoint{3.615849in}{1.436397in}}%
\pgfpathlineto{\pgfqpoint{3.581671in}{1.459200in}}%
\pgfpathlineto{\pgfqpoint{3.547596in}{1.481685in}}%
\pgfpathclose%
\pgfusepath{stroke,fill}%
\end{pgfscope}%
\begin{pgfscope}%
\pgfpathrectangle{\pgfqpoint{0.050000in}{0.050000in}}{\pgfqpoint{4.900000in}{4.900000in}}%
\pgfusepath{clip}%
\pgfsetbuttcap%
\pgfsetroundjoin%
\definecolor{currentfill}{rgb}{0.946298,0.944083,0.967889}%
\pgfsetfillcolor{currentfill}%
\pgfsetfillopacity{0.900000}%
\pgfsetlinewidth{0.200750pt}%
\definecolor{currentstroke}{rgb}{0.200000,0.200000,0.200000}%
\pgfsetstrokecolor{currentstroke}%
\pgfsetdash{}{0pt}%
\pgfpathmoveto{\pgfqpoint{3.204404in}{1.251272in}}%
\pgfpathlineto{\pgfqpoint{3.250278in}{1.285350in}}%
\pgfpathlineto{\pgfqpoint{3.296149in}{1.324453in}}%
\pgfpathlineto{\pgfqpoint{3.329832in}{1.305005in}}%
\pgfpathlineto{\pgfqpoint{3.363624in}{1.285389in}}%
\pgfpathlineto{\pgfqpoint{3.317648in}{1.245974in}}%
\pgfpathlineto{\pgfqpoint{3.271671in}{1.211297in}}%
\pgfpathlineto{\pgfqpoint{3.237983in}{1.231331in}}%
\pgfpathlineto{\pgfqpoint{3.204404in}{1.251272in}}%
\pgfpathclose%
\pgfusepath{stroke,fill}%
\end{pgfscope}%
\begin{pgfscope}%
\pgfpathrectangle{\pgfqpoint{0.050000in}{0.050000in}}{\pgfqpoint{4.900000in}{4.900000in}}%
\pgfusepath{clip}%
\pgfsetbuttcap%
\pgfsetroundjoin%
\definecolor{currentfill}{rgb}{1.000000,1.000000,1.000000}%
\pgfsetfillcolor{currentfill}%
\pgfsetfillopacity{0.900000}%
\pgfsetlinewidth{0.200750pt}%
\definecolor{currentstroke}{rgb}{0.200000,0.200000,0.200000}%
\pgfsetstrokecolor{currentstroke}%
\pgfsetdash{}{0pt}%
\pgfpathmoveto{\pgfqpoint{1.570442in}{1.189512in}}%
\pgfpathlineto{\pgfqpoint{1.617921in}{1.204764in}}%
\pgfpathlineto{\pgfqpoint{1.665294in}{1.219982in}}%
\pgfpathlineto{\pgfqpoint{1.696349in}{1.198238in}}%
\pgfpathlineto{\pgfqpoint{1.727503in}{1.176425in}}%
\pgfpathlineto{\pgfqpoint{1.680049in}{1.161109in}}%
\pgfpathlineto{\pgfqpoint{1.632488in}{1.145760in}}%
\pgfpathlineto{\pgfqpoint{1.601416in}{1.167671in}}%
\pgfpathlineto{\pgfqpoint{1.570442in}{1.189512in}}%
\pgfpathclose%
\pgfusepath{stroke,fill}%
\end{pgfscope}%
\begin{pgfscope}%
\pgfpathrectangle{\pgfqpoint{0.050000in}{0.050000in}}{\pgfqpoint{4.900000in}{4.900000in}}%
\pgfusepath{clip}%
\pgfsetbuttcap%
\pgfsetroundjoin%
\definecolor{currentfill}{rgb}{0.994033,0.993787,0.996432}%
\pgfsetfillcolor{currentfill}%
\pgfsetfillopacity{0.900000}%
\pgfsetlinewidth{0.200750pt}%
\definecolor{currentstroke}{rgb}{0.200000,0.200000,0.200000}%
\pgfsetstrokecolor{currentstroke}%
\pgfsetdash{}{0pt}%
\pgfpathmoveto{\pgfqpoint{2.954045in}{1.195899in}}%
\pgfpathlineto{\pgfqpoint{3.000025in}{1.215492in}}%
\pgfpathlineto{\pgfqpoint{3.045930in}{1.237330in}}%
\pgfpathlineto{\pgfqpoint{3.079198in}{1.216725in}}%
\pgfpathlineto{\pgfqpoint{3.112573in}{1.196068in}}%
\pgfpathlineto{\pgfqpoint{3.066581in}{1.173498in}}%
\pgfpathlineto{\pgfqpoint{3.020517in}{1.153288in}}%
\pgfpathlineto{\pgfqpoint{2.987228in}{1.174617in}}%
\pgfpathlineto{\pgfqpoint{2.954045in}{1.195899in}}%
\pgfpathclose%
\pgfusepath{stroke,fill}%
\end{pgfscope}%
\begin{pgfscope}%
\pgfpathrectangle{\pgfqpoint{0.050000in}{0.050000in}}{\pgfqpoint{4.900000in}{4.900000in}}%
\pgfusepath{clip}%
\pgfsetbuttcap%
\pgfsetroundjoin%
\definecolor{currentfill}{rgb}{1.000000,1.000000,1.000000}%
\pgfsetfillcolor{currentfill}%
\pgfsetfillopacity{0.900000}%
\pgfsetlinewidth{0.200750pt}%
\definecolor{currentstroke}{rgb}{0.200000,0.200000,0.200000}%
\pgfsetstrokecolor{currentstroke}%
\pgfsetdash{}{0pt}%
\pgfpathmoveto{\pgfqpoint{2.545430in}{1.185260in}}%
\pgfpathlineto{\pgfqpoint{2.591840in}{1.200590in}}%
\pgfpathlineto{\pgfqpoint{2.638147in}{1.215968in}}%
\pgfpathlineto{\pgfqpoint{2.670757in}{1.194258in}}%
\pgfpathlineto{\pgfqpoint{2.703471in}{1.172485in}}%
\pgfpathlineto{\pgfqpoint{2.657089in}{1.156958in}}%
\pgfpathlineto{\pgfqpoint{2.610604in}{1.141504in}}%
\pgfpathlineto{\pgfqpoint{2.577965in}{1.163416in}}%
\pgfpathlineto{\pgfqpoint{2.545430in}{1.185260in}}%
\pgfpathclose%
\pgfusepath{stroke,fill}%
\end{pgfscope}%
\begin{pgfscope}%
\pgfpathrectangle{\pgfqpoint{0.050000in}{0.050000in}}{\pgfqpoint{4.900000in}{4.900000in}}%
\pgfusepath{clip}%
\pgfsetbuttcap%
\pgfsetroundjoin%
\definecolor{currentfill}{rgb}{1.000000,1.000000,1.000000}%
\pgfsetfillcolor{currentfill}%
\pgfsetfillopacity{0.900000}%
\pgfsetlinewidth{0.200750pt}%
\definecolor{currentstroke}{rgb}{0.200000,0.200000,0.200000}%
\pgfsetstrokecolor{currentstroke}%
\pgfsetdash{}{0pt}%
\pgfpathmoveto{\pgfqpoint{2.136597in}{1.180820in}}%
\pgfpathlineto{\pgfqpoint{2.183476in}{1.196091in}}%
\pgfpathlineto{\pgfqpoint{2.230251in}{1.211329in}}%
\pgfpathlineto{\pgfqpoint{2.262224in}{1.189557in}}%
\pgfpathlineto{\pgfqpoint{2.294298in}{1.167716in}}%
\pgfpathlineto{\pgfqpoint{2.247446in}{1.152382in}}%
\pgfpathlineto{\pgfqpoint{2.200488in}{1.137013in}}%
\pgfpathlineto{\pgfqpoint{2.168492in}{1.158951in}}%
\pgfpathlineto{\pgfqpoint{2.136597in}{1.180820in}}%
\pgfpathclose%
\pgfusepath{stroke,fill}%
\end{pgfscope}%
\begin{pgfscope}%
\pgfpathrectangle{\pgfqpoint{0.050000in}{0.050000in}}{\pgfqpoint{4.900000in}{4.900000in}}%
\pgfusepath{clip}%
\pgfsetbuttcap%
\pgfsetroundjoin%
\definecolor{currentfill}{rgb}{0.844860,0.838462,0.907236}%
\pgfsetfillcolor{currentfill}%
\pgfsetfillopacity{0.900000}%
\pgfsetlinewidth{0.200750pt}%
\definecolor{currentstroke}{rgb}{0.200000,0.200000,0.200000}%
\pgfsetstrokecolor{currentstroke}%
\pgfsetdash{}{0pt}%
\pgfpathmoveto{\pgfqpoint{3.455608in}{1.377777in}}%
\pgfpathlineto{\pgfqpoint{3.501611in}{1.429088in}}%
\pgfpathlineto{\pgfqpoint{3.547596in}{1.481685in}}%
\pgfpathlineto{\pgfqpoint{3.581671in}{1.459200in}}%
\pgfpathlineto{\pgfqpoint{3.615849in}{1.436397in}}%
\pgfpathlineto{\pgfqpoint{3.569799in}{1.386338in}}%
\pgfpathlineto{\pgfqpoint{3.523714in}{1.336742in}}%
\pgfpathlineto{\pgfqpoint{3.489608in}{1.357390in}}%
\pgfpathlineto{\pgfqpoint{3.455608in}{1.377777in}}%
\pgfpathclose%
\pgfusepath{stroke,fill}%
\end{pgfscope}%
\begin{pgfscope}%
\pgfpathrectangle{\pgfqpoint{0.050000in}{0.050000in}}{\pgfqpoint{4.900000in}{4.900000in}}%
\pgfusepath{clip}%
\pgfsetbuttcap%
\pgfsetroundjoin%
\definecolor{currentfill}{rgb}{1.000000,1.000000,1.000000}%
\pgfsetfillcolor{currentfill}%
\pgfsetfillopacity{0.900000}%
\pgfsetlinewidth{0.200750pt}%
\definecolor{currentstroke}{rgb}{0.200000,0.200000,0.200000}%
\pgfsetstrokecolor{currentstroke}%
\pgfsetdash{}{0pt}%
\pgfpathmoveto{\pgfqpoint{1.727503in}{1.176425in}}%
\pgfpathlineto{\pgfqpoint{1.774852in}{1.191706in}}%
\pgfpathlineto{\pgfqpoint{1.822096in}{1.206953in}}%
\pgfpathlineto{\pgfqpoint{1.853430in}{1.185167in}}%
\pgfpathlineto{\pgfqpoint{1.884864in}{1.163312in}}%
\pgfpathlineto{\pgfqpoint{1.837540in}{1.147968in}}%
\pgfpathlineto{\pgfqpoint{1.790110in}{1.132589in}}%
\pgfpathlineto{\pgfqpoint{1.758756in}{1.154542in}}%
\pgfpathlineto{\pgfqpoint{1.727503in}{1.176425in}}%
\pgfpathclose%
\pgfusepath{stroke,fill}%
\end{pgfscope}%
\begin{pgfscope}%
\pgfpathrectangle{\pgfqpoint{0.050000in}{0.050000in}}{\pgfqpoint{4.900000in}{4.900000in}}%
\pgfusepath{clip}%
\pgfsetbuttcap%
\pgfsetroundjoin%
\definecolor{currentfill}{rgb}{1.000000,1.000000,1.000000}%
\pgfsetfillcolor{currentfill}%
\pgfsetfillopacity{0.900000}%
\pgfsetlinewidth{0.200750pt}%
\definecolor{currentstroke}{rgb}{0.200000,0.200000,0.200000}%
\pgfsetstrokecolor{currentstroke}%
\pgfsetdash{}{0pt}%
\pgfpathmoveto{\pgfqpoint{2.703471in}{1.172485in}}%
\pgfpathlineto{\pgfqpoint{2.749751in}{1.188178in}}%
\pgfpathlineto{\pgfqpoint{2.795930in}{1.204197in}}%
\pgfpathlineto{\pgfqpoint{2.828824in}{1.182593in}}%
\pgfpathlineto{\pgfqpoint{2.861822in}{1.160933in}}%
\pgfpathlineto{\pgfqpoint{2.815567in}{1.144645in}}%
\pgfpathlineto{\pgfqpoint{2.769212in}{1.128747in}}%
\pgfpathlineto{\pgfqpoint{2.736289in}{1.150648in}}%
\pgfpathlineto{\pgfqpoint{2.703471in}{1.172485in}}%
\pgfpathclose%
\pgfusepath{stroke,fill}%
\end{pgfscope}%
\begin{pgfscope}%
\pgfpathrectangle{\pgfqpoint{0.050000in}{0.050000in}}{\pgfqpoint{4.900000in}{4.900000in}}%
\pgfusepath{clip}%
\pgfsetbuttcap%
\pgfsetroundjoin%
\definecolor{currentfill}{rgb}{0.767290,0.757693,0.860854}%
\pgfsetfillcolor{currentfill}%
\pgfsetfillopacity{0.900000}%
\pgfsetlinewidth{0.200750pt}%
\definecolor{currentstroke}{rgb}{0.200000,0.200000,0.200000}%
\pgfsetstrokecolor{currentstroke}%
\pgfsetdash{}{0pt}%
\pgfpathmoveto{\pgfqpoint{3.867262in}{1.528436in}}%
\pgfpathlineto{\pgfqpoint{3.912188in}{1.538631in}}%
\pgfpathlineto{\pgfqpoint{3.956690in}{1.536124in}}%
\pgfpathlineto{\pgfqpoint{3.991105in}{1.501047in}}%
\pgfpathlineto{\pgfqpoint{3.946633in}{1.505189in}}%
\pgfpathlineto{\pgfqpoint{3.901735in}{1.496946in}}%
\pgfpathlineto{\pgfqpoint{3.867262in}{1.528436in}}%
\pgfpathclose%
\pgfusepath{stroke,fill}%
\end{pgfscope}%
\begin{pgfscope}%
\pgfpathrectangle{\pgfqpoint{0.050000in}{0.050000in}}{\pgfqpoint{4.900000in}{4.900000in}}%
\pgfusepath{clip}%
\pgfsetbuttcap%
\pgfsetroundjoin%
\definecolor{currentfill}{rgb}{0.755356,0.745267,0.853718}%
\pgfsetfillcolor{currentfill}%
\pgfsetfillopacity{0.900000}%
\pgfsetlinewidth{0.200750pt}%
\definecolor{currentstroke}{rgb}{0.200000,0.200000,0.200000}%
\pgfsetstrokecolor{currentstroke}%
\pgfsetdash{}{0pt}%
\pgfpathmoveto{\pgfqpoint{3.707648in}{1.528266in}}%
\pgfpathlineto{\pgfqpoint{3.753269in}{1.564617in}}%
\pgfpathlineto{\pgfqpoint{3.798605in}{1.591036in}}%
\pgfpathlineto{\pgfqpoint{3.832886in}{1.559803in}}%
\pgfpathlineto{\pgfqpoint{3.867262in}{1.528436in}}%
\pgfpathlineto{\pgfqpoint{3.821963in}{1.506454in}}%
\pgfpathlineto{\pgfqpoint{3.776357in}{1.474444in}}%
\pgfpathlineto{\pgfqpoint{3.741954in}{1.501489in}}%
\pgfpathlineto{\pgfqpoint{3.707648in}{1.528266in}}%
\pgfpathclose%
\pgfusepath{stroke,fill}%
\end{pgfscope}%
\begin{pgfscope}%
\pgfpathrectangle{\pgfqpoint{0.050000in}{0.050000in}}{\pgfqpoint{4.900000in}{4.900000in}}%
\pgfusepath{clip}%
\pgfsetbuttcap%
\pgfsetroundjoin%
\definecolor{currentfill}{rgb}{0.976132,0.975148,0.985729}%
\pgfsetfillcolor{currentfill}%
\pgfsetfillopacity{0.900000}%
\pgfsetlinewidth{0.200750pt}%
\definecolor{currentstroke}{rgb}{0.200000,0.200000,0.200000}%
\pgfsetstrokecolor{currentstroke}%
\pgfsetdash{}{0pt}%
\pgfpathmoveto{\pgfqpoint{3.112573in}{1.196068in}}%
\pgfpathlineto{\pgfqpoint{3.158509in}{1.221739in}}%
\pgfpathlineto{\pgfqpoint{3.204404in}{1.251272in}}%
\pgfpathlineto{\pgfqpoint{3.237983in}{1.231331in}}%
\pgfpathlineto{\pgfqpoint{3.271671in}{1.211297in}}%
\pgfpathlineto{\pgfqpoint{3.225676in}{1.181025in}}%
\pgfpathlineto{\pgfqpoint{3.179647in}{1.154591in}}%
\pgfpathlineto{\pgfqpoint{3.146056in}{1.175358in}}%
\pgfpathlineto{\pgfqpoint{3.112573in}{1.196068in}}%
\pgfpathclose%
\pgfusepath{stroke,fill}%
\end{pgfscope}%
\begin{pgfscope}%
\pgfpathrectangle{\pgfqpoint{0.050000in}{0.050000in}}{\pgfqpoint{4.900000in}{4.900000in}}%
\pgfusepath{clip}%
\pgfsetbuttcap%
\pgfsetroundjoin%
\definecolor{currentfill}{rgb}{1.000000,1.000000,1.000000}%
\pgfsetfillcolor{currentfill}%
\pgfsetfillopacity{0.900000}%
\pgfsetlinewidth{0.200750pt}%
\definecolor{currentstroke}{rgb}{0.200000,0.200000,0.200000}%
\pgfsetstrokecolor{currentstroke}%
\pgfsetdash{}{0pt}%
\pgfpathmoveto{\pgfqpoint{2.294298in}{1.167716in}}%
\pgfpathlineto{\pgfqpoint{2.341046in}{1.183018in}}%
\pgfpathlineto{\pgfqpoint{2.387691in}{1.198287in}}%
\pgfpathlineto{\pgfqpoint{2.419944in}{1.176475in}}%
\pgfpathlineto{\pgfqpoint{2.452300in}{1.154594in}}%
\pgfpathlineto{\pgfqpoint{2.405579in}{1.139225in}}%
\pgfpathlineto{\pgfqpoint{2.358754in}{1.123826in}}%
\pgfpathlineto{\pgfqpoint{2.326474in}{1.145806in}}%
\pgfpathlineto{\pgfqpoint{2.294298in}{1.167716in}}%
\pgfpathclose%
\pgfusepath{stroke,fill}%
\end{pgfscope}%
\begin{pgfscope}%
\pgfpathrectangle{\pgfqpoint{0.050000in}{0.050000in}}{\pgfqpoint{4.900000in}{4.900000in}}%
\pgfusepath{clip}%
\pgfsetbuttcap%
\pgfsetroundjoin%
\definecolor{currentfill}{rgb}{0.898562,0.894379,0.939346}%
\pgfsetfillcolor{currentfill}%
\pgfsetfillopacity{0.900000}%
\pgfsetlinewidth{0.200750pt}%
\definecolor{currentstroke}{rgb}{0.200000,0.200000,0.200000}%
\pgfsetstrokecolor{currentstroke}%
\pgfsetdash{}{0pt}%
\pgfpathmoveto{\pgfqpoint{3.363624in}{1.285389in}}%
\pgfpathlineto{\pgfqpoint{3.409610in}{1.329502in}}%
\pgfpathlineto{\pgfqpoint{3.455608in}{1.377777in}}%
\pgfpathlineto{\pgfqpoint{3.489608in}{1.357390in}}%
\pgfpathlineto{\pgfqpoint{3.523714in}{1.336742in}}%
\pgfpathlineto{\pgfqpoint{3.477621in}{1.289426in}}%
\pgfpathlineto{\pgfqpoint{3.431534in}{1.245640in}}%
\pgfpathlineto{\pgfqpoint{3.397524in}{1.265602in}}%
\pgfpathlineto{\pgfqpoint{3.363624in}{1.285389in}}%
\pgfpathclose%
\pgfusepath{stroke,fill}%
\end{pgfscope}%
\begin{pgfscope}%
\pgfpathrectangle{\pgfqpoint{0.050000in}{0.050000in}}{\pgfqpoint{4.900000in}{4.900000in}}%
\pgfusepath{clip}%
\pgfsetbuttcap%
\pgfsetroundjoin%
\definecolor{currentfill}{rgb}{1.000000,1.000000,1.000000}%
\pgfsetfillcolor{currentfill}%
\pgfsetfillopacity{0.900000}%
\pgfsetlinewidth{0.200750pt}%
\definecolor{currentstroke}{rgb}{0.200000,0.200000,0.200000}%
\pgfsetstrokecolor{currentstroke}%
\pgfsetdash{}{0pt}%
\pgfpathmoveto{\pgfqpoint{1.884864in}{1.163312in}}%
\pgfpathlineto{\pgfqpoint{1.932083in}{1.178623in}}%
\pgfpathlineto{\pgfqpoint{1.979197in}{1.193899in}}%
\pgfpathlineto{\pgfqpoint{2.010811in}{1.172072in}}%
\pgfpathlineto{\pgfqpoint{2.042525in}{1.150175in}}%
\pgfpathlineto{\pgfqpoint{1.995332in}{1.134801in}}%
\pgfpathlineto{\pgfqpoint{1.948033in}{1.119393in}}%
\pgfpathlineto{\pgfqpoint{1.916398in}{1.141388in}}%
\pgfpathlineto{\pgfqpoint{1.884864in}{1.163312in}}%
\pgfpathclose%
\pgfusepath{stroke,fill}%
\end{pgfscope}%
\begin{pgfscope}%
\pgfpathrectangle{\pgfqpoint{0.050000in}{0.050000in}}{\pgfqpoint{4.900000in}{4.900000in}}%
\pgfusepath{clip}%
\pgfsetbuttcap%
\pgfsetroundjoin%
\definecolor{currentfill}{rgb}{1.000000,1.000000,1.000000}%
\pgfsetfillcolor{currentfill}%
\pgfsetfillopacity{0.900000}%
\pgfsetlinewidth{0.200750pt}%
\definecolor{currentstroke}{rgb}{0.200000,0.200000,0.200000}%
\pgfsetstrokecolor{currentstroke}%
\pgfsetdash{}{0pt}%
\pgfpathmoveto{\pgfqpoint{1.475167in}{1.158906in}}%
\pgfpathlineto{\pgfqpoint{1.522858in}{1.174226in}}%
\pgfpathlineto{\pgfqpoint{1.570442in}{1.189512in}}%
\pgfpathlineto{\pgfqpoint{1.601416in}{1.167671in}}%
\pgfpathlineto{\pgfqpoint{1.632488in}{1.145760in}}%
\pgfpathlineto{\pgfqpoint{1.584821in}{1.130376in}}%
\pgfpathlineto{\pgfqpoint{1.537048in}{1.114958in}}%
\pgfpathlineto{\pgfqpoint{1.506058in}{1.136967in}}%
\pgfpathlineto{\pgfqpoint{1.475167in}{1.158906in}}%
\pgfpathclose%
\pgfusepath{stroke,fill}%
\end{pgfscope}%
\begin{pgfscope}%
\pgfpathrectangle{\pgfqpoint{0.050000in}{0.050000in}}{\pgfqpoint{4.900000in}{4.900000in}}%
\pgfusepath{clip}%
\pgfsetbuttcap%
\pgfsetroundjoin%
\definecolor{currentfill}{rgb}{1.000000,1.000000,1.000000}%
\pgfsetfillcolor{currentfill}%
\pgfsetfillopacity{0.900000}%
\pgfsetlinewidth{0.200750pt}%
\definecolor{currentstroke}{rgb}{0.200000,0.200000,0.200000}%
\pgfsetstrokecolor{currentstroke}%
\pgfsetdash{}{0pt}%
\pgfpathmoveto{\pgfqpoint{2.861822in}{1.160933in}}%
\pgfpathlineto{\pgfqpoint{2.907980in}{1.177888in}}%
\pgfpathlineto{\pgfqpoint{2.954045in}{1.195899in}}%
\pgfpathlineto{\pgfqpoint{2.987228in}{1.174617in}}%
\pgfpathlineto{\pgfqpoint{3.020517in}{1.153288in}}%
\pgfpathlineto{\pgfqpoint{2.974372in}{1.134785in}}%
\pgfpathlineto{\pgfqpoint{2.928137in}{1.117450in}}%
\pgfpathlineto{\pgfqpoint{2.894927in}{1.139219in}}%
\pgfpathlineto{\pgfqpoint{2.861822in}{1.160933in}}%
\pgfpathclose%
\pgfusepath{stroke,fill}%
\end{pgfscope}%
\begin{pgfscope}%
\pgfpathrectangle{\pgfqpoint{0.050000in}{0.050000in}}{\pgfqpoint{4.900000in}{4.900000in}}%
\pgfusepath{clip}%
\pgfsetbuttcap%
\pgfsetroundjoin%
\definecolor{currentfill}{rgb}{1.000000,1.000000,1.000000}%
\pgfsetfillcolor{currentfill}%
\pgfsetfillopacity{0.900000}%
\pgfsetlinewidth{0.200750pt}%
\definecolor{currentstroke}{rgb}{0.200000,0.200000,0.200000}%
\pgfsetstrokecolor{currentstroke}%
\pgfsetdash{}{0pt}%
\pgfpathmoveto{\pgfqpoint{2.452300in}{1.154594in}}%
\pgfpathlineto{\pgfqpoint{2.498917in}{1.169935in}}%
\pgfpathlineto{\pgfqpoint{2.545430in}{1.185260in}}%
\pgfpathlineto{\pgfqpoint{2.577965in}{1.163416in}}%
\pgfpathlineto{\pgfqpoint{2.610604in}{1.141504in}}%
\pgfpathlineto{\pgfqpoint{2.564015in}{1.126068in}}%
\pgfpathlineto{\pgfqpoint{2.517322in}{1.110623in}}%
\pgfpathlineto{\pgfqpoint{2.484759in}{1.132643in}}%
\pgfpathlineto{\pgfqpoint{2.452300in}{1.154594in}}%
\pgfpathclose%
\pgfusepath{stroke,fill}%
\end{pgfscope}%
\begin{pgfscope}%
\pgfpathrectangle{\pgfqpoint{0.050000in}{0.050000in}}{\pgfqpoint{4.900000in}{4.900000in}}%
\pgfusepath{clip}%
\pgfsetbuttcap%
\pgfsetroundjoin%
\definecolor{currentfill}{rgb}{1.000000,1.000000,1.000000}%
\pgfsetfillcolor{currentfill}%
\pgfsetfillopacity{0.900000}%
\pgfsetlinewidth{0.200750pt}%
\definecolor{currentstroke}{rgb}{0.200000,0.200000,0.200000}%
\pgfsetstrokecolor{currentstroke}%
\pgfsetdash{}{0pt}%
\pgfpathmoveto{\pgfqpoint{2.042525in}{1.150175in}}%
\pgfpathlineto{\pgfqpoint{2.089614in}{1.165515in}}%
\pgfpathlineto{\pgfqpoint{2.136597in}{1.180820in}}%
\pgfpathlineto{\pgfqpoint{2.168492in}{1.158951in}}%
\pgfpathlineto{\pgfqpoint{2.200488in}{1.137013in}}%
\pgfpathlineto{\pgfqpoint{2.153426in}{1.121609in}}%
\pgfpathlineto{\pgfqpoint{2.106259in}{1.106172in}}%
\pgfpathlineto{\pgfqpoint{2.074341in}{1.128208in}}%
\pgfpathlineto{\pgfqpoint{2.042525in}{1.150175in}}%
\pgfpathclose%
\pgfusepath{stroke,fill}%
\end{pgfscope}%
\begin{pgfscope}%
\pgfpathrectangle{\pgfqpoint{0.050000in}{0.050000in}}{\pgfqpoint{4.900000in}{4.900000in}}%
\pgfusepath{clip}%
\pgfsetbuttcap%
\pgfsetroundjoin%
\definecolor{currentfill}{rgb}{0.791157,0.782545,0.875125}%
\pgfsetfillcolor{currentfill}%
\pgfsetfillopacity{0.900000}%
\pgfsetlinewidth{0.200750pt}%
\definecolor{currentstroke}{rgb}{0.200000,0.200000,0.200000}%
\pgfsetstrokecolor{currentstroke}%
\pgfsetdash{}{0pt}%
\pgfpathmoveto{\pgfqpoint{3.615849in}{1.436397in}}%
\pgfpathlineto{\pgfqpoint{3.661819in}{1.484597in}}%
\pgfpathlineto{\pgfqpoint{3.707648in}{1.528266in}}%
\pgfpathlineto{\pgfqpoint{3.741954in}{1.501489in}}%
\pgfpathlineto{\pgfqpoint{3.776357in}{1.474444in}}%
\pgfpathlineto{\pgfqpoint{3.730517in}{1.434723in}}%
\pgfpathlineto{\pgfqpoint{3.684514in}{1.389866in}}%
\pgfpathlineto{\pgfqpoint{3.650130in}{1.413283in}}%
\pgfpathlineto{\pgfqpoint{3.615849in}{1.436397in}}%
\pgfpathclose%
\pgfusepath{stroke,fill}%
\end{pgfscope}%
\begin{pgfscope}%
\pgfpathrectangle{\pgfqpoint{0.050000in}{0.050000in}}{\pgfqpoint{4.900000in}{4.900000in}}%
\pgfusepath{clip}%
\pgfsetbuttcap%
\pgfsetroundjoin%
\definecolor{currentfill}{rgb}{0.946298,0.944083,0.967889}%
\pgfsetfillcolor{currentfill}%
\pgfsetfillopacity{0.900000}%
\pgfsetlinewidth{0.200750pt}%
\definecolor{currentstroke}{rgb}{0.200000,0.200000,0.200000}%
\pgfsetstrokecolor{currentstroke}%
\pgfsetdash{}{0pt}%
\pgfpathmoveto{\pgfqpoint{3.271671in}{1.211297in}}%
\pgfpathlineto{\pgfqpoint{3.317648in}{1.245974in}}%
\pgfpathlineto{\pgfqpoint{3.363624in}{1.285389in}}%
\pgfpathlineto{\pgfqpoint{3.397524in}{1.265602in}}%
\pgfpathlineto{\pgfqpoint{3.431534in}{1.245640in}}%
\pgfpathlineto{\pgfqpoint{3.385454in}{1.206075in}}%
\pgfpathlineto{\pgfqpoint{3.339374in}{1.170929in}}%
\pgfpathlineto{\pgfqpoint{3.305468in}{1.191164in}}%
\pgfpathlineto{\pgfqpoint{3.271671in}{1.211297in}}%
\pgfpathclose%
\pgfusepath{stroke,fill}%
\end{pgfscope}%
\begin{pgfscope}%
\pgfpathrectangle{\pgfqpoint{0.050000in}{0.050000in}}{\pgfqpoint{4.900000in}{4.900000in}}%
\pgfusepath{clip}%
\pgfsetbuttcap%
\pgfsetroundjoin%
\definecolor{currentfill}{rgb}{1.000000,1.000000,1.000000}%
\pgfsetfillcolor{currentfill}%
\pgfsetfillopacity{0.900000}%
\pgfsetlinewidth{0.200750pt}%
\definecolor{currentstroke}{rgb}{0.200000,0.200000,0.200000}%
\pgfsetstrokecolor{currentstroke}%
\pgfsetdash{}{0pt}%
\pgfpathmoveto{\pgfqpoint{1.632488in}{1.145760in}}%
\pgfpathlineto{\pgfqpoint{1.680049in}{1.161109in}}%
\pgfpathlineto{\pgfqpoint{1.727503in}{1.176425in}}%
\pgfpathlineto{\pgfqpoint{1.758756in}{1.154542in}}%
\pgfpathlineto{\pgfqpoint{1.790110in}{1.132589in}}%
\pgfpathlineto{\pgfqpoint{1.742573in}{1.117176in}}%
\pgfpathlineto{\pgfqpoint{1.694931in}{1.101728in}}%
\pgfpathlineto{\pgfqpoint{1.663660in}{1.123779in}}%
\pgfpathlineto{\pgfqpoint{1.632488in}{1.145760in}}%
\pgfpathclose%
\pgfusepath{stroke,fill}%
\end{pgfscope}%
\begin{pgfscope}%
\pgfpathrectangle{\pgfqpoint{0.050000in}{0.050000in}}{\pgfqpoint{4.900000in}{4.900000in}}%
\pgfusepath{clip}%
\pgfsetbuttcap%
\pgfsetroundjoin%
\definecolor{currentfill}{rgb}{0.994033,0.993787,0.996432}%
\pgfsetfillcolor{currentfill}%
\pgfsetfillopacity{0.900000}%
\pgfsetlinewidth{0.200750pt}%
\definecolor{currentstroke}{rgb}{0.200000,0.200000,0.200000}%
\pgfsetstrokecolor{currentstroke}%
\pgfsetdash{}{0pt}%
\pgfpathmoveto{\pgfqpoint{3.020517in}{1.153288in}}%
\pgfpathlineto{\pgfqpoint{3.066581in}{1.173498in}}%
\pgfpathlineto{\pgfqpoint{3.112573in}{1.196068in}}%
\pgfpathlineto{\pgfqpoint{3.146056in}{1.175358in}}%
\pgfpathlineto{\pgfqpoint{3.179647in}{1.154591in}}%
\pgfpathlineto{\pgfqpoint{3.133566in}{1.131309in}}%
\pgfpathlineto{\pgfqpoint{3.087418in}{1.110484in}}%
\pgfpathlineto{\pgfqpoint{3.053914in}{1.131911in}}%
\pgfpathlineto{\pgfqpoint{3.020517in}{1.153288in}}%
\pgfpathclose%
\pgfusepath{stroke,fill}%
\end{pgfscope}%
\begin{pgfscope}%
\pgfpathrectangle{\pgfqpoint{0.050000in}{0.050000in}}{\pgfqpoint{4.900000in}{4.900000in}}%
\pgfusepath{clip}%
\pgfsetbuttcap%
\pgfsetroundjoin%
\definecolor{currentfill}{rgb}{1.000000,1.000000,1.000000}%
\pgfsetfillcolor{currentfill}%
\pgfsetfillopacity{0.900000}%
\pgfsetlinewidth{0.200750pt}%
\definecolor{currentstroke}{rgb}{0.200000,0.200000,0.200000}%
\pgfsetstrokecolor{currentstroke}%
\pgfsetdash{}{0pt}%
\pgfpathmoveto{\pgfqpoint{2.610604in}{1.141504in}}%
\pgfpathlineto{\pgfqpoint{2.657089in}{1.156958in}}%
\pgfpathlineto{\pgfqpoint{2.703471in}{1.172485in}}%
\pgfpathlineto{\pgfqpoint{2.736289in}{1.150648in}}%
\pgfpathlineto{\pgfqpoint{2.769212in}{1.128747in}}%
\pgfpathlineto{\pgfqpoint{2.722755in}{1.113062in}}%
\pgfpathlineto{\pgfqpoint{2.676195in}{1.097476in}}%
\pgfpathlineto{\pgfqpoint{2.643347in}{1.119524in}}%
\pgfpathlineto{\pgfqpoint{2.610604in}{1.141504in}}%
\pgfpathclose%
\pgfusepath{stroke,fill}%
\end{pgfscope}%
\begin{pgfscope}%
\pgfpathrectangle{\pgfqpoint{0.050000in}{0.050000in}}{\pgfqpoint{4.900000in}{4.900000in}}%
\pgfusepath{clip}%
\pgfsetbuttcap%
\pgfsetroundjoin%
\definecolor{currentfill}{rgb}{0.767290,0.757693,0.860854}%
\pgfsetfillcolor{currentfill}%
\pgfsetfillopacity{0.900000}%
\pgfsetlinewidth{0.200750pt}%
\definecolor{currentstroke}{rgb}{0.200000,0.200000,0.200000}%
\pgfsetstrokecolor{currentstroke}%
\pgfsetdash{}{0pt}%
\pgfpathmoveto{\pgfqpoint{3.776357in}{1.474444in}}%
\pgfpathlineto{\pgfqpoint{3.821963in}{1.506454in}}%
\pgfpathlineto{\pgfqpoint{3.867262in}{1.528436in}}%
\pgfpathlineto{\pgfqpoint{3.901735in}{1.496946in}}%
\pgfpathlineto{\pgfqpoint{3.856456in}{1.477066in}}%
\pgfpathlineto{\pgfqpoint{3.810860in}{1.447144in}}%
\pgfpathlineto{\pgfqpoint{3.776357in}{1.474444in}}%
\pgfpathclose%
\pgfusepath{stroke,fill}%
\end{pgfscope}%
\begin{pgfscope}%
\pgfpathrectangle{\pgfqpoint{0.050000in}{0.050000in}}{\pgfqpoint{4.900000in}{4.900000in}}%
\pgfusepath{clip}%
\pgfsetbuttcap%
\pgfsetroundjoin%
\definecolor{currentfill}{rgb}{1.000000,1.000000,1.000000}%
\pgfsetfillcolor{currentfill}%
\pgfsetfillopacity{0.900000}%
\pgfsetlinewidth{0.200750pt}%
\definecolor{currentstroke}{rgb}{0.200000,0.200000,0.200000}%
\pgfsetstrokecolor{currentstroke}%
\pgfsetdash{}{0pt}%
\pgfpathmoveto{\pgfqpoint{2.200488in}{1.137013in}}%
\pgfpathlineto{\pgfqpoint{2.247446in}{1.152382in}}%
\pgfpathlineto{\pgfqpoint{2.294298in}{1.167716in}}%
\pgfpathlineto{\pgfqpoint{2.326474in}{1.145806in}}%
\pgfpathlineto{\pgfqpoint{2.358754in}{1.123826in}}%
\pgfpathlineto{\pgfqpoint{2.311823in}{1.108392in}}%
\pgfpathlineto{\pgfqpoint{2.264788in}{1.092925in}}%
\pgfpathlineto{\pgfqpoint{2.232587in}{1.115004in}}%
\pgfpathlineto{\pgfqpoint{2.200488in}{1.137013in}}%
\pgfpathclose%
\pgfusepath{stroke,fill}%
\end{pgfscope}%
\begin{pgfscope}%
\pgfpathrectangle{\pgfqpoint{0.050000in}{0.050000in}}{\pgfqpoint{4.900000in}{4.900000in}}%
\pgfusepath{clip}%
\pgfsetbuttcap%
\pgfsetroundjoin%
\definecolor{currentfill}{rgb}{0.844860,0.838462,0.907236}%
\pgfsetfillcolor{currentfill}%
\pgfsetfillopacity{0.900000}%
\pgfsetlinewidth{0.200750pt}%
\definecolor{currentstroke}{rgb}{0.200000,0.200000,0.200000}%
\pgfsetstrokecolor{currentstroke}%
\pgfsetdash{}{0pt}%
\pgfpathmoveto{\pgfqpoint{3.523714in}{1.336742in}}%
\pgfpathlineto{\pgfqpoint{3.569799in}{1.386338in}}%
\pgfpathlineto{\pgfqpoint{3.615849in}{1.436397in}}%
\pgfpathlineto{\pgfqpoint{3.650130in}{1.413283in}}%
\pgfpathlineto{\pgfqpoint{3.684514in}{1.389866in}}%
\pgfpathlineto{\pgfqpoint{3.638407in}{1.342418in}}%
\pgfpathlineto{\pgfqpoint{3.592246in}{1.294666in}}%
\pgfpathlineto{\pgfqpoint{3.557926in}{1.315833in}}%
\pgfpathlineto{\pgfqpoint{3.523714in}{1.336742in}}%
\pgfpathclose%
\pgfusepath{stroke,fill}%
\end{pgfscope}%
\begin{pgfscope}%
\pgfpathrectangle{\pgfqpoint{0.050000in}{0.050000in}}{\pgfqpoint{4.900000in}{4.900000in}}%
\pgfusepath{clip}%
\pgfsetbuttcap%
\pgfsetroundjoin%
\definecolor{currentfill}{rgb}{1.000000,1.000000,1.000000}%
\pgfsetfillcolor{currentfill}%
\pgfsetfillopacity{0.900000}%
\pgfsetlinewidth{0.200750pt}%
\definecolor{currentstroke}{rgb}{0.200000,0.200000,0.200000}%
\pgfsetstrokecolor{currentstroke}%
\pgfsetdash{}{0pt}%
\pgfpathmoveto{\pgfqpoint{1.790110in}{1.132589in}}%
\pgfpathlineto{\pgfqpoint{1.837540in}{1.147968in}}%
\pgfpathlineto{\pgfqpoint{1.884864in}{1.163312in}}%
\pgfpathlineto{\pgfqpoint{1.916398in}{1.141388in}}%
\pgfpathlineto{\pgfqpoint{1.948033in}{1.119393in}}%
\pgfpathlineto{\pgfqpoint{1.900628in}{1.103950in}}%
\pgfpathlineto{\pgfqpoint{1.853117in}{1.088473in}}%
\pgfpathlineto{\pgfqpoint{1.821563in}{1.110566in}}%
\pgfpathlineto{\pgfqpoint{1.790110in}{1.132589in}}%
\pgfpathclose%
\pgfusepath{stroke,fill}%
\end{pgfscope}%
\begin{pgfscope}%
\pgfpathrectangle{\pgfqpoint{0.050000in}{0.050000in}}{\pgfqpoint{4.900000in}{4.900000in}}%
\pgfusepath{clip}%
\pgfsetbuttcap%
\pgfsetroundjoin%
\definecolor{currentfill}{rgb}{1.000000,1.000000,1.000000}%
\pgfsetfillcolor{currentfill}%
\pgfsetfillopacity{0.900000}%
\pgfsetlinewidth{0.200750pt}%
\definecolor{currentstroke}{rgb}{0.200000,0.200000,0.200000}%
\pgfsetstrokecolor{currentstroke}%
\pgfsetdash{}{0pt}%
\pgfpathmoveto{\pgfqpoint{2.769212in}{1.128747in}}%
\pgfpathlineto{\pgfqpoint{2.815567in}{1.144645in}}%
\pgfpathlineto{\pgfqpoint{2.861822in}{1.160933in}}%
\pgfpathlineto{\pgfqpoint{2.894927in}{1.139219in}}%
\pgfpathlineto{\pgfqpoint{2.928137in}{1.117450in}}%
\pgfpathlineto{\pgfqpoint{2.881806in}{1.100873in}}%
\pgfpathlineto{\pgfqpoint{2.835375in}{1.084756in}}%
\pgfpathlineto{\pgfqpoint{2.802241in}{1.106784in}}%
\pgfpathlineto{\pgfqpoint{2.769212in}{1.128747in}}%
\pgfpathclose%
\pgfusepath{stroke,fill}%
\end{pgfscope}%
\begin{pgfscope}%
\pgfpathrectangle{\pgfqpoint{0.050000in}{0.050000in}}{\pgfqpoint{4.900000in}{4.900000in}}%
\pgfusepath{clip}%
\pgfsetbuttcap%
\pgfsetroundjoin%
\definecolor{currentfill}{rgb}{1.000000,1.000000,1.000000}%
\pgfsetfillcolor{currentfill}%
\pgfsetfillopacity{0.900000}%
\pgfsetlinewidth{0.200750pt}%
\definecolor{currentstroke}{rgb}{0.200000,0.200000,0.200000}%
\pgfsetstrokecolor{currentstroke}%
\pgfsetdash{}{0pt}%
\pgfpathmoveto{\pgfqpoint{2.358754in}{1.123826in}}%
\pgfpathlineto{\pgfqpoint{2.405579in}{1.139225in}}%
\pgfpathlineto{\pgfqpoint{2.452300in}{1.154594in}}%
\pgfpathlineto{\pgfqpoint{2.484759in}{1.132643in}}%
\pgfpathlineto{\pgfqpoint{2.517322in}{1.110623in}}%
\pgfpathlineto{\pgfqpoint{2.470524in}{1.095153in}}%
\pgfpathlineto{\pgfqpoint{2.423622in}{1.079654in}}%
\pgfpathlineto{\pgfqpoint{2.391136in}{1.101775in}}%
\pgfpathlineto{\pgfqpoint{2.358754in}{1.123826in}}%
\pgfpathclose%
\pgfusepath{stroke,fill}%
\end{pgfscope}%
\begin{pgfscope}%
\pgfpathrectangle{\pgfqpoint{0.050000in}{0.050000in}}{\pgfqpoint{4.900000in}{4.900000in}}%
\pgfusepath{clip}%
\pgfsetbuttcap%
\pgfsetroundjoin%
\definecolor{currentfill}{rgb}{0.970165,0.968935,0.982161}%
\pgfsetfillcolor{currentfill}%
\pgfsetfillopacity{0.900000}%
\pgfsetlinewidth{0.200750pt}%
\definecolor{currentstroke}{rgb}{0.200000,0.200000,0.200000}%
\pgfsetstrokecolor{currentstroke}%
\pgfsetdash{}{0pt}%
\pgfpathmoveto{\pgfqpoint{3.179647in}{1.154591in}}%
\pgfpathlineto{\pgfqpoint{3.225676in}{1.181025in}}%
\pgfpathlineto{\pgfqpoint{3.271671in}{1.211297in}}%
\pgfpathlineto{\pgfqpoint{3.305468in}{1.191164in}}%
\pgfpathlineto{\pgfqpoint{3.339374in}{1.170929in}}%
\pgfpathlineto{\pgfqpoint{3.293280in}{1.140012in}}%
\pgfpathlineto{\pgfqpoint{3.247156in}{1.112872in}}%
\pgfpathlineto{\pgfqpoint{3.213347in}{1.133763in}}%
\pgfpathlineto{\pgfqpoint{3.179647in}{1.154591in}}%
\pgfpathclose%
\pgfusepath{stroke,fill}%
\end{pgfscope}%
\begin{pgfscope}%
\pgfpathrectangle{\pgfqpoint{0.050000in}{0.050000in}}{\pgfqpoint{4.900000in}{4.900000in}}%
\pgfusepath{clip}%
\pgfsetbuttcap%
\pgfsetroundjoin%
\definecolor{currentfill}{rgb}{1.000000,1.000000,1.000000}%
\pgfsetfillcolor{currentfill}%
\pgfsetfillopacity{0.900000}%
\pgfsetlinewidth{0.200750pt}%
\definecolor{currentstroke}{rgb}{0.200000,0.200000,0.200000}%
\pgfsetstrokecolor{currentstroke}%
\pgfsetdash{}{0pt}%
\pgfpathmoveto{\pgfqpoint{1.948033in}{1.119393in}}%
\pgfpathlineto{\pgfqpoint{1.995332in}{1.134801in}}%
\pgfpathlineto{\pgfqpoint{2.042525in}{1.150175in}}%
\pgfpathlineto{\pgfqpoint{2.074341in}{1.128208in}}%
\pgfpathlineto{\pgfqpoint{2.106259in}{1.106172in}}%
\pgfpathlineto{\pgfqpoint{2.058986in}{1.090699in}}%
\pgfpathlineto{\pgfqpoint{2.011606in}{1.075192in}}%
\pgfpathlineto{\pgfqpoint{1.979769in}{1.097328in}}%
\pgfpathlineto{\pgfqpoint{1.948033in}{1.119393in}}%
\pgfpathclose%
\pgfusepath{stroke,fill}%
\end{pgfscope}%
\begin{pgfscope}%
\pgfpathrectangle{\pgfqpoint{0.050000in}{0.050000in}}{\pgfqpoint{4.900000in}{4.900000in}}%
\pgfusepath{clip}%
\pgfsetbuttcap%
\pgfsetroundjoin%
\definecolor{currentfill}{rgb}{0.898562,0.894379,0.939346}%
\pgfsetfillcolor{currentfill}%
\pgfsetfillopacity{0.900000}%
\pgfsetlinewidth{0.200750pt}%
\definecolor{currentstroke}{rgb}{0.200000,0.200000,0.200000}%
\pgfsetstrokecolor{currentstroke}%
\pgfsetdash{}{0pt}%
\pgfpathmoveto{\pgfqpoint{3.431534in}{1.245640in}}%
\pgfpathlineto{\pgfqpoint{3.477621in}{1.289426in}}%
\pgfpathlineto{\pgfqpoint{3.523714in}{1.336742in}}%
\pgfpathlineto{\pgfqpoint{3.557926in}{1.315833in}}%
\pgfpathlineto{\pgfqpoint{3.592246in}{1.294666in}}%
\pgfpathlineto{\pgfqpoint{3.546064in}{1.248470in}}%
\pgfpathlineto{\pgfqpoint{3.499880in}{1.205181in}}%
\pgfpathlineto{\pgfqpoint{3.465652in}{1.225501in}}%
\pgfpathlineto{\pgfqpoint{3.431534in}{1.245640in}}%
\pgfpathclose%
\pgfusepath{stroke,fill}%
\end{pgfscope}%
\begin{pgfscope}%
\pgfpathrectangle{\pgfqpoint{0.050000in}{0.050000in}}{\pgfqpoint{4.900000in}{4.900000in}}%
\pgfusepath{clip}%
\pgfsetbuttcap%
\pgfsetroundjoin%
\definecolor{currentfill}{rgb}{0.797124,0.788758,0.878693}%
\pgfsetfillcolor{currentfill}%
\pgfsetfillopacity{0.900000}%
\pgfsetlinewidth{0.200750pt}%
\definecolor{currentstroke}{rgb}{0.200000,0.200000,0.200000}%
\pgfsetstrokecolor{currentstroke}%
\pgfsetdash{}{0pt}%
\pgfpathmoveto{\pgfqpoint{3.684514in}{1.389866in}}%
\pgfpathlineto{\pgfqpoint{3.730517in}{1.434723in}}%
\pgfpathlineto{\pgfqpoint{3.776357in}{1.474444in}}%
\pgfpathlineto{\pgfqpoint{3.810860in}{1.447144in}}%
\pgfpathlineto{\pgfqpoint{3.765018in}{1.409351in}}%
\pgfpathlineto{\pgfqpoint{3.719000in}{1.366155in}}%
\pgfpathlineto{\pgfqpoint{3.684514in}{1.389866in}}%
\pgfpathclose%
\pgfusepath{stroke,fill}%
\end{pgfscope}%
\begin{pgfscope}%
\pgfpathrectangle{\pgfqpoint{0.050000in}{0.050000in}}{\pgfqpoint{4.900000in}{4.900000in}}%
\pgfusepath{clip}%
\pgfsetbuttcap%
\pgfsetroundjoin%
\definecolor{currentfill}{rgb}{1.000000,1.000000,1.000000}%
\pgfsetfillcolor{currentfill}%
\pgfsetfillopacity{0.900000}%
\pgfsetlinewidth{0.200750pt}%
\definecolor{currentstroke}{rgb}{0.200000,0.200000,0.200000}%
\pgfsetstrokecolor{currentstroke}%
\pgfsetdash{}{0pt}%
\pgfpathmoveto{\pgfqpoint{1.537048in}{1.114958in}}%
\pgfpathlineto{\pgfqpoint{1.584821in}{1.130376in}}%
\pgfpathlineto{\pgfqpoint{1.632488in}{1.145760in}}%
\pgfpathlineto{\pgfqpoint{1.663660in}{1.123779in}}%
\pgfpathlineto{\pgfqpoint{1.694931in}{1.101728in}}%
\pgfpathlineto{\pgfqpoint{1.647181in}{1.086246in}}%
\pgfpathlineto{\pgfqpoint{1.599325in}{1.070729in}}%
\pgfpathlineto{\pgfqpoint{1.568137in}{1.092879in}}%
\pgfpathlineto{\pgfqpoint{1.537048in}{1.114958in}}%
\pgfpathclose%
\pgfusepath{stroke,fill}%
\end{pgfscope}%
\begin{pgfscope}%
\pgfpathrectangle{\pgfqpoint{0.050000in}{0.050000in}}{\pgfqpoint{4.900000in}{4.900000in}}%
\pgfusepath{clip}%
\pgfsetbuttcap%
\pgfsetroundjoin%
\definecolor{currentfill}{rgb}{1.000000,1.000000,1.000000}%
\pgfsetfillcolor{currentfill}%
\pgfsetfillopacity{0.900000}%
\pgfsetlinewidth{0.200750pt}%
\definecolor{currentstroke}{rgb}{0.200000,0.200000,0.200000}%
\pgfsetstrokecolor{currentstroke}%
\pgfsetdash{}{0pt}%
\pgfpathmoveto{\pgfqpoint{2.928137in}{1.117450in}}%
\pgfpathlineto{\pgfqpoint{2.974372in}{1.134785in}}%
\pgfpathlineto{\pgfqpoint{3.020517in}{1.153288in}}%
\pgfpathlineto{\pgfqpoint{3.053914in}{1.131911in}}%
\pgfpathlineto{\pgfqpoint{3.087418in}{1.110484in}}%
\pgfpathlineto{\pgfqpoint{3.041192in}{1.091477in}}%
\pgfpathlineto{\pgfqpoint{2.994878in}{1.073745in}}%
\pgfpathlineto{\pgfqpoint{2.961454in}{1.095625in}}%
\pgfpathlineto{\pgfqpoint{2.928137in}{1.117450in}}%
\pgfpathclose%
\pgfusepath{stroke,fill}%
\end{pgfscope}%
\begin{pgfscope}%
\pgfpathrectangle{\pgfqpoint{0.050000in}{0.050000in}}{\pgfqpoint{4.900000in}{4.900000in}}%
\pgfusepath{clip}%
\pgfsetbuttcap%
\pgfsetroundjoin%
\definecolor{currentfill}{rgb}{1.000000,1.000000,1.000000}%
\pgfsetfillcolor{currentfill}%
\pgfsetfillopacity{0.900000}%
\pgfsetlinewidth{0.200750pt}%
\definecolor{currentstroke}{rgb}{0.200000,0.200000,0.200000}%
\pgfsetstrokecolor{currentstroke}%
\pgfsetdash{}{0pt}%
\pgfpathmoveto{\pgfqpoint{2.517322in}{1.110623in}}%
\pgfpathlineto{\pgfqpoint{2.564015in}{1.126068in}}%
\pgfpathlineto{\pgfqpoint{2.610604in}{1.141504in}}%
\pgfpathlineto{\pgfqpoint{2.643347in}{1.119524in}}%
\pgfpathlineto{\pgfqpoint{2.676195in}{1.097476in}}%
\pgfpathlineto{\pgfqpoint{2.629530in}{1.081924in}}%
\pgfpathlineto{\pgfqpoint{2.582761in}{1.066372in}}%
\pgfpathlineto{\pgfqpoint{2.549989in}{1.088532in}}%
\pgfpathlineto{\pgfqpoint{2.517322in}{1.110623in}}%
\pgfpathclose%
\pgfusepath{stroke,fill}%
\end{pgfscope}%
\begin{pgfscope}%
\pgfpathrectangle{\pgfqpoint{0.050000in}{0.050000in}}{\pgfqpoint{4.900000in}{4.900000in}}%
\pgfusepath{clip}%
\pgfsetbuttcap%
\pgfsetroundjoin%
\definecolor{currentfill}{rgb}{1.000000,1.000000,1.000000}%
\pgfsetfillcolor{currentfill}%
\pgfsetfillopacity{0.900000}%
\pgfsetlinewidth{0.200750pt}%
\definecolor{currentstroke}{rgb}{0.200000,0.200000,0.200000}%
\pgfsetstrokecolor{currentstroke}%
\pgfsetdash{}{0pt}%
\pgfpathmoveto{\pgfqpoint{2.106259in}{1.106172in}}%
\pgfpathlineto{\pgfqpoint{2.153426in}{1.121609in}}%
\pgfpathlineto{\pgfqpoint{2.200488in}{1.137013in}}%
\pgfpathlineto{\pgfqpoint{2.232587in}{1.115004in}}%
\pgfpathlineto{\pgfqpoint{2.264788in}{1.092925in}}%
\pgfpathlineto{\pgfqpoint{2.217647in}{1.077423in}}%
\pgfpathlineto{\pgfqpoint{2.170401in}{1.061886in}}%
\pgfpathlineto{\pgfqpoint{2.138278in}{1.084064in}}%
\pgfpathlineto{\pgfqpoint{2.106259in}{1.106172in}}%
\pgfpathclose%
\pgfusepath{stroke,fill}%
\end{pgfscope}%
\begin{pgfscope}%
\pgfpathrectangle{\pgfqpoint{0.050000in}{0.050000in}}{\pgfqpoint{4.900000in}{4.900000in}}%
\pgfusepath{clip}%
\pgfsetbuttcap%
\pgfsetroundjoin%
\definecolor{currentfill}{rgb}{1.000000,1.000000,1.000000}%
\pgfsetfillcolor{currentfill}%
\pgfsetfillopacity{0.900000}%
\pgfsetlinewidth{0.200750pt}%
\definecolor{currentstroke}{rgb}{0.200000,0.200000,0.200000}%
\pgfsetstrokecolor{currentstroke}%
\pgfsetdash{}{0pt}%
\pgfpathmoveto{\pgfqpoint{1.694931in}{1.101728in}}%
\pgfpathlineto{\pgfqpoint{1.742573in}{1.117176in}}%
\pgfpathlineto{\pgfqpoint{1.790110in}{1.132589in}}%
\pgfpathlineto{\pgfqpoint{1.821563in}{1.110566in}}%
\pgfpathlineto{\pgfqpoint{1.853117in}{1.088473in}}%
\pgfpathlineto{\pgfqpoint{1.805499in}{1.072961in}}%
\pgfpathlineto{\pgfqpoint{1.757774in}{1.057414in}}%
\pgfpathlineto{\pgfqpoint{1.726302in}{1.079607in}}%
\pgfpathlineto{\pgfqpoint{1.694931in}{1.101728in}}%
\pgfpathclose%
\pgfusepath{stroke,fill}%
\end{pgfscope}%
\begin{pgfscope}%
\pgfpathrectangle{\pgfqpoint{0.050000in}{0.050000in}}{\pgfqpoint{4.900000in}{4.900000in}}%
\pgfusepath{clip}%
\pgfsetbuttcap%
\pgfsetroundjoin%
\definecolor{currentfill}{rgb}{0.844860,0.838462,0.907236}%
\pgfsetfillcolor{currentfill}%
\pgfsetfillopacity{0.900000}%
\pgfsetlinewidth{0.200750pt}%
\definecolor{currentstroke}{rgb}{0.200000,0.200000,0.200000}%
\pgfsetstrokecolor{currentstroke}%
\pgfsetdash{}{0pt}%
\pgfpathmoveto{\pgfqpoint{3.592246in}{1.294666in}}%
\pgfpathlineto{\pgfqpoint{3.638407in}{1.342418in}}%
\pgfpathlineto{\pgfqpoint{3.684514in}{1.389866in}}%
\pgfpathlineto{\pgfqpoint{3.719000in}{1.366155in}}%
\pgfpathlineto{\pgfqpoint{3.672868in}{1.320032in}}%
\pgfpathlineto{\pgfqpoint{3.626672in}{1.273243in}}%
\pgfpathlineto{\pgfqpoint{3.592246in}{1.294666in}}%
\pgfpathclose%
\pgfusepath{stroke,fill}%
\end{pgfscope}%
\begin{pgfscope}%
\pgfpathrectangle{\pgfqpoint{0.050000in}{0.050000in}}{\pgfqpoint{4.900000in}{4.900000in}}%
\pgfusepath{clip}%
\pgfsetbuttcap%
\pgfsetroundjoin%
\definecolor{currentfill}{rgb}{0.940331,0.937870,0.964321}%
\pgfsetfillcolor{currentfill}%
\pgfsetfillopacity{0.900000}%
\pgfsetlinewidth{0.200750pt}%
\definecolor{currentstroke}{rgb}{0.200000,0.200000,0.200000}%
\pgfsetstrokecolor{currentstroke}%
\pgfsetdash{}{0pt}%
\pgfpathmoveto{\pgfqpoint{3.339374in}{1.170929in}}%
\pgfpathlineto{\pgfqpoint{3.385454in}{1.206075in}}%
\pgfpathlineto{\pgfqpoint{3.431534in}{1.245640in}}%
\pgfpathlineto{\pgfqpoint{3.465652in}{1.225501in}}%
\pgfpathlineto{\pgfqpoint{3.499880in}{1.205181in}}%
\pgfpathlineto{\pgfqpoint{3.453699in}{1.165621in}}%
\pgfpathlineto{\pgfqpoint{3.407518in}{1.130134in}}%
\pgfpathlineto{\pgfqpoint{3.373391in}{1.150587in}}%
\pgfpathlineto{\pgfqpoint{3.339374in}{1.170929in}}%
\pgfpathclose%
\pgfusepath{stroke,fill}%
\end{pgfscope}%
\begin{pgfscope}%
\pgfpathrectangle{\pgfqpoint{0.050000in}{0.050000in}}{\pgfqpoint{4.900000in}{4.900000in}}%
\pgfusepath{clip}%
\pgfsetbuttcap%
\pgfsetroundjoin%
\definecolor{currentfill}{rgb}{0.988066,0.987574,0.992864}%
\pgfsetfillcolor{currentfill}%
\pgfsetfillopacity{0.900000}%
\pgfsetlinewidth{0.200750pt}%
\definecolor{currentstroke}{rgb}{0.200000,0.200000,0.200000}%
\pgfsetstrokecolor{currentstroke}%
\pgfsetdash{}{0pt}%
\pgfpathmoveto{\pgfqpoint{3.087418in}{1.110484in}}%
\pgfpathlineto{\pgfqpoint{3.133566in}{1.131309in}}%
\pgfpathlineto{\pgfqpoint{3.179647in}{1.154591in}}%
\pgfpathlineto{\pgfqpoint{3.213347in}{1.133763in}}%
\pgfpathlineto{\pgfqpoint{3.247156in}{1.112872in}}%
\pgfpathlineto{\pgfqpoint{3.200985in}{1.088908in}}%
\pgfpathlineto{\pgfqpoint{3.154752in}{1.067476in}}%
\pgfpathlineto{\pgfqpoint{3.121030in}{1.089006in}}%
\pgfpathlineto{\pgfqpoint{3.087418in}{1.110484in}}%
\pgfpathclose%
\pgfusepath{stroke,fill}%
\end{pgfscope}%
\begin{pgfscope}%
\pgfpathrectangle{\pgfqpoint{0.050000in}{0.050000in}}{\pgfqpoint{4.900000in}{4.900000in}}%
\pgfusepath{clip}%
\pgfsetbuttcap%
\pgfsetroundjoin%
\definecolor{currentfill}{rgb}{1.000000,1.000000,1.000000}%
\pgfsetfillcolor{currentfill}%
\pgfsetfillopacity{0.900000}%
\pgfsetlinewidth{0.200750pt}%
\definecolor{currentstroke}{rgb}{0.200000,0.200000,0.200000}%
\pgfsetstrokecolor{currentstroke}%
\pgfsetdash{}{0pt}%
\pgfpathmoveto{\pgfqpoint{2.676195in}{1.097476in}}%
\pgfpathlineto{\pgfqpoint{2.722755in}{1.113062in}}%
\pgfpathlineto{\pgfqpoint{2.769212in}{1.128747in}}%
\pgfpathlineto{\pgfqpoint{2.802241in}{1.106784in}}%
\pgfpathlineto{\pgfqpoint{2.835375in}{1.084756in}}%
\pgfpathlineto{\pgfqpoint{2.788843in}{1.068900in}}%
\pgfpathlineto{\pgfqpoint{2.742207in}{1.053174in}}%
\pgfpathlineto{\pgfqpoint{2.709148in}{1.075359in}}%
\pgfpathlineto{\pgfqpoint{2.676195in}{1.097476in}}%
\pgfpathclose%
\pgfusepath{stroke,fill}%
\end{pgfscope}%
\begin{pgfscope}%
\pgfpathrectangle{\pgfqpoint{0.050000in}{0.050000in}}{\pgfqpoint{4.900000in}{4.900000in}}%
\pgfusepath{clip}%
\pgfsetbuttcap%
\pgfsetroundjoin%
\definecolor{currentfill}{rgb}{1.000000,1.000000,1.000000}%
\pgfsetfillcolor{currentfill}%
\pgfsetfillopacity{0.900000}%
\pgfsetlinewidth{0.200750pt}%
\definecolor{currentstroke}{rgb}{0.200000,0.200000,0.200000}%
\pgfsetstrokecolor{currentstroke}%
\pgfsetdash{}{0pt}%
\pgfpathmoveto{\pgfqpoint{2.264788in}{1.092925in}}%
\pgfpathlineto{\pgfqpoint{2.311823in}{1.108392in}}%
\pgfpathlineto{\pgfqpoint{2.358754in}{1.123826in}}%
\pgfpathlineto{\pgfqpoint{2.391136in}{1.101775in}}%
\pgfpathlineto{\pgfqpoint{2.423622in}{1.079654in}}%
\pgfpathlineto{\pgfqpoint{2.376614in}{1.064121in}}%
\pgfpathlineto{\pgfqpoint{2.329501in}{1.048554in}}%
\pgfpathlineto{\pgfqpoint{2.297093in}{1.070775in}}%
\pgfpathlineto{\pgfqpoint{2.264788in}{1.092925in}}%
\pgfpathclose%
\pgfusepath{stroke,fill}%
\end{pgfscope}%
\begin{pgfscope}%
\pgfpathrectangle{\pgfqpoint{0.050000in}{0.050000in}}{\pgfqpoint{4.900000in}{4.900000in}}%
\pgfusepath{clip}%
\pgfsetbuttcap%
\pgfsetroundjoin%
\definecolor{currentfill}{rgb}{1.000000,1.000000,1.000000}%
\pgfsetfillcolor{currentfill}%
\pgfsetfillopacity{0.900000}%
\pgfsetlinewidth{0.200750pt}%
\definecolor{currentstroke}{rgb}{0.200000,0.200000,0.200000}%
\pgfsetstrokecolor{currentstroke}%
\pgfsetdash{}{0pt}%
\pgfpathmoveto{\pgfqpoint{1.853117in}{1.088473in}}%
\pgfpathlineto{\pgfqpoint{1.900628in}{1.103950in}}%
\pgfpathlineto{\pgfqpoint{1.948033in}{1.119393in}}%
\pgfpathlineto{\pgfqpoint{1.979769in}{1.097328in}}%
\pgfpathlineto{\pgfqpoint{2.011606in}{1.075192in}}%
\pgfpathlineto{\pgfqpoint{1.964121in}{1.059650in}}%
\pgfpathlineto{\pgfqpoint{1.916528in}{1.044074in}}%
\pgfpathlineto{\pgfqpoint{1.884772in}{1.066309in}}%
\pgfpathlineto{\pgfqpoint{1.853117in}{1.088473in}}%
\pgfpathclose%
\pgfusepath{stroke,fill}%
\end{pgfscope}%
\begin{pgfscope}%
\pgfpathrectangle{\pgfqpoint{0.050000in}{0.050000in}}{\pgfqpoint{4.900000in}{4.900000in}}%
\pgfusepath{clip}%
\pgfsetbuttcap%
\pgfsetroundjoin%
\definecolor{currentfill}{rgb}{0.898562,0.894379,0.939346}%
\pgfsetfillcolor{currentfill}%
\pgfsetfillopacity{0.900000}%
\pgfsetlinewidth{0.200750pt}%
\definecolor{currentstroke}{rgb}{0.200000,0.200000,0.200000}%
\pgfsetstrokecolor{currentstroke}%
\pgfsetdash{}{0pt}%
\pgfpathmoveto{\pgfqpoint{3.499880in}{1.205181in}}%
\pgfpathlineto{\pgfqpoint{3.546064in}{1.248470in}}%
\pgfpathlineto{\pgfqpoint{3.592246in}{1.294666in}}%
\pgfpathlineto{\pgfqpoint{3.626672in}{1.273243in}}%
\pgfpathlineto{\pgfqpoint{3.580448in}{1.227661in}}%
\pgfpathlineto{\pgfqpoint{3.534217in}{1.184678in}}%
\pgfpathlineto{\pgfqpoint{3.499880in}{1.205181in}}%
\pgfpathclose%
\pgfusepath{stroke,fill}%
\end{pgfscope}%
\begin{pgfscope}%
\pgfpathrectangle{\pgfqpoint{0.050000in}{0.050000in}}{\pgfqpoint{4.900000in}{4.900000in}}%
\pgfusepath{clip}%
\pgfsetbuttcap%
\pgfsetroundjoin%
\definecolor{currentfill}{rgb}{1.000000,1.000000,1.000000}%
\pgfsetfillcolor{currentfill}%
\pgfsetfillopacity{0.900000}%
\pgfsetlinewidth{0.200750pt}%
\definecolor{currentstroke}{rgb}{0.200000,0.200000,0.200000}%
\pgfsetstrokecolor{currentstroke}%
\pgfsetdash{}{0pt}%
\pgfpathmoveto{\pgfqpoint{2.835375in}{1.084756in}}%
\pgfpathlineto{\pgfqpoint{2.881806in}{1.100873in}}%
\pgfpathlineto{\pgfqpoint{2.928137in}{1.117450in}}%
\pgfpathlineto{\pgfqpoint{2.961454in}{1.095625in}}%
\pgfpathlineto{\pgfqpoint{2.994878in}{1.073745in}}%
\pgfpathlineto{\pgfqpoint{2.948470in}{1.056862in}}%
\pgfpathlineto{\pgfqpoint{2.901964in}{1.040512in}}%
\pgfpathlineto{\pgfqpoint{2.868616in}{1.062666in}}%
\pgfpathlineto{\pgfqpoint{2.835375in}{1.084756in}}%
\pgfpathclose%
\pgfusepath{stroke,fill}%
\end{pgfscope}%
\begin{pgfscope}%
\pgfpathrectangle{\pgfqpoint{0.050000in}{0.050000in}}{\pgfqpoint{4.900000in}{4.900000in}}%
\pgfusepath{clip}%
\pgfsetbuttcap%
\pgfsetroundjoin%
\definecolor{currentfill}{rgb}{1.000000,1.000000,1.000000}%
\pgfsetfillcolor{currentfill}%
\pgfsetfillopacity{0.900000}%
\pgfsetlinewidth{0.200750pt}%
\definecolor{currentstroke}{rgb}{0.200000,0.200000,0.200000}%
\pgfsetstrokecolor{currentstroke}%
\pgfsetdash{}{0pt}%
\pgfpathmoveto{\pgfqpoint{2.423622in}{1.079654in}}%
\pgfpathlineto{\pgfqpoint{2.470524in}{1.095153in}}%
\pgfpathlineto{\pgfqpoint{2.517322in}{1.110623in}}%
\pgfpathlineto{\pgfqpoint{2.549989in}{1.088532in}}%
\pgfpathlineto{\pgfqpoint{2.582761in}{1.066372in}}%
\pgfpathlineto{\pgfqpoint{2.535887in}{1.050799in}}%
\pgfpathlineto{\pgfqpoint{2.488907in}{1.035199in}}%
\pgfpathlineto{\pgfqpoint{2.456212in}{1.057462in}}%
\pgfpathlineto{\pgfqpoint{2.423622in}{1.079654in}}%
\pgfpathclose%
\pgfusepath{stroke,fill}%
\end{pgfscope}%
\begin{pgfscope}%
\pgfpathrectangle{\pgfqpoint{0.050000in}{0.050000in}}{\pgfqpoint{4.900000in}{4.900000in}}%
\pgfusepath{clip}%
\pgfsetbuttcap%
\pgfsetroundjoin%
\definecolor{currentfill}{rgb}{0.970165,0.968935,0.982161}%
\pgfsetfillcolor{currentfill}%
\pgfsetfillopacity{0.900000}%
\pgfsetlinewidth{0.200750pt}%
\definecolor{currentstroke}{rgb}{0.200000,0.200000,0.200000}%
\pgfsetstrokecolor{currentstroke}%
\pgfsetdash{}{0pt}%
\pgfpathmoveto{\pgfqpoint{3.247156in}{1.112872in}}%
\pgfpathlineto{\pgfqpoint{3.293280in}{1.140012in}}%
\pgfpathlineto{\pgfqpoint{3.339374in}{1.170929in}}%
\pgfpathlineto{\pgfqpoint{3.373391in}{1.150587in}}%
\pgfpathlineto{\pgfqpoint{3.407518in}{1.130134in}}%
\pgfpathlineto{\pgfqpoint{3.361325in}{1.098668in}}%
\pgfpathlineto{\pgfqpoint{3.315105in}{1.070887in}}%
\pgfpathlineto{\pgfqpoint{3.281076in}{1.091914in}}%
\pgfpathlineto{\pgfqpoint{3.247156in}{1.112872in}}%
\pgfpathclose%
\pgfusepath{stroke,fill}%
\end{pgfscope}%
\begin{pgfscope}%
\pgfpathrectangle{\pgfqpoint{0.050000in}{0.050000in}}{\pgfqpoint{4.900000in}{4.900000in}}%
\pgfusepath{clip}%
\pgfsetbuttcap%
\pgfsetroundjoin%
\definecolor{currentfill}{rgb}{1.000000,1.000000,1.000000}%
\pgfsetfillcolor{currentfill}%
\pgfsetfillopacity{0.900000}%
\pgfsetlinewidth{0.200750pt}%
\definecolor{currentstroke}{rgb}{0.200000,0.200000,0.200000}%
\pgfsetstrokecolor{currentstroke}%
\pgfsetdash{}{0pt}%
\pgfpathmoveto{\pgfqpoint{2.011606in}{1.075192in}}%
\pgfpathlineto{\pgfqpoint{2.058986in}{1.090699in}}%
\pgfpathlineto{\pgfqpoint{2.106259in}{1.106172in}}%
\pgfpathlineto{\pgfqpoint{2.138278in}{1.084064in}}%
\pgfpathlineto{\pgfqpoint{2.170401in}{1.061886in}}%
\pgfpathlineto{\pgfqpoint{2.123048in}{1.046314in}}%
\pgfpathlineto{\pgfqpoint{2.075588in}{1.030708in}}%
\pgfpathlineto{\pgfqpoint{2.043546in}{1.052986in}}%
\pgfpathlineto{\pgfqpoint{2.011606in}{1.075192in}}%
\pgfpathclose%
\pgfusepath{stroke,fill}%
\end{pgfscope}%
\begin{pgfscope}%
\pgfpathrectangle{\pgfqpoint{0.050000in}{0.050000in}}{\pgfqpoint{4.900000in}{4.900000in}}%
\pgfusepath{clip}%
\pgfsetbuttcap%
\pgfsetroundjoin%
\definecolor{currentfill}{rgb}{1.000000,1.000000,1.000000}%
\pgfsetfillcolor{currentfill}%
\pgfsetfillopacity{0.900000}%
\pgfsetlinewidth{0.200750pt}%
\definecolor{currentstroke}{rgb}{0.200000,0.200000,0.200000}%
\pgfsetstrokecolor{currentstroke}%
\pgfsetdash{}{0pt}%
\pgfpathmoveto{\pgfqpoint{1.599325in}{1.070729in}}%
\pgfpathlineto{\pgfqpoint{1.647181in}{1.086246in}}%
\pgfpathlineto{\pgfqpoint{1.694931in}{1.101728in}}%
\pgfpathlineto{\pgfqpoint{1.726302in}{1.079607in}}%
\pgfpathlineto{\pgfqpoint{1.757774in}{1.057414in}}%
\pgfpathlineto{\pgfqpoint{1.709942in}{1.041832in}}%
\pgfpathlineto{\pgfqpoint{1.662002in}{1.026216in}}%
\pgfpathlineto{\pgfqpoint{1.630613in}{1.048508in}}%
\pgfpathlineto{\pgfqpoint{1.599325in}{1.070729in}}%
\pgfpathclose%
\pgfusepath{stroke,fill}%
\end{pgfscope}%
\begin{pgfscope}%
\pgfpathrectangle{\pgfqpoint{0.050000in}{0.050000in}}{\pgfqpoint{4.900000in}{4.900000in}}%
\pgfusepath{clip}%
\pgfsetbuttcap%
\pgfsetroundjoin%
\definecolor{currentfill}{rgb}{1.000000,1.000000,1.000000}%
\pgfsetfillcolor{currentfill}%
\pgfsetfillopacity{0.900000}%
\pgfsetlinewidth{0.200750pt}%
\definecolor{currentstroke}{rgb}{0.200000,0.200000,0.200000}%
\pgfsetstrokecolor{currentstroke}%
\pgfsetdash{}{0pt}%
\pgfpathmoveto{\pgfqpoint{2.994878in}{1.073745in}}%
\pgfpathlineto{\pgfqpoint{3.041192in}{1.091477in}}%
\pgfpathlineto{\pgfqpoint{3.087418in}{1.110484in}}%
\pgfpathlineto{\pgfqpoint{3.121030in}{1.089006in}}%
\pgfpathlineto{\pgfqpoint{3.154752in}{1.067476in}}%
\pgfpathlineto{\pgfqpoint{3.108444in}{1.047957in}}%
\pgfpathlineto{\pgfqpoint{3.062050in}{1.029815in}}%
\pgfpathlineto{\pgfqpoint{3.028410in}{1.051808in}}%
\pgfpathlineto{\pgfqpoint{2.994878in}{1.073745in}}%
\pgfpathclose%
\pgfusepath{stroke,fill}%
\end{pgfscope}%
\begin{pgfscope}%
\pgfpathrectangle{\pgfqpoint{0.050000in}{0.050000in}}{\pgfqpoint{4.900000in}{4.900000in}}%
\pgfusepath{clip}%
\pgfsetbuttcap%
\pgfsetroundjoin%
\definecolor{currentfill}{rgb}{0.940331,0.937870,0.964321}%
\pgfsetfillcolor{currentfill}%
\pgfsetfillopacity{0.900000}%
\pgfsetlinewidth{0.200750pt}%
\definecolor{currentstroke}{rgb}{0.200000,0.200000,0.200000}%
\pgfsetstrokecolor{currentstroke}%
\pgfsetdash{}{0pt}%
\pgfpathmoveto{\pgfqpoint{3.407518in}{1.130134in}}%
\pgfpathlineto{\pgfqpoint{3.453699in}{1.165621in}}%
\pgfpathlineto{\pgfqpoint{3.499880in}{1.205181in}}%
\pgfpathlineto{\pgfqpoint{3.534217in}{1.184678in}}%
\pgfpathlineto{\pgfqpoint{3.487987in}{1.145176in}}%
\pgfpathlineto{\pgfqpoint{3.441756in}{1.109567in}}%
\pgfpathlineto{\pgfqpoint{3.407518in}{1.130134in}}%
\pgfpathclose%
\pgfusepath{stroke,fill}%
\end{pgfscope}%
\begin{pgfscope}%
\pgfpathrectangle{\pgfqpoint{0.050000in}{0.050000in}}{\pgfqpoint{4.900000in}{4.900000in}}%
\pgfusepath{clip}%
\pgfsetbuttcap%
\pgfsetroundjoin%
\definecolor{currentfill}{rgb}{1.000000,1.000000,1.000000}%
\pgfsetfillcolor{currentfill}%
\pgfsetfillopacity{0.900000}%
\pgfsetlinewidth{0.200750pt}%
\definecolor{currentstroke}{rgb}{0.200000,0.200000,0.200000}%
\pgfsetstrokecolor{currentstroke}%
\pgfsetdash{}{0pt}%
\pgfpathmoveto{\pgfqpoint{2.582761in}{1.066372in}}%
\pgfpathlineto{\pgfqpoint{2.629530in}{1.081924in}}%
\pgfpathlineto{\pgfqpoint{2.676195in}{1.097476in}}%
\pgfpathlineto{\pgfqpoint{2.709148in}{1.075359in}}%
\pgfpathlineto{\pgfqpoint{2.742207in}{1.053174in}}%
\pgfpathlineto{\pgfqpoint{2.695466in}{1.037502in}}%
\pgfpathlineto{\pgfqpoint{2.648621in}{1.021838in}}%
\pgfpathlineto{\pgfqpoint{2.615638in}{1.044140in}}%
\pgfpathlineto{\pgfqpoint{2.582761in}{1.066372in}}%
\pgfpathclose%
\pgfusepath{stroke,fill}%
\end{pgfscope}%
\begin{pgfscope}%
\pgfpathrectangle{\pgfqpoint{0.050000in}{0.050000in}}{\pgfqpoint{4.900000in}{4.900000in}}%
\pgfusepath{clip}%
\pgfsetbuttcap%
\pgfsetroundjoin%
\definecolor{currentfill}{rgb}{1.000000,1.000000,1.000000}%
\pgfsetfillcolor{currentfill}%
\pgfsetfillopacity{0.900000}%
\pgfsetlinewidth{0.200750pt}%
\definecolor{currentstroke}{rgb}{0.200000,0.200000,0.200000}%
\pgfsetstrokecolor{currentstroke}%
\pgfsetdash{}{0pt}%
\pgfpathmoveto{\pgfqpoint{2.170401in}{1.061886in}}%
\pgfpathlineto{\pgfqpoint{2.217647in}{1.077423in}}%
\pgfpathlineto{\pgfqpoint{2.264788in}{1.092925in}}%
\pgfpathlineto{\pgfqpoint{2.297093in}{1.070775in}}%
\pgfpathlineto{\pgfqpoint{2.329501in}{1.048554in}}%
\pgfpathlineto{\pgfqpoint{2.282281in}{1.032952in}}%
\pgfpathlineto{\pgfqpoint{2.234955in}{1.017316in}}%
\pgfpathlineto{\pgfqpoint{2.202626in}{1.039637in}}%
\pgfpathlineto{\pgfqpoint{2.170401in}{1.061886in}}%
\pgfpathclose%
\pgfusepath{stroke,fill}%
\end{pgfscope}%
\begin{pgfscope}%
\pgfpathrectangle{\pgfqpoint{0.050000in}{0.050000in}}{\pgfqpoint{4.900000in}{4.900000in}}%
\pgfusepath{clip}%
\pgfsetbuttcap%
\pgfsetroundjoin%
\definecolor{currentfill}{rgb}{1.000000,1.000000,1.000000}%
\pgfsetfillcolor{currentfill}%
\pgfsetfillopacity{0.900000}%
\pgfsetlinewidth{0.200750pt}%
\definecolor{currentstroke}{rgb}{0.200000,0.200000,0.200000}%
\pgfsetstrokecolor{currentstroke}%
\pgfsetdash{}{0pt}%
\pgfpathmoveto{\pgfqpoint{1.757774in}{1.057414in}}%
\pgfpathlineto{\pgfqpoint{1.805499in}{1.072961in}}%
\pgfpathlineto{\pgfqpoint{1.853117in}{1.088473in}}%
\pgfpathlineto{\pgfqpoint{1.884772in}{1.066309in}}%
\pgfpathlineto{\pgfqpoint{1.916528in}{1.044074in}}%
\pgfpathlineto{\pgfqpoint{1.868829in}{1.028462in}}%
\pgfpathlineto{\pgfqpoint{1.821022in}{1.012815in}}%
\pgfpathlineto{\pgfqpoint{1.789347in}{1.035150in}}%
\pgfpathlineto{\pgfqpoint{1.757774in}{1.057414in}}%
\pgfpathclose%
\pgfusepath{stroke,fill}%
\end{pgfscope}%
\begin{pgfscope}%
\pgfpathrectangle{\pgfqpoint{0.050000in}{0.050000in}}{\pgfqpoint{4.900000in}{4.900000in}}%
\pgfusepath{clip}%
\pgfsetbuttcap%
\pgfsetroundjoin%
\definecolor{currentfill}{rgb}{0.988066,0.987574,0.992864}%
\pgfsetfillcolor{currentfill}%
\pgfsetfillopacity{0.900000}%
\pgfsetlinewidth{0.200750pt}%
\definecolor{currentstroke}{rgb}{0.200000,0.200000,0.200000}%
\pgfsetstrokecolor{currentstroke}%
\pgfsetdash{}{0pt}%
\pgfpathmoveto{\pgfqpoint{3.154752in}{1.067476in}}%
\pgfpathlineto{\pgfqpoint{3.200985in}{1.088908in}}%
\pgfpathlineto{\pgfqpoint{3.247156in}{1.112872in}}%
\pgfpathlineto{\pgfqpoint{3.281076in}{1.091914in}}%
\pgfpathlineto{\pgfqpoint{3.315105in}{1.070887in}}%
\pgfpathlineto{\pgfqpoint{3.268843in}{1.046277in}}%
\pgfpathlineto{\pgfqpoint{3.222523in}{1.024250in}}%
\pgfpathlineto{\pgfqpoint{3.188582in}{1.045891in}}%
\pgfpathlineto{\pgfqpoint{3.154752in}{1.067476in}}%
\pgfpathclose%
\pgfusepath{stroke,fill}%
\end{pgfscope}%
\begin{pgfscope}%
\pgfpathrectangle{\pgfqpoint{0.050000in}{0.050000in}}{\pgfqpoint{4.900000in}{4.900000in}}%
\pgfusepath{clip}%
\pgfsetbuttcap%
\pgfsetroundjoin%
\definecolor{currentfill}{rgb}{1.000000,1.000000,1.000000}%
\pgfsetfillcolor{currentfill}%
\pgfsetfillopacity{0.900000}%
\pgfsetlinewidth{0.200750pt}%
\definecolor{currentstroke}{rgb}{0.200000,0.200000,0.200000}%
\pgfsetstrokecolor{currentstroke}%
\pgfsetdash{}{0pt}%
\pgfpathmoveto{\pgfqpoint{2.742207in}{1.053174in}}%
\pgfpathlineto{\pgfqpoint{2.788843in}{1.068900in}}%
\pgfpathlineto{\pgfqpoint{2.835375in}{1.084756in}}%
\pgfpathlineto{\pgfqpoint{2.868616in}{1.062666in}}%
\pgfpathlineto{\pgfqpoint{2.901964in}{1.040512in}}%
\pgfpathlineto{\pgfqpoint{2.855356in}{1.024472in}}%
\pgfpathlineto{\pgfqpoint{2.808644in}{1.008598in}}%
\pgfpathlineto{\pgfqpoint{2.775372in}{1.030920in}}%
\pgfpathlineto{\pgfqpoint{2.742207in}{1.053174in}}%
\pgfpathclose%
\pgfusepath{stroke,fill}%
\end{pgfscope}%
\begin{pgfscope}%
\pgfpathrectangle{\pgfqpoint{0.050000in}{0.050000in}}{\pgfqpoint{4.900000in}{4.900000in}}%
\pgfusepath{clip}%
\pgfsetbuttcap%
\pgfsetroundjoin%
\definecolor{currentfill}{rgb}{1.000000,1.000000,1.000000}%
\pgfsetfillcolor{currentfill}%
\pgfsetfillopacity{0.900000}%
\pgfsetlinewidth{0.200750pt}%
\definecolor{currentstroke}{rgb}{0.200000,0.200000,0.200000}%
\pgfsetstrokecolor{currentstroke}%
\pgfsetdash{}{0pt}%
\pgfpathmoveto{\pgfqpoint{2.329501in}{1.048554in}}%
\pgfpathlineto{\pgfqpoint{2.376614in}{1.064121in}}%
\pgfpathlineto{\pgfqpoint{2.423622in}{1.079654in}}%
\pgfpathlineto{\pgfqpoint{2.456212in}{1.057462in}}%
\pgfpathlineto{\pgfqpoint{2.488907in}{1.035199in}}%
\pgfpathlineto{\pgfqpoint{2.441822in}{1.019565in}}%
\pgfpathlineto{\pgfqpoint{2.394630in}{1.003898in}}%
\pgfpathlineto{\pgfqpoint{2.362013in}{1.026262in}}%
\pgfpathlineto{\pgfqpoint{2.329501in}{1.048554in}}%
\pgfpathclose%
\pgfusepath{stroke,fill}%
\end{pgfscope}%
\begin{pgfscope}%
\pgfpathrectangle{\pgfqpoint{0.050000in}{0.050000in}}{\pgfqpoint{4.900000in}{4.900000in}}%
\pgfusepath{clip}%
\pgfsetbuttcap%
\pgfsetroundjoin%
\definecolor{currentfill}{rgb}{0.970165,0.968935,0.982161}%
\pgfsetfillcolor{currentfill}%
\pgfsetfillopacity{0.900000}%
\pgfsetlinewidth{0.200750pt}%
\definecolor{currentstroke}{rgb}{0.200000,0.200000,0.200000}%
\pgfsetstrokecolor{currentstroke}%
\pgfsetdash{}{0pt}%
\pgfpathmoveto{\pgfqpoint{3.315105in}{1.070887in}}%
\pgfpathlineto{\pgfqpoint{3.361325in}{1.098668in}}%
\pgfpathlineto{\pgfqpoint{3.407518in}{1.130134in}}%
\pgfpathlineto{\pgfqpoint{3.441756in}{1.109567in}}%
\pgfpathlineto{\pgfqpoint{3.395513in}{1.077864in}}%
\pgfpathlineto{\pgfqpoint{3.349246in}{1.049787in}}%
\pgfpathlineto{\pgfqpoint{3.315105in}{1.070887in}}%
\pgfpathclose%
\pgfusepath{stroke,fill}%
\end{pgfscope}%
\begin{pgfscope}%
\pgfpathrectangle{\pgfqpoint{0.050000in}{0.050000in}}{\pgfqpoint{4.900000in}{4.900000in}}%
\pgfusepath{clip}%
\pgfsetbuttcap%
\pgfsetroundjoin%
\definecolor{currentfill}{rgb}{1.000000,1.000000,1.000000}%
\pgfsetfillcolor{currentfill}%
\pgfsetfillopacity{0.900000}%
\pgfsetlinewidth{0.200750pt}%
\definecolor{currentstroke}{rgb}{0.200000,0.200000,0.200000}%
\pgfsetstrokecolor{currentstroke}%
\pgfsetdash{}{0pt}%
\pgfpathmoveto{\pgfqpoint{1.916528in}{1.044074in}}%
\pgfpathlineto{\pgfqpoint{1.964121in}{1.059650in}}%
\pgfpathlineto{\pgfqpoint{2.011606in}{1.075192in}}%
\pgfpathlineto{\pgfqpoint{2.043546in}{1.052986in}}%
\pgfpathlineto{\pgfqpoint{2.075588in}{1.030708in}}%
\pgfpathlineto{\pgfqpoint{2.028022in}{1.015066in}}%
\pgfpathlineto{\pgfqpoint{1.980349in}{0.999388in}}%
\pgfpathlineto{\pgfqpoint{1.948387in}{1.021767in}}%
\pgfpathlineto{\pgfqpoint{1.916528in}{1.044074in}}%
\pgfpathclose%
\pgfusepath{stroke,fill}%
\end{pgfscope}%
\begin{pgfscope}%
\pgfpathrectangle{\pgfqpoint{0.050000in}{0.050000in}}{\pgfqpoint{4.900000in}{4.900000in}}%
\pgfusepath{clip}%
\pgfsetbuttcap%
\pgfsetroundjoin%
\definecolor{currentfill}{rgb}{1.000000,1.000000,1.000000}%
\pgfsetfillcolor{currentfill}%
\pgfsetfillopacity{0.900000}%
\pgfsetlinewidth{0.200750pt}%
\definecolor{currentstroke}{rgb}{0.200000,0.200000,0.200000}%
\pgfsetstrokecolor{currentstroke}%
\pgfsetdash{}{0pt}%
\pgfpathmoveto{\pgfqpoint{2.901964in}{1.040512in}}%
\pgfpathlineto{\pgfqpoint{2.948470in}{1.056862in}}%
\pgfpathlineto{\pgfqpoint{2.994878in}{1.073745in}}%
\pgfpathlineto{\pgfqpoint{3.028410in}{1.051808in}}%
\pgfpathlineto{\pgfqpoint{3.062050in}{1.029815in}}%
\pgfpathlineto{\pgfqpoint{3.015565in}{1.012611in}}%
\pgfpathlineto{\pgfqpoint{2.968982in}{0.996012in}}%
\pgfpathlineto{\pgfqpoint{2.935419in}{1.018294in}}%
\pgfpathlineto{\pgfqpoint{2.901964in}{1.040512in}}%
\pgfpathclose%
\pgfusepath{stroke,fill}%
\end{pgfscope}%
\begin{pgfscope}%
\pgfpathrectangle{\pgfqpoint{0.050000in}{0.050000in}}{\pgfqpoint{4.900000in}{4.900000in}}%
\pgfusepath{clip}%
\pgfsetbuttcap%
\pgfsetroundjoin%
\definecolor{currentfill}{rgb}{1.000000,1.000000,1.000000}%
\pgfsetfillcolor{currentfill}%
\pgfsetfillopacity{0.900000}%
\pgfsetlinewidth{0.200750pt}%
\definecolor{currentstroke}{rgb}{0.200000,0.200000,0.200000}%
\pgfsetstrokecolor{currentstroke}%
\pgfsetdash{}{0pt}%
\pgfpathmoveto{\pgfqpoint{2.488907in}{1.035199in}}%
\pgfpathlineto{\pgfqpoint{2.535887in}{1.050799in}}%
\pgfpathlineto{\pgfqpoint{2.582761in}{1.066372in}}%
\pgfpathlineto{\pgfqpoint{2.615638in}{1.044140in}}%
\pgfpathlineto{\pgfqpoint{2.648621in}{1.021838in}}%
\pgfpathlineto{\pgfqpoint{2.601670in}{1.006161in}}%
\pgfpathlineto{\pgfqpoint{2.554613in}{0.990457in}}%
\pgfpathlineto{\pgfqpoint{2.521707in}{1.012864in}}%
\pgfpathlineto{\pgfqpoint{2.488907in}{1.035199in}}%
\pgfpathclose%
\pgfusepath{stroke,fill}%
\end{pgfscope}%
\begin{pgfscope}%
\pgfpathrectangle{\pgfqpoint{0.050000in}{0.050000in}}{\pgfqpoint{4.900000in}{4.900000in}}%
\pgfusepath{clip}%
\pgfsetbuttcap%
\pgfsetroundjoin%
\definecolor{currentfill}{rgb}{1.000000,1.000000,1.000000}%
\pgfsetfillcolor{currentfill}%
\pgfsetfillopacity{0.900000}%
\pgfsetlinewidth{0.200750pt}%
\definecolor{currentstroke}{rgb}{0.200000,0.200000,0.200000}%
\pgfsetstrokecolor{currentstroke}%
\pgfsetdash{}{0pt}%
\pgfpathmoveto{\pgfqpoint{2.075588in}{1.030708in}}%
\pgfpathlineto{\pgfqpoint{2.123048in}{1.046314in}}%
\pgfpathlineto{\pgfqpoint{2.170401in}{1.061886in}}%
\pgfpathlineto{\pgfqpoint{2.202626in}{1.039637in}}%
\pgfpathlineto{\pgfqpoint{2.234955in}{1.017316in}}%
\pgfpathlineto{\pgfqpoint{2.187523in}{1.001643in}}%
\pgfpathlineto{\pgfqpoint{2.139983in}{0.985936in}}%
\pgfpathlineto{\pgfqpoint{2.107734in}{1.008358in}}%
\pgfpathlineto{\pgfqpoint{2.075588in}{1.030708in}}%
\pgfpathclose%
\pgfusepath{stroke,fill}%
\end{pgfscope}%
\begin{pgfscope}%
\pgfpathrectangle{\pgfqpoint{0.050000in}{0.050000in}}{\pgfqpoint{4.900000in}{4.900000in}}%
\pgfusepath{clip}%
\pgfsetbuttcap%
\pgfsetroundjoin%
\definecolor{currentfill}{rgb}{1.000000,1.000000,1.000000}%
\pgfsetfillcolor{currentfill}%
\pgfsetfillopacity{0.900000}%
\pgfsetlinewidth{0.200750pt}%
\definecolor{currentstroke}{rgb}{0.200000,0.200000,0.200000}%
\pgfsetstrokecolor{currentstroke}%
\pgfsetdash{}{0pt}%
\pgfpathmoveto{\pgfqpoint{1.662002in}{1.026216in}}%
\pgfpathlineto{\pgfqpoint{1.709942in}{1.041832in}}%
\pgfpathlineto{\pgfqpoint{1.757774in}{1.057414in}}%
\pgfpathlineto{\pgfqpoint{1.789347in}{1.035150in}}%
\pgfpathlineto{\pgfqpoint{1.821022in}{1.012815in}}%
\pgfpathlineto{\pgfqpoint{1.773107in}{0.997133in}}%
\pgfpathlineto{\pgfqpoint{1.725084in}{0.981415in}}%
\pgfpathlineto{\pgfqpoint{1.693492in}{1.003851in}}%
\pgfpathlineto{\pgfqpoint{1.662002in}{1.026216in}}%
\pgfpathclose%
\pgfusepath{stroke,fill}%
\end{pgfscope}%
\begin{pgfscope}%
\pgfpathrectangle{\pgfqpoint{0.050000in}{0.050000in}}{\pgfqpoint{4.900000in}{4.900000in}}%
\pgfusepath{clip}%
\pgfsetbuttcap%
\pgfsetroundjoin%
\definecolor{currentfill}{rgb}{1.000000,1.000000,1.000000}%
\pgfsetfillcolor{currentfill}%
\pgfsetfillopacity{0.900000}%
\pgfsetlinewidth{0.200750pt}%
\definecolor{currentstroke}{rgb}{0.200000,0.200000,0.200000}%
\pgfsetstrokecolor{currentstroke}%
\pgfsetdash{}{0pt}%
\pgfpathmoveto{\pgfqpoint{3.062050in}{1.029815in}}%
\pgfpathlineto{\pgfqpoint{3.108444in}{1.047957in}}%
\pgfpathlineto{\pgfqpoint{3.154752in}{1.067476in}}%
\pgfpathlineto{\pgfqpoint{3.188582in}{1.045891in}}%
\pgfpathlineto{\pgfqpoint{3.222523in}{1.024250in}}%
\pgfpathlineto{\pgfqpoint{3.176131in}{1.004220in}}%
\pgfpathlineto{\pgfqpoint{3.129658in}{0.985657in}}%
\pgfpathlineto{\pgfqpoint{3.095799in}{1.007765in}}%
\pgfpathlineto{\pgfqpoint{3.062050in}{1.029815in}}%
\pgfpathclose%
\pgfusepath{stroke,fill}%
\end{pgfscope}%
\begin{pgfscope}%
\pgfpathrectangle{\pgfqpoint{0.050000in}{0.050000in}}{\pgfqpoint{4.900000in}{4.900000in}}%
\pgfusepath{clip}%
\pgfsetbuttcap%
\pgfsetroundjoin%
\definecolor{currentfill}{rgb}{1.000000,1.000000,1.000000}%
\pgfsetfillcolor{currentfill}%
\pgfsetfillopacity{0.900000}%
\pgfsetlinewidth{0.200750pt}%
\definecolor{currentstroke}{rgb}{0.200000,0.200000,0.200000}%
\pgfsetstrokecolor{currentstroke}%
\pgfsetdash{}{0pt}%
\pgfpathmoveto{\pgfqpoint{2.648621in}{1.021838in}}%
\pgfpathlineto{\pgfqpoint{2.695466in}{1.037502in}}%
\pgfpathlineto{\pgfqpoint{2.742207in}{1.053174in}}%
\pgfpathlineto{\pgfqpoint{2.775372in}{1.030920in}}%
\pgfpathlineto{\pgfqpoint{2.808644in}{1.008598in}}%
\pgfpathlineto{\pgfqpoint{2.761828in}{0.992798in}}%
\pgfpathlineto{\pgfqpoint{2.714906in}{0.977021in}}%
\pgfpathlineto{\pgfqpoint{2.681710in}{0.999465in}}%
\pgfpathlineto{\pgfqpoint{2.648621in}{1.021838in}}%
\pgfpathclose%
\pgfusepath{stroke,fill}%
\end{pgfscope}%
\begin{pgfscope}%
\pgfpathrectangle{\pgfqpoint{0.050000in}{0.050000in}}{\pgfqpoint{4.900000in}{4.900000in}}%
\pgfusepath{clip}%
\pgfsetbuttcap%
\pgfsetroundjoin%
\definecolor{currentfill}{rgb}{0.988066,0.987574,0.992864}%
\pgfsetfillcolor{currentfill}%
\pgfsetfillopacity{0.900000}%
\pgfsetlinewidth{0.200750pt}%
\definecolor{currentstroke}{rgb}{0.200000,0.200000,0.200000}%
\pgfsetstrokecolor{currentstroke}%
\pgfsetdash{}{0pt}%
\pgfpathmoveto{\pgfqpoint{3.222523in}{1.024250in}}%
\pgfpathlineto{\pgfqpoint{3.268843in}{1.046277in}}%
\pgfpathlineto{\pgfqpoint{3.315105in}{1.070887in}}%
\pgfpathlineto{\pgfqpoint{3.349246in}{1.049787in}}%
\pgfpathlineto{\pgfqpoint{3.302938in}{1.024869in}}%
\pgfpathlineto{\pgfqpoint{3.256574in}{1.002552in}}%
\pgfpathlineto{\pgfqpoint{3.222523in}{1.024250in}}%
\pgfpathclose%
\pgfusepath{stroke,fill}%
\end{pgfscope}%
\begin{pgfscope}%
\pgfpathrectangle{\pgfqpoint{0.050000in}{0.050000in}}{\pgfqpoint{4.900000in}{4.900000in}}%
\pgfusepath{clip}%
\pgfsetbuttcap%
\pgfsetroundjoin%
\definecolor{currentfill}{rgb}{1.000000,1.000000,1.000000}%
\pgfsetfillcolor{currentfill}%
\pgfsetfillopacity{0.900000}%
\pgfsetlinewidth{0.200750pt}%
\definecolor{currentstroke}{rgb}{0.200000,0.200000,0.200000}%
\pgfsetstrokecolor{currentstroke}%
\pgfsetdash{}{0pt}%
\pgfpathmoveto{\pgfqpoint{2.234955in}{1.017316in}}%
\pgfpathlineto{\pgfqpoint{2.282281in}{1.032952in}}%
\pgfpathlineto{\pgfqpoint{2.329501in}{1.048554in}}%
\pgfpathlineto{\pgfqpoint{2.362013in}{1.026262in}}%
\pgfpathlineto{\pgfqpoint{2.394630in}{1.003898in}}%
\pgfpathlineto{\pgfqpoint{2.347332in}{0.988195in}}%
\pgfpathlineto{\pgfqpoint{2.299927in}{0.972457in}}%
\pgfpathlineto{\pgfqpoint{2.267389in}{0.994923in}}%
\pgfpathlineto{\pgfqpoint{2.234955in}{1.017316in}}%
\pgfpathclose%
\pgfusepath{stroke,fill}%
\end{pgfscope}%
\begin{pgfscope}%
\pgfpathrectangle{\pgfqpoint{0.050000in}{0.050000in}}{\pgfqpoint{4.900000in}{4.900000in}}%
\pgfusepath{clip}%
\pgfsetbuttcap%
\pgfsetroundjoin%
\definecolor{currentfill}{rgb}{1.000000,1.000000,1.000000}%
\pgfsetfillcolor{currentfill}%
\pgfsetfillopacity{0.900000}%
\pgfsetlinewidth{0.200750pt}%
\definecolor{currentstroke}{rgb}{0.200000,0.200000,0.200000}%
\pgfsetstrokecolor{currentstroke}%
\pgfsetdash{}{0pt}%
\pgfpathmoveto{\pgfqpoint{1.821022in}{1.012815in}}%
\pgfpathlineto{\pgfqpoint{1.868829in}{1.028462in}}%
\pgfpathlineto{\pgfqpoint{1.916528in}{1.044074in}}%
\pgfpathlineto{\pgfqpoint{1.948387in}{1.021767in}}%
\pgfpathlineto{\pgfqpoint{1.980349in}{0.999388in}}%
\pgfpathlineto{\pgfqpoint{1.932567in}{0.983676in}}%
\pgfpathlineto{\pgfqpoint{1.884678in}{0.967928in}}%
\pgfpathlineto{\pgfqpoint{1.852799in}{0.990407in}}%
\pgfpathlineto{\pgfqpoint{1.821022in}{1.012815in}}%
\pgfpathclose%
\pgfusepath{stroke,fill}%
\end{pgfscope}%
\begin{pgfscope}%
\pgfpathrectangle{\pgfqpoint{0.050000in}{0.050000in}}{\pgfqpoint{4.900000in}{4.900000in}}%
\pgfusepath{clip}%
\pgfsetbuttcap%
\pgfsetroundjoin%
\definecolor{currentfill}{rgb}{1.000000,1.000000,1.000000}%
\pgfsetfillcolor{currentfill}%
\pgfsetfillopacity{0.900000}%
\pgfsetlinewidth{0.200750pt}%
\definecolor{currentstroke}{rgb}{0.200000,0.200000,0.200000}%
\pgfsetstrokecolor{currentstroke}%
\pgfsetdash{}{0pt}%
\pgfpathmoveto{\pgfqpoint{2.808644in}{1.008598in}}%
\pgfpathlineto{\pgfqpoint{2.855356in}{1.024472in}}%
\pgfpathlineto{\pgfqpoint{2.901964in}{1.040512in}}%
\pgfpathlineto{\pgfqpoint{2.935419in}{1.018294in}}%
\pgfpathlineto{\pgfqpoint{2.968982in}{0.996012in}}%
\pgfpathlineto{\pgfqpoint{2.922298in}{0.979777in}}%
\pgfpathlineto{\pgfqpoint{2.875510in}{0.963746in}}%
\pgfpathlineto{\pgfqpoint{2.842023in}{0.986206in}}%
\pgfpathlineto{\pgfqpoint{2.808644in}{1.008598in}}%
\pgfpathclose%
\pgfusepath{stroke,fill}%
\end{pgfscope}%
\begin{pgfscope}%
\pgfpathrectangle{\pgfqpoint{0.050000in}{0.050000in}}{\pgfqpoint{4.900000in}{4.900000in}}%
\pgfusepath{clip}%
\pgfsetbuttcap%
\pgfsetroundjoin%
\definecolor{currentfill}{rgb}{1.000000,1.000000,1.000000}%
\pgfsetfillcolor{currentfill}%
\pgfsetfillopacity{0.900000}%
\pgfsetlinewidth{0.200750pt}%
\definecolor{currentstroke}{rgb}{0.200000,0.200000,0.200000}%
\pgfsetstrokecolor{currentstroke}%
\pgfsetdash{}{0pt}%
\pgfpathmoveto{\pgfqpoint{2.394630in}{1.003898in}}%
\pgfpathlineto{\pgfqpoint{2.441822in}{1.019565in}}%
\pgfpathlineto{\pgfqpoint{2.488907in}{1.035199in}}%
\pgfpathlineto{\pgfqpoint{2.521707in}{1.012864in}}%
\pgfpathlineto{\pgfqpoint{2.554613in}{0.990457in}}%
\pgfpathlineto{\pgfqpoint{2.507450in}{0.974722in}}%
\pgfpathlineto{\pgfqpoint{2.460180in}{0.958953in}}%
\pgfpathlineto{\pgfqpoint{2.427352in}{0.981462in}}%
\pgfpathlineto{\pgfqpoint{2.394630in}{1.003898in}}%
\pgfpathclose%
\pgfusepath{stroke,fill}%
\end{pgfscope}%
\begin{pgfscope}%
\pgfpathrectangle{\pgfqpoint{0.050000in}{0.050000in}}{\pgfqpoint{4.900000in}{4.900000in}}%
\pgfusepath{clip}%
\pgfsetbuttcap%
\pgfsetroundjoin%
\definecolor{currentfill}{rgb}{1.000000,1.000000,1.000000}%
\pgfsetfillcolor{currentfill}%
\pgfsetfillopacity{0.900000}%
\pgfsetlinewidth{0.200750pt}%
\definecolor{currentstroke}{rgb}{0.200000,0.200000,0.200000}%
\pgfsetstrokecolor{currentstroke}%
\pgfsetdash{}{0pt}%
\pgfpathmoveto{\pgfqpoint{1.980349in}{0.999388in}}%
\pgfpathlineto{\pgfqpoint{2.028022in}{1.015066in}}%
\pgfpathlineto{\pgfqpoint{2.075588in}{1.030708in}}%
\pgfpathlineto{\pgfqpoint{2.107734in}{1.008358in}}%
\pgfpathlineto{\pgfqpoint{2.139983in}{0.985936in}}%
\pgfpathlineto{\pgfqpoint{2.092336in}{0.970193in}}%
\pgfpathlineto{\pgfqpoint{2.044581in}{0.954414in}}%
\pgfpathlineto{\pgfqpoint{2.012413in}{0.976938in}}%
\pgfpathlineto{\pgfqpoint{1.980349in}{0.999388in}}%
\pgfpathclose%
\pgfusepath{stroke,fill}%
\end{pgfscope}%
\begin{pgfscope}%
\pgfpathrectangle{\pgfqpoint{0.050000in}{0.050000in}}{\pgfqpoint{4.900000in}{4.900000in}}%
\pgfusepath{clip}%
\pgfsetbuttcap%
\pgfsetroundjoin%
\definecolor{currentfill}{rgb}{1.000000,1.000000,1.000000}%
\pgfsetfillcolor{currentfill}%
\pgfsetfillopacity{0.900000}%
\pgfsetlinewidth{0.200750pt}%
\definecolor{currentstroke}{rgb}{0.200000,0.200000,0.200000}%
\pgfsetstrokecolor{currentstroke}%
\pgfsetdash{}{0pt}%
\pgfpathmoveto{\pgfqpoint{2.968982in}{0.996012in}}%
\pgfpathlineto{\pgfqpoint{3.015565in}{1.012611in}}%
\pgfpathlineto{\pgfqpoint{3.062050in}{1.029815in}}%
\pgfpathlineto{\pgfqpoint{3.095799in}{1.007765in}}%
\pgfpathlineto{\pgfqpoint{3.129658in}{0.985657in}}%
\pgfpathlineto{\pgfqpoint{3.083094in}{0.968118in}}%
\pgfpathlineto{\pgfqpoint{3.036434in}{0.951256in}}%
\pgfpathlineto{\pgfqpoint{3.002653in}{0.973666in}}%
\pgfpathlineto{\pgfqpoint{2.968982in}{0.996012in}}%
\pgfpathclose%
\pgfusepath{stroke,fill}%
\end{pgfscope}%
\begin{pgfscope}%
\pgfpathrectangle{\pgfqpoint{0.050000in}{0.050000in}}{\pgfqpoint{4.900000in}{4.900000in}}%
\pgfusepath{clip}%
\pgfsetbuttcap%
\pgfsetroundjoin%
\definecolor{currentfill}{rgb}{0.994033,0.993787,0.996432}%
\pgfsetfillcolor{currentfill}%
\pgfsetfillopacity{0.900000}%
\pgfsetlinewidth{0.200750pt}%
\definecolor{currentstroke}{rgb}{0.200000,0.200000,0.200000}%
\pgfsetstrokecolor{currentstroke}%
\pgfsetdash{}{0pt}%
\pgfpathmoveto{\pgfqpoint{3.129658in}{0.985657in}}%
\pgfpathlineto{\pgfqpoint{3.176131in}{1.004220in}}%
\pgfpathlineto{\pgfqpoint{3.222523in}{1.024250in}}%
\pgfpathlineto{\pgfqpoint{3.256574in}{1.002552in}}%
\pgfpathlineto{\pgfqpoint{3.210140in}{0.982266in}}%
\pgfpathlineto{\pgfqpoint{3.163626in}{0.963490in}}%
\pgfpathlineto{\pgfqpoint{3.129658in}{0.985657in}}%
\pgfpathclose%
\pgfusepath{stroke,fill}%
\end{pgfscope}%
\begin{pgfscope}%
\pgfpathrectangle{\pgfqpoint{0.050000in}{0.050000in}}{\pgfqpoint{4.900000in}{4.900000in}}%
\pgfusepath{clip}%
\pgfsetbuttcap%
\pgfsetroundjoin%
\definecolor{currentfill}{rgb}{1.000000,1.000000,1.000000}%
\pgfsetfillcolor{currentfill}%
\pgfsetfillopacity{0.900000}%
\pgfsetlinewidth{0.200750pt}%
\definecolor{currentstroke}{rgb}{0.200000,0.200000,0.200000}%
\pgfsetstrokecolor{currentstroke}%
\pgfsetdash{}{0pt}%
\pgfpathmoveto{\pgfqpoint{2.554613in}{0.990457in}}%
\pgfpathlineto{\pgfqpoint{2.601670in}{1.006161in}}%
\pgfpathlineto{\pgfqpoint{2.648621in}{1.021838in}}%
\pgfpathlineto{\pgfqpoint{2.681710in}{0.999465in}}%
\pgfpathlineto{\pgfqpoint{2.714906in}{0.977021in}}%
\pgfpathlineto{\pgfqpoint{2.667878in}{0.961235in}}%
\pgfpathlineto{\pgfqpoint{2.620744in}{0.945427in}}%
\pgfpathlineto{\pgfqpoint{2.587625in}{0.967979in}}%
\pgfpathlineto{\pgfqpoint{2.554613in}{0.990457in}}%
\pgfpathclose%
\pgfusepath{stroke,fill}%
\end{pgfscope}%
\begin{pgfscope}%
\pgfpathrectangle{\pgfqpoint{0.050000in}{0.050000in}}{\pgfqpoint{4.900000in}{4.900000in}}%
\pgfusepath{clip}%
\pgfsetbuttcap%
\pgfsetroundjoin%
\definecolor{currentfill}{rgb}{1.000000,1.000000,1.000000}%
\pgfsetfillcolor{currentfill}%
\pgfsetfillopacity{0.900000}%
\pgfsetlinewidth{0.200750pt}%
\definecolor{currentstroke}{rgb}{0.200000,0.200000,0.200000}%
\pgfsetstrokecolor{currentstroke}%
\pgfsetdash{}{0pt}%
\pgfpathmoveto{\pgfqpoint{2.139983in}{0.985936in}}%
\pgfpathlineto{\pgfqpoint{2.187523in}{1.001643in}}%
\pgfpathlineto{\pgfqpoint{2.234955in}{1.017316in}}%
\pgfpathlineto{\pgfqpoint{2.267389in}{0.994923in}}%
\pgfpathlineto{\pgfqpoint{2.299927in}{0.972457in}}%
\pgfpathlineto{\pgfqpoint{2.252414in}{0.956684in}}%
\pgfpathlineto{\pgfqpoint{2.204795in}{0.940875in}}%
\pgfpathlineto{\pgfqpoint{2.172337in}{0.963442in}}%
\pgfpathlineto{\pgfqpoint{2.139983in}{0.985936in}}%
\pgfpathclose%
\pgfusepath{stroke,fill}%
\end{pgfscope}%
\begin{pgfscope}%
\pgfpathrectangle{\pgfqpoint{0.050000in}{0.050000in}}{\pgfqpoint{4.900000in}{4.900000in}}%
\pgfusepath{clip}%
\pgfsetbuttcap%
\pgfsetroundjoin%
\definecolor{currentfill}{rgb}{1.000000,1.000000,1.000000}%
\pgfsetfillcolor{currentfill}%
\pgfsetfillopacity{0.900000}%
\pgfsetlinewidth{0.200750pt}%
\definecolor{currentstroke}{rgb}{0.200000,0.200000,0.200000}%
\pgfsetstrokecolor{currentstroke}%
\pgfsetdash{}{0pt}%
\pgfpathmoveto{\pgfqpoint{1.725084in}{0.981415in}}%
\pgfpathlineto{\pgfqpoint{1.773107in}{0.997133in}}%
\pgfpathlineto{\pgfqpoint{1.821022in}{1.012815in}}%
\pgfpathlineto{\pgfqpoint{1.852799in}{0.990407in}}%
\pgfpathlineto{\pgfqpoint{1.884678in}{0.967928in}}%
\pgfpathlineto{\pgfqpoint{1.836680in}{0.952144in}}%
\pgfpathlineto{\pgfqpoint{1.788574in}{0.936324in}}%
\pgfpathlineto{\pgfqpoint{1.756778in}{0.958906in}}%
\pgfpathlineto{\pgfqpoint{1.725084in}{0.981415in}}%
\pgfpathclose%
\pgfusepath{stroke,fill}%
\end{pgfscope}%
\begin{pgfscope}%
\pgfpathrectangle{\pgfqpoint{0.050000in}{0.050000in}}{\pgfqpoint{4.900000in}{4.900000in}}%
\pgfusepath{clip}%
\pgfsetbuttcap%
\pgfsetroundjoin%
\definecolor{currentfill}{rgb}{1.000000,1.000000,1.000000}%
\pgfsetfillcolor{currentfill}%
\pgfsetfillopacity{0.900000}%
\pgfsetlinewidth{0.200750pt}%
\definecolor{currentstroke}{rgb}{0.200000,0.200000,0.200000}%
\pgfsetstrokecolor{currentstroke}%
\pgfsetdash{}{0pt}%
\pgfpathmoveto{\pgfqpoint{2.714906in}{0.977021in}}%
\pgfpathlineto{\pgfqpoint{2.761828in}{0.992798in}}%
\pgfpathlineto{\pgfqpoint{2.808644in}{1.008598in}}%
\pgfpathlineto{\pgfqpoint{2.842023in}{0.986206in}}%
\pgfpathlineto{\pgfqpoint{2.875510in}{0.963746in}}%
\pgfpathlineto{\pgfqpoint{2.828618in}{0.947813in}}%
\pgfpathlineto{\pgfqpoint{2.781620in}{0.931917in}}%
\pgfpathlineto{\pgfqpoint{2.748209in}{0.954505in}}%
\pgfpathlineto{\pgfqpoint{2.714906in}{0.977021in}}%
\pgfpathclose%
\pgfusepath{stroke,fill}%
\end{pgfscope}%
\begin{pgfscope}%
\pgfpathrectangle{\pgfqpoint{0.050000in}{0.050000in}}{\pgfqpoint{4.900000in}{4.900000in}}%
\pgfusepath{clip}%
\pgfsetbuttcap%
\pgfsetroundjoin%
\definecolor{currentfill}{rgb}{1.000000,1.000000,1.000000}%
\pgfsetfillcolor{currentfill}%
\pgfsetfillopacity{0.900000}%
\pgfsetlinewidth{0.200750pt}%
\definecolor{currentstroke}{rgb}{0.200000,0.200000,0.200000}%
\pgfsetstrokecolor{currentstroke}%
\pgfsetdash{}{0pt}%
\pgfpathmoveto{\pgfqpoint{2.299927in}{0.972457in}}%
\pgfpathlineto{\pgfqpoint{2.347332in}{0.988195in}}%
\pgfpathlineto{\pgfqpoint{2.394630in}{1.003898in}}%
\pgfpathlineto{\pgfqpoint{2.427352in}{0.981462in}}%
\pgfpathlineto{\pgfqpoint{2.460180in}{0.958953in}}%
\pgfpathlineto{\pgfqpoint{2.412803in}{0.943149in}}%
\pgfpathlineto{\pgfqpoint{2.365319in}{0.927309in}}%
\pgfpathlineto{\pgfqpoint{2.332570in}{0.949920in}}%
\pgfpathlineto{\pgfqpoint{2.299927in}{0.972457in}}%
\pgfpathclose%
\pgfusepath{stroke,fill}%
\end{pgfscope}%
\begin{pgfscope}%
\pgfpathrectangle{\pgfqpoint{0.050000in}{0.050000in}}{\pgfqpoint{4.900000in}{4.900000in}}%
\pgfusepath{clip}%
\pgfsetbuttcap%
\pgfsetroundjoin%
\definecolor{currentfill}{rgb}{1.000000,1.000000,1.000000}%
\pgfsetfillcolor{currentfill}%
\pgfsetfillopacity{0.900000}%
\pgfsetlinewidth{0.200750pt}%
\definecolor{currentstroke}{rgb}{0.200000,0.200000,0.200000}%
\pgfsetstrokecolor{currentstroke}%
\pgfsetdash{}{0pt}%
\pgfpathmoveto{\pgfqpoint{1.884678in}{0.967928in}}%
\pgfpathlineto{\pgfqpoint{1.932567in}{0.983676in}}%
\pgfpathlineto{\pgfqpoint{1.980349in}{0.999388in}}%
\pgfpathlineto{\pgfqpoint{2.012413in}{0.976938in}}%
\pgfpathlineto{\pgfqpoint{2.044581in}{0.954414in}}%
\pgfpathlineto{\pgfqpoint{1.996718in}{0.938600in}}%
\pgfpathlineto{\pgfqpoint{1.948747in}{0.922750in}}%
\pgfpathlineto{\pgfqpoint{1.916661in}{0.945375in}}%
\pgfpathlineto{\pgfqpoint{1.884678in}{0.967928in}}%
\pgfpathclose%
\pgfusepath{stroke,fill}%
\end{pgfscope}%
\begin{pgfscope}%
\pgfpathrectangle{\pgfqpoint{0.050000in}{0.050000in}}{\pgfqpoint{4.900000in}{4.900000in}}%
\pgfusepath{clip}%
\pgfsetbuttcap%
\pgfsetroundjoin%
\definecolor{currentfill}{rgb}{1.000000,1.000000,1.000000}%
\pgfsetfillcolor{currentfill}%
\pgfsetfillopacity{0.900000}%
\pgfsetlinewidth{0.200750pt}%
\definecolor{currentstroke}{rgb}{0.200000,0.200000,0.200000}%
\pgfsetstrokecolor{currentstroke}%
\pgfsetdash{}{0pt}%
\pgfpathmoveto{\pgfqpoint{2.875510in}{0.963746in}}%
\pgfpathlineto{\pgfqpoint{2.922298in}{0.979777in}}%
\pgfpathlineto{\pgfqpoint{2.968982in}{0.996012in}}%
\pgfpathlineto{\pgfqpoint{3.002653in}{0.973666in}}%
\pgfpathlineto{\pgfqpoint{3.036434in}{0.951256in}}%
\pgfpathlineto{\pgfqpoint{2.989674in}{0.934815in}}%
\pgfpathlineto{\pgfqpoint{2.942810in}{0.918616in}}%
\pgfpathlineto{\pgfqpoint{2.909106in}{0.941216in}}%
\pgfpathlineto{\pgfqpoint{2.875510in}{0.963746in}}%
\pgfpathclose%
\pgfusepath{stroke,fill}%
\end{pgfscope}%
\begin{pgfscope}%
\pgfpathrectangle{\pgfqpoint{0.050000in}{0.050000in}}{\pgfqpoint{4.900000in}{4.900000in}}%
\pgfusepath{clip}%
\pgfsetbuttcap%
\pgfsetroundjoin%
\definecolor{currentfill}{rgb}{1.000000,1.000000,1.000000}%
\pgfsetfillcolor{currentfill}%
\pgfsetfillopacity{0.900000}%
\pgfsetlinewidth{0.200750pt}%
\definecolor{currentstroke}{rgb}{0.200000,0.200000,0.200000}%
\pgfsetstrokecolor{currentstroke}%
\pgfsetdash{}{0pt}%
\pgfpathmoveto{\pgfqpoint{3.036434in}{0.951256in}}%
\pgfpathlineto{\pgfqpoint{3.083094in}{0.968118in}}%
\pgfpathlineto{\pgfqpoint{3.129658in}{0.985657in}}%
\pgfpathlineto{\pgfqpoint{3.163626in}{0.963490in}}%
\pgfpathlineto{\pgfqpoint{3.117023in}{0.945779in}}%
\pgfpathlineto{\pgfqpoint{3.070324in}{0.928781in}}%
\pgfpathlineto{\pgfqpoint{3.036434in}{0.951256in}}%
\pgfpathclose%
\pgfusepath{stroke,fill}%
\end{pgfscope}%
\begin{pgfscope}%
\pgfpathrectangle{\pgfqpoint{0.050000in}{0.050000in}}{\pgfqpoint{4.900000in}{4.900000in}}%
\pgfusepath{clip}%
\pgfsetbuttcap%
\pgfsetroundjoin%
\definecolor{currentfill}{rgb}{1.000000,1.000000,1.000000}%
\pgfsetfillcolor{currentfill}%
\pgfsetfillopacity{0.900000}%
\pgfsetlinewidth{0.200750pt}%
\definecolor{currentstroke}{rgb}{0.200000,0.200000,0.200000}%
\pgfsetstrokecolor{currentstroke}%
\pgfsetdash{}{0pt}%
\pgfpathmoveto{\pgfqpoint{2.460180in}{0.958953in}}%
\pgfpathlineto{\pgfqpoint{2.507450in}{0.974722in}}%
\pgfpathlineto{\pgfqpoint{2.554613in}{0.990457in}}%
\pgfpathlineto{\pgfqpoint{2.587625in}{0.967979in}}%
\pgfpathlineto{\pgfqpoint{2.620744in}{0.945427in}}%
\pgfpathlineto{\pgfqpoint{2.573503in}{0.929589in}}%
\pgfpathlineto{\pgfqpoint{2.526155in}{0.913717in}}%
\pgfpathlineto{\pgfqpoint{2.493114in}{0.936372in}}%
\pgfpathlineto{\pgfqpoint{2.460180in}{0.958953in}}%
\pgfpathclose%
\pgfusepath{stroke,fill}%
\end{pgfscope}%
\begin{pgfscope}%
\pgfpathrectangle{\pgfqpoint{0.050000in}{0.050000in}}{\pgfqpoint{4.900000in}{4.900000in}}%
\pgfusepath{clip}%
\pgfsetbuttcap%
\pgfsetroundjoin%
\definecolor{currentfill}{rgb}{1.000000,1.000000,1.000000}%
\pgfsetfillcolor{currentfill}%
\pgfsetfillopacity{0.900000}%
\pgfsetlinewidth{0.200750pt}%
\definecolor{currentstroke}{rgb}{0.200000,0.200000,0.200000}%
\pgfsetstrokecolor{currentstroke}%
\pgfsetdash{}{0pt}%
\pgfpathmoveto{\pgfqpoint{2.044581in}{0.954414in}}%
\pgfpathlineto{\pgfqpoint{2.092336in}{0.970193in}}%
\pgfpathlineto{\pgfqpoint{2.139983in}{0.985936in}}%
\pgfpathlineto{\pgfqpoint{2.172337in}{0.963442in}}%
\pgfpathlineto{\pgfqpoint{2.204795in}{0.940875in}}%
\pgfpathlineto{\pgfqpoint{2.157067in}{0.925030in}}%
\pgfpathlineto{\pgfqpoint{2.109230in}{0.909148in}}%
\pgfpathlineto{\pgfqpoint{2.076853in}{0.931818in}}%
\pgfpathlineto{\pgfqpoint{2.044581in}{0.954414in}}%
\pgfpathclose%
\pgfusepath{stroke,fill}%
\end{pgfscope}%
\begin{pgfscope}%
\pgfpathrectangle{\pgfqpoint{0.050000in}{0.050000in}}{\pgfqpoint{4.900000in}{4.900000in}}%
\pgfusepath{clip}%
\pgfsetbuttcap%
\pgfsetroundjoin%
\definecolor{currentfill}{rgb}{1.000000,1.000000,1.000000}%
\pgfsetfillcolor{currentfill}%
\pgfsetfillopacity{0.900000}%
\pgfsetlinewidth{0.200750pt}%
\definecolor{currentstroke}{rgb}{0.200000,0.200000,0.200000}%
\pgfsetstrokecolor{currentstroke}%
\pgfsetdash{}{0pt}%
\pgfpathmoveto{\pgfqpoint{2.620744in}{0.945427in}}%
\pgfpathlineto{\pgfqpoint{2.667878in}{0.961235in}}%
\pgfpathlineto{\pgfqpoint{2.714906in}{0.977021in}}%
\pgfpathlineto{\pgfqpoint{2.748209in}{0.954505in}}%
\pgfpathlineto{\pgfqpoint{2.781620in}{0.931917in}}%
\pgfpathlineto{\pgfqpoint{2.734515in}{0.916021in}}%
\pgfpathlineto{\pgfqpoint{2.687304in}{0.900106in}}%
\pgfpathlineto{\pgfqpoint{2.653970in}{0.922803in}}%
\pgfpathlineto{\pgfqpoint{2.620744in}{0.945427in}}%
\pgfpathclose%
\pgfusepath{stroke,fill}%
\end{pgfscope}%
\begin{pgfscope}%
\pgfpathrectangle{\pgfqpoint{0.050000in}{0.050000in}}{\pgfqpoint{4.900000in}{4.900000in}}%
\pgfusepath{clip}%
\pgfsetbuttcap%
\pgfsetroundjoin%
\definecolor{currentfill}{rgb}{1.000000,1.000000,1.000000}%
\pgfsetfillcolor{currentfill}%
\pgfsetfillopacity{0.900000}%
\pgfsetlinewidth{0.200750pt}%
\definecolor{currentstroke}{rgb}{0.200000,0.200000,0.200000}%
\pgfsetstrokecolor{currentstroke}%
\pgfsetdash{}{0pt}%
\pgfpathmoveto{\pgfqpoint{2.204795in}{0.940875in}}%
\pgfpathlineto{\pgfqpoint{2.252414in}{0.956684in}}%
\pgfpathlineto{\pgfqpoint{2.299927in}{0.972457in}}%
\pgfpathlineto{\pgfqpoint{2.332570in}{0.949920in}}%
\pgfpathlineto{\pgfqpoint{2.365319in}{0.927309in}}%
\pgfpathlineto{\pgfqpoint{2.317727in}{0.911433in}}%
\pgfpathlineto{\pgfqpoint{2.270026in}{0.895521in}}%
\pgfpathlineto{\pgfqpoint{2.237358in}{0.918235in}}%
\pgfpathlineto{\pgfqpoint{2.204795in}{0.940875in}}%
\pgfpathclose%
\pgfusepath{stroke,fill}%
\end{pgfscope}%
\begin{pgfscope}%
\pgfpathrectangle{\pgfqpoint{0.050000in}{0.050000in}}{\pgfqpoint{4.900000in}{4.900000in}}%
\pgfusepath{clip}%
\pgfsetbuttcap%
\pgfsetroundjoin%
\definecolor{currentfill}{rgb}{1.000000,1.000000,1.000000}%
\pgfsetfillcolor{currentfill}%
\pgfsetfillopacity{0.900000}%
\pgfsetlinewidth{0.200750pt}%
\definecolor{currentstroke}{rgb}{0.200000,0.200000,0.200000}%
\pgfsetstrokecolor{currentstroke}%
\pgfsetdash{}{0pt}%
\pgfpathmoveto{\pgfqpoint{1.788574in}{0.936324in}}%
\pgfpathlineto{\pgfqpoint{1.836680in}{0.952144in}}%
\pgfpathlineto{\pgfqpoint{1.884678in}{0.967928in}}%
\pgfpathlineto{\pgfqpoint{1.916661in}{0.945375in}}%
\pgfpathlineto{\pgfqpoint{1.948747in}{0.922750in}}%
\pgfpathlineto{\pgfqpoint{1.900666in}{0.906863in}}%
\pgfpathlineto{\pgfqpoint{1.852476in}{0.890941in}}%
\pgfpathlineto{\pgfqpoint{1.820473in}{0.913669in}}%
\pgfpathlineto{\pgfqpoint{1.788574in}{0.936324in}}%
\pgfpathclose%
\pgfusepath{stroke,fill}%
\end{pgfscope}%
\begin{pgfscope}%
\pgfpathrectangle{\pgfqpoint{0.050000in}{0.050000in}}{\pgfqpoint{4.900000in}{4.900000in}}%
\pgfusepath{clip}%
\pgfsetbuttcap%
\pgfsetroundjoin%
\definecolor{currentfill}{rgb}{1.000000,1.000000,1.000000}%
\pgfsetfillcolor{currentfill}%
\pgfsetfillopacity{0.900000}%
\pgfsetlinewidth{0.200750pt}%
\definecolor{currentstroke}{rgb}{0.200000,0.200000,0.200000}%
\pgfsetstrokecolor{currentstroke}%
\pgfsetdash{}{0pt}%
\pgfpathmoveto{\pgfqpoint{2.781620in}{0.931917in}}%
\pgfpathlineto{\pgfqpoint{2.828618in}{0.947813in}}%
\pgfpathlineto{\pgfqpoint{2.875510in}{0.963746in}}%
\pgfpathlineto{\pgfqpoint{2.909106in}{0.941216in}}%
\pgfpathlineto{\pgfqpoint{2.942810in}{0.918616in}}%
\pgfpathlineto{\pgfqpoint{2.895842in}{0.902544in}}%
\pgfpathlineto{\pgfqpoint{2.848767in}{0.886524in}}%
\pgfpathlineto{\pgfqpoint{2.815139in}{0.909256in}}%
\pgfpathlineto{\pgfqpoint{2.781620in}{0.931917in}}%
\pgfpathclose%
\pgfusepath{stroke,fill}%
\end{pgfscope}%
\begin{pgfscope}%
\pgfpathrectangle{\pgfqpoint{0.050000in}{0.050000in}}{\pgfqpoint{4.900000in}{4.900000in}}%
\pgfusepath{clip}%
\pgfsetbuttcap%
\pgfsetroundjoin%
\definecolor{currentfill}{rgb}{1.000000,1.000000,1.000000}%
\pgfsetfillcolor{currentfill}%
\pgfsetfillopacity{0.900000}%
\pgfsetlinewidth{0.200750pt}%
\definecolor{currentstroke}{rgb}{0.200000,0.200000,0.200000}%
\pgfsetstrokecolor{currentstroke}%
\pgfsetdash{}{0pt}%
\pgfpathmoveto{\pgfqpoint{2.942810in}{0.918616in}}%
\pgfpathlineto{\pgfqpoint{2.989674in}{0.934815in}}%
\pgfpathlineto{\pgfqpoint{3.036434in}{0.951256in}}%
\pgfpathlineto{\pgfqpoint{3.070324in}{0.928781in}}%
\pgfpathlineto{\pgfqpoint{3.023526in}{0.912232in}}%
\pgfpathlineto{\pgfqpoint{2.976624in}{0.895947in}}%
\pgfpathlineto{\pgfqpoint{2.942810in}{0.918616in}}%
\pgfpathclose%
\pgfusepath{stroke,fill}%
\end{pgfscope}%
\begin{pgfscope}%
\pgfpathrectangle{\pgfqpoint{0.050000in}{0.050000in}}{\pgfqpoint{4.900000in}{4.900000in}}%
\pgfusepath{clip}%
\pgfsetbuttcap%
\pgfsetroundjoin%
\definecolor{currentfill}{rgb}{1.000000,1.000000,1.000000}%
\pgfsetfillcolor{currentfill}%
\pgfsetfillopacity{0.900000}%
\pgfsetlinewidth{0.200750pt}%
\definecolor{currentstroke}{rgb}{0.200000,0.200000,0.200000}%
\pgfsetstrokecolor{currentstroke}%
\pgfsetdash{}{0pt}%
\pgfpathmoveto{\pgfqpoint{2.365319in}{0.927309in}}%
\pgfpathlineto{\pgfqpoint{2.412803in}{0.943149in}}%
\pgfpathlineto{\pgfqpoint{2.460180in}{0.958953in}}%
\pgfpathlineto{\pgfqpoint{2.493114in}{0.936372in}}%
\pgfpathlineto{\pgfqpoint{2.526155in}{0.913717in}}%
\pgfpathlineto{\pgfqpoint{2.478699in}{0.897810in}}%
\pgfpathlineto{\pgfqpoint{2.431136in}{0.881867in}}%
\pgfpathlineto{\pgfqpoint{2.398174in}{0.904625in}}%
\pgfpathlineto{\pgfqpoint{2.365319in}{0.927309in}}%
\pgfpathclose%
\pgfusepath{stroke,fill}%
\end{pgfscope}%
\begin{pgfscope}%
\pgfpathrectangle{\pgfqpoint{0.050000in}{0.050000in}}{\pgfqpoint{4.900000in}{4.900000in}}%
\pgfusepath{clip}%
\pgfsetbuttcap%
\pgfsetroundjoin%
\definecolor{currentfill}{rgb}{1.000000,1.000000,1.000000}%
\pgfsetfillcolor{currentfill}%
\pgfsetfillopacity{0.900000}%
\pgfsetlinewidth{0.200750pt}%
\definecolor{currentstroke}{rgb}{0.200000,0.200000,0.200000}%
\pgfsetstrokecolor{currentstroke}%
\pgfsetdash{}{0pt}%
\pgfpathmoveto{\pgfqpoint{1.948747in}{0.922750in}}%
\pgfpathlineto{\pgfqpoint{1.996718in}{0.938600in}}%
\pgfpathlineto{\pgfqpoint{2.044581in}{0.954414in}}%
\pgfpathlineto{\pgfqpoint{2.076853in}{0.931818in}}%
\pgfpathlineto{\pgfqpoint{2.109230in}{0.909148in}}%
\pgfpathlineto{\pgfqpoint{2.061285in}{0.893231in}}%
\pgfpathlineto{\pgfqpoint{2.013231in}{0.877278in}}%
\pgfpathlineto{\pgfqpoint{1.980937in}{0.900051in}}%
\pgfpathlineto{\pgfqpoint{1.948747in}{0.922750in}}%
\pgfpathclose%
\pgfusepath{stroke,fill}%
\end{pgfscope}%
\begin{pgfscope}%
\pgfpathrectangle{\pgfqpoint{0.050000in}{0.050000in}}{\pgfqpoint{4.900000in}{4.900000in}}%
\pgfusepath{clip}%
\pgfsetbuttcap%
\pgfsetroundjoin%
\definecolor{currentfill}{rgb}{1.000000,1.000000,1.000000}%
\pgfsetfillcolor{currentfill}%
\pgfsetfillopacity{0.900000}%
\pgfsetlinewidth{0.200750pt}%
\definecolor{currentstroke}{rgb}{0.200000,0.200000,0.200000}%
\pgfsetstrokecolor{currentstroke}%
\pgfsetdash{}{0pt}%
\pgfpathmoveto{\pgfqpoint{2.526155in}{0.913717in}}%
\pgfpathlineto{\pgfqpoint{2.573503in}{0.929589in}}%
\pgfpathlineto{\pgfqpoint{2.620744in}{0.945427in}}%
\pgfpathlineto{\pgfqpoint{2.653970in}{0.922803in}}%
\pgfpathlineto{\pgfqpoint{2.687304in}{0.900106in}}%
\pgfpathlineto{\pgfqpoint{2.639985in}{0.884163in}}%
\pgfpathlineto{\pgfqpoint{2.592559in}{0.868187in}}%
\pgfpathlineto{\pgfqpoint{2.559303in}{0.890989in}}%
\pgfpathlineto{\pgfqpoint{2.526155in}{0.913717in}}%
\pgfpathclose%
\pgfusepath{stroke,fill}%
\end{pgfscope}%
\begin{pgfscope}%
\pgfpathrectangle{\pgfqpoint{0.050000in}{0.050000in}}{\pgfqpoint{4.900000in}{4.900000in}}%
\pgfusepath{clip}%
\pgfsetbuttcap%
\pgfsetroundjoin%
\definecolor{currentfill}{rgb}{1.000000,1.000000,1.000000}%
\pgfsetfillcolor{currentfill}%
\pgfsetfillopacity{0.900000}%
\pgfsetlinewidth{0.200750pt}%
\definecolor{currentstroke}{rgb}{0.200000,0.200000,0.200000}%
\pgfsetstrokecolor{currentstroke}%
\pgfsetdash{}{0pt}%
\pgfpathmoveto{\pgfqpoint{2.109230in}{0.909148in}}%
\pgfpathlineto{\pgfqpoint{2.157067in}{0.925030in}}%
\pgfpathlineto{\pgfqpoint{2.204795in}{0.940875in}}%
\pgfpathlineto{\pgfqpoint{2.237358in}{0.918235in}}%
\pgfpathlineto{\pgfqpoint{2.270026in}{0.895521in}}%
\pgfpathlineto{\pgfqpoint{2.222218in}{0.879573in}}%
\pgfpathlineto{\pgfqpoint{2.174300in}{0.863588in}}%
\pgfpathlineto{\pgfqpoint{2.141712in}{0.886405in}}%
\pgfpathlineto{\pgfqpoint{2.109230in}{0.909148in}}%
\pgfpathclose%
\pgfusepath{stroke,fill}%
\end{pgfscope}%
\begin{pgfscope}%
\pgfpathrectangle{\pgfqpoint{0.050000in}{0.050000in}}{\pgfqpoint{4.900000in}{4.900000in}}%
\pgfusepath{clip}%
\pgfsetbuttcap%
\pgfsetroundjoin%
\definecolor{currentfill}{rgb}{1.000000,1.000000,1.000000}%
\pgfsetfillcolor{currentfill}%
\pgfsetfillopacity{0.900000}%
\pgfsetlinewidth{0.200750pt}%
\definecolor{currentstroke}{rgb}{0.200000,0.200000,0.200000}%
\pgfsetstrokecolor{currentstroke}%
\pgfsetdash{}{0pt}%
\pgfpathmoveto{\pgfqpoint{2.687304in}{0.900106in}}%
\pgfpathlineto{\pgfqpoint{2.734515in}{0.916021in}}%
\pgfpathlineto{\pgfqpoint{2.781620in}{0.931917in}}%
\pgfpathlineto{\pgfqpoint{2.815139in}{0.909256in}}%
\pgfpathlineto{\pgfqpoint{2.848767in}{0.886524in}}%
\pgfpathlineto{\pgfqpoint{2.801586in}{0.870514in}}%
\pgfpathlineto{\pgfqpoint{2.754297in}{0.854491in}}%
\pgfpathlineto{\pgfqpoint{2.720746in}{0.877335in}}%
\pgfpathlineto{\pgfqpoint{2.687304in}{0.900106in}}%
\pgfpathclose%
\pgfusepath{stroke,fill}%
\end{pgfscope}%
\begin{pgfscope}%
\pgfpathrectangle{\pgfqpoint{0.050000in}{0.050000in}}{\pgfqpoint{4.900000in}{4.900000in}}%
\pgfusepath{clip}%
\pgfsetbuttcap%
\pgfsetroundjoin%
\definecolor{currentfill}{rgb}{1.000000,1.000000,1.000000}%
\pgfsetfillcolor{currentfill}%
\pgfsetfillopacity{0.900000}%
\pgfsetlinewidth{0.200750pt}%
\definecolor{currentstroke}{rgb}{0.200000,0.200000,0.200000}%
\pgfsetstrokecolor{currentstroke}%
\pgfsetdash{}{0pt}%
\pgfpathmoveto{\pgfqpoint{2.848767in}{0.886524in}}%
\pgfpathlineto{\pgfqpoint{2.895842in}{0.902544in}}%
\pgfpathlineto{\pgfqpoint{2.942810in}{0.918616in}}%
\pgfpathlineto{\pgfqpoint{2.976624in}{0.895947in}}%
\pgfpathlineto{\pgfqpoint{2.929617in}{0.879802in}}%
\pgfpathlineto{\pgfqpoint{2.882505in}{0.863719in}}%
\pgfpathlineto{\pgfqpoint{2.848767in}{0.886524in}}%
\pgfpathclose%
\pgfusepath{stroke,fill}%
\end{pgfscope}%
\begin{pgfscope}%
\pgfpathrectangle{\pgfqpoint{0.050000in}{0.050000in}}{\pgfqpoint{4.900000in}{4.900000in}}%
\pgfusepath{clip}%
\pgfsetbuttcap%
\pgfsetroundjoin%
\definecolor{currentfill}{rgb}{1.000000,1.000000,1.000000}%
\pgfsetfillcolor{currentfill}%
\pgfsetfillopacity{0.900000}%
\pgfsetlinewidth{0.200750pt}%
\definecolor{currentstroke}{rgb}{0.200000,0.200000,0.200000}%
\pgfsetstrokecolor{currentstroke}%
\pgfsetdash{}{0pt}%
\pgfpathmoveto{\pgfqpoint{2.270026in}{0.895521in}}%
\pgfpathlineto{\pgfqpoint{2.317727in}{0.911433in}}%
\pgfpathlineto{\pgfqpoint{2.365319in}{0.927309in}}%
\pgfpathlineto{\pgfqpoint{2.398174in}{0.904625in}}%
\pgfpathlineto{\pgfqpoint{2.431136in}{0.881867in}}%
\pgfpathlineto{\pgfqpoint{2.383464in}{0.865888in}}%
\pgfpathlineto{\pgfqpoint{2.335683in}{0.849872in}}%
\pgfpathlineto{\pgfqpoint{2.302801in}{0.872733in}}%
\pgfpathlineto{\pgfqpoint{2.270026in}{0.895521in}}%
\pgfpathclose%
\pgfusepath{stroke,fill}%
\end{pgfscope}%
\begin{pgfscope}%
\pgfpathrectangle{\pgfqpoint{0.050000in}{0.050000in}}{\pgfqpoint{4.900000in}{4.900000in}}%
\pgfusepath{clip}%
\pgfsetbuttcap%
\pgfsetroundjoin%
\definecolor{currentfill}{rgb}{1.000000,1.000000,1.000000}%
\pgfsetfillcolor{currentfill}%
\pgfsetfillopacity{0.900000}%
\pgfsetlinewidth{0.200750pt}%
\definecolor{currentstroke}{rgb}{0.200000,0.200000,0.200000}%
\pgfsetstrokecolor{currentstroke}%
\pgfsetdash{}{0pt}%
\pgfpathmoveto{\pgfqpoint{1.852476in}{0.890941in}}%
\pgfpathlineto{\pgfqpoint{1.900666in}{0.906863in}}%
\pgfpathlineto{\pgfqpoint{1.948747in}{0.922750in}}%
\pgfpathlineto{\pgfqpoint{1.980937in}{0.900051in}}%
\pgfpathlineto{\pgfqpoint{2.013231in}{0.877278in}}%
\pgfpathlineto{\pgfqpoint{1.965068in}{0.861288in}}%
\pgfpathlineto{\pgfqpoint{1.916795in}{0.845262in}}%
\pgfpathlineto{\pgfqpoint{1.884583in}{0.868139in}}%
\pgfpathlineto{\pgfqpoint{1.852476in}{0.890941in}}%
\pgfpathclose%
\pgfusepath{stroke,fill}%
\end{pgfscope}%
\begin{pgfscope}%
\pgfpathrectangle{\pgfqpoint{0.050000in}{0.050000in}}{\pgfqpoint{4.900000in}{4.900000in}}%
\pgfusepath{clip}%
\pgfsetbuttcap%
\pgfsetroundjoin%
\definecolor{currentfill}{rgb}{1.000000,1.000000,1.000000}%
\pgfsetfillcolor{currentfill}%
\pgfsetfillopacity{0.900000}%
\pgfsetlinewidth{0.200750pt}%
\definecolor{currentstroke}{rgb}{0.200000,0.200000,0.200000}%
\pgfsetstrokecolor{currentstroke}%
\pgfsetdash{}{0pt}%
\pgfpathmoveto{\pgfqpoint{2.431136in}{0.881867in}}%
\pgfpathlineto{\pgfqpoint{2.478699in}{0.897810in}}%
\pgfpathlineto{\pgfqpoint{2.526155in}{0.913717in}}%
\pgfpathlineto{\pgfqpoint{2.559303in}{0.890989in}}%
\pgfpathlineto{\pgfqpoint{2.592559in}{0.868187in}}%
\pgfpathlineto{\pgfqpoint{2.545025in}{0.852176in}}%
\pgfpathlineto{\pgfqpoint{2.497382in}{0.836129in}}%
\pgfpathlineto{\pgfqpoint{2.464205in}{0.859035in}}%
\pgfpathlineto{\pgfqpoint{2.431136in}{0.881867in}}%
\pgfpathclose%
\pgfusepath{stroke,fill}%
\end{pgfscope}%
\begin{pgfscope}%
\pgfpathrectangle{\pgfqpoint{0.050000in}{0.050000in}}{\pgfqpoint{4.900000in}{4.900000in}}%
\pgfusepath{clip}%
\pgfsetbuttcap%
\pgfsetroundjoin%
\definecolor{currentfill}{rgb}{1.000000,1.000000,1.000000}%
\pgfsetfillcolor{currentfill}%
\pgfsetfillopacity{0.900000}%
\pgfsetlinewidth{0.200750pt}%
\definecolor{currentstroke}{rgb}{0.200000,0.200000,0.200000}%
\pgfsetstrokecolor{currentstroke}%
\pgfsetdash{}{0pt}%
\pgfpathmoveto{\pgfqpoint{2.013231in}{0.877278in}}%
\pgfpathlineto{\pgfqpoint{2.061285in}{0.893231in}}%
\pgfpathlineto{\pgfqpoint{2.109230in}{0.909148in}}%
\pgfpathlineto{\pgfqpoint{2.141712in}{0.886405in}}%
\pgfpathlineto{\pgfqpoint{2.174300in}{0.863588in}}%
\pgfpathlineto{\pgfqpoint{2.126273in}{0.847567in}}%
\pgfpathlineto{\pgfqpoint{2.078137in}{0.831510in}}%
\pgfpathlineto{\pgfqpoint{2.045631in}{0.854431in}}%
\pgfpathlineto{\pgfqpoint{2.013231in}{0.877278in}}%
\pgfpathclose%
\pgfusepath{stroke,fill}%
\end{pgfscope}%
\begin{pgfscope}%
\pgfpathrectangle{\pgfqpoint{0.050000in}{0.050000in}}{\pgfqpoint{4.900000in}{4.900000in}}%
\pgfusepath{clip}%
\pgfsetbuttcap%
\pgfsetroundjoin%
\definecolor{currentfill}{rgb}{1.000000,1.000000,1.000000}%
\pgfsetfillcolor{currentfill}%
\pgfsetfillopacity{0.900000}%
\pgfsetlinewidth{0.200750pt}%
\definecolor{currentstroke}{rgb}{0.200000,0.200000,0.200000}%
\pgfsetstrokecolor{currentstroke}%
\pgfsetdash{}{0pt}%
\pgfpathmoveto{\pgfqpoint{2.592559in}{0.868187in}}%
\pgfpathlineto{\pgfqpoint{2.639985in}{0.884163in}}%
\pgfpathlineto{\pgfqpoint{2.687304in}{0.900106in}}%
\pgfpathlineto{\pgfqpoint{2.720746in}{0.877335in}}%
\pgfpathlineto{\pgfqpoint{2.754297in}{0.854491in}}%
\pgfpathlineto{\pgfqpoint{2.706901in}{0.838442in}}%
\pgfpathlineto{\pgfqpoint{2.659396in}{0.822360in}}%
\pgfpathlineto{\pgfqpoint{2.625923in}{0.845311in}}%
\pgfpathlineto{\pgfqpoint{2.592559in}{0.868187in}}%
\pgfpathclose%
\pgfusepath{stroke,fill}%
\end{pgfscope}%
\begin{pgfscope}%
\pgfpathrectangle{\pgfqpoint{0.050000in}{0.050000in}}{\pgfqpoint{4.900000in}{4.900000in}}%
\pgfusepath{clip}%
\pgfsetbuttcap%
\pgfsetroundjoin%
\definecolor{currentfill}{rgb}{1.000000,1.000000,1.000000}%
\pgfsetfillcolor{currentfill}%
\pgfsetfillopacity{0.900000}%
\pgfsetlinewidth{0.200750pt}%
\definecolor{currentstroke}{rgb}{0.200000,0.200000,0.200000}%
\pgfsetstrokecolor{currentstroke}%
\pgfsetdash{}{0pt}%
\pgfpathmoveto{\pgfqpoint{2.754297in}{0.854491in}}%
\pgfpathlineto{\pgfqpoint{2.801586in}{0.870514in}}%
\pgfpathlineto{\pgfqpoint{2.848767in}{0.886524in}}%
\pgfpathlineto{\pgfqpoint{2.882505in}{0.863719in}}%
\pgfpathlineto{\pgfqpoint{2.835285in}{0.847651in}}%
\pgfpathlineto{\pgfqpoint{2.787958in}{0.831572in}}%
\pgfpathlineto{\pgfqpoint{2.754297in}{0.854491in}}%
\pgfpathclose%
\pgfusepath{stroke,fill}%
\end{pgfscope}%
\begin{pgfscope}%
\pgfpathrectangle{\pgfqpoint{0.050000in}{0.050000in}}{\pgfqpoint{4.900000in}{4.900000in}}%
\pgfusepath{clip}%
\pgfsetbuttcap%
\pgfsetroundjoin%
\definecolor{currentfill}{rgb}{1.000000,1.000000,1.000000}%
\pgfsetfillcolor{currentfill}%
\pgfsetfillopacity{0.900000}%
\pgfsetlinewidth{0.200750pt}%
\definecolor{currentstroke}{rgb}{0.200000,0.200000,0.200000}%
\pgfsetstrokecolor{currentstroke}%
\pgfsetdash{}{0pt}%
\pgfpathmoveto{\pgfqpoint{2.174300in}{0.863588in}}%
\pgfpathlineto{\pgfqpoint{2.222218in}{0.879573in}}%
\pgfpathlineto{\pgfqpoint{2.270026in}{0.895521in}}%
\pgfpathlineto{\pgfqpoint{2.302801in}{0.872733in}}%
\pgfpathlineto{\pgfqpoint{2.335683in}{0.849872in}}%
\pgfpathlineto{\pgfqpoint{2.287793in}{0.833820in}}%
\pgfpathlineto{\pgfqpoint{2.239794in}{0.817731in}}%
\pgfpathlineto{\pgfqpoint{2.206994in}{0.840697in}}%
\pgfpathlineto{\pgfqpoint{2.174300in}{0.863588in}}%
\pgfpathclose%
\pgfusepath{stroke,fill}%
\end{pgfscope}%
\begin{pgfscope}%
\pgfpathrectangle{\pgfqpoint{0.050000in}{0.050000in}}{\pgfqpoint{4.900000in}{4.900000in}}%
\pgfusepath{clip}%
\pgfsetbuttcap%
\pgfsetroundjoin%
\definecolor{currentfill}{rgb}{1.000000,1.000000,1.000000}%
\pgfsetfillcolor{currentfill}%
\pgfsetfillopacity{0.900000}%
\pgfsetlinewidth{0.200750pt}%
\definecolor{currentstroke}{rgb}{0.200000,0.200000,0.200000}%
\pgfsetstrokecolor{currentstroke}%
\pgfsetdash{}{0pt}%
\pgfpathmoveto{\pgfqpoint{2.335683in}{0.849872in}}%
\pgfpathlineto{\pgfqpoint{2.383464in}{0.865888in}}%
\pgfpathlineto{\pgfqpoint{2.431136in}{0.881867in}}%
\pgfpathlineto{\pgfqpoint{2.464205in}{0.859035in}}%
\pgfpathlineto{\pgfqpoint{2.497382in}{0.836129in}}%
\pgfpathlineto{\pgfqpoint{2.449630in}{0.820045in}}%
\pgfpathlineto{\pgfqpoint{2.401769in}{0.803925in}}%
\pgfpathlineto{\pgfqpoint{2.368672in}{0.826936in}}%
\pgfpathlineto{\pgfqpoint{2.335683in}{0.849872in}}%
\pgfpathclose%
\pgfusepath{stroke,fill}%
\end{pgfscope}%
\begin{pgfscope}%
\pgfpathrectangle{\pgfqpoint{0.050000in}{0.050000in}}{\pgfqpoint{4.900000in}{4.900000in}}%
\pgfusepath{clip}%
\pgfsetbuttcap%
\pgfsetroundjoin%
\definecolor{currentfill}{rgb}{1.000000,1.000000,1.000000}%
\pgfsetfillcolor{currentfill}%
\pgfsetfillopacity{0.900000}%
\pgfsetlinewidth{0.200750pt}%
\definecolor{currentstroke}{rgb}{0.200000,0.200000,0.200000}%
\pgfsetstrokecolor{currentstroke}%
\pgfsetdash{}{0pt}%
\pgfpathmoveto{\pgfqpoint{1.916795in}{0.845262in}}%
\pgfpathlineto{\pgfqpoint{1.965068in}{0.861288in}}%
\pgfpathlineto{\pgfqpoint{2.013231in}{0.877278in}}%
\pgfpathlineto{\pgfqpoint{2.045631in}{0.854431in}}%
\pgfpathlineto{\pgfqpoint{2.078137in}{0.831510in}}%
\pgfpathlineto{\pgfqpoint{2.029890in}{0.815416in}}%
\pgfpathlineto{\pgfqpoint{1.981533in}{0.799285in}}%
\pgfpathlineto{\pgfqpoint{1.949111in}{0.822311in}}%
\pgfpathlineto{\pgfqpoint{1.916795in}{0.845262in}}%
\pgfpathclose%
\pgfusepath{stroke,fill}%
\end{pgfscope}%
\begin{pgfscope}%
\pgfpathrectangle{\pgfqpoint{0.050000in}{0.050000in}}{\pgfqpoint{4.900000in}{4.900000in}}%
\pgfusepath{clip}%
\pgfsetbuttcap%
\pgfsetroundjoin%
\definecolor{currentfill}{rgb}{1.000000,1.000000,1.000000}%
\pgfsetfillcolor{currentfill}%
\pgfsetfillopacity{0.900000}%
\pgfsetlinewidth{0.200750pt}%
\definecolor{currentstroke}{rgb}{0.200000,0.200000,0.200000}%
\pgfsetstrokecolor{currentstroke}%
\pgfsetdash{}{0pt}%
\pgfpathmoveto{\pgfqpoint{2.497382in}{0.836129in}}%
\pgfpathlineto{\pgfqpoint{2.545025in}{0.852176in}}%
\pgfpathlineto{\pgfqpoint{2.592559in}{0.868187in}}%
\pgfpathlineto{\pgfqpoint{2.625923in}{0.845311in}}%
\pgfpathlineto{\pgfqpoint{2.659396in}{0.822360in}}%
\pgfpathlineto{\pgfqpoint{2.611783in}{0.806244in}}%
\pgfpathlineto{\pgfqpoint{2.564061in}{0.790092in}}%
\pgfpathlineto{\pgfqpoint{2.530667in}{0.813148in}}%
\pgfpathlineto{\pgfqpoint{2.497382in}{0.836129in}}%
\pgfpathclose%
\pgfusepath{stroke,fill}%
\end{pgfscope}%
\begin{pgfscope}%
\pgfpathrectangle{\pgfqpoint{0.050000in}{0.050000in}}{\pgfqpoint{4.900000in}{4.900000in}}%
\pgfusepath{clip}%
\pgfsetbuttcap%
\pgfsetroundjoin%
\definecolor{currentfill}{rgb}{1.000000,1.000000,1.000000}%
\pgfsetfillcolor{currentfill}%
\pgfsetfillopacity{0.900000}%
\pgfsetlinewidth{0.200750pt}%
\definecolor{currentstroke}{rgb}{0.200000,0.200000,0.200000}%
\pgfsetstrokecolor{currentstroke}%
\pgfsetdash{}{0pt}%
\pgfpathmoveto{\pgfqpoint{2.659396in}{0.822360in}}%
\pgfpathlineto{\pgfqpoint{2.706901in}{0.838442in}}%
\pgfpathlineto{\pgfqpoint{2.754297in}{0.854491in}}%
\pgfpathlineto{\pgfqpoint{2.787958in}{0.831572in}}%
\pgfpathlineto{\pgfqpoint{2.740522in}{0.815469in}}%
\pgfpathlineto{\pgfqpoint{2.692979in}{0.799335in}}%
\pgfpathlineto{\pgfqpoint{2.659396in}{0.822360in}}%
\pgfpathclose%
\pgfusepath{stroke,fill}%
\end{pgfscope}%
\begin{pgfscope}%
\pgfpathrectangle{\pgfqpoint{0.050000in}{0.050000in}}{\pgfqpoint{4.900000in}{4.900000in}}%
\pgfusepath{clip}%
\pgfsetbuttcap%
\pgfsetroundjoin%
\definecolor{currentfill}{rgb}{1.000000,1.000000,1.000000}%
\pgfsetfillcolor{currentfill}%
\pgfsetfillopacity{0.900000}%
\pgfsetlinewidth{0.200750pt}%
\definecolor{currentstroke}{rgb}{0.200000,0.200000,0.200000}%
\pgfsetstrokecolor{currentstroke}%
\pgfsetdash{}{0pt}%
\pgfpathmoveto{\pgfqpoint{2.078137in}{0.831510in}}%
\pgfpathlineto{\pgfqpoint{2.126273in}{0.847567in}}%
\pgfpathlineto{\pgfqpoint{2.174300in}{0.863588in}}%
\pgfpathlineto{\pgfqpoint{2.206994in}{0.840697in}}%
\pgfpathlineto{\pgfqpoint{2.239794in}{0.817731in}}%
\pgfpathlineto{\pgfqpoint{2.191686in}{0.801605in}}%
\pgfpathlineto{\pgfqpoint{2.143467in}{0.785442in}}%
\pgfpathlineto{\pgfqpoint{2.110748in}{0.808514in}}%
\pgfpathlineto{\pgfqpoint{2.078137in}{0.831510in}}%
\pgfpathclose%
\pgfusepath{stroke,fill}%
\end{pgfscope}%
\begin{pgfscope}%
\pgfpathrectangle{\pgfqpoint{0.050000in}{0.050000in}}{\pgfqpoint{4.900000in}{4.900000in}}%
\pgfusepath{clip}%
\pgfsetbuttcap%
\pgfsetroundjoin%
\definecolor{currentfill}{rgb}{1.000000,1.000000,1.000000}%
\pgfsetfillcolor{currentfill}%
\pgfsetfillopacity{0.900000}%
\pgfsetlinewidth{0.200750pt}%
\definecolor{currentstroke}{rgb}{0.200000,0.200000,0.200000}%
\pgfsetstrokecolor{currentstroke}%
\pgfsetdash{}{0pt}%
\pgfpathmoveto{\pgfqpoint{2.239794in}{0.817731in}}%
\pgfpathlineto{\pgfqpoint{2.287793in}{0.833820in}}%
\pgfpathlineto{\pgfqpoint{2.335683in}{0.849872in}}%
\pgfpathlineto{\pgfqpoint{2.368672in}{0.826936in}}%
\pgfpathlineto{\pgfqpoint{2.401769in}{0.803925in}}%
\pgfpathlineto{\pgfqpoint{2.353798in}{0.787767in}}%
\pgfpathlineto{\pgfqpoint{2.305718in}{0.771573in}}%
\pgfpathlineto{\pgfqpoint{2.272702in}{0.794690in}}%
\pgfpathlineto{\pgfqpoint{2.239794in}{0.817731in}}%
\pgfpathclose%
\pgfusepath{stroke,fill}%
\end{pgfscope}%
\begin{pgfscope}%
\pgfpathrectangle{\pgfqpoint{0.050000in}{0.050000in}}{\pgfqpoint{4.900000in}{4.900000in}}%
\pgfusepath{clip}%
\pgfsetbuttcap%
\pgfsetroundjoin%
\definecolor{currentfill}{rgb}{1.000000,1.000000,1.000000}%
\pgfsetfillcolor{currentfill}%
\pgfsetfillopacity{0.900000}%
\pgfsetlinewidth{0.200750pt}%
\definecolor{currentstroke}{rgb}{0.200000,0.200000,0.200000}%
\pgfsetstrokecolor{currentstroke}%
\pgfsetdash{}{0pt}%
\pgfpathmoveto{\pgfqpoint{2.401769in}{0.803925in}}%
\pgfpathlineto{\pgfqpoint{2.449630in}{0.820045in}}%
\pgfpathlineto{\pgfqpoint{2.497382in}{0.836129in}}%
\pgfpathlineto{\pgfqpoint{2.530667in}{0.813148in}}%
\pgfpathlineto{\pgfqpoint{2.564061in}{0.790092in}}%
\pgfpathlineto{\pgfqpoint{2.516229in}{0.773903in}}%
\pgfpathlineto{\pgfqpoint{2.468287in}{0.757676in}}%
\pgfpathlineto{\pgfqpoint{2.434974in}{0.780838in}}%
\pgfpathlineto{\pgfqpoint{2.401769in}{0.803925in}}%
\pgfpathclose%
\pgfusepath{stroke,fill}%
\end{pgfscope}%
\begin{pgfscope}%
\pgfpathrectangle{\pgfqpoint{0.050000in}{0.050000in}}{\pgfqpoint{4.900000in}{4.900000in}}%
\pgfusepath{clip}%
\pgfsetbuttcap%
\pgfsetroundjoin%
\definecolor{currentfill}{rgb}{1.000000,1.000000,1.000000}%
\pgfsetfillcolor{currentfill}%
\pgfsetfillopacity{0.900000}%
\pgfsetlinewidth{0.200750pt}%
\definecolor{currentstroke}{rgb}{0.200000,0.200000,0.200000}%
\pgfsetstrokecolor{currentstroke}%
\pgfsetdash{}{0pt}%
\pgfpathmoveto{\pgfqpoint{2.564061in}{0.790092in}}%
\pgfpathlineto{\pgfqpoint{2.611783in}{0.806244in}}%
\pgfpathlineto{\pgfqpoint{2.659396in}{0.822360in}}%
\pgfpathlineto{\pgfqpoint{2.692979in}{0.799335in}}%
\pgfpathlineto{\pgfqpoint{2.645326in}{0.783165in}}%
\pgfpathlineto{\pgfqpoint{2.597564in}{0.766960in}}%
\pgfpathlineto{\pgfqpoint{2.564061in}{0.790092in}}%
\pgfpathclose%
\pgfusepath{stroke,fill}%
\end{pgfscope}%
\begin{pgfscope}%
\pgfpathrectangle{\pgfqpoint{0.050000in}{0.050000in}}{\pgfqpoint{4.900000in}{4.900000in}}%
\pgfusepath{clip}%
\pgfsetbuttcap%
\pgfsetroundjoin%
\definecolor{currentfill}{rgb}{1.000000,1.000000,1.000000}%
\pgfsetfillcolor{currentfill}%
\pgfsetfillopacity{0.900000}%
\pgfsetlinewidth{0.200750pt}%
\definecolor{currentstroke}{rgb}{0.200000,0.200000,0.200000}%
\pgfsetstrokecolor{currentstroke}%
\pgfsetdash{}{0pt}%
\pgfpathmoveto{\pgfqpoint{1.981533in}{0.799285in}}%
\pgfpathlineto{\pgfqpoint{2.029890in}{0.815416in}}%
\pgfpathlineto{\pgfqpoint{2.078137in}{0.831510in}}%
\pgfpathlineto{\pgfqpoint{2.110748in}{0.808514in}}%
\pgfpathlineto{\pgfqpoint{2.143467in}{0.785442in}}%
\pgfpathlineto{\pgfqpoint{2.095137in}{0.769243in}}%
\pgfpathlineto{\pgfqpoint{2.046696in}{0.753006in}}%
\pgfpathlineto{\pgfqpoint{2.014062in}{0.776183in}}%
\pgfpathlineto{\pgfqpoint{1.981533in}{0.799285in}}%
\pgfpathclose%
\pgfusepath{stroke,fill}%
\end{pgfscope}%
\begin{pgfscope}%
\pgfpathrectangle{\pgfqpoint{0.050000in}{0.050000in}}{\pgfqpoint{4.900000in}{4.900000in}}%
\pgfusepath{clip}%
\pgfsetbuttcap%
\pgfsetroundjoin%
\definecolor{currentfill}{rgb}{1.000000,1.000000,1.000000}%
\pgfsetfillcolor{currentfill}%
\pgfsetfillopacity{0.900000}%
\pgfsetlinewidth{0.200750pt}%
\definecolor{currentstroke}{rgb}{0.200000,0.200000,0.200000}%
\pgfsetstrokecolor{currentstroke}%
\pgfsetdash{}{0pt}%
\pgfpathmoveto{\pgfqpoint{2.143467in}{0.785442in}}%
\pgfpathlineto{\pgfqpoint{2.191686in}{0.801605in}}%
\pgfpathlineto{\pgfqpoint{2.239794in}{0.817731in}}%
\pgfpathlineto{\pgfqpoint{2.272702in}{0.794690in}}%
\pgfpathlineto{\pgfqpoint{2.305718in}{0.771573in}}%
\pgfpathlineto{\pgfqpoint{2.257527in}{0.755341in}}%
\pgfpathlineto{\pgfqpoint{2.209225in}{0.739073in}}%
\pgfpathlineto{\pgfqpoint{2.176292in}{0.762296in}}%
\pgfpathlineto{\pgfqpoint{2.143467in}{0.785442in}}%
\pgfpathclose%
\pgfusepath{stroke,fill}%
\end{pgfscope}%
\begin{pgfscope}%
\pgfpathrectangle{\pgfqpoint{0.050000in}{0.050000in}}{\pgfqpoint{4.900000in}{4.900000in}}%
\pgfusepath{clip}%
\pgfsetbuttcap%
\pgfsetroundjoin%
\definecolor{currentfill}{rgb}{1.000000,1.000000,1.000000}%
\pgfsetfillcolor{currentfill}%
\pgfsetfillopacity{0.900000}%
\pgfsetlinewidth{0.200750pt}%
\definecolor{currentstroke}{rgb}{0.200000,0.200000,0.200000}%
\pgfsetstrokecolor{currentstroke}%
\pgfsetdash{}{0pt}%
\pgfpathmoveto{\pgfqpoint{2.305718in}{0.771573in}}%
\pgfpathlineto{\pgfqpoint{2.353798in}{0.787767in}}%
\pgfpathlineto{\pgfqpoint{2.401769in}{0.803925in}}%
\pgfpathlineto{\pgfqpoint{2.434974in}{0.780838in}}%
\pgfpathlineto{\pgfqpoint{2.468287in}{0.757676in}}%
\pgfpathlineto{\pgfqpoint{2.420236in}{0.741413in}}%
\pgfpathlineto{\pgfqpoint{2.372074in}{0.725112in}}%
\pgfpathlineto{\pgfqpoint{2.338841in}{0.748381in}}%
\pgfpathlineto{\pgfqpoint{2.305718in}{0.771573in}}%
\pgfpathclose%
\pgfusepath{stroke,fill}%
\end{pgfscope}%
\begin{pgfscope}%
\pgfpathrectangle{\pgfqpoint{0.050000in}{0.050000in}}{\pgfqpoint{4.900000in}{4.900000in}}%
\pgfusepath{clip}%
\pgfsetbuttcap%
\pgfsetroundjoin%
\definecolor{currentfill}{rgb}{1.000000,1.000000,1.000000}%
\pgfsetfillcolor{currentfill}%
\pgfsetfillopacity{0.900000}%
\pgfsetlinewidth{0.200750pt}%
\definecolor{currentstroke}{rgb}{0.200000,0.200000,0.200000}%
\pgfsetstrokecolor{currentstroke}%
\pgfsetdash{}{0pt}%
\pgfpathmoveto{\pgfqpoint{2.468287in}{0.757676in}}%
\pgfpathlineto{\pgfqpoint{2.516229in}{0.773903in}}%
\pgfpathlineto{\pgfqpoint{2.564061in}{0.790092in}}%
\pgfpathlineto{\pgfqpoint{2.597564in}{0.766960in}}%
\pgfpathlineto{\pgfqpoint{2.549692in}{0.750718in}}%
\pgfpathlineto{\pgfqpoint{2.501711in}{0.734438in}}%
\pgfpathlineto{\pgfqpoint{2.468287in}{0.757676in}}%
\pgfpathclose%
\pgfusepath{stroke,fill}%
\end{pgfscope}%
\begin{pgfscope}%
\pgfpathrectangle{\pgfqpoint{0.050000in}{0.050000in}}{\pgfqpoint{4.900000in}{4.900000in}}%
\pgfusepath{clip}%
\pgfsetbuttcap%
\pgfsetroundjoin%
\definecolor{currentfill}{rgb}{1.000000,1.000000,1.000000}%
\pgfsetfillcolor{currentfill}%
\pgfsetfillopacity{0.900000}%
\pgfsetlinewidth{0.200750pt}%
\definecolor{currentstroke}{rgb}{0.200000,0.200000,0.200000}%
\pgfsetstrokecolor{currentstroke}%
\pgfsetdash{}{0pt}%
\pgfpathmoveto{\pgfqpoint{2.046696in}{0.753006in}}%
\pgfpathlineto{\pgfqpoint{2.095137in}{0.769243in}}%
\pgfpathlineto{\pgfqpoint{2.143467in}{0.785442in}}%
\pgfpathlineto{\pgfqpoint{2.176292in}{0.762296in}}%
\pgfpathlineto{\pgfqpoint{2.209225in}{0.739073in}}%
\pgfpathlineto{\pgfqpoint{2.160812in}{0.722766in}}%
\pgfpathlineto{\pgfqpoint{2.112288in}{0.706423in}}%
\pgfpathlineto{\pgfqpoint{2.079438in}{0.729752in}}%
\pgfpathlineto{\pgfqpoint{2.046696in}{0.753006in}}%
\pgfpathclose%
\pgfusepath{stroke,fill}%
\end{pgfscope}%
\begin{pgfscope}%
\pgfpathrectangle{\pgfqpoint{0.050000in}{0.050000in}}{\pgfqpoint{4.900000in}{4.900000in}}%
\pgfusepath{clip}%
\pgfsetbuttcap%
\pgfsetroundjoin%
\definecolor{currentfill}{rgb}{1.000000,1.000000,1.000000}%
\pgfsetfillcolor{currentfill}%
\pgfsetfillopacity{0.900000}%
\pgfsetlinewidth{0.200750pt}%
\definecolor{currentstroke}{rgb}{0.200000,0.200000,0.200000}%
\pgfsetstrokecolor{currentstroke}%
\pgfsetdash{}{0pt}%
\pgfpathmoveto{\pgfqpoint{2.209225in}{0.739073in}}%
\pgfpathlineto{\pgfqpoint{2.257527in}{0.755341in}}%
\pgfpathlineto{\pgfqpoint{2.305718in}{0.771573in}}%
\pgfpathlineto{\pgfqpoint{2.338841in}{0.748381in}}%
\pgfpathlineto{\pgfqpoint{2.372074in}{0.725112in}}%
\pgfpathlineto{\pgfqpoint{2.323801in}{0.708774in}}%
\pgfpathlineto{\pgfqpoint{2.275417in}{0.692398in}}%
\pgfpathlineto{\pgfqpoint{2.242266in}{0.715773in}}%
\pgfpathlineto{\pgfqpoint{2.209225in}{0.739073in}}%
\pgfpathclose%
\pgfusepath{stroke,fill}%
\end{pgfscope}%
\begin{pgfscope}%
\pgfpathrectangle{\pgfqpoint{0.050000in}{0.050000in}}{\pgfqpoint{4.900000in}{4.900000in}}%
\pgfusepath{clip}%
\pgfsetbuttcap%
\pgfsetroundjoin%
\definecolor{currentfill}{rgb}{1.000000,1.000000,1.000000}%
\pgfsetfillcolor{currentfill}%
\pgfsetfillopacity{0.900000}%
\pgfsetlinewidth{0.200750pt}%
\definecolor{currentstroke}{rgb}{0.200000,0.200000,0.200000}%
\pgfsetstrokecolor{currentstroke}%
\pgfsetdash{}{0pt}%
\pgfpathmoveto{\pgfqpoint{2.372074in}{0.725112in}}%
\pgfpathlineto{\pgfqpoint{2.420236in}{0.741413in}}%
\pgfpathlineto{\pgfqpoint{2.468287in}{0.757676in}}%
\pgfpathlineto{\pgfqpoint{2.501711in}{0.734438in}}%
\pgfpathlineto{\pgfqpoint{2.453619in}{0.718121in}}%
\pgfpathlineto{\pgfqpoint{2.405416in}{0.701767in}}%
\pgfpathlineto{\pgfqpoint{2.372074in}{0.725112in}}%
\pgfpathclose%
\pgfusepath{stroke,fill}%
\end{pgfscope}%
\begin{pgfscope}%
\pgfpathrectangle{\pgfqpoint{0.050000in}{0.050000in}}{\pgfqpoint{4.900000in}{4.900000in}}%
\pgfusepath{clip}%
\pgfsetbuttcap%
\pgfsetroundjoin%
\definecolor{currentfill}{rgb}{1.000000,1.000000,1.000000}%
\pgfsetfillcolor{currentfill}%
\pgfsetfillopacity{0.900000}%
\pgfsetlinewidth{0.200750pt}%
\definecolor{currentstroke}{rgb}{0.200000,0.200000,0.200000}%
\pgfsetstrokecolor{currentstroke}%
\pgfsetdash{}{0pt}%
\pgfpathmoveto{\pgfqpoint{2.112288in}{0.706423in}}%
\pgfpathlineto{\pgfqpoint{2.160812in}{0.722766in}}%
\pgfpathlineto{\pgfqpoint{2.209225in}{0.739073in}}%
\pgfpathlineto{\pgfqpoint{2.242266in}{0.715773in}}%
\pgfpathlineto{\pgfqpoint{2.275417in}{0.692398in}}%
\pgfpathlineto{\pgfqpoint{2.226921in}{0.675984in}}%
\pgfpathlineto{\pgfqpoint{2.178313in}{0.659532in}}%
\pgfpathlineto{\pgfqpoint{2.145246in}{0.683016in}}%
\pgfpathlineto{\pgfqpoint{2.112288in}{0.706423in}}%
\pgfpathclose%
\pgfusepath{stroke,fill}%
\end{pgfscope}%
\begin{pgfscope}%
\pgfpathrectangle{\pgfqpoint{0.050000in}{0.050000in}}{\pgfqpoint{4.900000in}{4.900000in}}%
\pgfusepath{clip}%
\pgfsetbuttcap%
\pgfsetroundjoin%
\definecolor{currentfill}{rgb}{1.000000,1.000000,1.000000}%
\pgfsetfillcolor{currentfill}%
\pgfsetfillopacity{0.900000}%
\pgfsetlinewidth{0.200750pt}%
\definecolor{currentstroke}{rgb}{0.200000,0.200000,0.200000}%
\pgfsetstrokecolor{currentstroke}%
\pgfsetdash{}{0pt}%
\pgfpathmoveto{\pgfqpoint{2.275417in}{0.692398in}}%
\pgfpathlineto{\pgfqpoint{2.323801in}{0.708774in}}%
\pgfpathlineto{\pgfqpoint{2.372074in}{0.725112in}}%
\pgfpathlineto{\pgfqpoint{2.405416in}{0.701767in}}%
\pgfpathlineto{\pgfqpoint{2.357102in}{0.685375in}}%
\pgfpathlineto{\pgfqpoint{2.308676in}{0.668945in}}%
\pgfpathlineto{\pgfqpoint{2.275417in}{0.692398in}}%
\pgfpathclose%
\pgfusepath{stroke,fill}%
\end{pgfscope}%
\begin{pgfscope}%
\pgfpathrectangle{\pgfqpoint{0.050000in}{0.050000in}}{\pgfqpoint{4.900000in}{4.900000in}}%
\pgfusepath{clip}%
\pgfsetbuttcap%
\pgfsetroundjoin%
\definecolor{currentfill}{rgb}{1.000000,1.000000,1.000000}%
\pgfsetfillcolor{currentfill}%
\pgfsetfillopacity{0.900000}%
\pgfsetlinewidth{0.200750pt}%
\definecolor{currentstroke}{rgb}{0.200000,0.200000,0.200000}%
\pgfsetstrokecolor{currentstroke}%
\pgfsetdash{}{0pt}%
\pgfpathmoveto{\pgfqpoint{2.178313in}{0.659532in}}%
\pgfpathlineto{\pgfqpoint{2.226921in}{0.675984in}}%
\pgfpathlineto{\pgfqpoint{2.275417in}{0.692398in}}%
\pgfpathlineto{\pgfqpoint{2.308676in}{0.668945in}}%
\pgfpathlineto{\pgfqpoint{2.260139in}{0.652477in}}%
\pgfpathlineto{\pgfqpoint{2.211489in}{0.635970in}}%
\pgfpathlineto{\pgfqpoint{2.178313in}{0.659532in}}%
\pgfpathclose%
\pgfusepath{stroke,fill}%
\end{pgfscope}%
\end{pgfpicture}%
\makeatother%
\endgroup%

	\caption{Graph of the up-and-out call option price}\label{fig:upAndOutCallBarrierSurface}
\end{figure}

\clearpage
\subsubsection{Up-and-out barrier put option}
This option is also known as a \textit{Callable Bear Contract}, or a Bear CBBC with no residual value, or a turbo warrant with no rebate. In this context, \(B\) denotes the call level for pricing Bear CBBCs (up-and-out barrier put options). The payoff is given by

\begin{equation*}
	P = {(K - S_T)}^+ \I_{\left\{ \max\limits_{0 \le t \le T} S_t < B \right\}} =
	\begin{dcases}
		K - S_T & \text{if } \max\limits_{0 \le t \le T} S_t \le B, \\
		0       & \text{if } \max\limits_{0 \le t \le T} S_t > B.
	\end{dcases}
\end{equation*}

\begin{prop}{Up-and-Out Barrier Put Option Pricing}{Up-and-Out Barrier Put Option Pricing}
	The price of the up-and-out barrier put option with maturity \(T\), strike price \(K\), and barrier level \(B\) is given by the following formulas, where \(\tau = T-t\).

	If \(B < K\), the price is given by
	\begin{align}
		P_{UO}(t) & = S_t \I_{\{M_0^t < B\}} \Bigg\{
		-\Phi \left( -d_1{\left(S_t, B\right)} \right)
		+ {\left( \frac{B}{S_t} \right)}^{1+\frac{2r}{\sigma^2}} \Phi \left( -d_1{\left(\frac{B^2}{S_t}, B\right)} \right) \Bigg\} \nonumber \\
		          & \quad + K e^{-r\tau} \I_{\{M_0^t < B\}} \Bigg\{
		\Phi \left( -d_2{\left(S_t, B\right)} \right)
		- {\left( \frac{S_t}{B} \right)}^{1-\frac{2r}{\sigma^2}} \Phi \left( -d_2{\left(\frac{B^2}{S_t}, B\right)} \right)
		\Bigg\}.
	\end{align}

	If \(B \ge K\), the price is given by
	\begin{align}
		P_{UO}(t) & = S_t \I_{\{M_0^t < B\}} \Bigg\{
		-\Phi \left( -d_1{\left(S_t, K\right)} \right)
		+ {\left( \frac{B}{S_t} \right)}^{1+\frac{2r}{\sigma^2}} \Phi \left( -d_1{\left(\frac{B^2}{S_t}, K\right)} \right) \Bigg\} \nonumber \\
		          & \quad + K e^{-r\tau} \I_{\{M_0^t < B\}} \Bigg\{
		\Phi \left( -d_2{\left(S_t, K\right)} \right)
		- {\left( \frac{S_t}{B} \right)}^{1-\frac{2r}{\sigma^2}} \Phi \left( -d_2{\left(\frac{B^2}{S_t}, K\right)} \right)
		\Bigg\}.
	\end{align}
	This corresponds to the value of a vanilla put option minus the value of the ``knock-out'' component.
\end{prop}

\begin{figure}[H]
	\centering
	%% Creator: Matplotlib, PGF backend
%%
%% To include the figure in your LaTeX document, write
%%   \input{<filename>.pgf}
%%
%% Make sure the required packages are loaded in your preamble
%%   \usepackage{pgf}
%%
%% Also ensure that all the required font packages are loaded; for instance,
%% the lmodern package is sometimes necessary when using math font.
%%   \usepackage{lmodern}
%%
%% Figures using additional raster images can only be included by \input if
%% they are in the same directory as the main LaTeX file. For loading figures
%% from other directories you can use the `import` package
%%   \usepackage{import}
%%
%% and then include the figures with
%%   \import{<path to file>}{<filename>.pgf}
%%
%% Matplotlib used the following preamble
%%   \def\mathdefault#1{#1}
%%   \everymath=\expandafter{\the\everymath\displaystyle}
%%   \IfFileExists{scrextend.sty}{
%%     \usepackage[fontsize=10.000000pt]{scrextend}
%%   }{
%%     \renewcommand{\normalsize}{\fontsize{10.000000}{12.000000}\selectfont}
%%     \normalsize
%%   }
%%   
%%   \makeatletter\@ifpackageloaded{underscore}{}{\usepackage[strings]{underscore}}\makeatother
%%
\begingroup%
\makeatletter%
\begin{pgfpicture}%
\pgfpathrectangle{\pgfpointorigin}{\pgfqpoint{5.240711in}{5.000000in}}%
\pgfusepath{use as bounding box, clip}%
\begin{pgfscope}%
\pgfsetbuttcap%
\pgfsetmiterjoin%
\definecolor{currentfill}{rgb}{1.000000,1.000000,1.000000}%
\pgfsetfillcolor{currentfill}%
\pgfsetlinewidth{0.000000pt}%
\definecolor{currentstroke}{rgb}{1.000000,1.000000,1.000000}%
\pgfsetstrokecolor{currentstroke}%
\pgfsetdash{}{0pt}%
\pgfpathmoveto{\pgfqpoint{0.000000in}{0.000000in}}%
\pgfpathlineto{\pgfqpoint{5.240711in}{0.000000in}}%
\pgfpathlineto{\pgfqpoint{5.240711in}{5.000000in}}%
\pgfpathlineto{\pgfqpoint{0.000000in}{5.000000in}}%
\pgfpathlineto{\pgfqpoint{0.000000in}{0.000000in}}%
\pgfpathclose%
\pgfusepath{fill}%
\end{pgfscope}%
\begin{pgfscope}%
\pgfsetbuttcap%
\pgfsetmiterjoin%
\definecolor{currentfill}{rgb}{1.000000,1.000000,1.000000}%
\pgfsetfillcolor{currentfill}%
\pgfsetlinewidth{0.000000pt}%
\definecolor{currentstroke}{rgb}{0.000000,0.000000,0.000000}%
\pgfsetstrokecolor{currentstroke}%
\pgfsetstrokeopacity{0.000000}%
\pgfsetdash{}{0pt}%
\pgfpathmoveto{\pgfqpoint{0.050000in}{0.050000in}}%
\pgfpathlineto{\pgfqpoint{4.950000in}{0.050000in}}%
\pgfpathlineto{\pgfqpoint{4.950000in}{4.950000in}}%
\pgfpathlineto{\pgfqpoint{0.050000in}{4.950000in}}%
\pgfpathlineto{\pgfqpoint{0.050000in}{0.050000in}}%
\pgfpathclose%
\pgfusepath{fill}%
\end{pgfscope}%
\begin{pgfscope}%
\pgfsetbuttcap%
\pgfsetmiterjoin%
\pgfsetlinewidth{1.003750pt}%
\definecolor{currentstroke}{rgb}{0.950000,0.950000,0.950000}%
\pgfsetstrokecolor{currentstroke}%
\pgfsetstrokeopacity{0.500000}%
\pgfsetdash{}{0pt}%
\pgfpathmoveto{\pgfqpoint{0.438795in}{1.686447in}}%
\pgfpathlineto{\pgfqpoint{2.920118in}{2.434876in}}%
\pgfpathlineto{\pgfqpoint{2.932761in}{4.459308in}}%
\pgfpathlineto{\pgfqpoint{0.358066in}{3.808469in}}%
\pgfusepath{stroke}%
\end{pgfscope}%
\begin{pgfscope}%
\pgfsetbuttcap%
\pgfsetmiterjoin%
\pgfsetlinewidth{1.003750pt}%
\definecolor{currentstroke}{rgb}{0.900000,0.900000,0.900000}%
\pgfsetstrokecolor{currentstroke}%
\pgfsetstrokeopacity{0.500000}%
\pgfsetdash{}{0pt}%
\pgfpathmoveto{\pgfqpoint{2.920118in}{2.434876in}}%
\pgfpathlineto{\pgfqpoint{4.749691in}{1.337846in}}%
\pgfpathlineto{\pgfqpoint{4.834815in}{3.504368in}}%
\pgfpathlineto{\pgfqpoint{2.932761in}{4.459308in}}%
\pgfusepath{stroke}%
\end{pgfscope}%
\begin{pgfscope}%
\pgfsetbuttcap%
\pgfsetmiterjoin%
\pgfsetlinewidth{1.003750pt}%
\definecolor{currentstroke}{rgb}{0.925000,0.925000,0.925000}%
\pgfsetstrokecolor{currentstroke}%
\pgfsetstrokeopacity{0.500000}%
\pgfsetdash{}{0pt}%
\pgfpathmoveto{\pgfqpoint{0.438795in}{1.686447in}}%
\pgfpathlineto{\pgfqpoint{2.155964in}{0.447388in}}%
\pgfpathlineto{\pgfqpoint{4.749691in}{1.337846in}}%
\pgfpathlineto{\pgfqpoint{2.920118in}{2.434876in}}%
\pgfusepath{stroke}%
\end{pgfscope}%
\begin{pgfscope}%
\pgfsetbuttcap%
\pgfsetroundjoin%
\pgfsetlinewidth{0.501875pt}%
\definecolor{currentstroke}{rgb}{0.800000,0.800000,0.800000}%
\pgfsetstrokecolor{currentstroke}%
\pgfsetdash{}{0pt}%
\pgfpathmoveto{\pgfqpoint{0.459298in}{1.671653in}}%
\pgfpathlineto{\pgfqpoint{2.942077in}{2.421709in}}%
\pgfpathlineto{\pgfqpoint{2.955519in}{4.447882in}}%
\pgfusepath{stroke}%
\end{pgfscope}%
\begin{pgfscope}%
\pgfsetbuttcap%
\pgfsetroundjoin%
\pgfsetlinewidth{0.501875pt}%
\definecolor{currentstroke}{rgb}{0.800000,0.800000,0.800000}%
\pgfsetstrokecolor{currentstroke}%
\pgfsetdash{}{0pt}%
\pgfpathmoveto{\pgfqpoint{0.783380in}{1.437805in}}%
\pgfpathlineto{\pgfqpoint{3.288812in}{2.213803in}}%
\pgfpathlineto{\pgfqpoint{3.315100in}{4.267352in}}%
\pgfusepath{stroke}%
\end{pgfscope}%
\begin{pgfscope}%
\pgfsetbuttcap%
\pgfsetroundjoin%
\pgfsetlinewidth{0.501875pt}%
\definecolor{currentstroke}{rgb}{0.800000,0.800000,0.800000}%
\pgfsetstrokecolor{currentstroke}%
\pgfsetdash{}{0pt}%
\pgfpathmoveto{\pgfqpoint{1.119111in}{1.195551in}}%
\pgfpathlineto{\pgfqpoint{3.647281in}{1.998862in}}%
\pgfpathlineto{\pgfqpoint{3.687299in}{4.080487in}}%
\pgfusepath{stroke}%
\end{pgfscope}%
\begin{pgfscope}%
\pgfsetbuttcap%
\pgfsetroundjoin%
\pgfsetlinewidth{0.501875pt}%
\definecolor{currentstroke}{rgb}{0.800000,0.800000,0.800000}%
\pgfsetstrokecolor{currentstroke}%
\pgfsetdash{}{0pt}%
\pgfpathmoveto{\pgfqpoint{1.467131in}{0.944430in}}%
\pgfpathlineto{\pgfqpoint{4.018089in}{1.776522in}}%
\pgfpathlineto{\pgfqpoint{4.072794in}{3.886946in}}%
\pgfusepath{stroke}%
\end{pgfscope}%
\begin{pgfscope}%
\pgfsetbuttcap%
\pgfsetroundjoin%
\pgfsetlinewidth{0.501875pt}%
\definecolor{currentstroke}{rgb}{0.800000,0.800000,0.800000}%
\pgfsetstrokecolor{currentstroke}%
\pgfsetdash{}{0pt}%
\pgfpathmoveto{\pgfqpoint{1.828126in}{0.683946in}}%
\pgfpathlineto{\pgfqpoint{4.401884in}{1.546394in}}%
\pgfpathlineto{\pgfqpoint{4.472308in}{3.686368in}}%
\pgfusepath{stroke}%
\end{pgfscope}%
\begin{pgfscope}%
\pgfsetbuttcap%
\pgfsetroundjoin%
\pgfsetlinewidth{0.501875pt}%
\definecolor{currentstroke}{rgb}{0.800000,0.800000,0.800000}%
\pgfsetstrokecolor{currentstroke}%
\pgfsetdash{}{0pt}%
\pgfpathmoveto{\pgfqpoint{0.539260in}{1.613955in}}%
\pgfpathlineto{\pgfqpoint{3.027693in}{2.370373in}}%
\pgfpathlineto{\pgfqpoint{3.044267in}{4.403325in}}%
\pgfusepath{stroke}%
\end{pgfscope}%
\begin{pgfscope}%
\pgfsetbuttcap%
\pgfsetroundjoin%
\pgfsetlinewidth{0.501875pt}%
\definecolor{currentstroke}{rgb}{0.800000,0.800000,0.800000}%
\pgfsetstrokecolor{currentstroke}%
\pgfsetdash{}{0pt}%
\pgfpathmoveto{\pgfqpoint{0.619921in}{1.555752in}}%
\pgfpathlineto{\pgfqpoint{3.114014in}{2.318614in}}%
\pgfpathlineto{\pgfqpoint{3.133773in}{4.358388in}}%
\pgfusepath{stroke}%
\end{pgfscope}%
\begin{pgfscope}%
\pgfsetbuttcap%
\pgfsetroundjoin%
\pgfsetlinewidth{0.501875pt}%
\definecolor{currentstroke}{rgb}{0.800000,0.800000,0.800000}%
\pgfsetstrokecolor{currentstroke}%
\pgfsetdash{}{0pt}%
\pgfpathmoveto{\pgfqpoint{0.701291in}{1.497037in}}%
\pgfpathlineto{\pgfqpoint{3.201051in}{2.266426in}}%
\pgfpathlineto{\pgfqpoint{3.224047in}{4.313065in}}%
\pgfusepath{stroke}%
\end{pgfscope}%
\begin{pgfscope}%
\pgfsetbuttcap%
\pgfsetroundjoin%
\pgfsetlinewidth{0.501875pt}%
\definecolor{currentstroke}{rgb}{0.800000,0.800000,0.800000}%
\pgfsetstrokecolor{currentstroke}%
\pgfsetdash{}{0pt}%
\pgfpathmoveto{\pgfqpoint{0.866197in}{1.378047in}}%
\pgfpathlineto{\pgfqpoint{3.377306in}{2.160741in}}%
\pgfpathlineto{\pgfqpoint{3.406941in}{4.221242in}}%
\pgfusepath{stroke}%
\end{pgfscope}%
\begin{pgfscope}%
\pgfsetbuttcap%
\pgfsetroundjoin%
\pgfsetlinewidth{0.501875pt}%
\definecolor{currentstroke}{rgb}{0.800000,0.800000,0.800000}%
\pgfsetstrokecolor{currentstroke}%
\pgfsetdash{}{0pt}%
\pgfpathmoveto{\pgfqpoint{0.949751in}{1.317756in}}%
\pgfpathlineto{\pgfqpoint{3.466543in}{2.107234in}}%
\pgfpathlineto{\pgfqpoint{3.499581in}{4.174732in}}%
\pgfusepath{stroke}%
\end{pgfscope}%
\begin{pgfscope}%
\pgfsetbuttcap%
\pgfsetroundjoin%
\pgfsetlinewidth{0.501875pt}%
\definecolor{currentstroke}{rgb}{0.800000,0.800000,0.800000}%
\pgfsetstrokecolor{currentstroke}%
\pgfsetdash{}{0pt}%
\pgfpathmoveto{\pgfqpoint{1.034052in}{1.256927in}}%
\pgfpathlineto{\pgfqpoint{3.556531in}{2.053276in}}%
\pgfpathlineto{\pgfqpoint{3.593030in}{4.127815in}}%
\pgfusepath{stroke}%
\end{pgfscope}%
\begin{pgfscope}%
\pgfsetbuttcap%
\pgfsetroundjoin%
\pgfsetlinewidth{0.501875pt}%
\definecolor{currentstroke}{rgb}{0.800000,0.800000,0.800000}%
\pgfsetstrokecolor{currentstroke}%
\pgfsetdash{}{0pt}%
\pgfpathmoveto{\pgfqpoint{1.204938in}{1.133621in}}%
\pgfpathlineto{\pgfqpoint{3.738801in}{1.943985in}}%
\pgfpathlineto{\pgfqpoint{3.782399in}{4.032741in}}%
\pgfusepath{stroke}%
\end{pgfscope}%
\begin{pgfscope}%
\pgfsetbuttcap%
\pgfsetroundjoin%
\pgfsetlinewidth{0.501875pt}%
\definecolor{currentstroke}{rgb}{0.800000,0.800000,0.800000}%
\pgfsetstrokecolor{currentstroke}%
\pgfsetdash{}{0pt}%
\pgfpathmoveto{\pgfqpoint{1.291543in}{1.071129in}}%
\pgfpathlineto{\pgfqpoint{3.831103in}{1.888641in}}%
\pgfpathlineto{\pgfqpoint{3.878340in}{3.984573in}}%
\pgfusepath{stroke}%
\end{pgfscope}%
\begin{pgfscope}%
\pgfsetbuttcap%
\pgfsetroundjoin%
\pgfsetlinewidth{0.501875pt}%
\definecolor{currentstroke}{rgb}{0.800000,0.800000,0.800000}%
\pgfsetstrokecolor{currentstroke}%
\pgfsetdash{}{0pt}%
\pgfpathmoveto{\pgfqpoint{1.378937in}{1.008068in}}%
\pgfpathlineto{\pgfqpoint{3.924195in}{1.832821in}}%
\pgfpathlineto{\pgfqpoint{3.975135in}{3.935977in}}%
\pgfusepath{stroke}%
\end{pgfscope}%
\begin{pgfscope}%
\pgfsetbuttcap%
\pgfsetroundjoin%
\pgfsetlinewidth{0.501875pt}%
\definecolor{currentstroke}{rgb}{0.800000,0.800000,0.800000}%
\pgfsetstrokecolor{currentstroke}%
\pgfsetdash{}{0pt}%
\pgfpathmoveto{\pgfqpoint{1.556135in}{0.880207in}}%
\pgfpathlineto{\pgfqpoint{4.112794in}{1.719736in}}%
\pgfpathlineto{\pgfqpoint{4.171328in}{3.837476in}}%
\pgfusepath{stroke}%
\end{pgfscope}%
\begin{pgfscope}%
\pgfsetbuttcap%
\pgfsetroundjoin%
\pgfsetlinewidth{0.501875pt}%
\definecolor{currentstroke}{rgb}{0.800000,0.800000,0.800000}%
\pgfsetstrokecolor{currentstroke}%
\pgfsetdash{}{0pt}%
\pgfpathmoveto{\pgfqpoint{1.645962in}{0.815391in}}%
\pgfpathlineto{\pgfqpoint{4.208321in}{1.662457in}}%
\pgfpathlineto{\pgfqpoint{4.270751in}{3.787561in}}%
\pgfusepath{stroke}%
\end{pgfscope}%
\begin{pgfscope}%
\pgfsetbuttcap%
\pgfsetroundjoin%
\pgfsetlinewidth{0.501875pt}%
\definecolor{currentstroke}{rgb}{0.800000,0.800000,0.800000}%
\pgfsetstrokecolor{currentstroke}%
\pgfsetdash{}{0pt}%
\pgfpathmoveto{\pgfqpoint{1.736622in}{0.749974in}}%
\pgfpathlineto{\pgfqpoint{4.304681in}{1.604679in}}%
\pgfpathlineto{\pgfqpoint{4.371073in}{3.737193in}}%
\pgfusepath{stroke}%
\end{pgfscope}%
\begin{pgfscope}%
\pgfsetbuttcap%
\pgfsetroundjoin%
\pgfsetlinewidth{0.501875pt}%
\definecolor{currentstroke}{rgb}{0.800000,0.800000,0.800000}%
\pgfsetstrokecolor{currentstroke}%
\pgfsetdash{}{0pt}%
\pgfpathmoveto{\pgfqpoint{1.920488in}{0.617301in}}%
\pgfpathlineto{\pgfqpoint{4.499943in}{1.487597in}}%
\pgfpathlineto{\pgfqpoint{4.574467in}{3.635078in}}%
\pgfusepath{stroke}%
\end{pgfscope}%
\begin{pgfscope}%
\pgfsetbuttcap%
\pgfsetroundjoin%
\pgfsetlinewidth{0.501875pt}%
\definecolor{currentstroke}{rgb}{0.800000,0.800000,0.800000}%
\pgfsetstrokecolor{currentstroke}%
\pgfsetdash{}{0pt}%
\pgfpathmoveto{\pgfqpoint{2.013719in}{0.550028in}}%
\pgfpathlineto{\pgfqpoint{4.598868in}{1.428281in}}%
\pgfpathlineto{\pgfqpoint{4.677564in}{3.583317in}}%
\pgfusepath{stroke}%
\end{pgfscope}%
\begin{pgfscope}%
\pgfsetbuttcap%
\pgfsetroundjoin%
\pgfsetlinewidth{0.501875pt}%
\definecolor{currentstroke}{rgb}{0.800000,0.800000,0.800000}%
\pgfsetstrokecolor{currentstroke}%
\pgfsetdash{}{0pt}%
\pgfpathmoveto{\pgfqpoint{2.107831in}{0.482120in}}%
\pgfpathlineto{\pgfqpoint{4.698670in}{1.368439in}}%
\pgfpathlineto{\pgfqpoint{4.781611in}{3.531080in}}%
\pgfusepath{stroke}%
\end{pgfscope}%
\begin{pgfscope}%
\pgfsetbuttcap%
\pgfsetroundjoin%
\pgfsetlinewidth{0.501875pt}%
\definecolor{currentstroke}{rgb}{0.800000,0.800000,0.800000}%
\pgfsetstrokecolor{currentstroke}%
\pgfsetdash{}{0pt}%
\pgfpathmoveto{\pgfqpoint{0.361374in}{3.809305in}}%
\pgfpathlineto{\pgfqpoint{0.441976in}{1.687406in}}%
\pgfpathlineto{\pgfqpoint{2.159307in}{0.448536in}}%
\pgfusepath{stroke}%
\end{pgfscope}%
\begin{pgfscope}%
\pgfsetbuttcap%
\pgfsetroundjoin%
\pgfsetlinewidth{0.501875pt}%
\definecolor{currentstroke}{rgb}{0.800000,0.800000,0.800000}%
\pgfsetstrokecolor{currentstroke}%
\pgfsetdash{}{0pt}%
\pgfpathmoveto{\pgfqpoint{0.867596in}{3.937269in}}%
\pgfpathlineto{\pgfqpoint{0.929003in}{1.834306in}}%
\pgfpathlineto{\pgfqpoint{2.670554in}{0.624053in}}%
\pgfusepath{stroke}%
\end{pgfscope}%
\begin{pgfscope}%
\pgfsetbuttcap%
\pgfsetroundjoin%
\pgfsetlinewidth{0.501875pt}%
\definecolor{currentstroke}{rgb}{0.800000,0.800000,0.800000}%
\pgfsetstrokecolor{currentstroke}%
\pgfsetdash{}{0pt}%
\pgfpathmoveto{\pgfqpoint{1.362288in}{4.062319in}}%
\pgfpathlineto{\pgfqpoint{1.405333in}{1.977979in}}%
\pgfpathlineto{\pgfqpoint{3.169547in}{0.795364in}}%
\pgfusepath{stroke}%
\end{pgfscope}%
\begin{pgfscope}%
\pgfsetbuttcap%
\pgfsetroundjoin%
\pgfsetlinewidth{0.501875pt}%
\definecolor{currentstroke}{rgb}{0.800000,0.800000,0.800000}%
\pgfsetstrokecolor{currentstroke}%
\pgfsetdash{}{0pt}%
\pgfpathmoveto{\pgfqpoint{1.845838in}{4.184552in}}%
\pgfpathlineto{\pgfqpoint{1.871315in}{2.118531in}}%
\pgfpathlineto{\pgfqpoint{3.656720in}{0.962616in}}%
\pgfusepath{stroke}%
\end{pgfscope}%
\begin{pgfscope}%
\pgfsetbuttcap%
\pgfsetroundjoin%
\pgfsetlinewidth{0.501875pt}%
\definecolor{currentstroke}{rgb}{0.800000,0.800000,0.800000}%
\pgfsetstrokecolor{currentstroke}%
\pgfsetdash{}{0pt}%
\pgfpathmoveto{\pgfqpoint{2.318620in}{4.304063in}}%
\pgfpathlineto{\pgfqpoint{2.327282in}{2.256062in}}%
\pgfpathlineto{\pgfqpoint{4.132488in}{1.125953in}}%
\pgfusepath{stroke}%
\end{pgfscope}%
\begin{pgfscope}%
\pgfsetbuttcap%
\pgfsetroundjoin%
\pgfsetlinewidth{0.501875pt}%
\definecolor{currentstroke}{rgb}{0.800000,0.800000,0.800000}%
\pgfsetstrokecolor{currentstroke}%
\pgfsetdash{}{0pt}%
\pgfpathmoveto{\pgfqpoint{2.780988in}{4.420942in}}%
\pgfpathlineto{\pgfqpoint{2.773554in}{2.390668in}}%
\pgfpathlineto{\pgfqpoint{4.597247in}{1.285510in}}%
\pgfusepath{stroke}%
\end{pgfscope}%
\begin{pgfscope}%
\pgfsetbuttcap%
\pgfsetroundjoin%
\pgfsetlinewidth{0.501875pt}%
\definecolor{currentstroke}{rgb}{0.800000,0.800000,0.800000}%
\pgfsetstrokecolor{currentstroke}%
\pgfsetdash{}{0pt}%
\pgfpathmoveto{\pgfqpoint{0.489032in}{3.841575in}}%
\pgfpathlineto{\pgfqpoint{0.564755in}{1.724440in}}%
\pgfpathlineto{\pgfqpoint{2.288292in}{0.492818in}}%
\pgfusepath{stroke}%
\end{pgfscope}%
\begin{pgfscope}%
\pgfsetbuttcap%
\pgfsetroundjoin%
\pgfsetlinewidth{0.501875pt}%
\definecolor{currentstroke}{rgb}{0.800000,0.800000,0.800000}%
\pgfsetstrokecolor{currentstroke}%
\pgfsetdash{}{0pt}%
\pgfpathmoveto{\pgfqpoint{0.615951in}{3.873658in}}%
\pgfpathlineto{\pgfqpoint{0.686849in}{1.761266in}}%
\pgfpathlineto{\pgfqpoint{2.416491in}{0.536830in}}%
\pgfusepath{stroke}%
\end{pgfscope}%
\begin{pgfscope}%
\pgfsetbuttcap%
\pgfsetroundjoin%
\pgfsetlinewidth{0.501875pt}%
\definecolor{currentstroke}{rgb}{0.800000,0.800000,0.800000}%
\pgfsetstrokecolor{currentstroke}%
\pgfsetdash{}{0pt}%
\pgfpathmoveto{\pgfqpoint{0.742137in}{3.905555in}}%
\pgfpathlineto{\pgfqpoint{0.808263in}{1.797888in}}%
\pgfpathlineto{\pgfqpoint{2.543909in}{0.580574in}}%
\pgfusepath{stroke}%
\end{pgfscope}%
\begin{pgfscope}%
\pgfsetbuttcap%
\pgfsetroundjoin%
\pgfsetlinewidth{0.501875pt}%
\definecolor{currentstroke}{rgb}{0.800000,0.800000,0.800000}%
\pgfsetstrokecolor{currentstroke}%
\pgfsetdash{}{0pt}%
\pgfpathmoveto{\pgfqpoint{0.992334in}{3.968801in}}%
\pgfpathlineto{\pgfqpoint{1.049075in}{1.870522in}}%
\pgfpathlineto{\pgfqpoint{2.796434in}{0.667269in}}%
\pgfusepath{stroke}%
\end{pgfscope}%
\begin{pgfscope}%
\pgfsetbuttcap%
\pgfsetroundjoin%
\pgfsetlinewidth{0.501875pt}%
\definecolor{currentstroke}{rgb}{0.800000,0.800000,0.800000}%
\pgfsetstrokecolor{currentstroke}%
\pgfsetdash{}{0pt}%
\pgfpathmoveto{\pgfqpoint{1.116358in}{4.000152in}}%
\pgfpathlineto{\pgfqpoint{1.168483in}{1.906539in}}%
\pgfpathlineto{\pgfqpoint{2.921555in}{0.710225in}}%
\pgfusepath{stroke}%
\end{pgfscope}%
\begin{pgfscope}%
\pgfsetbuttcap%
\pgfsetroundjoin%
\pgfsetlinewidth{0.501875pt}%
\definecolor{currentstroke}{rgb}{0.800000,0.800000,0.800000}%
\pgfsetstrokecolor{currentstroke}%
\pgfsetdash{}{0pt}%
\pgfpathmoveto{\pgfqpoint{1.239674in}{4.031324in}}%
\pgfpathlineto{\pgfqpoint{1.287234in}{1.942357in}}%
\pgfpathlineto{\pgfqpoint{3.045924in}{0.752922in}}%
\pgfusepath{stroke}%
\end{pgfscope}%
\begin{pgfscope}%
\pgfsetbuttcap%
\pgfsetroundjoin%
\pgfsetlinewidth{0.501875pt}%
\definecolor{currentstroke}{rgb}{0.800000,0.800000,0.800000}%
\pgfsetstrokecolor{currentstroke}%
\pgfsetdash{}{0pt}%
\pgfpathmoveto{\pgfqpoint{1.484205in}{4.093137in}}%
\pgfpathlineto{\pgfqpoint{1.522785in}{2.013405in}}%
\pgfpathlineto{\pgfqpoint{3.292432in}{0.837551in}}%
\pgfusepath{stroke}%
\end{pgfscope}%
\begin{pgfscope}%
\pgfsetbuttcap%
\pgfsetroundjoin%
\pgfsetlinewidth{0.501875pt}%
\definecolor{currentstroke}{rgb}{0.800000,0.800000,0.800000}%
\pgfsetstrokecolor{currentstroke}%
\pgfsetdash{}{0pt}%
\pgfpathmoveto{\pgfqpoint{1.605432in}{4.123782in}}%
\pgfpathlineto{\pgfqpoint{1.639596in}{2.048638in}}%
\pgfpathlineto{\pgfqpoint{3.414584in}{0.879488in}}%
\pgfusepath{stroke}%
\end{pgfscope}%
\begin{pgfscope}%
\pgfsetbuttcap%
\pgfsetroundjoin%
\pgfsetlinewidth{0.501875pt}%
\definecolor{currentstroke}{rgb}{0.800000,0.800000,0.800000}%
\pgfsetstrokecolor{currentstroke}%
\pgfsetdash{}{0pt}%
\pgfpathmoveto{\pgfqpoint{1.725974in}{4.154253in}}%
\pgfpathlineto{\pgfqpoint{1.755771in}{2.083680in}}%
\pgfpathlineto{\pgfqpoint{3.536012in}{0.921175in}}%
\pgfusepath{stroke}%
\end{pgfscope}%
\begin{pgfscope}%
\pgfsetbuttcap%
\pgfsetroundjoin%
\pgfsetlinewidth{0.501875pt}%
\definecolor{currentstroke}{rgb}{0.800000,0.800000,0.800000}%
\pgfsetstrokecolor{currentstroke}%
\pgfsetdash{}{0pt}%
\pgfpathmoveto{\pgfqpoint{1.965029in}{4.214682in}}%
\pgfpathlineto{\pgfqpoint{1.986233in}{2.153193in}}%
\pgfpathlineto{\pgfqpoint{3.776715in}{1.003812in}}%
\pgfusepath{stroke}%
\end{pgfscope}%
\begin{pgfscope}%
\pgfsetbuttcap%
\pgfsetroundjoin%
\pgfsetlinewidth{0.501875pt}%
\definecolor{currentstroke}{rgb}{0.800000,0.800000,0.800000}%
\pgfsetstrokecolor{currentstroke}%
\pgfsetdash{}{0pt}%
\pgfpathmoveto{\pgfqpoint{2.083552in}{4.244642in}}%
\pgfpathlineto{\pgfqpoint{2.100530in}{2.187668in}}%
\pgfpathlineto{\pgfqpoint{3.896004in}{1.044765in}}%
\pgfusepath{stroke}%
\end{pgfscope}%
\begin{pgfscope}%
\pgfsetbuttcap%
\pgfsetroundjoin%
\pgfsetlinewidth{0.501875pt}%
\definecolor{currentstroke}{rgb}{0.800000,0.800000,0.800000}%
\pgfsetstrokecolor{currentstroke}%
\pgfsetdash{}{0pt}%
\pgfpathmoveto{\pgfqpoint{2.201414in}{4.274436in}}%
\pgfpathlineto{\pgfqpoint{2.214211in}{2.221957in}}%
\pgfpathlineto{\pgfqpoint{4.014593in}{1.085478in}}%
\pgfusepath{stroke}%
\end{pgfscope}%
\begin{pgfscope}%
\pgfsetbuttcap%
\pgfsetroundjoin%
\pgfsetlinewidth{0.501875pt}%
\definecolor{currentstroke}{rgb}{0.800000,0.800000,0.800000}%
\pgfsetstrokecolor{currentstroke}%
\pgfsetdash{}{0pt}%
\pgfpathmoveto{\pgfqpoint{2.435174in}{4.333526in}}%
\pgfpathlineto{\pgfqpoint{2.439747in}{2.289984in}}%
\pgfpathlineto{\pgfqpoint{4.249694in}{1.166191in}}%
\pgfusepath{stroke}%
\end{pgfscope}%
\begin{pgfscope}%
\pgfsetbuttcap%
\pgfsetroundjoin%
\pgfsetlinewidth{0.501875pt}%
\definecolor{currentstroke}{rgb}{0.800000,0.800000,0.800000}%
\pgfsetstrokecolor{currentstroke}%
\pgfsetdash{}{0pt}%
\pgfpathmoveto{\pgfqpoint{2.551084in}{4.362826in}}%
\pgfpathlineto{\pgfqpoint{2.551610in}{2.323725in}}%
\pgfpathlineto{\pgfqpoint{4.366219in}{1.206196in}}%
\pgfusepath{stroke}%
\end{pgfscope}%
\begin{pgfscope}%
\pgfsetbuttcap%
\pgfsetroundjoin%
\pgfsetlinewidth{0.501875pt}%
\definecolor{currentstroke}{rgb}{0.800000,0.800000,0.800000}%
\pgfsetstrokecolor{currentstroke}%
\pgfsetdash{}{0pt}%
\pgfpathmoveto{\pgfqpoint{2.666353in}{4.391964in}}%
\pgfpathlineto{\pgfqpoint{2.662878in}{2.357286in}}%
\pgfpathlineto{\pgfqpoint{4.482068in}{1.245968in}}%
\pgfusepath{stroke}%
\end{pgfscope}%
\begin{pgfscope}%
\pgfsetbuttcap%
\pgfsetroundjoin%
\pgfsetlinewidth{0.501875pt}%
\definecolor{currentstroke}{rgb}{0.800000,0.800000,0.800000}%
\pgfsetstrokecolor{currentstroke}%
\pgfsetdash{}{0pt}%
\pgfpathmoveto{\pgfqpoint{2.894993in}{4.449761in}}%
\pgfpathlineto{\pgfqpoint{2.883643in}{2.423874in}}%
\pgfpathlineto{\pgfqpoint{4.711761in}{1.324824in}}%
\pgfusepath{stroke}%
\end{pgfscope}%
\begin{pgfscope}%
\pgfsetbuttcap%
\pgfsetroundjoin%
\pgfsetlinewidth{0.501875pt}%
\definecolor{currentstroke}{rgb}{0.800000,0.800000,0.800000}%
\pgfsetstrokecolor{currentstroke}%
\pgfsetdash{}{0pt}%
\pgfpathmoveto{\pgfqpoint{4.751331in}{1.379582in}}%
\pgfpathlineto{\pgfqpoint{2.920363in}{2.473995in}}%
\pgfpathlineto{\pgfqpoint{0.437238in}{1.727366in}}%
\pgfusepath{stroke}%
\end{pgfscope}%
\begin{pgfscope}%
\pgfsetbuttcap%
\pgfsetroundjoin%
\pgfsetlinewidth{0.501875pt}%
\definecolor{currentstroke}{rgb}{0.800000,0.800000,0.800000}%
\pgfsetstrokecolor{currentstroke}%
\pgfsetdash{}{0pt}%
\pgfpathmoveto{\pgfqpoint{4.767217in}{1.783905in}}%
\pgfpathlineto{\pgfqpoint{2.922728in}{2.852722in}}%
\pgfpathlineto{\pgfqpoint{0.422161in}{2.123691in}}%
\pgfusepath{stroke}%
\end{pgfscope}%
\begin{pgfscope}%
\pgfsetbuttcap%
\pgfsetroundjoin%
\pgfsetlinewidth{0.501875pt}%
\definecolor{currentstroke}{rgb}{0.800000,0.800000,0.800000}%
\pgfsetstrokecolor{currentstroke}%
\pgfsetdash{}{0pt}%
\pgfpathmoveto{\pgfqpoint{4.783336in}{2.194150in}}%
\pgfpathlineto{\pgfqpoint{2.925125in}{3.236542in}}%
\pgfpathlineto{\pgfqpoint{0.406868in}{2.525670in}}%
\pgfusepath{stroke}%
\end{pgfscope}%
\begin{pgfscope}%
\pgfsetbuttcap%
\pgfsetroundjoin%
\pgfsetlinewidth{0.501875pt}%
\definecolor{currentstroke}{rgb}{0.800000,0.800000,0.800000}%
\pgfsetstrokecolor{currentstroke}%
\pgfsetdash{}{0pt}%
\pgfpathmoveto{\pgfqpoint{4.799693in}{2.610448in}}%
\pgfpathlineto{\pgfqpoint{2.927554in}{3.625558in}}%
\pgfpathlineto{\pgfqpoint{0.391356in}{2.933425in}}%
\pgfusepath{stroke}%
\end{pgfscope}%
\begin{pgfscope}%
\pgfsetbuttcap%
\pgfsetroundjoin%
\pgfsetlinewidth{0.501875pt}%
\definecolor{currentstroke}{rgb}{0.800000,0.800000,0.800000}%
\pgfsetstrokecolor{currentstroke}%
\pgfsetdash{}{0pt}%
\pgfpathmoveto{\pgfqpoint{4.816292in}{3.032935in}}%
\pgfpathlineto{\pgfqpoint{2.930017in}{4.019876in}}%
\pgfpathlineto{\pgfqpoint{0.375619in}{3.347080in}}%
\pgfusepath{stroke}%
\end{pgfscope}%
\begin{pgfscope}%
\pgfsetbuttcap%
\pgfsetroundjoin%
\pgfsetlinewidth{0.501875pt}%
\definecolor{currentstroke}{rgb}{0.800000,0.800000,0.800000}%
\pgfsetstrokecolor{currentstroke}%
\pgfsetdash{}{0pt}%
\pgfpathmoveto{\pgfqpoint{4.833141in}{3.461748in}}%
\pgfpathlineto{\pgfqpoint{2.932513in}{4.419606in}}%
\pgfpathlineto{\pgfqpoint{0.359652in}{3.766765in}}%
\pgfusepath{stroke}%
\end{pgfscope}%
\begin{pgfscope}%
\pgfsetrectcap%
\pgfsetroundjoin%
\pgfsetlinewidth{0.501875pt}%
\definecolor{currentstroke}{rgb}{0.000000,0.000000,0.000000}%
\pgfsetstrokecolor{currentstroke}%
\pgfsetdash{}{0pt}%
\pgfpathmoveto{\pgfqpoint{0.438795in}{1.686447in}}%
\pgfpathlineto{\pgfqpoint{2.155964in}{0.447388in}}%
\pgfusepath{stroke}%
\end{pgfscope}%
\begin{pgfscope}%
\pgfsetrectcap%
\pgfsetroundjoin%
\pgfsetlinewidth{0.501875pt}%
\definecolor{currentstroke}{rgb}{0.000000,0.000000,0.000000}%
\pgfsetstrokecolor{currentstroke}%
\pgfsetdash{}{0pt}%
\pgfpathmoveto{\pgfqpoint{0.480342in}{1.678010in}}%
\pgfpathlineto{\pgfqpoint{0.417148in}{1.658919in}}%
\pgfusepath{stroke}%
\end{pgfscope}%
\begin{pgfscope}%
\definecolor{textcolor}{rgb}{0.000000,0.000000,0.000000}%
\pgfsetstrokecolor{textcolor}%
\pgfsetfillcolor{textcolor}%
\pgftext[x=0.299177in,y=1.493571in,,top]{\color{textcolor}{\rmfamily\fontsize{10.000000}{12.000000}\selectfont\catcode`\^=\active\def^{\ifmmode\sp\else\^{}\fi}\catcode`\%=\active\def%{\%}$\mathdefault{0}$}}%
\end{pgfscope}%
\begin{pgfscope}%
\pgfsetrectcap%
\pgfsetroundjoin%
\pgfsetlinewidth{0.501875pt}%
\definecolor{currentstroke}{rgb}{0.000000,0.000000,0.000000}%
\pgfsetstrokecolor{currentstroke}%
\pgfsetdash{}{0pt}%
\pgfpathmoveto{\pgfqpoint{0.804637in}{1.444389in}}%
\pgfpathlineto{\pgfqpoint{0.740803in}{1.424617in}}%
\pgfusepath{stroke}%
\end{pgfscope}%
\begin{pgfscope}%
\definecolor{textcolor}{rgb}{0.000000,0.000000,0.000000}%
\pgfsetstrokecolor{textcolor}%
\pgfsetfillcolor{textcolor}%
\pgftext[x=0.620912in,y=1.256225in,,top]{\color{textcolor}{\rmfamily\fontsize{10.000000}{12.000000}\selectfont\catcode`\^=\active\def^{\ifmmode\sp\else\^{}\fi}\catcode`\%=\active\def%{\%}$\mathdefault{80}$}}%
\end{pgfscope}%
\begin{pgfscope}%
\pgfsetrectcap%
\pgfsetroundjoin%
\pgfsetlinewidth{0.501875pt}%
\definecolor{currentstroke}{rgb}{0.000000,0.000000,0.000000}%
\pgfsetstrokecolor{currentstroke}%
\pgfsetdash{}{0pt}%
\pgfpathmoveto{\pgfqpoint{1.140584in}{1.202374in}}%
\pgfpathlineto{\pgfqpoint{1.076102in}{1.181885in}}%
\pgfusepath{stroke}%
\end{pgfscope}%
\begin{pgfscope}%
\definecolor{textcolor}{rgb}{0.000000,0.000000,0.000000}%
\pgfsetstrokecolor{textcolor}%
\pgfsetfillcolor{textcolor}%
\pgftext[x=0.954230in,y=1.010334in,,top]{\color{textcolor}{\rmfamily\fontsize{10.000000}{12.000000}\selectfont\catcode`\^=\active\def^{\ifmmode\sp\else\^{}\fi}\catcode`\%=\active\def%{\%}$\mathdefault{160}$}}%
\end{pgfscope}%
\begin{pgfscope}%
\pgfsetrectcap%
\pgfsetroundjoin%
\pgfsetlinewidth{0.501875pt}%
\definecolor{currentstroke}{rgb}{0.000000,0.000000,0.000000}%
\pgfsetstrokecolor{currentstroke}%
\pgfsetdash{}{0pt}%
\pgfpathmoveto{\pgfqpoint{1.488820in}{0.951505in}}%
\pgfpathlineto{\pgfqpoint{1.423687in}{0.930259in}}%
\pgfusepath{stroke}%
\end{pgfscope}%
\begin{pgfscope}%
\definecolor{textcolor}{rgb}{0.000000,0.000000,0.000000}%
\pgfsetstrokecolor{textcolor}%
\pgfsetfillcolor{textcolor}%
\pgftext[x=1.299767in,y=0.755429in,,top]{\color{textcolor}{\rmfamily\fontsize{10.000000}{12.000000}\selectfont\catcode`\^=\active\def^{\ifmmode\sp\else\^{}\fi}\catcode`\%=\active\def%{\%}$\mathdefault{240}$}}%
\end{pgfscope}%
\begin{pgfscope}%
\pgfsetrectcap%
\pgfsetroundjoin%
\pgfsetlinewidth{0.501875pt}%
\definecolor{currentstroke}{rgb}{0.000000,0.000000,0.000000}%
\pgfsetstrokecolor{currentstroke}%
\pgfsetdash{}{0pt}%
\pgfpathmoveto{\pgfqpoint{1.850033in}{0.691287in}}%
\pgfpathlineto{\pgfqpoint{1.784245in}{0.669242in}}%
\pgfusepath{stroke}%
\end{pgfscope}%
\begin{pgfscope}%
\definecolor{textcolor}{rgb}{0.000000,0.000000,0.000000}%
\pgfsetstrokecolor{textcolor}%
\pgfsetfillcolor{textcolor}%
\pgftext[x=1.658207in,y=0.491006in,,top]{\color{textcolor}{\rmfamily\fontsize{10.000000}{12.000000}\selectfont\catcode`\^=\active\def^{\ifmmode\sp\else\^{}\fi}\catcode`\%=\active\def%{\%}$\mathdefault{320}$}}%
\end{pgfscope}%
\begin{pgfscope}%
\pgfsetrectcap%
\pgfsetroundjoin%
\pgfsetlinewidth{0.501875pt}%
\definecolor{currentstroke}{rgb}{0.000000,0.000000,0.000000}%
\pgfsetstrokecolor{currentstroke}%
\pgfsetdash{}{0pt}%
\pgfpathmoveto{\pgfqpoint{0.560357in}{1.620368in}}%
\pgfpathlineto{\pgfqpoint{0.497004in}{1.601110in}}%
\pgfusepath{stroke}%
\end{pgfscope}%
\begin{pgfscope}%
\pgfsetrectcap%
\pgfsetroundjoin%
\pgfsetlinewidth{0.501875pt}%
\definecolor{currentstroke}{rgb}{0.000000,0.000000,0.000000}%
\pgfsetstrokecolor{currentstroke}%
\pgfsetdash{}{0pt}%
\pgfpathmoveto{\pgfqpoint{0.641072in}{1.562221in}}%
\pgfpathlineto{\pgfqpoint{0.577558in}{1.542794in}}%
\pgfusepath{stroke}%
\end{pgfscope}%
\begin{pgfscope}%
\pgfsetrectcap%
\pgfsetroundjoin%
\pgfsetlinewidth{0.501875pt}%
\definecolor{currentstroke}{rgb}{0.000000,0.000000,0.000000}%
\pgfsetstrokecolor{currentstroke}%
\pgfsetdash{}{0pt}%
\pgfpathmoveto{\pgfqpoint{0.722495in}{1.503564in}}%
\pgfpathlineto{\pgfqpoint{0.658821in}{1.483966in}}%
\pgfusepath{stroke}%
\end{pgfscope}%
\begin{pgfscope}%
\pgfsetrectcap%
\pgfsetroundjoin%
\pgfsetlinewidth{0.501875pt}%
\definecolor{currentstroke}{rgb}{0.000000,0.000000,0.000000}%
\pgfsetstrokecolor{currentstroke}%
\pgfsetdash{}{0pt}%
\pgfpathmoveto{\pgfqpoint{0.887508in}{1.384689in}}%
\pgfpathlineto{\pgfqpoint{0.823512in}{1.364742in}}%
\pgfusepath{stroke}%
\end{pgfscope}%
\begin{pgfscope}%
\pgfsetrectcap%
\pgfsetroundjoin%
\pgfsetlinewidth{0.501875pt}%
\definecolor{currentstroke}{rgb}{0.000000,0.000000,0.000000}%
\pgfsetstrokecolor{currentstroke}%
\pgfsetdash{}{0pt}%
\pgfpathmoveto{\pgfqpoint{0.971115in}{1.324458in}}%
\pgfpathlineto{\pgfqpoint{0.906958in}{1.304333in}}%
\pgfusepath{stroke}%
\end{pgfscope}%
\begin{pgfscope}%
\pgfsetrectcap%
\pgfsetroundjoin%
\pgfsetlinewidth{0.501875pt}%
\definecolor{currentstroke}{rgb}{0.000000,0.000000,0.000000}%
\pgfsetstrokecolor{currentstroke}%
\pgfsetdash{}{0pt}%
\pgfpathmoveto{\pgfqpoint{1.055471in}{1.263689in}}%
\pgfpathlineto{\pgfqpoint{0.991151in}{1.243383in}}%
\pgfusepath{stroke}%
\end{pgfscope}%
\begin{pgfscope}%
\pgfsetrectcap%
\pgfsetroundjoin%
\pgfsetlinewidth{0.501875pt}%
\definecolor{currentstroke}{rgb}{0.000000,0.000000,0.000000}%
\pgfsetstrokecolor{currentstroke}%
\pgfsetdash{}{0pt}%
\pgfpathmoveto{\pgfqpoint{1.226464in}{1.140505in}}%
\pgfpathlineto{\pgfqpoint{1.161821in}{1.119831in}}%
\pgfusepath{stroke}%
\end{pgfscope}%
\begin{pgfscope}%
\pgfsetrectcap%
\pgfsetroundjoin%
\pgfsetlinewidth{0.501875pt}%
\definecolor{currentstroke}{rgb}{0.000000,0.000000,0.000000}%
\pgfsetstrokecolor{currentstroke}%
\pgfsetdash{}{0pt}%
\pgfpathmoveto{\pgfqpoint{1.313124in}{1.078076in}}%
\pgfpathlineto{\pgfqpoint{1.248317in}{1.057214in}}%
\pgfusepath{stroke}%
\end{pgfscope}%
\begin{pgfscope}%
\pgfsetrectcap%
\pgfsetroundjoin%
\pgfsetlinewidth{0.501875pt}%
\definecolor{currentstroke}{rgb}{0.000000,0.000000,0.000000}%
\pgfsetstrokecolor{currentstroke}%
\pgfsetdash{}{0pt}%
\pgfpathmoveto{\pgfqpoint{1.400572in}{1.015078in}}%
\pgfpathlineto{\pgfqpoint{1.335602in}{0.994026in}}%
\pgfusepath{stroke}%
\end{pgfscope}%
\begin{pgfscope}%
\pgfsetrectcap%
\pgfsetroundjoin%
\pgfsetlinewidth{0.501875pt}%
\definecolor{currentstroke}{rgb}{0.000000,0.000000,0.000000}%
\pgfsetstrokecolor{currentstroke}%
\pgfsetdash{}{0pt}%
\pgfpathmoveto{\pgfqpoint{1.577879in}{0.887347in}}%
\pgfpathlineto{\pgfqpoint{1.512582in}{0.865906in}}%
\pgfusepath{stroke}%
\end{pgfscope}%
\begin{pgfscope}%
\pgfsetrectcap%
\pgfsetroundjoin%
\pgfsetlinewidth{0.501875pt}%
\definecolor{currentstroke}{rgb}{0.000000,0.000000,0.000000}%
\pgfsetstrokecolor{currentstroke}%
\pgfsetdash{}{0pt}%
\pgfpathmoveto{\pgfqpoint{1.667760in}{0.822597in}}%
\pgfpathlineto{\pgfqpoint{1.602299in}{0.800957in}}%
\pgfusepath{stroke}%
\end{pgfscope}%
\begin{pgfscope}%
\pgfsetrectcap%
\pgfsetroundjoin%
\pgfsetlinewidth{0.501875pt}%
\definecolor{currentstroke}{rgb}{0.000000,0.000000,0.000000}%
\pgfsetstrokecolor{currentstroke}%
\pgfsetdash{}{0pt}%
\pgfpathmoveto{\pgfqpoint{1.758474in}{0.757246in}}%
\pgfpathlineto{\pgfqpoint{1.692850in}{0.735405in}}%
\pgfusepath{stroke}%
\end{pgfscope}%
\begin{pgfscope}%
\pgfsetrectcap%
\pgfsetroundjoin%
\pgfsetlinewidth{0.501875pt}%
\definecolor{currentstroke}{rgb}{0.000000,0.000000,0.000000}%
\pgfsetstrokecolor{currentstroke}%
\pgfsetdash{}{0pt}%
\pgfpathmoveto{\pgfqpoint{1.942450in}{0.624710in}}%
\pgfpathlineto{\pgfqpoint{1.876497in}{0.602458in}}%
\pgfusepath{stroke}%
\end{pgfscope}%
\begin{pgfscope}%
\pgfsetrectcap%
\pgfsetroundjoin%
\pgfsetlinewidth{0.501875pt}%
\definecolor{currentstroke}{rgb}{0.000000,0.000000,0.000000}%
\pgfsetstrokecolor{currentstroke}%
\pgfsetdash{}{0pt}%
\pgfpathmoveto{\pgfqpoint{2.035735in}{0.557508in}}%
\pgfpathlineto{\pgfqpoint{1.969618in}{0.535046in}}%
\pgfusepath{stroke}%
\end{pgfscope}%
\begin{pgfscope}%
\pgfsetrectcap%
\pgfsetroundjoin%
\pgfsetlinewidth{0.501875pt}%
\definecolor{currentstroke}{rgb}{0.000000,0.000000,0.000000}%
\pgfsetstrokecolor{currentstroke}%
\pgfsetdash{}{0pt}%
\pgfpathmoveto{\pgfqpoint{2.129902in}{0.489670in}}%
\pgfpathlineto{\pgfqpoint{2.063620in}{0.466995in}}%
\pgfusepath{stroke}%
\end{pgfscope}%
\begin{pgfscope}%
\definecolor{textcolor}{rgb}{0.000000,0.000000,0.000000}%
\pgfsetstrokecolor{textcolor}%
\pgfsetfillcolor{textcolor}%
\pgftext[x=0.883212in,y=0.671627in,,,rotate=324.186895]{\color{textcolor}{\rmfamily\fontsize{10.000000}{12.000000}\selectfont\catcode`\^=\active\def^{\ifmmode\sp\else\^{}\fi}\catcode`\%=\active\def%{\%}Time to Expiry (Days)}}%
\end{pgfscope}%
\begin{pgfscope}%
\pgfsetrectcap%
\pgfsetroundjoin%
\pgfsetlinewidth{0.501875pt}%
\definecolor{currentstroke}{rgb}{0.000000,0.000000,0.000000}%
\pgfsetstrokecolor{currentstroke}%
\pgfsetdash{}{0pt}%
\pgfpathmoveto{\pgfqpoint{4.749691in}{1.337846in}}%
\pgfpathlineto{\pgfqpoint{2.155964in}{0.447388in}}%
\pgfusepath{stroke}%
\end{pgfscope}%
\begin{pgfscope}%
\pgfsetrectcap%
\pgfsetroundjoin%
\pgfsetlinewidth{0.501875pt}%
\definecolor{currentstroke}{rgb}{0.000000,0.000000,0.000000}%
\pgfsetstrokecolor{currentstroke}%
\pgfsetdash{}{0pt}%
\pgfpathmoveto{\pgfqpoint{2.144293in}{0.459367in}}%
\pgfpathlineto{\pgfqpoint{2.189402in}{0.426825in}}%
\pgfusepath{stroke}%
\end{pgfscope}%
\begin{pgfscope}%
\definecolor{textcolor}{rgb}{0.000000,0.000000,0.000000}%
\pgfsetstrokecolor{textcolor}%
\pgfsetfillcolor{textcolor}%
\pgftext[x=2.278268in,y=0.226168in,,top]{\color{textcolor}{\rmfamily\fontsize{10.000000}{12.000000}\selectfont\catcode`\^=\active\def^{\ifmmode\sp\else\^{}\fi}\catcode`\%=\active\def%{\%}$\mathdefault{45}$}}%
\end{pgfscope}%
\begin{pgfscope}%
\pgfsetrectcap%
\pgfsetroundjoin%
\pgfsetlinewidth{0.501875pt}%
\definecolor{currentstroke}{rgb}{0.000000,0.000000,0.000000}%
\pgfsetstrokecolor{currentstroke}%
\pgfsetdash{}{0pt}%
\pgfpathmoveto{\pgfqpoint{2.655344in}{0.634623in}}%
\pgfpathlineto{\pgfqpoint{2.701042in}{0.602867in}}%
\pgfusepath{stroke}%
\end{pgfscope}%
\begin{pgfscope}%
\definecolor{textcolor}{rgb}{0.000000,0.000000,0.000000}%
\pgfsetstrokecolor{textcolor}%
\pgfsetfillcolor{textcolor}%
\pgftext[x=2.789770in,y=0.404976in,,top]{\color{textcolor}{\rmfamily\fontsize{10.000000}{12.000000}\selectfont\catcode`\^=\active\def^{\ifmmode\sp\else\^{}\fi}\catcode`\%=\active\def%{\%}$\mathdefault{60}$}}%
\end{pgfscope}%
\begin{pgfscope}%
\pgfsetrectcap%
\pgfsetroundjoin%
\pgfsetlinewidth{0.501875pt}%
\definecolor{currentstroke}{rgb}{0.000000,0.000000,0.000000}%
\pgfsetstrokecolor{currentstroke}%
\pgfsetdash{}{0pt}%
\pgfpathmoveto{\pgfqpoint{3.154154in}{0.805682in}}%
\pgfpathlineto{\pgfqpoint{3.200399in}{0.774682in}}%
\pgfusepath{stroke}%
\end{pgfscope}%
\begin{pgfscope}%
\definecolor{textcolor}{rgb}{0.000000,0.000000,0.000000}%
\pgfsetstrokecolor{textcolor}%
\pgfsetfillcolor{textcolor}%
\pgftext[x=3.288970in,y=0.579484in,,top]{\color{textcolor}{\rmfamily\fontsize{10.000000}{12.000000}\selectfont\catcode`\^=\active\def^{\ifmmode\sp\else\^{}\fi}\catcode`\%=\active\def%{\%}$\mathdefault{75}$}}%
\end{pgfscope}%
\begin{pgfscope}%
\pgfsetrectcap%
\pgfsetroundjoin%
\pgfsetlinewidth{0.501875pt}%
\definecolor{currentstroke}{rgb}{0.000000,0.000000,0.000000}%
\pgfsetstrokecolor{currentstroke}%
\pgfsetdash{}{0pt}%
\pgfpathmoveto{\pgfqpoint{3.641158in}{0.972691in}}%
\pgfpathlineto{\pgfqpoint{3.687911in}{0.942422in}}%
\pgfusepath{stroke}%
\end{pgfscope}%
\begin{pgfscope}%
\definecolor{textcolor}{rgb}{0.000000,0.000000,0.000000}%
\pgfsetstrokecolor{textcolor}%
\pgfsetfillcolor{textcolor}%
\pgftext[x=3.776307in,y=0.749845in,,top]{\color{textcolor}{\rmfamily\fontsize{10.000000}{12.000000}\selectfont\catcode`\^=\active\def^{\ifmmode\sp\else\^{}\fi}\catcode`\%=\active\def%{\%}$\mathdefault{90}$}}%
\end{pgfscope}%
\begin{pgfscope}%
\pgfsetrectcap%
\pgfsetroundjoin%
\pgfsetlinewidth{0.501875pt}%
\definecolor{currentstroke}{rgb}{0.000000,0.000000,0.000000}%
\pgfsetstrokecolor{currentstroke}%
\pgfsetdash{}{0pt}%
\pgfpathmoveto{\pgfqpoint{4.116768in}{1.135794in}}%
\pgfpathlineto{\pgfqpoint{4.163994in}{1.106229in}}%
\pgfusepath{stroke}%
\end{pgfscope}%
\begin{pgfscope}%
\definecolor{textcolor}{rgb}{0.000000,0.000000,0.000000}%
\pgfsetstrokecolor{textcolor}%
\pgfsetfillcolor{textcolor}%
\pgftext[x=4.252199in,y=0.916205in,,top]{\color{textcolor}{\rmfamily\fontsize{10.000000}{12.000000}\selectfont\catcode`\^=\active\def^{\ifmmode\sp\else\^{}\fi}\catcode`\%=\active\def%{\%}$\mathdefault{105}$}}%
\end{pgfscope}%
\begin{pgfscope}%
\pgfsetrectcap%
\pgfsetroundjoin%
\pgfsetlinewidth{0.501875pt}%
\definecolor{currentstroke}{rgb}{0.000000,0.000000,0.000000}%
\pgfsetstrokecolor{currentstroke}%
\pgfsetdash{}{0pt}%
\pgfpathmoveto{\pgfqpoint{4.581381in}{1.295125in}}%
\pgfpathlineto{\pgfqpoint{4.629045in}{1.266240in}}%
\pgfusepath{stroke}%
\end{pgfscope}%
\begin{pgfscope}%
\definecolor{textcolor}{rgb}{0.000000,0.000000,0.000000}%
\pgfsetstrokecolor{textcolor}%
\pgfsetfillcolor{textcolor}%
\pgftext[x=4.717044in,y=1.078703in,,top]{\color{textcolor}{\rmfamily\fontsize{10.000000}{12.000000}\selectfont\catcode`\^=\active\def^{\ifmmode\sp\else\^{}\fi}\catcode`\%=\active\def%{\%}$\mathdefault{120}$}}%
\end{pgfscope}%
\begin{pgfscope}%
\pgfsetrectcap%
\pgfsetroundjoin%
\pgfsetlinewidth{0.501875pt}%
\definecolor{currentstroke}{rgb}{0.000000,0.000000,0.000000}%
\pgfsetstrokecolor{currentstroke}%
\pgfsetdash{}{0pt}%
\pgfpathmoveto{\pgfqpoint{2.273228in}{0.503583in}}%
\pgfpathlineto{\pgfqpoint{2.318488in}{0.471240in}}%
\pgfusepath{stroke}%
\end{pgfscope}%
\begin{pgfscope}%
\pgfsetrectcap%
\pgfsetroundjoin%
\pgfsetlinewidth{0.501875pt}%
\definecolor{currentstroke}{rgb}{0.000000,0.000000,0.000000}%
\pgfsetstrokecolor{currentstroke}%
\pgfsetdash{}{0pt}%
\pgfpathmoveto{\pgfqpoint{2.401377in}{0.547529in}}%
\pgfpathlineto{\pgfqpoint{2.446786in}{0.515384in}}%
\pgfusepath{stroke}%
\end{pgfscope}%
\begin{pgfscope}%
\pgfsetrectcap%
\pgfsetroundjoin%
\pgfsetlinewidth{0.501875pt}%
\definecolor{currentstroke}{rgb}{0.000000,0.000000,0.000000}%
\pgfsetstrokecolor{currentstroke}%
\pgfsetdash{}{0pt}%
\pgfpathmoveto{\pgfqpoint{2.528747in}{0.591209in}}%
\pgfpathlineto{\pgfqpoint{2.574301in}{0.559259in}}%
\pgfusepath{stroke}%
\end{pgfscope}%
\begin{pgfscope}%
\pgfsetrectcap%
\pgfsetroundjoin%
\pgfsetlinewidth{0.501875pt}%
\definecolor{currentstroke}{rgb}{0.000000,0.000000,0.000000}%
\pgfsetstrokecolor{currentstroke}%
\pgfsetdash{}{0pt}%
\pgfpathmoveto{\pgfqpoint{2.781177in}{0.677775in}}%
\pgfpathlineto{\pgfqpoint{2.827015in}{0.646211in}}%
\pgfusepath{stroke}%
\end{pgfscope}%
\begin{pgfscope}%
\pgfsetrectcap%
\pgfsetroundjoin%
\pgfsetlinewidth{0.501875pt}%
\definecolor{currentstroke}{rgb}{0.000000,0.000000,0.000000}%
\pgfsetstrokecolor{currentstroke}%
\pgfsetdash{}{0pt}%
\pgfpathmoveto{\pgfqpoint{2.906252in}{0.720668in}}%
\pgfpathlineto{\pgfqpoint{2.952228in}{0.689293in}}%
\pgfusepath{stroke}%
\end{pgfscope}%
\begin{pgfscope}%
\pgfsetrectcap%
\pgfsetroundjoin%
\pgfsetlinewidth{0.501875pt}%
\definecolor{currentstroke}{rgb}{0.000000,0.000000,0.000000}%
\pgfsetstrokecolor{currentstroke}%
\pgfsetdash{}{0pt}%
\pgfpathmoveto{\pgfqpoint{3.030575in}{0.763302in}}%
\pgfpathlineto{\pgfqpoint{3.076687in}{0.732116in}}%
\pgfusepath{stroke}%
\end{pgfscope}%
\begin{pgfscope}%
\pgfsetrectcap%
\pgfsetroundjoin%
\pgfsetlinewidth{0.501875pt}%
\definecolor{currentstroke}{rgb}{0.000000,0.000000,0.000000}%
\pgfsetstrokecolor{currentstroke}%
\pgfsetdash{}{0pt}%
\pgfpathmoveto{\pgfqpoint{3.276996in}{0.847808in}}%
\pgfpathlineto{\pgfqpoint{3.323371in}{0.816993in}}%
\pgfusepath{stroke}%
\end{pgfscope}%
\begin{pgfscope}%
\pgfsetrectcap%
\pgfsetroundjoin%
\pgfsetlinewidth{0.501875pt}%
\definecolor{currentstroke}{rgb}{0.000000,0.000000,0.000000}%
\pgfsetstrokecolor{currentstroke}%
\pgfsetdash{}{0pt}%
\pgfpathmoveto{\pgfqpoint{3.399106in}{0.889683in}}%
\pgfpathlineto{\pgfqpoint{3.445610in}{0.859052in}}%
\pgfusepath{stroke}%
\end{pgfscope}%
\begin{pgfscope}%
\pgfsetrectcap%
\pgfsetroundjoin%
\pgfsetlinewidth{0.501875pt}%
\definecolor{currentstroke}{rgb}{0.000000,0.000000,0.000000}%
\pgfsetstrokecolor{currentstroke}%
\pgfsetdash{}{0pt}%
\pgfpathmoveto{\pgfqpoint{3.520491in}{0.931310in}}%
\pgfpathlineto{\pgfqpoint{3.567121in}{0.900861in}}%
\pgfusepath{stroke}%
\end{pgfscope}%
\begin{pgfscope}%
\pgfsetrectcap%
\pgfsetroundjoin%
\pgfsetlinewidth{0.501875pt}%
\definecolor{currentstroke}{rgb}{0.000000,0.000000,0.000000}%
\pgfsetstrokecolor{currentstroke}%
\pgfsetdash{}{0pt}%
\pgfpathmoveto{\pgfqpoint{3.761113in}{1.013828in}}%
\pgfpathlineto{\pgfqpoint{3.807988in}{0.983737in}}%
\pgfusepath{stroke}%
\end{pgfscope}%
\begin{pgfscope}%
\pgfsetrectcap%
\pgfsetroundjoin%
\pgfsetlinewidth{0.501875pt}%
\definecolor{currentstroke}{rgb}{0.000000,0.000000,0.000000}%
\pgfsetstrokecolor{currentstroke}%
\pgfsetdash{}{0pt}%
\pgfpathmoveto{\pgfqpoint{3.880362in}{1.054722in}}%
\pgfpathlineto{\pgfqpoint{3.927356in}{1.024808in}}%
\pgfusepath{stroke}%
\end{pgfscope}%
\begin{pgfscope}%
\pgfsetrectcap%
\pgfsetroundjoin%
\pgfsetlinewidth{0.501875pt}%
\definecolor{currentstroke}{rgb}{0.000000,0.000000,0.000000}%
\pgfsetstrokecolor{currentstroke}%
\pgfsetdash{}{0pt}%
\pgfpathmoveto{\pgfqpoint{3.998912in}{1.095377in}}%
\pgfpathlineto{\pgfqpoint{4.046023in}{1.065638in}}%
\pgfusepath{stroke}%
\end{pgfscope}%
\begin{pgfscope}%
\pgfsetrectcap%
\pgfsetroundjoin%
\pgfsetlinewidth{0.501875pt}%
\definecolor{currentstroke}{rgb}{0.000000,0.000000,0.000000}%
\pgfsetstrokecolor{currentstroke}%
\pgfsetdash{}{0pt}%
\pgfpathmoveto{\pgfqpoint{4.233937in}{1.175975in}}%
\pgfpathlineto{\pgfqpoint{4.281276in}{1.146582in}}%
\pgfusepath{stroke}%
\end{pgfscope}%
\begin{pgfscope}%
\pgfsetrectcap%
\pgfsetroundjoin%
\pgfsetlinewidth{0.501875pt}%
\definecolor{currentstroke}{rgb}{0.000000,0.000000,0.000000}%
\pgfsetstrokecolor{currentstroke}%
\pgfsetdash{}{0pt}%
\pgfpathmoveto{\pgfqpoint{4.350425in}{1.215922in}}%
\pgfpathlineto{\pgfqpoint{4.397874in}{1.186701in}}%
\pgfusepath{stroke}%
\end{pgfscope}%
\begin{pgfscope}%
\pgfsetrectcap%
\pgfsetroundjoin%
\pgfsetlinewidth{0.501875pt}%
\definecolor{currentstroke}{rgb}{0.000000,0.000000,0.000000}%
\pgfsetstrokecolor{currentstroke}%
\pgfsetdash{}{0pt}%
\pgfpathmoveto{\pgfqpoint{4.466238in}{1.255638in}}%
\pgfpathlineto{\pgfqpoint{4.513796in}{1.226586in}}%
\pgfusepath{stroke}%
\end{pgfscope}%
\begin{pgfscope}%
\pgfsetrectcap%
\pgfsetroundjoin%
\pgfsetlinewidth{0.501875pt}%
\definecolor{currentstroke}{rgb}{0.000000,0.000000,0.000000}%
\pgfsetstrokecolor{currentstroke}%
\pgfsetdash{}{0pt}%
\pgfpathmoveto{\pgfqpoint{4.695861in}{1.334383in}}%
\pgfpathlineto{\pgfqpoint{4.743629in}{1.305665in}}%
\pgfusepath{stroke}%
\end{pgfscope}%
\begin{pgfscope}%
\definecolor{textcolor}{rgb}{0.000000,0.000000,0.000000}%
\pgfsetstrokecolor{textcolor}%
\pgfsetfillcolor{textcolor}%
\pgftext[x=3.765372in,y=0.421364in,,,rotate=18.947963]{\color{textcolor}{\rmfamily\fontsize{10.000000}{12.000000}\selectfont\catcode`\^=\active\def^{\ifmmode\sp\else\^{}\fi}\catcode`\%=\active\def%{\%}Underlying (S)}}%
\end{pgfscope}%
\begin{pgfscope}%
\pgfsetrectcap%
\pgfsetroundjoin%
\pgfsetlinewidth{0.501875pt}%
\definecolor{currentstroke}{rgb}{0.000000,0.000000,0.000000}%
\pgfsetstrokecolor{currentstroke}%
\pgfsetdash{}{0pt}%
\pgfpathmoveto{\pgfqpoint{4.749691in}{1.337846in}}%
\pgfpathlineto{\pgfqpoint{4.834815in}{3.504368in}}%
\pgfusepath{stroke}%
\end{pgfscope}%
\begin{pgfscope}%
\pgfsetrectcap%
\pgfsetroundjoin%
\pgfsetlinewidth{0.501875pt}%
\definecolor{currentstroke}{rgb}{0.000000,0.000000,0.000000}%
\pgfsetstrokecolor{currentstroke}%
\pgfsetdash{}{0pt}%
\pgfpathmoveto{\pgfqpoint{4.735406in}{1.389101in}}%
\pgfpathlineto{\pgfqpoint{4.783248in}{1.360504in}}%
\pgfusepath{stroke}%
\end{pgfscope}%
\begin{pgfscope}%
\definecolor{textcolor}{rgb}{0.000000,0.000000,0.000000}%
\pgfsetstrokecolor{textcolor}%
\pgfsetfillcolor{textcolor}%
\pgftext[x=5.028906in,y=1.357205in,,top]{\color{textcolor}{\rmfamily\fontsize{10.000000}{12.000000}\selectfont\catcode`\^=\active\def^{\ifmmode\sp\else\^{}\fi}\catcode`\%=\active\def%{\%}$\mathdefault{0}$}}%
\end{pgfscope}%
\begin{pgfscope}%
\pgfsetrectcap%
\pgfsetroundjoin%
\pgfsetlinewidth{0.501875pt}%
\definecolor{currentstroke}{rgb}{0.000000,0.000000,0.000000}%
\pgfsetstrokecolor{currentstroke}%
\pgfsetdash{}{0pt}%
\pgfpathmoveto{\pgfqpoint{4.751165in}{1.793206in}}%
\pgfpathlineto{\pgfqpoint{4.799389in}{1.765262in}}%
\pgfusepath{stroke}%
\end{pgfscope}%
\begin{pgfscope}%
\definecolor{textcolor}{rgb}{0.000000,0.000000,0.000000}%
\pgfsetstrokecolor{textcolor}%
\pgfsetfillcolor{textcolor}%
\pgftext[x=5.046840in,y=1.762038in,,top]{\color{textcolor}{\rmfamily\fontsize{10.000000}{12.000000}\selectfont\catcode`\^=\active\def^{\ifmmode\sp\else\^{}\fi}\catcode`\%=\active\def%{\%}$\mathdefault{10}$}}%
\end{pgfscope}%
\begin{pgfscope}%
\pgfsetrectcap%
\pgfsetroundjoin%
\pgfsetlinewidth{0.501875pt}%
\definecolor{currentstroke}{rgb}{0.000000,0.000000,0.000000}%
\pgfsetstrokecolor{currentstroke}%
\pgfsetdash{}{0pt}%
\pgfpathmoveto{\pgfqpoint{4.767155in}{2.203227in}}%
\pgfpathlineto{\pgfqpoint{4.815767in}{2.175958in}}%
\pgfusepath{stroke}%
\end{pgfscope}%
\begin{pgfscope}%
\definecolor{textcolor}{rgb}{0.000000,0.000000,0.000000}%
\pgfsetstrokecolor{textcolor}%
\pgfsetfillcolor{textcolor}%
\pgftext[x=5.065037in,y=2.172811in,,top]{\color{textcolor}{\rmfamily\fontsize{10.000000}{12.000000}\selectfont\catcode`\^=\active\def^{\ifmmode\sp\else\^{}\fi}\catcode`\%=\active\def%{\%}$\mathdefault{20}$}}%
\end{pgfscope}%
\begin{pgfscope}%
\pgfsetrectcap%
\pgfsetroundjoin%
\pgfsetlinewidth{0.501875pt}%
\definecolor{currentstroke}{rgb}{0.000000,0.000000,0.000000}%
\pgfsetstrokecolor{currentstroke}%
\pgfsetdash{}{0pt}%
\pgfpathmoveto{\pgfqpoint{4.783381in}{2.619293in}}%
\pgfpathlineto{\pgfqpoint{4.832387in}{2.592721in}}%
\pgfusepath{stroke}%
\end{pgfscope}%
\begin{pgfscope}%
\definecolor{textcolor}{rgb}{0.000000,0.000000,0.000000}%
\pgfsetstrokecolor{textcolor}%
\pgfsetfillcolor{textcolor}%
\pgftext[x=5.083503in,y=2.589655in,,top]{\color{textcolor}{\rmfamily\fontsize{10.000000}{12.000000}\selectfont\catcode`\^=\active\def^{\ifmmode\sp\else\^{}\fi}\catcode`\%=\active\def%{\%}$\mathdefault{30}$}}%
\end{pgfscope}%
\begin{pgfscope}%
\pgfsetrectcap%
\pgfsetroundjoin%
\pgfsetlinewidth{0.501875pt}%
\definecolor{currentstroke}{rgb}{0.000000,0.000000,0.000000}%
\pgfsetstrokecolor{currentstroke}%
\pgfsetdash{}{0pt}%
\pgfpathmoveto{\pgfqpoint{4.799847in}{3.041539in}}%
\pgfpathlineto{\pgfqpoint{4.849254in}{3.015689in}}%
\pgfusepath{stroke}%
\end{pgfscope}%
\begin{pgfscope}%
\definecolor{textcolor}{rgb}{0.000000,0.000000,0.000000}%
\pgfsetstrokecolor{textcolor}%
\pgfsetfillcolor{textcolor}%
\pgftext[x=5.102244in,y=3.012706in,,top]{\color{textcolor}{\rmfamily\fontsize{10.000000}{12.000000}\selectfont\catcode`\^=\active\def^{\ifmmode\sp\else\^{}\fi}\catcode`\%=\active\def%{\%}$\mathdefault{40}$}}%
\end{pgfscope}%
\begin{pgfscope}%
\pgfsetrectcap%
\pgfsetroundjoin%
\pgfsetlinewidth{0.501875pt}%
\definecolor{currentstroke}{rgb}{0.000000,0.000000,0.000000}%
\pgfsetstrokecolor{currentstroke}%
\pgfsetdash{}{0pt}%
\pgfpathmoveto{\pgfqpoint{4.816561in}{3.470104in}}%
\pgfpathlineto{\pgfqpoint{4.866374in}{3.445000in}}%
\pgfusepath{stroke}%
\end{pgfscope}%
\begin{pgfscope}%
\definecolor{textcolor}{rgb}{0.000000,0.000000,0.000000}%
\pgfsetstrokecolor{textcolor}%
\pgfsetfillcolor{textcolor}%
\pgftext[x=5.121266in,y=3.442103in,,top]{\color{textcolor}{\rmfamily\fontsize{10.000000}{12.000000}\selectfont\catcode`\^=\active\def^{\ifmmode\sp\else\^{}\fi}\catcode`\%=\active\def%{\%}$\mathdefault{50}$}}%
\end{pgfscope}%
\begin{pgfscope}%
\definecolor{textcolor}{rgb}{0.000000,0.000000,0.000000}%
\pgfsetstrokecolor{textcolor}%
\pgfsetfillcolor{textcolor}%
\pgftext[x=5.432051in,y=2.352657in,,,rotate=87.749957]{\color{textcolor}{\rmfamily\fontsize{10.000000}{12.000000}\selectfont\catcode`\^=\active\def^{\ifmmode\sp\else\^{}\fi}\catcode`\%=\active\def%{\%}Option Price}}%
\end{pgfscope}%
\begin{pgfscope}%
\pgfpathrectangle{\pgfqpoint{0.050000in}{0.050000in}}{\pgfqpoint{4.900000in}{4.900000in}}%
\pgfusepath{clip}%
\pgfsetbuttcap%
\pgfsetroundjoin%
\definecolor{currentfill}{rgb}{1.000000,1.000000,1.000000}%
\pgfsetfillcolor{currentfill}%
\pgfsetfillopacity{0.900000}%
\pgfsetlinewidth{0.501875pt}%
\definecolor{currentstroke}{rgb}{0.200000,0.200000,0.200000}%
\pgfsetstrokecolor{currentstroke}%
\pgfsetdash{}{0pt}%
\pgfpathmoveto{\pgfqpoint{2.835840in}{2.351416in}}%
\pgfpathlineto{\pgfqpoint{2.878017in}{2.364197in}}%
\pgfpathlineto{\pgfqpoint{2.908238in}{2.345936in}}%
\pgfpathlineto{\pgfqpoint{2.866031in}{2.333119in}}%
\pgfpathlineto{\pgfqpoint{2.835840in}{2.351416in}}%
\pgfpathclose%
\pgfusepath{stroke,fill}%
\end{pgfscope}%
\begin{pgfscope}%
\pgfpathrectangle{\pgfqpoint{0.050000in}{0.050000in}}{\pgfqpoint{4.900000in}{4.900000in}}%
\pgfusepath{clip}%
\pgfsetbuttcap%
\pgfsetroundjoin%
\definecolor{currentfill}{rgb}{1.000000,1.000000,1.000000}%
\pgfsetfillcolor{currentfill}%
\pgfsetfillopacity{0.900000}%
\pgfsetlinewidth{0.501875pt}%
\definecolor{currentstroke}{rgb}{0.200000,0.200000,0.200000}%
\pgfsetstrokecolor{currentstroke}%
\pgfsetdash{}{0pt}%
\pgfpathmoveto{\pgfqpoint{2.793577in}{2.338610in}}%
\pgfpathlineto{\pgfqpoint{2.835840in}{2.351416in}}%
\pgfpathlineto{\pgfqpoint{2.866031in}{2.333119in}}%
\pgfpathlineto{\pgfqpoint{2.823737in}{2.320275in}}%
\pgfpathlineto{\pgfqpoint{2.793577in}{2.338610in}}%
\pgfpathclose%
\pgfusepath{stroke,fill}%
\end{pgfscope}%
\begin{pgfscope}%
\pgfpathrectangle{\pgfqpoint{0.050000in}{0.050000in}}{\pgfqpoint{4.900000in}{4.900000in}}%
\pgfusepath{clip}%
\pgfsetbuttcap%
\pgfsetroundjoin%
\definecolor{currentfill}{rgb}{1.000000,1.000000,1.000000}%
\pgfsetfillcolor{currentfill}%
\pgfsetfillopacity{0.900000}%
\pgfsetlinewidth{0.501875pt}%
\definecolor{currentstroke}{rgb}{0.200000,0.200000,0.200000}%
\pgfsetstrokecolor{currentstroke}%
\pgfsetdash{}{0pt}%
\pgfpathmoveto{\pgfqpoint{2.866031in}{2.333119in}}%
\pgfpathlineto{\pgfqpoint{2.908238in}{2.345936in}}%
\pgfpathlineto{\pgfqpoint{2.938548in}{2.327623in}}%
\pgfpathlineto{\pgfqpoint{2.896309in}{2.314767in}}%
\pgfpathlineto{\pgfqpoint{2.866031in}{2.333119in}}%
\pgfpathclose%
\pgfusepath{stroke,fill}%
\end{pgfscope}%
\begin{pgfscope}%
\pgfpathrectangle{\pgfqpoint{0.050000in}{0.050000in}}{\pgfqpoint{4.900000in}{4.900000in}}%
\pgfusepath{clip}%
\pgfsetbuttcap%
\pgfsetroundjoin%
\definecolor{currentfill}{rgb}{1.000000,1.000000,1.000000}%
\pgfsetfillcolor{currentfill}%
\pgfsetfillopacity{0.900000}%
\pgfsetlinewidth{0.501875pt}%
\definecolor{currentstroke}{rgb}{0.200000,0.200000,0.200000}%
\pgfsetstrokecolor{currentstroke}%
\pgfsetdash{}{0pt}%
\pgfpathmoveto{\pgfqpoint{2.751228in}{2.325777in}}%
\pgfpathlineto{\pgfqpoint{2.793577in}{2.338610in}}%
\pgfpathlineto{\pgfqpoint{2.823737in}{2.320275in}}%
\pgfpathlineto{\pgfqpoint{2.781357in}{2.307404in}}%
\pgfpathlineto{\pgfqpoint{2.751228in}{2.325777in}}%
\pgfpathclose%
\pgfusepath{stroke,fill}%
\end{pgfscope}%
\begin{pgfscope}%
\pgfpathrectangle{\pgfqpoint{0.050000in}{0.050000in}}{\pgfqpoint{4.900000in}{4.900000in}}%
\pgfusepath{clip}%
\pgfsetbuttcap%
\pgfsetroundjoin%
\definecolor{currentfill}{rgb}{1.000000,1.000000,1.000000}%
\pgfsetfillcolor{currentfill}%
\pgfsetfillopacity{0.900000}%
\pgfsetlinewidth{0.501875pt}%
\definecolor{currentstroke}{rgb}{0.200000,0.200000,0.200000}%
\pgfsetstrokecolor{currentstroke}%
\pgfsetdash{}{0pt}%
\pgfpathmoveto{\pgfqpoint{2.823737in}{2.320275in}}%
\pgfpathlineto{\pgfqpoint{2.866031in}{2.333119in}}%
\pgfpathlineto{\pgfqpoint{2.896309in}{2.314767in}}%
\pgfpathlineto{\pgfqpoint{2.853985in}{2.301886in}}%
\pgfpathlineto{\pgfqpoint{2.823737in}{2.320275in}}%
\pgfpathclose%
\pgfusepath{stroke,fill}%
\end{pgfscope}%
\begin{pgfscope}%
\pgfpathrectangle{\pgfqpoint{0.050000in}{0.050000in}}{\pgfqpoint{4.900000in}{4.900000in}}%
\pgfusepath{clip}%
\pgfsetbuttcap%
\pgfsetroundjoin%
\definecolor{currentfill}{rgb}{1.000000,1.000000,1.000000}%
\pgfsetfillcolor{currentfill}%
\pgfsetfillopacity{0.900000}%
\pgfsetlinewidth{0.501875pt}%
\definecolor{currentstroke}{rgb}{0.200000,0.200000,0.200000}%
\pgfsetstrokecolor{currentstroke}%
\pgfsetdash{}{0pt}%
\pgfpathmoveto{\pgfqpoint{2.896309in}{2.314767in}}%
\pgfpathlineto{\pgfqpoint{2.938548in}{2.327623in}}%
\pgfpathlineto{\pgfqpoint{2.968946in}{2.309255in}}%
\pgfpathlineto{\pgfqpoint{2.926677in}{2.296363in}}%
\pgfpathlineto{\pgfqpoint{2.896309in}{2.314767in}}%
\pgfpathclose%
\pgfusepath{stroke,fill}%
\end{pgfscope}%
\begin{pgfscope}%
\pgfpathrectangle{\pgfqpoint{0.050000in}{0.050000in}}{\pgfqpoint{4.900000in}{4.900000in}}%
\pgfusepath{clip}%
\pgfsetbuttcap%
\pgfsetroundjoin%
\definecolor{currentfill}{rgb}{1.000000,1.000000,1.000000}%
\pgfsetfillcolor{currentfill}%
\pgfsetfillopacity{0.900000}%
\pgfsetlinewidth{0.501875pt}%
\definecolor{currentstroke}{rgb}{0.200000,0.200000,0.200000}%
\pgfsetstrokecolor{currentstroke}%
\pgfsetdash{}{0pt}%
\pgfpathmoveto{\pgfqpoint{2.708792in}{2.312918in}}%
\pgfpathlineto{\pgfqpoint{2.751228in}{2.325777in}}%
\pgfpathlineto{\pgfqpoint{2.781357in}{2.307404in}}%
\pgfpathlineto{\pgfqpoint{2.738890in}{2.294508in}}%
\pgfpathlineto{\pgfqpoint{2.708792in}{2.312918in}}%
\pgfpathclose%
\pgfusepath{stroke,fill}%
\end{pgfscope}%
\begin{pgfscope}%
\pgfpathrectangle{\pgfqpoint{0.050000in}{0.050000in}}{\pgfqpoint{4.900000in}{4.900000in}}%
\pgfusepath{clip}%
\pgfsetbuttcap%
\pgfsetroundjoin%
\definecolor{currentfill}{rgb}{1.000000,1.000000,1.000000}%
\pgfsetfillcolor{currentfill}%
\pgfsetfillopacity{0.900000}%
\pgfsetlinewidth{0.501875pt}%
\definecolor{currentstroke}{rgb}{0.200000,0.200000,0.200000}%
\pgfsetstrokecolor{currentstroke}%
\pgfsetdash{}{0pt}%
\pgfpathmoveto{\pgfqpoint{2.781357in}{2.307404in}}%
\pgfpathlineto{\pgfqpoint{2.823737in}{2.320275in}}%
\pgfpathlineto{\pgfqpoint{2.853985in}{2.301886in}}%
\pgfpathlineto{\pgfqpoint{2.811574in}{2.288978in}}%
\pgfpathlineto{\pgfqpoint{2.781357in}{2.307404in}}%
\pgfpathclose%
\pgfusepath{stroke,fill}%
\end{pgfscope}%
\begin{pgfscope}%
\pgfpathrectangle{\pgfqpoint{0.050000in}{0.050000in}}{\pgfqpoint{4.900000in}{4.900000in}}%
\pgfusepath{clip}%
\pgfsetbuttcap%
\pgfsetroundjoin%
\definecolor{currentfill}{rgb}{1.000000,1.000000,1.000000}%
\pgfsetfillcolor{currentfill}%
\pgfsetfillopacity{0.900000}%
\pgfsetlinewidth{0.501875pt}%
\definecolor{currentstroke}{rgb}{0.200000,0.200000,0.200000}%
\pgfsetstrokecolor{currentstroke}%
\pgfsetdash{}{0pt}%
\pgfpathmoveto{\pgfqpoint{2.853985in}{2.301886in}}%
\pgfpathlineto{\pgfqpoint{2.896309in}{2.314767in}}%
\pgfpathlineto{\pgfqpoint{2.926677in}{2.296363in}}%
\pgfpathlineto{\pgfqpoint{2.884322in}{2.283443in}}%
\pgfpathlineto{\pgfqpoint{2.853985in}{2.301886in}}%
\pgfpathclose%
\pgfusepath{stroke,fill}%
\end{pgfscope}%
\begin{pgfscope}%
\pgfpathrectangle{\pgfqpoint{0.050000in}{0.050000in}}{\pgfqpoint{4.900000in}{4.900000in}}%
\pgfusepath{clip}%
\pgfsetbuttcap%
\pgfsetroundjoin%
\definecolor{currentfill}{rgb}{1.000000,1.000000,1.000000}%
\pgfsetfillcolor{currentfill}%
\pgfsetfillopacity{0.900000}%
\pgfsetlinewidth{0.501875pt}%
\definecolor{currentstroke}{rgb}{0.200000,0.200000,0.200000}%
\pgfsetstrokecolor{currentstroke}%
\pgfsetdash{}{0pt}%
\pgfpathmoveto{\pgfqpoint{2.666270in}{2.300033in}}%
\pgfpathlineto{\pgfqpoint{2.708792in}{2.312918in}}%
\pgfpathlineto{\pgfqpoint{2.738890in}{2.294508in}}%
\pgfpathlineto{\pgfqpoint{2.696336in}{2.281585in}}%
\pgfpathlineto{\pgfqpoint{2.666270in}{2.300033in}}%
\pgfpathclose%
\pgfusepath{stroke,fill}%
\end{pgfscope}%
\begin{pgfscope}%
\pgfpathrectangle{\pgfqpoint{0.050000in}{0.050000in}}{\pgfqpoint{4.900000in}{4.900000in}}%
\pgfusepath{clip}%
\pgfsetbuttcap%
\pgfsetroundjoin%
\definecolor{currentfill}{rgb}{1.000000,1.000000,1.000000}%
\pgfsetfillcolor{currentfill}%
\pgfsetfillopacity{0.900000}%
\pgfsetlinewidth{0.501875pt}%
\definecolor{currentstroke}{rgb}{0.200000,0.200000,0.200000}%
\pgfsetstrokecolor{currentstroke}%
\pgfsetdash{}{0pt}%
\pgfpathmoveto{\pgfqpoint{2.926677in}{2.296363in}}%
\pgfpathlineto{\pgfqpoint{2.968946in}{2.309255in}}%
\pgfpathlineto{\pgfqpoint{2.999432in}{2.290834in}}%
\pgfpathlineto{\pgfqpoint{2.957133in}{2.277904in}}%
\pgfpathlineto{\pgfqpoint{2.926677in}{2.296363in}}%
\pgfpathclose%
\pgfusepath{stroke,fill}%
\end{pgfscope}%
\begin{pgfscope}%
\pgfpathrectangle{\pgfqpoint{0.050000in}{0.050000in}}{\pgfqpoint{4.900000in}{4.900000in}}%
\pgfusepath{clip}%
\pgfsetbuttcap%
\pgfsetroundjoin%
\definecolor{currentfill}{rgb}{1.000000,1.000000,1.000000}%
\pgfsetfillcolor{currentfill}%
\pgfsetfillopacity{0.900000}%
\pgfsetlinewidth{0.501875pt}%
\definecolor{currentstroke}{rgb}{0.200000,0.200000,0.200000}%
\pgfsetstrokecolor{currentstroke}%
\pgfsetdash{}{0pt}%
\pgfpathmoveto{\pgfqpoint{2.738890in}{2.294508in}}%
\pgfpathlineto{\pgfqpoint{2.781357in}{2.307404in}}%
\pgfpathlineto{\pgfqpoint{2.811574in}{2.288978in}}%
\pgfpathlineto{\pgfqpoint{2.769076in}{2.276044in}}%
\pgfpathlineto{\pgfqpoint{2.738890in}{2.294508in}}%
\pgfpathclose%
\pgfusepath{stroke,fill}%
\end{pgfscope}%
\begin{pgfscope}%
\pgfpathrectangle{\pgfqpoint{0.050000in}{0.050000in}}{\pgfqpoint{4.900000in}{4.900000in}}%
\pgfusepath{clip}%
\pgfsetbuttcap%
\pgfsetroundjoin%
\definecolor{currentfill}{rgb}{1.000000,1.000000,1.000000}%
\pgfsetfillcolor{currentfill}%
\pgfsetfillopacity{0.900000}%
\pgfsetlinewidth{0.501875pt}%
\definecolor{currentstroke}{rgb}{0.200000,0.200000,0.200000}%
\pgfsetstrokecolor{currentstroke}%
\pgfsetdash{}{0pt}%
\pgfpathmoveto{\pgfqpoint{2.811574in}{2.288978in}}%
\pgfpathlineto{\pgfqpoint{2.853985in}{2.301886in}}%
\pgfpathlineto{\pgfqpoint{2.884322in}{2.283443in}}%
\pgfpathlineto{\pgfqpoint{2.841879in}{2.270498in}}%
\pgfpathlineto{\pgfqpoint{2.811574in}{2.288978in}}%
\pgfpathclose%
\pgfusepath{stroke,fill}%
\end{pgfscope}%
\begin{pgfscope}%
\pgfpathrectangle{\pgfqpoint{0.050000in}{0.050000in}}{\pgfqpoint{4.900000in}{4.900000in}}%
\pgfusepath{clip}%
\pgfsetbuttcap%
\pgfsetroundjoin%
\definecolor{currentfill}{rgb}{1.000000,1.000000,1.000000}%
\pgfsetfillcolor{currentfill}%
\pgfsetfillopacity{0.900000}%
\pgfsetlinewidth{0.501875pt}%
\definecolor{currentstroke}{rgb}{0.200000,0.200000,0.200000}%
\pgfsetstrokecolor{currentstroke}%
\pgfsetdash{}{0pt}%
\pgfpathmoveto{\pgfqpoint{2.623660in}{2.287121in}}%
\pgfpathlineto{\pgfqpoint{2.666270in}{2.300033in}}%
\pgfpathlineto{\pgfqpoint{2.696336in}{2.281585in}}%
\pgfpathlineto{\pgfqpoint{2.653695in}{2.268636in}}%
\pgfpathlineto{\pgfqpoint{2.623660in}{2.287121in}}%
\pgfpathclose%
\pgfusepath{stroke,fill}%
\end{pgfscope}%
\begin{pgfscope}%
\pgfpathrectangle{\pgfqpoint{0.050000in}{0.050000in}}{\pgfqpoint{4.900000in}{4.900000in}}%
\pgfusepath{clip}%
\pgfsetbuttcap%
\pgfsetroundjoin%
\definecolor{currentfill}{rgb}{1.000000,1.000000,1.000000}%
\pgfsetfillcolor{currentfill}%
\pgfsetfillopacity{0.900000}%
\pgfsetlinewidth{0.501875pt}%
\definecolor{currentstroke}{rgb}{0.200000,0.200000,0.200000}%
\pgfsetstrokecolor{currentstroke}%
\pgfsetdash{}{0pt}%
\pgfpathmoveto{\pgfqpoint{2.884322in}{2.283443in}}%
\pgfpathlineto{\pgfqpoint{2.926677in}{2.296363in}}%
\pgfpathlineto{\pgfqpoint{2.957133in}{2.277904in}}%
\pgfpathlineto{\pgfqpoint{2.914747in}{2.264947in}}%
\pgfpathlineto{\pgfqpoint{2.884322in}{2.283443in}}%
\pgfpathclose%
\pgfusepath{stroke,fill}%
\end{pgfscope}%
\begin{pgfscope}%
\pgfpathrectangle{\pgfqpoint{0.050000in}{0.050000in}}{\pgfqpoint{4.900000in}{4.900000in}}%
\pgfusepath{clip}%
\pgfsetbuttcap%
\pgfsetroundjoin%
\definecolor{currentfill}{rgb}{1.000000,1.000000,1.000000}%
\pgfsetfillcolor{currentfill}%
\pgfsetfillopacity{0.900000}%
\pgfsetlinewidth{0.501875pt}%
\definecolor{currentstroke}{rgb}{0.200000,0.200000,0.200000}%
\pgfsetstrokecolor{currentstroke}%
\pgfsetdash{}{0pt}%
\pgfpathmoveto{\pgfqpoint{2.696336in}{2.281585in}}%
\pgfpathlineto{\pgfqpoint{2.738890in}{2.294508in}}%
\pgfpathlineto{\pgfqpoint{2.769076in}{2.276044in}}%
\pgfpathlineto{\pgfqpoint{2.726491in}{2.263084in}}%
\pgfpathlineto{\pgfqpoint{2.696336in}{2.281585in}}%
\pgfpathclose%
\pgfusepath{stroke,fill}%
\end{pgfscope}%
\begin{pgfscope}%
\pgfpathrectangle{\pgfqpoint{0.050000in}{0.050000in}}{\pgfqpoint{4.900000in}{4.900000in}}%
\pgfusepath{clip}%
\pgfsetbuttcap%
\pgfsetroundjoin%
\definecolor{currentfill}{rgb}{1.000000,1.000000,1.000000}%
\pgfsetfillcolor{currentfill}%
\pgfsetfillopacity{0.900000}%
\pgfsetlinewidth{0.501875pt}%
\definecolor{currentstroke}{rgb}{0.200000,0.200000,0.200000}%
\pgfsetstrokecolor{currentstroke}%
\pgfsetdash{}{0pt}%
\pgfpathmoveto{\pgfqpoint{2.957133in}{2.277904in}}%
\pgfpathlineto{\pgfqpoint{2.999432in}{2.290834in}}%
\pgfpathlineto{\pgfqpoint{3.030008in}{2.272360in}}%
\pgfpathlineto{\pgfqpoint{2.987678in}{2.259391in}}%
\pgfpathlineto{\pgfqpoint{2.957133in}{2.277904in}}%
\pgfpathclose%
\pgfusepath{stroke,fill}%
\end{pgfscope}%
\begin{pgfscope}%
\pgfpathrectangle{\pgfqpoint{0.050000in}{0.050000in}}{\pgfqpoint{4.900000in}{4.900000in}}%
\pgfusepath{clip}%
\pgfsetbuttcap%
\pgfsetroundjoin%
\definecolor{currentfill}{rgb}{1.000000,1.000000,1.000000}%
\pgfsetfillcolor{currentfill}%
\pgfsetfillopacity{0.900000}%
\pgfsetlinewidth{0.501875pt}%
\definecolor{currentstroke}{rgb}{0.200000,0.200000,0.200000}%
\pgfsetstrokecolor{currentstroke}%
\pgfsetdash{}{0pt}%
\pgfpathmoveto{\pgfqpoint{2.769076in}{2.276044in}}%
\pgfpathlineto{\pgfqpoint{2.811574in}{2.288978in}}%
\pgfpathlineto{\pgfqpoint{2.841879in}{2.270498in}}%
\pgfpathlineto{\pgfqpoint{2.799350in}{2.257526in}}%
\pgfpathlineto{\pgfqpoint{2.769076in}{2.276044in}}%
\pgfpathclose%
\pgfusepath{stroke,fill}%
\end{pgfscope}%
\begin{pgfscope}%
\pgfpathrectangle{\pgfqpoint{0.050000in}{0.050000in}}{\pgfqpoint{4.900000in}{4.900000in}}%
\pgfusepath{clip}%
\pgfsetbuttcap%
\pgfsetroundjoin%
\definecolor{currentfill}{rgb}{1.000000,1.000000,1.000000}%
\pgfsetfillcolor{currentfill}%
\pgfsetfillopacity{0.900000}%
\pgfsetlinewidth{0.501875pt}%
\definecolor{currentstroke}{rgb}{0.200000,0.200000,0.200000}%
\pgfsetstrokecolor{currentstroke}%
\pgfsetdash{}{0pt}%
\pgfpathmoveto{\pgfqpoint{2.580963in}{2.274183in}}%
\pgfpathlineto{\pgfqpoint{2.623660in}{2.287121in}}%
\pgfpathlineto{\pgfqpoint{2.653695in}{2.268636in}}%
\pgfpathlineto{\pgfqpoint{2.610967in}{2.255665in}}%
\pgfpathlineto{\pgfqpoint{2.580963in}{2.274183in}}%
\pgfpathclose%
\pgfusepath{stroke,fill}%
\end{pgfscope}%
\begin{pgfscope}%
\pgfpathrectangle{\pgfqpoint{0.050000in}{0.050000in}}{\pgfqpoint{4.900000in}{4.900000in}}%
\pgfusepath{clip}%
\pgfsetbuttcap%
\pgfsetroundjoin%
\definecolor{currentfill}{rgb}{1.000000,1.000000,1.000000}%
\pgfsetfillcolor{currentfill}%
\pgfsetfillopacity{0.900000}%
\pgfsetlinewidth{0.501875pt}%
\definecolor{currentstroke}{rgb}{0.200000,0.200000,0.200000}%
\pgfsetstrokecolor{currentstroke}%
\pgfsetdash{}{0pt}%
\pgfpathmoveto{\pgfqpoint{2.841879in}{2.270498in}}%
\pgfpathlineto{\pgfqpoint{2.884322in}{2.283443in}}%
\pgfpathlineto{\pgfqpoint{2.914747in}{2.264947in}}%
\pgfpathlineto{\pgfqpoint{2.872274in}{2.251964in}}%
\pgfpathlineto{\pgfqpoint{2.841879in}{2.270498in}}%
\pgfpathclose%
\pgfusepath{stroke,fill}%
\end{pgfscope}%
\begin{pgfscope}%
\pgfpathrectangle{\pgfqpoint{0.050000in}{0.050000in}}{\pgfqpoint{4.900000in}{4.900000in}}%
\pgfusepath{clip}%
\pgfsetbuttcap%
\pgfsetroundjoin%
\definecolor{currentfill}{rgb}{1.000000,1.000000,1.000000}%
\pgfsetfillcolor{currentfill}%
\pgfsetfillopacity{0.900000}%
\pgfsetlinewidth{0.501875pt}%
\definecolor{currentstroke}{rgb}{0.200000,0.200000,0.200000}%
\pgfsetstrokecolor{currentstroke}%
\pgfsetdash{}{0pt}%
\pgfpathmoveto{\pgfqpoint{2.653695in}{2.268636in}}%
\pgfpathlineto{\pgfqpoint{2.696336in}{2.281585in}}%
\pgfpathlineto{\pgfqpoint{2.726491in}{2.263084in}}%
\pgfpathlineto{\pgfqpoint{2.683818in}{2.250101in}}%
\pgfpathlineto{\pgfqpoint{2.653695in}{2.268636in}}%
\pgfpathclose%
\pgfusepath{stroke,fill}%
\end{pgfscope}%
\begin{pgfscope}%
\pgfpathrectangle{\pgfqpoint{0.050000in}{0.050000in}}{\pgfqpoint{4.900000in}{4.900000in}}%
\pgfusepath{clip}%
\pgfsetbuttcap%
\pgfsetroundjoin%
\definecolor{currentfill}{rgb}{1.000000,1.000000,1.000000}%
\pgfsetfillcolor{currentfill}%
\pgfsetfillopacity{0.900000}%
\pgfsetlinewidth{0.501875pt}%
\definecolor{currentstroke}{rgb}{0.200000,0.200000,0.200000}%
\pgfsetstrokecolor{currentstroke}%
\pgfsetdash{}{0pt}%
\pgfpathmoveto{\pgfqpoint{2.914747in}{2.264947in}}%
\pgfpathlineto{\pgfqpoint{2.957133in}{2.277904in}}%
\pgfpathlineto{\pgfqpoint{2.987678in}{2.259391in}}%
\pgfpathlineto{\pgfqpoint{2.945261in}{2.246397in}}%
\pgfpathlineto{\pgfqpoint{2.914747in}{2.264947in}}%
\pgfpathclose%
\pgfusepath{stroke,fill}%
\end{pgfscope}%
\begin{pgfscope}%
\pgfpathrectangle{\pgfqpoint{0.050000in}{0.050000in}}{\pgfqpoint{4.900000in}{4.900000in}}%
\pgfusepath{clip}%
\pgfsetbuttcap%
\pgfsetroundjoin%
\definecolor{currentfill}{rgb}{1.000000,1.000000,1.000000}%
\pgfsetfillcolor{currentfill}%
\pgfsetfillopacity{0.900000}%
\pgfsetlinewidth{0.501875pt}%
\definecolor{currentstroke}{rgb}{0.200000,0.200000,0.200000}%
\pgfsetstrokecolor{currentstroke}%
\pgfsetdash{}{0pt}%
\pgfpathmoveto{\pgfqpoint{2.726491in}{2.263084in}}%
\pgfpathlineto{\pgfqpoint{2.769076in}{2.276044in}}%
\pgfpathlineto{\pgfqpoint{2.799350in}{2.257526in}}%
\pgfpathlineto{\pgfqpoint{2.756734in}{2.244531in}}%
\pgfpathlineto{\pgfqpoint{2.726491in}{2.263084in}}%
\pgfpathclose%
\pgfusepath{stroke,fill}%
\end{pgfscope}%
\begin{pgfscope}%
\pgfpathrectangle{\pgfqpoint{0.050000in}{0.050000in}}{\pgfqpoint{4.900000in}{4.900000in}}%
\pgfusepath{clip}%
\pgfsetbuttcap%
\pgfsetroundjoin%
\definecolor{currentfill}{rgb}{1.000000,1.000000,1.000000}%
\pgfsetfillcolor{currentfill}%
\pgfsetfillopacity{0.900000}%
\pgfsetlinewidth{0.501875pt}%
\definecolor{currentstroke}{rgb}{0.200000,0.200000,0.200000}%
\pgfsetstrokecolor{currentstroke}%
\pgfsetdash{}{0pt}%
\pgfpathmoveto{\pgfqpoint{2.538179in}{2.261223in}}%
\pgfpathlineto{\pgfqpoint{2.580963in}{2.274183in}}%
\pgfpathlineto{\pgfqpoint{2.610967in}{2.255665in}}%
\pgfpathlineto{\pgfqpoint{2.568150in}{2.242695in}}%
\pgfpathlineto{\pgfqpoint{2.538179in}{2.261223in}}%
\pgfpathclose%
\pgfusepath{stroke,fill}%
\end{pgfscope}%
\begin{pgfscope}%
\pgfpathrectangle{\pgfqpoint{0.050000in}{0.050000in}}{\pgfqpoint{4.900000in}{4.900000in}}%
\pgfusepath{clip}%
\pgfsetbuttcap%
\pgfsetroundjoin%
\definecolor{currentfill}{rgb}{1.000000,1.000000,1.000000}%
\pgfsetfillcolor{currentfill}%
\pgfsetfillopacity{0.900000}%
\pgfsetlinewidth{0.501875pt}%
\definecolor{currentstroke}{rgb}{0.200000,0.200000,0.200000}%
\pgfsetstrokecolor{currentstroke}%
\pgfsetdash{}{0pt}%
\pgfpathmoveto{\pgfqpoint{2.987678in}{2.259391in}}%
\pgfpathlineto{\pgfqpoint{3.030008in}{2.272360in}}%
\pgfpathlineto{\pgfqpoint{3.060674in}{2.253831in}}%
\pgfpathlineto{\pgfqpoint{3.018313in}{2.240825in}}%
\pgfpathlineto{\pgfqpoint{2.987678in}{2.259391in}}%
\pgfpathclose%
\pgfusepath{stroke,fill}%
\end{pgfscope}%
\begin{pgfscope}%
\pgfpathrectangle{\pgfqpoint{0.050000in}{0.050000in}}{\pgfqpoint{4.900000in}{4.900000in}}%
\pgfusepath{clip}%
\pgfsetbuttcap%
\pgfsetroundjoin%
\definecolor{currentfill}{rgb}{1.000000,1.000000,1.000000}%
\pgfsetfillcolor{currentfill}%
\pgfsetfillopacity{0.900000}%
\pgfsetlinewidth{0.501875pt}%
\definecolor{currentstroke}{rgb}{0.200000,0.200000,0.200000}%
\pgfsetstrokecolor{currentstroke}%
\pgfsetdash{}{0pt}%
\pgfpathmoveto{\pgfqpoint{2.799350in}{2.257526in}}%
\pgfpathlineto{\pgfqpoint{2.841879in}{2.270498in}}%
\pgfpathlineto{\pgfqpoint{2.872274in}{2.251964in}}%
\pgfpathlineto{\pgfqpoint{2.829714in}{2.238957in}}%
\pgfpathlineto{\pgfqpoint{2.799350in}{2.257526in}}%
\pgfpathclose%
\pgfusepath{stroke,fill}%
\end{pgfscope}%
\begin{pgfscope}%
\pgfpathrectangle{\pgfqpoint{0.050000in}{0.050000in}}{\pgfqpoint{4.900000in}{4.900000in}}%
\pgfusepath{clip}%
\pgfsetbuttcap%
\pgfsetroundjoin%
\definecolor{currentfill}{rgb}{1.000000,1.000000,1.000000}%
\pgfsetfillcolor{currentfill}%
\pgfsetfillopacity{0.900000}%
\pgfsetlinewidth{0.501875pt}%
\definecolor{currentstroke}{rgb}{0.200000,0.200000,0.200000}%
\pgfsetstrokecolor{currentstroke}%
\pgfsetdash{}{0pt}%
\pgfpathmoveto{\pgfqpoint{2.610967in}{2.255665in}}%
\pgfpathlineto{\pgfqpoint{2.653695in}{2.268636in}}%
\pgfpathlineto{\pgfqpoint{2.683818in}{2.250101in}}%
\pgfpathlineto{\pgfqpoint{2.641058in}{2.237111in}}%
\pgfpathlineto{\pgfqpoint{2.610967in}{2.255665in}}%
\pgfpathclose%
\pgfusepath{stroke,fill}%
\end{pgfscope}%
\begin{pgfscope}%
\pgfpathrectangle{\pgfqpoint{0.050000in}{0.050000in}}{\pgfqpoint{4.900000in}{4.900000in}}%
\pgfusepath{clip}%
\pgfsetbuttcap%
\pgfsetroundjoin%
\definecolor{currentfill}{rgb}{1.000000,1.000000,1.000000}%
\pgfsetfillcolor{currentfill}%
\pgfsetfillopacity{0.900000}%
\pgfsetlinewidth{0.501875pt}%
\definecolor{currentstroke}{rgb}{0.200000,0.200000,0.200000}%
\pgfsetstrokecolor{currentstroke}%
\pgfsetdash{}{0pt}%
\pgfpathmoveto{\pgfqpoint{2.872274in}{2.251964in}}%
\pgfpathlineto{\pgfqpoint{2.914747in}{2.264947in}}%
\pgfpathlineto{\pgfqpoint{2.945261in}{2.246397in}}%
\pgfpathlineto{\pgfqpoint{2.902757in}{2.233377in}}%
\pgfpathlineto{\pgfqpoint{2.872274in}{2.251964in}}%
\pgfpathclose%
\pgfusepath{stroke,fill}%
\end{pgfscope}%
\begin{pgfscope}%
\pgfpathrectangle{\pgfqpoint{0.050000in}{0.050000in}}{\pgfqpoint{4.900000in}{4.900000in}}%
\pgfusepath{clip}%
\pgfsetbuttcap%
\pgfsetroundjoin%
\definecolor{currentfill}{rgb}{1.000000,1.000000,1.000000}%
\pgfsetfillcolor{currentfill}%
\pgfsetfillopacity{0.900000}%
\pgfsetlinewidth{0.501875pt}%
\definecolor{currentstroke}{rgb}{0.200000,0.200000,0.200000}%
\pgfsetstrokecolor{currentstroke}%
\pgfsetdash{}{0pt}%
\pgfpathmoveto{\pgfqpoint{2.683818in}{2.250101in}}%
\pgfpathlineto{\pgfqpoint{2.726491in}{2.263084in}}%
\pgfpathlineto{\pgfqpoint{2.756734in}{2.244531in}}%
\pgfpathlineto{\pgfqpoint{2.714030in}{2.231523in}}%
\pgfpathlineto{\pgfqpoint{2.683818in}{2.250101in}}%
\pgfpathclose%
\pgfusepath{stroke,fill}%
\end{pgfscope}%
\begin{pgfscope}%
\pgfpathrectangle{\pgfqpoint{0.050000in}{0.050000in}}{\pgfqpoint{4.900000in}{4.900000in}}%
\pgfusepath{clip}%
\pgfsetbuttcap%
\pgfsetroundjoin%
\definecolor{currentfill}{rgb}{1.000000,1.000000,1.000000}%
\pgfsetfillcolor{currentfill}%
\pgfsetfillopacity{0.900000}%
\pgfsetlinewidth{0.501875pt}%
\definecolor{currentstroke}{rgb}{0.200000,0.200000,0.200000}%
\pgfsetstrokecolor{currentstroke}%
\pgfsetdash{}{0pt}%
\pgfpathmoveto{\pgfqpoint{2.495307in}{2.248270in}}%
\pgfpathlineto{\pgfqpoint{2.538179in}{2.261223in}}%
\pgfpathlineto{\pgfqpoint{2.568150in}{2.242695in}}%
\pgfpathlineto{\pgfqpoint{2.525246in}{2.229820in}}%
\pgfpathlineto{\pgfqpoint{2.495307in}{2.248270in}}%
\pgfpathclose%
\pgfusepath{stroke,fill}%
\end{pgfscope}%
\begin{pgfscope}%
\pgfpathrectangle{\pgfqpoint{0.050000in}{0.050000in}}{\pgfqpoint{4.900000in}{4.900000in}}%
\pgfusepath{clip}%
\pgfsetbuttcap%
\pgfsetroundjoin%
\definecolor{currentfill}{rgb}{1.000000,1.000000,1.000000}%
\pgfsetfillcolor{currentfill}%
\pgfsetfillopacity{0.900000}%
\pgfsetlinewidth{0.501875pt}%
\definecolor{currentstroke}{rgb}{0.200000,0.200000,0.200000}%
\pgfsetstrokecolor{currentstroke}%
\pgfsetdash{}{0pt}%
\pgfpathmoveto{\pgfqpoint{2.945261in}{2.246397in}}%
\pgfpathlineto{\pgfqpoint{2.987678in}{2.259391in}}%
\pgfpathlineto{\pgfqpoint{3.018313in}{2.240825in}}%
\pgfpathlineto{\pgfqpoint{2.975866in}{2.227793in}}%
\pgfpathlineto{\pgfqpoint{2.945261in}{2.246397in}}%
\pgfpathclose%
\pgfusepath{stroke,fill}%
\end{pgfscope}%
\begin{pgfscope}%
\pgfpathrectangle{\pgfqpoint{0.050000in}{0.050000in}}{\pgfqpoint{4.900000in}{4.900000in}}%
\pgfusepath{clip}%
\pgfsetbuttcap%
\pgfsetroundjoin%
\definecolor{currentfill}{rgb}{1.000000,1.000000,1.000000}%
\pgfsetfillcolor{currentfill}%
\pgfsetfillopacity{0.900000}%
\pgfsetlinewidth{0.501875pt}%
\definecolor{currentstroke}{rgb}{0.200000,0.200000,0.200000}%
\pgfsetstrokecolor{currentstroke}%
\pgfsetdash{}{0pt}%
\pgfpathmoveto{\pgfqpoint{2.756734in}{2.244531in}}%
\pgfpathlineto{\pgfqpoint{2.799350in}{2.257526in}}%
\pgfpathlineto{\pgfqpoint{2.829714in}{2.238957in}}%
\pgfpathlineto{\pgfqpoint{2.787066in}{2.225932in}}%
\pgfpathlineto{\pgfqpoint{2.756734in}{2.244531in}}%
\pgfpathclose%
\pgfusepath{stroke,fill}%
\end{pgfscope}%
\begin{pgfscope}%
\pgfpathrectangle{\pgfqpoint{0.050000in}{0.050000in}}{\pgfqpoint{4.900000in}{4.900000in}}%
\pgfusepath{clip}%
\pgfsetbuttcap%
\pgfsetroundjoin%
\definecolor{currentfill}{rgb}{1.000000,1.000000,1.000000}%
\pgfsetfillcolor{currentfill}%
\pgfsetfillopacity{0.900000}%
\pgfsetlinewidth{0.501875pt}%
\definecolor{currentstroke}{rgb}{0.200000,0.200000,0.200000}%
\pgfsetstrokecolor{currentstroke}%
\pgfsetdash{}{0pt}%
\pgfpathmoveto{\pgfqpoint{2.568150in}{2.242695in}}%
\pgfpathlineto{\pgfqpoint{2.610967in}{2.255665in}}%
\pgfpathlineto{\pgfqpoint{2.641058in}{2.237111in}}%
\pgfpathlineto{\pgfqpoint{2.598210in}{2.224170in}}%
\pgfpathlineto{\pgfqpoint{2.568150in}{2.242695in}}%
\pgfpathclose%
\pgfusepath{stroke,fill}%
\end{pgfscope}%
\begin{pgfscope}%
\pgfpathrectangle{\pgfqpoint{0.050000in}{0.050000in}}{\pgfqpoint{4.900000in}{4.900000in}}%
\pgfusepath{clip}%
\pgfsetbuttcap%
\pgfsetroundjoin%
\definecolor{currentfill}{rgb}{1.000000,1.000000,1.000000}%
\pgfsetfillcolor{currentfill}%
\pgfsetfillopacity{0.900000}%
\pgfsetlinewidth{0.501875pt}%
\definecolor{currentstroke}{rgb}{0.200000,0.200000,0.200000}%
\pgfsetstrokecolor{currentstroke}%
\pgfsetdash{}{0pt}%
\pgfpathmoveto{\pgfqpoint{3.018313in}{2.240825in}}%
\pgfpathlineto{\pgfqpoint{3.060674in}{2.253831in}}%
\pgfpathlineto{\pgfqpoint{3.091429in}{2.235247in}}%
\pgfpathlineto{\pgfqpoint{3.049038in}{2.222204in}}%
\pgfpathlineto{\pgfqpoint{3.018313in}{2.240825in}}%
\pgfpathclose%
\pgfusepath{stroke,fill}%
\end{pgfscope}%
\begin{pgfscope}%
\pgfpathrectangle{\pgfqpoint{0.050000in}{0.050000in}}{\pgfqpoint{4.900000in}{4.900000in}}%
\pgfusepath{clip}%
\pgfsetbuttcap%
\pgfsetroundjoin%
\definecolor{currentfill}{rgb}{1.000000,1.000000,1.000000}%
\pgfsetfillcolor{currentfill}%
\pgfsetfillopacity{0.900000}%
\pgfsetlinewidth{0.501875pt}%
\definecolor{currentstroke}{rgb}{0.200000,0.200000,0.200000}%
\pgfsetstrokecolor{currentstroke}%
\pgfsetdash{}{0pt}%
\pgfpathmoveto{\pgfqpoint{2.829714in}{2.238957in}}%
\pgfpathlineto{\pgfqpoint{2.872274in}{2.251964in}}%
\pgfpathlineto{\pgfqpoint{2.902757in}{2.233377in}}%
\pgfpathlineto{\pgfqpoint{2.860166in}{2.220337in}}%
\pgfpathlineto{\pgfqpoint{2.829714in}{2.238957in}}%
\pgfpathclose%
\pgfusepath{stroke,fill}%
\end{pgfscope}%
\begin{pgfscope}%
\pgfpathrectangle{\pgfqpoint{0.050000in}{0.050000in}}{\pgfqpoint{4.900000in}{4.900000in}}%
\pgfusepath{clip}%
\pgfsetbuttcap%
\pgfsetroundjoin%
\definecolor{currentfill}{rgb}{1.000000,1.000000,1.000000}%
\pgfsetfillcolor{currentfill}%
\pgfsetfillopacity{0.900000}%
\pgfsetlinewidth{0.501875pt}%
\definecolor{currentstroke}{rgb}{0.200000,0.200000,0.200000}%
\pgfsetstrokecolor{currentstroke}%
\pgfsetdash{}{0pt}%
\pgfpathmoveto{\pgfqpoint{2.641058in}{2.237111in}}%
\pgfpathlineto{\pgfqpoint{2.683818in}{2.250101in}}%
\pgfpathlineto{\pgfqpoint{2.714030in}{2.231523in}}%
\pgfpathlineto{\pgfqpoint{2.671238in}{2.218535in}}%
\pgfpathlineto{\pgfqpoint{2.641058in}{2.237111in}}%
\pgfpathclose%
\pgfusepath{stroke,fill}%
\end{pgfscope}%
\begin{pgfscope}%
\pgfpathrectangle{\pgfqpoint{0.050000in}{0.050000in}}{\pgfqpoint{4.900000in}{4.900000in}}%
\pgfusepath{clip}%
\pgfsetbuttcap%
\pgfsetroundjoin%
\definecolor{currentfill}{rgb}{1.000000,1.000000,1.000000}%
\pgfsetfillcolor{currentfill}%
\pgfsetfillopacity{0.900000}%
\pgfsetlinewidth{0.501875pt}%
\definecolor{currentstroke}{rgb}{0.200000,0.200000,0.200000}%
\pgfsetstrokecolor{currentstroke}%
\pgfsetdash{}{0pt}%
\pgfpathmoveto{\pgfqpoint{2.452346in}{2.235485in}}%
\pgfpathlineto{\pgfqpoint{2.495307in}{2.248270in}}%
\pgfpathlineto{\pgfqpoint{2.525246in}{2.229820in}}%
\pgfpathlineto{\pgfqpoint{2.482252in}{2.217381in}}%
\pgfpathlineto{\pgfqpoint{2.452346in}{2.235485in}}%
\pgfpathclose%
\pgfusepath{stroke,fill}%
\end{pgfscope}%
\begin{pgfscope}%
\pgfpathrectangle{\pgfqpoint{0.050000in}{0.050000in}}{\pgfqpoint{4.900000in}{4.900000in}}%
\pgfusepath{clip}%
\pgfsetbuttcap%
\pgfsetroundjoin%
\definecolor{currentfill}{rgb}{1.000000,1.000000,1.000000}%
\pgfsetfillcolor{currentfill}%
\pgfsetfillopacity{0.900000}%
\pgfsetlinewidth{0.501875pt}%
\definecolor{currentstroke}{rgb}{0.200000,0.200000,0.200000}%
\pgfsetstrokecolor{currentstroke}%
\pgfsetdash{}{0pt}%
\pgfpathmoveto{\pgfqpoint{2.902757in}{2.233377in}}%
\pgfpathlineto{\pgfqpoint{2.945261in}{2.246397in}}%
\pgfpathlineto{\pgfqpoint{2.975866in}{2.227793in}}%
\pgfpathlineto{\pgfqpoint{2.933331in}{2.214739in}}%
\pgfpathlineto{\pgfqpoint{2.902757in}{2.233377in}}%
\pgfpathclose%
\pgfusepath{stroke,fill}%
\end{pgfscope}%
\begin{pgfscope}%
\pgfpathrectangle{\pgfqpoint{0.050000in}{0.050000in}}{\pgfqpoint{4.900000in}{4.900000in}}%
\pgfusepath{clip}%
\pgfsetbuttcap%
\pgfsetroundjoin%
\definecolor{currentfill}{rgb}{1.000000,1.000000,1.000000}%
\pgfsetfillcolor{currentfill}%
\pgfsetfillopacity{0.900000}%
\pgfsetlinewidth{0.501875pt}%
\definecolor{currentstroke}{rgb}{0.200000,0.200000,0.200000}%
\pgfsetstrokecolor{currentstroke}%
\pgfsetdash{}{0pt}%
\pgfpathmoveto{\pgfqpoint{2.714030in}{2.231523in}}%
\pgfpathlineto{\pgfqpoint{2.756734in}{2.244531in}}%
\pgfpathlineto{\pgfqpoint{2.787066in}{2.225932in}}%
\pgfpathlineto{\pgfqpoint{2.744330in}{2.212909in}}%
\pgfpathlineto{\pgfqpoint{2.714030in}{2.231523in}}%
\pgfpathclose%
\pgfusepath{stroke,fill}%
\end{pgfscope}%
\begin{pgfscope}%
\pgfpathrectangle{\pgfqpoint{0.050000in}{0.050000in}}{\pgfqpoint{4.900000in}{4.900000in}}%
\pgfusepath{clip}%
\pgfsetbuttcap%
\pgfsetroundjoin%
\definecolor{currentfill}{rgb}{1.000000,1.000000,1.000000}%
\pgfsetfillcolor{currentfill}%
\pgfsetfillopacity{0.900000}%
\pgfsetlinewidth{0.501875pt}%
\definecolor{currentstroke}{rgb}{0.200000,0.200000,0.200000}%
\pgfsetstrokecolor{currentstroke}%
\pgfsetdash{}{0pt}%
\pgfpathmoveto{\pgfqpoint{2.525246in}{2.229820in}}%
\pgfpathlineto{\pgfqpoint{2.568150in}{2.242695in}}%
\pgfpathlineto{\pgfqpoint{2.598210in}{2.224170in}}%
\pgfpathlineto{\pgfqpoint{2.555273in}{2.211463in}}%
\pgfpathlineto{\pgfqpoint{2.525246in}{2.229820in}}%
\pgfpathclose%
\pgfusepath{stroke,fill}%
\end{pgfscope}%
\begin{pgfscope}%
\pgfpathrectangle{\pgfqpoint{0.050000in}{0.050000in}}{\pgfqpoint{4.900000in}{4.900000in}}%
\pgfusepath{clip}%
\pgfsetbuttcap%
\pgfsetroundjoin%
\definecolor{currentfill}{rgb}{1.000000,1.000000,1.000000}%
\pgfsetfillcolor{currentfill}%
\pgfsetfillopacity{0.900000}%
\pgfsetlinewidth{0.501875pt}%
\definecolor{currentstroke}{rgb}{0.200000,0.200000,0.200000}%
\pgfsetstrokecolor{currentstroke}%
\pgfsetdash{}{0pt}%
\pgfpathmoveto{\pgfqpoint{2.975866in}{2.227793in}}%
\pgfpathlineto{\pgfqpoint{3.018313in}{2.240825in}}%
\pgfpathlineto{\pgfqpoint{3.049038in}{2.222204in}}%
\pgfpathlineto{\pgfqpoint{3.006560in}{2.209137in}}%
\pgfpathlineto{\pgfqpoint{2.975866in}{2.227793in}}%
\pgfpathclose%
\pgfusepath{stroke,fill}%
\end{pgfscope}%
\begin{pgfscope}%
\pgfpathrectangle{\pgfqpoint{0.050000in}{0.050000in}}{\pgfqpoint{4.900000in}{4.900000in}}%
\pgfusepath{clip}%
\pgfsetbuttcap%
\pgfsetroundjoin%
\definecolor{currentfill}{rgb}{1.000000,1.000000,1.000000}%
\pgfsetfillcolor{currentfill}%
\pgfsetfillopacity{0.900000}%
\pgfsetlinewidth{0.501875pt}%
\definecolor{currentstroke}{rgb}{0.200000,0.200000,0.200000}%
\pgfsetstrokecolor{currentstroke}%
\pgfsetdash{}{0pt}%
\pgfpathmoveto{\pgfqpoint{2.787066in}{2.225932in}}%
\pgfpathlineto{\pgfqpoint{2.829714in}{2.238957in}}%
\pgfpathlineto{\pgfqpoint{2.860166in}{2.220337in}}%
\pgfpathlineto{\pgfqpoint{2.817487in}{2.207288in}}%
\pgfpathlineto{\pgfqpoint{2.787066in}{2.225932in}}%
\pgfpathclose%
\pgfusepath{stroke,fill}%
\end{pgfscope}%
\begin{pgfscope}%
\pgfpathrectangle{\pgfqpoint{0.050000in}{0.050000in}}{\pgfqpoint{4.900000in}{4.900000in}}%
\pgfusepath{clip}%
\pgfsetbuttcap%
\pgfsetroundjoin%
\definecolor{currentfill}{rgb}{1.000000,1.000000,1.000000}%
\pgfsetfillcolor{currentfill}%
\pgfsetfillopacity{0.900000}%
\pgfsetlinewidth{0.501875pt}%
\definecolor{currentstroke}{rgb}{0.200000,0.200000,0.200000}%
\pgfsetstrokecolor{currentstroke}%
\pgfsetdash{}{0pt}%
\pgfpathmoveto{\pgfqpoint{2.598210in}{2.224170in}}%
\pgfpathlineto{\pgfqpoint{2.641058in}{2.237111in}}%
\pgfpathlineto{\pgfqpoint{2.671238in}{2.218535in}}%
\pgfpathlineto{\pgfqpoint{2.628358in}{2.205671in}}%
\pgfpathlineto{\pgfqpoint{2.598210in}{2.224170in}}%
\pgfpathclose%
\pgfusepath{stroke,fill}%
\end{pgfscope}%
\begin{pgfscope}%
\pgfpathrectangle{\pgfqpoint{0.050000in}{0.050000in}}{\pgfqpoint{4.900000in}{4.900000in}}%
\pgfusepath{clip}%
\pgfsetbuttcap%
\pgfsetroundjoin%
\definecolor{currentfill}{rgb}{1.000000,1.000000,1.000000}%
\pgfsetfillcolor{currentfill}%
\pgfsetfillopacity{0.900000}%
\pgfsetlinewidth{0.501875pt}%
\definecolor{currentstroke}{rgb}{0.200000,0.200000,0.200000}%
\pgfsetstrokecolor{currentstroke}%
\pgfsetdash{}{0pt}%
\pgfpathmoveto{\pgfqpoint{3.049038in}{2.222204in}}%
\pgfpathlineto{\pgfqpoint{3.091429in}{2.235247in}}%
\pgfpathlineto{\pgfqpoint{3.122276in}{2.216609in}}%
\pgfpathlineto{\pgfqpoint{3.079854in}{2.203530in}}%
\pgfpathlineto{\pgfqpoint{3.049038in}{2.222204in}}%
\pgfpathclose%
\pgfusepath{stroke,fill}%
\end{pgfscope}%
\begin{pgfscope}%
\pgfpathrectangle{\pgfqpoint{0.050000in}{0.050000in}}{\pgfqpoint{4.900000in}{4.900000in}}%
\pgfusepath{clip}%
\pgfsetbuttcap%
\pgfsetroundjoin%
\definecolor{currentfill}{rgb}{1.000000,1.000000,1.000000}%
\pgfsetfillcolor{currentfill}%
\pgfsetfillopacity{0.900000}%
\pgfsetlinewidth{0.501875pt}%
\definecolor{currentstroke}{rgb}{0.200000,0.200000,0.200000}%
\pgfsetstrokecolor{currentstroke}%
\pgfsetdash{}{0pt}%
\pgfpathmoveto{\pgfqpoint{2.409294in}{2.223517in}}%
\pgfpathlineto{\pgfqpoint{2.452346in}{2.235485in}}%
\pgfpathlineto{\pgfqpoint{2.482252in}{2.217381in}}%
\pgfpathlineto{\pgfqpoint{2.439166in}{2.206313in}}%
\pgfpathlineto{\pgfqpoint{2.409294in}{2.223517in}}%
\pgfpathclose%
\pgfusepath{stroke,fill}%
\end{pgfscope}%
\begin{pgfscope}%
\pgfpathrectangle{\pgfqpoint{0.050000in}{0.050000in}}{\pgfqpoint{4.900000in}{4.900000in}}%
\pgfusepath{clip}%
\pgfsetbuttcap%
\pgfsetroundjoin%
\definecolor{currentfill}{rgb}{1.000000,1.000000,1.000000}%
\pgfsetfillcolor{currentfill}%
\pgfsetfillopacity{0.900000}%
\pgfsetlinewidth{0.501875pt}%
\definecolor{currentstroke}{rgb}{0.200000,0.200000,0.200000}%
\pgfsetstrokecolor{currentstroke}%
\pgfsetdash{}{0pt}%
\pgfpathmoveto{\pgfqpoint{2.860166in}{2.220337in}}%
\pgfpathlineto{\pgfqpoint{2.902757in}{2.233377in}}%
\pgfpathlineto{\pgfqpoint{2.933331in}{2.214739in}}%
\pgfpathlineto{\pgfqpoint{2.890708in}{2.201669in}}%
\pgfpathlineto{\pgfqpoint{2.860166in}{2.220337in}}%
\pgfpathclose%
\pgfusepath{stroke,fill}%
\end{pgfscope}%
\begin{pgfscope}%
\pgfpathrectangle{\pgfqpoint{0.050000in}{0.050000in}}{\pgfqpoint{4.900000in}{4.900000in}}%
\pgfusepath{clip}%
\pgfsetbuttcap%
\pgfsetroundjoin%
\definecolor{currentfill}{rgb}{1.000000,1.000000,1.000000}%
\pgfsetfillcolor{currentfill}%
\pgfsetfillopacity{0.900000}%
\pgfsetlinewidth{0.501875pt}%
\definecolor{currentstroke}{rgb}{0.200000,0.200000,0.200000}%
\pgfsetstrokecolor{currentstroke}%
\pgfsetdash{}{0pt}%
\pgfpathmoveto{\pgfqpoint{2.671238in}{2.218535in}}%
\pgfpathlineto{\pgfqpoint{2.714030in}{2.231523in}}%
\pgfpathlineto{\pgfqpoint{2.744330in}{2.212909in}}%
\pgfpathlineto{\pgfqpoint{2.701507in}{2.199949in}}%
\pgfpathlineto{\pgfqpoint{2.671238in}{2.218535in}}%
\pgfpathclose%
\pgfusepath{stroke,fill}%
\end{pgfscope}%
\begin{pgfscope}%
\pgfpathrectangle{\pgfqpoint{0.050000in}{0.050000in}}{\pgfqpoint{4.900000in}{4.900000in}}%
\pgfusepath{clip}%
\pgfsetbuttcap%
\pgfsetroundjoin%
\definecolor{currentfill}{rgb}{1.000000,1.000000,1.000000}%
\pgfsetfillcolor{currentfill}%
\pgfsetfillopacity{0.900000}%
\pgfsetlinewidth{0.501875pt}%
\definecolor{currentstroke}{rgb}{0.200000,0.200000,0.200000}%
\pgfsetstrokecolor{currentstroke}%
\pgfsetdash{}{0pt}%
\pgfpathmoveto{\pgfqpoint{2.482252in}{2.217381in}}%
\pgfpathlineto{\pgfqpoint{2.525246in}{2.229820in}}%
\pgfpathlineto{\pgfqpoint{2.555273in}{2.211463in}}%
\pgfpathlineto{\pgfqpoint{2.512247in}{2.199487in}}%
\pgfpathlineto{\pgfqpoint{2.482252in}{2.217381in}}%
\pgfpathclose%
\pgfusepath{stroke,fill}%
\end{pgfscope}%
\begin{pgfscope}%
\pgfpathrectangle{\pgfqpoint{0.050000in}{0.050000in}}{\pgfqpoint{4.900000in}{4.900000in}}%
\pgfusepath{clip}%
\pgfsetbuttcap%
\pgfsetroundjoin%
\definecolor{currentfill}{rgb}{1.000000,1.000000,1.000000}%
\pgfsetfillcolor{currentfill}%
\pgfsetfillopacity{0.900000}%
\pgfsetlinewidth{0.501875pt}%
\definecolor{currentstroke}{rgb}{0.200000,0.200000,0.200000}%
\pgfsetstrokecolor{currentstroke}%
\pgfsetdash{}{0pt}%
\pgfpathmoveto{\pgfqpoint{2.933331in}{2.214739in}}%
\pgfpathlineto{\pgfqpoint{2.975866in}{2.227793in}}%
\pgfpathlineto{\pgfqpoint{3.006560in}{2.209137in}}%
\pgfpathlineto{\pgfqpoint{2.963994in}{2.196049in}}%
\pgfpathlineto{\pgfqpoint{2.933331in}{2.214739in}}%
\pgfpathclose%
\pgfusepath{stroke,fill}%
\end{pgfscope}%
\begin{pgfscope}%
\pgfpathrectangle{\pgfqpoint{0.050000in}{0.050000in}}{\pgfqpoint{4.900000in}{4.900000in}}%
\pgfusepath{clip}%
\pgfsetbuttcap%
\pgfsetroundjoin%
\definecolor{currentfill}{rgb}{1.000000,1.000000,1.000000}%
\pgfsetfillcolor{currentfill}%
\pgfsetfillopacity{0.900000}%
\pgfsetlinewidth{0.501875pt}%
\definecolor{currentstroke}{rgb}{0.200000,0.200000,0.200000}%
\pgfsetstrokecolor{currentstroke}%
\pgfsetdash{}{0pt}%
\pgfpathmoveto{\pgfqpoint{2.744330in}{2.212909in}}%
\pgfpathlineto{\pgfqpoint{2.787066in}{2.225932in}}%
\pgfpathlineto{\pgfqpoint{2.817487in}{2.207288in}}%
\pgfpathlineto{\pgfqpoint{2.774720in}{2.194266in}}%
\pgfpathlineto{\pgfqpoint{2.744330in}{2.212909in}}%
\pgfpathclose%
\pgfusepath{stroke,fill}%
\end{pgfscope}%
\begin{pgfscope}%
\pgfpathrectangle{\pgfqpoint{0.050000in}{0.050000in}}{\pgfqpoint{4.900000in}{4.900000in}}%
\pgfusepath{clip}%
\pgfsetbuttcap%
\pgfsetroundjoin%
\definecolor{currentfill}{rgb}{1.000000,1.000000,1.000000}%
\pgfsetfillcolor{currentfill}%
\pgfsetfillopacity{0.900000}%
\pgfsetlinewidth{0.501875pt}%
\definecolor{currentstroke}{rgb}{0.200000,0.200000,0.200000}%
\pgfsetstrokecolor{currentstroke}%
\pgfsetdash{}{0pt}%
\pgfpathmoveto{\pgfqpoint{2.555273in}{2.211463in}}%
\pgfpathlineto{\pgfqpoint{2.598210in}{2.224170in}}%
\pgfpathlineto{\pgfqpoint{2.628358in}{2.205671in}}%
\pgfpathlineto{\pgfqpoint{2.585389in}{2.193204in}}%
\pgfpathlineto{\pgfqpoint{2.555273in}{2.211463in}}%
\pgfpathclose%
\pgfusepath{stroke,fill}%
\end{pgfscope}%
\begin{pgfscope}%
\pgfpathrectangle{\pgfqpoint{0.050000in}{0.050000in}}{\pgfqpoint{4.900000in}{4.900000in}}%
\pgfusepath{clip}%
\pgfsetbuttcap%
\pgfsetroundjoin%
\definecolor{currentfill}{rgb}{1.000000,1.000000,1.000000}%
\pgfsetfillcolor{currentfill}%
\pgfsetfillopacity{0.900000}%
\pgfsetlinewidth{0.501875pt}%
\definecolor{currentstroke}{rgb}{0.200000,0.200000,0.200000}%
\pgfsetstrokecolor{currentstroke}%
\pgfsetdash{}{0pt}%
\pgfpathmoveto{\pgfqpoint{3.006560in}{2.209137in}}%
\pgfpathlineto{\pgfqpoint{3.049038in}{2.222204in}}%
\pgfpathlineto{\pgfqpoint{3.079854in}{2.203530in}}%
\pgfpathlineto{\pgfqpoint{3.037345in}{2.190428in}}%
\pgfpathlineto{\pgfqpoint{3.006560in}{2.209137in}}%
\pgfpathclose%
\pgfusepath{stroke,fill}%
\end{pgfscope}%
\begin{pgfscope}%
\pgfpathrectangle{\pgfqpoint{0.050000in}{0.050000in}}{\pgfqpoint{4.900000in}{4.900000in}}%
\pgfusepath{clip}%
\pgfsetbuttcap%
\pgfsetroundjoin%
\definecolor{currentfill}{rgb}{1.000000,1.000000,1.000000}%
\pgfsetfillcolor{currentfill}%
\pgfsetfillopacity{0.900000}%
\pgfsetlinewidth{0.501875pt}%
\definecolor{currentstroke}{rgb}{0.200000,0.200000,0.200000}%
\pgfsetstrokecolor{currentstroke}%
\pgfsetdash{}{0pt}%
\pgfpathmoveto{\pgfqpoint{2.366143in}{2.214224in}}%
\pgfpathlineto{\pgfqpoint{2.409294in}{2.223517in}}%
\pgfpathlineto{\pgfqpoint{2.439166in}{2.206313in}}%
\pgfpathlineto{\pgfqpoint{2.395978in}{2.198594in}}%
\pgfpathlineto{\pgfqpoint{2.366143in}{2.214224in}}%
\pgfpathclose%
\pgfusepath{stroke,fill}%
\end{pgfscope}%
\begin{pgfscope}%
\pgfpathrectangle{\pgfqpoint{0.050000in}{0.050000in}}{\pgfqpoint{4.900000in}{4.900000in}}%
\pgfusepath{clip}%
\pgfsetbuttcap%
\pgfsetroundjoin%
\definecolor{currentfill}{rgb}{1.000000,1.000000,1.000000}%
\pgfsetfillcolor{currentfill}%
\pgfsetfillopacity{0.900000}%
\pgfsetlinewidth{0.501875pt}%
\definecolor{currentstroke}{rgb}{0.200000,0.200000,0.200000}%
\pgfsetstrokecolor{currentstroke}%
\pgfsetdash{}{0pt}%
\pgfpathmoveto{\pgfqpoint{2.817487in}{2.207288in}}%
\pgfpathlineto{\pgfqpoint{2.860166in}{2.220337in}}%
\pgfpathlineto{\pgfqpoint{2.890708in}{2.201669in}}%
\pgfpathlineto{\pgfqpoint{2.847998in}{2.188606in}}%
\pgfpathlineto{\pgfqpoint{2.817487in}{2.207288in}}%
\pgfpathclose%
\pgfusepath{stroke,fill}%
\end{pgfscope}%
\begin{pgfscope}%
\pgfpathrectangle{\pgfqpoint{0.050000in}{0.050000in}}{\pgfqpoint{4.900000in}{4.900000in}}%
\pgfusepath{clip}%
\pgfsetbuttcap%
\pgfsetroundjoin%
\definecolor{currentfill}{rgb}{1.000000,1.000000,1.000000}%
\pgfsetfillcolor{currentfill}%
\pgfsetfillopacity{0.900000}%
\pgfsetlinewidth{0.501875pt}%
\definecolor{currentstroke}{rgb}{0.200000,0.200000,0.200000}%
\pgfsetstrokecolor{currentstroke}%
\pgfsetdash{}{0pt}%
\pgfpathmoveto{\pgfqpoint{2.628358in}{2.205671in}}%
\pgfpathlineto{\pgfqpoint{2.671238in}{2.218535in}}%
\pgfpathlineto{\pgfqpoint{2.701507in}{2.199949in}}%
\pgfpathlineto{\pgfqpoint{2.658595in}{2.187205in}}%
\pgfpathlineto{\pgfqpoint{2.628358in}{2.205671in}}%
\pgfpathclose%
\pgfusepath{stroke,fill}%
\end{pgfscope}%
\begin{pgfscope}%
\pgfpathrectangle{\pgfqpoint{0.050000in}{0.050000in}}{\pgfqpoint{4.900000in}{4.900000in}}%
\pgfusepath{clip}%
\pgfsetbuttcap%
\pgfsetroundjoin%
\definecolor{currentfill}{rgb}{1.000000,1.000000,1.000000}%
\pgfsetfillcolor{currentfill}%
\pgfsetfillopacity{0.900000}%
\pgfsetlinewidth{0.501875pt}%
\definecolor{currentstroke}{rgb}{0.200000,0.200000,0.200000}%
\pgfsetstrokecolor{currentstroke}%
\pgfsetdash{}{0pt}%
\pgfpathmoveto{\pgfqpoint{3.079854in}{2.203530in}}%
\pgfpathlineto{\pgfqpoint{3.122276in}{2.216609in}}%
\pgfpathlineto{\pgfqpoint{3.153212in}{2.197917in}}%
\pgfpathlineto{\pgfqpoint{3.110760in}{2.184802in}}%
\pgfpathlineto{\pgfqpoint{3.079854in}{2.203530in}}%
\pgfpathclose%
\pgfusepath{stroke,fill}%
\end{pgfscope}%
\begin{pgfscope}%
\pgfpathrectangle{\pgfqpoint{0.050000in}{0.050000in}}{\pgfqpoint{4.900000in}{4.900000in}}%
\pgfusepath{clip}%
\pgfsetbuttcap%
\pgfsetroundjoin%
\definecolor{currentfill}{rgb}{1.000000,1.000000,1.000000}%
\pgfsetfillcolor{currentfill}%
\pgfsetfillopacity{0.900000}%
\pgfsetlinewidth{0.501875pt}%
\definecolor{currentstroke}{rgb}{0.200000,0.200000,0.200000}%
\pgfsetstrokecolor{currentstroke}%
\pgfsetdash{}{0pt}%
\pgfpathmoveto{\pgfqpoint{2.439166in}{2.206313in}}%
\pgfpathlineto{\pgfqpoint{2.482252in}{2.217381in}}%
\pgfpathlineto{\pgfqpoint{2.512247in}{2.199487in}}%
\pgfpathlineto{\pgfqpoint{2.469127in}{2.189330in}}%
\pgfpathlineto{\pgfqpoint{2.439166in}{2.206313in}}%
\pgfpathclose%
\pgfusepath{stroke,fill}%
\end{pgfscope}%
\begin{pgfscope}%
\pgfpathrectangle{\pgfqpoint{0.050000in}{0.050000in}}{\pgfqpoint{4.900000in}{4.900000in}}%
\pgfusepath{clip}%
\pgfsetbuttcap%
\pgfsetroundjoin%
\definecolor{currentfill}{rgb}{1.000000,1.000000,1.000000}%
\pgfsetfillcolor{currentfill}%
\pgfsetfillopacity{0.900000}%
\pgfsetlinewidth{0.501875pt}%
\definecolor{currentstroke}{rgb}{0.200000,0.200000,0.200000}%
\pgfsetstrokecolor{currentstroke}%
\pgfsetdash{}{0pt}%
\pgfpathmoveto{\pgfqpoint{2.890708in}{2.201669in}}%
\pgfpathlineto{\pgfqpoint{2.933331in}{2.214739in}}%
\pgfpathlineto{\pgfqpoint{2.963994in}{2.196049in}}%
\pgfpathlineto{\pgfqpoint{2.921340in}{2.182956in}}%
\pgfpathlineto{\pgfqpoint{2.890708in}{2.201669in}}%
\pgfpathclose%
\pgfusepath{stroke,fill}%
\end{pgfscope}%
\begin{pgfscope}%
\pgfpathrectangle{\pgfqpoint{0.050000in}{0.050000in}}{\pgfqpoint{4.900000in}{4.900000in}}%
\pgfusepath{clip}%
\pgfsetbuttcap%
\pgfsetroundjoin%
\definecolor{currentfill}{rgb}{1.000000,1.000000,1.000000}%
\pgfsetfillcolor{currentfill}%
\pgfsetfillopacity{0.900000}%
\pgfsetlinewidth{0.501875pt}%
\definecolor{currentstroke}{rgb}{0.200000,0.200000,0.200000}%
\pgfsetstrokecolor{currentstroke}%
\pgfsetdash{}{0pt}%
\pgfpathmoveto{\pgfqpoint{2.701507in}{2.199949in}}%
\pgfpathlineto{\pgfqpoint{2.744330in}{2.212909in}}%
\pgfpathlineto{\pgfqpoint{2.774720in}{2.194266in}}%
\pgfpathlineto{\pgfqpoint{2.731865in}{2.181361in}}%
\pgfpathlineto{\pgfqpoint{2.701507in}{2.199949in}}%
\pgfpathclose%
\pgfusepath{stroke,fill}%
\end{pgfscope}%
\begin{pgfscope}%
\pgfpathrectangle{\pgfqpoint{0.050000in}{0.050000in}}{\pgfqpoint{4.900000in}{4.900000in}}%
\pgfusepath{clip}%
\pgfsetbuttcap%
\pgfsetroundjoin%
\definecolor{currentfill}{rgb}{1.000000,1.000000,1.000000}%
\pgfsetfillcolor{currentfill}%
\pgfsetfillopacity{0.900000}%
\pgfsetlinewidth{0.501875pt}%
\definecolor{currentstroke}{rgb}{0.200000,0.200000,0.200000}%
\pgfsetstrokecolor{currentstroke}%
\pgfsetdash{}{0pt}%
\pgfpathmoveto{\pgfqpoint{2.512247in}{2.199487in}}%
\pgfpathlineto{\pgfqpoint{2.555273in}{2.211463in}}%
\pgfpathlineto{\pgfqpoint{2.585389in}{2.193204in}}%
\pgfpathlineto{\pgfqpoint{2.542331in}{2.181744in}}%
\pgfpathlineto{\pgfqpoint{2.512247in}{2.199487in}}%
\pgfpathclose%
\pgfusepath{stroke,fill}%
\end{pgfscope}%
\begin{pgfscope}%
\pgfpathrectangle{\pgfqpoint{0.050000in}{0.050000in}}{\pgfqpoint{4.900000in}{4.900000in}}%
\pgfusepath{clip}%
\pgfsetbuttcap%
\pgfsetroundjoin%
\definecolor{currentfill}{rgb}{1.000000,1.000000,1.000000}%
\pgfsetfillcolor{currentfill}%
\pgfsetfillopacity{0.900000}%
\pgfsetlinewidth{0.501875pt}%
\definecolor{currentstroke}{rgb}{0.200000,0.200000,0.200000}%
\pgfsetstrokecolor{currentstroke}%
\pgfsetdash{}{0pt}%
\pgfpathmoveto{\pgfqpoint{2.963994in}{2.196049in}}%
\pgfpathlineto{\pgfqpoint{3.006560in}{2.209137in}}%
\pgfpathlineto{\pgfqpoint{3.037345in}{2.190428in}}%
\pgfpathlineto{\pgfqpoint{2.994748in}{2.177311in}}%
\pgfpathlineto{\pgfqpoint{2.963994in}{2.196049in}}%
\pgfpathclose%
\pgfusepath{stroke,fill}%
\end{pgfscope}%
\begin{pgfscope}%
\pgfpathrectangle{\pgfqpoint{0.050000in}{0.050000in}}{\pgfqpoint{4.900000in}{4.900000in}}%
\pgfusepath{clip}%
\pgfsetbuttcap%
\pgfsetroundjoin%
\definecolor{currentfill}{rgb}{1.000000,1.000000,1.000000}%
\pgfsetfillcolor{currentfill}%
\pgfsetfillopacity{0.900000}%
\pgfsetlinewidth{0.501875pt}%
\definecolor{currentstroke}{rgb}{0.200000,0.200000,0.200000}%
\pgfsetstrokecolor{currentstroke}%
\pgfsetdash{}{0pt}%
\pgfpathmoveto{\pgfqpoint{2.774720in}{2.194266in}}%
\pgfpathlineto{\pgfqpoint{2.817487in}{2.207288in}}%
\pgfpathlineto{\pgfqpoint{2.847998in}{2.188606in}}%
\pgfpathlineto{\pgfqpoint{2.805200in}{2.175601in}}%
\pgfpathlineto{\pgfqpoint{2.774720in}{2.194266in}}%
\pgfpathclose%
\pgfusepath{stroke,fill}%
\end{pgfscope}%
\begin{pgfscope}%
\pgfpathrectangle{\pgfqpoint{0.050000in}{0.050000in}}{\pgfqpoint{4.900000in}{4.900000in}}%
\pgfusepath{clip}%
\pgfsetbuttcap%
\pgfsetroundjoin%
\definecolor{currentfill}{rgb}{0.998155,0.998155,0.998155}%
\pgfsetfillcolor{currentfill}%
\pgfsetfillopacity{0.900000}%
\pgfsetlinewidth{0.501875pt}%
\definecolor{currentstroke}{rgb}{0.200000,0.200000,0.200000}%
\pgfsetstrokecolor{currentstroke}%
\pgfsetdash{}{0pt}%
\pgfpathmoveto{\pgfqpoint{2.322869in}{2.211236in}}%
\pgfpathlineto{\pgfqpoint{2.366143in}{2.214224in}}%
\pgfpathlineto{\pgfqpoint{2.395978in}{2.198594in}}%
\pgfpathlineto{\pgfqpoint{2.352667in}{2.197299in}}%
\pgfpathlineto{\pgfqpoint{2.322869in}{2.211236in}}%
\pgfpathclose%
\pgfusepath{stroke,fill}%
\end{pgfscope}%
\begin{pgfscope}%
\pgfpathrectangle{\pgfqpoint{0.050000in}{0.050000in}}{\pgfqpoint{4.900000in}{4.900000in}}%
\pgfusepath{clip}%
\pgfsetbuttcap%
\pgfsetroundjoin%
\definecolor{currentfill}{rgb}{1.000000,1.000000,1.000000}%
\pgfsetfillcolor{currentfill}%
\pgfsetfillopacity{0.900000}%
\pgfsetlinewidth{0.501875pt}%
\definecolor{currentstroke}{rgb}{0.200000,0.200000,0.200000}%
\pgfsetstrokecolor{currentstroke}%
\pgfsetdash{}{0pt}%
\pgfpathmoveto{\pgfqpoint{2.585389in}{2.193204in}}%
\pgfpathlineto{\pgfqpoint{2.628358in}{2.205671in}}%
\pgfpathlineto{\pgfqpoint{2.658595in}{2.187205in}}%
\pgfpathlineto{\pgfqpoint{2.615595in}{2.175027in}}%
\pgfpathlineto{\pgfqpoint{2.585389in}{2.193204in}}%
\pgfpathclose%
\pgfusepath{stroke,fill}%
\end{pgfscope}%
\begin{pgfscope}%
\pgfpathrectangle{\pgfqpoint{0.050000in}{0.050000in}}{\pgfqpoint{4.900000in}{4.900000in}}%
\pgfusepath{clip}%
\pgfsetbuttcap%
\pgfsetroundjoin%
\definecolor{currentfill}{rgb}{1.000000,1.000000,1.000000}%
\pgfsetfillcolor{currentfill}%
\pgfsetfillopacity{0.900000}%
\pgfsetlinewidth{0.501875pt}%
\definecolor{currentstroke}{rgb}{0.200000,0.200000,0.200000}%
\pgfsetstrokecolor{currentstroke}%
\pgfsetdash{}{0pt}%
\pgfpathmoveto{\pgfqpoint{3.037345in}{2.190428in}}%
\pgfpathlineto{\pgfqpoint{3.079854in}{2.203530in}}%
\pgfpathlineto{\pgfqpoint{3.110760in}{2.184802in}}%
\pgfpathlineto{\pgfqpoint{3.068220in}{2.171667in}}%
\pgfpathlineto{\pgfqpoint{3.037345in}{2.190428in}}%
\pgfpathclose%
\pgfusepath{stroke,fill}%
\end{pgfscope}%
\begin{pgfscope}%
\pgfpathrectangle{\pgfqpoint{0.050000in}{0.050000in}}{\pgfqpoint{4.900000in}{4.900000in}}%
\pgfusepath{clip}%
\pgfsetbuttcap%
\pgfsetroundjoin%
\definecolor{currentfill}{rgb}{1.000000,1.000000,1.000000}%
\pgfsetfillcolor{currentfill}%
\pgfsetfillopacity{0.900000}%
\pgfsetlinewidth{0.501875pt}%
\definecolor{currentstroke}{rgb}{0.200000,0.200000,0.200000}%
\pgfsetstrokecolor{currentstroke}%
\pgfsetdash{}{0pt}%
\pgfpathmoveto{\pgfqpoint{2.395978in}{2.198594in}}%
\pgfpathlineto{\pgfqpoint{2.439166in}{2.206313in}}%
\pgfpathlineto{\pgfqpoint{2.469127in}{2.189330in}}%
\pgfpathlineto{\pgfqpoint{2.425904in}{2.182910in}}%
\pgfpathlineto{\pgfqpoint{2.395978in}{2.198594in}}%
\pgfpathclose%
\pgfusepath{stroke,fill}%
\end{pgfscope}%
\begin{pgfscope}%
\pgfpathrectangle{\pgfqpoint{0.050000in}{0.050000in}}{\pgfqpoint{4.900000in}{4.900000in}}%
\pgfusepath{clip}%
\pgfsetbuttcap%
\pgfsetroundjoin%
\definecolor{currentfill}{rgb}{1.000000,1.000000,1.000000}%
\pgfsetfillcolor{currentfill}%
\pgfsetfillopacity{0.900000}%
\pgfsetlinewidth{0.501875pt}%
\definecolor{currentstroke}{rgb}{0.200000,0.200000,0.200000}%
\pgfsetstrokecolor{currentstroke}%
\pgfsetdash{}{0pt}%
\pgfpathmoveto{\pgfqpoint{2.847998in}{2.188606in}}%
\pgfpathlineto{\pgfqpoint{2.890708in}{2.201669in}}%
\pgfpathlineto{\pgfqpoint{2.921340in}{2.182956in}}%
\pgfpathlineto{\pgfqpoint{2.878599in}{2.169888in}}%
\pgfpathlineto{\pgfqpoint{2.847998in}{2.188606in}}%
\pgfpathclose%
\pgfusepath{stroke,fill}%
\end{pgfscope}%
\begin{pgfscope}%
\pgfpathrectangle{\pgfqpoint{0.050000in}{0.050000in}}{\pgfqpoint{4.900000in}{4.900000in}}%
\pgfusepath{clip}%
\pgfsetbuttcap%
\pgfsetroundjoin%
\definecolor{currentfill}{rgb}{1.000000,1.000000,1.000000}%
\pgfsetfillcolor{currentfill}%
\pgfsetfillopacity{0.900000}%
\pgfsetlinewidth{0.501875pt}%
\definecolor{currentstroke}{rgb}{0.200000,0.200000,0.200000}%
\pgfsetstrokecolor{currentstroke}%
\pgfsetdash{}{0pt}%
\pgfpathmoveto{\pgfqpoint{2.658595in}{2.187205in}}%
\pgfpathlineto{\pgfqpoint{2.701507in}{2.199949in}}%
\pgfpathlineto{\pgfqpoint{2.731865in}{2.181361in}}%
\pgfpathlineto{\pgfqpoint{2.688922in}{2.168774in}}%
\pgfpathlineto{\pgfqpoint{2.658595in}{2.187205in}}%
\pgfpathclose%
\pgfusepath{stroke,fill}%
\end{pgfscope}%
\begin{pgfscope}%
\pgfpathrectangle{\pgfqpoint{0.050000in}{0.050000in}}{\pgfqpoint{4.900000in}{4.900000in}}%
\pgfusepath{clip}%
\pgfsetbuttcap%
\pgfsetroundjoin%
\definecolor{currentfill}{rgb}{1.000000,1.000000,1.000000}%
\pgfsetfillcolor{currentfill}%
\pgfsetfillopacity{0.900000}%
\pgfsetlinewidth{0.501875pt}%
\definecolor{currentstroke}{rgb}{0.200000,0.200000,0.200000}%
\pgfsetstrokecolor{currentstroke}%
\pgfsetdash{}{0pt}%
\pgfpathmoveto{\pgfqpoint{3.110760in}{2.184802in}}%
\pgfpathlineto{\pgfqpoint{3.153212in}{2.197917in}}%
\pgfpathlineto{\pgfqpoint{3.184240in}{2.179169in}}%
\pgfpathlineto{\pgfqpoint{3.141758in}{2.166021in}}%
\pgfpathlineto{\pgfqpoint{3.110760in}{2.184802in}}%
\pgfpathclose%
\pgfusepath{stroke,fill}%
\end{pgfscope}%
\begin{pgfscope}%
\pgfpathrectangle{\pgfqpoint{0.050000in}{0.050000in}}{\pgfqpoint{4.900000in}{4.900000in}}%
\pgfusepath{clip}%
\pgfsetbuttcap%
\pgfsetroundjoin%
\definecolor{currentfill}{rgb}{1.000000,1.000000,1.000000}%
\pgfsetfillcolor{currentfill}%
\pgfsetfillopacity{0.900000}%
\pgfsetlinewidth{0.501875pt}%
\definecolor{currentstroke}{rgb}{0.200000,0.200000,0.200000}%
\pgfsetstrokecolor{currentstroke}%
\pgfsetdash{}{0pt}%
\pgfpathmoveto{\pgfqpoint{2.469127in}{2.189330in}}%
\pgfpathlineto{\pgfqpoint{2.512247in}{2.199487in}}%
\pgfpathlineto{\pgfqpoint{2.542331in}{2.181744in}}%
\pgfpathlineto{\pgfqpoint{2.499178in}{2.172444in}}%
\pgfpathlineto{\pgfqpoint{2.469127in}{2.189330in}}%
\pgfpathclose%
\pgfusepath{stroke,fill}%
\end{pgfscope}%
\begin{pgfscope}%
\pgfpathrectangle{\pgfqpoint{0.050000in}{0.050000in}}{\pgfqpoint{4.900000in}{4.900000in}}%
\pgfusepath{clip}%
\pgfsetbuttcap%
\pgfsetroundjoin%
\definecolor{currentfill}{rgb}{1.000000,1.000000,1.000000}%
\pgfsetfillcolor{currentfill}%
\pgfsetfillopacity{0.900000}%
\pgfsetlinewidth{0.501875pt}%
\definecolor{currentstroke}{rgb}{0.200000,0.200000,0.200000}%
\pgfsetstrokecolor{currentstroke}%
\pgfsetdash{}{0pt}%
\pgfpathmoveto{\pgfqpoint{2.921340in}{2.182956in}}%
\pgfpathlineto{\pgfqpoint{2.963994in}{2.196049in}}%
\pgfpathlineto{\pgfqpoint{2.994748in}{2.177311in}}%
\pgfpathlineto{\pgfqpoint{2.952063in}{2.164200in}}%
\pgfpathlineto{\pgfqpoint{2.921340in}{2.182956in}}%
\pgfpathclose%
\pgfusepath{stroke,fill}%
\end{pgfscope}%
\begin{pgfscope}%
\pgfpathrectangle{\pgfqpoint{0.050000in}{0.050000in}}{\pgfqpoint{4.900000in}{4.900000in}}%
\pgfusepath{clip}%
\pgfsetbuttcap%
\pgfsetroundjoin%
\definecolor{currentfill}{rgb}{1.000000,1.000000,1.000000}%
\pgfsetfillcolor{currentfill}%
\pgfsetfillopacity{0.900000}%
\pgfsetlinewidth{0.501875pt}%
\definecolor{currentstroke}{rgb}{0.200000,0.200000,0.200000}%
\pgfsetstrokecolor{currentstroke}%
\pgfsetdash{}{0pt}%
\pgfpathmoveto{\pgfqpoint{2.731865in}{2.181361in}}%
\pgfpathlineto{\pgfqpoint{2.774720in}{2.194266in}}%
\pgfpathlineto{\pgfqpoint{2.805200in}{2.175601in}}%
\pgfpathlineto{\pgfqpoint{2.762313in}{2.162775in}}%
\pgfpathlineto{\pgfqpoint{2.731865in}{2.181361in}}%
\pgfpathclose%
\pgfusepath{stroke,fill}%
\end{pgfscope}%
\begin{pgfscope}%
\pgfpathrectangle{\pgfqpoint{0.050000in}{0.050000in}}{\pgfqpoint{4.900000in}{4.900000in}}%
\pgfusepath{clip}%
\pgfsetbuttcap%
\pgfsetroundjoin%
\definecolor{currentfill}{rgb}{1.000000,1.000000,1.000000}%
\pgfsetfillcolor{currentfill}%
\pgfsetfillopacity{0.900000}%
\pgfsetlinewidth{0.501875pt}%
\definecolor{currentstroke}{rgb}{0.200000,0.200000,0.200000}%
\pgfsetstrokecolor{currentstroke}%
\pgfsetdash{}{0pt}%
\pgfpathmoveto{\pgfqpoint{2.542331in}{2.181744in}}%
\pgfpathlineto{\pgfqpoint{2.585389in}{2.193204in}}%
\pgfpathlineto{\pgfqpoint{2.615595in}{2.175027in}}%
\pgfpathlineto{\pgfqpoint{2.572504in}{2.164099in}}%
\pgfpathlineto{\pgfqpoint{2.542331in}{2.181744in}}%
\pgfpathclose%
\pgfusepath{stroke,fill}%
\end{pgfscope}%
\begin{pgfscope}%
\pgfpathrectangle{\pgfqpoint{0.050000in}{0.050000in}}{\pgfqpoint{4.900000in}{4.900000in}}%
\pgfusepath{clip}%
\pgfsetbuttcap%
\pgfsetroundjoin%
\definecolor{currentfill}{rgb}{1.000000,1.000000,1.000000}%
\pgfsetfillcolor{currentfill}%
\pgfsetfillopacity{0.900000}%
\pgfsetlinewidth{0.501875pt}%
\definecolor{currentstroke}{rgb}{0.200000,0.200000,0.200000}%
\pgfsetstrokecolor{currentstroke}%
\pgfsetdash{}{0pt}%
\pgfpathmoveto{\pgfqpoint{2.994748in}{2.177311in}}%
\pgfpathlineto{\pgfqpoint{3.037345in}{2.190428in}}%
\pgfpathlineto{\pgfqpoint{3.068220in}{2.171667in}}%
\pgfpathlineto{\pgfqpoint{3.025593in}{2.158525in}}%
\pgfpathlineto{\pgfqpoint{2.994748in}{2.177311in}}%
\pgfpathclose%
\pgfusepath{stroke,fill}%
\end{pgfscope}%
\begin{pgfscope}%
\pgfpathrectangle{\pgfqpoint{0.050000in}{0.050000in}}{\pgfqpoint{4.900000in}{4.900000in}}%
\pgfusepath{clip}%
\pgfsetbuttcap%
\pgfsetroundjoin%
\definecolor{currentfill}{rgb}{0.994464,0.994464,0.994464}%
\pgfsetfillcolor{currentfill}%
\pgfsetfillopacity{0.900000}%
\pgfsetlinewidth{0.501875pt}%
\definecolor{currentstroke}{rgb}{0.200000,0.200000,0.200000}%
\pgfsetstrokecolor{currentstroke}%
\pgfsetdash{}{0pt}%
\pgfpathmoveto{\pgfqpoint{2.279436in}{2.218946in}}%
\pgfpathlineto{\pgfqpoint{2.322869in}{2.211236in}}%
\pgfpathlineto{\pgfqpoint{2.352667in}{2.197299in}}%
\pgfpathlineto{\pgfqpoint{2.309205in}{2.205651in}}%
\pgfpathlineto{\pgfqpoint{2.279436in}{2.218946in}}%
\pgfpathclose%
\pgfusepath{stroke,fill}%
\end{pgfscope}%
\begin{pgfscope}%
\pgfpathrectangle{\pgfqpoint{0.050000in}{0.050000in}}{\pgfqpoint{4.900000in}{4.900000in}}%
\pgfusepath{clip}%
\pgfsetbuttcap%
\pgfsetroundjoin%
\definecolor{currentfill}{rgb}{1.000000,1.000000,1.000000}%
\pgfsetfillcolor{currentfill}%
\pgfsetfillopacity{0.900000}%
\pgfsetlinewidth{0.501875pt}%
\definecolor{currentstroke}{rgb}{0.200000,0.200000,0.200000}%
\pgfsetstrokecolor{currentstroke}%
\pgfsetdash{}{0pt}%
\pgfpathmoveto{\pgfqpoint{2.805200in}{2.175601in}}%
\pgfpathlineto{\pgfqpoint{2.847998in}{2.188606in}}%
\pgfpathlineto{\pgfqpoint{2.878599in}{2.169888in}}%
\pgfpathlineto{\pgfqpoint{2.835769in}{2.156916in}}%
\pgfpathlineto{\pgfqpoint{2.805200in}{2.175601in}}%
\pgfpathclose%
\pgfusepath{stroke,fill}%
\end{pgfscope}%
\begin{pgfscope}%
\pgfpathrectangle{\pgfqpoint{0.050000in}{0.050000in}}{\pgfqpoint{4.900000in}{4.900000in}}%
\pgfusepath{clip}%
\pgfsetbuttcap%
\pgfsetroundjoin%
\definecolor{currentfill}{rgb}{0.996309,0.996309,0.996309}%
\pgfsetfillcolor{currentfill}%
\pgfsetfillopacity{0.900000}%
\pgfsetlinewidth{0.501875pt}%
\definecolor{currentstroke}{rgb}{0.200000,0.200000,0.200000}%
\pgfsetstrokecolor{currentstroke}%
\pgfsetdash{}{0pt}%
\pgfpathmoveto{\pgfqpoint{2.352667in}{2.197299in}}%
\pgfpathlineto{\pgfqpoint{2.395978in}{2.198594in}}%
\pgfpathlineto{\pgfqpoint{2.425904in}{2.182910in}}%
\pgfpathlineto{\pgfqpoint{2.382561in}{2.182829in}}%
\pgfpathlineto{\pgfqpoint{2.352667in}{2.197299in}}%
\pgfpathclose%
\pgfusepath{stroke,fill}%
\end{pgfscope}%
\begin{pgfscope}%
\pgfpathrectangle{\pgfqpoint{0.050000in}{0.050000in}}{\pgfqpoint{4.900000in}{4.900000in}}%
\pgfusepath{clip}%
\pgfsetbuttcap%
\pgfsetroundjoin%
\definecolor{currentfill}{rgb}{1.000000,1.000000,1.000000}%
\pgfsetfillcolor{currentfill}%
\pgfsetfillopacity{0.900000}%
\pgfsetlinewidth{0.501875pt}%
\definecolor{currentstroke}{rgb}{0.200000,0.200000,0.200000}%
\pgfsetstrokecolor{currentstroke}%
\pgfsetdash{}{0pt}%
\pgfpathmoveto{\pgfqpoint{2.615595in}{2.175027in}}%
\pgfpathlineto{\pgfqpoint{2.658595in}{2.187205in}}%
\pgfpathlineto{\pgfqpoint{2.688922in}{2.168774in}}%
\pgfpathlineto{\pgfqpoint{2.645890in}{2.156914in}}%
\pgfpathlineto{\pgfqpoint{2.615595in}{2.175027in}}%
\pgfpathclose%
\pgfusepath{stroke,fill}%
\end{pgfscope}%
\begin{pgfscope}%
\pgfpathrectangle{\pgfqpoint{0.050000in}{0.050000in}}{\pgfqpoint{4.900000in}{4.900000in}}%
\pgfusepath{clip}%
\pgfsetbuttcap%
\pgfsetroundjoin%
\definecolor{currentfill}{rgb}{1.000000,1.000000,1.000000}%
\pgfsetfillcolor{currentfill}%
\pgfsetfillopacity{0.900000}%
\pgfsetlinewidth{0.501875pt}%
\definecolor{currentstroke}{rgb}{0.200000,0.200000,0.200000}%
\pgfsetstrokecolor{currentstroke}%
\pgfsetdash{}{0pt}%
\pgfpathmoveto{\pgfqpoint{3.068220in}{2.171667in}}%
\pgfpathlineto{\pgfqpoint{3.110760in}{2.184802in}}%
\pgfpathlineto{\pgfqpoint{3.141758in}{2.166021in}}%
\pgfpathlineto{\pgfqpoint{3.099187in}{2.152856in}}%
\pgfpathlineto{\pgfqpoint{3.068220in}{2.171667in}}%
\pgfpathclose%
\pgfusepath{stroke,fill}%
\end{pgfscope}%
\begin{pgfscope}%
\pgfpathrectangle{\pgfqpoint{0.050000in}{0.050000in}}{\pgfqpoint{4.900000in}{4.900000in}}%
\pgfusepath{clip}%
\pgfsetbuttcap%
\pgfsetroundjoin%
\definecolor{currentfill}{rgb}{1.000000,1.000000,1.000000}%
\pgfsetfillcolor{currentfill}%
\pgfsetfillopacity{0.900000}%
\pgfsetlinewidth{0.501875pt}%
\definecolor{currentstroke}{rgb}{0.200000,0.200000,0.200000}%
\pgfsetstrokecolor{currentstroke}%
\pgfsetdash{}{0pt}%
\pgfpathmoveto{\pgfqpoint{2.878599in}{2.169888in}}%
\pgfpathlineto{\pgfqpoint{2.921340in}{2.182956in}}%
\pgfpathlineto{\pgfqpoint{2.952063in}{2.164200in}}%
\pgfpathlineto{\pgfqpoint{2.909291in}{2.151138in}}%
\pgfpathlineto{\pgfqpoint{2.878599in}{2.169888in}}%
\pgfpathclose%
\pgfusepath{stroke,fill}%
\end{pgfscope}%
\begin{pgfscope}%
\pgfpathrectangle{\pgfqpoint{0.050000in}{0.050000in}}{\pgfqpoint{4.900000in}{4.900000in}}%
\pgfusepath{clip}%
\pgfsetbuttcap%
\pgfsetroundjoin%
\definecolor{currentfill}{rgb}{0.998155,0.998155,0.998155}%
\pgfsetfillcolor{currentfill}%
\pgfsetfillopacity{0.900000}%
\pgfsetlinewidth{0.501875pt}%
\definecolor{currentstroke}{rgb}{0.200000,0.200000,0.200000}%
\pgfsetstrokecolor{currentstroke}%
\pgfsetdash{}{0pt}%
\pgfpathmoveto{\pgfqpoint{2.425904in}{2.182910in}}%
\pgfpathlineto{\pgfqpoint{2.469127in}{2.189330in}}%
\pgfpathlineto{\pgfqpoint{2.499178in}{2.172444in}}%
\pgfpathlineto{\pgfqpoint{2.455922in}{2.167099in}}%
\pgfpathlineto{\pgfqpoint{2.425904in}{2.182910in}}%
\pgfpathclose%
\pgfusepath{stroke,fill}%
\end{pgfscope}%
\begin{pgfscope}%
\pgfpathrectangle{\pgfqpoint{0.050000in}{0.050000in}}{\pgfqpoint{4.900000in}{4.900000in}}%
\pgfusepath{clip}%
\pgfsetbuttcap%
\pgfsetroundjoin%
\definecolor{currentfill}{rgb}{1.000000,1.000000,1.000000}%
\pgfsetfillcolor{currentfill}%
\pgfsetfillopacity{0.900000}%
\pgfsetlinewidth{0.501875pt}%
\definecolor{currentstroke}{rgb}{0.200000,0.200000,0.200000}%
\pgfsetstrokecolor{currentstroke}%
\pgfsetdash{}{0pt}%
\pgfpathmoveto{\pgfqpoint{2.688922in}{2.168774in}}%
\pgfpathlineto{\pgfqpoint{2.731865in}{2.181361in}}%
\pgfpathlineto{\pgfqpoint{2.762313in}{2.162775in}}%
\pgfpathlineto{\pgfqpoint{2.719339in}{2.150371in}}%
\pgfpathlineto{\pgfqpoint{2.688922in}{2.168774in}}%
\pgfpathclose%
\pgfusepath{stroke,fill}%
\end{pgfscope}%
\begin{pgfscope}%
\pgfpathrectangle{\pgfqpoint{0.050000in}{0.050000in}}{\pgfqpoint{4.900000in}{4.900000in}}%
\pgfusepath{clip}%
\pgfsetbuttcap%
\pgfsetroundjoin%
\definecolor{currentfill}{rgb}{1.000000,1.000000,1.000000}%
\pgfsetfillcolor{currentfill}%
\pgfsetfillopacity{0.900000}%
\pgfsetlinewidth{0.501875pt}%
\definecolor{currentstroke}{rgb}{0.200000,0.200000,0.200000}%
\pgfsetstrokecolor{currentstroke}%
\pgfsetdash{}{0pt}%
\pgfpathmoveto{\pgfqpoint{3.141758in}{2.166021in}}%
\pgfpathlineto{\pgfqpoint{3.184240in}{2.179169in}}%
\pgfpathlineto{\pgfqpoint{3.215360in}{2.160366in}}%
\pgfpathlineto{\pgfqpoint{3.172847in}{2.147187in}}%
\pgfpathlineto{\pgfqpoint{3.141758in}{2.166021in}}%
\pgfpathclose%
\pgfusepath{stroke,fill}%
\end{pgfscope}%
\begin{pgfscope}%
\pgfpathrectangle{\pgfqpoint{0.050000in}{0.050000in}}{\pgfqpoint{4.900000in}{4.900000in}}%
\pgfusepath{clip}%
\pgfsetbuttcap%
\pgfsetroundjoin%
\definecolor{currentfill}{rgb}{1.000000,1.000000,1.000000}%
\pgfsetfillcolor{currentfill}%
\pgfsetfillopacity{0.900000}%
\pgfsetlinewidth{0.501875pt}%
\definecolor{currentstroke}{rgb}{0.200000,0.200000,0.200000}%
\pgfsetstrokecolor{currentstroke}%
\pgfsetdash{}{0pt}%
\pgfpathmoveto{\pgfqpoint{2.499178in}{2.172444in}}%
\pgfpathlineto{\pgfqpoint{2.542331in}{2.181744in}}%
\pgfpathlineto{\pgfqpoint{2.572504in}{2.164099in}}%
\pgfpathlineto{\pgfqpoint{2.529320in}{2.155584in}}%
\pgfpathlineto{\pgfqpoint{2.499178in}{2.172444in}}%
\pgfpathclose%
\pgfusepath{stroke,fill}%
\end{pgfscope}%
\begin{pgfscope}%
\pgfpathrectangle{\pgfqpoint{0.050000in}{0.050000in}}{\pgfqpoint{4.900000in}{4.900000in}}%
\pgfusepath{clip}%
\pgfsetbuttcap%
\pgfsetroundjoin%
\definecolor{currentfill}{rgb}{1.000000,1.000000,1.000000}%
\pgfsetfillcolor{currentfill}%
\pgfsetfillopacity{0.900000}%
\pgfsetlinewidth{0.501875pt}%
\definecolor{currentstroke}{rgb}{0.200000,0.200000,0.200000}%
\pgfsetstrokecolor{currentstroke}%
\pgfsetdash{}{0pt}%
\pgfpathmoveto{\pgfqpoint{2.952063in}{2.164200in}}%
\pgfpathlineto{\pgfqpoint{2.994748in}{2.177311in}}%
\pgfpathlineto{\pgfqpoint{3.025593in}{2.158525in}}%
\pgfpathlineto{\pgfqpoint{2.982877in}{2.145403in}}%
\pgfpathlineto{\pgfqpoint{2.952063in}{2.164200in}}%
\pgfpathclose%
\pgfusepath{stroke,fill}%
\end{pgfscope}%
\begin{pgfscope}%
\pgfpathrectangle{\pgfqpoint{0.050000in}{0.050000in}}{\pgfqpoint{4.900000in}{4.900000in}}%
\pgfusepath{clip}%
\pgfsetbuttcap%
\pgfsetroundjoin%
\definecolor{currentfill}{rgb}{1.000000,1.000000,1.000000}%
\pgfsetfillcolor{currentfill}%
\pgfsetfillopacity{0.900000}%
\pgfsetlinewidth{0.501875pt}%
\definecolor{currentstroke}{rgb}{0.200000,0.200000,0.200000}%
\pgfsetstrokecolor{currentstroke}%
\pgfsetdash{}{0pt}%
\pgfpathmoveto{\pgfqpoint{2.762313in}{2.162775in}}%
\pgfpathlineto{\pgfqpoint{2.805200in}{2.175601in}}%
\pgfpathlineto{\pgfqpoint{2.835769in}{2.156916in}}%
\pgfpathlineto{\pgfqpoint{2.792852in}{2.144190in}}%
\pgfpathlineto{\pgfqpoint{2.762313in}{2.162775in}}%
\pgfpathclose%
\pgfusepath{stroke,fill}%
\end{pgfscope}%
\begin{pgfscope}%
\pgfpathrectangle{\pgfqpoint{0.050000in}{0.050000in}}{\pgfqpoint{4.900000in}{4.900000in}}%
\pgfusepath{clip}%
\pgfsetbuttcap%
\pgfsetroundjoin%
\definecolor{currentfill}{rgb}{1.000000,1.000000,1.000000}%
\pgfsetfillcolor{currentfill}%
\pgfsetfillopacity{0.900000}%
\pgfsetlinewidth{0.501875pt}%
\definecolor{currentstroke}{rgb}{0.200000,0.200000,0.200000}%
\pgfsetstrokecolor{currentstroke}%
\pgfsetdash{}{0pt}%
\pgfpathmoveto{\pgfqpoint{2.572504in}{2.164099in}}%
\pgfpathlineto{\pgfqpoint{2.615595in}{2.175027in}}%
\pgfpathlineto{\pgfqpoint{2.645890in}{2.156914in}}%
\pgfpathlineto{\pgfqpoint{2.602768in}{2.146510in}}%
\pgfpathlineto{\pgfqpoint{2.572504in}{2.164099in}}%
\pgfpathclose%
\pgfusepath{stroke,fill}%
\end{pgfscope}%
\begin{pgfscope}%
\pgfpathrectangle{\pgfqpoint{0.050000in}{0.050000in}}{\pgfqpoint{4.900000in}{4.900000in}}%
\pgfusepath{clip}%
\pgfsetbuttcap%
\pgfsetroundjoin%
\definecolor{currentfill}{rgb}{1.000000,1.000000,1.000000}%
\pgfsetfillcolor{currentfill}%
\pgfsetfillopacity{0.900000}%
\pgfsetlinewidth{0.501875pt}%
\definecolor{currentstroke}{rgb}{0.200000,0.200000,0.200000}%
\pgfsetstrokecolor{currentstroke}%
\pgfsetdash{}{0pt}%
\pgfpathmoveto{\pgfqpoint{3.025593in}{2.158525in}}%
\pgfpathlineto{\pgfqpoint{3.068220in}{2.171667in}}%
\pgfpathlineto{\pgfqpoint{3.099187in}{2.152856in}}%
\pgfpathlineto{\pgfqpoint{3.056529in}{2.139693in}}%
\pgfpathlineto{\pgfqpoint{3.025593in}{2.158525in}}%
\pgfpathclose%
\pgfusepath{stroke,fill}%
\end{pgfscope}%
\begin{pgfscope}%
\pgfpathrectangle{\pgfqpoint{0.050000in}{0.050000in}}{\pgfqpoint{4.900000in}{4.900000in}}%
\pgfusepath{clip}%
\pgfsetbuttcap%
\pgfsetroundjoin%
\definecolor{currentfill}{rgb}{0.987082,0.987082,0.987082}%
\pgfsetfillcolor{currentfill}%
\pgfsetfillopacity{0.900000}%
\pgfsetlinewidth{0.501875pt}%
\definecolor{currentstroke}{rgb}{0.200000,0.200000,0.200000}%
\pgfsetstrokecolor{currentstroke}%
\pgfsetdash{}{0pt}%
\pgfpathmoveto{\pgfqpoint{2.235805in}{2.239619in}}%
\pgfpathlineto{\pgfqpoint{2.279436in}{2.218946in}}%
\pgfpathlineto{\pgfqpoint{2.309205in}{2.205651in}}%
\pgfpathlineto{\pgfqpoint{2.265560in}{2.225257in}}%
\pgfpathlineto{\pgfqpoint{2.235805in}{2.239619in}}%
\pgfpathclose%
\pgfusepath{stroke,fill}%
\end{pgfscope}%
\begin{pgfscope}%
\pgfpathrectangle{\pgfqpoint{0.050000in}{0.050000in}}{\pgfqpoint{4.900000in}{4.900000in}}%
\pgfusepath{clip}%
\pgfsetbuttcap%
\pgfsetroundjoin%
\definecolor{currentfill}{rgb}{0.992618,0.992618,0.992618}%
\pgfsetfillcolor{currentfill}%
\pgfsetfillopacity{0.900000}%
\pgfsetlinewidth{0.501875pt}%
\definecolor{currentstroke}{rgb}{0.200000,0.200000,0.200000}%
\pgfsetstrokecolor{currentstroke}%
\pgfsetdash{}{0pt}%
\pgfpathmoveto{\pgfqpoint{2.309205in}{2.205651in}}%
\pgfpathlineto{\pgfqpoint{2.352667in}{2.197299in}}%
\pgfpathlineto{\pgfqpoint{2.382561in}{2.182829in}}%
\pgfpathlineto{\pgfqpoint{2.339072in}{2.191593in}}%
\pgfpathlineto{\pgfqpoint{2.309205in}{2.205651in}}%
\pgfpathclose%
\pgfusepath{stroke,fill}%
\end{pgfscope}%
\begin{pgfscope}%
\pgfpathrectangle{\pgfqpoint{0.050000in}{0.050000in}}{\pgfqpoint{4.900000in}{4.900000in}}%
\pgfusepath{clip}%
\pgfsetbuttcap%
\pgfsetroundjoin%
\definecolor{currentfill}{rgb}{1.000000,1.000000,1.000000}%
\pgfsetfillcolor{currentfill}%
\pgfsetfillopacity{0.900000}%
\pgfsetlinewidth{0.501875pt}%
\definecolor{currentstroke}{rgb}{0.200000,0.200000,0.200000}%
\pgfsetstrokecolor{currentstroke}%
\pgfsetdash{}{0pt}%
\pgfpathmoveto{\pgfqpoint{2.835769in}{2.156916in}}%
\pgfpathlineto{\pgfqpoint{2.878599in}{2.169888in}}%
\pgfpathlineto{\pgfqpoint{2.909291in}{2.151138in}}%
\pgfpathlineto{\pgfqpoint{2.866430in}{2.138214in}}%
\pgfpathlineto{\pgfqpoint{2.835769in}{2.156916in}}%
\pgfpathclose%
\pgfusepath{stroke,fill}%
\end{pgfscope}%
\begin{pgfscope}%
\pgfpathrectangle{\pgfqpoint{0.050000in}{0.050000in}}{\pgfqpoint{4.900000in}{4.900000in}}%
\pgfusepath{clip}%
\pgfsetbuttcap%
\pgfsetroundjoin%
\definecolor{currentfill}{rgb}{1.000000,1.000000,1.000000}%
\pgfsetfillcolor{currentfill}%
\pgfsetfillopacity{0.900000}%
\pgfsetlinewidth{0.501875pt}%
\definecolor{currentstroke}{rgb}{0.200000,0.200000,0.200000}%
\pgfsetstrokecolor{currentstroke}%
\pgfsetdash{}{0pt}%
\pgfpathmoveto{\pgfqpoint{2.645890in}{2.156914in}}%
\pgfpathlineto{\pgfqpoint{2.688922in}{2.168774in}}%
\pgfpathlineto{\pgfqpoint{2.719339in}{2.150371in}}%
\pgfpathlineto{\pgfqpoint{2.676276in}{2.138844in}}%
\pgfpathlineto{\pgfqpoint{2.645890in}{2.156914in}}%
\pgfpathclose%
\pgfusepath{stroke,fill}%
\end{pgfscope}%
\begin{pgfscope}%
\pgfpathrectangle{\pgfqpoint{0.050000in}{0.050000in}}{\pgfqpoint{4.900000in}{4.900000in}}%
\pgfusepath{clip}%
\pgfsetbuttcap%
\pgfsetroundjoin%
\definecolor{currentfill}{rgb}{0.996309,0.996309,0.996309}%
\pgfsetfillcolor{currentfill}%
\pgfsetfillopacity{0.900000}%
\pgfsetlinewidth{0.501875pt}%
\definecolor{currentstroke}{rgb}{0.200000,0.200000,0.200000}%
\pgfsetstrokecolor{currentstroke}%
\pgfsetdash{}{0pt}%
\pgfpathmoveto{\pgfqpoint{2.382561in}{2.182829in}}%
\pgfpathlineto{\pgfqpoint{2.425904in}{2.182910in}}%
\pgfpathlineto{\pgfqpoint{2.455922in}{2.167099in}}%
\pgfpathlineto{\pgfqpoint{2.412548in}{2.167936in}}%
\pgfpathlineto{\pgfqpoint{2.382561in}{2.182829in}}%
\pgfpathclose%
\pgfusepath{stroke,fill}%
\end{pgfscope}%
\begin{pgfscope}%
\pgfpathrectangle{\pgfqpoint{0.050000in}{0.050000in}}{\pgfqpoint{4.900000in}{4.900000in}}%
\pgfusepath{clip}%
\pgfsetbuttcap%
\pgfsetroundjoin%
\definecolor{currentfill}{rgb}{1.000000,1.000000,1.000000}%
\pgfsetfillcolor{currentfill}%
\pgfsetfillopacity{0.900000}%
\pgfsetlinewidth{0.501875pt}%
\definecolor{currentstroke}{rgb}{0.200000,0.200000,0.200000}%
\pgfsetstrokecolor{currentstroke}%
\pgfsetdash{}{0pt}%
\pgfpathmoveto{\pgfqpoint{3.099187in}{2.152856in}}%
\pgfpathlineto{\pgfqpoint{3.141758in}{2.166021in}}%
\pgfpathlineto{\pgfqpoint{3.172847in}{2.147187in}}%
\pgfpathlineto{\pgfqpoint{3.130245in}{2.133996in}}%
\pgfpathlineto{\pgfqpoint{3.099187in}{2.152856in}}%
\pgfpathclose%
\pgfusepath{stroke,fill}%
\end{pgfscope}%
\begin{pgfscope}%
\pgfpathrectangle{\pgfqpoint{0.050000in}{0.050000in}}{\pgfqpoint{4.900000in}{4.900000in}}%
\pgfusepath{clip}%
\pgfsetbuttcap%
\pgfsetroundjoin%
\definecolor{currentfill}{rgb}{1.000000,1.000000,1.000000}%
\pgfsetfillcolor{currentfill}%
\pgfsetfillopacity{0.900000}%
\pgfsetlinewidth{0.501875pt}%
\definecolor{currentstroke}{rgb}{0.200000,0.200000,0.200000}%
\pgfsetstrokecolor{currentstroke}%
\pgfsetdash{}{0pt}%
\pgfpathmoveto{\pgfqpoint{2.909291in}{2.151138in}}%
\pgfpathlineto{\pgfqpoint{2.952063in}{2.164200in}}%
\pgfpathlineto{\pgfqpoint{2.982877in}{2.145403in}}%
\pgfpathlineto{\pgfqpoint{2.940073in}{2.132357in}}%
\pgfpathlineto{\pgfqpoint{2.909291in}{2.151138in}}%
\pgfpathclose%
\pgfusepath{stroke,fill}%
\end{pgfscope}%
\begin{pgfscope}%
\pgfpathrectangle{\pgfqpoint{0.050000in}{0.050000in}}{\pgfqpoint{4.900000in}{4.900000in}}%
\pgfusepath{clip}%
\pgfsetbuttcap%
\pgfsetroundjoin%
\definecolor{currentfill}{rgb}{0.998155,0.998155,0.998155}%
\pgfsetfillcolor{currentfill}%
\pgfsetfillopacity{0.900000}%
\pgfsetlinewidth{0.501875pt}%
\definecolor{currentstroke}{rgb}{0.200000,0.200000,0.200000}%
\pgfsetstrokecolor{currentstroke}%
\pgfsetdash{}{0pt}%
\pgfpathmoveto{\pgfqpoint{2.455922in}{2.167099in}}%
\pgfpathlineto{\pgfqpoint{2.499178in}{2.172444in}}%
\pgfpathlineto{\pgfqpoint{2.529320in}{2.155584in}}%
\pgfpathlineto{\pgfqpoint{2.486032in}{2.151139in}}%
\pgfpathlineto{\pgfqpoint{2.455922in}{2.167099in}}%
\pgfpathclose%
\pgfusepath{stroke,fill}%
\end{pgfscope}%
\begin{pgfscope}%
\pgfpathrectangle{\pgfqpoint{0.050000in}{0.050000in}}{\pgfqpoint{4.900000in}{4.900000in}}%
\pgfusepath{clip}%
\pgfsetbuttcap%
\pgfsetroundjoin%
\definecolor{currentfill}{rgb}{1.000000,1.000000,1.000000}%
\pgfsetfillcolor{currentfill}%
\pgfsetfillopacity{0.900000}%
\pgfsetlinewidth{0.501875pt}%
\definecolor{currentstroke}{rgb}{0.200000,0.200000,0.200000}%
\pgfsetstrokecolor{currentstroke}%
\pgfsetdash{}{0pt}%
\pgfpathmoveto{\pgfqpoint{2.719339in}{2.150371in}}%
\pgfpathlineto{\pgfqpoint{2.762313in}{2.162775in}}%
\pgfpathlineto{\pgfqpoint{2.792852in}{2.144190in}}%
\pgfpathlineto{\pgfqpoint{2.749847in}{2.131988in}}%
\pgfpathlineto{\pgfqpoint{2.719339in}{2.150371in}}%
\pgfpathclose%
\pgfusepath{stroke,fill}%
\end{pgfscope}%
\begin{pgfscope}%
\pgfpathrectangle{\pgfqpoint{0.050000in}{0.050000in}}{\pgfqpoint{4.900000in}{4.900000in}}%
\pgfusepath{clip}%
\pgfsetbuttcap%
\pgfsetroundjoin%
\definecolor{currentfill}{rgb}{1.000000,1.000000,1.000000}%
\pgfsetfillcolor{currentfill}%
\pgfsetfillopacity{0.900000}%
\pgfsetlinewidth{0.501875pt}%
\definecolor{currentstroke}{rgb}{0.200000,0.200000,0.200000}%
\pgfsetstrokecolor{currentstroke}%
\pgfsetdash{}{0pt}%
\pgfpathmoveto{\pgfqpoint{3.172847in}{2.147187in}}%
\pgfpathlineto{\pgfqpoint{3.215360in}{2.160366in}}%
\pgfpathlineto{\pgfqpoint{3.246572in}{2.141508in}}%
\pgfpathlineto{\pgfqpoint{3.204028in}{2.128300in}}%
\pgfpathlineto{\pgfqpoint{3.172847in}{2.147187in}}%
\pgfpathclose%
\pgfusepath{stroke,fill}%
\end{pgfscope}%
\begin{pgfscope}%
\pgfpathrectangle{\pgfqpoint{0.050000in}{0.050000in}}{\pgfqpoint{4.900000in}{4.900000in}}%
\pgfusepath{clip}%
\pgfsetbuttcap%
\pgfsetroundjoin%
\definecolor{currentfill}{rgb}{1.000000,1.000000,1.000000}%
\pgfsetfillcolor{currentfill}%
\pgfsetfillopacity{0.900000}%
\pgfsetlinewidth{0.501875pt}%
\definecolor{currentstroke}{rgb}{0.200000,0.200000,0.200000}%
\pgfsetstrokecolor{currentstroke}%
\pgfsetdash{}{0pt}%
\pgfpathmoveto{\pgfqpoint{2.529320in}{2.155584in}}%
\pgfpathlineto{\pgfqpoint{2.572504in}{2.164099in}}%
\pgfpathlineto{\pgfqpoint{2.602768in}{2.146510in}}%
\pgfpathlineto{\pgfqpoint{2.559553in}{2.138705in}}%
\pgfpathlineto{\pgfqpoint{2.529320in}{2.155584in}}%
\pgfpathclose%
\pgfusepath{stroke,fill}%
\end{pgfscope}%
\begin{pgfscope}%
\pgfpathrectangle{\pgfqpoint{0.050000in}{0.050000in}}{\pgfqpoint{4.900000in}{4.900000in}}%
\pgfusepath{clip}%
\pgfsetbuttcap%
\pgfsetroundjoin%
\definecolor{currentfill}{rgb}{1.000000,1.000000,1.000000}%
\pgfsetfillcolor{currentfill}%
\pgfsetfillopacity{0.900000}%
\pgfsetlinewidth{0.501875pt}%
\definecolor{currentstroke}{rgb}{0.200000,0.200000,0.200000}%
\pgfsetstrokecolor{currentstroke}%
\pgfsetdash{}{0pt}%
\pgfpathmoveto{\pgfqpoint{2.982877in}{2.145403in}}%
\pgfpathlineto{\pgfqpoint{3.025593in}{2.158525in}}%
\pgfpathlineto{\pgfqpoint{3.056529in}{2.139693in}}%
\pgfpathlineto{\pgfqpoint{3.013782in}{2.126567in}}%
\pgfpathlineto{\pgfqpoint{2.982877in}{2.145403in}}%
\pgfpathclose%
\pgfusepath{stroke,fill}%
\end{pgfscope}%
\begin{pgfscope}%
\pgfpathrectangle{\pgfqpoint{0.050000in}{0.050000in}}{\pgfqpoint{4.900000in}{4.900000in}}%
\pgfusepath{clip}%
\pgfsetbuttcap%
\pgfsetroundjoin%
\definecolor{currentfill}{rgb}{1.000000,1.000000,1.000000}%
\pgfsetfillcolor{currentfill}%
\pgfsetfillopacity{0.900000}%
\pgfsetlinewidth{0.501875pt}%
\definecolor{currentstroke}{rgb}{0.200000,0.200000,0.200000}%
\pgfsetstrokecolor{currentstroke}%
\pgfsetdash{}{0pt}%
\pgfpathmoveto{\pgfqpoint{2.792852in}{2.144190in}}%
\pgfpathlineto{\pgfqpoint{2.835769in}{2.156916in}}%
\pgfpathlineto{\pgfqpoint{2.866430in}{2.138214in}}%
\pgfpathlineto{\pgfqpoint{2.823482in}{2.125604in}}%
\pgfpathlineto{\pgfqpoint{2.792852in}{2.144190in}}%
\pgfpathclose%
\pgfusepath{stroke,fill}%
\end{pgfscope}%
\begin{pgfscope}%
\pgfpathrectangle{\pgfqpoint{0.050000in}{0.050000in}}{\pgfqpoint{4.900000in}{4.900000in}}%
\pgfusepath{clip}%
\pgfsetbuttcap%
\pgfsetroundjoin%
\definecolor{currentfill}{rgb}{1.000000,1.000000,1.000000}%
\pgfsetfillcolor{currentfill}%
\pgfsetfillopacity{0.900000}%
\pgfsetlinewidth{0.501875pt}%
\definecolor{currentstroke}{rgb}{0.200000,0.200000,0.200000}%
\pgfsetstrokecolor{currentstroke}%
\pgfsetdash{}{0pt}%
\pgfpathmoveto{\pgfqpoint{2.602768in}{2.146510in}}%
\pgfpathlineto{\pgfqpoint{2.645890in}{2.156914in}}%
\pgfpathlineto{\pgfqpoint{2.676276in}{2.138844in}}%
\pgfpathlineto{\pgfqpoint{2.633124in}{2.128946in}}%
\pgfpathlineto{\pgfqpoint{2.602768in}{2.146510in}}%
\pgfpathclose%
\pgfusepath{stroke,fill}%
\end{pgfscope}%
\begin{pgfscope}%
\pgfpathrectangle{\pgfqpoint{0.050000in}{0.050000in}}{\pgfqpoint{4.900000in}{4.900000in}}%
\pgfusepath{clip}%
\pgfsetbuttcap%
\pgfsetroundjoin%
\definecolor{currentfill}{rgb}{1.000000,1.000000,1.000000}%
\pgfsetfillcolor{currentfill}%
\pgfsetfillopacity{0.900000}%
\pgfsetlinewidth{0.501875pt}%
\definecolor{currentstroke}{rgb}{0.200000,0.200000,0.200000}%
\pgfsetstrokecolor{currentstroke}%
\pgfsetdash{}{0pt}%
\pgfpathmoveto{\pgfqpoint{3.056529in}{2.139693in}}%
\pgfpathlineto{\pgfqpoint{3.099187in}{2.152856in}}%
\pgfpathlineto{\pgfqpoint{3.130245in}{2.133996in}}%
\pgfpathlineto{\pgfqpoint{3.087556in}{2.120816in}}%
\pgfpathlineto{\pgfqpoint{3.056529in}{2.139693in}}%
\pgfpathclose%
\pgfusepath{stroke,fill}%
\end{pgfscope}%
\begin{pgfscope}%
\pgfpathrectangle{\pgfqpoint{0.050000in}{0.050000in}}{\pgfqpoint{4.900000in}{4.900000in}}%
\pgfusepath{clip}%
\pgfsetbuttcap%
\pgfsetroundjoin%
\definecolor{currentfill}{rgb}{1.000000,1.000000,1.000000}%
\pgfsetfillcolor{currentfill}%
\pgfsetfillopacity{0.900000}%
\pgfsetlinewidth{0.501875pt}%
\definecolor{currentstroke}{rgb}{0.200000,0.200000,0.200000}%
\pgfsetstrokecolor{currentstroke}%
\pgfsetdash{}{0pt}%
\pgfpathmoveto{\pgfqpoint{2.866430in}{2.138214in}}%
\pgfpathlineto{\pgfqpoint{2.909291in}{2.151138in}}%
\pgfpathlineto{\pgfqpoint{2.940073in}{2.132357in}}%
\pgfpathlineto{\pgfqpoint{2.897182in}{2.119495in}}%
\pgfpathlineto{\pgfqpoint{2.866430in}{2.138214in}}%
\pgfpathclose%
\pgfusepath{stroke,fill}%
\end{pgfscope}%
\begin{pgfscope}%
\pgfpathrectangle{\pgfqpoint{0.050000in}{0.050000in}}{\pgfqpoint{4.900000in}{4.900000in}}%
\pgfusepath{clip}%
\pgfsetbuttcap%
\pgfsetroundjoin%
\definecolor{currentfill}{rgb}{0.985236,0.985236,0.985236}%
\pgfsetfillcolor{currentfill}%
\pgfsetfillopacity{0.900000}%
\pgfsetlinewidth{0.501875pt}%
\definecolor{currentstroke}{rgb}{0.200000,0.200000,0.200000}%
\pgfsetstrokecolor{currentstroke}%
\pgfsetdash{}{0pt}%
\pgfpathmoveto{\pgfqpoint{2.265560in}{2.225257in}}%
\pgfpathlineto{\pgfqpoint{2.309205in}{2.205651in}}%
\pgfpathlineto{\pgfqpoint{2.339072in}{2.191593in}}%
\pgfpathlineto{\pgfqpoint{2.295410in}{2.210443in}}%
\pgfpathlineto{\pgfqpoint{2.265560in}{2.225257in}}%
\pgfpathclose%
\pgfusepath{stroke,fill}%
\end{pgfscope}%
\begin{pgfscope}%
\pgfpathrectangle{\pgfqpoint{0.050000in}{0.050000in}}{\pgfqpoint{4.900000in}{4.900000in}}%
\pgfusepath{clip}%
\pgfsetbuttcap%
\pgfsetroundjoin%
\definecolor{currentfill}{rgb}{0.990773,0.990773,0.990773}%
\pgfsetfillcolor{currentfill}%
\pgfsetfillopacity{0.900000}%
\pgfsetlinewidth{0.501875pt}%
\definecolor{currentstroke}{rgb}{0.200000,0.200000,0.200000}%
\pgfsetstrokecolor{currentstroke}%
\pgfsetdash{}{0pt}%
\pgfpathmoveto{\pgfqpoint{2.339072in}{2.191593in}}%
\pgfpathlineto{\pgfqpoint{2.382561in}{2.182829in}}%
\pgfpathlineto{\pgfqpoint{2.412548in}{2.167936in}}%
\pgfpathlineto{\pgfqpoint{2.369033in}{2.176984in}}%
\pgfpathlineto{\pgfqpoint{2.339072in}{2.191593in}}%
\pgfpathclose%
\pgfusepath{stroke,fill}%
\end{pgfscope}%
\begin{pgfscope}%
\pgfpathrectangle{\pgfqpoint{0.050000in}{0.050000in}}{\pgfqpoint{4.900000in}{4.900000in}}%
\pgfusepath{clip}%
\pgfsetbuttcap%
\pgfsetroundjoin%
\definecolor{currentfill}{rgb}{0.977855,0.977855,0.977855}%
\pgfsetfillcolor{currentfill}%
\pgfsetfillopacity{0.900000}%
\pgfsetlinewidth{0.501875pt}%
\definecolor{currentstroke}{rgb}{0.200000,0.200000,0.200000}%
\pgfsetstrokecolor{currentstroke}%
\pgfsetdash{}{0pt}%
\pgfpathmoveto{\pgfqpoint{2.191959in}{2.271387in}}%
\pgfpathlineto{\pgfqpoint{2.235805in}{2.239619in}}%
\pgfpathlineto{\pgfqpoint{2.265560in}{2.225257in}}%
\pgfpathlineto{\pgfqpoint{2.221713in}{2.254979in}}%
\pgfpathlineto{\pgfqpoint{2.191959in}{2.271387in}}%
\pgfpathclose%
\pgfusepath{stroke,fill}%
\end{pgfscope}%
\begin{pgfscope}%
\pgfpathrectangle{\pgfqpoint{0.050000in}{0.050000in}}{\pgfqpoint{4.900000in}{4.900000in}}%
\pgfusepath{clip}%
\pgfsetbuttcap%
\pgfsetroundjoin%
\definecolor{currentfill}{rgb}{1.000000,1.000000,1.000000}%
\pgfsetfillcolor{currentfill}%
\pgfsetfillopacity{0.900000}%
\pgfsetlinewidth{0.501875pt}%
\definecolor{currentstroke}{rgb}{0.200000,0.200000,0.200000}%
\pgfsetstrokecolor{currentstroke}%
\pgfsetdash{}{0pt}%
\pgfpathmoveto{\pgfqpoint{2.676276in}{2.138844in}}%
\pgfpathlineto{\pgfqpoint{2.719339in}{2.150371in}}%
\pgfpathlineto{\pgfqpoint{2.749847in}{2.131988in}}%
\pgfpathlineto{\pgfqpoint{2.706754in}{2.120799in}}%
\pgfpathlineto{\pgfqpoint{2.676276in}{2.138844in}}%
\pgfpathclose%
\pgfusepath{stroke,fill}%
\end{pgfscope}%
\begin{pgfscope}%
\pgfpathrectangle{\pgfqpoint{0.050000in}{0.050000in}}{\pgfqpoint{4.900000in}{4.900000in}}%
\pgfusepath{clip}%
\pgfsetbuttcap%
\pgfsetroundjoin%
\definecolor{currentfill}{rgb}{0.994464,0.994464,0.994464}%
\pgfsetfillcolor{currentfill}%
\pgfsetfillopacity{0.900000}%
\pgfsetlinewidth{0.501875pt}%
\definecolor{currentstroke}{rgb}{0.200000,0.200000,0.200000}%
\pgfsetstrokecolor{currentstroke}%
\pgfsetdash{}{0pt}%
\pgfpathmoveto{\pgfqpoint{2.412548in}{2.167936in}}%
\pgfpathlineto{\pgfqpoint{2.455922in}{2.167099in}}%
\pgfpathlineto{\pgfqpoint{2.486032in}{2.151139in}}%
\pgfpathlineto{\pgfqpoint{2.442629in}{2.152695in}}%
\pgfpathlineto{\pgfqpoint{2.412548in}{2.167936in}}%
\pgfpathclose%
\pgfusepath{stroke,fill}%
\end{pgfscope}%
\begin{pgfscope}%
\pgfpathrectangle{\pgfqpoint{0.050000in}{0.050000in}}{\pgfqpoint{4.900000in}{4.900000in}}%
\pgfusepath{clip}%
\pgfsetbuttcap%
\pgfsetroundjoin%
\definecolor{currentfill}{rgb}{1.000000,1.000000,1.000000}%
\pgfsetfillcolor{currentfill}%
\pgfsetfillopacity{0.900000}%
\pgfsetlinewidth{0.501875pt}%
\definecolor{currentstroke}{rgb}{0.200000,0.200000,0.200000}%
\pgfsetstrokecolor{currentstroke}%
\pgfsetdash{}{0pt}%
\pgfpathmoveto{\pgfqpoint{3.130245in}{2.133996in}}%
\pgfpathlineto{\pgfqpoint{3.172847in}{2.147187in}}%
\pgfpathlineto{\pgfqpoint{3.204028in}{2.128300in}}%
\pgfpathlineto{\pgfqpoint{3.161396in}{2.115086in}}%
\pgfpathlineto{\pgfqpoint{3.130245in}{2.133996in}}%
\pgfpathclose%
\pgfusepath{stroke,fill}%
\end{pgfscope}%
\begin{pgfscope}%
\pgfpathrectangle{\pgfqpoint{0.050000in}{0.050000in}}{\pgfqpoint{4.900000in}{4.900000in}}%
\pgfusepath{clip}%
\pgfsetbuttcap%
\pgfsetroundjoin%
\definecolor{currentfill}{rgb}{1.000000,1.000000,1.000000}%
\pgfsetfillcolor{currentfill}%
\pgfsetfillopacity{0.900000}%
\pgfsetlinewidth{0.501875pt}%
\definecolor{currentstroke}{rgb}{0.200000,0.200000,0.200000}%
\pgfsetstrokecolor{currentstroke}%
\pgfsetdash{}{0pt}%
\pgfpathmoveto{\pgfqpoint{2.940073in}{2.132357in}}%
\pgfpathlineto{\pgfqpoint{2.982877in}{2.145403in}}%
\pgfpathlineto{\pgfqpoint{3.013782in}{2.126567in}}%
\pgfpathlineto{\pgfqpoint{2.970948in}{2.113547in}}%
\pgfpathlineto{\pgfqpoint{2.940073in}{2.132357in}}%
\pgfpathclose%
\pgfusepath{stroke,fill}%
\end{pgfscope}%
\begin{pgfscope}%
\pgfpathrectangle{\pgfqpoint{0.050000in}{0.050000in}}{\pgfqpoint{4.900000in}{4.900000in}}%
\pgfusepath{clip}%
\pgfsetbuttcap%
\pgfsetroundjoin%
\definecolor{currentfill}{rgb}{0.998155,0.998155,0.998155}%
\pgfsetfillcolor{currentfill}%
\pgfsetfillopacity{0.900000}%
\pgfsetlinewidth{0.501875pt}%
\definecolor{currentstroke}{rgb}{0.200000,0.200000,0.200000}%
\pgfsetstrokecolor{currentstroke}%
\pgfsetdash{}{0pt}%
\pgfpathmoveto{\pgfqpoint{2.486032in}{2.151139in}}%
\pgfpathlineto{\pgfqpoint{2.529320in}{2.155584in}}%
\pgfpathlineto{\pgfqpoint{2.559553in}{2.138705in}}%
\pgfpathlineto{\pgfqpoint{2.516235in}{2.135022in}}%
\pgfpathlineto{\pgfqpoint{2.486032in}{2.151139in}}%
\pgfpathclose%
\pgfusepath{stroke,fill}%
\end{pgfscope}%
\begin{pgfscope}%
\pgfpathrectangle{\pgfqpoint{0.050000in}{0.050000in}}{\pgfqpoint{4.900000in}{4.900000in}}%
\pgfusepath{clip}%
\pgfsetbuttcap%
\pgfsetroundjoin%
\definecolor{currentfill}{rgb}{1.000000,1.000000,1.000000}%
\pgfsetfillcolor{currentfill}%
\pgfsetfillopacity{0.900000}%
\pgfsetlinewidth{0.501875pt}%
\definecolor{currentstroke}{rgb}{0.200000,0.200000,0.200000}%
\pgfsetstrokecolor{currentstroke}%
\pgfsetdash{}{0pt}%
\pgfpathmoveto{\pgfqpoint{2.749847in}{2.131988in}}%
\pgfpathlineto{\pgfqpoint{2.792852in}{2.144190in}}%
\pgfpathlineto{\pgfqpoint{2.823482in}{2.125604in}}%
\pgfpathlineto{\pgfqpoint{2.780446in}{2.113618in}}%
\pgfpathlineto{\pgfqpoint{2.749847in}{2.131988in}}%
\pgfpathclose%
\pgfusepath{stroke,fill}%
\end{pgfscope}%
\begin{pgfscope}%
\pgfpathrectangle{\pgfqpoint{0.050000in}{0.050000in}}{\pgfqpoint{4.900000in}{4.900000in}}%
\pgfusepath{clip}%
\pgfsetbuttcap%
\pgfsetroundjoin%
\definecolor{currentfill}{rgb}{1.000000,1.000000,1.000000}%
\pgfsetfillcolor{currentfill}%
\pgfsetfillopacity{0.900000}%
\pgfsetlinewidth{0.501875pt}%
\definecolor{currentstroke}{rgb}{0.200000,0.200000,0.200000}%
\pgfsetstrokecolor{currentstroke}%
\pgfsetdash{}{0pt}%
\pgfpathmoveto{\pgfqpoint{3.204028in}{2.128300in}}%
\pgfpathlineto{\pgfqpoint{3.246572in}{2.141508in}}%
\pgfpathlineto{\pgfqpoint{3.277876in}{2.122594in}}%
\pgfpathlineto{\pgfqpoint{3.235302in}{2.109361in}}%
\pgfpathlineto{\pgfqpoint{3.204028in}{2.128300in}}%
\pgfpathclose%
\pgfusepath{stroke,fill}%
\end{pgfscope}%
\begin{pgfscope}%
\pgfpathrectangle{\pgfqpoint{0.050000in}{0.050000in}}{\pgfqpoint{4.900000in}{4.900000in}}%
\pgfusepath{clip}%
\pgfsetbuttcap%
\pgfsetroundjoin%
\definecolor{currentfill}{rgb}{1.000000,1.000000,1.000000}%
\pgfsetfillcolor{currentfill}%
\pgfsetfillopacity{0.900000}%
\pgfsetlinewidth{0.501875pt}%
\definecolor{currentstroke}{rgb}{0.200000,0.200000,0.200000}%
\pgfsetstrokecolor{currentstroke}%
\pgfsetdash{}{0pt}%
\pgfpathmoveto{\pgfqpoint{3.013782in}{2.126567in}}%
\pgfpathlineto{\pgfqpoint{3.056529in}{2.139693in}}%
\pgfpathlineto{\pgfqpoint{3.087556in}{2.120816in}}%
\pgfpathlineto{\pgfqpoint{3.044779in}{2.107693in}}%
\pgfpathlineto{\pgfqpoint{3.013782in}{2.126567in}}%
\pgfpathclose%
\pgfusepath{stroke,fill}%
\end{pgfscope}%
\begin{pgfscope}%
\pgfpathrectangle{\pgfqpoint{0.050000in}{0.050000in}}{\pgfqpoint{4.900000in}{4.900000in}}%
\pgfusepath{clip}%
\pgfsetbuttcap%
\pgfsetroundjoin%
\definecolor{currentfill}{rgb}{0.998155,0.998155,0.998155}%
\pgfsetfillcolor{currentfill}%
\pgfsetfillopacity{0.900000}%
\pgfsetlinewidth{0.501875pt}%
\definecolor{currentstroke}{rgb}{0.200000,0.200000,0.200000}%
\pgfsetstrokecolor{currentstroke}%
\pgfsetdash{}{0pt}%
\pgfpathmoveto{\pgfqpoint{2.559553in}{2.138705in}}%
\pgfpathlineto{\pgfqpoint{2.602768in}{2.146510in}}%
\pgfpathlineto{\pgfqpoint{2.633124in}{2.128946in}}%
\pgfpathlineto{\pgfqpoint{2.589878in}{2.121781in}}%
\pgfpathlineto{\pgfqpoint{2.559553in}{2.138705in}}%
\pgfpathclose%
\pgfusepath{stroke,fill}%
\end{pgfscope}%
\begin{pgfscope}%
\pgfpathrectangle{\pgfqpoint{0.050000in}{0.050000in}}{\pgfqpoint{4.900000in}{4.900000in}}%
\pgfusepath{clip}%
\pgfsetbuttcap%
\pgfsetroundjoin%
\definecolor{currentfill}{rgb}{1.000000,1.000000,1.000000}%
\pgfsetfillcolor{currentfill}%
\pgfsetfillopacity{0.900000}%
\pgfsetlinewidth{0.501875pt}%
\definecolor{currentstroke}{rgb}{0.200000,0.200000,0.200000}%
\pgfsetstrokecolor{currentstroke}%
\pgfsetdash{}{0pt}%
\pgfpathmoveto{\pgfqpoint{2.823482in}{2.125604in}}%
\pgfpathlineto{\pgfqpoint{2.866430in}{2.138214in}}%
\pgfpathlineto{\pgfqpoint{2.897182in}{2.119495in}}%
\pgfpathlineto{\pgfqpoint{2.854203in}{2.107015in}}%
\pgfpathlineto{\pgfqpoint{2.823482in}{2.125604in}}%
\pgfpathclose%
\pgfusepath{stroke,fill}%
\end{pgfscope}%
\begin{pgfscope}%
\pgfpathrectangle{\pgfqpoint{0.050000in}{0.050000in}}{\pgfqpoint{4.900000in}{4.900000in}}%
\pgfusepath{clip}%
\pgfsetbuttcap%
\pgfsetroundjoin%
\definecolor{currentfill}{rgb}{1.000000,1.000000,1.000000}%
\pgfsetfillcolor{currentfill}%
\pgfsetfillopacity{0.900000}%
\pgfsetlinewidth{0.501875pt}%
\definecolor{currentstroke}{rgb}{0.200000,0.200000,0.200000}%
\pgfsetstrokecolor{currentstroke}%
\pgfsetdash{}{0pt}%
\pgfpathmoveto{\pgfqpoint{2.633124in}{2.128946in}}%
\pgfpathlineto{\pgfqpoint{2.676276in}{2.138844in}}%
\pgfpathlineto{\pgfqpoint{2.706754in}{2.120799in}}%
\pgfpathlineto{\pgfqpoint{2.663571in}{2.111382in}}%
\pgfpathlineto{\pgfqpoint{2.633124in}{2.128946in}}%
\pgfpathclose%
\pgfusepath{stroke,fill}%
\end{pgfscope}%
\begin{pgfscope}%
\pgfpathrectangle{\pgfqpoint{0.050000in}{0.050000in}}{\pgfqpoint{4.900000in}{4.900000in}}%
\pgfusepath{clip}%
\pgfsetbuttcap%
\pgfsetroundjoin%
\definecolor{currentfill}{rgb}{1.000000,1.000000,1.000000}%
\pgfsetfillcolor{currentfill}%
\pgfsetfillopacity{0.900000}%
\pgfsetlinewidth{0.501875pt}%
\definecolor{currentstroke}{rgb}{0.200000,0.200000,0.200000}%
\pgfsetstrokecolor{currentstroke}%
\pgfsetdash{}{0pt}%
\pgfpathmoveto{\pgfqpoint{3.087556in}{2.120816in}}%
\pgfpathlineto{\pgfqpoint{3.130245in}{2.133996in}}%
\pgfpathlineto{\pgfqpoint{3.161396in}{2.115086in}}%
\pgfpathlineto{\pgfqpoint{3.118676in}{2.101896in}}%
\pgfpathlineto{\pgfqpoint{3.087556in}{2.120816in}}%
\pgfpathclose%
\pgfusepath{stroke,fill}%
\end{pgfscope}%
\begin{pgfscope}%
\pgfpathrectangle{\pgfqpoint{0.050000in}{0.050000in}}{\pgfqpoint{4.900000in}{4.900000in}}%
\pgfusepath{clip}%
\pgfsetbuttcap%
\pgfsetroundjoin%
\definecolor{currentfill}{rgb}{1.000000,1.000000,1.000000}%
\pgfsetfillcolor{currentfill}%
\pgfsetfillopacity{0.900000}%
\pgfsetlinewidth{0.501875pt}%
\definecolor{currentstroke}{rgb}{0.200000,0.200000,0.200000}%
\pgfsetstrokecolor{currentstroke}%
\pgfsetdash{}{0pt}%
\pgfpathmoveto{\pgfqpoint{2.897182in}{2.119495in}}%
\pgfpathlineto{\pgfqpoint{2.940073in}{2.132357in}}%
\pgfpathlineto{\pgfqpoint{2.970948in}{2.113547in}}%
\pgfpathlineto{\pgfqpoint{2.928026in}{2.100757in}}%
\pgfpathlineto{\pgfqpoint{2.897182in}{2.119495in}}%
\pgfpathclose%
\pgfusepath{stroke,fill}%
\end{pgfscope}%
\begin{pgfscope}%
\pgfpathrectangle{\pgfqpoint{0.050000in}{0.050000in}}{\pgfqpoint{4.900000in}{4.900000in}}%
\pgfusepath{clip}%
\pgfsetbuttcap%
\pgfsetroundjoin%
\definecolor{currentfill}{rgb}{0.985236,0.985236,0.985236}%
\pgfsetfillcolor{currentfill}%
\pgfsetfillopacity{0.900000}%
\pgfsetlinewidth{0.501875pt}%
\definecolor{currentstroke}{rgb}{0.200000,0.200000,0.200000}%
\pgfsetstrokecolor{currentstroke}%
\pgfsetdash{}{0pt}%
\pgfpathmoveto{\pgfqpoint{2.295410in}{2.210443in}}%
\pgfpathlineto{\pgfqpoint{2.339072in}{2.191593in}}%
\pgfpathlineto{\pgfqpoint{2.369033in}{2.176984in}}%
\pgfpathlineto{\pgfqpoint{2.325355in}{2.195256in}}%
\pgfpathlineto{\pgfqpoint{2.295410in}{2.210443in}}%
\pgfpathclose%
\pgfusepath{stroke,fill}%
\end{pgfscope}%
\begin{pgfscope}%
\pgfpathrectangle{\pgfqpoint{0.050000in}{0.050000in}}{\pgfqpoint{4.900000in}{4.900000in}}%
\pgfusepath{clip}%
\pgfsetbuttcap%
\pgfsetroundjoin%
\definecolor{currentfill}{rgb}{0.990773,0.990773,0.990773}%
\pgfsetfillcolor{currentfill}%
\pgfsetfillopacity{0.900000}%
\pgfsetlinewidth{0.501875pt}%
\definecolor{currentstroke}{rgb}{0.200000,0.200000,0.200000}%
\pgfsetstrokecolor{currentstroke}%
\pgfsetdash{}{0pt}%
\pgfpathmoveto{\pgfqpoint{2.369033in}{2.176984in}}%
\pgfpathlineto{\pgfqpoint{2.412548in}{2.167936in}}%
\pgfpathlineto{\pgfqpoint{2.442629in}{2.152695in}}%
\pgfpathlineto{\pgfqpoint{2.399089in}{2.161949in}}%
\pgfpathlineto{\pgfqpoint{2.369033in}{2.176984in}}%
\pgfpathclose%
\pgfusepath{stroke,fill}%
\end{pgfscope}%
\begin{pgfscope}%
\pgfpathrectangle{\pgfqpoint{0.050000in}{0.050000in}}{\pgfqpoint{4.900000in}{4.900000in}}%
\pgfusepath{clip}%
\pgfsetbuttcap%
\pgfsetroundjoin%
\definecolor{currentfill}{rgb}{1.000000,1.000000,1.000000}%
\pgfsetfillcolor{currentfill}%
\pgfsetfillopacity{0.900000}%
\pgfsetlinewidth{0.501875pt}%
\definecolor{currentstroke}{rgb}{0.200000,0.200000,0.200000}%
\pgfsetstrokecolor{currentstroke}%
\pgfsetdash{}{0pt}%
\pgfpathmoveto{\pgfqpoint{2.706754in}{2.120799in}}%
\pgfpathlineto{\pgfqpoint{2.749847in}{2.131988in}}%
\pgfpathlineto{\pgfqpoint{2.780446in}{2.113618in}}%
\pgfpathlineto{\pgfqpoint{2.737322in}{2.102766in}}%
\pgfpathlineto{\pgfqpoint{2.706754in}{2.120799in}}%
\pgfpathclose%
\pgfusepath{stroke,fill}%
\end{pgfscope}%
\begin{pgfscope}%
\pgfpathrectangle{\pgfqpoint{0.050000in}{0.050000in}}{\pgfqpoint{4.900000in}{4.900000in}}%
\pgfusepath{clip}%
\pgfsetbuttcap%
\pgfsetroundjoin%
\definecolor{currentfill}{rgb}{0.976009,0.976009,0.976009}%
\pgfsetfillcolor{currentfill}%
\pgfsetfillopacity{0.900000}%
\pgfsetlinewidth{0.501875pt}%
\definecolor{currentstroke}{rgb}{0.200000,0.200000,0.200000}%
\pgfsetstrokecolor{currentstroke}%
\pgfsetdash{}{0pt}%
\pgfpathmoveto{\pgfqpoint{2.221713in}{2.254979in}}%
\pgfpathlineto{\pgfqpoint{2.265560in}{2.225257in}}%
\pgfpathlineto{\pgfqpoint{2.295410in}{2.210443in}}%
\pgfpathlineto{\pgfqpoint{2.251557in}{2.238630in}}%
\pgfpathlineto{\pgfqpoint{2.221713in}{2.254979in}}%
\pgfpathclose%
\pgfusepath{stroke,fill}%
\end{pgfscope}%
\begin{pgfscope}%
\pgfpathrectangle{\pgfqpoint{0.050000in}{0.050000in}}{\pgfqpoint{4.900000in}{4.900000in}}%
\pgfusepath{clip}%
\pgfsetbuttcap%
\pgfsetroundjoin%
\definecolor{currentfill}{rgb}{1.000000,1.000000,1.000000}%
\pgfsetfillcolor{currentfill}%
\pgfsetfillopacity{0.900000}%
\pgfsetlinewidth{0.501875pt}%
\definecolor{currentstroke}{rgb}{0.200000,0.200000,0.200000}%
\pgfsetstrokecolor{currentstroke}%
\pgfsetdash{}{0pt}%
\pgfpathmoveto{\pgfqpoint{3.161396in}{2.115086in}}%
\pgfpathlineto{\pgfqpoint{3.204028in}{2.128300in}}%
\pgfpathlineto{\pgfqpoint{3.235302in}{2.109361in}}%
\pgfpathlineto{\pgfqpoint{3.192639in}{2.096128in}}%
\pgfpathlineto{\pgfqpoint{3.161396in}{2.115086in}}%
\pgfpathclose%
\pgfusepath{stroke,fill}%
\end{pgfscope}%
\begin{pgfscope}%
\pgfpathrectangle{\pgfqpoint{0.050000in}{0.050000in}}{\pgfqpoint{4.900000in}{4.900000in}}%
\pgfusepath{clip}%
\pgfsetbuttcap%
\pgfsetroundjoin%
\definecolor{currentfill}{rgb}{0.994464,0.994464,0.994464}%
\pgfsetfillcolor{currentfill}%
\pgfsetfillopacity{0.900000}%
\pgfsetlinewidth{0.501875pt}%
\definecolor{currentstroke}{rgb}{0.200000,0.200000,0.200000}%
\pgfsetstrokecolor{currentstroke}%
\pgfsetdash{}{0pt}%
\pgfpathmoveto{\pgfqpoint{2.442629in}{2.152695in}}%
\pgfpathlineto{\pgfqpoint{2.486032in}{2.151139in}}%
\pgfpathlineto{\pgfqpoint{2.516235in}{2.135022in}}%
\pgfpathlineto{\pgfqpoint{2.472803in}{2.137159in}}%
\pgfpathlineto{\pgfqpoint{2.442629in}{2.152695in}}%
\pgfpathclose%
\pgfusepath{stroke,fill}%
\end{pgfscope}%
\begin{pgfscope}%
\pgfpathrectangle{\pgfqpoint{0.050000in}{0.050000in}}{\pgfqpoint{4.900000in}{4.900000in}}%
\pgfusepath{clip}%
\pgfsetbuttcap%
\pgfsetroundjoin%
\definecolor{currentfill}{rgb}{1.000000,1.000000,1.000000}%
\pgfsetfillcolor{currentfill}%
\pgfsetfillopacity{0.900000}%
\pgfsetlinewidth{0.501875pt}%
\definecolor{currentstroke}{rgb}{0.200000,0.200000,0.200000}%
\pgfsetstrokecolor{currentstroke}%
\pgfsetdash{}{0pt}%
\pgfpathmoveto{\pgfqpoint{2.970948in}{2.113547in}}%
\pgfpathlineto{\pgfqpoint{3.013782in}{2.126567in}}%
\pgfpathlineto{\pgfqpoint{3.044779in}{2.107693in}}%
\pgfpathlineto{\pgfqpoint{3.001914in}{2.094707in}}%
\pgfpathlineto{\pgfqpoint{2.970948in}{2.113547in}}%
\pgfpathclose%
\pgfusepath{stroke,fill}%
\end{pgfscope}%
\begin{pgfscope}%
\pgfpathrectangle{\pgfqpoint{0.050000in}{0.050000in}}{\pgfqpoint{4.900000in}{4.900000in}}%
\pgfusepath{clip}%
\pgfsetbuttcap%
\pgfsetroundjoin%
\definecolor{currentfill}{rgb}{0.964937,0.964937,0.964937}%
\pgfsetfillcolor{currentfill}%
\pgfsetfillopacity{0.900000}%
\pgfsetlinewidth{0.501875pt}%
\definecolor{currentstroke}{rgb}{0.200000,0.200000,0.200000}%
\pgfsetstrokecolor{currentstroke}%
\pgfsetdash{}{0pt}%
\pgfpathmoveto{\pgfqpoint{2.147903in}{2.309807in}}%
\pgfpathlineto{\pgfqpoint{2.191959in}{2.271387in}}%
\pgfpathlineto{\pgfqpoint{2.221713in}{2.254979in}}%
\pgfpathlineto{\pgfqpoint{2.177660in}{2.291655in}}%
\pgfpathlineto{\pgfqpoint{2.147903in}{2.309807in}}%
\pgfpathclose%
\pgfusepath{stroke,fill}%
\end{pgfscope}%
\begin{pgfscope}%
\pgfpathrectangle{\pgfqpoint{0.050000in}{0.050000in}}{\pgfqpoint{4.900000in}{4.900000in}}%
\pgfusepath{clip}%
\pgfsetbuttcap%
\pgfsetroundjoin%
\definecolor{currentfill}{rgb}{0.996309,0.996309,0.996309}%
\pgfsetfillcolor{currentfill}%
\pgfsetfillopacity{0.900000}%
\pgfsetlinewidth{0.501875pt}%
\definecolor{currentstroke}{rgb}{0.200000,0.200000,0.200000}%
\pgfsetstrokecolor{currentstroke}%
\pgfsetdash{}{0pt}%
\pgfpathmoveto{\pgfqpoint{2.516235in}{2.135022in}}%
\pgfpathlineto{\pgfqpoint{2.559553in}{2.138705in}}%
\pgfpathlineto{\pgfqpoint{2.589878in}{2.121781in}}%
\pgfpathlineto{\pgfqpoint{2.546531in}{2.118753in}}%
\pgfpathlineto{\pgfqpoint{2.516235in}{2.135022in}}%
\pgfpathclose%
\pgfusepath{stroke,fill}%
\end{pgfscope}%
\begin{pgfscope}%
\pgfpathrectangle{\pgfqpoint{0.050000in}{0.050000in}}{\pgfqpoint{4.900000in}{4.900000in}}%
\pgfusepath{clip}%
\pgfsetbuttcap%
\pgfsetroundjoin%
\definecolor{currentfill}{rgb}{1.000000,1.000000,1.000000}%
\pgfsetfillcolor{currentfill}%
\pgfsetfillopacity{0.900000}%
\pgfsetlinewidth{0.501875pt}%
\definecolor{currentstroke}{rgb}{0.200000,0.200000,0.200000}%
\pgfsetstrokecolor{currentstroke}%
\pgfsetdash{}{0pt}%
\pgfpathmoveto{\pgfqpoint{2.780446in}{2.113618in}}%
\pgfpathlineto{\pgfqpoint{2.823482in}{2.125604in}}%
\pgfpathlineto{\pgfqpoint{2.854203in}{2.107015in}}%
\pgfpathlineto{\pgfqpoint{2.811137in}{2.095252in}}%
\pgfpathlineto{\pgfqpoint{2.780446in}{2.113618in}}%
\pgfpathclose%
\pgfusepath{stroke,fill}%
\end{pgfscope}%
\begin{pgfscope}%
\pgfpathrectangle{\pgfqpoint{0.050000in}{0.050000in}}{\pgfqpoint{4.900000in}{4.900000in}}%
\pgfusepath{clip}%
\pgfsetbuttcap%
\pgfsetroundjoin%
\definecolor{currentfill}{rgb}{1.000000,1.000000,1.000000}%
\pgfsetfillcolor{currentfill}%
\pgfsetfillopacity{0.900000}%
\pgfsetlinewidth{0.501875pt}%
\definecolor{currentstroke}{rgb}{0.200000,0.200000,0.200000}%
\pgfsetstrokecolor{currentstroke}%
\pgfsetdash{}{0pt}%
\pgfpathmoveto{\pgfqpoint{3.235302in}{2.109361in}}%
\pgfpathlineto{\pgfqpoint{3.277876in}{2.122594in}}%
\pgfpathlineto{\pgfqpoint{3.309272in}{2.103625in}}%
\pgfpathlineto{\pgfqpoint{3.266668in}{2.090370in}}%
\pgfpathlineto{\pgfqpoint{3.235302in}{2.109361in}}%
\pgfpathclose%
\pgfusepath{stroke,fill}%
\end{pgfscope}%
\begin{pgfscope}%
\pgfpathrectangle{\pgfqpoint{0.050000in}{0.050000in}}{\pgfqpoint{4.900000in}{4.900000in}}%
\pgfusepath{clip}%
\pgfsetbuttcap%
\pgfsetroundjoin%
\definecolor{currentfill}{rgb}{1.000000,1.000000,1.000000}%
\pgfsetfillcolor{currentfill}%
\pgfsetfillopacity{0.900000}%
\pgfsetlinewidth{0.501875pt}%
\definecolor{currentstroke}{rgb}{0.200000,0.200000,0.200000}%
\pgfsetstrokecolor{currentstroke}%
\pgfsetdash{}{0pt}%
\pgfpathmoveto{\pgfqpoint{3.044779in}{2.107693in}}%
\pgfpathlineto{\pgfqpoint{3.087556in}{2.120816in}}%
\pgfpathlineto{\pgfqpoint{3.118676in}{2.101896in}}%
\pgfpathlineto{\pgfqpoint{3.075868in}{2.088782in}}%
\pgfpathlineto{\pgfqpoint{3.044779in}{2.107693in}}%
\pgfpathclose%
\pgfusepath{stroke,fill}%
\end{pgfscope}%
\begin{pgfscope}%
\pgfpathrectangle{\pgfqpoint{0.050000in}{0.050000in}}{\pgfqpoint{4.900000in}{4.900000in}}%
\pgfusepath{clip}%
\pgfsetbuttcap%
\pgfsetroundjoin%
\definecolor{currentfill}{rgb}{0.998155,0.998155,0.998155}%
\pgfsetfillcolor{currentfill}%
\pgfsetfillopacity{0.900000}%
\pgfsetlinewidth{0.501875pt}%
\definecolor{currentstroke}{rgb}{0.200000,0.200000,0.200000}%
\pgfsetstrokecolor{currentstroke}%
\pgfsetdash{}{0pt}%
\pgfpathmoveto{\pgfqpoint{2.589878in}{2.121781in}}%
\pgfpathlineto{\pgfqpoint{2.633124in}{2.128946in}}%
\pgfpathlineto{\pgfqpoint{2.663571in}{2.111382in}}%
\pgfpathlineto{\pgfqpoint{2.620295in}{2.104795in}}%
\pgfpathlineto{\pgfqpoint{2.589878in}{2.121781in}}%
\pgfpathclose%
\pgfusepath{stroke,fill}%
\end{pgfscope}%
\begin{pgfscope}%
\pgfpathrectangle{\pgfqpoint{0.050000in}{0.050000in}}{\pgfqpoint{4.900000in}{4.900000in}}%
\pgfusepath{clip}%
\pgfsetbuttcap%
\pgfsetroundjoin%
\definecolor{currentfill}{rgb}{1.000000,1.000000,1.000000}%
\pgfsetfillcolor{currentfill}%
\pgfsetfillopacity{0.900000}%
\pgfsetlinewidth{0.501875pt}%
\definecolor{currentstroke}{rgb}{0.200000,0.200000,0.200000}%
\pgfsetstrokecolor{currentstroke}%
\pgfsetdash{}{0pt}%
\pgfpathmoveto{\pgfqpoint{2.854203in}{2.107015in}}%
\pgfpathlineto{\pgfqpoint{2.897182in}{2.119495in}}%
\pgfpathlineto{\pgfqpoint{2.928026in}{2.100757in}}%
\pgfpathlineto{\pgfqpoint{2.885016in}{2.088418in}}%
\pgfpathlineto{\pgfqpoint{2.854203in}{2.107015in}}%
\pgfpathclose%
\pgfusepath{stroke,fill}%
\end{pgfscope}%
\begin{pgfscope}%
\pgfpathrectangle{\pgfqpoint{0.050000in}{0.050000in}}{\pgfqpoint{4.900000in}{4.900000in}}%
\pgfusepath{clip}%
\pgfsetbuttcap%
\pgfsetroundjoin%
\definecolor{currentfill}{rgb}{1.000000,1.000000,1.000000}%
\pgfsetfillcolor{currentfill}%
\pgfsetfillopacity{0.900000}%
\pgfsetlinewidth{0.501875pt}%
\definecolor{currentstroke}{rgb}{0.200000,0.200000,0.200000}%
\pgfsetstrokecolor{currentstroke}%
\pgfsetdash{}{0pt}%
\pgfpathmoveto{\pgfqpoint{2.663571in}{2.111382in}}%
\pgfpathlineto{\pgfqpoint{2.706754in}{2.120799in}}%
\pgfpathlineto{\pgfqpoint{2.737322in}{2.102766in}}%
\pgfpathlineto{\pgfqpoint{2.694109in}{2.093801in}}%
\pgfpathlineto{\pgfqpoint{2.663571in}{2.111382in}}%
\pgfpathclose%
\pgfusepath{stroke,fill}%
\end{pgfscope}%
\begin{pgfscope}%
\pgfpathrectangle{\pgfqpoint{0.050000in}{0.050000in}}{\pgfqpoint{4.900000in}{4.900000in}}%
\pgfusepath{clip}%
\pgfsetbuttcap%
\pgfsetroundjoin%
\definecolor{currentfill}{rgb}{1.000000,1.000000,1.000000}%
\pgfsetfillcolor{currentfill}%
\pgfsetfillopacity{0.900000}%
\pgfsetlinewidth{0.501875pt}%
\definecolor{currentstroke}{rgb}{0.200000,0.200000,0.200000}%
\pgfsetstrokecolor{currentstroke}%
\pgfsetdash{}{0pt}%
\pgfpathmoveto{\pgfqpoint{3.118676in}{2.101896in}}%
\pgfpathlineto{\pgfqpoint{3.161396in}{2.115086in}}%
\pgfpathlineto{\pgfqpoint{3.192639in}{2.096128in}}%
\pgfpathlineto{\pgfqpoint{3.149889in}{2.082931in}}%
\pgfpathlineto{\pgfqpoint{3.118676in}{2.101896in}}%
\pgfpathclose%
\pgfusepath{stroke,fill}%
\end{pgfscope}%
\begin{pgfscope}%
\pgfpathrectangle{\pgfqpoint{0.050000in}{0.050000in}}{\pgfqpoint{4.900000in}{4.900000in}}%
\pgfusepath{clip}%
\pgfsetbuttcap%
\pgfsetroundjoin%
\definecolor{currentfill}{rgb}{1.000000,1.000000,1.000000}%
\pgfsetfillcolor{currentfill}%
\pgfsetfillopacity{0.900000}%
\pgfsetlinewidth{0.501875pt}%
\definecolor{currentstroke}{rgb}{0.200000,0.200000,0.200000}%
\pgfsetstrokecolor{currentstroke}%
\pgfsetdash{}{0pt}%
\pgfpathmoveto{\pgfqpoint{2.928026in}{2.100757in}}%
\pgfpathlineto{\pgfqpoint{2.970948in}{2.113547in}}%
\pgfpathlineto{\pgfqpoint{3.001914in}{2.094707in}}%
\pgfpathlineto{\pgfqpoint{2.958961in}{2.082001in}}%
\pgfpathlineto{\pgfqpoint{2.928026in}{2.100757in}}%
\pgfpathclose%
\pgfusepath{stroke,fill}%
\end{pgfscope}%
\begin{pgfscope}%
\pgfpathrectangle{\pgfqpoint{0.050000in}{0.050000in}}{\pgfqpoint{4.900000in}{4.900000in}}%
\pgfusepath{clip}%
\pgfsetbuttcap%
\pgfsetroundjoin%
\definecolor{currentfill}{rgb}{0.983391,0.983391,0.983391}%
\pgfsetfillcolor{currentfill}%
\pgfsetfillopacity{0.900000}%
\pgfsetlinewidth{0.501875pt}%
\definecolor{currentstroke}{rgb}{0.200000,0.200000,0.200000}%
\pgfsetstrokecolor{currentstroke}%
\pgfsetdash{}{0pt}%
\pgfpathmoveto{\pgfqpoint{2.325355in}{2.195256in}}%
\pgfpathlineto{\pgfqpoint{2.369033in}{2.176984in}}%
\pgfpathlineto{\pgfqpoint{2.399089in}{2.161949in}}%
\pgfpathlineto{\pgfqpoint{2.355393in}{2.179756in}}%
\pgfpathlineto{\pgfqpoint{2.325355in}{2.195256in}}%
\pgfpathclose%
\pgfusepath{stroke,fill}%
\end{pgfscope}%
\begin{pgfscope}%
\pgfpathrectangle{\pgfqpoint{0.050000in}{0.050000in}}{\pgfqpoint{4.900000in}{4.900000in}}%
\pgfusepath{clip}%
\pgfsetbuttcap%
\pgfsetroundjoin%
\definecolor{currentfill}{rgb}{0.988927,0.988927,0.988927}%
\pgfsetfillcolor{currentfill}%
\pgfsetfillopacity{0.900000}%
\pgfsetlinewidth{0.501875pt}%
\definecolor{currentstroke}{rgb}{0.200000,0.200000,0.200000}%
\pgfsetstrokecolor{currentstroke}%
\pgfsetdash{}{0pt}%
\pgfpathmoveto{\pgfqpoint{2.399089in}{2.161949in}}%
\pgfpathlineto{\pgfqpoint{2.442629in}{2.152695in}}%
\pgfpathlineto{\pgfqpoint{2.472803in}{2.137159in}}%
\pgfpathlineto{\pgfqpoint{2.429240in}{2.146566in}}%
\pgfpathlineto{\pgfqpoint{2.399089in}{2.161949in}}%
\pgfpathclose%
\pgfusepath{stroke,fill}%
\end{pgfscope}%
\begin{pgfscope}%
\pgfpathrectangle{\pgfqpoint{0.050000in}{0.050000in}}{\pgfqpoint{4.900000in}{4.900000in}}%
\pgfusepath{clip}%
\pgfsetbuttcap%
\pgfsetroundjoin%
\definecolor{currentfill}{rgb}{1.000000,1.000000,1.000000}%
\pgfsetfillcolor{currentfill}%
\pgfsetfillopacity{0.900000}%
\pgfsetlinewidth{0.501875pt}%
\definecolor{currentstroke}{rgb}{0.200000,0.200000,0.200000}%
\pgfsetstrokecolor{currentstroke}%
\pgfsetdash{}{0pt}%
\pgfpathmoveto{\pgfqpoint{2.737322in}{2.102766in}}%
\pgfpathlineto{\pgfqpoint{2.780446in}{2.113618in}}%
\pgfpathlineto{\pgfqpoint{2.811137in}{2.095252in}}%
\pgfpathlineto{\pgfqpoint{2.767983in}{2.084729in}}%
\pgfpathlineto{\pgfqpoint{2.737322in}{2.102766in}}%
\pgfpathclose%
\pgfusepath{stroke,fill}%
\end{pgfscope}%
\begin{pgfscope}%
\pgfpathrectangle{\pgfqpoint{0.050000in}{0.050000in}}{\pgfqpoint{4.900000in}{4.900000in}}%
\pgfusepath{clip}%
\pgfsetbuttcap%
\pgfsetroundjoin%
\definecolor{currentfill}{rgb}{1.000000,1.000000,1.000000}%
\pgfsetfillcolor{currentfill}%
\pgfsetfillopacity{0.900000}%
\pgfsetlinewidth{0.501875pt}%
\definecolor{currentstroke}{rgb}{0.200000,0.200000,0.200000}%
\pgfsetstrokecolor{currentstroke}%
\pgfsetdash{}{0pt}%
\pgfpathmoveto{\pgfqpoint{3.192639in}{2.096128in}}%
\pgfpathlineto{\pgfqpoint{3.235302in}{2.109361in}}%
\pgfpathlineto{\pgfqpoint{3.266668in}{2.090370in}}%
\pgfpathlineto{\pgfqpoint{3.223975in}{2.077121in}}%
\pgfpathlineto{\pgfqpoint{3.192639in}{2.096128in}}%
\pgfpathclose%
\pgfusepath{stroke,fill}%
\end{pgfscope}%
\begin{pgfscope}%
\pgfpathrectangle{\pgfqpoint{0.050000in}{0.050000in}}{\pgfqpoint{4.900000in}{4.900000in}}%
\pgfusepath{clip}%
\pgfsetbuttcap%
\pgfsetroundjoin%
\definecolor{currentfill}{rgb}{0.976009,0.976009,0.976009}%
\pgfsetfillcolor{currentfill}%
\pgfsetfillopacity{0.900000}%
\pgfsetlinewidth{0.501875pt}%
\definecolor{currentstroke}{rgb}{0.200000,0.200000,0.200000}%
\pgfsetstrokecolor{currentstroke}%
\pgfsetdash{}{0pt}%
\pgfpathmoveto{\pgfqpoint{2.251557in}{2.238630in}}%
\pgfpathlineto{\pgfqpoint{2.295410in}{2.210443in}}%
\pgfpathlineto{\pgfqpoint{2.325355in}{2.195256in}}%
\pgfpathlineto{\pgfqpoint{2.281492in}{2.222237in}}%
\pgfpathlineto{\pgfqpoint{2.251557in}{2.238630in}}%
\pgfpathclose%
\pgfusepath{stroke,fill}%
\end{pgfscope}%
\begin{pgfscope}%
\pgfpathrectangle{\pgfqpoint{0.050000in}{0.050000in}}{\pgfqpoint{4.900000in}{4.900000in}}%
\pgfusepath{clip}%
\pgfsetbuttcap%
\pgfsetroundjoin%
\definecolor{currentfill}{rgb}{0.994464,0.994464,0.994464}%
\pgfsetfillcolor{currentfill}%
\pgfsetfillopacity{0.900000}%
\pgfsetlinewidth{0.501875pt}%
\definecolor{currentstroke}{rgb}{0.200000,0.200000,0.200000}%
\pgfsetstrokecolor{currentstroke}%
\pgfsetdash{}{0pt}%
\pgfpathmoveto{\pgfqpoint{2.472803in}{2.137159in}}%
\pgfpathlineto{\pgfqpoint{2.516235in}{2.135022in}}%
\pgfpathlineto{\pgfqpoint{2.546531in}{2.118753in}}%
\pgfpathlineto{\pgfqpoint{2.503071in}{2.121368in}}%
\pgfpathlineto{\pgfqpoint{2.472803in}{2.137159in}}%
\pgfpathclose%
\pgfusepath{stroke,fill}%
\end{pgfscope}%
\begin{pgfscope}%
\pgfpathrectangle{\pgfqpoint{0.050000in}{0.050000in}}{\pgfqpoint{4.900000in}{4.900000in}}%
\pgfusepath{clip}%
\pgfsetbuttcap%
\pgfsetroundjoin%
\definecolor{currentfill}{rgb}{1.000000,1.000000,1.000000}%
\pgfsetfillcolor{currentfill}%
\pgfsetfillopacity{0.900000}%
\pgfsetlinewidth{0.501875pt}%
\definecolor{currentstroke}{rgb}{0.200000,0.200000,0.200000}%
\pgfsetstrokecolor{currentstroke}%
\pgfsetdash{}{0pt}%
\pgfpathmoveto{\pgfqpoint{3.001914in}{2.094707in}}%
\pgfpathlineto{\pgfqpoint{3.044779in}{2.107693in}}%
\pgfpathlineto{\pgfqpoint{3.075868in}{2.088782in}}%
\pgfpathlineto{\pgfqpoint{3.032973in}{2.075839in}}%
\pgfpathlineto{\pgfqpoint{3.001914in}{2.094707in}}%
\pgfpathclose%
\pgfusepath{stroke,fill}%
\end{pgfscope}%
\begin{pgfscope}%
\pgfpathrectangle{\pgfqpoint{0.050000in}{0.050000in}}{\pgfqpoint{4.900000in}{4.900000in}}%
\pgfusepath{clip}%
\pgfsetbuttcap%
\pgfsetroundjoin%
\definecolor{currentfill}{rgb}{0.964937,0.964937,0.964937}%
\pgfsetfillcolor{currentfill}%
\pgfsetfillopacity{0.900000}%
\pgfsetlinewidth{0.501875pt}%
\definecolor{currentstroke}{rgb}{0.200000,0.200000,0.200000}%
\pgfsetstrokecolor{currentstroke}%
\pgfsetdash{}{0pt}%
\pgfpathmoveto{\pgfqpoint{2.177660in}{2.291655in}}%
\pgfpathlineto{\pgfqpoint{2.221713in}{2.254979in}}%
\pgfpathlineto{\pgfqpoint{2.251557in}{2.238630in}}%
\pgfpathlineto{\pgfqpoint{2.207504in}{2.273790in}}%
\pgfpathlineto{\pgfqpoint{2.177660in}{2.291655in}}%
\pgfpathclose%
\pgfusepath{stroke,fill}%
\end{pgfscope}%
\begin{pgfscope}%
\pgfpathrectangle{\pgfqpoint{0.050000in}{0.050000in}}{\pgfqpoint{4.900000in}{4.900000in}}%
\pgfusepath{clip}%
\pgfsetbuttcap%
\pgfsetroundjoin%
\definecolor{currentfill}{rgb}{0.996309,0.996309,0.996309}%
\pgfsetfillcolor{currentfill}%
\pgfsetfillopacity{0.900000}%
\pgfsetlinewidth{0.501875pt}%
\definecolor{currentstroke}{rgb}{0.200000,0.200000,0.200000}%
\pgfsetstrokecolor{currentstroke}%
\pgfsetdash{}{0pt}%
\pgfpathmoveto{\pgfqpoint{2.546531in}{2.118753in}}%
\pgfpathlineto{\pgfqpoint{2.589878in}{2.121781in}}%
\pgfpathlineto{\pgfqpoint{2.620295in}{2.104795in}}%
\pgfpathlineto{\pgfqpoint{2.576919in}{2.102335in}}%
\pgfpathlineto{\pgfqpoint{2.546531in}{2.118753in}}%
\pgfpathclose%
\pgfusepath{stroke,fill}%
\end{pgfscope}%
\begin{pgfscope}%
\pgfpathrectangle{\pgfqpoint{0.050000in}{0.050000in}}{\pgfqpoint{4.900000in}{4.900000in}}%
\pgfusepath{clip}%
\pgfsetbuttcap%
\pgfsetroundjoin%
\definecolor{currentfill}{rgb}{1.000000,1.000000,1.000000}%
\pgfsetfillcolor{currentfill}%
\pgfsetfillopacity{0.900000}%
\pgfsetlinewidth{0.501875pt}%
\definecolor{currentstroke}{rgb}{0.200000,0.200000,0.200000}%
\pgfsetstrokecolor{currentstroke}%
\pgfsetdash{}{0pt}%
\pgfpathmoveto{\pgfqpoint{2.811137in}{2.095252in}}%
\pgfpathlineto{\pgfqpoint{2.854203in}{2.107015in}}%
\pgfpathlineto{\pgfqpoint{2.885016in}{2.088418in}}%
\pgfpathlineto{\pgfqpoint{2.841920in}{2.076883in}}%
\pgfpathlineto{\pgfqpoint{2.811137in}{2.095252in}}%
\pgfpathclose%
\pgfusepath{stroke,fill}%
\end{pgfscope}%
\begin{pgfscope}%
\pgfpathrectangle{\pgfqpoint{0.050000in}{0.050000in}}{\pgfqpoint{4.900000in}{4.900000in}}%
\pgfusepath{clip}%
\pgfsetbuttcap%
\pgfsetroundjoin%
\definecolor{currentfill}{rgb}{1.000000,1.000000,1.000000}%
\pgfsetfillcolor{currentfill}%
\pgfsetfillopacity{0.900000}%
\pgfsetlinewidth{0.501875pt}%
\definecolor{currentstroke}{rgb}{0.200000,0.200000,0.200000}%
\pgfsetstrokecolor{currentstroke}%
\pgfsetdash{}{0pt}%
\pgfpathmoveto{\pgfqpoint{3.266668in}{2.090370in}}%
\pgfpathlineto{\pgfqpoint{3.309272in}{2.103625in}}%
\pgfpathlineto{\pgfqpoint{3.340763in}{2.084599in}}%
\pgfpathlineto{\pgfqpoint{3.298128in}{2.071326in}}%
\pgfpathlineto{\pgfqpoint{3.266668in}{2.090370in}}%
\pgfpathclose%
\pgfusepath{stroke,fill}%
\end{pgfscope}%
\begin{pgfscope}%
\pgfpathrectangle{\pgfqpoint{0.050000in}{0.050000in}}{\pgfqpoint{4.900000in}{4.900000in}}%
\pgfusepath{clip}%
\pgfsetbuttcap%
\pgfsetroundjoin%
\definecolor{currentfill}{rgb}{1.000000,1.000000,1.000000}%
\pgfsetfillcolor{currentfill}%
\pgfsetfillopacity{0.900000}%
\pgfsetlinewidth{0.501875pt}%
\definecolor{currentstroke}{rgb}{0.200000,0.200000,0.200000}%
\pgfsetstrokecolor{currentstroke}%
\pgfsetdash{}{0pt}%
\pgfpathmoveto{\pgfqpoint{3.075868in}{2.088782in}}%
\pgfpathlineto{\pgfqpoint{3.118676in}{2.101896in}}%
\pgfpathlineto{\pgfqpoint{3.149889in}{2.082931in}}%
\pgfpathlineto{\pgfqpoint{3.107050in}{2.069834in}}%
\pgfpathlineto{\pgfqpoint{3.075868in}{2.088782in}}%
\pgfpathclose%
\pgfusepath{stroke,fill}%
\end{pgfscope}%
\begin{pgfscope}%
\pgfpathrectangle{\pgfqpoint{0.050000in}{0.050000in}}{\pgfqpoint{4.900000in}{4.900000in}}%
\pgfusepath{clip}%
\pgfsetbuttcap%
\pgfsetroundjoin%
\definecolor{currentfill}{rgb}{0.998155,0.998155,0.998155}%
\pgfsetfillcolor{currentfill}%
\pgfsetfillopacity{0.900000}%
\pgfsetlinewidth{0.501875pt}%
\definecolor{currentstroke}{rgb}{0.200000,0.200000,0.200000}%
\pgfsetstrokecolor{currentstroke}%
\pgfsetdash{}{0pt}%
\pgfpathmoveto{\pgfqpoint{2.620295in}{2.104795in}}%
\pgfpathlineto{\pgfqpoint{2.663571in}{2.111382in}}%
\pgfpathlineto{\pgfqpoint{2.694109in}{2.093801in}}%
\pgfpathlineto{\pgfqpoint{2.650804in}{2.087738in}}%
\pgfpathlineto{\pgfqpoint{2.620295in}{2.104795in}}%
\pgfpathclose%
\pgfusepath{stroke,fill}%
\end{pgfscope}%
\begin{pgfscope}%
\pgfpathrectangle{\pgfqpoint{0.050000in}{0.050000in}}{\pgfqpoint{4.900000in}{4.900000in}}%
\pgfusepath{clip}%
\pgfsetbuttcap%
\pgfsetroundjoin%
\definecolor{currentfill}{rgb}{0.952018,0.952018,0.952018}%
\pgfsetfillcolor{currentfill}%
\pgfsetfillopacity{0.900000}%
\pgfsetlinewidth{0.501875pt}%
\definecolor{currentstroke}{rgb}{0.200000,0.200000,0.200000}%
\pgfsetstrokecolor{currentstroke}%
\pgfsetdash{}{0pt}%
\pgfpathmoveto{\pgfqpoint{2.103654in}{2.351030in}}%
\pgfpathlineto{\pgfqpoint{2.147903in}{2.309807in}}%
\pgfpathlineto{\pgfqpoint{2.177660in}{2.291655in}}%
\pgfpathlineto{\pgfqpoint{2.133411in}{2.331983in}}%
\pgfpathlineto{\pgfqpoint{2.103654in}{2.351030in}}%
\pgfpathclose%
\pgfusepath{stroke,fill}%
\end{pgfscope}%
\begin{pgfscope}%
\pgfpathrectangle{\pgfqpoint{0.050000in}{0.050000in}}{\pgfqpoint{4.900000in}{4.900000in}}%
\pgfusepath{clip}%
\pgfsetbuttcap%
\pgfsetroundjoin%
\definecolor{currentfill}{rgb}{1.000000,1.000000,1.000000}%
\pgfsetfillcolor{currentfill}%
\pgfsetfillopacity{0.900000}%
\pgfsetlinewidth{0.501875pt}%
\definecolor{currentstroke}{rgb}{0.200000,0.200000,0.200000}%
\pgfsetstrokecolor{currentstroke}%
\pgfsetdash{}{0pt}%
\pgfpathmoveto{\pgfqpoint{2.885016in}{2.088418in}}%
\pgfpathlineto{\pgfqpoint{2.928026in}{2.100757in}}%
\pgfpathlineto{\pgfqpoint{2.958961in}{2.082001in}}%
\pgfpathlineto{\pgfqpoint{2.915922in}{2.069810in}}%
\pgfpathlineto{\pgfqpoint{2.885016in}{2.088418in}}%
\pgfpathclose%
\pgfusepath{stroke,fill}%
\end{pgfscope}%
\begin{pgfscope}%
\pgfpathrectangle{\pgfqpoint{0.050000in}{0.050000in}}{\pgfqpoint{4.900000in}{4.900000in}}%
\pgfusepath{clip}%
\pgfsetbuttcap%
\pgfsetroundjoin%
\definecolor{currentfill}{rgb}{0.998155,0.998155,0.998155}%
\pgfsetfillcolor{currentfill}%
\pgfsetfillopacity{0.900000}%
\pgfsetlinewidth{0.501875pt}%
\definecolor{currentstroke}{rgb}{0.200000,0.200000,0.200000}%
\pgfsetstrokecolor{currentstroke}%
\pgfsetdash{}{0pt}%
\pgfpathmoveto{\pgfqpoint{2.694109in}{2.093801in}}%
\pgfpathlineto{\pgfqpoint{2.737322in}{2.102766in}}%
\pgfpathlineto{\pgfqpoint{2.767983in}{2.084729in}}%
\pgfpathlineto{\pgfqpoint{2.724741in}{2.076191in}}%
\pgfpathlineto{\pgfqpoint{2.694109in}{2.093801in}}%
\pgfpathclose%
\pgfusepath{stroke,fill}%
\end{pgfscope}%
\begin{pgfscope}%
\pgfpathrectangle{\pgfqpoint{0.050000in}{0.050000in}}{\pgfqpoint{4.900000in}{4.900000in}}%
\pgfusepath{clip}%
\pgfsetbuttcap%
\pgfsetroundjoin%
\definecolor{currentfill}{rgb}{1.000000,1.000000,1.000000}%
\pgfsetfillcolor{currentfill}%
\pgfsetfillopacity{0.900000}%
\pgfsetlinewidth{0.501875pt}%
\definecolor{currentstroke}{rgb}{0.200000,0.200000,0.200000}%
\pgfsetstrokecolor{currentstroke}%
\pgfsetdash{}{0pt}%
\pgfpathmoveto{\pgfqpoint{3.149889in}{2.082931in}}%
\pgfpathlineto{\pgfqpoint{3.192639in}{2.096128in}}%
\pgfpathlineto{\pgfqpoint{3.223975in}{2.077121in}}%
\pgfpathlineto{\pgfqpoint{3.181194in}{2.063924in}}%
\pgfpathlineto{\pgfqpoint{3.149889in}{2.082931in}}%
\pgfpathclose%
\pgfusepath{stroke,fill}%
\end{pgfscope}%
\begin{pgfscope}%
\pgfpathrectangle{\pgfqpoint{0.050000in}{0.050000in}}{\pgfqpoint{4.900000in}{4.900000in}}%
\pgfusepath{clip}%
\pgfsetbuttcap%
\pgfsetroundjoin%
\definecolor{currentfill}{rgb}{1.000000,1.000000,1.000000}%
\pgfsetfillcolor{currentfill}%
\pgfsetfillopacity{0.900000}%
\pgfsetlinewidth{0.501875pt}%
\definecolor{currentstroke}{rgb}{0.200000,0.200000,0.200000}%
\pgfsetstrokecolor{currentstroke}%
\pgfsetdash{}{0pt}%
\pgfpathmoveto{\pgfqpoint{2.958961in}{2.082001in}}%
\pgfpathlineto{\pgfqpoint{3.001914in}{2.094707in}}%
\pgfpathlineto{\pgfqpoint{3.032973in}{2.075839in}}%
\pgfpathlineto{\pgfqpoint{2.989990in}{2.063223in}}%
\pgfpathlineto{\pgfqpoint{2.958961in}{2.082001in}}%
\pgfpathclose%
\pgfusepath{stroke,fill}%
\end{pgfscope}%
\begin{pgfscope}%
\pgfpathrectangle{\pgfqpoint{0.050000in}{0.050000in}}{\pgfqpoint{4.900000in}{4.900000in}}%
\pgfusepath{clip}%
\pgfsetbuttcap%
\pgfsetroundjoin%
\definecolor{currentfill}{rgb}{1.000000,1.000000,1.000000}%
\pgfsetfillcolor{currentfill}%
\pgfsetfillopacity{0.900000}%
\pgfsetlinewidth{0.501875pt}%
\definecolor{currentstroke}{rgb}{0.200000,0.200000,0.200000}%
\pgfsetstrokecolor{currentstroke}%
\pgfsetdash{}{0pt}%
\pgfpathmoveto{\pgfqpoint{2.767983in}{2.084729in}}%
\pgfpathlineto{\pgfqpoint{2.811137in}{2.095252in}}%
\pgfpathlineto{\pgfqpoint{2.841920in}{2.076883in}}%
\pgfpathlineto{\pgfqpoint{2.798736in}{2.066680in}}%
\pgfpathlineto{\pgfqpoint{2.767983in}{2.084729in}}%
\pgfpathclose%
\pgfusepath{stroke,fill}%
\end{pgfscope}%
\begin{pgfscope}%
\pgfpathrectangle{\pgfqpoint{0.050000in}{0.050000in}}{\pgfqpoint{4.900000in}{4.900000in}}%
\pgfusepath{clip}%
\pgfsetbuttcap%
\pgfsetroundjoin%
\definecolor{currentfill}{rgb}{0.983391,0.983391,0.983391}%
\pgfsetfillcolor{currentfill}%
\pgfsetfillopacity{0.900000}%
\pgfsetlinewidth{0.501875pt}%
\definecolor{currentstroke}{rgb}{0.200000,0.200000,0.200000}%
\pgfsetstrokecolor{currentstroke}%
\pgfsetdash{}{0pt}%
\pgfpathmoveto{\pgfqpoint{2.355393in}{2.179756in}}%
\pgfpathlineto{\pgfqpoint{2.399089in}{2.161949in}}%
\pgfpathlineto{\pgfqpoint{2.429240in}{2.146566in}}%
\pgfpathlineto{\pgfqpoint{2.385524in}{2.163987in}}%
\pgfpathlineto{\pgfqpoint{2.355393in}{2.179756in}}%
\pgfpathclose%
\pgfusepath{stroke,fill}%
\end{pgfscope}%
\begin{pgfscope}%
\pgfpathrectangle{\pgfqpoint{0.050000in}{0.050000in}}{\pgfqpoint{4.900000in}{4.900000in}}%
\pgfusepath{clip}%
\pgfsetbuttcap%
\pgfsetroundjoin%
\definecolor{currentfill}{rgb}{0.988927,0.988927,0.988927}%
\pgfsetfillcolor{currentfill}%
\pgfsetfillopacity{0.900000}%
\pgfsetlinewidth{0.501875pt}%
\definecolor{currentstroke}{rgb}{0.200000,0.200000,0.200000}%
\pgfsetstrokecolor{currentstroke}%
\pgfsetdash{}{0pt}%
\pgfpathmoveto{\pgfqpoint{2.429240in}{2.146566in}}%
\pgfpathlineto{\pgfqpoint{2.472803in}{2.137159in}}%
\pgfpathlineto{\pgfqpoint{2.503071in}{2.121368in}}%
\pgfpathlineto{\pgfqpoint{2.459484in}{2.130891in}}%
\pgfpathlineto{\pgfqpoint{2.429240in}{2.146566in}}%
\pgfpathclose%
\pgfusepath{stroke,fill}%
\end{pgfscope}%
\begin{pgfscope}%
\pgfpathrectangle{\pgfqpoint{0.050000in}{0.050000in}}{\pgfqpoint{4.900000in}{4.900000in}}%
\pgfusepath{clip}%
\pgfsetbuttcap%
\pgfsetroundjoin%
\definecolor{currentfill}{rgb}{1.000000,1.000000,1.000000}%
\pgfsetfillcolor{currentfill}%
\pgfsetfillopacity{0.900000}%
\pgfsetlinewidth{0.501875pt}%
\definecolor{currentstroke}{rgb}{0.200000,0.200000,0.200000}%
\pgfsetstrokecolor{currentstroke}%
\pgfsetdash{}{0pt}%
\pgfpathmoveto{\pgfqpoint{3.223975in}{2.077121in}}%
\pgfpathlineto{\pgfqpoint{3.266668in}{2.090370in}}%
\pgfpathlineto{\pgfqpoint{3.298128in}{2.071326in}}%
\pgfpathlineto{\pgfqpoint{3.255405in}{2.058067in}}%
\pgfpathlineto{\pgfqpoint{3.223975in}{2.077121in}}%
\pgfpathclose%
\pgfusepath{stroke,fill}%
\end{pgfscope}%
\begin{pgfscope}%
\pgfpathrectangle{\pgfqpoint{0.050000in}{0.050000in}}{\pgfqpoint{4.900000in}{4.900000in}}%
\pgfusepath{clip}%
\pgfsetbuttcap%
\pgfsetroundjoin%
\definecolor{currentfill}{rgb}{0.974164,0.974164,0.974164}%
\pgfsetfillcolor{currentfill}%
\pgfsetfillopacity{0.900000}%
\pgfsetlinewidth{0.501875pt}%
\definecolor{currentstroke}{rgb}{0.200000,0.200000,0.200000}%
\pgfsetstrokecolor{currentstroke}%
\pgfsetdash{}{0pt}%
\pgfpathmoveto{\pgfqpoint{2.281492in}{2.222237in}}%
\pgfpathlineto{\pgfqpoint{2.325355in}{2.195256in}}%
\pgfpathlineto{\pgfqpoint{2.355393in}{2.179756in}}%
\pgfpathlineto{\pgfqpoint{2.311520in}{2.205754in}}%
\pgfpathlineto{\pgfqpoint{2.281492in}{2.222237in}}%
\pgfpathclose%
\pgfusepath{stroke,fill}%
\end{pgfscope}%
\begin{pgfscope}%
\pgfpathrectangle{\pgfqpoint{0.050000in}{0.050000in}}{\pgfqpoint{4.900000in}{4.900000in}}%
\pgfusepath{clip}%
\pgfsetbuttcap%
\pgfsetroundjoin%
\definecolor{currentfill}{rgb}{0.992618,0.992618,0.992618}%
\pgfsetfillcolor{currentfill}%
\pgfsetfillopacity{0.900000}%
\pgfsetlinewidth{0.501875pt}%
\definecolor{currentstroke}{rgb}{0.200000,0.200000,0.200000}%
\pgfsetstrokecolor{currentstroke}%
\pgfsetdash{}{0pt}%
\pgfpathmoveto{\pgfqpoint{2.503071in}{2.121368in}}%
\pgfpathlineto{\pgfqpoint{2.546531in}{2.118753in}}%
\pgfpathlineto{\pgfqpoint{2.576919in}{2.102335in}}%
\pgfpathlineto{\pgfqpoint{2.533433in}{2.105351in}}%
\pgfpathlineto{\pgfqpoint{2.503071in}{2.121368in}}%
\pgfpathclose%
\pgfusepath{stroke,fill}%
\end{pgfscope}%
\begin{pgfscope}%
\pgfpathrectangle{\pgfqpoint{0.050000in}{0.050000in}}{\pgfqpoint{4.900000in}{4.900000in}}%
\pgfusepath{clip}%
\pgfsetbuttcap%
\pgfsetroundjoin%
\definecolor{currentfill}{rgb}{1.000000,1.000000,1.000000}%
\pgfsetfillcolor{currentfill}%
\pgfsetfillopacity{0.900000}%
\pgfsetlinewidth{0.501875pt}%
\definecolor{currentstroke}{rgb}{0.200000,0.200000,0.200000}%
\pgfsetstrokecolor{currentstroke}%
\pgfsetdash{}{0pt}%
\pgfpathmoveto{\pgfqpoint{3.032973in}{2.075839in}}%
\pgfpathlineto{\pgfqpoint{3.075868in}{2.088782in}}%
\pgfpathlineto{\pgfqpoint{3.107050in}{2.069834in}}%
\pgfpathlineto{\pgfqpoint{3.064124in}{2.056941in}}%
\pgfpathlineto{\pgfqpoint{3.032973in}{2.075839in}}%
\pgfpathclose%
\pgfusepath{stroke,fill}%
\end{pgfscope}%
\begin{pgfscope}%
\pgfpathrectangle{\pgfqpoint{0.050000in}{0.050000in}}{\pgfqpoint{4.900000in}{4.900000in}}%
\pgfusepath{clip}%
\pgfsetbuttcap%
\pgfsetroundjoin%
\definecolor{currentfill}{rgb}{0.964937,0.964937,0.964937}%
\pgfsetfillcolor{currentfill}%
\pgfsetfillopacity{0.900000}%
\pgfsetlinewidth{0.501875pt}%
\definecolor{currentstroke}{rgb}{0.200000,0.200000,0.200000}%
\pgfsetstrokecolor{currentstroke}%
\pgfsetdash{}{0pt}%
\pgfpathmoveto{\pgfqpoint{2.207504in}{2.273790in}}%
\pgfpathlineto{\pgfqpoint{2.251557in}{2.238630in}}%
\pgfpathlineto{\pgfqpoint{2.281492in}{2.222237in}}%
\pgfpathlineto{\pgfqpoint{2.237437in}{2.256091in}}%
\pgfpathlineto{\pgfqpoint{2.207504in}{2.273790in}}%
\pgfpathclose%
\pgfusepath{stroke,fill}%
\end{pgfscope}%
\begin{pgfscope}%
\pgfpathrectangle{\pgfqpoint{0.050000in}{0.050000in}}{\pgfqpoint{4.900000in}{4.900000in}}%
\pgfusepath{clip}%
\pgfsetbuttcap%
\pgfsetroundjoin%
\definecolor{currentfill}{rgb}{1.000000,1.000000,1.000000}%
\pgfsetfillcolor{currentfill}%
\pgfsetfillopacity{0.900000}%
\pgfsetlinewidth{0.501875pt}%
\definecolor{currentstroke}{rgb}{0.200000,0.200000,0.200000}%
\pgfsetstrokecolor{currentstroke}%
\pgfsetdash{}{0pt}%
\pgfpathmoveto{\pgfqpoint{2.841920in}{2.076883in}}%
\pgfpathlineto{\pgfqpoint{2.885016in}{2.088418in}}%
\pgfpathlineto{\pgfqpoint{2.915922in}{2.069810in}}%
\pgfpathlineto{\pgfqpoint{2.872795in}{2.058502in}}%
\pgfpathlineto{\pgfqpoint{2.841920in}{2.076883in}}%
\pgfpathclose%
\pgfusepath{stroke,fill}%
\end{pgfscope}%
\begin{pgfscope}%
\pgfpathrectangle{\pgfqpoint{0.050000in}{0.050000in}}{\pgfqpoint{4.900000in}{4.900000in}}%
\pgfusepath{clip}%
\pgfsetbuttcap%
\pgfsetroundjoin%
\definecolor{currentfill}{rgb}{0.996309,0.996309,0.996309}%
\pgfsetfillcolor{currentfill}%
\pgfsetfillopacity{0.900000}%
\pgfsetlinewidth{0.501875pt}%
\definecolor{currentstroke}{rgb}{0.200000,0.200000,0.200000}%
\pgfsetstrokecolor{currentstroke}%
\pgfsetdash{}{0pt}%
\pgfpathmoveto{\pgfqpoint{2.576919in}{2.102335in}}%
\pgfpathlineto{\pgfqpoint{2.620295in}{2.104795in}}%
\pgfpathlineto{\pgfqpoint{2.650804in}{2.087738in}}%
\pgfpathlineto{\pgfqpoint{2.607401in}{2.085774in}}%
\pgfpathlineto{\pgfqpoint{2.576919in}{2.102335in}}%
\pgfpathclose%
\pgfusepath{stroke,fill}%
\end{pgfscope}%
\begin{pgfscope}%
\pgfpathrectangle{\pgfqpoint{0.050000in}{0.050000in}}{\pgfqpoint{4.900000in}{4.900000in}}%
\pgfusepath{clip}%
\pgfsetbuttcap%
\pgfsetroundjoin%
\definecolor{currentfill}{rgb}{1.000000,1.000000,1.000000}%
\pgfsetfillcolor{currentfill}%
\pgfsetfillopacity{0.900000}%
\pgfsetlinewidth{0.501875pt}%
\definecolor{currentstroke}{rgb}{0.200000,0.200000,0.200000}%
\pgfsetstrokecolor{currentstroke}%
\pgfsetdash{}{0pt}%
\pgfpathmoveto{\pgfqpoint{3.298128in}{2.071326in}}%
\pgfpathlineto{\pgfqpoint{3.340763in}{2.084599in}}%
\pgfpathlineto{\pgfqpoint{3.372346in}{2.065517in}}%
\pgfpathlineto{\pgfqpoint{3.329682in}{2.052230in}}%
\pgfpathlineto{\pgfqpoint{3.298128in}{2.071326in}}%
\pgfpathclose%
\pgfusepath{stroke,fill}%
\end{pgfscope}%
\begin{pgfscope}%
\pgfpathrectangle{\pgfqpoint{0.050000in}{0.050000in}}{\pgfqpoint{4.900000in}{4.900000in}}%
\pgfusepath{clip}%
\pgfsetbuttcap%
\pgfsetroundjoin%
\definecolor{currentfill}{rgb}{1.000000,1.000000,1.000000}%
\pgfsetfillcolor{currentfill}%
\pgfsetfillopacity{0.900000}%
\pgfsetlinewidth{0.501875pt}%
\definecolor{currentstroke}{rgb}{0.200000,0.200000,0.200000}%
\pgfsetstrokecolor{currentstroke}%
\pgfsetdash{}{0pt}%
\pgfpathmoveto{\pgfqpoint{3.107050in}{2.069834in}}%
\pgfpathlineto{\pgfqpoint{3.149889in}{2.082931in}}%
\pgfpathlineto{\pgfqpoint{3.181194in}{2.063924in}}%
\pgfpathlineto{\pgfqpoint{3.138326in}{2.050850in}}%
\pgfpathlineto{\pgfqpoint{3.107050in}{2.069834in}}%
\pgfpathclose%
\pgfusepath{stroke,fill}%
\end{pgfscope}%
\begin{pgfscope}%
\pgfpathrectangle{\pgfqpoint{0.050000in}{0.050000in}}{\pgfqpoint{4.900000in}{4.900000in}}%
\pgfusepath{clip}%
\pgfsetbuttcap%
\pgfsetroundjoin%
\definecolor{currentfill}{rgb}{0.998155,0.998155,0.998155}%
\pgfsetfillcolor{currentfill}%
\pgfsetfillopacity{0.900000}%
\pgfsetlinewidth{0.501875pt}%
\definecolor{currentstroke}{rgb}{0.200000,0.200000,0.200000}%
\pgfsetstrokecolor{currentstroke}%
\pgfsetdash{}{0pt}%
\pgfpathmoveto{\pgfqpoint{2.650804in}{2.087738in}}%
\pgfpathlineto{\pgfqpoint{2.694109in}{2.093801in}}%
\pgfpathlineto{\pgfqpoint{2.724741in}{2.076191in}}%
\pgfpathlineto{\pgfqpoint{2.681407in}{2.070602in}}%
\pgfpathlineto{\pgfqpoint{2.650804in}{2.087738in}}%
\pgfpathclose%
\pgfusepath{stroke,fill}%
\end{pgfscope}%
\begin{pgfscope}%
\pgfpathrectangle{\pgfqpoint{0.050000in}{0.050000in}}{\pgfqpoint{4.900000in}{4.900000in}}%
\pgfusepath{clip}%
\pgfsetbuttcap%
\pgfsetroundjoin%
\definecolor{currentfill}{rgb}{0.952018,0.952018,0.952018}%
\pgfsetfillcolor{currentfill}%
\pgfsetfillopacity{0.900000}%
\pgfsetlinewidth{0.501875pt}%
\definecolor{currentstroke}{rgb}{0.200000,0.200000,0.200000}%
\pgfsetstrokecolor{currentstroke}%
\pgfsetdash{}{0pt}%
\pgfpathmoveto{\pgfqpoint{2.133411in}{2.331983in}}%
\pgfpathlineto{\pgfqpoint{2.177660in}{2.291655in}}%
\pgfpathlineto{\pgfqpoint{2.207504in}{2.273790in}}%
\pgfpathlineto{\pgfqpoint{2.163255in}{2.313138in}}%
\pgfpathlineto{\pgfqpoint{2.133411in}{2.331983in}}%
\pgfpathclose%
\pgfusepath{stroke,fill}%
\end{pgfscope}%
\begin{pgfscope}%
\pgfpathrectangle{\pgfqpoint{0.050000in}{0.050000in}}{\pgfqpoint{4.900000in}{4.900000in}}%
\pgfusepath{clip}%
\pgfsetbuttcap%
\pgfsetroundjoin%
\definecolor{currentfill}{rgb}{1.000000,1.000000,1.000000}%
\pgfsetfillcolor{currentfill}%
\pgfsetfillopacity{0.900000}%
\pgfsetlinewidth{0.501875pt}%
\definecolor{currentstroke}{rgb}{0.200000,0.200000,0.200000}%
\pgfsetstrokecolor{currentstroke}%
\pgfsetdash{}{0pt}%
\pgfpathmoveto{\pgfqpoint{2.915922in}{2.069810in}}%
\pgfpathlineto{\pgfqpoint{2.958961in}{2.082001in}}%
\pgfpathlineto{\pgfqpoint{2.989990in}{2.063223in}}%
\pgfpathlineto{\pgfqpoint{2.946920in}{2.051186in}}%
\pgfpathlineto{\pgfqpoint{2.915922in}{2.069810in}}%
\pgfpathclose%
\pgfusepath{stroke,fill}%
\end{pgfscope}%
\begin{pgfscope}%
\pgfpathrectangle{\pgfqpoint{0.050000in}{0.050000in}}{\pgfqpoint{4.900000in}{4.900000in}}%
\pgfusepath{clip}%
\pgfsetbuttcap%
\pgfsetroundjoin%
\definecolor{currentfill}{rgb}{1.000000,1.000000,1.000000}%
\pgfsetfillcolor{currentfill}%
\pgfsetfillopacity{0.900000}%
\pgfsetlinewidth{0.501875pt}%
\definecolor{currentstroke}{rgb}{0.200000,0.200000,0.200000}%
\pgfsetstrokecolor{currentstroke}%
\pgfsetdash{}{0pt}%
\pgfpathmoveto{\pgfqpoint{3.181194in}{2.063924in}}%
\pgfpathlineto{\pgfqpoint{3.223975in}{2.077121in}}%
\pgfpathlineto{\pgfqpoint{3.255405in}{2.058067in}}%
\pgfpathlineto{\pgfqpoint{3.212594in}{2.044874in}}%
\pgfpathlineto{\pgfqpoint{3.181194in}{2.063924in}}%
\pgfpathclose%
\pgfusepath{stroke,fill}%
\end{pgfscope}%
\begin{pgfscope}%
\pgfpathrectangle{\pgfqpoint{0.050000in}{0.050000in}}{\pgfqpoint{4.900000in}{4.900000in}}%
\pgfusepath{clip}%
\pgfsetbuttcap%
\pgfsetroundjoin%
\definecolor{currentfill}{rgb}{0.998155,0.998155,0.998155}%
\pgfsetfillcolor{currentfill}%
\pgfsetfillopacity{0.900000}%
\pgfsetlinewidth{0.501875pt}%
\definecolor{currentstroke}{rgb}{0.200000,0.200000,0.200000}%
\pgfsetstrokecolor{currentstroke}%
\pgfsetdash{}{0pt}%
\pgfpathmoveto{\pgfqpoint{2.724741in}{2.076191in}}%
\pgfpathlineto{\pgfqpoint{2.767983in}{2.084729in}}%
\pgfpathlineto{\pgfqpoint{2.798736in}{2.066680in}}%
\pgfpathlineto{\pgfqpoint{2.755465in}{2.058541in}}%
\pgfpathlineto{\pgfqpoint{2.724741in}{2.076191in}}%
\pgfpathclose%
\pgfusepath{stroke,fill}%
\end{pgfscope}%
\begin{pgfscope}%
\pgfpathrectangle{\pgfqpoint{0.050000in}{0.050000in}}{\pgfqpoint{4.900000in}{4.900000in}}%
\pgfusepath{clip}%
\pgfsetbuttcap%
\pgfsetroundjoin%
\definecolor{currentfill}{rgb}{1.000000,1.000000,1.000000}%
\pgfsetfillcolor{currentfill}%
\pgfsetfillopacity{0.900000}%
\pgfsetlinewidth{0.501875pt}%
\definecolor{currentstroke}{rgb}{0.200000,0.200000,0.200000}%
\pgfsetstrokecolor{currentstroke}%
\pgfsetdash{}{0pt}%
\pgfpathmoveto{\pgfqpoint{2.989990in}{2.063223in}}%
\pgfpathlineto{\pgfqpoint{3.032973in}{2.075839in}}%
\pgfpathlineto{\pgfqpoint{3.064124in}{2.056941in}}%
\pgfpathlineto{\pgfqpoint{3.021111in}{2.044423in}}%
\pgfpathlineto{\pgfqpoint{2.989990in}{2.063223in}}%
\pgfpathclose%
\pgfusepath{stroke,fill}%
\end{pgfscope}%
\begin{pgfscope}%
\pgfpathrectangle{\pgfqpoint{0.050000in}{0.050000in}}{\pgfqpoint{4.900000in}{4.900000in}}%
\pgfusepath{clip}%
\pgfsetbuttcap%
\pgfsetroundjoin%
\definecolor{currentfill}{rgb}{0.937993,0.937993,0.937993}%
\pgfsetfillcolor{currentfill}%
\pgfsetfillopacity{0.900000}%
\pgfsetlinewidth{0.501875pt}%
\definecolor{currentstroke}{rgb}{0.200000,0.200000,0.200000}%
\pgfsetstrokecolor{currentstroke}%
\pgfsetdash{}{0pt}%
\pgfpathmoveto{\pgfqpoint{2.059220in}{2.393144in}}%
\pgfpathlineto{\pgfqpoint{2.103654in}{2.351030in}}%
\pgfpathlineto{\pgfqpoint{2.133411in}{2.331983in}}%
\pgfpathlineto{\pgfqpoint{2.088975in}{2.373808in}}%
\pgfpathlineto{\pgfqpoint{2.059220in}{2.393144in}}%
\pgfpathclose%
\pgfusepath{stroke,fill}%
\end{pgfscope}%
\begin{pgfscope}%
\pgfpathrectangle{\pgfqpoint{0.050000in}{0.050000in}}{\pgfqpoint{4.900000in}{4.900000in}}%
\pgfusepath{clip}%
\pgfsetbuttcap%
\pgfsetroundjoin%
\definecolor{currentfill}{rgb}{1.000000,1.000000,1.000000}%
\pgfsetfillcolor{currentfill}%
\pgfsetfillopacity{0.900000}%
\pgfsetlinewidth{0.501875pt}%
\definecolor{currentstroke}{rgb}{0.200000,0.200000,0.200000}%
\pgfsetstrokecolor{currentstroke}%
\pgfsetdash{}{0pt}%
\pgfpathmoveto{\pgfqpoint{2.798736in}{2.066680in}}%
\pgfpathlineto{\pgfqpoint{2.841920in}{2.076883in}}%
\pgfpathlineto{\pgfqpoint{2.872795in}{2.058502in}}%
\pgfpathlineto{\pgfqpoint{2.829583in}{2.048608in}}%
\pgfpathlineto{\pgfqpoint{2.798736in}{2.066680in}}%
\pgfpathclose%
\pgfusepath{stroke,fill}%
\end{pgfscope}%
\begin{pgfscope}%
\pgfpathrectangle{\pgfqpoint{0.050000in}{0.050000in}}{\pgfqpoint{4.900000in}{4.900000in}}%
\pgfusepath{clip}%
\pgfsetbuttcap%
\pgfsetroundjoin%
\definecolor{currentfill}{rgb}{0.981546,0.981546,0.981546}%
\pgfsetfillcolor{currentfill}%
\pgfsetfillopacity{0.900000}%
\pgfsetlinewidth{0.501875pt}%
\definecolor{currentstroke}{rgb}{0.200000,0.200000,0.200000}%
\pgfsetstrokecolor{currentstroke}%
\pgfsetdash{}{0pt}%
\pgfpathmoveto{\pgfqpoint{2.385524in}{2.163987in}}%
\pgfpathlineto{\pgfqpoint{2.429240in}{2.146566in}}%
\pgfpathlineto{\pgfqpoint{2.459484in}{2.130891in}}%
\pgfpathlineto{\pgfqpoint{2.415750in}{2.147981in}}%
\pgfpathlineto{\pgfqpoint{2.385524in}{2.163987in}}%
\pgfpathclose%
\pgfusepath{stroke,fill}%
\end{pgfscope}%
\begin{pgfscope}%
\pgfpathrectangle{\pgfqpoint{0.050000in}{0.050000in}}{\pgfqpoint{4.900000in}{4.900000in}}%
\pgfusepath{clip}%
\pgfsetbuttcap%
\pgfsetroundjoin%
\definecolor{currentfill}{rgb}{1.000000,1.000000,1.000000}%
\pgfsetfillcolor{currentfill}%
\pgfsetfillopacity{0.900000}%
\pgfsetlinewidth{0.501875pt}%
\definecolor{currentstroke}{rgb}{0.200000,0.200000,0.200000}%
\pgfsetstrokecolor{currentstroke}%
\pgfsetdash{}{0pt}%
\pgfpathmoveto{\pgfqpoint{3.255405in}{2.058067in}}%
\pgfpathlineto{\pgfqpoint{3.298128in}{2.071326in}}%
\pgfpathlineto{\pgfqpoint{3.329682in}{2.052230in}}%
\pgfpathlineto{\pgfqpoint{3.286928in}{2.038965in}}%
\pgfpathlineto{\pgfqpoint{3.255405in}{2.058067in}}%
\pgfpathclose%
\pgfusepath{stroke,fill}%
\end{pgfscope}%
\begin{pgfscope}%
\pgfpathrectangle{\pgfqpoint{0.050000in}{0.050000in}}{\pgfqpoint{4.900000in}{4.900000in}}%
\pgfusepath{clip}%
\pgfsetbuttcap%
\pgfsetroundjoin%
\definecolor{currentfill}{rgb}{0.988927,0.988927,0.988927}%
\pgfsetfillcolor{currentfill}%
\pgfsetfillopacity{0.900000}%
\pgfsetlinewidth{0.501875pt}%
\definecolor{currentstroke}{rgb}{0.200000,0.200000,0.200000}%
\pgfsetstrokecolor{currentstroke}%
\pgfsetdash{}{0pt}%
\pgfpathmoveto{\pgfqpoint{2.459484in}{2.130891in}}%
\pgfpathlineto{\pgfqpoint{2.503071in}{2.121368in}}%
\pgfpathlineto{\pgfqpoint{2.533433in}{2.105351in}}%
\pgfpathlineto{\pgfqpoint{2.489822in}{2.114963in}}%
\pgfpathlineto{\pgfqpoint{2.459484in}{2.130891in}}%
\pgfpathclose%
\pgfusepath{stroke,fill}%
\end{pgfscope}%
\begin{pgfscope}%
\pgfpathrectangle{\pgfqpoint{0.050000in}{0.050000in}}{\pgfqpoint{4.900000in}{4.900000in}}%
\pgfusepath{clip}%
\pgfsetbuttcap%
\pgfsetroundjoin%
\definecolor{currentfill}{rgb}{0.974164,0.974164,0.974164}%
\pgfsetfillcolor{currentfill}%
\pgfsetfillopacity{0.900000}%
\pgfsetlinewidth{0.501875pt}%
\definecolor{currentstroke}{rgb}{0.200000,0.200000,0.200000}%
\pgfsetstrokecolor{currentstroke}%
\pgfsetdash{}{0pt}%
\pgfpathmoveto{\pgfqpoint{2.311520in}{2.205754in}}%
\pgfpathlineto{\pgfqpoint{2.355393in}{2.179756in}}%
\pgfpathlineto{\pgfqpoint{2.385524in}{2.163987in}}%
\pgfpathlineto{\pgfqpoint{2.341639in}{2.189161in}}%
\pgfpathlineto{\pgfqpoint{2.311520in}{2.205754in}}%
\pgfpathclose%
\pgfusepath{stroke,fill}%
\end{pgfscope}%
\begin{pgfscope}%
\pgfpathrectangle{\pgfqpoint{0.050000in}{0.050000in}}{\pgfqpoint{4.900000in}{4.900000in}}%
\pgfusepath{clip}%
\pgfsetbuttcap%
\pgfsetroundjoin%
\definecolor{currentfill}{rgb}{1.000000,1.000000,1.000000}%
\pgfsetfillcolor{currentfill}%
\pgfsetfillopacity{0.900000}%
\pgfsetlinewidth{0.501875pt}%
\definecolor{currentstroke}{rgb}{0.200000,0.200000,0.200000}%
\pgfsetstrokecolor{currentstroke}%
\pgfsetdash{}{0pt}%
\pgfpathmoveto{\pgfqpoint{3.064124in}{2.056941in}}%
\pgfpathlineto{\pgfqpoint{3.107050in}{2.069834in}}%
\pgfpathlineto{\pgfqpoint{3.138326in}{2.050850in}}%
\pgfpathlineto{\pgfqpoint{3.095369in}{2.038013in}}%
\pgfpathlineto{\pgfqpoint{3.064124in}{2.056941in}}%
\pgfpathclose%
\pgfusepath{stroke,fill}%
\end{pgfscope}%
\begin{pgfscope}%
\pgfpathrectangle{\pgfqpoint{0.050000in}{0.050000in}}{\pgfqpoint{4.900000in}{4.900000in}}%
\pgfusepath{clip}%
\pgfsetbuttcap%
\pgfsetroundjoin%
\definecolor{currentfill}{rgb}{0.992618,0.992618,0.992618}%
\pgfsetfillcolor{currentfill}%
\pgfsetfillopacity{0.900000}%
\pgfsetlinewidth{0.501875pt}%
\definecolor{currentstroke}{rgb}{0.200000,0.200000,0.200000}%
\pgfsetstrokecolor{currentstroke}%
\pgfsetdash{}{0pt}%
\pgfpathmoveto{\pgfqpoint{2.533433in}{2.105351in}}%
\pgfpathlineto{\pgfqpoint{2.576919in}{2.102335in}}%
\pgfpathlineto{\pgfqpoint{2.607401in}{2.085774in}}%
\pgfpathlineto{\pgfqpoint{2.563889in}{2.089131in}}%
\pgfpathlineto{\pgfqpoint{2.533433in}{2.105351in}}%
\pgfpathclose%
\pgfusepath{stroke,fill}%
\end{pgfscope}%
\begin{pgfscope}%
\pgfpathrectangle{\pgfqpoint{0.050000in}{0.050000in}}{\pgfqpoint{4.900000in}{4.900000in}}%
\pgfusepath{clip}%
\pgfsetbuttcap%
\pgfsetroundjoin%
\definecolor{currentfill}{rgb}{1.000000,1.000000,1.000000}%
\pgfsetfillcolor{currentfill}%
\pgfsetfillopacity{0.900000}%
\pgfsetlinewidth{0.501875pt}%
\definecolor{currentstroke}{rgb}{0.200000,0.200000,0.200000}%
\pgfsetstrokecolor{currentstroke}%
\pgfsetdash{}{0pt}%
\pgfpathmoveto{\pgfqpoint{2.872795in}{2.058502in}}%
\pgfpathlineto{\pgfqpoint{2.915922in}{2.069810in}}%
\pgfpathlineto{\pgfqpoint{2.946920in}{2.051186in}}%
\pgfpathlineto{\pgfqpoint{2.903764in}{2.040103in}}%
\pgfpathlineto{\pgfqpoint{2.872795in}{2.058502in}}%
\pgfpathclose%
\pgfusepath{stroke,fill}%
\end{pgfscope}%
\begin{pgfscope}%
\pgfpathrectangle{\pgfqpoint{0.050000in}{0.050000in}}{\pgfqpoint{4.900000in}{4.900000in}}%
\pgfusepath{clip}%
\pgfsetbuttcap%
\pgfsetroundjoin%
\definecolor{currentfill}{rgb}{0.964937,0.964937,0.964937}%
\pgfsetfillcolor{currentfill}%
\pgfsetfillopacity{0.900000}%
\pgfsetlinewidth{0.501875pt}%
\definecolor{currentstroke}{rgb}{0.200000,0.200000,0.200000}%
\pgfsetstrokecolor{currentstroke}%
\pgfsetdash{}{0pt}%
\pgfpathmoveto{\pgfqpoint{2.237437in}{2.256091in}}%
\pgfpathlineto{\pgfqpoint{2.281492in}{2.222237in}}%
\pgfpathlineto{\pgfqpoint{2.311520in}{2.205754in}}%
\pgfpathlineto{\pgfqpoint{2.267459in}{2.238476in}}%
\pgfpathlineto{\pgfqpoint{2.237437in}{2.256091in}}%
\pgfpathclose%
\pgfusepath{stroke,fill}%
\end{pgfscope}%
\begin{pgfscope}%
\pgfpathrectangle{\pgfqpoint{0.050000in}{0.050000in}}{\pgfqpoint{4.900000in}{4.900000in}}%
\pgfusepath{clip}%
\pgfsetbuttcap%
\pgfsetroundjoin%
\definecolor{currentfill}{rgb}{0.994464,0.994464,0.994464}%
\pgfsetfillcolor{currentfill}%
\pgfsetfillopacity{0.900000}%
\pgfsetlinewidth{0.501875pt}%
\definecolor{currentstroke}{rgb}{0.200000,0.200000,0.200000}%
\pgfsetstrokecolor{currentstroke}%
\pgfsetdash{}{0pt}%
\pgfpathmoveto{\pgfqpoint{2.607401in}{2.085774in}}%
\pgfpathlineto{\pgfqpoint{2.650804in}{2.087738in}}%
\pgfpathlineto{\pgfqpoint{2.681407in}{2.070602in}}%
\pgfpathlineto{\pgfqpoint{2.637976in}{2.069076in}}%
\pgfpathlineto{\pgfqpoint{2.607401in}{2.085774in}}%
\pgfpathclose%
\pgfusepath{stroke,fill}%
\end{pgfscope}%
\begin{pgfscope}%
\pgfpathrectangle{\pgfqpoint{0.050000in}{0.050000in}}{\pgfqpoint{4.900000in}{4.900000in}}%
\pgfusepath{clip}%
\pgfsetbuttcap%
\pgfsetroundjoin%
\definecolor{currentfill}{rgb}{1.000000,1.000000,1.000000}%
\pgfsetfillcolor{currentfill}%
\pgfsetfillopacity{0.900000}%
\pgfsetlinewidth{0.501875pt}%
\definecolor{currentstroke}{rgb}{0.200000,0.200000,0.200000}%
\pgfsetstrokecolor{currentstroke}%
\pgfsetdash{}{0pt}%
\pgfpathmoveto{\pgfqpoint{3.329682in}{2.052230in}}%
\pgfpathlineto{\pgfqpoint{3.372346in}{2.065517in}}%
\pgfpathlineto{\pgfqpoint{3.404024in}{2.046378in}}%
\pgfpathlineto{\pgfqpoint{3.361329in}{2.033082in}}%
\pgfpathlineto{\pgfqpoint{3.329682in}{2.052230in}}%
\pgfpathclose%
\pgfusepath{stroke,fill}%
\end{pgfscope}%
\begin{pgfscope}%
\pgfpathrectangle{\pgfqpoint{0.050000in}{0.050000in}}{\pgfqpoint{4.900000in}{4.900000in}}%
\pgfusepath{clip}%
\pgfsetbuttcap%
\pgfsetroundjoin%
\definecolor{currentfill}{rgb}{1.000000,1.000000,1.000000}%
\pgfsetfillcolor{currentfill}%
\pgfsetfillopacity{0.900000}%
\pgfsetlinewidth{0.501875pt}%
\definecolor{currentstroke}{rgb}{0.200000,0.200000,0.200000}%
\pgfsetstrokecolor{currentstroke}%
\pgfsetdash{}{0pt}%
\pgfpathmoveto{\pgfqpoint{3.138326in}{2.050850in}}%
\pgfpathlineto{\pgfqpoint{3.181194in}{2.063924in}}%
\pgfpathlineto{\pgfqpoint{3.212594in}{2.044874in}}%
\pgfpathlineto{\pgfqpoint{3.169694in}{2.031828in}}%
\pgfpathlineto{\pgfqpoint{3.138326in}{2.050850in}}%
\pgfpathclose%
\pgfusepath{stroke,fill}%
\end{pgfscope}%
\begin{pgfscope}%
\pgfpathrectangle{\pgfqpoint{0.050000in}{0.050000in}}{\pgfqpoint{4.900000in}{4.900000in}}%
\pgfusepath{clip}%
\pgfsetbuttcap%
\pgfsetroundjoin%
\definecolor{currentfill}{rgb}{0.996309,0.996309,0.996309}%
\pgfsetfillcolor{currentfill}%
\pgfsetfillopacity{0.900000}%
\pgfsetlinewidth{0.501875pt}%
\definecolor{currentstroke}{rgb}{0.200000,0.200000,0.200000}%
\pgfsetstrokecolor{currentstroke}%
\pgfsetdash{}{0pt}%
\pgfpathmoveto{\pgfqpoint{2.681407in}{2.070602in}}%
\pgfpathlineto{\pgfqpoint{2.724741in}{2.076191in}}%
\pgfpathlineto{\pgfqpoint{2.755465in}{2.058541in}}%
\pgfpathlineto{\pgfqpoint{2.712103in}{2.053384in}}%
\pgfpathlineto{\pgfqpoint{2.681407in}{2.070602in}}%
\pgfpathclose%
\pgfusepath{stroke,fill}%
\end{pgfscope}%
\begin{pgfscope}%
\pgfpathrectangle{\pgfqpoint{0.050000in}{0.050000in}}{\pgfqpoint{4.900000in}{4.900000in}}%
\pgfusepath{clip}%
\pgfsetbuttcap%
\pgfsetroundjoin%
\definecolor{currentfill}{rgb}{0.952018,0.952018,0.952018}%
\pgfsetfillcolor{currentfill}%
\pgfsetfillopacity{0.900000}%
\pgfsetlinewidth{0.501875pt}%
\definecolor{currentstroke}{rgb}{0.200000,0.200000,0.200000}%
\pgfsetstrokecolor{currentstroke}%
\pgfsetdash{}{0pt}%
\pgfpathmoveto{\pgfqpoint{2.163255in}{2.313138in}}%
\pgfpathlineto{\pgfqpoint{2.207504in}{2.273790in}}%
\pgfpathlineto{\pgfqpoint{2.237437in}{2.256091in}}%
\pgfpathlineto{\pgfqpoint{2.193187in}{2.294464in}}%
\pgfpathlineto{\pgfqpoint{2.163255in}{2.313138in}}%
\pgfpathclose%
\pgfusepath{stroke,fill}%
\end{pgfscope}%
\begin{pgfscope}%
\pgfpathrectangle{\pgfqpoint{0.050000in}{0.050000in}}{\pgfqpoint{4.900000in}{4.900000in}}%
\pgfusepath{clip}%
\pgfsetbuttcap%
\pgfsetroundjoin%
\definecolor{currentfill}{rgb}{1.000000,1.000000,1.000000}%
\pgfsetfillcolor{currentfill}%
\pgfsetfillopacity{0.900000}%
\pgfsetlinewidth{0.501875pt}%
\definecolor{currentstroke}{rgb}{0.200000,0.200000,0.200000}%
\pgfsetstrokecolor{currentstroke}%
\pgfsetdash{}{0pt}%
\pgfpathmoveto{\pgfqpoint{2.946920in}{2.051186in}}%
\pgfpathlineto{\pgfqpoint{2.989990in}{2.063223in}}%
\pgfpathlineto{\pgfqpoint{3.021111in}{2.044423in}}%
\pgfpathlineto{\pgfqpoint{2.978011in}{2.032542in}}%
\pgfpathlineto{\pgfqpoint{2.946920in}{2.051186in}}%
\pgfpathclose%
\pgfusepath{stroke,fill}%
\end{pgfscope}%
\begin{pgfscope}%
\pgfpathrectangle{\pgfqpoint{0.050000in}{0.050000in}}{\pgfqpoint{4.900000in}{4.900000in}}%
\pgfusepath{clip}%
\pgfsetbuttcap%
\pgfsetroundjoin%
\definecolor{currentfill}{rgb}{1.000000,1.000000,1.000000}%
\pgfsetfillcolor{currentfill}%
\pgfsetfillopacity{0.900000}%
\pgfsetlinewidth{0.501875pt}%
\definecolor{currentstroke}{rgb}{0.200000,0.200000,0.200000}%
\pgfsetstrokecolor{currentstroke}%
\pgfsetdash{}{0pt}%
\pgfpathmoveto{\pgfqpoint{3.212594in}{2.044874in}}%
\pgfpathlineto{\pgfqpoint{3.255405in}{2.058067in}}%
\pgfpathlineto{\pgfqpoint{3.286928in}{2.038965in}}%
\pgfpathlineto{\pgfqpoint{3.244087in}{2.025782in}}%
\pgfpathlineto{\pgfqpoint{3.212594in}{2.044874in}}%
\pgfpathclose%
\pgfusepath{stroke,fill}%
\end{pgfscope}%
\begin{pgfscope}%
\pgfpathrectangle{\pgfqpoint{0.050000in}{0.050000in}}{\pgfqpoint{4.900000in}{4.900000in}}%
\pgfusepath{clip}%
\pgfsetbuttcap%
\pgfsetroundjoin%
\definecolor{currentfill}{rgb}{0.998155,0.998155,0.998155}%
\pgfsetfillcolor{currentfill}%
\pgfsetfillopacity{0.900000}%
\pgfsetlinewidth{0.501875pt}%
\definecolor{currentstroke}{rgb}{0.200000,0.200000,0.200000}%
\pgfsetstrokecolor{currentstroke}%
\pgfsetdash{}{0pt}%
\pgfpathmoveto{\pgfqpoint{2.755465in}{2.058541in}}%
\pgfpathlineto{\pgfqpoint{2.798736in}{2.066680in}}%
\pgfpathlineto{\pgfqpoint{2.829583in}{2.048608in}}%
\pgfpathlineto{\pgfqpoint{2.786283in}{2.040843in}}%
\pgfpathlineto{\pgfqpoint{2.755465in}{2.058541in}}%
\pgfpathclose%
\pgfusepath{stroke,fill}%
\end{pgfscope}%
\begin{pgfscope}%
\pgfpathrectangle{\pgfqpoint{0.050000in}{0.050000in}}{\pgfqpoint{4.900000in}{4.900000in}}%
\pgfusepath{clip}%
\pgfsetbuttcap%
\pgfsetroundjoin%
\definecolor{currentfill}{rgb}{1.000000,1.000000,1.000000}%
\pgfsetfillcolor{currentfill}%
\pgfsetfillopacity{0.900000}%
\pgfsetlinewidth{0.501875pt}%
\definecolor{currentstroke}{rgb}{0.200000,0.200000,0.200000}%
\pgfsetstrokecolor{currentstroke}%
\pgfsetdash{}{0pt}%
\pgfpathmoveto{\pgfqpoint{3.021111in}{2.044423in}}%
\pgfpathlineto{\pgfqpoint{3.064124in}{2.056941in}}%
\pgfpathlineto{\pgfqpoint{3.095369in}{2.038013in}}%
\pgfpathlineto{\pgfqpoint{3.052326in}{2.025597in}}%
\pgfpathlineto{\pgfqpoint{3.021111in}{2.044423in}}%
\pgfpathclose%
\pgfusepath{stroke,fill}%
\end{pgfscope}%
\begin{pgfscope}%
\pgfpathrectangle{\pgfqpoint{0.050000in}{0.050000in}}{\pgfqpoint{4.900000in}{4.900000in}}%
\pgfusepath{clip}%
\pgfsetbuttcap%
\pgfsetroundjoin%
\definecolor{currentfill}{rgb}{0.937993,0.937993,0.937993}%
\pgfsetfillcolor{currentfill}%
\pgfsetfillopacity{0.900000}%
\pgfsetlinewidth{0.501875pt}%
\definecolor{currentstroke}{rgb}{0.200000,0.200000,0.200000}%
\pgfsetstrokecolor{currentstroke}%
\pgfsetdash{}{0pt}%
\pgfpathmoveto{\pgfqpoint{2.088975in}{2.373808in}}%
\pgfpathlineto{\pgfqpoint{2.133411in}{2.331983in}}%
\pgfpathlineto{\pgfqpoint{2.163255in}{2.313138in}}%
\pgfpathlineto{\pgfqpoint{2.118816in}{2.354528in}}%
\pgfpathlineto{\pgfqpoint{2.088975in}{2.373808in}}%
\pgfpathclose%
\pgfusepath{stroke,fill}%
\end{pgfscope}%
\begin{pgfscope}%
\pgfpathrectangle{\pgfqpoint{0.050000in}{0.050000in}}{\pgfqpoint{4.900000in}{4.900000in}}%
\pgfusepath{clip}%
\pgfsetbuttcap%
\pgfsetroundjoin%
\definecolor{currentfill}{rgb}{1.000000,1.000000,1.000000}%
\pgfsetfillcolor{currentfill}%
\pgfsetfillopacity{0.900000}%
\pgfsetlinewidth{0.501875pt}%
\definecolor{currentstroke}{rgb}{0.200000,0.200000,0.200000}%
\pgfsetstrokecolor{currentstroke}%
\pgfsetdash{}{0pt}%
\pgfpathmoveto{\pgfqpoint{2.829583in}{2.048608in}}%
\pgfpathlineto{\pgfqpoint{2.872795in}{2.058502in}}%
\pgfpathlineto{\pgfqpoint{2.903764in}{2.040103in}}%
\pgfpathlineto{\pgfqpoint{2.860522in}{2.030508in}}%
\pgfpathlineto{\pgfqpoint{2.829583in}{2.048608in}}%
\pgfpathclose%
\pgfusepath{stroke,fill}%
\end{pgfscope}%
\begin{pgfscope}%
\pgfpathrectangle{\pgfqpoint{0.050000in}{0.050000in}}{\pgfqpoint{4.900000in}{4.900000in}}%
\pgfusepath{clip}%
\pgfsetbuttcap%
\pgfsetroundjoin%
\definecolor{currentfill}{rgb}{1.000000,1.000000,1.000000}%
\pgfsetfillcolor{currentfill}%
\pgfsetfillopacity{0.900000}%
\pgfsetlinewidth{0.501875pt}%
\definecolor{currentstroke}{rgb}{0.200000,0.200000,0.200000}%
\pgfsetstrokecolor{currentstroke}%
\pgfsetdash{}{0pt}%
\pgfpathmoveto{\pgfqpoint{3.286928in}{2.038965in}}%
\pgfpathlineto{\pgfqpoint{3.329682in}{2.052230in}}%
\pgfpathlineto{\pgfqpoint{3.361329in}{2.033082in}}%
\pgfpathlineto{\pgfqpoint{3.318546in}{2.019815in}}%
\pgfpathlineto{\pgfqpoint{3.286928in}{2.038965in}}%
\pgfpathclose%
\pgfusepath{stroke,fill}%
\end{pgfscope}%
\begin{pgfscope}%
\pgfpathrectangle{\pgfqpoint{0.050000in}{0.050000in}}{\pgfqpoint{4.900000in}{4.900000in}}%
\pgfusepath{clip}%
\pgfsetbuttcap%
\pgfsetroundjoin%
\definecolor{currentfill}{rgb}{0.981546,0.981546,0.981546}%
\pgfsetfillcolor{currentfill}%
\pgfsetfillopacity{0.900000}%
\pgfsetlinewidth{0.501875pt}%
\definecolor{currentstroke}{rgb}{0.200000,0.200000,0.200000}%
\pgfsetstrokecolor{currentstroke}%
\pgfsetdash{}{0pt}%
\pgfpathmoveto{\pgfqpoint{2.415750in}{2.147981in}}%
\pgfpathlineto{\pgfqpoint{2.459484in}{2.130891in}}%
\pgfpathlineto{\pgfqpoint{2.489822in}{2.114963in}}%
\pgfpathlineto{\pgfqpoint{2.446069in}{2.131764in}}%
\pgfpathlineto{\pgfqpoint{2.415750in}{2.147981in}}%
\pgfpathclose%
\pgfusepath{stroke,fill}%
\end{pgfscope}%
\begin{pgfscope}%
\pgfpathrectangle{\pgfqpoint{0.050000in}{0.050000in}}{\pgfqpoint{4.900000in}{4.900000in}}%
\pgfusepath{clip}%
\pgfsetbuttcap%
\pgfsetroundjoin%
\definecolor{currentfill}{rgb}{0.987082,0.987082,0.987082}%
\pgfsetfillcolor{currentfill}%
\pgfsetfillopacity{0.900000}%
\pgfsetlinewidth{0.501875pt}%
\definecolor{currentstroke}{rgb}{0.200000,0.200000,0.200000}%
\pgfsetstrokecolor{currentstroke}%
\pgfsetdash{}{0pt}%
\pgfpathmoveto{\pgfqpoint{2.489822in}{2.114963in}}%
\pgfpathlineto{\pgfqpoint{2.533433in}{2.105351in}}%
\pgfpathlineto{\pgfqpoint{2.563889in}{2.089131in}}%
\pgfpathlineto{\pgfqpoint{2.520255in}{2.098812in}}%
\pgfpathlineto{\pgfqpoint{2.489822in}{2.114963in}}%
\pgfpathclose%
\pgfusepath{stroke,fill}%
\end{pgfscope}%
\begin{pgfscope}%
\pgfpathrectangle{\pgfqpoint{0.050000in}{0.050000in}}{\pgfqpoint{4.900000in}{4.900000in}}%
\pgfusepath{clip}%
\pgfsetbuttcap%
\pgfsetroundjoin%
\definecolor{currentfill}{rgb}{0.974164,0.974164,0.974164}%
\pgfsetfillcolor{currentfill}%
\pgfsetfillopacity{0.900000}%
\pgfsetlinewidth{0.501875pt}%
\definecolor{currentstroke}{rgb}{0.200000,0.200000,0.200000}%
\pgfsetstrokecolor{currentstroke}%
\pgfsetdash{}{0pt}%
\pgfpathmoveto{\pgfqpoint{2.341639in}{2.189161in}}%
\pgfpathlineto{\pgfqpoint{2.385524in}{2.163987in}}%
\pgfpathlineto{\pgfqpoint{2.415750in}{2.147981in}}%
\pgfpathlineto{\pgfqpoint{2.371852in}{2.172452in}}%
\pgfpathlineto{\pgfqpoint{2.341639in}{2.189161in}}%
\pgfpathclose%
\pgfusepath{stroke,fill}%
\end{pgfscope}%
\begin{pgfscope}%
\pgfpathrectangle{\pgfqpoint{0.050000in}{0.050000in}}{\pgfqpoint{4.900000in}{4.900000in}}%
\pgfusepath{clip}%
\pgfsetbuttcap%
\pgfsetroundjoin%
\definecolor{currentfill}{rgb}{1.000000,1.000000,1.000000}%
\pgfsetfillcolor{currentfill}%
\pgfsetfillopacity{0.900000}%
\pgfsetlinewidth{0.501875pt}%
\definecolor{currentstroke}{rgb}{0.200000,0.200000,0.200000}%
\pgfsetstrokecolor{currentstroke}%
\pgfsetdash{}{0pt}%
\pgfpathmoveto{\pgfqpoint{3.095369in}{2.038013in}}%
\pgfpathlineto{\pgfqpoint{3.138326in}{2.050850in}}%
\pgfpathlineto{\pgfqpoint{3.169694in}{2.031828in}}%
\pgfpathlineto{\pgfqpoint{3.126708in}{2.019053in}}%
\pgfpathlineto{\pgfqpoint{3.095369in}{2.038013in}}%
\pgfpathclose%
\pgfusepath{stroke,fill}%
\end{pgfscope}%
\begin{pgfscope}%
\pgfpathrectangle{\pgfqpoint{0.050000in}{0.050000in}}{\pgfqpoint{4.900000in}{4.900000in}}%
\pgfusepath{clip}%
\pgfsetbuttcap%
\pgfsetroundjoin%
\definecolor{currentfill}{rgb}{0.990773,0.990773,0.990773}%
\pgfsetfillcolor{currentfill}%
\pgfsetfillopacity{0.900000}%
\pgfsetlinewidth{0.501875pt}%
\definecolor{currentstroke}{rgb}{0.200000,0.200000,0.200000}%
\pgfsetstrokecolor{currentstroke}%
\pgfsetdash{}{0pt}%
\pgfpathmoveto{\pgfqpoint{2.563889in}{2.089131in}}%
\pgfpathlineto{\pgfqpoint{2.607401in}{2.085774in}}%
\pgfpathlineto{\pgfqpoint{2.637976in}{2.069076in}}%
\pgfpathlineto{\pgfqpoint{2.594438in}{2.072726in}}%
\pgfpathlineto{\pgfqpoint{2.563889in}{2.089131in}}%
\pgfpathclose%
\pgfusepath{stroke,fill}%
\end{pgfscope}%
\begin{pgfscope}%
\pgfpathrectangle{\pgfqpoint{0.050000in}{0.050000in}}{\pgfqpoint{4.900000in}{4.900000in}}%
\pgfusepath{clip}%
\pgfsetbuttcap%
\pgfsetroundjoin%
\definecolor{currentfill}{rgb}{0.918185,0.918185,0.918185}%
\pgfsetfillcolor{currentfill}%
\pgfsetfillopacity{0.900000}%
\pgfsetlinewidth{0.501875pt}%
\definecolor{currentstroke}{rgb}{0.200000,0.200000,0.200000}%
\pgfsetstrokecolor{currentstroke}%
\pgfsetdash{}{0pt}%
\pgfpathmoveto{\pgfqpoint{2.014606in}{2.435560in}}%
\pgfpathlineto{\pgfqpoint{2.059220in}{2.393144in}}%
\pgfpathlineto{\pgfqpoint{2.088975in}{2.373808in}}%
\pgfpathlineto{\pgfqpoint{2.044355in}{2.416177in}}%
\pgfpathlineto{\pgfqpoint{2.014606in}{2.435560in}}%
\pgfpathclose%
\pgfusepath{stroke,fill}%
\end{pgfscope}%
\begin{pgfscope}%
\pgfpathrectangle{\pgfqpoint{0.050000in}{0.050000in}}{\pgfqpoint{4.900000in}{4.900000in}}%
\pgfusepath{clip}%
\pgfsetbuttcap%
\pgfsetroundjoin%
\definecolor{currentfill}{rgb}{1.000000,1.000000,1.000000}%
\pgfsetfillcolor{currentfill}%
\pgfsetfillopacity{0.900000}%
\pgfsetlinewidth{0.501875pt}%
\definecolor{currentstroke}{rgb}{0.200000,0.200000,0.200000}%
\pgfsetstrokecolor{currentstroke}%
\pgfsetdash{}{0pt}%
\pgfpathmoveto{\pgfqpoint{2.903764in}{2.040103in}}%
\pgfpathlineto{\pgfqpoint{2.946920in}{2.051186in}}%
\pgfpathlineto{\pgfqpoint{2.978011in}{2.032542in}}%
\pgfpathlineto{\pgfqpoint{2.934826in}{2.021682in}}%
\pgfpathlineto{\pgfqpoint{2.903764in}{2.040103in}}%
\pgfpathclose%
\pgfusepath{stroke,fill}%
\end{pgfscope}%
\begin{pgfscope}%
\pgfpathrectangle{\pgfqpoint{0.050000in}{0.050000in}}{\pgfqpoint{4.900000in}{4.900000in}}%
\pgfusepath{clip}%
\pgfsetbuttcap%
\pgfsetroundjoin%
\definecolor{currentfill}{rgb}{0.963091,0.963091,0.963091}%
\pgfsetfillcolor{currentfill}%
\pgfsetfillopacity{0.900000}%
\pgfsetlinewidth{0.501875pt}%
\definecolor{currentstroke}{rgb}{0.200000,0.200000,0.200000}%
\pgfsetstrokecolor{currentstroke}%
\pgfsetdash{}{0pt}%
\pgfpathmoveto{\pgfqpoint{2.267459in}{2.238476in}}%
\pgfpathlineto{\pgfqpoint{2.311520in}{2.205754in}}%
\pgfpathlineto{\pgfqpoint{2.341639in}{2.189161in}}%
\pgfpathlineto{\pgfqpoint{2.297572in}{2.220896in}}%
\pgfpathlineto{\pgfqpoint{2.267459in}{2.238476in}}%
\pgfpathclose%
\pgfusepath{stroke,fill}%
\end{pgfscope}%
\begin{pgfscope}%
\pgfpathrectangle{\pgfqpoint{0.050000in}{0.050000in}}{\pgfqpoint{4.900000in}{4.900000in}}%
\pgfusepath{clip}%
\pgfsetbuttcap%
\pgfsetroundjoin%
\definecolor{currentfill}{rgb}{1.000000,1.000000,1.000000}%
\pgfsetfillcolor{currentfill}%
\pgfsetfillopacity{0.900000}%
\pgfsetlinewidth{0.501875pt}%
\definecolor{currentstroke}{rgb}{0.200000,0.200000,0.200000}%
\pgfsetstrokecolor{currentstroke}%
\pgfsetdash{}{0pt}%
\pgfpathmoveto{\pgfqpoint{3.361329in}{2.033082in}}%
\pgfpathlineto{\pgfqpoint{3.404024in}{2.046378in}}%
\pgfpathlineto{\pgfqpoint{3.435797in}{2.027183in}}%
\pgfpathlineto{\pgfqpoint{3.393072in}{2.013881in}}%
\pgfpathlineto{\pgfqpoint{3.361329in}{2.033082in}}%
\pgfpathclose%
\pgfusepath{stroke,fill}%
\end{pgfscope}%
\begin{pgfscope}%
\pgfpathrectangle{\pgfqpoint{0.050000in}{0.050000in}}{\pgfqpoint{4.900000in}{4.900000in}}%
\pgfusepath{clip}%
\pgfsetbuttcap%
\pgfsetroundjoin%
\definecolor{currentfill}{rgb}{0.994464,0.994464,0.994464}%
\pgfsetfillcolor{currentfill}%
\pgfsetfillopacity{0.900000}%
\pgfsetlinewidth{0.501875pt}%
\definecolor{currentstroke}{rgb}{0.200000,0.200000,0.200000}%
\pgfsetstrokecolor{currentstroke}%
\pgfsetdash{}{0pt}%
\pgfpathmoveto{\pgfqpoint{2.637976in}{2.069076in}}%
\pgfpathlineto{\pgfqpoint{2.681407in}{2.070602in}}%
\pgfpathlineto{\pgfqpoint{2.712103in}{2.053384in}}%
\pgfpathlineto{\pgfqpoint{2.668645in}{2.052247in}}%
\pgfpathlineto{\pgfqpoint{2.637976in}{2.069076in}}%
\pgfpathclose%
\pgfusepath{stroke,fill}%
\end{pgfscope}%
\begin{pgfscope}%
\pgfpathrectangle{\pgfqpoint{0.050000in}{0.050000in}}{\pgfqpoint{4.900000in}{4.900000in}}%
\pgfusepath{clip}%
\pgfsetbuttcap%
\pgfsetroundjoin%
\definecolor{currentfill}{rgb}{1.000000,1.000000,1.000000}%
\pgfsetfillcolor{currentfill}%
\pgfsetfillopacity{0.900000}%
\pgfsetlinewidth{0.501875pt}%
\definecolor{currentstroke}{rgb}{0.200000,0.200000,0.200000}%
\pgfsetstrokecolor{currentstroke}%
\pgfsetdash{}{0pt}%
\pgfpathmoveto{\pgfqpoint{3.169694in}{2.031828in}}%
\pgfpathlineto{\pgfqpoint{3.212594in}{2.044874in}}%
\pgfpathlineto{\pgfqpoint{3.244087in}{2.025782in}}%
\pgfpathlineto{\pgfqpoint{3.201157in}{2.012768in}}%
\pgfpathlineto{\pgfqpoint{3.169694in}{2.031828in}}%
\pgfpathclose%
\pgfusepath{stroke,fill}%
\end{pgfscope}%
\begin{pgfscope}%
\pgfpathrectangle{\pgfqpoint{0.050000in}{0.050000in}}{\pgfqpoint{4.900000in}{4.900000in}}%
\pgfusepath{clip}%
\pgfsetbuttcap%
\pgfsetroundjoin%
\definecolor{currentfill}{rgb}{0.996309,0.996309,0.996309}%
\pgfsetfillcolor{currentfill}%
\pgfsetfillopacity{0.900000}%
\pgfsetlinewidth{0.501875pt}%
\definecolor{currentstroke}{rgb}{0.200000,0.200000,0.200000}%
\pgfsetstrokecolor{currentstroke}%
\pgfsetdash{}{0pt}%
\pgfpathmoveto{\pgfqpoint{2.712103in}{2.053384in}}%
\pgfpathlineto{\pgfqpoint{2.755465in}{2.058541in}}%
\pgfpathlineto{\pgfqpoint{2.786283in}{2.040843in}}%
\pgfpathlineto{\pgfqpoint{2.742893in}{2.036081in}}%
\pgfpathlineto{\pgfqpoint{2.712103in}{2.053384in}}%
\pgfpathclose%
\pgfusepath{stroke,fill}%
\end{pgfscope}%
\begin{pgfscope}%
\pgfpathrectangle{\pgfqpoint{0.050000in}{0.050000in}}{\pgfqpoint{4.900000in}{4.900000in}}%
\pgfusepath{clip}%
\pgfsetbuttcap%
\pgfsetroundjoin%
\definecolor{currentfill}{rgb}{1.000000,1.000000,1.000000}%
\pgfsetfillcolor{currentfill}%
\pgfsetfillopacity{0.900000}%
\pgfsetlinewidth{0.501875pt}%
\definecolor{currentstroke}{rgb}{0.200000,0.200000,0.200000}%
\pgfsetstrokecolor{currentstroke}%
\pgfsetdash{}{0pt}%
\pgfpathmoveto{\pgfqpoint{2.978011in}{2.032542in}}%
\pgfpathlineto{\pgfqpoint{3.021111in}{2.044423in}}%
\pgfpathlineto{\pgfqpoint{3.052326in}{2.025597in}}%
\pgfpathlineto{\pgfqpoint{3.009197in}{2.013874in}}%
\pgfpathlineto{\pgfqpoint{2.978011in}{2.032542in}}%
\pgfpathclose%
\pgfusepath{stroke,fill}%
\end{pgfscope}%
\begin{pgfscope}%
\pgfpathrectangle{\pgfqpoint{0.050000in}{0.050000in}}{\pgfqpoint{4.900000in}{4.900000in}}%
\pgfusepath{clip}%
\pgfsetbuttcap%
\pgfsetroundjoin%
\definecolor{currentfill}{rgb}{0.952018,0.952018,0.952018}%
\pgfsetfillcolor{currentfill}%
\pgfsetfillopacity{0.900000}%
\pgfsetlinewidth{0.501875pt}%
\definecolor{currentstroke}{rgb}{0.200000,0.200000,0.200000}%
\pgfsetstrokecolor{currentstroke}%
\pgfsetdash{}{0pt}%
\pgfpathmoveto{\pgfqpoint{2.193187in}{2.294464in}}%
\pgfpathlineto{\pgfqpoint{2.237437in}{2.256091in}}%
\pgfpathlineto{\pgfqpoint{2.267459in}{2.238476in}}%
\pgfpathlineto{\pgfqpoint{2.223206in}{2.275920in}}%
\pgfpathlineto{\pgfqpoint{2.193187in}{2.294464in}}%
\pgfpathclose%
\pgfusepath{stroke,fill}%
\end{pgfscope}%
\begin{pgfscope}%
\pgfpathrectangle{\pgfqpoint{0.050000in}{0.050000in}}{\pgfqpoint{4.900000in}{4.900000in}}%
\pgfusepath{clip}%
\pgfsetbuttcap%
\pgfsetroundjoin%
\definecolor{currentfill}{rgb}{1.000000,1.000000,1.000000}%
\pgfsetfillcolor{currentfill}%
\pgfsetfillopacity{0.900000}%
\pgfsetlinewidth{0.501875pt}%
\definecolor{currentstroke}{rgb}{0.200000,0.200000,0.200000}%
\pgfsetstrokecolor{currentstroke}%
\pgfsetdash{}{0pt}%
\pgfpathmoveto{\pgfqpoint{3.244087in}{2.025782in}}%
\pgfpathlineto{\pgfqpoint{3.286928in}{2.038965in}}%
\pgfpathlineto{\pgfqpoint{3.318546in}{2.019815in}}%
\pgfpathlineto{\pgfqpoint{3.275674in}{2.006646in}}%
\pgfpathlineto{\pgfqpoint{3.244087in}{2.025782in}}%
\pgfpathclose%
\pgfusepath{stroke,fill}%
\end{pgfscope}%
\begin{pgfscope}%
\pgfpathrectangle{\pgfqpoint{0.050000in}{0.050000in}}{\pgfqpoint{4.900000in}{4.900000in}}%
\pgfusepath{clip}%
\pgfsetbuttcap%
\pgfsetroundjoin%
\definecolor{currentfill}{rgb}{0.998155,0.998155,0.998155}%
\pgfsetfillcolor{currentfill}%
\pgfsetfillopacity{0.900000}%
\pgfsetlinewidth{0.501875pt}%
\definecolor{currentstroke}{rgb}{0.200000,0.200000,0.200000}%
\pgfsetstrokecolor{currentstroke}%
\pgfsetdash{}{0pt}%
\pgfpathmoveto{\pgfqpoint{2.786283in}{2.040843in}}%
\pgfpathlineto{\pgfqpoint{2.829583in}{2.048608in}}%
\pgfpathlineto{\pgfqpoint{2.860522in}{2.030508in}}%
\pgfpathlineto{\pgfqpoint{2.817194in}{2.023093in}}%
\pgfpathlineto{\pgfqpoint{2.786283in}{2.040843in}}%
\pgfpathclose%
\pgfusepath{stroke,fill}%
\end{pgfscope}%
\begin{pgfscope}%
\pgfpathrectangle{\pgfqpoint{0.050000in}{0.050000in}}{\pgfqpoint{4.900000in}{4.900000in}}%
\pgfusepath{clip}%
\pgfsetbuttcap%
\pgfsetroundjoin%
\definecolor{currentfill}{rgb}{1.000000,1.000000,1.000000}%
\pgfsetfillcolor{currentfill}%
\pgfsetfillopacity{0.900000}%
\pgfsetlinewidth{0.501875pt}%
\definecolor{currentstroke}{rgb}{0.200000,0.200000,0.200000}%
\pgfsetstrokecolor{currentstroke}%
\pgfsetdash{}{0pt}%
\pgfpathmoveto{\pgfqpoint{3.052326in}{2.025597in}}%
\pgfpathlineto{\pgfqpoint{3.095369in}{2.038013in}}%
\pgfpathlineto{\pgfqpoint{3.126708in}{2.019053in}}%
\pgfpathlineto{\pgfqpoint{3.083635in}{2.006743in}}%
\pgfpathlineto{\pgfqpoint{3.052326in}{2.025597in}}%
\pgfpathclose%
\pgfusepath{stroke,fill}%
\end{pgfscope}%
\begin{pgfscope}%
\pgfpathrectangle{\pgfqpoint{0.050000in}{0.050000in}}{\pgfqpoint{4.900000in}{4.900000in}}%
\pgfusepath{clip}%
\pgfsetbuttcap%
\pgfsetroundjoin%
\definecolor{currentfill}{rgb}{0.937993,0.937993,0.937993}%
\pgfsetfillcolor{currentfill}%
\pgfsetfillopacity{0.900000}%
\pgfsetlinewidth{0.501875pt}%
\definecolor{currentstroke}{rgb}{0.200000,0.200000,0.200000}%
\pgfsetstrokecolor{currentstroke}%
\pgfsetdash{}{0pt}%
\pgfpathmoveto{\pgfqpoint{2.118816in}{2.354528in}}%
\pgfpathlineto{\pgfqpoint{2.163255in}{2.313138in}}%
\pgfpathlineto{\pgfqpoint{2.193187in}{2.294464in}}%
\pgfpathlineto{\pgfqpoint{2.148745in}{2.335323in}}%
\pgfpathlineto{\pgfqpoint{2.118816in}{2.354528in}}%
\pgfpathclose%
\pgfusepath{stroke,fill}%
\end{pgfscope}%
\begin{pgfscope}%
\pgfpathrectangle{\pgfqpoint{0.050000in}{0.050000in}}{\pgfqpoint{4.900000in}{4.900000in}}%
\pgfusepath{clip}%
\pgfsetbuttcap%
\pgfsetroundjoin%
\definecolor{currentfill}{rgb}{0.998155,0.998155,0.998155}%
\pgfsetfillcolor{currentfill}%
\pgfsetfillopacity{0.900000}%
\pgfsetlinewidth{0.501875pt}%
\definecolor{currentstroke}{rgb}{0.200000,0.200000,0.200000}%
\pgfsetstrokecolor{currentstroke}%
\pgfsetdash{}{0pt}%
\pgfpathmoveto{\pgfqpoint{2.860522in}{2.030508in}}%
\pgfpathlineto{\pgfqpoint{2.903764in}{2.040103in}}%
\pgfpathlineto{\pgfqpoint{2.934826in}{2.021682in}}%
\pgfpathlineto{\pgfqpoint{2.891556in}{2.012372in}}%
\pgfpathlineto{\pgfqpoint{2.860522in}{2.030508in}}%
\pgfpathclose%
\pgfusepath{stroke,fill}%
\end{pgfscope}%
\begin{pgfscope}%
\pgfpathrectangle{\pgfqpoint{0.050000in}{0.050000in}}{\pgfqpoint{4.900000in}{4.900000in}}%
\pgfusepath{clip}%
\pgfsetbuttcap%
\pgfsetroundjoin%
\definecolor{currentfill}{rgb}{1.000000,1.000000,1.000000}%
\pgfsetfillcolor{currentfill}%
\pgfsetfillopacity{0.900000}%
\pgfsetlinewidth{0.501875pt}%
\definecolor{currentstroke}{rgb}{0.200000,0.200000,0.200000}%
\pgfsetstrokecolor{currentstroke}%
\pgfsetdash{}{0pt}%
\pgfpathmoveto{\pgfqpoint{3.318546in}{2.019815in}}%
\pgfpathlineto{\pgfqpoint{3.361329in}{2.033082in}}%
\pgfpathlineto{\pgfqpoint{3.393072in}{2.013881in}}%
\pgfpathlineto{\pgfqpoint{3.350258in}{2.000618in}}%
\pgfpathlineto{\pgfqpoint{3.318546in}{2.019815in}}%
\pgfpathclose%
\pgfusepath{stroke,fill}%
\end{pgfscope}%
\begin{pgfscope}%
\pgfpathrectangle{\pgfqpoint{0.050000in}{0.050000in}}{\pgfqpoint{4.900000in}{4.900000in}}%
\pgfusepath{clip}%
\pgfsetbuttcap%
\pgfsetroundjoin%
\definecolor{currentfill}{rgb}{0.981546,0.981546,0.981546}%
\pgfsetfillcolor{currentfill}%
\pgfsetfillopacity{0.900000}%
\pgfsetlinewidth{0.501875pt}%
\definecolor{currentstroke}{rgb}{0.200000,0.200000,0.200000}%
\pgfsetstrokecolor{currentstroke}%
\pgfsetdash{}{0pt}%
\pgfpathmoveto{\pgfqpoint{2.446069in}{2.131764in}}%
\pgfpathlineto{\pgfqpoint{2.489822in}{2.114963in}}%
\pgfpathlineto{\pgfqpoint{2.520255in}{2.098812in}}%
\pgfpathlineto{\pgfqpoint{2.476483in}{2.115356in}}%
\pgfpathlineto{\pgfqpoint{2.446069in}{2.131764in}}%
\pgfpathclose%
\pgfusepath{stroke,fill}%
\end{pgfscope}%
\begin{pgfscope}%
\pgfpathrectangle{\pgfqpoint{0.050000in}{0.050000in}}{\pgfqpoint{4.900000in}{4.900000in}}%
\pgfusepath{clip}%
\pgfsetbuttcap%
\pgfsetroundjoin%
\definecolor{currentfill}{rgb}{0.987082,0.987082,0.987082}%
\pgfsetfillcolor{currentfill}%
\pgfsetfillopacity{0.900000}%
\pgfsetlinewidth{0.501875pt}%
\definecolor{currentstroke}{rgb}{0.200000,0.200000,0.200000}%
\pgfsetstrokecolor{currentstroke}%
\pgfsetdash{}{0pt}%
\pgfpathmoveto{\pgfqpoint{2.520255in}{2.098812in}}%
\pgfpathlineto{\pgfqpoint{2.563889in}{2.089131in}}%
\pgfpathlineto{\pgfqpoint{2.594438in}{2.072726in}}%
\pgfpathlineto{\pgfqpoint{2.550781in}{2.082460in}}%
\pgfpathlineto{\pgfqpoint{2.520255in}{2.098812in}}%
\pgfpathclose%
\pgfusepath{stroke,fill}%
\end{pgfscope}%
\begin{pgfscope}%
\pgfpathrectangle{\pgfqpoint{0.050000in}{0.050000in}}{\pgfqpoint{4.900000in}{4.900000in}}%
\pgfusepath{clip}%
\pgfsetbuttcap%
\pgfsetroundjoin%
\definecolor{currentfill}{rgb}{1.000000,1.000000,1.000000}%
\pgfsetfillcolor{currentfill}%
\pgfsetfillopacity{0.900000}%
\pgfsetlinewidth{0.501875pt}%
\definecolor{currentstroke}{rgb}{0.200000,0.200000,0.200000}%
\pgfsetstrokecolor{currentstroke}%
\pgfsetdash{}{0pt}%
\pgfpathmoveto{\pgfqpoint{3.126708in}{2.019053in}}%
\pgfpathlineto{\pgfqpoint{3.169694in}{2.031828in}}%
\pgfpathlineto{\pgfqpoint{3.201157in}{2.012768in}}%
\pgfpathlineto{\pgfqpoint{3.158141in}{2.000061in}}%
\pgfpathlineto{\pgfqpoint{3.126708in}{2.019053in}}%
\pgfpathclose%
\pgfusepath{stroke,fill}%
\end{pgfscope}%
\begin{pgfscope}%
\pgfpathrectangle{\pgfqpoint{0.050000in}{0.050000in}}{\pgfqpoint{4.900000in}{4.900000in}}%
\pgfusepath{clip}%
\pgfsetbuttcap%
\pgfsetroundjoin%
\definecolor{currentfill}{rgb}{0.972318,0.972318,0.972318}%
\pgfsetfillcolor{currentfill}%
\pgfsetfillopacity{0.900000}%
\pgfsetlinewidth{0.501875pt}%
\definecolor{currentstroke}{rgb}{0.200000,0.200000,0.200000}%
\pgfsetstrokecolor{currentstroke}%
\pgfsetdash{}{0pt}%
\pgfpathmoveto{\pgfqpoint{2.371852in}{2.172452in}}%
\pgfpathlineto{\pgfqpoint{2.415750in}{2.147981in}}%
\pgfpathlineto{\pgfqpoint{2.446069in}{2.131764in}}%
\pgfpathlineto{\pgfqpoint{2.402157in}{2.155623in}}%
\pgfpathlineto{\pgfqpoint{2.371852in}{2.172452in}}%
\pgfpathclose%
\pgfusepath{stroke,fill}%
\end{pgfscope}%
\begin{pgfscope}%
\pgfpathrectangle{\pgfqpoint{0.050000in}{0.050000in}}{\pgfqpoint{4.900000in}{4.900000in}}%
\pgfusepath{clip}%
\pgfsetbuttcap%
\pgfsetroundjoin%
\definecolor{currentfill}{rgb}{0.990773,0.990773,0.990773}%
\pgfsetfillcolor{currentfill}%
\pgfsetfillopacity{0.900000}%
\pgfsetlinewidth{0.501875pt}%
\definecolor{currentstroke}{rgb}{0.200000,0.200000,0.200000}%
\pgfsetstrokecolor{currentstroke}%
\pgfsetdash{}{0pt}%
\pgfpathmoveto{\pgfqpoint{2.594438in}{2.072726in}}%
\pgfpathlineto{\pgfqpoint{2.637976in}{2.069076in}}%
\pgfpathlineto{\pgfqpoint{2.668645in}{2.052247in}}%
\pgfpathlineto{\pgfqpoint{2.625082in}{2.056151in}}%
\pgfpathlineto{\pgfqpoint{2.594438in}{2.072726in}}%
\pgfpathclose%
\pgfusepath{stroke,fill}%
\end{pgfscope}%
\begin{pgfscope}%
\pgfpathrectangle{\pgfqpoint{0.050000in}{0.050000in}}{\pgfqpoint{4.900000in}{4.900000in}}%
\pgfusepath{clip}%
\pgfsetbuttcap%
\pgfsetroundjoin%
\definecolor{currentfill}{rgb}{0.918185,0.918185,0.918185}%
\pgfsetfillcolor{currentfill}%
\pgfsetfillopacity{0.900000}%
\pgfsetlinewidth{0.501875pt}%
\definecolor{currentstroke}{rgb}{0.200000,0.200000,0.200000}%
\pgfsetstrokecolor{currentstroke}%
\pgfsetdash{}{0pt}%
\pgfpathmoveto{\pgfqpoint{2.044355in}{2.416177in}}%
\pgfpathlineto{\pgfqpoint{2.088975in}{2.373808in}}%
\pgfpathlineto{\pgfqpoint{2.118816in}{2.354528in}}%
\pgfpathlineto{\pgfqpoint{2.074192in}{2.396768in}}%
\pgfpathlineto{\pgfqpoint{2.044355in}{2.416177in}}%
\pgfpathclose%
\pgfusepath{stroke,fill}%
\end{pgfscope}%
\begin{pgfscope}%
\pgfpathrectangle{\pgfqpoint{0.050000in}{0.050000in}}{\pgfqpoint{4.900000in}{4.900000in}}%
\pgfusepath{clip}%
\pgfsetbuttcap%
\pgfsetroundjoin%
\definecolor{currentfill}{rgb}{1.000000,1.000000,1.000000}%
\pgfsetfillcolor{currentfill}%
\pgfsetfillopacity{0.900000}%
\pgfsetlinewidth{0.501875pt}%
\definecolor{currentstroke}{rgb}{0.200000,0.200000,0.200000}%
\pgfsetstrokecolor{currentstroke}%
\pgfsetdash{}{0pt}%
\pgfpathmoveto{\pgfqpoint{2.934826in}{2.021682in}}%
\pgfpathlineto{\pgfqpoint{2.978011in}{2.032542in}}%
\pgfpathlineto{\pgfqpoint{3.009197in}{2.013874in}}%
\pgfpathlineto{\pgfqpoint{2.965983in}{2.003233in}}%
\pgfpathlineto{\pgfqpoint{2.934826in}{2.021682in}}%
\pgfpathclose%
\pgfusepath{stroke,fill}%
\end{pgfscope}%
\begin{pgfscope}%
\pgfpathrectangle{\pgfqpoint{0.050000in}{0.050000in}}{\pgfqpoint{4.900000in}{4.900000in}}%
\pgfusepath{clip}%
\pgfsetbuttcap%
\pgfsetroundjoin%
\definecolor{currentfill}{rgb}{1.000000,1.000000,1.000000}%
\pgfsetfillcolor{currentfill}%
\pgfsetfillopacity{0.900000}%
\pgfsetlinewidth{0.501875pt}%
\definecolor{currentstroke}{rgb}{0.200000,0.200000,0.200000}%
\pgfsetstrokecolor{currentstroke}%
\pgfsetdash{}{0pt}%
\pgfpathmoveto{\pgfqpoint{3.393072in}{2.013881in}}%
\pgfpathlineto{\pgfqpoint{3.435797in}{2.027183in}}%
\pgfpathlineto{\pgfqpoint{3.467664in}{2.007930in}}%
\pgfpathlineto{\pgfqpoint{3.424909in}{1.994628in}}%
\pgfpathlineto{\pgfqpoint{3.393072in}{2.013881in}}%
\pgfpathclose%
\pgfusepath{stroke,fill}%
\end{pgfscope}%
\begin{pgfscope}%
\pgfpathrectangle{\pgfqpoint{0.050000in}{0.050000in}}{\pgfqpoint{4.900000in}{4.900000in}}%
\pgfusepath{clip}%
\pgfsetbuttcap%
\pgfsetroundjoin%
\definecolor{currentfill}{rgb}{0.963091,0.963091,0.963091}%
\pgfsetfillcolor{currentfill}%
\pgfsetfillopacity{0.900000}%
\pgfsetlinewidth{0.501875pt}%
\definecolor{currentstroke}{rgb}{0.200000,0.200000,0.200000}%
\pgfsetstrokecolor{currentstroke}%
\pgfsetdash{}{0pt}%
\pgfpathmoveto{\pgfqpoint{2.297572in}{2.220896in}}%
\pgfpathlineto{\pgfqpoint{2.341639in}{2.189161in}}%
\pgfpathlineto{\pgfqpoint{2.371852in}{2.172452in}}%
\pgfpathlineto{\pgfqpoint{2.327776in}{2.203314in}}%
\pgfpathlineto{\pgfqpoint{2.297572in}{2.220896in}}%
\pgfpathclose%
\pgfusepath{stroke,fill}%
\end{pgfscope}%
\begin{pgfscope}%
\pgfpathrectangle{\pgfqpoint{0.050000in}{0.050000in}}{\pgfqpoint{4.900000in}{4.900000in}}%
\pgfusepath{clip}%
\pgfsetbuttcap%
\pgfsetroundjoin%
\definecolor{currentfill}{rgb}{0.994464,0.994464,0.994464}%
\pgfsetfillcolor{currentfill}%
\pgfsetfillopacity{0.900000}%
\pgfsetlinewidth{0.501875pt}%
\definecolor{currentstroke}{rgb}{0.200000,0.200000,0.200000}%
\pgfsetstrokecolor{currentstroke}%
\pgfsetdash{}{0pt}%
\pgfpathmoveto{\pgfqpoint{2.668645in}{2.052247in}}%
\pgfpathlineto{\pgfqpoint{2.712103in}{2.053384in}}%
\pgfpathlineto{\pgfqpoint{2.742893in}{2.036081in}}%
\pgfpathlineto{\pgfqpoint{2.699408in}{2.035291in}}%
\pgfpathlineto{\pgfqpoint{2.668645in}{2.052247in}}%
\pgfpathclose%
\pgfusepath{stroke,fill}%
\end{pgfscope}%
\begin{pgfscope}%
\pgfpathrectangle{\pgfqpoint{0.050000in}{0.050000in}}{\pgfqpoint{4.900000in}{4.900000in}}%
\pgfusepath{clip}%
\pgfsetbuttcap%
\pgfsetroundjoin%
\definecolor{currentfill}{rgb}{1.000000,1.000000,1.000000}%
\pgfsetfillcolor{currentfill}%
\pgfsetfillopacity{0.900000}%
\pgfsetlinewidth{0.501875pt}%
\definecolor{currentstroke}{rgb}{0.200000,0.200000,0.200000}%
\pgfsetstrokecolor{currentstroke}%
\pgfsetdash{}{0pt}%
\pgfpathmoveto{\pgfqpoint{3.201157in}{2.012768in}}%
\pgfpathlineto{\pgfqpoint{3.244087in}{2.025782in}}%
\pgfpathlineto{\pgfqpoint{3.275674in}{2.006646in}}%
\pgfpathlineto{\pgfqpoint{3.232715in}{1.993670in}}%
\pgfpathlineto{\pgfqpoint{3.201157in}{2.012768in}}%
\pgfpathclose%
\pgfusepath{stroke,fill}%
\end{pgfscope}%
\begin{pgfscope}%
\pgfpathrectangle{\pgfqpoint{0.050000in}{0.050000in}}{\pgfqpoint{4.900000in}{4.900000in}}%
\pgfusepath{clip}%
\pgfsetbuttcap%
\pgfsetroundjoin%
\definecolor{currentfill}{rgb}{0.996309,0.996309,0.996309}%
\pgfsetfillcolor{currentfill}%
\pgfsetfillopacity{0.900000}%
\pgfsetlinewidth{0.501875pt}%
\definecolor{currentstroke}{rgb}{0.200000,0.200000,0.200000}%
\pgfsetstrokecolor{currentstroke}%
\pgfsetdash{}{0pt}%
\pgfpathmoveto{\pgfqpoint{2.742893in}{2.036081in}}%
\pgfpathlineto{\pgfqpoint{2.786283in}{2.040843in}}%
\pgfpathlineto{\pgfqpoint{2.817194in}{2.023093in}}%
\pgfpathlineto{\pgfqpoint{2.773777in}{2.018692in}}%
\pgfpathlineto{\pgfqpoint{2.742893in}{2.036081in}}%
\pgfpathclose%
\pgfusepath{stroke,fill}%
\end{pgfscope}%
\begin{pgfscope}%
\pgfpathrectangle{\pgfqpoint{0.050000in}{0.050000in}}{\pgfqpoint{4.900000in}{4.900000in}}%
\pgfusepath{clip}%
\pgfsetbuttcap%
\pgfsetroundjoin%
\definecolor{currentfill}{rgb}{1.000000,1.000000,1.000000}%
\pgfsetfillcolor{currentfill}%
\pgfsetfillopacity{0.900000}%
\pgfsetlinewidth{0.501875pt}%
\definecolor{currentstroke}{rgb}{0.200000,0.200000,0.200000}%
\pgfsetstrokecolor{currentstroke}%
\pgfsetdash{}{0pt}%
\pgfpathmoveto{\pgfqpoint{3.009197in}{2.013874in}}%
\pgfpathlineto{\pgfqpoint{3.052326in}{2.025597in}}%
\pgfpathlineto{\pgfqpoint{3.083635in}{2.006743in}}%
\pgfpathlineto{\pgfqpoint{3.040476in}{1.995179in}}%
\pgfpathlineto{\pgfqpoint{3.009197in}{2.013874in}}%
\pgfpathclose%
\pgfusepath{stroke,fill}%
\end{pgfscope}%
\begin{pgfscope}%
\pgfpathrectangle{\pgfqpoint{0.050000in}{0.050000in}}{\pgfqpoint{4.900000in}{4.900000in}}%
\pgfusepath{clip}%
\pgfsetbuttcap%
\pgfsetroundjoin%
\definecolor{currentfill}{rgb}{0.952018,0.952018,0.952018}%
\pgfsetfillcolor{currentfill}%
\pgfsetfillopacity{0.900000}%
\pgfsetlinewidth{0.501875pt}%
\definecolor{currentstroke}{rgb}{0.200000,0.200000,0.200000}%
\pgfsetstrokecolor{currentstroke}%
\pgfsetdash{}{0pt}%
\pgfpathmoveto{\pgfqpoint{2.223206in}{2.275920in}}%
\pgfpathlineto{\pgfqpoint{2.267459in}{2.238476in}}%
\pgfpathlineto{\pgfqpoint{2.297572in}{2.220896in}}%
\pgfpathlineto{\pgfqpoint{2.253316in}{2.257466in}}%
\pgfpathlineto{\pgfqpoint{2.223206in}{2.275920in}}%
\pgfpathclose%
\pgfusepath{stroke,fill}%
\end{pgfscope}%
\begin{pgfscope}%
\pgfpathrectangle{\pgfqpoint{0.050000in}{0.050000in}}{\pgfqpoint{4.900000in}{4.900000in}}%
\pgfusepath{clip}%
\pgfsetbuttcap%
\pgfsetroundjoin%
\definecolor{currentfill}{rgb}{0.895548,0.895548,0.895548}%
\pgfsetfillcolor{currentfill}%
\pgfsetfillopacity{0.900000}%
\pgfsetlinewidth{0.501875pt}%
\definecolor{currentstroke}{rgb}{0.200000,0.200000,0.200000}%
\pgfsetstrokecolor{currentstroke}%
\pgfsetdash{}{0pt}%
\pgfpathmoveto{\pgfqpoint{1.969811in}{2.478164in}}%
\pgfpathlineto{\pgfqpoint{2.014606in}{2.435560in}}%
\pgfpathlineto{\pgfqpoint{2.044355in}{2.416177in}}%
\pgfpathlineto{\pgfqpoint{1.999554in}{2.458794in}}%
\pgfpathlineto{\pgfqpoint{1.969811in}{2.478164in}}%
\pgfpathclose%
\pgfusepath{stroke,fill}%
\end{pgfscope}%
\begin{pgfscope}%
\pgfpathrectangle{\pgfqpoint{0.050000in}{0.050000in}}{\pgfqpoint{4.900000in}{4.900000in}}%
\pgfusepath{clip}%
\pgfsetbuttcap%
\pgfsetroundjoin%
\definecolor{currentfill}{rgb}{1.000000,1.000000,1.000000}%
\pgfsetfillcolor{currentfill}%
\pgfsetfillopacity{0.900000}%
\pgfsetlinewidth{0.501875pt}%
\definecolor{currentstroke}{rgb}{0.200000,0.200000,0.200000}%
\pgfsetstrokecolor{currentstroke}%
\pgfsetdash{}{0pt}%
\pgfpathmoveto{\pgfqpoint{3.275674in}{2.006646in}}%
\pgfpathlineto{\pgfqpoint{3.318546in}{2.019815in}}%
\pgfpathlineto{\pgfqpoint{3.350258in}{2.000618in}}%
\pgfpathlineto{\pgfqpoint{3.307357in}{1.987467in}}%
\pgfpathlineto{\pgfqpoint{3.275674in}{2.006646in}}%
\pgfpathclose%
\pgfusepath{stroke,fill}%
\end{pgfscope}%
\begin{pgfscope}%
\pgfpathrectangle{\pgfqpoint{0.050000in}{0.050000in}}{\pgfqpoint{4.900000in}{4.900000in}}%
\pgfusepath{clip}%
\pgfsetbuttcap%
\pgfsetroundjoin%
\definecolor{currentfill}{rgb}{0.998155,0.998155,0.998155}%
\pgfsetfillcolor{currentfill}%
\pgfsetfillopacity{0.900000}%
\pgfsetlinewidth{0.501875pt}%
\definecolor{currentstroke}{rgb}{0.200000,0.200000,0.200000}%
\pgfsetstrokecolor{currentstroke}%
\pgfsetdash{}{0pt}%
\pgfpathmoveto{\pgfqpoint{2.817194in}{2.023093in}}%
\pgfpathlineto{\pgfqpoint{2.860522in}{2.030508in}}%
\pgfpathlineto{\pgfqpoint{2.891556in}{2.012372in}}%
\pgfpathlineto{\pgfqpoint{2.848199in}{2.005285in}}%
\pgfpathlineto{\pgfqpoint{2.817194in}{2.023093in}}%
\pgfpathclose%
\pgfusepath{stroke,fill}%
\end{pgfscope}%
\begin{pgfscope}%
\pgfpathrectangle{\pgfqpoint{0.050000in}{0.050000in}}{\pgfqpoint{4.900000in}{4.900000in}}%
\pgfusepath{clip}%
\pgfsetbuttcap%
\pgfsetroundjoin%
\definecolor{currentfill}{rgb}{1.000000,1.000000,1.000000}%
\pgfsetfillcolor{currentfill}%
\pgfsetfillopacity{0.900000}%
\pgfsetlinewidth{0.501875pt}%
\definecolor{currentstroke}{rgb}{0.200000,0.200000,0.200000}%
\pgfsetstrokecolor{currentstroke}%
\pgfsetdash{}{0pt}%
\pgfpathmoveto{\pgfqpoint{3.083635in}{2.006743in}}%
\pgfpathlineto{\pgfqpoint{3.126708in}{2.019053in}}%
\pgfpathlineto{\pgfqpoint{3.158141in}{2.000061in}}%
\pgfpathlineto{\pgfqpoint{3.115038in}{1.987860in}}%
\pgfpathlineto{\pgfqpoint{3.083635in}{2.006743in}}%
\pgfpathclose%
\pgfusepath{stroke,fill}%
\end{pgfscope}%
\begin{pgfscope}%
\pgfpathrectangle{\pgfqpoint{0.050000in}{0.050000in}}{\pgfqpoint{4.900000in}{4.900000in}}%
\pgfusepath{clip}%
\pgfsetbuttcap%
\pgfsetroundjoin%
\definecolor{currentfill}{rgb}{0.937993,0.937993,0.937993}%
\pgfsetfillcolor{currentfill}%
\pgfsetfillopacity{0.900000}%
\pgfsetlinewidth{0.501875pt}%
\definecolor{currentstroke}{rgb}{0.200000,0.200000,0.200000}%
\pgfsetstrokecolor{currentstroke}%
\pgfsetdash{}{0pt}%
\pgfpathmoveto{\pgfqpoint{2.148745in}{2.335323in}}%
\pgfpathlineto{\pgfqpoint{2.193187in}{2.294464in}}%
\pgfpathlineto{\pgfqpoint{2.223206in}{2.275920in}}%
\pgfpathlineto{\pgfqpoint{2.178763in}{2.316197in}}%
\pgfpathlineto{\pgfqpoint{2.148745in}{2.335323in}}%
\pgfpathclose%
\pgfusepath{stroke,fill}%
\end{pgfscope}%
\begin{pgfscope}%
\pgfpathrectangle{\pgfqpoint{0.050000in}{0.050000in}}{\pgfqpoint{4.900000in}{4.900000in}}%
\pgfusepath{clip}%
\pgfsetbuttcap%
\pgfsetroundjoin%
\definecolor{currentfill}{rgb}{1.000000,1.000000,1.000000}%
\pgfsetfillcolor{currentfill}%
\pgfsetfillopacity{0.900000}%
\pgfsetlinewidth{0.501875pt}%
\definecolor{currentstroke}{rgb}{0.200000,0.200000,0.200000}%
\pgfsetstrokecolor{currentstroke}%
\pgfsetdash{}{0pt}%
\pgfpathmoveto{\pgfqpoint{3.350258in}{2.000618in}}%
\pgfpathlineto{\pgfqpoint{3.393072in}{2.013881in}}%
\pgfpathlineto{\pgfqpoint{3.424909in}{1.994628in}}%
\pgfpathlineto{\pgfqpoint{3.382066in}{1.981372in}}%
\pgfpathlineto{\pgfqpoint{3.350258in}{2.000618in}}%
\pgfpathclose%
\pgfusepath{stroke,fill}%
\end{pgfscope}%
\begin{pgfscope}%
\pgfpathrectangle{\pgfqpoint{0.050000in}{0.050000in}}{\pgfqpoint{4.900000in}{4.900000in}}%
\pgfusepath{clip}%
\pgfsetbuttcap%
\pgfsetroundjoin%
\definecolor{currentfill}{rgb}{0.998155,0.998155,0.998155}%
\pgfsetfillcolor{currentfill}%
\pgfsetfillopacity{0.900000}%
\pgfsetlinewidth{0.501875pt}%
\definecolor{currentstroke}{rgb}{0.200000,0.200000,0.200000}%
\pgfsetstrokecolor{currentstroke}%
\pgfsetdash{}{0pt}%
\pgfpathmoveto{\pgfqpoint{2.891556in}{2.012372in}}%
\pgfpathlineto{\pgfqpoint{2.934826in}{2.021682in}}%
\pgfpathlineto{\pgfqpoint{2.965983in}{2.003233in}}%
\pgfpathlineto{\pgfqpoint{2.922684in}{1.994196in}}%
\pgfpathlineto{\pgfqpoint{2.891556in}{2.012372in}}%
\pgfpathclose%
\pgfusepath{stroke,fill}%
\end{pgfscope}%
\begin{pgfscope}%
\pgfpathrectangle{\pgfqpoint{0.050000in}{0.050000in}}{\pgfqpoint{4.900000in}{4.900000in}}%
\pgfusepath{clip}%
\pgfsetbuttcap%
\pgfsetroundjoin%
\definecolor{currentfill}{rgb}{0.979700,0.979700,0.979700}%
\pgfsetfillcolor{currentfill}%
\pgfsetfillopacity{0.900000}%
\pgfsetlinewidth{0.501875pt}%
\definecolor{currentstroke}{rgb}{0.200000,0.200000,0.200000}%
\pgfsetstrokecolor{currentstroke}%
\pgfsetdash{}{0pt}%
\pgfpathmoveto{\pgfqpoint{2.476483in}{2.115356in}}%
\pgfpathlineto{\pgfqpoint{2.520255in}{2.098812in}}%
\pgfpathlineto{\pgfqpoint{2.550781in}{2.082460in}}%
\pgfpathlineto{\pgfqpoint{2.506990in}{2.098774in}}%
\pgfpathlineto{\pgfqpoint{2.476483in}{2.115356in}}%
\pgfpathclose%
\pgfusepath{stroke,fill}%
\end{pgfscope}%
\begin{pgfscope}%
\pgfpathrectangle{\pgfqpoint{0.050000in}{0.050000in}}{\pgfqpoint{4.900000in}{4.900000in}}%
\pgfusepath{clip}%
\pgfsetbuttcap%
\pgfsetroundjoin%
\definecolor{currentfill}{rgb}{1.000000,1.000000,1.000000}%
\pgfsetfillcolor{currentfill}%
\pgfsetfillopacity{0.900000}%
\pgfsetlinewidth{0.501875pt}%
\definecolor{currentstroke}{rgb}{0.200000,0.200000,0.200000}%
\pgfsetstrokecolor{currentstroke}%
\pgfsetdash{}{0pt}%
\pgfpathmoveto{\pgfqpoint{3.158141in}{2.000061in}}%
\pgfpathlineto{\pgfqpoint{3.201157in}{2.012768in}}%
\pgfpathlineto{\pgfqpoint{3.232715in}{1.993670in}}%
\pgfpathlineto{\pgfqpoint{3.189669in}{1.981034in}}%
\pgfpathlineto{\pgfqpoint{3.158141in}{2.000061in}}%
\pgfpathclose%
\pgfusepath{stroke,fill}%
\end{pgfscope}%
\begin{pgfscope}%
\pgfpathrectangle{\pgfqpoint{0.050000in}{0.050000in}}{\pgfqpoint{4.900000in}{4.900000in}}%
\pgfusepath{clip}%
\pgfsetbuttcap%
\pgfsetroundjoin%
\definecolor{currentfill}{rgb}{0.985236,0.985236,0.985236}%
\pgfsetfillcolor{currentfill}%
\pgfsetfillopacity{0.900000}%
\pgfsetlinewidth{0.501875pt}%
\definecolor{currentstroke}{rgb}{0.200000,0.200000,0.200000}%
\pgfsetstrokecolor{currentstroke}%
\pgfsetdash{}{0pt}%
\pgfpathmoveto{\pgfqpoint{2.550781in}{2.082460in}}%
\pgfpathlineto{\pgfqpoint{2.594438in}{2.072726in}}%
\pgfpathlineto{\pgfqpoint{2.625082in}{2.056151in}}%
\pgfpathlineto{\pgfqpoint{2.581403in}{2.065925in}}%
\pgfpathlineto{\pgfqpoint{2.550781in}{2.082460in}}%
\pgfpathclose%
\pgfusepath{stroke,fill}%
\end{pgfscope}%
\begin{pgfscope}%
\pgfpathrectangle{\pgfqpoint{0.050000in}{0.050000in}}{\pgfqpoint{4.900000in}{4.900000in}}%
\pgfusepath{clip}%
\pgfsetbuttcap%
\pgfsetroundjoin%
\definecolor{currentfill}{rgb}{0.972318,0.972318,0.972318}%
\pgfsetfillcolor{currentfill}%
\pgfsetfillopacity{0.900000}%
\pgfsetlinewidth{0.501875pt}%
\definecolor{currentstroke}{rgb}{0.200000,0.200000,0.200000}%
\pgfsetstrokecolor{currentstroke}%
\pgfsetdash{}{0pt}%
\pgfpathmoveto{\pgfqpoint{2.402157in}{2.155623in}}%
\pgfpathlineto{\pgfqpoint{2.446069in}{2.131764in}}%
\pgfpathlineto{\pgfqpoint{2.476483in}{2.115356in}}%
\pgfpathlineto{\pgfqpoint{2.432556in}{2.138676in}}%
\pgfpathlineto{\pgfqpoint{2.402157in}{2.155623in}}%
\pgfpathclose%
\pgfusepath{stroke,fill}%
\end{pgfscope}%
\begin{pgfscope}%
\pgfpathrectangle{\pgfqpoint{0.050000in}{0.050000in}}{\pgfqpoint{4.900000in}{4.900000in}}%
\pgfusepath{clip}%
\pgfsetbuttcap%
\pgfsetroundjoin%
\definecolor{currentfill}{rgb}{0.918185,0.918185,0.918185}%
\pgfsetfillcolor{currentfill}%
\pgfsetfillopacity{0.900000}%
\pgfsetlinewidth{0.501875pt}%
\definecolor{currentstroke}{rgb}{0.200000,0.200000,0.200000}%
\pgfsetstrokecolor{currentstroke}%
\pgfsetdash{}{0pt}%
\pgfpathmoveto{\pgfqpoint{2.074192in}{2.396768in}}%
\pgfpathlineto{\pgfqpoint{2.118816in}{2.354528in}}%
\pgfpathlineto{\pgfqpoint{2.148745in}{2.335323in}}%
\pgfpathlineto{\pgfqpoint{2.104118in}{2.377355in}}%
\pgfpathlineto{\pgfqpoint{2.074192in}{2.396768in}}%
\pgfpathclose%
\pgfusepath{stroke,fill}%
\end{pgfscope}%
\begin{pgfscope}%
\pgfpathrectangle{\pgfqpoint{0.050000in}{0.050000in}}{\pgfqpoint{4.900000in}{4.900000in}}%
\pgfusepath{clip}%
\pgfsetbuttcap%
\pgfsetroundjoin%
\definecolor{currentfill}{rgb}{0.990773,0.990773,0.990773}%
\pgfsetfillcolor{currentfill}%
\pgfsetfillopacity{0.900000}%
\pgfsetlinewidth{0.501875pt}%
\definecolor{currentstroke}{rgb}{0.200000,0.200000,0.200000}%
\pgfsetstrokecolor{currentstroke}%
\pgfsetdash{}{0pt}%
\pgfpathmoveto{\pgfqpoint{2.625082in}{2.056151in}}%
\pgfpathlineto{\pgfqpoint{2.668645in}{2.052247in}}%
\pgfpathlineto{\pgfqpoint{2.699408in}{2.035291in}}%
\pgfpathlineto{\pgfqpoint{2.655820in}{2.039418in}}%
\pgfpathlineto{\pgfqpoint{2.625082in}{2.056151in}}%
\pgfpathclose%
\pgfusepath{stroke,fill}%
\end{pgfscope}%
\begin{pgfscope}%
\pgfpathrectangle{\pgfqpoint{0.050000in}{0.050000in}}{\pgfqpoint{4.900000in}{4.900000in}}%
\pgfusepath{clip}%
\pgfsetbuttcap%
\pgfsetroundjoin%
\definecolor{currentfill}{rgb}{1.000000,1.000000,1.000000}%
\pgfsetfillcolor{currentfill}%
\pgfsetfillopacity{0.900000}%
\pgfsetlinewidth{0.501875pt}%
\definecolor{currentstroke}{rgb}{0.200000,0.200000,0.200000}%
\pgfsetstrokecolor{currentstroke}%
\pgfsetdash{}{0pt}%
\pgfpathmoveto{\pgfqpoint{2.965983in}{2.003233in}}%
\pgfpathlineto{\pgfqpoint{3.009197in}{2.013874in}}%
\pgfpathlineto{\pgfqpoint{3.040476in}{1.995179in}}%
\pgfpathlineto{\pgfqpoint{2.997233in}{1.984752in}}%
\pgfpathlineto{\pgfqpoint{2.965983in}{2.003233in}}%
\pgfpathclose%
\pgfusepath{stroke,fill}%
\end{pgfscope}%
\begin{pgfscope}%
\pgfpathrectangle{\pgfqpoint{0.050000in}{0.050000in}}{\pgfqpoint{4.900000in}{4.900000in}}%
\pgfusepath{clip}%
\pgfsetbuttcap%
\pgfsetroundjoin%
\definecolor{currentfill}{rgb}{1.000000,1.000000,1.000000}%
\pgfsetfillcolor{currentfill}%
\pgfsetfillopacity{0.900000}%
\pgfsetlinewidth{0.501875pt}%
\definecolor{currentstroke}{rgb}{0.200000,0.200000,0.200000}%
\pgfsetstrokecolor{currentstroke}%
\pgfsetdash{}{0pt}%
\pgfpathmoveto{\pgfqpoint{3.424909in}{1.994628in}}%
\pgfpathlineto{\pgfqpoint{3.467664in}{2.007930in}}%
\pgfpathlineto{\pgfqpoint{3.499627in}{1.988621in}}%
\pgfpathlineto{\pgfqpoint{3.456842in}{1.975322in}}%
\pgfpathlineto{\pgfqpoint{3.424909in}{1.994628in}}%
\pgfpathclose%
\pgfusepath{stroke,fill}%
\end{pgfscope}%
\begin{pgfscope}%
\pgfpathrectangle{\pgfqpoint{0.050000in}{0.050000in}}{\pgfqpoint{4.900000in}{4.900000in}}%
\pgfusepath{clip}%
\pgfsetbuttcap%
\pgfsetroundjoin%
\definecolor{currentfill}{rgb}{0.963091,0.963091,0.963091}%
\pgfsetfillcolor{currentfill}%
\pgfsetfillopacity{0.900000}%
\pgfsetlinewidth{0.501875pt}%
\definecolor{currentstroke}{rgb}{0.200000,0.200000,0.200000}%
\pgfsetstrokecolor{currentstroke}%
\pgfsetdash{}{0pt}%
\pgfpathmoveto{\pgfqpoint{2.327776in}{2.203314in}}%
\pgfpathlineto{\pgfqpoint{2.371852in}{2.172452in}}%
\pgfpathlineto{\pgfqpoint{2.402157in}{2.155623in}}%
\pgfpathlineto{\pgfqpoint{2.358072in}{2.185709in}}%
\pgfpathlineto{\pgfqpoint{2.327776in}{2.203314in}}%
\pgfpathclose%
\pgfusepath{stroke,fill}%
\end{pgfscope}%
\begin{pgfscope}%
\pgfpathrectangle{\pgfqpoint{0.050000in}{0.050000in}}{\pgfqpoint{4.900000in}{4.900000in}}%
\pgfusepath{clip}%
\pgfsetbuttcap%
\pgfsetroundjoin%
\definecolor{currentfill}{rgb}{0.992618,0.992618,0.992618}%
\pgfsetfillcolor{currentfill}%
\pgfsetfillopacity{0.900000}%
\pgfsetlinewidth{0.501875pt}%
\definecolor{currentstroke}{rgb}{0.200000,0.200000,0.200000}%
\pgfsetstrokecolor{currentstroke}%
\pgfsetdash{}{0pt}%
\pgfpathmoveto{\pgfqpoint{2.699408in}{2.035291in}}%
\pgfpathlineto{\pgfqpoint{2.742893in}{2.036081in}}%
\pgfpathlineto{\pgfqpoint{2.773777in}{2.018692in}}%
\pgfpathlineto{\pgfqpoint{2.730266in}{2.018215in}}%
\pgfpathlineto{\pgfqpoint{2.699408in}{2.035291in}}%
\pgfpathclose%
\pgfusepath{stroke,fill}%
\end{pgfscope}%
\begin{pgfscope}%
\pgfpathrectangle{\pgfqpoint{0.050000in}{0.050000in}}{\pgfqpoint{4.900000in}{4.900000in}}%
\pgfusepath{clip}%
\pgfsetbuttcap%
\pgfsetroundjoin%
\definecolor{currentfill}{rgb}{1.000000,1.000000,1.000000}%
\pgfsetfillcolor{currentfill}%
\pgfsetfillopacity{0.900000}%
\pgfsetlinewidth{0.501875pt}%
\definecolor{currentstroke}{rgb}{0.200000,0.200000,0.200000}%
\pgfsetstrokecolor{currentstroke}%
\pgfsetdash{}{0pt}%
\pgfpathmoveto{\pgfqpoint{3.232715in}{1.993670in}}%
\pgfpathlineto{\pgfqpoint{3.275674in}{2.006646in}}%
\pgfpathlineto{\pgfqpoint{3.307357in}{1.987467in}}%
\pgfpathlineto{\pgfqpoint{3.264367in}{1.974534in}}%
\pgfpathlineto{\pgfqpoint{3.232715in}{1.993670in}}%
\pgfpathclose%
\pgfusepath{stroke,fill}%
\end{pgfscope}%
\begin{pgfscope}%
\pgfpathrectangle{\pgfqpoint{0.050000in}{0.050000in}}{\pgfqpoint{4.900000in}{4.900000in}}%
\pgfusepath{clip}%
\pgfsetbuttcap%
\pgfsetroundjoin%
\definecolor{currentfill}{rgb}{0.996309,0.996309,0.996309}%
\pgfsetfillcolor{currentfill}%
\pgfsetfillopacity{0.900000}%
\pgfsetlinewidth{0.501875pt}%
\definecolor{currentstroke}{rgb}{0.200000,0.200000,0.200000}%
\pgfsetstrokecolor{currentstroke}%
\pgfsetdash{}{0pt}%
\pgfpathmoveto{\pgfqpoint{2.773777in}{2.018692in}}%
\pgfpathlineto{\pgfqpoint{2.817194in}{2.023093in}}%
\pgfpathlineto{\pgfqpoint{2.848199in}{2.005285in}}%
\pgfpathlineto{\pgfqpoint{2.804755in}{2.001216in}}%
\pgfpathlineto{\pgfqpoint{2.773777in}{2.018692in}}%
\pgfpathclose%
\pgfusepath{stroke,fill}%
\end{pgfscope}%
\begin{pgfscope}%
\pgfpathrectangle{\pgfqpoint{0.050000in}{0.050000in}}{\pgfqpoint{4.900000in}{4.900000in}}%
\pgfusepath{clip}%
\pgfsetbuttcap%
\pgfsetroundjoin%
\definecolor{currentfill}{rgb}{1.000000,1.000000,1.000000}%
\pgfsetfillcolor{currentfill}%
\pgfsetfillopacity{0.900000}%
\pgfsetlinewidth{0.501875pt}%
\definecolor{currentstroke}{rgb}{0.200000,0.200000,0.200000}%
\pgfsetstrokecolor{currentstroke}%
\pgfsetdash{}{0pt}%
\pgfpathmoveto{\pgfqpoint{3.040476in}{1.995179in}}%
\pgfpathlineto{\pgfqpoint{3.083635in}{2.006743in}}%
\pgfpathlineto{\pgfqpoint{3.115038in}{1.987860in}}%
\pgfpathlineto{\pgfqpoint{3.071850in}{1.976454in}}%
\pgfpathlineto{\pgfqpoint{3.040476in}{1.995179in}}%
\pgfpathclose%
\pgfusepath{stroke,fill}%
\end{pgfscope}%
\begin{pgfscope}%
\pgfpathrectangle{\pgfqpoint{0.050000in}{0.050000in}}{\pgfqpoint{4.900000in}{4.900000in}}%
\pgfusepath{clip}%
\pgfsetbuttcap%
\pgfsetroundjoin%
\definecolor{currentfill}{rgb}{0.952018,0.952018,0.952018}%
\pgfsetfillcolor{currentfill}%
\pgfsetfillopacity{0.900000}%
\pgfsetlinewidth{0.501875pt}%
\definecolor{currentstroke}{rgb}{0.200000,0.200000,0.200000}%
\pgfsetstrokecolor{currentstroke}%
\pgfsetdash{}{0pt}%
\pgfpathmoveto{\pgfqpoint{2.253316in}{2.257466in}}%
\pgfpathlineto{\pgfqpoint{2.297572in}{2.220896in}}%
\pgfpathlineto{\pgfqpoint{2.327776in}{2.203314in}}%
\pgfpathlineto{\pgfqpoint{2.283515in}{2.239071in}}%
\pgfpathlineto{\pgfqpoint{2.253316in}{2.257466in}}%
\pgfpathclose%
\pgfusepath{stroke,fill}%
\end{pgfscope}%
\begin{pgfscope}%
\pgfpathrectangle{\pgfqpoint{0.050000in}{0.050000in}}{\pgfqpoint{4.900000in}{4.900000in}}%
\pgfusepath{clip}%
\pgfsetbuttcap%
\pgfsetroundjoin%
\definecolor{currentfill}{rgb}{0.895548,0.895548,0.895548}%
\pgfsetfillcolor{currentfill}%
\pgfsetfillopacity{0.900000}%
\pgfsetlinewidth{0.501875pt}%
\definecolor{currentstroke}{rgb}{0.200000,0.200000,0.200000}%
\pgfsetstrokecolor{currentstroke}%
\pgfsetdash{}{0pt}%
\pgfpathmoveto{\pgfqpoint{1.999554in}{2.458794in}}%
\pgfpathlineto{\pgfqpoint{2.044355in}{2.416177in}}%
\pgfpathlineto{\pgfqpoint{2.074192in}{2.396768in}}%
\pgfpathlineto{\pgfqpoint{2.029385in}{2.439373in}}%
\pgfpathlineto{\pgfqpoint{1.999554in}{2.458794in}}%
\pgfpathclose%
\pgfusepath{stroke,fill}%
\end{pgfscope}%
\begin{pgfscope}%
\pgfpathrectangle{\pgfqpoint{0.050000in}{0.050000in}}{\pgfqpoint{4.900000in}{4.900000in}}%
\pgfusepath{clip}%
\pgfsetbuttcap%
\pgfsetroundjoin%
\definecolor{currentfill}{rgb}{1.000000,1.000000,1.000000}%
\pgfsetfillcolor{currentfill}%
\pgfsetfillopacity{0.900000}%
\pgfsetlinewidth{0.501875pt}%
\definecolor{currentstroke}{rgb}{0.200000,0.200000,0.200000}%
\pgfsetstrokecolor{currentstroke}%
\pgfsetdash{}{0pt}%
\pgfpathmoveto{\pgfqpoint{3.307357in}{1.987467in}}%
\pgfpathlineto{\pgfqpoint{3.350258in}{2.000618in}}%
\pgfpathlineto{\pgfqpoint{3.382066in}{1.981372in}}%
\pgfpathlineto{\pgfqpoint{3.339134in}{1.968243in}}%
\pgfpathlineto{\pgfqpoint{3.307357in}{1.987467in}}%
\pgfpathclose%
\pgfusepath{stroke,fill}%
\end{pgfscope}%
\begin{pgfscope}%
\pgfpathrectangle{\pgfqpoint{0.050000in}{0.050000in}}{\pgfqpoint{4.900000in}{4.900000in}}%
\pgfusepath{clip}%
\pgfsetbuttcap%
\pgfsetroundjoin%
\definecolor{currentfill}{rgb}{0.998155,0.998155,0.998155}%
\pgfsetfillcolor{currentfill}%
\pgfsetfillopacity{0.900000}%
\pgfsetlinewidth{0.501875pt}%
\definecolor{currentstroke}{rgb}{0.200000,0.200000,0.200000}%
\pgfsetstrokecolor{currentstroke}%
\pgfsetdash{}{0pt}%
\pgfpathmoveto{\pgfqpoint{2.848199in}{2.005285in}}%
\pgfpathlineto{\pgfqpoint{2.891556in}{2.012372in}}%
\pgfpathlineto{\pgfqpoint{2.922684in}{1.994196in}}%
\pgfpathlineto{\pgfqpoint{2.879299in}{1.987417in}}%
\pgfpathlineto{\pgfqpoint{2.848199in}{2.005285in}}%
\pgfpathclose%
\pgfusepath{stroke,fill}%
\end{pgfscope}%
\begin{pgfscope}%
\pgfpathrectangle{\pgfqpoint{0.050000in}{0.050000in}}{\pgfqpoint{4.900000in}{4.900000in}}%
\pgfusepath{clip}%
\pgfsetbuttcap%
\pgfsetroundjoin%
\definecolor{currentfill}{rgb}{1.000000,1.000000,1.000000}%
\pgfsetfillcolor{currentfill}%
\pgfsetfillopacity{0.900000}%
\pgfsetlinewidth{0.501875pt}%
\definecolor{currentstroke}{rgb}{0.200000,0.200000,0.200000}%
\pgfsetstrokecolor{currentstroke}%
\pgfsetdash{}{0pt}%
\pgfpathmoveto{\pgfqpoint{3.115038in}{1.987860in}}%
\pgfpathlineto{\pgfqpoint{3.158141in}{2.000061in}}%
\pgfpathlineto{\pgfqpoint{3.189669in}{1.981034in}}%
\pgfpathlineto{\pgfqpoint{3.146536in}{1.968944in}}%
\pgfpathlineto{\pgfqpoint{3.115038in}{1.987860in}}%
\pgfpathclose%
\pgfusepath{stroke,fill}%
\end{pgfscope}%
\begin{pgfscope}%
\pgfpathrectangle{\pgfqpoint{0.050000in}{0.050000in}}{\pgfqpoint{4.900000in}{4.900000in}}%
\pgfusepath{clip}%
\pgfsetbuttcap%
\pgfsetroundjoin%
\definecolor{currentfill}{rgb}{0.937993,0.937993,0.937993}%
\pgfsetfillcolor{currentfill}%
\pgfsetfillopacity{0.900000}%
\pgfsetlinewidth{0.501875pt}%
\definecolor{currentstroke}{rgb}{0.200000,0.200000,0.200000}%
\pgfsetstrokecolor{currentstroke}%
\pgfsetdash{}{0pt}%
\pgfpathmoveto{\pgfqpoint{2.178763in}{2.316197in}}%
\pgfpathlineto{\pgfqpoint{2.223206in}{2.275920in}}%
\pgfpathlineto{\pgfqpoint{2.253316in}{2.257466in}}%
\pgfpathlineto{\pgfqpoint{2.208869in}{2.297141in}}%
\pgfpathlineto{\pgfqpoint{2.178763in}{2.316197in}}%
\pgfpathclose%
\pgfusepath{stroke,fill}%
\end{pgfscope}%
\begin{pgfscope}%
\pgfpathrectangle{\pgfqpoint{0.050000in}{0.050000in}}{\pgfqpoint{4.900000in}{4.900000in}}%
\pgfusepath{clip}%
\pgfsetbuttcap%
\pgfsetroundjoin%
\definecolor{currentfill}{rgb}{0.875740,0.875740,0.875740}%
\pgfsetfillcolor{currentfill}%
\pgfsetfillopacity{0.900000}%
\pgfsetlinewidth{0.501875pt}%
\definecolor{currentstroke}{rgb}{0.200000,0.200000,0.200000}%
\pgfsetstrokecolor{currentstroke}%
\pgfsetdash{}{0pt}%
\pgfpathmoveto{\pgfqpoint{1.924833in}{2.520942in}}%
\pgfpathlineto{\pgfqpoint{1.969811in}{2.478164in}}%
\pgfpathlineto{\pgfqpoint{1.999554in}{2.458794in}}%
\pgfpathlineto{\pgfqpoint{1.954570in}{2.501596in}}%
\pgfpathlineto{\pgfqpoint{1.924833in}{2.520942in}}%
\pgfpathclose%
\pgfusepath{stroke,fill}%
\end{pgfscope}%
\begin{pgfscope}%
\pgfpathrectangle{\pgfqpoint{0.050000in}{0.050000in}}{\pgfqpoint{4.900000in}{4.900000in}}%
\pgfusepath{clip}%
\pgfsetbuttcap%
\pgfsetroundjoin%
\definecolor{currentfill}{rgb}{1.000000,1.000000,1.000000}%
\pgfsetfillcolor{currentfill}%
\pgfsetfillopacity{0.900000}%
\pgfsetlinewidth{0.501875pt}%
\definecolor{currentstroke}{rgb}{0.200000,0.200000,0.200000}%
\pgfsetstrokecolor{currentstroke}%
\pgfsetdash{}{0pt}%
\pgfpathmoveto{\pgfqpoint{3.382066in}{1.981372in}}%
\pgfpathlineto{\pgfqpoint{3.424909in}{1.994628in}}%
\pgfpathlineto{\pgfqpoint{3.456842in}{1.975322in}}%
\pgfpathlineto{\pgfqpoint{3.413969in}{1.962077in}}%
\pgfpathlineto{\pgfqpoint{3.382066in}{1.981372in}}%
\pgfpathclose%
\pgfusepath{stroke,fill}%
\end{pgfscope}%
\begin{pgfscope}%
\pgfpathrectangle{\pgfqpoint{0.050000in}{0.050000in}}{\pgfqpoint{4.900000in}{4.900000in}}%
\pgfusepath{clip}%
\pgfsetbuttcap%
\pgfsetroundjoin%
\definecolor{currentfill}{rgb}{0.998155,0.998155,0.998155}%
\pgfsetfillcolor{currentfill}%
\pgfsetfillopacity{0.900000}%
\pgfsetlinewidth{0.501875pt}%
\definecolor{currentstroke}{rgb}{0.200000,0.200000,0.200000}%
\pgfsetstrokecolor{currentstroke}%
\pgfsetdash{}{0pt}%
\pgfpathmoveto{\pgfqpoint{2.922684in}{1.994196in}}%
\pgfpathlineto{\pgfqpoint{2.965983in}{2.003233in}}%
\pgfpathlineto{\pgfqpoint{2.997233in}{1.984752in}}%
\pgfpathlineto{\pgfqpoint{2.953906in}{1.975976in}}%
\pgfpathlineto{\pgfqpoint{2.922684in}{1.994196in}}%
\pgfpathclose%
\pgfusepath{stroke,fill}%
\end{pgfscope}%
\begin{pgfscope}%
\pgfpathrectangle{\pgfqpoint{0.050000in}{0.050000in}}{\pgfqpoint{4.900000in}{4.900000in}}%
\pgfusepath{clip}%
\pgfsetbuttcap%
\pgfsetroundjoin%
\definecolor{currentfill}{rgb}{1.000000,1.000000,1.000000}%
\pgfsetfillcolor{currentfill}%
\pgfsetfillopacity{0.900000}%
\pgfsetlinewidth{0.501875pt}%
\definecolor{currentstroke}{rgb}{0.200000,0.200000,0.200000}%
\pgfsetstrokecolor{currentstroke}%
\pgfsetdash{}{0pt}%
\pgfpathmoveto{\pgfqpoint{3.189669in}{1.981034in}}%
\pgfpathlineto{\pgfqpoint{3.232715in}{1.993670in}}%
\pgfpathlineto{\pgfqpoint{3.264367in}{1.974534in}}%
\pgfpathlineto{\pgfqpoint{3.221291in}{1.961971in}}%
\pgfpathlineto{\pgfqpoint{3.189669in}{1.981034in}}%
\pgfpathclose%
\pgfusepath{stroke,fill}%
\end{pgfscope}%
\begin{pgfscope}%
\pgfpathrectangle{\pgfqpoint{0.050000in}{0.050000in}}{\pgfqpoint{4.900000in}{4.900000in}}%
\pgfusepath{clip}%
\pgfsetbuttcap%
\pgfsetroundjoin%
\definecolor{currentfill}{rgb}{0.979700,0.979700,0.979700}%
\pgfsetfillcolor{currentfill}%
\pgfsetfillopacity{0.900000}%
\pgfsetlinewidth{0.501875pt}%
\definecolor{currentstroke}{rgb}{0.200000,0.200000,0.200000}%
\pgfsetstrokecolor{currentstroke}%
\pgfsetdash{}{0pt}%
\pgfpathmoveto{\pgfqpoint{2.506990in}{2.098774in}}%
\pgfpathlineto{\pgfqpoint{2.550781in}{2.082460in}}%
\pgfpathlineto{\pgfqpoint{2.581403in}{2.065925in}}%
\pgfpathlineto{\pgfqpoint{2.537592in}{2.082029in}}%
\pgfpathlineto{\pgfqpoint{2.506990in}{2.098774in}}%
\pgfpathclose%
\pgfusepath{stroke,fill}%
\end{pgfscope}%
\begin{pgfscope}%
\pgfpathrectangle{\pgfqpoint{0.050000in}{0.050000in}}{\pgfqpoint{4.900000in}{4.900000in}}%
\pgfusepath{clip}%
\pgfsetbuttcap%
\pgfsetroundjoin%
\definecolor{currentfill}{rgb}{0.985236,0.985236,0.985236}%
\pgfsetfillcolor{currentfill}%
\pgfsetfillopacity{0.900000}%
\pgfsetlinewidth{0.501875pt}%
\definecolor{currentstroke}{rgb}{0.200000,0.200000,0.200000}%
\pgfsetstrokecolor{currentstroke}%
\pgfsetdash{}{0pt}%
\pgfpathmoveto{\pgfqpoint{2.581403in}{2.065925in}}%
\pgfpathlineto{\pgfqpoint{2.625082in}{2.056151in}}%
\pgfpathlineto{\pgfqpoint{2.655820in}{2.039418in}}%
\pgfpathlineto{\pgfqpoint{2.612118in}{2.049222in}}%
\pgfpathlineto{\pgfqpoint{2.581403in}{2.065925in}}%
\pgfpathclose%
\pgfusepath{stroke,fill}%
\end{pgfscope}%
\begin{pgfscope}%
\pgfpathrectangle{\pgfqpoint{0.050000in}{0.050000in}}{\pgfqpoint{4.900000in}{4.900000in}}%
\pgfusepath{clip}%
\pgfsetbuttcap%
\pgfsetroundjoin%
\definecolor{currentfill}{rgb}{0.972318,0.972318,0.972318}%
\pgfsetfillcolor{currentfill}%
\pgfsetfillopacity{0.900000}%
\pgfsetlinewidth{0.501875pt}%
\definecolor{currentstroke}{rgb}{0.200000,0.200000,0.200000}%
\pgfsetstrokecolor{currentstroke}%
\pgfsetdash{}{0pt}%
\pgfpathmoveto{\pgfqpoint{2.432556in}{2.138676in}}%
\pgfpathlineto{\pgfqpoint{2.476483in}{2.115356in}}%
\pgfpathlineto{\pgfqpoint{2.506990in}{2.098774in}}%
\pgfpathlineto{\pgfqpoint{2.463048in}{2.121612in}}%
\pgfpathlineto{\pgfqpoint{2.432556in}{2.138676in}}%
\pgfpathclose%
\pgfusepath{stroke,fill}%
\end{pgfscope}%
\begin{pgfscope}%
\pgfpathrectangle{\pgfqpoint{0.050000in}{0.050000in}}{\pgfqpoint{4.900000in}{4.900000in}}%
\pgfusepath{clip}%
\pgfsetbuttcap%
\pgfsetroundjoin%
\definecolor{currentfill}{rgb}{0.918185,0.918185,0.918185}%
\pgfsetfillcolor{currentfill}%
\pgfsetfillopacity{0.900000}%
\pgfsetlinewidth{0.501875pt}%
\definecolor{currentstroke}{rgb}{0.200000,0.200000,0.200000}%
\pgfsetstrokecolor{currentstroke}%
\pgfsetdash{}{0pt}%
\pgfpathmoveto{\pgfqpoint{2.104118in}{2.377355in}}%
\pgfpathlineto{\pgfqpoint{2.148745in}{2.335323in}}%
\pgfpathlineto{\pgfqpoint{2.178763in}{2.316197in}}%
\pgfpathlineto{\pgfqpoint{2.134131in}{2.357952in}}%
\pgfpathlineto{\pgfqpoint{2.104118in}{2.377355in}}%
\pgfpathclose%
\pgfusepath{stroke,fill}%
\end{pgfscope}%
\begin{pgfscope}%
\pgfpathrectangle{\pgfqpoint{0.050000in}{0.050000in}}{\pgfqpoint{4.900000in}{4.900000in}}%
\pgfusepath{clip}%
\pgfsetbuttcap%
\pgfsetroundjoin%
\definecolor{currentfill}{rgb}{0.988927,0.988927,0.988927}%
\pgfsetfillcolor{currentfill}%
\pgfsetfillopacity{0.900000}%
\pgfsetlinewidth{0.501875pt}%
\definecolor{currentstroke}{rgb}{0.200000,0.200000,0.200000}%
\pgfsetstrokecolor{currentstroke}%
\pgfsetdash{}{0pt}%
\pgfpathmoveto{\pgfqpoint{2.655820in}{2.039418in}}%
\pgfpathlineto{\pgfqpoint{2.699408in}{2.035291in}}%
\pgfpathlineto{\pgfqpoint{2.730266in}{2.018215in}}%
\pgfpathlineto{\pgfqpoint{2.686654in}{2.022537in}}%
\pgfpathlineto{\pgfqpoint{2.655820in}{2.039418in}}%
\pgfpathclose%
\pgfusepath{stroke,fill}%
\end{pgfscope}%
\begin{pgfscope}%
\pgfpathrectangle{\pgfqpoint{0.050000in}{0.050000in}}{\pgfqpoint{4.900000in}{4.900000in}}%
\pgfusepath{clip}%
\pgfsetbuttcap%
\pgfsetroundjoin%
\definecolor{currentfill}{rgb}{1.000000,1.000000,1.000000}%
\pgfsetfillcolor{currentfill}%
\pgfsetfillopacity{0.900000}%
\pgfsetlinewidth{0.501875pt}%
\definecolor{currentstroke}{rgb}{0.200000,0.200000,0.200000}%
\pgfsetstrokecolor{currentstroke}%
\pgfsetdash{}{0pt}%
\pgfpathmoveto{\pgfqpoint{2.997233in}{1.984752in}}%
\pgfpathlineto{\pgfqpoint{3.040476in}{1.995179in}}%
\pgfpathlineto{\pgfqpoint{3.071850in}{1.976454in}}%
\pgfpathlineto{\pgfqpoint{3.028579in}{1.966235in}}%
\pgfpathlineto{\pgfqpoint{2.997233in}{1.984752in}}%
\pgfpathclose%
\pgfusepath{stroke,fill}%
\end{pgfscope}%
\begin{pgfscope}%
\pgfpathrectangle{\pgfqpoint{0.050000in}{0.050000in}}{\pgfqpoint{4.900000in}{4.900000in}}%
\pgfusepath{clip}%
\pgfsetbuttcap%
\pgfsetroundjoin%
\definecolor{currentfill}{rgb}{1.000000,1.000000,1.000000}%
\pgfsetfillcolor{currentfill}%
\pgfsetfillopacity{0.900000}%
\pgfsetlinewidth{0.501875pt}%
\definecolor{currentstroke}{rgb}{0.200000,0.200000,0.200000}%
\pgfsetstrokecolor{currentstroke}%
\pgfsetdash{}{0pt}%
\pgfpathmoveto{\pgfqpoint{3.456842in}{1.975322in}}%
\pgfpathlineto{\pgfqpoint{3.499627in}{1.988621in}}%
\pgfpathlineto{\pgfqpoint{3.531685in}{1.969253in}}%
\pgfpathlineto{\pgfqpoint{3.488871in}{1.955962in}}%
\pgfpathlineto{\pgfqpoint{3.456842in}{1.975322in}}%
\pgfpathclose%
\pgfusepath{stroke,fill}%
\end{pgfscope}%
\begin{pgfscope}%
\pgfpathrectangle{\pgfqpoint{0.050000in}{0.050000in}}{\pgfqpoint{4.900000in}{4.900000in}}%
\pgfusepath{clip}%
\pgfsetbuttcap%
\pgfsetroundjoin%
\definecolor{currentfill}{rgb}{0.963091,0.963091,0.963091}%
\pgfsetfillcolor{currentfill}%
\pgfsetfillopacity{0.900000}%
\pgfsetlinewidth{0.501875pt}%
\definecolor{currentstroke}{rgb}{0.200000,0.200000,0.200000}%
\pgfsetstrokecolor{currentstroke}%
\pgfsetdash{}{0pt}%
\pgfpathmoveto{\pgfqpoint{2.358072in}{2.185709in}}%
\pgfpathlineto{\pgfqpoint{2.402157in}{2.155623in}}%
\pgfpathlineto{\pgfqpoint{2.432556in}{2.138676in}}%
\pgfpathlineto{\pgfqpoint{2.388461in}{2.168064in}}%
\pgfpathlineto{\pgfqpoint{2.358072in}{2.185709in}}%
\pgfpathclose%
\pgfusepath{stroke,fill}%
\end{pgfscope}%
\begin{pgfscope}%
\pgfpathrectangle{\pgfqpoint{0.050000in}{0.050000in}}{\pgfqpoint{4.900000in}{4.900000in}}%
\pgfusepath{clip}%
\pgfsetbuttcap%
\pgfsetroundjoin%
\definecolor{currentfill}{rgb}{0.992618,0.992618,0.992618}%
\pgfsetfillcolor{currentfill}%
\pgfsetfillopacity{0.900000}%
\pgfsetlinewidth{0.501875pt}%
\definecolor{currentstroke}{rgb}{0.200000,0.200000,0.200000}%
\pgfsetstrokecolor{currentstroke}%
\pgfsetdash{}{0pt}%
\pgfpathmoveto{\pgfqpoint{2.730266in}{2.018215in}}%
\pgfpathlineto{\pgfqpoint{2.773777in}{2.018692in}}%
\pgfpathlineto{\pgfqpoint{2.804755in}{2.001216in}}%
\pgfpathlineto{\pgfqpoint{2.761218in}{2.001021in}}%
\pgfpathlineto{\pgfqpoint{2.730266in}{2.018215in}}%
\pgfpathclose%
\pgfusepath{stroke,fill}%
\end{pgfscope}%
\begin{pgfscope}%
\pgfpathrectangle{\pgfqpoint{0.050000in}{0.050000in}}{\pgfqpoint{4.900000in}{4.900000in}}%
\pgfusepath{clip}%
\pgfsetbuttcap%
\pgfsetroundjoin%
\definecolor{currentfill}{rgb}{1.000000,1.000000,1.000000}%
\pgfsetfillcolor{currentfill}%
\pgfsetfillopacity{0.900000}%
\pgfsetlinewidth{0.501875pt}%
\definecolor{currentstroke}{rgb}{0.200000,0.200000,0.200000}%
\pgfsetstrokecolor{currentstroke}%
\pgfsetdash{}{0pt}%
\pgfpathmoveto{\pgfqpoint{3.264367in}{1.974534in}}%
\pgfpathlineto{\pgfqpoint{3.307357in}{1.987467in}}%
\pgfpathlineto{\pgfqpoint{3.339134in}{1.968243in}}%
\pgfpathlineto{\pgfqpoint{3.296115in}{1.955357in}}%
\pgfpathlineto{\pgfqpoint{3.264367in}{1.974534in}}%
\pgfpathclose%
\pgfusepath{stroke,fill}%
\end{pgfscope}%
\begin{pgfscope}%
\pgfpathrectangle{\pgfqpoint{0.050000in}{0.050000in}}{\pgfqpoint{4.900000in}{4.900000in}}%
\pgfusepath{clip}%
\pgfsetbuttcap%
\pgfsetroundjoin%
\definecolor{currentfill}{rgb}{1.000000,1.000000,1.000000}%
\pgfsetfillcolor{currentfill}%
\pgfsetfillopacity{0.900000}%
\pgfsetlinewidth{0.501875pt}%
\definecolor{currentstroke}{rgb}{0.200000,0.200000,0.200000}%
\pgfsetstrokecolor{currentstroke}%
\pgfsetdash{}{0pt}%
\pgfpathmoveto{\pgfqpoint{3.071850in}{1.976454in}}%
\pgfpathlineto{\pgfqpoint{3.115038in}{1.987860in}}%
\pgfpathlineto{\pgfqpoint{3.146536in}{1.968944in}}%
\pgfpathlineto{\pgfqpoint{3.103320in}{1.957695in}}%
\pgfpathlineto{\pgfqpoint{3.071850in}{1.976454in}}%
\pgfpathclose%
\pgfusepath{stroke,fill}%
\end{pgfscope}%
\begin{pgfscope}%
\pgfpathrectangle{\pgfqpoint{0.050000in}{0.050000in}}{\pgfqpoint{4.900000in}{4.900000in}}%
\pgfusepath{clip}%
\pgfsetbuttcap%
\pgfsetroundjoin%
\definecolor{currentfill}{rgb}{0.895548,0.895548,0.895548}%
\pgfsetfillcolor{currentfill}%
\pgfsetfillopacity{0.900000}%
\pgfsetlinewidth{0.501875pt}%
\definecolor{currentstroke}{rgb}{0.200000,0.200000,0.200000}%
\pgfsetstrokecolor{currentstroke}%
\pgfsetdash{}{0pt}%
\pgfpathmoveto{\pgfqpoint{2.029385in}{2.439373in}}%
\pgfpathlineto{\pgfqpoint{2.074192in}{2.396768in}}%
\pgfpathlineto{\pgfqpoint{2.104118in}{2.377355in}}%
\pgfpathlineto{\pgfqpoint{2.059305in}{2.419909in}}%
\pgfpathlineto{\pgfqpoint{2.029385in}{2.439373in}}%
\pgfpathclose%
\pgfusepath{stroke,fill}%
\end{pgfscope}%
\begin{pgfscope}%
\pgfpathrectangle{\pgfqpoint{0.050000in}{0.050000in}}{\pgfqpoint{4.900000in}{4.900000in}}%
\pgfusepath{clip}%
\pgfsetbuttcap%
\pgfsetroundjoin%
\definecolor{currentfill}{rgb}{0.996309,0.996309,0.996309}%
\pgfsetfillcolor{currentfill}%
\pgfsetfillopacity{0.900000}%
\pgfsetlinewidth{0.501875pt}%
\definecolor{currentstroke}{rgb}{0.200000,0.200000,0.200000}%
\pgfsetstrokecolor{currentstroke}%
\pgfsetdash{}{0pt}%
\pgfpathmoveto{\pgfqpoint{2.804755in}{2.001216in}}%
\pgfpathlineto{\pgfqpoint{2.848199in}{2.005285in}}%
\pgfpathlineto{\pgfqpoint{2.879299in}{1.987417in}}%
\pgfpathlineto{\pgfqpoint{2.835828in}{1.983654in}}%
\pgfpathlineto{\pgfqpoint{2.804755in}{2.001216in}}%
\pgfpathclose%
\pgfusepath{stroke,fill}%
\end{pgfscope}%
\begin{pgfscope}%
\pgfpathrectangle{\pgfqpoint{0.050000in}{0.050000in}}{\pgfqpoint{4.900000in}{4.900000in}}%
\pgfusepath{clip}%
\pgfsetbuttcap%
\pgfsetroundjoin%
\definecolor{currentfill}{rgb}{0.952018,0.952018,0.952018}%
\pgfsetfillcolor{currentfill}%
\pgfsetfillopacity{0.900000}%
\pgfsetlinewidth{0.501875pt}%
\definecolor{currentstroke}{rgb}{0.200000,0.200000,0.200000}%
\pgfsetstrokecolor{currentstroke}%
\pgfsetdash{}{0pt}%
\pgfpathmoveto{\pgfqpoint{2.283515in}{2.239071in}}%
\pgfpathlineto{\pgfqpoint{2.327776in}{2.203314in}}%
\pgfpathlineto{\pgfqpoint{2.358072in}{2.185709in}}%
\pgfpathlineto{\pgfqpoint{2.313805in}{2.220712in}}%
\pgfpathlineto{\pgfqpoint{2.283515in}{2.239071in}}%
\pgfpathclose%
\pgfusepath{stroke,fill}%
\end{pgfscope}%
\begin{pgfscope}%
\pgfpathrectangle{\pgfqpoint{0.050000in}{0.050000in}}{\pgfqpoint{4.900000in}{4.900000in}}%
\pgfusepath{clip}%
\pgfsetbuttcap%
\pgfsetroundjoin%
\definecolor{currentfill}{rgb}{1.000000,1.000000,1.000000}%
\pgfsetfillcolor{currentfill}%
\pgfsetfillopacity{0.900000}%
\pgfsetlinewidth{0.501875pt}%
\definecolor{currentstroke}{rgb}{0.200000,0.200000,0.200000}%
\pgfsetstrokecolor{currentstroke}%
\pgfsetdash{}{0pt}%
\pgfpathmoveto{\pgfqpoint{3.339134in}{1.968243in}}%
\pgfpathlineto{\pgfqpoint{3.382066in}{1.981372in}}%
\pgfpathlineto{\pgfqpoint{3.413969in}{1.962077in}}%
\pgfpathlineto{\pgfqpoint{3.371008in}{1.948976in}}%
\pgfpathlineto{\pgfqpoint{3.339134in}{1.968243in}}%
\pgfpathclose%
\pgfusepath{stroke,fill}%
\end{pgfscope}%
\begin{pgfscope}%
\pgfpathrectangle{\pgfqpoint{0.050000in}{0.050000in}}{\pgfqpoint{4.900000in}{4.900000in}}%
\pgfusepath{clip}%
\pgfsetbuttcap%
\pgfsetroundjoin%
\definecolor{currentfill}{rgb}{0.996309,0.996309,0.996309}%
\pgfsetfillcolor{currentfill}%
\pgfsetfillopacity{0.900000}%
\pgfsetlinewidth{0.501875pt}%
\definecolor{currentstroke}{rgb}{0.200000,0.200000,0.200000}%
\pgfsetstrokecolor{currentstroke}%
\pgfsetdash{}{0pt}%
\pgfpathmoveto{\pgfqpoint{2.879299in}{1.987417in}}%
\pgfpathlineto{\pgfqpoint{2.922684in}{1.994196in}}%
\pgfpathlineto{\pgfqpoint{2.953906in}{1.975976in}}%
\pgfpathlineto{\pgfqpoint{2.910494in}{1.969485in}}%
\pgfpathlineto{\pgfqpoint{2.879299in}{1.987417in}}%
\pgfpathclose%
\pgfusepath{stroke,fill}%
\end{pgfscope}%
\begin{pgfscope}%
\pgfpathrectangle{\pgfqpoint{0.050000in}{0.050000in}}{\pgfqpoint{4.900000in}{4.900000in}}%
\pgfusepath{clip}%
\pgfsetbuttcap%
\pgfsetroundjoin%
\definecolor{currentfill}{rgb}{1.000000,1.000000,1.000000}%
\pgfsetfillcolor{currentfill}%
\pgfsetfillopacity{0.900000}%
\pgfsetlinewidth{0.501875pt}%
\definecolor{currentstroke}{rgb}{0.200000,0.200000,0.200000}%
\pgfsetstrokecolor{currentstroke}%
\pgfsetdash{}{0pt}%
\pgfpathmoveto{\pgfqpoint{3.146536in}{1.968944in}}%
\pgfpathlineto{\pgfqpoint{3.189669in}{1.981034in}}%
\pgfpathlineto{\pgfqpoint{3.221291in}{1.961971in}}%
\pgfpathlineto{\pgfqpoint{3.178130in}{1.949994in}}%
\pgfpathlineto{\pgfqpoint{3.146536in}{1.968944in}}%
\pgfpathclose%
\pgfusepath{stroke,fill}%
\end{pgfscope}%
\begin{pgfscope}%
\pgfpathrectangle{\pgfqpoint{0.050000in}{0.050000in}}{\pgfqpoint{4.900000in}{4.900000in}}%
\pgfusepath{clip}%
\pgfsetbuttcap%
\pgfsetroundjoin%
\definecolor{currentfill}{rgb}{0.937993,0.937993,0.937993}%
\pgfsetfillcolor{currentfill}%
\pgfsetfillopacity{0.900000}%
\pgfsetlinewidth{0.501875pt}%
\definecolor{currentstroke}{rgb}{0.200000,0.200000,0.200000}%
\pgfsetstrokecolor{currentstroke}%
\pgfsetdash{}{0pt}%
\pgfpathmoveto{\pgfqpoint{2.208869in}{2.297141in}}%
\pgfpathlineto{\pgfqpoint{2.253316in}{2.257466in}}%
\pgfpathlineto{\pgfqpoint{2.283515in}{2.239071in}}%
\pgfpathlineto{\pgfqpoint{2.239065in}{2.278142in}}%
\pgfpathlineto{\pgfqpoint{2.208869in}{2.297141in}}%
\pgfpathclose%
\pgfusepath{stroke,fill}%
\end{pgfscope}%
\begin{pgfscope}%
\pgfpathrectangle{\pgfqpoint{0.050000in}{0.050000in}}{\pgfqpoint{4.900000in}{4.900000in}}%
\pgfusepath{clip}%
\pgfsetbuttcap%
\pgfsetroundjoin%
\definecolor{currentfill}{rgb}{0.875740,0.875740,0.875740}%
\pgfsetfillcolor{currentfill}%
\pgfsetfillopacity{0.900000}%
\pgfsetlinewidth{0.501875pt}%
\definecolor{currentstroke}{rgb}{0.200000,0.200000,0.200000}%
\pgfsetstrokecolor{currentstroke}%
\pgfsetdash{}{0pt}%
\pgfpathmoveto{\pgfqpoint{1.954570in}{2.501596in}}%
\pgfpathlineto{\pgfqpoint{1.999554in}{2.458794in}}%
\pgfpathlineto{\pgfqpoint{2.029385in}{2.439373in}}%
\pgfpathlineto{\pgfqpoint{1.984395in}{2.482192in}}%
\pgfpathlineto{\pgfqpoint{1.954570in}{2.501596in}}%
\pgfpathclose%
\pgfusepath{stroke,fill}%
\end{pgfscope}%
\begin{pgfscope}%
\pgfpathrectangle{\pgfqpoint{0.050000in}{0.050000in}}{\pgfqpoint{4.900000in}{4.900000in}}%
\pgfusepath{clip}%
\pgfsetbuttcap%
\pgfsetroundjoin%
\definecolor{currentfill}{rgb}{1.000000,1.000000,1.000000}%
\pgfsetfillcolor{currentfill}%
\pgfsetfillopacity{0.900000}%
\pgfsetlinewidth{0.501875pt}%
\definecolor{currentstroke}{rgb}{0.200000,0.200000,0.200000}%
\pgfsetstrokecolor{currentstroke}%
\pgfsetdash{}{0pt}%
\pgfpathmoveto{\pgfqpoint{3.413969in}{1.962077in}}%
\pgfpathlineto{\pgfqpoint{3.456842in}{1.975322in}}%
\pgfpathlineto{\pgfqpoint{3.488871in}{1.955962in}}%
\pgfpathlineto{\pgfqpoint{3.445968in}{1.942733in}}%
\pgfpathlineto{\pgfqpoint{3.413969in}{1.962077in}}%
\pgfpathclose%
\pgfusepath{stroke,fill}%
\end{pgfscope}%
\begin{pgfscope}%
\pgfpathrectangle{\pgfqpoint{0.050000in}{0.050000in}}{\pgfqpoint{4.900000in}{4.900000in}}%
\pgfusepath{clip}%
\pgfsetbuttcap%
\pgfsetroundjoin%
\definecolor{currentfill}{rgb}{0.998155,0.998155,0.998155}%
\pgfsetfillcolor{currentfill}%
\pgfsetfillopacity{0.900000}%
\pgfsetlinewidth{0.501875pt}%
\definecolor{currentstroke}{rgb}{0.200000,0.200000,0.200000}%
\pgfsetstrokecolor{currentstroke}%
\pgfsetdash{}{0pt}%
\pgfpathmoveto{\pgfqpoint{2.953906in}{1.975976in}}%
\pgfpathlineto{\pgfqpoint{2.997233in}{1.984752in}}%
\pgfpathlineto{\pgfqpoint{3.028579in}{1.966235in}}%
\pgfpathlineto{\pgfqpoint{2.985223in}{1.957708in}}%
\pgfpathlineto{\pgfqpoint{2.953906in}{1.975976in}}%
\pgfpathclose%
\pgfusepath{stroke,fill}%
\end{pgfscope}%
\begin{pgfscope}%
\pgfpathrectangle{\pgfqpoint{0.050000in}{0.050000in}}{\pgfqpoint{4.900000in}{4.900000in}}%
\pgfusepath{clip}%
\pgfsetbuttcap%
\pgfsetroundjoin%
\definecolor{currentfill}{rgb}{1.000000,1.000000,1.000000}%
\pgfsetfillcolor{currentfill}%
\pgfsetfillopacity{0.900000}%
\pgfsetlinewidth{0.501875pt}%
\definecolor{currentstroke}{rgb}{0.200000,0.200000,0.200000}%
\pgfsetstrokecolor{currentstroke}%
\pgfsetdash{}{0pt}%
\pgfpathmoveto{\pgfqpoint{3.221291in}{1.961971in}}%
\pgfpathlineto{\pgfqpoint{3.264367in}{1.974534in}}%
\pgfpathlineto{\pgfqpoint{3.296115in}{1.955357in}}%
\pgfpathlineto{\pgfqpoint{3.253010in}{1.942871in}}%
\pgfpathlineto{\pgfqpoint{3.221291in}{1.961971in}}%
\pgfpathclose%
\pgfusepath{stroke,fill}%
\end{pgfscope}%
\begin{pgfscope}%
\pgfpathrectangle{\pgfqpoint{0.050000in}{0.050000in}}{\pgfqpoint{4.900000in}{4.900000in}}%
\pgfusepath{clip}%
\pgfsetbuttcap%
\pgfsetroundjoin%
\definecolor{currentfill}{rgb}{0.979700,0.979700,0.979700}%
\pgfsetfillcolor{currentfill}%
\pgfsetfillopacity{0.900000}%
\pgfsetlinewidth{0.501875pt}%
\definecolor{currentstroke}{rgb}{0.200000,0.200000,0.200000}%
\pgfsetstrokecolor{currentstroke}%
\pgfsetdash{}{0pt}%
\pgfpathmoveto{\pgfqpoint{2.537592in}{2.082029in}}%
\pgfpathlineto{\pgfqpoint{2.581403in}{2.065925in}}%
\pgfpathlineto{\pgfqpoint{2.612118in}{2.049222in}}%
\pgfpathlineto{\pgfqpoint{2.568288in}{2.065134in}}%
\pgfpathlineto{\pgfqpoint{2.537592in}{2.082029in}}%
\pgfpathclose%
\pgfusepath{stroke,fill}%
\end{pgfscope}%
\begin{pgfscope}%
\pgfpathrectangle{\pgfqpoint{0.050000in}{0.050000in}}{\pgfqpoint{4.900000in}{4.900000in}}%
\pgfusepath{clip}%
\pgfsetbuttcap%
\pgfsetroundjoin%
\definecolor{currentfill}{rgb}{0.985236,0.985236,0.985236}%
\pgfsetfillcolor{currentfill}%
\pgfsetfillopacity{0.900000}%
\pgfsetlinewidth{0.501875pt}%
\definecolor{currentstroke}{rgb}{0.200000,0.200000,0.200000}%
\pgfsetstrokecolor{currentstroke}%
\pgfsetdash{}{0pt}%
\pgfpathmoveto{\pgfqpoint{2.612118in}{2.049222in}}%
\pgfpathlineto{\pgfqpoint{2.655820in}{2.039418in}}%
\pgfpathlineto{\pgfqpoint{2.686654in}{2.022537in}}%
\pgfpathlineto{\pgfqpoint{2.642929in}{2.032363in}}%
\pgfpathlineto{\pgfqpoint{2.612118in}{2.049222in}}%
\pgfpathclose%
\pgfusepath{stroke,fill}%
\end{pgfscope}%
\begin{pgfscope}%
\pgfpathrectangle{\pgfqpoint{0.050000in}{0.050000in}}{\pgfqpoint{4.900000in}{4.900000in}}%
\pgfusepath{clip}%
\pgfsetbuttcap%
\pgfsetroundjoin%
\definecolor{currentfill}{rgb}{0.918185,0.918185,0.918185}%
\pgfsetfillcolor{currentfill}%
\pgfsetfillopacity{0.900000}%
\pgfsetlinewidth{0.501875pt}%
\definecolor{currentstroke}{rgb}{0.200000,0.200000,0.200000}%
\pgfsetstrokecolor{currentstroke}%
\pgfsetdash{}{0pt}%
\pgfpathmoveto{\pgfqpoint{2.134131in}{2.357952in}}%
\pgfpathlineto{\pgfqpoint{2.178763in}{2.316197in}}%
\pgfpathlineto{\pgfqpoint{2.208869in}{2.297141in}}%
\pgfpathlineto{\pgfqpoint{2.164234in}{2.338568in}}%
\pgfpathlineto{\pgfqpoint{2.134131in}{2.357952in}}%
\pgfpathclose%
\pgfusepath{stroke,fill}%
\end{pgfscope}%
\begin{pgfscope}%
\pgfpathrectangle{\pgfqpoint{0.050000in}{0.050000in}}{\pgfqpoint{4.900000in}{4.900000in}}%
\pgfusepath{clip}%
\pgfsetbuttcap%
\pgfsetroundjoin%
\definecolor{currentfill}{rgb}{0.970473,0.970473,0.970473}%
\pgfsetfillcolor{currentfill}%
\pgfsetfillopacity{0.900000}%
\pgfsetlinewidth{0.501875pt}%
\definecolor{currentstroke}{rgb}{0.200000,0.200000,0.200000}%
\pgfsetstrokecolor{currentstroke}%
\pgfsetdash{}{0pt}%
\pgfpathmoveto{\pgfqpoint{2.463048in}{2.121612in}}%
\pgfpathlineto{\pgfqpoint{2.506990in}{2.098774in}}%
\pgfpathlineto{\pgfqpoint{2.537592in}{2.082029in}}%
\pgfpathlineto{\pgfqpoint{2.493635in}{2.104434in}}%
\pgfpathlineto{\pgfqpoint{2.463048in}{2.121612in}}%
\pgfpathclose%
\pgfusepath{stroke,fill}%
\end{pgfscope}%
\begin{pgfscope}%
\pgfpathrectangle{\pgfqpoint{0.050000in}{0.050000in}}{\pgfqpoint{4.900000in}{4.900000in}}%
\pgfusepath{clip}%
\pgfsetbuttcap%
\pgfsetroundjoin%
\definecolor{currentfill}{rgb}{0.988927,0.988927,0.988927}%
\pgfsetfillcolor{currentfill}%
\pgfsetfillopacity{0.900000}%
\pgfsetlinewidth{0.501875pt}%
\definecolor{currentstroke}{rgb}{0.200000,0.200000,0.200000}%
\pgfsetstrokecolor{currentstroke}%
\pgfsetdash{}{0pt}%
\pgfpathmoveto{\pgfqpoint{2.686654in}{2.022537in}}%
\pgfpathlineto{\pgfqpoint{2.730266in}{2.018215in}}%
\pgfpathlineto{\pgfqpoint{2.761218in}{2.001021in}}%
\pgfpathlineto{\pgfqpoint{2.717582in}{2.005516in}}%
\pgfpathlineto{\pgfqpoint{2.686654in}{2.022537in}}%
\pgfpathclose%
\pgfusepath{stroke,fill}%
\end{pgfscope}%
\begin{pgfscope}%
\pgfpathrectangle{\pgfqpoint{0.050000in}{0.050000in}}{\pgfqpoint{4.900000in}{4.900000in}}%
\pgfusepath{clip}%
\pgfsetbuttcap%
\pgfsetroundjoin%
\definecolor{currentfill}{rgb}{1.000000,1.000000,1.000000}%
\pgfsetfillcolor{currentfill}%
\pgfsetfillopacity{0.900000}%
\pgfsetlinewidth{0.501875pt}%
\definecolor{currentstroke}{rgb}{0.200000,0.200000,0.200000}%
\pgfsetstrokecolor{currentstroke}%
\pgfsetdash{}{0pt}%
\pgfpathmoveto{\pgfqpoint{3.028579in}{1.966235in}}%
\pgfpathlineto{\pgfqpoint{3.071850in}{1.976454in}}%
\pgfpathlineto{\pgfqpoint{3.103320in}{1.957695in}}%
\pgfpathlineto{\pgfqpoint{3.060019in}{1.947678in}}%
\pgfpathlineto{\pgfqpoint{3.028579in}{1.966235in}}%
\pgfpathclose%
\pgfusepath{stroke,fill}%
\end{pgfscope}%
\begin{pgfscope}%
\pgfpathrectangle{\pgfqpoint{0.050000in}{0.050000in}}{\pgfqpoint{4.900000in}{4.900000in}}%
\pgfusepath{clip}%
\pgfsetbuttcap%
\pgfsetroundjoin%
\definecolor{currentfill}{rgb}{0.853103,0.853103,0.853103}%
\pgfsetfillcolor{currentfill}%
\pgfsetfillopacity{0.900000}%
\pgfsetlinewidth{0.501875pt}%
\definecolor{currentstroke}{rgb}{0.200000,0.200000,0.200000}%
\pgfsetstrokecolor{currentstroke}%
\pgfsetdash{}{0pt}%
\pgfpathmoveto{\pgfqpoint{1.879672in}{2.563895in}}%
\pgfpathlineto{\pgfqpoint{1.924833in}{2.520942in}}%
\pgfpathlineto{\pgfqpoint{1.954570in}{2.501596in}}%
\pgfpathlineto{\pgfqpoint{1.909402in}{2.544573in}}%
\pgfpathlineto{\pgfqpoint{1.879672in}{2.563895in}}%
\pgfpathclose%
\pgfusepath{stroke,fill}%
\end{pgfscope}%
\begin{pgfscope}%
\pgfpathrectangle{\pgfqpoint{0.050000in}{0.050000in}}{\pgfqpoint{4.900000in}{4.900000in}}%
\pgfusepath{clip}%
\pgfsetbuttcap%
\pgfsetroundjoin%
\definecolor{currentfill}{rgb}{1.000000,1.000000,1.000000}%
\pgfsetfillcolor{currentfill}%
\pgfsetfillopacity{0.900000}%
\pgfsetlinewidth{0.501875pt}%
\definecolor{currentstroke}{rgb}{0.200000,0.200000,0.200000}%
\pgfsetstrokecolor{currentstroke}%
\pgfsetdash{}{0pt}%
\pgfpathmoveto{\pgfqpoint{3.488871in}{1.955962in}}%
\pgfpathlineto{\pgfqpoint{3.531685in}{1.969253in}}%
\pgfpathlineto{\pgfqpoint{3.563840in}{1.949828in}}%
\pgfpathlineto{\pgfqpoint{3.520997in}{1.936549in}}%
\pgfpathlineto{\pgfqpoint{3.488871in}{1.955962in}}%
\pgfpathclose%
\pgfusepath{stroke,fill}%
\end{pgfscope}%
\begin{pgfscope}%
\pgfpathrectangle{\pgfqpoint{0.050000in}{0.050000in}}{\pgfqpoint{4.900000in}{4.900000in}}%
\pgfusepath{clip}%
\pgfsetbuttcap%
\pgfsetroundjoin%
\definecolor{currentfill}{rgb}{0.961246,0.961246,0.961246}%
\pgfsetfillcolor{currentfill}%
\pgfsetfillopacity{0.900000}%
\pgfsetlinewidth{0.501875pt}%
\definecolor{currentstroke}{rgb}{0.200000,0.200000,0.200000}%
\pgfsetstrokecolor{currentstroke}%
\pgfsetdash{}{0pt}%
\pgfpathmoveto{\pgfqpoint{2.388461in}{2.168064in}}%
\pgfpathlineto{\pgfqpoint{2.432556in}{2.138676in}}%
\pgfpathlineto{\pgfqpoint{2.463048in}{2.121612in}}%
\pgfpathlineto{\pgfqpoint{2.418942in}{2.150369in}}%
\pgfpathlineto{\pgfqpoint{2.388461in}{2.168064in}}%
\pgfpathclose%
\pgfusepath{stroke,fill}%
\end{pgfscope}%
\begin{pgfscope}%
\pgfpathrectangle{\pgfqpoint{0.050000in}{0.050000in}}{\pgfqpoint{4.900000in}{4.900000in}}%
\pgfusepath{clip}%
\pgfsetbuttcap%
\pgfsetroundjoin%
\definecolor{currentfill}{rgb}{0.992618,0.992618,0.992618}%
\pgfsetfillcolor{currentfill}%
\pgfsetfillopacity{0.900000}%
\pgfsetlinewidth{0.501875pt}%
\definecolor{currentstroke}{rgb}{0.200000,0.200000,0.200000}%
\pgfsetstrokecolor{currentstroke}%
\pgfsetdash{}{0pt}%
\pgfpathmoveto{\pgfqpoint{2.761218in}{2.001021in}}%
\pgfpathlineto{\pgfqpoint{2.804755in}{2.001216in}}%
\pgfpathlineto{\pgfqpoint{2.835828in}{1.983654in}}%
\pgfpathlineto{\pgfqpoint{2.792266in}{1.983714in}}%
\pgfpathlineto{\pgfqpoint{2.761218in}{2.001021in}}%
\pgfpathclose%
\pgfusepath{stroke,fill}%
\end{pgfscope}%
\begin{pgfscope}%
\pgfpathrectangle{\pgfqpoint{0.050000in}{0.050000in}}{\pgfqpoint{4.900000in}{4.900000in}}%
\pgfusepath{clip}%
\pgfsetbuttcap%
\pgfsetroundjoin%
\definecolor{currentfill}{rgb}{1.000000,1.000000,1.000000}%
\pgfsetfillcolor{currentfill}%
\pgfsetfillopacity{0.900000}%
\pgfsetlinewidth{0.501875pt}%
\definecolor{currentstroke}{rgb}{0.200000,0.200000,0.200000}%
\pgfsetstrokecolor{currentstroke}%
\pgfsetdash{}{0pt}%
\pgfpathmoveto{\pgfqpoint{3.296115in}{1.955357in}}%
\pgfpathlineto{\pgfqpoint{3.339134in}{1.968243in}}%
\pgfpathlineto{\pgfqpoint{3.371008in}{1.948976in}}%
\pgfpathlineto{\pgfqpoint{3.327959in}{1.936139in}}%
\pgfpathlineto{\pgfqpoint{3.296115in}{1.955357in}}%
\pgfpathclose%
\pgfusepath{stroke,fill}%
\end{pgfscope}%
\begin{pgfscope}%
\pgfpathrectangle{\pgfqpoint{0.050000in}{0.050000in}}{\pgfqpoint{4.900000in}{4.900000in}}%
\pgfusepath{clip}%
\pgfsetbuttcap%
\pgfsetroundjoin%
\definecolor{currentfill}{rgb}{0.898378,0.898378,0.898378}%
\pgfsetfillcolor{currentfill}%
\pgfsetfillopacity{0.900000}%
\pgfsetlinewidth{0.501875pt}%
\definecolor{currentstroke}{rgb}{0.200000,0.200000,0.200000}%
\pgfsetstrokecolor{currentstroke}%
\pgfsetdash{}{0pt}%
\pgfpathmoveto{\pgfqpoint{2.059305in}{2.419909in}}%
\pgfpathlineto{\pgfqpoint{2.104118in}{2.377355in}}%
\pgfpathlineto{\pgfqpoint{2.134131in}{2.357952in}}%
\pgfpathlineto{\pgfqpoint{2.089315in}{2.400412in}}%
\pgfpathlineto{\pgfqpoint{2.059305in}{2.419909in}}%
\pgfpathclose%
\pgfusepath{stroke,fill}%
\end{pgfscope}%
\begin{pgfscope}%
\pgfpathrectangle{\pgfqpoint{0.050000in}{0.050000in}}{\pgfqpoint{4.900000in}{4.900000in}}%
\pgfusepath{clip}%
\pgfsetbuttcap%
\pgfsetroundjoin%
\definecolor{currentfill}{rgb}{1.000000,1.000000,1.000000}%
\pgfsetfillcolor{currentfill}%
\pgfsetfillopacity{0.900000}%
\pgfsetlinewidth{0.501875pt}%
\definecolor{currentstroke}{rgb}{0.200000,0.200000,0.200000}%
\pgfsetstrokecolor{currentstroke}%
\pgfsetdash{}{0pt}%
\pgfpathmoveto{\pgfqpoint{3.103320in}{1.957695in}}%
\pgfpathlineto{\pgfqpoint{3.146536in}{1.968944in}}%
\pgfpathlineto{\pgfqpoint{3.178130in}{1.949994in}}%
\pgfpathlineto{\pgfqpoint{3.134884in}{1.938900in}}%
\pgfpathlineto{\pgfqpoint{3.103320in}{1.957695in}}%
\pgfpathclose%
\pgfusepath{stroke,fill}%
\end{pgfscope}%
\begin{pgfscope}%
\pgfpathrectangle{\pgfqpoint{0.050000in}{0.050000in}}{\pgfqpoint{4.900000in}{4.900000in}}%
\pgfusepath{clip}%
\pgfsetbuttcap%
\pgfsetroundjoin%
\definecolor{currentfill}{rgb}{0.952018,0.952018,0.952018}%
\pgfsetfillcolor{currentfill}%
\pgfsetfillopacity{0.900000}%
\pgfsetlinewidth{0.501875pt}%
\definecolor{currentstroke}{rgb}{0.200000,0.200000,0.200000}%
\pgfsetstrokecolor{currentstroke}%
\pgfsetdash{}{0pt}%
\pgfpathmoveto{\pgfqpoint{2.313805in}{2.220712in}}%
\pgfpathlineto{\pgfqpoint{2.358072in}{2.185709in}}%
\pgfpathlineto{\pgfqpoint{2.388461in}{2.168064in}}%
\pgfpathlineto{\pgfqpoint{2.344187in}{2.202366in}}%
\pgfpathlineto{\pgfqpoint{2.313805in}{2.220712in}}%
\pgfpathclose%
\pgfusepath{stroke,fill}%
\end{pgfscope}%
\begin{pgfscope}%
\pgfpathrectangle{\pgfqpoint{0.050000in}{0.050000in}}{\pgfqpoint{4.900000in}{4.900000in}}%
\pgfusepath{clip}%
\pgfsetbuttcap%
\pgfsetroundjoin%
\definecolor{currentfill}{rgb}{0.994464,0.994464,0.994464}%
\pgfsetfillcolor{currentfill}%
\pgfsetfillopacity{0.900000}%
\pgfsetlinewidth{0.501875pt}%
\definecolor{currentstroke}{rgb}{0.200000,0.200000,0.200000}%
\pgfsetstrokecolor{currentstroke}%
\pgfsetdash{}{0pt}%
\pgfpathmoveto{\pgfqpoint{2.835828in}{1.983654in}}%
\pgfpathlineto{\pgfqpoint{2.879299in}{1.987417in}}%
\pgfpathlineto{\pgfqpoint{2.910494in}{1.969485in}}%
\pgfpathlineto{\pgfqpoint{2.866997in}{1.966004in}}%
\pgfpathlineto{\pgfqpoint{2.835828in}{1.983654in}}%
\pgfpathclose%
\pgfusepath{stroke,fill}%
\end{pgfscope}%
\begin{pgfscope}%
\pgfpathrectangle{\pgfqpoint{0.050000in}{0.050000in}}{\pgfqpoint{4.900000in}{4.900000in}}%
\pgfusepath{clip}%
\pgfsetbuttcap%
\pgfsetroundjoin%
\definecolor{currentfill}{rgb}{1.000000,1.000000,1.000000}%
\pgfsetfillcolor{currentfill}%
\pgfsetfillopacity{0.900000}%
\pgfsetlinewidth{0.501875pt}%
\definecolor{currentstroke}{rgb}{0.200000,0.200000,0.200000}%
\pgfsetstrokecolor{currentstroke}%
\pgfsetdash{}{0pt}%
\pgfpathmoveto{\pgfqpoint{3.371008in}{1.948976in}}%
\pgfpathlineto{\pgfqpoint{3.413969in}{1.962077in}}%
\pgfpathlineto{\pgfqpoint{3.445968in}{1.942733in}}%
\pgfpathlineto{\pgfqpoint{3.402977in}{1.929663in}}%
\pgfpathlineto{\pgfqpoint{3.371008in}{1.948976in}}%
\pgfpathclose%
\pgfusepath{stroke,fill}%
\end{pgfscope}%
\begin{pgfscope}%
\pgfpathrectangle{\pgfqpoint{0.050000in}{0.050000in}}{\pgfqpoint{4.900000in}{4.900000in}}%
\pgfusepath{clip}%
\pgfsetbuttcap%
\pgfsetroundjoin%
\definecolor{currentfill}{rgb}{0.996309,0.996309,0.996309}%
\pgfsetfillcolor{currentfill}%
\pgfsetfillopacity{0.900000}%
\pgfsetlinewidth{0.501875pt}%
\definecolor{currentstroke}{rgb}{0.200000,0.200000,0.200000}%
\pgfsetstrokecolor{currentstroke}%
\pgfsetdash{}{0pt}%
\pgfpathmoveto{\pgfqpoint{2.910494in}{1.969485in}}%
\pgfpathlineto{\pgfqpoint{2.953906in}{1.975976in}}%
\pgfpathlineto{\pgfqpoint{2.985223in}{1.957708in}}%
\pgfpathlineto{\pgfqpoint{2.941785in}{1.951487in}}%
\pgfpathlineto{\pgfqpoint{2.910494in}{1.969485in}}%
\pgfpathclose%
\pgfusepath{stroke,fill}%
\end{pgfscope}%
\begin{pgfscope}%
\pgfpathrectangle{\pgfqpoint{0.050000in}{0.050000in}}{\pgfqpoint{4.900000in}{4.900000in}}%
\pgfusepath{clip}%
\pgfsetbuttcap%
\pgfsetroundjoin%
\definecolor{currentfill}{rgb}{1.000000,1.000000,1.000000}%
\pgfsetfillcolor{currentfill}%
\pgfsetfillopacity{0.900000}%
\pgfsetlinewidth{0.501875pt}%
\definecolor{currentstroke}{rgb}{0.200000,0.200000,0.200000}%
\pgfsetstrokecolor{currentstroke}%
\pgfsetdash{}{0pt}%
\pgfpathmoveto{\pgfqpoint{3.178130in}{1.949994in}}%
\pgfpathlineto{\pgfqpoint{3.221291in}{1.961971in}}%
\pgfpathlineto{\pgfqpoint{3.253010in}{1.942871in}}%
\pgfpathlineto{\pgfqpoint{3.209819in}{1.931008in}}%
\pgfpathlineto{\pgfqpoint{3.178130in}{1.949994in}}%
\pgfpathclose%
\pgfusepath{stroke,fill}%
\end{pgfscope}%
\begin{pgfscope}%
\pgfpathrectangle{\pgfqpoint{0.050000in}{0.050000in}}{\pgfqpoint{4.900000in}{4.900000in}}%
\pgfusepath{clip}%
\pgfsetbuttcap%
\pgfsetroundjoin%
\definecolor{currentfill}{rgb}{0.937993,0.937993,0.937993}%
\pgfsetfillcolor{currentfill}%
\pgfsetfillopacity{0.900000}%
\pgfsetlinewidth{0.501875pt}%
\definecolor{currentstroke}{rgb}{0.200000,0.200000,0.200000}%
\pgfsetstrokecolor{currentstroke}%
\pgfsetdash{}{0pt}%
\pgfpathmoveto{\pgfqpoint{2.239065in}{2.278142in}}%
\pgfpathlineto{\pgfqpoint{2.283515in}{2.239071in}}%
\pgfpathlineto{\pgfqpoint{2.313805in}{2.220712in}}%
\pgfpathlineto{\pgfqpoint{2.269351in}{2.259188in}}%
\pgfpathlineto{\pgfqpoint{2.239065in}{2.278142in}}%
\pgfpathclose%
\pgfusepath{stroke,fill}%
\end{pgfscope}%
\begin{pgfscope}%
\pgfpathrectangle{\pgfqpoint{0.050000in}{0.050000in}}{\pgfqpoint{4.900000in}{4.900000in}}%
\pgfusepath{clip}%
\pgfsetbuttcap%
\pgfsetroundjoin%
\definecolor{currentfill}{rgb}{0.875740,0.875740,0.875740}%
\pgfsetfillcolor{currentfill}%
\pgfsetfillopacity{0.900000}%
\pgfsetlinewidth{0.501875pt}%
\definecolor{currentstroke}{rgb}{0.200000,0.200000,0.200000}%
\pgfsetstrokecolor{currentstroke}%
\pgfsetdash{}{0pt}%
\pgfpathmoveto{\pgfqpoint{1.984395in}{2.482192in}}%
\pgfpathlineto{\pgfqpoint{2.029385in}{2.439373in}}%
\pgfpathlineto{\pgfqpoint{2.059305in}{2.419909in}}%
\pgfpathlineto{\pgfqpoint{2.014310in}{2.462734in}}%
\pgfpathlineto{\pgfqpoint{1.984395in}{2.482192in}}%
\pgfpathclose%
\pgfusepath{stroke,fill}%
\end{pgfscope}%
\begin{pgfscope}%
\pgfpathrectangle{\pgfqpoint{0.050000in}{0.050000in}}{\pgfqpoint{4.900000in}{4.900000in}}%
\pgfusepath{clip}%
\pgfsetbuttcap%
\pgfsetroundjoin%
\definecolor{currentfill}{rgb}{1.000000,1.000000,1.000000}%
\pgfsetfillcolor{currentfill}%
\pgfsetfillopacity{0.900000}%
\pgfsetlinewidth{0.501875pt}%
\definecolor{currentstroke}{rgb}{0.200000,0.200000,0.200000}%
\pgfsetstrokecolor{currentstroke}%
\pgfsetdash{}{0pt}%
\pgfpathmoveto{\pgfqpoint{3.445968in}{1.942733in}}%
\pgfpathlineto{\pgfqpoint{3.488871in}{1.955962in}}%
\pgfpathlineto{\pgfqpoint{3.520997in}{1.936549in}}%
\pgfpathlineto{\pgfqpoint{3.478064in}{1.923340in}}%
\pgfpathlineto{\pgfqpoint{3.445968in}{1.942733in}}%
\pgfpathclose%
\pgfusepath{stroke,fill}%
\end{pgfscope}%
\begin{pgfscope}%
\pgfpathrectangle{\pgfqpoint{0.050000in}{0.050000in}}{\pgfqpoint{4.900000in}{4.900000in}}%
\pgfusepath{clip}%
\pgfsetbuttcap%
\pgfsetroundjoin%
\definecolor{currentfill}{rgb}{0.998155,0.998155,0.998155}%
\pgfsetfillcolor{currentfill}%
\pgfsetfillopacity{0.900000}%
\pgfsetlinewidth{0.501875pt}%
\definecolor{currentstroke}{rgb}{0.200000,0.200000,0.200000}%
\pgfsetstrokecolor{currentstroke}%
\pgfsetdash{}{0pt}%
\pgfpathmoveto{\pgfqpoint{2.985223in}{1.957708in}}%
\pgfpathlineto{\pgfqpoint{3.028579in}{1.966235in}}%
\pgfpathlineto{\pgfqpoint{3.060019in}{1.947678in}}%
\pgfpathlineto{\pgfqpoint{3.016637in}{1.939388in}}%
\pgfpathlineto{\pgfqpoint{2.985223in}{1.957708in}}%
\pgfpathclose%
\pgfusepath{stroke,fill}%
\end{pgfscope}%
\begin{pgfscope}%
\pgfpathrectangle{\pgfqpoint{0.050000in}{0.050000in}}{\pgfqpoint{4.900000in}{4.900000in}}%
\pgfusepath{clip}%
\pgfsetbuttcap%
\pgfsetroundjoin%
\definecolor{currentfill}{rgb}{1.000000,1.000000,1.000000}%
\pgfsetfillcolor{currentfill}%
\pgfsetfillopacity{0.900000}%
\pgfsetlinewidth{0.501875pt}%
\definecolor{currentstroke}{rgb}{0.200000,0.200000,0.200000}%
\pgfsetstrokecolor{currentstroke}%
\pgfsetdash{}{0pt}%
\pgfpathmoveto{\pgfqpoint{3.253010in}{1.942871in}}%
\pgfpathlineto{\pgfqpoint{3.296115in}{1.955357in}}%
\pgfpathlineto{\pgfqpoint{3.327959in}{1.936139in}}%
\pgfpathlineto{\pgfqpoint{3.284824in}{1.923732in}}%
\pgfpathlineto{\pgfqpoint{3.253010in}{1.942871in}}%
\pgfpathclose%
\pgfusepath{stroke,fill}%
\end{pgfscope}%
\begin{pgfscope}%
\pgfpathrectangle{\pgfqpoint{0.050000in}{0.050000in}}{\pgfqpoint{4.900000in}{4.900000in}}%
\pgfusepath{clip}%
\pgfsetbuttcap%
\pgfsetroundjoin%
\definecolor{currentfill}{rgb}{0.918185,0.918185,0.918185}%
\pgfsetfillcolor{currentfill}%
\pgfsetfillopacity{0.900000}%
\pgfsetlinewidth{0.501875pt}%
\definecolor{currentstroke}{rgb}{0.200000,0.200000,0.200000}%
\pgfsetstrokecolor{currentstroke}%
\pgfsetdash{}{0pt}%
\pgfpathmoveto{\pgfqpoint{2.164234in}{2.338568in}}%
\pgfpathlineto{\pgfqpoint{2.208869in}{2.297141in}}%
\pgfpathlineto{\pgfqpoint{2.239065in}{2.278142in}}%
\pgfpathlineto{\pgfqpoint{2.194427in}{2.319206in}}%
\pgfpathlineto{\pgfqpoint{2.164234in}{2.338568in}}%
\pgfpathclose%
\pgfusepath{stroke,fill}%
\end{pgfscope}%
\begin{pgfscope}%
\pgfpathrectangle{\pgfqpoint{0.050000in}{0.050000in}}{\pgfqpoint{4.900000in}{4.900000in}}%
\pgfusepath{clip}%
\pgfsetbuttcap%
\pgfsetroundjoin%
\definecolor{currentfill}{rgb}{0.977855,0.977855,0.977855}%
\pgfsetfillcolor{currentfill}%
\pgfsetfillopacity{0.900000}%
\pgfsetlinewidth{0.501875pt}%
\definecolor{currentstroke}{rgb}{0.200000,0.200000,0.200000}%
\pgfsetstrokecolor{currentstroke}%
\pgfsetdash{}{0pt}%
\pgfpathmoveto{\pgfqpoint{2.568288in}{2.065134in}}%
\pgfpathlineto{\pgfqpoint{2.612118in}{2.049222in}}%
\pgfpathlineto{\pgfqpoint{2.642929in}{2.032363in}}%
\pgfpathlineto{\pgfqpoint{2.599080in}{2.048097in}}%
\pgfpathlineto{\pgfqpoint{2.568288in}{2.065134in}}%
\pgfpathclose%
\pgfusepath{stroke,fill}%
\end{pgfscope}%
\begin{pgfscope}%
\pgfpathrectangle{\pgfqpoint{0.050000in}{0.050000in}}{\pgfqpoint{4.900000in}{4.900000in}}%
\pgfusepath{clip}%
\pgfsetbuttcap%
\pgfsetroundjoin%
\definecolor{currentfill}{rgb}{0.983391,0.983391,0.983391}%
\pgfsetfillcolor{currentfill}%
\pgfsetfillopacity{0.900000}%
\pgfsetlinewidth{0.501875pt}%
\definecolor{currentstroke}{rgb}{0.200000,0.200000,0.200000}%
\pgfsetstrokecolor{currentstroke}%
\pgfsetdash{}{0pt}%
\pgfpathmoveto{\pgfqpoint{2.642929in}{2.032363in}}%
\pgfpathlineto{\pgfqpoint{2.686654in}{2.022537in}}%
\pgfpathlineto{\pgfqpoint{2.717582in}{2.005516in}}%
\pgfpathlineto{\pgfqpoint{2.673835in}{2.015357in}}%
\pgfpathlineto{\pgfqpoint{2.642929in}{2.032363in}}%
\pgfpathclose%
\pgfusepath{stroke,fill}%
\end{pgfscope}%
\begin{pgfscope}%
\pgfpathrectangle{\pgfqpoint{0.050000in}{0.050000in}}{\pgfqpoint{4.900000in}{4.900000in}}%
\pgfusepath{clip}%
\pgfsetbuttcap%
\pgfsetroundjoin%
\definecolor{currentfill}{rgb}{0.970473,0.970473,0.970473}%
\pgfsetfillcolor{currentfill}%
\pgfsetfillopacity{0.900000}%
\pgfsetlinewidth{0.501875pt}%
\definecolor{currentstroke}{rgb}{0.200000,0.200000,0.200000}%
\pgfsetstrokecolor{currentstroke}%
\pgfsetdash{}{0pt}%
\pgfpathmoveto{\pgfqpoint{2.493635in}{2.104434in}}%
\pgfpathlineto{\pgfqpoint{2.537592in}{2.082029in}}%
\pgfpathlineto{\pgfqpoint{2.568288in}{2.065134in}}%
\pgfpathlineto{\pgfqpoint{2.524316in}{2.087145in}}%
\pgfpathlineto{\pgfqpoint{2.493635in}{2.104434in}}%
\pgfpathclose%
\pgfusepath{stroke,fill}%
\end{pgfscope}%
\begin{pgfscope}%
\pgfpathrectangle{\pgfqpoint{0.050000in}{0.050000in}}{\pgfqpoint{4.900000in}{4.900000in}}%
\pgfusepath{clip}%
\pgfsetbuttcap%
\pgfsetroundjoin%
\definecolor{currentfill}{rgb}{0.855932,0.855932,0.855932}%
\pgfsetfillcolor{currentfill}%
\pgfsetfillopacity{0.900000}%
\pgfsetlinewidth{0.501875pt}%
\definecolor{currentstroke}{rgb}{0.200000,0.200000,0.200000}%
\pgfsetstrokecolor{currentstroke}%
\pgfsetdash{}{0pt}%
\pgfpathmoveto{\pgfqpoint{1.909402in}{2.544573in}}%
\pgfpathlineto{\pgfqpoint{1.954570in}{2.501596in}}%
\pgfpathlineto{\pgfqpoint{1.984395in}{2.482192in}}%
\pgfpathlineto{\pgfqpoint{1.939221in}{2.525193in}}%
\pgfpathlineto{\pgfqpoint{1.909402in}{2.544573in}}%
\pgfpathclose%
\pgfusepath{stroke,fill}%
\end{pgfscope}%
\begin{pgfscope}%
\pgfpathrectangle{\pgfqpoint{0.050000in}{0.050000in}}{\pgfqpoint{4.900000in}{4.900000in}}%
\pgfusepath{clip}%
\pgfsetbuttcap%
\pgfsetroundjoin%
\definecolor{currentfill}{rgb}{0.988927,0.988927,0.988927}%
\pgfsetfillcolor{currentfill}%
\pgfsetfillopacity{0.900000}%
\pgfsetlinewidth{0.501875pt}%
\definecolor{currentstroke}{rgb}{0.200000,0.200000,0.200000}%
\pgfsetstrokecolor{currentstroke}%
\pgfsetdash{}{0pt}%
\pgfpathmoveto{\pgfqpoint{2.717582in}{2.005516in}}%
\pgfpathlineto{\pgfqpoint{2.761218in}{2.001021in}}%
\pgfpathlineto{\pgfqpoint{2.792266in}{1.983714in}}%
\pgfpathlineto{\pgfqpoint{2.748605in}{1.988363in}}%
\pgfpathlineto{\pgfqpoint{2.717582in}{2.005516in}}%
\pgfpathclose%
\pgfusepath{stroke,fill}%
\end{pgfscope}%
\begin{pgfscope}%
\pgfpathrectangle{\pgfqpoint{0.050000in}{0.050000in}}{\pgfqpoint{4.900000in}{4.900000in}}%
\pgfusepath{clip}%
\pgfsetbuttcap%
\pgfsetroundjoin%
\definecolor{currentfill}{rgb}{0.998155,0.998155,0.998155}%
\pgfsetfillcolor{currentfill}%
\pgfsetfillopacity{0.900000}%
\pgfsetlinewidth{0.501875pt}%
\definecolor{currentstroke}{rgb}{0.200000,0.200000,0.200000}%
\pgfsetstrokecolor{currentstroke}%
\pgfsetdash{}{0pt}%
\pgfpathmoveto{\pgfqpoint{3.060019in}{1.947678in}}%
\pgfpathlineto{\pgfqpoint{3.103320in}{1.957695in}}%
\pgfpathlineto{\pgfqpoint{3.134884in}{1.938900in}}%
\pgfpathlineto{\pgfqpoint{3.091556in}{1.929078in}}%
\pgfpathlineto{\pgfqpoint{3.060019in}{1.947678in}}%
\pgfpathclose%
\pgfusepath{stroke,fill}%
\end{pgfscope}%
\begin{pgfscope}%
\pgfpathrectangle{\pgfqpoint{0.050000in}{0.050000in}}{\pgfqpoint{4.900000in}{4.900000in}}%
\pgfusepath{clip}%
\pgfsetbuttcap%
\pgfsetroundjoin%
\definecolor{currentfill}{rgb}{1.000000,1.000000,1.000000}%
\pgfsetfillcolor{currentfill}%
\pgfsetfillopacity{0.900000}%
\pgfsetlinewidth{0.501875pt}%
\definecolor{currentstroke}{rgb}{0.200000,0.200000,0.200000}%
\pgfsetstrokecolor{currentstroke}%
\pgfsetdash{}{0pt}%
\pgfpathmoveto{\pgfqpoint{3.520997in}{1.936549in}}%
\pgfpathlineto{\pgfqpoint{3.563840in}{1.949828in}}%
\pgfpathlineto{\pgfqpoint{3.596091in}{1.930345in}}%
\pgfpathlineto{\pgfqpoint{3.553219in}{1.917081in}}%
\pgfpathlineto{\pgfqpoint{3.520997in}{1.936549in}}%
\pgfpathclose%
\pgfusepath{stroke,fill}%
\end{pgfscope}%
\begin{pgfscope}%
\pgfpathrectangle{\pgfqpoint{0.050000in}{0.050000in}}{\pgfqpoint{4.900000in}{4.900000in}}%
\pgfusepath{clip}%
\pgfsetbuttcap%
\pgfsetroundjoin%
\definecolor{currentfill}{rgb}{0.961246,0.961246,0.961246}%
\pgfsetfillcolor{currentfill}%
\pgfsetfillopacity{0.900000}%
\pgfsetlinewidth{0.501875pt}%
\definecolor{currentstroke}{rgb}{0.200000,0.200000,0.200000}%
\pgfsetstrokecolor{currentstroke}%
\pgfsetdash{}{0pt}%
\pgfpathmoveto{\pgfqpoint{2.418942in}{2.150369in}}%
\pgfpathlineto{\pgfqpoint{2.463048in}{2.121612in}}%
\pgfpathlineto{\pgfqpoint{2.493635in}{2.104434in}}%
\pgfpathlineto{\pgfqpoint{2.449517in}{2.132615in}}%
\pgfpathlineto{\pgfqpoint{2.418942in}{2.150369in}}%
\pgfpathclose%
\pgfusepath{stroke,fill}%
\end{pgfscope}%
\begin{pgfscope}%
\pgfpathrectangle{\pgfqpoint{0.050000in}{0.050000in}}{\pgfqpoint{4.900000in}{4.900000in}}%
\pgfusepath{clip}%
\pgfsetbuttcap%
\pgfsetroundjoin%
\definecolor{currentfill}{rgb}{1.000000,1.000000,1.000000}%
\pgfsetfillcolor{currentfill}%
\pgfsetfillopacity{0.900000}%
\pgfsetlinewidth{0.501875pt}%
\definecolor{currentstroke}{rgb}{0.200000,0.200000,0.200000}%
\pgfsetstrokecolor{currentstroke}%
\pgfsetdash{}{0pt}%
\pgfpathmoveto{\pgfqpoint{3.327959in}{1.936139in}}%
\pgfpathlineto{\pgfqpoint{3.371008in}{1.948976in}}%
\pgfpathlineto{\pgfqpoint{3.402977in}{1.929663in}}%
\pgfpathlineto{\pgfqpoint{3.359899in}{1.916879in}}%
\pgfpathlineto{\pgfqpoint{3.327959in}{1.936139in}}%
\pgfpathclose%
\pgfusepath{stroke,fill}%
\end{pgfscope}%
\begin{pgfscope}%
\pgfpathrectangle{\pgfqpoint{0.050000in}{0.050000in}}{\pgfqpoint{4.900000in}{4.900000in}}%
\pgfusepath{clip}%
\pgfsetbuttcap%
\pgfsetroundjoin%
\definecolor{currentfill}{rgb}{0.992618,0.992618,0.992618}%
\pgfsetfillcolor{currentfill}%
\pgfsetfillopacity{0.900000}%
\pgfsetlinewidth{0.501875pt}%
\definecolor{currentstroke}{rgb}{0.200000,0.200000,0.200000}%
\pgfsetstrokecolor{currentstroke}%
\pgfsetdash{}{0pt}%
\pgfpathmoveto{\pgfqpoint{2.792266in}{1.983714in}}%
\pgfpathlineto{\pgfqpoint{2.835828in}{1.983654in}}%
\pgfpathlineto{\pgfqpoint{2.866997in}{1.966004in}}%
\pgfpathlineto{\pgfqpoint{2.823409in}{1.966297in}}%
\pgfpathlineto{\pgfqpoint{2.792266in}{1.983714in}}%
\pgfpathclose%
\pgfusepath{stroke,fill}%
\end{pgfscope}%
\begin{pgfscope}%
\pgfpathrectangle{\pgfqpoint{0.050000in}{0.050000in}}{\pgfqpoint{4.900000in}{4.900000in}}%
\pgfusepath{clip}%
\pgfsetbuttcap%
\pgfsetroundjoin%
\definecolor{currentfill}{rgb}{0.898378,0.898378,0.898378}%
\pgfsetfillcolor{currentfill}%
\pgfsetfillopacity{0.900000}%
\pgfsetlinewidth{0.501875pt}%
\definecolor{currentstroke}{rgb}{0.200000,0.200000,0.200000}%
\pgfsetstrokecolor{currentstroke}%
\pgfsetdash{}{0pt}%
\pgfpathmoveto{\pgfqpoint{2.089315in}{2.400412in}}%
\pgfpathlineto{\pgfqpoint{2.134131in}{2.357952in}}%
\pgfpathlineto{\pgfqpoint{2.164234in}{2.338568in}}%
\pgfpathlineto{\pgfqpoint{2.119413in}{2.380892in}}%
\pgfpathlineto{\pgfqpoint{2.089315in}{2.400412in}}%
\pgfpathclose%
\pgfusepath{stroke,fill}%
\end{pgfscope}%
\begin{pgfscope}%
\pgfpathrectangle{\pgfqpoint{0.050000in}{0.050000in}}{\pgfqpoint{4.900000in}{4.900000in}}%
\pgfusepath{clip}%
\pgfsetbuttcap%
\pgfsetroundjoin%
\definecolor{currentfill}{rgb}{1.000000,1.000000,1.000000}%
\pgfsetfillcolor{currentfill}%
\pgfsetfillopacity{0.900000}%
\pgfsetlinewidth{0.501875pt}%
\definecolor{currentstroke}{rgb}{0.200000,0.200000,0.200000}%
\pgfsetstrokecolor{currentstroke}%
\pgfsetdash{}{0pt}%
\pgfpathmoveto{\pgfqpoint{3.134884in}{1.938900in}}%
\pgfpathlineto{\pgfqpoint{3.178130in}{1.949994in}}%
\pgfpathlineto{\pgfqpoint{3.209819in}{1.931008in}}%
\pgfpathlineto{\pgfqpoint{3.166545in}{1.920066in}}%
\pgfpathlineto{\pgfqpoint{3.134884in}{1.938900in}}%
\pgfpathclose%
\pgfusepath{stroke,fill}%
\end{pgfscope}%
\begin{pgfscope}%
\pgfpathrectangle{\pgfqpoint{0.050000in}{0.050000in}}{\pgfqpoint{4.900000in}{4.900000in}}%
\pgfusepath{clip}%
\pgfsetbuttcap%
\pgfsetroundjoin%
\definecolor{currentfill}{rgb}{0.950173,0.950173,0.950173}%
\pgfsetfillcolor{currentfill}%
\pgfsetfillopacity{0.900000}%
\pgfsetlinewidth{0.501875pt}%
\definecolor{currentstroke}{rgb}{0.200000,0.200000,0.200000}%
\pgfsetstrokecolor{currentstroke}%
\pgfsetdash{}{0pt}%
\pgfpathmoveto{\pgfqpoint{2.344187in}{2.202366in}}%
\pgfpathlineto{\pgfqpoint{2.388461in}{2.168064in}}%
\pgfpathlineto{\pgfqpoint{2.418942in}{2.150369in}}%
\pgfpathlineto{\pgfqpoint{2.374661in}{2.184020in}}%
\pgfpathlineto{\pgfqpoint{2.344187in}{2.202366in}}%
\pgfpathclose%
\pgfusepath{stroke,fill}%
\end{pgfscope}%
\begin{pgfscope}%
\pgfpathrectangle{\pgfqpoint{0.050000in}{0.050000in}}{\pgfqpoint{4.900000in}{4.900000in}}%
\pgfusepath{clip}%
\pgfsetbuttcap%
\pgfsetroundjoin%
\definecolor{currentfill}{rgb}{0.994464,0.994464,0.994464}%
\pgfsetfillcolor{currentfill}%
\pgfsetfillopacity{0.900000}%
\pgfsetlinewidth{0.501875pt}%
\definecolor{currentstroke}{rgb}{0.200000,0.200000,0.200000}%
\pgfsetstrokecolor{currentstroke}%
\pgfsetdash{}{0pt}%
\pgfpathmoveto{\pgfqpoint{2.866997in}{1.966004in}}%
\pgfpathlineto{\pgfqpoint{2.910494in}{1.969485in}}%
\pgfpathlineto{\pgfqpoint{2.941785in}{1.951487in}}%
\pgfpathlineto{\pgfqpoint{2.898261in}{1.948268in}}%
\pgfpathlineto{\pgfqpoint{2.866997in}{1.966004in}}%
\pgfpathclose%
\pgfusepath{stroke,fill}%
\end{pgfscope}%
\begin{pgfscope}%
\pgfpathrectangle{\pgfqpoint{0.050000in}{0.050000in}}{\pgfqpoint{4.900000in}{4.900000in}}%
\pgfusepath{clip}%
\pgfsetbuttcap%
\pgfsetroundjoin%
\definecolor{currentfill}{rgb}{0.829450,0.829450,0.829450}%
\pgfsetfillcolor{currentfill}%
\pgfsetfillopacity{0.900000}%
\pgfsetlinewidth{0.501875pt}%
\definecolor{currentstroke}{rgb}{0.200000,0.200000,0.200000}%
\pgfsetstrokecolor{currentstroke}%
\pgfsetdash{}{0pt}%
\pgfpathmoveto{\pgfqpoint{1.834327in}{2.607023in}}%
\pgfpathlineto{\pgfqpoint{1.879672in}{2.563895in}}%
\pgfpathlineto{\pgfqpoint{1.909402in}{2.544573in}}%
\pgfpathlineto{\pgfqpoint{1.864050in}{2.587726in}}%
\pgfpathlineto{\pgfqpoint{1.834327in}{2.607023in}}%
\pgfpathclose%
\pgfusepath{stroke,fill}%
\end{pgfscope}%
\begin{pgfscope}%
\pgfpathrectangle{\pgfqpoint{0.050000in}{0.050000in}}{\pgfqpoint{4.900000in}{4.900000in}}%
\pgfusepath{clip}%
\pgfsetbuttcap%
\pgfsetroundjoin%
\definecolor{currentfill}{rgb}{1.000000,1.000000,1.000000}%
\pgfsetfillcolor{currentfill}%
\pgfsetfillopacity{0.900000}%
\pgfsetlinewidth{0.501875pt}%
\definecolor{currentstroke}{rgb}{0.200000,0.200000,0.200000}%
\pgfsetstrokecolor{currentstroke}%
\pgfsetdash{}{0pt}%
\pgfpathmoveto{\pgfqpoint{3.402977in}{1.929663in}}%
\pgfpathlineto{\pgfqpoint{3.445968in}{1.942733in}}%
\pgfpathlineto{\pgfqpoint{3.478064in}{1.923340in}}%
\pgfpathlineto{\pgfqpoint{3.435044in}{1.910304in}}%
\pgfpathlineto{\pgfqpoint{3.402977in}{1.929663in}}%
\pgfpathclose%
\pgfusepath{stroke,fill}%
\end{pgfscope}%
\begin{pgfscope}%
\pgfpathrectangle{\pgfqpoint{0.050000in}{0.050000in}}{\pgfqpoint{4.900000in}{4.900000in}}%
\pgfusepath{clip}%
\pgfsetbuttcap%
\pgfsetroundjoin%
\definecolor{currentfill}{rgb}{0.996309,0.996309,0.996309}%
\pgfsetfillcolor{currentfill}%
\pgfsetfillopacity{0.900000}%
\pgfsetlinewidth{0.501875pt}%
\definecolor{currentstroke}{rgb}{0.200000,0.200000,0.200000}%
\pgfsetstrokecolor{currentstroke}%
\pgfsetdash{}{0pt}%
\pgfpathmoveto{\pgfqpoint{2.941785in}{1.951487in}}%
\pgfpathlineto{\pgfqpoint{2.985223in}{1.957708in}}%
\pgfpathlineto{\pgfqpoint{3.016637in}{1.939388in}}%
\pgfpathlineto{\pgfqpoint{2.973171in}{1.933423in}}%
\pgfpathlineto{\pgfqpoint{2.941785in}{1.951487in}}%
\pgfpathclose%
\pgfusepath{stroke,fill}%
\end{pgfscope}%
\begin{pgfscope}%
\pgfpathrectangle{\pgfqpoint{0.050000in}{0.050000in}}{\pgfqpoint{4.900000in}{4.900000in}}%
\pgfusepath{clip}%
\pgfsetbuttcap%
\pgfsetroundjoin%
\definecolor{currentfill}{rgb}{1.000000,1.000000,1.000000}%
\pgfsetfillcolor{currentfill}%
\pgfsetfillopacity{0.900000}%
\pgfsetlinewidth{0.501875pt}%
\definecolor{currentstroke}{rgb}{0.200000,0.200000,0.200000}%
\pgfsetstrokecolor{currentstroke}%
\pgfsetdash{}{0pt}%
\pgfpathmoveto{\pgfqpoint{3.209819in}{1.931008in}}%
\pgfpathlineto{\pgfqpoint{3.253010in}{1.942871in}}%
\pgfpathlineto{\pgfqpoint{3.284824in}{1.923732in}}%
\pgfpathlineto{\pgfqpoint{3.241605in}{1.911983in}}%
\pgfpathlineto{\pgfqpoint{3.209819in}{1.931008in}}%
\pgfpathclose%
\pgfusepath{stroke,fill}%
\end{pgfscope}%
\begin{pgfscope}%
\pgfpathrectangle{\pgfqpoint{0.050000in}{0.050000in}}{\pgfqpoint{4.900000in}{4.900000in}}%
\pgfusepath{clip}%
\pgfsetbuttcap%
\pgfsetroundjoin%
\definecolor{currentfill}{rgb}{0.937993,0.937993,0.937993}%
\pgfsetfillcolor{currentfill}%
\pgfsetfillopacity{0.900000}%
\pgfsetlinewidth{0.501875pt}%
\definecolor{currentstroke}{rgb}{0.200000,0.200000,0.200000}%
\pgfsetstrokecolor{currentstroke}%
\pgfsetdash{}{0pt}%
\pgfpathmoveto{\pgfqpoint{2.269351in}{2.259188in}}%
\pgfpathlineto{\pgfqpoint{2.313805in}{2.220712in}}%
\pgfpathlineto{\pgfqpoint{2.344187in}{2.202366in}}%
\pgfpathlineto{\pgfqpoint{2.299728in}{2.240263in}}%
\pgfpathlineto{\pgfqpoint{2.269351in}{2.259188in}}%
\pgfpathclose%
\pgfusepath{stroke,fill}%
\end{pgfscope}%
\begin{pgfscope}%
\pgfpathrectangle{\pgfqpoint{0.050000in}{0.050000in}}{\pgfqpoint{4.900000in}{4.900000in}}%
\pgfusepath{clip}%
\pgfsetbuttcap%
\pgfsetroundjoin%
\definecolor{currentfill}{rgb}{0.875740,0.875740,0.875740}%
\pgfsetfillcolor{currentfill}%
\pgfsetfillopacity{0.900000}%
\pgfsetlinewidth{0.501875pt}%
\definecolor{currentstroke}{rgb}{0.200000,0.200000,0.200000}%
\pgfsetstrokecolor{currentstroke}%
\pgfsetdash{}{0pt}%
\pgfpathmoveto{\pgfqpoint{2.014310in}{2.462734in}}%
\pgfpathlineto{\pgfqpoint{2.059305in}{2.419909in}}%
\pgfpathlineto{\pgfqpoint{2.089315in}{2.400412in}}%
\pgfpathlineto{\pgfqpoint{2.044313in}{2.443225in}}%
\pgfpathlineto{\pgfqpoint{2.014310in}{2.462734in}}%
\pgfpathclose%
\pgfusepath{stroke,fill}%
\end{pgfscope}%
\begin{pgfscope}%
\pgfpathrectangle{\pgfqpoint{0.050000in}{0.050000in}}{\pgfqpoint{4.900000in}{4.900000in}}%
\pgfusepath{clip}%
\pgfsetbuttcap%
\pgfsetroundjoin%
\definecolor{currentfill}{rgb}{1.000000,1.000000,1.000000}%
\pgfsetfillcolor{currentfill}%
\pgfsetfillopacity{0.900000}%
\pgfsetlinewidth{0.501875pt}%
\definecolor{currentstroke}{rgb}{0.200000,0.200000,0.200000}%
\pgfsetstrokecolor{currentstroke}%
\pgfsetdash{}{0pt}%
\pgfpathmoveto{\pgfqpoint{3.478064in}{1.923340in}}%
\pgfpathlineto{\pgfqpoint{3.520997in}{1.936549in}}%
\pgfpathlineto{\pgfqpoint{3.553219in}{1.917081in}}%
\pgfpathlineto{\pgfqpoint{3.510257in}{1.903896in}}%
\pgfpathlineto{\pgfqpoint{3.478064in}{1.923340in}}%
\pgfpathclose%
\pgfusepath{stroke,fill}%
\end{pgfscope}%
\begin{pgfscope}%
\pgfpathrectangle{\pgfqpoint{0.050000in}{0.050000in}}{\pgfqpoint{4.900000in}{4.900000in}}%
\pgfusepath{clip}%
\pgfsetbuttcap%
\pgfsetroundjoin%
\definecolor{currentfill}{rgb}{0.998155,0.998155,0.998155}%
\pgfsetfillcolor{currentfill}%
\pgfsetfillopacity{0.900000}%
\pgfsetlinewidth{0.501875pt}%
\definecolor{currentstroke}{rgb}{0.200000,0.200000,0.200000}%
\pgfsetstrokecolor{currentstroke}%
\pgfsetdash{}{0pt}%
\pgfpathmoveto{\pgfqpoint{3.016637in}{1.939388in}}%
\pgfpathlineto{\pgfqpoint{3.060019in}{1.947678in}}%
\pgfpathlineto{\pgfqpoint{3.091556in}{1.929078in}}%
\pgfpathlineto{\pgfqpoint{3.048146in}{1.921016in}}%
\pgfpathlineto{\pgfqpoint{3.016637in}{1.939388in}}%
\pgfpathclose%
\pgfusepath{stroke,fill}%
\end{pgfscope}%
\begin{pgfscope}%
\pgfpathrectangle{\pgfqpoint{0.050000in}{0.050000in}}{\pgfqpoint{4.900000in}{4.900000in}}%
\pgfusepath{clip}%
\pgfsetbuttcap%
\pgfsetroundjoin%
\definecolor{currentfill}{rgb}{1.000000,1.000000,1.000000}%
\pgfsetfillcolor{currentfill}%
\pgfsetfillopacity{0.900000}%
\pgfsetlinewidth{0.501875pt}%
\definecolor{currentstroke}{rgb}{0.200000,0.200000,0.200000}%
\pgfsetstrokecolor{currentstroke}%
\pgfsetdash{}{0pt}%
\pgfpathmoveto{\pgfqpoint{3.284824in}{1.923732in}}%
\pgfpathlineto{\pgfqpoint{3.327959in}{1.936139in}}%
\pgfpathlineto{\pgfqpoint{3.359899in}{1.916879in}}%
\pgfpathlineto{\pgfqpoint{3.316735in}{1.904553in}}%
\pgfpathlineto{\pgfqpoint{3.284824in}{1.923732in}}%
\pgfpathclose%
\pgfusepath{stroke,fill}%
\end{pgfscope}%
\begin{pgfscope}%
\pgfpathrectangle{\pgfqpoint{0.050000in}{0.050000in}}{\pgfqpoint{4.900000in}{4.900000in}}%
\pgfusepath{clip}%
\pgfsetbuttcap%
\pgfsetroundjoin%
\definecolor{currentfill}{rgb}{0.918185,0.918185,0.918185}%
\pgfsetfillcolor{currentfill}%
\pgfsetfillopacity{0.900000}%
\pgfsetlinewidth{0.501875pt}%
\definecolor{currentstroke}{rgb}{0.200000,0.200000,0.200000}%
\pgfsetstrokecolor{currentstroke}%
\pgfsetdash{}{0pt}%
\pgfpathmoveto{\pgfqpoint{2.194427in}{2.319206in}}%
\pgfpathlineto{\pgfqpoint{2.239065in}{2.278142in}}%
\pgfpathlineto{\pgfqpoint{2.269351in}{2.259188in}}%
\pgfpathlineto{\pgfqpoint{2.224709in}{2.299864in}}%
\pgfpathlineto{\pgfqpoint{2.194427in}{2.319206in}}%
\pgfpathclose%
\pgfusepath{stroke,fill}%
\end{pgfscope}%
\begin{pgfscope}%
\pgfpathrectangle{\pgfqpoint{0.050000in}{0.050000in}}{\pgfqpoint{4.900000in}{4.900000in}}%
\pgfusepath{clip}%
\pgfsetbuttcap%
\pgfsetroundjoin%
\definecolor{currentfill}{rgb}{0.977855,0.977855,0.977855}%
\pgfsetfillcolor{currentfill}%
\pgfsetfillopacity{0.900000}%
\pgfsetlinewidth{0.501875pt}%
\definecolor{currentstroke}{rgb}{0.200000,0.200000,0.200000}%
\pgfsetstrokecolor{currentstroke}%
\pgfsetdash{}{0pt}%
\pgfpathmoveto{\pgfqpoint{2.599080in}{2.048097in}}%
\pgfpathlineto{\pgfqpoint{2.642929in}{2.032363in}}%
\pgfpathlineto{\pgfqpoint{2.673835in}{2.015357in}}%
\pgfpathlineto{\pgfqpoint{2.629967in}{2.030925in}}%
\pgfpathlineto{\pgfqpoint{2.599080in}{2.048097in}}%
\pgfpathclose%
\pgfusepath{stroke,fill}%
\end{pgfscope}%
\begin{pgfscope}%
\pgfpathrectangle{\pgfqpoint{0.050000in}{0.050000in}}{\pgfqpoint{4.900000in}{4.900000in}}%
\pgfusepath{clip}%
\pgfsetbuttcap%
\pgfsetroundjoin%
\definecolor{currentfill}{rgb}{0.983391,0.983391,0.983391}%
\pgfsetfillcolor{currentfill}%
\pgfsetfillopacity{0.900000}%
\pgfsetlinewidth{0.501875pt}%
\definecolor{currentstroke}{rgb}{0.200000,0.200000,0.200000}%
\pgfsetstrokecolor{currentstroke}%
\pgfsetdash{}{0pt}%
\pgfpathmoveto{\pgfqpoint{2.673835in}{2.015357in}}%
\pgfpathlineto{\pgfqpoint{2.717582in}{2.005516in}}%
\pgfpathlineto{\pgfqpoint{2.748605in}{1.988363in}}%
\pgfpathlineto{\pgfqpoint{2.704837in}{1.998213in}}%
\pgfpathlineto{\pgfqpoint{2.673835in}{2.015357in}}%
\pgfpathclose%
\pgfusepath{stroke,fill}%
\end{pgfscope}%
\begin{pgfscope}%
\pgfpathrectangle{\pgfqpoint{0.050000in}{0.050000in}}{\pgfqpoint{4.900000in}{4.900000in}}%
\pgfusepath{clip}%
\pgfsetbuttcap%
\pgfsetroundjoin%
\definecolor{currentfill}{rgb}{0.855932,0.855932,0.855932}%
\pgfsetfillcolor{currentfill}%
\pgfsetfillopacity{0.900000}%
\pgfsetlinewidth{0.501875pt}%
\definecolor{currentstroke}{rgb}{0.200000,0.200000,0.200000}%
\pgfsetstrokecolor{currentstroke}%
\pgfsetdash{}{0pt}%
\pgfpathmoveto{\pgfqpoint{1.939221in}{2.525193in}}%
\pgfpathlineto{\pgfqpoint{1.984395in}{2.482192in}}%
\pgfpathlineto{\pgfqpoint{2.014310in}{2.462734in}}%
\pgfpathlineto{\pgfqpoint{1.969129in}{2.505756in}}%
\pgfpathlineto{\pgfqpoint{1.939221in}{2.525193in}}%
\pgfpathclose%
\pgfusepath{stroke,fill}%
\end{pgfscope}%
\begin{pgfscope}%
\pgfpathrectangle{\pgfqpoint{0.050000in}{0.050000in}}{\pgfqpoint{4.900000in}{4.900000in}}%
\pgfusepath{clip}%
\pgfsetbuttcap%
\pgfsetroundjoin%
\definecolor{currentfill}{rgb}{0.970473,0.970473,0.970473}%
\pgfsetfillcolor{currentfill}%
\pgfsetfillopacity{0.900000}%
\pgfsetlinewidth{0.501875pt}%
\definecolor{currentstroke}{rgb}{0.200000,0.200000,0.200000}%
\pgfsetstrokecolor{currentstroke}%
\pgfsetdash{}{0pt}%
\pgfpathmoveto{\pgfqpoint{2.524316in}{2.087145in}}%
\pgfpathlineto{\pgfqpoint{2.568288in}{2.065134in}}%
\pgfpathlineto{\pgfqpoint{2.599080in}{2.048097in}}%
\pgfpathlineto{\pgfqpoint{2.555091in}{2.069747in}}%
\pgfpathlineto{\pgfqpoint{2.524316in}{2.087145in}}%
\pgfpathclose%
\pgfusepath{stroke,fill}%
\end{pgfscope}%
\begin{pgfscope}%
\pgfpathrectangle{\pgfqpoint{0.050000in}{0.050000in}}{\pgfqpoint{4.900000in}{4.900000in}}%
\pgfusepath{clip}%
\pgfsetbuttcap%
\pgfsetroundjoin%
\definecolor{currentfill}{rgb}{1.000000,1.000000,1.000000}%
\pgfsetfillcolor{currentfill}%
\pgfsetfillopacity{0.900000}%
\pgfsetlinewidth{0.501875pt}%
\definecolor{currentstroke}{rgb}{0.200000,0.200000,0.200000}%
\pgfsetstrokecolor{currentstroke}%
\pgfsetdash{}{0pt}%
\pgfpathmoveto{\pgfqpoint{3.553219in}{1.917081in}}%
\pgfpathlineto{\pgfqpoint{3.596091in}{1.930345in}}%
\pgfpathlineto{\pgfqpoint{3.628440in}{1.910803in}}%
\pgfpathlineto{\pgfqpoint{3.585538in}{1.897559in}}%
\pgfpathlineto{\pgfqpoint{3.553219in}{1.917081in}}%
\pgfpathclose%
\pgfusepath{stroke,fill}%
\end{pgfscope}%
\begin{pgfscope}%
\pgfpathrectangle{\pgfqpoint{0.050000in}{0.050000in}}{\pgfqpoint{4.900000in}{4.900000in}}%
\pgfusepath{clip}%
\pgfsetbuttcap%
\pgfsetroundjoin%
\definecolor{currentfill}{rgb}{0.998155,0.998155,0.998155}%
\pgfsetfillcolor{currentfill}%
\pgfsetfillopacity{0.900000}%
\pgfsetlinewidth{0.501875pt}%
\definecolor{currentstroke}{rgb}{0.200000,0.200000,0.200000}%
\pgfsetstrokecolor{currentstroke}%
\pgfsetdash{}{0pt}%
\pgfpathmoveto{\pgfqpoint{3.091556in}{1.929078in}}%
\pgfpathlineto{\pgfqpoint{3.134884in}{1.938900in}}%
\pgfpathlineto{\pgfqpoint{3.166545in}{1.920066in}}%
\pgfpathlineto{\pgfqpoint{3.123189in}{1.910433in}}%
\pgfpathlineto{\pgfqpoint{3.091556in}{1.929078in}}%
\pgfpathclose%
\pgfusepath{stroke,fill}%
\end{pgfscope}%
\begin{pgfscope}%
\pgfpathrectangle{\pgfqpoint{0.050000in}{0.050000in}}{\pgfqpoint{4.900000in}{4.900000in}}%
\pgfusepath{clip}%
\pgfsetbuttcap%
\pgfsetroundjoin%
\definecolor{currentfill}{rgb}{0.988927,0.988927,0.988927}%
\pgfsetfillcolor{currentfill}%
\pgfsetfillopacity{0.900000}%
\pgfsetlinewidth{0.501875pt}%
\definecolor{currentstroke}{rgb}{0.200000,0.200000,0.200000}%
\pgfsetstrokecolor{currentstroke}%
\pgfsetdash{}{0pt}%
\pgfpathmoveto{\pgfqpoint{2.748605in}{1.988363in}}%
\pgfpathlineto{\pgfqpoint{2.792266in}{1.983714in}}%
\pgfpathlineto{\pgfqpoint{2.823409in}{1.966297in}}%
\pgfpathlineto{\pgfqpoint{2.779724in}{1.971083in}}%
\pgfpathlineto{\pgfqpoint{2.748605in}{1.988363in}}%
\pgfpathclose%
\pgfusepath{stroke,fill}%
\end{pgfscope}%
\begin{pgfscope}%
\pgfpathrectangle{\pgfqpoint{0.050000in}{0.050000in}}{\pgfqpoint{4.900000in}{4.900000in}}%
\pgfusepath{clip}%
\pgfsetbuttcap%
\pgfsetroundjoin%
\definecolor{currentfill}{rgb}{0.961246,0.961246,0.961246}%
\pgfsetfillcolor{currentfill}%
\pgfsetfillopacity{0.900000}%
\pgfsetlinewidth{0.501875pt}%
\definecolor{currentstroke}{rgb}{0.200000,0.200000,0.200000}%
\pgfsetstrokecolor{currentstroke}%
\pgfsetdash{}{0pt}%
\pgfpathmoveto{\pgfqpoint{2.449517in}{2.132615in}}%
\pgfpathlineto{\pgfqpoint{2.493635in}{2.104434in}}%
\pgfpathlineto{\pgfqpoint{2.524316in}{2.087145in}}%
\pgfpathlineto{\pgfqpoint{2.480186in}{2.114799in}}%
\pgfpathlineto{\pgfqpoint{2.449517in}{2.132615in}}%
\pgfpathclose%
\pgfusepath{stroke,fill}%
\end{pgfscope}%
\begin{pgfscope}%
\pgfpathrectangle{\pgfqpoint{0.050000in}{0.050000in}}{\pgfqpoint{4.900000in}{4.900000in}}%
\pgfusepath{clip}%
\pgfsetbuttcap%
\pgfsetroundjoin%
\definecolor{currentfill}{rgb}{1.000000,1.000000,1.000000}%
\pgfsetfillcolor{currentfill}%
\pgfsetfillopacity{0.900000}%
\pgfsetlinewidth{0.501875pt}%
\definecolor{currentstroke}{rgb}{0.200000,0.200000,0.200000}%
\pgfsetstrokecolor{currentstroke}%
\pgfsetdash{}{0pt}%
\pgfpathmoveto{\pgfqpoint{3.359899in}{1.916879in}}%
\pgfpathlineto{\pgfqpoint{3.402977in}{1.929663in}}%
\pgfpathlineto{\pgfqpoint{3.435044in}{1.910304in}}%
\pgfpathlineto{\pgfqpoint{3.391936in}{1.897575in}}%
\pgfpathlineto{\pgfqpoint{3.359899in}{1.916879in}}%
\pgfpathclose%
\pgfusepath{stroke,fill}%
\end{pgfscope}%
\begin{pgfscope}%
\pgfpathrectangle{\pgfqpoint{0.050000in}{0.050000in}}{\pgfqpoint{4.900000in}{4.900000in}}%
\pgfusepath{clip}%
\pgfsetbuttcap%
\pgfsetroundjoin%
\definecolor{currentfill}{rgb}{0.990773,0.990773,0.990773}%
\pgfsetfillcolor{currentfill}%
\pgfsetfillopacity{0.900000}%
\pgfsetlinewidth{0.501875pt}%
\definecolor{currentstroke}{rgb}{0.200000,0.200000,0.200000}%
\pgfsetstrokecolor{currentstroke}%
\pgfsetdash{}{0pt}%
\pgfpathmoveto{\pgfqpoint{2.823409in}{1.966297in}}%
\pgfpathlineto{\pgfqpoint{2.866997in}{1.966004in}}%
\pgfpathlineto{\pgfqpoint{2.898261in}{1.948268in}}%
\pgfpathlineto{\pgfqpoint{2.854648in}{1.948774in}}%
\pgfpathlineto{\pgfqpoint{2.823409in}{1.966297in}}%
\pgfpathclose%
\pgfusepath{stroke,fill}%
\end{pgfscope}%
\begin{pgfscope}%
\pgfpathrectangle{\pgfqpoint{0.050000in}{0.050000in}}{\pgfqpoint{4.900000in}{4.900000in}}%
\pgfusepath{clip}%
\pgfsetbuttcap%
\pgfsetroundjoin%
\definecolor{currentfill}{rgb}{0.898378,0.898378,0.898378}%
\pgfsetfillcolor{currentfill}%
\pgfsetfillopacity{0.900000}%
\pgfsetlinewidth{0.501875pt}%
\definecolor{currentstroke}{rgb}{0.200000,0.200000,0.200000}%
\pgfsetstrokecolor{currentstroke}%
\pgfsetdash{}{0pt}%
\pgfpathmoveto{\pgfqpoint{2.119413in}{2.380892in}}%
\pgfpathlineto{\pgfqpoint{2.164234in}{2.338568in}}%
\pgfpathlineto{\pgfqpoint{2.194427in}{2.319206in}}%
\pgfpathlineto{\pgfqpoint{2.149601in}{2.361355in}}%
\pgfpathlineto{\pgfqpoint{2.119413in}{2.380892in}}%
\pgfpathclose%
\pgfusepath{stroke,fill}%
\end{pgfscope}%
\begin{pgfscope}%
\pgfpathrectangle{\pgfqpoint{0.050000in}{0.050000in}}{\pgfqpoint{4.900000in}{4.900000in}}%
\pgfusepath{clip}%
\pgfsetbuttcap%
\pgfsetroundjoin%
\definecolor{currentfill}{rgb}{1.000000,1.000000,1.000000}%
\pgfsetfillcolor{currentfill}%
\pgfsetfillopacity{0.900000}%
\pgfsetlinewidth{0.501875pt}%
\definecolor{currentstroke}{rgb}{0.200000,0.200000,0.200000}%
\pgfsetstrokecolor{currentstroke}%
\pgfsetdash{}{0pt}%
\pgfpathmoveto{\pgfqpoint{3.166545in}{1.920066in}}%
\pgfpathlineto{\pgfqpoint{3.209819in}{1.931008in}}%
\pgfpathlineto{\pgfqpoint{3.241605in}{1.911983in}}%
\pgfpathlineto{\pgfqpoint{3.198302in}{1.901191in}}%
\pgfpathlineto{\pgfqpoint{3.166545in}{1.920066in}}%
\pgfpathclose%
\pgfusepath{stroke,fill}%
\end{pgfscope}%
\begin{pgfscope}%
\pgfpathrectangle{\pgfqpoint{0.050000in}{0.050000in}}{\pgfqpoint{4.900000in}{4.900000in}}%
\pgfusepath{clip}%
\pgfsetbuttcap%
\pgfsetroundjoin%
\definecolor{currentfill}{rgb}{0.950173,0.950173,0.950173}%
\pgfsetfillcolor{currentfill}%
\pgfsetfillopacity{0.900000}%
\pgfsetlinewidth{0.501875pt}%
\definecolor{currentstroke}{rgb}{0.200000,0.200000,0.200000}%
\pgfsetstrokecolor{currentstroke}%
\pgfsetdash{}{0pt}%
\pgfpathmoveto{\pgfqpoint{2.374661in}{2.184020in}}%
\pgfpathlineto{\pgfqpoint{2.418942in}{2.150369in}}%
\pgfpathlineto{\pgfqpoint{2.449517in}{2.132615in}}%
\pgfpathlineto{\pgfqpoint{2.405228in}{2.165660in}}%
\pgfpathlineto{\pgfqpoint{2.374661in}{2.184020in}}%
\pgfpathclose%
\pgfusepath{stroke,fill}%
\end{pgfscope}%
\begin{pgfscope}%
\pgfpathrectangle{\pgfqpoint{0.050000in}{0.050000in}}{\pgfqpoint{4.900000in}{4.900000in}}%
\pgfusepath{clip}%
\pgfsetbuttcap%
\pgfsetroundjoin%
\definecolor{currentfill}{rgb}{0.994464,0.994464,0.994464}%
\pgfsetfillcolor{currentfill}%
\pgfsetfillopacity{0.900000}%
\pgfsetlinewidth{0.501875pt}%
\definecolor{currentstroke}{rgb}{0.200000,0.200000,0.200000}%
\pgfsetstrokecolor{currentstroke}%
\pgfsetdash{}{0pt}%
\pgfpathmoveto{\pgfqpoint{2.898261in}{1.948268in}}%
\pgfpathlineto{\pgfqpoint{2.941785in}{1.951487in}}%
\pgfpathlineto{\pgfqpoint{2.973171in}{1.933423in}}%
\pgfpathlineto{\pgfqpoint{2.929621in}{1.930446in}}%
\pgfpathlineto{\pgfqpoint{2.898261in}{1.948268in}}%
\pgfpathclose%
\pgfusepath{stroke,fill}%
\end{pgfscope}%
\begin{pgfscope}%
\pgfpathrectangle{\pgfqpoint{0.050000in}{0.050000in}}{\pgfqpoint{4.900000in}{4.900000in}}%
\pgfusepath{clip}%
\pgfsetbuttcap%
\pgfsetroundjoin%
\definecolor{currentfill}{rgb}{0.829450,0.829450,0.829450}%
\pgfsetfillcolor{currentfill}%
\pgfsetfillopacity{0.900000}%
\pgfsetlinewidth{0.501875pt}%
\definecolor{currentstroke}{rgb}{0.200000,0.200000,0.200000}%
\pgfsetstrokecolor{currentstroke}%
\pgfsetdash{}{0pt}%
\pgfpathmoveto{\pgfqpoint{1.864050in}{2.587726in}}%
\pgfpathlineto{\pgfqpoint{1.909402in}{2.544573in}}%
\pgfpathlineto{\pgfqpoint{1.939221in}{2.525193in}}%
\pgfpathlineto{\pgfqpoint{1.893862in}{2.568371in}}%
\pgfpathlineto{\pgfqpoint{1.864050in}{2.587726in}}%
\pgfpathclose%
\pgfusepath{stroke,fill}%
\end{pgfscope}%
\begin{pgfscope}%
\pgfpathrectangle{\pgfqpoint{0.050000in}{0.050000in}}{\pgfqpoint{4.900000in}{4.900000in}}%
\pgfusepath{clip}%
\pgfsetbuttcap%
\pgfsetroundjoin%
\definecolor{currentfill}{rgb}{1.000000,1.000000,1.000000}%
\pgfsetfillcolor{currentfill}%
\pgfsetfillopacity{0.900000}%
\pgfsetlinewidth{0.501875pt}%
\definecolor{currentstroke}{rgb}{0.200000,0.200000,0.200000}%
\pgfsetstrokecolor{currentstroke}%
\pgfsetdash{}{0pt}%
\pgfpathmoveto{\pgfqpoint{3.435044in}{1.910304in}}%
\pgfpathlineto{\pgfqpoint{3.478064in}{1.923340in}}%
\pgfpathlineto{\pgfqpoint{3.510257in}{1.903896in}}%
\pgfpathlineto{\pgfqpoint{3.467207in}{1.890898in}}%
\pgfpathlineto{\pgfqpoint{3.435044in}{1.910304in}}%
\pgfpathclose%
\pgfusepath{stroke,fill}%
\end{pgfscope}%
\begin{pgfscope}%
\pgfpathrectangle{\pgfqpoint{0.050000in}{0.050000in}}{\pgfqpoint{4.900000in}{4.900000in}}%
\pgfusepath{clip}%
\pgfsetbuttcap%
\pgfsetroundjoin%
\definecolor{currentfill}{rgb}{0.996309,0.996309,0.996309}%
\pgfsetfillcolor{currentfill}%
\pgfsetfillopacity{0.900000}%
\pgfsetlinewidth{0.501875pt}%
\definecolor{currentstroke}{rgb}{0.200000,0.200000,0.200000}%
\pgfsetstrokecolor{currentstroke}%
\pgfsetdash{}{0pt}%
\pgfpathmoveto{\pgfqpoint{2.973171in}{1.933423in}}%
\pgfpathlineto{\pgfqpoint{3.016637in}{1.939388in}}%
\pgfpathlineto{\pgfqpoint{3.048146in}{1.921016in}}%
\pgfpathlineto{\pgfqpoint{3.004653in}{1.915289in}}%
\pgfpathlineto{\pgfqpoint{2.973171in}{1.933423in}}%
\pgfpathclose%
\pgfusepath{stroke,fill}%
\end{pgfscope}%
\begin{pgfscope}%
\pgfpathrectangle{\pgfqpoint{0.050000in}{0.050000in}}{\pgfqpoint{4.900000in}{4.900000in}}%
\pgfusepath{clip}%
\pgfsetbuttcap%
\pgfsetroundjoin%
\definecolor{currentfill}{rgb}{1.000000,1.000000,1.000000}%
\pgfsetfillcolor{currentfill}%
\pgfsetfillopacity{0.900000}%
\pgfsetlinewidth{0.501875pt}%
\definecolor{currentstroke}{rgb}{0.200000,0.200000,0.200000}%
\pgfsetstrokecolor{currentstroke}%
\pgfsetdash{}{0pt}%
\pgfpathmoveto{\pgfqpoint{3.241605in}{1.911983in}}%
\pgfpathlineto{\pgfqpoint{3.284824in}{1.923732in}}%
\pgfpathlineto{\pgfqpoint{3.316735in}{1.904553in}}%
\pgfpathlineto{\pgfqpoint{3.273487in}{1.892917in}}%
\pgfpathlineto{\pgfqpoint{3.241605in}{1.911983in}}%
\pgfpathclose%
\pgfusepath{stroke,fill}%
\end{pgfscope}%
\begin{pgfscope}%
\pgfpathrectangle{\pgfqpoint{0.050000in}{0.050000in}}{\pgfqpoint{4.900000in}{4.900000in}}%
\pgfusepath{clip}%
\pgfsetbuttcap%
\pgfsetroundjoin%
\definecolor{currentfill}{rgb}{0.937993,0.937993,0.937993}%
\pgfsetfillcolor{currentfill}%
\pgfsetfillopacity{0.900000}%
\pgfsetlinewidth{0.501875pt}%
\definecolor{currentstroke}{rgb}{0.200000,0.200000,0.200000}%
\pgfsetstrokecolor{currentstroke}%
\pgfsetdash{}{0pt}%
\pgfpathmoveto{\pgfqpoint{2.299728in}{2.240263in}}%
\pgfpathlineto{\pgfqpoint{2.344187in}{2.202366in}}%
\pgfpathlineto{\pgfqpoint{2.374661in}{2.184020in}}%
\pgfpathlineto{\pgfqpoint{2.330197in}{2.221357in}}%
\pgfpathlineto{\pgfqpoint{2.299728in}{2.240263in}}%
\pgfpathclose%
\pgfusepath{stroke,fill}%
\end{pgfscope}%
\begin{pgfscope}%
\pgfpathrectangle{\pgfqpoint{0.050000in}{0.050000in}}{\pgfqpoint{4.900000in}{4.900000in}}%
\pgfusepath{clip}%
\pgfsetbuttcap%
\pgfsetroundjoin%
\definecolor{currentfill}{rgb}{0.875740,0.875740,0.875740}%
\pgfsetfillcolor{currentfill}%
\pgfsetfillopacity{0.900000}%
\pgfsetlinewidth{0.501875pt}%
\definecolor{currentstroke}{rgb}{0.200000,0.200000,0.200000}%
\pgfsetstrokecolor{currentstroke}%
\pgfsetdash{}{0pt}%
\pgfpathmoveto{\pgfqpoint{2.044313in}{2.443225in}}%
\pgfpathlineto{\pgfqpoint{2.089315in}{2.400412in}}%
\pgfpathlineto{\pgfqpoint{2.119413in}{2.380892in}}%
\pgfpathlineto{\pgfqpoint{2.074406in}{2.423669in}}%
\pgfpathlineto{\pgfqpoint{2.044313in}{2.443225in}}%
\pgfpathclose%
\pgfusepath{stroke,fill}%
\end{pgfscope}%
\begin{pgfscope}%
\pgfpathrectangle{\pgfqpoint{0.050000in}{0.050000in}}{\pgfqpoint{4.900000in}{4.900000in}}%
\pgfusepath{clip}%
\pgfsetbuttcap%
\pgfsetroundjoin%
\definecolor{currentfill}{rgb}{1.000000,1.000000,1.000000}%
\pgfsetfillcolor{currentfill}%
\pgfsetfillopacity{0.900000}%
\pgfsetlinewidth{0.501875pt}%
\definecolor{currentstroke}{rgb}{0.200000,0.200000,0.200000}%
\pgfsetstrokecolor{currentstroke}%
\pgfsetdash{}{0pt}%
\pgfpathmoveto{\pgfqpoint{3.510257in}{1.903896in}}%
\pgfpathlineto{\pgfqpoint{3.553219in}{1.917081in}}%
\pgfpathlineto{\pgfqpoint{3.585538in}{1.897559in}}%
\pgfpathlineto{\pgfqpoint{3.542547in}{1.884401in}}%
\pgfpathlineto{\pgfqpoint{3.510257in}{1.903896in}}%
\pgfpathclose%
\pgfusepath{stroke,fill}%
\end{pgfscope}%
\begin{pgfscope}%
\pgfpathrectangle{\pgfqpoint{0.050000in}{0.050000in}}{\pgfqpoint{4.900000in}{4.900000in}}%
\pgfusepath{clip}%
\pgfsetbuttcap%
\pgfsetroundjoin%
\definecolor{currentfill}{rgb}{0.805336,0.805336,0.805336}%
\pgfsetfillcolor{currentfill}%
\pgfsetfillopacity{0.900000}%
\pgfsetlinewidth{0.501875pt}%
\definecolor{currentstroke}{rgb}{0.200000,0.200000,0.200000}%
\pgfsetstrokecolor{currentstroke}%
\pgfsetdash{}{0pt}%
\pgfpathmoveto{\pgfqpoint{1.788796in}{2.650328in}}%
\pgfpathlineto{\pgfqpoint{1.834327in}{2.607023in}}%
\pgfpathlineto{\pgfqpoint{1.864050in}{2.587726in}}%
\pgfpathlineto{\pgfqpoint{1.818512in}{2.631055in}}%
\pgfpathlineto{\pgfqpoint{1.788796in}{2.650328in}}%
\pgfpathclose%
\pgfusepath{stroke,fill}%
\end{pgfscope}%
\begin{pgfscope}%
\pgfpathrectangle{\pgfqpoint{0.050000in}{0.050000in}}{\pgfqpoint{4.900000in}{4.900000in}}%
\pgfusepath{clip}%
\pgfsetbuttcap%
\pgfsetroundjoin%
\definecolor{currentfill}{rgb}{0.998155,0.998155,0.998155}%
\pgfsetfillcolor{currentfill}%
\pgfsetfillopacity{0.900000}%
\pgfsetlinewidth{0.501875pt}%
\definecolor{currentstroke}{rgb}{0.200000,0.200000,0.200000}%
\pgfsetstrokecolor{currentstroke}%
\pgfsetdash{}{0pt}%
\pgfpathmoveto{\pgfqpoint{3.048146in}{1.921016in}}%
\pgfpathlineto{\pgfqpoint{3.091556in}{1.929078in}}%
\pgfpathlineto{\pgfqpoint{3.123189in}{1.910433in}}%
\pgfpathlineto{\pgfqpoint{3.079751in}{1.902587in}}%
\pgfpathlineto{\pgfqpoint{3.048146in}{1.921016in}}%
\pgfpathclose%
\pgfusepath{stroke,fill}%
\end{pgfscope}%
\begin{pgfscope}%
\pgfpathrectangle{\pgfqpoint{0.050000in}{0.050000in}}{\pgfqpoint{4.900000in}{4.900000in}}%
\pgfusepath{clip}%
\pgfsetbuttcap%
\pgfsetroundjoin%
\definecolor{currentfill}{rgb}{1.000000,1.000000,1.000000}%
\pgfsetfillcolor{currentfill}%
\pgfsetfillopacity{0.900000}%
\pgfsetlinewidth{0.501875pt}%
\definecolor{currentstroke}{rgb}{0.200000,0.200000,0.200000}%
\pgfsetstrokecolor{currentstroke}%
\pgfsetdash{}{0pt}%
\pgfpathmoveto{\pgfqpoint{3.316735in}{1.904553in}}%
\pgfpathlineto{\pgfqpoint{3.359899in}{1.916879in}}%
\pgfpathlineto{\pgfqpoint{3.391936in}{1.897575in}}%
\pgfpathlineto{\pgfqpoint{3.348743in}{1.885332in}}%
\pgfpathlineto{\pgfqpoint{3.316735in}{1.904553in}}%
\pgfpathclose%
\pgfusepath{stroke,fill}%
\end{pgfscope}%
\begin{pgfscope}%
\pgfpathrectangle{\pgfqpoint{0.050000in}{0.050000in}}{\pgfqpoint{4.900000in}{4.900000in}}%
\pgfusepath{clip}%
\pgfsetbuttcap%
\pgfsetroundjoin%
\definecolor{currentfill}{rgb}{0.918185,0.918185,0.918185}%
\pgfsetfillcolor{currentfill}%
\pgfsetfillopacity{0.900000}%
\pgfsetlinewidth{0.501875pt}%
\definecolor{currentstroke}{rgb}{0.200000,0.200000,0.200000}%
\pgfsetstrokecolor{currentstroke}%
\pgfsetdash{}{0pt}%
\pgfpathmoveto{\pgfqpoint{2.224709in}{2.299864in}}%
\pgfpathlineto{\pgfqpoint{2.269351in}{2.259188in}}%
\pgfpathlineto{\pgfqpoint{2.299728in}{2.240263in}}%
\pgfpathlineto{\pgfqpoint{2.255082in}{2.280538in}}%
\pgfpathlineto{\pgfqpoint{2.224709in}{2.299864in}}%
\pgfpathclose%
\pgfusepath{stroke,fill}%
\end{pgfscope}%
\begin{pgfscope}%
\pgfpathrectangle{\pgfqpoint{0.050000in}{0.050000in}}{\pgfqpoint{4.900000in}{4.900000in}}%
\pgfusepath{clip}%
\pgfsetbuttcap%
\pgfsetroundjoin%
\definecolor{currentfill}{rgb}{0.977855,0.977855,0.977855}%
\pgfsetfillcolor{currentfill}%
\pgfsetfillopacity{0.900000}%
\pgfsetlinewidth{0.501875pt}%
\definecolor{currentstroke}{rgb}{0.200000,0.200000,0.200000}%
\pgfsetstrokecolor{currentstroke}%
\pgfsetdash{}{0pt}%
\pgfpathmoveto{\pgfqpoint{2.629967in}{2.030925in}}%
\pgfpathlineto{\pgfqpoint{2.673835in}{2.015357in}}%
\pgfpathlineto{\pgfqpoint{2.704837in}{1.998213in}}%
\pgfpathlineto{\pgfqpoint{2.660949in}{2.013626in}}%
\pgfpathlineto{\pgfqpoint{2.629967in}{2.030925in}}%
\pgfpathclose%
\pgfusepath{stroke,fill}%
\end{pgfscope}%
\begin{pgfscope}%
\pgfpathrectangle{\pgfqpoint{0.050000in}{0.050000in}}{\pgfqpoint{4.900000in}{4.900000in}}%
\pgfusepath{clip}%
\pgfsetbuttcap%
\pgfsetroundjoin%
\definecolor{currentfill}{rgb}{0.983391,0.983391,0.983391}%
\pgfsetfillcolor{currentfill}%
\pgfsetfillopacity{0.900000}%
\pgfsetlinewidth{0.501875pt}%
\definecolor{currentstroke}{rgb}{0.200000,0.200000,0.200000}%
\pgfsetstrokecolor{currentstroke}%
\pgfsetdash{}{0pt}%
\pgfpathmoveto{\pgfqpoint{2.704837in}{1.998213in}}%
\pgfpathlineto{\pgfqpoint{2.748605in}{1.988363in}}%
\pgfpathlineto{\pgfqpoint{2.779724in}{1.971083in}}%
\pgfpathlineto{\pgfqpoint{2.735934in}{1.980937in}}%
\pgfpathlineto{\pgfqpoint{2.704837in}{1.998213in}}%
\pgfpathclose%
\pgfusepath{stroke,fill}%
\end{pgfscope}%
\begin{pgfscope}%
\pgfpathrectangle{\pgfqpoint{0.050000in}{0.050000in}}{\pgfqpoint{4.900000in}{4.900000in}}%
\pgfusepath{clip}%
\pgfsetbuttcap%
\pgfsetroundjoin%
\definecolor{currentfill}{rgb}{0.855932,0.855932,0.855932}%
\pgfsetfillcolor{currentfill}%
\pgfsetfillopacity{0.900000}%
\pgfsetlinewidth{0.501875pt}%
\definecolor{currentstroke}{rgb}{0.200000,0.200000,0.200000}%
\pgfsetstrokecolor{currentstroke}%
\pgfsetdash{}{0pt}%
\pgfpathmoveto{\pgfqpoint{1.969129in}{2.505756in}}%
\pgfpathlineto{\pgfqpoint{2.014310in}{2.462734in}}%
\pgfpathlineto{\pgfqpoint{2.044313in}{2.443225in}}%
\pgfpathlineto{\pgfqpoint{1.999127in}{2.486262in}}%
\pgfpathlineto{\pgfqpoint{1.969129in}{2.505756in}}%
\pgfpathclose%
\pgfusepath{stroke,fill}%
\end{pgfscope}%
\begin{pgfscope}%
\pgfpathrectangle{\pgfqpoint{0.050000in}{0.050000in}}{\pgfqpoint{4.900000in}{4.900000in}}%
\pgfusepath{clip}%
\pgfsetbuttcap%
\pgfsetroundjoin%
\definecolor{currentfill}{rgb}{0.970473,0.970473,0.970473}%
\pgfsetfillcolor{currentfill}%
\pgfsetfillopacity{0.900000}%
\pgfsetlinewidth{0.501875pt}%
\definecolor{currentstroke}{rgb}{0.200000,0.200000,0.200000}%
\pgfsetstrokecolor{currentstroke}%
\pgfsetdash{}{0pt}%
\pgfpathmoveto{\pgfqpoint{2.555091in}{2.069747in}}%
\pgfpathlineto{\pgfqpoint{2.599080in}{2.048097in}}%
\pgfpathlineto{\pgfqpoint{2.629967in}{2.030925in}}%
\pgfpathlineto{\pgfqpoint{2.585961in}{2.052243in}}%
\pgfpathlineto{\pgfqpoint{2.555091in}{2.069747in}}%
\pgfpathclose%
\pgfusepath{stroke,fill}%
\end{pgfscope}%
\begin{pgfscope}%
\pgfpathrectangle{\pgfqpoint{0.050000in}{0.050000in}}{\pgfqpoint{4.900000in}{4.900000in}}%
\pgfusepath{clip}%
\pgfsetbuttcap%
\pgfsetroundjoin%
\definecolor{currentfill}{rgb}{1.000000,1.000000,1.000000}%
\pgfsetfillcolor{currentfill}%
\pgfsetfillopacity{0.900000}%
\pgfsetlinewidth{0.501875pt}%
\definecolor{currentstroke}{rgb}{0.200000,0.200000,0.200000}%
\pgfsetstrokecolor{currentstroke}%
\pgfsetdash{}{0pt}%
\pgfpathmoveto{\pgfqpoint{3.585538in}{1.897559in}}%
\pgfpathlineto{\pgfqpoint{3.628440in}{1.910803in}}%
\pgfpathlineto{\pgfqpoint{3.660886in}{1.891203in}}%
\pgfpathlineto{\pgfqpoint{3.617955in}{1.877982in}}%
\pgfpathlineto{\pgfqpoint{3.585538in}{1.897559in}}%
\pgfpathclose%
\pgfusepath{stroke,fill}%
\end{pgfscope}%
\begin{pgfscope}%
\pgfpathrectangle{\pgfqpoint{0.050000in}{0.050000in}}{\pgfqpoint{4.900000in}{4.900000in}}%
\pgfusepath{clip}%
\pgfsetbuttcap%
\pgfsetroundjoin%
\definecolor{currentfill}{rgb}{0.998155,0.998155,0.998155}%
\pgfsetfillcolor{currentfill}%
\pgfsetfillopacity{0.900000}%
\pgfsetlinewidth{0.501875pt}%
\definecolor{currentstroke}{rgb}{0.200000,0.200000,0.200000}%
\pgfsetstrokecolor{currentstroke}%
\pgfsetdash{}{0pt}%
\pgfpathmoveto{\pgfqpoint{3.123189in}{1.910433in}}%
\pgfpathlineto{\pgfqpoint{3.166545in}{1.920066in}}%
\pgfpathlineto{\pgfqpoint{3.198302in}{1.901191in}}%
\pgfpathlineto{\pgfqpoint{3.154918in}{1.891741in}}%
\pgfpathlineto{\pgfqpoint{3.123189in}{1.910433in}}%
\pgfpathclose%
\pgfusepath{stroke,fill}%
\end{pgfscope}%
\begin{pgfscope}%
\pgfpathrectangle{\pgfqpoint{0.050000in}{0.050000in}}{\pgfqpoint{4.900000in}{4.900000in}}%
\pgfusepath{clip}%
\pgfsetbuttcap%
\pgfsetroundjoin%
\definecolor{currentfill}{rgb}{0.987082,0.987082,0.987082}%
\pgfsetfillcolor{currentfill}%
\pgfsetfillopacity{0.900000}%
\pgfsetlinewidth{0.501875pt}%
\definecolor{currentstroke}{rgb}{0.200000,0.200000,0.200000}%
\pgfsetstrokecolor{currentstroke}%
\pgfsetdash{}{0pt}%
\pgfpathmoveto{\pgfqpoint{2.779724in}{1.971083in}}%
\pgfpathlineto{\pgfqpoint{2.823409in}{1.966297in}}%
\pgfpathlineto{\pgfqpoint{2.854648in}{1.948774in}}%
\pgfpathlineto{\pgfqpoint{2.810940in}{1.953682in}}%
\pgfpathlineto{\pgfqpoint{2.779724in}{1.971083in}}%
\pgfpathclose%
\pgfusepath{stroke,fill}%
\end{pgfscope}%
\begin{pgfscope}%
\pgfpathrectangle{\pgfqpoint{0.050000in}{0.050000in}}{\pgfqpoint{4.900000in}{4.900000in}}%
\pgfusepath{clip}%
\pgfsetbuttcap%
\pgfsetroundjoin%
\definecolor{currentfill}{rgb}{1.000000,1.000000,1.000000}%
\pgfsetfillcolor{currentfill}%
\pgfsetfillopacity{0.900000}%
\pgfsetlinewidth{0.501875pt}%
\definecolor{currentstroke}{rgb}{0.200000,0.200000,0.200000}%
\pgfsetstrokecolor{currentstroke}%
\pgfsetdash{}{0pt}%
\pgfpathmoveto{\pgfqpoint{3.391936in}{1.897575in}}%
\pgfpathlineto{\pgfqpoint{3.435044in}{1.910304in}}%
\pgfpathlineto{\pgfqpoint{3.467207in}{1.890898in}}%
\pgfpathlineto{\pgfqpoint{3.424070in}{1.878228in}}%
\pgfpathlineto{\pgfqpoint{3.391936in}{1.897575in}}%
\pgfpathclose%
\pgfusepath{stroke,fill}%
\end{pgfscope}%
\begin{pgfscope}%
\pgfpathrectangle{\pgfqpoint{0.050000in}{0.050000in}}{\pgfqpoint{4.900000in}{4.900000in}}%
\pgfusepath{clip}%
\pgfsetbuttcap%
\pgfsetroundjoin%
\definecolor{currentfill}{rgb}{0.961246,0.961246,0.961246}%
\pgfsetfillcolor{currentfill}%
\pgfsetfillopacity{0.900000}%
\pgfsetlinewidth{0.501875pt}%
\definecolor{currentstroke}{rgb}{0.200000,0.200000,0.200000}%
\pgfsetstrokecolor{currentstroke}%
\pgfsetdash{}{0pt}%
\pgfpathmoveto{\pgfqpoint{2.480186in}{2.114799in}}%
\pgfpathlineto{\pgfqpoint{2.524316in}{2.087145in}}%
\pgfpathlineto{\pgfqpoint{2.555091in}{2.069747in}}%
\pgfpathlineto{\pgfqpoint{2.510949in}{2.096915in}}%
\pgfpathlineto{\pgfqpoint{2.480186in}{2.114799in}}%
\pgfpathclose%
\pgfusepath{stroke,fill}%
\end{pgfscope}%
\begin{pgfscope}%
\pgfpathrectangle{\pgfqpoint{0.050000in}{0.050000in}}{\pgfqpoint{4.900000in}{4.900000in}}%
\pgfusepath{clip}%
\pgfsetbuttcap%
\pgfsetroundjoin%
\definecolor{currentfill}{rgb}{0.990773,0.990773,0.990773}%
\pgfsetfillcolor{currentfill}%
\pgfsetfillopacity{0.900000}%
\pgfsetlinewidth{0.501875pt}%
\definecolor{currentstroke}{rgb}{0.200000,0.200000,0.200000}%
\pgfsetstrokecolor{currentstroke}%
\pgfsetdash{}{0pt}%
\pgfpathmoveto{\pgfqpoint{2.854648in}{1.948774in}}%
\pgfpathlineto{\pgfqpoint{2.898261in}{1.948268in}}%
\pgfpathlineto{\pgfqpoint{2.929621in}{1.930446in}}%
\pgfpathlineto{\pgfqpoint{2.885983in}{1.931146in}}%
\pgfpathlineto{\pgfqpoint{2.854648in}{1.948774in}}%
\pgfpathclose%
\pgfusepath{stroke,fill}%
\end{pgfscope}%
\begin{pgfscope}%
\pgfpathrectangle{\pgfqpoint{0.050000in}{0.050000in}}{\pgfqpoint{4.900000in}{4.900000in}}%
\pgfusepath{clip}%
\pgfsetbuttcap%
\pgfsetroundjoin%
\definecolor{currentfill}{rgb}{0.898378,0.898378,0.898378}%
\pgfsetfillcolor{currentfill}%
\pgfsetfillopacity{0.900000}%
\pgfsetlinewidth{0.501875pt}%
\definecolor{currentstroke}{rgb}{0.200000,0.200000,0.200000}%
\pgfsetstrokecolor{currentstroke}%
\pgfsetdash{}{0pt}%
\pgfpathmoveto{\pgfqpoint{2.149601in}{2.361355in}}%
\pgfpathlineto{\pgfqpoint{2.194427in}{2.319206in}}%
\pgfpathlineto{\pgfqpoint{2.224709in}{2.299864in}}%
\pgfpathlineto{\pgfqpoint{2.179880in}{2.341807in}}%
\pgfpathlineto{\pgfqpoint{2.149601in}{2.361355in}}%
\pgfpathclose%
\pgfusepath{stroke,fill}%
\end{pgfscope}%
\begin{pgfscope}%
\pgfpathrectangle{\pgfqpoint{0.050000in}{0.050000in}}{\pgfqpoint{4.900000in}{4.900000in}}%
\pgfusepath{clip}%
\pgfsetbuttcap%
\pgfsetroundjoin%
\definecolor{currentfill}{rgb}{1.000000,1.000000,1.000000}%
\pgfsetfillcolor{currentfill}%
\pgfsetfillopacity{0.900000}%
\pgfsetlinewidth{0.501875pt}%
\definecolor{currentstroke}{rgb}{0.200000,0.200000,0.200000}%
\pgfsetstrokecolor{currentstroke}%
\pgfsetdash{}{0pt}%
\pgfpathmoveto{\pgfqpoint{3.198302in}{1.901191in}}%
\pgfpathlineto{\pgfqpoint{3.241605in}{1.911983in}}%
\pgfpathlineto{\pgfqpoint{3.273487in}{1.892917in}}%
\pgfpathlineto{\pgfqpoint{3.230156in}{1.882272in}}%
\pgfpathlineto{\pgfqpoint{3.198302in}{1.901191in}}%
\pgfpathclose%
\pgfusepath{stroke,fill}%
\end{pgfscope}%
\begin{pgfscope}%
\pgfpathrectangle{\pgfqpoint{0.050000in}{0.050000in}}{\pgfqpoint{4.900000in}{4.900000in}}%
\pgfusepath{clip}%
\pgfsetbuttcap%
\pgfsetroundjoin%
\definecolor{currentfill}{rgb}{0.950173,0.950173,0.950173}%
\pgfsetfillcolor{currentfill}%
\pgfsetfillopacity{0.900000}%
\pgfsetlinewidth{0.501875pt}%
\definecolor{currentstroke}{rgb}{0.200000,0.200000,0.200000}%
\pgfsetstrokecolor{currentstroke}%
\pgfsetdash{}{0pt}%
\pgfpathmoveto{\pgfqpoint{2.405228in}{2.165660in}}%
\pgfpathlineto{\pgfqpoint{2.449517in}{2.132615in}}%
\pgfpathlineto{\pgfqpoint{2.480186in}{2.114799in}}%
\pgfpathlineto{\pgfqpoint{2.435887in}{2.147277in}}%
\pgfpathlineto{\pgfqpoint{2.405228in}{2.165660in}}%
\pgfpathclose%
\pgfusepath{stroke,fill}%
\end{pgfscope}%
\begin{pgfscope}%
\pgfpathrectangle{\pgfqpoint{0.050000in}{0.050000in}}{\pgfqpoint{4.900000in}{4.900000in}}%
\pgfusepath{clip}%
\pgfsetbuttcap%
\pgfsetroundjoin%
\definecolor{currentfill}{rgb}{0.829450,0.829450,0.829450}%
\pgfsetfillcolor{currentfill}%
\pgfsetfillopacity{0.900000}%
\pgfsetlinewidth{0.501875pt}%
\definecolor{currentstroke}{rgb}{0.200000,0.200000,0.200000}%
\pgfsetstrokecolor{currentstroke}%
\pgfsetdash{}{0pt}%
\pgfpathmoveto{\pgfqpoint{1.893862in}{2.568371in}}%
\pgfpathlineto{\pgfqpoint{1.939221in}{2.525193in}}%
\pgfpathlineto{\pgfqpoint{1.969129in}{2.505756in}}%
\pgfpathlineto{\pgfqpoint{1.923764in}{2.548958in}}%
\pgfpathlineto{\pgfqpoint{1.893862in}{2.568371in}}%
\pgfpathclose%
\pgfusepath{stroke,fill}%
\end{pgfscope}%
\begin{pgfscope}%
\pgfpathrectangle{\pgfqpoint{0.050000in}{0.050000in}}{\pgfqpoint{4.900000in}{4.900000in}}%
\pgfusepath{clip}%
\pgfsetbuttcap%
\pgfsetroundjoin%
\definecolor{currentfill}{rgb}{0.994464,0.994464,0.994464}%
\pgfsetfillcolor{currentfill}%
\pgfsetfillopacity{0.900000}%
\pgfsetlinewidth{0.501875pt}%
\definecolor{currentstroke}{rgb}{0.200000,0.200000,0.200000}%
\pgfsetstrokecolor{currentstroke}%
\pgfsetdash{}{0pt}%
\pgfpathmoveto{\pgfqpoint{2.929621in}{1.930446in}}%
\pgfpathlineto{\pgfqpoint{2.973171in}{1.933423in}}%
\pgfpathlineto{\pgfqpoint{3.004653in}{1.915289in}}%
\pgfpathlineto{\pgfqpoint{2.961077in}{1.912538in}}%
\pgfpathlineto{\pgfqpoint{2.929621in}{1.930446in}}%
\pgfpathclose%
\pgfusepath{stroke,fill}%
\end{pgfscope}%
\begin{pgfscope}%
\pgfpathrectangle{\pgfqpoint{0.050000in}{0.050000in}}{\pgfqpoint{4.900000in}{4.900000in}}%
\pgfusepath{clip}%
\pgfsetbuttcap%
\pgfsetroundjoin%
\definecolor{currentfill}{rgb}{1.000000,1.000000,1.000000}%
\pgfsetfillcolor{currentfill}%
\pgfsetfillopacity{0.900000}%
\pgfsetlinewidth{0.501875pt}%
\definecolor{currentstroke}{rgb}{0.200000,0.200000,0.200000}%
\pgfsetstrokecolor{currentstroke}%
\pgfsetdash{}{0pt}%
\pgfpathmoveto{\pgfqpoint{3.467207in}{1.890898in}}%
\pgfpathlineto{\pgfqpoint{3.510257in}{1.903896in}}%
\pgfpathlineto{\pgfqpoint{3.542547in}{1.884401in}}%
\pgfpathlineto{\pgfqpoint{3.499468in}{1.871444in}}%
\pgfpathlineto{\pgfqpoint{3.467207in}{1.890898in}}%
\pgfpathclose%
\pgfusepath{stroke,fill}%
\end{pgfscope}%
\begin{pgfscope}%
\pgfpathrectangle{\pgfqpoint{0.050000in}{0.050000in}}{\pgfqpoint{4.900000in}{4.900000in}}%
\pgfusepath{clip}%
\pgfsetbuttcap%
\pgfsetroundjoin%
\definecolor{currentfill}{rgb}{0.996309,0.996309,0.996309}%
\pgfsetfillcolor{currentfill}%
\pgfsetfillopacity{0.900000}%
\pgfsetlinewidth{0.501875pt}%
\definecolor{currentstroke}{rgb}{0.200000,0.200000,0.200000}%
\pgfsetstrokecolor{currentstroke}%
\pgfsetdash{}{0pt}%
\pgfpathmoveto{\pgfqpoint{3.004653in}{1.915289in}}%
\pgfpathlineto{\pgfqpoint{3.048146in}{1.921016in}}%
\pgfpathlineto{\pgfqpoint{3.079751in}{1.902587in}}%
\pgfpathlineto{\pgfqpoint{3.036232in}{1.897087in}}%
\pgfpathlineto{\pgfqpoint{3.004653in}{1.915289in}}%
\pgfpathclose%
\pgfusepath{stroke,fill}%
\end{pgfscope}%
\begin{pgfscope}%
\pgfpathrectangle{\pgfqpoint{0.050000in}{0.050000in}}{\pgfqpoint{4.900000in}{4.900000in}}%
\pgfusepath{clip}%
\pgfsetbuttcap%
\pgfsetroundjoin%
\definecolor{currentfill}{rgb}{1.000000,1.000000,1.000000}%
\pgfsetfillcolor{currentfill}%
\pgfsetfillopacity{0.900000}%
\pgfsetlinewidth{0.501875pt}%
\definecolor{currentstroke}{rgb}{0.200000,0.200000,0.200000}%
\pgfsetstrokecolor{currentstroke}%
\pgfsetdash{}{0pt}%
\pgfpathmoveto{\pgfqpoint{3.273487in}{1.892917in}}%
\pgfpathlineto{\pgfqpoint{3.316735in}{1.904553in}}%
\pgfpathlineto{\pgfqpoint{3.348743in}{1.885332in}}%
\pgfpathlineto{\pgfqpoint{3.305467in}{1.873809in}}%
\pgfpathlineto{\pgfqpoint{3.273487in}{1.892917in}}%
\pgfpathclose%
\pgfusepath{stroke,fill}%
\end{pgfscope}%
\begin{pgfscope}%
\pgfpathrectangle{\pgfqpoint{0.050000in}{0.050000in}}{\pgfqpoint{4.900000in}{4.900000in}}%
\pgfusepath{clip}%
\pgfsetbuttcap%
\pgfsetroundjoin%
\definecolor{currentfill}{rgb}{0.937993,0.937993,0.937993}%
\pgfsetfillcolor{currentfill}%
\pgfsetfillopacity{0.900000}%
\pgfsetlinewidth{0.501875pt}%
\definecolor{currentstroke}{rgb}{0.200000,0.200000,0.200000}%
\pgfsetstrokecolor{currentstroke}%
\pgfsetdash{}{0pt}%
\pgfpathmoveto{\pgfqpoint{2.330197in}{2.221357in}}%
\pgfpathlineto{\pgfqpoint{2.374661in}{2.184020in}}%
\pgfpathlineto{\pgfqpoint{2.405228in}{2.165660in}}%
\pgfpathlineto{\pgfqpoint{2.360758in}{2.202459in}}%
\pgfpathlineto{\pgfqpoint{2.330197in}{2.221357in}}%
\pgfpathclose%
\pgfusepath{stroke,fill}%
\end{pgfscope}%
\begin{pgfscope}%
\pgfpathrectangle{\pgfqpoint{0.050000in}{0.050000in}}{\pgfqpoint{4.900000in}{4.900000in}}%
\pgfusepath{clip}%
\pgfsetbuttcap%
\pgfsetroundjoin%
\definecolor{currentfill}{rgb}{0.875740,0.875740,0.875740}%
\pgfsetfillcolor{currentfill}%
\pgfsetfillopacity{0.900000}%
\pgfsetlinewidth{0.501875pt}%
\definecolor{currentstroke}{rgb}{0.200000,0.200000,0.200000}%
\pgfsetstrokecolor{currentstroke}%
\pgfsetdash{}{0pt}%
\pgfpathmoveto{\pgfqpoint{2.074406in}{2.423669in}}%
\pgfpathlineto{\pgfqpoint{2.119413in}{2.380892in}}%
\pgfpathlineto{\pgfqpoint{2.149601in}{2.361355in}}%
\pgfpathlineto{\pgfqpoint{2.104590in}{2.404073in}}%
\pgfpathlineto{\pgfqpoint{2.074406in}{2.423669in}}%
\pgfpathclose%
\pgfusepath{stroke,fill}%
\end{pgfscope}%
\begin{pgfscope}%
\pgfpathrectangle{\pgfqpoint{0.050000in}{0.050000in}}{\pgfqpoint{4.900000in}{4.900000in}}%
\pgfusepath{clip}%
\pgfsetbuttcap%
\pgfsetroundjoin%
\definecolor{currentfill}{rgb}{0.805336,0.805336,0.805336}%
\pgfsetfillcolor{currentfill}%
\pgfsetfillopacity{0.900000}%
\pgfsetlinewidth{0.501875pt}%
\definecolor{currentstroke}{rgb}{0.200000,0.200000,0.200000}%
\pgfsetstrokecolor{currentstroke}%
\pgfsetdash{}{0pt}%
\pgfpathmoveto{\pgfqpoint{1.818512in}{2.631055in}}%
\pgfpathlineto{\pgfqpoint{1.864050in}{2.587726in}}%
\pgfpathlineto{\pgfqpoint{1.893862in}{2.568371in}}%
\pgfpathlineto{\pgfqpoint{1.848317in}{2.611726in}}%
\pgfpathlineto{\pgfqpoint{1.818512in}{2.631055in}}%
\pgfpathclose%
\pgfusepath{stroke,fill}%
\end{pgfscope}%
\begin{pgfscope}%
\pgfpathrectangle{\pgfqpoint{0.050000in}{0.050000in}}{\pgfqpoint{4.900000in}{4.900000in}}%
\pgfusepath{clip}%
\pgfsetbuttcap%
\pgfsetroundjoin%
\definecolor{currentfill}{rgb}{1.000000,1.000000,1.000000}%
\pgfsetfillcolor{currentfill}%
\pgfsetfillopacity{0.900000}%
\pgfsetlinewidth{0.501875pt}%
\definecolor{currentstroke}{rgb}{0.200000,0.200000,0.200000}%
\pgfsetstrokecolor{currentstroke}%
\pgfsetdash{}{0pt}%
\pgfpathmoveto{\pgfqpoint{3.542547in}{1.884401in}}%
\pgfpathlineto{\pgfqpoint{3.585538in}{1.897559in}}%
\pgfpathlineto{\pgfqpoint{3.617955in}{1.877982in}}%
\pgfpathlineto{\pgfqpoint{3.574935in}{1.864854in}}%
\pgfpathlineto{\pgfqpoint{3.542547in}{1.884401in}}%
\pgfpathclose%
\pgfusepath{stroke,fill}%
\end{pgfscope}%
\begin{pgfscope}%
\pgfpathrectangle{\pgfqpoint{0.050000in}{0.050000in}}{\pgfqpoint{4.900000in}{4.900000in}}%
\pgfusepath{clip}%
\pgfsetbuttcap%
\pgfsetroundjoin%
\definecolor{currentfill}{rgb}{0.998155,0.998155,0.998155}%
\pgfsetfillcolor{currentfill}%
\pgfsetfillopacity{0.900000}%
\pgfsetlinewidth{0.501875pt}%
\definecolor{currentstroke}{rgb}{0.200000,0.200000,0.200000}%
\pgfsetstrokecolor{currentstroke}%
\pgfsetdash{}{0pt}%
\pgfpathmoveto{\pgfqpoint{3.079751in}{1.902587in}}%
\pgfpathlineto{\pgfqpoint{3.123189in}{1.910433in}}%
\pgfpathlineto{\pgfqpoint{3.154918in}{1.891741in}}%
\pgfpathlineto{\pgfqpoint{3.111453in}{1.884101in}}%
\pgfpathlineto{\pgfqpoint{3.079751in}{1.902587in}}%
\pgfpathclose%
\pgfusepath{stroke,fill}%
\end{pgfscope}%
\begin{pgfscope}%
\pgfpathrectangle{\pgfqpoint{0.050000in}{0.050000in}}{\pgfqpoint{4.900000in}{4.900000in}}%
\pgfusepath{clip}%
\pgfsetbuttcap%
\pgfsetroundjoin%
\definecolor{currentfill}{rgb}{1.000000,1.000000,1.000000}%
\pgfsetfillcolor{currentfill}%
\pgfsetfillopacity{0.900000}%
\pgfsetlinewidth{0.501875pt}%
\definecolor{currentstroke}{rgb}{0.200000,0.200000,0.200000}%
\pgfsetstrokecolor{currentstroke}%
\pgfsetdash{}{0pt}%
\pgfpathmoveto{\pgfqpoint{3.348743in}{1.885332in}}%
\pgfpathlineto{\pgfqpoint{3.391936in}{1.897575in}}%
\pgfpathlineto{\pgfqpoint{3.424070in}{1.878228in}}%
\pgfpathlineto{\pgfqpoint{3.380849in}{1.866068in}}%
\pgfpathlineto{\pgfqpoint{3.348743in}{1.885332in}}%
\pgfpathclose%
\pgfusepath{stroke,fill}%
\end{pgfscope}%
\begin{pgfscope}%
\pgfpathrectangle{\pgfqpoint{0.050000in}{0.050000in}}{\pgfqpoint{4.900000in}{4.900000in}}%
\pgfusepath{clip}%
\pgfsetbuttcap%
\pgfsetroundjoin%
\definecolor{currentfill}{rgb}{0.918185,0.918185,0.918185}%
\pgfsetfillcolor{currentfill}%
\pgfsetfillopacity{0.900000}%
\pgfsetlinewidth{0.501875pt}%
\definecolor{currentstroke}{rgb}{0.200000,0.200000,0.200000}%
\pgfsetstrokecolor{currentstroke}%
\pgfsetdash{}{0pt}%
\pgfpathmoveto{\pgfqpoint{2.255082in}{2.280538in}}%
\pgfpathlineto{\pgfqpoint{2.299728in}{2.240263in}}%
\pgfpathlineto{\pgfqpoint{2.330197in}{2.221357in}}%
\pgfpathlineto{\pgfqpoint{2.285547in}{2.261224in}}%
\pgfpathlineto{\pgfqpoint{2.255082in}{2.280538in}}%
\pgfpathclose%
\pgfusepath{stroke,fill}%
\end{pgfscope}%
\begin{pgfscope}%
\pgfpathrectangle{\pgfqpoint{0.050000in}{0.050000in}}{\pgfqpoint{4.900000in}{4.900000in}}%
\pgfusepath{clip}%
\pgfsetbuttcap%
\pgfsetroundjoin%
\definecolor{currentfill}{rgb}{0.855932,0.855932,0.855932}%
\pgfsetfillcolor{currentfill}%
\pgfsetfillopacity{0.900000}%
\pgfsetlinewidth{0.501875pt}%
\definecolor{currentstroke}{rgb}{0.200000,0.200000,0.200000}%
\pgfsetstrokecolor{currentstroke}%
\pgfsetdash{}{0pt}%
\pgfpathmoveto{\pgfqpoint{1.999127in}{2.486262in}}%
\pgfpathlineto{\pgfqpoint{2.044313in}{2.443225in}}%
\pgfpathlineto{\pgfqpoint{2.074406in}{2.423669in}}%
\pgfpathlineto{\pgfqpoint{2.029214in}{2.466713in}}%
\pgfpathlineto{\pgfqpoint{1.999127in}{2.486262in}}%
\pgfpathclose%
\pgfusepath{stroke,fill}%
\end{pgfscope}%
\begin{pgfscope}%
\pgfpathrectangle{\pgfqpoint{0.050000in}{0.050000in}}{\pgfqpoint{4.900000in}{4.900000in}}%
\pgfusepath{clip}%
\pgfsetbuttcap%
\pgfsetroundjoin%
\definecolor{currentfill}{rgb}{0.976009,0.976009,0.976009}%
\pgfsetfillcolor{currentfill}%
\pgfsetfillopacity{0.900000}%
\pgfsetlinewidth{0.501875pt}%
\definecolor{currentstroke}{rgb}{0.200000,0.200000,0.200000}%
\pgfsetstrokecolor{currentstroke}%
\pgfsetdash{}{0pt}%
\pgfpathmoveto{\pgfqpoint{2.660949in}{2.013626in}}%
\pgfpathlineto{\pgfqpoint{2.704837in}{1.998213in}}%
\pgfpathlineto{\pgfqpoint{2.735934in}{1.980937in}}%
\pgfpathlineto{\pgfqpoint{2.692027in}{1.996204in}}%
\pgfpathlineto{\pgfqpoint{2.660949in}{2.013626in}}%
\pgfpathclose%
\pgfusepath{stroke,fill}%
\end{pgfscope}%
\begin{pgfscope}%
\pgfpathrectangle{\pgfqpoint{0.050000in}{0.050000in}}{\pgfqpoint{4.900000in}{4.900000in}}%
\pgfusepath{clip}%
\pgfsetbuttcap%
\pgfsetroundjoin%
\definecolor{currentfill}{rgb}{0.981546,0.981546,0.981546}%
\pgfsetfillcolor{currentfill}%
\pgfsetfillopacity{0.900000}%
\pgfsetlinewidth{0.501875pt}%
\definecolor{currentstroke}{rgb}{0.200000,0.200000,0.200000}%
\pgfsetstrokecolor{currentstroke}%
\pgfsetdash{}{0pt}%
\pgfpathmoveto{\pgfqpoint{2.735934in}{1.980937in}}%
\pgfpathlineto{\pgfqpoint{2.779724in}{1.971083in}}%
\pgfpathlineto{\pgfqpoint{2.810940in}{1.953682in}}%
\pgfpathlineto{\pgfqpoint{2.767128in}{1.963536in}}%
\pgfpathlineto{\pgfqpoint{2.735934in}{1.980937in}}%
\pgfpathclose%
\pgfusepath{stroke,fill}%
\end{pgfscope}%
\begin{pgfscope}%
\pgfpathrectangle{\pgfqpoint{0.050000in}{0.050000in}}{\pgfqpoint{4.900000in}{4.900000in}}%
\pgfusepath{clip}%
\pgfsetbuttcap%
\pgfsetroundjoin%
\definecolor{currentfill}{rgb}{1.000000,1.000000,1.000000}%
\pgfsetfillcolor{currentfill}%
\pgfsetfillopacity{0.900000}%
\pgfsetlinewidth{0.501875pt}%
\definecolor{currentstroke}{rgb}{0.200000,0.200000,0.200000}%
\pgfsetstrokecolor{currentstroke}%
\pgfsetdash{}{0pt}%
\pgfpathmoveto{\pgfqpoint{3.617955in}{1.877982in}}%
\pgfpathlineto{\pgfqpoint{3.660886in}{1.891203in}}%
\pgfpathlineto{\pgfqpoint{3.693430in}{1.871544in}}%
\pgfpathlineto{\pgfqpoint{3.650470in}{1.858349in}}%
\pgfpathlineto{\pgfqpoint{3.617955in}{1.877982in}}%
\pgfpathclose%
\pgfusepath{stroke,fill}%
\end{pgfscope}%
\begin{pgfscope}%
\pgfpathrectangle{\pgfqpoint{0.050000in}{0.050000in}}{\pgfqpoint{4.900000in}{4.900000in}}%
\pgfusepath{clip}%
\pgfsetbuttcap%
\pgfsetroundjoin%
\definecolor{currentfill}{rgb}{0.968627,0.968627,0.968627}%
\pgfsetfillcolor{currentfill}%
\pgfsetfillopacity{0.900000}%
\pgfsetlinewidth{0.501875pt}%
\definecolor{currentstroke}{rgb}{0.200000,0.200000,0.200000}%
\pgfsetstrokecolor{currentstroke}%
\pgfsetdash{}{0pt}%
\pgfpathmoveto{\pgfqpoint{2.585961in}{2.052243in}}%
\pgfpathlineto{\pgfqpoint{2.629967in}{2.030925in}}%
\pgfpathlineto{\pgfqpoint{2.660949in}{2.013626in}}%
\pgfpathlineto{\pgfqpoint{2.616927in}{2.034636in}}%
\pgfpathlineto{\pgfqpoint{2.585961in}{2.052243in}}%
\pgfpathclose%
\pgfusepath{stroke,fill}%
\end{pgfscope}%
\begin{pgfscope}%
\pgfpathrectangle{\pgfqpoint{0.050000in}{0.050000in}}{\pgfqpoint{4.900000in}{4.900000in}}%
\pgfusepath{clip}%
\pgfsetbuttcap%
\pgfsetroundjoin%
\definecolor{currentfill}{rgb}{0.998155,0.998155,0.998155}%
\pgfsetfillcolor{currentfill}%
\pgfsetfillopacity{0.900000}%
\pgfsetlinewidth{0.501875pt}%
\definecolor{currentstroke}{rgb}{0.200000,0.200000,0.200000}%
\pgfsetstrokecolor{currentstroke}%
\pgfsetdash{}{0pt}%
\pgfpathmoveto{\pgfqpoint{3.154918in}{1.891741in}}%
\pgfpathlineto{\pgfqpoint{3.198302in}{1.901191in}}%
\pgfpathlineto{\pgfqpoint{3.230156in}{1.882272in}}%
\pgfpathlineto{\pgfqpoint{3.186744in}{1.872999in}}%
\pgfpathlineto{\pgfqpoint{3.154918in}{1.891741in}}%
\pgfpathclose%
\pgfusepath{stroke,fill}%
\end{pgfscope}%
\begin{pgfscope}%
\pgfpathrectangle{\pgfqpoint{0.050000in}{0.050000in}}{\pgfqpoint{4.900000in}{4.900000in}}%
\pgfusepath{clip}%
\pgfsetbuttcap%
\pgfsetroundjoin%
\definecolor{currentfill}{rgb}{0.987082,0.987082,0.987082}%
\pgfsetfillcolor{currentfill}%
\pgfsetfillopacity{0.900000}%
\pgfsetlinewidth{0.501875pt}%
\definecolor{currentstroke}{rgb}{0.200000,0.200000,0.200000}%
\pgfsetstrokecolor{currentstroke}%
\pgfsetdash{}{0pt}%
\pgfpathmoveto{\pgfqpoint{2.810940in}{1.953682in}}%
\pgfpathlineto{\pgfqpoint{2.854648in}{1.948774in}}%
\pgfpathlineto{\pgfqpoint{2.885983in}{1.931146in}}%
\pgfpathlineto{\pgfqpoint{2.842252in}{1.936164in}}%
\pgfpathlineto{\pgfqpoint{2.810940in}{1.953682in}}%
\pgfpathclose%
\pgfusepath{stroke,fill}%
\end{pgfscope}%
\begin{pgfscope}%
\pgfpathrectangle{\pgfqpoint{0.050000in}{0.050000in}}{\pgfqpoint{4.900000in}{4.900000in}}%
\pgfusepath{clip}%
\pgfsetbuttcap%
\pgfsetroundjoin%
\definecolor{currentfill}{rgb}{0.777778,0.777778,0.777778}%
\pgfsetfillcolor{currentfill}%
\pgfsetfillopacity{0.900000}%
\pgfsetlinewidth{0.501875pt}%
\definecolor{currentstroke}{rgb}{0.200000,0.200000,0.200000}%
\pgfsetstrokecolor{currentstroke}%
\pgfsetdash{}{0pt}%
\pgfpathmoveto{\pgfqpoint{1.743078in}{2.693810in}}%
\pgfpathlineto{\pgfqpoint{1.788796in}{2.650328in}}%
\pgfpathlineto{\pgfqpoint{1.818512in}{2.631055in}}%
\pgfpathlineto{\pgfqpoint{1.772787in}{2.674563in}}%
\pgfpathlineto{\pgfqpoint{1.743078in}{2.693810in}}%
\pgfpathclose%
\pgfusepath{stroke,fill}%
\end{pgfscope}%
\begin{pgfscope}%
\pgfpathrectangle{\pgfqpoint{0.050000in}{0.050000in}}{\pgfqpoint{4.900000in}{4.900000in}}%
\pgfusepath{clip}%
\pgfsetbuttcap%
\pgfsetroundjoin%
\definecolor{currentfill}{rgb}{1.000000,1.000000,1.000000}%
\pgfsetfillcolor{currentfill}%
\pgfsetfillopacity{0.900000}%
\pgfsetlinewidth{0.501875pt}%
\definecolor{currentstroke}{rgb}{0.200000,0.200000,0.200000}%
\pgfsetstrokecolor{currentstroke}%
\pgfsetdash{}{0pt}%
\pgfpathmoveto{\pgfqpoint{3.424070in}{1.878228in}}%
\pgfpathlineto{\pgfqpoint{3.467207in}{1.890898in}}%
\pgfpathlineto{\pgfqpoint{3.499468in}{1.871444in}}%
\pgfpathlineto{\pgfqpoint{3.456302in}{1.858834in}}%
\pgfpathlineto{\pgfqpoint{3.424070in}{1.878228in}}%
\pgfpathclose%
\pgfusepath{stroke,fill}%
\end{pgfscope}%
\begin{pgfscope}%
\pgfpathrectangle{\pgfqpoint{0.050000in}{0.050000in}}{\pgfqpoint{4.900000in}{4.900000in}}%
\pgfusepath{clip}%
\pgfsetbuttcap%
\pgfsetroundjoin%
\definecolor{currentfill}{rgb}{0.961246,0.961246,0.961246}%
\pgfsetfillcolor{currentfill}%
\pgfsetfillopacity{0.900000}%
\pgfsetlinewidth{0.501875pt}%
\definecolor{currentstroke}{rgb}{0.200000,0.200000,0.200000}%
\pgfsetstrokecolor{currentstroke}%
\pgfsetdash{}{0pt}%
\pgfpathmoveto{\pgfqpoint{2.510949in}{2.096915in}}%
\pgfpathlineto{\pgfqpoint{2.555091in}{2.069747in}}%
\pgfpathlineto{\pgfqpoint{2.585961in}{2.052243in}}%
\pgfpathlineto{\pgfqpoint{2.541806in}{2.078961in}}%
\pgfpathlineto{\pgfqpoint{2.510949in}{2.096915in}}%
\pgfpathclose%
\pgfusepath{stroke,fill}%
\end{pgfscope}%
\begin{pgfscope}%
\pgfpathrectangle{\pgfqpoint{0.050000in}{0.050000in}}{\pgfqpoint{4.900000in}{4.900000in}}%
\pgfusepath{clip}%
\pgfsetbuttcap%
\pgfsetroundjoin%
\definecolor{currentfill}{rgb}{0.990773,0.990773,0.990773}%
\pgfsetfillcolor{currentfill}%
\pgfsetfillopacity{0.900000}%
\pgfsetlinewidth{0.501875pt}%
\definecolor{currentstroke}{rgb}{0.200000,0.200000,0.200000}%
\pgfsetstrokecolor{currentstroke}%
\pgfsetdash{}{0pt}%
\pgfpathmoveto{\pgfqpoint{2.885983in}{1.931146in}}%
\pgfpathlineto{\pgfqpoint{2.929621in}{1.930446in}}%
\pgfpathlineto{\pgfqpoint{2.961077in}{1.912538in}}%
\pgfpathlineto{\pgfqpoint{2.917415in}{1.913417in}}%
\pgfpathlineto{\pgfqpoint{2.885983in}{1.931146in}}%
\pgfpathclose%
\pgfusepath{stroke,fill}%
\end{pgfscope}%
\begin{pgfscope}%
\pgfpathrectangle{\pgfqpoint{0.050000in}{0.050000in}}{\pgfqpoint{4.900000in}{4.900000in}}%
\pgfusepath{clip}%
\pgfsetbuttcap%
\pgfsetroundjoin%
\definecolor{currentfill}{rgb}{0.898378,0.898378,0.898378}%
\pgfsetfillcolor{currentfill}%
\pgfsetfillopacity{0.900000}%
\pgfsetlinewidth{0.501875pt}%
\definecolor{currentstroke}{rgb}{0.200000,0.200000,0.200000}%
\pgfsetstrokecolor{currentstroke}%
\pgfsetdash{}{0pt}%
\pgfpathmoveto{\pgfqpoint{2.179880in}{2.341807in}}%
\pgfpathlineto{\pgfqpoint{2.224709in}{2.299864in}}%
\pgfpathlineto{\pgfqpoint{2.255082in}{2.280538in}}%
\pgfpathlineto{\pgfqpoint{2.210249in}{2.322250in}}%
\pgfpathlineto{\pgfqpoint{2.179880in}{2.341807in}}%
\pgfpathclose%
\pgfusepath{stroke,fill}%
\end{pgfscope}%
\begin{pgfscope}%
\pgfpathrectangle{\pgfqpoint{0.050000in}{0.050000in}}{\pgfqpoint{4.900000in}{4.900000in}}%
\pgfusepath{clip}%
\pgfsetbuttcap%
\pgfsetroundjoin%
\definecolor{currentfill}{rgb}{0.998155,0.998155,0.998155}%
\pgfsetfillcolor{currentfill}%
\pgfsetfillopacity{0.900000}%
\pgfsetlinewidth{0.501875pt}%
\definecolor{currentstroke}{rgb}{0.200000,0.200000,0.200000}%
\pgfsetstrokecolor{currentstroke}%
\pgfsetdash{}{0pt}%
\pgfpathmoveto{\pgfqpoint{3.230156in}{1.882272in}}%
\pgfpathlineto{\pgfqpoint{3.273487in}{1.892917in}}%
\pgfpathlineto{\pgfqpoint{3.305467in}{1.873809in}}%
\pgfpathlineto{\pgfqpoint{3.262108in}{1.863308in}}%
\pgfpathlineto{\pgfqpoint{3.230156in}{1.882272in}}%
\pgfpathclose%
\pgfusepath{stroke,fill}%
\end{pgfscope}%
\begin{pgfscope}%
\pgfpathrectangle{\pgfqpoint{0.050000in}{0.050000in}}{\pgfqpoint{4.900000in}{4.900000in}}%
\pgfusepath{clip}%
\pgfsetbuttcap%
\pgfsetroundjoin%
\definecolor{currentfill}{rgb}{0.829450,0.829450,0.829450}%
\pgfsetfillcolor{currentfill}%
\pgfsetfillopacity{0.900000}%
\pgfsetlinewidth{0.501875pt}%
\definecolor{currentstroke}{rgb}{0.200000,0.200000,0.200000}%
\pgfsetstrokecolor{currentstroke}%
\pgfsetdash{}{0pt}%
\pgfpathmoveto{\pgfqpoint{1.923764in}{2.548958in}}%
\pgfpathlineto{\pgfqpoint{1.969129in}{2.505756in}}%
\pgfpathlineto{\pgfqpoint{1.999127in}{2.486262in}}%
\pgfpathlineto{\pgfqpoint{1.953755in}{2.529487in}}%
\pgfpathlineto{\pgfqpoint{1.923764in}{2.548958in}}%
\pgfpathclose%
\pgfusepath{stroke,fill}%
\end{pgfscope}%
\begin{pgfscope}%
\pgfpathrectangle{\pgfqpoint{0.050000in}{0.050000in}}{\pgfqpoint{4.900000in}{4.900000in}}%
\pgfusepath{clip}%
\pgfsetbuttcap%
\pgfsetroundjoin%
\definecolor{currentfill}{rgb}{0.950173,0.950173,0.950173}%
\pgfsetfillcolor{currentfill}%
\pgfsetfillopacity{0.900000}%
\pgfsetlinewidth{0.501875pt}%
\definecolor{currentstroke}{rgb}{0.200000,0.200000,0.200000}%
\pgfsetstrokecolor{currentstroke}%
\pgfsetdash{}{0pt}%
\pgfpathmoveto{\pgfqpoint{2.435887in}{2.147277in}}%
\pgfpathlineto{\pgfqpoint{2.480186in}{2.114799in}}%
\pgfpathlineto{\pgfqpoint{2.510949in}{2.096915in}}%
\pgfpathlineto{\pgfqpoint{2.466641in}{2.128863in}}%
\pgfpathlineto{\pgfqpoint{2.435887in}{2.147277in}}%
\pgfpathclose%
\pgfusepath{stroke,fill}%
\end{pgfscope}%
\begin{pgfscope}%
\pgfpathrectangle{\pgfqpoint{0.050000in}{0.050000in}}{\pgfqpoint{4.900000in}{4.900000in}}%
\pgfusepath{clip}%
\pgfsetbuttcap%
\pgfsetroundjoin%
\definecolor{currentfill}{rgb}{0.992618,0.992618,0.992618}%
\pgfsetfillcolor{currentfill}%
\pgfsetfillopacity{0.900000}%
\pgfsetlinewidth{0.501875pt}%
\definecolor{currentstroke}{rgb}{0.200000,0.200000,0.200000}%
\pgfsetstrokecolor{currentstroke}%
\pgfsetdash{}{0pt}%
\pgfpathmoveto{\pgfqpoint{2.961077in}{1.912538in}}%
\pgfpathlineto{\pgfqpoint{3.004653in}{1.915289in}}%
\pgfpathlineto{\pgfqpoint{3.036232in}{1.897087in}}%
\pgfpathlineto{\pgfqpoint{2.992631in}{1.894545in}}%
\pgfpathlineto{\pgfqpoint{2.961077in}{1.912538in}}%
\pgfpathclose%
\pgfusepath{stroke,fill}%
\end{pgfscope}%
\begin{pgfscope}%
\pgfpathrectangle{\pgfqpoint{0.050000in}{0.050000in}}{\pgfqpoint{4.900000in}{4.900000in}}%
\pgfusepath{clip}%
\pgfsetbuttcap%
\pgfsetroundjoin%
\definecolor{currentfill}{rgb}{1.000000,1.000000,1.000000}%
\pgfsetfillcolor{currentfill}%
\pgfsetfillopacity{0.900000}%
\pgfsetlinewidth{0.501875pt}%
\definecolor{currentstroke}{rgb}{0.200000,0.200000,0.200000}%
\pgfsetstrokecolor{currentstroke}%
\pgfsetdash{}{0pt}%
\pgfpathmoveto{\pgfqpoint{3.499468in}{1.871444in}}%
\pgfpathlineto{\pgfqpoint{3.542547in}{1.884401in}}%
\pgfpathlineto{\pgfqpoint{3.574935in}{1.864854in}}%
\pgfpathlineto{\pgfqpoint{3.531827in}{1.851942in}}%
\pgfpathlineto{\pgfqpoint{3.499468in}{1.871444in}}%
\pgfpathclose%
\pgfusepath{stroke,fill}%
\end{pgfscope}%
\begin{pgfscope}%
\pgfpathrectangle{\pgfqpoint{0.050000in}{0.050000in}}{\pgfqpoint{4.900000in}{4.900000in}}%
\pgfusepath{clip}%
\pgfsetbuttcap%
\pgfsetroundjoin%
\definecolor{currentfill}{rgb}{0.996309,0.996309,0.996309}%
\pgfsetfillcolor{currentfill}%
\pgfsetfillopacity{0.900000}%
\pgfsetlinewidth{0.501875pt}%
\definecolor{currentstroke}{rgb}{0.200000,0.200000,0.200000}%
\pgfsetstrokecolor{currentstroke}%
\pgfsetdash{}{0pt}%
\pgfpathmoveto{\pgfqpoint{3.036232in}{1.897087in}}%
\pgfpathlineto{\pgfqpoint{3.079751in}{1.902587in}}%
\pgfpathlineto{\pgfqpoint{3.111453in}{1.884101in}}%
\pgfpathlineto{\pgfqpoint{3.067908in}{1.878814in}}%
\pgfpathlineto{\pgfqpoint{3.036232in}{1.897087in}}%
\pgfpathclose%
\pgfusepath{stroke,fill}%
\end{pgfscope}%
\begin{pgfscope}%
\pgfpathrectangle{\pgfqpoint{0.050000in}{0.050000in}}{\pgfqpoint{4.900000in}{4.900000in}}%
\pgfusepath{clip}%
\pgfsetbuttcap%
\pgfsetroundjoin%
\definecolor{currentfill}{rgb}{1.000000,1.000000,1.000000}%
\pgfsetfillcolor{currentfill}%
\pgfsetfillopacity{0.900000}%
\pgfsetlinewidth{0.501875pt}%
\definecolor{currentstroke}{rgb}{0.200000,0.200000,0.200000}%
\pgfsetstrokecolor{currentstroke}%
\pgfsetdash{}{0pt}%
\pgfpathmoveto{\pgfqpoint{3.305467in}{1.873809in}}%
\pgfpathlineto{\pgfqpoint{3.348743in}{1.885332in}}%
\pgfpathlineto{\pgfqpoint{3.380849in}{1.866068in}}%
\pgfpathlineto{\pgfqpoint{3.337543in}{1.854657in}}%
\pgfpathlineto{\pgfqpoint{3.305467in}{1.873809in}}%
\pgfpathclose%
\pgfusepath{stroke,fill}%
\end{pgfscope}%
\begin{pgfscope}%
\pgfpathrectangle{\pgfqpoint{0.050000in}{0.050000in}}{\pgfqpoint{4.900000in}{4.900000in}}%
\pgfusepath{clip}%
\pgfsetbuttcap%
\pgfsetroundjoin%
\definecolor{currentfill}{rgb}{0.875740,0.875740,0.875740}%
\pgfsetfillcolor{currentfill}%
\pgfsetfillopacity{0.900000}%
\pgfsetlinewidth{0.501875pt}%
\definecolor{currentstroke}{rgb}{0.200000,0.200000,0.200000}%
\pgfsetstrokecolor{currentstroke}%
\pgfsetdash{}{0pt}%
\pgfpathmoveto{\pgfqpoint{2.104590in}{2.404073in}}%
\pgfpathlineto{\pgfqpoint{2.149601in}{2.361355in}}%
\pgfpathlineto{\pgfqpoint{2.179880in}{2.341807in}}%
\pgfpathlineto{\pgfqpoint{2.134863in}{2.384440in}}%
\pgfpathlineto{\pgfqpoint{2.104590in}{2.404073in}}%
\pgfpathclose%
\pgfusepath{stroke,fill}%
\end{pgfscope}%
\begin{pgfscope}%
\pgfpathrectangle{\pgfqpoint{0.050000in}{0.050000in}}{\pgfqpoint{4.900000in}{4.900000in}}%
\pgfusepath{clip}%
\pgfsetbuttcap%
\pgfsetroundjoin%
\definecolor{currentfill}{rgb}{0.937993,0.937993,0.937993}%
\pgfsetfillcolor{currentfill}%
\pgfsetfillopacity{0.900000}%
\pgfsetlinewidth{0.501875pt}%
\definecolor{currentstroke}{rgb}{0.200000,0.200000,0.200000}%
\pgfsetstrokecolor{currentstroke}%
\pgfsetdash{}{0pt}%
\pgfpathmoveto{\pgfqpoint{2.360758in}{2.202459in}}%
\pgfpathlineto{\pgfqpoint{2.405228in}{2.165660in}}%
\pgfpathlineto{\pgfqpoint{2.435887in}{2.147277in}}%
\pgfpathlineto{\pgfqpoint{2.391411in}{2.183560in}}%
\pgfpathlineto{\pgfqpoint{2.360758in}{2.202459in}}%
\pgfpathclose%
\pgfusepath{stroke,fill}%
\end{pgfscope}%
\begin{pgfscope}%
\pgfpathrectangle{\pgfqpoint{0.050000in}{0.050000in}}{\pgfqpoint{4.900000in}{4.900000in}}%
\pgfusepath{clip}%
\pgfsetbuttcap%
\pgfsetroundjoin%
\definecolor{currentfill}{rgb}{0.805336,0.805336,0.805336}%
\pgfsetfillcolor{currentfill}%
\pgfsetfillopacity{0.900000}%
\pgfsetlinewidth{0.501875pt}%
\definecolor{currentstroke}{rgb}{0.200000,0.200000,0.200000}%
\pgfsetstrokecolor{currentstroke}%
\pgfsetdash{}{0pt}%
\pgfpathmoveto{\pgfqpoint{1.848317in}{2.611726in}}%
\pgfpathlineto{\pgfqpoint{1.893862in}{2.568371in}}%
\pgfpathlineto{\pgfqpoint{1.923764in}{2.548958in}}%
\pgfpathlineto{\pgfqpoint{1.878212in}{2.592338in}}%
\pgfpathlineto{\pgfqpoint{1.848317in}{2.611726in}}%
\pgfpathclose%
\pgfusepath{stroke,fill}%
\end{pgfscope}%
\begin{pgfscope}%
\pgfpathrectangle{\pgfqpoint{0.050000in}{0.050000in}}{\pgfqpoint{4.900000in}{4.900000in}}%
\pgfusepath{clip}%
\pgfsetbuttcap%
\pgfsetroundjoin%
\definecolor{currentfill}{rgb}{1.000000,1.000000,1.000000}%
\pgfsetfillcolor{currentfill}%
\pgfsetfillopacity{0.900000}%
\pgfsetlinewidth{0.501875pt}%
\definecolor{currentstroke}{rgb}{0.200000,0.200000,0.200000}%
\pgfsetstrokecolor{currentstroke}%
\pgfsetdash{}{0pt}%
\pgfpathmoveto{\pgfqpoint{3.574935in}{1.864854in}}%
\pgfpathlineto{\pgfqpoint{3.617955in}{1.877982in}}%
\pgfpathlineto{\pgfqpoint{3.650470in}{1.858349in}}%
\pgfpathlineto{\pgfqpoint{3.607421in}{1.845256in}}%
\pgfpathlineto{\pgfqpoint{3.574935in}{1.864854in}}%
\pgfpathclose%
\pgfusepath{stroke,fill}%
\end{pgfscope}%
\begin{pgfscope}%
\pgfpathrectangle{\pgfqpoint{0.050000in}{0.050000in}}{\pgfqpoint{4.900000in}{4.900000in}}%
\pgfusepath{clip}%
\pgfsetbuttcap%
\pgfsetroundjoin%
\definecolor{currentfill}{rgb}{0.996309,0.996309,0.996309}%
\pgfsetfillcolor{currentfill}%
\pgfsetfillopacity{0.900000}%
\pgfsetlinewidth{0.501875pt}%
\definecolor{currentstroke}{rgb}{0.200000,0.200000,0.200000}%
\pgfsetstrokecolor{currentstroke}%
\pgfsetdash{}{0pt}%
\pgfpathmoveto{\pgfqpoint{3.111453in}{1.884101in}}%
\pgfpathlineto{\pgfqpoint{3.154918in}{1.891741in}}%
\pgfpathlineto{\pgfqpoint{3.186744in}{1.872999in}}%
\pgfpathlineto{\pgfqpoint{3.143253in}{1.865556in}}%
\pgfpathlineto{\pgfqpoint{3.111453in}{1.884101in}}%
\pgfpathclose%
\pgfusepath{stroke,fill}%
\end{pgfscope}%
\begin{pgfscope}%
\pgfpathrectangle{\pgfqpoint{0.050000in}{0.050000in}}{\pgfqpoint{4.900000in}{4.900000in}}%
\pgfusepath{clip}%
\pgfsetbuttcap%
\pgfsetroundjoin%
\definecolor{currentfill}{rgb}{1.000000,1.000000,1.000000}%
\pgfsetfillcolor{currentfill}%
\pgfsetfillopacity{0.900000}%
\pgfsetlinewidth{0.501875pt}%
\definecolor{currentstroke}{rgb}{0.200000,0.200000,0.200000}%
\pgfsetstrokecolor{currentstroke}%
\pgfsetdash{}{0pt}%
\pgfpathmoveto{\pgfqpoint{3.380849in}{1.866068in}}%
\pgfpathlineto{\pgfqpoint{3.424070in}{1.878228in}}%
\pgfpathlineto{\pgfqpoint{3.456302in}{1.858834in}}%
\pgfpathlineto{\pgfqpoint{3.413052in}{1.846760in}}%
\pgfpathlineto{\pgfqpoint{3.380849in}{1.866068in}}%
\pgfpathclose%
\pgfusepath{stroke,fill}%
\end{pgfscope}%
\begin{pgfscope}%
\pgfpathrectangle{\pgfqpoint{0.050000in}{0.050000in}}{\pgfqpoint{4.900000in}{4.900000in}}%
\pgfusepath{clip}%
\pgfsetbuttcap%
\pgfsetroundjoin%
\definecolor{currentfill}{rgb}{0.918185,0.918185,0.918185}%
\pgfsetfillcolor{currentfill}%
\pgfsetfillopacity{0.900000}%
\pgfsetlinewidth{0.501875pt}%
\definecolor{currentstroke}{rgb}{0.200000,0.200000,0.200000}%
\pgfsetstrokecolor{currentstroke}%
\pgfsetdash{}{0pt}%
\pgfpathmoveto{\pgfqpoint{2.285547in}{2.261224in}}%
\pgfpathlineto{\pgfqpoint{2.330197in}{2.221357in}}%
\pgfpathlineto{\pgfqpoint{2.360758in}{2.202459in}}%
\pgfpathlineto{\pgfqpoint{2.316103in}{2.241917in}}%
\pgfpathlineto{\pgfqpoint{2.285547in}{2.261224in}}%
\pgfpathclose%
\pgfusepath{stroke,fill}%
\end{pgfscope}%
\begin{pgfscope}%
\pgfpathrectangle{\pgfqpoint{0.050000in}{0.050000in}}{\pgfqpoint{4.900000in}{4.900000in}}%
\pgfusepath{clip}%
\pgfsetbuttcap%
\pgfsetroundjoin%
\definecolor{currentfill}{rgb}{0.855932,0.855932,0.855932}%
\pgfsetfillcolor{currentfill}%
\pgfsetfillopacity{0.900000}%
\pgfsetlinewidth{0.501875pt}%
\definecolor{currentstroke}{rgb}{0.200000,0.200000,0.200000}%
\pgfsetstrokecolor{currentstroke}%
\pgfsetdash{}{0pt}%
\pgfpathmoveto{\pgfqpoint{2.029214in}{2.466713in}}%
\pgfpathlineto{\pgfqpoint{2.074406in}{2.423669in}}%
\pgfpathlineto{\pgfqpoint{2.104590in}{2.404073in}}%
\pgfpathlineto{\pgfqpoint{2.059392in}{2.447112in}}%
\pgfpathlineto{\pgfqpoint{2.029214in}{2.466713in}}%
\pgfpathclose%
\pgfusepath{stroke,fill}%
\end{pgfscope}%
\begin{pgfscope}%
\pgfpathrectangle{\pgfqpoint{0.050000in}{0.050000in}}{\pgfqpoint{4.900000in}{4.900000in}}%
\pgfusepath{clip}%
\pgfsetbuttcap%
\pgfsetroundjoin%
\definecolor{currentfill}{rgb}{0.976009,0.976009,0.976009}%
\pgfsetfillcolor{currentfill}%
\pgfsetfillopacity{0.900000}%
\pgfsetlinewidth{0.501875pt}%
\definecolor{currentstroke}{rgb}{0.200000,0.200000,0.200000}%
\pgfsetstrokecolor{currentstroke}%
\pgfsetdash{}{0pt}%
\pgfpathmoveto{\pgfqpoint{2.692027in}{1.996204in}}%
\pgfpathlineto{\pgfqpoint{2.735934in}{1.980937in}}%
\pgfpathlineto{\pgfqpoint{2.767128in}{1.963536in}}%
\pgfpathlineto{\pgfqpoint{2.723201in}{1.978664in}}%
\pgfpathlineto{\pgfqpoint{2.692027in}{1.996204in}}%
\pgfpathclose%
\pgfusepath{stroke,fill}%
\end{pgfscope}%
\begin{pgfscope}%
\pgfpathrectangle{\pgfqpoint{0.050000in}{0.050000in}}{\pgfqpoint{4.900000in}{4.900000in}}%
\pgfusepath{clip}%
\pgfsetbuttcap%
\pgfsetroundjoin%
\definecolor{currentfill}{rgb}{1.000000,1.000000,1.000000}%
\pgfsetfillcolor{currentfill}%
\pgfsetfillopacity{0.900000}%
\pgfsetlinewidth{0.501875pt}%
\definecolor{currentstroke}{rgb}{0.200000,0.200000,0.200000}%
\pgfsetstrokecolor{currentstroke}%
\pgfsetdash{}{0pt}%
\pgfpathmoveto{\pgfqpoint{3.650470in}{1.858349in}}%
\pgfpathlineto{\pgfqpoint{3.693430in}{1.871544in}}%
\pgfpathlineto{\pgfqpoint{3.726072in}{1.851825in}}%
\pgfpathlineto{\pgfqpoint{3.683084in}{1.838661in}}%
\pgfpathlineto{\pgfqpoint{3.650470in}{1.858349in}}%
\pgfpathclose%
\pgfusepath{stroke,fill}%
\end{pgfscope}%
\begin{pgfscope}%
\pgfpathrectangle{\pgfqpoint{0.050000in}{0.050000in}}{\pgfqpoint{4.900000in}{4.900000in}}%
\pgfusepath{clip}%
\pgfsetbuttcap%
\pgfsetroundjoin%
\definecolor{currentfill}{rgb}{0.981546,0.981546,0.981546}%
\pgfsetfillcolor{currentfill}%
\pgfsetfillopacity{0.900000}%
\pgfsetlinewidth{0.501875pt}%
\definecolor{currentstroke}{rgb}{0.200000,0.200000,0.200000}%
\pgfsetstrokecolor{currentstroke}%
\pgfsetdash{}{0pt}%
\pgfpathmoveto{\pgfqpoint{2.767128in}{1.963536in}}%
\pgfpathlineto{\pgfqpoint{2.810940in}{1.953682in}}%
\pgfpathlineto{\pgfqpoint{2.842252in}{1.936164in}}%
\pgfpathlineto{\pgfqpoint{2.798418in}{1.946014in}}%
\pgfpathlineto{\pgfqpoint{2.767128in}{1.963536in}}%
\pgfpathclose%
\pgfusepath{stroke,fill}%
\end{pgfscope}%
\begin{pgfscope}%
\pgfpathrectangle{\pgfqpoint{0.050000in}{0.050000in}}{\pgfqpoint{4.900000in}{4.900000in}}%
\pgfusepath{clip}%
\pgfsetbuttcap%
\pgfsetroundjoin%
\definecolor{currentfill}{rgb}{0.998155,0.998155,0.998155}%
\pgfsetfillcolor{currentfill}%
\pgfsetfillopacity{0.900000}%
\pgfsetlinewidth{0.501875pt}%
\definecolor{currentstroke}{rgb}{0.200000,0.200000,0.200000}%
\pgfsetstrokecolor{currentstroke}%
\pgfsetdash{}{0pt}%
\pgfpathmoveto{\pgfqpoint{3.186744in}{1.872999in}}%
\pgfpathlineto{\pgfqpoint{3.230156in}{1.882272in}}%
\pgfpathlineto{\pgfqpoint{3.262108in}{1.863308in}}%
\pgfpathlineto{\pgfqpoint{3.218668in}{1.854206in}}%
\pgfpathlineto{\pgfqpoint{3.186744in}{1.872999in}}%
\pgfpathclose%
\pgfusepath{stroke,fill}%
\end{pgfscope}%
\begin{pgfscope}%
\pgfpathrectangle{\pgfqpoint{0.050000in}{0.050000in}}{\pgfqpoint{4.900000in}{4.900000in}}%
\pgfusepath{clip}%
\pgfsetbuttcap%
\pgfsetroundjoin%
\definecolor{currentfill}{rgb}{0.968627,0.968627,0.968627}%
\pgfsetfillcolor{currentfill}%
\pgfsetfillopacity{0.900000}%
\pgfsetlinewidth{0.501875pt}%
\definecolor{currentstroke}{rgb}{0.200000,0.200000,0.200000}%
\pgfsetstrokecolor{currentstroke}%
\pgfsetdash{}{0pt}%
\pgfpathmoveto{\pgfqpoint{2.616927in}{2.034636in}}%
\pgfpathlineto{\pgfqpoint{2.660949in}{2.013626in}}%
\pgfpathlineto{\pgfqpoint{2.692027in}{1.996204in}}%
\pgfpathlineto{\pgfqpoint{2.647989in}{2.016927in}}%
\pgfpathlineto{\pgfqpoint{2.616927in}{2.034636in}}%
\pgfpathclose%
\pgfusepath{stroke,fill}%
\end{pgfscope}%
\begin{pgfscope}%
\pgfpathrectangle{\pgfqpoint{0.050000in}{0.050000in}}{\pgfqpoint{4.900000in}{4.900000in}}%
\pgfusepath{clip}%
\pgfsetbuttcap%
\pgfsetroundjoin%
\definecolor{currentfill}{rgb}{0.777778,0.777778,0.777778}%
\pgfsetfillcolor{currentfill}%
\pgfsetfillopacity{0.900000}%
\pgfsetlinewidth{0.501875pt}%
\definecolor{currentstroke}{rgb}{0.200000,0.200000,0.200000}%
\pgfsetstrokecolor{currentstroke}%
\pgfsetdash{}{0pt}%
\pgfpathmoveto{\pgfqpoint{1.772787in}{2.674563in}}%
\pgfpathlineto{\pgfqpoint{1.818512in}{2.631055in}}%
\pgfpathlineto{\pgfqpoint{1.848317in}{2.611726in}}%
\pgfpathlineto{\pgfqpoint{1.802585in}{2.655259in}}%
\pgfpathlineto{\pgfqpoint{1.772787in}{2.674563in}}%
\pgfpathclose%
\pgfusepath{stroke,fill}%
\end{pgfscope}%
\begin{pgfscope}%
\pgfpathrectangle{\pgfqpoint{0.050000in}{0.050000in}}{\pgfqpoint{4.900000in}{4.900000in}}%
\pgfusepath{clip}%
\pgfsetbuttcap%
\pgfsetroundjoin%
\definecolor{currentfill}{rgb}{0.987082,0.987082,0.987082}%
\pgfsetfillcolor{currentfill}%
\pgfsetfillopacity{0.900000}%
\pgfsetlinewidth{0.501875pt}%
\definecolor{currentstroke}{rgb}{0.200000,0.200000,0.200000}%
\pgfsetstrokecolor{currentstroke}%
\pgfsetdash{}{0pt}%
\pgfpathmoveto{\pgfqpoint{2.842252in}{1.936164in}}%
\pgfpathlineto{\pgfqpoint{2.885983in}{1.931146in}}%
\pgfpathlineto{\pgfqpoint{2.917415in}{1.913417in}}%
\pgfpathlineto{\pgfqpoint{2.873660in}{1.918533in}}%
\pgfpathlineto{\pgfqpoint{2.842252in}{1.936164in}}%
\pgfpathclose%
\pgfusepath{stroke,fill}%
\end{pgfscope}%
\begin{pgfscope}%
\pgfpathrectangle{\pgfqpoint{0.050000in}{0.050000in}}{\pgfqpoint{4.900000in}{4.900000in}}%
\pgfusepath{clip}%
\pgfsetbuttcap%
\pgfsetroundjoin%
\definecolor{currentfill}{rgb}{1.000000,1.000000,1.000000}%
\pgfsetfillcolor{currentfill}%
\pgfsetfillopacity{0.900000}%
\pgfsetlinewidth{0.501875pt}%
\definecolor{currentstroke}{rgb}{0.200000,0.200000,0.200000}%
\pgfsetstrokecolor{currentstroke}%
\pgfsetdash{}{0pt}%
\pgfpathmoveto{\pgfqpoint{3.456302in}{1.858834in}}%
\pgfpathlineto{\pgfqpoint{3.499468in}{1.871444in}}%
\pgfpathlineto{\pgfqpoint{3.531827in}{1.851942in}}%
\pgfpathlineto{\pgfqpoint{3.488632in}{1.839395in}}%
\pgfpathlineto{\pgfqpoint{3.456302in}{1.858834in}}%
\pgfpathclose%
\pgfusepath{stroke,fill}%
\end{pgfscope}%
\begin{pgfscope}%
\pgfpathrectangle{\pgfqpoint{0.050000in}{0.050000in}}{\pgfqpoint{4.900000in}{4.900000in}}%
\pgfusepath{clip}%
\pgfsetbuttcap%
\pgfsetroundjoin%
\definecolor{currentfill}{rgb}{0.959400,0.959400,0.959400}%
\pgfsetfillcolor{currentfill}%
\pgfsetfillopacity{0.900000}%
\pgfsetlinewidth{0.501875pt}%
\definecolor{currentstroke}{rgb}{0.200000,0.200000,0.200000}%
\pgfsetstrokecolor{currentstroke}%
\pgfsetdash{}{0pt}%
\pgfpathmoveto{\pgfqpoint{2.541806in}{2.078961in}}%
\pgfpathlineto{\pgfqpoint{2.585961in}{2.052243in}}%
\pgfpathlineto{\pgfqpoint{2.616927in}{2.034636in}}%
\pgfpathlineto{\pgfqpoint{2.572759in}{2.060935in}}%
\pgfpathlineto{\pgfqpoint{2.541806in}{2.078961in}}%
\pgfpathclose%
\pgfusepath{stroke,fill}%
\end{pgfscope}%
\begin{pgfscope}%
\pgfpathrectangle{\pgfqpoint{0.050000in}{0.050000in}}{\pgfqpoint{4.900000in}{4.900000in}}%
\pgfusepath{clip}%
\pgfsetbuttcap%
\pgfsetroundjoin%
\definecolor{currentfill}{rgb}{0.990773,0.990773,0.990773}%
\pgfsetfillcolor{currentfill}%
\pgfsetfillopacity{0.900000}%
\pgfsetlinewidth{0.501875pt}%
\definecolor{currentstroke}{rgb}{0.200000,0.200000,0.200000}%
\pgfsetstrokecolor{currentstroke}%
\pgfsetdash{}{0pt}%
\pgfpathmoveto{\pgfqpoint{2.917415in}{1.913417in}}%
\pgfpathlineto{\pgfqpoint{2.961077in}{1.912538in}}%
\pgfpathlineto{\pgfqpoint{2.992631in}{1.894545in}}%
\pgfpathlineto{\pgfqpoint{2.948944in}{1.895588in}}%
\pgfpathlineto{\pgfqpoint{2.917415in}{1.913417in}}%
\pgfpathclose%
\pgfusepath{stroke,fill}%
\end{pgfscope}%
\begin{pgfscope}%
\pgfpathrectangle{\pgfqpoint{0.050000in}{0.050000in}}{\pgfqpoint{4.900000in}{4.900000in}}%
\pgfusepath{clip}%
\pgfsetbuttcap%
\pgfsetroundjoin%
\definecolor{currentfill}{rgb}{0.898378,0.898378,0.898378}%
\pgfsetfillcolor{currentfill}%
\pgfsetfillopacity{0.900000}%
\pgfsetlinewidth{0.501875pt}%
\definecolor{currentstroke}{rgb}{0.200000,0.200000,0.200000}%
\pgfsetstrokecolor{currentstroke}%
\pgfsetdash{}{0pt}%
\pgfpathmoveto{\pgfqpoint{2.210249in}{2.322250in}}%
\pgfpathlineto{\pgfqpoint{2.255082in}{2.280538in}}%
\pgfpathlineto{\pgfqpoint{2.285547in}{2.261224in}}%
\pgfpathlineto{\pgfqpoint{2.240709in}{2.302683in}}%
\pgfpathlineto{\pgfqpoint{2.210249in}{2.322250in}}%
\pgfpathclose%
\pgfusepath{stroke,fill}%
\end{pgfscope}%
\begin{pgfscope}%
\pgfpathrectangle{\pgfqpoint{0.050000in}{0.050000in}}{\pgfqpoint{4.900000in}{4.900000in}}%
\pgfusepath{clip}%
\pgfsetbuttcap%
\pgfsetroundjoin%
\definecolor{currentfill}{rgb}{0.998155,0.998155,0.998155}%
\pgfsetfillcolor{currentfill}%
\pgfsetfillopacity{0.900000}%
\pgfsetlinewidth{0.501875pt}%
\definecolor{currentstroke}{rgb}{0.200000,0.200000,0.200000}%
\pgfsetstrokecolor{currentstroke}%
\pgfsetdash{}{0pt}%
\pgfpathmoveto{\pgfqpoint{3.262108in}{1.863308in}}%
\pgfpathlineto{\pgfqpoint{3.305467in}{1.873809in}}%
\pgfpathlineto{\pgfqpoint{3.337543in}{1.854657in}}%
\pgfpathlineto{\pgfqpoint{3.294157in}{1.844297in}}%
\pgfpathlineto{\pgfqpoint{3.262108in}{1.863308in}}%
\pgfpathclose%
\pgfusepath{stroke,fill}%
\end{pgfscope}%
\begin{pgfscope}%
\pgfpathrectangle{\pgfqpoint{0.050000in}{0.050000in}}{\pgfqpoint{4.900000in}{4.900000in}}%
\pgfusepath{clip}%
\pgfsetbuttcap%
\pgfsetroundjoin%
\definecolor{currentfill}{rgb}{0.829450,0.829450,0.829450}%
\pgfsetfillcolor{currentfill}%
\pgfsetfillopacity{0.900000}%
\pgfsetlinewidth{0.501875pt}%
\definecolor{currentstroke}{rgb}{0.200000,0.200000,0.200000}%
\pgfsetstrokecolor{currentstroke}%
\pgfsetdash{}{0pt}%
\pgfpathmoveto{\pgfqpoint{1.953755in}{2.529487in}}%
\pgfpathlineto{\pgfqpoint{1.999127in}{2.486262in}}%
\pgfpathlineto{\pgfqpoint{2.029214in}{2.466713in}}%
\pgfpathlineto{\pgfqpoint{1.983836in}{2.509959in}}%
\pgfpathlineto{\pgfqpoint{1.953755in}{2.529487in}}%
\pgfpathclose%
\pgfusepath{stroke,fill}%
\end{pgfscope}%
\begin{pgfscope}%
\pgfpathrectangle{\pgfqpoint{0.050000in}{0.050000in}}{\pgfqpoint{4.900000in}{4.900000in}}%
\pgfusepath{clip}%
\pgfsetbuttcap%
\pgfsetroundjoin%
\definecolor{currentfill}{rgb}{0.950173,0.950173,0.950173}%
\pgfsetfillcolor{currentfill}%
\pgfsetfillopacity{0.900000}%
\pgfsetlinewidth{0.501875pt}%
\definecolor{currentstroke}{rgb}{0.200000,0.200000,0.200000}%
\pgfsetstrokecolor{currentstroke}%
\pgfsetdash{}{0pt}%
\pgfpathmoveto{\pgfqpoint{2.466641in}{2.128863in}}%
\pgfpathlineto{\pgfqpoint{2.510949in}{2.096915in}}%
\pgfpathlineto{\pgfqpoint{2.541806in}{2.078961in}}%
\pgfpathlineto{\pgfqpoint{2.497488in}{2.110411in}}%
\pgfpathlineto{\pgfqpoint{2.466641in}{2.128863in}}%
\pgfpathclose%
\pgfusepath{stroke,fill}%
\end{pgfscope}%
\begin{pgfscope}%
\pgfpathrectangle{\pgfqpoint{0.050000in}{0.050000in}}{\pgfqpoint{4.900000in}{4.900000in}}%
\pgfusepath{clip}%
\pgfsetbuttcap%
\pgfsetroundjoin%
\definecolor{currentfill}{rgb}{0.992618,0.992618,0.992618}%
\pgfsetfillcolor{currentfill}%
\pgfsetfillopacity{0.900000}%
\pgfsetlinewidth{0.501875pt}%
\definecolor{currentstroke}{rgb}{0.200000,0.200000,0.200000}%
\pgfsetstrokecolor{currentstroke}%
\pgfsetdash{}{0pt}%
\pgfpathmoveto{\pgfqpoint{2.992631in}{1.894545in}}%
\pgfpathlineto{\pgfqpoint{3.036232in}{1.897087in}}%
\pgfpathlineto{\pgfqpoint{3.067908in}{1.878814in}}%
\pgfpathlineto{\pgfqpoint{3.024281in}{1.876468in}}%
\pgfpathlineto{\pgfqpoint{2.992631in}{1.894545in}}%
\pgfpathclose%
\pgfusepath{stroke,fill}%
\end{pgfscope}%
\begin{pgfscope}%
\pgfpathrectangle{\pgfqpoint{0.050000in}{0.050000in}}{\pgfqpoint{4.900000in}{4.900000in}}%
\pgfusepath{clip}%
\pgfsetbuttcap%
\pgfsetroundjoin%
\definecolor{currentfill}{rgb}{1.000000,1.000000,1.000000}%
\pgfsetfillcolor{currentfill}%
\pgfsetfillopacity{0.900000}%
\pgfsetlinewidth{0.501875pt}%
\definecolor{currentstroke}{rgb}{0.200000,0.200000,0.200000}%
\pgfsetstrokecolor{currentstroke}%
\pgfsetdash{}{0pt}%
\pgfpathmoveto{\pgfqpoint{3.531827in}{1.851942in}}%
\pgfpathlineto{\pgfqpoint{3.574935in}{1.864854in}}%
\pgfpathlineto{\pgfqpoint{3.607421in}{1.845256in}}%
\pgfpathlineto{\pgfqpoint{3.564284in}{1.832390in}}%
\pgfpathlineto{\pgfqpoint{3.531827in}{1.851942in}}%
\pgfpathclose%
\pgfusepath{stroke,fill}%
\end{pgfscope}%
\begin{pgfscope}%
\pgfpathrectangle{\pgfqpoint{0.050000in}{0.050000in}}{\pgfqpoint{4.900000in}{4.900000in}}%
\pgfusepath{clip}%
\pgfsetbuttcap%
\pgfsetroundjoin%
\definecolor{currentfill}{rgb}{0.753664,0.753664,0.753664}%
\pgfsetfillcolor{currentfill}%
\pgfsetfillopacity{0.900000}%
\pgfsetlinewidth{0.501875pt}%
\definecolor{currentstroke}{rgb}{0.200000,0.200000,0.200000}%
\pgfsetstrokecolor{currentstroke}%
\pgfsetdash{}{0pt}%
\pgfpathmoveto{\pgfqpoint{1.697173in}{2.737470in}}%
\pgfpathlineto{\pgfqpoint{1.743078in}{2.693810in}}%
\pgfpathlineto{\pgfqpoint{1.772787in}{2.674563in}}%
\pgfpathlineto{\pgfqpoint{1.726874in}{2.718250in}}%
\pgfpathlineto{\pgfqpoint{1.697173in}{2.737470in}}%
\pgfpathclose%
\pgfusepath{stroke,fill}%
\end{pgfscope}%
\begin{pgfscope}%
\pgfpathrectangle{\pgfqpoint{0.050000in}{0.050000in}}{\pgfqpoint{4.900000in}{4.900000in}}%
\pgfusepath{clip}%
\pgfsetbuttcap%
\pgfsetroundjoin%
\definecolor{currentfill}{rgb}{1.000000,1.000000,1.000000}%
\pgfsetfillcolor{currentfill}%
\pgfsetfillopacity{0.900000}%
\pgfsetlinewidth{0.501875pt}%
\definecolor{currentstroke}{rgb}{0.200000,0.200000,0.200000}%
\pgfsetstrokecolor{currentstroke}%
\pgfsetdash{}{0pt}%
\pgfpathmoveto{\pgfqpoint{3.337543in}{1.854657in}}%
\pgfpathlineto{\pgfqpoint{3.380849in}{1.866068in}}%
\pgfpathlineto{\pgfqpoint{3.413052in}{1.846760in}}%
\pgfpathlineto{\pgfqpoint{3.369718in}{1.835460in}}%
\pgfpathlineto{\pgfqpoint{3.337543in}{1.854657in}}%
\pgfpathclose%
\pgfusepath{stroke,fill}%
\end{pgfscope}%
\begin{pgfscope}%
\pgfpathrectangle{\pgfqpoint{0.050000in}{0.050000in}}{\pgfqpoint{4.900000in}{4.900000in}}%
\pgfusepath{clip}%
\pgfsetbuttcap%
\pgfsetroundjoin%
\definecolor{currentfill}{rgb}{0.994464,0.994464,0.994464}%
\pgfsetfillcolor{currentfill}%
\pgfsetfillopacity{0.900000}%
\pgfsetlinewidth{0.501875pt}%
\definecolor{currentstroke}{rgb}{0.200000,0.200000,0.200000}%
\pgfsetstrokecolor{currentstroke}%
\pgfsetdash{}{0pt}%
\pgfpathmoveto{\pgfqpoint{3.067908in}{1.878814in}}%
\pgfpathlineto{\pgfqpoint{3.111453in}{1.884101in}}%
\pgfpathlineto{\pgfqpoint{3.143253in}{1.865556in}}%
\pgfpathlineto{\pgfqpoint{3.099681in}{1.860469in}}%
\pgfpathlineto{\pgfqpoint{3.067908in}{1.878814in}}%
\pgfpathclose%
\pgfusepath{stroke,fill}%
\end{pgfscope}%
\begin{pgfscope}%
\pgfpathrectangle{\pgfqpoint{0.050000in}{0.050000in}}{\pgfqpoint{4.900000in}{4.900000in}}%
\pgfusepath{clip}%
\pgfsetbuttcap%
\pgfsetroundjoin%
\definecolor{currentfill}{rgb}{0.875740,0.875740,0.875740}%
\pgfsetfillcolor{currentfill}%
\pgfsetfillopacity{0.900000}%
\pgfsetlinewidth{0.501875pt}%
\definecolor{currentstroke}{rgb}{0.200000,0.200000,0.200000}%
\pgfsetstrokecolor{currentstroke}%
\pgfsetdash{}{0pt}%
\pgfpathmoveto{\pgfqpoint{2.134863in}{2.384440in}}%
\pgfpathlineto{\pgfqpoint{2.179880in}{2.341807in}}%
\pgfpathlineto{\pgfqpoint{2.210249in}{2.322250in}}%
\pgfpathlineto{\pgfqpoint{2.165228in}{2.364775in}}%
\pgfpathlineto{\pgfqpoint{2.134863in}{2.384440in}}%
\pgfpathclose%
\pgfusepath{stroke,fill}%
\end{pgfscope}%
\begin{pgfscope}%
\pgfpathrectangle{\pgfqpoint{0.050000in}{0.050000in}}{\pgfqpoint{4.900000in}{4.900000in}}%
\pgfusepath{clip}%
\pgfsetbuttcap%
\pgfsetroundjoin%
\definecolor{currentfill}{rgb}{0.937993,0.937993,0.937993}%
\pgfsetfillcolor{currentfill}%
\pgfsetfillopacity{0.900000}%
\pgfsetlinewidth{0.501875pt}%
\definecolor{currentstroke}{rgb}{0.200000,0.200000,0.200000}%
\pgfsetstrokecolor{currentstroke}%
\pgfsetdash{}{0pt}%
\pgfpathmoveto{\pgfqpoint{2.391411in}{2.183560in}}%
\pgfpathlineto{\pgfqpoint{2.435887in}{2.147277in}}%
\pgfpathlineto{\pgfqpoint{2.466641in}{2.128863in}}%
\pgfpathlineto{\pgfqpoint{2.422157in}{2.164651in}}%
\pgfpathlineto{\pgfqpoint{2.391411in}{2.183560in}}%
\pgfpathclose%
\pgfusepath{stroke,fill}%
\end{pgfscope}%
\begin{pgfscope}%
\pgfpathrectangle{\pgfqpoint{0.050000in}{0.050000in}}{\pgfqpoint{4.900000in}{4.900000in}}%
\pgfusepath{clip}%
\pgfsetbuttcap%
\pgfsetroundjoin%
\definecolor{currentfill}{rgb}{0.805336,0.805336,0.805336}%
\pgfsetfillcolor{currentfill}%
\pgfsetfillopacity{0.900000}%
\pgfsetlinewidth{0.501875pt}%
\definecolor{currentstroke}{rgb}{0.200000,0.200000,0.200000}%
\pgfsetstrokecolor{currentstroke}%
\pgfsetdash{}{0pt}%
\pgfpathmoveto{\pgfqpoint{1.878212in}{2.592338in}}%
\pgfpathlineto{\pgfqpoint{1.923764in}{2.548958in}}%
\pgfpathlineto{\pgfqpoint{1.953755in}{2.529487in}}%
\pgfpathlineto{\pgfqpoint{1.908196in}{2.572892in}}%
\pgfpathlineto{\pgfqpoint{1.878212in}{2.592338in}}%
\pgfpathclose%
\pgfusepath{stroke,fill}%
\end{pgfscope}%
\begin{pgfscope}%
\pgfpathrectangle{\pgfqpoint{0.050000in}{0.050000in}}{\pgfqpoint{4.900000in}{4.900000in}}%
\pgfusepath{clip}%
\pgfsetbuttcap%
\pgfsetroundjoin%
\definecolor{currentfill}{rgb}{1.000000,1.000000,1.000000}%
\pgfsetfillcolor{currentfill}%
\pgfsetfillopacity{0.900000}%
\pgfsetlinewidth{0.501875pt}%
\definecolor{currentstroke}{rgb}{0.200000,0.200000,0.200000}%
\pgfsetstrokecolor{currentstroke}%
\pgfsetdash{}{0pt}%
\pgfpathmoveto{\pgfqpoint{3.607421in}{1.845256in}}%
\pgfpathlineto{\pgfqpoint{3.650470in}{1.858349in}}%
\pgfpathlineto{\pgfqpoint{3.683084in}{1.838661in}}%
\pgfpathlineto{\pgfqpoint{3.640006in}{1.825604in}}%
\pgfpathlineto{\pgfqpoint{3.607421in}{1.845256in}}%
\pgfpathclose%
\pgfusepath{stroke,fill}%
\end{pgfscope}%
\begin{pgfscope}%
\pgfpathrectangle{\pgfqpoint{0.050000in}{0.050000in}}{\pgfqpoint{4.900000in}{4.900000in}}%
\pgfusepath{clip}%
\pgfsetbuttcap%
\pgfsetroundjoin%
\definecolor{currentfill}{rgb}{0.996309,0.996309,0.996309}%
\pgfsetfillcolor{currentfill}%
\pgfsetfillopacity{0.900000}%
\pgfsetlinewidth{0.501875pt}%
\definecolor{currentstroke}{rgb}{0.200000,0.200000,0.200000}%
\pgfsetstrokecolor{currentstroke}%
\pgfsetdash{}{0pt}%
\pgfpathmoveto{\pgfqpoint{3.143253in}{1.865556in}}%
\pgfpathlineto{\pgfqpoint{3.186744in}{1.872999in}}%
\pgfpathlineto{\pgfqpoint{3.218668in}{1.854206in}}%
\pgfpathlineto{\pgfqpoint{3.175150in}{1.846951in}}%
\pgfpathlineto{\pgfqpoint{3.143253in}{1.865556in}}%
\pgfpathclose%
\pgfusepath{stroke,fill}%
\end{pgfscope}%
\begin{pgfscope}%
\pgfpathrectangle{\pgfqpoint{0.050000in}{0.050000in}}{\pgfqpoint{4.900000in}{4.900000in}}%
\pgfusepath{clip}%
\pgfsetbuttcap%
\pgfsetroundjoin%
\definecolor{currentfill}{rgb}{1.000000,1.000000,1.000000}%
\pgfsetfillcolor{currentfill}%
\pgfsetfillopacity{0.900000}%
\pgfsetlinewidth{0.501875pt}%
\definecolor{currentstroke}{rgb}{0.200000,0.200000,0.200000}%
\pgfsetstrokecolor{currentstroke}%
\pgfsetdash{}{0pt}%
\pgfpathmoveto{\pgfqpoint{3.413052in}{1.846760in}}%
\pgfpathlineto{\pgfqpoint{3.456302in}{1.858834in}}%
\pgfpathlineto{\pgfqpoint{3.488632in}{1.839395in}}%
\pgfpathlineto{\pgfqpoint{3.445353in}{1.827405in}}%
\pgfpathlineto{\pgfqpoint{3.413052in}{1.846760in}}%
\pgfpathclose%
\pgfusepath{stroke,fill}%
\end{pgfscope}%
\begin{pgfscope}%
\pgfpathrectangle{\pgfqpoint{0.050000in}{0.050000in}}{\pgfqpoint{4.900000in}{4.900000in}}%
\pgfusepath{clip}%
\pgfsetbuttcap%
\pgfsetroundjoin%
\definecolor{currentfill}{rgb}{0.918185,0.918185,0.918185}%
\pgfsetfillcolor{currentfill}%
\pgfsetfillopacity{0.900000}%
\pgfsetlinewidth{0.501875pt}%
\definecolor{currentstroke}{rgb}{0.200000,0.200000,0.200000}%
\pgfsetstrokecolor{currentstroke}%
\pgfsetdash{}{0pt}%
\pgfpathmoveto{\pgfqpoint{2.316103in}{2.241917in}}%
\pgfpathlineto{\pgfqpoint{2.360758in}{2.202459in}}%
\pgfpathlineto{\pgfqpoint{2.391411in}{2.183560in}}%
\pgfpathlineto{\pgfqpoint{2.346751in}{2.222610in}}%
\pgfpathlineto{\pgfqpoint{2.316103in}{2.241917in}}%
\pgfpathclose%
\pgfusepath{stroke,fill}%
\end{pgfscope}%
\begin{pgfscope}%
\pgfpathrectangle{\pgfqpoint{0.050000in}{0.050000in}}{\pgfqpoint{4.900000in}{4.900000in}}%
\pgfusepath{clip}%
\pgfsetbuttcap%
\pgfsetroundjoin%
\definecolor{currentfill}{rgb}{0.855932,0.855932,0.855932}%
\pgfsetfillcolor{currentfill}%
\pgfsetfillopacity{0.900000}%
\pgfsetlinewidth{0.501875pt}%
\definecolor{currentstroke}{rgb}{0.200000,0.200000,0.200000}%
\pgfsetstrokecolor{currentstroke}%
\pgfsetdash{}{0pt}%
\pgfpathmoveto{\pgfqpoint{2.059392in}{2.447112in}}%
\pgfpathlineto{\pgfqpoint{2.104590in}{2.404073in}}%
\pgfpathlineto{\pgfqpoint{2.134863in}{2.384440in}}%
\pgfpathlineto{\pgfqpoint{2.089660in}{2.427460in}}%
\pgfpathlineto{\pgfqpoint{2.059392in}{2.447112in}}%
\pgfpathclose%
\pgfusepath{stroke,fill}%
\end{pgfscope}%
\begin{pgfscope}%
\pgfpathrectangle{\pgfqpoint{0.050000in}{0.050000in}}{\pgfqpoint{4.900000in}{4.900000in}}%
\pgfusepath{clip}%
\pgfsetbuttcap%
\pgfsetroundjoin%
\definecolor{currentfill}{rgb}{1.000000,1.000000,1.000000}%
\pgfsetfillcolor{currentfill}%
\pgfsetfillopacity{0.900000}%
\pgfsetlinewidth{0.501875pt}%
\definecolor{currentstroke}{rgb}{0.200000,0.200000,0.200000}%
\pgfsetstrokecolor{currentstroke}%
\pgfsetdash{}{0pt}%
\pgfpathmoveto{\pgfqpoint{3.683084in}{1.838661in}}%
\pgfpathlineto{\pgfqpoint{3.726072in}{1.851825in}}%
\pgfpathlineto{\pgfqpoint{3.758813in}{1.832047in}}%
\pgfpathlineto{\pgfqpoint{3.715796in}{1.818916in}}%
\pgfpathlineto{\pgfqpoint{3.683084in}{1.838661in}}%
\pgfpathclose%
\pgfusepath{stroke,fill}%
\end{pgfscope}%
\begin{pgfscope}%
\pgfpathrectangle{\pgfqpoint{0.050000in}{0.050000in}}{\pgfqpoint{4.900000in}{4.900000in}}%
\pgfusepath{clip}%
\pgfsetbuttcap%
\pgfsetroundjoin%
\definecolor{currentfill}{rgb}{0.976009,0.976009,0.976009}%
\pgfsetfillcolor{currentfill}%
\pgfsetfillopacity{0.900000}%
\pgfsetlinewidth{0.501875pt}%
\definecolor{currentstroke}{rgb}{0.200000,0.200000,0.200000}%
\pgfsetstrokecolor{currentstroke}%
\pgfsetdash{}{0pt}%
\pgfpathmoveto{\pgfqpoint{2.723201in}{1.978664in}}%
\pgfpathlineto{\pgfqpoint{2.767128in}{1.963536in}}%
\pgfpathlineto{\pgfqpoint{2.798418in}{1.946014in}}%
\pgfpathlineto{\pgfqpoint{2.754472in}{1.961012in}}%
\pgfpathlineto{\pgfqpoint{2.723201in}{1.978664in}}%
\pgfpathclose%
\pgfusepath{stroke,fill}%
\end{pgfscope}%
\begin{pgfscope}%
\pgfpathrectangle{\pgfqpoint{0.050000in}{0.050000in}}{\pgfqpoint{4.900000in}{4.900000in}}%
\pgfusepath{clip}%
\pgfsetbuttcap%
\pgfsetroundjoin%
\definecolor{currentfill}{rgb}{0.981546,0.981546,0.981546}%
\pgfsetfillcolor{currentfill}%
\pgfsetfillopacity{0.900000}%
\pgfsetlinewidth{0.501875pt}%
\definecolor{currentstroke}{rgb}{0.200000,0.200000,0.200000}%
\pgfsetstrokecolor{currentstroke}%
\pgfsetdash{}{0pt}%
\pgfpathmoveto{\pgfqpoint{2.798418in}{1.946014in}}%
\pgfpathlineto{\pgfqpoint{2.842252in}{1.936164in}}%
\pgfpathlineto{\pgfqpoint{2.873660in}{1.918533in}}%
\pgfpathlineto{\pgfqpoint{2.829805in}{1.928377in}}%
\pgfpathlineto{\pgfqpoint{2.798418in}{1.946014in}}%
\pgfpathclose%
\pgfusepath{stroke,fill}%
\end{pgfscope}%
\begin{pgfscope}%
\pgfpathrectangle{\pgfqpoint{0.050000in}{0.050000in}}{\pgfqpoint{4.900000in}{4.900000in}}%
\pgfusepath{clip}%
\pgfsetbuttcap%
\pgfsetroundjoin%
\definecolor{currentfill}{rgb}{0.998155,0.998155,0.998155}%
\pgfsetfillcolor{currentfill}%
\pgfsetfillopacity{0.900000}%
\pgfsetlinewidth{0.501875pt}%
\definecolor{currentstroke}{rgb}{0.200000,0.200000,0.200000}%
\pgfsetstrokecolor{currentstroke}%
\pgfsetdash{}{0pt}%
\pgfpathmoveto{\pgfqpoint{3.218668in}{1.854206in}}%
\pgfpathlineto{\pgfqpoint{3.262108in}{1.863308in}}%
\pgfpathlineto{\pgfqpoint{3.294157in}{1.844297in}}%
\pgfpathlineto{\pgfqpoint{3.250690in}{1.835360in}}%
\pgfpathlineto{\pgfqpoint{3.218668in}{1.854206in}}%
\pgfpathclose%
\pgfusepath{stroke,fill}%
\end{pgfscope}%
\begin{pgfscope}%
\pgfpathrectangle{\pgfqpoint{0.050000in}{0.050000in}}{\pgfqpoint{4.900000in}{4.900000in}}%
\pgfusepath{clip}%
\pgfsetbuttcap%
\pgfsetroundjoin%
\definecolor{currentfill}{rgb}{0.777778,0.777778,0.777778}%
\pgfsetfillcolor{currentfill}%
\pgfsetfillopacity{0.900000}%
\pgfsetlinewidth{0.501875pt}%
\definecolor{currentstroke}{rgb}{0.200000,0.200000,0.200000}%
\pgfsetstrokecolor{currentstroke}%
\pgfsetdash{}{0pt}%
\pgfpathmoveto{\pgfqpoint{1.802585in}{2.655259in}}%
\pgfpathlineto{\pgfqpoint{1.848317in}{2.611726in}}%
\pgfpathlineto{\pgfqpoint{1.878212in}{2.592338in}}%
\pgfpathlineto{\pgfqpoint{1.832472in}{2.635896in}}%
\pgfpathlineto{\pgfqpoint{1.802585in}{2.655259in}}%
\pgfpathclose%
\pgfusepath{stroke,fill}%
\end{pgfscope}%
\begin{pgfscope}%
\pgfpathrectangle{\pgfqpoint{0.050000in}{0.050000in}}{\pgfqpoint{4.900000in}{4.900000in}}%
\pgfusepath{clip}%
\pgfsetbuttcap%
\pgfsetroundjoin%
\definecolor{currentfill}{rgb}{0.968627,0.968627,0.968627}%
\pgfsetfillcolor{currentfill}%
\pgfsetfillopacity{0.900000}%
\pgfsetlinewidth{0.501875pt}%
\definecolor{currentstroke}{rgb}{0.200000,0.200000,0.200000}%
\pgfsetstrokecolor{currentstroke}%
\pgfsetdash{}{0pt}%
\pgfpathmoveto{\pgfqpoint{2.647989in}{2.016927in}}%
\pgfpathlineto{\pgfqpoint{2.692027in}{1.996204in}}%
\pgfpathlineto{\pgfqpoint{2.723201in}{1.978664in}}%
\pgfpathlineto{\pgfqpoint{2.679146in}{1.999120in}}%
\pgfpathlineto{\pgfqpoint{2.647989in}{2.016927in}}%
\pgfpathclose%
\pgfusepath{stroke,fill}%
\end{pgfscope}%
\begin{pgfscope}%
\pgfpathrectangle{\pgfqpoint{0.050000in}{0.050000in}}{\pgfqpoint{4.900000in}{4.900000in}}%
\pgfusepath{clip}%
\pgfsetbuttcap%
\pgfsetroundjoin%
\definecolor{currentfill}{rgb}{0.985236,0.985236,0.985236}%
\pgfsetfillcolor{currentfill}%
\pgfsetfillopacity{0.900000}%
\pgfsetlinewidth{0.501875pt}%
\definecolor{currentstroke}{rgb}{0.200000,0.200000,0.200000}%
\pgfsetstrokecolor{currentstroke}%
\pgfsetdash{}{0pt}%
\pgfpathmoveto{\pgfqpoint{2.873660in}{1.918533in}}%
\pgfpathlineto{\pgfqpoint{2.917415in}{1.913417in}}%
\pgfpathlineto{\pgfqpoint{2.948944in}{1.895588in}}%
\pgfpathlineto{\pgfqpoint{2.905166in}{1.900794in}}%
\pgfpathlineto{\pgfqpoint{2.873660in}{1.918533in}}%
\pgfpathclose%
\pgfusepath{stroke,fill}%
\end{pgfscope}%
\begin{pgfscope}%
\pgfpathrectangle{\pgfqpoint{0.050000in}{0.050000in}}{\pgfqpoint{4.900000in}{4.900000in}}%
\pgfusepath{clip}%
\pgfsetbuttcap%
\pgfsetroundjoin%
\definecolor{currentfill}{rgb}{1.000000,1.000000,1.000000}%
\pgfsetfillcolor{currentfill}%
\pgfsetfillopacity{0.900000}%
\pgfsetlinewidth{0.501875pt}%
\definecolor{currentstroke}{rgb}{0.200000,0.200000,0.200000}%
\pgfsetstrokecolor{currentstroke}%
\pgfsetdash{}{0pt}%
\pgfpathmoveto{\pgfqpoint{3.488632in}{1.839395in}}%
\pgfpathlineto{\pgfqpoint{3.531827in}{1.851942in}}%
\pgfpathlineto{\pgfqpoint{3.564284in}{1.832390in}}%
\pgfpathlineto{\pgfqpoint{3.521061in}{1.819907in}}%
\pgfpathlineto{\pgfqpoint{3.488632in}{1.839395in}}%
\pgfpathclose%
\pgfusepath{stroke,fill}%
\end{pgfscope}%
\begin{pgfscope}%
\pgfpathrectangle{\pgfqpoint{0.050000in}{0.050000in}}{\pgfqpoint{4.900000in}{4.900000in}}%
\pgfusepath{clip}%
\pgfsetbuttcap%
\pgfsetroundjoin%
\definecolor{currentfill}{rgb}{0.959400,0.959400,0.959400}%
\pgfsetfillcolor{currentfill}%
\pgfsetfillopacity{0.900000}%
\pgfsetlinewidth{0.501875pt}%
\definecolor{currentstroke}{rgb}{0.200000,0.200000,0.200000}%
\pgfsetstrokecolor{currentstroke}%
\pgfsetdash{}{0pt}%
\pgfpathmoveto{\pgfqpoint{2.572759in}{2.060935in}}%
\pgfpathlineto{\pgfqpoint{2.616927in}{2.034636in}}%
\pgfpathlineto{\pgfqpoint{2.647989in}{2.016927in}}%
\pgfpathlineto{\pgfqpoint{2.603807in}{2.042835in}}%
\pgfpathlineto{\pgfqpoint{2.572759in}{2.060935in}}%
\pgfpathclose%
\pgfusepath{stroke,fill}%
\end{pgfscope}%
\begin{pgfscope}%
\pgfpathrectangle{\pgfqpoint{0.050000in}{0.050000in}}{\pgfqpoint{4.900000in}{4.900000in}}%
\pgfusepath{clip}%
\pgfsetbuttcap%
\pgfsetroundjoin%
\definecolor{currentfill}{rgb}{0.898378,0.898378,0.898378}%
\pgfsetfillcolor{currentfill}%
\pgfsetfillopacity{0.900000}%
\pgfsetlinewidth{0.501875pt}%
\definecolor{currentstroke}{rgb}{0.200000,0.200000,0.200000}%
\pgfsetstrokecolor{currentstroke}%
\pgfsetdash{}{0pt}%
\pgfpathmoveto{\pgfqpoint{2.240709in}{2.302683in}}%
\pgfpathlineto{\pgfqpoint{2.285547in}{2.261224in}}%
\pgfpathlineto{\pgfqpoint{2.316103in}{2.241917in}}%
\pgfpathlineto{\pgfqpoint{2.271261in}{2.283108in}}%
\pgfpathlineto{\pgfqpoint{2.240709in}{2.302683in}}%
\pgfpathclose%
\pgfusepath{stroke,fill}%
\end{pgfscope}%
\begin{pgfscope}%
\pgfpathrectangle{\pgfqpoint{0.050000in}{0.050000in}}{\pgfqpoint{4.900000in}{4.900000in}}%
\pgfusepath{clip}%
\pgfsetbuttcap%
\pgfsetroundjoin%
\definecolor{currentfill}{rgb}{0.988927,0.988927,0.988927}%
\pgfsetfillcolor{currentfill}%
\pgfsetfillopacity{0.900000}%
\pgfsetlinewidth{0.501875pt}%
\definecolor{currentstroke}{rgb}{0.200000,0.200000,0.200000}%
\pgfsetstrokecolor{currentstroke}%
\pgfsetdash{}{0pt}%
\pgfpathmoveto{\pgfqpoint{2.948944in}{1.895588in}}%
\pgfpathlineto{\pgfqpoint{2.992631in}{1.894545in}}%
\pgfpathlineto{\pgfqpoint{3.024281in}{1.876468in}}%
\pgfpathlineto{\pgfqpoint{2.980570in}{1.877662in}}%
\pgfpathlineto{\pgfqpoint{2.948944in}{1.895588in}}%
\pgfpathclose%
\pgfusepath{stroke,fill}%
\end{pgfscope}%
\begin{pgfscope}%
\pgfpathrectangle{\pgfqpoint{0.050000in}{0.050000in}}{\pgfqpoint{4.900000in}{4.900000in}}%
\pgfusepath{clip}%
\pgfsetbuttcap%
\pgfsetroundjoin%
\definecolor{currentfill}{rgb}{0.998155,0.998155,0.998155}%
\pgfsetfillcolor{currentfill}%
\pgfsetfillopacity{0.900000}%
\pgfsetlinewidth{0.501875pt}%
\definecolor{currentstroke}{rgb}{0.200000,0.200000,0.200000}%
\pgfsetstrokecolor{currentstroke}%
\pgfsetdash{}{0pt}%
\pgfpathmoveto{\pgfqpoint{3.294157in}{1.844297in}}%
\pgfpathlineto{\pgfqpoint{3.337543in}{1.854657in}}%
\pgfpathlineto{\pgfqpoint{3.369718in}{1.835460in}}%
\pgfpathlineto{\pgfqpoint{3.326304in}{1.825237in}}%
\pgfpathlineto{\pgfqpoint{3.294157in}{1.844297in}}%
\pgfpathclose%
\pgfusepath{stroke,fill}%
\end{pgfscope}%
\begin{pgfscope}%
\pgfpathrectangle{\pgfqpoint{0.050000in}{0.050000in}}{\pgfqpoint{4.900000in}{4.900000in}}%
\pgfusepath{clip}%
\pgfsetbuttcap%
\pgfsetroundjoin%
\definecolor{currentfill}{rgb}{0.829450,0.829450,0.829450}%
\pgfsetfillcolor{currentfill}%
\pgfsetfillopacity{0.900000}%
\pgfsetlinewidth{0.501875pt}%
\definecolor{currentstroke}{rgb}{0.200000,0.200000,0.200000}%
\pgfsetstrokecolor{currentstroke}%
\pgfsetdash{}{0pt}%
\pgfpathmoveto{\pgfqpoint{1.983836in}{2.509959in}}%
\pgfpathlineto{\pgfqpoint{2.029214in}{2.466713in}}%
\pgfpathlineto{\pgfqpoint{2.059392in}{2.447112in}}%
\pgfpathlineto{\pgfqpoint{2.014008in}{2.490373in}}%
\pgfpathlineto{\pgfqpoint{1.983836in}{2.509959in}}%
\pgfpathclose%
\pgfusepath{stroke,fill}%
\end{pgfscope}%
\begin{pgfscope}%
\pgfpathrectangle{\pgfqpoint{0.050000in}{0.050000in}}{\pgfqpoint{4.900000in}{4.900000in}}%
\pgfusepath{clip}%
\pgfsetbuttcap%
\pgfsetroundjoin%
\definecolor{currentfill}{rgb}{0.950173,0.950173,0.950173}%
\pgfsetfillcolor{currentfill}%
\pgfsetfillopacity{0.900000}%
\pgfsetlinewidth{0.501875pt}%
\definecolor{currentstroke}{rgb}{0.200000,0.200000,0.200000}%
\pgfsetstrokecolor{currentstroke}%
\pgfsetdash{}{0pt}%
\pgfpathmoveto{\pgfqpoint{2.497488in}{2.110411in}}%
\pgfpathlineto{\pgfqpoint{2.541806in}{2.078961in}}%
\pgfpathlineto{\pgfqpoint{2.572759in}{2.060935in}}%
\pgfpathlineto{\pgfqpoint{2.528430in}{2.091916in}}%
\pgfpathlineto{\pgfqpoint{2.497488in}{2.110411in}}%
\pgfpathclose%
\pgfusepath{stroke,fill}%
\end{pgfscope}%
\begin{pgfscope}%
\pgfpathrectangle{\pgfqpoint{0.050000in}{0.050000in}}{\pgfqpoint{4.900000in}{4.900000in}}%
\pgfusepath{clip}%
\pgfsetbuttcap%
\pgfsetroundjoin%
\definecolor{currentfill}{rgb}{0.992618,0.992618,0.992618}%
\pgfsetfillcolor{currentfill}%
\pgfsetfillopacity{0.900000}%
\pgfsetlinewidth{0.501875pt}%
\definecolor{currentstroke}{rgb}{0.200000,0.200000,0.200000}%
\pgfsetstrokecolor{currentstroke}%
\pgfsetdash{}{0pt}%
\pgfpathmoveto{\pgfqpoint{3.024281in}{1.876468in}}%
\pgfpathlineto{\pgfqpoint{3.067908in}{1.878814in}}%
\pgfpathlineto{\pgfqpoint{3.099681in}{1.860469in}}%
\pgfpathlineto{\pgfqpoint{3.056029in}{1.858307in}}%
\pgfpathlineto{\pgfqpoint{3.024281in}{1.876468in}}%
\pgfpathclose%
\pgfusepath{stroke,fill}%
\end{pgfscope}%
\begin{pgfscope}%
\pgfpathrectangle{\pgfqpoint{0.050000in}{0.050000in}}{\pgfqpoint{4.900000in}{4.900000in}}%
\pgfusepath{clip}%
\pgfsetbuttcap%
\pgfsetroundjoin%
\definecolor{currentfill}{rgb}{0.753664,0.753664,0.753664}%
\pgfsetfillcolor{currentfill}%
\pgfsetfillopacity{0.900000}%
\pgfsetlinewidth{0.501875pt}%
\definecolor{currentstroke}{rgb}{0.200000,0.200000,0.200000}%
\pgfsetstrokecolor{currentstroke}%
\pgfsetdash{}{0pt}%
\pgfpathmoveto{\pgfqpoint{1.726874in}{2.718250in}}%
\pgfpathlineto{\pgfqpoint{1.772787in}{2.674563in}}%
\pgfpathlineto{\pgfqpoint{1.802585in}{2.655259in}}%
\pgfpathlineto{\pgfqpoint{1.756664in}{2.698971in}}%
\pgfpathlineto{\pgfqpoint{1.726874in}{2.718250in}}%
\pgfpathclose%
\pgfusepath{stroke,fill}%
\end{pgfscope}%
\begin{pgfscope}%
\pgfpathrectangle{\pgfqpoint{0.050000in}{0.050000in}}{\pgfqpoint{4.900000in}{4.900000in}}%
\pgfusepath{clip}%
\pgfsetbuttcap%
\pgfsetroundjoin%
\definecolor{currentfill}{rgb}{1.000000,1.000000,1.000000}%
\pgfsetfillcolor{currentfill}%
\pgfsetfillopacity{0.900000}%
\pgfsetlinewidth{0.501875pt}%
\definecolor{currentstroke}{rgb}{0.200000,0.200000,0.200000}%
\pgfsetstrokecolor{currentstroke}%
\pgfsetdash{}{0pt}%
\pgfpathmoveto{\pgfqpoint{3.564284in}{1.832390in}}%
\pgfpathlineto{\pgfqpoint{3.607421in}{1.845256in}}%
\pgfpathlineto{\pgfqpoint{3.640006in}{1.825604in}}%
\pgfpathlineto{\pgfqpoint{3.596840in}{1.812787in}}%
\pgfpathlineto{\pgfqpoint{3.564284in}{1.832390in}}%
\pgfpathclose%
\pgfusepath{stroke,fill}%
\end{pgfscope}%
\begin{pgfscope}%
\pgfpathrectangle{\pgfqpoint{0.050000in}{0.050000in}}{\pgfqpoint{4.900000in}{4.900000in}}%
\pgfusepath{clip}%
\pgfsetbuttcap%
\pgfsetroundjoin%
\definecolor{currentfill}{rgb}{0.878570,0.878570,0.878570}%
\pgfsetfillcolor{currentfill}%
\pgfsetfillopacity{0.900000}%
\pgfsetlinewidth{0.501875pt}%
\definecolor{currentstroke}{rgb}{0.200000,0.200000,0.200000}%
\pgfsetstrokecolor{currentstroke}%
\pgfsetdash{}{0pt}%
\pgfpathmoveto{\pgfqpoint{2.165228in}{2.364775in}}%
\pgfpathlineto{\pgfqpoint{2.210249in}{2.322250in}}%
\pgfpathlineto{\pgfqpoint{2.240709in}{2.302683in}}%
\pgfpathlineto{\pgfqpoint{2.195684in}{2.345080in}}%
\pgfpathlineto{\pgfqpoint{2.165228in}{2.364775in}}%
\pgfpathclose%
\pgfusepath{stroke,fill}%
\end{pgfscope}%
\begin{pgfscope}%
\pgfpathrectangle{\pgfqpoint{0.050000in}{0.050000in}}{\pgfqpoint{4.900000in}{4.900000in}}%
\pgfusepath{clip}%
\pgfsetbuttcap%
\pgfsetroundjoin%
\definecolor{currentfill}{rgb}{1.000000,1.000000,1.000000}%
\pgfsetfillcolor{currentfill}%
\pgfsetfillopacity{0.900000}%
\pgfsetlinewidth{0.501875pt}%
\definecolor{currentstroke}{rgb}{0.200000,0.200000,0.200000}%
\pgfsetstrokecolor{currentstroke}%
\pgfsetdash{}{0pt}%
\pgfpathmoveto{\pgfqpoint{3.369718in}{1.835460in}}%
\pgfpathlineto{\pgfqpoint{3.413052in}{1.846760in}}%
\pgfpathlineto{\pgfqpoint{3.445353in}{1.827405in}}%
\pgfpathlineto{\pgfqpoint{3.401992in}{1.816215in}}%
\pgfpathlineto{\pgfqpoint{3.369718in}{1.835460in}}%
\pgfpathclose%
\pgfusepath{stroke,fill}%
\end{pgfscope}%
\begin{pgfscope}%
\pgfpathrectangle{\pgfqpoint{0.050000in}{0.050000in}}{\pgfqpoint{4.900000in}{4.900000in}}%
\pgfusepath{clip}%
\pgfsetbuttcap%
\pgfsetroundjoin%
\definecolor{currentfill}{rgb}{0.994464,0.994464,0.994464}%
\pgfsetfillcolor{currentfill}%
\pgfsetfillopacity{0.900000}%
\pgfsetlinewidth{0.501875pt}%
\definecolor{currentstroke}{rgb}{0.200000,0.200000,0.200000}%
\pgfsetstrokecolor{currentstroke}%
\pgfsetdash{}{0pt}%
\pgfpathmoveto{\pgfqpoint{3.099681in}{1.860469in}}%
\pgfpathlineto{\pgfqpoint{3.143253in}{1.865556in}}%
\pgfpathlineto{\pgfqpoint{3.175150in}{1.846951in}}%
\pgfpathlineto{\pgfqpoint{3.131552in}{1.842054in}}%
\pgfpathlineto{\pgfqpoint{3.099681in}{1.860469in}}%
\pgfpathclose%
\pgfusepath{stroke,fill}%
\end{pgfscope}%
\begin{pgfscope}%
\pgfpathrectangle{\pgfqpoint{0.050000in}{0.050000in}}{\pgfqpoint{4.900000in}{4.900000in}}%
\pgfusepath{clip}%
\pgfsetbuttcap%
\pgfsetroundjoin%
\definecolor{currentfill}{rgb}{0.937993,0.937993,0.937993}%
\pgfsetfillcolor{currentfill}%
\pgfsetfillopacity{0.900000}%
\pgfsetlinewidth{0.501875pt}%
\definecolor{currentstroke}{rgb}{0.200000,0.200000,0.200000}%
\pgfsetstrokecolor{currentstroke}%
\pgfsetdash{}{0pt}%
\pgfpathmoveto{\pgfqpoint{2.422157in}{2.164651in}}%
\pgfpathlineto{\pgfqpoint{2.466641in}{2.128863in}}%
\pgfpathlineto{\pgfqpoint{2.497488in}{2.110411in}}%
\pgfpathlineto{\pgfqpoint{2.452997in}{2.145725in}}%
\pgfpathlineto{\pgfqpoint{2.422157in}{2.164651in}}%
\pgfpathclose%
\pgfusepath{stroke,fill}%
\end{pgfscope}%
\begin{pgfscope}%
\pgfpathrectangle{\pgfqpoint{0.050000in}{0.050000in}}{\pgfqpoint{4.900000in}{4.900000in}}%
\pgfusepath{clip}%
\pgfsetbuttcap%
\pgfsetroundjoin%
\definecolor{currentfill}{rgb}{0.805336,0.805336,0.805336}%
\pgfsetfillcolor{currentfill}%
\pgfsetfillopacity{0.900000}%
\pgfsetlinewidth{0.501875pt}%
\definecolor{currentstroke}{rgb}{0.200000,0.200000,0.200000}%
\pgfsetstrokecolor{currentstroke}%
\pgfsetdash{}{0pt}%
\pgfpathmoveto{\pgfqpoint{1.908196in}{2.572892in}}%
\pgfpathlineto{\pgfqpoint{1.953755in}{2.529487in}}%
\pgfpathlineto{\pgfqpoint{1.983836in}{2.509959in}}%
\pgfpathlineto{\pgfqpoint{1.938271in}{2.553388in}}%
\pgfpathlineto{\pgfqpoint{1.908196in}{2.572892in}}%
\pgfpathclose%
\pgfusepath{stroke,fill}%
\end{pgfscope}%
\begin{pgfscope}%
\pgfpathrectangle{\pgfqpoint{0.050000in}{0.050000in}}{\pgfqpoint{4.900000in}{4.900000in}}%
\pgfusepath{clip}%
\pgfsetbuttcap%
\pgfsetroundjoin%
\definecolor{currentfill}{rgb}{1.000000,1.000000,1.000000}%
\pgfsetfillcolor{currentfill}%
\pgfsetfillopacity{0.900000}%
\pgfsetlinewidth{0.501875pt}%
\definecolor{currentstroke}{rgb}{0.200000,0.200000,0.200000}%
\pgfsetstrokecolor{currentstroke}%
\pgfsetdash{}{0pt}%
\pgfpathmoveto{\pgfqpoint{3.640006in}{1.825604in}}%
\pgfpathlineto{\pgfqpoint{3.683084in}{1.838661in}}%
\pgfpathlineto{\pgfqpoint{3.715796in}{1.818916in}}%
\pgfpathlineto{\pgfqpoint{3.672690in}{1.805899in}}%
\pgfpathlineto{\pgfqpoint{3.640006in}{1.825604in}}%
\pgfpathclose%
\pgfusepath{stroke,fill}%
\end{pgfscope}%
\begin{pgfscope}%
\pgfpathrectangle{\pgfqpoint{0.050000in}{0.050000in}}{\pgfqpoint{4.900000in}{4.900000in}}%
\pgfusepath{clip}%
\pgfsetbuttcap%
\pgfsetroundjoin%
\definecolor{currentfill}{rgb}{0.720185,0.720185,0.720185}%
\pgfsetfillcolor{currentfill}%
\pgfsetfillopacity{0.900000}%
\pgfsetlinewidth{0.501875pt}%
\definecolor{currentstroke}{rgb}{0.200000,0.200000,0.200000}%
\pgfsetstrokecolor{currentstroke}%
\pgfsetdash{}{0pt}%
\pgfpathmoveto{\pgfqpoint{1.651079in}{2.781311in}}%
\pgfpathlineto{\pgfqpoint{1.697173in}{2.737470in}}%
\pgfpathlineto{\pgfqpoint{1.726874in}{2.718250in}}%
\pgfpathlineto{\pgfqpoint{1.680772in}{2.762116in}}%
\pgfpathlineto{\pgfqpoint{1.651079in}{2.781311in}}%
\pgfpathclose%
\pgfusepath{stroke,fill}%
\end{pgfscope}%
\begin{pgfscope}%
\pgfpathrectangle{\pgfqpoint{0.050000in}{0.050000in}}{\pgfqpoint{4.900000in}{4.900000in}}%
\pgfusepath{clip}%
\pgfsetbuttcap%
\pgfsetroundjoin%
\definecolor{currentfill}{rgb}{0.996309,0.996309,0.996309}%
\pgfsetfillcolor{currentfill}%
\pgfsetfillopacity{0.900000}%
\pgfsetlinewidth{0.501875pt}%
\definecolor{currentstroke}{rgb}{0.200000,0.200000,0.200000}%
\pgfsetstrokecolor{currentstroke}%
\pgfsetdash{}{0pt}%
\pgfpathmoveto{\pgfqpoint{3.175150in}{1.846951in}}%
\pgfpathlineto{\pgfqpoint{3.218668in}{1.854206in}}%
\pgfpathlineto{\pgfqpoint{3.250690in}{1.835360in}}%
\pgfpathlineto{\pgfqpoint{3.207145in}{1.828284in}}%
\pgfpathlineto{\pgfqpoint{3.175150in}{1.846951in}}%
\pgfpathclose%
\pgfusepath{stroke,fill}%
\end{pgfscope}%
\begin{pgfscope}%
\pgfpathrectangle{\pgfqpoint{0.050000in}{0.050000in}}{\pgfqpoint{4.900000in}{4.900000in}}%
\pgfusepath{clip}%
\pgfsetbuttcap%
\pgfsetroundjoin%
\definecolor{currentfill}{rgb}{1.000000,1.000000,1.000000}%
\pgfsetfillcolor{currentfill}%
\pgfsetfillopacity{0.900000}%
\pgfsetlinewidth{0.501875pt}%
\definecolor{currentstroke}{rgb}{0.200000,0.200000,0.200000}%
\pgfsetstrokecolor{currentstroke}%
\pgfsetdash{}{0pt}%
\pgfpathmoveto{\pgfqpoint{3.445353in}{1.827405in}}%
\pgfpathlineto{\pgfqpoint{3.488632in}{1.839395in}}%
\pgfpathlineto{\pgfqpoint{3.521061in}{1.819907in}}%
\pgfpathlineto{\pgfqpoint{3.477754in}{1.808003in}}%
\pgfpathlineto{\pgfqpoint{3.445353in}{1.827405in}}%
\pgfpathclose%
\pgfusepath{stroke,fill}%
\end{pgfscope}%
\begin{pgfscope}%
\pgfpathrectangle{\pgfqpoint{0.050000in}{0.050000in}}{\pgfqpoint{4.900000in}{4.900000in}}%
\pgfusepath{clip}%
\pgfsetbuttcap%
\pgfsetroundjoin%
\definecolor{currentfill}{rgb}{0.918185,0.918185,0.918185}%
\pgfsetfillcolor{currentfill}%
\pgfsetfillopacity{0.900000}%
\pgfsetlinewidth{0.501875pt}%
\definecolor{currentstroke}{rgb}{0.200000,0.200000,0.200000}%
\pgfsetstrokecolor{currentstroke}%
\pgfsetdash{}{0pt}%
\pgfpathmoveto{\pgfqpoint{2.346751in}{2.222610in}}%
\pgfpathlineto{\pgfqpoint{2.391411in}{2.183560in}}%
\pgfpathlineto{\pgfqpoint{2.422157in}{2.164651in}}%
\pgfpathlineto{\pgfqpoint{2.377492in}{2.203298in}}%
\pgfpathlineto{\pgfqpoint{2.346751in}{2.222610in}}%
\pgfpathclose%
\pgfusepath{stroke,fill}%
\end{pgfscope}%
\begin{pgfscope}%
\pgfpathrectangle{\pgfqpoint{0.050000in}{0.050000in}}{\pgfqpoint{4.900000in}{4.900000in}}%
\pgfusepath{clip}%
\pgfsetbuttcap%
\pgfsetroundjoin%
\definecolor{currentfill}{rgb}{0.855932,0.855932,0.855932}%
\pgfsetfillcolor{currentfill}%
\pgfsetfillopacity{0.900000}%
\pgfsetlinewidth{0.501875pt}%
\definecolor{currentstroke}{rgb}{0.200000,0.200000,0.200000}%
\pgfsetstrokecolor{currentstroke}%
\pgfsetdash{}{0pt}%
\pgfpathmoveto{\pgfqpoint{2.089660in}{2.427460in}}%
\pgfpathlineto{\pgfqpoint{2.134863in}{2.384440in}}%
\pgfpathlineto{\pgfqpoint{2.165228in}{2.364775in}}%
\pgfpathlineto{\pgfqpoint{2.120020in}{2.407761in}}%
\pgfpathlineto{\pgfqpoint{2.089660in}{2.427460in}}%
\pgfpathclose%
\pgfusepath{stroke,fill}%
\end{pgfscope}%
\begin{pgfscope}%
\pgfpathrectangle{\pgfqpoint{0.050000in}{0.050000in}}{\pgfqpoint{4.900000in}{4.900000in}}%
\pgfusepath{clip}%
\pgfsetbuttcap%
\pgfsetroundjoin%
\definecolor{currentfill}{rgb}{1.000000,1.000000,1.000000}%
\pgfsetfillcolor{currentfill}%
\pgfsetfillopacity{0.900000}%
\pgfsetlinewidth{0.501875pt}%
\definecolor{currentstroke}{rgb}{0.200000,0.200000,0.200000}%
\pgfsetstrokecolor{currentstroke}%
\pgfsetdash{}{0pt}%
\pgfpathmoveto{\pgfqpoint{3.715796in}{1.818916in}}%
\pgfpathlineto{\pgfqpoint{3.758813in}{1.832047in}}%
\pgfpathlineto{\pgfqpoint{3.791654in}{1.812210in}}%
\pgfpathlineto{\pgfqpoint{3.748609in}{1.799114in}}%
\pgfpathlineto{\pgfqpoint{3.715796in}{1.818916in}}%
\pgfpathclose%
\pgfusepath{stroke,fill}%
\end{pgfscope}%
\begin{pgfscope}%
\pgfpathrectangle{\pgfqpoint{0.050000in}{0.050000in}}{\pgfqpoint{4.900000in}{4.900000in}}%
\pgfusepath{clip}%
\pgfsetbuttcap%
\pgfsetroundjoin%
\definecolor{currentfill}{rgb}{0.777778,0.777778,0.777778}%
\pgfsetfillcolor{currentfill}%
\pgfsetfillopacity{0.900000}%
\pgfsetlinewidth{0.501875pt}%
\definecolor{currentstroke}{rgb}{0.200000,0.200000,0.200000}%
\pgfsetstrokecolor{currentstroke}%
\pgfsetdash{}{0pt}%
\pgfpathmoveto{\pgfqpoint{1.832472in}{2.635896in}}%
\pgfpathlineto{\pgfqpoint{1.878212in}{2.592338in}}%
\pgfpathlineto{\pgfqpoint{1.908196in}{2.572892in}}%
\pgfpathlineto{\pgfqpoint{1.862450in}{2.616476in}}%
\pgfpathlineto{\pgfqpoint{1.832472in}{2.635896in}}%
\pgfpathclose%
\pgfusepath{stroke,fill}%
\end{pgfscope}%
\begin{pgfscope}%
\pgfpathrectangle{\pgfqpoint{0.050000in}{0.050000in}}{\pgfqpoint{4.900000in}{4.900000in}}%
\pgfusepath{clip}%
\pgfsetbuttcap%
\pgfsetroundjoin%
\definecolor{currentfill}{rgb}{0.974164,0.974164,0.974164}%
\pgfsetfillcolor{currentfill}%
\pgfsetfillopacity{0.900000}%
\pgfsetlinewidth{0.501875pt}%
\definecolor{currentstroke}{rgb}{0.200000,0.200000,0.200000}%
\pgfsetstrokecolor{currentstroke}%
\pgfsetdash{}{0pt}%
\pgfpathmoveto{\pgfqpoint{2.754472in}{1.961012in}}%
\pgfpathlineto{\pgfqpoint{2.798418in}{1.946014in}}%
\pgfpathlineto{\pgfqpoint{2.829805in}{1.928377in}}%
\pgfpathlineto{\pgfqpoint{2.785840in}{1.943249in}}%
\pgfpathlineto{\pgfqpoint{2.754472in}{1.961012in}}%
\pgfpathclose%
\pgfusepath{stroke,fill}%
\end{pgfscope}%
\begin{pgfscope}%
\pgfpathrectangle{\pgfqpoint{0.050000in}{0.050000in}}{\pgfqpoint{4.900000in}{4.900000in}}%
\pgfusepath{clip}%
\pgfsetbuttcap%
\pgfsetroundjoin%
\definecolor{currentfill}{rgb}{0.998155,0.998155,0.998155}%
\pgfsetfillcolor{currentfill}%
\pgfsetfillopacity{0.900000}%
\pgfsetlinewidth{0.501875pt}%
\definecolor{currentstroke}{rgb}{0.200000,0.200000,0.200000}%
\pgfsetstrokecolor{currentstroke}%
\pgfsetdash{}{0pt}%
\pgfpathmoveto{\pgfqpoint{3.250690in}{1.835360in}}%
\pgfpathlineto{\pgfqpoint{3.294157in}{1.844297in}}%
\pgfpathlineto{\pgfqpoint{3.326304in}{1.825237in}}%
\pgfpathlineto{\pgfqpoint{3.282810in}{1.816458in}}%
\pgfpathlineto{\pgfqpoint{3.250690in}{1.835360in}}%
\pgfpathclose%
\pgfusepath{stroke,fill}%
\end{pgfscope}%
\begin{pgfscope}%
\pgfpathrectangle{\pgfqpoint{0.050000in}{0.050000in}}{\pgfqpoint{4.900000in}{4.900000in}}%
\pgfusepath{clip}%
\pgfsetbuttcap%
\pgfsetroundjoin%
\definecolor{currentfill}{rgb}{0.981546,0.981546,0.981546}%
\pgfsetfillcolor{currentfill}%
\pgfsetfillopacity{0.900000}%
\pgfsetlinewidth{0.501875pt}%
\definecolor{currentstroke}{rgb}{0.200000,0.200000,0.200000}%
\pgfsetstrokecolor{currentstroke}%
\pgfsetdash{}{0pt}%
\pgfpathmoveto{\pgfqpoint{2.829805in}{1.928377in}}%
\pgfpathlineto{\pgfqpoint{2.873660in}{1.918533in}}%
\pgfpathlineto{\pgfqpoint{2.905166in}{1.900794in}}%
\pgfpathlineto{\pgfqpoint{2.861290in}{1.910627in}}%
\pgfpathlineto{\pgfqpoint{2.829805in}{1.928377in}}%
\pgfpathclose%
\pgfusepath{stroke,fill}%
\end{pgfscope}%
\begin{pgfscope}%
\pgfpathrectangle{\pgfqpoint{0.050000in}{0.050000in}}{\pgfqpoint{4.900000in}{4.900000in}}%
\pgfusepath{clip}%
\pgfsetbuttcap%
\pgfsetroundjoin%
\definecolor{currentfill}{rgb}{0.968627,0.968627,0.968627}%
\pgfsetfillcolor{currentfill}%
\pgfsetfillopacity{0.900000}%
\pgfsetlinewidth{0.501875pt}%
\definecolor{currentstroke}{rgb}{0.200000,0.200000,0.200000}%
\pgfsetstrokecolor{currentstroke}%
\pgfsetdash{}{0pt}%
\pgfpathmoveto{\pgfqpoint{2.679146in}{1.999120in}}%
\pgfpathlineto{\pgfqpoint{2.723201in}{1.978664in}}%
\pgfpathlineto{\pgfqpoint{2.754472in}{1.961012in}}%
\pgfpathlineto{\pgfqpoint{2.710400in}{1.981215in}}%
\pgfpathlineto{\pgfqpoint{2.679146in}{1.999120in}}%
\pgfpathclose%
\pgfusepath{stroke,fill}%
\end{pgfscope}%
\begin{pgfscope}%
\pgfpathrectangle{\pgfqpoint{0.050000in}{0.050000in}}{\pgfqpoint{4.900000in}{4.900000in}}%
\pgfusepath{clip}%
\pgfsetbuttcap%
\pgfsetroundjoin%
\definecolor{currentfill}{rgb}{0.985236,0.985236,0.985236}%
\pgfsetfillcolor{currentfill}%
\pgfsetfillopacity{0.900000}%
\pgfsetlinewidth{0.501875pt}%
\definecolor{currentstroke}{rgb}{0.200000,0.200000,0.200000}%
\pgfsetstrokecolor{currentstroke}%
\pgfsetdash{}{0pt}%
\pgfpathmoveto{\pgfqpoint{2.905166in}{1.900794in}}%
\pgfpathlineto{\pgfqpoint{2.948944in}{1.895588in}}%
\pgfpathlineto{\pgfqpoint{2.980570in}{1.877662in}}%
\pgfpathlineto{\pgfqpoint{2.936769in}{1.882948in}}%
\pgfpathlineto{\pgfqpoint{2.905166in}{1.900794in}}%
\pgfpathclose%
\pgfusepath{stroke,fill}%
\end{pgfscope}%
\begin{pgfscope}%
\pgfpathrectangle{\pgfqpoint{0.050000in}{0.050000in}}{\pgfqpoint{4.900000in}{4.900000in}}%
\pgfusepath{clip}%
\pgfsetbuttcap%
\pgfsetroundjoin%
\definecolor{currentfill}{rgb}{1.000000,1.000000,1.000000}%
\pgfsetfillcolor{currentfill}%
\pgfsetfillopacity{0.900000}%
\pgfsetlinewidth{0.501875pt}%
\definecolor{currentstroke}{rgb}{0.200000,0.200000,0.200000}%
\pgfsetstrokecolor{currentstroke}%
\pgfsetdash{}{0pt}%
\pgfpathmoveto{\pgfqpoint{3.521061in}{1.819907in}}%
\pgfpathlineto{\pgfqpoint{3.564284in}{1.832390in}}%
\pgfpathlineto{\pgfqpoint{3.596840in}{1.812787in}}%
\pgfpathlineto{\pgfqpoint{3.553589in}{1.800371in}}%
\pgfpathlineto{\pgfqpoint{3.521061in}{1.819907in}}%
\pgfpathclose%
\pgfusepath{stroke,fill}%
\end{pgfscope}%
\begin{pgfscope}%
\pgfpathrectangle{\pgfqpoint{0.050000in}{0.050000in}}{\pgfqpoint{4.900000in}{4.900000in}}%
\pgfusepath{clip}%
\pgfsetbuttcap%
\pgfsetroundjoin%
\definecolor{currentfill}{rgb}{0.959400,0.959400,0.959400}%
\pgfsetfillcolor{currentfill}%
\pgfsetfillopacity{0.900000}%
\pgfsetlinewidth{0.501875pt}%
\definecolor{currentstroke}{rgb}{0.200000,0.200000,0.200000}%
\pgfsetstrokecolor{currentstroke}%
\pgfsetdash{}{0pt}%
\pgfpathmoveto{\pgfqpoint{2.603807in}{2.042835in}}%
\pgfpathlineto{\pgfqpoint{2.647989in}{2.016927in}}%
\pgfpathlineto{\pgfqpoint{2.679146in}{1.999120in}}%
\pgfpathlineto{\pgfqpoint{2.634950in}{2.024661in}}%
\pgfpathlineto{\pgfqpoint{2.603807in}{2.042835in}}%
\pgfpathclose%
\pgfusepath{stroke,fill}%
\end{pgfscope}%
\begin{pgfscope}%
\pgfpathrectangle{\pgfqpoint{0.050000in}{0.050000in}}{\pgfqpoint{4.900000in}{4.900000in}}%
\pgfusepath{clip}%
\pgfsetbuttcap%
\pgfsetroundjoin%
\definecolor{currentfill}{rgb}{0.898378,0.898378,0.898378}%
\pgfsetfillcolor{currentfill}%
\pgfsetfillopacity{0.900000}%
\pgfsetlinewidth{0.501875pt}%
\definecolor{currentstroke}{rgb}{0.200000,0.200000,0.200000}%
\pgfsetstrokecolor{currentstroke}%
\pgfsetdash{}{0pt}%
\pgfpathmoveto{\pgfqpoint{2.271261in}{2.283108in}}%
\pgfpathlineto{\pgfqpoint{2.316103in}{2.241917in}}%
\pgfpathlineto{\pgfqpoint{2.346751in}{2.222610in}}%
\pgfpathlineto{\pgfqpoint{2.301905in}{2.263522in}}%
\pgfpathlineto{\pgfqpoint{2.271261in}{2.283108in}}%
\pgfpathclose%
\pgfusepath{stroke,fill}%
\end{pgfscope}%
\begin{pgfscope}%
\pgfpathrectangle{\pgfqpoint{0.050000in}{0.050000in}}{\pgfqpoint{4.900000in}{4.900000in}}%
\pgfusepath{clip}%
\pgfsetbuttcap%
\pgfsetroundjoin%
\definecolor{currentfill}{rgb}{0.988927,0.988927,0.988927}%
\pgfsetfillcolor{currentfill}%
\pgfsetfillopacity{0.900000}%
\pgfsetlinewidth{0.501875pt}%
\definecolor{currentstroke}{rgb}{0.200000,0.200000,0.200000}%
\pgfsetstrokecolor{currentstroke}%
\pgfsetdash{}{0pt}%
\pgfpathmoveto{\pgfqpoint{2.980570in}{1.877662in}}%
\pgfpathlineto{\pgfqpoint{3.024281in}{1.876468in}}%
\pgfpathlineto{\pgfqpoint{3.056029in}{1.858307in}}%
\pgfpathlineto{\pgfqpoint{3.012294in}{1.859641in}}%
\pgfpathlineto{\pgfqpoint{2.980570in}{1.877662in}}%
\pgfpathclose%
\pgfusepath{stroke,fill}%
\end{pgfscope}%
\begin{pgfscope}%
\pgfpathrectangle{\pgfqpoint{0.050000in}{0.050000in}}{\pgfqpoint{4.900000in}{4.900000in}}%
\pgfusepath{clip}%
\pgfsetbuttcap%
\pgfsetroundjoin%
\definecolor{currentfill}{rgb}{0.832895,0.832895,0.832895}%
\pgfsetfillcolor{currentfill}%
\pgfsetfillopacity{0.900000}%
\pgfsetlinewidth{0.501875pt}%
\definecolor{currentstroke}{rgb}{0.200000,0.200000,0.200000}%
\pgfsetstrokecolor{currentstroke}%
\pgfsetdash{}{0pt}%
\pgfpathmoveto{\pgfqpoint{2.014008in}{2.490373in}}%
\pgfpathlineto{\pgfqpoint{2.059392in}{2.447112in}}%
\pgfpathlineto{\pgfqpoint{2.089660in}{2.427460in}}%
\pgfpathlineto{\pgfqpoint{2.044271in}{2.470732in}}%
\pgfpathlineto{\pgfqpoint{2.014008in}{2.490373in}}%
\pgfpathclose%
\pgfusepath{stroke,fill}%
\end{pgfscope}%
\begin{pgfscope}%
\pgfpathrectangle{\pgfqpoint{0.050000in}{0.050000in}}{\pgfqpoint{4.900000in}{4.900000in}}%
\pgfusepath{clip}%
\pgfsetbuttcap%
\pgfsetroundjoin%
\definecolor{currentfill}{rgb}{0.998155,0.998155,0.998155}%
\pgfsetfillcolor{currentfill}%
\pgfsetfillopacity{0.900000}%
\pgfsetlinewidth{0.501875pt}%
\definecolor{currentstroke}{rgb}{0.200000,0.200000,0.200000}%
\pgfsetstrokecolor{currentstroke}%
\pgfsetdash{}{0pt}%
\pgfpathmoveto{\pgfqpoint{3.326304in}{1.825237in}}%
\pgfpathlineto{\pgfqpoint{3.369718in}{1.835460in}}%
\pgfpathlineto{\pgfqpoint{3.401992in}{1.816215in}}%
\pgfpathlineto{\pgfqpoint{3.358550in}{1.806127in}}%
\pgfpathlineto{\pgfqpoint{3.326304in}{1.825237in}}%
\pgfpathclose%
\pgfusepath{stroke,fill}%
\end{pgfscope}%
\begin{pgfscope}%
\pgfpathrectangle{\pgfqpoint{0.050000in}{0.050000in}}{\pgfqpoint{4.900000in}{4.900000in}}%
\pgfusepath{clip}%
\pgfsetbuttcap%
\pgfsetroundjoin%
\definecolor{currentfill}{rgb}{0.950173,0.950173,0.950173}%
\pgfsetfillcolor{currentfill}%
\pgfsetfillopacity{0.900000}%
\pgfsetlinewidth{0.501875pt}%
\definecolor{currentstroke}{rgb}{0.200000,0.200000,0.200000}%
\pgfsetstrokecolor{currentstroke}%
\pgfsetdash{}{0pt}%
\pgfpathmoveto{\pgfqpoint{2.528430in}{2.091916in}}%
\pgfpathlineto{\pgfqpoint{2.572759in}{2.060935in}}%
\pgfpathlineto{\pgfqpoint{2.603807in}{2.042835in}}%
\pgfpathlineto{\pgfqpoint{2.559467in}{2.073375in}}%
\pgfpathlineto{\pgfqpoint{2.528430in}{2.091916in}}%
\pgfpathclose%
\pgfusepath{stroke,fill}%
\end{pgfscope}%
\begin{pgfscope}%
\pgfpathrectangle{\pgfqpoint{0.050000in}{0.050000in}}{\pgfqpoint{4.900000in}{4.900000in}}%
\pgfusepath{clip}%
\pgfsetbuttcap%
\pgfsetroundjoin%
\definecolor{currentfill}{rgb}{0.753664,0.753664,0.753664}%
\pgfsetfillcolor{currentfill}%
\pgfsetfillopacity{0.900000}%
\pgfsetlinewidth{0.501875pt}%
\definecolor{currentstroke}{rgb}{0.200000,0.200000,0.200000}%
\pgfsetstrokecolor{currentstroke}%
\pgfsetdash{}{0pt}%
\pgfpathmoveto{\pgfqpoint{1.756664in}{2.698971in}}%
\pgfpathlineto{\pgfqpoint{1.802585in}{2.655259in}}%
\pgfpathlineto{\pgfqpoint{1.832472in}{2.635896in}}%
\pgfpathlineto{\pgfqpoint{1.786544in}{2.679635in}}%
\pgfpathlineto{\pgfqpoint{1.756664in}{2.698971in}}%
\pgfpathclose%
\pgfusepath{stroke,fill}%
\end{pgfscope}%
\begin{pgfscope}%
\pgfpathrectangle{\pgfqpoint{0.050000in}{0.050000in}}{\pgfqpoint{4.900000in}{4.900000in}}%
\pgfusepath{clip}%
\pgfsetbuttcap%
\pgfsetroundjoin%
\definecolor{currentfill}{rgb}{0.992618,0.992618,0.992618}%
\pgfsetfillcolor{currentfill}%
\pgfsetfillopacity{0.900000}%
\pgfsetlinewidth{0.501875pt}%
\definecolor{currentstroke}{rgb}{0.200000,0.200000,0.200000}%
\pgfsetstrokecolor{currentstroke}%
\pgfsetdash{}{0pt}%
\pgfpathmoveto{\pgfqpoint{3.056029in}{1.858307in}}%
\pgfpathlineto{\pgfqpoint{3.099681in}{1.860469in}}%
\pgfpathlineto{\pgfqpoint{3.131552in}{1.842054in}}%
\pgfpathlineto{\pgfqpoint{3.087875in}{1.840062in}}%
\pgfpathlineto{\pgfqpoint{3.056029in}{1.858307in}}%
\pgfpathclose%
\pgfusepath{stroke,fill}%
\end{pgfscope}%
\begin{pgfscope}%
\pgfpathrectangle{\pgfqpoint{0.050000in}{0.050000in}}{\pgfqpoint{4.900000in}{4.900000in}}%
\pgfusepath{clip}%
\pgfsetbuttcap%
\pgfsetroundjoin%
\definecolor{currentfill}{rgb}{1.000000,1.000000,1.000000}%
\pgfsetfillcolor{currentfill}%
\pgfsetfillopacity{0.900000}%
\pgfsetlinewidth{0.501875pt}%
\definecolor{currentstroke}{rgb}{0.200000,0.200000,0.200000}%
\pgfsetstrokecolor{currentstroke}%
\pgfsetdash{}{0pt}%
\pgfpathmoveto{\pgfqpoint{3.596840in}{1.812787in}}%
\pgfpathlineto{\pgfqpoint{3.640006in}{1.825604in}}%
\pgfpathlineto{\pgfqpoint{3.672690in}{1.805899in}}%
\pgfpathlineto{\pgfqpoint{3.629496in}{1.793134in}}%
\pgfpathlineto{\pgfqpoint{3.596840in}{1.812787in}}%
\pgfpathclose%
\pgfusepath{stroke,fill}%
\end{pgfscope}%
\begin{pgfscope}%
\pgfpathrectangle{\pgfqpoint{0.050000in}{0.050000in}}{\pgfqpoint{4.900000in}{4.900000in}}%
\pgfusepath{clip}%
\pgfsetbuttcap%
\pgfsetroundjoin%
\definecolor{currentfill}{rgb}{0.878570,0.878570,0.878570}%
\pgfsetfillcolor{currentfill}%
\pgfsetfillopacity{0.900000}%
\pgfsetlinewidth{0.501875pt}%
\definecolor{currentstroke}{rgb}{0.200000,0.200000,0.200000}%
\pgfsetstrokecolor{currentstroke}%
\pgfsetdash{}{0pt}%
\pgfpathmoveto{\pgfqpoint{2.195684in}{2.345080in}}%
\pgfpathlineto{\pgfqpoint{2.240709in}{2.302683in}}%
\pgfpathlineto{\pgfqpoint{2.271261in}{2.283108in}}%
\pgfpathlineto{\pgfqpoint{2.226231in}{2.325357in}}%
\pgfpathlineto{\pgfqpoint{2.195684in}{2.345080in}}%
\pgfpathclose%
\pgfusepath{stroke,fill}%
\end{pgfscope}%
\begin{pgfscope}%
\pgfpathrectangle{\pgfqpoint{0.050000in}{0.050000in}}{\pgfqpoint{4.900000in}{4.900000in}}%
\pgfusepath{clip}%
\pgfsetbuttcap%
\pgfsetroundjoin%
\definecolor{currentfill}{rgb}{1.000000,1.000000,1.000000}%
\pgfsetfillcolor{currentfill}%
\pgfsetfillopacity{0.900000}%
\pgfsetlinewidth{0.501875pt}%
\definecolor{currentstroke}{rgb}{0.200000,0.200000,0.200000}%
\pgfsetstrokecolor{currentstroke}%
\pgfsetdash{}{0pt}%
\pgfpathmoveto{\pgfqpoint{3.401992in}{1.816215in}}%
\pgfpathlineto{\pgfqpoint{3.445353in}{1.827405in}}%
\pgfpathlineto{\pgfqpoint{3.477754in}{1.808003in}}%
\pgfpathlineto{\pgfqpoint{3.434364in}{1.796922in}}%
\pgfpathlineto{\pgfqpoint{3.401992in}{1.816215in}}%
\pgfpathclose%
\pgfusepath{stroke,fill}%
\end{pgfscope}%
\begin{pgfscope}%
\pgfpathrectangle{\pgfqpoint{0.050000in}{0.050000in}}{\pgfqpoint{4.900000in}{4.900000in}}%
\pgfusepath{clip}%
\pgfsetbuttcap%
\pgfsetroundjoin%
\definecolor{currentfill}{rgb}{0.994464,0.994464,0.994464}%
\pgfsetfillcolor{currentfill}%
\pgfsetfillopacity{0.900000}%
\pgfsetlinewidth{0.501875pt}%
\definecolor{currentstroke}{rgb}{0.200000,0.200000,0.200000}%
\pgfsetstrokecolor{currentstroke}%
\pgfsetdash{}{0pt}%
\pgfpathmoveto{\pgfqpoint{3.131552in}{1.842054in}}%
\pgfpathlineto{\pgfqpoint{3.175150in}{1.846951in}}%
\pgfpathlineto{\pgfqpoint{3.207145in}{1.828284in}}%
\pgfpathlineto{\pgfqpoint{3.163522in}{1.823566in}}%
\pgfpathlineto{\pgfqpoint{3.131552in}{1.842054in}}%
\pgfpathclose%
\pgfusepath{stroke,fill}%
\end{pgfscope}%
\begin{pgfscope}%
\pgfpathrectangle{\pgfqpoint{0.050000in}{0.050000in}}{\pgfqpoint{4.900000in}{4.900000in}}%
\pgfusepath{clip}%
\pgfsetbuttcap%
\pgfsetroundjoin%
\definecolor{currentfill}{rgb}{0.937993,0.937993,0.937993}%
\pgfsetfillcolor{currentfill}%
\pgfsetfillopacity{0.900000}%
\pgfsetlinewidth{0.501875pt}%
\definecolor{currentstroke}{rgb}{0.200000,0.200000,0.200000}%
\pgfsetstrokecolor{currentstroke}%
\pgfsetdash{}{0pt}%
\pgfpathmoveto{\pgfqpoint{2.452997in}{2.145725in}}%
\pgfpathlineto{\pgfqpoint{2.497488in}{2.110411in}}%
\pgfpathlineto{\pgfqpoint{2.528430in}{2.091916in}}%
\pgfpathlineto{\pgfqpoint{2.483932in}{2.126777in}}%
\pgfpathlineto{\pgfqpoint{2.452997in}{2.145725in}}%
\pgfpathclose%
\pgfusepath{stroke,fill}%
\end{pgfscope}%
\begin{pgfscope}%
\pgfpathrectangle{\pgfqpoint{0.050000in}{0.050000in}}{\pgfqpoint{4.900000in}{4.900000in}}%
\pgfusepath{clip}%
\pgfsetbuttcap%
\pgfsetroundjoin%
\definecolor{currentfill}{rgb}{0.805336,0.805336,0.805336}%
\pgfsetfillcolor{currentfill}%
\pgfsetfillopacity{0.900000}%
\pgfsetlinewidth{0.501875pt}%
\definecolor{currentstroke}{rgb}{0.200000,0.200000,0.200000}%
\pgfsetstrokecolor{currentstroke}%
\pgfsetdash{}{0pt}%
\pgfpathmoveto{\pgfqpoint{1.938271in}{2.553388in}}%
\pgfpathlineto{\pgfqpoint{1.983836in}{2.509959in}}%
\pgfpathlineto{\pgfqpoint{2.014008in}{2.490373in}}%
\pgfpathlineto{\pgfqpoint{1.968436in}{2.533826in}}%
\pgfpathlineto{\pgfqpoint{1.938271in}{2.553388in}}%
\pgfpathclose%
\pgfusepath{stroke,fill}%
\end{pgfscope}%
\begin{pgfscope}%
\pgfpathrectangle{\pgfqpoint{0.050000in}{0.050000in}}{\pgfqpoint{4.900000in}{4.900000in}}%
\pgfusepath{clip}%
\pgfsetbuttcap%
\pgfsetroundjoin%
\definecolor{currentfill}{rgb}{1.000000,1.000000,1.000000}%
\pgfsetfillcolor{currentfill}%
\pgfsetfillopacity{0.900000}%
\pgfsetlinewidth{0.501875pt}%
\definecolor{currentstroke}{rgb}{0.200000,0.200000,0.200000}%
\pgfsetstrokecolor{currentstroke}%
\pgfsetdash{}{0pt}%
\pgfpathmoveto{\pgfqpoint{3.672690in}{1.805899in}}%
\pgfpathlineto{\pgfqpoint{3.715796in}{1.818916in}}%
\pgfpathlineto{\pgfqpoint{3.748609in}{1.799114in}}%
\pgfpathlineto{\pgfqpoint{3.705474in}{1.786140in}}%
\pgfpathlineto{\pgfqpoint{3.672690in}{1.805899in}}%
\pgfpathclose%
\pgfusepath{stroke,fill}%
\end{pgfscope}%
\begin{pgfscope}%
\pgfpathrectangle{\pgfqpoint{0.050000in}{0.050000in}}{\pgfqpoint{4.900000in}{4.900000in}}%
\pgfusepath{clip}%
\pgfsetbuttcap%
\pgfsetroundjoin%
\definecolor{currentfill}{rgb}{0.720185,0.720185,0.720185}%
\pgfsetfillcolor{currentfill}%
\pgfsetfillopacity{0.900000}%
\pgfsetlinewidth{0.501875pt}%
\definecolor{currentstroke}{rgb}{0.200000,0.200000,0.200000}%
\pgfsetstrokecolor{currentstroke}%
\pgfsetdash{}{0pt}%
\pgfpathmoveto{\pgfqpoint{1.680772in}{2.762116in}}%
\pgfpathlineto{\pgfqpoint{1.726874in}{2.718250in}}%
\pgfpathlineto{\pgfqpoint{1.756664in}{2.698971in}}%
\pgfpathlineto{\pgfqpoint{1.710554in}{2.742864in}}%
\pgfpathlineto{\pgfqpoint{1.680772in}{2.762116in}}%
\pgfpathclose%
\pgfusepath{stroke,fill}%
\end{pgfscope}%
\begin{pgfscope}%
\pgfpathrectangle{\pgfqpoint{0.050000in}{0.050000in}}{\pgfqpoint{4.900000in}{4.900000in}}%
\pgfusepath{clip}%
\pgfsetbuttcap%
\pgfsetroundjoin%
\definecolor{currentfill}{rgb}{0.996309,0.996309,0.996309}%
\pgfsetfillcolor{currentfill}%
\pgfsetfillopacity{0.900000}%
\pgfsetlinewidth{0.501875pt}%
\definecolor{currentstroke}{rgb}{0.200000,0.200000,0.200000}%
\pgfsetstrokecolor{currentstroke}%
\pgfsetdash{}{0pt}%
\pgfpathmoveto{\pgfqpoint{3.207145in}{1.828284in}}%
\pgfpathlineto{\pgfqpoint{3.250690in}{1.835360in}}%
\pgfpathlineto{\pgfqpoint{3.282810in}{1.816458in}}%
\pgfpathlineto{\pgfqpoint{3.239239in}{1.809554in}}%
\pgfpathlineto{\pgfqpoint{3.207145in}{1.828284in}}%
\pgfpathclose%
\pgfusepath{stroke,fill}%
\end{pgfscope}%
\begin{pgfscope}%
\pgfpathrectangle{\pgfqpoint{0.050000in}{0.050000in}}{\pgfqpoint{4.900000in}{4.900000in}}%
\pgfusepath{clip}%
\pgfsetbuttcap%
\pgfsetroundjoin%
\definecolor{currentfill}{rgb}{1.000000,1.000000,1.000000}%
\pgfsetfillcolor{currentfill}%
\pgfsetfillopacity{0.900000}%
\pgfsetlinewidth{0.501875pt}%
\definecolor{currentstroke}{rgb}{0.200000,0.200000,0.200000}%
\pgfsetstrokecolor{currentstroke}%
\pgfsetdash{}{0pt}%
\pgfpathmoveto{\pgfqpoint{3.477754in}{1.808003in}}%
\pgfpathlineto{\pgfqpoint{3.521061in}{1.819907in}}%
\pgfpathlineto{\pgfqpoint{3.553589in}{1.800371in}}%
\pgfpathlineto{\pgfqpoint{3.510253in}{1.788553in}}%
\pgfpathlineto{\pgfqpoint{3.477754in}{1.808003in}}%
\pgfpathclose%
\pgfusepath{stroke,fill}%
\end{pgfscope}%
\begin{pgfscope}%
\pgfpathrectangle{\pgfqpoint{0.050000in}{0.050000in}}{\pgfqpoint{4.900000in}{4.900000in}}%
\pgfusepath{clip}%
\pgfsetbuttcap%
\pgfsetroundjoin%
\definecolor{currentfill}{rgb}{0.918185,0.918185,0.918185}%
\pgfsetfillcolor{currentfill}%
\pgfsetfillopacity{0.900000}%
\pgfsetlinewidth{0.501875pt}%
\definecolor{currentstroke}{rgb}{0.200000,0.200000,0.200000}%
\pgfsetstrokecolor{currentstroke}%
\pgfsetdash{}{0pt}%
\pgfpathmoveto{\pgfqpoint{2.377492in}{2.203298in}}%
\pgfpathlineto{\pgfqpoint{2.422157in}{2.164651in}}%
\pgfpathlineto{\pgfqpoint{2.452997in}{2.145725in}}%
\pgfpathlineto{\pgfqpoint{2.408327in}{2.183976in}}%
\pgfpathlineto{\pgfqpoint{2.377492in}{2.203298in}}%
\pgfpathclose%
\pgfusepath{stroke,fill}%
\end{pgfscope}%
\begin{pgfscope}%
\pgfpathrectangle{\pgfqpoint{0.050000in}{0.050000in}}{\pgfqpoint{4.900000in}{4.900000in}}%
\pgfusepath{clip}%
\pgfsetbuttcap%
\pgfsetroundjoin%
\definecolor{currentfill}{rgb}{0.855932,0.855932,0.855932}%
\pgfsetfillcolor{currentfill}%
\pgfsetfillopacity{0.900000}%
\pgfsetlinewidth{0.501875pt}%
\definecolor{currentstroke}{rgb}{0.200000,0.200000,0.200000}%
\pgfsetstrokecolor{currentstroke}%
\pgfsetdash{}{0pt}%
\pgfpathmoveto{\pgfqpoint{2.120020in}{2.407761in}}%
\pgfpathlineto{\pgfqpoint{2.165228in}{2.364775in}}%
\pgfpathlineto{\pgfqpoint{2.195684in}{2.345080in}}%
\pgfpathlineto{\pgfqpoint{2.150471in}{2.388017in}}%
\pgfpathlineto{\pgfqpoint{2.120020in}{2.407761in}}%
\pgfpathclose%
\pgfusepath{stroke,fill}%
\end{pgfscope}%
\begin{pgfscope}%
\pgfpathrectangle{\pgfqpoint{0.050000in}{0.050000in}}{\pgfqpoint{4.900000in}{4.900000in}}%
\pgfusepath{clip}%
\pgfsetbuttcap%
\pgfsetroundjoin%
\definecolor{currentfill}{rgb}{1.000000,1.000000,1.000000}%
\pgfsetfillcolor{currentfill}%
\pgfsetfillopacity{0.900000}%
\pgfsetlinewidth{0.501875pt}%
\definecolor{currentstroke}{rgb}{0.200000,0.200000,0.200000}%
\pgfsetstrokecolor{currentstroke}%
\pgfsetdash{}{0pt}%
\pgfpathmoveto{\pgfqpoint{3.748609in}{1.799114in}}%
\pgfpathlineto{\pgfqpoint{3.791654in}{1.812210in}}%
\pgfpathlineto{\pgfqpoint{3.824594in}{1.792312in}}%
\pgfpathlineto{\pgfqpoint{3.781521in}{1.779255in}}%
\pgfpathlineto{\pgfqpoint{3.748609in}{1.799114in}}%
\pgfpathclose%
\pgfusepath{stroke,fill}%
\end{pgfscope}%
\begin{pgfscope}%
\pgfpathrectangle{\pgfqpoint{0.050000in}{0.050000in}}{\pgfqpoint{4.900000in}{4.900000in}}%
\pgfusepath{clip}%
\pgfsetbuttcap%
\pgfsetroundjoin%
\definecolor{currentfill}{rgb}{0.781223,0.781223,0.781223}%
\pgfsetfillcolor{currentfill}%
\pgfsetfillopacity{0.900000}%
\pgfsetlinewidth{0.501875pt}%
\definecolor{currentstroke}{rgb}{0.200000,0.200000,0.200000}%
\pgfsetstrokecolor{currentstroke}%
\pgfsetdash{}{0pt}%
\pgfpathmoveto{\pgfqpoint{1.862450in}{2.616476in}}%
\pgfpathlineto{\pgfqpoint{1.908196in}{2.572892in}}%
\pgfpathlineto{\pgfqpoint{1.938271in}{2.553388in}}%
\pgfpathlineto{\pgfqpoint{1.892518in}{2.596997in}}%
\pgfpathlineto{\pgfqpoint{1.862450in}{2.616476in}}%
\pgfpathclose%
\pgfusepath{stroke,fill}%
\end{pgfscope}%
\begin{pgfscope}%
\pgfpathrectangle{\pgfqpoint{0.050000in}{0.050000in}}{\pgfqpoint{4.900000in}{4.900000in}}%
\pgfusepath{clip}%
\pgfsetbuttcap%
\pgfsetroundjoin%
\definecolor{currentfill}{rgb}{0.998155,0.998155,0.998155}%
\pgfsetfillcolor{currentfill}%
\pgfsetfillopacity{0.900000}%
\pgfsetlinewidth{0.501875pt}%
\definecolor{currentstroke}{rgb}{0.200000,0.200000,0.200000}%
\pgfsetstrokecolor{currentstroke}%
\pgfsetdash{}{0pt}%
\pgfpathmoveto{\pgfqpoint{3.282810in}{1.816458in}}%
\pgfpathlineto{\pgfqpoint{3.326304in}{1.825237in}}%
\pgfpathlineto{\pgfqpoint{3.358550in}{1.806127in}}%
\pgfpathlineto{\pgfqpoint{3.315029in}{1.797501in}}%
\pgfpathlineto{\pgfqpoint{3.282810in}{1.816458in}}%
\pgfpathclose%
\pgfusepath{stroke,fill}%
\end{pgfscope}%
\begin{pgfscope}%
\pgfpathrectangle{\pgfqpoint{0.050000in}{0.050000in}}{\pgfqpoint{4.900000in}{4.900000in}}%
\pgfusepath{clip}%
\pgfsetbuttcap%
\pgfsetroundjoin%
\definecolor{currentfill}{rgb}{0.974164,0.974164,0.974164}%
\pgfsetfillcolor{currentfill}%
\pgfsetfillopacity{0.900000}%
\pgfsetlinewidth{0.501875pt}%
\definecolor{currentstroke}{rgb}{0.200000,0.200000,0.200000}%
\pgfsetstrokecolor{currentstroke}%
\pgfsetdash{}{0pt}%
\pgfpathmoveto{\pgfqpoint{2.785840in}{1.943249in}}%
\pgfpathlineto{\pgfqpoint{2.829805in}{1.928377in}}%
\pgfpathlineto{\pgfqpoint{2.861290in}{1.910627in}}%
\pgfpathlineto{\pgfqpoint{2.817305in}{1.925380in}}%
\pgfpathlineto{\pgfqpoint{2.785840in}{1.943249in}}%
\pgfpathclose%
\pgfusepath{stroke,fill}%
\end{pgfscope}%
\begin{pgfscope}%
\pgfpathrectangle{\pgfqpoint{0.050000in}{0.050000in}}{\pgfqpoint{4.900000in}{4.900000in}}%
\pgfusepath{clip}%
\pgfsetbuttcap%
\pgfsetroundjoin%
\definecolor{currentfill}{rgb}{0.979700,0.979700,0.979700}%
\pgfsetfillcolor{currentfill}%
\pgfsetfillopacity{0.900000}%
\pgfsetlinewidth{0.501875pt}%
\definecolor{currentstroke}{rgb}{0.200000,0.200000,0.200000}%
\pgfsetstrokecolor{currentstroke}%
\pgfsetdash{}{0pt}%
\pgfpathmoveto{\pgfqpoint{2.861290in}{1.910627in}}%
\pgfpathlineto{\pgfqpoint{2.905166in}{1.900794in}}%
\pgfpathlineto{\pgfqpoint{2.936769in}{1.882948in}}%
\pgfpathlineto{\pgfqpoint{2.892872in}{1.892770in}}%
\pgfpathlineto{\pgfqpoint{2.861290in}{1.910627in}}%
\pgfpathclose%
\pgfusepath{stroke,fill}%
\end{pgfscope}%
\begin{pgfscope}%
\pgfpathrectangle{\pgfqpoint{0.050000in}{0.050000in}}{\pgfqpoint{4.900000in}{4.900000in}}%
\pgfusepath{clip}%
\pgfsetbuttcap%
\pgfsetroundjoin%
\definecolor{currentfill}{rgb}{0.966782,0.966782,0.966782}%
\pgfsetfillcolor{currentfill}%
\pgfsetfillopacity{0.900000}%
\pgfsetlinewidth{0.501875pt}%
\definecolor{currentstroke}{rgb}{0.200000,0.200000,0.200000}%
\pgfsetstrokecolor{currentstroke}%
\pgfsetdash{}{0pt}%
\pgfpathmoveto{\pgfqpoint{2.710400in}{1.981215in}}%
\pgfpathlineto{\pgfqpoint{2.754472in}{1.961012in}}%
\pgfpathlineto{\pgfqpoint{2.785840in}{1.943249in}}%
\pgfpathlineto{\pgfqpoint{2.741751in}{1.963216in}}%
\pgfpathlineto{\pgfqpoint{2.710400in}{1.981215in}}%
\pgfpathclose%
\pgfusepath{stroke,fill}%
\end{pgfscope}%
\begin{pgfscope}%
\pgfpathrectangle{\pgfqpoint{0.050000in}{0.050000in}}{\pgfqpoint{4.900000in}{4.900000in}}%
\pgfusepath{clip}%
\pgfsetbuttcap%
\pgfsetroundjoin%
\definecolor{currentfill}{rgb}{0.686597,0.686597,0.686597}%
\pgfsetfillcolor{currentfill}%
\pgfsetfillopacity{0.900000}%
\pgfsetlinewidth{0.501875pt}%
\definecolor{currentstroke}{rgb}{0.200000,0.200000,0.200000}%
\pgfsetstrokecolor{currentstroke}%
\pgfsetdash{}{0pt}%
\pgfpathmoveto{\pgfqpoint{1.604795in}{2.825332in}}%
\pgfpathlineto{\pgfqpoint{1.651079in}{2.781311in}}%
\pgfpathlineto{\pgfqpoint{1.680772in}{2.762116in}}%
\pgfpathlineto{\pgfqpoint{1.634479in}{2.806164in}}%
\pgfpathlineto{\pgfqpoint{1.604795in}{2.825332in}}%
\pgfpathclose%
\pgfusepath{stroke,fill}%
\end{pgfscope}%
\begin{pgfscope}%
\pgfpathrectangle{\pgfqpoint{0.050000in}{0.050000in}}{\pgfqpoint{4.900000in}{4.900000in}}%
\pgfusepath{clip}%
\pgfsetbuttcap%
\pgfsetroundjoin%
\definecolor{currentfill}{rgb}{1.000000,1.000000,1.000000}%
\pgfsetfillcolor{currentfill}%
\pgfsetfillopacity{0.900000}%
\pgfsetlinewidth{0.501875pt}%
\definecolor{currentstroke}{rgb}{0.200000,0.200000,0.200000}%
\pgfsetstrokecolor{currentstroke}%
\pgfsetdash{}{0pt}%
\pgfpathmoveto{\pgfqpoint{3.553589in}{1.800371in}}%
\pgfpathlineto{\pgfqpoint{3.596840in}{1.812787in}}%
\pgfpathlineto{\pgfqpoint{3.629496in}{1.793134in}}%
\pgfpathlineto{\pgfqpoint{3.586216in}{1.780784in}}%
\pgfpathlineto{\pgfqpoint{3.553589in}{1.800371in}}%
\pgfpathclose%
\pgfusepath{stroke,fill}%
\end{pgfscope}%
\begin{pgfscope}%
\pgfpathrectangle{\pgfqpoint{0.050000in}{0.050000in}}{\pgfqpoint{4.900000in}{4.900000in}}%
\pgfusepath{clip}%
\pgfsetbuttcap%
\pgfsetroundjoin%
\definecolor{currentfill}{rgb}{0.985236,0.985236,0.985236}%
\pgfsetfillcolor{currentfill}%
\pgfsetfillopacity{0.900000}%
\pgfsetlinewidth{0.501875pt}%
\definecolor{currentstroke}{rgb}{0.200000,0.200000,0.200000}%
\pgfsetstrokecolor{currentstroke}%
\pgfsetdash{}{0pt}%
\pgfpathmoveto{\pgfqpoint{2.936769in}{1.882948in}}%
\pgfpathlineto{\pgfqpoint{2.980570in}{1.877662in}}%
\pgfpathlineto{\pgfqpoint{3.012294in}{1.859641in}}%
\pgfpathlineto{\pgfqpoint{2.968471in}{1.864999in}}%
\pgfpathlineto{\pgfqpoint{2.936769in}{1.882948in}}%
\pgfpathclose%
\pgfusepath{stroke,fill}%
\end{pgfscope}%
\begin{pgfscope}%
\pgfpathrectangle{\pgfqpoint{0.050000in}{0.050000in}}{\pgfqpoint{4.900000in}{4.900000in}}%
\pgfusepath{clip}%
\pgfsetbuttcap%
\pgfsetroundjoin%
\definecolor{currentfill}{rgb}{0.898378,0.898378,0.898378}%
\pgfsetfillcolor{currentfill}%
\pgfsetfillopacity{0.900000}%
\pgfsetlinewidth{0.501875pt}%
\definecolor{currentstroke}{rgb}{0.200000,0.200000,0.200000}%
\pgfsetstrokecolor{currentstroke}%
\pgfsetdash{}{0pt}%
\pgfpathmoveto{\pgfqpoint{2.301905in}{2.263522in}}%
\pgfpathlineto{\pgfqpoint{2.346751in}{2.222610in}}%
\pgfpathlineto{\pgfqpoint{2.377492in}{2.203298in}}%
\pgfpathlineto{\pgfqpoint{2.332641in}{2.243924in}}%
\pgfpathlineto{\pgfqpoint{2.301905in}{2.263522in}}%
\pgfpathclose%
\pgfusepath{stroke,fill}%
\end{pgfscope}%
\begin{pgfscope}%
\pgfpathrectangle{\pgfqpoint{0.050000in}{0.050000in}}{\pgfqpoint{4.900000in}{4.900000in}}%
\pgfusepath{clip}%
\pgfsetbuttcap%
\pgfsetroundjoin%
\definecolor{currentfill}{rgb}{0.959400,0.959400,0.959400}%
\pgfsetfillcolor{currentfill}%
\pgfsetfillopacity{0.900000}%
\pgfsetlinewidth{0.501875pt}%
\definecolor{currentstroke}{rgb}{0.200000,0.200000,0.200000}%
\pgfsetstrokecolor{currentstroke}%
\pgfsetdash{}{0pt}%
\pgfpathmoveto{\pgfqpoint{2.634950in}{2.024661in}}%
\pgfpathlineto{\pgfqpoint{2.679146in}{1.999120in}}%
\pgfpathlineto{\pgfqpoint{2.710400in}{1.981215in}}%
\pgfpathlineto{\pgfqpoint{2.666190in}{2.006411in}}%
\pgfpathlineto{\pgfqpoint{2.634950in}{2.024661in}}%
\pgfpathclose%
\pgfusepath{stroke,fill}%
\end{pgfscope}%
\begin{pgfscope}%
\pgfpathrectangle{\pgfqpoint{0.050000in}{0.050000in}}{\pgfqpoint{4.900000in}{4.900000in}}%
\pgfusepath{clip}%
\pgfsetbuttcap%
\pgfsetroundjoin%
\definecolor{currentfill}{rgb}{0.988927,0.988927,0.988927}%
\pgfsetfillcolor{currentfill}%
\pgfsetfillopacity{0.900000}%
\pgfsetlinewidth{0.501875pt}%
\definecolor{currentstroke}{rgb}{0.200000,0.200000,0.200000}%
\pgfsetstrokecolor{currentstroke}%
\pgfsetdash{}{0pt}%
\pgfpathmoveto{\pgfqpoint{3.012294in}{1.859641in}}%
\pgfpathlineto{\pgfqpoint{3.056029in}{1.858307in}}%
\pgfpathlineto{\pgfqpoint{3.087875in}{1.840062in}}%
\pgfpathlineto{\pgfqpoint{3.044116in}{1.841525in}}%
\pgfpathlineto{\pgfqpoint{3.012294in}{1.859641in}}%
\pgfpathclose%
\pgfusepath{stroke,fill}%
\end{pgfscope}%
\begin{pgfscope}%
\pgfpathrectangle{\pgfqpoint{0.050000in}{0.050000in}}{\pgfqpoint{4.900000in}{4.900000in}}%
\pgfusepath{clip}%
\pgfsetbuttcap%
\pgfsetroundjoin%
\definecolor{currentfill}{rgb}{0.832895,0.832895,0.832895}%
\pgfsetfillcolor{currentfill}%
\pgfsetfillopacity{0.900000}%
\pgfsetlinewidth{0.501875pt}%
\definecolor{currentstroke}{rgb}{0.200000,0.200000,0.200000}%
\pgfsetstrokecolor{currentstroke}%
\pgfsetdash{}{0pt}%
\pgfpathmoveto{\pgfqpoint{2.044271in}{2.470732in}}%
\pgfpathlineto{\pgfqpoint{2.089660in}{2.427460in}}%
\pgfpathlineto{\pgfqpoint{2.120020in}{2.407761in}}%
\pgfpathlineto{\pgfqpoint{2.074624in}{2.451035in}}%
\pgfpathlineto{\pgfqpoint{2.044271in}{2.470732in}}%
\pgfpathclose%
\pgfusepath{stroke,fill}%
\end{pgfscope}%
\begin{pgfscope}%
\pgfpathrectangle{\pgfqpoint{0.050000in}{0.050000in}}{\pgfqpoint{4.900000in}{4.900000in}}%
\pgfusepath{clip}%
\pgfsetbuttcap%
\pgfsetroundjoin%
\definecolor{currentfill}{rgb}{0.998155,0.998155,0.998155}%
\pgfsetfillcolor{currentfill}%
\pgfsetfillopacity{0.900000}%
\pgfsetlinewidth{0.501875pt}%
\definecolor{currentstroke}{rgb}{0.200000,0.200000,0.200000}%
\pgfsetstrokecolor{currentstroke}%
\pgfsetdash{}{0pt}%
\pgfpathmoveto{\pgfqpoint{3.358550in}{1.806127in}}%
\pgfpathlineto{\pgfqpoint{3.401992in}{1.816215in}}%
\pgfpathlineto{\pgfqpoint{3.434364in}{1.796922in}}%
\pgfpathlineto{\pgfqpoint{3.390895in}{1.786964in}}%
\pgfpathlineto{\pgfqpoint{3.358550in}{1.806127in}}%
\pgfpathclose%
\pgfusepath{stroke,fill}%
\end{pgfscope}%
\begin{pgfscope}%
\pgfpathrectangle{\pgfqpoint{0.050000in}{0.050000in}}{\pgfqpoint{4.900000in}{4.900000in}}%
\pgfusepath{clip}%
\pgfsetbuttcap%
\pgfsetroundjoin%
\definecolor{currentfill}{rgb}{0.948328,0.948328,0.948328}%
\pgfsetfillcolor{currentfill}%
\pgfsetfillopacity{0.900000}%
\pgfsetlinewidth{0.501875pt}%
\definecolor{currentstroke}{rgb}{0.200000,0.200000,0.200000}%
\pgfsetstrokecolor{currentstroke}%
\pgfsetdash{}{0pt}%
\pgfpathmoveto{\pgfqpoint{2.559467in}{2.073375in}}%
\pgfpathlineto{\pgfqpoint{2.603807in}{2.042835in}}%
\pgfpathlineto{\pgfqpoint{2.634950in}{2.024661in}}%
\pgfpathlineto{\pgfqpoint{2.590600in}{2.054783in}}%
\pgfpathlineto{\pgfqpoint{2.559467in}{2.073375in}}%
\pgfpathclose%
\pgfusepath{stroke,fill}%
\end{pgfscope}%
\begin{pgfscope}%
\pgfpathrectangle{\pgfqpoint{0.050000in}{0.050000in}}{\pgfqpoint{4.900000in}{4.900000in}}%
\pgfusepath{clip}%
\pgfsetbuttcap%
\pgfsetroundjoin%
\definecolor{currentfill}{rgb}{0.753664,0.753664,0.753664}%
\pgfsetfillcolor{currentfill}%
\pgfsetfillopacity{0.900000}%
\pgfsetlinewidth{0.501875pt}%
\definecolor{currentstroke}{rgb}{0.200000,0.200000,0.200000}%
\pgfsetstrokecolor{currentstroke}%
\pgfsetdash{}{0pt}%
\pgfpathmoveto{\pgfqpoint{1.786544in}{2.679635in}}%
\pgfpathlineto{\pgfqpoint{1.832472in}{2.635896in}}%
\pgfpathlineto{\pgfqpoint{1.862450in}{2.616476in}}%
\pgfpathlineto{\pgfqpoint{1.816514in}{2.660240in}}%
\pgfpathlineto{\pgfqpoint{1.786544in}{2.679635in}}%
\pgfpathclose%
\pgfusepath{stroke,fill}%
\end{pgfscope}%
\begin{pgfscope}%
\pgfpathrectangle{\pgfqpoint{0.050000in}{0.050000in}}{\pgfqpoint{4.900000in}{4.900000in}}%
\pgfusepath{clip}%
\pgfsetbuttcap%
\pgfsetroundjoin%
\definecolor{currentfill}{rgb}{1.000000,1.000000,1.000000}%
\pgfsetfillcolor{currentfill}%
\pgfsetfillopacity{0.900000}%
\pgfsetlinewidth{0.501875pt}%
\definecolor{currentstroke}{rgb}{0.200000,0.200000,0.200000}%
\pgfsetstrokecolor{currentstroke}%
\pgfsetdash{}{0pt}%
\pgfpathmoveto{\pgfqpoint{3.629496in}{1.793134in}}%
\pgfpathlineto{\pgfqpoint{3.672690in}{1.805899in}}%
\pgfpathlineto{\pgfqpoint{3.705474in}{1.786140in}}%
\pgfpathlineto{\pgfqpoint{3.662251in}{1.773428in}}%
\pgfpathlineto{\pgfqpoint{3.629496in}{1.793134in}}%
\pgfpathclose%
\pgfusepath{stroke,fill}%
\end{pgfscope}%
\begin{pgfscope}%
\pgfpathrectangle{\pgfqpoint{0.050000in}{0.050000in}}{\pgfqpoint{4.900000in}{4.900000in}}%
\pgfusepath{clip}%
\pgfsetbuttcap%
\pgfsetroundjoin%
\definecolor{currentfill}{rgb}{0.992618,0.992618,0.992618}%
\pgfsetfillcolor{currentfill}%
\pgfsetfillopacity{0.900000}%
\pgfsetlinewidth{0.501875pt}%
\definecolor{currentstroke}{rgb}{0.200000,0.200000,0.200000}%
\pgfsetstrokecolor{currentstroke}%
\pgfsetdash{}{0pt}%
\pgfpathmoveto{\pgfqpoint{3.087875in}{1.840062in}}%
\pgfpathlineto{\pgfqpoint{3.131552in}{1.842054in}}%
\pgfpathlineto{\pgfqpoint{3.163522in}{1.823566in}}%
\pgfpathlineto{\pgfqpoint{3.119820in}{1.821734in}}%
\pgfpathlineto{\pgfqpoint{3.087875in}{1.840062in}}%
\pgfpathclose%
\pgfusepath{stroke,fill}%
\end{pgfscope}%
\begin{pgfscope}%
\pgfpathrectangle{\pgfqpoint{0.050000in}{0.050000in}}{\pgfqpoint{4.900000in}{4.900000in}}%
\pgfusepath{clip}%
\pgfsetbuttcap%
\pgfsetroundjoin%
\definecolor{currentfill}{rgb}{0.878570,0.878570,0.878570}%
\pgfsetfillcolor{currentfill}%
\pgfsetfillopacity{0.900000}%
\pgfsetlinewidth{0.501875pt}%
\definecolor{currentstroke}{rgb}{0.200000,0.200000,0.200000}%
\pgfsetstrokecolor{currentstroke}%
\pgfsetdash{}{0pt}%
\pgfpathmoveto{\pgfqpoint{2.226231in}{2.325357in}}%
\pgfpathlineto{\pgfqpoint{2.271261in}{2.283108in}}%
\pgfpathlineto{\pgfqpoint{2.301905in}{2.263522in}}%
\pgfpathlineto{\pgfqpoint{2.256871in}{2.305608in}}%
\pgfpathlineto{\pgfqpoint{2.226231in}{2.325357in}}%
\pgfpathclose%
\pgfusepath{stroke,fill}%
\end{pgfscope}%
\begin{pgfscope}%
\pgfpathrectangle{\pgfqpoint{0.050000in}{0.050000in}}{\pgfqpoint{4.900000in}{4.900000in}}%
\pgfusepath{clip}%
\pgfsetbuttcap%
\pgfsetroundjoin%
\definecolor{currentfill}{rgb}{1.000000,1.000000,1.000000}%
\pgfsetfillcolor{currentfill}%
\pgfsetfillopacity{0.900000}%
\pgfsetlinewidth{0.501875pt}%
\definecolor{currentstroke}{rgb}{0.200000,0.200000,0.200000}%
\pgfsetstrokecolor{currentstroke}%
\pgfsetdash{}{0pt}%
\pgfpathmoveto{\pgfqpoint{3.434364in}{1.796922in}}%
\pgfpathlineto{\pgfqpoint{3.477754in}{1.808003in}}%
\pgfpathlineto{\pgfqpoint{3.510253in}{1.788553in}}%
\pgfpathlineto{\pgfqpoint{3.466836in}{1.777579in}}%
\pgfpathlineto{\pgfqpoint{3.434364in}{1.796922in}}%
\pgfpathclose%
\pgfusepath{stroke,fill}%
\end{pgfscope}%
\begin{pgfscope}%
\pgfpathrectangle{\pgfqpoint{0.050000in}{0.050000in}}{\pgfqpoint{4.900000in}{4.900000in}}%
\pgfusepath{clip}%
\pgfsetbuttcap%
\pgfsetroundjoin%
\definecolor{currentfill}{rgb}{0.935163,0.935163,0.935163}%
\pgfsetfillcolor{currentfill}%
\pgfsetfillopacity{0.900000}%
\pgfsetlinewidth{0.501875pt}%
\definecolor{currentstroke}{rgb}{0.200000,0.200000,0.200000}%
\pgfsetstrokecolor{currentstroke}%
\pgfsetdash{}{0pt}%
\pgfpathmoveto{\pgfqpoint{2.483932in}{2.126777in}}%
\pgfpathlineto{\pgfqpoint{2.528430in}{2.091916in}}%
\pgfpathlineto{\pgfqpoint{2.559467in}{2.073375in}}%
\pgfpathlineto{\pgfqpoint{2.514961in}{2.107801in}}%
\pgfpathlineto{\pgfqpoint{2.483932in}{2.126777in}}%
\pgfpathclose%
\pgfusepath{stroke,fill}%
\end{pgfscope}%
\begin{pgfscope}%
\pgfpathrectangle{\pgfqpoint{0.050000in}{0.050000in}}{\pgfqpoint{4.900000in}{4.900000in}}%
\pgfusepath{clip}%
\pgfsetbuttcap%
\pgfsetroundjoin%
\definecolor{currentfill}{rgb}{0.994464,0.994464,0.994464}%
\pgfsetfillcolor{currentfill}%
\pgfsetfillopacity{0.900000}%
\pgfsetlinewidth{0.501875pt}%
\definecolor{currentstroke}{rgb}{0.200000,0.200000,0.200000}%
\pgfsetstrokecolor{currentstroke}%
\pgfsetdash{}{0pt}%
\pgfpathmoveto{\pgfqpoint{3.163522in}{1.823566in}}%
\pgfpathlineto{\pgfqpoint{3.207145in}{1.828284in}}%
\pgfpathlineto{\pgfqpoint{3.239239in}{1.809554in}}%
\pgfpathlineto{\pgfqpoint{3.195590in}{1.805005in}}%
\pgfpathlineto{\pgfqpoint{3.163522in}{1.823566in}}%
\pgfpathclose%
\pgfusepath{stroke,fill}%
\end{pgfscope}%
\begin{pgfscope}%
\pgfpathrectangle{\pgfqpoint{0.050000in}{0.050000in}}{\pgfqpoint{4.900000in}{4.900000in}}%
\pgfusepath{clip}%
\pgfsetbuttcap%
\pgfsetroundjoin%
\definecolor{currentfill}{rgb}{0.805336,0.805336,0.805336}%
\pgfsetfillcolor{currentfill}%
\pgfsetfillopacity{0.900000}%
\pgfsetlinewidth{0.501875pt}%
\definecolor{currentstroke}{rgb}{0.200000,0.200000,0.200000}%
\pgfsetstrokecolor{currentstroke}%
\pgfsetdash{}{0pt}%
\pgfpathmoveto{\pgfqpoint{1.968436in}{2.533826in}}%
\pgfpathlineto{\pgfqpoint{2.014008in}{2.490373in}}%
\pgfpathlineto{\pgfqpoint{2.044271in}{2.470732in}}%
\pgfpathlineto{\pgfqpoint{1.998693in}{2.514205in}}%
\pgfpathlineto{\pgfqpoint{1.968436in}{2.533826in}}%
\pgfpathclose%
\pgfusepath{stroke,fill}%
\end{pgfscope}%
\begin{pgfscope}%
\pgfpathrectangle{\pgfqpoint{0.050000in}{0.050000in}}{\pgfqpoint{4.900000in}{4.900000in}}%
\pgfusepath{clip}%
\pgfsetbuttcap%
\pgfsetroundjoin%
\definecolor{currentfill}{rgb}{1.000000,1.000000,1.000000}%
\pgfsetfillcolor{currentfill}%
\pgfsetfillopacity{0.900000}%
\pgfsetlinewidth{0.501875pt}%
\definecolor{currentstroke}{rgb}{0.200000,0.200000,0.200000}%
\pgfsetstrokecolor{currentstroke}%
\pgfsetdash{}{0pt}%
\pgfpathmoveto{\pgfqpoint{3.705474in}{1.786140in}}%
\pgfpathlineto{\pgfqpoint{3.748609in}{1.799114in}}%
\pgfpathlineto{\pgfqpoint{3.781521in}{1.779255in}}%
\pgfpathlineto{\pgfqpoint{3.738358in}{1.766326in}}%
\pgfpathlineto{\pgfqpoint{3.705474in}{1.786140in}}%
\pgfpathclose%
\pgfusepath{stroke,fill}%
\end{pgfscope}%
\begin{pgfscope}%
\pgfpathrectangle{\pgfqpoint{0.050000in}{0.050000in}}{\pgfqpoint{4.900000in}{4.900000in}}%
\pgfusepath{clip}%
\pgfsetbuttcap%
\pgfsetroundjoin%
\definecolor{currentfill}{rgb}{0.724983,0.724983,0.724983}%
\pgfsetfillcolor{currentfill}%
\pgfsetfillopacity{0.900000}%
\pgfsetlinewidth{0.501875pt}%
\definecolor{currentstroke}{rgb}{0.200000,0.200000,0.200000}%
\pgfsetstrokecolor{currentstroke}%
\pgfsetdash{}{0pt}%
\pgfpathmoveto{\pgfqpoint{1.710554in}{2.742864in}}%
\pgfpathlineto{\pgfqpoint{1.756664in}{2.698971in}}%
\pgfpathlineto{\pgfqpoint{1.786544in}{2.679635in}}%
\pgfpathlineto{\pgfqpoint{1.740426in}{2.723554in}}%
\pgfpathlineto{\pgfqpoint{1.710554in}{2.742864in}}%
\pgfpathclose%
\pgfusepath{stroke,fill}%
\end{pgfscope}%
\begin{pgfscope}%
\pgfpathrectangle{\pgfqpoint{0.050000in}{0.050000in}}{\pgfqpoint{4.900000in}{4.900000in}}%
\pgfusepath{clip}%
\pgfsetbuttcap%
\pgfsetroundjoin%
\definecolor{currentfill}{rgb}{0.996309,0.996309,0.996309}%
\pgfsetfillcolor{currentfill}%
\pgfsetfillopacity{0.900000}%
\pgfsetlinewidth{0.501875pt}%
\definecolor{currentstroke}{rgb}{0.200000,0.200000,0.200000}%
\pgfsetstrokecolor{currentstroke}%
\pgfsetdash{}{0pt}%
\pgfpathmoveto{\pgfqpoint{3.239239in}{1.809554in}}%
\pgfpathlineto{\pgfqpoint{3.282810in}{1.816458in}}%
\pgfpathlineto{\pgfqpoint{3.315029in}{1.797501in}}%
\pgfpathlineto{\pgfqpoint{3.271432in}{1.790760in}}%
\pgfpathlineto{\pgfqpoint{3.239239in}{1.809554in}}%
\pgfpathclose%
\pgfusepath{stroke,fill}%
\end{pgfscope}%
\begin{pgfscope}%
\pgfpathrectangle{\pgfqpoint{0.050000in}{0.050000in}}{\pgfqpoint{4.900000in}{4.900000in}}%
\pgfusepath{clip}%
\pgfsetbuttcap%
\pgfsetroundjoin%
\definecolor{currentfill}{rgb}{1.000000,1.000000,1.000000}%
\pgfsetfillcolor{currentfill}%
\pgfsetfillopacity{0.900000}%
\pgfsetlinewidth{0.501875pt}%
\definecolor{currentstroke}{rgb}{0.200000,0.200000,0.200000}%
\pgfsetstrokecolor{currentstroke}%
\pgfsetdash{}{0pt}%
\pgfpathmoveto{\pgfqpoint{3.510253in}{1.788553in}}%
\pgfpathlineto{\pgfqpoint{3.553589in}{1.800371in}}%
\pgfpathlineto{\pgfqpoint{3.586216in}{1.780784in}}%
\pgfpathlineto{\pgfqpoint{3.542852in}{1.769052in}}%
\pgfpathlineto{\pgfqpoint{3.510253in}{1.788553in}}%
\pgfpathclose%
\pgfusepath{stroke,fill}%
\end{pgfscope}%
\begin{pgfscope}%
\pgfpathrectangle{\pgfqpoint{0.050000in}{0.050000in}}{\pgfqpoint{4.900000in}{4.900000in}}%
\pgfusepath{clip}%
\pgfsetbuttcap%
\pgfsetroundjoin%
\definecolor{currentfill}{rgb}{0.918185,0.918185,0.918185}%
\pgfsetfillcolor{currentfill}%
\pgfsetfillopacity{0.900000}%
\pgfsetlinewidth{0.501875pt}%
\definecolor{currentstroke}{rgb}{0.200000,0.200000,0.200000}%
\pgfsetstrokecolor{currentstroke}%
\pgfsetdash{}{0pt}%
\pgfpathmoveto{\pgfqpoint{2.408327in}{2.183976in}}%
\pgfpathlineto{\pgfqpoint{2.452997in}{2.145725in}}%
\pgfpathlineto{\pgfqpoint{2.483932in}{2.126777in}}%
\pgfpathlineto{\pgfqpoint{2.439255in}{2.164640in}}%
\pgfpathlineto{\pgfqpoint{2.408327in}{2.183976in}}%
\pgfpathclose%
\pgfusepath{stroke,fill}%
\end{pgfscope}%
\begin{pgfscope}%
\pgfpathrectangle{\pgfqpoint{0.050000in}{0.050000in}}{\pgfqpoint{4.900000in}{4.900000in}}%
\pgfusepath{clip}%
\pgfsetbuttcap%
\pgfsetroundjoin%
\definecolor{currentfill}{rgb}{0.855932,0.855932,0.855932}%
\pgfsetfillcolor{currentfill}%
\pgfsetfillopacity{0.900000}%
\pgfsetlinewidth{0.501875pt}%
\definecolor{currentstroke}{rgb}{0.200000,0.200000,0.200000}%
\pgfsetstrokecolor{currentstroke}%
\pgfsetdash{}{0pt}%
\pgfpathmoveto{\pgfqpoint{2.150471in}{2.388017in}}%
\pgfpathlineto{\pgfqpoint{2.195684in}{2.345080in}}%
\pgfpathlineto{\pgfqpoint{2.226231in}{2.325357in}}%
\pgfpathlineto{\pgfqpoint{2.181013in}{2.368230in}}%
\pgfpathlineto{\pgfqpoint{2.150471in}{2.388017in}}%
\pgfpathclose%
\pgfusepath{stroke,fill}%
\end{pgfscope}%
\begin{pgfscope}%
\pgfpathrectangle{\pgfqpoint{0.050000in}{0.050000in}}{\pgfqpoint{4.900000in}{4.900000in}}%
\pgfusepath{clip}%
\pgfsetbuttcap%
\pgfsetroundjoin%
\definecolor{currentfill}{rgb}{1.000000,1.000000,1.000000}%
\pgfsetfillcolor{currentfill}%
\pgfsetfillopacity{0.900000}%
\pgfsetlinewidth{0.501875pt}%
\definecolor{currentstroke}{rgb}{0.200000,0.200000,0.200000}%
\pgfsetstrokecolor{currentstroke}%
\pgfsetdash{}{0pt}%
\pgfpathmoveto{\pgfqpoint{3.781521in}{1.779255in}}%
\pgfpathlineto{\pgfqpoint{3.824594in}{1.792312in}}%
\pgfpathlineto{\pgfqpoint{3.857634in}{1.772354in}}%
\pgfpathlineto{\pgfqpoint{3.814533in}{1.759338in}}%
\pgfpathlineto{\pgfqpoint{3.781521in}{1.779255in}}%
\pgfpathclose%
\pgfusepath{stroke,fill}%
\end{pgfscope}%
\begin{pgfscope}%
\pgfpathrectangle{\pgfqpoint{0.050000in}{0.050000in}}{\pgfqpoint{4.900000in}{4.900000in}}%
\pgfusepath{clip}%
\pgfsetbuttcap%
\pgfsetroundjoin%
\definecolor{currentfill}{rgb}{0.781223,0.781223,0.781223}%
\pgfsetfillcolor{currentfill}%
\pgfsetfillopacity{0.900000}%
\pgfsetlinewidth{0.501875pt}%
\definecolor{currentstroke}{rgb}{0.200000,0.200000,0.200000}%
\pgfsetstrokecolor{currentstroke}%
\pgfsetdash{}{0pt}%
\pgfpathmoveto{\pgfqpoint{1.892518in}{2.596997in}}%
\pgfpathlineto{\pgfqpoint{1.938271in}{2.553388in}}%
\pgfpathlineto{\pgfqpoint{1.968436in}{2.533826in}}%
\pgfpathlineto{\pgfqpoint{1.922676in}{2.577460in}}%
\pgfpathlineto{\pgfqpoint{1.892518in}{2.596997in}}%
\pgfpathclose%
\pgfusepath{stroke,fill}%
\end{pgfscope}%
\begin{pgfscope}%
\pgfpathrectangle{\pgfqpoint{0.050000in}{0.050000in}}{\pgfqpoint{4.900000in}{4.900000in}}%
\pgfusepath{clip}%
\pgfsetbuttcap%
\pgfsetroundjoin%
\definecolor{currentfill}{rgb}{0.998155,0.998155,0.998155}%
\pgfsetfillcolor{currentfill}%
\pgfsetfillopacity{0.900000}%
\pgfsetlinewidth{0.501875pt}%
\definecolor{currentstroke}{rgb}{0.200000,0.200000,0.200000}%
\pgfsetstrokecolor{currentstroke}%
\pgfsetdash{}{0pt}%
\pgfpathmoveto{\pgfqpoint{3.315029in}{1.797501in}}%
\pgfpathlineto{\pgfqpoint{3.358550in}{1.806127in}}%
\pgfpathlineto{\pgfqpoint{3.390895in}{1.786964in}}%
\pgfpathlineto{\pgfqpoint{3.347348in}{1.778486in}}%
\pgfpathlineto{\pgfqpoint{3.315029in}{1.797501in}}%
\pgfpathclose%
\pgfusepath{stroke,fill}%
\end{pgfscope}%
\begin{pgfscope}%
\pgfpathrectangle{\pgfqpoint{0.050000in}{0.050000in}}{\pgfqpoint{4.900000in}{4.900000in}}%
\pgfusepath{clip}%
\pgfsetbuttcap%
\pgfsetroundjoin%
\definecolor{currentfill}{rgb}{0.974164,0.974164,0.974164}%
\pgfsetfillcolor{currentfill}%
\pgfsetfillopacity{0.900000}%
\pgfsetlinewidth{0.501875pt}%
\definecolor{currentstroke}{rgb}{0.200000,0.200000,0.200000}%
\pgfsetstrokecolor{currentstroke}%
\pgfsetdash{}{0pt}%
\pgfpathmoveto{\pgfqpoint{2.817305in}{1.925380in}}%
\pgfpathlineto{\pgfqpoint{2.861290in}{1.910627in}}%
\pgfpathlineto{\pgfqpoint{2.892872in}{1.892770in}}%
\pgfpathlineto{\pgfqpoint{2.848868in}{1.907408in}}%
\pgfpathlineto{\pgfqpoint{2.817305in}{1.925380in}}%
\pgfpathclose%
\pgfusepath{stroke,fill}%
\end{pgfscope}%
\begin{pgfscope}%
\pgfpathrectangle{\pgfqpoint{0.050000in}{0.050000in}}{\pgfqpoint{4.900000in}{4.900000in}}%
\pgfusepath{clip}%
\pgfsetbuttcap%
\pgfsetroundjoin%
\definecolor{currentfill}{rgb}{0.979700,0.979700,0.979700}%
\pgfsetfillcolor{currentfill}%
\pgfsetfillopacity{0.900000}%
\pgfsetlinewidth{0.501875pt}%
\definecolor{currentstroke}{rgb}{0.200000,0.200000,0.200000}%
\pgfsetstrokecolor{currentstroke}%
\pgfsetdash{}{0pt}%
\pgfpathmoveto{\pgfqpoint{2.892872in}{1.892770in}}%
\pgfpathlineto{\pgfqpoint{2.936769in}{1.882948in}}%
\pgfpathlineto{\pgfqpoint{2.968471in}{1.864999in}}%
\pgfpathlineto{\pgfqpoint{2.924552in}{1.874806in}}%
\pgfpathlineto{\pgfqpoint{2.892872in}{1.892770in}}%
\pgfpathclose%
\pgfusepath{stroke,fill}%
\end{pgfscope}%
\begin{pgfscope}%
\pgfpathrectangle{\pgfqpoint{0.050000in}{0.050000in}}{\pgfqpoint{4.900000in}{4.900000in}}%
\pgfusepath{clip}%
\pgfsetbuttcap%
\pgfsetroundjoin%
\definecolor{currentfill}{rgb}{0.686597,0.686597,0.686597}%
\pgfsetfillcolor{currentfill}%
\pgfsetfillopacity{0.900000}%
\pgfsetlinewidth{0.501875pt}%
\definecolor{currentstroke}{rgb}{0.200000,0.200000,0.200000}%
\pgfsetstrokecolor{currentstroke}%
\pgfsetdash{}{0pt}%
\pgfpathmoveto{\pgfqpoint{1.634479in}{2.806164in}}%
\pgfpathlineto{\pgfqpoint{1.680772in}{2.762116in}}%
\pgfpathlineto{\pgfqpoint{1.710554in}{2.742864in}}%
\pgfpathlineto{\pgfqpoint{1.664253in}{2.786939in}}%
\pgfpathlineto{\pgfqpoint{1.634479in}{2.806164in}}%
\pgfpathclose%
\pgfusepath{stroke,fill}%
\end{pgfscope}%
\begin{pgfscope}%
\pgfpathrectangle{\pgfqpoint{0.050000in}{0.050000in}}{\pgfqpoint{4.900000in}{4.900000in}}%
\pgfusepath{clip}%
\pgfsetbuttcap%
\pgfsetroundjoin%
\definecolor{currentfill}{rgb}{0.966782,0.966782,0.966782}%
\pgfsetfillcolor{currentfill}%
\pgfsetfillopacity{0.900000}%
\pgfsetlinewidth{0.501875pt}%
\definecolor{currentstroke}{rgb}{0.200000,0.200000,0.200000}%
\pgfsetstrokecolor{currentstroke}%
\pgfsetdash{}{0pt}%
\pgfpathmoveto{\pgfqpoint{2.741751in}{1.963216in}}%
\pgfpathlineto{\pgfqpoint{2.785840in}{1.943249in}}%
\pgfpathlineto{\pgfqpoint{2.817305in}{1.925380in}}%
\pgfpathlineto{\pgfqpoint{2.773200in}{1.945122in}}%
\pgfpathlineto{\pgfqpoint{2.741751in}{1.963216in}}%
\pgfpathclose%
\pgfusepath{stroke,fill}%
\end{pgfscope}%
\begin{pgfscope}%
\pgfpathrectangle{\pgfqpoint{0.050000in}{0.050000in}}{\pgfqpoint{4.900000in}{4.900000in}}%
\pgfusepath{clip}%
\pgfsetbuttcap%
\pgfsetroundjoin%
\definecolor{currentfill}{rgb}{1.000000,1.000000,1.000000}%
\pgfsetfillcolor{currentfill}%
\pgfsetfillopacity{0.900000}%
\pgfsetlinewidth{0.501875pt}%
\definecolor{currentstroke}{rgb}{0.200000,0.200000,0.200000}%
\pgfsetstrokecolor{currentstroke}%
\pgfsetdash{}{0pt}%
\pgfpathmoveto{\pgfqpoint{3.586216in}{1.780784in}}%
\pgfpathlineto{\pgfqpoint{3.629496in}{1.793134in}}%
\pgfpathlineto{\pgfqpoint{3.662251in}{1.773428in}}%
\pgfpathlineto{\pgfqpoint{3.618943in}{1.761147in}}%
\pgfpathlineto{\pgfqpoint{3.586216in}{1.780784in}}%
\pgfpathclose%
\pgfusepath{stroke,fill}%
\end{pgfscope}%
\begin{pgfscope}%
\pgfpathrectangle{\pgfqpoint{0.050000in}{0.050000in}}{\pgfqpoint{4.900000in}{4.900000in}}%
\pgfusepath{clip}%
\pgfsetbuttcap%
\pgfsetroundjoin%
\definecolor{currentfill}{rgb}{0.985236,0.985236,0.985236}%
\pgfsetfillcolor{currentfill}%
\pgfsetfillopacity{0.900000}%
\pgfsetlinewidth{0.501875pt}%
\definecolor{currentstroke}{rgb}{0.200000,0.200000,0.200000}%
\pgfsetstrokecolor{currentstroke}%
\pgfsetdash{}{0pt}%
\pgfpathmoveto{\pgfqpoint{2.968471in}{1.864999in}}%
\pgfpathlineto{\pgfqpoint{3.012294in}{1.859641in}}%
\pgfpathlineto{\pgfqpoint{3.044116in}{1.841525in}}%
\pgfpathlineto{\pgfqpoint{3.000270in}{1.846949in}}%
\pgfpathlineto{\pgfqpoint{2.968471in}{1.864999in}}%
\pgfpathclose%
\pgfusepath{stroke,fill}%
\end{pgfscope}%
\begin{pgfscope}%
\pgfpathrectangle{\pgfqpoint{0.050000in}{0.050000in}}{\pgfqpoint{4.900000in}{4.900000in}}%
\pgfusepath{clip}%
\pgfsetbuttcap%
\pgfsetroundjoin%
\definecolor{currentfill}{rgb}{0.898378,0.898378,0.898378}%
\pgfsetfillcolor{currentfill}%
\pgfsetfillopacity{0.900000}%
\pgfsetlinewidth{0.501875pt}%
\definecolor{currentstroke}{rgb}{0.200000,0.200000,0.200000}%
\pgfsetstrokecolor{currentstroke}%
\pgfsetdash{}{0pt}%
\pgfpathmoveto{\pgfqpoint{2.332641in}{2.243924in}}%
\pgfpathlineto{\pgfqpoint{2.377492in}{2.203298in}}%
\pgfpathlineto{\pgfqpoint{2.408327in}{2.183976in}}%
\pgfpathlineto{\pgfqpoint{2.363471in}{2.224310in}}%
\pgfpathlineto{\pgfqpoint{2.332641in}{2.243924in}}%
\pgfpathclose%
\pgfusepath{stroke,fill}%
\end{pgfscope}%
\begin{pgfscope}%
\pgfpathrectangle{\pgfqpoint{0.050000in}{0.050000in}}{\pgfqpoint{4.900000in}{4.900000in}}%
\pgfusepath{clip}%
\pgfsetbuttcap%
\pgfsetroundjoin%
\definecolor{currentfill}{rgb}{0.959400,0.959400,0.959400}%
\pgfsetfillcolor{currentfill}%
\pgfsetfillopacity{0.900000}%
\pgfsetlinewidth{0.501875pt}%
\definecolor{currentstroke}{rgb}{0.200000,0.200000,0.200000}%
\pgfsetstrokecolor{currentstroke}%
\pgfsetdash{}{0pt}%
\pgfpathmoveto{\pgfqpoint{2.666190in}{2.006411in}}%
\pgfpathlineto{\pgfqpoint{2.710400in}{1.981215in}}%
\pgfpathlineto{\pgfqpoint{2.741751in}{1.963216in}}%
\pgfpathlineto{\pgfqpoint{2.697527in}{1.988086in}}%
\pgfpathlineto{\pgfqpoint{2.666190in}{2.006411in}}%
\pgfpathclose%
\pgfusepath{stroke,fill}%
\end{pgfscope}%
\begin{pgfscope}%
\pgfpathrectangle{\pgfqpoint{0.050000in}{0.050000in}}{\pgfqpoint{4.900000in}{4.900000in}}%
\pgfusepath{clip}%
\pgfsetbuttcap%
\pgfsetroundjoin%
\definecolor{currentfill}{rgb}{0.832895,0.832895,0.832895}%
\pgfsetfillcolor{currentfill}%
\pgfsetfillopacity{0.900000}%
\pgfsetlinewidth{0.501875pt}%
\definecolor{currentstroke}{rgb}{0.200000,0.200000,0.200000}%
\pgfsetstrokecolor{currentstroke}%
\pgfsetdash{}{0pt}%
\pgfpathmoveto{\pgfqpoint{2.074624in}{2.451035in}}%
\pgfpathlineto{\pgfqpoint{2.120020in}{2.407761in}}%
\pgfpathlineto{\pgfqpoint{2.150471in}{2.388017in}}%
\pgfpathlineto{\pgfqpoint{2.105070in}{2.431286in}}%
\pgfpathlineto{\pgfqpoint{2.074624in}{2.451035in}}%
\pgfpathclose%
\pgfusepath{stroke,fill}%
\end{pgfscope}%
\begin{pgfscope}%
\pgfpathrectangle{\pgfqpoint{0.050000in}{0.050000in}}{\pgfqpoint{4.900000in}{4.900000in}}%
\pgfusepath{clip}%
\pgfsetbuttcap%
\pgfsetroundjoin%
\definecolor{currentfill}{rgb}{0.988927,0.988927,0.988927}%
\pgfsetfillcolor{currentfill}%
\pgfsetfillopacity{0.900000}%
\pgfsetlinewidth{0.501875pt}%
\definecolor{currentstroke}{rgb}{0.200000,0.200000,0.200000}%
\pgfsetstrokecolor{currentstroke}%
\pgfsetdash{}{0pt}%
\pgfpathmoveto{\pgfqpoint{3.044116in}{1.841525in}}%
\pgfpathlineto{\pgfqpoint{3.087875in}{1.840062in}}%
\pgfpathlineto{\pgfqpoint{3.119820in}{1.821734in}}%
\pgfpathlineto{\pgfqpoint{3.076037in}{1.823318in}}%
\pgfpathlineto{\pgfqpoint{3.044116in}{1.841525in}}%
\pgfpathclose%
\pgfusepath{stroke,fill}%
\end{pgfscope}%
\begin{pgfscope}%
\pgfpathrectangle{\pgfqpoint{0.050000in}{0.050000in}}{\pgfqpoint{4.900000in}{4.900000in}}%
\pgfusepath{clip}%
\pgfsetbuttcap%
\pgfsetroundjoin%
\definecolor{currentfill}{rgb}{0.998155,0.998155,0.998155}%
\pgfsetfillcolor{currentfill}%
\pgfsetfillopacity{0.900000}%
\pgfsetlinewidth{0.501875pt}%
\definecolor{currentstroke}{rgb}{0.200000,0.200000,0.200000}%
\pgfsetstrokecolor{currentstroke}%
\pgfsetdash{}{0pt}%
\pgfpathmoveto{\pgfqpoint{3.390895in}{1.786964in}}%
\pgfpathlineto{\pgfqpoint{3.434364in}{1.796922in}}%
\pgfpathlineto{\pgfqpoint{3.466836in}{1.777579in}}%
\pgfpathlineto{\pgfqpoint{3.423339in}{1.767748in}}%
\pgfpathlineto{\pgfqpoint{3.390895in}{1.786964in}}%
\pgfpathclose%
\pgfusepath{stroke,fill}%
\end{pgfscope}%
\begin{pgfscope}%
\pgfpathrectangle{\pgfqpoint{0.050000in}{0.050000in}}{\pgfqpoint{4.900000in}{4.900000in}}%
\pgfusepath{clip}%
\pgfsetbuttcap%
\pgfsetroundjoin%
\definecolor{currentfill}{rgb}{0.753664,0.753664,0.753664}%
\pgfsetfillcolor{currentfill}%
\pgfsetfillopacity{0.900000}%
\pgfsetlinewidth{0.501875pt}%
\definecolor{currentstroke}{rgb}{0.200000,0.200000,0.200000}%
\pgfsetstrokecolor{currentstroke}%
\pgfsetdash{}{0pt}%
\pgfpathmoveto{\pgfqpoint{1.816514in}{2.660240in}}%
\pgfpathlineto{\pgfqpoint{1.862450in}{2.616476in}}%
\pgfpathlineto{\pgfqpoint{1.892518in}{2.596997in}}%
\pgfpathlineto{\pgfqpoint{1.846575in}{2.640787in}}%
\pgfpathlineto{\pgfqpoint{1.816514in}{2.660240in}}%
\pgfpathclose%
\pgfusepath{stroke,fill}%
\end{pgfscope}%
\begin{pgfscope}%
\pgfpathrectangle{\pgfqpoint{0.050000in}{0.050000in}}{\pgfqpoint{4.900000in}{4.900000in}}%
\pgfusepath{clip}%
\pgfsetbuttcap%
\pgfsetroundjoin%
\definecolor{currentfill}{rgb}{0.948328,0.948328,0.948328}%
\pgfsetfillcolor{currentfill}%
\pgfsetfillopacity{0.900000}%
\pgfsetlinewidth{0.501875pt}%
\definecolor{currentstroke}{rgb}{0.200000,0.200000,0.200000}%
\pgfsetstrokecolor{currentstroke}%
\pgfsetdash{}{0pt}%
\pgfpathmoveto{\pgfqpoint{2.590600in}{2.054783in}}%
\pgfpathlineto{\pgfqpoint{2.634950in}{2.024661in}}%
\pgfpathlineto{\pgfqpoint{2.666190in}{2.006411in}}%
\pgfpathlineto{\pgfqpoint{2.621829in}{2.036138in}}%
\pgfpathlineto{\pgfqpoint{2.590600in}{2.054783in}}%
\pgfpathclose%
\pgfusepath{stroke,fill}%
\end{pgfscope}%
\begin{pgfscope}%
\pgfpathrectangle{\pgfqpoint{0.050000in}{0.050000in}}{\pgfqpoint{4.900000in}{4.900000in}}%
\pgfusepath{clip}%
\pgfsetbuttcap%
\pgfsetroundjoin%
\definecolor{currentfill}{rgb}{1.000000,1.000000,1.000000}%
\pgfsetfillcolor{currentfill}%
\pgfsetfillopacity{0.900000}%
\pgfsetlinewidth{0.501875pt}%
\definecolor{currentstroke}{rgb}{0.200000,0.200000,0.200000}%
\pgfsetstrokecolor{currentstroke}%
\pgfsetdash{}{0pt}%
\pgfpathmoveto{\pgfqpoint{3.662251in}{1.773428in}}%
\pgfpathlineto{\pgfqpoint{3.705474in}{1.786140in}}%
\pgfpathlineto{\pgfqpoint{3.738358in}{1.766326in}}%
\pgfpathlineto{\pgfqpoint{3.695107in}{1.753669in}}%
\pgfpathlineto{\pgfqpoint{3.662251in}{1.773428in}}%
\pgfpathclose%
\pgfusepath{stroke,fill}%
\end{pgfscope}%
\begin{pgfscope}%
\pgfpathrectangle{\pgfqpoint{0.050000in}{0.050000in}}{\pgfqpoint{4.900000in}{4.900000in}}%
\pgfusepath{clip}%
\pgfsetbuttcap%
\pgfsetroundjoin%
\definecolor{currentfill}{rgb}{0.990773,0.990773,0.990773}%
\pgfsetfillcolor{currentfill}%
\pgfsetfillopacity{0.900000}%
\pgfsetlinewidth{0.501875pt}%
\definecolor{currentstroke}{rgb}{0.200000,0.200000,0.200000}%
\pgfsetstrokecolor{currentstroke}%
\pgfsetdash{}{0pt}%
\pgfpathmoveto{\pgfqpoint{3.119820in}{1.821734in}}%
\pgfpathlineto{\pgfqpoint{3.163522in}{1.823566in}}%
\pgfpathlineto{\pgfqpoint{3.195590in}{1.805005in}}%
\pgfpathlineto{\pgfqpoint{3.151864in}{1.803324in}}%
\pgfpathlineto{\pgfqpoint{3.119820in}{1.821734in}}%
\pgfpathclose%
\pgfusepath{stroke,fill}%
\end{pgfscope}%
\begin{pgfscope}%
\pgfpathrectangle{\pgfqpoint{0.050000in}{0.050000in}}{\pgfqpoint{4.900000in}{4.900000in}}%
\pgfusepath{clip}%
\pgfsetbuttcap%
\pgfsetroundjoin%
\definecolor{currentfill}{rgb}{0.648212,0.648212,0.648212}%
\pgfsetfillcolor{currentfill}%
\pgfsetfillopacity{0.900000}%
\pgfsetlinewidth{0.501875pt}%
\definecolor{currentstroke}{rgb}{0.200000,0.200000,0.200000}%
\pgfsetstrokecolor{currentstroke}%
\pgfsetdash{}{0pt}%
\pgfpathmoveto{\pgfqpoint{1.558319in}{2.869535in}}%
\pgfpathlineto{\pgfqpoint{1.604795in}{2.825332in}}%
\pgfpathlineto{\pgfqpoint{1.634479in}{2.806164in}}%
\pgfpathlineto{\pgfqpoint{1.587995in}{2.850394in}}%
\pgfpathlineto{\pgfqpoint{1.558319in}{2.869535in}}%
\pgfpathclose%
\pgfusepath{stroke,fill}%
\end{pgfscope}%
\begin{pgfscope}%
\pgfpathrectangle{\pgfqpoint{0.050000in}{0.050000in}}{\pgfqpoint{4.900000in}{4.900000in}}%
\pgfusepath{clip}%
\pgfsetbuttcap%
\pgfsetroundjoin%
\definecolor{currentfill}{rgb}{0.878570,0.878570,0.878570}%
\pgfsetfillcolor{currentfill}%
\pgfsetfillopacity{0.900000}%
\pgfsetlinewidth{0.501875pt}%
\definecolor{currentstroke}{rgb}{0.200000,0.200000,0.200000}%
\pgfsetstrokecolor{currentstroke}%
\pgfsetdash{}{0pt}%
\pgfpathmoveto{\pgfqpoint{2.256871in}{2.305608in}}%
\pgfpathlineto{\pgfqpoint{2.301905in}{2.263522in}}%
\pgfpathlineto{\pgfqpoint{2.332641in}{2.243924in}}%
\pgfpathlineto{\pgfqpoint{2.287603in}{2.285832in}}%
\pgfpathlineto{\pgfqpoint{2.256871in}{2.305608in}}%
\pgfpathclose%
\pgfusepath{stroke,fill}%
\end{pgfscope}%
\begin{pgfscope}%
\pgfpathrectangle{\pgfqpoint{0.050000in}{0.050000in}}{\pgfqpoint{4.900000in}{4.900000in}}%
\pgfusepath{clip}%
\pgfsetbuttcap%
\pgfsetroundjoin%
\definecolor{currentfill}{rgb}{0.998155,0.998155,0.998155}%
\pgfsetfillcolor{currentfill}%
\pgfsetfillopacity{0.900000}%
\pgfsetlinewidth{0.501875pt}%
\definecolor{currentstroke}{rgb}{0.200000,0.200000,0.200000}%
\pgfsetstrokecolor{currentstroke}%
\pgfsetdash{}{0pt}%
\pgfpathmoveto{\pgfqpoint{3.466836in}{1.777579in}}%
\pgfpathlineto{\pgfqpoint{3.510253in}{1.788553in}}%
\pgfpathlineto{\pgfqpoint{3.542852in}{1.769052in}}%
\pgfpathlineto{\pgfqpoint{3.499408in}{1.758184in}}%
\pgfpathlineto{\pgfqpoint{3.466836in}{1.777579in}}%
\pgfpathclose%
\pgfusepath{stroke,fill}%
\end{pgfscope}%
\begin{pgfscope}%
\pgfpathrectangle{\pgfqpoint{0.050000in}{0.050000in}}{\pgfqpoint{4.900000in}{4.900000in}}%
\pgfusepath{clip}%
\pgfsetbuttcap%
\pgfsetroundjoin%
\definecolor{currentfill}{rgb}{0.935163,0.935163,0.935163}%
\pgfsetfillcolor{currentfill}%
\pgfsetfillopacity{0.900000}%
\pgfsetlinewidth{0.501875pt}%
\definecolor{currentstroke}{rgb}{0.200000,0.200000,0.200000}%
\pgfsetstrokecolor{currentstroke}%
\pgfsetdash{}{0pt}%
\pgfpathmoveto{\pgfqpoint{2.514961in}{2.107801in}}%
\pgfpathlineto{\pgfqpoint{2.559467in}{2.073375in}}%
\pgfpathlineto{\pgfqpoint{2.590600in}{2.054783in}}%
\pgfpathlineto{\pgfqpoint{2.546085in}{2.088792in}}%
\pgfpathlineto{\pgfqpoint{2.514961in}{2.107801in}}%
\pgfpathclose%
\pgfusepath{stroke,fill}%
\end{pgfscope}%
\begin{pgfscope}%
\pgfpathrectangle{\pgfqpoint{0.050000in}{0.050000in}}{\pgfqpoint{4.900000in}{4.900000in}}%
\pgfusepath{clip}%
\pgfsetbuttcap%
\pgfsetroundjoin%
\definecolor{currentfill}{rgb}{0.994464,0.994464,0.994464}%
\pgfsetfillcolor{currentfill}%
\pgfsetfillopacity{0.900000}%
\pgfsetlinewidth{0.501875pt}%
\definecolor{currentstroke}{rgb}{0.200000,0.200000,0.200000}%
\pgfsetstrokecolor{currentstroke}%
\pgfsetdash{}{0pt}%
\pgfpathmoveto{\pgfqpoint{3.195590in}{1.805005in}}%
\pgfpathlineto{\pgfqpoint{3.239239in}{1.809554in}}%
\pgfpathlineto{\pgfqpoint{3.271432in}{1.790760in}}%
\pgfpathlineto{\pgfqpoint{3.227757in}{1.786373in}}%
\pgfpathlineto{\pgfqpoint{3.195590in}{1.805005in}}%
\pgfpathclose%
\pgfusepath{stroke,fill}%
\end{pgfscope}%
\begin{pgfscope}%
\pgfpathrectangle{\pgfqpoint{0.050000in}{0.050000in}}{\pgfqpoint{4.900000in}{4.900000in}}%
\pgfusepath{clip}%
\pgfsetbuttcap%
\pgfsetroundjoin%
\definecolor{currentfill}{rgb}{0.805336,0.805336,0.805336}%
\pgfsetfillcolor{currentfill}%
\pgfsetfillopacity{0.900000}%
\pgfsetlinewidth{0.501875pt}%
\definecolor{currentstroke}{rgb}{0.200000,0.200000,0.200000}%
\pgfsetstrokecolor{currentstroke}%
\pgfsetdash{}{0pt}%
\pgfpathmoveto{\pgfqpoint{1.998693in}{2.514205in}}%
\pgfpathlineto{\pgfqpoint{2.044271in}{2.470732in}}%
\pgfpathlineto{\pgfqpoint{2.074624in}{2.451035in}}%
\pgfpathlineto{\pgfqpoint{2.029041in}{2.494526in}}%
\pgfpathlineto{\pgfqpoint{1.998693in}{2.514205in}}%
\pgfpathclose%
\pgfusepath{stroke,fill}%
\end{pgfscope}%
\begin{pgfscope}%
\pgfpathrectangle{\pgfqpoint{0.050000in}{0.050000in}}{\pgfqpoint{4.900000in}{4.900000in}}%
\pgfusepath{clip}%
\pgfsetbuttcap%
\pgfsetroundjoin%
\definecolor{currentfill}{rgb}{1.000000,1.000000,1.000000}%
\pgfsetfillcolor{currentfill}%
\pgfsetfillopacity{0.900000}%
\pgfsetlinewidth{0.501875pt}%
\definecolor{currentstroke}{rgb}{0.200000,0.200000,0.200000}%
\pgfsetstrokecolor{currentstroke}%
\pgfsetdash{}{0pt}%
\pgfpathmoveto{\pgfqpoint{3.738358in}{1.766326in}}%
\pgfpathlineto{\pgfqpoint{3.781521in}{1.779255in}}%
\pgfpathlineto{\pgfqpoint{3.814533in}{1.759338in}}%
\pgfpathlineto{\pgfqpoint{3.771342in}{1.746456in}}%
\pgfpathlineto{\pgfqpoint{3.738358in}{1.766326in}}%
\pgfpathclose%
\pgfusepath{stroke,fill}%
\end{pgfscope}%
\begin{pgfscope}%
\pgfpathrectangle{\pgfqpoint{0.050000in}{0.050000in}}{\pgfqpoint{4.900000in}{4.900000in}}%
\pgfusepath{clip}%
\pgfsetbuttcap%
\pgfsetroundjoin%
\definecolor{currentfill}{rgb}{0.724983,0.724983,0.724983}%
\pgfsetfillcolor{currentfill}%
\pgfsetfillopacity{0.900000}%
\pgfsetlinewidth{0.501875pt}%
\definecolor{currentstroke}{rgb}{0.200000,0.200000,0.200000}%
\pgfsetstrokecolor{currentstroke}%
\pgfsetdash{}{0pt}%
\pgfpathmoveto{\pgfqpoint{1.740426in}{2.723554in}}%
\pgfpathlineto{\pgfqpoint{1.786544in}{2.679635in}}%
\pgfpathlineto{\pgfqpoint{1.816514in}{2.660240in}}%
\pgfpathlineto{\pgfqpoint{1.770389in}{2.704186in}}%
\pgfpathlineto{\pgfqpoint{1.740426in}{2.723554in}}%
\pgfpathclose%
\pgfusepath{stroke,fill}%
\end{pgfscope}%
\begin{pgfscope}%
\pgfpathrectangle{\pgfqpoint{0.050000in}{0.050000in}}{\pgfqpoint{4.900000in}{4.900000in}}%
\pgfusepath{clip}%
\pgfsetbuttcap%
\pgfsetroundjoin%
\definecolor{currentfill}{rgb}{0.996309,0.996309,0.996309}%
\pgfsetfillcolor{currentfill}%
\pgfsetfillopacity{0.900000}%
\pgfsetlinewidth{0.501875pt}%
\definecolor{currentstroke}{rgb}{0.200000,0.200000,0.200000}%
\pgfsetstrokecolor{currentstroke}%
\pgfsetdash{}{0pt}%
\pgfpathmoveto{\pgfqpoint{3.271432in}{1.790760in}}%
\pgfpathlineto{\pgfqpoint{3.315029in}{1.797501in}}%
\pgfpathlineto{\pgfqpoint{3.347348in}{1.778486in}}%
\pgfpathlineto{\pgfqpoint{3.303724in}{1.771902in}}%
\pgfpathlineto{\pgfqpoint{3.271432in}{1.790760in}}%
\pgfpathclose%
\pgfusepath{stroke,fill}%
\end{pgfscope}%
\begin{pgfscope}%
\pgfpathrectangle{\pgfqpoint{0.050000in}{0.050000in}}{\pgfqpoint{4.900000in}{4.900000in}}%
\pgfusepath{clip}%
\pgfsetbuttcap%
\pgfsetroundjoin%
\definecolor{currentfill}{rgb}{1.000000,1.000000,1.000000}%
\pgfsetfillcolor{currentfill}%
\pgfsetfillopacity{0.900000}%
\pgfsetlinewidth{0.501875pt}%
\definecolor{currentstroke}{rgb}{0.200000,0.200000,0.200000}%
\pgfsetstrokecolor{currentstroke}%
\pgfsetdash{}{0pt}%
\pgfpathmoveto{\pgfqpoint{3.542852in}{1.769052in}}%
\pgfpathlineto{\pgfqpoint{3.586216in}{1.780784in}}%
\pgfpathlineto{\pgfqpoint{3.618943in}{1.761147in}}%
\pgfpathlineto{\pgfqpoint{3.575552in}{1.749501in}}%
\pgfpathlineto{\pgfqpoint{3.542852in}{1.769052in}}%
\pgfpathclose%
\pgfusepath{stroke,fill}%
\end{pgfscope}%
\begin{pgfscope}%
\pgfpathrectangle{\pgfqpoint{0.050000in}{0.050000in}}{\pgfqpoint{4.900000in}{4.900000in}}%
\pgfusepath{clip}%
\pgfsetbuttcap%
\pgfsetroundjoin%
\definecolor{currentfill}{rgb}{0.918185,0.918185,0.918185}%
\pgfsetfillcolor{currentfill}%
\pgfsetfillopacity{0.900000}%
\pgfsetlinewidth{0.501875pt}%
\definecolor{currentstroke}{rgb}{0.200000,0.200000,0.200000}%
\pgfsetstrokecolor{currentstroke}%
\pgfsetdash{}{0pt}%
\pgfpathmoveto{\pgfqpoint{2.439255in}{2.164640in}}%
\pgfpathlineto{\pgfqpoint{2.483932in}{2.126777in}}%
\pgfpathlineto{\pgfqpoint{2.514961in}{2.107801in}}%
\pgfpathlineto{\pgfqpoint{2.470277in}{2.145285in}}%
\pgfpathlineto{\pgfqpoint{2.439255in}{2.164640in}}%
\pgfpathclose%
\pgfusepath{stroke,fill}%
\end{pgfscope}%
\begin{pgfscope}%
\pgfpathrectangle{\pgfqpoint{0.050000in}{0.050000in}}{\pgfqpoint{4.900000in}{4.900000in}}%
\pgfusepath{clip}%
\pgfsetbuttcap%
\pgfsetroundjoin%
\definecolor{currentfill}{rgb}{0.855932,0.855932,0.855932}%
\pgfsetfillcolor{currentfill}%
\pgfsetfillopacity{0.900000}%
\pgfsetlinewidth{0.501875pt}%
\definecolor{currentstroke}{rgb}{0.200000,0.200000,0.200000}%
\pgfsetstrokecolor{currentstroke}%
\pgfsetdash{}{0pt}%
\pgfpathmoveto{\pgfqpoint{2.181013in}{2.368230in}}%
\pgfpathlineto{\pgfqpoint{2.226231in}{2.325357in}}%
\pgfpathlineto{\pgfqpoint{2.256871in}{2.305608in}}%
\pgfpathlineto{\pgfqpoint{2.211649in}{2.348402in}}%
\pgfpathlineto{\pgfqpoint{2.181013in}{2.368230in}}%
\pgfpathclose%
\pgfusepath{stroke,fill}%
\end{pgfscope}%
\begin{pgfscope}%
\pgfpathrectangle{\pgfqpoint{0.050000in}{0.050000in}}{\pgfqpoint{4.900000in}{4.900000in}}%
\pgfusepath{clip}%
\pgfsetbuttcap%
\pgfsetroundjoin%
\definecolor{currentfill}{rgb}{1.000000,1.000000,1.000000}%
\pgfsetfillcolor{currentfill}%
\pgfsetfillopacity{0.900000}%
\pgfsetlinewidth{0.501875pt}%
\definecolor{currentstroke}{rgb}{0.200000,0.200000,0.200000}%
\pgfsetstrokecolor{currentstroke}%
\pgfsetdash{}{0pt}%
\pgfpathmoveto{\pgfqpoint{3.814533in}{1.759338in}}%
\pgfpathlineto{\pgfqpoint{3.857634in}{1.772354in}}%
\pgfpathlineto{\pgfqpoint{3.890776in}{1.752335in}}%
\pgfpathlineto{\pgfqpoint{3.847647in}{1.739362in}}%
\pgfpathlineto{\pgfqpoint{3.814533in}{1.759338in}}%
\pgfpathclose%
\pgfusepath{stroke,fill}%
\end{pgfscope}%
\begin{pgfscope}%
\pgfpathrectangle{\pgfqpoint{0.050000in}{0.050000in}}{\pgfqpoint{4.900000in}{4.900000in}}%
\pgfusepath{clip}%
\pgfsetbuttcap%
\pgfsetroundjoin%
\definecolor{currentfill}{rgb}{0.781223,0.781223,0.781223}%
\pgfsetfillcolor{currentfill}%
\pgfsetfillopacity{0.900000}%
\pgfsetlinewidth{0.501875pt}%
\definecolor{currentstroke}{rgb}{0.200000,0.200000,0.200000}%
\pgfsetstrokecolor{currentstroke}%
\pgfsetdash{}{0pt}%
\pgfpathmoveto{\pgfqpoint{1.922676in}{2.577460in}}%
\pgfpathlineto{\pgfqpoint{1.968436in}{2.533826in}}%
\pgfpathlineto{\pgfqpoint{1.998693in}{2.514205in}}%
\pgfpathlineto{\pgfqpoint{1.952926in}{2.557864in}}%
\pgfpathlineto{\pgfqpoint{1.922676in}{2.577460in}}%
\pgfpathclose%
\pgfusepath{stroke,fill}%
\end{pgfscope}%
\begin{pgfscope}%
\pgfpathrectangle{\pgfqpoint{0.050000in}{0.050000in}}{\pgfqpoint{4.900000in}{4.900000in}}%
\pgfusepath{clip}%
\pgfsetbuttcap%
\pgfsetroundjoin%
\definecolor{currentfill}{rgb}{0.996309,0.996309,0.996309}%
\pgfsetfillcolor{currentfill}%
\pgfsetfillopacity{0.900000}%
\pgfsetlinewidth{0.501875pt}%
\definecolor{currentstroke}{rgb}{0.200000,0.200000,0.200000}%
\pgfsetstrokecolor{currentstroke}%
\pgfsetdash{}{0pt}%
\pgfpathmoveto{\pgfqpoint{3.347348in}{1.778486in}}%
\pgfpathlineto{\pgfqpoint{3.390895in}{1.786964in}}%
\pgfpathlineto{\pgfqpoint{3.423339in}{1.767748in}}%
\pgfpathlineto{\pgfqpoint{3.379765in}{1.759413in}}%
\pgfpathlineto{\pgfqpoint{3.347348in}{1.778486in}}%
\pgfpathclose%
\pgfusepath{stroke,fill}%
\end{pgfscope}%
\begin{pgfscope}%
\pgfpathrectangle{\pgfqpoint{0.050000in}{0.050000in}}{\pgfqpoint{4.900000in}{4.900000in}}%
\pgfusepath{clip}%
\pgfsetbuttcap%
\pgfsetroundjoin%
\definecolor{currentfill}{rgb}{0.974164,0.974164,0.974164}%
\pgfsetfillcolor{currentfill}%
\pgfsetfillopacity{0.900000}%
\pgfsetlinewidth{0.501875pt}%
\definecolor{currentstroke}{rgb}{0.200000,0.200000,0.200000}%
\pgfsetstrokecolor{currentstroke}%
\pgfsetdash{}{0pt}%
\pgfpathmoveto{\pgfqpoint{2.848868in}{1.907408in}}%
\pgfpathlineto{\pgfqpoint{2.892872in}{1.892770in}}%
\pgfpathlineto{\pgfqpoint{2.924552in}{1.874806in}}%
\pgfpathlineto{\pgfqpoint{2.880529in}{1.889334in}}%
\pgfpathlineto{\pgfqpoint{2.848868in}{1.907408in}}%
\pgfpathclose%
\pgfusepath{stroke,fill}%
\end{pgfscope}%
\begin{pgfscope}%
\pgfpathrectangle{\pgfqpoint{0.050000in}{0.050000in}}{\pgfqpoint{4.900000in}{4.900000in}}%
\pgfusepath{clip}%
\pgfsetbuttcap%
\pgfsetroundjoin%
\definecolor{currentfill}{rgb}{0.686597,0.686597,0.686597}%
\pgfsetfillcolor{currentfill}%
\pgfsetfillopacity{0.900000}%
\pgfsetlinewidth{0.501875pt}%
\definecolor{currentstroke}{rgb}{0.200000,0.200000,0.200000}%
\pgfsetstrokecolor{currentstroke}%
\pgfsetdash{}{0pt}%
\pgfpathmoveto{\pgfqpoint{1.664253in}{2.786939in}}%
\pgfpathlineto{\pgfqpoint{1.710554in}{2.742864in}}%
\pgfpathlineto{\pgfqpoint{1.740426in}{2.723554in}}%
\pgfpathlineto{\pgfqpoint{1.694117in}{2.767655in}}%
\pgfpathlineto{\pgfqpoint{1.664253in}{2.786939in}}%
\pgfpathclose%
\pgfusepath{stroke,fill}%
\end{pgfscope}%
\begin{pgfscope}%
\pgfpathrectangle{\pgfqpoint{0.050000in}{0.050000in}}{\pgfqpoint{4.900000in}{4.900000in}}%
\pgfusepath{clip}%
\pgfsetbuttcap%
\pgfsetroundjoin%
\definecolor{currentfill}{rgb}{0.979700,0.979700,0.979700}%
\pgfsetfillcolor{currentfill}%
\pgfsetfillopacity{0.900000}%
\pgfsetlinewidth{0.501875pt}%
\definecolor{currentstroke}{rgb}{0.200000,0.200000,0.200000}%
\pgfsetstrokecolor{currentstroke}%
\pgfsetdash{}{0pt}%
\pgfpathmoveto{\pgfqpoint{2.924552in}{1.874806in}}%
\pgfpathlineto{\pgfqpoint{2.968471in}{1.864999in}}%
\pgfpathlineto{\pgfqpoint{3.000270in}{1.846949in}}%
\pgfpathlineto{\pgfqpoint{2.956331in}{1.856740in}}%
\pgfpathlineto{\pgfqpoint{2.924552in}{1.874806in}}%
\pgfpathclose%
\pgfusepath{stroke,fill}%
\end{pgfscope}%
\begin{pgfscope}%
\pgfpathrectangle{\pgfqpoint{0.050000in}{0.050000in}}{\pgfqpoint{4.900000in}{4.900000in}}%
\pgfusepath{clip}%
\pgfsetbuttcap%
\pgfsetroundjoin%
\definecolor{currentfill}{rgb}{1.000000,1.000000,1.000000}%
\pgfsetfillcolor{currentfill}%
\pgfsetfillopacity{0.900000}%
\pgfsetlinewidth{0.501875pt}%
\definecolor{currentstroke}{rgb}{0.200000,0.200000,0.200000}%
\pgfsetstrokecolor{currentstroke}%
\pgfsetdash{}{0pt}%
\pgfpathmoveto{\pgfqpoint{3.618943in}{1.761147in}}%
\pgfpathlineto{\pgfqpoint{3.662251in}{1.773428in}}%
\pgfpathlineto{\pgfqpoint{3.695107in}{1.753669in}}%
\pgfpathlineto{\pgfqpoint{3.651771in}{1.741457in}}%
\pgfpathlineto{\pgfqpoint{3.618943in}{1.761147in}}%
\pgfpathclose%
\pgfusepath{stroke,fill}%
\end{pgfscope}%
\begin{pgfscope}%
\pgfpathrectangle{\pgfqpoint{0.050000in}{0.050000in}}{\pgfqpoint{4.900000in}{4.900000in}}%
\pgfusepath{clip}%
\pgfsetbuttcap%
\pgfsetroundjoin%
\definecolor{currentfill}{rgb}{0.966782,0.966782,0.966782}%
\pgfsetfillcolor{currentfill}%
\pgfsetfillopacity{0.900000}%
\pgfsetlinewidth{0.501875pt}%
\definecolor{currentstroke}{rgb}{0.200000,0.200000,0.200000}%
\pgfsetstrokecolor{currentstroke}%
\pgfsetdash{}{0pt}%
\pgfpathmoveto{\pgfqpoint{2.773200in}{1.945122in}}%
\pgfpathlineto{\pgfqpoint{2.817305in}{1.925380in}}%
\pgfpathlineto{\pgfqpoint{2.848868in}{1.907408in}}%
\pgfpathlineto{\pgfqpoint{2.804745in}{1.926937in}}%
\pgfpathlineto{\pgfqpoint{2.773200in}{1.945122in}}%
\pgfpathclose%
\pgfusepath{stroke,fill}%
\end{pgfscope}%
\begin{pgfscope}%
\pgfpathrectangle{\pgfqpoint{0.050000in}{0.050000in}}{\pgfqpoint{4.900000in}{4.900000in}}%
\pgfusepath{clip}%
\pgfsetbuttcap%
\pgfsetroundjoin%
\definecolor{currentfill}{rgb}{0.985236,0.985236,0.985236}%
\pgfsetfillcolor{currentfill}%
\pgfsetfillopacity{0.900000}%
\pgfsetlinewidth{0.501875pt}%
\definecolor{currentstroke}{rgb}{0.200000,0.200000,0.200000}%
\pgfsetstrokecolor{currentstroke}%
\pgfsetdash{}{0pt}%
\pgfpathmoveto{\pgfqpoint{3.000270in}{1.846949in}}%
\pgfpathlineto{\pgfqpoint{3.044116in}{1.841525in}}%
\pgfpathlineto{\pgfqpoint{3.076037in}{1.823318in}}%
\pgfpathlineto{\pgfqpoint{3.032169in}{1.828800in}}%
\pgfpathlineto{\pgfqpoint{3.000270in}{1.846949in}}%
\pgfpathclose%
\pgfusepath{stroke,fill}%
\end{pgfscope}%
\begin{pgfscope}%
\pgfpathrectangle{\pgfqpoint{0.050000in}{0.050000in}}{\pgfqpoint{4.900000in}{4.900000in}}%
\pgfusepath{clip}%
\pgfsetbuttcap%
\pgfsetroundjoin%
\definecolor{currentfill}{rgb}{0.898378,0.898378,0.898378}%
\pgfsetfillcolor{currentfill}%
\pgfsetfillopacity{0.900000}%
\pgfsetlinewidth{0.501875pt}%
\definecolor{currentstroke}{rgb}{0.200000,0.200000,0.200000}%
\pgfsetstrokecolor{currentstroke}%
\pgfsetdash{}{0pt}%
\pgfpathmoveto{\pgfqpoint{2.363471in}{2.224310in}}%
\pgfpathlineto{\pgfqpoint{2.408327in}{2.183976in}}%
\pgfpathlineto{\pgfqpoint{2.439255in}{2.164640in}}%
\pgfpathlineto{\pgfqpoint{2.394395in}{2.204679in}}%
\pgfpathlineto{\pgfqpoint{2.363471in}{2.224310in}}%
\pgfpathclose%
\pgfusepath{stroke,fill}%
\end{pgfscope}%
\begin{pgfscope}%
\pgfpathrectangle{\pgfqpoint{0.050000in}{0.050000in}}{\pgfqpoint{4.900000in}{4.900000in}}%
\pgfusepath{clip}%
\pgfsetbuttcap%
\pgfsetroundjoin%
\definecolor{currentfill}{rgb}{0.957555,0.957555,0.957555}%
\pgfsetfillcolor{currentfill}%
\pgfsetfillopacity{0.900000}%
\pgfsetlinewidth{0.501875pt}%
\definecolor{currentstroke}{rgb}{0.200000,0.200000,0.200000}%
\pgfsetstrokecolor{currentstroke}%
\pgfsetdash{}{0pt}%
\pgfpathmoveto{\pgfqpoint{2.697527in}{1.988086in}}%
\pgfpathlineto{\pgfqpoint{2.741751in}{1.963216in}}%
\pgfpathlineto{\pgfqpoint{2.773200in}{1.945122in}}%
\pgfpathlineto{\pgfqpoint{2.728961in}{1.969685in}}%
\pgfpathlineto{\pgfqpoint{2.697527in}{1.988086in}}%
\pgfpathclose%
\pgfusepath{stroke,fill}%
\end{pgfscope}%
\begin{pgfscope}%
\pgfpathrectangle{\pgfqpoint{0.050000in}{0.050000in}}{\pgfqpoint{4.900000in}{4.900000in}}%
\pgfusepath{clip}%
\pgfsetbuttcap%
\pgfsetroundjoin%
\definecolor{currentfill}{rgb}{0.832895,0.832895,0.832895}%
\pgfsetfillcolor{currentfill}%
\pgfsetfillopacity{0.900000}%
\pgfsetlinewidth{0.501875pt}%
\definecolor{currentstroke}{rgb}{0.200000,0.200000,0.200000}%
\pgfsetstrokecolor{currentstroke}%
\pgfsetdash{}{0pt}%
\pgfpathmoveto{\pgfqpoint{2.105070in}{2.431286in}}%
\pgfpathlineto{\pgfqpoint{2.150471in}{2.388017in}}%
\pgfpathlineto{\pgfqpoint{2.181013in}{2.368230in}}%
\pgfpathlineto{\pgfqpoint{2.135607in}{2.411483in}}%
\pgfpathlineto{\pgfqpoint{2.105070in}{2.431286in}}%
\pgfpathclose%
\pgfusepath{stroke,fill}%
\end{pgfscope}%
\begin{pgfscope}%
\pgfpathrectangle{\pgfqpoint{0.050000in}{0.050000in}}{\pgfqpoint{4.900000in}{4.900000in}}%
\pgfusepath{clip}%
\pgfsetbuttcap%
\pgfsetroundjoin%
\definecolor{currentfill}{rgb}{0.988927,0.988927,0.988927}%
\pgfsetfillcolor{currentfill}%
\pgfsetfillopacity{0.900000}%
\pgfsetlinewidth{0.501875pt}%
\definecolor{currentstroke}{rgb}{0.200000,0.200000,0.200000}%
\pgfsetstrokecolor{currentstroke}%
\pgfsetdash{}{0pt}%
\pgfpathmoveto{\pgfqpoint{3.076037in}{1.823318in}}%
\pgfpathlineto{\pgfqpoint{3.119820in}{1.821734in}}%
\pgfpathlineto{\pgfqpoint{3.151864in}{1.803324in}}%
\pgfpathlineto{\pgfqpoint{3.108057in}{1.805019in}}%
\pgfpathlineto{\pgfqpoint{3.076037in}{1.823318in}}%
\pgfpathclose%
\pgfusepath{stroke,fill}%
\end{pgfscope}%
\begin{pgfscope}%
\pgfpathrectangle{\pgfqpoint{0.050000in}{0.050000in}}{\pgfqpoint{4.900000in}{4.900000in}}%
\pgfusepath{clip}%
\pgfsetbuttcap%
\pgfsetroundjoin%
\definecolor{currentfill}{rgb}{0.998155,0.998155,0.998155}%
\pgfsetfillcolor{currentfill}%
\pgfsetfillopacity{0.900000}%
\pgfsetlinewidth{0.501875pt}%
\definecolor{currentstroke}{rgb}{0.200000,0.200000,0.200000}%
\pgfsetstrokecolor{currentstroke}%
\pgfsetdash{}{0pt}%
\pgfpathmoveto{\pgfqpoint{3.423339in}{1.767748in}}%
\pgfpathlineto{\pgfqpoint{3.466836in}{1.777579in}}%
\pgfpathlineto{\pgfqpoint{3.499408in}{1.758184in}}%
\pgfpathlineto{\pgfqpoint{3.455884in}{1.748478in}}%
\pgfpathlineto{\pgfqpoint{3.423339in}{1.767748in}}%
\pgfpathclose%
\pgfusepath{stroke,fill}%
\end{pgfscope}%
\begin{pgfscope}%
\pgfpathrectangle{\pgfqpoint{0.050000in}{0.050000in}}{\pgfqpoint{4.900000in}{4.900000in}}%
\pgfusepath{clip}%
\pgfsetbuttcap%
\pgfsetroundjoin%
\definecolor{currentfill}{rgb}{0.753664,0.753664,0.753664}%
\pgfsetfillcolor{currentfill}%
\pgfsetfillopacity{0.900000}%
\pgfsetlinewidth{0.501875pt}%
\definecolor{currentstroke}{rgb}{0.200000,0.200000,0.200000}%
\pgfsetstrokecolor{currentstroke}%
\pgfsetdash{}{0pt}%
\pgfpathmoveto{\pgfqpoint{1.846575in}{2.640787in}}%
\pgfpathlineto{\pgfqpoint{1.892518in}{2.596997in}}%
\pgfpathlineto{\pgfqpoint{1.922676in}{2.577460in}}%
\pgfpathlineto{\pgfqpoint{1.876727in}{2.621276in}}%
\pgfpathlineto{\pgfqpoint{1.846575in}{2.640787in}}%
\pgfpathclose%
\pgfusepath{stroke,fill}%
\end{pgfscope}%
\begin{pgfscope}%
\pgfpathrectangle{\pgfqpoint{0.050000in}{0.050000in}}{\pgfqpoint{4.900000in}{4.900000in}}%
\pgfusepath{clip}%
\pgfsetbuttcap%
\pgfsetroundjoin%
\definecolor{currentfill}{rgb}{0.948328,0.948328,0.948328}%
\pgfsetfillcolor{currentfill}%
\pgfsetfillopacity{0.900000}%
\pgfsetlinewidth{0.501875pt}%
\definecolor{currentstroke}{rgb}{0.200000,0.200000,0.200000}%
\pgfsetstrokecolor{currentstroke}%
\pgfsetdash{}{0pt}%
\pgfpathmoveto{\pgfqpoint{2.621829in}{2.036138in}}%
\pgfpathlineto{\pgfqpoint{2.666190in}{2.006411in}}%
\pgfpathlineto{\pgfqpoint{2.697527in}{1.988086in}}%
\pgfpathlineto{\pgfqpoint{2.653154in}{2.017436in}}%
\pgfpathlineto{\pgfqpoint{2.621829in}{2.036138in}}%
\pgfpathclose%
\pgfusepath{stroke,fill}%
\end{pgfscope}%
\begin{pgfscope}%
\pgfpathrectangle{\pgfqpoint{0.050000in}{0.050000in}}{\pgfqpoint{4.900000in}{4.900000in}}%
\pgfusepath{clip}%
\pgfsetbuttcap%
\pgfsetroundjoin%
\definecolor{currentfill}{rgb}{1.000000,1.000000,1.000000}%
\pgfsetfillcolor{currentfill}%
\pgfsetfillopacity{0.900000}%
\pgfsetlinewidth{0.501875pt}%
\definecolor{currentstroke}{rgb}{0.200000,0.200000,0.200000}%
\pgfsetstrokecolor{currentstroke}%
\pgfsetdash{}{0pt}%
\pgfpathmoveto{\pgfqpoint{3.695107in}{1.753669in}}%
\pgfpathlineto{\pgfqpoint{3.738358in}{1.766326in}}%
\pgfpathlineto{\pgfqpoint{3.771342in}{1.746456in}}%
\pgfpathlineto{\pgfqpoint{3.728063in}{1.733856in}}%
\pgfpathlineto{\pgfqpoint{3.695107in}{1.753669in}}%
\pgfpathclose%
\pgfusepath{stroke,fill}%
\end{pgfscope}%
\begin{pgfscope}%
\pgfpathrectangle{\pgfqpoint{0.050000in}{0.050000in}}{\pgfqpoint{4.900000in}{4.900000in}}%
\pgfusepath{clip}%
\pgfsetbuttcap%
\pgfsetroundjoin%
\definecolor{currentfill}{rgb}{0.990773,0.990773,0.990773}%
\pgfsetfillcolor{currentfill}%
\pgfsetfillopacity{0.900000}%
\pgfsetlinewidth{0.501875pt}%
\definecolor{currentstroke}{rgb}{0.200000,0.200000,0.200000}%
\pgfsetstrokecolor{currentstroke}%
\pgfsetdash{}{0pt}%
\pgfpathmoveto{\pgfqpoint{3.151864in}{1.803324in}}%
\pgfpathlineto{\pgfqpoint{3.195590in}{1.805005in}}%
\pgfpathlineto{\pgfqpoint{3.227757in}{1.786373in}}%
\pgfpathlineto{\pgfqpoint{3.184007in}{1.784833in}}%
\pgfpathlineto{\pgfqpoint{3.151864in}{1.803324in}}%
\pgfpathclose%
\pgfusepath{stroke,fill}%
\end{pgfscope}%
\begin{pgfscope}%
\pgfpathrectangle{\pgfqpoint{0.050000in}{0.050000in}}{\pgfqpoint{4.900000in}{4.900000in}}%
\pgfusepath{clip}%
\pgfsetbuttcap%
\pgfsetroundjoin%
\definecolor{currentfill}{rgb}{0.653010,0.653010,0.653010}%
\pgfsetfillcolor{currentfill}%
\pgfsetfillopacity{0.900000}%
\pgfsetlinewidth{0.501875pt}%
\definecolor{currentstroke}{rgb}{0.200000,0.200000,0.200000}%
\pgfsetstrokecolor{currentstroke}%
\pgfsetdash{}{0pt}%
\pgfpathmoveto{\pgfqpoint{1.587995in}{2.850394in}}%
\pgfpathlineto{\pgfqpoint{1.634479in}{2.806164in}}%
\pgfpathlineto{\pgfqpoint{1.664253in}{2.786939in}}%
\pgfpathlineto{\pgfqpoint{1.617760in}{2.831196in}}%
\pgfpathlineto{\pgfqpoint{1.587995in}{2.850394in}}%
\pgfpathclose%
\pgfusepath{stroke,fill}%
\end{pgfscope}%
\begin{pgfscope}%
\pgfpathrectangle{\pgfqpoint{0.050000in}{0.050000in}}{\pgfqpoint{4.900000in}{4.900000in}}%
\pgfusepath{clip}%
\pgfsetbuttcap%
\pgfsetroundjoin%
\definecolor{currentfill}{rgb}{0.878570,0.878570,0.878570}%
\pgfsetfillcolor{currentfill}%
\pgfsetfillopacity{0.900000}%
\pgfsetlinewidth{0.501875pt}%
\definecolor{currentstroke}{rgb}{0.200000,0.200000,0.200000}%
\pgfsetstrokecolor{currentstroke}%
\pgfsetdash{}{0pt}%
\pgfpathmoveto{\pgfqpoint{2.287603in}{2.285832in}}%
\pgfpathlineto{\pgfqpoint{2.332641in}{2.243924in}}%
\pgfpathlineto{\pgfqpoint{2.363471in}{2.224310in}}%
\pgfpathlineto{\pgfqpoint{2.318429in}{2.266030in}}%
\pgfpathlineto{\pgfqpoint{2.287603in}{2.285832in}}%
\pgfpathclose%
\pgfusepath{stroke,fill}%
\end{pgfscope}%
\begin{pgfscope}%
\pgfpathrectangle{\pgfqpoint{0.050000in}{0.050000in}}{\pgfqpoint{4.900000in}{4.900000in}}%
\pgfusepath{clip}%
\pgfsetbuttcap%
\pgfsetroundjoin%
\definecolor{currentfill}{rgb}{0.998155,0.998155,0.998155}%
\pgfsetfillcolor{currentfill}%
\pgfsetfillopacity{0.900000}%
\pgfsetlinewidth{0.501875pt}%
\definecolor{currentstroke}{rgb}{0.200000,0.200000,0.200000}%
\pgfsetstrokecolor{currentstroke}%
\pgfsetdash{}{0pt}%
\pgfpathmoveto{\pgfqpoint{3.499408in}{1.758184in}}%
\pgfpathlineto{\pgfqpoint{3.542852in}{1.769052in}}%
\pgfpathlineto{\pgfqpoint{3.575552in}{1.749501in}}%
\pgfpathlineto{\pgfqpoint{3.532080in}{1.738737in}}%
\pgfpathlineto{\pgfqpoint{3.499408in}{1.758184in}}%
\pgfpathclose%
\pgfusepath{stroke,fill}%
\end{pgfscope}%
\begin{pgfscope}%
\pgfpathrectangle{\pgfqpoint{0.050000in}{0.050000in}}{\pgfqpoint{4.900000in}{4.900000in}}%
\pgfusepath{clip}%
\pgfsetbuttcap%
\pgfsetroundjoin%
\definecolor{currentfill}{rgb}{0.935163,0.935163,0.935163}%
\pgfsetfillcolor{currentfill}%
\pgfsetfillopacity{0.900000}%
\pgfsetlinewidth{0.501875pt}%
\definecolor{currentstroke}{rgb}{0.200000,0.200000,0.200000}%
\pgfsetstrokecolor{currentstroke}%
\pgfsetdash{}{0pt}%
\pgfpathmoveto{\pgfqpoint{2.546085in}{2.088792in}}%
\pgfpathlineto{\pgfqpoint{2.590600in}{2.054783in}}%
\pgfpathlineto{\pgfqpoint{2.621829in}{2.036138in}}%
\pgfpathlineto{\pgfqpoint{2.577305in}{2.069748in}}%
\pgfpathlineto{\pgfqpoint{2.546085in}{2.088792in}}%
\pgfpathclose%
\pgfusepath{stroke,fill}%
\end{pgfscope}%
\begin{pgfscope}%
\pgfpathrectangle{\pgfqpoint{0.050000in}{0.050000in}}{\pgfqpoint{4.900000in}{4.900000in}}%
\pgfusepath{clip}%
\pgfsetbuttcap%
\pgfsetroundjoin%
\definecolor{currentfill}{rgb}{0.805336,0.805336,0.805336}%
\pgfsetfillcolor{currentfill}%
\pgfsetfillopacity{0.900000}%
\pgfsetlinewidth{0.501875pt}%
\definecolor{currentstroke}{rgb}{0.200000,0.200000,0.200000}%
\pgfsetstrokecolor{currentstroke}%
\pgfsetdash{}{0pt}%
\pgfpathmoveto{\pgfqpoint{2.029041in}{2.494526in}}%
\pgfpathlineto{\pgfqpoint{2.074624in}{2.451035in}}%
\pgfpathlineto{\pgfqpoint{2.105070in}{2.431286in}}%
\pgfpathlineto{\pgfqpoint{2.059480in}{2.474791in}}%
\pgfpathlineto{\pgfqpoint{2.029041in}{2.494526in}}%
\pgfpathclose%
\pgfusepath{stroke,fill}%
\end{pgfscope}%
\begin{pgfscope}%
\pgfpathrectangle{\pgfqpoint{0.050000in}{0.050000in}}{\pgfqpoint{4.900000in}{4.900000in}}%
\pgfusepath{clip}%
\pgfsetbuttcap%
\pgfsetroundjoin%
\definecolor{currentfill}{rgb}{0.994464,0.994464,0.994464}%
\pgfsetfillcolor{currentfill}%
\pgfsetfillopacity{0.900000}%
\pgfsetlinewidth{0.501875pt}%
\definecolor{currentstroke}{rgb}{0.200000,0.200000,0.200000}%
\pgfsetstrokecolor{currentstroke}%
\pgfsetdash{}{0pt}%
\pgfpathmoveto{\pgfqpoint{3.227757in}{1.786373in}}%
\pgfpathlineto{\pgfqpoint{3.271432in}{1.790760in}}%
\pgfpathlineto{\pgfqpoint{3.303724in}{1.771902in}}%
\pgfpathlineto{\pgfqpoint{3.260024in}{1.767667in}}%
\pgfpathlineto{\pgfqpoint{3.227757in}{1.786373in}}%
\pgfpathclose%
\pgfusepath{stroke,fill}%
\end{pgfscope}%
\begin{pgfscope}%
\pgfpathrectangle{\pgfqpoint{0.050000in}{0.050000in}}{\pgfqpoint{4.900000in}{4.900000in}}%
\pgfusepath{clip}%
\pgfsetbuttcap%
\pgfsetroundjoin%
\definecolor{currentfill}{rgb}{1.000000,1.000000,1.000000}%
\pgfsetfillcolor{currentfill}%
\pgfsetfillopacity{0.900000}%
\pgfsetlinewidth{0.501875pt}%
\definecolor{currentstroke}{rgb}{0.200000,0.200000,0.200000}%
\pgfsetstrokecolor{currentstroke}%
\pgfsetdash{}{0pt}%
\pgfpathmoveto{\pgfqpoint{3.771342in}{1.746456in}}%
\pgfpathlineto{\pgfqpoint{3.814533in}{1.759338in}}%
\pgfpathlineto{\pgfqpoint{3.847647in}{1.739362in}}%
\pgfpathlineto{\pgfqpoint{3.804428in}{1.726530in}}%
\pgfpathlineto{\pgfqpoint{3.771342in}{1.746456in}}%
\pgfpathclose%
\pgfusepath{stroke,fill}%
\end{pgfscope}%
\begin{pgfscope}%
\pgfpathrectangle{\pgfqpoint{0.050000in}{0.050000in}}{\pgfqpoint{4.900000in}{4.900000in}}%
\pgfusepath{clip}%
\pgfsetbuttcap%
\pgfsetroundjoin%
\definecolor{currentfill}{rgb}{0.724983,0.724983,0.724983}%
\pgfsetfillcolor{currentfill}%
\pgfsetfillopacity{0.900000}%
\pgfsetlinewidth{0.501875pt}%
\definecolor{currentstroke}{rgb}{0.200000,0.200000,0.200000}%
\pgfsetstrokecolor{currentstroke}%
\pgfsetdash{}{0pt}%
\pgfpathmoveto{\pgfqpoint{1.770389in}{2.704186in}}%
\pgfpathlineto{\pgfqpoint{1.816514in}{2.660240in}}%
\pgfpathlineto{\pgfqpoint{1.846575in}{2.640787in}}%
\pgfpathlineto{\pgfqpoint{1.800442in}{2.684759in}}%
\pgfpathlineto{\pgfqpoint{1.770389in}{2.704186in}}%
\pgfpathclose%
\pgfusepath{stroke,fill}%
\end{pgfscope}%
\begin{pgfscope}%
\pgfpathrectangle{\pgfqpoint{0.050000in}{0.050000in}}{\pgfqpoint{4.900000in}{4.900000in}}%
\pgfusepath{clip}%
\pgfsetbuttcap%
\pgfsetroundjoin%
\definecolor{currentfill}{rgb}{0.614625,0.614625,0.614625}%
\pgfsetfillcolor{currentfill}%
\pgfsetfillopacity{0.900000}%
\pgfsetlinewidth{0.501875pt}%
\definecolor{currentstroke}{rgb}{0.200000,0.200000,0.200000}%
\pgfsetstrokecolor{currentstroke}%
\pgfsetdash{}{0pt}%
\pgfpathmoveto{\pgfqpoint{1.511651in}{2.913921in}}%
\pgfpathlineto{\pgfqpoint{1.558319in}{2.869535in}}%
\pgfpathlineto{\pgfqpoint{1.587995in}{2.850394in}}%
\pgfpathlineto{\pgfqpoint{1.541318in}{2.894808in}}%
\pgfpathlineto{\pgfqpoint{1.511651in}{2.913921in}}%
\pgfpathclose%
\pgfusepath{stroke,fill}%
\end{pgfscope}%
\begin{pgfscope}%
\pgfpathrectangle{\pgfqpoint{0.050000in}{0.050000in}}{\pgfqpoint{4.900000in}{4.900000in}}%
\pgfusepath{clip}%
\pgfsetbuttcap%
\pgfsetroundjoin%
\definecolor{currentfill}{rgb}{0.996309,0.996309,0.996309}%
\pgfsetfillcolor{currentfill}%
\pgfsetfillopacity{0.900000}%
\pgfsetlinewidth{0.501875pt}%
\definecolor{currentstroke}{rgb}{0.200000,0.200000,0.200000}%
\pgfsetstrokecolor{currentstroke}%
\pgfsetdash{}{0pt}%
\pgfpathmoveto{\pgfqpoint{3.303724in}{1.771902in}}%
\pgfpathlineto{\pgfqpoint{3.347348in}{1.778486in}}%
\pgfpathlineto{\pgfqpoint{3.379765in}{1.759413in}}%
\pgfpathlineto{\pgfqpoint{3.336116in}{1.752978in}}%
\pgfpathlineto{\pgfqpoint{3.303724in}{1.771902in}}%
\pgfpathclose%
\pgfusepath{stroke,fill}%
\end{pgfscope}%
\begin{pgfscope}%
\pgfpathrectangle{\pgfqpoint{0.050000in}{0.050000in}}{\pgfqpoint{4.900000in}{4.900000in}}%
\pgfusepath{clip}%
\pgfsetbuttcap%
\pgfsetroundjoin%
\definecolor{currentfill}{rgb}{1.000000,1.000000,1.000000}%
\pgfsetfillcolor{currentfill}%
\pgfsetfillopacity{0.900000}%
\pgfsetlinewidth{0.501875pt}%
\definecolor{currentstroke}{rgb}{0.200000,0.200000,0.200000}%
\pgfsetstrokecolor{currentstroke}%
\pgfsetdash{}{0pt}%
\pgfpathmoveto{\pgfqpoint{3.575552in}{1.749501in}}%
\pgfpathlineto{\pgfqpoint{3.618943in}{1.761147in}}%
\pgfpathlineto{\pgfqpoint{3.651771in}{1.741457in}}%
\pgfpathlineto{\pgfqpoint{3.608352in}{1.729897in}}%
\pgfpathlineto{\pgfqpoint{3.575552in}{1.749501in}}%
\pgfpathclose%
\pgfusepath{stroke,fill}%
\end{pgfscope}%
\begin{pgfscope}%
\pgfpathrectangle{\pgfqpoint{0.050000in}{0.050000in}}{\pgfqpoint{4.900000in}{4.900000in}}%
\pgfusepath{clip}%
\pgfsetbuttcap%
\pgfsetroundjoin%
\definecolor{currentfill}{rgb}{0.855932,0.855932,0.855932}%
\pgfsetfillcolor{currentfill}%
\pgfsetfillopacity{0.900000}%
\pgfsetlinewidth{0.501875pt}%
\definecolor{currentstroke}{rgb}{0.200000,0.200000,0.200000}%
\pgfsetstrokecolor{currentstroke}%
\pgfsetdash{}{0pt}%
\pgfpathmoveto{\pgfqpoint{2.211649in}{2.348402in}}%
\pgfpathlineto{\pgfqpoint{2.256871in}{2.305608in}}%
\pgfpathlineto{\pgfqpoint{2.287603in}{2.285832in}}%
\pgfpathlineto{\pgfqpoint{2.242376in}{2.328535in}}%
\pgfpathlineto{\pgfqpoint{2.211649in}{2.348402in}}%
\pgfpathclose%
\pgfusepath{stroke,fill}%
\end{pgfscope}%
\begin{pgfscope}%
\pgfpathrectangle{\pgfqpoint{0.050000in}{0.050000in}}{\pgfqpoint{4.900000in}{4.900000in}}%
\pgfusepath{clip}%
\pgfsetbuttcap%
\pgfsetroundjoin%
\definecolor{currentfill}{rgb}{0.918185,0.918185,0.918185}%
\pgfsetfillcolor{currentfill}%
\pgfsetfillopacity{0.900000}%
\pgfsetlinewidth{0.501875pt}%
\definecolor{currentstroke}{rgb}{0.200000,0.200000,0.200000}%
\pgfsetstrokecolor{currentstroke}%
\pgfsetdash{}{0pt}%
\pgfpathmoveto{\pgfqpoint{2.470277in}{2.145285in}}%
\pgfpathlineto{\pgfqpoint{2.514961in}{2.107801in}}%
\pgfpathlineto{\pgfqpoint{2.546085in}{2.088792in}}%
\pgfpathlineto{\pgfqpoint{2.501395in}{2.125906in}}%
\pgfpathlineto{\pgfqpoint{2.470277in}{2.145285in}}%
\pgfpathclose%
\pgfusepath{stroke,fill}%
\end{pgfscope}%
\begin{pgfscope}%
\pgfpathrectangle{\pgfqpoint{0.050000in}{0.050000in}}{\pgfqpoint{4.900000in}{4.900000in}}%
\pgfusepath{clip}%
\pgfsetbuttcap%
\pgfsetroundjoin%
\definecolor{currentfill}{rgb}{0.781223,0.781223,0.781223}%
\pgfsetfillcolor{currentfill}%
\pgfsetfillopacity{0.900000}%
\pgfsetlinewidth{0.501875pt}%
\definecolor{currentstroke}{rgb}{0.200000,0.200000,0.200000}%
\pgfsetstrokecolor{currentstroke}%
\pgfsetdash{}{0pt}%
\pgfpathmoveto{\pgfqpoint{1.952926in}{2.557864in}}%
\pgfpathlineto{\pgfqpoint{1.998693in}{2.514205in}}%
\pgfpathlineto{\pgfqpoint{2.029041in}{2.494526in}}%
\pgfpathlineto{\pgfqpoint{1.983268in}{2.538209in}}%
\pgfpathlineto{\pgfqpoint{1.952926in}{2.557864in}}%
\pgfpathclose%
\pgfusepath{stroke,fill}%
\end{pgfscope}%
\begin{pgfscope}%
\pgfpathrectangle{\pgfqpoint{0.050000in}{0.050000in}}{\pgfqpoint{4.900000in}{4.900000in}}%
\pgfusepath{clip}%
\pgfsetbuttcap%
\pgfsetroundjoin%
\definecolor{currentfill}{rgb}{1.000000,1.000000,1.000000}%
\pgfsetfillcolor{currentfill}%
\pgfsetfillopacity{0.900000}%
\pgfsetlinewidth{0.501875pt}%
\definecolor{currentstroke}{rgb}{0.200000,0.200000,0.200000}%
\pgfsetstrokecolor{currentstroke}%
\pgfsetdash{}{0pt}%
\pgfpathmoveto{\pgfqpoint{3.847647in}{1.739362in}}%
\pgfpathlineto{\pgfqpoint{3.890776in}{1.752335in}}%
\pgfpathlineto{\pgfqpoint{3.924018in}{1.732255in}}%
\pgfpathlineto{\pgfqpoint{3.880861in}{1.719328in}}%
\pgfpathlineto{\pgfqpoint{3.847647in}{1.739362in}}%
\pgfpathclose%
\pgfusepath{stroke,fill}%
\end{pgfscope}%
\begin{pgfscope}%
\pgfpathrectangle{\pgfqpoint{0.050000in}{0.050000in}}{\pgfqpoint{4.900000in}{4.900000in}}%
\pgfusepath{clip}%
\pgfsetbuttcap%
\pgfsetroundjoin%
\definecolor{currentfill}{rgb}{0.996309,0.996309,0.996309}%
\pgfsetfillcolor{currentfill}%
\pgfsetfillopacity{0.900000}%
\pgfsetlinewidth{0.501875pt}%
\definecolor{currentstroke}{rgb}{0.200000,0.200000,0.200000}%
\pgfsetstrokecolor{currentstroke}%
\pgfsetdash{}{0pt}%
\pgfpathmoveto{\pgfqpoint{3.379765in}{1.759413in}}%
\pgfpathlineto{\pgfqpoint{3.423339in}{1.767748in}}%
\pgfpathlineto{\pgfqpoint{3.455884in}{1.748478in}}%
\pgfpathlineto{\pgfqpoint{3.412284in}{1.740280in}}%
\pgfpathlineto{\pgfqpoint{3.379765in}{1.759413in}}%
\pgfpathclose%
\pgfusepath{stroke,fill}%
\end{pgfscope}%
\begin{pgfscope}%
\pgfpathrectangle{\pgfqpoint{0.050000in}{0.050000in}}{\pgfqpoint{4.900000in}{4.900000in}}%
\pgfusepath{clip}%
\pgfsetbuttcap%
\pgfsetroundjoin%
\definecolor{currentfill}{rgb}{0.686597,0.686597,0.686597}%
\pgfsetfillcolor{currentfill}%
\pgfsetfillopacity{0.900000}%
\pgfsetlinewidth{0.501875pt}%
\definecolor{currentstroke}{rgb}{0.200000,0.200000,0.200000}%
\pgfsetstrokecolor{currentstroke}%
\pgfsetdash{}{0pt}%
\pgfpathmoveto{\pgfqpoint{1.694117in}{2.767655in}}%
\pgfpathlineto{\pgfqpoint{1.740426in}{2.723554in}}%
\pgfpathlineto{\pgfqpoint{1.770389in}{2.704186in}}%
\pgfpathlineto{\pgfqpoint{1.724071in}{2.748314in}}%
\pgfpathlineto{\pgfqpoint{1.694117in}{2.767655in}}%
\pgfpathclose%
\pgfusepath{stroke,fill}%
\end{pgfscope}%
\begin{pgfscope}%
\pgfpathrectangle{\pgfqpoint{0.050000in}{0.050000in}}{\pgfqpoint{4.900000in}{4.900000in}}%
\pgfusepath{clip}%
\pgfsetbuttcap%
\pgfsetroundjoin%
\definecolor{currentfill}{rgb}{0.972318,0.972318,0.972318}%
\pgfsetfillcolor{currentfill}%
\pgfsetfillopacity{0.900000}%
\pgfsetlinewidth{0.501875pt}%
\definecolor{currentstroke}{rgb}{0.200000,0.200000,0.200000}%
\pgfsetstrokecolor{currentstroke}%
\pgfsetdash{}{0pt}%
\pgfpathmoveto{\pgfqpoint{2.880529in}{1.889334in}}%
\pgfpathlineto{\pgfqpoint{2.924552in}{1.874806in}}%
\pgfpathlineto{\pgfqpoint{2.956331in}{1.856740in}}%
\pgfpathlineto{\pgfqpoint{2.912288in}{1.871162in}}%
\pgfpathlineto{\pgfqpoint{2.880529in}{1.889334in}}%
\pgfpathclose%
\pgfusepath{stroke,fill}%
\end{pgfscope}%
\begin{pgfscope}%
\pgfpathrectangle{\pgfqpoint{0.050000in}{0.050000in}}{\pgfqpoint{4.900000in}{4.900000in}}%
\pgfusepath{clip}%
\pgfsetbuttcap%
\pgfsetroundjoin%
\definecolor{currentfill}{rgb}{0.979700,0.979700,0.979700}%
\pgfsetfillcolor{currentfill}%
\pgfsetfillopacity{0.900000}%
\pgfsetlinewidth{0.501875pt}%
\definecolor{currentstroke}{rgb}{0.200000,0.200000,0.200000}%
\pgfsetstrokecolor{currentstroke}%
\pgfsetdash{}{0pt}%
\pgfpathmoveto{\pgfqpoint{2.956331in}{1.856740in}}%
\pgfpathlineto{\pgfqpoint{3.000270in}{1.846949in}}%
\pgfpathlineto{\pgfqpoint{3.032169in}{1.828800in}}%
\pgfpathlineto{\pgfqpoint{2.988208in}{1.838573in}}%
\pgfpathlineto{\pgfqpoint{2.956331in}{1.856740in}}%
\pgfpathclose%
\pgfusepath{stroke,fill}%
\end{pgfscope}%
\begin{pgfscope}%
\pgfpathrectangle{\pgfqpoint{0.050000in}{0.050000in}}{\pgfqpoint{4.900000in}{4.900000in}}%
\pgfusepath{clip}%
\pgfsetbuttcap%
\pgfsetroundjoin%
\definecolor{currentfill}{rgb}{1.000000,1.000000,1.000000}%
\pgfsetfillcolor{currentfill}%
\pgfsetfillopacity{0.900000}%
\pgfsetlinewidth{0.501875pt}%
\definecolor{currentstroke}{rgb}{0.200000,0.200000,0.200000}%
\pgfsetstrokecolor{currentstroke}%
\pgfsetdash{}{0pt}%
\pgfpathmoveto{\pgfqpoint{3.651771in}{1.741457in}}%
\pgfpathlineto{\pgfqpoint{3.695107in}{1.753669in}}%
\pgfpathlineto{\pgfqpoint{3.728063in}{1.733856in}}%
\pgfpathlineto{\pgfqpoint{3.684699in}{1.721715in}}%
\pgfpathlineto{\pgfqpoint{3.651771in}{1.741457in}}%
\pgfpathclose%
\pgfusepath{stroke,fill}%
\end{pgfscope}%
\begin{pgfscope}%
\pgfpathrectangle{\pgfqpoint{0.050000in}{0.050000in}}{\pgfqpoint{4.900000in}{4.900000in}}%
\pgfusepath{clip}%
\pgfsetbuttcap%
\pgfsetroundjoin%
\definecolor{currentfill}{rgb}{0.966782,0.966782,0.966782}%
\pgfsetfillcolor{currentfill}%
\pgfsetfillopacity{0.900000}%
\pgfsetlinewidth{0.501875pt}%
\definecolor{currentstroke}{rgb}{0.200000,0.200000,0.200000}%
\pgfsetstrokecolor{currentstroke}%
\pgfsetdash{}{0pt}%
\pgfpathmoveto{\pgfqpoint{2.804745in}{1.926937in}}%
\pgfpathlineto{\pgfqpoint{2.848868in}{1.907408in}}%
\pgfpathlineto{\pgfqpoint{2.880529in}{1.889334in}}%
\pgfpathlineto{\pgfqpoint{2.836389in}{1.908661in}}%
\pgfpathlineto{\pgfqpoint{2.804745in}{1.926937in}}%
\pgfpathclose%
\pgfusepath{stroke,fill}%
\end{pgfscope}%
\begin{pgfscope}%
\pgfpathrectangle{\pgfqpoint{0.050000in}{0.050000in}}{\pgfqpoint{4.900000in}{4.900000in}}%
\pgfusepath{clip}%
\pgfsetbuttcap%
\pgfsetroundjoin%
\definecolor{currentfill}{rgb}{0.898378,0.898378,0.898378}%
\pgfsetfillcolor{currentfill}%
\pgfsetfillopacity{0.900000}%
\pgfsetlinewidth{0.501875pt}%
\definecolor{currentstroke}{rgb}{0.200000,0.200000,0.200000}%
\pgfsetstrokecolor{currentstroke}%
\pgfsetdash{}{0pt}%
\pgfpathmoveto{\pgfqpoint{2.394395in}{2.204679in}}%
\pgfpathlineto{\pgfqpoint{2.439255in}{2.164640in}}%
\pgfpathlineto{\pgfqpoint{2.470277in}{2.145285in}}%
\pgfpathlineto{\pgfqpoint{2.425412in}{2.185028in}}%
\pgfpathlineto{\pgfqpoint{2.394395in}{2.204679in}}%
\pgfpathclose%
\pgfusepath{stroke,fill}%
\end{pgfscope}%
\begin{pgfscope}%
\pgfpathrectangle{\pgfqpoint{0.050000in}{0.050000in}}{\pgfqpoint{4.900000in}{4.900000in}}%
\pgfusepath{clip}%
\pgfsetbuttcap%
\pgfsetroundjoin%
\definecolor{currentfill}{rgb}{0.983391,0.983391,0.983391}%
\pgfsetfillcolor{currentfill}%
\pgfsetfillopacity{0.900000}%
\pgfsetlinewidth{0.501875pt}%
\definecolor{currentstroke}{rgb}{0.200000,0.200000,0.200000}%
\pgfsetstrokecolor{currentstroke}%
\pgfsetdash{}{0pt}%
\pgfpathmoveto{\pgfqpoint{3.032169in}{1.828800in}}%
\pgfpathlineto{\pgfqpoint{3.076037in}{1.823318in}}%
\pgfpathlineto{\pgfqpoint{3.108057in}{1.805019in}}%
\pgfpathlineto{\pgfqpoint{3.064167in}{1.810555in}}%
\pgfpathlineto{\pgfqpoint{3.032169in}{1.828800in}}%
\pgfpathclose%
\pgfusepath{stroke,fill}%
\end{pgfscope}%
\begin{pgfscope}%
\pgfpathrectangle{\pgfqpoint{0.050000in}{0.050000in}}{\pgfqpoint{4.900000in}{4.900000in}}%
\pgfusepath{clip}%
\pgfsetbuttcap%
\pgfsetroundjoin%
\definecolor{currentfill}{rgb}{0.832895,0.832895,0.832895}%
\pgfsetfillcolor{currentfill}%
\pgfsetfillopacity{0.900000}%
\pgfsetlinewidth{0.501875pt}%
\definecolor{currentstroke}{rgb}{0.200000,0.200000,0.200000}%
\pgfsetstrokecolor{currentstroke}%
\pgfsetdash{}{0pt}%
\pgfpathmoveto{\pgfqpoint{2.135607in}{2.411483in}}%
\pgfpathlineto{\pgfqpoint{2.181013in}{2.368230in}}%
\pgfpathlineto{\pgfqpoint{2.211649in}{2.348402in}}%
\pgfpathlineto{\pgfqpoint{2.166237in}{2.391631in}}%
\pgfpathlineto{\pgfqpoint{2.135607in}{2.411483in}}%
\pgfpathclose%
\pgfusepath{stroke,fill}%
\end{pgfscope}%
\begin{pgfscope}%
\pgfpathrectangle{\pgfqpoint{0.050000in}{0.050000in}}{\pgfqpoint{4.900000in}{4.900000in}}%
\pgfusepath{clip}%
\pgfsetbuttcap%
\pgfsetroundjoin%
\definecolor{currentfill}{rgb}{0.957555,0.957555,0.957555}%
\pgfsetfillcolor{currentfill}%
\pgfsetfillopacity{0.900000}%
\pgfsetlinewidth{0.501875pt}%
\definecolor{currentstroke}{rgb}{0.200000,0.200000,0.200000}%
\pgfsetstrokecolor{currentstroke}%
\pgfsetdash{}{0pt}%
\pgfpathmoveto{\pgfqpoint{2.728961in}{1.969685in}}%
\pgfpathlineto{\pgfqpoint{2.773200in}{1.945122in}}%
\pgfpathlineto{\pgfqpoint{2.804745in}{1.926937in}}%
\pgfpathlineto{\pgfqpoint{2.760492in}{1.951207in}}%
\pgfpathlineto{\pgfqpoint{2.728961in}{1.969685in}}%
\pgfpathclose%
\pgfusepath{stroke,fill}%
\end{pgfscope}%
\begin{pgfscope}%
\pgfpathrectangle{\pgfqpoint{0.050000in}{0.050000in}}{\pgfqpoint{4.900000in}{4.900000in}}%
\pgfusepath{clip}%
\pgfsetbuttcap%
\pgfsetroundjoin%
\definecolor{currentfill}{rgb}{0.987082,0.987082,0.987082}%
\pgfsetfillcolor{currentfill}%
\pgfsetfillopacity{0.900000}%
\pgfsetlinewidth{0.501875pt}%
\definecolor{currentstroke}{rgb}{0.200000,0.200000,0.200000}%
\pgfsetstrokecolor{currentstroke}%
\pgfsetdash{}{0pt}%
\pgfpathmoveto{\pgfqpoint{3.108057in}{1.805019in}}%
\pgfpathlineto{\pgfqpoint{3.151864in}{1.803324in}}%
\pgfpathlineto{\pgfqpoint{3.184007in}{1.784833in}}%
\pgfpathlineto{\pgfqpoint{3.140177in}{1.786630in}}%
\pgfpathlineto{\pgfqpoint{3.108057in}{1.805019in}}%
\pgfpathclose%
\pgfusepath{stroke,fill}%
\end{pgfscope}%
\begin{pgfscope}%
\pgfpathrectangle{\pgfqpoint{0.050000in}{0.050000in}}{\pgfqpoint{4.900000in}{4.900000in}}%
\pgfusepath{clip}%
\pgfsetbuttcap%
\pgfsetroundjoin%
\definecolor{currentfill}{rgb}{0.998155,0.998155,0.998155}%
\pgfsetfillcolor{currentfill}%
\pgfsetfillopacity{0.900000}%
\pgfsetlinewidth{0.501875pt}%
\definecolor{currentstroke}{rgb}{0.200000,0.200000,0.200000}%
\pgfsetstrokecolor{currentstroke}%
\pgfsetdash{}{0pt}%
\pgfpathmoveto{\pgfqpoint{3.455884in}{1.748478in}}%
\pgfpathlineto{\pgfqpoint{3.499408in}{1.758184in}}%
\pgfpathlineto{\pgfqpoint{3.532080in}{1.738737in}}%
\pgfpathlineto{\pgfqpoint{3.488529in}{1.729151in}}%
\pgfpathlineto{\pgfqpoint{3.455884in}{1.748478in}}%
\pgfpathclose%
\pgfusepath{stroke,fill}%
\end{pgfscope}%
\begin{pgfscope}%
\pgfpathrectangle{\pgfqpoint{0.050000in}{0.050000in}}{\pgfqpoint{4.900000in}{4.900000in}}%
\pgfusepath{clip}%
\pgfsetbuttcap%
\pgfsetroundjoin%
\definecolor{currentfill}{rgb}{0.753664,0.753664,0.753664}%
\pgfsetfillcolor{currentfill}%
\pgfsetfillopacity{0.900000}%
\pgfsetlinewidth{0.501875pt}%
\definecolor{currentstroke}{rgb}{0.200000,0.200000,0.200000}%
\pgfsetstrokecolor{currentstroke}%
\pgfsetdash{}{0pt}%
\pgfpathmoveto{\pgfqpoint{1.876727in}{2.621276in}}%
\pgfpathlineto{\pgfqpoint{1.922676in}{2.577460in}}%
\pgfpathlineto{\pgfqpoint{1.952926in}{2.557864in}}%
\pgfpathlineto{\pgfqpoint{1.906970in}{2.601705in}}%
\pgfpathlineto{\pgfqpoint{1.876727in}{2.621276in}}%
\pgfpathclose%
\pgfusepath{stroke,fill}%
\end{pgfscope}%
\begin{pgfscope}%
\pgfpathrectangle{\pgfqpoint{0.050000in}{0.050000in}}{\pgfqpoint{4.900000in}{4.900000in}}%
\pgfusepath{clip}%
\pgfsetbuttcap%
\pgfsetroundjoin%
\definecolor{currentfill}{rgb}{0.948328,0.948328,0.948328}%
\pgfsetfillcolor{currentfill}%
\pgfsetfillopacity{0.900000}%
\pgfsetlinewidth{0.501875pt}%
\definecolor{currentstroke}{rgb}{0.200000,0.200000,0.200000}%
\pgfsetstrokecolor{currentstroke}%
\pgfsetdash{}{0pt}%
\pgfpathmoveto{\pgfqpoint{2.653154in}{2.017436in}}%
\pgfpathlineto{\pgfqpoint{2.697527in}{1.988086in}}%
\pgfpathlineto{\pgfqpoint{2.728961in}{1.969685in}}%
\pgfpathlineto{\pgfqpoint{2.684576in}{1.998678in}}%
\pgfpathlineto{\pgfqpoint{2.653154in}{2.017436in}}%
\pgfpathclose%
\pgfusepath{stroke,fill}%
\end{pgfscope}%
\begin{pgfscope}%
\pgfpathrectangle{\pgfqpoint{0.050000in}{0.050000in}}{\pgfqpoint{4.900000in}{4.900000in}}%
\pgfusepath{clip}%
\pgfsetbuttcap%
\pgfsetroundjoin%
\definecolor{currentfill}{rgb}{1.000000,1.000000,1.000000}%
\pgfsetfillcolor{currentfill}%
\pgfsetfillopacity{0.900000}%
\pgfsetlinewidth{0.501875pt}%
\definecolor{currentstroke}{rgb}{0.200000,0.200000,0.200000}%
\pgfsetstrokecolor{currentstroke}%
\pgfsetdash{}{0pt}%
\pgfpathmoveto{\pgfqpoint{3.728063in}{1.733856in}}%
\pgfpathlineto{\pgfqpoint{3.771342in}{1.746456in}}%
\pgfpathlineto{\pgfqpoint{3.804428in}{1.726530in}}%
\pgfpathlineto{\pgfqpoint{3.761121in}{1.713989in}}%
\pgfpathlineto{\pgfqpoint{3.728063in}{1.733856in}}%
\pgfpathclose%
\pgfusepath{stroke,fill}%
\end{pgfscope}%
\begin{pgfscope}%
\pgfpathrectangle{\pgfqpoint{0.050000in}{0.050000in}}{\pgfqpoint{4.900000in}{4.900000in}}%
\pgfusepath{clip}%
\pgfsetbuttcap%
\pgfsetroundjoin%
\definecolor{currentfill}{rgb}{0.653010,0.653010,0.653010}%
\pgfsetfillcolor{currentfill}%
\pgfsetfillopacity{0.900000}%
\pgfsetlinewidth{0.501875pt}%
\definecolor{currentstroke}{rgb}{0.200000,0.200000,0.200000}%
\pgfsetstrokecolor{currentstroke}%
\pgfsetdash{}{0pt}%
\pgfpathmoveto{\pgfqpoint{1.617760in}{2.831196in}}%
\pgfpathlineto{\pgfqpoint{1.664253in}{2.786939in}}%
\pgfpathlineto{\pgfqpoint{1.694117in}{2.767655in}}%
\pgfpathlineto{\pgfqpoint{1.647615in}{2.811940in}}%
\pgfpathlineto{\pgfqpoint{1.617760in}{2.831196in}}%
\pgfpathclose%
\pgfusepath{stroke,fill}%
\end{pgfscope}%
\begin{pgfscope}%
\pgfpathrectangle{\pgfqpoint{0.050000in}{0.050000in}}{\pgfqpoint{4.900000in}{4.900000in}}%
\pgfusepath{clip}%
\pgfsetbuttcap%
\pgfsetroundjoin%
\definecolor{currentfill}{rgb}{0.990773,0.990773,0.990773}%
\pgfsetfillcolor{currentfill}%
\pgfsetfillopacity{0.900000}%
\pgfsetlinewidth{0.501875pt}%
\definecolor{currentstroke}{rgb}{0.200000,0.200000,0.200000}%
\pgfsetstrokecolor{currentstroke}%
\pgfsetdash{}{0pt}%
\pgfpathmoveto{\pgfqpoint{3.184007in}{1.784833in}}%
\pgfpathlineto{\pgfqpoint{3.227757in}{1.786373in}}%
\pgfpathlineto{\pgfqpoint{3.260024in}{1.767667in}}%
\pgfpathlineto{\pgfqpoint{3.216249in}{1.766259in}}%
\pgfpathlineto{\pgfqpoint{3.184007in}{1.784833in}}%
\pgfpathclose%
\pgfusepath{stroke,fill}%
\end{pgfscope}%
\begin{pgfscope}%
\pgfpathrectangle{\pgfqpoint{0.050000in}{0.050000in}}{\pgfqpoint{4.900000in}{4.900000in}}%
\pgfusepath{clip}%
\pgfsetbuttcap%
\pgfsetroundjoin%
\definecolor{currentfill}{rgb}{0.878570,0.878570,0.878570}%
\pgfsetfillcolor{currentfill}%
\pgfsetfillopacity{0.900000}%
\pgfsetlinewidth{0.501875pt}%
\definecolor{currentstroke}{rgb}{0.200000,0.200000,0.200000}%
\pgfsetstrokecolor{currentstroke}%
\pgfsetdash{}{0pt}%
\pgfpathmoveto{\pgfqpoint{2.318429in}{2.266030in}}%
\pgfpathlineto{\pgfqpoint{2.363471in}{2.224310in}}%
\pgfpathlineto{\pgfqpoint{2.394395in}{2.204679in}}%
\pgfpathlineto{\pgfqpoint{2.349348in}{2.246201in}}%
\pgfpathlineto{\pgfqpoint{2.318429in}{2.266030in}}%
\pgfpathclose%
\pgfusepath{stroke,fill}%
\end{pgfscope}%
\begin{pgfscope}%
\pgfpathrectangle{\pgfqpoint{0.050000in}{0.050000in}}{\pgfqpoint{4.900000in}{4.900000in}}%
\pgfusepath{clip}%
\pgfsetbuttcap%
\pgfsetroundjoin%
\definecolor{currentfill}{rgb}{0.998155,0.998155,0.998155}%
\pgfsetfillcolor{currentfill}%
\pgfsetfillopacity{0.900000}%
\pgfsetlinewidth{0.501875pt}%
\definecolor{currentstroke}{rgb}{0.200000,0.200000,0.200000}%
\pgfsetstrokecolor{currentstroke}%
\pgfsetdash{}{0pt}%
\pgfpathmoveto{\pgfqpoint{3.532080in}{1.738737in}}%
\pgfpathlineto{\pgfqpoint{3.575552in}{1.749501in}}%
\pgfpathlineto{\pgfqpoint{3.608352in}{1.729897in}}%
\pgfpathlineto{\pgfqpoint{3.564852in}{1.719236in}}%
\pgfpathlineto{\pgfqpoint{3.532080in}{1.738737in}}%
\pgfpathclose%
\pgfusepath{stroke,fill}%
\end{pgfscope}%
\begin{pgfscope}%
\pgfpathrectangle{\pgfqpoint{0.050000in}{0.050000in}}{\pgfqpoint{4.900000in}{4.900000in}}%
\pgfusepath{clip}%
\pgfsetbuttcap%
\pgfsetroundjoin%
\definecolor{currentfill}{rgb}{0.805336,0.805336,0.805336}%
\pgfsetfillcolor{currentfill}%
\pgfsetfillopacity{0.900000}%
\pgfsetlinewidth{0.501875pt}%
\definecolor{currentstroke}{rgb}{0.200000,0.200000,0.200000}%
\pgfsetstrokecolor{currentstroke}%
\pgfsetdash{}{0pt}%
\pgfpathmoveto{\pgfqpoint{2.059480in}{2.474791in}}%
\pgfpathlineto{\pgfqpoint{2.105070in}{2.431286in}}%
\pgfpathlineto{\pgfqpoint{2.135607in}{2.411483in}}%
\pgfpathlineto{\pgfqpoint{2.090012in}{2.454998in}}%
\pgfpathlineto{\pgfqpoint{2.059480in}{2.474791in}}%
\pgfpathclose%
\pgfusepath{stroke,fill}%
\end{pgfscope}%
\begin{pgfscope}%
\pgfpathrectangle{\pgfqpoint{0.050000in}{0.050000in}}{\pgfqpoint{4.900000in}{4.900000in}}%
\pgfusepath{clip}%
\pgfsetbuttcap%
\pgfsetroundjoin%
\definecolor{currentfill}{rgb}{0.935163,0.935163,0.935163}%
\pgfsetfillcolor{currentfill}%
\pgfsetfillopacity{0.900000}%
\pgfsetlinewidth{0.501875pt}%
\definecolor{currentstroke}{rgb}{0.200000,0.200000,0.200000}%
\pgfsetstrokecolor{currentstroke}%
\pgfsetdash{}{0pt}%
\pgfpathmoveto{\pgfqpoint{2.577305in}{2.069748in}}%
\pgfpathlineto{\pgfqpoint{2.621829in}{2.036138in}}%
\pgfpathlineto{\pgfqpoint{2.653154in}{2.017436in}}%
\pgfpathlineto{\pgfqpoint{2.608620in}{2.050663in}}%
\pgfpathlineto{\pgfqpoint{2.577305in}{2.069748in}}%
\pgfpathclose%
\pgfusepath{stroke,fill}%
\end{pgfscope}%
\begin{pgfscope}%
\pgfpathrectangle{\pgfqpoint{0.050000in}{0.050000in}}{\pgfqpoint{4.900000in}{4.900000in}}%
\pgfusepath{clip}%
\pgfsetbuttcap%
\pgfsetroundjoin%
\definecolor{currentfill}{rgb}{0.992618,0.992618,0.992618}%
\pgfsetfillcolor{currentfill}%
\pgfsetfillopacity{0.900000}%
\pgfsetlinewidth{0.501875pt}%
\definecolor{currentstroke}{rgb}{0.200000,0.200000,0.200000}%
\pgfsetstrokecolor{currentstroke}%
\pgfsetdash{}{0pt}%
\pgfpathmoveto{\pgfqpoint{3.260024in}{1.767667in}}%
\pgfpathlineto{\pgfqpoint{3.303724in}{1.771902in}}%
\pgfpathlineto{\pgfqpoint{3.336116in}{1.752978in}}%
\pgfpathlineto{\pgfqpoint{3.292391in}{1.748888in}}%
\pgfpathlineto{\pgfqpoint{3.260024in}{1.767667in}}%
\pgfpathclose%
\pgfusepath{stroke,fill}%
\end{pgfscope}%
\begin{pgfscope}%
\pgfpathrectangle{\pgfqpoint{0.050000in}{0.050000in}}{\pgfqpoint{4.900000in}{4.900000in}}%
\pgfusepath{clip}%
\pgfsetbuttcap%
\pgfsetroundjoin%
\definecolor{currentfill}{rgb}{1.000000,1.000000,1.000000}%
\pgfsetfillcolor{currentfill}%
\pgfsetfillopacity{0.900000}%
\pgfsetlinewidth{0.501875pt}%
\definecolor{currentstroke}{rgb}{0.200000,0.200000,0.200000}%
\pgfsetstrokecolor{currentstroke}%
\pgfsetdash{}{0pt}%
\pgfpathmoveto{\pgfqpoint{3.804428in}{1.726530in}}%
\pgfpathlineto{\pgfqpoint{3.847647in}{1.739362in}}%
\pgfpathlineto{\pgfqpoint{3.880861in}{1.719328in}}%
\pgfpathlineto{\pgfqpoint{3.837615in}{1.706547in}}%
\pgfpathlineto{\pgfqpoint{3.804428in}{1.726530in}}%
\pgfpathclose%
\pgfusepath{stroke,fill}%
\end{pgfscope}%
\begin{pgfscope}%
\pgfpathrectangle{\pgfqpoint{0.050000in}{0.050000in}}{\pgfqpoint{4.900000in}{4.900000in}}%
\pgfusepath{clip}%
\pgfsetbuttcap%
\pgfsetroundjoin%
\definecolor{currentfill}{rgb}{0.724983,0.724983,0.724983}%
\pgfsetfillcolor{currentfill}%
\pgfsetfillopacity{0.900000}%
\pgfsetlinewidth{0.501875pt}%
\definecolor{currentstroke}{rgb}{0.200000,0.200000,0.200000}%
\pgfsetstrokecolor{currentstroke}%
\pgfsetdash{}{0pt}%
\pgfpathmoveto{\pgfqpoint{1.800442in}{2.684759in}}%
\pgfpathlineto{\pgfqpoint{1.846575in}{2.640787in}}%
\pgfpathlineto{\pgfqpoint{1.876727in}{2.621276in}}%
\pgfpathlineto{\pgfqpoint{1.830586in}{2.665274in}}%
\pgfpathlineto{\pgfqpoint{1.800442in}{2.684759in}}%
\pgfpathclose%
\pgfusepath{stroke,fill}%
\end{pgfscope}%
\begin{pgfscope}%
\pgfpathrectangle{\pgfqpoint{0.050000in}{0.050000in}}{\pgfqpoint{4.900000in}{4.900000in}}%
\pgfusepath{clip}%
\pgfsetbuttcap%
\pgfsetroundjoin%
\definecolor{currentfill}{rgb}{0.614625,0.614625,0.614625}%
\pgfsetfillcolor{currentfill}%
\pgfsetfillopacity{0.900000}%
\pgfsetlinewidth{0.501875pt}%
\definecolor{currentstroke}{rgb}{0.200000,0.200000,0.200000}%
\pgfsetstrokecolor{currentstroke}%
\pgfsetdash{}{0pt}%
\pgfpathmoveto{\pgfqpoint{1.541318in}{2.894808in}}%
\pgfpathlineto{\pgfqpoint{1.587995in}{2.850394in}}%
\pgfpathlineto{\pgfqpoint{1.617760in}{2.831196in}}%
\pgfpathlineto{\pgfqpoint{1.571074in}{2.875637in}}%
\pgfpathlineto{\pgfqpoint{1.541318in}{2.894808in}}%
\pgfpathclose%
\pgfusepath{stroke,fill}%
\end{pgfscope}%
\begin{pgfscope}%
\pgfpathrectangle{\pgfqpoint{0.050000in}{0.050000in}}{\pgfqpoint{4.900000in}{4.900000in}}%
\pgfusepath{clip}%
\pgfsetbuttcap%
\pgfsetroundjoin%
\definecolor{currentfill}{rgb}{0.994464,0.994464,0.994464}%
\pgfsetfillcolor{currentfill}%
\pgfsetfillopacity{0.900000}%
\pgfsetlinewidth{0.501875pt}%
\definecolor{currentstroke}{rgb}{0.200000,0.200000,0.200000}%
\pgfsetstrokecolor{currentstroke}%
\pgfsetdash{}{0pt}%
\pgfpathmoveto{\pgfqpoint{3.336116in}{1.752978in}}%
\pgfpathlineto{\pgfqpoint{3.379765in}{1.759413in}}%
\pgfpathlineto{\pgfqpoint{3.412284in}{1.740280in}}%
\pgfpathlineto{\pgfqpoint{3.368608in}{1.733989in}}%
\pgfpathlineto{\pgfqpoint{3.336116in}{1.752978in}}%
\pgfpathclose%
\pgfusepath{stroke,fill}%
\end{pgfscope}%
\begin{pgfscope}%
\pgfpathrectangle{\pgfqpoint{0.050000in}{0.050000in}}{\pgfqpoint{4.900000in}{4.900000in}}%
\pgfusepath{clip}%
\pgfsetbuttcap%
\pgfsetroundjoin%
\definecolor{currentfill}{rgb}{1.000000,1.000000,1.000000}%
\pgfsetfillcolor{currentfill}%
\pgfsetfillopacity{0.900000}%
\pgfsetlinewidth{0.501875pt}%
\definecolor{currentstroke}{rgb}{0.200000,0.200000,0.200000}%
\pgfsetstrokecolor{currentstroke}%
\pgfsetdash{}{0pt}%
\pgfpathmoveto{\pgfqpoint{3.608352in}{1.729897in}}%
\pgfpathlineto{\pgfqpoint{3.651771in}{1.741457in}}%
\pgfpathlineto{\pgfqpoint{3.684699in}{1.721715in}}%
\pgfpathlineto{\pgfqpoint{3.641253in}{1.710240in}}%
\pgfpathlineto{\pgfqpoint{3.608352in}{1.729897in}}%
\pgfpathclose%
\pgfusepath{stroke,fill}%
\end{pgfscope}%
\begin{pgfscope}%
\pgfpathrectangle{\pgfqpoint{0.050000in}{0.050000in}}{\pgfqpoint{4.900000in}{4.900000in}}%
\pgfusepath{clip}%
\pgfsetbuttcap%
\pgfsetroundjoin%
\definecolor{currentfill}{rgb}{0.858762,0.858762,0.858762}%
\pgfsetfillcolor{currentfill}%
\pgfsetfillopacity{0.900000}%
\pgfsetlinewidth{0.501875pt}%
\definecolor{currentstroke}{rgb}{0.200000,0.200000,0.200000}%
\pgfsetstrokecolor{currentstroke}%
\pgfsetdash{}{0pt}%
\pgfpathmoveto{\pgfqpoint{2.242376in}{2.328535in}}%
\pgfpathlineto{\pgfqpoint{2.287603in}{2.285832in}}%
\pgfpathlineto{\pgfqpoint{2.318429in}{2.266030in}}%
\pgfpathlineto{\pgfqpoint{2.273197in}{2.308629in}}%
\pgfpathlineto{\pgfqpoint{2.242376in}{2.328535in}}%
\pgfpathclose%
\pgfusepath{stroke,fill}%
\end{pgfscope}%
\begin{pgfscope}%
\pgfpathrectangle{\pgfqpoint{0.050000in}{0.050000in}}{\pgfqpoint{4.900000in}{4.900000in}}%
\pgfusepath{clip}%
\pgfsetbuttcap%
\pgfsetroundjoin%
\definecolor{currentfill}{rgb}{0.918185,0.918185,0.918185}%
\pgfsetfillcolor{currentfill}%
\pgfsetfillopacity{0.900000}%
\pgfsetlinewidth{0.501875pt}%
\definecolor{currentstroke}{rgb}{0.200000,0.200000,0.200000}%
\pgfsetstrokecolor{currentstroke}%
\pgfsetdash{}{0pt}%
\pgfpathmoveto{\pgfqpoint{2.501395in}{2.125906in}}%
\pgfpathlineto{\pgfqpoint{2.546085in}{2.088792in}}%
\pgfpathlineto{\pgfqpoint{2.577305in}{2.069748in}}%
\pgfpathlineto{\pgfqpoint{2.532608in}{2.106501in}}%
\pgfpathlineto{\pgfqpoint{2.501395in}{2.125906in}}%
\pgfpathclose%
\pgfusepath{stroke,fill}%
\end{pgfscope}%
\begin{pgfscope}%
\pgfpathrectangle{\pgfqpoint{0.050000in}{0.050000in}}{\pgfqpoint{4.900000in}{4.900000in}}%
\pgfusepath{clip}%
\pgfsetbuttcap%
\pgfsetroundjoin%
\definecolor{currentfill}{rgb}{0.781223,0.781223,0.781223}%
\pgfsetfillcolor{currentfill}%
\pgfsetfillopacity{0.900000}%
\pgfsetlinewidth{0.501875pt}%
\definecolor{currentstroke}{rgb}{0.200000,0.200000,0.200000}%
\pgfsetstrokecolor{currentstroke}%
\pgfsetdash{}{0pt}%
\pgfpathmoveto{\pgfqpoint{1.983268in}{2.538209in}}%
\pgfpathlineto{\pgfqpoint{2.029041in}{2.494526in}}%
\pgfpathlineto{\pgfqpoint{2.059480in}{2.474791in}}%
\pgfpathlineto{\pgfqpoint{2.013701in}{2.518495in}}%
\pgfpathlineto{\pgfqpoint{1.983268in}{2.538209in}}%
\pgfpathclose%
\pgfusepath{stroke,fill}%
\end{pgfscope}%
\begin{pgfscope}%
\pgfpathrectangle{\pgfqpoint{0.050000in}{0.050000in}}{\pgfqpoint{4.900000in}{4.900000in}}%
\pgfusepath{clip}%
\pgfsetbuttcap%
\pgfsetroundjoin%
\definecolor{currentfill}{rgb}{1.000000,1.000000,1.000000}%
\pgfsetfillcolor{currentfill}%
\pgfsetfillopacity{0.900000}%
\pgfsetlinewidth{0.501875pt}%
\definecolor{currentstroke}{rgb}{0.200000,0.200000,0.200000}%
\pgfsetstrokecolor{currentstroke}%
\pgfsetdash{}{0pt}%
\pgfpathmoveto{\pgfqpoint{3.880861in}{1.719328in}}%
\pgfpathlineto{\pgfqpoint{3.924018in}{1.732255in}}%
\pgfpathlineto{\pgfqpoint{3.957362in}{1.712114in}}%
\pgfpathlineto{\pgfqpoint{3.914178in}{1.699235in}}%
\pgfpathlineto{\pgfqpoint{3.880861in}{1.719328in}}%
\pgfpathclose%
\pgfusepath{stroke,fill}%
\end{pgfscope}%
\begin{pgfscope}%
\pgfpathrectangle{\pgfqpoint{0.050000in}{0.050000in}}{\pgfqpoint{4.900000in}{4.900000in}}%
\pgfusepath{clip}%
\pgfsetbuttcap%
\pgfsetroundjoin%
\definecolor{currentfill}{rgb}{0.686597,0.686597,0.686597}%
\pgfsetfillcolor{currentfill}%
\pgfsetfillopacity{0.900000}%
\pgfsetlinewidth{0.501875pt}%
\definecolor{currentstroke}{rgb}{0.200000,0.200000,0.200000}%
\pgfsetstrokecolor{currentstroke}%
\pgfsetdash{}{0pt}%
\pgfpathmoveto{\pgfqpoint{1.724071in}{2.748314in}}%
\pgfpathlineto{\pgfqpoint{1.770389in}{2.704186in}}%
\pgfpathlineto{\pgfqpoint{1.800442in}{2.684759in}}%
\pgfpathlineto{\pgfqpoint{1.754117in}{2.728914in}}%
\pgfpathlineto{\pgfqpoint{1.724071in}{2.748314in}}%
\pgfpathclose%
\pgfusepath{stroke,fill}%
\end{pgfscope}%
\begin{pgfscope}%
\pgfpathrectangle{\pgfqpoint{0.050000in}{0.050000in}}{\pgfqpoint{4.900000in}{4.900000in}}%
\pgfusepath{clip}%
\pgfsetbuttcap%
\pgfsetroundjoin%
\definecolor{currentfill}{rgb}{0.996309,0.996309,0.996309}%
\pgfsetfillcolor{currentfill}%
\pgfsetfillopacity{0.900000}%
\pgfsetlinewidth{0.501875pt}%
\definecolor{currentstroke}{rgb}{0.200000,0.200000,0.200000}%
\pgfsetstrokecolor{currentstroke}%
\pgfsetdash{}{0pt}%
\pgfpathmoveto{\pgfqpoint{3.412284in}{1.740280in}}%
\pgfpathlineto{\pgfqpoint{3.455884in}{1.748478in}}%
\pgfpathlineto{\pgfqpoint{3.488529in}{1.729151in}}%
\pgfpathlineto{\pgfqpoint{3.444902in}{1.721087in}}%
\pgfpathlineto{\pgfqpoint{3.412284in}{1.740280in}}%
\pgfpathclose%
\pgfusepath{stroke,fill}%
\end{pgfscope}%
\begin{pgfscope}%
\pgfpathrectangle{\pgfqpoint{0.050000in}{0.050000in}}{\pgfqpoint{4.900000in}{4.900000in}}%
\pgfusepath{clip}%
\pgfsetbuttcap%
\pgfsetroundjoin%
\definecolor{currentfill}{rgb}{0.972318,0.972318,0.972318}%
\pgfsetfillcolor{currentfill}%
\pgfsetfillopacity{0.900000}%
\pgfsetlinewidth{0.501875pt}%
\definecolor{currentstroke}{rgb}{0.200000,0.200000,0.200000}%
\pgfsetstrokecolor{currentstroke}%
\pgfsetdash{}{0pt}%
\pgfpathmoveto{\pgfqpoint{2.912288in}{1.871162in}}%
\pgfpathlineto{\pgfqpoint{2.956331in}{1.856740in}}%
\pgfpathlineto{\pgfqpoint{2.988208in}{1.838573in}}%
\pgfpathlineto{\pgfqpoint{2.944147in}{1.852892in}}%
\pgfpathlineto{\pgfqpoint{2.912288in}{1.871162in}}%
\pgfpathclose%
\pgfusepath{stroke,fill}%
\end{pgfscope}%
\begin{pgfscope}%
\pgfpathrectangle{\pgfqpoint{0.050000in}{0.050000in}}{\pgfqpoint{4.900000in}{4.900000in}}%
\pgfusepath{clip}%
\pgfsetbuttcap%
\pgfsetroundjoin%
\definecolor{currentfill}{rgb}{1.000000,1.000000,1.000000}%
\pgfsetfillcolor{currentfill}%
\pgfsetfillopacity{0.900000}%
\pgfsetlinewidth{0.501875pt}%
\definecolor{currentstroke}{rgb}{0.200000,0.200000,0.200000}%
\pgfsetstrokecolor{currentstroke}%
\pgfsetdash{}{0pt}%
\pgfpathmoveto{\pgfqpoint{3.684699in}{1.721715in}}%
\pgfpathlineto{\pgfqpoint{3.728063in}{1.733856in}}%
\pgfpathlineto{\pgfqpoint{3.761121in}{1.713989in}}%
\pgfpathlineto{\pgfqpoint{3.717729in}{1.701919in}}%
\pgfpathlineto{\pgfqpoint{3.684699in}{1.721715in}}%
\pgfpathclose%
\pgfusepath{stroke,fill}%
\end{pgfscope}%
\begin{pgfscope}%
\pgfpathrectangle{\pgfqpoint{0.050000in}{0.050000in}}{\pgfqpoint{4.900000in}{4.900000in}}%
\pgfusepath{clip}%
\pgfsetbuttcap%
\pgfsetroundjoin%
\definecolor{currentfill}{rgb}{0.977855,0.977855,0.977855}%
\pgfsetfillcolor{currentfill}%
\pgfsetfillopacity{0.900000}%
\pgfsetlinewidth{0.501875pt}%
\definecolor{currentstroke}{rgb}{0.200000,0.200000,0.200000}%
\pgfsetstrokecolor{currentstroke}%
\pgfsetdash{}{0pt}%
\pgfpathmoveto{\pgfqpoint{2.988208in}{1.838573in}}%
\pgfpathlineto{\pgfqpoint{3.032169in}{1.828800in}}%
\pgfpathlineto{\pgfqpoint{3.064167in}{1.810555in}}%
\pgfpathlineto{\pgfqpoint{3.020185in}{1.820308in}}%
\pgfpathlineto{\pgfqpoint{2.988208in}{1.838573in}}%
\pgfpathclose%
\pgfusepath{stroke,fill}%
\end{pgfscope}%
\begin{pgfscope}%
\pgfpathrectangle{\pgfqpoint{0.050000in}{0.050000in}}{\pgfqpoint{4.900000in}{4.900000in}}%
\pgfusepath{clip}%
\pgfsetbuttcap%
\pgfsetroundjoin%
\definecolor{currentfill}{rgb}{0.964937,0.964937,0.964937}%
\pgfsetfillcolor{currentfill}%
\pgfsetfillopacity{0.900000}%
\pgfsetlinewidth{0.501875pt}%
\definecolor{currentstroke}{rgb}{0.200000,0.200000,0.200000}%
\pgfsetstrokecolor{currentstroke}%
\pgfsetdash{}{0pt}%
\pgfpathmoveto{\pgfqpoint{2.836389in}{1.908661in}}%
\pgfpathlineto{\pgfqpoint{2.880529in}{1.889334in}}%
\pgfpathlineto{\pgfqpoint{2.912288in}{1.871162in}}%
\pgfpathlineto{\pgfqpoint{2.868132in}{1.890295in}}%
\pgfpathlineto{\pgfqpoint{2.836389in}{1.908661in}}%
\pgfpathclose%
\pgfusepath{stroke,fill}%
\end{pgfscope}%
\begin{pgfscope}%
\pgfpathrectangle{\pgfqpoint{0.050000in}{0.050000in}}{\pgfqpoint{4.900000in}{4.900000in}}%
\pgfusepath{clip}%
\pgfsetbuttcap%
\pgfsetroundjoin%
\definecolor{currentfill}{rgb}{0.898378,0.898378,0.898378}%
\pgfsetfillcolor{currentfill}%
\pgfsetfillopacity{0.900000}%
\pgfsetlinewidth{0.501875pt}%
\definecolor{currentstroke}{rgb}{0.200000,0.200000,0.200000}%
\pgfsetstrokecolor{currentstroke}%
\pgfsetdash{}{0pt}%
\pgfpathmoveto{\pgfqpoint{2.425412in}{2.185028in}}%
\pgfpathlineto{\pgfqpoint{2.470277in}{2.145285in}}%
\pgfpathlineto{\pgfqpoint{2.501395in}{2.125906in}}%
\pgfpathlineto{\pgfqpoint{2.456524in}{2.165355in}}%
\pgfpathlineto{\pgfqpoint{2.425412in}{2.185028in}}%
\pgfpathclose%
\pgfusepath{stroke,fill}%
\end{pgfscope}%
\begin{pgfscope}%
\pgfpathrectangle{\pgfqpoint{0.050000in}{0.050000in}}{\pgfqpoint{4.900000in}{4.900000in}}%
\pgfusepath{clip}%
\pgfsetbuttcap%
\pgfsetroundjoin%
\definecolor{currentfill}{rgb}{0.581776,0.581776,0.581776}%
\pgfsetfillcolor{currentfill}%
\pgfsetfillopacity{0.900000}%
\pgfsetlinewidth{0.501875pt}%
\definecolor{currentstroke}{rgb}{0.200000,0.200000,0.200000}%
\pgfsetstrokecolor{currentstroke}%
\pgfsetdash{}{0pt}%
\pgfpathmoveto{\pgfqpoint{1.464790in}{2.958491in}}%
\pgfpathlineto{\pgfqpoint{1.511651in}{2.913921in}}%
\pgfpathlineto{\pgfqpoint{1.541318in}{2.894808in}}%
\pgfpathlineto{\pgfqpoint{1.494446in}{2.939406in}}%
\pgfpathlineto{\pgfqpoint{1.464790in}{2.958491in}}%
\pgfpathclose%
\pgfusepath{stroke,fill}%
\end{pgfscope}%
\begin{pgfscope}%
\pgfpathrectangle{\pgfqpoint{0.050000in}{0.050000in}}{\pgfqpoint{4.900000in}{4.900000in}}%
\pgfusepath{clip}%
\pgfsetbuttcap%
\pgfsetroundjoin%
\definecolor{currentfill}{rgb}{0.983391,0.983391,0.983391}%
\pgfsetfillcolor{currentfill}%
\pgfsetfillopacity{0.900000}%
\pgfsetlinewidth{0.501875pt}%
\definecolor{currentstroke}{rgb}{0.200000,0.200000,0.200000}%
\pgfsetstrokecolor{currentstroke}%
\pgfsetdash{}{0pt}%
\pgfpathmoveto{\pgfqpoint{3.064167in}{1.810555in}}%
\pgfpathlineto{\pgfqpoint{3.108057in}{1.805019in}}%
\pgfpathlineto{\pgfqpoint{3.140177in}{1.786630in}}%
\pgfpathlineto{\pgfqpoint{3.096264in}{1.792214in}}%
\pgfpathlineto{\pgfqpoint{3.064167in}{1.810555in}}%
\pgfpathclose%
\pgfusepath{stroke,fill}%
\end{pgfscope}%
\begin{pgfscope}%
\pgfpathrectangle{\pgfqpoint{0.050000in}{0.050000in}}{\pgfqpoint{4.900000in}{4.900000in}}%
\pgfusepath{clip}%
\pgfsetbuttcap%
\pgfsetroundjoin%
\definecolor{currentfill}{rgb}{0.832895,0.832895,0.832895}%
\pgfsetfillcolor{currentfill}%
\pgfsetfillopacity{0.900000}%
\pgfsetlinewidth{0.501875pt}%
\definecolor{currentstroke}{rgb}{0.200000,0.200000,0.200000}%
\pgfsetstrokecolor{currentstroke}%
\pgfsetdash{}{0pt}%
\pgfpathmoveto{\pgfqpoint{2.166237in}{2.391631in}}%
\pgfpathlineto{\pgfqpoint{2.211649in}{2.348402in}}%
\pgfpathlineto{\pgfqpoint{2.242376in}{2.328535in}}%
\pgfpathlineto{\pgfqpoint{2.196960in}{2.371729in}}%
\pgfpathlineto{\pgfqpoint{2.166237in}{2.391631in}}%
\pgfpathclose%
\pgfusepath{stroke,fill}%
\end{pgfscope}%
\begin{pgfscope}%
\pgfpathrectangle{\pgfqpoint{0.050000in}{0.050000in}}{\pgfqpoint{4.900000in}{4.900000in}}%
\pgfusepath{clip}%
\pgfsetbuttcap%
\pgfsetroundjoin%
\definecolor{currentfill}{rgb}{0.957555,0.957555,0.957555}%
\pgfsetfillcolor{currentfill}%
\pgfsetfillopacity{0.900000}%
\pgfsetlinewidth{0.501875pt}%
\definecolor{currentstroke}{rgb}{0.200000,0.200000,0.200000}%
\pgfsetstrokecolor{currentstroke}%
\pgfsetdash{}{0pt}%
\pgfpathmoveto{\pgfqpoint{2.760492in}{1.951207in}}%
\pgfpathlineto{\pgfqpoint{2.804745in}{1.926937in}}%
\pgfpathlineto{\pgfqpoint{2.836389in}{1.908661in}}%
\pgfpathlineto{\pgfqpoint{2.792121in}{1.932653in}}%
\pgfpathlineto{\pgfqpoint{2.760492in}{1.951207in}}%
\pgfpathclose%
\pgfusepath{stroke,fill}%
\end{pgfscope}%
\begin{pgfscope}%
\pgfpathrectangle{\pgfqpoint{0.050000in}{0.050000in}}{\pgfqpoint{4.900000in}{4.900000in}}%
\pgfusepath{clip}%
\pgfsetbuttcap%
\pgfsetroundjoin%
\definecolor{currentfill}{rgb}{0.987082,0.987082,0.987082}%
\pgfsetfillcolor{currentfill}%
\pgfsetfillopacity{0.900000}%
\pgfsetlinewidth{0.501875pt}%
\definecolor{currentstroke}{rgb}{0.200000,0.200000,0.200000}%
\pgfsetstrokecolor{currentstroke}%
\pgfsetdash{}{0pt}%
\pgfpathmoveto{\pgfqpoint{3.140177in}{1.786630in}}%
\pgfpathlineto{\pgfqpoint{3.184007in}{1.784833in}}%
\pgfpathlineto{\pgfqpoint{3.216249in}{1.766259in}}%
\pgfpathlineto{\pgfqpoint{3.172396in}{1.768153in}}%
\pgfpathlineto{\pgfqpoint{3.140177in}{1.786630in}}%
\pgfpathclose%
\pgfusepath{stroke,fill}%
\end{pgfscope}%
\begin{pgfscope}%
\pgfpathrectangle{\pgfqpoint{0.050000in}{0.050000in}}{\pgfqpoint{4.900000in}{4.900000in}}%
\pgfusepath{clip}%
\pgfsetbuttcap%
\pgfsetroundjoin%
\definecolor{currentfill}{rgb}{0.998155,0.998155,0.998155}%
\pgfsetfillcolor{currentfill}%
\pgfsetfillopacity{0.900000}%
\pgfsetlinewidth{0.501875pt}%
\definecolor{currentstroke}{rgb}{0.200000,0.200000,0.200000}%
\pgfsetstrokecolor{currentstroke}%
\pgfsetdash{}{0pt}%
\pgfpathmoveto{\pgfqpoint{3.488529in}{1.729151in}}%
\pgfpathlineto{\pgfqpoint{3.532080in}{1.738737in}}%
\pgfpathlineto{\pgfqpoint{3.564852in}{1.719236in}}%
\pgfpathlineto{\pgfqpoint{3.521275in}{1.709768in}}%
\pgfpathlineto{\pgfqpoint{3.488529in}{1.729151in}}%
\pgfpathclose%
\pgfusepath{stroke,fill}%
\end{pgfscope}%
\begin{pgfscope}%
\pgfpathrectangle{\pgfqpoint{0.050000in}{0.050000in}}{\pgfqpoint{4.900000in}{4.900000in}}%
\pgfusepath{clip}%
\pgfsetbuttcap%
\pgfsetroundjoin%
\definecolor{currentfill}{rgb}{0.753664,0.753664,0.753664}%
\pgfsetfillcolor{currentfill}%
\pgfsetfillopacity{0.900000}%
\pgfsetlinewidth{0.501875pt}%
\definecolor{currentstroke}{rgb}{0.200000,0.200000,0.200000}%
\pgfsetstrokecolor{currentstroke}%
\pgfsetdash{}{0pt}%
\pgfpathmoveto{\pgfqpoint{1.906970in}{2.601705in}}%
\pgfpathlineto{\pgfqpoint{1.952926in}{2.557864in}}%
\pgfpathlineto{\pgfqpoint{1.983268in}{2.538209in}}%
\pgfpathlineto{\pgfqpoint{1.937304in}{2.582075in}}%
\pgfpathlineto{\pgfqpoint{1.906970in}{2.601705in}}%
\pgfpathclose%
\pgfusepath{stroke,fill}%
\end{pgfscope}%
\begin{pgfscope}%
\pgfpathrectangle{\pgfqpoint{0.050000in}{0.050000in}}{\pgfqpoint{4.900000in}{4.900000in}}%
\pgfusepath{clip}%
\pgfsetbuttcap%
\pgfsetroundjoin%
\definecolor{currentfill}{rgb}{1.000000,1.000000,1.000000}%
\pgfsetfillcolor{currentfill}%
\pgfsetfillopacity{0.900000}%
\pgfsetlinewidth{0.501875pt}%
\definecolor{currentstroke}{rgb}{0.200000,0.200000,0.200000}%
\pgfsetstrokecolor{currentstroke}%
\pgfsetdash{}{0pt}%
\pgfpathmoveto{\pgfqpoint{3.761121in}{1.713989in}}%
\pgfpathlineto{\pgfqpoint{3.804428in}{1.726530in}}%
\pgfpathlineto{\pgfqpoint{3.837615in}{1.706547in}}%
\pgfpathlineto{\pgfqpoint{3.794280in}{1.694066in}}%
\pgfpathlineto{\pgfqpoint{3.761121in}{1.713989in}}%
\pgfpathclose%
\pgfusepath{stroke,fill}%
\end{pgfscope}%
\begin{pgfscope}%
\pgfpathrectangle{\pgfqpoint{0.050000in}{0.050000in}}{\pgfqpoint{4.900000in}{4.900000in}}%
\pgfusepath{clip}%
\pgfsetbuttcap%
\pgfsetroundjoin%
\definecolor{currentfill}{rgb}{0.948328,0.948328,0.948328}%
\pgfsetfillcolor{currentfill}%
\pgfsetfillopacity{0.900000}%
\pgfsetlinewidth{0.501875pt}%
\definecolor{currentstroke}{rgb}{0.200000,0.200000,0.200000}%
\pgfsetstrokecolor{currentstroke}%
\pgfsetdash{}{0pt}%
\pgfpathmoveto{\pgfqpoint{2.684576in}{1.998678in}}%
\pgfpathlineto{\pgfqpoint{2.728961in}{1.969685in}}%
\pgfpathlineto{\pgfqpoint{2.760492in}{1.951207in}}%
\pgfpathlineto{\pgfqpoint{2.716095in}{1.979859in}}%
\pgfpathlineto{\pgfqpoint{2.684576in}{1.998678in}}%
\pgfpathclose%
\pgfusepath{stroke,fill}%
\end{pgfscope}%
\begin{pgfscope}%
\pgfpathrectangle{\pgfqpoint{0.050000in}{0.050000in}}{\pgfqpoint{4.900000in}{4.900000in}}%
\pgfusepath{clip}%
\pgfsetbuttcap%
\pgfsetroundjoin%
\definecolor{currentfill}{rgb}{0.653010,0.653010,0.653010}%
\pgfsetfillcolor{currentfill}%
\pgfsetfillopacity{0.900000}%
\pgfsetlinewidth{0.501875pt}%
\definecolor{currentstroke}{rgb}{0.200000,0.200000,0.200000}%
\pgfsetstrokecolor{currentstroke}%
\pgfsetdash{}{0pt}%
\pgfpathmoveto{\pgfqpoint{1.647615in}{2.811940in}}%
\pgfpathlineto{\pgfqpoint{1.694117in}{2.767655in}}%
\pgfpathlineto{\pgfqpoint{1.724071in}{2.748314in}}%
\pgfpathlineto{\pgfqpoint{1.677562in}{2.792625in}}%
\pgfpathlineto{\pgfqpoint{1.647615in}{2.811940in}}%
\pgfpathclose%
\pgfusepath{stroke,fill}%
\end{pgfscope}%
\begin{pgfscope}%
\pgfpathrectangle{\pgfqpoint{0.050000in}{0.050000in}}{\pgfqpoint{4.900000in}{4.900000in}}%
\pgfusepath{clip}%
\pgfsetbuttcap%
\pgfsetroundjoin%
\definecolor{currentfill}{rgb}{0.990773,0.990773,0.990773}%
\pgfsetfillcolor{currentfill}%
\pgfsetfillopacity{0.900000}%
\pgfsetlinewidth{0.501875pt}%
\definecolor{currentstroke}{rgb}{0.200000,0.200000,0.200000}%
\pgfsetstrokecolor{currentstroke}%
\pgfsetdash{}{0pt}%
\pgfpathmoveto{\pgfqpoint{3.216249in}{1.766259in}}%
\pgfpathlineto{\pgfqpoint{3.260024in}{1.767667in}}%
\pgfpathlineto{\pgfqpoint{3.292391in}{1.748888in}}%
\pgfpathlineto{\pgfqpoint{3.248592in}{1.747605in}}%
\pgfpathlineto{\pgfqpoint{3.216249in}{1.766259in}}%
\pgfpathclose%
\pgfusepath{stroke,fill}%
\end{pgfscope}%
\begin{pgfscope}%
\pgfpathrectangle{\pgfqpoint{0.050000in}{0.050000in}}{\pgfqpoint{4.900000in}{4.900000in}}%
\pgfusepath{clip}%
\pgfsetbuttcap%
\pgfsetroundjoin%
\definecolor{currentfill}{rgb}{0.878570,0.878570,0.878570}%
\pgfsetfillcolor{currentfill}%
\pgfsetfillopacity{0.900000}%
\pgfsetlinewidth{0.501875pt}%
\definecolor{currentstroke}{rgb}{0.200000,0.200000,0.200000}%
\pgfsetstrokecolor{currentstroke}%
\pgfsetdash{}{0pt}%
\pgfpathmoveto{\pgfqpoint{2.349348in}{2.246201in}}%
\pgfpathlineto{\pgfqpoint{2.394395in}{2.204679in}}%
\pgfpathlineto{\pgfqpoint{2.425412in}{2.185028in}}%
\pgfpathlineto{\pgfqpoint{2.380361in}{2.226345in}}%
\pgfpathlineto{\pgfqpoint{2.349348in}{2.246201in}}%
\pgfpathclose%
\pgfusepath{stroke,fill}%
\end{pgfscope}%
\begin{pgfscope}%
\pgfpathrectangle{\pgfqpoint{0.050000in}{0.050000in}}{\pgfqpoint{4.900000in}{4.900000in}}%
\pgfusepath{clip}%
\pgfsetbuttcap%
\pgfsetroundjoin%
\definecolor{currentfill}{rgb}{0.998155,0.998155,0.998155}%
\pgfsetfillcolor{currentfill}%
\pgfsetfillopacity{0.900000}%
\pgfsetlinewidth{0.501875pt}%
\definecolor{currentstroke}{rgb}{0.200000,0.200000,0.200000}%
\pgfsetstrokecolor{currentstroke}%
\pgfsetdash{}{0pt}%
\pgfpathmoveto{\pgfqpoint{3.564852in}{1.719236in}}%
\pgfpathlineto{\pgfqpoint{3.608352in}{1.729897in}}%
\pgfpathlineto{\pgfqpoint{3.641253in}{1.710240in}}%
\pgfpathlineto{\pgfqpoint{3.597726in}{1.699681in}}%
\pgfpathlineto{\pgfqpoint{3.564852in}{1.719236in}}%
\pgfpathclose%
\pgfusepath{stroke,fill}%
\end{pgfscope}%
\begin{pgfscope}%
\pgfpathrectangle{\pgfqpoint{0.050000in}{0.050000in}}{\pgfqpoint{4.900000in}{4.900000in}}%
\pgfusepath{clip}%
\pgfsetbuttcap%
\pgfsetroundjoin%
\definecolor{currentfill}{rgb}{0.808781,0.808781,0.808781}%
\pgfsetfillcolor{currentfill}%
\pgfsetfillopacity{0.900000}%
\pgfsetlinewidth{0.501875pt}%
\definecolor{currentstroke}{rgb}{0.200000,0.200000,0.200000}%
\pgfsetstrokecolor{currentstroke}%
\pgfsetdash{}{0pt}%
\pgfpathmoveto{\pgfqpoint{2.090012in}{2.454998in}}%
\pgfpathlineto{\pgfqpoint{2.135607in}{2.411483in}}%
\pgfpathlineto{\pgfqpoint{2.166237in}{2.391631in}}%
\pgfpathlineto{\pgfqpoint{2.120637in}{2.435149in}}%
\pgfpathlineto{\pgfqpoint{2.090012in}{2.454998in}}%
\pgfpathclose%
\pgfusepath{stroke,fill}%
\end{pgfscope}%
\begin{pgfscope}%
\pgfpathrectangle{\pgfqpoint{0.050000in}{0.050000in}}{\pgfqpoint{4.900000in}{4.900000in}}%
\pgfusepath{clip}%
\pgfsetbuttcap%
\pgfsetroundjoin%
\definecolor{currentfill}{rgb}{0.935163,0.935163,0.935163}%
\pgfsetfillcolor{currentfill}%
\pgfsetfillopacity{0.900000}%
\pgfsetlinewidth{0.501875pt}%
\definecolor{currentstroke}{rgb}{0.200000,0.200000,0.200000}%
\pgfsetstrokecolor{currentstroke}%
\pgfsetdash{}{0pt}%
\pgfpathmoveto{\pgfqpoint{2.608620in}{2.050663in}}%
\pgfpathlineto{\pgfqpoint{2.653154in}{2.017436in}}%
\pgfpathlineto{\pgfqpoint{2.684576in}{1.998678in}}%
\pgfpathlineto{\pgfqpoint{2.640033in}{2.031536in}}%
\pgfpathlineto{\pgfqpoint{2.608620in}{2.050663in}}%
\pgfpathclose%
\pgfusepath{stroke,fill}%
\end{pgfscope}%
\begin{pgfscope}%
\pgfpathrectangle{\pgfqpoint{0.050000in}{0.050000in}}{\pgfqpoint{4.900000in}{4.900000in}}%
\pgfusepath{clip}%
\pgfsetbuttcap%
\pgfsetroundjoin%
\definecolor{currentfill}{rgb}{0.992618,0.992618,0.992618}%
\pgfsetfillcolor{currentfill}%
\pgfsetfillopacity{0.900000}%
\pgfsetlinewidth{0.501875pt}%
\definecolor{currentstroke}{rgb}{0.200000,0.200000,0.200000}%
\pgfsetstrokecolor{currentstroke}%
\pgfsetdash{}{0pt}%
\pgfpathmoveto{\pgfqpoint{3.292391in}{1.748888in}}%
\pgfpathlineto{\pgfqpoint{3.336116in}{1.752978in}}%
\pgfpathlineto{\pgfqpoint{3.368608in}{1.733989in}}%
\pgfpathlineto{\pgfqpoint{3.324859in}{1.730035in}}%
\pgfpathlineto{\pgfqpoint{3.292391in}{1.748888in}}%
\pgfpathclose%
\pgfusepath{stroke,fill}%
\end{pgfscope}%
\begin{pgfscope}%
\pgfpathrectangle{\pgfqpoint{0.050000in}{0.050000in}}{\pgfqpoint{4.900000in}{4.900000in}}%
\pgfusepath{clip}%
\pgfsetbuttcap%
\pgfsetroundjoin%
\definecolor{currentfill}{rgb}{1.000000,1.000000,1.000000}%
\pgfsetfillcolor{currentfill}%
\pgfsetfillopacity{0.900000}%
\pgfsetlinewidth{0.501875pt}%
\definecolor{currentstroke}{rgb}{0.200000,0.200000,0.200000}%
\pgfsetstrokecolor{currentstroke}%
\pgfsetdash{}{0pt}%
\pgfpathmoveto{\pgfqpoint{3.837615in}{1.706547in}}%
\pgfpathlineto{\pgfqpoint{3.880861in}{1.719328in}}%
\pgfpathlineto{\pgfqpoint{3.914178in}{1.699235in}}%
\pgfpathlineto{\pgfqpoint{3.870903in}{1.686506in}}%
\pgfpathlineto{\pgfqpoint{3.837615in}{1.706547in}}%
\pgfpathclose%
\pgfusepath{stroke,fill}%
\end{pgfscope}%
\begin{pgfscope}%
\pgfpathrectangle{\pgfqpoint{0.050000in}{0.050000in}}{\pgfqpoint{4.900000in}{4.900000in}}%
\pgfusepath{clip}%
\pgfsetbuttcap%
\pgfsetroundjoin%
\definecolor{currentfill}{rgb}{0.724983,0.724983,0.724983}%
\pgfsetfillcolor{currentfill}%
\pgfsetfillopacity{0.900000}%
\pgfsetlinewidth{0.501875pt}%
\definecolor{currentstroke}{rgb}{0.200000,0.200000,0.200000}%
\pgfsetstrokecolor{currentstroke}%
\pgfsetdash{}{0pt}%
\pgfpathmoveto{\pgfqpoint{1.830586in}{2.665274in}}%
\pgfpathlineto{\pgfqpoint{1.876727in}{2.621276in}}%
\pgfpathlineto{\pgfqpoint{1.906970in}{2.601705in}}%
\pgfpathlineto{\pgfqpoint{1.860822in}{2.645729in}}%
\pgfpathlineto{\pgfqpoint{1.830586in}{2.665274in}}%
\pgfpathclose%
\pgfusepath{stroke,fill}%
\end{pgfscope}%
\begin{pgfscope}%
\pgfpathrectangle{\pgfqpoint{0.050000in}{0.050000in}}{\pgfqpoint{4.900000in}{4.900000in}}%
\pgfusepath{clip}%
\pgfsetbuttcap%
\pgfsetroundjoin%
\definecolor{currentfill}{rgb}{0.614625,0.614625,0.614625}%
\pgfsetfillcolor{currentfill}%
\pgfsetfillopacity{0.900000}%
\pgfsetlinewidth{0.501875pt}%
\definecolor{currentstroke}{rgb}{0.200000,0.200000,0.200000}%
\pgfsetstrokecolor{currentstroke}%
\pgfsetdash{}{0pt}%
\pgfpathmoveto{\pgfqpoint{1.571074in}{2.875637in}}%
\pgfpathlineto{\pgfqpoint{1.617760in}{2.831196in}}%
\pgfpathlineto{\pgfqpoint{1.647615in}{2.811940in}}%
\pgfpathlineto{\pgfqpoint{1.600921in}{2.856408in}}%
\pgfpathlineto{\pgfqpoint{1.571074in}{2.875637in}}%
\pgfpathclose%
\pgfusepath{stroke,fill}%
\end{pgfscope}%
\begin{pgfscope}%
\pgfpathrectangle{\pgfqpoint{0.050000in}{0.050000in}}{\pgfqpoint{4.900000in}{4.900000in}}%
\pgfusepath{clip}%
\pgfsetbuttcap%
\pgfsetroundjoin%
\definecolor{currentfill}{rgb}{1.000000,1.000000,1.000000}%
\pgfsetfillcolor{currentfill}%
\pgfsetfillopacity{0.900000}%
\pgfsetlinewidth{0.501875pt}%
\definecolor{currentstroke}{rgb}{0.200000,0.200000,0.200000}%
\pgfsetstrokecolor{currentstroke}%
\pgfsetdash{}{0pt}%
\pgfpathmoveto{\pgfqpoint{3.641253in}{1.710240in}}%
\pgfpathlineto{\pgfqpoint{3.684699in}{1.721715in}}%
\pgfpathlineto{\pgfqpoint{3.717729in}{1.701919in}}%
\pgfpathlineto{\pgfqpoint{3.674255in}{1.690529in}}%
\pgfpathlineto{\pgfqpoint{3.641253in}{1.710240in}}%
\pgfpathclose%
\pgfusepath{stroke,fill}%
\end{pgfscope}%
\begin{pgfscope}%
\pgfpathrectangle{\pgfqpoint{0.050000in}{0.050000in}}{\pgfqpoint{4.900000in}{4.900000in}}%
\pgfusepath{clip}%
\pgfsetbuttcap%
\pgfsetroundjoin%
\definecolor{currentfill}{rgb}{0.994464,0.994464,0.994464}%
\pgfsetfillcolor{currentfill}%
\pgfsetfillopacity{0.900000}%
\pgfsetlinewidth{0.501875pt}%
\definecolor{currentstroke}{rgb}{0.200000,0.200000,0.200000}%
\pgfsetstrokecolor{currentstroke}%
\pgfsetdash{}{0pt}%
\pgfpathmoveto{\pgfqpoint{3.368608in}{1.733989in}}%
\pgfpathlineto{\pgfqpoint{3.412284in}{1.740280in}}%
\pgfpathlineto{\pgfqpoint{3.444902in}{1.721087in}}%
\pgfpathlineto{\pgfqpoint{3.401201in}{1.714932in}}%
\pgfpathlineto{\pgfqpoint{3.368608in}{1.733989in}}%
\pgfpathclose%
\pgfusepath{stroke,fill}%
\end{pgfscope}%
\begin{pgfscope}%
\pgfpathrectangle{\pgfqpoint{0.050000in}{0.050000in}}{\pgfqpoint{4.900000in}{4.900000in}}%
\pgfusepath{clip}%
\pgfsetbuttcap%
\pgfsetroundjoin%
\definecolor{currentfill}{rgb}{0.858762,0.858762,0.858762}%
\pgfsetfillcolor{currentfill}%
\pgfsetfillopacity{0.900000}%
\pgfsetlinewidth{0.501875pt}%
\definecolor{currentstroke}{rgb}{0.200000,0.200000,0.200000}%
\pgfsetstrokecolor{currentstroke}%
\pgfsetdash{}{0pt}%
\pgfpathmoveto{\pgfqpoint{2.273197in}{2.308629in}}%
\pgfpathlineto{\pgfqpoint{2.318429in}{2.266030in}}%
\pgfpathlineto{\pgfqpoint{2.349348in}{2.246201in}}%
\pgfpathlineto{\pgfqpoint{2.304112in}{2.288686in}}%
\pgfpathlineto{\pgfqpoint{2.273197in}{2.308629in}}%
\pgfpathclose%
\pgfusepath{stroke,fill}%
\end{pgfscope}%
\begin{pgfscope}%
\pgfpathrectangle{\pgfqpoint{0.050000in}{0.050000in}}{\pgfqpoint{4.900000in}{4.900000in}}%
\pgfusepath{clip}%
\pgfsetbuttcap%
\pgfsetroundjoin%
\definecolor{currentfill}{rgb}{0.918185,0.918185,0.918185}%
\pgfsetfillcolor{currentfill}%
\pgfsetfillopacity{0.900000}%
\pgfsetlinewidth{0.501875pt}%
\definecolor{currentstroke}{rgb}{0.200000,0.200000,0.200000}%
\pgfsetstrokecolor{currentstroke}%
\pgfsetdash{}{0pt}%
\pgfpathmoveto{\pgfqpoint{2.532608in}{2.106501in}}%
\pgfpathlineto{\pgfqpoint{2.577305in}{2.069748in}}%
\pgfpathlineto{\pgfqpoint{2.608620in}{2.050663in}}%
\pgfpathlineto{\pgfqpoint{2.563917in}{2.087065in}}%
\pgfpathlineto{\pgfqpoint{2.532608in}{2.106501in}}%
\pgfpathclose%
\pgfusepath{stroke,fill}%
\end{pgfscope}%
\begin{pgfscope}%
\pgfpathrectangle{\pgfqpoint{0.050000in}{0.050000in}}{\pgfqpoint{4.900000in}{4.900000in}}%
\pgfusepath{clip}%
\pgfsetbuttcap%
\pgfsetroundjoin%
\definecolor{currentfill}{rgb}{0.781223,0.781223,0.781223}%
\pgfsetfillcolor{currentfill}%
\pgfsetfillopacity{0.900000}%
\pgfsetlinewidth{0.501875pt}%
\definecolor{currentstroke}{rgb}{0.200000,0.200000,0.200000}%
\pgfsetstrokecolor{currentstroke}%
\pgfsetdash{}{0pt}%
\pgfpathmoveto{\pgfqpoint{2.013701in}{2.518495in}}%
\pgfpathlineto{\pgfqpoint{2.059480in}{2.474791in}}%
\pgfpathlineto{\pgfqpoint{2.090012in}{2.454998in}}%
\pgfpathlineto{\pgfqpoint{2.044227in}{2.498722in}}%
\pgfpathlineto{\pgfqpoint{2.013701in}{2.518495in}}%
\pgfpathclose%
\pgfusepath{stroke,fill}%
\end{pgfscope}%
\begin{pgfscope}%
\pgfpathrectangle{\pgfqpoint{0.050000in}{0.050000in}}{\pgfqpoint{4.900000in}{4.900000in}}%
\pgfusepath{clip}%
\pgfsetbuttcap%
\pgfsetroundjoin%
\definecolor{currentfill}{rgb}{1.000000,1.000000,1.000000}%
\pgfsetfillcolor{currentfill}%
\pgfsetfillopacity{0.900000}%
\pgfsetlinewidth{0.501875pt}%
\definecolor{currentstroke}{rgb}{0.200000,0.200000,0.200000}%
\pgfsetstrokecolor{currentstroke}%
\pgfsetdash{}{0pt}%
\pgfpathmoveto{\pgfqpoint{3.914178in}{1.699235in}}%
\pgfpathlineto{\pgfqpoint{3.957362in}{1.712114in}}%
\pgfpathlineto{\pgfqpoint{3.990808in}{1.691912in}}%
\pgfpathlineto{\pgfqpoint{3.947597in}{1.679082in}}%
\pgfpathlineto{\pgfqpoint{3.914178in}{1.699235in}}%
\pgfpathclose%
\pgfusepath{stroke,fill}%
\end{pgfscope}%
\begin{pgfscope}%
\pgfpathrectangle{\pgfqpoint{0.050000in}{0.050000in}}{\pgfqpoint{4.900000in}{4.900000in}}%
\pgfusepath{clip}%
\pgfsetbuttcap%
\pgfsetroundjoin%
\definecolor{currentfill}{rgb}{0.686597,0.686597,0.686597}%
\pgfsetfillcolor{currentfill}%
\pgfsetfillopacity{0.900000}%
\pgfsetlinewidth{0.501875pt}%
\definecolor{currentstroke}{rgb}{0.200000,0.200000,0.200000}%
\pgfsetstrokecolor{currentstroke}%
\pgfsetdash{}{0pt}%
\pgfpathmoveto{\pgfqpoint{1.754117in}{2.728914in}}%
\pgfpathlineto{\pgfqpoint{1.800442in}{2.684759in}}%
\pgfpathlineto{\pgfqpoint{1.830586in}{2.665274in}}%
\pgfpathlineto{\pgfqpoint{1.784253in}{2.709455in}}%
\pgfpathlineto{\pgfqpoint{1.754117in}{2.728914in}}%
\pgfpathclose%
\pgfusepath{stroke,fill}%
\end{pgfscope}%
\begin{pgfscope}%
\pgfpathrectangle{\pgfqpoint{0.050000in}{0.050000in}}{\pgfqpoint{4.900000in}{4.900000in}}%
\pgfusepath{clip}%
\pgfsetbuttcap%
\pgfsetroundjoin%
\definecolor{currentfill}{rgb}{0.996309,0.996309,0.996309}%
\pgfsetfillcolor{currentfill}%
\pgfsetfillopacity{0.900000}%
\pgfsetlinewidth{0.501875pt}%
\definecolor{currentstroke}{rgb}{0.200000,0.200000,0.200000}%
\pgfsetstrokecolor{currentstroke}%
\pgfsetdash{}{0pt}%
\pgfpathmoveto{\pgfqpoint{3.444902in}{1.721087in}}%
\pgfpathlineto{\pgfqpoint{3.488529in}{1.729151in}}%
\pgfpathlineto{\pgfqpoint{3.521275in}{1.709768in}}%
\pgfpathlineto{\pgfqpoint{3.477622in}{1.701832in}}%
\pgfpathlineto{\pgfqpoint{3.444902in}{1.721087in}}%
\pgfpathclose%
\pgfusepath{stroke,fill}%
\end{pgfscope}%
\begin{pgfscope}%
\pgfpathrectangle{\pgfqpoint{0.050000in}{0.050000in}}{\pgfqpoint{4.900000in}{4.900000in}}%
\pgfusepath{clip}%
\pgfsetbuttcap%
\pgfsetroundjoin%
\definecolor{currentfill}{rgb}{1.000000,1.000000,1.000000}%
\pgfsetfillcolor{currentfill}%
\pgfsetfillopacity{0.900000}%
\pgfsetlinewidth{0.501875pt}%
\definecolor{currentstroke}{rgb}{0.200000,0.200000,0.200000}%
\pgfsetstrokecolor{currentstroke}%
\pgfsetdash{}{0pt}%
\pgfpathmoveto{\pgfqpoint{3.717729in}{1.701919in}}%
\pgfpathlineto{\pgfqpoint{3.761121in}{1.713989in}}%
\pgfpathlineto{\pgfqpoint{3.794280in}{1.694066in}}%
\pgfpathlineto{\pgfqpoint{3.750861in}{1.682067in}}%
\pgfpathlineto{\pgfqpoint{3.717729in}{1.701919in}}%
\pgfpathclose%
\pgfusepath{stroke,fill}%
\end{pgfscope}%
\begin{pgfscope}%
\pgfpathrectangle{\pgfqpoint{0.050000in}{0.050000in}}{\pgfqpoint{4.900000in}{4.900000in}}%
\pgfusepath{clip}%
\pgfsetbuttcap%
\pgfsetroundjoin%
\definecolor{currentfill}{rgb}{0.972318,0.972318,0.972318}%
\pgfsetfillcolor{currentfill}%
\pgfsetfillopacity{0.900000}%
\pgfsetlinewidth{0.501875pt}%
\definecolor{currentstroke}{rgb}{0.200000,0.200000,0.200000}%
\pgfsetstrokecolor{currentstroke}%
\pgfsetdash{}{0pt}%
\pgfpathmoveto{\pgfqpoint{2.944147in}{1.852892in}}%
\pgfpathlineto{\pgfqpoint{2.988208in}{1.838573in}}%
\pgfpathlineto{\pgfqpoint{3.020185in}{1.820308in}}%
\pgfpathlineto{\pgfqpoint{2.976104in}{1.834528in}}%
\pgfpathlineto{\pgfqpoint{2.944147in}{1.852892in}}%
\pgfpathclose%
\pgfusepath{stroke,fill}%
\end{pgfscope}%
\begin{pgfscope}%
\pgfpathrectangle{\pgfqpoint{0.050000in}{0.050000in}}{\pgfqpoint{4.900000in}{4.900000in}}%
\pgfusepath{clip}%
\pgfsetbuttcap%
\pgfsetroundjoin%
\definecolor{currentfill}{rgb}{0.977855,0.977855,0.977855}%
\pgfsetfillcolor{currentfill}%
\pgfsetfillopacity{0.900000}%
\pgfsetlinewidth{0.501875pt}%
\definecolor{currentstroke}{rgb}{0.200000,0.200000,0.200000}%
\pgfsetstrokecolor{currentstroke}%
\pgfsetdash{}{0pt}%
\pgfpathmoveto{\pgfqpoint{3.020185in}{1.820308in}}%
\pgfpathlineto{\pgfqpoint{3.064167in}{1.810555in}}%
\pgfpathlineto{\pgfqpoint{3.096264in}{1.792214in}}%
\pgfpathlineto{\pgfqpoint{3.052262in}{1.801947in}}%
\pgfpathlineto{\pgfqpoint{3.020185in}{1.820308in}}%
\pgfpathclose%
\pgfusepath{stroke,fill}%
\end{pgfscope}%
\begin{pgfscope}%
\pgfpathrectangle{\pgfqpoint{0.050000in}{0.050000in}}{\pgfqpoint{4.900000in}{4.900000in}}%
\pgfusepath{clip}%
\pgfsetbuttcap%
\pgfsetroundjoin%
\definecolor{currentfill}{rgb}{0.581776,0.581776,0.581776}%
\pgfsetfillcolor{currentfill}%
\pgfsetfillopacity{0.900000}%
\pgfsetlinewidth{0.501875pt}%
\definecolor{currentstroke}{rgb}{0.200000,0.200000,0.200000}%
\pgfsetstrokecolor{currentstroke}%
\pgfsetdash{}{0pt}%
\pgfpathmoveto{\pgfqpoint{1.494446in}{2.939406in}}%
\pgfpathlineto{\pgfqpoint{1.541318in}{2.894808in}}%
\pgfpathlineto{\pgfqpoint{1.571074in}{2.875637in}}%
\pgfpathlineto{\pgfqpoint{1.524193in}{2.920263in}}%
\pgfpathlineto{\pgfqpoint{1.494446in}{2.939406in}}%
\pgfpathclose%
\pgfusepath{stroke,fill}%
\end{pgfscope}%
\begin{pgfscope}%
\pgfpathrectangle{\pgfqpoint{0.050000in}{0.050000in}}{\pgfqpoint{4.900000in}{4.900000in}}%
\pgfusepath{clip}%
\pgfsetbuttcap%
\pgfsetroundjoin%
\definecolor{currentfill}{rgb}{0.898378,0.898378,0.898378}%
\pgfsetfillcolor{currentfill}%
\pgfsetfillopacity{0.900000}%
\pgfsetlinewidth{0.501875pt}%
\definecolor{currentstroke}{rgb}{0.200000,0.200000,0.200000}%
\pgfsetstrokecolor{currentstroke}%
\pgfsetdash{}{0pt}%
\pgfpathmoveto{\pgfqpoint{2.456524in}{2.165355in}}%
\pgfpathlineto{\pgfqpoint{2.501395in}{2.125906in}}%
\pgfpathlineto{\pgfqpoint{2.532608in}{2.106501in}}%
\pgfpathlineto{\pgfqpoint{2.487732in}{2.145655in}}%
\pgfpathlineto{\pgfqpoint{2.456524in}{2.165355in}}%
\pgfpathclose%
\pgfusepath{stroke,fill}%
\end{pgfscope}%
\begin{pgfscope}%
\pgfpathrectangle{\pgfqpoint{0.050000in}{0.050000in}}{\pgfqpoint{4.900000in}{4.900000in}}%
\pgfusepath{clip}%
\pgfsetbuttcap%
\pgfsetroundjoin%
\definecolor{currentfill}{rgb}{0.964937,0.964937,0.964937}%
\pgfsetfillcolor{currentfill}%
\pgfsetfillopacity{0.900000}%
\pgfsetlinewidth{0.501875pt}%
\definecolor{currentstroke}{rgb}{0.200000,0.200000,0.200000}%
\pgfsetstrokecolor{currentstroke}%
\pgfsetdash{}{0pt}%
\pgfpathmoveto{\pgfqpoint{2.868132in}{1.890295in}}%
\pgfpathlineto{\pgfqpoint{2.912288in}{1.871162in}}%
\pgfpathlineto{\pgfqpoint{2.944147in}{1.852892in}}%
\pgfpathlineto{\pgfqpoint{2.899973in}{1.871841in}}%
\pgfpathlineto{\pgfqpoint{2.868132in}{1.890295in}}%
\pgfpathclose%
\pgfusepath{stroke,fill}%
\end{pgfscope}%
\begin{pgfscope}%
\pgfpathrectangle{\pgfqpoint{0.050000in}{0.050000in}}{\pgfqpoint{4.900000in}{4.900000in}}%
\pgfusepath{clip}%
\pgfsetbuttcap%
\pgfsetroundjoin%
\definecolor{currentfill}{rgb}{0.832895,0.832895,0.832895}%
\pgfsetfillcolor{currentfill}%
\pgfsetfillopacity{0.900000}%
\pgfsetlinewidth{0.501875pt}%
\definecolor{currentstroke}{rgb}{0.200000,0.200000,0.200000}%
\pgfsetstrokecolor{currentstroke}%
\pgfsetdash{}{0pt}%
\pgfpathmoveto{\pgfqpoint{2.196960in}{2.371729in}}%
\pgfpathlineto{\pgfqpoint{2.242376in}{2.328535in}}%
\pgfpathlineto{\pgfqpoint{2.273197in}{2.308629in}}%
\pgfpathlineto{\pgfqpoint{2.227777in}{2.351778in}}%
\pgfpathlineto{\pgfqpoint{2.196960in}{2.371729in}}%
\pgfpathclose%
\pgfusepath{stroke,fill}%
\end{pgfscope}%
\begin{pgfscope}%
\pgfpathrectangle{\pgfqpoint{0.050000in}{0.050000in}}{\pgfqpoint{4.900000in}{4.900000in}}%
\pgfusepath{clip}%
\pgfsetbuttcap%
\pgfsetroundjoin%
\definecolor{currentfill}{rgb}{0.983391,0.983391,0.983391}%
\pgfsetfillcolor{currentfill}%
\pgfsetfillopacity{0.900000}%
\pgfsetlinewidth{0.501875pt}%
\definecolor{currentstroke}{rgb}{0.200000,0.200000,0.200000}%
\pgfsetstrokecolor{currentstroke}%
\pgfsetdash{}{0pt}%
\pgfpathmoveto{\pgfqpoint{3.096264in}{1.792214in}}%
\pgfpathlineto{\pgfqpoint{3.140177in}{1.786630in}}%
\pgfpathlineto{\pgfqpoint{3.172396in}{1.768153in}}%
\pgfpathlineto{\pgfqpoint{3.128461in}{1.773780in}}%
\pgfpathlineto{\pgfqpoint{3.096264in}{1.792214in}}%
\pgfpathclose%
\pgfusepath{stroke,fill}%
\end{pgfscope}%
\begin{pgfscope}%
\pgfpathrectangle{\pgfqpoint{0.050000in}{0.050000in}}{\pgfqpoint{4.900000in}{4.900000in}}%
\pgfusepath{clip}%
\pgfsetbuttcap%
\pgfsetroundjoin%
\definecolor{currentfill}{rgb}{0.957555,0.957555,0.957555}%
\pgfsetfillcolor{currentfill}%
\pgfsetfillopacity{0.900000}%
\pgfsetlinewidth{0.501875pt}%
\definecolor{currentstroke}{rgb}{0.200000,0.200000,0.200000}%
\pgfsetstrokecolor{currentstroke}%
\pgfsetdash{}{0pt}%
\pgfpathmoveto{\pgfqpoint{2.792121in}{1.932653in}}%
\pgfpathlineto{\pgfqpoint{2.836389in}{1.908661in}}%
\pgfpathlineto{\pgfqpoint{2.868132in}{1.890295in}}%
\pgfpathlineto{\pgfqpoint{2.823849in}{1.914023in}}%
\pgfpathlineto{\pgfqpoint{2.792121in}{1.932653in}}%
\pgfpathclose%
\pgfusepath{stroke,fill}%
\end{pgfscope}%
\begin{pgfscope}%
\pgfpathrectangle{\pgfqpoint{0.050000in}{0.050000in}}{\pgfqpoint{4.900000in}{4.900000in}}%
\pgfusepath{clip}%
\pgfsetbuttcap%
\pgfsetroundjoin%
\definecolor{currentfill}{rgb}{0.757109,0.757109,0.757109}%
\pgfsetfillcolor{currentfill}%
\pgfsetfillopacity{0.900000}%
\pgfsetlinewidth{0.501875pt}%
\definecolor{currentstroke}{rgb}{0.200000,0.200000,0.200000}%
\pgfsetstrokecolor{currentstroke}%
\pgfsetdash{}{0pt}%
\pgfpathmoveto{\pgfqpoint{1.937304in}{2.582075in}}%
\pgfpathlineto{\pgfqpoint{1.983268in}{2.538209in}}%
\pgfpathlineto{\pgfqpoint{2.013701in}{2.518495in}}%
\pgfpathlineto{\pgfqpoint{1.967731in}{2.562386in}}%
\pgfpathlineto{\pgfqpoint{1.937304in}{2.582075in}}%
\pgfpathclose%
\pgfusepath{stroke,fill}%
\end{pgfscope}%
\begin{pgfscope}%
\pgfpathrectangle{\pgfqpoint{0.050000in}{0.050000in}}{\pgfqpoint{4.900000in}{4.900000in}}%
\pgfusepath{clip}%
\pgfsetbuttcap%
\pgfsetroundjoin%
\definecolor{currentfill}{rgb}{0.998155,0.998155,0.998155}%
\pgfsetfillcolor{currentfill}%
\pgfsetfillopacity{0.900000}%
\pgfsetlinewidth{0.501875pt}%
\definecolor{currentstroke}{rgb}{0.200000,0.200000,0.200000}%
\pgfsetstrokecolor{currentstroke}%
\pgfsetdash{}{0pt}%
\pgfpathmoveto{\pgfqpoint{3.521275in}{1.709768in}}%
\pgfpathlineto{\pgfqpoint{3.564852in}{1.719236in}}%
\pgfpathlineto{\pgfqpoint{3.597726in}{1.699681in}}%
\pgfpathlineto{\pgfqpoint{3.554122in}{1.690327in}}%
\pgfpathlineto{\pgfqpoint{3.521275in}{1.709768in}}%
\pgfpathclose%
\pgfusepath{stroke,fill}%
\end{pgfscope}%
\begin{pgfscope}%
\pgfpathrectangle{\pgfqpoint{0.050000in}{0.050000in}}{\pgfqpoint{4.900000in}{4.900000in}}%
\pgfusepath{clip}%
\pgfsetbuttcap%
\pgfsetroundjoin%
\definecolor{currentfill}{rgb}{0.987082,0.987082,0.987082}%
\pgfsetfillcolor{currentfill}%
\pgfsetfillopacity{0.900000}%
\pgfsetlinewidth{0.501875pt}%
\definecolor{currentstroke}{rgb}{0.200000,0.200000,0.200000}%
\pgfsetstrokecolor{currentstroke}%
\pgfsetdash{}{0pt}%
\pgfpathmoveto{\pgfqpoint{3.172396in}{1.768153in}}%
\pgfpathlineto{\pgfqpoint{3.216249in}{1.766259in}}%
\pgfpathlineto{\pgfqpoint{3.248592in}{1.747605in}}%
\pgfpathlineto{\pgfqpoint{3.204716in}{1.749588in}}%
\pgfpathlineto{\pgfqpoint{3.172396in}{1.768153in}}%
\pgfpathclose%
\pgfusepath{stroke,fill}%
\end{pgfscope}%
\begin{pgfscope}%
\pgfpathrectangle{\pgfqpoint{0.050000in}{0.050000in}}{\pgfqpoint{4.900000in}{4.900000in}}%
\pgfusepath{clip}%
\pgfsetbuttcap%
\pgfsetroundjoin%
\definecolor{currentfill}{rgb}{1.000000,1.000000,1.000000}%
\pgfsetfillcolor{currentfill}%
\pgfsetfillopacity{0.900000}%
\pgfsetlinewidth{0.501875pt}%
\definecolor{currentstroke}{rgb}{0.200000,0.200000,0.200000}%
\pgfsetstrokecolor{currentstroke}%
\pgfsetdash{}{0pt}%
\pgfpathmoveto{\pgfqpoint{3.794280in}{1.694066in}}%
\pgfpathlineto{\pgfqpoint{3.837615in}{1.706547in}}%
\pgfpathlineto{\pgfqpoint{3.870903in}{1.686506in}}%
\pgfpathlineto{\pgfqpoint{3.827541in}{1.674087in}}%
\pgfpathlineto{\pgfqpoint{3.794280in}{1.694066in}}%
\pgfpathclose%
\pgfusepath{stroke,fill}%
\end{pgfscope}%
\begin{pgfscope}%
\pgfpathrectangle{\pgfqpoint{0.050000in}{0.050000in}}{\pgfqpoint{4.900000in}{4.900000in}}%
\pgfusepath{clip}%
\pgfsetbuttcap%
\pgfsetroundjoin%
\definecolor{currentfill}{rgb}{0.948328,0.948328,0.948328}%
\pgfsetfillcolor{currentfill}%
\pgfsetfillopacity{0.900000}%
\pgfsetlinewidth{0.501875pt}%
\definecolor{currentstroke}{rgb}{0.200000,0.200000,0.200000}%
\pgfsetstrokecolor{currentstroke}%
\pgfsetdash{}{0pt}%
\pgfpathmoveto{\pgfqpoint{2.716095in}{1.979859in}}%
\pgfpathlineto{\pgfqpoint{2.760492in}{1.951207in}}%
\pgfpathlineto{\pgfqpoint{2.792121in}{1.932653in}}%
\pgfpathlineto{\pgfqpoint{2.747712in}{1.960979in}}%
\pgfpathlineto{\pgfqpoint{2.716095in}{1.979859in}}%
\pgfpathclose%
\pgfusepath{stroke,fill}%
\end{pgfscope}%
\begin{pgfscope}%
\pgfpathrectangle{\pgfqpoint{0.050000in}{0.050000in}}{\pgfqpoint{4.900000in}{4.900000in}}%
\pgfusepath{clip}%
\pgfsetbuttcap%
\pgfsetroundjoin%
\definecolor{currentfill}{rgb}{0.653010,0.653010,0.653010}%
\pgfsetfillcolor{currentfill}%
\pgfsetfillopacity{0.900000}%
\pgfsetlinewidth{0.501875pt}%
\definecolor{currentstroke}{rgb}{0.200000,0.200000,0.200000}%
\pgfsetstrokecolor{currentstroke}%
\pgfsetdash{}{0pt}%
\pgfpathmoveto{\pgfqpoint{1.677562in}{2.792625in}}%
\pgfpathlineto{\pgfqpoint{1.724071in}{2.748314in}}%
\pgfpathlineto{\pgfqpoint{1.754117in}{2.728914in}}%
\pgfpathlineto{\pgfqpoint{1.707599in}{2.773252in}}%
\pgfpathlineto{\pgfqpoint{1.677562in}{2.792625in}}%
\pgfpathclose%
\pgfusepath{stroke,fill}%
\end{pgfscope}%
\begin{pgfscope}%
\pgfpathrectangle{\pgfqpoint{0.050000in}{0.050000in}}{\pgfqpoint{4.900000in}{4.900000in}}%
\pgfusepath{clip}%
\pgfsetbuttcap%
\pgfsetroundjoin%
\definecolor{currentfill}{rgb}{0.990773,0.990773,0.990773}%
\pgfsetfillcolor{currentfill}%
\pgfsetfillopacity{0.900000}%
\pgfsetlinewidth{0.501875pt}%
\definecolor{currentstroke}{rgb}{0.200000,0.200000,0.200000}%
\pgfsetstrokecolor{currentstroke}%
\pgfsetdash{}{0pt}%
\pgfpathmoveto{\pgfqpoint{3.248592in}{1.747605in}}%
\pgfpathlineto{\pgfqpoint{3.292391in}{1.748888in}}%
\pgfpathlineto{\pgfqpoint{3.324859in}{1.730035in}}%
\pgfpathlineto{\pgfqpoint{3.281035in}{1.728870in}}%
\pgfpathlineto{\pgfqpoint{3.248592in}{1.747605in}}%
\pgfpathclose%
\pgfusepath{stroke,fill}%
\end{pgfscope}%
\begin{pgfscope}%
\pgfpathrectangle{\pgfqpoint{0.050000in}{0.050000in}}{\pgfqpoint{4.900000in}{4.900000in}}%
\pgfusepath{clip}%
\pgfsetbuttcap%
\pgfsetroundjoin%
\definecolor{currentfill}{rgb}{0.878570,0.878570,0.878570}%
\pgfsetfillcolor{currentfill}%
\pgfsetfillopacity{0.900000}%
\pgfsetlinewidth{0.501875pt}%
\definecolor{currentstroke}{rgb}{0.200000,0.200000,0.200000}%
\pgfsetstrokecolor{currentstroke}%
\pgfsetdash{}{0pt}%
\pgfpathmoveto{\pgfqpoint{2.380361in}{2.226345in}}%
\pgfpathlineto{\pgfqpoint{2.425412in}{2.185028in}}%
\pgfpathlineto{\pgfqpoint{2.456524in}{2.165355in}}%
\pgfpathlineto{\pgfqpoint{2.411468in}{2.206460in}}%
\pgfpathlineto{\pgfqpoint{2.380361in}{2.226345in}}%
\pgfpathclose%
\pgfusepath{stroke,fill}%
\end{pgfscope}%
\begin{pgfscope}%
\pgfpathrectangle{\pgfqpoint{0.050000in}{0.050000in}}{\pgfqpoint{4.900000in}{4.900000in}}%
\pgfusepath{clip}%
\pgfsetbuttcap%
\pgfsetroundjoin%
\definecolor{currentfill}{rgb}{0.547328,0.547328,0.547328}%
\pgfsetfillcolor{currentfill}%
\pgfsetfillopacity{0.900000}%
\pgfsetlinewidth{0.501875pt}%
\definecolor{currentstroke}{rgb}{0.200000,0.200000,0.200000}%
\pgfsetstrokecolor{currentstroke}%
\pgfsetdash{}{0pt}%
\pgfpathmoveto{\pgfqpoint{1.417733in}{3.003247in}}%
\pgfpathlineto{\pgfqpoint{1.464790in}{2.958491in}}%
\pgfpathlineto{\pgfqpoint{1.494446in}{2.939406in}}%
\pgfpathlineto{\pgfqpoint{1.447380in}{2.984190in}}%
\pgfpathlineto{\pgfqpoint{1.417733in}{3.003247in}}%
\pgfpathclose%
\pgfusepath{stroke,fill}%
\end{pgfscope}%
\begin{pgfscope}%
\pgfpathrectangle{\pgfqpoint{0.050000in}{0.050000in}}{\pgfqpoint{4.900000in}{4.900000in}}%
\pgfusepath{clip}%
\pgfsetbuttcap%
\pgfsetroundjoin%
\definecolor{currentfill}{rgb}{0.808781,0.808781,0.808781}%
\pgfsetfillcolor{currentfill}%
\pgfsetfillopacity{0.900000}%
\pgfsetlinewidth{0.501875pt}%
\definecolor{currentstroke}{rgb}{0.200000,0.200000,0.200000}%
\pgfsetstrokecolor{currentstroke}%
\pgfsetdash{}{0pt}%
\pgfpathmoveto{\pgfqpoint{2.120637in}{2.435149in}}%
\pgfpathlineto{\pgfqpoint{2.166237in}{2.391631in}}%
\pgfpathlineto{\pgfqpoint{2.196960in}{2.371729in}}%
\pgfpathlineto{\pgfqpoint{2.151355in}{2.415245in}}%
\pgfpathlineto{\pgfqpoint{2.120637in}{2.435149in}}%
\pgfpathclose%
\pgfusepath{stroke,fill}%
\end{pgfscope}%
\begin{pgfscope}%
\pgfpathrectangle{\pgfqpoint{0.050000in}{0.050000in}}{\pgfqpoint{4.900000in}{4.900000in}}%
\pgfusepath{clip}%
\pgfsetbuttcap%
\pgfsetroundjoin%
\definecolor{currentfill}{rgb}{0.998155,0.998155,0.998155}%
\pgfsetfillcolor{currentfill}%
\pgfsetfillopacity{0.900000}%
\pgfsetlinewidth{0.501875pt}%
\definecolor{currentstroke}{rgb}{0.200000,0.200000,0.200000}%
\pgfsetstrokecolor{currentstroke}%
\pgfsetdash{}{0pt}%
\pgfpathmoveto{\pgfqpoint{3.597726in}{1.699681in}}%
\pgfpathlineto{\pgfqpoint{3.641253in}{1.710240in}}%
\pgfpathlineto{\pgfqpoint{3.674255in}{1.690529in}}%
\pgfpathlineto{\pgfqpoint{3.630702in}{1.680069in}}%
\pgfpathlineto{\pgfqpoint{3.597726in}{1.699681in}}%
\pgfpathclose%
\pgfusepath{stroke,fill}%
\end{pgfscope}%
\begin{pgfscope}%
\pgfpathrectangle{\pgfqpoint{0.050000in}{0.050000in}}{\pgfqpoint{4.900000in}{4.900000in}}%
\pgfusepath{clip}%
\pgfsetbuttcap%
\pgfsetroundjoin%
\definecolor{currentfill}{rgb}{0.935163,0.935163,0.935163}%
\pgfsetfillcolor{currentfill}%
\pgfsetfillopacity{0.900000}%
\pgfsetlinewidth{0.501875pt}%
\definecolor{currentstroke}{rgb}{0.200000,0.200000,0.200000}%
\pgfsetstrokecolor{currentstroke}%
\pgfsetdash{}{0pt}%
\pgfpathmoveto{\pgfqpoint{2.640033in}{2.031536in}}%
\pgfpathlineto{\pgfqpoint{2.684576in}{1.998678in}}%
\pgfpathlineto{\pgfqpoint{2.716095in}{1.979859in}}%
\pgfpathlineto{\pgfqpoint{2.671542in}{2.012364in}}%
\pgfpathlineto{\pgfqpoint{2.640033in}{2.031536in}}%
\pgfpathclose%
\pgfusepath{stroke,fill}%
\end{pgfscope}%
\begin{pgfscope}%
\pgfpathrectangle{\pgfqpoint{0.050000in}{0.050000in}}{\pgfqpoint{4.900000in}{4.900000in}}%
\pgfusepath{clip}%
\pgfsetbuttcap%
\pgfsetroundjoin%
\definecolor{currentfill}{rgb}{0.992618,0.992618,0.992618}%
\pgfsetfillcolor{currentfill}%
\pgfsetfillopacity{0.900000}%
\pgfsetlinewidth{0.501875pt}%
\definecolor{currentstroke}{rgb}{0.200000,0.200000,0.200000}%
\pgfsetstrokecolor{currentstroke}%
\pgfsetdash{}{0pt}%
\pgfpathmoveto{\pgfqpoint{3.324859in}{1.730035in}}%
\pgfpathlineto{\pgfqpoint{3.368608in}{1.733989in}}%
\pgfpathlineto{\pgfqpoint{3.401201in}{1.714932in}}%
\pgfpathlineto{\pgfqpoint{3.357427in}{1.711109in}}%
\pgfpathlineto{\pgfqpoint{3.324859in}{1.730035in}}%
\pgfpathclose%
\pgfusepath{stroke,fill}%
\end{pgfscope}%
\begin{pgfscope}%
\pgfpathrectangle{\pgfqpoint{0.050000in}{0.050000in}}{\pgfqpoint{4.900000in}{4.900000in}}%
\pgfusepath{clip}%
\pgfsetbuttcap%
\pgfsetroundjoin%
\definecolor{currentfill}{rgb}{0.724983,0.724983,0.724983}%
\pgfsetfillcolor{currentfill}%
\pgfsetfillopacity{0.900000}%
\pgfsetlinewidth{0.501875pt}%
\definecolor{currentstroke}{rgb}{0.200000,0.200000,0.200000}%
\pgfsetstrokecolor{currentstroke}%
\pgfsetdash{}{0pt}%
\pgfpathmoveto{\pgfqpoint{1.860822in}{2.645729in}}%
\pgfpathlineto{\pgfqpoint{1.906970in}{2.601705in}}%
\pgfpathlineto{\pgfqpoint{1.937304in}{2.582075in}}%
\pgfpathlineto{\pgfqpoint{1.891150in}{2.626125in}}%
\pgfpathlineto{\pgfqpoint{1.860822in}{2.645729in}}%
\pgfpathclose%
\pgfusepath{stroke,fill}%
\end{pgfscope}%
\begin{pgfscope}%
\pgfpathrectangle{\pgfqpoint{0.050000in}{0.050000in}}{\pgfqpoint{4.900000in}{4.900000in}}%
\pgfusepath{clip}%
\pgfsetbuttcap%
\pgfsetroundjoin%
\definecolor{currentfill}{rgb}{1.000000,1.000000,1.000000}%
\pgfsetfillcolor{currentfill}%
\pgfsetfillopacity{0.900000}%
\pgfsetlinewidth{0.501875pt}%
\definecolor{currentstroke}{rgb}{0.200000,0.200000,0.200000}%
\pgfsetstrokecolor{currentstroke}%
\pgfsetdash{}{0pt}%
\pgfpathmoveto{\pgfqpoint{3.870903in}{1.686506in}}%
\pgfpathlineto{\pgfqpoint{3.914178in}{1.699235in}}%
\pgfpathlineto{\pgfqpoint{3.947597in}{1.679082in}}%
\pgfpathlineto{\pgfqpoint{3.904295in}{1.666408in}}%
\pgfpathlineto{\pgfqpoint{3.870903in}{1.686506in}}%
\pgfpathclose%
\pgfusepath{stroke,fill}%
\end{pgfscope}%
\begin{pgfscope}%
\pgfpathrectangle{\pgfqpoint{0.050000in}{0.050000in}}{\pgfqpoint{4.900000in}{4.900000in}}%
\pgfusepath{clip}%
\pgfsetbuttcap%
\pgfsetroundjoin%
\definecolor{currentfill}{rgb}{0.614625,0.614625,0.614625}%
\pgfsetfillcolor{currentfill}%
\pgfsetfillopacity{0.900000}%
\pgfsetlinewidth{0.501875pt}%
\definecolor{currentstroke}{rgb}{0.200000,0.200000,0.200000}%
\pgfsetstrokecolor{currentstroke}%
\pgfsetdash{}{0pt}%
\pgfpathmoveto{\pgfqpoint{1.600921in}{2.856408in}}%
\pgfpathlineto{\pgfqpoint{1.647615in}{2.811940in}}%
\pgfpathlineto{\pgfqpoint{1.677562in}{2.792625in}}%
\pgfpathlineto{\pgfqpoint{1.630858in}{2.837121in}}%
\pgfpathlineto{\pgfqpoint{1.600921in}{2.856408in}}%
\pgfpathclose%
\pgfusepath{stroke,fill}%
\end{pgfscope}%
\begin{pgfscope}%
\pgfpathrectangle{\pgfqpoint{0.050000in}{0.050000in}}{\pgfqpoint{4.900000in}{4.900000in}}%
\pgfusepath{clip}%
\pgfsetbuttcap%
\pgfsetroundjoin%
\definecolor{currentfill}{rgb}{0.858762,0.858762,0.858762}%
\pgfsetfillcolor{currentfill}%
\pgfsetfillopacity{0.900000}%
\pgfsetlinewidth{0.501875pt}%
\definecolor{currentstroke}{rgb}{0.200000,0.200000,0.200000}%
\pgfsetstrokecolor{currentstroke}%
\pgfsetdash{}{0pt}%
\pgfpathmoveto{\pgfqpoint{2.304112in}{2.288686in}}%
\pgfpathlineto{\pgfqpoint{2.349348in}{2.246201in}}%
\pgfpathlineto{\pgfqpoint{2.380361in}{2.226345in}}%
\pgfpathlineto{\pgfqpoint{2.335121in}{2.268705in}}%
\pgfpathlineto{\pgfqpoint{2.304112in}{2.288686in}}%
\pgfpathclose%
\pgfusepath{stroke,fill}%
\end{pgfscope}%
\begin{pgfscope}%
\pgfpathrectangle{\pgfqpoint{0.050000in}{0.050000in}}{\pgfqpoint{4.900000in}{4.900000in}}%
\pgfusepath{clip}%
\pgfsetbuttcap%
\pgfsetroundjoin%
\definecolor{currentfill}{rgb}{0.998155,0.998155,0.998155}%
\pgfsetfillcolor{currentfill}%
\pgfsetfillopacity{0.900000}%
\pgfsetlinewidth{0.501875pt}%
\definecolor{currentstroke}{rgb}{0.200000,0.200000,0.200000}%
\pgfsetstrokecolor{currentstroke}%
\pgfsetdash{}{0pt}%
\pgfpathmoveto{\pgfqpoint{3.674255in}{1.690529in}}%
\pgfpathlineto{\pgfqpoint{3.717729in}{1.701919in}}%
\pgfpathlineto{\pgfqpoint{3.750861in}{1.682067in}}%
\pgfpathlineto{\pgfqpoint{3.707360in}{1.670763in}}%
\pgfpathlineto{\pgfqpoint{3.674255in}{1.690529in}}%
\pgfpathclose%
\pgfusepath{stroke,fill}%
\end{pgfscope}%
\begin{pgfscope}%
\pgfpathrectangle{\pgfqpoint{0.050000in}{0.050000in}}{\pgfqpoint{4.900000in}{4.900000in}}%
\pgfusepath{clip}%
\pgfsetbuttcap%
\pgfsetroundjoin%
\definecolor{currentfill}{rgb}{0.994464,0.994464,0.994464}%
\pgfsetfillcolor{currentfill}%
\pgfsetfillopacity{0.900000}%
\pgfsetlinewidth{0.501875pt}%
\definecolor{currentstroke}{rgb}{0.200000,0.200000,0.200000}%
\pgfsetstrokecolor{currentstroke}%
\pgfsetdash{}{0pt}%
\pgfpathmoveto{\pgfqpoint{3.401201in}{1.714932in}}%
\pgfpathlineto{\pgfqpoint{3.444902in}{1.721087in}}%
\pgfpathlineto{\pgfqpoint{3.477622in}{1.701832in}}%
\pgfpathlineto{\pgfqpoint{3.433896in}{1.695809in}}%
\pgfpathlineto{\pgfqpoint{3.401201in}{1.714932in}}%
\pgfpathclose%
\pgfusepath{stroke,fill}%
\end{pgfscope}%
\begin{pgfscope}%
\pgfpathrectangle{\pgfqpoint{0.050000in}{0.050000in}}{\pgfqpoint{4.900000in}{4.900000in}}%
\pgfusepath{clip}%
\pgfsetbuttcap%
\pgfsetroundjoin%
\definecolor{currentfill}{rgb}{0.918185,0.918185,0.918185}%
\pgfsetfillcolor{currentfill}%
\pgfsetfillopacity{0.900000}%
\pgfsetlinewidth{0.501875pt}%
\definecolor{currentstroke}{rgb}{0.200000,0.200000,0.200000}%
\pgfsetstrokecolor{currentstroke}%
\pgfsetdash{}{0pt}%
\pgfpathmoveto{\pgfqpoint{2.563917in}{2.087065in}}%
\pgfpathlineto{\pgfqpoint{2.608620in}{2.050663in}}%
\pgfpathlineto{\pgfqpoint{2.640033in}{2.031536in}}%
\pgfpathlineto{\pgfqpoint{2.595322in}{2.067597in}}%
\pgfpathlineto{\pgfqpoint{2.563917in}{2.087065in}}%
\pgfpathclose%
\pgfusepath{stroke,fill}%
\end{pgfscope}%
\begin{pgfscope}%
\pgfpathrectangle{\pgfqpoint{0.050000in}{0.050000in}}{\pgfqpoint{4.900000in}{4.900000in}}%
\pgfusepath{clip}%
\pgfsetbuttcap%
\pgfsetroundjoin%
\definecolor{currentfill}{rgb}{0.781223,0.781223,0.781223}%
\pgfsetfillcolor{currentfill}%
\pgfsetfillopacity{0.900000}%
\pgfsetlinewidth{0.501875pt}%
\definecolor{currentstroke}{rgb}{0.200000,0.200000,0.200000}%
\pgfsetstrokecolor{currentstroke}%
\pgfsetdash{}{0pt}%
\pgfpathmoveto{\pgfqpoint{2.044227in}{2.498722in}}%
\pgfpathlineto{\pgfqpoint{2.090012in}{2.454998in}}%
\pgfpathlineto{\pgfqpoint{2.120637in}{2.435149in}}%
\pgfpathlineto{\pgfqpoint{2.074846in}{2.478890in}}%
\pgfpathlineto{\pgfqpoint{2.044227in}{2.498722in}}%
\pgfpathclose%
\pgfusepath{stroke,fill}%
\end{pgfscope}%
\begin{pgfscope}%
\pgfpathrectangle{\pgfqpoint{0.050000in}{0.050000in}}{\pgfqpoint{4.900000in}{4.900000in}}%
\pgfusepath{clip}%
\pgfsetbuttcap%
\pgfsetroundjoin%
\definecolor{currentfill}{rgb}{1.000000,1.000000,1.000000}%
\pgfsetfillcolor{currentfill}%
\pgfsetfillopacity{0.900000}%
\pgfsetlinewidth{0.501875pt}%
\definecolor{currentstroke}{rgb}{0.200000,0.200000,0.200000}%
\pgfsetstrokecolor{currentstroke}%
\pgfsetdash{}{0pt}%
\pgfpathmoveto{\pgfqpoint{3.947597in}{1.679082in}}%
\pgfpathlineto{\pgfqpoint{3.990808in}{1.691912in}}%
\pgfpathlineto{\pgfqpoint{4.024356in}{1.671648in}}%
\pgfpathlineto{\pgfqpoint{3.981118in}{1.658868in}}%
\pgfpathlineto{\pgfqpoint{3.947597in}{1.679082in}}%
\pgfpathclose%
\pgfusepath{stroke,fill}%
\end{pgfscope}%
\begin{pgfscope}%
\pgfpathrectangle{\pgfqpoint{0.050000in}{0.050000in}}{\pgfqpoint{4.900000in}{4.900000in}}%
\pgfusepath{clip}%
\pgfsetbuttcap%
\pgfsetroundjoin%
\definecolor{currentfill}{rgb}{0.686597,0.686597,0.686597}%
\pgfsetfillcolor{currentfill}%
\pgfsetfillopacity{0.900000}%
\pgfsetlinewidth{0.501875pt}%
\definecolor{currentstroke}{rgb}{0.200000,0.200000,0.200000}%
\pgfsetstrokecolor{currentstroke}%
\pgfsetdash{}{0pt}%
\pgfpathmoveto{\pgfqpoint{1.784253in}{2.709455in}}%
\pgfpathlineto{\pgfqpoint{1.830586in}{2.665274in}}%
\pgfpathlineto{\pgfqpoint{1.860822in}{2.645729in}}%
\pgfpathlineto{\pgfqpoint{1.814482in}{2.689936in}}%
\pgfpathlineto{\pgfqpoint{1.784253in}{2.709455in}}%
\pgfpathclose%
\pgfusepath{stroke,fill}%
\end{pgfscope}%
\begin{pgfscope}%
\pgfpathrectangle{\pgfqpoint{0.050000in}{0.050000in}}{\pgfqpoint{4.900000in}{4.900000in}}%
\pgfusepath{clip}%
\pgfsetbuttcap%
\pgfsetroundjoin%
\definecolor{currentfill}{rgb}{0.996309,0.996309,0.996309}%
\pgfsetfillcolor{currentfill}%
\pgfsetfillopacity{0.900000}%
\pgfsetlinewidth{0.501875pt}%
\definecolor{currentstroke}{rgb}{0.200000,0.200000,0.200000}%
\pgfsetstrokecolor{currentstroke}%
\pgfsetdash{}{0pt}%
\pgfpathmoveto{\pgfqpoint{3.477622in}{1.701832in}}%
\pgfpathlineto{\pgfqpoint{3.521275in}{1.709768in}}%
\pgfpathlineto{\pgfqpoint{3.554122in}{1.690327in}}%
\pgfpathlineto{\pgfqpoint{3.510443in}{1.682515in}}%
\pgfpathlineto{\pgfqpoint{3.477622in}{1.701832in}}%
\pgfpathclose%
\pgfusepath{stroke,fill}%
\end{pgfscope}%
\begin{pgfscope}%
\pgfpathrectangle{\pgfqpoint{0.050000in}{0.050000in}}{\pgfqpoint{4.900000in}{4.900000in}}%
\pgfusepath{clip}%
\pgfsetbuttcap%
\pgfsetroundjoin%
\definecolor{currentfill}{rgb}{1.000000,1.000000,1.000000}%
\pgfsetfillcolor{currentfill}%
\pgfsetfillopacity{0.900000}%
\pgfsetlinewidth{0.501875pt}%
\definecolor{currentstroke}{rgb}{0.200000,0.200000,0.200000}%
\pgfsetstrokecolor{currentstroke}%
\pgfsetdash{}{0pt}%
\pgfpathmoveto{\pgfqpoint{3.750861in}{1.682067in}}%
\pgfpathlineto{\pgfqpoint{3.794280in}{1.694066in}}%
\pgfpathlineto{\pgfqpoint{3.827541in}{1.674087in}}%
\pgfpathlineto{\pgfqpoint{3.784095in}{1.662160in}}%
\pgfpathlineto{\pgfqpoint{3.750861in}{1.682067in}}%
\pgfpathclose%
\pgfusepath{stroke,fill}%
\end{pgfscope}%
\begin{pgfscope}%
\pgfpathrectangle{\pgfqpoint{0.050000in}{0.050000in}}{\pgfqpoint{4.900000in}{4.900000in}}%
\pgfusepath{clip}%
\pgfsetbuttcap%
\pgfsetroundjoin%
\definecolor{currentfill}{rgb}{0.581776,0.581776,0.581776}%
\pgfsetfillcolor{currentfill}%
\pgfsetfillopacity{0.900000}%
\pgfsetlinewidth{0.501875pt}%
\definecolor{currentstroke}{rgb}{0.200000,0.200000,0.200000}%
\pgfsetstrokecolor{currentstroke}%
\pgfsetdash{}{0pt}%
\pgfpathmoveto{\pgfqpoint{1.524193in}{2.920263in}}%
\pgfpathlineto{\pgfqpoint{1.571074in}{2.875637in}}%
\pgfpathlineto{\pgfqpoint{1.600921in}{2.856408in}}%
\pgfpathlineto{\pgfqpoint{1.554031in}{2.901062in}}%
\pgfpathlineto{\pgfqpoint{1.524193in}{2.920263in}}%
\pgfpathclose%
\pgfusepath{stroke,fill}%
\end{pgfscope}%
\begin{pgfscope}%
\pgfpathrectangle{\pgfqpoint{0.050000in}{0.050000in}}{\pgfqpoint{4.900000in}{4.900000in}}%
\pgfusepath{clip}%
\pgfsetbuttcap%
\pgfsetroundjoin%
\definecolor{currentfill}{rgb}{0.972318,0.972318,0.972318}%
\pgfsetfillcolor{currentfill}%
\pgfsetfillopacity{0.900000}%
\pgfsetlinewidth{0.501875pt}%
\definecolor{currentstroke}{rgb}{0.200000,0.200000,0.200000}%
\pgfsetstrokecolor{currentstroke}%
\pgfsetdash{}{0pt}%
\pgfpathmoveto{\pgfqpoint{2.976104in}{1.834528in}}%
\pgfpathlineto{\pgfqpoint{3.020185in}{1.820308in}}%
\pgfpathlineto{\pgfqpoint{3.052262in}{1.801947in}}%
\pgfpathlineto{\pgfqpoint{3.008162in}{1.816070in}}%
\pgfpathlineto{\pgfqpoint{2.976104in}{1.834528in}}%
\pgfpathclose%
\pgfusepath{stroke,fill}%
\end{pgfscope}%
\begin{pgfscope}%
\pgfpathrectangle{\pgfqpoint{0.050000in}{0.050000in}}{\pgfqpoint{4.900000in}{4.900000in}}%
\pgfusepath{clip}%
\pgfsetbuttcap%
\pgfsetroundjoin%
\definecolor{currentfill}{rgb}{0.898378,0.898378,0.898378}%
\pgfsetfillcolor{currentfill}%
\pgfsetfillopacity{0.900000}%
\pgfsetlinewidth{0.501875pt}%
\definecolor{currentstroke}{rgb}{0.200000,0.200000,0.200000}%
\pgfsetstrokecolor{currentstroke}%
\pgfsetdash{}{0pt}%
\pgfpathmoveto{\pgfqpoint{2.487732in}{2.145655in}}%
\pgfpathlineto{\pgfqpoint{2.532608in}{2.106501in}}%
\pgfpathlineto{\pgfqpoint{2.563917in}{2.087065in}}%
\pgfpathlineto{\pgfqpoint{2.519035in}{2.125928in}}%
\pgfpathlineto{\pgfqpoint{2.487732in}{2.145655in}}%
\pgfpathclose%
\pgfusepath{stroke,fill}%
\end{pgfscope}%
\begin{pgfscope}%
\pgfpathrectangle{\pgfqpoint{0.050000in}{0.050000in}}{\pgfqpoint{4.900000in}{4.900000in}}%
\pgfusepath{clip}%
\pgfsetbuttcap%
\pgfsetroundjoin%
\definecolor{currentfill}{rgb}{0.977855,0.977855,0.977855}%
\pgfsetfillcolor{currentfill}%
\pgfsetfillopacity{0.900000}%
\pgfsetlinewidth{0.501875pt}%
\definecolor{currentstroke}{rgb}{0.200000,0.200000,0.200000}%
\pgfsetstrokecolor{currentstroke}%
\pgfsetdash{}{0pt}%
\pgfpathmoveto{\pgfqpoint{3.052262in}{1.801947in}}%
\pgfpathlineto{\pgfqpoint{3.096264in}{1.792214in}}%
\pgfpathlineto{\pgfqpoint{3.128461in}{1.773780in}}%
\pgfpathlineto{\pgfqpoint{3.084438in}{1.783490in}}%
\pgfpathlineto{\pgfqpoint{3.052262in}{1.801947in}}%
\pgfpathclose%
\pgfusepath{stroke,fill}%
\end{pgfscope}%
\begin{pgfscope}%
\pgfpathrectangle{\pgfqpoint{0.050000in}{0.050000in}}{\pgfqpoint{4.900000in}{4.900000in}}%
\pgfusepath{clip}%
\pgfsetbuttcap%
\pgfsetroundjoin%
\definecolor{currentfill}{rgb}{0.964937,0.964937,0.964937}%
\pgfsetfillcolor{currentfill}%
\pgfsetfillopacity{0.900000}%
\pgfsetlinewidth{0.501875pt}%
\definecolor{currentstroke}{rgb}{0.200000,0.200000,0.200000}%
\pgfsetstrokecolor{currentstroke}%
\pgfsetdash{}{0pt}%
\pgfpathmoveto{\pgfqpoint{2.899973in}{1.871841in}}%
\pgfpathlineto{\pgfqpoint{2.944147in}{1.852892in}}%
\pgfpathlineto{\pgfqpoint{2.976104in}{1.834528in}}%
\pgfpathlineto{\pgfqpoint{2.931914in}{1.853300in}}%
\pgfpathlineto{\pgfqpoint{2.899973in}{1.871841in}}%
\pgfpathclose%
\pgfusepath{stroke,fill}%
\end{pgfscope}%
\begin{pgfscope}%
\pgfpathrectangle{\pgfqpoint{0.050000in}{0.050000in}}{\pgfqpoint{4.900000in}{4.900000in}}%
\pgfusepath{clip}%
\pgfsetbuttcap%
\pgfsetroundjoin%
\definecolor{currentfill}{rgb}{0.832895,0.832895,0.832895}%
\pgfsetfillcolor{currentfill}%
\pgfsetfillopacity{0.900000}%
\pgfsetlinewidth{0.501875pt}%
\definecolor{currentstroke}{rgb}{0.200000,0.200000,0.200000}%
\pgfsetstrokecolor{currentstroke}%
\pgfsetdash{}{0pt}%
\pgfpathmoveto{\pgfqpoint{2.227777in}{2.351778in}}%
\pgfpathlineto{\pgfqpoint{2.273197in}{2.308629in}}%
\pgfpathlineto{\pgfqpoint{2.304112in}{2.288686in}}%
\pgfpathlineto{\pgfqpoint{2.258687in}{2.331781in}}%
\pgfpathlineto{\pgfqpoint{2.227777in}{2.351778in}}%
\pgfpathclose%
\pgfusepath{stroke,fill}%
\end{pgfscope}%
\begin{pgfscope}%
\pgfpathrectangle{\pgfqpoint{0.050000in}{0.050000in}}{\pgfqpoint{4.900000in}{4.900000in}}%
\pgfusepath{clip}%
\pgfsetbuttcap%
\pgfsetroundjoin%
\definecolor{currentfill}{rgb}{0.983391,0.983391,0.983391}%
\pgfsetfillcolor{currentfill}%
\pgfsetfillopacity{0.900000}%
\pgfsetlinewidth{0.501875pt}%
\definecolor{currentstroke}{rgb}{0.200000,0.200000,0.200000}%
\pgfsetstrokecolor{currentstroke}%
\pgfsetdash{}{0pt}%
\pgfpathmoveto{\pgfqpoint{3.128461in}{1.773780in}}%
\pgfpathlineto{\pgfqpoint{3.172396in}{1.768153in}}%
\pgfpathlineto{\pgfqpoint{3.204716in}{1.749588in}}%
\pgfpathlineto{\pgfqpoint{3.160759in}{1.755254in}}%
\pgfpathlineto{\pgfqpoint{3.128461in}{1.773780in}}%
\pgfpathclose%
\pgfusepath{stroke,fill}%
\end{pgfscope}%
\begin{pgfscope}%
\pgfpathrectangle{\pgfqpoint{0.050000in}{0.050000in}}{\pgfqpoint{4.900000in}{4.900000in}}%
\pgfusepath{clip}%
\pgfsetbuttcap%
\pgfsetroundjoin%
\definecolor{currentfill}{rgb}{0.957555,0.957555,0.957555}%
\pgfsetfillcolor{currentfill}%
\pgfsetfillopacity{0.900000}%
\pgfsetlinewidth{0.501875pt}%
\definecolor{currentstroke}{rgb}{0.200000,0.200000,0.200000}%
\pgfsetstrokecolor{currentstroke}%
\pgfsetdash{}{0pt}%
\pgfpathmoveto{\pgfqpoint{2.823849in}{1.914023in}}%
\pgfpathlineto{\pgfqpoint{2.868132in}{1.890295in}}%
\pgfpathlineto{\pgfqpoint{2.899973in}{1.871841in}}%
\pgfpathlineto{\pgfqpoint{2.855676in}{1.895316in}}%
\pgfpathlineto{\pgfqpoint{2.823849in}{1.914023in}}%
\pgfpathclose%
\pgfusepath{stroke,fill}%
\end{pgfscope}%
\begin{pgfscope}%
\pgfpathrectangle{\pgfqpoint{0.050000in}{0.050000in}}{\pgfqpoint{4.900000in}{4.900000in}}%
\pgfusepath{clip}%
\pgfsetbuttcap%
\pgfsetroundjoin%
\definecolor{currentfill}{rgb}{0.757109,0.757109,0.757109}%
\pgfsetfillcolor{currentfill}%
\pgfsetfillopacity{0.900000}%
\pgfsetlinewidth{0.501875pt}%
\definecolor{currentstroke}{rgb}{0.200000,0.200000,0.200000}%
\pgfsetstrokecolor{currentstroke}%
\pgfsetdash{}{0pt}%
\pgfpathmoveto{\pgfqpoint{1.967731in}{2.562386in}}%
\pgfpathlineto{\pgfqpoint{2.013701in}{2.518495in}}%
\pgfpathlineto{\pgfqpoint{2.044227in}{2.498722in}}%
\pgfpathlineto{\pgfqpoint{1.998251in}{2.542637in}}%
\pgfpathlineto{\pgfqpoint{1.967731in}{2.562386in}}%
\pgfpathclose%
\pgfusepath{stroke,fill}%
\end{pgfscope}%
\begin{pgfscope}%
\pgfpathrectangle{\pgfqpoint{0.050000in}{0.050000in}}{\pgfqpoint{4.900000in}{4.900000in}}%
\pgfusepath{clip}%
\pgfsetbuttcap%
\pgfsetroundjoin%
\definecolor{currentfill}{rgb}{0.998155,0.998155,0.998155}%
\pgfsetfillcolor{currentfill}%
\pgfsetfillopacity{0.900000}%
\pgfsetlinewidth{0.501875pt}%
\definecolor{currentstroke}{rgb}{0.200000,0.200000,0.200000}%
\pgfsetstrokecolor{currentstroke}%
\pgfsetdash{}{0pt}%
\pgfpathmoveto{\pgfqpoint{3.554122in}{1.690327in}}%
\pgfpathlineto{\pgfqpoint{3.597726in}{1.699681in}}%
\pgfpathlineto{\pgfqpoint{3.630702in}{1.680069in}}%
\pgfpathlineto{\pgfqpoint{3.587071in}{1.670827in}}%
\pgfpathlineto{\pgfqpoint{3.554122in}{1.690327in}}%
\pgfpathclose%
\pgfusepath{stroke,fill}%
\end{pgfscope}%
\begin{pgfscope}%
\pgfpathrectangle{\pgfqpoint{0.050000in}{0.050000in}}{\pgfqpoint{4.900000in}{4.900000in}}%
\pgfusepath{clip}%
\pgfsetbuttcap%
\pgfsetroundjoin%
\definecolor{currentfill}{rgb}{0.987082,0.987082,0.987082}%
\pgfsetfillcolor{currentfill}%
\pgfsetfillopacity{0.900000}%
\pgfsetlinewidth{0.501875pt}%
\definecolor{currentstroke}{rgb}{0.200000,0.200000,0.200000}%
\pgfsetstrokecolor{currentstroke}%
\pgfsetdash{}{0pt}%
\pgfpathmoveto{\pgfqpoint{3.204716in}{1.749588in}}%
\pgfpathlineto{\pgfqpoint{3.248592in}{1.747605in}}%
\pgfpathlineto{\pgfqpoint{3.281035in}{1.728870in}}%
\pgfpathlineto{\pgfqpoint{3.237136in}{1.730937in}}%
\pgfpathlineto{\pgfqpoint{3.204716in}{1.749588in}}%
\pgfpathclose%
\pgfusepath{stroke,fill}%
\end{pgfscope}%
\begin{pgfscope}%
\pgfpathrectangle{\pgfqpoint{0.050000in}{0.050000in}}{\pgfqpoint{4.900000in}{4.900000in}}%
\pgfusepath{clip}%
\pgfsetbuttcap%
\pgfsetroundjoin%
\definecolor{currentfill}{rgb}{1.000000,1.000000,1.000000}%
\pgfsetfillcolor{currentfill}%
\pgfsetfillopacity{0.900000}%
\pgfsetlinewidth{0.501875pt}%
\definecolor{currentstroke}{rgb}{0.200000,0.200000,0.200000}%
\pgfsetstrokecolor{currentstroke}%
\pgfsetdash{}{0pt}%
\pgfpathmoveto{\pgfqpoint{3.827541in}{1.674087in}}%
\pgfpathlineto{\pgfqpoint{3.870903in}{1.686506in}}%
\pgfpathlineto{\pgfqpoint{3.904295in}{1.666408in}}%
\pgfpathlineto{\pgfqpoint{3.860905in}{1.654051in}}%
\pgfpathlineto{\pgfqpoint{3.827541in}{1.674087in}}%
\pgfpathclose%
\pgfusepath{stroke,fill}%
\end{pgfscope}%
\begin{pgfscope}%
\pgfpathrectangle{\pgfqpoint{0.050000in}{0.050000in}}{\pgfqpoint{4.900000in}{4.900000in}}%
\pgfusepath{clip}%
\pgfsetbuttcap%
\pgfsetroundjoin%
\definecolor{currentfill}{rgb}{0.653010,0.653010,0.653010}%
\pgfsetfillcolor{currentfill}%
\pgfsetfillopacity{0.900000}%
\pgfsetlinewidth{0.501875pt}%
\definecolor{currentstroke}{rgb}{0.200000,0.200000,0.200000}%
\pgfsetstrokecolor{currentstroke}%
\pgfsetdash{}{0pt}%
\pgfpathmoveto{\pgfqpoint{1.707599in}{2.773252in}}%
\pgfpathlineto{\pgfqpoint{1.754117in}{2.728914in}}%
\pgfpathlineto{\pgfqpoint{1.784253in}{2.709455in}}%
\pgfpathlineto{\pgfqpoint{1.737727in}{2.753820in}}%
\pgfpathlineto{\pgfqpoint{1.707599in}{2.773252in}}%
\pgfpathclose%
\pgfusepath{stroke,fill}%
\end{pgfscope}%
\begin{pgfscope}%
\pgfpathrectangle{\pgfqpoint{0.050000in}{0.050000in}}{\pgfqpoint{4.900000in}{4.900000in}}%
\pgfusepath{clip}%
\pgfsetbuttcap%
\pgfsetroundjoin%
\definecolor{currentfill}{rgb}{0.948328,0.948328,0.948328}%
\pgfsetfillcolor{currentfill}%
\pgfsetfillopacity{0.900000}%
\pgfsetlinewidth{0.501875pt}%
\definecolor{currentstroke}{rgb}{0.200000,0.200000,0.200000}%
\pgfsetstrokecolor{currentstroke}%
\pgfsetdash{}{0pt}%
\pgfpathmoveto{\pgfqpoint{2.747712in}{1.960979in}}%
\pgfpathlineto{\pgfqpoint{2.792121in}{1.932653in}}%
\pgfpathlineto{\pgfqpoint{2.823849in}{1.914023in}}%
\pgfpathlineto{\pgfqpoint{2.779427in}{1.942037in}}%
\pgfpathlineto{\pgfqpoint{2.747712in}{1.960979in}}%
\pgfpathclose%
\pgfusepath{stroke,fill}%
\end{pgfscope}%
\begin{pgfscope}%
\pgfpathrectangle{\pgfqpoint{0.050000in}{0.050000in}}{\pgfqpoint{4.900000in}{4.900000in}}%
\pgfusepath{clip}%
\pgfsetbuttcap%
\pgfsetroundjoin%
\definecolor{currentfill}{rgb}{0.990773,0.990773,0.990773}%
\pgfsetfillcolor{currentfill}%
\pgfsetfillopacity{0.900000}%
\pgfsetlinewidth{0.501875pt}%
\definecolor{currentstroke}{rgb}{0.200000,0.200000,0.200000}%
\pgfsetstrokecolor{currentstroke}%
\pgfsetdash{}{0pt}%
\pgfpathmoveto{\pgfqpoint{3.281035in}{1.728870in}}%
\pgfpathlineto{\pgfqpoint{3.324859in}{1.730035in}}%
\pgfpathlineto{\pgfqpoint{3.357427in}{1.711109in}}%
\pgfpathlineto{\pgfqpoint{3.313580in}{1.710055in}}%
\pgfpathlineto{\pgfqpoint{3.281035in}{1.728870in}}%
\pgfpathclose%
\pgfusepath{stroke,fill}%
\end{pgfscope}%
\begin{pgfscope}%
\pgfpathrectangle{\pgfqpoint{0.050000in}{0.050000in}}{\pgfqpoint{4.900000in}{4.900000in}}%
\pgfusepath{clip}%
\pgfsetbuttcap%
\pgfsetroundjoin%
\definecolor{currentfill}{rgb}{0.878570,0.878570,0.878570}%
\pgfsetfillcolor{currentfill}%
\pgfsetfillopacity{0.900000}%
\pgfsetlinewidth{0.501875pt}%
\definecolor{currentstroke}{rgb}{0.200000,0.200000,0.200000}%
\pgfsetstrokecolor{currentstroke}%
\pgfsetdash{}{0pt}%
\pgfpathmoveto{\pgfqpoint{2.411468in}{2.206460in}}%
\pgfpathlineto{\pgfqpoint{2.456524in}{2.165355in}}%
\pgfpathlineto{\pgfqpoint{2.487732in}{2.145655in}}%
\pgfpathlineto{\pgfqpoint{2.442671in}{2.186545in}}%
\pgfpathlineto{\pgfqpoint{2.411468in}{2.206460in}}%
\pgfpathclose%
\pgfusepath{stroke,fill}%
\end{pgfscope}%
\begin{pgfscope}%
\pgfpathrectangle{\pgfqpoint{0.050000in}{0.050000in}}{\pgfqpoint{4.900000in}{4.900000in}}%
\pgfusepath{clip}%
\pgfsetbuttcap%
\pgfsetroundjoin%
\definecolor{currentfill}{rgb}{0.547328,0.547328,0.547328}%
\pgfsetfillcolor{currentfill}%
\pgfsetfillopacity{0.900000}%
\pgfsetlinewidth{0.501875pt}%
\definecolor{currentstroke}{rgb}{0.200000,0.200000,0.200000}%
\pgfsetstrokecolor{currentstroke}%
\pgfsetdash{}{0pt}%
\pgfpathmoveto{\pgfqpoint{1.447380in}{2.984190in}}%
\pgfpathlineto{\pgfqpoint{1.494446in}{2.939406in}}%
\pgfpathlineto{\pgfqpoint{1.524193in}{2.920263in}}%
\pgfpathlineto{\pgfqpoint{1.477117in}{2.965075in}}%
\pgfpathlineto{\pgfqpoint{1.447380in}{2.984190in}}%
\pgfpathclose%
\pgfusepath{stroke,fill}%
\end{pgfscope}%
\begin{pgfscope}%
\pgfpathrectangle{\pgfqpoint{0.050000in}{0.050000in}}{\pgfqpoint{4.900000in}{4.900000in}}%
\pgfusepath{clip}%
\pgfsetbuttcap%
\pgfsetroundjoin%
\definecolor{currentfill}{rgb}{0.808781,0.808781,0.808781}%
\pgfsetfillcolor{currentfill}%
\pgfsetfillopacity{0.900000}%
\pgfsetlinewidth{0.501875pt}%
\definecolor{currentstroke}{rgb}{0.200000,0.200000,0.200000}%
\pgfsetstrokecolor{currentstroke}%
\pgfsetdash{}{0pt}%
\pgfpathmoveto{\pgfqpoint{2.151355in}{2.415245in}}%
\pgfpathlineto{\pgfqpoint{2.196960in}{2.371729in}}%
\pgfpathlineto{\pgfqpoint{2.227777in}{2.351778in}}%
\pgfpathlineto{\pgfqpoint{2.182166in}{2.395286in}}%
\pgfpathlineto{\pgfqpoint{2.151355in}{2.415245in}}%
\pgfpathclose%
\pgfusepath{stroke,fill}%
\end{pgfscope}%
\begin{pgfscope}%
\pgfpathrectangle{\pgfqpoint{0.050000in}{0.050000in}}{\pgfqpoint{4.900000in}{4.900000in}}%
\pgfusepath{clip}%
\pgfsetbuttcap%
\pgfsetroundjoin%
\definecolor{currentfill}{rgb}{0.998155,0.998155,0.998155}%
\pgfsetfillcolor{currentfill}%
\pgfsetfillopacity{0.900000}%
\pgfsetlinewidth{0.501875pt}%
\definecolor{currentstroke}{rgb}{0.200000,0.200000,0.200000}%
\pgfsetstrokecolor{currentstroke}%
\pgfsetdash{}{0pt}%
\pgfpathmoveto{\pgfqpoint{3.630702in}{1.680069in}}%
\pgfpathlineto{\pgfqpoint{3.674255in}{1.690529in}}%
\pgfpathlineto{\pgfqpoint{3.707360in}{1.670763in}}%
\pgfpathlineto{\pgfqpoint{3.663779in}{1.660400in}}%
\pgfpathlineto{\pgfqpoint{3.630702in}{1.680069in}}%
\pgfpathclose%
\pgfusepath{stroke,fill}%
\end{pgfscope}%
\begin{pgfscope}%
\pgfpathrectangle{\pgfqpoint{0.050000in}{0.050000in}}{\pgfqpoint{4.900000in}{4.900000in}}%
\pgfusepath{clip}%
\pgfsetbuttcap%
\pgfsetroundjoin%
\definecolor{currentfill}{rgb}{0.935163,0.935163,0.935163}%
\pgfsetfillcolor{currentfill}%
\pgfsetfillopacity{0.900000}%
\pgfsetlinewidth{0.501875pt}%
\definecolor{currentstroke}{rgb}{0.200000,0.200000,0.200000}%
\pgfsetstrokecolor{currentstroke}%
\pgfsetdash{}{0pt}%
\pgfpathmoveto{\pgfqpoint{2.671542in}{2.012364in}}%
\pgfpathlineto{\pgfqpoint{2.716095in}{1.979859in}}%
\pgfpathlineto{\pgfqpoint{2.747712in}{1.960979in}}%
\pgfpathlineto{\pgfqpoint{2.703150in}{1.993144in}}%
\pgfpathlineto{\pgfqpoint{2.671542in}{2.012364in}}%
\pgfpathclose%
\pgfusepath{stroke,fill}%
\end{pgfscope}%
\begin{pgfscope}%
\pgfpathrectangle{\pgfqpoint{0.050000in}{0.050000in}}{\pgfqpoint{4.900000in}{4.900000in}}%
\pgfusepath{clip}%
\pgfsetbuttcap%
\pgfsetroundjoin%
\definecolor{currentfill}{rgb}{0.992618,0.992618,0.992618}%
\pgfsetfillcolor{currentfill}%
\pgfsetfillopacity{0.900000}%
\pgfsetlinewidth{0.501875pt}%
\definecolor{currentstroke}{rgb}{0.200000,0.200000,0.200000}%
\pgfsetstrokecolor{currentstroke}%
\pgfsetdash{}{0pt}%
\pgfpathmoveto{\pgfqpoint{3.357427in}{1.711109in}}%
\pgfpathlineto{\pgfqpoint{3.401201in}{1.714932in}}%
\pgfpathlineto{\pgfqpoint{3.433896in}{1.695809in}}%
\pgfpathlineto{\pgfqpoint{3.390097in}{1.692110in}}%
\pgfpathlineto{\pgfqpoint{3.357427in}{1.711109in}}%
\pgfpathclose%
\pgfusepath{stroke,fill}%
\end{pgfscope}%
\begin{pgfscope}%
\pgfpathrectangle{\pgfqpoint{0.050000in}{0.050000in}}{\pgfqpoint{4.900000in}{4.900000in}}%
\pgfusepath{clip}%
\pgfsetbuttcap%
\pgfsetroundjoin%
\definecolor{currentfill}{rgb}{0.724983,0.724983,0.724983}%
\pgfsetfillcolor{currentfill}%
\pgfsetfillopacity{0.900000}%
\pgfsetlinewidth{0.501875pt}%
\definecolor{currentstroke}{rgb}{0.200000,0.200000,0.200000}%
\pgfsetstrokecolor{currentstroke}%
\pgfsetdash{}{0pt}%
\pgfpathmoveto{\pgfqpoint{1.891150in}{2.626125in}}%
\pgfpathlineto{\pgfqpoint{1.937304in}{2.582075in}}%
\pgfpathlineto{\pgfqpoint{1.967731in}{2.562386in}}%
\pgfpathlineto{\pgfqpoint{1.921570in}{2.606462in}}%
\pgfpathlineto{\pgfqpoint{1.891150in}{2.626125in}}%
\pgfpathclose%
\pgfusepath{stroke,fill}%
\end{pgfscope}%
\begin{pgfscope}%
\pgfpathrectangle{\pgfqpoint{0.050000in}{0.050000in}}{\pgfqpoint{4.900000in}{4.900000in}}%
\pgfusepath{clip}%
\pgfsetbuttcap%
\pgfsetroundjoin%
\definecolor{currentfill}{rgb}{1.000000,1.000000,1.000000}%
\pgfsetfillcolor{currentfill}%
\pgfsetfillopacity{0.900000}%
\pgfsetlinewidth{0.501875pt}%
\definecolor{currentstroke}{rgb}{0.200000,0.200000,0.200000}%
\pgfsetstrokecolor{currentstroke}%
\pgfsetdash{}{0pt}%
\pgfpathmoveto{\pgfqpoint{3.904295in}{1.666408in}}%
\pgfpathlineto{\pgfqpoint{3.947597in}{1.679082in}}%
\pgfpathlineto{\pgfqpoint{3.981118in}{1.658868in}}%
\pgfpathlineto{\pgfqpoint{3.937789in}{1.646250in}}%
\pgfpathlineto{\pgfqpoint{3.904295in}{1.666408in}}%
\pgfpathclose%
\pgfusepath{stroke,fill}%
\end{pgfscope}%
\begin{pgfscope}%
\pgfpathrectangle{\pgfqpoint{0.050000in}{0.050000in}}{\pgfqpoint{4.900000in}{4.900000in}}%
\pgfusepath{clip}%
\pgfsetbuttcap%
\pgfsetroundjoin%
\definecolor{currentfill}{rgb}{0.614625,0.614625,0.614625}%
\pgfsetfillcolor{currentfill}%
\pgfsetfillopacity{0.900000}%
\pgfsetlinewidth{0.501875pt}%
\definecolor{currentstroke}{rgb}{0.200000,0.200000,0.200000}%
\pgfsetstrokecolor{currentstroke}%
\pgfsetdash{}{0pt}%
\pgfpathmoveto{\pgfqpoint{1.630858in}{2.837121in}}%
\pgfpathlineto{\pgfqpoint{1.677562in}{2.792625in}}%
\pgfpathlineto{\pgfqpoint{1.707599in}{2.773252in}}%
\pgfpathlineto{\pgfqpoint{1.660887in}{2.817775in}}%
\pgfpathlineto{\pgfqpoint{1.630858in}{2.837121in}}%
\pgfpathclose%
\pgfusepath{stroke,fill}%
\end{pgfscope}%
\begin{pgfscope}%
\pgfpathrectangle{\pgfqpoint{0.050000in}{0.050000in}}{\pgfqpoint{4.900000in}{4.900000in}}%
\pgfusepath{clip}%
\pgfsetbuttcap%
\pgfsetroundjoin%
\definecolor{currentfill}{rgb}{0.858762,0.858762,0.858762}%
\pgfsetfillcolor{currentfill}%
\pgfsetfillopacity{0.900000}%
\pgfsetlinewidth{0.501875pt}%
\definecolor{currentstroke}{rgb}{0.200000,0.200000,0.200000}%
\pgfsetstrokecolor{currentstroke}%
\pgfsetdash{}{0pt}%
\pgfpathmoveto{\pgfqpoint{2.335121in}{2.268705in}}%
\pgfpathlineto{\pgfqpoint{2.380361in}{2.226345in}}%
\pgfpathlineto{\pgfqpoint{2.411468in}{2.206460in}}%
\pgfpathlineto{\pgfqpoint{2.366224in}{2.248686in}}%
\pgfpathlineto{\pgfqpoint{2.335121in}{2.268705in}}%
\pgfpathclose%
\pgfusepath{stroke,fill}%
\end{pgfscope}%
\begin{pgfscope}%
\pgfpathrectangle{\pgfqpoint{0.050000in}{0.050000in}}{\pgfqpoint{4.900000in}{4.900000in}}%
\pgfusepath{clip}%
\pgfsetbuttcap%
\pgfsetroundjoin%
\definecolor{currentfill}{rgb}{0.998155,0.998155,0.998155}%
\pgfsetfillcolor{currentfill}%
\pgfsetfillopacity{0.900000}%
\pgfsetlinewidth{0.501875pt}%
\definecolor{currentstroke}{rgb}{0.200000,0.200000,0.200000}%
\pgfsetstrokecolor{currentstroke}%
\pgfsetdash{}{0pt}%
\pgfpathmoveto{\pgfqpoint{3.707360in}{1.670763in}}%
\pgfpathlineto{\pgfqpoint{3.750861in}{1.682067in}}%
\pgfpathlineto{\pgfqpoint{3.784095in}{1.662160in}}%
\pgfpathlineto{\pgfqpoint{3.740567in}{1.650940in}}%
\pgfpathlineto{\pgfqpoint{3.707360in}{1.670763in}}%
\pgfpathclose%
\pgfusepath{stroke,fill}%
\end{pgfscope}%
\begin{pgfscope}%
\pgfpathrectangle{\pgfqpoint{0.050000in}{0.050000in}}{\pgfqpoint{4.900000in}{4.900000in}}%
\pgfusepath{clip}%
\pgfsetbuttcap%
\pgfsetroundjoin%
\definecolor{currentfill}{rgb}{0.994464,0.994464,0.994464}%
\pgfsetfillcolor{currentfill}%
\pgfsetfillopacity{0.900000}%
\pgfsetlinewidth{0.501875pt}%
\definecolor{currentstroke}{rgb}{0.200000,0.200000,0.200000}%
\pgfsetstrokecolor{currentstroke}%
\pgfsetdash{}{0pt}%
\pgfpathmoveto{\pgfqpoint{3.433896in}{1.695809in}}%
\pgfpathlineto{\pgfqpoint{3.477622in}{1.701832in}}%
\pgfpathlineto{\pgfqpoint{3.510443in}{1.682515in}}%
\pgfpathlineto{\pgfqpoint{3.466692in}{1.676618in}}%
\pgfpathlineto{\pgfqpoint{3.433896in}{1.695809in}}%
\pgfpathclose%
\pgfusepath{stroke,fill}%
\end{pgfscope}%
\begin{pgfscope}%
\pgfpathrectangle{\pgfqpoint{0.050000in}{0.050000in}}{\pgfqpoint{4.900000in}{4.900000in}}%
\pgfusepath{clip}%
\pgfsetbuttcap%
\pgfsetroundjoin%
\definecolor{currentfill}{rgb}{0.918185,0.918185,0.918185}%
\pgfsetfillcolor{currentfill}%
\pgfsetfillopacity{0.900000}%
\pgfsetlinewidth{0.501875pt}%
\definecolor{currentstroke}{rgb}{0.200000,0.200000,0.200000}%
\pgfsetstrokecolor{currentstroke}%
\pgfsetdash{}{0pt}%
\pgfpathmoveto{\pgfqpoint{2.595322in}{2.067597in}}%
\pgfpathlineto{\pgfqpoint{2.640033in}{2.031536in}}%
\pgfpathlineto{\pgfqpoint{2.671542in}{2.012364in}}%
\pgfpathlineto{\pgfqpoint{2.626823in}{2.048092in}}%
\pgfpathlineto{\pgfqpoint{2.595322in}{2.067597in}}%
\pgfpathclose%
\pgfusepath{stroke,fill}%
\end{pgfscope}%
\begin{pgfscope}%
\pgfpathrectangle{\pgfqpoint{0.050000in}{0.050000in}}{\pgfqpoint{4.900000in}{4.900000in}}%
\pgfusepath{clip}%
\pgfsetbuttcap%
\pgfsetroundjoin%
\definecolor{currentfill}{rgb}{0.517186,0.517186,0.517186}%
\pgfsetfillcolor{currentfill}%
\pgfsetfillopacity{0.900000}%
\pgfsetlinewidth{0.501875pt}%
\definecolor{currentstroke}{rgb}{0.200000,0.200000,0.200000}%
\pgfsetstrokecolor{currentstroke}%
\pgfsetdash{}{0pt}%
\pgfpathmoveto{\pgfqpoint{1.370480in}{3.048189in}}%
\pgfpathlineto{\pgfqpoint{1.417733in}{3.003247in}}%
\pgfpathlineto{\pgfqpoint{1.447380in}{2.984190in}}%
\pgfpathlineto{\pgfqpoint{1.400117in}{3.029161in}}%
\pgfpathlineto{\pgfqpoint{1.370480in}{3.048189in}}%
\pgfpathclose%
\pgfusepath{stroke,fill}%
\end{pgfscope}%
\begin{pgfscope}%
\pgfpathrectangle{\pgfqpoint{0.050000in}{0.050000in}}{\pgfqpoint{4.900000in}{4.900000in}}%
\pgfusepath{clip}%
\pgfsetbuttcap%
\pgfsetroundjoin%
\definecolor{currentfill}{rgb}{0.781223,0.781223,0.781223}%
\pgfsetfillcolor{currentfill}%
\pgfsetfillopacity{0.900000}%
\pgfsetlinewidth{0.501875pt}%
\definecolor{currentstroke}{rgb}{0.200000,0.200000,0.200000}%
\pgfsetstrokecolor{currentstroke}%
\pgfsetdash{}{0pt}%
\pgfpathmoveto{\pgfqpoint{2.074846in}{2.478890in}}%
\pgfpathlineto{\pgfqpoint{2.120637in}{2.435149in}}%
\pgfpathlineto{\pgfqpoint{2.151355in}{2.415245in}}%
\pgfpathlineto{\pgfqpoint{2.105558in}{2.459001in}}%
\pgfpathlineto{\pgfqpoint{2.074846in}{2.478890in}}%
\pgfpathclose%
\pgfusepath{stroke,fill}%
\end{pgfscope}%
\begin{pgfscope}%
\pgfpathrectangle{\pgfqpoint{0.050000in}{0.050000in}}{\pgfqpoint{4.900000in}{4.900000in}}%
\pgfusepath{clip}%
\pgfsetbuttcap%
\pgfsetroundjoin%
\definecolor{currentfill}{rgb}{1.000000,1.000000,1.000000}%
\pgfsetfillcolor{currentfill}%
\pgfsetfillopacity{0.900000}%
\pgfsetlinewidth{0.501875pt}%
\definecolor{currentstroke}{rgb}{0.200000,0.200000,0.200000}%
\pgfsetstrokecolor{currentstroke}%
\pgfsetdash{}{0pt}%
\pgfpathmoveto{\pgfqpoint{3.981118in}{1.658868in}}%
\pgfpathlineto{\pgfqpoint{4.024356in}{1.671648in}}%
\pgfpathlineto{\pgfqpoint{4.058008in}{1.651321in}}%
\pgfpathlineto{\pgfqpoint{4.014743in}{1.638594in}}%
\pgfpathlineto{\pgfqpoint{3.981118in}{1.658868in}}%
\pgfpathclose%
\pgfusepath{stroke,fill}%
\end{pgfscope}%
\begin{pgfscope}%
\pgfpathrectangle{\pgfqpoint{0.050000in}{0.050000in}}{\pgfqpoint{4.900000in}{4.900000in}}%
\pgfusepath{clip}%
\pgfsetbuttcap%
\pgfsetroundjoin%
\definecolor{currentfill}{rgb}{0.691396,0.691396,0.691396}%
\pgfsetfillcolor{currentfill}%
\pgfsetfillopacity{0.900000}%
\pgfsetlinewidth{0.501875pt}%
\definecolor{currentstroke}{rgb}{0.200000,0.200000,0.200000}%
\pgfsetstrokecolor{currentstroke}%
\pgfsetdash{}{0pt}%
\pgfpathmoveto{\pgfqpoint{1.814482in}{2.689936in}}%
\pgfpathlineto{\pgfqpoint{1.860822in}{2.645729in}}%
\pgfpathlineto{\pgfqpoint{1.891150in}{2.626125in}}%
\pgfpathlineto{\pgfqpoint{1.844802in}{2.670359in}}%
\pgfpathlineto{\pgfqpoint{1.814482in}{2.689936in}}%
\pgfpathclose%
\pgfusepath{stroke,fill}%
\end{pgfscope}%
\begin{pgfscope}%
\pgfpathrectangle{\pgfqpoint{0.050000in}{0.050000in}}{\pgfqpoint{4.900000in}{4.900000in}}%
\pgfusepath{clip}%
\pgfsetbuttcap%
\pgfsetroundjoin%
\definecolor{currentfill}{rgb}{0.996309,0.996309,0.996309}%
\pgfsetfillcolor{currentfill}%
\pgfsetfillopacity{0.900000}%
\pgfsetlinewidth{0.501875pt}%
\definecolor{currentstroke}{rgb}{0.200000,0.200000,0.200000}%
\pgfsetstrokecolor{currentstroke}%
\pgfsetdash{}{0pt}%
\pgfpathmoveto{\pgfqpoint{3.510443in}{1.682515in}}%
\pgfpathlineto{\pgfqpoint{3.554122in}{1.690327in}}%
\pgfpathlineto{\pgfqpoint{3.587071in}{1.670827in}}%
\pgfpathlineto{\pgfqpoint{3.543367in}{1.663135in}}%
\pgfpathlineto{\pgfqpoint{3.510443in}{1.682515in}}%
\pgfpathclose%
\pgfusepath{stroke,fill}%
\end{pgfscope}%
\begin{pgfscope}%
\pgfpathrectangle{\pgfqpoint{0.050000in}{0.050000in}}{\pgfqpoint{4.900000in}{4.900000in}}%
\pgfusepath{clip}%
\pgfsetbuttcap%
\pgfsetroundjoin%
\definecolor{currentfill}{rgb}{1.000000,1.000000,1.000000}%
\pgfsetfillcolor{currentfill}%
\pgfsetfillopacity{0.900000}%
\pgfsetlinewidth{0.501875pt}%
\definecolor{currentstroke}{rgb}{0.200000,0.200000,0.200000}%
\pgfsetstrokecolor{currentstroke}%
\pgfsetdash{}{0pt}%
\pgfpathmoveto{\pgfqpoint{3.784095in}{1.662160in}}%
\pgfpathlineto{\pgfqpoint{3.827541in}{1.674087in}}%
\pgfpathlineto{\pgfqpoint{3.860905in}{1.654051in}}%
\pgfpathlineto{\pgfqpoint{3.817432in}{1.642197in}}%
\pgfpathlineto{\pgfqpoint{3.784095in}{1.662160in}}%
\pgfpathclose%
\pgfusepath{stroke,fill}%
\end{pgfscope}%
\begin{pgfscope}%
\pgfpathrectangle{\pgfqpoint{0.050000in}{0.050000in}}{\pgfqpoint{4.900000in}{4.900000in}}%
\pgfusepath{clip}%
\pgfsetbuttcap%
\pgfsetroundjoin%
\definecolor{currentfill}{rgb}{0.581776,0.581776,0.581776}%
\pgfsetfillcolor{currentfill}%
\pgfsetfillopacity{0.900000}%
\pgfsetlinewidth{0.501875pt}%
\definecolor{currentstroke}{rgb}{0.200000,0.200000,0.200000}%
\pgfsetstrokecolor{currentstroke}%
\pgfsetdash{}{0pt}%
\pgfpathmoveto{\pgfqpoint{1.554031in}{2.901062in}}%
\pgfpathlineto{\pgfqpoint{1.600921in}{2.856408in}}%
\pgfpathlineto{\pgfqpoint{1.630858in}{2.837121in}}%
\pgfpathlineto{\pgfqpoint{1.583959in}{2.881803in}}%
\pgfpathlineto{\pgfqpoint{1.554031in}{2.901062in}}%
\pgfpathclose%
\pgfusepath{stroke,fill}%
\end{pgfscope}%
\begin{pgfscope}%
\pgfpathrectangle{\pgfqpoint{0.050000in}{0.050000in}}{\pgfqpoint{4.900000in}{4.900000in}}%
\pgfusepath{clip}%
\pgfsetbuttcap%
\pgfsetroundjoin%
\definecolor{currentfill}{rgb}{0.898378,0.898378,0.898378}%
\pgfsetfillcolor{currentfill}%
\pgfsetfillopacity{0.900000}%
\pgfsetlinewidth{0.501875pt}%
\definecolor{currentstroke}{rgb}{0.200000,0.200000,0.200000}%
\pgfsetstrokecolor{currentstroke}%
\pgfsetdash{}{0pt}%
\pgfpathmoveto{\pgfqpoint{2.519035in}{2.125928in}}%
\pgfpathlineto{\pgfqpoint{2.563917in}{2.087065in}}%
\pgfpathlineto{\pgfqpoint{2.595322in}{2.067597in}}%
\pgfpathlineto{\pgfqpoint{2.550434in}{2.106171in}}%
\pgfpathlineto{\pgfqpoint{2.519035in}{2.125928in}}%
\pgfpathclose%
\pgfusepath{stroke,fill}%
\end{pgfscope}%
\begin{pgfscope}%
\pgfpathrectangle{\pgfqpoint{0.050000in}{0.050000in}}{\pgfqpoint{4.900000in}{4.900000in}}%
\pgfusepath{clip}%
\pgfsetbuttcap%
\pgfsetroundjoin%
\definecolor{currentfill}{rgb}{0.972318,0.972318,0.972318}%
\pgfsetfillcolor{currentfill}%
\pgfsetfillopacity{0.900000}%
\pgfsetlinewidth{0.501875pt}%
\definecolor{currentstroke}{rgb}{0.200000,0.200000,0.200000}%
\pgfsetstrokecolor{currentstroke}%
\pgfsetdash{}{0pt}%
\pgfpathmoveto{\pgfqpoint{3.008162in}{1.816070in}}%
\pgfpathlineto{\pgfqpoint{3.052262in}{1.801947in}}%
\pgfpathlineto{\pgfqpoint{3.084438in}{1.783490in}}%
\pgfpathlineto{\pgfqpoint{3.040320in}{1.797520in}}%
\pgfpathlineto{\pgfqpoint{3.008162in}{1.816070in}}%
\pgfpathclose%
\pgfusepath{stroke,fill}%
\end{pgfscope}%
\begin{pgfscope}%
\pgfpathrectangle{\pgfqpoint{0.050000in}{0.050000in}}{\pgfqpoint{4.900000in}{4.900000in}}%
\pgfusepath{clip}%
\pgfsetbuttcap%
\pgfsetroundjoin%
\definecolor{currentfill}{rgb}{0.977855,0.977855,0.977855}%
\pgfsetfillcolor{currentfill}%
\pgfsetfillopacity{0.900000}%
\pgfsetlinewidth{0.501875pt}%
\definecolor{currentstroke}{rgb}{0.200000,0.200000,0.200000}%
\pgfsetstrokecolor{currentstroke}%
\pgfsetdash{}{0pt}%
\pgfpathmoveto{\pgfqpoint{3.084438in}{1.783490in}}%
\pgfpathlineto{\pgfqpoint{3.128461in}{1.773780in}}%
\pgfpathlineto{\pgfqpoint{3.160759in}{1.755254in}}%
\pgfpathlineto{\pgfqpoint{3.116715in}{1.764941in}}%
\pgfpathlineto{\pgfqpoint{3.084438in}{1.783490in}}%
\pgfpathclose%
\pgfusepath{stroke,fill}%
\end{pgfscope}%
\begin{pgfscope}%
\pgfpathrectangle{\pgfqpoint{0.050000in}{0.050000in}}{\pgfqpoint{4.900000in}{4.900000in}}%
\pgfusepath{clip}%
\pgfsetbuttcap%
\pgfsetroundjoin%
\definecolor{currentfill}{rgb}{0.832895,0.832895,0.832895}%
\pgfsetfillcolor{currentfill}%
\pgfsetfillopacity{0.900000}%
\pgfsetlinewidth{0.501875pt}%
\definecolor{currentstroke}{rgb}{0.200000,0.200000,0.200000}%
\pgfsetstrokecolor{currentstroke}%
\pgfsetdash{}{0pt}%
\pgfpathmoveto{\pgfqpoint{2.258687in}{2.331781in}}%
\pgfpathlineto{\pgfqpoint{2.304112in}{2.288686in}}%
\pgfpathlineto{\pgfqpoint{2.335121in}{2.268705in}}%
\pgfpathlineto{\pgfqpoint{2.289691in}{2.311737in}}%
\pgfpathlineto{\pgfqpoint{2.258687in}{2.331781in}}%
\pgfpathclose%
\pgfusepath{stroke,fill}%
\end{pgfscope}%
\begin{pgfscope}%
\pgfpathrectangle{\pgfqpoint{0.050000in}{0.050000in}}{\pgfqpoint{4.900000in}{4.900000in}}%
\pgfusepath{clip}%
\pgfsetbuttcap%
\pgfsetroundjoin%
\definecolor{currentfill}{rgb}{0.964937,0.964937,0.964937}%
\pgfsetfillcolor{currentfill}%
\pgfsetfillopacity{0.900000}%
\pgfsetlinewidth{0.501875pt}%
\definecolor{currentstroke}{rgb}{0.200000,0.200000,0.200000}%
\pgfsetstrokecolor{currentstroke}%
\pgfsetdash{}{0pt}%
\pgfpathmoveto{\pgfqpoint{2.931914in}{1.853300in}}%
\pgfpathlineto{\pgfqpoint{2.976104in}{1.834528in}}%
\pgfpathlineto{\pgfqpoint{3.008162in}{1.816070in}}%
\pgfpathlineto{\pgfqpoint{2.963954in}{1.834672in}}%
\pgfpathlineto{\pgfqpoint{2.931914in}{1.853300in}}%
\pgfpathclose%
\pgfusepath{stroke,fill}%
\end{pgfscope}%
\begin{pgfscope}%
\pgfpathrectangle{\pgfqpoint{0.050000in}{0.050000in}}{\pgfqpoint{4.900000in}{4.900000in}}%
\pgfusepath{clip}%
\pgfsetbuttcap%
\pgfsetroundjoin%
\definecolor{currentfill}{rgb}{0.981546,0.981546,0.981546}%
\pgfsetfillcolor{currentfill}%
\pgfsetfillopacity{0.900000}%
\pgfsetlinewidth{0.501875pt}%
\definecolor{currentstroke}{rgb}{0.200000,0.200000,0.200000}%
\pgfsetstrokecolor{currentstroke}%
\pgfsetdash{}{0pt}%
\pgfpathmoveto{\pgfqpoint{3.160759in}{1.755254in}}%
\pgfpathlineto{\pgfqpoint{3.204716in}{1.749588in}}%
\pgfpathlineto{\pgfqpoint{3.237136in}{1.730937in}}%
\pgfpathlineto{\pgfqpoint{3.193157in}{1.736636in}}%
\pgfpathlineto{\pgfqpoint{3.160759in}{1.755254in}}%
\pgfpathclose%
\pgfusepath{stroke,fill}%
\end{pgfscope}%
\begin{pgfscope}%
\pgfpathrectangle{\pgfqpoint{0.050000in}{0.050000in}}{\pgfqpoint{4.900000in}{4.900000in}}%
\pgfusepath{clip}%
\pgfsetbuttcap%
\pgfsetroundjoin%
\definecolor{currentfill}{rgb}{0.757109,0.757109,0.757109}%
\pgfsetfillcolor{currentfill}%
\pgfsetfillopacity{0.900000}%
\pgfsetlinewidth{0.501875pt}%
\definecolor{currentstroke}{rgb}{0.200000,0.200000,0.200000}%
\pgfsetstrokecolor{currentstroke}%
\pgfsetdash{}{0pt}%
\pgfpathmoveto{\pgfqpoint{1.998251in}{2.542637in}}%
\pgfpathlineto{\pgfqpoint{2.044227in}{2.498722in}}%
\pgfpathlineto{\pgfqpoint{2.074846in}{2.478890in}}%
\pgfpathlineto{\pgfqpoint{2.028864in}{2.522829in}}%
\pgfpathlineto{\pgfqpoint{1.998251in}{2.542637in}}%
\pgfpathclose%
\pgfusepath{stroke,fill}%
\end{pgfscope}%
\begin{pgfscope}%
\pgfpathrectangle{\pgfqpoint{0.050000in}{0.050000in}}{\pgfqpoint{4.900000in}{4.900000in}}%
\pgfusepath{clip}%
\pgfsetbuttcap%
\pgfsetroundjoin%
\definecolor{currentfill}{rgb}{0.955709,0.955709,0.955709}%
\pgfsetfillcolor{currentfill}%
\pgfsetfillopacity{0.900000}%
\pgfsetlinewidth{0.501875pt}%
\definecolor{currentstroke}{rgb}{0.200000,0.200000,0.200000}%
\pgfsetstrokecolor{currentstroke}%
\pgfsetdash{}{0pt}%
\pgfpathmoveto{\pgfqpoint{2.855676in}{1.895316in}}%
\pgfpathlineto{\pgfqpoint{2.899973in}{1.871841in}}%
\pgfpathlineto{\pgfqpoint{2.931914in}{1.853300in}}%
\pgfpathlineto{\pgfqpoint{2.887601in}{1.876533in}}%
\pgfpathlineto{\pgfqpoint{2.855676in}{1.895316in}}%
\pgfpathclose%
\pgfusepath{stroke,fill}%
\end{pgfscope}%
\begin{pgfscope}%
\pgfpathrectangle{\pgfqpoint{0.050000in}{0.050000in}}{\pgfqpoint{4.900000in}{4.900000in}}%
\pgfusepath{clip}%
\pgfsetbuttcap%
\pgfsetroundjoin%
\definecolor{currentfill}{rgb}{0.996309,0.996309,0.996309}%
\pgfsetfillcolor{currentfill}%
\pgfsetfillopacity{0.900000}%
\pgfsetlinewidth{0.501875pt}%
\definecolor{currentstroke}{rgb}{0.200000,0.200000,0.200000}%
\pgfsetstrokecolor{currentstroke}%
\pgfsetdash{}{0pt}%
\pgfpathmoveto{\pgfqpoint{3.587071in}{1.670827in}}%
\pgfpathlineto{\pgfqpoint{3.630702in}{1.680069in}}%
\pgfpathlineto{\pgfqpoint{3.663779in}{1.660400in}}%
\pgfpathlineto{\pgfqpoint{3.620123in}{1.651268in}}%
\pgfpathlineto{\pgfqpoint{3.587071in}{1.670827in}}%
\pgfpathclose%
\pgfusepath{stroke,fill}%
\end{pgfscope}%
\begin{pgfscope}%
\pgfpathrectangle{\pgfqpoint{0.050000in}{0.050000in}}{\pgfqpoint{4.900000in}{4.900000in}}%
\pgfusepath{clip}%
\pgfsetbuttcap%
\pgfsetroundjoin%
\definecolor{currentfill}{rgb}{0.987082,0.987082,0.987082}%
\pgfsetfillcolor{currentfill}%
\pgfsetfillopacity{0.900000}%
\pgfsetlinewidth{0.501875pt}%
\definecolor{currentstroke}{rgb}{0.200000,0.200000,0.200000}%
\pgfsetstrokecolor{currentstroke}%
\pgfsetdash{}{0pt}%
\pgfpathmoveto{\pgfqpoint{3.237136in}{1.730937in}}%
\pgfpathlineto{\pgfqpoint{3.281035in}{1.728870in}}%
\pgfpathlineto{\pgfqpoint{3.313580in}{1.710055in}}%
\pgfpathlineto{\pgfqpoint{3.269658in}{1.712199in}}%
\pgfpathlineto{\pgfqpoint{3.237136in}{1.730937in}}%
\pgfpathclose%
\pgfusepath{stroke,fill}%
\end{pgfscope}%
\begin{pgfscope}%
\pgfpathrectangle{\pgfqpoint{0.050000in}{0.050000in}}{\pgfqpoint{4.900000in}{4.900000in}}%
\pgfusepath{clip}%
\pgfsetbuttcap%
\pgfsetroundjoin%
\definecolor{currentfill}{rgb}{0.653010,0.653010,0.653010}%
\pgfsetfillcolor{currentfill}%
\pgfsetfillopacity{0.900000}%
\pgfsetlinewidth{0.501875pt}%
\definecolor{currentstroke}{rgb}{0.200000,0.200000,0.200000}%
\pgfsetstrokecolor{currentstroke}%
\pgfsetdash{}{0pt}%
\pgfpathmoveto{\pgfqpoint{1.737727in}{2.753820in}}%
\pgfpathlineto{\pgfqpoint{1.784253in}{2.709455in}}%
\pgfpathlineto{\pgfqpoint{1.814482in}{2.689936in}}%
\pgfpathlineto{\pgfqpoint{1.767948in}{2.734329in}}%
\pgfpathlineto{\pgfqpoint{1.737727in}{2.753820in}}%
\pgfpathclose%
\pgfusepath{stroke,fill}%
\end{pgfscope}%
\begin{pgfscope}%
\pgfpathrectangle{\pgfqpoint{0.050000in}{0.050000in}}{\pgfqpoint{4.900000in}{4.900000in}}%
\pgfusepath{clip}%
\pgfsetbuttcap%
\pgfsetroundjoin%
\definecolor{currentfill}{rgb}{1.000000,1.000000,1.000000}%
\pgfsetfillcolor{currentfill}%
\pgfsetfillopacity{0.900000}%
\pgfsetlinewidth{0.501875pt}%
\definecolor{currentstroke}{rgb}{0.200000,0.200000,0.200000}%
\pgfsetstrokecolor{currentstroke}%
\pgfsetdash{}{0pt}%
\pgfpathmoveto{\pgfqpoint{3.860905in}{1.654051in}}%
\pgfpathlineto{\pgfqpoint{3.904295in}{1.666408in}}%
\pgfpathlineto{\pgfqpoint{3.937789in}{1.646250in}}%
\pgfpathlineto{\pgfqpoint{3.894372in}{1.633957in}}%
\pgfpathlineto{\pgfqpoint{3.860905in}{1.654051in}}%
\pgfpathclose%
\pgfusepath{stroke,fill}%
\end{pgfscope}%
\begin{pgfscope}%
\pgfpathrectangle{\pgfqpoint{0.050000in}{0.050000in}}{\pgfqpoint{4.900000in}{4.900000in}}%
\pgfusepath{clip}%
\pgfsetbuttcap%
\pgfsetroundjoin%
\definecolor{currentfill}{rgb}{0.946482,0.946482,0.946482}%
\pgfsetfillcolor{currentfill}%
\pgfsetfillopacity{0.900000}%
\pgfsetlinewidth{0.501875pt}%
\definecolor{currentstroke}{rgb}{0.200000,0.200000,0.200000}%
\pgfsetstrokecolor{currentstroke}%
\pgfsetdash{}{0pt}%
\pgfpathmoveto{\pgfqpoint{2.779427in}{1.942037in}}%
\pgfpathlineto{\pgfqpoint{2.823849in}{1.914023in}}%
\pgfpathlineto{\pgfqpoint{2.855676in}{1.895316in}}%
\pgfpathlineto{\pgfqpoint{2.811241in}{1.923032in}}%
\pgfpathlineto{\pgfqpoint{2.779427in}{1.942037in}}%
\pgfpathclose%
\pgfusepath{stroke,fill}%
\end{pgfscope}%
\begin{pgfscope}%
\pgfpathrectangle{\pgfqpoint{0.050000in}{0.050000in}}{\pgfqpoint{4.900000in}{4.900000in}}%
\pgfusepath{clip}%
\pgfsetbuttcap%
\pgfsetroundjoin%
\definecolor{currentfill}{rgb}{0.878570,0.878570,0.878570}%
\pgfsetfillcolor{currentfill}%
\pgfsetfillopacity{0.900000}%
\pgfsetlinewidth{0.501875pt}%
\definecolor{currentstroke}{rgb}{0.200000,0.200000,0.200000}%
\pgfsetstrokecolor{currentstroke}%
\pgfsetdash{}{0pt}%
\pgfpathmoveto{\pgfqpoint{2.442671in}{2.186545in}}%
\pgfpathlineto{\pgfqpoint{2.487732in}{2.145655in}}%
\pgfpathlineto{\pgfqpoint{2.519035in}{2.125928in}}%
\pgfpathlineto{\pgfqpoint{2.473969in}{2.166599in}}%
\pgfpathlineto{\pgfqpoint{2.442671in}{2.186545in}}%
\pgfpathclose%
\pgfusepath{stroke,fill}%
\end{pgfscope}%
\begin{pgfscope}%
\pgfpathrectangle{\pgfqpoint{0.050000in}{0.050000in}}{\pgfqpoint{4.900000in}{4.900000in}}%
\pgfusepath{clip}%
\pgfsetbuttcap%
\pgfsetroundjoin%
\definecolor{currentfill}{rgb}{0.988927,0.988927,0.988927}%
\pgfsetfillcolor{currentfill}%
\pgfsetfillopacity{0.900000}%
\pgfsetlinewidth{0.501875pt}%
\definecolor{currentstroke}{rgb}{0.200000,0.200000,0.200000}%
\pgfsetstrokecolor{currentstroke}%
\pgfsetdash{}{0pt}%
\pgfpathmoveto{\pgfqpoint{3.313580in}{1.710055in}}%
\pgfpathlineto{\pgfqpoint{3.357427in}{1.711109in}}%
\pgfpathlineto{\pgfqpoint{3.390097in}{1.692110in}}%
\pgfpathlineto{\pgfqpoint{3.346226in}{1.691160in}}%
\pgfpathlineto{\pgfqpoint{3.313580in}{1.710055in}}%
\pgfpathclose%
\pgfusepath{stroke,fill}%
\end{pgfscope}%
\begin{pgfscope}%
\pgfpathrectangle{\pgfqpoint{0.050000in}{0.050000in}}{\pgfqpoint{4.900000in}{4.900000in}}%
\pgfusepath{clip}%
\pgfsetbuttcap%
\pgfsetroundjoin%
\definecolor{currentfill}{rgb}{0.547328,0.547328,0.547328}%
\pgfsetfillcolor{currentfill}%
\pgfsetfillopacity{0.900000}%
\pgfsetlinewidth{0.501875pt}%
\definecolor{currentstroke}{rgb}{0.200000,0.200000,0.200000}%
\pgfsetstrokecolor{currentstroke}%
\pgfsetdash{}{0pt}%
\pgfpathmoveto{\pgfqpoint{1.477117in}{2.965075in}}%
\pgfpathlineto{\pgfqpoint{1.524193in}{2.920263in}}%
\pgfpathlineto{\pgfqpoint{1.554031in}{2.901062in}}%
\pgfpathlineto{\pgfqpoint{1.506945in}{2.945903in}}%
\pgfpathlineto{\pgfqpoint{1.477117in}{2.965075in}}%
\pgfpathclose%
\pgfusepath{stroke,fill}%
\end{pgfscope}%
\begin{pgfscope}%
\pgfpathrectangle{\pgfqpoint{0.050000in}{0.050000in}}{\pgfqpoint{4.900000in}{4.900000in}}%
\pgfusepath{clip}%
\pgfsetbuttcap%
\pgfsetroundjoin%
\definecolor{currentfill}{rgb}{0.808781,0.808781,0.808781}%
\pgfsetfillcolor{currentfill}%
\pgfsetfillopacity{0.900000}%
\pgfsetlinewidth{0.501875pt}%
\definecolor{currentstroke}{rgb}{0.200000,0.200000,0.200000}%
\pgfsetstrokecolor{currentstroke}%
\pgfsetdash{}{0pt}%
\pgfpathmoveto{\pgfqpoint{2.182166in}{2.395286in}}%
\pgfpathlineto{\pgfqpoint{2.227777in}{2.351778in}}%
\pgfpathlineto{\pgfqpoint{2.258687in}{2.331781in}}%
\pgfpathlineto{\pgfqpoint{2.213071in}{2.375273in}}%
\pgfpathlineto{\pgfqpoint{2.182166in}{2.395286in}}%
\pgfpathclose%
\pgfusepath{stroke,fill}%
\end{pgfscope}%
\begin{pgfscope}%
\pgfpathrectangle{\pgfqpoint{0.050000in}{0.050000in}}{\pgfqpoint{4.900000in}{4.900000in}}%
\pgfusepath{clip}%
\pgfsetbuttcap%
\pgfsetroundjoin%
\definecolor{currentfill}{rgb}{0.998155,0.998155,0.998155}%
\pgfsetfillcolor{currentfill}%
\pgfsetfillopacity{0.900000}%
\pgfsetlinewidth{0.501875pt}%
\definecolor{currentstroke}{rgb}{0.200000,0.200000,0.200000}%
\pgfsetstrokecolor{currentstroke}%
\pgfsetdash{}{0pt}%
\pgfpathmoveto{\pgfqpoint{3.663779in}{1.660400in}}%
\pgfpathlineto{\pgfqpoint{3.707360in}{1.670763in}}%
\pgfpathlineto{\pgfqpoint{3.740567in}{1.650940in}}%
\pgfpathlineto{\pgfqpoint{3.696960in}{1.640674in}}%
\pgfpathlineto{\pgfqpoint{3.663779in}{1.660400in}}%
\pgfpathclose%
\pgfusepath{stroke,fill}%
\end{pgfscope}%
\begin{pgfscope}%
\pgfpathrectangle{\pgfqpoint{0.050000in}{0.050000in}}{\pgfqpoint{4.900000in}{4.900000in}}%
\pgfusepath{clip}%
\pgfsetbuttcap%
\pgfsetroundjoin%
\definecolor{currentfill}{rgb}{0.935163,0.935163,0.935163}%
\pgfsetfillcolor{currentfill}%
\pgfsetfillopacity{0.900000}%
\pgfsetlinewidth{0.501875pt}%
\definecolor{currentstroke}{rgb}{0.200000,0.200000,0.200000}%
\pgfsetstrokecolor{currentstroke}%
\pgfsetdash{}{0pt}%
\pgfpathmoveto{\pgfqpoint{2.703150in}{1.993144in}}%
\pgfpathlineto{\pgfqpoint{2.747712in}{1.960979in}}%
\pgfpathlineto{\pgfqpoint{2.779427in}{1.942037in}}%
\pgfpathlineto{\pgfqpoint{2.734855in}{1.973874in}}%
\pgfpathlineto{\pgfqpoint{2.703150in}{1.993144in}}%
\pgfpathclose%
\pgfusepath{stroke,fill}%
\end{pgfscope}%
\begin{pgfscope}%
\pgfpathrectangle{\pgfqpoint{0.050000in}{0.050000in}}{\pgfqpoint{4.900000in}{4.900000in}}%
\pgfusepath{clip}%
\pgfsetbuttcap%
\pgfsetroundjoin%
\definecolor{currentfill}{rgb}{0.724983,0.724983,0.724983}%
\pgfsetfillcolor{currentfill}%
\pgfsetfillopacity{0.900000}%
\pgfsetlinewidth{0.501875pt}%
\definecolor{currentstroke}{rgb}{0.200000,0.200000,0.200000}%
\pgfsetstrokecolor{currentstroke}%
\pgfsetdash{}{0pt}%
\pgfpathmoveto{\pgfqpoint{1.921570in}{2.606462in}}%
\pgfpathlineto{\pgfqpoint{1.967731in}{2.562386in}}%
\pgfpathlineto{\pgfqpoint{1.998251in}{2.542637in}}%
\pgfpathlineto{\pgfqpoint{1.952083in}{2.586739in}}%
\pgfpathlineto{\pgfqpoint{1.921570in}{2.606462in}}%
\pgfpathclose%
\pgfusepath{stroke,fill}%
\end{pgfscope}%
\begin{pgfscope}%
\pgfpathrectangle{\pgfqpoint{0.050000in}{0.050000in}}{\pgfqpoint{4.900000in}{4.900000in}}%
\pgfusepath{clip}%
\pgfsetbuttcap%
\pgfsetroundjoin%
\definecolor{currentfill}{rgb}{0.992618,0.992618,0.992618}%
\pgfsetfillcolor{currentfill}%
\pgfsetfillopacity{0.900000}%
\pgfsetlinewidth{0.501875pt}%
\definecolor{currentstroke}{rgb}{0.200000,0.200000,0.200000}%
\pgfsetstrokecolor{currentstroke}%
\pgfsetdash{}{0pt}%
\pgfpathmoveto{\pgfqpoint{3.390097in}{1.692110in}}%
\pgfpathlineto{\pgfqpoint{3.433896in}{1.695809in}}%
\pgfpathlineto{\pgfqpoint{3.466692in}{1.676618in}}%
\pgfpathlineto{\pgfqpoint{3.422869in}{1.673037in}}%
\pgfpathlineto{\pgfqpoint{3.390097in}{1.692110in}}%
\pgfpathclose%
\pgfusepath{stroke,fill}%
\end{pgfscope}%
\begin{pgfscope}%
\pgfpathrectangle{\pgfqpoint{0.050000in}{0.050000in}}{\pgfqpoint{4.900000in}{4.900000in}}%
\pgfusepath{clip}%
\pgfsetbuttcap%
\pgfsetroundjoin%
\definecolor{currentfill}{rgb}{1.000000,1.000000,1.000000}%
\pgfsetfillcolor{currentfill}%
\pgfsetfillopacity{0.900000}%
\pgfsetlinewidth{0.501875pt}%
\definecolor{currentstroke}{rgb}{0.200000,0.200000,0.200000}%
\pgfsetstrokecolor{currentstroke}%
\pgfsetdash{}{0pt}%
\pgfpathmoveto{\pgfqpoint{3.937789in}{1.646250in}}%
\pgfpathlineto{\pgfqpoint{3.981118in}{1.658868in}}%
\pgfpathlineto{\pgfqpoint{4.014743in}{1.638594in}}%
\pgfpathlineto{\pgfqpoint{3.971387in}{1.626033in}}%
\pgfpathlineto{\pgfqpoint{3.937789in}{1.646250in}}%
\pgfpathclose%
\pgfusepath{stroke,fill}%
\end{pgfscope}%
\begin{pgfscope}%
\pgfpathrectangle{\pgfqpoint{0.050000in}{0.050000in}}{\pgfqpoint{4.900000in}{4.900000in}}%
\pgfusepath{clip}%
\pgfsetbuttcap%
\pgfsetroundjoin%
\definecolor{currentfill}{rgb}{0.619423,0.619423,0.619423}%
\pgfsetfillcolor{currentfill}%
\pgfsetfillopacity{0.900000}%
\pgfsetlinewidth{0.501875pt}%
\definecolor{currentstroke}{rgb}{0.200000,0.200000,0.200000}%
\pgfsetstrokecolor{currentstroke}%
\pgfsetdash{}{0pt}%
\pgfpathmoveto{\pgfqpoint{1.660887in}{2.817775in}}%
\pgfpathlineto{\pgfqpoint{1.707599in}{2.773252in}}%
\pgfpathlineto{\pgfqpoint{1.737727in}{2.753820in}}%
\pgfpathlineto{\pgfqpoint{1.691007in}{2.798370in}}%
\pgfpathlineto{\pgfqpoint{1.660887in}{2.817775in}}%
\pgfpathclose%
\pgfusepath{stroke,fill}%
\end{pgfscope}%
\begin{pgfscope}%
\pgfpathrectangle{\pgfqpoint{0.050000in}{0.050000in}}{\pgfqpoint{4.900000in}{4.900000in}}%
\pgfusepath{clip}%
\pgfsetbuttcap%
\pgfsetroundjoin%
\definecolor{currentfill}{rgb}{0.858762,0.858762,0.858762}%
\pgfsetfillcolor{currentfill}%
\pgfsetfillopacity{0.900000}%
\pgfsetlinewidth{0.501875pt}%
\definecolor{currentstroke}{rgb}{0.200000,0.200000,0.200000}%
\pgfsetstrokecolor{currentstroke}%
\pgfsetdash{}{0pt}%
\pgfpathmoveto{\pgfqpoint{2.366224in}{2.248686in}}%
\pgfpathlineto{\pgfqpoint{2.411468in}{2.206460in}}%
\pgfpathlineto{\pgfqpoint{2.442671in}{2.186545in}}%
\pgfpathlineto{\pgfqpoint{2.397423in}{2.228631in}}%
\pgfpathlineto{\pgfqpoint{2.366224in}{2.248686in}}%
\pgfpathclose%
\pgfusepath{stroke,fill}%
\end{pgfscope}%
\begin{pgfscope}%
\pgfpathrectangle{\pgfqpoint{0.050000in}{0.050000in}}{\pgfqpoint{4.900000in}{4.900000in}}%
\pgfusepath{clip}%
\pgfsetbuttcap%
\pgfsetroundjoin%
\definecolor{currentfill}{rgb}{0.998155,0.998155,0.998155}%
\pgfsetfillcolor{currentfill}%
\pgfsetfillopacity{0.900000}%
\pgfsetlinewidth{0.501875pt}%
\definecolor{currentstroke}{rgb}{0.200000,0.200000,0.200000}%
\pgfsetstrokecolor{currentstroke}%
\pgfsetdash{}{0pt}%
\pgfpathmoveto{\pgfqpoint{3.740567in}{1.650940in}}%
\pgfpathlineto{\pgfqpoint{3.784095in}{1.662160in}}%
\pgfpathlineto{\pgfqpoint{3.817432in}{1.642197in}}%
\pgfpathlineto{\pgfqpoint{3.773876in}{1.631060in}}%
\pgfpathlineto{\pgfqpoint{3.740567in}{1.650940in}}%
\pgfpathclose%
\pgfusepath{stroke,fill}%
\end{pgfscope}%
\begin{pgfscope}%
\pgfpathrectangle{\pgfqpoint{0.050000in}{0.050000in}}{\pgfqpoint{4.900000in}{4.900000in}}%
\pgfusepath{clip}%
\pgfsetbuttcap%
\pgfsetroundjoin%
\definecolor{currentfill}{rgb}{0.517186,0.517186,0.517186}%
\pgfsetfillcolor{currentfill}%
\pgfsetfillopacity{0.900000}%
\pgfsetlinewidth{0.501875pt}%
\definecolor{currentstroke}{rgb}{0.200000,0.200000,0.200000}%
\pgfsetstrokecolor{currentstroke}%
\pgfsetdash{}{0pt}%
\pgfpathmoveto{\pgfqpoint{1.400117in}{3.029161in}}%
\pgfpathlineto{\pgfqpoint{1.447380in}{2.984190in}}%
\pgfpathlineto{\pgfqpoint{1.477117in}{2.965075in}}%
\pgfpathlineto{\pgfqpoint{1.429845in}{3.010075in}}%
\pgfpathlineto{\pgfqpoint{1.400117in}{3.029161in}}%
\pgfpathclose%
\pgfusepath{stroke,fill}%
\end{pgfscope}%
\begin{pgfscope}%
\pgfpathrectangle{\pgfqpoint{0.050000in}{0.050000in}}{\pgfqpoint{4.900000in}{4.900000in}}%
\pgfusepath{clip}%
\pgfsetbuttcap%
\pgfsetroundjoin%
\definecolor{currentfill}{rgb}{0.994464,0.994464,0.994464}%
\pgfsetfillcolor{currentfill}%
\pgfsetfillopacity{0.900000}%
\pgfsetlinewidth{0.501875pt}%
\definecolor{currentstroke}{rgb}{0.200000,0.200000,0.200000}%
\pgfsetstrokecolor{currentstroke}%
\pgfsetdash{}{0pt}%
\pgfpathmoveto{\pgfqpoint{3.466692in}{1.676618in}}%
\pgfpathlineto{\pgfqpoint{3.510443in}{1.682515in}}%
\pgfpathlineto{\pgfqpoint{3.543367in}{1.663135in}}%
\pgfpathlineto{\pgfqpoint{3.499590in}{1.657358in}}%
\pgfpathlineto{\pgfqpoint{3.466692in}{1.676618in}}%
\pgfpathclose%
\pgfusepath{stroke,fill}%
\end{pgfscope}%
\begin{pgfscope}%
\pgfpathrectangle{\pgfqpoint{0.050000in}{0.050000in}}{\pgfqpoint{4.900000in}{4.900000in}}%
\pgfusepath{clip}%
\pgfsetbuttcap%
\pgfsetroundjoin%
\definecolor{currentfill}{rgb}{0.918185,0.918185,0.918185}%
\pgfsetfillcolor{currentfill}%
\pgfsetfillopacity{0.900000}%
\pgfsetlinewidth{0.501875pt}%
\definecolor{currentstroke}{rgb}{0.200000,0.200000,0.200000}%
\pgfsetstrokecolor{currentstroke}%
\pgfsetdash{}{0pt}%
\pgfpathmoveto{\pgfqpoint{2.626823in}{2.048092in}}%
\pgfpathlineto{\pgfqpoint{2.671542in}{2.012364in}}%
\pgfpathlineto{\pgfqpoint{2.703150in}{1.993144in}}%
\pgfpathlineto{\pgfqpoint{2.658423in}{2.028548in}}%
\pgfpathlineto{\pgfqpoint{2.626823in}{2.048092in}}%
\pgfpathclose%
\pgfusepath{stroke,fill}%
\end{pgfscope}%
\begin{pgfscope}%
\pgfpathrectangle{\pgfqpoint{0.050000in}{0.050000in}}{\pgfqpoint{4.900000in}{4.900000in}}%
\pgfusepath{clip}%
\pgfsetbuttcap%
\pgfsetroundjoin%
\definecolor{currentfill}{rgb}{0.781223,0.781223,0.781223}%
\pgfsetfillcolor{currentfill}%
\pgfsetfillopacity{0.900000}%
\pgfsetlinewidth{0.501875pt}%
\definecolor{currentstroke}{rgb}{0.200000,0.200000,0.200000}%
\pgfsetstrokecolor{currentstroke}%
\pgfsetdash{}{0pt}%
\pgfpathmoveto{\pgfqpoint{2.105558in}{2.459001in}}%
\pgfpathlineto{\pgfqpoint{2.151355in}{2.415245in}}%
\pgfpathlineto{\pgfqpoint{2.182166in}{2.395286in}}%
\pgfpathlineto{\pgfqpoint{2.136364in}{2.439053in}}%
\pgfpathlineto{\pgfqpoint{2.105558in}{2.459001in}}%
\pgfpathclose%
\pgfusepath{stroke,fill}%
\end{pgfscope}%
\begin{pgfscope}%
\pgfpathrectangle{\pgfqpoint{0.050000in}{0.050000in}}{\pgfqpoint{4.900000in}{4.900000in}}%
\pgfusepath{clip}%
\pgfsetbuttcap%
\pgfsetroundjoin%
\definecolor{currentfill}{rgb}{1.000000,1.000000,1.000000}%
\pgfsetfillcolor{currentfill}%
\pgfsetfillopacity{0.900000}%
\pgfsetlinewidth{0.501875pt}%
\definecolor{currentstroke}{rgb}{0.200000,0.200000,0.200000}%
\pgfsetstrokecolor{currentstroke}%
\pgfsetdash{}{0pt}%
\pgfpathmoveto{\pgfqpoint{4.014743in}{1.638594in}}%
\pgfpathlineto{\pgfqpoint{4.058008in}{1.651321in}}%
\pgfpathlineto{\pgfqpoint{4.091763in}{1.630932in}}%
\pgfpathlineto{\pgfqpoint{4.048471in}{1.618259in}}%
\pgfpathlineto{\pgfqpoint{4.014743in}{1.638594in}}%
\pgfpathclose%
\pgfusepath{stroke,fill}%
\end{pgfscope}%
\begin{pgfscope}%
\pgfpathrectangle{\pgfqpoint{0.050000in}{0.050000in}}{\pgfqpoint{4.900000in}{4.900000in}}%
\pgfusepath{clip}%
\pgfsetbuttcap%
\pgfsetroundjoin%
\definecolor{currentfill}{rgb}{0.691396,0.691396,0.691396}%
\pgfsetfillcolor{currentfill}%
\pgfsetfillopacity{0.900000}%
\pgfsetlinewidth{0.501875pt}%
\definecolor{currentstroke}{rgb}{0.200000,0.200000,0.200000}%
\pgfsetstrokecolor{currentstroke}%
\pgfsetdash{}{0pt}%
\pgfpathmoveto{\pgfqpoint{1.844802in}{2.670359in}}%
\pgfpathlineto{\pgfqpoint{1.891150in}{2.626125in}}%
\pgfpathlineto{\pgfqpoint{1.921570in}{2.606462in}}%
\pgfpathlineto{\pgfqpoint{1.875215in}{2.650722in}}%
\pgfpathlineto{\pgfqpoint{1.844802in}{2.670359in}}%
\pgfpathclose%
\pgfusepath{stroke,fill}%
\end{pgfscope}%
\begin{pgfscope}%
\pgfpathrectangle{\pgfqpoint{0.050000in}{0.050000in}}{\pgfqpoint{4.900000in}{4.900000in}}%
\pgfusepath{clip}%
\pgfsetbuttcap%
\pgfsetroundjoin%
\definecolor{currentfill}{rgb}{0.996309,0.996309,0.996309}%
\pgfsetfillcolor{currentfill}%
\pgfsetfillopacity{0.900000}%
\pgfsetlinewidth{0.501875pt}%
\definecolor{currentstroke}{rgb}{0.200000,0.200000,0.200000}%
\pgfsetstrokecolor{currentstroke}%
\pgfsetdash{}{0pt}%
\pgfpathmoveto{\pgfqpoint{3.543367in}{1.663135in}}%
\pgfpathlineto{\pgfqpoint{3.587071in}{1.670827in}}%
\pgfpathlineto{\pgfqpoint{3.620123in}{1.651268in}}%
\pgfpathlineto{\pgfqpoint{3.576392in}{1.643691in}}%
\pgfpathlineto{\pgfqpoint{3.543367in}{1.663135in}}%
\pgfpathclose%
\pgfusepath{stroke,fill}%
\end{pgfscope}%
\begin{pgfscope}%
\pgfpathrectangle{\pgfqpoint{0.050000in}{0.050000in}}{\pgfqpoint{4.900000in}{4.900000in}}%
\pgfusepath{clip}%
\pgfsetbuttcap%
\pgfsetroundjoin%
\definecolor{currentfill}{rgb}{1.000000,1.000000,1.000000}%
\pgfsetfillcolor{currentfill}%
\pgfsetfillopacity{0.900000}%
\pgfsetlinewidth{0.501875pt}%
\definecolor{currentstroke}{rgb}{0.200000,0.200000,0.200000}%
\pgfsetstrokecolor{currentstroke}%
\pgfsetdash{}{0pt}%
\pgfpathmoveto{\pgfqpoint{3.817432in}{1.642197in}}%
\pgfpathlineto{\pgfqpoint{3.860905in}{1.654051in}}%
\pgfpathlineto{\pgfqpoint{3.894372in}{1.633957in}}%
\pgfpathlineto{\pgfqpoint{3.850872in}{1.622175in}}%
\pgfpathlineto{\pgfqpoint{3.817432in}{1.642197in}}%
\pgfpathclose%
\pgfusepath{stroke,fill}%
\end{pgfscope}%
\begin{pgfscope}%
\pgfpathrectangle{\pgfqpoint{0.050000in}{0.050000in}}{\pgfqpoint{4.900000in}{4.900000in}}%
\pgfusepath{clip}%
\pgfsetbuttcap%
\pgfsetroundjoin%
\definecolor{currentfill}{rgb}{0.581776,0.581776,0.581776}%
\pgfsetfillcolor{currentfill}%
\pgfsetfillopacity{0.900000}%
\pgfsetlinewidth{0.501875pt}%
\definecolor{currentstroke}{rgb}{0.200000,0.200000,0.200000}%
\pgfsetstrokecolor{currentstroke}%
\pgfsetdash{}{0pt}%
\pgfpathmoveto{\pgfqpoint{1.583959in}{2.881803in}}%
\pgfpathlineto{\pgfqpoint{1.630858in}{2.837121in}}%
\pgfpathlineto{\pgfqpoint{1.660887in}{2.817775in}}%
\pgfpathlineto{\pgfqpoint{1.613979in}{2.862485in}}%
\pgfpathlineto{\pgfqpoint{1.583959in}{2.881803in}}%
\pgfpathclose%
\pgfusepath{stroke,fill}%
\end{pgfscope}%
\begin{pgfscope}%
\pgfpathrectangle{\pgfqpoint{0.050000in}{0.050000in}}{\pgfqpoint{4.900000in}{4.900000in}}%
\pgfusepath{clip}%
\pgfsetbuttcap%
\pgfsetroundjoin%
\definecolor{currentfill}{rgb}{0.898378,0.898378,0.898378}%
\pgfsetfillcolor{currentfill}%
\pgfsetfillopacity{0.900000}%
\pgfsetlinewidth{0.501875pt}%
\definecolor{currentstroke}{rgb}{0.200000,0.200000,0.200000}%
\pgfsetstrokecolor{currentstroke}%
\pgfsetdash{}{0pt}%
\pgfpathmoveto{\pgfqpoint{2.550434in}{2.106171in}}%
\pgfpathlineto{\pgfqpoint{2.595322in}{2.067597in}}%
\pgfpathlineto{\pgfqpoint{2.626823in}{2.048092in}}%
\pgfpathlineto{\pgfqpoint{2.581930in}{2.086381in}}%
\pgfpathlineto{\pgfqpoint{2.550434in}{2.106171in}}%
\pgfpathclose%
\pgfusepath{stroke,fill}%
\end{pgfscope}%
\begin{pgfscope}%
\pgfpathrectangle{\pgfqpoint{0.050000in}{0.050000in}}{\pgfqpoint{4.900000in}{4.900000in}}%
\pgfusepath{clip}%
\pgfsetbuttcap%
\pgfsetroundjoin%
\definecolor{currentfill}{rgb}{0.832895,0.832895,0.832895}%
\pgfsetfillcolor{currentfill}%
\pgfsetfillopacity{0.900000}%
\pgfsetlinewidth{0.501875pt}%
\definecolor{currentstroke}{rgb}{0.200000,0.200000,0.200000}%
\pgfsetstrokecolor{currentstroke}%
\pgfsetdash{}{0pt}%
\pgfpathmoveto{\pgfqpoint{2.289691in}{2.311737in}}%
\pgfpathlineto{\pgfqpoint{2.335121in}{2.268705in}}%
\pgfpathlineto{\pgfqpoint{2.366224in}{2.248686in}}%
\pgfpathlineto{\pgfqpoint{2.320790in}{2.291647in}}%
\pgfpathlineto{\pgfqpoint{2.289691in}{2.311737in}}%
\pgfpathclose%
\pgfusepath{stroke,fill}%
\end{pgfscope}%
\begin{pgfscope}%
\pgfpathrectangle{\pgfqpoint{0.050000in}{0.050000in}}{\pgfqpoint{4.900000in}{4.900000in}}%
\pgfusepath{clip}%
\pgfsetbuttcap%
\pgfsetroundjoin%
\definecolor{currentfill}{rgb}{0.970473,0.970473,0.970473}%
\pgfsetfillcolor{currentfill}%
\pgfsetfillopacity{0.900000}%
\pgfsetlinewidth{0.501875pt}%
\definecolor{currentstroke}{rgb}{0.200000,0.200000,0.200000}%
\pgfsetstrokecolor{currentstroke}%
\pgfsetdash{}{0pt}%
\pgfpathmoveto{\pgfqpoint{3.040320in}{1.797520in}}%
\pgfpathlineto{\pgfqpoint{3.084438in}{1.783490in}}%
\pgfpathlineto{\pgfqpoint{3.116715in}{1.764941in}}%
\pgfpathlineto{\pgfqpoint{3.072578in}{1.778880in}}%
\pgfpathlineto{\pgfqpoint{3.040320in}{1.797520in}}%
\pgfpathclose%
\pgfusepath{stroke,fill}%
\end{pgfscope}%
\begin{pgfscope}%
\pgfpathrectangle{\pgfqpoint{0.050000in}{0.050000in}}{\pgfqpoint{4.900000in}{4.900000in}}%
\pgfusepath{clip}%
\pgfsetbuttcap%
\pgfsetroundjoin%
\definecolor{currentfill}{rgb}{0.977855,0.977855,0.977855}%
\pgfsetfillcolor{currentfill}%
\pgfsetfillopacity{0.900000}%
\pgfsetlinewidth{0.501875pt}%
\definecolor{currentstroke}{rgb}{0.200000,0.200000,0.200000}%
\pgfsetstrokecolor{currentstroke}%
\pgfsetdash{}{0pt}%
\pgfpathmoveto{\pgfqpoint{3.116715in}{1.764941in}}%
\pgfpathlineto{\pgfqpoint{3.160759in}{1.755254in}}%
\pgfpathlineto{\pgfqpoint{3.193157in}{1.736636in}}%
\pgfpathlineto{\pgfqpoint{3.149094in}{1.746299in}}%
\pgfpathlineto{\pgfqpoint{3.116715in}{1.764941in}}%
\pgfpathclose%
\pgfusepath{stroke,fill}%
\end{pgfscope}%
\begin{pgfscope}%
\pgfpathrectangle{\pgfqpoint{0.050000in}{0.050000in}}{\pgfqpoint{4.900000in}{4.900000in}}%
\pgfusepath{clip}%
\pgfsetbuttcap%
\pgfsetroundjoin%
\definecolor{currentfill}{rgb}{0.964937,0.964937,0.964937}%
\pgfsetfillcolor{currentfill}%
\pgfsetfillopacity{0.900000}%
\pgfsetlinewidth{0.501875pt}%
\definecolor{currentstroke}{rgb}{0.200000,0.200000,0.200000}%
\pgfsetstrokecolor{currentstroke}%
\pgfsetdash{}{0pt}%
\pgfpathmoveto{\pgfqpoint{2.963954in}{1.834672in}}%
\pgfpathlineto{\pgfqpoint{3.008162in}{1.816070in}}%
\pgfpathlineto{\pgfqpoint{3.040320in}{1.797520in}}%
\pgfpathlineto{\pgfqpoint{2.996095in}{1.815959in}}%
\pgfpathlineto{\pgfqpoint{2.963954in}{1.834672in}}%
\pgfpathclose%
\pgfusepath{stroke,fill}%
\end{pgfscope}%
\begin{pgfscope}%
\pgfpathrectangle{\pgfqpoint{0.050000in}{0.050000in}}{\pgfqpoint{4.900000in}{4.900000in}}%
\pgfusepath{clip}%
\pgfsetbuttcap%
\pgfsetroundjoin%
\definecolor{currentfill}{rgb}{0.482737,0.482737,0.482737}%
\pgfsetfillcolor{currentfill}%
\pgfsetfillopacity{0.900000}%
\pgfsetlinewidth{0.501875pt}%
\definecolor{currentstroke}{rgb}{0.200000,0.200000,0.200000}%
\pgfsetstrokecolor{currentstroke}%
\pgfsetdash{}{0pt}%
\pgfpathmoveto{\pgfqpoint{1.323030in}{3.093319in}}%
\pgfpathlineto{\pgfqpoint{1.370480in}{3.048189in}}%
\pgfpathlineto{\pgfqpoint{1.400117in}{3.029161in}}%
\pgfpathlineto{\pgfqpoint{1.352657in}{3.074320in}}%
\pgfpathlineto{\pgfqpoint{1.323030in}{3.093319in}}%
\pgfpathclose%
\pgfusepath{stroke,fill}%
\end{pgfscope}%
\begin{pgfscope}%
\pgfpathrectangle{\pgfqpoint{0.050000in}{0.050000in}}{\pgfqpoint{4.900000in}{4.900000in}}%
\pgfusepath{clip}%
\pgfsetbuttcap%
\pgfsetroundjoin%
\definecolor{currentfill}{rgb}{0.981546,0.981546,0.981546}%
\pgfsetfillcolor{currentfill}%
\pgfsetfillopacity{0.900000}%
\pgfsetlinewidth{0.501875pt}%
\definecolor{currentstroke}{rgb}{0.200000,0.200000,0.200000}%
\pgfsetstrokecolor{currentstroke}%
\pgfsetdash{}{0pt}%
\pgfpathmoveto{\pgfqpoint{3.193157in}{1.736636in}}%
\pgfpathlineto{\pgfqpoint{3.237136in}{1.730937in}}%
\pgfpathlineto{\pgfqpoint{3.269658in}{1.712199in}}%
\pgfpathlineto{\pgfqpoint{3.225657in}{1.717930in}}%
\pgfpathlineto{\pgfqpoint{3.193157in}{1.736636in}}%
\pgfpathclose%
\pgfusepath{stroke,fill}%
\end{pgfscope}%
\begin{pgfscope}%
\pgfpathrectangle{\pgfqpoint{0.050000in}{0.050000in}}{\pgfqpoint{4.900000in}{4.900000in}}%
\pgfusepath{clip}%
\pgfsetbuttcap%
\pgfsetroundjoin%
\definecolor{currentfill}{rgb}{0.757109,0.757109,0.757109}%
\pgfsetfillcolor{currentfill}%
\pgfsetfillopacity{0.900000}%
\pgfsetlinewidth{0.501875pt}%
\definecolor{currentstroke}{rgb}{0.200000,0.200000,0.200000}%
\pgfsetstrokecolor{currentstroke}%
\pgfsetdash{}{0pt}%
\pgfpathmoveto{\pgfqpoint{2.028864in}{2.522829in}}%
\pgfpathlineto{\pgfqpoint{2.074846in}{2.478890in}}%
\pgfpathlineto{\pgfqpoint{2.105558in}{2.459001in}}%
\pgfpathlineto{\pgfqpoint{2.059570in}{2.502961in}}%
\pgfpathlineto{\pgfqpoint{2.028864in}{2.522829in}}%
\pgfpathclose%
\pgfusepath{stroke,fill}%
\end{pgfscope}%
\begin{pgfscope}%
\pgfpathrectangle{\pgfqpoint{0.050000in}{0.050000in}}{\pgfqpoint{4.900000in}{4.900000in}}%
\pgfusepath{clip}%
\pgfsetbuttcap%
\pgfsetroundjoin%
\definecolor{currentfill}{rgb}{0.955709,0.955709,0.955709}%
\pgfsetfillcolor{currentfill}%
\pgfsetfillopacity{0.900000}%
\pgfsetlinewidth{0.501875pt}%
\definecolor{currentstroke}{rgb}{0.200000,0.200000,0.200000}%
\pgfsetstrokecolor{currentstroke}%
\pgfsetdash{}{0pt}%
\pgfpathmoveto{\pgfqpoint{2.887601in}{1.876533in}}%
\pgfpathlineto{\pgfqpoint{2.931914in}{1.853300in}}%
\pgfpathlineto{\pgfqpoint{2.963954in}{1.834672in}}%
\pgfpathlineto{\pgfqpoint{2.919627in}{1.857673in}}%
\pgfpathlineto{\pgfqpoint{2.887601in}{1.876533in}}%
\pgfpathclose%
\pgfusepath{stroke,fill}%
\end{pgfscope}%
\begin{pgfscope}%
\pgfpathrectangle{\pgfqpoint{0.050000in}{0.050000in}}{\pgfqpoint{4.900000in}{4.900000in}}%
\pgfusepath{clip}%
\pgfsetbuttcap%
\pgfsetroundjoin%
\definecolor{currentfill}{rgb}{0.996309,0.996309,0.996309}%
\pgfsetfillcolor{currentfill}%
\pgfsetfillopacity{0.900000}%
\pgfsetlinewidth{0.501875pt}%
\definecolor{currentstroke}{rgb}{0.200000,0.200000,0.200000}%
\pgfsetstrokecolor{currentstroke}%
\pgfsetdash{}{0pt}%
\pgfpathmoveto{\pgfqpoint{3.620123in}{1.651268in}}%
\pgfpathlineto{\pgfqpoint{3.663779in}{1.660400in}}%
\pgfpathlineto{\pgfqpoint{3.696960in}{1.640674in}}%
\pgfpathlineto{\pgfqpoint{3.653277in}{1.631647in}}%
\pgfpathlineto{\pgfqpoint{3.620123in}{1.651268in}}%
\pgfpathclose%
\pgfusepath{stroke,fill}%
\end{pgfscope}%
\begin{pgfscope}%
\pgfpathrectangle{\pgfqpoint{0.050000in}{0.050000in}}{\pgfqpoint{4.900000in}{4.900000in}}%
\pgfusepath{clip}%
\pgfsetbuttcap%
\pgfsetroundjoin%
\definecolor{currentfill}{rgb}{0.985236,0.985236,0.985236}%
\pgfsetfillcolor{currentfill}%
\pgfsetfillopacity{0.900000}%
\pgfsetlinewidth{0.501875pt}%
\definecolor{currentstroke}{rgb}{0.200000,0.200000,0.200000}%
\pgfsetstrokecolor{currentstroke}%
\pgfsetdash{}{0pt}%
\pgfpathmoveto{\pgfqpoint{3.269658in}{1.712199in}}%
\pgfpathlineto{\pgfqpoint{3.313580in}{1.710055in}}%
\pgfpathlineto{\pgfqpoint{3.346226in}{1.691160in}}%
\pgfpathlineto{\pgfqpoint{3.302281in}{1.693376in}}%
\pgfpathlineto{\pgfqpoint{3.269658in}{1.712199in}}%
\pgfpathclose%
\pgfusepath{stroke,fill}%
\end{pgfscope}%
\begin{pgfscope}%
\pgfpathrectangle{\pgfqpoint{0.050000in}{0.050000in}}{\pgfqpoint{4.900000in}{4.900000in}}%
\pgfusepath{clip}%
\pgfsetbuttcap%
\pgfsetroundjoin%
\definecolor{currentfill}{rgb}{0.653010,0.653010,0.653010}%
\pgfsetfillcolor{currentfill}%
\pgfsetfillopacity{0.900000}%
\pgfsetlinewidth{0.501875pt}%
\definecolor{currentstroke}{rgb}{0.200000,0.200000,0.200000}%
\pgfsetstrokecolor{currentstroke}%
\pgfsetdash{}{0pt}%
\pgfpathmoveto{\pgfqpoint{1.767948in}{2.734329in}}%
\pgfpathlineto{\pgfqpoint{1.814482in}{2.689936in}}%
\pgfpathlineto{\pgfqpoint{1.844802in}{2.670359in}}%
\pgfpathlineto{\pgfqpoint{1.798261in}{2.714778in}}%
\pgfpathlineto{\pgfqpoint{1.767948in}{2.734329in}}%
\pgfpathclose%
\pgfusepath{stroke,fill}%
\end{pgfscope}%
\begin{pgfscope}%
\pgfpathrectangle{\pgfqpoint{0.050000in}{0.050000in}}{\pgfqpoint{4.900000in}{4.900000in}}%
\pgfusepath{clip}%
\pgfsetbuttcap%
\pgfsetroundjoin%
\definecolor{currentfill}{rgb}{1.000000,1.000000,1.000000}%
\pgfsetfillcolor{currentfill}%
\pgfsetfillopacity{0.900000}%
\pgfsetlinewidth{0.501875pt}%
\definecolor{currentstroke}{rgb}{0.200000,0.200000,0.200000}%
\pgfsetstrokecolor{currentstroke}%
\pgfsetdash{}{0pt}%
\pgfpathmoveto{\pgfqpoint{3.894372in}{1.633957in}}%
\pgfpathlineto{\pgfqpoint{3.937789in}{1.646250in}}%
\pgfpathlineto{\pgfqpoint{3.971387in}{1.626033in}}%
\pgfpathlineto{\pgfqpoint{3.927943in}{1.613804in}}%
\pgfpathlineto{\pgfqpoint{3.894372in}{1.633957in}}%
\pgfpathclose%
\pgfusepath{stroke,fill}%
\end{pgfscope}%
\begin{pgfscope}%
\pgfpathrectangle{\pgfqpoint{0.050000in}{0.050000in}}{\pgfqpoint{4.900000in}{4.900000in}}%
\pgfusepath{clip}%
\pgfsetbuttcap%
\pgfsetroundjoin%
\definecolor{currentfill}{rgb}{0.946482,0.946482,0.946482}%
\pgfsetfillcolor{currentfill}%
\pgfsetfillopacity{0.900000}%
\pgfsetlinewidth{0.501875pt}%
\definecolor{currentstroke}{rgb}{0.200000,0.200000,0.200000}%
\pgfsetstrokecolor{currentstroke}%
\pgfsetdash{}{0pt}%
\pgfpathmoveto{\pgfqpoint{2.811241in}{1.923032in}}%
\pgfpathlineto{\pgfqpoint{2.855676in}{1.895316in}}%
\pgfpathlineto{\pgfqpoint{2.887601in}{1.876533in}}%
\pgfpathlineto{\pgfqpoint{2.843154in}{1.903962in}}%
\pgfpathlineto{\pgfqpoint{2.811241in}{1.923032in}}%
\pgfpathclose%
\pgfusepath{stroke,fill}%
\end{pgfscope}%
\begin{pgfscope}%
\pgfpathrectangle{\pgfqpoint{0.050000in}{0.050000in}}{\pgfqpoint{4.900000in}{4.900000in}}%
\pgfusepath{clip}%
\pgfsetbuttcap%
\pgfsetroundjoin%
\definecolor{currentfill}{rgb}{0.878570,0.878570,0.878570}%
\pgfsetfillcolor{currentfill}%
\pgfsetfillopacity{0.900000}%
\pgfsetlinewidth{0.501875pt}%
\definecolor{currentstroke}{rgb}{0.200000,0.200000,0.200000}%
\pgfsetstrokecolor{currentstroke}%
\pgfsetdash{}{0pt}%
\pgfpathmoveto{\pgfqpoint{2.473969in}{2.166599in}}%
\pgfpathlineto{\pgfqpoint{2.519035in}{2.125928in}}%
\pgfpathlineto{\pgfqpoint{2.550434in}{2.106171in}}%
\pgfpathlineto{\pgfqpoint{2.505364in}{2.146621in}}%
\pgfpathlineto{\pgfqpoint{2.473969in}{2.166599in}}%
\pgfpathclose%
\pgfusepath{stroke,fill}%
\end{pgfscope}%
\begin{pgfscope}%
\pgfpathrectangle{\pgfqpoint{0.050000in}{0.050000in}}{\pgfqpoint{4.900000in}{4.900000in}}%
\pgfusepath{clip}%
\pgfsetbuttcap%
\pgfsetroundjoin%
\definecolor{currentfill}{rgb}{0.551634,0.551634,0.551634}%
\pgfsetfillcolor{currentfill}%
\pgfsetfillopacity{0.900000}%
\pgfsetlinewidth{0.501875pt}%
\definecolor{currentstroke}{rgb}{0.200000,0.200000,0.200000}%
\pgfsetstrokecolor{currentstroke}%
\pgfsetdash{}{0pt}%
\pgfpathmoveto{\pgfqpoint{1.506945in}{2.945903in}}%
\pgfpathlineto{\pgfqpoint{1.554031in}{2.901062in}}%
\pgfpathlineto{\pgfqpoint{1.583959in}{2.881803in}}%
\pgfpathlineto{\pgfqpoint{1.536865in}{2.926671in}}%
\pgfpathlineto{\pgfqpoint{1.506945in}{2.945903in}}%
\pgfpathclose%
\pgfusepath{stroke,fill}%
\end{pgfscope}%
\begin{pgfscope}%
\pgfpathrectangle{\pgfqpoint{0.050000in}{0.050000in}}{\pgfqpoint{4.900000in}{4.900000in}}%
\pgfusepath{clip}%
\pgfsetbuttcap%
\pgfsetroundjoin%
\definecolor{currentfill}{rgb}{0.988927,0.988927,0.988927}%
\pgfsetfillcolor{currentfill}%
\pgfsetfillopacity{0.900000}%
\pgfsetlinewidth{0.501875pt}%
\definecolor{currentstroke}{rgb}{0.200000,0.200000,0.200000}%
\pgfsetstrokecolor{currentstroke}%
\pgfsetdash{}{0pt}%
\pgfpathmoveto{\pgfqpoint{3.346226in}{1.691160in}}%
\pgfpathlineto{\pgfqpoint{3.390097in}{1.692110in}}%
\pgfpathlineto{\pgfqpoint{3.422869in}{1.673037in}}%
\pgfpathlineto{\pgfqpoint{3.378974in}{1.672186in}}%
\pgfpathlineto{\pgfqpoint{3.346226in}{1.691160in}}%
\pgfpathclose%
\pgfusepath{stroke,fill}%
\end{pgfscope}%
\begin{pgfscope}%
\pgfpathrectangle{\pgfqpoint{0.050000in}{0.050000in}}{\pgfqpoint{4.900000in}{4.900000in}}%
\pgfusepath{clip}%
\pgfsetbuttcap%
\pgfsetroundjoin%
\definecolor{currentfill}{rgb}{0.808781,0.808781,0.808781}%
\pgfsetfillcolor{currentfill}%
\pgfsetfillopacity{0.900000}%
\pgfsetlinewidth{0.501875pt}%
\definecolor{currentstroke}{rgb}{0.200000,0.200000,0.200000}%
\pgfsetstrokecolor{currentstroke}%
\pgfsetdash{}{0pt}%
\pgfpathmoveto{\pgfqpoint{2.213071in}{2.375273in}}%
\pgfpathlineto{\pgfqpoint{2.258687in}{2.331781in}}%
\pgfpathlineto{\pgfqpoint{2.289691in}{2.311737in}}%
\pgfpathlineto{\pgfqpoint{2.244071in}{2.355208in}}%
\pgfpathlineto{\pgfqpoint{2.213071in}{2.375273in}}%
\pgfpathclose%
\pgfusepath{stroke,fill}%
\end{pgfscope}%
\begin{pgfscope}%
\pgfpathrectangle{\pgfqpoint{0.050000in}{0.050000in}}{\pgfqpoint{4.900000in}{4.900000in}}%
\pgfusepath{clip}%
\pgfsetbuttcap%
\pgfsetroundjoin%
\definecolor{currentfill}{rgb}{0.998155,0.998155,0.998155}%
\pgfsetfillcolor{currentfill}%
\pgfsetfillopacity{0.900000}%
\pgfsetlinewidth{0.501875pt}%
\definecolor{currentstroke}{rgb}{0.200000,0.200000,0.200000}%
\pgfsetstrokecolor{currentstroke}%
\pgfsetdash{}{0pt}%
\pgfpathmoveto{\pgfqpoint{3.696960in}{1.640674in}}%
\pgfpathlineto{\pgfqpoint{3.740567in}{1.650940in}}%
\pgfpathlineto{\pgfqpoint{3.773876in}{1.631060in}}%
\pgfpathlineto{\pgfqpoint{3.730243in}{1.620888in}}%
\pgfpathlineto{\pgfqpoint{3.696960in}{1.640674in}}%
\pgfpathclose%
\pgfusepath{stroke,fill}%
\end{pgfscope}%
\begin{pgfscope}%
\pgfpathrectangle{\pgfqpoint{0.050000in}{0.050000in}}{\pgfqpoint{4.900000in}{4.900000in}}%
\pgfusepath{clip}%
\pgfsetbuttcap%
\pgfsetroundjoin%
\definecolor{currentfill}{rgb}{0.935163,0.935163,0.935163}%
\pgfsetfillcolor{currentfill}%
\pgfsetfillopacity{0.900000}%
\pgfsetlinewidth{0.501875pt}%
\definecolor{currentstroke}{rgb}{0.200000,0.200000,0.200000}%
\pgfsetstrokecolor{currentstroke}%
\pgfsetdash{}{0pt}%
\pgfpathmoveto{\pgfqpoint{2.734855in}{1.973874in}}%
\pgfpathlineto{\pgfqpoint{2.779427in}{1.942037in}}%
\pgfpathlineto{\pgfqpoint{2.811241in}{1.923032in}}%
\pgfpathlineto{\pgfqpoint{2.766658in}{1.954552in}}%
\pgfpathlineto{\pgfqpoint{2.734855in}{1.973874in}}%
\pgfpathclose%
\pgfusepath{stroke,fill}%
\end{pgfscope}%
\begin{pgfscope}%
\pgfpathrectangle{\pgfqpoint{0.050000in}{0.050000in}}{\pgfqpoint{4.900000in}{4.900000in}}%
\pgfusepath{clip}%
\pgfsetbuttcap%
\pgfsetroundjoin%
\definecolor{currentfill}{rgb}{0.724983,0.724983,0.724983}%
\pgfsetfillcolor{currentfill}%
\pgfsetfillopacity{0.900000}%
\pgfsetlinewidth{0.501875pt}%
\definecolor{currentstroke}{rgb}{0.200000,0.200000,0.200000}%
\pgfsetstrokecolor{currentstroke}%
\pgfsetdash{}{0pt}%
\pgfpathmoveto{\pgfqpoint{1.952083in}{2.586739in}}%
\pgfpathlineto{\pgfqpoint{1.998251in}{2.542637in}}%
\pgfpathlineto{\pgfqpoint{2.028864in}{2.522829in}}%
\pgfpathlineto{\pgfqpoint{1.982689in}{2.566955in}}%
\pgfpathlineto{\pgfqpoint{1.952083in}{2.586739in}}%
\pgfpathclose%
\pgfusepath{stroke,fill}%
\end{pgfscope}%
\begin{pgfscope}%
\pgfpathrectangle{\pgfqpoint{0.050000in}{0.050000in}}{\pgfqpoint{4.900000in}{4.900000in}}%
\pgfusepath{clip}%
\pgfsetbuttcap%
\pgfsetroundjoin%
\definecolor{currentfill}{rgb}{1.000000,1.000000,1.000000}%
\pgfsetfillcolor{currentfill}%
\pgfsetfillopacity{0.900000}%
\pgfsetlinewidth{0.501875pt}%
\definecolor{currentstroke}{rgb}{0.200000,0.200000,0.200000}%
\pgfsetstrokecolor{currentstroke}%
\pgfsetdash{}{0pt}%
\pgfpathmoveto{\pgfqpoint{3.971387in}{1.626033in}}%
\pgfpathlineto{\pgfqpoint{4.014743in}{1.638594in}}%
\pgfpathlineto{\pgfqpoint{4.048471in}{1.618259in}}%
\pgfpathlineto{\pgfqpoint{4.005088in}{1.605756in}}%
\pgfpathlineto{\pgfqpoint{3.971387in}{1.626033in}}%
\pgfpathclose%
\pgfusepath{stroke,fill}%
\end{pgfscope}%
\begin{pgfscope}%
\pgfpathrectangle{\pgfqpoint{0.050000in}{0.050000in}}{\pgfqpoint{4.900000in}{4.900000in}}%
\pgfusepath{clip}%
\pgfsetbuttcap%
\pgfsetroundjoin%
\definecolor{currentfill}{rgb}{0.992618,0.992618,0.992618}%
\pgfsetfillcolor{currentfill}%
\pgfsetfillopacity{0.900000}%
\pgfsetlinewidth{0.501875pt}%
\definecolor{currentstroke}{rgb}{0.200000,0.200000,0.200000}%
\pgfsetstrokecolor{currentstroke}%
\pgfsetdash{}{0pt}%
\pgfpathmoveto{\pgfqpoint{3.422869in}{1.673037in}}%
\pgfpathlineto{\pgfqpoint{3.466692in}{1.676618in}}%
\pgfpathlineto{\pgfqpoint{3.499590in}{1.657358in}}%
\pgfpathlineto{\pgfqpoint{3.455742in}{1.653889in}}%
\pgfpathlineto{\pgfqpoint{3.422869in}{1.673037in}}%
\pgfpathclose%
\pgfusepath{stroke,fill}%
\end{pgfscope}%
\begin{pgfscope}%
\pgfpathrectangle{\pgfqpoint{0.050000in}{0.050000in}}{\pgfqpoint{4.900000in}{4.900000in}}%
\pgfusepath{clip}%
\pgfsetbuttcap%
\pgfsetroundjoin%
\definecolor{currentfill}{rgb}{0.619423,0.619423,0.619423}%
\pgfsetfillcolor{currentfill}%
\pgfsetfillopacity{0.900000}%
\pgfsetlinewidth{0.501875pt}%
\definecolor{currentstroke}{rgb}{0.200000,0.200000,0.200000}%
\pgfsetstrokecolor{currentstroke}%
\pgfsetdash{}{0pt}%
\pgfpathmoveto{\pgfqpoint{1.691007in}{2.798370in}}%
\pgfpathlineto{\pgfqpoint{1.737727in}{2.753820in}}%
\pgfpathlineto{\pgfqpoint{1.767948in}{2.734329in}}%
\pgfpathlineto{\pgfqpoint{1.721219in}{2.778907in}}%
\pgfpathlineto{\pgfqpoint{1.691007in}{2.798370in}}%
\pgfpathclose%
\pgfusepath{stroke,fill}%
\end{pgfscope}%
\begin{pgfscope}%
\pgfpathrectangle{\pgfqpoint{0.050000in}{0.050000in}}{\pgfqpoint{4.900000in}{4.900000in}}%
\pgfusepath{clip}%
\pgfsetbuttcap%
\pgfsetroundjoin%
\definecolor{currentfill}{rgb}{0.858762,0.858762,0.858762}%
\pgfsetfillcolor{currentfill}%
\pgfsetfillopacity{0.900000}%
\pgfsetlinewidth{0.501875pt}%
\definecolor{currentstroke}{rgb}{0.200000,0.200000,0.200000}%
\pgfsetstrokecolor{currentstroke}%
\pgfsetdash{}{0pt}%
\pgfpathmoveto{\pgfqpoint{2.397423in}{2.228631in}}%
\pgfpathlineto{\pgfqpoint{2.442671in}{2.186545in}}%
\pgfpathlineto{\pgfqpoint{2.473969in}{2.166599in}}%
\pgfpathlineto{\pgfqpoint{2.428717in}{2.208537in}}%
\pgfpathlineto{\pgfqpoint{2.397423in}{2.228631in}}%
\pgfpathclose%
\pgfusepath{stroke,fill}%
\end{pgfscope}%
\begin{pgfscope}%
\pgfpathrectangle{\pgfqpoint{0.050000in}{0.050000in}}{\pgfqpoint{4.900000in}{4.900000in}}%
\pgfusepath{clip}%
\pgfsetbuttcap%
\pgfsetroundjoin%
\definecolor{currentfill}{rgb}{0.998155,0.998155,0.998155}%
\pgfsetfillcolor{currentfill}%
\pgfsetfillopacity{0.900000}%
\pgfsetlinewidth{0.501875pt}%
\definecolor{currentstroke}{rgb}{0.200000,0.200000,0.200000}%
\pgfsetstrokecolor{currentstroke}%
\pgfsetdash{}{0pt}%
\pgfpathmoveto{\pgfqpoint{3.773876in}{1.631060in}}%
\pgfpathlineto{\pgfqpoint{3.817432in}{1.642197in}}%
\pgfpathlineto{\pgfqpoint{3.850872in}{1.622175in}}%
\pgfpathlineto{\pgfqpoint{3.807290in}{1.611122in}}%
\pgfpathlineto{\pgfqpoint{3.773876in}{1.631060in}}%
\pgfpathclose%
\pgfusepath{stroke,fill}%
\end{pgfscope}%
\begin{pgfscope}%
\pgfpathrectangle{\pgfqpoint{0.050000in}{0.050000in}}{\pgfqpoint{4.900000in}{4.900000in}}%
\pgfusepath{clip}%
\pgfsetbuttcap%
\pgfsetroundjoin%
\definecolor{currentfill}{rgb}{0.517186,0.517186,0.517186}%
\pgfsetfillcolor{currentfill}%
\pgfsetfillopacity{0.900000}%
\pgfsetlinewidth{0.501875pt}%
\definecolor{currentstroke}{rgb}{0.200000,0.200000,0.200000}%
\pgfsetstrokecolor{currentstroke}%
\pgfsetdash{}{0pt}%
\pgfpathmoveto{\pgfqpoint{1.429845in}{3.010075in}}%
\pgfpathlineto{\pgfqpoint{1.477117in}{2.965075in}}%
\pgfpathlineto{\pgfqpoint{1.506945in}{2.945903in}}%
\pgfpathlineto{\pgfqpoint{1.459663in}{2.990931in}}%
\pgfpathlineto{\pgfqpoint{1.429845in}{3.010075in}}%
\pgfpathclose%
\pgfusepath{stroke,fill}%
\end{pgfscope}%
\begin{pgfscope}%
\pgfpathrectangle{\pgfqpoint{0.050000in}{0.050000in}}{\pgfqpoint{4.900000in}{4.900000in}}%
\pgfusepath{clip}%
\pgfsetbuttcap%
\pgfsetroundjoin%
\definecolor{currentfill}{rgb}{0.918185,0.918185,0.918185}%
\pgfsetfillcolor{currentfill}%
\pgfsetfillopacity{0.900000}%
\pgfsetlinewidth{0.501875pt}%
\definecolor{currentstroke}{rgb}{0.200000,0.200000,0.200000}%
\pgfsetstrokecolor{currentstroke}%
\pgfsetdash{}{0pt}%
\pgfpathmoveto{\pgfqpoint{2.658423in}{2.028548in}}%
\pgfpathlineto{\pgfqpoint{2.703150in}{1.993144in}}%
\pgfpathlineto{\pgfqpoint{2.734855in}{1.973874in}}%
\pgfpathlineto{\pgfqpoint{2.690120in}{2.008964in}}%
\pgfpathlineto{\pgfqpoint{2.658423in}{2.028548in}}%
\pgfpathclose%
\pgfusepath{stroke,fill}%
\end{pgfscope}%
\begin{pgfscope}%
\pgfpathrectangle{\pgfqpoint{0.050000in}{0.050000in}}{\pgfqpoint{4.900000in}{4.900000in}}%
\pgfusepath{clip}%
\pgfsetbuttcap%
\pgfsetroundjoin%
\definecolor{currentfill}{rgb}{0.994464,0.994464,0.994464}%
\pgfsetfillcolor{currentfill}%
\pgfsetfillopacity{0.900000}%
\pgfsetlinewidth{0.501875pt}%
\definecolor{currentstroke}{rgb}{0.200000,0.200000,0.200000}%
\pgfsetstrokecolor{currentstroke}%
\pgfsetdash{}{0pt}%
\pgfpathmoveto{\pgfqpoint{3.499590in}{1.657358in}}%
\pgfpathlineto{\pgfqpoint{3.543367in}{1.663135in}}%
\pgfpathlineto{\pgfqpoint{3.576392in}{1.643691in}}%
\pgfpathlineto{\pgfqpoint{3.532591in}{1.638030in}}%
\pgfpathlineto{\pgfqpoint{3.499590in}{1.657358in}}%
\pgfpathclose%
\pgfusepath{stroke,fill}%
\end{pgfscope}%
\begin{pgfscope}%
\pgfpathrectangle{\pgfqpoint{0.050000in}{0.050000in}}{\pgfqpoint{4.900000in}{4.900000in}}%
\pgfusepath{clip}%
\pgfsetbuttcap%
\pgfsetroundjoin%
\definecolor{currentfill}{rgb}{0.781223,0.781223,0.781223}%
\pgfsetfillcolor{currentfill}%
\pgfsetfillopacity{0.900000}%
\pgfsetlinewidth{0.501875pt}%
\definecolor{currentstroke}{rgb}{0.200000,0.200000,0.200000}%
\pgfsetstrokecolor{currentstroke}%
\pgfsetdash{}{0pt}%
\pgfpathmoveto{\pgfqpoint{2.136364in}{2.439053in}}%
\pgfpathlineto{\pgfqpoint{2.182166in}{2.395286in}}%
\pgfpathlineto{\pgfqpoint{2.213071in}{2.375273in}}%
\pgfpathlineto{\pgfqpoint{2.167264in}{2.419047in}}%
\pgfpathlineto{\pgfqpoint{2.136364in}{2.439053in}}%
\pgfpathclose%
\pgfusepath{stroke,fill}%
\end{pgfscope}%
\begin{pgfscope}%
\pgfpathrectangle{\pgfqpoint{0.050000in}{0.050000in}}{\pgfqpoint{4.900000in}{4.900000in}}%
\pgfusepath{clip}%
\pgfsetbuttcap%
\pgfsetroundjoin%
\definecolor{currentfill}{rgb}{1.000000,1.000000,1.000000}%
\pgfsetfillcolor{currentfill}%
\pgfsetfillopacity{0.900000}%
\pgfsetlinewidth{0.501875pt}%
\definecolor{currentstroke}{rgb}{0.200000,0.200000,0.200000}%
\pgfsetstrokecolor{currentstroke}%
\pgfsetdash{}{0pt}%
\pgfpathmoveto{\pgfqpoint{4.048471in}{1.618259in}}%
\pgfpathlineto{\pgfqpoint{4.091763in}{1.630932in}}%
\pgfpathlineto{\pgfqpoint{4.125622in}{1.610481in}}%
\pgfpathlineto{\pgfqpoint{4.082303in}{1.597862in}}%
\pgfpathlineto{\pgfqpoint{4.048471in}{1.618259in}}%
\pgfpathclose%
\pgfusepath{stroke,fill}%
\end{pgfscope}%
\begin{pgfscope}%
\pgfpathrectangle{\pgfqpoint{0.050000in}{0.050000in}}{\pgfqpoint{4.900000in}{4.900000in}}%
\pgfusepath{clip}%
\pgfsetbuttcap%
\pgfsetroundjoin%
\definecolor{currentfill}{rgb}{0.691396,0.691396,0.691396}%
\pgfsetfillcolor{currentfill}%
\pgfsetfillopacity{0.900000}%
\pgfsetlinewidth{0.501875pt}%
\definecolor{currentstroke}{rgb}{0.200000,0.200000,0.200000}%
\pgfsetstrokecolor{currentstroke}%
\pgfsetdash{}{0pt}%
\pgfpathmoveto{\pgfqpoint{1.875215in}{2.650722in}}%
\pgfpathlineto{\pgfqpoint{1.921570in}{2.606462in}}%
\pgfpathlineto{\pgfqpoint{1.952083in}{2.586739in}}%
\pgfpathlineto{\pgfqpoint{1.905721in}{2.631025in}}%
\pgfpathlineto{\pgfqpoint{1.875215in}{2.650722in}}%
\pgfpathclose%
\pgfusepath{stroke,fill}%
\end{pgfscope}%
\begin{pgfscope}%
\pgfpathrectangle{\pgfqpoint{0.050000in}{0.050000in}}{\pgfqpoint{4.900000in}{4.900000in}}%
\pgfusepath{clip}%
\pgfsetbuttcap%
\pgfsetroundjoin%
\definecolor{currentfill}{rgb}{0.996309,0.996309,0.996309}%
\pgfsetfillcolor{currentfill}%
\pgfsetfillopacity{0.900000}%
\pgfsetlinewidth{0.501875pt}%
\definecolor{currentstroke}{rgb}{0.200000,0.200000,0.200000}%
\pgfsetstrokecolor{currentstroke}%
\pgfsetdash{}{0pt}%
\pgfpathmoveto{\pgfqpoint{3.576392in}{1.643691in}}%
\pgfpathlineto{\pgfqpoint{3.620123in}{1.651268in}}%
\pgfpathlineto{\pgfqpoint{3.653277in}{1.631647in}}%
\pgfpathlineto{\pgfqpoint{3.609521in}{1.624182in}}%
\pgfpathlineto{\pgfqpoint{3.576392in}{1.643691in}}%
\pgfpathclose%
\pgfusepath{stroke,fill}%
\end{pgfscope}%
\begin{pgfscope}%
\pgfpathrectangle{\pgfqpoint{0.050000in}{0.050000in}}{\pgfqpoint{4.900000in}{4.900000in}}%
\pgfusepath{clip}%
\pgfsetbuttcap%
\pgfsetroundjoin%
\definecolor{currentfill}{rgb}{0.581776,0.581776,0.581776}%
\pgfsetfillcolor{currentfill}%
\pgfsetfillopacity{0.900000}%
\pgfsetlinewidth{0.501875pt}%
\definecolor{currentstroke}{rgb}{0.200000,0.200000,0.200000}%
\pgfsetstrokecolor{currentstroke}%
\pgfsetdash{}{0pt}%
\pgfpathmoveto{\pgfqpoint{1.613979in}{2.862485in}}%
\pgfpathlineto{\pgfqpoint{1.660887in}{2.817775in}}%
\pgfpathlineto{\pgfqpoint{1.691007in}{2.798370in}}%
\pgfpathlineto{\pgfqpoint{1.644091in}{2.843108in}}%
\pgfpathlineto{\pgfqpoint{1.613979in}{2.862485in}}%
\pgfpathclose%
\pgfusepath{stroke,fill}%
\end{pgfscope}%
\begin{pgfscope}%
\pgfpathrectangle{\pgfqpoint{0.050000in}{0.050000in}}{\pgfqpoint{4.900000in}{4.900000in}}%
\pgfusepath{clip}%
\pgfsetbuttcap%
\pgfsetroundjoin%
\definecolor{currentfill}{rgb}{1.000000,1.000000,1.000000}%
\pgfsetfillcolor{currentfill}%
\pgfsetfillopacity{0.900000}%
\pgfsetlinewidth{0.501875pt}%
\definecolor{currentstroke}{rgb}{0.200000,0.200000,0.200000}%
\pgfsetstrokecolor{currentstroke}%
\pgfsetdash{}{0pt}%
\pgfpathmoveto{\pgfqpoint{3.850872in}{1.622175in}}%
\pgfpathlineto{\pgfqpoint{3.894372in}{1.633957in}}%
\pgfpathlineto{\pgfqpoint{3.927943in}{1.613804in}}%
\pgfpathlineto{\pgfqpoint{3.884415in}{1.602096in}}%
\pgfpathlineto{\pgfqpoint{3.850872in}{1.622175in}}%
\pgfpathclose%
\pgfusepath{stroke,fill}%
\end{pgfscope}%
\begin{pgfscope}%
\pgfpathrectangle{\pgfqpoint{0.050000in}{0.050000in}}{\pgfqpoint{4.900000in}{4.900000in}}%
\pgfusepath{clip}%
\pgfsetbuttcap%
\pgfsetroundjoin%
\definecolor{currentfill}{rgb}{0.898378,0.898378,0.898378}%
\pgfsetfillcolor{currentfill}%
\pgfsetfillopacity{0.900000}%
\pgfsetlinewidth{0.501875pt}%
\definecolor{currentstroke}{rgb}{0.200000,0.200000,0.200000}%
\pgfsetstrokecolor{currentstroke}%
\pgfsetdash{}{0pt}%
\pgfpathmoveto{\pgfqpoint{2.581930in}{2.086381in}}%
\pgfpathlineto{\pgfqpoint{2.626823in}{2.048092in}}%
\pgfpathlineto{\pgfqpoint{2.658423in}{2.028548in}}%
\pgfpathlineto{\pgfqpoint{2.613523in}{2.066556in}}%
\pgfpathlineto{\pgfqpoint{2.581930in}{2.086381in}}%
\pgfpathclose%
\pgfusepath{stroke,fill}%
\end{pgfscope}%
\begin{pgfscope}%
\pgfpathrectangle{\pgfqpoint{0.050000in}{0.050000in}}{\pgfqpoint{4.900000in}{4.900000in}}%
\pgfusepath{clip}%
\pgfsetbuttcap%
\pgfsetroundjoin%
\definecolor{currentfill}{rgb}{0.832895,0.832895,0.832895}%
\pgfsetfillcolor{currentfill}%
\pgfsetfillopacity{0.900000}%
\pgfsetlinewidth{0.501875pt}%
\definecolor{currentstroke}{rgb}{0.200000,0.200000,0.200000}%
\pgfsetstrokecolor{currentstroke}%
\pgfsetdash{}{0pt}%
\pgfpathmoveto{\pgfqpoint{2.320790in}{2.291647in}}%
\pgfpathlineto{\pgfqpoint{2.366224in}{2.248686in}}%
\pgfpathlineto{\pgfqpoint{2.397423in}{2.228631in}}%
\pgfpathlineto{\pgfqpoint{2.351985in}{2.271512in}}%
\pgfpathlineto{\pgfqpoint{2.320790in}{2.291647in}}%
\pgfpathclose%
\pgfusepath{stroke,fill}%
\end{pgfscope}%
\begin{pgfscope}%
\pgfpathrectangle{\pgfqpoint{0.050000in}{0.050000in}}{\pgfqpoint{4.900000in}{4.900000in}}%
\pgfusepath{clip}%
\pgfsetbuttcap%
\pgfsetroundjoin%
\definecolor{currentfill}{rgb}{0.970473,0.970473,0.970473}%
\pgfsetfillcolor{currentfill}%
\pgfsetfillopacity{0.900000}%
\pgfsetlinewidth{0.501875pt}%
\definecolor{currentstroke}{rgb}{0.200000,0.200000,0.200000}%
\pgfsetstrokecolor{currentstroke}%
\pgfsetdash{}{0pt}%
\pgfpathmoveto{\pgfqpoint{3.072578in}{1.778880in}}%
\pgfpathlineto{\pgfqpoint{3.116715in}{1.764941in}}%
\pgfpathlineto{\pgfqpoint{3.149094in}{1.746299in}}%
\pgfpathlineto{\pgfqpoint{3.104937in}{1.760149in}}%
\pgfpathlineto{\pgfqpoint{3.072578in}{1.778880in}}%
\pgfpathclose%
\pgfusepath{stroke,fill}%
\end{pgfscope}%
\begin{pgfscope}%
\pgfpathrectangle{\pgfqpoint{0.050000in}{0.050000in}}{\pgfqpoint{4.900000in}{4.900000in}}%
\pgfusepath{clip}%
\pgfsetbuttcap%
\pgfsetroundjoin%
\definecolor{currentfill}{rgb}{0.976009,0.976009,0.976009}%
\pgfsetfillcolor{currentfill}%
\pgfsetfillopacity{0.900000}%
\pgfsetlinewidth{0.501875pt}%
\definecolor{currentstroke}{rgb}{0.200000,0.200000,0.200000}%
\pgfsetstrokecolor{currentstroke}%
\pgfsetdash{}{0pt}%
\pgfpathmoveto{\pgfqpoint{3.149094in}{1.746299in}}%
\pgfpathlineto{\pgfqpoint{3.193157in}{1.736636in}}%
\pgfpathlineto{\pgfqpoint{3.225657in}{1.717930in}}%
\pgfpathlineto{\pgfqpoint{3.181573in}{1.727567in}}%
\pgfpathlineto{\pgfqpoint{3.149094in}{1.746299in}}%
\pgfpathclose%
\pgfusepath{stroke,fill}%
\end{pgfscope}%
\begin{pgfscope}%
\pgfpathrectangle{\pgfqpoint{0.050000in}{0.050000in}}{\pgfqpoint{4.900000in}{4.900000in}}%
\pgfusepath{clip}%
\pgfsetbuttcap%
\pgfsetroundjoin%
\definecolor{currentfill}{rgb}{0.482737,0.482737,0.482737}%
\pgfsetfillcolor{currentfill}%
\pgfsetfillopacity{0.900000}%
\pgfsetlinewidth{0.501875pt}%
\definecolor{currentstroke}{rgb}{0.200000,0.200000,0.200000}%
\pgfsetstrokecolor{currentstroke}%
\pgfsetdash{}{0pt}%
\pgfpathmoveto{\pgfqpoint{1.352657in}{3.074320in}}%
\pgfpathlineto{\pgfqpoint{1.400117in}{3.029161in}}%
\pgfpathlineto{\pgfqpoint{1.429845in}{3.010075in}}%
\pgfpathlineto{\pgfqpoint{1.382374in}{3.055263in}}%
\pgfpathlineto{\pgfqpoint{1.352657in}{3.074320in}}%
\pgfpathclose%
\pgfusepath{stroke,fill}%
\end{pgfscope}%
\begin{pgfscope}%
\pgfpathrectangle{\pgfqpoint{0.050000in}{0.050000in}}{\pgfqpoint{4.900000in}{4.900000in}}%
\pgfusepath{clip}%
\pgfsetbuttcap%
\pgfsetroundjoin%
\definecolor{currentfill}{rgb}{0.963091,0.963091,0.963091}%
\pgfsetfillcolor{currentfill}%
\pgfsetfillopacity{0.900000}%
\pgfsetlinewidth{0.501875pt}%
\definecolor{currentstroke}{rgb}{0.200000,0.200000,0.200000}%
\pgfsetstrokecolor{currentstroke}%
\pgfsetdash{}{0pt}%
\pgfpathmoveto{\pgfqpoint{2.996095in}{1.815959in}}%
\pgfpathlineto{\pgfqpoint{3.040320in}{1.797520in}}%
\pgfpathlineto{\pgfqpoint{3.072578in}{1.778880in}}%
\pgfpathlineto{\pgfqpoint{3.028336in}{1.797161in}}%
\pgfpathlineto{\pgfqpoint{2.996095in}{1.815959in}}%
\pgfpathclose%
\pgfusepath{stroke,fill}%
\end{pgfscope}%
\begin{pgfscope}%
\pgfpathrectangle{\pgfqpoint{0.050000in}{0.050000in}}{\pgfqpoint{4.900000in}{4.900000in}}%
\pgfusepath{clip}%
\pgfsetbuttcap%
\pgfsetroundjoin%
\definecolor{currentfill}{rgb}{0.757109,0.757109,0.757109}%
\pgfsetfillcolor{currentfill}%
\pgfsetfillopacity{0.900000}%
\pgfsetlinewidth{0.501875pt}%
\definecolor{currentstroke}{rgb}{0.200000,0.200000,0.200000}%
\pgfsetstrokecolor{currentstroke}%
\pgfsetdash{}{0pt}%
\pgfpathmoveto{\pgfqpoint{2.059570in}{2.502961in}}%
\pgfpathlineto{\pgfqpoint{2.105558in}{2.459001in}}%
\pgfpathlineto{\pgfqpoint{2.136364in}{2.439053in}}%
\pgfpathlineto{\pgfqpoint{2.090370in}{2.483033in}}%
\pgfpathlineto{\pgfqpoint{2.059570in}{2.502961in}}%
\pgfpathclose%
\pgfusepath{stroke,fill}%
\end{pgfscope}%
\begin{pgfscope}%
\pgfpathrectangle{\pgfqpoint{0.050000in}{0.050000in}}{\pgfqpoint{4.900000in}{4.900000in}}%
\pgfusepath{clip}%
\pgfsetbuttcap%
\pgfsetroundjoin%
\definecolor{currentfill}{rgb}{0.981546,0.981546,0.981546}%
\pgfsetfillcolor{currentfill}%
\pgfsetfillopacity{0.900000}%
\pgfsetlinewidth{0.501875pt}%
\definecolor{currentstroke}{rgb}{0.200000,0.200000,0.200000}%
\pgfsetstrokecolor{currentstroke}%
\pgfsetdash{}{0pt}%
\pgfpathmoveto{\pgfqpoint{3.225657in}{1.717930in}}%
\pgfpathlineto{\pgfqpoint{3.269658in}{1.712199in}}%
\pgfpathlineto{\pgfqpoint{3.302281in}{1.693376in}}%
\pgfpathlineto{\pgfqpoint{3.258259in}{1.699134in}}%
\pgfpathlineto{\pgfqpoint{3.225657in}{1.717930in}}%
\pgfpathclose%
\pgfusepath{stroke,fill}%
\end{pgfscope}%
\begin{pgfscope}%
\pgfpathrectangle{\pgfqpoint{0.050000in}{0.050000in}}{\pgfqpoint{4.900000in}{4.900000in}}%
\pgfusepath{clip}%
\pgfsetbuttcap%
\pgfsetroundjoin%
\definecolor{currentfill}{rgb}{0.955709,0.955709,0.955709}%
\pgfsetfillcolor{currentfill}%
\pgfsetfillopacity{0.900000}%
\pgfsetlinewidth{0.501875pt}%
\definecolor{currentstroke}{rgb}{0.200000,0.200000,0.200000}%
\pgfsetstrokecolor{currentstroke}%
\pgfsetdash{}{0pt}%
\pgfpathmoveto{\pgfqpoint{2.919627in}{1.857673in}}%
\pgfpathlineto{\pgfqpoint{2.963954in}{1.834672in}}%
\pgfpathlineto{\pgfqpoint{2.996095in}{1.815959in}}%
\pgfpathlineto{\pgfqpoint{2.951752in}{1.838737in}}%
\pgfpathlineto{\pgfqpoint{2.919627in}{1.857673in}}%
\pgfpathclose%
\pgfusepath{stroke,fill}%
\end{pgfscope}%
\begin{pgfscope}%
\pgfpathrectangle{\pgfqpoint{0.050000in}{0.050000in}}{\pgfqpoint{4.900000in}{4.900000in}}%
\pgfusepath{clip}%
\pgfsetbuttcap%
\pgfsetroundjoin%
\definecolor{currentfill}{rgb}{0.996309,0.996309,0.996309}%
\pgfsetfillcolor{currentfill}%
\pgfsetfillopacity{0.900000}%
\pgfsetlinewidth{0.501875pt}%
\definecolor{currentstroke}{rgb}{0.200000,0.200000,0.200000}%
\pgfsetstrokecolor{currentstroke}%
\pgfsetdash{}{0pt}%
\pgfpathmoveto{\pgfqpoint{3.653277in}{1.631647in}}%
\pgfpathlineto{\pgfqpoint{3.696960in}{1.640674in}}%
\pgfpathlineto{\pgfqpoint{3.730243in}{1.620888in}}%
\pgfpathlineto{\pgfqpoint{3.686534in}{1.611965in}}%
\pgfpathlineto{\pgfqpoint{3.653277in}{1.631647in}}%
\pgfpathclose%
\pgfusepath{stroke,fill}%
\end{pgfscope}%
\begin{pgfscope}%
\pgfpathrectangle{\pgfqpoint{0.050000in}{0.050000in}}{\pgfqpoint{4.900000in}{4.900000in}}%
\pgfusepath{clip}%
\pgfsetbuttcap%
\pgfsetroundjoin%
\definecolor{currentfill}{rgb}{0.653010,0.653010,0.653010}%
\pgfsetfillcolor{currentfill}%
\pgfsetfillopacity{0.900000}%
\pgfsetlinewidth{0.501875pt}%
\definecolor{currentstroke}{rgb}{0.200000,0.200000,0.200000}%
\pgfsetstrokecolor{currentstroke}%
\pgfsetdash{}{0pt}%
\pgfpathmoveto{\pgfqpoint{1.798261in}{2.714778in}}%
\pgfpathlineto{\pgfqpoint{1.844802in}{2.670359in}}%
\pgfpathlineto{\pgfqpoint{1.875215in}{2.650722in}}%
\pgfpathlineto{\pgfqpoint{1.828666in}{2.695168in}}%
\pgfpathlineto{\pgfqpoint{1.798261in}{2.714778in}}%
\pgfpathclose%
\pgfusepath{stroke,fill}%
\end{pgfscope}%
\begin{pgfscope}%
\pgfpathrectangle{\pgfqpoint{0.050000in}{0.050000in}}{\pgfqpoint{4.900000in}{4.900000in}}%
\pgfusepath{clip}%
\pgfsetbuttcap%
\pgfsetroundjoin%
\definecolor{currentfill}{rgb}{0.985236,0.985236,0.985236}%
\pgfsetfillcolor{currentfill}%
\pgfsetfillopacity{0.900000}%
\pgfsetlinewidth{0.501875pt}%
\definecolor{currentstroke}{rgb}{0.200000,0.200000,0.200000}%
\pgfsetstrokecolor{currentstroke}%
\pgfsetdash{}{0pt}%
\pgfpathmoveto{\pgfqpoint{3.302281in}{1.693376in}}%
\pgfpathlineto{\pgfqpoint{3.346226in}{1.691160in}}%
\pgfpathlineto{\pgfqpoint{3.378974in}{1.672186in}}%
\pgfpathlineto{\pgfqpoint{3.335006in}{1.674468in}}%
\pgfpathlineto{\pgfqpoint{3.302281in}{1.693376in}}%
\pgfpathclose%
\pgfusepath{stroke,fill}%
\end{pgfscope}%
\begin{pgfscope}%
\pgfpathrectangle{\pgfqpoint{0.050000in}{0.050000in}}{\pgfqpoint{4.900000in}{4.900000in}}%
\pgfusepath{clip}%
\pgfsetbuttcap%
\pgfsetroundjoin%
\definecolor{currentfill}{rgb}{1.000000,1.000000,1.000000}%
\pgfsetfillcolor{currentfill}%
\pgfsetfillopacity{0.900000}%
\pgfsetlinewidth{0.501875pt}%
\definecolor{currentstroke}{rgb}{0.200000,0.200000,0.200000}%
\pgfsetstrokecolor{currentstroke}%
\pgfsetdash{}{0pt}%
\pgfpathmoveto{\pgfqpoint{3.927943in}{1.613804in}}%
\pgfpathlineto{\pgfqpoint{3.971387in}{1.626033in}}%
\pgfpathlineto{\pgfqpoint{4.005088in}{1.605756in}}%
\pgfpathlineto{\pgfqpoint{3.961618in}{1.593592in}}%
\pgfpathlineto{\pgfqpoint{3.927943in}{1.613804in}}%
\pgfpathclose%
\pgfusepath{stroke,fill}%
\end{pgfscope}%
\begin{pgfscope}%
\pgfpathrectangle{\pgfqpoint{0.050000in}{0.050000in}}{\pgfqpoint{4.900000in}{4.900000in}}%
\pgfusepath{clip}%
\pgfsetbuttcap%
\pgfsetroundjoin%
\definecolor{currentfill}{rgb}{0.946482,0.946482,0.946482}%
\pgfsetfillcolor{currentfill}%
\pgfsetfillopacity{0.900000}%
\pgfsetlinewidth{0.501875pt}%
\definecolor{currentstroke}{rgb}{0.200000,0.200000,0.200000}%
\pgfsetstrokecolor{currentstroke}%
\pgfsetdash{}{0pt}%
\pgfpathmoveto{\pgfqpoint{2.843154in}{1.903962in}}%
\pgfpathlineto{\pgfqpoint{2.887601in}{1.876533in}}%
\pgfpathlineto{\pgfqpoint{2.919627in}{1.857673in}}%
\pgfpathlineto{\pgfqpoint{2.875167in}{1.884827in}}%
\pgfpathlineto{\pgfqpoint{2.843154in}{1.903962in}}%
\pgfpathclose%
\pgfusepath{stroke,fill}%
\end{pgfscope}%
\begin{pgfscope}%
\pgfpathrectangle{\pgfqpoint{0.050000in}{0.050000in}}{\pgfqpoint{4.900000in}{4.900000in}}%
\pgfusepath{clip}%
\pgfsetbuttcap%
\pgfsetroundjoin%
\definecolor{currentfill}{rgb}{0.878570,0.878570,0.878570}%
\pgfsetfillcolor{currentfill}%
\pgfsetfillopacity{0.900000}%
\pgfsetlinewidth{0.501875pt}%
\definecolor{currentstroke}{rgb}{0.200000,0.200000,0.200000}%
\pgfsetstrokecolor{currentstroke}%
\pgfsetdash{}{0pt}%
\pgfpathmoveto{\pgfqpoint{2.505364in}{2.146621in}}%
\pgfpathlineto{\pgfqpoint{2.550434in}{2.106171in}}%
\pgfpathlineto{\pgfqpoint{2.581930in}{2.086381in}}%
\pgfpathlineto{\pgfqpoint{2.536854in}{2.126608in}}%
\pgfpathlineto{\pgfqpoint{2.505364in}{2.146621in}}%
\pgfpathclose%
\pgfusepath{stroke,fill}%
\end{pgfscope}%
\begin{pgfscope}%
\pgfpathrectangle{\pgfqpoint{0.050000in}{0.050000in}}{\pgfqpoint{4.900000in}{4.900000in}}%
\pgfusepath{clip}%
\pgfsetbuttcap%
\pgfsetroundjoin%
\definecolor{currentfill}{rgb}{0.551634,0.551634,0.551634}%
\pgfsetfillcolor{currentfill}%
\pgfsetfillopacity{0.900000}%
\pgfsetlinewidth{0.501875pt}%
\definecolor{currentstroke}{rgb}{0.200000,0.200000,0.200000}%
\pgfsetstrokecolor{currentstroke}%
\pgfsetdash{}{0pt}%
\pgfpathmoveto{\pgfqpoint{1.536865in}{2.926671in}}%
\pgfpathlineto{\pgfqpoint{1.583959in}{2.881803in}}%
\pgfpathlineto{\pgfqpoint{1.613979in}{2.862485in}}%
\pgfpathlineto{\pgfqpoint{1.566875in}{2.907381in}}%
\pgfpathlineto{\pgfqpoint{1.536865in}{2.926671in}}%
\pgfpathclose%
\pgfusepath{stroke,fill}%
\end{pgfscope}%
\begin{pgfscope}%
\pgfpathrectangle{\pgfqpoint{0.050000in}{0.050000in}}{\pgfqpoint{4.900000in}{4.900000in}}%
\pgfusepath{clip}%
\pgfsetbuttcap%
\pgfsetroundjoin%
\definecolor{currentfill}{rgb}{0.988927,0.988927,0.988927}%
\pgfsetfillcolor{currentfill}%
\pgfsetfillopacity{0.900000}%
\pgfsetlinewidth{0.501875pt}%
\definecolor{currentstroke}{rgb}{0.200000,0.200000,0.200000}%
\pgfsetstrokecolor{currentstroke}%
\pgfsetdash{}{0pt}%
\pgfpathmoveto{\pgfqpoint{3.378974in}{1.672186in}}%
\pgfpathlineto{\pgfqpoint{3.422869in}{1.673037in}}%
\pgfpathlineto{\pgfqpoint{3.455742in}{1.653889in}}%
\pgfpathlineto{\pgfqpoint{3.411824in}{1.653131in}}%
\pgfpathlineto{\pgfqpoint{3.378974in}{1.672186in}}%
\pgfpathclose%
\pgfusepath{stroke,fill}%
\end{pgfscope}%
\begin{pgfscope}%
\pgfpathrectangle{\pgfqpoint{0.050000in}{0.050000in}}{\pgfqpoint{4.900000in}{4.900000in}}%
\pgfusepath{clip}%
\pgfsetbuttcap%
\pgfsetroundjoin%
\definecolor{currentfill}{rgb}{0.808781,0.808781,0.808781}%
\pgfsetfillcolor{currentfill}%
\pgfsetfillopacity{0.900000}%
\pgfsetlinewidth{0.501875pt}%
\definecolor{currentstroke}{rgb}{0.200000,0.200000,0.200000}%
\pgfsetstrokecolor{currentstroke}%
\pgfsetdash{}{0pt}%
\pgfpathmoveto{\pgfqpoint{2.244071in}{2.355208in}}%
\pgfpathlineto{\pgfqpoint{2.289691in}{2.311737in}}%
\pgfpathlineto{\pgfqpoint{2.320790in}{2.291647in}}%
\pgfpathlineto{\pgfqpoint{2.275165in}{2.335089in}}%
\pgfpathlineto{\pgfqpoint{2.244071in}{2.355208in}}%
\pgfpathclose%
\pgfusepath{stroke,fill}%
\end{pgfscope}%
\begin{pgfscope}%
\pgfpathrectangle{\pgfqpoint{0.050000in}{0.050000in}}{\pgfqpoint{4.900000in}{4.900000in}}%
\pgfusepath{clip}%
\pgfsetbuttcap%
\pgfsetroundjoin%
\definecolor{currentfill}{rgb}{0.452595,0.452595,0.452595}%
\pgfsetfillcolor{currentfill}%
\pgfsetfillopacity{0.900000}%
\pgfsetlinewidth{0.501875pt}%
\definecolor{currentstroke}{rgb}{0.200000,0.200000,0.200000}%
\pgfsetstrokecolor{currentstroke}%
\pgfsetdash{}{0pt}%
\pgfpathmoveto{\pgfqpoint{1.275382in}{3.138638in}}%
\pgfpathlineto{\pgfqpoint{1.323030in}{3.093319in}}%
\pgfpathlineto{\pgfqpoint{1.352657in}{3.074320in}}%
\pgfpathlineto{\pgfqpoint{1.304997in}{3.119668in}}%
\pgfpathlineto{\pgfqpoint{1.275382in}{3.138638in}}%
\pgfpathclose%
\pgfusepath{stroke,fill}%
\end{pgfscope}%
\begin{pgfscope}%
\pgfpathrectangle{\pgfqpoint{0.050000in}{0.050000in}}{\pgfqpoint{4.900000in}{4.900000in}}%
\pgfusepath{clip}%
\pgfsetbuttcap%
\pgfsetroundjoin%
\definecolor{currentfill}{rgb}{0.998155,0.998155,0.998155}%
\pgfsetfillcolor{currentfill}%
\pgfsetfillopacity{0.900000}%
\pgfsetlinewidth{0.501875pt}%
\definecolor{currentstroke}{rgb}{0.200000,0.200000,0.200000}%
\pgfsetstrokecolor{currentstroke}%
\pgfsetdash{}{0pt}%
\pgfpathmoveto{\pgfqpoint{3.730243in}{1.620888in}}%
\pgfpathlineto{\pgfqpoint{3.773876in}{1.631060in}}%
\pgfpathlineto{\pgfqpoint{3.807290in}{1.611122in}}%
\pgfpathlineto{\pgfqpoint{3.763630in}{1.601043in}}%
\pgfpathlineto{\pgfqpoint{3.730243in}{1.620888in}}%
\pgfpathclose%
\pgfusepath{stroke,fill}%
\end{pgfscope}%
\begin{pgfscope}%
\pgfpathrectangle{\pgfqpoint{0.050000in}{0.050000in}}{\pgfqpoint{4.900000in}{4.900000in}}%
\pgfusepath{clip}%
\pgfsetbuttcap%
\pgfsetroundjoin%
\definecolor{currentfill}{rgb}{0.724983,0.724983,0.724983}%
\pgfsetfillcolor{currentfill}%
\pgfsetfillopacity{0.900000}%
\pgfsetlinewidth{0.501875pt}%
\definecolor{currentstroke}{rgb}{0.200000,0.200000,0.200000}%
\pgfsetstrokecolor{currentstroke}%
\pgfsetdash{}{0pt}%
\pgfpathmoveto{\pgfqpoint{1.982689in}{2.566955in}}%
\pgfpathlineto{\pgfqpoint{2.028864in}{2.522829in}}%
\pgfpathlineto{\pgfqpoint{2.059570in}{2.502961in}}%
\pgfpathlineto{\pgfqpoint{2.013389in}{2.547112in}}%
\pgfpathlineto{\pgfqpoint{1.982689in}{2.566955in}}%
\pgfpathclose%
\pgfusepath{stroke,fill}%
\end{pgfscope}%
\begin{pgfscope}%
\pgfpathrectangle{\pgfqpoint{0.050000in}{0.050000in}}{\pgfqpoint{4.900000in}{4.900000in}}%
\pgfusepath{clip}%
\pgfsetbuttcap%
\pgfsetroundjoin%
\definecolor{currentfill}{rgb}{0.935163,0.935163,0.935163}%
\pgfsetfillcolor{currentfill}%
\pgfsetfillopacity{0.900000}%
\pgfsetlinewidth{0.501875pt}%
\definecolor{currentstroke}{rgb}{0.200000,0.200000,0.200000}%
\pgfsetstrokecolor{currentstroke}%
\pgfsetdash{}{0pt}%
\pgfpathmoveto{\pgfqpoint{2.766658in}{1.954552in}}%
\pgfpathlineto{\pgfqpoint{2.811241in}{1.923032in}}%
\pgfpathlineto{\pgfqpoint{2.843154in}{1.903962in}}%
\pgfpathlineto{\pgfqpoint{2.798561in}{1.935177in}}%
\pgfpathlineto{\pgfqpoint{2.766658in}{1.954552in}}%
\pgfpathclose%
\pgfusepath{stroke,fill}%
\end{pgfscope}%
\begin{pgfscope}%
\pgfpathrectangle{\pgfqpoint{0.050000in}{0.050000in}}{\pgfqpoint{4.900000in}{4.900000in}}%
\pgfusepath{clip}%
\pgfsetbuttcap%
\pgfsetroundjoin%
\definecolor{currentfill}{rgb}{1.000000,1.000000,1.000000}%
\pgfsetfillcolor{currentfill}%
\pgfsetfillopacity{0.900000}%
\pgfsetlinewidth{0.501875pt}%
\definecolor{currentstroke}{rgb}{0.200000,0.200000,0.200000}%
\pgfsetstrokecolor{currentstroke}%
\pgfsetdash{}{0pt}%
\pgfpathmoveto{\pgfqpoint{4.005088in}{1.605756in}}%
\pgfpathlineto{\pgfqpoint{4.048471in}{1.618259in}}%
\pgfpathlineto{\pgfqpoint{4.082303in}{1.597862in}}%
\pgfpathlineto{\pgfqpoint{4.038894in}{1.585419in}}%
\pgfpathlineto{\pgfqpoint{4.005088in}{1.605756in}}%
\pgfpathclose%
\pgfusepath{stroke,fill}%
\end{pgfscope}%
\begin{pgfscope}%
\pgfpathrectangle{\pgfqpoint{0.050000in}{0.050000in}}{\pgfqpoint{4.900000in}{4.900000in}}%
\pgfusepath{clip}%
\pgfsetbuttcap%
\pgfsetroundjoin%
\definecolor{currentfill}{rgb}{0.990773,0.990773,0.990773}%
\pgfsetfillcolor{currentfill}%
\pgfsetfillopacity{0.900000}%
\pgfsetlinewidth{0.501875pt}%
\definecolor{currentstroke}{rgb}{0.200000,0.200000,0.200000}%
\pgfsetstrokecolor{currentstroke}%
\pgfsetdash{}{0pt}%
\pgfpathmoveto{\pgfqpoint{3.455742in}{1.653889in}}%
\pgfpathlineto{\pgfqpoint{3.499590in}{1.657358in}}%
\pgfpathlineto{\pgfqpoint{3.532591in}{1.638030in}}%
\pgfpathlineto{\pgfqpoint{3.488719in}{1.634668in}}%
\pgfpathlineto{\pgfqpoint{3.455742in}{1.653889in}}%
\pgfpathclose%
\pgfusepath{stroke,fill}%
\end{pgfscope}%
\begin{pgfscope}%
\pgfpathrectangle{\pgfqpoint{0.050000in}{0.050000in}}{\pgfqpoint{4.900000in}{4.900000in}}%
\pgfusepath{clip}%
\pgfsetbuttcap%
\pgfsetroundjoin%
\definecolor{currentfill}{rgb}{0.619423,0.619423,0.619423}%
\pgfsetfillcolor{currentfill}%
\pgfsetfillopacity{0.900000}%
\pgfsetlinewidth{0.501875pt}%
\definecolor{currentstroke}{rgb}{0.200000,0.200000,0.200000}%
\pgfsetstrokecolor{currentstroke}%
\pgfsetdash{}{0pt}%
\pgfpathmoveto{\pgfqpoint{1.721219in}{2.778907in}}%
\pgfpathlineto{\pgfqpoint{1.767948in}{2.734329in}}%
\pgfpathlineto{\pgfqpoint{1.798261in}{2.714778in}}%
\pgfpathlineto{\pgfqpoint{1.751524in}{2.759383in}}%
\pgfpathlineto{\pgfqpoint{1.721219in}{2.778907in}}%
\pgfpathclose%
\pgfusepath{stroke,fill}%
\end{pgfscope}%
\begin{pgfscope}%
\pgfpathrectangle{\pgfqpoint{0.050000in}{0.050000in}}{\pgfqpoint{4.900000in}{4.900000in}}%
\pgfusepath{clip}%
\pgfsetbuttcap%
\pgfsetroundjoin%
\definecolor{currentfill}{rgb}{0.858762,0.858762,0.858762}%
\pgfsetfillcolor{currentfill}%
\pgfsetfillopacity{0.900000}%
\pgfsetlinewidth{0.501875pt}%
\definecolor{currentstroke}{rgb}{0.200000,0.200000,0.200000}%
\pgfsetstrokecolor{currentstroke}%
\pgfsetdash{}{0pt}%
\pgfpathmoveto{\pgfqpoint{2.428717in}{2.208537in}}%
\pgfpathlineto{\pgfqpoint{2.473969in}{2.166599in}}%
\pgfpathlineto{\pgfqpoint{2.505364in}{2.146621in}}%
\pgfpathlineto{\pgfqpoint{2.460107in}{2.188405in}}%
\pgfpathlineto{\pgfqpoint{2.428717in}{2.208537in}}%
\pgfpathclose%
\pgfusepath{stroke,fill}%
\end{pgfscope}%
\begin{pgfscope}%
\pgfpathrectangle{\pgfqpoint{0.050000in}{0.050000in}}{\pgfqpoint{4.900000in}{4.900000in}}%
\pgfusepath{clip}%
\pgfsetbuttcap%
\pgfsetroundjoin%
\definecolor{currentfill}{rgb}{0.517186,0.517186,0.517186}%
\pgfsetfillcolor{currentfill}%
\pgfsetfillopacity{0.900000}%
\pgfsetlinewidth{0.501875pt}%
\definecolor{currentstroke}{rgb}{0.200000,0.200000,0.200000}%
\pgfsetstrokecolor{currentstroke}%
\pgfsetdash{}{0pt}%
\pgfpathmoveto{\pgfqpoint{1.459663in}{2.990931in}}%
\pgfpathlineto{\pgfqpoint{1.506945in}{2.945903in}}%
\pgfpathlineto{\pgfqpoint{1.536865in}{2.926671in}}%
\pgfpathlineto{\pgfqpoint{1.489572in}{2.971728in}}%
\pgfpathlineto{\pgfqpoint{1.459663in}{2.990931in}}%
\pgfpathclose%
\pgfusepath{stroke,fill}%
\end{pgfscope}%
\begin{pgfscope}%
\pgfpathrectangle{\pgfqpoint{0.050000in}{0.050000in}}{\pgfqpoint{4.900000in}{4.900000in}}%
\pgfusepath{clip}%
\pgfsetbuttcap%
\pgfsetroundjoin%
\definecolor{currentfill}{rgb}{0.998155,0.998155,0.998155}%
\pgfsetfillcolor{currentfill}%
\pgfsetfillopacity{0.900000}%
\pgfsetlinewidth{0.501875pt}%
\definecolor{currentstroke}{rgb}{0.200000,0.200000,0.200000}%
\pgfsetstrokecolor{currentstroke}%
\pgfsetdash{}{0pt}%
\pgfpathmoveto{\pgfqpoint{3.807290in}{1.611122in}}%
\pgfpathlineto{\pgfqpoint{3.850872in}{1.622175in}}%
\pgfpathlineto{\pgfqpoint{3.884415in}{1.602096in}}%
\pgfpathlineto{\pgfqpoint{3.840807in}{1.591125in}}%
\pgfpathlineto{\pgfqpoint{3.807290in}{1.611122in}}%
\pgfpathclose%
\pgfusepath{stroke,fill}%
\end{pgfscope}%
\begin{pgfscope}%
\pgfpathrectangle{\pgfqpoint{0.050000in}{0.050000in}}{\pgfqpoint{4.900000in}{4.900000in}}%
\pgfusepath{clip}%
\pgfsetbuttcap%
\pgfsetroundjoin%
\definecolor{currentfill}{rgb}{0.915356,0.915356,0.915356}%
\pgfsetfillcolor{currentfill}%
\pgfsetfillopacity{0.900000}%
\pgfsetlinewidth{0.501875pt}%
\definecolor{currentstroke}{rgb}{0.200000,0.200000,0.200000}%
\pgfsetstrokecolor{currentstroke}%
\pgfsetdash{}{0pt}%
\pgfpathmoveto{\pgfqpoint{2.690120in}{2.008964in}}%
\pgfpathlineto{\pgfqpoint{2.734855in}{1.973874in}}%
\pgfpathlineto{\pgfqpoint{2.766658in}{1.954552in}}%
\pgfpathlineto{\pgfqpoint{2.721915in}{1.989336in}}%
\pgfpathlineto{\pgfqpoint{2.690120in}{2.008964in}}%
\pgfpathclose%
\pgfusepath{stroke,fill}%
\end{pgfscope}%
\begin{pgfscope}%
\pgfpathrectangle{\pgfqpoint{0.050000in}{0.050000in}}{\pgfqpoint{4.900000in}{4.900000in}}%
\pgfusepath{clip}%
\pgfsetbuttcap%
\pgfsetroundjoin%
\definecolor{currentfill}{rgb}{0.784667,0.784667,0.784667}%
\pgfsetfillcolor{currentfill}%
\pgfsetfillopacity{0.900000}%
\pgfsetlinewidth{0.501875pt}%
\definecolor{currentstroke}{rgb}{0.200000,0.200000,0.200000}%
\pgfsetstrokecolor{currentstroke}%
\pgfsetdash{}{0pt}%
\pgfpathmoveto{\pgfqpoint{2.167264in}{2.419047in}}%
\pgfpathlineto{\pgfqpoint{2.213071in}{2.375273in}}%
\pgfpathlineto{\pgfqpoint{2.244071in}{2.355208in}}%
\pgfpathlineto{\pgfqpoint{2.198259in}{2.398984in}}%
\pgfpathlineto{\pgfqpoint{2.167264in}{2.419047in}}%
\pgfpathclose%
\pgfusepath{stroke,fill}%
\end{pgfscope}%
\begin{pgfscope}%
\pgfpathrectangle{\pgfqpoint{0.050000in}{0.050000in}}{\pgfqpoint{4.900000in}{4.900000in}}%
\pgfusepath{clip}%
\pgfsetbuttcap%
\pgfsetroundjoin%
\definecolor{currentfill}{rgb}{0.994464,0.994464,0.994464}%
\pgfsetfillcolor{currentfill}%
\pgfsetfillopacity{0.900000}%
\pgfsetlinewidth{0.501875pt}%
\definecolor{currentstroke}{rgb}{0.200000,0.200000,0.200000}%
\pgfsetstrokecolor{currentstroke}%
\pgfsetdash{}{0pt}%
\pgfpathmoveto{\pgfqpoint{3.532591in}{1.638030in}}%
\pgfpathlineto{\pgfqpoint{3.576392in}{1.643691in}}%
\pgfpathlineto{\pgfqpoint{3.609521in}{1.624182in}}%
\pgfpathlineto{\pgfqpoint{3.565695in}{1.618632in}}%
\pgfpathlineto{\pgfqpoint{3.532591in}{1.638030in}}%
\pgfpathclose%
\pgfusepath{stroke,fill}%
\end{pgfscope}%
\begin{pgfscope}%
\pgfpathrectangle{\pgfqpoint{0.050000in}{0.050000in}}{\pgfqpoint{4.900000in}{4.900000in}}%
\pgfusepath{clip}%
\pgfsetbuttcap%
\pgfsetroundjoin%
\definecolor{currentfill}{rgb}{1.000000,1.000000,1.000000}%
\pgfsetfillcolor{currentfill}%
\pgfsetfillopacity{0.900000}%
\pgfsetlinewidth{0.501875pt}%
\definecolor{currentstroke}{rgb}{0.200000,0.200000,0.200000}%
\pgfsetstrokecolor{currentstroke}%
\pgfsetdash{}{0pt}%
\pgfpathmoveto{\pgfqpoint{4.082303in}{1.597862in}}%
\pgfpathlineto{\pgfqpoint{4.125622in}{1.610481in}}%
\pgfpathlineto{\pgfqpoint{4.159585in}{1.589966in}}%
\pgfpathlineto{\pgfqpoint{4.116240in}{1.577404in}}%
\pgfpathlineto{\pgfqpoint{4.082303in}{1.597862in}}%
\pgfpathclose%
\pgfusepath{stroke,fill}%
\end{pgfscope}%
\begin{pgfscope}%
\pgfpathrectangle{\pgfqpoint{0.050000in}{0.050000in}}{\pgfqpoint{4.900000in}{4.900000in}}%
\pgfusepath{clip}%
\pgfsetbuttcap%
\pgfsetroundjoin%
\definecolor{currentfill}{rgb}{0.691396,0.691396,0.691396}%
\pgfsetfillcolor{currentfill}%
\pgfsetfillopacity{0.900000}%
\pgfsetlinewidth{0.501875pt}%
\definecolor{currentstroke}{rgb}{0.200000,0.200000,0.200000}%
\pgfsetstrokecolor{currentstroke}%
\pgfsetdash{}{0pt}%
\pgfpathmoveto{\pgfqpoint{1.905721in}{2.631025in}}%
\pgfpathlineto{\pgfqpoint{1.952083in}{2.586739in}}%
\pgfpathlineto{\pgfqpoint{1.982689in}{2.566955in}}%
\pgfpathlineto{\pgfqpoint{1.936321in}{2.611268in}}%
\pgfpathlineto{\pgfqpoint{1.905721in}{2.631025in}}%
\pgfpathclose%
\pgfusepath{stroke,fill}%
\end{pgfscope}%
\begin{pgfscope}%
\pgfpathrectangle{\pgfqpoint{0.050000in}{0.050000in}}{\pgfqpoint{4.900000in}{4.900000in}}%
\pgfusepath{clip}%
\pgfsetbuttcap%
\pgfsetroundjoin%
\definecolor{currentfill}{rgb}{0.581776,0.581776,0.581776}%
\pgfsetfillcolor{currentfill}%
\pgfsetfillopacity{0.900000}%
\pgfsetlinewidth{0.501875pt}%
\definecolor{currentstroke}{rgb}{0.200000,0.200000,0.200000}%
\pgfsetstrokecolor{currentstroke}%
\pgfsetdash{}{0pt}%
\pgfpathmoveto{\pgfqpoint{1.644091in}{2.843108in}}%
\pgfpathlineto{\pgfqpoint{1.691007in}{2.798370in}}%
\pgfpathlineto{\pgfqpoint{1.721219in}{2.778907in}}%
\pgfpathlineto{\pgfqpoint{1.674295in}{2.823671in}}%
\pgfpathlineto{\pgfqpoint{1.644091in}{2.843108in}}%
\pgfpathclose%
\pgfusepath{stroke,fill}%
\end{pgfscope}%
\begin{pgfscope}%
\pgfpathrectangle{\pgfqpoint{0.050000in}{0.050000in}}{\pgfqpoint{4.900000in}{4.900000in}}%
\pgfusepath{clip}%
\pgfsetbuttcap%
\pgfsetroundjoin%
\definecolor{currentfill}{rgb}{0.994464,0.994464,0.994464}%
\pgfsetfillcolor{currentfill}%
\pgfsetfillopacity{0.900000}%
\pgfsetlinewidth{0.501875pt}%
\definecolor{currentstroke}{rgb}{0.200000,0.200000,0.200000}%
\pgfsetstrokecolor{currentstroke}%
\pgfsetdash{}{0pt}%
\pgfpathmoveto{\pgfqpoint{3.609521in}{1.624182in}}%
\pgfpathlineto{\pgfqpoint{3.653277in}{1.631647in}}%
\pgfpathlineto{\pgfqpoint{3.686534in}{1.611965in}}%
\pgfpathlineto{\pgfqpoint{3.642753in}{1.604609in}}%
\pgfpathlineto{\pgfqpoint{3.609521in}{1.624182in}}%
\pgfpathclose%
\pgfusepath{stroke,fill}%
\end{pgfscope}%
\begin{pgfscope}%
\pgfpathrectangle{\pgfqpoint{0.050000in}{0.050000in}}{\pgfqpoint{4.900000in}{4.900000in}}%
\pgfusepath{clip}%
\pgfsetbuttcap%
\pgfsetroundjoin%
\definecolor{currentfill}{rgb}{1.000000,1.000000,1.000000}%
\pgfsetfillcolor{currentfill}%
\pgfsetfillopacity{0.900000}%
\pgfsetlinewidth{0.501875pt}%
\definecolor{currentstroke}{rgb}{0.200000,0.200000,0.200000}%
\pgfsetstrokecolor{currentstroke}%
\pgfsetdash{}{0pt}%
\pgfpathmoveto{\pgfqpoint{3.884415in}{1.602096in}}%
\pgfpathlineto{\pgfqpoint{3.927943in}{1.613804in}}%
\pgfpathlineto{\pgfqpoint{3.961618in}{1.593592in}}%
\pgfpathlineto{\pgfqpoint{3.918063in}{1.581957in}}%
\pgfpathlineto{\pgfqpoint{3.884415in}{1.602096in}}%
\pgfpathclose%
\pgfusepath{stroke,fill}%
\end{pgfscope}%
\begin{pgfscope}%
\pgfpathrectangle{\pgfqpoint{0.050000in}{0.050000in}}{\pgfqpoint{4.900000in}{4.900000in}}%
\pgfusepath{clip}%
\pgfsetbuttcap%
\pgfsetroundjoin%
\definecolor{currentfill}{rgb}{0.836340,0.836340,0.836340}%
\pgfsetfillcolor{currentfill}%
\pgfsetfillopacity{0.900000}%
\pgfsetlinewidth{0.501875pt}%
\definecolor{currentstroke}{rgb}{0.200000,0.200000,0.200000}%
\pgfsetstrokecolor{currentstroke}%
\pgfsetdash{}{0pt}%
\pgfpathmoveto{\pgfqpoint{2.351985in}{2.271512in}}%
\pgfpathlineto{\pgfqpoint{2.397423in}{2.228631in}}%
\pgfpathlineto{\pgfqpoint{2.428717in}{2.208537in}}%
\pgfpathlineto{\pgfqpoint{2.383274in}{2.251331in}}%
\pgfpathlineto{\pgfqpoint{2.351985in}{2.271512in}}%
\pgfpathclose%
\pgfusepath{stroke,fill}%
\end{pgfscope}%
\begin{pgfscope}%
\pgfpathrectangle{\pgfqpoint{0.050000in}{0.050000in}}{\pgfqpoint{4.900000in}{4.900000in}}%
\pgfusepath{clip}%
\pgfsetbuttcap%
\pgfsetroundjoin%
\definecolor{currentfill}{rgb}{0.898378,0.898378,0.898378}%
\pgfsetfillcolor{currentfill}%
\pgfsetfillopacity{0.900000}%
\pgfsetlinewidth{0.501875pt}%
\definecolor{currentstroke}{rgb}{0.200000,0.200000,0.200000}%
\pgfsetstrokecolor{currentstroke}%
\pgfsetdash{}{0pt}%
\pgfpathmoveto{\pgfqpoint{2.613523in}{2.066556in}}%
\pgfpathlineto{\pgfqpoint{2.658423in}{2.028548in}}%
\pgfpathlineto{\pgfqpoint{2.690120in}{2.008964in}}%
\pgfpathlineto{\pgfqpoint{2.645213in}{2.046695in}}%
\pgfpathlineto{\pgfqpoint{2.613523in}{2.066556in}}%
\pgfpathclose%
\pgfusepath{stroke,fill}%
\end{pgfscope}%
\begin{pgfscope}%
\pgfpathrectangle{\pgfqpoint{0.050000in}{0.050000in}}{\pgfqpoint{4.900000in}{4.900000in}}%
\pgfusepath{clip}%
\pgfsetbuttcap%
\pgfsetroundjoin%
\definecolor{currentfill}{rgb}{0.487043,0.487043,0.487043}%
\pgfsetfillcolor{currentfill}%
\pgfsetfillopacity{0.900000}%
\pgfsetlinewidth{0.501875pt}%
\definecolor{currentstroke}{rgb}{0.200000,0.200000,0.200000}%
\pgfsetstrokecolor{currentstroke}%
\pgfsetdash{}{0pt}%
\pgfpathmoveto{\pgfqpoint{1.382374in}{3.055263in}}%
\pgfpathlineto{\pgfqpoint{1.429845in}{3.010075in}}%
\pgfpathlineto{\pgfqpoint{1.459663in}{2.990931in}}%
\pgfpathlineto{\pgfqpoint{1.412182in}{3.036148in}}%
\pgfpathlineto{\pgfqpoint{1.382374in}{3.055263in}}%
\pgfpathclose%
\pgfusepath{stroke,fill}%
\end{pgfscope}%
\begin{pgfscope}%
\pgfpathrectangle{\pgfqpoint{0.050000in}{0.050000in}}{\pgfqpoint{4.900000in}{4.900000in}}%
\pgfusepath{clip}%
\pgfsetbuttcap%
\pgfsetroundjoin%
\definecolor{currentfill}{rgb}{0.970473,0.970473,0.970473}%
\pgfsetfillcolor{currentfill}%
\pgfsetfillopacity{0.900000}%
\pgfsetlinewidth{0.501875pt}%
\definecolor{currentstroke}{rgb}{0.200000,0.200000,0.200000}%
\pgfsetstrokecolor{currentstroke}%
\pgfsetdash{}{0pt}%
\pgfpathmoveto{\pgfqpoint{3.104937in}{1.760149in}}%
\pgfpathlineto{\pgfqpoint{3.149094in}{1.746299in}}%
\pgfpathlineto{\pgfqpoint{3.181573in}{1.727567in}}%
\pgfpathlineto{\pgfqpoint{3.137398in}{1.741331in}}%
\pgfpathlineto{\pgfqpoint{3.104937in}{1.760149in}}%
\pgfpathclose%
\pgfusepath{stroke,fill}%
\end{pgfscope}%
\begin{pgfscope}%
\pgfpathrectangle{\pgfqpoint{0.050000in}{0.050000in}}{\pgfqpoint{4.900000in}{4.900000in}}%
\pgfusepath{clip}%
\pgfsetbuttcap%
\pgfsetroundjoin%
\definecolor{currentfill}{rgb}{0.976009,0.976009,0.976009}%
\pgfsetfillcolor{currentfill}%
\pgfsetfillopacity{0.900000}%
\pgfsetlinewidth{0.501875pt}%
\definecolor{currentstroke}{rgb}{0.200000,0.200000,0.200000}%
\pgfsetstrokecolor{currentstroke}%
\pgfsetdash{}{0pt}%
\pgfpathmoveto{\pgfqpoint{3.181573in}{1.727567in}}%
\pgfpathlineto{\pgfqpoint{3.225657in}{1.717930in}}%
\pgfpathlineto{\pgfqpoint{3.258259in}{1.699134in}}%
\pgfpathlineto{\pgfqpoint{3.214154in}{1.708746in}}%
\pgfpathlineto{\pgfqpoint{3.181573in}{1.727567in}}%
\pgfpathclose%
\pgfusepath{stroke,fill}%
\end{pgfscope}%
\begin{pgfscope}%
\pgfpathrectangle{\pgfqpoint{0.050000in}{0.050000in}}{\pgfqpoint{4.900000in}{4.900000in}}%
\pgfusepath{clip}%
\pgfsetbuttcap%
\pgfsetroundjoin%
\definecolor{currentfill}{rgb}{0.963091,0.963091,0.963091}%
\pgfsetfillcolor{currentfill}%
\pgfsetfillopacity{0.900000}%
\pgfsetlinewidth{0.501875pt}%
\definecolor{currentstroke}{rgb}{0.200000,0.200000,0.200000}%
\pgfsetstrokecolor{currentstroke}%
\pgfsetdash{}{0pt}%
\pgfpathmoveto{\pgfqpoint{3.028336in}{1.797161in}}%
\pgfpathlineto{\pgfqpoint{3.072578in}{1.778880in}}%
\pgfpathlineto{\pgfqpoint{3.104937in}{1.760149in}}%
\pgfpathlineto{\pgfqpoint{3.060678in}{1.778279in}}%
\pgfpathlineto{\pgfqpoint{3.028336in}{1.797161in}}%
\pgfpathclose%
\pgfusepath{stroke,fill}%
\end{pgfscope}%
\begin{pgfscope}%
\pgfpathrectangle{\pgfqpoint{0.050000in}{0.050000in}}{\pgfqpoint{4.900000in}{4.900000in}}%
\pgfusepath{clip}%
\pgfsetbuttcap%
\pgfsetroundjoin%
\definecolor{currentfill}{rgb}{0.757109,0.757109,0.757109}%
\pgfsetfillcolor{currentfill}%
\pgfsetfillopacity{0.900000}%
\pgfsetlinewidth{0.501875pt}%
\definecolor{currentstroke}{rgb}{0.200000,0.200000,0.200000}%
\pgfsetstrokecolor{currentstroke}%
\pgfsetdash{}{0pt}%
\pgfpathmoveto{\pgfqpoint{2.090370in}{2.483033in}}%
\pgfpathlineto{\pgfqpoint{2.136364in}{2.439053in}}%
\pgfpathlineto{\pgfqpoint{2.167264in}{2.419047in}}%
\pgfpathlineto{\pgfqpoint{2.121264in}{2.463046in}}%
\pgfpathlineto{\pgfqpoint{2.090370in}{2.483033in}}%
\pgfpathclose%
\pgfusepath{stroke,fill}%
\end{pgfscope}%
\begin{pgfscope}%
\pgfpathrectangle{\pgfqpoint{0.050000in}{0.050000in}}{\pgfqpoint{4.900000in}{4.900000in}}%
\pgfusepath{clip}%
\pgfsetbuttcap%
\pgfsetroundjoin%
\definecolor{currentfill}{rgb}{0.981546,0.981546,0.981546}%
\pgfsetfillcolor{currentfill}%
\pgfsetfillopacity{0.900000}%
\pgfsetlinewidth{0.501875pt}%
\definecolor{currentstroke}{rgb}{0.200000,0.200000,0.200000}%
\pgfsetstrokecolor{currentstroke}%
\pgfsetdash{}{0pt}%
\pgfpathmoveto{\pgfqpoint{3.258259in}{1.699134in}}%
\pgfpathlineto{\pgfqpoint{3.302281in}{1.693376in}}%
\pgfpathlineto{\pgfqpoint{3.335006in}{1.674468in}}%
\pgfpathlineto{\pgfqpoint{3.290963in}{1.680251in}}%
\pgfpathlineto{\pgfqpoint{3.258259in}{1.699134in}}%
\pgfpathclose%
\pgfusepath{stroke,fill}%
\end{pgfscope}%
\begin{pgfscope}%
\pgfpathrectangle{\pgfqpoint{0.050000in}{0.050000in}}{\pgfqpoint{4.900000in}{4.900000in}}%
\pgfusepath{clip}%
\pgfsetbuttcap%
\pgfsetroundjoin%
\definecolor{currentfill}{rgb}{0.955709,0.955709,0.955709}%
\pgfsetfillcolor{currentfill}%
\pgfsetfillopacity{0.900000}%
\pgfsetlinewidth{0.501875pt}%
\definecolor{currentstroke}{rgb}{0.200000,0.200000,0.200000}%
\pgfsetstrokecolor{currentstroke}%
\pgfsetdash{}{0pt}%
\pgfpathmoveto{\pgfqpoint{2.951752in}{1.838737in}}%
\pgfpathlineto{\pgfqpoint{2.996095in}{1.815959in}}%
\pgfpathlineto{\pgfqpoint{3.028336in}{1.797161in}}%
\pgfpathlineto{\pgfqpoint{2.983978in}{1.819725in}}%
\pgfpathlineto{\pgfqpoint{2.951752in}{1.838737in}}%
\pgfpathclose%
\pgfusepath{stroke,fill}%
\end{pgfscope}%
\begin{pgfscope}%
\pgfpathrectangle{\pgfqpoint{0.050000in}{0.050000in}}{\pgfqpoint{4.900000in}{4.900000in}}%
\pgfusepath{clip}%
\pgfsetbuttcap%
\pgfsetroundjoin%
\definecolor{currentfill}{rgb}{0.996309,0.996309,0.996309}%
\pgfsetfillcolor{currentfill}%
\pgfsetfillopacity{0.900000}%
\pgfsetlinewidth{0.501875pt}%
\definecolor{currentstroke}{rgb}{0.200000,0.200000,0.200000}%
\pgfsetstrokecolor{currentstroke}%
\pgfsetdash{}{0pt}%
\pgfpathmoveto{\pgfqpoint{3.686534in}{1.611965in}}%
\pgfpathlineto{\pgfqpoint{3.730243in}{1.620888in}}%
\pgfpathlineto{\pgfqpoint{3.763630in}{1.601043in}}%
\pgfpathlineto{\pgfqpoint{3.719895in}{1.592221in}}%
\pgfpathlineto{\pgfqpoint{3.686534in}{1.611965in}}%
\pgfpathclose%
\pgfusepath{stroke,fill}%
\end{pgfscope}%
\begin{pgfscope}%
\pgfpathrectangle{\pgfqpoint{0.050000in}{0.050000in}}{\pgfqpoint{4.900000in}{4.900000in}}%
\pgfusepath{clip}%
\pgfsetbuttcap%
\pgfsetroundjoin%
\definecolor{currentfill}{rgb}{0.653010,0.653010,0.653010}%
\pgfsetfillcolor{currentfill}%
\pgfsetfillopacity{0.900000}%
\pgfsetlinewidth{0.501875pt}%
\definecolor{currentstroke}{rgb}{0.200000,0.200000,0.200000}%
\pgfsetstrokecolor{currentstroke}%
\pgfsetdash{}{0pt}%
\pgfpathmoveto{\pgfqpoint{1.828666in}{2.695168in}}%
\pgfpathlineto{\pgfqpoint{1.875215in}{2.650722in}}%
\pgfpathlineto{\pgfqpoint{1.905721in}{2.631025in}}%
\pgfpathlineto{\pgfqpoint{1.859165in}{2.675498in}}%
\pgfpathlineto{\pgfqpoint{1.828666in}{2.695168in}}%
\pgfpathclose%
\pgfusepath{stroke,fill}%
\end{pgfscope}%
\begin{pgfscope}%
\pgfpathrectangle{\pgfqpoint{0.050000in}{0.050000in}}{\pgfqpoint{4.900000in}{4.900000in}}%
\pgfusepath{clip}%
\pgfsetbuttcap%
\pgfsetroundjoin%
\definecolor{currentfill}{rgb}{0.985236,0.985236,0.985236}%
\pgfsetfillcolor{currentfill}%
\pgfsetfillopacity{0.900000}%
\pgfsetlinewidth{0.501875pt}%
\definecolor{currentstroke}{rgb}{0.200000,0.200000,0.200000}%
\pgfsetstrokecolor{currentstroke}%
\pgfsetdash{}{0pt}%
\pgfpathmoveto{\pgfqpoint{3.335006in}{1.674468in}}%
\pgfpathlineto{\pgfqpoint{3.378974in}{1.672186in}}%
\pgfpathlineto{\pgfqpoint{3.411824in}{1.653131in}}%
\pgfpathlineto{\pgfqpoint{3.367834in}{1.655477in}}%
\pgfpathlineto{\pgfqpoint{3.335006in}{1.674468in}}%
\pgfpathclose%
\pgfusepath{stroke,fill}%
\end{pgfscope}%
\begin{pgfscope}%
\pgfpathrectangle{\pgfqpoint{0.050000in}{0.050000in}}{\pgfqpoint{4.900000in}{4.900000in}}%
\pgfusepath{clip}%
\pgfsetbuttcap%
\pgfsetroundjoin%
\definecolor{currentfill}{rgb}{1.000000,1.000000,1.000000}%
\pgfsetfillcolor{currentfill}%
\pgfsetfillopacity{0.900000}%
\pgfsetlinewidth{0.501875pt}%
\definecolor{currentstroke}{rgb}{0.200000,0.200000,0.200000}%
\pgfsetstrokecolor{currentstroke}%
\pgfsetdash{}{0pt}%
\pgfpathmoveto{\pgfqpoint{3.961618in}{1.593592in}}%
\pgfpathlineto{\pgfqpoint{4.005088in}{1.605756in}}%
\pgfpathlineto{\pgfqpoint{4.038894in}{1.585419in}}%
\pgfpathlineto{\pgfqpoint{3.995397in}{1.573320in}}%
\pgfpathlineto{\pgfqpoint{3.961618in}{1.593592in}}%
\pgfpathclose%
\pgfusepath{stroke,fill}%
\end{pgfscope}%
\begin{pgfscope}%
\pgfpathrectangle{\pgfqpoint{0.050000in}{0.050000in}}{\pgfqpoint{4.900000in}{4.900000in}}%
\pgfusepath{clip}%
\pgfsetbuttcap%
\pgfsetroundjoin%
\definecolor{currentfill}{rgb}{0.946482,0.946482,0.946482}%
\pgfsetfillcolor{currentfill}%
\pgfsetfillopacity{0.900000}%
\pgfsetlinewidth{0.501875pt}%
\definecolor{currentstroke}{rgb}{0.200000,0.200000,0.200000}%
\pgfsetstrokecolor{currentstroke}%
\pgfsetdash{}{0pt}%
\pgfpathmoveto{\pgfqpoint{2.875167in}{1.884827in}}%
\pgfpathlineto{\pgfqpoint{2.919627in}{1.857673in}}%
\pgfpathlineto{\pgfqpoint{2.951752in}{1.838737in}}%
\pgfpathlineto{\pgfqpoint{2.907279in}{1.865626in}}%
\pgfpathlineto{\pgfqpoint{2.875167in}{1.884827in}}%
\pgfpathclose%
\pgfusepath{stroke,fill}%
\end{pgfscope}%
\begin{pgfscope}%
\pgfpathrectangle{\pgfqpoint{0.050000in}{0.050000in}}{\pgfqpoint{4.900000in}{4.900000in}}%
\pgfusepath{clip}%
\pgfsetbuttcap%
\pgfsetroundjoin%
\definecolor{currentfill}{rgb}{0.878570,0.878570,0.878570}%
\pgfsetfillcolor{currentfill}%
\pgfsetfillopacity{0.900000}%
\pgfsetlinewidth{0.501875pt}%
\definecolor{currentstroke}{rgb}{0.200000,0.200000,0.200000}%
\pgfsetstrokecolor{currentstroke}%
\pgfsetdash{}{0pt}%
\pgfpathmoveto{\pgfqpoint{2.536854in}{2.126608in}}%
\pgfpathlineto{\pgfqpoint{2.581930in}{2.086381in}}%
\pgfpathlineto{\pgfqpoint{2.613523in}{2.066556in}}%
\pgfpathlineto{\pgfqpoint{2.568442in}{2.106560in}}%
\pgfpathlineto{\pgfqpoint{2.536854in}{2.126608in}}%
\pgfpathclose%
\pgfusepath{stroke,fill}%
\end{pgfscope}%
\begin{pgfscope}%
\pgfpathrectangle{\pgfqpoint{0.050000in}{0.050000in}}{\pgfqpoint{4.900000in}{4.900000in}}%
\pgfusepath{clip}%
\pgfsetbuttcap%
\pgfsetroundjoin%
\definecolor{currentfill}{rgb}{0.551634,0.551634,0.551634}%
\pgfsetfillcolor{currentfill}%
\pgfsetfillopacity{0.900000}%
\pgfsetlinewidth{0.501875pt}%
\definecolor{currentstroke}{rgb}{0.200000,0.200000,0.200000}%
\pgfsetstrokecolor{currentstroke}%
\pgfsetdash{}{0pt}%
\pgfpathmoveto{\pgfqpoint{1.566875in}{2.907381in}}%
\pgfpathlineto{\pgfqpoint{1.613979in}{2.862485in}}%
\pgfpathlineto{\pgfqpoint{1.644091in}{2.843108in}}%
\pgfpathlineto{\pgfqpoint{1.596978in}{2.888033in}}%
\pgfpathlineto{\pgfqpoint{1.566875in}{2.907381in}}%
\pgfpathclose%
\pgfusepath{stroke,fill}%
\end{pgfscope}%
\begin{pgfscope}%
\pgfpathrectangle{\pgfqpoint{0.050000in}{0.050000in}}{\pgfqpoint{4.900000in}{4.900000in}}%
\pgfusepath{clip}%
\pgfsetbuttcap%
\pgfsetroundjoin%
\definecolor{currentfill}{rgb}{0.808781,0.808781,0.808781}%
\pgfsetfillcolor{currentfill}%
\pgfsetfillopacity{0.900000}%
\pgfsetlinewidth{0.501875pt}%
\definecolor{currentstroke}{rgb}{0.200000,0.200000,0.200000}%
\pgfsetstrokecolor{currentstroke}%
\pgfsetdash{}{0pt}%
\pgfpathmoveto{\pgfqpoint{2.275165in}{2.335089in}}%
\pgfpathlineto{\pgfqpoint{2.320790in}{2.291647in}}%
\pgfpathlineto{\pgfqpoint{2.351985in}{2.271512in}}%
\pgfpathlineto{\pgfqpoint{2.306355in}{2.314919in}}%
\pgfpathlineto{\pgfqpoint{2.275165in}{2.335089in}}%
\pgfpathclose%
\pgfusepath{stroke,fill}%
\end{pgfscope}%
\begin{pgfscope}%
\pgfpathrectangle{\pgfqpoint{0.050000in}{0.050000in}}{\pgfqpoint{4.900000in}{4.900000in}}%
\pgfusepath{clip}%
\pgfsetbuttcap%
\pgfsetroundjoin%
\definecolor{currentfill}{rgb}{0.988927,0.988927,0.988927}%
\pgfsetfillcolor{currentfill}%
\pgfsetfillopacity{0.900000}%
\pgfsetlinewidth{0.501875pt}%
\definecolor{currentstroke}{rgb}{0.200000,0.200000,0.200000}%
\pgfsetstrokecolor{currentstroke}%
\pgfsetdash{}{0pt}%
\pgfpathmoveto{\pgfqpoint{3.411824in}{1.653131in}}%
\pgfpathlineto{\pgfqpoint{3.455742in}{1.653889in}}%
\pgfpathlineto{\pgfqpoint{3.488719in}{1.634668in}}%
\pgfpathlineto{\pgfqpoint{3.444777in}{1.633998in}}%
\pgfpathlineto{\pgfqpoint{3.411824in}{1.653131in}}%
\pgfpathclose%
\pgfusepath{stroke,fill}%
\end{pgfscope}%
\begin{pgfscope}%
\pgfpathrectangle{\pgfqpoint{0.050000in}{0.050000in}}{\pgfqpoint{4.900000in}{4.900000in}}%
\pgfusepath{clip}%
\pgfsetbuttcap%
\pgfsetroundjoin%
\definecolor{currentfill}{rgb}{0.452595,0.452595,0.452595}%
\pgfsetfillcolor{currentfill}%
\pgfsetfillopacity{0.900000}%
\pgfsetlinewidth{0.501875pt}%
\definecolor{currentstroke}{rgb}{0.200000,0.200000,0.200000}%
\pgfsetstrokecolor{currentstroke}%
\pgfsetdash{}{0pt}%
\pgfpathmoveto{\pgfqpoint{1.304997in}{3.119668in}}%
\pgfpathlineto{\pgfqpoint{1.352657in}{3.074320in}}%
\pgfpathlineto{\pgfqpoint{1.382374in}{3.055263in}}%
\pgfpathlineto{\pgfqpoint{1.334704in}{3.100641in}}%
\pgfpathlineto{\pgfqpoint{1.304997in}{3.119668in}}%
\pgfpathclose%
\pgfusepath{stroke,fill}%
\end{pgfscope}%
\begin{pgfscope}%
\pgfpathrectangle{\pgfqpoint{0.050000in}{0.050000in}}{\pgfqpoint{4.900000in}{4.900000in}}%
\pgfusepath{clip}%
\pgfsetbuttcap%
\pgfsetroundjoin%
\definecolor{currentfill}{rgb}{0.998155,0.998155,0.998155}%
\pgfsetfillcolor{currentfill}%
\pgfsetfillopacity{0.900000}%
\pgfsetlinewidth{0.501875pt}%
\definecolor{currentstroke}{rgb}{0.200000,0.200000,0.200000}%
\pgfsetstrokecolor{currentstroke}%
\pgfsetdash{}{0pt}%
\pgfpathmoveto{\pgfqpoint{3.763630in}{1.601043in}}%
\pgfpathlineto{\pgfqpoint{3.807290in}{1.611122in}}%
\pgfpathlineto{\pgfqpoint{3.840807in}{1.591125in}}%
\pgfpathlineto{\pgfqpoint{3.797121in}{1.581138in}}%
\pgfpathlineto{\pgfqpoint{3.763630in}{1.601043in}}%
\pgfpathclose%
\pgfusepath{stroke,fill}%
\end{pgfscope}%
\begin{pgfscope}%
\pgfpathrectangle{\pgfqpoint{0.050000in}{0.050000in}}{\pgfqpoint{4.900000in}{4.900000in}}%
\pgfusepath{clip}%
\pgfsetbuttcap%
\pgfsetroundjoin%
\definecolor{currentfill}{rgb}{0.729781,0.729781,0.729781}%
\pgfsetfillcolor{currentfill}%
\pgfsetfillopacity{0.900000}%
\pgfsetlinewidth{0.501875pt}%
\definecolor{currentstroke}{rgb}{0.200000,0.200000,0.200000}%
\pgfsetstrokecolor{currentstroke}%
\pgfsetdash{}{0pt}%
\pgfpathmoveto{\pgfqpoint{2.013389in}{2.547112in}}%
\pgfpathlineto{\pgfqpoint{2.059570in}{2.502961in}}%
\pgfpathlineto{\pgfqpoint{2.090370in}{2.483033in}}%
\pgfpathlineto{\pgfqpoint{2.044183in}{2.527208in}}%
\pgfpathlineto{\pgfqpoint{2.013389in}{2.547112in}}%
\pgfpathclose%
\pgfusepath{stroke,fill}%
\end{pgfscope}%
\begin{pgfscope}%
\pgfpathrectangle{\pgfqpoint{0.050000in}{0.050000in}}{\pgfqpoint{4.900000in}{4.900000in}}%
\pgfusepath{clip}%
\pgfsetbuttcap%
\pgfsetroundjoin%
\definecolor{currentfill}{rgb}{0.932334,0.932334,0.932334}%
\pgfsetfillcolor{currentfill}%
\pgfsetfillopacity{0.900000}%
\pgfsetlinewidth{0.501875pt}%
\definecolor{currentstroke}{rgb}{0.200000,0.200000,0.200000}%
\pgfsetstrokecolor{currentstroke}%
\pgfsetdash{}{0pt}%
\pgfpathmoveto{\pgfqpoint{2.798561in}{1.935177in}}%
\pgfpathlineto{\pgfqpoint{2.843154in}{1.903962in}}%
\pgfpathlineto{\pgfqpoint{2.875167in}{1.884827in}}%
\pgfpathlineto{\pgfqpoint{2.830563in}{1.915747in}}%
\pgfpathlineto{\pgfqpoint{2.798561in}{1.935177in}}%
\pgfpathclose%
\pgfusepath{stroke,fill}%
\end{pgfscope}%
\begin{pgfscope}%
\pgfpathrectangle{\pgfqpoint{0.050000in}{0.050000in}}{\pgfqpoint{4.900000in}{4.900000in}}%
\pgfusepath{clip}%
\pgfsetbuttcap%
\pgfsetroundjoin%
\definecolor{currentfill}{rgb}{1.000000,1.000000,1.000000}%
\pgfsetfillcolor{currentfill}%
\pgfsetfillopacity{0.900000}%
\pgfsetlinewidth{0.501875pt}%
\definecolor{currentstroke}{rgb}{0.200000,0.200000,0.200000}%
\pgfsetstrokecolor{currentstroke}%
\pgfsetdash{}{0pt}%
\pgfpathmoveto{\pgfqpoint{4.038894in}{1.585419in}}%
\pgfpathlineto{\pgfqpoint{4.082303in}{1.597862in}}%
\pgfpathlineto{\pgfqpoint{4.116240in}{1.577404in}}%
\pgfpathlineto{\pgfqpoint{4.072804in}{1.565021in}}%
\pgfpathlineto{\pgfqpoint{4.038894in}{1.585419in}}%
\pgfpathclose%
\pgfusepath{stroke,fill}%
\end{pgfscope}%
\begin{pgfscope}%
\pgfpathrectangle{\pgfqpoint{0.050000in}{0.050000in}}{\pgfqpoint{4.900000in}{4.900000in}}%
\pgfusepath{clip}%
\pgfsetbuttcap%
\pgfsetroundjoin%
\definecolor{currentfill}{rgb}{0.990773,0.990773,0.990773}%
\pgfsetfillcolor{currentfill}%
\pgfsetfillopacity{0.900000}%
\pgfsetlinewidth{0.501875pt}%
\definecolor{currentstroke}{rgb}{0.200000,0.200000,0.200000}%
\pgfsetstrokecolor{currentstroke}%
\pgfsetdash{}{0pt}%
\pgfpathmoveto{\pgfqpoint{3.488719in}{1.634668in}}%
\pgfpathlineto{\pgfqpoint{3.532591in}{1.638030in}}%
\pgfpathlineto{\pgfqpoint{3.565695in}{1.618632in}}%
\pgfpathlineto{\pgfqpoint{3.521799in}{1.615373in}}%
\pgfpathlineto{\pgfqpoint{3.488719in}{1.634668in}}%
\pgfpathclose%
\pgfusepath{stroke,fill}%
\end{pgfscope}%
\begin{pgfscope}%
\pgfpathrectangle{\pgfqpoint{0.050000in}{0.050000in}}{\pgfqpoint{4.900000in}{4.900000in}}%
\pgfusepath{clip}%
\pgfsetbuttcap%
\pgfsetroundjoin%
\definecolor{currentfill}{rgb}{0.619423,0.619423,0.619423}%
\pgfsetfillcolor{currentfill}%
\pgfsetfillopacity{0.900000}%
\pgfsetlinewidth{0.501875pt}%
\definecolor{currentstroke}{rgb}{0.200000,0.200000,0.200000}%
\pgfsetstrokecolor{currentstroke}%
\pgfsetdash{}{0pt}%
\pgfpathmoveto{\pgfqpoint{1.751524in}{2.759383in}}%
\pgfpathlineto{\pgfqpoint{1.798261in}{2.714778in}}%
\pgfpathlineto{\pgfqpoint{1.828666in}{2.695168in}}%
\pgfpathlineto{\pgfqpoint{1.781922in}{2.739800in}}%
\pgfpathlineto{\pgfqpoint{1.751524in}{2.759383in}}%
\pgfpathclose%
\pgfusepath{stroke,fill}%
\end{pgfscope}%
\begin{pgfscope}%
\pgfpathrectangle{\pgfqpoint{0.050000in}{0.050000in}}{\pgfqpoint{4.900000in}{4.900000in}}%
\pgfusepath{clip}%
\pgfsetbuttcap%
\pgfsetroundjoin%
\definecolor{currentfill}{rgb}{0.858762,0.858762,0.858762}%
\pgfsetfillcolor{currentfill}%
\pgfsetfillopacity{0.900000}%
\pgfsetlinewidth{0.501875pt}%
\definecolor{currentstroke}{rgb}{0.200000,0.200000,0.200000}%
\pgfsetstrokecolor{currentstroke}%
\pgfsetdash{}{0pt}%
\pgfpathmoveto{\pgfqpoint{2.460107in}{2.188405in}}%
\pgfpathlineto{\pgfqpoint{2.505364in}{2.146621in}}%
\pgfpathlineto{\pgfqpoint{2.536854in}{2.126608in}}%
\pgfpathlineto{\pgfqpoint{2.491593in}{2.168235in}}%
\pgfpathlineto{\pgfqpoint{2.460107in}{2.188405in}}%
\pgfpathclose%
\pgfusepath{stroke,fill}%
\end{pgfscope}%
\begin{pgfscope}%
\pgfpathrectangle{\pgfqpoint{0.050000in}{0.050000in}}{\pgfqpoint{4.900000in}{4.900000in}}%
\pgfusepath{clip}%
\pgfsetbuttcap%
\pgfsetroundjoin%
\definecolor{currentfill}{rgb}{0.517186,0.517186,0.517186}%
\pgfsetfillcolor{currentfill}%
\pgfsetfillopacity{0.900000}%
\pgfsetlinewidth{0.501875pt}%
\definecolor{currentstroke}{rgb}{0.200000,0.200000,0.200000}%
\pgfsetstrokecolor{currentstroke}%
\pgfsetdash{}{0pt}%
\pgfpathmoveto{\pgfqpoint{1.489572in}{2.971728in}}%
\pgfpathlineto{\pgfqpoint{1.536865in}{2.926671in}}%
\pgfpathlineto{\pgfqpoint{1.566875in}{2.907381in}}%
\pgfpathlineto{\pgfqpoint{1.519574in}{2.952467in}}%
\pgfpathlineto{\pgfqpoint{1.489572in}{2.971728in}}%
\pgfpathclose%
\pgfusepath{stroke,fill}%
\end{pgfscope}%
\begin{pgfscope}%
\pgfpathrectangle{\pgfqpoint{0.050000in}{0.050000in}}{\pgfqpoint{4.900000in}{4.900000in}}%
\pgfusepath{clip}%
\pgfsetbuttcap%
\pgfsetroundjoin%
\definecolor{currentfill}{rgb}{0.998155,0.998155,0.998155}%
\pgfsetfillcolor{currentfill}%
\pgfsetfillopacity{0.900000}%
\pgfsetlinewidth{0.501875pt}%
\definecolor{currentstroke}{rgb}{0.200000,0.200000,0.200000}%
\pgfsetstrokecolor{currentstroke}%
\pgfsetdash{}{0pt}%
\pgfpathmoveto{\pgfqpoint{3.840807in}{1.591125in}}%
\pgfpathlineto{\pgfqpoint{3.884415in}{1.602096in}}%
\pgfpathlineto{\pgfqpoint{3.918063in}{1.581957in}}%
\pgfpathlineto{\pgfqpoint{3.874428in}{1.571068in}}%
\pgfpathlineto{\pgfqpoint{3.840807in}{1.591125in}}%
\pgfpathclose%
\pgfusepath{stroke,fill}%
\end{pgfscope}%
\begin{pgfscope}%
\pgfpathrectangle{\pgfqpoint{0.050000in}{0.050000in}}{\pgfqpoint{4.900000in}{4.900000in}}%
\pgfusepath{clip}%
\pgfsetbuttcap%
\pgfsetroundjoin%
\definecolor{currentfill}{rgb}{0.784667,0.784667,0.784667}%
\pgfsetfillcolor{currentfill}%
\pgfsetfillopacity{0.900000}%
\pgfsetlinewidth{0.501875pt}%
\definecolor{currentstroke}{rgb}{0.200000,0.200000,0.200000}%
\pgfsetstrokecolor{currentstroke}%
\pgfsetdash{}{0pt}%
\pgfpathmoveto{\pgfqpoint{2.198259in}{2.398984in}}%
\pgfpathlineto{\pgfqpoint{2.244071in}{2.355208in}}%
\pgfpathlineto{\pgfqpoint{2.275165in}{2.335089in}}%
\pgfpathlineto{\pgfqpoint{2.229348in}{2.378865in}}%
\pgfpathlineto{\pgfqpoint{2.198259in}{2.398984in}}%
\pgfpathclose%
\pgfusepath{stroke,fill}%
\end{pgfscope}%
\begin{pgfscope}%
\pgfpathrectangle{\pgfqpoint{0.050000in}{0.050000in}}{\pgfqpoint{4.900000in}{4.900000in}}%
\pgfusepath{clip}%
\pgfsetbuttcap%
\pgfsetroundjoin%
\definecolor{currentfill}{rgb}{0.915356,0.915356,0.915356}%
\pgfsetfillcolor{currentfill}%
\pgfsetfillopacity{0.900000}%
\pgfsetlinewidth{0.501875pt}%
\definecolor{currentstroke}{rgb}{0.200000,0.200000,0.200000}%
\pgfsetstrokecolor{currentstroke}%
\pgfsetdash{}{0pt}%
\pgfpathmoveto{\pgfqpoint{2.721915in}{1.989336in}}%
\pgfpathlineto{\pgfqpoint{2.766658in}{1.954552in}}%
\pgfpathlineto{\pgfqpoint{2.798561in}{1.935177in}}%
\pgfpathlineto{\pgfqpoint{2.753809in}{1.969663in}}%
\pgfpathlineto{\pgfqpoint{2.721915in}{1.989336in}}%
\pgfpathclose%
\pgfusepath{stroke,fill}%
\end{pgfscope}%
\begin{pgfscope}%
\pgfpathrectangle{\pgfqpoint{0.050000in}{0.050000in}}{\pgfqpoint{4.900000in}{4.900000in}}%
\pgfusepath{clip}%
\pgfsetbuttcap%
\pgfsetroundjoin%
\definecolor{currentfill}{rgb}{0.992618,0.992618,0.992618}%
\pgfsetfillcolor{currentfill}%
\pgfsetfillopacity{0.900000}%
\pgfsetlinewidth{0.501875pt}%
\definecolor{currentstroke}{rgb}{0.200000,0.200000,0.200000}%
\pgfsetstrokecolor{currentstroke}%
\pgfsetdash{}{0pt}%
\pgfpathmoveto{\pgfqpoint{3.565695in}{1.618632in}}%
\pgfpathlineto{\pgfqpoint{3.609521in}{1.624182in}}%
\pgfpathlineto{\pgfqpoint{3.642753in}{1.604609in}}%
\pgfpathlineto{\pgfqpoint{3.598902in}{1.599165in}}%
\pgfpathlineto{\pgfqpoint{3.565695in}{1.618632in}}%
\pgfpathclose%
\pgfusepath{stroke,fill}%
\end{pgfscope}%
\begin{pgfscope}%
\pgfpathrectangle{\pgfqpoint{0.050000in}{0.050000in}}{\pgfqpoint{4.900000in}{4.900000in}}%
\pgfusepath{clip}%
\pgfsetbuttcap%
\pgfsetroundjoin%
\definecolor{currentfill}{rgb}{0.424083,0.424083,0.424083}%
\pgfsetfillcolor{currentfill}%
\pgfsetfillopacity{0.900000}%
\pgfsetlinewidth{0.501875pt}%
\definecolor{currentstroke}{rgb}{0.200000,0.200000,0.200000}%
\pgfsetstrokecolor{currentstroke}%
\pgfsetdash{}{0pt}%
\pgfpathmoveto{\pgfqpoint{1.227533in}{3.184147in}}%
\pgfpathlineto{\pgfqpoint{1.275382in}{3.138638in}}%
\pgfpathlineto{\pgfqpoint{1.304997in}{3.119668in}}%
\pgfpathlineto{\pgfqpoint{1.257138in}{3.165207in}}%
\pgfpathlineto{\pgfqpoint{1.227533in}{3.184147in}}%
\pgfpathclose%
\pgfusepath{stroke,fill}%
\end{pgfscope}%
\begin{pgfscope}%
\pgfpathrectangle{\pgfqpoint{0.050000in}{0.050000in}}{\pgfqpoint{4.900000in}{4.900000in}}%
\pgfusepath{clip}%
\pgfsetbuttcap%
\pgfsetroundjoin%
\definecolor{currentfill}{rgb}{1.000000,1.000000,1.000000}%
\pgfsetfillcolor{currentfill}%
\pgfsetfillopacity{0.900000}%
\pgfsetlinewidth{0.501875pt}%
\definecolor{currentstroke}{rgb}{0.200000,0.200000,0.200000}%
\pgfsetstrokecolor{currentstroke}%
\pgfsetdash{}{0pt}%
\pgfpathmoveto{\pgfqpoint{4.116240in}{1.577404in}}%
\pgfpathlineto{\pgfqpoint{4.159585in}{1.589966in}}%
\pgfpathlineto{\pgfqpoint{4.193654in}{1.569388in}}%
\pgfpathlineto{\pgfqpoint{4.150283in}{1.556883in}}%
\pgfpathlineto{\pgfqpoint{4.116240in}{1.577404in}}%
\pgfpathclose%
\pgfusepath{stroke,fill}%
\end{pgfscope}%
\begin{pgfscope}%
\pgfpathrectangle{\pgfqpoint{0.050000in}{0.050000in}}{\pgfqpoint{4.900000in}{4.900000in}}%
\pgfusepath{clip}%
\pgfsetbuttcap%
\pgfsetroundjoin%
\definecolor{currentfill}{rgb}{0.691396,0.691396,0.691396}%
\pgfsetfillcolor{currentfill}%
\pgfsetfillopacity{0.900000}%
\pgfsetlinewidth{0.501875pt}%
\definecolor{currentstroke}{rgb}{0.200000,0.200000,0.200000}%
\pgfsetstrokecolor{currentstroke}%
\pgfsetdash{}{0pt}%
\pgfpathmoveto{\pgfqpoint{1.936321in}{2.611268in}}%
\pgfpathlineto{\pgfqpoint{1.982689in}{2.566955in}}%
\pgfpathlineto{\pgfqpoint{2.013389in}{2.547112in}}%
\pgfpathlineto{\pgfqpoint{1.967014in}{2.591451in}}%
\pgfpathlineto{\pgfqpoint{1.936321in}{2.611268in}}%
\pgfpathclose%
\pgfusepath{stroke,fill}%
\end{pgfscope}%
\begin{pgfscope}%
\pgfpathrectangle{\pgfqpoint{0.050000in}{0.050000in}}{\pgfqpoint{4.900000in}{4.900000in}}%
\pgfusepath{clip}%
\pgfsetbuttcap%
\pgfsetroundjoin%
\definecolor{currentfill}{rgb}{0.581776,0.581776,0.581776}%
\pgfsetfillcolor{currentfill}%
\pgfsetfillopacity{0.900000}%
\pgfsetlinewidth{0.501875pt}%
\definecolor{currentstroke}{rgb}{0.200000,0.200000,0.200000}%
\pgfsetstrokecolor{currentstroke}%
\pgfsetdash{}{0pt}%
\pgfpathmoveto{\pgfqpoint{1.674295in}{2.823671in}}%
\pgfpathlineto{\pgfqpoint{1.721219in}{2.778907in}}%
\pgfpathlineto{\pgfqpoint{1.751524in}{2.759383in}}%
\pgfpathlineto{\pgfqpoint{1.704591in}{2.804176in}}%
\pgfpathlineto{\pgfqpoint{1.674295in}{2.823671in}}%
\pgfpathclose%
\pgfusepath{stroke,fill}%
\end{pgfscope}%
\begin{pgfscope}%
\pgfpathrectangle{\pgfqpoint{0.050000in}{0.050000in}}{\pgfqpoint{4.900000in}{4.900000in}}%
\pgfusepath{clip}%
\pgfsetbuttcap%
\pgfsetroundjoin%
\definecolor{currentfill}{rgb}{0.994464,0.994464,0.994464}%
\pgfsetfillcolor{currentfill}%
\pgfsetfillopacity{0.900000}%
\pgfsetlinewidth{0.501875pt}%
\definecolor{currentstroke}{rgb}{0.200000,0.200000,0.200000}%
\pgfsetstrokecolor{currentstroke}%
\pgfsetdash{}{0pt}%
\pgfpathmoveto{\pgfqpoint{3.642753in}{1.604609in}}%
\pgfpathlineto{\pgfqpoint{3.686534in}{1.611965in}}%
\pgfpathlineto{\pgfqpoint{3.719895in}{1.592221in}}%
\pgfpathlineto{\pgfqpoint{3.676089in}{1.584969in}}%
\pgfpathlineto{\pgfqpoint{3.642753in}{1.604609in}}%
\pgfpathclose%
\pgfusepath{stroke,fill}%
\end{pgfscope}%
\begin{pgfscope}%
\pgfpathrectangle{\pgfqpoint{0.050000in}{0.050000in}}{\pgfqpoint{4.900000in}{4.900000in}}%
\pgfusepath{clip}%
\pgfsetbuttcap%
\pgfsetroundjoin%
\definecolor{currentfill}{rgb}{1.000000,1.000000,1.000000}%
\pgfsetfillcolor{currentfill}%
\pgfsetfillopacity{0.900000}%
\pgfsetlinewidth{0.501875pt}%
\definecolor{currentstroke}{rgb}{0.200000,0.200000,0.200000}%
\pgfsetstrokecolor{currentstroke}%
\pgfsetdash{}{0pt}%
\pgfpathmoveto{\pgfqpoint{3.918063in}{1.581957in}}%
\pgfpathlineto{\pgfqpoint{3.961618in}{1.593592in}}%
\pgfpathlineto{\pgfqpoint{3.995397in}{1.573320in}}%
\pgfpathlineto{\pgfqpoint{3.951816in}{1.561758in}}%
\pgfpathlineto{\pgfqpoint{3.918063in}{1.581957in}}%
\pgfpathclose%
\pgfusepath{stroke,fill}%
\end{pgfscope}%
\begin{pgfscope}%
\pgfpathrectangle{\pgfqpoint{0.050000in}{0.050000in}}{\pgfqpoint{4.900000in}{4.900000in}}%
\pgfusepath{clip}%
\pgfsetbuttcap%
\pgfsetroundjoin%
\definecolor{currentfill}{rgb}{0.836340,0.836340,0.836340}%
\pgfsetfillcolor{currentfill}%
\pgfsetfillopacity{0.900000}%
\pgfsetlinewidth{0.501875pt}%
\definecolor{currentstroke}{rgb}{0.200000,0.200000,0.200000}%
\pgfsetstrokecolor{currentstroke}%
\pgfsetdash{}{0pt}%
\pgfpathmoveto{\pgfqpoint{2.383274in}{2.251331in}}%
\pgfpathlineto{\pgfqpoint{2.428717in}{2.208537in}}%
\pgfpathlineto{\pgfqpoint{2.460107in}{2.188405in}}%
\pgfpathlineto{\pgfqpoint{2.414660in}{2.231106in}}%
\pgfpathlineto{\pgfqpoint{2.383274in}{2.251331in}}%
\pgfpathclose%
\pgfusepath{stroke,fill}%
\end{pgfscope}%
\begin{pgfscope}%
\pgfpathrectangle{\pgfqpoint{0.050000in}{0.050000in}}{\pgfqpoint{4.900000in}{4.900000in}}%
\pgfusepath{clip}%
\pgfsetbuttcap%
\pgfsetroundjoin%
\definecolor{currentfill}{rgb}{0.898378,0.898378,0.898378}%
\pgfsetfillcolor{currentfill}%
\pgfsetfillopacity{0.900000}%
\pgfsetlinewidth{0.501875pt}%
\definecolor{currentstroke}{rgb}{0.200000,0.200000,0.200000}%
\pgfsetstrokecolor{currentstroke}%
\pgfsetdash{}{0pt}%
\pgfpathmoveto{\pgfqpoint{2.645213in}{2.046695in}}%
\pgfpathlineto{\pgfqpoint{2.690120in}{2.008964in}}%
\pgfpathlineto{\pgfqpoint{2.721915in}{1.989336in}}%
\pgfpathlineto{\pgfqpoint{2.677002in}{2.026795in}}%
\pgfpathlineto{\pgfqpoint{2.645213in}{2.046695in}}%
\pgfpathclose%
\pgfusepath{stroke,fill}%
\end{pgfscope}%
\begin{pgfscope}%
\pgfpathrectangle{\pgfqpoint{0.050000in}{0.050000in}}{\pgfqpoint{4.900000in}{4.900000in}}%
\pgfusepath{clip}%
\pgfsetbuttcap%
\pgfsetroundjoin%
\definecolor{currentfill}{rgb}{0.487043,0.487043,0.487043}%
\pgfsetfillcolor{currentfill}%
\pgfsetfillopacity{0.900000}%
\pgfsetlinewidth{0.501875pt}%
\definecolor{currentstroke}{rgb}{0.200000,0.200000,0.200000}%
\pgfsetstrokecolor{currentstroke}%
\pgfsetdash{}{0pt}%
\pgfpathmoveto{\pgfqpoint{1.412182in}{3.036148in}}%
\pgfpathlineto{\pgfqpoint{1.459663in}{2.990931in}}%
\pgfpathlineto{\pgfqpoint{1.489572in}{2.971728in}}%
\pgfpathlineto{\pgfqpoint{1.442081in}{3.016974in}}%
\pgfpathlineto{\pgfqpoint{1.412182in}{3.036148in}}%
\pgfpathclose%
\pgfusepath{stroke,fill}%
\end{pgfscope}%
\begin{pgfscope}%
\pgfpathrectangle{\pgfqpoint{0.050000in}{0.050000in}}{\pgfqpoint{4.900000in}{4.900000in}}%
\pgfusepath{clip}%
\pgfsetbuttcap%
\pgfsetroundjoin%
\definecolor{currentfill}{rgb}{0.970473,0.970473,0.970473}%
\pgfsetfillcolor{currentfill}%
\pgfsetfillopacity{0.900000}%
\pgfsetlinewidth{0.501875pt}%
\definecolor{currentstroke}{rgb}{0.200000,0.200000,0.200000}%
\pgfsetstrokecolor{currentstroke}%
\pgfsetdash{}{0pt}%
\pgfpathmoveto{\pgfqpoint{3.137398in}{1.741331in}}%
\pgfpathlineto{\pgfqpoint{3.181573in}{1.727567in}}%
\pgfpathlineto{\pgfqpoint{3.214154in}{1.708746in}}%
\pgfpathlineto{\pgfqpoint{3.169960in}{1.722424in}}%
\pgfpathlineto{\pgfqpoint{3.137398in}{1.741331in}}%
\pgfpathclose%
\pgfusepath{stroke,fill}%
\end{pgfscope}%
\begin{pgfscope}%
\pgfpathrectangle{\pgfqpoint{0.050000in}{0.050000in}}{\pgfqpoint{4.900000in}{4.900000in}}%
\pgfusepath{clip}%
\pgfsetbuttcap%
\pgfsetroundjoin%
\definecolor{currentfill}{rgb}{0.757109,0.757109,0.757109}%
\pgfsetfillcolor{currentfill}%
\pgfsetfillopacity{0.900000}%
\pgfsetlinewidth{0.501875pt}%
\definecolor{currentstroke}{rgb}{0.200000,0.200000,0.200000}%
\pgfsetstrokecolor{currentstroke}%
\pgfsetdash{}{0pt}%
\pgfpathmoveto{\pgfqpoint{2.121264in}{2.463046in}}%
\pgfpathlineto{\pgfqpoint{2.167264in}{2.419047in}}%
\pgfpathlineto{\pgfqpoint{2.198259in}{2.398984in}}%
\pgfpathlineto{\pgfqpoint{2.152254in}{2.442999in}}%
\pgfpathlineto{\pgfqpoint{2.121264in}{2.463046in}}%
\pgfpathclose%
\pgfusepath{stroke,fill}%
\end{pgfscope}%
\begin{pgfscope}%
\pgfpathrectangle{\pgfqpoint{0.050000in}{0.050000in}}{\pgfqpoint{4.900000in}{4.900000in}}%
\pgfusepath{clip}%
\pgfsetbuttcap%
\pgfsetroundjoin%
\definecolor{currentfill}{rgb}{0.976009,0.976009,0.976009}%
\pgfsetfillcolor{currentfill}%
\pgfsetfillopacity{0.900000}%
\pgfsetlinewidth{0.501875pt}%
\definecolor{currentstroke}{rgb}{0.200000,0.200000,0.200000}%
\pgfsetstrokecolor{currentstroke}%
\pgfsetdash{}{0pt}%
\pgfpathmoveto{\pgfqpoint{3.214154in}{1.708746in}}%
\pgfpathlineto{\pgfqpoint{3.258259in}{1.699134in}}%
\pgfpathlineto{\pgfqpoint{3.290963in}{1.680251in}}%
\pgfpathlineto{\pgfqpoint{3.246838in}{1.689835in}}%
\pgfpathlineto{\pgfqpoint{3.214154in}{1.708746in}}%
\pgfpathclose%
\pgfusepath{stroke,fill}%
\end{pgfscope}%
\begin{pgfscope}%
\pgfpathrectangle{\pgfqpoint{0.050000in}{0.050000in}}{\pgfqpoint{4.900000in}{4.900000in}}%
\pgfusepath{clip}%
\pgfsetbuttcap%
\pgfsetroundjoin%
\definecolor{currentfill}{rgb}{0.963091,0.963091,0.963091}%
\pgfsetfillcolor{currentfill}%
\pgfsetfillopacity{0.900000}%
\pgfsetlinewidth{0.501875pt}%
\definecolor{currentstroke}{rgb}{0.200000,0.200000,0.200000}%
\pgfsetstrokecolor{currentstroke}%
\pgfsetdash{}{0pt}%
\pgfpathmoveto{\pgfqpoint{3.060678in}{1.778279in}}%
\pgfpathlineto{\pgfqpoint{3.104937in}{1.760149in}}%
\pgfpathlineto{\pgfqpoint{3.137398in}{1.741331in}}%
\pgfpathlineto{\pgfqpoint{3.093121in}{1.759313in}}%
\pgfpathlineto{\pgfqpoint{3.060678in}{1.778279in}}%
\pgfpathclose%
\pgfusepath{stroke,fill}%
\end{pgfscope}%
\begin{pgfscope}%
\pgfpathrectangle{\pgfqpoint{0.050000in}{0.050000in}}{\pgfqpoint{4.900000in}{4.900000in}}%
\pgfusepath{clip}%
\pgfsetbuttcap%
\pgfsetroundjoin%
\definecolor{currentfill}{rgb}{0.981546,0.981546,0.981546}%
\pgfsetfillcolor{currentfill}%
\pgfsetfillopacity{0.900000}%
\pgfsetlinewidth{0.501875pt}%
\definecolor{currentstroke}{rgb}{0.200000,0.200000,0.200000}%
\pgfsetstrokecolor{currentstroke}%
\pgfsetdash{}{0pt}%
\pgfpathmoveto{\pgfqpoint{3.290963in}{1.680251in}}%
\pgfpathlineto{\pgfqpoint{3.335006in}{1.674468in}}%
\pgfpathlineto{\pgfqpoint{3.367834in}{1.655477in}}%
\pgfpathlineto{\pgfqpoint{3.323769in}{1.661281in}}%
\pgfpathlineto{\pgfqpoint{3.290963in}{1.680251in}}%
\pgfpathclose%
\pgfusepath{stroke,fill}%
\end{pgfscope}%
\begin{pgfscope}%
\pgfpathrectangle{\pgfqpoint{0.050000in}{0.050000in}}{\pgfqpoint{4.900000in}{4.900000in}}%
\pgfusepath{clip}%
\pgfsetbuttcap%
\pgfsetroundjoin%
\definecolor{currentfill}{rgb}{0.955709,0.955709,0.955709}%
\pgfsetfillcolor{currentfill}%
\pgfsetfillopacity{0.900000}%
\pgfsetlinewidth{0.501875pt}%
\definecolor{currentstroke}{rgb}{0.200000,0.200000,0.200000}%
\pgfsetstrokecolor{currentstroke}%
\pgfsetdash{}{0pt}%
\pgfpathmoveto{\pgfqpoint{2.983978in}{1.819725in}}%
\pgfpathlineto{\pgfqpoint{3.028336in}{1.797161in}}%
\pgfpathlineto{\pgfqpoint{3.060678in}{1.778279in}}%
\pgfpathlineto{\pgfqpoint{3.016305in}{1.800636in}}%
\pgfpathlineto{\pgfqpoint{2.983978in}{1.819725in}}%
\pgfpathclose%
\pgfusepath{stroke,fill}%
\end{pgfscope}%
\begin{pgfscope}%
\pgfpathrectangle{\pgfqpoint{0.050000in}{0.050000in}}{\pgfqpoint{4.900000in}{4.900000in}}%
\pgfusepath{clip}%
\pgfsetbuttcap%
\pgfsetroundjoin%
\definecolor{currentfill}{rgb}{0.653010,0.653010,0.653010}%
\pgfsetfillcolor{currentfill}%
\pgfsetfillopacity{0.900000}%
\pgfsetlinewidth{0.501875pt}%
\definecolor{currentstroke}{rgb}{0.200000,0.200000,0.200000}%
\pgfsetstrokecolor{currentstroke}%
\pgfsetdash{}{0pt}%
\pgfpathmoveto{\pgfqpoint{1.859165in}{2.675498in}}%
\pgfpathlineto{\pgfqpoint{1.905721in}{2.631025in}}%
\pgfpathlineto{\pgfqpoint{1.936321in}{2.611268in}}%
\pgfpathlineto{\pgfqpoint{1.889757in}{2.655767in}}%
\pgfpathlineto{\pgfqpoint{1.859165in}{2.675498in}}%
\pgfpathclose%
\pgfusepath{stroke,fill}%
\end{pgfscope}%
\begin{pgfscope}%
\pgfpathrectangle{\pgfqpoint{0.050000in}{0.050000in}}{\pgfqpoint{4.900000in}{4.900000in}}%
\pgfusepath{clip}%
\pgfsetbuttcap%
\pgfsetroundjoin%
\definecolor{currentfill}{rgb}{0.996309,0.996309,0.996309}%
\pgfsetfillcolor{currentfill}%
\pgfsetfillopacity{0.900000}%
\pgfsetlinewidth{0.501875pt}%
\definecolor{currentstroke}{rgb}{0.200000,0.200000,0.200000}%
\pgfsetstrokecolor{currentstroke}%
\pgfsetdash{}{0pt}%
\pgfpathmoveto{\pgfqpoint{3.719895in}{1.592221in}}%
\pgfpathlineto{\pgfqpoint{3.763630in}{1.601043in}}%
\pgfpathlineto{\pgfqpoint{3.797121in}{1.581138in}}%
\pgfpathlineto{\pgfqpoint{3.753361in}{1.572414in}}%
\pgfpathlineto{\pgfqpoint{3.719895in}{1.592221in}}%
\pgfpathclose%
\pgfusepath{stroke,fill}%
\end{pgfscope}%
\begin{pgfscope}%
\pgfpathrectangle{\pgfqpoint{0.050000in}{0.050000in}}{\pgfqpoint{4.900000in}{4.900000in}}%
\pgfusepath{clip}%
\pgfsetbuttcap%
\pgfsetroundjoin%
\definecolor{currentfill}{rgb}{1.000000,1.000000,1.000000}%
\pgfsetfillcolor{currentfill}%
\pgfsetfillopacity{0.900000}%
\pgfsetlinewidth{0.501875pt}%
\definecolor{currentstroke}{rgb}{0.200000,0.200000,0.200000}%
\pgfsetstrokecolor{currentstroke}%
\pgfsetdash{}{0pt}%
\pgfpathmoveto{\pgfqpoint{3.995397in}{1.573320in}}%
\pgfpathlineto{\pgfqpoint{4.038894in}{1.585419in}}%
\pgfpathlineto{\pgfqpoint{4.072804in}{1.565021in}}%
\pgfpathlineto{\pgfqpoint{4.029281in}{1.552988in}}%
\pgfpathlineto{\pgfqpoint{3.995397in}{1.573320in}}%
\pgfpathclose%
\pgfusepath{stroke,fill}%
\end{pgfscope}%
\begin{pgfscope}%
\pgfpathrectangle{\pgfqpoint{0.050000in}{0.050000in}}{\pgfqpoint{4.900000in}{4.900000in}}%
\pgfusepath{clip}%
\pgfsetbuttcap%
\pgfsetroundjoin%
\definecolor{currentfill}{rgb}{0.985236,0.985236,0.985236}%
\pgfsetfillcolor{currentfill}%
\pgfsetfillopacity{0.900000}%
\pgfsetlinewidth{0.501875pt}%
\definecolor{currentstroke}{rgb}{0.200000,0.200000,0.200000}%
\pgfsetstrokecolor{currentstroke}%
\pgfsetdash{}{0pt}%
\pgfpathmoveto{\pgfqpoint{3.367834in}{1.655477in}}%
\pgfpathlineto{\pgfqpoint{3.411824in}{1.653131in}}%
\pgfpathlineto{\pgfqpoint{3.444777in}{1.633998in}}%
\pgfpathlineto{\pgfqpoint{3.400765in}{1.636402in}}%
\pgfpathlineto{\pgfqpoint{3.367834in}{1.655477in}}%
\pgfpathclose%
\pgfusepath{stroke,fill}%
\end{pgfscope}%
\begin{pgfscope}%
\pgfpathrectangle{\pgfqpoint{0.050000in}{0.050000in}}{\pgfqpoint{4.900000in}{4.900000in}}%
\pgfusepath{clip}%
\pgfsetbuttcap%
\pgfsetroundjoin%
\definecolor{currentfill}{rgb}{0.551634,0.551634,0.551634}%
\pgfsetfillcolor{currentfill}%
\pgfsetfillopacity{0.900000}%
\pgfsetlinewidth{0.501875pt}%
\definecolor{currentstroke}{rgb}{0.200000,0.200000,0.200000}%
\pgfsetstrokecolor{currentstroke}%
\pgfsetdash{}{0pt}%
\pgfpathmoveto{\pgfqpoint{1.596978in}{2.888033in}}%
\pgfpathlineto{\pgfqpoint{1.644091in}{2.843108in}}%
\pgfpathlineto{\pgfqpoint{1.674295in}{2.823671in}}%
\pgfpathlineto{\pgfqpoint{1.627173in}{2.868625in}}%
\pgfpathlineto{\pgfqpoint{1.596978in}{2.888033in}}%
\pgfpathclose%
\pgfusepath{stroke,fill}%
\end{pgfscope}%
\begin{pgfscope}%
\pgfpathrectangle{\pgfqpoint{0.050000in}{0.050000in}}{\pgfqpoint{4.900000in}{4.900000in}}%
\pgfusepath{clip}%
\pgfsetbuttcap%
\pgfsetroundjoin%
\definecolor{currentfill}{rgb}{0.878570,0.878570,0.878570}%
\pgfsetfillcolor{currentfill}%
\pgfsetfillopacity{0.900000}%
\pgfsetlinewidth{0.501875pt}%
\definecolor{currentstroke}{rgb}{0.200000,0.200000,0.200000}%
\pgfsetstrokecolor{currentstroke}%
\pgfsetdash{}{0pt}%
\pgfpathmoveto{\pgfqpoint{2.568442in}{2.106560in}}%
\pgfpathlineto{\pgfqpoint{2.613523in}{2.066556in}}%
\pgfpathlineto{\pgfqpoint{2.645213in}{2.046695in}}%
\pgfpathlineto{\pgfqpoint{2.600128in}{2.086475in}}%
\pgfpathlineto{\pgfqpoint{2.568442in}{2.106560in}}%
\pgfpathclose%
\pgfusepath{stroke,fill}%
\end{pgfscope}%
\begin{pgfscope}%
\pgfpathrectangle{\pgfqpoint{0.050000in}{0.050000in}}{\pgfqpoint{4.900000in}{4.900000in}}%
\pgfusepath{clip}%
\pgfsetbuttcap%
\pgfsetroundjoin%
\definecolor{currentfill}{rgb}{0.946482,0.946482,0.946482}%
\pgfsetfillcolor{currentfill}%
\pgfsetfillopacity{0.900000}%
\pgfsetlinewidth{0.501875pt}%
\definecolor{currentstroke}{rgb}{0.200000,0.200000,0.200000}%
\pgfsetstrokecolor{currentstroke}%
\pgfsetdash{}{0pt}%
\pgfpathmoveto{\pgfqpoint{2.907279in}{1.865626in}}%
\pgfpathlineto{\pgfqpoint{2.951752in}{1.838737in}}%
\pgfpathlineto{\pgfqpoint{2.983978in}{1.819725in}}%
\pgfpathlineto{\pgfqpoint{2.939492in}{1.846358in}}%
\pgfpathlineto{\pgfqpoint{2.907279in}{1.865626in}}%
\pgfpathclose%
\pgfusepath{stroke,fill}%
\end{pgfscope}%
\begin{pgfscope}%
\pgfpathrectangle{\pgfqpoint{0.050000in}{0.050000in}}{\pgfqpoint{4.900000in}{4.900000in}}%
\pgfusepath{clip}%
\pgfsetbuttcap%
\pgfsetroundjoin%
\definecolor{currentfill}{rgb}{0.808781,0.808781,0.808781}%
\pgfsetfillcolor{currentfill}%
\pgfsetfillopacity{0.900000}%
\pgfsetlinewidth{0.501875pt}%
\definecolor{currentstroke}{rgb}{0.200000,0.200000,0.200000}%
\pgfsetstrokecolor{currentstroke}%
\pgfsetdash{}{0pt}%
\pgfpathmoveto{\pgfqpoint{2.306355in}{2.314919in}}%
\pgfpathlineto{\pgfqpoint{2.351985in}{2.271512in}}%
\pgfpathlineto{\pgfqpoint{2.383274in}{2.251331in}}%
\pgfpathlineto{\pgfqpoint{2.337641in}{2.294698in}}%
\pgfpathlineto{\pgfqpoint{2.306355in}{2.314919in}}%
\pgfpathclose%
\pgfusepath{stroke,fill}%
\end{pgfscope}%
\begin{pgfscope}%
\pgfpathrectangle{\pgfqpoint{0.050000in}{0.050000in}}{\pgfqpoint{4.900000in}{4.900000in}}%
\pgfusepath{clip}%
\pgfsetbuttcap%
\pgfsetroundjoin%
\definecolor{currentfill}{rgb}{0.988927,0.988927,0.988927}%
\pgfsetfillcolor{currentfill}%
\pgfsetfillopacity{0.900000}%
\pgfsetlinewidth{0.501875pt}%
\definecolor{currentstroke}{rgb}{0.200000,0.200000,0.200000}%
\pgfsetstrokecolor{currentstroke}%
\pgfsetdash{}{0pt}%
\pgfpathmoveto{\pgfqpoint{3.444777in}{1.633998in}}%
\pgfpathlineto{\pgfqpoint{3.488719in}{1.634668in}}%
\pgfpathlineto{\pgfqpoint{3.521799in}{1.615373in}}%
\pgfpathlineto{\pgfqpoint{3.477834in}{1.614786in}}%
\pgfpathlineto{\pgfqpoint{3.444777in}{1.633998in}}%
\pgfpathclose%
\pgfusepath{stroke,fill}%
\end{pgfscope}%
\begin{pgfscope}%
\pgfpathrectangle{\pgfqpoint{0.050000in}{0.050000in}}{\pgfqpoint{4.900000in}{4.900000in}}%
\pgfusepath{clip}%
\pgfsetbuttcap%
\pgfsetroundjoin%
\definecolor{currentfill}{rgb}{0.452595,0.452595,0.452595}%
\pgfsetfillcolor{currentfill}%
\pgfsetfillopacity{0.900000}%
\pgfsetlinewidth{0.501875pt}%
\definecolor{currentstroke}{rgb}{0.200000,0.200000,0.200000}%
\pgfsetstrokecolor{currentstroke}%
\pgfsetdash{}{0pt}%
\pgfpathmoveto{\pgfqpoint{1.334704in}{3.100641in}}%
\pgfpathlineto{\pgfqpoint{1.382374in}{3.055263in}}%
\pgfpathlineto{\pgfqpoint{1.412182in}{3.036148in}}%
\pgfpathlineto{\pgfqpoint{1.364501in}{3.081555in}}%
\pgfpathlineto{\pgfqpoint{1.334704in}{3.100641in}}%
\pgfpathclose%
\pgfusepath{stroke,fill}%
\end{pgfscope}%
\begin{pgfscope}%
\pgfpathrectangle{\pgfqpoint{0.050000in}{0.050000in}}{\pgfqpoint{4.900000in}{4.900000in}}%
\pgfusepath{clip}%
\pgfsetbuttcap%
\pgfsetroundjoin%
\definecolor{currentfill}{rgb}{0.998155,0.998155,0.998155}%
\pgfsetfillcolor{currentfill}%
\pgfsetfillopacity{0.900000}%
\pgfsetlinewidth{0.501875pt}%
\definecolor{currentstroke}{rgb}{0.200000,0.200000,0.200000}%
\pgfsetstrokecolor{currentstroke}%
\pgfsetdash{}{0pt}%
\pgfpathmoveto{\pgfqpoint{3.797121in}{1.581138in}}%
\pgfpathlineto{\pgfqpoint{3.840807in}{1.591125in}}%
\pgfpathlineto{\pgfqpoint{3.874428in}{1.571068in}}%
\pgfpathlineto{\pgfqpoint{3.830716in}{1.561171in}}%
\pgfpathlineto{\pgfqpoint{3.797121in}{1.581138in}}%
\pgfpathclose%
\pgfusepath{stroke,fill}%
\end{pgfscope}%
\begin{pgfscope}%
\pgfpathrectangle{\pgfqpoint{0.050000in}{0.050000in}}{\pgfqpoint{4.900000in}{4.900000in}}%
\pgfusepath{clip}%
\pgfsetbuttcap%
\pgfsetroundjoin%
\definecolor{currentfill}{rgb}{0.729781,0.729781,0.729781}%
\pgfsetfillcolor{currentfill}%
\pgfsetfillopacity{0.900000}%
\pgfsetlinewidth{0.501875pt}%
\definecolor{currentstroke}{rgb}{0.200000,0.200000,0.200000}%
\pgfsetstrokecolor{currentstroke}%
\pgfsetdash{}{0pt}%
\pgfpathmoveto{\pgfqpoint{2.044183in}{2.527208in}}%
\pgfpathlineto{\pgfqpoint{2.090370in}{2.483033in}}%
\pgfpathlineto{\pgfqpoint{2.121264in}{2.463046in}}%
\pgfpathlineto{\pgfqpoint{2.075071in}{2.507244in}}%
\pgfpathlineto{\pgfqpoint{2.044183in}{2.527208in}}%
\pgfpathclose%
\pgfusepath{stroke,fill}%
\end{pgfscope}%
\begin{pgfscope}%
\pgfpathrectangle{\pgfqpoint{0.050000in}{0.050000in}}{\pgfqpoint{4.900000in}{4.900000in}}%
\pgfusepath{clip}%
\pgfsetbuttcap%
\pgfsetroundjoin%
\definecolor{currentfill}{rgb}{0.932334,0.932334,0.932334}%
\pgfsetfillcolor{currentfill}%
\pgfsetfillopacity{0.900000}%
\pgfsetlinewidth{0.501875pt}%
\definecolor{currentstroke}{rgb}{0.200000,0.200000,0.200000}%
\pgfsetstrokecolor{currentstroke}%
\pgfsetdash{}{0pt}%
\pgfpathmoveto{\pgfqpoint{2.830563in}{1.915747in}}%
\pgfpathlineto{\pgfqpoint{2.875167in}{1.884827in}}%
\pgfpathlineto{\pgfqpoint{2.907279in}{1.865626in}}%
\pgfpathlineto{\pgfqpoint{2.862664in}{1.896261in}}%
\pgfpathlineto{\pgfqpoint{2.830563in}{1.915747in}}%
\pgfpathclose%
\pgfusepath{stroke,fill}%
\end{pgfscope}%
\begin{pgfscope}%
\pgfpathrectangle{\pgfqpoint{0.050000in}{0.050000in}}{\pgfqpoint{4.900000in}{4.900000in}}%
\pgfusepath{clip}%
\pgfsetbuttcap%
\pgfsetroundjoin%
\definecolor{currentfill}{rgb}{1.000000,1.000000,1.000000}%
\pgfsetfillcolor{currentfill}%
\pgfsetfillopacity{0.900000}%
\pgfsetlinewidth{0.501875pt}%
\definecolor{currentstroke}{rgb}{0.200000,0.200000,0.200000}%
\pgfsetstrokecolor{currentstroke}%
\pgfsetdash{}{0pt}%
\pgfpathmoveto{\pgfqpoint{4.072804in}{1.565021in}}%
\pgfpathlineto{\pgfqpoint{4.116240in}{1.577404in}}%
\pgfpathlineto{\pgfqpoint{4.150283in}{1.556883in}}%
\pgfpathlineto{\pgfqpoint{4.106820in}{1.544561in}}%
\pgfpathlineto{\pgfqpoint{4.072804in}{1.565021in}}%
\pgfpathclose%
\pgfusepath{stroke,fill}%
\end{pgfscope}%
\begin{pgfscope}%
\pgfpathrectangle{\pgfqpoint{0.050000in}{0.050000in}}{\pgfqpoint{4.900000in}{4.900000in}}%
\pgfusepath{clip}%
\pgfsetbuttcap%
\pgfsetroundjoin%
\definecolor{currentfill}{rgb}{0.619423,0.619423,0.619423}%
\pgfsetfillcolor{currentfill}%
\pgfsetfillopacity{0.900000}%
\pgfsetlinewidth{0.501875pt}%
\definecolor{currentstroke}{rgb}{0.200000,0.200000,0.200000}%
\pgfsetstrokecolor{currentstroke}%
\pgfsetdash{}{0pt}%
\pgfpathmoveto{\pgfqpoint{1.781922in}{2.739800in}}%
\pgfpathlineto{\pgfqpoint{1.828666in}{2.695168in}}%
\pgfpathlineto{\pgfqpoint{1.859165in}{2.675498in}}%
\pgfpathlineto{\pgfqpoint{1.812413in}{2.720157in}}%
\pgfpathlineto{\pgfqpoint{1.781922in}{2.739800in}}%
\pgfpathclose%
\pgfusepath{stroke,fill}%
\end{pgfscope}%
\begin{pgfscope}%
\pgfpathrectangle{\pgfqpoint{0.050000in}{0.050000in}}{\pgfqpoint{4.900000in}{4.900000in}}%
\pgfusepath{clip}%
\pgfsetbuttcap%
\pgfsetroundjoin%
\definecolor{currentfill}{rgb}{0.990773,0.990773,0.990773}%
\pgfsetfillcolor{currentfill}%
\pgfsetfillopacity{0.900000}%
\pgfsetlinewidth{0.501875pt}%
\definecolor{currentstroke}{rgb}{0.200000,0.200000,0.200000}%
\pgfsetstrokecolor{currentstroke}%
\pgfsetdash{}{0pt}%
\pgfpathmoveto{\pgfqpoint{3.521799in}{1.615373in}}%
\pgfpathlineto{\pgfqpoint{3.565695in}{1.618632in}}%
\pgfpathlineto{\pgfqpoint{3.598902in}{1.599165in}}%
\pgfpathlineto{\pgfqpoint{3.554982in}{1.596003in}}%
\pgfpathlineto{\pgfqpoint{3.521799in}{1.615373in}}%
\pgfpathclose%
\pgfusepath{stroke,fill}%
\end{pgfscope}%
\begin{pgfscope}%
\pgfpathrectangle{\pgfqpoint{0.050000in}{0.050000in}}{\pgfqpoint{4.900000in}{4.900000in}}%
\pgfusepath{clip}%
\pgfsetbuttcap%
\pgfsetroundjoin%
\definecolor{currentfill}{rgb}{0.858762,0.858762,0.858762}%
\pgfsetfillcolor{currentfill}%
\pgfsetfillopacity{0.900000}%
\pgfsetlinewidth{0.501875pt}%
\definecolor{currentstroke}{rgb}{0.200000,0.200000,0.200000}%
\pgfsetstrokecolor{currentstroke}%
\pgfsetdash{}{0pt}%
\pgfpathmoveto{\pgfqpoint{2.491593in}{2.168235in}}%
\pgfpathlineto{\pgfqpoint{2.536854in}{2.126608in}}%
\pgfpathlineto{\pgfqpoint{2.568442in}{2.106560in}}%
\pgfpathlineto{\pgfqpoint{2.523177in}{2.148024in}}%
\pgfpathlineto{\pgfqpoint{2.491593in}{2.168235in}}%
\pgfpathclose%
\pgfusepath{stroke,fill}%
\end{pgfscope}%
\begin{pgfscope}%
\pgfpathrectangle{\pgfqpoint{0.050000in}{0.050000in}}{\pgfqpoint{4.900000in}{4.900000in}}%
\pgfusepath{clip}%
\pgfsetbuttcap%
\pgfsetroundjoin%
\definecolor{currentfill}{rgb}{0.517186,0.517186,0.517186}%
\pgfsetfillcolor{currentfill}%
\pgfsetfillopacity{0.900000}%
\pgfsetlinewidth{0.501875pt}%
\definecolor{currentstroke}{rgb}{0.200000,0.200000,0.200000}%
\pgfsetstrokecolor{currentstroke}%
\pgfsetdash{}{0pt}%
\pgfpathmoveto{\pgfqpoint{1.519574in}{2.952467in}}%
\pgfpathlineto{\pgfqpoint{1.566875in}{2.907381in}}%
\pgfpathlineto{\pgfqpoint{1.596978in}{2.888033in}}%
\pgfpathlineto{\pgfqpoint{1.549667in}{2.933147in}}%
\pgfpathlineto{\pgfqpoint{1.519574in}{2.952467in}}%
\pgfpathclose%
\pgfusepath{stroke,fill}%
\end{pgfscope}%
\begin{pgfscope}%
\pgfpathrectangle{\pgfqpoint{0.050000in}{0.050000in}}{\pgfqpoint{4.900000in}{4.900000in}}%
\pgfusepath{clip}%
\pgfsetbuttcap%
\pgfsetroundjoin%
\definecolor{currentfill}{rgb}{0.998155,0.998155,0.998155}%
\pgfsetfillcolor{currentfill}%
\pgfsetfillopacity{0.900000}%
\pgfsetlinewidth{0.501875pt}%
\definecolor{currentstroke}{rgb}{0.200000,0.200000,0.200000}%
\pgfsetstrokecolor{currentstroke}%
\pgfsetdash{}{0pt}%
\pgfpathmoveto{\pgfqpoint{3.874428in}{1.571068in}}%
\pgfpathlineto{\pgfqpoint{3.918063in}{1.581957in}}%
\pgfpathlineto{\pgfqpoint{3.951816in}{1.561758in}}%
\pgfpathlineto{\pgfqpoint{3.908155in}{1.550950in}}%
\pgfpathlineto{\pgfqpoint{3.874428in}{1.571068in}}%
\pgfpathclose%
\pgfusepath{stroke,fill}%
\end{pgfscope}%
\begin{pgfscope}%
\pgfpathrectangle{\pgfqpoint{0.050000in}{0.050000in}}{\pgfqpoint{4.900000in}{4.900000in}}%
\pgfusepath{clip}%
\pgfsetbuttcap%
\pgfsetroundjoin%
\definecolor{currentfill}{rgb}{0.784667,0.784667,0.784667}%
\pgfsetfillcolor{currentfill}%
\pgfsetfillopacity{0.900000}%
\pgfsetlinewidth{0.501875pt}%
\definecolor{currentstroke}{rgb}{0.200000,0.200000,0.200000}%
\pgfsetstrokecolor{currentstroke}%
\pgfsetdash{}{0pt}%
\pgfpathmoveto{\pgfqpoint{2.229348in}{2.378865in}}%
\pgfpathlineto{\pgfqpoint{2.275165in}{2.335089in}}%
\pgfpathlineto{\pgfqpoint{2.306355in}{2.314919in}}%
\pgfpathlineto{\pgfqpoint{2.260533in}{2.358688in}}%
\pgfpathlineto{\pgfqpoint{2.229348in}{2.378865in}}%
\pgfpathclose%
\pgfusepath{stroke,fill}%
\end{pgfscope}%
\begin{pgfscope}%
\pgfpathrectangle{\pgfqpoint{0.050000in}{0.050000in}}{\pgfqpoint{4.900000in}{4.900000in}}%
\pgfusepath{clip}%
\pgfsetbuttcap%
\pgfsetroundjoin%
\definecolor{currentfill}{rgb}{0.915356,0.915356,0.915356}%
\pgfsetfillcolor{currentfill}%
\pgfsetfillopacity{0.900000}%
\pgfsetlinewidth{0.501875pt}%
\definecolor{currentstroke}{rgb}{0.200000,0.200000,0.200000}%
\pgfsetstrokecolor{currentstroke}%
\pgfsetdash{}{0pt}%
\pgfpathmoveto{\pgfqpoint{2.753809in}{1.969663in}}%
\pgfpathlineto{\pgfqpoint{2.798561in}{1.935177in}}%
\pgfpathlineto{\pgfqpoint{2.830563in}{1.915747in}}%
\pgfpathlineto{\pgfqpoint{2.785802in}{1.949944in}}%
\pgfpathlineto{\pgfqpoint{2.753809in}{1.969663in}}%
\pgfpathclose%
\pgfusepath{stroke,fill}%
\end{pgfscope}%
\begin{pgfscope}%
\pgfpathrectangle{\pgfqpoint{0.050000in}{0.050000in}}{\pgfqpoint{4.900000in}{4.900000in}}%
\pgfusepath{clip}%
\pgfsetbuttcap%
\pgfsetroundjoin%
\definecolor{currentfill}{rgb}{0.992618,0.992618,0.992618}%
\pgfsetfillcolor{currentfill}%
\pgfsetfillopacity{0.900000}%
\pgfsetlinewidth{0.501875pt}%
\definecolor{currentstroke}{rgb}{0.200000,0.200000,0.200000}%
\pgfsetstrokecolor{currentstroke}%
\pgfsetdash{}{0pt}%
\pgfpathmoveto{\pgfqpoint{3.598902in}{1.599165in}}%
\pgfpathlineto{\pgfqpoint{3.642753in}{1.604609in}}%
\pgfpathlineto{\pgfqpoint{3.676089in}{1.584969in}}%
\pgfpathlineto{\pgfqpoint{3.632213in}{1.579628in}}%
\pgfpathlineto{\pgfqpoint{3.598902in}{1.599165in}}%
\pgfpathclose%
\pgfusepath{stroke,fill}%
\end{pgfscope}%
\begin{pgfscope}%
\pgfpathrectangle{\pgfqpoint{0.050000in}{0.050000in}}{\pgfqpoint{4.900000in}{4.900000in}}%
\pgfusepath{clip}%
\pgfsetbuttcap%
\pgfsetroundjoin%
\definecolor{currentfill}{rgb}{0.424083,0.424083,0.424083}%
\pgfsetfillcolor{currentfill}%
\pgfsetfillopacity{0.900000}%
\pgfsetlinewidth{0.501875pt}%
\definecolor{currentstroke}{rgb}{0.200000,0.200000,0.200000}%
\pgfsetstrokecolor{currentstroke}%
\pgfsetdash{}{0pt}%
\pgfpathmoveto{\pgfqpoint{1.257138in}{3.165207in}}%
\pgfpathlineto{\pgfqpoint{1.304997in}{3.119668in}}%
\pgfpathlineto{\pgfqpoint{1.334704in}{3.100641in}}%
\pgfpathlineto{\pgfqpoint{1.286833in}{3.146210in}}%
\pgfpathlineto{\pgfqpoint{1.257138in}{3.165207in}}%
\pgfpathclose%
\pgfusepath{stroke,fill}%
\end{pgfscope}%
\begin{pgfscope}%
\pgfpathrectangle{\pgfqpoint{0.050000in}{0.050000in}}{\pgfqpoint{4.900000in}{4.900000in}}%
\pgfusepath{clip}%
\pgfsetbuttcap%
\pgfsetroundjoin%
\definecolor{currentfill}{rgb}{1.000000,1.000000,1.000000}%
\pgfsetfillcolor{currentfill}%
\pgfsetfillopacity{0.900000}%
\pgfsetlinewidth{0.501875pt}%
\definecolor{currentstroke}{rgb}{0.200000,0.200000,0.200000}%
\pgfsetstrokecolor{currentstroke}%
\pgfsetdash{}{0pt}%
\pgfpathmoveto{\pgfqpoint{4.150283in}{1.556883in}}%
\pgfpathlineto{\pgfqpoint{4.193654in}{1.569388in}}%
\pgfpathlineto{\pgfqpoint{4.227828in}{1.548747in}}%
\pgfpathlineto{\pgfqpoint{4.184430in}{1.536299in}}%
\pgfpathlineto{\pgfqpoint{4.150283in}{1.556883in}}%
\pgfpathclose%
\pgfusepath{stroke,fill}%
\end{pgfscope}%
\begin{pgfscope}%
\pgfpathrectangle{\pgfqpoint{0.050000in}{0.050000in}}{\pgfqpoint{4.900000in}{4.900000in}}%
\pgfusepath{clip}%
\pgfsetbuttcap%
\pgfsetroundjoin%
\definecolor{currentfill}{rgb}{0.691396,0.691396,0.691396}%
\pgfsetfillcolor{currentfill}%
\pgfsetfillopacity{0.900000}%
\pgfsetlinewidth{0.501875pt}%
\definecolor{currentstroke}{rgb}{0.200000,0.200000,0.200000}%
\pgfsetstrokecolor{currentstroke}%
\pgfsetdash{}{0pt}%
\pgfpathmoveto{\pgfqpoint{1.967014in}{2.591451in}}%
\pgfpathlineto{\pgfqpoint{2.013389in}{2.547112in}}%
\pgfpathlineto{\pgfqpoint{2.044183in}{2.527208in}}%
\pgfpathlineto{\pgfqpoint{1.997801in}{2.571573in}}%
\pgfpathlineto{\pgfqpoint{1.967014in}{2.591451in}}%
\pgfpathclose%
\pgfusepath{stroke,fill}%
\end{pgfscope}%
\begin{pgfscope}%
\pgfpathrectangle{\pgfqpoint{0.050000in}{0.050000in}}{\pgfqpoint{4.900000in}{4.900000in}}%
\pgfusepath{clip}%
\pgfsetbuttcap%
\pgfsetroundjoin%
\definecolor{currentfill}{rgb}{0.581776,0.581776,0.581776}%
\pgfsetfillcolor{currentfill}%
\pgfsetfillopacity{0.900000}%
\pgfsetlinewidth{0.501875pt}%
\definecolor{currentstroke}{rgb}{0.200000,0.200000,0.200000}%
\pgfsetstrokecolor{currentstroke}%
\pgfsetdash{}{0pt}%
\pgfpathmoveto{\pgfqpoint{1.704591in}{2.804176in}}%
\pgfpathlineto{\pgfqpoint{1.751524in}{2.759383in}}%
\pgfpathlineto{\pgfqpoint{1.781922in}{2.739800in}}%
\pgfpathlineto{\pgfqpoint{1.734981in}{2.784620in}}%
\pgfpathlineto{\pgfqpoint{1.704591in}{2.804176in}}%
\pgfpathclose%
\pgfusepath{stroke,fill}%
\end{pgfscope}%
\begin{pgfscope}%
\pgfpathrectangle{\pgfqpoint{0.050000in}{0.050000in}}{\pgfqpoint{4.900000in}{4.900000in}}%
\pgfusepath{clip}%
\pgfsetbuttcap%
\pgfsetroundjoin%
\definecolor{currentfill}{rgb}{0.998155,0.998155,0.998155}%
\pgfsetfillcolor{currentfill}%
\pgfsetfillopacity{0.900000}%
\pgfsetlinewidth{0.501875pt}%
\definecolor{currentstroke}{rgb}{0.200000,0.200000,0.200000}%
\pgfsetstrokecolor{currentstroke}%
\pgfsetdash{}{0pt}%
\pgfpathmoveto{\pgfqpoint{3.951816in}{1.561758in}}%
\pgfpathlineto{\pgfqpoint{3.995397in}{1.573320in}}%
\pgfpathlineto{\pgfqpoint{4.029281in}{1.552988in}}%
\pgfpathlineto{\pgfqpoint{3.985674in}{1.541499in}}%
\pgfpathlineto{\pgfqpoint{3.951816in}{1.561758in}}%
\pgfpathclose%
\pgfusepath{stroke,fill}%
\end{pgfscope}%
\begin{pgfscope}%
\pgfpathrectangle{\pgfqpoint{0.050000in}{0.050000in}}{\pgfqpoint{4.900000in}{4.900000in}}%
\pgfusepath{clip}%
\pgfsetbuttcap%
\pgfsetroundjoin%
\definecolor{currentfill}{rgb}{0.994464,0.994464,0.994464}%
\pgfsetfillcolor{currentfill}%
\pgfsetfillopacity{0.900000}%
\pgfsetlinewidth{0.501875pt}%
\definecolor{currentstroke}{rgb}{0.200000,0.200000,0.200000}%
\pgfsetstrokecolor{currentstroke}%
\pgfsetdash{}{0pt}%
\pgfpathmoveto{\pgfqpoint{3.676089in}{1.584969in}}%
\pgfpathlineto{\pgfqpoint{3.719895in}{1.592221in}}%
\pgfpathlineto{\pgfqpoint{3.753361in}{1.572414in}}%
\pgfpathlineto{\pgfqpoint{3.709529in}{1.565263in}}%
\pgfpathlineto{\pgfqpoint{3.676089in}{1.584969in}}%
\pgfpathclose%
\pgfusepath{stroke,fill}%
\end{pgfscope}%
\begin{pgfscope}%
\pgfpathrectangle{\pgfqpoint{0.050000in}{0.050000in}}{\pgfqpoint{4.900000in}{4.900000in}}%
\pgfusepath{clip}%
\pgfsetbuttcap%
\pgfsetroundjoin%
\definecolor{currentfill}{rgb}{0.836340,0.836340,0.836340}%
\pgfsetfillcolor{currentfill}%
\pgfsetfillopacity{0.900000}%
\pgfsetlinewidth{0.501875pt}%
\definecolor{currentstroke}{rgb}{0.200000,0.200000,0.200000}%
\pgfsetstrokecolor{currentstroke}%
\pgfsetdash{}{0pt}%
\pgfpathmoveto{\pgfqpoint{2.414660in}{2.231106in}}%
\pgfpathlineto{\pgfqpoint{2.460107in}{2.188405in}}%
\pgfpathlineto{\pgfqpoint{2.491593in}{2.168235in}}%
\pgfpathlineto{\pgfqpoint{2.446143in}{2.210835in}}%
\pgfpathlineto{\pgfqpoint{2.414660in}{2.231106in}}%
\pgfpathclose%
\pgfusepath{stroke,fill}%
\end{pgfscope}%
\begin{pgfscope}%
\pgfpathrectangle{\pgfqpoint{0.050000in}{0.050000in}}{\pgfqpoint{4.900000in}{4.900000in}}%
\pgfusepath{clip}%
\pgfsetbuttcap%
\pgfsetroundjoin%
\definecolor{currentfill}{rgb}{0.898378,0.898378,0.898378}%
\pgfsetfillcolor{currentfill}%
\pgfsetfillopacity{0.900000}%
\pgfsetlinewidth{0.501875pt}%
\definecolor{currentstroke}{rgb}{0.200000,0.200000,0.200000}%
\pgfsetstrokecolor{currentstroke}%
\pgfsetdash{}{0pt}%
\pgfpathmoveto{\pgfqpoint{2.677002in}{2.026795in}}%
\pgfpathlineto{\pgfqpoint{2.721915in}{1.989336in}}%
\pgfpathlineto{\pgfqpoint{2.753809in}{1.969663in}}%
\pgfpathlineto{\pgfqpoint{2.708889in}{2.006854in}}%
\pgfpathlineto{\pgfqpoint{2.677002in}{2.026795in}}%
\pgfpathclose%
\pgfusepath{stroke,fill}%
\end{pgfscope}%
\begin{pgfscope}%
\pgfpathrectangle{\pgfqpoint{0.050000in}{0.050000in}}{\pgfqpoint{4.900000in}{4.900000in}}%
\pgfusepath{clip}%
\pgfsetbuttcap%
\pgfsetroundjoin%
\definecolor{currentfill}{rgb}{0.487043,0.487043,0.487043}%
\pgfsetfillcolor{currentfill}%
\pgfsetfillopacity{0.900000}%
\pgfsetlinewidth{0.501875pt}%
\definecolor{currentstroke}{rgb}{0.200000,0.200000,0.200000}%
\pgfsetstrokecolor{currentstroke}%
\pgfsetdash{}{0pt}%
\pgfpathmoveto{\pgfqpoint{1.442081in}{3.016974in}}%
\pgfpathlineto{\pgfqpoint{1.489572in}{2.971728in}}%
\pgfpathlineto{\pgfqpoint{1.519574in}{2.952467in}}%
\pgfpathlineto{\pgfqpoint{1.472073in}{2.997742in}}%
\pgfpathlineto{\pgfqpoint{1.442081in}{3.016974in}}%
\pgfpathclose%
\pgfusepath{stroke,fill}%
\end{pgfscope}%
\begin{pgfscope}%
\pgfpathrectangle{\pgfqpoint{0.050000in}{0.050000in}}{\pgfqpoint{4.900000in}{4.900000in}}%
\pgfusepath{clip}%
\pgfsetbuttcap%
\pgfsetroundjoin%
\definecolor{currentfill}{rgb}{0.757109,0.757109,0.757109}%
\pgfsetfillcolor{currentfill}%
\pgfsetfillopacity{0.900000}%
\pgfsetlinewidth{0.501875pt}%
\definecolor{currentstroke}{rgb}{0.200000,0.200000,0.200000}%
\pgfsetstrokecolor{currentstroke}%
\pgfsetdash{}{0pt}%
\pgfpathmoveto{\pgfqpoint{2.152254in}{2.442999in}}%
\pgfpathlineto{\pgfqpoint{2.198259in}{2.398984in}}%
\pgfpathlineto{\pgfqpoint{2.229348in}{2.378865in}}%
\pgfpathlineto{\pgfqpoint{2.183338in}{2.422892in}}%
\pgfpathlineto{\pgfqpoint{2.152254in}{2.442999in}}%
\pgfpathclose%
\pgfusepath{stroke,fill}%
\end{pgfscope}%
\begin{pgfscope}%
\pgfpathrectangle{\pgfqpoint{0.050000in}{0.050000in}}{\pgfqpoint{4.900000in}{4.900000in}}%
\pgfusepath{clip}%
\pgfsetbuttcap%
\pgfsetroundjoin%
\definecolor{currentfill}{rgb}{0.970473,0.970473,0.970473}%
\pgfsetfillcolor{currentfill}%
\pgfsetfillopacity{0.900000}%
\pgfsetlinewidth{0.501875pt}%
\definecolor{currentstroke}{rgb}{0.200000,0.200000,0.200000}%
\pgfsetstrokecolor{currentstroke}%
\pgfsetdash{}{0pt}%
\pgfpathmoveto{\pgfqpoint{3.169960in}{1.722424in}}%
\pgfpathlineto{\pgfqpoint{3.214154in}{1.708746in}}%
\pgfpathlineto{\pgfqpoint{3.246838in}{1.689835in}}%
\pgfpathlineto{\pgfqpoint{3.202625in}{1.703431in}}%
\pgfpathlineto{\pgfqpoint{3.169960in}{1.722424in}}%
\pgfpathclose%
\pgfusepath{stroke,fill}%
\end{pgfscope}%
\begin{pgfscope}%
\pgfpathrectangle{\pgfqpoint{0.050000in}{0.050000in}}{\pgfqpoint{4.900000in}{4.900000in}}%
\pgfusepath{clip}%
\pgfsetbuttcap%
\pgfsetroundjoin%
\definecolor{currentfill}{rgb}{0.976009,0.976009,0.976009}%
\pgfsetfillcolor{currentfill}%
\pgfsetfillopacity{0.900000}%
\pgfsetlinewidth{0.501875pt}%
\definecolor{currentstroke}{rgb}{0.200000,0.200000,0.200000}%
\pgfsetstrokecolor{currentstroke}%
\pgfsetdash{}{0pt}%
\pgfpathmoveto{\pgfqpoint{3.246838in}{1.689835in}}%
\pgfpathlineto{\pgfqpoint{3.290963in}{1.680251in}}%
\pgfpathlineto{\pgfqpoint{3.323769in}{1.661281in}}%
\pgfpathlineto{\pgfqpoint{3.279624in}{1.670837in}}%
\pgfpathlineto{\pgfqpoint{3.246838in}{1.689835in}}%
\pgfpathclose%
\pgfusepath{stroke,fill}%
\end{pgfscope}%
\begin{pgfscope}%
\pgfpathrectangle{\pgfqpoint{0.050000in}{0.050000in}}{\pgfqpoint{4.900000in}{4.900000in}}%
\pgfusepath{clip}%
\pgfsetbuttcap%
\pgfsetroundjoin%
\definecolor{currentfill}{rgb}{0.963091,0.963091,0.963091}%
\pgfsetfillcolor{currentfill}%
\pgfsetfillopacity{0.900000}%
\pgfsetlinewidth{0.501875pt}%
\definecolor{currentstroke}{rgb}{0.200000,0.200000,0.200000}%
\pgfsetstrokecolor{currentstroke}%
\pgfsetdash{}{0pt}%
\pgfpathmoveto{\pgfqpoint{3.093121in}{1.759313in}}%
\pgfpathlineto{\pgfqpoint{3.137398in}{1.741331in}}%
\pgfpathlineto{\pgfqpoint{3.169960in}{1.722424in}}%
\pgfpathlineto{\pgfqpoint{3.125667in}{1.740265in}}%
\pgfpathlineto{\pgfqpoint{3.093121in}{1.759313in}}%
\pgfpathclose%
\pgfusepath{stroke,fill}%
\end{pgfscope}%
\begin{pgfscope}%
\pgfpathrectangle{\pgfqpoint{0.050000in}{0.050000in}}{\pgfqpoint{4.900000in}{4.900000in}}%
\pgfusepath{clip}%
\pgfsetbuttcap%
\pgfsetroundjoin%
\definecolor{currentfill}{rgb}{0.391603,0.391603,0.391603}%
\pgfsetfillcolor{currentfill}%
\pgfsetfillopacity{0.900000}%
\pgfsetlinewidth{0.501875pt}%
\definecolor{currentstroke}{rgb}{0.200000,0.200000,0.200000}%
\pgfsetstrokecolor{currentstroke}%
\pgfsetdash{}{0pt}%
\pgfpathmoveto{\pgfqpoint{1.179484in}{3.229847in}}%
\pgfpathlineto{\pgfqpoint{1.227533in}{3.184147in}}%
\pgfpathlineto{\pgfqpoint{1.257138in}{3.165207in}}%
\pgfpathlineto{\pgfqpoint{1.209077in}{3.210937in}}%
\pgfpathlineto{\pgfqpoint{1.179484in}{3.229847in}}%
\pgfpathclose%
\pgfusepath{stroke,fill}%
\end{pgfscope}%
\begin{pgfscope}%
\pgfpathrectangle{\pgfqpoint{0.050000in}{0.050000in}}{\pgfqpoint{4.900000in}{4.900000in}}%
\pgfusepath{clip}%
\pgfsetbuttcap%
\pgfsetroundjoin%
\definecolor{currentfill}{rgb}{0.979700,0.979700,0.979700}%
\pgfsetfillcolor{currentfill}%
\pgfsetfillopacity{0.900000}%
\pgfsetlinewidth{0.501875pt}%
\definecolor{currentstroke}{rgb}{0.200000,0.200000,0.200000}%
\pgfsetstrokecolor{currentstroke}%
\pgfsetdash{}{0pt}%
\pgfpathmoveto{\pgfqpoint{3.323769in}{1.661281in}}%
\pgfpathlineto{\pgfqpoint{3.367834in}{1.655477in}}%
\pgfpathlineto{\pgfqpoint{3.400765in}{1.636402in}}%
\pgfpathlineto{\pgfqpoint{3.356679in}{1.642224in}}%
\pgfpathlineto{\pgfqpoint{3.323769in}{1.661281in}}%
\pgfpathclose%
\pgfusepath{stroke,fill}%
\end{pgfscope}%
\begin{pgfscope}%
\pgfpathrectangle{\pgfqpoint{0.050000in}{0.050000in}}{\pgfqpoint{4.900000in}{4.900000in}}%
\pgfusepath{clip}%
\pgfsetbuttcap%
\pgfsetroundjoin%
\definecolor{currentfill}{rgb}{0.657809,0.657809,0.657809}%
\pgfsetfillcolor{currentfill}%
\pgfsetfillopacity{0.900000}%
\pgfsetlinewidth{0.501875pt}%
\definecolor{currentstroke}{rgb}{0.200000,0.200000,0.200000}%
\pgfsetstrokecolor{currentstroke}%
\pgfsetdash{}{0pt}%
\pgfpathmoveto{\pgfqpoint{1.889757in}{2.655767in}}%
\pgfpathlineto{\pgfqpoint{1.936321in}{2.611268in}}%
\pgfpathlineto{\pgfqpoint{1.967014in}{2.591451in}}%
\pgfpathlineto{\pgfqpoint{1.920444in}{2.635976in}}%
\pgfpathlineto{\pgfqpoint{1.889757in}{2.655767in}}%
\pgfpathclose%
\pgfusepath{stroke,fill}%
\end{pgfscope}%
\begin{pgfscope}%
\pgfpathrectangle{\pgfqpoint{0.050000in}{0.050000in}}{\pgfqpoint{4.900000in}{4.900000in}}%
\pgfusepath{clip}%
\pgfsetbuttcap%
\pgfsetroundjoin%
\definecolor{currentfill}{rgb}{0.955709,0.955709,0.955709}%
\pgfsetfillcolor{currentfill}%
\pgfsetfillopacity{0.900000}%
\pgfsetlinewidth{0.501875pt}%
\definecolor{currentstroke}{rgb}{0.200000,0.200000,0.200000}%
\pgfsetstrokecolor{currentstroke}%
\pgfsetdash{}{0pt}%
\pgfpathmoveto{\pgfqpoint{3.016305in}{1.800636in}}%
\pgfpathlineto{\pgfqpoint{3.060678in}{1.778279in}}%
\pgfpathlineto{\pgfqpoint{3.093121in}{1.759313in}}%
\pgfpathlineto{\pgfqpoint{3.048733in}{1.781471in}}%
\pgfpathlineto{\pgfqpoint{3.016305in}{1.800636in}}%
\pgfpathclose%
\pgfusepath{stroke,fill}%
\end{pgfscope}%
\begin{pgfscope}%
\pgfpathrectangle{\pgfqpoint{0.050000in}{0.050000in}}{\pgfqpoint{4.900000in}{4.900000in}}%
\pgfusepath{clip}%
\pgfsetbuttcap%
\pgfsetroundjoin%
\definecolor{currentfill}{rgb}{0.996309,0.996309,0.996309}%
\pgfsetfillcolor{currentfill}%
\pgfsetfillopacity{0.900000}%
\pgfsetlinewidth{0.501875pt}%
\definecolor{currentstroke}{rgb}{0.200000,0.200000,0.200000}%
\pgfsetstrokecolor{currentstroke}%
\pgfsetdash{}{0pt}%
\pgfpathmoveto{\pgfqpoint{3.753361in}{1.572414in}}%
\pgfpathlineto{\pgfqpoint{3.797121in}{1.581138in}}%
\pgfpathlineto{\pgfqpoint{3.830716in}{1.561171in}}%
\pgfpathlineto{\pgfqpoint{3.786931in}{1.552543in}}%
\pgfpathlineto{\pgfqpoint{3.753361in}{1.572414in}}%
\pgfpathclose%
\pgfusepath{stroke,fill}%
\end{pgfscope}%
\begin{pgfscope}%
\pgfpathrectangle{\pgfqpoint{0.050000in}{0.050000in}}{\pgfqpoint{4.900000in}{4.900000in}}%
\pgfusepath{clip}%
\pgfsetbuttcap%
\pgfsetroundjoin%
\definecolor{currentfill}{rgb}{1.000000,1.000000,1.000000}%
\pgfsetfillcolor{currentfill}%
\pgfsetfillopacity{0.900000}%
\pgfsetlinewidth{0.501875pt}%
\definecolor{currentstroke}{rgb}{0.200000,0.200000,0.200000}%
\pgfsetstrokecolor{currentstroke}%
\pgfsetdash{}{0pt}%
\pgfpathmoveto{\pgfqpoint{4.029281in}{1.552988in}}%
\pgfpathlineto{\pgfqpoint{4.072804in}{1.565021in}}%
\pgfpathlineto{\pgfqpoint{4.106820in}{1.544561in}}%
\pgfpathlineto{\pgfqpoint{4.063270in}{1.532595in}}%
\pgfpathlineto{\pgfqpoint{4.029281in}{1.552988in}}%
\pgfpathclose%
\pgfusepath{stroke,fill}%
\end{pgfscope}%
\begin{pgfscope}%
\pgfpathrectangle{\pgfqpoint{0.050000in}{0.050000in}}{\pgfqpoint{4.900000in}{4.900000in}}%
\pgfusepath{clip}%
\pgfsetbuttcap%
\pgfsetroundjoin%
\definecolor{currentfill}{rgb}{0.985236,0.985236,0.985236}%
\pgfsetfillcolor{currentfill}%
\pgfsetfillopacity{0.900000}%
\pgfsetlinewidth{0.501875pt}%
\definecolor{currentstroke}{rgb}{0.200000,0.200000,0.200000}%
\pgfsetstrokecolor{currentstroke}%
\pgfsetdash{}{0pt}%
\pgfpathmoveto{\pgfqpoint{3.400765in}{1.636402in}}%
\pgfpathlineto{\pgfqpoint{3.444777in}{1.633998in}}%
\pgfpathlineto{\pgfqpoint{3.477834in}{1.614786in}}%
\pgfpathlineto{\pgfqpoint{3.433799in}{1.617244in}}%
\pgfpathlineto{\pgfqpoint{3.400765in}{1.636402in}}%
\pgfpathclose%
\pgfusepath{stroke,fill}%
\end{pgfscope}%
\begin{pgfscope}%
\pgfpathrectangle{\pgfqpoint{0.050000in}{0.050000in}}{\pgfqpoint{4.900000in}{4.900000in}}%
\pgfusepath{clip}%
\pgfsetbuttcap%
\pgfsetroundjoin%
\definecolor{currentfill}{rgb}{0.551634,0.551634,0.551634}%
\pgfsetfillcolor{currentfill}%
\pgfsetfillopacity{0.900000}%
\pgfsetlinewidth{0.501875pt}%
\definecolor{currentstroke}{rgb}{0.200000,0.200000,0.200000}%
\pgfsetstrokecolor{currentstroke}%
\pgfsetdash{}{0pt}%
\pgfpathmoveto{\pgfqpoint{1.627173in}{2.868625in}}%
\pgfpathlineto{\pgfqpoint{1.674295in}{2.823671in}}%
\pgfpathlineto{\pgfqpoint{1.704591in}{2.804176in}}%
\pgfpathlineto{\pgfqpoint{1.657461in}{2.849157in}}%
\pgfpathlineto{\pgfqpoint{1.627173in}{2.868625in}}%
\pgfpathclose%
\pgfusepath{stroke,fill}%
\end{pgfscope}%
\begin{pgfscope}%
\pgfpathrectangle{\pgfqpoint{0.050000in}{0.050000in}}{\pgfqpoint{4.900000in}{4.900000in}}%
\pgfusepath{clip}%
\pgfsetbuttcap%
\pgfsetroundjoin%
\definecolor{currentfill}{rgb}{0.878570,0.878570,0.878570}%
\pgfsetfillcolor{currentfill}%
\pgfsetfillopacity{0.900000}%
\pgfsetlinewidth{0.501875pt}%
\definecolor{currentstroke}{rgb}{0.200000,0.200000,0.200000}%
\pgfsetstrokecolor{currentstroke}%
\pgfsetdash{}{0pt}%
\pgfpathmoveto{\pgfqpoint{2.600128in}{2.086475in}}%
\pgfpathlineto{\pgfqpoint{2.645213in}{2.046695in}}%
\pgfpathlineto{\pgfqpoint{2.677002in}{2.026795in}}%
\pgfpathlineto{\pgfqpoint{2.631911in}{2.066351in}}%
\pgfpathlineto{\pgfqpoint{2.600128in}{2.086475in}}%
\pgfpathclose%
\pgfusepath{stroke,fill}%
\end{pgfscope}%
\begin{pgfscope}%
\pgfpathrectangle{\pgfqpoint{0.050000in}{0.050000in}}{\pgfqpoint{4.900000in}{4.900000in}}%
\pgfusepath{clip}%
\pgfsetbuttcap%
\pgfsetroundjoin%
\definecolor{currentfill}{rgb}{0.946482,0.946482,0.946482}%
\pgfsetfillcolor{currentfill}%
\pgfsetfillopacity{0.900000}%
\pgfsetlinewidth{0.501875pt}%
\definecolor{currentstroke}{rgb}{0.200000,0.200000,0.200000}%
\pgfsetstrokecolor{currentstroke}%
\pgfsetdash{}{0pt}%
\pgfpathmoveto{\pgfqpoint{2.939492in}{1.846358in}}%
\pgfpathlineto{\pgfqpoint{2.983978in}{1.819725in}}%
\pgfpathlineto{\pgfqpoint{3.016305in}{1.800636in}}%
\pgfpathlineto{\pgfqpoint{2.971806in}{1.827023in}}%
\pgfpathlineto{\pgfqpoint{2.939492in}{1.846358in}}%
\pgfpathclose%
\pgfusepath{stroke,fill}%
\end{pgfscope}%
\begin{pgfscope}%
\pgfpathrectangle{\pgfqpoint{0.050000in}{0.050000in}}{\pgfqpoint{4.900000in}{4.900000in}}%
\pgfusepath{clip}%
\pgfsetbuttcap%
\pgfsetroundjoin%
\definecolor{currentfill}{rgb}{0.808781,0.808781,0.808781}%
\pgfsetfillcolor{currentfill}%
\pgfsetfillopacity{0.900000}%
\pgfsetlinewidth{0.501875pt}%
\definecolor{currentstroke}{rgb}{0.200000,0.200000,0.200000}%
\pgfsetstrokecolor{currentstroke}%
\pgfsetdash{}{0pt}%
\pgfpathmoveto{\pgfqpoint{2.337641in}{2.294698in}}%
\pgfpathlineto{\pgfqpoint{2.383274in}{2.251331in}}%
\pgfpathlineto{\pgfqpoint{2.414660in}{2.231106in}}%
\pgfpathlineto{\pgfqpoint{2.369023in}{2.274425in}}%
\pgfpathlineto{\pgfqpoint{2.337641in}{2.294698in}}%
\pgfpathclose%
\pgfusepath{stroke,fill}%
\end{pgfscope}%
\begin{pgfscope}%
\pgfpathrectangle{\pgfqpoint{0.050000in}{0.050000in}}{\pgfqpoint{4.900000in}{4.900000in}}%
\pgfusepath{clip}%
\pgfsetbuttcap%
\pgfsetroundjoin%
\definecolor{currentfill}{rgb}{0.452595,0.452595,0.452595}%
\pgfsetfillcolor{currentfill}%
\pgfsetfillopacity{0.900000}%
\pgfsetlinewidth{0.501875pt}%
\definecolor{currentstroke}{rgb}{0.200000,0.200000,0.200000}%
\pgfsetstrokecolor{currentstroke}%
\pgfsetdash{}{0pt}%
\pgfpathmoveto{\pgfqpoint{1.364501in}{3.081555in}}%
\pgfpathlineto{\pgfqpoint{1.412182in}{3.036148in}}%
\pgfpathlineto{\pgfqpoint{1.442081in}{3.016974in}}%
\pgfpathlineto{\pgfqpoint{1.394390in}{3.062411in}}%
\pgfpathlineto{\pgfqpoint{1.364501in}{3.081555in}}%
\pgfpathclose%
\pgfusepath{stroke,fill}%
\end{pgfscope}%
\begin{pgfscope}%
\pgfpathrectangle{\pgfqpoint{0.050000in}{0.050000in}}{\pgfqpoint{4.900000in}{4.900000in}}%
\pgfusepath{clip}%
\pgfsetbuttcap%
\pgfsetroundjoin%
\definecolor{currentfill}{rgb}{0.987082,0.987082,0.987082}%
\pgfsetfillcolor{currentfill}%
\pgfsetfillopacity{0.900000}%
\pgfsetlinewidth{0.501875pt}%
\definecolor{currentstroke}{rgb}{0.200000,0.200000,0.200000}%
\pgfsetstrokecolor{currentstroke}%
\pgfsetdash{}{0pt}%
\pgfpathmoveto{\pgfqpoint{3.477834in}{1.614786in}}%
\pgfpathlineto{\pgfqpoint{3.521799in}{1.615373in}}%
\pgfpathlineto{\pgfqpoint{3.554982in}{1.596003in}}%
\pgfpathlineto{\pgfqpoint{3.510994in}{1.595495in}}%
\pgfpathlineto{\pgfqpoint{3.477834in}{1.614786in}}%
\pgfpathclose%
\pgfusepath{stroke,fill}%
\end{pgfscope}%
\begin{pgfscope}%
\pgfpathrectangle{\pgfqpoint{0.050000in}{0.050000in}}{\pgfqpoint{4.900000in}{4.900000in}}%
\pgfusepath{clip}%
\pgfsetbuttcap%
\pgfsetroundjoin%
\definecolor{currentfill}{rgb}{0.729781,0.729781,0.729781}%
\pgfsetfillcolor{currentfill}%
\pgfsetfillopacity{0.900000}%
\pgfsetlinewidth{0.501875pt}%
\definecolor{currentstroke}{rgb}{0.200000,0.200000,0.200000}%
\pgfsetstrokecolor{currentstroke}%
\pgfsetdash{}{0pt}%
\pgfpathmoveto{\pgfqpoint{2.075071in}{2.507244in}}%
\pgfpathlineto{\pgfqpoint{2.121264in}{2.463046in}}%
\pgfpathlineto{\pgfqpoint{2.152254in}{2.442999in}}%
\pgfpathlineto{\pgfqpoint{2.106055in}{2.487219in}}%
\pgfpathlineto{\pgfqpoint{2.075071in}{2.507244in}}%
\pgfpathclose%
\pgfusepath{stroke,fill}%
\end{pgfscope}%
\begin{pgfscope}%
\pgfpathrectangle{\pgfqpoint{0.050000in}{0.050000in}}{\pgfqpoint{4.900000in}{4.900000in}}%
\pgfusepath{clip}%
\pgfsetbuttcap%
\pgfsetroundjoin%
\definecolor{currentfill}{rgb}{0.998155,0.998155,0.998155}%
\pgfsetfillcolor{currentfill}%
\pgfsetfillopacity{0.900000}%
\pgfsetlinewidth{0.501875pt}%
\definecolor{currentstroke}{rgb}{0.200000,0.200000,0.200000}%
\pgfsetstrokecolor{currentstroke}%
\pgfsetdash{}{0pt}%
\pgfpathmoveto{\pgfqpoint{3.830716in}{1.561171in}}%
\pgfpathlineto{\pgfqpoint{3.874428in}{1.571068in}}%
\pgfpathlineto{\pgfqpoint{3.908155in}{1.550950in}}%
\pgfpathlineto{\pgfqpoint{3.864417in}{1.541141in}}%
\pgfpathlineto{\pgfqpoint{3.830716in}{1.561171in}}%
\pgfpathclose%
\pgfusepath{stroke,fill}%
\end{pgfscope}%
\begin{pgfscope}%
\pgfpathrectangle{\pgfqpoint{0.050000in}{0.050000in}}{\pgfqpoint{4.900000in}{4.900000in}}%
\pgfusepath{clip}%
\pgfsetbuttcap%
\pgfsetroundjoin%
\definecolor{currentfill}{rgb}{0.932334,0.932334,0.932334}%
\pgfsetfillcolor{currentfill}%
\pgfsetfillopacity{0.900000}%
\pgfsetlinewidth{0.501875pt}%
\definecolor{currentstroke}{rgb}{0.200000,0.200000,0.200000}%
\pgfsetstrokecolor{currentstroke}%
\pgfsetdash{}{0pt}%
\pgfpathmoveto{\pgfqpoint{2.862664in}{1.896261in}}%
\pgfpathlineto{\pgfqpoint{2.907279in}{1.865626in}}%
\pgfpathlineto{\pgfqpoint{2.939492in}{1.846358in}}%
\pgfpathlineto{\pgfqpoint{2.894867in}{1.876718in}}%
\pgfpathlineto{\pgfqpoint{2.862664in}{1.896261in}}%
\pgfpathclose%
\pgfusepath{stroke,fill}%
\end{pgfscope}%
\begin{pgfscope}%
\pgfpathrectangle{\pgfqpoint{0.050000in}{0.050000in}}{\pgfqpoint{4.900000in}{4.900000in}}%
\pgfusepath{clip}%
\pgfsetbuttcap%
\pgfsetroundjoin%
\definecolor{currentfill}{rgb}{1.000000,1.000000,1.000000}%
\pgfsetfillcolor{currentfill}%
\pgfsetfillopacity{0.900000}%
\pgfsetlinewidth{0.501875pt}%
\definecolor{currentstroke}{rgb}{0.200000,0.200000,0.200000}%
\pgfsetstrokecolor{currentstroke}%
\pgfsetdash{}{0pt}%
\pgfpathmoveto{\pgfqpoint{4.106820in}{1.544561in}}%
\pgfpathlineto{\pgfqpoint{4.150283in}{1.556883in}}%
\pgfpathlineto{\pgfqpoint{4.184430in}{1.536299in}}%
\pgfpathlineto{\pgfqpoint{4.140942in}{1.524038in}}%
\pgfpathlineto{\pgfqpoint{4.106820in}{1.544561in}}%
\pgfpathclose%
\pgfusepath{stroke,fill}%
\end{pgfscope}%
\begin{pgfscope}%
\pgfpathrectangle{\pgfqpoint{0.050000in}{0.050000in}}{\pgfqpoint{4.900000in}{4.900000in}}%
\pgfusepath{clip}%
\pgfsetbuttcap%
\pgfsetroundjoin%
\definecolor{currentfill}{rgb}{0.619423,0.619423,0.619423}%
\pgfsetfillcolor{currentfill}%
\pgfsetfillopacity{0.900000}%
\pgfsetlinewidth{0.501875pt}%
\definecolor{currentstroke}{rgb}{0.200000,0.200000,0.200000}%
\pgfsetstrokecolor{currentstroke}%
\pgfsetdash{}{0pt}%
\pgfpathmoveto{\pgfqpoint{1.812413in}{2.720157in}}%
\pgfpathlineto{\pgfqpoint{1.859165in}{2.675498in}}%
\pgfpathlineto{\pgfqpoint{1.889757in}{2.655767in}}%
\pgfpathlineto{\pgfqpoint{1.842998in}{2.700454in}}%
\pgfpathlineto{\pgfqpoint{1.812413in}{2.720157in}}%
\pgfpathclose%
\pgfusepath{stroke,fill}%
\end{pgfscope}%
\begin{pgfscope}%
\pgfpathrectangle{\pgfqpoint{0.050000in}{0.050000in}}{\pgfqpoint{4.900000in}{4.900000in}}%
\pgfusepath{clip}%
\pgfsetbuttcap%
\pgfsetroundjoin%
\definecolor{currentfill}{rgb}{0.990773,0.990773,0.990773}%
\pgfsetfillcolor{currentfill}%
\pgfsetfillopacity{0.900000}%
\pgfsetlinewidth{0.501875pt}%
\definecolor{currentstroke}{rgb}{0.200000,0.200000,0.200000}%
\pgfsetstrokecolor{currentstroke}%
\pgfsetdash{}{0pt}%
\pgfpathmoveto{\pgfqpoint{3.554982in}{1.596003in}}%
\pgfpathlineto{\pgfqpoint{3.598902in}{1.599165in}}%
\pgfpathlineto{\pgfqpoint{3.632213in}{1.579628in}}%
\pgfpathlineto{\pgfqpoint{3.588270in}{1.576559in}}%
\pgfpathlineto{\pgfqpoint{3.554982in}{1.596003in}}%
\pgfpathclose%
\pgfusepath{stroke,fill}%
\end{pgfscope}%
\begin{pgfscope}%
\pgfpathrectangle{\pgfqpoint{0.050000in}{0.050000in}}{\pgfqpoint{4.900000in}{4.900000in}}%
\pgfusepath{clip}%
\pgfsetbuttcap%
\pgfsetroundjoin%
\definecolor{currentfill}{rgb}{0.858762,0.858762,0.858762}%
\pgfsetfillcolor{currentfill}%
\pgfsetfillopacity{0.900000}%
\pgfsetlinewidth{0.501875pt}%
\definecolor{currentstroke}{rgb}{0.200000,0.200000,0.200000}%
\pgfsetstrokecolor{currentstroke}%
\pgfsetdash{}{0pt}%
\pgfpathmoveto{\pgfqpoint{2.523177in}{2.148024in}}%
\pgfpathlineto{\pgfqpoint{2.568442in}{2.106560in}}%
\pgfpathlineto{\pgfqpoint{2.600128in}{2.086475in}}%
\pgfpathlineto{\pgfqpoint{2.554858in}{2.127774in}}%
\pgfpathlineto{\pgfqpoint{2.523177in}{2.148024in}}%
\pgfpathclose%
\pgfusepath{stroke,fill}%
\end{pgfscope}%
\begin{pgfscope}%
\pgfpathrectangle{\pgfqpoint{0.050000in}{0.050000in}}{\pgfqpoint{4.900000in}{4.900000in}}%
\pgfusepath{clip}%
\pgfsetbuttcap%
\pgfsetroundjoin%
\definecolor{currentfill}{rgb}{0.517186,0.517186,0.517186}%
\pgfsetfillcolor{currentfill}%
\pgfsetfillopacity{0.900000}%
\pgfsetlinewidth{0.501875pt}%
\definecolor{currentstroke}{rgb}{0.200000,0.200000,0.200000}%
\pgfsetstrokecolor{currentstroke}%
\pgfsetdash{}{0pt}%
\pgfpathmoveto{\pgfqpoint{1.549667in}{2.933147in}}%
\pgfpathlineto{\pgfqpoint{1.596978in}{2.888033in}}%
\pgfpathlineto{\pgfqpoint{1.627173in}{2.868625in}}%
\pgfpathlineto{\pgfqpoint{1.579853in}{2.913767in}}%
\pgfpathlineto{\pgfqpoint{1.549667in}{2.933147in}}%
\pgfpathclose%
\pgfusepath{stroke,fill}%
\end{pgfscope}%
\begin{pgfscope}%
\pgfpathrectangle{\pgfqpoint{0.050000in}{0.050000in}}{\pgfqpoint{4.900000in}{4.900000in}}%
\pgfusepath{clip}%
\pgfsetbuttcap%
\pgfsetroundjoin%
\definecolor{currentfill}{rgb}{0.784667,0.784667,0.784667}%
\pgfsetfillcolor{currentfill}%
\pgfsetfillopacity{0.900000}%
\pgfsetlinewidth{0.501875pt}%
\definecolor{currentstroke}{rgb}{0.200000,0.200000,0.200000}%
\pgfsetstrokecolor{currentstroke}%
\pgfsetdash{}{0pt}%
\pgfpathmoveto{\pgfqpoint{2.260533in}{2.358688in}}%
\pgfpathlineto{\pgfqpoint{2.306355in}{2.314919in}}%
\pgfpathlineto{\pgfqpoint{2.337641in}{2.294698in}}%
\pgfpathlineto{\pgfqpoint{2.291814in}{2.338456in}}%
\pgfpathlineto{\pgfqpoint{2.260533in}{2.358688in}}%
\pgfpathclose%
\pgfusepath{stroke,fill}%
\end{pgfscope}%
\begin{pgfscope}%
\pgfpathrectangle{\pgfqpoint{0.050000in}{0.050000in}}{\pgfqpoint{4.900000in}{4.900000in}}%
\pgfusepath{clip}%
\pgfsetbuttcap%
\pgfsetroundjoin%
\definecolor{currentfill}{rgb}{0.998155,0.998155,0.998155}%
\pgfsetfillcolor{currentfill}%
\pgfsetfillopacity{0.900000}%
\pgfsetlinewidth{0.501875pt}%
\definecolor{currentstroke}{rgb}{0.200000,0.200000,0.200000}%
\pgfsetstrokecolor{currentstroke}%
\pgfsetdash{}{0pt}%
\pgfpathmoveto{\pgfqpoint{3.908155in}{1.550950in}}%
\pgfpathlineto{\pgfqpoint{3.951816in}{1.561758in}}%
\pgfpathlineto{\pgfqpoint{3.985674in}{1.541499in}}%
\pgfpathlineto{\pgfqpoint{3.941986in}{1.530771in}}%
\pgfpathlineto{\pgfqpoint{3.908155in}{1.550950in}}%
\pgfpathclose%
\pgfusepath{stroke,fill}%
\end{pgfscope}%
\begin{pgfscope}%
\pgfpathrectangle{\pgfqpoint{0.050000in}{0.050000in}}{\pgfqpoint{4.900000in}{4.900000in}}%
\pgfusepath{clip}%
\pgfsetbuttcap%
\pgfsetroundjoin%
\definecolor{currentfill}{rgb}{0.915356,0.915356,0.915356}%
\pgfsetfillcolor{currentfill}%
\pgfsetfillopacity{0.900000}%
\pgfsetlinewidth{0.501875pt}%
\definecolor{currentstroke}{rgb}{0.200000,0.200000,0.200000}%
\pgfsetstrokecolor{currentstroke}%
\pgfsetdash{}{0pt}%
\pgfpathmoveto{\pgfqpoint{2.785802in}{1.949944in}}%
\pgfpathlineto{\pgfqpoint{2.830563in}{1.915747in}}%
\pgfpathlineto{\pgfqpoint{2.862664in}{1.896261in}}%
\pgfpathlineto{\pgfqpoint{2.817895in}{1.930175in}}%
\pgfpathlineto{\pgfqpoint{2.785802in}{1.949944in}}%
\pgfpathclose%
\pgfusepath{stroke,fill}%
\end{pgfscope}%
\begin{pgfscope}%
\pgfpathrectangle{\pgfqpoint{0.050000in}{0.050000in}}{\pgfqpoint{4.900000in}{4.900000in}}%
\pgfusepath{clip}%
\pgfsetbuttcap%
\pgfsetroundjoin%
\definecolor{currentfill}{rgb}{0.424083,0.424083,0.424083}%
\pgfsetfillcolor{currentfill}%
\pgfsetfillopacity{0.900000}%
\pgfsetlinewidth{0.501875pt}%
\definecolor{currentstroke}{rgb}{0.200000,0.200000,0.200000}%
\pgfsetstrokecolor{currentstroke}%
\pgfsetdash{}{0pt}%
\pgfpathmoveto{\pgfqpoint{1.286833in}{3.146210in}}%
\pgfpathlineto{\pgfqpoint{1.334704in}{3.100641in}}%
\pgfpathlineto{\pgfqpoint{1.364501in}{3.081555in}}%
\pgfpathlineto{\pgfqpoint{1.316620in}{3.127154in}}%
\pgfpathlineto{\pgfqpoint{1.286833in}{3.146210in}}%
\pgfpathclose%
\pgfusepath{stroke,fill}%
\end{pgfscope}%
\begin{pgfscope}%
\pgfpathrectangle{\pgfqpoint{0.050000in}{0.050000in}}{\pgfqpoint{4.900000in}{4.900000in}}%
\pgfusepath{clip}%
\pgfsetbuttcap%
\pgfsetroundjoin%
\definecolor{currentfill}{rgb}{0.992618,0.992618,0.992618}%
\pgfsetfillcolor{currentfill}%
\pgfsetfillopacity{0.900000}%
\pgfsetlinewidth{0.501875pt}%
\definecolor{currentstroke}{rgb}{0.200000,0.200000,0.200000}%
\pgfsetstrokecolor{currentstroke}%
\pgfsetdash{}{0pt}%
\pgfpathmoveto{\pgfqpoint{3.632213in}{1.579628in}}%
\pgfpathlineto{\pgfqpoint{3.676089in}{1.584969in}}%
\pgfpathlineto{\pgfqpoint{3.709529in}{1.565263in}}%
\pgfpathlineto{\pgfqpoint{3.665629in}{1.560021in}}%
\pgfpathlineto{\pgfqpoint{3.632213in}{1.579628in}}%
\pgfpathclose%
\pgfusepath{stroke,fill}%
\end{pgfscope}%
\begin{pgfscope}%
\pgfpathrectangle{\pgfqpoint{0.050000in}{0.050000in}}{\pgfqpoint{4.900000in}{4.900000in}}%
\pgfusepath{clip}%
\pgfsetbuttcap%
\pgfsetroundjoin%
\definecolor{currentfill}{rgb}{1.000000,1.000000,1.000000}%
\pgfsetfillcolor{currentfill}%
\pgfsetfillopacity{0.900000}%
\pgfsetlinewidth{0.501875pt}%
\definecolor{currentstroke}{rgb}{0.200000,0.200000,0.200000}%
\pgfsetstrokecolor{currentstroke}%
\pgfsetdash{}{0pt}%
\pgfpathmoveto{\pgfqpoint{4.184430in}{1.536299in}}%
\pgfpathlineto{\pgfqpoint{4.227828in}{1.548747in}}%
\pgfpathlineto{\pgfqpoint{4.262107in}{1.528041in}}%
\pgfpathlineto{\pgfqpoint{4.218684in}{1.515652in}}%
\pgfpathlineto{\pgfqpoint{4.184430in}{1.536299in}}%
\pgfpathclose%
\pgfusepath{stroke,fill}%
\end{pgfscope}%
\begin{pgfscope}%
\pgfpathrectangle{\pgfqpoint{0.050000in}{0.050000in}}{\pgfqpoint{4.900000in}{4.900000in}}%
\pgfusepath{clip}%
\pgfsetbuttcap%
\pgfsetroundjoin%
\definecolor{currentfill}{rgb}{0.691396,0.691396,0.691396}%
\pgfsetfillcolor{currentfill}%
\pgfsetfillopacity{0.900000}%
\pgfsetlinewidth{0.501875pt}%
\definecolor{currentstroke}{rgb}{0.200000,0.200000,0.200000}%
\pgfsetstrokecolor{currentstroke}%
\pgfsetdash{}{0pt}%
\pgfpathmoveto{\pgfqpoint{1.997801in}{2.571573in}}%
\pgfpathlineto{\pgfqpoint{2.044183in}{2.527208in}}%
\pgfpathlineto{\pgfqpoint{2.075071in}{2.507244in}}%
\pgfpathlineto{\pgfqpoint{2.028684in}{2.551634in}}%
\pgfpathlineto{\pgfqpoint{1.997801in}{2.571573in}}%
\pgfpathclose%
\pgfusepath{stroke,fill}%
\end{pgfscope}%
\begin{pgfscope}%
\pgfpathrectangle{\pgfqpoint{0.050000in}{0.050000in}}{\pgfqpoint{4.900000in}{4.900000in}}%
\pgfusepath{clip}%
\pgfsetbuttcap%
\pgfsetroundjoin%
\definecolor{currentfill}{rgb}{0.586082,0.586082,0.586082}%
\pgfsetfillcolor{currentfill}%
\pgfsetfillopacity{0.900000}%
\pgfsetlinewidth{0.501875pt}%
\definecolor{currentstroke}{rgb}{0.200000,0.200000,0.200000}%
\pgfsetstrokecolor{currentstroke}%
\pgfsetdash{}{0pt}%
\pgfpathmoveto{\pgfqpoint{1.734981in}{2.784620in}}%
\pgfpathlineto{\pgfqpoint{1.781922in}{2.739800in}}%
\pgfpathlineto{\pgfqpoint{1.812413in}{2.720157in}}%
\pgfpathlineto{\pgfqpoint{1.765464in}{2.765005in}}%
\pgfpathlineto{\pgfqpoint{1.734981in}{2.784620in}}%
\pgfpathclose%
\pgfusepath{stroke,fill}%
\end{pgfscope}%
\begin{pgfscope}%
\pgfpathrectangle{\pgfqpoint{0.050000in}{0.050000in}}{\pgfqpoint{4.900000in}{4.900000in}}%
\pgfusepath{clip}%
\pgfsetbuttcap%
\pgfsetroundjoin%
\definecolor{currentfill}{rgb}{0.836340,0.836340,0.836340}%
\pgfsetfillcolor{currentfill}%
\pgfsetfillopacity{0.900000}%
\pgfsetlinewidth{0.501875pt}%
\definecolor{currentstroke}{rgb}{0.200000,0.200000,0.200000}%
\pgfsetstrokecolor{currentstroke}%
\pgfsetdash{}{0pt}%
\pgfpathmoveto{\pgfqpoint{2.446143in}{2.210835in}}%
\pgfpathlineto{\pgfqpoint{2.491593in}{2.168235in}}%
\pgfpathlineto{\pgfqpoint{2.523177in}{2.148024in}}%
\pgfpathlineto{\pgfqpoint{2.477722in}{2.190520in}}%
\pgfpathlineto{\pgfqpoint{2.446143in}{2.210835in}}%
\pgfpathclose%
\pgfusepath{stroke,fill}%
\end{pgfscope}%
\begin{pgfscope}%
\pgfpathrectangle{\pgfqpoint{0.050000in}{0.050000in}}{\pgfqpoint{4.900000in}{4.900000in}}%
\pgfusepath{clip}%
\pgfsetbuttcap%
\pgfsetroundjoin%
\definecolor{currentfill}{rgb}{0.998155,0.998155,0.998155}%
\pgfsetfillcolor{currentfill}%
\pgfsetfillopacity{0.900000}%
\pgfsetlinewidth{0.501875pt}%
\definecolor{currentstroke}{rgb}{0.200000,0.200000,0.200000}%
\pgfsetstrokecolor{currentstroke}%
\pgfsetdash{}{0pt}%
\pgfpathmoveto{\pgfqpoint{3.985674in}{1.541499in}}%
\pgfpathlineto{\pgfqpoint{4.029281in}{1.552988in}}%
\pgfpathlineto{\pgfqpoint{4.063270in}{1.532595in}}%
\pgfpathlineto{\pgfqpoint{4.019637in}{1.521178in}}%
\pgfpathlineto{\pgfqpoint{3.985674in}{1.541499in}}%
\pgfpathclose%
\pgfusepath{stroke,fill}%
\end{pgfscope}%
\begin{pgfscope}%
\pgfpathrectangle{\pgfqpoint{0.050000in}{0.050000in}}{\pgfqpoint{4.900000in}{4.900000in}}%
\pgfusepath{clip}%
\pgfsetbuttcap%
\pgfsetroundjoin%
\definecolor{currentfill}{rgb}{0.994464,0.994464,0.994464}%
\pgfsetfillcolor{currentfill}%
\pgfsetfillopacity{0.900000}%
\pgfsetlinewidth{0.501875pt}%
\definecolor{currentstroke}{rgb}{0.200000,0.200000,0.200000}%
\pgfsetstrokecolor{currentstroke}%
\pgfsetdash{}{0pt}%
\pgfpathmoveto{\pgfqpoint{3.709529in}{1.565263in}}%
\pgfpathlineto{\pgfqpoint{3.753361in}{1.572414in}}%
\pgfpathlineto{\pgfqpoint{3.786931in}{1.552543in}}%
\pgfpathlineto{\pgfqpoint{3.743074in}{1.545490in}}%
\pgfpathlineto{\pgfqpoint{3.709529in}{1.565263in}}%
\pgfpathclose%
\pgfusepath{stroke,fill}%
\end{pgfscope}%
\begin{pgfscope}%
\pgfpathrectangle{\pgfqpoint{0.050000in}{0.050000in}}{\pgfqpoint{4.900000in}{4.900000in}}%
\pgfusepath{clip}%
\pgfsetbuttcap%
\pgfsetroundjoin%
\definecolor{currentfill}{rgb}{0.898378,0.898378,0.898378}%
\pgfsetfillcolor{currentfill}%
\pgfsetfillopacity{0.900000}%
\pgfsetlinewidth{0.501875pt}%
\definecolor{currentstroke}{rgb}{0.200000,0.200000,0.200000}%
\pgfsetstrokecolor{currentstroke}%
\pgfsetdash{}{0pt}%
\pgfpathmoveto{\pgfqpoint{2.708889in}{2.006854in}}%
\pgfpathlineto{\pgfqpoint{2.753809in}{1.969663in}}%
\pgfpathlineto{\pgfqpoint{2.785802in}{1.949944in}}%
\pgfpathlineto{\pgfqpoint{2.740876in}{1.986872in}}%
\pgfpathlineto{\pgfqpoint{2.708889in}{2.006854in}}%
\pgfpathclose%
\pgfusepath{stroke,fill}%
\end{pgfscope}%
\begin{pgfscope}%
\pgfpathrectangle{\pgfqpoint{0.050000in}{0.050000in}}{\pgfqpoint{4.900000in}{4.900000in}}%
\pgfusepath{clip}%
\pgfsetbuttcap%
\pgfsetroundjoin%
\definecolor{currentfill}{rgb}{0.487043,0.487043,0.487043}%
\pgfsetfillcolor{currentfill}%
\pgfsetfillopacity{0.900000}%
\pgfsetlinewidth{0.501875pt}%
\definecolor{currentstroke}{rgb}{0.200000,0.200000,0.200000}%
\pgfsetstrokecolor{currentstroke}%
\pgfsetdash{}{0pt}%
\pgfpathmoveto{\pgfqpoint{1.472073in}{2.997742in}}%
\pgfpathlineto{\pgfqpoint{1.519574in}{2.952467in}}%
\pgfpathlineto{\pgfqpoint{1.549667in}{2.933147in}}%
\pgfpathlineto{\pgfqpoint{1.502157in}{2.978451in}}%
\pgfpathlineto{\pgfqpoint{1.472073in}{2.997742in}}%
\pgfpathclose%
\pgfusepath{stroke,fill}%
\end{pgfscope}%
\begin{pgfscope}%
\pgfpathrectangle{\pgfqpoint{0.050000in}{0.050000in}}{\pgfqpoint{4.900000in}{4.900000in}}%
\pgfusepath{clip}%
\pgfsetbuttcap%
\pgfsetroundjoin%
\definecolor{currentfill}{rgb}{0.757109,0.757109,0.757109}%
\pgfsetfillcolor{currentfill}%
\pgfsetfillopacity{0.900000}%
\pgfsetlinewidth{0.501875pt}%
\definecolor{currentstroke}{rgb}{0.200000,0.200000,0.200000}%
\pgfsetstrokecolor{currentstroke}%
\pgfsetdash{}{0pt}%
\pgfpathmoveto{\pgfqpoint{2.183338in}{2.422892in}}%
\pgfpathlineto{\pgfqpoint{2.229348in}{2.378865in}}%
\pgfpathlineto{\pgfqpoint{2.260533in}{2.358688in}}%
\pgfpathlineto{\pgfqpoint{2.214518in}{2.402726in}}%
\pgfpathlineto{\pgfqpoint{2.183338in}{2.422892in}}%
\pgfpathclose%
\pgfusepath{stroke,fill}%
\end{pgfscope}%
\begin{pgfscope}%
\pgfpathrectangle{\pgfqpoint{0.050000in}{0.050000in}}{\pgfqpoint{4.900000in}{4.900000in}}%
\pgfusepath{clip}%
\pgfsetbuttcap%
\pgfsetroundjoin%
\definecolor{currentfill}{rgb}{0.968627,0.968627,0.968627}%
\pgfsetfillcolor{currentfill}%
\pgfsetfillopacity{0.900000}%
\pgfsetlinewidth{0.501875pt}%
\definecolor{currentstroke}{rgb}{0.200000,0.200000,0.200000}%
\pgfsetstrokecolor{currentstroke}%
\pgfsetdash{}{0pt}%
\pgfpathmoveto{\pgfqpoint{3.202625in}{1.703431in}}%
\pgfpathlineto{\pgfqpoint{3.246838in}{1.689835in}}%
\pgfpathlineto{\pgfqpoint{3.279624in}{1.670837in}}%
\pgfpathlineto{\pgfqpoint{3.235393in}{1.684352in}}%
\pgfpathlineto{\pgfqpoint{3.202625in}{1.703431in}}%
\pgfpathclose%
\pgfusepath{stroke,fill}%
\end{pgfscope}%
\begin{pgfscope}%
\pgfpathrectangle{\pgfqpoint{0.050000in}{0.050000in}}{\pgfqpoint{4.900000in}{4.900000in}}%
\pgfusepath{clip}%
\pgfsetbuttcap%
\pgfsetroundjoin%
\definecolor{currentfill}{rgb}{0.976009,0.976009,0.976009}%
\pgfsetfillcolor{currentfill}%
\pgfsetfillopacity{0.900000}%
\pgfsetlinewidth{0.501875pt}%
\definecolor{currentstroke}{rgb}{0.200000,0.200000,0.200000}%
\pgfsetstrokecolor{currentstroke}%
\pgfsetdash{}{0pt}%
\pgfpathmoveto{\pgfqpoint{3.279624in}{1.670837in}}%
\pgfpathlineto{\pgfqpoint{3.323769in}{1.661281in}}%
\pgfpathlineto{\pgfqpoint{3.356679in}{1.642224in}}%
\pgfpathlineto{\pgfqpoint{3.312514in}{1.651753in}}%
\pgfpathlineto{\pgfqpoint{3.279624in}{1.670837in}}%
\pgfpathclose%
\pgfusepath{stroke,fill}%
\end{pgfscope}%
\begin{pgfscope}%
\pgfpathrectangle{\pgfqpoint{0.050000in}{0.050000in}}{\pgfqpoint{4.900000in}{4.900000in}}%
\pgfusepath{clip}%
\pgfsetbuttcap%
\pgfsetroundjoin%
\definecolor{currentfill}{rgb}{0.391603,0.391603,0.391603}%
\pgfsetfillcolor{currentfill}%
\pgfsetfillopacity{0.900000}%
\pgfsetlinewidth{0.501875pt}%
\definecolor{currentstroke}{rgb}{0.200000,0.200000,0.200000}%
\pgfsetstrokecolor{currentstroke}%
\pgfsetdash{}{0pt}%
\pgfpathmoveto{\pgfqpoint{1.209077in}{3.210937in}}%
\pgfpathlineto{\pgfqpoint{1.257138in}{3.165207in}}%
\pgfpathlineto{\pgfqpoint{1.286833in}{3.146210in}}%
\pgfpathlineto{\pgfqpoint{1.238761in}{3.191970in}}%
\pgfpathlineto{\pgfqpoint{1.209077in}{3.210937in}}%
\pgfpathclose%
\pgfusepath{stroke,fill}%
\end{pgfscope}%
\begin{pgfscope}%
\pgfpathrectangle{\pgfqpoint{0.050000in}{0.050000in}}{\pgfqpoint{4.900000in}{4.900000in}}%
\pgfusepath{clip}%
\pgfsetbuttcap%
\pgfsetroundjoin%
\definecolor{currentfill}{rgb}{0.963091,0.963091,0.963091}%
\pgfsetfillcolor{currentfill}%
\pgfsetfillopacity{0.900000}%
\pgfsetlinewidth{0.501875pt}%
\definecolor{currentstroke}{rgb}{0.200000,0.200000,0.200000}%
\pgfsetstrokecolor{currentstroke}%
\pgfsetdash{}{0pt}%
\pgfpathmoveto{\pgfqpoint{3.125667in}{1.740265in}}%
\pgfpathlineto{\pgfqpoint{3.169960in}{1.722424in}}%
\pgfpathlineto{\pgfqpoint{3.202625in}{1.703431in}}%
\pgfpathlineto{\pgfqpoint{3.158315in}{1.721134in}}%
\pgfpathlineto{\pgfqpoint{3.125667in}{1.740265in}}%
\pgfpathclose%
\pgfusepath{stroke,fill}%
\end{pgfscope}%
\begin{pgfscope}%
\pgfpathrectangle{\pgfqpoint{0.050000in}{0.050000in}}{\pgfqpoint{4.900000in}{4.900000in}}%
\pgfusepath{clip}%
\pgfsetbuttcap%
\pgfsetroundjoin%
\definecolor{currentfill}{rgb}{0.657809,0.657809,0.657809}%
\pgfsetfillcolor{currentfill}%
\pgfsetfillopacity{0.900000}%
\pgfsetlinewidth{0.501875pt}%
\definecolor{currentstroke}{rgb}{0.200000,0.200000,0.200000}%
\pgfsetstrokecolor{currentstroke}%
\pgfsetdash{}{0pt}%
\pgfpathmoveto{\pgfqpoint{1.920444in}{2.635976in}}%
\pgfpathlineto{\pgfqpoint{1.967014in}{2.591451in}}%
\pgfpathlineto{\pgfqpoint{1.997801in}{2.571573in}}%
\pgfpathlineto{\pgfqpoint{1.951224in}{2.616125in}}%
\pgfpathlineto{\pgfqpoint{1.920444in}{2.635976in}}%
\pgfpathclose%
\pgfusepath{stroke,fill}%
\end{pgfscope}%
\begin{pgfscope}%
\pgfpathrectangle{\pgfqpoint{0.050000in}{0.050000in}}{\pgfqpoint{4.900000in}{4.900000in}}%
\pgfusepath{clip}%
\pgfsetbuttcap%
\pgfsetroundjoin%
\definecolor{currentfill}{rgb}{0.979700,0.979700,0.979700}%
\pgfsetfillcolor{currentfill}%
\pgfsetfillopacity{0.900000}%
\pgfsetlinewidth{0.501875pt}%
\definecolor{currentstroke}{rgb}{0.200000,0.200000,0.200000}%
\pgfsetstrokecolor{currentstroke}%
\pgfsetdash{}{0pt}%
\pgfpathmoveto{\pgfqpoint{3.356679in}{1.642224in}}%
\pgfpathlineto{\pgfqpoint{3.400765in}{1.636402in}}%
\pgfpathlineto{\pgfqpoint{3.433799in}{1.617244in}}%
\pgfpathlineto{\pgfqpoint{3.389692in}{1.623083in}}%
\pgfpathlineto{\pgfqpoint{3.356679in}{1.642224in}}%
\pgfpathclose%
\pgfusepath{stroke,fill}%
\end{pgfscope}%
\begin{pgfscope}%
\pgfpathrectangle{\pgfqpoint{0.050000in}{0.050000in}}{\pgfqpoint{4.900000in}{4.900000in}}%
\pgfusepath{clip}%
\pgfsetbuttcap%
\pgfsetroundjoin%
\definecolor{currentfill}{rgb}{0.953864,0.953864,0.953864}%
\pgfsetfillcolor{currentfill}%
\pgfsetfillopacity{0.900000}%
\pgfsetlinewidth{0.501875pt}%
\definecolor{currentstroke}{rgb}{0.200000,0.200000,0.200000}%
\pgfsetstrokecolor{currentstroke}%
\pgfsetdash{}{0pt}%
\pgfpathmoveto{\pgfqpoint{3.048733in}{1.781471in}}%
\pgfpathlineto{\pgfqpoint{3.093121in}{1.759313in}}%
\pgfpathlineto{\pgfqpoint{3.125667in}{1.740265in}}%
\pgfpathlineto{\pgfqpoint{3.081264in}{1.762230in}}%
\pgfpathlineto{\pgfqpoint{3.048733in}{1.781471in}}%
\pgfpathclose%
\pgfusepath{stroke,fill}%
\end{pgfscope}%
\begin{pgfscope}%
\pgfpathrectangle{\pgfqpoint{0.050000in}{0.050000in}}{\pgfqpoint{4.900000in}{4.900000in}}%
\pgfusepath{clip}%
\pgfsetbuttcap%
\pgfsetroundjoin%
\definecolor{currentfill}{rgb}{0.996309,0.996309,0.996309}%
\pgfsetfillcolor{currentfill}%
\pgfsetfillopacity{0.900000}%
\pgfsetlinewidth{0.501875pt}%
\definecolor{currentstroke}{rgb}{0.200000,0.200000,0.200000}%
\pgfsetstrokecolor{currentstroke}%
\pgfsetdash{}{0pt}%
\pgfpathmoveto{\pgfqpoint{3.786931in}{1.552543in}}%
\pgfpathlineto{\pgfqpoint{3.830716in}{1.561171in}}%
\pgfpathlineto{\pgfqpoint{3.864417in}{1.541141in}}%
\pgfpathlineto{\pgfqpoint{3.820606in}{1.532607in}}%
\pgfpathlineto{\pgfqpoint{3.786931in}{1.552543in}}%
\pgfpathclose%
\pgfusepath{stroke,fill}%
\end{pgfscope}%
\begin{pgfscope}%
\pgfpathrectangle{\pgfqpoint{0.050000in}{0.050000in}}{\pgfqpoint{4.900000in}{4.900000in}}%
\pgfusepath{clip}%
\pgfsetbuttcap%
\pgfsetroundjoin%
\definecolor{currentfill}{rgb}{1.000000,1.000000,1.000000}%
\pgfsetfillcolor{currentfill}%
\pgfsetfillopacity{0.900000}%
\pgfsetlinewidth{0.501875pt}%
\definecolor{currentstroke}{rgb}{0.200000,0.200000,0.200000}%
\pgfsetstrokecolor{currentstroke}%
\pgfsetdash{}{0pt}%
\pgfpathmoveto{\pgfqpoint{4.063270in}{1.532595in}}%
\pgfpathlineto{\pgfqpoint{4.106820in}{1.544561in}}%
\pgfpathlineto{\pgfqpoint{4.140942in}{1.524038in}}%
\pgfpathlineto{\pgfqpoint{4.097366in}{1.512139in}}%
\pgfpathlineto{\pgfqpoint{4.063270in}{1.532595in}}%
\pgfpathclose%
\pgfusepath{stroke,fill}%
\end{pgfscope}%
\begin{pgfscope}%
\pgfpathrectangle{\pgfqpoint{0.050000in}{0.050000in}}{\pgfqpoint{4.900000in}{4.900000in}}%
\pgfusepath{clip}%
\pgfsetbuttcap%
\pgfsetroundjoin%
\definecolor{currentfill}{rgb}{0.551634,0.551634,0.551634}%
\pgfsetfillcolor{currentfill}%
\pgfsetfillopacity{0.900000}%
\pgfsetlinewidth{0.501875pt}%
\definecolor{currentstroke}{rgb}{0.200000,0.200000,0.200000}%
\pgfsetstrokecolor{currentstroke}%
\pgfsetdash{}{0pt}%
\pgfpathmoveto{\pgfqpoint{1.657461in}{2.849157in}}%
\pgfpathlineto{\pgfqpoint{1.704591in}{2.804176in}}%
\pgfpathlineto{\pgfqpoint{1.734981in}{2.784620in}}%
\pgfpathlineto{\pgfqpoint{1.687842in}{2.829629in}}%
\pgfpathlineto{\pgfqpoint{1.657461in}{2.849157in}}%
\pgfpathclose%
\pgfusepath{stroke,fill}%
\end{pgfscope}%
\begin{pgfscope}%
\pgfpathrectangle{\pgfqpoint{0.050000in}{0.050000in}}{\pgfqpoint{4.900000in}{4.900000in}}%
\pgfusepath{clip}%
\pgfsetbuttcap%
\pgfsetroundjoin%
\definecolor{currentfill}{rgb}{0.983391,0.983391,0.983391}%
\pgfsetfillcolor{currentfill}%
\pgfsetfillopacity{0.900000}%
\pgfsetlinewidth{0.501875pt}%
\definecolor{currentstroke}{rgb}{0.200000,0.200000,0.200000}%
\pgfsetstrokecolor{currentstroke}%
\pgfsetdash{}{0pt}%
\pgfpathmoveto{\pgfqpoint{3.433799in}{1.617244in}}%
\pgfpathlineto{\pgfqpoint{3.477834in}{1.614786in}}%
\pgfpathlineto{\pgfqpoint{3.510994in}{1.595495in}}%
\pgfpathlineto{\pgfqpoint{3.466937in}{1.598003in}}%
\pgfpathlineto{\pgfqpoint{3.433799in}{1.617244in}}%
\pgfpathclose%
\pgfusepath{stroke,fill}%
\end{pgfscope}%
\begin{pgfscope}%
\pgfpathrectangle{\pgfqpoint{0.050000in}{0.050000in}}{\pgfqpoint{4.900000in}{4.900000in}}%
\pgfusepath{clip}%
\pgfsetbuttcap%
\pgfsetroundjoin%
\definecolor{currentfill}{rgb}{0.878570,0.878570,0.878570}%
\pgfsetfillcolor{currentfill}%
\pgfsetfillopacity{0.900000}%
\pgfsetlinewidth{0.501875pt}%
\definecolor{currentstroke}{rgb}{0.200000,0.200000,0.200000}%
\pgfsetstrokecolor{currentstroke}%
\pgfsetdash{}{0pt}%
\pgfpathmoveto{\pgfqpoint{2.631911in}{2.066351in}}%
\pgfpathlineto{\pgfqpoint{2.677002in}{2.026795in}}%
\pgfpathlineto{\pgfqpoint{2.708889in}{2.006854in}}%
\pgfpathlineto{\pgfqpoint{2.663793in}{2.046188in}}%
\pgfpathlineto{\pgfqpoint{2.631911in}{2.066351in}}%
\pgfpathclose%
\pgfusepath{stroke,fill}%
\end{pgfscope}%
\begin{pgfscope}%
\pgfpathrectangle{\pgfqpoint{0.050000in}{0.050000in}}{\pgfqpoint{4.900000in}{4.900000in}}%
\pgfusepath{clip}%
\pgfsetbuttcap%
\pgfsetroundjoin%
\definecolor{currentfill}{rgb}{0.808781,0.808781,0.808781}%
\pgfsetfillcolor{currentfill}%
\pgfsetfillopacity{0.900000}%
\pgfsetlinewidth{0.501875pt}%
\definecolor{currentstroke}{rgb}{0.200000,0.200000,0.200000}%
\pgfsetstrokecolor{currentstroke}%
\pgfsetdash{}{0pt}%
\pgfpathmoveto{\pgfqpoint{2.369023in}{2.274425in}}%
\pgfpathlineto{\pgfqpoint{2.414660in}{2.231106in}}%
\pgfpathlineto{\pgfqpoint{2.446143in}{2.210835in}}%
\pgfpathlineto{\pgfqpoint{2.400501in}{2.254101in}}%
\pgfpathlineto{\pgfqpoint{2.369023in}{2.274425in}}%
\pgfpathclose%
\pgfusepath{stroke,fill}%
\end{pgfscope}%
\begin{pgfscope}%
\pgfpathrectangle{\pgfqpoint{0.050000in}{0.050000in}}{\pgfqpoint{4.900000in}{4.900000in}}%
\pgfusepath{clip}%
\pgfsetbuttcap%
\pgfsetroundjoin%
\definecolor{currentfill}{rgb}{0.946482,0.946482,0.946482}%
\pgfsetfillcolor{currentfill}%
\pgfsetfillopacity{0.900000}%
\pgfsetlinewidth{0.501875pt}%
\definecolor{currentstroke}{rgb}{0.200000,0.200000,0.200000}%
\pgfsetstrokecolor{currentstroke}%
\pgfsetdash{}{0pt}%
\pgfpathmoveto{\pgfqpoint{2.971806in}{1.827023in}}%
\pgfpathlineto{\pgfqpoint{3.016305in}{1.800636in}}%
\pgfpathlineto{\pgfqpoint{3.048733in}{1.781471in}}%
\pgfpathlineto{\pgfqpoint{3.004222in}{1.807621in}}%
\pgfpathlineto{\pgfqpoint{2.971806in}{1.827023in}}%
\pgfpathclose%
\pgfusepath{stroke,fill}%
\end{pgfscope}%
\begin{pgfscope}%
\pgfpathrectangle{\pgfqpoint{0.050000in}{0.050000in}}{\pgfqpoint{4.900000in}{4.900000in}}%
\pgfusepath{clip}%
\pgfsetbuttcap%
\pgfsetroundjoin%
\definecolor{currentfill}{rgb}{0.452595,0.452595,0.452595}%
\pgfsetfillcolor{currentfill}%
\pgfsetfillopacity{0.900000}%
\pgfsetlinewidth{0.501875pt}%
\definecolor{currentstroke}{rgb}{0.200000,0.200000,0.200000}%
\pgfsetstrokecolor{currentstroke}%
\pgfsetdash{}{0pt}%
\pgfpathmoveto{\pgfqpoint{1.394390in}{3.062411in}}%
\pgfpathlineto{\pgfqpoint{1.442081in}{3.016974in}}%
\pgfpathlineto{\pgfqpoint{1.472073in}{2.997742in}}%
\pgfpathlineto{\pgfqpoint{1.424372in}{3.043208in}}%
\pgfpathlineto{\pgfqpoint{1.394390in}{3.062411in}}%
\pgfpathclose%
\pgfusepath{stroke,fill}%
\end{pgfscope}%
\begin{pgfscope}%
\pgfpathrectangle{\pgfqpoint{0.050000in}{0.050000in}}{\pgfqpoint{4.900000in}{4.900000in}}%
\pgfusepath{clip}%
\pgfsetbuttcap%
\pgfsetroundjoin%
\definecolor{currentfill}{rgb}{0.987082,0.987082,0.987082}%
\pgfsetfillcolor{currentfill}%
\pgfsetfillopacity{0.900000}%
\pgfsetlinewidth{0.501875pt}%
\definecolor{currentstroke}{rgb}{0.200000,0.200000,0.200000}%
\pgfsetstrokecolor{currentstroke}%
\pgfsetdash{}{0pt}%
\pgfpathmoveto{\pgfqpoint{3.510994in}{1.595495in}}%
\pgfpathlineto{\pgfqpoint{3.554982in}{1.596003in}}%
\pgfpathlineto{\pgfqpoint{3.588270in}{1.576559in}}%
\pgfpathlineto{\pgfqpoint{3.544259in}{1.576125in}}%
\pgfpathlineto{\pgfqpoint{3.510994in}{1.595495in}}%
\pgfpathclose%
\pgfusepath{stroke,fill}%
\end{pgfscope}%
\begin{pgfscope}%
\pgfpathrectangle{\pgfqpoint{0.050000in}{0.050000in}}{\pgfqpoint{4.900000in}{4.900000in}}%
\pgfusepath{clip}%
\pgfsetbuttcap%
\pgfsetroundjoin%
\definecolor{currentfill}{rgb}{0.729781,0.729781,0.729781}%
\pgfsetfillcolor{currentfill}%
\pgfsetfillopacity{0.900000}%
\pgfsetlinewidth{0.501875pt}%
\definecolor{currentstroke}{rgb}{0.200000,0.200000,0.200000}%
\pgfsetstrokecolor{currentstroke}%
\pgfsetdash{}{0pt}%
\pgfpathmoveto{\pgfqpoint{2.106055in}{2.487219in}}%
\pgfpathlineto{\pgfqpoint{2.152254in}{2.442999in}}%
\pgfpathlineto{\pgfqpoint{2.183338in}{2.422892in}}%
\pgfpathlineto{\pgfqpoint{2.137134in}{2.467133in}}%
\pgfpathlineto{\pgfqpoint{2.106055in}{2.487219in}}%
\pgfpathclose%
\pgfusepath{stroke,fill}%
\end{pgfscope}%
\begin{pgfscope}%
\pgfpathrectangle{\pgfqpoint{0.050000in}{0.050000in}}{\pgfqpoint{4.900000in}{4.900000in}}%
\pgfusepath{clip}%
\pgfsetbuttcap%
\pgfsetroundjoin%
\definecolor{currentfill}{rgb}{0.996309,0.996309,0.996309}%
\pgfsetfillcolor{currentfill}%
\pgfsetfillopacity{0.900000}%
\pgfsetlinewidth{0.501875pt}%
\definecolor{currentstroke}{rgb}{0.200000,0.200000,0.200000}%
\pgfsetstrokecolor{currentstroke}%
\pgfsetdash{}{0pt}%
\pgfpathmoveto{\pgfqpoint{3.864417in}{1.541141in}}%
\pgfpathlineto{\pgfqpoint{3.908155in}{1.550950in}}%
\pgfpathlineto{\pgfqpoint{3.941986in}{1.530771in}}%
\pgfpathlineto{\pgfqpoint{3.898223in}{1.521049in}}%
\pgfpathlineto{\pgfqpoint{3.864417in}{1.541141in}}%
\pgfpathclose%
\pgfusepath{stroke,fill}%
\end{pgfscope}%
\begin{pgfscope}%
\pgfpathrectangle{\pgfqpoint{0.050000in}{0.050000in}}{\pgfqpoint{4.900000in}{4.900000in}}%
\pgfusepath{clip}%
\pgfsetbuttcap%
\pgfsetroundjoin%
\definecolor{currentfill}{rgb}{0.363183,0.363183,0.363183}%
\pgfsetfillcolor{currentfill}%
\pgfsetfillopacity{0.900000}%
\pgfsetlinewidth{0.501875pt}%
\definecolor{currentstroke}{rgb}{0.200000,0.200000,0.200000}%
\pgfsetstrokecolor{currentstroke}%
\pgfsetdash{}{0pt}%
\pgfpathmoveto{\pgfqpoint{1.131232in}{3.275739in}}%
\pgfpathlineto{\pgfqpoint{1.179484in}{3.229847in}}%
\pgfpathlineto{\pgfqpoint{1.209077in}{3.210937in}}%
\pgfpathlineto{\pgfqpoint{1.160813in}{3.256861in}}%
\pgfpathlineto{\pgfqpoint{1.131232in}{3.275739in}}%
\pgfpathclose%
\pgfusepath{stroke,fill}%
\end{pgfscope}%
\begin{pgfscope}%
\pgfpathrectangle{\pgfqpoint{0.050000in}{0.050000in}}{\pgfqpoint{4.900000in}{4.900000in}}%
\pgfusepath{clip}%
\pgfsetbuttcap%
\pgfsetroundjoin%
\definecolor{currentfill}{rgb}{0.932334,0.932334,0.932334}%
\pgfsetfillcolor{currentfill}%
\pgfsetfillopacity{0.900000}%
\pgfsetlinewidth{0.501875pt}%
\definecolor{currentstroke}{rgb}{0.200000,0.200000,0.200000}%
\pgfsetstrokecolor{currentstroke}%
\pgfsetdash{}{0pt}%
\pgfpathmoveto{\pgfqpoint{2.894867in}{1.876718in}}%
\pgfpathlineto{\pgfqpoint{2.939492in}{1.846358in}}%
\pgfpathlineto{\pgfqpoint{2.971806in}{1.827023in}}%
\pgfpathlineto{\pgfqpoint{2.927169in}{1.857116in}}%
\pgfpathlineto{\pgfqpoint{2.894867in}{1.876718in}}%
\pgfpathclose%
\pgfusepath{stroke,fill}%
\end{pgfscope}%
\begin{pgfscope}%
\pgfpathrectangle{\pgfqpoint{0.050000in}{0.050000in}}{\pgfqpoint{4.900000in}{4.900000in}}%
\pgfusepath{clip}%
\pgfsetbuttcap%
\pgfsetroundjoin%
\definecolor{currentfill}{rgb}{1.000000,1.000000,1.000000}%
\pgfsetfillcolor{currentfill}%
\pgfsetfillopacity{0.900000}%
\pgfsetlinewidth{0.501875pt}%
\definecolor{currentstroke}{rgb}{0.200000,0.200000,0.200000}%
\pgfsetstrokecolor{currentstroke}%
\pgfsetdash{}{0pt}%
\pgfpathmoveto{\pgfqpoint{4.140942in}{1.524038in}}%
\pgfpathlineto{\pgfqpoint{4.184430in}{1.536299in}}%
\pgfpathlineto{\pgfqpoint{4.218684in}{1.515652in}}%
\pgfpathlineto{\pgfqpoint{4.175170in}{1.503454in}}%
\pgfpathlineto{\pgfqpoint{4.140942in}{1.524038in}}%
\pgfpathclose%
\pgfusepath{stroke,fill}%
\end{pgfscope}%
\begin{pgfscope}%
\pgfpathrectangle{\pgfqpoint{0.050000in}{0.050000in}}{\pgfqpoint{4.900000in}{4.900000in}}%
\pgfusepath{clip}%
\pgfsetbuttcap%
\pgfsetroundjoin%
\definecolor{currentfill}{rgb}{0.619423,0.619423,0.619423}%
\pgfsetfillcolor{currentfill}%
\pgfsetfillopacity{0.900000}%
\pgfsetlinewidth{0.501875pt}%
\definecolor{currentstroke}{rgb}{0.200000,0.200000,0.200000}%
\pgfsetstrokecolor{currentstroke}%
\pgfsetdash{}{0pt}%
\pgfpathmoveto{\pgfqpoint{1.842998in}{2.700454in}}%
\pgfpathlineto{\pgfqpoint{1.889757in}{2.655767in}}%
\pgfpathlineto{\pgfqpoint{1.920444in}{2.635976in}}%
\pgfpathlineto{\pgfqpoint{1.873677in}{2.680690in}}%
\pgfpathlineto{\pgfqpoint{1.842998in}{2.700454in}}%
\pgfpathclose%
\pgfusepath{stroke,fill}%
\end{pgfscope}%
\begin{pgfscope}%
\pgfpathrectangle{\pgfqpoint{0.050000in}{0.050000in}}{\pgfqpoint{4.900000in}{4.900000in}}%
\pgfusepath{clip}%
\pgfsetbuttcap%
\pgfsetroundjoin%
\definecolor{currentfill}{rgb}{0.990773,0.990773,0.990773}%
\pgfsetfillcolor{currentfill}%
\pgfsetfillopacity{0.900000}%
\pgfsetlinewidth{0.501875pt}%
\definecolor{currentstroke}{rgb}{0.200000,0.200000,0.200000}%
\pgfsetstrokecolor{currentstroke}%
\pgfsetdash{}{0pt}%
\pgfpathmoveto{\pgfqpoint{3.588270in}{1.576559in}}%
\pgfpathlineto{\pgfqpoint{3.632213in}{1.579628in}}%
\pgfpathlineto{\pgfqpoint{3.665629in}{1.560021in}}%
\pgfpathlineto{\pgfqpoint{3.621662in}{1.557040in}}%
\pgfpathlineto{\pgfqpoint{3.588270in}{1.576559in}}%
\pgfpathclose%
\pgfusepath{stroke,fill}%
\end{pgfscope}%
\begin{pgfscope}%
\pgfpathrectangle{\pgfqpoint{0.050000in}{0.050000in}}{\pgfqpoint{4.900000in}{4.900000in}}%
\pgfusepath{clip}%
\pgfsetbuttcap%
\pgfsetroundjoin%
\definecolor{currentfill}{rgb}{0.858762,0.858762,0.858762}%
\pgfsetfillcolor{currentfill}%
\pgfsetfillopacity{0.900000}%
\pgfsetlinewidth{0.501875pt}%
\definecolor{currentstroke}{rgb}{0.200000,0.200000,0.200000}%
\pgfsetstrokecolor{currentstroke}%
\pgfsetdash{}{0pt}%
\pgfpathmoveto{\pgfqpoint{2.554858in}{2.127774in}}%
\pgfpathlineto{\pgfqpoint{2.600128in}{2.086475in}}%
\pgfpathlineto{\pgfqpoint{2.631911in}{2.066351in}}%
\pgfpathlineto{\pgfqpoint{2.586637in}{2.107482in}}%
\pgfpathlineto{\pgfqpoint{2.554858in}{2.127774in}}%
\pgfpathclose%
\pgfusepath{stroke,fill}%
\end{pgfscope}%
\begin{pgfscope}%
\pgfpathrectangle{\pgfqpoint{0.050000in}{0.050000in}}{\pgfqpoint{4.900000in}{4.900000in}}%
\pgfusepath{clip}%
\pgfsetbuttcap%
\pgfsetroundjoin%
\definecolor{currentfill}{rgb}{0.517186,0.517186,0.517186}%
\pgfsetfillcolor{currentfill}%
\pgfsetfillopacity{0.900000}%
\pgfsetlinewidth{0.501875pt}%
\definecolor{currentstroke}{rgb}{0.200000,0.200000,0.200000}%
\pgfsetstrokecolor{currentstroke}%
\pgfsetdash{}{0pt}%
\pgfpathmoveto{\pgfqpoint{1.579853in}{2.913767in}}%
\pgfpathlineto{\pgfqpoint{1.627173in}{2.868625in}}%
\pgfpathlineto{\pgfqpoint{1.657461in}{2.849157in}}%
\pgfpathlineto{\pgfqpoint{1.610132in}{2.894328in}}%
\pgfpathlineto{\pgfqpoint{1.579853in}{2.913767in}}%
\pgfpathclose%
\pgfusepath{stroke,fill}%
\end{pgfscope}%
\begin{pgfscope}%
\pgfpathrectangle{\pgfqpoint{0.050000in}{0.050000in}}{\pgfqpoint{4.900000in}{4.900000in}}%
\pgfusepath{clip}%
\pgfsetbuttcap%
\pgfsetroundjoin%
\definecolor{currentfill}{rgb}{0.784667,0.784667,0.784667}%
\pgfsetfillcolor{currentfill}%
\pgfsetfillopacity{0.900000}%
\pgfsetlinewidth{0.501875pt}%
\definecolor{currentstroke}{rgb}{0.200000,0.200000,0.200000}%
\pgfsetstrokecolor{currentstroke}%
\pgfsetdash{}{0pt}%
\pgfpathmoveto{\pgfqpoint{2.291814in}{2.338456in}}%
\pgfpathlineto{\pgfqpoint{2.337641in}{2.294698in}}%
\pgfpathlineto{\pgfqpoint{2.369023in}{2.274425in}}%
\pgfpathlineto{\pgfqpoint{2.323192in}{2.318167in}}%
\pgfpathlineto{\pgfqpoint{2.291814in}{2.338456in}}%
\pgfpathclose%
\pgfusepath{stroke,fill}%
\end{pgfscope}%
\begin{pgfscope}%
\pgfpathrectangle{\pgfqpoint{0.050000in}{0.050000in}}{\pgfqpoint{4.900000in}{4.900000in}}%
\pgfusepath{clip}%
\pgfsetbuttcap%
\pgfsetroundjoin%
\definecolor{currentfill}{rgb}{0.998155,0.998155,0.998155}%
\pgfsetfillcolor{currentfill}%
\pgfsetfillopacity{0.900000}%
\pgfsetlinewidth{0.501875pt}%
\definecolor{currentstroke}{rgb}{0.200000,0.200000,0.200000}%
\pgfsetstrokecolor{currentstroke}%
\pgfsetdash{}{0pt}%
\pgfpathmoveto{\pgfqpoint{3.941986in}{1.530771in}}%
\pgfpathlineto{\pgfqpoint{3.985674in}{1.541499in}}%
\pgfpathlineto{\pgfqpoint{4.019637in}{1.521178in}}%
\pgfpathlineto{\pgfqpoint{3.975924in}{1.510530in}}%
\pgfpathlineto{\pgfqpoint{3.941986in}{1.530771in}}%
\pgfpathclose%
\pgfusepath{stroke,fill}%
\end{pgfscope}%
\begin{pgfscope}%
\pgfpathrectangle{\pgfqpoint{0.050000in}{0.050000in}}{\pgfqpoint{4.900000in}{4.900000in}}%
\pgfusepath{clip}%
\pgfsetbuttcap%
\pgfsetroundjoin%
\definecolor{currentfill}{rgb}{0.915356,0.915356,0.915356}%
\pgfsetfillcolor{currentfill}%
\pgfsetfillopacity{0.900000}%
\pgfsetlinewidth{0.501875pt}%
\definecolor{currentstroke}{rgb}{0.200000,0.200000,0.200000}%
\pgfsetstrokecolor{currentstroke}%
\pgfsetdash{}{0pt}%
\pgfpathmoveto{\pgfqpoint{2.817895in}{1.930175in}}%
\pgfpathlineto{\pgfqpoint{2.862664in}{1.896261in}}%
\pgfpathlineto{\pgfqpoint{2.894867in}{1.876718in}}%
\pgfpathlineto{\pgfqpoint{2.850089in}{1.910357in}}%
\pgfpathlineto{\pgfqpoint{2.817895in}{1.930175in}}%
\pgfpathclose%
\pgfusepath{stroke,fill}%
\end{pgfscope}%
\begin{pgfscope}%
\pgfpathrectangle{\pgfqpoint{0.050000in}{0.050000in}}{\pgfqpoint{4.900000in}{4.900000in}}%
\pgfusepath{clip}%
\pgfsetbuttcap%
\pgfsetroundjoin%
\definecolor{currentfill}{rgb}{0.424083,0.424083,0.424083}%
\pgfsetfillcolor{currentfill}%
\pgfsetfillopacity{0.900000}%
\pgfsetlinewidth{0.501875pt}%
\definecolor{currentstroke}{rgb}{0.200000,0.200000,0.200000}%
\pgfsetstrokecolor{currentstroke}%
\pgfsetdash{}{0pt}%
\pgfpathmoveto{\pgfqpoint{1.316620in}{3.127154in}}%
\pgfpathlineto{\pgfqpoint{1.364501in}{3.081555in}}%
\pgfpathlineto{\pgfqpoint{1.394390in}{3.062411in}}%
\pgfpathlineto{\pgfqpoint{1.346498in}{3.108039in}}%
\pgfpathlineto{\pgfqpoint{1.316620in}{3.127154in}}%
\pgfpathclose%
\pgfusepath{stroke,fill}%
\end{pgfscope}%
\begin{pgfscope}%
\pgfpathrectangle{\pgfqpoint{0.050000in}{0.050000in}}{\pgfqpoint{4.900000in}{4.900000in}}%
\pgfusepath{clip}%
\pgfsetbuttcap%
\pgfsetroundjoin%
\definecolor{currentfill}{rgb}{0.992618,0.992618,0.992618}%
\pgfsetfillcolor{currentfill}%
\pgfsetfillopacity{0.900000}%
\pgfsetlinewidth{0.501875pt}%
\definecolor{currentstroke}{rgb}{0.200000,0.200000,0.200000}%
\pgfsetstrokecolor{currentstroke}%
\pgfsetdash{}{0pt}%
\pgfpathmoveto{\pgfqpoint{3.665629in}{1.560021in}}%
\pgfpathlineto{\pgfqpoint{3.709529in}{1.565263in}}%
\pgfpathlineto{\pgfqpoint{3.743074in}{1.545490in}}%
\pgfpathlineto{\pgfqpoint{3.699150in}{1.540343in}}%
\pgfpathlineto{\pgfqpoint{3.665629in}{1.560021in}}%
\pgfpathclose%
\pgfusepath{stroke,fill}%
\end{pgfscope}%
\begin{pgfscope}%
\pgfpathrectangle{\pgfqpoint{0.050000in}{0.050000in}}{\pgfqpoint{4.900000in}{4.900000in}}%
\pgfusepath{clip}%
\pgfsetbuttcap%
\pgfsetroundjoin%
\definecolor{currentfill}{rgb}{1.000000,1.000000,1.000000}%
\pgfsetfillcolor{currentfill}%
\pgfsetfillopacity{0.900000}%
\pgfsetlinewidth{0.501875pt}%
\definecolor{currentstroke}{rgb}{0.200000,0.200000,0.200000}%
\pgfsetstrokecolor{currentstroke}%
\pgfsetdash{}{0pt}%
\pgfpathmoveto{\pgfqpoint{4.218684in}{1.515652in}}%
\pgfpathlineto{\pgfqpoint{4.262107in}{1.528041in}}%
\pgfpathlineto{\pgfqpoint{4.296493in}{1.507272in}}%
\pgfpathlineto{\pgfqpoint{4.253045in}{1.494941in}}%
\pgfpathlineto{\pgfqpoint{4.218684in}{1.515652in}}%
\pgfpathclose%
\pgfusepath{stroke,fill}%
\end{pgfscope}%
\begin{pgfscope}%
\pgfpathrectangle{\pgfqpoint{0.050000in}{0.050000in}}{\pgfqpoint{4.900000in}{4.900000in}}%
\pgfusepath{clip}%
\pgfsetbuttcap%
\pgfsetroundjoin%
\definecolor{currentfill}{rgb}{0.691396,0.691396,0.691396}%
\pgfsetfillcolor{currentfill}%
\pgfsetfillopacity{0.900000}%
\pgfsetlinewidth{0.501875pt}%
\definecolor{currentstroke}{rgb}{0.200000,0.200000,0.200000}%
\pgfsetstrokecolor{currentstroke}%
\pgfsetdash{}{0pt}%
\pgfpathmoveto{\pgfqpoint{2.028684in}{2.551634in}}%
\pgfpathlineto{\pgfqpoint{2.075071in}{2.507244in}}%
\pgfpathlineto{\pgfqpoint{2.106055in}{2.487219in}}%
\pgfpathlineto{\pgfqpoint{2.059661in}{2.531634in}}%
\pgfpathlineto{\pgfqpoint{2.028684in}{2.551634in}}%
\pgfpathclose%
\pgfusepath{stroke,fill}%
\end{pgfscope}%
\begin{pgfscope}%
\pgfpathrectangle{\pgfqpoint{0.050000in}{0.050000in}}{\pgfqpoint{4.900000in}{4.900000in}}%
\pgfusepath{clip}%
\pgfsetbuttcap%
\pgfsetroundjoin%
\definecolor{currentfill}{rgb}{0.586082,0.586082,0.586082}%
\pgfsetfillcolor{currentfill}%
\pgfsetfillopacity{0.900000}%
\pgfsetlinewidth{0.501875pt}%
\definecolor{currentstroke}{rgb}{0.200000,0.200000,0.200000}%
\pgfsetstrokecolor{currentstroke}%
\pgfsetdash{}{0pt}%
\pgfpathmoveto{\pgfqpoint{1.765464in}{2.765005in}}%
\pgfpathlineto{\pgfqpoint{1.812413in}{2.720157in}}%
\pgfpathlineto{\pgfqpoint{1.842998in}{2.700454in}}%
\pgfpathlineto{\pgfqpoint{1.796041in}{2.745329in}}%
\pgfpathlineto{\pgfqpoint{1.765464in}{2.765005in}}%
\pgfpathclose%
\pgfusepath{stroke,fill}%
\end{pgfscope}%
\begin{pgfscope}%
\pgfpathrectangle{\pgfqpoint{0.050000in}{0.050000in}}{\pgfqpoint{4.900000in}{4.900000in}}%
\pgfusepath{clip}%
\pgfsetbuttcap%
\pgfsetroundjoin%
\definecolor{currentfill}{rgb}{0.836340,0.836340,0.836340}%
\pgfsetfillcolor{currentfill}%
\pgfsetfillopacity{0.900000}%
\pgfsetlinewidth{0.501875pt}%
\definecolor{currentstroke}{rgb}{0.200000,0.200000,0.200000}%
\pgfsetstrokecolor{currentstroke}%
\pgfsetdash{}{0pt}%
\pgfpathmoveto{\pgfqpoint{2.477722in}{2.190520in}}%
\pgfpathlineto{\pgfqpoint{2.523177in}{2.148024in}}%
\pgfpathlineto{\pgfqpoint{2.554858in}{2.127774in}}%
\pgfpathlineto{\pgfqpoint{2.509399in}{2.170159in}}%
\pgfpathlineto{\pgfqpoint{2.477722in}{2.190520in}}%
\pgfpathclose%
\pgfusepath{stroke,fill}%
\end{pgfscope}%
\begin{pgfscope}%
\pgfpathrectangle{\pgfqpoint{0.050000in}{0.050000in}}{\pgfqpoint{4.900000in}{4.900000in}}%
\pgfusepath{clip}%
\pgfsetbuttcap%
\pgfsetroundjoin%
\definecolor{currentfill}{rgb}{0.998155,0.998155,0.998155}%
\pgfsetfillcolor{currentfill}%
\pgfsetfillopacity{0.900000}%
\pgfsetlinewidth{0.501875pt}%
\definecolor{currentstroke}{rgb}{0.200000,0.200000,0.200000}%
\pgfsetstrokecolor{currentstroke}%
\pgfsetdash{}{0pt}%
\pgfpathmoveto{\pgfqpoint{4.019637in}{1.521178in}}%
\pgfpathlineto{\pgfqpoint{4.063270in}{1.532595in}}%
\pgfpathlineto{\pgfqpoint{4.097366in}{1.512139in}}%
\pgfpathlineto{\pgfqpoint{4.053706in}{1.500796in}}%
\pgfpathlineto{\pgfqpoint{4.019637in}{1.521178in}}%
\pgfpathclose%
\pgfusepath{stroke,fill}%
\end{pgfscope}%
\begin{pgfscope}%
\pgfpathrectangle{\pgfqpoint{0.050000in}{0.050000in}}{\pgfqpoint{4.900000in}{4.900000in}}%
\pgfusepath{clip}%
\pgfsetbuttcap%
\pgfsetroundjoin%
\definecolor{currentfill}{rgb}{0.994464,0.994464,0.994464}%
\pgfsetfillcolor{currentfill}%
\pgfsetfillopacity{0.900000}%
\pgfsetlinewidth{0.501875pt}%
\definecolor{currentstroke}{rgb}{0.200000,0.200000,0.200000}%
\pgfsetstrokecolor{currentstroke}%
\pgfsetdash{}{0pt}%
\pgfpathmoveto{\pgfqpoint{3.743074in}{1.545490in}}%
\pgfpathlineto{\pgfqpoint{3.786931in}{1.552543in}}%
\pgfpathlineto{\pgfqpoint{3.820606in}{1.532607in}}%
\pgfpathlineto{\pgfqpoint{3.776725in}{1.525650in}}%
\pgfpathlineto{\pgfqpoint{3.743074in}{1.545490in}}%
\pgfpathclose%
\pgfusepath{stroke,fill}%
\end{pgfscope}%
\begin{pgfscope}%
\pgfpathrectangle{\pgfqpoint{0.050000in}{0.050000in}}{\pgfqpoint{4.900000in}{4.900000in}}%
\pgfusepath{clip}%
\pgfsetbuttcap%
\pgfsetroundjoin%
\definecolor{currentfill}{rgb}{0.898378,0.898378,0.898378}%
\pgfsetfillcolor{currentfill}%
\pgfsetfillopacity{0.900000}%
\pgfsetlinewidth{0.501875pt}%
\definecolor{currentstroke}{rgb}{0.200000,0.200000,0.200000}%
\pgfsetstrokecolor{currentstroke}%
\pgfsetdash{}{0pt}%
\pgfpathmoveto{\pgfqpoint{2.740876in}{1.986872in}}%
\pgfpathlineto{\pgfqpoint{2.785802in}{1.949944in}}%
\pgfpathlineto{\pgfqpoint{2.817895in}{1.930175in}}%
\pgfpathlineto{\pgfqpoint{2.772962in}{1.966845in}}%
\pgfpathlineto{\pgfqpoint{2.740876in}{1.986872in}}%
\pgfpathclose%
\pgfusepath{stroke,fill}%
\end{pgfscope}%
\begin{pgfscope}%
\pgfpathrectangle{\pgfqpoint{0.050000in}{0.050000in}}{\pgfqpoint{4.900000in}{4.900000in}}%
\pgfusepath{clip}%
\pgfsetbuttcap%
\pgfsetroundjoin%
\definecolor{currentfill}{rgb}{0.487043,0.487043,0.487043}%
\pgfsetfillcolor{currentfill}%
\pgfsetfillopacity{0.900000}%
\pgfsetlinewidth{0.501875pt}%
\definecolor{currentstroke}{rgb}{0.200000,0.200000,0.200000}%
\pgfsetstrokecolor{currentstroke}%
\pgfsetdash{}{0pt}%
\pgfpathmoveto{\pgfqpoint{1.502157in}{2.978451in}}%
\pgfpathlineto{\pgfqpoint{1.549667in}{2.933147in}}%
\pgfpathlineto{\pgfqpoint{1.579853in}{2.913767in}}%
\pgfpathlineto{\pgfqpoint{1.532333in}{2.959100in}}%
\pgfpathlineto{\pgfqpoint{1.502157in}{2.978451in}}%
\pgfpathclose%
\pgfusepath{stroke,fill}%
\end{pgfscope}%
\begin{pgfscope}%
\pgfpathrectangle{\pgfqpoint{0.050000in}{0.050000in}}{\pgfqpoint{4.900000in}{4.900000in}}%
\pgfusepath{clip}%
\pgfsetbuttcap%
\pgfsetroundjoin%
\definecolor{currentfill}{rgb}{0.757109,0.757109,0.757109}%
\pgfsetfillcolor{currentfill}%
\pgfsetfillopacity{0.900000}%
\pgfsetlinewidth{0.501875pt}%
\definecolor{currentstroke}{rgb}{0.200000,0.200000,0.200000}%
\pgfsetstrokecolor{currentstroke}%
\pgfsetdash{}{0pt}%
\pgfpathmoveto{\pgfqpoint{2.214518in}{2.402726in}}%
\pgfpathlineto{\pgfqpoint{2.260533in}{2.358688in}}%
\pgfpathlineto{\pgfqpoint{2.291814in}{2.338456in}}%
\pgfpathlineto{\pgfqpoint{2.245794in}{2.382502in}}%
\pgfpathlineto{\pgfqpoint{2.214518in}{2.402726in}}%
\pgfpathclose%
\pgfusepath{stroke,fill}%
\end{pgfscope}%
\begin{pgfscope}%
\pgfpathrectangle{\pgfqpoint{0.050000in}{0.050000in}}{\pgfqpoint{4.900000in}{4.900000in}}%
\pgfusepath{clip}%
\pgfsetbuttcap%
\pgfsetroundjoin%
\definecolor{currentfill}{rgb}{0.968627,0.968627,0.968627}%
\pgfsetfillcolor{currentfill}%
\pgfsetfillopacity{0.900000}%
\pgfsetlinewidth{0.501875pt}%
\definecolor{currentstroke}{rgb}{0.200000,0.200000,0.200000}%
\pgfsetstrokecolor{currentstroke}%
\pgfsetdash{}{0pt}%
\pgfpathmoveto{\pgfqpoint{3.235393in}{1.684352in}}%
\pgfpathlineto{\pgfqpoint{3.279624in}{1.670837in}}%
\pgfpathlineto{\pgfqpoint{3.312514in}{1.651753in}}%
\pgfpathlineto{\pgfqpoint{3.268264in}{1.665188in}}%
\pgfpathlineto{\pgfqpoint{3.235393in}{1.684352in}}%
\pgfpathclose%
\pgfusepath{stroke,fill}%
\end{pgfscope}%
\begin{pgfscope}%
\pgfpathrectangle{\pgfqpoint{0.050000in}{0.050000in}}{\pgfqpoint{4.900000in}{4.900000in}}%
\pgfusepath{clip}%
\pgfsetbuttcap%
\pgfsetroundjoin%
\definecolor{currentfill}{rgb}{0.391603,0.391603,0.391603}%
\pgfsetfillcolor{currentfill}%
\pgfsetfillopacity{0.900000}%
\pgfsetlinewidth{0.501875pt}%
\definecolor{currentstroke}{rgb}{0.200000,0.200000,0.200000}%
\pgfsetstrokecolor{currentstroke}%
\pgfsetdash{}{0pt}%
\pgfpathmoveto{\pgfqpoint{1.238761in}{3.191970in}}%
\pgfpathlineto{\pgfqpoint{1.286833in}{3.146210in}}%
\pgfpathlineto{\pgfqpoint{1.316620in}{3.127154in}}%
\pgfpathlineto{\pgfqpoint{1.268536in}{3.172945in}}%
\pgfpathlineto{\pgfqpoint{1.238761in}{3.191970in}}%
\pgfpathclose%
\pgfusepath{stroke,fill}%
\end{pgfscope}%
\begin{pgfscope}%
\pgfpathrectangle{\pgfqpoint{0.050000in}{0.050000in}}{\pgfqpoint{4.900000in}{4.900000in}}%
\pgfusepath{clip}%
\pgfsetbuttcap%
\pgfsetroundjoin%
\definecolor{currentfill}{rgb}{0.974164,0.974164,0.974164}%
\pgfsetfillcolor{currentfill}%
\pgfsetfillopacity{0.900000}%
\pgfsetlinewidth{0.501875pt}%
\definecolor{currentstroke}{rgb}{0.200000,0.200000,0.200000}%
\pgfsetstrokecolor{currentstroke}%
\pgfsetdash{}{0pt}%
\pgfpathmoveto{\pgfqpoint{3.312514in}{1.651753in}}%
\pgfpathlineto{\pgfqpoint{3.356679in}{1.642224in}}%
\pgfpathlineto{\pgfqpoint{3.389692in}{1.623083in}}%
\pgfpathlineto{\pgfqpoint{3.345507in}{1.632582in}}%
\pgfpathlineto{\pgfqpoint{3.312514in}{1.651753in}}%
\pgfpathclose%
\pgfusepath{stroke,fill}%
\end{pgfscope}%
\begin{pgfscope}%
\pgfpathrectangle{\pgfqpoint{0.050000in}{0.050000in}}{\pgfqpoint{4.900000in}{4.900000in}}%
\pgfusepath{clip}%
\pgfsetbuttcap%
\pgfsetroundjoin%
\definecolor{currentfill}{rgb}{0.961246,0.961246,0.961246}%
\pgfsetfillcolor{currentfill}%
\pgfsetfillopacity{0.900000}%
\pgfsetlinewidth{0.501875pt}%
\definecolor{currentstroke}{rgb}{0.200000,0.200000,0.200000}%
\pgfsetstrokecolor{currentstroke}%
\pgfsetdash{}{0pt}%
\pgfpathmoveto{\pgfqpoint{3.158315in}{1.721134in}}%
\pgfpathlineto{\pgfqpoint{3.202625in}{1.703431in}}%
\pgfpathlineto{\pgfqpoint{3.235393in}{1.684352in}}%
\pgfpathlineto{\pgfqpoint{3.191066in}{1.701921in}}%
\pgfpathlineto{\pgfqpoint{3.158315in}{1.721134in}}%
\pgfpathclose%
\pgfusepath{stroke,fill}%
\end{pgfscope}%
\begin{pgfscope}%
\pgfpathrectangle{\pgfqpoint{0.050000in}{0.050000in}}{\pgfqpoint{4.900000in}{4.900000in}}%
\pgfusepath{clip}%
\pgfsetbuttcap%
\pgfsetroundjoin%
\definecolor{currentfill}{rgb}{0.657809,0.657809,0.657809}%
\pgfsetfillcolor{currentfill}%
\pgfsetfillopacity{0.900000}%
\pgfsetlinewidth{0.501875pt}%
\definecolor{currentstroke}{rgb}{0.200000,0.200000,0.200000}%
\pgfsetstrokecolor{currentstroke}%
\pgfsetdash{}{0pt}%
\pgfpathmoveto{\pgfqpoint{1.951224in}{2.616125in}}%
\pgfpathlineto{\pgfqpoint{1.997801in}{2.571573in}}%
\pgfpathlineto{\pgfqpoint{2.028684in}{2.551634in}}%
\pgfpathlineto{\pgfqpoint{1.982100in}{2.596212in}}%
\pgfpathlineto{\pgfqpoint{1.951224in}{2.616125in}}%
\pgfpathclose%
\pgfusepath{stroke,fill}%
\end{pgfscope}%
\begin{pgfscope}%
\pgfpathrectangle{\pgfqpoint{0.050000in}{0.050000in}}{\pgfqpoint{4.900000in}{4.900000in}}%
\pgfusepath{clip}%
\pgfsetbuttcap%
\pgfsetroundjoin%
\definecolor{currentfill}{rgb}{0.979700,0.979700,0.979700}%
\pgfsetfillcolor{currentfill}%
\pgfsetfillopacity{0.900000}%
\pgfsetlinewidth{0.501875pt}%
\definecolor{currentstroke}{rgb}{0.200000,0.200000,0.200000}%
\pgfsetstrokecolor{currentstroke}%
\pgfsetdash{}{0pt}%
\pgfpathmoveto{\pgfqpoint{3.389692in}{1.623083in}}%
\pgfpathlineto{\pgfqpoint{3.433799in}{1.617244in}}%
\pgfpathlineto{\pgfqpoint{3.466937in}{1.598003in}}%
\pgfpathlineto{\pgfqpoint{3.422809in}{1.603856in}}%
\pgfpathlineto{\pgfqpoint{3.389692in}{1.623083in}}%
\pgfpathclose%
\pgfusepath{stroke,fill}%
\end{pgfscope}%
\begin{pgfscope}%
\pgfpathrectangle{\pgfqpoint{0.050000in}{0.050000in}}{\pgfqpoint{4.900000in}{4.900000in}}%
\pgfusepath{clip}%
\pgfsetbuttcap%
\pgfsetroundjoin%
\definecolor{currentfill}{rgb}{0.996309,0.996309,0.996309}%
\pgfsetfillcolor{currentfill}%
\pgfsetfillopacity{0.900000}%
\pgfsetlinewidth{0.501875pt}%
\definecolor{currentstroke}{rgb}{0.200000,0.200000,0.200000}%
\pgfsetstrokecolor{currentstroke}%
\pgfsetdash{}{0pt}%
\pgfpathmoveto{\pgfqpoint{3.820606in}{1.532607in}}%
\pgfpathlineto{\pgfqpoint{3.864417in}{1.541141in}}%
\pgfpathlineto{\pgfqpoint{3.898223in}{1.521049in}}%
\pgfpathlineto{\pgfqpoint{3.854387in}{1.512607in}}%
\pgfpathlineto{\pgfqpoint{3.820606in}{1.532607in}}%
\pgfpathclose%
\pgfusepath{stroke,fill}%
\end{pgfscope}%
\begin{pgfscope}%
\pgfpathrectangle{\pgfqpoint{0.050000in}{0.050000in}}{\pgfqpoint{4.900000in}{4.900000in}}%
\pgfusepath{clip}%
\pgfsetbuttcap%
\pgfsetroundjoin%
\definecolor{currentfill}{rgb}{0.953864,0.953864,0.953864}%
\pgfsetfillcolor{currentfill}%
\pgfsetfillopacity{0.900000}%
\pgfsetlinewidth{0.501875pt}%
\definecolor{currentstroke}{rgb}{0.200000,0.200000,0.200000}%
\pgfsetstrokecolor{currentstroke}%
\pgfsetdash{}{0pt}%
\pgfpathmoveto{\pgfqpoint{3.081264in}{1.762230in}}%
\pgfpathlineto{\pgfqpoint{3.125667in}{1.740265in}}%
\pgfpathlineto{\pgfqpoint{3.158315in}{1.721134in}}%
\pgfpathlineto{\pgfqpoint{3.113897in}{1.742913in}}%
\pgfpathlineto{\pgfqpoint{3.081264in}{1.762230in}}%
\pgfpathclose%
\pgfusepath{stroke,fill}%
\end{pgfscope}%
\begin{pgfscope}%
\pgfpathrectangle{\pgfqpoint{0.050000in}{0.050000in}}{\pgfqpoint{4.900000in}{4.900000in}}%
\pgfusepath{clip}%
\pgfsetbuttcap%
\pgfsetroundjoin%
\definecolor{currentfill}{rgb}{1.000000,1.000000,1.000000}%
\pgfsetfillcolor{currentfill}%
\pgfsetfillopacity{0.900000}%
\pgfsetlinewidth{0.501875pt}%
\definecolor{currentstroke}{rgb}{0.200000,0.200000,0.200000}%
\pgfsetstrokecolor{currentstroke}%
\pgfsetdash{}{0pt}%
\pgfpathmoveto{\pgfqpoint{4.097366in}{1.512139in}}%
\pgfpathlineto{\pgfqpoint{4.140942in}{1.524038in}}%
\pgfpathlineto{\pgfqpoint{4.175170in}{1.503454in}}%
\pgfpathlineto{\pgfqpoint{4.131568in}{1.491621in}}%
\pgfpathlineto{\pgfqpoint{4.097366in}{1.512139in}}%
\pgfpathclose%
\pgfusepath{stroke,fill}%
\end{pgfscope}%
\begin{pgfscope}%
\pgfpathrectangle{\pgfqpoint{0.050000in}{0.050000in}}{\pgfqpoint{4.900000in}{4.900000in}}%
\pgfusepath{clip}%
\pgfsetbuttcap%
\pgfsetroundjoin%
\definecolor{currentfill}{rgb}{0.551634,0.551634,0.551634}%
\pgfsetfillcolor{currentfill}%
\pgfsetfillopacity{0.900000}%
\pgfsetlinewidth{0.501875pt}%
\definecolor{currentstroke}{rgb}{0.200000,0.200000,0.200000}%
\pgfsetstrokecolor{currentstroke}%
\pgfsetdash{}{0pt}%
\pgfpathmoveto{\pgfqpoint{1.687842in}{2.829629in}}%
\pgfpathlineto{\pgfqpoint{1.734981in}{2.784620in}}%
\pgfpathlineto{\pgfqpoint{1.765464in}{2.765005in}}%
\pgfpathlineto{\pgfqpoint{1.718317in}{2.810042in}}%
\pgfpathlineto{\pgfqpoint{1.687842in}{2.829629in}}%
\pgfpathclose%
\pgfusepath{stroke,fill}%
\end{pgfscope}%
\begin{pgfscope}%
\pgfpathrectangle{\pgfqpoint{0.050000in}{0.050000in}}{\pgfqpoint{4.900000in}{4.900000in}}%
\pgfusepath{clip}%
\pgfsetbuttcap%
\pgfsetroundjoin%
\definecolor{currentfill}{rgb}{0.983391,0.983391,0.983391}%
\pgfsetfillcolor{currentfill}%
\pgfsetfillopacity{0.900000}%
\pgfsetlinewidth{0.501875pt}%
\definecolor{currentstroke}{rgb}{0.200000,0.200000,0.200000}%
\pgfsetstrokecolor{currentstroke}%
\pgfsetdash{}{0pt}%
\pgfpathmoveto{\pgfqpoint{3.466937in}{1.598003in}}%
\pgfpathlineto{\pgfqpoint{3.510994in}{1.595495in}}%
\pgfpathlineto{\pgfqpoint{3.544259in}{1.576125in}}%
\pgfpathlineto{\pgfqpoint{3.500180in}{1.578680in}}%
\pgfpathlineto{\pgfqpoint{3.466937in}{1.598003in}}%
\pgfpathclose%
\pgfusepath{stroke,fill}%
\end{pgfscope}%
\begin{pgfscope}%
\pgfpathrectangle{\pgfqpoint{0.050000in}{0.050000in}}{\pgfqpoint{4.900000in}{4.900000in}}%
\pgfusepath{clip}%
\pgfsetbuttcap%
\pgfsetroundjoin%
\definecolor{currentfill}{rgb}{0.878570,0.878570,0.878570}%
\pgfsetfillcolor{currentfill}%
\pgfsetfillopacity{0.900000}%
\pgfsetlinewidth{0.501875pt}%
\definecolor{currentstroke}{rgb}{0.200000,0.200000,0.200000}%
\pgfsetstrokecolor{currentstroke}%
\pgfsetdash{}{0pt}%
\pgfpathmoveto{\pgfqpoint{2.663793in}{2.046188in}}%
\pgfpathlineto{\pgfqpoint{2.708889in}{2.006854in}}%
\pgfpathlineto{\pgfqpoint{2.740876in}{1.986872in}}%
\pgfpathlineto{\pgfqpoint{2.695774in}{2.025983in}}%
\pgfpathlineto{\pgfqpoint{2.663793in}{2.046188in}}%
\pgfpathclose%
\pgfusepath{stroke,fill}%
\end{pgfscope}%
\begin{pgfscope}%
\pgfpathrectangle{\pgfqpoint{0.050000in}{0.050000in}}{\pgfqpoint{4.900000in}{4.900000in}}%
\pgfusepath{clip}%
\pgfsetbuttcap%
\pgfsetroundjoin%
\definecolor{currentfill}{rgb}{0.808781,0.808781,0.808781}%
\pgfsetfillcolor{currentfill}%
\pgfsetfillopacity{0.900000}%
\pgfsetlinewidth{0.501875pt}%
\definecolor{currentstroke}{rgb}{0.200000,0.200000,0.200000}%
\pgfsetstrokecolor{currentstroke}%
\pgfsetdash{}{0pt}%
\pgfpathmoveto{\pgfqpoint{2.400501in}{2.254101in}}%
\pgfpathlineto{\pgfqpoint{2.446143in}{2.210835in}}%
\pgfpathlineto{\pgfqpoint{2.477722in}{2.190520in}}%
\pgfpathlineto{\pgfqpoint{2.432077in}{2.233727in}}%
\pgfpathlineto{\pgfqpoint{2.400501in}{2.254101in}}%
\pgfpathclose%
\pgfusepath{stroke,fill}%
\end{pgfscope}%
\begin{pgfscope}%
\pgfpathrectangle{\pgfqpoint{0.050000in}{0.050000in}}{\pgfqpoint{4.900000in}{4.900000in}}%
\pgfusepath{clip}%
\pgfsetbuttcap%
\pgfsetroundjoin%
\definecolor{currentfill}{rgb}{0.944637,0.944637,0.944637}%
\pgfsetfillcolor{currentfill}%
\pgfsetfillopacity{0.900000}%
\pgfsetlinewidth{0.501875pt}%
\definecolor{currentstroke}{rgb}{0.200000,0.200000,0.200000}%
\pgfsetstrokecolor{currentstroke}%
\pgfsetdash{}{0pt}%
\pgfpathmoveto{\pgfqpoint{3.004222in}{1.807621in}}%
\pgfpathlineto{\pgfqpoint{3.048733in}{1.781471in}}%
\pgfpathlineto{\pgfqpoint{3.081264in}{1.762230in}}%
\pgfpathlineto{\pgfqpoint{3.036739in}{1.788150in}}%
\pgfpathlineto{\pgfqpoint{3.004222in}{1.807621in}}%
\pgfpathclose%
\pgfusepath{stroke,fill}%
\end{pgfscope}%
\begin{pgfscope}%
\pgfpathrectangle{\pgfqpoint{0.050000in}{0.050000in}}{\pgfqpoint{4.900000in}{4.900000in}}%
\pgfusepath{clip}%
\pgfsetbuttcap%
\pgfsetroundjoin%
\definecolor{currentfill}{rgb}{0.452595,0.452595,0.452595}%
\pgfsetfillcolor{currentfill}%
\pgfsetfillopacity{0.900000}%
\pgfsetlinewidth{0.501875pt}%
\definecolor{currentstroke}{rgb}{0.200000,0.200000,0.200000}%
\pgfsetstrokecolor{currentstroke}%
\pgfsetdash{}{0pt}%
\pgfpathmoveto{\pgfqpoint{1.424372in}{3.043208in}}%
\pgfpathlineto{\pgfqpoint{1.472073in}{2.997742in}}%
\pgfpathlineto{\pgfqpoint{1.502157in}{2.978451in}}%
\pgfpathlineto{\pgfqpoint{1.454445in}{3.023946in}}%
\pgfpathlineto{\pgfqpoint{1.424372in}{3.043208in}}%
\pgfpathclose%
\pgfusepath{stroke,fill}%
\end{pgfscope}%
\begin{pgfscope}%
\pgfpathrectangle{\pgfqpoint{0.050000in}{0.050000in}}{\pgfqpoint{4.900000in}{4.900000in}}%
\pgfusepath{clip}%
\pgfsetbuttcap%
\pgfsetroundjoin%
\definecolor{currentfill}{rgb}{0.729781,0.729781,0.729781}%
\pgfsetfillcolor{currentfill}%
\pgfsetfillopacity{0.900000}%
\pgfsetlinewidth{0.501875pt}%
\definecolor{currentstroke}{rgb}{0.200000,0.200000,0.200000}%
\pgfsetstrokecolor{currentstroke}%
\pgfsetdash{}{0pt}%
\pgfpathmoveto{\pgfqpoint{2.137134in}{2.467133in}}%
\pgfpathlineto{\pgfqpoint{2.183338in}{2.422892in}}%
\pgfpathlineto{\pgfqpoint{2.214518in}{2.402726in}}%
\pgfpathlineto{\pgfqpoint{2.168308in}{2.446987in}}%
\pgfpathlineto{\pgfqpoint{2.137134in}{2.467133in}}%
\pgfpathclose%
\pgfusepath{stroke,fill}%
\end{pgfscope}%
\begin{pgfscope}%
\pgfpathrectangle{\pgfqpoint{0.050000in}{0.050000in}}{\pgfqpoint{4.900000in}{4.900000in}}%
\pgfusepath{clip}%
\pgfsetbuttcap%
\pgfsetroundjoin%
\definecolor{currentfill}{rgb}{0.987082,0.987082,0.987082}%
\pgfsetfillcolor{currentfill}%
\pgfsetfillopacity{0.900000}%
\pgfsetlinewidth{0.501875pt}%
\definecolor{currentstroke}{rgb}{0.200000,0.200000,0.200000}%
\pgfsetstrokecolor{currentstroke}%
\pgfsetdash{}{0pt}%
\pgfpathmoveto{\pgfqpoint{3.544259in}{1.576125in}}%
\pgfpathlineto{\pgfqpoint{3.588270in}{1.576559in}}%
\pgfpathlineto{\pgfqpoint{3.621662in}{1.557040in}}%
\pgfpathlineto{\pgfqpoint{3.577628in}{1.556677in}}%
\pgfpathlineto{\pgfqpoint{3.544259in}{1.576125in}}%
\pgfpathclose%
\pgfusepath{stroke,fill}%
\end{pgfscope}%
\begin{pgfscope}%
\pgfpathrectangle{\pgfqpoint{0.050000in}{0.050000in}}{\pgfqpoint{4.900000in}{4.900000in}}%
\pgfusepath{clip}%
\pgfsetbuttcap%
\pgfsetroundjoin%
\definecolor{currentfill}{rgb}{0.996309,0.996309,0.996309}%
\pgfsetfillcolor{currentfill}%
\pgfsetfillopacity{0.900000}%
\pgfsetlinewidth{0.501875pt}%
\definecolor{currentstroke}{rgb}{0.200000,0.200000,0.200000}%
\pgfsetstrokecolor{currentstroke}%
\pgfsetdash{}{0pt}%
\pgfpathmoveto{\pgfqpoint{3.898223in}{1.521049in}}%
\pgfpathlineto{\pgfqpoint{3.941986in}{1.530771in}}%
\pgfpathlineto{\pgfqpoint{3.975924in}{1.510530in}}%
\pgfpathlineto{\pgfqpoint{3.932135in}{1.500894in}}%
\pgfpathlineto{\pgfqpoint{3.898223in}{1.521049in}}%
\pgfpathclose%
\pgfusepath{stroke,fill}%
\end{pgfscope}%
\begin{pgfscope}%
\pgfpathrectangle{\pgfqpoint{0.050000in}{0.050000in}}{\pgfqpoint{4.900000in}{4.900000in}}%
\pgfusepath{clip}%
\pgfsetbuttcap%
\pgfsetroundjoin%
\definecolor{currentfill}{rgb}{0.363183,0.363183,0.363183}%
\pgfsetfillcolor{currentfill}%
\pgfsetfillopacity{0.900000}%
\pgfsetlinewidth{0.501875pt}%
\definecolor{currentstroke}{rgb}{0.200000,0.200000,0.200000}%
\pgfsetstrokecolor{currentstroke}%
\pgfsetdash{}{0pt}%
\pgfpathmoveto{\pgfqpoint{1.160813in}{3.256861in}}%
\pgfpathlineto{\pgfqpoint{1.209077in}{3.210937in}}%
\pgfpathlineto{\pgfqpoint{1.238761in}{3.191970in}}%
\pgfpathlineto{\pgfqpoint{1.190485in}{3.237924in}}%
\pgfpathlineto{\pgfqpoint{1.160813in}{3.256861in}}%
\pgfpathclose%
\pgfusepath{stroke,fill}%
\end{pgfscope}%
\begin{pgfscope}%
\pgfpathrectangle{\pgfqpoint{0.050000in}{0.050000in}}{\pgfqpoint{4.900000in}{4.900000in}}%
\pgfusepath{clip}%
\pgfsetbuttcap%
\pgfsetroundjoin%
\definecolor{currentfill}{rgb}{1.000000,1.000000,1.000000}%
\pgfsetfillcolor{currentfill}%
\pgfsetfillopacity{0.900000}%
\pgfsetlinewidth{0.501875pt}%
\definecolor{currentstroke}{rgb}{0.200000,0.200000,0.200000}%
\pgfsetstrokecolor{currentstroke}%
\pgfsetdash{}{0pt}%
\pgfpathmoveto{\pgfqpoint{4.175170in}{1.503454in}}%
\pgfpathlineto{\pgfqpoint{4.218684in}{1.515652in}}%
\pgfpathlineto{\pgfqpoint{4.253045in}{1.494941in}}%
\pgfpathlineto{\pgfqpoint{4.209504in}{1.482806in}}%
\pgfpathlineto{\pgfqpoint{4.175170in}{1.503454in}}%
\pgfpathclose%
\pgfusepath{stroke,fill}%
\end{pgfscope}%
\begin{pgfscope}%
\pgfpathrectangle{\pgfqpoint{0.050000in}{0.050000in}}{\pgfqpoint{4.900000in}{4.900000in}}%
\pgfusepath{clip}%
\pgfsetbuttcap%
\pgfsetroundjoin%
\definecolor{currentfill}{rgb}{0.932334,0.932334,0.932334}%
\pgfsetfillcolor{currentfill}%
\pgfsetfillopacity{0.900000}%
\pgfsetlinewidth{0.501875pt}%
\definecolor{currentstroke}{rgb}{0.200000,0.200000,0.200000}%
\pgfsetstrokecolor{currentstroke}%
\pgfsetdash{}{0pt}%
\pgfpathmoveto{\pgfqpoint{2.927169in}{1.857116in}}%
\pgfpathlineto{\pgfqpoint{2.971806in}{1.827023in}}%
\pgfpathlineto{\pgfqpoint{3.004222in}{1.807621in}}%
\pgfpathlineto{\pgfqpoint{2.959574in}{1.837454in}}%
\pgfpathlineto{\pgfqpoint{2.927169in}{1.857116in}}%
\pgfpathclose%
\pgfusepath{stroke,fill}%
\end{pgfscope}%
\begin{pgfscope}%
\pgfpathrectangle{\pgfqpoint{0.050000in}{0.050000in}}{\pgfqpoint{4.900000in}{4.900000in}}%
\pgfusepath{clip}%
\pgfsetbuttcap%
\pgfsetroundjoin%
\definecolor{currentfill}{rgb}{0.619423,0.619423,0.619423}%
\pgfsetfillcolor{currentfill}%
\pgfsetfillopacity{0.900000}%
\pgfsetlinewidth{0.501875pt}%
\definecolor{currentstroke}{rgb}{0.200000,0.200000,0.200000}%
\pgfsetstrokecolor{currentstroke}%
\pgfsetdash{}{0pt}%
\pgfpathmoveto{\pgfqpoint{1.873677in}{2.680690in}}%
\pgfpathlineto{\pgfqpoint{1.920444in}{2.635976in}}%
\pgfpathlineto{\pgfqpoint{1.951224in}{2.616125in}}%
\pgfpathlineto{\pgfqpoint{1.904451in}{2.660866in}}%
\pgfpathlineto{\pgfqpoint{1.873677in}{2.680690in}}%
\pgfpathclose%
\pgfusepath{stroke,fill}%
\end{pgfscope}%
\begin{pgfscope}%
\pgfpathrectangle{\pgfqpoint{0.050000in}{0.050000in}}{\pgfqpoint{4.900000in}{4.900000in}}%
\pgfusepath{clip}%
\pgfsetbuttcap%
\pgfsetroundjoin%
\definecolor{currentfill}{rgb}{0.990773,0.990773,0.990773}%
\pgfsetfillcolor{currentfill}%
\pgfsetfillopacity{0.900000}%
\pgfsetlinewidth{0.501875pt}%
\definecolor{currentstroke}{rgb}{0.200000,0.200000,0.200000}%
\pgfsetstrokecolor{currentstroke}%
\pgfsetdash{}{0pt}%
\pgfpathmoveto{\pgfqpoint{3.621662in}{1.557040in}}%
\pgfpathlineto{\pgfqpoint{3.665629in}{1.560021in}}%
\pgfpathlineto{\pgfqpoint{3.699150in}{1.540343in}}%
\pgfpathlineto{\pgfqpoint{3.655159in}{1.537447in}}%
\pgfpathlineto{\pgfqpoint{3.621662in}{1.557040in}}%
\pgfpathclose%
\pgfusepath{stroke,fill}%
\end{pgfscope}%
\begin{pgfscope}%
\pgfpathrectangle{\pgfqpoint{0.050000in}{0.050000in}}{\pgfqpoint{4.900000in}{4.900000in}}%
\pgfusepath{clip}%
\pgfsetbuttcap%
\pgfsetroundjoin%
\definecolor{currentfill}{rgb}{0.858762,0.858762,0.858762}%
\pgfsetfillcolor{currentfill}%
\pgfsetfillopacity{0.900000}%
\pgfsetlinewidth{0.501875pt}%
\definecolor{currentstroke}{rgb}{0.200000,0.200000,0.200000}%
\pgfsetstrokecolor{currentstroke}%
\pgfsetdash{}{0pt}%
\pgfpathmoveto{\pgfqpoint{2.586637in}{2.107482in}}%
\pgfpathlineto{\pgfqpoint{2.631911in}{2.066351in}}%
\pgfpathlineto{\pgfqpoint{2.663793in}{2.046188in}}%
\pgfpathlineto{\pgfqpoint{2.618514in}{2.087148in}}%
\pgfpathlineto{\pgfqpoint{2.586637in}{2.107482in}}%
\pgfpathclose%
\pgfusepath{stroke,fill}%
\end{pgfscope}%
\begin{pgfscope}%
\pgfpathrectangle{\pgfqpoint{0.050000in}{0.050000in}}{\pgfqpoint{4.900000in}{4.900000in}}%
\pgfusepath{clip}%
\pgfsetbuttcap%
\pgfsetroundjoin%
\definecolor{currentfill}{rgb}{0.521492,0.521492,0.521492}%
\pgfsetfillcolor{currentfill}%
\pgfsetfillopacity{0.900000}%
\pgfsetlinewidth{0.501875pt}%
\definecolor{currentstroke}{rgb}{0.200000,0.200000,0.200000}%
\pgfsetstrokecolor{currentstroke}%
\pgfsetdash{}{0pt}%
\pgfpathmoveto{\pgfqpoint{1.610132in}{2.894328in}}%
\pgfpathlineto{\pgfqpoint{1.657461in}{2.849157in}}%
\pgfpathlineto{\pgfqpoint{1.687842in}{2.829629in}}%
\pgfpathlineto{\pgfqpoint{1.640504in}{2.874829in}}%
\pgfpathlineto{\pgfqpoint{1.610132in}{2.894328in}}%
\pgfpathclose%
\pgfusepath{stroke,fill}%
\end{pgfscope}%
\begin{pgfscope}%
\pgfpathrectangle{\pgfqpoint{0.050000in}{0.050000in}}{\pgfqpoint{4.900000in}{4.900000in}}%
\pgfusepath{clip}%
\pgfsetbuttcap%
\pgfsetroundjoin%
\definecolor{currentfill}{rgb}{0.784667,0.784667,0.784667}%
\pgfsetfillcolor{currentfill}%
\pgfsetfillopacity{0.900000}%
\pgfsetlinewidth{0.501875pt}%
\definecolor{currentstroke}{rgb}{0.200000,0.200000,0.200000}%
\pgfsetstrokecolor{currentstroke}%
\pgfsetdash{}{0pt}%
\pgfpathmoveto{\pgfqpoint{2.323192in}{2.318167in}}%
\pgfpathlineto{\pgfqpoint{2.369023in}{2.274425in}}%
\pgfpathlineto{\pgfqpoint{2.400501in}{2.254101in}}%
\pgfpathlineto{\pgfqpoint{2.354666in}{2.297823in}}%
\pgfpathlineto{\pgfqpoint{2.323192in}{2.318167in}}%
\pgfpathclose%
\pgfusepath{stroke,fill}%
\end{pgfscope}%
\begin{pgfscope}%
\pgfpathrectangle{\pgfqpoint{0.050000in}{0.050000in}}{\pgfqpoint{4.900000in}{4.900000in}}%
\pgfusepath{clip}%
\pgfsetbuttcap%
\pgfsetroundjoin%
\definecolor{currentfill}{rgb}{0.998155,0.998155,0.998155}%
\pgfsetfillcolor{currentfill}%
\pgfsetfillopacity{0.900000}%
\pgfsetlinewidth{0.501875pt}%
\definecolor{currentstroke}{rgb}{0.200000,0.200000,0.200000}%
\pgfsetstrokecolor{currentstroke}%
\pgfsetdash{}{0pt}%
\pgfpathmoveto{\pgfqpoint{3.975924in}{1.510530in}}%
\pgfpathlineto{\pgfqpoint{4.019637in}{1.521178in}}%
\pgfpathlineto{\pgfqpoint{4.053706in}{1.500796in}}%
\pgfpathlineto{\pgfqpoint{4.009968in}{1.490226in}}%
\pgfpathlineto{\pgfqpoint{3.975924in}{1.510530in}}%
\pgfpathclose%
\pgfusepath{stroke,fill}%
\end{pgfscope}%
\begin{pgfscope}%
\pgfpathrectangle{\pgfqpoint{0.050000in}{0.050000in}}{\pgfqpoint{4.900000in}{4.900000in}}%
\pgfusepath{clip}%
\pgfsetbuttcap%
\pgfsetroundjoin%
\definecolor{currentfill}{rgb}{0.424083,0.424083,0.424083}%
\pgfsetfillcolor{currentfill}%
\pgfsetfillopacity{0.900000}%
\pgfsetlinewidth{0.501875pt}%
\definecolor{currentstroke}{rgb}{0.200000,0.200000,0.200000}%
\pgfsetstrokecolor{currentstroke}%
\pgfsetdash{}{0pt}%
\pgfpathmoveto{\pgfqpoint{1.346498in}{3.108039in}}%
\pgfpathlineto{\pgfqpoint{1.394390in}{3.062411in}}%
\pgfpathlineto{\pgfqpoint{1.424372in}{3.043208in}}%
\pgfpathlineto{\pgfqpoint{1.376469in}{3.088866in}}%
\pgfpathlineto{\pgfqpoint{1.346498in}{3.108039in}}%
\pgfpathclose%
\pgfusepath{stroke,fill}%
\end{pgfscope}%
\begin{pgfscope}%
\pgfpathrectangle{\pgfqpoint{0.050000in}{0.050000in}}{\pgfqpoint{4.900000in}{4.900000in}}%
\pgfusepath{clip}%
\pgfsetbuttcap%
\pgfsetroundjoin%
\definecolor{currentfill}{rgb}{0.915356,0.915356,0.915356}%
\pgfsetfillcolor{currentfill}%
\pgfsetfillopacity{0.900000}%
\pgfsetlinewidth{0.501875pt}%
\definecolor{currentstroke}{rgb}{0.200000,0.200000,0.200000}%
\pgfsetstrokecolor{currentstroke}%
\pgfsetdash{}{0pt}%
\pgfpathmoveto{\pgfqpoint{2.850089in}{1.910357in}}%
\pgfpathlineto{\pgfqpoint{2.894867in}{1.876718in}}%
\pgfpathlineto{\pgfqpoint{2.927169in}{1.857116in}}%
\pgfpathlineto{\pgfqpoint{2.882382in}{1.890487in}}%
\pgfpathlineto{\pgfqpoint{2.850089in}{1.910357in}}%
\pgfpathclose%
\pgfusepath{stroke,fill}%
\end{pgfscope}%
\begin{pgfscope}%
\pgfpathrectangle{\pgfqpoint{0.050000in}{0.050000in}}{\pgfqpoint{4.900000in}{4.900000in}}%
\pgfusepath{clip}%
\pgfsetbuttcap%
\pgfsetroundjoin%
\definecolor{currentfill}{rgb}{0.691396,0.691396,0.691396}%
\pgfsetfillcolor{currentfill}%
\pgfsetfillopacity{0.900000}%
\pgfsetlinewidth{0.501875pt}%
\definecolor{currentstroke}{rgb}{0.200000,0.200000,0.200000}%
\pgfsetstrokecolor{currentstroke}%
\pgfsetdash{}{0pt}%
\pgfpathmoveto{\pgfqpoint{2.059661in}{2.531634in}}%
\pgfpathlineto{\pgfqpoint{2.106055in}{2.487219in}}%
\pgfpathlineto{\pgfqpoint{2.137134in}{2.467133in}}%
\pgfpathlineto{\pgfqpoint{2.090734in}{2.511572in}}%
\pgfpathlineto{\pgfqpoint{2.059661in}{2.531634in}}%
\pgfpathclose%
\pgfusepath{stroke,fill}%
\end{pgfscope}%
\begin{pgfscope}%
\pgfpathrectangle{\pgfqpoint{0.050000in}{0.050000in}}{\pgfqpoint{4.900000in}{4.900000in}}%
\pgfusepath{clip}%
\pgfsetbuttcap%
\pgfsetroundjoin%
\definecolor{currentfill}{rgb}{0.992618,0.992618,0.992618}%
\pgfsetfillcolor{currentfill}%
\pgfsetfillopacity{0.900000}%
\pgfsetlinewidth{0.501875pt}%
\definecolor{currentstroke}{rgb}{0.200000,0.200000,0.200000}%
\pgfsetstrokecolor{currentstroke}%
\pgfsetdash{}{0pt}%
\pgfpathmoveto{\pgfqpoint{3.699150in}{1.540343in}}%
\pgfpathlineto{\pgfqpoint{3.743074in}{1.545490in}}%
\pgfpathlineto{\pgfqpoint{3.776725in}{1.525650in}}%
\pgfpathlineto{\pgfqpoint{3.732776in}{1.520594in}}%
\pgfpathlineto{\pgfqpoint{3.699150in}{1.540343in}}%
\pgfpathclose%
\pgfusepath{stroke,fill}%
\end{pgfscope}%
\begin{pgfscope}%
\pgfpathrectangle{\pgfqpoint{0.050000in}{0.050000in}}{\pgfqpoint{4.900000in}{4.900000in}}%
\pgfusepath{clip}%
\pgfsetbuttcap%
\pgfsetroundjoin%
\definecolor{currentfill}{rgb}{1.000000,1.000000,1.000000}%
\pgfsetfillcolor{currentfill}%
\pgfsetfillopacity{0.900000}%
\pgfsetlinewidth{0.501875pt}%
\definecolor{currentstroke}{rgb}{0.200000,0.200000,0.200000}%
\pgfsetstrokecolor{currentstroke}%
\pgfsetdash{}{0pt}%
\pgfpathmoveto{\pgfqpoint{4.253045in}{1.494941in}}%
\pgfpathlineto{\pgfqpoint{4.296493in}{1.507272in}}%
\pgfpathlineto{\pgfqpoint{4.330986in}{1.486438in}}%
\pgfpathlineto{\pgfqpoint{4.287512in}{1.474167in}}%
\pgfpathlineto{\pgfqpoint{4.253045in}{1.494941in}}%
\pgfpathclose%
\pgfusepath{stroke,fill}%
\end{pgfscope}%
\begin{pgfscope}%
\pgfpathrectangle{\pgfqpoint{0.050000in}{0.050000in}}{\pgfqpoint{4.900000in}{4.900000in}}%
\pgfusepath{clip}%
\pgfsetbuttcap%
\pgfsetroundjoin%
\definecolor{currentfill}{rgb}{0.330704,0.330704,0.330704}%
\pgfsetfillcolor{currentfill}%
\pgfsetfillopacity{0.900000}%
\pgfsetlinewidth{0.501875pt}%
\definecolor{currentstroke}{rgb}{0.200000,0.200000,0.200000}%
\pgfsetstrokecolor{currentstroke}%
\pgfsetdash{}{0pt}%
\pgfpathmoveto{\pgfqpoint{1.082776in}{3.321825in}}%
\pgfpathlineto{\pgfqpoint{1.131232in}{3.275739in}}%
\pgfpathlineto{\pgfqpoint{1.160813in}{3.256861in}}%
\pgfpathlineto{\pgfqpoint{1.112345in}{3.302978in}}%
\pgfpathlineto{\pgfqpoint{1.082776in}{3.321825in}}%
\pgfpathclose%
\pgfusepath{stroke,fill}%
\end{pgfscope}%
\begin{pgfscope}%
\pgfpathrectangle{\pgfqpoint{0.050000in}{0.050000in}}{\pgfqpoint{4.900000in}{4.900000in}}%
\pgfusepath{clip}%
\pgfsetbuttcap%
\pgfsetroundjoin%
\definecolor{currentfill}{rgb}{0.586082,0.586082,0.586082}%
\pgfsetfillcolor{currentfill}%
\pgfsetfillopacity{0.900000}%
\pgfsetlinewidth{0.501875pt}%
\definecolor{currentstroke}{rgb}{0.200000,0.200000,0.200000}%
\pgfsetstrokecolor{currentstroke}%
\pgfsetdash{}{0pt}%
\pgfpathmoveto{\pgfqpoint{1.796041in}{2.745329in}}%
\pgfpathlineto{\pgfqpoint{1.842998in}{2.700454in}}%
\pgfpathlineto{\pgfqpoint{1.873677in}{2.680690in}}%
\pgfpathlineto{\pgfqpoint{1.826713in}{2.725593in}}%
\pgfpathlineto{\pgfqpoint{1.796041in}{2.745329in}}%
\pgfpathclose%
\pgfusepath{stroke,fill}%
\end{pgfscope}%
\begin{pgfscope}%
\pgfpathrectangle{\pgfqpoint{0.050000in}{0.050000in}}{\pgfqpoint{4.900000in}{4.900000in}}%
\pgfusepath{clip}%
\pgfsetbuttcap%
\pgfsetroundjoin%
\definecolor{currentfill}{rgb}{0.836340,0.836340,0.836340}%
\pgfsetfillcolor{currentfill}%
\pgfsetfillopacity{0.900000}%
\pgfsetlinewidth{0.501875pt}%
\definecolor{currentstroke}{rgb}{0.200000,0.200000,0.200000}%
\pgfsetstrokecolor{currentstroke}%
\pgfsetdash{}{0pt}%
\pgfpathmoveto{\pgfqpoint{2.509399in}{2.170159in}}%
\pgfpathlineto{\pgfqpoint{2.554858in}{2.127774in}}%
\pgfpathlineto{\pgfqpoint{2.586637in}{2.107482in}}%
\pgfpathlineto{\pgfqpoint{2.541174in}{2.149752in}}%
\pgfpathlineto{\pgfqpoint{2.509399in}{2.170159in}}%
\pgfpathclose%
\pgfusepath{stroke,fill}%
\end{pgfscope}%
\begin{pgfscope}%
\pgfpathrectangle{\pgfqpoint{0.050000in}{0.050000in}}{\pgfqpoint{4.900000in}{4.900000in}}%
\pgfusepath{clip}%
\pgfsetbuttcap%
\pgfsetroundjoin%
\definecolor{currentfill}{rgb}{0.998155,0.998155,0.998155}%
\pgfsetfillcolor{currentfill}%
\pgfsetfillopacity{0.900000}%
\pgfsetlinewidth{0.501875pt}%
\definecolor{currentstroke}{rgb}{0.200000,0.200000,0.200000}%
\pgfsetstrokecolor{currentstroke}%
\pgfsetdash{}{0pt}%
\pgfpathmoveto{\pgfqpoint{4.053706in}{1.500796in}}%
\pgfpathlineto{\pgfqpoint{4.097366in}{1.512139in}}%
\pgfpathlineto{\pgfqpoint{4.131568in}{1.491621in}}%
\pgfpathlineto{\pgfqpoint{4.087883in}{1.480350in}}%
\pgfpathlineto{\pgfqpoint{4.053706in}{1.500796in}}%
\pgfpathclose%
\pgfusepath{stroke,fill}%
\end{pgfscope}%
\begin{pgfscope}%
\pgfpathrectangle{\pgfqpoint{0.050000in}{0.050000in}}{\pgfqpoint{4.900000in}{4.900000in}}%
\pgfusepath{clip}%
\pgfsetbuttcap%
\pgfsetroundjoin%
\definecolor{currentfill}{rgb}{0.994464,0.994464,0.994464}%
\pgfsetfillcolor{currentfill}%
\pgfsetfillopacity{0.900000}%
\pgfsetlinewidth{0.501875pt}%
\definecolor{currentstroke}{rgb}{0.200000,0.200000,0.200000}%
\pgfsetstrokecolor{currentstroke}%
\pgfsetdash{}{0pt}%
\pgfpathmoveto{\pgfqpoint{3.776725in}{1.525650in}}%
\pgfpathlineto{\pgfqpoint{3.820606in}{1.532607in}}%
\pgfpathlineto{\pgfqpoint{3.854387in}{1.512607in}}%
\pgfpathlineto{\pgfqpoint{3.810481in}{1.505741in}}%
\pgfpathlineto{\pgfqpoint{3.776725in}{1.525650in}}%
\pgfpathclose%
\pgfusepath{stroke,fill}%
\end{pgfscope}%
\begin{pgfscope}%
\pgfpathrectangle{\pgfqpoint{0.050000in}{0.050000in}}{\pgfqpoint{4.900000in}{4.900000in}}%
\pgfusepath{clip}%
\pgfsetbuttcap%
\pgfsetroundjoin%
\definecolor{currentfill}{rgb}{0.487043,0.487043,0.487043}%
\pgfsetfillcolor{currentfill}%
\pgfsetfillopacity{0.900000}%
\pgfsetlinewidth{0.501875pt}%
\definecolor{currentstroke}{rgb}{0.200000,0.200000,0.200000}%
\pgfsetstrokecolor{currentstroke}%
\pgfsetdash{}{0pt}%
\pgfpathmoveto{\pgfqpoint{1.532333in}{2.959100in}}%
\pgfpathlineto{\pgfqpoint{1.579853in}{2.913767in}}%
\pgfpathlineto{\pgfqpoint{1.610132in}{2.894328in}}%
\pgfpathlineto{\pgfqpoint{1.562603in}{2.939690in}}%
\pgfpathlineto{\pgfqpoint{1.532333in}{2.959100in}}%
\pgfpathclose%
\pgfusepath{stroke,fill}%
\end{pgfscope}%
\begin{pgfscope}%
\pgfpathrectangle{\pgfqpoint{0.050000in}{0.050000in}}{\pgfqpoint{4.900000in}{4.900000in}}%
\pgfusepath{clip}%
\pgfsetbuttcap%
\pgfsetroundjoin%
\definecolor{currentfill}{rgb}{0.898378,0.898378,0.898378}%
\pgfsetfillcolor{currentfill}%
\pgfsetfillopacity{0.900000}%
\pgfsetlinewidth{0.501875pt}%
\definecolor{currentstroke}{rgb}{0.200000,0.200000,0.200000}%
\pgfsetstrokecolor{currentstroke}%
\pgfsetdash{}{0pt}%
\pgfpathmoveto{\pgfqpoint{2.772962in}{1.966845in}}%
\pgfpathlineto{\pgfqpoint{2.817895in}{1.930175in}}%
\pgfpathlineto{\pgfqpoint{2.850089in}{1.910357in}}%
\pgfpathlineto{\pgfqpoint{2.805148in}{1.946772in}}%
\pgfpathlineto{\pgfqpoint{2.772962in}{1.966845in}}%
\pgfpathclose%
\pgfusepath{stroke,fill}%
\end{pgfscope}%
\begin{pgfscope}%
\pgfpathrectangle{\pgfqpoint{0.050000in}{0.050000in}}{\pgfqpoint{4.900000in}{4.900000in}}%
\pgfusepath{clip}%
\pgfsetbuttcap%
\pgfsetroundjoin%
\definecolor{currentfill}{rgb}{0.760554,0.760554,0.760554}%
\pgfsetfillcolor{currentfill}%
\pgfsetfillopacity{0.900000}%
\pgfsetlinewidth{0.501875pt}%
\definecolor{currentstroke}{rgb}{0.200000,0.200000,0.200000}%
\pgfsetstrokecolor{currentstroke}%
\pgfsetdash{}{0pt}%
\pgfpathmoveto{\pgfqpoint{2.245794in}{2.382502in}}%
\pgfpathlineto{\pgfqpoint{2.291814in}{2.338456in}}%
\pgfpathlineto{\pgfqpoint{2.323192in}{2.318167in}}%
\pgfpathlineto{\pgfqpoint{2.277167in}{2.362218in}}%
\pgfpathlineto{\pgfqpoint{2.245794in}{2.382502in}}%
\pgfpathclose%
\pgfusepath{stroke,fill}%
\end{pgfscope}%
\begin{pgfscope}%
\pgfpathrectangle{\pgfqpoint{0.050000in}{0.050000in}}{\pgfqpoint{4.900000in}{4.900000in}}%
\pgfusepath{clip}%
\pgfsetbuttcap%
\pgfsetroundjoin%
\definecolor{currentfill}{rgb}{0.391603,0.391603,0.391603}%
\pgfsetfillcolor{currentfill}%
\pgfsetfillopacity{0.900000}%
\pgfsetlinewidth{0.501875pt}%
\definecolor{currentstroke}{rgb}{0.200000,0.200000,0.200000}%
\pgfsetstrokecolor{currentstroke}%
\pgfsetdash{}{0pt}%
\pgfpathmoveto{\pgfqpoint{1.268536in}{3.172945in}}%
\pgfpathlineto{\pgfqpoint{1.316620in}{3.127154in}}%
\pgfpathlineto{\pgfqpoint{1.346498in}{3.108039in}}%
\pgfpathlineto{\pgfqpoint{1.298404in}{3.153861in}}%
\pgfpathlineto{\pgfqpoint{1.268536in}{3.172945in}}%
\pgfpathclose%
\pgfusepath{stroke,fill}%
\end{pgfscope}%
\begin{pgfscope}%
\pgfpathrectangle{\pgfqpoint{0.050000in}{0.050000in}}{\pgfqpoint{4.900000in}{4.900000in}}%
\pgfusepath{clip}%
\pgfsetbuttcap%
\pgfsetroundjoin%
\definecolor{currentfill}{rgb}{0.968627,0.968627,0.968627}%
\pgfsetfillcolor{currentfill}%
\pgfsetfillopacity{0.900000}%
\pgfsetlinewidth{0.501875pt}%
\definecolor{currentstroke}{rgb}{0.200000,0.200000,0.200000}%
\pgfsetstrokecolor{currentstroke}%
\pgfsetdash{}{0pt}%
\pgfpathmoveto{\pgfqpoint{3.268264in}{1.665188in}}%
\pgfpathlineto{\pgfqpoint{3.312514in}{1.651753in}}%
\pgfpathlineto{\pgfqpoint{3.345507in}{1.632582in}}%
\pgfpathlineto{\pgfqpoint{3.301239in}{1.645939in}}%
\pgfpathlineto{\pgfqpoint{3.268264in}{1.665188in}}%
\pgfpathclose%
\pgfusepath{stroke,fill}%
\end{pgfscope}%
\begin{pgfscope}%
\pgfpathrectangle{\pgfqpoint{0.050000in}{0.050000in}}{\pgfqpoint{4.900000in}{4.900000in}}%
\pgfusepath{clip}%
\pgfsetbuttcap%
\pgfsetroundjoin%
\definecolor{currentfill}{rgb}{0.657809,0.657809,0.657809}%
\pgfsetfillcolor{currentfill}%
\pgfsetfillopacity{0.900000}%
\pgfsetlinewidth{0.501875pt}%
\definecolor{currentstroke}{rgb}{0.200000,0.200000,0.200000}%
\pgfsetstrokecolor{currentstroke}%
\pgfsetdash{}{0pt}%
\pgfpathmoveto{\pgfqpoint{1.982100in}{2.596212in}}%
\pgfpathlineto{\pgfqpoint{2.028684in}{2.551634in}}%
\pgfpathlineto{\pgfqpoint{2.059661in}{2.531634in}}%
\pgfpathlineto{\pgfqpoint{2.013071in}{2.576238in}}%
\pgfpathlineto{\pgfqpoint{1.982100in}{2.596212in}}%
\pgfpathclose%
\pgfusepath{stroke,fill}%
\end{pgfscope}%
\begin{pgfscope}%
\pgfpathrectangle{\pgfqpoint{0.050000in}{0.050000in}}{\pgfqpoint{4.900000in}{4.900000in}}%
\pgfusepath{clip}%
\pgfsetbuttcap%
\pgfsetroundjoin%
\definecolor{currentfill}{rgb}{0.974164,0.974164,0.974164}%
\pgfsetfillcolor{currentfill}%
\pgfsetfillopacity{0.900000}%
\pgfsetlinewidth{0.501875pt}%
\definecolor{currentstroke}{rgb}{0.200000,0.200000,0.200000}%
\pgfsetstrokecolor{currentstroke}%
\pgfsetdash{}{0pt}%
\pgfpathmoveto{\pgfqpoint{3.345507in}{1.632582in}}%
\pgfpathlineto{\pgfqpoint{3.389692in}{1.623083in}}%
\pgfpathlineto{\pgfqpoint{3.422809in}{1.603856in}}%
\pgfpathlineto{\pgfqpoint{3.378604in}{1.613325in}}%
\pgfpathlineto{\pgfqpoint{3.345507in}{1.632582in}}%
\pgfpathclose%
\pgfusepath{stroke,fill}%
\end{pgfscope}%
\begin{pgfscope}%
\pgfpathrectangle{\pgfqpoint{0.050000in}{0.050000in}}{\pgfqpoint{4.900000in}{4.900000in}}%
\pgfusepath{clip}%
\pgfsetbuttcap%
\pgfsetroundjoin%
\definecolor{currentfill}{rgb}{0.961246,0.961246,0.961246}%
\pgfsetfillcolor{currentfill}%
\pgfsetfillopacity{0.900000}%
\pgfsetlinewidth{0.501875pt}%
\definecolor{currentstroke}{rgb}{0.200000,0.200000,0.200000}%
\pgfsetstrokecolor{currentstroke}%
\pgfsetdash{}{0pt}%
\pgfpathmoveto{\pgfqpoint{3.191066in}{1.701921in}}%
\pgfpathlineto{\pgfqpoint{3.235393in}{1.684352in}}%
\pgfpathlineto{\pgfqpoint{3.268264in}{1.665188in}}%
\pgfpathlineto{\pgfqpoint{3.223920in}{1.682627in}}%
\pgfpathlineto{\pgfqpoint{3.191066in}{1.701921in}}%
\pgfpathclose%
\pgfusepath{stroke,fill}%
\end{pgfscope}%
\begin{pgfscope}%
\pgfpathrectangle{\pgfqpoint{0.050000in}{0.050000in}}{\pgfqpoint{4.900000in}{4.900000in}}%
\pgfusepath{clip}%
\pgfsetbuttcap%
\pgfsetroundjoin%
\definecolor{currentfill}{rgb}{0.979700,0.979700,0.979700}%
\pgfsetfillcolor{currentfill}%
\pgfsetfillopacity{0.900000}%
\pgfsetlinewidth{0.501875pt}%
\definecolor{currentstroke}{rgb}{0.200000,0.200000,0.200000}%
\pgfsetstrokecolor{currentstroke}%
\pgfsetdash{}{0pt}%
\pgfpathmoveto{\pgfqpoint{3.422809in}{1.603856in}}%
\pgfpathlineto{\pgfqpoint{3.466937in}{1.598003in}}%
\pgfpathlineto{\pgfqpoint{3.500180in}{1.578680in}}%
\pgfpathlineto{\pgfqpoint{3.456031in}{1.584545in}}%
\pgfpathlineto{\pgfqpoint{3.422809in}{1.603856in}}%
\pgfpathclose%
\pgfusepath{stroke,fill}%
\end{pgfscope}%
\begin{pgfscope}%
\pgfpathrectangle{\pgfqpoint{0.050000in}{0.050000in}}{\pgfqpoint{4.900000in}{4.900000in}}%
\pgfusepath{clip}%
\pgfsetbuttcap%
\pgfsetroundjoin%
\definecolor{currentfill}{rgb}{0.996309,0.996309,0.996309}%
\pgfsetfillcolor{currentfill}%
\pgfsetfillopacity{0.900000}%
\pgfsetlinewidth{0.501875pt}%
\definecolor{currentstroke}{rgb}{0.200000,0.200000,0.200000}%
\pgfsetstrokecolor{currentstroke}%
\pgfsetdash{}{0pt}%
\pgfpathmoveto{\pgfqpoint{3.854387in}{1.512607in}}%
\pgfpathlineto{\pgfqpoint{3.898223in}{1.521049in}}%
\pgfpathlineto{\pgfqpoint{3.932135in}{1.500894in}}%
\pgfpathlineto{\pgfqpoint{3.888273in}{1.492540in}}%
\pgfpathlineto{\pgfqpoint{3.854387in}{1.512607in}}%
\pgfpathclose%
\pgfusepath{stroke,fill}%
\end{pgfscope}%
\begin{pgfscope}%
\pgfpathrectangle{\pgfqpoint{0.050000in}{0.050000in}}{\pgfqpoint{4.900000in}{4.900000in}}%
\pgfusepath{clip}%
\pgfsetbuttcap%
\pgfsetroundjoin%
\definecolor{currentfill}{rgb}{0.953864,0.953864,0.953864}%
\pgfsetfillcolor{currentfill}%
\pgfsetfillopacity{0.900000}%
\pgfsetlinewidth{0.501875pt}%
\definecolor{currentstroke}{rgb}{0.200000,0.200000,0.200000}%
\pgfsetstrokecolor{currentstroke}%
\pgfsetdash{}{0pt}%
\pgfpathmoveto{\pgfqpoint{3.113897in}{1.742913in}}%
\pgfpathlineto{\pgfqpoint{3.158315in}{1.721134in}}%
\pgfpathlineto{\pgfqpoint{3.191066in}{1.701921in}}%
\pgfpathlineto{\pgfqpoint{3.146632in}{1.723520in}}%
\pgfpathlineto{\pgfqpoint{3.113897in}{1.742913in}}%
\pgfpathclose%
\pgfusepath{stroke,fill}%
\end{pgfscope}%
\begin{pgfscope}%
\pgfpathrectangle{\pgfqpoint{0.050000in}{0.050000in}}{\pgfqpoint{4.900000in}{4.900000in}}%
\pgfusepath{clip}%
\pgfsetbuttcap%
\pgfsetroundjoin%
\definecolor{currentfill}{rgb}{1.000000,1.000000,1.000000}%
\pgfsetfillcolor{currentfill}%
\pgfsetfillopacity{0.900000}%
\pgfsetlinewidth{0.501875pt}%
\definecolor{currentstroke}{rgb}{0.200000,0.200000,0.200000}%
\pgfsetstrokecolor{currentstroke}%
\pgfsetdash{}{0pt}%
\pgfpathmoveto{\pgfqpoint{4.131568in}{1.491621in}}%
\pgfpathlineto{\pgfqpoint{4.175170in}{1.503454in}}%
\pgfpathlineto{\pgfqpoint{4.209504in}{1.482806in}}%
\pgfpathlineto{\pgfqpoint{4.165877in}{1.471040in}}%
\pgfpathlineto{\pgfqpoint{4.131568in}{1.491621in}}%
\pgfpathclose%
\pgfusepath{stroke,fill}%
\end{pgfscope}%
\begin{pgfscope}%
\pgfpathrectangle{\pgfqpoint{0.050000in}{0.050000in}}{\pgfqpoint{4.900000in}{4.900000in}}%
\pgfusepath{clip}%
\pgfsetbuttcap%
\pgfsetroundjoin%
\definecolor{currentfill}{rgb}{0.551634,0.551634,0.551634}%
\pgfsetfillcolor{currentfill}%
\pgfsetfillopacity{0.900000}%
\pgfsetlinewidth{0.501875pt}%
\definecolor{currentstroke}{rgb}{0.200000,0.200000,0.200000}%
\pgfsetstrokecolor{currentstroke}%
\pgfsetdash{}{0pt}%
\pgfpathmoveto{\pgfqpoint{1.718317in}{2.810042in}}%
\pgfpathlineto{\pgfqpoint{1.765464in}{2.765005in}}%
\pgfpathlineto{\pgfqpoint{1.796041in}{2.745329in}}%
\pgfpathlineto{\pgfqpoint{1.748886in}{2.790394in}}%
\pgfpathlineto{\pgfqpoint{1.718317in}{2.810042in}}%
\pgfpathclose%
\pgfusepath{stroke,fill}%
\end{pgfscope}%
\begin{pgfscope}%
\pgfpathrectangle{\pgfqpoint{0.050000in}{0.050000in}}{\pgfqpoint{4.900000in}{4.900000in}}%
\pgfusepath{clip}%
\pgfsetbuttcap%
\pgfsetroundjoin%
\definecolor{currentfill}{rgb}{0.878570,0.878570,0.878570}%
\pgfsetfillcolor{currentfill}%
\pgfsetfillopacity{0.900000}%
\pgfsetlinewidth{0.501875pt}%
\definecolor{currentstroke}{rgb}{0.200000,0.200000,0.200000}%
\pgfsetstrokecolor{currentstroke}%
\pgfsetdash{}{0pt}%
\pgfpathmoveto{\pgfqpoint{2.695774in}{2.025983in}}%
\pgfpathlineto{\pgfqpoint{2.740876in}{1.986872in}}%
\pgfpathlineto{\pgfqpoint{2.772962in}{1.966845in}}%
\pgfpathlineto{\pgfqpoint{2.727854in}{2.005736in}}%
\pgfpathlineto{\pgfqpoint{2.695774in}{2.025983in}}%
\pgfpathclose%
\pgfusepath{stroke,fill}%
\end{pgfscope}%
\begin{pgfscope}%
\pgfpathrectangle{\pgfqpoint{0.050000in}{0.050000in}}{\pgfqpoint{4.900000in}{4.900000in}}%
\pgfusepath{clip}%
\pgfsetbuttcap%
\pgfsetroundjoin%
\definecolor{currentfill}{rgb}{0.983391,0.983391,0.983391}%
\pgfsetfillcolor{currentfill}%
\pgfsetfillopacity{0.900000}%
\pgfsetlinewidth{0.501875pt}%
\definecolor{currentstroke}{rgb}{0.200000,0.200000,0.200000}%
\pgfsetstrokecolor{currentstroke}%
\pgfsetdash{}{0pt}%
\pgfpathmoveto{\pgfqpoint{3.500180in}{1.578680in}}%
\pgfpathlineto{\pgfqpoint{3.544259in}{1.576125in}}%
\pgfpathlineto{\pgfqpoint{3.577628in}{1.556677in}}%
\pgfpathlineto{\pgfqpoint{3.533527in}{1.559276in}}%
\pgfpathlineto{\pgfqpoint{3.500180in}{1.578680in}}%
\pgfpathclose%
\pgfusepath{stroke,fill}%
\end{pgfscope}%
\begin{pgfscope}%
\pgfpathrectangle{\pgfqpoint{0.050000in}{0.050000in}}{\pgfqpoint{4.900000in}{4.900000in}}%
\pgfusepath{clip}%
\pgfsetbuttcap%
\pgfsetroundjoin%
\definecolor{currentfill}{rgb}{0.812226,0.812226,0.812226}%
\pgfsetfillcolor{currentfill}%
\pgfsetfillopacity{0.900000}%
\pgfsetlinewidth{0.501875pt}%
\definecolor{currentstroke}{rgb}{0.200000,0.200000,0.200000}%
\pgfsetstrokecolor{currentstroke}%
\pgfsetdash{}{0pt}%
\pgfpathmoveto{\pgfqpoint{2.432077in}{2.233727in}}%
\pgfpathlineto{\pgfqpoint{2.477722in}{2.190520in}}%
\pgfpathlineto{\pgfqpoint{2.509399in}{2.170159in}}%
\pgfpathlineto{\pgfqpoint{2.463750in}{2.213301in}}%
\pgfpathlineto{\pgfqpoint{2.432077in}{2.233727in}}%
\pgfpathclose%
\pgfusepath{stroke,fill}%
\end{pgfscope}%
\begin{pgfscope}%
\pgfpathrectangle{\pgfqpoint{0.050000in}{0.050000in}}{\pgfqpoint{4.900000in}{4.900000in}}%
\pgfusepath{clip}%
\pgfsetbuttcap%
\pgfsetroundjoin%
\definecolor{currentfill}{rgb}{0.944637,0.944637,0.944637}%
\pgfsetfillcolor{currentfill}%
\pgfsetfillopacity{0.900000}%
\pgfsetlinewidth{0.501875pt}%
\definecolor{currentstroke}{rgb}{0.200000,0.200000,0.200000}%
\pgfsetstrokecolor{currentstroke}%
\pgfsetdash{}{0pt}%
\pgfpathmoveto{\pgfqpoint{3.036739in}{1.788150in}}%
\pgfpathlineto{\pgfqpoint{3.081264in}{1.762230in}}%
\pgfpathlineto{\pgfqpoint{3.113897in}{1.742913in}}%
\pgfpathlineto{\pgfqpoint{3.069358in}{1.768610in}}%
\pgfpathlineto{\pgfqpoint{3.036739in}{1.788150in}}%
\pgfpathclose%
\pgfusepath{stroke,fill}%
\end{pgfscope}%
\begin{pgfscope}%
\pgfpathrectangle{\pgfqpoint{0.050000in}{0.050000in}}{\pgfqpoint{4.900000in}{4.900000in}}%
\pgfusepath{clip}%
\pgfsetbuttcap%
\pgfsetroundjoin%
\definecolor{currentfill}{rgb}{0.456901,0.456901,0.456901}%
\pgfsetfillcolor{currentfill}%
\pgfsetfillopacity{0.900000}%
\pgfsetlinewidth{0.501875pt}%
\definecolor{currentstroke}{rgb}{0.200000,0.200000,0.200000}%
\pgfsetstrokecolor{currentstroke}%
\pgfsetdash{}{0pt}%
\pgfpathmoveto{\pgfqpoint{1.454445in}{3.023946in}}%
\pgfpathlineto{\pgfqpoint{1.502157in}{2.978451in}}%
\pgfpathlineto{\pgfqpoint{1.532333in}{2.959100in}}%
\pgfpathlineto{\pgfqpoint{1.484612in}{3.004625in}}%
\pgfpathlineto{\pgfqpoint{1.454445in}{3.023946in}}%
\pgfpathclose%
\pgfusepath{stroke,fill}%
\end{pgfscope}%
\begin{pgfscope}%
\pgfpathrectangle{\pgfqpoint{0.050000in}{0.050000in}}{\pgfqpoint{4.900000in}{4.900000in}}%
\pgfusepath{clip}%
\pgfsetbuttcap%
\pgfsetroundjoin%
\definecolor{currentfill}{rgb}{0.729781,0.729781,0.729781}%
\pgfsetfillcolor{currentfill}%
\pgfsetfillopacity{0.900000}%
\pgfsetlinewidth{0.501875pt}%
\definecolor{currentstroke}{rgb}{0.200000,0.200000,0.200000}%
\pgfsetstrokecolor{currentstroke}%
\pgfsetdash{}{0pt}%
\pgfpathmoveto{\pgfqpoint{2.168308in}{2.446987in}}%
\pgfpathlineto{\pgfqpoint{2.214518in}{2.402726in}}%
\pgfpathlineto{\pgfqpoint{2.245794in}{2.382502in}}%
\pgfpathlineto{\pgfqpoint{2.199579in}{2.426780in}}%
\pgfpathlineto{\pgfqpoint{2.168308in}{2.446987in}}%
\pgfpathclose%
\pgfusepath{stroke,fill}%
\end{pgfscope}%
\begin{pgfscope}%
\pgfpathrectangle{\pgfqpoint{0.050000in}{0.050000in}}{\pgfqpoint{4.900000in}{4.900000in}}%
\pgfusepath{clip}%
\pgfsetbuttcap%
\pgfsetroundjoin%
\definecolor{currentfill}{rgb}{0.987082,0.987082,0.987082}%
\pgfsetfillcolor{currentfill}%
\pgfsetfillopacity{0.900000}%
\pgfsetlinewidth{0.501875pt}%
\definecolor{currentstroke}{rgb}{0.200000,0.200000,0.200000}%
\pgfsetstrokecolor{currentstroke}%
\pgfsetdash{}{0pt}%
\pgfpathmoveto{\pgfqpoint{3.577628in}{1.556677in}}%
\pgfpathlineto{\pgfqpoint{3.621662in}{1.557040in}}%
\pgfpathlineto{\pgfqpoint{3.655159in}{1.537447in}}%
\pgfpathlineto{\pgfqpoint{3.611103in}{1.537150in}}%
\pgfpathlineto{\pgfqpoint{3.577628in}{1.556677in}}%
\pgfpathclose%
\pgfusepath{stroke,fill}%
\end{pgfscope}%
\begin{pgfscope}%
\pgfpathrectangle{\pgfqpoint{0.050000in}{0.050000in}}{\pgfqpoint{4.900000in}{4.900000in}}%
\pgfusepath{clip}%
\pgfsetbuttcap%
\pgfsetroundjoin%
\definecolor{currentfill}{rgb}{0.996309,0.996309,0.996309}%
\pgfsetfillcolor{currentfill}%
\pgfsetfillopacity{0.900000}%
\pgfsetlinewidth{0.501875pt}%
\definecolor{currentstroke}{rgb}{0.200000,0.200000,0.200000}%
\pgfsetstrokecolor{currentstroke}%
\pgfsetdash{}{0pt}%
\pgfpathmoveto{\pgfqpoint{3.932135in}{1.500894in}}%
\pgfpathlineto{\pgfqpoint{3.975924in}{1.510530in}}%
\pgfpathlineto{\pgfqpoint{4.009968in}{1.490226in}}%
\pgfpathlineto{\pgfqpoint{3.966153in}{1.480674in}}%
\pgfpathlineto{\pgfqpoint{3.932135in}{1.500894in}}%
\pgfpathclose%
\pgfusepath{stroke,fill}%
\end{pgfscope}%
\begin{pgfscope}%
\pgfpathrectangle{\pgfqpoint{0.050000in}{0.050000in}}{\pgfqpoint{4.900000in}{4.900000in}}%
\pgfusepath{clip}%
\pgfsetbuttcap%
\pgfsetroundjoin%
\definecolor{currentfill}{rgb}{0.363183,0.363183,0.363183}%
\pgfsetfillcolor{currentfill}%
\pgfsetfillopacity{0.900000}%
\pgfsetlinewidth{0.501875pt}%
\definecolor{currentstroke}{rgb}{0.200000,0.200000,0.200000}%
\pgfsetstrokecolor{currentstroke}%
\pgfsetdash{}{0pt}%
\pgfpathmoveto{\pgfqpoint{1.190485in}{3.237924in}}%
\pgfpathlineto{\pgfqpoint{1.238761in}{3.191970in}}%
\pgfpathlineto{\pgfqpoint{1.268536in}{3.172945in}}%
\pgfpathlineto{\pgfqpoint{1.220249in}{3.218930in}}%
\pgfpathlineto{\pgfqpoint{1.190485in}{3.237924in}}%
\pgfpathclose%
\pgfusepath{stroke,fill}%
\end{pgfscope}%
\begin{pgfscope}%
\pgfpathrectangle{\pgfqpoint{0.050000in}{0.050000in}}{\pgfqpoint{4.900000in}{4.900000in}}%
\pgfusepath{clip}%
\pgfsetbuttcap%
\pgfsetroundjoin%
\definecolor{currentfill}{rgb}{0.619423,0.619423,0.619423}%
\pgfsetfillcolor{currentfill}%
\pgfsetfillopacity{0.900000}%
\pgfsetlinewidth{0.501875pt}%
\definecolor{currentstroke}{rgb}{0.200000,0.200000,0.200000}%
\pgfsetstrokecolor{currentstroke}%
\pgfsetdash{}{0pt}%
\pgfpathmoveto{\pgfqpoint{1.904451in}{2.660866in}}%
\pgfpathlineto{\pgfqpoint{1.951224in}{2.616125in}}%
\pgfpathlineto{\pgfqpoint{1.982100in}{2.596212in}}%
\pgfpathlineto{\pgfqpoint{1.935319in}{2.640980in}}%
\pgfpathlineto{\pgfqpoint{1.904451in}{2.660866in}}%
\pgfpathclose%
\pgfusepath{stroke,fill}%
\end{pgfscope}%
\begin{pgfscope}%
\pgfpathrectangle{\pgfqpoint{0.050000in}{0.050000in}}{\pgfqpoint{4.900000in}{4.900000in}}%
\pgfusepath{clip}%
\pgfsetbuttcap%
\pgfsetroundjoin%
\definecolor{currentfill}{rgb}{1.000000,1.000000,1.000000}%
\pgfsetfillcolor{currentfill}%
\pgfsetfillopacity{0.900000}%
\pgfsetlinewidth{0.501875pt}%
\definecolor{currentstroke}{rgb}{0.200000,0.200000,0.200000}%
\pgfsetstrokecolor{currentstroke}%
\pgfsetdash{}{0pt}%
\pgfpathmoveto{\pgfqpoint{4.209504in}{1.482806in}}%
\pgfpathlineto{\pgfqpoint{4.253045in}{1.494941in}}%
\pgfpathlineto{\pgfqpoint{4.287512in}{1.474167in}}%
\pgfpathlineto{\pgfqpoint{4.243946in}{1.462094in}}%
\pgfpathlineto{\pgfqpoint{4.209504in}{1.482806in}}%
\pgfpathclose%
\pgfusepath{stroke,fill}%
\end{pgfscope}%
\begin{pgfscope}%
\pgfpathrectangle{\pgfqpoint{0.050000in}{0.050000in}}{\pgfqpoint{4.900000in}{4.900000in}}%
\pgfusepath{clip}%
\pgfsetbuttcap%
\pgfsetroundjoin%
\definecolor{currentfill}{rgb}{0.932334,0.932334,0.932334}%
\pgfsetfillcolor{currentfill}%
\pgfsetfillopacity{0.900000}%
\pgfsetlinewidth{0.501875pt}%
\definecolor{currentstroke}{rgb}{0.200000,0.200000,0.200000}%
\pgfsetstrokecolor{currentstroke}%
\pgfsetdash{}{0pt}%
\pgfpathmoveto{\pgfqpoint{2.959574in}{1.837454in}}%
\pgfpathlineto{\pgfqpoint{3.004222in}{1.807621in}}%
\pgfpathlineto{\pgfqpoint{3.036739in}{1.788150in}}%
\pgfpathlineto{\pgfqpoint{2.992080in}{1.817732in}}%
\pgfpathlineto{\pgfqpoint{2.959574in}{1.837454in}}%
\pgfpathclose%
\pgfusepath{stroke,fill}%
\end{pgfscope}%
\begin{pgfscope}%
\pgfpathrectangle{\pgfqpoint{0.050000in}{0.050000in}}{\pgfqpoint{4.900000in}{4.900000in}}%
\pgfusepath{clip}%
\pgfsetbuttcap%
\pgfsetroundjoin%
\definecolor{currentfill}{rgb}{0.990773,0.990773,0.990773}%
\pgfsetfillcolor{currentfill}%
\pgfsetfillopacity{0.900000}%
\pgfsetlinewidth{0.501875pt}%
\definecolor{currentstroke}{rgb}{0.200000,0.200000,0.200000}%
\pgfsetstrokecolor{currentstroke}%
\pgfsetdash{}{0pt}%
\pgfpathmoveto{\pgfqpoint{3.655159in}{1.537447in}}%
\pgfpathlineto{\pgfqpoint{3.699150in}{1.540343in}}%
\pgfpathlineto{\pgfqpoint{3.732776in}{1.520594in}}%
\pgfpathlineto{\pgfqpoint{3.688762in}{1.517778in}}%
\pgfpathlineto{\pgfqpoint{3.655159in}{1.537447in}}%
\pgfpathclose%
\pgfusepath{stroke,fill}%
\end{pgfscope}%
\begin{pgfscope}%
\pgfpathrectangle{\pgfqpoint{0.050000in}{0.050000in}}{\pgfqpoint{4.900000in}{4.900000in}}%
\pgfusepath{clip}%
\pgfsetbuttcap%
\pgfsetroundjoin%
\definecolor{currentfill}{rgb}{0.858762,0.858762,0.858762}%
\pgfsetfillcolor{currentfill}%
\pgfsetfillopacity{0.900000}%
\pgfsetlinewidth{0.501875pt}%
\definecolor{currentstroke}{rgb}{0.200000,0.200000,0.200000}%
\pgfsetstrokecolor{currentstroke}%
\pgfsetdash{}{0pt}%
\pgfpathmoveto{\pgfqpoint{2.618514in}{2.087148in}}%
\pgfpathlineto{\pgfqpoint{2.663793in}{2.046188in}}%
\pgfpathlineto{\pgfqpoint{2.695774in}{2.025983in}}%
\pgfpathlineto{\pgfqpoint{2.650491in}{2.066772in}}%
\pgfpathlineto{\pgfqpoint{2.618514in}{2.087148in}}%
\pgfpathclose%
\pgfusepath{stroke,fill}%
\end{pgfscope}%
\begin{pgfscope}%
\pgfpathrectangle{\pgfqpoint{0.050000in}{0.050000in}}{\pgfqpoint{4.900000in}{4.900000in}}%
\pgfusepath{clip}%
\pgfsetbuttcap%
\pgfsetroundjoin%
\definecolor{currentfill}{rgb}{0.521492,0.521492,0.521492}%
\pgfsetfillcolor{currentfill}%
\pgfsetfillopacity{0.900000}%
\pgfsetlinewidth{0.501875pt}%
\definecolor{currentstroke}{rgb}{0.200000,0.200000,0.200000}%
\pgfsetstrokecolor{currentstroke}%
\pgfsetdash{}{0pt}%
\pgfpathmoveto{\pgfqpoint{1.640504in}{2.874829in}}%
\pgfpathlineto{\pgfqpoint{1.687842in}{2.829629in}}%
\pgfpathlineto{\pgfqpoint{1.718317in}{2.810042in}}%
\pgfpathlineto{\pgfqpoint{1.670970in}{2.855270in}}%
\pgfpathlineto{\pgfqpoint{1.640504in}{2.874829in}}%
\pgfpathclose%
\pgfusepath{stroke,fill}%
\end{pgfscope}%
\begin{pgfscope}%
\pgfpathrectangle{\pgfqpoint{0.050000in}{0.050000in}}{\pgfqpoint{4.900000in}{4.900000in}}%
\pgfusepath{clip}%
\pgfsetbuttcap%
\pgfsetroundjoin%
\definecolor{currentfill}{rgb}{0.784667,0.784667,0.784667}%
\pgfsetfillcolor{currentfill}%
\pgfsetfillopacity{0.900000}%
\pgfsetlinewidth{0.501875pt}%
\definecolor{currentstroke}{rgb}{0.200000,0.200000,0.200000}%
\pgfsetstrokecolor{currentstroke}%
\pgfsetdash{}{0pt}%
\pgfpathmoveto{\pgfqpoint{2.354666in}{2.297823in}}%
\pgfpathlineto{\pgfqpoint{2.400501in}{2.254101in}}%
\pgfpathlineto{\pgfqpoint{2.432077in}{2.233727in}}%
\pgfpathlineto{\pgfqpoint{2.386238in}{2.277424in}}%
\pgfpathlineto{\pgfqpoint{2.354666in}{2.297823in}}%
\pgfpathclose%
\pgfusepath{stroke,fill}%
\end{pgfscope}%
\begin{pgfscope}%
\pgfpathrectangle{\pgfqpoint{0.050000in}{0.050000in}}{\pgfqpoint{4.900000in}{4.900000in}}%
\pgfusepath{clip}%
\pgfsetbuttcap%
\pgfsetroundjoin%
\definecolor{currentfill}{rgb}{0.998155,0.998155,0.998155}%
\pgfsetfillcolor{currentfill}%
\pgfsetfillopacity{0.900000}%
\pgfsetlinewidth{0.501875pt}%
\definecolor{currentstroke}{rgb}{0.200000,0.200000,0.200000}%
\pgfsetstrokecolor{currentstroke}%
\pgfsetdash{}{0pt}%
\pgfpathmoveto{\pgfqpoint{4.009968in}{1.490226in}}%
\pgfpathlineto{\pgfqpoint{4.053706in}{1.500796in}}%
\pgfpathlineto{\pgfqpoint{4.087883in}{1.480350in}}%
\pgfpathlineto{\pgfqpoint{4.044118in}{1.469859in}}%
\pgfpathlineto{\pgfqpoint{4.009968in}{1.490226in}}%
\pgfpathclose%
\pgfusepath{stroke,fill}%
\end{pgfscope}%
\begin{pgfscope}%
\pgfpathrectangle{\pgfqpoint{0.050000in}{0.050000in}}{\pgfqpoint{4.900000in}{4.900000in}}%
\pgfusepath{clip}%
\pgfsetbuttcap%
\pgfsetroundjoin%
\definecolor{currentfill}{rgb}{0.424083,0.424083,0.424083}%
\pgfsetfillcolor{currentfill}%
\pgfsetfillopacity{0.900000}%
\pgfsetlinewidth{0.501875pt}%
\definecolor{currentstroke}{rgb}{0.200000,0.200000,0.200000}%
\pgfsetstrokecolor{currentstroke}%
\pgfsetdash{}{0pt}%
\pgfpathmoveto{\pgfqpoint{1.376469in}{3.088866in}}%
\pgfpathlineto{\pgfqpoint{1.424372in}{3.043208in}}%
\pgfpathlineto{\pgfqpoint{1.454445in}{3.023946in}}%
\pgfpathlineto{\pgfqpoint{1.406532in}{3.069634in}}%
\pgfpathlineto{\pgfqpoint{1.376469in}{3.088866in}}%
\pgfpathclose%
\pgfusepath{stroke,fill}%
\end{pgfscope}%
\begin{pgfscope}%
\pgfpathrectangle{\pgfqpoint{0.050000in}{0.050000in}}{\pgfqpoint{4.900000in}{4.900000in}}%
\pgfusepath{clip}%
\pgfsetbuttcap%
\pgfsetroundjoin%
\definecolor{currentfill}{rgb}{0.915356,0.915356,0.915356}%
\pgfsetfillcolor{currentfill}%
\pgfsetfillopacity{0.900000}%
\pgfsetlinewidth{0.501875pt}%
\definecolor{currentstroke}{rgb}{0.200000,0.200000,0.200000}%
\pgfsetstrokecolor{currentstroke}%
\pgfsetdash{}{0pt}%
\pgfpathmoveto{\pgfqpoint{2.882382in}{1.890487in}}%
\pgfpathlineto{\pgfqpoint{2.927169in}{1.857116in}}%
\pgfpathlineto{\pgfqpoint{2.959574in}{1.837454in}}%
\pgfpathlineto{\pgfqpoint{2.914778in}{1.870565in}}%
\pgfpathlineto{\pgfqpoint{2.882382in}{1.890487in}}%
\pgfpathclose%
\pgfusepath{stroke,fill}%
\end{pgfscope}%
\begin{pgfscope}%
\pgfpathrectangle{\pgfqpoint{0.050000in}{0.050000in}}{\pgfqpoint{4.900000in}{4.900000in}}%
\pgfusepath{clip}%
\pgfsetbuttcap%
\pgfsetroundjoin%
\definecolor{currentfill}{rgb}{0.691396,0.691396,0.691396}%
\pgfsetfillcolor{currentfill}%
\pgfsetfillopacity{0.900000}%
\pgfsetlinewidth{0.501875pt}%
\definecolor{currentstroke}{rgb}{0.200000,0.200000,0.200000}%
\pgfsetstrokecolor{currentstroke}%
\pgfsetdash{}{0pt}%
\pgfpathmoveto{\pgfqpoint{2.090734in}{2.511572in}}%
\pgfpathlineto{\pgfqpoint{2.137134in}{2.467133in}}%
\pgfpathlineto{\pgfqpoint{2.168308in}{2.446987in}}%
\pgfpathlineto{\pgfqpoint{2.121903in}{2.491450in}}%
\pgfpathlineto{\pgfqpoint{2.090734in}{2.511572in}}%
\pgfpathclose%
\pgfusepath{stroke,fill}%
\end{pgfscope}%
\begin{pgfscope}%
\pgfpathrectangle{\pgfqpoint{0.050000in}{0.050000in}}{\pgfqpoint{4.900000in}{4.900000in}}%
\pgfusepath{clip}%
\pgfsetbuttcap%
\pgfsetroundjoin%
\definecolor{currentfill}{rgb}{1.000000,1.000000,1.000000}%
\pgfsetfillcolor{currentfill}%
\pgfsetfillopacity{0.900000}%
\pgfsetlinewidth{0.501875pt}%
\definecolor{currentstroke}{rgb}{0.200000,0.200000,0.200000}%
\pgfsetstrokecolor{currentstroke}%
\pgfsetdash{}{0pt}%
\pgfpathmoveto{\pgfqpoint{4.287512in}{1.474167in}}%
\pgfpathlineto{\pgfqpoint{4.330986in}{1.486438in}}%
\pgfpathlineto{\pgfqpoint{4.365587in}{1.465539in}}%
\pgfpathlineto{\pgfqpoint{4.322087in}{1.453328in}}%
\pgfpathlineto{\pgfqpoint{4.287512in}{1.474167in}}%
\pgfpathclose%
\pgfusepath{stroke,fill}%
\end{pgfscope}%
\begin{pgfscope}%
\pgfpathrectangle{\pgfqpoint{0.050000in}{0.050000in}}{\pgfqpoint{4.900000in}{4.900000in}}%
\pgfusepath{clip}%
\pgfsetbuttcap%
\pgfsetroundjoin%
\definecolor{currentfill}{rgb}{0.992618,0.992618,0.992618}%
\pgfsetfillcolor{currentfill}%
\pgfsetfillopacity{0.900000}%
\pgfsetlinewidth{0.501875pt}%
\definecolor{currentstroke}{rgb}{0.200000,0.200000,0.200000}%
\pgfsetstrokecolor{currentstroke}%
\pgfsetdash{}{0pt}%
\pgfpathmoveto{\pgfqpoint{3.732776in}{1.520594in}}%
\pgfpathlineto{\pgfqpoint{3.776725in}{1.525650in}}%
\pgfpathlineto{\pgfqpoint{3.810481in}{1.505741in}}%
\pgfpathlineto{\pgfqpoint{3.766508in}{1.500773in}}%
\pgfpathlineto{\pgfqpoint{3.732776in}{1.520594in}}%
\pgfpathclose%
\pgfusepath{stroke,fill}%
\end{pgfscope}%
\begin{pgfscope}%
\pgfpathrectangle{\pgfqpoint{0.050000in}{0.050000in}}{\pgfqpoint{4.900000in}{4.900000in}}%
\pgfusepath{clip}%
\pgfsetbuttcap%
\pgfsetroundjoin%
\definecolor{currentfill}{rgb}{0.330704,0.330704,0.330704}%
\pgfsetfillcolor{currentfill}%
\pgfsetfillopacity{0.900000}%
\pgfsetlinewidth{0.501875pt}%
\definecolor{currentstroke}{rgb}{0.200000,0.200000,0.200000}%
\pgfsetstrokecolor{currentstroke}%
\pgfsetdash{}{0pt}%
\pgfpathmoveto{\pgfqpoint{1.112345in}{3.302978in}}%
\pgfpathlineto{\pgfqpoint{1.160813in}{3.256861in}}%
\pgfpathlineto{\pgfqpoint{1.190485in}{3.237924in}}%
\pgfpathlineto{\pgfqpoint{1.142005in}{3.284073in}}%
\pgfpathlineto{\pgfqpoint{1.112345in}{3.302978in}}%
\pgfpathclose%
\pgfusepath{stroke,fill}%
\end{pgfscope}%
\begin{pgfscope}%
\pgfpathrectangle{\pgfqpoint{0.050000in}{0.050000in}}{\pgfqpoint{4.900000in}{4.900000in}}%
\pgfusepath{clip}%
\pgfsetbuttcap%
\pgfsetroundjoin%
\definecolor{currentfill}{rgb}{0.586082,0.586082,0.586082}%
\pgfsetfillcolor{currentfill}%
\pgfsetfillopacity{0.900000}%
\pgfsetlinewidth{0.501875pt}%
\definecolor{currentstroke}{rgb}{0.200000,0.200000,0.200000}%
\pgfsetstrokecolor{currentstroke}%
\pgfsetdash{}{0pt}%
\pgfpathmoveto{\pgfqpoint{1.826713in}{2.725593in}}%
\pgfpathlineto{\pgfqpoint{1.873677in}{2.680690in}}%
\pgfpathlineto{\pgfqpoint{1.904451in}{2.660866in}}%
\pgfpathlineto{\pgfqpoint{1.857479in}{2.705796in}}%
\pgfpathlineto{\pgfqpoint{1.826713in}{2.725593in}}%
\pgfpathclose%
\pgfusepath{stroke,fill}%
\end{pgfscope}%
\begin{pgfscope}%
\pgfpathrectangle{\pgfqpoint{0.050000in}{0.050000in}}{\pgfqpoint{4.900000in}{4.900000in}}%
\pgfusepath{clip}%
\pgfsetbuttcap%
\pgfsetroundjoin%
\definecolor{currentfill}{rgb}{0.836340,0.836340,0.836340}%
\pgfsetfillcolor{currentfill}%
\pgfsetfillopacity{0.900000}%
\pgfsetlinewidth{0.501875pt}%
\definecolor{currentstroke}{rgb}{0.200000,0.200000,0.200000}%
\pgfsetstrokecolor{currentstroke}%
\pgfsetdash{}{0pt}%
\pgfpathmoveto{\pgfqpoint{2.541174in}{2.149752in}}%
\pgfpathlineto{\pgfqpoint{2.586637in}{2.107482in}}%
\pgfpathlineto{\pgfqpoint{2.618514in}{2.087148in}}%
\pgfpathlineto{\pgfqpoint{2.573048in}{2.129300in}}%
\pgfpathlineto{\pgfqpoint{2.541174in}{2.149752in}}%
\pgfpathclose%
\pgfusepath{stroke,fill}%
\end{pgfscope}%
\begin{pgfscope}%
\pgfpathrectangle{\pgfqpoint{0.050000in}{0.050000in}}{\pgfqpoint{4.900000in}{4.900000in}}%
\pgfusepath{clip}%
\pgfsetbuttcap%
\pgfsetroundjoin%
\definecolor{currentfill}{rgb}{0.998155,0.998155,0.998155}%
\pgfsetfillcolor{currentfill}%
\pgfsetfillopacity{0.900000}%
\pgfsetlinewidth{0.501875pt}%
\definecolor{currentstroke}{rgb}{0.200000,0.200000,0.200000}%
\pgfsetstrokecolor{currentstroke}%
\pgfsetdash{}{0pt}%
\pgfpathmoveto{\pgfqpoint{4.087883in}{1.480350in}}%
\pgfpathlineto{\pgfqpoint{4.131568in}{1.491621in}}%
\pgfpathlineto{\pgfqpoint{4.165877in}{1.471040in}}%
\pgfpathlineto{\pgfqpoint{4.122166in}{1.459842in}}%
\pgfpathlineto{\pgfqpoint{4.087883in}{1.480350in}}%
\pgfpathclose%
\pgfusepath{stroke,fill}%
\end{pgfscope}%
\begin{pgfscope}%
\pgfpathrectangle{\pgfqpoint{0.050000in}{0.050000in}}{\pgfqpoint{4.900000in}{4.900000in}}%
\pgfusepath{clip}%
\pgfsetbuttcap%
\pgfsetroundjoin%
\definecolor{currentfill}{rgb}{0.487043,0.487043,0.487043}%
\pgfsetfillcolor{currentfill}%
\pgfsetfillopacity{0.900000}%
\pgfsetlinewidth{0.501875pt}%
\definecolor{currentstroke}{rgb}{0.200000,0.200000,0.200000}%
\pgfsetstrokecolor{currentstroke}%
\pgfsetdash{}{0pt}%
\pgfpathmoveto{\pgfqpoint{1.562603in}{2.939690in}}%
\pgfpathlineto{\pgfqpoint{1.610132in}{2.894328in}}%
\pgfpathlineto{\pgfqpoint{1.640504in}{2.874829in}}%
\pgfpathlineto{\pgfqpoint{1.592966in}{2.920220in}}%
\pgfpathlineto{\pgfqpoint{1.562603in}{2.939690in}}%
\pgfpathclose%
\pgfusepath{stroke,fill}%
\end{pgfscope}%
\begin{pgfscope}%
\pgfpathrectangle{\pgfqpoint{0.050000in}{0.050000in}}{\pgfqpoint{4.900000in}{4.900000in}}%
\pgfusepath{clip}%
\pgfsetbuttcap%
\pgfsetroundjoin%
\definecolor{currentfill}{rgb}{0.898378,0.898378,0.898378}%
\pgfsetfillcolor{currentfill}%
\pgfsetfillopacity{0.900000}%
\pgfsetlinewidth{0.501875pt}%
\definecolor{currentstroke}{rgb}{0.200000,0.200000,0.200000}%
\pgfsetstrokecolor{currentstroke}%
\pgfsetdash{}{0pt}%
\pgfpathmoveto{\pgfqpoint{2.805148in}{1.946772in}}%
\pgfpathlineto{\pgfqpoint{2.850089in}{1.910357in}}%
\pgfpathlineto{\pgfqpoint{2.882382in}{1.890487in}}%
\pgfpathlineto{\pgfqpoint{2.837434in}{1.926653in}}%
\pgfpathlineto{\pgfqpoint{2.805148in}{1.946772in}}%
\pgfpathclose%
\pgfusepath{stroke,fill}%
\end{pgfscope}%
\begin{pgfscope}%
\pgfpathrectangle{\pgfqpoint{0.050000in}{0.050000in}}{\pgfqpoint{4.900000in}{4.900000in}}%
\pgfusepath{clip}%
\pgfsetbuttcap%
\pgfsetroundjoin%
\definecolor{currentfill}{rgb}{0.994464,0.994464,0.994464}%
\pgfsetfillcolor{currentfill}%
\pgfsetfillopacity{0.900000}%
\pgfsetlinewidth{0.501875pt}%
\definecolor{currentstroke}{rgb}{0.200000,0.200000,0.200000}%
\pgfsetstrokecolor{currentstroke}%
\pgfsetdash{}{0pt}%
\pgfpathmoveto{\pgfqpoint{3.810481in}{1.505741in}}%
\pgfpathlineto{\pgfqpoint{3.854387in}{1.512607in}}%
\pgfpathlineto{\pgfqpoint{3.888273in}{1.492540in}}%
\pgfpathlineto{\pgfqpoint{3.844343in}{1.485764in}}%
\pgfpathlineto{\pgfqpoint{3.810481in}{1.505741in}}%
\pgfpathclose%
\pgfusepath{stroke,fill}%
\end{pgfscope}%
\begin{pgfscope}%
\pgfpathrectangle{\pgfqpoint{0.050000in}{0.050000in}}{\pgfqpoint{4.900000in}{4.900000in}}%
\pgfusepath{clip}%
\pgfsetbuttcap%
\pgfsetroundjoin%
\definecolor{currentfill}{rgb}{0.760554,0.760554,0.760554}%
\pgfsetfillcolor{currentfill}%
\pgfsetfillopacity{0.900000}%
\pgfsetlinewidth{0.501875pt}%
\definecolor{currentstroke}{rgb}{0.200000,0.200000,0.200000}%
\pgfsetstrokecolor{currentstroke}%
\pgfsetdash{}{0pt}%
\pgfpathmoveto{\pgfqpoint{2.277167in}{2.362218in}}%
\pgfpathlineto{\pgfqpoint{2.323192in}{2.318167in}}%
\pgfpathlineto{\pgfqpoint{2.354666in}{2.297823in}}%
\pgfpathlineto{\pgfqpoint{2.308637in}{2.341875in}}%
\pgfpathlineto{\pgfqpoint{2.277167in}{2.362218in}}%
\pgfpathclose%
\pgfusepath{stroke,fill}%
\end{pgfscope}%
\begin{pgfscope}%
\pgfpathrectangle{\pgfqpoint{0.050000in}{0.050000in}}{\pgfqpoint{4.900000in}{4.900000in}}%
\pgfusepath{clip}%
\pgfsetbuttcap%
\pgfsetroundjoin%
\definecolor{currentfill}{rgb}{0.395663,0.395663,0.395663}%
\pgfsetfillcolor{currentfill}%
\pgfsetfillopacity{0.900000}%
\pgfsetlinewidth{0.501875pt}%
\definecolor{currentstroke}{rgb}{0.200000,0.200000,0.200000}%
\pgfsetstrokecolor{currentstroke}%
\pgfsetdash{}{0pt}%
\pgfpathmoveto{\pgfqpoint{1.298404in}{3.153861in}}%
\pgfpathlineto{\pgfqpoint{1.346498in}{3.108039in}}%
\pgfpathlineto{\pgfqpoint{1.376469in}{3.088866in}}%
\pgfpathlineto{\pgfqpoint{1.328363in}{3.134718in}}%
\pgfpathlineto{\pgfqpoint{1.298404in}{3.153861in}}%
\pgfpathclose%
\pgfusepath{stroke,fill}%
\end{pgfscope}%
\begin{pgfscope}%
\pgfpathrectangle{\pgfqpoint{0.050000in}{0.050000in}}{\pgfqpoint{4.900000in}{4.900000in}}%
\pgfusepath{clip}%
\pgfsetbuttcap%
\pgfsetroundjoin%
\definecolor{currentfill}{rgb}{0.657809,0.657809,0.657809}%
\pgfsetfillcolor{currentfill}%
\pgfsetfillopacity{0.900000}%
\pgfsetlinewidth{0.501875pt}%
\definecolor{currentstroke}{rgb}{0.200000,0.200000,0.200000}%
\pgfsetstrokecolor{currentstroke}%
\pgfsetdash{}{0pt}%
\pgfpathmoveto{\pgfqpoint{2.013071in}{2.576238in}}%
\pgfpathlineto{\pgfqpoint{2.059661in}{2.531634in}}%
\pgfpathlineto{\pgfqpoint{2.090734in}{2.511572in}}%
\pgfpathlineto{\pgfqpoint{2.044138in}{2.556203in}}%
\pgfpathlineto{\pgfqpoint{2.013071in}{2.576238in}}%
\pgfpathclose%
\pgfusepath{stroke,fill}%
\end{pgfscope}%
\begin{pgfscope}%
\pgfpathrectangle{\pgfqpoint{0.050000in}{0.050000in}}{\pgfqpoint{4.900000in}{4.900000in}}%
\pgfusepath{clip}%
\pgfsetbuttcap%
\pgfsetroundjoin%
\definecolor{currentfill}{rgb}{0.968627,0.968627,0.968627}%
\pgfsetfillcolor{currentfill}%
\pgfsetfillopacity{0.900000}%
\pgfsetlinewidth{0.501875pt}%
\definecolor{currentstroke}{rgb}{0.200000,0.200000,0.200000}%
\pgfsetstrokecolor{currentstroke}%
\pgfsetdash{}{0pt}%
\pgfpathmoveto{\pgfqpoint{3.301239in}{1.645939in}}%
\pgfpathlineto{\pgfqpoint{3.345507in}{1.632582in}}%
\pgfpathlineto{\pgfqpoint{3.378604in}{1.613325in}}%
\pgfpathlineto{\pgfqpoint{3.334317in}{1.626606in}}%
\pgfpathlineto{\pgfqpoint{3.301239in}{1.645939in}}%
\pgfpathclose%
\pgfusepath{stroke,fill}%
\end{pgfscope}%
\begin{pgfscope}%
\pgfpathrectangle{\pgfqpoint{0.050000in}{0.050000in}}{\pgfqpoint{4.900000in}{4.900000in}}%
\pgfusepath{clip}%
\pgfsetbuttcap%
\pgfsetroundjoin%
\definecolor{currentfill}{rgb}{0.974164,0.974164,0.974164}%
\pgfsetfillcolor{currentfill}%
\pgfsetfillopacity{0.900000}%
\pgfsetlinewidth{0.501875pt}%
\definecolor{currentstroke}{rgb}{0.200000,0.200000,0.200000}%
\pgfsetstrokecolor{currentstroke}%
\pgfsetdash{}{0pt}%
\pgfpathmoveto{\pgfqpoint{3.378604in}{1.613325in}}%
\pgfpathlineto{\pgfqpoint{3.422809in}{1.603856in}}%
\pgfpathlineto{\pgfqpoint{3.456031in}{1.584545in}}%
\pgfpathlineto{\pgfqpoint{3.411806in}{1.593984in}}%
\pgfpathlineto{\pgfqpoint{3.378604in}{1.613325in}}%
\pgfpathclose%
\pgfusepath{stroke,fill}%
\end{pgfscope}%
\begin{pgfscope}%
\pgfpathrectangle{\pgfqpoint{0.050000in}{0.050000in}}{\pgfqpoint{4.900000in}{4.900000in}}%
\pgfusepath{clip}%
\pgfsetbuttcap%
\pgfsetroundjoin%
\definecolor{currentfill}{rgb}{0.961246,0.961246,0.961246}%
\pgfsetfillcolor{currentfill}%
\pgfsetfillopacity{0.900000}%
\pgfsetlinewidth{0.501875pt}%
\definecolor{currentstroke}{rgb}{0.200000,0.200000,0.200000}%
\pgfsetstrokecolor{currentstroke}%
\pgfsetdash{}{0pt}%
\pgfpathmoveto{\pgfqpoint{3.223920in}{1.682627in}}%
\pgfpathlineto{\pgfqpoint{3.268264in}{1.665188in}}%
\pgfpathlineto{\pgfqpoint{3.301239in}{1.645939in}}%
\pgfpathlineto{\pgfqpoint{3.256877in}{1.663252in}}%
\pgfpathlineto{\pgfqpoint{3.223920in}{1.682627in}}%
\pgfpathclose%
\pgfusepath{stroke,fill}%
\end{pgfscope}%
\begin{pgfscope}%
\pgfpathrectangle{\pgfqpoint{0.050000in}{0.050000in}}{\pgfqpoint{4.900000in}{4.900000in}}%
\pgfusepath{clip}%
\pgfsetbuttcap%
\pgfsetroundjoin%
\definecolor{currentfill}{rgb}{0.295271,0.295271,0.295271}%
\pgfsetfillcolor{currentfill}%
\pgfsetfillopacity{0.900000}%
\pgfsetlinewidth{0.501875pt}%
\definecolor{currentstroke}{rgb}{0.200000,0.200000,0.200000}%
\pgfsetstrokecolor{currentstroke}%
\pgfsetdash{}{0pt}%
\pgfpathmoveto{\pgfqpoint{1.034116in}{3.368106in}}%
\pgfpathlineto{\pgfqpoint{1.082776in}{3.321825in}}%
\pgfpathlineto{\pgfqpoint{1.112345in}{3.302978in}}%
\pgfpathlineto{\pgfqpoint{1.063672in}{3.349291in}}%
\pgfpathlineto{\pgfqpoint{1.034116in}{3.368106in}}%
\pgfpathclose%
\pgfusepath{stroke,fill}%
\end{pgfscope}%
\begin{pgfscope}%
\pgfpathrectangle{\pgfqpoint{0.050000in}{0.050000in}}{\pgfqpoint{4.900000in}{4.900000in}}%
\pgfusepath{clip}%
\pgfsetbuttcap%
\pgfsetroundjoin%
\definecolor{currentfill}{rgb}{0.979700,0.979700,0.979700}%
\pgfsetfillcolor{currentfill}%
\pgfsetfillopacity{0.900000}%
\pgfsetlinewidth{0.501875pt}%
\definecolor{currentstroke}{rgb}{0.200000,0.200000,0.200000}%
\pgfsetstrokecolor{currentstroke}%
\pgfsetdash{}{0pt}%
\pgfpathmoveto{\pgfqpoint{3.456031in}{1.584545in}}%
\pgfpathlineto{\pgfqpoint{3.500180in}{1.578680in}}%
\pgfpathlineto{\pgfqpoint{3.533527in}{1.559276in}}%
\pgfpathlineto{\pgfqpoint{3.489357in}{1.565150in}}%
\pgfpathlineto{\pgfqpoint{3.456031in}{1.584545in}}%
\pgfpathclose%
\pgfusepath{stroke,fill}%
\end{pgfscope}%
\begin{pgfscope}%
\pgfpathrectangle{\pgfqpoint{0.050000in}{0.050000in}}{\pgfqpoint{4.900000in}{4.900000in}}%
\pgfusepath{clip}%
\pgfsetbuttcap%
\pgfsetroundjoin%
\definecolor{currentfill}{rgb}{0.996309,0.996309,0.996309}%
\pgfsetfillcolor{currentfill}%
\pgfsetfillopacity{0.900000}%
\pgfsetlinewidth{0.501875pt}%
\definecolor{currentstroke}{rgb}{0.200000,0.200000,0.200000}%
\pgfsetstrokecolor{currentstroke}%
\pgfsetdash{}{0pt}%
\pgfpathmoveto{\pgfqpoint{3.888273in}{1.492540in}}%
\pgfpathlineto{\pgfqpoint{3.932135in}{1.500894in}}%
\pgfpathlineto{\pgfqpoint{3.966153in}{1.480674in}}%
\pgfpathlineto{\pgfqpoint{3.922267in}{1.472407in}}%
\pgfpathlineto{\pgfqpoint{3.888273in}{1.492540in}}%
\pgfpathclose%
\pgfusepath{stroke,fill}%
\end{pgfscope}%
\begin{pgfscope}%
\pgfpathrectangle{\pgfqpoint{0.050000in}{0.050000in}}{\pgfqpoint{4.900000in}{4.900000in}}%
\pgfusepath{clip}%
\pgfsetbuttcap%
\pgfsetroundjoin%
\definecolor{currentfill}{rgb}{0.551634,0.551634,0.551634}%
\pgfsetfillcolor{currentfill}%
\pgfsetfillopacity{0.900000}%
\pgfsetlinewidth{0.501875pt}%
\definecolor{currentstroke}{rgb}{0.200000,0.200000,0.200000}%
\pgfsetstrokecolor{currentstroke}%
\pgfsetdash{}{0pt}%
\pgfpathmoveto{\pgfqpoint{1.748886in}{2.790394in}}%
\pgfpathlineto{\pgfqpoint{1.796041in}{2.745329in}}%
\pgfpathlineto{\pgfqpoint{1.826713in}{2.725593in}}%
\pgfpathlineto{\pgfqpoint{1.779549in}{2.770686in}}%
\pgfpathlineto{\pgfqpoint{1.748886in}{2.790394in}}%
\pgfpathclose%
\pgfusepath{stroke,fill}%
\end{pgfscope}%
\begin{pgfscope}%
\pgfpathrectangle{\pgfqpoint{0.050000in}{0.050000in}}{\pgfqpoint{4.900000in}{4.900000in}}%
\pgfusepath{clip}%
\pgfsetbuttcap%
\pgfsetroundjoin%
\definecolor{currentfill}{rgb}{0.953864,0.953864,0.953864}%
\pgfsetfillcolor{currentfill}%
\pgfsetfillopacity{0.900000}%
\pgfsetlinewidth{0.501875pt}%
\definecolor{currentstroke}{rgb}{0.200000,0.200000,0.200000}%
\pgfsetstrokecolor{currentstroke}%
\pgfsetdash{}{0pt}%
\pgfpathmoveto{\pgfqpoint{3.146632in}{1.723520in}}%
\pgfpathlineto{\pgfqpoint{3.191066in}{1.701921in}}%
\pgfpathlineto{\pgfqpoint{3.223920in}{1.682627in}}%
\pgfpathlineto{\pgfqpoint{3.179471in}{1.704051in}}%
\pgfpathlineto{\pgfqpoint{3.146632in}{1.723520in}}%
\pgfpathclose%
\pgfusepath{stroke,fill}%
\end{pgfscope}%
\begin{pgfscope}%
\pgfpathrectangle{\pgfqpoint{0.050000in}{0.050000in}}{\pgfqpoint{4.900000in}{4.900000in}}%
\pgfusepath{clip}%
\pgfsetbuttcap%
\pgfsetroundjoin%
\definecolor{currentfill}{rgb}{1.000000,1.000000,1.000000}%
\pgfsetfillcolor{currentfill}%
\pgfsetfillopacity{0.900000}%
\pgfsetlinewidth{0.501875pt}%
\definecolor{currentstroke}{rgb}{0.200000,0.200000,0.200000}%
\pgfsetstrokecolor{currentstroke}%
\pgfsetdash{}{0pt}%
\pgfpathmoveto{\pgfqpoint{4.165877in}{1.471040in}}%
\pgfpathlineto{\pgfqpoint{4.209504in}{1.482806in}}%
\pgfpathlineto{\pgfqpoint{4.243946in}{1.462094in}}%
\pgfpathlineto{\pgfqpoint{4.200293in}{1.450396in}}%
\pgfpathlineto{\pgfqpoint{4.165877in}{1.471040in}}%
\pgfpathclose%
\pgfusepath{stroke,fill}%
\end{pgfscope}%
\begin{pgfscope}%
\pgfpathrectangle{\pgfqpoint{0.050000in}{0.050000in}}{\pgfqpoint{4.900000in}{4.900000in}}%
\pgfusepath{clip}%
\pgfsetbuttcap%
\pgfsetroundjoin%
\definecolor{currentfill}{rgb}{0.878570,0.878570,0.878570}%
\pgfsetfillcolor{currentfill}%
\pgfsetfillopacity{0.900000}%
\pgfsetlinewidth{0.501875pt}%
\definecolor{currentstroke}{rgb}{0.200000,0.200000,0.200000}%
\pgfsetstrokecolor{currentstroke}%
\pgfsetdash{}{0pt}%
\pgfpathmoveto{\pgfqpoint{2.727854in}{2.005736in}}%
\pgfpathlineto{\pgfqpoint{2.772962in}{1.966845in}}%
\pgfpathlineto{\pgfqpoint{2.805148in}{1.946772in}}%
\pgfpathlineto{\pgfqpoint{2.760035in}{1.985446in}}%
\pgfpathlineto{\pgfqpoint{2.727854in}{2.005736in}}%
\pgfpathclose%
\pgfusepath{stroke,fill}%
\end{pgfscope}%
\begin{pgfscope}%
\pgfpathrectangle{\pgfqpoint{0.050000in}{0.050000in}}{\pgfqpoint{4.900000in}{4.900000in}}%
\pgfusepath{clip}%
\pgfsetbuttcap%
\pgfsetroundjoin%
\definecolor{currentfill}{rgb}{0.812226,0.812226,0.812226}%
\pgfsetfillcolor{currentfill}%
\pgfsetfillopacity{0.900000}%
\pgfsetlinewidth{0.501875pt}%
\definecolor{currentstroke}{rgb}{0.200000,0.200000,0.200000}%
\pgfsetstrokecolor{currentstroke}%
\pgfsetdash{}{0pt}%
\pgfpathmoveto{\pgfqpoint{2.463750in}{2.213301in}}%
\pgfpathlineto{\pgfqpoint{2.509399in}{2.170159in}}%
\pgfpathlineto{\pgfqpoint{2.541174in}{2.149752in}}%
\pgfpathlineto{\pgfqpoint{2.495521in}{2.192826in}}%
\pgfpathlineto{\pgfqpoint{2.463750in}{2.213301in}}%
\pgfpathclose%
\pgfusepath{stroke,fill}%
\end{pgfscope}%
\begin{pgfscope}%
\pgfpathrectangle{\pgfqpoint{0.050000in}{0.050000in}}{\pgfqpoint{4.900000in}{4.900000in}}%
\pgfusepath{clip}%
\pgfsetbuttcap%
\pgfsetroundjoin%
\definecolor{currentfill}{rgb}{0.983391,0.983391,0.983391}%
\pgfsetfillcolor{currentfill}%
\pgfsetfillopacity{0.900000}%
\pgfsetlinewidth{0.501875pt}%
\definecolor{currentstroke}{rgb}{0.200000,0.200000,0.200000}%
\pgfsetstrokecolor{currentstroke}%
\pgfsetdash{}{0pt}%
\pgfpathmoveto{\pgfqpoint{3.533527in}{1.559276in}}%
\pgfpathlineto{\pgfqpoint{3.577628in}{1.556677in}}%
\pgfpathlineto{\pgfqpoint{3.611103in}{1.537150in}}%
\pgfpathlineto{\pgfqpoint{3.566980in}{1.539790in}}%
\pgfpathlineto{\pgfqpoint{3.533527in}{1.559276in}}%
\pgfpathclose%
\pgfusepath{stroke,fill}%
\end{pgfscope}%
\begin{pgfscope}%
\pgfpathrectangle{\pgfqpoint{0.050000in}{0.050000in}}{\pgfqpoint{4.900000in}{4.900000in}}%
\pgfusepath{clip}%
\pgfsetbuttcap%
\pgfsetroundjoin%
\definecolor{currentfill}{rgb}{0.456901,0.456901,0.456901}%
\pgfsetfillcolor{currentfill}%
\pgfsetfillopacity{0.900000}%
\pgfsetlinewidth{0.501875pt}%
\definecolor{currentstroke}{rgb}{0.200000,0.200000,0.200000}%
\pgfsetstrokecolor{currentstroke}%
\pgfsetdash{}{0pt}%
\pgfpathmoveto{\pgfqpoint{1.484612in}{3.004625in}}%
\pgfpathlineto{\pgfqpoint{1.532333in}{2.959100in}}%
\pgfpathlineto{\pgfqpoint{1.562603in}{2.939690in}}%
\pgfpathlineto{\pgfqpoint{1.514872in}{2.985244in}}%
\pgfpathlineto{\pgfqpoint{1.484612in}{3.004625in}}%
\pgfpathclose%
\pgfusepath{stroke,fill}%
\end{pgfscope}%
\begin{pgfscope}%
\pgfpathrectangle{\pgfqpoint{0.050000in}{0.050000in}}{\pgfqpoint{4.900000in}{4.900000in}}%
\pgfusepath{clip}%
\pgfsetbuttcap%
\pgfsetroundjoin%
\definecolor{currentfill}{rgb}{0.944637,0.944637,0.944637}%
\pgfsetfillcolor{currentfill}%
\pgfsetfillopacity{0.900000}%
\pgfsetlinewidth{0.501875pt}%
\definecolor{currentstroke}{rgb}{0.200000,0.200000,0.200000}%
\pgfsetstrokecolor{currentstroke}%
\pgfsetdash{}{0pt}%
\pgfpathmoveto{\pgfqpoint{3.069358in}{1.768610in}}%
\pgfpathlineto{\pgfqpoint{3.113897in}{1.742913in}}%
\pgfpathlineto{\pgfqpoint{3.146632in}{1.723520in}}%
\pgfpathlineto{\pgfqpoint{3.102081in}{1.749001in}}%
\pgfpathlineto{\pgfqpoint{3.069358in}{1.768610in}}%
\pgfpathclose%
\pgfusepath{stroke,fill}%
\end{pgfscope}%
\begin{pgfscope}%
\pgfpathrectangle{\pgfqpoint{0.050000in}{0.050000in}}{\pgfqpoint{4.900000in}{4.900000in}}%
\pgfusepath{clip}%
\pgfsetbuttcap%
\pgfsetroundjoin%
\definecolor{currentfill}{rgb}{0.729781,0.729781,0.729781}%
\pgfsetfillcolor{currentfill}%
\pgfsetfillopacity{0.900000}%
\pgfsetlinewidth{0.501875pt}%
\definecolor{currentstroke}{rgb}{0.200000,0.200000,0.200000}%
\pgfsetstrokecolor{currentstroke}%
\pgfsetdash{}{0pt}%
\pgfpathmoveto{\pgfqpoint{2.199579in}{2.426780in}}%
\pgfpathlineto{\pgfqpoint{2.245794in}{2.382502in}}%
\pgfpathlineto{\pgfqpoint{2.277167in}{2.362218in}}%
\pgfpathlineto{\pgfqpoint{2.230947in}{2.406512in}}%
\pgfpathlineto{\pgfqpoint{2.199579in}{2.426780in}}%
\pgfpathclose%
\pgfusepath{stroke,fill}%
\end{pgfscope}%
\begin{pgfscope}%
\pgfpathrectangle{\pgfqpoint{0.050000in}{0.050000in}}{\pgfqpoint{4.900000in}{4.900000in}}%
\pgfusepath{clip}%
\pgfsetbuttcap%
\pgfsetroundjoin%
\definecolor{currentfill}{rgb}{0.987082,0.987082,0.987082}%
\pgfsetfillcolor{currentfill}%
\pgfsetfillopacity{0.900000}%
\pgfsetlinewidth{0.501875pt}%
\definecolor{currentstroke}{rgb}{0.200000,0.200000,0.200000}%
\pgfsetstrokecolor{currentstroke}%
\pgfsetdash{}{0pt}%
\pgfpathmoveto{\pgfqpoint{3.611103in}{1.537150in}}%
\pgfpathlineto{\pgfqpoint{3.655159in}{1.537447in}}%
\pgfpathlineto{\pgfqpoint{3.688762in}{1.517778in}}%
\pgfpathlineto{\pgfqpoint{3.644683in}{1.517545in}}%
\pgfpathlineto{\pgfqpoint{3.611103in}{1.537150in}}%
\pgfpathclose%
\pgfusepath{stroke,fill}%
\end{pgfscope}%
\begin{pgfscope}%
\pgfpathrectangle{\pgfqpoint{0.050000in}{0.050000in}}{\pgfqpoint{4.900000in}{4.900000in}}%
\pgfusepath{clip}%
\pgfsetbuttcap%
\pgfsetroundjoin%
\definecolor{currentfill}{rgb}{0.363183,0.363183,0.363183}%
\pgfsetfillcolor{currentfill}%
\pgfsetfillopacity{0.900000}%
\pgfsetlinewidth{0.501875pt}%
\definecolor{currentstroke}{rgb}{0.200000,0.200000,0.200000}%
\pgfsetstrokecolor{currentstroke}%
\pgfsetdash{}{0pt}%
\pgfpathmoveto{\pgfqpoint{1.220249in}{3.218930in}}%
\pgfpathlineto{\pgfqpoint{1.268536in}{3.172945in}}%
\pgfpathlineto{\pgfqpoint{1.298404in}{3.153861in}}%
\pgfpathlineto{\pgfqpoint{1.250105in}{3.199876in}}%
\pgfpathlineto{\pgfqpoint{1.220249in}{3.218930in}}%
\pgfpathclose%
\pgfusepath{stroke,fill}%
\end{pgfscope}%
\begin{pgfscope}%
\pgfpathrectangle{\pgfqpoint{0.050000in}{0.050000in}}{\pgfqpoint{4.900000in}{4.900000in}}%
\pgfusepath{clip}%
\pgfsetbuttcap%
\pgfsetroundjoin%
\definecolor{currentfill}{rgb}{0.996309,0.996309,0.996309}%
\pgfsetfillcolor{currentfill}%
\pgfsetfillopacity{0.900000}%
\pgfsetlinewidth{0.501875pt}%
\definecolor{currentstroke}{rgb}{0.200000,0.200000,0.200000}%
\pgfsetstrokecolor{currentstroke}%
\pgfsetdash{}{0pt}%
\pgfpathmoveto{\pgfqpoint{3.966153in}{1.480674in}}%
\pgfpathlineto{\pgfqpoint{4.009968in}{1.490226in}}%
\pgfpathlineto{\pgfqpoint{4.044118in}{1.469859in}}%
\pgfpathlineto{\pgfqpoint{4.000278in}{1.460389in}}%
\pgfpathlineto{\pgfqpoint{3.966153in}{1.480674in}}%
\pgfpathclose%
\pgfusepath{stroke,fill}%
\end{pgfscope}%
\begin{pgfscope}%
\pgfpathrectangle{\pgfqpoint{0.050000in}{0.050000in}}{\pgfqpoint{4.900000in}{4.900000in}}%
\pgfusepath{clip}%
\pgfsetbuttcap%
\pgfsetroundjoin%
\definecolor{currentfill}{rgb}{0.619423,0.619423,0.619423}%
\pgfsetfillcolor{currentfill}%
\pgfsetfillopacity{0.900000}%
\pgfsetlinewidth{0.501875pt}%
\definecolor{currentstroke}{rgb}{0.200000,0.200000,0.200000}%
\pgfsetstrokecolor{currentstroke}%
\pgfsetdash{}{0pt}%
\pgfpathmoveto{\pgfqpoint{1.935319in}{2.640980in}}%
\pgfpathlineto{\pgfqpoint{1.982100in}{2.596212in}}%
\pgfpathlineto{\pgfqpoint{2.013071in}{2.576238in}}%
\pgfpathlineto{\pgfqpoint{1.966284in}{2.621033in}}%
\pgfpathlineto{\pgfqpoint{1.935319in}{2.640980in}}%
\pgfpathclose%
\pgfusepath{stroke,fill}%
\end{pgfscope}%
\begin{pgfscope}%
\pgfpathrectangle{\pgfqpoint{0.050000in}{0.050000in}}{\pgfqpoint{4.900000in}{4.900000in}}%
\pgfusepath{clip}%
\pgfsetbuttcap%
\pgfsetroundjoin%
\definecolor{currentfill}{rgb}{1.000000,1.000000,1.000000}%
\pgfsetfillcolor{currentfill}%
\pgfsetfillopacity{0.900000}%
\pgfsetlinewidth{0.501875pt}%
\definecolor{currentstroke}{rgb}{0.200000,0.200000,0.200000}%
\pgfsetstrokecolor{currentstroke}%
\pgfsetdash{}{0pt}%
\pgfpathmoveto{\pgfqpoint{4.243946in}{1.462094in}}%
\pgfpathlineto{\pgfqpoint{4.287512in}{1.474167in}}%
\pgfpathlineto{\pgfqpoint{4.322087in}{1.453328in}}%
\pgfpathlineto{\pgfqpoint{4.278496in}{1.441318in}}%
\pgfpathlineto{\pgfqpoint{4.243946in}{1.462094in}}%
\pgfpathclose%
\pgfusepath{stroke,fill}%
\end{pgfscope}%
\begin{pgfscope}%
\pgfpathrectangle{\pgfqpoint{0.050000in}{0.050000in}}{\pgfqpoint{4.900000in}{4.900000in}}%
\pgfusepath{clip}%
\pgfsetbuttcap%
\pgfsetroundjoin%
\definecolor{currentfill}{rgb}{0.932334,0.932334,0.932334}%
\pgfsetfillcolor{currentfill}%
\pgfsetfillopacity{0.900000}%
\pgfsetlinewidth{0.501875pt}%
\definecolor{currentstroke}{rgb}{0.200000,0.200000,0.200000}%
\pgfsetstrokecolor{currentstroke}%
\pgfsetdash{}{0pt}%
\pgfpathmoveto{\pgfqpoint{2.992080in}{1.817732in}}%
\pgfpathlineto{\pgfqpoint{3.036739in}{1.788150in}}%
\pgfpathlineto{\pgfqpoint{3.069358in}{1.768610in}}%
\pgfpathlineto{\pgfqpoint{3.024688in}{1.797948in}}%
\pgfpathlineto{\pgfqpoint{2.992080in}{1.817732in}}%
\pgfpathclose%
\pgfusepath{stroke,fill}%
\end{pgfscope}%
\begin{pgfscope}%
\pgfpathrectangle{\pgfqpoint{0.050000in}{0.050000in}}{\pgfqpoint{4.900000in}{4.900000in}}%
\pgfusepath{clip}%
\pgfsetbuttcap%
\pgfsetroundjoin%
\definecolor{currentfill}{rgb}{0.858762,0.858762,0.858762}%
\pgfsetfillcolor{currentfill}%
\pgfsetfillopacity{0.900000}%
\pgfsetlinewidth{0.501875pt}%
\definecolor{currentstroke}{rgb}{0.200000,0.200000,0.200000}%
\pgfsetstrokecolor{currentstroke}%
\pgfsetdash{}{0pt}%
\pgfpathmoveto{\pgfqpoint{2.650491in}{2.066772in}}%
\pgfpathlineto{\pgfqpoint{2.695774in}{2.025983in}}%
\pgfpathlineto{\pgfqpoint{2.727854in}{2.005736in}}%
\pgfpathlineto{\pgfqpoint{2.682566in}{2.046351in}}%
\pgfpathlineto{\pgfqpoint{2.650491in}{2.066772in}}%
\pgfpathclose%
\pgfusepath{stroke,fill}%
\end{pgfscope}%
\begin{pgfscope}%
\pgfpathrectangle{\pgfqpoint{0.050000in}{0.050000in}}{\pgfqpoint{4.900000in}{4.900000in}}%
\pgfusepath{clip}%
\pgfsetbuttcap%
\pgfsetroundjoin%
\definecolor{currentfill}{rgb}{0.521492,0.521492,0.521492}%
\pgfsetfillcolor{currentfill}%
\pgfsetfillopacity{0.900000}%
\pgfsetlinewidth{0.501875pt}%
\definecolor{currentstroke}{rgb}{0.200000,0.200000,0.200000}%
\pgfsetstrokecolor{currentstroke}%
\pgfsetdash{}{0pt}%
\pgfpathmoveto{\pgfqpoint{1.670970in}{2.855270in}}%
\pgfpathlineto{\pgfqpoint{1.718317in}{2.810042in}}%
\pgfpathlineto{\pgfqpoint{1.748886in}{2.790394in}}%
\pgfpathlineto{\pgfqpoint{1.701531in}{2.835650in}}%
\pgfpathlineto{\pgfqpoint{1.670970in}{2.855270in}}%
\pgfpathclose%
\pgfusepath{stroke,fill}%
\end{pgfscope}%
\begin{pgfscope}%
\pgfpathrectangle{\pgfqpoint{0.050000in}{0.050000in}}{\pgfqpoint{4.900000in}{4.900000in}}%
\pgfusepath{clip}%
\pgfsetbuttcap%
\pgfsetroundjoin%
\definecolor{currentfill}{rgb}{0.988927,0.988927,0.988927}%
\pgfsetfillcolor{currentfill}%
\pgfsetfillopacity{0.900000}%
\pgfsetlinewidth{0.501875pt}%
\definecolor{currentstroke}{rgb}{0.200000,0.200000,0.200000}%
\pgfsetstrokecolor{currentstroke}%
\pgfsetdash{}{0pt}%
\pgfpathmoveto{\pgfqpoint{3.688762in}{1.517778in}}%
\pgfpathlineto{\pgfqpoint{3.732776in}{1.520594in}}%
\pgfpathlineto{\pgfqpoint{3.766508in}{1.500773in}}%
\pgfpathlineto{\pgfqpoint{3.722471in}{1.498035in}}%
\pgfpathlineto{\pgfqpoint{3.688762in}{1.517778in}}%
\pgfpathclose%
\pgfusepath{stroke,fill}%
\end{pgfscope}%
\begin{pgfscope}%
\pgfpathrectangle{\pgfqpoint{0.050000in}{0.050000in}}{\pgfqpoint{4.900000in}{4.900000in}}%
\pgfusepath{clip}%
\pgfsetbuttcap%
\pgfsetroundjoin%
\definecolor{currentfill}{rgb}{0.784667,0.784667,0.784667}%
\pgfsetfillcolor{currentfill}%
\pgfsetfillopacity{0.900000}%
\pgfsetlinewidth{0.501875pt}%
\definecolor{currentstroke}{rgb}{0.200000,0.200000,0.200000}%
\pgfsetstrokecolor{currentstroke}%
\pgfsetdash{}{0pt}%
\pgfpathmoveto{\pgfqpoint{2.386238in}{2.277424in}}%
\pgfpathlineto{\pgfqpoint{2.432077in}{2.233727in}}%
\pgfpathlineto{\pgfqpoint{2.463750in}{2.213301in}}%
\pgfpathlineto{\pgfqpoint{2.417907in}{2.256970in}}%
\pgfpathlineto{\pgfqpoint{2.386238in}{2.277424in}}%
\pgfpathclose%
\pgfusepath{stroke,fill}%
\end{pgfscope}%
\begin{pgfscope}%
\pgfpathrectangle{\pgfqpoint{0.050000in}{0.050000in}}{\pgfqpoint{4.900000in}{4.900000in}}%
\pgfusepath{clip}%
\pgfsetbuttcap%
\pgfsetroundjoin%
\definecolor{currentfill}{rgb}{0.424083,0.424083,0.424083}%
\pgfsetfillcolor{currentfill}%
\pgfsetfillopacity{0.900000}%
\pgfsetlinewidth{0.501875pt}%
\definecolor{currentstroke}{rgb}{0.200000,0.200000,0.200000}%
\pgfsetstrokecolor{currentstroke}%
\pgfsetdash{}{0pt}%
\pgfpathmoveto{\pgfqpoint{1.406532in}{3.069634in}}%
\pgfpathlineto{\pgfqpoint{1.454445in}{3.023946in}}%
\pgfpathlineto{\pgfqpoint{1.484612in}{3.004625in}}%
\pgfpathlineto{\pgfqpoint{1.436689in}{3.050342in}}%
\pgfpathlineto{\pgfqpoint{1.406532in}{3.069634in}}%
\pgfpathclose%
\pgfusepath{stroke,fill}%
\end{pgfscope}%
\begin{pgfscope}%
\pgfpathrectangle{\pgfqpoint{0.050000in}{0.050000in}}{\pgfqpoint{4.900000in}{4.900000in}}%
\pgfusepath{clip}%
\pgfsetbuttcap%
\pgfsetroundjoin%
\definecolor{currentfill}{rgb}{0.998155,0.998155,0.998155}%
\pgfsetfillcolor{currentfill}%
\pgfsetfillopacity{0.900000}%
\pgfsetlinewidth{0.501875pt}%
\definecolor{currentstroke}{rgb}{0.200000,0.200000,0.200000}%
\pgfsetstrokecolor{currentstroke}%
\pgfsetdash{}{0pt}%
\pgfpathmoveto{\pgfqpoint{4.044118in}{1.469859in}}%
\pgfpathlineto{\pgfqpoint{4.087883in}{1.480350in}}%
\pgfpathlineto{\pgfqpoint{4.122166in}{1.459842in}}%
\pgfpathlineto{\pgfqpoint{4.078376in}{1.449428in}}%
\pgfpathlineto{\pgfqpoint{4.044118in}{1.469859in}}%
\pgfpathclose%
\pgfusepath{stroke,fill}%
\end{pgfscope}%
\begin{pgfscope}%
\pgfpathrectangle{\pgfqpoint{0.050000in}{0.050000in}}{\pgfqpoint{4.900000in}{4.900000in}}%
\pgfusepath{clip}%
\pgfsetbuttcap%
\pgfsetroundjoin%
\definecolor{currentfill}{rgb}{0.915356,0.915356,0.915356}%
\pgfsetfillcolor{currentfill}%
\pgfsetfillopacity{0.900000}%
\pgfsetlinewidth{0.501875pt}%
\definecolor{currentstroke}{rgb}{0.200000,0.200000,0.200000}%
\pgfsetstrokecolor{currentstroke}%
\pgfsetdash{}{0pt}%
\pgfpathmoveto{\pgfqpoint{2.914778in}{1.870565in}}%
\pgfpathlineto{\pgfqpoint{2.959574in}{1.837454in}}%
\pgfpathlineto{\pgfqpoint{2.992080in}{1.817732in}}%
\pgfpathlineto{\pgfqpoint{2.947274in}{1.850588in}}%
\pgfpathlineto{\pgfqpoint{2.914778in}{1.870565in}}%
\pgfpathclose%
\pgfusepath{stroke,fill}%
\end{pgfscope}%
\begin{pgfscope}%
\pgfpathrectangle{\pgfqpoint{0.050000in}{0.050000in}}{\pgfqpoint{4.900000in}{4.900000in}}%
\pgfusepath{clip}%
\pgfsetbuttcap%
\pgfsetroundjoin%
\definecolor{currentfill}{rgb}{0.696194,0.696194,0.696194}%
\pgfsetfillcolor{currentfill}%
\pgfsetfillopacity{0.900000}%
\pgfsetlinewidth{0.501875pt}%
\definecolor{currentstroke}{rgb}{0.200000,0.200000,0.200000}%
\pgfsetstrokecolor{currentstroke}%
\pgfsetdash{}{0pt}%
\pgfpathmoveto{\pgfqpoint{2.121903in}{2.491450in}}%
\pgfpathlineto{\pgfqpoint{2.168308in}{2.446987in}}%
\pgfpathlineto{\pgfqpoint{2.199579in}{2.426780in}}%
\pgfpathlineto{\pgfqpoint{2.153169in}{2.471265in}}%
\pgfpathlineto{\pgfqpoint{2.121903in}{2.491450in}}%
\pgfpathclose%
\pgfusepath{stroke,fill}%
\end{pgfscope}%
\begin{pgfscope}%
\pgfpathrectangle{\pgfqpoint{0.050000in}{0.050000in}}{\pgfqpoint{4.900000in}{4.900000in}}%
\pgfusepath{clip}%
\pgfsetbuttcap%
\pgfsetroundjoin%
\definecolor{currentfill}{rgb}{1.000000,1.000000,1.000000}%
\pgfsetfillcolor{currentfill}%
\pgfsetfillopacity{0.900000}%
\pgfsetlinewidth{0.501875pt}%
\definecolor{currentstroke}{rgb}{0.200000,0.200000,0.200000}%
\pgfsetstrokecolor{currentstroke}%
\pgfsetdash{}{0pt}%
\pgfpathmoveto{\pgfqpoint{4.322087in}{1.453328in}}%
\pgfpathlineto{\pgfqpoint{4.365587in}{1.465539in}}%
\pgfpathlineto{\pgfqpoint{4.400295in}{1.444575in}}%
\pgfpathlineto{\pgfqpoint{4.356770in}{1.432424in}}%
\pgfpathlineto{\pgfqpoint{4.322087in}{1.453328in}}%
\pgfpathclose%
\pgfusepath{stroke,fill}%
\end{pgfscope}%
\begin{pgfscope}%
\pgfpathrectangle{\pgfqpoint{0.050000in}{0.050000in}}{\pgfqpoint{4.900000in}{4.900000in}}%
\pgfusepath{clip}%
\pgfsetbuttcap%
\pgfsetroundjoin%
\definecolor{currentfill}{rgb}{0.992618,0.992618,0.992618}%
\pgfsetfillcolor{currentfill}%
\pgfsetfillopacity{0.900000}%
\pgfsetlinewidth{0.501875pt}%
\definecolor{currentstroke}{rgb}{0.200000,0.200000,0.200000}%
\pgfsetstrokecolor{currentstroke}%
\pgfsetdash{}{0pt}%
\pgfpathmoveto{\pgfqpoint{3.766508in}{1.500773in}}%
\pgfpathlineto{\pgfqpoint{3.810481in}{1.505741in}}%
\pgfpathlineto{\pgfqpoint{3.844343in}{1.485764in}}%
\pgfpathlineto{\pgfqpoint{3.800346in}{1.480881in}}%
\pgfpathlineto{\pgfqpoint{3.766508in}{1.500773in}}%
\pgfpathclose%
\pgfusepath{stroke,fill}%
\end{pgfscope}%
\begin{pgfscope}%
\pgfpathrectangle{\pgfqpoint{0.050000in}{0.050000in}}{\pgfqpoint{4.900000in}{4.900000in}}%
\pgfusepath{clip}%
\pgfsetbuttcap%
\pgfsetroundjoin%
\definecolor{currentfill}{rgb}{0.330704,0.330704,0.330704}%
\pgfsetfillcolor{currentfill}%
\pgfsetfillopacity{0.900000}%
\pgfsetlinewidth{0.501875pt}%
\definecolor{currentstroke}{rgb}{0.200000,0.200000,0.200000}%
\pgfsetstrokecolor{currentstroke}%
\pgfsetdash{}{0pt}%
\pgfpathmoveto{\pgfqpoint{1.142005in}{3.284073in}}%
\pgfpathlineto{\pgfqpoint{1.190485in}{3.237924in}}%
\pgfpathlineto{\pgfqpoint{1.220249in}{3.218930in}}%
\pgfpathlineto{\pgfqpoint{1.171757in}{3.265109in}}%
\pgfpathlineto{\pgfqpoint{1.142005in}{3.284073in}}%
\pgfpathclose%
\pgfusepath{stroke,fill}%
\end{pgfscope}%
\begin{pgfscope}%
\pgfpathrectangle{\pgfqpoint{0.050000in}{0.050000in}}{\pgfqpoint{4.900000in}{4.900000in}}%
\pgfusepath{clip}%
\pgfsetbuttcap%
\pgfsetroundjoin%
\definecolor{currentfill}{rgb}{0.586082,0.586082,0.586082}%
\pgfsetfillcolor{currentfill}%
\pgfsetfillopacity{0.900000}%
\pgfsetlinewidth{0.501875pt}%
\definecolor{currentstroke}{rgb}{0.200000,0.200000,0.200000}%
\pgfsetstrokecolor{currentstroke}%
\pgfsetdash{}{0pt}%
\pgfpathmoveto{\pgfqpoint{1.857479in}{2.705796in}}%
\pgfpathlineto{\pgfqpoint{1.904451in}{2.660866in}}%
\pgfpathlineto{\pgfqpoint{1.935319in}{2.640980in}}%
\pgfpathlineto{\pgfqpoint{1.888340in}{2.685937in}}%
\pgfpathlineto{\pgfqpoint{1.857479in}{2.705796in}}%
\pgfpathclose%
\pgfusepath{stroke,fill}%
\end{pgfscope}%
\begin{pgfscope}%
\pgfpathrectangle{\pgfqpoint{0.050000in}{0.050000in}}{\pgfqpoint{4.900000in}{4.900000in}}%
\pgfusepath{clip}%
\pgfsetbuttcap%
\pgfsetroundjoin%
\definecolor{currentfill}{rgb}{0.836340,0.836340,0.836340}%
\pgfsetfillcolor{currentfill}%
\pgfsetfillopacity{0.900000}%
\pgfsetlinewidth{0.501875pt}%
\definecolor{currentstroke}{rgb}{0.200000,0.200000,0.200000}%
\pgfsetstrokecolor{currentstroke}%
\pgfsetdash{}{0pt}%
\pgfpathmoveto{\pgfqpoint{2.573048in}{2.129300in}}%
\pgfpathlineto{\pgfqpoint{2.618514in}{2.087148in}}%
\pgfpathlineto{\pgfqpoint{2.650491in}{2.066772in}}%
\pgfpathlineto{\pgfqpoint{2.605020in}{2.108801in}}%
\pgfpathlineto{\pgfqpoint{2.573048in}{2.129300in}}%
\pgfpathclose%
\pgfusepath{stroke,fill}%
\end{pgfscope}%
\begin{pgfscope}%
\pgfpathrectangle{\pgfqpoint{0.050000in}{0.050000in}}{\pgfqpoint{4.900000in}{4.900000in}}%
\pgfusepath{clip}%
\pgfsetbuttcap%
\pgfsetroundjoin%
\definecolor{currentfill}{rgb}{0.487043,0.487043,0.487043}%
\pgfsetfillcolor{currentfill}%
\pgfsetfillopacity{0.900000}%
\pgfsetlinewidth{0.501875pt}%
\definecolor{currentstroke}{rgb}{0.200000,0.200000,0.200000}%
\pgfsetstrokecolor{currentstroke}%
\pgfsetdash{}{0pt}%
\pgfpathmoveto{\pgfqpoint{1.592966in}{2.920220in}}%
\pgfpathlineto{\pgfqpoint{1.640504in}{2.874829in}}%
\pgfpathlineto{\pgfqpoint{1.670970in}{2.855270in}}%
\pgfpathlineto{\pgfqpoint{1.623423in}{2.900689in}}%
\pgfpathlineto{\pgfqpoint{1.592966in}{2.920220in}}%
\pgfpathclose%
\pgfusepath{stroke,fill}%
\end{pgfscope}%
\begin{pgfscope}%
\pgfpathrectangle{\pgfqpoint{0.050000in}{0.050000in}}{\pgfqpoint{4.900000in}{4.900000in}}%
\pgfusepath{clip}%
\pgfsetbuttcap%
\pgfsetroundjoin%
\definecolor{currentfill}{rgb}{0.998155,0.998155,0.998155}%
\pgfsetfillcolor{currentfill}%
\pgfsetfillopacity{0.900000}%
\pgfsetlinewidth{0.501875pt}%
\definecolor{currentstroke}{rgb}{0.200000,0.200000,0.200000}%
\pgfsetstrokecolor{currentstroke}%
\pgfsetdash{}{0pt}%
\pgfpathmoveto{\pgfqpoint{4.122166in}{1.459842in}}%
\pgfpathlineto{\pgfqpoint{4.165877in}{1.471040in}}%
\pgfpathlineto{\pgfqpoint{4.200293in}{1.450396in}}%
\pgfpathlineto{\pgfqpoint{4.156557in}{1.439269in}}%
\pgfpathlineto{\pgfqpoint{4.122166in}{1.459842in}}%
\pgfpathclose%
\pgfusepath{stroke,fill}%
\end{pgfscope}%
\begin{pgfscope}%
\pgfpathrectangle{\pgfqpoint{0.050000in}{0.050000in}}{\pgfqpoint{4.900000in}{4.900000in}}%
\pgfusepath{clip}%
\pgfsetbuttcap%
\pgfsetroundjoin%
\definecolor{currentfill}{rgb}{0.898378,0.898378,0.898378}%
\pgfsetfillcolor{currentfill}%
\pgfsetfillopacity{0.900000}%
\pgfsetlinewidth{0.501875pt}%
\definecolor{currentstroke}{rgb}{0.200000,0.200000,0.200000}%
\pgfsetstrokecolor{currentstroke}%
\pgfsetdash{}{0pt}%
\pgfpathmoveto{\pgfqpoint{2.837434in}{1.926653in}}%
\pgfpathlineto{\pgfqpoint{2.882382in}{1.890487in}}%
\pgfpathlineto{\pgfqpoint{2.914778in}{1.870565in}}%
\pgfpathlineto{\pgfqpoint{2.869822in}{1.906486in}}%
\pgfpathlineto{\pgfqpoint{2.837434in}{1.926653in}}%
\pgfpathclose%
\pgfusepath{stroke,fill}%
\end{pgfscope}%
\begin{pgfscope}%
\pgfpathrectangle{\pgfqpoint{0.050000in}{0.050000in}}{\pgfqpoint{4.900000in}{4.900000in}}%
\pgfusepath{clip}%
\pgfsetbuttcap%
\pgfsetroundjoin%
\definecolor{currentfill}{rgb}{0.994464,0.994464,0.994464}%
\pgfsetfillcolor{currentfill}%
\pgfsetfillopacity{0.900000}%
\pgfsetlinewidth{0.501875pt}%
\definecolor{currentstroke}{rgb}{0.200000,0.200000,0.200000}%
\pgfsetstrokecolor{currentstroke}%
\pgfsetdash{}{0pt}%
\pgfpathmoveto{\pgfqpoint{3.844343in}{1.485764in}}%
\pgfpathlineto{\pgfqpoint{3.888273in}{1.492540in}}%
\pgfpathlineto{\pgfqpoint{3.922267in}{1.472407in}}%
\pgfpathlineto{\pgfqpoint{3.878312in}{1.465718in}}%
\pgfpathlineto{\pgfqpoint{3.844343in}{1.485764in}}%
\pgfpathclose%
\pgfusepath{stroke,fill}%
\end{pgfscope}%
\begin{pgfscope}%
\pgfpathrectangle{\pgfqpoint{0.050000in}{0.050000in}}{\pgfqpoint{4.900000in}{4.900000in}}%
\pgfusepath{clip}%
\pgfsetbuttcap%
\pgfsetroundjoin%
\definecolor{currentfill}{rgb}{0.760554,0.760554,0.760554}%
\pgfsetfillcolor{currentfill}%
\pgfsetfillopacity{0.900000}%
\pgfsetlinewidth{0.501875pt}%
\definecolor{currentstroke}{rgb}{0.200000,0.200000,0.200000}%
\pgfsetstrokecolor{currentstroke}%
\pgfsetdash{}{0pt}%
\pgfpathmoveto{\pgfqpoint{2.308637in}{2.341875in}}%
\pgfpathlineto{\pgfqpoint{2.354666in}{2.297823in}}%
\pgfpathlineto{\pgfqpoint{2.386238in}{2.277424in}}%
\pgfpathlineto{\pgfqpoint{2.340204in}{2.321474in}}%
\pgfpathlineto{\pgfqpoint{2.308637in}{2.341875in}}%
\pgfpathclose%
\pgfusepath{stroke,fill}%
\end{pgfscope}%
\begin{pgfscope}%
\pgfpathrectangle{\pgfqpoint{0.050000in}{0.050000in}}{\pgfqpoint{4.900000in}{4.900000in}}%
\pgfusepath{clip}%
\pgfsetbuttcap%
\pgfsetroundjoin%
\definecolor{currentfill}{rgb}{0.395663,0.395663,0.395663}%
\pgfsetfillcolor{currentfill}%
\pgfsetfillopacity{0.900000}%
\pgfsetlinewidth{0.501875pt}%
\definecolor{currentstroke}{rgb}{0.200000,0.200000,0.200000}%
\pgfsetstrokecolor{currentstroke}%
\pgfsetdash{}{0pt}%
\pgfpathmoveto{\pgfqpoint{1.328363in}{3.134718in}}%
\pgfpathlineto{\pgfqpoint{1.376469in}{3.088866in}}%
\pgfpathlineto{\pgfqpoint{1.406532in}{3.069634in}}%
\pgfpathlineto{\pgfqpoint{1.358416in}{3.115516in}}%
\pgfpathlineto{\pgfqpoint{1.328363in}{3.134718in}}%
\pgfpathclose%
\pgfusepath{stroke,fill}%
\end{pgfscope}%
\begin{pgfscope}%
\pgfpathrectangle{\pgfqpoint{0.050000in}{0.050000in}}{\pgfqpoint{4.900000in}{4.900000in}}%
\pgfusepath{clip}%
\pgfsetbuttcap%
\pgfsetroundjoin%
\definecolor{currentfill}{rgb}{0.657809,0.657809,0.657809}%
\pgfsetfillcolor{currentfill}%
\pgfsetfillopacity{0.900000}%
\pgfsetlinewidth{0.501875pt}%
\definecolor{currentstroke}{rgb}{0.200000,0.200000,0.200000}%
\pgfsetstrokecolor{currentstroke}%
\pgfsetdash{}{0pt}%
\pgfpathmoveto{\pgfqpoint{2.044138in}{2.556203in}}%
\pgfpathlineto{\pgfqpoint{2.090734in}{2.511572in}}%
\pgfpathlineto{\pgfqpoint{2.121903in}{2.491450in}}%
\pgfpathlineto{\pgfqpoint{2.075301in}{2.536106in}}%
\pgfpathlineto{\pgfqpoint{2.044138in}{2.556203in}}%
\pgfpathclose%
\pgfusepath{stroke,fill}%
\end{pgfscope}%
\begin{pgfscope}%
\pgfpathrectangle{\pgfqpoint{0.050000in}{0.050000in}}{\pgfqpoint{4.900000in}{4.900000in}}%
\pgfusepath{clip}%
\pgfsetbuttcap%
\pgfsetroundjoin%
\definecolor{currentfill}{rgb}{0.968627,0.968627,0.968627}%
\pgfsetfillcolor{currentfill}%
\pgfsetfillopacity{0.900000}%
\pgfsetlinewidth{0.501875pt}%
\definecolor{currentstroke}{rgb}{0.200000,0.200000,0.200000}%
\pgfsetstrokecolor{currentstroke}%
\pgfsetdash{}{0pt}%
\pgfpathmoveto{\pgfqpoint{3.334317in}{1.626606in}}%
\pgfpathlineto{\pgfqpoint{3.378604in}{1.613325in}}%
\pgfpathlineto{\pgfqpoint{3.411806in}{1.593984in}}%
\pgfpathlineto{\pgfqpoint{3.367501in}{1.607189in}}%
\pgfpathlineto{\pgfqpoint{3.334317in}{1.626606in}}%
\pgfpathclose%
\pgfusepath{stroke,fill}%
\end{pgfscope}%
\begin{pgfscope}%
\pgfpathrectangle{\pgfqpoint{0.050000in}{0.050000in}}{\pgfqpoint{4.900000in}{4.900000in}}%
\pgfusepath{clip}%
\pgfsetbuttcap%
\pgfsetroundjoin%
\definecolor{currentfill}{rgb}{0.974164,0.974164,0.974164}%
\pgfsetfillcolor{currentfill}%
\pgfsetfillopacity{0.900000}%
\pgfsetlinewidth{0.501875pt}%
\definecolor{currentstroke}{rgb}{0.200000,0.200000,0.200000}%
\pgfsetstrokecolor{currentstroke}%
\pgfsetdash{}{0pt}%
\pgfpathmoveto{\pgfqpoint{3.411806in}{1.593984in}}%
\pgfpathlineto{\pgfqpoint{3.456031in}{1.584545in}}%
\pgfpathlineto{\pgfqpoint{3.489357in}{1.565150in}}%
\pgfpathlineto{\pgfqpoint{3.445113in}{1.574558in}}%
\pgfpathlineto{\pgfqpoint{3.411806in}{1.593984in}}%
\pgfpathclose%
\pgfusepath{stroke,fill}%
\end{pgfscope}%
\begin{pgfscope}%
\pgfpathrectangle{\pgfqpoint{0.050000in}{0.050000in}}{\pgfqpoint{4.900000in}{4.900000in}}%
\pgfusepath{clip}%
\pgfsetbuttcap%
\pgfsetroundjoin%
\definecolor{currentfill}{rgb}{0.961246,0.961246,0.961246}%
\pgfsetfillcolor{currentfill}%
\pgfsetfillopacity{0.900000}%
\pgfsetlinewidth{0.501875pt}%
\definecolor{currentstroke}{rgb}{0.200000,0.200000,0.200000}%
\pgfsetstrokecolor{currentstroke}%
\pgfsetdash{}{0pt}%
\pgfpathmoveto{\pgfqpoint{3.256877in}{1.663252in}}%
\pgfpathlineto{\pgfqpoint{3.301239in}{1.645939in}}%
\pgfpathlineto{\pgfqpoint{3.334317in}{1.626606in}}%
\pgfpathlineto{\pgfqpoint{3.289939in}{1.643795in}}%
\pgfpathlineto{\pgfqpoint{3.256877in}{1.663252in}}%
\pgfpathclose%
\pgfusepath{stroke,fill}%
\end{pgfscope}%
\begin{pgfscope}%
\pgfpathrectangle{\pgfqpoint{0.050000in}{0.050000in}}{\pgfqpoint{4.900000in}{4.900000in}}%
\pgfusepath{clip}%
\pgfsetbuttcap%
\pgfsetroundjoin%
\definecolor{currentfill}{rgb}{0.295271,0.295271,0.295271}%
\pgfsetfillcolor{currentfill}%
\pgfsetfillopacity{0.900000}%
\pgfsetlinewidth{0.501875pt}%
\definecolor{currentstroke}{rgb}{0.200000,0.200000,0.200000}%
\pgfsetstrokecolor{currentstroke}%
\pgfsetdash{}{0pt}%
\pgfpathmoveto{\pgfqpoint{1.063672in}{3.349291in}}%
\pgfpathlineto{\pgfqpoint{1.112345in}{3.302978in}}%
\pgfpathlineto{\pgfqpoint{1.142005in}{3.284073in}}%
\pgfpathlineto{\pgfqpoint{1.093320in}{3.330417in}}%
\pgfpathlineto{\pgfqpoint{1.063672in}{3.349291in}}%
\pgfpathclose%
\pgfusepath{stroke,fill}%
\end{pgfscope}%
\begin{pgfscope}%
\pgfpathrectangle{\pgfqpoint{0.050000in}{0.050000in}}{\pgfqpoint{4.900000in}{4.900000in}}%
\pgfusepath{clip}%
\pgfsetbuttcap%
\pgfsetroundjoin%
\definecolor{currentfill}{rgb}{0.551634,0.551634,0.551634}%
\pgfsetfillcolor{currentfill}%
\pgfsetfillopacity{0.900000}%
\pgfsetlinewidth{0.501875pt}%
\definecolor{currentstroke}{rgb}{0.200000,0.200000,0.200000}%
\pgfsetstrokecolor{currentstroke}%
\pgfsetdash{}{0pt}%
\pgfpathmoveto{\pgfqpoint{1.779549in}{2.770686in}}%
\pgfpathlineto{\pgfqpoint{1.826713in}{2.725593in}}%
\pgfpathlineto{\pgfqpoint{1.857479in}{2.705796in}}%
\pgfpathlineto{\pgfqpoint{1.810308in}{2.750917in}}%
\pgfpathlineto{\pgfqpoint{1.779549in}{2.770686in}}%
\pgfpathclose%
\pgfusepath{stroke,fill}%
\end{pgfscope}%
\begin{pgfscope}%
\pgfpathrectangle{\pgfqpoint{0.050000in}{0.050000in}}{\pgfqpoint{4.900000in}{4.900000in}}%
\pgfusepath{clip}%
\pgfsetbuttcap%
\pgfsetroundjoin%
\definecolor{currentfill}{rgb}{0.977855,0.977855,0.977855}%
\pgfsetfillcolor{currentfill}%
\pgfsetfillopacity{0.900000}%
\pgfsetlinewidth{0.501875pt}%
\definecolor{currentstroke}{rgb}{0.200000,0.200000,0.200000}%
\pgfsetstrokecolor{currentstroke}%
\pgfsetdash{}{0pt}%
\pgfpathmoveto{\pgfqpoint{3.489357in}{1.565150in}}%
\pgfpathlineto{\pgfqpoint{3.533527in}{1.559276in}}%
\pgfpathlineto{\pgfqpoint{3.566980in}{1.539790in}}%
\pgfpathlineto{\pgfqpoint{3.522789in}{1.545672in}}%
\pgfpathlineto{\pgfqpoint{3.489357in}{1.565150in}}%
\pgfpathclose%
\pgfusepath{stroke,fill}%
\end{pgfscope}%
\begin{pgfscope}%
\pgfpathrectangle{\pgfqpoint{0.050000in}{0.050000in}}{\pgfqpoint{4.900000in}{4.900000in}}%
\pgfusepath{clip}%
\pgfsetbuttcap%
\pgfsetroundjoin%
\definecolor{currentfill}{rgb}{0.994464,0.994464,0.994464}%
\pgfsetfillcolor{currentfill}%
\pgfsetfillopacity{0.900000}%
\pgfsetlinewidth{0.501875pt}%
\definecolor{currentstroke}{rgb}{0.200000,0.200000,0.200000}%
\pgfsetstrokecolor{currentstroke}%
\pgfsetdash{}{0pt}%
\pgfpathmoveto{\pgfqpoint{3.922267in}{1.472407in}}%
\pgfpathlineto{\pgfqpoint{3.966153in}{1.480674in}}%
\pgfpathlineto{\pgfqpoint{4.000278in}{1.460389in}}%
\pgfpathlineto{\pgfqpoint{3.956367in}{1.452208in}}%
\pgfpathlineto{\pgfqpoint{3.922267in}{1.472407in}}%
\pgfpathclose%
\pgfusepath{stroke,fill}%
\end{pgfscope}%
\begin{pgfscope}%
\pgfpathrectangle{\pgfqpoint{0.050000in}{0.050000in}}{\pgfqpoint{4.900000in}{4.900000in}}%
\pgfusepath{clip}%
\pgfsetbuttcap%
\pgfsetroundjoin%
\definecolor{currentfill}{rgb}{1.000000,1.000000,1.000000}%
\pgfsetfillcolor{currentfill}%
\pgfsetfillopacity{0.900000}%
\pgfsetlinewidth{0.501875pt}%
\definecolor{currentstroke}{rgb}{0.200000,0.200000,0.200000}%
\pgfsetstrokecolor{currentstroke}%
\pgfsetdash{}{0pt}%
\pgfpathmoveto{\pgfqpoint{4.200293in}{1.450396in}}%
\pgfpathlineto{\pgfqpoint{4.243946in}{1.462094in}}%
\pgfpathlineto{\pgfqpoint{4.278496in}{1.441318in}}%
\pgfpathlineto{\pgfqpoint{4.234817in}{1.429687in}}%
\pgfpathlineto{\pgfqpoint{4.200293in}{1.450396in}}%
\pgfpathclose%
\pgfusepath{stroke,fill}%
\end{pgfscope}%
\begin{pgfscope}%
\pgfpathrectangle{\pgfqpoint{0.050000in}{0.050000in}}{\pgfqpoint{4.900000in}{4.900000in}}%
\pgfusepath{clip}%
\pgfsetbuttcap%
\pgfsetroundjoin%
\definecolor{currentfill}{rgb}{0.953864,0.953864,0.953864}%
\pgfsetfillcolor{currentfill}%
\pgfsetfillopacity{0.900000}%
\pgfsetlinewidth{0.501875pt}%
\definecolor{currentstroke}{rgb}{0.200000,0.200000,0.200000}%
\pgfsetstrokecolor{currentstroke}%
\pgfsetdash{}{0pt}%
\pgfpathmoveto{\pgfqpoint{3.179471in}{1.704051in}}%
\pgfpathlineto{\pgfqpoint{3.223920in}{1.682627in}}%
\pgfpathlineto{\pgfqpoint{3.256877in}{1.663252in}}%
\pgfpathlineto{\pgfqpoint{3.212413in}{1.684506in}}%
\pgfpathlineto{\pgfqpoint{3.179471in}{1.704051in}}%
\pgfpathclose%
\pgfusepath{stroke,fill}%
\end{pgfscope}%
\begin{pgfscope}%
\pgfpathrectangle{\pgfqpoint{0.050000in}{0.050000in}}{\pgfqpoint{4.900000in}{4.900000in}}%
\pgfusepath{clip}%
\pgfsetbuttcap%
\pgfsetroundjoin%
\definecolor{currentfill}{rgb}{0.878570,0.878570,0.878570}%
\pgfsetfillcolor{currentfill}%
\pgfsetfillopacity{0.900000}%
\pgfsetlinewidth{0.501875pt}%
\definecolor{currentstroke}{rgb}{0.200000,0.200000,0.200000}%
\pgfsetstrokecolor{currentstroke}%
\pgfsetdash{}{0pt}%
\pgfpathmoveto{\pgfqpoint{2.760035in}{1.985446in}}%
\pgfpathlineto{\pgfqpoint{2.805148in}{1.946772in}}%
\pgfpathlineto{\pgfqpoint{2.837434in}{1.926653in}}%
\pgfpathlineto{\pgfqpoint{2.792316in}{1.965110in}}%
\pgfpathlineto{\pgfqpoint{2.760035in}{1.985446in}}%
\pgfpathclose%
\pgfusepath{stroke,fill}%
\end{pgfscope}%
\begin{pgfscope}%
\pgfpathrectangle{\pgfqpoint{0.050000in}{0.050000in}}{\pgfqpoint{4.900000in}{4.900000in}}%
\pgfusepath{clip}%
\pgfsetbuttcap%
\pgfsetroundjoin%
\definecolor{currentfill}{rgb}{0.812226,0.812226,0.812226}%
\pgfsetfillcolor{currentfill}%
\pgfsetfillopacity{0.900000}%
\pgfsetlinewidth{0.501875pt}%
\definecolor{currentstroke}{rgb}{0.200000,0.200000,0.200000}%
\pgfsetstrokecolor{currentstroke}%
\pgfsetdash{}{0pt}%
\pgfpathmoveto{\pgfqpoint{2.495521in}{2.192826in}}%
\pgfpathlineto{\pgfqpoint{2.541174in}{2.149752in}}%
\pgfpathlineto{\pgfqpoint{2.573048in}{2.129300in}}%
\pgfpathlineto{\pgfqpoint{2.527391in}{2.172299in}}%
\pgfpathlineto{\pgfqpoint{2.495521in}{2.192826in}}%
\pgfpathclose%
\pgfusepath{stroke,fill}%
\end{pgfscope}%
\begin{pgfscope}%
\pgfpathrectangle{\pgfqpoint{0.050000in}{0.050000in}}{\pgfqpoint{4.900000in}{4.900000in}}%
\pgfusepath{clip}%
\pgfsetbuttcap%
\pgfsetroundjoin%
\definecolor{currentfill}{rgb}{0.983391,0.983391,0.983391}%
\pgfsetfillcolor{currentfill}%
\pgfsetfillopacity{0.900000}%
\pgfsetlinewidth{0.501875pt}%
\definecolor{currentstroke}{rgb}{0.200000,0.200000,0.200000}%
\pgfsetstrokecolor{currentstroke}%
\pgfsetdash{}{0pt}%
\pgfpathmoveto{\pgfqpoint{3.566980in}{1.539790in}}%
\pgfpathlineto{\pgfqpoint{3.611103in}{1.537150in}}%
\pgfpathlineto{\pgfqpoint{3.644683in}{1.517545in}}%
\pgfpathlineto{\pgfqpoint{3.600539in}{1.520223in}}%
\pgfpathlineto{\pgfqpoint{3.566980in}{1.539790in}}%
\pgfpathclose%
\pgfusepath{stroke,fill}%
\end{pgfscope}%
\begin{pgfscope}%
\pgfpathrectangle{\pgfqpoint{0.050000in}{0.050000in}}{\pgfqpoint{4.900000in}{4.900000in}}%
\pgfusepath{clip}%
\pgfsetbuttcap%
\pgfsetroundjoin%
\definecolor{currentfill}{rgb}{0.456901,0.456901,0.456901}%
\pgfsetfillcolor{currentfill}%
\pgfsetfillopacity{0.900000}%
\pgfsetlinewidth{0.501875pt}%
\definecolor{currentstroke}{rgb}{0.200000,0.200000,0.200000}%
\pgfsetstrokecolor{currentstroke}%
\pgfsetdash{}{0pt}%
\pgfpathmoveto{\pgfqpoint{1.514872in}{2.985244in}}%
\pgfpathlineto{\pgfqpoint{1.562603in}{2.939690in}}%
\pgfpathlineto{\pgfqpoint{1.592966in}{2.920220in}}%
\pgfpathlineto{\pgfqpoint{1.545225in}{2.965803in}}%
\pgfpathlineto{\pgfqpoint{1.514872in}{2.985244in}}%
\pgfpathclose%
\pgfusepath{stroke,fill}%
\end{pgfscope}%
\begin{pgfscope}%
\pgfpathrectangle{\pgfqpoint{0.050000in}{0.050000in}}{\pgfqpoint{4.900000in}{4.900000in}}%
\pgfusepath{clip}%
\pgfsetbuttcap%
\pgfsetroundjoin%
\definecolor{currentfill}{rgb}{0.944637,0.944637,0.944637}%
\pgfsetfillcolor{currentfill}%
\pgfsetfillopacity{0.900000}%
\pgfsetlinewidth{0.501875pt}%
\definecolor{currentstroke}{rgb}{0.200000,0.200000,0.200000}%
\pgfsetstrokecolor{currentstroke}%
\pgfsetdash{}{0pt}%
\pgfpathmoveto{\pgfqpoint{3.102081in}{1.749001in}}%
\pgfpathlineto{\pgfqpoint{3.146632in}{1.723520in}}%
\pgfpathlineto{\pgfqpoint{3.179471in}{1.704051in}}%
\pgfpathlineto{\pgfqpoint{3.134906in}{1.729323in}}%
\pgfpathlineto{\pgfqpoint{3.102081in}{1.749001in}}%
\pgfpathclose%
\pgfusepath{stroke,fill}%
\end{pgfscope}%
\begin{pgfscope}%
\pgfpathrectangle{\pgfqpoint{0.050000in}{0.050000in}}{\pgfqpoint{4.900000in}{4.900000in}}%
\pgfusepath{clip}%
\pgfsetbuttcap%
\pgfsetroundjoin%
\definecolor{currentfill}{rgb}{0.729781,0.729781,0.729781}%
\pgfsetfillcolor{currentfill}%
\pgfsetfillopacity{0.900000}%
\pgfsetlinewidth{0.501875pt}%
\definecolor{currentstroke}{rgb}{0.200000,0.200000,0.200000}%
\pgfsetstrokecolor{currentstroke}%
\pgfsetdash{}{0pt}%
\pgfpathmoveto{\pgfqpoint{2.230947in}{2.406512in}}%
\pgfpathlineto{\pgfqpoint{2.277167in}{2.362218in}}%
\pgfpathlineto{\pgfqpoint{2.308637in}{2.341875in}}%
\pgfpathlineto{\pgfqpoint{2.262412in}{2.386183in}}%
\pgfpathlineto{\pgfqpoint{2.230947in}{2.406512in}}%
\pgfpathclose%
\pgfusepath{stroke,fill}%
\end{pgfscope}%
\begin{pgfscope}%
\pgfpathrectangle{\pgfqpoint{0.050000in}{0.050000in}}{\pgfqpoint{4.900000in}{4.900000in}}%
\pgfusepath{clip}%
\pgfsetbuttcap%
\pgfsetroundjoin%
\definecolor{currentfill}{rgb}{0.363183,0.363183,0.363183}%
\pgfsetfillcolor{currentfill}%
\pgfsetfillopacity{0.900000}%
\pgfsetlinewidth{0.501875pt}%
\definecolor{currentstroke}{rgb}{0.200000,0.200000,0.200000}%
\pgfsetstrokecolor{currentstroke}%
\pgfsetdash{}{0pt}%
\pgfpathmoveto{\pgfqpoint{1.250105in}{3.199876in}}%
\pgfpathlineto{\pgfqpoint{1.298404in}{3.153861in}}%
\pgfpathlineto{\pgfqpoint{1.328363in}{3.134718in}}%
\pgfpathlineto{\pgfqpoint{1.280054in}{3.180764in}}%
\pgfpathlineto{\pgfqpoint{1.250105in}{3.199876in}}%
\pgfpathclose%
\pgfusepath{stroke,fill}%
\end{pgfscope}%
\begin{pgfscope}%
\pgfpathrectangle{\pgfqpoint{0.050000in}{0.050000in}}{\pgfqpoint{4.900000in}{4.900000in}}%
\pgfusepath{clip}%
\pgfsetbuttcap%
\pgfsetroundjoin%
\definecolor{currentfill}{rgb}{0.987082,0.987082,0.987082}%
\pgfsetfillcolor{currentfill}%
\pgfsetfillopacity{0.900000}%
\pgfsetlinewidth{0.501875pt}%
\definecolor{currentstroke}{rgb}{0.200000,0.200000,0.200000}%
\pgfsetstrokecolor{currentstroke}%
\pgfsetdash{}{0pt}%
\pgfpathmoveto{\pgfqpoint{3.644683in}{1.517545in}}%
\pgfpathlineto{\pgfqpoint{3.688762in}{1.517778in}}%
\pgfpathlineto{\pgfqpoint{3.722471in}{1.498035in}}%
\pgfpathlineto{\pgfqpoint{3.678369in}{1.497862in}}%
\pgfpathlineto{\pgfqpoint{3.644683in}{1.517545in}}%
\pgfpathclose%
\pgfusepath{stroke,fill}%
\end{pgfscope}%
\begin{pgfscope}%
\pgfpathrectangle{\pgfqpoint{0.050000in}{0.050000in}}{\pgfqpoint{4.900000in}{4.900000in}}%
\pgfusepath{clip}%
\pgfsetbuttcap%
\pgfsetroundjoin%
\definecolor{currentfill}{rgb}{0.996309,0.996309,0.996309}%
\pgfsetfillcolor{currentfill}%
\pgfsetfillopacity{0.900000}%
\pgfsetlinewidth{0.501875pt}%
\definecolor{currentstroke}{rgb}{0.200000,0.200000,0.200000}%
\pgfsetstrokecolor{currentstroke}%
\pgfsetdash{}{0pt}%
\pgfpathmoveto{\pgfqpoint{4.000278in}{1.460389in}}%
\pgfpathlineto{\pgfqpoint{4.044118in}{1.469859in}}%
\pgfpathlineto{\pgfqpoint{4.078376in}{1.449428in}}%
\pgfpathlineto{\pgfqpoint{4.034511in}{1.440039in}}%
\pgfpathlineto{\pgfqpoint{4.000278in}{1.460389in}}%
\pgfpathclose%
\pgfusepath{stroke,fill}%
\end{pgfscope}%
\begin{pgfscope}%
\pgfpathrectangle{\pgfqpoint{0.050000in}{0.050000in}}{\pgfqpoint{4.900000in}{4.900000in}}%
\pgfusepath{clip}%
\pgfsetbuttcap%
\pgfsetroundjoin%
\definecolor{currentfill}{rgb}{0.624221,0.624221,0.624221}%
\pgfsetfillcolor{currentfill}%
\pgfsetfillopacity{0.900000}%
\pgfsetlinewidth{0.501875pt}%
\definecolor{currentstroke}{rgb}{0.200000,0.200000,0.200000}%
\pgfsetstrokecolor{currentstroke}%
\pgfsetdash{}{0pt}%
\pgfpathmoveto{\pgfqpoint{1.966284in}{2.621033in}}%
\pgfpathlineto{\pgfqpoint{2.013071in}{2.576238in}}%
\pgfpathlineto{\pgfqpoint{2.044138in}{2.556203in}}%
\pgfpathlineto{\pgfqpoint{1.997344in}{2.601024in}}%
\pgfpathlineto{\pgfqpoint{1.966284in}{2.621033in}}%
\pgfpathclose%
\pgfusepath{stroke,fill}%
\end{pgfscope}%
\begin{pgfscope}%
\pgfpathrectangle{\pgfqpoint{0.050000in}{0.050000in}}{\pgfqpoint{4.900000in}{4.900000in}}%
\pgfusepath{clip}%
\pgfsetbuttcap%
\pgfsetroundjoin%
\definecolor{currentfill}{rgb}{1.000000,1.000000,1.000000}%
\pgfsetfillcolor{currentfill}%
\pgfsetfillopacity{0.900000}%
\pgfsetlinewidth{0.501875pt}%
\definecolor{currentstroke}{rgb}{0.200000,0.200000,0.200000}%
\pgfsetstrokecolor{currentstroke}%
\pgfsetdash{}{0pt}%
\pgfpathmoveto{\pgfqpoint{4.278496in}{1.441318in}}%
\pgfpathlineto{\pgfqpoint{4.322087in}{1.453328in}}%
\pgfpathlineto{\pgfqpoint{4.356770in}{1.432424in}}%
\pgfpathlineto{\pgfqpoint{4.313154in}{1.420478in}}%
\pgfpathlineto{\pgfqpoint{4.278496in}{1.441318in}}%
\pgfpathclose%
\pgfusepath{stroke,fill}%
\end{pgfscope}%
\begin{pgfscope}%
\pgfpathrectangle{\pgfqpoint{0.050000in}{0.050000in}}{\pgfqpoint{4.900000in}{4.900000in}}%
\pgfusepath{clip}%
\pgfsetbuttcap%
\pgfsetroundjoin%
\definecolor{currentfill}{rgb}{0.932334,0.932334,0.932334}%
\pgfsetfillcolor{currentfill}%
\pgfsetfillopacity{0.900000}%
\pgfsetlinewidth{0.501875pt}%
\definecolor{currentstroke}{rgb}{0.200000,0.200000,0.200000}%
\pgfsetstrokecolor{currentstroke}%
\pgfsetdash{}{0pt}%
\pgfpathmoveto{\pgfqpoint{3.024688in}{1.797948in}}%
\pgfpathlineto{\pgfqpoint{3.069358in}{1.768610in}}%
\pgfpathlineto{\pgfqpoint{3.102081in}{1.749001in}}%
\pgfpathlineto{\pgfqpoint{3.057399in}{1.778102in}}%
\pgfpathlineto{\pgfqpoint{3.024688in}{1.797948in}}%
\pgfpathclose%
\pgfusepath{stroke,fill}%
\end{pgfscope}%
\begin{pgfscope}%
\pgfpathrectangle{\pgfqpoint{0.050000in}{0.050000in}}{\pgfqpoint{4.900000in}{4.900000in}}%
\pgfusepath{clip}%
\pgfsetbuttcap%
\pgfsetroundjoin%
\definecolor{currentfill}{rgb}{0.250980,0.250980,0.250980}%
\pgfsetfillcolor{currentfill}%
\pgfsetfillopacity{0.900000}%
\pgfsetlinewidth{0.501875pt}%
\definecolor{currentstroke}{rgb}{0.200000,0.200000,0.200000}%
\pgfsetstrokecolor{currentstroke}%
\pgfsetdash{}{0pt}%
\pgfpathmoveto{\pgfqpoint{0.985249in}{3.414584in}}%
\pgfpathlineto{\pgfqpoint{1.034116in}{3.368106in}}%
\pgfpathlineto{\pgfqpoint{1.063672in}{3.349291in}}%
\pgfpathlineto{\pgfqpoint{1.014792in}{3.395800in}}%
\pgfpathlineto{\pgfqpoint{0.985249in}{3.414584in}}%
\pgfpathclose%
\pgfusepath{stroke,fill}%
\end{pgfscope}%
\begin{pgfscope}%
\pgfpathrectangle{\pgfqpoint{0.050000in}{0.050000in}}{\pgfqpoint{4.900000in}{4.900000in}}%
\pgfusepath{clip}%
\pgfsetbuttcap%
\pgfsetroundjoin%
\definecolor{currentfill}{rgb}{0.521492,0.521492,0.521492}%
\pgfsetfillcolor{currentfill}%
\pgfsetfillopacity{0.900000}%
\pgfsetlinewidth{0.501875pt}%
\definecolor{currentstroke}{rgb}{0.200000,0.200000,0.200000}%
\pgfsetstrokecolor{currentstroke}%
\pgfsetdash{}{0pt}%
\pgfpathmoveto{\pgfqpoint{1.701531in}{2.835650in}}%
\pgfpathlineto{\pgfqpoint{1.748886in}{2.790394in}}%
\pgfpathlineto{\pgfqpoint{1.779549in}{2.770686in}}%
\pgfpathlineto{\pgfqpoint{1.732186in}{2.815970in}}%
\pgfpathlineto{\pgfqpoint{1.701531in}{2.835650in}}%
\pgfpathclose%
\pgfusepath{stroke,fill}%
\end{pgfscope}%
\begin{pgfscope}%
\pgfpathrectangle{\pgfqpoint{0.050000in}{0.050000in}}{\pgfqpoint{4.900000in}{4.900000in}}%
\pgfusepath{clip}%
\pgfsetbuttcap%
\pgfsetroundjoin%
\definecolor{currentfill}{rgb}{0.858762,0.858762,0.858762}%
\pgfsetfillcolor{currentfill}%
\pgfsetfillopacity{0.900000}%
\pgfsetlinewidth{0.501875pt}%
\definecolor{currentstroke}{rgb}{0.200000,0.200000,0.200000}%
\pgfsetstrokecolor{currentstroke}%
\pgfsetdash{}{0pt}%
\pgfpathmoveto{\pgfqpoint{2.682566in}{2.046351in}}%
\pgfpathlineto{\pgfqpoint{2.727854in}{2.005736in}}%
\pgfpathlineto{\pgfqpoint{2.760035in}{1.985446in}}%
\pgfpathlineto{\pgfqpoint{2.714742in}{2.025886in}}%
\pgfpathlineto{\pgfqpoint{2.682566in}{2.046351in}}%
\pgfpathclose%
\pgfusepath{stroke,fill}%
\end{pgfscope}%
\begin{pgfscope}%
\pgfpathrectangle{\pgfqpoint{0.050000in}{0.050000in}}{\pgfqpoint{4.900000in}{4.900000in}}%
\pgfusepath{clip}%
\pgfsetbuttcap%
\pgfsetroundjoin%
\definecolor{currentfill}{rgb}{0.988927,0.988927,0.988927}%
\pgfsetfillcolor{currentfill}%
\pgfsetfillopacity{0.900000}%
\pgfsetlinewidth{0.501875pt}%
\definecolor{currentstroke}{rgb}{0.200000,0.200000,0.200000}%
\pgfsetstrokecolor{currentstroke}%
\pgfsetdash{}{0pt}%
\pgfpathmoveto{\pgfqpoint{3.722471in}{1.498035in}}%
\pgfpathlineto{\pgfqpoint{3.766508in}{1.500773in}}%
\pgfpathlineto{\pgfqpoint{3.800346in}{1.480881in}}%
\pgfpathlineto{\pgfqpoint{3.756286in}{1.478217in}}%
\pgfpathlineto{\pgfqpoint{3.722471in}{1.498035in}}%
\pgfpathclose%
\pgfusepath{stroke,fill}%
\end{pgfscope}%
\begin{pgfscope}%
\pgfpathrectangle{\pgfqpoint{0.050000in}{0.050000in}}{\pgfqpoint{4.900000in}{4.900000in}}%
\pgfusepath{clip}%
\pgfsetbuttcap%
\pgfsetroundjoin%
\definecolor{currentfill}{rgb}{0.784667,0.784667,0.784667}%
\pgfsetfillcolor{currentfill}%
\pgfsetfillopacity{0.900000}%
\pgfsetlinewidth{0.501875pt}%
\definecolor{currentstroke}{rgb}{0.200000,0.200000,0.200000}%
\pgfsetstrokecolor{currentstroke}%
\pgfsetdash{}{0pt}%
\pgfpathmoveto{\pgfqpoint{2.417907in}{2.256970in}}%
\pgfpathlineto{\pgfqpoint{2.463750in}{2.213301in}}%
\pgfpathlineto{\pgfqpoint{2.495521in}{2.192826in}}%
\pgfpathlineto{\pgfqpoint{2.449674in}{2.236460in}}%
\pgfpathlineto{\pgfqpoint{2.417907in}{2.256970in}}%
\pgfpathclose%
\pgfusepath{stroke,fill}%
\end{pgfscope}%
\begin{pgfscope}%
\pgfpathrectangle{\pgfqpoint{0.050000in}{0.050000in}}{\pgfqpoint{4.900000in}{4.900000in}}%
\pgfusepath{clip}%
\pgfsetbuttcap%
\pgfsetroundjoin%
\definecolor{currentfill}{rgb}{0.424083,0.424083,0.424083}%
\pgfsetfillcolor{currentfill}%
\pgfsetfillopacity{0.900000}%
\pgfsetlinewidth{0.501875pt}%
\definecolor{currentstroke}{rgb}{0.200000,0.200000,0.200000}%
\pgfsetstrokecolor{currentstroke}%
\pgfsetdash{}{0pt}%
\pgfpathmoveto{\pgfqpoint{1.436689in}{3.050342in}}%
\pgfpathlineto{\pgfqpoint{1.484612in}{3.004625in}}%
\pgfpathlineto{\pgfqpoint{1.514872in}{2.985244in}}%
\pgfpathlineto{\pgfqpoint{1.466939in}{3.030991in}}%
\pgfpathlineto{\pgfqpoint{1.436689in}{3.050342in}}%
\pgfpathclose%
\pgfusepath{stroke,fill}%
\end{pgfscope}%
\begin{pgfscope}%
\pgfpathrectangle{\pgfqpoint{0.050000in}{0.050000in}}{\pgfqpoint{4.900000in}{4.900000in}}%
\pgfusepath{clip}%
\pgfsetbuttcap%
\pgfsetroundjoin%
\definecolor{currentfill}{rgb}{0.998155,0.998155,0.998155}%
\pgfsetfillcolor{currentfill}%
\pgfsetfillopacity{0.900000}%
\pgfsetlinewidth{0.501875pt}%
\definecolor{currentstroke}{rgb}{0.200000,0.200000,0.200000}%
\pgfsetstrokecolor{currentstroke}%
\pgfsetdash{}{0pt}%
\pgfpathmoveto{\pgfqpoint{4.078376in}{1.449428in}}%
\pgfpathlineto{\pgfqpoint{4.122166in}{1.459842in}}%
\pgfpathlineto{\pgfqpoint{4.156557in}{1.439269in}}%
\pgfpathlineto{\pgfqpoint{4.112742in}{1.428932in}}%
\pgfpathlineto{\pgfqpoint{4.078376in}{1.449428in}}%
\pgfpathclose%
\pgfusepath{stroke,fill}%
\end{pgfscope}%
\begin{pgfscope}%
\pgfpathrectangle{\pgfqpoint{0.050000in}{0.050000in}}{\pgfqpoint{4.900000in}{4.900000in}}%
\pgfusepath{clip}%
\pgfsetbuttcap%
\pgfsetroundjoin%
\definecolor{currentfill}{rgb}{0.696194,0.696194,0.696194}%
\pgfsetfillcolor{currentfill}%
\pgfsetfillopacity{0.900000}%
\pgfsetlinewidth{0.501875pt}%
\definecolor{currentstroke}{rgb}{0.200000,0.200000,0.200000}%
\pgfsetstrokecolor{currentstroke}%
\pgfsetdash{}{0pt}%
\pgfpathmoveto{\pgfqpoint{2.153169in}{2.471265in}}%
\pgfpathlineto{\pgfqpoint{2.199579in}{2.426780in}}%
\pgfpathlineto{\pgfqpoint{2.230947in}{2.406512in}}%
\pgfpathlineto{\pgfqpoint{2.184531in}{2.451019in}}%
\pgfpathlineto{\pgfqpoint{2.153169in}{2.471265in}}%
\pgfpathclose%
\pgfusepath{stroke,fill}%
\end{pgfscope}%
\begin{pgfscope}%
\pgfpathrectangle{\pgfqpoint{0.050000in}{0.050000in}}{\pgfqpoint{4.900000in}{4.900000in}}%
\pgfusepath{clip}%
\pgfsetbuttcap%
\pgfsetroundjoin%
\definecolor{currentfill}{rgb}{0.915356,0.915356,0.915356}%
\pgfsetfillcolor{currentfill}%
\pgfsetfillopacity{0.900000}%
\pgfsetlinewidth{0.501875pt}%
\definecolor{currentstroke}{rgb}{0.200000,0.200000,0.200000}%
\pgfsetstrokecolor{currentstroke}%
\pgfsetdash{}{0pt}%
\pgfpathmoveto{\pgfqpoint{2.947274in}{1.850588in}}%
\pgfpathlineto{\pgfqpoint{2.992080in}{1.817732in}}%
\pgfpathlineto{\pgfqpoint{3.024688in}{1.797948in}}%
\pgfpathlineto{\pgfqpoint{2.979873in}{1.830557in}}%
\pgfpathlineto{\pgfqpoint{2.947274in}{1.850588in}}%
\pgfpathclose%
\pgfusepath{stroke,fill}%
\end{pgfscope}%
\begin{pgfscope}%
\pgfpathrectangle{\pgfqpoint{0.050000in}{0.050000in}}{\pgfqpoint{4.900000in}{4.900000in}}%
\pgfusepath{clip}%
\pgfsetbuttcap%
\pgfsetroundjoin%
\definecolor{currentfill}{rgb}{1.000000,1.000000,1.000000}%
\pgfsetfillcolor{currentfill}%
\pgfsetfillopacity{0.900000}%
\pgfsetlinewidth{0.501875pt}%
\definecolor{currentstroke}{rgb}{0.200000,0.200000,0.200000}%
\pgfsetstrokecolor{currentstroke}%
\pgfsetdash{}{0pt}%
\pgfpathmoveto{\pgfqpoint{4.356770in}{1.432424in}}%
\pgfpathlineto{\pgfqpoint{4.400295in}{1.444575in}}%
\pgfpathlineto{\pgfqpoint{4.435112in}{1.423545in}}%
\pgfpathlineto{\pgfqpoint{4.391562in}{1.411454in}}%
\pgfpathlineto{\pgfqpoint{4.356770in}{1.432424in}}%
\pgfpathclose%
\pgfusepath{stroke,fill}%
\end{pgfscope}%
\begin{pgfscope}%
\pgfpathrectangle{\pgfqpoint{0.050000in}{0.050000in}}{\pgfqpoint{4.900000in}{4.900000in}}%
\pgfusepath{clip}%
\pgfsetbuttcap%
\pgfsetroundjoin%
\definecolor{currentfill}{rgb}{0.992618,0.992618,0.992618}%
\pgfsetfillcolor{currentfill}%
\pgfsetfillopacity{0.900000}%
\pgfsetlinewidth{0.501875pt}%
\definecolor{currentstroke}{rgb}{0.200000,0.200000,0.200000}%
\pgfsetstrokecolor{currentstroke}%
\pgfsetdash{}{0pt}%
\pgfpathmoveto{\pgfqpoint{3.800346in}{1.480881in}}%
\pgfpathlineto{\pgfqpoint{3.844343in}{1.485764in}}%
\pgfpathlineto{\pgfqpoint{3.878312in}{1.465718in}}%
\pgfpathlineto{\pgfqpoint{3.834292in}{1.460916in}}%
\pgfpathlineto{\pgfqpoint{3.800346in}{1.480881in}}%
\pgfpathclose%
\pgfusepath{stroke,fill}%
\end{pgfscope}%
\begin{pgfscope}%
\pgfpathrectangle{\pgfqpoint{0.050000in}{0.050000in}}{\pgfqpoint{4.900000in}{4.900000in}}%
\pgfusepath{clip}%
\pgfsetbuttcap%
\pgfsetroundjoin%
\definecolor{currentfill}{rgb}{0.334764,0.334764,0.334764}%
\pgfsetfillcolor{currentfill}%
\pgfsetfillopacity{0.900000}%
\pgfsetlinewidth{0.501875pt}%
\definecolor{currentstroke}{rgb}{0.200000,0.200000,0.200000}%
\pgfsetstrokecolor{currentstroke}%
\pgfsetdash{}{0pt}%
\pgfpathmoveto{\pgfqpoint{1.171757in}{3.265109in}}%
\pgfpathlineto{\pgfqpoint{1.220249in}{3.218930in}}%
\pgfpathlineto{\pgfqpoint{1.250105in}{3.199876in}}%
\pgfpathlineto{\pgfqpoint{1.201602in}{3.246087in}}%
\pgfpathlineto{\pgfqpoint{1.171757in}{3.265109in}}%
\pgfpathclose%
\pgfusepath{stroke,fill}%
\end{pgfscope}%
\begin{pgfscope}%
\pgfpathrectangle{\pgfqpoint{0.050000in}{0.050000in}}{\pgfqpoint{4.900000in}{4.900000in}}%
\pgfusepath{clip}%
\pgfsetbuttcap%
\pgfsetroundjoin%
\definecolor{currentfill}{rgb}{0.586082,0.586082,0.586082}%
\pgfsetfillcolor{currentfill}%
\pgfsetfillopacity{0.900000}%
\pgfsetlinewidth{0.501875pt}%
\definecolor{currentstroke}{rgb}{0.200000,0.200000,0.200000}%
\pgfsetstrokecolor{currentstroke}%
\pgfsetdash{}{0pt}%
\pgfpathmoveto{\pgfqpoint{1.888340in}{2.685937in}}%
\pgfpathlineto{\pgfqpoint{1.935319in}{2.640980in}}%
\pgfpathlineto{\pgfqpoint{1.966284in}{2.621033in}}%
\pgfpathlineto{\pgfqpoint{1.919297in}{2.666018in}}%
\pgfpathlineto{\pgfqpoint{1.888340in}{2.685937in}}%
\pgfpathclose%
\pgfusepath{stroke,fill}%
\end{pgfscope}%
\begin{pgfscope}%
\pgfpathrectangle{\pgfqpoint{0.050000in}{0.050000in}}{\pgfqpoint{4.900000in}{4.900000in}}%
\pgfusepath{clip}%
\pgfsetbuttcap%
\pgfsetroundjoin%
\definecolor{currentfill}{rgb}{0.836340,0.836340,0.836340}%
\pgfsetfillcolor{currentfill}%
\pgfsetfillopacity{0.900000}%
\pgfsetlinewidth{0.501875pt}%
\definecolor{currentstroke}{rgb}{0.200000,0.200000,0.200000}%
\pgfsetstrokecolor{currentstroke}%
\pgfsetdash{}{0pt}%
\pgfpathmoveto{\pgfqpoint{2.605020in}{2.108801in}}%
\pgfpathlineto{\pgfqpoint{2.650491in}{2.066772in}}%
\pgfpathlineto{\pgfqpoint{2.682566in}{2.046351in}}%
\pgfpathlineto{\pgfqpoint{2.637092in}{2.088254in}}%
\pgfpathlineto{\pgfqpoint{2.605020in}{2.108801in}}%
\pgfpathclose%
\pgfusepath{stroke,fill}%
\end{pgfscope}%
\begin{pgfscope}%
\pgfpathrectangle{\pgfqpoint{0.050000in}{0.050000in}}{\pgfqpoint{4.900000in}{4.900000in}}%
\pgfusepath{clip}%
\pgfsetbuttcap%
\pgfsetroundjoin%
\definecolor{currentfill}{rgb}{0.487043,0.487043,0.487043}%
\pgfsetfillcolor{currentfill}%
\pgfsetfillopacity{0.900000}%
\pgfsetlinewidth{0.501875pt}%
\definecolor{currentstroke}{rgb}{0.200000,0.200000,0.200000}%
\pgfsetstrokecolor{currentstroke}%
\pgfsetdash{}{0pt}%
\pgfpathmoveto{\pgfqpoint{1.623423in}{2.900689in}}%
\pgfpathlineto{\pgfqpoint{1.670970in}{2.855270in}}%
\pgfpathlineto{\pgfqpoint{1.701531in}{2.835650in}}%
\pgfpathlineto{\pgfqpoint{1.653974in}{2.881099in}}%
\pgfpathlineto{\pgfqpoint{1.623423in}{2.900689in}}%
\pgfpathclose%
\pgfusepath{stroke,fill}%
\end{pgfscope}%
\begin{pgfscope}%
\pgfpathrectangle{\pgfqpoint{0.050000in}{0.050000in}}{\pgfqpoint{4.900000in}{4.900000in}}%
\pgfusepath{clip}%
\pgfsetbuttcap%
\pgfsetroundjoin%
\definecolor{currentfill}{rgb}{0.998155,0.998155,0.998155}%
\pgfsetfillcolor{currentfill}%
\pgfsetfillopacity{0.900000}%
\pgfsetlinewidth{0.501875pt}%
\definecolor{currentstroke}{rgb}{0.200000,0.200000,0.200000}%
\pgfsetstrokecolor{currentstroke}%
\pgfsetdash{}{0pt}%
\pgfpathmoveto{\pgfqpoint{4.156557in}{1.439269in}}%
\pgfpathlineto{\pgfqpoint{4.200293in}{1.450396in}}%
\pgfpathlineto{\pgfqpoint{4.234817in}{1.429687in}}%
\pgfpathlineto{\pgfqpoint{4.191056in}{1.418632in}}%
\pgfpathlineto{\pgfqpoint{4.156557in}{1.439269in}}%
\pgfpathclose%
\pgfusepath{stroke,fill}%
\end{pgfscope}%
\begin{pgfscope}%
\pgfpathrectangle{\pgfqpoint{0.050000in}{0.050000in}}{\pgfqpoint{4.900000in}{4.900000in}}%
\pgfusepath{clip}%
\pgfsetbuttcap%
\pgfsetroundjoin%
\definecolor{currentfill}{rgb}{0.898378,0.898378,0.898378}%
\pgfsetfillcolor{currentfill}%
\pgfsetfillopacity{0.900000}%
\pgfsetlinewidth{0.501875pt}%
\definecolor{currentstroke}{rgb}{0.200000,0.200000,0.200000}%
\pgfsetstrokecolor{currentstroke}%
\pgfsetdash{}{0pt}%
\pgfpathmoveto{\pgfqpoint{2.869822in}{1.906486in}}%
\pgfpathlineto{\pgfqpoint{2.914778in}{1.870565in}}%
\pgfpathlineto{\pgfqpoint{2.947274in}{1.850588in}}%
\pgfpathlineto{\pgfqpoint{2.902311in}{1.886269in}}%
\pgfpathlineto{\pgfqpoint{2.869822in}{1.906486in}}%
\pgfpathclose%
\pgfusepath{stroke,fill}%
\end{pgfscope}%
\begin{pgfscope}%
\pgfpathrectangle{\pgfqpoint{0.050000in}{0.050000in}}{\pgfqpoint{4.900000in}{4.900000in}}%
\pgfusepath{clip}%
\pgfsetbuttcap%
\pgfsetroundjoin%
\definecolor{currentfill}{rgb}{0.760554,0.760554,0.760554}%
\pgfsetfillcolor{currentfill}%
\pgfsetfillopacity{0.900000}%
\pgfsetlinewidth{0.501875pt}%
\definecolor{currentstroke}{rgb}{0.200000,0.200000,0.200000}%
\pgfsetstrokecolor{currentstroke}%
\pgfsetdash{}{0pt}%
\pgfpathmoveto{\pgfqpoint{2.340204in}{2.321474in}}%
\pgfpathlineto{\pgfqpoint{2.386238in}{2.277424in}}%
\pgfpathlineto{\pgfqpoint{2.417907in}{2.256970in}}%
\pgfpathlineto{\pgfqpoint{2.371869in}{2.301014in}}%
\pgfpathlineto{\pgfqpoint{2.340204in}{2.321474in}}%
\pgfpathclose%
\pgfusepath{stroke,fill}%
\end{pgfscope}%
\begin{pgfscope}%
\pgfpathrectangle{\pgfqpoint{0.050000in}{0.050000in}}{\pgfqpoint{4.900000in}{4.900000in}}%
\pgfusepath{clip}%
\pgfsetbuttcap%
\pgfsetroundjoin%
\definecolor{currentfill}{rgb}{0.994464,0.994464,0.994464}%
\pgfsetfillcolor{currentfill}%
\pgfsetfillopacity{0.900000}%
\pgfsetlinewidth{0.501875pt}%
\definecolor{currentstroke}{rgb}{0.200000,0.200000,0.200000}%
\pgfsetstrokecolor{currentstroke}%
\pgfsetdash{}{0pt}%
\pgfpathmoveto{\pgfqpoint{3.878312in}{1.465718in}}%
\pgfpathlineto{\pgfqpoint{3.922267in}{1.472407in}}%
\pgfpathlineto{\pgfqpoint{3.956367in}{1.452208in}}%
\pgfpathlineto{\pgfqpoint{3.912388in}{1.445603in}}%
\pgfpathlineto{\pgfqpoint{3.878312in}{1.465718in}}%
\pgfpathclose%
\pgfusepath{stroke,fill}%
\end{pgfscope}%
\begin{pgfscope}%
\pgfpathrectangle{\pgfqpoint{0.050000in}{0.050000in}}{\pgfqpoint{4.900000in}{4.900000in}}%
\pgfusepath{clip}%
\pgfsetbuttcap%
\pgfsetroundjoin%
\definecolor{currentfill}{rgb}{0.395663,0.395663,0.395663}%
\pgfsetfillcolor{currentfill}%
\pgfsetfillopacity{0.900000}%
\pgfsetlinewidth{0.501875pt}%
\definecolor{currentstroke}{rgb}{0.200000,0.200000,0.200000}%
\pgfsetstrokecolor{currentstroke}%
\pgfsetdash{}{0pt}%
\pgfpathmoveto{\pgfqpoint{1.358416in}{3.115516in}}%
\pgfpathlineto{\pgfqpoint{1.406532in}{3.069634in}}%
\pgfpathlineto{\pgfqpoint{1.436689in}{3.050342in}}%
\pgfpathlineto{\pgfqpoint{1.388562in}{3.096254in}}%
\pgfpathlineto{\pgfqpoint{1.358416in}{3.115516in}}%
\pgfpathclose%
\pgfusepath{stroke,fill}%
\end{pgfscope}%
\begin{pgfscope}%
\pgfpathrectangle{\pgfqpoint{0.050000in}{0.050000in}}{\pgfqpoint{4.900000in}{4.900000in}}%
\pgfusepath{clip}%
\pgfsetbuttcap%
\pgfsetroundjoin%
\definecolor{currentfill}{rgb}{0.657809,0.657809,0.657809}%
\pgfsetfillcolor{currentfill}%
\pgfsetfillopacity{0.900000}%
\pgfsetlinewidth{0.501875pt}%
\definecolor{currentstroke}{rgb}{0.200000,0.200000,0.200000}%
\pgfsetstrokecolor{currentstroke}%
\pgfsetdash{}{0pt}%
\pgfpathmoveto{\pgfqpoint{2.075301in}{2.536106in}}%
\pgfpathlineto{\pgfqpoint{2.121903in}{2.491450in}}%
\pgfpathlineto{\pgfqpoint{2.153169in}{2.471265in}}%
\pgfpathlineto{\pgfqpoint{2.106560in}{2.515947in}}%
\pgfpathlineto{\pgfqpoint{2.075301in}{2.536106in}}%
\pgfpathclose%
\pgfusepath{stroke,fill}%
\end{pgfscope}%
\begin{pgfscope}%
\pgfpathrectangle{\pgfqpoint{0.050000in}{0.050000in}}{\pgfqpoint{4.900000in}{4.900000in}}%
\pgfusepath{clip}%
\pgfsetbuttcap%
\pgfsetroundjoin%
\definecolor{currentfill}{rgb}{0.966782,0.966782,0.966782}%
\pgfsetfillcolor{currentfill}%
\pgfsetfillopacity{0.900000}%
\pgfsetlinewidth{0.501875pt}%
\definecolor{currentstroke}{rgb}{0.200000,0.200000,0.200000}%
\pgfsetstrokecolor{currentstroke}%
\pgfsetdash{}{0pt}%
\pgfpathmoveto{\pgfqpoint{3.367501in}{1.607189in}}%
\pgfpathlineto{\pgfqpoint{3.411806in}{1.593984in}}%
\pgfpathlineto{\pgfqpoint{3.445113in}{1.574558in}}%
\pgfpathlineto{\pgfqpoint{3.400790in}{1.587689in}}%
\pgfpathlineto{\pgfqpoint{3.367501in}{1.607189in}}%
\pgfpathclose%
\pgfusepath{stroke,fill}%
\end{pgfscope}%
\begin{pgfscope}%
\pgfpathrectangle{\pgfqpoint{0.050000in}{0.050000in}}{\pgfqpoint{4.900000in}{4.900000in}}%
\pgfusepath{clip}%
\pgfsetbuttcap%
\pgfsetroundjoin%
\definecolor{currentfill}{rgb}{0.295271,0.295271,0.295271}%
\pgfsetfillcolor{currentfill}%
\pgfsetfillopacity{0.900000}%
\pgfsetlinewidth{0.501875pt}%
\definecolor{currentstroke}{rgb}{0.200000,0.200000,0.200000}%
\pgfsetstrokecolor{currentstroke}%
\pgfsetdash{}{0pt}%
\pgfpathmoveto{\pgfqpoint{1.093320in}{3.330417in}}%
\pgfpathlineto{\pgfqpoint{1.142005in}{3.284073in}}%
\pgfpathlineto{\pgfqpoint{1.171757in}{3.265109in}}%
\pgfpathlineto{\pgfqpoint{1.123059in}{3.311485in}}%
\pgfpathlineto{\pgfqpoint{1.093320in}{3.330417in}}%
\pgfpathclose%
\pgfusepath{stroke,fill}%
\end{pgfscope}%
\begin{pgfscope}%
\pgfpathrectangle{\pgfqpoint{0.050000in}{0.050000in}}{\pgfqpoint{4.900000in}{4.900000in}}%
\pgfusepath{clip}%
\pgfsetbuttcap%
\pgfsetroundjoin%
\definecolor{currentfill}{rgb}{0.974164,0.974164,0.974164}%
\pgfsetfillcolor{currentfill}%
\pgfsetfillopacity{0.900000}%
\pgfsetlinewidth{0.501875pt}%
\definecolor{currentstroke}{rgb}{0.200000,0.200000,0.200000}%
\pgfsetstrokecolor{currentstroke}%
\pgfsetdash{}{0pt}%
\pgfpathmoveto{\pgfqpoint{3.445113in}{1.574558in}}%
\pgfpathlineto{\pgfqpoint{3.489357in}{1.565150in}}%
\pgfpathlineto{\pgfqpoint{3.522789in}{1.545672in}}%
\pgfpathlineto{\pgfqpoint{3.478526in}{1.555049in}}%
\pgfpathlineto{\pgfqpoint{3.445113in}{1.574558in}}%
\pgfpathclose%
\pgfusepath{stroke,fill}%
\end{pgfscope}%
\begin{pgfscope}%
\pgfpathrectangle{\pgfqpoint{0.050000in}{0.050000in}}{\pgfqpoint{4.900000in}{4.900000in}}%
\pgfusepath{clip}%
\pgfsetbuttcap%
\pgfsetroundjoin%
\definecolor{currentfill}{rgb}{0.961246,0.961246,0.961246}%
\pgfsetfillcolor{currentfill}%
\pgfsetfillopacity{0.900000}%
\pgfsetlinewidth{0.501875pt}%
\definecolor{currentstroke}{rgb}{0.200000,0.200000,0.200000}%
\pgfsetstrokecolor{currentstroke}%
\pgfsetdash{}{0pt}%
\pgfpathmoveto{\pgfqpoint{3.289939in}{1.643795in}}%
\pgfpathlineto{\pgfqpoint{3.334317in}{1.626606in}}%
\pgfpathlineto{\pgfqpoint{3.367501in}{1.607189in}}%
\pgfpathlineto{\pgfqpoint{3.323106in}{1.624259in}}%
\pgfpathlineto{\pgfqpoint{3.289939in}{1.643795in}}%
\pgfpathclose%
\pgfusepath{stroke,fill}%
\end{pgfscope}%
\begin{pgfscope}%
\pgfpathrectangle{\pgfqpoint{0.050000in}{0.050000in}}{\pgfqpoint{4.900000in}{4.900000in}}%
\pgfusepath{clip}%
\pgfsetbuttcap%
\pgfsetroundjoin%
\definecolor{currentfill}{rgb}{0.555940,0.555940,0.555940}%
\pgfsetfillcolor{currentfill}%
\pgfsetfillopacity{0.900000}%
\pgfsetlinewidth{0.501875pt}%
\definecolor{currentstroke}{rgb}{0.200000,0.200000,0.200000}%
\pgfsetstrokecolor{currentstroke}%
\pgfsetdash{}{0pt}%
\pgfpathmoveto{\pgfqpoint{1.810308in}{2.750917in}}%
\pgfpathlineto{\pgfqpoint{1.857479in}{2.705796in}}%
\pgfpathlineto{\pgfqpoint{1.888340in}{2.685937in}}%
\pgfpathlineto{\pgfqpoint{1.841162in}{2.731086in}}%
\pgfpathlineto{\pgfqpoint{1.810308in}{2.750917in}}%
\pgfpathclose%
\pgfusepath{stroke,fill}%
\end{pgfscope}%
\begin{pgfscope}%
\pgfpathrectangle{\pgfqpoint{0.050000in}{0.050000in}}{\pgfqpoint{4.900000in}{4.900000in}}%
\pgfusepath{clip}%
\pgfsetbuttcap%
\pgfsetroundjoin%
\definecolor{currentfill}{rgb}{0.994464,0.994464,0.994464}%
\pgfsetfillcolor{currentfill}%
\pgfsetfillopacity{0.900000}%
\pgfsetlinewidth{0.501875pt}%
\definecolor{currentstroke}{rgb}{0.200000,0.200000,0.200000}%
\pgfsetstrokecolor{currentstroke}%
\pgfsetdash{}{0pt}%
\pgfpathmoveto{\pgfqpoint{3.956367in}{1.452208in}}%
\pgfpathlineto{\pgfqpoint{4.000278in}{1.460389in}}%
\pgfpathlineto{\pgfqpoint{4.034511in}{1.440039in}}%
\pgfpathlineto{\pgfqpoint{3.990575in}{1.431941in}}%
\pgfpathlineto{\pgfqpoint{3.956367in}{1.452208in}}%
\pgfpathclose%
\pgfusepath{stroke,fill}%
\end{pgfscope}%
\begin{pgfscope}%
\pgfpathrectangle{\pgfqpoint{0.050000in}{0.050000in}}{\pgfqpoint{4.900000in}{4.900000in}}%
\pgfusepath{clip}%
\pgfsetbuttcap%
\pgfsetroundjoin%
\definecolor{currentfill}{rgb}{1.000000,1.000000,1.000000}%
\pgfsetfillcolor{currentfill}%
\pgfsetfillopacity{0.900000}%
\pgfsetlinewidth{0.501875pt}%
\definecolor{currentstroke}{rgb}{0.200000,0.200000,0.200000}%
\pgfsetstrokecolor{currentstroke}%
\pgfsetdash{}{0pt}%
\pgfpathmoveto{\pgfqpoint{4.234817in}{1.429687in}}%
\pgfpathlineto{\pgfqpoint{4.278496in}{1.441318in}}%
\pgfpathlineto{\pgfqpoint{4.313154in}{1.420478in}}%
\pgfpathlineto{\pgfqpoint{4.269450in}{1.408914in}}%
\pgfpathlineto{\pgfqpoint{4.234817in}{1.429687in}}%
\pgfpathclose%
\pgfusepath{stroke,fill}%
\end{pgfscope}%
\begin{pgfscope}%
\pgfpathrectangle{\pgfqpoint{0.050000in}{0.050000in}}{\pgfqpoint{4.900000in}{4.900000in}}%
\pgfusepath{clip}%
\pgfsetbuttcap%
\pgfsetroundjoin%
\definecolor{currentfill}{rgb}{0.977855,0.977855,0.977855}%
\pgfsetfillcolor{currentfill}%
\pgfsetfillopacity{0.900000}%
\pgfsetlinewidth{0.501875pt}%
\definecolor{currentstroke}{rgb}{0.200000,0.200000,0.200000}%
\pgfsetstrokecolor{currentstroke}%
\pgfsetdash{}{0pt}%
\pgfpathmoveto{\pgfqpoint{3.522789in}{1.545672in}}%
\pgfpathlineto{\pgfqpoint{3.566980in}{1.539790in}}%
\pgfpathlineto{\pgfqpoint{3.600539in}{1.520223in}}%
\pgfpathlineto{\pgfqpoint{3.556327in}{1.526110in}}%
\pgfpathlineto{\pgfqpoint{3.522789in}{1.545672in}}%
\pgfpathclose%
\pgfusepath{stroke,fill}%
\end{pgfscope}%
\begin{pgfscope}%
\pgfpathrectangle{\pgfqpoint{0.050000in}{0.050000in}}{\pgfqpoint{4.900000in}{4.900000in}}%
\pgfusepath{clip}%
\pgfsetbuttcap%
\pgfsetroundjoin%
\definecolor{currentfill}{rgb}{0.953864,0.953864,0.953864}%
\pgfsetfillcolor{currentfill}%
\pgfsetfillopacity{0.900000}%
\pgfsetlinewidth{0.501875pt}%
\definecolor{currentstroke}{rgb}{0.200000,0.200000,0.200000}%
\pgfsetstrokecolor{currentstroke}%
\pgfsetdash{}{0pt}%
\pgfpathmoveto{\pgfqpoint{3.212413in}{1.684506in}}%
\pgfpathlineto{\pgfqpoint{3.256877in}{1.663252in}}%
\pgfpathlineto{\pgfqpoint{3.289939in}{1.643795in}}%
\pgfpathlineto{\pgfqpoint{3.245460in}{1.664885in}}%
\pgfpathlineto{\pgfqpoint{3.212413in}{1.684506in}}%
\pgfpathclose%
\pgfusepath{stroke,fill}%
\end{pgfscope}%
\begin{pgfscope}%
\pgfpathrectangle{\pgfqpoint{0.050000in}{0.050000in}}{\pgfqpoint{4.900000in}{4.900000in}}%
\pgfusepath{clip}%
\pgfsetbuttcap%
\pgfsetroundjoin%
\definecolor{currentfill}{rgb}{0.812226,0.812226,0.812226}%
\pgfsetfillcolor{currentfill}%
\pgfsetfillopacity{0.900000}%
\pgfsetlinewidth{0.501875pt}%
\definecolor{currentstroke}{rgb}{0.200000,0.200000,0.200000}%
\pgfsetstrokecolor{currentstroke}%
\pgfsetdash{}{0pt}%
\pgfpathmoveto{\pgfqpoint{2.527391in}{2.172299in}}%
\pgfpathlineto{\pgfqpoint{2.573048in}{2.129300in}}%
\pgfpathlineto{\pgfqpoint{2.605020in}{2.108801in}}%
\pgfpathlineto{\pgfqpoint{2.559360in}{2.151722in}}%
\pgfpathlineto{\pgfqpoint{2.527391in}{2.172299in}}%
\pgfpathclose%
\pgfusepath{stroke,fill}%
\end{pgfscope}%
\begin{pgfscope}%
\pgfpathrectangle{\pgfqpoint{0.050000in}{0.050000in}}{\pgfqpoint{4.900000in}{4.900000in}}%
\pgfusepath{clip}%
\pgfsetbuttcap%
\pgfsetroundjoin%
\definecolor{currentfill}{rgb}{0.878570,0.878570,0.878570}%
\pgfsetfillcolor{currentfill}%
\pgfsetfillopacity{0.900000}%
\pgfsetlinewidth{0.501875pt}%
\definecolor{currentstroke}{rgb}{0.200000,0.200000,0.200000}%
\pgfsetstrokecolor{currentstroke}%
\pgfsetdash{}{0pt}%
\pgfpathmoveto{\pgfqpoint{2.792316in}{1.965110in}}%
\pgfpathlineto{\pgfqpoint{2.837434in}{1.926653in}}%
\pgfpathlineto{\pgfqpoint{2.869822in}{1.906486in}}%
\pgfpathlineto{\pgfqpoint{2.824697in}{1.944728in}}%
\pgfpathlineto{\pgfqpoint{2.792316in}{1.965110in}}%
\pgfpathclose%
\pgfusepath{stroke,fill}%
\end{pgfscope}%
\begin{pgfscope}%
\pgfpathrectangle{\pgfqpoint{0.050000in}{0.050000in}}{\pgfqpoint{4.900000in}{4.900000in}}%
\pgfusepath{clip}%
\pgfsetbuttcap%
\pgfsetroundjoin%
\definecolor{currentfill}{rgb}{0.456901,0.456901,0.456901}%
\pgfsetfillcolor{currentfill}%
\pgfsetfillopacity{0.900000}%
\pgfsetlinewidth{0.501875pt}%
\definecolor{currentstroke}{rgb}{0.200000,0.200000,0.200000}%
\pgfsetstrokecolor{currentstroke}%
\pgfsetdash{}{0pt}%
\pgfpathmoveto{\pgfqpoint{1.545225in}{2.965803in}}%
\pgfpathlineto{\pgfqpoint{1.592966in}{2.920220in}}%
\pgfpathlineto{\pgfqpoint{1.623423in}{2.900689in}}%
\pgfpathlineto{\pgfqpoint{1.575673in}{2.946302in}}%
\pgfpathlineto{\pgfqpoint{1.545225in}{2.965803in}}%
\pgfpathclose%
\pgfusepath{stroke,fill}%
\end{pgfscope}%
\begin{pgfscope}%
\pgfpathrectangle{\pgfqpoint{0.050000in}{0.050000in}}{\pgfqpoint{4.900000in}{4.900000in}}%
\pgfusepath{clip}%
\pgfsetbuttcap%
\pgfsetroundjoin%
\definecolor{currentfill}{rgb}{0.983391,0.983391,0.983391}%
\pgfsetfillcolor{currentfill}%
\pgfsetfillopacity{0.900000}%
\pgfsetlinewidth{0.501875pt}%
\definecolor{currentstroke}{rgb}{0.200000,0.200000,0.200000}%
\pgfsetstrokecolor{currentstroke}%
\pgfsetdash{}{0pt}%
\pgfpathmoveto{\pgfqpoint{3.600539in}{1.520223in}}%
\pgfpathlineto{\pgfqpoint{3.644683in}{1.517545in}}%
\pgfpathlineto{\pgfqpoint{3.678369in}{1.497862in}}%
\pgfpathlineto{\pgfqpoint{3.634203in}{1.500575in}}%
\pgfpathlineto{\pgfqpoint{3.600539in}{1.520223in}}%
\pgfpathclose%
\pgfusepath{stroke,fill}%
\end{pgfscope}%
\begin{pgfscope}%
\pgfpathrectangle{\pgfqpoint{0.050000in}{0.050000in}}{\pgfqpoint{4.900000in}{4.900000in}}%
\pgfusepath{clip}%
\pgfsetbuttcap%
\pgfsetroundjoin%
\definecolor{currentfill}{rgb}{0.729781,0.729781,0.729781}%
\pgfsetfillcolor{currentfill}%
\pgfsetfillopacity{0.900000}%
\pgfsetlinewidth{0.501875pt}%
\definecolor{currentstroke}{rgb}{0.200000,0.200000,0.200000}%
\pgfsetstrokecolor{currentstroke}%
\pgfsetdash{}{0pt}%
\pgfpathmoveto{\pgfqpoint{2.262412in}{2.386183in}}%
\pgfpathlineto{\pgfqpoint{2.308637in}{2.341875in}}%
\pgfpathlineto{\pgfqpoint{2.340204in}{2.321474in}}%
\pgfpathlineto{\pgfqpoint{2.293975in}{2.365794in}}%
\pgfpathlineto{\pgfqpoint{2.262412in}{2.386183in}}%
\pgfpathclose%
\pgfusepath{stroke,fill}%
\end{pgfscope}%
\begin{pgfscope}%
\pgfpathrectangle{\pgfqpoint{0.050000in}{0.050000in}}{\pgfqpoint{4.900000in}{4.900000in}}%
\pgfusepath{clip}%
\pgfsetbuttcap%
\pgfsetroundjoin%
\definecolor{currentfill}{rgb}{0.944637,0.944637,0.944637}%
\pgfsetfillcolor{currentfill}%
\pgfsetfillopacity{0.900000}%
\pgfsetlinewidth{0.501875pt}%
\definecolor{currentstroke}{rgb}{0.200000,0.200000,0.200000}%
\pgfsetstrokecolor{currentstroke}%
\pgfsetdash{}{0pt}%
\pgfpathmoveto{\pgfqpoint{3.134906in}{1.729323in}}%
\pgfpathlineto{\pgfqpoint{3.179471in}{1.704051in}}%
\pgfpathlineto{\pgfqpoint{3.212413in}{1.684506in}}%
\pgfpathlineto{\pgfqpoint{3.167835in}{1.709574in}}%
\pgfpathlineto{\pgfqpoint{3.134906in}{1.729323in}}%
\pgfpathclose%
\pgfusepath{stroke,fill}%
\end{pgfscope}%
\begin{pgfscope}%
\pgfpathrectangle{\pgfqpoint{0.050000in}{0.050000in}}{\pgfqpoint{4.900000in}{4.900000in}}%
\pgfusepath{clip}%
\pgfsetbuttcap%
\pgfsetroundjoin%
\definecolor{currentfill}{rgb}{0.363183,0.363183,0.363183}%
\pgfsetfillcolor{currentfill}%
\pgfsetfillopacity{0.900000}%
\pgfsetlinewidth{0.501875pt}%
\definecolor{currentstroke}{rgb}{0.200000,0.200000,0.200000}%
\pgfsetstrokecolor{currentstroke}%
\pgfsetdash{}{0pt}%
\pgfpathmoveto{\pgfqpoint{1.280054in}{3.180764in}}%
\pgfpathlineto{\pgfqpoint{1.328363in}{3.134718in}}%
\pgfpathlineto{\pgfqpoint{1.358416in}{3.115516in}}%
\pgfpathlineto{\pgfqpoint{1.310095in}{3.161593in}}%
\pgfpathlineto{\pgfqpoint{1.280054in}{3.180764in}}%
\pgfpathclose%
\pgfusepath{stroke,fill}%
\end{pgfscope}%
\begin{pgfscope}%
\pgfpathrectangle{\pgfqpoint{0.050000in}{0.050000in}}{\pgfqpoint{4.900000in}{4.900000in}}%
\pgfusepath{clip}%
\pgfsetbuttcap%
\pgfsetroundjoin%
\definecolor{currentfill}{rgb}{0.987082,0.987082,0.987082}%
\pgfsetfillcolor{currentfill}%
\pgfsetfillopacity{0.900000}%
\pgfsetlinewidth{0.501875pt}%
\definecolor{currentstroke}{rgb}{0.200000,0.200000,0.200000}%
\pgfsetstrokecolor{currentstroke}%
\pgfsetdash{}{0pt}%
\pgfpathmoveto{\pgfqpoint{3.678369in}{1.497862in}}%
\pgfpathlineto{\pgfqpoint{3.722471in}{1.498035in}}%
\pgfpathlineto{\pgfqpoint{3.756286in}{1.478217in}}%
\pgfpathlineto{\pgfqpoint{3.712162in}{1.478100in}}%
\pgfpathlineto{\pgfqpoint{3.678369in}{1.497862in}}%
\pgfpathclose%
\pgfusepath{stroke,fill}%
\end{pgfscope}%
\begin{pgfscope}%
\pgfpathrectangle{\pgfqpoint{0.050000in}{0.050000in}}{\pgfqpoint{4.900000in}{4.900000in}}%
\pgfusepath{clip}%
\pgfsetbuttcap%
\pgfsetroundjoin%
\definecolor{currentfill}{rgb}{0.624221,0.624221,0.624221}%
\pgfsetfillcolor{currentfill}%
\pgfsetfillopacity{0.900000}%
\pgfsetlinewidth{0.501875pt}%
\definecolor{currentstroke}{rgb}{0.200000,0.200000,0.200000}%
\pgfsetstrokecolor{currentstroke}%
\pgfsetdash{}{0pt}%
\pgfpathmoveto{\pgfqpoint{1.997344in}{2.601024in}}%
\pgfpathlineto{\pgfqpoint{2.044138in}{2.556203in}}%
\pgfpathlineto{\pgfqpoint{2.075301in}{2.536106in}}%
\pgfpathlineto{\pgfqpoint{2.028500in}{2.580954in}}%
\pgfpathlineto{\pgfqpoint{1.997344in}{2.601024in}}%
\pgfpathclose%
\pgfusepath{stroke,fill}%
\end{pgfscope}%
\begin{pgfscope}%
\pgfpathrectangle{\pgfqpoint{0.050000in}{0.050000in}}{\pgfqpoint{4.900000in}{4.900000in}}%
\pgfusepath{clip}%
\pgfsetbuttcap%
\pgfsetroundjoin%
\definecolor{currentfill}{rgb}{0.996309,0.996309,0.996309}%
\pgfsetfillcolor{currentfill}%
\pgfsetfillopacity{0.900000}%
\pgfsetlinewidth{0.501875pt}%
\definecolor{currentstroke}{rgb}{0.200000,0.200000,0.200000}%
\pgfsetstrokecolor{currentstroke}%
\pgfsetdash{}{0pt}%
\pgfpathmoveto{\pgfqpoint{4.034511in}{1.440039in}}%
\pgfpathlineto{\pgfqpoint{4.078376in}{1.449428in}}%
\pgfpathlineto{\pgfqpoint{4.112742in}{1.428932in}}%
\pgfpathlineto{\pgfqpoint{4.068852in}{1.419623in}}%
\pgfpathlineto{\pgfqpoint{4.034511in}{1.440039in}}%
\pgfpathclose%
\pgfusepath{stroke,fill}%
\end{pgfscope}%
\begin{pgfscope}%
\pgfpathrectangle{\pgfqpoint{0.050000in}{0.050000in}}{\pgfqpoint{4.900000in}{4.900000in}}%
\pgfusepath{clip}%
\pgfsetbuttcap%
\pgfsetroundjoin%
\definecolor{currentfill}{rgb}{1.000000,1.000000,1.000000}%
\pgfsetfillcolor{currentfill}%
\pgfsetfillopacity{0.900000}%
\pgfsetlinewidth{0.501875pt}%
\definecolor{currentstroke}{rgb}{0.200000,0.200000,0.200000}%
\pgfsetstrokecolor{currentstroke}%
\pgfsetdash{}{0pt}%
\pgfpathmoveto{\pgfqpoint{4.313154in}{1.420478in}}%
\pgfpathlineto{\pgfqpoint{4.356770in}{1.432424in}}%
\pgfpathlineto{\pgfqpoint{4.391562in}{1.411454in}}%
\pgfpathlineto{\pgfqpoint{4.347921in}{1.399572in}}%
\pgfpathlineto{\pgfqpoint{4.313154in}{1.420478in}}%
\pgfpathclose%
\pgfusepath{stroke,fill}%
\end{pgfscope}%
\begin{pgfscope}%
\pgfpathrectangle{\pgfqpoint{0.050000in}{0.050000in}}{\pgfqpoint{4.900000in}{4.900000in}}%
\pgfusepath{clip}%
\pgfsetbuttcap%
\pgfsetroundjoin%
\definecolor{currentfill}{rgb}{0.256517,0.256517,0.256517}%
\pgfsetfillcolor{currentfill}%
\pgfsetfillopacity{0.900000}%
\pgfsetlinewidth{0.501875pt}%
\definecolor{currentstroke}{rgb}{0.200000,0.200000,0.200000}%
\pgfsetstrokecolor{currentstroke}%
\pgfsetdash{}{0pt}%
\pgfpathmoveto{\pgfqpoint{1.014792in}{3.395800in}}%
\pgfpathlineto{\pgfqpoint{1.063672in}{3.349291in}}%
\pgfpathlineto{\pgfqpoint{1.093320in}{3.330417in}}%
\pgfpathlineto{\pgfqpoint{1.044427in}{3.376959in}}%
\pgfpathlineto{\pgfqpoint{1.014792in}{3.395800in}}%
\pgfpathclose%
\pgfusepath{stroke,fill}%
\end{pgfscope}%
\begin{pgfscope}%
\pgfpathrectangle{\pgfqpoint{0.050000in}{0.050000in}}{\pgfqpoint{4.900000in}{4.900000in}}%
\pgfusepath{clip}%
\pgfsetbuttcap%
\pgfsetroundjoin%
\definecolor{currentfill}{rgb}{0.932334,0.932334,0.932334}%
\pgfsetfillcolor{currentfill}%
\pgfsetfillopacity{0.900000}%
\pgfsetlinewidth{0.501875pt}%
\definecolor{currentstroke}{rgb}{0.200000,0.200000,0.200000}%
\pgfsetstrokecolor{currentstroke}%
\pgfsetdash{}{0pt}%
\pgfpathmoveto{\pgfqpoint{3.057399in}{1.778102in}}%
\pgfpathlineto{\pgfqpoint{3.102081in}{1.749001in}}%
\pgfpathlineto{\pgfqpoint{3.134906in}{1.729323in}}%
\pgfpathlineto{\pgfqpoint{3.090213in}{1.758193in}}%
\pgfpathlineto{\pgfqpoint{3.057399in}{1.778102in}}%
\pgfpathclose%
\pgfusepath{stroke,fill}%
\end{pgfscope}%
\begin{pgfscope}%
\pgfpathrectangle{\pgfqpoint{0.050000in}{0.050000in}}{\pgfqpoint{4.900000in}{4.900000in}}%
\pgfusepath{clip}%
\pgfsetbuttcap%
\pgfsetroundjoin%
\definecolor{currentfill}{rgb}{0.521492,0.521492,0.521492}%
\pgfsetfillcolor{currentfill}%
\pgfsetfillopacity{0.900000}%
\pgfsetlinewidth{0.501875pt}%
\definecolor{currentstroke}{rgb}{0.200000,0.200000,0.200000}%
\pgfsetstrokecolor{currentstroke}%
\pgfsetdash{}{0pt}%
\pgfpathmoveto{\pgfqpoint{1.732186in}{2.815970in}}%
\pgfpathlineto{\pgfqpoint{1.779549in}{2.770686in}}%
\pgfpathlineto{\pgfqpoint{1.810308in}{2.750917in}}%
\pgfpathlineto{\pgfqpoint{1.762936in}{2.796229in}}%
\pgfpathlineto{\pgfqpoint{1.732186in}{2.815970in}}%
\pgfpathclose%
\pgfusepath{stroke,fill}%
\end{pgfscope}%
\begin{pgfscope}%
\pgfpathrectangle{\pgfqpoint{0.050000in}{0.050000in}}{\pgfqpoint{4.900000in}{4.900000in}}%
\pgfusepath{clip}%
\pgfsetbuttcap%
\pgfsetroundjoin%
\definecolor{currentfill}{rgb}{0.858762,0.858762,0.858762}%
\pgfsetfillcolor{currentfill}%
\pgfsetfillopacity{0.900000}%
\pgfsetlinewidth{0.501875pt}%
\definecolor{currentstroke}{rgb}{0.200000,0.200000,0.200000}%
\pgfsetstrokecolor{currentstroke}%
\pgfsetdash{}{0pt}%
\pgfpathmoveto{\pgfqpoint{2.714742in}{2.025886in}}%
\pgfpathlineto{\pgfqpoint{2.760035in}{1.985446in}}%
\pgfpathlineto{\pgfqpoint{2.792316in}{1.965110in}}%
\pgfpathlineto{\pgfqpoint{2.747018in}{2.005375in}}%
\pgfpathlineto{\pgfqpoint{2.714742in}{2.025886in}}%
\pgfpathclose%
\pgfusepath{stroke,fill}%
\end{pgfscope}%
\begin{pgfscope}%
\pgfpathrectangle{\pgfqpoint{0.050000in}{0.050000in}}{\pgfqpoint{4.900000in}{4.900000in}}%
\pgfusepath{clip}%
\pgfsetbuttcap%
\pgfsetroundjoin%
\definecolor{currentfill}{rgb}{0.988927,0.988927,0.988927}%
\pgfsetfillcolor{currentfill}%
\pgfsetfillopacity{0.900000}%
\pgfsetlinewidth{0.501875pt}%
\definecolor{currentstroke}{rgb}{0.200000,0.200000,0.200000}%
\pgfsetstrokecolor{currentstroke}%
\pgfsetdash{}{0pt}%
\pgfpathmoveto{\pgfqpoint{3.756286in}{1.478217in}}%
\pgfpathlineto{\pgfqpoint{3.800346in}{1.480881in}}%
\pgfpathlineto{\pgfqpoint{3.834292in}{1.460916in}}%
\pgfpathlineto{\pgfqpoint{3.790208in}{1.458323in}}%
\pgfpathlineto{\pgfqpoint{3.756286in}{1.478217in}}%
\pgfpathclose%
\pgfusepath{stroke,fill}%
\end{pgfscope}%
\begin{pgfscope}%
\pgfpathrectangle{\pgfqpoint{0.050000in}{0.050000in}}{\pgfqpoint{4.900000in}{4.900000in}}%
\pgfusepath{clip}%
\pgfsetbuttcap%
\pgfsetroundjoin%
\definecolor{currentfill}{rgb}{0.784667,0.784667,0.784667}%
\pgfsetfillcolor{currentfill}%
\pgfsetfillopacity{0.900000}%
\pgfsetlinewidth{0.501875pt}%
\definecolor{currentstroke}{rgb}{0.200000,0.200000,0.200000}%
\pgfsetstrokecolor{currentstroke}%
\pgfsetdash{}{0pt}%
\pgfpathmoveto{\pgfqpoint{2.449674in}{2.236460in}}%
\pgfpathlineto{\pgfqpoint{2.495521in}{2.192826in}}%
\pgfpathlineto{\pgfqpoint{2.527391in}{2.172299in}}%
\pgfpathlineto{\pgfqpoint{2.481540in}{2.215896in}}%
\pgfpathlineto{\pgfqpoint{2.449674in}{2.236460in}}%
\pgfpathclose%
\pgfusepath{stroke,fill}%
\end{pgfscope}%
\begin{pgfscope}%
\pgfpathrectangle{\pgfqpoint{0.050000in}{0.050000in}}{\pgfqpoint{4.900000in}{4.900000in}}%
\pgfusepath{clip}%
\pgfsetbuttcap%
\pgfsetroundjoin%
\definecolor{currentfill}{rgb}{0.424083,0.424083,0.424083}%
\pgfsetfillcolor{currentfill}%
\pgfsetfillopacity{0.900000}%
\pgfsetlinewidth{0.501875pt}%
\definecolor{currentstroke}{rgb}{0.200000,0.200000,0.200000}%
\pgfsetstrokecolor{currentstroke}%
\pgfsetdash{}{0pt}%
\pgfpathmoveto{\pgfqpoint{1.466939in}{3.030991in}}%
\pgfpathlineto{\pgfqpoint{1.514872in}{2.985244in}}%
\pgfpathlineto{\pgfqpoint{1.545225in}{2.965803in}}%
\pgfpathlineto{\pgfqpoint{1.497282in}{3.011580in}}%
\pgfpathlineto{\pgfqpoint{1.466939in}{3.030991in}}%
\pgfpathclose%
\pgfusepath{stroke,fill}%
\end{pgfscope}%
\begin{pgfscope}%
\pgfpathrectangle{\pgfqpoint{0.050000in}{0.050000in}}{\pgfqpoint{4.900000in}{4.900000in}}%
\pgfusepath{clip}%
\pgfsetbuttcap%
\pgfsetroundjoin%
\definecolor{currentfill}{rgb}{0.998155,0.998155,0.998155}%
\pgfsetfillcolor{currentfill}%
\pgfsetfillopacity{0.900000}%
\pgfsetlinewidth{0.501875pt}%
\definecolor{currentstroke}{rgb}{0.200000,0.200000,0.200000}%
\pgfsetstrokecolor{currentstroke}%
\pgfsetdash{}{0pt}%
\pgfpathmoveto{\pgfqpoint{4.112742in}{1.428932in}}%
\pgfpathlineto{\pgfqpoint{4.156557in}{1.439269in}}%
\pgfpathlineto{\pgfqpoint{4.191056in}{1.418632in}}%
\pgfpathlineto{\pgfqpoint{4.147215in}{1.408370in}}%
\pgfpathlineto{\pgfqpoint{4.112742in}{1.428932in}}%
\pgfpathclose%
\pgfusepath{stroke,fill}%
\end{pgfscope}%
\begin{pgfscope}%
\pgfpathrectangle{\pgfqpoint{0.050000in}{0.050000in}}{\pgfqpoint{4.900000in}{4.900000in}}%
\pgfusepath{clip}%
\pgfsetbuttcap%
\pgfsetroundjoin%
\definecolor{currentfill}{rgb}{0.696194,0.696194,0.696194}%
\pgfsetfillcolor{currentfill}%
\pgfsetfillopacity{0.900000}%
\pgfsetlinewidth{0.501875pt}%
\definecolor{currentstroke}{rgb}{0.200000,0.200000,0.200000}%
\pgfsetstrokecolor{currentstroke}%
\pgfsetdash{}{0pt}%
\pgfpathmoveto{\pgfqpoint{2.184531in}{2.451019in}}%
\pgfpathlineto{\pgfqpoint{2.230947in}{2.406512in}}%
\pgfpathlineto{\pgfqpoint{2.262412in}{2.386183in}}%
\pgfpathlineto{\pgfqpoint{2.215991in}{2.430711in}}%
\pgfpathlineto{\pgfqpoint{2.184531in}{2.451019in}}%
\pgfpathclose%
\pgfusepath{stroke,fill}%
\end{pgfscope}%
\begin{pgfscope}%
\pgfpathrectangle{\pgfqpoint{0.050000in}{0.050000in}}{\pgfqpoint{4.900000in}{4.900000in}}%
\pgfusepath{clip}%
\pgfsetbuttcap%
\pgfsetroundjoin%
\definecolor{currentfill}{rgb}{0.915356,0.915356,0.915356}%
\pgfsetfillcolor{currentfill}%
\pgfsetfillopacity{0.900000}%
\pgfsetlinewidth{0.501875pt}%
\definecolor{currentstroke}{rgb}{0.200000,0.200000,0.200000}%
\pgfsetstrokecolor{currentstroke}%
\pgfsetdash{}{0pt}%
\pgfpathmoveto{\pgfqpoint{2.979873in}{1.830557in}}%
\pgfpathlineto{\pgfqpoint{3.024688in}{1.797948in}}%
\pgfpathlineto{\pgfqpoint{3.057399in}{1.778102in}}%
\pgfpathlineto{\pgfqpoint{3.012575in}{1.810469in}}%
\pgfpathlineto{\pgfqpoint{2.979873in}{1.830557in}}%
\pgfpathclose%
\pgfusepath{stroke,fill}%
\end{pgfscope}%
\begin{pgfscope}%
\pgfpathrectangle{\pgfqpoint{0.050000in}{0.050000in}}{\pgfqpoint{4.900000in}{4.900000in}}%
\pgfusepath{clip}%
\pgfsetbuttcap%
\pgfsetroundjoin%
\definecolor{currentfill}{rgb}{1.000000,1.000000,1.000000}%
\pgfsetfillcolor{currentfill}%
\pgfsetfillopacity{0.900000}%
\pgfsetlinewidth{0.501875pt}%
\definecolor{currentstroke}{rgb}{0.200000,0.200000,0.200000}%
\pgfsetstrokecolor{currentstroke}%
\pgfsetdash{}{0pt}%
\pgfpathmoveto{\pgfqpoint{4.391562in}{1.411454in}}%
\pgfpathlineto{\pgfqpoint{4.435112in}{1.423545in}}%
\pgfpathlineto{\pgfqpoint{4.470037in}{1.402450in}}%
\pgfpathlineto{\pgfqpoint{4.426463in}{1.390420in}}%
\pgfpathlineto{\pgfqpoint{4.391562in}{1.411454in}}%
\pgfpathclose%
\pgfusepath{stroke,fill}%
\end{pgfscope}%
\begin{pgfscope}%
\pgfpathrectangle{\pgfqpoint{0.050000in}{0.050000in}}{\pgfqpoint{4.900000in}{4.900000in}}%
\pgfusepath{clip}%
\pgfsetbuttcap%
\pgfsetroundjoin%
\definecolor{currentfill}{rgb}{0.334764,0.334764,0.334764}%
\pgfsetfillcolor{currentfill}%
\pgfsetfillopacity{0.900000}%
\pgfsetlinewidth{0.501875pt}%
\definecolor{currentstroke}{rgb}{0.200000,0.200000,0.200000}%
\pgfsetstrokecolor{currentstroke}%
\pgfsetdash{}{0pt}%
\pgfpathmoveto{\pgfqpoint{1.201602in}{3.246087in}}%
\pgfpathlineto{\pgfqpoint{1.250105in}{3.199876in}}%
\pgfpathlineto{\pgfqpoint{1.280054in}{3.180764in}}%
\pgfpathlineto{\pgfqpoint{1.231538in}{3.227006in}}%
\pgfpathlineto{\pgfqpoint{1.201602in}{3.246087in}}%
\pgfpathclose%
\pgfusepath{stroke,fill}%
\end{pgfscope}%
\begin{pgfscope}%
\pgfpathrectangle{\pgfqpoint{0.050000in}{0.050000in}}{\pgfqpoint{4.900000in}{4.900000in}}%
\pgfusepath{clip}%
\pgfsetbuttcap%
\pgfsetroundjoin%
\definecolor{currentfill}{rgb}{0.990773,0.990773,0.990773}%
\pgfsetfillcolor{currentfill}%
\pgfsetfillopacity{0.900000}%
\pgfsetlinewidth{0.501875pt}%
\definecolor{currentstroke}{rgb}{0.200000,0.200000,0.200000}%
\pgfsetstrokecolor{currentstroke}%
\pgfsetdash{}{0pt}%
\pgfpathmoveto{\pgfqpoint{3.834292in}{1.460916in}}%
\pgfpathlineto{\pgfqpoint{3.878312in}{1.465718in}}%
\pgfpathlineto{\pgfqpoint{3.912388in}{1.445603in}}%
\pgfpathlineto{\pgfqpoint{3.868345in}{1.440879in}}%
\pgfpathlineto{\pgfqpoint{3.834292in}{1.460916in}}%
\pgfpathclose%
\pgfusepath{stroke,fill}%
\end{pgfscope}%
\begin{pgfscope}%
\pgfpathrectangle{\pgfqpoint{0.050000in}{0.050000in}}{\pgfqpoint{4.900000in}{4.900000in}}%
\pgfusepath{clip}%
\pgfsetbuttcap%
\pgfsetroundjoin%
\definecolor{currentfill}{rgb}{0.586082,0.586082,0.586082}%
\pgfsetfillcolor{currentfill}%
\pgfsetfillopacity{0.900000}%
\pgfsetlinewidth{0.501875pt}%
\definecolor{currentstroke}{rgb}{0.200000,0.200000,0.200000}%
\pgfsetstrokecolor{currentstroke}%
\pgfsetdash{}{0pt}%
\pgfpathmoveto{\pgfqpoint{1.919297in}{2.666018in}}%
\pgfpathlineto{\pgfqpoint{1.966284in}{2.621033in}}%
\pgfpathlineto{\pgfqpoint{1.997344in}{2.601024in}}%
\pgfpathlineto{\pgfqpoint{1.950351in}{2.646036in}}%
\pgfpathlineto{\pgfqpoint{1.919297in}{2.666018in}}%
\pgfpathclose%
\pgfusepath{stroke,fill}%
\end{pgfscope}%
\begin{pgfscope}%
\pgfpathrectangle{\pgfqpoint{0.050000in}{0.050000in}}{\pgfqpoint{4.900000in}{4.900000in}}%
\pgfusepath{clip}%
\pgfsetbuttcap%
\pgfsetroundjoin%
\definecolor{currentfill}{rgb}{0.212226,0.212226,0.212226}%
\pgfsetfillcolor{currentfill}%
\pgfsetfillopacity{0.900000}%
\pgfsetlinewidth{0.501875pt}%
\definecolor{currentstroke}{rgb}{0.200000,0.200000,0.200000}%
\pgfsetstrokecolor{currentstroke}%
\pgfsetdash{}{0pt}%
\pgfpathmoveto{\pgfqpoint{0.936175in}{3.461258in}}%
\pgfpathlineto{\pgfqpoint{0.985249in}{3.414584in}}%
\pgfpathlineto{\pgfqpoint{1.014792in}{3.395800in}}%
\pgfpathlineto{\pgfqpoint{0.965705in}{3.442507in}}%
\pgfpathlineto{\pgfqpoint{0.936175in}{3.461258in}}%
\pgfpathclose%
\pgfusepath{stroke,fill}%
\end{pgfscope}%
\begin{pgfscope}%
\pgfpathrectangle{\pgfqpoint{0.050000in}{0.050000in}}{\pgfqpoint{4.900000in}{4.900000in}}%
\pgfusepath{clip}%
\pgfsetbuttcap%
\pgfsetroundjoin%
\definecolor{currentfill}{rgb}{0.836340,0.836340,0.836340}%
\pgfsetfillcolor{currentfill}%
\pgfsetfillopacity{0.900000}%
\pgfsetlinewidth{0.501875pt}%
\definecolor{currentstroke}{rgb}{0.200000,0.200000,0.200000}%
\pgfsetstrokecolor{currentstroke}%
\pgfsetdash{}{0pt}%
\pgfpathmoveto{\pgfqpoint{2.637092in}{2.088254in}}%
\pgfpathlineto{\pgfqpoint{2.682566in}{2.046351in}}%
\pgfpathlineto{\pgfqpoint{2.714742in}{2.025886in}}%
\pgfpathlineto{\pgfqpoint{2.669264in}{2.067661in}}%
\pgfpathlineto{\pgfqpoint{2.637092in}{2.088254in}}%
\pgfpathclose%
\pgfusepath{stroke,fill}%
\end{pgfscope}%
\begin{pgfscope}%
\pgfpathrectangle{\pgfqpoint{0.050000in}{0.050000in}}{\pgfqpoint{4.900000in}{4.900000in}}%
\pgfusepath{clip}%
\pgfsetbuttcap%
\pgfsetroundjoin%
\definecolor{currentfill}{rgb}{0.487043,0.487043,0.487043}%
\pgfsetfillcolor{currentfill}%
\pgfsetfillopacity{0.900000}%
\pgfsetlinewidth{0.501875pt}%
\definecolor{currentstroke}{rgb}{0.200000,0.200000,0.200000}%
\pgfsetstrokecolor{currentstroke}%
\pgfsetdash{}{0pt}%
\pgfpathmoveto{\pgfqpoint{1.653974in}{2.881099in}}%
\pgfpathlineto{\pgfqpoint{1.701531in}{2.835650in}}%
\pgfpathlineto{\pgfqpoint{1.732186in}{2.815970in}}%
\pgfpathlineto{\pgfqpoint{1.684621in}{2.861447in}}%
\pgfpathlineto{\pgfqpoint{1.653974in}{2.881099in}}%
\pgfpathclose%
\pgfusepath{stroke,fill}%
\end{pgfscope}%
\begin{pgfscope}%
\pgfpathrectangle{\pgfqpoint{0.050000in}{0.050000in}}{\pgfqpoint{4.900000in}{4.900000in}}%
\pgfusepath{clip}%
\pgfsetbuttcap%
\pgfsetroundjoin%
\definecolor{currentfill}{rgb}{0.998155,0.998155,0.998155}%
\pgfsetfillcolor{currentfill}%
\pgfsetfillopacity{0.900000}%
\pgfsetlinewidth{0.501875pt}%
\definecolor{currentstroke}{rgb}{0.200000,0.200000,0.200000}%
\pgfsetstrokecolor{currentstroke}%
\pgfsetdash{}{0pt}%
\pgfpathmoveto{\pgfqpoint{4.191056in}{1.418632in}}%
\pgfpathlineto{\pgfqpoint{4.234817in}{1.429687in}}%
\pgfpathlineto{\pgfqpoint{4.269450in}{1.408914in}}%
\pgfpathlineto{\pgfqpoint{4.225663in}{1.397930in}}%
\pgfpathlineto{\pgfqpoint{4.191056in}{1.418632in}}%
\pgfpathclose%
\pgfusepath{stroke,fill}%
\end{pgfscope}%
\begin{pgfscope}%
\pgfpathrectangle{\pgfqpoint{0.050000in}{0.050000in}}{\pgfqpoint{4.900000in}{4.900000in}}%
\pgfusepath{clip}%
\pgfsetbuttcap%
\pgfsetroundjoin%
\definecolor{currentfill}{rgb}{0.760554,0.760554,0.760554}%
\pgfsetfillcolor{currentfill}%
\pgfsetfillopacity{0.900000}%
\pgfsetlinewidth{0.501875pt}%
\definecolor{currentstroke}{rgb}{0.200000,0.200000,0.200000}%
\pgfsetstrokecolor{currentstroke}%
\pgfsetdash{}{0pt}%
\pgfpathmoveto{\pgfqpoint{2.371869in}{2.301014in}}%
\pgfpathlineto{\pgfqpoint{2.417907in}{2.256970in}}%
\pgfpathlineto{\pgfqpoint{2.449674in}{2.236460in}}%
\pgfpathlineto{\pgfqpoint{2.403633in}{2.280495in}}%
\pgfpathlineto{\pgfqpoint{2.371869in}{2.301014in}}%
\pgfpathclose%
\pgfusepath{stroke,fill}%
\end{pgfscope}%
\begin{pgfscope}%
\pgfpathrectangle{\pgfqpoint{0.050000in}{0.050000in}}{\pgfqpoint{4.900000in}{4.900000in}}%
\pgfusepath{clip}%
\pgfsetbuttcap%
\pgfsetroundjoin%
\definecolor{currentfill}{rgb}{0.898378,0.898378,0.898378}%
\pgfsetfillcolor{currentfill}%
\pgfsetfillopacity{0.900000}%
\pgfsetlinewidth{0.501875pt}%
\definecolor{currentstroke}{rgb}{0.200000,0.200000,0.200000}%
\pgfsetstrokecolor{currentstroke}%
\pgfsetdash{}{0pt}%
\pgfpathmoveto{\pgfqpoint{2.902311in}{1.886269in}}%
\pgfpathlineto{\pgfqpoint{2.947274in}{1.850588in}}%
\pgfpathlineto{\pgfqpoint{2.979873in}{1.830557in}}%
\pgfpathlineto{\pgfqpoint{2.934903in}{1.866002in}}%
\pgfpathlineto{\pgfqpoint{2.902311in}{1.886269in}}%
\pgfpathclose%
\pgfusepath{stroke,fill}%
\end{pgfscope}%
\begin{pgfscope}%
\pgfpathrectangle{\pgfqpoint{0.050000in}{0.050000in}}{\pgfqpoint{4.900000in}{4.900000in}}%
\pgfusepath{clip}%
\pgfsetbuttcap%
\pgfsetroundjoin%
\definecolor{currentfill}{rgb}{0.992618,0.992618,0.992618}%
\pgfsetfillcolor{currentfill}%
\pgfsetfillopacity{0.900000}%
\pgfsetlinewidth{0.501875pt}%
\definecolor{currentstroke}{rgb}{0.200000,0.200000,0.200000}%
\pgfsetstrokecolor{currentstroke}%
\pgfsetdash{}{0pt}%
\pgfpathmoveto{\pgfqpoint{3.912388in}{1.445603in}}%
\pgfpathlineto{\pgfqpoint{3.956367in}{1.452208in}}%
\pgfpathlineto{\pgfqpoint{3.990575in}{1.431941in}}%
\pgfpathlineto{\pgfqpoint{3.946572in}{1.425417in}}%
\pgfpathlineto{\pgfqpoint{3.912388in}{1.445603in}}%
\pgfpathclose%
\pgfusepath{stroke,fill}%
\end{pgfscope}%
\begin{pgfscope}%
\pgfpathrectangle{\pgfqpoint{0.050000in}{0.050000in}}{\pgfqpoint{4.900000in}{4.900000in}}%
\pgfusepath{clip}%
\pgfsetbuttcap%
\pgfsetroundjoin%
\definecolor{currentfill}{rgb}{0.395663,0.395663,0.395663}%
\pgfsetfillcolor{currentfill}%
\pgfsetfillopacity{0.900000}%
\pgfsetlinewidth{0.501875pt}%
\definecolor{currentstroke}{rgb}{0.200000,0.200000,0.200000}%
\pgfsetstrokecolor{currentstroke}%
\pgfsetdash{}{0pt}%
\pgfpathmoveto{\pgfqpoint{1.388562in}{3.096254in}}%
\pgfpathlineto{\pgfqpoint{1.436689in}{3.050342in}}%
\pgfpathlineto{\pgfqpoint{1.466939in}{3.030991in}}%
\pgfpathlineto{\pgfqpoint{1.418801in}{3.076933in}}%
\pgfpathlineto{\pgfqpoint{1.388562in}{3.096254in}}%
\pgfpathclose%
\pgfusepath{stroke,fill}%
\end{pgfscope}%
\begin{pgfscope}%
\pgfpathrectangle{\pgfqpoint{0.050000in}{0.050000in}}{\pgfqpoint{4.900000in}{4.900000in}}%
\pgfusepath{clip}%
\pgfsetbuttcap%
\pgfsetroundjoin%
\definecolor{currentfill}{rgb}{0.657809,0.657809,0.657809}%
\pgfsetfillcolor{currentfill}%
\pgfsetfillopacity{0.900000}%
\pgfsetlinewidth{0.501875pt}%
\definecolor{currentstroke}{rgb}{0.200000,0.200000,0.200000}%
\pgfsetstrokecolor{currentstroke}%
\pgfsetdash{}{0pt}%
\pgfpathmoveto{\pgfqpoint{2.106560in}{2.515947in}}%
\pgfpathlineto{\pgfqpoint{2.153169in}{2.471265in}}%
\pgfpathlineto{\pgfqpoint{2.184531in}{2.451019in}}%
\pgfpathlineto{\pgfqpoint{2.137917in}{2.495725in}}%
\pgfpathlineto{\pgfqpoint{2.106560in}{2.515947in}}%
\pgfpathclose%
\pgfusepath{stroke,fill}%
\end{pgfscope}%
\begin{pgfscope}%
\pgfpathrectangle{\pgfqpoint{0.050000in}{0.050000in}}{\pgfqpoint{4.900000in}{4.900000in}}%
\pgfusepath{clip}%
\pgfsetbuttcap%
\pgfsetroundjoin%
\definecolor{currentfill}{rgb}{0.295271,0.295271,0.295271}%
\pgfsetfillcolor{currentfill}%
\pgfsetfillopacity{0.900000}%
\pgfsetlinewidth{0.501875pt}%
\definecolor{currentstroke}{rgb}{0.200000,0.200000,0.200000}%
\pgfsetstrokecolor{currentstroke}%
\pgfsetdash{}{0pt}%
\pgfpathmoveto{\pgfqpoint{1.123059in}{3.311485in}}%
\pgfpathlineto{\pgfqpoint{1.171757in}{3.265109in}}%
\pgfpathlineto{\pgfqpoint{1.201602in}{3.246087in}}%
\pgfpathlineto{\pgfqpoint{1.152892in}{3.292495in}}%
\pgfpathlineto{\pgfqpoint{1.123059in}{3.311485in}}%
\pgfpathclose%
\pgfusepath{stroke,fill}%
\end{pgfscope}%
\begin{pgfscope}%
\pgfpathrectangle{\pgfqpoint{0.050000in}{0.050000in}}{\pgfqpoint{4.900000in}{4.900000in}}%
\pgfusepath{clip}%
\pgfsetbuttcap%
\pgfsetroundjoin%
\definecolor{currentfill}{rgb}{0.966782,0.966782,0.966782}%
\pgfsetfillcolor{currentfill}%
\pgfsetfillopacity{0.900000}%
\pgfsetlinewidth{0.501875pt}%
\definecolor{currentstroke}{rgb}{0.200000,0.200000,0.200000}%
\pgfsetstrokecolor{currentstroke}%
\pgfsetdash{}{0pt}%
\pgfpathmoveto{\pgfqpoint{3.400790in}{1.587689in}}%
\pgfpathlineto{\pgfqpoint{3.445113in}{1.574558in}}%
\pgfpathlineto{\pgfqpoint{3.478526in}{1.555049in}}%
\pgfpathlineto{\pgfqpoint{3.434184in}{1.568107in}}%
\pgfpathlineto{\pgfqpoint{3.400790in}{1.587689in}}%
\pgfpathclose%
\pgfusepath{stroke,fill}%
\end{pgfscope}%
\begin{pgfscope}%
\pgfpathrectangle{\pgfqpoint{0.050000in}{0.050000in}}{\pgfqpoint{4.900000in}{4.900000in}}%
\pgfusepath{clip}%
\pgfsetbuttcap%
\pgfsetroundjoin%
\definecolor{currentfill}{rgb}{0.972318,0.972318,0.972318}%
\pgfsetfillcolor{currentfill}%
\pgfsetfillopacity{0.900000}%
\pgfsetlinewidth{0.501875pt}%
\definecolor{currentstroke}{rgb}{0.200000,0.200000,0.200000}%
\pgfsetstrokecolor{currentstroke}%
\pgfsetdash{}{0pt}%
\pgfpathmoveto{\pgfqpoint{3.478526in}{1.555049in}}%
\pgfpathlineto{\pgfqpoint{3.522789in}{1.545672in}}%
\pgfpathlineto{\pgfqpoint{3.556327in}{1.526110in}}%
\pgfpathlineto{\pgfqpoint{3.512044in}{1.535456in}}%
\pgfpathlineto{\pgfqpoint{3.478526in}{1.555049in}}%
\pgfpathclose%
\pgfusepath{stroke,fill}%
\end{pgfscope}%
\begin{pgfscope}%
\pgfpathrectangle{\pgfqpoint{0.050000in}{0.050000in}}{\pgfqpoint{4.900000in}{4.900000in}}%
\pgfusepath{clip}%
\pgfsetbuttcap%
\pgfsetroundjoin%
\definecolor{currentfill}{rgb}{0.961246,0.961246,0.961246}%
\pgfsetfillcolor{currentfill}%
\pgfsetfillopacity{0.900000}%
\pgfsetlinewidth{0.501875pt}%
\definecolor{currentstroke}{rgb}{0.200000,0.200000,0.200000}%
\pgfsetstrokecolor{currentstroke}%
\pgfsetdash{}{0pt}%
\pgfpathmoveto{\pgfqpoint{3.323106in}{1.624259in}}%
\pgfpathlineto{\pgfqpoint{3.367501in}{1.607189in}}%
\pgfpathlineto{\pgfqpoint{3.400790in}{1.587689in}}%
\pgfpathlineto{\pgfqpoint{3.356378in}{1.604642in}}%
\pgfpathlineto{\pgfqpoint{3.323106in}{1.624259in}}%
\pgfpathclose%
\pgfusepath{stroke,fill}%
\end{pgfscope}%
\begin{pgfscope}%
\pgfpathrectangle{\pgfqpoint{0.050000in}{0.050000in}}{\pgfqpoint{4.900000in}{4.900000in}}%
\pgfusepath{clip}%
\pgfsetbuttcap%
\pgfsetroundjoin%
\definecolor{currentfill}{rgb}{0.555940,0.555940,0.555940}%
\pgfsetfillcolor{currentfill}%
\pgfsetfillopacity{0.900000}%
\pgfsetlinewidth{0.501875pt}%
\definecolor{currentstroke}{rgb}{0.200000,0.200000,0.200000}%
\pgfsetstrokecolor{currentstroke}%
\pgfsetdash{}{0pt}%
\pgfpathmoveto{\pgfqpoint{1.841162in}{2.731086in}}%
\pgfpathlineto{\pgfqpoint{1.888340in}{2.685937in}}%
\pgfpathlineto{\pgfqpoint{1.919297in}{2.666018in}}%
\pgfpathlineto{\pgfqpoint{1.872111in}{2.711194in}}%
\pgfpathlineto{\pgfqpoint{1.841162in}{2.731086in}}%
\pgfpathclose%
\pgfusepath{stroke,fill}%
\end{pgfscope}%
\begin{pgfscope}%
\pgfpathrectangle{\pgfqpoint{0.050000in}{0.050000in}}{\pgfqpoint{4.900000in}{4.900000in}}%
\pgfusepath{clip}%
\pgfsetbuttcap%
\pgfsetroundjoin%
\definecolor{currentfill}{rgb}{0.994464,0.994464,0.994464}%
\pgfsetfillcolor{currentfill}%
\pgfsetfillopacity{0.900000}%
\pgfsetlinewidth{0.501875pt}%
\definecolor{currentstroke}{rgb}{0.200000,0.200000,0.200000}%
\pgfsetstrokecolor{currentstroke}%
\pgfsetdash{}{0pt}%
\pgfpathmoveto{\pgfqpoint{3.990575in}{1.431941in}}%
\pgfpathlineto{\pgfqpoint{4.034511in}{1.440039in}}%
\pgfpathlineto{\pgfqpoint{4.068852in}{1.419623in}}%
\pgfpathlineto{\pgfqpoint{4.024892in}{1.411605in}}%
\pgfpathlineto{\pgfqpoint{3.990575in}{1.431941in}}%
\pgfpathclose%
\pgfusepath{stroke,fill}%
\end{pgfscope}%
\begin{pgfscope}%
\pgfpathrectangle{\pgfqpoint{0.050000in}{0.050000in}}{\pgfqpoint{4.900000in}{4.900000in}}%
\pgfusepath{clip}%
\pgfsetbuttcap%
\pgfsetroundjoin%
\definecolor{currentfill}{rgb}{1.000000,1.000000,1.000000}%
\pgfsetfillcolor{currentfill}%
\pgfsetfillopacity{0.900000}%
\pgfsetlinewidth{0.501875pt}%
\definecolor{currentstroke}{rgb}{0.200000,0.200000,0.200000}%
\pgfsetstrokecolor{currentstroke}%
\pgfsetdash{}{0pt}%
\pgfpathmoveto{\pgfqpoint{4.269450in}{1.408914in}}%
\pgfpathlineto{\pgfqpoint{4.313154in}{1.420478in}}%
\pgfpathlineto{\pgfqpoint{4.347921in}{1.399572in}}%
\pgfpathlineto{\pgfqpoint{4.304192in}{1.388076in}}%
\pgfpathlineto{\pgfqpoint{4.269450in}{1.408914in}}%
\pgfpathclose%
\pgfusepath{stroke,fill}%
\end{pgfscope}%
\begin{pgfscope}%
\pgfpathrectangle{\pgfqpoint{0.050000in}{0.050000in}}{\pgfqpoint{4.900000in}{4.900000in}}%
\pgfusepath{clip}%
\pgfsetbuttcap%
\pgfsetroundjoin%
\definecolor{currentfill}{rgb}{0.977855,0.977855,0.977855}%
\pgfsetfillcolor{currentfill}%
\pgfsetfillopacity{0.900000}%
\pgfsetlinewidth{0.501875pt}%
\definecolor{currentstroke}{rgb}{0.200000,0.200000,0.200000}%
\pgfsetstrokecolor{currentstroke}%
\pgfsetdash{}{0pt}%
\pgfpathmoveto{\pgfqpoint{3.556327in}{1.526110in}}%
\pgfpathlineto{\pgfqpoint{3.600539in}{1.520223in}}%
\pgfpathlineto{\pgfqpoint{3.634203in}{1.500575in}}%
\pgfpathlineto{\pgfqpoint{3.589971in}{1.506466in}}%
\pgfpathlineto{\pgfqpoint{3.556327in}{1.526110in}}%
\pgfpathclose%
\pgfusepath{stroke,fill}%
\end{pgfscope}%
\begin{pgfscope}%
\pgfpathrectangle{\pgfqpoint{0.050000in}{0.050000in}}{\pgfqpoint{4.900000in}{4.900000in}}%
\pgfusepath{clip}%
\pgfsetbuttcap%
\pgfsetroundjoin%
\definecolor{currentfill}{rgb}{0.812226,0.812226,0.812226}%
\pgfsetfillcolor{currentfill}%
\pgfsetfillopacity{0.900000}%
\pgfsetlinewidth{0.501875pt}%
\definecolor{currentstroke}{rgb}{0.200000,0.200000,0.200000}%
\pgfsetstrokecolor{currentstroke}%
\pgfsetdash{}{0pt}%
\pgfpathmoveto{\pgfqpoint{2.559360in}{2.151722in}}%
\pgfpathlineto{\pgfqpoint{2.605020in}{2.108801in}}%
\pgfpathlineto{\pgfqpoint{2.637092in}{2.088254in}}%
\pgfpathlineto{\pgfqpoint{2.591428in}{2.131094in}}%
\pgfpathlineto{\pgfqpoint{2.559360in}{2.151722in}}%
\pgfpathclose%
\pgfusepath{stroke,fill}%
\end{pgfscope}%
\begin{pgfscope}%
\pgfpathrectangle{\pgfqpoint{0.050000in}{0.050000in}}{\pgfqpoint{4.900000in}{4.900000in}}%
\pgfusepath{clip}%
\pgfsetbuttcap%
\pgfsetroundjoin%
\definecolor{currentfill}{rgb}{0.952018,0.952018,0.952018}%
\pgfsetfillcolor{currentfill}%
\pgfsetfillopacity{0.900000}%
\pgfsetlinewidth{0.501875pt}%
\definecolor{currentstroke}{rgb}{0.200000,0.200000,0.200000}%
\pgfsetstrokecolor{currentstroke}%
\pgfsetdash{}{0pt}%
\pgfpathmoveto{\pgfqpoint{3.245460in}{1.664885in}}%
\pgfpathlineto{\pgfqpoint{3.289939in}{1.643795in}}%
\pgfpathlineto{\pgfqpoint{3.323106in}{1.624259in}}%
\pgfpathlineto{\pgfqpoint{3.278612in}{1.645188in}}%
\pgfpathlineto{\pgfqpoint{3.245460in}{1.664885in}}%
\pgfpathclose%
\pgfusepath{stroke,fill}%
\end{pgfscope}%
\begin{pgfscope}%
\pgfpathrectangle{\pgfqpoint{0.050000in}{0.050000in}}{\pgfqpoint{4.900000in}{4.900000in}}%
\pgfusepath{clip}%
\pgfsetbuttcap%
\pgfsetroundjoin%
\definecolor{currentfill}{rgb}{0.878570,0.878570,0.878570}%
\pgfsetfillcolor{currentfill}%
\pgfsetfillopacity{0.900000}%
\pgfsetlinewidth{0.501875pt}%
\definecolor{currentstroke}{rgb}{0.200000,0.200000,0.200000}%
\pgfsetstrokecolor{currentstroke}%
\pgfsetdash{}{0pt}%
\pgfpathmoveto{\pgfqpoint{2.824697in}{1.944728in}}%
\pgfpathlineto{\pgfqpoint{2.869822in}{1.906486in}}%
\pgfpathlineto{\pgfqpoint{2.902311in}{1.886269in}}%
\pgfpathlineto{\pgfqpoint{2.857181in}{1.924299in}}%
\pgfpathlineto{\pgfqpoint{2.824697in}{1.944728in}}%
\pgfpathclose%
\pgfusepath{stroke,fill}%
\end{pgfscope}%
\begin{pgfscope}%
\pgfpathrectangle{\pgfqpoint{0.050000in}{0.050000in}}{\pgfqpoint{4.900000in}{4.900000in}}%
\pgfusepath{clip}%
\pgfsetbuttcap%
\pgfsetroundjoin%
\definecolor{currentfill}{rgb}{0.456901,0.456901,0.456901}%
\pgfsetfillcolor{currentfill}%
\pgfsetfillopacity{0.900000}%
\pgfsetlinewidth{0.501875pt}%
\definecolor{currentstroke}{rgb}{0.200000,0.200000,0.200000}%
\pgfsetstrokecolor{currentstroke}%
\pgfsetdash{}{0pt}%
\pgfpathmoveto{\pgfqpoint{1.575673in}{2.946302in}}%
\pgfpathlineto{\pgfqpoint{1.623423in}{2.900689in}}%
\pgfpathlineto{\pgfqpoint{1.653974in}{2.881099in}}%
\pgfpathlineto{\pgfqpoint{1.606216in}{2.926740in}}%
\pgfpathlineto{\pgfqpoint{1.575673in}{2.946302in}}%
\pgfpathclose%
\pgfusepath{stroke,fill}%
\end{pgfscope}%
\begin{pgfscope}%
\pgfpathrectangle{\pgfqpoint{0.050000in}{0.050000in}}{\pgfqpoint{4.900000in}{4.900000in}}%
\pgfusepath{clip}%
\pgfsetbuttcap%
\pgfsetroundjoin%
\definecolor{currentfill}{rgb}{0.981546,0.981546,0.981546}%
\pgfsetfillcolor{currentfill}%
\pgfsetfillopacity{0.900000}%
\pgfsetlinewidth{0.501875pt}%
\definecolor{currentstroke}{rgb}{0.200000,0.200000,0.200000}%
\pgfsetstrokecolor{currentstroke}%
\pgfsetdash{}{0pt}%
\pgfpathmoveto{\pgfqpoint{3.634203in}{1.500575in}}%
\pgfpathlineto{\pgfqpoint{3.678369in}{1.497862in}}%
\pgfpathlineto{\pgfqpoint{3.712162in}{1.478100in}}%
\pgfpathlineto{\pgfqpoint{3.667975in}{1.480846in}}%
\pgfpathlineto{\pgfqpoint{3.634203in}{1.500575in}}%
\pgfpathclose%
\pgfusepath{stroke,fill}%
\end{pgfscope}%
\begin{pgfscope}%
\pgfpathrectangle{\pgfqpoint{0.050000in}{0.050000in}}{\pgfqpoint{4.900000in}{4.900000in}}%
\pgfusepath{clip}%
\pgfsetbuttcap%
\pgfsetroundjoin%
\definecolor{currentfill}{rgb}{0.729781,0.729781,0.729781}%
\pgfsetfillcolor{currentfill}%
\pgfsetfillopacity{0.900000}%
\pgfsetlinewidth{0.501875pt}%
\definecolor{currentstroke}{rgb}{0.200000,0.200000,0.200000}%
\pgfsetstrokecolor{currentstroke}%
\pgfsetdash{}{0pt}%
\pgfpathmoveto{\pgfqpoint{2.293975in}{2.365794in}}%
\pgfpathlineto{\pgfqpoint{2.340204in}{2.321474in}}%
\pgfpathlineto{\pgfqpoint{2.371869in}{2.301014in}}%
\pgfpathlineto{\pgfqpoint{2.325635in}{2.345343in}}%
\pgfpathlineto{\pgfqpoint{2.293975in}{2.365794in}}%
\pgfpathclose%
\pgfusepath{stroke,fill}%
\end{pgfscope}%
\begin{pgfscope}%
\pgfpathrectangle{\pgfqpoint{0.050000in}{0.050000in}}{\pgfqpoint{4.900000in}{4.900000in}}%
\pgfusepath{clip}%
\pgfsetbuttcap%
\pgfsetroundjoin%
\definecolor{currentfill}{rgb}{0.944637,0.944637,0.944637}%
\pgfsetfillcolor{currentfill}%
\pgfsetfillopacity{0.900000}%
\pgfsetlinewidth{0.501875pt}%
\definecolor{currentstroke}{rgb}{0.200000,0.200000,0.200000}%
\pgfsetstrokecolor{currentstroke}%
\pgfsetdash{}{0pt}%
\pgfpathmoveto{\pgfqpoint{3.167835in}{1.709574in}}%
\pgfpathlineto{\pgfqpoint{3.212413in}{1.684506in}}%
\pgfpathlineto{\pgfqpoint{3.245460in}{1.664885in}}%
\pgfpathlineto{\pgfqpoint{3.200869in}{1.689756in}}%
\pgfpathlineto{\pgfqpoint{3.167835in}{1.709574in}}%
\pgfpathclose%
\pgfusepath{stroke,fill}%
\end{pgfscope}%
\begin{pgfscope}%
\pgfpathrectangle{\pgfqpoint{0.050000in}{0.050000in}}{\pgfqpoint{4.900000in}{4.900000in}}%
\pgfusepath{clip}%
\pgfsetbuttcap%
\pgfsetroundjoin%
\definecolor{currentfill}{rgb}{0.363183,0.363183,0.363183}%
\pgfsetfillcolor{currentfill}%
\pgfsetfillopacity{0.900000}%
\pgfsetlinewidth{0.501875pt}%
\definecolor{currentstroke}{rgb}{0.200000,0.200000,0.200000}%
\pgfsetstrokecolor{currentstroke}%
\pgfsetdash{}{0pt}%
\pgfpathmoveto{\pgfqpoint{1.310095in}{3.161593in}}%
\pgfpathlineto{\pgfqpoint{1.358416in}{3.115516in}}%
\pgfpathlineto{\pgfqpoint{1.388562in}{3.096254in}}%
\pgfpathlineto{\pgfqpoint{1.340230in}{3.142362in}}%
\pgfpathlineto{\pgfqpoint{1.310095in}{3.161593in}}%
\pgfpathclose%
\pgfusepath{stroke,fill}%
\end{pgfscope}%
\begin{pgfscope}%
\pgfpathrectangle{\pgfqpoint{0.050000in}{0.050000in}}{\pgfqpoint{4.900000in}{4.900000in}}%
\pgfusepath{clip}%
\pgfsetbuttcap%
\pgfsetroundjoin%
\definecolor{currentfill}{rgb}{0.624221,0.624221,0.624221}%
\pgfsetfillcolor{currentfill}%
\pgfsetfillopacity{0.900000}%
\pgfsetlinewidth{0.501875pt}%
\definecolor{currentstroke}{rgb}{0.200000,0.200000,0.200000}%
\pgfsetstrokecolor{currentstroke}%
\pgfsetdash{}{0pt}%
\pgfpathmoveto{\pgfqpoint{2.028500in}{2.580954in}}%
\pgfpathlineto{\pgfqpoint{2.075301in}{2.536106in}}%
\pgfpathlineto{\pgfqpoint{2.106560in}{2.515947in}}%
\pgfpathlineto{\pgfqpoint{2.059754in}{2.560821in}}%
\pgfpathlineto{\pgfqpoint{2.028500in}{2.580954in}}%
\pgfpathclose%
\pgfusepath{stroke,fill}%
\end{pgfscope}%
\begin{pgfscope}%
\pgfpathrectangle{\pgfqpoint{0.050000in}{0.050000in}}{\pgfqpoint{4.900000in}{4.900000in}}%
\pgfusepath{clip}%
\pgfsetbuttcap%
\pgfsetroundjoin%
\definecolor{currentfill}{rgb}{0.996309,0.996309,0.996309}%
\pgfsetfillcolor{currentfill}%
\pgfsetfillopacity{0.900000}%
\pgfsetlinewidth{0.501875pt}%
\definecolor{currentstroke}{rgb}{0.200000,0.200000,0.200000}%
\pgfsetstrokecolor{currentstroke}%
\pgfsetdash{}{0pt}%
\pgfpathmoveto{\pgfqpoint{4.068852in}{1.419623in}}%
\pgfpathlineto{\pgfqpoint{4.112742in}{1.428932in}}%
\pgfpathlineto{\pgfqpoint{4.147215in}{1.408370in}}%
\pgfpathlineto{\pgfqpoint{4.103301in}{1.399140in}}%
\pgfpathlineto{\pgfqpoint{4.068852in}{1.419623in}}%
\pgfpathclose%
\pgfusepath{stroke,fill}%
\end{pgfscope}%
\begin{pgfscope}%
\pgfpathrectangle{\pgfqpoint{0.050000in}{0.050000in}}{\pgfqpoint{4.900000in}{4.900000in}}%
\pgfusepath{clip}%
\pgfsetbuttcap%
\pgfsetroundjoin%
\definecolor{currentfill}{rgb}{0.985236,0.985236,0.985236}%
\pgfsetfillcolor{currentfill}%
\pgfsetfillopacity{0.900000}%
\pgfsetlinewidth{0.501875pt}%
\definecolor{currentstroke}{rgb}{0.200000,0.200000,0.200000}%
\pgfsetstrokecolor{currentstroke}%
\pgfsetdash{}{0pt}%
\pgfpathmoveto{\pgfqpoint{3.712162in}{1.478100in}}%
\pgfpathlineto{\pgfqpoint{3.756286in}{1.478217in}}%
\pgfpathlineto{\pgfqpoint{3.790208in}{1.458323in}}%
\pgfpathlineto{\pgfqpoint{3.746062in}{1.458261in}}%
\pgfpathlineto{\pgfqpoint{3.712162in}{1.478100in}}%
\pgfpathclose%
\pgfusepath{stroke,fill}%
\end{pgfscope}%
\begin{pgfscope}%
\pgfpathrectangle{\pgfqpoint{0.050000in}{0.050000in}}{\pgfqpoint{4.900000in}{4.900000in}}%
\pgfusepath{clip}%
\pgfsetbuttcap%
\pgfsetroundjoin%
\definecolor{currentfill}{rgb}{1.000000,1.000000,1.000000}%
\pgfsetfillcolor{currentfill}%
\pgfsetfillopacity{0.900000}%
\pgfsetlinewidth{0.501875pt}%
\definecolor{currentstroke}{rgb}{0.200000,0.200000,0.200000}%
\pgfsetstrokecolor{currentstroke}%
\pgfsetdash{}{0pt}%
\pgfpathmoveto{\pgfqpoint{4.347921in}{1.399572in}}%
\pgfpathlineto{\pgfqpoint{4.391562in}{1.411454in}}%
\pgfpathlineto{\pgfqpoint{4.426463in}{1.390420in}}%
\pgfpathlineto{\pgfqpoint{4.382797in}{1.378601in}}%
\pgfpathlineto{\pgfqpoint{4.347921in}{1.399572in}}%
\pgfpathclose%
\pgfusepath{stroke,fill}%
\end{pgfscope}%
\begin{pgfscope}%
\pgfpathrectangle{\pgfqpoint{0.050000in}{0.050000in}}{\pgfqpoint{4.900000in}{4.900000in}}%
\pgfusepath{clip}%
\pgfsetbuttcap%
\pgfsetroundjoin%
\definecolor{currentfill}{rgb}{0.256517,0.256517,0.256517}%
\pgfsetfillcolor{currentfill}%
\pgfsetfillopacity{0.900000}%
\pgfsetlinewidth{0.501875pt}%
\definecolor{currentstroke}{rgb}{0.200000,0.200000,0.200000}%
\pgfsetstrokecolor{currentstroke}%
\pgfsetdash{}{0pt}%
\pgfpathmoveto{\pgfqpoint{1.044427in}{3.376959in}}%
\pgfpathlineto{\pgfqpoint{1.093320in}{3.330417in}}%
\pgfpathlineto{\pgfqpoint{1.123059in}{3.311485in}}%
\pgfpathlineto{\pgfqpoint{1.074154in}{3.358059in}}%
\pgfpathlineto{\pgfqpoint{1.044427in}{3.376959in}}%
\pgfpathclose%
\pgfusepath{stroke,fill}%
\end{pgfscope}%
\begin{pgfscope}%
\pgfpathrectangle{\pgfqpoint{0.050000in}{0.050000in}}{\pgfqpoint{4.900000in}{4.900000in}}%
\pgfusepath{clip}%
\pgfsetbuttcap%
\pgfsetroundjoin%
\definecolor{currentfill}{rgb}{0.932334,0.932334,0.932334}%
\pgfsetfillcolor{currentfill}%
\pgfsetfillopacity{0.900000}%
\pgfsetlinewidth{0.501875pt}%
\definecolor{currentstroke}{rgb}{0.200000,0.200000,0.200000}%
\pgfsetstrokecolor{currentstroke}%
\pgfsetdash{}{0pt}%
\pgfpathmoveto{\pgfqpoint{3.090213in}{1.758193in}}%
\pgfpathlineto{\pgfqpoint{3.134906in}{1.729323in}}%
\pgfpathlineto{\pgfqpoint{3.167835in}{1.709574in}}%
\pgfpathlineto{\pgfqpoint{3.123131in}{1.738221in}}%
\pgfpathlineto{\pgfqpoint{3.090213in}{1.758193in}}%
\pgfpathclose%
\pgfusepath{stroke,fill}%
\end{pgfscope}%
\begin{pgfscope}%
\pgfpathrectangle{\pgfqpoint{0.050000in}{0.050000in}}{\pgfqpoint{4.900000in}{4.900000in}}%
\pgfusepath{clip}%
\pgfsetbuttcap%
\pgfsetroundjoin%
\definecolor{currentfill}{rgb}{0.521492,0.521492,0.521492}%
\pgfsetfillcolor{currentfill}%
\pgfsetfillopacity{0.900000}%
\pgfsetlinewidth{0.501875pt}%
\definecolor{currentstroke}{rgb}{0.200000,0.200000,0.200000}%
\pgfsetstrokecolor{currentstroke}%
\pgfsetdash{}{0pt}%
\pgfpathmoveto{\pgfqpoint{1.762936in}{2.796229in}}%
\pgfpathlineto{\pgfqpoint{1.810308in}{2.750917in}}%
\pgfpathlineto{\pgfqpoint{1.841162in}{2.731086in}}%
\pgfpathlineto{\pgfqpoint{1.793782in}{2.776427in}}%
\pgfpathlineto{\pgfqpoint{1.762936in}{2.796229in}}%
\pgfpathclose%
\pgfusepath{stroke,fill}%
\end{pgfscope}%
\begin{pgfscope}%
\pgfpathrectangle{\pgfqpoint{0.050000in}{0.050000in}}{\pgfqpoint{4.900000in}{4.900000in}}%
\pgfusepath{clip}%
\pgfsetbuttcap%
\pgfsetroundjoin%
\definecolor{currentfill}{rgb}{0.858762,0.858762,0.858762}%
\pgfsetfillcolor{currentfill}%
\pgfsetfillopacity{0.900000}%
\pgfsetlinewidth{0.501875pt}%
\definecolor{currentstroke}{rgb}{0.200000,0.200000,0.200000}%
\pgfsetstrokecolor{currentstroke}%
\pgfsetdash{}{0pt}%
\pgfpathmoveto{\pgfqpoint{2.747018in}{2.005375in}}%
\pgfpathlineto{\pgfqpoint{2.792316in}{1.965110in}}%
\pgfpathlineto{\pgfqpoint{2.824697in}{1.944728in}}%
\pgfpathlineto{\pgfqpoint{2.779395in}{1.984818in}}%
\pgfpathlineto{\pgfqpoint{2.747018in}{2.005375in}}%
\pgfpathclose%
\pgfusepath{stroke,fill}%
\end{pgfscope}%
\begin{pgfscope}%
\pgfpathrectangle{\pgfqpoint{0.050000in}{0.050000in}}{\pgfqpoint{4.900000in}{4.900000in}}%
\pgfusepath{clip}%
\pgfsetbuttcap%
\pgfsetroundjoin%
\definecolor{currentfill}{rgb}{0.784667,0.784667,0.784667}%
\pgfsetfillcolor{currentfill}%
\pgfsetfillopacity{0.900000}%
\pgfsetlinewidth{0.501875pt}%
\definecolor{currentstroke}{rgb}{0.200000,0.200000,0.200000}%
\pgfsetstrokecolor{currentstroke}%
\pgfsetdash{}{0pt}%
\pgfpathmoveto{\pgfqpoint{2.481540in}{2.215896in}}%
\pgfpathlineto{\pgfqpoint{2.527391in}{2.172299in}}%
\pgfpathlineto{\pgfqpoint{2.559360in}{2.151722in}}%
\pgfpathlineto{\pgfqpoint{2.513506in}{2.195276in}}%
\pgfpathlineto{\pgfqpoint{2.481540in}{2.215896in}}%
\pgfpathclose%
\pgfusepath{stroke,fill}%
\end{pgfscope}%
\begin{pgfscope}%
\pgfpathrectangle{\pgfqpoint{0.050000in}{0.050000in}}{\pgfqpoint{4.900000in}{4.900000in}}%
\pgfusepath{clip}%
\pgfsetbuttcap%
\pgfsetroundjoin%
\definecolor{currentfill}{rgb}{0.988927,0.988927,0.988927}%
\pgfsetfillcolor{currentfill}%
\pgfsetfillopacity{0.900000}%
\pgfsetlinewidth{0.501875pt}%
\definecolor{currentstroke}{rgb}{0.200000,0.200000,0.200000}%
\pgfsetstrokecolor{currentstroke}%
\pgfsetdash{}{0pt}%
\pgfpathmoveto{\pgfqpoint{3.790208in}{1.458323in}}%
\pgfpathlineto{\pgfqpoint{3.834292in}{1.460916in}}%
\pgfpathlineto{\pgfqpoint{3.868345in}{1.440879in}}%
\pgfpathlineto{\pgfqpoint{3.824238in}{1.438354in}}%
\pgfpathlineto{\pgfqpoint{3.790208in}{1.458323in}}%
\pgfpathclose%
\pgfusepath{stroke,fill}%
\end{pgfscope}%
\begin{pgfscope}%
\pgfpathrectangle{\pgfqpoint{0.050000in}{0.050000in}}{\pgfqpoint{4.900000in}{4.900000in}}%
\pgfusepath{clip}%
\pgfsetbuttcap%
\pgfsetroundjoin%
\definecolor{currentfill}{rgb}{0.424083,0.424083,0.424083}%
\pgfsetfillcolor{currentfill}%
\pgfsetfillopacity{0.900000}%
\pgfsetlinewidth{0.501875pt}%
\definecolor{currentstroke}{rgb}{0.200000,0.200000,0.200000}%
\pgfsetstrokecolor{currentstroke}%
\pgfsetdash{}{0pt}%
\pgfpathmoveto{\pgfqpoint{1.497282in}{3.011580in}}%
\pgfpathlineto{\pgfqpoint{1.545225in}{2.965803in}}%
\pgfpathlineto{\pgfqpoint{1.575673in}{2.946302in}}%
\pgfpathlineto{\pgfqpoint{1.527720in}{2.992109in}}%
\pgfpathlineto{\pgfqpoint{1.497282in}{3.011580in}}%
\pgfpathclose%
\pgfusepath{stroke,fill}%
\end{pgfscope}%
\begin{pgfscope}%
\pgfpathrectangle{\pgfqpoint{0.050000in}{0.050000in}}{\pgfqpoint{4.900000in}{4.900000in}}%
\pgfusepath{clip}%
\pgfsetbuttcap%
\pgfsetroundjoin%
\definecolor{currentfill}{rgb}{0.696194,0.696194,0.696194}%
\pgfsetfillcolor{currentfill}%
\pgfsetfillopacity{0.900000}%
\pgfsetlinewidth{0.501875pt}%
\definecolor{currentstroke}{rgb}{0.200000,0.200000,0.200000}%
\pgfsetstrokecolor{currentstroke}%
\pgfsetdash{}{0pt}%
\pgfpathmoveto{\pgfqpoint{2.215991in}{2.430711in}}%
\pgfpathlineto{\pgfqpoint{2.262412in}{2.386183in}}%
\pgfpathlineto{\pgfqpoint{2.293975in}{2.365794in}}%
\pgfpathlineto{\pgfqpoint{2.247548in}{2.410341in}}%
\pgfpathlineto{\pgfqpoint{2.215991in}{2.430711in}}%
\pgfpathclose%
\pgfusepath{stroke,fill}%
\end{pgfscope}%
\begin{pgfscope}%
\pgfpathrectangle{\pgfqpoint{0.050000in}{0.050000in}}{\pgfqpoint{4.900000in}{4.900000in}}%
\pgfusepath{clip}%
\pgfsetbuttcap%
\pgfsetroundjoin%
\definecolor{currentfill}{rgb}{0.998155,0.998155,0.998155}%
\pgfsetfillcolor{currentfill}%
\pgfsetfillopacity{0.900000}%
\pgfsetlinewidth{0.501875pt}%
\definecolor{currentstroke}{rgb}{0.200000,0.200000,0.200000}%
\pgfsetstrokecolor{currentstroke}%
\pgfsetdash{}{0pt}%
\pgfpathmoveto{\pgfqpoint{4.147215in}{1.408370in}}%
\pgfpathlineto{\pgfqpoint{4.191056in}{1.418632in}}%
\pgfpathlineto{\pgfqpoint{4.225663in}{1.397930in}}%
\pgfpathlineto{\pgfqpoint{4.181798in}{1.387742in}}%
\pgfpathlineto{\pgfqpoint{4.147215in}{1.408370in}}%
\pgfpathclose%
\pgfusepath{stroke,fill}%
\end{pgfscope}%
\begin{pgfscope}%
\pgfpathrectangle{\pgfqpoint{0.050000in}{0.050000in}}{\pgfqpoint{4.900000in}{4.900000in}}%
\pgfusepath{clip}%
\pgfsetbuttcap%
\pgfsetroundjoin%
\definecolor{currentfill}{rgb}{0.915356,0.915356,0.915356}%
\pgfsetfillcolor{currentfill}%
\pgfsetfillopacity{0.900000}%
\pgfsetlinewidth{0.501875pt}%
\definecolor{currentstroke}{rgb}{0.200000,0.200000,0.200000}%
\pgfsetstrokecolor{currentstroke}%
\pgfsetdash{}{0pt}%
\pgfpathmoveto{\pgfqpoint{3.012575in}{1.810469in}}%
\pgfpathlineto{\pgfqpoint{3.057399in}{1.778102in}}%
\pgfpathlineto{\pgfqpoint{3.090213in}{1.758193in}}%
\pgfpathlineto{\pgfqpoint{3.045379in}{1.790323in}}%
\pgfpathlineto{\pgfqpoint{3.012575in}{1.810469in}}%
\pgfpathclose%
\pgfusepath{stroke,fill}%
\end{pgfscope}%
\begin{pgfscope}%
\pgfpathrectangle{\pgfqpoint{0.050000in}{0.050000in}}{\pgfqpoint{4.900000in}{4.900000in}}%
\pgfusepath{clip}%
\pgfsetbuttcap%
\pgfsetroundjoin%
\definecolor{currentfill}{rgb}{0.334764,0.334764,0.334764}%
\pgfsetfillcolor{currentfill}%
\pgfsetfillopacity{0.900000}%
\pgfsetlinewidth{0.501875pt}%
\definecolor{currentstroke}{rgb}{0.200000,0.200000,0.200000}%
\pgfsetstrokecolor{currentstroke}%
\pgfsetdash{}{0pt}%
\pgfpathmoveto{\pgfqpoint{1.231538in}{3.227006in}}%
\pgfpathlineto{\pgfqpoint{1.280054in}{3.180764in}}%
\pgfpathlineto{\pgfqpoint{1.310095in}{3.161593in}}%
\pgfpathlineto{\pgfqpoint{1.261569in}{3.207866in}}%
\pgfpathlineto{\pgfqpoint{1.231538in}{3.227006in}}%
\pgfpathclose%
\pgfusepath{stroke,fill}%
\end{pgfscope}%
\begin{pgfscope}%
\pgfpathrectangle{\pgfqpoint{0.050000in}{0.050000in}}{\pgfqpoint{4.900000in}{4.900000in}}%
\pgfusepath{clip}%
\pgfsetbuttcap%
\pgfsetroundjoin%
\definecolor{currentfill}{rgb}{0.990773,0.990773,0.990773}%
\pgfsetfillcolor{currentfill}%
\pgfsetfillopacity{0.900000}%
\pgfsetlinewidth{0.501875pt}%
\definecolor{currentstroke}{rgb}{0.200000,0.200000,0.200000}%
\pgfsetstrokecolor{currentstroke}%
\pgfsetdash{}{0pt}%
\pgfpathmoveto{\pgfqpoint{3.868345in}{1.440879in}}%
\pgfpathlineto{\pgfqpoint{3.912388in}{1.445603in}}%
\pgfpathlineto{\pgfqpoint{3.946572in}{1.425417in}}%
\pgfpathlineto{\pgfqpoint{3.902505in}{1.420770in}}%
\pgfpathlineto{\pgfqpoint{3.868345in}{1.440879in}}%
\pgfpathclose%
\pgfusepath{stroke,fill}%
\end{pgfscope}%
\begin{pgfscope}%
\pgfpathrectangle{\pgfqpoint{0.050000in}{0.050000in}}{\pgfqpoint{4.900000in}{4.900000in}}%
\pgfusepath{clip}%
\pgfsetbuttcap%
\pgfsetroundjoin%
\definecolor{currentfill}{rgb}{0.586082,0.586082,0.586082}%
\pgfsetfillcolor{currentfill}%
\pgfsetfillopacity{0.900000}%
\pgfsetlinewidth{0.501875pt}%
\definecolor{currentstroke}{rgb}{0.200000,0.200000,0.200000}%
\pgfsetstrokecolor{currentstroke}%
\pgfsetdash{}{0pt}%
\pgfpathmoveto{\pgfqpoint{1.950351in}{2.646036in}}%
\pgfpathlineto{\pgfqpoint{1.997344in}{2.601024in}}%
\pgfpathlineto{\pgfqpoint{2.028500in}{2.580954in}}%
\pgfpathlineto{\pgfqpoint{1.981501in}{2.625993in}}%
\pgfpathlineto{\pgfqpoint{1.950351in}{2.646036in}}%
\pgfpathclose%
\pgfusepath{stroke,fill}%
\end{pgfscope}%
\begin{pgfscope}%
\pgfpathrectangle{\pgfqpoint{0.050000in}{0.050000in}}{\pgfqpoint{4.900000in}{4.900000in}}%
\pgfusepath{clip}%
\pgfsetbuttcap%
\pgfsetroundjoin%
\definecolor{currentfill}{rgb}{0.212226,0.212226,0.212226}%
\pgfsetfillcolor{currentfill}%
\pgfsetfillopacity{0.900000}%
\pgfsetlinewidth{0.501875pt}%
\definecolor{currentstroke}{rgb}{0.200000,0.200000,0.200000}%
\pgfsetstrokecolor{currentstroke}%
\pgfsetdash{}{0pt}%
\pgfpathmoveto{\pgfqpoint{0.965705in}{3.442507in}}%
\pgfpathlineto{\pgfqpoint{1.014792in}{3.395800in}}%
\pgfpathlineto{\pgfqpoint{1.044427in}{3.376959in}}%
\pgfpathlineto{\pgfqpoint{0.995326in}{3.423698in}}%
\pgfpathlineto{\pgfqpoint{0.965705in}{3.442507in}}%
\pgfpathclose%
\pgfusepath{stroke,fill}%
\end{pgfscope}%
\begin{pgfscope}%
\pgfpathrectangle{\pgfqpoint{0.050000in}{0.050000in}}{\pgfqpoint{4.900000in}{4.900000in}}%
\pgfusepath{clip}%
\pgfsetbuttcap%
\pgfsetroundjoin%
\definecolor{currentfill}{rgb}{0.836340,0.836340,0.836340}%
\pgfsetfillcolor{currentfill}%
\pgfsetfillopacity{0.900000}%
\pgfsetlinewidth{0.501875pt}%
\definecolor{currentstroke}{rgb}{0.200000,0.200000,0.200000}%
\pgfsetstrokecolor{currentstroke}%
\pgfsetdash{}{0pt}%
\pgfpathmoveto{\pgfqpoint{2.669264in}{2.067661in}}%
\pgfpathlineto{\pgfqpoint{2.714742in}{2.025886in}}%
\pgfpathlineto{\pgfqpoint{2.747018in}{2.005375in}}%
\pgfpathlineto{\pgfqpoint{2.701537in}{2.047019in}}%
\pgfpathlineto{\pgfqpoint{2.669264in}{2.067661in}}%
\pgfpathclose%
\pgfusepath{stroke,fill}%
\end{pgfscope}%
\begin{pgfscope}%
\pgfpathrectangle{\pgfqpoint{0.050000in}{0.050000in}}{\pgfqpoint{4.900000in}{4.900000in}}%
\pgfusepath{clip}%
\pgfsetbuttcap%
\pgfsetroundjoin%
\definecolor{currentfill}{rgb}{0.491349,0.491349,0.491349}%
\pgfsetfillcolor{currentfill}%
\pgfsetfillopacity{0.900000}%
\pgfsetlinewidth{0.501875pt}%
\definecolor{currentstroke}{rgb}{0.200000,0.200000,0.200000}%
\pgfsetstrokecolor{currentstroke}%
\pgfsetdash{}{0pt}%
\pgfpathmoveto{\pgfqpoint{1.684621in}{2.861447in}}%
\pgfpathlineto{\pgfqpoint{1.732186in}{2.815970in}}%
\pgfpathlineto{\pgfqpoint{1.762936in}{2.796229in}}%
\pgfpathlineto{\pgfqpoint{1.715363in}{2.841735in}}%
\pgfpathlineto{\pgfqpoint{1.684621in}{2.861447in}}%
\pgfpathclose%
\pgfusepath{stroke,fill}%
\end{pgfscope}%
\begin{pgfscope}%
\pgfpathrectangle{\pgfqpoint{0.050000in}{0.050000in}}{\pgfqpoint{4.900000in}{4.900000in}}%
\pgfusepath{clip}%
\pgfsetbuttcap%
\pgfsetroundjoin%
\definecolor{currentfill}{rgb}{0.760554,0.760554,0.760554}%
\pgfsetfillcolor{currentfill}%
\pgfsetfillopacity{0.900000}%
\pgfsetlinewidth{0.501875pt}%
\definecolor{currentstroke}{rgb}{0.200000,0.200000,0.200000}%
\pgfsetstrokecolor{currentstroke}%
\pgfsetdash{}{0pt}%
\pgfpathmoveto{\pgfqpoint{2.403633in}{2.280495in}}%
\pgfpathlineto{\pgfqpoint{2.449674in}{2.236460in}}%
\pgfpathlineto{\pgfqpoint{2.481540in}{2.215896in}}%
\pgfpathlineto{\pgfqpoint{2.435495in}{2.259918in}}%
\pgfpathlineto{\pgfqpoint{2.403633in}{2.280495in}}%
\pgfpathclose%
\pgfusepath{stroke,fill}%
\end{pgfscope}%
\begin{pgfscope}%
\pgfpathrectangle{\pgfqpoint{0.050000in}{0.050000in}}{\pgfqpoint{4.900000in}{4.900000in}}%
\pgfusepath{clip}%
\pgfsetbuttcap%
\pgfsetroundjoin%
\definecolor{currentfill}{rgb}{0.998155,0.998155,0.998155}%
\pgfsetfillcolor{currentfill}%
\pgfsetfillopacity{0.900000}%
\pgfsetlinewidth{0.501875pt}%
\definecolor{currentstroke}{rgb}{0.200000,0.200000,0.200000}%
\pgfsetstrokecolor{currentstroke}%
\pgfsetdash{}{0pt}%
\pgfpathmoveto{\pgfqpoint{4.225663in}{1.397930in}}%
\pgfpathlineto{\pgfqpoint{4.269450in}{1.408914in}}%
\pgfpathlineto{\pgfqpoint{4.304192in}{1.388076in}}%
\pgfpathlineto{\pgfqpoint{4.260380in}{1.377162in}}%
\pgfpathlineto{\pgfqpoint{4.225663in}{1.397930in}}%
\pgfpathclose%
\pgfusepath{stroke,fill}%
\end{pgfscope}%
\begin{pgfscope}%
\pgfpathrectangle{\pgfqpoint{0.050000in}{0.050000in}}{\pgfqpoint{4.900000in}{4.900000in}}%
\pgfusepath{clip}%
\pgfsetbuttcap%
\pgfsetroundjoin%
\definecolor{currentfill}{rgb}{0.898378,0.898378,0.898378}%
\pgfsetfillcolor{currentfill}%
\pgfsetfillopacity{0.900000}%
\pgfsetlinewidth{0.501875pt}%
\definecolor{currentstroke}{rgb}{0.200000,0.200000,0.200000}%
\pgfsetstrokecolor{currentstroke}%
\pgfsetdash{}{0pt}%
\pgfpathmoveto{\pgfqpoint{2.934903in}{1.866002in}}%
\pgfpathlineto{\pgfqpoint{2.979873in}{1.830557in}}%
\pgfpathlineto{\pgfqpoint{3.012575in}{1.810469in}}%
\pgfpathlineto{\pgfqpoint{2.967596in}{1.845682in}}%
\pgfpathlineto{\pgfqpoint{2.934903in}{1.866002in}}%
\pgfpathclose%
\pgfusepath{stroke,fill}%
\end{pgfscope}%
\begin{pgfscope}%
\pgfpathrectangle{\pgfqpoint{0.050000in}{0.050000in}}{\pgfqpoint{4.900000in}{4.900000in}}%
\pgfusepath{clip}%
\pgfsetbuttcap%
\pgfsetroundjoin%
\definecolor{currentfill}{rgb}{0.992618,0.992618,0.992618}%
\pgfsetfillcolor{currentfill}%
\pgfsetfillopacity{0.900000}%
\pgfsetlinewidth{0.501875pt}%
\definecolor{currentstroke}{rgb}{0.200000,0.200000,0.200000}%
\pgfsetstrokecolor{currentstroke}%
\pgfsetdash{}{0pt}%
\pgfpathmoveto{\pgfqpoint{3.946572in}{1.425417in}}%
\pgfpathlineto{\pgfqpoint{3.990575in}{1.431941in}}%
\pgfpathlineto{\pgfqpoint{4.024892in}{1.411605in}}%
\pgfpathlineto{\pgfqpoint{3.980865in}{1.405161in}}%
\pgfpathlineto{\pgfqpoint{3.946572in}{1.425417in}}%
\pgfpathclose%
\pgfusepath{stroke,fill}%
\end{pgfscope}%
\begin{pgfscope}%
\pgfpathrectangle{\pgfqpoint{0.050000in}{0.050000in}}{\pgfqpoint{4.900000in}{4.900000in}}%
\pgfusepath{clip}%
\pgfsetbuttcap%
\pgfsetroundjoin%
\definecolor{currentfill}{rgb}{0.395663,0.395663,0.395663}%
\pgfsetfillcolor{currentfill}%
\pgfsetfillopacity{0.900000}%
\pgfsetlinewidth{0.501875pt}%
\definecolor{currentstroke}{rgb}{0.200000,0.200000,0.200000}%
\pgfsetstrokecolor{currentstroke}%
\pgfsetdash{}{0pt}%
\pgfpathmoveto{\pgfqpoint{1.418801in}{3.076933in}}%
\pgfpathlineto{\pgfqpoint{1.466939in}{3.030991in}}%
\pgfpathlineto{\pgfqpoint{1.497282in}{3.011580in}}%
\pgfpathlineto{\pgfqpoint{1.449135in}{3.057552in}}%
\pgfpathlineto{\pgfqpoint{1.418801in}{3.076933in}}%
\pgfpathclose%
\pgfusepath{stroke,fill}%
\end{pgfscope}%
\begin{pgfscope}%
\pgfpathrectangle{\pgfqpoint{0.050000in}{0.050000in}}{\pgfqpoint{4.900000in}{4.900000in}}%
\pgfusepath{clip}%
\pgfsetbuttcap%
\pgfsetroundjoin%
\definecolor{currentfill}{rgb}{0.657809,0.657809,0.657809}%
\pgfsetfillcolor{currentfill}%
\pgfsetfillopacity{0.900000}%
\pgfsetlinewidth{0.501875pt}%
\definecolor{currentstroke}{rgb}{0.200000,0.200000,0.200000}%
\pgfsetstrokecolor{currentstroke}%
\pgfsetdash{}{0pt}%
\pgfpathmoveto{\pgfqpoint{2.137917in}{2.495725in}}%
\pgfpathlineto{\pgfqpoint{2.184531in}{2.451019in}}%
\pgfpathlineto{\pgfqpoint{2.215991in}{2.430711in}}%
\pgfpathlineto{\pgfqpoint{2.169371in}{2.475442in}}%
\pgfpathlineto{\pgfqpoint{2.137917in}{2.495725in}}%
\pgfpathclose%
\pgfusepath{stroke,fill}%
\end{pgfscope}%
\begin{pgfscope}%
\pgfpathrectangle{\pgfqpoint{0.050000in}{0.050000in}}{\pgfqpoint{4.900000in}{4.900000in}}%
\pgfusepath{clip}%
\pgfsetbuttcap%
\pgfsetroundjoin%
\definecolor{currentfill}{rgb}{0.295271,0.295271,0.295271}%
\pgfsetfillcolor{currentfill}%
\pgfsetfillopacity{0.900000}%
\pgfsetlinewidth{0.501875pt}%
\definecolor{currentstroke}{rgb}{0.200000,0.200000,0.200000}%
\pgfsetstrokecolor{currentstroke}%
\pgfsetdash{}{0pt}%
\pgfpathmoveto{\pgfqpoint{1.152892in}{3.292495in}}%
\pgfpathlineto{\pgfqpoint{1.201602in}{3.246087in}}%
\pgfpathlineto{\pgfqpoint{1.231538in}{3.227006in}}%
\pgfpathlineto{\pgfqpoint{1.182817in}{3.273445in}}%
\pgfpathlineto{\pgfqpoint{1.152892in}{3.292495in}}%
\pgfpathclose%
\pgfusepath{stroke,fill}%
\end{pgfscope}%
\begin{pgfscope}%
\pgfpathrectangle{\pgfqpoint{0.050000in}{0.050000in}}{\pgfqpoint{4.900000in}{4.900000in}}%
\pgfusepath{clip}%
\pgfsetbuttcap%
\pgfsetroundjoin%
\definecolor{currentfill}{rgb}{0.966782,0.966782,0.966782}%
\pgfsetfillcolor{currentfill}%
\pgfsetfillopacity{0.900000}%
\pgfsetlinewidth{0.501875pt}%
\definecolor{currentstroke}{rgb}{0.200000,0.200000,0.200000}%
\pgfsetstrokecolor{currentstroke}%
\pgfsetdash{}{0pt}%
\pgfpathmoveto{\pgfqpoint{3.434184in}{1.568107in}}%
\pgfpathlineto{\pgfqpoint{3.478526in}{1.555049in}}%
\pgfpathlineto{\pgfqpoint{3.512044in}{1.535456in}}%
\pgfpathlineto{\pgfqpoint{3.467684in}{1.548442in}}%
\pgfpathlineto{\pgfqpoint{3.434184in}{1.568107in}}%
\pgfpathclose%
\pgfusepath{stroke,fill}%
\end{pgfscope}%
\begin{pgfscope}%
\pgfpathrectangle{\pgfqpoint{0.050000in}{0.050000in}}{\pgfqpoint{4.900000in}{4.900000in}}%
\pgfusepath{clip}%
\pgfsetbuttcap%
\pgfsetroundjoin%
\definecolor{currentfill}{rgb}{0.555940,0.555940,0.555940}%
\pgfsetfillcolor{currentfill}%
\pgfsetfillopacity{0.900000}%
\pgfsetlinewidth{0.501875pt}%
\definecolor{currentstroke}{rgb}{0.200000,0.200000,0.200000}%
\pgfsetstrokecolor{currentstroke}%
\pgfsetdash{}{0pt}%
\pgfpathmoveto{\pgfqpoint{1.872111in}{2.711194in}}%
\pgfpathlineto{\pgfqpoint{1.919297in}{2.666018in}}%
\pgfpathlineto{\pgfqpoint{1.950351in}{2.646036in}}%
\pgfpathlineto{\pgfqpoint{1.903157in}{2.691240in}}%
\pgfpathlineto{\pgfqpoint{1.872111in}{2.711194in}}%
\pgfpathclose%
\pgfusepath{stroke,fill}%
\end{pgfscope}%
\begin{pgfscope}%
\pgfpathrectangle{\pgfqpoint{0.050000in}{0.050000in}}{\pgfqpoint{4.900000in}{4.900000in}}%
\pgfusepath{clip}%
\pgfsetbuttcap%
\pgfsetroundjoin%
\definecolor{currentfill}{rgb}{0.972318,0.972318,0.972318}%
\pgfsetfillcolor{currentfill}%
\pgfsetfillopacity{0.900000}%
\pgfsetlinewidth{0.501875pt}%
\definecolor{currentstroke}{rgb}{0.200000,0.200000,0.200000}%
\pgfsetstrokecolor{currentstroke}%
\pgfsetdash{}{0pt}%
\pgfpathmoveto{\pgfqpoint{3.512044in}{1.535456in}}%
\pgfpathlineto{\pgfqpoint{3.556327in}{1.526110in}}%
\pgfpathlineto{\pgfqpoint{3.589971in}{1.506466in}}%
\pgfpathlineto{\pgfqpoint{3.545669in}{1.515780in}}%
\pgfpathlineto{\pgfqpoint{3.512044in}{1.535456in}}%
\pgfpathclose%
\pgfusepath{stroke,fill}%
\end{pgfscope}%
\begin{pgfscope}%
\pgfpathrectangle{\pgfqpoint{0.050000in}{0.050000in}}{\pgfqpoint{4.900000in}{4.900000in}}%
\pgfusepath{clip}%
\pgfsetbuttcap%
\pgfsetroundjoin%
\definecolor{currentfill}{rgb}{0.959400,0.959400,0.959400}%
\pgfsetfillcolor{currentfill}%
\pgfsetfillopacity{0.900000}%
\pgfsetlinewidth{0.501875pt}%
\definecolor{currentstroke}{rgb}{0.200000,0.200000,0.200000}%
\pgfsetstrokecolor{currentstroke}%
\pgfsetdash{}{0pt}%
\pgfpathmoveto{\pgfqpoint{3.356378in}{1.604642in}}%
\pgfpathlineto{\pgfqpoint{3.400790in}{1.587689in}}%
\pgfpathlineto{\pgfqpoint{3.434184in}{1.568107in}}%
\pgfpathlineto{\pgfqpoint{3.389755in}{1.584945in}}%
\pgfpathlineto{\pgfqpoint{3.356378in}{1.604642in}}%
\pgfpathclose%
\pgfusepath{stroke,fill}%
\end{pgfscope}%
\begin{pgfscope}%
\pgfpathrectangle{\pgfqpoint{0.050000in}{0.050000in}}{\pgfqpoint{4.900000in}{4.900000in}}%
\pgfusepath{clip}%
\pgfsetbuttcap%
\pgfsetroundjoin%
\definecolor{currentfill}{rgb}{1.000000,1.000000,1.000000}%
\pgfsetfillcolor{currentfill}%
\pgfsetfillopacity{0.900000}%
\pgfsetlinewidth{0.501875pt}%
\definecolor{currentstroke}{rgb}{0.200000,0.200000,0.200000}%
\pgfsetstrokecolor{currentstroke}%
\pgfsetdash{}{0pt}%
\pgfpathmoveto{\pgfqpoint{4.304192in}{1.388076in}}%
\pgfpathlineto{\pgfqpoint{4.347921in}{1.399572in}}%
\pgfpathlineto{\pgfqpoint{4.382797in}{1.378601in}}%
\pgfpathlineto{\pgfqpoint{4.339043in}{1.367171in}}%
\pgfpathlineto{\pgfqpoint{4.304192in}{1.388076in}}%
\pgfpathclose%
\pgfusepath{stroke,fill}%
\end{pgfscope}%
\begin{pgfscope}%
\pgfpathrectangle{\pgfqpoint{0.050000in}{0.050000in}}{\pgfqpoint{4.900000in}{4.900000in}}%
\pgfusepath{clip}%
\pgfsetbuttcap%
\pgfsetroundjoin%
\definecolor{currentfill}{rgb}{0.994464,0.994464,0.994464}%
\pgfsetfillcolor{currentfill}%
\pgfsetfillopacity{0.900000}%
\pgfsetlinewidth{0.501875pt}%
\definecolor{currentstroke}{rgb}{0.200000,0.200000,0.200000}%
\pgfsetstrokecolor{currentstroke}%
\pgfsetdash{}{0pt}%
\pgfpathmoveto{\pgfqpoint{4.024892in}{1.411605in}}%
\pgfpathlineto{\pgfqpoint{4.068852in}{1.419623in}}%
\pgfpathlineto{\pgfqpoint{4.103301in}{1.399140in}}%
\pgfpathlineto{\pgfqpoint{4.059316in}{1.391201in}}%
\pgfpathlineto{\pgfqpoint{4.024892in}{1.411605in}}%
\pgfpathclose%
\pgfusepath{stroke,fill}%
\end{pgfscope}%
\begin{pgfscope}%
\pgfpathrectangle{\pgfqpoint{0.050000in}{0.050000in}}{\pgfqpoint{4.900000in}{4.900000in}}%
\pgfusepath{clip}%
\pgfsetbuttcap%
\pgfsetroundjoin%
\definecolor{currentfill}{rgb}{0.812226,0.812226,0.812226}%
\pgfsetfillcolor{currentfill}%
\pgfsetfillopacity{0.900000}%
\pgfsetlinewidth{0.501875pt}%
\definecolor{currentstroke}{rgb}{0.200000,0.200000,0.200000}%
\pgfsetstrokecolor{currentstroke}%
\pgfsetdash{}{0pt}%
\pgfpathmoveto{\pgfqpoint{2.591428in}{2.131094in}}%
\pgfpathlineto{\pgfqpoint{2.637092in}{2.088254in}}%
\pgfpathlineto{\pgfqpoint{2.669264in}{2.067661in}}%
\pgfpathlineto{\pgfqpoint{2.623596in}{2.110414in}}%
\pgfpathlineto{\pgfqpoint{2.591428in}{2.131094in}}%
\pgfpathclose%
\pgfusepath{stroke,fill}%
\end{pgfscope}%
\begin{pgfscope}%
\pgfpathrectangle{\pgfqpoint{0.050000in}{0.050000in}}{\pgfqpoint{4.900000in}{4.900000in}}%
\pgfusepath{clip}%
\pgfsetbuttcap%
\pgfsetroundjoin%
\definecolor{currentfill}{rgb}{0.173472,0.173472,0.173472}%
\pgfsetfillcolor{currentfill}%
\pgfsetfillopacity{0.900000}%
\pgfsetlinewidth{0.501875pt}%
\definecolor{currentstroke}{rgb}{0.200000,0.200000,0.200000}%
\pgfsetstrokecolor{currentstroke}%
\pgfsetdash{}{0pt}%
\pgfpathmoveto{\pgfqpoint{0.886892in}{3.508132in}}%
\pgfpathlineto{\pgfqpoint{0.936175in}{3.461258in}}%
\pgfpathlineto{\pgfqpoint{0.965705in}{3.442507in}}%
\pgfpathlineto{\pgfqpoint{0.916408in}{3.489414in}}%
\pgfpathlineto{\pgfqpoint{0.886892in}{3.508132in}}%
\pgfpathclose%
\pgfusepath{stroke,fill}%
\end{pgfscope}%
\begin{pgfscope}%
\pgfpathrectangle{\pgfqpoint{0.050000in}{0.050000in}}{\pgfqpoint{4.900000in}{4.900000in}}%
\pgfusepath{clip}%
\pgfsetbuttcap%
\pgfsetroundjoin%
\definecolor{currentfill}{rgb}{0.977855,0.977855,0.977855}%
\pgfsetfillcolor{currentfill}%
\pgfsetfillopacity{0.900000}%
\pgfsetlinewidth{0.501875pt}%
\definecolor{currentstroke}{rgb}{0.200000,0.200000,0.200000}%
\pgfsetstrokecolor{currentstroke}%
\pgfsetdash{}{0pt}%
\pgfpathmoveto{\pgfqpoint{3.589971in}{1.506466in}}%
\pgfpathlineto{\pgfqpoint{3.634203in}{1.500575in}}%
\pgfpathlineto{\pgfqpoint{3.667975in}{1.480846in}}%
\pgfpathlineto{\pgfqpoint{3.623722in}{1.486740in}}%
\pgfpathlineto{\pgfqpoint{3.589971in}{1.506466in}}%
\pgfpathclose%
\pgfusepath{stroke,fill}%
\end{pgfscope}%
\begin{pgfscope}%
\pgfpathrectangle{\pgfqpoint{0.050000in}{0.050000in}}{\pgfqpoint{4.900000in}{4.900000in}}%
\pgfusepath{clip}%
\pgfsetbuttcap%
\pgfsetroundjoin%
\definecolor{currentfill}{rgb}{0.878570,0.878570,0.878570}%
\pgfsetfillcolor{currentfill}%
\pgfsetfillopacity{0.900000}%
\pgfsetlinewidth{0.501875pt}%
\definecolor{currentstroke}{rgb}{0.200000,0.200000,0.200000}%
\pgfsetstrokecolor{currentstroke}%
\pgfsetdash{}{0pt}%
\pgfpathmoveto{\pgfqpoint{2.857181in}{1.924299in}}%
\pgfpathlineto{\pgfqpoint{2.902311in}{1.886269in}}%
\pgfpathlineto{\pgfqpoint{2.934903in}{1.866002in}}%
\pgfpathlineto{\pgfqpoint{2.889766in}{1.903821in}}%
\pgfpathlineto{\pgfqpoint{2.857181in}{1.924299in}}%
\pgfpathclose%
\pgfusepath{stroke,fill}%
\end{pgfscope}%
\begin{pgfscope}%
\pgfpathrectangle{\pgfqpoint{0.050000in}{0.050000in}}{\pgfqpoint{4.900000in}{4.900000in}}%
\pgfusepath{clip}%
\pgfsetbuttcap%
\pgfsetroundjoin%
\definecolor{currentfill}{rgb}{0.952018,0.952018,0.952018}%
\pgfsetfillcolor{currentfill}%
\pgfsetfillopacity{0.900000}%
\pgfsetlinewidth{0.501875pt}%
\definecolor{currentstroke}{rgb}{0.200000,0.200000,0.200000}%
\pgfsetstrokecolor{currentstroke}%
\pgfsetdash{}{0pt}%
\pgfpathmoveto{\pgfqpoint{3.278612in}{1.645188in}}%
\pgfpathlineto{\pgfqpoint{3.323106in}{1.624259in}}%
\pgfpathlineto{\pgfqpoint{3.356378in}{1.604642in}}%
\pgfpathlineto{\pgfqpoint{3.311869in}{1.625415in}}%
\pgfpathlineto{\pgfqpoint{3.278612in}{1.645188in}}%
\pgfpathclose%
\pgfusepath{stroke,fill}%
\end{pgfscope}%
\begin{pgfscope}%
\pgfpathrectangle{\pgfqpoint{0.050000in}{0.050000in}}{\pgfqpoint{4.900000in}{4.900000in}}%
\pgfusepath{clip}%
\pgfsetbuttcap%
\pgfsetroundjoin%
\definecolor{currentfill}{rgb}{0.456901,0.456901,0.456901}%
\pgfsetfillcolor{currentfill}%
\pgfsetfillopacity{0.900000}%
\pgfsetlinewidth{0.501875pt}%
\definecolor{currentstroke}{rgb}{0.200000,0.200000,0.200000}%
\pgfsetstrokecolor{currentstroke}%
\pgfsetdash{}{0pt}%
\pgfpathmoveto{\pgfqpoint{1.606216in}{2.926740in}}%
\pgfpathlineto{\pgfqpoint{1.653974in}{2.881099in}}%
\pgfpathlineto{\pgfqpoint{1.684621in}{2.861447in}}%
\pgfpathlineto{\pgfqpoint{1.636853in}{2.907118in}}%
\pgfpathlineto{\pgfqpoint{1.606216in}{2.926740in}}%
\pgfpathclose%
\pgfusepath{stroke,fill}%
\end{pgfscope}%
\begin{pgfscope}%
\pgfpathrectangle{\pgfqpoint{0.050000in}{0.050000in}}{\pgfqpoint{4.900000in}{4.900000in}}%
\pgfusepath{clip}%
\pgfsetbuttcap%
\pgfsetroundjoin%
\definecolor{currentfill}{rgb}{0.729781,0.729781,0.729781}%
\pgfsetfillcolor{currentfill}%
\pgfsetfillopacity{0.900000}%
\pgfsetlinewidth{0.501875pt}%
\definecolor{currentstroke}{rgb}{0.200000,0.200000,0.200000}%
\pgfsetstrokecolor{currentstroke}%
\pgfsetdash{}{0pt}%
\pgfpathmoveto{\pgfqpoint{2.325635in}{2.345343in}}%
\pgfpathlineto{\pgfqpoint{2.371869in}{2.301014in}}%
\pgfpathlineto{\pgfqpoint{2.403633in}{2.280495in}}%
\pgfpathlineto{\pgfqpoint{2.357395in}{2.324832in}}%
\pgfpathlineto{\pgfqpoint{2.325635in}{2.345343in}}%
\pgfpathclose%
\pgfusepath{stroke,fill}%
\end{pgfscope}%
\begin{pgfscope}%
\pgfpathrectangle{\pgfqpoint{0.050000in}{0.050000in}}{\pgfqpoint{4.900000in}{4.900000in}}%
\pgfusepath{clip}%
\pgfsetbuttcap%
\pgfsetroundjoin%
\definecolor{currentfill}{rgb}{0.981546,0.981546,0.981546}%
\pgfsetfillcolor{currentfill}%
\pgfsetfillopacity{0.900000}%
\pgfsetlinewidth{0.501875pt}%
\definecolor{currentstroke}{rgb}{0.200000,0.200000,0.200000}%
\pgfsetstrokecolor{currentstroke}%
\pgfsetdash{}{0pt}%
\pgfpathmoveto{\pgfqpoint{3.667975in}{1.480846in}}%
\pgfpathlineto{\pgfqpoint{3.712162in}{1.478100in}}%
\pgfpathlineto{\pgfqpoint{3.746062in}{1.458261in}}%
\pgfpathlineto{\pgfqpoint{3.701854in}{1.461037in}}%
\pgfpathlineto{\pgfqpoint{3.667975in}{1.480846in}}%
\pgfpathclose%
\pgfusepath{stroke,fill}%
\end{pgfscope}%
\begin{pgfscope}%
\pgfpathrectangle{\pgfqpoint{0.050000in}{0.050000in}}{\pgfqpoint{4.900000in}{4.900000in}}%
\pgfusepath{clip}%
\pgfsetbuttcap%
\pgfsetroundjoin%
\definecolor{currentfill}{rgb}{0.944637,0.944637,0.944637}%
\pgfsetfillcolor{currentfill}%
\pgfsetfillopacity{0.900000}%
\pgfsetlinewidth{0.501875pt}%
\definecolor{currentstroke}{rgb}{0.200000,0.200000,0.200000}%
\pgfsetstrokecolor{currentstroke}%
\pgfsetdash{}{0pt}%
\pgfpathmoveto{\pgfqpoint{3.200869in}{1.689756in}}%
\pgfpathlineto{\pgfqpoint{3.245460in}{1.664885in}}%
\pgfpathlineto{\pgfqpoint{3.278612in}{1.645188in}}%
\pgfpathlineto{\pgfqpoint{3.234007in}{1.669867in}}%
\pgfpathlineto{\pgfqpoint{3.200869in}{1.689756in}}%
\pgfpathclose%
\pgfusepath{stroke,fill}%
\end{pgfscope}%
\begin{pgfscope}%
\pgfpathrectangle{\pgfqpoint{0.050000in}{0.050000in}}{\pgfqpoint{4.900000in}{4.900000in}}%
\pgfusepath{clip}%
\pgfsetbuttcap%
\pgfsetroundjoin%
\definecolor{currentfill}{rgb}{0.363183,0.363183,0.363183}%
\pgfsetfillcolor{currentfill}%
\pgfsetfillopacity{0.900000}%
\pgfsetlinewidth{0.501875pt}%
\definecolor{currentstroke}{rgb}{0.200000,0.200000,0.200000}%
\pgfsetstrokecolor{currentstroke}%
\pgfsetdash{}{0pt}%
\pgfpathmoveto{\pgfqpoint{1.340230in}{3.142362in}}%
\pgfpathlineto{\pgfqpoint{1.388562in}{3.096254in}}%
\pgfpathlineto{\pgfqpoint{1.418801in}{3.076933in}}%
\pgfpathlineto{\pgfqpoint{1.370459in}{3.123071in}}%
\pgfpathlineto{\pgfqpoint{1.340230in}{3.142362in}}%
\pgfpathclose%
\pgfusepath{stroke,fill}%
\end{pgfscope}%
\begin{pgfscope}%
\pgfpathrectangle{\pgfqpoint{0.050000in}{0.050000in}}{\pgfqpoint{4.900000in}{4.900000in}}%
\pgfusepath{clip}%
\pgfsetbuttcap%
\pgfsetroundjoin%
\definecolor{currentfill}{rgb}{0.624221,0.624221,0.624221}%
\pgfsetfillcolor{currentfill}%
\pgfsetfillopacity{0.900000}%
\pgfsetlinewidth{0.501875pt}%
\definecolor{currentstroke}{rgb}{0.200000,0.200000,0.200000}%
\pgfsetstrokecolor{currentstroke}%
\pgfsetdash{}{0pt}%
\pgfpathmoveto{\pgfqpoint{2.059754in}{2.560821in}}%
\pgfpathlineto{\pgfqpoint{2.106560in}{2.515947in}}%
\pgfpathlineto{\pgfqpoint{2.137917in}{2.495725in}}%
\pgfpathlineto{\pgfqpoint{2.091105in}{2.540626in}}%
\pgfpathlineto{\pgfqpoint{2.059754in}{2.560821in}}%
\pgfpathclose%
\pgfusepath{stroke,fill}%
\end{pgfscope}%
\begin{pgfscope}%
\pgfpathrectangle{\pgfqpoint{0.050000in}{0.050000in}}{\pgfqpoint{4.900000in}{4.900000in}}%
\pgfusepath{clip}%
\pgfsetbuttcap%
\pgfsetroundjoin%
\definecolor{currentfill}{rgb}{0.996309,0.996309,0.996309}%
\pgfsetfillcolor{currentfill}%
\pgfsetfillopacity{0.900000}%
\pgfsetlinewidth{0.501875pt}%
\definecolor{currentstroke}{rgb}{0.200000,0.200000,0.200000}%
\pgfsetstrokecolor{currentstroke}%
\pgfsetdash{}{0pt}%
\pgfpathmoveto{\pgfqpoint{4.103301in}{1.399140in}}%
\pgfpathlineto{\pgfqpoint{4.147215in}{1.408370in}}%
\pgfpathlineto{\pgfqpoint{4.181798in}{1.387742in}}%
\pgfpathlineto{\pgfqpoint{4.137859in}{1.378589in}}%
\pgfpathlineto{\pgfqpoint{4.103301in}{1.399140in}}%
\pgfpathclose%
\pgfusepath{stroke,fill}%
\end{pgfscope}%
\begin{pgfscope}%
\pgfpathrectangle{\pgfqpoint{0.050000in}{0.050000in}}{\pgfqpoint{4.900000in}{4.900000in}}%
\pgfusepath{clip}%
\pgfsetbuttcap%
\pgfsetroundjoin%
\definecolor{currentfill}{rgb}{0.985236,0.985236,0.985236}%
\pgfsetfillcolor{currentfill}%
\pgfsetfillopacity{0.900000}%
\pgfsetlinewidth{0.501875pt}%
\definecolor{currentstroke}{rgb}{0.200000,0.200000,0.200000}%
\pgfsetstrokecolor{currentstroke}%
\pgfsetdash{}{0pt}%
\pgfpathmoveto{\pgfqpoint{3.746062in}{1.458261in}}%
\pgfpathlineto{\pgfqpoint{3.790208in}{1.458323in}}%
\pgfpathlineto{\pgfqpoint{3.824238in}{1.438354in}}%
\pgfpathlineto{\pgfqpoint{3.780070in}{1.438343in}}%
\pgfpathlineto{\pgfqpoint{3.746062in}{1.458261in}}%
\pgfpathclose%
\pgfusepath{stroke,fill}%
\end{pgfscope}%
\begin{pgfscope}%
\pgfpathrectangle{\pgfqpoint{0.050000in}{0.050000in}}{\pgfqpoint{4.900000in}{4.900000in}}%
\pgfusepath{clip}%
\pgfsetbuttcap%
\pgfsetroundjoin%
\definecolor{currentfill}{rgb}{0.256517,0.256517,0.256517}%
\pgfsetfillcolor{currentfill}%
\pgfsetfillopacity{0.900000}%
\pgfsetlinewidth{0.501875pt}%
\definecolor{currentstroke}{rgb}{0.200000,0.200000,0.200000}%
\pgfsetstrokecolor{currentstroke}%
\pgfsetdash{}{0pt}%
\pgfpathmoveto{\pgfqpoint{1.074154in}{3.358059in}}%
\pgfpathlineto{\pgfqpoint{1.123059in}{3.311485in}}%
\pgfpathlineto{\pgfqpoint{1.152892in}{3.292495in}}%
\pgfpathlineto{\pgfqpoint{1.103974in}{3.339100in}}%
\pgfpathlineto{\pgfqpoint{1.074154in}{3.358059in}}%
\pgfpathclose%
\pgfusepath{stroke,fill}%
\end{pgfscope}%
\begin{pgfscope}%
\pgfpathrectangle{\pgfqpoint{0.050000in}{0.050000in}}{\pgfqpoint{4.900000in}{4.900000in}}%
\pgfusepath{clip}%
\pgfsetbuttcap%
\pgfsetroundjoin%
\definecolor{currentfill}{rgb}{0.929504,0.929504,0.929504}%
\pgfsetfillcolor{currentfill}%
\pgfsetfillopacity{0.900000}%
\pgfsetlinewidth{0.501875pt}%
\definecolor{currentstroke}{rgb}{0.200000,0.200000,0.200000}%
\pgfsetstrokecolor{currentstroke}%
\pgfsetdash{}{0pt}%
\pgfpathmoveto{\pgfqpoint{3.123131in}{1.738221in}}%
\pgfpathlineto{\pgfqpoint{3.167835in}{1.709574in}}%
\pgfpathlineto{\pgfqpoint{3.200869in}{1.689756in}}%
\pgfpathlineto{\pgfqpoint{3.156153in}{1.718184in}}%
\pgfpathlineto{\pgfqpoint{3.123131in}{1.738221in}}%
\pgfpathclose%
\pgfusepath{stroke,fill}%
\end{pgfscope}%
\begin{pgfscope}%
\pgfpathrectangle{\pgfqpoint{0.050000in}{0.050000in}}{\pgfqpoint{4.900000in}{4.900000in}}%
\pgfusepath{clip}%
\pgfsetbuttcap%
\pgfsetroundjoin%
\definecolor{currentfill}{rgb}{0.521492,0.521492,0.521492}%
\pgfsetfillcolor{currentfill}%
\pgfsetfillopacity{0.900000}%
\pgfsetlinewidth{0.501875pt}%
\definecolor{currentstroke}{rgb}{0.200000,0.200000,0.200000}%
\pgfsetstrokecolor{currentstroke}%
\pgfsetdash{}{0pt}%
\pgfpathmoveto{\pgfqpoint{1.793782in}{2.776427in}}%
\pgfpathlineto{\pgfqpoint{1.841162in}{2.731086in}}%
\pgfpathlineto{\pgfqpoint{1.872111in}{2.711194in}}%
\pgfpathlineto{\pgfqpoint{1.824724in}{2.756563in}}%
\pgfpathlineto{\pgfqpoint{1.793782in}{2.776427in}}%
\pgfpathclose%
\pgfusepath{stroke,fill}%
\end{pgfscope}%
\begin{pgfscope}%
\pgfpathrectangle{\pgfqpoint{0.050000in}{0.050000in}}{\pgfqpoint{4.900000in}{4.900000in}}%
\pgfusepath{clip}%
\pgfsetbuttcap%
\pgfsetroundjoin%
\definecolor{currentfill}{rgb}{0.858762,0.858762,0.858762}%
\pgfsetfillcolor{currentfill}%
\pgfsetfillopacity{0.900000}%
\pgfsetlinewidth{0.501875pt}%
\definecolor{currentstroke}{rgb}{0.200000,0.200000,0.200000}%
\pgfsetstrokecolor{currentstroke}%
\pgfsetdash{}{0pt}%
\pgfpathmoveto{\pgfqpoint{2.779395in}{1.984818in}}%
\pgfpathlineto{\pgfqpoint{2.824697in}{1.944728in}}%
\pgfpathlineto{\pgfqpoint{2.857181in}{1.924299in}}%
\pgfpathlineto{\pgfqpoint{2.811874in}{1.964213in}}%
\pgfpathlineto{\pgfqpoint{2.779395in}{1.984818in}}%
\pgfpathclose%
\pgfusepath{stroke,fill}%
\end{pgfscope}%
\begin{pgfscope}%
\pgfpathrectangle{\pgfqpoint{0.050000in}{0.050000in}}{\pgfqpoint{4.900000in}{4.900000in}}%
\pgfusepath{clip}%
\pgfsetbuttcap%
\pgfsetroundjoin%
\definecolor{currentfill}{rgb}{0.788112,0.788112,0.788112}%
\pgfsetfillcolor{currentfill}%
\pgfsetfillopacity{0.900000}%
\pgfsetlinewidth{0.501875pt}%
\definecolor{currentstroke}{rgb}{0.200000,0.200000,0.200000}%
\pgfsetstrokecolor{currentstroke}%
\pgfsetdash{}{0pt}%
\pgfpathmoveto{\pgfqpoint{2.513506in}{2.195276in}}%
\pgfpathlineto{\pgfqpoint{2.559360in}{2.151722in}}%
\pgfpathlineto{\pgfqpoint{2.591428in}{2.131094in}}%
\pgfpathlineto{\pgfqpoint{2.545571in}{2.174602in}}%
\pgfpathlineto{\pgfqpoint{2.513506in}{2.195276in}}%
\pgfpathclose%
\pgfusepath{stroke,fill}%
\end{pgfscope}%
\begin{pgfscope}%
\pgfpathrectangle{\pgfqpoint{0.050000in}{0.050000in}}{\pgfqpoint{4.900000in}{4.900000in}}%
\pgfusepath{clip}%
\pgfsetbuttcap%
\pgfsetroundjoin%
\definecolor{currentfill}{rgb}{0.988927,0.988927,0.988927}%
\pgfsetfillcolor{currentfill}%
\pgfsetfillopacity{0.900000}%
\pgfsetlinewidth{0.501875pt}%
\definecolor{currentstroke}{rgb}{0.200000,0.200000,0.200000}%
\pgfsetstrokecolor{currentstroke}%
\pgfsetdash{}{0pt}%
\pgfpathmoveto{\pgfqpoint{3.824238in}{1.438354in}}%
\pgfpathlineto{\pgfqpoint{3.868345in}{1.440879in}}%
\pgfpathlineto{\pgfqpoint{3.902505in}{1.420770in}}%
\pgfpathlineto{\pgfqpoint{3.858376in}{1.418310in}}%
\pgfpathlineto{\pgfqpoint{3.824238in}{1.438354in}}%
\pgfpathclose%
\pgfusepath{stroke,fill}%
\end{pgfscope}%
\begin{pgfscope}%
\pgfpathrectangle{\pgfqpoint{0.050000in}{0.050000in}}{\pgfqpoint{4.900000in}{4.900000in}}%
\pgfusepath{clip}%
\pgfsetbuttcap%
\pgfsetroundjoin%
\definecolor{currentfill}{rgb}{0.428143,0.428143,0.428143}%
\pgfsetfillcolor{currentfill}%
\pgfsetfillopacity{0.900000}%
\pgfsetlinewidth{0.501875pt}%
\definecolor{currentstroke}{rgb}{0.200000,0.200000,0.200000}%
\pgfsetstrokecolor{currentstroke}%
\pgfsetdash{}{0pt}%
\pgfpathmoveto{\pgfqpoint{1.527720in}{2.992109in}}%
\pgfpathlineto{\pgfqpoint{1.575673in}{2.946302in}}%
\pgfpathlineto{\pgfqpoint{1.606216in}{2.926740in}}%
\pgfpathlineto{\pgfqpoint{1.558253in}{2.972577in}}%
\pgfpathlineto{\pgfqpoint{1.527720in}{2.992109in}}%
\pgfpathclose%
\pgfusepath{stroke,fill}%
\end{pgfscope}%
\begin{pgfscope}%
\pgfpathrectangle{\pgfqpoint{0.050000in}{0.050000in}}{\pgfqpoint{4.900000in}{4.900000in}}%
\pgfusepath{clip}%
\pgfsetbuttcap%
\pgfsetroundjoin%
\definecolor{currentfill}{rgb}{0.696194,0.696194,0.696194}%
\pgfsetfillcolor{currentfill}%
\pgfsetfillopacity{0.900000}%
\pgfsetlinewidth{0.501875pt}%
\definecolor{currentstroke}{rgb}{0.200000,0.200000,0.200000}%
\pgfsetstrokecolor{currentstroke}%
\pgfsetdash{}{0pt}%
\pgfpathmoveto{\pgfqpoint{2.247548in}{2.410341in}}%
\pgfpathlineto{\pgfqpoint{2.293975in}{2.365794in}}%
\pgfpathlineto{\pgfqpoint{2.325635in}{2.345343in}}%
\pgfpathlineto{\pgfqpoint{2.279204in}{2.389909in}}%
\pgfpathlineto{\pgfqpoint{2.247548in}{2.410341in}}%
\pgfpathclose%
\pgfusepath{stroke,fill}%
\end{pgfscope}%
\begin{pgfscope}%
\pgfpathrectangle{\pgfqpoint{0.050000in}{0.050000in}}{\pgfqpoint{4.900000in}{4.900000in}}%
\pgfusepath{clip}%
\pgfsetbuttcap%
\pgfsetroundjoin%
\definecolor{currentfill}{rgb}{0.998155,0.998155,0.998155}%
\pgfsetfillcolor{currentfill}%
\pgfsetfillopacity{0.900000}%
\pgfsetlinewidth{0.501875pt}%
\definecolor{currentstroke}{rgb}{0.200000,0.200000,0.200000}%
\pgfsetstrokecolor{currentstroke}%
\pgfsetdash{}{0pt}%
\pgfpathmoveto{\pgfqpoint{4.181798in}{1.387742in}}%
\pgfpathlineto{\pgfqpoint{4.225663in}{1.397930in}}%
\pgfpathlineto{\pgfqpoint{4.260380in}{1.377162in}}%
\pgfpathlineto{\pgfqpoint{4.216490in}{1.367048in}}%
\pgfpathlineto{\pgfqpoint{4.181798in}{1.387742in}}%
\pgfpathclose%
\pgfusepath{stroke,fill}%
\end{pgfscope}%
\begin{pgfscope}%
\pgfpathrectangle{\pgfqpoint{0.050000in}{0.050000in}}{\pgfqpoint{4.900000in}{4.900000in}}%
\pgfusepath{clip}%
\pgfsetbuttcap%
\pgfsetroundjoin%
\definecolor{currentfill}{rgb}{0.915356,0.915356,0.915356}%
\pgfsetfillcolor{currentfill}%
\pgfsetfillopacity{0.900000}%
\pgfsetlinewidth{0.501875pt}%
\definecolor{currentstroke}{rgb}{0.200000,0.200000,0.200000}%
\pgfsetstrokecolor{currentstroke}%
\pgfsetdash{}{0pt}%
\pgfpathmoveto{\pgfqpoint{3.045379in}{1.790323in}}%
\pgfpathlineto{\pgfqpoint{3.090213in}{1.758193in}}%
\pgfpathlineto{\pgfqpoint{3.123131in}{1.738221in}}%
\pgfpathlineto{\pgfqpoint{3.078288in}{1.770120in}}%
\pgfpathlineto{\pgfqpoint{3.045379in}{1.790323in}}%
\pgfpathclose%
\pgfusepath{stroke,fill}%
\end{pgfscope}%
\begin{pgfscope}%
\pgfpathrectangle{\pgfqpoint{0.050000in}{0.050000in}}{\pgfqpoint{4.900000in}{4.900000in}}%
\pgfusepath{clip}%
\pgfsetbuttcap%
\pgfsetroundjoin%
\definecolor{currentfill}{rgb}{0.334764,0.334764,0.334764}%
\pgfsetfillcolor{currentfill}%
\pgfsetfillopacity{0.900000}%
\pgfsetlinewidth{0.501875pt}%
\definecolor{currentstroke}{rgb}{0.200000,0.200000,0.200000}%
\pgfsetstrokecolor{currentstroke}%
\pgfsetdash{}{0pt}%
\pgfpathmoveto{\pgfqpoint{1.261569in}{3.207866in}}%
\pgfpathlineto{\pgfqpoint{1.310095in}{3.161593in}}%
\pgfpathlineto{\pgfqpoint{1.340230in}{3.142362in}}%
\pgfpathlineto{\pgfqpoint{1.291692in}{3.188666in}}%
\pgfpathlineto{\pgfqpoint{1.261569in}{3.207866in}}%
\pgfpathclose%
\pgfusepath{stroke,fill}%
\end{pgfscope}%
\begin{pgfscope}%
\pgfpathrectangle{\pgfqpoint{0.050000in}{0.050000in}}{\pgfqpoint{4.900000in}{4.900000in}}%
\pgfusepath{clip}%
\pgfsetbuttcap%
\pgfsetroundjoin%
\definecolor{currentfill}{rgb}{0.586082,0.586082,0.586082}%
\pgfsetfillcolor{currentfill}%
\pgfsetfillopacity{0.900000}%
\pgfsetlinewidth{0.501875pt}%
\definecolor{currentstroke}{rgb}{0.200000,0.200000,0.200000}%
\pgfsetstrokecolor{currentstroke}%
\pgfsetdash{}{0pt}%
\pgfpathmoveto{\pgfqpoint{1.981501in}{2.625993in}}%
\pgfpathlineto{\pgfqpoint{2.028500in}{2.580954in}}%
\pgfpathlineto{\pgfqpoint{2.059754in}{2.560821in}}%
\pgfpathlineto{\pgfqpoint{2.012748in}{2.605887in}}%
\pgfpathlineto{\pgfqpoint{1.981501in}{2.625993in}}%
\pgfpathclose%
\pgfusepath{stroke,fill}%
\end{pgfscope}%
\begin{pgfscope}%
\pgfpathrectangle{\pgfqpoint{0.050000in}{0.050000in}}{\pgfqpoint{4.900000in}{4.900000in}}%
\pgfusepath{clip}%
\pgfsetbuttcap%
\pgfsetroundjoin%
\definecolor{currentfill}{rgb}{0.990773,0.990773,0.990773}%
\pgfsetfillcolor{currentfill}%
\pgfsetfillopacity{0.900000}%
\pgfsetlinewidth{0.501875pt}%
\definecolor{currentstroke}{rgb}{0.200000,0.200000,0.200000}%
\pgfsetstrokecolor{currentstroke}%
\pgfsetdash{}{0pt}%
\pgfpathmoveto{\pgfqpoint{3.902505in}{1.420770in}}%
\pgfpathlineto{\pgfqpoint{3.946572in}{1.425417in}}%
\pgfpathlineto{\pgfqpoint{3.980865in}{1.405161in}}%
\pgfpathlineto{\pgfqpoint{3.936774in}{1.400587in}}%
\pgfpathlineto{\pgfqpoint{3.902505in}{1.420770in}}%
\pgfpathclose%
\pgfusepath{stroke,fill}%
\end{pgfscope}%
\begin{pgfscope}%
\pgfpathrectangle{\pgfqpoint{0.050000in}{0.050000in}}{\pgfqpoint{4.900000in}{4.900000in}}%
\pgfusepath{clip}%
\pgfsetbuttcap%
\pgfsetroundjoin%
\definecolor{currentfill}{rgb}{0.212226,0.212226,0.212226}%
\pgfsetfillcolor{currentfill}%
\pgfsetfillopacity{0.900000}%
\pgfsetlinewidth{0.501875pt}%
\definecolor{currentstroke}{rgb}{0.200000,0.200000,0.200000}%
\pgfsetstrokecolor{currentstroke}%
\pgfsetdash{}{0pt}%
\pgfpathmoveto{\pgfqpoint{0.995326in}{3.423698in}}%
\pgfpathlineto{\pgfqpoint{1.044427in}{3.376959in}}%
\pgfpathlineto{\pgfqpoint{1.074154in}{3.358059in}}%
\pgfpathlineto{\pgfqpoint{1.025040in}{3.404831in}}%
\pgfpathlineto{\pgfqpoint{0.995326in}{3.423698in}}%
\pgfpathclose%
\pgfusepath{stroke,fill}%
\end{pgfscope}%
\begin{pgfscope}%
\pgfpathrectangle{\pgfqpoint{0.050000in}{0.050000in}}{\pgfqpoint{4.900000in}{4.900000in}}%
\pgfusepath{clip}%
\pgfsetbuttcap%
\pgfsetroundjoin%
\definecolor{currentfill}{rgb}{0.836340,0.836340,0.836340}%
\pgfsetfillcolor{currentfill}%
\pgfsetfillopacity{0.900000}%
\pgfsetlinewidth{0.501875pt}%
\definecolor{currentstroke}{rgb}{0.200000,0.200000,0.200000}%
\pgfsetstrokecolor{currentstroke}%
\pgfsetdash{}{0pt}%
\pgfpathmoveto{\pgfqpoint{2.701537in}{2.047019in}}%
\pgfpathlineto{\pgfqpoint{2.747018in}{2.005375in}}%
\pgfpathlineto{\pgfqpoint{2.779395in}{1.984818in}}%
\pgfpathlineto{\pgfqpoint{2.733910in}{2.026329in}}%
\pgfpathlineto{\pgfqpoint{2.701537in}{2.047019in}}%
\pgfpathclose%
\pgfusepath{stroke,fill}%
\end{pgfscope}%
\begin{pgfscope}%
\pgfpathrectangle{\pgfqpoint{0.050000in}{0.050000in}}{\pgfqpoint{4.900000in}{4.900000in}}%
\pgfusepath{clip}%
\pgfsetbuttcap%
\pgfsetroundjoin%
\definecolor{currentfill}{rgb}{0.491349,0.491349,0.491349}%
\pgfsetfillcolor{currentfill}%
\pgfsetfillopacity{0.900000}%
\pgfsetlinewidth{0.501875pt}%
\definecolor{currentstroke}{rgb}{0.200000,0.200000,0.200000}%
\pgfsetstrokecolor{currentstroke}%
\pgfsetdash{}{0pt}%
\pgfpathmoveto{\pgfqpoint{1.715363in}{2.841735in}}%
\pgfpathlineto{\pgfqpoint{1.762936in}{2.796229in}}%
\pgfpathlineto{\pgfqpoint{1.793782in}{2.776427in}}%
\pgfpathlineto{\pgfqpoint{1.746200in}{2.821961in}}%
\pgfpathlineto{\pgfqpoint{1.715363in}{2.841735in}}%
\pgfpathclose%
\pgfusepath{stroke,fill}%
\end{pgfscope}%
\begin{pgfscope}%
\pgfpathrectangle{\pgfqpoint{0.050000in}{0.050000in}}{\pgfqpoint{4.900000in}{4.900000in}}%
\pgfusepath{clip}%
\pgfsetbuttcap%
\pgfsetroundjoin%
\definecolor{currentfill}{rgb}{0.760554,0.760554,0.760554}%
\pgfsetfillcolor{currentfill}%
\pgfsetfillopacity{0.900000}%
\pgfsetlinewidth{0.501875pt}%
\definecolor{currentstroke}{rgb}{0.200000,0.200000,0.200000}%
\pgfsetstrokecolor{currentstroke}%
\pgfsetdash{}{0pt}%
\pgfpathmoveto{\pgfqpoint{2.435495in}{2.259918in}}%
\pgfpathlineto{\pgfqpoint{2.481540in}{2.215896in}}%
\pgfpathlineto{\pgfqpoint{2.513506in}{2.195276in}}%
\pgfpathlineto{\pgfqpoint{2.467457in}{2.239283in}}%
\pgfpathlineto{\pgfqpoint{2.435495in}{2.259918in}}%
\pgfpathclose%
\pgfusepath{stroke,fill}%
\end{pgfscope}%
\begin{pgfscope}%
\pgfpathrectangle{\pgfqpoint{0.050000in}{0.050000in}}{\pgfqpoint{4.900000in}{4.900000in}}%
\pgfusepath{clip}%
\pgfsetbuttcap%
\pgfsetroundjoin%
\definecolor{currentfill}{rgb}{0.998155,0.998155,0.998155}%
\pgfsetfillcolor{currentfill}%
\pgfsetfillopacity{0.900000}%
\pgfsetlinewidth{0.501875pt}%
\definecolor{currentstroke}{rgb}{0.200000,0.200000,0.200000}%
\pgfsetstrokecolor{currentstroke}%
\pgfsetdash{}{0pt}%
\pgfpathmoveto{\pgfqpoint{4.260380in}{1.377162in}}%
\pgfpathlineto{\pgfqpoint{4.304192in}{1.388076in}}%
\pgfpathlineto{\pgfqpoint{4.339043in}{1.367171in}}%
\pgfpathlineto{\pgfqpoint{4.295206in}{1.356328in}}%
\pgfpathlineto{\pgfqpoint{4.260380in}{1.377162in}}%
\pgfpathclose%
\pgfusepath{stroke,fill}%
\end{pgfscope}%
\begin{pgfscope}%
\pgfpathrectangle{\pgfqpoint{0.050000in}{0.050000in}}{\pgfqpoint{4.900000in}{4.900000in}}%
\pgfusepath{clip}%
\pgfsetbuttcap%
\pgfsetroundjoin%
\definecolor{currentfill}{rgb}{0.898378,0.898378,0.898378}%
\pgfsetfillcolor{currentfill}%
\pgfsetfillopacity{0.900000}%
\pgfsetlinewidth{0.501875pt}%
\definecolor{currentstroke}{rgb}{0.200000,0.200000,0.200000}%
\pgfsetstrokecolor{currentstroke}%
\pgfsetdash{}{0pt}%
\pgfpathmoveto{\pgfqpoint{2.967596in}{1.845682in}}%
\pgfpathlineto{\pgfqpoint{3.012575in}{1.810469in}}%
\pgfpathlineto{\pgfqpoint{3.045379in}{1.790323in}}%
\pgfpathlineto{\pgfqpoint{3.000393in}{1.825310in}}%
\pgfpathlineto{\pgfqpoint{2.967596in}{1.845682in}}%
\pgfpathclose%
\pgfusepath{stroke,fill}%
\end{pgfscope}%
\begin{pgfscope}%
\pgfpathrectangle{\pgfqpoint{0.050000in}{0.050000in}}{\pgfqpoint{4.900000in}{4.900000in}}%
\pgfusepath{clip}%
\pgfsetbuttcap%
\pgfsetroundjoin%
\definecolor{currentfill}{rgb}{0.395663,0.395663,0.395663}%
\pgfsetfillcolor{currentfill}%
\pgfsetfillopacity{0.900000}%
\pgfsetlinewidth{0.501875pt}%
\definecolor{currentstroke}{rgb}{0.200000,0.200000,0.200000}%
\pgfsetstrokecolor{currentstroke}%
\pgfsetdash{}{0pt}%
\pgfpathmoveto{\pgfqpoint{1.449135in}{3.057552in}}%
\pgfpathlineto{\pgfqpoint{1.497282in}{3.011580in}}%
\pgfpathlineto{\pgfqpoint{1.527720in}{2.992109in}}%
\pgfpathlineto{\pgfqpoint{1.479563in}{3.038111in}}%
\pgfpathlineto{\pgfqpoint{1.449135in}{3.057552in}}%
\pgfpathclose%
\pgfusepath{stroke,fill}%
\end{pgfscope}%
\begin{pgfscope}%
\pgfpathrectangle{\pgfqpoint{0.050000in}{0.050000in}}{\pgfqpoint{4.900000in}{4.900000in}}%
\pgfusepath{clip}%
\pgfsetbuttcap%
\pgfsetroundjoin%
\definecolor{currentfill}{rgb}{0.992618,0.992618,0.992618}%
\pgfsetfillcolor{currentfill}%
\pgfsetfillopacity{0.900000}%
\pgfsetlinewidth{0.501875pt}%
\definecolor{currentstroke}{rgb}{0.200000,0.200000,0.200000}%
\pgfsetstrokecolor{currentstroke}%
\pgfsetdash{}{0pt}%
\pgfpathmoveto{\pgfqpoint{3.980865in}{1.405161in}}%
\pgfpathlineto{\pgfqpoint{4.024892in}{1.411605in}}%
\pgfpathlineto{\pgfqpoint{4.059316in}{1.391201in}}%
\pgfpathlineto{\pgfqpoint{4.015266in}{1.384834in}}%
\pgfpathlineto{\pgfqpoint{3.980865in}{1.405161in}}%
\pgfpathclose%
\pgfusepath{stroke,fill}%
\end{pgfscope}%
\begin{pgfscope}%
\pgfpathrectangle{\pgfqpoint{0.050000in}{0.050000in}}{\pgfqpoint{4.900000in}{4.900000in}}%
\pgfusepath{clip}%
\pgfsetbuttcap%
\pgfsetroundjoin%
\definecolor{currentfill}{rgb}{0.657809,0.657809,0.657809}%
\pgfsetfillcolor{currentfill}%
\pgfsetfillopacity{0.900000}%
\pgfsetlinewidth{0.501875pt}%
\definecolor{currentstroke}{rgb}{0.200000,0.200000,0.200000}%
\pgfsetstrokecolor{currentstroke}%
\pgfsetdash{}{0pt}%
\pgfpathmoveto{\pgfqpoint{2.169371in}{2.475442in}}%
\pgfpathlineto{\pgfqpoint{2.215991in}{2.430711in}}%
\pgfpathlineto{\pgfqpoint{2.247548in}{2.410341in}}%
\pgfpathlineto{\pgfqpoint{2.200924in}{2.455095in}}%
\pgfpathlineto{\pgfqpoint{2.169371in}{2.475442in}}%
\pgfpathclose%
\pgfusepath{stroke,fill}%
\end{pgfscope}%
\begin{pgfscope}%
\pgfpathrectangle{\pgfqpoint{0.050000in}{0.050000in}}{\pgfqpoint{4.900000in}{4.900000in}}%
\pgfusepath{clip}%
\pgfsetbuttcap%
\pgfsetroundjoin%
\definecolor{currentfill}{rgb}{0.295271,0.295271,0.295271}%
\pgfsetfillcolor{currentfill}%
\pgfsetfillopacity{0.900000}%
\pgfsetlinewidth{0.501875pt}%
\definecolor{currentstroke}{rgb}{0.200000,0.200000,0.200000}%
\pgfsetstrokecolor{currentstroke}%
\pgfsetdash{}{0pt}%
\pgfpathmoveto{\pgfqpoint{1.182817in}{3.273445in}}%
\pgfpathlineto{\pgfqpoint{1.231538in}{3.227006in}}%
\pgfpathlineto{\pgfqpoint{1.261569in}{3.207866in}}%
\pgfpathlineto{\pgfqpoint{1.212835in}{3.254336in}}%
\pgfpathlineto{\pgfqpoint{1.182817in}{3.273445in}}%
\pgfpathclose%
\pgfusepath{stroke,fill}%
\end{pgfscope}%
\begin{pgfscope}%
\pgfpathrectangle{\pgfqpoint{0.050000in}{0.050000in}}{\pgfqpoint{4.900000in}{4.900000in}}%
\pgfusepath{clip}%
\pgfsetbuttcap%
\pgfsetroundjoin%
\definecolor{currentfill}{rgb}{0.555940,0.555940,0.555940}%
\pgfsetfillcolor{currentfill}%
\pgfsetfillopacity{0.900000}%
\pgfsetlinewidth{0.501875pt}%
\definecolor{currentstroke}{rgb}{0.200000,0.200000,0.200000}%
\pgfsetstrokecolor{currentstroke}%
\pgfsetdash{}{0pt}%
\pgfpathmoveto{\pgfqpoint{1.903157in}{2.691240in}}%
\pgfpathlineto{\pgfqpoint{1.950351in}{2.646036in}}%
\pgfpathlineto{\pgfqpoint{1.981501in}{2.625993in}}%
\pgfpathlineto{\pgfqpoint{1.934300in}{2.671225in}}%
\pgfpathlineto{\pgfqpoint{1.903157in}{2.691240in}}%
\pgfpathclose%
\pgfusepath{stroke,fill}%
\end{pgfscope}%
\begin{pgfscope}%
\pgfpathrectangle{\pgfqpoint{0.050000in}{0.050000in}}{\pgfqpoint{4.900000in}{4.900000in}}%
\pgfusepath{clip}%
\pgfsetbuttcap%
\pgfsetroundjoin%
\definecolor{currentfill}{rgb}{0.966782,0.966782,0.966782}%
\pgfsetfillcolor{currentfill}%
\pgfsetfillopacity{0.900000}%
\pgfsetlinewidth{0.501875pt}%
\definecolor{currentstroke}{rgb}{0.200000,0.200000,0.200000}%
\pgfsetstrokecolor{currentstroke}%
\pgfsetdash{}{0pt}%
\pgfpathmoveto{\pgfqpoint{3.467684in}{1.548442in}}%
\pgfpathlineto{\pgfqpoint{3.512044in}{1.535456in}}%
\pgfpathlineto{\pgfqpoint{3.545669in}{1.515780in}}%
\pgfpathlineto{\pgfqpoint{3.501291in}{1.528695in}}%
\pgfpathlineto{\pgfqpoint{3.467684in}{1.548442in}}%
\pgfpathclose%
\pgfusepath{stroke,fill}%
\end{pgfscope}%
\begin{pgfscope}%
\pgfpathrectangle{\pgfqpoint{0.050000in}{0.050000in}}{\pgfqpoint{4.900000in}{4.900000in}}%
\pgfusepath{clip}%
\pgfsetbuttcap%
\pgfsetroundjoin%
\definecolor{currentfill}{rgb}{0.972318,0.972318,0.972318}%
\pgfsetfillcolor{currentfill}%
\pgfsetfillopacity{0.900000}%
\pgfsetlinewidth{0.501875pt}%
\definecolor{currentstroke}{rgb}{0.200000,0.200000,0.200000}%
\pgfsetstrokecolor{currentstroke}%
\pgfsetdash{}{0pt}%
\pgfpathmoveto{\pgfqpoint{3.545669in}{1.515780in}}%
\pgfpathlineto{\pgfqpoint{3.589971in}{1.506466in}}%
\pgfpathlineto{\pgfqpoint{3.623722in}{1.486740in}}%
\pgfpathlineto{\pgfqpoint{3.579401in}{1.496021in}}%
\pgfpathlineto{\pgfqpoint{3.545669in}{1.515780in}}%
\pgfpathclose%
\pgfusepath{stroke,fill}%
\end{pgfscope}%
\begin{pgfscope}%
\pgfpathrectangle{\pgfqpoint{0.050000in}{0.050000in}}{\pgfqpoint{4.900000in}{4.900000in}}%
\pgfusepath{clip}%
\pgfsetbuttcap%
\pgfsetroundjoin%
\definecolor{currentfill}{rgb}{0.959400,0.959400,0.959400}%
\pgfsetfillcolor{currentfill}%
\pgfsetfillopacity{0.900000}%
\pgfsetlinewidth{0.501875pt}%
\definecolor{currentstroke}{rgb}{0.200000,0.200000,0.200000}%
\pgfsetstrokecolor{currentstroke}%
\pgfsetdash{}{0pt}%
\pgfpathmoveto{\pgfqpoint{3.389755in}{1.584945in}}%
\pgfpathlineto{\pgfqpoint{3.434184in}{1.568107in}}%
\pgfpathlineto{\pgfqpoint{3.467684in}{1.548442in}}%
\pgfpathlineto{\pgfqpoint{3.423239in}{1.565168in}}%
\pgfpathlineto{\pgfqpoint{3.389755in}{1.584945in}}%
\pgfpathclose%
\pgfusepath{stroke,fill}%
\end{pgfscope}%
\begin{pgfscope}%
\pgfpathrectangle{\pgfqpoint{0.050000in}{0.050000in}}{\pgfqpoint{4.900000in}{4.900000in}}%
\pgfusepath{clip}%
\pgfsetbuttcap%
\pgfsetroundjoin%
\definecolor{currentfill}{rgb}{0.173472,0.173472,0.173472}%
\pgfsetfillcolor{currentfill}%
\pgfsetfillopacity{0.900000}%
\pgfsetlinewidth{0.501875pt}%
\definecolor{currentstroke}{rgb}{0.200000,0.200000,0.200000}%
\pgfsetstrokecolor{currentstroke}%
\pgfsetdash{}{0pt}%
\pgfpathmoveto{\pgfqpoint{0.916408in}{3.489414in}}%
\pgfpathlineto{\pgfqpoint{0.965705in}{3.442507in}}%
\pgfpathlineto{\pgfqpoint{0.995326in}{3.423698in}}%
\pgfpathlineto{\pgfqpoint{0.946016in}{3.470638in}}%
\pgfpathlineto{\pgfqpoint{0.916408in}{3.489414in}}%
\pgfpathclose%
\pgfusepath{stroke,fill}%
\end{pgfscope}%
\begin{pgfscope}%
\pgfpathrectangle{\pgfqpoint{0.050000in}{0.050000in}}{\pgfqpoint{4.900000in}{4.900000in}}%
\pgfusepath{clip}%
\pgfsetbuttcap%
\pgfsetroundjoin%
\definecolor{currentfill}{rgb}{0.812226,0.812226,0.812226}%
\pgfsetfillcolor{currentfill}%
\pgfsetfillopacity{0.900000}%
\pgfsetlinewidth{0.501875pt}%
\definecolor{currentstroke}{rgb}{0.200000,0.200000,0.200000}%
\pgfsetstrokecolor{currentstroke}%
\pgfsetdash{}{0pt}%
\pgfpathmoveto{\pgfqpoint{2.623596in}{2.110414in}}%
\pgfpathlineto{\pgfqpoint{2.669264in}{2.067661in}}%
\pgfpathlineto{\pgfqpoint{2.701537in}{2.047019in}}%
\pgfpathlineto{\pgfqpoint{2.655865in}{2.089683in}}%
\pgfpathlineto{\pgfqpoint{2.623596in}{2.110414in}}%
\pgfpathclose%
\pgfusepath{stroke,fill}%
\end{pgfscope}%
\begin{pgfscope}%
\pgfpathrectangle{\pgfqpoint{0.050000in}{0.050000in}}{\pgfqpoint{4.900000in}{4.900000in}}%
\pgfusepath{clip}%
\pgfsetbuttcap%
\pgfsetroundjoin%
\definecolor{currentfill}{rgb}{0.994464,0.994464,0.994464}%
\pgfsetfillcolor{currentfill}%
\pgfsetfillopacity{0.900000}%
\pgfsetlinewidth{0.501875pt}%
\definecolor{currentstroke}{rgb}{0.200000,0.200000,0.200000}%
\pgfsetstrokecolor{currentstroke}%
\pgfsetdash{}{0pt}%
\pgfpathmoveto{\pgfqpoint{4.059316in}{1.391201in}}%
\pgfpathlineto{\pgfqpoint{4.103301in}{1.399140in}}%
\pgfpathlineto{\pgfqpoint{4.137859in}{1.378589in}}%
\pgfpathlineto{\pgfqpoint{4.093850in}{1.370728in}}%
\pgfpathlineto{\pgfqpoint{4.059316in}{1.391201in}}%
\pgfpathclose%
\pgfusepath{stroke,fill}%
\end{pgfscope}%
\begin{pgfscope}%
\pgfpathrectangle{\pgfqpoint{0.050000in}{0.050000in}}{\pgfqpoint{4.900000in}{4.900000in}}%
\pgfusepath{clip}%
\pgfsetbuttcap%
\pgfsetroundjoin%
\definecolor{currentfill}{rgb}{0.878570,0.878570,0.878570}%
\pgfsetfillcolor{currentfill}%
\pgfsetfillopacity{0.900000}%
\pgfsetlinewidth{0.501875pt}%
\definecolor{currentstroke}{rgb}{0.200000,0.200000,0.200000}%
\pgfsetstrokecolor{currentstroke}%
\pgfsetdash{}{0pt}%
\pgfpathmoveto{\pgfqpoint{2.889766in}{1.903821in}}%
\pgfpathlineto{\pgfqpoint{2.934903in}{1.866002in}}%
\pgfpathlineto{\pgfqpoint{2.967596in}{1.845682in}}%
\pgfpathlineto{\pgfqpoint{2.922454in}{1.883294in}}%
\pgfpathlineto{\pgfqpoint{2.889766in}{1.903821in}}%
\pgfpathclose%
\pgfusepath{stroke,fill}%
\end{pgfscope}%
\begin{pgfscope}%
\pgfpathrectangle{\pgfqpoint{0.050000in}{0.050000in}}{\pgfqpoint{4.900000in}{4.900000in}}%
\pgfusepath{clip}%
\pgfsetbuttcap%
\pgfsetroundjoin%
\definecolor{currentfill}{rgb}{0.977855,0.977855,0.977855}%
\pgfsetfillcolor{currentfill}%
\pgfsetfillopacity{0.900000}%
\pgfsetlinewidth{0.501875pt}%
\definecolor{currentstroke}{rgb}{0.200000,0.200000,0.200000}%
\pgfsetstrokecolor{currentstroke}%
\pgfsetdash{}{0pt}%
\pgfpathmoveto{\pgfqpoint{3.623722in}{1.486740in}}%
\pgfpathlineto{\pgfqpoint{3.667975in}{1.480846in}}%
\pgfpathlineto{\pgfqpoint{3.701854in}{1.461037in}}%
\pgfpathlineto{\pgfqpoint{3.657581in}{1.466931in}}%
\pgfpathlineto{\pgfqpoint{3.623722in}{1.486740in}}%
\pgfpathclose%
\pgfusepath{stroke,fill}%
\end{pgfscope}%
\begin{pgfscope}%
\pgfpathrectangle{\pgfqpoint{0.050000in}{0.050000in}}{\pgfqpoint{4.900000in}{4.900000in}}%
\pgfusepath{clip}%
\pgfsetbuttcap%
\pgfsetroundjoin%
\definecolor{currentfill}{rgb}{0.456901,0.456901,0.456901}%
\pgfsetfillcolor{currentfill}%
\pgfsetfillopacity{0.900000}%
\pgfsetlinewidth{0.501875pt}%
\definecolor{currentstroke}{rgb}{0.200000,0.200000,0.200000}%
\pgfsetstrokecolor{currentstroke}%
\pgfsetdash{}{0pt}%
\pgfpathmoveto{\pgfqpoint{1.636853in}{2.907118in}}%
\pgfpathlineto{\pgfqpoint{1.684621in}{2.861447in}}%
\pgfpathlineto{\pgfqpoint{1.715363in}{2.841735in}}%
\pgfpathlineto{\pgfqpoint{1.667586in}{2.887435in}}%
\pgfpathlineto{\pgfqpoint{1.636853in}{2.907118in}}%
\pgfpathclose%
\pgfusepath{stroke,fill}%
\end{pgfscope}%
\begin{pgfscope}%
\pgfpathrectangle{\pgfqpoint{0.050000in}{0.050000in}}{\pgfqpoint{4.900000in}{4.900000in}}%
\pgfusepath{clip}%
\pgfsetbuttcap%
\pgfsetroundjoin%
\definecolor{currentfill}{rgb}{0.952018,0.952018,0.952018}%
\pgfsetfillcolor{currentfill}%
\pgfsetfillopacity{0.900000}%
\pgfsetlinewidth{0.501875pt}%
\definecolor{currentstroke}{rgb}{0.200000,0.200000,0.200000}%
\pgfsetstrokecolor{currentstroke}%
\pgfsetdash{}{0pt}%
\pgfpathmoveto{\pgfqpoint{3.311869in}{1.625415in}}%
\pgfpathlineto{\pgfqpoint{3.356378in}{1.604642in}}%
\pgfpathlineto{\pgfqpoint{3.389755in}{1.584945in}}%
\pgfpathlineto{\pgfqpoint{3.345231in}{1.605566in}}%
\pgfpathlineto{\pgfqpoint{3.311869in}{1.625415in}}%
\pgfpathclose%
\pgfusepath{stroke,fill}%
\end{pgfscope}%
\begin{pgfscope}%
\pgfpathrectangle{\pgfqpoint{0.050000in}{0.050000in}}{\pgfqpoint{4.900000in}{4.900000in}}%
\pgfusepath{clip}%
\pgfsetbuttcap%
\pgfsetroundjoin%
\definecolor{currentfill}{rgb}{0.734579,0.734579,0.734579}%
\pgfsetfillcolor{currentfill}%
\pgfsetfillopacity{0.900000}%
\pgfsetlinewidth{0.501875pt}%
\definecolor{currentstroke}{rgb}{0.200000,0.200000,0.200000}%
\pgfsetstrokecolor{currentstroke}%
\pgfsetdash{}{0pt}%
\pgfpathmoveto{\pgfqpoint{2.357395in}{2.324832in}}%
\pgfpathlineto{\pgfqpoint{2.403633in}{2.280495in}}%
\pgfpathlineto{\pgfqpoint{2.435495in}{2.259918in}}%
\pgfpathlineto{\pgfqpoint{2.389253in}{2.304260in}}%
\pgfpathlineto{\pgfqpoint{2.357395in}{2.324832in}}%
\pgfpathclose%
\pgfusepath{stroke,fill}%
\end{pgfscope}%
\begin{pgfscope}%
\pgfpathrectangle{\pgfqpoint{0.050000in}{0.050000in}}{\pgfqpoint{4.900000in}{4.900000in}}%
\pgfusepath{clip}%
\pgfsetbuttcap%
\pgfsetroundjoin%
\definecolor{currentfill}{rgb}{0.981546,0.981546,0.981546}%
\pgfsetfillcolor{currentfill}%
\pgfsetfillopacity{0.900000}%
\pgfsetlinewidth{0.501875pt}%
\definecolor{currentstroke}{rgb}{0.200000,0.200000,0.200000}%
\pgfsetstrokecolor{currentstroke}%
\pgfsetdash{}{0pt}%
\pgfpathmoveto{\pgfqpoint{3.701854in}{1.461037in}}%
\pgfpathlineto{\pgfqpoint{3.746062in}{1.458261in}}%
\pgfpathlineto{\pgfqpoint{3.780070in}{1.438343in}}%
\pgfpathlineto{\pgfqpoint{3.735840in}{1.441147in}}%
\pgfpathlineto{\pgfqpoint{3.701854in}{1.461037in}}%
\pgfpathclose%
\pgfusepath{stroke,fill}%
\end{pgfscope}%
\begin{pgfscope}%
\pgfpathrectangle{\pgfqpoint{0.050000in}{0.050000in}}{\pgfqpoint{4.900000in}{4.900000in}}%
\pgfusepath{clip}%
\pgfsetbuttcap%
\pgfsetroundjoin%
\definecolor{currentfill}{rgb}{0.363183,0.363183,0.363183}%
\pgfsetfillcolor{currentfill}%
\pgfsetfillopacity{0.900000}%
\pgfsetlinewidth{0.501875pt}%
\definecolor{currentstroke}{rgb}{0.200000,0.200000,0.200000}%
\pgfsetstrokecolor{currentstroke}%
\pgfsetdash{}{0pt}%
\pgfpathmoveto{\pgfqpoint{1.370459in}{3.123071in}}%
\pgfpathlineto{\pgfqpoint{1.418801in}{3.076933in}}%
\pgfpathlineto{\pgfqpoint{1.449135in}{3.057552in}}%
\pgfpathlineto{\pgfqpoint{1.400782in}{3.103721in}}%
\pgfpathlineto{\pgfqpoint{1.370459in}{3.123071in}}%
\pgfpathclose%
\pgfusepath{stroke,fill}%
\end{pgfscope}%
\begin{pgfscope}%
\pgfpathrectangle{\pgfqpoint{0.050000in}{0.050000in}}{\pgfqpoint{4.900000in}{4.900000in}}%
\pgfusepath{clip}%
\pgfsetbuttcap%
\pgfsetroundjoin%
\definecolor{currentfill}{rgb}{0.942791,0.942791,0.942791}%
\pgfsetfillcolor{currentfill}%
\pgfsetfillopacity{0.900000}%
\pgfsetlinewidth{0.501875pt}%
\definecolor{currentstroke}{rgb}{0.200000,0.200000,0.200000}%
\pgfsetstrokecolor{currentstroke}%
\pgfsetdash{}{0pt}%
\pgfpathmoveto{\pgfqpoint{3.234007in}{1.669867in}}%
\pgfpathlineto{\pgfqpoint{3.278612in}{1.645188in}}%
\pgfpathlineto{\pgfqpoint{3.311869in}{1.625415in}}%
\pgfpathlineto{\pgfqpoint{3.267250in}{1.649907in}}%
\pgfpathlineto{\pgfqpoint{3.234007in}{1.669867in}}%
\pgfpathclose%
\pgfusepath{stroke,fill}%
\end{pgfscope}%
\begin{pgfscope}%
\pgfpathrectangle{\pgfqpoint{0.050000in}{0.050000in}}{\pgfqpoint{4.900000in}{4.900000in}}%
\pgfusepath{clip}%
\pgfsetbuttcap%
\pgfsetroundjoin%
\definecolor{currentfill}{rgb}{0.624221,0.624221,0.624221}%
\pgfsetfillcolor{currentfill}%
\pgfsetfillopacity{0.900000}%
\pgfsetlinewidth{0.501875pt}%
\definecolor{currentstroke}{rgb}{0.200000,0.200000,0.200000}%
\pgfsetstrokecolor{currentstroke}%
\pgfsetdash{}{0pt}%
\pgfpathmoveto{\pgfqpoint{2.091105in}{2.540626in}}%
\pgfpathlineto{\pgfqpoint{2.137917in}{2.495725in}}%
\pgfpathlineto{\pgfqpoint{2.169371in}{2.475442in}}%
\pgfpathlineto{\pgfqpoint{2.122553in}{2.520368in}}%
\pgfpathlineto{\pgfqpoint{2.091105in}{2.540626in}}%
\pgfpathclose%
\pgfusepath{stroke,fill}%
\end{pgfscope}%
\begin{pgfscope}%
\pgfpathrectangle{\pgfqpoint{0.050000in}{0.050000in}}{\pgfqpoint{4.900000in}{4.900000in}}%
\pgfusepath{clip}%
\pgfsetbuttcap%
\pgfsetroundjoin%
\definecolor{currentfill}{rgb}{0.996309,0.996309,0.996309}%
\pgfsetfillcolor{currentfill}%
\pgfsetfillopacity{0.900000}%
\pgfsetlinewidth{0.501875pt}%
\definecolor{currentstroke}{rgb}{0.200000,0.200000,0.200000}%
\pgfsetstrokecolor{currentstroke}%
\pgfsetdash{}{0pt}%
\pgfpathmoveto{\pgfqpoint{4.137859in}{1.378589in}}%
\pgfpathlineto{\pgfqpoint{4.181798in}{1.387742in}}%
\pgfpathlineto{\pgfqpoint{4.216490in}{1.367048in}}%
\pgfpathlineto{\pgfqpoint{4.172526in}{1.357971in}}%
\pgfpathlineto{\pgfqpoint{4.137859in}{1.378589in}}%
\pgfpathclose%
\pgfusepath{stroke,fill}%
\end{pgfscope}%
\begin{pgfscope}%
\pgfpathrectangle{\pgfqpoint{0.050000in}{0.050000in}}{\pgfqpoint{4.900000in}{4.900000in}}%
\pgfusepath{clip}%
\pgfsetbuttcap%
\pgfsetroundjoin%
\definecolor{currentfill}{rgb}{0.985236,0.985236,0.985236}%
\pgfsetfillcolor{currentfill}%
\pgfsetfillopacity{0.900000}%
\pgfsetlinewidth{0.501875pt}%
\definecolor{currentstroke}{rgb}{0.200000,0.200000,0.200000}%
\pgfsetstrokecolor{currentstroke}%
\pgfsetdash{}{0pt}%
\pgfpathmoveto{\pgfqpoint{3.780070in}{1.438343in}}%
\pgfpathlineto{\pgfqpoint{3.824238in}{1.438354in}}%
\pgfpathlineto{\pgfqpoint{3.858376in}{1.418310in}}%
\pgfpathlineto{\pgfqpoint{3.814186in}{1.418346in}}%
\pgfpathlineto{\pgfqpoint{3.780070in}{1.438343in}}%
\pgfpathclose%
\pgfusepath{stroke,fill}%
\end{pgfscope}%
\begin{pgfscope}%
\pgfpathrectangle{\pgfqpoint{0.050000in}{0.050000in}}{\pgfqpoint{4.900000in}{4.900000in}}%
\pgfusepath{clip}%
\pgfsetbuttcap%
\pgfsetroundjoin%
\definecolor{currentfill}{rgb}{0.256517,0.256517,0.256517}%
\pgfsetfillcolor{currentfill}%
\pgfsetfillopacity{0.900000}%
\pgfsetlinewidth{0.501875pt}%
\definecolor{currentstroke}{rgb}{0.200000,0.200000,0.200000}%
\pgfsetstrokecolor{currentstroke}%
\pgfsetdash{}{0pt}%
\pgfpathmoveto{\pgfqpoint{1.103974in}{3.339100in}}%
\pgfpathlineto{\pgfqpoint{1.152892in}{3.292495in}}%
\pgfpathlineto{\pgfqpoint{1.182817in}{3.273445in}}%
\pgfpathlineto{\pgfqpoint{1.133886in}{3.320083in}}%
\pgfpathlineto{\pgfqpoint{1.103974in}{3.339100in}}%
\pgfpathclose%
\pgfusepath{stroke,fill}%
\end{pgfscope}%
\begin{pgfscope}%
\pgfpathrectangle{\pgfqpoint{0.050000in}{0.050000in}}{\pgfqpoint{4.900000in}{4.900000in}}%
\pgfusepath{clip}%
\pgfsetbuttcap%
\pgfsetroundjoin%
\definecolor{currentfill}{rgb}{0.521492,0.521492,0.521492}%
\pgfsetfillcolor{currentfill}%
\pgfsetfillopacity{0.900000}%
\pgfsetlinewidth{0.501875pt}%
\definecolor{currentstroke}{rgb}{0.200000,0.200000,0.200000}%
\pgfsetstrokecolor{currentstroke}%
\pgfsetdash{}{0pt}%
\pgfpathmoveto{\pgfqpoint{1.824724in}{2.756563in}}%
\pgfpathlineto{\pgfqpoint{1.872111in}{2.711194in}}%
\pgfpathlineto{\pgfqpoint{1.903157in}{2.691240in}}%
\pgfpathlineto{\pgfqpoint{1.855763in}{2.736638in}}%
\pgfpathlineto{\pgfqpoint{1.824724in}{2.756563in}}%
\pgfpathclose%
\pgfusepath{stroke,fill}%
\end{pgfscope}%
\begin{pgfscope}%
\pgfpathrectangle{\pgfqpoint{0.050000in}{0.050000in}}{\pgfqpoint{4.900000in}{4.900000in}}%
\pgfusepath{clip}%
\pgfsetbuttcap%
\pgfsetroundjoin%
\definecolor{currentfill}{rgb}{0.929504,0.929504,0.929504}%
\pgfsetfillcolor{currentfill}%
\pgfsetfillopacity{0.900000}%
\pgfsetlinewidth{0.501875pt}%
\definecolor{currentstroke}{rgb}{0.200000,0.200000,0.200000}%
\pgfsetstrokecolor{currentstroke}%
\pgfsetdash{}{0pt}%
\pgfpathmoveto{\pgfqpoint{3.156153in}{1.718184in}}%
\pgfpathlineto{\pgfqpoint{3.200869in}{1.689756in}}%
\pgfpathlineto{\pgfqpoint{3.234007in}{1.669867in}}%
\pgfpathlineto{\pgfqpoint{3.189280in}{1.698082in}}%
\pgfpathlineto{\pgfqpoint{3.156153in}{1.718184in}}%
\pgfpathclose%
\pgfusepath{stroke,fill}%
\end{pgfscope}%
\begin{pgfscope}%
\pgfpathrectangle{\pgfqpoint{0.050000in}{0.050000in}}{\pgfqpoint{4.900000in}{4.900000in}}%
\pgfusepath{clip}%
\pgfsetbuttcap%
\pgfsetroundjoin%
\definecolor{currentfill}{rgb}{0.858762,0.858762,0.858762}%
\pgfsetfillcolor{currentfill}%
\pgfsetfillopacity{0.900000}%
\pgfsetlinewidth{0.501875pt}%
\definecolor{currentstroke}{rgb}{0.200000,0.200000,0.200000}%
\pgfsetstrokecolor{currentstroke}%
\pgfsetdash{}{0pt}%
\pgfpathmoveto{\pgfqpoint{2.811874in}{1.964213in}}%
\pgfpathlineto{\pgfqpoint{2.857181in}{1.924299in}}%
\pgfpathlineto{\pgfqpoint{2.889766in}{1.903821in}}%
\pgfpathlineto{\pgfqpoint{2.844455in}{1.943560in}}%
\pgfpathlineto{\pgfqpoint{2.811874in}{1.964213in}}%
\pgfpathclose%
\pgfusepath{stroke,fill}%
\end{pgfscope}%
\begin{pgfscope}%
\pgfpathrectangle{\pgfqpoint{0.050000in}{0.050000in}}{\pgfqpoint{4.900000in}{4.900000in}}%
\pgfusepath{clip}%
\pgfsetbuttcap%
\pgfsetroundjoin%
\definecolor{currentfill}{rgb}{0.788112,0.788112,0.788112}%
\pgfsetfillcolor{currentfill}%
\pgfsetfillopacity{0.900000}%
\pgfsetlinewidth{0.501875pt}%
\definecolor{currentstroke}{rgb}{0.200000,0.200000,0.200000}%
\pgfsetstrokecolor{currentstroke}%
\pgfsetdash{}{0pt}%
\pgfpathmoveto{\pgfqpoint{2.545571in}{2.174602in}}%
\pgfpathlineto{\pgfqpoint{2.591428in}{2.131094in}}%
\pgfpathlineto{\pgfqpoint{2.623596in}{2.110414in}}%
\pgfpathlineto{\pgfqpoint{2.577736in}{2.153872in}}%
\pgfpathlineto{\pgfqpoint{2.545571in}{2.174602in}}%
\pgfpathclose%
\pgfusepath{stroke,fill}%
\end{pgfscope}%
\begin{pgfscope}%
\pgfpathrectangle{\pgfqpoint{0.050000in}{0.050000in}}{\pgfqpoint{4.900000in}{4.900000in}}%
\pgfusepath{clip}%
\pgfsetbuttcap%
\pgfsetroundjoin%
\definecolor{currentfill}{rgb}{0.132011,0.132011,0.132011}%
\pgfsetfillcolor{currentfill}%
\pgfsetfillopacity{0.900000}%
\pgfsetlinewidth{0.501875pt}%
\definecolor{currentstroke}{rgb}{0.200000,0.200000,0.200000}%
\pgfsetstrokecolor{currentstroke}%
\pgfsetdash{}{0pt}%
\pgfpathmoveto{\pgfqpoint{0.837399in}{3.555205in}}%
\pgfpathlineto{\pgfqpoint{0.886892in}{3.508132in}}%
\pgfpathlineto{\pgfqpoint{0.916408in}{3.489414in}}%
\pgfpathlineto{\pgfqpoint{0.866901in}{3.536520in}}%
\pgfpathlineto{\pgfqpoint{0.837399in}{3.555205in}}%
\pgfpathclose%
\pgfusepath{stroke,fill}%
\end{pgfscope}%
\begin{pgfscope}%
\pgfpathrectangle{\pgfqpoint{0.050000in}{0.050000in}}{\pgfqpoint{4.900000in}{4.900000in}}%
\pgfusepath{clip}%
\pgfsetbuttcap%
\pgfsetroundjoin%
\definecolor{currentfill}{rgb}{0.428143,0.428143,0.428143}%
\pgfsetfillcolor{currentfill}%
\pgfsetfillopacity{0.900000}%
\pgfsetlinewidth{0.501875pt}%
\definecolor{currentstroke}{rgb}{0.200000,0.200000,0.200000}%
\pgfsetstrokecolor{currentstroke}%
\pgfsetdash{}{0pt}%
\pgfpathmoveto{\pgfqpoint{1.558253in}{2.972577in}}%
\pgfpathlineto{\pgfqpoint{1.606216in}{2.926740in}}%
\pgfpathlineto{\pgfqpoint{1.636853in}{2.907118in}}%
\pgfpathlineto{\pgfqpoint{1.588881in}{2.952984in}}%
\pgfpathlineto{\pgfqpoint{1.558253in}{2.972577in}}%
\pgfpathclose%
\pgfusepath{stroke,fill}%
\end{pgfscope}%
\begin{pgfscope}%
\pgfpathrectangle{\pgfqpoint{0.050000in}{0.050000in}}{\pgfqpoint{4.900000in}{4.900000in}}%
\pgfusepath{clip}%
\pgfsetbuttcap%
\pgfsetroundjoin%
\definecolor{currentfill}{rgb}{0.988927,0.988927,0.988927}%
\pgfsetfillcolor{currentfill}%
\pgfsetfillopacity{0.900000}%
\pgfsetlinewidth{0.501875pt}%
\definecolor{currentstroke}{rgb}{0.200000,0.200000,0.200000}%
\pgfsetstrokecolor{currentstroke}%
\pgfsetdash{}{0pt}%
\pgfpathmoveto{\pgfqpoint{3.858376in}{1.418310in}}%
\pgfpathlineto{\pgfqpoint{3.902505in}{1.420770in}}%
\pgfpathlineto{\pgfqpoint{3.936774in}{1.400587in}}%
\pgfpathlineto{\pgfqpoint{3.892622in}{1.398189in}}%
\pgfpathlineto{\pgfqpoint{3.858376in}{1.418310in}}%
\pgfpathclose%
\pgfusepath{stroke,fill}%
\end{pgfscope}%
\begin{pgfscope}%
\pgfpathrectangle{\pgfqpoint{0.050000in}{0.050000in}}{\pgfqpoint{4.900000in}{4.900000in}}%
\pgfusepath{clip}%
\pgfsetbuttcap%
\pgfsetroundjoin%
\definecolor{currentfill}{rgb}{0.696194,0.696194,0.696194}%
\pgfsetfillcolor{currentfill}%
\pgfsetfillopacity{0.900000}%
\pgfsetlinewidth{0.501875pt}%
\definecolor{currentstroke}{rgb}{0.200000,0.200000,0.200000}%
\pgfsetstrokecolor{currentstroke}%
\pgfsetdash{}{0pt}%
\pgfpathmoveto{\pgfqpoint{2.279204in}{2.389909in}}%
\pgfpathlineto{\pgfqpoint{2.325635in}{2.345343in}}%
\pgfpathlineto{\pgfqpoint{2.357395in}{2.324832in}}%
\pgfpathlineto{\pgfqpoint{2.310959in}{2.369414in}}%
\pgfpathlineto{\pgfqpoint{2.279204in}{2.389909in}}%
\pgfpathclose%
\pgfusepath{stroke,fill}%
\end{pgfscope}%
\begin{pgfscope}%
\pgfpathrectangle{\pgfqpoint{0.050000in}{0.050000in}}{\pgfqpoint{4.900000in}{4.900000in}}%
\pgfusepath{clip}%
\pgfsetbuttcap%
\pgfsetroundjoin%
\definecolor{currentfill}{rgb}{0.996309,0.996309,0.996309}%
\pgfsetfillcolor{currentfill}%
\pgfsetfillopacity{0.900000}%
\pgfsetlinewidth{0.501875pt}%
\definecolor{currentstroke}{rgb}{0.200000,0.200000,0.200000}%
\pgfsetstrokecolor{currentstroke}%
\pgfsetdash{}{0pt}%
\pgfpathmoveto{\pgfqpoint{4.216490in}{1.367048in}}%
\pgfpathlineto{\pgfqpoint{4.260380in}{1.377162in}}%
\pgfpathlineto{\pgfqpoint{4.295206in}{1.356328in}}%
\pgfpathlineto{\pgfqpoint{4.251292in}{1.346287in}}%
\pgfpathlineto{\pgfqpoint{4.216490in}{1.367048in}}%
\pgfpathclose%
\pgfusepath{stroke,fill}%
\end{pgfscope}%
\begin{pgfscope}%
\pgfpathrectangle{\pgfqpoint{0.050000in}{0.050000in}}{\pgfqpoint{4.900000in}{4.900000in}}%
\pgfusepath{clip}%
\pgfsetbuttcap%
\pgfsetroundjoin%
\definecolor{currentfill}{rgb}{0.915356,0.915356,0.915356}%
\pgfsetfillcolor{currentfill}%
\pgfsetfillopacity{0.900000}%
\pgfsetlinewidth{0.501875pt}%
\definecolor{currentstroke}{rgb}{0.200000,0.200000,0.200000}%
\pgfsetstrokecolor{currentstroke}%
\pgfsetdash{}{0pt}%
\pgfpathmoveto{\pgfqpoint{3.078288in}{1.770120in}}%
\pgfpathlineto{\pgfqpoint{3.123131in}{1.738221in}}%
\pgfpathlineto{\pgfqpoint{3.156153in}{1.718184in}}%
\pgfpathlineto{\pgfqpoint{3.111300in}{1.749857in}}%
\pgfpathlineto{\pgfqpoint{3.078288in}{1.770120in}}%
\pgfpathclose%
\pgfusepath{stroke,fill}%
\end{pgfscope}%
\begin{pgfscope}%
\pgfpathrectangle{\pgfqpoint{0.050000in}{0.050000in}}{\pgfqpoint{4.900000in}{4.900000in}}%
\pgfusepath{clip}%
\pgfsetbuttcap%
\pgfsetroundjoin%
\definecolor{currentfill}{rgb}{0.334764,0.334764,0.334764}%
\pgfsetfillcolor{currentfill}%
\pgfsetfillopacity{0.900000}%
\pgfsetlinewidth{0.501875pt}%
\definecolor{currentstroke}{rgb}{0.200000,0.200000,0.200000}%
\pgfsetstrokecolor{currentstroke}%
\pgfsetdash{}{0pt}%
\pgfpathmoveto{\pgfqpoint{1.291692in}{3.188666in}}%
\pgfpathlineto{\pgfqpoint{1.340230in}{3.142362in}}%
\pgfpathlineto{\pgfqpoint{1.370459in}{3.123071in}}%
\pgfpathlineto{\pgfqpoint{1.321910in}{3.169406in}}%
\pgfpathlineto{\pgfqpoint{1.291692in}{3.188666in}}%
\pgfpathclose%
\pgfusepath{stroke,fill}%
\end{pgfscope}%
\begin{pgfscope}%
\pgfpathrectangle{\pgfqpoint{0.050000in}{0.050000in}}{\pgfqpoint{4.900000in}{4.900000in}}%
\pgfusepath{clip}%
\pgfsetbuttcap%
\pgfsetroundjoin%
\definecolor{currentfill}{rgb}{0.586082,0.586082,0.586082}%
\pgfsetfillcolor{currentfill}%
\pgfsetfillopacity{0.900000}%
\pgfsetlinewidth{0.501875pt}%
\definecolor{currentstroke}{rgb}{0.200000,0.200000,0.200000}%
\pgfsetstrokecolor{currentstroke}%
\pgfsetdash{}{0pt}%
\pgfpathmoveto{\pgfqpoint{2.012748in}{2.605887in}}%
\pgfpathlineto{\pgfqpoint{2.059754in}{2.560821in}}%
\pgfpathlineto{\pgfqpoint{2.091105in}{2.540626in}}%
\pgfpathlineto{\pgfqpoint{2.044092in}{2.585719in}}%
\pgfpathlineto{\pgfqpoint{2.012748in}{2.605887in}}%
\pgfpathclose%
\pgfusepath{stroke,fill}%
\end{pgfscope}%
\begin{pgfscope}%
\pgfpathrectangle{\pgfqpoint{0.050000in}{0.050000in}}{\pgfqpoint{4.900000in}{4.900000in}}%
\pgfusepath{clip}%
\pgfsetbuttcap%
\pgfsetroundjoin%
\definecolor{currentfill}{rgb}{0.990773,0.990773,0.990773}%
\pgfsetfillcolor{currentfill}%
\pgfsetfillopacity{0.900000}%
\pgfsetlinewidth{0.501875pt}%
\definecolor{currentstroke}{rgb}{0.200000,0.200000,0.200000}%
\pgfsetstrokecolor{currentstroke}%
\pgfsetdash{}{0pt}%
\pgfpathmoveto{\pgfqpoint{3.936774in}{1.400587in}}%
\pgfpathlineto{\pgfqpoint{3.980865in}{1.405161in}}%
\pgfpathlineto{\pgfqpoint{4.015266in}{1.384834in}}%
\pgfpathlineto{\pgfqpoint{3.971152in}{1.380331in}}%
\pgfpathlineto{\pgfqpoint{3.936774in}{1.400587in}}%
\pgfpathclose%
\pgfusepath{stroke,fill}%
\end{pgfscope}%
\begin{pgfscope}%
\pgfpathrectangle{\pgfqpoint{0.050000in}{0.050000in}}{\pgfqpoint{4.900000in}{4.900000in}}%
\pgfusepath{clip}%
\pgfsetbuttcap%
\pgfsetroundjoin%
\definecolor{currentfill}{rgb}{0.212226,0.212226,0.212226}%
\pgfsetfillcolor{currentfill}%
\pgfsetfillopacity{0.900000}%
\pgfsetlinewidth{0.501875pt}%
\definecolor{currentstroke}{rgb}{0.200000,0.200000,0.200000}%
\pgfsetstrokecolor{currentstroke}%
\pgfsetdash{}{0pt}%
\pgfpathmoveto{\pgfqpoint{1.025040in}{3.404831in}}%
\pgfpathlineto{\pgfqpoint{1.074154in}{3.358059in}}%
\pgfpathlineto{\pgfqpoint{1.103974in}{3.339100in}}%
\pgfpathlineto{\pgfqpoint{1.054847in}{3.385905in}}%
\pgfpathlineto{\pgfqpoint{1.025040in}{3.404831in}}%
\pgfpathclose%
\pgfusepath{stroke,fill}%
\end{pgfscope}%
\begin{pgfscope}%
\pgfpathrectangle{\pgfqpoint{0.050000in}{0.050000in}}{\pgfqpoint{4.900000in}{4.900000in}}%
\pgfusepath{clip}%
\pgfsetbuttcap%
\pgfsetroundjoin%
\definecolor{currentfill}{rgb}{0.836340,0.836340,0.836340}%
\pgfsetfillcolor{currentfill}%
\pgfsetfillopacity{0.900000}%
\pgfsetlinewidth{0.501875pt}%
\definecolor{currentstroke}{rgb}{0.200000,0.200000,0.200000}%
\pgfsetstrokecolor{currentstroke}%
\pgfsetdash{}{0pt}%
\pgfpathmoveto{\pgfqpoint{2.733910in}{2.026329in}}%
\pgfpathlineto{\pgfqpoint{2.779395in}{1.984818in}}%
\pgfpathlineto{\pgfqpoint{2.811874in}{1.964213in}}%
\pgfpathlineto{\pgfqpoint{2.766384in}{2.005589in}}%
\pgfpathlineto{\pgfqpoint{2.733910in}{2.026329in}}%
\pgfpathclose%
\pgfusepath{stroke,fill}%
\end{pgfscope}%
\begin{pgfscope}%
\pgfpathrectangle{\pgfqpoint{0.050000in}{0.050000in}}{\pgfqpoint{4.900000in}{4.900000in}}%
\pgfusepath{clip}%
\pgfsetbuttcap%
\pgfsetroundjoin%
\definecolor{currentfill}{rgb}{0.491349,0.491349,0.491349}%
\pgfsetfillcolor{currentfill}%
\pgfsetfillopacity{0.900000}%
\pgfsetlinewidth{0.501875pt}%
\definecolor{currentstroke}{rgb}{0.200000,0.200000,0.200000}%
\pgfsetstrokecolor{currentstroke}%
\pgfsetdash{}{0pt}%
\pgfpathmoveto{\pgfqpoint{1.746200in}{2.821961in}}%
\pgfpathlineto{\pgfqpoint{1.793782in}{2.776427in}}%
\pgfpathlineto{\pgfqpoint{1.824724in}{2.756563in}}%
\pgfpathlineto{\pgfqpoint{1.777134in}{2.802126in}}%
\pgfpathlineto{\pgfqpoint{1.746200in}{2.821961in}}%
\pgfpathclose%
\pgfusepath{stroke,fill}%
\end{pgfscope}%
\begin{pgfscope}%
\pgfpathrectangle{\pgfqpoint{0.050000in}{0.050000in}}{\pgfqpoint{4.900000in}{4.900000in}}%
\pgfusepath{clip}%
\pgfsetbuttcap%
\pgfsetroundjoin%
\definecolor{currentfill}{rgb}{0.760554,0.760554,0.760554}%
\pgfsetfillcolor{currentfill}%
\pgfsetfillopacity{0.900000}%
\pgfsetlinewidth{0.501875pt}%
\definecolor{currentstroke}{rgb}{0.200000,0.200000,0.200000}%
\pgfsetstrokecolor{currentstroke}%
\pgfsetdash{}{0pt}%
\pgfpathmoveto{\pgfqpoint{2.467457in}{2.239283in}}%
\pgfpathlineto{\pgfqpoint{2.513506in}{2.195276in}}%
\pgfpathlineto{\pgfqpoint{2.545571in}{2.174602in}}%
\pgfpathlineto{\pgfqpoint{2.499518in}{2.218590in}}%
\pgfpathlineto{\pgfqpoint{2.467457in}{2.239283in}}%
\pgfpathclose%
\pgfusepath{stroke,fill}%
\end{pgfscope}%
\begin{pgfscope}%
\pgfpathrectangle{\pgfqpoint{0.050000in}{0.050000in}}{\pgfqpoint{4.900000in}{4.900000in}}%
\pgfusepath{clip}%
\pgfsetbuttcap%
\pgfsetroundjoin%
\definecolor{currentfill}{rgb}{0.898378,0.898378,0.898378}%
\pgfsetfillcolor{currentfill}%
\pgfsetfillopacity{0.900000}%
\pgfsetlinewidth{0.501875pt}%
\definecolor{currentstroke}{rgb}{0.200000,0.200000,0.200000}%
\pgfsetstrokecolor{currentstroke}%
\pgfsetdash{}{0pt}%
\pgfpathmoveto{\pgfqpoint{3.000393in}{1.825310in}}%
\pgfpathlineto{\pgfqpoint{3.045379in}{1.790323in}}%
\pgfpathlineto{\pgfqpoint{3.078288in}{1.770120in}}%
\pgfpathlineto{\pgfqpoint{3.033294in}{1.804883in}}%
\pgfpathlineto{\pgfqpoint{3.000393in}{1.825310in}}%
\pgfpathclose%
\pgfusepath{stroke,fill}%
\end{pgfscope}%
\begin{pgfscope}%
\pgfpathrectangle{\pgfqpoint{0.050000in}{0.050000in}}{\pgfqpoint{4.900000in}{4.900000in}}%
\pgfusepath{clip}%
\pgfsetbuttcap%
\pgfsetroundjoin%
\definecolor{currentfill}{rgb}{0.395663,0.395663,0.395663}%
\pgfsetfillcolor{currentfill}%
\pgfsetfillopacity{0.900000}%
\pgfsetlinewidth{0.501875pt}%
\definecolor{currentstroke}{rgb}{0.200000,0.200000,0.200000}%
\pgfsetstrokecolor{currentstroke}%
\pgfsetdash{}{0pt}%
\pgfpathmoveto{\pgfqpoint{1.479563in}{3.038111in}}%
\pgfpathlineto{\pgfqpoint{1.527720in}{2.992109in}}%
\pgfpathlineto{\pgfqpoint{1.558253in}{2.972577in}}%
\pgfpathlineto{\pgfqpoint{1.510086in}{3.018609in}}%
\pgfpathlineto{\pgfqpoint{1.479563in}{3.038111in}}%
\pgfpathclose%
\pgfusepath{stroke,fill}%
\end{pgfscope}%
\begin{pgfscope}%
\pgfpathrectangle{\pgfqpoint{0.050000in}{0.050000in}}{\pgfqpoint{4.900000in}{4.900000in}}%
\pgfusepath{clip}%
\pgfsetbuttcap%
\pgfsetroundjoin%
\definecolor{currentfill}{rgb}{0.992618,0.992618,0.992618}%
\pgfsetfillcolor{currentfill}%
\pgfsetfillopacity{0.900000}%
\pgfsetlinewidth{0.501875pt}%
\definecolor{currentstroke}{rgb}{0.200000,0.200000,0.200000}%
\pgfsetstrokecolor{currentstroke}%
\pgfsetdash{}{0pt}%
\pgfpathmoveto{\pgfqpoint{4.015266in}{1.384834in}}%
\pgfpathlineto{\pgfqpoint{4.059316in}{1.391201in}}%
\pgfpathlineto{\pgfqpoint{4.093850in}{1.370728in}}%
\pgfpathlineto{\pgfqpoint{4.049776in}{1.364436in}}%
\pgfpathlineto{\pgfqpoint{4.015266in}{1.384834in}}%
\pgfpathclose%
\pgfusepath{stroke,fill}%
\end{pgfscope}%
\begin{pgfscope}%
\pgfpathrectangle{\pgfqpoint{0.050000in}{0.050000in}}{\pgfqpoint{4.900000in}{4.900000in}}%
\pgfusepath{clip}%
\pgfsetbuttcap%
\pgfsetroundjoin%
\definecolor{currentfill}{rgb}{0.662607,0.662607,0.662607}%
\pgfsetfillcolor{currentfill}%
\pgfsetfillopacity{0.900000}%
\pgfsetlinewidth{0.501875pt}%
\definecolor{currentstroke}{rgb}{0.200000,0.200000,0.200000}%
\pgfsetstrokecolor{currentstroke}%
\pgfsetdash{}{0pt}%
\pgfpathmoveto{\pgfqpoint{2.200924in}{2.455095in}}%
\pgfpathlineto{\pgfqpoint{2.247548in}{2.410341in}}%
\pgfpathlineto{\pgfqpoint{2.279204in}{2.389909in}}%
\pgfpathlineto{\pgfqpoint{2.232575in}{2.434686in}}%
\pgfpathlineto{\pgfqpoint{2.200924in}{2.455095in}}%
\pgfpathclose%
\pgfusepath{stroke,fill}%
\end{pgfscope}%
\begin{pgfscope}%
\pgfpathrectangle{\pgfqpoint{0.050000in}{0.050000in}}{\pgfqpoint{4.900000in}{4.900000in}}%
\pgfusepath{clip}%
\pgfsetbuttcap%
\pgfsetroundjoin%
\definecolor{currentfill}{rgb}{0.295271,0.295271,0.295271}%
\pgfsetfillcolor{currentfill}%
\pgfsetfillopacity{0.900000}%
\pgfsetlinewidth{0.501875pt}%
\definecolor{currentstroke}{rgb}{0.200000,0.200000,0.200000}%
\pgfsetstrokecolor{currentstroke}%
\pgfsetdash{}{0pt}%
\pgfpathmoveto{\pgfqpoint{1.212835in}{3.254336in}}%
\pgfpathlineto{\pgfqpoint{1.261569in}{3.207866in}}%
\pgfpathlineto{\pgfqpoint{1.291692in}{3.188666in}}%
\pgfpathlineto{\pgfqpoint{1.242947in}{3.235168in}}%
\pgfpathlineto{\pgfqpoint{1.212835in}{3.254336in}}%
\pgfpathclose%
\pgfusepath{stroke,fill}%
\end{pgfscope}%
\begin{pgfscope}%
\pgfpathrectangle{\pgfqpoint{0.050000in}{0.050000in}}{\pgfqpoint{4.900000in}{4.900000in}}%
\pgfusepath{clip}%
\pgfsetbuttcap%
\pgfsetroundjoin%
\definecolor{currentfill}{rgb}{0.555940,0.555940,0.555940}%
\pgfsetfillcolor{currentfill}%
\pgfsetfillopacity{0.900000}%
\pgfsetlinewidth{0.501875pt}%
\definecolor{currentstroke}{rgb}{0.200000,0.200000,0.200000}%
\pgfsetstrokecolor{currentstroke}%
\pgfsetdash{}{0pt}%
\pgfpathmoveto{\pgfqpoint{1.934300in}{2.671225in}}%
\pgfpathlineto{\pgfqpoint{1.981501in}{2.625993in}}%
\pgfpathlineto{\pgfqpoint{2.012748in}{2.605887in}}%
\pgfpathlineto{\pgfqpoint{1.965540in}{2.651147in}}%
\pgfpathlineto{\pgfqpoint{1.934300in}{2.671225in}}%
\pgfpathclose%
\pgfusepath{stroke,fill}%
\end{pgfscope}%
\begin{pgfscope}%
\pgfpathrectangle{\pgfqpoint{0.050000in}{0.050000in}}{\pgfqpoint{4.900000in}{4.900000in}}%
\pgfusepath{clip}%
\pgfsetbuttcap%
\pgfsetroundjoin%
\definecolor{currentfill}{rgb}{0.966782,0.966782,0.966782}%
\pgfsetfillcolor{currentfill}%
\pgfsetfillopacity{0.900000}%
\pgfsetlinewidth{0.501875pt}%
\definecolor{currentstroke}{rgb}{0.200000,0.200000,0.200000}%
\pgfsetstrokecolor{currentstroke}%
\pgfsetdash{}{0pt}%
\pgfpathmoveto{\pgfqpoint{3.501291in}{1.528695in}}%
\pgfpathlineto{\pgfqpoint{3.545669in}{1.515780in}}%
\pgfpathlineto{\pgfqpoint{3.579401in}{1.496021in}}%
\pgfpathlineto{\pgfqpoint{3.535005in}{1.508866in}}%
\pgfpathlineto{\pgfqpoint{3.501291in}{1.528695in}}%
\pgfpathclose%
\pgfusepath{stroke,fill}%
\end{pgfscope}%
\begin{pgfscope}%
\pgfpathrectangle{\pgfqpoint{0.050000in}{0.050000in}}{\pgfqpoint{4.900000in}{4.900000in}}%
\pgfusepath{clip}%
\pgfsetbuttcap%
\pgfsetroundjoin%
\definecolor{currentfill}{rgb}{0.972318,0.972318,0.972318}%
\pgfsetfillcolor{currentfill}%
\pgfsetfillopacity{0.900000}%
\pgfsetlinewidth{0.501875pt}%
\definecolor{currentstroke}{rgb}{0.200000,0.200000,0.200000}%
\pgfsetstrokecolor{currentstroke}%
\pgfsetdash{}{0pt}%
\pgfpathmoveto{\pgfqpoint{3.579401in}{1.496021in}}%
\pgfpathlineto{\pgfqpoint{3.623722in}{1.486740in}}%
\pgfpathlineto{\pgfqpoint{3.657581in}{1.466931in}}%
\pgfpathlineto{\pgfqpoint{3.613240in}{1.476180in}}%
\pgfpathlineto{\pgfqpoint{3.579401in}{1.496021in}}%
\pgfpathclose%
\pgfusepath{stroke,fill}%
\end{pgfscope}%
\begin{pgfscope}%
\pgfpathrectangle{\pgfqpoint{0.050000in}{0.050000in}}{\pgfqpoint{4.900000in}{4.900000in}}%
\pgfusepath{clip}%
\pgfsetbuttcap%
\pgfsetroundjoin%
\definecolor{currentfill}{rgb}{0.173472,0.173472,0.173472}%
\pgfsetfillcolor{currentfill}%
\pgfsetfillopacity{0.900000}%
\pgfsetlinewidth{0.501875pt}%
\definecolor{currentstroke}{rgb}{0.200000,0.200000,0.200000}%
\pgfsetstrokecolor{currentstroke}%
\pgfsetdash{}{0pt}%
\pgfpathmoveto{\pgfqpoint{0.946016in}{3.470638in}}%
\pgfpathlineto{\pgfqpoint{0.995326in}{3.423698in}}%
\pgfpathlineto{\pgfqpoint{1.025040in}{3.404831in}}%
\pgfpathlineto{\pgfqpoint{0.975717in}{3.451803in}}%
\pgfpathlineto{\pgfqpoint{0.946016in}{3.470638in}}%
\pgfpathclose%
\pgfusepath{stroke,fill}%
\end{pgfscope}%
\begin{pgfscope}%
\pgfpathrectangle{\pgfqpoint{0.050000in}{0.050000in}}{\pgfqpoint{4.900000in}{4.900000in}}%
\pgfusepath{clip}%
\pgfsetbuttcap%
\pgfsetroundjoin%
\definecolor{currentfill}{rgb}{0.959400,0.959400,0.959400}%
\pgfsetfillcolor{currentfill}%
\pgfsetfillopacity{0.900000}%
\pgfsetlinewidth{0.501875pt}%
\definecolor{currentstroke}{rgb}{0.200000,0.200000,0.200000}%
\pgfsetstrokecolor{currentstroke}%
\pgfsetdash{}{0pt}%
\pgfpathmoveto{\pgfqpoint{3.423239in}{1.565168in}}%
\pgfpathlineto{\pgfqpoint{3.467684in}{1.548442in}}%
\pgfpathlineto{\pgfqpoint{3.501291in}{1.528695in}}%
\pgfpathlineto{\pgfqpoint{3.456829in}{1.545312in}}%
\pgfpathlineto{\pgfqpoint{3.423239in}{1.565168in}}%
\pgfpathclose%
\pgfusepath{stroke,fill}%
\end{pgfscope}%
\begin{pgfscope}%
\pgfpathrectangle{\pgfqpoint{0.050000in}{0.050000in}}{\pgfqpoint{4.900000in}{4.900000in}}%
\pgfusepath{clip}%
\pgfsetbuttcap%
\pgfsetroundjoin%
\definecolor{currentfill}{rgb}{0.812226,0.812226,0.812226}%
\pgfsetfillcolor{currentfill}%
\pgfsetfillopacity{0.900000}%
\pgfsetlinewidth{0.501875pt}%
\definecolor{currentstroke}{rgb}{0.200000,0.200000,0.200000}%
\pgfsetstrokecolor{currentstroke}%
\pgfsetdash{}{0pt}%
\pgfpathmoveto{\pgfqpoint{2.655865in}{2.089683in}}%
\pgfpathlineto{\pgfqpoint{2.701537in}{2.047019in}}%
\pgfpathlineto{\pgfqpoint{2.733910in}{2.026329in}}%
\pgfpathlineto{\pgfqpoint{2.688235in}{2.068900in}}%
\pgfpathlineto{\pgfqpoint{2.655865in}{2.089683in}}%
\pgfpathclose%
\pgfusepath{stroke,fill}%
\end{pgfscope}%
\begin{pgfscope}%
\pgfpathrectangle{\pgfqpoint{0.050000in}{0.050000in}}{\pgfqpoint{4.900000in}{4.900000in}}%
\pgfusepath{clip}%
\pgfsetbuttcap%
\pgfsetroundjoin%
\definecolor{currentfill}{rgb}{0.994464,0.994464,0.994464}%
\pgfsetfillcolor{currentfill}%
\pgfsetfillopacity{0.900000}%
\pgfsetlinewidth{0.501875pt}%
\definecolor{currentstroke}{rgb}{0.200000,0.200000,0.200000}%
\pgfsetstrokecolor{currentstroke}%
\pgfsetdash{}{0pt}%
\pgfpathmoveto{\pgfqpoint{4.093850in}{1.370728in}}%
\pgfpathlineto{\pgfqpoint{4.137859in}{1.378589in}}%
\pgfpathlineto{\pgfqpoint{4.172526in}{1.357971in}}%
\pgfpathlineto{\pgfqpoint{4.128494in}{1.350186in}}%
\pgfpathlineto{\pgfqpoint{4.093850in}{1.370728in}}%
\pgfpathclose%
\pgfusepath{stroke,fill}%
\end{pgfscope}%
\begin{pgfscope}%
\pgfpathrectangle{\pgfqpoint{0.050000in}{0.050000in}}{\pgfqpoint{4.900000in}{4.900000in}}%
\pgfusepath{clip}%
\pgfsetbuttcap%
\pgfsetroundjoin%
\definecolor{currentfill}{rgb}{0.878570,0.878570,0.878570}%
\pgfsetfillcolor{currentfill}%
\pgfsetfillopacity{0.900000}%
\pgfsetlinewidth{0.501875pt}%
\definecolor{currentstroke}{rgb}{0.200000,0.200000,0.200000}%
\pgfsetstrokecolor{currentstroke}%
\pgfsetdash{}{0pt}%
\pgfpathmoveto{\pgfqpoint{2.922454in}{1.883294in}}%
\pgfpathlineto{\pgfqpoint{2.967596in}{1.845682in}}%
\pgfpathlineto{\pgfqpoint{3.000393in}{1.825310in}}%
\pgfpathlineto{\pgfqpoint{2.955245in}{1.862716in}}%
\pgfpathlineto{\pgfqpoint{2.922454in}{1.883294in}}%
\pgfpathclose%
\pgfusepath{stroke,fill}%
\end{pgfscope}%
\begin{pgfscope}%
\pgfpathrectangle{\pgfqpoint{0.050000in}{0.050000in}}{\pgfqpoint{4.900000in}{4.900000in}}%
\pgfusepath{clip}%
\pgfsetbuttcap%
\pgfsetroundjoin%
\definecolor{currentfill}{rgb}{0.456901,0.456901,0.456901}%
\pgfsetfillcolor{currentfill}%
\pgfsetfillopacity{0.900000}%
\pgfsetlinewidth{0.501875pt}%
\definecolor{currentstroke}{rgb}{0.200000,0.200000,0.200000}%
\pgfsetstrokecolor{currentstroke}%
\pgfsetdash{}{0pt}%
\pgfpathmoveto{\pgfqpoint{1.667586in}{2.887435in}}%
\pgfpathlineto{\pgfqpoint{1.715363in}{2.841735in}}%
\pgfpathlineto{\pgfqpoint{1.746200in}{2.821961in}}%
\pgfpathlineto{\pgfqpoint{1.698415in}{2.867690in}}%
\pgfpathlineto{\pgfqpoint{1.667586in}{2.887435in}}%
\pgfpathclose%
\pgfusepath{stroke,fill}%
\end{pgfscope}%
\begin{pgfscope}%
\pgfpathrectangle{\pgfqpoint{0.050000in}{0.050000in}}{\pgfqpoint{4.900000in}{4.900000in}}%
\pgfusepath{clip}%
\pgfsetbuttcap%
\pgfsetroundjoin%
\definecolor{currentfill}{rgb}{0.977855,0.977855,0.977855}%
\pgfsetfillcolor{currentfill}%
\pgfsetfillopacity{0.900000}%
\pgfsetlinewidth{0.501875pt}%
\definecolor{currentstroke}{rgb}{0.200000,0.200000,0.200000}%
\pgfsetstrokecolor{currentstroke}%
\pgfsetdash{}{0pt}%
\pgfpathmoveto{\pgfqpoint{3.657581in}{1.466931in}}%
\pgfpathlineto{\pgfqpoint{3.701854in}{1.461037in}}%
\pgfpathlineto{\pgfqpoint{3.735840in}{1.441147in}}%
\pgfpathlineto{\pgfqpoint{3.691547in}{1.447041in}}%
\pgfpathlineto{\pgfqpoint{3.657581in}{1.466931in}}%
\pgfpathclose%
\pgfusepath{stroke,fill}%
\end{pgfscope}%
\begin{pgfscope}%
\pgfpathrectangle{\pgfqpoint{0.050000in}{0.050000in}}{\pgfqpoint{4.900000in}{4.900000in}}%
\pgfusepath{clip}%
\pgfsetbuttcap%
\pgfsetroundjoin%
\definecolor{currentfill}{rgb}{0.952018,0.952018,0.952018}%
\pgfsetfillcolor{currentfill}%
\pgfsetfillopacity{0.900000}%
\pgfsetlinewidth{0.501875pt}%
\definecolor{currentstroke}{rgb}{0.200000,0.200000,0.200000}%
\pgfsetstrokecolor{currentstroke}%
\pgfsetdash{}{0pt}%
\pgfpathmoveto{\pgfqpoint{3.345231in}{1.605566in}}%
\pgfpathlineto{\pgfqpoint{3.389755in}{1.584945in}}%
\pgfpathlineto{\pgfqpoint{3.423239in}{1.565168in}}%
\pgfpathlineto{\pgfqpoint{3.378699in}{1.585642in}}%
\pgfpathlineto{\pgfqpoint{3.345231in}{1.605566in}}%
\pgfpathclose%
\pgfusepath{stroke,fill}%
\end{pgfscope}%
\begin{pgfscope}%
\pgfpathrectangle{\pgfqpoint{0.050000in}{0.050000in}}{\pgfqpoint{4.900000in}{4.900000in}}%
\pgfusepath{clip}%
\pgfsetbuttcap%
\pgfsetroundjoin%
\definecolor{currentfill}{rgb}{0.734579,0.734579,0.734579}%
\pgfsetfillcolor{currentfill}%
\pgfsetfillopacity{0.900000}%
\pgfsetlinewidth{0.501875pt}%
\definecolor{currentstroke}{rgb}{0.200000,0.200000,0.200000}%
\pgfsetstrokecolor{currentstroke}%
\pgfsetdash{}{0pt}%
\pgfpathmoveto{\pgfqpoint{2.389253in}{2.304260in}}%
\pgfpathlineto{\pgfqpoint{2.435495in}{2.259918in}}%
\pgfpathlineto{\pgfqpoint{2.467457in}{2.239283in}}%
\pgfpathlineto{\pgfqpoint{2.421211in}{2.283627in}}%
\pgfpathlineto{\pgfqpoint{2.389253in}{2.304260in}}%
\pgfpathclose%
\pgfusepath{stroke,fill}%
\end{pgfscope}%
\begin{pgfscope}%
\pgfpathrectangle{\pgfqpoint{0.050000in}{0.050000in}}{\pgfqpoint{4.900000in}{4.900000in}}%
\pgfusepath{clip}%
\pgfsetbuttcap%
\pgfsetroundjoin%
\definecolor{currentfill}{rgb}{0.367243,0.367243,0.367243}%
\pgfsetfillcolor{currentfill}%
\pgfsetfillopacity{0.900000}%
\pgfsetlinewidth{0.501875pt}%
\definecolor{currentstroke}{rgb}{0.200000,0.200000,0.200000}%
\pgfsetstrokecolor{currentstroke}%
\pgfsetdash{}{0pt}%
\pgfpathmoveto{\pgfqpoint{1.400782in}{3.103721in}}%
\pgfpathlineto{\pgfqpoint{1.449135in}{3.057552in}}%
\pgfpathlineto{\pgfqpoint{1.479563in}{3.038111in}}%
\pgfpathlineto{\pgfqpoint{1.431199in}{3.084310in}}%
\pgfpathlineto{\pgfqpoint{1.400782in}{3.103721in}}%
\pgfpathclose%
\pgfusepath{stroke,fill}%
\end{pgfscope}%
\begin{pgfscope}%
\pgfpathrectangle{\pgfqpoint{0.050000in}{0.050000in}}{\pgfqpoint{4.900000in}{4.900000in}}%
\pgfusepath{clip}%
\pgfsetbuttcap%
\pgfsetroundjoin%
\definecolor{currentfill}{rgb}{0.981546,0.981546,0.981546}%
\pgfsetfillcolor{currentfill}%
\pgfsetfillopacity{0.900000}%
\pgfsetlinewidth{0.501875pt}%
\definecolor{currentstroke}{rgb}{0.200000,0.200000,0.200000}%
\pgfsetstrokecolor{currentstroke}%
\pgfsetdash{}{0pt}%
\pgfpathmoveto{\pgfqpoint{3.735840in}{1.441147in}}%
\pgfpathlineto{\pgfqpoint{3.780070in}{1.438343in}}%
\pgfpathlineto{\pgfqpoint{3.814186in}{1.418346in}}%
\pgfpathlineto{\pgfqpoint{3.769935in}{1.421177in}}%
\pgfpathlineto{\pgfqpoint{3.735840in}{1.441147in}}%
\pgfpathclose%
\pgfusepath{stroke,fill}%
\end{pgfscope}%
\begin{pgfscope}%
\pgfpathrectangle{\pgfqpoint{0.050000in}{0.050000in}}{\pgfqpoint{4.900000in}{4.900000in}}%
\pgfusepath{clip}%
\pgfsetbuttcap%
\pgfsetroundjoin%
\definecolor{currentfill}{rgb}{0.942791,0.942791,0.942791}%
\pgfsetfillcolor{currentfill}%
\pgfsetfillopacity{0.900000}%
\pgfsetlinewidth{0.501875pt}%
\definecolor{currentstroke}{rgb}{0.200000,0.200000,0.200000}%
\pgfsetstrokecolor{currentstroke}%
\pgfsetdash{}{0pt}%
\pgfpathmoveto{\pgfqpoint{3.267250in}{1.649907in}}%
\pgfpathlineto{\pgfqpoint{3.311869in}{1.625415in}}%
\pgfpathlineto{\pgfqpoint{3.345231in}{1.605566in}}%
\pgfpathlineto{\pgfqpoint{3.300599in}{1.629876in}}%
\pgfpathlineto{\pgfqpoint{3.267250in}{1.649907in}}%
\pgfpathclose%
\pgfusepath{stroke,fill}%
\end{pgfscope}%
\begin{pgfscope}%
\pgfpathrectangle{\pgfqpoint{0.050000in}{0.050000in}}{\pgfqpoint{4.900000in}{4.900000in}}%
\pgfusepath{clip}%
\pgfsetbuttcap%
\pgfsetroundjoin%
\definecolor{currentfill}{rgb}{0.624221,0.624221,0.624221}%
\pgfsetfillcolor{currentfill}%
\pgfsetfillopacity{0.900000}%
\pgfsetlinewidth{0.501875pt}%
\definecolor{currentstroke}{rgb}{0.200000,0.200000,0.200000}%
\pgfsetstrokecolor{currentstroke}%
\pgfsetdash{}{0pt}%
\pgfpathmoveto{\pgfqpoint{2.122553in}{2.520368in}}%
\pgfpathlineto{\pgfqpoint{2.169371in}{2.475442in}}%
\pgfpathlineto{\pgfqpoint{2.200924in}{2.455095in}}%
\pgfpathlineto{\pgfqpoint{2.154100in}{2.500047in}}%
\pgfpathlineto{\pgfqpoint{2.122553in}{2.520368in}}%
\pgfpathclose%
\pgfusepath{stroke,fill}%
\end{pgfscope}%
\begin{pgfscope}%
\pgfpathrectangle{\pgfqpoint{0.050000in}{0.050000in}}{\pgfqpoint{4.900000in}{4.900000in}}%
\pgfusepath{clip}%
\pgfsetbuttcap%
\pgfsetroundjoin%
\definecolor{currentfill}{rgb}{0.996309,0.996309,0.996309}%
\pgfsetfillcolor{currentfill}%
\pgfsetfillopacity{0.900000}%
\pgfsetlinewidth{0.501875pt}%
\definecolor{currentstroke}{rgb}{0.200000,0.200000,0.200000}%
\pgfsetstrokecolor{currentstroke}%
\pgfsetdash{}{0pt}%
\pgfpathmoveto{\pgfqpoint{4.172526in}{1.357971in}}%
\pgfpathlineto{\pgfqpoint{4.216490in}{1.367048in}}%
\pgfpathlineto{\pgfqpoint{4.251292in}{1.346287in}}%
\pgfpathlineto{\pgfqpoint{4.207304in}{1.337284in}}%
\pgfpathlineto{\pgfqpoint{4.172526in}{1.357971in}}%
\pgfpathclose%
\pgfusepath{stroke,fill}%
\end{pgfscope}%
\begin{pgfscope}%
\pgfpathrectangle{\pgfqpoint{0.050000in}{0.050000in}}{\pgfqpoint{4.900000in}{4.900000in}}%
\pgfusepath{clip}%
\pgfsetbuttcap%
\pgfsetroundjoin%
\definecolor{currentfill}{rgb}{0.256517,0.256517,0.256517}%
\pgfsetfillcolor{currentfill}%
\pgfsetfillopacity{0.900000}%
\pgfsetlinewidth{0.501875pt}%
\definecolor{currentstroke}{rgb}{0.200000,0.200000,0.200000}%
\pgfsetstrokecolor{currentstroke}%
\pgfsetdash{}{0pt}%
\pgfpathmoveto{\pgfqpoint{1.133886in}{3.320083in}}%
\pgfpathlineto{\pgfqpoint{1.182817in}{3.273445in}}%
\pgfpathlineto{\pgfqpoint{1.212835in}{3.254336in}}%
\pgfpathlineto{\pgfqpoint{1.163892in}{3.301006in}}%
\pgfpathlineto{\pgfqpoint{1.133886in}{3.320083in}}%
\pgfpathclose%
\pgfusepath{stroke,fill}%
\end{pgfscope}%
\begin{pgfscope}%
\pgfpathrectangle{\pgfqpoint{0.050000in}{0.050000in}}{\pgfqpoint{4.900000in}{4.900000in}}%
\pgfusepath{clip}%
\pgfsetbuttcap%
\pgfsetroundjoin%
\definecolor{currentfill}{rgb}{0.985236,0.985236,0.985236}%
\pgfsetfillcolor{currentfill}%
\pgfsetfillopacity{0.900000}%
\pgfsetlinewidth{0.501875pt}%
\definecolor{currentstroke}{rgb}{0.200000,0.200000,0.200000}%
\pgfsetstrokecolor{currentstroke}%
\pgfsetdash{}{0pt}%
\pgfpathmoveto{\pgfqpoint{3.814186in}{1.418346in}}%
\pgfpathlineto{\pgfqpoint{3.858376in}{1.418310in}}%
\pgfpathlineto{\pgfqpoint{3.892622in}{1.398189in}}%
\pgfpathlineto{\pgfqpoint{3.848410in}{1.398272in}}%
\pgfpathlineto{\pgfqpoint{3.814186in}{1.418346in}}%
\pgfpathclose%
\pgfusepath{stroke,fill}%
\end{pgfscope}%
\begin{pgfscope}%
\pgfpathrectangle{\pgfqpoint{0.050000in}{0.050000in}}{\pgfqpoint{4.900000in}{4.900000in}}%
\pgfusepath{clip}%
\pgfsetbuttcap%
\pgfsetroundjoin%
\definecolor{currentfill}{rgb}{0.521492,0.521492,0.521492}%
\pgfsetfillcolor{currentfill}%
\pgfsetfillopacity{0.900000}%
\pgfsetlinewidth{0.501875pt}%
\definecolor{currentstroke}{rgb}{0.200000,0.200000,0.200000}%
\pgfsetstrokecolor{currentstroke}%
\pgfsetdash{}{0pt}%
\pgfpathmoveto{\pgfqpoint{1.855763in}{2.736638in}}%
\pgfpathlineto{\pgfqpoint{1.903157in}{2.691240in}}%
\pgfpathlineto{\pgfqpoint{1.934300in}{2.671225in}}%
\pgfpathlineto{\pgfqpoint{1.886898in}{2.716650in}}%
\pgfpathlineto{\pgfqpoint{1.855763in}{2.736638in}}%
\pgfpathclose%
\pgfusepath{stroke,fill}%
\end{pgfscope}%
\begin{pgfscope}%
\pgfpathrectangle{\pgfqpoint{0.050000in}{0.050000in}}{\pgfqpoint{4.900000in}{4.900000in}}%
\pgfusepath{clip}%
\pgfsetbuttcap%
\pgfsetroundjoin%
\definecolor{currentfill}{rgb}{0.929504,0.929504,0.929504}%
\pgfsetfillcolor{currentfill}%
\pgfsetfillopacity{0.900000}%
\pgfsetlinewidth{0.501875pt}%
\definecolor{currentstroke}{rgb}{0.200000,0.200000,0.200000}%
\pgfsetstrokecolor{currentstroke}%
\pgfsetdash{}{0pt}%
\pgfpathmoveto{\pgfqpoint{3.189280in}{1.698082in}}%
\pgfpathlineto{\pgfqpoint{3.234007in}{1.669867in}}%
\pgfpathlineto{\pgfqpoint{3.267250in}{1.649907in}}%
\pgfpathlineto{\pgfqpoint{3.222511in}{1.677914in}}%
\pgfpathlineto{\pgfqpoint{3.189280in}{1.698082in}}%
\pgfpathclose%
\pgfusepath{stroke,fill}%
\end{pgfscope}%
\begin{pgfscope}%
\pgfpathrectangle{\pgfqpoint{0.050000in}{0.050000in}}{\pgfqpoint{4.900000in}{4.900000in}}%
\pgfusepath{clip}%
\pgfsetbuttcap%
\pgfsetroundjoin%
\definecolor{currentfill}{rgb}{0.858762,0.858762,0.858762}%
\pgfsetfillcolor{currentfill}%
\pgfsetfillopacity{0.900000}%
\pgfsetlinewidth{0.501875pt}%
\definecolor{currentstroke}{rgb}{0.200000,0.200000,0.200000}%
\pgfsetstrokecolor{currentstroke}%
\pgfsetdash{}{0pt}%
\pgfpathmoveto{\pgfqpoint{2.844455in}{1.943560in}}%
\pgfpathlineto{\pgfqpoint{2.889766in}{1.903821in}}%
\pgfpathlineto{\pgfqpoint{2.922454in}{1.883294in}}%
\pgfpathlineto{\pgfqpoint{2.877137in}{1.922857in}}%
\pgfpathlineto{\pgfqpoint{2.844455in}{1.943560in}}%
\pgfpathclose%
\pgfusepath{stroke,fill}%
\end{pgfscope}%
\begin{pgfscope}%
\pgfpathrectangle{\pgfqpoint{0.050000in}{0.050000in}}{\pgfqpoint{4.900000in}{4.900000in}}%
\pgfusepath{clip}%
\pgfsetbuttcap%
\pgfsetroundjoin%
\definecolor{currentfill}{rgb}{0.788112,0.788112,0.788112}%
\pgfsetfillcolor{currentfill}%
\pgfsetfillopacity{0.900000}%
\pgfsetlinewidth{0.501875pt}%
\definecolor{currentstroke}{rgb}{0.200000,0.200000,0.200000}%
\pgfsetstrokecolor{currentstroke}%
\pgfsetdash{}{0pt}%
\pgfpathmoveto{\pgfqpoint{2.577736in}{2.153872in}}%
\pgfpathlineto{\pgfqpoint{2.623596in}{2.110414in}}%
\pgfpathlineto{\pgfqpoint{2.655865in}{2.089683in}}%
\pgfpathlineto{\pgfqpoint{2.610002in}{2.133088in}}%
\pgfpathlineto{\pgfqpoint{2.577736in}{2.153872in}}%
\pgfpathclose%
\pgfusepath{stroke,fill}%
\end{pgfscope}%
\begin{pgfscope}%
\pgfpathrectangle{\pgfqpoint{0.050000in}{0.050000in}}{\pgfqpoint{4.900000in}{4.900000in}}%
\pgfusepath{clip}%
\pgfsetbuttcap%
\pgfsetroundjoin%
\definecolor{currentfill}{rgb}{0.132011,0.132011,0.132011}%
\pgfsetfillcolor{currentfill}%
\pgfsetfillopacity{0.900000}%
\pgfsetlinewidth{0.501875pt}%
\definecolor{currentstroke}{rgb}{0.200000,0.200000,0.200000}%
\pgfsetstrokecolor{currentstroke}%
\pgfsetdash{}{0pt}%
\pgfpathmoveto{\pgfqpoint{0.866901in}{3.536520in}}%
\pgfpathlineto{\pgfqpoint{0.916408in}{3.489414in}}%
\pgfpathlineto{\pgfqpoint{0.946016in}{3.470638in}}%
\pgfpathlineto{\pgfqpoint{0.896495in}{3.517778in}}%
\pgfpathlineto{\pgfqpoint{0.866901in}{3.536520in}}%
\pgfpathclose%
\pgfusepath{stroke,fill}%
\end{pgfscope}%
\begin{pgfscope}%
\pgfpathrectangle{\pgfqpoint{0.050000in}{0.050000in}}{\pgfqpoint{4.900000in}{4.900000in}}%
\pgfusepath{clip}%
\pgfsetbuttcap%
\pgfsetroundjoin%
\definecolor{currentfill}{rgb}{0.428143,0.428143,0.428143}%
\pgfsetfillcolor{currentfill}%
\pgfsetfillopacity{0.900000}%
\pgfsetlinewidth{0.501875pt}%
\definecolor{currentstroke}{rgb}{0.200000,0.200000,0.200000}%
\pgfsetstrokecolor{currentstroke}%
\pgfsetdash{}{0pt}%
\pgfpathmoveto{\pgfqpoint{1.588881in}{2.952984in}}%
\pgfpathlineto{\pgfqpoint{1.636853in}{2.907118in}}%
\pgfpathlineto{\pgfqpoint{1.667586in}{2.887435in}}%
\pgfpathlineto{\pgfqpoint{1.619605in}{2.933330in}}%
\pgfpathlineto{\pgfqpoint{1.588881in}{2.952984in}}%
\pgfpathclose%
\pgfusepath{stroke,fill}%
\end{pgfscope}%
\begin{pgfscope}%
\pgfpathrectangle{\pgfqpoint{0.050000in}{0.050000in}}{\pgfqpoint{4.900000in}{4.900000in}}%
\pgfusepath{clip}%
\pgfsetbuttcap%
\pgfsetroundjoin%
\definecolor{currentfill}{rgb}{0.988927,0.988927,0.988927}%
\pgfsetfillcolor{currentfill}%
\pgfsetfillopacity{0.900000}%
\pgfsetlinewidth{0.501875pt}%
\definecolor{currentstroke}{rgb}{0.200000,0.200000,0.200000}%
\pgfsetstrokecolor{currentstroke}%
\pgfsetdash{}{0pt}%
\pgfpathmoveto{\pgfqpoint{3.892622in}{1.398189in}}%
\pgfpathlineto{\pgfqpoint{3.936774in}{1.400587in}}%
\pgfpathlineto{\pgfqpoint{3.971152in}{1.380331in}}%
\pgfpathlineto{\pgfqpoint{3.926977in}{1.377993in}}%
\pgfpathlineto{\pgfqpoint{3.892622in}{1.398189in}}%
\pgfpathclose%
\pgfusepath{stroke,fill}%
\end{pgfscope}%
\begin{pgfscope}%
\pgfpathrectangle{\pgfqpoint{0.050000in}{0.050000in}}{\pgfqpoint{4.900000in}{4.900000in}}%
\pgfusepath{clip}%
\pgfsetbuttcap%
\pgfsetroundjoin%
\definecolor{currentfill}{rgb}{0.696194,0.696194,0.696194}%
\pgfsetfillcolor{currentfill}%
\pgfsetfillopacity{0.900000}%
\pgfsetlinewidth{0.501875pt}%
\definecolor{currentstroke}{rgb}{0.200000,0.200000,0.200000}%
\pgfsetstrokecolor{currentstroke}%
\pgfsetdash{}{0pt}%
\pgfpathmoveto{\pgfqpoint{2.310959in}{2.369414in}}%
\pgfpathlineto{\pgfqpoint{2.357395in}{2.324832in}}%
\pgfpathlineto{\pgfqpoint{2.389253in}{2.304260in}}%
\pgfpathlineto{\pgfqpoint{2.342813in}{2.348858in}}%
\pgfpathlineto{\pgfqpoint{2.310959in}{2.369414in}}%
\pgfpathclose%
\pgfusepath{stroke,fill}%
\end{pgfscope}%
\begin{pgfscope}%
\pgfpathrectangle{\pgfqpoint{0.050000in}{0.050000in}}{\pgfqpoint{4.900000in}{4.900000in}}%
\pgfusepath{clip}%
\pgfsetbuttcap%
\pgfsetroundjoin%
\definecolor{currentfill}{rgb}{0.334764,0.334764,0.334764}%
\pgfsetfillcolor{currentfill}%
\pgfsetfillopacity{0.900000}%
\pgfsetlinewidth{0.501875pt}%
\definecolor{currentstroke}{rgb}{0.200000,0.200000,0.200000}%
\pgfsetstrokecolor{currentstroke}%
\pgfsetdash{}{0pt}%
\pgfpathmoveto{\pgfqpoint{1.321910in}{3.169406in}}%
\pgfpathlineto{\pgfqpoint{1.370459in}{3.123071in}}%
\pgfpathlineto{\pgfqpoint{1.400782in}{3.103721in}}%
\pgfpathlineto{\pgfqpoint{1.352222in}{3.150087in}}%
\pgfpathlineto{\pgfqpoint{1.321910in}{3.169406in}}%
\pgfpathclose%
\pgfusepath{stroke,fill}%
\end{pgfscope}%
\begin{pgfscope}%
\pgfpathrectangle{\pgfqpoint{0.050000in}{0.050000in}}{\pgfqpoint{4.900000in}{4.900000in}}%
\pgfusepath{clip}%
\pgfsetbuttcap%
\pgfsetroundjoin%
\definecolor{currentfill}{rgb}{0.915356,0.915356,0.915356}%
\pgfsetfillcolor{currentfill}%
\pgfsetfillopacity{0.900000}%
\pgfsetlinewidth{0.501875pt}%
\definecolor{currentstroke}{rgb}{0.200000,0.200000,0.200000}%
\pgfsetstrokecolor{currentstroke}%
\pgfsetdash{}{0pt}%
\pgfpathmoveto{\pgfqpoint{3.111300in}{1.749857in}}%
\pgfpathlineto{\pgfqpoint{3.156153in}{1.718184in}}%
\pgfpathlineto{\pgfqpoint{3.189280in}{1.698082in}}%
\pgfpathlineto{\pgfqpoint{3.144417in}{1.729535in}}%
\pgfpathlineto{\pgfqpoint{3.111300in}{1.749857in}}%
\pgfpathclose%
\pgfusepath{stroke,fill}%
\end{pgfscope}%
\begin{pgfscope}%
\pgfpathrectangle{\pgfqpoint{0.050000in}{0.050000in}}{\pgfqpoint{4.900000in}{4.900000in}}%
\pgfusepath{clip}%
\pgfsetbuttcap%
\pgfsetroundjoin%
\definecolor{currentfill}{rgb}{0.590634,0.590634,0.590634}%
\pgfsetfillcolor{currentfill}%
\pgfsetfillopacity{0.900000}%
\pgfsetlinewidth{0.501875pt}%
\definecolor{currentstroke}{rgb}{0.200000,0.200000,0.200000}%
\pgfsetstrokecolor{currentstroke}%
\pgfsetdash{}{0pt}%
\pgfpathmoveto{\pgfqpoint{2.044092in}{2.585719in}}%
\pgfpathlineto{\pgfqpoint{2.091105in}{2.540626in}}%
\pgfpathlineto{\pgfqpoint{2.122553in}{2.520368in}}%
\pgfpathlineto{\pgfqpoint{2.075534in}{2.565488in}}%
\pgfpathlineto{\pgfqpoint{2.044092in}{2.585719in}}%
\pgfpathclose%
\pgfusepath{stroke,fill}%
\end{pgfscope}%
\begin{pgfscope}%
\pgfpathrectangle{\pgfqpoint{0.050000in}{0.050000in}}{\pgfqpoint{4.900000in}{4.900000in}}%
\pgfusepath{clip}%
\pgfsetbuttcap%
\pgfsetroundjoin%
\definecolor{currentfill}{rgb}{0.990773,0.990773,0.990773}%
\pgfsetfillcolor{currentfill}%
\pgfsetfillopacity{0.900000}%
\pgfsetlinewidth{0.501875pt}%
\definecolor{currentstroke}{rgb}{0.200000,0.200000,0.200000}%
\pgfsetstrokecolor{currentstroke}%
\pgfsetdash{}{0pt}%
\pgfpathmoveto{\pgfqpoint{3.971152in}{1.380331in}}%
\pgfpathlineto{\pgfqpoint{4.015266in}{1.384834in}}%
\pgfpathlineto{\pgfqpoint{4.049776in}{1.364436in}}%
\pgfpathlineto{\pgfqpoint{4.005639in}{1.360001in}}%
\pgfpathlineto{\pgfqpoint{3.971152in}{1.380331in}}%
\pgfpathclose%
\pgfusepath{stroke,fill}%
\end{pgfscope}%
\begin{pgfscope}%
\pgfpathrectangle{\pgfqpoint{0.050000in}{0.050000in}}{\pgfqpoint{4.900000in}{4.900000in}}%
\pgfusepath{clip}%
\pgfsetbuttcap%
\pgfsetroundjoin%
\definecolor{currentfill}{rgb}{0.212226,0.212226,0.212226}%
\pgfsetfillcolor{currentfill}%
\pgfsetfillopacity{0.900000}%
\pgfsetlinewidth{0.501875pt}%
\definecolor{currentstroke}{rgb}{0.200000,0.200000,0.200000}%
\pgfsetstrokecolor{currentstroke}%
\pgfsetdash{}{0pt}%
\pgfpathmoveto{\pgfqpoint{1.054847in}{3.385905in}}%
\pgfpathlineto{\pgfqpoint{1.103974in}{3.339100in}}%
\pgfpathlineto{\pgfqpoint{1.133886in}{3.320083in}}%
\pgfpathlineto{\pgfqpoint{1.084747in}{3.366920in}}%
\pgfpathlineto{\pgfqpoint{1.054847in}{3.385905in}}%
\pgfpathclose%
\pgfusepath{stroke,fill}%
\end{pgfscope}%
\begin{pgfscope}%
\pgfpathrectangle{\pgfqpoint{0.050000in}{0.050000in}}{\pgfqpoint{4.900000in}{4.900000in}}%
\pgfusepath{clip}%
\pgfsetbuttcap%
\pgfsetroundjoin%
\definecolor{currentfill}{rgb}{0.836340,0.836340,0.836340}%
\pgfsetfillcolor{currentfill}%
\pgfsetfillopacity{0.900000}%
\pgfsetlinewidth{0.501875pt}%
\definecolor{currentstroke}{rgb}{0.200000,0.200000,0.200000}%
\pgfsetstrokecolor{currentstroke}%
\pgfsetdash{}{0pt}%
\pgfpathmoveto{\pgfqpoint{2.766384in}{2.005589in}}%
\pgfpathlineto{\pgfqpoint{2.811874in}{1.964213in}}%
\pgfpathlineto{\pgfqpoint{2.844455in}{1.943560in}}%
\pgfpathlineto{\pgfqpoint{2.798961in}{1.984799in}}%
\pgfpathlineto{\pgfqpoint{2.766384in}{2.005589in}}%
\pgfpathclose%
\pgfusepath{stroke,fill}%
\end{pgfscope}%
\begin{pgfscope}%
\pgfpathrectangle{\pgfqpoint{0.050000in}{0.050000in}}{\pgfqpoint{4.900000in}{4.900000in}}%
\pgfusepath{clip}%
\pgfsetbuttcap%
\pgfsetroundjoin%
\definecolor{currentfill}{rgb}{0.491349,0.491349,0.491349}%
\pgfsetfillcolor{currentfill}%
\pgfsetfillopacity{0.900000}%
\pgfsetlinewidth{0.501875pt}%
\definecolor{currentstroke}{rgb}{0.200000,0.200000,0.200000}%
\pgfsetstrokecolor{currentstroke}%
\pgfsetdash{}{0pt}%
\pgfpathmoveto{\pgfqpoint{1.777134in}{2.802126in}}%
\pgfpathlineto{\pgfqpoint{1.824724in}{2.756563in}}%
\pgfpathlineto{\pgfqpoint{1.855763in}{2.736638in}}%
\pgfpathlineto{\pgfqpoint{1.808165in}{2.782229in}}%
\pgfpathlineto{\pgfqpoint{1.777134in}{2.802126in}}%
\pgfpathclose%
\pgfusepath{stroke,fill}%
\end{pgfscope}%
\begin{pgfscope}%
\pgfpathrectangle{\pgfqpoint{0.050000in}{0.050000in}}{\pgfqpoint{4.900000in}{4.900000in}}%
\pgfusepath{clip}%
\pgfsetbuttcap%
\pgfsetroundjoin%
\definecolor{currentfill}{rgb}{0.760554,0.760554,0.760554}%
\pgfsetfillcolor{currentfill}%
\pgfsetfillopacity{0.900000}%
\pgfsetlinewidth{0.501875pt}%
\definecolor{currentstroke}{rgb}{0.200000,0.200000,0.200000}%
\pgfsetstrokecolor{currentstroke}%
\pgfsetdash{}{0pt}%
\pgfpathmoveto{\pgfqpoint{2.499518in}{2.218590in}}%
\pgfpathlineto{\pgfqpoint{2.545571in}{2.174602in}}%
\pgfpathlineto{\pgfqpoint{2.577736in}{2.153872in}}%
\pgfpathlineto{\pgfqpoint{2.531680in}{2.197838in}}%
\pgfpathlineto{\pgfqpoint{2.499518in}{2.218590in}}%
\pgfpathclose%
\pgfusepath{stroke,fill}%
\end{pgfscope}%
\begin{pgfscope}%
\pgfpathrectangle{\pgfqpoint{0.050000in}{0.050000in}}{\pgfqpoint{4.900000in}{4.900000in}}%
\pgfusepath{clip}%
\pgfsetbuttcap%
\pgfsetroundjoin%
\definecolor{currentfill}{rgb}{0.100146,0.100146,0.100146}%
\pgfsetfillcolor{currentfill}%
\pgfsetfillopacity{0.900000}%
\pgfsetlinewidth{0.501875pt}%
\definecolor{currentstroke}{rgb}{0.200000,0.200000,0.200000}%
\pgfsetstrokecolor{currentstroke}%
\pgfsetdash{}{0pt}%
\pgfpathmoveto{\pgfqpoint{0.787694in}{3.602479in}}%
\pgfpathlineto{\pgfqpoint{0.837399in}{3.555205in}}%
\pgfpathlineto{\pgfqpoint{0.866901in}{3.536520in}}%
\pgfpathlineto{\pgfqpoint{0.817181in}{3.583829in}}%
\pgfpathlineto{\pgfqpoint{0.787694in}{3.602479in}}%
\pgfpathclose%
\pgfusepath{stroke,fill}%
\end{pgfscope}%
\begin{pgfscope}%
\pgfpathrectangle{\pgfqpoint{0.050000in}{0.050000in}}{\pgfqpoint{4.900000in}{4.900000in}}%
\pgfusepath{clip}%
\pgfsetbuttcap%
\pgfsetroundjoin%
\definecolor{currentfill}{rgb}{0.898378,0.898378,0.898378}%
\pgfsetfillcolor{currentfill}%
\pgfsetfillopacity{0.900000}%
\pgfsetlinewidth{0.501875pt}%
\definecolor{currentstroke}{rgb}{0.200000,0.200000,0.200000}%
\pgfsetstrokecolor{currentstroke}%
\pgfsetdash{}{0pt}%
\pgfpathmoveto{\pgfqpoint{3.033294in}{1.804883in}}%
\pgfpathlineto{\pgfqpoint{3.078288in}{1.770120in}}%
\pgfpathlineto{\pgfqpoint{3.111300in}{1.749857in}}%
\pgfpathlineto{\pgfqpoint{3.066298in}{1.784402in}}%
\pgfpathlineto{\pgfqpoint{3.033294in}{1.804883in}}%
\pgfpathclose%
\pgfusepath{stroke,fill}%
\end{pgfscope}%
\begin{pgfscope}%
\pgfpathrectangle{\pgfqpoint{0.050000in}{0.050000in}}{\pgfqpoint{4.900000in}{4.900000in}}%
\pgfusepath{clip}%
\pgfsetbuttcap%
\pgfsetroundjoin%
\definecolor{currentfill}{rgb}{0.395663,0.395663,0.395663}%
\pgfsetfillcolor{currentfill}%
\pgfsetfillopacity{0.900000}%
\pgfsetlinewidth{0.501875pt}%
\definecolor{currentstroke}{rgb}{0.200000,0.200000,0.200000}%
\pgfsetstrokecolor{currentstroke}%
\pgfsetdash{}{0pt}%
\pgfpathmoveto{\pgfqpoint{1.510086in}{3.018609in}}%
\pgfpathlineto{\pgfqpoint{1.558253in}{2.972577in}}%
\pgfpathlineto{\pgfqpoint{1.588881in}{2.952984in}}%
\pgfpathlineto{\pgfqpoint{1.540704in}{2.999046in}}%
\pgfpathlineto{\pgfqpoint{1.510086in}{3.018609in}}%
\pgfpathclose%
\pgfusepath{stroke,fill}%
\end{pgfscope}%
\begin{pgfscope}%
\pgfpathrectangle{\pgfqpoint{0.050000in}{0.050000in}}{\pgfqpoint{4.900000in}{4.900000in}}%
\pgfusepath{clip}%
\pgfsetbuttcap%
\pgfsetroundjoin%
\definecolor{currentfill}{rgb}{0.992618,0.992618,0.992618}%
\pgfsetfillcolor{currentfill}%
\pgfsetfillopacity{0.900000}%
\pgfsetlinewidth{0.501875pt}%
\definecolor{currentstroke}{rgb}{0.200000,0.200000,0.200000}%
\pgfsetstrokecolor{currentstroke}%
\pgfsetdash{}{0pt}%
\pgfpathmoveto{\pgfqpoint{4.049776in}{1.364436in}}%
\pgfpathlineto{\pgfqpoint{4.093850in}{1.370728in}}%
\pgfpathlineto{\pgfqpoint{4.128494in}{1.350186in}}%
\pgfpathlineto{\pgfqpoint{4.084396in}{1.343966in}}%
\pgfpathlineto{\pgfqpoint{4.049776in}{1.364436in}}%
\pgfpathclose%
\pgfusepath{stroke,fill}%
\end{pgfscope}%
\begin{pgfscope}%
\pgfpathrectangle{\pgfqpoint{0.050000in}{0.050000in}}{\pgfqpoint{4.900000in}{4.900000in}}%
\pgfusepath{clip}%
\pgfsetbuttcap%
\pgfsetroundjoin%
\definecolor{currentfill}{rgb}{0.662607,0.662607,0.662607}%
\pgfsetfillcolor{currentfill}%
\pgfsetfillopacity{0.900000}%
\pgfsetlinewidth{0.501875pt}%
\definecolor{currentstroke}{rgb}{0.200000,0.200000,0.200000}%
\pgfsetstrokecolor{currentstroke}%
\pgfsetdash{}{0pt}%
\pgfpathmoveto{\pgfqpoint{2.232575in}{2.434686in}}%
\pgfpathlineto{\pgfqpoint{2.279204in}{2.389909in}}%
\pgfpathlineto{\pgfqpoint{2.310959in}{2.369414in}}%
\pgfpathlineto{\pgfqpoint{2.264324in}{2.414213in}}%
\pgfpathlineto{\pgfqpoint{2.232575in}{2.434686in}}%
\pgfpathclose%
\pgfusepath{stroke,fill}%
\end{pgfscope}%
\begin{pgfscope}%
\pgfpathrectangle{\pgfqpoint{0.050000in}{0.050000in}}{\pgfqpoint{4.900000in}{4.900000in}}%
\pgfusepath{clip}%
\pgfsetbuttcap%
\pgfsetroundjoin%
\definecolor{currentfill}{rgb}{0.300807,0.300807,0.300807}%
\pgfsetfillcolor{currentfill}%
\pgfsetfillopacity{0.900000}%
\pgfsetlinewidth{0.501875pt}%
\definecolor{currentstroke}{rgb}{0.200000,0.200000,0.200000}%
\pgfsetstrokecolor{currentstroke}%
\pgfsetdash{}{0pt}%
\pgfpathmoveto{\pgfqpoint{1.242947in}{3.235168in}}%
\pgfpathlineto{\pgfqpoint{1.291692in}{3.188666in}}%
\pgfpathlineto{\pgfqpoint{1.321910in}{3.169406in}}%
\pgfpathlineto{\pgfqpoint{1.273153in}{3.215940in}}%
\pgfpathlineto{\pgfqpoint{1.242947in}{3.235168in}}%
\pgfpathclose%
\pgfusepath{stroke,fill}%
\end{pgfscope}%
\begin{pgfscope}%
\pgfpathrectangle{\pgfqpoint{0.050000in}{0.050000in}}{\pgfqpoint{4.900000in}{4.900000in}}%
\pgfusepath{clip}%
\pgfsetbuttcap%
\pgfsetroundjoin%
\definecolor{currentfill}{rgb}{0.555940,0.555940,0.555940}%
\pgfsetfillcolor{currentfill}%
\pgfsetfillopacity{0.900000}%
\pgfsetlinewidth{0.501875pt}%
\definecolor{currentstroke}{rgb}{0.200000,0.200000,0.200000}%
\pgfsetstrokecolor{currentstroke}%
\pgfsetdash{}{0pt}%
\pgfpathmoveto{\pgfqpoint{1.965540in}{2.651147in}}%
\pgfpathlineto{\pgfqpoint{2.012748in}{2.605887in}}%
\pgfpathlineto{\pgfqpoint{2.044092in}{2.585719in}}%
\pgfpathlineto{\pgfqpoint{1.996878in}{2.631006in}}%
\pgfpathlineto{\pgfqpoint{1.965540in}{2.651147in}}%
\pgfpathclose%
\pgfusepath{stroke,fill}%
\end{pgfscope}%
\begin{pgfscope}%
\pgfpathrectangle{\pgfqpoint{0.050000in}{0.050000in}}{\pgfqpoint{4.900000in}{4.900000in}}%
\pgfusepath{clip}%
\pgfsetbuttcap%
\pgfsetroundjoin%
\definecolor{currentfill}{rgb}{0.173472,0.173472,0.173472}%
\pgfsetfillcolor{currentfill}%
\pgfsetfillopacity{0.900000}%
\pgfsetlinewidth{0.501875pt}%
\definecolor{currentstroke}{rgb}{0.200000,0.200000,0.200000}%
\pgfsetstrokecolor{currentstroke}%
\pgfsetdash{}{0pt}%
\pgfpathmoveto{\pgfqpoint{0.975717in}{3.451803in}}%
\pgfpathlineto{\pgfqpoint{1.025040in}{3.404831in}}%
\pgfpathlineto{\pgfqpoint{1.054847in}{3.385905in}}%
\pgfpathlineto{\pgfqpoint{1.005510in}{3.432910in}}%
\pgfpathlineto{\pgfqpoint{0.975717in}{3.451803in}}%
\pgfpathclose%
\pgfusepath{stroke,fill}%
\end{pgfscope}%
\begin{pgfscope}%
\pgfpathrectangle{\pgfqpoint{0.050000in}{0.050000in}}{\pgfqpoint{4.900000in}{4.900000in}}%
\pgfusepath{clip}%
\pgfsetbuttcap%
\pgfsetroundjoin%
\definecolor{currentfill}{rgb}{0.966782,0.966782,0.966782}%
\pgfsetfillcolor{currentfill}%
\pgfsetfillopacity{0.900000}%
\pgfsetlinewidth{0.501875pt}%
\definecolor{currentstroke}{rgb}{0.200000,0.200000,0.200000}%
\pgfsetstrokecolor{currentstroke}%
\pgfsetdash{}{0pt}%
\pgfpathmoveto{\pgfqpoint{3.535005in}{1.508866in}}%
\pgfpathlineto{\pgfqpoint{3.579401in}{1.496021in}}%
\pgfpathlineto{\pgfqpoint{3.613240in}{1.476180in}}%
\pgfpathlineto{\pgfqpoint{3.568826in}{1.488956in}}%
\pgfpathlineto{\pgfqpoint{3.535005in}{1.508866in}}%
\pgfpathclose%
\pgfusepath{stroke,fill}%
\end{pgfscope}%
\begin{pgfscope}%
\pgfpathrectangle{\pgfqpoint{0.050000in}{0.050000in}}{\pgfqpoint{4.900000in}{4.900000in}}%
\pgfusepath{clip}%
\pgfsetbuttcap%
\pgfsetroundjoin%
\definecolor{currentfill}{rgb}{0.812226,0.812226,0.812226}%
\pgfsetfillcolor{currentfill}%
\pgfsetfillopacity{0.900000}%
\pgfsetlinewidth{0.501875pt}%
\definecolor{currentstroke}{rgb}{0.200000,0.200000,0.200000}%
\pgfsetstrokecolor{currentstroke}%
\pgfsetdash{}{0pt}%
\pgfpathmoveto{\pgfqpoint{2.688235in}{2.068900in}}%
\pgfpathlineto{\pgfqpoint{2.733910in}{2.026329in}}%
\pgfpathlineto{\pgfqpoint{2.766384in}{2.005589in}}%
\pgfpathlineto{\pgfqpoint{2.720707in}{2.048065in}}%
\pgfpathlineto{\pgfqpoint{2.688235in}{2.068900in}}%
\pgfpathclose%
\pgfusepath{stroke,fill}%
\end{pgfscope}%
\begin{pgfscope}%
\pgfpathrectangle{\pgfqpoint{0.050000in}{0.050000in}}{\pgfqpoint{4.900000in}{4.900000in}}%
\pgfusepath{clip}%
\pgfsetbuttcap%
\pgfsetroundjoin%
\definecolor{currentfill}{rgb}{0.972318,0.972318,0.972318}%
\pgfsetfillcolor{currentfill}%
\pgfsetfillopacity{0.900000}%
\pgfsetlinewidth{0.501875pt}%
\definecolor{currentstroke}{rgb}{0.200000,0.200000,0.200000}%
\pgfsetstrokecolor{currentstroke}%
\pgfsetdash{}{0pt}%
\pgfpathmoveto{\pgfqpoint{3.613240in}{1.476180in}}%
\pgfpathlineto{\pgfqpoint{3.657581in}{1.466931in}}%
\pgfpathlineto{\pgfqpoint{3.691547in}{1.447041in}}%
\pgfpathlineto{\pgfqpoint{3.647187in}{1.456256in}}%
\pgfpathlineto{\pgfqpoint{3.613240in}{1.476180in}}%
\pgfpathclose%
\pgfusepath{stroke,fill}%
\end{pgfscope}%
\begin{pgfscope}%
\pgfpathrectangle{\pgfqpoint{0.050000in}{0.050000in}}{\pgfqpoint{4.900000in}{4.900000in}}%
\pgfusepath{clip}%
\pgfsetbuttcap%
\pgfsetroundjoin%
\definecolor{currentfill}{rgb}{0.959400,0.959400,0.959400}%
\pgfsetfillcolor{currentfill}%
\pgfsetfillopacity{0.900000}%
\pgfsetlinewidth{0.501875pt}%
\definecolor{currentstroke}{rgb}{0.200000,0.200000,0.200000}%
\pgfsetstrokecolor{currentstroke}%
\pgfsetdash{}{0pt}%
\pgfpathmoveto{\pgfqpoint{3.456829in}{1.545312in}}%
\pgfpathlineto{\pgfqpoint{3.501291in}{1.528695in}}%
\pgfpathlineto{\pgfqpoint{3.535005in}{1.508866in}}%
\pgfpathlineto{\pgfqpoint{3.490527in}{1.525376in}}%
\pgfpathlineto{\pgfqpoint{3.456829in}{1.545312in}}%
\pgfpathclose%
\pgfusepath{stroke,fill}%
\end{pgfscope}%
\begin{pgfscope}%
\pgfpathrectangle{\pgfqpoint{0.050000in}{0.050000in}}{\pgfqpoint{4.900000in}{4.900000in}}%
\pgfusepath{clip}%
\pgfsetbuttcap%
\pgfsetroundjoin%
\definecolor{currentfill}{rgb}{0.994464,0.994464,0.994464}%
\pgfsetfillcolor{currentfill}%
\pgfsetfillopacity{0.900000}%
\pgfsetlinewidth{0.501875pt}%
\definecolor{currentstroke}{rgb}{0.200000,0.200000,0.200000}%
\pgfsetstrokecolor{currentstroke}%
\pgfsetdash{}{0pt}%
\pgfpathmoveto{\pgfqpoint{4.128494in}{1.350186in}}%
\pgfpathlineto{\pgfqpoint{4.172526in}{1.357971in}}%
\pgfpathlineto{\pgfqpoint{4.207304in}{1.337284in}}%
\pgfpathlineto{\pgfqpoint{4.163247in}{1.329573in}}%
\pgfpathlineto{\pgfqpoint{4.128494in}{1.350186in}}%
\pgfpathclose%
\pgfusepath{stroke,fill}%
\end{pgfscope}%
\begin{pgfscope}%
\pgfpathrectangle{\pgfqpoint{0.050000in}{0.050000in}}{\pgfqpoint{4.900000in}{4.900000in}}%
\pgfusepath{clip}%
\pgfsetbuttcap%
\pgfsetroundjoin%
\definecolor{currentfill}{rgb}{0.456901,0.456901,0.456901}%
\pgfsetfillcolor{currentfill}%
\pgfsetfillopacity{0.900000}%
\pgfsetlinewidth{0.501875pt}%
\definecolor{currentstroke}{rgb}{0.200000,0.200000,0.200000}%
\pgfsetstrokecolor{currentstroke}%
\pgfsetdash{}{0pt}%
\pgfpathmoveto{\pgfqpoint{1.698415in}{2.867690in}}%
\pgfpathlineto{\pgfqpoint{1.746200in}{2.821961in}}%
\pgfpathlineto{\pgfqpoint{1.777134in}{2.802126in}}%
\pgfpathlineto{\pgfqpoint{1.729341in}{2.847884in}}%
\pgfpathlineto{\pgfqpoint{1.698415in}{2.867690in}}%
\pgfpathclose%
\pgfusepath{stroke,fill}%
\end{pgfscope}%
\begin{pgfscope}%
\pgfpathrectangle{\pgfqpoint{0.050000in}{0.050000in}}{\pgfqpoint{4.900000in}{4.900000in}}%
\pgfusepath{clip}%
\pgfsetbuttcap%
\pgfsetroundjoin%
\definecolor{currentfill}{rgb}{0.878570,0.878570,0.878570}%
\pgfsetfillcolor{currentfill}%
\pgfsetfillopacity{0.900000}%
\pgfsetlinewidth{0.501875pt}%
\definecolor{currentstroke}{rgb}{0.200000,0.200000,0.200000}%
\pgfsetstrokecolor{currentstroke}%
\pgfsetdash{}{0pt}%
\pgfpathmoveto{\pgfqpoint{2.955245in}{1.862716in}}%
\pgfpathlineto{\pgfqpoint{3.000393in}{1.825310in}}%
\pgfpathlineto{\pgfqpoint{3.033294in}{1.804883in}}%
\pgfpathlineto{\pgfqpoint{2.988139in}{1.842086in}}%
\pgfpathlineto{\pgfqpoint{2.955245in}{1.862716in}}%
\pgfpathclose%
\pgfusepath{stroke,fill}%
\end{pgfscope}%
\begin{pgfscope}%
\pgfpathrectangle{\pgfqpoint{0.050000in}{0.050000in}}{\pgfqpoint{4.900000in}{4.900000in}}%
\pgfusepath{clip}%
\pgfsetbuttcap%
\pgfsetroundjoin%
\definecolor{currentfill}{rgb}{0.976009,0.976009,0.976009}%
\pgfsetfillcolor{currentfill}%
\pgfsetfillopacity{0.900000}%
\pgfsetlinewidth{0.501875pt}%
\definecolor{currentstroke}{rgb}{0.200000,0.200000,0.200000}%
\pgfsetstrokecolor{currentstroke}%
\pgfsetdash{}{0pt}%
\pgfpathmoveto{\pgfqpoint{3.691547in}{1.447041in}}%
\pgfpathlineto{\pgfqpoint{3.735840in}{1.441147in}}%
\pgfpathlineto{\pgfqpoint{3.769935in}{1.421177in}}%
\pgfpathlineto{\pgfqpoint{3.725622in}{1.427069in}}%
\pgfpathlineto{\pgfqpoint{3.691547in}{1.447041in}}%
\pgfpathclose%
\pgfusepath{stroke,fill}%
\end{pgfscope}%
\begin{pgfscope}%
\pgfpathrectangle{\pgfqpoint{0.050000in}{0.050000in}}{\pgfqpoint{4.900000in}{4.900000in}}%
\pgfusepath{clip}%
\pgfsetbuttcap%
\pgfsetroundjoin%
\definecolor{currentfill}{rgb}{0.734579,0.734579,0.734579}%
\pgfsetfillcolor{currentfill}%
\pgfsetfillopacity{0.900000}%
\pgfsetlinewidth{0.501875pt}%
\definecolor{currentstroke}{rgb}{0.200000,0.200000,0.200000}%
\pgfsetstrokecolor{currentstroke}%
\pgfsetdash{}{0pt}%
\pgfpathmoveto{\pgfqpoint{2.421211in}{2.283627in}}%
\pgfpathlineto{\pgfqpoint{2.467457in}{2.239283in}}%
\pgfpathlineto{\pgfqpoint{2.499518in}{2.218590in}}%
\pgfpathlineto{\pgfqpoint{2.453268in}{2.262933in}}%
\pgfpathlineto{\pgfqpoint{2.421211in}{2.283627in}}%
\pgfpathclose%
\pgfusepath{stroke,fill}%
\end{pgfscope}%
\begin{pgfscope}%
\pgfpathrectangle{\pgfqpoint{0.050000in}{0.050000in}}{\pgfqpoint{4.900000in}{4.900000in}}%
\pgfusepath{clip}%
\pgfsetbuttcap%
\pgfsetroundjoin%
\definecolor{currentfill}{rgb}{0.952018,0.952018,0.952018}%
\pgfsetfillcolor{currentfill}%
\pgfsetfillopacity{0.900000}%
\pgfsetlinewidth{0.501875pt}%
\definecolor{currentstroke}{rgb}{0.200000,0.200000,0.200000}%
\pgfsetstrokecolor{currentstroke}%
\pgfsetdash{}{0pt}%
\pgfpathmoveto{\pgfqpoint{3.378699in}{1.585642in}}%
\pgfpathlineto{\pgfqpoint{3.423239in}{1.565168in}}%
\pgfpathlineto{\pgfqpoint{3.456829in}{1.545312in}}%
\pgfpathlineto{\pgfqpoint{3.412275in}{1.565640in}}%
\pgfpathlineto{\pgfqpoint{3.378699in}{1.585642in}}%
\pgfpathclose%
\pgfusepath{stroke,fill}%
\end{pgfscope}%
\begin{pgfscope}%
\pgfpathrectangle{\pgfqpoint{0.050000in}{0.050000in}}{\pgfqpoint{4.900000in}{4.900000in}}%
\pgfusepath{clip}%
\pgfsetbuttcap%
\pgfsetroundjoin%
\definecolor{currentfill}{rgb}{0.367243,0.367243,0.367243}%
\pgfsetfillcolor{currentfill}%
\pgfsetfillopacity{0.900000}%
\pgfsetlinewidth{0.501875pt}%
\definecolor{currentstroke}{rgb}{0.200000,0.200000,0.200000}%
\pgfsetstrokecolor{currentstroke}%
\pgfsetdash{}{0pt}%
\pgfpathmoveto{\pgfqpoint{1.431199in}{3.084310in}}%
\pgfpathlineto{\pgfqpoint{1.479563in}{3.038111in}}%
\pgfpathlineto{\pgfqpoint{1.510086in}{3.018609in}}%
\pgfpathlineto{\pgfqpoint{1.461712in}{3.064839in}}%
\pgfpathlineto{\pgfqpoint{1.431199in}{3.084310in}}%
\pgfpathclose%
\pgfusepath{stroke,fill}%
\end{pgfscope}%
\begin{pgfscope}%
\pgfpathrectangle{\pgfqpoint{0.050000in}{0.050000in}}{\pgfqpoint{4.900000in}{4.900000in}}%
\pgfusepath{clip}%
\pgfsetbuttcap%
\pgfsetroundjoin%
\definecolor{currentfill}{rgb}{0.981546,0.981546,0.981546}%
\pgfsetfillcolor{currentfill}%
\pgfsetfillopacity{0.900000}%
\pgfsetlinewidth{0.501875pt}%
\definecolor{currentstroke}{rgb}{0.200000,0.200000,0.200000}%
\pgfsetstrokecolor{currentstroke}%
\pgfsetdash{}{0pt}%
\pgfpathmoveto{\pgfqpoint{3.769935in}{1.421177in}}%
\pgfpathlineto{\pgfqpoint{3.814186in}{1.418346in}}%
\pgfpathlineto{\pgfqpoint{3.848410in}{1.398272in}}%
\pgfpathlineto{\pgfqpoint{3.804138in}{1.401126in}}%
\pgfpathlineto{\pgfqpoint{3.769935in}{1.421177in}}%
\pgfpathclose%
\pgfusepath{stroke,fill}%
\end{pgfscope}%
\begin{pgfscope}%
\pgfpathrectangle{\pgfqpoint{0.050000in}{0.050000in}}{\pgfqpoint{4.900000in}{4.900000in}}%
\pgfusepath{clip}%
\pgfsetbuttcap%
\pgfsetroundjoin%
\definecolor{currentfill}{rgb}{0.624221,0.624221,0.624221}%
\pgfsetfillcolor{currentfill}%
\pgfsetfillopacity{0.900000}%
\pgfsetlinewidth{0.501875pt}%
\definecolor{currentstroke}{rgb}{0.200000,0.200000,0.200000}%
\pgfsetstrokecolor{currentstroke}%
\pgfsetdash{}{0pt}%
\pgfpathmoveto{\pgfqpoint{2.154100in}{2.500047in}}%
\pgfpathlineto{\pgfqpoint{2.200924in}{2.455095in}}%
\pgfpathlineto{\pgfqpoint{2.232575in}{2.434686in}}%
\pgfpathlineto{\pgfqpoint{2.185745in}{2.479663in}}%
\pgfpathlineto{\pgfqpoint{2.154100in}{2.500047in}}%
\pgfpathclose%
\pgfusepath{stroke,fill}%
\end{pgfscope}%
\begin{pgfscope}%
\pgfpathrectangle{\pgfqpoint{0.050000in}{0.050000in}}{\pgfqpoint{4.900000in}{4.900000in}}%
\pgfusepath{clip}%
\pgfsetbuttcap%
\pgfsetroundjoin%
\definecolor{currentfill}{rgb}{0.942791,0.942791,0.942791}%
\pgfsetfillcolor{currentfill}%
\pgfsetfillopacity{0.900000}%
\pgfsetlinewidth{0.501875pt}%
\definecolor{currentstroke}{rgb}{0.200000,0.200000,0.200000}%
\pgfsetstrokecolor{currentstroke}%
\pgfsetdash{}{0pt}%
\pgfpathmoveto{\pgfqpoint{3.300599in}{1.629876in}}%
\pgfpathlineto{\pgfqpoint{3.345231in}{1.605566in}}%
\pgfpathlineto{\pgfqpoint{3.378699in}{1.585642in}}%
\pgfpathlineto{\pgfqpoint{3.334055in}{1.609773in}}%
\pgfpathlineto{\pgfqpoint{3.300599in}{1.629876in}}%
\pgfpathclose%
\pgfusepath{stroke,fill}%
\end{pgfscope}%
\begin{pgfscope}%
\pgfpathrectangle{\pgfqpoint{0.050000in}{0.050000in}}{\pgfqpoint{4.900000in}{4.900000in}}%
\pgfusepath{clip}%
\pgfsetbuttcap%
\pgfsetroundjoin%
\definecolor{currentfill}{rgb}{0.256517,0.256517,0.256517}%
\pgfsetfillcolor{currentfill}%
\pgfsetfillopacity{0.900000}%
\pgfsetlinewidth{0.501875pt}%
\definecolor{currentstroke}{rgb}{0.200000,0.200000,0.200000}%
\pgfsetstrokecolor{currentstroke}%
\pgfsetdash{}{0pt}%
\pgfpathmoveto{\pgfqpoint{1.163892in}{3.301006in}}%
\pgfpathlineto{\pgfqpoint{1.212835in}{3.254336in}}%
\pgfpathlineto{\pgfqpoint{1.242947in}{3.235168in}}%
\pgfpathlineto{\pgfqpoint{1.193992in}{3.281869in}}%
\pgfpathlineto{\pgfqpoint{1.163892in}{3.301006in}}%
\pgfpathclose%
\pgfusepath{stroke,fill}%
\end{pgfscope}%
\begin{pgfscope}%
\pgfpathrectangle{\pgfqpoint{0.050000in}{0.050000in}}{\pgfqpoint{4.900000in}{4.900000in}}%
\pgfusepath{clip}%
\pgfsetbuttcap%
\pgfsetroundjoin%
\definecolor{currentfill}{rgb}{0.521492,0.521492,0.521492}%
\pgfsetfillcolor{currentfill}%
\pgfsetfillopacity{0.900000}%
\pgfsetlinewidth{0.501875pt}%
\definecolor{currentstroke}{rgb}{0.200000,0.200000,0.200000}%
\pgfsetstrokecolor{currentstroke}%
\pgfsetdash{}{0pt}%
\pgfpathmoveto{\pgfqpoint{1.886898in}{2.716650in}}%
\pgfpathlineto{\pgfqpoint{1.934300in}{2.671225in}}%
\pgfpathlineto{\pgfqpoint{1.965540in}{2.651147in}}%
\pgfpathlineto{\pgfqpoint{1.918131in}{2.696600in}}%
\pgfpathlineto{\pgfqpoint{1.886898in}{2.716650in}}%
\pgfpathclose%
\pgfusepath{stroke,fill}%
\end{pgfscope}%
\begin{pgfscope}%
\pgfpathrectangle{\pgfqpoint{0.050000in}{0.050000in}}{\pgfqpoint{4.900000in}{4.900000in}}%
\pgfusepath{clip}%
\pgfsetbuttcap%
\pgfsetroundjoin%
\definecolor{currentfill}{rgb}{0.985236,0.985236,0.985236}%
\pgfsetfillcolor{currentfill}%
\pgfsetfillopacity{0.900000}%
\pgfsetlinewidth{0.501875pt}%
\definecolor{currentstroke}{rgb}{0.200000,0.200000,0.200000}%
\pgfsetstrokecolor{currentstroke}%
\pgfsetdash{}{0pt}%
\pgfpathmoveto{\pgfqpoint{3.848410in}{1.398272in}}%
\pgfpathlineto{\pgfqpoint{3.892622in}{1.398189in}}%
\pgfpathlineto{\pgfqpoint{3.926977in}{1.377993in}}%
\pgfpathlineto{\pgfqpoint{3.882744in}{1.378119in}}%
\pgfpathlineto{\pgfqpoint{3.848410in}{1.398272in}}%
\pgfpathclose%
\pgfusepath{stroke,fill}%
\end{pgfscope}%
\begin{pgfscope}%
\pgfpathrectangle{\pgfqpoint{0.050000in}{0.050000in}}{\pgfqpoint{4.900000in}{4.900000in}}%
\pgfusepath{clip}%
\pgfsetbuttcap%
\pgfsetroundjoin%
\definecolor{currentfill}{rgb}{0.858762,0.858762,0.858762}%
\pgfsetfillcolor{currentfill}%
\pgfsetfillopacity{0.900000}%
\pgfsetlinewidth{0.501875pt}%
\definecolor{currentstroke}{rgb}{0.200000,0.200000,0.200000}%
\pgfsetstrokecolor{currentstroke}%
\pgfsetdash{}{0pt}%
\pgfpathmoveto{\pgfqpoint{2.877137in}{1.922857in}}%
\pgfpathlineto{\pgfqpoint{2.922454in}{1.883294in}}%
\pgfpathlineto{\pgfqpoint{2.955245in}{1.862716in}}%
\pgfpathlineto{\pgfqpoint{2.909923in}{1.902105in}}%
\pgfpathlineto{\pgfqpoint{2.877137in}{1.922857in}}%
\pgfpathclose%
\pgfusepath{stroke,fill}%
\end{pgfscope}%
\begin{pgfscope}%
\pgfpathrectangle{\pgfqpoint{0.050000in}{0.050000in}}{\pgfqpoint{4.900000in}{4.900000in}}%
\pgfusepath{clip}%
\pgfsetbuttcap%
\pgfsetroundjoin%
\definecolor{currentfill}{rgb}{0.929504,0.929504,0.929504}%
\pgfsetfillcolor{currentfill}%
\pgfsetfillopacity{0.900000}%
\pgfsetlinewidth{0.501875pt}%
\definecolor{currentstroke}{rgb}{0.200000,0.200000,0.200000}%
\pgfsetstrokecolor{currentstroke}%
\pgfsetdash{}{0pt}%
\pgfpathmoveto{\pgfqpoint{3.222511in}{1.677914in}}%
\pgfpathlineto{\pgfqpoint{3.267250in}{1.649907in}}%
\pgfpathlineto{\pgfqpoint{3.300599in}{1.629876in}}%
\pgfpathlineto{\pgfqpoint{3.255849in}{1.657681in}}%
\pgfpathlineto{\pgfqpoint{3.222511in}{1.677914in}}%
\pgfpathclose%
\pgfusepath{stroke,fill}%
\end{pgfscope}%
\begin{pgfscope}%
\pgfpathrectangle{\pgfqpoint{0.050000in}{0.050000in}}{\pgfqpoint{4.900000in}{4.900000in}}%
\pgfusepath{clip}%
\pgfsetbuttcap%
\pgfsetroundjoin%
\definecolor{currentfill}{rgb}{0.788112,0.788112,0.788112}%
\pgfsetfillcolor{currentfill}%
\pgfsetfillopacity{0.900000}%
\pgfsetlinewidth{0.501875pt}%
\definecolor{currentstroke}{rgb}{0.200000,0.200000,0.200000}%
\pgfsetstrokecolor{currentstroke}%
\pgfsetdash{}{0pt}%
\pgfpathmoveto{\pgfqpoint{2.610002in}{2.133088in}}%
\pgfpathlineto{\pgfqpoint{2.655865in}{2.089683in}}%
\pgfpathlineto{\pgfqpoint{2.688235in}{2.068900in}}%
\pgfpathlineto{\pgfqpoint{2.642368in}{2.112248in}}%
\pgfpathlineto{\pgfqpoint{2.610002in}{2.133088in}}%
\pgfpathclose%
\pgfusepath{stroke,fill}%
\end{pgfscope}%
\begin{pgfscope}%
\pgfpathrectangle{\pgfqpoint{0.050000in}{0.050000in}}{\pgfqpoint{4.900000in}{4.900000in}}%
\pgfusepath{clip}%
\pgfsetbuttcap%
\pgfsetroundjoin%
\definecolor{currentfill}{rgb}{0.132011,0.132011,0.132011}%
\pgfsetfillcolor{currentfill}%
\pgfsetfillopacity{0.900000}%
\pgfsetlinewidth{0.501875pt}%
\definecolor{currentstroke}{rgb}{0.200000,0.200000,0.200000}%
\pgfsetstrokecolor{currentstroke}%
\pgfsetdash{}{0pt}%
\pgfpathmoveto{\pgfqpoint{0.896495in}{3.517778in}}%
\pgfpathlineto{\pgfqpoint{0.946016in}{3.470638in}}%
\pgfpathlineto{\pgfqpoint{0.975717in}{3.451803in}}%
\pgfpathlineto{\pgfqpoint{0.926181in}{3.498977in}}%
\pgfpathlineto{\pgfqpoint{0.896495in}{3.517778in}}%
\pgfpathclose%
\pgfusepath{stroke,fill}%
\end{pgfscope}%
\begin{pgfscope}%
\pgfpathrectangle{\pgfqpoint{0.050000in}{0.050000in}}{\pgfqpoint{4.900000in}{4.900000in}}%
\pgfusepath{clip}%
\pgfsetbuttcap%
\pgfsetroundjoin%
\definecolor{currentfill}{rgb}{0.428143,0.428143,0.428143}%
\pgfsetfillcolor{currentfill}%
\pgfsetfillopacity{0.900000}%
\pgfsetlinewidth{0.501875pt}%
\definecolor{currentstroke}{rgb}{0.200000,0.200000,0.200000}%
\pgfsetstrokecolor{currentstroke}%
\pgfsetdash{}{0pt}%
\pgfpathmoveto{\pgfqpoint{1.619605in}{2.933330in}}%
\pgfpathlineto{\pgfqpoint{1.667586in}{2.887435in}}%
\pgfpathlineto{\pgfqpoint{1.698415in}{2.867690in}}%
\pgfpathlineto{\pgfqpoint{1.650425in}{2.913615in}}%
\pgfpathlineto{\pgfqpoint{1.619605in}{2.933330in}}%
\pgfpathclose%
\pgfusepath{stroke,fill}%
\end{pgfscope}%
\begin{pgfscope}%
\pgfpathrectangle{\pgfqpoint{0.050000in}{0.050000in}}{\pgfqpoint{4.900000in}{4.900000in}}%
\pgfusepath{clip}%
\pgfsetbuttcap%
\pgfsetroundjoin%
\definecolor{currentfill}{rgb}{0.696194,0.696194,0.696194}%
\pgfsetfillcolor{currentfill}%
\pgfsetfillopacity{0.900000}%
\pgfsetlinewidth{0.501875pt}%
\definecolor{currentstroke}{rgb}{0.200000,0.200000,0.200000}%
\pgfsetstrokecolor{currentstroke}%
\pgfsetdash{}{0pt}%
\pgfpathmoveto{\pgfqpoint{2.342813in}{2.348858in}}%
\pgfpathlineto{\pgfqpoint{2.389253in}{2.304260in}}%
\pgfpathlineto{\pgfqpoint{2.421211in}{2.283627in}}%
\pgfpathlineto{\pgfqpoint{2.374766in}{2.328238in}}%
\pgfpathlineto{\pgfqpoint{2.342813in}{2.348858in}}%
\pgfpathclose%
\pgfusepath{stroke,fill}%
\end{pgfscope}%
\begin{pgfscope}%
\pgfpathrectangle{\pgfqpoint{0.050000in}{0.050000in}}{\pgfqpoint{4.900000in}{4.900000in}}%
\pgfusepath{clip}%
\pgfsetbuttcap%
\pgfsetroundjoin%
\definecolor{currentfill}{rgb}{0.987082,0.987082,0.987082}%
\pgfsetfillcolor{currentfill}%
\pgfsetfillopacity{0.900000}%
\pgfsetlinewidth{0.501875pt}%
\definecolor{currentstroke}{rgb}{0.200000,0.200000,0.200000}%
\pgfsetstrokecolor{currentstroke}%
\pgfsetdash{}{0pt}%
\pgfpathmoveto{\pgfqpoint{3.926977in}{1.377993in}}%
\pgfpathlineto{\pgfqpoint{3.971152in}{1.380331in}}%
\pgfpathlineto{\pgfqpoint{4.005639in}{1.360001in}}%
\pgfpathlineto{\pgfqpoint{3.961442in}{1.357721in}}%
\pgfpathlineto{\pgfqpoint{3.926977in}{1.377993in}}%
\pgfpathclose%
\pgfusepath{stroke,fill}%
\end{pgfscope}%
\begin{pgfscope}%
\pgfpathrectangle{\pgfqpoint{0.050000in}{0.050000in}}{\pgfqpoint{4.900000in}{4.900000in}}%
\pgfusepath{clip}%
\pgfsetbuttcap%
\pgfsetroundjoin%
\definecolor{currentfill}{rgb}{0.334764,0.334764,0.334764}%
\pgfsetfillcolor{currentfill}%
\pgfsetfillopacity{0.900000}%
\pgfsetlinewidth{0.501875pt}%
\definecolor{currentstroke}{rgb}{0.200000,0.200000,0.200000}%
\pgfsetstrokecolor{currentstroke}%
\pgfsetdash{}{0pt}%
\pgfpathmoveto{\pgfqpoint{1.352222in}{3.150087in}}%
\pgfpathlineto{\pgfqpoint{1.400782in}{3.103721in}}%
\pgfpathlineto{\pgfqpoint{1.431199in}{3.084310in}}%
\pgfpathlineto{\pgfqpoint{1.382628in}{3.130707in}}%
\pgfpathlineto{\pgfqpoint{1.352222in}{3.150087in}}%
\pgfpathclose%
\pgfusepath{stroke,fill}%
\end{pgfscope}%
\begin{pgfscope}%
\pgfpathrectangle{\pgfqpoint{0.050000in}{0.050000in}}{\pgfqpoint{4.900000in}{4.900000in}}%
\pgfusepath{clip}%
\pgfsetbuttcap%
\pgfsetroundjoin%
\definecolor{currentfill}{rgb}{0.915356,0.915356,0.915356}%
\pgfsetfillcolor{currentfill}%
\pgfsetfillopacity{0.900000}%
\pgfsetlinewidth{0.501875pt}%
\definecolor{currentstroke}{rgb}{0.200000,0.200000,0.200000}%
\pgfsetstrokecolor{currentstroke}%
\pgfsetdash{}{0pt}%
\pgfpathmoveto{\pgfqpoint{3.144417in}{1.729535in}}%
\pgfpathlineto{\pgfqpoint{3.189280in}{1.698082in}}%
\pgfpathlineto{\pgfqpoint{3.222511in}{1.677914in}}%
\pgfpathlineto{\pgfqpoint{3.177639in}{1.709152in}}%
\pgfpathlineto{\pgfqpoint{3.144417in}{1.729535in}}%
\pgfpathclose%
\pgfusepath{stroke,fill}%
\end{pgfscope}%
\begin{pgfscope}%
\pgfpathrectangle{\pgfqpoint{0.050000in}{0.050000in}}{\pgfqpoint{4.900000in}{4.900000in}}%
\pgfusepath{clip}%
\pgfsetbuttcap%
\pgfsetroundjoin%
\definecolor{currentfill}{rgb}{0.590634,0.590634,0.590634}%
\pgfsetfillcolor{currentfill}%
\pgfsetfillopacity{0.900000}%
\pgfsetlinewidth{0.501875pt}%
\definecolor{currentstroke}{rgb}{0.200000,0.200000,0.200000}%
\pgfsetstrokecolor{currentstroke}%
\pgfsetdash{}{0pt}%
\pgfpathmoveto{\pgfqpoint{2.075534in}{2.565488in}}%
\pgfpathlineto{\pgfqpoint{2.122553in}{2.520368in}}%
\pgfpathlineto{\pgfqpoint{2.154100in}{2.500047in}}%
\pgfpathlineto{\pgfqpoint{2.107075in}{2.545194in}}%
\pgfpathlineto{\pgfqpoint{2.075534in}{2.565488in}}%
\pgfpathclose%
\pgfusepath{stroke,fill}%
\end{pgfscope}%
\begin{pgfscope}%
\pgfpathrectangle{\pgfqpoint{0.050000in}{0.050000in}}{\pgfqpoint{4.900000in}{4.900000in}}%
\pgfusepath{clip}%
\pgfsetbuttcap%
\pgfsetroundjoin%
\definecolor{currentfill}{rgb}{0.217762,0.217762,0.217762}%
\pgfsetfillcolor{currentfill}%
\pgfsetfillopacity{0.900000}%
\pgfsetlinewidth{0.501875pt}%
\definecolor{currentstroke}{rgb}{0.200000,0.200000,0.200000}%
\pgfsetstrokecolor{currentstroke}%
\pgfsetdash{}{0pt}%
\pgfpathmoveto{\pgfqpoint{1.084747in}{3.366920in}}%
\pgfpathlineto{\pgfqpoint{1.133886in}{3.320083in}}%
\pgfpathlineto{\pgfqpoint{1.163892in}{3.301006in}}%
\pgfpathlineto{\pgfqpoint{1.114740in}{3.347875in}}%
\pgfpathlineto{\pgfqpoint{1.084747in}{3.366920in}}%
\pgfpathclose%
\pgfusepath{stroke,fill}%
\end{pgfscope}%
\begin{pgfscope}%
\pgfpathrectangle{\pgfqpoint{0.050000in}{0.050000in}}{\pgfqpoint{4.900000in}{4.900000in}}%
\pgfusepath{clip}%
\pgfsetbuttcap%
\pgfsetroundjoin%
\definecolor{currentfill}{rgb}{0.990773,0.990773,0.990773}%
\pgfsetfillcolor{currentfill}%
\pgfsetfillopacity{0.900000}%
\pgfsetlinewidth{0.501875pt}%
\definecolor{currentstroke}{rgb}{0.200000,0.200000,0.200000}%
\pgfsetstrokecolor{currentstroke}%
\pgfsetdash{}{0pt}%
\pgfpathmoveto{\pgfqpoint{4.005639in}{1.360001in}}%
\pgfpathlineto{\pgfqpoint{4.049776in}{1.364436in}}%
\pgfpathlineto{\pgfqpoint{4.084396in}{1.343966in}}%
\pgfpathlineto{\pgfqpoint{4.040236in}{1.339598in}}%
\pgfpathlineto{\pgfqpoint{4.005639in}{1.360001in}}%
\pgfpathclose%
\pgfusepath{stroke,fill}%
\end{pgfscope}%
\begin{pgfscope}%
\pgfpathrectangle{\pgfqpoint{0.050000in}{0.050000in}}{\pgfqpoint{4.900000in}{4.900000in}}%
\pgfusepath{clip}%
\pgfsetbuttcap%
\pgfsetroundjoin%
\definecolor{currentfill}{rgb}{0.836340,0.836340,0.836340}%
\pgfsetfillcolor{currentfill}%
\pgfsetfillopacity{0.900000}%
\pgfsetlinewidth{0.501875pt}%
\definecolor{currentstroke}{rgb}{0.200000,0.200000,0.200000}%
\pgfsetstrokecolor{currentstroke}%
\pgfsetdash{}{0pt}%
\pgfpathmoveto{\pgfqpoint{2.798961in}{1.984799in}}%
\pgfpathlineto{\pgfqpoint{2.844455in}{1.943560in}}%
\pgfpathlineto{\pgfqpoint{2.877137in}{1.922857in}}%
\pgfpathlineto{\pgfqpoint{2.831640in}{1.963959in}}%
\pgfpathlineto{\pgfqpoint{2.798961in}{1.984799in}}%
\pgfpathclose%
\pgfusepath{stroke,fill}%
\end{pgfscope}%
\begin{pgfscope}%
\pgfpathrectangle{\pgfqpoint{0.050000in}{0.050000in}}{\pgfqpoint{4.900000in}{4.900000in}}%
\pgfusepath{clip}%
\pgfsetbuttcap%
\pgfsetroundjoin%
\definecolor{currentfill}{rgb}{0.491349,0.491349,0.491349}%
\pgfsetfillcolor{currentfill}%
\pgfsetfillopacity{0.900000}%
\pgfsetlinewidth{0.501875pt}%
\definecolor{currentstroke}{rgb}{0.200000,0.200000,0.200000}%
\pgfsetstrokecolor{currentstroke}%
\pgfsetdash{}{0pt}%
\pgfpathmoveto{\pgfqpoint{1.808165in}{2.782229in}}%
\pgfpathlineto{\pgfqpoint{1.855763in}{2.736638in}}%
\pgfpathlineto{\pgfqpoint{1.886898in}{2.716650in}}%
\pgfpathlineto{\pgfqpoint{1.839292in}{2.762270in}}%
\pgfpathlineto{\pgfqpoint{1.808165in}{2.782229in}}%
\pgfpathclose%
\pgfusepath{stroke,fill}%
\end{pgfscope}%
\begin{pgfscope}%
\pgfpathrectangle{\pgfqpoint{0.050000in}{0.050000in}}{\pgfqpoint{4.900000in}{4.900000in}}%
\pgfusepath{clip}%
\pgfsetbuttcap%
\pgfsetroundjoin%
\definecolor{currentfill}{rgb}{0.760554,0.760554,0.760554}%
\pgfsetfillcolor{currentfill}%
\pgfsetfillopacity{0.900000}%
\pgfsetlinewidth{0.501875pt}%
\definecolor{currentstroke}{rgb}{0.200000,0.200000,0.200000}%
\pgfsetstrokecolor{currentstroke}%
\pgfsetdash{}{0pt}%
\pgfpathmoveto{\pgfqpoint{2.531680in}{2.197838in}}%
\pgfpathlineto{\pgfqpoint{2.577736in}{2.153872in}}%
\pgfpathlineto{\pgfqpoint{2.610002in}{2.133088in}}%
\pgfpathlineto{\pgfqpoint{2.563942in}{2.177028in}}%
\pgfpathlineto{\pgfqpoint{2.531680in}{2.197838in}}%
\pgfpathclose%
\pgfusepath{stroke,fill}%
\end{pgfscope}%
\begin{pgfscope}%
\pgfpathrectangle{\pgfqpoint{0.050000in}{0.050000in}}{\pgfqpoint{4.900000in}{4.900000in}}%
\pgfusepath{clip}%
\pgfsetbuttcap%
\pgfsetroundjoin%
\definecolor{currentfill}{rgb}{0.100146,0.100146,0.100146}%
\pgfsetfillcolor{currentfill}%
\pgfsetfillopacity{0.900000}%
\pgfsetlinewidth{0.501875pt}%
\definecolor{currentstroke}{rgb}{0.200000,0.200000,0.200000}%
\pgfsetstrokecolor{currentstroke}%
\pgfsetdash{}{0pt}%
\pgfpathmoveto{\pgfqpoint{0.817181in}{3.583829in}}%
\pgfpathlineto{\pgfqpoint{0.866901in}{3.536520in}}%
\pgfpathlineto{\pgfqpoint{0.896495in}{3.517778in}}%
\pgfpathlineto{\pgfqpoint{0.846761in}{3.565120in}}%
\pgfpathlineto{\pgfqpoint{0.817181in}{3.583829in}}%
\pgfpathclose%
\pgfusepath{stroke,fill}%
\end{pgfscope}%
\begin{pgfscope}%
\pgfpathrectangle{\pgfqpoint{0.050000in}{0.050000in}}{\pgfqpoint{4.900000in}{4.900000in}}%
\pgfusepath{clip}%
\pgfsetbuttcap%
\pgfsetroundjoin%
\definecolor{currentfill}{rgb}{0.395663,0.395663,0.395663}%
\pgfsetfillcolor{currentfill}%
\pgfsetfillopacity{0.900000}%
\pgfsetlinewidth{0.501875pt}%
\definecolor{currentstroke}{rgb}{0.200000,0.200000,0.200000}%
\pgfsetstrokecolor{currentstroke}%
\pgfsetdash{}{0pt}%
\pgfpathmoveto{\pgfqpoint{1.540704in}{2.999046in}}%
\pgfpathlineto{\pgfqpoint{1.588881in}{2.952984in}}%
\pgfpathlineto{\pgfqpoint{1.619605in}{2.933330in}}%
\pgfpathlineto{\pgfqpoint{1.571418in}{2.979423in}}%
\pgfpathlineto{\pgfqpoint{1.540704in}{2.999046in}}%
\pgfpathclose%
\pgfusepath{stroke,fill}%
\end{pgfscope}%
\begin{pgfscope}%
\pgfpathrectangle{\pgfqpoint{0.050000in}{0.050000in}}{\pgfqpoint{4.900000in}{4.900000in}}%
\pgfusepath{clip}%
\pgfsetbuttcap%
\pgfsetroundjoin%
\definecolor{currentfill}{rgb}{0.898378,0.898378,0.898378}%
\pgfsetfillcolor{currentfill}%
\pgfsetfillopacity{0.900000}%
\pgfsetlinewidth{0.501875pt}%
\definecolor{currentstroke}{rgb}{0.200000,0.200000,0.200000}%
\pgfsetstrokecolor{currentstroke}%
\pgfsetdash{}{0pt}%
\pgfpathmoveto{\pgfqpoint{3.066298in}{1.784402in}}%
\pgfpathlineto{\pgfqpoint{3.111300in}{1.749857in}}%
\pgfpathlineto{\pgfqpoint{3.144417in}{1.729535in}}%
\pgfpathlineto{\pgfqpoint{3.099408in}{1.763864in}}%
\pgfpathlineto{\pgfqpoint{3.066298in}{1.784402in}}%
\pgfpathclose%
\pgfusepath{stroke,fill}%
\end{pgfscope}%
\begin{pgfscope}%
\pgfpathrectangle{\pgfqpoint{0.050000in}{0.050000in}}{\pgfqpoint{4.900000in}{4.900000in}}%
\pgfusepath{clip}%
\pgfsetbuttcap%
\pgfsetroundjoin%
\definecolor{currentfill}{rgb}{0.662607,0.662607,0.662607}%
\pgfsetfillcolor{currentfill}%
\pgfsetfillopacity{0.900000}%
\pgfsetlinewidth{0.501875pt}%
\definecolor{currentstroke}{rgb}{0.200000,0.200000,0.200000}%
\pgfsetstrokecolor{currentstroke}%
\pgfsetdash{}{0pt}%
\pgfpathmoveto{\pgfqpoint{2.264324in}{2.414213in}}%
\pgfpathlineto{\pgfqpoint{2.310959in}{2.369414in}}%
\pgfpathlineto{\pgfqpoint{2.342813in}{2.348858in}}%
\pgfpathlineto{\pgfqpoint{2.296174in}{2.393678in}}%
\pgfpathlineto{\pgfqpoint{2.264324in}{2.414213in}}%
\pgfpathclose%
\pgfusepath{stroke,fill}%
\end{pgfscope}%
\begin{pgfscope}%
\pgfpathrectangle{\pgfqpoint{0.050000in}{0.050000in}}{\pgfqpoint{4.900000in}{4.900000in}}%
\pgfusepath{clip}%
\pgfsetbuttcap%
\pgfsetroundjoin%
\definecolor{currentfill}{rgb}{0.992618,0.992618,0.992618}%
\pgfsetfillcolor{currentfill}%
\pgfsetfillopacity{0.900000}%
\pgfsetlinewidth{0.501875pt}%
\definecolor{currentstroke}{rgb}{0.200000,0.200000,0.200000}%
\pgfsetstrokecolor{currentstroke}%
\pgfsetdash{}{0pt}%
\pgfpathmoveto{\pgfqpoint{4.084396in}{1.343966in}}%
\pgfpathlineto{\pgfqpoint{4.128494in}{1.350186in}}%
\pgfpathlineto{\pgfqpoint{4.163247in}{1.329573in}}%
\pgfpathlineto{\pgfqpoint{4.119126in}{1.323424in}}%
\pgfpathlineto{\pgfqpoint{4.084396in}{1.343966in}}%
\pgfpathclose%
\pgfusepath{stroke,fill}%
\end{pgfscope}%
\begin{pgfscope}%
\pgfpathrectangle{\pgfqpoint{0.050000in}{0.050000in}}{\pgfqpoint{4.900000in}{4.900000in}}%
\pgfusepath{clip}%
\pgfsetbuttcap%
\pgfsetroundjoin%
\definecolor{currentfill}{rgb}{0.300807,0.300807,0.300807}%
\pgfsetfillcolor{currentfill}%
\pgfsetfillopacity{0.900000}%
\pgfsetlinewidth{0.501875pt}%
\definecolor{currentstroke}{rgb}{0.200000,0.200000,0.200000}%
\pgfsetstrokecolor{currentstroke}%
\pgfsetdash{}{0pt}%
\pgfpathmoveto{\pgfqpoint{1.273153in}{3.215940in}}%
\pgfpathlineto{\pgfqpoint{1.321910in}{3.169406in}}%
\pgfpathlineto{\pgfqpoint{1.352222in}{3.150087in}}%
\pgfpathlineto{\pgfqpoint{1.303453in}{3.196652in}}%
\pgfpathlineto{\pgfqpoint{1.273153in}{3.215940in}}%
\pgfpathclose%
\pgfusepath{stroke,fill}%
\end{pgfscope}%
\begin{pgfscope}%
\pgfpathrectangle{\pgfqpoint{0.050000in}{0.050000in}}{\pgfqpoint{4.900000in}{4.900000in}}%
\pgfusepath{clip}%
\pgfsetbuttcap%
\pgfsetroundjoin%
\definecolor{currentfill}{rgb}{0.555940,0.555940,0.555940}%
\pgfsetfillcolor{currentfill}%
\pgfsetfillopacity{0.900000}%
\pgfsetlinewidth{0.501875pt}%
\definecolor{currentstroke}{rgb}{0.200000,0.200000,0.200000}%
\pgfsetstrokecolor{currentstroke}%
\pgfsetdash{}{0pt}%
\pgfpathmoveto{\pgfqpoint{1.996878in}{2.631006in}}%
\pgfpathlineto{\pgfqpoint{2.044092in}{2.585719in}}%
\pgfpathlineto{\pgfqpoint{2.075534in}{2.565488in}}%
\pgfpathlineto{\pgfqpoint{2.028314in}{2.610802in}}%
\pgfpathlineto{\pgfqpoint{1.996878in}{2.631006in}}%
\pgfpathclose%
\pgfusepath{stroke,fill}%
\end{pgfscope}%
\begin{pgfscope}%
\pgfpathrectangle{\pgfqpoint{0.050000in}{0.050000in}}{\pgfqpoint{4.900000in}{4.900000in}}%
\pgfusepath{clip}%
\pgfsetbuttcap%
\pgfsetroundjoin%
\definecolor{currentfill}{rgb}{0.173472,0.173472,0.173472}%
\pgfsetfillcolor{currentfill}%
\pgfsetfillopacity{0.900000}%
\pgfsetlinewidth{0.501875pt}%
\definecolor{currentstroke}{rgb}{0.200000,0.200000,0.200000}%
\pgfsetstrokecolor{currentstroke}%
\pgfsetdash{}{0pt}%
\pgfpathmoveto{\pgfqpoint{1.005510in}{3.432910in}}%
\pgfpathlineto{\pgfqpoint{1.054847in}{3.385905in}}%
\pgfpathlineto{\pgfqpoint{1.084747in}{3.366920in}}%
\pgfpathlineto{\pgfqpoint{1.035397in}{3.413958in}}%
\pgfpathlineto{\pgfqpoint{1.005510in}{3.432910in}}%
\pgfpathclose%
\pgfusepath{stroke,fill}%
\end{pgfscope}%
\begin{pgfscope}%
\pgfpathrectangle{\pgfqpoint{0.050000in}{0.050000in}}{\pgfqpoint{4.900000in}{4.900000in}}%
\pgfusepath{clip}%
\pgfsetbuttcap%
\pgfsetroundjoin%
\definecolor{currentfill}{rgb}{0.812226,0.812226,0.812226}%
\pgfsetfillcolor{currentfill}%
\pgfsetfillopacity{0.900000}%
\pgfsetlinewidth{0.501875pt}%
\definecolor{currentstroke}{rgb}{0.200000,0.200000,0.200000}%
\pgfsetstrokecolor{currentstroke}%
\pgfsetdash{}{0pt}%
\pgfpathmoveto{\pgfqpoint{2.720707in}{2.048065in}}%
\pgfpathlineto{\pgfqpoint{2.766384in}{2.005589in}}%
\pgfpathlineto{\pgfqpoint{2.798961in}{1.984799in}}%
\pgfpathlineto{\pgfqpoint{2.753280in}{2.027178in}}%
\pgfpathlineto{\pgfqpoint{2.720707in}{2.048065in}}%
\pgfpathclose%
\pgfusepath{stroke,fill}%
\end{pgfscope}%
\begin{pgfscope}%
\pgfpathrectangle{\pgfqpoint{0.050000in}{0.050000in}}{\pgfqpoint{4.900000in}{4.900000in}}%
\pgfusepath{clip}%
\pgfsetbuttcap%
\pgfsetroundjoin%
\definecolor{currentfill}{rgb}{0.964937,0.964937,0.964937}%
\pgfsetfillcolor{currentfill}%
\pgfsetfillopacity{0.900000}%
\pgfsetlinewidth{0.501875pt}%
\definecolor{currentstroke}{rgb}{0.200000,0.200000,0.200000}%
\pgfsetstrokecolor{currentstroke}%
\pgfsetdash{}{0pt}%
\pgfpathmoveto{\pgfqpoint{3.568826in}{1.488956in}}%
\pgfpathlineto{\pgfqpoint{3.613240in}{1.476180in}}%
\pgfpathlineto{\pgfqpoint{3.647187in}{1.456256in}}%
\pgfpathlineto{\pgfqpoint{3.602756in}{1.468964in}}%
\pgfpathlineto{\pgfqpoint{3.568826in}{1.488956in}}%
\pgfpathclose%
\pgfusepath{stroke,fill}%
\end{pgfscope}%
\begin{pgfscope}%
\pgfpathrectangle{\pgfqpoint{0.050000in}{0.050000in}}{\pgfqpoint{4.900000in}{4.900000in}}%
\pgfusepath{clip}%
\pgfsetbuttcap%
\pgfsetroundjoin%
\definecolor{currentfill}{rgb}{0.456901,0.456901,0.456901}%
\pgfsetfillcolor{currentfill}%
\pgfsetfillopacity{0.900000}%
\pgfsetlinewidth{0.501875pt}%
\definecolor{currentstroke}{rgb}{0.200000,0.200000,0.200000}%
\pgfsetstrokecolor{currentstroke}%
\pgfsetdash{}{0pt}%
\pgfpathmoveto{\pgfqpoint{1.729341in}{2.847884in}}%
\pgfpathlineto{\pgfqpoint{1.777134in}{2.802126in}}%
\pgfpathlineto{\pgfqpoint{1.808165in}{2.782229in}}%
\pgfpathlineto{\pgfqpoint{1.760363in}{2.828016in}}%
\pgfpathlineto{\pgfqpoint{1.729341in}{2.847884in}}%
\pgfpathclose%
\pgfusepath{stroke,fill}%
\end{pgfscope}%
\begin{pgfscope}%
\pgfpathrectangle{\pgfqpoint{0.050000in}{0.050000in}}{\pgfqpoint{4.900000in}{4.900000in}}%
\pgfusepath{clip}%
\pgfsetbuttcap%
\pgfsetroundjoin%
\definecolor{currentfill}{rgb}{0.972318,0.972318,0.972318}%
\pgfsetfillcolor{currentfill}%
\pgfsetfillopacity{0.900000}%
\pgfsetlinewidth{0.501875pt}%
\definecolor{currentstroke}{rgb}{0.200000,0.200000,0.200000}%
\pgfsetstrokecolor{currentstroke}%
\pgfsetdash{}{0pt}%
\pgfpathmoveto{\pgfqpoint{3.647187in}{1.456256in}}%
\pgfpathlineto{\pgfqpoint{3.691547in}{1.447041in}}%
\pgfpathlineto{\pgfqpoint{3.725622in}{1.427069in}}%
\pgfpathlineto{\pgfqpoint{3.681243in}{1.436251in}}%
\pgfpathlineto{\pgfqpoint{3.647187in}{1.456256in}}%
\pgfpathclose%
\pgfusepath{stroke,fill}%
\end{pgfscope}%
\begin{pgfscope}%
\pgfpathrectangle{\pgfqpoint{0.050000in}{0.050000in}}{\pgfqpoint{4.900000in}{4.900000in}}%
\pgfusepath{clip}%
\pgfsetbuttcap%
\pgfsetroundjoin%
\definecolor{currentfill}{rgb}{0.959400,0.959400,0.959400}%
\pgfsetfillcolor{currentfill}%
\pgfsetfillopacity{0.900000}%
\pgfsetlinewidth{0.501875pt}%
\definecolor{currentstroke}{rgb}{0.200000,0.200000,0.200000}%
\pgfsetstrokecolor{currentstroke}%
\pgfsetdash{}{0pt}%
\pgfpathmoveto{\pgfqpoint{3.490527in}{1.525376in}}%
\pgfpathlineto{\pgfqpoint{3.535005in}{1.508866in}}%
\pgfpathlineto{\pgfqpoint{3.568826in}{1.488956in}}%
\pgfpathlineto{\pgfqpoint{3.524331in}{1.505360in}}%
\pgfpathlineto{\pgfqpoint{3.490527in}{1.525376in}}%
\pgfpathclose%
\pgfusepath{stroke,fill}%
\end{pgfscope}%
\begin{pgfscope}%
\pgfpathrectangle{\pgfqpoint{0.050000in}{0.050000in}}{\pgfqpoint{4.900000in}{4.900000in}}%
\pgfusepath{clip}%
\pgfsetbuttcap%
\pgfsetroundjoin%
\definecolor{currentfill}{rgb}{0.878570,0.878570,0.878570}%
\pgfsetfillcolor{currentfill}%
\pgfsetfillopacity{0.900000}%
\pgfsetlinewidth{0.501875pt}%
\definecolor{currentstroke}{rgb}{0.200000,0.200000,0.200000}%
\pgfsetstrokecolor{currentstroke}%
\pgfsetdash{}{0pt}%
\pgfpathmoveto{\pgfqpoint{2.988139in}{1.842086in}}%
\pgfpathlineto{\pgfqpoint{3.033294in}{1.804883in}}%
\pgfpathlineto{\pgfqpoint{3.066298in}{1.784402in}}%
\pgfpathlineto{\pgfqpoint{3.021137in}{1.821404in}}%
\pgfpathlineto{\pgfqpoint{2.988139in}{1.842086in}}%
\pgfpathclose%
\pgfusepath{stroke,fill}%
\end{pgfscope}%
\begin{pgfscope}%
\pgfpathrectangle{\pgfqpoint{0.050000in}{0.050000in}}{\pgfqpoint{4.900000in}{4.900000in}}%
\pgfusepath{clip}%
\pgfsetbuttcap%
\pgfsetroundjoin%
\definecolor{currentfill}{rgb}{0.734579,0.734579,0.734579}%
\pgfsetfillcolor{currentfill}%
\pgfsetfillopacity{0.900000}%
\pgfsetlinewidth{0.501875pt}%
\definecolor{currentstroke}{rgb}{0.200000,0.200000,0.200000}%
\pgfsetstrokecolor{currentstroke}%
\pgfsetdash{}{0pt}%
\pgfpathmoveto{\pgfqpoint{2.453268in}{2.262933in}}%
\pgfpathlineto{\pgfqpoint{2.499518in}{2.218590in}}%
\pgfpathlineto{\pgfqpoint{2.531680in}{2.197838in}}%
\pgfpathlineto{\pgfqpoint{2.485426in}{2.242179in}}%
\pgfpathlineto{\pgfqpoint{2.453268in}{2.262933in}}%
\pgfpathclose%
\pgfusepath{stroke,fill}%
\end{pgfscope}%
\begin{pgfscope}%
\pgfpathrectangle{\pgfqpoint{0.050000in}{0.050000in}}{\pgfqpoint{4.900000in}{4.900000in}}%
\pgfusepath{clip}%
\pgfsetbuttcap%
\pgfsetroundjoin%
\definecolor{currentfill}{rgb}{0.976009,0.976009,0.976009}%
\pgfsetfillcolor{currentfill}%
\pgfsetfillopacity{0.900000}%
\pgfsetlinewidth{0.501875pt}%
\definecolor{currentstroke}{rgb}{0.200000,0.200000,0.200000}%
\pgfsetstrokecolor{currentstroke}%
\pgfsetdash{}{0pt}%
\pgfpathmoveto{\pgfqpoint{3.725622in}{1.427069in}}%
\pgfpathlineto{\pgfqpoint{3.769935in}{1.421177in}}%
\pgfpathlineto{\pgfqpoint{3.804138in}{1.401126in}}%
\pgfpathlineto{\pgfqpoint{3.759805in}{1.407015in}}%
\pgfpathlineto{\pgfqpoint{3.725622in}{1.427069in}}%
\pgfpathclose%
\pgfusepath{stroke,fill}%
\end{pgfscope}%
\begin{pgfscope}%
\pgfpathrectangle{\pgfqpoint{0.050000in}{0.050000in}}{\pgfqpoint{4.900000in}{4.900000in}}%
\pgfusepath{clip}%
\pgfsetbuttcap%
\pgfsetroundjoin%
\definecolor{currentfill}{rgb}{0.063729,0.063729,0.063729}%
\pgfsetfillcolor{currentfill}%
\pgfsetfillopacity{0.900000}%
\pgfsetlinewidth{0.501875pt}%
\definecolor{currentstroke}{rgb}{0.200000,0.200000,0.200000}%
\pgfsetstrokecolor{currentstroke}%
\pgfsetdash{}{0pt}%
\pgfpathmoveto{\pgfqpoint{0.737776in}{3.649956in}}%
\pgfpathlineto{\pgfqpoint{0.787694in}{3.602479in}}%
\pgfpathlineto{\pgfqpoint{0.817181in}{3.583829in}}%
\pgfpathlineto{\pgfqpoint{0.767249in}{3.631340in}}%
\pgfpathlineto{\pgfqpoint{0.737776in}{3.649956in}}%
\pgfpathclose%
\pgfusepath{stroke,fill}%
\end{pgfscope}%
\begin{pgfscope}%
\pgfpathrectangle{\pgfqpoint{0.050000in}{0.050000in}}{\pgfqpoint{4.900000in}{4.900000in}}%
\pgfusepath{clip}%
\pgfsetbuttcap%
\pgfsetroundjoin%
\definecolor{currentfill}{rgb}{0.952018,0.952018,0.952018}%
\pgfsetfillcolor{currentfill}%
\pgfsetfillopacity{0.900000}%
\pgfsetlinewidth{0.501875pt}%
\definecolor{currentstroke}{rgb}{0.200000,0.200000,0.200000}%
\pgfsetstrokecolor{currentstroke}%
\pgfsetdash{}{0pt}%
\pgfpathmoveto{\pgfqpoint{3.412275in}{1.565640in}}%
\pgfpathlineto{\pgfqpoint{3.456829in}{1.545312in}}%
\pgfpathlineto{\pgfqpoint{3.490527in}{1.525376in}}%
\pgfpathlineto{\pgfqpoint{3.445957in}{1.545563in}}%
\pgfpathlineto{\pgfqpoint{3.412275in}{1.565640in}}%
\pgfpathclose%
\pgfusepath{stroke,fill}%
\end{pgfscope}%
\begin{pgfscope}%
\pgfpathrectangle{\pgfqpoint{0.050000in}{0.050000in}}{\pgfqpoint{4.900000in}{4.900000in}}%
\pgfusepath{clip}%
\pgfsetbuttcap%
\pgfsetroundjoin%
\definecolor{currentfill}{rgb}{0.367243,0.367243,0.367243}%
\pgfsetfillcolor{currentfill}%
\pgfsetfillopacity{0.900000}%
\pgfsetlinewidth{0.501875pt}%
\definecolor{currentstroke}{rgb}{0.200000,0.200000,0.200000}%
\pgfsetstrokecolor{currentstroke}%
\pgfsetdash{}{0pt}%
\pgfpathmoveto{\pgfqpoint{1.461712in}{3.064839in}}%
\pgfpathlineto{\pgfqpoint{1.510086in}{3.018609in}}%
\pgfpathlineto{\pgfqpoint{1.540704in}{2.999046in}}%
\pgfpathlineto{\pgfqpoint{1.492320in}{3.045306in}}%
\pgfpathlineto{\pgfqpoint{1.461712in}{3.064839in}}%
\pgfpathclose%
\pgfusepath{stroke,fill}%
\end{pgfscope}%
\begin{pgfscope}%
\pgfpathrectangle{\pgfqpoint{0.050000in}{0.050000in}}{\pgfqpoint{4.900000in}{4.900000in}}%
\pgfusepath{clip}%
\pgfsetbuttcap%
\pgfsetroundjoin%
\definecolor{currentfill}{rgb}{0.624221,0.624221,0.624221}%
\pgfsetfillcolor{currentfill}%
\pgfsetfillopacity{0.900000}%
\pgfsetlinewidth{0.501875pt}%
\definecolor{currentstroke}{rgb}{0.200000,0.200000,0.200000}%
\pgfsetstrokecolor{currentstroke}%
\pgfsetdash{}{0pt}%
\pgfpathmoveto{\pgfqpoint{2.185745in}{2.479663in}}%
\pgfpathlineto{\pgfqpoint{2.232575in}{2.434686in}}%
\pgfpathlineto{\pgfqpoint{2.264324in}{2.414213in}}%
\pgfpathlineto{\pgfqpoint{2.217490in}{2.459215in}}%
\pgfpathlineto{\pgfqpoint{2.185745in}{2.479663in}}%
\pgfpathclose%
\pgfusepath{stroke,fill}%
\end{pgfscope}%
\begin{pgfscope}%
\pgfpathrectangle{\pgfqpoint{0.050000in}{0.050000in}}{\pgfqpoint{4.900000in}{4.900000in}}%
\pgfusepath{clip}%
\pgfsetbuttcap%
\pgfsetroundjoin%
\definecolor{currentfill}{rgb}{0.981546,0.981546,0.981546}%
\pgfsetfillcolor{currentfill}%
\pgfsetfillopacity{0.900000}%
\pgfsetlinewidth{0.501875pt}%
\definecolor{currentstroke}{rgb}{0.200000,0.200000,0.200000}%
\pgfsetstrokecolor{currentstroke}%
\pgfsetdash{}{0pt}%
\pgfpathmoveto{\pgfqpoint{3.804138in}{1.401126in}}%
\pgfpathlineto{\pgfqpoint{3.848410in}{1.398272in}}%
\pgfpathlineto{\pgfqpoint{3.882744in}{1.378119in}}%
\pgfpathlineto{\pgfqpoint{3.838451in}{1.380995in}}%
\pgfpathlineto{\pgfqpoint{3.804138in}{1.401126in}}%
\pgfpathclose%
\pgfusepath{stroke,fill}%
\end{pgfscope}%
\begin{pgfscope}%
\pgfpathrectangle{\pgfqpoint{0.050000in}{0.050000in}}{\pgfqpoint{4.900000in}{4.900000in}}%
\pgfusepath{clip}%
\pgfsetbuttcap%
\pgfsetroundjoin%
\definecolor{currentfill}{rgb}{0.942791,0.942791,0.942791}%
\pgfsetfillcolor{currentfill}%
\pgfsetfillopacity{0.900000}%
\pgfsetlinewidth{0.501875pt}%
\definecolor{currentstroke}{rgb}{0.200000,0.200000,0.200000}%
\pgfsetstrokecolor{currentstroke}%
\pgfsetdash{}{0pt}%
\pgfpathmoveto{\pgfqpoint{3.334055in}{1.609773in}}%
\pgfpathlineto{\pgfqpoint{3.378699in}{1.585642in}}%
\pgfpathlineto{\pgfqpoint{3.412275in}{1.565640in}}%
\pgfpathlineto{\pgfqpoint{3.367616in}{1.589599in}}%
\pgfpathlineto{\pgfqpoint{3.334055in}{1.609773in}}%
\pgfpathclose%
\pgfusepath{stroke,fill}%
\end{pgfscope}%
\begin{pgfscope}%
\pgfpathrectangle{\pgfqpoint{0.050000in}{0.050000in}}{\pgfqpoint{4.900000in}{4.900000in}}%
\pgfusepath{clip}%
\pgfsetbuttcap%
\pgfsetroundjoin%
\definecolor{currentfill}{rgb}{0.256517,0.256517,0.256517}%
\pgfsetfillcolor{currentfill}%
\pgfsetfillopacity{0.900000}%
\pgfsetlinewidth{0.501875pt}%
\definecolor{currentstroke}{rgb}{0.200000,0.200000,0.200000}%
\pgfsetstrokecolor{currentstroke}%
\pgfsetdash{}{0pt}%
\pgfpathmoveto{\pgfqpoint{1.193992in}{3.281869in}}%
\pgfpathlineto{\pgfqpoint{1.242947in}{3.235168in}}%
\pgfpathlineto{\pgfqpoint{1.273153in}{3.215940in}}%
\pgfpathlineto{\pgfqpoint{1.224186in}{3.262673in}}%
\pgfpathlineto{\pgfqpoint{1.193992in}{3.281869in}}%
\pgfpathclose%
\pgfusepath{stroke,fill}%
\end{pgfscope}%
\begin{pgfscope}%
\pgfpathrectangle{\pgfqpoint{0.050000in}{0.050000in}}{\pgfqpoint{4.900000in}{4.900000in}}%
\pgfusepath{clip}%
\pgfsetbuttcap%
\pgfsetroundjoin%
\definecolor{currentfill}{rgb}{0.525798,0.525798,0.525798}%
\pgfsetfillcolor{currentfill}%
\pgfsetfillopacity{0.900000}%
\pgfsetlinewidth{0.501875pt}%
\definecolor{currentstroke}{rgb}{0.200000,0.200000,0.200000}%
\pgfsetstrokecolor{currentstroke}%
\pgfsetdash{}{0pt}%
\pgfpathmoveto{\pgfqpoint{1.918131in}{2.696600in}}%
\pgfpathlineto{\pgfqpoint{1.965540in}{2.651147in}}%
\pgfpathlineto{\pgfqpoint{1.996878in}{2.631006in}}%
\pgfpathlineto{\pgfqpoint{1.949461in}{2.676487in}}%
\pgfpathlineto{\pgfqpoint{1.918131in}{2.696600in}}%
\pgfpathclose%
\pgfusepath{stroke,fill}%
\end{pgfscope}%
\begin{pgfscope}%
\pgfpathrectangle{\pgfqpoint{0.050000in}{0.050000in}}{\pgfqpoint{4.900000in}{4.900000in}}%
\pgfusepath{clip}%
\pgfsetbuttcap%
\pgfsetroundjoin%
\definecolor{currentfill}{rgb}{0.985236,0.985236,0.985236}%
\pgfsetfillcolor{currentfill}%
\pgfsetfillopacity{0.900000}%
\pgfsetlinewidth{0.501875pt}%
\definecolor{currentstroke}{rgb}{0.200000,0.200000,0.200000}%
\pgfsetstrokecolor{currentstroke}%
\pgfsetdash{}{0pt}%
\pgfpathmoveto{\pgfqpoint{3.882744in}{1.378119in}}%
\pgfpathlineto{\pgfqpoint{3.926977in}{1.377993in}}%
\pgfpathlineto{\pgfqpoint{3.961442in}{1.357721in}}%
\pgfpathlineto{\pgfqpoint{3.917187in}{1.357888in}}%
\pgfpathlineto{\pgfqpoint{3.882744in}{1.378119in}}%
\pgfpathclose%
\pgfusepath{stroke,fill}%
\end{pgfscope}%
\begin{pgfscope}%
\pgfpathrectangle{\pgfqpoint{0.050000in}{0.050000in}}{\pgfqpoint{4.900000in}{4.900000in}}%
\pgfusepath{clip}%
\pgfsetbuttcap%
\pgfsetroundjoin%
\definecolor{currentfill}{rgb}{0.858762,0.858762,0.858762}%
\pgfsetfillcolor{currentfill}%
\pgfsetfillopacity{0.900000}%
\pgfsetlinewidth{0.501875pt}%
\definecolor{currentstroke}{rgb}{0.200000,0.200000,0.200000}%
\pgfsetstrokecolor{currentstroke}%
\pgfsetdash{}{0pt}%
\pgfpathmoveto{\pgfqpoint{2.909923in}{1.902105in}}%
\pgfpathlineto{\pgfqpoint{2.955245in}{1.862716in}}%
\pgfpathlineto{\pgfqpoint{2.988139in}{1.842086in}}%
\pgfpathlineto{\pgfqpoint{2.942813in}{1.881301in}}%
\pgfpathlineto{\pgfqpoint{2.909923in}{1.902105in}}%
\pgfpathclose%
\pgfusepath{stroke,fill}%
\end{pgfscope}%
\begin{pgfscope}%
\pgfpathrectangle{\pgfqpoint{0.050000in}{0.050000in}}{\pgfqpoint{4.900000in}{4.900000in}}%
\pgfusepath{clip}%
\pgfsetbuttcap%
\pgfsetroundjoin%
\definecolor{currentfill}{rgb}{0.788112,0.788112,0.788112}%
\pgfsetfillcolor{currentfill}%
\pgfsetfillopacity{0.900000}%
\pgfsetlinewidth{0.501875pt}%
\definecolor{currentstroke}{rgb}{0.200000,0.200000,0.200000}%
\pgfsetstrokecolor{currentstroke}%
\pgfsetdash{}{0pt}%
\pgfpathmoveto{\pgfqpoint{2.642368in}{2.112248in}}%
\pgfpathlineto{\pgfqpoint{2.688235in}{2.068900in}}%
\pgfpathlineto{\pgfqpoint{2.720707in}{2.048065in}}%
\pgfpathlineto{\pgfqpoint{2.674837in}{2.091352in}}%
\pgfpathlineto{\pgfqpoint{2.642368in}{2.112248in}}%
\pgfpathclose%
\pgfusepath{stroke,fill}%
\end{pgfscope}%
\begin{pgfscope}%
\pgfpathrectangle{\pgfqpoint{0.050000in}{0.050000in}}{\pgfqpoint{4.900000in}{4.900000in}}%
\pgfusepath{clip}%
\pgfsetbuttcap%
\pgfsetroundjoin%
\definecolor{currentfill}{rgb}{0.132011,0.132011,0.132011}%
\pgfsetfillcolor{currentfill}%
\pgfsetfillopacity{0.900000}%
\pgfsetlinewidth{0.501875pt}%
\definecolor{currentstroke}{rgb}{0.200000,0.200000,0.200000}%
\pgfsetstrokecolor{currentstroke}%
\pgfsetdash{}{0pt}%
\pgfpathmoveto{\pgfqpoint{0.926181in}{3.498977in}}%
\pgfpathlineto{\pgfqpoint{0.975717in}{3.451803in}}%
\pgfpathlineto{\pgfqpoint{1.005510in}{3.432910in}}%
\pgfpathlineto{\pgfqpoint{0.955961in}{3.480117in}}%
\pgfpathlineto{\pgfqpoint{0.926181in}{3.498977in}}%
\pgfpathclose%
\pgfusepath{stroke,fill}%
\end{pgfscope}%
\begin{pgfscope}%
\pgfpathrectangle{\pgfqpoint{0.050000in}{0.050000in}}{\pgfqpoint{4.900000in}{4.900000in}}%
\pgfusepath{clip}%
\pgfsetbuttcap%
\pgfsetroundjoin%
\definecolor{currentfill}{rgb}{0.929504,0.929504,0.929504}%
\pgfsetfillcolor{currentfill}%
\pgfsetfillopacity{0.900000}%
\pgfsetlinewidth{0.501875pt}%
\definecolor{currentstroke}{rgb}{0.200000,0.200000,0.200000}%
\pgfsetstrokecolor{currentstroke}%
\pgfsetdash{}{0pt}%
\pgfpathmoveto{\pgfqpoint{3.255849in}{1.657681in}}%
\pgfpathlineto{\pgfqpoint{3.300599in}{1.629876in}}%
\pgfpathlineto{\pgfqpoint{3.334055in}{1.609773in}}%
\pgfpathlineto{\pgfqpoint{3.289292in}{1.637380in}}%
\pgfpathlineto{\pgfqpoint{3.255849in}{1.657681in}}%
\pgfpathclose%
\pgfusepath{stroke,fill}%
\end{pgfscope}%
\begin{pgfscope}%
\pgfpathrectangle{\pgfqpoint{0.050000in}{0.050000in}}{\pgfqpoint{4.900000in}{4.900000in}}%
\pgfusepath{clip}%
\pgfsetbuttcap%
\pgfsetroundjoin%
\definecolor{currentfill}{rgb}{0.428143,0.428143,0.428143}%
\pgfsetfillcolor{currentfill}%
\pgfsetfillopacity{0.900000}%
\pgfsetlinewidth{0.501875pt}%
\definecolor{currentstroke}{rgb}{0.200000,0.200000,0.200000}%
\pgfsetstrokecolor{currentstroke}%
\pgfsetdash{}{0pt}%
\pgfpathmoveto{\pgfqpoint{1.650425in}{2.913615in}}%
\pgfpathlineto{\pgfqpoint{1.698415in}{2.867690in}}%
\pgfpathlineto{\pgfqpoint{1.729341in}{2.847884in}}%
\pgfpathlineto{\pgfqpoint{1.681342in}{2.893839in}}%
\pgfpathlineto{\pgfqpoint{1.650425in}{2.913615in}}%
\pgfpathclose%
\pgfusepath{stroke,fill}%
\end{pgfscope}%
\begin{pgfscope}%
\pgfpathrectangle{\pgfqpoint{0.050000in}{0.050000in}}{\pgfqpoint{4.900000in}{4.900000in}}%
\pgfusepath{clip}%
\pgfsetbuttcap%
\pgfsetroundjoin%
\definecolor{currentfill}{rgb}{0.696194,0.696194,0.696194}%
\pgfsetfillcolor{currentfill}%
\pgfsetfillopacity{0.900000}%
\pgfsetlinewidth{0.501875pt}%
\definecolor{currentstroke}{rgb}{0.200000,0.200000,0.200000}%
\pgfsetstrokecolor{currentstroke}%
\pgfsetdash{}{0pt}%
\pgfpathmoveto{\pgfqpoint{2.374766in}{2.328238in}}%
\pgfpathlineto{\pgfqpoint{2.421211in}{2.283627in}}%
\pgfpathlineto{\pgfqpoint{2.453268in}{2.262933in}}%
\pgfpathlineto{\pgfqpoint{2.406820in}{2.307557in}}%
\pgfpathlineto{\pgfqpoint{2.374766in}{2.328238in}}%
\pgfpathclose%
\pgfusepath{stroke,fill}%
\end{pgfscope}%
\begin{pgfscope}%
\pgfpathrectangle{\pgfqpoint{0.050000in}{0.050000in}}{\pgfqpoint{4.900000in}{4.900000in}}%
\pgfusepath{clip}%
\pgfsetbuttcap%
\pgfsetroundjoin%
\definecolor{currentfill}{rgb}{0.987082,0.987082,0.987082}%
\pgfsetfillcolor{currentfill}%
\pgfsetfillopacity{0.900000}%
\pgfsetlinewidth{0.501875pt}%
\definecolor{currentstroke}{rgb}{0.200000,0.200000,0.200000}%
\pgfsetstrokecolor{currentstroke}%
\pgfsetdash{}{0pt}%
\pgfpathmoveto{\pgfqpoint{3.961442in}{1.357721in}}%
\pgfpathlineto{\pgfqpoint{4.005639in}{1.360001in}}%
\pgfpathlineto{\pgfqpoint{4.040236in}{1.339598in}}%
\pgfpathlineto{\pgfqpoint{3.996017in}{1.337372in}}%
\pgfpathlineto{\pgfqpoint{3.961442in}{1.357721in}}%
\pgfpathclose%
\pgfusepath{stroke,fill}%
\end{pgfscope}%
\begin{pgfscope}%
\pgfpathrectangle{\pgfqpoint{0.050000in}{0.050000in}}{\pgfqpoint{4.900000in}{4.900000in}}%
\pgfusepath{clip}%
\pgfsetbuttcap%
\pgfsetroundjoin%
\definecolor{currentfill}{rgb}{0.334764,0.334764,0.334764}%
\pgfsetfillcolor{currentfill}%
\pgfsetfillopacity{0.900000}%
\pgfsetlinewidth{0.501875pt}%
\definecolor{currentstroke}{rgb}{0.200000,0.200000,0.200000}%
\pgfsetstrokecolor{currentstroke}%
\pgfsetdash{}{0pt}%
\pgfpathmoveto{\pgfqpoint{1.382628in}{3.130707in}}%
\pgfpathlineto{\pgfqpoint{1.431199in}{3.084310in}}%
\pgfpathlineto{\pgfqpoint{1.461712in}{3.064839in}}%
\pgfpathlineto{\pgfqpoint{1.413131in}{3.111267in}}%
\pgfpathlineto{\pgfqpoint{1.382628in}{3.130707in}}%
\pgfpathclose%
\pgfusepath{stroke,fill}%
\end{pgfscope}%
\begin{pgfscope}%
\pgfpathrectangle{\pgfqpoint{0.050000in}{0.050000in}}{\pgfqpoint{4.900000in}{4.900000in}}%
\pgfusepath{clip}%
\pgfsetbuttcap%
\pgfsetroundjoin%
\definecolor{currentfill}{rgb}{0.590634,0.590634,0.590634}%
\pgfsetfillcolor{currentfill}%
\pgfsetfillopacity{0.900000}%
\pgfsetlinewidth{0.501875pt}%
\definecolor{currentstroke}{rgb}{0.200000,0.200000,0.200000}%
\pgfsetstrokecolor{currentstroke}%
\pgfsetdash{}{0pt}%
\pgfpathmoveto{\pgfqpoint{2.107075in}{2.545194in}}%
\pgfpathlineto{\pgfqpoint{2.154100in}{2.500047in}}%
\pgfpathlineto{\pgfqpoint{2.185745in}{2.479663in}}%
\pgfpathlineto{\pgfqpoint{2.138715in}{2.524836in}}%
\pgfpathlineto{\pgfqpoint{2.107075in}{2.545194in}}%
\pgfpathclose%
\pgfusepath{stroke,fill}%
\end{pgfscope}%
\begin{pgfscope}%
\pgfpathrectangle{\pgfqpoint{0.050000in}{0.050000in}}{\pgfqpoint{4.900000in}{4.900000in}}%
\pgfusepath{clip}%
\pgfsetbuttcap%
\pgfsetroundjoin%
\definecolor{currentfill}{rgb}{0.912526,0.912526,0.912526}%
\pgfsetfillcolor{currentfill}%
\pgfsetfillopacity{0.900000}%
\pgfsetlinewidth{0.501875pt}%
\definecolor{currentstroke}{rgb}{0.200000,0.200000,0.200000}%
\pgfsetstrokecolor{currentstroke}%
\pgfsetdash{}{0pt}%
\pgfpathmoveto{\pgfqpoint{3.177639in}{1.709152in}}%
\pgfpathlineto{\pgfqpoint{3.222511in}{1.677914in}}%
\pgfpathlineto{\pgfqpoint{3.255849in}{1.657681in}}%
\pgfpathlineto{\pgfqpoint{3.210967in}{1.688707in}}%
\pgfpathlineto{\pgfqpoint{3.177639in}{1.709152in}}%
\pgfpathclose%
\pgfusepath{stroke,fill}%
\end{pgfscope}%
\begin{pgfscope}%
\pgfpathrectangle{\pgfqpoint{0.050000in}{0.050000in}}{\pgfqpoint{4.900000in}{4.900000in}}%
\pgfusepath{clip}%
\pgfsetbuttcap%
\pgfsetroundjoin%
\definecolor{currentfill}{rgb}{0.217762,0.217762,0.217762}%
\pgfsetfillcolor{currentfill}%
\pgfsetfillopacity{0.900000}%
\pgfsetlinewidth{0.501875pt}%
\definecolor{currentstroke}{rgb}{0.200000,0.200000,0.200000}%
\pgfsetstrokecolor{currentstroke}%
\pgfsetdash{}{0pt}%
\pgfpathmoveto{\pgfqpoint{1.114740in}{3.347875in}}%
\pgfpathlineto{\pgfqpoint{1.163892in}{3.301006in}}%
\pgfpathlineto{\pgfqpoint{1.193992in}{3.281869in}}%
\pgfpathlineto{\pgfqpoint{1.144828in}{3.328771in}}%
\pgfpathlineto{\pgfqpoint{1.114740in}{3.347875in}}%
\pgfpathclose%
\pgfusepath{stroke,fill}%
\end{pgfscope}%
\begin{pgfscope}%
\pgfpathrectangle{\pgfqpoint{0.050000in}{0.050000in}}{\pgfqpoint{4.900000in}{4.900000in}}%
\pgfusepath{clip}%
\pgfsetbuttcap%
\pgfsetroundjoin%
\definecolor{currentfill}{rgb}{0.990773,0.990773,0.990773}%
\pgfsetfillcolor{currentfill}%
\pgfsetfillopacity{0.900000}%
\pgfsetlinewidth{0.501875pt}%
\definecolor{currentstroke}{rgb}{0.200000,0.200000,0.200000}%
\pgfsetstrokecolor{currentstroke}%
\pgfsetdash{}{0pt}%
\pgfpathmoveto{\pgfqpoint{4.040236in}{1.339598in}}%
\pgfpathlineto{\pgfqpoint{4.084396in}{1.343966in}}%
\pgfpathlineto{\pgfqpoint{4.119126in}{1.323424in}}%
\pgfpathlineto{\pgfqpoint{4.074943in}{1.319119in}}%
\pgfpathlineto{\pgfqpoint{4.040236in}{1.339598in}}%
\pgfpathclose%
\pgfusepath{stroke,fill}%
\end{pgfscope}%
\begin{pgfscope}%
\pgfpathrectangle{\pgfqpoint{0.050000in}{0.050000in}}{\pgfqpoint{4.900000in}{4.900000in}}%
\pgfusepath{clip}%
\pgfsetbuttcap%
\pgfsetroundjoin%
\definecolor{currentfill}{rgb}{0.491349,0.491349,0.491349}%
\pgfsetfillcolor{currentfill}%
\pgfsetfillopacity{0.900000}%
\pgfsetlinewidth{0.501875pt}%
\definecolor{currentstroke}{rgb}{0.200000,0.200000,0.200000}%
\pgfsetstrokecolor{currentstroke}%
\pgfsetdash{}{0pt}%
\pgfpathmoveto{\pgfqpoint{1.839292in}{2.762270in}}%
\pgfpathlineto{\pgfqpoint{1.886898in}{2.716650in}}%
\pgfpathlineto{\pgfqpoint{1.918131in}{2.696600in}}%
\pgfpathlineto{\pgfqpoint{1.870518in}{2.742248in}}%
\pgfpathlineto{\pgfqpoint{1.839292in}{2.762270in}}%
\pgfpathclose%
\pgfusepath{stroke,fill}%
\end{pgfscope}%
\begin{pgfscope}%
\pgfpathrectangle{\pgfqpoint{0.050000in}{0.050000in}}{\pgfqpoint{4.900000in}{4.900000in}}%
\pgfusepath{clip}%
\pgfsetbuttcap%
\pgfsetroundjoin%
\definecolor{currentfill}{rgb}{0.836340,0.836340,0.836340}%
\pgfsetfillcolor{currentfill}%
\pgfsetfillopacity{0.900000}%
\pgfsetlinewidth{0.501875pt}%
\definecolor{currentstroke}{rgb}{0.200000,0.200000,0.200000}%
\pgfsetstrokecolor{currentstroke}%
\pgfsetdash{}{0pt}%
\pgfpathmoveto{\pgfqpoint{2.831640in}{1.963959in}}%
\pgfpathlineto{\pgfqpoint{2.877137in}{1.922857in}}%
\pgfpathlineto{\pgfqpoint{2.909923in}{1.902105in}}%
\pgfpathlineto{\pgfqpoint{2.864422in}{1.943068in}}%
\pgfpathlineto{\pgfqpoint{2.831640in}{1.963959in}}%
\pgfpathclose%
\pgfusepath{stroke,fill}%
\end{pgfscope}%
\begin{pgfscope}%
\pgfpathrectangle{\pgfqpoint{0.050000in}{0.050000in}}{\pgfqpoint{4.900000in}{4.900000in}}%
\pgfusepath{clip}%
\pgfsetbuttcap%
\pgfsetroundjoin%
\definecolor{currentfill}{rgb}{0.760554,0.760554,0.760554}%
\pgfsetfillcolor{currentfill}%
\pgfsetfillopacity{0.900000}%
\pgfsetlinewidth{0.501875pt}%
\definecolor{currentstroke}{rgb}{0.200000,0.200000,0.200000}%
\pgfsetstrokecolor{currentstroke}%
\pgfsetdash{}{0pt}%
\pgfpathmoveto{\pgfqpoint{2.563942in}{2.177028in}}%
\pgfpathlineto{\pgfqpoint{2.610002in}{2.133088in}}%
\pgfpathlineto{\pgfqpoint{2.642368in}{2.112248in}}%
\pgfpathlineto{\pgfqpoint{2.596306in}{2.156159in}}%
\pgfpathlineto{\pgfqpoint{2.563942in}{2.177028in}}%
\pgfpathclose%
\pgfusepath{stroke,fill}%
\end{pgfscope}%
\begin{pgfscope}%
\pgfpathrectangle{\pgfqpoint{0.050000in}{0.050000in}}{\pgfqpoint{4.900000in}{4.900000in}}%
\pgfusepath{clip}%
\pgfsetbuttcap%
\pgfsetroundjoin%
\definecolor{currentfill}{rgb}{0.100146,0.100146,0.100146}%
\pgfsetfillcolor{currentfill}%
\pgfsetfillopacity{0.900000}%
\pgfsetlinewidth{0.501875pt}%
\definecolor{currentstroke}{rgb}{0.200000,0.200000,0.200000}%
\pgfsetstrokecolor{currentstroke}%
\pgfsetdash{}{0pt}%
\pgfpathmoveto{\pgfqpoint{0.846761in}{3.565120in}}%
\pgfpathlineto{\pgfqpoint{0.896495in}{3.517778in}}%
\pgfpathlineto{\pgfqpoint{0.926181in}{3.498977in}}%
\pgfpathlineto{\pgfqpoint{0.876433in}{3.546353in}}%
\pgfpathlineto{\pgfqpoint{0.846761in}{3.565120in}}%
\pgfpathclose%
\pgfusepath{stroke,fill}%
\end{pgfscope}%
\begin{pgfscope}%
\pgfpathrectangle{\pgfqpoint{0.050000in}{0.050000in}}{\pgfqpoint{4.900000in}{4.900000in}}%
\pgfusepath{clip}%
\pgfsetbuttcap%
\pgfsetroundjoin%
\definecolor{currentfill}{rgb}{0.395663,0.395663,0.395663}%
\pgfsetfillcolor{currentfill}%
\pgfsetfillopacity{0.900000}%
\pgfsetlinewidth{0.501875pt}%
\definecolor{currentstroke}{rgb}{0.200000,0.200000,0.200000}%
\pgfsetstrokecolor{currentstroke}%
\pgfsetdash{}{0pt}%
\pgfpathmoveto{\pgfqpoint{1.571418in}{2.979423in}}%
\pgfpathlineto{\pgfqpoint{1.619605in}{2.933330in}}%
\pgfpathlineto{\pgfqpoint{1.650425in}{2.913615in}}%
\pgfpathlineto{\pgfqpoint{1.602229in}{2.959737in}}%
\pgfpathlineto{\pgfqpoint{1.571418in}{2.979423in}}%
\pgfpathclose%
\pgfusepath{stroke,fill}%
\end{pgfscope}%
\begin{pgfscope}%
\pgfpathrectangle{\pgfqpoint{0.050000in}{0.050000in}}{\pgfqpoint{4.900000in}{4.900000in}}%
\pgfusepath{clip}%
\pgfsetbuttcap%
\pgfsetroundjoin%
\definecolor{currentfill}{rgb}{0.895548,0.895548,0.895548}%
\pgfsetfillcolor{currentfill}%
\pgfsetfillopacity{0.900000}%
\pgfsetlinewidth{0.501875pt}%
\definecolor{currentstroke}{rgb}{0.200000,0.200000,0.200000}%
\pgfsetstrokecolor{currentstroke}%
\pgfsetdash{}{0pt}%
\pgfpathmoveto{\pgfqpoint{3.099408in}{1.763864in}}%
\pgfpathlineto{\pgfqpoint{3.144417in}{1.729535in}}%
\pgfpathlineto{\pgfqpoint{3.177639in}{1.709152in}}%
\pgfpathlineto{\pgfqpoint{3.132622in}{1.743270in}}%
\pgfpathlineto{\pgfqpoint{3.099408in}{1.763864in}}%
\pgfpathclose%
\pgfusepath{stroke,fill}%
\end{pgfscope}%
\begin{pgfscope}%
\pgfpathrectangle{\pgfqpoint{0.050000in}{0.050000in}}{\pgfqpoint{4.900000in}{4.900000in}}%
\pgfusepath{clip}%
\pgfsetbuttcap%
\pgfsetroundjoin%
\definecolor{currentfill}{rgb}{0.662607,0.662607,0.662607}%
\pgfsetfillcolor{currentfill}%
\pgfsetfillopacity{0.900000}%
\pgfsetlinewidth{0.501875pt}%
\definecolor{currentstroke}{rgb}{0.200000,0.200000,0.200000}%
\pgfsetstrokecolor{currentstroke}%
\pgfsetdash{}{0pt}%
\pgfpathmoveto{\pgfqpoint{2.296174in}{2.393678in}}%
\pgfpathlineto{\pgfqpoint{2.342813in}{2.348858in}}%
\pgfpathlineto{\pgfqpoint{2.374766in}{2.328238in}}%
\pgfpathlineto{\pgfqpoint{2.328123in}{2.373079in}}%
\pgfpathlineto{\pgfqpoint{2.296174in}{2.393678in}}%
\pgfpathclose%
\pgfusepath{stroke,fill}%
\end{pgfscope}%
\begin{pgfscope}%
\pgfpathrectangle{\pgfqpoint{0.050000in}{0.050000in}}{\pgfqpoint{4.900000in}{4.900000in}}%
\pgfusepath{clip}%
\pgfsetbuttcap%
\pgfsetroundjoin%
\definecolor{currentfill}{rgb}{0.300807,0.300807,0.300807}%
\pgfsetfillcolor{currentfill}%
\pgfsetfillopacity{0.900000}%
\pgfsetlinewidth{0.501875pt}%
\definecolor{currentstroke}{rgb}{0.200000,0.200000,0.200000}%
\pgfsetstrokecolor{currentstroke}%
\pgfsetdash{}{0pt}%
\pgfpathmoveto{\pgfqpoint{1.303453in}{3.196652in}}%
\pgfpathlineto{\pgfqpoint{1.352222in}{3.150087in}}%
\pgfpathlineto{\pgfqpoint{1.382628in}{3.130707in}}%
\pgfpathlineto{\pgfqpoint{1.333849in}{3.177303in}}%
\pgfpathlineto{\pgfqpoint{1.303453in}{3.196652in}}%
\pgfpathclose%
\pgfusepath{stroke,fill}%
\end{pgfscope}%
\begin{pgfscope}%
\pgfpathrectangle{\pgfqpoint{0.050000in}{0.050000in}}{\pgfqpoint{4.900000in}{4.900000in}}%
\pgfusepath{clip}%
\pgfsetbuttcap%
\pgfsetroundjoin%
\definecolor{currentfill}{rgb}{0.555940,0.555940,0.555940}%
\pgfsetfillcolor{currentfill}%
\pgfsetfillopacity{0.900000}%
\pgfsetlinewidth{0.501875pt}%
\definecolor{currentstroke}{rgb}{0.200000,0.200000,0.200000}%
\pgfsetstrokecolor{currentstroke}%
\pgfsetdash{}{0pt}%
\pgfpathmoveto{\pgfqpoint{2.028314in}{2.610802in}}%
\pgfpathlineto{\pgfqpoint{2.075534in}{2.565488in}}%
\pgfpathlineto{\pgfqpoint{2.107075in}{2.545194in}}%
\pgfpathlineto{\pgfqpoint{2.059848in}{2.590536in}}%
\pgfpathlineto{\pgfqpoint{2.028314in}{2.610802in}}%
\pgfpathclose%
\pgfusepath{stroke,fill}%
\end{pgfscope}%
\begin{pgfscope}%
\pgfpathrectangle{\pgfqpoint{0.050000in}{0.050000in}}{\pgfqpoint{4.900000in}{4.900000in}}%
\pgfusepath{clip}%
\pgfsetbuttcap%
\pgfsetroundjoin%
\definecolor{currentfill}{rgb}{0.173472,0.173472,0.173472}%
\pgfsetfillcolor{currentfill}%
\pgfsetfillopacity{0.900000}%
\pgfsetlinewidth{0.501875pt}%
\definecolor{currentstroke}{rgb}{0.200000,0.200000,0.200000}%
\pgfsetstrokecolor{currentstroke}%
\pgfsetdash{}{0pt}%
\pgfpathmoveto{\pgfqpoint{1.035397in}{3.413958in}}%
\pgfpathlineto{\pgfqpoint{1.084747in}{3.366920in}}%
\pgfpathlineto{\pgfqpoint{1.114740in}{3.347875in}}%
\pgfpathlineto{\pgfqpoint{1.065377in}{3.394946in}}%
\pgfpathlineto{\pgfqpoint{1.035397in}{3.413958in}}%
\pgfpathclose%
\pgfusepath{stroke,fill}%
\end{pgfscope}%
\begin{pgfscope}%
\pgfpathrectangle{\pgfqpoint{0.050000in}{0.050000in}}{\pgfqpoint{4.900000in}{4.900000in}}%
\pgfusepath{clip}%
\pgfsetbuttcap%
\pgfsetroundjoin%
\definecolor{currentfill}{rgb}{0.812226,0.812226,0.812226}%
\pgfsetfillcolor{currentfill}%
\pgfsetfillopacity{0.900000}%
\pgfsetlinewidth{0.501875pt}%
\definecolor{currentstroke}{rgb}{0.200000,0.200000,0.200000}%
\pgfsetstrokecolor{currentstroke}%
\pgfsetdash{}{0pt}%
\pgfpathmoveto{\pgfqpoint{2.753280in}{2.027178in}}%
\pgfpathlineto{\pgfqpoint{2.798961in}{1.984799in}}%
\pgfpathlineto{\pgfqpoint{2.831640in}{1.963959in}}%
\pgfpathlineto{\pgfqpoint{2.785956in}{2.006237in}}%
\pgfpathlineto{\pgfqpoint{2.753280in}{2.027178in}}%
\pgfpathclose%
\pgfusepath{stroke,fill}%
\end{pgfscope}%
\begin{pgfscope}%
\pgfpathrectangle{\pgfqpoint{0.050000in}{0.050000in}}{\pgfqpoint{4.900000in}{4.900000in}}%
\pgfusepath{clip}%
\pgfsetbuttcap%
\pgfsetroundjoin%
\definecolor{currentfill}{rgb}{0.461207,0.461207,0.461207}%
\pgfsetfillcolor{currentfill}%
\pgfsetfillopacity{0.900000}%
\pgfsetlinewidth{0.501875pt}%
\definecolor{currentstroke}{rgb}{0.200000,0.200000,0.200000}%
\pgfsetstrokecolor{currentstroke}%
\pgfsetdash{}{0pt}%
\pgfpathmoveto{\pgfqpoint{1.760363in}{2.828016in}}%
\pgfpathlineto{\pgfqpoint{1.808165in}{2.782229in}}%
\pgfpathlineto{\pgfqpoint{1.839292in}{2.762270in}}%
\pgfpathlineto{\pgfqpoint{1.791482in}{2.808086in}}%
\pgfpathlineto{\pgfqpoint{1.760363in}{2.828016in}}%
\pgfpathclose%
\pgfusepath{stroke,fill}%
\end{pgfscope}%
\begin{pgfscope}%
\pgfpathrectangle{\pgfqpoint{0.050000in}{0.050000in}}{\pgfqpoint{4.900000in}{4.900000in}}%
\pgfusepath{clip}%
\pgfsetbuttcap%
\pgfsetroundjoin%
\definecolor{currentfill}{rgb}{0.964937,0.964937,0.964937}%
\pgfsetfillcolor{currentfill}%
\pgfsetfillopacity{0.900000}%
\pgfsetlinewidth{0.501875pt}%
\definecolor{currentstroke}{rgb}{0.200000,0.200000,0.200000}%
\pgfsetstrokecolor{currentstroke}%
\pgfsetdash{}{0pt}%
\pgfpathmoveto{\pgfqpoint{3.602756in}{1.468964in}}%
\pgfpathlineto{\pgfqpoint{3.647187in}{1.456256in}}%
\pgfpathlineto{\pgfqpoint{3.681243in}{1.436251in}}%
\pgfpathlineto{\pgfqpoint{3.636794in}{1.448892in}}%
\pgfpathlineto{\pgfqpoint{3.602756in}{1.468964in}}%
\pgfpathclose%
\pgfusepath{stroke,fill}%
\end{pgfscope}%
\begin{pgfscope}%
\pgfpathrectangle{\pgfqpoint{0.050000in}{0.050000in}}{\pgfqpoint{4.900000in}{4.900000in}}%
\pgfusepath{clip}%
\pgfsetbuttcap%
\pgfsetroundjoin%
\definecolor{currentfill}{rgb}{0.878570,0.878570,0.878570}%
\pgfsetfillcolor{currentfill}%
\pgfsetfillopacity{0.900000}%
\pgfsetlinewidth{0.501875pt}%
\definecolor{currentstroke}{rgb}{0.200000,0.200000,0.200000}%
\pgfsetstrokecolor{currentstroke}%
\pgfsetdash{}{0pt}%
\pgfpathmoveto{\pgfqpoint{3.021137in}{1.821404in}}%
\pgfpathlineto{\pgfqpoint{3.066298in}{1.784402in}}%
\pgfpathlineto{\pgfqpoint{3.099408in}{1.763864in}}%
\pgfpathlineto{\pgfqpoint{3.054240in}{1.800668in}}%
\pgfpathlineto{\pgfqpoint{3.021137in}{1.821404in}}%
\pgfpathclose%
\pgfusepath{stroke,fill}%
\end{pgfscope}%
\begin{pgfscope}%
\pgfpathrectangle{\pgfqpoint{0.050000in}{0.050000in}}{\pgfqpoint{4.900000in}{4.900000in}}%
\pgfusepath{clip}%
\pgfsetbuttcap%
\pgfsetroundjoin%
\definecolor{currentfill}{rgb}{0.970473,0.970473,0.970473}%
\pgfsetfillcolor{currentfill}%
\pgfsetfillopacity{0.900000}%
\pgfsetlinewidth{0.501875pt}%
\definecolor{currentstroke}{rgb}{0.200000,0.200000,0.200000}%
\pgfsetstrokecolor{currentstroke}%
\pgfsetdash{}{0pt}%
\pgfpathmoveto{\pgfqpoint{3.681243in}{1.436251in}}%
\pgfpathlineto{\pgfqpoint{3.725622in}{1.427069in}}%
\pgfpathlineto{\pgfqpoint{3.759805in}{1.407015in}}%
\pgfpathlineto{\pgfqpoint{3.715408in}{1.416163in}}%
\pgfpathlineto{\pgfqpoint{3.681243in}{1.436251in}}%
\pgfpathclose%
\pgfusepath{stroke,fill}%
\end{pgfscope}%
\begin{pgfscope}%
\pgfpathrectangle{\pgfqpoint{0.050000in}{0.050000in}}{\pgfqpoint{4.900000in}{4.900000in}}%
\pgfusepath{clip}%
\pgfsetbuttcap%
\pgfsetroundjoin%
\definecolor{currentfill}{rgb}{0.959400,0.959400,0.959400}%
\pgfsetfillcolor{currentfill}%
\pgfsetfillopacity{0.900000}%
\pgfsetlinewidth{0.501875pt}%
\definecolor{currentstroke}{rgb}{0.200000,0.200000,0.200000}%
\pgfsetstrokecolor{currentstroke}%
\pgfsetdash{}{0pt}%
\pgfpathmoveto{\pgfqpoint{3.524331in}{1.505360in}}%
\pgfpathlineto{\pgfqpoint{3.568826in}{1.488956in}}%
\pgfpathlineto{\pgfqpoint{3.602756in}{1.468964in}}%
\pgfpathlineto{\pgfqpoint{3.558244in}{1.485266in}}%
\pgfpathlineto{\pgfqpoint{3.524331in}{1.505360in}}%
\pgfpathclose%
\pgfusepath{stroke,fill}%
\end{pgfscope}%
\begin{pgfscope}%
\pgfpathrectangle{\pgfqpoint{0.050000in}{0.050000in}}{\pgfqpoint{4.900000in}{4.900000in}}%
\pgfusepath{clip}%
\pgfsetbuttcap%
\pgfsetroundjoin%
\definecolor{currentfill}{rgb}{0.734579,0.734579,0.734579}%
\pgfsetfillcolor{currentfill}%
\pgfsetfillopacity{0.900000}%
\pgfsetlinewidth{0.501875pt}%
\definecolor{currentstroke}{rgb}{0.200000,0.200000,0.200000}%
\pgfsetstrokecolor{currentstroke}%
\pgfsetdash{}{0pt}%
\pgfpathmoveto{\pgfqpoint{2.485426in}{2.242179in}}%
\pgfpathlineto{\pgfqpoint{2.531680in}{2.197838in}}%
\pgfpathlineto{\pgfqpoint{2.563942in}{2.177028in}}%
\pgfpathlineto{\pgfqpoint{2.517686in}{2.221363in}}%
\pgfpathlineto{\pgfqpoint{2.485426in}{2.242179in}}%
\pgfpathclose%
\pgfusepath{stroke,fill}%
\end{pgfscope}%
\begin{pgfscope}%
\pgfpathrectangle{\pgfqpoint{0.050000in}{0.050000in}}{\pgfqpoint{4.900000in}{4.900000in}}%
\pgfusepath{clip}%
\pgfsetbuttcap%
\pgfsetroundjoin%
\definecolor{currentfill}{rgb}{0.063729,0.063729,0.063729}%
\pgfsetfillcolor{currentfill}%
\pgfsetfillopacity{0.900000}%
\pgfsetlinewidth{0.501875pt}%
\definecolor{currentstroke}{rgb}{0.200000,0.200000,0.200000}%
\pgfsetstrokecolor{currentstroke}%
\pgfsetdash{}{0pt}%
\pgfpathmoveto{\pgfqpoint{0.767249in}{3.631340in}}%
\pgfpathlineto{\pgfqpoint{0.817181in}{3.583829in}}%
\pgfpathlineto{\pgfqpoint{0.846761in}{3.565120in}}%
\pgfpathlineto{\pgfqpoint{0.796813in}{3.612665in}}%
\pgfpathlineto{\pgfqpoint{0.767249in}{3.631340in}}%
\pgfpathclose%
\pgfusepath{stroke,fill}%
\end{pgfscope}%
\begin{pgfscope}%
\pgfpathrectangle{\pgfqpoint{0.050000in}{0.050000in}}{\pgfqpoint{4.900000in}{4.900000in}}%
\pgfusepath{clip}%
\pgfsetbuttcap%
\pgfsetroundjoin%
\definecolor{currentfill}{rgb}{0.976009,0.976009,0.976009}%
\pgfsetfillcolor{currentfill}%
\pgfsetfillopacity{0.900000}%
\pgfsetlinewidth{0.501875pt}%
\definecolor{currentstroke}{rgb}{0.200000,0.200000,0.200000}%
\pgfsetstrokecolor{currentstroke}%
\pgfsetdash{}{0pt}%
\pgfpathmoveto{\pgfqpoint{3.759805in}{1.407015in}}%
\pgfpathlineto{\pgfqpoint{3.804138in}{1.401126in}}%
\pgfpathlineto{\pgfqpoint{3.838451in}{1.380995in}}%
\pgfpathlineto{\pgfqpoint{3.794098in}{1.386881in}}%
\pgfpathlineto{\pgfqpoint{3.759805in}{1.407015in}}%
\pgfpathclose%
\pgfusepath{stroke,fill}%
\end{pgfscope}%
\begin{pgfscope}%
\pgfpathrectangle{\pgfqpoint{0.050000in}{0.050000in}}{\pgfqpoint{4.900000in}{4.900000in}}%
\pgfusepath{clip}%
\pgfsetbuttcap%
\pgfsetroundjoin%
\definecolor{currentfill}{rgb}{0.952018,0.952018,0.952018}%
\pgfsetfillcolor{currentfill}%
\pgfsetfillopacity{0.900000}%
\pgfsetlinewidth{0.501875pt}%
\definecolor{currentstroke}{rgb}{0.200000,0.200000,0.200000}%
\pgfsetstrokecolor{currentstroke}%
\pgfsetdash{}{0pt}%
\pgfpathmoveto{\pgfqpoint{3.445957in}{1.545563in}}%
\pgfpathlineto{\pgfqpoint{3.490527in}{1.525376in}}%
\pgfpathlineto{\pgfqpoint{3.524331in}{1.505360in}}%
\pgfpathlineto{\pgfqpoint{3.479747in}{1.525410in}}%
\pgfpathlineto{\pgfqpoint{3.445957in}{1.545563in}}%
\pgfpathclose%
\pgfusepath{stroke,fill}%
\end{pgfscope}%
\begin{pgfscope}%
\pgfpathrectangle{\pgfqpoint{0.050000in}{0.050000in}}{\pgfqpoint{4.900000in}{4.900000in}}%
\pgfusepath{clip}%
\pgfsetbuttcap%
\pgfsetroundjoin%
\definecolor{currentfill}{rgb}{0.367243,0.367243,0.367243}%
\pgfsetfillcolor{currentfill}%
\pgfsetfillopacity{0.900000}%
\pgfsetlinewidth{0.501875pt}%
\definecolor{currentstroke}{rgb}{0.200000,0.200000,0.200000}%
\pgfsetstrokecolor{currentstroke}%
\pgfsetdash{}{0pt}%
\pgfpathmoveto{\pgfqpoint{1.492320in}{3.045306in}}%
\pgfpathlineto{\pgfqpoint{1.540704in}{2.999046in}}%
\pgfpathlineto{\pgfqpoint{1.571418in}{2.979423in}}%
\pgfpathlineto{\pgfqpoint{1.523025in}{3.025713in}}%
\pgfpathlineto{\pgfqpoint{1.492320in}{3.045306in}}%
\pgfpathclose%
\pgfusepath{stroke,fill}%
\end{pgfscope}%
\begin{pgfscope}%
\pgfpathrectangle{\pgfqpoint{0.050000in}{0.050000in}}{\pgfqpoint{4.900000in}{4.900000in}}%
\pgfusepath{clip}%
\pgfsetbuttcap%
\pgfsetroundjoin%
\definecolor{currentfill}{rgb}{0.624221,0.624221,0.624221}%
\pgfsetfillcolor{currentfill}%
\pgfsetfillopacity{0.900000}%
\pgfsetlinewidth{0.501875pt}%
\definecolor{currentstroke}{rgb}{0.200000,0.200000,0.200000}%
\pgfsetstrokecolor{currentstroke}%
\pgfsetdash{}{0pt}%
\pgfpathmoveto{\pgfqpoint{2.217490in}{2.459215in}}%
\pgfpathlineto{\pgfqpoint{2.264324in}{2.414213in}}%
\pgfpathlineto{\pgfqpoint{2.296174in}{2.393678in}}%
\pgfpathlineto{\pgfqpoint{2.249334in}{2.438704in}}%
\pgfpathlineto{\pgfqpoint{2.217490in}{2.459215in}}%
\pgfpathclose%
\pgfusepath{stroke,fill}%
\end{pgfscope}%
\begin{pgfscope}%
\pgfpathrectangle{\pgfqpoint{0.050000in}{0.050000in}}{\pgfqpoint{4.900000in}{4.900000in}}%
\pgfusepath{clip}%
\pgfsetbuttcap%
\pgfsetroundjoin%
\definecolor{currentfill}{rgb}{0.979700,0.979700,0.979700}%
\pgfsetfillcolor{currentfill}%
\pgfsetfillopacity{0.900000}%
\pgfsetlinewidth{0.501875pt}%
\definecolor{currentstroke}{rgb}{0.200000,0.200000,0.200000}%
\pgfsetstrokecolor{currentstroke}%
\pgfsetdash{}{0pt}%
\pgfpathmoveto{\pgfqpoint{3.838451in}{1.380995in}}%
\pgfpathlineto{\pgfqpoint{3.882744in}{1.378119in}}%
\pgfpathlineto{\pgfqpoint{3.917187in}{1.357888in}}%
\pgfpathlineto{\pgfqpoint{3.872874in}{1.360784in}}%
\pgfpathlineto{\pgfqpoint{3.838451in}{1.380995in}}%
\pgfpathclose%
\pgfusepath{stroke,fill}%
\end{pgfscope}%
\begin{pgfscope}%
\pgfpathrectangle{\pgfqpoint{0.050000in}{0.050000in}}{\pgfqpoint{4.900000in}{4.900000in}}%
\pgfusepath{clip}%
\pgfsetbuttcap%
\pgfsetroundjoin%
\definecolor{currentfill}{rgb}{0.942791,0.942791,0.942791}%
\pgfsetfillcolor{currentfill}%
\pgfsetfillopacity{0.900000}%
\pgfsetlinewidth{0.501875pt}%
\definecolor{currentstroke}{rgb}{0.200000,0.200000,0.200000}%
\pgfsetstrokecolor{currentstroke}%
\pgfsetdash{}{0pt}%
\pgfpathmoveto{\pgfqpoint{3.367616in}{1.589599in}}%
\pgfpathlineto{\pgfqpoint{3.412275in}{1.565640in}}%
\pgfpathlineto{\pgfqpoint{3.445957in}{1.545563in}}%
\pgfpathlineto{\pgfqpoint{3.401285in}{1.569352in}}%
\pgfpathlineto{\pgfqpoint{3.367616in}{1.589599in}}%
\pgfpathclose%
\pgfusepath{stroke,fill}%
\end{pgfscope}%
\begin{pgfscope}%
\pgfpathrectangle{\pgfqpoint{0.050000in}{0.050000in}}{\pgfqpoint{4.900000in}{4.900000in}}%
\pgfusepath{clip}%
\pgfsetbuttcap%
\pgfsetroundjoin%
\definecolor{currentfill}{rgb}{0.256517,0.256517,0.256517}%
\pgfsetfillcolor{currentfill}%
\pgfsetfillopacity{0.900000}%
\pgfsetlinewidth{0.501875pt}%
\definecolor{currentstroke}{rgb}{0.200000,0.200000,0.200000}%
\pgfsetstrokecolor{currentstroke}%
\pgfsetdash{}{0pt}%
\pgfpathmoveto{\pgfqpoint{1.224186in}{3.262673in}}%
\pgfpathlineto{\pgfqpoint{1.273153in}{3.215940in}}%
\pgfpathlineto{\pgfqpoint{1.303453in}{3.196652in}}%
\pgfpathlineto{\pgfqpoint{1.254475in}{3.243417in}}%
\pgfpathlineto{\pgfqpoint{1.224186in}{3.262673in}}%
\pgfpathclose%
\pgfusepath{stroke,fill}%
\end{pgfscope}%
\begin{pgfscope}%
\pgfpathrectangle{\pgfqpoint{0.050000in}{0.050000in}}{\pgfqpoint{4.900000in}{4.900000in}}%
\pgfusepath{clip}%
\pgfsetbuttcap%
\pgfsetroundjoin%
\definecolor{currentfill}{rgb}{0.525798,0.525798,0.525798}%
\pgfsetfillcolor{currentfill}%
\pgfsetfillopacity{0.900000}%
\pgfsetlinewidth{0.501875pt}%
\definecolor{currentstroke}{rgb}{0.200000,0.200000,0.200000}%
\pgfsetstrokecolor{currentstroke}%
\pgfsetdash{}{0pt}%
\pgfpathmoveto{\pgfqpoint{1.949461in}{2.676487in}}%
\pgfpathlineto{\pgfqpoint{1.996878in}{2.631006in}}%
\pgfpathlineto{\pgfqpoint{2.028314in}{2.610802in}}%
\pgfpathlineto{\pgfqpoint{1.980890in}{2.656312in}}%
\pgfpathlineto{\pgfqpoint{1.949461in}{2.676487in}}%
\pgfpathclose%
\pgfusepath{stroke,fill}%
\end{pgfscope}%
\begin{pgfscope}%
\pgfpathrectangle{\pgfqpoint{0.050000in}{0.050000in}}{\pgfqpoint{4.900000in}{4.900000in}}%
\pgfusepath{clip}%
\pgfsetbuttcap%
\pgfsetroundjoin%
\definecolor{currentfill}{rgb}{0.983391,0.983391,0.983391}%
\pgfsetfillcolor{currentfill}%
\pgfsetfillopacity{0.900000}%
\pgfsetlinewidth{0.501875pt}%
\definecolor{currentstroke}{rgb}{0.200000,0.200000,0.200000}%
\pgfsetstrokecolor{currentstroke}%
\pgfsetdash{}{0pt}%
\pgfpathmoveto{\pgfqpoint{3.917187in}{1.357888in}}%
\pgfpathlineto{\pgfqpoint{3.961442in}{1.357721in}}%
\pgfpathlineto{\pgfqpoint{3.996017in}{1.337372in}}%
\pgfpathlineto{\pgfqpoint{3.951740in}{1.337579in}}%
\pgfpathlineto{\pgfqpoint{3.917187in}{1.357888in}}%
\pgfpathclose%
\pgfusepath{stroke,fill}%
\end{pgfscope}%
\begin{pgfscope}%
\pgfpathrectangle{\pgfqpoint{0.050000in}{0.050000in}}{\pgfqpoint{4.900000in}{4.900000in}}%
\pgfusepath{clip}%
\pgfsetbuttcap%
\pgfsetroundjoin%
\definecolor{currentfill}{rgb}{0.858762,0.858762,0.858762}%
\pgfsetfillcolor{currentfill}%
\pgfsetfillopacity{0.900000}%
\pgfsetlinewidth{0.501875pt}%
\definecolor{currentstroke}{rgb}{0.200000,0.200000,0.200000}%
\pgfsetstrokecolor{currentstroke}%
\pgfsetdash{}{0pt}%
\pgfpathmoveto{\pgfqpoint{2.942813in}{1.881301in}}%
\pgfpathlineto{\pgfqpoint{2.988139in}{1.842086in}}%
\pgfpathlineto{\pgfqpoint{3.021137in}{1.821404in}}%
\pgfpathlineto{\pgfqpoint{2.975806in}{1.860446in}}%
\pgfpathlineto{\pgfqpoint{2.942813in}{1.881301in}}%
\pgfpathclose%
\pgfusepath{stroke,fill}%
\end{pgfscope}%
\begin{pgfscope}%
\pgfpathrectangle{\pgfqpoint{0.050000in}{0.050000in}}{\pgfqpoint{4.900000in}{4.900000in}}%
\pgfusepath{clip}%
\pgfsetbuttcap%
\pgfsetroundjoin%
\definecolor{currentfill}{rgb}{0.788112,0.788112,0.788112}%
\pgfsetfillcolor{currentfill}%
\pgfsetfillopacity{0.900000}%
\pgfsetlinewidth{0.501875pt}%
\definecolor{currentstroke}{rgb}{0.200000,0.200000,0.200000}%
\pgfsetstrokecolor{currentstroke}%
\pgfsetdash{}{0pt}%
\pgfpathmoveto{\pgfqpoint{2.674837in}{2.091352in}}%
\pgfpathlineto{\pgfqpoint{2.720707in}{2.048065in}}%
\pgfpathlineto{\pgfqpoint{2.753280in}{2.027178in}}%
\pgfpathlineto{\pgfqpoint{2.707407in}{2.070402in}}%
\pgfpathlineto{\pgfqpoint{2.674837in}{2.091352in}}%
\pgfpathclose%
\pgfusepath{stroke,fill}%
\end{pgfscope}%
\begin{pgfscope}%
\pgfpathrectangle{\pgfqpoint{0.050000in}{0.050000in}}{\pgfqpoint{4.900000in}{4.900000in}}%
\pgfusepath{clip}%
\pgfsetbuttcap%
\pgfsetroundjoin%
\definecolor{currentfill}{rgb}{0.136563,0.136563,0.136563}%
\pgfsetfillcolor{currentfill}%
\pgfsetfillopacity{0.900000}%
\pgfsetlinewidth{0.501875pt}%
\definecolor{currentstroke}{rgb}{0.200000,0.200000,0.200000}%
\pgfsetstrokecolor{currentstroke}%
\pgfsetdash{}{0pt}%
\pgfpathmoveto{\pgfqpoint{0.955961in}{3.480117in}}%
\pgfpathlineto{\pgfqpoint{1.005510in}{3.432910in}}%
\pgfpathlineto{\pgfqpoint{1.035397in}{3.413958in}}%
\pgfpathlineto{\pgfqpoint{0.985834in}{3.461198in}}%
\pgfpathlineto{\pgfqpoint{0.955961in}{3.480117in}}%
\pgfpathclose%
\pgfusepath{stroke,fill}%
\end{pgfscope}%
\begin{pgfscope}%
\pgfpathrectangle{\pgfqpoint{0.050000in}{0.050000in}}{\pgfqpoint{4.900000in}{4.900000in}}%
\pgfusepath{clip}%
\pgfsetbuttcap%
\pgfsetroundjoin%
\definecolor{currentfill}{rgb}{0.929504,0.929504,0.929504}%
\pgfsetfillcolor{currentfill}%
\pgfsetfillopacity{0.900000}%
\pgfsetlinewidth{0.501875pt}%
\definecolor{currentstroke}{rgb}{0.200000,0.200000,0.200000}%
\pgfsetstrokecolor{currentstroke}%
\pgfsetdash{}{0pt}%
\pgfpathmoveto{\pgfqpoint{3.289292in}{1.637380in}}%
\pgfpathlineto{\pgfqpoint{3.334055in}{1.609773in}}%
\pgfpathlineto{\pgfqpoint{3.367616in}{1.589599in}}%
\pgfpathlineto{\pgfqpoint{3.322843in}{1.617013in}}%
\pgfpathlineto{\pgfqpoint{3.289292in}{1.637380in}}%
\pgfpathclose%
\pgfusepath{stroke,fill}%
\end{pgfscope}%
\begin{pgfscope}%
\pgfpathrectangle{\pgfqpoint{0.050000in}{0.050000in}}{\pgfqpoint{4.900000in}{4.900000in}}%
\pgfusepath{clip}%
\pgfsetbuttcap%
\pgfsetroundjoin%
\definecolor{currentfill}{rgb}{0.428143,0.428143,0.428143}%
\pgfsetfillcolor{currentfill}%
\pgfsetfillopacity{0.900000}%
\pgfsetlinewidth{0.501875pt}%
\definecolor{currentstroke}{rgb}{0.200000,0.200000,0.200000}%
\pgfsetstrokecolor{currentstroke}%
\pgfsetdash{}{0pt}%
\pgfpathmoveto{\pgfqpoint{1.681342in}{2.893839in}}%
\pgfpathlineto{\pgfqpoint{1.729341in}{2.847884in}}%
\pgfpathlineto{\pgfqpoint{1.760363in}{2.828016in}}%
\pgfpathlineto{\pgfqpoint{1.712355in}{2.874000in}}%
\pgfpathlineto{\pgfqpoint{1.681342in}{2.893839in}}%
\pgfpathclose%
\pgfusepath{stroke,fill}%
\end{pgfscope}%
\begin{pgfscope}%
\pgfpathrectangle{\pgfqpoint{0.050000in}{0.050000in}}{\pgfqpoint{4.900000in}{4.900000in}}%
\pgfusepath{clip}%
\pgfsetbuttcap%
\pgfsetroundjoin%
\definecolor{currentfill}{rgb}{0.696194,0.696194,0.696194}%
\pgfsetfillcolor{currentfill}%
\pgfsetfillopacity{0.900000}%
\pgfsetlinewidth{0.501875pt}%
\definecolor{currentstroke}{rgb}{0.200000,0.200000,0.200000}%
\pgfsetstrokecolor{currentstroke}%
\pgfsetdash{}{0pt}%
\pgfpathmoveto{\pgfqpoint{2.406820in}{2.307557in}}%
\pgfpathlineto{\pgfqpoint{2.453268in}{2.262933in}}%
\pgfpathlineto{\pgfqpoint{2.485426in}{2.242179in}}%
\pgfpathlineto{\pgfqpoint{2.438974in}{2.286812in}}%
\pgfpathlineto{\pgfqpoint{2.406820in}{2.307557in}}%
\pgfpathclose%
\pgfusepath{stroke,fill}%
\end{pgfscope}%
\begin{pgfscope}%
\pgfpathrectangle{\pgfqpoint{0.050000in}{0.050000in}}{\pgfqpoint{4.900000in}{4.900000in}}%
\pgfusepath{clip}%
\pgfsetbuttcap%
\pgfsetroundjoin%
\definecolor{currentfill}{rgb}{0.031865,0.031865,0.031865}%
\pgfsetfillcolor{currentfill}%
\pgfsetfillopacity{0.900000}%
\pgfsetlinewidth{0.501875pt}%
\definecolor{currentstroke}{rgb}{0.200000,0.200000,0.200000}%
\pgfsetstrokecolor{currentstroke}%
\pgfsetdash{}{0pt}%
\pgfpathmoveto{\pgfqpoint{0.687644in}{3.697637in}}%
\pgfpathlineto{\pgfqpoint{0.737776in}{3.649956in}}%
\pgfpathlineto{\pgfqpoint{0.767249in}{3.631340in}}%
\pgfpathlineto{\pgfqpoint{0.717101in}{3.679055in}}%
\pgfpathlineto{\pgfqpoint{0.687644in}{3.697637in}}%
\pgfpathclose%
\pgfusepath{stroke,fill}%
\end{pgfscope}%
\begin{pgfscope}%
\pgfpathrectangle{\pgfqpoint{0.050000in}{0.050000in}}{\pgfqpoint{4.900000in}{4.900000in}}%
\pgfusepath{clip}%
\pgfsetbuttcap%
\pgfsetroundjoin%
\definecolor{currentfill}{rgb}{0.987082,0.987082,0.987082}%
\pgfsetfillcolor{currentfill}%
\pgfsetfillopacity{0.900000}%
\pgfsetlinewidth{0.501875pt}%
\definecolor{currentstroke}{rgb}{0.200000,0.200000,0.200000}%
\pgfsetstrokecolor{currentstroke}%
\pgfsetdash{}{0pt}%
\pgfpathmoveto{\pgfqpoint{3.996017in}{1.337372in}}%
\pgfpathlineto{\pgfqpoint{4.040236in}{1.339598in}}%
\pgfpathlineto{\pgfqpoint{4.074943in}{1.319119in}}%
\pgfpathlineto{\pgfqpoint{4.030702in}{1.316947in}}%
\pgfpathlineto{\pgfqpoint{3.996017in}{1.337372in}}%
\pgfpathclose%
\pgfusepath{stroke,fill}%
\end{pgfscope}%
\begin{pgfscope}%
\pgfpathrectangle{\pgfqpoint{0.050000in}{0.050000in}}{\pgfqpoint{4.900000in}{4.900000in}}%
\pgfusepath{clip}%
\pgfsetbuttcap%
\pgfsetroundjoin%
\definecolor{currentfill}{rgb}{0.334764,0.334764,0.334764}%
\pgfsetfillcolor{currentfill}%
\pgfsetfillopacity{0.900000}%
\pgfsetlinewidth{0.501875pt}%
\definecolor{currentstroke}{rgb}{0.200000,0.200000,0.200000}%
\pgfsetstrokecolor{currentstroke}%
\pgfsetdash{}{0pt}%
\pgfpathmoveto{\pgfqpoint{1.413131in}{3.111267in}}%
\pgfpathlineto{\pgfqpoint{1.461712in}{3.064839in}}%
\pgfpathlineto{\pgfqpoint{1.492320in}{3.045306in}}%
\pgfpathlineto{\pgfqpoint{1.443728in}{3.091765in}}%
\pgfpathlineto{\pgfqpoint{1.413131in}{3.111267in}}%
\pgfpathclose%
\pgfusepath{stroke,fill}%
\end{pgfscope}%
\begin{pgfscope}%
\pgfpathrectangle{\pgfqpoint{0.050000in}{0.050000in}}{\pgfqpoint{4.900000in}{4.900000in}}%
\pgfusepath{clip}%
\pgfsetbuttcap%
\pgfsetroundjoin%
\definecolor{currentfill}{rgb}{0.590634,0.590634,0.590634}%
\pgfsetfillcolor{currentfill}%
\pgfsetfillopacity{0.900000}%
\pgfsetlinewidth{0.501875pt}%
\definecolor{currentstroke}{rgb}{0.200000,0.200000,0.200000}%
\pgfsetstrokecolor{currentstroke}%
\pgfsetdash{}{0pt}%
\pgfpathmoveto{\pgfqpoint{2.138715in}{2.524836in}}%
\pgfpathlineto{\pgfqpoint{2.185745in}{2.479663in}}%
\pgfpathlineto{\pgfqpoint{2.217490in}{2.459215in}}%
\pgfpathlineto{\pgfqpoint{2.170454in}{2.504415in}}%
\pgfpathlineto{\pgfqpoint{2.138715in}{2.524836in}}%
\pgfpathclose%
\pgfusepath{stroke,fill}%
\end{pgfscope}%
\begin{pgfscope}%
\pgfpathrectangle{\pgfqpoint{0.050000in}{0.050000in}}{\pgfqpoint{4.900000in}{4.900000in}}%
\pgfusepath{clip}%
\pgfsetbuttcap%
\pgfsetroundjoin%
\definecolor{currentfill}{rgb}{0.912526,0.912526,0.912526}%
\pgfsetfillcolor{currentfill}%
\pgfsetfillopacity{0.900000}%
\pgfsetlinewidth{0.501875pt}%
\definecolor{currentstroke}{rgb}{0.200000,0.200000,0.200000}%
\pgfsetstrokecolor{currentstroke}%
\pgfsetdash{}{0pt}%
\pgfpathmoveto{\pgfqpoint{3.210967in}{1.688707in}}%
\pgfpathlineto{\pgfqpoint{3.255849in}{1.657681in}}%
\pgfpathlineto{\pgfqpoint{3.289292in}{1.637380in}}%
\pgfpathlineto{\pgfqpoint{3.244401in}{1.668201in}}%
\pgfpathlineto{\pgfqpoint{3.210967in}{1.688707in}}%
\pgfpathclose%
\pgfusepath{stroke,fill}%
\end{pgfscope}%
\begin{pgfscope}%
\pgfpathrectangle{\pgfqpoint{0.050000in}{0.050000in}}{\pgfqpoint{4.900000in}{4.900000in}}%
\pgfusepath{clip}%
\pgfsetbuttcap%
\pgfsetroundjoin%
\definecolor{currentfill}{rgb}{0.217762,0.217762,0.217762}%
\pgfsetfillcolor{currentfill}%
\pgfsetfillopacity{0.900000}%
\pgfsetlinewidth{0.501875pt}%
\definecolor{currentstroke}{rgb}{0.200000,0.200000,0.200000}%
\pgfsetstrokecolor{currentstroke}%
\pgfsetdash{}{0pt}%
\pgfpathmoveto{\pgfqpoint{1.144828in}{3.328771in}}%
\pgfpathlineto{\pgfqpoint{1.193992in}{3.281869in}}%
\pgfpathlineto{\pgfqpoint{1.224186in}{3.262673in}}%
\pgfpathlineto{\pgfqpoint{1.175010in}{3.309607in}}%
\pgfpathlineto{\pgfqpoint{1.144828in}{3.328771in}}%
\pgfpathclose%
\pgfusepath{stroke,fill}%
\end{pgfscope}%
\begin{pgfscope}%
\pgfpathrectangle{\pgfqpoint{0.050000in}{0.050000in}}{\pgfqpoint{4.900000in}{4.900000in}}%
\pgfusepath{clip}%
\pgfsetbuttcap%
\pgfsetroundjoin%
\definecolor{currentfill}{rgb}{0.491349,0.491349,0.491349}%
\pgfsetfillcolor{currentfill}%
\pgfsetfillopacity{0.900000}%
\pgfsetlinewidth{0.501875pt}%
\definecolor{currentstroke}{rgb}{0.200000,0.200000,0.200000}%
\pgfsetstrokecolor{currentstroke}%
\pgfsetdash{}{0pt}%
\pgfpathmoveto{\pgfqpoint{1.870518in}{2.742248in}}%
\pgfpathlineto{\pgfqpoint{1.918131in}{2.696600in}}%
\pgfpathlineto{\pgfqpoint{1.949461in}{2.676487in}}%
\pgfpathlineto{\pgfqpoint{1.901841in}{2.722164in}}%
\pgfpathlineto{\pgfqpoint{1.870518in}{2.742248in}}%
\pgfpathclose%
\pgfusepath{stroke,fill}%
\end{pgfscope}%
\begin{pgfscope}%
\pgfpathrectangle{\pgfqpoint{0.050000in}{0.050000in}}{\pgfqpoint{4.900000in}{4.900000in}}%
\pgfusepath{clip}%
\pgfsetbuttcap%
\pgfsetroundjoin%
\definecolor{currentfill}{rgb}{0.836340,0.836340,0.836340}%
\pgfsetfillcolor{currentfill}%
\pgfsetfillopacity{0.900000}%
\pgfsetlinewidth{0.501875pt}%
\definecolor{currentstroke}{rgb}{0.200000,0.200000,0.200000}%
\pgfsetstrokecolor{currentstroke}%
\pgfsetdash{}{0pt}%
\pgfpathmoveto{\pgfqpoint{2.864422in}{1.943068in}}%
\pgfpathlineto{\pgfqpoint{2.909923in}{1.902105in}}%
\pgfpathlineto{\pgfqpoint{2.942813in}{1.881301in}}%
\pgfpathlineto{\pgfqpoint{2.897308in}{1.922124in}}%
\pgfpathlineto{\pgfqpoint{2.864422in}{1.943068in}}%
\pgfpathclose%
\pgfusepath{stroke,fill}%
\end{pgfscope}%
\begin{pgfscope}%
\pgfpathrectangle{\pgfqpoint{0.050000in}{0.050000in}}{\pgfqpoint{4.900000in}{4.900000in}}%
\pgfusepath{clip}%
\pgfsetbuttcap%
\pgfsetroundjoin%
\definecolor{currentfill}{rgb}{0.763998,0.763998,0.763998}%
\pgfsetfillcolor{currentfill}%
\pgfsetfillopacity{0.900000}%
\pgfsetlinewidth{0.501875pt}%
\definecolor{currentstroke}{rgb}{0.200000,0.200000,0.200000}%
\pgfsetstrokecolor{currentstroke}%
\pgfsetdash{}{0pt}%
\pgfpathmoveto{\pgfqpoint{2.596306in}{2.156159in}}%
\pgfpathlineto{\pgfqpoint{2.642368in}{2.112248in}}%
\pgfpathlineto{\pgfqpoint{2.674837in}{2.091352in}}%
\pgfpathlineto{\pgfqpoint{2.628771in}{2.135232in}}%
\pgfpathlineto{\pgfqpoint{2.596306in}{2.156159in}}%
\pgfpathclose%
\pgfusepath{stroke,fill}%
\end{pgfscope}%
\begin{pgfscope}%
\pgfpathrectangle{\pgfqpoint{0.050000in}{0.050000in}}{\pgfqpoint{4.900000in}{4.900000in}}%
\pgfusepath{clip}%
\pgfsetbuttcap%
\pgfsetroundjoin%
\definecolor{currentfill}{rgb}{0.100146,0.100146,0.100146}%
\pgfsetfillcolor{currentfill}%
\pgfsetfillopacity{0.900000}%
\pgfsetlinewidth{0.501875pt}%
\definecolor{currentstroke}{rgb}{0.200000,0.200000,0.200000}%
\pgfsetstrokecolor{currentstroke}%
\pgfsetdash{}{0pt}%
\pgfpathmoveto{\pgfqpoint{0.876433in}{3.546353in}}%
\pgfpathlineto{\pgfqpoint{0.926181in}{3.498977in}}%
\pgfpathlineto{\pgfqpoint{0.955961in}{3.480117in}}%
\pgfpathlineto{\pgfqpoint{0.906199in}{3.527526in}}%
\pgfpathlineto{\pgfqpoint{0.876433in}{3.546353in}}%
\pgfpathclose%
\pgfusepath{stroke,fill}%
\end{pgfscope}%
\begin{pgfscope}%
\pgfpathrectangle{\pgfqpoint{0.050000in}{0.050000in}}{\pgfqpoint{4.900000in}{4.900000in}}%
\pgfusepath{clip}%
\pgfsetbuttcap%
\pgfsetroundjoin%
\definecolor{currentfill}{rgb}{0.399723,0.399723,0.399723}%
\pgfsetfillcolor{currentfill}%
\pgfsetfillopacity{0.900000}%
\pgfsetlinewidth{0.501875pt}%
\definecolor{currentstroke}{rgb}{0.200000,0.200000,0.200000}%
\pgfsetstrokecolor{currentstroke}%
\pgfsetdash{}{0pt}%
\pgfpathmoveto{\pgfqpoint{1.602229in}{2.959737in}}%
\pgfpathlineto{\pgfqpoint{1.650425in}{2.913615in}}%
\pgfpathlineto{\pgfqpoint{1.681342in}{2.893839in}}%
\pgfpathlineto{\pgfqpoint{1.633136in}{2.939991in}}%
\pgfpathlineto{\pgfqpoint{1.602229in}{2.959737in}}%
\pgfpathclose%
\pgfusepath{stroke,fill}%
\end{pgfscope}%
\begin{pgfscope}%
\pgfpathrectangle{\pgfqpoint{0.050000in}{0.050000in}}{\pgfqpoint{4.900000in}{4.900000in}}%
\pgfusepath{clip}%
\pgfsetbuttcap%
\pgfsetroundjoin%
\definecolor{currentfill}{rgb}{0.895548,0.895548,0.895548}%
\pgfsetfillcolor{currentfill}%
\pgfsetfillopacity{0.900000}%
\pgfsetlinewidth{0.501875pt}%
\definecolor{currentstroke}{rgb}{0.200000,0.200000,0.200000}%
\pgfsetstrokecolor{currentstroke}%
\pgfsetdash{}{0pt}%
\pgfpathmoveto{\pgfqpoint{3.132622in}{1.743270in}}%
\pgfpathlineto{\pgfqpoint{3.177639in}{1.709152in}}%
\pgfpathlineto{\pgfqpoint{3.210967in}{1.688707in}}%
\pgfpathlineto{\pgfqpoint{3.165942in}{1.722618in}}%
\pgfpathlineto{\pgfqpoint{3.132622in}{1.743270in}}%
\pgfpathclose%
\pgfusepath{stroke,fill}%
\end{pgfscope}%
\begin{pgfscope}%
\pgfpathrectangle{\pgfqpoint{0.050000in}{0.050000in}}{\pgfqpoint{4.900000in}{4.900000in}}%
\pgfusepath{clip}%
\pgfsetbuttcap%
\pgfsetroundjoin%
\definecolor{currentfill}{rgb}{0.662607,0.662607,0.662607}%
\pgfsetfillcolor{currentfill}%
\pgfsetfillopacity{0.900000}%
\pgfsetlinewidth{0.501875pt}%
\definecolor{currentstroke}{rgb}{0.200000,0.200000,0.200000}%
\pgfsetstrokecolor{currentstroke}%
\pgfsetdash{}{0pt}%
\pgfpathmoveto{\pgfqpoint{2.328123in}{2.373079in}}%
\pgfpathlineto{\pgfqpoint{2.374766in}{2.328238in}}%
\pgfpathlineto{\pgfqpoint{2.406820in}{2.307557in}}%
\pgfpathlineto{\pgfqpoint{2.360172in}{2.352416in}}%
\pgfpathlineto{\pgfqpoint{2.328123in}{2.373079in}}%
\pgfpathclose%
\pgfusepath{stroke,fill}%
\end{pgfscope}%
\begin{pgfscope}%
\pgfpathrectangle{\pgfqpoint{0.050000in}{0.050000in}}{\pgfqpoint{4.900000in}{4.900000in}}%
\pgfusepath{clip}%
\pgfsetbuttcap%
\pgfsetroundjoin%
\definecolor{currentfill}{rgb}{0.300807,0.300807,0.300807}%
\pgfsetfillcolor{currentfill}%
\pgfsetfillopacity{0.900000}%
\pgfsetlinewidth{0.501875pt}%
\definecolor{currentstroke}{rgb}{0.200000,0.200000,0.200000}%
\pgfsetstrokecolor{currentstroke}%
\pgfsetdash{}{0pt}%
\pgfpathmoveto{\pgfqpoint{1.333849in}{3.177303in}}%
\pgfpathlineto{\pgfqpoint{1.382628in}{3.130707in}}%
\pgfpathlineto{\pgfqpoint{1.413131in}{3.111267in}}%
\pgfpathlineto{\pgfqpoint{1.364340in}{3.157894in}}%
\pgfpathlineto{\pgfqpoint{1.333849in}{3.177303in}}%
\pgfpathclose%
\pgfusepath{stroke,fill}%
\end{pgfscope}%
\begin{pgfscope}%
\pgfpathrectangle{\pgfqpoint{0.050000in}{0.050000in}}{\pgfqpoint{4.900000in}{4.900000in}}%
\pgfusepath{clip}%
\pgfsetbuttcap%
\pgfsetroundjoin%
\definecolor{currentfill}{rgb}{0.555940,0.555940,0.555940}%
\pgfsetfillcolor{currentfill}%
\pgfsetfillopacity{0.900000}%
\pgfsetlinewidth{0.501875pt}%
\definecolor{currentstroke}{rgb}{0.200000,0.200000,0.200000}%
\pgfsetstrokecolor{currentstroke}%
\pgfsetdash{}{0pt}%
\pgfpathmoveto{\pgfqpoint{2.059848in}{2.590536in}}%
\pgfpathlineto{\pgfqpoint{2.107075in}{2.545194in}}%
\pgfpathlineto{\pgfqpoint{2.138715in}{2.524836in}}%
\pgfpathlineto{\pgfqpoint{2.091482in}{2.570205in}}%
\pgfpathlineto{\pgfqpoint{2.059848in}{2.590536in}}%
\pgfpathclose%
\pgfusepath{stroke,fill}%
\end{pgfscope}%
\begin{pgfscope}%
\pgfpathrectangle{\pgfqpoint{0.050000in}{0.050000in}}{\pgfqpoint{4.900000in}{4.900000in}}%
\pgfusepath{clip}%
\pgfsetbuttcap%
\pgfsetroundjoin%
\definecolor{currentfill}{rgb}{0.173472,0.173472,0.173472}%
\pgfsetfillcolor{currentfill}%
\pgfsetfillopacity{0.900000}%
\pgfsetlinewidth{0.501875pt}%
\definecolor{currentstroke}{rgb}{0.200000,0.200000,0.200000}%
\pgfsetstrokecolor{currentstroke}%
\pgfsetdash{}{0pt}%
\pgfpathmoveto{\pgfqpoint{1.065377in}{3.394946in}}%
\pgfpathlineto{\pgfqpoint{1.114740in}{3.347875in}}%
\pgfpathlineto{\pgfqpoint{1.144828in}{3.328771in}}%
\pgfpathlineto{\pgfqpoint{1.095452in}{3.375875in}}%
\pgfpathlineto{\pgfqpoint{1.065377in}{3.394946in}}%
\pgfpathclose%
\pgfusepath{stroke,fill}%
\end{pgfscope}%
\begin{pgfscope}%
\pgfpathrectangle{\pgfqpoint{0.050000in}{0.050000in}}{\pgfqpoint{4.900000in}{4.900000in}}%
\pgfusepath{clip}%
\pgfsetbuttcap%
\pgfsetroundjoin%
\definecolor{currentfill}{rgb}{0.812226,0.812226,0.812226}%
\pgfsetfillcolor{currentfill}%
\pgfsetfillopacity{0.900000}%
\pgfsetlinewidth{0.501875pt}%
\definecolor{currentstroke}{rgb}{0.200000,0.200000,0.200000}%
\pgfsetstrokecolor{currentstroke}%
\pgfsetdash{}{0pt}%
\pgfpathmoveto{\pgfqpoint{2.785956in}{2.006237in}}%
\pgfpathlineto{\pgfqpoint{2.831640in}{1.963959in}}%
\pgfpathlineto{\pgfqpoint{2.864422in}{1.943068in}}%
\pgfpathlineto{\pgfqpoint{2.818735in}{1.985243in}}%
\pgfpathlineto{\pgfqpoint{2.785956in}{2.006237in}}%
\pgfpathclose%
\pgfusepath{stroke,fill}%
\end{pgfscope}%
\begin{pgfscope}%
\pgfpathrectangle{\pgfqpoint{0.050000in}{0.050000in}}{\pgfqpoint{4.900000in}{4.900000in}}%
\pgfusepath{clip}%
\pgfsetbuttcap%
\pgfsetroundjoin%
\definecolor{currentfill}{rgb}{0.461207,0.461207,0.461207}%
\pgfsetfillcolor{currentfill}%
\pgfsetfillopacity{0.900000}%
\pgfsetlinewidth{0.501875pt}%
\definecolor{currentstroke}{rgb}{0.200000,0.200000,0.200000}%
\pgfsetstrokecolor{currentstroke}%
\pgfsetdash{}{0pt}%
\pgfpathmoveto{\pgfqpoint{1.791482in}{2.808086in}}%
\pgfpathlineto{\pgfqpoint{1.839292in}{2.762270in}}%
\pgfpathlineto{\pgfqpoint{1.870518in}{2.742248in}}%
\pgfpathlineto{\pgfqpoint{1.822700in}{2.788093in}}%
\pgfpathlineto{\pgfqpoint{1.791482in}{2.808086in}}%
\pgfpathclose%
\pgfusepath{stroke,fill}%
\end{pgfscope}%
\begin{pgfscope}%
\pgfpathrectangle{\pgfqpoint{0.050000in}{0.050000in}}{\pgfqpoint{4.900000in}{4.900000in}}%
\pgfusepath{clip}%
\pgfsetbuttcap%
\pgfsetroundjoin%
\definecolor{currentfill}{rgb}{0.964937,0.964937,0.964937}%
\pgfsetfillcolor{currentfill}%
\pgfsetfillopacity{0.900000}%
\pgfsetlinewidth{0.501875pt}%
\definecolor{currentstroke}{rgb}{0.200000,0.200000,0.200000}%
\pgfsetstrokecolor{currentstroke}%
\pgfsetdash{}{0pt}%
\pgfpathmoveto{\pgfqpoint{3.636794in}{1.448892in}}%
\pgfpathlineto{\pgfqpoint{3.681243in}{1.436251in}}%
\pgfpathlineto{\pgfqpoint{3.715408in}{1.416163in}}%
\pgfpathlineto{\pgfqpoint{3.670940in}{1.428738in}}%
\pgfpathlineto{\pgfqpoint{3.636794in}{1.448892in}}%
\pgfpathclose%
\pgfusepath{stroke,fill}%
\end{pgfscope}%
\begin{pgfscope}%
\pgfpathrectangle{\pgfqpoint{0.050000in}{0.050000in}}{\pgfqpoint{4.900000in}{4.900000in}}%
\pgfusepath{clip}%
\pgfsetbuttcap%
\pgfsetroundjoin%
\definecolor{currentfill}{rgb}{0.878570,0.878570,0.878570}%
\pgfsetfillcolor{currentfill}%
\pgfsetfillopacity{0.900000}%
\pgfsetlinewidth{0.501875pt}%
\definecolor{currentstroke}{rgb}{0.200000,0.200000,0.200000}%
\pgfsetstrokecolor{currentstroke}%
\pgfsetdash{}{0pt}%
\pgfpathmoveto{\pgfqpoint{3.054240in}{1.800668in}}%
\pgfpathlineto{\pgfqpoint{3.099408in}{1.763864in}}%
\pgfpathlineto{\pgfqpoint{3.132622in}{1.743270in}}%
\pgfpathlineto{\pgfqpoint{3.087448in}{1.779878in}}%
\pgfpathlineto{\pgfqpoint{3.054240in}{1.800668in}}%
\pgfpathclose%
\pgfusepath{stroke,fill}%
\end{pgfscope}%
\begin{pgfscope}%
\pgfpathrectangle{\pgfqpoint{0.050000in}{0.050000in}}{\pgfqpoint{4.900000in}{4.900000in}}%
\pgfusepath{clip}%
\pgfsetbuttcap%
\pgfsetroundjoin%
\definecolor{currentfill}{rgb}{0.734579,0.734579,0.734579}%
\pgfsetfillcolor{currentfill}%
\pgfsetfillopacity{0.900000}%
\pgfsetlinewidth{0.501875pt}%
\definecolor{currentstroke}{rgb}{0.200000,0.200000,0.200000}%
\pgfsetstrokecolor{currentstroke}%
\pgfsetdash{}{0pt}%
\pgfpathmoveto{\pgfqpoint{2.517686in}{2.221363in}}%
\pgfpathlineto{\pgfqpoint{2.563942in}{2.177028in}}%
\pgfpathlineto{\pgfqpoint{2.596306in}{2.156159in}}%
\pgfpathlineto{\pgfqpoint{2.550046in}{2.200487in}}%
\pgfpathlineto{\pgfqpoint{2.517686in}{2.221363in}}%
\pgfpathclose%
\pgfusepath{stroke,fill}%
\end{pgfscope}%
\begin{pgfscope}%
\pgfpathrectangle{\pgfqpoint{0.050000in}{0.050000in}}{\pgfqpoint{4.900000in}{4.900000in}}%
\pgfusepath{clip}%
\pgfsetbuttcap%
\pgfsetroundjoin%
\definecolor{currentfill}{rgb}{0.959400,0.959400,0.959400}%
\pgfsetfillcolor{currentfill}%
\pgfsetfillopacity{0.900000}%
\pgfsetlinewidth{0.501875pt}%
\definecolor{currentstroke}{rgb}{0.200000,0.200000,0.200000}%
\pgfsetstrokecolor{currentstroke}%
\pgfsetdash{}{0pt}%
\pgfpathmoveto{\pgfqpoint{3.558244in}{1.485266in}}%
\pgfpathlineto{\pgfqpoint{3.602756in}{1.468964in}}%
\pgfpathlineto{\pgfqpoint{3.636794in}{1.448892in}}%
\pgfpathlineto{\pgfqpoint{3.592266in}{1.465092in}}%
\pgfpathlineto{\pgfqpoint{3.558244in}{1.485266in}}%
\pgfpathclose%
\pgfusepath{stroke,fill}%
\end{pgfscope}%
\begin{pgfscope}%
\pgfpathrectangle{\pgfqpoint{0.050000in}{0.050000in}}{\pgfqpoint{4.900000in}{4.900000in}}%
\pgfusepath{clip}%
\pgfsetbuttcap%
\pgfsetroundjoin%
\definecolor{currentfill}{rgb}{0.970473,0.970473,0.970473}%
\pgfsetfillcolor{currentfill}%
\pgfsetfillopacity{0.900000}%
\pgfsetlinewidth{0.501875pt}%
\definecolor{currentstroke}{rgb}{0.200000,0.200000,0.200000}%
\pgfsetstrokecolor{currentstroke}%
\pgfsetdash{}{0pt}%
\pgfpathmoveto{\pgfqpoint{3.715408in}{1.416163in}}%
\pgfpathlineto{\pgfqpoint{3.759805in}{1.407015in}}%
\pgfpathlineto{\pgfqpoint{3.794098in}{1.386881in}}%
\pgfpathlineto{\pgfqpoint{3.749682in}{1.395995in}}%
\pgfpathlineto{\pgfqpoint{3.715408in}{1.416163in}}%
\pgfpathclose%
\pgfusepath{stroke,fill}%
\end{pgfscope}%
\begin{pgfscope}%
\pgfpathrectangle{\pgfqpoint{0.050000in}{0.050000in}}{\pgfqpoint{4.900000in}{4.900000in}}%
\pgfusepath{clip}%
\pgfsetbuttcap%
\pgfsetroundjoin%
\definecolor{currentfill}{rgb}{0.068281,0.068281,0.068281}%
\pgfsetfillcolor{currentfill}%
\pgfsetfillopacity{0.900000}%
\pgfsetlinewidth{0.501875pt}%
\definecolor{currentstroke}{rgb}{0.200000,0.200000,0.200000}%
\pgfsetstrokecolor{currentstroke}%
\pgfsetdash{}{0pt}%
\pgfpathmoveto{\pgfqpoint{0.796813in}{3.612665in}}%
\pgfpathlineto{\pgfqpoint{0.846761in}{3.565120in}}%
\pgfpathlineto{\pgfqpoint{0.876433in}{3.546353in}}%
\pgfpathlineto{\pgfqpoint{0.826471in}{3.593932in}}%
\pgfpathlineto{\pgfqpoint{0.796813in}{3.612665in}}%
\pgfpathclose%
\pgfusepath{stroke,fill}%
\end{pgfscope}%
\begin{pgfscope}%
\pgfpathrectangle{\pgfqpoint{0.050000in}{0.050000in}}{\pgfqpoint{4.900000in}{4.900000in}}%
\pgfusepath{clip}%
\pgfsetbuttcap%
\pgfsetroundjoin%
\definecolor{currentfill}{rgb}{0.976009,0.976009,0.976009}%
\pgfsetfillcolor{currentfill}%
\pgfsetfillopacity{0.900000}%
\pgfsetlinewidth{0.501875pt}%
\definecolor{currentstroke}{rgb}{0.200000,0.200000,0.200000}%
\pgfsetstrokecolor{currentstroke}%
\pgfsetdash{}{0pt}%
\pgfpathmoveto{\pgfqpoint{3.794098in}{1.386881in}}%
\pgfpathlineto{\pgfqpoint{3.838451in}{1.380995in}}%
\pgfpathlineto{\pgfqpoint{3.872874in}{1.360784in}}%
\pgfpathlineto{\pgfqpoint{3.828501in}{1.366664in}}%
\pgfpathlineto{\pgfqpoint{3.794098in}{1.386881in}}%
\pgfpathclose%
\pgfusepath{stroke,fill}%
\end{pgfscope}%
\begin{pgfscope}%
\pgfpathrectangle{\pgfqpoint{0.050000in}{0.050000in}}{\pgfqpoint{4.900000in}{4.900000in}}%
\pgfusepath{clip}%
\pgfsetbuttcap%
\pgfsetroundjoin%
\definecolor{currentfill}{rgb}{0.367243,0.367243,0.367243}%
\pgfsetfillcolor{currentfill}%
\pgfsetfillopacity{0.900000}%
\pgfsetlinewidth{0.501875pt}%
\definecolor{currentstroke}{rgb}{0.200000,0.200000,0.200000}%
\pgfsetstrokecolor{currentstroke}%
\pgfsetdash{}{0pt}%
\pgfpathmoveto{\pgfqpoint{1.523025in}{3.025713in}}%
\pgfpathlineto{\pgfqpoint{1.571418in}{2.979423in}}%
\pgfpathlineto{\pgfqpoint{1.602229in}{2.959737in}}%
\pgfpathlineto{\pgfqpoint{1.553826in}{3.006058in}}%
\pgfpathlineto{\pgfqpoint{1.523025in}{3.025713in}}%
\pgfpathclose%
\pgfusepath{stroke,fill}%
\end{pgfscope}%
\begin{pgfscope}%
\pgfpathrectangle{\pgfqpoint{0.050000in}{0.050000in}}{\pgfqpoint{4.900000in}{4.900000in}}%
\pgfusepath{clip}%
\pgfsetbuttcap%
\pgfsetroundjoin%
\definecolor{currentfill}{rgb}{0.950173,0.950173,0.950173}%
\pgfsetfillcolor{currentfill}%
\pgfsetfillopacity{0.900000}%
\pgfsetlinewidth{0.501875pt}%
\definecolor{currentstroke}{rgb}{0.200000,0.200000,0.200000}%
\pgfsetstrokecolor{currentstroke}%
\pgfsetdash{}{0pt}%
\pgfpathmoveto{\pgfqpoint{3.479747in}{1.525410in}}%
\pgfpathlineto{\pgfqpoint{3.524331in}{1.505360in}}%
\pgfpathlineto{\pgfqpoint{3.558244in}{1.485266in}}%
\pgfpathlineto{\pgfqpoint{3.513645in}{1.505180in}}%
\pgfpathlineto{\pgfqpoint{3.479747in}{1.525410in}}%
\pgfpathclose%
\pgfusepath{stroke,fill}%
\end{pgfscope}%
\begin{pgfscope}%
\pgfpathrectangle{\pgfqpoint{0.050000in}{0.050000in}}{\pgfqpoint{4.900000in}{4.900000in}}%
\pgfusepath{clip}%
\pgfsetbuttcap%
\pgfsetroundjoin%
\definecolor{currentfill}{rgb}{0.624221,0.624221,0.624221}%
\pgfsetfillcolor{currentfill}%
\pgfsetfillopacity{0.900000}%
\pgfsetlinewidth{0.501875pt}%
\definecolor{currentstroke}{rgb}{0.200000,0.200000,0.200000}%
\pgfsetstrokecolor{currentstroke}%
\pgfsetdash{}{0pt}%
\pgfpathmoveto{\pgfqpoint{2.249334in}{2.438704in}}%
\pgfpathlineto{\pgfqpoint{2.296174in}{2.393678in}}%
\pgfpathlineto{\pgfqpoint{2.328123in}{2.373079in}}%
\pgfpathlineto{\pgfqpoint{2.281278in}{2.418129in}}%
\pgfpathlineto{\pgfqpoint{2.249334in}{2.438704in}}%
\pgfpathclose%
\pgfusepath{stroke,fill}%
\end{pgfscope}%
\begin{pgfscope}%
\pgfpathrectangle{\pgfqpoint{0.050000in}{0.050000in}}{\pgfqpoint{4.900000in}{4.900000in}}%
\pgfusepath{clip}%
\pgfsetbuttcap%
\pgfsetroundjoin%
\definecolor{currentfill}{rgb}{0.979700,0.979700,0.979700}%
\pgfsetfillcolor{currentfill}%
\pgfsetfillopacity{0.900000}%
\pgfsetlinewidth{0.501875pt}%
\definecolor{currentstroke}{rgb}{0.200000,0.200000,0.200000}%
\pgfsetstrokecolor{currentstroke}%
\pgfsetdash{}{0pt}%
\pgfpathmoveto{\pgfqpoint{3.872874in}{1.360784in}}%
\pgfpathlineto{\pgfqpoint{3.917187in}{1.357888in}}%
\pgfpathlineto{\pgfqpoint{3.951740in}{1.337579in}}%
\pgfpathlineto{\pgfqpoint{3.907406in}{1.340493in}}%
\pgfpathlineto{\pgfqpoint{3.872874in}{1.360784in}}%
\pgfpathclose%
\pgfusepath{stroke,fill}%
\end{pgfscope}%
\begin{pgfscope}%
\pgfpathrectangle{\pgfqpoint{0.050000in}{0.050000in}}{\pgfqpoint{4.900000in}{4.900000in}}%
\pgfusepath{clip}%
\pgfsetbuttcap%
\pgfsetroundjoin%
\definecolor{currentfill}{rgb}{0.942791,0.942791,0.942791}%
\pgfsetfillcolor{currentfill}%
\pgfsetfillopacity{0.900000}%
\pgfsetlinewidth{0.501875pt}%
\definecolor{currentstroke}{rgb}{0.200000,0.200000,0.200000}%
\pgfsetstrokecolor{currentstroke}%
\pgfsetdash{}{0pt}%
\pgfpathmoveto{\pgfqpoint{3.401285in}{1.569352in}}%
\pgfpathlineto{\pgfqpoint{3.445957in}{1.545563in}}%
\pgfpathlineto{\pgfqpoint{3.479747in}{1.525410in}}%
\pgfpathlineto{\pgfqpoint{3.435062in}{1.549033in}}%
\pgfpathlineto{\pgfqpoint{3.401285in}{1.569352in}}%
\pgfpathclose%
\pgfusepath{stroke,fill}%
\end{pgfscope}%
\begin{pgfscope}%
\pgfpathrectangle{\pgfqpoint{0.050000in}{0.050000in}}{\pgfqpoint{4.900000in}{4.900000in}}%
\pgfusepath{clip}%
\pgfsetbuttcap%
\pgfsetroundjoin%
\definecolor{currentfill}{rgb}{0.256517,0.256517,0.256517}%
\pgfsetfillcolor{currentfill}%
\pgfsetfillopacity{0.900000}%
\pgfsetlinewidth{0.501875pt}%
\definecolor{currentstroke}{rgb}{0.200000,0.200000,0.200000}%
\pgfsetstrokecolor{currentstroke}%
\pgfsetdash{}{0pt}%
\pgfpathmoveto{\pgfqpoint{1.254475in}{3.243417in}}%
\pgfpathlineto{\pgfqpoint{1.303453in}{3.196652in}}%
\pgfpathlineto{\pgfqpoint{1.333849in}{3.177303in}}%
\pgfpathlineto{\pgfqpoint{1.284859in}{3.224100in}}%
\pgfpathlineto{\pgfqpoint{1.254475in}{3.243417in}}%
\pgfpathclose%
\pgfusepath{stroke,fill}%
\end{pgfscope}%
\begin{pgfscope}%
\pgfpathrectangle{\pgfqpoint{0.050000in}{0.050000in}}{\pgfqpoint{4.900000in}{4.900000in}}%
\pgfusepath{clip}%
\pgfsetbuttcap%
\pgfsetroundjoin%
\definecolor{currentfill}{rgb}{0.525798,0.525798,0.525798}%
\pgfsetfillcolor{currentfill}%
\pgfsetfillopacity{0.900000}%
\pgfsetlinewidth{0.501875pt}%
\definecolor{currentstroke}{rgb}{0.200000,0.200000,0.200000}%
\pgfsetstrokecolor{currentstroke}%
\pgfsetdash{}{0pt}%
\pgfpathmoveto{\pgfqpoint{1.980890in}{2.656312in}}%
\pgfpathlineto{\pgfqpoint{2.028314in}{2.610802in}}%
\pgfpathlineto{\pgfqpoint{2.059848in}{2.590536in}}%
\pgfpathlineto{\pgfqpoint{2.012418in}{2.636073in}}%
\pgfpathlineto{\pgfqpoint{1.980890in}{2.656312in}}%
\pgfpathclose%
\pgfusepath{stroke,fill}%
\end{pgfscope}%
\begin{pgfscope}%
\pgfpathrectangle{\pgfqpoint{0.050000in}{0.050000in}}{\pgfqpoint{4.900000in}{4.900000in}}%
\pgfusepath{clip}%
\pgfsetbuttcap%
\pgfsetroundjoin%
\definecolor{currentfill}{rgb}{0.788112,0.788112,0.788112}%
\pgfsetfillcolor{currentfill}%
\pgfsetfillopacity{0.900000}%
\pgfsetlinewidth{0.501875pt}%
\definecolor{currentstroke}{rgb}{0.200000,0.200000,0.200000}%
\pgfsetstrokecolor{currentstroke}%
\pgfsetdash{}{0pt}%
\pgfpathmoveto{\pgfqpoint{2.707407in}{2.070402in}}%
\pgfpathlineto{\pgfqpoint{2.753280in}{2.027178in}}%
\pgfpathlineto{\pgfqpoint{2.785956in}{2.006237in}}%
\pgfpathlineto{\pgfqpoint{2.740080in}{2.049395in}}%
\pgfpathlineto{\pgfqpoint{2.707407in}{2.070402in}}%
\pgfpathclose%
\pgfusepath{stroke,fill}%
\end{pgfscope}%
\begin{pgfscope}%
\pgfpathrectangle{\pgfqpoint{0.050000in}{0.050000in}}{\pgfqpoint{4.900000in}{4.900000in}}%
\pgfusepath{clip}%
\pgfsetbuttcap%
\pgfsetroundjoin%
\definecolor{currentfill}{rgb}{0.983391,0.983391,0.983391}%
\pgfsetfillcolor{currentfill}%
\pgfsetfillopacity{0.900000}%
\pgfsetlinewidth{0.501875pt}%
\definecolor{currentstroke}{rgb}{0.200000,0.200000,0.200000}%
\pgfsetstrokecolor{currentstroke}%
\pgfsetdash{}{0pt}%
\pgfpathmoveto{\pgfqpoint{3.951740in}{1.337579in}}%
\pgfpathlineto{\pgfqpoint{3.996017in}{1.337372in}}%
\pgfpathlineto{\pgfqpoint{4.030702in}{1.316947in}}%
\pgfpathlineto{\pgfqpoint{3.986404in}{1.317190in}}%
\pgfpathlineto{\pgfqpoint{3.951740in}{1.337579in}}%
\pgfpathclose%
\pgfusepath{stroke,fill}%
\end{pgfscope}%
\begin{pgfscope}%
\pgfpathrectangle{\pgfqpoint{0.050000in}{0.050000in}}{\pgfqpoint{4.900000in}{4.900000in}}%
\pgfusepath{clip}%
\pgfsetbuttcap%
\pgfsetroundjoin%
\definecolor{currentfill}{rgb}{0.136563,0.136563,0.136563}%
\pgfsetfillcolor{currentfill}%
\pgfsetfillopacity{0.900000}%
\pgfsetlinewidth{0.501875pt}%
\definecolor{currentstroke}{rgb}{0.200000,0.200000,0.200000}%
\pgfsetstrokecolor{currentstroke}%
\pgfsetdash{}{0pt}%
\pgfpathmoveto{\pgfqpoint{0.985834in}{3.461198in}}%
\pgfpathlineto{\pgfqpoint{1.035397in}{3.413958in}}%
\pgfpathlineto{\pgfqpoint{1.065377in}{3.394946in}}%
\pgfpathlineto{\pgfqpoint{1.015801in}{3.442219in}}%
\pgfpathlineto{\pgfqpoint{0.985834in}{3.461198in}}%
\pgfpathclose%
\pgfusepath{stroke,fill}%
\end{pgfscope}%
\begin{pgfscope}%
\pgfpathrectangle{\pgfqpoint{0.050000in}{0.050000in}}{\pgfqpoint{4.900000in}{4.900000in}}%
\pgfusepath{clip}%
\pgfsetbuttcap%
\pgfsetroundjoin%
\definecolor{currentfill}{rgb}{0.858762,0.858762,0.858762}%
\pgfsetfillcolor{currentfill}%
\pgfsetfillopacity{0.900000}%
\pgfsetlinewidth{0.501875pt}%
\definecolor{currentstroke}{rgb}{0.200000,0.200000,0.200000}%
\pgfsetstrokecolor{currentstroke}%
\pgfsetdash{}{0pt}%
\pgfpathmoveto{\pgfqpoint{2.975806in}{1.860446in}}%
\pgfpathlineto{\pgfqpoint{3.021137in}{1.821404in}}%
\pgfpathlineto{\pgfqpoint{3.054240in}{1.800668in}}%
\pgfpathlineto{\pgfqpoint{3.008904in}{1.839538in}}%
\pgfpathlineto{\pgfqpoint{2.975806in}{1.860446in}}%
\pgfpathclose%
\pgfusepath{stroke,fill}%
\end{pgfscope}%
\begin{pgfscope}%
\pgfpathrectangle{\pgfqpoint{0.050000in}{0.050000in}}{\pgfqpoint{4.900000in}{4.900000in}}%
\pgfusepath{clip}%
\pgfsetbuttcap%
\pgfsetroundjoin%
\definecolor{currentfill}{rgb}{0.929504,0.929504,0.929504}%
\pgfsetfillcolor{currentfill}%
\pgfsetfillopacity{0.900000}%
\pgfsetlinewidth{0.501875pt}%
\definecolor{currentstroke}{rgb}{0.200000,0.200000,0.200000}%
\pgfsetstrokecolor{currentstroke}%
\pgfsetdash{}{0pt}%
\pgfpathmoveto{\pgfqpoint{3.322843in}{1.617013in}}%
\pgfpathlineto{\pgfqpoint{3.367616in}{1.589599in}}%
\pgfpathlineto{\pgfqpoint{3.401285in}{1.569352in}}%
\pgfpathlineto{\pgfqpoint{3.356500in}{1.596578in}}%
\pgfpathlineto{\pgfqpoint{3.322843in}{1.617013in}}%
\pgfpathclose%
\pgfusepath{stroke,fill}%
\end{pgfscope}%
\begin{pgfscope}%
\pgfpathrectangle{\pgfqpoint{0.050000in}{0.050000in}}{\pgfqpoint{4.900000in}{4.900000in}}%
\pgfusepath{clip}%
\pgfsetbuttcap%
\pgfsetroundjoin%
\definecolor{currentfill}{rgb}{0.428143,0.428143,0.428143}%
\pgfsetfillcolor{currentfill}%
\pgfsetfillopacity{0.900000}%
\pgfsetlinewidth{0.501875pt}%
\definecolor{currentstroke}{rgb}{0.200000,0.200000,0.200000}%
\pgfsetstrokecolor{currentstroke}%
\pgfsetdash{}{0pt}%
\pgfpathmoveto{\pgfqpoint{1.712355in}{2.874000in}}%
\pgfpathlineto{\pgfqpoint{1.760363in}{2.828016in}}%
\pgfpathlineto{\pgfqpoint{1.791482in}{2.808086in}}%
\pgfpathlineto{\pgfqpoint{1.743466in}{2.854099in}}%
\pgfpathlineto{\pgfqpoint{1.712355in}{2.874000in}}%
\pgfpathclose%
\pgfusepath{stroke,fill}%
\end{pgfscope}%
\begin{pgfscope}%
\pgfpathrectangle{\pgfqpoint{0.050000in}{0.050000in}}{\pgfqpoint{4.900000in}{4.900000in}}%
\pgfusepath{clip}%
\pgfsetbuttcap%
\pgfsetroundjoin%
\definecolor{currentfill}{rgb}{0.700992,0.700992,0.700992}%
\pgfsetfillcolor{currentfill}%
\pgfsetfillopacity{0.900000}%
\pgfsetlinewidth{0.501875pt}%
\definecolor{currentstroke}{rgb}{0.200000,0.200000,0.200000}%
\pgfsetstrokecolor{currentstroke}%
\pgfsetdash{}{0pt}%
\pgfpathmoveto{\pgfqpoint{2.438974in}{2.286812in}}%
\pgfpathlineto{\pgfqpoint{2.485426in}{2.242179in}}%
\pgfpathlineto{\pgfqpoint{2.517686in}{2.221363in}}%
\pgfpathlineto{\pgfqpoint{2.471230in}{2.266005in}}%
\pgfpathlineto{\pgfqpoint{2.438974in}{2.286812in}}%
\pgfpathclose%
\pgfusepath{stroke,fill}%
\end{pgfscope}%
\begin{pgfscope}%
\pgfpathrectangle{\pgfqpoint{0.050000in}{0.050000in}}{\pgfqpoint{4.900000in}{4.900000in}}%
\pgfusepath{clip}%
\pgfsetbuttcap%
\pgfsetroundjoin%
\definecolor{currentfill}{rgb}{0.031865,0.031865,0.031865}%
\pgfsetfillcolor{currentfill}%
\pgfsetfillopacity{0.900000}%
\pgfsetlinewidth{0.501875pt}%
\definecolor{currentstroke}{rgb}{0.200000,0.200000,0.200000}%
\pgfsetstrokecolor{currentstroke}%
\pgfsetdash{}{0pt}%
\pgfpathmoveto{\pgfqpoint{0.717101in}{3.679055in}}%
\pgfpathlineto{\pgfqpoint{0.767249in}{3.631340in}}%
\pgfpathlineto{\pgfqpoint{0.796813in}{3.612665in}}%
\pgfpathlineto{\pgfqpoint{0.746651in}{3.660416in}}%
\pgfpathlineto{\pgfqpoint{0.717101in}{3.679055in}}%
\pgfpathclose%
\pgfusepath{stroke,fill}%
\end{pgfscope}%
\begin{pgfscope}%
\pgfpathrectangle{\pgfqpoint{0.050000in}{0.050000in}}{\pgfqpoint{4.900000in}{4.900000in}}%
\pgfusepath{clip}%
\pgfsetbuttcap%
\pgfsetroundjoin%
\definecolor{currentfill}{rgb}{0.334764,0.334764,0.334764}%
\pgfsetfillcolor{currentfill}%
\pgfsetfillopacity{0.900000}%
\pgfsetlinewidth{0.501875pt}%
\definecolor{currentstroke}{rgb}{0.200000,0.200000,0.200000}%
\pgfsetstrokecolor{currentstroke}%
\pgfsetdash{}{0pt}%
\pgfpathmoveto{\pgfqpoint{1.443728in}{3.091765in}}%
\pgfpathlineto{\pgfqpoint{1.492320in}{3.045306in}}%
\pgfpathlineto{\pgfqpoint{1.523025in}{3.025713in}}%
\pgfpathlineto{\pgfqpoint{1.474422in}{3.072203in}}%
\pgfpathlineto{\pgfqpoint{1.443728in}{3.091765in}}%
\pgfpathclose%
\pgfusepath{stroke,fill}%
\end{pgfscope}%
\begin{pgfscope}%
\pgfpathrectangle{\pgfqpoint{0.050000in}{0.050000in}}{\pgfqpoint{4.900000in}{4.900000in}}%
\pgfusepath{clip}%
\pgfsetbuttcap%
\pgfsetroundjoin%
\definecolor{currentfill}{rgb}{0.590634,0.590634,0.590634}%
\pgfsetfillcolor{currentfill}%
\pgfsetfillopacity{0.900000}%
\pgfsetlinewidth{0.501875pt}%
\definecolor{currentstroke}{rgb}{0.200000,0.200000,0.200000}%
\pgfsetstrokecolor{currentstroke}%
\pgfsetdash{}{0pt}%
\pgfpathmoveto{\pgfqpoint{2.170454in}{2.504415in}}%
\pgfpathlineto{\pgfqpoint{2.217490in}{2.459215in}}%
\pgfpathlineto{\pgfqpoint{2.249334in}{2.438704in}}%
\pgfpathlineto{\pgfqpoint{2.202292in}{2.483930in}}%
\pgfpathlineto{\pgfqpoint{2.170454in}{2.504415in}}%
\pgfpathclose%
\pgfusepath{stroke,fill}%
\end{pgfscope}%
\begin{pgfscope}%
\pgfpathrectangle{\pgfqpoint{0.050000in}{0.050000in}}{\pgfqpoint{4.900000in}{4.900000in}}%
\pgfusepath{clip}%
\pgfsetbuttcap%
\pgfsetroundjoin%
\definecolor{currentfill}{rgb}{0.912526,0.912526,0.912526}%
\pgfsetfillcolor{currentfill}%
\pgfsetfillopacity{0.900000}%
\pgfsetlinewidth{0.501875pt}%
\definecolor{currentstroke}{rgb}{0.200000,0.200000,0.200000}%
\pgfsetstrokecolor{currentstroke}%
\pgfsetdash{}{0pt}%
\pgfpathmoveto{\pgfqpoint{3.244401in}{1.668201in}}%
\pgfpathlineto{\pgfqpoint{3.289292in}{1.637380in}}%
\pgfpathlineto{\pgfqpoint{3.322843in}{1.617013in}}%
\pgfpathlineto{\pgfqpoint{3.277941in}{1.647631in}}%
\pgfpathlineto{\pgfqpoint{3.244401in}{1.668201in}}%
\pgfpathclose%
\pgfusepath{stroke,fill}%
\end{pgfscope}%
\begin{pgfscope}%
\pgfpathrectangle{\pgfqpoint{0.050000in}{0.050000in}}{\pgfqpoint{4.900000in}{4.900000in}}%
\pgfusepath{clip}%
\pgfsetbuttcap%
\pgfsetroundjoin%
\definecolor{currentfill}{rgb}{0.217762,0.217762,0.217762}%
\pgfsetfillcolor{currentfill}%
\pgfsetfillopacity{0.900000}%
\pgfsetlinewidth{0.501875pt}%
\definecolor{currentstroke}{rgb}{0.200000,0.200000,0.200000}%
\pgfsetstrokecolor{currentstroke}%
\pgfsetdash{}{0pt}%
\pgfpathmoveto{\pgfqpoint{1.175010in}{3.309607in}}%
\pgfpathlineto{\pgfqpoint{1.224186in}{3.262673in}}%
\pgfpathlineto{\pgfqpoint{1.254475in}{3.243417in}}%
\pgfpathlineto{\pgfqpoint{1.205286in}{3.290383in}}%
\pgfpathlineto{\pgfqpoint{1.175010in}{3.309607in}}%
\pgfpathclose%
\pgfusepath{stroke,fill}%
\end{pgfscope}%
\begin{pgfscope}%
\pgfpathrectangle{\pgfqpoint{0.050000in}{0.050000in}}{\pgfqpoint{4.900000in}{4.900000in}}%
\pgfusepath{clip}%
\pgfsetbuttcap%
\pgfsetroundjoin%
\definecolor{currentfill}{rgb}{0.491349,0.491349,0.491349}%
\pgfsetfillcolor{currentfill}%
\pgfsetfillopacity{0.900000}%
\pgfsetlinewidth{0.501875pt}%
\definecolor{currentstroke}{rgb}{0.200000,0.200000,0.200000}%
\pgfsetstrokecolor{currentstroke}%
\pgfsetdash{}{0pt}%
\pgfpathmoveto{\pgfqpoint{1.901841in}{2.722164in}}%
\pgfpathlineto{\pgfqpoint{1.949461in}{2.676487in}}%
\pgfpathlineto{\pgfqpoint{1.980890in}{2.656312in}}%
\pgfpathlineto{\pgfqpoint{1.933263in}{2.702017in}}%
\pgfpathlineto{\pgfqpoint{1.901841in}{2.722164in}}%
\pgfpathclose%
\pgfusepath{stroke,fill}%
\end{pgfscope}%
\begin{pgfscope}%
\pgfpathrectangle{\pgfqpoint{0.050000in}{0.050000in}}{\pgfqpoint{4.900000in}{4.900000in}}%
\pgfusepath{clip}%
\pgfsetbuttcap%
\pgfsetroundjoin%
\definecolor{currentfill}{rgb}{0.836340,0.836340,0.836340}%
\pgfsetfillcolor{currentfill}%
\pgfsetfillopacity{0.900000}%
\pgfsetlinewidth{0.501875pt}%
\definecolor{currentstroke}{rgb}{0.200000,0.200000,0.200000}%
\pgfsetstrokecolor{currentstroke}%
\pgfsetdash{}{0pt}%
\pgfpathmoveto{\pgfqpoint{2.897308in}{1.922124in}}%
\pgfpathlineto{\pgfqpoint{2.942813in}{1.881301in}}%
\pgfpathlineto{\pgfqpoint{2.975806in}{1.860446in}}%
\pgfpathlineto{\pgfqpoint{2.930297in}{1.901129in}}%
\pgfpathlineto{\pgfqpoint{2.897308in}{1.922124in}}%
\pgfpathclose%
\pgfusepath{stroke,fill}%
\end{pgfscope}%
\begin{pgfscope}%
\pgfpathrectangle{\pgfqpoint{0.050000in}{0.050000in}}{\pgfqpoint{4.900000in}{4.900000in}}%
\pgfusepath{clip}%
\pgfsetbuttcap%
\pgfsetroundjoin%
\definecolor{currentfill}{rgb}{0.763998,0.763998,0.763998}%
\pgfsetfillcolor{currentfill}%
\pgfsetfillopacity{0.900000}%
\pgfsetlinewidth{0.501875pt}%
\definecolor{currentstroke}{rgb}{0.200000,0.200000,0.200000}%
\pgfsetstrokecolor{currentstroke}%
\pgfsetdash{}{0pt}%
\pgfpathmoveto{\pgfqpoint{2.628771in}{2.135232in}}%
\pgfpathlineto{\pgfqpoint{2.674837in}{2.091352in}}%
\pgfpathlineto{\pgfqpoint{2.707407in}{2.070402in}}%
\pgfpathlineto{\pgfqpoint{2.661339in}{2.114246in}}%
\pgfpathlineto{\pgfqpoint{2.628771in}{2.135232in}}%
\pgfpathclose%
\pgfusepath{stroke,fill}%
\end{pgfscope}%
\begin{pgfscope}%
\pgfpathrectangle{\pgfqpoint{0.050000in}{0.050000in}}{\pgfqpoint{4.900000in}{4.900000in}}%
\pgfusepath{clip}%
\pgfsetbuttcap%
\pgfsetroundjoin%
\definecolor{currentfill}{rgb}{0.100146,0.100146,0.100146}%
\pgfsetfillcolor{currentfill}%
\pgfsetfillopacity{0.900000}%
\pgfsetlinewidth{0.501875pt}%
\definecolor{currentstroke}{rgb}{0.200000,0.200000,0.200000}%
\pgfsetstrokecolor{currentstroke}%
\pgfsetdash{}{0pt}%
\pgfpathmoveto{\pgfqpoint{0.906199in}{3.527526in}}%
\pgfpathlineto{\pgfqpoint{0.955961in}{3.480117in}}%
\pgfpathlineto{\pgfqpoint{0.985834in}{3.461198in}}%
\pgfpathlineto{\pgfqpoint{0.936058in}{3.508641in}}%
\pgfpathlineto{\pgfqpoint{0.906199in}{3.527526in}}%
\pgfpathclose%
\pgfusepath{stroke,fill}%
\end{pgfscope}%
\begin{pgfscope}%
\pgfpathrectangle{\pgfqpoint{0.050000in}{0.050000in}}{\pgfqpoint{4.900000in}{4.900000in}}%
\pgfusepath{clip}%
\pgfsetbuttcap%
\pgfsetroundjoin%
\definecolor{currentfill}{rgb}{0.399723,0.399723,0.399723}%
\pgfsetfillcolor{currentfill}%
\pgfsetfillopacity{0.900000}%
\pgfsetlinewidth{0.501875pt}%
\definecolor{currentstroke}{rgb}{0.200000,0.200000,0.200000}%
\pgfsetstrokecolor{currentstroke}%
\pgfsetdash{}{0pt}%
\pgfpathmoveto{\pgfqpoint{1.633136in}{2.939991in}}%
\pgfpathlineto{\pgfqpoint{1.681342in}{2.893839in}}%
\pgfpathlineto{\pgfqpoint{1.712355in}{2.874000in}}%
\pgfpathlineto{\pgfqpoint{1.664141in}{2.920182in}}%
\pgfpathlineto{\pgfqpoint{1.633136in}{2.939991in}}%
\pgfpathclose%
\pgfusepath{stroke,fill}%
\end{pgfscope}%
\begin{pgfscope}%
\pgfpathrectangle{\pgfqpoint{0.050000in}{0.050000in}}{\pgfqpoint{4.900000in}{4.900000in}}%
\pgfusepath{clip}%
\pgfsetbuttcap%
\pgfsetroundjoin%
\definecolor{currentfill}{rgb}{0.895548,0.895548,0.895548}%
\pgfsetfillcolor{currentfill}%
\pgfsetfillopacity{0.900000}%
\pgfsetlinewidth{0.501875pt}%
\definecolor{currentstroke}{rgb}{0.200000,0.200000,0.200000}%
\pgfsetstrokecolor{currentstroke}%
\pgfsetdash{}{0pt}%
\pgfpathmoveto{\pgfqpoint{3.165942in}{1.722618in}}%
\pgfpathlineto{\pgfqpoint{3.210967in}{1.688707in}}%
\pgfpathlineto{\pgfqpoint{3.244401in}{1.668201in}}%
\pgfpathlineto{\pgfqpoint{3.199367in}{1.701907in}}%
\pgfpathlineto{\pgfqpoint{3.165942in}{1.722618in}}%
\pgfpathclose%
\pgfusepath{stroke,fill}%
\end{pgfscope}%
\begin{pgfscope}%
\pgfpathrectangle{\pgfqpoint{0.050000in}{0.050000in}}{\pgfqpoint{4.900000in}{4.900000in}}%
\pgfusepath{clip}%
\pgfsetbuttcap%
\pgfsetroundjoin%
\definecolor{currentfill}{rgb}{0.662607,0.662607,0.662607}%
\pgfsetfillcolor{currentfill}%
\pgfsetfillopacity{0.900000}%
\pgfsetlinewidth{0.501875pt}%
\definecolor{currentstroke}{rgb}{0.200000,0.200000,0.200000}%
\pgfsetstrokecolor{currentstroke}%
\pgfsetdash{}{0pt}%
\pgfpathmoveto{\pgfqpoint{2.360172in}{2.352416in}}%
\pgfpathlineto{\pgfqpoint{2.406820in}{2.307557in}}%
\pgfpathlineto{\pgfqpoint{2.438974in}{2.286812in}}%
\pgfpathlineto{\pgfqpoint{2.392322in}{2.331690in}}%
\pgfpathlineto{\pgfqpoint{2.360172in}{2.352416in}}%
\pgfpathclose%
\pgfusepath{stroke,fill}%
\end{pgfscope}%
\begin{pgfscope}%
\pgfpathrectangle{\pgfqpoint{0.050000in}{0.050000in}}{\pgfqpoint{4.900000in}{4.900000in}}%
\pgfusepath{clip}%
\pgfsetbuttcap%
\pgfsetroundjoin%
\definecolor{currentfill}{rgb}{0.000000,0.000000,0.000000}%
\pgfsetfillcolor{currentfill}%
\pgfsetfillopacity{0.900000}%
\pgfsetlinewidth{0.501875pt}%
\definecolor{currentstroke}{rgb}{0.200000,0.200000,0.200000}%
\pgfsetstrokecolor{currentstroke}%
\pgfsetdash{}{0pt}%
\pgfpathmoveto{\pgfqpoint{0.637296in}{3.745523in}}%
\pgfpathlineto{\pgfqpoint{0.687644in}{3.697637in}}%
\pgfpathlineto{\pgfqpoint{0.717101in}{3.679055in}}%
\pgfpathlineto{\pgfqpoint{0.666738in}{3.726977in}}%
\pgfpathlineto{\pgfqpoint{0.637296in}{3.745523in}}%
\pgfpathclose%
\pgfusepath{stroke,fill}%
\end{pgfscope}%
\begin{pgfscope}%
\pgfpathrectangle{\pgfqpoint{0.050000in}{0.050000in}}{\pgfqpoint{4.900000in}{4.900000in}}%
\pgfusepath{clip}%
\pgfsetbuttcap%
\pgfsetroundjoin%
\definecolor{currentfill}{rgb}{0.300807,0.300807,0.300807}%
\pgfsetfillcolor{currentfill}%
\pgfsetfillopacity{0.900000}%
\pgfsetlinewidth{0.501875pt}%
\definecolor{currentstroke}{rgb}{0.200000,0.200000,0.200000}%
\pgfsetstrokecolor{currentstroke}%
\pgfsetdash{}{0pt}%
\pgfpathmoveto{\pgfqpoint{1.364340in}{3.157894in}}%
\pgfpathlineto{\pgfqpoint{1.413131in}{3.111267in}}%
\pgfpathlineto{\pgfqpoint{1.443728in}{3.091765in}}%
\pgfpathlineto{\pgfqpoint{1.394927in}{3.138424in}}%
\pgfpathlineto{\pgfqpoint{1.364340in}{3.157894in}}%
\pgfpathclose%
\pgfusepath{stroke,fill}%
\end{pgfscope}%
\begin{pgfscope}%
\pgfpathrectangle{\pgfqpoint{0.050000in}{0.050000in}}{\pgfqpoint{4.900000in}{4.900000in}}%
\pgfusepath{clip}%
\pgfsetbuttcap%
\pgfsetroundjoin%
\definecolor{currentfill}{rgb}{0.555940,0.555940,0.555940}%
\pgfsetfillcolor{currentfill}%
\pgfsetfillopacity{0.900000}%
\pgfsetlinewidth{0.501875pt}%
\definecolor{currentstroke}{rgb}{0.200000,0.200000,0.200000}%
\pgfsetstrokecolor{currentstroke}%
\pgfsetdash{}{0pt}%
\pgfpathmoveto{\pgfqpoint{2.091482in}{2.570205in}}%
\pgfpathlineto{\pgfqpoint{2.138715in}{2.524836in}}%
\pgfpathlineto{\pgfqpoint{2.170454in}{2.504415in}}%
\pgfpathlineto{\pgfqpoint{2.123215in}{2.549811in}}%
\pgfpathlineto{\pgfqpoint{2.091482in}{2.570205in}}%
\pgfpathclose%
\pgfusepath{stroke,fill}%
\end{pgfscope}%
\begin{pgfscope}%
\pgfpathrectangle{\pgfqpoint{0.050000in}{0.050000in}}{\pgfqpoint{4.900000in}{4.900000in}}%
\pgfusepath{clip}%
\pgfsetbuttcap%
\pgfsetroundjoin%
\definecolor{currentfill}{rgb}{0.173472,0.173472,0.173472}%
\pgfsetfillcolor{currentfill}%
\pgfsetfillopacity{0.900000}%
\pgfsetlinewidth{0.501875pt}%
\definecolor{currentstroke}{rgb}{0.200000,0.200000,0.200000}%
\pgfsetstrokecolor{currentstroke}%
\pgfsetdash{}{0pt}%
\pgfpathmoveto{\pgfqpoint{1.095452in}{3.375875in}}%
\pgfpathlineto{\pgfqpoint{1.144828in}{3.328771in}}%
\pgfpathlineto{\pgfqpoint{1.175010in}{3.309607in}}%
\pgfpathlineto{\pgfqpoint{1.125621in}{3.356744in}}%
\pgfpathlineto{\pgfqpoint{1.095452in}{3.375875in}}%
\pgfpathclose%
\pgfusepath{stroke,fill}%
\end{pgfscope}%
\begin{pgfscope}%
\pgfpathrectangle{\pgfqpoint{0.050000in}{0.050000in}}{\pgfqpoint{4.900000in}{4.900000in}}%
\pgfusepath{clip}%
\pgfsetbuttcap%
\pgfsetroundjoin%
\definecolor{currentfill}{rgb}{0.812226,0.812226,0.812226}%
\pgfsetfillcolor{currentfill}%
\pgfsetfillopacity{0.900000}%
\pgfsetlinewidth{0.501875pt}%
\definecolor{currentstroke}{rgb}{0.200000,0.200000,0.200000}%
\pgfsetstrokecolor{currentstroke}%
\pgfsetdash{}{0pt}%
\pgfpathmoveto{\pgfqpoint{2.818735in}{1.985243in}}%
\pgfpathlineto{\pgfqpoint{2.864422in}{1.943068in}}%
\pgfpathlineto{\pgfqpoint{2.897308in}{1.922124in}}%
\pgfpathlineto{\pgfqpoint{2.851617in}{1.964195in}}%
\pgfpathlineto{\pgfqpoint{2.818735in}{1.985243in}}%
\pgfpathclose%
\pgfusepath{stroke,fill}%
\end{pgfscope}%
\begin{pgfscope}%
\pgfpathrectangle{\pgfqpoint{0.050000in}{0.050000in}}{\pgfqpoint{4.900000in}{4.900000in}}%
\pgfusepath{clip}%
\pgfsetbuttcap%
\pgfsetroundjoin%
\definecolor{currentfill}{rgb}{0.461207,0.461207,0.461207}%
\pgfsetfillcolor{currentfill}%
\pgfsetfillopacity{0.900000}%
\pgfsetlinewidth{0.501875pt}%
\definecolor{currentstroke}{rgb}{0.200000,0.200000,0.200000}%
\pgfsetstrokecolor{currentstroke}%
\pgfsetdash{}{0pt}%
\pgfpathmoveto{\pgfqpoint{1.822700in}{2.788093in}}%
\pgfpathlineto{\pgfqpoint{1.870518in}{2.742248in}}%
\pgfpathlineto{\pgfqpoint{1.901841in}{2.722164in}}%
\pgfpathlineto{\pgfqpoint{1.854015in}{2.768038in}}%
\pgfpathlineto{\pgfqpoint{1.822700in}{2.788093in}}%
\pgfpathclose%
\pgfusepath{stroke,fill}%
\end{pgfscope}%
\begin{pgfscope}%
\pgfpathrectangle{\pgfqpoint{0.050000in}{0.050000in}}{\pgfqpoint{4.900000in}{4.900000in}}%
\pgfusepath{clip}%
\pgfsetbuttcap%
\pgfsetroundjoin%
\definecolor{currentfill}{rgb}{0.878570,0.878570,0.878570}%
\pgfsetfillcolor{currentfill}%
\pgfsetfillopacity{0.900000}%
\pgfsetlinewidth{0.501875pt}%
\definecolor{currentstroke}{rgb}{0.200000,0.200000,0.200000}%
\pgfsetstrokecolor{currentstroke}%
\pgfsetdash{}{0pt}%
\pgfpathmoveto{\pgfqpoint{3.087448in}{1.779878in}}%
\pgfpathlineto{\pgfqpoint{3.132622in}{1.743270in}}%
\pgfpathlineto{\pgfqpoint{3.165942in}{1.722618in}}%
\pgfpathlineto{\pgfqpoint{3.120761in}{1.759033in}}%
\pgfpathlineto{\pgfqpoint{3.087448in}{1.779878in}}%
\pgfpathclose%
\pgfusepath{stroke,fill}%
\end{pgfscope}%
\begin{pgfscope}%
\pgfpathrectangle{\pgfqpoint{0.050000in}{0.050000in}}{\pgfqpoint{4.900000in}{4.900000in}}%
\pgfusepath{clip}%
\pgfsetbuttcap%
\pgfsetroundjoin%
\definecolor{currentfill}{rgb}{0.734579,0.734579,0.734579}%
\pgfsetfillcolor{currentfill}%
\pgfsetfillopacity{0.900000}%
\pgfsetlinewidth{0.501875pt}%
\definecolor{currentstroke}{rgb}{0.200000,0.200000,0.200000}%
\pgfsetstrokecolor{currentstroke}%
\pgfsetdash{}{0pt}%
\pgfpathmoveto{\pgfqpoint{2.550046in}{2.200487in}}%
\pgfpathlineto{\pgfqpoint{2.596306in}{2.156159in}}%
\pgfpathlineto{\pgfqpoint{2.628771in}{2.135232in}}%
\pgfpathlineto{\pgfqpoint{2.582508in}{2.179550in}}%
\pgfpathlineto{\pgfqpoint{2.550046in}{2.200487in}}%
\pgfpathclose%
\pgfusepath{stroke,fill}%
\end{pgfscope}%
\begin{pgfscope}%
\pgfpathrectangle{\pgfqpoint{0.050000in}{0.050000in}}{\pgfqpoint{4.900000in}{4.900000in}}%
\pgfusepath{clip}%
\pgfsetbuttcap%
\pgfsetroundjoin%
\definecolor{currentfill}{rgb}{0.964937,0.964937,0.964937}%
\pgfsetfillcolor{currentfill}%
\pgfsetfillopacity{0.900000}%
\pgfsetlinewidth{0.501875pt}%
\definecolor{currentstroke}{rgb}{0.200000,0.200000,0.200000}%
\pgfsetstrokecolor{currentstroke}%
\pgfsetdash{}{0pt}%
\pgfpathmoveto{\pgfqpoint{3.670940in}{1.428738in}}%
\pgfpathlineto{\pgfqpoint{3.715408in}{1.416163in}}%
\pgfpathlineto{\pgfqpoint{3.749682in}{1.395995in}}%
\pgfpathlineto{\pgfqpoint{3.705197in}{1.408503in}}%
\pgfpathlineto{\pgfqpoint{3.670940in}{1.428738in}}%
\pgfpathclose%
\pgfusepath{stroke,fill}%
\end{pgfscope}%
\begin{pgfscope}%
\pgfpathrectangle{\pgfqpoint{0.050000in}{0.050000in}}{\pgfqpoint{4.900000in}{4.900000in}}%
\pgfusepath{clip}%
\pgfsetbuttcap%
\pgfsetroundjoin%
\definecolor{currentfill}{rgb}{0.068281,0.068281,0.068281}%
\pgfsetfillcolor{currentfill}%
\pgfsetfillopacity{0.900000}%
\pgfsetlinewidth{0.501875pt}%
\definecolor{currentstroke}{rgb}{0.200000,0.200000,0.200000}%
\pgfsetstrokecolor{currentstroke}%
\pgfsetdash{}{0pt}%
\pgfpathmoveto{\pgfqpoint{0.826471in}{3.593932in}}%
\pgfpathlineto{\pgfqpoint{0.876433in}{3.546353in}}%
\pgfpathlineto{\pgfqpoint{0.906199in}{3.527526in}}%
\pgfpathlineto{\pgfqpoint{0.856223in}{3.575141in}}%
\pgfpathlineto{\pgfqpoint{0.826471in}{3.593932in}}%
\pgfpathclose%
\pgfusepath{stroke,fill}%
\end{pgfscope}%
\begin{pgfscope}%
\pgfpathrectangle{\pgfqpoint{0.050000in}{0.050000in}}{\pgfqpoint{4.900000in}{4.900000in}}%
\pgfusepath{clip}%
\pgfsetbuttcap%
\pgfsetroundjoin%
\definecolor{currentfill}{rgb}{0.957555,0.957555,0.957555}%
\pgfsetfillcolor{currentfill}%
\pgfsetfillopacity{0.900000}%
\pgfsetlinewidth{0.501875pt}%
\definecolor{currentstroke}{rgb}{0.200000,0.200000,0.200000}%
\pgfsetstrokecolor{currentstroke}%
\pgfsetdash{}{0pt}%
\pgfpathmoveto{\pgfqpoint{3.592266in}{1.465092in}}%
\pgfpathlineto{\pgfqpoint{3.636794in}{1.448892in}}%
\pgfpathlineto{\pgfqpoint{3.670940in}{1.428738in}}%
\pgfpathlineto{\pgfqpoint{3.626397in}{1.444839in}}%
\pgfpathlineto{\pgfqpoint{3.592266in}{1.465092in}}%
\pgfpathclose%
\pgfusepath{stroke,fill}%
\end{pgfscope}%
\begin{pgfscope}%
\pgfpathrectangle{\pgfqpoint{0.050000in}{0.050000in}}{\pgfqpoint{4.900000in}{4.900000in}}%
\pgfusepath{clip}%
\pgfsetbuttcap%
\pgfsetroundjoin%
\definecolor{currentfill}{rgb}{0.970473,0.970473,0.970473}%
\pgfsetfillcolor{currentfill}%
\pgfsetfillopacity{0.900000}%
\pgfsetlinewidth{0.501875pt}%
\definecolor{currentstroke}{rgb}{0.200000,0.200000,0.200000}%
\pgfsetstrokecolor{currentstroke}%
\pgfsetdash{}{0pt}%
\pgfpathmoveto{\pgfqpoint{3.749682in}{1.395995in}}%
\pgfpathlineto{\pgfqpoint{3.794098in}{1.386881in}}%
\pgfpathlineto{\pgfqpoint{3.828501in}{1.366664in}}%
\pgfpathlineto{\pgfqpoint{3.784066in}{1.375744in}}%
\pgfpathlineto{\pgfqpoint{3.749682in}{1.395995in}}%
\pgfpathclose%
\pgfusepath{stroke,fill}%
\end{pgfscope}%
\begin{pgfscope}%
\pgfpathrectangle{\pgfqpoint{0.050000in}{0.050000in}}{\pgfqpoint{4.900000in}{4.900000in}}%
\pgfusepath{clip}%
\pgfsetbuttcap%
\pgfsetroundjoin%
\definecolor{currentfill}{rgb}{0.367243,0.367243,0.367243}%
\pgfsetfillcolor{currentfill}%
\pgfsetfillopacity{0.900000}%
\pgfsetlinewidth{0.501875pt}%
\definecolor{currentstroke}{rgb}{0.200000,0.200000,0.200000}%
\pgfsetstrokecolor{currentstroke}%
\pgfsetdash{}{0pt}%
\pgfpathmoveto{\pgfqpoint{1.553826in}{3.006058in}}%
\pgfpathlineto{\pgfqpoint{1.602229in}{2.959737in}}%
\pgfpathlineto{\pgfqpoint{1.633136in}{2.939991in}}%
\pgfpathlineto{\pgfqpoint{1.584723in}{2.986342in}}%
\pgfpathlineto{\pgfqpoint{1.553826in}{3.006058in}}%
\pgfpathclose%
\pgfusepath{stroke,fill}%
\end{pgfscope}%
\begin{pgfscope}%
\pgfpathrectangle{\pgfqpoint{0.050000in}{0.050000in}}{\pgfqpoint{4.900000in}{4.900000in}}%
\pgfusepath{clip}%
\pgfsetbuttcap%
\pgfsetroundjoin%
\definecolor{currentfill}{rgb}{0.976009,0.976009,0.976009}%
\pgfsetfillcolor{currentfill}%
\pgfsetfillopacity{0.900000}%
\pgfsetlinewidth{0.501875pt}%
\definecolor{currentstroke}{rgb}{0.200000,0.200000,0.200000}%
\pgfsetstrokecolor{currentstroke}%
\pgfsetdash{}{0pt}%
\pgfpathmoveto{\pgfqpoint{3.828501in}{1.366664in}}%
\pgfpathlineto{\pgfqpoint{3.872874in}{1.360784in}}%
\pgfpathlineto{\pgfqpoint{3.907406in}{1.340493in}}%
\pgfpathlineto{\pgfqpoint{3.863014in}{1.346367in}}%
\pgfpathlineto{\pgfqpoint{3.828501in}{1.366664in}}%
\pgfpathclose%
\pgfusepath{stroke,fill}%
\end{pgfscope}%
\begin{pgfscope}%
\pgfpathrectangle{\pgfqpoint{0.050000in}{0.050000in}}{\pgfqpoint{4.900000in}{4.900000in}}%
\pgfusepath{clip}%
\pgfsetbuttcap%
\pgfsetroundjoin%
\definecolor{currentfill}{rgb}{0.950173,0.950173,0.950173}%
\pgfsetfillcolor{currentfill}%
\pgfsetfillopacity{0.900000}%
\pgfsetlinewidth{0.501875pt}%
\definecolor{currentstroke}{rgb}{0.200000,0.200000,0.200000}%
\pgfsetstrokecolor{currentstroke}%
\pgfsetdash{}{0pt}%
\pgfpathmoveto{\pgfqpoint{3.513645in}{1.505180in}}%
\pgfpathlineto{\pgfqpoint{3.558244in}{1.485266in}}%
\pgfpathlineto{\pgfqpoint{3.592266in}{1.465092in}}%
\pgfpathlineto{\pgfqpoint{3.547652in}{1.484874in}}%
\pgfpathlineto{\pgfqpoint{3.513645in}{1.505180in}}%
\pgfpathclose%
\pgfusepath{stroke,fill}%
\end{pgfscope}%
\begin{pgfscope}%
\pgfpathrectangle{\pgfqpoint{0.050000in}{0.050000in}}{\pgfqpoint{4.900000in}{4.900000in}}%
\pgfusepath{clip}%
\pgfsetbuttcap%
\pgfsetroundjoin%
\definecolor{currentfill}{rgb}{0.629020,0.629020,0.629020}%
\pgfsetfillcolor{currentfill}%
\pgfsetfillopacity{0.900000}%
\pgfsetlinewidth{0.501875pt}%
\definecolor{currentstroke}{rgb}{0.200000,0.200000,0.200000}%
\pgfsetstrokecolor{currentstroke}%
\pgfsetdash{}{0pt}%
\pgfpathmoveto{\pgfqpoint{2.281278in}{2.418129in}}%
\pgfpathlineto{\pgfqpoint{2.328123in}{2.373079in}}%
\pgfpathlineto{\pgfqpoint{2.360172in}{2.352416in}}%
\pgfpathlineto{\pgfqpoint{2.313323in}{2.397490in}}%
\pgfpathlineto{\pgfqpoint{2.281278in}{2.418129in}}%
\pgfpathclose%
\pgfusepath{stroke,fill}%
\end{pgfscope}%
\begin{pgfscope}%
\pgfpathrectangle{\pgfqpoint{0.050000in}{0.050000in}}{\pgfqpoint{4.900000in}{4.900000in}}%
\pgfusepath{clip}%
\pgfsetbuttcap%
\pgfsetroundjoin%
\definecolor{currentfill}{rgb}{0.256517,0.256517,0.256517}%
\pgfsetfillcolor{currentfill}%
\pgfsetfillopacity{0.900000}%
\pgfsetlinewidth{0.501875pt}%
\definecolor{currentstroke}{rgb}{0.200000,0.200000,0.200000}%
\pgfsetstrokecolor{currentstroke}%
\pgfsetdash{}{0pt}%
\pgfpathmoveto{\pgfqpoint{1.284859in}{3.224100in}}%
\pgfpathlineto{\pgfqpoint{1.333849in}{3.177303in}}%
\pgfpathlineto{\pgfqpoint{1.364340in}{3.157894in}}%
\pgfpathlineto{\pgfqpoint{1.315339in}{3.204723in}}%
\pgfpathlineto{\pgfqpoint{1.284859in}{3.224100in}}%
\pgfpathclose%
\pgfusepath{stroke,fill}%
\end{pgfscope}%
\begin{pgfscope}%
\pgfpathrectangle{\pgfqpoint{0.050000in}{0.050000in}}{\pgfqpoint{4.900000in}{4.900000in}}%
\pgfusepath{clip}%
\pgfsetbuttcap%
\pgfsetroundjoin%
\definecolor{currentfill}{rgb}{0.979700,0.979700,0.979700}%
\pgfsetfillcolor{currentfill}%
\pgfsetfillopacity{0.900000}%
\pgfsetlinewidth{0.501875pt}%
\definecolor{currentstroke}{rgb}{0.200000,0.200000,0.200000}%
\pgfsetstrokecolor{currentstroke}%
\pgfsetdash{}{0pt}%
\pgfpathmoveto{\pgfqpoint{3.907406in}{1.340493in}}%
\pgfpathlineto{\pgfqpoint{3.951740in}{1.337579in}}%
\pgfpathlineto{\pgfqpoint{3.986404in}{1.317190in}}%
\pgfpathlineto{\pgfqpoint{3.942050in}{1.320122in}}%
\pgfpathlineto{\pgfqpoint{3.907406in}{1.340493in}}%
\pgfpathclose%
\pgfusepath{stroke,fill}%
\end{pgfscope}%
\begin{pgfscope}%
\pgfpathrectangle{\pgfqpoint{0.050000in}{0.050000in}}{\pgfqpoint{4.900000in}{4.900000in}}%
\pgfusepath{clip}%
\pgfsetbuttcap%
\pgfsetroundjoin%
\definecolor{currentfill}{rgb}{0.942791,0.942791,0.942791}%
\pgfsetfillcolor{currentfill}%
\pgfsetfillopacity{0.900000}%
\pgfsetlinewidth{0.501875pt}%
\definecolor{currentstroke}{rgb}{0.200000,0.200000,0.200000}%
\pgfsetstrokecolor{currentstroke}%
\pgfsetdash{}{0pt}%
\pgfpathmoveto{\pgfqpoint{3.435062in}{1.549033in}}%
\pgfpathlineto{\pgfqpoint{3.479747in}{1.525410in}}%
\pgfpathlineto{\pgfqpoint{3.513645in}{1.505180in}}%
\pgfpathlineto{\pgfqpoint{3.468947in}{1.528642in}}%
\pgfpathlineto{\pgfqpoint{3.435062in}{1.549033in}}%
\pgfpathclose%
\pgfusepath{stroke,fill}%
\end{pgfscope}%
\begin{pgfscope}%
\pgfpathrectangle{\pgfqpoint{0.050000in}{0.050000in}}{\pgfqpoint{4.900000in}{4.900000in}}%
\pgfusepath{clip}%
\pgfsetbuttcap%
\pgfsetroundjoin%
\definecolor{currentfill}{rgb}{0.525798,0.525798,0.525798}%
\pgfsetfillcolor{currentfill}%
\pgfsetfillopacity{0.900000}%
\pgfsetlinewidth{0.501875pt}%
\definecolor{currentstroke}{rgb}{0.200000,0.200000,0.200000}%
\pgfsetstrokecolor{currentstroke}%
\pgfsetdash{}{0pt}%
\pgfpathmoveto{\pgfqpoint{2.012418in}{2.636073in}}%
\pgfpathlineto{\pgfqpoint{2.059848in}{2.590536in}}%
\pgfpathlineto{\pgfqpoint{2.091482in}{2.570205in}}%
\pgfpathlineto{\pgfqpoint{2.044045in}{2.615770in}}%
\pgfpathlineto{\pgfqpoint{2.012418in}{2.636073in}}%
\pgfpathclose%
\pgfusepath{stroke,fill}%
\end{pgfscope}%
\begin{pgfscope}%
\pgfpathrectangle{\pgfqpoint{0.050000in}{0.050000in}}{\pgfqpoint{4.900000in}{4.900000in}}%
\pgfusepath{clip}%
\pgfsetbuttcap%
\pgfsetroundjoin%
\definecolor{currentfill}{rgb}{0.136563,0.136563,0.136563}%
\pgfsetfillcolor{currentfill}%
\pgfsetfillopacity{0.900000}%
\pgfsetlinewidth{0.501875pt}%
\definecolor{currentstroke}{rgb}{0.200000,0.200000,0.200000}%
\pgfsetstrokecolor{currentstroke}%
\pgfsetdash{}{0pt}%
\pgfpathmoveto{\pgfqpoint{1.015801in}{3.442219in}}%
\pgfpathlineto{\pgfqpoint{1.065377in}{3.394946in}}%
\pgfpathlineto{\pgfqpoint{1.095452in}{3.375875in}}%
\pgfpathlineto{\pgfqpoint{1.045863in}{3.423182in}}%
\pgfpathlineto{\pgfqpoint{1.015801in}{3.442219in}}%
\pgfpathclose%
\pgfusepath{stroke,fill}%
\end{pgfscope}%
\begin{pgfscope}%
\pgfpathrectangle{\pgfqpoint{0.050000in}{0.050000in}}{\pgfqpoint{4.900000in}{4.900000in}}%
\pgfusepath{clip}%
\pgfsetbuttcap%
\pgfsetroundjoin%
\definecolor{currentfill}{rgb}{0.788112,0.788112,0.788112}%
\pgfsetfillcolor{currentfill}%
\pgfsetfillopacity{0.900000}%
\pgfsetlinewidth{0.501875pt}%
\definecolor{currentstroke}{rgb}{0.200000,0.200000,0.200000}%
\pgfsetstrokecolor{currentstroke}%
\pgfsetdash{}{0pt}%
\pgfpathmoveto{\pgfqpoint{2.740080in}{2.049395in}}%
\pgfpathlineto{\pgfqpoint{2.785956in}{2.006237in}}%
\pgfpathlineto{\pgfqpoint{2.818735in}{1.985243in}}%
\pgfpathlineto{\pgfqpoint{2.772856in}{2.028332in}}%
\pgfpathlineto{\pgfqpoint{2.740080in}{2.049395in}}%
\pgfpathclose%
\pgfusepath{stroke,fill}%
\end{pgfscope}%
\begin{pgfscope}%
\pgfpathrectangle{\pgfqpoint{0.050000in}{0.050000in}}{\pgfqpoint{4.900000in}{4.900000in}}%
\pgfusepath{clip}%
\pgfsetbuttcap%
\pgfsetroundjoin%
\definecolor{currentfill}{rgb}{0.858762,0.858762,0.858762}%
\pgfsetfillcolor{currentfill}%
\pgfsetfillopacity{0.900000}%
\pgfsetlinewidth{0.501875pt}%
\definecolor{currentstroke}{rgb}{0.200000,0.200000,0.200000}%
\pgfsetstrokecolor{currentstroke}%
\pgfsetdash{}{0pt}%
\pgfpathmoveto{\pgfqpoint{3.008904in}{1.839538in}}%
\pgfpathlineto{\pgfqpoint{3.054240in}{1.800668in}}%
\pgfpathlineto{\pgfqpoint{3.087448in}{1.779878in}}%
\pgfpathlineto{\pgfqpoint{3.042107in}{1.818577in}}%
\pgfpathlineto{\pgfqpoint{3.008904in}{1.839538in}}%
\pgfpathclose%
\pgfusepath{stroke,fill}%
\end{pgfscope}%
\begin{pgfscope}%
\pgfpathrectangle{\pgfqpoint{0.050000in}{0.050000in}}{\pgfqpoint{4.900000in}{4.900000in}}%
\pgfusepath{clip}%
\pgfsetbuttcap%
\pgfsetroundjoin%
\definecolor{currentfill}{rgb}{0.428143,0.428143,0.428143}%
\pgfsetfillcolor{currentfill}%
\pgfsetfillopacity{0.900000}%
\pgfsetlinewidth{0.501875pt}%
\definecolor{currentstroke}{rgb}{0.200000,0.200000,0.200000}%
\pgfsetstrokecolor{currentstroke}%
\pgfsetdash{}{0pt}%
\pgfpathmoveto{\pgfqpoint{1.743466in}{2.854099in}}%
\pgfpathlineto{\pgfqpoint{1.791482in}{2.808086in}}%
\pgfpathlineto{\pgfqpoint{1.822700in}{2.788093in}}%
\pgfpathlineto{\pgfqpoint{1.774675in}{2.834136in}}%
\pgfpathlineto{\pgfqpoint{1.743466in}{2.854099in}}%
\pgfpathclose%
\pgfusepath{stroke,fill}%
\end{pgfscope}%
\begin{pgfscope}%
\pgfpathrectangle{\pgfqpoint{0.050000in}{0.050000in}}{\pgfqpoint{4.900000in}{4.900000in}}%
\pgfusepath{clip}%
\pgfsetbuttcap%
\pgfsetroundjoin%
\definecolor{currentfill}{rgb}{0.929504,0.929504,0.929504}%
\pgfsetfillcolor{currentfill}%
\pgfsetfillopacity{0.900000}%
\pgfsetlinewidth{0.501875pt}%
\definecolor{currentstroke}{rgb}{0.200000,0.200000,0.200000}%
\pgfsetstrokecolor{currentstroke}%
\pgfsetdash{}{0pt}%
\pgfpathmoveto{\pgfqpoint{3.356500in}{1.596578in}}%
\pgfpathlineto{\pgfqpoint{3.401285in}{1.569352in}}%
\pgfpathlineto{\pgfqpoint{3.435062in}{1.549033in}}%
\pgfpathlineto{\pgfqpoint{3.390265in}{1.576075in}}%
\pgfpathlineto{\pgfqpoint{3.356500in}{1.596578in}}%
\pgfpathclose%
\pgfusepath{stroke,fill}%
\end{pgfscope}%
\begin{pgfscope}%
\pgfpathrectangle{\pgfqpoint{0.050000in}{0.050000in}}{\pgfqpoint{4.900000in}{4.900000in}}%
\pgfusepath{clip}%
\pgfsetbuttcap%
\pgfsetroundjoin%
\definecolor{currentfill}{rgb}{0.700992,0.700992,0.700992}%
\pgfsetfillcolor{currentfill}%
\pgfsetfillopacity{0.900000}%
\pgfsetlinewidth{0.501875pt}%
\definecolor{currentstroke}{rgb}{0.200000,0.200000,0.200000}%
\pgfsetstrokecolor{currentstroke}%
\pgfsetdash{}{0pt}%
\pgfpathmoveto{\pgfqpoint{2.471230in}{2.266005in}}%
\pgfpathlineto{\pgfqpoint{2.517686in}{2.221363in}}%
\pgfpathlineto{\pgfqpoint{2.550046in}{2.200487in}}%
\pgfpathlineto{\pgfqpoint{2.503587in}{2.245135in}}%
\pgfpathlineto{\pgfqpoint{2.471230in}{2.266005in}}%
\pgfpathclose%
\pgfusepath{stroke,fill}%
\end{pgfscope}%
\begin{pgfscope}%
\pgfpathrectangle{\pgfqpoint{0.050000in}{0.050000in}}{\pgfqpoint{4.900000in}{4.900000in}}%
\pgfusepath{clip}%
\pgfsetbuttcap%
\pgfsetroundjoin%
\definecolor{currentfill}{rgb}{0.031865,0.031865,0.031865}%
\pgfsetfillcolor{currentfill}%
\pgfsetfillopacity{0.900000}%
\pgfsetlinewidth{0.501875pt}%
\definecolor{currentstroke}{rgb}{0.200000,0.200000,0.200000}%
\pgfsetstrokecolor{currentstroke}%
\pgfsetdash{}{0pt}%
\pgfpathmoveto{\pgfqpoint{0.746651in}{3.660416in}}%
\pgfpathlineto{\pgfqpoint{0.796813in}{3.612665in}}%
\pgfpathlineto{\pgfqpoint{0.826471in}{3.593932in}}%
\pgfpathlineto{\pgfqpoint{0.776294in}{3.641718in}}%
\pgfpathlineto{\pgfqpoint{0.746651in}{3.660416in}}%
\pgfpathclose%
\pgfusepath{stroke,fill}%
\end{pgfscope}%
\begin{pgfscope}%
\pgfpathrectangle{\pgfqpoint{0.050000in}{0.050000in}}{\pgfqpoint{4.900000in}{4.900000in}}%
\pgfusepath{clip}%
\pgfsetbuttcap%
\pgfsetroundjoin%
\definecolor{currentfill}{rgb}{0.338824,0.338824,0.338824}%
\pgfsetfillcolor{currentfill}%
\pgfsetfillopacity{0.900000}%
\pgfsetlinewidth{0.501875pt}%
\definecolor{currentstroke}{rgb}{0.200000,0.200000,0.200000}%
\pgfsetstrokecolor{currentstroke}%
\pgfsetdash{}{0pt}%
\pgfpathmoveto{\pgfqpoint{1.474422in}{3.072203in}}%
\pgfpathlineto{\pgfqpoint{1.523025in}{3.025713in}}%
\pgfpathlineto{\pgfqpoint{1.553826in}{3.006058in}}%
\pgfpathlineto{\pgfqpoint{1.505213in}{3.052579in}}%
\pgfpathlineto{\pgfqpoint{1.474422in}{3.072203in}}%
\pgfpathclose%
\pgfusepath{stroke,fill}%
\end{pgfscope}%
\begin{pgfscope}%
\pgfpathrectangle{\pgfqpoint{0.050000in}{0.050000in}}{\pgfqpoint{4.900000in}{4.900000in}}%
\pgfusepath{clip}%
\pgfsetbuttcap%
\pgfsetroundjoin%
\definecolor{currentfill}{rgb}{0.590634,0.590634,0.590634}%
\pgfsetfillcolor{currentfill}%
\pgfsetfillopacity{0.900000}%
\pgfsetlinewidth{0.501875pt}%
\definecolor{currentstroke}{rgb}{0.200000,0.200000,0.200000}%
\pgfsetstrokecolor{currentstroke}%
\pgfsetdash{}{0pt}%
\pgfpathmoveto{\pgfqpoint{2.202292in}{2.483930in}}%
\pgfpathlineto{\pgfqpoint{2.249334in}{2.438704in}}%
\pgfpathlineto{\pgfqpoint{2.281278in}{2.418129in}}%
\pgfpathlineto{\pgfqpoint{2.234231in}{2.463381in}}%
\pgfpathlineto{\pgfqpoint{2.202292in}{2.483930in}}%
\pgfpathclose%
\pgfusepath{stroke,fill}%
\end{pgfscope}%
\begin{pgfscope}%
\pgfpathrectangle{\pgfqpoint{0.050000in}{0.050000in}}{\pgfqpoint{4.900000in}{4.900000in}}%
\pgfusepath{clip}%
\pgfsetbuttcap%
\pgfsetroundjoin%
\definecolor{currentfill}{rgb}{0.912526,0.912526,0.912526}%
\pgfsetfillcolor{currentfill}%
\pgfsetfillopacity{0.900000}%
\pgfsetlinewidth{0.501875pt}%
\definecolor{currentstroke}{rgb}{0.200000,0.200000,0.200000}%
\pgfsetstrokecolor{currentstroke}%
\pgfsetdash{}{0pt}%
\pgfpathmoveto{\pgfqpoint{3.277941in}{1.647631in}}%
\pgfpathlineto{\pgfqpoint{3.322843in}{1.617013in}}%
\pgfpathlineto{\pgfqpoint{3.356500in}{1.596578in}}%
\pgfpathlineto{\pgfqpoint{3.311589in}{1.626998in}}%
\pgfpathlineto{\pgfqpoint{3.277941in}{1.647631in}}%
\pgfpathclose%
\pgfusepath{stroke,fill}%
\end{pgfscope}%
\begin{pgfscope}%
\pgfpathrectangle{\pgfqpoint{0.050000in}{0.050000in}}{\pgfqpoint{4.900000in}{4.900000in}}%
\pgfusepath{clip}%
\pgfsetbuttcap%
\pgfsetroundjoin%
\definecolor{currentfill}{rgb}{0.217762,0.217762,0.217762}%
\pgfsetfillcolor{currentfill}%
\pgfsetfillopacity{0.900000}%
\pgfsetlinewidth{0.501875pt}%
\definecolor{currentstroke}{rgb}{0.200000,0.200000,0.200000}%
\pgfsetstrokecolor{currentstroke}%
\pgfsetdash{}{0pt}%
\pgfpathmoveto{\pgfqpoint{1.205286in}{3.290383in}}%
\pgfpathlineto{\pgfqpoint{1.254475in}{3.243417in}}%
\pgfpathlineto{\pgfqpoint{1.284859in}{3.224100in}}%
\pgfpathlineto{\pgfqpoint{1.235659in}{3.271099in}}%
\pgfpathlineto{\pgfqpoint{1.205286in}{3.290383in}}%
\pgfpathclose%
\pgfusepath{stroke,fill}%
\end{pgfscope}%
\begin{pgfscope}%
\pgfpathrectangle{\pgfqpoint{0.050000in}{0.050000in}}{\pgfqpoint{4.900000in}{4.900000in}}%
\pgfusepath{clip}%
\pgfsetbuttcap%
\pgfsetroundjoin%
\definecolor{currentfill}{rgb}{0.491349,0.491349,0.491349}%
\pgfsetfillcolor{currentfill}%
\pgfsetfillopacity{0.900000}%
\pgfsetlinewidth{0.501875pt}%
\definecolor{currentstroke}{rgb}{0.200000,0.200000,0.200000}%
\pgfsetstrokecolor{currentstroke}%
\pgfsetdash{}{0pt}%
\pgfpathmoveto{\pgfqpoint{1.933263in}{2.702017in}}%
\pgfpathlineto{\pgfqpoint{1.980890in}{2.656312in}}%
\pgfpathlineto{\pgfqpoint{2.012418in}{2.636073in}}%
\pgfpathlineto{\pgfqpoint{1.964783in}{2.681806in}}%
\pgfpathlineto{\pgfqpoint{1.933263in}{2.702017in}}%
\pgfpathclose%
\pgfusepath{stroke,fill}%
\end{pgfscope}%
\begin{pgfscope}%
\pgfpathrectangle{\pgfqpoint{0.050000in}{0.050000in}}{\pgfqpoint{4.900000in}{4.900000in}}%
\pgfusepath{clip}%
\pgfsetbuttcap%
\pgfsetroundjoin%
\definecolor{currentfill}{rgb}{0.839785,0.839785,0.839785}%
\pgfsetfillcolor{currentfill}%
\pgfsetfillopacity{0.900000}%
\pgfsetlinewidth{0.501875pt}%
\definecolor{currentstroke}{rgb}{0.200000,0.200000,0.200000}%
\pgfsetstrokecolor{currentstroke}%
\pgfsetdash{}{0pt}%
\pgfpathmoveto{\pgfqpoint{2.930297in}{1.901129in}}%
\pgfpathlineto{\pgfqpoint{2.975806in}{1.860446in}}%
\pgfpathlineto{\pgfqpoint{3.008904in}{1.839538in}}%
\pgfpathlineto{\pgfqpoint{2.963391in}{1.880080in}}%
\pgfpathlineto{\pgfqpoint{2.930297in}{1.901129in}}%
\pgfpathclose%
\pgfusepath{stroke,fill}%
\end{pgfscope}%
\begin{pgfscope}%
\pgfpathrectangle{\pgfqpoint{0.050000in}{0.050000in}}{\pgfqpoint{4.900000in}{4.900000in}}%
\pgfusepath{clip}%
\pgfsetbuttcap%
\pgfsetroundjoin%
\definecolor{currentfill}{rgb}{0.763998,0.763998,0.763998}%
\pgfsetfillcolor{currentfill}%
\pgfsetfillopacity{0.900000}%
\pgfsetlinewidth{0.501875pt}%
\definecolor{currentstroke}{rgb}{0.200000,0.200000,0.200000}%
\pgfsetstrokecolor{currentstroke}%
\pgfsetdash{}{0pt}%
\pgfpathmoveto{\pgfqpoint{2.661339in}{2.114246in}}%
\pgfpathlineto{\pgfqpoint{2.707407in}{2.070402in}}%
\pgfpathlineto{\pgfqpoint{2.740080in}{2.049395in}}%
\pgfpathlineto{\pgfqpoint{2.694009in}{2.093202in}}%
\pgfpathlineto{\pgfqpoint{2.661339in}{2.114246in}}%
\pgfpathclose%
\pgfusepath{stroke,fill}%
\end{pgfscope}%
\begin{pgfscope}%
\pgfpathrectangle{\pgfqpoint{0.050000in}{0.050000in}}{\pgfqpoint{4.900000in}{4.900000in}}%
\pgfusepath{clip}%
\pgfsetbuttcap%
\pgfsetroundjoin%
\definecolor{currentfill}{rgb}{0.100146,0.100146,0.100146}%
\pgfsetfillcolor{currentfill}%
\pgfsetfillopacity{0.900000}%
\pgfsetlinewidth{0.501875pt}%
\definecolor{currentstroke}{rgb}{0.200000,0.200000,0.200000}%
\pgfsetstrokecolor{currentstroke}%
\pgfsetdash{}{0pt}%
\pgfpathmoveto{\pgfqpoint{0.936058in}{3.508641in}}%
\pgfpathlineto{\pgfqpoint{0.985834in}{3.461198in}}%
\pgfpathlineto{\pgfqpoint{1.015801in}{3.442219in}}%
\pgfpathlineto{\pgfqpoint{0.966012in}{3.489697in}}%
\pgfpathlineto{\pgfqpoint{0.936058in}{3.508641in}}%
\pgfpathclose%
\pgfusepath{stroke,fill}%
\end{pgfscope}%
\begin{pgfscope}%
\pgfpathrectangle{\pgfqpoint{0.050000in}{0.050000in}}{\pgfqpoint{4.900000in}{4.900000in}}%
\pgfusepath{clip}%
\pgfsetbuttcap%
\pgfsetroundjoin%
\definecolor{currentfill}{rgb}{0.399723,0.399723,0.399723}%
\pgfsetfillcolor{currentfill}%
\pgfsetfillopacity{0.900000}%
\pgfsetlinewidth{0.501875pt}%
\definecolor{currentstroke}{rgb}{0.200000,0.200000,0.200000}%
\pgfsetstrokecolor{currentstroke}%
\pgfsetdash{}{0pt}%
\pgfpathmoveto{\pgfqpoint{1.664141in}{2.920182in}}%
\pgfpathlineto{\pgfqpoint{1.712355in}{2.874000in}}%
\pgfpathlineto{\pgfqpoint{1.743466in}{2.854099in}}%
\pgfpathlineto{\pgfqpoint{1.695243in}{2.900311in}}%
\pgfpathlineto{\pgfqpoint{1.664141in}{2.920182in}}%
\pgfpathclose%
\pgfusepath{stroke,fill}%
\end{pgfscope}%
\begin{pgfscope}%
\pgfpathrectangle{\pgfqpoint{0.050000in}{0.050000in}}{\pgfqpoint{4.900000in}{4.900000in}}%
\pgfusepath{clip}%
\pgfsetbuttcap%
\pgfsetroundjoin%
\definecolor{currentfill}{rgb}{0.662607,0.662607,0.662607}%
\pgfsetfillcolor{currentfill}%
\pgfsetfillopacity{0.900000}%
\pgfsetlinewidth{0.501875pt}%
\definecolor{currentstroke}{rgb}{0.200000,0.200000,0.200000}%
\pgfsetstrokecolor{currentstroke}%
\pgfsetdash{}{0pt}%
\pgfpathmoveto{\pgfqpoint{2.392322in}{2.331690in}}%
\pgfpathlineto{\pgfqpoint{2.438974in}{2.286812in}}%
\pgfpathlineto{\pgfqpoint{2.471230in}{2.266005in}}%
\pgfpathlineto{\pgfqpoint{2.424574in}{2.310900in}}%
\pgfpathlineto{\pgfqpoint{2.392322in}{2.331690in}}%
\pgfpathclose%
\pgfusepath{stroke,fill}%
\end{pgfscope}%
\begin{pgfscope}%
\pgfpathrectangle{\pgfqpoint{0.050000in}{0.050000in}}{\pgfqpoint{4.900000in}{4.900000in}}%
\pgfusepath{clip}%
\pgfsetbuttcap%
\pgfsetroundjoin%
\definecolor{currentfill}{rgb}{0.895548,0.895548,0.895548}%
\pgfsetfillcolor{currentfill}%
\pgfsetfillopacity{0.900000}%
\pgfsetlinewidth{0.501875pt}%
\definecolor{currentstroke}{rgb}{0.200000,0.200000,0.200000}%
\pgfsetstrokecolor{currentstroke}%
\pgfsetdash{}{0pt}%
\pgfpathmoveto{\pgfqpoint{3.199367in}{1.701907in}}%
\pgfpathlineto{\pgfqpoint{3.244401in}{1.668201in}}%
\pgfpathlineto{\pgfqpoint{3.277941in}{1.647631in}}%
\pgfpathlineto{\pgfqpoint{3.232900in}{1.681137in}}%
\pgfpathlineto{\pgfqpoint{3.199367in}{1.701907in}}%
\pgfpathclose%
\pgfusepath{stroke,fill}%
\end{pgfscope}%
\begin{pgfscope}%
\pgfpathrectangle{\pgfqpoint{0.050000in}{0.050000in}}{\pgfqpoint{4.900000in}{4.900000in}}%
\pgfusepath{clip}%
\pgfsetbuttcap%
\pgfsetroundjoin%
\definecolor{currentfill}{rgb}{0.000000,0.000000,0.000000}%
\pgfsetfillcolor{currentfill}%
\pgfsetfillopacity{0.900000}%
\pgfsetlinewidth{0.501875pt}%
\definecolor{currentstroke}{rgb}{0.200000,0.200000,0.200000}%
\pgfsetstrokecolor{currentstroke}%
\pgfsetdash{}{0pt}%
\pgfpathmoveto{\pgfqpoint{0.666738in}{3.726977in}}%
\pgfpathlineto{\pgfqpoint{0.717101in}{3.679055in}}%
\pgfpathlineto{\pgfqpoint{0.746651in}{3.660416in}}%
\pgfpathlineto{\pgfqpoint{0.696272in}{3.708372in}}%
\pgfpathlineto{\pgfqpoint{0.666738in}{3.726977in}}%
\pgfpathclose%
\pgfusepath{stroke,fill}%
\end{pgfscope}%
\begin{pgfscope}%
\pgfpathrectangle{\pgfqpoint{0.050000in}{0.050000in}}{\pgfqpoint{4.900000in}{4.900000in}}%
\pgfusepath{clip}%
\pgfsetbuttcap%
\pgfsetroundjoin%
\definecolor{currentfill}{rgb}{0.300807,0.300807,0.300807}%
\pgfsetfillcolor{currentfill}%
\pgfsetfillopacity{0.900000}%
\pgfsetlinewidth{0.501875pt}%
\definecolor{currentstroke}{rgb}{0.200000,0.200000,0.200000}%
\pgfsetstrokecolor{currentstroke}%
\pgfsetdash{}{0pt}%
\pgfpathmoveto{\pgfqpoint{1.394927in}{3.138424in}}%
\pgfpathlineto{\pgfqpoint{1.443728in}{3.091765in}}%
\pgfpathlineto{\pgfqpoint{1.474422in}{3.072203in}}%
\pgfpathlineto{\pgfqpoint{1.425611in}{3.118893in}}%
\pgfpathlineto{\pgfqpoint{1.394927in}{3.138424in}}%
\pgfpathclose%
\pgfusepath{stroke,fill}%
\end{pgfscope}%
\begin{pgfscope}%
\pgfpathrectangle{\pgfqpoint{0.050000in}{0.050000in}}{\pgfqpoint{4.900000in}{4.900000in}}%
\pgfusepath{clip}%
\pgfsetbuttcap%
\pgfsetroundjoin%
\definecolor{currentfill}{rgb}{0.560246,0.560246,0.560246}%
\pgfsetfillcolor{currentfill}%
\pgfsetfillopacity{0.900000}%
\pgfsetlinewidth{0.501875pt}%
\definecolor{currentstroke}{rgb}{0.200000,0.200000,0.200000}%
\pgfsetstrokecolor{currentstroke}%
\pgfsetdash{}{0pt}%
\pgfpathmoveto{\pgfqpoint{2.123215in}{2.549811in}}%
\pgfpathlineto{\pgfqpoint{2.170454in}{2.504415in}}%
\pgfpathlineto{\pgfqpoint{2.202292in}{2.483930in}}%
\pgfpathlineto{\pgfqpoint{2.155048in}{2.529353in}}%
\pgfpathlineto{\pgfqpoint{2.123215in}{2.549811in}}%
\pgfpathclose%
\pgfusepath{stroke,fill}%
\end{pgfscope}%
\begin{pgfscope}%
\pgfpathrectangle{\pgfqpoint{0.050000in}{0.050000in}}{\pgfqpoint{4.900000in}{4.900000in}}%
\pgfusepath{clip}%
\pgfsetbuttcap%
\pgfsetroundjoin%
\definecolor{currentfill}{rgb}{0.173472,0.173472,0.173472}%
\pgfsetfillcolor{currentfill}%
\pgfsetfillopacity{0.900000}%
\pgfsetlinewidth{0.501875pt}%
\definecolor{currentstroke}{rgb}{0.200000,0.200000,0.200000}%
\pgfsetstrokecolor{currentstroke}%
\pgfsetdash{}{0pt}%
\pgfpathmoveto{\pgfqpoint{1.125621in}{3.356744in}}%
\pgfpathlineto{\pgfqpoint{1.175010in}{3.309607in}}%
\pgfpathlineto{\pgfqpoint{1.205286in}{3.290383in}}%
\pgfpathlineto{\pgfqpoint{1.155885in}{3.337553in}}%
\pgfpathlineto{\pgfqpoint{1.125621in}{3.356744in}}%
\pgfpathclose%
\pgfusepath{stroke,fill}%
\end{pgfscope}%
\begin{pgfscope}%
\pgfpathrectangle{\pgfqpoint{0.050000in}{0.050000in}}{\pgfqpoint{4.900000in}{4.900000in}}%
\pgfusepath{clip}%
\pgfsetbuttcap%
\pgfsetroundjoin%
\definecolor{currentfill}{rgb}{0.812226,0.812226,0.812226}%
\pgfsetfillcolor{currentfill}%
\pgfsetfillopacity{0.900000}%
\pgfsetlinewidth{0.501875pt}%
\definecolor{currentstroke}{rgb}{0.200000,0.200000,0.200000}%
\pgfsetstrokecolor{currentstroke}%
\pgfsetdash{}{0pt}%
\pgfpathmoveto{\pgfqpoint{2.851617in}{1.964195in}}%
\pgfpathlineto{\pgfqpoint{2.897308in}{1.922124in}}%
\pgfpathlineto{\pgfqpoint{2.930297in}{1.901129in}}%
\pgfpathlineto{\pgfqpoint{2.884604in}{1.943093in}}%
\pgfpathlineto{\pgfqpoint{2.851617in}{1.964195in}}%
\pgfpathclose%
\pgfusepath{stroke,fill}%
\end{pgfscope}%
\begin{pgfscope}%
\pgfpathrectangle{\pgfqpoint{0.050000in}{0.050000in}}{\pgfqpoint{4.900000in}{4.900000in}}%
\pgfusepath{clip}%
\pgfsetbuttcap%
\pgfsetroundjoin%
\definecolor{currentfill}{rgb}{0.461207,0.461207,0.461207}%
\pgfsetfillcolor{currentfill}%
\pgfsetfillopacity{0.900000}%
\pgfsetlinewidth{0.501875pt}%
\definecolor{currentstroke}{rgb}{0.200000,0.200000,0.200000}%
\pgfsetstrokecolor{currentstroke}%
\pgfsetdash{}{0pt}%
\pgfpathmoveto{\pgfqpoint{1.854015in}{2.768038in}}%
\pgfpathlineto{\pgfqpoint{1.901841in}{2.722164in}}%
\pgfpathlineto{\pgfqpoint{1.933263in}{2.702017in}}%
\pgfpathlineto{\pgfqpoint{1.885429in}{2.747919in}}%
\pgfpathlineto{\pgfqpoint{1.854015in}{2.768038in}}%
\pgfpathclose%
\pgfusepath{stroke,fill}%
\end{pgfscope}%
\begin{pgfscope}%
\pgfpathrectangle{\pgfqpoint{0.050000in}{0.050000in}}{\pgfqpoint{4.900000in}{4.900000in}}%
\pgfusepath{clip}%
\pgfsetbuttcap%
\pgfsetroundjoin%
\definecolor{currentfill}{rgb}{0.734579,0.734579,0.734579}%
\pgfsetfillcolor{currentfill}%
\pgfsetfillopacity{0.900000}%
\pgfsetlinewidth{0.501875pt}%
\definecolor{currentstroke}{rgb}{0.200000,0.200000,0.200000}%
\pgfsetstrokecolor{currentstroke}%
\pgfsetdash{}{0pt}%
\pgfpathmoveto{\pgfqpoint{2.582508in}{2.179550in}}%
\pgfpathlineto{\pgfqpoint{2.628771in}{2.135232in}}%
\pgfpathlineto{\pgfqpoint{2.661339in}{2.114246in}}%
\pgfpathlineto{\pgfqpoint{2.615073in}{2.158551in}}%
\pgfpathlineto{\pgfqpoint{2.582508in}{2.179550in}}%
\pgfpathclose%
\pgfusepath{stroke,fill}%
\end{pgfscope}%
\begin{pgfscope}%
\pgfpathrectangle{\pgfqpoint{0.050000in}{0.050000in}}{\pgfqpoint{4.900000in}{4.900000in}}%
\pgfusepath{clip}%
\pgfsetbuttcap%
\pgfsetroundjoin%
\definecolor{currentfill}{rgb}{0.878570,0.878570,0.878570}%
\pgfsetfillcolor{currentfill}%
\pgfsetfillopacity{0.900000}%
\pgfsetlinewidth{0.501875pt}%
\definecolor{currentstroke}{rgb}{0.200000,0.200000,0.200000}%
\pgfsetstrokecolor{currentstroke}%
\pgfsetdash{}{0pt}%
\pgfpathmoveto{\pgfqpoint{3.120761in}{1.759033in}}%
\pgfpathlineto{\pgfqpoint{3.165942in}{1.722618in}}%
\pgfpathlineto{\pgfqpoint{3.199367in}{1.701907in}}%
\pgfpathlineto{\pgfqpoint{3.154181in}{1.738131in}}%
\pgfpathlineto{\pgfqpoint{3.120761in}{1.759033in}}%
\pgfpathclose%
\pgfusepath{stroke,fill}%
\end{pgfscope}%
\begin{pgfscope}%
\pgfpathrectangle{\pgfqpoint{0.050000in}{0.050000in}}{\pgfqpoint{4.900000in}{4.900000in}}%
\pgfusepath{clip}%
\pgfsetbuttcap%
\pgfsetroundjoin%
\definecolor{currentfill}{rgb}{0.068281,0.068281,0.068281}%
\pgfsetfillcolor{currentfill}%
\pgfsetfillopacity{0.900000}%
\pgfsetlinewidth{0.501875pt}%
\definecolor{currentstroke}{rgb}{0.200000,0.200000,0.200000}%
\pgfsetstrokecolor{currentstroke}%
\pgfsetdash{}{0pt}%
\pgfpathmoveto{\pgfqpoint{0.856223in}{3.575141in}}%
\pgfpathlineto{\pgfqpoint{0.906199in}{3.527526in}}%
\pgfpathlineto{\pgfqpoint{0.936058in}{3.508641in}}%
\pgfpathlineto{\pgfqpoint{0.886067in}{3.556290in}}%
\pgfpathlineto{\pgfqpoint{0.856223in}{3.575141in}}%
\pgfpathclose%
\pgfusepath{stroke,fill}%
\end{pgfscope}%
\begin{pgfscope}%
\pgfpathrectangle{\pgfqpoint{0.050000in}{0.050000in}}{\pgfqpoint{4.900000in}{4.900000in}}%
\pgfusepath{clip}%
\pgfsetbuttcap%
\pgfsetroundjoin%
\definecolor{currentfill}{rgb}{0.964937,0.964937,0.964937}%
\pgfsetfillcolor{currentfill}%
\pgfsetfillopacity{0.900000}%
\pgfsetlinewidth{0.501875pt}%
\definecolor{currentstroke}{rgb}{0.200000,0.200000,0.200000}%
\pgfsetstrokecolor{currentstroke}%
\pgfsetdash{}{0pt}%
\pgfpathmoveto{\pgfqpoint{3.705197in}{1.408503in}}%
\pgfpathlineto{\pgfqpoint{3.749682in}{1.395995in}}%
\pgfpathlineto{\pgfqpoint{3.784066in}{1.375744in}}%
\pgfpathlineto{\pgfqpoint{3.739564in}{1.388188in}}%
\pgfpathlineto{\pgfqpoint{3.705197in}{1.408503in}}%
\pgfpathclose%
\pgfusepath{stroke,fill}%
\end{pgfscope}%
\begin{pgfscope}%
\pgfpathrectangle{\pgfqpoint{0.050000in}{0.050000in}}{\pgfqpoint{4.900000in}{4.900000in}}%
\pgfusepath{clip}%
\pgfsetbuttcap%
\pgfsetroundjoin%
\definecolor{currentfill}{rgb}{0.957555,0.957555,0.957555}%
\pgfsetfillcolor{currentfill}%
\pgfsetfillopacity{0.900000}%
\pgfsetlinewidth{0.501875pt}%
\definecolor{currentstroke}{rgb}{0.200000,0.200000,0.200000}%
\pgfsetstrokecolor{currentstroke}%
\pgfsetdash{}{0pt}%
\pgfpathmoveto{\pgfqpoint{3.626397in}{1.444839in}}%
\pgfpathlineto{\pgfqpoint{3.670940in}{1.428738in}}%
\pgfpathlineto{\pgfqpoint{3.705197in}{1.408503in}}%
\pgfpathlineto{\pgfqpoint{3.660637in}{1.424507in}}%
\pgfpathlineto{\pgfqpoint{3.626397in}{1.444839in}}%
\pgfpathclose%
\pgfusepath{stroke,fill}%
\end{pgfscope}%
\begin{pgfscope}%
\pgfpathrectangle{\pgfqpoint{0.050000in}{0.050000in}}{\pgfqpoint{4.900000in}{4.900000in}}%
\pgfusepath{clip}%
\pgfsetbuttcap%
\pgfsetroundjoin%
\definecolor{currentfill}{rgb}{0.970473,0.970473,0.970473}%
\pgfsetfillcolor{currentfill}%
\pgfsetfillopacity{0.900000}%
\pgfsetlinewidth{0.501875pt}%
\definecolor{currentstroke}{rgb}{0.200000,0.200000,0.200000}%
\pgfsetstrokecolor{currentstroke}%
\pgfsetdash{}{0pt}%
\pgfpathmoveto{\pgfqpoint{3.784066in}{1.375744in}}%
\pgfpathlineto{\pgfqpoint{3.828501in}{1.366664in}}%
\pgfpathlineto{\pgfqpoint{3.863014in}{1.346367in}}%
\pgfpathlineto{\pgfqpoint{3.818560in}{1.355412in}}%
\pgfpathlineto{\pgfqpoint{3.784066in}{1.375744in}}%
\pgfpathclose%
\pgfusepath{stroke,fill}%
\end{pgfscope}%
\begin{pgfscope}%
\pgfpathrectangle{\pgfqpoint{0.050000in}{0.050000in}}{\pgfqpoint{4.900000in}{4.900000in}}%
\pgfusepath{clip}%
\pgfsetbuttcap%
\pgfsetroundjoin%
\definecolor{currentfill}{rgb}{0.367243,0.367243,0.367243}%
\pgfsetfillcolor{currentfill}%
\pgfsetfillopacity{0.900000}%
\pgfsetlinewidth{0.501875pt}%
\definecolor{currentstroke}{rgb}{0.200000,0.200000,0.200000}%
\pgfsetstrokecolor{currentstroke}%
\pgfsetdash{}{0pt}%
\pgfpathmoveto{\pgfqpoint{1.584723in}{2.986342in}}%
\pgfpathlineto{\pgfqpoint{1.633136in}{2.939991in}}%
\pgfpathlineto{\pgfqpoint{1.664141in}{2.920182in}}%
\pgfpathlineto{\pgfqpoint{1.615719in}{2.966563in}}%
\pgfpathlineto{\pgfqpoint{1.584723in}{2.986342in}}%
\pgfpathclose%
\pgfusepath{stroke,fill}%
\end{pgfscope}%
\begin{pgfscope}%
\pgfpathrectangle{\pgfqpoint{0.050000in}{0.050000in}}{\pgfqpoint{4.900000in}{4.900000in}}%
\pgfusepath{clip}%
\pgfsetbuttcap%
\pgfsetroundjoin%
\definecolor{currentfill}{rgb}{0.629020,0.629020,0.629020}%
\pgfsetfillcolor{currentfill}%
\pgfsetfillopacity{0.900000}%
\pgfsetlinewidth{0.501875pt}%
\definecolor{currentstroke}{rgb}{0.200000,0.200000,0.200000}%
\pgfsetstrokecolor{currentstroke}%
\pgfsetdash{}{0pt}%
\pgfpathmoveto{\pgfqpoint{2.313323in}{2.397490in}}%
\pgfpathlineto{\pgfqpoint{2.360172in}{2.352416in}}%
\pgfpathlineto{\pgfqpoint{2.392322in}{2.331690in}}%
\pgfpathlineto{\pgfqpoint{2.345469in}{2.376786in}}%
\pgfpathlineto{\pgfqpoint{2.313323in}{2.397490in}}%
\pgfpathclose%
\pgfusepath{stroke,fill}%
\end{pgfscope}%
\begin{pgfscope}%
\pgfpathrectangle{\pgfqpoint{0.050000in}{0.050000in}}{\pgfqpoint{4.900000in}{4.900000in}}%
\pgfusepath{clip}%
\pgfsetbuttcap%
\pgfsetroundjoin%
\definecolor{currentfill}{rgb}{0.950173,0.950173,0.950173}%
\pgfsetfillcolor{currentfill}%
\pgfsetfillopacity{0.900000}%
\pgfsetlinewidth{0.501875pt}%
\definecolor{currentstroke}{rgb}{0.200000,0.200000,0.200000}%
\pgfsetstrokecolor{currentstroke}%
\pgfsetdash{}{0pt}%
\pgfpathmoveto{\pgfqpoint{3.547652in}{1.484874in}}%
\pgfpathlineto{\pgfqpoint{3.592266in}{1.465092in}}%
\pgfpathlineto{\pgfqpoint{3.626397in}{1.444839in}}%
\pgfpathlineto{\pgfqpoint{3.581768in}{1.464492in}}%
\pgfpathlineto{\pgfqpoint{3.547652in}{1.484874in}}%
\pgfpathclose%
\pgfusepath{stroke,fill}%
\end{pgfscope}%
\begin{pgfscope}%
\pgfpathrectangle{\pgfqpoint{0.050000in}{0.050000in}}{\pgfqpoint{4.900000in}{4.900000in}}%
\pgfusepath{clip}%
\pgfsetbuttcap%
\pgfsetroundjoin%
\definecolor{currentfill}{rgb}{0.976009,0.976009,0.976009}%
\pgfsetfillcolor{currentfill}%
\pgfsetfillopacity{0.900000}%
\pgfsetlinewidth{0.501875pt}%
\definecolor{currentstroke}{rgb}{0.200000,0.200000,0.200000}%
\pgfsetstrokecolor{currentstroke}%
\pgfsetdash{}{0pt}%
\pgfpathmoveto{\pgfqpoint{3.863014in}{1.346367in}}%
\pgfpathlineto{\pgfqpoint{3.907406in}{1.340493in}}%
\pgfpathlineto{\pgfqpoint{3.942050in}{1.320122in}}%
\pgfpathlineto{\pgfqpoint{3.897638in}{1.325989in}}%
\pgfpathlineto{\pgfqpoint{3.863014in}{1.346367in}}%
\pgfpathclose%
\pgfusepath{stroke,fill}%
\end{pgfscope}%
\begin{pgfscope}%
\pgfpathrectangle{\pgfqpoint{0.050000in}{0.050000in}}{\pgfqpoint{4.900000in}{4.900000in}}%
\pgfusepath{clip}%
\pgfsetbuttcap%
\pgfsetroundjoin%
\definecolor{currentfill}{rgb}{0.262053,0.262053,0.262053}%
\pgfsetfillcolor{currentfill}%
\pgfsetfillopacity{0.900000}%
\pgfsetlinewidth{0.501875pt}%
\definecolor{currentstroke}{rgb}{0.200000,0.200000,0.200000}%
\pgfsetstrokecolor{currentstroke}%
\pgfsetdash{}{0pt}%
\pgfpathmoveto{\pgfqpoint{1.315339in}{3.204723in}}%
\pgfpathlineto{\pgfqpoint{1.364340in}{3.157894in}}%
\pgfpathlineto{\pgfqpoint{1.394927in}{3.138424in}}%
\pgfpathlineto{\pgfqpoint{1.345915in}{3.185285in}}%
\pgfpathlineto{\pgfqpoint{1.315339in}{3.204723in}}%
\pgfpathclose%
\pgfusepath{stroke,fill}%
\end{pgfscope}%
\begin{pgfscope}%
\pgfpathrectangle{\pgfqpoint{0.050000in}{0.050000in}}{\pgfqpoint{4.900000in}{4.900000in}}%
\pgfusepath{clip}%
\pgfsetbuttcap%
\pgfsetroundjoin%
\definecolor{currentfill}{rgb}{0.942791,0.942791,0.942791}%
\pgfsetfillcolor{currentfill}%
\pgfsetfillopacity{0.900000}%
\pgfsetlinewidth{0.501875pt}%
\definecolor{currentstroke}{rgb}{0.200000,0.200000,0.200000}%
\pgfsetstrokecolor{currentstroke}%
\pgfsetdash{}{0pt}%
\pgfpathmoveto{\pgfqpoint{3.468947in}{1.528642in}}%
\pgfpathlineto{\pgfqpoint{3.513645in}{1.505180in}}%
\pgfpathlineto{\pgfqpoint{3.547652in}{1.484874in}}%
\pgfpathlineto{\pgfqpoint{3.502940in}{1.508178in}}%
\pgfpathlineto{\pgfqpoint{3.468947in}{1.528642in}}%
\pgfpathclose%
\pgfusepath{stroke,fill}%
\end{pgfscope}%
\begin{pgfscope}%
\pgfpathrectangle{\pgfqpoint{0.050000in}{0.050000in}}{\pgfqpoint{4.900000in}{4.900000in}}%
\pgfusepath{clip}%
\pgfsetbuttcap%
\pgfsetroundjoin%
\definecolor{currentfill}{rgb}{0.525798,0.525798,0.525798}%
\pgfsetfillcolor{currentfill}%
\pgfsetfillopacity{0.900000}%
\pgfsetlinewidth{0.501875pt}%
\definecolor{currentstroke}{rgb}{0.200000,0.200000,0.200000}%
\pgfsetstrokecolor{currentstroke}%
\pgfsetdash{}{0pt}%
\pgfpathmoveto{\pgfqpoint{2.044045in}{2.615770in}}%
\pgfpathlineto{\pgfqpoint{2.091482in}{2.570205in}}%
\pgfpathlineto{\pgfqpoint{2.123215in}{2.549811in}}%
\pgfpathlineto{\pgfqpoint{2.075772in}{2.595404in}}%
\pgfpathlineto{\pgfqpoint{2.044045in}{2.615770in}}%
\pgfpathclose%
\pgfusepath{stroke,fill}%
\end{pgfscope}%
\begin{pgfscope}%
\pgfpathrectangle{\pgfqpoint{0.050000in}{0.050000in}}{\pgfqpoint{4.900000in}{4.900000in}}%
\pgfusepath{clip}%
\pgfsetbuttcap%
\pgfsetroundjoin%
\definecolor{currentfill}{rgb}{0.136563,0.136563,0.136563}%
\pgfsetfillcolor{currentfill}%
\pgfsetfillopacity{0.900000}%
\pgfsetlinewidth{0.501875pt}%
\definecolor{currentstroke}{rgb}{0.200000,0.200000,0.200000}%
\pgfsetstrokecolor{currentstroke}%
\pgfsetdash{}{0pt}%
\pgfpathmoveto{\pgfqpoint{1.045863in}{3.423182in}}%
\pgfpathlineto{\pgfqpoint{1.095452in}{3.375875in}}%
\pgfpathlineto{\pgfqpoint{1.125621in}{3.356744in}}%
\pgfpathlineto{\pgfqpoint{1.076019in}{3.404084in}}%
\pgfpathlineto{\pgfqpoint{1.045863in}{3.423182in}}%
\pgfpathclose%
\pgfusepath{stroke,fill}%
\end{pgfscope}%
\begin{pgfscope}%
\pgfpathrectangle{\pgfqpoint{0.050000in}{0.050000in}}{\pgfqpoint{4.900000in}{4.900000in}}%
\pgfusepath{clip}%
\pgfsetbuttcap%
\pgfsetroundjoin%
\definecolor{currentfill}{rgb}{0.788112,0.788112,0.788112}%
\pgfsetfillcolor{currentfill}%
\pgfsetfillopacity{0.900000}%
\pgfsetlinewidth{0.501875pt}%
\definecolor{currentstroke}{rgb}{0.200000,0.200000,0.200000}%
\pgfsetstrokecolor{currentstroke}%
\pgfsetdash{}{0pt}%
\pgfpathmoveto{\pgfqpoint{2.772856in}{2.028332in}}%
\pgfpathlineto{\pgfqpoint{2.818735in}{1.985243in}}%
\pgfpathlineto{\pgfqpoint{2.851617in}{1.964195in}}%
\pgfpathlineto{\pgfqpoint{2.805736in}{2.007213in}}%
\pgfpathlineto{\pgfqpoint{2.772856in}{2.028332in}}%
\pgfpathclose%
\pgfusepath{stroke,fill}%
\end{pgfscope}%
\begin{pgfscope}%
\pgfpathrectangle{\pgfqpoint{0.050000in}{0.050000in}}{\pgfqpoint{4.900000in}{4.900000in}}%
\pgfusepath{clip}%
\pgfsetbuttcap%
\pgfsetroundjoin%
\definecolor{currentfill}{rgb}{0.858762,0.858762,0.858762}%
\pgfsetfillcolor{currentfill}%
\pgfsetfillopacity{0.900000}%
\pgfsetlinewidth{0.501875pt}%
\definecolor{currentstroke}{rgb}{0.200000,0.200000,0.200000}%
\pgfsetstrokecolor{currentstroke}%
\pgfsetdash{}{0pt}%
\pgfpathmoveto{\pgfqpoint{3.042107in}{1.818577in}}%
\pgfpathlineto{\pgfqpoint{3.087448in}{1.779878in}}%
\pgfpathlineto{\pgfqpoint{3.120761in}{1.759033in}}%
\pgfpathlineto{\pgfqpoint{3.075415in}{1.797561in}}%
\pgfpathlineto{\pgfqpoint{3.042107in}{1.818577in}}%
\pgfpathclose%
\pgfusepath{stroke,fill}%
\end{pgfscope}%
\begin{pgfscope}%
\pgfpathrectangle{\pgfqpoint{0.050000in}{0.050000in}}{\pgfqpoint{4.900000in}{4.900000in}}%
\pgfusepath{clip}%
\pgfsetbuttcap%
\pgfsetroundjoin%
\definecolor{currentfill}{rgb}{0.428143,0.428143,0.428143}%
\pgfsetfillcolor{currentfill}%
\pgfsetfillopacity{0.900000}%
\pgfsetlinewidth{0.501875pt}%
\definecolor{currentstroke}{rgb}{0.200000,0.200000,0.200000}%
\pgfsetstrokecolor{currentstroke}%
\pgfsetdash{}{0pt}%
\pgfpathmoveto{\pgfqpoint{1.774675in}{2.834136in}}%
\pgfpathlineto{\pgfqpoint{1.822700in}{2.788093in}}%
\pgfpathlineto{\pgfqpoint{1.854015in}{2.768038in}}%
\pgfpathlineto{\pgfqpoint{1.805983in}{2.814110in}}%
\pgfpathlineto{\pgfqpoint{1.774675in}{2.834136in}}%
\pgfpathclose%
\pgfusepath{stroke,fill}%
\end{pgfscope}%
\begin{pgfscope}%
\pgfpathrectangle{\pgfqpoint{0.050000in}{0.050000in}}{\pgfqpoint{4.900000in}{4.900000in}}%
\pgfusepath{clip}%
\pgfsetbuttcap%
\pgfsetroundjoin%
\definecolor{currentfill}{rgb}{0.929504,0.929504,0.929504}%
\pgfsetfillcolor{currentfill}%
\pgfsetfillopacity{0.900000}%
\pgfsetlinewidth{0.501875pt}%
\definecolor{currentstroke}{rgb}{0.200000,0.200000,0.200000}%
\pgfsetstrokecolor{currentstroke}%
\pgfsetdash{}{0pt}%
\pgfpathmoveto{\pgfqpoint{3.390265in}{1.576075in}}%
\pgfpathlineto{\pgfqpoint{3.435062in}{1.549033in}}%
\pgfpathlineto{\pgfqpoint{3.468947in}{1.528642in}}%
\pgfpathlineto{\pgfqpoint{3.424139in}{1.555503in}}%
\pgfpathlineto{\pgfqpoint{3.390265in}{1.576075in}}%
\pgfpathclose%
\pgfusepath{stroke,fill}%
\end{pgfscope}%
\begin{pgfscope}%
\pgfpathrectangle{\pgfqpoint{0.050000in}{0.050000in}}{\pgfqpoint{4.900000in}{4.900000in}}%
\pgfusepath{clip}%
\pgfsetbuttcap%
\pgfsetroundjoin%
\definecolor{currentfill}{rgb}{0.700992,0.700992,0.700992}%
\pgfsetfillcolor{currentfill}%
\pgfsetfillopacity{0.900000}%
\pgfsetlinewidth{0.501875pt}%
\definecolor{currentstroke}{rgb}{0.200000,0.200000,0.200000}%
\pgfsetstrokecolor{currentstroke}%
\pgfsetdash{}{0pt}%
\pgfpathmoveto{\pgfqpoint{2.503587in}{2.245135in}}%
\pgfpathlineto{\pgfqpoint{2.550046in}{2.200487in}}%
\pgfpathlineto{\pgfqpoint{2.582508in}{2.179550in}}%
\pgfpathlineto{\pgfqpoint{2.536046in}{2.224202in}}%
\pgfpathlineto{\pgfqpoint{2.503587in}{2.245135in}}%
\pgfpathclose%
\pgfusepath{stroke,fill}%
\end{pgfscope}%
\begin{pgfscope}%
\pgfpathrectangle{\pgfqpoint{0.050000in}{0.050000in}}{\pgfqpoint{4.900000in}{4.900000in}}%
\pgfusepath{clip}%
\pgfsetbuttcap%
\pgfsetroundjoin%
\definecolor{currentfill}{rgb}{0.031865,0.031865,0.031865}%
\pgfsetfillcolor{currentfill}%
\pgfsetfillopacity{0.900000}%
\pgfsetlinewidth{0.501875pt}%
\definecolor{currentstroke}{rgb}{0.200000,0.200000,0.200000}%
\pgfsetstrokecolor{currentstroke}%
\pgfsetdash{}{0pt}%
\pgfpathmoveto{\pgfqpoint{0.776294in}{3.641718in}}%
\pgfpathlineto{\pgfqpoint{0.826471in}{3.593932in}}%
\pgfpathlineto{\pgfqpoint{0.856223in}{3.575141in}}%
\pgfpathlineto{\pgfqpoint{0.806030in}{3.622961in}}%
\pgfpathlineto{\pgfqpoint{0.776294in}{3.641718in}}%
\pgfpathclose%
\pgfusepath{stroke,fill}%
\end{pgfscope}%
\begin{pgfscope}%
\pgfpathrectangle{\pgfqpoint{0.050000in}{0.050000in}}{\pgfqpoint{4.900000in}{4.900000in}}%
\pgfusepath{clip}%
\pgfsetbuttcap%
\pgfsetroundjoin%
\definecolor{currentfill}{rgb}{0.338824,0.338824,0.338824}%
\pgfsetfillcolor{currentfill}%
\pgfsetfillopacity{0.900000}%
\pgfsetlinewidth{0.501875pt}%
\definecolor{currentstroke}{rgb}{0.200000,0.200000,0.200000}%
\pgfsetstrokecolor{currentstroke}%
\pgfsetdash{}{0pt}%
\pgfpathmoveto{\pgfqpoint{1.505213in}{3.052579in}}%
\pgfpathlineto{\pgfqpoint{1.553826in}{3.006058in}}%
\pgfpathlineto{\pgfqpoint{1.584723in}{2.986342in}}%
\pgfpathlineto{\pgfqpoint{1.536101in}{3.032893in}}%
\pgfpathlineto{\pgfqpoint{1.505213in}{3.052579in}}%
\pgfpathclose%
\pgfusepath{stroke,fill}%
\end{pgfscope}%
\begin{pgfscope}%
\pgfpathrectangle{\pgfqpoint{0.050000in}{0.050000in}}{\pgfqpoint{4.900000in}{4.900000in}}%
\pgfusepath{clip}%
\pgfsetbuttcap%
\pgfsetroundjoin%
\definecolor{currentfill}{rgb}{0.590634,0.590634,0.590634}%
\pgfsetfillcolor{currentfill}%
\pgfsetfillopacity{0.900000}%
\pgfsetlinewidth{0.501875pt}%
\definecolor{currentstroke}{rgb}{0.200000,0.200000,0.200000}%
\pgfsetstrokecolor{currentstroke}%
\pgfsetdash{}{0pt}%
\pgfpathmoveto{\pgfqpoint{2.234231in}{2.463381in}}%
\pgfpathlineto{\pgfqpoint{2.281278in}{2.418129in}}%
\pgfpathlineto{\pgfqpoint{2.313323in}{2.397490in}}%
\pgfpathlineto{\pgfqpoint{2.266271in}{2.442767in}}%
\pgfpathlineto{\pgfqpoint{2.234231in}{2.463381in}}%
\pgfpathclose%
\pgfusepath{stroke,fill}%
\end{pgfscope}%
\begin{pgfscope}%
\pgfpathrectangle{\pgfqpoint{0.050000in}{0.050000in}}{\pgfqpoint{4.900000in}{4.900000in}}%
\pgfusepath{clip}%
\pgfsetbuttcap%
\pgfsetroundjoin%
\definecolor{currentfill}{rgb}{0.912526,0.912526,0.912526}%
\pgfsetfillcolor{currentfill}%
\pgfsetfillopacity{0.900000}%
\pgfsetlinewidth{0.501875pt}%
\definecolor{currentstroke}{rgb}{0.200000,0.200000,0.200000}%
\pgfsetstrokecolor{currentstroke}%
\pgfsetdash{}{0pt}%
\pgfpathmoveto{\pgfqpoint{3.311589in}{1.626998in}}%
\pgfpathlineto{\pgfqpoint{3.356500in}{1.596578in}}%
\pgfpathlineto{\pgfqpoint{3.390265in}{1.576075in}}%
\pgfpathlineto{\pgfqpoint{3.345344in}{1.606300in}}%
\pgfpathlineto{\pgfqpoint{3.311589in}{1.626998in}}%
\pgfpathclose%
\pgfusepath{stroke,fill}%
\end{pgfscope}%
\begin{pgfscope}%
\pgfpathrectangle{\pgfqpoint{0.050000in}{0.050000in}}{\pgfqpoint{4.900000in}{4.900000in}}%
\pgfusepath{clip}%
\pgfsetbuttcap%
\pgfsetroundjoin%
\definecolor{currentfill}{rgb}{0.217762,0.217762,0.217762}%
\pgfsetfillcolor{currentfill}%
\pgfsetfillopacity{0.900000}%
\pgfsetlinewidth{0.501875pt}%
\definecolor{currentstroke}{rgb}{0.200000,0.200000,0.200000}%
\pgfsetstrokecolor{currentstroke}%
\pgfsetdash{}{0pt}%
\pgfpathmoveto{\pgfqpoint{1.235659in}{3.271099in}}%
\pgfpathlineto{\pgfqpoint{1.284859in}{3.224100in}}%
\pgfpathlineto{\pgfqpoint{1.315339in}{3.204723in}}%
\pgfpathlineto{\pgfqpoint{1.266127in}{3.251754in}}%
\pgfpathlineto{\pgfqpoint{1.235659in}{3.271099in}}%
\pgfpathclose%
\pgfusepath{stroke,fill}%
\end{pgfscope}%
\begin{pgfscope}%
\pgfpathrectangle{\pgfqpoint{0.050000in}{0.050000in}}{\pgfqpoint{4.900000in}{4.900000in}}%
\pgfusepath{clip}%
\pgfsetbuttcap%
\pgfsetroundjoin%
\definecolor{currentfill}{rgb}{0.491349,0.491349,0.491349}%
\pgfsetfillcolor{currentfill}%
\pgfsetfillopacity{0.900000}%
\pgfsetlinewidth{0.501875pt}%
\definecolor{currentstroke}{rgb}{0.200000,0.200000,0.200000}%
\pgfsetstrokecolor{currentstroke}%
\pgfsetdash{}{0pt}%
\pgfpathmoveto{\pgfqpoint{1.964783in}{2.681806in}}%
\pgfpathlineto{\pgfqpoint{2.012418in}{2.636073in}}%
\pgfpathlineto{\pgfqpoint{2.044045in}{2.615770in}}%
\pgfpathlineto{\pgfqpoint{1.996404in}{2.661532in}}%
\pgfpathlineto{\pgfqpoint{1.964783in}{2.681806in}}%
\pgfpathclose%
\pgfusepath{stroke,fill}%
\end{pgfscope}%
\begin{pgfscope}%
\pgfpathrectangle{\pgfqpoint{0.050000in}{0.050000in}}{\pgfqpoint{4.900000in}{4.900000in}}%
\pgfusepath{clip}%
\pgfsetbuttcap%
\pgfsetroundjoin%
\definecolor{currentfill}{rgb}{0.839785,0.839785,0.839785}%
\pgfsetfillcolor{currentfill}%
\pgfsetfillopacity{0.900000}%
\pgfsetlinewidth{0.501875pt}%
\definecolor{currentstroke}{rgb}{0.200000,0.200000,0.200000}%
\pgfsetstrokecolor{currentstroke}%
\pgfsetdash{}{0pt}%
\pgfpathmoveto{\pgfqpoint{2.963391in}{1.880080in}}%
\pgfpathlineto{\pgfqpoint{3.008904in}{1.839538in}}%
\pgfpathlineto{\pgfqpoint{3.042107in}{1.818577in}}%
\pgfpathlineto{\pgfqpoint{2.996590in}{1.858977in}}%
\pgfpathlineto{\pgfqpoint{2.963391in}{1.880080in}}%
\pgfpathclose%
\pgfusepath{stroke,fill}%
\end{pgfscope}%
\begin{pgfscope}%
\pgfpathrectangle{\pgfqpoint{0.050000in}{0.050000in}}{\pgfqpoint{4.900000in}{4.900000in}}%
\pgfusepath{clip}%
\pgfsetbuttcap%
\pgfsetroundjoin%
\definecolor{currentfill}{rgb}{0.763998,0.763998,0.763998}%
\pgfsetfillcolor{currentfill}%
\pgfsetfillopacity{0.900000}%
\pgfsetlinewidth{0.501875pt}%
\definecolor{currentstroke}{rgb}{0.200000,0.200000,0.200000}%
\pgfsetstrokecolor{currentstroke}%
\pgfsetdash{}{0pt}%
\pgfpathmoveto{\pgfqpoint{2.694009in}{2.093202in}}%
\pgfpathlineto{\pgfqpoint{2.740080in}{2.049395in}}%
\pgfpathlineto{\pgfqpoint{2.772856in}{2.028332in}}%
\pgfpathlineto{\pgfqpoint{2.726783in}{2.072099in}}%
\pgfpathlineto{\pgfqpoint{2.694009in}{2.093202in}}%
\pgfpathclose%
\pgfusepath{stroke,fill}%
\end{pgfscope}%
\begin{pgfscope}%
\pgfpathrectangle{\pgfqpoint{0.050000in}{0.050000in}}{\pgfqpoint{4.900000in}{4.900000in}}%
\pgfusepath{clip}%
\pgfsetbuttcap%
\pgfsetroundjoin%
\definecolor{currentfill}{rgb}{0.100146,0.100146,0.100146}%
\pgfsetfillcolor{currentfill}%
\pgfsetfillopacity{0.900000}%
\pgfsetlinewidth{0.501875pt}%
\definecolor{currentstroke}{rgb}{0.200000,0.200000,0.200000}%
\pgfsetstrokecolor{currentstroke}%
\pgfsetdash{}{0pt}%
\pgfpathmoveto{\pgfqpoint{0.966012in}{3.489697in}}%
\pgfpathlineto{\pgfqpoint{1.015801in}{3.442219in}}%
\pgfpathlineto{\pgfqpoint{1.045863in}{3.423182in}}%
\pgfpathlineto{\pgfqpoint{0.996059in}{3.470693in}}%
\pgfpathlineto{\pgfqpoint{0.966012in}{3.489697in}}%
\pgfpathclose%
\pgfusepath{stroke,fill}%
\end{pgfscope}%
\begin{pgfscope}%
\pgfpathrectangle{\pgfqpoint{0.050000in}{0.050000in}}{\pgfqpoint{4.900000in}{4.900000in}}%
\pgfusepath{clip}%
\pgfsetbuttcap%
\pgfsetroundjoin%
\definecolor{currentfill}{rgb}{0.399723,0.399723,0.399723}%
\pgfsetfillcolor{currentfill}%
\pgfsetfillopacity{0.900000}%
\pgfsetlinewidth{0.501875pt}%
\definecolor{currentstroke}{rgb}{0.200000,0.200000,0.200000}%
\pgfsetstrokecolor{currentstroke}%
\pgfsetdash{}{0pt}%
\pgfpathmoveto{\pgfqpoint{1.695243in}{2.900311in}}%
\pgfpathlineto{\pgfqpoint{1.743466in}{2.854099in}}%
\pgfpathlineto{\pgfqpoint{1.774675in}{2.834136in}}%
\pgfpathlineto{\pgfqpoint{1.726444in}{2.880377in}}%
\pgfpathlineto{\pgfqpoint{1.695243in}{2.900311in}}%
\pgfpathclose%
\pgfusepath{stroke,fill}%
\end{pgfscope}%
\begin{pgfscope}%
\pgfpathrectangle{\pgfqpoint{0.050000in}{0.050000in}}{\pgfqpoint{4.900000in}{4.900000in}}%
\pgfusepath{clip}%
\pgfsetbuttcap%
\pgfsetroundjoin%
\definecolor{currentfill}{rgb}{0.662607,0.662607,0.662607}%
\pgfsetfillcolor{currentfill}%
\pgfsetfillopacity{0.900000}%
\pgfsetlinewidth{0.501875pt}%
\definecolor{currentstroke}{rgb}{0.200000,0.200000,0.200000}%
\pgfsetstrokecolor{currentstroke}%
\pgfsetdash{}{0pt}%
\pgfpathmoveto{\pgfqpoint{2.424574in}{2.310900in}}%
\pgfpathlineto{\pgfqpoint{2.471230in}{2.266005in}}%
\pgfpathlineto{\pgfqpoint{2.503587in}{2.245135in}}%
\pgfpathlineto{\pgfqpoint{2.456927in}{2.290045in}}%
\pgfpathlineto{\pgfqpoint{2.424574in}{2.310900in}}%
\pgfpathclose%
\pgfusepath{stroke,fill}%
\end{pgfscope}%
\begin{pgfscope}%
\pgfpathrectangle{\pgfqpoint{0.050000in}{0.050000in}}{\pgfqpoint{4.900000in}{4.900000in}}%
\pgfusepath{clip}%
\pgfsetbuttcap%
\pgfsetroundjoin%
\definecolor{currentfill}{rgb}{0.000000,0.000000,0.000000}%
\pgfsetfillcolor{currentfill}%
\pgfsetfillopacity{0.900000}%
\pgfsetlinewidth{0.501875pt}%
\definecolor{currentstroke}{rgb}{0.200000,0.200000,0.200000}%
\pgfsetstrokecolor{currentstroke}%
\pgfsetdash{}{0pt}%
\pgfpathmoveto{\pgfqpoint{0.696272in}{3.708372in}}%
\pgfpathlineto{\pgfqpoint{0.746651in}{3.660416in}}%
\pgfpathlineto{\pgfqpoint{0.776294in}{3.641718in}}%
\pgfpathlineto{\pgfqpoint{0.725899in}{3.689709in}}%
\pgfpathlineto{\pgfqpoint{0.696272in}{3.708372in}}%
\pgfpathclose%
\pgfusepath{stroke,fill}%
\end{pgfscope}%
\begin{pgfscope}%
\pgfpathrectangle{\pgfqpoint{0.050000in}{0.050000in}}{\pgfqpoint{4.900000in}{4.900000in}}%
\pgfusepath{clip}%
\pgfsetbuttcap%
\pgfsetroundjoin%
\definecolor{currentfill}{rgb}{0.895548,0.895548,0.895548}%
\pgfsetfillcolor{currentfill}%
\pgfsetfillopacity{0.900000}%
\pgfsetlinewidth{0.501875pt}%
\definecolor{currentstroke}{rgb}{0.200000,0.200000,0.200000}%
\pgfsetstrokecolor{currentstroke}%
\pgfsetdash{}{0pt}%
\pgfpathmoveto{\pgfqpoint{3.232900in}{1.681137in}}%
\pgfpathlineto{\pgfqpoint{3.277941in}{1.647631in}}%
\pgfpathlineto{\pgfqpoint{3.311589in}{1.626998in}}%
\pgfpathlineto{\pgfqpoint{3.266539in}{1.660306in}}%
\pgfpathlineto{\pgfqpoint{3.232900in}{1.681137in}}%
\pgfpathclose%
\pgfusepath{stroke,fill}%
\end{pgfscope}%
\begin{pgfscope}%
\pgfpathrectangle{\pgfqpoint{0.050000in}{0.050000in}}{\pgfqpoint{4.900000in}{4.900000in}}%
\pgfusepath{clip}%
\pgfsetbuttcap%
\pgfsetroundjoin%
\definecolor{currentfill}{rgb}{0.300807,0.300807,0.300807}%
\pgfsetfillcolor{currentfill}%
\pgfsetfillopacity{0.900000}%
\pgfsetlinewidth{0.501875pt}%
\definecolor{currentstroke}{rgb}{0.200000,0.200000,0.200000}%
\pgfsetstrokecolor{currentstroke}%
\pgfsetdash{}{0pt}%
\pgfpathmoveto{\pgfqpoint{1.425611in}{3.118893in}}%
\pgfpathlineto{\pgfqpoint{1.474422in}{3.072203in}}%
\pgfpathlineto{\pgfqpoint{1.505213in}{3.052579in}}%
\pgfpathlineto{\pgfqpoint{1.456391in}{3.099300in}}%
\pgfpathlineto{\pgfqpoint{1.425611in}{3.118893in}}%
\pgfpathclose%
\pgfusepath{stroke,fill}%
\end{pgfscope}%
\begin{pgfscope}%
\pgfpathrectangle{\pgfqpoint{0.050000in}{0.050000in}}{\pgfqpoint{4.900000in}{4.900000in}}%
\pgfusepath{clip}%
\pgfsetbuttcap%
\pgfsetroundjoin%
\definecolor{currentfill}{rgb}{0.560246,0.560246,0.560246}%
\pgfsetfillcolor{currentfill}%
\pgfsetfillopacity{0.900000}%
\pgfsetlinewidth{0.501875pt}%
\definecolor{currentstroke}{rgb}{0.200000,0.200000,0.200000}%
\pgfsetstrokecolor{currentstroke}%
\pgfsetdash{}{0pt}%
\pgfpathmoveto{\pgfqpoint{2.155048in}{2.529353in}}%
\pgfpathlineto{\pgfqpoint{2.202292in}{2.483930in}}%
\pgfpathlineto{\pgfqpoint{2.234231in}{2.463381in}}%
\pgfpathlineto{\pgfqpoint{2.186981in}{2.508831in}}%
\pgfpathlineto{\pgfqpoint{2.155048in}{2.529353in}}%
\pgfpathclose%
\pgfusepath{stroke,fill}%
\end{pgfscope}%
\begin{pgfscope}%
\pgfpathrectangle{\pgfqpoint{0.050000in}{0.050000in}}{\pgfqpoint{4.900000in}{4.900000in}}%
\pgfusepath{clip}%
\pgfsetbuttcap%
\pgfsetroundjoin%
\definecolor{currentfill}{rgb}{0.173472,0.173472,0.173472}%
\pgfsetfillcolor{currentfill}%
\pgfsetfillopacity{0.900000}%
\pgfsetlinewidth{0.501875pt}%
\definecolor{currentstroke}{rgb}{0.200000,0.200000,0.200000}%
\pgfsetstrokecolor{currentstroke}%
\pgfsetdash{}{0pt}%
\pgfpathmoveto{\pgfqpoint{1.155885in}{3.337553in}}%
\pgfpathlineto{\pgfqpoint{1.205286in}{3.290383in}}%
\pgfpathlineto{\pgfqpoint{1.235659in}{3.271099in}}%
\pgfpathlineto{\pgfqpoint{1.186245in}{3.318301in}}%
\pgfpathlineto{\pgfqpoint{1.155885in}{3.337553in}}%
\pgfpathclose%
\pgfusepath{stroke,fill}%
\end{pgfscope}%
\begin{pgfscope}%
\pgfpathrectangle{\pgfqpoint{0.050000in}{0.050000in}}{\pgfqpoint{4.900000in}{4.900000in}}%
\pgfusepath{clip}%
\pgfsetbuttcap%
\pgfsetroundjoin%
\definecolor{currentfill}{rgb}{0.812226,0.812226,0.812226}%
\pgfsetfillcolor{currentfill}%
\pgfsetfillopacity{0.900000}%
\pgfsetlinewidth{0.501875pt}%
\definecolor{currentstroke}{rgb}{0.200000,0.200000,0.200000}%
\pgfsetstrokecolor{currentstroke}%
\pgfsetdash{}{0pt}%
\pgfpathmoveto{\pgfqpoint{2.884604in}{1.943093in}}%
\pgfpathlineto{\pgfqpoint{2.930297in}{1.901129in}}%
\pgfpathlineto{\pgfqpoint{2.963391in}{1.880080in}}%
\pgfpathlineto{\pgfqpoint{2.917695in}{1.921936in}}%
\pgfpathlineto{\pgfqpoint{2.884604in}{1.943093in}}%
\pgfpathclose%
\pgfusepath{stroke,fill}%
\end{pgfscope}%
\begin{pgfscope}%
\pgfpathrectangle{\pgfqpoint{0.050000in}{0.050000in}}{\pgfqpoint{4.900000in}{4.900000in}}%
\pgfusepath{clip}%
\pgfsetbuttcap%
\pgfsetroundjoin%
\definecolor{currentfill}{rgb}{0.461207,0.461207,0.461207}%
\pgfsetfillcolor{currentfill}%
\pgfsetfillopacity{0.900000}%
\pgfsetlinewidth{0.501875pt}%
\definecolor{currentstroke}{rgb}{0.200000,0.200000,0.200000}%
\pgfsetstrokecolor{currentstroke}%
\pgfsetdash{}{0pt}%
\pgfpathmoveto{\pgfqpoint{1.885429in}{2.747919in}}%
\pgfpathlineto{\pgfqpoint{1.933263in}{2.702017in}}%
\pgfpathlineto{\pgfqpoint{1.964783in}{2.681806in}}%
\pgfpathlineto{\pgfqpoint{1.916943in}{2.727737in}}%
\pgfpathlineto{\pgfqpoint{1.885429in}{2.747919in}}%
\pgfpathclose%
\pgfusepath{stroke,fill}%
\end{pgfscope}%
\begin{pgfscope}%
\pgfpathrectangle{\pgfqpoint{0.050000in}{0.050000in}}{\pgfqpoint{4.900000in}{4.900000in}}%
\pgfusepath{clip}%
\pgfsetbuttcap%
\pgfsetroundjoin%
\definecolor{currentfill}{rgb}{0.734579,0.734579,0.734579}%
\pgfsetfillcolor{currentfill}%
\pgfsetfillopacity{0.900000}%
\pgfsetlinewidth{0.501875pt}%
\definecolor{currentstroke}{rgb}{0.200000,0.200000,0.200000}%
\pgfsetstrokecolor{currentstroke}%
\pgfsetdash{}{0pt}%
\pgfpathmoveto{\pgfqpoint{2.615073in}{2.158551in}}%
\pgfpathlineto{\pgfqpoint{2.661339in}{2.114246in}}%
\pgfpathlineto{\pgfqpoint{2.694009in}{2.093202in}}%
\pgfpathlineto{\pgfqpoint{2.647741in}{2.137491in}}%
\pgfpathlineto{\pgfqpoint{2.615073in}{2.158551in}}%
\pgfpathclose%
\pgfusepath{stroke,fill}%
\end{pgfscope}%
\begin{pgfscope}%
\pgfpathrectangle{\pgfqpoint{0.050000in}{0.050000in}}{\pgfqpoint{4.900000in}{4.900000in}}%
\pgfusepath{clip}%
\pgfsetbuttcap%
\pgfsetroundjoin%
\definecolor{currentfill}{rgb}{0.878570,0.878570,0.878570}%
\pgfsetfillcolor{currentfill}%
\pgfsetfillopacity{0.900000}%
\pgfsetlinewidth{0.501875pt}%
\definecolor{currentstroke}{rgb}{0.200000,0.200000,0.200000}%
\pgfsetstrokecolor{currentstroke}%
\pgfsetdash{}{0pt}%
\pgfpathmoveto{\pgfqpoint{3.154181in}{1.738131in}}%
\pgfpathlineto{\pgfqpoint{3.199367in}{1.701907in}}%
\pgfpathlineto{\pgfqpoint{3.232900in}{1.681137in}}%
\pgfpathlineto{\pgfqpoint{3.187707in}{1.717173in}}%
\pgfpathlineto{\pgfqpoint{3.154181in}{1.738131in}}%
\pgfpathclose%
\pgfusepath{stroke,fill}%
\end{pgfscope}%
\begin{pgfscope}%
\pgfpathrectangle{\pgfqpoint{0.050000in}{0.050000in}}{\pgfqpoint{4.900000in}{4.900000in}}%
\pgfusepath{clip}%
\pgfsetbuttcap%
\pgfsetroundjoin%
\definecolor{currentfill}{rgb}{0.068281,0.068281,0.068281}%
\pgfsetfillcolor{currentfill}%
\pgfsetfillopacity{0.900000}%
\pgfsetlinewidth{0.501875pt}%
\definecolor{currentstroke}{rgb}{0.200000,0.200000,0.200000}%
\pgfsetstrokecolor{currentstroke}%
\pgfsetdash{}{0pt}%
\pgfpathmoveto{\pgfqpoint{0.886067in}{3.556290in}}%
\pgfpathlineto{\pgfqpoint{0.936058in}{3.508641in}}%
\pgfpathlineto{\pgfqpoint{0.966012in}{3.489697in}}%
\pgfpathlineto{\pgfqpoint{0.916007in}{3.537380in}}%
\pgfpathlineto{\pgfqpoint{0.886067in}{3.556290in}}%
\pgfpathclose%
\pgfusepath{stroke,fill}%
\end{pgfscope}%
\begin{pgfscope}%
\pgfpathrectangle{\pgfqpoint{0.050000in}{0.050000in}}{\pgfqpoint{4.900000in}{4.900000in}}%
\pgfusepath{clip}%
\pgfsetbuttcap%
\pgfsetroundjoin%
\definecolor{currentfill}{rgb}{0.964937,0.964937,0.964937}%
\pgfsetfillcolor{currentfill}%
\pgfsetfillopacity{0.900000}%
\pgfsetlinewidth{0.501875pt}%
\definecolor{currentstroke}{rgb}{0.200000,0.200000,0.200000}%
\pgfsetstrokecolor{currentstroke}%
\pgfsetdash{}{0pt}%
\pgfpathmoveto{\pgfqpoint{3.739564in}{1.388188in}}%
\pgfpathlineto{\pgfqpoint{3.784066in}{1.375744in}}%
\pgfpathlineto{\pgfqpoint{3.818560in}{1.355412in}}%
\pgfpathlineto{\pgfqpoint{3.774041in}{1.367791in}}%
\pgfpathlineto{\pgfqpoint{3.739564in}{1.388188in}}%
\pgfpathclose%
\pgfusepath{stroke,fill}%
\end{pgfscope}%
\begin{pgfscope}%
\pgfpathrectangle{\pgfqpoint{0.050000in}{0.050000in}}{\pgfqpoint{4.900000in}{4.900000in}}%
\pgfusepath{clip}%
\pgfsetbuttcap%
\pgfsetroundjoin%
\definecolor{currentfill}{rgb}{0.367243,0.367243,0.367243}%
\pgfsetfillcolor{currentfill}%
\pgfsetfillopacity{0.900000}%
\pgfsetlinewidth{0.501875pt}%
\definecolor{currentstroke}{rgb}{0.200000,0.200000,0.200000}%
\pgfsetstrokecolor{currentstroke}%
\pgfsetdash{}{0pt}%
\pgfpathmoveto{\pgfqpoint{1.615719in}{2.966563in}}%
\pgfpathlineto{\pgfqpoint{1.664141in}{2.920182in}}%
\pgfpathlineto{\pgfqpoint{1.695243in}{2.900311in}}%
\pgfpathlineto{\pgfqpoint{1.646812in}{2.946723in}}%
\pgfpathlineto{\pgfqpoint{1.615719in}{2.966563in}}%
\pgfpathclose%
\pgfusepath{stroke,fill}%
\end{pgfscope}%
\begin{pgfscope}%
\pgfpathrectangle{\pgfqpoint{0.050000in}{0.050000in}}{\pgfqpoint{4.900000in}{4.900000in}}%
\pgfusepath{clip}%
\pgfsetbuttcap%
\pgfsetroundjoin%
\definecolor{currentfill}{rgb}{0.957555,0.957555,0.957555}%
\pgfsetfillcolor{currentfill}%
\pgfsetfillopacity{0.900000}%
\pgfsetlinewidth{0.501875pt}%
\definecolor{currentstroke}{rgb}{0.200000,0.200000,0.200000}%
\pgfsetstrokecolor{currentstroke}%
\pgfsetdash{}{0pt}%
\pgfpathmoveto{\pgfqpoint{3.660637in}{1.424507in}}%
\pgfpathlineto{\pgfqpoint{3.705197in}{1.408503in}}%
\pgfpathlineto{\pgfqpoint{3.739564in}{1.388188in}}%
\pgfpathlineto{\pgfqpoint{3.694987in}{1.404095in}}%
\pgfpathlineto{\pgfqpoint{3.660637in}{1.424507in}}%
\pgfpathclose%
\pgfusepath{stroke,fill}%
\end{pgfscope}%
\begin{pgfscope}%
\pgfpathrectangle{\pgfqpoint{0.050000in}{0.050000in}}{\pgfqpoint{4.900000in}{4.900000in}}%
\pgfusepath{clip}%
\pgfsetbuttcap%
\pgfsetroundjoin%
\definecolor{currentfill}{rgb}{0.970473,0.970473,0.970473}%
\pgfsetfillcolor{currentfill}%
\pgfsetfillopacity{0.900000}%
\pgfsetlinewidth{0.501875pt}%
\definecolor{currentstroke}{rgb}{0.200000,0.200000,0.200000}%
\pgfsetstrokecolor{currentstroke}%
\pgfsetdash{}{0pt}%
\pgfpathmoveto{\pgfqpoint{3.818560in}{1.355412in}}%
\pgfpathlineto{\pgfqpoint{3.863014in}{1.346367in}}%
\pgfpathlineto{\pgfqpoint{3.897638in}{1.325989in}}%
\pgfpathlineto{\pgfqpoint{3.853166in}{1.334999in}}%
\pgfpathlineto{\pgfqpoint{3.818560in}{1.355412in}}%
\pgfpathclose%
\pgfusepath{stroke,fill}%
\end{pgfscope}%
\begin{pgfscope}%
\pgfpathrectangle{\pgfqpoint{0.050000in}{0.050000in}}{\pgfqpoint{4.900000in}{4.900000in}}%
\pgfusepath{clip}%
\pgfsetbuttcap%
\pgfsetroundjoin%
\definecolor{currentfill}{rgb}{0.629020,0.629020,0.629020}%
\pgfsetfillcolor{currentfill}%
\pgfsetfillopacity{0.900000}%
\pgfsetlinewidth{0.501875pt}%
\definecolor{currentstroke}{rgb}{0.200000,0.200000,0.200000}%
\pgfsetstrokecolor{currentstroke}%
\pgfsetdash{}{0pt}%
\pgfpathmoveto{\pgfqpoint{2.345469in}{2.376786in}}%
\pgfpathlineto{\pgfqpoint{2.392322in}{2.331690in}}%
\pgfpathlineto{\pgfqpoint{2.424574in}{2.310900in}}%
\pgfpathlineto{\pgfqpoint{2.377716in}{2.356018in}}%
\pgfpathlineto{\pgfqpoint{2.345469in}{2.376786in}}%
\pgfpathclose%
\pgfusepath{stroke,fill}%
\end{pgfscope}%
\begin{pgfscope}%
\pgfpathrectangle{\pgfqpoint{0.050000in}{0.050000in}}{\pgfqpoint{4.900000in}{4.900000in}}%
\pgfusepath{clip}%
\pgfsetbuttcap%
\pgfsetroundjoin%
\definecolor{currentfill}{rgb}{0.950173,0.950173,0.950173}%
\pgfsetfillcolor{currentfill}%
\pgfsetfillopacity{0.900000}%
\pgfsetlinewidth{0.501875pt}%
\definecolor{currentstroke}{rgb}{0.200000,0.200000,0.200000}%
\pgfsetstrokecolor{currentstroke}%
\pgfsetdash{}{0pt}%
\pgfpathmoveto{\pgfqpoint{3.581768in}{1.464492in}}%
\pgfpathlineto{\pgfqpoint{3.626397in}{1.444839in}}%
\pgfpathlineto{\pgfqpoint{3.660637in}{1.424507in}}%
\pgfpathlineto{\pgfqpoint{3.615993in}{1.444033in}}%
\pgfpathlineto{\pgfqpoint{3.581768in}{1.464492in}}%
\pgfpathclose%
\pgfusepath{stroke,fill}%
\end{pgfscope}%
\begin{pgfscope}%
\pgfpathrectangle{\pgfqpoint{0.050000in}{0.050000in}}{\pgfqpoint{4.900000in}{4.900000in}}%
\pgfusepath{clip}%
\pgfsetbuttcap%
\pgfsetroundjoin%
\definecolor{currentfill}{rgb}{0.262053,0.262053,0.262053}%
\pgfsetfillcolor{currentfill}%
\pgfsetfillopacity{0.900000}%
\pgfsetlinewidth{0.501875pt}%
\definecolor{currentstroke}{rgb}{0.200000,0.200000,0.200000}%
\pgfsetstrokecolor{currentstroke}%
\pgfsetdash{}{0pt}%
\pgfpathmoveto{\pgfqpoint{1.345915in}{3.185285in}}%
\pgfpathlineto{\pgfqpoint{1.394927in}{3.138424in}}%
\pgfpathlineto{\pgfqpoint{1.425611in}{3.118893in}}%
\pgfpathlineto{\pgfqpoint{1.376587in}{3.165785in}}%
\pgfpathlineto{\pgfqpoint{1.345915in}{3.185285in}}%
\pgfpathclose%
\pgfusepath{stroke,fill}%
\end{pgfscope}%
\begin{pgfscope}%
\pgfpathrectangle{\pgfqpoint{0.050000in}{0.050000in}}{\pgfqpoint{4.900000in}{4.900000in}}%
\pgfusepath{clip}%
\pgfsetbuttcap%
\pgfsetroundjoin%
\definecolor{currentfill}{rgb}{0.525798,0.525798,0.525798}%
\pgfsetfillcolor{currentfill}%
\pgfsetfillopacity{0.900000}%
\pgfsetlinewidth{0.501875pt}%
\definecolor{currentstroke}{rgb}{0.200000,0.200000,0.200000}%
\pgfsetstrokecolor{currentstroke}%
\pgfsetdash{}{0pt}%
\pgfpathmoveto{\pgfqpoint{2.075772in}{2.595404in}}%
\pgfpathlineto{\pgfqpoint{2.123215in}{2.549811in}}%
\pgfpathlineto{\pgfqpoint{2.155048in}{2.529353in}}%
\pgfpathlineto{\pgfqpoint{2.107599in}{2.574973in}}%
\pgfpathlineto{\pgfqpoint{2.075772in}{2.595404in}}%
\pgfpathclose%
\pgfusepath{stroke,fill}%
\end{pgfscope}%
\begin{pgfscope}%
\pgfpathrectangle{\pgfqpoint{0.050000in}{0.050000in}}{\pgfqpoint{4.900000in}{4.900000in}}%
\pgfusepath{clip}%
\pgfsetbuttcap%
\pgfsetroundjoin%
\definecolor{currentfill}{rgb}{0.942791,0.942791,0.942791}%
\pgfsetfillcolor{currentfill}%
\pgfsetfillopacity{0.900000}%
\pgfsetlinewidth{0.501875pt}%
\definecolor{currentstroke}{rgb}{0.200000,0.200000,0.200000}%
\pgfsetstrokecolor{currentstroke}%
\pgfsetdash{}{0pt}%
\pgfpathmoveto{\pgfqpoint{3.502940in}{1.508178in}}%
\pgfpathlineto{\pgfqpoint{3.547652in}{1.484874in}}%
\pgfpathlineto{\pgfqpoint{3.581768in}{1.464492in}}%
\pgfpathlineto{\pgfqpoint{3.537043in}{1.487641in}}%
\pgfpathlineto{\pgfqpoint{3.502940in}{1.508178in}}%
\pgfpathclose%
\pgfusepath{stroke,fill}%
\end{pgfscope}%
\begin{pgfscope}%
\pgfpathrectangle{\pgfqpoint{0.050000in}{0.050000in}}{\pgfqpoint{4.900000in}{4.900000in}}%
\pgfusepath{clip}%
\pgfsetbuttcap%
\pgfsetroundjoin%
\definecolor{currentfill}{rgb}{0.136563,0.136563,0.136563}%
\pgfsetfillcolor{currentfill}%
\pgfsetfillopacity{0.900000}%
\pgfsetlinewidth{0.501875pt}%
\definecolor{currentstroke}{rgb}{0.200000,0.200000,0.200000}%
\pgfsetstrokecolor{currentstroke}%
\pgfsetdash{}{0pt}%
\pgfpathmoveto{\pgfqpoint{1.076019in}{3.404084in}}%
\pgfpathlineto{\pgfqpoint{1.125621in}{3.356744in}}%
\pgfpathlineto{\pgfqpoint{1.155885in}{3.337553in}}%
\pgfpathlineto{\pgfqpoint{1.106270in}{3.384926in}}%
\pgfpathlineto{\pgfqpoint{1.076019in}{3.404084in}}%
\pgfpathclose%
\pgfusepath{stroke,fill}%
\end{pgfscope}%
\begin{pgfscope}%
\pgfpathrectangle{\pgfqpoint{0.050000in}{0.050000in}}{\pgfqpoint{4.900000in}{4.900000in}}%
\pgfusepath{clip}%
\pgfsetbuttcap%
\pgfsetroundjoin%
\definecolor{currentfill}{rgb}{0.788112,0.788112,0.788112}%
\pgfsetfillcolor{currentfill}%
\pgfsetfillopacity{0.900000}%
\pgfsetlinewidth{0.501875pt}%
\definecolor{currentstroke}{rgb}{0.200000,0.200000,0.200000}%
\pgfsetstrokecolor{currentstroke}%
\pgfsetdash{}{0pt}%
\pgfpathmoveto{\pgfqpoint{2.805736in}{2.007213in}}%
\pgfpathlineto{\pgfqpoint{2.851617in}{1.964195in}}%
\pgfpathlineto{\pgfqpoint{2.884604in}{1.943093in}}%
\pgfpathlineto{\pgfqpoint{2.838720in}{1.986037in}}%
\pgfpathlineto{\pgfqpoint{2.805736in}{2.007213in}}%
\pgfpathclose%
\pgfusepath{stroke,fill}%
\end{pgfscope}%
\begin{pgfscope}%
\pgfpathrectangle{\pgfqpoint{0.050000in}{0.050000in}}{\pgfqpoint{4.900000in}{4.900000in}}%
\pgfusepath{clip}%
\pgfsetbuttcap%
\pgfsetroundjoin%
\definecolor{currentfill}{rgb}{0.858762,0.858762,0.858762}%
\pgfsetfillcolor{currentfill}%
\pgfsetfillopacity{0.900000}%
\pgfsetlinewidth{0.501875pt}%
\definecolor{currentstroke}{rgb}{0.200000,0.200000,0.200000}%
\pgfsetstrokecolor{currentstroke}%
\pgfsetdash{}{0pt}%
\pgfpathmoveto{\pgfqpoint{3.075415in}{1.797561in}}%
\pgfpathlineto{\pgfqpoint{3.120761in}{1.759033in}}%
\pgfpathlineto{\pgfqpoint{3.154181in}{1.738131in}}%
\pgfpathlineto{\pgfqpoint{3.108830in}{1.776490in}}%
\pgfpathlineto{\pgfqpoint{3.075415in}{1.797561in}}%
\pgfpathclose%
\pgfusepath{stroke,fill}%
\end{pgfscope}%
\begin{pgfscope}%
\pgfpathrectangle{\pgfqpoint{0.050000in}{0.050000in}}{\pgfqpoint{4.900000in}{4.900000in}}%
\pgfusepath{clip}%
\pgfsetbuttcap%
\pgfsetroundjoin%
\definecolor{currentfill}{rgb}{0.428143,0.428143,0.428143}%
\pgfsetfillcolor{currentfill}%
\pgfsetfillopacity{0.900000}%
\pgfsetlinewidth{0.501875pt}%
\definecolor{currentstroke}{rgb}{0.200000,0.200000,0.200000}%
\pgfsetstrokecolor{currentstroke}%
\pgfsetdash{}{0pt}%
\pgfpathmoveto{\pgfqpoint{1.805983in}{2.814110in}}%
\pgfpathlineto{\pgfqpoint{1.854015in}{2.768038in}}%
\pgfpathlineto{\pgfqpoint{1.885429in}{2.747919in}}%
\pgfpathlineto{\pgfqpoint{1.837389in}{2.794020in}}%
\pgfpathlineto{\pgfqpoint{1.805983in}{2.814110in}}%
\pgfpathclose%
\pgfusepath{stroke,fill}%
\end{pgfscope}%
\begin{pgfscope}%
\pgfpathrectangle{\pgfqpoint{0.050000in}{0.050000in}}{\pgfqpoint{4.900000in}{4.900000in}}%
\pgfusepath{clip}%
\pgfsetbuttcap%
\pgfsetroundjoin%
\definecolor{currentfill}{rgb}{0.700992,0.700992,0.700992}%
\pgfsetfillcolor{currentfill}%
\pgfsetfillopacity{0.900000}%
\pgfsetlinewidth{0.501875pt}%
\definecolor{currentstroke}{rgb}{0.200000,0.200000,0.200000}%
\pgfsetstrokecolor{currentstroke}%
\pgfsetdash{}{0pt}%
\pgfpathmoveto{\pgfqpoint{2.536046in}{2.224202in}}%
\pgfpathlineto{\pgfqpoint{2.582508in}{2.179550in}}%
\pgfpathlineto{\pgfqpoint{2.615073in}{2.158551in}}%
\pgfpathlineto{\pgfqpoint{2.568608in}{2.203207in}}%
\pgfpathlineto{\pgfqpoint{2.536046in}{2.224202in}}%
\pgfpathclose%
\pgfusepath{stroke,fill}%
\end{pgfscope}%
\begin{pgfscope}%
\pgfpathrectangle{\pgfqpoint{0.050000in}{0.050000in}}{\pgfqpoint{4.900000in}{4.900000in}}%
\pgfusepath{clip}%
\pgfsetbuttcap%
\pgfsetroundjoin%
\definecolor{currentfill}{rgb}{0.929504,0.929504,0.929504}%
\pgfsetfillcolor{currentfill}%
\pgfsetfillopacity{0.900000}%
\pgfsetlinewidth{0.501875pt}%
\definecolor{currentstroke}{rgb}{0.200000,0.200000,0.200000}%
\pgfsetstrokecolor{currentstroke}%
\pgfsetdash{}{0pt}%
\pgfpathmoveto{\pgfqpoint{3.424139in}{1.555503in}}%
\pgfpathlineto{\pgfqpoint{3.468947in}{1.528642in}}%
\pgfpathlineto{\pgfqpoint{3.502940in}{1.508178in}}%
\pgfpathlineto{\pgfqpoint{3.458120in}{1.534862in}}%
\pgfpathlineto{\pgfqpoint{3.424139in}{1.555503in}}%
\pgfpathclose%
\pgfusepath{stroke,fill}%
\end{pgfscope}%
\begin{pgfscope}%
\pgfpathrectangle{\pgfqpoint{0.050000in}{0.050000in}}{\pgfqpoint{4.900000in}{4.900000in}}%
\pgfusepath{clip}%
\pgfsetbuttcap%
\pgfsetroundjoin%
\definecolor{currentfill}{rgb}{0.031865,0.031865,0.031865}%
\pgfsetfillcolor{currentfill}%
\pgfsetfillopacity{0.900000}%
\pgfsetlinewidth{0.501875pt}%
\definecolor{currentstroke}{rgb}{0.200000,0.200000,0.200000}%
\pgfsetstrokecolor{currentstroke}%
\pgfsetdash{}{0pt}%
\pgfpathmoveto{\pgfqpoint{0.806030in}{3.622961in}}%
\pgfpathlineto{\pgfqpoint{0.856223in}{3.575141in}}%
\pgfpathlineto{\pgfqpoint{0.886067in}{3.556290in}}%
\pgfpathlineto{\pgfqpoint{0.835860in}{3.604144in}}%
\pgfpathlineto{\pgfqpoint{0.806030in}{3.622961in}}%
\pgfpathclose%
\pgfusepath{stroke,fill}%
\end{pgfscope}%
\begin{pgfscope}%
\pgfpathrectangle{\pgfqpoint{0.050000in}{0.050000in}}{\pgfqpoint{4.900000in}{4.900000in}}%
\pgfusepath{clip}%
\pgfsetbuttcap%
\pgfsetroundjoin%
\definecolor{currentfill}{rgb}{0.338824,0.338824,0.338824}%
\pgfsetfillcolor{currentfill}%
\pgfsetfillopacity{0.900000}%
\pgfsetlinewidth{0.501875pt}%
\definecolor{currentstroke}{rgb}{0.200000,0.200000,0.200000}%
\pgfsetstrokecolor{currentstroke}%
\pgfsetdash{}{0pt}%
\pgfpathmoveto{\pgfqpoint{1.536101in}{3.032893in}}%
\pgfpathlineto{\pgfqpoint{1.584723in}{2.986342in}}%
\pgfpathlineto{\pgfqpoint{1.615719in}{2.966563in}}%
\pgfpathlineto{\pgfqpoint{1.567087in}{3.013145in}}%
\pgfpathlineto{\pgfqpoint{1.536101in}{3.032893in}}%
\pgfpathclose%
\pgfusepath{stroke,fill}%
\end{pgfscope}%
\begin{pgfscope}%
\pgfpathrectangle{\pgfqpoint{0.050000in}{0.050000in}}{\pgfqpoint{4.900000in}{4.900000in}}%
\pgfusepath{clip}%
\pgfsetbuttcap%
\pgfsetroundjoin%
\definecolor{currentfill}{rgb}{0.590634,0.590634,0.590634}%
\pgfsetfillcolor{currentfill}%
\pgfsetfillopacity{0.900000}%
\pgfsetlinewidth{0.501875pt}%
\definecolor{currentstroke}{rgb}{0.200000,0.200000,0.200000}%
\pgfsetstrokecolor{currentstroke}%
\pgfsetdash{}{0pt}%
\pgfpathmoveto{\pgfqpoint{2.266271in}{2.442767in}}%
\pgfpathlineto{\pgfqpoint{2.313323in}{2.397490in}}%
\pgfpathlineto{\pgfqpoint{2.345469in}{2.376786in}}%
\pgfpathlineto{\pgfqpoint{2.298412in}{2.422088in}}%
\pgfpathlineto{\pgfqpoint{2.266271in}{2.442767in}}%
\pgfpathclose%
\pgfusepath{stroke,fill}%
\end{pgfscope}%
\begin{pgfscope}%
\pgfpathrectangle{\pgfqpoint{0.050000in}{0.050000in}}{\pgfqpoint{4.900000in}{4.900000in}}%
\pgfusepath{clip}%
\pgfsetbuttcap%
\pgfsetroundjoin%
\definecolor{currentfill}{rgb}{0.217762,0.217762,0.217762}%
\pgfsetfillcolor{currentfill}%
\pgfsetfillopacity{0.900000}%
\pgfsetlinewidth{0.501875pt}%
\definecolor{currentstroke}{rgb}{0.200000,0.200000,0.200000}%
\pgfsetstrokecolor{currentstroke}%
\pgfsetdash{}{0pt}%
\pgfpathmoveto{\pgfqpoint{1.266127in}{3.251754in}}%
\pgfpathlineto{\pgfqpoint{1.315339in}{3.204723in}}%
\pgfpathlineto{\pgfqpoint{1.345915in}{3.185285in}}%
\pgfpathlineto{\pgfqpoint{1.296691in}{3.232348in}}%
\pgfpathlineto{\pgfqpoint{1.266127in}{3.251754in}}%
\pgfpathclose%
\pgfusepath{stroke,fill}%
\end{pgfscope}%
\begin{pgfscope}%
\pgfpathrectangle{\pgfqpoint{0.050000in}{0.050000in}}{\pgfqpoint{4.900000in}{4.900000in}}%
\pgfusepath{clip}%
\pgfsetbuttcap%
\pgfsetroundjoin%
\definecolor{currentfill}{rgb}{0.912526,0.912526,0.912526}%
\pgfsetfillcolor{currentfill}%
\pgfsetfillopacity{0.900000}%
\pgfsetlinewidth{0.501875pt}%
\definecolor{currentstroke}{rgb}{0.200000,0.200000,0.200000}%
\pgfsetstrokecolor{currentstroke}%
\pgfsetdash{}{0pt}%
\pgfpathmoveto{\pgfqpoint{3.345344in}{1.606300in}}%
\pgfpathlineto{\pgfqpoint{3.390265in}{1.576075in}}%
\pgfpathlineto{\pgfqpoint{3.424139in}{1.555503in}}%
\pgfpathlineto{\pgfqpoint{3.379208in}{1.585538in}}%
\pgfpathlineto{\pgfqpoint{3.345344in}{1.606300in}}%
\pgfpathclose%
\pgfusepath{stroke,fill}%
\end{pgfscope}%
\begin{pgfscope}%
\pgfpathrectangle{\pgfqpoint{0.050000in}{0.050000in}}{\pgfqpoint{4.900000in}{4.900000in}}%
\pgfusepath{clip}%
\pgfsetbuttcap%
\pgfsetroundjoin%
\definecolor{currentfill}{rgb}{0.495656,0.495656,0.495656}%
\pgfsetfillcolor{currentfill}%
\pgfsetfillopacity{0.900000}%
\pgfsetlinewidth{0.501875pt}%
\definecolor{currentstroke}{rgb}{0.200000,0.200000,0.200000}%
\pgfsetstrokecolor{currentstroke}%
\pgfsetdash{}{0pt}%
\pgfpathmoveto{\pgfqpoint{1.996404in}{2.661532in}}%
\pgfpathlineto{\pgfqpoint{2.044045in}{2.615770in}}%
\pgfpathlineto{\pgfqpoint{2.075772in}{2.595404in}}%
\pgfpathlineto{\pgfqpoint{2.028124in}{2.641194in}}%
\pgfpathlineto{\pgfqpoint{1.996404in}{2.661532in}}%
\pgfpathclose%
\pgfusepath{stroke,fill}%
\end{pgfscope}%
\begin{pgfscope}%
\pgfpathrectangle{\pgfqpoint{0.050000in}{0.050000in}}{\pgfqpoint{4.900000in}{4.900000in}}%
\pgfusepath{clip}%
\pgfsetbuttcap%
\pgfsetroundjoin%
\definecolor{currentfill}{rgb}{0.839785,0.839785,0.839785}%
\pgfsetfillcolor{currentfill}%
\pgfsetfillopacity{0.900000}%
\pgfsetlinewidth{0.501875pt}%
\definecolor{currentstroke}{rgb}{0.200000,0.200000,0.200000}%
\pgfsetstrokecolor{currentstroke}%
\pgfsetdash{}{0pt}%
\pgfpathmoveto{\pgfqpoint{2.996590in}{1.858977in}}%
\pgfpathlineto{\pgfqpoint{3.042107in}{1.818577in}}%
\pgfpathlineto{\pgfqpoint{3.075415in}{1.797561in}}%
\pgfpathlineto{\pgfqpoint{3.029895in}{1.837819in}}%
\pgfpathlineto{\pgfqpoint{2.996590in}{1.858977in}}%
\pgfpathclose%
\pgfusepath{stroke,fill}%
\end{pgfscope}%
\begin{pgfscope}%
\pgfpathrectangle{\pgfqpoint{0.050000in}{0.050000in}}{\pgfqpoint{4.900000in}{4.900000in}}%
\pgfusepath{clip}%
\pgfsetbuttcap%
\pgfsetroundjoin%
\definecolor{currentfill}{rgb}{0.763998,0.763998,0.763998}%
\pgfsetfillcolor{currentfill}%
\pgfsetfillopacity{0.900000}%
\pgfsetlinewidth{0.501875pt}%
\definecolor{currentstroke}{rgb}{0.200000,0.200000,0.200000}%
\pgfsetstrokecolor{currentstroke}%
\pgfsetdash{}{0pt}%
\pgfpathmoveto{\pgfqpoint{2.726783in}{2.072099in}}%
\pgfpathlineto{\pgfqpoint{2.772856in}{2.028332in}}%
\pgfpathlineto{\pgfqpoint{2.805736in}{2.007213in}}%
\pgfpathlineto{\pgfqpoint{2.759660in}{2.050937in}}%
\pgfpathlineto{\pgfqpoint{2.726783in}{2.072099in}}%
\pgfpathclose%
\pgfusepath{stroke,fill}%
\end{pgfscope}%
\begin{pgfscope}%
\pgfpathrectangle{\pgfqpoint{0.050000in}{0.050000in}}{\pgfqpoint{4.900000in}{4.900000in}}%
\pgfusepath{clip}%
\pgfsetbuttcap%
\pgfsetroundjoin%
\definecolor{currentfill}{rgb}{0.100146,0.100146,0.100146}%
\pgfsetfillcolor{currentfill}%
\pgfsetfillopacity{0.900000}%
\pgfsetlinewidth{0.501875pt}%
\definecolor{currentstroke}{rgb}{0.200000,0.200000,0.200000}%
\pgfsetstrokecolor{currentstroke}%
\pgfsetdash{}{0pt}%
\pgfpathmoveto{\pgfqpoint{0.996059in}{3.470693in}}%
\pgfpathlineto{\pgfqpoint{1.045863in}{3.423182in}}%
\pgfpathlineto{\pgfqpoint{1.076019in}{3.404084in}}%
\pgfpathlineto{\pgfqpoint{1.026202in}{3.451629in}}%
\pgfpathlineto{\pgfqpoint{0.996059in}{3.470693in}}%
\pgfpathclose%
\pgfusepath{stroke,fill}%
\end{pgfscope}%
\begin{pgfscope}%
\pgfpathrectangle{\pgfqpoint{0.050000in}{0.050000in}}{\pgfqpoint{4.900000in}{4.900000in}}%
\pgfusepath{clip}%
\pgfsetbuttcap%
\pgfsetroundjoin%
\definecolor{currentfill}{rgb}{0.399723,0.399723,0.399723}%
\pgfsetfillcolor{currentfill}%
\pgfsetfillopacity{0.900000}%
\pgfsetlinewidth{0.501875pt}%
\definecolor{currentstroke}{rgb}{0.200000,0.200000,0.200000}%
\pgfsetstrokecolor{currentstroke}%
\pgfsetdash{}{0pt}%
\pgfpathmoveto{\pgfqpoint{1.726444in}{2.880377in}}%
\pgfpathlineto{\pgfqpoint{1.774675in}{2.834136in}}%
\pgfpathlineto{\pgfqpoint{1.805983in}{2.814110in}}%
\pgfpathlineto{\pgfqpoint{1.757743in}{2.860381in}}%
\pgfpathlineto{\pgfqpoint{1.726444in}{2.880377in}}%
\pgfpathclose%
\pgfusepath{stroke,fill}%
\end{pgfscope}%
\begin{pgfscope}%
\pgfpathrectangle{\pgfqpoint{0.050000in}{0.050000in}}{\pgfqpoint{4.900000in}{4.900000in}}%
\pgfusepath{clip}%
\pgfsetbuttcap%
\pgfsetroundjoin%
\definecolor{currentfill}{rgb}{0.662607,0.662607,0.662607}%
\pgfsetfillcolor{currentfill}%
\pgfsetfillopacity{0.900000}%
\pgfsetlinewidth{0.501875pt}%
\definecolor{currentstroke}{rgb}{0.200000,0.200000,0.200000}%
\pgfsetstrokecolor{currentstroke}%
\pgfsetdash{}{0pt}%
\pgfpathmoveto{\pgfqpoint{2.456927in}{2.290045in}}%
\pgfpathlineto{\pgfqpoint{2.503587in}{2.245135in}}%
\pgfpathlineto{\pgfqpoint{2.536046in}{2.224202in}}%
\pgfpathlineto{\pgfqpoint{2.489383in}{2.269127in}}%
\pgfpathlineto{\pgfqpoint{2.456927in}{2.290045in}}%
\pgfpathclose%
\pgfusepath{stroke,fill}%
\end{pgfscope}%
\begin{pgfscope}%
\pgfpathrectangle{\pgfqpoint{0.050000in}{0.050000in}}{\pgfqpoint{4.900000in}{4.900000in}}%
\pgfusepath{clip}%
\pgfsetbuttcap%
\pgfsetroundjoin%
\definecolor{currentfill}{rgb}{0.000000,0.000000,0.000000}%
\pgfsetfillcolor{currentfill}%
\pgfsetfillopacity{0.900000}%
\pgfsetlinewidth{0.501875pt}%
\definecolor{currentstroke}{rgb}{0.200000,0.200000,0.200000}%
\pgfsetstrokecolor{currentstroke}%
\pgfsetdash{}{0pt}%
\pgfpathmoveto{\pgfqpoint{0.725899in}{3.689709in}}%
\pgfpathlineto{\pgfqpoint{0.776294in}{3.641718in}}%
\pgfpathlineto{\pgfqpoint{0.806030in}{3.622961in}}%
\pgfpathlineto{\pgfqpoint{0.755620in}{3.670987in}}%
\pgfpathlineto{\pgfqpoint{0.725899in}{3.689709in}}%
\pgfpathclose%
\pgfusepath{stroke,fill}%
\end{pgfscope}%
\begin{pgfscope}%
\pgfpathrectangle{\pgfqpoint{0.050000in}{0.050000in}}{\pgfqpoint{4.900000in}{4.900000in}}%
\pgfusepath{clip}%
\pgfsetbuttcap%
\pgfsetroundjoin%
\definecolor{currentfill}{rgb}{0.895548,0.895548,0.895548}%
\pgfsetfillcolor{currentfill}%
\pgfsetfillopacity{0.900000}%
\pgfsetlinewidth{0.501875pt}%
\definecolor{currentstroke}{rgb}{0.200000,0.200000,0.200000}%
\pgfsetstrokecolor{currentstroke}%
\pgfsetdash{}{0pt}%
\pgfpathmoveto{\pgfqpoint{3.266539in}{1.660306in}}%
\pgfpathlineto{\pgfqpoint{3.311589in}{1.626998in}}%
\pgfpathlineto{\pgfqpoint{3.345344in}{1.606300in}}%
\pgfpathlineto{\pgfqpoint{3.300287in}{1.639415in}}%
\pgfpathlineto{\pgfqpoint{3.266539in}{1.660306in}}%
\pgfpathclose%
\pgfusepath{stroke,fill}%
\end{pgfscope}%
\begin{pgfscope}%
\pgfpathrectangle{\pgfqpoint{0.050000in}{0.050000in}}{\pgfqpoint{4.900000in}{4.900000in}}%
\pgfusepath{clip}%
\pgfsetbuttcap%
\pgfsetroundjoin%
\definecolor{currentfill}{rgb}{0.300807,0.300807,0.300807}%
\pgfsetfillcolor{currentfill}%
\pgfsetfillopacity{0.900000}%
\pgfsetlinewidth{0.501875pt}%
\definecolor{currentstroke}{rgb}{0.200000,0.200000,0.200000}%
\pgfsetstrokecolor{currentstroke}%
\pgfsetdash{}{0pt}%
\pgfpathmoveto{\pgfqpoint{1.456391in}{3.099300in}}%
\pgfpathlineto{\pgfqpoint{1.505213in}{3.052579in}}%
\pgfpathlineto{\pgfqpoint{1.536101in}{3.032893in}}%
\pgfpathlineto{\pgfqpoint{1.487268in}{3.079646in}}%
\pgfpathlineto{\pgfqpoint{1.456391in}{3.099300in}}%
\pgfpathclose%
\pgfusepath{stroke,fill}%
\end{pgfscope}%
\begin{pgfscope}%
\pgfpathrectangle{\pgfqpoint{0.050000in}{0.050000in}}{\pgfqpoint{4.900000in}{4.900000in}}%
\pgfusepath{clip}%
\pgfsetbuttcap%
\pgfsetroundjoin%
\definecolor{currentfill}{rgb}{0.560246,0.560246,0.560246}%
\pgfsetfillcolor{currentfill}%
\pgfsetfillopacity{0.900000}%
\pgfsetlinewidth{0.501875pt}%
\definecolor{currentstroke}{rgb}{0.200000,0.200000,0.200000}%
\pgfsetstrokecolor{currentstroke}%
\pgfsetdash{}{0pt}%
\pgfpathmoveto{\pgfqpoint{2.186981in}{2.508831in}}%
\pgfpathlineto{\pgfqpoint{2.234231in}{2.463381in}}%
\pgfpathlineto{\pgfqpoint{2.266271in}{2.442767in}}%
\pgfpathlineto{\pgfqpoint{2.219016in}{2.488243in}}%
\pgfpathlineto{\pgfqpoint{2.186981in}{2.508831in}}%
\pgfpathclose%
\pgfusepath{stroke,fill}%
\end{pgfscope}%
\begin{pgfscope}%
\pgfpathrectangle{\pgfqpoint{0.050000in}{0.050000in}}{\pgfqpoint{4.900000in}{4.900000in}}%
\pgfusepath{clip}%
\pgfsetbuttcap%
\pgfsetroundjoin%
\definecolor{currentfill}{rgb}{0.179008,0.179008,0.179008}%
\pgfsetfillcolor{currentfill}%
\pgfsetfillopacity{0.900000}%
\pgfsetlinewidth{0.501875pt}%
\definecolor{currentstroke}{rgb}{0.200000,0.200000,0.200000}%
\pgfsetstrokecolor{currentstroke}%
\pgfsetdash{}{0pt}%
\pgfpathmoveto{\pgfqpoint{1.186245in}{3.318301in}}%
\pgfpathlineto{\pgfqpoint{1.235659in}{3.271099in}}%
\pgfpathlineto{\pgfqpoint{1.266127in}{3.251754in}}%
\pgfpathlineto{\pgfqpoint{1.216701in}{3.298989in}}%
\pgfpathlineto{\pgfqpoint{1.186245in}{3.318301in}}%
\pgfpathclose%
\pgfusepath{stroke,fill}%
\end{pgfscope}%
\begin{pgfscope}%
\pgfpathrectangle{\pgfqpoint{0.050000in}{0.050000in}}{\pgfqpoint{4.900000in}{4.900000in}}%
\pgfusepath{clip}%
\pgfsetbuttcap%
\pgfsetroundjoin%
\definecolor{currentfill}{rgb}{0.815671,0.815671,0.815671}%
\pgfsetfillcolor{currentfill}%
\pgfsetfillopacity{0.900000}%
\pgfsetlinewidth{0.501875pt}%
\definecolor{currentstroke}{rgb}{0.200000,0.200000,0.200000}%
\pgfsetstrokecolor{currentstroke}%
\pgfsetdash{}{0pt}%
\pgfpathmoveto{\pgfqpoint{2.917695in}{1.921936in}}%
\pgfpathlineto{\pgfqpoint{2.963391in}{1.880080in}}%
\pgfpathlineto{\pgfqpoint{2.996590in}{1.858977in}}%
\pgfpathlineto{\pgfqpoint{2.950891in}{1.900724in}}%
\pgfpathlineto{\pgfqpoint{2.917695in}{1.921936in}}%
\pgfpathclose%
\pgfusepath{stroke,fill}%
\end{pgfscope}%
\begin{pgfscope}%
\pgfpathrectangle{\pgfqpoint{0.050000in}{0.050000in}}{\pgfqpoint{4.900000in}{4.900000in}}%
\pgfusepath{clip}%
\pgfsetbuttcap%
\pgfsetroundjoin%
\definecolor{currentfill}{rgb}{0.461207,0.461207,0.461207}%
\pgfsetfillcolor{currentfill}%
\pgfsetfillopacity{0.900000}%
\pgfsetlinewidth{0.501875pt}%
\definecolor{currentstroke}{rgb}{0.200000,0.200000,0.200000}%
\pgfsetstrokecolor{currentstroke}%
\pgfsetdash{}{0pt}%
\pgfpathmoveto{\pgfqpoint{1.916943in}{2.727737in}}%
\pgfpathlineto{\pgfqpoint{1.964783in}{2.681806in}}%
\pgfpathlineto{\pgfqpoint{1.996404in}{2.661532in}}%
\pgfpathlineto{\pgfqpoint{1.948556in}{2.707492in}}%
\pgfpathlineto{\pgfqpoint{1.916943in}{2.727737in}}%
\pgfpathclose%
\pgfusepath{stroke,fill}%
\end{pgfscope}%
\begin{pgfscope}%
\pgfpathrectangle{\pgfqpoint{0.050000in}{0.050000in}}{\pgfqpoint{4.900000in}{4.900000in}}%
\pgfusepath{clip}%
\pgfsetbuttcap%
\pgfsetroundjoin%
\definecolor{currentfill}{rgb}{0.734579,0.734579,0.734579}%
\pgfsetfillcolor{currentfill}%
\pgfsetfillopacity{0.900000}%
\pgfsetlinewidth{0.501875pt}%
\definecolor{currentstroke}{rgb}{0.200000,0.200000,0.200000}%
\pgfsetstrokecolor{currentstroke}%
\pgfsetdash{}{0pt}%
\pgfpathmoveto{\pgfqpoint{2.647741in}{2.137491in}}%
\pgfpathlineto{\pgfqpoint{2.694009in}{2.093202in}}%
\pgfpathlineto{\pgfqpoint{2.726783in}{2.072099in}}%
\pgfpathlineto{\pgfqpoint{2.680512in}{2.116370in}}%
\pgfpathlineto{\pgfqpoint{2.647741in}{2.137491in}}%
\pgfpathclose%
\pgfusepath{stroke,fill}%
\end{pgfscope}%
\begin{pgfscope}%
\pgfpathrectangle{\pgfqpoint{0.050000in}{0.050000in}}{\pgfqpoint{4.900000in}{4.900000in}}%
\pgfusepath{clip}%
\pgfsetbuttcap%
\pgfsetroundjoin%
\definecolor{currentfill}{rgb}{0.068281,0.068281,0.068281}%
\pgfsetfillcolor{currentfill}%
\pgfsetfillopacity{0.900000}%
\pgfsetlinewidth{0.501875pt}%
\definecolor{currentstroke}{rgb}{0.200000,0.200000,0.200000}%
\pgfsetstrokecolor{currentstroke}%
\pgfsetdash{}{0pt}%
\pgfpathmoveto{\pgfqpoint{0.916007in}{3.537380in}}%
\pgfpathlineto{\pgfqpoint{0.966012in}{3.489697in}}%
\pgfpathlineto{\pgfqpoint{0.996059in}{3.470693in}}%
\pgfpathlineto{\pgfqpoint{0.946040in}{3.518410in}}%
\pgfpathlineto{\pgfqpoint{0.916007in}{3.537380in}}%
\pgfpathclose%
\pgfusepath{stroke,fill}%
\end{pgfscope}%
\begin{pgfscope}%
\pgfpathrectangle{\pgfqpoint{0.050000in}{0.050000in}}{\pgfqpoint{4.900000in}{4.900000in}}%
\pgfusepath{clip}%
\pgfsetbuttcap%
\pgfsetroundjoin%
\definecolor{currentfill}{rgb}{0.878570,0.878570,0.878570}%
\pgfsetfillcolor{currentfill}%
\pgfsetfillopacity{0.900000}%
\pgfsetlinewidth{0.501875pt}%
\definecolor{currentstroke}{rgb}{0.200000,0.200000,0.200000}%
\pgfsetstrokecolor{currentstroke}%
\pgfsetdash{}{0pt}%
\pgfpathmoveto{\pgfqpoint{3.187707in}{1.717173in}}%
\pgfpathlineto{\pgfqpoint{3.232900in}{1.681137in}}%
\pgfpathlineto{\pgfqpoint{3.266539in}{1.660306in}}%
\pgfpathlineto{\pgfqpoint{3.221340in}{1.696157in}}%
\pgfpathlineto{\pgfqpoint{3.187707in}{1.717173in}}%
\pgfpathclose%
\pgfusepath{stroke,fill}%
\end{pgfscope}%
\begin{pgfscope}%
\pgfpathrectangle{\pgfqpoint{0.050000in}{0.050000in}}{\pgfqpoint{4.900000in}{4.900000in}}%
\pgfusepath{clip}%
\pgfsetbuttcap%
\pgfsetroundjoin%
\definecolor{currentfill}{rgb}{0.367243,0.367243,0.367243}%
\pgfsetfillcolor{currentfill}%
\pgfsetfillopacity{0.900000}%
\pgfsetlinewidth{0.501875pt}%
\definecolor{currentstroke}{rgb}{0.200000,0.200000,0.200000}%
\pgfsetstrokecolor{currentstroke}%
\pgfsetdash{}{0pt}%
\pgfpathmoveto{\pgfqpoint{1.646812in}{2.946723in}}%
\pgfpathlineto{\pgfqpoint{1.695243in}{2.900311in}}%
\pgfpathlineto{\pgfqpoint{1.726444in}{2.880377in}}%
\pgfpathlineto{\pgfqpoint{1.678003in}{2.926819in}}%
\pgfpathlineto{\pgfqpoint{1.646812in}{2.946723in}}%
\pgfpathclose%
\pgfusepath{stroke,fill}%
\end{pgfscope}%
\begin{pgfscope}%
\pgfpathrectangle{\pgfqpoint{0.050000in}{0.050000in}}{\pgfqpoint{4.900000in}{4.900000in}}%
\pgfusepath{clip}%
\pgfsetbuttcap%
\pgfsetroundjoin%
\definecolor{currentfill}{rgb}{0.964937,0.964937,0.964937}%
\pgfsetfillcolor{currentfill}%
\pgfsetfillopacity{0.900000}%
\pgfsetlinewidth{0.501875pt}%
\definecolor{currentstroke}{rgb}{0.200000,0.200000,0.200000}%
\pgfsetstrokecolor{currentstroke}%
\pgfsetdash{}{0pt}%
\pgfpathmoveto{\pgfqpoint{3.774041in}{1.367791in}}%
\pgfpathlineto{\pgfqpoint{3.818560in}{1.355412in}}%
\pgfpathlineto{\pgfqpoint{3.853166in}{1.334999in}}%
\pgfpathlineto{\pgfqpoint{3.808629in}{1.347314in}}%
\pgfpathlineto{\pgfqpoint{3.774041in}{1.367791in}}%
\pgfpathclose%
\pgfusepath{stroke,fill}%
\end{pgfscope}%
\begin{pgfscope}%
\pgfpathrectangle{\pgfqpoint{0.050000in}{0.050000in}}{\pgfqpoint{4.900000in}{4.900000in}}%
\pgfusepath{clip}%
\pgfsetbuttcap%
\pgfsetroundjoin%
\definecolor{currentfill}{rgb}{0.957555,0.957555,0.957555}%
\pgfsetfillcolor{currentfill}%
\pgfsetfillopacity{0.900000}%
\pgfsetlinewidth{0.501875pt}%
\definecolor{currentstroke}{rgb}{0.200000,0.200000,0.200000}%
\pgfsetstrokecolor{currentstroke}%
\pgfsetdash{}{0pt}%
\pgfpathmoveto{\pgfqpoint{3.694987in}{1.404095in}}%
\pgfpathlineto{\pgfqpoint{3.739564in}{1.388188in}}%
\pgfpathlineto{\pgfqpoint{3.774041in}{1.367791in}}%
\pgfpathlineto{\pgfqpoint{3.729448in}{1.383605in}}%
\pgfpathlineto{\pgfqpoint{3.694987in}{1.404095in}}%
\pgfpathclose%
\pgfusepath{stroke,fill}%
\end{pgfscope}%
\begin{pgfscope}%
\pgfpathrectangle{\pgfqpoint{0.050000in}{0.050000in}}{\pgfqpoint{4.900000in}{4.900000in}}%
\pgfusepath{clip}%
\pgfsetbuttcap%
\pgfsetroundjoin%
\definecolor{currentfill}{rgb}{0.629020,0.629020,0.629020}%
\pgfsetfillcolor{currentfill}%
\pgfsetfillopacity{0.900000}%
\pgfsetlinewidth{0.501875pt}%
\definecolor{currentstroke}{rgb}{0.200000,0.200000,0.200000}%
\pgfsetstrokecolor{currentstroke}%
\pgfsetdash{}{0pt}%
\pgfpathmoveto{\pgfqpoint{2.377716in}{2.356018in}}%
\pgfpathlineto{\pgfqpoint{2.424574in}{2.310900in}}%
\pgfpathlineto{\pgfqpoint{2.456927in}{2.290045in}}%
\pgfpathlineto{\pgfqpoint{2.410065in}{2.335184in}}%
\pgfpathlineto{\pgfqpoint{2.377716in}{2.356018in}}%
\pgfpathclose%
\pgfusepath{stroke,fill}%
\end{pgfscope}%
\begin{pgfscope}%
\pgfpathrectangle{\pgfqpoint{0.050000in}{0.050000in}}{\pgfqpoint{4.900000in}{4.900000in}}%
\pgfusepath{clip}%
\pgfsetbuttcap%
\pgfsetroundjoin%
\definecolor{currentfill}{rgb}{0.950173,0.950173,0.950173}%
\pgfsetfillcolor{currentfill}%
\pgfsetfillopacity{0.900000}%
\pgfsetlinewidth{0.501875pt}%
\definecolor{currentstroke}{rgb}{0.200000,0.200000,0.200000}%
\pgfsetstrokecolor{currentstroke}%
\pgfsetdash{}{0pt}%
\pgfpathmoveto{\pgfqpoint{3.615993in}{1.444033in}}%
\pgfpathlineto{\pgfqpoint{3.660637in}{1.424507in}}%
\pgfpathlineto{\pgfqpoint{3.694987in}{1.404095in}}%
\pgfpathlineto{\pgfqpoint{3.650329in}{1.423497in}}%
\pgfpathlineto{\pgfqpoint{3.615993in}{1.444033in}}%
\pgfpathclose%
\pgfusepath{stroke,fill}%
\end{pgfscope}%
\begin{pgfscope}%
\pgfpathrectangle{\pgfqpoint{0.050000in}{0.050000in}}{\pgfqpoint{4.900000in}{4.900000in}}%
\pgfusepath{clip}%
\pgfsetbuttcap%
\pgfsetroundjoin%
\definecolor{currentfill}{rgb}{0.262053,0.262053,0.262053}%
\pgfsetfillcolor{currentfill}%
\pgfsetfillopacity{0.900000}%
\pgfsetlinewidth{0.501875pt}%
\definecolor{currentstroke}{rgb}{0.200000,0.200000,0.200000}%
\pgfsetstrokecolor{currentstroke}%
\pgfsetdash{}{0pt}%
\pgfpathmoveto{\pgfqpoint{1.376587in}{3.165785in}}%
\pgfpathlineto{\pgfqpoint{1.425611in}{3.118893in}}%
\pgfpathlineto{\pgfqpoint{1.456391in}{3.099300in}}%
\pgfpathlineto{\pgfqpoint{1.407357in}{3.146224in}}%
\pgfpathlineto{\pgfqpoint{1.376587in}{3.165785in}}%
\pgfpathclose%
\pgfusepath{stroke,fill}%
\end{pgfscope}%
\begin{pgfscope}%
\pgfpathrectangle{\pgfqpoint{0.050000in}{0.050000in}}{\pgfqpoint{4.900000in}{4.900000in}}%
\pgfusepath{clip}%
\pgfsetbuttcap%
\pgfsetroundjoin%
\definecolor{currentfill}{rgb}{0.525798,0.525798,0.525798}%
\pgfsetfillcolor{currentfill}%
\pgfsetfillopacity{0.900000}%
\pgfsetlinewidth{0.501875pt}%
\definecolor{currentstroke}{rgb}{0.200000,0.200000,0.200000}%
\pgfsetstrokecolor{currentstroke}%
\pgfsetdash{}{0pt}%
\pgfpathmoveto{\pgfqpoint{2.107599in}{2.574973in}}%
\pgfpathlineto{\pgfqpoint{2.155048in}{2.529353in}}%
\pgfpathlineto{\pgfqpoint{2.186981in}{2.508831in}}%
\pgfpathlineto{\pgfqpoint{2.139527in}{2.554478in}}%
\pgfpathlineto{\pgfqpoint{2.107599in}{2.574973in}}%
\pgfpathclose%
\pgfusepath{stroke,fill}%
\end{pgfscope}%
\begin{pgfscope}%
\pgfpathrectangle{\pgfqpoint{0.050000in}{0.050000in}}{\pgfqpoint{4.900000in}{4.900000in}}%
\pgfusepath{clip}%
\pgfsetbuttcap%
\pgfsetroundjoin%
\definecolor{currentfill}{rgb}{0.940823,0.940823,0.940823}%
\pgfsetfillcolor{currentfill}%
\pgfsetfillopacity{0.900000}%
\pgfsetlinewidth{0.501875pt}%
\definecolor{currentstroke}{rgb}{0.200000,0.200000,0.200000}%
\pgfsetstrokecolor{currentstroke}%
\pgfsetdash{}{0pt}%
\pgfpathmoveto{\pgfqpoint{3.537043in}{1.487641in}}%
\pgfpathlineto{\pgfqpoint{3.581768in}{1.464492in}}%
\pgfpathlineto{\pgfqpoint{3.615993in}{1.444033in}}%
\pgfpathlineto{\pgfqpoint{3.571255in}{1.467030in}}%
\pgfpathlineto{\pgfqpoint{3.537043in}{1.487641in}}%
\pgfpathclose%
\pgfusepath{stroke,fill}%
\end{pgfscope}%
\begin{pgfscope}%
\pgfpathrectangle{\pgfqpoint{0.050000in}{0.050000in}}{\pgfqpoint{4.900000in}{4.900000in}}%
\pgfusepath{clip}%
\pgfsetbuttcap%
\pgfsetroundjoin%
\definecolor{currentfill}{rgb}{0.136563,0.136563,0.136563}%
\pgfsetfillcolor{currentfill}%
\pgfsetfillopacity{0.900000}%
\pgfsetlinewidth{0.501875pt}%
\definecolor{currentstroke}{rgb}{0.200000,0.200000,0.200000}%
\pgfsetstrokecolor{currentstroke}%
\pgfsetdash{}{0pt}%
\pgfpathmoveto{\pgfqpoint{1.106270in}{3.384926in}}%
\pgfpathlineto{\pgfqpoint{1.155885in}{3.337553in}}%
\pgfpathlineto{\pgfqpoint{1.186245in}{3.318301in}}%
\pgfpathlineto{\pgfqpoint{1.136618in}{3.365708in}}%
\pgfpathlineto{\pgfqpoint{1.106270in}{3.384926in}}%
\pgfpathclose%
\pgfusepath{stroke,fill}%
\end{pgfscope}%
\begin{pgfscope}%
\pgfpathrectangle{\pgfqpoint{0.050000in}{0.050000in}}{\pgfqpoint{4.900000in}{4.900000in}}%
\pgfusepath{clip}%
\pgfsetbuttcap%
\pgfsetroundjoin%
\definecolor{currentfill}{rgb}{0.788112,0.788112,0.788112}%
\pgfsetfillcolor{currentfill}%
\pgfsetfillopacity{0.900000}%
\pgfsetlinewidth{0.501875pt}%
\definecolor{currentstroke}{rgb}{0.200000,0.200000,0.200000}%
\pgfsetstrokecolor{currentstroke}%
\pgfsetdash{}{0pt}%
\pgfpathmoveto{\pgfqpoint{2.838720in}{1.986037in}}%
\pgfpathlineto{\pgfqpoint{2.884604in}{1.943093in}}%
\pgfpathlineto{\pgfqpoint{2.917695in}{1.921936in}}%
\pgfpathlineto{\pgfqpoint{2.871808in}{1.964804in}}%
\pgfpathlineto{\pgfqpoint{2.838720in}{1.986037in}}%
\pgfpathclose%
\pgfusepath{stroke,fill}%
\end{pgfscope}%
\begin{pgfscope}%
\pgfpathrectangle{\pgfqpoint{0.050000in}{0.050000in}}{\pgfqpoint{4.900000in}{4.900000in}}%
\pgfusepath{clip}%
\pgfsetbuttcap%
\pgfsetroundjoin%
\definecolor{currentfill}{rgb}{0.858762,0.858762,0.858762}%
\pgfsetfillcolor{currentfill}%
\pgfsetfillopacity{0.900000}%
\pgfsetlinewidth{0.501875pt}%
\definecolor{currentstroke}{rgb}{0.200000,0.200000,0.200000}%
\pgfsetstrokecolor{currentstroke}%
\pgfsetdash{}{0pt}%
\pgfpathmoveto{\pgfqpoint{3.108830in}{1.776490in}}%
\pgfpathlineto{\pgfqpoint{3.154181in}{1.738131in}}%
\pgfpathlineto{\pgfqpoint{3.187707in}{1.717173in}}%
\pgfpathlineto{\pgfqpoint{3.142351in}{1.755364in}}%
\pgfpathlineto{\pgfqpoint{3.108830in}{1.776490in}}%
\pgfpathclose%
\pgfusepath{stroke,fill}%
\end{pgfscope}%
\begin{pgfscope}%
\pgfpathrectangle{\pgfqpoint{0.050000in}{0.050000in}}{\pgfqpoint{4.900000in}{4.900000in}}%
\pgfusepath{clip}%
\pgfsetbuttcap%
\pgfsetroundjoin%
\definecolor{currentfill}{rgb}{0.432203,0.432203,0.432203}%
\pgfsetfillcolor{currentfill}%
\pgfsetfillopacity{0.900000}%
\pgfsetlinewidth{0.501875pt}%
\definecolor{currentstroke}{rgb}{0.200000,0.200000,0.200000}%
\pgfsetstrokecolor{currentstroke}%
\pgfsetdash{}{0pt}%
\pgfpathmoveto{\pgfqpoint{1.837389in}{2.794020in}}%
\pgfpathlineto{\pgfqpoint{1.885429in}{2.747919in}}%
\pgfpathlineto{\pgfqpoint{1.916943in}{2.727737in}}%
\pgfpathlineto{\pgfqpoint{1.868895in}{2.773868in}}%
\pgfpathlineto{\pgfqpoint{1.837389in}{2.794020in}}%
\pgfpathclose%
\pgfusepath{stroke,fill}%
\end{pgfscope}%
\begin{pgfscope}%
\pgfpathrectangle{\pgfqpoint{0.050000in}{0.050000in}}{\pgfqpoint{4.900000in}{4.900000in}}%
\pgfusepath{clip}%
\pgfsetbuttcap%
\pgfsetroundjoin%
\definecolor{currentfill}{rgb}{0.700992,0.700992,0.700992}%
\pgfsetfillcolor{currentfill}%
\pgfsetfillopacity{0.900000}%
\pgfsetlinewidth{0.501875pt}%
\definecolor{currentstroke}{rgb}{0.200000,0.200000,0.200000}%
\pgfsetstrokecolor{currentstroke}%
\pgfsetdash{}{0pt}%
\pgfpathmoveto{\pgfqpoint{2.568608in}{2.203207in}}%
\pgfpathlineto{\pgfqpoint{2.615073in}{2.158551in}}%
\pgfpathlineto{\pgfqpoint{2.647741in}{2.137491in}}%
\pgfpathlineto{\pgfqpoint{2.601272in}{2.182147in}}%
\pgfpathlineto{\pgfqpoint{2.568608in}{2.203207in}}%
\pgfpathclose%
\pgfusepath{stroke,fill}%
\end{pgfscope}%
\begin{pgfscope}%
\pgfpathrectangle{\pgfqpoint{0.050000in}{0.050000in}}{\pgfqpoint{4.900000in}{4.900000in}}%
\pgfusepath{clip}%
\pgfsetbuttcap%
\pgfsetroundjoin%
\definecolor{currentfill}{rgb}{0.926674,0.926674,0.926674}%
\pgfsetfillcolor{currentfill}%
\pgfsetfillopacity{0.900000}%
\pgfsetlinewidth{0.501875pt}%
\definecolor{currentstroke}{rgb}{0.200000,0.200000,0.200000}%
\pgfsetstrokecolor{currentstroke}%
\pgfsetdash{}{0pt}%
\pgfpathmoveto{\pgfqpoint{3.458120in}{1.534862in}}%
\pgfpathlineto{\pgfqpoint{3.502940in}{1.508178in}}%
\pgfpathlineto{\pgfqpoint{3.537043in}{1.487641in}}%
\pgfpathlineto{\pgfqpoint{3.492212in}{1.514151in}}%
\pgfpathlineto{\pgfqpoint{3.458120in}{1.534862in}}%
\pgfpathclose%
\pgfusepath{stroke,fill}%
\end{pgfscope}%
\begin{pgfscope}%
\pgfpathrectangle{\pgfqpoint{0.050000in}{0.050000in}}{\pgfqpoint{4.900000in}{4.900000in}}%
\pgfusepath{clip}%
\pgfsetbuttcap%
\pgfsetroundjoin%
\definecolor{currentfill}{rgb}{0.031865,0.031865,0.031865}%
\pgfsetfillcolor{currentfill}%
\pgfsetfillopacity{0.900000}%
\pgfsetlinewidth{0.501875pt}%
\definecolor{currentstroke}{rgb}{0.200000,0.200000,0.200000}%
\pgfsetstrokecolor{currentstroke}%
\pgfsetdash{}{0pt}%
\pgfpathmoveto{\pgfqpoint{0.835860in}{3.604144in}}%
\pgfpathlineto{\pgfqpoint{0.886067in}{3.556290in}}%
\pgfpathlineto{\pgfqpoint{0.916007in}{3.537380in}}%
\pgfpathlineto{\pgfqpoint{0.865785in}{3.585269in}}%
\pgfpathlineto{\pgfqpoint{0.835860in}{3.604144in}}%
\pgfpathclose%
\pgfusepath{stroke,fill}%
\end{pgfscope}%
\begin{pgfscope}%
\pgfpathrectangle{\pgfqpoint{0.050000in}{0.050000in}}{\pgfqpoint{4.900000in}{4.900000in}}%
\pgfusepath{clip}%
\pgfsetbuttcap%
\pgfsetroundjoin%
\definecolor{currentfill}{rgb}{0.338824,0.338824,0.338824}%
\pgfsetfillcolor{currentfill}%
\pgfsetfillopacity{0.900000}%
\pgfsetlinewidth{0.501875pt}%
\definecolor{currentstroke}{rgb}{0.200000,0.200000,0.200000}%
\pgfsetstrokecolor{currentstroke}%
\pgfsetdash{}{0pt}%
\pgfpathmoveto{\pgfqpoint{1.567087in}{3.013145in}}%
\pgfpathlineto{\pgfqpoint{1.615719in}{2.966563in}}%
\pgfpathlineto{\pgfqpoint{1.646812in}{2.946723in}}%
\pgfpathlineto{\pgfqpoint{1.598170in}{2.993335in}}%
\pgfpathlineto{\pgfqpoint{1.567087in}{3.013145in}}%
\pgfpathclose%
\pgfusepath{stroke,fill}%
\end{pgfscope}%
\begin{pgfscope}%
\pgfpathrectangle{\pgfqpoint{0.050000in}{0.050000in}}{\pgfqpoint{4.900000in}{4.900000in}}%
\pgfusepath{clip}%
\pgfsetbuttcap%
\pgfsetroundjoin%
\definecolor{currentfill}{rgb}{0.590634,0.590634,0.590634}%
\pgfsetfillcolor{currentfill}%
\pgfsetfillopacity{0.900000}%
\pgfsetlinewidth{0.501875pt}%
\definecolor{currentstroke}{rgb}{0.200000,0.200000,0.200000}%
\pgfsetstrokecolor{currentstroke}%
\pgfsetdash{}{0pt}%
\pgfpathmoveto{\pgfqpoint{2.298412in}{2.422088in}}%
\pgfpathlineto{\pgfqpoint{2.345469in}{2.376786in}}%
\pgfpathlineto{\pgfqpoint{2.377716in}{2.356018in}}%
\pgfpathlineto{\pgfqpoint{2.330655in}{2.401345in}}%
\pgfpathlineto{\pgfqpoint{2.298412in}{2.422088in}}%
\pgfpathclose%
\pgfusepath{stroke,fill}%
\end{pgfscope}%
\begin{pgfscope}%
\pgfpathrectangle{\pgfqpoint{0.050000in}{0.050000in}}{\pgfqpoint{4.900000in}{4.900000in}}%
\pgfusepath{clip}%
\pgfsetbuttcap%
\pgfsetroundjoin%
\definecolor{currentfill}{rgb}{0.217762,0.217762,0.217762}%
\pgfsetfillcolor{currentfill}%
\pgfsetfillopacity{0.900000}%
\pgfsetlinewidth{0.501875pt}%
\definecolor{currentstroke}{rgb}{0.200000,0.200000,0.200000}%
\pgfsetstrokecolor{currentstroke}%
\pgfsetdash{}{0pt}%
\pgfpathmoveto{\pgfqpoint{1.296691in}{3.232348in}}%
\pgfpathlineto{\pgfqpoint{1.345915in}{3.185285in}}%
\pgfpathlineto{\pgfqpoint{1.376587in}{3.165785in}}%
\pgfpathlineto{\pgfqpoint{1.327352in}{3.212881in}}%
\pgfpathlineto{\pgfqpoint{1.296691in}{3.232348in}}%
\pgfpathclose%
\pgfusepath{stroke,fill}%
\end{pgfscope}%
\begin{pgfscope}%
\pgfpathrectangle{\pgfqpoint{0.050000in}{0.050000in}}{\pgfqpoint{4.900000in}{4.900000in}}%
\pgfusepath{clip}%
\pgfsetbuttcap%
\pgfsetroundjoin%
\definecolor{currentfill}{rgb}{0.912526,0.912526,0.912526}%
\pgfsetfillcolor{currentfill}%
\pgfsetfillopacity{0.900000}%
\pgfsetlinewidth{0.501875pt}%
\definecolor{currentstroke}{rgb}{0.200000,0.200000,0.200000}%
\pgfsetstrokecolor{currentstroke}%
\pgfsetdash{}{0pt}%
\pgfpathmoveto{\pgfqpoint{3.379208in}{1.585538in}}%
\pgfpathlineto{\pgfqpoint{3.424139in}{1.555503in}}%
\pgfpathlineto{\pgfqpoint{3.458120in}{1.534862in}}%
\pgfpathlineto{\pgfqpoint{3.413180in}{1.564710in}}%
\pgfpathlineto{\pgfqpoint{3.379208in}{1.585538in}}%
\pgfpathclose%
\pgfusepath{stroke,fill}%
\end{pgfscope}%
\begin{pgfscope}%
\pgfpathrectangle{\pgfqpoint{0.050000in}{0.050000in}}{\pgfqpoint{4.900000in}{4.900000in}}%
\pgfusepath{clip}%
\pgfsetbuttcap%
\pgfsetroundjoin%
\definecolor{currentfill}{rgb}{0.495656,0.495656,0.495656}%
\pgfsetfillcolor{currentfill}%
\pgfsetfillopacity{0.900000}%
\pgfsetlinewidth{0.501875pt}%
\definecolor{currentstroke}{rgb}{0.200000,0.200000,0.200000}%
\pgfsetstrokecolor{currentstroke}%
\pgfsetdash{}{0pt}%
\pgfpathmoveto{\pgfqpoint{2.028124in}{2.641194in}}%
\pgfpathlineto{\pgfqpoint{2.075772in}{2.595404in}}%
\pgfpathlineto{\pgfqpoint{2.107599in}{2.574973in}}%
\pgfpathlineto{\pgfqpoint{2.059944in}{2.620791in}}%
\pgfpathlineto{\pgfqpoint{2.028124in}{2.641194in}}%
\pgfpathclose%
\pgfusepath{stroke,fill}%
\end{pgfscope}%
\begin{pgfscope}%
\pgfpathrectangle{\pgfqpoint{0.050000in}{0.050000in}}{\pgfqpoint{4.900000in}{4.900000in}}%
\pgfusepath{clip}%
\pgfsetbuttcap%
\pgfsetroundjoin%
\definecolor{currentfill}{rgb}{0.839785,0.839785,0.839785}%
\pgfsetfillcolor{currentfill}%
\pgfsetfillopacity{0.900000}%
\pgfsetlinewidth{0.501875pt}%
\definecolor{currentstroke}{rgb}{0.200000,0.200000,0.200000}%
\pgfsetstrokecolor{currentstroke}%
\pgfsetdash{}{0pt}%
\pgfpathmoveto{\pgfqpoint{3.029895in}{1.837819in}}%
\pgfpathlineto{\pgfqpoint{3.075415in}{1.797561in}}%
\pgfpathlineto{\pgfqpoint{3.108830in}{1.776490in}}%
\pgfpathlineto{\pgfqpoint{3.063305in}{1.816607in}}%
\pgfpathlineto{\pgfqpoint{3.029895in}{1.837819in}}%
\pgfpathclose%
\pgfusepath{stroke,fill}%
\end{pgfscope}%
\begin{pgfscope}%
\pgfpathrectangle{\pgfqpoint{0.050000in}{0.050000in}}{\pgfqpoint{4.900000in}{4.900000in}}%
\pgfusepath{clip}%
\pgfsetbuttcap%
\pgfsetroundjoin%
\definecolor{currentfill}{rgb}{0.763998,0.763998,0.763998}%
\pgfsetfillcolor{currentfill}%
\pgfsetfillopacity{0.900000}%
\pgfsetlinewidth{0.501875pt}%
\definecolor{currentstroke}{rgb}{0.200000,0.200000,0.200000}%
\pgfsetstrokecolor{currentstroke}%
\pgfsetdash{}{0pt}%
\pgfpathmoveto{\pgfqpoint{2.759660in}{2.050937in}}%
\pgfpathlineto{\pgfqpoint{2.805736in}{2.007213in}}%
\pgfpathlineto{\pgfqpoint{2.838720in}{1.986037in}}%
\pgfpathlineto{\pgfqpoint{2.792642in}{2.029715in}}%
\pgfpathlineto{\pgfqpoint{2.759660in}{2.050937in}}%
\pgfpathclose%
\pgfusepath{stroke,fill}%
\end{pgfscope}%
\begin{pgfscope}%
\pgfpathrectangle{\pgfqpoint{0.050000in}{0.050000in}}{\pgfqpoint{4.900000in}{4.900000in}}%
\pgfusepath{clip}%
\pgfsetbuttcap%
\pgfsetroundjoin%
\definecolor{currentfill}{rgb}{0.104698,0.104698,0.104698}%
\pgfsetfillcolor{currentfill}%
\pgfsetfillopacity{0.900000}%
\pgfsetlinewidth{0.501875pt}%
\definecolor{currentstroke}{rgb}{0.200000,0.200000,0.200000}%
\pgfsetstrokecolor{currentstroke}%
\pgfsetdash{}{0pt}%
\pgfpathmoveto{\pgfqpoint{1.026202in}{3.451629in}}%
\pgfpathlineto{\pgfqpoint{1.076019in}{3.404084in}}%
\pgfpathlineto{\pgfqpoint{1.106270in}{3.384926in}}%
\pgfpathlineto{\pgfqpoint{1.056440in}{3.432505in}}%
\pgfpathlineto{\pgfqpoint{1.026202in}{3.451629in}}%
\pgfpathclose%
\pgfusepath{stroke,fill}%
\end{pgfscope}%
\begin{pgfscope}%
\pgfpathrectangle{\pgfqpoint{0.050000in}{0.050000in}}{\pgfqpoint{4.900000in}{4.900000in}}%
\pgfusepath{clip}%
\pgfsetbuttcap%
\pgfsetroundjoin%
\definecolor{currentfill}{rgb}{0.399723,0.399723,0.399723}%
\pgfsetfillcolor{currentfill}%
\pgfsetfillopacity{0.900000}%
\pgfsetlinewidth{0.501875pt}%
\definecolor{currentstroke}{rgb}{0.200000,0.200000,0.200000}%
\pgfsetstrokecolor{currentstroke}%
\pgfsetdash{}{0pt}%
\pgfpathmoveto{\pgfqpoint{1.757743in}{2.860381in}}%
\pgfpathlineto{\pgfqpoint{1.805983in}{2.814110in}}%
\pgfpathlineto{\pgfqpoint{1.837389in}{2.794020in}}%
\pgfpathlineto{\pgfqpoint{1.789141in}{2.840321in}}%
\pgfpathlineto{\pgfqpoint{1.757743in}{2.860381in}}%
\pgfpathclose%
\pgfusepath{stroke,fill}%
\end{pgfscope}%
\begin{pgfscope}%
\pgfpathrectangle{\pgfqpoint{0.050000in}{0.050000in}}{\pgfqpoint{4.900000in}{4.900000in}}%
\pgfusepath{clip}%
\pgfsetbuttcap%
\pgfsetroundjoin%
\definecolor{currentfill}{rgb}{0.662607,0.662607,0.662607}%
\pgfsetfillcolor{currentfill}%
\pgfsetfillopacity{0.900000}%
\pgfsetlinewidth{0.501875pt}%
\definecolor{currentstroke}{rgb}{0.200000,0.200000,0.200000}%
\pgfsetstrokecolor{currentstroke}%
\pgfsetdash{}{0pt}%
\pgfpathmoveto{\pgfqpoint{2.489383in}{2.269127in}}%
\pgfpathlineto{\pgfqpoint{2.536046in}{2.224202in}}%
\pgfpathlineto{\pgfqpoint{2.568608in}{2.203207in}}%
\pgfpathlineto{\pgfqpoint{2.521941in}{2.248144in}}%
\pgfpathlineto{\pgfqpoint{2.489383in}{2.269127in}}%
\pgfpathclose%
\pgfusepath{stroke,fill}%
\end{pgfscope}%
\begin{pgfscope}%
\pgfpathrectangle{\pgfqpoint{0.050000in}{0.050000in}}{\pgfqpoint{4.900000in}{4.900000in}}%
\pgfusepath{clip}%
\pgfsetbuttcap%
\pgfsetroundjoin%
\definecolor{currentfill}{rgb}{0.000000,0.000000,0.000000}%
\pgfsetfillcolor{currentfill}%
\pgfsetfillopacity{0.900000}%
\pgfsetlinewidth{0.501875pt}%
\definecolor{currentstroke}{rgb}{0.200000,0.200000,0.200000}%
\pgfsetstrokecolor{currentstroke}%
\pgfsetdash{}{0pt}%
\pgfpathmoveto{\pgfqpoint{0.755620in}{3.670987in}}%
\pgfpathlineto{\pgfqpoint{0.806030in}{3.622961in}}%
\pgfpathlineto{\pgfqpoint{0.835860in}{3.604144in}}%
\pgfpathlineto{\pgfqpoint{0.785436in}{3.652207in}}%
\pgfpathlineto{\pgfqpoint{0.755620in}{3.670987in}}%
\pgfpathclose%
\pgfusepath{stroke,fill}%
\end{pgfscope}%
\begin{pgfscope}%
\pgfpathrectangle{\pgfqpoint{0.050000in}{0.050000in}}{\pgfqpoint{4.900000in}{4.900000in}}%
\pgfusepath{clip}%
\pgfsetbuttcap%
\pgfsetroundjoin%
\definecolor{currentfill}{rgb}{0.895548,0.895548,0.895548}%
\pgfsetfillcolor{currentfill}%
\pgfsetfillopacity{0.900000}%
\pgfsetlinewidth{0.501875pt}%
\definecolor{currentstroke}{rgb}{0.200000,0.200000,0.200000}%
\pgfsetstrokecolor{currentstroke}%
\pgfsetdash{}{0pt}%
\pgfpathmoveto{\pgfqpoint{3.300287in}{1.639415in}}%
\pgfpathlineto{\pgfqpoint{3.345344in}{1.606300in}}%
\pgfpathlineto{\pgfqpoint{3.379208in}{1.585538in}}%
\pgfpathlineto{\pgfqpoint{3.334142in}{1.618463in}}%
\pgfpathlineto{\pgfqpoint{3.300287in}{1.639415in}}%
\pgfpathclose%
\pgfusepath{stroke,fill}%
\end{pgfscope}%
\begin{pgfscope}%
\pgfpathrectangle{\pgfqpoint{0.050000in}{0.050000in}}{\pgfqpoint{4.900000in}{4.900000in}}%
\pgfusepath{clip}%
\pgfsetbuttcap%
\pgfsetroundjoin%
\definecolor{currentfill}{rgb}{0.300807,0.300807,0.300807}%
\pgfsetfillcolor{currentfill}%
\pgfsetfillopacity{0.900000}%
\pgfsetlinewidth{0.501875pt}%
\definecolor{currentstroke}{rgb}{0.200000,0.200000,0.200000}%
\pgfsetstrokecolor{currentstroke}%
\pgfsetdash{}{0pt}%
\pgfpathmoveto{\pgfqpoint{1.487268in}{3.079646in}}%
\pgfpathlineto{\pgfqpoint{1.536101in}{3.032893in}}%
\pgfpathlineto{\pgfqpoint{1.567087in}{3.013145in}}%
\pgfpathlineto{\pgfqpoint{1.518244in}{3.059929in}}%
\pgfpathlineto{\pgfqpoint{1.487268in}{3.079646in}}%
\pgfpathclose%
\pgfusepath{stroke,fill}%
\end{pgfscope}%
\begin{pgfscope}%
\pgfpathrectangle{\pgfqpoint{0.050000in}{0.050000in}}{\pgfqpoint{4.900000in}{4.900000in}}%
\pgfusepath{clip}%
\pgfsetbuttcap%
\pgfsetroundjoin%
\definecolor{currentfill}{rgb}{0.560246,0.560246,0.560246}%
\pgfsetfillcolor{currentfill}%
\pgfsetfillopacity{0.900000}%
\pgfsetlinewidth{0.501875pt}%
\definecolor{currentstroke}{rgb}{0.200000,0.200000,0.200000}%
\pgfsetstrokecolor{currentstroke}%
\pgfsetdash{}{0pt}%
\pgfpathmoveto{\pgfqpoint{2.219016in}{2.488243in}}%
\pgfpathlineto{\pgfqpoint{2.266271in}{2.442767in}}%
\pgfpathlineto{\pgfqpoint{2.298412in}{2.422088in}}%
\pgfpathlineto{\pgfqpoint{2.251152in}{2.467591in}}%
\pgfpathlineto{\pgfqpoint{2.219016in}{2.488243in}}%
\pgfpathclose%
\pgfusepath{stroke,fill}%
\end{pgfscope}%
\begin{pgfscope}%
\pgfpathrectangle{\pgfqpoint{0.050000in}{0.050000in}}{\pgfqpoint{4.900000in}{4.900000in}}%
\pgfusepath{clip}%
\pgfsetbuttcap%
\pgfsetroundjoin%
\definecolor{currentfill}{rgb}{0.179008,0.179008,0.179008}%
\pgfsetfillcolor{currentfill}%
\pgfsetfillopacity{0.900000}%
\pgfsetlinewidth{0.501875pt}%
\definecolor{currentstroke}{rgb}{0.200000,0.200000,0.200000}%
\pgfsetstrokecolor{currentstroke}%
\pgfsetdash{}{0pt}%
\pgfpathmoveto{\pgfqpoint{1.216701in}{3.298989in}}%
\pgfpathlineto{\pgfqpoint{1.266127in}{3.251754in}}%
\pgfpathlineto{\pgfqpoint{1.296691in}{3.232348in}}%
\pgfpathlineto{\pgfqpoint{1.247253in}{3.279616in}}%
\pgfpathlineto{\pgfqpoint{1.216701in}{3.298989in}}%
\pgfpathclose%
\pgfusepath{stroke,fill}%
\end{pgfscope}%
\begin{pgfscope}%
\pgfpathrectangle{\pgfqpoint{0.050000in}{0.050000in}}{\pgfqpoint{4.900000in}{4.900000in}}%
\pgfusepath{clip}%
\pgfsetbuttcap%
\pgfsetroundjoin%
\definecolor{currentfill}{rgb}{0.815671,0.815671,0.815671}%
\pgfsetfillcolor{currentfill}%
\pgfsetfillopacity{0.900000}%
\pgfsetlinewidth{0.501875pt}%
\definecolor{currentstroke}{rgb}{0.200000,0.200000,0.200000}%
\pgfsetstrokecolor{currentstroke}%
\pgfsetdash{}{0pt}%
\pgfpathmoveto{\pgfqpoint{2.950891in}{1.900724in}}%
\pgfpathlineto{\pgfqpoint{2.996590in}{1.858977in}}%
\pgfpathlineto{\pgfqpoint{3.029895in}{1.837819in}}%
\pgfpathlineto{\pgfqpoint{2.984192in}{1.879456in}}%
\pgfpathlineto{\pgfqpoint{2.950891in}{1.900724in}}%
\pgfpathclose%
\pgfusepath{stroke,fill}%
\end{pgfscope}%
\begin{pgfscope}%
\pgfpathrectangle{\pgfqpoint{0.050000in}{0.050000in}}{\pgfqpoint{4.900000in}{4.900000in}}%
\pgfusepath{clip}%
\pgfsetbuttcap%
\pgfsetroundjoin%
\definecolor{currentfill}{rgb}{0.461207,0.461207,0.461207}%
\pgfsetfillcolor{currentfill}%
\pgfsetfillopacity{0.900000}%
\pgfsetlinewidth{0.501875pt}%
\definecolor{currentstroke}{rgb}{0.200000,0.200000,0.200000}%
\pgfsetstrokecolor{currentstroke}%
\pgfsetdash{}{0pt}%
\pgfpathmoveto{\pgfqpoint{1.948556in}{2.707492in}}%
\pgfpathlineto{\pgfqpoint{1.996404in}{2.661532in}}%
\pgfpathlineto{\pgfqpoint{2.028124in}{2.641194in}}%
\pgfpathlineto{\pgfqpoint{1.980269in}{2.687182in}}%
\pgfpathlineto{\pgfqpoint{1.948556in}{2.707492in}}%
\pgfpathclose%
\pgfusepath{stroke,fill}%
\end{pgfscope}%
\begin{pgfscope}%
\pgfpathrectangle{\pgfqpoint{0.050000in}{0.050000in}}{\pgfqpoint{4.900000in}{4.900000in}}%
\pgfusepath{clip}%
\pgfsetbuttcap%
\pgfsetroundjoin%
\definecolor{currentfill}{rgb}{0.739377,0.739377,0.739377}%
\pgfsetfillcolor{currentfill}%
\pgfsetfillopacity{0.900000}%
\pgfsetlinewidth{0.501875pt}%
\definecolor{currentstroke}{rgb}{0.200000,0.200000,0.200000}%
\pgfsetstrokecolor{currentstroke}%
\pgfsetdash{}{0pt}%
\pgfpathmoveto{\pgfqpoint{2.680512in}{2.116370in}}%
\pgfpathlineto{\pgfqpoint{2.726783in}{2.072099in}}%
\pgfpathlineto{\pgfqpoint{2.759660in}{2.050937in}}%
\pgfpathlineto{\pgfqpoint{2.713387in}{2.095188in}}%
\pgfpathlineto{\pgfqpoint{2.680512in}{2.116370in}}%
\pgfpathclose%
\pgfusepath{stroke,fill}%
\end{pgfscope}%
\begin{pgfscope}%
\pgfpathrectangle{\pgfqpoint{0.050000in}{0.050000in}}{\pgfqpoint{4.900000in}{4.900000in}}%
\pgfusepath{clip}%
\pgfsetbuttcap%
\pgfsetroundjoin%
\definecolor{currentfill}{rgb}{0.068281,0.068281,0.068281}%
\pgfsetfillcolor{currentfill}%
\pgfsetfillopacity{0.900000}%
\pgfsetlinewidth{0.501875pt}%
\definecolor{currentstroke}{rgb}{0.200000,0.200000,0.200000}%
\pgfsetstrokecolor{currentstroke}%
\pgfsetdash{}{0pt}%
\pgfpathmoveto{\pgfqpoint{0.946040in}{3.518410in}}%
\pgfpathlineto{\pgfqpoint{0.996059in}{3.470693in}}%
\pgfpathlineto{\pgfqpoint{1.026202in}{3.451629in}}%
\pgfpathlineto{\pgfqpoint{0.976169in}{3.499380in}}%
\pgfpathlineto{\pgfqpoint{0.946040in}{3.518410in}}%
\pgfpathclose%
\pgfusepath{stroke,fill}%
\end{pgfscope}%
\begin{pgfscope}%
\pgfpathrectangle{\pgfqpoint{0.050000in}{0.050000in}}{\pgfqpoint{4.900000in}{4.900000in}}%
\pgfusepath{clip}%
\pgfsetbuttcap%
\pgfsetroundjoin%
\definecolor{currentfill}{rgb}{0.878570,0.878570,0.878570}%
\pgfsetfillcolor{currentfill}%
\pgfsetfillopacity{0.900000}%
\pgfsetlinewidth{0.501875pt}%
\definecolor{currentstroke}{rgb}{0.200000,0.200000,0.200000}%
\pgfsetstrokecolor{currentstroke}%
\pgfsetdash{}{0pt}%
\pgfpathmoveto{\pgfqpoint{3.221340in}{1.696157in}}%
\pgfpathlineto{\pgfqpoint{3.266539in}{1.660306in}}%
\pgfpathlineto{\pgfqpoint{3.300287in}{1.639415in}}%
\pgfpathlineto{\pgfqpoint{3.255080in}{1.675082in}}%
\pgfpathlineto{\pgfqpoint{3.221340in}{1.696157in}}%
\pgfpathclose%
\pgfusepath{stroke,fill}%
\end{pgfscope}%
\begin{pgfscope}%
\pgfpathrectangle{\pgfqpoint{0.050000in}{0.050000in}}{\pgfqpoint{4.900000in}{4.900000in}}%
\pgfusepath{clip}%
\pgfsetbuttcap%
\pgfsetroundjoin%
\definecolor{currentfill}{rgb}{0.367243,0.367243,0.367243}%
\pgfsetfillcolor{currentfill}%
\pgfsetfillopacity{0.900000}%
\pgfsetlinewidth{0.501875pt}%
\definecolor{currentstroke}{rgb}{0.200000,0.200000,0.200000}%
\pgfsetstrokecolor{currentstroke}%
\pgfsetdash{}{0pt}%
\pgfpathmoveto{\pgfqpoint{1.678003in}{2.926819in}}%
\pgfpathlineto{\pgfqpoint{1.726444in}{2.880377in}}%
\pgfpathlineto{\pgfqpoint{1.757743in}{2.860381in}}%
\pgfpathlineto{\pgfqpoint{1.709293in}{2.906853in}}%
\pgfpathlineto{\pgfqpoint{1.678003in}{2.926819in}}%
\pgfpathclose%
\pgfusepath{stroke,fill}%
\end{pgfscope}%
\begin{pgfscope}%
\pgfpathrectangle{\pgfqpoint{0.050000in}{0.050000in}}{\pgfqpoint{4.900000in}{4.900000in}}%
\pgfusepath{clip}%
\pgfsetbuttcap%
\pgfsetroundjoin%
\definecolor{currentfill}{rgb}{0.629020,0.629020,0.629020}%
\pgfsetfillcolor{currentfill}%
\pgfsetfillopacity{0.900000}%
\pgfsetlinewidth{0.501875pt}%
\definecolor{currentstroke}{rgb}{0.200000,0.200000,0.200000}%
\pgfsetstrokecolor{currentstroke}%
\pgfsetdash{}{0pt}%
\pgfpathmoveto{\pgfqpoint{2.410065in}{2.335184in}}%
\pgfpathlineto{\pgfqpoint{2.456927in}{2.290045in}}%
\pgfpathlineto{\pgfqpoint{2.489383in}{2.269127in}}%
\pgfpathlineto{\pgfqpoint{2.442517in}{2.314286in}}%
\pgfpathlineto{\pgfqpoint{2.410065in}{2.335184in}}%
\pgfpathclose%
\pgfusepath{stroke,fill}%
\end{pgfscope}%
\begin{pgfscope}%
\pgfpathrectangle{\pgfqpoint{0.050000in}{0.050000in}}{\pgfqpoint{4.900000in}{4.900000in}}%
\pgfusepath{clip}%
\pgfsetbuttcap%
\pgfsetroundjoin%
\definecolor{currentfill}{rgb}{0.957555,0.957555,0.957555}%
\pgfsetfillcolor{currentfill}%
\pgfsetfillopacity{0.900000}%
\pgfsetlinewidth{0.501875pt}%
\definecolor{currentstroke}{rgb}{0.200000,0.200000,0.200000}%
\pgfsetstrokecolor{currentstroke}%
\pgfsetdash{}{0pt}%
\pgfpathmoveto{\pgfqpoint{3.729448in}{1.383605in}}%
\pgfpathlineto{\pgfqpoint{3.774041in}{1.367791in}}%
\pgfpathlineto{\pgfqpoint{3.808629in}{1.347314in}}%
\pgfpathlineto{\pgfqpoint{3.764021in}{1.363035in}}%
\pgfpathlineto{\pgfqpoint{3.729448in}{1.383605in}}%
\pgfpathclose%
\pgfusepath{stroke,fill}%
\end{pgfscope}%
\begin{pgfscope}%
\pgfpathrectangle{\pgfqpoint{0.050000in}{0.050000in}}{\pgfqpoint{4.900000in}{4.900000in}}%
\pgfusepath{clip}%
\pgfsetbuttcap%
\pgfsetroundjoin%
\definecolor{currentfill}{rgb}{0.950173,0.950173,0.950173}%
\pgfsetfillcolor{currentfill}%
\pgfsetfillopacity{0.900000}%
\pgfsetlinewidth{0.501875pt}%
\definecolor{currentstroke}{rgb}{0.200000,0.200000,0.200000}%
\pgfsetstrokecolor{currentstroke}%
\pgfsetdash{}{0pt}%
\pgfpathmoveto{\pgfqpoint{3.650329in}{1.423497in}}%
\pgfpathlineto{\pgfqpoint{3.694987in}{1.404095in}}%
\pgfpathlineto{\pgfqpoint{3.729448in}{1.383605in}}%
\pgfpathlineto{\pgfqpoint{3.684775in}{1.402884in}}%
\pgfpathlineto{\pgfqpoint{3.650329in}{1.423497in}}%
\pgfpathclose%
\pgfusepath{stroke,fill}%
\end{pgfscope}%
\begin{pgfscope}%
\pgfpathrectangle{\pgfqpoint{0.050000in}{0.050000in}}{\pgfqpoint{4.900000in}{4.900000in}}%
\pgfusepath{clip}%
\pgfsetbuttcap%
\pgfsetroundjoin%
\definecolor{currentfill}{rgb}{0.262053,0.262053,0.262053}%
\pgfsetfillcolor{currentfill}%
\pgfsetfillopacity{0.900000}%
\pgfsetlinewidth{0.501875pt}%
\definecolor{currentstroke}{rgb}{0.200000,0.200000,0.200000}%
\pgfsetstrokecolor{currentstroke}%
\pgfsetdash{}{0pt}%
\pgfpathmoveto{\pgfqpoint{1.407357in}{3.146224in}}%
\pgfpathlineto{\pgfqpoint{1.456391in}{3.099300in}}%
\pgfpathlineto{\pgfqpoint{1.487268in}{3.079646in}}%
\pgfpathlineto{\pgfqpoint{1.438224in}{3.126602in}}%
\pgfpathlineto{\pgfqpoint{1.407357in}{3.146224in}}%
\pgfpathclose%
\pgfusepath{stroke,fill}%
\end{pgfscope}%
\begin{pgfscope}%
\pgfpathrectangle{\pgfqpoint{0.050000in}{0.050000in}}{\pgfqpoint{4.900000in}{4.900000in}}%
\pgfusepath{clip}%
\pgfsetbuttcap%
\pgfsetroundjoin%
\definecolor{currentfill}{rgb}{0.525798,0.525798,0.525798}%
\pgfsetfillcolor{currentfill}%
\pgfsetfillopacity{0.900000}%
\pgfsetlinewidth{0.501875pt}%
\definecolor{currentstroke}{rgb}{0.200000,0.200000,0.200000}%
\pgfsetstrokecolor{currentstroke}%
\pgfsetdash{}{0pt}%
\pgfpathmoveto{\pgfqpoint{2.139527in}{2.554478in}}%
\pgfpathlineto{\pgfqpoint{2.186981in}{2.508831in}}%
\pgfpathlineto{\pgfqpoint{2.219016in}{2.488243in}}%
\pgfpathlineto{\pgfqpoint{2.171555in}{2.533919in}}%
\pgfpathlineto{\pgfqpoint{2.139527in}{2.554478in}}%
\pgfpathclose%
\pgfusepath{stroke,fill}%
\end{pgfscope}%
\begin{pgfscope}%
\pgfpathrectangle{\pgfqpoint{0.050000in}{0.050000in}}{\pgfqpoint{4.900000in}{4.900000in}}%
\pgfusepath{clip}%
\pgfsetbuttcap%
\pgfsetroundjoin%
\definecolor{currentfill}{rgb}{0.940823,0.940823,0.940823}%
\pgfsetfillcolor{currentfill}%
\pgfsetfillopacity{0.900000}%
\pgfsetlinewidth{0.501875pt}%
\definecolor{currentstroke}{rgb}{0.200000,0.200000,0.200000}%
\pgfsetstrokecolor{currentstroke}%
\pgfsetdash{}{0pt}%
\pgfpathmoveto{\pgfqpoint{3.571255in}{1.467030in}}%
\pgfpathlineto{\pgfqpoint{3.615993in}{1.444033in}}%
\pgfpathlineto{\pgfqpoint{3.650329in}{1.423497in}}%
\pgfpathlineto{\pgfqpoint{3.605578in}{1.446346in}}%
\pgfpathlineto{\pgfqpoint{3.571255in}{1.467030in}}%
\pgfpathclose%
\pgfusepath{stroke,fill}%
\end{pgfscope}%
\begin{pgfscope}%
\pgfpathrectangle{\pgfqpoint{0.050000in}{0.050000in}}{\pgfqpoint{4.900000in}{4.900000in}}%
\pgfusepath{clip}%
\pgfsetbuttcap%
\pgfsetroundjoin%
\definecolor{currentfill}{rgb}{0.136563,0.136563,0.136563}%
\pgfsetfillcolor{currentfill}%
\pgfsetfillopacity{0.900000}%
\pgfsetlinewidth{0.501875pt}%
\definecolor{currentstroke}{rgb}{0.200000,0.200000,0.200000}%
\pgfsetstrokecolor{currentstroke}%
\pgfsetdash{}{0pt}%
\pgfpathmoveto{\pgfqpoint{1.136618in}{3.365708in}}%
\pgfpathlineto{\pgfqpoint{1.186245in}{3.318301in}}%
\pgfpathlineto{\pgfqpoint{1.216701in}{3.298989in}}%
\pgfpathlineto{\pgfqpoint{1.167061in}{3.346429in}}%
\pgfpathlineto{\pgfqpoint{1.136618in}{3.365708in}}%
\pgfpathclose%
\pgfusepath{stroke,fill}%
\end{pgfscope}%
\begin{pgfscope}%
\pgfpathrectangle{\pgfqpoint{0.050000in}{0.050000in}}{\pgfqpoint{4.900000in}{4.900000in}}%
\pgfusepath{clip}%
\pgfsetbuttcap%
\pgfsetroundjoin%
\definecolor{currentfill}{rgb}{0.788112,0.788112,0.788112}%
\pgfsetfillcolor{currentfill}%
\pgfsetfillopacity{0.900000}%
\pgfsetlinewidth{0.501875pt}%
\definecolor{currentstroke}{rgb}{0.200000,0.200000,0.200000}%
\pgfsetstrokecolor{currentstroke}%
\pgfsetdash{}{0pt}%
\pgfpathmoveto{\pgfqpoint{2.871808in}{1.964804in}}%
\pgfpathlineto{\pgfqpoint{2.917695in}{1.921936in}}%
\pgfpathlineto{\pgfqpoint{2.950891in}{1.900724in}}%
\pgfpathlineto{\pgfqpoint{2.905002in}{1.943513in}}%
\pgfpathlineto{\pgfqpoint{2.871808in}{1.964804in}}%
\pgfpathclose%
\pgfusepath{stroke,fill}%
\end{pgfscope}%
\begin{pgfscope}%
\pgfpathrectangle{\pgfqpoint{0.050000in}{0.050000in}}{\pgfqpoint{4.900000in}{4.900000in}}%
\pgfusepath{clip}%
\pgfsetbuttcap%
\pgfsetroundjoin%
\definecolor{currentfill}{rgb}{0.858762,0.858762,0.858762}%
\pgfsetfillcolor{currentfill}%
\pgfsetfillopacity{0.900000}%
\pgfsetlinewidth{0.501875pt}%
\definecolor{currentstroke}{rgb}{0.200000,0.200000,0.200000}%
\pgfsetstrokecolor{currentstroke}%
\pgfsetdash{}{0pt}%
\pgfpathmoveto{\pgfqpoint{3.142351in}{1.755364in}}%
\pgfpathlineto{\pgfqpoint{3.187707in}{1.717173in}}%
\pgfpathlineto{\pgfqpoint{3.221340in}{1.696157in}}%
\pgfpathlineto{\pgfqpoint{3.175979in}{1.734181in}}%
\pgfpathlineto{\pgfqpoint{3.142351in}{1.755364in}}%
\pgfpathclose%
\pgfusepath{stroke,fill}%
\end{pgfscope}%
\begin{pgfscope}%
\pgfpathrectangle{\pgfqpoint{0.050000in}{0.050000in}}{\pgfqpoint{4.900000in}{4.900000in}}%
\pgfusepath{clip}%
\pgfsetbuttcap%
\pgfsetroundjoin%
\definecolor{currentfill}{rgb}{0.432203,0.432203,0.432203}%
\pgfsetfillcolor{currentfill}%
\pgfsetfillopacity{0.900000}%
\pgfsetlinewidth{0.501875pt}%
\definecolor{currentstroke}{rgb}{0.200000,0.200000,0.200000}%
\pgfsetstrokecolor{currentstroke}%
\pgfsetdash{}{0pt}%
\pgfpathmoveto{\pgfqpoint{1.868895in}{2.773868in}}%
\pgfpathlineto{\pgfqpoint{1.916943in}{2.727737in}}%
\pgfpathlineto{\pgfqpoint{1.948556in}{2.707492in}}%
\pgfpathlineto{\pgfqpoint{1.900500in}{2.753651in}}%
\pgfpathlineto{\pgfqpoint{1.868895in}{2.773868in}}%
\pgfpathclose%
\pgfusepath{stroke,fill}%
\end{pgfscope}%
\begin{pgfscope}%
\pgfpathrectangle{\pgfqpoint{0.050000in}{0.050000in}}{\pgfqpoint{4.900000in}{4.900000in}}%
\pgfusepath{clip}%
\pgfsetbuttcap%
\pgfsetroundjoin%
\definecolor{currentfill}{rgb}{0.700992,0.700992,0.700992}%
\pgfsetfillcolor{currentfill}%
\pgfsetfillopacity{0.900000}%
\pgfsetlinewidth{0.501875pt}%
\definecolor{currentstroke}{rgb}{0.200000,0.200000,0.200000}%
\pgfsetstrokecolor{currentstroke}%
\pgfsetdash{}{0pt}%
\pgfpathmoveto{\pgfqpoint{2.601272in}{2.182147in}}%
\pgfpathlineto{\pgfqpoint{2.647741in}{2.137491in}}%
\pgfpathlineto{\pgfqpoint{2.680512in}{2.116370in}}%
\pgfpathlineto{\pgfqpoint{2.634041in}{2.161025in}}%
\pgfpathlineto{\pgfqpoint{2.601272in}{2.182147in}}%
\pgfpathclose%
\pgfusepath{stroke,fill}%
\end{pgfscope}%
\begin{pgfscope}%
\pgfpathrectangle{\pgfqpoint{0.050000in}{0.050000in}}{\pgfqpoint{4.900000in}{4.900000in}}%
\pgfusepath{clip}%
\pgfsetbuttcap%
\pgfsetroundjoin%
\definecolor{currentfill}{rgb}{0.036417,0.036417,0.036417}%
\pgfsetfillcolor{currentfill}%
\pgfsetfillopacity{0.900000}%
\pgfsetlinewidth{0.501875pt}%
\definecolor{currentstroke}{rgb}{0.200000,0.200000,0.200000}%
\pgfsetstrokecolor{currentstroke}%
\pgfsetdash{}{0pt}%
\pgfpathmoveto{\pgfqpoint{0.865785in}{3.585269in}}%
\pgfpathlineto{\pgfqpoint{0.916007in}{3.537380in}}%
\pgfpathlineto{\pgfqpoint{0.946040in}{3.518410in}}%
\pgfpathlineto{\pgfqpoint{0.895804in}{3.566334in}}%
\pgfpathlineto{\pgfqpoint{0.865785in}{3.585269in}}%
\pgfpathclose%
\pgfusepath{stroke,fill}%
\end{pgfscope}%
\begin{pgfscope}%
\pgfpathrectangle{\pgfqpoint{0.050000in}{0.050000in}}{\pgfqpoint{4.900000in}{4.900000in}}%
\pgfusepath{clip}%
\pgfsetbuttcap%
\pgfsetroundjoin%
\definecolor{currentfill}{rgb}{0.926674,0.926674,0.926674}%
\pgfsetfillcolor{currentfill}%
\pgfsetfillopacity{0.900000}%
\pgfsetlinewidth{0.501875pt}%
\definecolor{currentstroke}{rgb}{0.200000,0.200000,0.200000}%
\pgfsetstrokecolor{currentstroke}%
\pgfsetdash{}{0pt}%
\pgfpathmoveto{\pgfqpoint{3.492212in}{1.514151in}}%
\pgfpathlineto{\pgfqpoint{3.537043in}{1.487641in}}%
\pgfpathlineto{\pgfqpoint{3.571255in}{1.467030in}}%
\pgfpathlineto{\pgfqpoint{3.526413in}{1.493371in}}%
\pgfpathlineto{\pgfqpoint{3.492212in}{1.514151in}}%
\pgfpathclose%
\pgfusepath{stroke,fill}%
\end{pgfscope}%
\begin{pgfscope}%
\pgfpathrectangle{\pgfqpoint{0.050000in}{0.050000in}}{\pgfqpoint{4.900000in}{4.900000in}}%
\pgfusepath{clip}%
\pgfsetbuttcap%
\pgfsetroundjoin%
\definecolor{currentfill}{rgb}{0.338824,0.338824,0.338824}%
\pgfsetfillcolor{currentfill}%
\pgfsetfillopacity{0.900000}%
\pgfsetlinewidth{0.501875pt}%
\definecolor{currentstroke}{rgb}{0.200000,0.200000,0.200000}%
\pgfsetstrokecolor{currentstroke}%
\pgfsetdash{}{0pt}%
\pgfpathmoveto{\pgfqpoint{1.598170in}{2.993335in}}%
\pgfpathlineto{\pgfqpoint{1.646812in}{2.946723in}}%
\pgfpathlineto{\pgfqpoint{1.678003in}{2.926819in}}%
\pgfpathlineto{\pgfqpoint{1.629352in}{2.973463in}}%
\pgfpathlineto{\pgfqpoint{1.598170in}{2.993335in}}%
\pgfpathclose%
\pgfusepath{stroke,fill}%
\end{pgfscope}%
\begin{pgfscope}%
\pgfpathrectangle{\pgfqpoint{0.050000in}{0.050000in}}{\pgfqpoint{4.900000in}{4.900000in}}%
\pgfusepath{clip}%
\pgfsetbuttcap%
\pgfsetroundjoin%
\definecolor{currentfill}{rgb}{0.590634,0.590634,0.590634}%
\pgfsetfillcolor{currentfill}%
\pgfsetfillopacity{0.900000}%
\pgfsetlinewidth{0.501875pt}%
\definecolor{currentstroke}{rgb}{0.200000,0.200000,0.200000}%
\pgfsetstrokecolor{currentstroke}%
\pgfsetdash{}{0pt}%
\pgfpathmoveto{\pgfqpoint{2.330655in}{2.401345in}}%
\pgfpathlineto{\pgfqpoint{2.377716in}{2.356018in}}%
\pgfpathlineto{\pgfqpoint{2.410065in}{2.335184in}}%
\pgfpathlineto{\pgfqpoint{2.363000in}{2.380536in}}%
\pgfpathlineto{\pgfqpoint{2.330655in}{2.401345in}}%
\pgfpathclose%
\pgfusepath{stroke,fill}%
\end{pgfscope}%
\begin{pgfscope}%
\pgfpathrectangle{\pgfqpoint{0.050000in}{0.050000in}}{\pgfqpoint{4.900000in}{4.900000in}}%
\pgfusepath{clip}%
\pgfsetbuttcap%
\pgfsetroundjoin%
\definecolor{currentfill}{rgb}{0.217762,0.217762,0.217762}%
\pgfsetfillcolor{currentfill}%
\pgfsetfillopacity{0.900000}%
\pgfsetlinewidth{0.501875pt}%
\definecolor{currentstroke}{rgb}{0.200000,0.200000,0.200000}%
\pgfsetstrokecolor{currentstroke}%
\pgfsetdash{}{0pt}%
\pgfpathmoveto{\pgfqpoint{1.327352in}{3.212881in}}%
\pgfpathlineto{\pgfqpoint{1.376587in}{3.165785in}}%
\pgfpathlineto{\pgfqpoint{1.407357in}{3.146224in}}%
\pgfpathlineto{\pgfqpoint{1.358110in}{3.193352in}}%
\pgfpathlineto{\pgfqpoint{1.327352in}{3.212881in}}%
\pgfpathclose%
\pgfusepath{stroke,fill}%
\end{pgfscope}%
\begin{pgfscope}%
\pgfpathrectangle{\pgfqpoint{0.050000in}{0.050000in}}{\pgfqpoint{4.900000in}{4.900000in}}%
\pgfusepath{clip}%
\pgfsetbuttcap%
\pgfsetroundjoin%
\definecolor{currentfill}{rgb}{0.912526,0.912526,0.912526}%
\pgfsetfillcolor{currentfill}%
\pgfsetfillopacity{0.900000}%
\pgfsetlinewidth{0.501875pt}%
\definecolor{currentstroke}{rgb}{0.200000,0.200000,0.200000}%
\pgfsetstrokecolor{currentstroke}%
\pgfsetdash{}{0pt}%
\pgfpathmoveto{\pgfqpoint{3.413180in}{1.564710in}}%
\pgfpathlineto{\pgfqpoint{3.458120in}{1.534862in}}%
\pgfpathlineto{\pgfqpoint{3.492212in}{1.514151in}}%
\pgfpathlineto{\pgfqpoint{3.447261in}{1.543817in}}%
\pgfpathlineto{\pgfqpoint{3.413180in}{1.564710in}}%
\pgfpathclose%
\pgfusepath{stroke,fill}%
\end{pgfscope}%
\begin{pgfscope}%
\pgfpathrectangle{\pgfqpoint{0.050000in}{0.050000in}}{\pgfqpoint{4.900000in}{4.900000in}}%
\pgfusepath{clip}%
\pgfsetbuttcap%
\pgfsetroundjoin%
\definecolor{currentfill}{rgb}{0.495656,0.495656,0.495656}%
\pgfsetfillcolor{currentfill}%
\pgfsetfillopacity{0.900000}%
\pgfsetlinewidth{0.501875pt}%
\definecolor{currentstroke}{rgb}{0.200000,0.200000,0.200000}%
\pgfsetstrokecolor{currentstroke}%
\pgfsetdash{}{0pt}%
\pgfpathmoveto{\pgfqpoint{2.059944in}{2.620791in}}%
\pgfpathlineto{\pgfqpoint{2.107599in}{2.574973in}}%
\pgfpathlineto{\pgfqpoint{2.139527in}{2.554478in}}%
\pgfpathlineto{\pgfqpoint{2.091866in}{2.600325in}}%
\pgfpathlineto{\pgfqpoint{2.059944in}{2.620791in}}%
\pgfpathclose%
\pgfusepath{stroke,fill}%
\end{pgfscope}%
\begin{pgfscope}%
\pgfpathrectangle{\pgfqpoint{0.050000in}{0.050000in}}{\pgfqpoint{4.900000in}{4.900000in}}%
\pgfusepath{clip}%
\pgfsetbuttcap%
\pgfsetroundjoin%
\definecolor{currentfill}{rgb}{0.839785,0.839785,0.839785}%
\pgfsetfillcolor{currentfill}%
\pgfsetfillopacity{0.900000}%
\pgfsetlinewidth{0.501875pt}%
\definecolor{currentstroke}{rgb}{0.200000,0.200000,0.200000}%
\pgfsetstrokecolor{currentstroke}%
\pgfsetdash{}{0pt}%
\pgfpathmoveto{\pgfqpoint{3.063305in}{1.816607in}}%
\pgfpathlineto{\pgfqpoint{3.108830in}{1.776490in}}%
\pgfpathlineto{\pgfqpoint{3.142351in}{1.755364in}}%
\pgfpathlineto{\pgfqpoint{3.096823in}{1.795338in}}%
\pgfpathlineto{\pgfqpoint{3.063305in}{1.816607in}}%
\pgfpathclose%
\pgfusepath{stroke,fill}%
\end{pgfscope}%
\begin{pgfscope}%
\pgfpathrectangle{\pgfqpoint{0.050000in}{0.050000in}}{\pgfqpoint{4.900000in}{4.900000in}}%
\pgfusepath{clip}%
\pgfsetbuttcap%
\pgfsetroundjoin%
\definecolor{currentfill}{rgb}{0.763998,0.763998,0.763998}%
\pgfsetfillcolor{currentfill}%
\pgfsetfillopacity{0.900000}%
\pgfsetlinewidth{0.501875pt}%
\definecolor{currentstroke}{rgb}{0.200000,0.200000,0.200000}%
\pgfsetstrokecolor{currentstroke}%
\pgfsetdash{}{0pt}%
\pgfpathmoveto{\pgfqpoint{2.792642in}{2.029715in}}%
\pgfpathlineto{\pgfqpoint{2.838720in}{1.986037in}}%
\pgfpathlineto{\pgfqpoint{2.871808in}{1.964804in}}%
\pgfpathlineto{\pgfqpoint{2.825728in}{2.008434in}}%
\pgfpathlineto{\pgfqpoint{2.792642in}{2.029715in}}%
\pgfpathclose%
\pgfusepath{stroke,fill}%
\end{pgfscope}%
\begin{pgfscope}%
\pgfpathrectangle{\pgfqpoint{0.050000in}{0.050000in}}{\pgfqpoint{4.900000in}{4.900000in}}%
\pgfusepath{clip}%
\pgfsetbuttcap%
\pgfsetroundjoin%
\definecolor{currentfill}{rgb}{0.104698,0.104698,0.104698}%
\pgfsetfillcolor{currentfill}%
\pgfsetfillopacity{0.900000}%
\pgfsetlinewidth{0.501875pt}%
\definecolor{currentstroke}{rgb}{0.200000,0.200000,0.200000}%
\pgfsetstrokecolor{currentstroke}%
\pgfsetdash{}{0pt}%
\pgfpathmoveto{\pgfqpoint{1.056440in}{3.432505in}}%
\pgfpathlineto{\pgfqpoint{1.106270in}{3.384926in}}%
\pgfpathlineto{\pgfqpoint{1.136618in}{3.365708in}}%
\pgfpathlineto{\pgfqpoint{1.086774in}{3.413320in}}%
\pgfpathlineto{\pgfqpoint{1.056440in}{3.432505in}}%
\pgfpathclose%
\pgfusepath{stroke,fill}%
\end{pgfscope}%
\begin{pgfscope}%
\pgfpathrectangle{\pgfqpoint{0.050000in}{0.050000in}}{\pgfqpoint{4.900000in}{4.900000in}}%
\pgfusepath{clip}%
\pgfsetbuttcap%
\pgfsetroundjoin%
\definecolor{currentfill}{rgb}{0.399723,0.399723,0.399723}%
\pgfsetfillcolor{currentfill}%
\pgfsetfillopacity{0.900000}%
\pgfsetlinewidth{0.501875pt}%
\definecolor{currentstroke}{rgb}{0.200000,0.200000,0.200000}%
\pgfsetstrokecolor{currentstroke}%
\pgfsetdash{}{0pt}%
\pgfpathmoveto{\pgfqpoint{1.789141in}{2.840321in}}%
\pgfpathlineto{\pgfqpoint{1.837389in}{2.794020in}}%
\pgfpathlineto{\pgfqpoint{1.868895in}{2.773868in}}%
\pgfpathlineto{\pgfqpoint{1.820638in}{2.820198in}}%
\pgfpathlineto{\pgfqpoint{1.789141in}{2.840321in}}%
\pgfpathclose%
\pgfusepath{stroke,fill}%
\end{pgfscope}%
\begin{pgfscope}%
\pgfpathrectangle{\pgfqpoint{0.050000in}{0.050000in}}{\pgfqpoint{4.900000in}{4.900000in}}%
\pgfusepath{clip}%
\pgfsetbuttcap%
\pgfsetroundjoin%
\definecolor{currentfill}{rgb}{0.667405,0.667405,0.667405}%
\pgfsetfillcolor{currentfill}%
\pgfsetfillopacity{0.900000}%
\pgfsetlinewidth{0.501875pt}%
\definecolor{currentstroke}{rgb}{0.200000,0.200000,0.200000}%
\pgfsetstrokecolor{currentstroke}%
\pgfsetdash{}{0pt}%
\pgfpathmoveto{\pgfqpoint{2.521941in}{2.248144in}}%
\pgfpathlineto{\pgfqpoint{2.568608in}{2.203207in}}%
\pgfpathlineto{\pgfqpoint{2.601272in}{2.182147in}}%
\pgfpathlineto{\pgfqpoint{2.554603in}{2.227097in}}%
\pgfpathlineto{\pgfqpoint{2.521941in}{2.248144in}}%
\pgfpathclose%
\pgfusepath{stroke,fill}%
\end{pgfscope}%
\begin{pgfscope}%
\pgfpathrectangle{\pgfqpoint{0.050000in}{0.050000in}}{\pgfqpoint{4.900000in}{4.900000in}}%
\pgfusepath{clip}%
\pgfsetbuttcap%
\pgfsetroundjoin%
\definecolor{currentfill}{rgb}{0.000000,0.000000,0.000000}%
\pgfsetfillcolor{currentfill}%
\pgfsetfillopacity{0.900000}%
\pgfsetlinewidth{0.501875pt}%
\definecolor{currentstroke}{rgb}{0.200000,0.200000,0.200000}%
\pgfsetstrokecolor{currentstroke}%
\pgfsetdash{}{0pt}%
\pgfpathmoveto{\pgfqpoint{0.785436in}{3.652207in}}%
\pgfpathlineto{\pgfqpoint{0.835860in}{3.604144in}}%
\pgfpathlineto{\pgfqpoint{0.865785in}{3.585269in}}%
\pgfpathlineto{\pgfqpoint{0.815345in}{3.633366in}}%
\pgfpathlineto{\pgfqpoint{0.785436in}{3.652207in}}%
\pgfpathclose%
\pgfusepath{stroke,fill}%
\end{pgfscope}%
\begin{pgfscope}%
\pgfpathrectangle{\pgfqpoint{0.050000in}{0.050000in}}{\pgfqpoint{4.900000in}{4.900000in}}%
\pgfusepath{clip}%
\pgfsetbuttcap%
\pgfsetroundjoin%
\definecolor{currentfill}{rgb}{0.895548,0.895548,0.895548}%
\pgfsetfillcolor{currentfill}%
\pgfsetfillopacity{0.900000}%
\pgfsetlinewidth{0.501875pt}%
\definecolor{currentstroke}{rgb}{0.200000,0.200000,0.200000}%
\pgfsetstrokecolor{currentstroke}%
\pgfsetdash{}{0pt}%
\pgfpathmoveto{\pgfqpoint{3.334142in}{1.618463in}}%
\pgfpathlineto{\pgfqpoint{3.379208in}{1.585538in}}%
\pgfpathlineto{\pgfqpoint{3.413180in}{1.564710in}}%
\pgfpathlineto{\pgfqpoint{3.368106in}{1.597448in}}%
\pgfpathlineto{\pgfqpoint{3.334142in}{1.618463in}}%
\pgfpathclose%
\pgfusepath{stroke,fill}%
\end{pgfscope}%
\begin{pgfscope}%
\pgfpathrectangle{\pgfqpoint{0.050000in}{0.050000in}}{\pgfqpoint{4.900000in}{4.900000in}}%
\pgfusepath{clip}%
\pgfsetbuttcap%
\pgfsetroundjoin%
\definecolor{currentfill}{rgb}{0.300807,0.300807,0.300807}%
\pgfsetfillcolor{currentfill}%
\pgfsetfillopacity{0.900000}%
\pgfsetlinewidth{0.501875pt}%
\definecolor{currentstroke}{rgb}{0.200000,0.200000,0.200000}%
\pgfsetstrokecolor{currentstroke}%
\pgfsetdash{}{0pt}%
\pgfpathmoveto{\pgfqpoint{1.518244in}{3.059929in}}%
\pgfpathlineto{\pgfqpoint{1.567087in}{3.013145in}}%
\pgfpathlineto{\pgfqpoint{1.598170in}{2.993335in}}%
\pgfpathlineto{\pgfqpoint{1.549317in}{3.040151in}}%
\pgfpathlineto{\pgfqpoint{1.518244in}{3.059929in}}%
\pgfpathclose%
\pgfusepath{stroke,fill}%
\end{pgfscope}%
\begin{pgfscope}%
\pgfpathrectangle{\pgfqpoint{0.050000in}{0.050000in}}{\pgfqpoint{4.900000in}{4.900000in}}%
\pgfusepath{clip}%
\pgfsetbuttcap%
\pgfsetroundjoin%
\definecolor{currentfill}{rgb}{0.560246,0.560246,0.560246}%
\pgfsetfillcolor{currentfill}%
\pgfsetfillopacity{0.900000}%
\pgfsetlinewidth{0.501875pt}%
\definecolor{currentstroke}{rgb}{0.200000,0.200000,0.200000}%
\pgfsetstrokecolor{currentstroke}%
\pgfsetdash{}{0pt}%
\pgfpathmoveto{\pgfqpoint{2.251152in}{2.467591in}}%
\pgfpathlineto{\pgfqpoint{2.298412in}{2.422088in}}%
\pgfpathlineto{\pgfqpoint{2.330655in}{2.401345in}}%
\pgfpathlineto{\pgfqpoint{2.283389in}{2.446874in}}%
\pgfpathlineto{\pgfqpoint{2.251152in}{2.467591in}}%
\pgfpathclose%
\pgfusepath{stroke,fill}%
\end{pgfscope}%
\begin{pgfscope}%
\pgfpathrectangle{\pgfqpoint{0.050000in}{0.050000in}}{\pgfqpoint{4.900000in}{4.900000in}}%
\pgfusepath{clip}%
\pgfsetbuttcap%
\pgfsetroundjoin%
\definecolor{currentfill}{rgb}{0.179008,0.179008,0.179008}%
\pgfsetfillcolor{currentfill}%
\pgfsetfillopacity{0.900000}%
\pgfsetlinewidth{0.501875pt}%
\definecolor{currentstroke}{rgb}{0.200000,0.200000,0.200000}%
\pgfsetstrokecolor{currentstroke}%
\pgfsetdash{}{0pt}%
\pgfpathmoveto{\pgfqpoint{1.247253in}{3.279616in}}%
\pgfpathlineto{\pgfqpoint{1.296691in}{3.232348in}}%
\pgfpathlineto{\pgfqpoint{1.327352in}{3.212881in}}%
\pgfpathlineto{\pgfqpoint{1.277902in}{3.260181in}}%
\pgfpathlineto{\pgfqpoint{1.247253in}{3.279616in}}%
\pgfpathclose%
\pgfusepath{stroke,fill}%
\end{pgfscope}%
\begin{pgfscope}%
\pgfpathrectangle{\pgfqpoint{0.050000in}{0.050000in}}{\pgfqpoint{4.900000in}{4.900000in}}%
\pgfusepath{clip}%
\pgfsetbuttcap%
\pgfsetroundjoin%
\definecolor{currentfill}{rgb}{0.815671,0.815671,0.815671}%
\pgfsetfillcolor{currentfill}%
\pgfsetfillopacity{0.900000}%
\pgfsetlinewidth{0.501875pt}%
\definecolor{currentstroke}{rgb}{0.200000,0.200000,0.200000}%
\pgfsetstrokecolor{currentstroke}%
\pgfsetdash{}{0pt}%
\pgfpathmoveto{\pgfqpoint{2.984192in}{1.879456in}}%
\pgfpathlineto{\pgfqpoint{3.029895in}{1.837819in}}%
\pgfpathlineto{\pgfqpoint{3.063305in}{1.816607in}}%
\pgfpathlineto{\pgfqpoint{3.017600in}{1.858132in}}%
\pgfpathlineto{\pgfqpoint{2.984192in}{1.879456in}}%
\pgfpathclose%
\pgfusepath{stroke,fill}%
\end{pgfscope}%
\begin{pgfscope}%
\pgfpathrectangle{\pgfqpoint{0.050000in}{0.050000in}}{\pgfqpoint{4.900000in}{4.900000in}}%
\pgfusepath{clip}%
\pgfsetbuttcap%
\pgfsetroundjoin%
\definecolor{currentfill}{rgb}{0.461207,0.461207,0.461207}%
\pgfsetfillcolor{currentfill}%
\pgfsetfillopacity{0.900000}%
\pgfsetlinewidth{0.501875pt}%
\definecolor{currentstroke}{rgb}{0.200000,0.200000,0.200000}%
\pgfsetstrokecolor{currentstroke}%
\pgfsetdash{}{0pt}%
\pgfpathmoveto{\pgfqpoint{1.980269in}{2.687182in}}%
\pgfpathlineto{\pgfqpoint{2.028124in}{2.641194in}}%
\pgfpathlineto{\pgfqpoint{2.059944in}{2.620791in}}%
\pgfpathlineto{\pgfqpoint{2.012083in}{2.666809in}}%
\pgfpathlineto{\pgfqpoint{1.980269in}{2.687182in}}%
\pgfpathclose%
\pgfusepath{stroke,fill}%
\end{pgfscope}%
\begin{pgfscope}%
\pgfpathrectangle{\pgfqpoint{0.050000in}{0.050000in}}{\pgfqpoint{4.900000in}{4.900000in}}%
\pgfusepath{clip}%
\pgfsetbuttcap%
\pgfsetroundjoin%
\definecolor{currentfill}{rgb}{0.739377,0.739377,0.739377}%
\pgfsetfillcolor{currentfill}%
\pgfsetfillopacity{0.900000}%
\pgfsetlinewidth{0.501875pt}%
\definecolor{currentstroke}{rgb}{0.200000,0.200000,0.200000}%
\pgfsetstrokecolor{currentstroke}%
\pgfsetdash{}{0pt}%
\pgfpathmoveto{\pgfqpoint{2.713387in}{2.095188in}}%
\pgfpathlineto{\pgfqpoint{2.759660in}{2.050937in}}%
\pgfpathlineto{\pgfqpoint{2.792642in}{2.029715in}}%
\pgfpathlineto{\pgfqpoint{2.746366in}{2.073943in}}%
\pgfpathlineto{\pgfqpoint{2.713387in}{2.095188in}}%
\pgfpathclose%
\pgfusepath{stroke,fill}%
\end{pgfscope}%
\begin{pgfscope}%
\pgfpathrectangle{\pgfqpoint{0.050000in}{0.050000in}}{\pgfqpoint{4.900000in}{4.900000in}}%
\pgfusepath{clip}%
\pgfsetbuttcap%
\pgfsetroundjoin%
\definecolor{currentfill}{rgb}{0.068281,0.068281,0.068281}%
\pgfsetfillcolor{currentfill}%
\pgfsetfillopacity{0.900000}%
\pgfsetlinewidth{0.501875pt}%
\definecolor{currentstroke}{rgb}{0.200000,0.200000,0.200000}%
\pgfsetstrokecolor{currentstroke}%
\pgfsetdash{}{0pt}%
\pgfpathmoveto{\pgfqpoint{0.976169in}{3.499380in}}%
\pgfpathlineto{\pgfqpoint{1.026202in}{3.451629in}}%
\pgfpathlineto{\pgfqpoint{1.056440in}{3.432505in}}%
\pgfpathlineto{\pgfqpoint{1.006394in}{3.480290in}}%
\pgfpathlineto{\pgfqpoint{0.976169in}{3.499380in}}%
\pgfpathclose%
\pgfusepath{stroke,fill}%
\end{pgfscope}%
\begin{pgfscope}%
\pgfpathrectangle{\pgfqpoint{0.050000in}{0.050000in}}{\pgfqpoint{4.900000in}{4.900000in}}%
\pgfusepath{clip}%
\pgfsetbuttcap%
\pgfsetroundjoin%
\definecolor{currentfill}{rgb}{0.878570,0.878570,0.878570}%
\pgfsetfillcolor{currentfill}%
\pgfsetfillopacity{0.900000}%
\pgfsetlinewidth{0.501875pt}%
\definecolor{currentstroke}{rgb}{0.200000,0.200000,0.200000}%
\pgfsetstrokecolor{currentstroke}%
\pgfsetdash{}{0pt}%
\pgfpathmoveto{\pgfqpoint{3.255080in}{1.675082in}}%
\pgfpathlineto{\pgfqpoint{3.300287in}{1.639415in}}%
\pgfpathlineto{\pgfqpoint{3.334142in}{1.618463in}}%
\pgfpathlineto{\pgfqpoint{3.288929in}{1.653948in}}%
\pgfpathlineto{\pgfqpoint{3.255080in}{1.675082in}}%
\pgfpathclose%
\pgfusepath{stroke,fill}%
\end{pgfscope}%
\begin{pgfscope}%
\pgfpathrectangle{\pgfqpoint{0.050000in}{0.050000in}}{\pgfqpoint{4.900000in}{4.900000in}}%
\pgfusepath{clip}%
\pgfsetbuttcap%
\pgfsetroundjoin%
\definecolor{currentfill}{rgb}{0.371303,0.371303,0.371303}%
\pgfsetfillcolor{currentfill}%
\pgfsetfillopacity{0.900000}%
\pgfsetlinewidth{0.501875pt}%
\definecolor{currentstroke}{rgb}{0.200000,0.200000,0.200000}%
\pgfsetstrokecolor{currentstroke}%
\pgfsetdash{}{0pt}%
\pgfpathmoveto{\pgfqpoint{1.709293in}{2.906853in}}%
\pgfpathlineto{\pgfqpoint{1.757743in}{2.860381in}}%
\pgfpathlineto{\pgfqpoint{1.789141in}{2.840321in}}%
\pgfpathlineto{\pgfqpoint{1.740683in}{2.886824in}}%
\pgfpathlineto{\pgfqpoint{1.709293in}{2.906853in}}%
\pgfpathclose%
\pgfusepath{stroke,fill}%
\end{pgfscope}%
\begin{pgfscope}%
\pgfpathrectangle{\pgfqpoint{0.050000in}{0.050000in}}{\pgfqpoint{4.900000in}{4.900000in}}%
\pgfusepath{clip}%
\pgfsetbuttcap%
\pgfsetroundjoin%
\definecolor{currentfill}{rgb}{0.629020,0.629020,0.629020}%
\pgfsetfillcolor{currentfill}%
\pgfsetfillopacity{0.900000}%
\pgfsetlinewidth{0.501875pt}%
\definecolor{currentstroke}{rgb}{0.200000,0.200000,0.200000}%
\pgfsetstrokecolor{currentstroke}%
\pgfsetdash{}{0pt}%
\pgfpathmoveto{\pgfqpoint{2.442517in}{2.314286in}}%
\pgfpathlineto{\pgfqpoint{2.489383in}{2.269127in}}%
\pgfpathlineto{\pgfqpoint{2.521941in}{2.248144in}}%
\pgfpathlineto{\pgfqpoint{2.475072in}{2.293323in}}%
\pgfpathlineto{\pgfqpoint{2.442517in}{2.314286in}}%
\pgfpathclose%
\pgfusepath{stroke,fill}%
\end{pgfscope}%
\begin{pgfscope}%
\pgfpathrectangle{\pgfqpoint{0.050000in}{0.050000in}}{\pgfqpoint{4.900000in}{4.900000in}}%
\pgfusepath{clip}%
\pgfsetbuttcap%
\pgfsetroundjoin%
\definecolor{currentfill}{rgb}{0.262053,0.262053,0.262053}%
\pgfsetfillcolor{currentfill}%
\pgfsetfillopacity{0.900000}%
\pgfsetlinewidth{0.501875pt}%
\definecolor{currentstroke}{rgb}{0.200000,0.200000,0.200000}%
\pgfsetstrokecolor{currentstroke}%
\pgfsetdash{}{0pt}%
\pgfpathmoveto{\pgfqpoint{1.438224in}{3.126602in}}%
\pgfpathlineto{\pgfqpoint{1.487268in}{3.079646in}}%
\pgfpathlineto{\pgfqpoint{1.518244in}{3.059929in}}%
\pgfpathlineto{\pgfqpoint{1.469189in}{3.106917in}}%
\pgfpathlineto{\pgfqpoint{1.438224in}{3.126602in}}%
\pgfpathclose%
\pgfusepath{stroke,fill}%
\end{pgfscope}%
\begin{pgfscope}%
\pgfpathrectangle{\pgfqpoint{0.050000in}{0.050000in}}{\pgfqpoint{4.900000in}{4.900000in}}%
\pgfusepath{clip}%
\pgfsetbuttcap%
\pgfsetroundjoin%
\definecolor{currentfill}{rgb}{0.950173,0.950173,0.950173}%
\pgfsetfillcolor{currentfill}%
\pgfsetfillopacity{0.900000}%
\pgfsetlinewidth{0.501875pt}%
\definecolor{currentstroke}{rgb}{0.200000,0.200000,0.200000}%
\pgfsetstrokecolor{currentstroke}%
\pgfsetdash{}{0pt}%
\pgfpathmoveto{\pgfqpoint{3.684775in}{1.402884in}}%
\pgfpathlineto{\pgfqpoint{3.729448in}{1.383605in}}%
\pgfpathlineto{\pgfqpoint{3.764021in}{1.363035in}}%
\pgfpathlineto{\pgfqpoint{3.719333in}{1.382195in}}%
\pgfpathlineto{\pgfqpoint{3.684775in}{1.402884in}}%
\pgfpathclose%
\pgfusepath{stroke,fill}%
\end{pgfscope}%
\begin{pgfscope}%
\pgfpathrectangle{\pgfqpoint{0.050000in}{0.050000in}}{\pgfqpoint{4.900000in}{4.900000in}}%
\pgfusepath{clip}%
\pgfsetbuttcap%
\pgfsetroundjoin%
\definecolor{currentfill}{rgb}{0.525798,0.525798,0.525798}%
\pgfsetfillcolor{currentfill}%
\pgfsetfillopacity{0.900000}%
\pgfsetlinewidth{0.501875pt}%
\definecolor{currentstroke}{rgb}{0.200000,0.200000,0.200000}%
\pgfsetstrokecolor{currentstroke}%
\pgfsetdash{}{0pt}%
\pgfpathmoveto{\pgfqpoint{2.171555in}{2.533919in}}%
\pgfpathlineto{\pgfqpoint{2.219016in}{2.488243in}}%
\pgfpathlineto{\pgfqpoint{2.251152in}{2.467591in}}%
\pgfpathlineto{\pgfqpoint{2.203686in}{2.513294in}}%
\pgfpathlineto{\pgfqpoint{2.171555in}{2.533919in}}%
\pgfpathclose%
\pgfusepath{stroke,fill}%
\end{pgfscope}%
\begin{pgfscope}%
\pgfpathrectangle{\pgfqpoint{0.050000in}{0.050000in}}{\pgfqpoint{4.900000in}{4.900000in}}%
\pgfusepath{clip}%
\pgfsetbuttcap%
\pgfsetroundjoin%
\definecolor{currentfill}{rgb}{0.136563,0.136563,0.136563}%
\pgfsetfillcolor{currentfill}%
\pgfsetfillopacity{0.900000}%
\pgfsetlinewidth{0.501875pt}%
\definecolor{currentstroke}{rgb}{0.200000,0.200000,0.200000}%
\pgfsetstrokecolor{currentstroke}%
\pgfsetdash{}{0pt}%
\pgfpathmoveto{\pgfqpoint{1.167061in}{3.346429in}}%
\pgfpathlineto{\pgfqpoint{1.216701in}{3.298989in}}%
\pgfpathlineto{\pgfqpoint{1.247253in}{3.279616in}}%
\pgfpathlineto{\pgfqpoint{1.197601in}{3.327089in}}%
\pgfpathlineto{\pgfqpoint{1.167061in}{3.346429in}}%
\pgfpathclose%
\pgfusepath{stroke,fill}%
\end{pgfscope}%
\begin{pgfscope}%
\pgfpathrectangle{\pgfqpoint{0.050000in}{0.050000in}}{\pgfqpoint{4.900000in}{4.900000in}}%
\pgfusepath{clip}%
\pgfsetbuttcap%
\pgfsetroundjoin%
\definecolor{currentfill}{rgb}{0.940823,0.940823,0.940823}%
\pgfsetfillcolor{currentfill}%
\pgfsetfillopacity{0.900000}%
\pgfsetlinewidth{0.501875pt}%
\definecolor{currentstroke}{rgb}{0.200000,0.200000,0.200000}%
\pgfsetstrokecolor{currentstroke}%
\pgfsetdash{}{0pt}%
\pgfpathmoveto{\pgfqpoint{3.605578in}{1.446346in}}%
\pgfpathlineto{\pgfqpoint{3.650329in}{1.423497in}}%
\pgfpathlineto{\pgfqpoint{3.684775in}{1.402884in}}%
\pgfpathlineto{\pgfqpoint{3.640011in}{1.425588in}}%
\pgfpathlineto{\pgfqpoint{3.605578in}{1.446346in}}%
\pgfpathclose%
\pgfusepath{stroke,fill}%
\end{pgfscope}%
\begin{pgfscope}%
\pgfpathrectangle{\pgfqpoint{0.050000in}{0.050000in}}{\pgfqpoint{4.900000in}{4.900000in}}%
\pgfusepath{clip}%
\pgfsetbuttcap%
\pgfsetroundjoin%
\definecolor{currentfill}{rgb}{0.788112,0.788112,0.788112}%
\pgfsetfillcolor{currentfill}%
\pgfsetfillopacity{0.900000}%
\pgfsetlinewidth{0.501875pt}%
\definecolor{currentstroke}{rgb}{0.200000,0.200000,0.200000}%
\pgfsetstrokecolor{currentstroke}%
\pgfsetdash{}{0pt}%
\pgfpathmoveto{\pgfqpoint{2.905002in}{1.943513in}}%
\pgfpathlineto{\pgfqpoint{2.950891in}{1.900724in}}%
\pgfpathlineto{\pgfqpoint{2.984192in}{1.879456in}}%
\pgfpathlineto{\pgfqpoint{2.938301in}{1.922165in}}%
\pgfpathlineto{\pgfqpoint{2.905002in}{1.943513in}}%
\pgfpathclose%
\pgfusepath{stroke,fill}%
\end{pgfscope}%
\begin{pgfscope}%
\pgfpathrectangle{\pgfqpoint{0.050000in}{0.050000in}}{\pgfqpoint{4.900000in}{4.900000in}}%
\pgfusepath{clip}%
\pgfsetbuttcap%
\pgfsetroundjoin%
\definecolor{currentfill}{rgb}{0.432203,0.432203,0.432203}%
\pgfsetfillcolor{currentfill}%
\pgfsetfillopacity{0.900000}%
\pgfsetlinewidth{0.501875pt}%
\definecolor{currentstroke}{rgb}{0.200000,0.200000,0.200000}%
\pgfsetstrokecolor{currentstroke}%
\pgfsetdash{}{0pt}%
\pgfpathmoveto{\pgfqpoint{1.900500in}{2.753651in}}%
\pgfpathlineto{\pgfqpoint{1.948556in}{2.707492in}}%
\pgfpathlineto{\pgfqpoint{1.980269in}{2.687182in}}%
\pgfpathlineto{\pgfqpoint{1.932206in}{2.733371in}}%
\pgfpathlineto{\pgfqpoint{1.900500in}{2.753651in}}%
\pgfpathclose%
\pgfusepath{stroke,fill}%
\end{pgfscope}%
\begin{pgfscope}%
\pgfpathrectangle{\pgfqpoint{0.050000in}{0.050000in}}{\pgfqpoint{4.900000in}{4.900000in}}%
\pgfusepath{clip}%
\pgfsetbuttcap%
\pgfsetroundjoin%
\definecolor{currentfill}{rgb}{0.858762,0.858762,0.858762}%
\pgfsetfillcolor{currentfill}%
\pgfsetfillopacity{0.900000}%
\pgfsetlinewidth{0.501875pt}%
\definecolor{currentstroke}{rgb}{0.200000,0.200000,0.200000}%
\pgfsetstrokecolor{currentstroke}%
\pgfsetdash{}{0pt}%
\pgfpathmoveto{\pgfqpoint{3.175979in}{1.734181in}}%
\pgfpathlineto{\pgfqpoint{3.221340in}{1.696157in}}%
\pgfpathlineto{\pgfqpoint{3.255080in}{1.675082in}}%
\pgfpathlineto{\pgfqpoint{3.209714in}{1.712940in}}%
\pgfpathlineto{\pgfqpoint{3.175979in}{1.734181in}}%
\pgfpathclose%
\pgfusepath{stroke,fill}%
\end{pgfscope}%
\begin{pgfscope}%
\pgfpathrectangle{\pgfqpoint{0.050000in}{0.050000in}}{\pgfqpoint{4.900000in}{4.900000in}}%
\pgfusepath{clip}%
\pgfsetbuttcap%
\pgfsetroundjoin%
\definecolor{currentfill}{rgb}{0.700992,0.700992,0.700992}%
\pgfsetfillcolor{currentfill}%
\pgfsetfillopacity{0.900000}%
\pgfsetlinewidth{0.501875pt}%
\definecolor{currentstroke}{rgb}{0.200000,0.200000,0.200000}%
\pgfsetstrokecolor{currentstroke}%
\pgfsetdash{}{0pt}%
\pgfpathmoveto{\pgfqpoint{2.634041in}{2.161025in}}%
\pgfpathlineto{\pgfqpoint{2.680512in}{2.116370in}}%
\pgfpathlineto{\pgfqpoint{2.713387in}{2.095188in}}%
\pgfpathlineto{\pgfqpoint{2.666913in}{2.139839in}}%
\pgfpathlineto{\pgfqpoint{2.634041in}{2.161025in}}%
\pgfpathclose%
\pgfusepath{stroke,fill}%
\end{pgfscope}%
\begin{pgfscope}%
\pgfpathrectangle{\pgfqpoint{0.050000in}{0.050000in}}{\pgfqpoint{4.900000in}{4.900000in}}%
\pgfusepath{clip}%
\pgfsetbuttcap%
\pgfsetroundjoin%
\definecolor{currentfill}{rgb}{0.036417,0.036417,0.036417}%
\pgfsetfillcolor{currentfill}%
\pgfsetfillopacity{0.900000}%
\pgfsetlinewidth{0.501875pt}%
\definecolor{currentstroke}{rgb}{0.200000,0.200000,0.200000}%
\pgfsetstrokecolor{currentstroke}%
\pgfsetdash{}{0pt}%
\pgfpathmoveto{\pgfqpoint{0.895804in}{3.566334in}}%
\pgfpathlineto{\pgfqpoint{0.946040in}{3.518410in}}%
\pgfpathlineto{\pgfqpoint{0.976169in}{3.499380in}}%
\pgfpathlineto{\pgfqpoint{0.925919in}{3.547339in}}%
\pgfpathlineto{\pgfqpoint{0.895804in}{3.566334in}}%
\pgfpathclose%
\pgfusepath{stroke,fill}%
\end{pgfscope}%
\begin{pgfscope}%
\pgfpathrectangle{\pgfqpoint{0.050000in}{0.050000in}}{\pgfqpoint{4.900000in}{4.900000in}}%
\pgfusepath{clip}%
\pgfsetbuttcap%
\pgfsetroundjoin%
\definecolor{currentfill}{rgb}{0.926674,0.926674,0.926674}%
\pgfsetfillcolor{currentfill}%
\pgfsetfillopacity{0.900000}%
\pgfsetlinewidth{0.501875pt}%
\definecolor{currentstroke}{rgb}{0.200000,0.200000,0.200000}%
\pgfsetstrokecolor{currentstroke}%
\pgfsetdash{}{0pt}%
\pgfpathmoveto{\pgfqpoint{3.526413in}{1.493371in}}%
\pgfpathlineto{\pgfqpoint{3.571255in}{1.467030in}}%
\pgfpathlineto{\pgfqpoint{3.605578in}{1.446346in}}%
\pgfpathlineto{\pgfqpoint{3.560724in}{1.472520in}}%
\pgfpathlineto{\pgfqpoint{3.526413in}{1.493371in}}%
\pgfpathclose%
\pgfusepath{stroke,fill}%
\end{pgfscope}%
\begin{pgfscope}%
\pgfpathrectangle{\pgfqpoint{0.050000in}{0.050000in}}{\pgfqpoint{4.900000in}{4.900000in}}%
\pgfusepath{clip}%
\pgfsetbuttcap%
\pgfsetroundjoin%
\definecolor{currentfill}{rgb}{0.338824,0.338824,0.338824}%
\pgfsetfillcolor{currentfill}%
\pgfsetfillopacity{0.900000}%
\pgfsetlinewidth{0.501875pt}%
\definecolor{currentstroke}{rgb}{0.200000,0.200000,0.200000}%
\pgfsetstrokecolor{currentstroke}%
\pgfsetdash{}{0pt}%
\pgfpathmoveto{\pgfqpoint{1.629352in}{2.973463in}}%
\pgfpathlineto{\pgfqpoint{1.678003in}{2.926819in}}%
\pgfpathlineto{\pgfqpoint{1.709293in}{2.906853in}}%
\pgfpathlineto{\pgfqpoint{1.660633in}{2.953527in}}%
\pgfpathlineto{\pgfqpoint{1.629352in}{2.973463in}}%
\pgfpathclose%
\pgfusepath{stroke,fill}%
\end{pgfscope}%
\begin{pgfscope}%
\pgfpathrectangle{\pgfqpoint{0.050000in}{0.050000in}}{\pgfqpoint{4.900000in}{4.900000in}}%
\pgfusepath{clip}%
\pgfsetbuttcap%
\pgfsetroundjoin%
\definecolor{currentfill}{rgb}{0.595433,0.595433,0.595433}%
\pgfsetfillcolor{currentfill}%
\pgfsetfillopacity{0.900000}%
\pgfsetlinewidth{0.501875pt}%
\definecolor{currentstroke}{rgb}{0.200000,0.200000,0.200000}%
\pgfsetstrokecolor{currentstroke}%
\pgfsetdash{}{0pt}%
\pgfpathmoveto{\pgfqpoint{2.363000in}{2.380536in}}%
\pgfpathlineto{\pgfqpoint{2.410065in}{2.335184in}}%
\pgfpathlineto{\pgfqpoint{2.442517in}{2.314286in}}%
\pgfpathlineto{\pgfqpoint{2.395447in}{2.359661in}}%
\pgfpathlineto{\pgfqpoint{2.363000in}{2.380536in}}%
\pgfpathclose%
\pgfusepath{stroke,fill}%
\end{pgfscope}%
\begin{pgfscope}%
\pgfpathrectangle{\pgfqpoint{0.050000in}{0.050000in}}{\pgfqpoint{4.900000in}{4.900000in}}%
\pgfusepath{clip}%
\pgfsetbuttcap%
\pgfsetroundjoin%
\definecolor{currentfill}{rgb}{0.217762,0.217762,0.217762}%
\pgfsetfillcolor{currentfill}%
\pgfsetfillopacity{0.900000}%
\pgfsetlinewidth{0.501875pt}%
\definecolor{currentstroke}{rgb}{0.200000,0.200000,0.200000}%
\pgfsetstrokecolor{currentstroke}%
\pgfsetdash{}{0pt}%
\pgfpathmoveto{\pgfqpoint{1.358110in}{3.193352in}}%
\pgfpathlineto{\pgfqpoint{1.407357in}{3.146224in}}%
\pgfpathlineto{\pgfqpoint{1.438224in}{3.126602in}}%
\pgfpathlineto{\pgfqpoint{1.388966in}{3.173762in}}%
\pgfpathlineto{\pgfqpoint{1.358110in}{3.193352in}}%
\pgfpathclose%
\pgfusepath{stroke,fill}%
\end{pgfscope}%
\begin{pgfscope}%
\pgfpathrectangle{\pgfqpoint{0.050000in}{0.050000in}}{\pgfqpoint{4.900000in}{4.900000in}}%
\pgfusepath{clip}%
\pgfsetbuttcap%
\pgfsetroundjoin%
\definecolor{currentfill}{rgb}{0.495656,0.495656,0.495656}%
\pgfsetfillcolor{currentfill}%
\pgfsetfillopacity{0.900000}%
\pgfsetlinewidth{0.501875pt}%
\definecolor{currentstroke}{rgb}{0.200000,0.200000,0.200000}%
\pgfsetstrokecolor{currentstroke}%
\pgfsetdash{}{0pt}%
\pgfpathmoveto{\pgfqpoint{2.091866in}{2.600325in}}%
\pgfpathlineto{\pgfqpoint{2.139527in}{2.554478in}}%
\pgfpathlineto{\pgfqpoint{2.171555in}{2.533919in}}%
\pgfpathlineto{\pgfqpoint{2.123889in}{2.579793in}}%
\pgfpathlineto{\pgfqpoint{2.091866in}{2.600325in}}%
\pgfpathclose%
\pgfusepath{stroke,fill}%
\end{pgfscope}%
\begin{pgfscope}%
\pgfpathrectangle{\pgfqpoint{0.050000in}{0.050000in}}{\pgfqpoint{4.900000in}{4.900000in}}%
\pgfusepath{clip}%
\pgfsetbuttcap%
\pgfsetroundjoin%
\definecolor{currentfill}{rgb}{0.912526,0.912526,0.912526}%
\pgfsetfillcolor{currentfill}%
\pgfsetfillopacity{0.900000}%
\pgfsetlinewidth{0.501875pt}%
\definecolor{currentstroke}{rgb}{0.200000,0.200000,0.200000}%
\pgfsetstrokecolor{currentstroke}%
\pgfsetdash{}{0pt}%
\pgfpathmoveto{\pgfqpoint{3.447261in}{1.543817in}}%
\pgfpathlineto{\pgfqpoint{3.492212in}{1.514151in}}%
\pgfpathlineto{\pgfqpoint{3.526413in}{1.493371in}}%
\pgfpathlineto{\pgfqpoint{3.481452in}{1.522856in}}%
\pgfpathlineto{\pgfqpoint{3.447261in}{1.543817in}}%
\pgfpathclose%
\pgfusepath{stroke,fill}%
\end{pgfscope}%
\begin{pgfscope}%
\pgfpathrectangle{\pgfqpoint{0.050000in}{0.050000in}}{\pgfqpoint{4.900000in}{4.900000in}}%
\pgfusepath{clip}%
\pgfsetbuttcap%
\pgfsetroundjoin%
\definecolor{currentfill}{rgb}{0.839785,0.839785,0.839785}%
\pgfsetfillcolor{currentfill}%
\pgfsetfillopacity{0.900000}%
\pgfsetlinewidth{0.501875pt}%
\definecolor{currentstroke}{rgb}{0.200000,0.200000,0.200000}%
\pgfsetstrokecolor{currentstroke}%
\pgfsetdash{}{0pt}%
\pgfpathmoveto{\pgfqpoint{3.096823in}{1.795338in}}%
\pgfpathlineto{\pgfqpoint{3.142351in}{1.755364in}}%
\pgfpathlineto{\pgfqpoint{3.175979in}{1.734181in}}%
\pgfpathlineto{\pgfqpoint{3.130447in}{1.774014in}}%
\pgfpathlineto{\pgfqpoint{3.096823in}{1.795338in}}%
\pgfpathclose%
\pgfusepath{stroke,fill}%
\end{pgfscope}%
\begin{pgfscope}%
\pgfpathrectangle{\pgfqpoint{0.050000in}{0.050000in}}{\pgfqpoint{4.900000in}{4.900000in}}%
\pgfusepath{clip}%
\pgfsetbuttcap%
\pgfsetroundjoin%
\definecolor{currentfill}{rgb}{0.763998,0.763998,0.763998}%
\pgfsetfillcolor{currentfill}%
\pgfsetfillopacity{0.900000}%
\pgfsetlinewidth{0.501875pt}%
\definecolor{currentstroke}{rgb}{0.200000,0.200000,0.200000}%
\pgfsetstrokecolor{currentstroke}%
\pgfsetdash{}{0pt}%
\pgfpathmoveto{\pgfqpoint{2.825728in}{2.008434in}}%
\pgfpathlineto{\pgfqpoint{2.871808in}{1.964804in}}%
\pgfpathlineto{\pgfqpoint{2.905002in}{1.943513in}}%
\pgfpathlineto{\pgfqpoint{2.858919in}{1.987093in}}%
\pgfpathlineto{\pgfqpoint{2.825728in}{2.008434in}}%
\pgfpathclose%
\pgfusepath{stroke,fill}%
\end{pgfscope}%
\begin{pgfscope}%
\pgfpathrectangle{\pgfqpoint{0.050000in}{0.050000in}}{\pgfqpoint{4.900000in}{4.900000in}}%
\pgfusepath{clip}%
\pgfsetbuttcap%
\pgfsetroundjoin%
\definecolor{currentfill}{rgb}{0.104698,0.104698,0.104698}%
\pgfsetfillcolor{currentfill}%
\pgfsetfillopacity{0.900000}%
\pgfsetlinewidth{0.501875pt}%
\definecolor{currentstroke}{rgb}{0.200000,0.200000,0.200000}%
\pgfsetstrokecolor{currentstroke}%
\pgfsetdash{}{0pt}%
\pgfpathmoveto{\pgfqpoint{1.086774in}{3.413320in}}%
\pgfpathlineto{\pgfqpoint{1.136618in}{3.365708in}}%
\pgfpathlineto{\pgfqpoint{1.167061in}{3.346429in}}%
\pgfpathlineto{\pgfqpoint{1.117205in}{3.394075in}}%
\pgfpathlineto{\pgfqpoint{1.086774in}{3.413320in}}%
\pgfpathclose%
\pgfusepath{stroke,fill}%
\end{pgfscope}%
\begin{pgfscope}%
\pgfpathrectangle{\pgfqpoint{0.050000in}{0.050000in}}{\pgfqpoint{4.900000in}{4.900000in}}%
\pgfusepath{clip}%
\pgfsetbuttcap%
\pgfsetroundjoin%
\definecolor{currentfill}{rgb}{0.399723,0.399723,0.399723}%
\pgfsetfillcolor{currentfill}%
\pgfsetfillopacity{0.900000}%
\pgfsetlinewidth{0.501875pt}%
\definecolor{currentstroke}{rgb}{0.200000,0.200000,0.200000}%
\pgfsetstrokecolor{currentstroke}%
\pgfsetdash{}{0pt}%
\pgfpathmoveto{\pgfqpoint{1.820638in}{2.820198in}}%
\pgfpathlineto{\pgfqpoint{1.868895in}{2.773868in}}%
\pgfpathlineto{\pgfqpoint{1.900500in}{2.753651in}}%
\pgfpathlineto{\pgfqpoint{1.852236in}{2.800012in}}%
\pgfpathlineto{\pgfqpoint{1.820638in}{2.820198in}}%
\pgfpathclose%
\pgfusepath{stroke,fill}%
\end{pgfscope}%
\begin{pgfscope}%
\pgfpathrectangle{\pgfqpoint{0.050000in}{0.050000in}}{\pgfqpoint{4.900000in}{4.900000in}}%
\pgfusepath{clip}%
\pgfsetbuttcap%
\pgfsetroundjoin%
\definecolor{currentfill}{rgb}{0.667405,0.667405,0.667405}%
\pgfsetfillcolor{currentfill}%
\pgfsetfillopacity{0.900000}%
\pgfsetlinewidth{0.501875pt}%
\definecolor{currentstroke}{rgb}{0.200000,0.200000,0.200000}%
\pgfsetstrokecolor{currentstroke}%
\pgfsetdash{}{0pt}%
\pgfpathmoveto{\pgfqpoint{2.554603in}{2.227097in}}%
\pgfpathlineto{\pgfqpoint{2.601272in}{2.182147in}}%
\pgfpathlineto{\pgfqpoint{2.634041in}{2.161025in}}%
\pgfpathlineto{\pgfqpoint{2.587368in}{2.205984in}}%
\pgfpathlineto{\pgfqpoint{2.554603in}{2.227097in}}%
\pgfpathclose%
\pgfusepath{stroke,fill}%
\end{pgfscope}%
\begin{pgfscope}%
\pgfpathrectangle{\pgfqpoint{0.050000in}{0.050000in}}{\pgfqpoint{4.900000in}{4.900000in}}%
\pgfusepath{clip}%
\pgfsetbuttcap%
\pgfsetroundjoin%
\definecolor{currentfill}{rgb}{0.000000,0.000000,0.000000}%
\pgfsetfillcolor{currentfill}%
\pgfsetfillopacity{0.900000}%
\pgfsetlinewidth{0.501875pt}%
\definecolor{currentstroke}{rgb}{0.200000,0.200000,0.200000}%
\pgfsetstrokecolor{currentstroke}%
\pgfsetdash{}{0pt}%
\pgfpathmoveto{\pgfqpoint{0.815345in}{3.633366in}}%
\pgfpathlineto{\pgfqpoint{0.865785in}{3.585269in}}%
\pgfpathlineto{\pgfqpoint{0.895804in}{3.566334in}}%
\pgfpathlineto{\pgfqpoint{0.845350in}{3.614466in}}%
\pgfpathlineto{\pgfqpoint{0.815345in}{3.633366in}}%
\pgfpathclose%
\pgfusepath{stroke,fill}%
\end{pgfscope}%
\begin{pgfscope}%
\pgfpathrectangle{\pgfqpoint{0.050000in}{0.050000in}}{\pgfqpoint{4.900000in}{4.900000in}}%
\pgfusepath{clip}%
\pgfsetbuttcap%
\pgfsetroundjoin%
\definecolor{currentfill}{rgb}{0.306344,0.306344,0.306344}%
\pgfsetfillcolor{currentfill}%
\pgfsetfillopacity{0.900000}%
\pgfsetlinewidth{0.501875pt}%
\definecolor{currentstroke}{rgb}{0.200000,0.200000,0.200000}%
\pgfsetstrokecolor{currentstroke}%
\pgfsetdash{}{0pt}%
\pgfpathmoveto{\pgfqpoint{1.549317in}{3.040151in}}%
\pgfpathlineto{\pgfqpoint{1.598170in}{2.993335in}}%
\pgfpathlineto{\pgfqpoint{1.629352in}{2.973463in}}%
\pgfpathlineto{\pgfqpoint{1.580489in}{3.020309in}}%
\pgfpathlineto{\pgfqpoint{1.549317in}{3.040151in}}%
\pgfpathclose%
\pgfusepath{stroke,fill}%
\end{pgfscope}%
\begin{pgfscope}%
\pgfpathrectangle{\pgfqpoint{0.050000in}{0.050000in}}{\pgfqpoint{4.900000in}{4.900000in}}%
\pgfusepath{clip}%
\pgfsetbuttcap%
\pgfsetroundjoin%
\definecolor{currentfill}{rgb}{0.895548,0.895548,0.895548}%
\pgfsetfillcolor{currentfill}%
\pgfsetfillopacity{0.900000}%
\pgfsetlinewidth{0.501875pt}%
\definecolor{currentstroke}{rgb}{0.200000,0.200000,0.200000}%
\pgfsetstrokecolor{currentstroke}%
\pgfsetdash{}{0pt}%
\pgfpathmoveto{\pgfqpoint{3.368106in}{1.597448in}}%
\pgfpathlineto{\pgfqpoint{3.413180in}{1.564710in}}%
\pgfpathlineto{\pgfqpoint{3.447261in}{1.543817in}}%
\pgfpathlineto{\pgfqpoint{3.402179in}{1.576370in}}%
\pgfpathlineto{\pgfqpoint{3.368106in}{1.597448in}}%
\pgfpathclose%
\pgfusepath{stroke,fill}%
\end{pgfscope}%
\begin{pgfscope}%
\pgfpathrectangle{\pgfqpoint{0.050000in}{0.050000in}}{\pgfqpoint{4.900000in}{4.900000in}}%
\pgfusepath{clip}%
\pgfsetbuttcap%
\pgfsetroundjoin%
\definecolor{currentfill}{rgb}{0.560246,0.560246,0.560246}%
\pgfsetfillcolor{currentfill}%
\pgfsetfillopacity{0.900000}%
\pgfsetlinewidth{0.501875pt}%
\definecolor{currentstroke}{rgb}{0.200000,0.200000,0.200000}%
\pgfsetstrokecolor{currentstroke}%
\pgfsetdash{}{0pt}%
\pgfpathmoveto{\pgfqpoint{2.283389in}{2.446874in}}%
\pgfpathlineto{\pgfqpoint{2.330655in}{2.401345in}}%
\pgfpathlineto{\pgfqpoint{2.363000in}{2.380536in}}%
\pgfpathlineto{\pgfqpoint{2.315730in}{2.426091in}}%
\pgfpathlineto{\pgfqpoint{2.283389in}{2.446874in}}%
\pgfpathclose%
\pgfusepath{stroke,fill}%
\end{pgfscope}%
\begin{pgfscope}%
\pgfpathrectangle{\pgfqpoint{0.050000in}{0.050000in}}{\pgfqpoint{4.900000in}{4.900000in}}%
\pgfusepath{clip}%
\pgfsetbuttcap%
\pgfsetroundjoin%
\definecolor{currentfill}{rgb}{0.179008,0.179008,0.179008}%
\pgfsetfillcolor{currentfill}%
\pgfsetfillopacity{0.900000}%
\pgfsetlinewidth{0.501875pt}%
\definecolor{currentstroke}{rgb}{0.200000,0.200000,0.200000}%
\pgfsetstrokecolor{currentstroke}%
\pgfsetdash{}{0pt}%
\pgfpathmoveto{\pgfqpoint{1.277902in}{3.260181in}}%
\pgfpathlineto{\pgfqpoint{1.327352in}{3.212881in}}%
\pgfpathlineto{\pgfqpoint{1.358110in}{3.193352in}}%
\pgfpathlineto{\pgfqpoint{1.308649in}{3.240685in}}%
\pgfpathlineto{\pgfqpoint{1.277902in}{3.260181in}}%
\pgfpathclose%
\pgfusepath{stroke,fill}%
\end{pgfscope}%
\begin{pgfscope}%
\pgfpathrectangle{\pgfqpoint{0.050000in}{0.050000in}}{\pgfqpoint{4.900000in}{4.900000in}}%
\pgfusepath{clip}%
\pgfsetbuttcap%
\pgfsetroundjoin%
\definecolor{currentfill}{rgb}{0.461207,0.461207,0.461207}%
\pgfsetfillcolor{currentfill}%
\pgfsetfillopacity{0.900000}%
\pgfsetlinewidth{0.501875pt}%
\definecolor{currentstroke}{rgb}{0.200000,0.200000,0.200000}%
\pgfsetstrokecolor{currentstroke}%
\pgfsetdash{}{0pt}%
\pgfpathmoveto{\pgfqpoint{2.012083in}{2.666809in}}%
\pgfpathlineto{\pgfqpoint{2.059944in}{2.620791in}}%
\pgfpathlineto{\pgfqpoint{2.091866in}{2.600325in}}%
\pgfpathlineto{\pgfqpoint{2.043998in}{2.646370in}}%
\pgfpathlineto{\pgfqpoint{2.012083in}{2.666809in}}%
\pgfpathclose%
\pgfusepath{stroke,fill}%
\end{pgfscope}%
\begin{pgfscope}%
\pgfpathrectangle{\pgfqpoint{0.050000in}{0.050000in}}{\pgfqpoint{4.900000in}{4.900000in}}%
\pgfusepath{clip}%
\pgfsetbuttcap%
\pgfsetroundjoin%
\definecolor{currentfill}{rgb}{0.815671,0.815671,0.815671}%
\pgfsetfillcolor{currentfill}%
\pgfsetfillopacity{0.900000}%
\pgfsetlinewidth{0.501875pt}%
\definecolor{currentstroke}{rgb}{0.200000,0.200000,0.200000}%
\pgfsetstrokecolor{currentstroke}%
\pgfsetdash{}{0pt}%
\pgfpathmoveto{\pgfqpoint{3.017600in}{1.858132in}}%
\pgfpathlineto{\pgfqpoint{3.063305in}{1.816607in}}%
\pgfpathlineto{\pgfqpoint{3.096823in}{1.795338in}}%
\pgfpathlineto{\pgfqpoint{3.051114in}{1.836751in}}%
\pgfpathlineto{\pgfqpoint{3.017600in}{1.858132in}}%
\pgfpathclose%
\pgfusepath{stroke,fill}%
\end{pgfscope}%
\begin{pgfscope}%
\pgfpathrectangle{\pgfqpoint{0.050000in}{0.050000in}}{\pgfqpoint{4.900000in}{4.900000in}}%
\pgfusepath{clip}%
\pgfsetbuttcap%
\pgfsetroundjoin%
\definecolor{currentfill}{rgb}{0.739377,0.739377,0.739377}%
\pgfsetfillcolor{currentfill}%
\pgfsetfillopacity{0.900000}%
\pgfsetlinewidth{0.501875pt}%
\definecolor{currentstroke}{rgb}{0.200000,0.200000,0.200000}%
\pgfsetstrokecolor{currentstroke}%
\pgfsetdash{}{0pt}%
\pgfpathmoveto{\pgfqpoint{2.746366in}{2.073943in}}%
\pgfpathlineto{\pgfqpoint{2.792642in}{2.029715in}}%
\pgfpathlineto{\pgfqpoint{2.825728in}{2.008434in}}%
\pgfpathlineto{\pgfqpoint{2.779450in}{2.052637in}}%
\pgfpathlineto{\pgfqpoint{2.746366in}{2.073943in}}%
\pgfpathclose%
\pgfusepath{stroke,fill}%
\end{pgfscope}%
\begin{pgfscope}%
\pgfpathrectangle{\pgfqpoint{0.050000in}{0.050000in}}{\pgfqpoint{4.900000in}{4.900000in}}%
\pgfusepath{clip}%
\pgfsetbuttcap%
\pgfsetroundjoin%
\definecolor{currentfill}{rgb}{0.068281,0.068281,0.068281}%
\pgfsetfillcolor{currentfill}%
\pgfsetfillopacity{0.900000}%
\pgfsetlinewidth{0.501875pt}%
\definecolor{currentstroke}{rgb}{0.200000,0.200000,0.200000}%
\pgfsetstrokecolor{currentstroke}%
\pgfsetdash{}{0pt}%
\pgfpathmoveto{\pgfqpoint{1.006394in}{3.480290in}}%
\pgfpathlineto{\pgfqpoint{1.056440in}{3.432505in}}%
\pgfpathlineto{\pgfqpoint{1.086774in}{3.413320in}}%
\pgfpathlineto{\pgfqpoint{1.036714in}{3.461140in}}%
\pgfpathlineto{\pgfqpoint{1.006394in}{3.480290in}}%
\pgfpathclose%
\pgfusepath{stroke,fill}%
\end{pgfscope}%
\begin{pgfscope}%
\pgfpathrectangle{\pgfqpoint{0.050000in}{0.050000in}}{\pgfqpoint{4.900000in}{4.900000in}}%
\pgfusepath{clip}%
\pgfsetbuttcap%
\pgfsetroundjoin%
\definecolor{currentfill}{rgb}{0.878570,0.878570,0.878570}%
\pgfsetfillcolor{currentfill}%
\pgfsetfillopacity{0.900000}%
\pgfsetlinewidth{0.501875pt}%
\definecolor{currentstroke}{rgb}{0.200000,0.200000,0.200000}%
\pgfsetstrokecolor{currentstroke}%
\pgfsetdash{}{0pt}%
\pgfpathmoveto{\pgfqpoint{3.288929in}{1.653948in}}%
\pgfpathlineto{\pgfqpoint{3.334142in}{1.618463in}}%
\pgfpathlineto{\pgfqpoint{3.368106in}{1.597448in}}%
\pgfpathlineto{\pgfqpoint{3.322887in}{1.632754in}}%
\pgfpathlineto{\pgfqpoint{3.288929in}{1.653948in}}%
\pgfpathclose%
\pgfusepath{stroke,fill}%
\end{pgfscope}%
\begin{pgfscope}%
\pgfpathrectangle{\pgfqpoint{0.050000in}{0.050000in}}{\pgfqpoint{4.900000in}{4.900000in}}%
\pgfusepath{clip}%
\pgfsetbuttcap%
\pgfsetroundjoin%
\definecolor{currentfill}{rgb}{0.371303,0.371303,0.371303}%
\pgfsetfillcolor{currentfill}%
\pgfsetfillopacity{0.900000}%
\pgfsetlinewidth{0.501875pt}%
\definecolor{currentstroke}{rgb}{0.200000,0.200000,0.200000}%
\pgfsetstrokecolor{currentstroke}%
\pgfsetdash{}{0pt}%
\pgfpathmoveto{\pgfqpoint{1.740683in}{2.886824in}}%
\pgfpathlineto{\pgfqpoint{1.789141in}{2.840321in}}%
\pgfpathlineto{\pgfqpoint{1.820638in}{2.820198in}}%
\pgfpathlineto{\pgfqpoint{1.772172in}{2.866731in}}%
\pgfpathlineto{\pgfqpoint{1.740683in}{2.886824in}}%
\pgfpathclose%
\pgfusepath{stroke,fill}%
\end{pgfscope}%
\begin{pgfscope}%
\pgfpathrectangle{\pgfqpoint{0.050000in}{0.050000in}}{\pgfqpoint{4.900000in}{4.900000in}}%
\pgfusepath{clip}%
\pgfsetbuttcap%
\pgfsetroundjoin%
\definecolor{currentfill}{rgb}{0.629020,0.629020,0.629020}%
\pgfsetfillcolor{currentfill}%
\pgfsetfillopacity{0.900000}%
\pgfsetlinewidth{0.501875pt}%
\definecolor{currentstroke}{rgb}{0.200000,0.200000,0.200000}%
\pgfsetstrokecolor{currentstroke}%
\pgfsetdash{}{0pt}%
\pgfpathmoveto{\pgfqpoint{2.475072in}{2.293323in}}%
\pgfpathlineto{\pgfqpoint{2.521941in}{2.248144in}}%
\pgfpathlineto{\pgfqpoint{2.554603in}{2.227097in}}%
\pgfpathlineto{\pgfqpoint{2.507730in}{2.272294in}}%
\pgfpathlineto{\pgfqpoint{2.475072in}{2.293323in}}%
\pgfpathclose%
\pgfusepath{stroke,fill}%
\end{pgfscope}%
\begin{pgfscope}%
\pgfpathrectangle{\pgfqpoint{0.050000in}{0.050000in}}{\pgfqpoint{4.900000in}{4.900000in}}%
\pgfusepath{clip}%
\pgfsetbuttcap%
\pgfsetroundjoin%
\definecolor{currentfill}{rgb}{0.262053,0.262053,0.262053}%
\pgfsetfillcolor{currentfill}%
\pgfsetfillopacity{0.900000}%
\pgfsetlinewidth{0.501875pt}%
\definecolor{currentstroke}{rgb}{0.200000,0.200000,0.200000}%
\pgfsetstrokecolor{currentstroke}%
\pgfsetdash{}{0pt}%
\pgfpathmoveto{\pgfqpoint{1.469189in}{3.106917in}}%
\pgfpathlineto{\pgfqpoint{1.518244in}{3.059929in}}%
\pgfpathlineto{\pgfqpoint{1.549317in}{3.040151in}}%
\pgfpathlineto{\pgfqpoint{1.500252in}{3.087170in}}%
\pgfpathlineto{\pgfqpoint{1.469189in}{3.106917in}}%
\pgfpathclose%
\pgfusepath{stroke,fill}%
\end{pgfscope}%
\begin{pgfscope}%
\pgfpathrectangle{\pgfqpoint{0.050000in}{0.050000in}}{\pgfqpoint{4.900000in}{4.900000in}}%
\pgfusepath{clip}%
\pgfsetbuttcap%
\pgfsetroundjoin%
\definecolor{currentfill}{rgb}{0.525798,0.525798,0.525798}%
\pgfsetfillcolor{currentfill}%
\pgfsetfillopacity{0.900000}%
\pgfsetlinewidth{0.501875pt}%
\definecolor{currentstroke}{rgb}{0.200000,0.200000,0.200000}%
\pgfsetstrokecolor{currentstroke}%
\pgfsetdash{}{0pt}%
\pgfpathmoveto{\pgfqpoint{2.203686in}{2.513294in}}%
\pgfpathlineto{\pgfqpoint{2.251152in}{2.467591in}}%
\pgfpathlineto{\pgfqpoint{2.283389in}{2.446874in}}%
\pgfpathlineto{\pgfqpoint{2.235918in}{2.492604in}}%
\pgfpathlineto{\pgfqpoint{2.203686in}{2.513294in}}%
\pgfpathclose%
\pgfusepath{stroke,fill}%
\end{pgfscope}%
\begin{pgfscope}%
\pgfpathrectangle{\pgfqpoint{0.050000in}{0.050000in}}{\pgfqpoint{4.900000in}{4.900000in}}%
\pgfusepath{clip}%
\pgfsetbuttcap%
\pgfsetroundjoin%
\definecolor{currentfill}{rgb}{0.136563,0.136563,0.136563}%
\pgfsetfillcolor{currentfill}%
\pgfsetfillopacity{0.900000}%
\pgfsetlinewidth{0.501875pt}%
\definecolor{currentstroke}{rgb}{0.200000,0.200000,0.200000}%
\pgfsetstrokecolor{currentstroke}%
\pgfsetdash{}{0pt}%
\pgfpathmoveto{\pgfqpoint{1.197601in}{3.327089in}}%
\pgfpathlineto{\pgfqpoint{1.247253in}{3.279616in}}%
\pgfpathlineto{\pgfqpoint{1.277902in}{3.260181in}}%
\pgfpathlineto{\pgfqpoint{1.228238in}{3.307687in}}%
\pgfpathlineto{\pgfqpoint{1.197601in}{3.327089in}}%
\pgfpathclose%
\pgfusepath{stroke,fill}%
\end{pgfscope}%
\begin{pgfscope}%
\pgfpathrectangle{\pgfqpoint{0.050000in}{0.050000in}}{\pgfqpoint{4.900000in}{4.900000in}}%
\pgfusepath{clip}%
\pgfsetbuttcap%
\pgfsetroundjoin%
\definecolor{currentfill}{rgb}{0.788112,0.788112,0.788112}%
\pgfsetfillcolor{currentfill}%
\pgfsetfillopacity{0.900000}%
\pgfsetlinewidth{0.501875pt}%
\definecolor{currentstroke}{rgb}{0.200000,0.200000,0.200000}%
\pgfsetstrokecolor{currentstroke}%
\pgfsetdash{}{0pt}%
\pgfpathmoveto{\pgfqpoint{2.938301in}{1.922165in}}%
\pgfpathlineto{\pgfqpoint{2.984192in}{1.879456in}}%
\pgfpathlineto{\pgfqpoint{3.017600in}{1.858132in}}%
\pgfpathlineto{\pgfqpoint{2.971706in}{1.900758in}}%
\pgfpathlineto{\pgfqpoint{2.938301in}{1.922165in}}%
\pgfpathclose%
\pgfusepath{stroke,fill}%
\end{pgfscope}%
\begin{pgfscope}%
\pgfpathrectangle{\pgfqpoint{0.050000in}{0.050000in}}{\pgfqpoint{4.900000in}{4.900000in}}%
\pgfusepath{clip}%
\pgfsetbuttcap%
\pgfsetroundjoin%
\definecolor{currentfill}{rgb}{0.940823,0.940823,0.940823}%
\pgfsetfillcolor{currentfill}%
\pgfsetfillopacity{0.900000}%
\pgfsetlinewidth{0.501875pt}%
\definecolor{currentstroke}{rgb}{0.200000,0.200000,0.200000}%
\pgfsetstrokecolor{currentstroke}%
\pgfsetdash{}{0pt}%
\pgfpathmoveto{\pgfqpoint{3.640011in}{1.425588in}}%
\pgfpathlineto{\pgfqpoint{3.684775in}{1.402884in}}%
\pgfpathlineto{\pgfqpoint{3.719333in}{1.382195in}}%
\pgfpathlineto{\pgfqpoint{3.674556in}{1.404755in}}%
\pgfpathlineto{\pgfqpoint{3.640011in}{1.425588in}}%
\pgfpathclose%
\pgfusepath{stroke,fill}%
\end{pgfscope}%
\begin{pgfscope}%
\pgfpathrectangle{\pgfqpoint{0.050000in}{0.050000in}}{\pgfqpoint{4.900000in}{4.900000in}}%
\pgfusepath{clip}%
\pgfsetbuttcap%
\pgfsetroundjoin%
\definecolor{currentfill}{rgb}{0.432203,0.432203,0.432203}%
\pgfsetfillcolor{currentfill}%
\pgfsetfillopacity{0.900000}%
\pgfsetlinewidth{0.501875pt}%
\definecolor{currentstroke}{rgb}{0.200000,0.200000,0.200000}%
\pgfsetstrokecolor{currentstroke}%
\pgfsetdash{}{0pt}%
\pgfpathmoveto{\pgfqpoint{1.932206in}{2.733371in}}%
\pgfpathlineto{\pgfqpoint{1.980269in}{2.687182in}}%
\pgfpathlineto{\pgfqpoint{2.012083in}{2.666809in}}%
\pgfpathlineto{\pgfqpoint{1.964013in}{2.713026in}}%
\pgfpathlineto{\pgfqpoint{1.932206in}{2.733371in}}%
\pgfpathclose%
\pgfusepath{stroke,fill}%
\end{pgfscope}%
\begin{pgfscope}%
\pgfpathrectangle{\pgfqpoint{0.050000in}{0.050000in}}{\pgfqpoint{4.900000in}{4.900000in}}%
\pgfusepath{clip}%
\pgfsetbuttcap%
\pgfsetroundjoin%
\definecolor{currentfill}{rgb}{0.858762,0.858762,0.858762}%
\pgfsetfillcolor{currentfill}%
\pgfsetfillopacity{0.900000}%
\pgfsetlinewidth{0.501875pt}%
\definecolor{currentstroke}{rgb}{0.200000,0.200000,0.200000}%
\pgfsetstrokecolor{currentstroke}%
\pgfsetdash{}{0pt}%
\pgfpathmoveto{\pgfqpoint{3.209714in}{1.712940in}}%
\pgfpathlineto{\pgfqpoint{3.255080in}{1.675082in}}%
\pgfpathlineto{\pgfqpoint{3.288929in}{1.653948in}}%
\pgfpathlineto{\pgfqpoint{3.243558in}{1.691642in}}%
\pgfpathlineto{\pgfqpoint{3.209714in}{1.712940in}}%
\pgfpathclose%
\pgfusepath{stroke,fill}%
\end{pgfscope}%
\begin{pgfscope}%
\pgfpathrectangle{\pgfqpoint{0.050000in}{0.050000in}}{\pgfqpoint{4.900000in}{4.900000in}}%
\pgfusepath{clip}%
\pgfsetbuttcap%
\pgfsetroundjoin%
\definecolor{currentfill}{rgb}{0.700992,0.700992,0.700992}%
\pgfsetfillcolor{currentfill}%
\pgfsetfillopacity{0.900000}%
\pgfsetlinewidth{0.501875pt}%
\definecolor{currentstroke}{rgb}{0.200000,0.200000,0.200000}%
\pgfsetstrokecolor{currentstroke}%
\pgfsetdash{}{0pt}%
\pgfpathmoveto{\pgfqpoint{2.666913in}{2.139839in}}%
\pgfpathlineto{\pgfqpoint{2.713387in}{2.095188in}}%
\pgfpathlineto{\pgfqpoint{2.746366in}{2.073943in}}%
\pgfpathlineto{\pgfqpoint{2.699890in}{2.118590in}}%
\pgfpathlineto{\pgfqpoint{2.666913in}{2.139839in}}%
\pgfpathclose%
\pgfusepath{stroke,fill}%
\end{pgfscope}%
\begin{pgfscope}%
\pgfpathrectangle{\pgfqpoint{0.050000in}{0.050000in}}{\pgfqpoint{4.900000in}{4.900000in}}%
\pgfusepath{clip}%
\pgfsetbuttcap%
\pgfsetroundjoin%
\definecolor{currentfill}{rgb}{0.036417,0.036417,0.036417}%
\pgfsetfillcolor{currentfill}%
\pgfsetfillopacity{0.900000}%
\pgfsetlinewidth{0.501875pt}%
\definecolor{currentstroke}{rgb}{0.200000,0.200000,0.200000}%
\pgfsetstrokecolor{currentstroke}%
\pgfsetdash{}{0pt}%
\pgfpathmoveto{\pgfqpoint{0.925919in}{3.547339in}}%
\pgfpathlineto{\pgfqpoint{0.976169in}{3.499380in}}%
\pgfpathlineto{\pgfqpoint{1.006394in}{3.480290in}}%
\pgfpathlineto{\pgfqpoint{0.956130in}{3.528284in}}%
\pgfpathlineto{\pgfqpoint{0.925919in}{3.547339in}}%
\pgfpathclose%
\pgfusepath{stroke,fill}%
\end{pgfscope}%
\begin{pgfscope}%
\pgfpathrectangle{\pgfqpoint{0.050000in}{0.050000in}}{\pgfqpoint{4.900000in}{4.900000in}}%
\pgfusepath{clip}%
\pgfsetbuttcap%
\pgfsetroundjoin%
\definecolor{currentfill}{rgb}{0.338824,0.338824,0.338824}%
\pgfsetfillcolor{currentfill}%
\pgfsetfillopacity{0.900000}%
\pgfsetlinewidth{0.501875pt}%
\definecolor{currentstroke}{rgb}{0.200000,0.200000,0.200000}%
\pgfsetstrokecolor{currentstroke}%
\pgfsetdash{}{0pt}%
\pgfpathmoveto{\pgfqpoint{1.660633in}{2.953527in}}%
\pgfpathlineto{\pgfqpoint{1.709293in}{2.906853in}}%
\pgfpathlineto{\pgfqpoint{1.740683in}{2.886824in}}%
\pgfpathlineto{\pgfqpoint{1.692013in}{2.933528in}}%
\pgfpathlineto{\pgfqpoint{1.660633in}{2.953527in}}%
\pgfpathclose%
\pgfusepath{stroke,fill}%
\end{pgfscope}%
\begin{pgfscope}%
\pgfpathrectangle{\pgfqpoint{0.050000in}{0.050000in}}{\pgfqpoint{4.900000in}{4.900000in}}%
\pgfusepath{clip}%
\pgfsetbuttcap%
\pgfsetroundjoin%
\definecolor{currentfill}{rgb}{0.926674,0.926674,0.926674}%
\pgfsetfillcolor{currentfill}%
\pgfsetfillopacity{0.900000}%
\pgfsetlinewidth{0.501875pt}%
\definecolor{currentstroke}{rgb}{0.200000,0.200000,0.200000}%
\pgfsetstrokecolor{currentstroke}%
\pgfsetdash{}{0pt}%
\pgfpathmoveto{\pgfqpoint{3.560724in}{1.472520in}}%
\pgfpathlineto{\pgfqpoint{3.605578in}{1.446346in}}%
\pgfpathlineto{\pgfqpoint{3.640011in}{1.425588in}}%
\pgfpathlineto{\pgfqpoint{3.595146in}{1.451599in}}%
\pgfpathlineto{\pgfqpoint{3.560724in}{1.472520in}}%
\pgfpathclose%
\pgfusepath{stroke,fill}%
\end{pgfscope}%
\begin{pgfscope}%
\pgfpathrectangle{\pgfqpoint{0.050000in}{0.050000in}}{\pgfqpoint{4.900000in}{4.900000in}}%
\pgfusepath{clip}%
\pgfsetbuttcap%
\pgfsetroundjoin%
\definecolor{currentfill}{rgb}{0.595433,0.595433,0.595433}%
\pgfsetfillcolor{currentfill}%
\pgfsetfillopacity{0.900000}%
\pgfsetlinewidth{0.501875pt}%
\definecolor{currentstroke}{rgb}{0.200000,0.200000,0.200000}%
\pgfsetstrokecolor{currentstroke}%
\pgfsetdash{}{0pt}%
\pgfpathmoveto{\pgfqpoint{2.395447in}{2.359661in}}%
\pgfpathlineto{\pgfqpoint{2.442517in}{2.314286in}}%
\pgfpathlineto{\pgfqpoint{2.475072in}{2.293323in}}%
\pgfpathlineto{\pgfqpoint{2.427998in}{2.338721in}}%
\pgfpathlineto{\pgfqpoint{2.395447in}{2.359661in}}%
\pgfpathclose%
\pgfusepath{stroke,fill}%
\end{pgfscope}%
\begin{pgfscope}%
\pgfpathrectangle{\pgfqpoint{0.050000in}{0.050000in}}{\pgfqpoint{4.900000in}{4.900000in}}%
\pgfusepath{clip}%
\pgfsetbuttcap%
\pgfsetroundjoin%
\definecolor{currentfill}{rgb}{0.223299,0.223299,0.223299}%
\pgfsetfillcolor{currentfill}%
\pgfsetfillopacity{0.900000}%
\pgfsetlinewidth{0.501875pt}%
\definecolor{currentstroke}{rgb}{0.200000,0.200000,0.200000}%
\pgfsetstrokecolor{currentstroke}%
\pgfsetdash{}{0pt}%
\pgfpathmoveto{\pgfqpoint{1.388966in}{3.173762in}}%
\pgfpathlineto{\pgfqpoint{1.438224in}{3.126602in}}%
\pgfpathlineto{\pgfqpoint{1.469189in}{3.106917in}}%
\pgfpathlineto{\pgfqpoint{1.419920in}{3.154109in}}%
\pgfpathlineto{\pgfqpoint{1.388966in}{3.173762in}}%
\pgfpathclose%
\pgfusepath{stroke,fill}%
\end{pgfscope}%
\begin{pgfscope}%
\pgfpathrectangle{\pgfqpoint{0.050000in}{0.050000in}}{\pgfqpoint{4.900000in}{4.900000in}}%
\pgfusepath{clip}%
\pgfsetbuttcap%
\pgfsetroundjoin%
\definecolor{currentfill}{rgb}{0.495656,0.495656,0.495656}%
\pgfsetfillcolor{currentfill}%
\pgfsetfillopacity{0.900000}%
\pgfsetlinewidth{0.501875pt}%
\definecolor{currentstroke}{rgb}{0.200000,0.200000,0.200000}%
\pgfsetstrokecolor{currentstroke}%
\pgfsetdash{}{0pt}%
\pgfpathmoveto{\pgfqpoint{2.123889in}{2.579793in}}%
\pgfpathlineto{\pgfqpoint{2.171555in}{2.533919in}}%
\pgfpathlineto{\pgfqpoint{2.203686in}{2.513294in}}%
\pgfpathlineto{\pgfqpoint{2.156013in}{2.559196in}}%
\pgfpathlineto{\pgfqpoint{2.123889in}{2.579793in}}%
\pgfpathclose%
\pgfusepath{stroke,fill}%
\end{pgfscope}%
\begin{pgfscope}%
\pgfpathrectangle{\pgfqpoint{0.050000in}{0.050000in}}{\pgfqpoint{4.900000in}{4.900000in}}%
\pgfusepath{clip}%
\pgfsetbuttcap%
\pgfsetroundjoin%
\definecolor{currentfill}{rgb}{0.912526,0.912526,0.912526}%
\pgfsetfillcolor{currentfill}%
\pgfsetfillopacity{0.900000}%
\pgfsetlinewidth{0.501875pt}%
\definecolor{currentstroke}{rgb}{0.200000,0.200000,0.200000}%
\pgfsetstrokecolor{currentstroke}%
\pgfsetdash{}{0pt}%
\pgfpathmoveto{\pgfqpoint{3.481452in}{1.522856in}}%
\pgfpathlineto{\pgfqpoint{3.526413in}{1.493371in}}%
\pgfpathlineto{\pgfqpoint{3.560724in}{1.472520in}}%
\pgfpathlineto{\pgfqpoint{3.515754in}{1.501829in}}%
\pgfpathlineto{\pgfqpoint{3.481452in}{1.522856in}}%
\pgfpathclose%
\pgfusepath{stroke,fill}%
\end{pgfscope}%
\begin{pgfscope}%
\pgfpathrectangle{\pgfqpoint{0.050000in}{0.050000in}}{\pgfqpoint{4.900000in}{4.900000in}}%
\pgfusepath{clip}%
\pgfsetbuttcap%
\pgfsetroundjoin%
\definecolor{currentfill}{rgb}{0.839785,0.839785,0.839785}%
\pgfsetfillcolor{currentfill}%
\pgfsetfillopacity{0.900000}%
\pgfsetlinewidth{0.501875pt}%
\definecolor{currentstroke}{rgb}{0.200000,0.200000,0.200000}%
\pgfsetstrokecolor{currentstroke}%
\pgfsetdash{}{0pt}%
\pgfpathmoveto{\pgfqpoint{3.130447in}{1.774014in}}%
\pgfpathlineto{\pgfqpoint{3.175979in}{1.734181in}}%
\pgfpathlineto{\pgfqpoint{3.209714in}{1.712940in}}%
\pgfpathlineto{\pgfqpoint{3.164179in}{1.752632in}}%
\pgfpathlineto{\pgfqpoint{3.130447in}{1.774014in}}%
\pgfpathclose%
\pgfusepath{stroke,fill}%
\end{pgfscope}%
\begin{pgfscope}%
\pgfpathrectangle{\pgfqpoint{0.050000in}{0.050000in}}{\pgfqpoint{4.900000in}{4.900000in}}%
\pgfusepath{clip}%
\pgfsetbuttcap%
\pgfsetroundjoin%
\definecolor{currentfill}{rgb}{0.763998,0.763998,0.763998}%
\pgfsetfillcolor{currentfill}%
\pgfsetfillopacity{0.900000}%
\pgfsetlinewidth{0.501875pt}%
\definecolor{currentstroke}{rgb}{0.200000,0.200000,0.200000}%
\pgfsetstrokecolor{currentstroke}%
\pgfsetdash{}{0pt}%
\pgfpathmoveto{\pgfqpoint{2.858919in}{1.987093in}}%
\pgfpathlineto{\pgfqpoint{2.905002in}{1.943513in}}%
\pgfpathlineto{\pgfqpoint{2.938301in}{1.922165in}}%
\pgfpathlineto{\pgfqpoint{2.892217in}{1.965692in}}%
\pgfpathlineto{\pgfqpoint{2.858919in}{1.987093in}}%
\pgfpathclose%
\pgfusepath{stroke,fill}%
\end{pgfscope}%
\begin{pgfscope}%
\pgfpathrectangle{\pgfqpoint{0.050000in}{0.050000in}}{\pgfqpoint{4.900000in}{4.900000in}}%
\pgfusepath{clip}%
\pgfsetbuttcap%
\pgfsetroundjoin%
\definecolor{currentfill}{rgb}{0.104698,0.104698,0.104698}%
\pgfsetfillcolor{currentfill}%
\pgfsetfillopacity{0.900000}%
\pgfsetlinewidth{0.501875pt}%
\definecolor{currentstroke}{rgb}{0.200000,0.200000,0.200000}%
\pgfsetstrokecolor{currentstroke}%
\pgfsetdash{}{0pt}%
\pgfpathmoveto{\pgfqpoint{1.117205in}{3.394075in}}%
\pgfpathlineto{\pgfqpoint{1.167061in}{3.346429in}}%
\pgfpathlineto{\pgfqpoint{1.197601in}{3.327089in}}%
\pgfpathlineto{\pgfqpoint{1.147732in}{3.374768in}}%
\pgfpathlineto{\pgfqpoint{1.117205in}{3.394075in}}%
\pgfpathclose%
\pgfusepath{stroke,fill}%
\end{pgfscope}%
\begin{pgfscope}%
\pgfpathrectangle{\pgfqpoint{0.050000in}{0.050000in}}{\pgfqpoint{4.900000in}{4.900000in}}%
\pgfusepath{clip}%
\pgfsetbuttcap%
\pgfsetroundjoin%
\definecolor{currentfill}{rgb}{0.399723,0.399723,0.399723}%
\pgfsetfillcolor{currentfill}%
\pgfsetfillopacity{0.900000}%
\pgfsetlinewidth{0.501875pt}%
\definecolor{currentstroke}{rgb}{0.200000,0.200000,0.200000}%
\pgfsetstrokecolor{currentstroke}%
\pgfsetdash{}{0pt}%
\pgfpathmoveto{\pgfqpoint{1.852236in}{2.800012in}}%
\pgfpathlineto{\pgfqpoint{1.900500in}{2.753651in}}%
\pgfpathlineto{\pgfqpoint{1.932206in}{2.733371in}}%
\pgfpathlineto{\pgfqpoint{1.883934in}{2.779761in}}%
\pgfpathlineto{\pgfqpoint{1.852236in}{2.800012in}}%
\pgfpathclose%
\pgfusepath{stroke,fill}%
\end{pgfscope}%
\begin{pgfscope}%
\pgfpathrectangle{\pgfqpoint{0.050000in}{0.050000in}}{\pgfqpoint{4.900000in}{4.900000in}}%
\pgfusepath{clip}%
\pgfsetbuttcap%
\pgfsetroundjoin%
\definecolor{currentfill}{rgb}{0.667405,0.667405,0.667405}%
\pgfsetfillcolor{currentfill}%
\pgfsetfillopacity{0.900000}%
\pgfsetlinewidth{0.501875pt}%
\definecolor{currentstroke}{rgb}{0.200000,0.200000,0.200000}%
\pgfsetstrokecolor{currentstroke}%
\pgfsetdash{}{0pt}%
\pgfpathmoveto{\pgfqpoint{2.587368in}{2.205984in}}%
\pgfpathlineto{\pgfqpoint{2.634041in}{2.161025in}}%
\pgfpathlineto{\pgfqpoint{2.666913in}{2.139839in}}%
\pgfpathlineto{\pgfqpoint{2.620238in}{2.184807in}}%
\pgfpathlineto{\pgfqpoint{2.587368in}{2.205984in}}%
\pgfpathclose%
\pgfusepath{stroke,fill}%
\end{pgfscope}%
\begin{pgfscope}%
\pgfpathrectangle{\pgfqpoint{0.050000in}{0.050000in}}{\pgfqpoint{4.900000in}{4.900000in}}%
\pgfusepath{clip}%
\pgfsetbuttcap%
\pgfsetroundjoin%
\definecolor{currentfill}{rgb}{0.000000,0.000000,0.000000}%
\pgfsetfillcolor{currentfill}%
\pgfsetfillopacity{0.900000}%
\pgfsetlinewidth{0.501875pt}%
\definecolor{currentstroke}{rgb}{0.200000,0.200000,0.200000}%
\pgfsetstrokecolor{currentstroke}%
\pgfsetdash{}{0pt}%
\pgfpathmoveto{\pgfqpoint{0.845350in}{3.614466in}}%
\pgfpathlineto{\pgfqpoint{0.895804in}{3.566334in}}%
\pgfpathlineto{\pgfqpoint{0.925919in}{3.547339in}}%
\pgfpathlineto{\pgfqpoint{0.875450in}{3.595507in}}%
\pgfpathlineto{\pgfqpoint{0.845350in}{3.614466in}}%
\pgfpathclose%
\pgfusepath{stroke,fill}%
\end{pgfscope}%
\begin{pgfscope}%
\pgfpathrectangle{\pgfqpoint{0.050000in}{0.050000in}}{\pgfqpoint{4.900000in}{4.900000in}}%
\pgfusepath{clip}%
\pgfsetbuttcap%
\pgfsetroundjoin%
\definecolor{currentfill}{rgb}{0.306344,0.306344,0.306344}%
\pgfsetfillcolor{currentfill}%
\pgfsetfillopacity{0.900000}%
\pgfsetlinewidth{0.501875pt}%
\definecolor{currentstroke}{rgb}{0.200000,0.200000,0.200000}%
\pgfsetstrokecolor{currentstroke}%
\pgfsetdash{}{0pt}%
\pgfpathmoveto{\pgfqpoint{1.580489in}{3.020309in}}%
\pgfpathlineto{\pgfqpoint{1.629352in}{2.973463in}}%
\pgfpathlineto{\pgfqpoint{1.660633in}{2.953527in}}%
\pgfpathlineto{\pgfqpoint{1.611761in}{3.000404in}}%
\pgfpathlineto{\pgfqpoint{1.580489in}{3.020309in}}%
\pgfpathclose%
\pgfusepath{stroke,fill}%
\end{pgfscope}%
\begin{pgfscope}%
\pgfpathrectangle{\pgfqpoint{0.050000in}{0.050000in}}{\pgfqpoint{4.900000in}{4.900000in}}%
\pgfusepath{clip}%
\pgfsetbuttcap%
\pgfsetroundjoin%
\definecolor{currentfill}{rgb}{0.895548,0.895548,0.895548}%
\pgfsetfillcolor{currentfill}%
\pgfsetfillopacity{0.900000}%
\pgfsetlinewidth{0.501875pt}%
\definecolor{currentstroke}{rgb}{0.200000,0.200000,0.200000}%
\pgfsetstrokecolor{currentstroke}%
\pgfsetdash{}{0pt}%
\pgfpathmoveto{\pgfqpoint{3.402179in}{1.576370in}}%
\pgfpathlineto{\pgfqpoint{3.447261in}{1.543817in}}%
\pgfpathlineto{\pgfqpoint{3.481452in}{1.522856in}}%
\pgfpathlineto{\pgfqpoint{3.436362in}{1.555228in}}%
\pgfpathlineto{\pgfqpoint{3.402179in}{1.576370in}}%
\pgfpathclose%
\pgfusepath{stroke,fill}%
\end{pgfscope}%
\begin{pgfscope}%
\pgfpathrectangle{\pgfqpoint{0.050000in}{0.050000in}}{\pgfqpoint{4.900000in}{4.900000in}}%
\pgfusepath{clip}%
\pgfsetbuttcap%
\pgfsetroundjoin%
\definecolor{currentfill}{rgb}{0.560246,0.560246,0.560246}%
\pgfsetfillcolor{currentfill}%
\pgfsetfillopacity{0.900000}%
\pgfsetlinewidth{0.501875pt}%
\definecolor{currentstroke}{rgb}{0.200000,0.200000,0.200000}%
\pgfsetstrokecolor{currentstroke}%
\pgfsetdash{}{0pt}%
\pgfpathmoveto{\pgfqpoint{2.315730in}{2.426091in}}%
\pgfpathlineto{\pgfqpoint{2.363000in}{2.380536in}}%
\pgfpathlineto{\pgfqpoint{2.395447in}{2.359661in}}%
\pgfpathlineto{\pgfqpoint{2.348173in}{2.405243in}}%
\pgfpathlineto{\pgfqpoint{2.315730in}{2.426091in}}%
\pgfpathclose%
\pgfusepath{stroke,fill}%
\end{pgfscope}%
\begin{pgfscope}%
\pgfpathrectangle{\pgfqpoint{0.050000in}{0.050000in}}{\pgfqpoint{4.900000in}{4.900000in}}%
\pgfusepath{clip}%
\pgfsetbuttcap%
\pgfsetroundjoin%
\definecolor{currentfill}{rgb}{0.179008,0.179008,0.179008}%
\pgfsetfillcolor{currentfill}%
\pgfsetfillopacity{0.900000}%
\pgfsetlinewidth{0.501875pt}%
\definecolor{currentstroke}{rgb}{0.200000,0.200000,0.200000}%
\pgfsetstrokecolor{currentstroke}%
\pgfsetdash{}{0pt}%
\pgfpathmoveto{\pgfqpoint{1.308649in}{3.240685in}}%
\pgfpathlineto{\pgfqpoint{1.358110in}{3.193352in}}%
\pgfpathlineto{\pgfqpoint{1.388966in}{3.173762in}}%
\pgfpathlineto{\pgfqpoint{1.339493in}{3.221127in}}%
\pgfpathlineto{\pgfqpoint{1.308649in}{3.240685in}}%
\pgfpathclose%
\pgfusepath{stroke,fill}%
\end{pgfscope}%
\begin{pgfscope}%
\pgfpathrectangle{\pgfqpoint{0.050000in}{0.050000in}}{\pgfqpoint{4.900000in}{4.900000in}}%
\pgfusepath{clip}%
\pgfsetbuttcap%
\pgfsetroundjoin%
\definecolor{currentfill}{rgb}{0.461207,0.461207,0.461207}%
\pgfsetfillcolor{currentfill}%
\pgfsetfillopacity{0.900000}%
\pgfsetlinewidth{0.501875pt}%
\definecolor{currentstroke}{rgb}{0.200000,0.200000,0.200000}%
\pgfsetstrokecolor{currentstroke}%
\pgfsetdash{}{0pt}%
\pgfpathmoveto{\pgfqpoint{2.043998in}{2.646370in}}%
\pgfpathlineto{\pgfqpoint{2.091866in}{2.600325in}}%
\pgfpathlineto{\pgfqpoint{2.123889in}{2.579793in}}%
\pgfpathlineto{\pgfqpoint{2.076014in}{2.625867in}}%
\pgfpathlineto{\pgfqpoint{2.043998in}{2.646370in}}%
\pgfpathclose%
\pgfusepath{stroke,fill}%
\end{pgfscope}%
\begin{pgfscope}%
\pgfpathrectangle{\pgfqpoint{0.050000in}{0.050000in}}{\pgfqpoint{4.900000in}{4.900000in}}%
\pgfusepath{clip}%
\pgfsetbuttcap%
\pgfsetroundjoin%
\definecolor{currentfill}{rgb}{0.815671,0.815671,0.815671}%
\pgfsetfillcolor{currentfill}%
\pgfsetfillopacity{0.900000}%
\pgfsetlinewidth{0.501875pt}%
\definecolor{currentstroke}{rgb}{0.200000,0.200000,0.200000}%
\pgfsetstrokecolor{currentstroke}%
\pgfsetdash{}{0pt}%
\pgfpathmoveto{\pgfqpoint{3.051114in}{1.836751in}}%
\pgfpathlineto{\pgfqpoint{3.096823in}{1.795338in}}%
\pgfpathlineto{\pgfqpoint{3.130447in}{1.774014in}}%
\pgfpathlineto{\pgfqpoint{3.084736in}{1.815312in}}%
\pgfpathlineto{\pgfqpoint{3.051114in}{1.836751in}}%
\pgfpathclose%
\pgfusepath{stroke,fill}%
\end{pgfscope}%
\begin{pgfscope}%
\pgfpathrectangle{\pgfqpoint{0.050000in}{0.050000in}}{\pgfqpoint{4.900000in}{4.900000in}}%
\pgfusepath{clip}%
\pgfsetbuttcap%
\pgfsetroundjoin%
\definecolor{currentfill}{rgb}{0.739377,0.739377,0.739377}%
\pgfsetfillcolor{currentfill}%
\pgfsetfillopacity{0.900000}%
\pgfsetlinewidth{0.501875pt}%
\definecolor{currentstroke}{rgb}{0.200000,0.200000,0.200000}%
\pgfsetstrokecolor{currentstroke}%
\pgfsetdash{}{0pt}%
\pgfpathmoveto{\pgfqpoint{2.779450in}{2.052637in}}%
\pgfpathlineto{\pgfqpoint{2.825728in}{2.008434in}}%
\pgfpathlineto{\pgfqpoint{2.858919in}{1.987093in}}%
\pgfpathlineto{\pgfqpoint{2.812639in}{2.031269in}}%
\pgfpathlineto{\pgfqpoint{2.779450in}{2.052637in}}%
\pgfpathclose%
\pgfusepath{stroke,fill}%
\end{pgfscope}%
\begin{pgfscope}%
\pgfpathrectangle{\pgfqpoint{0.050000in}{0.050000in}}{\pgfqpoint{4.900000in}{4.900000in}}%
\pgfusepath{clip}%
\pgfsetbuttcap%
\pgfsetroundjoin%
\definecolor{currentfill}{rgb}{0.068281,0.068281,0.068281}%
\pgfsetfillcolor{currentfill}%
\pgfsetfillopacity{0.900000}%
\pgfsetlinewidth{0.501875pt}%
\definecolor{currentstroke}{rgb}{0.200000,0.200000,0.200000}%
\pgfsetstrokecolor{currentstroke}%
\pgfsetdash{}{0pt}%
\pgfpathmoveto{\pgfqpoint{1.036714in}{3.461140in}}%
\pgfpathlineto{\pgfqpoint{1.086774in}{3.413320in}}%
\pgfpathlineto{\pgfqpoint{1.117205in}{3.394075in}}%
\pgfpathlineto{\pgfqpoint{1.067132in}{3.441928in}}%
\pgfpathlineto{\pgfqpoint{1.036714in}{3.461140in}}%
\pgfpathclose%
\pgfusepath{stroke,fill}%
\end{pgfscope}%
\begin{pgfscope}%
\pgfpathrectangle{\pgfqpoint{0.050000in}{0.050000in}}{\pgfqpoint{4.900000in}{4.900000in}}%
\pgfusepath{clip}%
\pgfsetbuttcap%
\pgfsetroundjoin%
\definecolor{currentfill}{rgb}{0.371303,0.371303,0.371303}%
\pgfsetfillcolor{currentfill}%
\pgfsetfillopacity{0.900000}%
\pgfsetlinewidth{0.501875pt}%
\definecolor{currentstroke}{rgb}{0.200000,0.200000,0.200000}%
\pgfsetstrokecolor{currentstroke}%
\pgfsetdash{}{0pt}%
\pgfpathmoveto{\pgfqpoint{1.772172in}{2.866731in}}%
\pgfpathlineto{\pgfqpoint{1.820638in}{2.820198in}}%
\pgfpathlineto{\pgfqpoint{1.852236in}{2.800012in}}%
\pgfpathlineto{\pgfqpoint{1.803761in}{2.846574in}}%
\pgfpathlineto{\pgfqpoint{1.772172in}{2.866731in}}%
\pgfpathclose%
\pgfusepath{stroke,fill}%
\end{pgfscope}%
\begin{pgfscope}%
\pgfpathrectangle{\pgfqpoint{0.050000in}{0.050000in}}{\pgfqpoint{4.900000in}{4.900000in}}%
\pgfusepath{clip}%
\pgfsetbuttcap%
\pgfsetroundjoin%
\definecolor{currentfill}{rgb}{0.878570,0.878570,0.878570}%
\pgfsetfillcolor{currentfill}%
\pgfsetfillopacity{0.900000}%
\pgfsetlinewidth{0.501875pt}%
\definecolor{currentstroke}{rgb}{0.200000,0.200000,0.200000}%
\pgfsetstrokecolor{currentstroke}%
\pgfsetdash{}{0pt}%
\pgfpathmoveto{\pgfqpoint{3.322887in}{1.632754in}}%
\pgfpathlineto{\pgfqpoint{3.368106in}{1.597448in}}%
\pgfpathlineto{\pgfqpoint{3.402179in}{1.576370in}}%
\pgfpathlineto{\pgfqpoint{3.356954in}{1.611499in}}%
\pgfpathlineto{\pgfqpoint{3.322887in}{1.632754in}}%
\pgfpathclose%
\pgfusepath{stroke,fill}%
\end{pgfscope}%
\begin{pgfscope}%
\pgfpathrectangle{\pgfqpoint{0.050000in}{0.050000in}}{\pgfqpoint{4.900000in}{4.900000in}}%
\pgfusepath{clip}%
\pgfsetbuttcap%
\pgfsetroundjoin%
\definecolor{currentfill}{rgb}{0.629020,0.629020,0.629020}%
\pgfsetfillcolor{currentfill}%
\pgfsetfillopacity{0.900000}%
\pgfsetlinewidth{0.501875pt}%
\definecolor{currentstroke}{rgb}{0.200000,0.200000,0.200000}%
\pgfsetstrokecolor{currentstroke}%
\pgfsetdash{}{0pt}%
\pgfpathmoveto{\pgfqpoint{2.507730in}{2.272294in}}%
\pgfpathlineto{\pgfqpoint{2.554603in}{2.227097in}}%
\pgfpathlineto{\pgfqpoint{2.587368in}{2.205984in}}%
\pgfpathlineto{\pgfqpoint{2.540492in}{2.251199in}}%
\pgfpathlineto{\pgfqpoint{2.507730in}{2.272294in}}%
\pgfpathclose%
\pgfusepath{stroke,fill}%
\end{pgfscope}%
\begin{pgfscope}%
\pgfpathrectangle{\pgfqpoint{0.050000in}{0.050000in}}{\pgfqpoint{4.900000in}{4.900000in}}%
\pgfusepath{clip}%
\pgfsetbuttcap%
\pgfsetroundjoin%
\definecolor{currentfill}{rgb}{0.262053,0.262053,0.262053}%
\pgfsetfillcolor{currentfill}%
\pgfsetfillopacity{0.900000}%
\pgfsetlinewidth{0.501875pt}%
\definecolor{currentstroke}{rgb}{0.200000,0.200000,0.200000}%
\pgfsetstrokecolor{currentstroke}%
\pgfsetdash{}{0pt}%
\pgfpathmoveto{\pgfqpoint{1.500252in}{3.087170in}}%
\pgfpathlineto{\pgfqpoint{1.549317in}{3.040151in}}%
\pgfpathlineto{\pgfqpoint{1.580489in}{3.020309in}}%
\pgfpathlineto{\pgfqpoint{1.531414in}{3.067360in}}%
\pgfpathlineto{\pgfqpoint{1.500252in}{3.087170in}}%
\pgfpathclose%
\pgfusepath{stroke,fill}%
\end{pgfscope}%
\begin{pgfscope}%
\pgfpathrectangle{\pgfqpoint{0.050000in}{0.050000in}}{\pgfqpoint{4.900000in}{4.900000in}}%
\pgfusepath{clip}%
\pgfsetbuttcap%
\pgfsetroundjoin%
\definecolor{currentfill}{rgb}{0.530104,0.530104,0.530104}%
\pgfsetfillcolor{currentfill}%
\pgfsetfillopacity{0.900000}%
\pgfsetlinewidth{0.501875pt}%
\definecolor{currentstroke}{rgb}{0.200000,0.200000,0.200000}%
\pgfsetstrokecolor{currentstroke}%
\pgfsetdash{}{0pt}%
\pgfpathmoveto{\pgfqpoint{2.235918in}{2.492604in}}%
\pgfpathlineto{\pgfqpoint{2.283389in}{2.446874in}}%
\pgfpathlineto{\pgfqpoint{2.315730in}{2.426091in}}%
\pgfpathlineto{\pgfqpoint{2.268253in}{2.471848in}}%
\pgfpathlineto{\pgfqpoint{2.235918in}{2.492604in}}%
\pgfpathclose%
\pgfusepath{stroke,fill}%
\end{pgfscope}%
\begin{pgfscope}%
\pgfpathrectangle{\pgfqpoint{0.050000in}{0.050000in}}{\pgfqpoint{4.900000in}{4.900000in}}%
\pgfusepath{clip}%
\pgfsetbuttcap%
\pgfsetroundjoin%
\definecolor{currentfill}{rgb}{0.136563,0.136563,0.136563}%
\pgfsetfillcolor{currentfill}%
\pgfsetfillopacity{0.900000}%
\pgfsetlinewidth{0.501875pt}%
\definecolor{currentstroke}{rgb}{0.200000,0.200000,0.200000}%
\pgfsetstrokecolor{currentstroke}%
\pgfsetdash{}{0pt}%
\pgfpathmoveto{\pgfqpoint{1.228238in}{3.307687in}}%
\pgfpathlineto{\pgfqpoint{1.277902in}{3.260181in}}%
\pgfpathlineto{\pgfqpoint{1.308649in}{3.240685in}}%
\pgfpathlineto{\pgfqpoint{1.258972in}{3.288224in}}%
\pgfpathlineto{\pgfqpoint{1.228238in}{3.307687in}}%
\pgfpathclose%
\pgfusepath{stroke,fill}%
\end{pgfscope}%
\begin{pgfscope}%
\pgfpathrectangle{\pgfqpoint{0.050000in}{0.050000in}}{\pgfqpoint{4.900000in}{4.900000in}}%
\pgfusepath{clip}%
\pgfsetbuttcap%
\pgfsetroundjoin%
\definecolor{currentfill}{rgb}{0.791557,0.791557,0.791557}%
\pgfsetfillcolor{currentfill}%
\pgfsetfillopacity{0.900000}%
\pgfsetlinewidth{0.501875pt}%
\definecolor{currentstroke}{rgb}{0.200000,0.200000,0.200000}%
\pgfsetstrokecolor{currentstroke}%
\pgfsetdash{}{0pt}%
\pgfpathmoveto{\pgfqpoint{2.971706in}{1.900758in}}%
\pgfpathlineto{\pgfqpoint{3.017600in}{1.858132in}}%
\pgfpathlineto{\pgfqpoint{3.051114in}{1.836751in}}%
\pgfpathlineto{\pgfqpoint{3.005219in}{1.879294in}}%
\pgfpathlineto{\pgfqpoint{2.971706in}{1.900758in}}%
\pgfpathclose%
\pgfusepath{stroke,fill}%
\end{pgfscope}%
\begin{pgfscope}%
\pgfpathrectangle{\pgfqpoint{0.050000in}{0.050000in}}{\pgfqpoint{4.900000in}{4.900000in}}%
\pgfusepath{clip}%
\pgfsetbuttcap%
\pgfsetroundjoin%
\definecolor{currentfill}{rgb}{0.432203,0.432203,0.432203}%
\pgfsetfillcolor{currentfill}%
\pgfsetfillopacity{0.900000}%
\pgfsetlinewidth{0.501875pt}%
\definecolor{currentstroke}{rgb}{0.200000,0.200000,0.200000}%
\pgfsetstrokecolor{currentstroke}%
\pgfsetdash{}{0pt}%
\pgfpathmoveto{\pgfqpoint{1.964013in}{2.713026in}}%
\pgfpathlineto{\pgfqpoint{2.012083in}{2.666809in}}%
\pgfpathlineto{\pgfqpoint{2.043998in}{2.646370in}}%
\pgfpathlineto{\pgfqpoint{1.995921in}{2.692617in}}%
\pgfpathlineto{\pgfqpoint{1.964013in}{2.713026in}}%
\pgfpathclose%
\pgfusepath{stroke,fill}%
\end{pgfscope}%
\begin{pgfscope}%
\pgfpathrectangle{\pgfqpoint{0.050000in}{0.050000in}}{\pgfqpoint{4.900000in}{4.900000in}}%
\pgfusepath{clip}%
\pgfsetbuttcap%
\pgfsetroundjoin%
\definecolor{currentfill}{rgb}{0.858762,0.858762,0.858762}%
\pgfsetfillcolor{currentfill}%
\pgfsetfillopacity{0.900000}%
\pgfsetlinewidth{0.501875pt}%
\definecolor{currentstroke}{rgb}{0.200000,0.200000,0.200000}%
\pgfsetstrokecolor{currentstroke}%
\pgfsetdash{}{0pt}%
\pgfpathmoveto{\pgfqpoint{3.243558in}{1.691642in}}%
\pgfpathlineto{\pgfqpoint{3.288929in}{1.653948in}}%
\pgfpathlineto{\pgfqpoint{3.322887in}{1.632754in}}%
\pgfpathlineto{\pgfqpoint{3.277511in}{1.670285in}}%
\pgfpathlineto{\pgfqpoint{3.243558in}{1.691642in}}%
\pgfpathclose%
\pgfusepath{stroke,fill}%
\end{pgfscope}%
\begin{pgfscope}%
\pgfpathrectangle{\pgfqpoint{0.050000in}{0.050000in}}{\pgfqpoint{4.900000in}{4.900000in}}%
\pgfusepath{clip}%
\pgfsetbuttcap%
\pgfsetroundjoin%
\definecolor{currentfill}{rgb}{0.700992,0.700992,0.700992}%
\pgfsetfillcolor{currentfill}%
\pgfsetfillopacity{0.900000}%
\pgfsetlinewidth{0.501875pt}%
\definecolor{currentstroke}{rgb}{0.200000,0.200000,0.200000}%
\pgfsetstrokecolor{currentstroke}%
\pgfsetdash{}{0pt}%
\pgfpathmoveto{\pgfqpoint{2.699890in}{2.118590in}}%
\pgfpathlineto{\pgfqpoint{2.746366in}{2.073943in}}%
\pgfpathlineto{\pgfqpoint{2.779450in}{2.052637in}}%
\pgfpathlineto{\pgfqpoint{2.732972in}{2.097277in}}%
\pgfpathlineto{\pgfqpoint{2.699890in}{2.118590in}}%
\pgfpathclose%
\pgfusepath{stroke,fill}%
\end{pgfscope}%
\begin{pgfscope}%
\pgfpathrectangle{\pgfqpoint{0.050000in}{0.050000in}}{\pgfqpoint{4.900000in}{4.900000in}}%
\pgfusepath{clip}%
\pgfsetbuttcap%
\pgfsetroundjoin%
\definecolor{currentfill}{rgb}{0.036417,0.036417,0.036417}%
\pgfsetfillcolor{currentfill}%
\pgfsetfillopacity{0.900000}%
\pgfsetlinewidth{0.501875pt}%
\definecolor{currentstroke}{rgb}{0.200000,0.200000,0.200000}%
\pgfsetstrokecolor{currentstroke}%
\pgfsetdash{}{0pt}%
\pgfpathmoveto{\pgfqpoint{0.956130in}{3.528284in}}%
\pgfpathlineto{\pgfqpoint{1.006394in}{3.480290in}}%
\pgfpathlineto{\pgfqpoint{1.036714in}{3.461140in}}%
\pgfpathlineto{\pgfqpoint{0.986436in}{3.509168in}}%
\pgfpathlineto{\pgfqpoint{0.956130in}{3.528284in}}%
\pgfpathclose%
\pgfusepath{stroke,fill}%
\end{pgfscope}%
\begin{pgfscope}%
\pgfpathrectangle{\pgfqpoint{0.050000in}{0.050000in}}{\pgfqpoint{4.900000in}{4.900000in}}%
\pgfusepath{clip}%
\pgfsetbuttcap%
\pgfsetroundjoin%
\definecolor{currentfill}{rgb}{0.338824,0.338824,0.338824}%
\pgfsetfillcolor{currentfill}%
\pgfsetfillopacity{0.900000}%
\pgfsetlinewidth{0.501875pt}%
\definecolor{currentstroke}{rgb}{0.200000,0.200000,0.200000}%
\pgfsetstrokecolor{currentstroke}%
\pgfsetdash{}{0pt}%
\pgfpathmoveto{\pgfqpoint{1.692013in}{2.933528in}}%
\pgfpathlineto{\pgfqpoint{1.740683in}{2.886824in}}%
\pgfpathlineto{\pgfqpoint{1.772172in}{2.866731in}}%
\pgfpathlineto{\pgfqpoint{1.723494in}{2.913466in}}%
\pgfpathlineto{\pgfqpoint{1.692013in}{2.933528in}}%
\pgfpathclose%
\pgfusepath{stroke,fill}%
\end{pgfscope}%
\begin{pgfscope}%
\pgfpathrectangle{\pgfqpoint{0.050000in}{0.050000in}}{\pgfqpoint{4.900000in}{4.900000in}}%
\pgfusepath{clip}%
\pgfsetbuttcap%
\pgfsetroundjoin%
\definecolor{currentfill}{rgb}{0.926674,0.926674,0.926674}%
\pgfsetfillcolor{currentfill}%
\pgfsetfillopacity{0.900000}%
\pgfsetlinewidth{0.501875pt}%
\definecolor{currentstroke}{rgb}{0.200000,0.200000,0.200000}%
\pgfsetstrokecolor{currentstroke}%
\pgfsetdash{}{0pt}%
\pgfpathmoveto{\pgfqpoint{3.595146in}{1.451599in}}%
\pgfpathlineto{\pgfqpoint{3.640011in}{1.425588in}}%
\pgfpathlineto{\pgfqpoint{3.674556in}{1.404755in}}%
\pgfpathlineto{\pgfqpoint{3.629679in}{1.430606in}}%
\pgfpathlineto{\pgfqpoint{3.595146in}{1.451599in}}%
\pgfpathclose%
\pgfusepath{stroke,fill}%
\end{pgfscope}%
\begin{pgfscope}%
\pgfpathrectangle{\pgfqpoint{0.050000in}{0.050000in}}{\pgfqpoint{4.900000in}{4.900000in}}%
\pgfusepath{clip}%
\pgfsetbuttcap%
\pgfsetroundjoin%
\definecolor{currentfill}{rgb}{0.595433,0.595433,0.595433}%
\pgfsetfillcolor{currentfill}%
\pgfsetfillopacity{0.900000}%
\pgfsetlinewidth{0.501875pt}%
\definecolor{currentstroke}{rgb}{0.200000,0.200000,0.200000}%
\pgfsetstrokecolor{currentstroke}%
\pgfsetdash{}{0pt}%
\pgfpathmoveto{\pgfqpoint{2.427998in}{2.338721in}}%
\pgfpathlineto{\pgfqpoint{2.475072in}{2.293323in}}%
\pgfpathlineto{\pgfqpoint{2.507730in}{2.272294in}}%
\pgfpathlineto{\pgfqpoint{2.460653in}{2.317715in}}%
\pgfpathlineto{\pgfqpoint{2.427998in}{2.338721in}}%
\pgfpathclose%
\pgfusepath{stroke,fill}%
\end{pgfscope}%
\begin{pgfscope}%
\pgfpathrectangle{\pgfqpoint{0.050000in}{0.050000in}}{\pgfqpoint{4.900000in}{4.900000in}}%
\pgfusepath{clip}%
\pgfsetbuttcap%
\pgfsetroundjoin%
\definecolor{currentfill}{rgb}{0.223299,0.223299,0.223299}%
\pgfsetfillcolor{currentfill}%
\pgfsetfillopacity{0.900000}%
\pgfsetlinewidth{0.501875pt}%
\definecolor{currentstroke}{rgb}{0.200000,0.200000,0.200000}%
\pgfsetstrokecolor{currentstroke}%
\pgfsetdash{}{0pt}%
\pgfpathmoveto{\pgfqpoint{1.419920in}{3.154109in}}%
\pgfpathlineto{\pgfqpoint{1.469189in}{3.106917in}}%
\pgfpathlineto{\pgfqpoint{1.500252in}{3.087170in}}%
\pgfpathlineto{\pgfqpoint{1.450972in}{3.134394in}}%
\pgfpathlineto{\pgfqpoint{1.419920in}{3.154109in}}%
\pgfpathclose%
\pgfusepath{stroke,fill}%
\end{pgfscope}%
\begin{pgfscope}%
\pgfpathrectangle{\pgfqpoint{0.050000in}{0.050000in}}{\pgfqpoint{4.900000in}{4.900000in}}%
\pgfusepath{clip}%
\pgfsetbuttcap%
\pgfsetroundjoin%
\definecolor{currentfill}{rgb}{0.495656,0.495656,0.495656}%
\pgfsetfillcolor{currentfill}%
\pgfsetfillopacity{0.900000}%
\pgfsetlinewidth{0.501875pt}%
\definecolor{currentstroke}{rgb}{0.200000,0.200000,0.200000}%
\pgfsetstrokecolor{currentstroke}%
\pgfsetdash{}{0pt}%
\pgfpathmoveto{\pgfqpoint{2.156013in}{2.559196in}}%
\pgfpathlineto{\pgfqpoint{2.203686in}{2.513294in}}%
\pgfpathlineto{\pgfqpoint{2.235918in}{2.492604in}}%
\pgfpathlineto{\pgfqpoint{2.188240in}{2.538534in}}%
\pgfpathlineto{\pgfqpoint{2.156013in}{2.559196in}}%
\pgfpathclose%
\pgfusepath{stroke,fill}%
\end{pgfscope}%
\begin{pgfscope}%
\pgfpathrectangle{\pgfqpoint{0.050000in}{0.050000in}}{\pgfqpoint{4.900000in}{4.900000in}}%
\pgfusepath{clip}%
\pgfsetbuttcap%
\pgfsetroundjoin%
\definecolor{currentfill}{rgb}{0.912526,0.912526,0.912526}%
\pgfsetfillcolor{currentfill}%
\pgfsetfillopacity{0.900000}%
\pgfsetlinewidth{0.501875pt}%
\definecolor{currentstroke}{rgb}{0.200000,0.200000,0.200000}%
\pgfsetstrokecolor{currentstroke}%
\pgfsetdash{}{0pt}%
\pgfpathmoveto{\pgfqpoint{3.515754in}{1.501829in}}%
\pgfpathlineto{\pgfqpoint{3.560724in}{1.472520in}}%
\pgfpathlineto{\pgfqpoint{3.595146in}{1.451599in}}%
\pgfpathlineto{\pgfqpoint{3.550166in}{1.480734in}}%
\pgfpathlineto{\pgfqpoint{3.515754in}{1.501829in}}%
\pgfpathclose%
\pgfusepath{stroke,fill}%
\end{pgfscope}%
\begin{pgfscope}%
\pgfpathrectangle{\pgfqpoint{0.050000in}{0.050000in}}{\pgfqpoint{4.900000in}{4.900000in}}%
\pgfusepath{clip}%
\pgfsetbuttcap%
\pgfsetroundjoin%
\definecolor{currentfill}{rgb}{0.839785,0.839785,0.839785}%
\pgfsetfillcolor{currentfill}%
\pgfsetfillopacity{0.900000}%
\pgfsetlinewidth{0.501875pt}%
\definecolor{currentstroke}{rgb}{0.200000,0.200000,0.200000}%
\pgfsetstrokecolor{currentstroke}%
\pgfsetdash{}{0pt}%
\pgfpathmoveto{\pgfqpoint{3.164179in}{1.752632in}}%
\pgfpathlineto{\pgfqpoint{3.209714in}{1.712940in}}%
\pgfpathlineto{\pgfqpoint{3.243558in}{1.691642in}}%
\pgfpathlineto{\pgfqpoint{3.198019in}{1.731192in}}%
\pgfpathlineto{\pgfqpoint{3.164179in}{1.752632in}}%
\pgfpathclose%
\pgfusepath{stroke,fill}%
\end{pgfscope}%
\begin{pgfscope}%
\pgfpathrectangle{\pgfqpoint{0.050000in}{0.050000in}}{\pgfqpoint{4.900000in}{4.900000in}}%
\pgfusepath{clip}%
\pgfsetbuttcap%
\pgfsetroundjoin%
\definecolor{currentfill}{rgb}{0.104698,0.104698,0.104698}%
\pgfsetfillcolor{currentfill}%
\pgfsetfillopacity{0.900000}%
\pgfsetlinewidth{0.501875pt}%
\definecolor{currentstroke}{rgb}{0.200000,0.200000,0.200000}%
\pgfsetstrokecolor{currentstroke}%
\pgfsetdash{}{0pt}%
\pgfpathmoveto{\pgfqpoint{1.147732in}{3.374768in}}%
\pgfpathlineto{\pgfqpoint{1.197601in}{3.327089in}}%
\pgfpathlineto{\pgfqpoint{1.228238in}{3.307687in}}%
\pgfpathlineto{\pgfqpoint{1.178356in}{3.355400in}}%
\pgfpathlineto{\pgfqpoint{1.147732in}{3.374768in}}%
\pgfpathclose%
\pgfusepath{stroke,fill}%
\end{pgfscope}%
\begin{pgfscope}%
\pgfpathrectangle{\pgfqpoint{0.050000in}{0.050000in}}{\pgfqpoint{4.900000in}{4.900000in}}%
\pgfusepath{clip}%
\pgfsetbuttcap%
\pgfsetroundjoin%
\definecolor{currentfill}{rgb}{0.763998,0.763998,0.763998}%
\pgfsetfillcolor{currentfill}%
\pgfsetfillopacity{0.900000}%
\pgfsetlinewidth{0.501875pt}%
\definecolor{currentstroke}{rgb}{0.200000,0.200000,0.200000}%
\pgfsetstrokecolor{currentstroke}%
\pgfsetdash{}{0pt}%
\pgfpathmoveto{\pgfqpoint{2.892217in}{1.965692in}}%
\pgfpathlineto{\pgfqpoint{2.938301in}{1.922165in}}%
\pgfpathlineto{\pgfqpoint{2.971706in}{1.900758in}}%
\pgfpathlineto{\pgfqpoint{2.925620in}{1.944231in}}%
\pgfpathlineto{\pgfqpoint{2.892217in}{1.965692in}}%
\pgfpathclose%
\pgfusepath{stroke,fill}%
\end{pgfscope}%
\begin{pgfscope}%
\pgfpathrectangle{\pgfqpoint{0.050000in}{0.050000in}}{\pgfqpoint{4.900000in}{4.900000in}}%
\pgfusepath{clip}%
\pgfsetbuttcap%
\pgfsetroundjoin%
\definecolor{currentfill}{rgb}{0.399723,0.399723,0.399723}%
\pgfsetfillcolor{currentfill}%
\pgfsetfillopacity{0.900000}%
\pgfsetlinewidth{0.501875pt}%
\definecolor{currentstroke}{rgb}{0.200000,0.200000,0.200000}%
\pgfsetstrokecolor{currentstroke}%
\pgfsetdash{}{0pt}%
\pgfpathmoveto{\pgfqpoint{1.883934in}{2.779761in}}%
\pgfpathlineto{\pgfqpoint{1.932206in}{2.733371in}}%
\pgfpathlineto{\pgfqpoint{1.964013in}{2.713026in}}%
\pgfpathlineto{\pgfqpoint{1.915733in}{2.759446in}}%
\pgfpathlineto{\pgfqpoint{1.883934in}{2.779761in}}%
\pgfpathclose%
\pgfusepath{stroke,fill}%
\end{pgfscope}%
\begin{pgfscope}%
\pgfpathrectangle{\pgfqpoint{0.050000in}{0.050000in}}{\pgfqpoint{4.900000in}{4.900000in}}%
\pgfusepath{clip}%
\pgfsetbuttcap%
\pgfsetroundjoin%
\definecolor{currentfill}{rgb}{0.667405,0.667405,0.667405}%
\pgfsetfillcolor{currentfill}%
\pgfsetfillopacity{0.900000}%
\pgfsetlinewidth{0.501875pt}%
\definecolor{currentstroke}{rgb}{0.200000,0.200000,0.200000}%
\pgfsetstrokecolor{currentstroke}%
\pgfsetdash{}{0pt}%
\pgfpathmoveto{\pgfqpoint{2.620238in}{2.184807in}}%
\pgfpathlineto{\pgfqpoint{2.666913in}{2.139839in}}%
\pgfpathlineto{\pgfqpoint{2.699890in}{2.118590in}}%
\pgfpathlineto{\pgfqpoint{2.653212in}{2.163565in}}%
\pgfpathlineto{\pgfqpoint{2.620238in}{2.184807in}}%
\pgfpathclose%
\pgfusepath{stroke,fill}%
\end{pgfscope}%
\begin{pgfscope}%
\pgfpathrectangle{\pgfqpoint{0.050000in}{0.050000in}}{\pgfqpoint{4.900000in}{4.900000in}}%
\pgfusepath{clip}%
\pgfsetbuttcap%
\pgfsetroundjoin%
\definecolor{currentfill}{rgb}{0.000000,0.000000,0.000000}%
\pgfsetfillcolor{currentfill}%
\pgfsetfillopacity{0.900000}%
\pgfsetlinewidth{0.501875pt}%
\definecolor{currentstroke}{rgb}{0.200000,0.200000,0.200000}%
\pgfsetstrokecolor{currentstroke}%
\pgfsetdash{}{0pt}%
\pgfpathmoveto{\pgfqpoint{0.875450in}{3.595507in}}%
\pgfpathlineto{\pgfqpoint{0.925919in}{3.547339in}}%
\pgfpathlineto{\pgfqpoint{0.956130in}{3.528284in}}%
\pgfpathlineto{\pgfqpoint{0.905646in}{3.576486in}}%
\pgfpathlineto{\pgfqpoint{0.875450in}{3.595507in}}%
\pgfpathclose%
\pgfusepath{stroke,fill}%
\end{pgfscope}%
\begin{pgfscope}%
\pgfpathrectangle{\pgfqpoint{0.050000in}{0.050000in}}{\pgfqpoint{4.900000in}{4.900000in}}%
\pgfusepath{clip}%
\pgfsetbuttcap%
\pgfsetroundjoin%
\definecolor{currentfill}{rgb}{0.306344,0.306344,0.306344}%
\pgfsetfillcolor{currentfill}%
\pgfsetfillopacity{0.900000}%
\pgfsetlinewidth{0.501875pt}%
\definecolor{currentstroke}{rgb}{0.200000,0.200000,0.200000}%
\pgfsetstrokecolor{currentstroke}%
\pgfsetdash{}{0pt}%
\pgfpathmoveto{\pgfqpoint{1.611761in}{3.000404in}}%
\pgfpathlineto{\pgfqpoint{1.660633in}{2.953527in}}%
\pgfpathlineto{\pgfqpoint{1.692013in}{2.933528in}}%
\pgfpathlineto{\pgfqpoint{1.643132in}{2.980437in}}%
\pgfpathlineto{\pgfqpoint{1.611761in}{3.000404in}}%
\pgfpathclose%
\pgfusepath{stroke,fill}%
\end{pgfscope}%
\begin{pgfscope}%
\pgfpathrectangle{\pgfqpoint{0.050000in}{0.050000in}}{\pgfqpoint{4.900000in}{4.900000in}}%
\pgfusepath{clip}%
\pgfsetbuttcap%
\pgfsetroundjoin%
\definecolor{currentfill}{rgb}{0.895548,0.895548,0.895548}%
\pgfsetfillcolor{currentfill}%
\pgfsetfillopacity{0.900000}%
\pgfsetlinewidth{0.501875pt}%
\definecolor{currentstroke}{rgb}{0.200000,0.200000,0.200000}%
\pgfsetstrokecolor{currentstroke}%
\pgfsetdash{}{0pt}%
\pgfpathmoveto{\pgfqpoint{3.436362in}{1.555228in}}%
\pgfpathlineto{\pgfqpoint{3.481452in}{1.522856in}}%
\pgfpathlineto{\pgfqpoint{3.515754in}{1.501829in}}%
\pgfpathlineto{\pgfqpoint{3.470655in}{1.534022in}}%
\pgfpathlineto{\pgfqpoint{3.436362in}{1.555228in}}%
\pgfpathclose%
\pgfusepath{stroke,fill}%
\end{pgfscope}%
\begin{pgfscope}%
\pgfpathrectangle{\pgfqpoint{0.050000in}{0.050000in}}{\pgfqpoint{4.900000in}{4.900000in}}%
\pgfusepath{clip}%
\pgfsetbuttcap%
\pgfsetroundjoin%
\definecolor{currentfill}{rgb}{0.560246,0.560246,0.560246}%
\pgfsetfillcolor{currentfill}%
\pgfsetfillopacity{0.900000}%
\pgfsetlinewidth{0.501875pt}%
\definecolor{currentstroke}{rgb}{0.200000,0.200000,0.200000}%
\pgfsetstrokecolor{currentstroke}%
\pgfsetdash{}{0pt}%
\pgfpathmoveto{\pgfqpoint{2.348173in}{2.405243in}}%
\pgfpathlineto{\pgfqpoint{2.395447in}{2.359661in}}%
\pgfpathlineto{\pgfqpoint{2.427998in}{2.338721in}}%
\pgfpathlineto{\pgfqpoint{2.380719in}{2.384328in}}%
\pgfpathlineto{\pgfqpoint{2.348173in}{2.405243in}}%
\pgfpathclose%
\pgfusepath{stroke,fill}%
\end{pgfscope}%
\begin{pgfscope}%
\pgfpathrectangle{\pgfqpoint{0.050000in}{0.050000in}}{\pgfqpoint{4.900000in}{4.900000in}}%
\pgfusepath{clip}%
\pgfsetbuttcap%
\pgfsetroundjoin%
\definecolor{currentfill}{rgb}{0.179008,0.179008,0.179008}%
\pgfsetfillcolor{currentfill}%
\pgfsetfillopacity{0.900000}%
\pgfsetlinewidth{0.501875pt}%
\definecolor{currentstroke}{rgb}{0.200000,0.200000,0.200000}%
\pgfsetstrokecolor{currentstroke}%
\pgfsetdash{}{0pt}%
\pgfpathmoveto{\pgfqpoint{1.339493in}{3.221127in}}%
\pgfpathlineto{\pgfqpoint{1.388966in}{3.173762in}}%
\pgfpathlineto{\pgfqpoint{1.419920in}{3.154109in}}%
\pgfpathlineto{\pgfqpoint{1.370436in}{3.201507in}}%
\pgfpathlineto{\pgfqpoint{1.339493in}{3.221127in}}%
\pgfpathclose%
\pgfusepath{stroke,fill}%
\end{pgfscope}%
\begin{pgfscope}%
\pgfpathrectangle{\pgfqpoint{0.050000in}{0.050000in}}{\pgfqpoint{4.900000in}{4.900000in}}%
\pgfusepath{clip}%
\pgfsetbuttcap%
\pgfsetroundjoin%
\definecolor{currentfill}{rgb}{0.465513,0.465513,0.465513}%
\pgfsetfillcolor{currentfill}%
\pgfsetfillopacity{0.900000}%
\pgfsetlinewidth{0.501875pt}%
\definecolor{currentstroke}{rgb}{0.200000,0.200000,0.200000}%
\pgfsetstrokecolor{currentstroke}%
\pgfsetdash{}{0pt}%
\pgfpathmoveto{\pgfqpoint{2.076014in}{2.625867in}}%
\pgfpathlineto{\pgfqpoint{2.123889in}{2.579793in}}%
\pgfpathlineto{\pgfqpoint{2.156013in}{2.559196in}}%
\pgfpathlineto{\pgfqpoint{2.108132in}{2.605299in}}%
\pgfpathlineto{\pgfqpoint{2.076014in}{2.625867in}}%
\pgfpathclose%
\pgfusepath{stroke,fill}%
\end{pgfscope}%
\begin{pgfscope}%
\pgfpathrectangle{\pgfqpoint{0.050000in}{0.050000in}}{\pgfqpoint{4.900000in}{4.900000in}}%
\pgfusepath{clip}%
\pgfsetbuttcap%
\pgfsetroundjoin%
\definecolor{currentfill}{rgb}{0.815671,0.815671,0.815671}%
\pgfsetfillcolor{currentfill}%
\pgfsetfillopacity{0.900000}%
\pgfsetlinewidth{0.501875pt}%
\definecolor{currentstroke}{rgb}{0.200000,0.200000,0.200000}%
\pgfsetstrokecolor{currentstroke}%
\pgfsetdash{}{0pt}%
\pgfpathmoveto{\pgfqpoint{3.084736in}{1.815312in}}%
\pgfpathlineto{\pgfqpoint{3.130447in}{1.774014in}}%
\pgfpathlineto{\pgfqpoint{3.164179in}{1.752632in}}%
\pgfpathlineto{\pgfqpoint{3.118465in}{1.793815in}}%
\pgfpathlineto{\pgfqpoint{3.084736in}{1.815312in}}%
\pgfpathclose%
\pgfusepath{stroke,fill}%
\end{pgfscope}%
\begin{pgfscope}%
\pgfpathrectangle{\pgfqpoint{0.050000in}{0.050000in}}{\pgfqpoint{4.900000in}{4.900000in}}%
\pgfusepath{clip}%
\pgfsetbuttcap%
\pgfsetroundjoin%
\definecolor{currentfill}{rgb}{0.739377,0.739377,0.739377}%
\pgfsetfillcolor{currentfill}%
\pgfsetfillopacity{0.900000}%
\pgfsetlinewidth{0.501875pt}%
\definecolor{currentstroke}{rgb}{0.200000,0.200000,0.200000}%
\pgfsetstrokecolor{currentstroke}%
\pgfsetdash{}{0pt}%
\pgfpathmoveto{\pgfqpoint{2.812639in}{2.031269in}}%
\pgfpathlineto{\pgfqpoint{2.858919in}{1.987093in}}%
\pgfpathlineto{\pgfqpoint{2.892217in}{1.965692in}}%
\pgfpathlineto{\pgfqpoint{2.845935in}{2.009839in}}%
\pgfpathlineto{\pgfqpoint{2.812639in}{2.031269in}}%
\pgfpathclose%
\pgfusepath{stroke,fill}%
\end{pgfscope}%
\begin{pgfscope}%
\pgfpathrectangle{\pgfqpoint{0.050000in}{0.050000in}}{\pgfqpoint{4.900000in}{4.900000in}}%
\pgfusepath{clip}%
\pgfsetbuttcap%
\pgfsetroundjoin%
\definecolor{currentfill}{rgb}{0.068281,0.068281,0.068281}%
\pgfsetfillcolor{currentfill}%
\pgfsetfillopacity{0.900000}%
\pgfsetlinewidth{0.501875pt}%
\definecolor{currentstroke}{rgb}{0.200000,0.200000,0.200000}%
\pgfsetstrokecolor{currentstroke}%
\pgfsetdash{}{0pt}%
\pgfpathmoveto{\pgfqpoint{1.067132in}{3.441928in}}%
\pgfpathlineto{\pgfqpoint{1.117205in}{3.394075in}}%
\pgfpathlineto{\pgfqpoint{1.147732in}{3.374768in}}%
\pgfpathlineto{\pgfqpoint{1.097645in}{3.422656in}}%
\pgfpathlineto{\pgfqpoint{1.067132in}{3.441928in}}%
\pgfpathclose%
\pgfusepath{stroke,fill}%
\end{pgfscope}%
\begin{pgfscope}%
\pgfpathrectangle{\pgfqpoint{0.050000in}{0.050000in}}{\pgfqpoint{4.900000in}{4.900000in}}%
\pgfusepath{clip}%
\pgfsetbuttcap%
\pgfsetroundjoin%
\definecolor{currentfill}{rgb}{0.371303,0.371303,0.371303}%
\pgfsetfillcolor{currentfill}%
\pgfsetfillopacity{0.900000}%
\pgfsetlinewidth{0.501875pt}%
\definecolor{currentstroke}{rgb}{0.200000,0.200000,0.200000}%
\pgfsetstrokecolor{currentstroke}%
\pgfsetdash{}{0pt}%
\pgfpathmoveto{\pgfqpoint{1.803761in}{2.846574in}}%
\pgfpathlineto{\pgfqpoint{1.852236in}{2.800012in}}%
\pgfpathlineto{\pgfqpoint{1.883934in}{2.779761in}}%
\pgfpathlineto{\pgfqpoint{1.835451in}{2.826353in}}%
\pgfpathlineto{\pgfqpoint{1.803761in}{2.846574in}}%
\pgfpathclose%
\pgfusepath{stroke,fill}%
\end{pgfscope}%
\begin{pgfscope}%
\pgfpathrectangle{\pgfqpoint{0.050000in}{0.050000in}}{\pgfqpoint{4.900000in}{4.900000in}}%
\pgfusepath{clip}%
\pgfsetbuttcap%
\pgfsetroundjoin%
\definecolor{currentfill}{rgb}{0.878570,0.878570,0.878570}%
\pgfsetfillcolor{currentfill}%
\pgfsetfillopacity{0.900000}%
\pgfsetlinewidth{0.501875pt}%
\definecolor{currentstroke}{rgb}{0.200000,0.200000,0.200000}%
\pgfsetstrokecolor{currentstroke}%
\pgfsetdash{}{0pt}%
\pgfpathmoveto{\pgfqpoint{3.356954in}{1.611499in}}%
\pgfpathlineto{\pgfqpoint{3.402179in}{1.576370in}}%
\pgfpathlineto{\pgfqpoint{3.436362in}{1.555228in}}%
\pgfpathlineto{\pgfqpoint{3.391130in}{1.590183in}}%
\pgfpathlineto{\pgfqpoint{3.356954in}{1.611499in}}%
\pgfpathclose%
\pgfusepath{stroke,fill}%
\end{pgfscope}%
\begin{pgfscope}%
\pgfpathrectangle{\pgfqpoint{0.050000in}{0.050000in}}{\pgfqpoint{4.900000in}{4.900000in}}%
\pgfusepath{clip}%
\pgfsetbuttcap%
\pgfsetroundjoin%
\definecolor{currentfill}{rgb}{0.629020,0.629020,0.629020}%
\pgfsetfillcolor{currentfill}%
\pgfsetfillopacity{0.900000}%
\pgfsetlinewidth{0.501875pt}%
\definecolor{currentstroke}{rgb}{0.200000,0.200000,0.200000}%
\pgfsetstrokecolor{currentstroke}%
\pgfsetdash{}{0pt}%
\pgfpathmoveto{\pgfqpoint{2.540492in}{2.251199in}}%
\pgfpathlineto{\pgfqpoint{2.587368in}{2.205984in}}%
\pgfpathlineto{\pgfqpoint{2.620238in}{2.184807in}}%
\pgfpathlineto{\pgfqpoint{2.573359in}{2.230039in}}%
\pgfpathlineto{\pgfqpoint{2.540492in}{2.251199in}}%
\pgfpathclose%
\pgfusepath{stroke,fill}%
\end{pgfscope}%
\begin{pgfscope}%
\pgfpathrectangle{\pgfqpoint{0.050000in}{0.050000in}}{\pgfqpoint{4.900000in}{4.900000in}}%
\pgfusepath{clip}%
\pgfsetbuttcap%
\pgfsetroundjoin%
\definecolor{currentfill}{rgb}{0.262053,0.262053,0.262053}%
\pgfsetfillcolor{currentfill}%
\pgfsetfillopacity{0.900000}%
\pgfsetlinewidth{0.501875pt}%
\definecolor{currentstroke}{rgb}{0.200000,0.200000,0.200000}%
\pgfsetstrokecolor{currentstroke}%
\pgfsetdash{}{0pt}%
\pgfpathmoveto{\pgfqpoint{1.531414in}{3.067360in}}%
\pgfpathlineto{\pgfqpoint{1.580489in}{3.020309in}}%
\pgfpathlineto{\pgfqpoint{1.611761in}{3.000404in}}%
\pgfpathlineto{\pgfqpoint{1.562675in}{3.047487in}}%
\pgfpathlineto{\pgfqpoint{1.531414in}{3.067360in}}%
\pgfpathclose%
\pgfusepath{stroke,fill}%
\end{pgfscope}%
\begin{pgfscope}%
\pgfpathrectangle{\pgfqpoint{0.050000in}{0.050000in}}{\pgfqpoint{4.900000in}{4.900000in}}%
\pgfusepath{clip}%
\pgfsetbuttcap%
\pgfsetroundjoin%
\definecolor{currentfill}{rgb}{0.530104,0.530104,0.530104}%
\pgfsetfillcolor{currentfill}%
\pgfsetfillopacity{0.900000}%
\pgfsetlinewidth{0.501875pt}%
\definecolor{currentstroke}{rgb}{0.200000,0.200000,0.200000}%
\pgfsetstrokecolor{currentstroke}%
\pgfsetdash{}{0pt}%
\pgfpathmoveto{\pgfqpoint{2.268253in}{2.471848in}}%
\pgfpathlineto{\pgfqpoint{2.315730in}{2.426091in}}%
\pgfpathlineto{\pgfqpoint{2.348173in}{2.405243in}}%
\pgfpathlineto{\pgfqpoint{2.300692in}{2.451026in}}%
\pgfpathlineto{\pgfqpoint{2.268253in}{2.471848in}}%
\pgfpathclose%
\pgfusepath{stroke,fill}%
\end{pgfscope}%
\begin{pgfscope}%
\pgfpathrectangle{\pgfqpoint{0.050000in}{0.050000in}}{\pgfqpoint{4.900000in}{4.900000in}}%
\pgfusepath{clip}%
\pgfsetbuttcap%
\pgfsetroundjoin%
\definecolor{currentfill}{rgb}{0.141115,0.141115,0.141115}%
\pgfsetfillcolor{currentfill}%
\pgfsetfillopacity{0.900000}%
\pgfsetlinewidth{0.501875pt}%
\definecolor{currentstroke}{rgb}{0.200000,0.200000,0.200000}%
\pgfsetstrokecolor{currentstroke}%
\pgfsetdash{}{0pt}%
\pgfpathmoveto{\pgfqpoint{1.258972in}{3.288224in}}%
\pgfpathlineto{\pgfqpoint{1.308649in}{3.240685in}}%
\pgfpathlineto{\pgfqpoint{1.339493in}{3.221127in}}%
\pgfpathlineto{\pgfqpoint{1.289805in}{3.268699in}}%
\pgfpathlineto{\pgfqpoint{1.258972in}{3.288224in}}%
\pgfpathclose%
\pgfusepath{stroke,fill}%
\end{pgfscope}%
\begin{pgfscope}%
\pgfpathrectangle{\pgfqpoint{0.050000in}{0.050000in}}{\pgfqpoint{4.900000in}{4.900000in}}%
\pgfusepath{clip}%
\pgfsetbuttcap%
\pgfsetroundjoin%
\definecolor{currentfill}{rgb}{0.791557,0.791557,0.791557}%
\pgfsetfillcolor{currentfill}%
\pgfsetfillopacity{0.900000}%
\pgfsetlinewidth{0.501875pt}%
\definecolor{currentstroke}{rgb}{0.200000,0.200000,0.200000}%
\pgfsetstrokecolor{currentstroke}%
\pgfsetdash{}{0pt}%
\pgfpathmoveto{\pgfqpoint{3.005219in}{1.879294in}}%
\pgfpathlineto{\pgfqpoint{3.051114in}{1.836751in}}%
\pgfpathlineto{\pgfqpoint{3.084736in}{1.815312in}}%
\pgfpathlineto{\pgfqpoint{3.038838in}{1.857770in}}%
\pgfpathlineto{\pgfqpoint{3.005219in}{1.879294in}}%
\pgfpathclose%
\pgfusepath{stroke,fill}%
\end{pgfscope}%
\begin{pgfscope}%
\pgfpathrectangle{\pgfqpoint{0.050000in}{0.050000in}}{\pgfqpoint{4.900000in}{4.900000in}}%
\pgfusepath{clip}%
\pgfsetbuttcap%
\pgfsetroundjoin%
\definecolor{currentfill}{rgb}{0.432203,0.432203,0.432203}%
\pgfsetfillcolor{currentfill}%
\pgfsetfillopacity{0.900000}%
\pgfsetlinewidth{0.501875pt}%
\definecolor{currentstroke}{rgb}{0.200000,0.200000,0.200000}%
\pgfsetstrokecolor{currentstroke}%
\pgfsetdash{}{0pt}%
\pgfpathmoveto{\pgfqpoint{1.995921in}{2.692617in}}%
\pgfpathlineto{\pgfqpoint{2.043998in}{2.646370in}}%
\pgfpathlineto{\pgfqpoint{2.076014in}{2.625867in}}%
\pgfpathlineto{\pgfqpoint{2.027930in}{2.672143in}}%
\pgfpathlineto{\pgfqpoint{1.995921in}{2.692617in}}%
\pgfpathclose%
\pgfusepath{stroke,fill}%
\end{pgfscope}%
\begin{pgfscope}%
\pgfpathrectangle{\pgfqpoint{0.050000in}{0.050000in}}{\pgfqpoint{4.900000in}{4.900000in}}%
\pgfusepath{clip}%
\pgfsetbuttcap%
\pgfsetroundjoin%
\definecolor{currentfill}{rgb}{0.858762,0.858762,0.858762}%
\pgfsetfillcolor{currentfill}%
\pgfsetfillopacity{0.900000}%
\pgfsetlinewidth{0.501875pt}%
\definecolor{currentstroke}{rgb}{0.200000,0.200000,0.200000}%
\pgfsetstrokecolor{currentstroke}%
\pgfsetdash{}{0pt}%
\pgfpathmoveto{\pgfqpoint{3.277511in}{1.670285in}}%
\pgfpathlineto{\pgfqpoint{3.322887in}{1.632754in}}%
\pgfpathlineto{\pgfqpoint{3.356954in}{1.611499in}}%
\pgfpathlineto{\pgfqpoint{3.311573in}{1.648869in}}%
\pgfpathlineto{\pgfqpoint{3.277511in}{1.670285in}}%
\pgfpathclose%
\pgfusepath{stroke,fill}%
\end{pgfscope}%
\begin{pgfscope}%
\pgfpathrectangle{\pgfqpoint{0.050000in}{0.050000in}}{\pgfqpoint{4.900000in}{4.900000in}}%
\pgfusepath{clip}%
\pgfsetbuttcap%
\pgfsetroundjoin%
\definecolor{currentfill}{rgb}{0.700992,0.700992,0.700992}%
\pgfsetfillcolor{currentfill}%
\pgfsetfillopacity{0.900000}%
\pgfsetlinewidth{0.501875pt}%
\definecolor{currentstroke}{rgb}{0.200000,0.200000,0.200000}%
\pgfsetstrokecolor{currentstroke}%
\pgfsetdash{}{0pt}%
\pgfpathmoveto{\pgfqpoint{2.732972in}{2.097277in}}%
\pgfpathlineto{\pgfqpoint{2.779450in}{2.052637in}}%
\pgfpathlineto{\pgfqpoint{2.812639in}{2.031269in}}%
\pgfpathlineto{\pgfqpoint{2.766159in}{2.075900in}}%
\pgfpathlineto{\pgfqpoint{2.732972in}{2.097277in}}%
\pgfpathclose%
\pgfusepath{stroke,fill}%
\end{pgfscope}%
\begin{pgfscope}%
\pgfpathrectangle{\pgfqpoint{0.050000in}{0.050000in}}{\pgfqpoint{4.900000in}{4.900000in}}%
\pgfusepath{clip}%
\pgfsetbuttcap%
\pgfsetroundjoin%
\definecolor{currentfill}{rgb}{0.036417,0.036417,0.036417}%
\pgfsetfillcolor{currentfill}%
\pgfsetfillopacity{0.900000}%
\pgfsetlinewidth{0.501875pt}%
\definecolor{currentstroke}{rgb}{0.200000,0.200000,0.200000}%
\pgfsetstrokecolor{currentstroke}%
\pgfsetdash{}{0pt}%
\pgfpathmoveto{\pgfqpoint{0.986436in}{3.509168in}}%
\pgfpathlineto{\pgfqpoint{1.036714in}{3.461140in}}%
\pgfpathlineto{\pgfqpoint{1.067132in}{3.441928in}}%
\pgfpathlineto{\pgfqpoint{1.016839in}{3.489991in}}%
\pgfpathlineto{\pgfqpoint{0.986436in}{3.509168in}}%
\pgfpathclose%
\pgfusepath{stroke,fill}%
\end{pgfscope}%
\begin{pgfscope}%
\pgfpathrectangle{\pgfqpoint{0.050000in}{0.050000in}}{\pgfqpoint{4.900000in}{4.900000in}}%
\pgfusepath{clip}%
\pgfsetbuttcap%
\pgfsetroundjoin%
\definecolor{currentfill}{rgb}{0.338824,0.338824,0.338824}%
\pgfsetfillcolor{currentfill}%
\pgfsetfillopacity{0.900000}%
\pgfsetlinewidth{0.501875pt}%
\definecolor{currentstroke}{rgb}{0.200000,0.200000,0.200000}%
\pgfsetstrokecolor{currentstroke}%
\pgfsetdash{}{0pt}%
\pgfpathmoveto{\pgfqpoint{1.723494in}{2.913466in}}%
\pgfpathlineto{\pgfqpoint{1.772172in}{2.866731in}}%
\pgfpathlineto{\pgfqpoint{1.803761in}{2.846574in}}%
\pgfpathlineto{\pgfqpoint{1.755074in}{2.893340in}}%
\pgfpathlineto{\pgfqpoint{1.723494in}{2.913466in}}%
\pgfpathclose%
\pgfusepath{stroke,fill}%
\end{pgfscope}%
\begin{pgfscope}%
\pgfpathrectangle{\pgfqpoint{0.050000in}{0.050000in}}{\pgfqpoint{4.900000in}{4.900000in}}%
\pgfusepath{clip}%
\pgfsetbuttcap%
\pgfsetroundjoin%
\definecolor{currentfill}{rgb}{0.595433,0.595433,0.595433}%
\pgfsetfillcolor{currentfill}%
\pgfsetfillopacity{0.900000}%
\pgfsetlinewidth{0.501875pt}%
\definecolor{currentstroke}{rgb}{0.200000,0.200000,0.200000}%
\pgfsetstrokecolor{currentstroke}%
\pgfsetdash{}{0pt}%
\pgfpathmoveto{\pgfqpoint{2.460653in}{2.317715in}}%
\pgfpathlineto{\pgfqpoint{2.507730in}{2.272294in}}%
\pgfpathlineto{\pgfqpoint{2.540492in}{2.251199in}}%
\pgfpathlineto{\pgfqpoint{2.493411in}{2.296642in}}%
\pgfpathlineto{\pgfqpoint{2.460653in}{2.317715in}}%
\pgfpathclose%
\pgfusepath{stroke,fill}%
\end{pgfscope}%
\begin{pgfscope}%
\pgfpathrectangle{\pgfqpoint{0.050000in}{0.050000in}}{\pgfqpoint{4.900000in}{4.900000in}}%
\pgfusepath{clip}%
\pgfsetbuttcap%
\pgfsetroundjoin%
\definecolor{currentfill}{rgb}{0.223299,0.223299,0.223299}%
\pgfsetfillcolor{currentfill}%
\pgfsetfillopacity{0.900000}%
\pgfsetlinewidth{0.501875pt}%
\definecolor{currentstroke}{rgb}{0.200000,0.200000,0.200000}%
\pgfsetstrokecolor{currentstroke}%
\pgfsetdash{}{0pt}%
\pgfpathmoveto{\pgfqpoint{1.450972in}{3.134394in}}%
\pgfpathlineto{\pgfqpoint{1.500252in}{3.087170in}}%
\pgfpathlineto{\pgfqpoint{1.531414in}{3.067360in}}%
\pgfpathlineto{\pgfqpoint{1.482124in}{3.114616in}}%
\pgfpathlineto{\pgfqpoint{1.450972in}{3.134394in}}%
\pgfpathclose%
\pgfusepath{stroke,fill}%
\end{pgfscope}%
\begin{pgfscope}%
\pgfpathrectangle{\pgfqpoint{0.050000in}{0.050000in}}{\pgfqpoint{4.900000in}{4.900000in}}%
\pgfusepath{clip}%
\pgfsetbuttcap%
\pgfsetroundjoin%
\definecolor{currentfill}{rgb}{0.495656,0.495656,0.495656}%
\pgfsetfillcolor{currentfill}%
\pgfsetfillopacity{0.900000}%
\pgfsetlinewidth{0.501875pt}%
\definecolor{currentstroke}{rgb}{0.200000,0.200000,0.200000}%
\pgfsetstrokecolor{currentstroke}%
\pgfsetdash{}{0pt}%
\pgfpathmoveto{\pgfqpoint{2.188240in}{2.538534in}}%
\pgfpathlineto{\pgfqpoint{2.235918in}{2.492604in}}%
\pgfpathlineto{\pgfqpoint{2.268253in}{2.471848in}}%
\pgfpathlineto{\pgfqpoint{2.220570in}{2.517806in}}%
\pgfpathlineto{\pgfqpoint{2.188240in}{2.538534in}}%
\pgfpathclose%
\pgfusepath{stroke,fill}%
\end{pgfscope}%
\begin{pgfscope}%
\pgfpathrectangle{\pgfqpoint{0.050000in}{0.050000in}}{\pgfqpoint{4.900000in}{4.900000in}}%
\pgfusepath{clip}%
\pgfsetbuttcap%
\pgfsetroundjoin%
\definecolor{currentfill}{rgb}{0.912526,0.912526,0.912526}%
\pgfsetfillcolor{currentfill}%
\pgfsetfillopacity{0.900000}%
\pgfsetlinewidth{0.501875pt}%
\definecolor{currentstroke}{rgb}{0.200000,0.200000,0.200000}%
\pgfsetstrokecolor{currentstroke}%
\pgfsetdash{}{0pt}%
\pgfpathmoveto{\pgfqpoint{3.550166in}{1.480734in}}%
\pgfpathlineto{\pgfqpoint{3.595146in}{1.451599in}}%
\pgfpathlineto{\pgfqpoint{3.629679in}{1.430606in}}%
\pgfpathlineto{\pgfqpoint{3.584689in}{1.459571in}}%
\pgfpathlineto{\pgfqpoint{3.550166in}{1.480734in}}%
\pgfpathclose%
\pgfusepath{stroke,fill}%
\end{pgfscope}%
\begin{pgfscope}%
\pgfpathrectangle{\pgfqpoint{0.050000in}{0.050000in}}{\pgfqpoint{4.900000in}{4.900000in}}%
\pgfusepath{clip}%
\pgfsetbuttcap%
\pgfsetroundjoin%
\definecolor{currentfill}{rgb}{0.104698,0.104698,0.104698}%
\pgfsetfillcolor{currentfill}%
\pgfsetfillopacity{0.900000}%
\pgfsetlinewidth{0.501875pt}%
\definecolor{currentstroke}{rgb}{0.200000,0.200000,0.200000}%
\pgfsetstrokecolor{currentstroke}%
\pgfsetdash{}{0pt}%
\pgfpathmoveto{\pgfqpoint{1.178356in}{3.355400in}}%
\pgfpathlineto{\pgfqpoint{1.228238in}{3.307687in}}%
\pgfpathlineto{\pgfqpoint{1.258972in}{3.288224in}}%
\pgfpathlineto{\pgfqpoint{1.209078in}{3.335971in}}%
\pgfpathlineto{\pgfqpoint{1.178356in}{3.355400in}}%
\pgfpathclose%
\pgfusepath{stroke,fill}%
\end{pgfscope}%
\begin{pgfscope}%
\pgfpathrectangle{\pgfqpoint{0.050000in}{0.050000in}}{\pgfqpoint{4.900000in}{4.900000in}}%
\pgfusepath{clip}%
\pgfsetbuttcap%
\pgfsetroundjoin%
\definecolor{currentfill}{rgb}{0.839785,0.839785,0.839785}%
\pgfsetfillcolor{currentfill}%
\pgfsetfillopacity{0.900000}%
\pgfsetlinewidth{0.501875pt}%
\definecolor{currentstroke}{rgb}{0.200000,0.200000,0.200000}%
\pgfsetstrokecolor{currentstroke}%
\pgfsetdash{}{0pt}%
\pgfpathmoveto{\pgfqpoint{3.198019in}{1.731192in}}%
\pgfpathlineto{\pgfqpoint{3.243558in}{1.691642in}}%
\pgfpathlineto{\pgfqpoint{3.277511in}{1.670285in}}%
\pgfpathlineto{\pgfqpoint{3.231968in}{1.709694in}}%
\pgfpathlineto{\pgfqpoint{3.198019in}{1.731192in}}%
\pgfpathclose%
\pgfusepath{stroke,fill}%
\end{pgfscope}%
\begin{pgfscope}%
\pgfpathrectangle{\pgfqpoint{0.050000in}{0.050000in}}{\pgfqpoint{4.900000in}{4.900000in}}%
\pgfusepath{clip}%
\pgfsetbuttcap%
\pgfsetroundjoin%
\definecolor{currentfill}{rgb}{0.763998,0.763998,0.763998}%
\pgfsetfillcolor{currentfill}%
\pgfsetfillopacity{0.900000}%
\pgfsetlinewidth{0.501875pt}%
\definecolor{currentstroke}{rgb}{0.200000,0.200000,0.200000}%
\pgfsetstrokecolor{currentstroke}%
\pgfsetdash{}{0pt}%
\pgfpathmoveto{\pgfqpoint{2.925620in}{1.944231in}}%
\pgfpathlineto{\pgfqpoint{2.971706in}{1.900758in}}%
\pgfpathlineto{\pgfqpoint{3.005219in}{1.879294in}}%
\pgfpathlineto{\pgfqpoint{2.959130in}{1.922709in}}%
\pgfpathlineto{\pgfqpoint{2.925620in}{1.944231in}}%
\pgfpathclose%
\pgfusepath{stroke,fill}%
\end{pgfscope}%
\begin{pgfscope}%
\pgfpathrectangle{\pgfqpoint{0.050000in}{0.050000in}}{\pgfqpoint{4.900000in}{4.900000in}}%
\pgfusepath{clip}%
\pgfsetbuttcap%
\pgfsetroundjoin%
\definecolor{currentfill}{rgb}{0.403783,0.403783,0.403783}%
\pgfsetfillcolor{currentfill}%
\pgfsetfillopacity{0.900000}%
\pgfsetlinewidth{0.501875pt}%
\definecolor{currentstroke}{rgb}{0.200000,0.200000,0.200000}%
\pgfsetstrokecolor{currentstroke}%
\pgfsetdash{}{0pt}%
\pgfpathmoveto{\pgfqpoint{1.915733in}{2.759446in}}%
\pgfpathlineto{\pgfqpoint{1.964013in}{2.713026in}}%
\pgfpathlineto{\pgfqpoint{1.995921in}{2.692617in}}%
\pgfpathlineto{\pgfqpoint{1.947634in}{2.739066in}}%
\pgfpathlineto{\pgfqpoint{1.915733in}{2.759446in}}%
\pgfpathclose%
\pgfusepath{stroke,fill}%
\end{pgfscope}%
\begin{pgfscope}%
\pgfpathrectangle{\pgfqpoint{0.050000in}{0.050000in}}{\pgfqpoint{4.900000in}{4.900000in}}%
\pgfusepath{clip}%
\pgfsetbuttcap%
\pgfsetroundjoin%
\definecolor{currentfill}{rgb}{0.667405,0.667405,0.667405}%
\pgfsetfillcolor{currentfill}%
\pgfsetfillopacity{0.900000}%
\pgfsetlinewidth{0.501875pt}%
\definecolor{currentstroke}{rgb}{0.200000,0.200000,0.200000}%
\pgfsetstrokecolor{currentstroke}%
\pgfsetdash{}{0pt}%
\pgfpathmoveto{\pgfqpoint{2.653212in}{2.163565in}}%
\pgfpathlineto{\pgfqpoint{2.699890in}{2.118590in}}%
\pgfpathlineto{\pgfqpoint{2.732972in}{2.097277in}}%
\pgfpathlineto{\pgfqpoint{2.686291in}{2.142258in}}%
\pgfpathlineto{\pgfqpoint{2.653212in}{2.163565in}}%
\pgfpathclose%
\pgfusepath{stroke,fill}%
\end{pgfscope}%
\begin{pgfscope}%
\pgfpathrectangle{\pgfqpoint{0.050000in}{0.050000in}}{\pgfqpoint{4.900000in}{4.900000in}}%
\pgfusepath{clip}%
\pgfsetbuttcap%
\pgfsetroundjoin%
\definecolor{currentfill}{rgb}{0.000000,0.000000,0.000000}%
\pgfsetfillcolor{currentfill}%
\pgfsetfillopacity{0.900000}%
\pgfsetlinewidth{0.501875pt}%
\definecolor{currentstroke}{rgb}{0.200000,0.200000,0.200000}%
\pgfsetstrokecolor{currentstroke}%
\pgfsetdash{}{0pt}%
\pgfpathmoveto{\pgfqpoint{0.905646in}{3.576486in}}%
\pgfpathlineto{\pgfqpoint{0.956130in}{3.528284in}}%
\pgfpathlineto{\pgfqpoint{0.986436in}{3.509168in}}%
\pgfpathlineto{\pgfqpoint{0.935938in}{3.557406in}}%
\pgfpathlineto{\pgfqpoint{0.905646in}{3.576486in}}%
\pgfpathclose%
\pgfusepath{stroke,fill}%
\end{pgfscope}%
\begin{pgfscope}%
\pgfpathrectangle{\pgfqpoint{0.050000in}{0.050000in}}{\pgfqpoint{4.900000in}{4.900000in}}%
\pgfusepath{clip}%
\pgfsetbuttcap%
\pgfsetroundjoin%
\definecolor{currentfill}{rgb}{0.306344,0.306344,0.306344}%
\pgfsetfillcolor{currentfill}%
\pgfsetfillopacity{0.900000}%
\pgfsetlinewidth{0.501875pt}%
\definecolor{currentstroke}{rgb}{0.200000,0.200000,0.200000}%
\pgfsetstrokecolor{currentstroke}%
\pgfsetdash{}{0pt}%
\pgfpathmoveto{\pgfqpoint{1.643132in}{2.980437in}}%
\pgfpathlineto{\pgfqpoint{1.692013in}{2.933528in}}%
\pgfpathlineto{\pgfqpoint{1.723494in}{2.913466in}}%
\pgfpathlineto{\pgfqpoint{1.674603in}{2.960405in}}%
\pgfpathlineto{\pgfqpoint{1.643132in}{2.980437in}}%
\pgfpathclose%
\pgfusepath{stroke,fill}%
\end{pgfscope}%
\begin{pgfscope}%
\pgfpathrectangle{\pgfqpoint{0.050000in}{0.050000in}}{\pgfqpoint{4.900000in}{4.900000in}}%
\pgfusepath{clip}%
\pgfsetbuttcap%
\pgfsetroundjoin%
\definecolor{currentfill}{rgb}{0.560246,0.560246,0.560246}%
\pgfsetfillcolor{currentfill}%
\pgfsetfillopacity{0.900000}%
\pgfsetlinewidth{0.501875pt}%
\definecolor{currentstroke}{rgb}{0.200000,0.200000,0.200000}%
\pgfsetstrokecolor{currentstroke}%
\pgfsetdash{}{0pt}%
\pgfpathmoveto{\pgfqpoint{2.380719in}{2.384328in}}%
\pgfpathlineto{\pgfqpoint{2.427998in}{2.338721in}}%
\pgfpathlineto{\pgfqpoint{2.460653in}{2.317715in}}%
\pgfpathlineto{\pgfqpoint{2.413370in}{2.363347in}}%
\pgfpathlineto{\pgfqpoint{2.380719in}{2.384328in}}%
\pgfpathclose%
\pgfusepath{stroke,fill}%
\end{pgfscope}%
\begin{pgfscope}%
\pgfpathrectangle{\pgfqpoint{0.050000in}{0.050000in}}{\pgfqpoint{4.900000in}{4.900000in}}%
\pgfusepath{clip}%
\pgfsetbuttcap%
\pgfsetroundjoin%
\definecolor{currentfill}{rgb}{0.895548,0.895548,0.895548}%
\pgfsetfillcolor{currentfill}%
\pgfsetfillopacity{0.900000}%
\pgfsetlinewidth{0.501875pt}%
\definecolor{currentstroke}{rgb}{0.200000,0.200000,0.200000}%
\pgfsetstrokecolor{currentstroke}%
\pgfsetdash{}{0pt}%
\pgfpathmoveto{\pgfqpoint{3.470655in}{1.534022in}}%
\pgfpathlineto{\pgfqpoint{3.515754in}{1.501829in}}%
\pgfpathlineto{\pgfqpoint{3.550166in}{1.480734in}}%
\pgfpathlineto{\pgfqpoint{3.505060in}{1.512751in}}%
\pgfpathlineto{\pgfqpoint{3.470655in}{1.534022in}}%
\pgfpathclose%
\pgfusepath{stroke,fill}%
\end{pgfscope}%
\begin{pgfscope}%
\pgfpathrectangle{\pgfqpoint{0.050000in}{0.050000in}}{\pgfqpoint{4.900000in}{4.900000in}}%
\pgfusepath{clip}%
\pgfsetbuttcap%
\pgfsetroundjoin%
\definecolor{currentfill}{rgb}{0.179008,0.179008,0.179008}%
\pgfsetfillcolor{currentfill}%
\pgfsetfillopacity{0.900000}%
\pgfsetlinewidth{0.501875pt}%
\definecolor{currentstroke}{rgb}{0.200000,0.200000,0.200000}%
\pgfsetstrokecolor{currentstroke}%
\pgfsetdash{}{0pt}%
\pgfpathmoveto{\pgfqpoint{1.370436in}{3.201507in}}%
\pgfpathlineto{\pgfqpoint{1.419920in}{3.154109in}}%
\pgfpathlineto{\pgfqpoint{1.450972in}{3.134394in}}%
\pgfpathlineto{\pgfqpoint{1.401477in}{3.181824in}}%
\pgfpathlineto{\pgfqpoint{1.370436in}{3.201507in}}%
\pgfpathclose%
\pgfusepath{stroke,fill}%
\end{pgfscope}%
\begin{pgfscope}%
\pgfpathrectangle{\pgfqpoint{0.050000in}{0.050000in}}{\pgfqpoint{4.900000in}{4.900000in}}%
\pgfusepath{clip}%
\pgfsetbuttcap%
\pgfsetroundjoin%
\definecolor{currentfill}{rgb}{0.465513,0.465513,0.465513}%
\pgfsetfillcolor{currentfill}%
\pgfsetfillopacity{0.900000}%
\pgfsetlinewidth{0.501875pt}%
\definecolor{currentstroke}{rgb}{0.200000,0.200000,0.200000}%
\pgfsetstrokecolor{currentstroke}%
\pgfsetdash{}{0pt}%
\pgfpathmoveto{\pgfqpoint{2.108132in}{2.605299in}}%
\pgfpathlineto{\pgfqpoint{2.156013in}{2.559196in}}%
\pgfpathlineto{\pgfqpoint{2.188240in}{2.538534in}}%
\pgfpathlineto{\pgfqpoint{2.140353in}{2.584665in}}%
\pgfpathlineto{\pgfqpoint{2.108132in}{2.605299in}}%
\pgfpathclose%
\pgfusepath{stroke,fill}%
\end{pgfscope}%
\begin{pgfscope}%
\pgfpathrectangle{\pgfqpoint{0.050000in}{0.050000in}}{\pgfqpoint{4.900000in}{4.900000in}}%
\pgfusepath{clip}%
\pgfsetbuttcap%
\pgfsetroundjoin%
\definecolor{currentfill}{rgb}{0.815671,0.815671,0.815671}%
\pgfsetfillcolor{currentfill}%
\pgfsetfillopacity{0.900000}%
\pgfsetlinewidth{0.501875pt}%
\definecolor{currentstroke}{rgb}{0.200000,0.200000,0.200000}%
\pgfsetstrokecolor{currentstroke}%
\pgfsetdash{}{0pt}%
\pgfpathmoveto{\pgfqpoint{3.118465in}{1.793815in}}%
\pgfpathlineto{\pgfqpoint{3.164179in}{1.752632in}}%
\pgfpathlineto{\pgfqpoint{3.198019in}{1.731192in}}%
\pgfpathlineto{\pgfqpoint{3.152302in}{1.772261in}}%
\pgfpathlineto{\pgfqpoint{3.118465in}{1.793815in}}%
\pgfpathclose%
\pgfusepath{stroke,fill}%
\end{pgfscope}%
\begin{pgfscope}%
\pgfpathrectangle{\pgfqpoint{0.050000in}{0.050000in}}{\pgfqpoint{4.900000in}{4.900000in}}%
\pgfusepath{clip}%
\pgfsetbuttcap%
\pgfsetroundjoin%
\definecolor{currentfill}{rgb}{0.739377,0.739377,0.739377}%
\pgfsetfillcolor{currentfill}%
\pgfsetfillopacity{0.900000}%
\pgfsetlinewidth{0.501875pt}%
\definecolor{currentstroke}{rgb}{0.200000,0.200000,0.200000}%
\pgfsetstrokecolor{currentstroke}%
\pgfsetdash{}{0pt}%
\pgfpathmoveto{\pgfqpoint{2.845935in}{2.009839in}}%
\pgfpathlineto{\pgfqpoint{2.892217in}{1.965692in}}%
\pgfpathlineto{\pgfqpoint{2.925620in}{1.944231in}}%
\pgfpathlineto{\pgfqpoint{2.879337in}{1.988347in}}%
\pgfpathlineto{\pgfqpoint{2.845935in}{2.009839in}}%
\pgfpathclose%
\pgfusepath{stroke,fill}%
\end{pgfscope}%
\begin{pgfscope}%
\pgfpathrectangle{\pgfqpoint{0.050000in}{0.050000in}}{\pgfqpoint{4.900000in}{4.900000in}}%
\pgfusepath{clip}%
\pgfsetbuttcap%
\pgfsetroundjoin%
\definecolor{currentfill}{rgb}{0.072834,0.072834,0.072834}%
\pgfsetfillcolor{currentfill}%
\pgfsetfillopacity{0.900000}%
\pgfsetlinewidth{0.501875pt}%
\definecolor{currentstroke}{rgb}{0.200000,0.200000,0.200000}%
\pgfsetstrokecolor{currentstroke}%
\pgfsetdash{}{0pt}%
\pgfpathmoveto{\pgfqpoint{1.097645in}{3.422656in}}%
\pgfpathlineto{\pgfqpoint{1.147732in}{3.374768in}}%
\pgfpathlineto{\pgfqpoint{1.178356in}{3.355400in}}%
\pgfpathlineto{\pgfqpoint{1.128257in}{3.403322in}}%
\pgfpathlineto{\pgfqpoint{1.097645in}{3.422656in}}%
\pgfpathclose%
\pgfusepath{stroke,fill}%
\end{pgfscope}%
\begin{pgfscope}%
\pgfpathrectangle{\pgfqpoint{0.050000in}{0.050000in}}{\pgfqpoint{4.900000in}{4.900000in}}%
\pgfusepath{clip}%
\pgfsetbuttcap%
\pgfsetroundjoin%
\definecolor{currentfill}{rgb}{0.371303,0.371303,0.371303}%
\pgfsetfillcolor{currentfill}%
\pgfsetfillopacity{0.900000}%
\pgfsetlinewidth{0.501875pt}%
\definecolor{currentstroke}{rgb}{0.200000,0.200000,0.200000}%
\pgfsetstrokecolor{currentstroke}%
\pgfsetdash{}{0pt}%
\pgfpathmoveto{\pgfqpoint{1.835451in}{2.826353in}}%
\pgfpathlineto{\pgfqpoint{1.883934in}{2.779761in}}%
\pgfpathlineto{\pgfqpoint{1.915733in}{2.759446in}}%
\pgfpathlineto{\pgfqpoint{1.867242in}{2.806068in}}%
\pgfpathlineto{\pgfqpoint{1.835451in}{2.826353in}}%
\pgfpathclose%
\pgfusepath{stroke,fill}%
\end{pgfscope}%
\begin{pgfscope}%
\pgfpathrectangle{\pgfqpoint{0.050000in}{0.050000in}}{\pgfqpoint{4.900000in}{4.900000in}}%
\pgfusepath{clip}%
\pgfsetbuttcap%
\pgfsetroundjoin%
\definecolor{currentfill}{rgb}{0.878570,0.878570,0.878570}%
\pgfsetfillcolor{currentfill}%
\pgfsetfillopacity{0.900000}%
\pgfsetlinewidth{0.501875pt}%
\definecolor{currentstroke}{rgb}{0.200000,0.200000,0.200000}%
\pgfsetstrokecolor{currentstroke}%
\pgfsetdash{}{0pt}%
\pgfpathmoveto{\pgfqpoint{3.391130in}{1.590183in}}%
\pgfpathlineto{\pgfqpoint{3.436362in}{1.555228in}}%
\pgfpathlineto{\pgfqpoint{3.470655in}{1.534022in}}%
\pgfpathlineto{\pgfqpoint{3.425417in}{1.568805in}}%
\pgfpathlineto{\pgfqpoint{3.391130in}{1.590183in}}%
\pgfpathclose%
\pgfusepath{stroke,fill}%
\end{pgfscope}%
\begin{pgfscope}%
\pgfpathrectangle{\pgfqpoint{0.050000in}{0.050000in}}{\pgfqpoint{4.900000in}{4.900000in}}%
\pgfusepath{clip}%
\pgfsetbuttcap%
\pgfsetroundjoin%
\definecolor{currentfill}{rgb}{0.629020,0.629020,0.629020}%
\pgfsetfillcolor{currentfill}%
\pgfsetfillopacity{0.900000}%
\pgfsetlinewidth{0.501875pt}%
\definecolor{currentstroke}{rgb}{0.200000,0.200000,0.200000}%
\pgfsetstrokecolor{currentstroke}%
\pgfsetdash{}{0pt}%
\pgfpathmoveto{\pgfqpoint{2.573359in}{2.230039in}}%
\pgfpathlineto{\pgfqpoint{2.620238in}{2.184807in}}%
\pgfpathlineto{\pgfqpoint{2.653212in}{2.163565in}}%
\pgfpathlineto{\pgfqpoint{2.606330in}{2.208812in}}%
\pgfpathlineto{\pgfqpoint{2.573359in}{2.230039in}}%
\pgfpathclose%
\pgfusepath{stroke,fill}%
\end{pgfscope}%
\begin{pgfscope}%
\pgfpathrectangle{\pgfqpoint{0.050000in}{0.050000in}}{\pgfqpoint{4.900000in}{4.900000in}}%
\pgfusepath{clip}%
\pgfsetbuttcap%
\pgfsetroundjoin%
\definecolor{currentfill}{rgb}{0.262053,0.262053,0.262053}%
\pgfsetfillcolor{currentfill}%
\pgfsetfillopacity{0.900000}%
\pgfsetlinewidth{0.501875pt}%
\definecolor{currentstroke}{rgb}{0.200000,0.200000,0.200000}%
\pgfsetstrokecolor{currentstroke}%
\pgfsetdash{}{0pt}%
\pgfpathmoveto{\pgfqpoint{1.562675in}{3.047487in}}%
\pgfpathlineto{\pgfqpoint{1.611761in}{3.000404in}}%
\pgfpathlineto{\pgfqpoint{1.643132in}{2.980437in}}%
\pgfpathlineto{\pgfqpoint{1.594037in}{3.027550in}}%
\pgfpathlineto{\pgfqpoint{1.562675in}{3.047487in}}%
\pgfpathclose%
\pgfusepath{stroke,fill}%
\end{pgfscope}%
\begin{pgfscope}%
\pgfpathrectangle{\pgfqpoint{0.050000in}{0.050000in}}{\pgfqpoint{4.900000in}{4.900000in}}%
\pgfusepath{clip}%
\pgfsetbuttcap%
\pgfsetroundjoin%
\definecolor{currentfill}{rgb}{0.530104,0.530104,0.530104}%
\pgfsetfillcolor{currentfill}%
\pgfsetfillopacity{0.900000}%
\pgfsetlinewidth{0.501875pt}%
\definecolor{currentstroke}{rgb}{0.200000,0.200000,0.200000}%
\pgfsetstrokecolor{currentstroke}%
\pgfsetdash{}{0pt}%
\pgfpathmoveto{\pgfqpoint{2.300692in}{2.451026in}}%
\pgfpathlineto{\pgfqpoint{2.348173in}{2.405243in}}%
\pgfpathlineto{\pgfqpoint{2.380719in}{2.384328in}}%
\pgfpathlineto{\pgfqpoint{2.333234in}{2.430138in}}%
\pgfpathlineto{\pgfqpoint{2.300692in}{2.451026in}}%
\pgfpathclose%
\pgfusepath{stroke,fill}%
\end{pgfscope}%
\begin{pgfscope}%
\pgfpathrectangle{\pgfqpoint{0.050000in}{0.050000in}}{\pgfqpoint{4.900000in}{4.900000in}}%
\pgfusepath{clip}%
\pgfsetbuttcap%
\pgfsetroundjoin%
\definecolor{currentfill}{rgb}{0.141115,0.141115,0.141115}%
\pgfsetfillcolor{currentfill}%
\pgfsetfillopacity{0.900000}%
\pgfsetlinewidth{0.501875pt}%
\definecolor{currentstroke}{rgb}{0.200000,0.200000,0.200000}%
\pgfsetstrokecolor{currentstroke}%
\pgfsetdash{}{0pt}%
\pgfpathmoveto{\pgfqpoint{1.289805in}{3.268699in}}%
\pgfpathlineto{\pgfqpoint{1.339493in}{3.221127in}}%
\pgfpathlineto{\pgfqpoint{1.370436in}{3.201507in}}%
\pgfpathlineto{\pgfqpoint{1.320736in}{3.249112in}}%
\pgfpathlineto{\pgfqpoint{1.289805in}{3.268699in}}%
\pgfpathclose%
\pgfusepath{stroke,fill}%
\end{pgfscope}%
\begin{pgfscope}%
\pgfpathrectangle{\pgfqpoint{0.050000in}{0.050000in}}{\pgfqpoint{4.900000in}{4.900000in}}%
\pgfusepath{clip}%
\pgfsetbuttcap%
\pgfsetroundjoin%
\definecolor{currentfill}{rgb}{0.791557,0.791557,0.791557}%
\pgfsetfillcolor{currentfill}%
\pgfsetfillopacity{0.900000}%
\pgfsetlinewidth{0.501875pt}%
\definecolor{currentstroke}{rgb}{0.200000,0.200000,0.200000}%
\pgfsetstrokecolor{currentstroke}%
\pgfsetdash{}{0pt}%
\pgfpathmoveto{\pgfqpoint{3.038838in}{1.857770in}}%
\pgfpathlineto{\pgfqpoint{3.084736in}{1.815312in}}%
\pgfpathlineto{\pgfqpoint{3.118465in}{1.793815in}}%
\pgfpathlineto{\pgfqpoint{3.072565in}{1.836186in}}%
\pgfpathlineto{\pgfqpoint{3.038838in}{1.857770in}}%
\pgfpathclose%
\pgfusepath{stroke,fill}%
\end{pgfscope}%
\begin{pgfscope}%
\pgfpathrectangle{\pgfqpoint{0.050000in}{0.050000in}}{\pgfqpoint{4.900000in}{4.900000in}}%
\pgfusepath{clip}%
\pgfsetbuttcap%
\pgfsetroundjoin%
\definecolor{currentfill}{rgb}{0.432203,0.432203,0.432203}%
\pgfsetfillcolor{currentfill}%
\pgfsetfillopacity{0.900000}%
\pgfsetlinewidth{0.501875pt}%
\definecolor{currentstroke}{rgb}{0.200000,0.200000,0.200000}%
\pgfsetstrokecolor{currentstroke}%
\pgfsetdash{}{0pt}%
\pgfpathmoveto{\pgfqpoint{2.027930in}{2.672143in}}%
\pgfpathlineto{\pgfqpoint{2.076014in}{2.625867in}}%
\pgfpathlineto{\pgfqpoint{2.108132in}{2.605299in}}%
\pgfpathlineto{\pgfqpoint{2.060042in}{2.651603in}}%
\pgfpathlineto{\pgfqpoint{2.027930in}{2.672143in}}%
\pgfpathclose%
\pgfusepath{stroke,fill}%
\end{pgfscope}%
\begin{pgfscope}%
\pgfpathrectangle{\pgfqpoint{0.050000in}{0.050000in}}{\pgfqpoint{4.900000in}{4.900000in}}%
\pgfusepath{clip}%
\pgfsetbuttcap%
\pgfsetroundjoin%
\definecolor{currentfill}{rgb}{0.858762,0.858762,0.858762}%
\pgfsetfillcolor{currentfill}%
\pgfsetfillopacity{0.900000}%
\pgfsetlinewidth{0.501875pt}%
\definecolor{currentstroke}{rgb}{0.200000,0.200000,0.200000}%
\pgfsetstrokecolor{currentstroke}%
\pgfsetdash{}{0pt}%
\pgfpathmoveto{\pgfqpoint{3.311573in}{1.648869in}}%
\pgfpathlineto{\pgfqpoint{3.356954in}{1.611499in}}%
\pgfpathlineto{\pgfqpoint{3.391130in}{1.590183in}}%
\pgfpathlineto{\pgfqpoint{3.345744in}{1.627392in}}%
\pgfpathlineto{\pgfqpoint{3.311573in}{1.648869in}}%
\pgfpathclose%
\pgfusepath{stroke,fill}%
\end{pgfscope}%
\begin{pgfscope}%
\pgfpathrectangle{\pgfqpoint{0.050000in}{0.050000in}}{\pgfqpoint{4.900000in}{4.900000in}}%
\pgfusepath{clip}%
\pgfsetbuttcap%
\pgfsetroundjoin%
\definecolor{currentfill}{rgb}{0.705790,0.705790,0.705790}%
\pgfsetfillcolor{currentfill}%
\pgfsetfillopacity{0.900000}%
\pgfsetlinewidth{0.501875pt}%
\definecolor{currentstroke}{rgb}{0.200000,0.200000,0.200000}%
\pgfsetstrokecolor{currentstroke}%
\pgfsetdash{}{0pt}%
\pgfpathmoveto{\pgfqpoint{2.766159in}{2.075900in}}%
\pgfpathlineto{\pgfqpoint{2.812639in}{2.031269in}}%
\pgfpathlineto{\pgfqpoint{2.845935in}{2.009839in}}%
\pgfpathlineto{\pgfqpoint{2.799453in}{2.054459in}}%
\pgfpathlineto{\pgfqpoint{2.766159in}{2.075900in}}%
\pgfpathclose%
\pgfusepath{stroke,fill}%
\end{pgfscope}%
\begin{pgfscope}%
\pgfpathrectangle{\pgfqpoint{0.050000in}{0.050000in}}{\pgfqpoint{4.900000in}{4.900000in}}%
\pgfusepath{clip}%
\pgfsetbuttcap%
\pgfsetroundjoin%
\definecolor{currentfill}{rgb}{0.036417,0.036417,0.036417}%
\pgfsetfillcolor{currentfill}%
\pgfsetfillopacity{0.900000}%
\pgfsetlinewidth{0.501875pt}%
\definecolor{currentstroke}{rgb}{0.200000,0.200000,0.200000}%
\pgfsetstrokecolor{currentstroke}%
\pgfsetdash{}{0pt}%
\pgfpathmoveto{\pgfqpoint{1.016839in}{3.489991in}}%
\pgfpathlineto{\pgfqpoint{1.067132in}{3.441928in}}%
\pgfpathlineto{\pgfqpoint{1.097645in}{3.422656in}}%
\pgfpathlineto{\pgfqpoint{1.047340in}{3.470753in}}%
\pgfpathlineto{\pgfqpoint{1.016839in}{3.489991in}}%
\pgfpathclose%
\pgfusepath{stroke,fill}%
\end{pgfscope}%
\begin{pgfscope}%
\pgfpathrectangle{\pgfqpoint{0.050000in}{0.050000in}}{\pgfqpoint{4.900000in}{4.900000in}}%
\pgfusepath{clip}%
\pgfsetbuttcap%
\pgfsetroundjoin%
\definecolor{currentfill}{rgb}{0.338824,0.338824,0.338824}%
\pgfsetfillcolor{currentfill}%
\pgfsetfillopacity{0.900000}%
\pgfsetlinewidth{0.501875pt}%
\definecolor{currentstroke}{rgb}{0.200000,0.200000,0.200000}%
\pgfsetstrokecolor{currentstroke}%
\pgfsetdash{}{0pt}%
\pgfpathmoveto{\pgfqpoint{1.755074in}{2.893340in}}%
\pgfpathlineto{\pgfqpoint{1.803761in}{2.846574in}}%
\pgfpathlineto{\pgfqpoint{1.835451in}{2.826353in}}%
\pgfpathlineto{\pgfqpoint{1.786756in}{2.873149in}}%
\pgfpathlineto{\pgfqpoint{1.755074in}{2.893340in}}%
\pgfpathclose%
\pgfusepath{stroke,fill}%
\end{pgfscope}%
\begin{pgfscope}%
\pgfpathrectangle{\pgfqpoint{0.050000in}{0.050000in}}{\pgfqpoint{4.900000in}{4.900000in}}%
\pgfusepath{clip}%
\pgfsetbuttcap%
\pgfsetroundjoin%
\definecolor{currentfill}{rgb}{0.595433,0.595433,0.595433}%
\pgfsetfillcolor{currentfill}%
\pgfsetfillopacity{0.900000}%
\pgfsetlinewidth{0.501875pt}%
\definecolor{currentstroke}{rgb}{0.200000,0.200000,0.200000}%
\pgfsetstrokecolor{currentstroke}%
\pgfsetdash{}{0pt}%
\pgfpathmoveto{\pgfqpoint{2.493411in}{2.296642in}}%
\pgfpathlineto{\pgfqpoint{2.540492in}{2.251199in}}%
\pgfpathlineto{\pgfqpoint{2.573359in}{2.230039in}}%
\pgfpathlineto{\pgfqpoint{2.526274in}{2.275503in}}%
\pgfpathlineto{\pgfqpoint{2.493411in}{2.296642in}}%
\pgfpathclose%
\pgfusepath{stroke,fill}%
\end{pgfscope}%
\begin{pgfscope}%
\pgfpathrectangle{\pgfqpoint{0.050000in}{0.050000in}}{\pgfqpoint{4.900000in}{4.900000in}}%
\pgfusepath{clip}%
\pgfsetbuttcap%
\pgfsetroundjoin%
\definecolor{currentfill}{rgb}{0.223299,0.223299,0.223299}%
\pgfsetfillcolor{currentfill}%
\pgfsetfillopacity{0.900000}%
\pgfsetlinewidth{0.501875pt}%
\definecolor{currentstroke}{rgb}{0.200000,0.200000,0.200000}%
\pgfsetstrokecolor{currentstroke}%
\pgfsetdash{}{0pt}%
\pgfpathmoveto{\pgfqpoint{1.482124in}{3.114616in}}%
\pgfpathlineto{\pgfqpoint{1.531414in}{3.067360in}}%
\pgfpathlineto{\pgfqpoint{1.562675in}{3.047487in}}%
\pgfpathlineto{\pgfqpoint{1.513375in}{3.094775in}}%
\pgfpathlineto{\pgfqpoint{1.482124in}{3.114616in}}%
\pgfpathclose%
\pgfusepath{stroke,fill}%
\end{pgfscope}%
\begin{pgfscope}%
\pgfpathrectangle{\pgfqpoint{0.050000in}{0.050000in}}{\pgfqpoint{4.900000in}{4.900000in}}%
\pgfusepath{clip}%
\pgfsetbuttcap%
\pgfsetroundjoin%
\definecolor{currentfill}{rgb}{0.495656,0.495656,0.495656}%
\pgfsetfillcolor{currentfill}%
\pgfsetfillopacity{0.900000}%
\pgfsetlinewidth{0.501875pt}%
\definecolor{currentstroke}{rgb}{0.200000,0.200000,0.200000}%
\pgfsetstrokecolor{currentstroke}%
\pgfsetdash{}{0pt}%
\pgfpathmoveto{\pgfqpoint{2.220570in}{2.517806in}}%
\pgfpathlineto{\pgfqpoint{2.268253in}{2.471848in}}%
\pgfpathlineto{\pgfqpoint{2.300692in}{2.451026in}}%
\pgfpathlineto{\pgfqpoint{2.253003in}{2.497012in}}%
\pgfpathlineto{\pgfqpoint{2.220570in}{2.517806in}}%
\pgfpathclose%
\pgfusepath{stroke,fill}%
\end{pgfscope}%
\begin{pgfscope}%
\pgfpathrectangle{\pgfqpoint{0.050000in}{0.050000in}}{\pgfqpoint{4.900000in}{4.900000in}}%
\pgfusepath{clip}%
\pgfsetbuttcap%
\pgfsetroundjoin%
\definecolor{currentfill}{rgb}{0.104698,0.104698,0.104698}%
\pgfsetfillcolor{currentfill}%
\pgfsetfillopacity{0.900000}%
\pgfsetlinewidth{0.501875pt}%
\definecolor{currentstroke}{rgb}{0.200000,0.200000,0.200000}%
\pgfsetstrokecolor{currentstroke}%
\pgfsetdash{}{0pt}%
\pgfpathmoveto{\pgfqpoint{1.209078in}{3.335971in}}%
\pgfpathlineto{\pgfqpoint{1.258972in}{3.288224in}}%
\pgfpathlineto{\pgfqpoint{1.289805in}{3.268699in}}%
\pgfpathlineto{\pgfqpoint{1.239899in}{3.316480in}}%
\pgfpathlineto{\pgfqpoint{1.209078in}{3.335971in}}%
\pgfpathclose%
\pgfusepath{stroke,fill}%
\end{pgfscope}%
\begin{pgfscope}%
\pgfpathrectangle{\pgfqpoint{0.050000in}{0.050000in}}{\pgfqpoint{4.900000in}{4.900000in}}%
\pgfusepath{clip}%
\pgfsetbuttcap%
\pgfsetroundjoin%
\definecolor{currentfill}{rgb}{0.763998,0.763998,0.763998}%
\pgfsetfillcolor{currentfill}%
\pgfsetfillopacity{0.900000}%
\pgfsetlinewidth{0.501875pt}%
\definecolor{currentstroke}{rgb}{0.200000,0.200000,0.200000}%
\pgfsetstrokecolor{currentstroke}%
\pgfsetdash{}{0pt}%
\pgfpathmoveto{\pgfqpoint{2.959130in}{1.922709in}}%
\pgfpathlineto{\pgfqpoint{3.005219in}{1.879294in}}%
\pgfpathlineto{\pgfqpoint{3.038838in}{1.857770in}}%
\pgfpathlineto{\pgfqpoint{2.992748in}{1.901127in}}%
\pgfpathlineto{\pgfqpoint{2.959130in}{1.922709in}}%
\pgfpathclose%
\pgfusepath{stroke,fill}%
\end{pgfscope}%
\begin{pgfscope}%
\pgfpathrectangle{\pgfqpoint{0.050000in}{0.050000in}}{\pgfqpoint{4.900000in}{4.900000in}}%
\pgfusepath{clip}%
\pgfsetbuttcap%
\pgfsetroundjoin%
\definecolor{currentfill}{rgb}{0.839785,0.839785,0.839785}%
\pgfsetfillcolor{currentfill}%
\pgfsetfillopacity{0.900000}%
\pgfsetlinewidth{0.501875pt}%
\definecolor{currentstroke}{rgb}{0.200000,0.200000,0.200000}%
\pgfsetstrokecolor{currentstroke}%
\pgfsetdash{}{0pt}%
\pgfpathmoveto{\pgfqpoint{3.231968in}{1.709694in}}%
\pgfpathlineto{\pgfqpoint{3.277511in}{1.670285in}}%
\pgfpathlineto{\pgfqpoint{3.311573in}{1.648869in}}%
\pgfpathlineto{\pgfqpoint{3.266026in}{1.688137in}}%
\pgfpathlineto{\pgfqpoint{3.231968in}{1.709694in}}%
\pgfpathclose%
\pgfusepath{stroke,fill}%
\end{pgfscope}%
\begin{pgfscope}%
\pgfpathrectangle{\pgfqpoint{0.050000in}{0.050000in}}{\pgfqpoint{4.900000in}{4.900000in}}%
\pgfusepath{clip}%
\pgfsetbuttcap%
\pgfsetroundjoin%
\definecolor{currentfill}{rgb}{0.403783,0.403783,0.403783}%
\pgfsetfillcolor{currentfill}%
\pgfsetfillopacity{0.900000}%
\pgfsetlinewidth{0.501875pt}%
\definecolor{currentstroke}{rgb}{0.200000,0.200000,0.200000}%
\pgfsetstrokecolor{currentstroke}%
\pgfsetdash{}{0pt}%
\pgfpathmoveto{\pgfqpoint{1.947634in}{2.739066in}}%
\pgfpathlineto{\pgfqpoint{1.995921in}{2.692617in}}%
\pgfpathlineto{\pgfqpoint{2.027930in}{2.672143in}}%
\pgfpathlineto{\pgfqpoint{1.979636in}{2.718621in}}%
\pgfpathlineto{\pgfqpoint{1.947634in}{2.739066in}}%
\pgfpathclose%
\pgfusepath{stroke,fill}%
\end{pgfscope}%
\begin{pgfscope}%
\pgfpathrectangle{\pgfqpoint{0.050000in}{0.050000in}}{\pgfqpoint{4.900000in}{4.900000in}}%
\pgfusepath{clip}%
\pgfsetbuttcap%
\pgfsetroundjoin%
\definecolor{currentfill}{rgb}{0.667405,0.667405,0.667405}%
\pgfsetfillcolor{currentfill}%
\pgfsetfillopacity{0.900000}%
\pgfsetlinewidth{0.501875pt}%
\definecolor{currentstroke}{rgb}{0.200000,0.200000,0.200000}%
\pgfsetstrokecolor{currentstroke}%
\pgfsetdash{}{0pt}%
\pgfpathmoveto{\pgfqpoint{2.686291in}{2.142258in}}%
\pgfpathlineto{\pgfqpoint{2.732972in}{2.097277in}}%
\pgfpathlineto{\pgfqpoint{2.766159in}{2.075900in}}%
\pgfpathlineto{\pgfqpoint{2.719477in}{2.120885in}}%
\pgfpathlineto{\pgfqpoint{2.686291in}{2.142258in}}%
\pgfpathclose%
\pgfusepath{stroke,fill}%
\end{pgfscope}%
\begin{pgfscope}%
\pgfpathrectangle{\pgfqpoint{0.050000in}{0.050000in}}{\pgfqpoint{4.900000in}{4.900000in}}%
\pgfusepath{clip}%
\pgfsetbuttcap%
\pgfsetroundjoin%
\definecolor{currentfill}{rgb}{0.004552,0.004552,0.004552}%
\pgfsetfillcolor{currentfill}%
\pgfsetfillopacity{0.900000}%
\pgfsetlinewidth{0.501875pt}%
\definecolor{currentstroke}{rgb}{0.200000,0.200000,0.200000}%
\pgfsetstrokecolor{currentstroke}%
\pgfsetdash{}{0pt}%
\pgfpathmoveto{\pgfqpoint{0.935938in}{3.557406in}}%
\pgfpathlineto{\pgfqpoint{0.986436in}{3.509168in}}%
\pgfpathlineto{\pgfqpoint{1.016839in}{3.489991in}}%
\pgfpathlineto{\pgfqpoint{0.966327in}{3.538264in}}%
\pgfpathlineto{\pgfqpoint{0.935938in}{3.557406in}}%
\pgfpathclose%
\pgfusepath{stroke,fill}%
\end{pgfscope}%
\begin{pgfscope}%
\pgfpathrectangle{\pgfqpoint{0.050000in}{0.050000in}}{\pgfqpoint{4.900000in}{4.900000in}}%
\pgfusepath{clip}%
\pgfsetbuttcap%
\pgfsetroundjoin%
\definecolor{currentfill}{rgb}{0.306344,0.306344,0.306344}%
\pgfsetfillcolor{currentfill}%
\pgfsetfillopacity{0.900000}%
\pgfsetlinewidth{0.501875pt}%
\definecolor{currentstroke}{rgb}{0.200000,0.200000,0.200000}%
\pgfsetstrokecolor{currentstroke}%
\pgfsetdash{}{0pt}%
\pgfpathmoveto{\pgfqpoint{1.674603in}{2.960405in}}%
\pgfpathlineto{\pgfqpoint{1.723494in}{2.913466in}}%
\pgfpathlineto{\pgfqpoint{1.755074in}{2.893340in}}%
\pgfpathlineto{\pgfqpoint{1.706175in}{2.940310in}}%
\pgfpathlineto{\pgfqpoint{1.674603in}{2.960405in}}%
\pgfpathclose%
\pgfusepath{stroke,fill}%
\end{pgfscope}%
\begin{pgfscope}%
\pgfpathrectangle{\pgfqpoint{0.050000in}{0.050000in}}{\pgfqpoint{4.900000in}{4.900000in}}%
\pgfusepath{clip}%
\pgfsetbuttcap%
\pgfsetroundjoin%
\definecolor{currentfill}{rgb}{0.560246,0.560246,0.560246}%
\pgfsetfillcolor{currentfill}%
\pgfsetfillopacity{0.900000}%
\pgfsetlinewidth{0.501875pt}%
\definecolor{currentstroke}{rgb}{0.200000,0.200000,0.200000}%
\pgfsetstrokecolor{currentstroke}%
\pgfsetdash{}{0pt}%
\pgfpathmoveto{\pgfqpoint{2.413370in}{2.363347in}}%
\pgfpathlineto{\pgfqpoint{2.460653in}{2.317715in}}%
\pgfpathlineto{\pgfqpoint{2.493411in}{2.296642in}}%
\pgfpathlineto{\pgfqpoint{2.446124in}{2.342299in}}%
\pgfpathlineto{\pgfqpoint{2.413370in}{2.363347in}}%
\pgfpathclose%
\pgfusepath{stroke,fill}%
\end{pgfscope}%
\begin{pgfscope}%
\pgfpathrectangle{\pgfqpoint{0.050000in}{0.050000in}}{\pgfqpoint{4.900000in}{4.900000in}}%
\pgfusepath{clip}%
\pgfsetbuttcap%
\pgfsetroundjoin%
\definecolor{currentfill}{rgb}{0.895548,0.895548,0.895548}%
\pgfsetfillcolor{currentfill}%
\pgfsetfillopacity{0.900000}%
\pgfsetlinewidth{0.501875pt}%
\definecolor{currentstroke}{rgb}{0.200000,0.200000,0.200000}%
\pgfsetstrokecolor{currentstroke}%
\pgfsetdash{}{0pt}%
\pgfpathmoveto{\pgfqpoint{3.505060in}{1.512751in}}%
\pgfpathlineto{\pgfqpoint{3.550166in}{1.480734in}}%
\pgfpathlineto{\pgfqpoint{3.584689in}{1.459571in}}%
\pgfpathlineto{\pgfqpoint{3.539575in}{1.491415in}}%
\pgfpathlineto{\pgfqpoint{3.505060in}{1.512751in}}%
\pgfpathclose%
\pgfusepath{stroke,fill}%
\end{pgfscope}%
\begin{pgfscope}%
\pgfpathrectangle{\pgfqpoint{0.050000in}{0.050000in}}{\pgfqpoint{4.900000in}{4.900000in}}%
\pgfusepath{clip}%
\pgfsetbuttcap%
\pgfsetroundjoin%
\definecolor{currentfill}{rgb}{0.179008,0.179008,0.179008}%
\pgfsetfillcolor{currentfill}%
\pgfsetfillopacity{0.900000}%
\pgfsetlinewidth{0.501875pt}%
\definecolor{currentstroke}{rgb}{0.200000,0.200000,0.200000}%
\pgfsetstrokecolor{currentstroke}%
\pgfsetdash{}{0pt}%
\pgfpathmoveto{\pgfqpoint{1.401477in}{3.181824in}}%
\pgfpathlineto{\pgfqpoint{1.450972in}{3.134394in}}%
\pgfpathlineto{\pgfqpoint{1.482124in}{3.114616in}}%
\pgfpathlineto{\pgfqpoint{1.432618in}{3.162079in}}%
\pgfpathlineto{\pgfqpoint{1.401477in}{3.181824in}}%
\pgfpathclose%
\pgfusepath{stroke,fill}%
\end{pgfscope}%
\begin{pgfscope}%
\pgfpathrectangle{\pgfqpoint{0.050000in}{0.050000in}}{\pgfqpoint{4.900000in}{4.900000in}}%
\pgfusepath{clip}%
\pgfsetbuttcap%
\pgfsetroundjoin%
\definecolor{currentfill}{rgb}{0.465513,0.465513,0.465513}%
\pgfsetfillcolor{currentfill}%
\pgfsetfillopacity{0.900000}%
\pgfsetlinewidth{0.501875pt}%
\definecolor{currentstroke}{rgb}{0.200000,0.200000,0.200000}%
\pgfsetstrokecolor{currentstroke}%
\pgfsetdash{}{0pt}%
\pgfpathmoveto{\pgfqpoint{2.140353in}{2.584665in}}%
\pgfpathlineto{\pgfqpoint{2.188240in}{2.538534in}}%
\pgfpathlineto{\pgfqpoint{2.220570in}{2.517806in}}%
\pgfpathlineto{\pgfqpoint{2.172677in}{2.563965in}}%
\pgfpathlineto{\pgfqpoint{2.140353in}{2.584665in}}%
\pgfpathclose%
\pgfusepath{stroke,fill}%
\end{pgfscope}%
\begin{pgfscope}%
\pgfpathrectangle{\pgfqpoint{0.050000in}{0.050000in}}{\pgfqpoint{4.900000in}{4.900000in}}%
\pgfusepath{clip}%
\pgfsetbuttcap%
\pgfsetroundjoin%
\definecolor{currentfill}{rgb}{0.815671,0.815671,0.815671}%
\pgfsetfillcolor{currentfill}%
\pgfsetfillopacity{0.900000}%
\pgfsetlinewidth{0.501875pt}%
\definecolor{currentstroke}{rgb}{0.200000,0.200000,0.200000}%
\pgfsetstrokecolor{currentstroke}%
\pgfsetdash{}{0pt}%
\pgfpathmoveto{\pgfqpoint{3.152302in}{1.772261in}}%
\pgfpathlineto{\pgfqpoint{3.198019in}{1.731192in}}%
\pgfpathlineto{\pgfqpoint{3.231968in}{1.709694in}}%
\pgfpathlineto{\pgfqpoint{3.186248in}{1.750647in}}%
\pgfpathlineto{\pgfqpoint{3.152302in}{1.772261in}}%
\pgfpathclose%
\pgfusepath{stroke,fill}%
\end{pgfscope}%
\begin{pgfscope}%
\pgfpathrectangle{\pgfqpoint{0.050000in}{0.050000in}}{\pgfqpoint{4.900000in}{4.900000in}}%
\pgfusepath{clip}%
\pgfsetbuttcap%
\pgfsetroundjoin%
\definecolor{currentfill}{rgb}{0.739377,0.739377,0.739377}%
\pgfsetfillcolor{currentfill}%
\pgfsetfillopacity{0.900000}%
\pgfsetlinewidth{0.501875pt}%
\definecolor{currentstroke}{rgb}{0.200000,0.200000,0.200000}%
\pgfsetstrokecolor{currentstroke}%
\pgfsetdash{}{0pt}%
\pgfpathmoveto{\pgfqpoint{2.879337in}{1.988347in}}%
\pgfpathlineto{\pgfqpoint{2.925620in}{1.944231in}}%
\pgfpathlineto{\pgfqpoint{2.959130in}{1.922709in}}%
\pgfpathlineto{\pgfqpoint{2.912846in}{1.966792in}}%
\pgfpathlineto{\pgfqpoint{2.879337in}{1.988347in}}%
\pgfpathclose%
\pgfusepath{stroke,fill}%
\end{pgfscope}%
\begin{pgfscope}%
\pgfpathrectangle{\pgfqpoint{0.050000in}{0.050000in}}{\pgfqpoint{4.900000in}{4.900000in}}%
\pgfusepath{clip}%
\pgfsetbuttcap%
\pgfsetroundjoin%
\definecolor{currentfill}{rgb}{0.072834,0.072834,0.072834}%
\pgfsetfillcolor{currentfill}%
\pgfsetfillopacity{0.900000}%
\pgfsetlinewidth{0.501875pt}%
\definecolor{currentstroke}{rgb}{0.200000,0.200000,0.200000}%
\pgfsetstrokecolor{currentstroke}%
\pgfsetdash{}{0pt}%
\pgfpathmoveto{\pgfqpoint{1.128257in}{3.403322in}}%
\pgfpathlineto{\pgfqpoint{1.178356in}{3.355400in}}%
\pgfpathlineto{\pgfqpoint{1.209078in}{3.335971in}}%
\pgfpathlineto{\pgfqpoint{1.158966in}{3.383927in}}%
\pgfpathlineto{\pgfqpoint{1.128257in}{3.403322in}}%
\pgfpathclose%
\pgfusepath{stroke,fill}%
\end{pgfscope}%
\begin{pgfscope}%
\pgfpathrectangle{\pgfqpoint{0.050000in}{0.050000in}}{\pgfqpoint{4.900000in}{4.900000in}}%
\pgfusepath{clip}%
\pgfsetbuttcap%
\pgfsetroundjoin%
\definecolor{currentfill}{rgb}{0.371303,0.371303,0.371303}%
\pgfsetfillcolor{currentfill}%
\pgfsetfillopacity{0.900000}%
\pgfsetlinewidth{0.501875pt}%
\definecolor{currentstroke}{rgb}{0.200000,0.200000,0.200000}%
\pgfsetstrokecolor{currentstroke}%
\pgfsetdash{}{0pt}%
\pgfpathmoveto{\pgfqpoint{1.867242in}{2.806068in}}%
\pgfpathlineto{\pgfqpoint{1.915733in}{2.759446in}}%
\pgfpathlineto{\pgfqpoint{1.947634in}{2.739066in}}%
\pgfpathlineto{\pgfqpoint{1.899135in}{2.785718in}}%
\pgfpathlineto{\pgfqpoint{1.867242in}{2.806068in}}%
\pgfpathclose%
\pgfusepath{stroke,fill}%
\end{pgfscope}%
\begin{pgfscope}%
\pgfpathrectangle{\pgfqpoint{0.050000in}{0.050000in}}{\pgfqpoint{4.900000in}{4.900000in}}%
\pgfusepath{clip}%
\pgfsetbuttcap%
\pgfsetroundjoin%
\definecolor{currentfill}{rgb}{0.878570,0.878570,0.878570}%
\pgfsetfillcolor{currentfill}%
\pgfsetfillopacity{0.900000}%
\pgfsetlinewidth{0.501875pt}%
\definecolor{currentstroke}{rgb}{0.200000,0.200000,0.200000}%
\pgfsetstrokecolor{currentstroke}%
\pgfsetdash{}{0pt}%
\pgfpathmoveto{\pgfqpoint{3.425417in}{1.568805in}}%
\pgfpathlineto{\pgfqpoint{3.470655in}{1.534022in}}%
\pgfpathlineto{\pgfqpoint{3.505060in}{1.512751in}}%
\pgfpathlineto{\pgfqpoint{3.459814in}{1.547364in}}%
\pgfpathlineto{\pgfqpoint{3.425417in}{1.568805in}}%
\pgfpathclose%
\pgfusepath{stroke,fill}%
\end{pgfscope}%
\begin{pgfscope}%
\pgfpathrectangle{\pgfqpoint{0.050000in}{0.050000in}}{\pgfqpoint{4.900000in}{4.900000in}}%
\pgfusepath{clip}%
\pgfsetbuttcap%
\pgfsetroundjoin%
\definecolor{currentfill}{rgb}{0.633818,0.633818,0.633818}%
\pgfsetfillcolor{currentfill}%
\pgfsetfillopacity{0.900000}%
\pgfsetlinewidth{0.501875pt}%
\definecolor{currentstroke}{rgb}{0.200000,0.200000,0.200000}%
\pgfsetstrokecolor{currentstroke}%
\pgfsetdash{}{0pt}%
\pgfpathmoveto{\pgfqpoint{2.606330in}{2.208812in}}%
\pgfpathlineto{\pgfqpoint{2.653212in}{2.163565in}}%
\pgfpathlineto{\pgfqpoint{2.686291in}{2.142258in}}%
\pgfpathlineto{\pgfqpoint{2.639407in}{2.187519in}}%
\pgfpathlineto{\pgfqpoint{2.606330in}{2.208812in}}%
\pgfpathclose%
\pgfusepath{stroke,fill}%
\end{pgfscope}%
\begin{pgfscope}%
\pgfpathrectangle{\pgfqpoint{0.050000in}{0.050000in}}{\pgfqpoint{4.900000in}{4.900000in}}%
\pgfusepath{clip}%
\pgfsetbuttcap%
\pgfsetroundjoin%
\definecolor{currentfill}{rgb}{0.262053,0.262053,0.262053}%
\pgfsetfillcolor{currentfill}%
\pgfsetfillopacity{0.900000}%
\pgfsetlinewidth{0.501875pt}%
\definecolor{currentstroke}{rgb}{0.200000,0.200000,0.200000}%
\pgfsetstrokecolor{currentstroke}%
\pgfsetdash{}{0pt}%
\pgfpathmoveto{\pgfqpoint{1.594037in}{3.027550in}}%
\pgfpathlineto{\pgfqpoint{1.643132in}{2.980437in}}%
\pgfpathlineto{\pgfqpoint{1.674603in}{2.960405in}}%
\pgfpathlineto{\pgfqpoint{1.625498in}{3.007550in}}%
\pgfpathlineto{\pgfqpoint{1.594037in}{3.027550in}}%
\pgfpathclose%
\pgfusepath{stroke,fill}%
\end{pgfscope}%
\begin{pgfscope}%
\pgfpathrectangle{\pgfqpoint{0.050000in}{0.050000in}}{\pgfqpoint{4.900000in}{4.900000in}}%
\pgfusepath{clip}%
\pgfsetbuttcap%
\pgfsetroundjoin%
\definecolor{currentfill}{rgb}{0.530104,0.530104,0.530104}%
\pgfsetfillcolor{currentfill}%
\pgfsetfillopacity{0.900000}%
\pgfsetlinewidth{0.501875pt}%
\definecolor{currentstroke}{rgb}{0.200000,0.200000,0.200000}%
\pgfsetstrokecolor{currentstroke}%
\pgfsetdash{}{0pt}%
\pgfpathmoveto{\pgfqpoint{2.333234in}{2.430138in}}%
\pgfpathlineto{\pgfqpoint{2.380719in}{2.384328in}}%
\pgfpathlineto{\pgfqpoint{2.413370in}{2.363347in}}%
\pgfpathlineto{\pgfqpoint{2.365879in}{2.409184in}}%
\pgfpathlineto{\pgfqpoint{2.333234in}{2.430138in}}%
\pgfpathclose%
\pgfusepath{stroke,fill}%
\end{pgfscope}%
\begin{pgfscope}%
\pgfpathrectangle{\pgfqpoint{0.050000in}{0.050000in}}{\pgfqpoint{4.900000in}{4.900000in}}%
\pgfusepath{clip}%
\pgfsetbuttcap%
\pgfsetroundjoin%
\definecolor{currentfill}{rgb}{0.141115,0.141115,0.141115}%
\pgfsetfillcolor{currentfill}%
\pgfsetfillopacity{0.900000}%
\pgfsetlinewidth{0.501875pt}%
\definecolor{currentstroke}{rgb}{0.200000,0.200000,0.200000}%
\pgfsetstrokecolor{currentstroke}%
\pgfsetdash{}{0pt}%
\pgfpathmoveto{\pgfqpoint{1.320736in}{3.249112in}}%
\pgfpathlineto{\pgfqpoint{1.370436in}{3.201507in}}%
\pgfpathlineto{\pgfqpoint{1.401477in}{3.181824in}}%
\pgfpathlineto{\pgfqpoint{1.351766in}{3.229462in}}%
\pgfpathlineto{\pgfqpoint{1.320736in}{3.249112in}}%
\pgfpathclose%
\pgfusepath{stroke,fill}%
\end{pgfscope}%
\begin{pgfscope}%
\pgfpathrectangle{\pgfqpoint{0.050000in}{0.050000in}}{\pgfqpoint{4.900000in}{4.900000in}}%
\pgfusepath{clip}%
\pgfsetbuttcap%
\pgfsetroundjoin%
\definecolor{currentfill}{rgb}{0.791557,0.791557,0.791557}%
\pgfsetfillcolor{currentfill}%
\pgfsetfillopacity{0.900000}%
\pgfsetlinewidth{0.501875pt}%
\definecolor{currentstroke}{rgb}{0.200000,0.200000,0.200000}%
\pgfsetstrokecolor{currentstroke}%
\pgfsetdash{}{0pt}%
\pgfpathmoveto{\pgfqpoint{3.072565in}{1.836186in}}%
\pgfpathlineto{\pgfqpoint{3.118465in}{1.793815in}}%
\pgfpathlineto{\pgfqpoint{3.152302in}{1.772261in}}%
\pgfpathlineto{\pgfqpoint{3.106400in}{1.814543in}}%
\pgfpathlineto{\pgfqpoint{3.072565in}{1.836186in}}%
\pgfpathclose%
\pgfusepath{stroke,fill}%
\end{pgfscope}%
\begin{pgfscope}%
\pgfpathrectangle{\pgfqpoint{0.050000in}{0.050000in}}{\pgfqpoint{4.900000in}{4.900000in}}%
\pgfusepath{clip}%
\pgfsetbuttcap%
\pgfsetroundjoin%
\definecolor{currentfill}{rgb}{0.432203,0.432203,0.432203}%
\pgfsetfillcolor{currentfill}%
\pgfsetfillopacity{0.900000}%
\pgfsetlinewidth{0.501875pt}%
\definecolor{currentstroke}{rgb}{0.200000,0.200000,0.200000}%
\pgfsetstrokecolor{currentstroke}%
\pgfsetdash{}{0pt}%
\pgfpathmoveto{\pgfqpoint{2.060042in}{2.651603in}}%
\pgfpathlineto{\pgfqpoint{2.108132in}{2.605299in}}%
\pgfpathlineto{\pgfqpoint{2.140353in}{2.584665in}}%
\pgfpathlineto{\pgfqpoint{2.092257in}{2.630998in}}%
\pgfpathlineto{\pgfqpoint{2.060042in}{2.651603in}}%
\pgfpathclose%
\pgfusepath{stroke,fill}%
\end{pgfscope}%
\begin{pgfscope}%
\pgfpathrectangle{\pgfqpoint{0.050000in}{0.050000in}}{\pgfqpoint{4.900000in}{4.900000in}}%
\pgfusepath{clip}%
\pgfsetbuttcap%
\pgfsetroundjoin%
\definecolor{currentfill}{rgb}{0.858762,0.858762,0.858762}%
\pgfsetfillcolor{currentfill}%
\pgfsetfillopacity{0.900000}%
\pgfsetlinewidth{0.501875pt}%
\definecolor{currentstroke}{rgb}{0.200000,0.200000,0.200000}%
\pgfsetstrokecolor{currentstroke}%
\pgfsetdash{}{0pt}%
\pgfpathmoveto{\pgfqpoint{3.345744in}{1.627392in}}%
\pgfpathlineto{\pgfqpoint{3.391130in}{1.590183in}}%
\pgfpathlineto{\pgfqpoint{3.425417in}{1.568805in}}%
\pgfpathlineto{\pgfqpoint{3.380026in}{1.605855in}}%
\pgfpathlineto{\pgfqpoint{3.345744in}{1.627392in}}%
\pgfpathclose%
\pgfusepath{stroke,fill}%
\end{pgfscope}%
\begin{pgfscope}%
\pgfpathrectangle{\pgfqpoint{0.050000in}{0.050000in}}{\pgfqpoint{4.900000in}{4.900000in}}%
\pgfusepath{clip}%
\pgfsetbuttcap%
\pgfsetroundjoin%
\definecolor{currentfill}{rgb}{0.705790,0.705790,0.705790}%
\pgfsetfillcolor{currentfill}%
\pgfsetfillopacity{0.900000}%
\pgfsetlinewidth{0.501875pt}%
\definecolor{currentstroke}{rgb}{0.200000,0.200000,0.200000}%
\pgfsetstrokecolor{currentstroke}%
\pgfsetdash{}{0pt}%
\pgfpathmoveto{\pgfqpoint{2.799453in}{2.054459in}}%
\pgfpathlineto{\pgfqpoint{2.845935in}{2.009839in}}%
\pgfpathlineto{\pgfqpoint{2.879337in}{1.988347in}}%
\pgfpathlineto{\pgfqpoint{2.832853in}{2.032953in}}%
\pgfpathlineto{\pgfqpoint{2.799453in}{2.054459in}}%
\pgfpathclose%
\pgfusepath{stroke,fill}%
\end{pgfscope}%
\begin{pgfscope}%
\pgfpathrectangle{\pgfqpoint{0.050000in}{0.050000in}}{\pgfqpoint{4.900000in}{4.900000in}}%
\pgfusepath{clip}%
\pgfsetbuttcap%
\pgfsetroundjoin%
\definecolor{currentfill}{rgb}{0.036417,0.036417,0.036417}%
\pgfsetfillcolor{currentfill}%
\pgfsetfillopacity{0.900000}%
\pgfsetlinewidth{0.501875pt}%
\definecolor{currentstroke}{rgb}{0.200000,0.200000,0.200000}%
\pgfsetstrokecolor{currentstroke}%
\pgfsetdash{}{0pt}%
\pgfpathmoveto{\pgfqpoint{1.047340in}{3.470753in}}%
\pgfpathlineto{\pgfqpoint{1.097645in}{3.422656in}}%
\pgfpathlineto{\pgfqpoint{1.128257in}{3.403322in}}%
\pgfpathlineto{\pgfqpoint{1.077938in}{3.451454in}}%
\pgfpathlineto{\pgfqpoint{1.047340in}{3.470753in}}%
\pgfpathclose%
\pgfusepath{stroke,fill}%
\end{pgfscope}%
\begin{pgfscope}%
\pgfpathrectangle{\pgfqpoint{0.050000in}{0.050000in}}{\pgfqpoint{4.900000in}{4.900000in}}%
\pgfusepath{clip}%
\pgfsetbuttcap%
\pgfsetroundjoin%
\definecolor{currentfill}{rgb}{0.342884,0.342884,0.342884}%
\pgfsetfillcolor{currentfill}%
\pgfsetfillopacity{0.900000}%
\pgfsetlinewidth{0.501875pt}%
\definecolor{currentstroke}{rgb}{0.200000,0.200000,0.200000}%
\pgfsetstrokecolor{currentstroke}%
\pgfsetdash{}{0pt}%
\pgfpathmoveto{\pgfqpoint{1.786756in}{2.873149in}}%
\pgfpathlineto{\pgfqpoint{1.835451in}{2.826353in}}%
\pgfpathlineto{\pgfqpoint{1.867242in}{2.806068in}}%
\pgfpathlineto{\pgfqpoint{1.818539in}{2.852895in}}%
\pgfpathlineto{\pgfqpoint{1.786756in}{2.873149in}}%
\pgfpathclose%
\pgfusepath{stroke,fill}%
\end{pgfscope}%
\begin{pgfscope}%
\pgfpathrectangle{\pgfqpoint{0.050000in}{0.050000in}}{\pgfqpoint{4.900000in}{4.900000in}}%
\pgfusepath{clip}%
\pgfsetbuttcap%
\pgfsetroundjoin%
\definecolor{currentfill}{rgb}{0.595433,0.595433,0.595433}%
\pgfsetfillcolor{currentfill}%
\pgfsetfillopacity{0.900000}%
\pgfsetlinewidth{0.501875pt}%
\definecolor{currentstroke}{rgb}{0.200000,0.200000,0.200000}%
\pgfsetstrokecolor{currentstroke}%
\pgfsetdash{}{0pt}%
\pgfpathmoveto{\pgfqpoint{2.526274in}{2.275503in}}%
\pgfpathlineto{\pgfqpoint{2.573359in}{2.230039in}}%
\pgfpathlineto{\pgfqpoint{2.606330in}{2.208812in}}%
\pgfpathlineto{\pgfqpoint{2.559243in}{2.254297in}}%
\pgfpathlineto{\pgfqpoint{2.526274in}{2.275503in}}%
\pgfpathclose%
\pgfusepath{stroke,fill}%
\end{pgfscope}%
\begin{pgfscope}%
\pgfpathrectangle{\pgfqpoint{0.050000in}{0.050000in}}{\pgfqpoint{4.900000in}{4.900000in}}%
\pgfusepath{clip}%
\pgfsetbuttcap%
\pgfsetroundjoin%
\definecolor{currentfill}{rgb}{0.223299,0.223299,0.223299}%
\pgfsetfillcolor{currentfill}%
\pgfsetfillopacity{0.900000}%
\pgfsetlinewidth{0.501875pt}%
\definecolor{currentstroke}{rgb}{0.200000,0.200000,0.200000}%
\pgfsetstrokecolor{currentstroke}%
\pgfsetdash{}{0pt}%
\pgfpathmoveto{\pgfqpoint{1.513375in}{3.094775in}}%
\pgfpathlineto{\pgfqpoint{1.562675in}{3.047487in}}%
\pgfpathlineto{\pgfqpoint{1.594037in}{3.027550in}}%
\pgfpathlineto{\pgfqpoint{1.544726in}{3.074870in}}%
\pgfpathlineto{\pgfqpoint{1.513375in}{3.094775in}}%
\pgfpathclose%
\pgfusepath{stroke,fill}%
\end{pgfscope}%
\begin{pgfscope}%
\pgfpathrectangle{\pgfqpoint{0.050000in}{0.050000in}}{\pgfqpoint{4.900000in}{4.900000in}}%
\pgfusepath{clip}%
\pgfsetbuttcap%
\pgfsetroundjoin%
\definecolor{currentfill}{rgb}{0.495656,0.495656,0.495656}%
\pgfsetfillcolor{currentfill}%
\pgfsetfillopacity{0.900000}%
\pgfsetlinewidth{0.501875pt}%
\definecolor{currentstroke}{rgb}{0.200000,0.200000,0.200000}%
\pgfsetstrokecolor{currentstroke}%
\pgfsetdash{}{0pt}%
\pgfpathmoveto{\pgfqpoint{2.253003in}{2.497012in}}%
\pgfpathlineto{\pgfqpoint{2.300692in}{2.451026in}}%
\pgfpathlineto{\pgfqpoint{2.333234in}{2.430138in}}%
\pgfpathlineto{\pgfqpoint{2.285540in}{2.476152in}}%
\pgfpathlineto{\pgfqpoint{2.253003in}{2.497012in}}%
\pgfpathclose%
\pgfusepath{stroke,fill}%
\end{pgfscope}%
\begin{pgfscope}%
\pgfpathrectangle{\pgfqpoint{0.050000in}{0.050000in}}{\pgfqpoint{4.900000in}{4.900000in}}%
\pgfusepath{clip}%
\pgfsetbuttcap%
\pgfsetroundjoin%
\definecolor{currentfill}{rgb}{0.104698,0.104698,0.104698}%
\pgfsetfillcolor{currentfill}%
\pgfsetfillopacity{0.900000}%
\pgfsetlinewidth{0.501875pt}%
\definecolor{currentstroke}{rgb}{0.200000,0.200000,0.200000}%
\pgfsetstrokecolor{currentstroke}%
\pgfsetdash{}{0pt}%
\pgfpathmoveto{\pgfqpoint{1.239899in}{3.316480in}}%
\pgfpathlineto{\pgfqpoint{1.289805in}{3.268699in}}%
\pgfpathlineto{\pgfqpoint{1.320736in}{3.249112in}}%
\pgfpathlineto{\pgfqpoint{1.270818in}{3.296926in}}%
\pgfpathlineto{\pgfqpoint{1.239899in}{3.316480in}}%
\pgfpathclose%
\pgfusepath{stroke,fill}%
\end{pgfscope}%
\begin{pgfscope}%
\pgfpathrectangle{\pgfqpoint{0.050000in}{0.050000in}}{\pgfqpoint{4.900000in}{4.900000in}}%
\pgfusepath{clip}%
\pgfsetbuttcap%
\pgfsetroundjoin%
\definecolor{currentfill}{rgb}{0.763998,0.763998,0.763998}%
\pgfsetfillcolor{currentfill}%
\pgfsetfillopacity{0.900000}%
\pgfsetlinewidth{0.501875pt}%
\definecolor{currentstroke}{rgb}{0.200000,0.200000,0.200000}%
\pgfsetstrokecolor{currentstroke}%
\pgfsetdash{}{0pt}%
\pgfpathmoveto{\pgfqpoint{2.992748in}{1.901127in}}%
\pgfpathlineto{\pgfqpoint{3.038838in}{1.857770in}}%
\pgfpathlineto{\pgfqpoint{3.072565in}{1.836186in}}%
\pgfpathlineto{\pgfqpoint{3.026474in}{1.879483in}}%
\pgfpathlineto{\pgfqpoint{2.992748in}{1.901127in}}%
\pgfpathclose%
\pgfusepath{stroke,fill}%
\end{pgfscope}%
\begin{pgfscope}%
\pgfpathrectangle{\pgfqpoint{0.050000in}{0.050000in}}{\pgfqpoint{4.900000in}{4.900000in}}%
\pgfusepath{clip}%
\pgfsetbuttcap%
\pgfsetroundjoin%
\definecolor{currentfill}{rgb}{0.839785,0.839785,0.839785}%
\pgfsetfillcolor{currentfill}%
\pgfsetfillopacity{0.900000}%
\pgfsetlinewidth{0.501875pt}%
\definecolor{currentstroke}{rgb}{0.200000,0.200000,0.200000}%
\pgfsetstrokecolor{currentstroke}%
\pgfsetdash{}{0pt}%
\pgfpathmoveto{\pgfqpoint{3.266026in}{1.688137in}}%
\pgfpathlineto{\pgfqpoint{3.311573in}{1.648869in}}%
\pgfpathlineto{\pgfqpoint{3.345744in}{1.627392in}}%
\pgfpathlineto{\pgfqpoint{3.300194in}{1.666520in}}%
\pgfpathlineto{\pgfqpoint{3.266026in}{1.688137in}}%
\pgfpathclose%
\pgfusepath{stroke,fill}%
\end{pgfscope}%
\begin{pgfscope}%
\pgfpathrectangle{\pgfqpoint{0.050000in}{0.050000in}}{\pgfqpoint{4.900000in}{4.900000in}}%
\pgfusepath{clip}%
\pgfsetbuttcap%
\pgfsetroundjoin%
\definecolor{currentfill}{rgb}{0.403783,0.403783,0.403783}%
\pgfsetfillcolor{currentfill}%
\pgfsetfillopacity{0.900000}%
\pgfsetlinewidth{0.501875pt}%
\definecolor{currentstroke}{rgb}{0.200000,0.200000,0.200000}%
\pgfsetstrokecolor{currentstroke}%
\pgfsetdash{}{0pt}%
\pgfpathmoveto{\pgfqpoint{1.979636in}{2.718621in}}%
\pgfpathlineto{\pgfqpoint{2.027930in}{2.672143in}}%
\pgfpathlineto{\pgfqpoint{2.060042in}{2.651603in}}%
\pgfpathlineto{\pgfqpoint{2.011741in}{2.698111in}}%
\pgfpathlineto{\pgfqpoint{1.979636in}{2.718621in}}%
\pgfpathclose%
\pgfusepath{stroke,fill}%
\end{pgfscope}%
\begin{pgfscope}%
\pgfpathrectangle{\pgfqpoint{0.050000in}{0.050000in}}{\pgfqpoint{4.900000in}{4.900000in}}%
\pgfusepath{clip}%
\pgfsetbuttcap%
\pgfsetroundjoin%
\definecolor{currentfill}{rgb}{0.667405,0.667405,0.667405}%
\pgfsetfillcolor{currentfill}%
\pgfsetfillopacity{0.900000}%
\pgfsetlinewidth{0.501875pt}%
\definecolor{currentstroke}{rgb}{0.200000,0.200000,0.200000}%
\pgfsetstrokecolor{currentstroke}%
\pgfsetdash{}{0pt}%
\pgfpathmoveto{\pgfqpoint{2.719477in}{2.120885in}}%
\pgfpathlineto{\pgfqpoint{2.766159in}{2.075900in}}%
\pgfpathlineto{\pgfqpoint{2.799453in}{2.054459in}}%
\pgfpathlineto{\pgfqpoint{2.752768in}{2.099447in}}%
\pgfpathlineto{\pgfqpoint{2.719477in}{2.120885in}}%
\pgfpathclose%
\pgfusepath{stroke,fill}%
\end{pgfscope}%
\begin{pgfscope}%
\pgfpathrectangle{\pgfqpoint{0.050000in}{0.050000in}}{\pgfqpoint{4.900000in}{4.900000in}}%
\pgfusepath{clip}%
\pgfsetbuttcap%
\pgfsetroundjoin%
\definecolor{currentfill}{rgb}{0.004552,0.004552,0.004552}%
\pgfsetfillcolor{currentfill}%
\pgfsetfillopacity{0.900000}%
\pgfsetlinewidth{0.501875pt}%
\definecolor{currentstroke}{rgb}{0.200000,0.200000,0.200000}%
\pgfsetstrokecolor{currentstroke}%
\pgfsetdash{}{0pt}%
\pgfpathmoveto{\pgfqpoint{0.966327in}{3.538264in}}%
\pgfpathlineto{\pgfqpoint{1.016839in}{3.489991in}}%
\pgfpathlineto{\pgfqpoint{1.047340in}{3.470753in}}%
\pgfpathlineto{\pgfqpoint{0.996814in}{3.519062in}}%
\pgfpathlineto{\pgfqpoint{0.966327in}{3.538264in}}%
\pgfpathclose%
\pgfusepath{stroke,fill}%
\end{pgfscope}%
\begin{pgfscope}%
\pgfpathrectangle{\pgfqpoint{0.050000in}{0.050000in}}{\pgfqpoint{4.900000in}{4.900000in}}%
\pgfusepath{clip}%
\pgfsetbuttcap%
\pgfsetroundjoin%
\definecolor{currentfill}{rgb}{0.306344,0.306344,0.306344}%
\pgfsetfillcolor{currentfill}%
\pgfsetfillopacity{0.900000}%
\pgfsetlinewidth{0.501875pt}%
\definecolor{currentstroke}{rgb}{0.200000,0.200000,0.200000}%
\pgfsetstrokecolor{currentstroke}%
\pgfsetdash{}{0pt}%
\pgfpathmoveto{\pgfqpoint{1.706175in}{2.940310in}}%
\pgfpathlineto{\pgfqpoint{1.755074in}{2.893340in}}%
\pgfpathlineto{\pgfqpoint{1.786756in}{2.873149in}}%
\pgfpathlineto{\pgfqpoint{1.737848in}{2.920151in}}%
\pgfpathlineto{\pgfqpoint{1.706175in}{2.940310in}}%
\pgfpathclose%
\pgfusepath{stroke,fill}%
\end{pgfscope}%
\begin{pgfscope}%
\pgfpathrectangle{\pgfqpoint{0.050000in}{0.050000in}}{\pgfqpoint{4.900000in}{4.900000in}}%
\pgfusepath{clip}%
\pgfsetbuttcap%
\pgfsetroundjoin%
\definecolor{currentfill}{rgb}{0.564552,0.564552,0.564552}%
\pgfsetfillcolor{currentfill}%
\pgfsetfillopacity{0.900000}%
\pgfsetlinewidth{0.501875pt}%
\definecolor{currentstroke}{rgb}{0.200000,0.200000,0.200000}%
\pgfsetstrokecolor{currentstroke}%
\pgfsetdash{}{0pt}%
\pgfpathmoveto{\pgfqpoint{2.446124in}{2.342299in}}%
\pgfpathlineto{\pgfqpoint{2.493411in}{2.296642in}}%
\pgfpathlineto{\pgfqpoint{2.526274in}{2.275503in}}%
\pgfpathlineto{\pgfqpoint{2.478984in}{2.321184in}}%
\pgfpathlineto{\pgfqpoint{2.446124in}{2.342299in}}%
\pgfpathclose%
\pgfusepath{stroke,fill}%
\end{pgfscope}%
\begin{pgfscope}%
\pgfpathrectangle{\pgfqpoint{0.050000in}{0.050000in}}{\pgfqpoint{4.900000in}{4.900000in}}%
\pgfusepath{clip}%
\pgfsetbuttcap%
\pgfsetroundjoin%
\definecolor{currentfill}{rgb}{0.179008,0.179008,0.179008}%
\pgfsetfillcolor{currentfill}%
\pgfsetfillopacity{0.900000}%
\pgfsetlinewidth{0.501875pt}%
\definecolor{currentstroke}{rgb}{0.200000,0.200000,0.200000}%
\pgfsetstrokecolor{currentstroke}%
\pgfsetdash{}{0pt}%
\pgfpathmoveto{\pgfqpoint{1.432618in}{3.162079in}}%
\pgfpathlineto{\pgfqpoint{1.482124in}{3.114616in}}%
\pgfpathlineto{\pgfqpoint{1.513375in}{3.094775in}}%
\pgfpathlineto{\pgfqpoint{1.463858in}{3.142270in}}%
\pgfpathlineto{\pgfqpoint{1.432618in}{3.162079in}}%
\pgfpathclose%
\pgfusepath{stroke,fill}%
\end{pgfscope}%
\begin{pgfscope}%
\pgfpathrectangle{\pgfqpoint{0.050000in}{0.050000in}}{\pgfqpoint{4.900000in}{4.900000in}}%
\pgfusepath{clip}%
\pgfsetbuttcap%
\pgfsetroundjoin%
\definecolor{currentfill}{rgb}{0.465513,0.465513,0.465513}%
\pgfsetfillcolor{currentfill}%
\pgfsetfillopacity{0.900000}%
\pgfsetlinewidth{0.501875pt}%
\definecolor{currentstroke}{rgb}{0.200000,0.200000,0.200000}%
\pgfsetstrokecolor{currentstroke}%
\pgfsetdash{}{0pt}%
\pgfpathmoveto{\pgfqpoint{2.172677in}{2.563965in}}%
\pgfpathlineto{\pgfqpoint{2.220570in}{2.517806in}}%
\pgfpathlineto{\pgfqpoint{2.253003in}{2.497012in}}%
\pgfpathlineto{\pgfqpoint{2.205105in}{2.543200in}}%
\pgfpathlineto{\pgfqpoint{2.172677in}{2.563965in}}%
\pgfpathclose%
\pgfusepath{stroke,fill}%
\end{pgfscope}%
\begin{pgfscope}%
\pgfpathrectangle{\pgfqpoint{0.050000in}{0.050000in}}{\pgfqpoint{4.900000in}{4.900000in}}%
\pgfusepath{clip}%
\pgfsetbuttcap%
\pgfsetroundjoin%
\definecolor{currentfill}{rgb}{0.815671,0.815671,0.815671}%
\pgfsetfillcolor{currentfill}%
\pgfsetfillopacity{0.900000}%
\pgfsetlinewidth{0.501875pt}%
\definecolor{currentstroke}{rgb}{0.200000,0.200000,0.200000}%
\pgfsetstrokecolor{currentstroke}%
\pgfsetdash{}{0pt}%
\pgfpathmoveto{\pgfqpoint{3.186248in}{1.750647in}}%
\pgfpathlineto{\pgfqpoint{3.231968in}{1.709694in}}%
\pgfpathlineto{\pgfqpoint{3.266026in}{1.688137in}}%
\pgfpathlineto{\pgfqpoint{3.220304in}{1.728973in}}%
\pgfpathlineto{\pgfqpoint{3.186248in}{1.750647in}}%
\pgfpathclose%
\pgfusepath{stroke,fill}%
\end{pgfscope}%
\begin{pgfscope}%
\pgfpathrectangle{\pgfqpoint{0.050000in}{0.050000in}}{\pgfqpoint{4.900000in}{4.900000in}}%
\pgfusepath{clip}%
\pgfsetbuttcap%
\pgfsetroundjoin%
\definecolor{currentfill}{rgb}{0.739377,0.739377,0.739377}%
\pgfsetfillcolor{currentfill}%
\pgfsetfillopacity{0.900000}%
\pgfsetlinewidth{0.501875pt}%
\definecolor{currentstroke}{rgb}{0.200000,0.200000,0.200000}%
\pgfsetstrokecolor{currentstroke}%
\pgfsetdash{}{0pt}%
\pgfpathmoveto{\pgfqpoint{2.912846in}{1.966792in}}%
\pgfpathlineto{\pgfqpoint{2.959130in}{1.922709in}}%
\pgfpathlineto{\pgfqpoint{2.992748in}{1.901127in}}%
\pgfpathlineto{\pgfqpoint{2.946462in}{1.945173in}}%
\pgfpathlineto{\pgfqpoint{2.912846in}{1.966792in}}%
\pgfpathclose%
\pgfusepath{stroke,fill}%
\end{pgfscope}%
\begin{pgfscope}%
\pgfpathrectangle{\pgfqpoint{0.050000in}{0.050000in}}{\pgfqpoint{4.900000in}{4.900000in}}%
\pgfusepath{clip}%
\pgfsetbuttcap%
\pgfsetroundjoin%
\definecolor{currentfill}{rgb}{0.072834,0.072834,0.072834}%
\pgfsetfillcolor{currentfill}%
\pgfsetfillopacity{0.900000}%
\pgfsetlinewidth{0.501875pt}%
\definecolor{currentstroke}{rgb}{0.200000,0.200000,0.200000}%
\pgfsetstrokecolor{currentstroke}%
\pgfsetdash{}{0pt}%
\pgfpathmoveto{\pgfqpoint{1.158966in}{3.383927in}}%
\pgfpathlineto{\pgfqpoint{1.209078in}{3.335971in}}%
\pgfpathlineto{\pgfqpoint{1.239899in}{3.316480in}}%
\pgfpathlineto{\pgfqpoint{1.189774in}{3.364470in}}%
\pgfpathlineto{\pgfqpoint{1.158966in}{3.383927in}}%
\pgfpathclose%
\pgfusepath{stroke,fill}%
\end{pgfscope}%
\begin{pgfscope}%
\pgfpathrectangle{\pgfqpoint{0.050000in}{0.050000in}}{\pgfqpoint{4.900000in}{4.900000in}}%
\pgfusepath{clip}%
\pgfsetbuttcap%
\pgfsetroundjoin%
\definecolor{currentfill}{rgb}{0.371303,0.371303,0.371303}%
\pgfsetfillcolor{currentfill}%
\pgfsetfillopacity{0.900000}%
\pgfsetlinewidth{0.501875pt}%
\definecolor{currentstroke}{rgb}{0.200000,0.200000,0.200000}%
\pgfsetstrokecolor{currentstroke}%
\pgfsetdash{}{0pt}%
\pgfpathmoveto{\pgfqpoint{1.899135in}{2.785718in}}%
\pgfpathlineto{\pgfqpoint{1.947634in}{2.739066in}}%
\pgfpathlineto{\pgfqpoint{1.979636in}{2.718621in}}%
\pgfpathlineto{\pgfqpoint{1.931130in}{2.765303in}}%
\pgfpathlineto{\pgfqpoint{1.899135in}{2.785718in}}%
\pgfpathclose%
\pgfusepath{stroke,fill}%
\end{pgfscope}%
\begin{pgfscope}%
\pgfpathrectangle{\pgfqpoint{0.050000in}{0.050000in}}{\pgfqpoint{4.900000in}{4.900000in}}%
\pgfusepath{clip}%
\pgfsetbuttcap%
\pgfsetroundjoin%
\definecolor{currentfill}{rgb}{0.633818,0.633818,0.633818}%
\pgfsetfillcolor{currentfill}%
\pgfsetfillopacity{0.900000}%
\pgfsetlinewidth{0.501875pt}%
\definecolor{currentstroke}{rgb}{0.200000,0.200000,0.200000}%
\pgfsetstrokecolor{currentstroke}%
\pgfsetdash{}{0pt}%
\pgfpathmoveto{\pgfqpoint{2.639407in}{2.187519in}}%
\pgfpathlineto{\pgfqpoint{2.686291in}{2.142258in}}%
\pgfpathlineto{\pgfqpoint{2.719477in}{2.120885in}}%
\pgfpathlineto{\pgfqpoint{2.672590in}{2.166160in}}%
\pgfpathlineto{\pgfqpoint{2.639407in}{2.187519in}}%
\pgfpathclose%
\pgfusepath{stroke,fill}%
\end{pgfscope}%
\begin{pgfscope}%
\pgfpathrectangle{\pgfqpoint{0.050000in}{0.050000in}}{\pgfqpoint{4.900000in}{4.900000in}}%
\pgfusepath{clip}%
\pgfsetbuttcap%
\pgfsetroundjoin%
\definecolor{currentfill}{rgb}{0.878570,0.878570,0.878570}%
\pgfsetfillcolor{currentfill}%
\pgfsetfillopacity{0.900000}%
\pgfsetlinewidth{0.501875pt}%
\definecolor{currentstroke}{rgb}{0.200000,0.200000,0.200000}%
\pgfsetstrokecolor{currentstroke}%
\pgfsetdash{}{0pt}%
\pgfpathmoveto{\pgfqpoint{3.459814in}{1.547364in}}%
\pgfpathlineto{\pgfqpoint{3.505060in}{1.512751in}}%
\pgfpathlineto{\pgfqpoint{3.539575in}{1.491415in}}%
\pgfpathlineto{\pgfqpoint{3.494323in}{1.525859in}}%
\pgfpathlineto{\pgfqpoint{3.459814in}{1.547364in}}%
\pgfpathclose%
\pgfusepath{stroke,fill}%
\end{pgfscope}%
\begin{pgfscope}%
\pgfpathrectangle{\pgfqpoint{0.050000in}{0.050000in}}{\pgfqpoint{4.900000in}{4.900000in}}%
\pgfusepath{clip}%
\pgfsetbuttcap%
\pgfsetroundjoin%
\definecolor{currentfill}{rgb}{0.267589,0.267589,0.267589}%
\pgfsetfillcolor{currentfill}%
\pgfsetfillopacity{0.900000}%
\pgfsetlinewidth{0.501875pt}%
\definecolor{currentstroke}{rgb}{0.200000,0.200000,0.200000}%
\pgfsetstrokecolor{currentstroke}%
\pgfsetdash{}{0pt}%
\pgfpathmoveto{\pgfqpoint{1.625498in}{3.007550in}}%
\pgfpathlineto{\pgfqpoint{1.674603in}{2.960405in}}%
\pgfpathlineto{\pgfqpoint{1.706175in}{2.940310in}}%
\pgfpathlineto{\pgfqpoint{1.657060in}{2.987487in}}%
\pgfpathlineto{\pgfqpoint{1.625498in}{3.007550in}}%
\pgfpathclose%
\pgfusepath{stroke,fill}%
\end{pgfscope}%
\begin{pgfscope}%
\pgfpathrectangle{\pgfqpoint{0.050000in}{0.050000in}}{\pgfqpoint{4.900000in}{4.900000in}}%
\pgfusepath{clip}%
\pgfsetbuttcap%
\pgfsetroundjoin%
\definecolor{currentfill}{rgb}{0.530104,0.530104,0.530104}%
\pgfsetfillcolor{currentfill}%
\pgfsetfillopacity{0.900000}%
\pgfsetlinewidth{0.501875pt}%
\definecolor{currentstroke}{rgb}{0.200000,0.200000,0.200000}%
\pgfsetstrokecolor{currentstroke}%
\pgfsetdash{}{0pt}%
\pgfpathmoveto{\pgfqpoint{2.365879in}{2.409184in}}%
\pgfpathlineto{\pgfqpoint{2.413370in}{2.363347in}}%
\pgfpathlineto{\pgfqpoint{2.446124in}{2.342299in}}%
\pgfpathlineto{\pgfqpoint{2.398630in}{2.388162in}}%
\pgfpathlineto{\pgfqpoint{2.365879in}{2.409184in}}%
\pgfpathclose%
\pgfusepath{stroke,fill}%
\end{pgfscope}%
\begin{pgfscope}%
\pgfpathrectangle{\pgfqpoint{0.050000in}{0.050000in}}{\pgfqpoint{4.900000in}{4.900000in}}%
\pgfusepath{clip}%
\pgfsetbuttcap%
\pgfsetroundjoin%
\definecolor{currentfill}{rgb}{0.141115,0.141115,0.141115}%
\pgfsetfillcolor{currentfill}%
\pgfsetfillopacity{0.900000}%
\pgfsetlinewidth{0.501875pt}%
\definecolor{currentstroke}{rgb}{0.200000,0.200000,0.200000}%
\pgfsetstrokecolor{currentstroke}%
\pgfsetdash{}{0pt}%
\pgfpathmoveto{\pgfqpoint{1.351766in}{3.229462in}}%
\pgfpathlineto{\pgfqpoint{1.401477in}{3.181824in}}%
\pgfpathlineto{\pgfqpoint{1.432618in}{3.162079in}}%
\pgfpathlineto{\pgfqpoint{1.382895in}{3.209750in}}%
\pgfpathlineto{\pgfqpoint{1.351766in}{3.229462in}}%
\pgfpathclose%
\pgfusepath{stroke,fill}%
\end{pgfscope}%
\begin{pgfscope}%
\pgfpathrectangle{\pgfqpoint{0.050000in}{0.050000in}}{\pgfqpoint{4.900000in}{4.900000in}}%
\pgfusepath{clip}%
\pgfsetbuttcap%
\pgfsetroundjoin%
\definecolor{currentfill}{rgb}{0.791557,0.791557,0.791557}%
\pgfsetfillcolor{currentfill}%
\pgfsetfillopacity{0.900000}%
\pgfsetlinewidth{0.501875pt}%
\definecolor{currentstroke}{rgb}{0.200000,0.200000,0.200000}%
\pgfsetstrokecolor{currentstroke}%
\pgfsetdash{}{0pt}%
\pgfpathmoveto{\pgfqpoint{3.106400in}{1.814543in}}%
\pgfpathlineto{\pgfqpoint{3.152302in}{1.772261in}}%
\pgfpathlineto{\pgfqpoint{3.186248in}{1.750647in}}%
\pgfpathlineto{\pgfqpoint{3.140345in}{1.792840in}}%
\pgfpathlineto{\pgfqpoint{3.106400in}{1.814543in}}%
\pgfpathclose%
\pgfusepath{stroke,fill}%
\end{pgfscope}%
\begin{pgfscope}%
\pgfpathrectangle{\pgfqpoint{0.050000in}{0.050000in}}{\pgfqpoint{4.900000in}{4.900000in}}%
\pgfusepath{clip}%
\pgfsetbuttcap%
\pgfsetroundjoin%
\definecolor{currentfill}{rgb}{0.432203,0.432203,0.432203}%
\pgfsetfillcolor{currentfill}%
\pgfsetfillopacity{0.900000}%
\pgfsetlinewidth{0.501875pt}%
\definecolor{currentstroke}{rgb}{0.200000,0.200000,0.200000}%
\pgfsetstrokecolor{currentstroke}%
\pgfsetdash{}{0pt}%
\pgfpathmoveto{\pgfqpoint{2.092257in}{2.630998in}}%
\pgfpathlineto{\pgfqpoint{2.140353in}{2.584665in}}%
\pgfpathlineto{\pgfqpoint{2.172677in}{2.563965in}}%
\pgfpathlineto{\pgfqpoint{2.124575in}{2.610328in}}%
\pgfpathlineto{\pgfqpoint{2.092257in}{2.630998in}}%
\pgfpathclose%
\pgfusepath{stroke,fill}%
\end{pgfscope}%
\begin{pgfscope}%
\pgfpathrectangle{\pgfqpoint{0.050000in}{0.050000in}}{\pgfqpoint{4.900000in}{4.900000in}}%
\pgfusepath{clip}%
\pgfsetbuttcap%
\pgfsetroundjoin%
\definecolor{currentfill}{rgb}{0.705790,0.705790,0.705790}%
\pgfsetfillcolor{currentfill}%
\pgfsetfillopacity{0.900000}%
\pgfsetlinewidth{0.501875pt}%
\definecolor{currentstroke}{rgb}{0.200000,0.200000,0.200000}%
\pgfsetstrokecolor{currentstroke}%
\pgfsetdash{}{0pt}%
\pgfpathmoveto{\pgfqpoint{2.832853in}{2.032953in}}%
\pgfpathlineto{\pgfqpoint{2.879337in}{1.988347in}}%
\pgfpathlineto{\pgfqpoint{2.912846in}{1.966792in}}%
\pgfpathlineto{\pgfqpoint{2.866361in}{2.011383in}}%
\pgfpathlineto{\pgfqpoint{2.832853in}{2.032953in}}%
\pgfpathclose%
\pgfusepath{stroke,fill}%
\end{pgfscope}%
\begin{pgfscope}%
\pgfpathrectangle{\pgfqpoint{0.050000in}{0.050000in}}{\pgfqpoint{4.900000in}{4.900000in}}%
\pgfusepath{clip}%
\pgfsetbuttcap%
\pgfsetroundjoin%
\definecolor{currentfill}{rgb}{0.858762,0.858762,0.858762}%
\pgfsetfillcolor{currentfill}%
\pgfsetfillopacity{0.900000}%
\pgfsetlinewidth{0.501875pt}%
\definecolor{currentstroke}{rgb}{0.200000,0.200000,0.200000}%
\pgfsetstrokecolor{currentstroke}%
\pgfsetdash{}{0pt}%
\pgfpathmoveto{\pgfqpoint{3.380026in}{1.605855in}}%
\pgfpathlineto{\pgfqpoint{3.425417in}{1.568805in}}%
\pgfpathlineto{\pgfqpoint{3.459814in}{1.547364in}}%
\pgfpathlineto{\pgfqpoint{3.414419in}{1.584257in}}%
\pgfpathlineto{\pgfqpoint{3.380026in}{1.605855in}}%
\pgfpathclose%
\pgfusepath{stroke,fill}%
\end{pgfscope}%
\begin{pgfscope}%
\pgfpathrectangle{\pgfqpoint{0.050000in}{0.050000in}}{\pgfqpoint{4.900000in}{4.900000in}}%
\pgfusepath{clip}%
\pgfsetbuttcap%
\pgfsetroundjoin%
\definecolor{currentfill}{rgb}{0.036417,0.036417,0.036417}%
\pgfsetfillcolor{currentfill}%
\pgfsetfillopacity{0.900000}%
\pgfsetlinewidth{0.501875pt}%
\definecolor{currentstroke}{rgb}{0.200000,0.200000,0.200000}%
\pgfsetstrokecolor{currentstroke}%
\pgfsetdash{}{0pt}%
\pgfpathmoveto{\pgfqpoint{1.077938in}{3.451454in}}%
\pgfpathlineto{\pgfqpoint{1.128257in}{3.403322in}}%
\pgfpathlineto{\pgfqpoint{1.158966in}{3.383927in}}%
\pgfpathlineto{\pgfqpoint{1.108634in}{3.432093in}}%
\pgfpathlineto{\pgfqpoint{1.077938in}{3.451454in}}%
\pgfpathclose%
\pgfusepath{stroke,fill}%
\end{pgfscope}%
\begin{pgfscope}%
\pgfpathrectangle{\pgfqpoint{0.050000in}{0.050000in}}{\pgfqpoint{4.900000in}{4.900000in}}%
\pgfusepath{clip}%
\pgfsetbuttcap%
\pgfsetroundjoin%
\definecolor{currentfill}{rgb}{0.342884,0.342884,0.342884}%
\pgfsetfillcolor{currentfill}%
\pgfsetfillopacity{0.900000}%
\pgfsetlinewidth{0.501875pt}%
\definecolor{currentstroke}{rgb}{0.200000,0.200000,0.200000}%
\pgfsetstrokecolor{currentstroke}%
\pgfsetdash{}{0pt}%
\pgfpathmoveto{\pgfqpoint{1.818539in}{2.852895in}}%
\pgfpathlineto{\pgfqpoint{1.867242in}{2.806068in}}%
\pgfpathlineto{\pgfqpoint{1.899135in}{2.785718in}}%
\pgfpathlineto{\pgfqpoint{1.850424in}{2.832575in}}%
\pgfpathlineto{\pgfqpoint{1.818539in}{2.852895in}}%
\pgfpathclose%
\pgfusepath{stroke,fill}%
\end{pgfscope}%
\begin{pgfscope}%
\pgfpathrectangle{\pgfqpoint{0.050000in}{0.050000in}}{\pgfqpoint{4.900000in}{4.900000in}}%
\pgfusepath{clip}%
\pgfsetbuttcap%
\pgfsetroundjoin%
\definecolor{currentfill}{rgb}{0.595433,0.595433,0.595433}%
\pgfsetfillcolor{currentfill}%
\pgfsetfillopacity{0.900000}%
\pgfsetlinewidth{0.501875pt}%
\definecolor{currentstroke}{rgb}{0.200000,0.200000,0.200000}%
\pgfsetstrokecolor{currentstroke}%
\pgfsetdash{}{0pt}%
\pgfpathmoveto{\pgfqpoint{2.559243in}{2.254297in}}%
\pgfpathlineto{\pgfqpoint{2.606330in}{2.208812in}}%
\pgfpathlineto{\pgfqpoint{2.639407in}{2.187519in}}%
\pgfpathlineto{\pgfqpoint{2.592317in}{2.233025in}}%
\pgfpathlineto{\pgfqpoint{2.559243in}{2.254297in}}%
\pgfpathclose%
\pgfusepath{stroke,fill}%
\end{pgfscope}%
\begin{pgfscope}%
\pgfpathrectangle{\pgfqpoint{0.050000in}{0.050000in}}{\pgfqpoint{4.900000in}{4.900000in}}%
\pgfusepath{clip}%
\pgfsetbuttcap%
\pgfsetroundjoin%
\definecolor{currentfill}{rgb}{0.223299,0.223299,0.223299}%
\pgfsetfillcolor{currentfill}%
\pgfsetfillopacity{0.900000}%
\pgfsetlinewidth{0.501875pt}%
\definecolor{currentstroke}{rgb}{0.200000,0.200000,0.200000}%
\pgfsetstrokecolor{currentstroke}%
\pgfsetdash{}{0pt}%
\pgfpathmoveto{\pgfqpoint{1.544726in}{3.074870in}}%
\pgfpathlineto{\pgfqpoint{1.594037in}{3.027550in}}%
\pgfpathlineto{\pgfqpoint{1.625498in}{3.007550in}}%
\pgfpathlineto{\pgfqpoint{1.576178in}{3.054902in}}%
\pgfpathlineto{\pgfqpoint{1.544726in}{3.074870in}}%
\pgfpathclose%
\pgfusepath{stroke,fill}%
\end{pgfscope}%
\begin{pgfscope}%
\pgfpathrectangle{\pgfqpoint{0.050000in}{0.050000in}}{\pgfqpoint{4.900000in}{4.900000in}}%
\pgfusepath{clip}%
\pgfsetbuttcap%
\pgfsetroundjoin%
\definecolor{currentfill}{rgb}{0.495656,0.495656,0.495656}%
\pgfsetfillcolor{currentfill}%
\pgfsetfillopacity{0.900000}%
\pgfsetlinewidth{0.501875pt}%
\definecolor{currentstroke}{rgb}{0.200000,0.200000,0.200000}%
\pgfsetstrokecolor{currentstroke}%
\pgfsetdash{}{0pt}%
\pgfpathmoveto{\pgfqpoint{2.285540in}{2.476152in}}%
\pgfpathlineto{\pgfqpoint{2.333234in}{2.430138in}}%
\pgfpathlineto{\pgfqpoint{2.365879in}{2.409184in}}%
\pgfpathlineto{\pgfqpoint{2.318181in}{2.455224in}}%
\pgfpathlineto{\pgfqpoint{2.285540in}{2.476152in}}%
\pgfpathclose%
\pgfusepath{stroke,fill}%
\end{pgfscope}%
\begin{pgfscope}%
\pgfpathrectangle{\pgfqpoint{0.050000in}{0.050000in}}{\pgfqpoint{4.900000in}{4.900000in}}%
\pgfusepath{clip}%
\pgfsetbuttcap%
\pgfsetroundjoin%
\definecolor{currentfill}{rgb}{0.104698,0.104698,0.104698}%
\pgfsetfillcolor{currentfill}%
\pgfsetfillopacity{0.900000}%
\pgfsetlinewidth{0.501875pt}%
\definecolor{currentstroke}{rgb}{0.200000,0.200000,0.200000}%
\pgfsetstrokecolor{currentstroke}%
\pgfsetdash{}{0pt}%
\pgfpathmoveto{\pgfqpoint{1.270818in}{3.296926in}}%
\pgfpathlineto{\pgfqpoint{1.320736in}{3.249112in}}%
\pgfpathlineto{\pgfqpoint{1.351766in}{3.229462in}}%
\pgfpathlineto{\pgfqpoint{1.301836in}{3.277310in}}%
\pgfpathlineto{\pgfqpoint{1.270818in}{3.296926in}}%
\pgfpathclose%
\pgfusepath{stroke,fill}%
\end{pgfscope}%
\begin{pgfscope}%
\pgfpathrectangle{\pgfqpoint{0.050000in}{0.050000in}}{\pgfqpoint{4.900000in}{4.900000in}}%
\pgfusepath{clip}%
\pgfsetbuttcap%
\pgfsetroundjoin%
\definecolor{currentfill}{rgb}{0.767443,0.767443,0.767443}%
\pgfsetfillcolor{currentfill}%
\pgfsetfillopacity{0.900000}%
\pgfsetlinewidth{0.501875pt}%
\definecolor{currentstroke}{rgb}{0.200000,0.200000,0.200000}%
\pgfsetstrokecolor{currentstroke}%
\pgfsetdash{}{0pt}%
\pgfpathmoveto{\pgfqpoint{3.026474in}{1.879483in}}%
\pgfpathlineto{\pgfqpoint{3.072565in}{1.836186in}}%
\pgfpathlineto{\pgfqpoint{3.106400in}{1.814543in}}%
\pgfpathlineto{\pgfqpoint{3.060308in}{1.857778in}}%
\pgfpathlineto{\pgfqpoint{3.026474in}{1.879483in}}%
\pgfpathclose%
\pgfusepath{stroke,fill}%
\end{pgfscope}%
\begin{pgfscope}%
\pgfpathrectangle{\pgfqpoint{0.050000in}{0.050000in}}{\pgfqpoint{4.900000in}{4.900000in}}%
\pgfusepath{clip}%
\pgfsetbuttcap%
\pgfsetroundjoin%
\definecolor{currentfill}{rgb}{0.839785,0.839785,0.839785}%
\pgfsetfillcolor{currentfill}%
\pgfsetfillopacity{0.900000}%
\pgfsetlinewidth{0.501875pt}%
\definecolor{currentstroke}{rgb}{0.200000,0.200000,0.200000}%
\pgfsetstrokecolor{currentstroke}%
\pgfsetdash{}{0pt}%
\pgfpathmoveto{\pgfqpoint{3.300194in}{1.666520in}}%
\pgfpathlineto{\pgfqpoint{3.345744in}{1.627392in}}%
\pgfpathlineto{\pgfqpoint{3.380026in}{1.605855in}}%
\pgfpathlineto{\pgfqpoint{3.334472in}{1.644843in}}%
\pgfpathlineto{\pgfqpoint{3.300194in}{1.666520in}}%
\pgfpathclose%
\pgfusepath{stroke,fill}%
\end{pgfscope}%
\begin{pgfscope}%
\pgfpathrectangle{\pgfqpoint{0.050000in}{0.050000in}}{\pgfqpoint{4.900000in}{4.900000in}}%
\pgfusepath{clip}%
\pgfsetbuttcap%
\pgfsetroundjoin%
\definecolor{currentfill}{rgb}{0.403783,0.403783,0.403783}%
\pgfsetfillcolor{currentfill}%
\pgfsetfillopacity{0.900000}%
\pgfsetlinewidth{0.501875pt}%
\definecolor{currentstroke}{rgb}{0.200000,0.200000,0.200000}%
\pgfsetstrokecolor{currentstroke}%
\pgfsetdash{}{0pt}%
\pgfpathmoveto{\pgfqpoint{2.011741in}{2.698111in}}%
\pgfpathlineto{\pgfqpoint{2.060042in}{2.651603in}}%
\pgfpathlineto{\pgfqpoint{2.092257in}{2.630998in}}%
\pgfpathlineto{\pgfqpoint{2.043949in}{2.677535in}}%
\pgfpathlineto{\pgfqpoint{2.011741in}{2.698111in}}%
\pgfpathclose%
\pgfusepath{stroke,fill}%
\end{pgfscope}%
\begin{pgfscope}%
\pgfpathrectangle{\pgfqpoint{0.050000in}{0.050000in}}{\pgfqpoint{4.900000in}{4.900000in}}%
\pgfusepath{clip}%
\pgfsetbuttcap%
\pgfsetroundjoin%
\definecolor{currentfill}{rgb}{0.667405,0.667405,0.667405}%
\pgfsetfillcolor{currentfill}%
\pgfsetfillopacity{0.900000}%
\pgfsetlinewidth{0.501875pt}%
\definecolor{currentstroke}{rgb}{0.200000,0.200000,0.200000}%
\pgfsetstrokecolor{currentstroke}%
\pgfsetdash{}{0pt}%
\pgfpathmoveto{\pgfqpoint{2.752768in}{2.099447in}}%
\pgfpathlineto{\pgfqpoint{2.799453in}{2.054459in}}%
\pgfpathlineto{\pgfqpoint{2.832853in}{2.032953in}}%
\pgfpathlineto{\pgfqpoint{2.786167in}{2.077942in}}%
\pgfpathlineto{\pgfqpoint{2.752768in}{2.099447in}}%
\pgfpathclose%
\pgfusepath{stroke,fill}%
\end{pgfscope}%
\begin{pgfscope}%
\pgfpathrectangle{\pgfqpoint{0.050000in}{0.050000in}}{\pgfqpoint{4.900000in}{4.900000in}}%
\pgfusepath{clip}%
\pgfsetbuttcap%
\pgfsetroundjoin%
\definecolor{currentfill}{rgb}{0.004552,0.004552,0.004552}%
\pgfsetfillcolor{currentfill}%
\pgfsetfillopacity{0.900000}%
\pgfsetlinewidth{0.501875pt}%
\definecolor{currentstroke}{rgb}{0.200000,0.200000,0.200000}%
\pgfsetstrokecolor{currentstroke}%
\pgfsetdash{}{0pt}%
\pgfpathmoveto{\pgfqpoint{0.996814in}{3.519062in}}%
\pgfpathlineto{\pgfqpoint{1.047340in}{3.470753in}}%
\pgfpathlineto{\pgfqpoint{1.077938in}{3.451454in}}%
\pgfpathlineto{\pgfqpoint{1.027398in}{3.499798in}}%
\pgfpathlineto{\pgfqpoint{0.996814in}{3.519062in}}%
\pgfpathclose%
\pgfusepath{stroke,fill}%
\end{pgfscope}%
\begin{pgfscope}%
\pgfpathrectangle{\pgfqpoint{0.050000in}{0.050000in}}{\pgfqpoint{4.900000in}{4.900000in}}%
\pgfusepath{clip}%
\pgfsetbuttcap%
\pgfsetroundjoin%
\definecolor{currentfill}{rgb}{0.306344,0.306344,0.306344}%
\pgfsetfillcolor{currentfill}%
\pgfsetfillopacity{0.900000}%
\pgfsetlinewidth{0.501875pt}%
\definecolor{currentstroke}{rgb}{0.200000,0.200000,0.200000}%
\pgfsetstrokecolor{currentstroke}%
\pgfsetdash{}{0pt}%
\pgfpathmoveto{\pgfqpoint{1.737848in}{2.920151in}}%
\pgfpathlineto{\pgfqpoint{1.786756in}{2.873149in}}%
\pgfpathlineto{\pgfqpoint{1.818539in}{2.852895in}}%
\pgfpathlineto{\pgfqpoint{1.769622in}{2.899927in}}%
\pgfpathlineto{\pgfqpoint{1.737848in}{2.920151in}}%
\pgfpathclose%
\pgfusepath{stroke,fill}%
\end{pgfscope}%
\begin{pgfscope}%
\pgfpathrectangle{\pgfqpoint{0.050000in}{0.050000in}}{\pgfqpoint{4.900000in}{4.900000in}}%
\pgfusepath{clip}%
\pgfsetbuttcap%
\pgfsetroundjoin%
\definecolor{currentfill}{rgb}{0.564552,0.564552,0.564552}%
\pgfsetfillcolor{currentfill}%
\pgfsetfillopacity{0.900000}%
\pgfsetlinewidth{0.501875pt}%
\definecolor{currentstroke}{rgb}{0.200000,0.200000,0.200000}%
\pgfsetstrokecolor{currentstroke}%
\pgfsetdash{}{0pt}%
\pgfpathmoveto{\pgfqpoint{2.478984in}{2.321184in}}%
\pgfpathlineto{\pgfqpoint{2.526274in}{2.275503in}}%
\pgfpathlineto{\pgfqpoint{2.559243in}{2.254297in}}%
\pgfpathlineto{\pgfqpoint{2.511949in}{2.300002in}}%
\pgfpathlineto{\pgfqpoint{2.478984in}{2.321184in}}%
\pgfpathclose%
\pgfusepath{stroke,fill}%
\end{pgfscope}%
\begin{pgfscope}%
\pgfpathrectangle{\pgfqpoint{0.050000in}{0.050000in}}{\pgfqpoint{4.900000in}{4.900000in}}%
\pgfusepath{clip}%
\pgfsetbuttcap%
\pgfsetroundjoin%
\definecolor{currentfill}{rgb}{0.179008,0.179008,0.179008}%
\pgfsetfillcolor{currentfill}%
\pgfsetfillopacity{0.900000}%
\pgfsetlinewidth{0.501875pt}%
\definecolor{currentstroke}{rgb}{0.200000,0.200000,0.200000}%
\pgfsetstrokecolor{currentstroke}%
\pgfsetdash{}{0pt}%
\pgfpathmoveto{\pgfqpoint{1.463858in}{3.142270in}}%
\pgfpathlineto{\pgfqpoint{1.513375in}{3.094775in}}%
\pgfpathlineto{\pgfqpoint{1.544726in}{3.074870in}}%
\pgfpathlineto{\pgfqpoint{1.495199in}{3.122398in}}%
\pgfpathlineto{\pgfqpoint{1.463858in}{3.142270in}}%
\pgfpathclose%
\pgfusepath{stroke,fill}%
\end{pgfscope}%
\begin{pgfscope}%
\pgfpathrectangle{\pgfqpoint{0.050000in}{0.050000in}}{\pgfqpoint{4.900000in}{4.900000in}}%
\pgfusepath{clip}%
\pgfsetbuttcap%
\pgfsetroundjoin%
\definecolor{currentfill}{rgb}{0.465513,0.465513,0.465513}%
\pgfsetfillcolor{currentfill}%
\pgfsetfillopacity{0.900000}%
\pgfsetlinewidth{0.501875pt}%
\definecolor{currentstroke}{rgb}{0.200000,0.200000,0.200000}%
\pgfsetstrokecolor{currentstroke}%
\pgfsetdash{}{0pt}%
\pgfpathmoveto{\pgfqpoint{2.205105in}{2.543200in}}%
\pgfpathlineto{\pgfqpoint{2.253003in}{2.497012in}}%
\pgfpathlineto{\pgfqpoint{2.285540in}{2.476152in}}%
\pgfpathlineto{\pgfqpoint{2.237636in}{2.522367in}}%
\pgfpathlineto{\pgfqpoint{2.205105in}{2.543200in}}%
\pgfpathclose%
\pgfusepath{stroke,fill}%
\end{pgfscope}%
\begin{pgfscope}%
\pgfpathrectangle{\pgfqpoint{0.050000in}{0.050000in}}{\pgfqpoint{4.900000in}{4.900000in}}%
\pgfusepath{clip}%
\pgfsetbuttcap%
\pgfsetroundjoin%
\definecolor{currentfill}{rgb}{0.815671,0.815671,0.815671}%
\pgfsetfillcolor{currentfill}%
\pgfsetfillopacity{0.900000}%
\pgfsetlinewidth{0.501875pt}%
\definecolor{currentstroke}{rgb}{0.200000,0.200000,0.200000}%
\pgfsetstrokecolor{currentstroke}%
\pgfsetdash{}{0pt}%
\pgfpathmoveto{\pgfqpoint{3.220304in}{1.728973in}}%
\pgfpathlineto{\pgfqpoint{3.266026in}{1.688137in}}%
\pgfpathlineto{\pgfqpoint{3.300194in}{1.666520in}}%
\pgfpathlineto{\pgfqpoint{3.254469in}{1.707240in}}%
\pgfpathlineto{\pgfqpoint{3.220304in}{1.728973in}}%
\pgfpathclose%
\pgfusepath{stroke,fill}%
\end{pgfscope}%
\begin{pgfscope}%
\pgfpathrectangle{\pgfqpoint{0.050000in}{0.050000in}}{\pgfqpoint{4.900000in}{4.900000in}}%
\pgfusepath{clip}%
\pgfsetbuttcap%
\pgfsetroundjoin%
\definecolor{currentfill}{rgb}{0.739377,0.739377,0.739377}%
\pgfsetfillcolor{currentfill}%
\pgfsetfillopacity{0.900000}%
\pgfsetlinewidth{0.501875pt}%
\definecolor{currentstroke}{rgb}{0.200000,0.200000,0.200000}%
\pgfsetstrokecolor{currentstroke}%
\pgfsetdash{}{0pt}%
\pgfpathmoveto{\pgfqpoint{2.946462in}{1.945173in}}%
\pgfpathlineto{\pgfqpoint{2.992748in}{1.901127in}}%
\pgfpathlineto{\pgfqpoint{3.026474in}{1.879483in}}%
\pgfpathlineto{\pgfqpoint{2.980186in}{1.923492in}}%
\pgfpathlineto{\pgfqpoint{2.946462in}{1.945173in}}%
\pgfpathclose%
\pgfusepath{stroke,fill}%
\end{pgfscope}%
\begin{pgfscope}%
\pgfpathrectangle{\pgfqpoint{0.050000in}{0.050000in}}{\pgfqpoint{4.900000in}{4.900000in}}%
\pgfusepath{clip}%
\pgfsetbuttcap%
\pgfsetroundjoin%
\definecolor{currentfill}{rgb}{0.072834,0.072834,0.072834}%
\pgfsetfillcolor{currentfill}%
\pgfsetfillopacity{0.900000}%
\pgfsetlinewidth{0.501875pt}%
\definecolor{currentstroke}{rgb}{0.200000,0.200000,0.200000}%
\pgfsetstrokecolor{currentstroke}%
\pgfsetdash{}{0pt}%
\pgfpathmoveto{\pgfqpoint{1.189774in}{3.364470in}}%
\pgfpathlineto{\pgfqpoint{1.239899in}{3.316480in}}%
\pgfpathlineto{\pgfqpoint{1.270818in}{3.296926in}}%
\pgfpathlineto{\pgfqpoint{1.220680in}{3.344950in}}%
\pgfpathlineto{\pgfqpoint{1.189774in}{3.364470in}}%
\pgfpathclose%
\pgfusepath{stroke,fill}%
\end{pgfscope}%
\begin{pgfscope}%
\pgfpathrectangle{\pgfqpoint{0.050000in}{0.050000in}}{\pgfqpoint{4.900000in}{4.900000in}}%
\pgfusepath{clip}%
\pgfsetbuttcap%
\pgfsetroundjoin%
\definecolor{currentfill}{rgb}{0.371303,0.371303,0.371303}%
\pgfsetfillcolor{currentfill}%
\pgfsetfillopacity{0.900000}%
\pgfsetlinewidth{0.501875pt}%
\definecolor{currentstroke}{rgb}{0.200000,0.200000,0.200000}%
\pgfsetstrokecolor{currentstroke}%
\pgfsetdash{}{0pt}%
\pgfpathmoveto{\pgfqpoint{1.931130in}{2.765303in}}%
\pgfpathlineto{\pgfqpoint{1.979636in}{2.718621in}}%
\pgfpathlineto{\pgfqpoint{2.011741in}{2.698111in}}%
\pgfpathlineto{\pgfqpoint{1.963228in}{2.744823in}}%
\pgfpathlineto{\pgfqpoint{1.931130in}{2.765303in}}%
\pgfpathclose%
\pgfusepath{stroke,fill}%
\end{pgfscope}%
\begin{pgfscope}%
\pgfpathrectangle{\pgfqpoint{0.050000in}{0.050000in}}{\pgfqpoint{4.900000in}{4.900000in}}%
\pgfusepath{clip}%
\pgfsetbuttcap%
\pgfsetroundjoin%
\definecolor{currentfill}{rgb}{0.633818,0.633818,0.633818}%
\pgfsetfillcolor{currentfill}%
\pgfsetfillopacity{0.900000}%
\pgfsetlinewidth{0.501875pt}%
\definecolor{currentstroke}{rgb}{0.200000,0.200000,0.200000}%
\pgfsetstrokecolor{currentstroke}%
\pgfsetdash{}{0pt}%
\pgfpathmoveto{\pgfqpoint{2.672590in}{2.166160in}}%
\pgfpathlineto{\pgfqpoint{2.719477in}{2.120885in}}%
\pgfpathlineto{\pgfqpoint{2.752768in}{2.099447in}}%
\pgfpathlineto{\pgfqpoint{2.705879in}{2.144734in}}%
\pgfpathlineto{\pgfqpoint{2.672590in}{2.166160in}}%
\pgfpathclose%
\pgfusepath{stroke,fill}%
\end{pgfscope}%
\begin{pgfscope}%
\pgfpathrectangle{\pgfqpoint{0.050000in}{0.050000in}}{\pgfqpoint{4.900000in}{4.900000in}}%
\pgfusepath{clip}%
\pgfsetbuttcap%
\pgfsetroundjoin%
\definecolor{currentfill}{rgb}{0.267589,0.267589,0.267589}%
\pgfsetfillcolor{currentfill}%
\pgfsetfillopacity{0.900000}%
\pgfsetlinewidth{0.501875pt}%
\definecolor{currentstroke}{rgb}{0.200000,0.200000,0.200000}%
\pgfsetstrokecolor{currentstroke}%
\pgfsetdash{}{0pt}%
\pgfpathmoveto{\pgfqpoint{1.657060in}{2.987487in}}%
\pgfpathlineto{\pgfqpoint{1.706175in}{2.940310in}}%
\pgfpathlineto{\pgfqpoint{1.737848in}{2.920151in}}%
\pgfpathlineto{\pgfqpoint{1.688724in}{2.967359in}}%
\pgfpathlineto{\pgfqpoint{1.657060in}{2.987487in}}%
\pgfpathclose%
\pgfusepath{stroke,fill}%
\end{pgfscope}%
\begin{pgfscope}%
\pgfpathrectangle{\pgfqpoint{0.050000in}{0.050000in}}{\pgfqpoint{4.900000in}{4.900000in}}%
\pgfusepath{clip}%
\pgfsetbuttcap%
\pgfsetroundjoin%
\definecolor{currentfill}{rgb}{0.530104,0.530104,0.530104}%
\pgfsetfillcolor{currentfill}%
\pgfsetfillopacity{0.900000}%
\pgfsetlinewidth{0.501875pt}%
\definecolor{currentstroke}{rgb}{0.200000,0.200000,0.200000}%
\pgfsetstrokecolor{currentstroke}%
\pgfsetdash{}{0pt}%
\pgfpathmoveto{\pgfqpoint{2.398630in}{2.388162in}}%
\pgfpathlineto{\pgfqpoint{2.446124in}{2.342299in}}%
\pgfpathlineto{\pgfqpoint{2.478984in}{2.321184in}}%
\pgfpathlineto{\pgfqpoint{2.431486in}{2.367074in}}%
\pgfpathlineto{\pgfqpoint{2.398630in}{2.388162in}}%
\pgfpathclose%
\pgfusepath{stroke,fill}%
\end{pgfscope}%
\begin{pgfscope}%
\pgfpathrectangle{\pgfqpoint{0.050000in}{0.050000in}}{\pgfqpoint{4.900000in}{4.900000in}}%
\pgfusepath{clip}%
\pgfsetbuttcap%
\pgfsetroundjoin%
\definecolor{currentfill}{rgb}{0.141115,0.141115,0.141115}%
\pgfsetfillcolor{currentfill}%
\pgfsetfillopacity{0.900000}%
\pgfsetlinewidth{0.501875pt}%
\definecolor{currentstroke}{rgb}{0.200000,0.200000,0.200000}%
\pgfsetstrokecolor{currentstroke}%
\pgfsetdash{}{0pt}%
\pgfpathmoveto{\pgfqpoint{1.382895in}{3.209750in}}%
\pgfpathlineto{\pgfqpoint{1.432618in}{3.162079in}}%
\pgfpathlineto{\pgfqpoint{1.463858in}{3.142270in}}%
\pgfpathlineto{\pgfqpoint{1.414124in}{3.189975in}}%
\pgfpathlineto{\pgfqpoint{1.382895in}{3.209750in}}%
\pgfpathclose%
\pgfusepath{stroke,fill}%
\end{pgfscope}%
\begin{pgfscope}%
\pgfpathrectangle{\pgfqpoint{0.050000in}{0.050000in}}{\pgfqpoint{4.900000in}{4.900000in}}%
\pgfusepath{clip}%
\pgfsetbuttcap%
\pgfsetroundjoin%
\definecolor{currentfill}{rgb}{0.791557,0.791557,0.791557}%
\pgfsetfillcolor{currentfill}%
\pgfsetfillopacity{0.900000}%
\pgfsetlinewidth{0.501875pt}%
\definecolor{currentstroke}{rgb}{0.200000,0.200000,0.200000}%
\pgfsetstrokecolor{currentstroke}%
\pgfsetdash{}{0pt}%
\pgfpathmoveto{\pgfqpoint{3.140345in}{1.792840in}}%
\pgfpathlineto{\pgfqpoint{3.186248in}{1.750647in}}%
\pgfpathlineto{\pgfqpoint{3.220304in}{1.728973in}}%
\pgfpathlineto{\pgfqpoint{3.174398in}{1.771076in}}%
\pgfpathlineto{\pgfqpoint{3.140345in}{1.792840in}}%
\pgfpathclose%
\pgfusepath{stroke,fill}%
\end{pgfscope}%
\begin{pgfscope}%
\pgfpathrectangle{\pgfqpoint{0.050000in}{0.050000in}}{\pgfqpoint{4.900000in}{4.900000in}}%
\pgfusepath{clip}%
\pgfsetbuttcap%
\pgfsetroundjoin%
\definecolor{currentfill}{rgb}{0.432203,0.432203,0.432203}%
\pgfsetfillcolor{currentfill}%
\pgfsetfillopacity{0.900000}%
\pgfsetlinewidth{0.501875pt}%
\definecolor{currentstroke}{rgb}{0.200000,0.200000,0.200000}%
\pgfsetstrokecolor{currentstroke}%
\pgfsetdash{}{0pt}%
\pgfpathmoveto{\pgfqpoint{2.124575in}{2.610328in}}%
\pgfpathlineto{\pgfqpoint{2.172677in}{2.563965in}}%
\pgfpathlineto{\pgfqpoint{2.205105in}{2.543200in}}%
\pgfpathlineto{\pgfqpoint{2.156996in}{2.589591in}}%
\pgfpathlineto{\pgfqpoint{2.124575in}{2.610328in}}%
\pgfpathclose%
\pgfusepath{stroke,fill}%
\end{pgfscope}%
\begin{pgfscope}%
\pgfpathrectangle{\pgfqpoint{0.050000in}{0.050000in}}{\pgfqpoint{4.900000in}{4.900000in}}%
\pgfusepath{clip}%
\pgfsetbuttcap%
\pgfsetroundjoin%
\definecolor{currentfill}{rgb}{0.705790,0.705790,0.705790}%
\pgfsetfillcolor{currentfill}%
\pgfsetfillopacity{0.900000}%
\pgfsetlinewidth{0.501875pt}%
\definecolor{currentstroke}{rgb}{0.200000,0.200000,0.200000}%
\pgfsetstrokecolor{currentstroke}%
\pgfsetdash{}{0pt}%
\pgfpathmoveto{\pgfqpoint{2.866361in}{2.011383in}}%
\pgfpathlineto{\pgfqpoint{2.912846in}{1.966792in}}%
\pgfpathlineto{\pgfqpoint{2.946462in}{1.945173in}}%
\pgfpathlineto{\pgfqpoint{2.899975in}{1.989749in}}%
\pgfpathlineto{\pgfqpoint{2.866361in}{2.011383in}}%
\pgfpathclose%
\pgfusepath{stroke,fill}%
\end{pgfscope}%
\begin{pgfscope}%
\pgfpathrectangle{\pgfqpoint{0.050000in}{0.050000in}}{\pgfqpoint{4.900000in}{4.900000in}}%
\pgfusepath{clip}%
\pgfsetbuttcap%
\pgfsetroundjoin%
\definecolor{currentfill}{rgb}{0.036417,0.036417,0.036417}%
\pgfsetfillcolor{currentfill}%
\pgfsetfillopacity{0.900000}%
\pgfsetlinewidth{0.501875pt}%
\definecolor{currentstroke}{rgb}{0.200000,0.200000,0.200000}%
\pgfsetstrokecolor{currentstroke}%
\pgfsetdash{}{0pt}%
\pgfpathmoveto{\pgfqpoint{1.108634in}{3.432093in}}%
\pgfpathlineto{\pgfqpoint{1.158966in}{3.383927in}}%
\pgfpathlineto{\pgfqpoint{1.189774in}{3.364470in}}%
\pgfpathlineto{\pgfqpoint{1.139428in}{3.412671in}}%
\pgfpathlineto{\pgfqpoint{1.108634in}{3.432093in}}%
\pgfpathclose%
\pgfusepath{stroke,fill}%
\end{pgfscope}%
\begin{pgfscope}%
\pgfpathrectangle{\pgfqpoint{0.050000in}{0.050000in}}{\pgfqpoint{4.900000in}{4.900000in}}%
\pgfusepath{clip}%
\pgfsetbuttcap%
\pgfsetroundjoin%
\definecolor{currentfill}{rgb}{0.858762,0.858762,0.858762}%
\pgfsetfillcolor{currentfill}%
\pgfsetfillopacity{0.900000}%
\pgfsetlinewidth{0.501875pt}%
\definecolor{currentstroke}{rgb}{0.200000,0.200000,0.200000}%
\pgfsetstrokecolor{currentstroke}%
\pgfsetdash{}{0pt}%
\pgfpathmoveto{\pgfqpoint{3.414419in}{1.584257in}}%
\pgfpathlineto{\pgfqpoint{3.459814in}{1.547364in}}%
\pgfpathlineto{\pgfqpoint{3.494323in}{1.525859in}}%
\pgfpathlineto{\pgfqpoint{3.448923in}{1.562596in}}%
\pgfpathlineto{\pgfqpoint{3.414419in}{1.584257in}}%
\pgfpathclose%
\pgfusepath{stroke,fill}%
\end{pgfscope}%
\begin{pgfscope}%
\pgfpathrectangle{\pgfqpoint{0.050000in}{0.050000in}}{\pgfqpoint{4.900000in}{4.900000in}}%
\pgfusepath{clip}%
\pgfsetbuttcap%
\pgfsetroundjoin%
\definecolor{currentfill}{rgb}{0.342884,0.342884,0.342884}%
\pgfsetfillcolor{currentfill}%
\pgfsetfillopacity{0.900000}%
\pgfsetlinewidth{0.501875pt}%
\definecolor{currentstroke}{rgb}{0.200000,0.200000,0.200000}%
\pgfsetstrokecolor{currentstroke}%
\pgfsetdash{}{0pt}%
\pgfpathmoveto{\pgfqpoint{1.850424in}{2.832575in}}%
\pgfpathlineto{\pgfqpoint{1.899135in}{2.785718in}}%
\pgfpathlineto{\pgfqpoint{1.931130in}{2.765303in}}%
\pgfpathlineto{\pgfqpoint{1.882411in}{2.812190in}}%
\pgfpathlineto{\pgfqpoint{1.850424in}{2.832575in}}%
\pgfpathclose%
\pgfusepath{stroke,fill}%
\end{pgfscope}%
\begin{pgfscope}%
\pgfpathrectangle{\pgfqpoint{0.050000in}{0.050000in}}{\pgfqpoint{4.900000in}{4.900000in}}%
\pgfusepath{clip}%
\pgfsetbuttcap%
\pgfsetroundjoin%
\definecolor{currentfill}{rgb}{0.595433,0.595433,0.595433}%
\pgfsetfillcolor{currentfill}%
\pgfsetfillopacity{0.900000}%
\pgfsetlinewidth{0.501875pt}%
\definecolor{currentstroke}{rgb}{0.200000,0.200000,0.200000}%
\pgfsetstrokecolor{currentstroke}%
\pgfsetdash{}{0pt}%
\pgfpathmoveto{\pgfqpoint{2.592317in}{2.233025in}}%
\pgfpathlineto{\pgfqpoint{2.639407in}{2.187519in}}%
\pgfpathlineto{\pgfqpoint{2.672590in}{2.166160in}}%
\pgfpathlineto{\pgfqpoint{2.625497in}{2.211684in}}%
\pgfpathlineto{\pgfqpoint{2.592317in}{2.233025in}}%
\pgfpathclose%
\pgfusepath{stroke,fill}%
\end{pgfscope}%
\begin{pgfscope}%
\pgfpathrectangle{\pgfqpoint{0.050000in}{0.050000in}}{\pgfqpoint{4.900000in}{4.900000in}}%
\pgfusepath{clip}%
\pgfsetbuttcap%
\pgfsetroundjoin%
\definecolor{currentfill}{rgb}{0.223299,0.223299,0.223299}%
\pgfsetfillcolor{currentfill}%
\pgfsetfillopacity{0.900000}%
\pgfsetlinewidth{0.501875pt}%
\definecolor{currentstroke}{rgb}{0.200000,0.200000,0.200000}%
\pgfsetstrokecolor{currentstroke}%
\pgfsetdash{}{0pt}%
\pgfpathmoveto{\pgfqpoint{1.576178in}{3.054902in}}%
\pgfpathlineto{\pgfqpoint{1.625498in}{3.007550in}}%
\pgfpathlineto{\pgfqpoint{1.657060in}{2.987487in}}%
\pgfpathlineto{\pgfqpoint{1.607730in}{3.034871in}}%
\pgfpathlineto{\pgfqpoint{1.576178in}{3.054902in}}%
\pgfpathclose%
\pgfusepath{stroke,fill}%
\end{pgfscope}%
\begin{pgfscope}%
\pgfpathrectangle{\pgfqpoint{0.050000in}{0.050000in}}{\pgfqpoint{4.900000in}{4.900000in}}%
\pgfusepath{clip}%
\pgfsetbuttcap%
\pgfsetroundjoin%
\definecolor{currentfill}{rgb}{0.499962,0.499962,0.499962}%
\pgfsetfillcolor{currentfill}%
\pgfsetfillopacity{0.900000}%
\pgfsetlinewidth{0.501875pt}%
\definecolor{currentstroke}{rgb}{0.200000,0.200000,0.200000}%
\pgfsetstrokecolor{currentstroke}%
\pgfsetdash{}{0pt}%
\pgfpathmoveto{\pgfqpoint{2.318181in}{2.455224in}}%
\pgfpathlineto{\pgfqpoint{2.365879in}{2.409184in}}%
\pgfpathlineto{\pgfqpoint{2.398630in}{2.388162in}}%
\pgfpathlineto{\pgfqpoint{2.350927in}{2.434230in}}%
\pgfpathlineto{\pgfqpoint{2.318181in}{2.455224in}}%
\pgfpathclose%
\pgfusepath{stroke,fill}%
\end{pgfscope}%
\begin{pgfscope}%
\pgfpathrectangle{\pgfqpoint{0.050000in}{0.050000in}}{\pgfqpoint{4.900000in}{4.900000in}}%
\pgfusepath{clip}%
\pgfsetbuttcap%
\pgfsetroundjoin%
\definecolor{currentfill}{rgb}{0.104698,0.104698,0.104698}%
\pgfsetfillcolor{currentfill}%
\pgfsetfillopacity{0.900000}%
\pgfsetlinewidth{0.501875pt}%
\definecolor{currentstroke}{rgb}{0.200000,0.200000,0.200000}%
\pgfsetstrokecolor{currentstroke}%
\pgfsetdash{}{0pt}%
\pgfpathmoveto{\pgfqpoint{1.301836in}{3.277310in}}%
\pgfpathlineto{\pgfqpoint{1.351766in}{3.229462in}}%
\pgfpathlineto{\pgfqpoint{1.382895in}{3.209750in}}%
\pgfpathlineto{\pgfqpoint{1.332953in}{3.257631in}}%
\pgfpathlineto{\pgfqpoint{1.301836in}{3.277310in}}%
\pgfpathclose%
\pgfusepath{stroke,fill}%
\end{pgfscope}%
\begin{pgfscope}%
\pgfpathrectangle{\pgfqpoint{0.050000in}{0.050000in}}{\pgfqpoint{4.900000in}{4.900000in}}%
\pgfusepath{clip}%
\pgfsetbuttcap%
\pgfsetroundjoin%
\definecolor{currentfill}{rgb}{0.767443,0.767443,0.767443}%
\pgfsetfillcolor{currentfill}%
\pgfsetfillopacity{0.900000}%
\pgfsetlinewidth{0.501875pt}%
\definecolor{currentstroke}{rgb}{0.200000,0.200000,0.200000}%
\pgfsetstrokecolor{currentstroke}%
\pgfsetdash{}{0pt}%
\pgfpathmoveto{\pgfqpoint{3.060308in}{1.857778in}}%
\pgfpathlineto{\pgfqpoint{3.106400in}{1.814543in}}%
\pgfpathlineto{\pgfqpoint{3.140345in}{1.792840in}}%
\pgfpathlineto{\pgfqpoint{3.094251in}{1.836010in}}%
\pgfpathlineto{\pgfqpoint{3.060308in}{1.857778in}}%
\pgfpathclose%
\pgfusepath{stroke,fill}%
\end{pgfscope}%
\begin{pgfscope}%
\pgfpathrectangle{\pgfqpoint{0.050000in}{0.050000in}}{\pgfqpoint{4.900000in}{4.900000in}}%
\pgfusepath{clip}%
\pgfsetbuttcap%
\pgfsetroundjoin%
\definecolor{currentfill}{rgb}{0.839785,0.839785,0.839785}%
\pgfsetfillcolor{currentfill}%
\pgfsetfillopacity{0.900000}%
\pgfsetlinewidth{0.501875pt}%
\definecolor{currentstroke}{rgb}{0.200000,0.200000,0.200000}%
\pgfsetstrokecolor{currentstroke}%
\pgfsetdash{}{0pt}%
\pgfpathmoveto{\pgfqpoint{3.334472in}{1.644843in}}%
\pgfpathlineto{\pgfqpoint{3.380026in}{1.605855in}}%
\pgfpathlineto{\pgfqpoint{3.414419in}{1.584257in}}%
\pgfpathlineto{\pgfqpoint{3.368861in}{1.623105in}}%
\pgfpathlineto{\pgfqpoint{3.334472in}{1.644843in}}%
\pgfpathclose%
\pgfusepath{stroke,fill}%
\end{pgfscope}%
\begin{pgfscope}%
\pgfpathrectangle{\pgfqpoint{0.050000in}{0.050000in}}{\pgfqpoint{4.900000in}{4.900000in}}%
\pgfusepath{clip}%
\pgfsetbuttcap%
\pgfsetroundjoin%
\definecolor{currentfill}{rgb}{0.403783,0.403783,0.403783}%
\pgfsetfillcolor{currentfill}%
\pgfsetfillopacity{0.900000}%
\pgfsetlinewidth{0.501875pt}%
\definecolor{currentstroke}{rgb}{0.200000,0.200000,0.200000}%
\pgfsetstrokecolor{currentstroke}%
\pgfsetdash{}{0pt}%
\pgfpathmoveto{\pgfqpoint{2.043949in}{2.677535in}}%
\pgfpathlineto{\pgfqpoint{2.092257in}{2.630998in}}%
\pgfpathlineto{\pgfqpoint{2.124575in}{2.610328in}}%
\pgfpathlineto{\pgfqpoint{2.076261in}{2.656894in}}%
\pgfpathlineto{\pgfqpoint{2.043949in}{2.677535in}}%
\pgfpathclose%
\pgfusepath{stroke,fill}%
\end{pgfscope}%
\begin{pgfscope}%
\pgfpathrectangle{\pgfqpoint{0.050000in}{0.050000in}}{\pgfqpoint{4.900000in}{4.900000in}}%
\pgfusepath{clip}%
\pgfsetbuttcap%
\pgfsetroundjoin%
\definecolor{currentfill}{rgb}{0.667405,0.667405,0.667405}%
\pgfsetfillcolor{currentfill}%
\pgfsetfillopacity{0.900000}%
\pgfsetlinewidth{0.501875pt}%
\definecolor{currentstroke}{rgb}{0.200000,0.200000,0.200000}%
\pgfsetstrokecolor{currentstroke}%
\pgfsetdash{}{0pt}%
\pgfpathmoveto{\pgfqpoint{2.786167in}{2.077942in}}%
\pgfpathlineto{\pgfqpoint{2.832853in}{2.032953in}}%
\pgfpathlineto{\pgfqpoint{2.866361in}{2.011383in}}%
\pgfpathlineto{\pgfqpoint{2.819673in}{2.056372in}}%
\pgfpathlineto{\pgfqpoint{2.786167in}{2.077942in}}%
\pgfpathclose%
\pgfusepath{stroke,fill}%
\end{pgfscope}%
\begin{pgfscope}%
\pgfpathrectangle{\pgfqpoint{0.050000in}{0.050000in}}{\pgfqpoint{4.900000in}{4.900000in}}%
\pgfusepath{clip}%
\pgfsetbuttcap%
\pgfsetroundjoin%
\definecolor{currentfill}{rgb}{0.004552,0.004552,0.004552}%
\pgfsetfillcolor{currentfill}%
\pgfsetfillopacity{0.900000}%
\pgfsetlinewidth{0.501875pt}%
\definecolor{currentstroke}{rgb}{0.200000,0.200000,0.200000}%
\pgfsetstrokecolor{currentstroke}%
\pgfsetdash{}{0pt}%
\pgfpathmoveto{\pgfqpoint{1.027398in}{3.499798in}}%
\pgfpathlineto{\pgfqpoint{1.077938in}{3.451454in}}%
\pgfpathlineto{\pgfqpoint{1.108634in}{3.432093in}}%
\pgfpathlineto{\pgfqpoint{1.058080in}{3.480472in}}%
\pgfpathlineto{\pgfqpoint{1.027398in}{3.499798in}}%
\pgfpathclose%
\pgfusepath{stroke,fill}%
\end{pgfscope}%
\begin{pgfscope}%
\pgfpathrectangle{\pgfqpoint{0.050000in}{0.050000in}}{\pgfqpoint{4.900000in}{4.900000in}}%
\pgfusepath{clip}%
\pgfsetbuttcap%
\pgfsetroundjoin%
\definecolor{currentfill}{rgb}{0.306344,0.306344,0.306344}%
\pgfsetfillcolor{currentfill}%
\pgfsetfillopacity{0.900000}%
\pgfsetlinewidth{0.501875pt}%
\definecolor{currentstroke}{rgb}{0.200000,0.200000,0.200000}%
\pgfsetstrokecolor{currentstroke}%
\pgfsetdash{}{0pt}%
\pgfpathmoveto{\pgfqpoint{1.769622in}{2.899927in}}%
\pgfpathlineto{\pgfqpoint{1.818539in}{2.852895in}}%
\pgfpathlineto{\pgfqpoint{1.850424in}{2.832575in}}%
\pgfpathlineto{\pgfqpoint{1.801498in}{2.879638in}}%
\pgfpathlineto{\pgfqpoint{1.769622in}{2.899927in}}%
\pgfpathclose%
\pgfusepath{stroke,fill}%
\end{pgfscope}%
\begin{pgfscope}%
\pgfpathrectangle{\pgfqpoint{0.050000in}{0.050000in}}{\pgfqpoint{4.900000in}{4.900000in}}%
\pgfusepath{clip}%
\pgfsetbuttcap%
\pgfsetroundjoin%
\definecolor{currentfill}{rgb}{0.564552,0.564552,0.564552}%
\pgfsetfillcolor{currentfill}%
\pgfsetfillopacity{0.900000}%
\pgfsetlinewidth{0.501875pt}%
\definecolor{currentstroke}{rgb}{0.200000,0.200000,0.200000}%
\pgfsetstrokecolor{currentstroke}%
\pgfsetdash{}{0pt}%
\pgfpathmoveto{\pgfqpoint{2.511949in}{2.300002in}}%
\pgfpathlineto{\pgfqpoint{2.559243in}{2.254297in}}%
\pgfpathlineto{\pgfqpoint{2.592317in}{2.233025in}}%
\pgfpathlineto{\pgfqpoint{2.545020in}{2.278753in}}%
\pgfpathlineto{\pgfqpoint{2.511949in}{2.300002in}}%
\pgfpathclose%
\pgfusepath{stroke,fill}%
\end{pgfscope}%
\begin{pgfscope}%
\pgfpathrectangle{\pgfqpoint{0.050000in}{0.050000in}}{\pgfqpoint{4.900000in}{4.900000in}}%
\pgfusepath{clip}%
\pgfsetbuttcap%
\pgfsetroundjoin%
\definecolor{currentfill}{rgb}{0.184544,0.184544,0.184544}%
\pgfsetfillcolor{currentfill}%
\pgfsetfillopacity{0.900000}%
\pgfsetlinewidth{0.501875pt}%
\definecolor{currentstroke}{rgb}{0.200000,0.200000,0.200000}%
\pgfsetstrokecolor{currentstroke}%
\pgfsetdash{}{0pt}%
\pgfpathmoveto{\pgfqpoint{1.495199in}{3.122398in}}%
\pgfpathlineto{\pgfqpoint{1.544726in}{3.074870in}}%
\pgfpathlineto{\pgfqpoint{1.576178in}{3.054902in}}%
\pgfpathlineto{\pgfqpoint{1.526640in}{3.102463in}}%
\pgfpathlineto{\pgfqpoint{1.495199in}{3.122398in}}%
\pgfpathclose%
\pgfusepath{stroke,fill}%
\end{pgfscope}%
\begin{pgfscope}%
\pgfpathrectangle{\pgfqpoint{0.050000in}{0.050000in}}{\pgfqpoint{4.900000in}{4.900000in}}%
\pgfusepath{clip}%
\pgfsetbuttcap%
\pgfsetroundjoin%
\definecolor{currentfill}{rgb}{0.465513,0.465513,0.465513}%
\pgfsetfillcolor{currentfill}%
\pgfsetfillopacity{0.900000}%
\pgfsetlinewidth{0.501875pt}%
\definecolor{currentstroke}{rgb}{0.200000,0.200000,0.200000}%
\pgfsetstrokecolor{currentstroke}%
\pgfsetdash{}{0pt}%
\pgfpathmoveto{\pgfqpoint{2.237636in}{2.522367in}}%
\pgfpathlineto{\pgfqpoint{2.285540in}{2.476152in}}%
\pgfpathlineto{\pgfqpoint{2.318181in}{2.455224in}}%
\pgfpathlineto{\pgfqpoint{2.270272in}{2.501468in}}%
\pgfpathlineto{\pgfqpoint{2.237636in}{2.522367in}}%
\pgfpathclose%
\pgfusepath{stroke,fill}%
\end{pgfscope}%
\begin{pgfscope}%
\pgfpathrectangle{\pgfqpoint{0.050000in}{0.050000in}}{\pgfqpoint{4.900000in}{4.900000in}}%
\pgfusepath{clip}%
\pgfsetbuttcap%
\pgfsetroundjoin%
\definecolor{currentfill}{rgb}{0.815671,0.815671,0.815671}%
\pgfsetfillcolor{currentfill}%
\pgfsetfillopacity{0.900000}%
\pgfsetlinewidth{0.501875pt}%
\definecolor{currentstroke}{rgb}{0.200000,0.200000,0.200000}%
\pgfsetstrokecolor{currentstroke}%
\pgfsetdash{}{0pt}%
\pgfpathmoveto{\pgfqpoint{3.254469in}{1.707240in}}%
\pgfpathlineto{\pgfqpoint{3.300194in}{1.666520in}}%
\pgfpathlineto{\pgfqpoint{3.334472in}{1.644843in}}%
\pgfpathlineto{\pgfqpoint{3.288744in}{1.685446in}}%
\pgfpathlineto{\pgfqpoint{3.254469in}{1.707240in}}%
\pgfpathclose%
\pgfusepath{stroke,fill}%
\end{pgfscope}%
\begin{pgfscope}%
\pgfpathrectangle{\pgfqpoint{0.050000in}{0.050000in}}{\pgfqpoint{4.900000in}{4.900000in}}%
\pgfusepath{clip}%
\pgfsetbuttcap%
\pgfsetroundjoin%
\definecolor{currentfill}{rgb}{0.739377,0.739377,0.739377}%
\pgfsetfillcolor{currentfill}%
\pgfsetfillopacity{0.900000}%
\pgfsetlinewidth{0.501875pt}%
\definecolor{currentstroke}{rgb}{0.200000,0.200000,0.200000}%
\pgfsetstrokecolor{currentstroke}%
\pgfsetdash{}{0pt}%
\pgfpathmoveto{\pgfqpoint{2.980186in}{1.923492in}}%
\pgfpathlineto{\pgfqpoint{3.026474in}{1.879483in}}%
\pgfpathlineto{\pgfqpoint{3.060308in}{1.857778in}}%
\pgfpathlineto{\pgfqpoint{3.014019in}{1.901748in}}%
\pgfpathlineto{\pgfqpoint{2.980186in}{1.923492in}}%
\pgfpathclose%
\pgfusepath{stroke,fill}%
\end{pgfscope}%
\begin{pgfscope}%
\pgfpathrectangle{\pgfqpoint{0.050000in}{0.050000in}}{\pgfqpoint{4.900000in}{4.900000in}}%
\pgfusepath{clip}%
\pgfsetbuttcap%
\pgfsetroundjoin%
\definecolor{currentfill}{rgb}{0.072834,0.072834,0.072834}%
\pgfsetfillcolor{currentfill}%
\pgfsetfillopacity{0.900000}%
\pgfsetlinewidth{0.501875pt}%
\definecolor{currentstroke}{rgb}{0.200000,0.200000,0.200000}%
\pgfsetstrokecolor{currentstroke}%
\pgfsetdash{}{0pt}%
\pgfpathmoveto{\pgfqpoint{1.220680in}{3.344950in}}%
\pgfpathlineto{\pgfqpoint{1.270818in}{3.296926in}}%
\pgfpathlineto{\pgfqpoint{1.301836in}{3.277310in}}%
\pgfpathlineto{\pgfqpoint{1.251686in}{3.325368in}}%
\pgfpathlineto{\pgfqpoint{1.220680in}{3.344950in}}%
\pgfpathclose%
\pgfusepath{stroke,fill}%
\end{pgfscope}%
\begin{pgfscope}%
\pgfpathrectangle{\pgfqpoint{0.050000in}{0.050000in}}{\pgfqpoint{4.900000in}{4.900000in}}%
\pgfusepath{clip}%
\pgfsetbuttcap%
\pgfsetroundjoin%
\definecolor{currentfill}{rgb}{0.371303,0.371303,0.371303}%
\pgfsetfillcolor{currentfill}%
\pgfsetfillopacity{0.900000}%
\pgfsetlinewidth{0.501875pt}%
\definecolor{currentstroke}{rgb}{0.200000,0.200000,0.200000}%
\pgfsetstrokecolor{currentstroke}%
\pgfsetdash{}{0pt}%
\pgfpathmoveto{\pgfqpoint{1.963228in}{2.744823in}}%
\pgfpathlineto{\pgfqpoint{2.011741in}{2.698111in}}%
\pgfpathlineto{\pgfqpoint{2.043949in}{2.677535in}}%
\pgfpathlineto{\pgfqpoint{1.995429in}{2.724277in}}%
\pgfpathlineto{\pgfqpoint{1.963228in}{2.744823in}}%
\pgfpathclose%
\pgfusepath{stroke,fill}%
\end{pgfscope}%
\begin{pgfscope}%
\pgfpathrectangle{\pgfqpoint{0.050000in}{0.050000in}}{\pgfqpoint{4.900000in}{4.900000in}}%
\pgfusepath{clip}%
\pgfsetbuttcap%
\pgfsetroundjoin%
\definecolor{currentfill}{rgb}{0.633818,0.633818,0.633818}%
\pgfsetfillcolor{currentfill}%
\pgfsetfillopacity{0.900000}%
\pgfsetlinewidth{0.501875pt}%
\definecolor{currentstroke}{rgb}{0.200000,0.200000,0.200000}%
\pgfsetstrokecolor{currentstroke}%
\pgfsetdash{}{0pt}%
\pgfpathmoveto{\pgfqpoint{2.705879in}{2.144734in}}%
\pgfpathlineto{\pgfqpoint{2.752768in}{2.099447in}}%
\pgfpathlineto{\pgfqpoint{2.786167in}{2.077942in}}%
\pgfpathlineto{\pgfqpoint{2.739276in}{2.123240in}}%
\pgfpathlineto{\pgfqpoint{2.705879in}{2.144734in}}%
\pgfpathclose%
\pgfusepath{stroke,fill}%
\end{pgfscope}%
\begin{pgfscope}%
\pgfpathrectangle{\pgfqpoint{0.050000in}{0.050000in}}{\pgfqpoint{4.900000in}{4.900000in}}%
\pgfusepath{clip}%
\pgfsetbuttcap%
\pgfsetroundjoin%
\definecolor{currentfill}{rgb}{0.267589,0.267589,0.267589}%
\pgfsetfillcolor{currentfill}%
\pgfsetfillopacity{0.900000}%
\pgfsetlinewidth{0.501875pt}%
\definecolor{currentstroke}{rgb}{0.200000,0.200000,0.200000}%
\pgfsetstrokecolor{currentstroke}%
\pgfsetdash{}{0pt}%
\pgfpathmoveto{\pgfqpoint{1.688724in}{2.967359in}}%
\pgfpathlineto{\pgfqpoint{1.737848in}{2.920151in}}%
\pgfpathlineto{\pgfqpoint{1.769622in}{2.899927in}}%
\pgfpathlineto{\pgfqpoint{1.720489in}{2.947166in}}%
\pgfpathlineto{\pgfqpoint{1.688724in}{2.967359in}}%
\pgfpathclose%
\pgfusepath{stroke,fill}%
\end{pgfscope}%
\begin{pgfscope}%
\pgfpathrectangle{\pgfqpoint{0.050000in}{0.050000in}}{\pgfqpoint{4.900000in}{4.900000in}}%
\pgfusepath{clip}%
\pgfsetbuttcap%
\pgfsetroundjoin%
\definecolor{currentfill}{rgb}{0.530104,0.530104,0.530104}%
\pgfsetfillcolor{currentfill}%
\pgfsetfillopacity{0.900000}%
\pgfsetlinewidth{0.501875pt}%
\definecolor{currentstroke}{rgb}{0.200000,0.200000,0.200000}%
\pgfsetstrokecolor{currentstroke}%
\pgfsetdash{}{0pt}%
\pgfpathmoveto{\pgfqpoint{2.431486in}{2.367074in}}%
\pgfpathlineto{\pgfqpoint{2.478984in}{2.321184in}}%
\pgfpathlineto{\pgfqpoint{2.511949in}{2.300002in}}%
\pgfpathlineto{\pgfqpoint{2.464447in}{2.345918in}}%
\pgfpathlineto{\pgfqpoint{2.431486in}{2.367074in}}%
\pgfpathclose%
\pgfusepath{stroke,fill}%
\end{pgfscope}%
\begin{pgfscope}%
\pgfpathrectangle{\pgfqpoint{0.050000in}{0.050000in}}{\pgfqpoint{4.900000in}{4.900000in}}%
\pgfusepath{clip}%
\pgfsetbuttcap%
\pgfsetroundjoin%
\definecolor{currentfill}{rgb}{0.141115,0.141115,0.141115}%
\pgfsetfillcolor{currentfill}%
\pgfsetfillopacity{0.900000}%
\pgfsetlinewidth{0.501875pt}%
\definecolor{currentstroke}{rgb}{0.200000,0.200000,0.200000}%
\pgfsetstrokecolor{currentstroke}%
\pgfsetdash{}{0pt}%
\pgfpathmoveto{\pgfqpoint{1.414124in}{3.189975in}}%
\pgfpathlineto{\pgfqpoint{1.463858in}{3.142270in}}%
\pgfpathlineto{\pgfqpoint{1.495199in}{3.122398in}}%
\pgfpathlineto{\pgfqpoint{1.445454in}{3.170136in}}%
\pgfpathlineto{\pgfqpoint{1.414124in}{3.189975in}}%
\pgfpathclose%
\pgfusepath{stroke,fill}%
\end{pgfscope}%
\begin{pgfscope}%
\pgfpathrectangle{\pgfqpoint{0.050000in}{0.050000in}}{\pgfqpoint{4.900000in}{4.900000in}}%
\pgfusepath{clip}%
\pgfsetbuttcap%
\pgfsetroundjoin%
\definecolor{currentfill}{rgb}{0.791557,0.791557,0.791557}%
\pgfsetfillcolor{currentfill}%
\pgfsetfillopacity{0.900000}%
\pgfsetlinewidth{0.501875pt}%
\definecolor{currentstroke}{rgb}{0.200000,0.200000,0.200000}%
\pgfsetstrokecolor{currentstroke}%
\pgfsetdash{}{0pt}%
\pgfpathmoveto{\pgfqpoint{3.174398in}{1.771076in}}%
\pgfpathlineto{\pgfqpoint{3.220304in}{1.728973in}}%
\pgfpathlineto{\pgfqpoint{3.254469in}{1.707240in}}%
\pgfpathlineto{\pgfqpoint{3.208562in}{1.749251in}}%
\pgfpathlineto{\pgfqpoint{3.174398in}{1.771076in}}%
\pgfpathclose%
\pgfusepath{stroke,fill}%
\end{pgfscope}%
\begin{pgfscope}%
\pgfpathrectangle{\pgfqpoint{0.050000in}{0.050000in}}{\pgfqpoint{4.900000in}{4.900000in}}%
\pgfusepath{clip}%
\pgfsetbuttcap%
\pgfsetroundjoin%
\definecolor{currentfill}{rgb}{0.436263,0.436263,0.436263}%
\pgfsetfillcolor{currentfill}%
\pgfsetfillopacity{0.900000}%
\pgfsetlinewidth{0.501875pt}%
\definecolor{currentstroke}{rgb}{0.200000,0.200000,0.200000}%
\pgfsetstrokecolor{currentstroke}%
\pgfsetdash{}{0pt}%
\pgfpathmoveto{\pgfqpoint{2.156996in}{2.589591in}}%
\pgfpathlineto{\pgfqpoint{2.205105in}{2.543200in}}%
\pgfpathlineto{\pgfqpoint{2.237636in}{2.522367in}}%
\pgfpathlineto{\pgfqpoint{2.189522in}{2.568787in}}%
\pgfpathlineto{\pgfqpoint{2.156996in}{2.589591in}}%
\pgfpathclose%
\pgfusepath{stroke,fill}%
\end{pgfscope}%
\begin{pgfscope}%
\pgfpathrectangle{\pgfqpoint{0.050000in}{0.050000in}}{\pgfqpoint{4.900000in}{4.900000in}}%
\pgfusepath{clip}%
\pgfsetbuttcap%
\pgfsetroundjoin%
\definecolor{currentfill}{rgb}{0.705790,0.705790,0.705790}%
\pgfsetfillcolor{currentfill}%
\pgfsetfillopacity{0.900000}%
\pgfsetlinewidth{0.501875pt}%
\definecolor{currentstroke}{rgb}{0.200000,0.200000,0.200000}%
\pgfsetstrokecolor{currentstroke}%
\pgfsetdash{}{0pt}%
\pgfpathmoveto{\pgfqpoint{2.899975in}{1.989749in}}%
\pgfpathlineto{\pgfqpoint{2.946462in}{1.945173in}}%
\pgfpathlineto{\pgfqpoint{2.980186in}{1.923492in}}%
\pgfpathlineto{\pgfqpoint{2.933699in}{1.968049in}}%
\pgfpathlineto{\pgfqpoint{2.899975in}{1.989749in}}%
\pgfpathclose%
\pgfusepath{stroke,fill}%
\end{pgfscope}%
\begin{pgfscope}%
\pgfpathrectangle{\pgfqpoint{0.050000in}{0.050000in}}{\pgfqpoint{4.900000in}{4.900000in}}%
\pgfusepath{clip}%
\pgfsetbuttcap%
\pgfsetroundjoin%
\definecolor{currentfill}{rgb}{0.036417,0.036417,0.036417}%
\pgfsetfillcolor{currentfill}%
\pgfsetfillopacity{0.900000}%
\pgfsetlinewidth{0.501875pt}%
\definecolor{currentstroke}{rgb}{0.200000,0.200000,0.200000}%
\pgfsetstrokecolor{currentstroke}%
\pgfsetdash{}{0pt}%
\pgfpathmoveto{\pgfqpoint{1.139428in}{3.412671in}}%
\pgfpathlineto{\pgfqpoint{1.189774in}{3.364470in}}%
\pgfpathlineto{\pgfqpoint{1.220680in}{3.344950in}}%
\pgfpathlineto{\pgfqpoint{1.170322in}{3.393186in}}%
\pgfpathlineto{\pgfqpoint{1.139428in}{3.412671in}}%
\pgfpathclose%
\pgfusepath{stroke,fill}%
\end{pgfscope}%
\begin{pgfscope}%
\pgfpathrectangle{\pgfqpoint{0.050000in}{0.050000in}}{\pgfqpoint{4.900000in}{4.900000in}}%
\pgfusepath{clip}%
\pgfsetbuttcap%
\pgfsetroundjoin%
\definecolor{currentfill}{rgb}{0.342884,0.342884,0.342884}%
\pgfsetfillcolor{currentfill}%
\pgfsetfillopacity{0.900000}%
\pgfsetlinewidth{0.501875pt}%
\definecolor{currentstroke}{rgb}{0.200000,0.200000,0.200000}%
\pgfsetstrokecolor{currentstroke}%
\pgfsetdash{}{0pt}%
\pgfpathmoveto{\pgfqpoint{1.882411in}{2.812190in}}%
\pgfpathlineto{\pgfqpoint{1.931130in}{2.765303in}}%
\pgfpathlineto{\pgfqpoint{1.963228in}{2.744823in}}%
\pgfpathlineto{\pgfqpoint{1.914501in}{2.791740in}}%
\pgfpathlineto{\pgfqpoint{1.882411in}{2.812190in}}%
\pgfpathclose%
\pgfusepath{stroke,fill}%
\end{pgfscope}%
\begin{pgfscope}%
\pgfpathrectangle{\pgfqpoint{0.050000in}{0.050000in}}{\pgfqpoint{4.900000in}{4.900000in}}%
\pgfusepath{clip}%
\pgfsetbuttcap%
\pgfsetroundjoin%
\definecolor{currentfill}{rgb}{0.595433,0.595433,0.595433}%
\pgfsetfillcolor{currentfill}%
\pgfsetfillopacity{0.900000}%
\pgfsetlinewidth{0.501875pt}%
\definecolor{currentstroke}{rgb}{0.200000,0.200000,0.200000}%
\pgfsetstrokecolor{currentstroke}%
\pgfsetdash{}{0pt}%
\pgfpathmoveto{\pgfqpoint{2.625497in}{2.211684in}}%
\pgfpathlineto{\pgfqpoint{2.672590in}{2.166160in}}%
\pgfpathlineto{\pgfqpoint{2.705879in}{2.144734in}}%
\pgfpathlineto{\pgfqpoint{2.658784in}{2.190277in}}%
\pgfpathlineto{\pgfqpoint{2.625497in}{2.211684in}}%
\pgfpathclose%
\pgfusepath{stroke,fill}%
\end{pgfscope}%
\begin{pgfscope}%
\pgfpathrectangle{\pgfqpoint{0.050000in}{0.050000in}}{\pgfqpoint{4.900000in}{4.900000in}}%
\pgfusepath{clip}%
\pgfsetbuttcap%
\pgfsetroundjoin%
\definecolor{currentfill}{rgb}{0.223299,0.223299,0.223299}%
\pgfsetfillcolor{currentfill}%
\pgfsetfillopacity{0.900000}%
\pgfsetlinewidth{0.501875pt}%
\definecolor{currentstroke}{rgb}{0.200000,0.200000,0.200000}%
\pgfsetstrokecolor{currentstroke}%
\pgfsetdash{}{0pt}%
\pgfpathmoveto{\pgfqpoint{1.607730in}{3.034871in}}%
\pgfpathlineto{\pgfqpoint{1.657060in}{2.987487in}}%
\pgfpathlineto{\pgfqpoint{1.688724in}{2.967359in}}%
\pgfpathlineto{\pgfqpoint{1.639384in}{3.014774in}}%
\pgfpathlineto{\pgfqpoint{1.607730in}{3.034871in}}%
\pgfpathclose%
\pgfusepath{stroke,fill}%
\end{pgfscope}%
\begin{pgfscope}%
\pgfpathrectangle{\pgfqpoint{0.050000in}{0.050000in}}{\pgfqpoint{4.900000in}{4.900000in}}%
\pgfusepath{clip}%
\pgfsetbuttcap%
\pgfsetroundjoin%
\definecolor{currentfill}{rgb}{0.499962,0.499962,0.499962}%
\pgfsetfillcolor{currentfill}%
\pgfsetfillopacity{0.900000}%
\pgfsetlinewidth{0.501875pt}%
\definecolor{currentstroke}{rgb}{0.200000,0.200000,0.200000}%
\pgfsetstrokecolor{currentstroke}%
\pgfsetdash{}{0pt}%
\pgfpathmoveto{\pgfqpoint{2.350927in}{2.434230in}}%
\pgfpathlineto{\pgfqpoint{2.398630in}{2.388162in}}%
\pgfpathlineto{\pgfqpoint{2.431486in}{2.367074in}}%
\pgfpathlineto{\pgfqpoint{2.383778in}{2.413169in}}%
\pgfpathlineto{\pgfqpoint{2.350927in}{2.434230in}}%
\pgfpathclose%
\pgfusepath{stroke,fill}%
\end{pgfscope}%
\begin{pgfscope}%
\pgfpathrectangle{\pgfqpoint{0.050000in}{0.050000in}}{\pgfqpoint{4.900000in}{4.900000in}}%
\pgfusepath{clip}%
\pgfsetbuttcap%
\pgfsetroundjoin%
\definecolor{currentfill}{rgb}{0.109250,0.109250,0.109250}%
\pgfsetfillcolor{currentfill}%
\pgfsetfillopacity{0.900000}%
\pgfsetlinewidth{0.501875pt}%
\definecolor{currentstroke}{rgb}{0.200000,0.200000,0.200000}%
\pgfsetstrokecolor{currentstroke}%
\pgfsetdash{}{0pt}%
\pgfpathmoveto{\pgfqpoint{1.332953in}{3.257631in}}%
\pgfpathlineto{\pgfqpoint{1.382895in}{3.209750in}}%
\pgfpathlineto{\pgfqpoint{1.414124in}{3.189975in}}%
\pgfpathlineto{\pgfqpoint{1.364171in}{3.237889in}}%
\pgfpathlineto{\pgfqpoint{1.332953in}{3.257631in}}%
\pgfpathclose%
\pgfusepath{stroke,fill}%
\end{pgfscope}%
\begin{pgfscope}%
\pgfpathrectangle{\pgfqpoint{0.050000in}{0.050000in}}{\pgfqpoint{4.900000in}{4.900000in}}%
\pgfusepath{clip}%
\pgfsetbuttcap%
\pgfsetroundjoin%
\definecolor{currentfill}{rgb}{0.767443,0.767443,0.767443}%
\pgfsetfillcolor{currentfill}%
\pgfsetfillopacity{0.900000}%
\pgfsetlinewidth{0.501875pt}%
\definecolor{currentstroke}{rgb}{0.200000,0.200000,0.200000}%
\pgfsetstrokecolor{currentstroke}%
\pgfsetdash{}{0pt}%
\pgfpathmoveto{\pgfqpoint{3.094251in}{1.836010in}}%
\pgfpathlineto{\pgfqpoint{3.140345in}{1.792840in}}%
\pgfpathlineto{\pgfqpoint{3.174398in}{1.771076in}}%
\pgfpathlineto{\pgfqpoint{3.128303in}{1.814181in}}%
\pgfpathlineto{\pgfqpoint{3.094251in}{1.836010in}}%
\pgfpathclose%
\pgfusepath{stroke,fill}%
\end{pgfscope}%
\begin{pgfscope}%
\pgfpathrectangle{\pgfqpoint{0.050000in}{0.050000in}}{\pgfqpoint{4.900000in}{4.900000in}}%
\pgfusepath{clip}%
\pgfsetbuttcap%
\pgfsetroundjoin%
\definecolor{currentfill}{rgb}{0.839785,0.839785,0.839785}%
\pgfsetfillcolor{currentfill}%
\pgfsetfillopacity{0.900000}%
\pgfsetlinewidth{0.501875pt}%
\definecolor{currentstroke}{rgb}{0.200000,0.200000,0.200000}%
\pgfsetstrokecolor{currentstroke}%
\pgfsetdash{}{0pt}%
\pgfpathmoveto{\pgfqpoint{3.368861in}{1.623105in}}%
\pgfpathlineto{\pgfqpoint{3.414419in}{1.584257in}}%
\pgfpathlineto{\pgfqpoint{3.448923in}{1.562596in}}%
\pgfpathlineto{\pgfqpoint{3.403361in}{1.601306in}}%
\pgfpathlineto{\pgfqpoint{3.368861in}{1.623105in}}%
\pgfpathclose%
\pgfusepath{stroke,fill}%
\end{pgfscope}%
\begin{pgfscope}%
\pgfpathrectangle{\pgfqpoint{0.050000in}{0.050000in}}{\pgfqpoint{4.900000in}{4.900000in}}%
\pgfusepath{clip}%
\pgfsetbuttcap%
\pgfsetroundjoin%
\definecolor{currentfill}{rgb}{0.403783,0.403783,0.403783}%
\pgfsetfillcolor{currentfill}%
\pgfsetfillopacity{0.900000}%
\pgfsetlinewidth{0.501875pt}%
\definecolor{currentstroke}{rgb}{0.200000,0.200000,0.200000}%
\pgfsetstrokecolor{currentstroke}%
\pgfsetdash{}{0pt}%
\pgfpathmoveto{\pgfqpoint{2.076261in}{2.656894in}}%
\pgfpathlineto{\pgfqpoint{2.124575in}{2.610328in}}%
\pgfpathlineto{\pgfqpoint{2.156996in}{2.589591in}}%
\pgfpathlineto{\pgfqpoint{2.108676in}{2.636186in}}%
\pgfpathlineto{\pgfqpoint{2.076261in}{2.656894in}}%
\pgfpathclose%
\pgfusepath{stroke,fill}%
\end{pgfscope}%
\begin{pgfscope}%
\pgfpathrectangle{\pgfqpoint{0.050000in}{0.050000in}}{\pgfqpoint{4.900000in}{4.900000in}}%
\pgfusepath{clip}%
\pgfsetbuttcap%
\pgfsetroundjoin%
\definecolor{currentfill}{rgb}{0.667405,0.667405,0.667405}%
\pgfsetfillcolor{currentfill}%
\pgfsetfillopacity{0.900000}%
\pgfsetlinewidth{0.501875pt}%
\definecolor{currentstroke}{rgb}{0.200000,0.200000,0.200000}%
\pgfsetstrokecolor{currentstroke}%
\pgfsetdash{}{0pt}%
\pgfpathmoveto{\pgfqpoint{2.819673in}{2.056372in}}%
\pgfpathlineto{\pgfqpoint{2.866361in}{2.011383in}}%
\pgfpathlineto{\pgfqpoint{2.899975in}{1.989749in}}%
\pgfpathlineto{\pgfqpoint{2.853286in}{2.034735in}}%
\pgfpathlineto{\pgfqpoint{2.819673in}{2.056372in}}%
\pgfpathclose%
\pgfusepath{stroke,fill}%
\end{pgfscope}%
\begin{pgfscope}%
\pgfpathrectangle{\pgfqpoint{0.050000in}{0.050000in}}{\pgfqpoint{4.900000in}{4.900000in}}%
\pgfusepath{clip}%
\pgfsetbuttcap%
\pgfsetroundjoin%
\definecolor{currentfill}{rgb}{0.004552,0.004552,0.004552}%
\pgfsetfillcolor{currentfill}%
\pgfsetfillopacity{0.900000}%
\pgfsetlinewidth{0.501875pt}%
\definecolor{currentstroke}{rgb}{0.200000,0.200000,0.200000}%
\pgfsetstrokecolor{currentstroke}%
\pgfsetdash{}{0pt}%
\pgfpathmoveto{\pgfqpoint{1.058080in}{3.480472in}}%
\pgfpathlineto{\pgfqpoint{1.108634in}{3.432093in}}%
\pgfpathlineto{\pgfqpoint{1.139428in}{3.412671in}}%
\pgfpathlineto{\pgfqpoint{1.088861in}{3.461084in}}%
\pgfpathlineto{\pgfqpoint{1.058080in}{3.480472in}}%
\pgfpathclose%
\pgfusepath{stroke,fill}%
\end{pgfscope}%
\begin{pgfscope}%
\pgfpathrectangle{\pgfqpoint{0.050000in}{0.050000in}}{\pgfqpoint{4.900000in}{4.900000in}}%
\pgfusepath{clip}%
\pgfsetbuttcap%
\pgfsetroundjoin%
\definecolor{currentfill}{rgb}{0.306344,0.306344,0.306344}%
\pgfsetfillcolor{currentfill}%
\pgfsetfillopacity{0.900000}%
\pgfsetlinewidth{0.501875pt}%
\definecolor{currentstroke}{rgb}{0.200000,0.200000,0.200000}%
\pgfsetstrokecolor{currentstroke}%
\pgfsetdash{}{0pt}%
\pgfpathmoveto{\pgfqpoint{1.801498in}{2.879638in}}%
\pgfpathlineto{\pgfqpoint{1.850424in}{2.832575in}}%
\pgfpathlineto{\pgfqpoint{1.882411in}{2.812190in}}%
\pgfpathlineto{\pgfqpoint{1.833477in}{2.859284in}}%
\pgfpathlineto{\pgfqpoint{1.801498in}{2.879638in}}%
\pgfpathclose%
\pgfusepath{stroke,fill}%
\end{pgfscope}%
\begin{pgfscope}%
\pgfpathrectangle{\pgfqpoint{0.050000in}{0.050000in}}{\pgfqpoint{4.900000in}{4.900000in}}%
\pgfusepath{clip}%
\pgfsetbuttcap%
\pgfsetroundjoin%
\definecolor{currentfill}{rgb}{0.564552,0.564552,0.564552}%
\pgfsetfillcolor{currentfill}%
\pgfsetfillopacity{0.900000}%
\pgfsetlinewidth{0.501875pt}%
\definecolor{currentstroke}{rgb}{0.200000,0.200000,0.200000}%
\pgfsetstrokecolor{currentstroke}%
\pgfsetdash{}{0pt}%
\pgfpathmoveto{\pgfqpoint{2.545020in}{2.278753in}}%
\pgfpathlineto{\pgfqpoint{2.592317in}{2.233025in}}%
\pgfpathlineto{\pgfqpoint{2.625497in}{2.211684in}}%
\pgfpathlineto{\pgfqpoint{2.578197in}{2.257436in}}%
\pgfpathlineto{\pgfqpoint{2.545020in}{2.278753in}}%
\pgfpathclose%
\pgfusepath{stroke,fill}%
\end{pgfscope}%
\begin{pgfscope}%
\pgfpathrectangle{\pgfqpoint{0.050000in}{0.050000in}}{\pgfqpoint{4.900000in}{4.900000in}}%
\pgfusepath{clip}%
\pgfsetbuttcap%
\pgfsetroundjoin%
\definecolor{currentfill}{rgb}{0.184544,0.184544,0.184544}%
\pgfsetfillcolor{currentfill}%
\pgfsetfillopacity{0.900000}%
\pgfsetlinewidth{0.501875pt}%
\definecolor{currentstroke}{rgb}{0.200000,0.200000,0.200000}%
\pgfsetstrokecolor{currentstroke}%
\pgfsetdash{}{0pt}%
\pgfpathmoveto{\pgfqpoint{1.526640in}{3.102463in}}%
\pgfpathlineto{\pgfqpoint{1.576178in}{3.054902in}}%
\pgfpathlineto{\pgfqpoint{1.607730in}{3.034871in}}%
\pgfpathlineto{\pgfqpoint{1.558183in}{3.082463in}}%
\pgfpathlineto{\pgfqpoint{1.526640in}{3.102463in}}%
\pgfpathclose%
\pgfusepath{stroke,fill}%
\end{pgfscope}%
\begin{pgfscope}%
\pgfpathrectangle{\pgfqpoint{0.050000in}{0.050000in}}{\pgfqpoint{4.900000in}{4.900000in}}%
\pgfusepath{clip}%
\pgfsetbuttcap%
\pgfsetroundjoin%
\definecolor{currentfill}{rgb}{0.465513,0.465513,0.465513}%
\pgfsetfillcolor{currentfill}%
\pgfsetfillopacity{0.900000}%
\pgfsetlinewidth{0.501875pt}%
\definecolor{currentstroke}{rgb}{0.200000,0.200000,0.200000}%
\pgfsetstrokecolor{currentstroke}%
\pgfsetdash{}{0pt}%
\pgfpathmoveto{\pgfqpoint{2.270272in}{2.501468in}}%
\pgfpathlineto{\pgfqpoint{2.318181in}{2.455224in}}%
\pgfpathlineto{\pgfqpoint{2.350927in}{2.434230in}}%
\pgfpathlineto{\pgfqpoint{2.303013in}{2.480502in}}%
\pgfpathlineto{\pgfqpoint{2.270272in}{2.501468in}}%
\pgfpathclose%
\pgfusepath{stroke,fill}%
\end{pgfscope}%
\begin{pgfscope}%
\pgfpathrectangle{\pgfqpoint{0.050000in}{0.050000in}}{\pgfqpoint{4.900000in}{4.900000in}}%
\pgfusepath{clip}%
\pgfsetbuttcap%
\pgfsetroundjoin%
\definecolor{currentfill}{rgb}{0.815671,0.815671,0.815671}%
\pgfsetfillcolor{currentfill}%
\pgfsetfillopacity{0.900000}%
\pgfsetlinewidth{0.501875pt}%
\definecolor{currentstroke}{rgb}{0.200000,0.200000,0.200000}%
\pgfsetstrokecolor{currentstroke}%
\pgfsetdash{}{0pt}%
\pgfpathmoveto{\pgfqpoint{3.288744in}{1.685446in}}%
\pgfpathlineto{\pgfqpoint{3.334472in}{1.644843in}}%
\pgfpathlineto{\pgfqpoint{3.368861in}{1.623105in}}%
\pgfpathlineto{\pgfqpoint{3.323131in}{1.663591in}}%
\pgfpathlineto{\pgfqpoint{3.288744in}{1.685446in}}%
\pgfpathclose%
\pgfusepath{stroke,fill}%
\end{pgfscope}%
\begin{pgfscope}%
\pgfpathrectangle{\pgfqpoint{0.050000in}{0.050000in}}{\pgfqpoint{4.900000in}{4.900000in}}%
\pgfusepath{clip}%
\pgfsetbuttcap%
\pgfsetroundjoin%
\definecolor{currentfill}{rgb}{0.739377,0.739377,0.739377}%
\pgfsetfillcolor{currentfill}%
\pgfsetfillopacity{0.900000}%
\pgfsetlinewidth{0.501875pt}%
\definecolor{currentstroke}{rgb}{0.200000,0.200000,0.200000}%
\pgfsetstrokecolor{currentstroke}%
\pgfsetdash{}{0pt}%
\pgfpathmoveto{\pgfqpoint{3.014019in}{1.901748in}}%
\pgfpathlineto{\pgfqpoint{3.060308in}{1.857778in}}%
\pgfpathlineto{\pgfqpoint{3.094251in}{1.836010in}}%
\pgfpathlineto{\pgfqpoint{3.047961in}{1.879939in}}%
\pgfpathlineto{\pgfqpoint{3.014019in}{1.901748in}}%
\pgfpathclose%
\pgfusepath{stroke,fill}%
\end{pgfscope}%
\begin{pgfscope}%
\pgfpathrectangle{\pgfqpoint{0.050000in}{0.050000in}}{\pgfqpoint{4.900000in}{4.900000in}}%
\pgfusepath{clip}%
\pgfsetbuttcap%
\pgfsetroundjoin%
\definecolor{currentfill}{rgb}{0.072834,0.072834,0.072834}%
\pgfsetfillcolor{currentfill}%
\pgfsetfillopacity{0.900000}%
\pgfsetlinewidth{0.501875pt}%
\definecolor{currentstroke}{rgb}{0.200000,0.200000,0.200000}%
\pgfsetstrokecolor{currentstroke}%
\pgfsetdash{}{0pt}%
\pgfpathmoveto{\pgfqpoint{1.251686in}{3.325368in}}%
\pgfpathlineto{\pgfqpoint{1.301836in}{3.277310in}}%
\pgfpathlineto{\pgfqpoint{1.332953in}{3.257631in}}%
\pgfpathlineto{\pgfqpoint{1.282791in}{3.305723in}}%
\pgfpathlineto{\pgfqpoint{1.251686in}{3.325368in}}%
\pgfpathclose%
\pgfusepath{stroke,fill}%
\end{pgfscope}%
\begin{pgfscope}%
\pgfpathrectangle{\pgfqpoint{0.050000in}{0.050000in}}{\pgfqpoint{4.900000in}{4.900000in}}%
\pgfusepath{clip}%
\pgfsetbuttcap%
\pgfsetroundjoin%
\definecolor{currentfill}{rgb}{0.371303,0.371303,0.371303}%
\pgfsetfillcolor{currentfill}%
\pgfsetfillopacity{0.900000}%
\pgfsetlinewidth{0.501875pt}%
\definecolor{currentstroke}{rgb}{0.200000,0.200000,0.200000}%
\pgfsetstrokecolor{currentstroke}%
\pgfsetdash{}{0pt}%
\pgfpathmoveto{\pgfqpoint{1.995429in}{2.724277in}}%
\pgfpathlineto{\pgfqpoint{2.043949in}{2.677535in}}%
\pgfpathlineto{\pgfqpoint{2.076261in}{2.656894in}}%
\pgfpathlineto{\pgfqpoint{2.027733in}{2.703665in}}%
\pgfpathlineto{\pgfqpoint{1.995429in}{2.724277in}}%
\pgfpathclose%
\pgfusepath{stroke,fill}%
\end{pgfscope}%
\begin{pgfscope}%
\pgfpathrectangle{\pgfqpoint{0.050000in}{0.050000in}}{\pgfqpoint{4.900000in}{4.900000in}}%
\pgfusepath{clip}%
\pgfsetbuttcap%
\pgfsetroundjoin%
\definecolor{currentfill}{rgb}{0.633818,0.633818,0.633818}%
\pgfsetfillcolor{currentfill}%
\pgfsetfillopacity{0.900000}%
\pgfsetlinewidth{0.501875pt}%
\definecolor{currentstroke}{rgb}{0.200000,0.200000,0.200000}%
\pgfsetstrokecolor{currentstroke}%
\pgfsetdash{}{0pt}%
\pgfpathmoveto{\pgfqpoint{2.739276in}{2.123240in}}%
\pgfpathlineto{\pgfqpoint{2.786167in}{2.077942in}}%
\pgfpathlineto{\pgfqpoint{2.819673in}{2.056372in}}%
\pgfpathlineto{\pgfqpoint{2.772780in}{2.101680in}}%
\pgfpathlineto{\pgfqpoint{2.739276in}{2.123240in}}%
\pgfpathclose%
\pgfusepath{stroke,fill}%
\end{pgfscope}%
\begin{pgfscope}%
\pgfpathrectangle{\pgfqpoint{0.050000in}{0.050000in}}{\pgfqpoint{4.900000in}{4.900000in}}%
\pgfusepath{clip}%
\pgfsetbuttcap%
\pgfsetroundjoin%
\definecolor{currentfill}{rgb}{0.267589,0.267589,0.267589}%
\pgfsetfillcolor{currentfill}%
\pgfsetfillopacity{0.900000}%
\pgfsetlinewidth{0.501875pt}%
\definecolor{currentstroke}{rgb}{0.200000,0.200000,0.200000}%
\pgfsetstrokecolor{currentstroke}%
\pgfsetdash{}{0pt}%
\pgfpathmoveto{\pgfqpoint{1.720489in}{2.947166in}}%
\pgfpathlineto{\pgfqpoint{1.769622in}{2.899927in}}%
\pgfpathlineto{\pgfqpoint{1.801498in}{2.879638in}}%
\pgfpathlineto{\pgfqpoint{1.752357in}{2.926909in}}%
\pgfpathlineto{\pgfqpoint{1.720489in}{2.947166in}}%
\pgfpathclose%
\pgfusepath{stroke,fill}%
\end{pgfscope}%
\begin{pgfscope}%
\pgfpathrectangle{\pgfqpoint{0.050000in}{0.050000in}}{\pgfqpoint{4.900000in}{4.900000in}}%
\pgfusepath{clip}%
\pgfsetbuttcap%
\pgfsetroundjoin%
\definecolor{currentfill}{rgb}{0.530104,0.530104,0.530104}%
\pgfsetfillcolor{currentfill}%
\pgfsetfillopacity{0.900000}%
\pgfsetlinewidth{0.501875pt}%
\definecolor{currentstroke}{rgb}{0.200000,0.200000,0.200000}%
\pgfsetstrokecolor{currentstroke}%
\pgfsetdash{}{0pt}%
\pgfpathmoveto{\pgfqpoint{2.464447in}{2.345918in}}%
\pgfpathlineto{\pgfqpoint{2.511949in}{2.300002in}}%
\pgfpathlineto{\pgfqpoint{2.545020in}{2.278753in}}%
\pgfpathlineto{\pgfqpoint{2.497514in}{2.324694in}}%
\pgfpathlineto{\pgfqpoint{2.464447in}{2.345918in}}%
\pgfpathclose%
\pgfusepath{stroke,fill}%
\end{pgfscope}%
\begin{pgfscope}%
\pgfpathrectangle{\pgfqpoint{0.050000in}{0.050000in}}{\pgfqpoint{4.900000in}{4.900000in}}%
\pgfusepath{clip}%
\pgfsetbuttcap%
\pgfsetroundjoin%
\definecolor{currentfill}{rgb}{0.141115,0.141115,0.141115}%
\pgfsetfillcolor{currentfill}%
\pgfsetfillopacity{0.900000}%
\pgfsetlinewidth{0.501875pt}%
\definecolor{currentstroke}{rgb}{0.200000,0.200000,0.200000}%
\pgfsetstrokecolor{currentstroke}%
\pgfsetdash{}{0pt}%
\pgfpathmoveto{\pgfqpoint{1.445454in}{3.170136in}}%
\pgfpathlineto{\pgfqpoint{1.495199in}{3.122398in}}%
\pgfpathlineto{\pgfqpoint{1.526640in}{3.102463in}}%
\pgfpathlineto{\pgfqpoint{1.476884in}{3.150233in}}%
\pgfpathlineto{\pgfqpoint{1.445454in}{3.170136in}}%
\pgfpathclose%
\pgfusepath{stroke,fill}%
\end{pgfscope}%
\begin{pgfscope}%
\pgfpathrectangle{\pgfqpoint{0.050000in}{0.050000in}}{\pgfqpoint{4.900000in}{4.900000in}}%
\pgfusepath{clip}%
\pgfsetbuttcap%
\pgfsetroundjoin%
\definecolor{currentfill}{rgb}{0.436263,0.436263,0.436263}%
\pgfsetfillcolor{currentfill}%
\pgfsetfillopacity{0.900000}%
\pgfsetlinewidth{0.501875pt}%
\definecolor{currentstroke}{rgb}{0.200000,0.200000,0.200000}%
\pgfsetstrokecolor{currentstroke}%
\pgfsetdash{}{0pt}%
\pgfpathmoveto{\pgfqpoint{2.189522in}{2.568787in}}%
\pgfpathlineto{\pgfqpoint{2.237636in}{2.522367in}}%
\pgfpathlineto{\pgfqpoint{2.270272in}{2.501468in}}%
\pgfpathlineto{\pgfqpoint{2.222153in}{2.547917in}}%
\pgfpathlineto{\pgfqpoint{2.189522in}{2.568787in}}%
\pgfpathclose%
\pgfusepath{stroke,fill}%
\end{pgfscope}%
\begin{pgfscope}%
\pgfpathrectangle{\pgfqpoint{0.050000in}{0.050000in}}{\pgfqpoint{4.900000in}{4.900000in}}%
\pgfusepath{clip}%
\pgfsetbuttcap%
\pgfsetroundjoin%
\definecolor{currentfill}{rgb}{0.791557,0.791557,0.791557}%
\pgfsetfillcolor{currentfill}%
\pgfsetfillopacity{0.900000}%
\pgfsetlinewidth{0.501875pt}%
\definecolor{currentstroke}{rgb}{0.200000,0.200000,0.200000}%
\pgfsetstrokecolor{currentstroke}%
\pgfsetdash{}{0pt}%
\pgfpathmoveto{\pgfqpoint{3.208562in}{1.749251in}}%
\pgfpathlineto{\pgfqpoint{3.254469in}{1.707240in}}%
\pgfpathlineto{\pgfqpoint{3.288744in}{1.685446in}}%
\pgfpathlineto{\pgfqpoint{3.242836in}{1.727364in}}%
\pgfpathlineto{\pgfqpoint{3.208562in}{1.749251in}}%
\pgfpathclose%
\pgfusepath{stroke,fill}%
\end{pgfscope}%
\begin{pgfscope}%
\pgfpathrectangle{\pgfqpoint{0.050000in}{0.050000in}}{\pgfqpoint{4.900000in}{4.900000in}}%
\pgfusepath{clip}%
\pgfsetbuttcap%
\pgfsetroundjoin%
\definecolor{currentfill}{rgb}{0.705790,0.705790,0.705790}%
\pgfsetfillcolor{currentfill}%
\pgfsetfillopacity{0.900000}%
\pgfsetlinewidth{0.501875pt}%
\definecolor{currentstroke}{rgb}{0.200000,0.200000,0.200000}%
\pgfsetstrokecolor{currentstroke}%
\pgfsetdash{}{0pt}%
\pgfpathmoveto{\pgfqpoint{2.933699in}{1.968049in}}%
\pgfpathlineto{\pgfqpoint{2.980186in}{1.923492in}}%
\pgfpathlineto{\pgfqpoint{3.014019in}{1.901748in}}%
\pgfpathlineto{\pgfqpoint{2.967530in}{1.946284in}}%
\pgfpathlineto{\pgfqpoint{2.933699in}{1.968049in}}%
\pgfpathclose%
\pgfusepath{stroke,fill}%
\end{pgfscope}%
\begin{pgfscope}%
\pgfpathrectangle{\pgfqpoint{0.050000in}{0.050000in}}{\pgfqpoint{4.900000in}{4.900000in}}%
\pgfusepath{clip}%
\pgfsetbuttcap%
\pgfsetroundjoin%
\definecolor{currentfill}{rgb}{0.040969,0.040969,0.040969}%
\pgfsetfillcolor{currentfill}%
\pgfsetfillopacity{0.900000}%
\pgfsetlinewidth{0.501875pt}%
\definecolor{currentstroke}{rgb}{0.200000,0.200000,0.200000}%
\pgfsetstrokecolor{currentstroke}%
\pgfsetdash{}{0pt}%
\pgfpathmoveto{\pgfqpoint{1.170322in}{3.393186in}}%
\pgfpathlineto{\pgfqpoint{1.220680in}{3.344950in}}%
\pgfpathlineto{\pgfqpoint{1.251686in}{3.325368in}}%
\pgfpathlineto{\pgfqpoint{1.201315in}{3.373638in}}%
\pgfpathlineto{\pgfqpoint{1.170322in}{3.393186in}}%
\pgfpathclose%
\pgfusepath{stroke,fill}%
\end{pgfscope}%
\begin{pgfscope}%
\pgfpathrectangle{\pgfqpoint{0.050000in}{0.050000in}}{\pgfqpoint{4.900000in}{4.900000in}}%
\pgfusepath{clip}%
\pgfsetbuttcap%
\pgfsetroundjoin%
\definecolor{currentfill}{rgb}{0.342884,0.342884,0.342884}%
\pgfsetfillcolor{currentfill}%
\pgfsetfillopacity{0.900000}%
\pgfsetlinewidth{0.501875pt}%
\definecolor{currentstroke}{rgb}{0.200000,0.200000,0.200000}%
\pgfsetstrokecolor{currentstroke}%
\pgfsetdash{}{0pt}%
\pgfpathmoveto{\pgfqpoint{1.914501in}{2.791740in}}%
\pgfpathlineto{\pgfqpoint{1.963228in}{2.744823in}}%
\pgfpathlineto{\pgfqpoint{1.995429in}{2.724277in}}%
\pgfpathlineto{\pgfqpoint{1.946695in}{2.771225in}}%
\pgfpathlineto{\pgfqpoint{1.914501in}{2.791740in}}%
\pgfpathclose%
\pgfusepath{stroke,fill}%
\end{pgfscope}%
\begin{pgfscope}%
\pgfpathrectangle{\pgfqpoint{0.050000in}{0.050000in}}{\pgfqpoint{4.900000in}{4.900000in}}%
\pgfusepath{clip}%
\pgfsetbuttcap%
\pgfsetroundjoin%
\definecolor{currentfill}{rgb}{0.595433,0.595433,0.595433}%
\pgfsetfillcolor{currentfill}%
\pgfsetfillopacity{0.900000}%
\pgfsetlinewidth{0.501875pt}%
\definecolor{currentstroke}{rgb}{0.200000,0.200000,0.200000}%
\pgfsetstrokecolor{currentstroke}%
\pgfsetdash{}{0pt}%
\pgfpathmoveto{\pgfqpoint{2.658784in}{2.190277in}}%
\pgfpathlineto{\pgfqpoint{2.705879in}{2.144734in}}%
\pgfpathlineto{\pgfqpoint{2.739276in}{2.123240in}}%
\pgfpathlineto{\pgfqpoint{2.692178in}{2.168801in}}%
\pgfpathlineto{\pgfqpoint{2.658784in}{2.190277in}}%
\pgfpathclose%
\pgfusepath{stroke,fill}%
\end{pgfscope}%
\begin{pgfscope}%
\pgfpathrectangle{\pgfqpoint{0.050000in}{0.050000in}}{\pgfqpoint{4.900000in}{4.900000in}}%
\pgfusepath{clip}%
\pgfsetbuttcap%
\pgfsetroundjoin%
\definecolor{currentfill}{rgb}{0.223299,0.223299,0.223299}%
\pgfsetfillcolor{currentfill}%
\pgfsetfillopacity{0.900000}%
\pgfsetlinewidth{0.501875pt}%
\definecolor{currentstroke}{rgb}{0.200000,0.200000,0.200000}%
\pgfsetstrokecolor{currentstroke}%
\pgfsetdash{}{0pt}%
\pgfpathmoveto{\pgfqpoint{1.639384in}{3.014774in}}%
\pgfpathlineto{\pgfqpoint{1.688724in}{2.967359in}}%
\pgfpathlineto{\pgfqpoint{1.720489in}{2.947166in}}%
\pgfpathlineto{\pgfqpoint{1.671140in}{2.994614in}}%
\pgfpathlineto{\pgfqpoint{1.639384in}{3.014774in}}%
\pgfpathclose%
\pgfusepath{stroke,fill}%
\end{pgfscope}%
\begin{pgfscope}%
\pgfpathrectangle{\pgfqpoint{0.050000in}{0.050000in}}{\pgfqpoint{4.900000in}{4.900000in}}%
\pgfusepath{clip}%
\pgfsetbuttcap%
\pgfsetroundjoin%
\definecolor{currentfill}{rgb}{0.499962,0.499962,0.499962}%
\pgfsetfillcolor{currentfill}%
\pgfsetfillopacity{0.900000}%
\pgfsetlinewidth{0.501875pt}%
\definecolor{currentstroke}{rgb}{0.200000,0.200000,0.200000}%
\pgfsetstrokecolor{currentstroke}%
\pgfsetdash{}{0pt}%
\pgfpathmoveto{\pgfqpoint{2.383778in}{2.413169in}}%
\pgfpathlineto{\pgfqpoint{2.431486in}{2.367074in}}%
\pgfpathlineto{\pgfqpoint{2.464447in}{2.345918in}}%
\pgfpathlineto{\pgfqpoint{2.416735in}{2.392040in}}%
\pgfpathlineto{\pgfqpoint{2.383778in}{2.413169in}}%
\pgfpathclose%
\pgfusepath{stroke,fill}%
\end{pgfscope}%
\begin{pgfscope}%
\pgfpathrectangle{\pgfqpoint{0.050000in}{0.050000in}}{\pgfqpoint{4.900000in}{4.900000in}}%
\pgfusepath{clip}%
\pgfsetbuttcap%
\pgfsetroundjoin%
\definecolor{currentfill}{rgb}{0.109250,0.109250,0.109250}%
\pgfsetfillcolor{currentfill}%
\pgfsetfillopacity{0.900000}%
\pgfsetlinewidth{0.501875pt}%
\definecolor{currentstroke}{rgb}{0.200000,0.200000,0.200000}%
\pgfsetstrokecolor{currentstroke}%
\pgfsetdash{}{0pt}%
\pgfpathmoveto{\pgfqpoint{1.364171in}{3.237889in}}%
\pgfpathlineto{\pgfqpoint{1.414124in}{3.189975in}}%
\pgfpathlineto{\pgfqpoint{1.445454in}{3.170136in}}%
\pgfpathlineto{\pgfqpoint{1.395489in}{3.218083in}}%
\pgfpathlineto{\pgfqpoint{1.364171in}{3.237889in}}%
\pgfpathclose%
\pgfusepath{stroke,fill}%
\end{pgfscope}%
\begin{pgfscope}%
\pgfpathrectangle{\pgfqpoint{0.050000in}{0.050000in}}{\pgfqpoint{4.900000in}{4.900000in}}%
\pgfusepath{clip}%
\pgfsetbuttcap%
\pgfsetroundjoin%
\definecolor{currentfill}{rgb}{0.767443,0.767443,0.767443}%
\pgfsetfillcolor{currentfill}%
\pgfsetfillopacity{0.900000}%
\pgfsetlinewidth{0.501875pt}%
\definecolor{currentstroke}{rgb}{0.200000,0.200000,0.200000}%
\pgfsetstrokecolor{currentstroke}%
\pgfsetdash{}{0pt}%
\pgfpathmoveto{\pgfqpoint{3.128303in}{1.814181in}}%
\pgfpathlineto{\pgfqpoint{3.174398in}{1.771076in}}%
\pgfpathlineto{\pgfqpoint{3.208562in}{1.749251in}}%
\pgfpathlineto{\pgfqpoint{3.162465in}{1.792289in}}%
\pgfpathlineto{\pgfqpoint{3.128303in}{1.814181in}}%
\pgfpathclose%
\pgfusepath{stroke,fill}%
\end{pgfscope}%
\begin{pgfscope}%
\pgfpathrectangle{\pgfqpoint{0.050000in}{0.050000in}}{\pgfqpoint{4.900000in}{4.900000in}}%
\pgfusepath{clip}%
\pgfsetbuttcap%
\pgfsetroundjoin%
\definecolor{currentfill}{rgb}{0.403783,0.403783,0.403783}%
\pgfsetfillcolor{currentfill}%
\pgfsetfillopacity{0.900000}%
\pgfsetlinewidth{0.501875pt}%
\definecolor{currentstroke}{rgb}{0.200000,0.200000,0.200000}%
\pgfsetstrokecolor{currentstroke}%
\pgfsetdash{}{0pt}%
\pgfpathmoveto{\pgfqpoint{2.108676in}{2.636186in}}%
\pgfpathlineto{\pgfqpoint{2.156996in}{2.589591in}}%
\pgfpathlineto{\pgfqpoint{2.189522in}{2.568787in}}%
\pgfpathlineto{\pgfqpoint{2.141195in}{2.615412in}}%
\pgfpathlineto{\pgfqpoint{2.108676in}{2.636186in}}%
\pgfpathclose%
\pgfusepath{stroke,fill}%
\end{pgfscope}%
\begin{pgfscope}%
\pgfpathrectangle{\pgfqpoint{0.050000in}{0.050000in}}{\pgfqpoint{4.900000in}{4.900000in}}%
\pgfusepath{clip}%
\pgfsetbuttcap%
\pgfsetroundjoin%
\definecolor{currentfill}{rgb}{0.667405,0.667405,0.667405}%
\pgfsetfillcolor{currentfill}%
\pgfsetfillopacity{0.900000}%
\pgfsetlinewidth{0.501875pt}%
\definecolor{currentstroke}{rgb}{0.200000,0.200000,0.200000}%
\pgfsetstrokecolor{currentstroke}%
\pgfsetdash{}{0pt}%
\pgfpathmoveto{\pgfqpoint{2.853286in}{2.034735in}}%
\pgfpathlineto{\pgfqpoint{2.899975in}{1.989749in}}%
\pgfpathlineto{\pgfqpoint{2.933699in}{1.968049in}}%
\pgfpathlineto{\pgfqpoint{2.887008in}{2.013032in}}%
\pgfpathlineto{\pgfqpoint{2.853286in}{2.034735in}}%
\pgfpathclose%
\pgfusepath{stroke,fill}%
\end{pgfscope}%
\begin{pgfscope}%
\pgfpathrectangle{\pgfqpoint{0.050000in}{0.050000in}}{\pgfqpoint{4.900000in}{4.900000in}}%
\pgfusepath{clip}%
\pgfsetbuttcap%
\pgfsetroundjoin%
\definecolor{currentfill}{rgb}{0.004552,0.004552,0.004552}%
\pgfsetfillcolor{currentfill}%
\pgfsetfillopacity{0.900000}%
\pgfsetlinewidth{0.501875pt}%
\definecolor{currentstroke}{rgb}{0.200000,0.200000,0.200000}%
\pgfsetstrokecolor{currentstroke}%
\pgfsetdash{}{0pt}%
\pgfpathmoveto{\pgfqpoint{1.088861in}{3.461084in}}%
\pgfpathlineto{\pgfqpoint{1.139428in}{3.412671in}}%
\pgfpathlineto{\pgfqpoint{1.170322in}{3.393186in}}%
\pgfpathlineto{\pgfqpoint{1.119742in}{3.441634in}}%
\pgfpathlineto{\pgfqpoint{1.088861in}{3.461084in}}%
\pgfpathclose%
\pgfusepath{stroke,fill}%
\end{pgfscope}%
\begin{pgfscope}%
\pgfpathrectangle{\pgfqpoint{0.050000in}{0.050000in}}{\pgfqpoint{4.900000in}{4.900000in}}%
\pgfusepath{clip}%
\pgfsetbuttcap%
\pgfsetroundjoin%
\definecolor{currentfill}{rgb}{0.306344,0.306344,0.306344}%
\pgfsetfillcolor{currentfill}%
\pgfsetfillopacity{0.900000}%
\pgfsetlinewidth{0.501875pt}%
\definecolor{currentstroke}{rgb}{0.200000,0.200000,0.200000}%
\pgfsetstrokecolor{currentstroke}%
\pgfsetdash{}{0pt}%
\pgfpathmoveto{\pgfqpoint{1.833477in}{2.859284in}}%
\pgfpathlineto{\pgfqpoint{1.882411in}{2.812190in}}%
\pgfpathlineto{\pgfqpoint{1.914501in}{2.791740in}}%
\pgfpathlineto{\pgfqpoint{1.865559in}{2.838865in}}%
\pgfpathlineto{\pgfqpoint{1.833477in}{2.859284in}}%
\pgfpathclose%
\pgfusepath{stroke,fill}%
\end{pgfscope}%
\begin{pgfscope}%
\pgfpathrectangle{\pgfqpoint{0.050000in}{0.050000in}}{\pgfqpoint{4.900000in}{4.900000in}}%
\pgfusepath{clip}%
\pgfsetbuttcap%
\pgfsetroundjoin%
\definecolor{currentfill}{rgb}{0.564552,0.564552,0.564552}%
\pgfsetfillcolor{currentfill}%
\pgfsetfillopacity{0.900000}%
\pgfsetlinewidth{0.501875pt}%
\definecolor{currentstroke}{rgb}{0.200000,0.200000,0.200000}%
\pgfsetstrokecolor{currentstroke}%
\pgfsetdash{}{0pt}%
\pgfpathmoveto{\pgfqpoint{2.578197in}{2.257436in}}%
\pgfpathlineto{\pgfqpoint{2.625497in}{2.211684in}}%
\pgfpathlineto{\pgfqpoint{2.658784in}{2.190277in}}%
\pgfpathlineto{\pgfqpoint{2.611481in}{2.236051in}}%
\pgfpathlineto{\pgfqpoint{2.578197in}{2.257436in}}%
\pgfpathclose%
\pgfusepath{stroke,fill}%
\end{pgfscope}%
\begin{pgfscope}%
\pgfpathrectangle{\pgfqpoint{0.050000in}{0.050000in}}{\pgfqpoint{4.900000in}{4.900000in}}%
\pgfusepath{clip}%
\pgfsetbuttcap%
\pgfsetroundjoin%
\definecolor{currentfill}{rgb}{0.184544,0.184544,0.184544}%
\pgfsetfillcolor{currentfill}%
\pgfsetfillopacity{0.900000}%
\pgfsetlinewidth{0.501875pt}%
\definecolor{currentstroke}{rgb}{0.200000,0.200000,0.200000}%
\pgfsetstrokecolor{currentstroke}%
\pgfsetdash{}{0pt}%
\pgfpathmoveto{\pgfqpoint{1.558183in}{3.082463in}}%
\pgfpathlineto{\pgfqpoint{1.607730in}{3.034871in}}%
\pgfpathlineto{\pgfqpoint{1.639384in}{3.014774in}}%
\pgfpathlineto{\pgfqpoint{1.589827in}{3.062400in}}%
\pgfpathlineto{\pgfqpoint{1.558183in}{3.082463in}}%
\pgfpathclose%
\pgfusepath{stroke,fill}%
\end{pgfscope}%
\begin{pgfscope}%
\pgfpathrectangle{\pgfqpoint{0.050000in}{0.050000in}}{\pgfqpoint{4.900000in}{4.900000in}}%
\pgfusepath{clip}%
\pgfsetbuttcap%
\pgfsetroundjoin%
\definecolor{currentfill}{rgb}{0.465513,0.465513,0.465513}%
\pgfsetfillcolor{currentfill}%
\pgfsetfillopacity{0.900000}%
\pgfsetlinewidth{0.501875pt}%
\definecolor{currentstroke}{rgb}{0.200000,0.200000,0.200000}%
\pgfsetstrokecolor{currentstroke}%
\pgfsetdash{}{0pt}%
\pgfpathmoveto{\pgfqpoint{2.303013in}{2.480502in}}%
\pgfpathlineto{\pgfqpoint{2.350927in}{2.434230in}}%
\pgfpathlineto{\pgfqpoint{2.383778in}{2.413169in}}%
\pgfpathlineto{\pgfqpoint{2.335860in}{2.459469in}}%
\pgfpathlineto{\pgfqpoint{2.303013in}{2.480502in}}%
\pgfpathclose%
\pgfusepath{stroke,fill}%
\end{pgfscope}%
\begin{pgfscope}%
\pgfpathrectangle{\pgfqpoint{0.050000in}{0.050000in}}{\pgfqpoint{4.900000in}{4.900000in}}%
\pgfusepath{clip}%
\pgfsetbuttcap%
\pgfsetroundjoin%
\definecolor{currentfill}{rgb}{0.815671,0.815671,0.815671}%
\pgfsetfillcolor{currentfill}%
\pgfsetfillopacity{0.900000}%
\pgfsetlinewidth{0.501875pt}%
\definecolor{currentstroke}{rgb}{0.200000,0.200000,0.200000}%
\pgfsetstrokecolor{currentstroke}%
\pgfsetdash{}{0pt}%
\pgfpathmoveto{\pgfqpoint{3.323131in}{1.663591in}}%
\pgfpathlineto{\pgfqpoint{3.368861in}{1.623105in}}%
\pgfpathlineto{\pgfqpoint{3.403361in}{1.601306in}}%
\pgfpathlineto{\pgfqpoint{3.357629in}{1.641674in}}%
\pgfpathlineto{\pgfqpoint{3.323131in}{1.663591in}}%
\pgfpathclose%
\pgfusepath{stroke,fill}%
\end{pgfscope}%
\begin{pgfscope}%
\pgfpathrectangle{\pgfqpoint{0.050000in}{0.050000in}}{\pgfqpoint{4.900000in}{4.900000in}}%
\pgfusepath{clip}%
\pgfsetbuttcap%
\pgfsetroundjoin%
\definecolor{currentfill}{rgb}{0.072834,0.072834,0.072834}%
\pgfsetfillcolor{currentfill}%
\pgfsetfillopacity{0.900000}%
\pgfsetlinewidth{0.501875pt}%
\definecolor{currentstroke}{rgb}{0.200000,0.200000,0.200000}%
\pgfsetstrokecolor{currentstroke}%
\pgfsetdash{}{0pt}%
\pgfpathmoveto{\pgfqpoint{1.282791in}{3.305723in}}%
\pgfpathlineto{\pgfqpoint{1.332953in}{3.257631in}}%
\pgfpathlineto{\pgfqpoint{1.364171in}{3.237889in}}%
\pgfpathlineto{\pgfqpoint{1.313997in}{3.286015in}}%
\pgfpathlineto{\pgfqpoint{1.282791in}{3.305723in}}%
\pgfpathclose%
\pgfusepath{stroke,fill}%
\end{pgfscope}%
\begin{pgfscope}%
\pgfpathrectangle{\pgfqpoint{0.050000in}{0.050000in}}{\pgfqpoint{4.900000in}{4.900000in}}%
\pgfusepath{clip}%
\pgfsetbuttcap%
\pgfsetroundjoin%
\definecolor{currentfill}{rgb}{0.739377,0.739377,0.739377}%
\pgfsetfillcolor{currentfill}%
\pgfsetfillopacity{0.900000}%
\pgfsetlinewidth{0.501875pt}%
\definecolor{currentstroke}{rgb}{0.200000,0.200000,0.200000}%
\pgfsetstrokecolor{currentstroke}%
\pgfsetdash{}{0pt}%
\pgfpathmoveto{\pgfqpoint{3.047961in}{1.879939in}}%
\pgfpathlineto{\pgfqpoint{3.094251in}{1.836010in}}%
\pgfpathlineto{\pgfqpoint{3.128303in}{1.814181in}}%
\pgfpathlineto{\pgfqpoint{3.082012in}{1.858067in}}%
\pgfpathlineto{\pgfqpoint{3.047961in}{1.879939in}}%
\pgfpathclose%
\pgfusepath{stroke,fill}%
\end{pgfscope}%
\begin{pgfscope}%
\pgfpathrectangle{\pgfqpoint{0.050000in}{0.050000in}}{\pgfqpoint{4.900000in}{4.900000in}}%
\pgfusepath{clip}%
\pgfsetbuttcap%
\pgfsetroundjoin%
\definecolor{currentfill}{rgb}{0.375363,0.375363,0.375363}%
\pgfsetfillcolor{currentfill}%
\pgfsetfillopacity{0.900000}%
\pgfsetlinewidth{0.501875pt}%
\definecolor{currentstroke}{rgb}{0.200000,0.200000,0.200000}%
\pgfsetstrokecolor{currentstroke}%
\pgfsetdash{}{0pt}%
\pgfpathmoveto{\pgfqpoint{2.027733in}{2.703665in}}%
\pgfpathlineto{\pgfqpoint{2.076261in}{2.656894in}}%
\pgfpathlineto{\pgfqpoint{2.108676in}{2.636186in}}%
\pgfpathlineto{\pgfqpoint{2.060142in}{2.682987in}}%
\pgfpathlineto{\pgfqpoint{2.027733in}{2.703665in}}%
\pgfpathclose%
\pgfusepath{stroke,fill}%
\end{pgfscope}%
\begin{pgfscope}%
\pgfpathrectangle{\pgfqpoint{0.050000in}{0.050000in}}{\pgfqpoint{4.900000in}{4.900000in}}%
\pgfusepath{clip}%
\pgfsetbuttcap%
\pgfsetroundjoin%
\definecolor{currentfill}{rgb}{0.633818,0.633818,0.633818}%
\pgfsetfillcolor{currentfill}%
\pgfsetfillopacity{0.900000}%
\pgfsetlinewidth{0.501875pt}%
\definecolor{currentstroke}{rgb}{0.200000,0.200000,0.200000}%
\pgfsetstrokecolor{currentstroke}%
\pgfsetdash{}{0pt}%
\pgfpathmoveto{\pgfqpoint{2.772780in}{2.101680in}}%
\pgfpathlineto{\pgfqpoint{2.819673in}{2.056372in}}%
\pgfpathlineto{\pgfqpoint{2.853286in}{2.034735in}}%
\pgfpathlineto{\pgfqpoint{2.806392in}{2.080052in}}%
\pgfpathlineto{\pgfqpoint{2.772780in}{2.101680in}}%
\pgfpathclose%
\pgfusepath{stroke,fill}%
\end{pgfscope}%
\begin{pgfscope}%
\pgfpathrectangle{\pgfqpoint{0.050000in}{0.050000in}}{\pgfqpoint{4.900000in}{4.900000in}}%
\pgfusepath{clip}%
\pgfsetbuttcap%
\pgfsetroundjoin%
\definecolor{currentfill}{rgb}{0.267589,0.267589,0.267589}%
\pgfsetfillcolor{currentfill}%
\pgfsetfillopacity{0.900000}%
\pgfsetlinewidth{0.501875pt}%
\definecolor{currentstroke}{rgb}{0.200000,0.200000,0.200000}%
\pgfsetstrokecolor{currentstroke}%
\pgfsetdash{}{0pt}%
\pgfpathmoveto{\pgfqpoint{1.752357in}{2.926909in}}%
\pgfpathlineto{\pgfqpoint{1.801498in}{2.879638in}}%
\pgfpathlineto{\pgfqpoint{1.833477in}{2.859284in}}%
\pgfpathlineto{\pgfqpoint{1.784327in}{2.906586in}}%
\pgfpathlineto{\pgfqpoint{1.752357in}{2.926909in}}%
\pgfpathclose%
\pgfusepath{stroke,fill}%
\end{pgfscope}%
\begin{pgfscope}%
\pgfpathrectangle{\pgfqpoint{0.050000in}{0.050000in}}{\pgfqpoint{4.900000in}{4.900000in}}%
\pgfusepath{clip}%
\pgfsetbuttcap%
\pgfsetroundjoin%
\definecolor{currentfill}{rgb}{0.530104,0.530104,0.530104}%
\pgfsetfillcolor{currentfill}%
\pgfsetfillopacity{0.900000}%
\pgfsetlinewidth{0.501875pt}%
\definecolor{currentstroke}{rgb}{0.200000,0.200000,0.200000}%
\pgfsetstrokecolor{currentstroke}%
\pgfsetdash{}{0pt}%
\pgfpathmoveto{\pgfqpoint{2.497514in}{2.324694in}}%
\pgfpathlineto{\pgfqpoint{2.545020in}{2.278753in}}%
\pgfpathlineto{\pgfqpoint{2.578197in}{2.257436in}}%
\pgfpathlineto{\pgfqpoint{2.530688in}{2.303402in}}%
\pgfpathlineto{\pgfqpoint{2.497514in}{2.324694in}}%
\pgfpathclose%
\pgfusepath{stroke,fill}%
\end{pgfscope}%
\begin{pgfscope}%
\pgfpathrectangle{\pgfqpoint{0.050000in}{0.050000in}}{\pgfqpoint{4.900000in}{4.900000in}}%
\pgfusepath{clip}%
\pgfsetbuttcap%
\pgfsetroundjoin%
\definecolor{currentfill}{rgb}{0.141115,0.141115,0.141115}%
\pgfsetfillcolor{currentfill}%
\pgfsetfillopacity{0.900000}%
\pgfsetlinewidth{0.501875pt}%
\definecolor{currentstroke}{rgb}{0.200000,0.200000,0.200000}%
\pgfsetstrokecolor{currentstroke}%
\pgfsetdash{}{0pt}%
\pgfpathmoveto{\pgfqpoint{1.476884in}{3.150233in}}%
\pgfpathlineto{\pgfqpoint{1.526640in}{3.102463in}}%
\pgfpathlineto{\pgfqpoint{1.558183in}{3.082463in}}%
\pgfpathlineto{\pgfqpoint{1.508416in}{3.130266in}}%
\pgfpathlineto{\pgfqpoint{1.476884in}{3.150233in}}%
\pgfpathclose%
\pgfusepath{stroke,fill}%
\end{pgfscope}%
\begin{pgfscope}%
\pgfpathrectangle{\pgfqpoint{0.050000in}{0.050000in}}{\pgfqpoint{4.900000in}{4.900000in}}%
\pgfusepath{clip}%
\pgfsetbuttcap%
\pgfsetroundjoin%
\definecolor{currentfill}{rgb}{0.436263,0.436263,0.436263}%
\pgfsetfillcolor{currentfill}%
\pgfsetfillopacity{0.900000}%
\pgfsetlinewidth{0.501875pt}%
\definecolor{currentstroke}{rgb}{0.200000,0.200000,0.200000}%
\pgfsetstrokecolor{currentstroke}%
\pgfsetdash{}{0pt}%
\pgfpathmoveto{\pgfqpoint{2.222153in}{2.547917in}}%
\pgfpathlineto{\pgfqpoint{2.270272in}{2.501468in}}%
\pgfpathlineto{\pgfqpoint{2.303013in}{2.480502in}}%
\pgfpathlineto{\pgfqpoint{2.254888in}{2.526979in}}%
\pgfpathlineto{\pgfqpoint{2.222153in}{2.547917in}}%
\pgfpathclose%
\pgfusepath{stroke,fill}%
\end{pgfscope}%
\begin{pgfscope}%
\pgfpathrectangle{\pgfqpoint{0.050000in}{0.050000in}}{\pgfqpoint{4.900000in}{4.900000in}}%
\pgfusepath{clip}%
\pgfsetbuttcap%
\pgfsetroundjoin%
\definecolor{currentfill}{rgb}{0.791557,0.791557,0.791557}%
\pgfsetfillcolor{currentfill}%
\pgfsetfillopacity{0.900000}%
\pgfsetlinewidth{0.501875pt}%
\definecolor{currentstroke}{rgb}{0.200000,0.200000,0.200000}%
\pgfsetstrokecolor{currentstroke}%
\pgfsetdash{}{0pt}%
\pgfpathmoveto{\pgfqpoint{3.242836in}{1.727364in}}%
\pgfpathlineto{\pgfqpoint{3.288744in}{1.685446in}}%
\pgfpathlineto{\pgfqpoint{3.323131in}{1.663591in}}%
\pgfpathlineto{\pgfqpoint{3.277221in}{1.705416in}}%
\pgfpathlineto{\pgfqpoint{3.242836in}{1.727364in}}%
\pgfpathclose%
\pgfusepath{stroke,fill}%
\end{pgfscope}%
\begin{pgfscope}%
\pgfpathrectangle{\pgfqpoint{0.050000in}{0.050000in}}{\pgfqpoint{4.900000in}{4.900000in}}%
\pgfusepath{clip}%
\pgfsetbuttcap%
\pgfsetroundjoin%
\definecolor{currentfill}{rgb}{0.705790,0.705790,0.705790}%
\pgfsetfillcolor{currentfill}%
\pgfsetfillopacity{0.900000}%
\pgfsetlinewidth{0.501875pt}%
\definecolor{currentstroke}{rgb}{0.200000,0.200000,0.200000}%
\pgfsetstrokecolor{currentstroke}%
\pgfsetdash{}{0pt}%
\pgfpathmoveto{\pgfqpoint{2.967530in}{1.946284in}}%
\pgfpathlineto{\pgfqpoint{3.014019in}{1.901748in}}%
\pgfpathlineto{\pgfqpoint{3.047961in}{1.879939in}}%
\pgfpathlineto{\pgfqpoint{3.001472in}{1.924454in}}%
\pgfpathlineto{\pgfqpoint{2.967530in}{1.946284in}}%
\pgfpathclose%
\pgfusepath{stroke,fill}%
\end{pgfscope}%
\begin{pgfscope}%
\pgfpathrectangle{\pgfqpoint{0.050000in}{0.050000in}}{\pgfqpoint{4.900000in}{4.900000in}}%
\pgfusepath{clip}%
\pgfsetbuttcap%
\pgfsetroundjoin%
\definecolor{currentfill}{rgb}{0.040969,0.040969,0.040969}%
\pgfsetfillcolor{currentfill}%
\pgfsetfillopacity{0.900000}%
\pgfsetlinewidth{0.501875pt}%
\definecolor{currentstroke}{rgb}{0.200000,0.200000,0.200000}%
\pgfsetstrokecolor{currentstroke}%
\pgfsetdash{}{0pt}%
\pgfpathmoveto{\pgfqpoint{1.201315in}{3.373638in}}%
\pgfpathlineto{\pgfqpoint{1.251686in}{3.325368in}}%
\pgfpathlineto{\pgfqpoint{1.282791in}{3.305723in}}%
\pgfpathlineto{\pgfqpoint{1.232408in}{3.354027in}}%
\pgfpathlineto{\pgfqpoint{1.201315in}{3.373638in}}%
\pgfpathclose%
\pgfusepath{stroke,fill}%
\end{pgfscope}%
\begin{pgfscope}%
\pgfpathrectangle{\pgfqpoint{0.050000in}{0.050000in}}{\pgfqpoint{4.900000in}{4.900000in}}%
\pgfusepath{clip}%
\pgfsetbuttcap%
\pgfsetroundjoin%
\definecolor{currentfill}{rgb}{0.342884,0.342884,0.342884}%
\pgfsetfillcolor{currentfill}%
\pgfsetfillopacity{0.900000}%
\pgfsetlinewidth{0.501875pt}%
\definecolor{currentstroke}{rgb}{0.200000,0.200000,0.200000}%
\pgfsetstrokecolor{currentstroke}%
\pgfsetdash{}{0pt}%
\pgfpathmoveto{\pgfqpoint{1.946695in}{2.771225in}}%
\pgfpathlineto{\pgfqpoint{1.995429in}{2.724277in}}%
\pgfpathlineto{\pgfqpoint{2.027733in}{2.703665in}}%
\pgfpathlineto{\pgfqpoint{1.978992in}{2.750643in}}%
\pgfpathlineto{\pgfqpoint{1.946695in}{2.771225in}}%
\pgfpathclose%
\pgfusepath{stroke,fill}%
\end{pgfscope}%
\begin{pgfscope}%
\pgfpathrectangle{\pgfqpoint{0.050000in}{0.050000in}}{\pgfqpoint{4.900000in}{4.900000in}}%
\pgfusepath{clip}%
\pgfsetbuttcap%
\pgfsetroundjoin%
\definecolor{currentfill}{rgb}{0.600231,0.600231,0.600231}%
\pgfsetfillcolor{currentfill}%
\pgfsetfillopacity{0.900000}%
\pgfsetlinewidth{0.501875pt}%
\definecolor{currentstroke}{rgb}{0.200000,0.200000,0.200000}%
\pgfsetstrokecolor{currentstroke}%
\pgfsetdash{}{0pt}%
\pgfpathmoveto{\pgfqpoint{2.692178in}{2.168801in}}%
\pgfpathlineto{\pgfqpoint{2.739276in}{2.123240in}}%
\pgfpathlineto{\pgfqpoint{2.772780in}{2.101680in}}%
\pgfpathlineto{\pgfqpoint{2.725680in}{2.147258in}}%
\pgfpathlineto{\pgfqpoint{2.692178in}{2.168801in}}%
\pgfpathclose%
\pgfusepath{stroke,fill}%
\end{pgfscope}%
\begin{pgfscope}%
\pgfpathrectangle{\pgfqpoint{0.050000in}{0.050000in}}{\pgfqpoint{4.900000in}{4.900000in}}%
\pgfusepath{clip}%
\pgfsetbuttcap%
\pgfsetroundjoin%
\definecolor{currentfill}{rgb}{0.223299,0.223299,0.223299}%
\pgfsetfillcolor{currentfill}%
\pgfsetfillopacity{0.900000}%
\pgfsetlinewidth{0.501875pt}%
\definecolor{currentstroke}{rgb}{0.200000,0.200000,0.200000}%
\pgfsetstrokecolor{currentstroke}%
\pgfsetdash{}{0pt}%
\pgfpathmoveto{\pgfqpoint{1.671140in}{2.994614in}}%
\pgfpathlineto{\pgfqpoint{1.720489in}{2.947166in}}%
\pgfpathlineto{\pgfqpoint{1.752357in}{2.926909in}}%
\pgfpathlineto{\pgfqpoint{1.702999in}{2.974388in}}%
\pgfpathlineto{\pgfqpoint{1.671140in}{2.994614in}}%
\pgfpathclose%
\pgfusepath{stroke,fill}%
\end{pgfscope}%
\begin{pgfscope}%
\pgfpathrectangle{\pgfqpoint{0.050000in}{0.050000in}}{\pgfqpoint{4.900000in}{4.900000in}}%
\pgfusepath{clip}%
\pgfsetbuttcap%
\pgfsetroundjoin%
\definecolor{currentfill}{rgb}{0.499962,0.499962,0.499962}%
\pgfsetfillcolor{currentfill}%
\pgfsetfillopacity{0.900000}%
\pgfsetlinewidth{0.501875pt}%
\definecolor{currentstroke}{rgb}{0.200000,0.200000,0.200000}%
\pgfsetstrokecolor{currentstroke}%
\pgfsetdash{}{0pt}%
\pgfpathmoveto{\pgfqpoint{2.416735in}{2.392040in}}%
\pgfpathlineto{\pgfqpoint{2.464447in}{2.345918in}}%
\pgfpathlineto{\pgfqpoint{2.497514in}{2.324694in}}%
\pgfpathlineto{\pgfqpoint{2.449799in}{2.370843in}}%
\pgfpathlineto{\pgfqpoint{2.416735in}{2.392040in}}%
\pgfpathclose%
\pgfusepath{stroke,fill}%
\end{pgfscope}%
\begin{pgfscope}%
\pgfpathrectangle{\pgfqpoint{0.050000in}{0.050000in}}{\pgfqpoint{4.900000in}{4.900000in}}%
\pgfusepath{clip}%
\pgfsetbuttcap%
\pgfsetroundjoin%
\definecolor{currentfill}{rgb}{0.109250,0.109250,0.109250}%
\pgfsetfillcolor{currentfill}%
\pgfsetfillopacity{0.900000}%
\pgfsetlinewidth{0.501875pt}%
\definecolor{currentstroke}{rgb}{0.200000,0.200000,0.200000}%
\pgfsetstrokecolor{currentstroke}%
\pgfsetdash{}{0pt}%
\pgfpathmoveto{\pgfqpoint{1.395489in}{3.218083in}}%
\pgfpathlineto{\pgfqpoint{1.445454in}{3.170136in}}%
\pgfpathlineto{\pgfqpoint{1.476884in}{3.150233in}}%
\pgfpathlineto{\pgfqpoint{1.426908in}{3.198214in}}%
\pgfpathlineto{\pgfqpoint{1.395489in}{3.218083in}}%
\pgfpathclose%
\pgfusepath{stroke,fill}%
\end{pgfscope}%
\begin{pgfscope}%
\pgfpathrectangle{\pgfqpoint{0.050000in}{0.050000in}}{\pgfqpoint{4.900000in}{4.900000in}}%
\pgfusepath{clip}%
\pgfsetbuttcap%
\pgfsetroundjoin%
\definecolor{currentfill}{rgb}{0.767443,0.767443,0.767443}%
\pgfsetfillcolor{currentfill}%
\pgfsetfillopacity{0.900000}%
\pgfsetlinewidth{0.501875pt}%
\definecolor{currentstroke}{rgb}{0.200000,0.200000,0.200000}%
\pgfsetstrokecolor{currentstroke}%
\pgfsetdash{}{0pt}%
\pgfpathmoveto{\pgfqpoint{3.162465in}{1.792289in}}%
\pgfpathlineto{\pgfqpoint{3.208562in}{1.749251in}}%
\pgfpathlineto{\pgfqpoint{3.242836in}{1.727364in}}%
\pgfpathlineto{\pgfqpoint{3.196738in}{1.770334in}}%
\pgfpathlineto{\pgfqpoint{3.162465in}{1.792289in}}%
\pgfpathclose%
\pgfusepath{stroke,fill}%
\end{pgfscope}%
\begin{pgfscope}%
\pgfpathrectangle{\pgfqpoint{0.050000in}{0.050000in}}{\pgfqpoint{4.900000in}{4.900000in}}%
\pgfusepath{clip}%
\pgfsetbuttcap%
\pgfsetroundjoin%
\definecolor{currentfill}{rgb}{0.403783,0.403783,0.403783}%
\pgfsetfillcolor{currentfill}%
\pgfsetfillopacity{0.900000}%
\pgfsetlinewidth{0.501875pt}%
\definecolor{currentstroke}{rgb}{0.200000,0.200000,0.200000}%
\pgfsetstrokecolor{currentstroke}%
\pgfsetdash{}{0pt}%
\pgfpathmoveto{\pgfqpoint{2.141195in}{2.615412in}}%
\pgfpathlineto{\pgfqpoint{2.189522in}{2.568787in}}%
\pgfpathlineto{\pgfqpoint{2.222153in}{2.547917in}}%
\pgfpathlineto{\pgfqpoint{2.173820in}{2.594570in}}%
\pgfpathlineto{\pgfqpoint{2.141195in}{2.615412in}}%
\pgfpathclose%
\pgfusepath{stroke,fill}%
\end{pgfscope}%
\begin{pgfscope}%
\pgfpathrectangle{\pgfqpoint{0.050000in}{0.050000in}}{\pgfqpoint{4.900000in}{4.900000in}}%
\pgfusepath{clip}%
\pgfsetbuttcap%
\pgfsetroundjoin%
\definecolor{currentfill}{rgb}{0.672203,0.672203,0.672203}%
\pgfsetfillcolor{currentfill}%
\pgfsetfillopacity{0.900000}%
\pgfsetlinewidth{0.501875pt}%
\definecolor{currentstroke}{rgb}{0.200000,0.200000,0.200000}%
\pgfsetstrokecolor{currentstroke}%
\pgfsetdash{}{0pt}%
\pgfpathmoveto{\pgfqpoint{2.887008in}{2.013032in}}%
\pgfpathlineto{\pgfqpoint{2.933699in}{1.968049in}}%
\pgfpathlineto{\pgfqpoint{2.967530in}{1.946284in}}%
\pgfpathlineto{\pgfqpoint{2.920839in}{1.991263in}}%
\pgfpathlineto{\pgfqpoint{2.887008in}{2.013032in}}%
\pgfpathclose%
\pgfusepath{stroke,fill}%
\end{pgfscope}%
\begin{pgfscope}%
\pgfpathrectangle{\pgfqpoint{0.050000in}{0.050000in}}{\pgfqpoint{4.900000in}{4.900000in}}%
\pgfusepath{clip}%
\pgfsetbuttcap%
\pgfsetroundjoin%
\definecolor{currentfill}{rgb}{0.004552,0.004552,0.004552}%
\pgfsetfillcolor{currentfill}%
\pgfsetfillopacity{0.900000}%
\pgfsetlinewidth{0.501875pt}%
\definecolor{currentstroke}{rgb}{0.200000,0.200000,0.200000}%
\pgfsetstrokecolor{currentstroke}%
\pgfsetdash{}{0pt}%
\pgfpathmoveto{\pgfqpoint{1.119742in}{3.441634in}}%
\pgfpathlineto{\pgfqpoint{1.170322in}{3.393186in}}%
\pgfpathlineto{\pgfqpoint{1.201315in}{3.373638in}}%
\pgfpathlineto{\pgfqpoint{1.150721in}{3.422121in}}%
\pgfpathlineto{\pgfqpoint{1.119742in}{3.441634in}}%
\pgfpathclose%
\pgfusepath{stroke,fill}%
\end{pgfscope}%
\begin{pgfscope}%
\pgfpathrectangle{\pgfqpoint{0.050000in}{0.050000in}}{\pgfqpoint{4.900000in}{4.900000in}}%
\pgfusepath{clip}%
\pgfsetbuttcap%
\pgfsetroundjoin%
\definecolor{currentfill}{rgb}{0.311880,0.311880,0.311880}%
\pgfsetfillcolor{currentfill}%
\pgfsetfillopacity{0.900000}%
\pgfsetlinewidth{0.501875pt}%
\definecolor{currentstroke}{rgb}{0.200000,0.200000,0.200000}%
\pgfsetstrokecolor{currentstroke}%
\pgfsetdash{}{0pt}%
\pgfpathmoveto{\pgfqpoint{1.865559in}{2.838865in}}%
\pgfpathlineto{\pgfqpoint{1.914501in}{2.791740in}}%
\pgfpathlineto{\pgfqpoint{1.946695in}{2.771225in}}%
\pgfpathlineto{\pgfqpoint{1.897745in}{2.818380in}}%
\pgfpathlineto{\pgfqpoint{1.865559in}{2.838865in}}%
\pgfpathclose%
\pgfusepath{stroke,fill}%
\end{pgfscope}%
\begin{pgfscope}%
\pgfpathrectangle{\pgfqpoint{0.050000in}{0.050000in}}{\pgfqpoint{4.900000in}{4.900000in}}%
\pgfusepath{clip}%
\pgfsetbuttcap%
\pgfsetroundjoin%
\definecolor{currentfill}{rgb}{0.564552,0.564552,0.564552}%
\pgfsetfillcolor{currentfill}%
\pgfsetfillopacity{0.900000}%
\pgfsetlinewidth{0.501875pt}%
\definecolor{currentstroke}{rgb}{0.200000,0.200000,0.200000}%
\pgfsetstrokecolor{currentstroke}%
\pgfsetdash{}{0pt}%
\pgfpathmoveto{\pgfqpoint{2.611481in}{2.236051in}}%
\pgfpathlineto{\pgfqpoint{2.658784in}{2.190277in}}%
\pgfpathlineto{\pgfqpoint{2.692178in}{2.168801in}}%
\pgfpathlineto{\pgfqpoint{2.644873in}{2.214597in}}%
\pgfpathlineto{\pgfqpoint{2.611481in}{2.236051in}}%
\pgfpathclose%
\pgfusepath{stroke,fill}%
\end{pgfscope}%
\begin{pgfscope}%
\pgfpathrectangle{\pgfqpoint{0.050000in}{0.050000in}}{\pgfqpoint{4.900000in}{4.900000in}}%
\pgfusepath{clip}%
\pgfsetbuttcap%
\pgfsetroundjoin%
\definecolor{currentfill}{rgb}{0.184544,0.184544,0.184544}%
\pgfsetfillcolor{currentfill}%
\pgfsetfillopacity{0.900000}%
\pgfsetlinewidth{0.501875pt}%
\definecolor{currentstroke}{rgb}{0.200000,0.200000,0.200000}%
\pgfsetstrokecolor{currentstroke}%
\pgfsetdash{}{0pt}%
\pgfpathmoveto{\pgfqpoint{1.589827in}{3.062400in}}%
\pgfpathlineto{\pgfqpoint{1.639384in}{3.014774in}}%
\pgfpathlineto{\pgfqpoint{1.671140in}{2.994614in}}%
\pgfpathlineto{\pgfqpoint{1.621573in}{3.042271in}}%
\pgfpathlineto{\pgfqpoint{1.589827in}{3.062400in}}%
\pgfpathclose%
\pgfusepath{stroke,fill}%
\end{pgfscope}%
\begin{pgfscope}%
\pgfpathrectangle{\pgfqpoint{0.050000in}{0.050000in}}{\pgfqpoint{4.900000in}{4.900000in}}%
\pgfusepath{clip}%
\pgfsetbuttcap%
\pgfsetroundjoin%
\definecolor{currentfill}{rgb}{0.465513,0.465513,0.465513}%
\pgfsetfillcolor{currentfill}%
\pgfsetfillopacity{0.900000}%
\pgfsetlinewidth{0.501875pt}%
\definecolor{currentstroke}{rgb}{0.200000,0.200000,0.200000}%
\pgfsetstrokecolor{currentstroke}%
\pgfsetdash{}{0pt}%
\pgfpathmoveto{\pgfqpoint{2.335860in}{2.459469in}}%
\pgfpathlineto{\pgfqpoint{2.383778in}{2.413169in}}%
\pgfpathlineto{\pgfqpoint{2.416735in}{2.392040in}}%
\pgfpathlineto{\pgfqpoint{2.368812in}{2.438368in}}%
\pgfpathlineto{\pgfqpoint{2.335860in}{2.459469in}}%
\pgfpathclose%
\pgfusepath{stroke,fill}%
\end{pgfscope}%
\begin{pgfscope}%
\pgfpathrectangle{\pgfqpoint{0.050000in}{0.050000in}}{\pgfqpoint{4.900000in}{4.900000in}}%
\pgfusepath{clip}%
\pgfsetbuttcap%
\pgfsetroundjoin%
\definecolor{currentfill}{rgb}{0.072834,0.072834,0.072834}%
\pgfsetfillcolor{currentfill}%
\pgfsetfillopacity{0.900000}%
\pgfsetlinewidth{0.501875pt}%
\definecolor{currentstroke}{rgb}{0.200000,0.200000,0.200000}%
\pgfsetstrokecolor{currentstroke}%
\pgfsetdash{}{0pt}%
\pgfpathmoveto{\pgfqpoint{1.313997in}{3.286015in}}%
\pgfpathlineto{\pgfqpoint{1.364171in}{3.237889in}}%
\pgfpathlineto{\pgfqpoint{1.395489in}{3.218083in}}%
\pgfpathlineto{\pgfqpoint{1.345304in}{3.266243in}}%
\pgfpathlineto{\pgfqpoint{1.313997in}{3.286015in}}%
\pgfpathclose%
\pgfusepath{stroke,fill}%
\end{pgfscope}%
\begin{pgfscope}%
\pgfpathrectangle{\pgfqpoint{0.050000in}{0.050000in}}{\pgfqpoint{4.900000in}{4.900000in}}%
\pgfusepath{clip}%
\pgfsetbuttcap%
\pgfsetroundjoin%
\definecolor{currentfill}{rgb}{0.743329,0.743329,0.743329}%
\pgfsetfillcolor{currentfill}%
\pgfsetfillopacity{0.900000}%
\pgfsetlinewidth{0.501875pt}%
\definecolor{currentstroke}{rgb}{0.200000,0.200000,0.200000}%
\pgfsetstrokecolor{currentstroke}%
\pgfsetdash{}{0pt}%
\pgfpathmoveto{\pgfqpoint{3.082012in}{1.858067in}}%
\pgfpathlineto{\pgfqpoint{3.128303in}{1.814181in}}%
\pgfpathlineto{\pgfqpoint{3.162465in}{1.792289in}}%
\pgfpathlineto{\pgfqpoint{3.116174in}{1.836130in}}%
\pgfpathlineto{\pgfqpoint{3.082012in}{1.858067in}}%
\pgfpathclose%
\pgfusepath{stroke,fill}%
\end{pgfscope}%
\begin{pgfscope}%
\pgfpathrectangle{\pgfqpoint{0.050000in}{0.050000in}}{\pgfqpoint{4.900000in}{4.900000in}}%
\pgfusepath{clip}%
\pgfsetbuttcap%
\pgfsetroundjoin%
\definecolor{currentfill}{rgb}{0.375363,0.375363,0.375363}%
\pgfsetfillcolor{currentfill}%
\pgfsetfillopacity{0.900000}%
\pgfsetlinewidth{0.501875pt}%
\definecolor{currentstroke}{rgb}{0.200000,0.200000,0.200000}%
\pgfsetstrokecolor{currentstroke}%
\pgfsetdash{}{0pt}%
\pgfpathmoveto{\pgfqpoint{2.060142in}{2.682987in}}%
\pgfpathlineto{\pgfqpoint{2.108676in}{2.636186in}}%
\pgfpathlineto{\pgfqpoint{2.141195in}{2.615412in}}%
\pgfpathlineto{\pgfqpoint{2.092655in}{2.662242in}}%
\pgfpathlineto{\pgfqpoint{2.060142in}{2.682987in}}%
\pgfpathclose%
\pgfusepath{stroke,fill}%
\end{pgfscope}%
\begin{pgfscope}%
\pgfpathrectangle{\pgfqpoint{0.050000in}{0.050000in}}{\pgfqpoint{4.900000in}{4.900000in}}%
\pgfusepath{clip}%
\pgfsetbuttcap%
\pgfsetroundjoin%
\definecolor{currentfill}{rgb}{0.633818,0.633818,0.633818}%
\pgfsetfillcolor{currentfill}%
\pgfsetfillopacity{0.900000}%
\pgfsetlinewidth{0.501875pt}%
\definecolor{currentstroke}{rgb}{0.200000,0.200000,0.200000}%
\pgfsetstrokecolor{currentstroke}%
\pgfsetdash{}{0pt}%
\pgfpathmoveto{\pgfqpoint{2.806392in}{2.080052in}}%
\pgfpathlineto{\pgfqpoint{2.853286in}{2.034735in}}%
\pgfpathlineto{\pgfqpoint{2.887008in}{2.013032in}}%
\pgfpathlineto{\pgfqpoint{2.840112in}{2.058357in}}%
\pgfpathlineto{\pgfqpoint{2.806392in}{2.080052in}}%
\pgfpathclose%
\pgfusepath{stroke,fill}%
\end{pgfscope}%
\begin{pgfscope}%
\pgfpathrectangle{\pgfqpoint{0.050000in}{0.050000in}}{\pgfqpoint{4.900000in}{4.900000in}}%
\pgfusepath{clip}%
\pgfsetbuttcap%
\pgfsetroundjoin%
\definecolor{currentfill}{rgb}{0.267589,0.267589,0.267589}%
\pgfsetfillcolor{currentfill}%
\pgfsetfillopacity{0.900000}%
\pgfsetlinewidth{0.501875pt}%
\definecolor{currentstroke}{rgb}{0.200000,0.200000,0.200000}%
\pgfsetstrokecolor{currentstroke}%
\pgfsetdash{}{0pt}%
\pgfpathmoveto{\pgfqpoint{1.784327in}{2.906586in}}%
\pgfpathlineto{\pgfqpoint{1.833477in}{2.859284in}}%
\pgfpathlineto{\pgfqpoint{1.865559in}{2.838865in}}%
\pgfpathlineto{\pgfqpoint{1.816401in}{2.886198in}}%
\pgfpathlineto{\pgfqpoint{1.784327in}{2.906586in}}%
\pgfpathclose%
\pgfusepath{stroke,fill}%
\end{pgfscope}%
\begin{pgfscope}%
\pgfpathrectangle{\pgfqpoint{0.050000in}{0.050000in}}{\pgfqpoint{4.900000in}{4.900000in}}%
\pgfusepath{clip}%
\pgfsetbuttcap%
\pgfsetroundjoin%
\definecolor{currentfill}{rgb}{0.530104,0.530104,0.530104}%
\pgfsetfillcolor{currentfill}%
\pgfsetfillopacity{0.900000}%
\pgfsetlinewidth{0.501875pt}%
\definecolor{currentstroke}{rgb}{0.200000,0.200000,0.200000}%
\pgfsetstrokecolor{currentstroke}%
\pgfsetdash{}{0pt}%
\pgfpathmoveto{\pgfqpoint{2.530688in}{2.303402in}}%
\pgfpathlineto{\pgfqpoint{2.578197in}{2.257436in}}%
\pgfpathlineto{\pgfqpoint{2.611481in}{2.236051in}}%
\pgfpathlineto{\pgfqpoint{2.563969in}{2.282042in}}%
\pgfpathlineto{\pgfqpoint{2.530688in}{2.303402in}}%
\pgfpathclose%
\pgfusepath{stroke,fill}%
\end{pgfscope}%
\begin{pgfscope}%
\pgfpathrectangle{\pgfqpoint{0.050000in}{0.050000in}}{\pgfqpoint{4.900000in}{4.900000in}}%
\pgfusepath{clip}%
\pgfsetbuttcap%
\pgfsetroundjoin%
\definecolor{currentfill}{rgb}{0.141115,0.141115,0.141115}%
\pgfsetfillcolor{currentfill}%
\pgfsetfillopacity{0.900000}%
\pgfsetlinewidth{0.501875pt}%
\definecolor{currentstroke}{rgb}{0.200000,0.200000,0.200000}%
\pgfsetstrokecolor{currentstroke}%
\pgfsetdash{}{0pt}%
\pgfpathmoveto{\pgfqpoint{1.508416in}{3.130266in}}%
\pgfpathlineto{\pgfqpoint{1.558183in}{3.082463in}}%
\pgfpathlineto{\pgfqpoint{1.589827in}{3.062400in}}%
\pgfpathlineto{\pgfqpoint{1.540050in}{3.110235in}}%
\pgfpathlineto{\pgfqpoint{1.508416in}{3.130266in}}%
\pgfpathclose%
\pgfusepath{stroke,fill}%
\end{pgfscope}%
\begin{pgfscope}%
\pgfpathrectangle{\pgfqpoint{0.050000in}{0.050000in}}{\pgfqpoint{4.900000in}{4.900000in}}%
\pgfusepath{clip}%
\pgfsetbuttcap%
\pgfsetroundjoin%
\definecolor{currentfill}{rgb}{0.436263,0.436263,0.436263}%
\pgfsetfillcolor{currentfill}%
\pgfsetfillopacity{0.900000}%
\pgfsetlinewidth{0.501875pt}%
\definecolor{currentstroke}{rgb}{0.200000,0.200000,0.200000}%
\pgfsetstrokecolor{currentstroke}%
\pgfsetdash{}{0pt}%
\pgfpathmoveto{\pgfqpoint{2.254888in}{2.526979in}}%
\pgfpathlineto{\pgfqpoint{2.303013in}{2.480502in}}%
\pgfpathlineto{\pgfqpoint{2.335860in}{2.459469in}}%
\pgfpathlineto{\pgfqpoint{2.287730in}{2.505974in}}%
\pgfpathlineto{\pgfqpoint{2.254888in}{2.526979in}}%
\pgfpathclose%
\pgfusepath{stroke,fill}%
\end{pgfscope}%
\begin{pgfscope}%
\pgfpathrectangle{\pgfqpoint{0.050000in}{0.050000in}}{\pgfqpoint{4.900000in}{4.900000in}}%
\pgfusepath{clip}%
\pgfsetbuttcap%
\pgfsetroundjoin%
\definecolor{currentfill}{rgb}{0.791557,0.791557,0.791557}%
\pgfsetfillcolor{currentfill}%
\pgfsetfillopacity{0.900000}%
\pgfsetlinewidth{0.501875pt}%
\definecolor{currentstroke}{rgb}{0.200000,0.200000,0.200000}%
\pgfsetstrokecolor{currentstroke}%
\pgfsetdash{}{0pt}%
\pgfpathmoveto{\pgfqpoint{3.277221in}{1.705416in}}%
\pgfpathlineto{\pgfqpoint{3.323131in}{1.663591in}}%
\pgfpathlineto{\pgfqpoint{3.357629in}{1.641674in}}%
\pgfpathlineto{\pgfqpoint{3.311717in}{1.683405in}}%
\pgfpathlineto{\pgfqpoint{3.277221in}{1.705416in}}%
\pgfpathclose%
\pgfusepath{stroke,fill}%
\end{pgfscope}%
\begin{pgfscope}%
\pgfpathrectangle{\pgfqpoint{0.050000in}{0.050000in}}{\pgfqpoint{4.900000in}{4.900000in}}%
\pgfusepath{clip}%
\pgfsetbuttcap%
\pgfsetroundjoin%
\definecolor{currentfill}{rgb}{0.705790,0.705790,0.705790}%
\pgfsetfillcolor{currentfill}%
\pgfsetfillopacity{0.900000}%
\pgfsetlinewidth{0.501875pt}%
\definecolor{currentstroke}{rgb}{0.200000,0.200000,0.200000}%
\pgfsetstrokecolor{currentstroke}%
\pgfsetdash{}{0pt}%
\pgfpathmoveto{\pgfqpoint{3.001472in}{1.924454in}}%
\pgfpathlineto{\pgfqpoint{3.047961in}{1.879939in}}%
\pgfpathlineto{\pgfqpoint{3.082012in}{1.858067in}}%
\pgfpathlineto{\pgfqpoint{3.035522in}{1.902558in}}%
\pgfpathlineto{\pgfqpoint{3.001472in}{1.924454in}}%
\pgfpathclose%
\pgfusepath{stroke,fill}%
\end{pgfscope}%
\begin{pgfscope}%
\pgfpathrectangle{\pgfqpoint{0.050000in}{0.050000in}}{\pgfqpoint{4.900000in}{4.900000in}}%
\pgfusepath{clip}%
\pgfsetbuttcap%
\pgfsetroundjoin%
\definecolor{currentfill}{rgb}{0.040969,0.040969,0.040969}%
\pgfsetfillcolor{currentfill}%
\pgfsetfillopacity{0.900000}%
\pgfsetlinewidth{0.501875pt}%
\definecolor{currentstroke}{rgb}{0.200000,0.200000,0.200000}%
\pgfsetstrokecolor{currentstroke}%
\pgfsetdash{}{0pt}%
\pgfpathmoveto{\pgfqpoint{1.232408in}{3.354027in}}%
\pgfpathlineto{\pgfqpoint{1.282791in}{3.305723in}}%
\pgfpathlineto{\pgfqpoint{1.313997in}{3.286015in}}%
\pgfpathlineto{\pgfqpoint{1.263601in}{3.334354in}}%
\pgfpathlineto{\pgfqpoint{1.232408in}{3.354027in}}%
\pgfpathclose%
\pgfusepath{stroke,fill}%
\end{pgfscope}%
\begin{pgfscope}%
\pgfpathrectangle{\pgfqpoint{0.050000in}{0.050000in}}{\pgfqpoint{4.900000in}{4.900000in}}%
\pgfusepath{clip}%
\pgfsetbuttcap%
\pgfsetroundjoin%
\definecolor{currentfill}{rgb}{0.342884,0.342884,0.342884}%
\pgfsetfillcolor{currentfill}%
\pgfsetfillopacity{0.900000}%
\pgfsetlinewidth{0.501875pt}%
\definecolor{currentstroke}{rgb}{0.200000,0.200000,0.200000}%
\pgfsetstrokecolor{currentstroke}%
\pgfsetdash{}{0pt}%
\pgfpathmoveto{\pgfqpoint{1.978992in}{2.750643in}}%
\pgfpathlineto{\pgfqpoint{2.027733in}{2.703665in}}%
\pgfpathlineto{\pgfqpoint{2.060142in}{2.682987in}}%
\pgfpathlineto{\pgfqpoint{2.011393in}{2.729995in}}%
\pgfpathlineto{\pgfqpoint{1.978992in}{2.750643in}}%
\pgfpathclose%
\pgfusepath{stroke,fill}%
\end{pgfscope}%
\begin{pgfscope}%
\pgfpathrectangle{\pgfqpoint{0.050000in}{0.050000in}}{\pgfqpoint{4.900000in}{4.900000in}}%
\pgfusepath{clip}%
\pgfsetbuttcap%
\pgfsetroundjoin%
\definecolor{currentfill}{rgb}{0.600231,0.600231,0.600231}%
\pgfsetfillcolor{currentfill}%
\pgfsetfillopacity{0.900000}%
\pgfsetlinewidth{0.501875pt}%
\definecolor{currentstroke}{rgb}{0.200000,0.200000,0.200000}%
\pgfsetstrokecolor{currentstroke}%
\pgfsetdash{}{0pt}%
\pgfpathmoveto{\pgfqpoint{2.725680in}{2.147258in}}%
\pgfpathlineto{\pgfqpoint{2.772780in}{2.101680in}}%
\pgfpathlineto{\pgfqpoint{2.806392in}{2.080052in}}%
\pgfpathlineto{\pgfqpoint{2.759291in}{2.125646in}}%
\pgfpathlineto{\pgfqpoint{2.725680in}{2.147258in}}%
\pgfpathclose%
\pgfusepath{stroke,fill}%
\end{pgfscope}%
\begin{pgfscope}%
\pgfpathrectangle{\pgfqpoint{0.050000in}{0.050000in}}{\pgfqpoint{4.900000in}{4.900000in}}%
\pgfusepath{clip}%
\pgfsetbuttcap%
\pgfsetroundjoin%
\definecolor{currentfill}{rgb}{0.228835,0.228835,0.228835}%
\pgfsetfillcolor{currentfill}%
\pgfsetfillopacity{0.900000}%
\pgfsetlinewidth{0.501875pt}%
\definecolor{currentstroke}{rgb}{0.200000,0.200000,0.200000}%
\pgfsetstrokecolor{currentstroke}%
\pgfsetdash{}{0pt}%
\pgfpathmoveto{\pgfqpoint{1.702999in}{2.974388in}}%
\pgfpathlineto{\pgfqpoint{1.752357in}{2.926909in}}%
\pgfpathlineto{\pgfqpoint{1.784327in}{2.906586in}}%
\pgfpathlineto{\pgfqpoint{1.734960in}{2.954097in}}%
\pgfpathlineto{\pgfqpoint{1.702999in}{2.974388in}}%
\pgfpathclose%
\pgfusepath{stroke,fill}%
\end{pgfscope}%
\begin{pgfscope}%
\pgfpathrectangle{\pgfqpoint{0.050000in}{0.050000in}}{\pgfqpoint{4.900000in}{4.900000in}}%
\pgfusepath{clip}%
\pgfsetbuttcap%
\pgfsetroundjoin%
\definecolor{currentfill}{rgb}{0.499962,0.499962,0.499962}%
\pgfsetfillcolor{currentfill}%
\pgfsetfillopacity{0.900000}%
\pgfsetlinewidth{0.501875pt}%
\definecolor{currentstroke}{rgb}{0.200000,0.200000,0.200000}%
\pgfsetstrokecolor{currentstroke}%
\pgfsetdash{}{0pt}%
\pgfpathmoveto{\pgfqpoint{2.449799in}{2.370843in}}%
\pgfpathlineto{\pgfqpoint{2.497514in}{2.324694in}}%
\pgfpathlineto{\pgfqpoint{2.530688in}{2.303402in}}%
\pgfpathlineto{\pgfqpoint{2.482969in}{2.349578in}}%
\pgfpathlineto{\pgfqpoint{2.449799in}{2.370843in}}%
\pgfpathclose%
\pgfusepath{stroke,fill}%
\end{pgfscope}%
\begin{pgfscope}%
\pgfpathrectangle{\pgfqpoint{0.050000in}{0.050000in}}{\pgfqpoint{4.900000in}{4.900000in}}%
\pgfusepath{clip}%
\pgfsetbuttcap%
\pgfsetroundjoin%
\definecolor{currentfill}{rgb}{0.109250,0.109250,0.109250}%
\pgfsetfillcolor{currentfill}%
\pgfsetfillopacity{0.900000}%
\pgfsetlinewidth{0.501875pt}%
\definecolor{currentstroke}{rgb}{0.200000,0.200000,0.200000}%
\pgfsetstrokecolor{currentstroke}%
\pgfsetdash{}{0pt}%
\pgfpathmoveto{\pgfqpoint{1.426908in}{3.198214in}}%
\pgfpathlineto{\pgfqpoint{1.476884in}{3.150233in}}%
\pgfpathlineto{\pgfqpoint{1.508416in}{3.130266in}}%
\pgfpathlineto{\pgfqpoint{1.458429in}{3.178281in}}%
\pgfpathlineto{\pgfqpoint{1.426908in}{3.198214in}}%
\pgfpathclose%
\pgfusepath{stroke,fill}%
\end{pgfscope}%
\begin{pgfscope}%
\pgfpathrectangle{\pgfqpoint{0.050000in}{0.050000in}}{\pgfqpoint{4.900000in}{4.900000in}}%
\pgfusepath{clip}%
\pgfsetbuttcap%
\pgfsetroundjoin%
\definecolor{currentfill}{rgb}{0.767443,0.767443,0.767443}%
\pgfsetfillcolor{currentfill}%
\pgfsetfillopacity{0.900000}%
\pgfsetlinewidth{0.501875pt}%
\definecolor{currentstroke}{rgb}{0.200000,0.200000,0.200000}%
\pgfsetstrokecolor{currentstroke}%
\pgfsetdash{}{0pt}%
\pgfpathmoveto{\pgfqpoint{3.196738in}{1.770334in}}%
\pgfpathlineto{\pgfqpoint{3.242836in}{1.727364in}}%
\pgfpathlineto{\pgfqpoint{3.277221in}{1.705416in}}%
\pgfpathlineto{\pgfqpoint{3.231122in}{1.748315in}}%
\pgfpathlineto{\pgfqpoint{3.196738in}{1.770334in}}%
\pgfpathclose%
\pgfusepath{stroke,fill}%
\end{pgfscope}%
\begin{pgfscope}%
\pgfpathrectangle{\pgfqpoint{0.050000in}{0.050000in}}{\pgfqpoint{4.900000in}{4.900000in}}%
\pgfusepath{clip}%
\pgfsetbuttcap%
\pgfsetroundjoin%
\definecolor{currentfill}{rgb}{0.403783,0.403783,0.403783}%
\pgfsetfillcolor{currentfill}%
\pgfsetfillopacity{0.900000}%
\pgfsetlinewidth{0.501875pt}%
\definecolor{currentstroke}{rgb}{0.200000,0.200000,0.200000}%
\pgfsetstrokecolor{currentstroke}%
\pgfsetdash{}{0pt}%
\pgfpathmoveto{\pgfqpoint{2.173820in}{2.594570in}}%
\pgfpathlineto{\pgfqpoint{2.222153in}{2.547917in}}%
\pgfpathlineto{\pgfqpoint{2.254888in}{2.526979in}}%
\pgfpathlineto{\pgfqpoint{2.206550in}{2.573662in}}%
\pgfpathlineto{\pgfqpoint{2.173820in}{2.594570in}}%
\pgfpathclose%
\pgfusepath{stroke,fill}%
\end{pgfscope}%
\begin{pgfscope}%
\pgfpathrectangle{\pgfqpoint{0.050000in}{0.050000in}}{\pgfqpoint{4.900000in}{4.900000in}}%
\pgfusepath{clip}%
\pgfsetbuttcap%
\pgfsetroundjoin%
\definecolor{currentfill}{rgb}{0.672203,0.672203,0.672203}%
\pgfsetfillcolor{currentfill}%
\pgfsetfillopacity{0.900000}%
\pgfsetlinewidth{0.501875pt}%
\definecolor{currentstroke}{rgb}{0.200000,0.200000,0.200000}%
\pgfsetstrokecolor{currentstroke}%
\pgfsetdash{}{0pt}%
\pgfpathmoveto{\pgfqpoint{2.920839in}{1.991263in}}%
\pgfpathlineto{\pgfqpoint{2.967530in}{1.946284in}}%
\pgfpathlineto{\pgfqpoint{3.001472in}{1.924454in}}%
\pgfpathlineto{\pgfqpoint{2.954779in}{1.969426in}}%
\pgfpathlineto{\pgfqpoint{2.920839in}{1.991263in}}%
\pgfpathclose%
\pgfusepath{stroke,fill}%
\end{pgfscope}%
\begin{pgfscope}%
\pgfpathrectangle{\pgfqpoint{0.050000in}{0.050000in}}{\pgfqpoint{4.900000in}{4.900000in}}%
\pgfusepath{clip}%
\pgfsetbuttcap%
\pgfsetroundjoin%
\definecolor{currentfill}{rgb}{0.004552,0.004552,0.004552}%
\pgfsetfillcolor{currentfill}%
\pgfsetfillopacity{0.900000}%
\pgfsetlinewidth{0.501875pt}%
\definecolor{currentstroke}{rgb}{0.200000,0.200000,0.200000}%
\pgfsetstrokecolor{currentstroke}%
\pgfsetdash{}{0pt}%
\pgfpathmoveto{\pgfqpoint{1.150721in}{3.422121in}}%
\pgfpathlineto{\pgfqpoint{1.201315in}{3.373638in}}%
\pgfpathlineto{\pgfqpoint{1.232408in}{3.354027in}}%
\pgfpathlineto{\pgfqpoint{1.181801in}{3.402546in}}%
\pgfpathlineto{\pgfqpoint{1.150721in}{3.422121in}}%
\pgfpathclose%
\pgfusepath{stroke,fill}%
\end{pgfscope}%
\begin{pgfscope}%
\pgfpathrectangle{\pgfqpoint{0.050000in}{0.050000in}}{\pgfqpoint{4.900000in}{4.900000in}}%
\pgfusepath{clip}%
\pgfsetbuttcap%
\pgfsetroundjoin%
\definecolor{currentfill}{rgb}{0.311880,0.311880,0.311880}%
\pgfsetfillcolor{currentfill}%
\pgfsetfillopacity{0.900000}%
\pgfsetlinewidth{0.501875pt}%
\definecolor{currentstroke}{rgb}{0.200000,0.200000,0.200000}%
\pgfsetstrokecolor{currentstroke}%
\pgfsetdash{}{0pt}%
\pgfpathmoveto{\pgfqpoint{1.897745in}{2.818380in}}%
\pgfpathlineto{\pgfqpoint{1.946695in}{2.771225in}}%
\pgfpathlineto{\pgfqpoint{1.978992in}{2.750643in}}%
\pgfpathlineto{\pgfqpoint{1.930035in}{2.797829in}}%
\pgfpathlineto{\pgfqpoint{1.897745in}{2.818380in}}%
\pgfpathclose%
\pgfusepath{stroke,fill}%
\end{pgfscope}%
\begin{pgfscope}%
\pgfpathrectangle{\pgfqpoint{0.050000in}{0.050000in}}{\pgfqpoint{4.900000in}{4.900000in}}%
\pgfusepath{clip}%
\pgfsetbuttcap%
\pgfsetroundjoin%
\definecolor{currentfill}{rgb}{0.564552,0.564552,0.564552}%
\pgfsetfillcolor{currentfill}%
\pgfsetfillopacity{0.900000}%
\pgfsetlinewidth{0.501875pt}%
\definecolor{currentstroke}{rgb}{0.200000,0.200000,0.200000}%
\pgfsetstrokecolor{currentstroke}%
\pgfsetdash{}{0pt}%
\pgfpathmoveto{\pgfqpoint{2.644873in}{2.214597in}}%
\pgfpathlineto{\pgfqpoint{2.692178in}{2.168801in}}%
\pgfpathlineto{\pgfqpoint{2.725680in}{2.147258in}}%
\pgfpathlineto{\pgfqpoint{2.678372in}{2.193075in}}%
\pgfpathlineto{\pgfqpoint{2.644873in}{2.214597in}}%
\pgfpathclose%
\pgfusepath{stroke,fill}%
\end{pgfscope}%
\begin{pgfscope}%
\pgfpathrectangle{\pgfqpoint{0.050000in}{0.050000in}}{\pgfqpoint{4.900000in}{4.900000in}}%
\pgfusepath{clip}%
\pgfsetbuttcap%
\pgfsetroundjoin%
\definecolor{currentfill}{rgb}{0.184544,0.184544,0.184544}%
\pgfsetfillcolor{currentfill}%
\pgfsetfillopacity{0.900000}%
\pgfsetlinewidth{0.501875pt}%
\definecolor{currentstroke}{rgb}{0.200000,0.200000,0.200000}%
\pgfsetstrokecolor{currentstroke}%
\pgfsetdash{}{0pt}%
\pgfpathmoveto{\pgfqpoint{1.621573in}{3.042271in}}%
\pgfpathlineto{\pgfqpoint{1.671140in}{2.994614in}}%
\pgfpathlineto{\pgfqpoint{1.702999in}{2.974388in}}%
\pgfpathlineto{\pgfqpoint{1.653422in}{3.022078in}}%
\pgfpathlineto{\pgfqpoint{1.621573in}{3.042271in}}%
\pgfpathclose%
\pgfusepath{stroke,fill}%
\end{pgfscope}%
\begin{pgfscope}%
\pgfpathrectangle{\pgfqpoint{0.050000in}{0.050000in}}{\pgfqpoint{4.900000in}{4.900000in}}%
\pgfusepath{clip}%
\pgfsetbuttcap%
\pgfsetroundjoin%
\definecolor{currentfill}{rgb}{0.465513,0.465513,0.465513}%
\pgfsetfillcolor{currentfill}%
\pgfsetfillopacity{0.900000}%
\pgfsetlinewidth{0.501875pt}%
\definecolor{currentstroke}{rgb}{0.200000,0.200000,0.200000}%
\pgfsetstrokecolor{currentstroke}%
\pgfsetdash{}{0pt}%
\pgfpathmoveto{\pgfqpoint{2.368812in}{2.438368in}}%
\pgfpathlineto{\pgfqpoint{2.416735in}{2.392040in}}%
\pgfpathlineto{\pgfqpoint{2.449799in}{2.370843in}}%
\pgfpathlineto{\pgfqpoint{2.401872in}{2.417199in}}%
\pgfpathlineto{\pgfqpoint{2.368812in}{2.438368in}}%
\pgfpathclose%
\pgfusepath{stroke,fill}%
\end{pgfscope}%
\begin{pgfscope}%
\pgfpathrectangle{\pgfqpoint{0.050000in}{0.050000in}}{\pgfqpoint{4.900000in}{4.900000in}}%
\pgfusepath{clip}%
\pgfsetbuttcap%
\pgfsetroundjoin%
\definecolor{currentfill}{rgb}{0.072834,0.072834,0.072834}%
\pgfsetfillcolor{currentfill}%
\pgfsetfillopacity{0.900000}%
\pgfsetlinewidth{0.501875pt}%
\definecolor{currentstroke}{rgb}{0.200000,0.200000,0.200000}%
\pgfsetstrokecolor{currentstroke}%
\pgfsetdash{}{0pt}%
\pgfpathmoveto{\pgfqpoint{1.345304in}{3.266243in}}%
\pgfpathlineto{\pgfqpoint{1.395489in}{3.218083in}}%
\pgfpathlineto{\pgfqpoint{1.426908in}{3.198214in}}%
\pgfpathlineto{\pgfqpoint{1.376711in}{3.246408in}}%
\pgfpathlineto{\pgfqpoint{1.345304in}{3.266243in}}%
\pgfpathclose%
\pgfusepath{stroke,fill}%
\end{pgfscope}%
\begin{pgfscope}%
\pgfpathrectangle{\pgfqpoint{0.050000in}{0.050000in}}{\pgfqpoint{4.900000in}{4.900000in}}%
\pgfusepath{clip}%
\pgfsetbuttcap%
\pgfsetroundjoin%
\definecolor{currentfill}{rgb}{0.743329,0.743329,0.743329}%
\pgfsetfillcolor{currentfill}%
\pgfsetfillopacity{0.900000}%
\pgfsetlinewidth{0.501875pt}%
\definecolor{currentstroke}{rgb}{0.200000,0.200000,0.200000}%
\pgfsetstrokecolor{currentstroke}%
\pgfsetdash{}{0pt}%
\pgfpathmoveto{\pgfqpoint{3.116174in}{1.836130in}}%
\pgfpathlineto{\pgfqpoint{3.162465in}{1.792289in}}%
\pgfpathlineto{\pgfqpoint{3.196738in}{1.770334in}}%
\pgfpathlineto{\pgfqpoint{3.150447in}{1.814129in}}%
\pgfpathlineto{\pgfqpoint{3.116174in}{1.836130in}}%
\pgfpathclose%
\pgfusepath{stroke,fill}%
\end{pgfscope}%
\begin{pgfscope}%
\pgfpathrectangle{\pgfqpoint{0.050000in}{0.050000in}}{\pgfqpoint{4.900000in}{4.900000in}}%
\pgfusepath{clip}%
\pgfsetbuttcap%
\pgfsetroundjoin%
\definecolor{currentfill}{rgb}{0.375363,0.375363,0.375363}%
\pgfsetfillcolor{currentfill}%
\pgfsetfillopacity{0.900000}%
\pgfsetlinewidth{0.501875pt}%
\definecolor{currentstroke}{rgb}{0.200000,0.200000,0.200000}%
\pgfsetstrokecolor{currentstroke}%
\pgfsetdash{}{0pt}%
\pgfpathmoveto{\pgfqpoint{2.092655in}{2.662242in}}%
\pgfpathlineto{\pgfqpoint{2.141195in}{2.615412in}}%
\pgfpathlineto{\pgfqpoint{2.173820in}{2.594570in}}%
\pgfpathlineto{\pgfqpoint{2.125274in}{2.641430in}}%
\pgfpathlineto{\pgfqpoint{2.092655in}{2.662242in}}%
\pgfpathclose%
\pgfusepath{stroke,fill}%
\end{pgfscope}%
\begin{pgfscope}%
\pgfpathrectangle{\pgfqpoint{0.050000in}{0.050000in}}{\pgfqpoint{4.900000in}{4.900000in}}%
\pgfusepath{clip}%
\pgfsetbuttcap%
\pgfsetroundjoin%
\definecolor{currentfill}{rgb}{0.633818,0.633818,0.633818}%
\pgfsetfillcolor{currentfill}%
\pgfsetfillopacity{0.900000}%
\pgfsetlinewidth{0.501875pt}%
\definecolor{currentstroke}{rgb}{0.200000,0.200000,0.200000}%
\pgfsetstrokecolor{currentstroke}%
\pgfsetdash{}{0pt}%
\pgfpathmoveto{\pgfqpoint{2.840112in}{2.058357in}}%
\pgfpathlineto{\pgfqpoint{2.887008in}{2.013032in}}%
\pgfpathlineto{\pgfqpoint{2.920839in}{1.991263in}}%
\pgfpathlineto{\pgfqpoint{2.873942in}{2.036593in}}%
\pgfpathlineto{\pgfqpoint{2.840112in}{2.058357in}}%
\pgfpathclose%
\pgfusepath{stroke,fill}%
\end{pgfscope}%
\begin{pgfscope}%
\pgfpathrectangle{\pgfqpoint{0.050000in}{0.050000in}}{\pgfqpoint{4.900000in}{4.900000in}}%
\pgfusepath{clip}%
\pgfsetbuttcap%
\pgfsetroundjoin%
\definecolor{currentfill}{rgb}{0.267589,0.267589,0.267589}%
\pgfsetfillcolor{currentfill}%
\pgfsetfillopacity{0.900000}%
\pgfsetlinewidth{0.501875pt}%
\definecolor{currentstroke}{rgb}{0.200000,0.200000,0.200000}%
\pgfsetstrokecolor{currentstroke}%
\pgfsetdash{}{0pt}%
\pgfpathmoveto{\pgfqpoint{1.816401in}{2.886198in}}%
\pgfpathlineto{\pgfqpoint{1.865559in}{2.838865in}}%
\pgfpathlineto{\pgfqpoint{1.897745in}{2.818380in}}%
\pgfpathlineto{\pgfqpoint{1.848578in}{2.865744in}}%
\pgfpathlineto{\pgfqpoint{1.816401in}{2.886198in}}%
\pgfpathclose%
\pgfusepath{stroke,fill}%
\end{pgfscope}%
\begin{pgfscope}%
\pgfpathrectangle{\pgfqpoint{0.050000in}{0.050000in}}{\pgfqpoint{4.900000in}{4.900000in}}%
\pgfusepath{clip}%
\pgfsetbuttcap%
\pgfsetroundjoin%
\definecolor{currentfill}{rgb}{0.534410,0.534410,0.534410}%
\pgfsetfillcolor{currentfill}%
\pgfsetfillopacity{0.900000}%
\pgfsetlinewidth{0.501875pt}%
\definecolor{currentstroke}{rgb}{0.200000,0.200000,0.200000}%
\pgfsetstrokecolor{currentstroke}%
\pgfsetdash{}{0pt}%
\pgfpathmoveto{\pgfqpoint{2.563969in}{2.282042in}}%
\pgfpathlineto{\pgfqpoint{2.611481in}{2.236051in}}%
\pgfpathlineto{\pgfqpoint{2.644873in}{2.214597in}}%
\pgfpathlineto{\pgfqpoint{2.597358in}{2.260613in}}%
\pgfpathlineto{\pgfqpoint{2.563969in}{2.282042in}}%
\pgfpathclose%
\pgfusepath{stroke,fill}%
\end{pgfscope}%
\begin{pgfscope}%
\pgfpathrectangle{\pgfqpoint{0.050000in}{0.050000in}}{\pgfqpoint{4.900000in}{4.900000in}}%
\pgfusepath{clip}%
\pgfsetbuttcap%
\pgfsetroundjoin%
\definecolor{currentfill}{rgb}{0.141115,0.141115,0.141115}%
\pgfsetfillcolor{currentfill}%
\pgfsetfillopacity{0.900000}%
\pgfsetlinewidth{0.501875pt}%
\definecolor{currentstroke}{rgb}{0.200000,0.200000,0.200000}%
\pgfsetstrokecolor{currentstroke}%
\pgfsetdash{}{0pt}%
\pgfpathmoveto{\pgfqpoint{1.540050in}{3.110235in}}%
\pgfpathlineto{\pgfqpoint{1.589827in}{3.062400in}}%
\pgfpathlineto{\pgfqpoint{1.621573in}{3.042271in}}%
\pgfpathlineto{\pgfqpoint{1.571786in}{3.090140in}}%
\pgfpathlineto{\pgfqpoint{1.540050in}{3.110235in}}%
\pgfpathclose%
\pgfusepath{stroke,fill}%
\end{pgfscope}%
\begin{pgfscope}%
\pgfpathrectangle{\pgfqpoint{0.050000in}{0.050000in}}{\pgfqpoint{4.900000in}{4.900000in}}%
\pgfusepath{clip}%
\pgfsetbuttcap%
\pgfsetroundjoin%
\definecolor{currentfill}{rgb}{0.436263,0.436263,0.436263}%
\pgfsetfillcolor{currentfill}%
\pgfsetfillopacity{0.900000}%
\pgfsetlinewidth{0.501875pt}%
\definecolor{currentstroke}{rgb}{0.200000,0.200000,0.200000}%
\pgfsetstrokecolor{currentstroke}%
\pgfsetdash{}{0pt}%
\pgfpathmoveto{\pgfqpoint{2.287730in}{2.505974in}}%
\pgfpathlineto{\pgfqpoint{2.335860in}{2.459469in}}%
\pgfpathlineto{\pgfqpoint{2.368812in}{2.438368in}}%
\pgfpathlineto{\pgfqpoint{2.320677in}{2.484901in}}%
\pgfpathlineto{\pgfqpoint{2.287730in}{2.505974in}}%
\pgfpathclose%
\pgfusepath{stroke,fill}%
\end{pgfscope}%
\begin{pgfscope}%
\pgfpathrectangle{\pgfqpoint{0.050000in}{0.050000in}}{\pgfqpoint{4.900000in}{4.900000in}}%
\pgfusepath{clip}%
\pgfsetbuttcap%
\pgfsetroundjoin%
\definecolor{currentfill}{rgb}{0.705790,0.705790,0.705790}%
\pgfsetfillcolor{currentfill}%
\pgfsetfillopacity{0.900000}%
\pgfsetlinewidth{0.501875pt}%
\definecolor{currentstroke}{rgb}{0.200000,0.200000,0.200000}%
\pgfsetstrokecolor{currentstroke}%
\pgfsetdash{}{0pt}%
\pgfpathmoveto{\pgfqpoint{3.035522in}{1.902558in}}%
\pgfpathlineto{\pgfqpoint{3.082012in}{1.858067in}}%
\pgfpathlineto{\pgfqpoint{3.116174in}{1.836130in}}%
\pgfpathlineto{\pgfqpoint{3.069684in}{1.880597in}}%
\pgfpathlineto{\pgfqpoint{3.035522in}{1.902558in}}%
\pgfpathclose%
\pgfusepath{stroke,fill}%
\end{pgfscope}%
\begin{pgfscope}%
\pgfpathrectangle{\pgfqpoint{0.050000in}{0.050000in}}{\pgfqpoint{4.900000in}{4.900000in}}%
\pgfusepath{clip}%
\pgfsetbuttcap%
\pgfsetroundjoin%
\definecolor{currentfill}{rgb}{0.040969,0.040969,0.040969}%
\pgfsetfillcolor{currentfill}%
\pgfsetfillopacity{0.900000}%
\pgfsetlinewidth{0.501875pt}%
\definecolor{currentstroke}{rgb}{0.200000,0.200000,0.200000}%
\pgfsetstrokecolor{currentstroke}%
\pgfsetdash{}{0pt}%
\pgfpathmoveto{\pgfqpoint{1.263601in}{3.334354in}}%
\pgfpathlineto{\pgfqpoint{1.313997in}{3.286015in}}%
\pgfpathlineto{\pgfqpoint{1.345304in}{3.266243in}}%
\pgfpathlineto{\pgfqpoint{1.294896in}{3.314617in}}%
\pgfpathlineto{\pgfqpoint{1.263601in}{3.334354in}}%
\pgfpathclose%
\pgfusepath{stroke,fill}%
\end{pgfscope}%
\begin{pgfscope}%
\pgfpathrectangle{\pgfqpoint{0.050000in}{0.050000in}}{\pgfqpoint{4.900000in}{4.900000in}}%
\pgfusepath{clip}%
\pgfsetbuttcap%
\pgfsetroundjoin%
\definecolor{currentfill}{rgb}{0.342884,0.342884,0.342884}%
\pgfsetfillcolor{currentfill}%
\pgfsetfillopacity{0.900000}%
\pgfsetlinewidth{0.501875pt}%
\definecolor{currentstroke}{rgb}{0.200000,0.200000,0.200000}%
\pgfsetstrokecolor{currentstroke}%
\pgfsetdash{}{0pt}%
\pgfpathmoveto{\pgfqpoint{2.011393in}{2.729995in}}%
\pgfpathlineto{\pgfqpoint{2.060142in}{2.682987in}}%
\pgfpathlineto{\pgfqpoint{2.092655in}{2.662242in}}%
\pgfpathlineto{\pgfqpoint{2.043900in}{2.709280in}}%
\pgfpathlineto{\pgfqpoint{2.011393in}{2.729995in}}%
\pgfpathclose%
\pgfusepath{stroke,fill}%
\end{pgfscope}%
\begin{pgfscope}%
\pgfpathrectangle{\pgfqpoint{0.050000in}{0.050000in}}{\pgfqpoint{4.900000in}{4.900000in}}%
\pgfusepath{clip}%
\pgfsetbuttcap%
\pgfsetroundjoin%
\definecolor{currentfill}{rgb}{0.600231,0.600231,0.600231}%
\pgfsetfillcolor{currentfill}%
\pgfsetfillopacity{0.900000}%
\pgfsetlinewidth{0.501875pt}%
\definecolor{currentstroke}{rgb}{0.200000,0.200000,0.200000}%
\pgfsetstrokecolor{currentstroke}%
\pgfsetdash{}{0pt}%
\pgfpathmoveto{\pgfqpoint{2.759291in}{2.125646in}}%
\pgfpathlineto{\pgfqpoint{2.806392in}{2.080052in}}%
\pgfpathlineto{\pgfqpoint{2.840112in}{2.058357in}}%
\pgfpathlineto{\pgfqpoint{2.793010in}{2.103966in}}%
\pgfpathlineto{\pgfqpoint{2.759291in}{2.125646in}}%
\pgfpathclose%
\pgfusepath{stroke,fill}%
\end{pgfscope}%
\begin{pgfscope}%
\pgfpathrectangle{\pgfqpoint{0.050000in}{0.050000in}}{\pgfqpoint{4.900000in}{4.900000in}}%
\pgfusepath{clip}%
\pgfsetbuttcap%
\pgfsetroundjoin%
\definecolor{currentfill}{rgb}{0.228835,0.228835,0.228835}%
\pgfsetfillcolor{currentfill}%
\pgfsetfillopacity{0.900000}%
\pgfsetlinewidth{0.501875pt}%
\definecolor{currentstroke}{rgb}{0.200000,0.200000,0.200000}%
\pgfsetstrokecolor{currentstroke}%
\pgfsetdash{}{0pt}%
\pgfpathmoveto{\pgfqpoint{1.734960in}{2.954097in}}%
\pgfpathlineto{\pgfqpoint{1.784327in}{2.906586in}}%
\pgfpathlineto{\pgfqpoint{1.816401in}{2.886198in}}%
\pgfpathlineto{\pgfqpoint{1.767025in}{2.933741in}}%
\pgfpathlineto{\pgfqpoint{1.734960in}{2.954097in}}%
\pgfpathclose%
\pgfusepath{stroke,fill}%
\end{pgfscope}%
\begin{pgfscope}%
\pgfpathrectangle{\pgfqpoint{0.050000in}{0.050000in}}{\pgfqpoint{4.900000in}{4.900000in}}%
\pgfusepath{clip}%
\pgfsetbuttcap%
\pgfsetroundjoin%
\definecolor{currentfill}{rgb}{0.499962,0.499962,0.499962}%
\pgfsetfillcolor{currentfill}%
\pgfsetfillopacity{0.900000}%
\pgfsetlinewidth{0.501875pt}%
\definecolor{currentstroke}{rgb}{0.200000,0.200000,0.200000}%
\pgfsetstrokecolor{currentstroke}%
\pgfsetdash{}{0pt}%
\pgfpathmoveto{\pgfqpoint{2.482969in}{2.349578in}}%
\pgfpathlineto{\pgfqpoint{2.530688in}{2.303402in}}%
\pgfpathlineto{\pgfqpoint{2.563969in}{2.282042in}}%
\pgfpathlineto{\pgfqpoint{2.516246in}{2.328244in}}%
\pgfpathlineto{\pgfqpoint{2.482969in}{2.349578in}}%
\pgfpathclose%
\pgfusepath{stroke,fill}%
\end{pgfscope}%
\begin{pgfscope}%
\pgfpathrectangle{\pgfqpoint{0.050000in}{0.050000in}}{\pgfqpoint{4.900000in}{4.900000in}}%
\pgfusepath{clip}%
\pgfsetbuttcap%
\pgfsetroundjoin%
\definecolor{currentfill}{rgb}{0.109250,0.109250,0.109250}%
\pgfsetfillcolor{currentfill}%
\pgfsetfillopacity{0.900000}%
\pgfsetlinewidth{0.501875pt}%
\definecolor{currentstroke}{rgb}{0.200000,0.200000,0.200000}%
\pgfsetstrokecolor{currentstroke}%
\pgfsetdash{}{0pt}%
\pgfpathmoveto{\pgfqpoint{1.458429in}{3.178281in}}%
\pgfpathlineto{\pgfqpoint{1.508416in}{3.130266in}}%
\pgfpathlineto{\pgfqpoint{1.540050in}{3.110235in}}%
\pgfpathlineto{\pgfqpoint{1.490052in}{3.158283in}}%
\pgfpathlineto{\pgfqpoint{1.458429in}{3.178281in}}%
\pgfpathclose%
\pgfusepath{stroke,fill}%
\end{pgfscope}%
\begin{pgfscope}%
\pgfpathrectangle{\pgfqpoint{0.050000in}{0.050000in}}{\pgfqpoint{4.900000in}{4.900000in}}%
\pgfusepath{clip}%
\pgfsetbuttcap%
\pgfsetroundjoin%
\definecolor{currentfill}{rgb}{0.767443,0.767443,0.767443}%
\pgfsetfillcolor{currentfill}%
\pgfsetfillopacity{0.900000}%
\pgfsetlinewidth{0.501875pt}%
\definecolor{currentstroke}{rgb}{0.200000,0.200000,0.200000}%
\pgfsetstrokecolor{currentstroke}%
\pgfsetdash{}{0pt}%
\pgfpathmoveto{\pgfqpoint{3.231122in}{1.748315in}}%
\pgfpathlineto{\pgfqpoint{3.277221in}{1.705416in}}%
\pgfpathlineto{\pgfqpoint{3.311717in}{1.683405in}}%
\pgfpathlineto{\pgfqpoint{3.265618in}{1.726233in}}%
\pgfpathlineto{\pgfqpoint{3.231122in}{1.748315in}}%
\pgfpathclose%
\pgfusepath{stroke,fill}%
\end{pgfscope}%
\begin{pgfscope}%
\pgfpathrectangle{\pgfqpoint{0.050000in}{0.050000in}}{\pgfqpoint{4.900000in}{4.900000in}}%
\pgfusepath{clip}%
\pgfsetbuttcap%
\pgfsetroundjoin%
\definecolor{currentfill}{rgb}{0.403783,0.403783,0.403783}%
\pgfsetfillcolor{currentfill}%
\pgfsetfillopacity{0.900000}%
\pgfsetlinewidth{0.501875pt}%
\definecolor{currentstroke}{rgb}{0.200000,0.200000,0.200000}%
\pgfsetstrokecolor{currentstroke}%
\pgfsetdash{}{0pt}%
\pgfpathmoveto{\pgfqpoint{2.206550in}{2.573662in}}%
\pgfpathlineto{\pgfqpoint{2.254888in}{2.526979in}}%
\pgfpathlineto{\pgfqpoint{2.287730in}{2.505974in}}%
\pgfpathlineto{\pgfqpoint{2.239386in}{2.552686in}}%
\pgfpathlineto{\pgfqpoint{2.206550in}{2.573662in}}%
\pgfpathclose%
\pgfusepath{stroke,fill}%
\end{pgfscope}%
\begin{pgfscope}%
\pgfpathrectangle{\pgfqpoint{0.050000in}{0.050000in}}{\pgfqpoint{4.900000in}{4.900000in}}%
\pgfusepath{clip}%
\pgfsetbuttcap%
\pgfsetroundjoin%
\definecolor{currentfill}{rgb}{0.672203,0.672203,0.672203}%
\pgfsetfillcolor{currentfill}%
\pgfsetfillopacity{0.900000}%
\pgfsetlinewidth{0.501875pt}%
\definecolor{currentstroke}{rgb}{0.200000,0.200000,0.200000}%
\pgfsetstrokecolor{currentstroke}%
\pgfsetdash{}{0pt}%
\pgfpathmoveto{\pgfqpoint{2.954779in}{1.969426in}}%
\pgfpathlineto{\pgfqpoint{3.001472in}{1.924454in}}%
\pgfpathlineto{\pgfqpoint{3.035522in}{1.902558in}}%
\pgfpathlineto{\pgfqpoint{2.988830in}{1.947522in}}%
\pgfpathlineto{\pgfqpoint{2.954779in}{1.969426in}}%
\pgfpathclose%
\pgfusepath{stroke,fill}%
\end{pgfscope}%
\begin{pgfscope}%
\pgfpathrectangle{\pgfqpoint{0.050000in}{0.050000in}}{\pgfqpoint{4.900000in}{4.900000in}}%
\pgfusepath{clip}%
\pgfsetbuttcap%
\pgfsetroundjoin%
\definecolor{currentfill}{rgb}{0.004552,0.004552,0.004552}%
\pgfsetfillcolor{currentfill}%
\pgfsetfillopacity{0.900000}%
\pgfsetlinewidth{0.501875pt}%
\definecolor{currentstroke}{rgb}{0.200000,0.200000,0.200000}%
\pgfsetstrokecolor{currentstroke}%
\pgfsetdash{}{0pt}%
\pgfpathmoveto{\pgfqpoint{1.181801in}{3.402546in}}%
\pgfpathlineto{\pgfqpoint{1.232408in}{3.354027in}}%
\pgfpathlineto{\pgfqpoint{1.263601in}{3.334354in}}%
\pgfpathlineto{\pgfqpoint{1.212982in}{3.382907in}}%
\pgfpathlineto{\pgfqpoint{1.181801in}{3.402546in}}%
\pgfpathclose%
\pgfusepath{stroke,fill}%
\end{pgfscope}%
\begin{pgfscope}%
\pgfpathrectangle{\pgfqpoint{0.050000in}{0.050000in}}{\pgfqpoint{4.900000in}{4.900000in}}%
\pgfusepath{clip}%
\pgfsetbuttcap%
\pgfsetroundjoin%
\definecolor{currentfill}{rgb}{0.311880,0.311880,0.311880}%
\pgfsetfillcolor{currentfill}%
\pgfsetfillopacity{0.900000}%
\pgfsetlinewidth{0.501875pt}%
\definecolor{currentstroke}{rgb}{0.200000,0.200000,0.200000}%
\pgfsetstrokecolor{currentstroke}%
\pgfsetdash{}{0pt}%
\pgfpathmoveto{\pgfqpoint{1.930035in}{2.797829in}}%
\pgfpathlineto{\pgfqpoint{1.978992in}{2.750643in}}%
\pgfpathlineto{\pgfqpoint{2.011393in}{2.729995in}}%
\pgfpathlineto{\pgfqpoint{1.962429in}{2.777212in}}%
\pgfpathlineto{\pgfqpoint{1.930035in}{2.797829in}}%
\pgfpathclose%
\pgfusepath{stroke,fill}%
\end{pgfscope}%
\begin{pgfscope}%
\pgfpathrectangle{\pgfqpoint{0.050000in}{0.050000in}}{\pgfqpoint{4.900000in}{4.900000in}}%
\pgfusepath{clip}%
\pgfsetbuttcap%
\pgfsetroundjoin%
\definecolor{currentfill}{rgb}{0.564552,0.564552,0.564552}%
\pgfsetfillcolor{currentfill}%
\pgfsetfillopacity{0.900000}%
\pgfsetlinewidth{0.501875pt}%
\definecolor{currentstroke}{rgb}{0.200000,0.200000,0.200000}%
\pgfsetstrokecolor{currentstroke}%
\pgfsetdash{}{0pt}%
\pgfpathmoveto{\pgfqpoint{2.678372in}{2.193075in}}%
\pgfpathlineto{\pgfqpoint{2.725680in}{2.147258in}}%
\pgfpathlineto{\pgfqpoint{2.759291in}{2.125646in}}%
\pgfpathlineto{\pgfqpoint{2.711981in}{2.171484in}}%
\pgfpathlineto{\pgfqpoint{2.678372in}{2.193075in}}%
\pgfpathclose%
\pgfusepath{stroke,fill}%
\end{pgfscope}%
\begin{pgfscope}%
\pgfpathrectangle{\pgfqpoint{0.050000in}{0.050000in}}{\pgfqpoint{4.900000in}{4.900000in}}%
\pgfusepath{clip}%
\pgfsetbuttcap%
\pgfsetroundjoin%
\definecolor{currentfill}{rgb}{0.184544,0.184544,0.184544}%
\pgfsetfillcolor{currentfill}%
\pgfsetfillopacity{0.900000}%
\pgfsetlinewidth{0.501875pt}%
\definecolor{currentstroke}{rgb}{0.200000,0.200000,0.200000}%
\pgfsetstrokecolor{currentstroke}%
\pgfsetdash{}{0pt}%
\pgfpathmoveto{\pgfqpoint{1.653422in}{3.022078in}}%
\pgfpathlineto{\pgfqpoint{1.702999in}{2.974388in}}%
\pgfpathlineto{\pgfqpoint{1.734960in}{2.954097in}}%
\pgfpathlineto{\pgfqpoint{1.685374in}{3.001819in}}%
\pgfpathlineto{\pgfqpoint{1.653422in}{3.022078in}}%
\pgfpathclose%
\pgfusepath{stroke,fill}%
\end{pgfscope}%
\begin{pgfscope}%
\pgfpathrectangle{\pgfqpoint{0.050000in}{0.050000in}}{\pgfqpoint{4.900000in}{4.900000in}}%
\pgfusepath{clip}%
\pgfsetbuttcap%
\pgfsetroundjoin%
\definecolor{currentfill}{rgb}{0.469819,0.469819,0.469819}%
\pgfsetfillcolor{currentfill}%
\pgfsetfillopacity{0.900000}%
\pgfsetlinewidth{0.501875pt}%
\definecolor{currentstroke}{rgb}{0.200000,0.200000,0.200000}%
\pgfsetstrokecolor{currentstroke}%
\pgfsetdash{}{0pt}%
\pgfpathmoveto{\pgfqpoint{2.401872in}{2.417199in}}%
\pgfpathlineto{\pgfqpoint{2.449799in}{2.370843in}}%
\pgfpathlineto{\pgfqpoint{2.482969in}{2.349578in}}%
\pgfpathlineto{\pgfqpoint{2.435038in}{2.395961in}}%
\pgfpathlineto{\pgfqpoint{2.401872in}{2.417199in}}%
\pgfpathclose%
\pgfusepath{stroke,fill}%
\end{pgfscope}%
\begin{pgfscope}%
\pgfpathrectangle{\pgfqpoint{0.050000in}{0.050000in}}{\pgfqpoint{4.900000in}{4.900000in}}%
\pgfusepath{clip}%
\pgfsetbuttcap%
\pgfsetroundjoin%
\definecolor{currentfill}{rgb}{0.072834,0.072834,0.072834}%
\pgfsetfillcolor{currentfill}%
\pgfsetfillopacity{0.900000}%
\pgfsetlinewidth{0.501875pt}%
\definecolor{currentstroke}{rgb}{0.200000,0.200000,0.200000}%
\pgfsetstrokecolor{currentstroke}%
\pgfsetdash{}{0pt}%
\pgfpathmoveto{\pgfqpoint{1.376711in}{3.246408in}}%
\pgfpathlineto{\pgfqpoint{1.426908in}{3.198214in}}%
\pgfpathlineto{\pgfqpoint{1.458429in}{3.178281in}}%
\pgfpathlineto{\pgfqpoint{1.408221in}{3.226508in}}%
\pgfpathlineto{\pgfqpoint{1.376711in}{3.246408in}}%
\pgfpathclose%
\pgfusepath{stroke,fill}%
\end{pgfscope}%
\begin{pgfscope}%
\pgfpathrectangle{\pgfqpoint{0.050000in}{0.050000in}}{\pgfqpoint{4.900000in}{4.900000in}}%
\pgfusepath{clip}%
\pgfsetbuttcap%
\pgfsetroundjoin%
\definecolor{currentfill}{rgb}{0.743329,0.743329,0.743329}%
\pgfsetfillcolor{currentfill}%
\pgfsetfillopacity{0.900000}%
\pgfsetlinewidth{0.501875pt}%
\definecolor{currentstroke}{rgb}{0.200000,0.200000,0.200000}%
\pgfsetstrokecolor{currentstroke}%
\pgfsetdash{}{0pt}%
\pgfpathmoveto{\pgfqpoint{3.150447in}{1.814129in}}%
\pgfpathlineto{\pgfqpoint{3.196738in}{1.770334in}}%
\pgfpathlineto{\pgfqpoint{3.231122in}{1.748315in}}%
\pgfpathlineto{\pgfqpoint{3.184830in}{1.792063in}}%
\pgfpathlineto{\pgfqpoint{3.150447in}{1.814129in}}%
\pgfpathclose%
\pgfusepath{stroke,fill}%
\end{pgfscope}%
\begin{pgfscope}%
\pgfpathrectangle{\pgfqpoint{0.050000in}{0.050000in}}{\pgfqpoint{4.900000in}{4.900000in}}%
\pgfusepath{clip}%
\pgfsetbuttcap%
\pgfsetroundjoin%
\definecolor{currentfill}{rgb}{0.375363,0.375363,0.375363}%
\pgfsetfillcolor{currentfill}%
\pgfsetfillopacity{0.900000}%
\pgfsetlinewidth{0.501875pt}%
\definecolor{currentstroke}{rgb}{0.200000,0.200000,0.200000}%
\pgfsetstrokecolor{currentstroke}%
\pgfsetdash{}{0pt}%
\pgfpathmoveto{\pgfqpoint{2.125274in}{2.641430in}}%
\pgfpathlineto{\pgfqpoint{2.173820in}{2.594570in}}%
\pgfpathlineto{\pgfqpoint{2.206550in}{2.573662in}}%
\pgfpathlineto{\pgfqpoint{2.157998in}{2.620552in}}%
\pgfpathlineto{\pgfqpoint{2.125274in}{2.641430in}}%
\pgfpathclose%
\pgfusepath{stroke,fill}%
\end{pgfscope}%
\begin{pgfscope}%
\pgfpathrectangle{\pgfqpoint{0.050000in}{0.050000in}}{\pgfqpoint{4.900000in}{4.900000in}}%
\pgfusepath{clip}%
\pgfsetbuttcap%
\pgfsetroundjoin%
\definecolor{currentfill}{rgb}{0.633818,0.633818,0.633818}%
\pgfsetfillcolor{currentfill}%
\pgfsetfillopacity{0.900000}%
\pgfsetlinewidth{0.501875pt}%
\definecolor{currentstroke}{rgb}{0.200000,0.200000,0.200000}%
\pgfsetstrokecolor{currentstroke}%
\pgfsetdash{}{0pt}%
\pgfpathmoveto{\pgfqpoint{2.873942in}{2.036593in}}%
\pgfpathlineto{\pgfqpoint{2.920839in}{1.991263in}}%
\pgfpathlineto{\pgfqpoint{2.954779in}{1.969426in}}%
\pgfpathlineto{\pgfqpoint{2.907881in}{2.014762in}}%
\pgfpathlineto{\pgfqpoint{2.873942in}{2.036593in}}%
\pgfpathclose%
\pgfusepath{stroke,fill}%
\end{pgfscope}%
\begin{pgfscope}%
\pgfpathrectangle{\pgfqpoint{0.050000in}{0.050000in}}{\pgfqpoint{4.900000in}{4.900000in}}%
\pgfusepath{clip}%
\pgfsetbuttcap%
\pgfsetroundjoin%
\definecolor{currentfill}{rgb}{0.267589,0.267589,0.267589}%
\pgfsetfillcolor{currentfill}%
\pgfsetfillopacity{0.900000}%
\pgfsetlinewidth{0.501875pt}%
\definecolor{currentstroke}{rgb}{0.200000,0.200000,0.200000}%
\pgfsetstrokecolor{currentstroke}%
\pgfsetdash{}{0pt}%
\pgfpathmoveto{\pgfqpoint{1.848578in}{2.865744in}}%
\pgfpathlineto{\pgfqpoint{1.897745in}{2.818380in}}%
\pgfpathlineto{\pgfqpoint{1.930035in}{2.797829in}}%
\pgfpathlineto{\pgfqpoint{1.880860in}{2.845224in}}%
\pgfpathlineto{\pgfqpoint{1.848578in}{2.865744in}}%
\pgfpathclose%
\pgfusepath{stroke,fill}%
\end{pgfscope}%
\begin{pgfscope}%
\pgfpathrectangle{\pgfqpoint{0.050000in}{0.050000in}}{\pgfqpoint{4.900000in}{4.900000in}}%
\pgfusepath{clip}%
\pgfsetbuttcap%
\pgfsetroundjoin%
\definecolor{currentfill}{rgb}{0.534410,0.534410,0.534410}%
\pgfsetfillcolor{currentfill}%
\pgfsetfillopacity{0.900000}%
\pgfsetlinewidth{0.501875pt}%
\definecolor{currentstroke}{rgb}{0.200000,0.200000,0.200000}%
\pgfsetstrokecolor{currentstroke}%
\pgfsetdash{}{0pt}%
\pgfpathmoveto{\pgfqpoint{2.597358in}{2.260613in}}%
\pgfpathlineto{\pgfqpoint{2.644873in}{2.214597in}}%
\pgfpathlineto{\pgfqpoint{2.678372in}{2.193075in}}%
\pgfpathlineto{\pgfqpoint{2.630855in}{2.239115in}}%
\pgfpathlineto{\pgfqpoint{2.597358in}{2.260613in}}%
\pgfpathclose%
\pgfusepath{stroke,fill}%
\end{pgfscope}%
\begin{pgfscope}%
\pgfpathrectangle{\pgfqpoint{0.050000in}{0.050000in}}{\pgfqpoint{4.900000in}{4.900000in}}%
\pgfusepath{clip}%
\pgfsetbuttcap%
\pgfsetroundjoin%
\definecolor{currentfill}{rgb}{0.145790,0.145790,0.145790}%
\pgfsetfillcolor{currentfill}%
\pgfsetfillopacity{0.900000}%
\pgfsetlinewidth{0.501875pt}%
\definecolor{currentstroke}{rgb}{0.200000,0.200000,0.200000}%
\pgfsetstrokecolor{currentstroke}%
\pgfsetdash{}{0pt}%
\pgfpathmoveto{\pgfqpoint{1.571786in}{3.090140in}}%
\pgfpathlineto{\pgfqpoint{1.621573in}{3.042271in}}%
\pgfpathlineto{\pgfqpoint{1.653422in}{3.022078in}}%
\pgfpathlineto{\pgfqpoint{1.603625in}{3.069979in}}%
\pgfpathlineto{\pgfqpoint{1.571786in}{3.090140in}}%
\pgfpathclose%
\pgfusepath{stroke,fill}%
\end{pgfscope}%
\begin{pgfscope}%
\pgfpathrectangle{\pgfqpoint{0.050000in}{0.050000in}}{\pgfqpoint{4.900000in}{4.900000in}}%
\pgfusepath{clip}%
\pgfsetbuttcap%
\pgfsetroundjoin%
\definecolor{currentfill}{rgb}{0.436263,0.436263,0.436263}%
\pgfsetfillcolor{currentfill}%
\pgfsetfillopacity{0.900000}%
\pgfsetlinewidth{0.501875pt}%
\definecolor{currentstroke}{rgb}{0.200000,0.200000,0.200000}%
\pgfsetstrokecolor{currentstroke}%
\pgfsetdash{}{0pt}%
\pgfpathmoveto{\pgfqpoint{2.320677in}{2.484901in}}%
\pgfpathlineto{\pgfqpoint{2.368812in}{2.438368in}}%
\pgfpathlineto{\pgfqpoint{2.401872in}{2.417199in}}%
\pgfpathlineto{\pgfqpoint{2.353732in}{2.463761in}}%
\pgfpathlineto{\pgfqpoint{2.320677in}{2.484901in}}%
\pgfpathclose%
\pgfusepath{stroke,fill}%
\end{pgfscope}%
\begin{pgfscope}%
\pgfpathrectangle{\pgfqpoint{0.050000in}{0.050000in}}{\pgfqpoint{4.900000in}{4.900000in}}%
\pgfusepath{clip}%
\pgfsetbuttcap%
\pgfsetroundjoin%
\definecolor{currentfill}{rgb}{0.705790,0.705790,0.705790}%
\pgfsetfillcolor{currentfill}%
\pgfsetfillopacity{0.900000}%
\pgfsetlinewidth{0.501875pt}%
\definecolor{currentstroke}{rgb}{0.200000,0.200000,0.200000}%
\pgfsetstrokecolor{currentstroke}%
\pgfsetdash{}{0pt}%
\pgfpathmoveto{\pgfqpoint{3.069684in}{1.880597in}}%
\pgfpathlineto{\pgfqpoint{3.116174in}{1.836130in}}%
\pgfpathlineto{\pgfqpoint{3.150447in}{1.814129in}}%
\pgfpathlineto{\pgfqpoint{3.103956in}{1.858569in}}%
\pgfpathlineto{\pgfqpoint{3.069684in}{1.880597in}}%
\pgfpathclose%
\pgfusepath{stroke,fill}%
\end{pgfscope}%
\begin{pgfscope}%
\pgfpathrectangle{\pgfqpoint{0.050000in}{0.050000in}}{\pgfqpoint{4.900000in}{4.900000in}}%
\pgfusepath{clip}%
\pgfsetbuttcap%
\pgfsetroundjoin%
\definecolor{currentfill}{rgb}{0.040969,0.040969,0.040969}%
\pgfsetfillcolor{currentfill}%
\pgfsetfillopacity{0.900000}%
\pgfsetlinewidth{0.501875pt}%
\definecolor{currentstroke}{rgb}{0.200000,0.200000,0.200000}%
\pgfsetstrokecolor{currentstroke}%
\pgfsetdash{}{0pt}%
\pgfpathmoveto{\pgfqpoint{1.294896in}{3.314617in}}%
\pgfpathlineto{\pgfqpoint{1.345304in}{3.266243in}}%
\pgfpathlineto{\pgfqpoint{1.376711in}{3.246408in}}%
\pgfpathlineto{\pgfqpoint{1.326291in}{3.294816in}}%
\pgfpathlineto{\pgfqpoint{1.294896in}{3.314617in}}%
\pgfpathclose%
\pgfusepath{stroke,fill}%
\end{pgfscope}%
\begin{pgfscope}%
\pgfpathrectangle{\pgfqpoint{0.050000in}{0.050000in}}{\pgfqpoint{4.900000in}{4.900000in}}%
\pgfusepath{clip}%
\pgfsetbuttcap%
\pgfsetroundjoin%
\definecolor{currentfill}{rgb}{0.342884,0.342884,0.342884}%
\pgfsetfillcolor{currentfill}%
\pgfsetfillopacity{0.900000}%
\pgfsetlinewidth{0.501875pt}%
\definecolor{currentstroke}{rgb}{0.200000,0.200000,0.200000}%
\pgfsetstrokecolor{currentstroke}%
\pgfsetdash{}{0pt}%
\pgfpathmoveto{\pgfqpoint{2.043900in}{2.709280in}}%
\pgfpathlineto{\pgfqpoint{2.092655in}{2.662242in}}%
\pgfpathlineto{\pgfqpoint{2.125274in}{2.641430in}}%
\pgfpathlineto{\pgfqpoint{2.076512in}{2.688499in}}%
\pgfpathlineto{\pgfqpoint{2.043900in}{2.709280in}}%
\pgfpathclose%
\pgfusepath{stroke,fill}%
\end{pgfscope}%
\begin{pgfscope}%
\pgfpathrectangle{\pgfqpoint{0.050000in}{0.050000in}}{\pgfqpoint{4.900000in}{4.900000in}}%
\pgfusepath{clip}%
\pgfsetbuttcap%
\pgfsetroundjoin%
\definecolor{currentfill}{rgb}{0.600231,0.600231,0.600231}%
\pgfsetfillcolor{currentfill}%
\pgfsetfillopacity{0.900000}%
\pgfsetlinewidth{0.501875pt}%
\definecolor{currentstroke}{rgb}{0.200000,0.200000,0.200000}%
\pgfsetstrokecolor{currentstroke}%
\pgfsetdash{}{0pt}%
\pgfpathmoveto{\pgfqpoint{2.793010in}{2.103966in}}%
\pgfpathlineto{\pgfqpoint{2.840112in}{2.058357in}}%
\pgfpathlineto{\pgfqpoint{2.873942in}{2.036593in}}%
\pgfpathlineto{\pgfqpoint{2.826838in}{2.082216in}}%
\pgfpathlineto{\pgfqpoint{2.793010in}{2.103966in}}%
\pgfpathclose%
\pgfusepath{stroke,fill}%
\end{pgfscope}%
\begin{pgfscope}%
\pgfpathrectangle{\pgfqpoint{0.050000in}{0.050000in}}{\pgfqpoint{4.900000in}{4.900000in}}%
\pgfusepath{clip}%
\pgfsetbuttcap%
\pgfsetroundjoin%
\definecolor{currentfill}{rgb}{0.228835,0.228835,0.228835}%
\pgfsetfillcolor{currentfill}%
\pgfsetfillopacity{0.900000}%
\pgfsetlinewidth{0.501875pt}%
\definecolor{currentstroke}{rgb}{0.200000,0.200000,0.200000}%
\pgfsetstrokecolor{currentstroke}%
\pgfsetdash{}{0pt}%
\pgfpathmoveto{\pgfqpoint{1.767025in}{2.933741in}}%
\pgfpathlineto{\pgfqpoint{1.816401in}{2.886198in}}%
\pgfpathlineto{\pgfqpoint{1.848578in}{2.865744in}}%
\pgfpathlineto{\pgfqpoint{1.799194in}{2.913319in}}%
\pgfpathlineto{\pgfqpoint{1.767025in}{2.933741in}}%
\pgfpathclose%
\pgfusepath{stroke,fill}%
\end{pgfscope}%
\begin{pgfscope}%
\pgfpathrectangle{\pgfqpoint{0.050000in}{0.050000in}}{\pgfqpoint{4.900000in}{4.900000in}}%
\pgfusepath{clip}%
\pgfsetbuttcap%
\pgfsetroundjoin%
\definecolor{currentfill}{rgb}{0.499962,0.499962,0.499962}%
\pgfsetfillcolor{currentfill}%
\pgfsetfillopacity{0.900000}%
\pgfsetlinewidth{0.501875pt}%
\definecolor{currentstroke}{rgb}{0.200000,0.200000,0.200000}%
\pgfsetstrokecolor{currentstroke}%
\pgfsetdash{}{0pt}%
\pgfpathmoveto{\pgfqpoint{2.516246in}{2.328244in}}%
\pgfpathlineto{\pgfqpoint{2.563969in}{2.282042in}}%
\pgfpathlineto{\pgfqpoint{2.597358in}{2.260613in}}%
\pgfpathlineto{\pgfqpoint{2.549632in}{2.306842in}}%
\pgfpathlineto{\pgfqpoint{2.516246in}{2.328244in}}%
\pgfpathclose%
\pgfusepath{stroke,fill}%
\end{pgfscope}%
\begin{pgfscope}%
\pgfpathrectangle{\pgfqpoint{0.050000in}{0.050000in}}{\pgfqpoint{4.900000in}{4.900000in}}%
\pgfusepath{clip}%
\pgfsetbuttcap%
\pgfsetroundjoin%
\definecolor{currentfill}{rgb}{0.109250,0.109250,0.109250}%
\pgfsetfillcolor{currentfill}%
\pgfsetfillopacity{0.900000}%
\pgfsetlinewidth{0.501875pt}%
\definecolor{currentstroke}{rgb}{0.200000,0.200000,0.200000}%
\pgfsetstrokecolor{currentstroke}%
\pgfsetdash{}{0pt}%
\pgfpathmoveto{\pgfqpoint{1.490052in}{3.158283in}}%
\pgfpathlineto{\pgfqpoint{1.540050in}{3.110235in}}%
\pgfpathlineto{\pgfqpoint{1.571786in}{3.090140in}}%
\pgfpathlineto{\pgfqpoint{1.521778in}{3.138221in}}%
\pgfpathlineto{\pgfqpoint{1.490052in}{3.158283in}}%
\pgfpathclose%
\pgfusepath{stroke,fill}%
\end{pgfscope}%
\begin{pgfscope}%
\pgfpathrectangle{\pgfqpoint{0.050000in}{0.050000in}}{\pgfqpoint{4.900000in}{4.900000in}}%
\pgfusepath{clip}%
\pgfsetbuttcap%
\pgfsetroundjoin%
\definecolor{currentfill}{rgb}{0.407843,0.407843,0.407843}%
\pgfsetfillcolor{currentfill}%
\pgfsetfillopacity{0.900000}%
\pgfsetlinewidth{0.501875pt}%
\definecolor{currentstroke}{rgb}{0.200000,0.200000,0.200000}%
\pgfsetstrokecolor{currentstroke}%
\pgfsetdash{}{0pt}%
\pgfpathmoveto{\pgfqpoint{2.239386in}{2.552686in}}%
\pgfpathlineto{\pgfqpoint{2.287730in}{2.505974in}}%
\pgfpathlineto{\pgfqpoint{2.320677in}{2.484901in}}%
\pgfpathlineto{\pgfqpoint{2.272329in}{2.531642in}}%
\pgfpathlineto{\pgfqpoint{2.239386in}{2.552686in}}%
\pgfpathclose%
\pgfusepath{stroke,fill}%
\end{pgfscope}%
\begin{pgfscope}%
\pgfpathrectangle{\pgfqpoint{0.050000in}{0.050000in}}{\pgfqpoint{4.900000in}{4.900000in}}%
\pgfusepath{clip}%
\pgfsetbuttcap%
\pgfsetroundjoin%
\definecolor{currentfill}{rgb}{0.672203,0.672203,0.672203}%
\pgfsetfillcolor{currentfill}%
\pgfsetfillopacity{0.900000}%
\pgfsetlinewidth{0.501875pt}%
\definecolor{currentstroke}{rgb}{0.200000,0.200000,0.200000}%
\pgfsetstrokecolor{currentstroke}%
\pgfsetdash{}{0pt}%
\pgfpathmoveto{\pgfqpoint{2.988830in}{1.947522in}}%
\pgfpathlineto{\pgfqpoint{3.035522in}{1.902558in}}%
\pgfpathlineto{\pgfqpoint{3.069684in}{1.880597in}}%
\pgfpathlineto{\pgfqpoint{3.022990in}{1.925551in}}%
\pgfpathlineto{\pgfqpoint{2.988830in}{1.947522in}}%
\pgfpathclose%
\pgfusepath{stroke,fill}%
\end{pgfscope}%
\begin{pgfscope}%
\pgfpathrectangle{\pgfqpoint{0.050000in}{0.050000in}}{\pgfqpoint{4.900000in}{4.900000in}}%
\pgfusepath{clip}%
\pgfsetbuttcap%
\pgfsetroundjoin%
\definecolor{currentfill}{rgb}{0.004552,0.004552,0.004552}%
\pgfsetfillcolor{currentfill}%
\pgfsetfillopacity{0.900000}%
\pgfsetlinewidth{0.501875pt}%
\definecolor{currentstroke}{rgb}{0.200000,0.200000,0.200000}%
\pgfsetstrokecolor{currentstroke}%
\pgfsetdash{}{0pt}%
\pgfpathmoveto{\pgfqpoint{1.212982in}{3.382907in}}%
\pgfpathlineto{\pgfqpoint{1.263601in}{3.334354in}}%
\pgfpathlineto{\pgfqpoint{1.294896in}{3.314617in}}%
\pgfpathlineto{\pgfqpoint{1.244264in}{3.363205in}}%
\pgfpathlineto{\pgfqpoint{1.212982in}{3.382907in}}%
\pgfpathclose%
\pgfusepath{stroke,fill}%
\end{pgfscope}%
\begin{pgfscope}%
\pgfpathrectangle{\pgfqpoint{0.050000in}{0.050000in}}{\pgfqpoint{4.900000in}{4.900000in}}%
\pgfusepath{clip}%
\pgfsetbuttcap%
\pgfsetroundjoin%
\definecolor{currentfill}{rgb}{0.311880,0.311880,0.311880}%
\pgfsetfillcolor{currentfill}%
\pgfsetfillopacity{0.900000}%
\pgfsetlinewidth{0.501875pt}%
\definecolor{currentstroke}{rgb}{0.200000,0.200000,0.200000}%
\pgfsetstrokecolor{currentstroke}%
\pgfsetdash{}{0pt}%
\pgfpathmoveto{\pgfqpoint{1.962429in}{2.777212in}}%
\pgfpathlineto{\pgfqpoint{2.011393in}{2.729995in}}%
\pgfpathlineto{\pgfqpoint{2.043900in}{2.709280in}}%
\pgfpathlineto{\pgfqpoint{1.994928in}{2.756528in}}%
\pgfpathlineto{\pgfqpoint{1.962429in}{2.777212in}}%
\pgfpathclose%
\pgfusepath{stroke,fill}%
\end{pgfscope}%
\begin{pgfscope}%
\pgfpathrectangle{\pgfqpoint{0.050000in}{0.050000in}}{\pgfqpoint{4.900000in}{4.900000in}}%
\pgfusepath{clip}%
\pgfsetbuttcap%
\pgfsetroundjoin%
\definecolor{currentfill}{rgb}{0.564552,0.564552,0.564552}%
\pgfsetfillcolor{currentfill}%
\pgfsetfillopacity{0.900000}%
\pgfsetlinewidth{0.501875pt}%
\definecolor{currentstroke}{rgb}{0.200000,0.200000,0.200000}%
\pgfsetstrokecolor{currentstroke}%
\pgfsetdash{}{0pt}%
\pgfpathmoveto{\pgfqpoint{2.711981in}{2.171484in}}%
\pgfpathlineto{\pgfqpoint{2.759291in}{2.125646in}}%
\pgfpathlineto{\pgfqpoint{2.793010in}{2.103966in}}%
\pgfpathlineto{\pgfqpoint{2.745698in}{2.149824in}}%
\pgfpathlineto{\pgfqpoint{2.711981in}{2.171484in}}%
\pgfpathclose%
\pgfusepath{stroke,fill}%
\end{pgfscope}%
\begin{pgfscope}%
\pgfpathrectangle{\pgfqpoint{0.050000in}{0.050000in}}{\pgfqpoint{4.900000in}{4.900000in}}%
\pgfusepath{clip}%
\pgfsetbuttcap%
\pgfsetroundjoin%
\definecolor{currentfill}{rgb}{0.184544,0.184544,0.184544}%
\pgfsetfillcolor{currentfill}%
\pgfsetfillopacity{0.900000}%
\pgfsetlinewidth{0.501875pt}%
\definecolor{currentstroke}{rgb}{0.200000,0.200000,0.200000}%
\pgfsetstrokecolor{currentstroke}%
\pgfsetdash{}{0pt}%
\pgfpathmoveto{\pgfqpoint{1.685374in}{3.001819in}}%
\pgfpathlineto{\pgfqpoint{1.734960in}{2.954097in}}%
\pgfpathlineto{\pgfqpoint{1.767025in}{2.933741in}}%
\pgfpathlineto{\pgfqpoint{1.717429in}{2.981495in}}%
\pgfpathlineto{\pgfqpoint{1.685374in}{3.001819in}}%
\pgfpathclose%
\pgfusepath{stroke,fill}%
\end{pgfscope}%
\begin{pgfscope}%
\pgfpathrectangle{\pgfqpoint{0.050000in}{0.050000in}}{\pgfqpoint{4.900000in}{4.900000in}}%
\pgfusepath{clip}%
\pgfsetbuttcap%
\pgfsetroundjoin%
\definecolor{currentfill}{rgb}{0.469819,0.469819,0.469819}%
\pgfsetfillcolor{currentfill}%
\pgfsetfillopacity{0.900000}%
\pgfsetlinewidth{0.501875pt}%
\definecolor{currentstroke}{rgb}{0.200000,0.200000,0.200000}%
\pgfsetstrokecolor{currentstroke}%
\pgfsetdash{}{0pt}%
\pgfpathmoveto{\pgfqpoint{2.435038in}{2.395961in}}%
\pgfpathlineto{\pgfqpoint{2.482969in}{2.349578in}}%
\pgfpathlineto{\pgfqpoint{2.516246in}{2.328244in}}%
\pgfpathlineto{\pgfqpoint{2.468312in}{2.374655in}}%
\pgfpathlineto{\pgfqpoint{2.435038in}{2.395961in}}%
\pgfpathclose%
\pgfusepath{stroke,fill}%
\end{pgfscope}%
\begin{pgfscope}%
\pgfpathrectangle{\pgfqpoint{0.050000in}{0.050000in}}{\pgfqpoint{4.900000in}{4.900000in}}%
\pgfusepath{clip}%
\pgfsetbuttcap%
\pgfsetroundjoin%
\definecolor{currentfill}{rgb}{0.077386,0.077386,0.077386}%
\pgfsetfillcolor{currentfill}%
\pgfsetfillopacity{0.900000}%
\pgfsetlinewidth{0.501875pt}%
\definecolor{currentstroke}{rgb}{0.200000,0.200000,0.200000}%
\pgfsetstrokecolor{currentstroke}%
\pgfsetdash{}{0pt}%
\pgfpathmoveto{\pgfqpoint{1.408221in}{3.226508in}}%
\pgfpathlineto{\pgfqpoint{1.458429in}{3.178281in}}%
\pgfpathlineto{\pgfqpoint{1.490052in}{3.158283in}}%
\pgfpathlineto{\pgfqpoint{1.439833in}{3.206545in}}%
\pgfpathlineto{\pgfqpoint{1.408221in}{3.226508in}}%
\pgfpathclose%
\pgfusepath{stroke,fill}%
\end{pgfscope}%
\begin{pgfscope}%
\pgfpathrectangle{\pgfqpoint{0.050000in}{0.050000in}}{\pgfqpoint{4.900000in}{4.900000in}}%
\pgfusepath{clip}%
\pgfsetbuttcap%
\pgfsetroundjoin%
\definecolor{currentfill}{rgb}{0.743329,0.743329,0.743329}%
\pgfsetfillcolor{currentfill}%
\pgfsetfillopacity{0.900000}%
\pgfsetlinewidth{0.501875pt}%
\definecolor{currentstroke}{rgb}{0.200000,0.200000,0.200000}%
\pgfsetstrokecolor{currentstroke}%
\pgfsetdash{}{0pt}%
\pgfpathmoveto{\pgfqpoint{3.184830in}{1.792063in}}%
\pgfpathlineto{\pgfqpoint{3.231122in}{1.748315in}}%
\pgfpathlineto{\pgfqpoint{3.265618in}{1.726233in}}%
\pgfpathlineto{\pgfqpoint{3.219326in}{1.769932in}}%
\pgfpathlineto{\pgfqpoint{3.184830in}{1.792063in}}%
\pgfpathclose%
\pgfusepath{stroke,fill}%
\end{pgfscope}%
\begin{pgfscope}%
\pgfpathrectangle{\pgfqpoint{0.050000in}{0.050000in}}{\pgfqpoint{4.900000in}{4.900000in}}%
\pgfusepath{clip}%
\pgfsetbuttcap%
\pgfsetroundjoin%
\definecolor{currentfill}{rgb}{0.375363,0.375363,0.375363}%
\pgfsetfillcolor{currentfill}%
\pgfsetfillopacity{0.900000}%
\pgfsetlinewidth{0.501875pt}%
\definecolor{currentstroke}{rgb}{0.200000,0.200000,0.200000}%
\pgfsetstrokecolor{currentstroke}%
\pgfsetdash{}{0pt}%
\pgfpathmoveto{\pgfqpoint{2.157998in}{2.620552in}}%
\pgfpathlineto{\pgfqpoint{2.206550in}{2.573662in}}%
\pgfpathlineto{\pgfqpoint{2.239386in}{2.552686in}}%
\pgfpathlineto{\pgfqpoint{2.190828in}{2.599605in}}%
\pgfpathlineto{\pgfqpoint{2.157998in}{2.620552in}}%
\pgfpathclose%
\pgfusepath{stroke,fill}%
\end{pgfscope}%
\begin{pgfscope}%
\pgfpathrectangle{\pgfqpoint{0.050000in}{0.050000in}}{\pgfqpoint{4.900000in}{4.900000in}}%
\pgfusepath{clip}%
\pgfsetbuttcap%
\pgfsetroundjoin%
\definecolor{currentfill}{rgb}{0.633818,0.633818,0.633818}%
\pgfsetfillcolor{currentfill}%
\pgfsetfillopacity{0.900000}%
\pgfsetlinewidth{0.501875pt}%
\definecolor{currentstroke}{rgb}{0.200000,0.200000,0.200000}%
\pgfsetstrokecolor{currentstroke}%
\pgfsetdash{}{0pt}%
\pgfpathmoveto{\pgfqpoint{2.907881in}{2.014762in}}%
\pgfpathlineto{\pgfqpoint{2.954779in}{1.969426in}}%
\pgfpathlineto{\pgfqpoint{2.988830in}{1.947522in}}%
\pgfpathlineto{\pgfqpoint{2.941931in}{1.992862in}}%
\pgfpathlineto{\pgfqpoint{2.907881in}{2.014762in}}%
\pgfpathclose%
\pgfusepath{stroke,fill}%
\end{pgfscope}%
\begin{pgfscope}%
\pgfpathrectangle{\pgfqpoint{0.050000in}{0.050000in}}{\pgfqpoint{4.900000in}{4.900000in}}%
\pgfusepath{clip}%
\pgfsetbuttcap%
\pgfsetroundjoin%
\definecolor{currentfill}{rgb}{0.267589,0.267589,0.267589}%
\pgfsetfillcolor{currentfill}%
\pgfsetfillopacity{0.900000}%
\pgfsetlinewidth{0.501875pt}%
\definecolor{currentstroke}{rgb}{0.200000,0.200000,0.200000}%
\pgfsetstrokecolor{currentstroke}%
\pgfsetdash{}{0pt}%
\pgfpathmoveto{\pgfqpoint{1.880860in}{2.845224in}}%
\pgfpathlineto{\pgfqpoint{1.930035in}{2.797829in}}%
\pgfpathlineto{\pgfqpoint{1.962429in}{2.777212in}}%
\pgfpathlineto{\pgfqpoint{1.913246in}{2.824638in}}%
\pgfpathlineto{\pgfqpoint{1.880860in}{2.845224in}}%
\pgfpathclose%
\pgfusepath{stroke,fill}%
\end{pgfscope}%
\begin{pgfscope}%
\pgfpathrectangle{\pgfqpoint{0.050000in}{0.050000in}}{\pgfqpoint{4.900000in}{4.900000in}}%
\pgfusepath{clip}%
\pgfsetbuttcap%
\pgfsetroundjoin%
\definecolor{currentfill}{rgb}{0.534410,0.534410,0.534410}%
\pgfsetfillcolor{currentfill}%
\pgfsetfillopacity{0.900000}%
\pgfsetlinewidth{0.501875pt}%
\definecolor{currentstroke}{rgb}{0.200000,0.200000,0.200000}%
\pgfsetstrokecolor{currentstroke}%
\pgfsetdash{}{0pt}%
\pgfpathmoveto{\pgfqpoint{2.630855in}{2.239115in}}%
\pgfpathlineto{\pgfqpoint{2.678372in}{2.193075in}}%
\pgfpathlineto{\pgfqpoint{2.711981in}{2.171484in}}%
\pgfpathlineto{\pgfqpoint{2.664460in}{2.217548in}}%
\pgfpathlineto{\pgfqpoint{2.630855in}{2.239115in}}%
\pgfpathclose%
\pgfusepath{stroke,fill}%
\end{pgfscope}%
\begin{pgfscope}%
\pgfpathrectangle{\pgfqpoint{0.050000in}{0.050000in}}{\pgfqpoint{4.900000in}{4.900000in}}%
\pgfusepath{clip}%
\pgfsetbuttcap%
\pgfsetroundjoin%
\definecolor{currentfill}{rgb}{0.145790,0.145790,0.145790}%
\pgfsetfillcolor{currentfill}%
\pgfsetfillopacity{0.900000}%
\pgfsetlinewidth{0.501875pt}%
\definecolor{currentstroke}{rgb}{0.200000,0.200000,0.200000}%
\pgfsetstrokecolor{currentstroke}%
\pgfsetdash{}{0pt}%
\pgfpathmoveto{\pgfqpoint{1.603625in}{3.069979in}}%
\pgfpathlineto{\pgfqpoint{1.653422in}{3.022078in}}%
\pgfpathlineto{\pgfqpoint{1.685374in}{3.001819in}}%
\pgfpathlineto{\pgfqpoint{1.635567in}{3.049753in}}%
\pgfpathlineto{\pgfqpoint{1.603625in}{3.069979in}}%
\pgfpathclose%
\pgfusepath{stroke,fill}%
\end{pgfscope}%
\begin{pgfscope}%
\pgfpathrectangle{\pgfqpoint{0.050000in}{0.050000in}}{\pgfqpoint{4.900000in}{4.900000in}}%
\pgfusepath{clip}%
\pgfsetbuttcap%
\pgfsetroundjoin%
\definecolor{currentfill}{rgb}{0.436263,0.436263,0.436263}%
\pgfsetfillcolor{currentfill}%
\pgfsetfillopacity{0.900000}%
\pgfsetlinewidth{0.501875pt}%
\definecolor{currentstroke}{rgb}{0.200000,0.200000,0.200000}%
\pgfsetstrokecolor{currentstroke}%
\pgfsetdash{}{0pt}%
\pgfpathmoveto{\pgfqpoint{2.353732in}{2.463761in}}%
\pgfpathlineto{\pgfqpoint{2.401872in}{2.417199in}}%
\pgfpathlineto{\pgfqpoint{2.435038in}{2.395961in}}%
\pgfpathlineto{\pgfqpoint{2.386894in}{2.442551in}}%
\pgfpathlineto{\pgfqpoint{2.353732in}{2.463761in}}%
\pgfpathclose%
\pgfusepath{stroke,fill}%
\end{pgfscope}%
\begin{pgfscope}%
\pgfpathrectangle{\pgfqpoint{0.050000in}{0.050000in}}{\pgfqpoint{4.900000in}{4.900000in}}%
\pgfusepath{clip}%
\pgfsetbuttcap%
\pgfsetroundjoin%
\definecolor{currentfill}{rgb}{0.705790,0.705790,0.705790}%
\pgfsetfillcolor{currentfill}%
\pgfsetfillopacity{0.900000}%
\pgfsetlinewidth{0.501875pt}%
\definecolor{currentstroke}{rgb}{0.200000,0.200000,0.200000}%
\pgfsetstrokecolor{currentstroke}%
\pgfsetdash{}{0pt}%
\pgfpathmoveto{\pgfqpoint{3.103956in}{1.858569in}}%
\pgfpathlineto{\pgfqpoint{3.150447in}{1.814129in}}%
\pgfpathlineto{\pgfqpoint{3.184830in}{1.792063in}}%
\pgfpathlineto{\pgfqpoint{3.138339in}{1.836474in}}%
\pgfpathlineto{\pgfqpoint{3.103956in}{1.858569in}}%
\pgfpathclose%
\pgfusepath{stroke,fill}%
\end{pgfscope}%
\begin{pgfscope}%
\pgfpathrectangle{\pgfqpoint{0.050000in}{0.050000in}}{\pgfqpoint{4.900000in}{4.900000in}}%
\pgfusepath{clip}%
\pgfsetbuttcap%
\pgfsetroundjoin%
\definecolor{currentfill}{rgb}{0.040969,0.040969,0.040969}%
\pgfsetfillcolor{currentfill}%
\pgfsetfillopacity{0.900000}%
\pgfsetlinewidth{0.501875pt}%
\definecolor{currentstroke}{rgb}{0.200000,0.200000,0.200000}%
\pgfsetstrokecolor{currentstroke}%
\pgfsetdash{}{0pt}%
\pgfpathmoveto{\pgfqpoint{1.326291in}{3.294816in}}%
\pgfpathlineto{\pgfqpoint{1.376711in}{3.246408in}}%
\pgfpathlineto{\pgfqpoint{1.408221in}{3.226508in}}%
\pgfpathlineto{\pgfqpoint{1.357789in}{3.274950in}}%
\pgfpathlineto{\pgfqpoint{1.326291in}{3.294816in}}%
\pgfpathclose%
\pgfusepath{stroke,fill}%
\end{pgfscope}%
\begin{pgfscope}%
\pgfpathrectangle{\pgfqpoint{0.050000in}{0.050000in}}{\pgfqpoint{4.900000in}{4.900000in}}%
\pgfusepath{clip}%
\pgfsetbuttcap%
\pgfsetroundjoin%
\definecolor{currentfill}{rgb}{0.342884,0.342884,0.342884}%
\pgfsetfillcolor{currentfill}%
\pgfsetfillopacity{0.900000}%
\pgfsetlinewidth{0.501875pt}%
\definecolor{currentstroke}{rgb}{0.200000,0.200000,0.200000}%
\pgfsetstrokecolor{currentstroke}%
\pgfsetdash{}{0pt}%
\pgfpathmoveto{\pgfqpoint{2.076512in}{2.688499in}}%
\pgfpathlineto{\pgfqpoint{2.125274in}{2.641430in}}%
\pgfpathlineto{\pgfqpoint{2.157998in}{2.620552in}}%
\pgfpathlineto{\pgfqpoint{2.109229in}{2.667650in}}%
\pgfpathlineto{\pgfqpoint{2.076512in}{2.688499in}}%
\pgfpathclose%
\pgfusepath{stroke,fill}%
\end{pgfscope}%
\begin{pgfscope}%
\pgfpathrectangle{\pgfqpoint{0.050000in}{0.050000in}}{\pgfqpoint{4.900000in}{4.900000in}}%
\pgfusepath{clip}%
\pgfsetbuttcap%
\pgfsetroundjoin%
\definecolor{currentfill}{rgb}{0.600231,0.600231,0.600231}%
\pgfsetfillcolor{currentfill}%
\pgfsetfillopacity{0.900000}%
\pgfsetlinewidth{0.501875pt}%
\definecolor{currentstroke}{rgb}{0.200000,0.200000,0.200000}%
\pgfsetstrokecolor{currentstroke}%
\pgfsetdash{}{0pt}%
\pgfpathmoveto{\pgfqpoint{2.826838in}{2.082216in}}%
\pgfpathlineto{\pgfqpoint{2.873942in}{2.036593in}}%
\pgfpathlineto{\pgfqpoint{2.907881in}{2.014762in}}%
\pgfpathlineto{\pgfqpoint{2.860776in}{2.060398in}}%
\pgfpathlineto{\pgfqpoint{2.826838in}{2.082216in}}%
\pgfpathclose%
\pgfusepath{stroke,fill}%
\end{pgfscope}%
\begin{pgfscope}%
\pgfpathrectangle{\pgfqpoint{0.050000in}{0.050000in}}{\pgfqpoint{4.900000in}{4.900000in}}%
\pgfusepath{clip}%
\pgfsetbuttcap%
\pgfsetroundjoin%
\definecolor{currentfill}{rgb}{0.228835,0.228835,0.228835}%
\pgfsetfillcolor{currentfill}%
\pgfsetfillopacity{0.900000}%
\pgfsetlinewidth{0.501875pt}%
\definecolor{currentstroke}{rgb}{0.200000,0.200000,0.200000}%
\pgfsetstrokecolor{currentstroke}%
\pgfsetdash{}{0pt}%
\pgfpathmoveto{\pgfqpoint{1.799194in}{2.913319in}}%
\pgfpathlineto{\pgfqpoint{1.848578in}{2.865744in}}%
\pgfpathlineto{\pgfqpoint{1.880860in}{2.845224in}}%
\pgfpathlineto{\pgfqpoint{1.831467in}{2.892831in}}%
\pgfpathlineto{\pgfqpoint{1.799194in}{2.913319in}}%
\pgfpathclose%
\pgfusepath{stroke,fill}%
\end{pgfscope}%
\begin{pgfscope}%
\pgfpathrectangle{\pgfqpoint{0.050000in}{0.050000in}}{\pgfqpoint{4.900000in}{4.900000in}}%
\pgfusepath{clip}%
\pgfsetbuttcap%
\pgfsetroundjoin%
\definecolor{currentfill}{rgb}{0.499962,0.499962,0.499962}%
\pgfsetfillcolor{currentfill}%
\pgfsetfillopacity{0.900000}%
\pgfsetlinewidth{0.501875pt}%
\definecolor{currentstroke}{rgb}{0.200000,0.200000,0.200000}%
\pgfsetstrokecolor{currentstroke}%
\pgfsetdash{}{0pt}%
\pgfpathmoveto{\pgfqpoint{2.549632in}{2.306842in}}%
\pgfpathlineto{\pgfqpoint{2.597358in}{2.260613in}}%
\pgfpathlineto{\pgfqpoint{2.630855in}{2.239115in}}%
\pgfpathlineto{\pgfqpoint{2.583126in}{2.285370in}}%
\pgfpathlineto{\pgfqpoint{2.549632in}{2.306842in}}%
\pgfpathclose%
\pgfusepath{stroke,fill}%
\end{pgfscope}%
\begin{pgfscope}%
\pgfpathrectangle{\pgfqpoint{0.050000in}{0.050000in}}{\pgfqpoint{4.900000in}{4.900000in}}%
\pgfusepath{clip}%
\pgfsetbuttcap%
\pgfsetroundjoin%
\definecolor{currentfill}{rgb}{0.109250,0.109250,0.109250}%
\pgfsetfillcolor{currentfill}%
\pgfsetfillopacity{0.900000}%
\pgfsetlinewidth{0.501875pt}%
\definecolor{currentstroke}{rgb}{0.200000,0.200000,0.200000}%
\pgfsetstrokecolor{currentstroke}%
\pgfsetdash{}{0pt}%
\pgfpathmoveto{\pgfqpoint{1.521778in}{3.138221in}}%
\pgfpathlineto{\pgfqpoint{1.571786in}{3.090140in}}%
\pgfpathlineto{\pgfqpoint{1.603625in}{3.069979in}}%
\pgfpathlineto{\pgfqpoint{1.553606in}{3.118093in}}%
\pgfpathlineto{\pgfqpoint{1.521778in}{3.138221in}}%
\pgfpathclose%
\pgfusepath{stroke,fill}%
\end{pgfscope}%
\begin{pgfscope}%
\pgfpathrectangle{\pgfqpoint{0.050000in}{0.050000in}}{\pgfqpoint{4.900000in}{4.900000in}}%
\pgfusepath{clip}%
\pgfsetbuttcap%
\pgfsetroundjoin%
\definecolor{currentfill}{rgb}{0.407843,0.407843,0.407843}%
\pgfsetfillcolor{currentfill}%
\pgfsetfillopacity{0.900000}%
\pgfsetlinewidth{0.501875pt}%
\definecolor{currentstroke}{rgb}{0.200000,0.200000,0.200000}%
\pgfsetstrokecolor{currentstroke}%
\pgfsetdash{}{0pt}%
\pgfpathmoveto{\pgfqpoint{2.272329in}{2.531642in}}%
\pgfpathlineto{\pgfqpoint{2.320677in}{2.484901in}}%
\pgfpathlineto{\pgfqpoint{2.353732in}{2.463761in}}%
\pgfpathlineto{\pgfqpoint{2.305378in}{2.510530in}}%
\pgfpathlineto{\pgfqpoint{2.272329in}{2.531642in}}%
\pgfpathclose%
\pgfusepath{stroke,fill}%
\end{pgfscope}%
\begin{pgfscope}%
\pgfpathrectangle{\pgfqpoint{0.050000in}{0.050000in}}{\pgfqpoint{4.900000in}{4.900000in}}%
\pgfusepath{clip}%
\pgfsetbuttcap%
\pgfsetroundjoin%
\definecolor{currentfill}{rgb}{0.672203,0.672203,0.672203}%
\pgfsetfillcolor{currentfill}%
\pgfsetfillopacity{0.900000}%
\pgfsetlinewidth{0.501875pt}%
\definecolor{currentstroke}{rgb}{0.200000,0.200000,0.200000}%
\pgfsetstrokecolor{currentstroke}%
\pgfsetdash{}{0pt}%
\pgfpathmoveto{\pgfqpoint{3.022990in}{1.925551in}}%
\pgfpathlineto{\pgfqpoint{3.069684in}{1.880597in}}%
\pgfpathlineto{\pgfqpoint{3.103956in}{1.858569in}}%
\pgfpathlineto{\pgfqpoint{3.057262in}{1.903512in}}%
\pgfpathlineto{\pgfqpoint{3.022990in}{1.925551in}}%
\pgfpathclose%
\pgfusepath{stroke,fill}%
\end{pgfscope}%
\begin{pgfscope}%
\pgfpathrectangle{\pgfqpoint{0.050000in}{0.050000in}}{\pgfqpoint{4.900000in}{4.900000in}}%
\pgfusepath{clip}%
\pgfsetbuttcap%
\pgfsetroundjoin%
\definecolor{currentfill}{rgb}{0.004552,0.004552,0.004552}%
\pgfsetfillcolor{currentfill}%
\pgfsetfillopacity{0.900000}%
\pgfsetlinewidth{0.501875pt}%
\definecolor{currentstroke}{rgb}{0.200000,0.200000,0.200000}%
\pgfsetstrokecolor{currentstroke}%
\pgfsetdash{}{0pt}%
\pgfpathmoveto{\pgfqpoint{1.244264in}{3.363205in}}%
\pgfpathlineto{\pgfqpoint{1.294896in}{3.314617in}}%
\pgfpathlineto{\pgfqpoint{1.326291in}{3.294816in}}%
\pgfpathlineto{\pgfqpoint{1.275647in}{3.343438in}}%
\pgfpathlineto{\pgfqpoint{1.244264in}{3.363205in}}%
\pgfpathclose%
\pgfusepath{stroke,fill}%
\end{pgfscope}%
\begin{pgfscope}%
\pgfpathrectangle{\pgfqpoint{0.050000in}{0.050000in}}{\pgfqpoint{4.900000in}{4.900000in}}%
\pgfusepath{clip}%
\pgfsetbuttcap%
\pgfsetroundjoin%
\definecolor{currentfill}{rgb}{0.311880,0.311880,0.311880}%
\pgfsetfillcolor{currentfill}%
\pgfsetfillopacity{0.900000}%
\pgfsetlinewidth{0.501875pt}%
\definecolor{currentstroke}{rgb}{0.200000,0.200000,0.200000}%
\pgfsetstrokecolor{currentstroke}%
\pgfsetdash{}{0pt}%
\pgfpathmoveto{\pgfqpoint{1.994928in}{2.756528in}}%
\pgfpathlineto{\pgfqpoint{2.043900in}{2.709280in}}%
\pgfpathlineto{\pgfqpoint{2.076512in}{2.688499in}}%
\pgfpathlineto{\pgfqpoint{2.027533in}{2.735776in}}%
\pgfpathlineto{\pgfqpoint{1.994928in}{2.756528in}}%
\pgfpathclose%
\pgfusepath{stroke,fill}%
\end{pgfscope}%
\begin{pgfscope}%
\pgfpathrectangle{\pgfqpoint{0.050000in}{0.050000in}}{\pgfqpoint{4.900000in}{4.900000in}}%
\pgfusepath{clip}%
\pgfsetbuttcap%
\pgfsetroundjoin%
\definecolor{currentfill}{rgb}{0.564552,0.564552,0.564552}%
\pgfsetfillcolor{currentfill}%
\pgfsetfillopacity{0.900000}%
\pgfsetlinewidth{0.501875pt}%
\definecolor{currentstroke}{rgb}{0.200000,0.200000,0.200000}%
\pgfsetstrokecolor{currentstroke}%
\pgfsetdash{}{0pt}%
\pgfpathmoveto{\pgfqpoint{2.745698in}{2.149824in}}%
\pgfpathlineto{\pgfqpoint{2.793010in}{2.103966in}}%
\pgfpathlineto{\pgfqpoint{2.826838in}{2.082216in}}%
\pgfpathlineto{\pgfqpoint{2.779524in}{2.128094in}}%
\pgfpathlineto{\pgfqpoint{2.745698in}{2.149824in}}%
\pgfpathclose%
\pgfusepath{stroke,fill}%
\end{pgfscope}%
\begin{pgfscope}%
\pgfpathrectangle{\pgfqpoint{0.050000in}{0.050000in}}{\pgfqpoint{4.900000in}{4.900000in}}%
\pgfusepath{clip}%
\pgfsetbuttcap%
\pgfsetroundjoin%
\definecolor{currentfill}{rgb}{0.184544,0.184544,0.184544}%
\pgfsetfillcolor{currentfill}%
\pgfsetfillopacity{0.900000}%
\pgfsetlinewidth{0.501875pt}%
\definecolor{currentstroke}{rgb}{0.200000,0.200000,0.200000}%
\pgfsetstrokecolor{currentstroke}%
\pgfsetdash{}{0pt}%
\pgfpathmoveto{\pgfqpoint{1.717429in}{2.981495in}}%
\pgfpathlineto{\pgfqpoint{1.767025in}{2.933741in}}%
\pgfpathlineto{\pgfqpoint{1.799194in}{2.913319in}}%
\pgfpathlineto{\pgfqpoint{1.749589in}{2.961105in}}%
\pgfpathlineto{\pgfqpoint{1.717429in}{2.981495in}}%
\pgfpathclose%
\pgfusepath{stroke,fill}%
\end{pgfscope}%
\begin{pgfscope}%
\pgfpathrectangle{\pgfqpoint{0.050000in}{0.050000in}}{\pgfqpoint{4.900000in}{4.900000in}}%
\pgfusepath{clip}%
\pgfsetbuttcap%
\pgfsetroundjoin%
\definecolor{currentfill}{rgb}{0.469819,0.469819,0.469819}%
\pgfsetfillcolor{currentfill}%
\pgfsetfillopacity{0.900000}%
\pgfsetlinewidth{0.501875pt}%
\definecolor{currentstroke}{rgb}{0.200000,0.200000,0.200000}%
\pgfsetstrokecolor{currentstroke}%
\pgfsetdash{}{0pt}%
\pgfpathmoveto{\pgfqpoint{2.468312in}{2.374655in}}%
\pgfpathlineto{\pgfqpoint{2.516246in}{2.328244in}}%
\pgfpathlineto{\pgfqpoint{2.549632in}{2.306842in}}%
\pgfpathlineto{\pgfqpoint{2.501694in}{2.353279in}}%
\pgfpathlineto{\pgfqpoint{2.468312in}{2.374655in}}%
\pgfpathclose%
\pgfusepath{stroke,fill}%
\end{pgfscope}%
\begin{pgfscope}%
\pgfpathrectangle{\pgfqpoint{0.050000in}{0.050000in}}{\pgfqpoint{4.900000in}{4.900000in}}%
\pgfusepath{clip}%
\pgfsetbuttcap%
\pgfsetroundjoin%
\definecolor{currentfill}{rgb}{0.077386,0.077386,0.077386}%
\pgfsetfillcolor{currentfill}%
\pgfsetfillopacity{0.900000}%
\pgfsetlinewidth{0.501875pt}%
\definecolor{currentstroke}{rgb}{0.200000,0.200000,0.200000}%
\pgfsetstrokecolor{currentstroke}%
\pgfsetdash{}{0pt}%
\pgfpathmoveto{\pgfqpoint{1.439833in}{3.206545in}}%
\pgfpathlineto{\pgfqpoint{1.490052in}{3.158283in}}%
\pgfpathlineto{\pgfqpoint{1.521778in}{3.138221in}}%
\pgfpathlineto{\pgfqpoint{1.471547in}{3.186516in}}%
\pgfpathlineto{\pgfqpoint{1.439833in}{3.206545in}}%
\pgfpathclose%
\pgfusepath{stroke,fill}%
\end{pgfscope}%
\begin{pgfscope}%
\pgfpathrectangle{\pgfqpoint{0.050000in}{0.050000in}}{\pgfqpoint{4.900000in}{4.900000in}}%
\pgfusepath{clip}%
\pgfsetbuttcap%
\pgfsetroundjoin%
\definecolor{currentfill}{rgb}{0.375363,0.375363,0.375363}%
\pgfsetfillcolor{currentfill}%
\pgfsetfillopacity{0.900000}%
\pgfsetlinewidth{0.501875pt}%
\definecolor{currentstroke}{rgb}{0.200000,0.200000,0.200000}%
\pgfsetstrokecolor{currentstroke}%
\pgfsetdash{}{0pt}%
\pgfpathmoveto{\pgfqpoint{2.190828in}{2.599605in}}%
\pgfpathlineto{\pgfqpoint{2.239386in}{2.552686in}}%
\pgfpathlineto{\pgfqpoint{2.272329in}{2.531642in}}%
\pgfpathlineto{\pgfqpoint{2.223765in}{2.578591in}}%
\pgfpathlineto{\pgfqpoint{2.190828in}{2.599605in}}%
\pgfpathclose%
\pgfusepath{stroke,fill}%
\end{pgfscope}%
\begin{pgfscope}%
\pgfpathrectangle{\pgfqpoint{0.050000in}{0.050000in}}{\pgfqpoint{4.900000in}{4.900000in}}%
\pgfusepath{clip}%
\pgfsetbuttcap%
\pgfsetroundjoin%
\definecolor{currentfill}{rgb}{0.633818,0.633818,0.633818}%
\pgfsetfillcolor{currentfill}%
\pgfsetfillopacity{0.900000}%
\pgfsetlinewidth{0.501875pt}%
\definecolor{currentstroke}{rgb}{0.200000,0.200000,0.200000}%
\pgfsetstrokecolor{currentstroke}%
\pgfsetdash{}{0pt}%
\pgfpathmoveto{\pgfqpoint{2.941931in}{1.992862in}}%
\pgfpathlineto{\pgfqpoint{2.988830in}{1.947522in}}%
\pgfpathlineto{\pgfqpoint{3.022990in}{1.925551in}}%
\pgfpathlineto{\pgfqpoint{2.976091in}{1.970893in}}%
\pgfpathlineto{\pgfqpoint{2.941931in}{1.992862in}}%
\pgfpathclose%
\pgfusepath{stroke,fill}%
\end{pgfscope}%
\begin{pgfscope}%
\pgfpathrectangle{\pgfqpoint{0.050000in}{0.050000in}}{\pgfqpoint{4.900000in}{4.900000in}}%
\pgfusepath{clip}%
\pgfsetbuttcap%
\pgfsetroundjoin%
\definecolor{currentfill}{rgb}{0.267589,0.267589,0.267589}%
\pgfsetfillcolor{currentfill}%
\pgfsetfillopacity{0.900000}%
\pgfsetlinewidth{0.501875pt}%
\definecolor{currentstroke}{rgb}{0.200000,0.200000,0.200000}%
\pgfsetstrokecolor{currentstroke}%
\pgfsetdash{}{0pt}%
\pgfpathmoveto{\pgfqpoint{1.913246in}{2.824638in}}%
\pgfpathlineto{\pgfqpoint{1.962429in}{2.777212in}}%
\pgfpathlineto{\pgfqpoint{1.994928in}{2.756528in}}%
\pgfpathlineto{\pgfqpoint{1.945738in}{2.803985in}}%
\pgfpathlineto{\pgfqpoint{1.913246in}{2.824638in}}%
\pgfpathclose%
\pgfusepath{stroke,fill}%
\end{pgfscope}%
\begin{pgfscope}%
\pgfpathrectangle{\pgfqpoint{0.050000in}{0.050000in}}{\pgfqpoint{4.900000in}{4.900000in}}%
\pgfusepath{clip}%
\pgfsetbuttcap%
\pgfsetroundjoin%
\definecolor{currentfill}{rgb}{0.534410,0.534410,0.534410}%
\pgfsetfillcolor{currentfill}%
\pgfsetfillopacity{0.900000}%
\pgfsetlinewidth{0.501875pt}%
\definecolor{currentstroke}{rgb}{0.200000,0.200000,0.200000}%
\pgfsetstrokecolor{currentstroke}%
\pgfsetdash{}{0pt}%
\pgfpathmoveto{\pgfqpoint{2.664460in}{2.217548in}}%
\pgfpathlineto{\pgfqpoint{2.711981in}{2.171484in}}%
\pgfpathlineto{\pgfqpoint{2.745698in}{2.149824in}}%
\pgfpathlineto{\pgfqpoint{2.698175in}{2.195912in}}%
\pgfpathlineto{\pgfqpoint{2.664460in}{2.217548in}}%
\pgfpathclose%
\pgfusepath{stroke,fill}%
\end{pgfscope}%
\begin{pgfscope}%
\pgfpathrectangle{\pgfqpoint{0.050000in}{0.050000in}}{\pgfqpoint{4.900000in}{4.900000in}}%
\pgfusepath{clip}%
\pgfsetbuttcap%
\pgfsetroundjoin%
\definecolor{currentfill}{rgb}{0.145790,0.145790,0.145790}%
\pgfsetfillcolor{currentfill}%
\pgfsetfillopacity{0.900000}%
\pgfsetlinewidth{0.501875pt}%
\definecolor{currentstroke}{rgb}{0.200000,0.200000,0.200000}%
\pgfsetstrokecolor{currentstroke}%
\pgfsetdash{}{0pt}%
\pgfpathmoveto{\pgfqpoint{1.635567in}{3.049753in}}%
\pgfpathlineto{\pgfqpoint{1.685374in}{3.001819in}}%
\pgfpathlineto{\pgfqpoint{1.717429in}{2.981495in}}%
\pgfpathlineto{\pgfqpoint{1.667613in}{3.029462in}}%
\pgfpathlineto{\pgfqpoint{1.635567in}{3.049753in}}%
\pgfpathclose%
\pgfusepath{stroke,fill}%
\end{pgfscope}%
\begin{pgfscope}%
\pgfpathrectangle{\pgfqpoint{0.050000in}{0.050000in}}{\pgfqpoint{4.900000in}{4.900000in}}%
\pgfusepath{clip}%
\pgfsetbuttcap%
\pgfsetroundjoin%
\definecolor{currentfill}{rgb}{0.436263,0.436263,0.436263}%
\pgfsetfillcolor{currentfill}%
\pgfsetfillopacity{0.900000}%
\pgfsetlinewidth{0.501875pt}%
\definecolor{currentstroke}{rgb}{0.200000,0.200000,0.200000}%
\pgfsetstrokecolor{currentstroke}%
\pgfsetdash{}{0pt}%
\pgfpathmoveto{\pgfqpoint{2.386894in}{2.442551in}}%
\pgfpathlineto{\pgfqpoint{2.435038in}{2.395961in}}%
\pgfpathlineto{\pgfqpoint{2.468312in}{2.374655in}}%
\pgfpathlineto{\pgfqpoint{2.420163in}{2.421273in}}%
\pgfpathlineto{\pgfqpoint{2.386894in}{2.442551in}}%
\pgfpathclose%
\pgfusepath{stroke,fill}%
\end{pgfscope}%
\begin{pgfscope}%
\pgfpathrectangle{\pgfqpoint{0.050000in}{0.050000in}}{\pgfqpoint{4.900000in}{4.900000in}}%
\pgfusepath{clip}%
\pgfsetbuttcap%
\pgfsetroundjoin%
\definecolor{currentfill}{rgb}{0.705790,0.705790,0.705790}%
\pgfsetfillcolor{currentfill}%
\pgfsetfillopacity{0.900000}%
\pgfsetlinewidth{0.501875pt}%
\definecolor{currentstroke}{rgb}{0.200000,0.200000,0.200000}%
\pgfsetstrokecolor{currentstroke}%
\pgfsetdash{}{0pt}%
\pgfpathmoveto{\pgfqpoint{3.138339in}{1.836474in}}%
\pgfpathlineto{\pgfqpoint{3.184830in}{1.792063in}}%
\pgfpathlineto{\pgfqpoint{3.219326in}{1.769932in}}%
\pgfpathlineto{\pgfqpoint{3.172835in}{1.814313in}}%
\pgfpathlineto{\pgfqpoint{3.138339in}{1.836474in}}%
\pgfpathclose%
\pgfusepath{stroke,fill}%
\end{pgfscope}%
\begin{pgfscope}%
\pgfpathrectangle{\pgfqpoint{0.050000in}{0.050000in}}{\pgfqpoint{4.900000in}{4.900000in}}%
\pgfusepath{clip}%
\pgfsetbuttcap%
\pgfsetroundjoin%
\definecolor{currentfill}{rgb}{0.040969,0.040969,0.040969}%
\pgfsetfillcolor{currentfill}%
\pgfsetfillopacity{0.900000}%
\pgfsetlinewidth{0.501875pt}%
\definecolor{currentstroke}{rgb}{0.200000,0.200000,0.200000}%
\pgfsetstrokecolor{currentstroke}%
\pgfsetdash{}{0pt}%
\pgfpathmoveto{\pgfqpoint{1.357789in}{3.274950in}}%
\pgfpathlineto{\pgfqpoint{1.408221in}{3.226508in}}%
\pgfpathlineto{\pgfqpoint{1.439833in}{3.206545in}}%
\pgfpathlineto{\pgfqpoint{1.389389in}{3.255021in}}%
\pgfpathlineto{\pgfqpoint{1.357789in}{3.274950in}}%
\pgfpathclose%
\pgfusepath{stroke,fill}%
\end{pgfscope}%
\begin{pgfscope}%
\pgfpathrectangle{\pgfqpoint{0.050000in}{0.050000in}}{\pgfqpoint{4.900000in}{4.900000in}}%
\pgfusepath{clip}%
\pgfsetbuttcap%
\pgfsetroundjoin%
\definecolor{currentfill}{rgb}{0.346943,0.346943,0.346943}%
\pgfsetfillcolor{currentfill}%
\pgfsetfillopacity{0.900000}%
\pgfsetlinewidth{0.501875pt}%
\definecolor{currentstroke}{rgb}{0.200000,0.200000,0.200000}%
\pgfsetstrokecolor{currentstroke}%
\pgfsetdash{}{0pt}%
\pgfpathmoveto{\pgfqpoint{2.109229in}{2.667650in}}%
\pgfpathlineto{\pgfqpoint{2.157998in}{2.620552in}}%
\pgfpathlineto{\pgfqpoint{2.190828in}{2.599605in}}%
\pgfpathlineto{\pgfqpoint{2.142053in}{2.646733in}}%
\pgfpathlineto{\pgfqpoint{2.109229in}{2.667650in}}%
\pgfpathclose%
\pgfusepath{stroke,fill}%
\end{pgfscope}%
\begin{pgfscope}%
\pgfpathrectangle{\pgfqpoint{0.050000in}{0.050000in}}{\pgfqpoint{4.900000in}{4.900000in}}%
\pgfusepath{clip}%
\pgfsetbuttcap%
\pgfsetroundjoin%
\definecolor{currentfill}{rgb}{0.600231,0.600231,0.600231}%
\pgfsetfillcolor{currentfill}%
\pgfsetfillopacity{0.900000}%
\pgfsetlinewidth{0.501875pt}%
\definecolor{currentstroke}{rgb}{0.200000,0.200000,0.200000}%
\pgfsetstrokecolor{currentstroke}%
\pgfsetdash{}{0pt}%
\pgfpathmoveto{\pgfqpoint{2.860776in}{2.060398in}}%
\pgfpathlineto{\pgfqpoint{2.907881in}{2.014762in}}%
\pgfpathlineto{\pgfqpoint{2.941931in}{1.992862in}}%
\pgfpathlineto{\pgfqpoint{2.894825in}{2.038510in}}%
\pgfpathlineto{\pgfqpoint{2.860776in}{2.060398in}}%
\pgfpathclose%
\pgfusepath{stroke,fill}%
\end{pgfscope}%
\begin{pgfscope}%
\pgfpathrectangle{\pgfqpoint{0.050000in}{0.050000in}}{\pgfqpoint{4.900000in}{4.900000in}}%
\pgfusepath{clip}%
\pgfsetbuttcap%
\pgfsetroundjoin%
\definecolor{currentfill}{rgb}{0.228835,0.228835,0.228835}%
\pgfsetfillcolor{currentfill}%
\pgfsetfillopacity{0.900000}%
\pgfsetlinewidth{0.501875pt}%
\definecolor{currentstroke}{rgb}{0.200000,0.200000,0.200000}%
\pgfsetstrokecolor{currentstroke}%
\pgfsetdash{}{0pt}%
\pgfpathmoveto{\pgfqpoint{1.831467in}{2.892831in}}%
\pgfpathlineto{\pgfqpoint{1.880860in}{2.845224in}}%
\pgfpathlineto{\pgfqpoint{1.913246in}{2.824638in}}%
\pgfpathlineto{\pgfqpoint{1.863845in}{2.872276in}}%
\pgfpathlineto{\pgfqpoint{1.831467in}{2.892831in}}%
\pgfpathclose%
\pgfusepath{stroke,fill}%
\end{pgfscope}%
\begin{pgfscope}%
\pgfpathrectangle{\pgfqpoint{0.050000in}{0.050000in}}{\pgfqpoint{4.900000in}{4.900000in}}%
\pgfusepath{clip}%
\pgfsetbuttcap%
\pgfsetroundjoin%
\definecolor{currentfill}{rgb}{0.499962,0.499962,0.499962}%
\pgfsetfillcolor{currentfill}%
\pgfsetfillopacity{0.900000}%
\pgfsetlinewidth{0.501875pt}%
\definecolor{currentstroke}{rgb}{0.200000,0.200000,0.200000}%
\pgfsetstrokecolor{currentstroke}%
\pgfsetdash{}{0pt}%
\pgfpathmoveto{\pgfqpoint{2.583126in}{2.285370in}}%
\pgfpathlineto{\pgfqpoint{2.630855in}{2.239115in}}%
\pgfpathlineto{\pgfqpoint{2.664460in}{2.217548in}}%
\pgfpathlineto{\pgfqpoint{2.616729in}{2.263829in}}%
\pgfpathlineto{\pgfqpoint{2.583126in}{2.285370in}}%
\pgfpathclose%
\pgfusepath{stroke,fill}%
\end{pgfscope}%
\begin{pgfscope}%
\pgfpathrectangle{\pgfqpoint{0.050000in}{0.050000in}}{\pgfqpoint{4.900000in}{4.900000in}}%
\pgfusepath{clip}%
\pgfsetbuttcap%
\pgfsetroundjoin%
\definecolor{currentfill}{rgb}{0.109250,0.109250,0.109250}%
\pgfsetfillcolor{currentfill}%
\pgfsetfillopacity{0.900000}%
\pgfsetlinewidth{0.501875pt}%
\definecolor{currentstroke}{rgb}{0.200000,0.200000,0.200000}%
\pgfsetstrokecolor{currentstroke}%
\pgfsetdash{}{0pt}%
\pgfpathmoveto{\pgfqpoint{1.553606in}{3.118093in}}%
\pgfpathlineto{\pgfqpoint{1.603625in}{3.069979in}}%
\pgfpathlineto{\pgfqpoint{1.635567in}{3.049753in}}%
\pgfpathlineto{\pgfqpoint{1.585538in}{3.097901in}}%
\pgfpathlineto{\pgfqpoint{1.553606in}{3.118093in}}%
\pgfpathclose%
\pgfusepath{stroke,fill}%
\end{pgfscope}%
\begin{pgfscope}%
\pgfpathrectangle{\pgfqpoint{0.050000in}{0.050000in}}{\pgfqpoint{4.900000in}{4.900000in}}%
\pgfusepath{clip}%
\pgfsetbuttcap%
\pgfsetroundjoin%
\definecolor{currentfill}{rgb}{0.407843,0.407843,0.407843}%
\pgfsetfillcolor{currentfill}%
\pgfsetfillopacity{0.900000}%
\pgfsetlinewidth{0.501875pt}%
\definecolor{currentstroke}{rgb}{0.200000,0.200000,0.200000}%
\pgfsetstrokecolor{currentstroke}%
\pgfsetdash{}{0pt}%
\pgfpathmoveto{\pgfqpoint{2.305378in}{2.510530in}}%
\pgfpathlineto{\pgfqpoint{2.353732in}{2.463761in}}%
\pgfpathlineto{\pgfqpoint{2.386894in}{2.442551in}}%
\pgfpathlineto{\pgfqpoint{2.338535in}{2.489349in}}%
\pgfpathlineto{\pgfqpoint{2.305378in}{2.510530in}}%
\pgfpathclose%
\pgfusepath{stroke,fill}%
\end{pgfscope}%
\begin{pgfscope}%
\pgfpathrectangle{\pgfqpoint{0.050000in}{0.050000in}}{\pgfqpoint{4.900000in}{4.900000in}}%
\pgfusepath{clip}%
\pgfsetbuttcap%
\pgfsetroundjoin%
\definecolor{currentfill}{rgb}{0.672203,0.672203,0.672203}%
\pgfsetfillcolor{currentfill}%
\pgfsetfillopacity{0.900000}%
\pgfsetlinewidth{0.501875pt}%
\definecolor{currentstroke}{rgb}{0.200000,0.200000,0.200000}%
\pgfsetstrokecolor{currentstroke}%
\pgfsetdash{}{0pt}%
\pgfpathmoveto{\pgfqpoint{3.057262in}{1.903512in}}%
\pgfpathlineto{\pgfqpoint{3.103956in}{1.858569in}}%
\pgfpathlineto{\pgfqpoint{3.138339in}{1.836474in}}%
\pgfpathlineto{\pgfqpoint{3.091646in}{1.881405in}}%
\pgfpathlineto{\pgfqpoint{3.057262in}{1.903512in}}%
\pgfpathclose%
\pgfusepath{stroke,fill}%
\end{pgfscope}%
\begin{pgfscope}%
\pgfpathrectangle{\pgfqpoint{0.050000in}{0.050000in}}{\pgfqpoint{4.900000in}{4.900000in}}%
\pgfusepath{clip}%
\pgfsetbuttcap%
\pgfsetroundjoin%
\definecolor{currentfill}{rgb}{0.009104,0.009104,0.009104}%
\pgfsetfillcolor{currentfill}%
\pgfsetfillopacity{0.900000}%
\pgfsetlinewidth{0.501875pt}%
\definecolor{currentstroke}{rgb}{0.200000,0.200000,0.200000}%
\pgfsetstrokecolor{currentstroke}%
\pgfsetdash{}{0pt}%
\pgfpathmoveto{\pgfqpoint{1.275647in}{3.343438in}}%
\pgfpathlineto{\pgfqpoint{1.326291in}{3.294816in}}%
\pgfpathlineto{\pgfqpoint{1.357789in}{3.274950in}}%
\pgfpathlineto{\pgfqpoint{1.307133in}{3.323608in}}%
\pgfpathlineto{\pgfqpoint{1.275647in}{3.343438in}}%
\pgfpathclose%
\pgfusepath{stroke,fill}%
\end{pgfscope}%
\begin{pgfscope}%
\pgfpathrectangle{\pgfqpoint{0.050000in}{0.050000in}}{\pgfqpoint{4.900000in}{4.900000in}}%
\pgfusepath{clip}%
\pgfsetbuttcap%
\pgfsetroundjoin%
\definecolor{currentfill}{rgb}{0.311880,0.311880,0.311880}%
\pgfsetfillcolor{currentfill}%
\pgfsetfillopacity{0.900000}%
\pgfsetlinewidth{0.501875pt}%
\definecolor{currentstroke}{rgb}{0.200000,0.200000,0.200000}%
\pgfsetstrokecolor{currentstroke}%
\pgfsetdash{}{0pt}%
\pgfpathmoveto{\pgfqpoint{2.027533in}{2.735776in}}%
\pgfpathlineto{\pgfqpoint{2.076512in}{2.688499in}}%
\pgfpathlineto{\pgfqpoint{2.109229in}{2.667650in}}%
\pgfpathlineto{\pgfqpoint{2.060244in}{2.714958in}}%
\pgfpathlineto{\pgfqpoint{2.027533in}{2.735776in}}%
\pgfpathclose%
\pgfusepath{stroke,fill}%
\end{pgfscope}%
\begin{pgfscope}%
\pgfpathrectangle{\pgfqpoint{0.050000in}{0.050000in}}{\pgfqpoint{4.900000in}{4.900000in}}%
\pgfusepath{clip}%
\pgfsetbuttcap%
\pgfsetroundjoin%
\definecolor{currentfill}{rgb}{0.564552,0.564552,0.564552}%
\pgfsetfillcolor{currentfill}%
\pgfsetfillopacity{0.900000}%
\pgfsetlinewidth{0.501875pt}%
\definecolor{currentstroke}{rgb}{0.200000,0.200000,0.200000}%
\pgfsetstrokecolor{currentstroke}%
\pgfsetdash{}{0pt}%
\pgfpathmoveto{\pgfqpoint{2.779524in}{2.128094in}}%
\pgfpathlineto{\pgfqpoint{2.826838in}{2.082216in}}%
\pgfpathlineto{\pgfqpoint{2.860776in}{2.060398in}}%
\pgfpathlineto{\pgfqpoint{2.813461in}{2.106295in}}%
\pgfpathlineto{\pgfqpoint{2.779524in}{2.128094in}}%
\pgfpathclose%
\pgfusepath{stroke,fill}%
\end{pgfscope}%
\begin{pgfscope}%
\pgfpathrectangle{\pgfqpoint{0.050000in}{0.050000in}}{\pgfqpoint{4.900000in}{4.900000in}}%
\pgfusepath{clip}%
\pgfsetbuttcap%
\pgfsetroundjoin%
\definecolor{currentfill}{rgb}{0.184544,0.184544,0.184544}%
\pgfsetfillcolor{currentfill}%
\pgfsetfillopacity{0.900000}%
\pgfsetlinewidth{0.501875pt}%
\definecolor{currentstroke}{rgb}{0.200000,0.200000,0.200000}%
\pgfsetstrokecolor{currentstroke}%
\pgfsetdash{}{0pt}%
\pgfpathmoveto{\pgfqpoint{1.749589in}{2.961105in}}%
\pgfpathlineto{\pgfqpoint{1.799194in}{2.913319in}}%
\pgfpathlineto{\pgfqpoint{1.831467in}{2.892831in}}%
\pgfpathlineto{\pgfqpoint{1.781853in}{2.940649in}}%
\pgfpathlineto{\pgfqpoint{1.749589in}{2.961105in}}%
\pgfpathclose%
\pgfusepath{stroke,fill}%
\end{pgfscope}%
\begin{pgfscope}%
\pgfpathrectangle{\pgfqpoint{0.050000in}{0.050000in}}{\pgfqpoint{4.900000in}{4.900000in}}%
\pgfusepath{clip}%
\pgfsetbuttcap%
\pgfsetroundjoin%
\definecolor{currentfill}{rgb}{0.469819,0.469819,0.469819}%
\pgfsetfillcolor{currentfill}%
\pgfsetfillopacity{0.900000}%
\pgfsetlinewidth{0.501875pt}%
\definecolor{currentstroke}{rgb}{0.200000,0.200000,0.200000}%
\pgfsetstrokecolor{currentstroke}%
\pgfsetdash{}{0pt}%
\pgfpathmoveto{\pgfqpoint{2.501694in}{2.353279in}}%
\pgfpathlineto{\pgfqpoint{2.549632in}{2.306842in}}%
\pgfpathlineto{\pgfqpoint{2.583126in}{2.285370in}}%
\pgfpathlineto{\pgfqpoint{2.535184in}{2.331835in}}%
\pgfpathlineto{\pgfqpoint{2.501694in}{2.353279in}}%
\pgfpathclose%
\pgfusepath{stroke,fill}%
\end{pgfscope}%
\begin{pgfscope}%
\pgfpathrectangle{\pgfqpoint{0.050000in}{0.050000in}}{\pgfqpoint{4.900000in}{4.900000in}}%
\pgfusepath{clip}%
\pgfsetbuttcap%
\pgfsetroundjoin%
\definecolor{currentfill}{rgb}{0.077386,0.077386,0.077386}%
\pgfsetfillcolor{currentfill}%
\pgfsetfillopacity{0.900000}%
\pgfsetlinewidth{0.501875pt}%
\definecolor{currentstroke}{rgb}{0.200000,0.200000,0.200000}%
\pgfsetstrokecolor{currentstroke}%
\pgfsetdash{}{0pt}%
\pgfpathmoveto{\pgfqpoint{1.471547in}{3.186516in}}%
\pgfpathlineto{\pgfqpoint{1.521778in}{3.138221in}}%
\pgfpathlineto{\pgfqpoint{1.553606in}{3.118093in}}%
\pgfpathlineto{\pgfqpoint{1.503365in}{3.166422in}}%
\pgfpathlineto{\pgfqpoint{1.471547in}{3.186516in}}%
\pgfpathclose%
\pgfusepath{stroke,fill}%
\end{pgfscope}%
\begin{pgfscope}%
\pgfpathrectangle{\pgfqpoint{0.050000in}{0.050000in}}{\pgfqpoint{4.900000in}{4.900000in}}%
\pgfusepath{clip}%
\pgfsetbuttcap%
\pgfsetroundjoin%
\definecolor{currentfill}{rgb}{0.375363,0.375363,0.375363}%
\pgfsetfillcolor{currentfill}%
\pgfsetfillopacity{0.900000}%
\pgfsetlinewidth{0.501875pt}%
\definecolor{currentstroke}{rgb}{0.200000,0.200000,0.200000}%
\pgfsetstrokecolor{currentstroke}%
\pgfsetdash{}{0pt}%
\pgfpathmoveto{\pgfqpoint{2.223765in}{2.578591in}}%
\pgfpathlineto{\pgfqpoint{2.272329in}{2.531642in}}%
\pgfpathlineto{\pgfqpoint{2.305378in}{2.510530in}}%
\pgfpathlineto{\pgfqpoint{2.256809in}{2.557508in}}%
\pgfpathlineto{\pgfqpoint{2.223765in}{2.578591in}}%
\pgfpathclose%
\pgfusepath{stroke,fill}%
\end{pgfscope}%
\begin{pgfscope}%
\pgfpathrectangle{\pgfqpoint{0.050000in}{0.050000in}}{\pgfqpoint{4.900000in}{4.900000in}}%
\pgfusepath{clip}%
\pgfsetbuttcap%
\pgfsetroundjoin%
\definecolor{currentfill}{rgb}{0.638616,0.638616,0.638616}%
\pgfsetfillcolor{currentfill}%
\pgfsetfillopacity{0.900000}%
\pgfsetlinewidth{0.501875pt}%
\definecolor{currentstroke}{rgb}{0.200000,0.200000,0.200000}%
\pgfsetstrokecolor{currentstroke}%
\pgfsetdash{}{0pt}%
\pgfpathmoveto{\pgfqpoint{2.976091in}{1.970893in}}%
\pgfpathlineto{\pgfqpoint{3.022990in}{1.925551in}}%
\pgfpathlineto{\pgfqpoint{3.057262in}{1.903512in}}%
\pgfpathlineto{\pgfqpoint{3.010363in}{1.948855in}}%
\pgfpathlineto{\pgfqpoint{2.976091in}{1.970893in}}%
\pgfpathclose%
\pgfusepath{stroke,fill}%
\end{pgfscope}%
\begin{pgfscope}%
\pgfpathrectangle{\pgfqpoint{0.050000in}{0.050000in}}{\pgfqpoint{4.900000in}{4.900000in}}%
\pgfusepath{clip}%
\pgfsetbuttcap%
\pgfsetroundjoin%
\definecolor{currentfill}{rgb}{0.273126,0.273126,0.273126}%
\pgfsetfillcolor{currentfill}%
\pgfsetfillopacity{0.900000}%
\pgfsetlinewidth{0.501875pt}%
\definecolor{currentstroke}{rgb}{0.200000,0.200000,0.200000}%
\pgfsetstrokecolor{currentstroke}%
\pgfsetdash{}{0pt}%
\pgfpathmoveto{\pgfqpoint{1.945738in}{2.803985in}}%
\pgfpathlineto{\pgfqpoint{1.994928in}{2.756528in}}%
\pgfpathlineto{\pgfqpoint{2.027533in}{2.735776in}}%
\pgfpathlineto{\pgfqpoint{1.978335in}{2.783265in}}%
\pgfpathlineto{\pgfqpoint{1.945738in}{2.803985in}}%
\pgfpathclose%
\pgfusepath{stroke,fill}%
\end{pgfscope}%
\begin{pgfscope}%
\pgfpathrectangle{\pgfqpoint{0.050000in}{0.050000in}}{\pgfqpoint{4.900000in}{4.900000in}}%
\pgfusepath{clip}%
\pgfsetbuttcap%
\pgfsetroundjoin%
\definecolor{currentfill}{rgb}{0.534410,0.534410,0.534410}%
\pgfsetfillcolor{currentfill}%
\pgfsetfillopacity{0.900000}%
\pgfsetlinewidth{0.501875pt}%
\definecolor{currentstroke}{rgb}{0.200000,0.200000,0.200000}%
\pgfsetstrokecolor{currentstroke}%
\pgfsetdash{}{0pt}%
\pgfpathmoveto{\pgfqpoint{2.698175in}{2.195912in}}%
\pgfpathlineto{\pgfqpoint{2.745698in}{2.149824in}}%
\pgfpathlineto{\pgfqpoint{2.779524in}{2.128094in}}%
\pgfpathlineto{\pgfqpoint{2.732000in}{2.174205in}}%
\pgfpathlineto{\pgfqpoint{2.698175in}{2.195912in}}%
\pgfpathclose%
\pgfusepath{stroke,fill}%
\end{pgfscope}%
\begin{pgfscope}%
\pgfpathrectangle{\pgfqpoint{0.050000in}{0.050000in}}{\pgfqpoint{4.900000in}{4.900000in}}%
\pgfusepath{clip}%
\pgfsetbuttcap%
\pgfsetroundjoin%
\definecolor{currentfill}{rgb}{0.145790,0.145790,0.145790}%
\pgfsetfillcolor{currentfill}%
\pgfsetfillopacity{0.900000}%
\pgfsetlinewidth{0.501875pt}%
\definecolor{currentstroke}{rgb}{0.200000,0.200000,0.200000}%
\pgfsetstrokecolor{currentstroke}%
\pgfsetdash{}{0pt}%
\pgfpathmoveto{\pgfqpoint{1.667613in}{3.029462in}}%
\pgfpathlineto{\pgfqpoint{1.717429in}{2.981495in}}%
\pgfpathlineto{\pgfqpoint{1.749589in}{2.961105in}}%
\pgfpathlineto{\pgfqpoint{1.699763in}{3.009104in}}%
\pgfpathlineto{\pgfqpoint{1.667613in}{3.029462in}}%
\pgfpathclose%
\pgfusepath{stroke,fill}%
\end{pgfscope}%
\begin{pgfscope}%
\pgfpathrectangle{\pgfqpoint{0.050000in}{0.050000in}}{\pgfqpoint{4.900000in}{4.900000in}}%
\pgfusepath{clip}%
\pgfsetbuttcap%
\pgfsetroundjoin%
\definecolor{currentfill}{rgb}{0.436263,0.436263,0.436263}%
\pgfsetfillcolor{currentfill}%
\pgfsetfillopacity{0.900000}%
\pgfsetlinewidth{0.501875pt}%
\definecolor{currentstroke}{rgb}{0.200000,0.200000,0.200000}%
\pgfsetstrokecolor{currentstroke}%
\pgfsetdash{}{0pt}%
\pgfpathmoveto{\pgfqpoint{2.420163in}{2.421273in}}%
\pgfpathlineto{\pgfqpoint{2.468312in}{2.374655in}}%
\pgfpathlineto{\pgfqpoint{2.501694in}{2.353279in}}%
\pgfpathlineto{\pgfqpoint{2.453541in}{2.399926in}}%
\pgfpathlineto{\pgfqpoint{2.420163in}{2.421273in}}%
\pgfpathclose%
\pgfusepath{stroke,fill}%
\end{pgfscope}%
\begin{pgfscope}%
\pgfpathrectangle{\pgfqpoint{0.050000in}{0.050000in}}{\pgfqpoint{4.900000in}{4.900000in}}%
\pgfusepath{clip}%
\pgfsetbuttcap%
\pgfsetroundjoin%
\definecolor{currentfill}{rgb}{0.040969,0.040969,0.040969}%
\pgfsetfillcolor{currentfill}%
\pgfsetfillopacity{0.900000}%
\pgfsetlinewidth{0.501875pt}%
\definecolor{currentstroke}{rgb}{0.200000,0.200000,0.200000}%
\pgfsetstrokecolor{currentstroke}%
\pgfsetdash{}{0pt}%
\pgfpathmoveto{\pgfqpoint{1.389389in}{3.255021in}}%
\pgfpathlineto{\pgfqpoint{1.439833in}{3.206545in}}%
\pgfpathlineto{\pgfqpoint{1.471547in}{3.186516in}}%
\pgfpathlineto{\pgfqpoint{1.421092in}{3.235026in}}%
\pgfpathlineto{\pgfqpoint{1.389389in}{3.255021in}}%
\pgfpathclose%
\pgfusepath{stroke,fill}%
\end{pgfscope}%
\begin{pgfscope}%
\pgfpathrectangle{\pgfqpoint{0.050000in}{0.050000in}}{\pgfqpoint{4.900000in}{4.900000in}}%
\pgfusepath{clip}%
\pgfsetbuttcap%
\pgfsetroundjoin%
\definecolor{currentfill}{rgb}{0.346943,0.346943,0.346943}%
\pgfsetfillcolor{currentfill}%
\pgfsetfillopacity{0.900000}%
\pgfsetlinewidth{0.501875pt}%
\definecolor{currentstroke}{rgb}{0.200000,0.200000,0.200000}%
\pgfsetstrokecolor{currentstroke}%
\pgfsetdash{}{0pt}%
\pgfpathmoveto{\pgfqpoint{2.142053in}{2.646733in}}%
\pgfpathlineto{\pgfqpoint{2.190828in}{2.599605in}}%
\pgfpathlineto{\pgfqpoint{2.223765in}{2.578591in}}%
\pgfpathlineto{\pgfqpoint{2.174984in}{2.625749in}}%
\pgfpathlineto{\pgfqpoint{2.142053in}{2.646733in}}%
\pgfpathclose%
\pgfusepath{stroke,fill}%
\end{pgfscope}%
\begin{pgfscope}%
\pgfpathrectangle{\pgfqpoint{0.050000in}{0.050000in}}{\pgfqpoint{4.900000in}{4.900000in}}%
\pgfusepath{clip}%
\pgfsetbuttcap%
\pgfsetroundjoin%
\definecolor{currentfill}{rgb}{0.600231,0.600231,0.600231}%
\pgfsetfillcolor{currentfill}%
\pgfsetfillopacity{0.900000}%
\pgfsetlinewidth{0.501875pt}%
\definecolor{currentstroke}{rgb}{0.200000,0.200000,0.200000}%
\pgfsetstrokecolor{currentstroke}%
\pgfsetdash{}{0pt}%
\pgfpathmoveto{\pgfqpoint{2.894825in}{2.038510in}}%
\pgfpathlineto{\pgfqpoint{2.941931in}{1.992862in}}%
\pgfpathlineto{\pgfqpoint{2.976091in}{1.970893in}}%
\pgfpathlineto{\pgfqpoint{2.928984in}{2.016553in}}%
\pgfpathlineto{\pgfqpoint{2.894825in}{2.038510in}}%
\pgfpathclose%
\pgfusepath{stroke,fill}%
\end{pgfscope}%
\begin{pgfscope}%
\pgfpathrectangle{\pgfqpoint{0.050000in}{0.050000in}}{\pgfqpoint{4.900000in}{4.900000in}}%
\pgfusepath{clip}%
\pgfsetbuttcap%
\pgfsetroundjoin%
\definecolor{currentfill}{rgb}{0.228835,0.228835,0.228835}%
\pgfsetfillcolor{currentfill}%
\pgfsetfillopacity{0.900000}%
\pgfsetlinewidth{0.501875pt}%
\definecolor{currentstroke}{rgb}{0.200000,0.200000,0.200000}%
\pgfsetstrokecolor{currentstroke}%
\pgfsetdash{}{0pt}%
\pgfpathmoveto{\pgfqpoint{1.863845in}{2.872276in}}%
\pgfpathlineto{\pgfqpoint{1.913246in}{2.824638in}}%
\pgfpathlineto{\pgfqpoint{1.945738in}{2.803985in}}%
\pgfpathlineto{\pgfqpoint{1.896328in}{2.851654in}}%
\pgfpathlineto{\pgfqpoint{1.863845in}{2.872276in}}%
\pgfpathclose%
\pgfusepath{stroke,fill}%
\end{pgfscope}%
\begin{pgfscope}%
\pgfpathrectangle{\pgfqpoint{0.050000in}{0.050000in}}{\pgfqpoint{4.900000in}{4.900000in}}%
\pgfusepath{clip}%
\pgfsetbuttcap%
\pgfsetroundjoin%
\definecolor{currentfill}{rgb}{0.499962,0.499962,0.499962}%
\pgfsetfillcolor{currentfill}%
\pgfsetfillopacity{0.900000}%
\pgfsetlinewidth{0.501875pt}%
\definecolor{currentstroke}{rgb}{0.200000,0.200000,0.200000}%
\pgfsetstrokecolor{currentstroke}%
\pgfsetdash{}{0pt}%
\pgfpathmoveto{\pgfqpoint{2.616729in}{2.263829in}}%
\pgfpathlineto{\pgfqpoint{2.664460in}{2.217548in}}%
\pgfpathlineto{\pgfqpoint{2.698175in}{2.195912in}}%
\pgfpathlineto{\pgfqpoint{2.650441in}{2.242217in}}%
\pgfpathlineto{\pgfqpoint{2.616729in}{2.263829in}}%
\pgfpathclose%
\pgfusepath{stroke,fill}%
\end{pgfscope}%
\begin{pgfscope}%
\pgfpathrectangle{\pgfqpoint{0.050000in}{0.050000in}}{\pgfqpoint{4.900000in}{4.900000in}}%
\pgfusepath{clip}%
\pgfsetbuttcap%
\pgfsetroundjoin%
\definecolor{currentfill}{rgb}{0.109250,0.109250,0.109250}%
\pgfsetfillcolor{currentfill}%
\pgfsetfillopacity{0.900000}%
\pgfsetlinewidth{0.501875pt}%
\definecolor{currentstroke}{rgb}{0.200000,0.200000,0.200000}%
\pgfsetstrokecolor{currentstroke}%
\pgfsetdash{}{0pt}%
\pgfpathmoveto{\pgfqpoint{1.585538in}{3.097901in}}%
\pgfpathlineto{\pgfqpoint{1.635567in}{3.049753in}}%
\pgfpathlineto{\pgfqpoint{1.667613in}{3.029462in}}%
\pgfpathlineto{\pgfqpoint{1.617574in}{3.077643in}}%
\pgfpathlineto{\pgfqpoint{1.585538in}{3.097901in}}%
\pgfpathclose%
\pgfusepath{stroke,fill}%
\end{pgfscope}%
\begin{pgfscope}%
\pgfpathrectangle{\pgfqpoint{0.050000in}{0.050000in}}{\pgfqpoint{4.900000in}{4.900000in}}%
\pgfusepath{clip}%
\pgfsetbuttcap%
\pgfsetroundjoin%
\definecolor{currentfill}{rgb}{0.407843,0.407843,0.407843}%
\pgfsetfillcolor{currentfill}%
\pgfsetfillopacity{0.900000}%
\pgfsetlinewidth{0.501875pt}%
\definecolor{currentstroke}{rgb}{0.200000,0.200000,0.200000}%
\pgfsetstrokecolor{currentstroke}%
\pgfsetdash{}{0pt}%
\pgfpathmoveto{\pgfqpoint{2.338535in}{2.489349in}}%
\pgfpathlineto{\pgfqpoint{2.386894in}{2.442551in}}%
\pgfpathlineto{\pgfqpoint{2.420163in}{2.421273in}}%
\pgfpathlineto{\pgfqpoint{2.371800in}{2.468100in}}%
\pgfpathlineto{\pgfqpoint{2.338535in}{2.489349in}}%
\pgfpathclose%
\pgfusepath{stroke,fill}%
\end{pgfscope}%
\begin{pgfscope}%
\pgfpathrectangle{\pgfqpoint{0.050000in}{0.050000in}}{\pgfqpoint{4.900000in}{4.900000in}}%
\pgfusepath{clip}%
\pgfsetbuttcap%
\pgfsetroundjoin%
\definecolor{currentfill}{rgb}{0.672203,0.672203,0.672203}%
\pgfsetfillcolor{currentfill}%
\pgfsetfillopacity{0.900000}%
\pgfsetlinewidth{0.501875pt}%
\definecolor{currentstroke}{rgb}{0.200000,0.200000,0.200000}%
\pgfsetstrokecolor{currentstroke}%
\pgfsetdash{}{0pt}%
\pgfpathmoveto{\pgfqpoint{3.091646in}{1.881405in}}%
\pgfpathlineto{\pgfqpoint{3.138339in}{1.836474in}}%
\pgfpathlineto{\pgfqpoint{3.172835in}{1.814313in}}%
\pgfpathlineto{\pgfqpoint{3.126141in}{1.859230in}}%
\pgfpathlineto{\pgfqpoint{3.091646in}{1.881405in}}%
\pgfpathclose%
\pgfusepath{stroke,fill}%
\end{pgfscope}%
\begin{pgfscope}%
\pgfpathrectangle{\pgfqpoint{0.050000in}{0.050000in}}{\pgfqpoint{4.900000in}{4.900000in}}%
\pgfusepath{clip}%
\pgfsetbuttcap%
\pgfsetroundjoin%
\definecolor{currentfill}{rgb}{0.009104,0.009104,0.009104}%
\pgfsetfillcolor{currentfill}%
\pgfsetfillopacity{0.900000}%
\pgfsetlinewidth{0.501875pt}%
\definecolor{currentstroke}{rgb}{0.200000,0.200000,0.200000}%
\pgfsetstrokecolor{currentstroke}%
\pgfsetdash{}{0pt}%
\pgfpathmoveto{\pgfqpoint{1.307133in}{3.323608in}}%
\pgfpathlineto{\pgfqpoint{1.357789in}{3.274950in}}%
\pgfpathlineto{\pgfqpoint{1.389389in}{3.255021in}}%
\pgfpathlineto{\pgfqpoint{1.338721in}{3.303713in}}%
\pgfpathlineto{\pgfqpoint{1.307133in}{3.323608in}}%
\pgfpathclose%
\pgfusepath{stroke,fill}%
\end{pgfscope}%
\begin{pgfscope}%
\pgfpathrectangle{\pgfqpoint{0.050000in}{0.050000in}}{\pgfqpoint{4.900000in}{4.900000in}}%
\pgfusepath{clip}%
\pgfsetbuttcap%
\pgfsetroundjoin%
\definecolor{currentfill}{rgb}{0.311880,0.311880,0.311880}%
\pgfsetfillcolor{currentfill}%
\pgfsetfillopacity{0.900000}%
\pgfsetlinewidth{0.501875pt}%
\definecolor{currentstroke}{rgb}{0.200000,0.200000,0.200000}%
\pgfsetstrokecolor{currentstroke}%
\pgfsetdash{}{0pt}%
\pgfpathmoveto{\pgfqpoint{2.060244in}{2.714958in}}%
\pgfpathlineto{\pgfqpoint{2.109229in}{2.667650in}}%
\pgfpathlineto{\pgfqpoint{2.142053in}{2.646733in}}%
\pgfpathlineto{\pgfqpoint{2.093061in}{2.694072in}}%
\pgfpathlineto{\pgfqpoint{2.060244in}{2.714958in}}%
\pgfpathclose%
\pgfusepath{stroke,fill}%
\end{pgfscope}%
\begin{pgfscope}%
\pgfpathrectangle{\pgfqpoint{0.050000in}{0.050000in}}{\pgfqpoint{4.900000in}{4.900000in}}%
\pgfusepath{clip}%
\pgfsetbuttcap%
\pgfsetroundjoin%
\definecolor{currentfill}{rgb}{0.568858,0.568858,0.568858}%
\pgfsetfillcolor{currentfill}%
\pgfsetfillopacity{0.900000}%
\pgfsetlinewidth{0.501875pt}%
\definecolor{currentstroke}{rgb}{0.200000,0.200000,0.200000}%
\pgfsetstrokecolor{currentstroke}%
\pgfsetdash{}{0pt}%
\pgfpathmoveto{\pgfqpoint{2.813461in}{2.106295in}}%
\pgfpathlineto{\pgfqpoint{2.860776in}{2.060398in}}%
\pgfpathlineto{\pgfqpoint{2.894825in}{2.038510in}}%
\pgfpathlineto{\pgfqpoint{2.847508in}{2.084426in}}%
\pgfpathlineto{\pgfqpoint{2.813461in}{2.106295in}}%
\pgfpathclose%
\pgfusepath{stroke,fill}%
\end{pgfscope}%
\begin{pgfscope}%
\pgfpathrectangle{\pgfqpoint{0.050000in}{0.050000in}}{\pgfqpoint{4.900000in}{4.900000in}}%
\pgfusepath{clip}%
\pgfsetbuttcap%
\pgfsetroundjoin%
\definecolor{currentfill}{rgb}{0.184544,0.184544,0.184544}%
\pgfsetfillcolor{currentfill}%
\pgfsetfillopacity{0.900000}%
\pgfsetlinewidth{0.501875pt}%
\definecolor{currentstroke}{rgb}{0.200000,0.200000,0.200000}%
\pgfsetstrokecolor{currentstroke}%
\pgfsetdash{}{0pt}%
\pgfpathmoveto{\pgfqpoint{1.781853in}{2.940649in}}%
\pgfpathlineto{\pgfqpoint{1.831467in}{2.892831in}}%
\pgfpathlineto{\pgfqpoint{1.863845in}{2.872276in}}%
\pgfpathlineto{\pgfqpoint{1.814223in}{2.920126in}}%
\pgfpathlineto{\pgfqpoint{1.781853in}{2.940649in}}%
\pgfpathclose%
\pgfusepath{stroke,fill}%
\end{pgfscope}%
\begin{pgfscope}%
\pgfpathrectangle{\pgfqpoint{0.050000in}{0.050000in}}{\pgfqpoint{4.900000in}{4.900000in}}%
\pgfusepath{clip}%
\pgfsetbuttcap%
\pgfsetroundjoin%
\definecolor{currentfill}{rgb}{0.469819,0.469819,0.469819}%
\pgfsetfillcolor{currentfill}%
\pgfsetfillopacity{0.900000}%
\pgfsetlinewidth{0.501875pt}%
\definecolor{currentstroke}{rgb}{0.200000,0.200000,0.200000}%
\pgfsetstrokecolor{currentstroke}%
\pgfsetdash{}{0pt}%
\pgfpathmoveto{\pgfqpoint{2.535184in}{2.331835in}}%
\pgfpathlineto{\pgfqpoint{2.583126in}{2.285370in}}%
\pgfpathlineto{\pgfqpoint{2.616729in}{2.263829in}}%
\pgfpathlineto{\pgfqpoint{2.568784in}{2.310320in}}%
\pgfpathlineto{\pgfqpoint{2.535184in}{2.331835in}}%
\pgfpathclose%
\pgfusepath{stroke,fill}%
\end{pgfscope}%
\begin{pgfscope}%
\pgfpathrectangle{\pgfqpoint{0.050000in}{0.050000in}}{\pgfqpoint{4.900000in}{4.900000in}}%
\pgfusepath{clip}%
\pgfsetbuttcap%
\pgfsetroundjoin%
\definecolor{currentfill}{rgb}{0.077386,0.077386,0.077386}%
\pgfsetfillcolor{currentfill}%
\pgfsetfillopacity{0.900000}%
\pgfsetlinewidth{0.501875pt}%
\definecolor{currentstroke}{rgb}{0.200000,0.200000,0.200000}%
\pgfsetstrokecolor{currentstroke}%
\pgfsetdash{}{0pt}%
\pgfpathmoveto{\pgfqpoint{1.503365in}{3.166422in}}%
\pgfpathlineto{\pgfqpoint{1.553606in}{3.118093in}}%
\pgfpathlineto{\pgfqpoint{1.585538in}{3.097901in}}%
\pgfpathlineto{\pgfqpoint{1.535286in}{3.146263in}}%
\pgfpathlineto{\pgfqpoint{1.503365in}{3.166422in}}%
\pgfpathclose%
\pgfusepath{stroke,fill}%
\end{pgfscope}%
\begin{pgfscope}%
\pgfpathrectangle{\pgfqpoint{0.050000in}{0.050000in}}{\pgfqpoint{4.900000in}{4.900000in}}%
\pgfusepath{clip}%
\pgfsetbuttcap%
\pgfsetroundjoin%
\definecolor{currentfill}{rgb}{0.375363,0.375363,0.375363}%
\pgfsetfillcolor{currentfill}%
\pgfsetfillopacity{0.900000}%
\pgfsetlinewidth{0.501875pt}%
\definecolor{currentstroke}{rgb}{0.200000,0.200000,0.200000}%
\pgfsetstrokecolor{currentstroke}%
\pgfsetdash{}{0pt}%
\pgfpathmoveto{\pgfqpoint{2.256809in}{2.557508in}}%
\pgfpathlineto{\pgfqpoint{2.305378in}{2.510530in}}%
\pgfpathlineto{\pgfqpoint{2.338535in}{2.489349in}}%
\pgfpathlineto{\pgfqpoint{2.289961in}{2.536357in}}%
\pgfpathlineto{\pgfqpoint{2.256809in}{2.557508in}}%
\pgfpathclose%
\pgfusepath{stroke,fill}%
\end{pgfscope}%
\begin{pgfscope}%
\pgfpathrectangle{\pgfqpoint{0.050000in}{0.050000in}}{\pgfqpoint{4.900000in}{4.900000in}}%
\pgfusepath{clip}%
\pgfsetbuttcap%
\pgfsetroundjoin%
\definecolor{currentfill}{rgb}{0.638616,0.638616,0.638616}%
\pgfsetfillcolor{currentfill}%
\pgfsetfillopacity{0.900000}%
\pgfsetlinewidth{0.501875pt}%
\definecolor{currentstroke}{rgb}{0.200000,0.200000,0.200000}%
\pgfsetstrokecolor{currentstroke}%
\pgfsetdash{}{0pt}%
\pgfpathmoveto{\pgfqpoint{3.010363in}{1.948855in}}%
\pgfpathlineto{\pgfqpoint{3.057262in}{1.903512in}}%
\pgfpathlineto{\pgfqpoint{3.091646in}{1.881405in}}%
\pgfpathlineto{\pgfqpoint{3.044746in}{1.926749in}}%
\pgfpathlineto{\pgfqpoint{3.010363in}{1.948855in}}%
\pgfpathclose%
\pgfusepath{stroke,fill}%
\end{pgfscope}%
\begin{pgfscope}%
\pgfpathrectangle{\pgfqpoint{0.050000in}{0.050000in}}{\pgfqpoint{4.900000in}{4.900000in}}%
\pgfusepath{clip}%
\pgfsetbuttcap%
\pgfsetroundjoin%
\definecolor{currentfill}{rgb}{0.273126,0.273126,0.273126}%
\pgfsetfillcolor{currentfill}%
\pgfsetfillopacity{0.900000}%
\pgfsetlinewidth{0.501875pt}%
\definecolor{currentstroke}{rgb}{0.200000,0.200000,0.200000}%
\pgfsetstrokecolor{currentstroke}%
\pgfsetdash{}{0pt}%
\pgfpathmoveto{\pgfqpoint{1.978335in}{2.783265in}}%
\pgfpathlineto{\pgfqpoint{2.027533in}{2.735776in}}%
\pgfpathlineto{\pgfqpoint{2.060244in}{2.714958in}}%
\pgfpathlineto{\pgfqpoint{2.011039in}{2.762477in}}%
\pgfpathlineto{\pgfqpoint{1.978335in}{2.783265in}}%
\pgfpathclose%
\pgfusepath{stroke,fill}%
\end{pgfscope}%
\begin{pgfscope}%
\pgfpathrectangle{\pgfqpoint{0.050000in}{0.050000in}}{\pgfqpoint{4.900000in}{4.900000in}}%
\pgfusepath{clip}%
\pgfsetbuttcap%
\pgfsetroundjoin%
\definecolor{currentfill}{rgb}{0.534410,0.534410,0.534410}%
\pgfsetfillcolor{currentfill}%
\pgfsetfillopacity{0.900000}%
\pgfsetlinewidth{0.501875pt}%
\definecolor{currentstroke}{rgb}{0.200000,0.200000,0.200000}%
\pgfsetstrokecolor{currentstroke}%
\pgfsetdash{}{0pt}%
\pgfpathmoveto{\pgfqpoint{2.732000in}{2.174205in}}%
\pgfpathlineto{\pgfqpoint{2.779524in}{2.128094in}}%
\pgfpathlineto{\pgfqpoint{2.813461in}{2.106295in}}%
\pgfpathlineto{\pgfqpoint{2.765935in}{2.152428in}}%
\pgfpathlineto{\pgfqpoint{2.732000in}{2.174205in}}%
\pgfpathclose%
\pgfusepath{stroke,fill}%
\end{pgfscope}%
\begin{pgfscope}%
\pgfpathrectangle{\pgfqpoint{0.050000in}{0.050000in}}{\pgfqpoint{4.900000in}{4.900000in}}%
\pgfusepath{clip}%
\pgfsetbuttcap%
\pgfsetroundjoin%
\definecolor{currentfill}{rgb}{0.145790,0.145790,0.145790}%
\pgfsetfillcolor{currentfill}%
\pgfsetfillopacity{0.900000}%
\pgfsetlinewidth{0.501875pt}%
\definecolor{currentstroke}{rgb}{0.200000,0.200000,0.200000}%
\pgfsetstrokecolor{currentstroke}%
\pgfsetdash{}{0pt}%
\pgfpathmoveto{\pgfqpoint{1.699763in}{3.009104in}}%
\pgfpathlineto{\pgfqpoint{1.749589in}{2.961105in}}%
\pgfpathlineto{\pgfqpoint{1.781853in}{2.940649in}}%
\pgfpathlineto{\pgfqpoint{1.732018in}{2.988681in}}%
\pgfpathlineto{\pgfqpoint{1.699763in}{3.009104in}}%
\pgfpathclose%
\pgfusepath{stroke,fill}%
\end{pgfscope}%
\begin{pgfscope}%
\pgfpathrectangle{\pgfqpoint{0.050000in}{0.050000in}}{\pgfqpoint{4.900000in}{4.900000in}}%
\pgfusepath{clip}%
\pgfsetbuttcap%
\pgfsetroundjoin%
\definecolor{currentfill}{rgb}{0.436263,0.436263,0.436263}%
\pgfsetfillcolor{currentfill}%
\pgfsetfillopacity{0.900000}%
\pgfsetlinewidth{0.501875pt}%
\definecolor{currentstroke}{rgb}{0.200000,0.200000,0.200000}%
\pgfsetstrokecolor{currentstroke}%
\pgfsetdash{}{0pt}%
\pgfpathmoveto{\pgfqpoint{2.453541in}{2.399926in}}%
\pgfpathlineto{\pgfqpoint{2.501694in}{2.353279in}}%
\pgfpathlineto{\pgfqpoint{2.535184in}{2.331835in}}%
\pgfpathlineto{\pgfqpoint{2.487028in}{2.378509in}}%
\pgfpathlineto{\pgfqpoint{2.453541in}{2.399926in}}%
\pgfpathclose%
\pgfusepath{stroke,fill}%
\end{pgfscope}%
\begin{pgfscope}%
\pgfpathrectangle{\pgfqpoint{0.050000in}{0.050000in}}{\pgfqpoint{4.900000in}{4.900000in}}%
\pgfusepath{clip}%
\pgfsetbuttcap%
\pgfsetroundjoin%
\definecolor{currentfill}{rgb}{0.040969,0.040969,0.040969}%
\pgfsetfillcolor{currentfill}%
\pgfsetfillopacity{0.900000}%
\pgfsetlinewidth{0.501875pt}%
\definecolor{currentstroke}{rgb}{0.200000,0.200000,0.200000}%
\pgfsetstrokecolor{currentstroke}%
\pgfsetdash{}{0pt}%
\pgfpathmoveto{\pgfqpoint{1.421092in}{3.235026in}}%
\pgfpathlineto{\pgfqpoint{1.471547in}{3.186516in}}%
\pgfpathlineto{\pgfqpoint{1.503365in}{3.166422in}}%
\pgfpathlineto{\pgfqpoint{1.452899in}{3.214967in}}%
\pgfpathlineto{\pgfqpoint{1.421092in}{3.235026in}}%
\pgfpathclose%
\pgfusepath{stroke,fill}%
\end{pgfscope}%
\begin{pgfscope}%
\pgfpathrectangle{\pgfqpoint{0.050000in}{0.050000in}}{\pgfqpoint{4.900000in}{4.900000in}}%
\pgfusepath{clip}%
\pgfsetbuttcap%
\pgfsetroundjoin%
\definecolor{currentfill}{rgb}{0.346943,0.346943,0.346943}%
\pgfsetfillcolor{currentfill}%
\pgfsetfillopacity{0.900000}%
\pgfsetlinewidth{0.501875pt}%
\definecolor{currentstroke}{rgb}{0.200000,0.200000,0.200000}%
\pgfsetstrokecolor{currentstroke}%
\pgfsetdash{}{0pt}%
\pgfpathmoveto{\pgfqpoint{2.174984in}{2.625749in}}%
\pgfpathlineto{\pgfqpoint{2.223765in}{2.578591in}}%
\pgfpathlineto{\pgfqpoint{2.256809in}{2.557508in}}%
\pgfpathlineto{\pgfqpoint{2.208023in}{2.604696in}}%
\pgfpathlineto{\pgfqpoint{2.174984in}{2.625749in}}%
\pgfpathclose%
\pgfusepath{stroke,fill}%
\end{pgfscope}%
\begin{pgfscope}%
\pgfpathrectangle{\pgfqpoint{0.050000in}{0.050000in}}{\pgfqpoint{4.900000in}{4.900000in}}%
\pgfusepath{clip}%
\pgfsetbuttcap%
\pgfsetroundjoin%
\definecolor{currentfill}{rgb}{0.600231,0.600231,0.600231}%
\pgfsetfillcolor{currentfill}%
\pgfsetfillopacity{0.900000}%
\pgfsetlinewidth{0.501875pt}%
\definecolor{currentstroke}{rgb}{0.200000,0.200000,0.200000}%
\pgfsetstrokecolor{currentstroke}%
\pgfsetdash{}{0pt}%
\pgfpathmoveto{\pgfqpoint{2.928984in}{2.016553in}}%
\pgfpathlineto{\pgfqpoint{2.976091in}{1.970893in}}%
\pgfpathlineto{\pgfqpoint{3.010363in}{1.948855in}}%
\pgfpathlineto{\pgfqpoint{2.963255in}{1.994526in}}%
\pgfpathlineto{\pgfqpoint{2.928984in}{2.016553in}}%
\pgfpathclose%
\pgfusepath{stroke,fill}%
\end{pgfscope}%
\begin{pgfscope}%
\pgfpathrectangle{\pgfqpoint{0.050000in}{0.050000in}}{\pgfqpoint{4.900000in}{4.900000in}}%
\pgfusepath{clip}%
\pgfsetbuttcap%
\pgfsetroundjoin%
\definecolor{currentfill}{rgb}{0.228835,0.228835,0.228835}%
\pgfsetfillcolor{currentfill}%
\pgfsetfillopacity{0.900000}%
\pgfsetlinewidth{0.501875pt}%
\definecolor{currentstroke}{rgb}{0.200000,0.200000,0.200000}%
\pgfsetstrokecolor{currentstroke}%
\pgfsetdash{}{0pt}%
\pgfpathmoveto{\pgfqpoint{1.896328in}{2.851654in}}%
\pgfpathlineto{\pgfqpoint{1.945738in}{2.803985in}}%
\pgfpathlineto{\pgfqpoint{1.978335in}{2.783265in}}%
\pgfpathlineto{\pgfqpoint{1.928918in}{2.830966in}}%
\pgfpathlineto{\pgfqpoint{1.896328in}{2.851654in}}%
\pgfpathclose%
\pgfusepath{stroke,fill}%
\end{pgfscope}%
\begin{pgfscope}%
\pgfpathrectangle{\pgfqpoint{0.050000in}{0.050000in}}{\pgfqpoint{4.900000in}{4.900000in}}%
\pgfusepath{clip}%
\pgfsetbuttcap%
\pgfsetroundjoin%
\definecolor{currentfill}{rgb}{0.504268,0.504268,0.504268}%
\pgfsetfillcolor{currentfill}%
\pgfsetfillopacity{0.900000}%
\pgfsetlinewidth{0.501875pt}%
\definecolor{currentstroke}{rgb}{0.200000,0.200000,0.200000}%
\pgfsetstrokecolor{currentstroke}%
\pgfsetdash{}{0pt}%
\pgfpathmoveto{\pgfqpoint{2.650441in}{2.242217in}}%
\pgfpathlineto{\pgfqpoint{2.698175in}{2.195912in}}%
\pgfpathlineto{\pgfqpoint{2.732000in}{2.174205in}}%
\pgfpathlineto{\pgfqpoint{2.684263in}{2.220536in}}%
\pgfpathlineto{\pgfqpoint{2.650441in}{2.242217in}}%
\pgfpathclose%
\pgfusepath{stroke,fill}%
\end{pgfscope}%
\begin{pgfscope}%
\pgfpathrectangle{\pgfqpoint{0.050000in}{0.050000in}}{\pgfqpoint{4.900000in}{4.900000in}}%
\pgfusepath{clip}%
\pgfsetbuttcap%
\pgfsetroundjoin%
\definecolor{currentfill}{rgb}{0.109250,0.109250,0.109250}%
\pgfsetfillcolor{currentfill}%
\pgfsetfillopacity{0.900000}%
\pgfsetlinewidth{0.501875pt}%
\definecolor{currentstroke}{rgb}{0.200000,0.200000,0.200000}%
\pgfsetstrokecolor{currentstroke}%
\pgfsetdash{}{0pt}%
\pgfpathmoveto{\pgfqpoint{1.617574in}{3.077643in}}%
\pgfpathlineto{\pgfqpoint{1.667613in}{3.029462in}}%
\pgfpathlineto{\pgfqpoint{1.699763in}{3.009104in}}%
\pgfpathlineto{\pgfqpoint{1.649715in}{3.057318in}}%
\pgfpathlineto{\pgfqpoint{1.617574in}{3.077643in}}%
\pgfpathclose%
\pgfusepath{stroke,fill}%
\end{pgfscope}%
\begin{pgfscope}%
\pgfpathrectangle{\pgfqpoint{0.050000in}{0.050000in}}{\pgfqpoint{4.900000in}{4.900000in}}%
\pgfusepath{clip}%
\pgfsetbuttcap%
\pgfsetroundjoin%
\definecolor{currentfill}{rgb}{0.407843,0.407843,0.407843}%
\pgfsetfillcolor{currentfill}%
\pgfsetfillopacity{0.900000}%
\pgfsetlinewidth{0.501875pt}%
\definecolor{currentstroke}{rgb}{0.200000,0.200000,0.200000}%
\pgfsetstrokecolor{currentstroke}%
\pgfsetdash{}{0pt}%
\pgfpathmoveto{\pgfqpoint{2.371800in}{2.468100in}}%
\pgfpathlineto{\pgfqpoint{2.420163in}{2.421273in}}%
\pgfpathlineto{\pgfqpoint{2.453541in}{2.399926in}}%
\pgfpathlineto{\pgfqpoint{2.405174in}{2.446781in}}%
\pgfpathlineto{\pgfqpoint{2.371800in}{2.468100in}}%
\pgfpathclose%
\pgfusepath{stroke,fill}%
\end{pgfscope}%
\begin{pgfscope}%
\pgfpathrectangle{\pgfqpoint{0.050000in}{0.050000in}}{\pgfqpoint{4.900000in}{4.900000in}}%
\pgfusepath{clip}%
\pgfsetbuttcap%
\pgfsetroundjoin%
\definecolor{currentfill}{rgb}{0.009104,0.009104,0.009104}%
\pgfsetfillcolor{currentfill}%
\pgfsetfillopacity{0.900000}%
\pgfsetlinewidth{0.501875pt}%
\definecolor{currentstroke}{rgb}{0.200000,0.200000,0.200000}%
\pgfsetstrokecolor{currentstroke}%
\pgfsetdash{}{0pt}%
\pgfpathmoveto{\pgfqpoint{1.338721in}{3.303713in}}%
\pgfpathlineto{\pgfqpoint{1.389389in}{3.255021in}}%
\pgfpathlineto{\pgfqpoint{1.421092in}{3.235026in}}%
\pgfpathlineto{\pgfqpoint{1.370412in}{3.283754in}}%
\pgfpathlineto{\pgfqpoint{1.338721in}{3.303713in}}%
\pgfpathclose%
\pgfusepath{stroke,fill}%
\end{pgfscope}%
\begin{pgfscope}%
\pgfpathrectangle{\pgfqpoint{0.050000in}{0.050000in}}{\pgfqpoint{4.900000in}{4.900000in}}%
\pgfusepath{clip}%
\pgfsetbuttcap%
\pgfsetroundjoin%
\definecolor{currentfill}{rgb}{0.311880,0.311880,0.311880}%
\pgfsetfillcolor{currentfill}%
\pgfsetfillopacity{0.900000}%
\pgfsetlinewidth{0.501875pt}%
\definecolor{currentstroke}{rgb}{0.200000,0.200000,0.200000}%
\pgfsetstrokecolor{currentstroke}%
\pgfsetdash{}{0pt}%
\pgfpathmoveto{\pgfqpoint{2.093061in}{2.694072in}}%
\pgfpathlineto{\pgfqpoint{2.142053in}{2.646733in}}%
\pgfpathlineto{\pgfqpoint{2.174984in}{2.625749in}}%
\pgfpathlineto{\pgfqpoint{2.125986in}{2.673117in}}%
\pgfpathlineto{\pgfqpoint{2.093061in}{2.694072in}}%
\pgfpathclose%
\pgfusepath{stroke,fill}%
\end{pgfscope}%
\begin{pgfscope}%
\pgfpathrectangle{\pgfqpoint{0.050000in}{0.050000in}}{\pgfqpoint{4.900000in}{4.900000in}}%
\pgfusepath{clip}%
\pgfsetbuttcap%
\pgfsetroundjoin%
\definecolor{currentfill}{rgb}{0.568858,0.568858,0.568858}%
\pgfsetfillcolor{currentfill}%
\pgfsetfillopacity{0.900000}%
\pgfsetlinewidth{0.501875pt}%
\definecolor{currentstroke}{rgb}{0.200000,0.200000,0.200000}%
\pgfsetstrokecolor{currentstroke}%
\pgfsetdash{}{0pt}%
\pgfpathmoveto{\pgfqpoint{2.847508in}{2.084426in}}%
\pgfpathlineto{\pgfqpoint{2.894825in}{2.038510in}}%
\pgfpathlineto{\pgfqpoint{2.928984in}{2.016553in}}%
\pgfpathlineto{\pgfqpoint{2.881667in}{2.062486in}}%
\pgfpathlineto{\pgfqpoint{2.847508in}{2.084426in}}%
\pgfpathclose%
\pgfusepath{stroke,fill}%
\end{pgfscope}%
\begin{pgfscope}%
\pgfpathrectangle{\pgfqpoint{0.050000in}{0.050000in}}{\pgfqpoint{4.900000in}{4.900000in}}%
\pgfusepath{clip}%
\pgfsetbuttcap%
\pgfsetroundjoin%
\definecolor{currentfill}{rgb}{0.190081,0.190081,0.190081}%
\pgfsetfillcolor{currentfill}%
\pgfsetfillopacity{0.900000}%
\pgfsetlinewidth{0.501875pt}%
\definecolor{currentstroke}{rgb}{0.200000,0.200000,0.200000}%
\pgfsetstrokecolor{currentstroke}%
\pgfsetdash{}{0pt}%
\pgfpathmoveto{\pgfqpoint{1.814223in}{2.920126in}}%
\pgfpathlineto{\pgfqpoint{1.863845in}{2.872276in}}%
\pgfpathlineto{\pgfqpoint{1.896328in}{2.851654in}}%
\pgfpathlineto{\pgfqpoint{1.846698in}{2.899537in}}%
\pgfpathlineto{\pgfqpoint{1.814223in}{2.920126in}}%
\pgfpathclose%
\pgfusepath{stroke,fill}%
\end{pgfscope}%
\begin{pgfscope}%
\pgfpathrectangle{\pgfqpoint{0.050000in}{0.050000in}}{\pgfqpoint{4.900000in}{4.900000in}}%
\pgfusepath{clip}%
\pgfsetbuttcap%
\pgfsetroundjoin%
\definecolor{currentfill}{rgb}{0.469819,0.469819,0.469819}%
\pgfsetfillcolor{currentfill}%
\pgfsetfillopacity{0.900000}%
\pgfsetlinewidth{0.501875pt}%
\definecolor{currentstroke}{rgb}{0.200000,0.200000,0.200000}%
\pgfsetstrokecolor{currentstroke}%
\pgfsetdash{}{0pt}%
\pgfpathmoveto{\pgfqpoint{2.568784in}{2.310320in}}%
\pgfpathlineto{\pgfqpoint{2.616729in}{2.263829in}}%
\pgfpathlineto{\pgfqpoint{2.650441in}{2.242217in}}%
\pgfpathlineto{\pgfqpoint{2.602493in}{2.288736in}}%
\pgfpathlineto{\pgfqpoint{2.568784in}{2.310320in}}%
\pgfpathclose%
\pgfusepath{stroke,fill}%
\end{pgfscope}%
\begin{pgfscope}%
\pgfpathrectangle{\pgfqpoint{0.050000in}{0.050000in}}{\pgfqpoint{4.900000in}{4.900000in}}%
\pgfusepath{clip}%
\pgfsetbuttcap%
\pgfsetroundjoin%
\definecolor{currentfill}{rgb}{0.077386,0.077386,0.077386}%
\pgfsetfillcolor{currentfill}%
\pgfsetfillopacity{0.900000}%
\pgfsetlinewidth{0.501875pt}%
\definecolor{currentstroke}{rgb}{0.200000,0.200000,0.200000}%
\pgfsetstrokecolor{currentstroke}%
\pgfsetdash{}{0pt}%
\pgfpathmoveto{\pgfqpoint{1.535286in}{3.146263in}}%
\pgfpathlineto{\pgfqpoint{1.585538in}{3.097901in}}%
\pgfpathlineto{\pgfqpoint{1.617574in}{3.077643in}}%
\pgfpathlineto{\pgfqpoint{1.567312in}{3.126039in}}%
\pgfpathlineto{\pgfqpoint{1.535286in}{3.146263in}}%
\pgfpathclose%
\pgfusepath{stroke,fill}%
\end{pgfscope}%
\begin{pgfscope}%
\pgfpathrectangle{\pgfqpoint{0.050000in}{0.050000in}}{\pgfqpoint{4.900000in}{4.900000in}}%
\pgfusepath{clip}%
\pgfsetbuttcap%
\pgfsetroundjoin%
\definecolor{currentfill}{rgb}{0.375363,0.375363,0.375363}%
\pgfsetfillcolor{currentfill}%
\pgfsetfillopacity{0.900000}%
\pgfsetlinewidth{0.501875pt}%
\definecolor{currentstroke}{rgb}{0.200000,0.200000,0.200000}%
\pgfsetstrokecolor{currentstroke}%
\pgfsetdash{}{0pt}%
\pgfpathmoveto{\pgfqpoint{2.289961in}{2.536357in}}%
\pgfpathlineto{\pgfqpoint{2.338535in}{2.489349in}}%
\pgfpathlineto{\pgfqpoint{2.371800in}{2.468100in}}%
\pgfpathlineto{\pgfqpoint{2.323221in}{2.515136in}}%
\pgfpathlineto{\pgfqpoint{2.289961in}{2.536357in}}%
\pgfpathclose%
\pgfusepath{stroke,fill}%
\end{pgfscope}%
\begin{pgfscope}%
\pgfpathrectangle{\pgfqpoint{0.050000in}{0.050000in}}{\pgfqpoint{4.900000in}{4.900000in}}%
\pgfusepath{clip}%
\pgfsetbuttcap%
\pgfsetroundjoin%
\definecolor{currentfill}{rgb}{0.638616,0.638616,0.638616}%
\pgfsetfillcolor{currentfill}%
\pgfsetfillopacity{0.900000}%
\pgfsetlinewidth{0.501875pt}%
\definecolor{currentstroke}{rgb}{0.200000,0.200000,0.200000}%
\pgfsetstrokecolor{currentstroke}%
\pgfsetdash{}{0pt}%
\pgfpathmoveto{\pgfqpoint{3.044746in}{1.926749in}}%
\pgfpathlineto{\pgfqpoint{3.091646in}{1.881405in}}%
\pgfpathlineto{\pgfqpoint{3.126141in}{1.859230in}}%
\pgfpathlineto{\pgfqpoint{3.079242in}{1.904572in}}%
\pgfpathlineto{\pgfqpoint{3.044746in}{1.926749in}}%
\pgfpathclose%
\pgfusepath{stroke,fill}%
\end{pgfscope}%
\begin{pgfscope}%
\pgfpathrectangle{\pgfqpoint{0.050000in}{0.050000in}}{\pgfqpoint{4.900000in}{4.900000in}}%
\pgfusepath{clip}%
\pgfsetbuttcap%
\pgfsetroundjoin%
\definecolor{currentfill}{rgb}{0.273126,0.273126,0.273126}%
\pgfsetfillcolor{currentfill}%
\pgfsetfillopacity{0.900000}%
\pgfsetlinewidth{0.501875pt}%
\definecolor{currentstroke}{rgb}{0.200000,0.200000,0.200000}%
\pgfsetstrokecolor{currentstroke}%
\pgfsetdash{}{0pt}%
\pgfpathmoveto{\pgfqpoint{2.011039in}{2.762477in}}%
\pgfpathlineto{\pgfqpoint{2.060244in}{2.714958in}}%
\pgfpathlineto{\pgfqpoint{2.093061in}{2.694072in}}%
\pgfpathlineto{\pgfqpoint{2.043849in}{2.741622in}}%
\pgfpathlineto{\pgfqpoint{2.011039in}{2.762477in}}%
\pgfpathclose%
\pgfusepath{stroke,fill}%
\end{pgfscope}%
\begin{pgfscope}%
\pgfpathrectangle{\pgfqpoint{0.050000in}{0.050000in}}{\pgfqpoint{4.900000in}{4.900000in}}%
\pgfusepath{clip}%
\pgfsetbuttcap%
\pgfsetroundjoin%
\definecolor{currentfill}{rgb}{0.534410,0.534410,0.534410}%
\pgfsetfillcolor{currentfill}%
\pgfsetfillopacity{0.900000}%
\pgfsetlinewidth{0.501875pt}%
\definecolor{currentstroke}{rgb}{0.200000,0.200000,0.200000}%
\pgfsetstrokecolor{currentstroke}%
\pgfsetdash{}{0pt}%
\pgfpathmoveto{\pgfqpoint{2.765935in}{2.152428in}}%
\pgfpathlineto{\pgfqpoint{2.813461in}{2.106295in}}%
\pgfpathlineto{\pgfqpoint{2.847508in}{2.084426in}}%
\pgfpathlineto{\pgfqpoint{2.799981in}{2.130581in}}%
\pgfpathlineto{\pgfqpoint{2.765935in}{2.152428in}}%
\pgfpathclose%
\pgfusepath{stroke,fill}%
\end{pgfscope}%
\begin{pgfscope}%
\pgfpathrectangle{\pgfqpoint{0.050000in}{0.050000in}}{\pgfqpoint{4.900000in}{4.900000in}}%
\pgfusepath{clip}%
\pgfsetbuttcap%
\pgfsetroundjoin%
\definecolor{currentfill}{rgb}{0.145790,0.145790,0.145790}%
\pgfsetfillcolor{currentfill}%
\pgfsetfillopacity{0.900000}%
\pgfsetlinewidth{0.501875pt}%
\definecolor{currentstroke}{rgb}{0.200000,0.200000,0.200000}%
\pgfsetstrokecolor{currentstroke}%
\pgfsetdash{}{0pt}%
\pgfpathmoveto{\pgfqpoint{1.732018in}{2.988681in}}%
\pgfpathlineto{\pgfqpoint{1.781853in}{2.940649in}}%
\pgfpathlineto{\pgfqpoint{1.814223in}{2.920126in}}%
\pgfpathlineto{\pgfqpoint{1.764379in}{2.968191in}}%
\pgfpathlineto{\pgfqpoint{1.732018in}{2.988681in}}%
\pgfpathclose%
\pgfusepath{stroke,fill}%
\end{pgfscope}%
\begin{pgfscope}%
\pgfpathrectangle{\pgfqpoint{0.050000in}{0.050000in}}{\pgfqpoint{4.900000in}{4.900000in}}%
\pgfusepath{clip}%
\pgfsetbuttcap%
\pgfsetroundjoin%
\definecolor{currentfill}{rgb}{0.440323,0.440323,0.440323}%
\pgfsetfillcolor{currentfill}%
\pgfsetfillopacity{0.900000}%
\pgfsetlinewidth{0.501875pt}%
\definecolor{currentstroke}{rgb}{0.200000,0.200000,0.200000}%
\pgfsetstrokecolor{currentstroke}%
\pgfsetdash{}{0pt}%
\pgfpathmoveto{\pgfqpoint{2.487028in}{2.378509in}}%
\pgfpathlineto{\pgfqpoint{2.535184in}{2.331835in}}%
\pgfpathlineto{\pgfqpoint{2.568784in}{2.310320in}}%
\pgfpathlineto{\pgfqpoint{2.520624in}{2.357023in}}%
\pgfpathlineto{\pgfqpoint{2.487028in}{2.378509in}}%
\pgfpathclose%
\pgfusepath{stroke,fill}%
\end{pgfscope}%
\begin{pgfscope}%
\pgfpathrectangle{\pgfqpoint{0.050000in}{0.050000in}}{\pgfqpoint{4.900000in}{4.900000in}}%
\pgfusepath{clip}%
\pgfsetbuttcap%
\pgfsetroundjoin%
\definecolor{currentfill}{rgb}{0.040969,0.040969,0.040969}%
\pgfsetfillcolor{currentfill}%
\pgfsetfillopacity{0.900000}%
\pgfsetlinewidth{0.501875pt}%
\definecolor{currentstroke}{rgb}{0.200000,0.200000,0.200000}%
\pgfsetstrokecolor{currentstroke}%
\pgfsetdash{}{0pt}%
\pgfpathmoveto{\pgfqpoint{1.452899in}{3.214967in}}%
\pgfpathlineto{\pgfqpoint{1.503365in}{3.166422in}}%
\pgfpathlineto{\pgfqpoint{1.535286in}{3.146263in}}%
\pgfpathlineto{\pgfqpoint{1.484809in}{3.194842in}}%
\pgfpathlineto{\pgfqpoint{1.452899in}{3.214967in}}%
\pgfpathclose%
\pgfusepath{stroke,fill}%
\end{pgfscope}%
\begin{pgfscope}%
\pgfpathrectangle{\pgfqpoint{0.050000in}{0.050000in}}{\pgfqpoint{4.900000in}{4.900000in}}%
\pgfusepath{clip}%
\pgfsetbuttcap%
\pgfsetroundjoin%
\definecolor{currentfill}{rgb}{0.346943,0.346943,0.346943}%
\pgfsetfillcolor{currentfill}%
\pgfsetfillopacity{0.900000}%
\pgfsetlinewidth{0.501875pt}%
\definecolor{currentstroke}{rgb}{0.200000,0.200000,0.200000}%
\pgfsetstrokecolor{currentstroke}%
\pgfsetdash{}{0pt}%
\pgfpathmoveto{\pgfqpoint{2.208023in}{2.604696in}}%
\pgfpathlineto{\pgfqpoint{2.256809in}{2.557508in}}%
\pgfpathlineto{\pgfqpoint{2.289961in}{2.536357in}}%
\pgfpathlineto{\pgfqpoint{2.241169in}{2.583574in}}%
\pgfpathlineto{\pgfqpoint{2.208023in}{2.604696in}}%
\pgfpathclose%
\pgfusepath{stroke,fill}%
\end{pgfscope}%
\begin{pgfscope}%
\pgfpathrectangle{\pgfqpoint{0.050000in}{0.050000in}}{\pgfqpoint{4.900000in}{4.900000in}}%
\pgfusepath{clip}%
\pgfsetbuttcap%
\pgfsetroundjoin%
\definecolor{currentfill}{rgb}{0.600231,0.600231,0.600231}%
\pgfsetfillcolor{currentfill}%
\pgfsetfillopacity{0.900000}%
\pgfsetlinewidth{0.501875pt}%
\definecolor{currentstroke}{rgb}{0.200000,0.200000,0.200000}%
\pgfsetstrokecolor{currentstroke}%
\pgfsetdash{}{0pt}%
\pgfpathmoveto{\pgfqpoint{2.963255in}{1.994526in}}%
\pgfpathlineto{\pgfqpoint{3.010363in}{1.948855in}}%
\pgfpathlineto{\pgfqpoint{3.044746in}{1.926749in}}%
\pgfpathlineto{\pgfqpoint{2.997638in}{1.972429in}}%
\pgfpathlineto{\pgfqpoint{2.963255in}{1.994526in}}%
\pgfpathclose%
\pgfusepath{stroke,fill}%
\end{pgfscope}%
\begin{pgfscope}%
\pgfpathrectangle{\pgfqpoint{0.050000in}{0.050000in}}{\pgfqpoint{4.900000in}{4.900000in}}%
\pgfusepath{clip}%
\pgfsetbuttcap%
\pgfsetroundjoin%
\definecolor{currentfill}{rgb}{0.228835,0.228835,0.228835}%
\pgfsetfillcolor{currentfill}%
\pgfsetfillopacity{0.900000}%
\pgfsetlinewidth{0.501875pt}%
\definecolor{currentstroke}{rgb}{0.200000,0.200000,0.200000}%
\pgfsetstrokecolor{currentstroke}%
\pgfsetdash{}{0pt}%
\pgfpathmoveto{\pgfqpoint{1.928918in}{2.830966in}}%
\pgfpathlineto{\pgfqpoint{1.978335in}{2.783265in}}%
\pgfpathlineto{\pgfqpoint{2.011039in}{2.762477in}}%
\pgfpathlineto{\pgfqpoint{1.961614in}{2.810209in}}%
\pgfpathlineto{\pgfqpoint{1.928918in}{2.830966in}}%
\pgfpathclose%
\pgfusepath{stroke,fill}%
\end{pgfscope}%
\begin{pgfscope}%
\pgfpathrectangle{\pgfqpoint{0.050000in}{0.050000in}}{\pgfqpoint{4.900000in}{4.900000in}}%
\pgfusepath{clip}%
\pgfsetbuttcap%
\pgfsetroundjoin%
\definecolor{currentfill}{rgb}{0.504268,0.504268,0.504268}%
\pgfsetfillcolor{currentfill}%
\pgfsetfillopacity{0.900000}%
\pgfsetlinewidth{0.501875pt}%
\definecolor{currentstroke}{rgb}{0.200000,0.200000,0.200000}%
\pgfsetstrokecolor{currentstroke}%
\pgfsetdash{}{0pt}%
\pgfpathmoveto{\pgfqpoint{2.684263in}{2.220536in}}%
\pgfpathlineto{\pgfqpoint{2.732000in}{2.174205in}}%
\pgfpathlineto{\pgfqpoint{2.765935in}{2.152428in}}%
\pgfpathlineto{\pgfqpoint{2.718196in}{2.198784in}}%
\pgfpathlineto{\pgfqpoint{2.684263in}{2.220536in}}%
\pgfpathclose%
\pgfusepath{stroke,fill}%
\end{pgfscope}%
\begin{pgfscope}%
\pgfpathrectangle{\pgfqpoint{0.050000in}{0.050000in}}{\pgfqpoint{4.900000in}{4.900000in}}%
\pgfusepath{clip}%
\pgfsetbuttcap%
\pgfsetroundjoin%
\definecolor{currentfill}{rgb}{0.113802,0.113802,0.113802}%
\pgfsetfillcolor{currentfill}%
\pgfsetfillopacity{0.900000}%
\pgfsetlinewidth{0.501875pt}%
\definecolor{currentstroke}{rgb}{0.200000,0.200000,0.200000}%
\pgfsetstrokecolor{currentstroke}%
\pgfsetdash{}{0pt}%
\pgfpathmoveto{\pgfqpoint{1.649715in}{3.057318in}}%
\pgfpathlineto{\pgfqpoint{1.699763in}{3.009104in}}%
\pgfpathlineto{\pgfqpoint{1.732018in}{2.988681in}}%
\pgfpathlineto{\pgfqpoint{1.681960in}{3.036928in}}%
\pgfpathlineto{\pgfqpoint{1.649715in}{3.057318in}}%
\pgfpathclose%
\pgfusepath{stroke,fill}%
\end{pgfscope}%
\begin{pgfscope}%
\pgfpathrectangle{\pgfqpoint{0.050000in}{0.050000in}}{\pgfqpoint{4.900000in}{4.900000in}}%
\pgfusepath{clip}%
\pgfsetbuttcap%
\pgfsetroundjoin%
\definecolor{currentfill}{rgb}{0.407843,0.407843,0.407843}%
\pgfsetfillcolor{currentfill}%
\pgfsetfillopacity{0.900000}%
\pgfsetlinewidth{0.501875pt}%
\definecolor{currentstroke}{rgb}{0.200000,0.200000,0.200000}%
\pgfsetstrokecolor{currentstroke}%
\pgfsetdash{}{0pt}%
\pgfpathmoveto{\pgfqpoint{2.405174in}{2.446781in}}%
\pgfpathlineto{\pgfqpoint{2.453541in}{2.399926in}}%
\pgfpathlineto{\pgfqpoint{2.487028in}{2.378509in}}%
\pgfpathlineto{\pgfqpoint{2.438657in}{2.425393in}}%
\pgfpathlineto{\pgfqpoint{2.405174in}{2.446781in}}%
\pgfpathclose%
\pgfusepath{stroke,fill}%
\end{pgfscope}%
\begin{pgfscope}%
\pgfpathrectangle{\pgfqpoint{0.050000in}{0.050000in}}{\pgfqpoint{4.900000in}{4.900000in}}%
\pgfusepath{clip}%
\pgfsetbuttcap%
\pgfsetroundjoin%
\definecolor{currentfill}{rgb}{0.009104,0.009104,0.009104}%
\pgfsetfillcolor{currentfill}%
\pgfsetfillopacity{0.900000}%
\pgfsetlinewidth{0.501875pt}%
\definecolor{currentstroke}{rgb}{0.200000,0.200000,0.200000}%
\pgfsetstrokecolor{currentstroke}%
\pgfsetdash{}{0pt}%
\pgfpathmoveto{\pgfqpoint{1.370412in}{3.283754in}}%
\pgfpathlineto{\pgfqpoint{1.421092in}{3.235026in}}%
\pgfpathlineto{\pgfqpoint{1.452899in}{3.214967in}}%
\pgfpathlineto{\pgfqpoint{1.402207in}{3.263729in}}%
\pgfpathlineto{\pgfqpoint{1.370412in}{3.283754in}}%
\pgfpathclose%
\pgfusepath{stroke,fill}%
\end{pgfscope}%
\begin{pgfscope}%
\pgfpathrectangle{\pgfqpoint{0.050000in}{0.050000in}}{\pgfqpoint{4.900000in}{4.900000in}}%
\pgfusepath{clip}%
\pgfsetbuttcap%
\pgfsetroundjoin%
\definecolor{currentfill}{rgb}{0.311880,0.311880,0.311880}%
\pgfsetfillcolor{currentfill}%
\pgfsetfillopacity{0.900000}%
\pgfsetlinewidth{0.501875pt}%
\definecolor{currentstroke}{rgb}{0.200000,0.200000,0.200000}%
\pgfsetstrokecolor{currentstroke}%
\pgfsetdash{}{0pt}%
\pgfpathmoveto{\pgfqpoint{2.125986in}{2.673117in}}%
\pgfpathlineto{\pgfqpoint{2.174984in}{2.625749in}}%
\pgfpathlineto{\pgfqpoint{2.208023in}{2.604696in}}%
\pgfpathlineto{\pgfqpoint{2.159018in}{2.652095in}}%
\pgfpathlineto{\pgfqpoint{2.125986in}{2.673117in}}%
\pgfpathclose%
\pgfusepath{stroke,fill}%
\end{pgfscope}%
\begin{pgfscope}%
\pgfpathrectangle{\pgfqpoint{0.050000in}{0.050000in}}{\pgfqpoint{4.900000in}{4.900000in}}%
\pgfusepath{clip}%
\pgfsetbuttcap%
\pgfsetroundjoin%
\definecolor{currentfill}{rgb}{0.568858,0.568858,0.568858}%
\pgfsetfillcolor{currentfill}%
\pgfsetfillopacity{0.900000}%
\pgfsetlinewidth{0.501875pt}%
\definecolor{currentstroke}{rgb}{0.200000,0.200000,0.200000}%
\pgfsetstrokecolor{currentstroke}%
\pgfsetdash{}{0pt}%
\pgfpathmoveto{\pgfqpoint{2.881667in}{2.062486in}}%
\pgfpathlineto{\pgfqpoint{2.928984in}{2.016553in}}%
\pgfpathlineto{\pgfqpoint{2.963255in}{1.994526in}}%
\pgfpathlineto{\pgfqpoint{2.915937in}{2.040476in}}%
\pgfpathlineto{\pgfqpoint{2.881667in}{2.062486in}}%
\pgfpathclose%
\pgfusepath{stroke,fill}%
\end{pgfscope}%
\begin{pgfscope}%
\pgfpathrectangle{\pgfqpoint{0.050000in}{0.050000in}}{\pgfqpoint{4.900000in}{4.900000in}}%
\pgfusepath{clip}%
\pgfsetbuttcap%
\pgfsetroundjoin%
\definecolor{currentfill}{rgb}{0.190081,0.190081,0.190081}%
\pgfsetfillcolor{currentfill}%
\pgfsetfillopacity{0.900000}%
\pgfsetlinewidth{0.501875pt}%
\definecolor{currentstroke}{rgb}{0.200000,0.200000,0.200000}%
\pgfsetstrokecolor{currentstroke}%
\pgfsetdash{}{0pt}%
\pgfpathmoveto{\pgfqpoint{1.846698in}{2.899537in}}%
\pgfpathlineto{\pgfqpoint{1.896328in}{2.851654in}}%
\pgfpathlineto{\pgfqpoint{1.928918in}{2.830966in}}%
\pgfpathlineto{\pgfqpoint{1.879280in}{2.878880in}}%
\pgfpathlineto{\pgfqpoint{1.846698in}{2.899537in}}%
\pgfpathclose%
\pgfusepath{stroke,fill}%
\end{pgfscope}%
\begin{pgfscope}%
\pgfpathrectangle{\pgfqpoint{0.050000in}{0.050000in}}{\pgfqpoint{4.900000in}{4.900000in}}%
\pgfusepath{clip}%
\pgfsetbuttcap%
\pgfsetroundjoin%
\definecolor{currentfill}{rgb}{0.469819,0.469819,0.469819}%
\pgfsetfillcolor{currentfill}%
\pgfsetfillopacity{0.900000}%
\pgfsetlinewidth{0.501875pt}%
\definecolor{currentstroke}{rgb}{0.200000,0.200000,0.200000}%
\pgfsetstrokecolor{currentstroke}%
\pgfsetdash{}{0pt}%
\pgfpathmoveto{\pgfqpoint{2.602493in}{2.288736in}}%
\pgfpathlineto{\pgfqpoint{2.650441in}{2.242217in}}%
\pgfpathlineto{\pgfqpoint{2.684263in}{2.220536in}}%
\pgfpathlineto{\pgfqpoint{2.636313in}{2.267082in}}%
\pgfpathlineto{\pgfqpoint{2.602493in}{2.288736in}}%
\pgfpathclose%
\pgfusepath{stroke,fill}%
\end{pgfscope}%
\begin{pgfscope}%
\pgfpathrectangle{\pgfqpoint{0.050000in}{0.050000in}}{\pgfqpoint{4.900000in}{4.900000in}}%
\pgfusepath{clip}%
\pgfsetbuttcap%
\pgfsetroundjoin%
\definecolor{currentfill}{rgb}{0.077386,0.077386,0.077386}%
\pgfsetfillcolor{currentfill}%
\pgfsetfillopacity{0.900000}%
\pgfsetlinewidth{0.501875pt}%
\definecolor{currentstroke}{rgb}{0.200000,0.200000,0.200000}%
\pgfsetstrokecolor{currentstroke}%
\pgfsetdash{}{0pt}%
\pgfpathmoveto{\pgfqpoint{1.567312in}{3.126039in}}%
\pgfpathlineto{\pgfqpoint{1.617574in}{3.077643in}}%
\pgfpathlineto{\pgfqpoint{1.649715in}{3.057318in}}%
\pgfpathlineto{\pgfqpoint{1.599442in}{3.105748in}}%
\pgfpathlineto{\pgfqpoint{1.567312in}{3.126039in}}%
\pgfpathclose%
\pgfusepath{stroke,fill}%
\end{pgfscope}%
\begin{pgfscope}%
\pgfpathrectangle{\pgfqpoint{0.050000in}{0.050000in}}{\pgfqpoint{4.900000in}{4.900000in}}%
\pgfusepath{clip}%
\pgfsetbuttcap%
\pgfsetroundjoin%
\definecolor{currentfill}{rgb}{0.375363,0.375363,0.375363}%
\pgfsetfillcolor{currentfill}%
\pgfsetfillopacity{0.900000}%
\pgfsetlinewidth{0.501875pt}%
\definecolor{currentstroke}{rgb}{0.200000,0.200000,0.200000}%
\pgfsetstrokecolor{currentstroke}%
\pgfsetdash{}{0pt}%
\pgfpathmoveto{\pgfqpoint{2.323221in}{2.515136in}}%
\pgfpathlineto{\pgfqpoint{2.371800in}{2.468100in}}%
\pgfpathlineto{\pgfqpoint{2.405174in}{2.446781in}}%
\pgfpathlineto{\pgfqpoint{2.356590in}{2.493847in}}%
\pgfpathlineto{\pgfqpoint{2.323221in}{2.515136in}}%
\pgfpathclose%
\pgfusepath{stroke,fill}%
\end{pgfscope}%
\begin{pgfscope}%
\pgfpathrectangle{\pgfqpoint{0.050000in}{0.050000in}}{\pgfqpoint{4.900000in}{4.900000in}}%
\pgfusepath{clip}%
\pgfsetbuttcap%
\pgfsetroundjoin%
\definecolor{currentfill}{rgb}{0.273126,0.273126,0.273126}%
\pgfsetfillcolor{currentfill}%
\pgfsetfillopacity{0.900000}%
\pgfsetlinewidth{0.501875pt}%
\definecolor{currentstroke}{rgb}{0.200000,0.200000,0.200000}%
\pgfsetstrokecolor{currentstroke}%
\pgfsetdash{}{0pt}%
\pgfpathmoveto{\pgfqpoint{2.043849in}{2.741622in}}%
\pgfpathlineto{\pgfqpoint{2.093061in}{2.694072in}}%
\pgfpathlineto{\pgfqpoint{2.125986in}{2.673117in}}%
\pgfpathlineto{\pgfqpoint{2.076767in}{2.720699in}}%
\pgfpathlineto{\pgfqpoint{2.043849in}{2.741622in}}%
\pgfpathclose%
\pgfusepath{stroke,fill}%
\end{pgfscope}%
\begin{pgfscope}%
\pgfpathrectangle{\pgfqpoint{0.050000in}{0.050000in}}{\pgfqpoint{4.900000in}{4.900000in}}%
\pgfusepath{clip}%
\pgfsetbuttcap%
\pgfsetroundjoin%
\definecolor{currentfill}{rgb}{0.534410,0.534410,0.534410}%
\pgfsetfillcolor{currentfill}%
\pgfsetfillopacity{0.900000}%
\pgfsetlinewidth{0.501875pt}%
\definecolor{currentstroke}{rgb}{0.200000,0.200000,0.200000}%
\pgfsetstrokecolor{currentstroke}%
\pgfsetdash{}{0pt}%
\pgfpathmoveto{\pgfqpoint{2.799981in}{2.130581in}}%
\pgfpathlineto{\pgfqpoint{2.847508in}{2.084426in}}%
\pgfpathlineto{\pgfqpoint{2.881667in}{2.062486in}}%
\pgfpathlineto{\pgfqpoint{2.834138in}{2.108663in}}%
\pgfpathlineto{\pgfqpoint{2.799981in}{2.130581in}}%
\pgfpathclose%
\pgfusepath{stroke,fill}%
\end{pgfscope}%
\begin{pgfscope}%
\pgfpathrectangle{\pgfqpoint{0.050000in}{0.050000in}}{\pgfqpoint{4.900000in}{4.900000in}}%
\pgfusepath{clip}%
\pgfsetbuttcap%
\pgfsetroundjoin%
\definecolor{currentfill}{rgb}{0.145790,0.145790,0.145790}%
\pgfsetfillcolor{currentfill}%
\pgfsetfillopacity{0.900000}%
\pgfsetlinewidth{0.501875pt}%
\definecolor{currentstroke}{rgb}{0.200000,0.200000,0.200000}%
\pgfsetstrokecolor{currentstroke}%
\pgfsetdash{}{0pt}%
\pgfpathmoveto{\pgfqpoint{1.764379in}{2.968191in}}%
\pgfpathlineto{\pgfqpoint{1.814223in}{2.920126in}}%
\pgfpathlineto{\pgfqpoint{1.846698in}{2.899537in}}%
\pgfpathlineto{\pgfqpoint{1.796845in}{2.947634in}}%
\pgfpathlineto{\pgfqpoint{1.764379in}{2.968191in}}%
\pgfpathclose%
\pgfusepath{stroke,fill}%
\end{pgfscope}%
\begin{pgfscope}%
\pgfpathrectangle{\pgfqpoint{0.050000in}{0.050000in}}{\pgfqpoint{4.900000in}{4.900000in}}%
\pgfusepath{clip}%
\pgfsetbuttcap%
\pgfsetroundjoin%
\definecolor{currentfill}{rgb}{0.440323,0.440323,0.440323}%
\pgfsetfillcolor{currentfill}%
\pgfsetfillopacity{0.900000}%
\pgfsetlinewidth{0.501875pt}%
\definecolor{currentstroke}{rgb}{0.200000,0.200000,0.200000}%
\pgfsetstrokecolor{currentstroke}%
\pgfsetdash{}{0pt}%
\pgfpathmoveto{\pgfqpoint{2.520624in}{2.357023in}}%
\pgfpathlineto{\pgfqpoint{2.568784in}{2.310320in}}%
\pgfpathlineto{\pgfqpoint{2.602493in}{2.288736in}}%
\pgfpathlineto{\pgfqpoint{2.554331in}{2.335466in}}%
\pgfpathlineto{\pgfqpoint{2.520624in}{2.357023in}}%
\pgfpathclose%
\pgfusepath{stroke,fill}%
\end{pgfscope}%
\begin{pgfscope}%
\pgfpathrectangle{\pgfqpoint{0.050000in}{0.050000in}}{\pgfqpoint{4.900000in}{4.900000in}}%
\pgfusepath{clip}%
\pgfsetbuttcap%
\pgfsetroundjoin%
\definecolor{currentfill}{rgb}{0.045521,0.045521,0.045521}%
\pgfsetfillcolor{currentfill}%
\pgfsetfillopacity{0.900000}%
\pgfsetlinewidth{0.501875pt}%
\definecolor{currentstroke}{rgb}{0.200000,0.200000,0.200000}%
\pgfsetstrokecolor{currentstroke}%
\pgfsetdash{}{0pt}%
\pgfpathmoveto{\pgfqpoint{1.484809in}{3.194842in}}%
\pgfpathlineto{\pgfqpoint{1.535286in}{3.146263in}}%
\pgfpathlineto{\pgfqpoint{1.567312in}{3.126039in}}%
\pgfpathlineto{\pgfqpoint{1.516824in}{3.174652in}}%
\pgfpathlineto{\pgfqpoint{1.484809in}{3.194842in}}%
\pgfpathclose%
\pgfusepath{stroke,fill}%
\end{pgfscope}%
\begin{pgfscope}%
\pgfpathrectangle{\pgfqpoint{0.050000in}{0.050000in}}{\pgfqpoint{4.900000in}{4.900000in}}%
\pgfusepath{clip}%
\pgfsetbuttcap%
\pgfsetroundjoin%
\definecolor{currentfill}{rgb}{0.346943,0.346943,0.346943}%
\pgfsetfillcolor{currentfill}%
\pgfsetfillopacity{0.900000}%
\pgfsetlinewidth{0.501875pt}%
\definecolor{currentstroke}{rgb}{0.200000,0.200000,0.200000}%
\pgfsetstrokecolor{currentstroke}%
\pgfsetdash{}{0pt}%
\pgfpathmoveto{\pgfqpoint{2.241169in}{2.583574in}}%
\pgfpathlineto{\pgfqpoint{2.289961in}{2.536357in}}%
\pgfpathlineto{\pgfqpoint{2.323221in}{2.515136in}}%
\pgfpathlineto{\pgfqpoint{2.274424in}{2.562384in}}%
\pgfpathlineto{\pgfqpoint{2.241169in}{2.583574in}}%
\pgfpathclose%
\pgfusepath{stroke,fill}%
\end{pgfscope}%
\begin{pgfscope}%
\pgfpathrectangle{\pgfqpoint{0.050000in}{0.050000in}}{\pgfqpoint{4.900000in}{4.900000in}}%
\pgfusepath{clip}%
\pgfsetbuttcap%
\pgfsetroundjoin%
\definecolor{currentfill}{rgb}{0.600231,0.600231,0.600231}%
\pgfsetfillcolor{currentfill}%
\pgfsetfillopacity{0.900000}%
\pgfsetlinewidth{0.501875pt}%
\definecolor{currentstroke}{rgb}{0.200000,0.200000,0.200000}%
\pgfsetstrokecolor{currentstroke}%
\pgfsetdash{}{0pt}%
\pgfpathmoveto{\pgfqpoint{2.997638in}{1.972429in}}%
\pgfpathlineto{\pgfqpoint{3.044746in}{1.926749in}}%
\pgfpathlineto{\pgfqpoint{3.079242in}{1.904572in}}%
\pgfpathlineto{\pgfqpoint{3.032134in}{1.950261in}}%
\pgfpathlineto{\pgfqpoint{2.997638in}{1.972429in}}%
\pgfpathclose%
\pgfusepath{stroke,fill}%
\end{pgfscope}%
\begin{pgfscope}%
\pgfpathrectangle{\pgfqpoint{0.050000in}{0.050000in}}{\pgfqpoint{4.900000in}{4.900000in}}%
\pgfusepath{clip}%
\pgfsetbuttcap%
\pgfsetroundjoin%
\definecolor{currentfill}{rgb}{0.228835,0.228835,0.228835}%
\pgfsetfillcolor{currentfill}%
\pgfsetfillopacity{0.900000}%
\pgfsetlinewidth{0.501875pt}%
\definecolor{currentstroke}{rgb}{0.200000,0.200000,0.200000}%
\pgfsetstrokecolor{currentstroke}%
\pgfsetdash{}{0pt}%
\pgfpathmoveto{\pgfqpoint{1.961614in}{2.810209in}}%
\pgfpathlineto{\pgfqpoint{2.011039in}{2.762477in}}%
\pgfpathlineto{\pgfqpoint{2.043849in}{2.741622in}}%
\pgfpathlineto{\pgfqpoint{1.994417in}{2.789386in}}%
\pgfpathlineto{\pgfqpoint{1.961614in}{2.810209in}}%
\pgfpathclose%
\pgfusepath{stroke,fill}%
\end{pgfscope}%
\begin{pgfscope}%
\pgfpathrectangle{\pgfqpoint{0.050000in}{0.050000in}}{\pgfqpoint{4.900000in}{4.900000in}}%
\pgfusepath{clip}%
\pgfsetbuttcap%
\pgfsetroundjoin%
\definecolor{currentfill}{rgb}{0.504268,0.504268,0.504268}%
\pgfsetfillcolor{currentfill}%
\pgfsetfillopacity{0.900000}%
\pgfsetlinewidth{0.501875pt}%
\definecolor{currentstroke}{rgb}{0.200000,0.200000,0.200000}%
\pgfsetstrokecolor{currentstroke}%
\pgfsetdash{}{0pt}%
\pgfpathmoveto{\pgfqpoint{2.718196in}{2.198784in}}%
\pgfpathlineto{\pgfqpoint{2.765935in}{2.152428in}}%
\pgfpathlineto{\pgfqpoint{2.799981in}{2.130581in}}%
\pgfpathlineto{\pgfqpoint{2.752240in}{2.176962in}}%
\pgfpathlineto{\pgfqpoint{2.718196in}{2.198784in}}%
\pgfpathclose%
\pgfusepath{stroke,fill}%
\end{pgfscope}%
\begin{pgfscope}%
\pgfpathrectangle{\pgfqpoint{0.050000in}{0.050000in}}{\pgfqpoint{4.900000in}{4.900000in}}%
\pgfusepath{clip}%
\pgfsetbuttcap%
\pgfsetroundjoin%
\definecolor{currentfill}{rgb}{0.113802,0.113802,0.113802}%
\pgfsetfillcolor{currentfill}%
\pgfsetfillopacity{0.900000}%
\pgfsetlinewidth{0.501875pt}%
\definecolor{currentstroke}{rgb}{0.200000,0.200000,0.200000}%
\pgfsetstrokecolor{currentstroke}%
\pgfsetdash{}{0pt}%
\pgfpathmoveto{\pgfqpoint{1.681960in}{3.036928in}}%
\pgfpathlineto{\pgfqpoint{1.732018in}{2.988681in}}%
\pgfpathlineto{\pgfqpoint{1.764379in}{2.968191in}}%
\pgfpathlineto{\pgfqpoint{1.714311in}{3.016471in}}%
\pgfpathlineto{\pgfqpoint{1.681960in}{3.036928in}}%
\pgfpathclose%
\pgfusepath{stroke,fill}%
\end{pgfscope}%
\begin{pgfscope}%
\pgfpathrectangle{\pgfqpoint{0.050000in}{0.050000in}}{\pgfqpoint{4.900000in}{4.900000in}}%
\pgfusepath{clip}%
\pgfsetbuttcap%
\pgfsetroundjoin%
\definecolor{currentfill}{rgb}{0.407843,0.407843,0.407843}%
\pgfsetfillcolor{currentfill}%
\pgfsetfillopacity{0.900000}%
\pgfsetlinewidth{0.501875pt}%
\definecolor{currentstroke}{rgb}{0.200000,0.200000,0.200000}%
\pgfsetstrokecolor{currentstroke}%
\pgfsetdash{}{0pt}%
\pgfpathmoveto{\pgfqpoint{2.438657in}{2.425393in}}%
\pgfpathlineto{\pgfqpoint{2.487028in}{2.378509in}}%
\pgfpathlineto{\pgfqpoint{2.520624in}{2.357023in}}%
\pgfpathlineto{\pgfqpoint{2.472249in}{2.403935in}}%
\pgfpathlineto{\pgfqpoint{2.438657in}{2.425393in}}%
\pgfpathclose%
\pgfusepath{stroke,fill}%
\end{pgfscope}%
\begin{pgfscope}%
\pgfpathrectangle{\pgfqpoint{0.050000in}{0.050000in}}{\pgfqpoint{4.900000in}{4.900000in}}%
\pgfusepath{clip}%
\pgfsetbuttcap%
\pgfsetroundjoin%
\definecolor{currentfill}{rgb}{0.009104,0.009104,0.009104}%
\pgfsetfillcolor{currentfill}%
\pgfsetfillopacity{0.900000}%
\pgfsetlinewidth{0.501875pt}%
\definecolor{currentstroke}{rgb}{0.200000,0.200000,0.200000}%
\pgfsetstrokecolor{currentstroke}%
\pgfsetdash{}{0pt}%
\pgfpathmoveto{\pgfqpoint{1.402207in}{3.263729in}}%
\pgfpathlineto{\pgfqpoint{1.452899in}{3.214967in}}%
\pgfpathlineto{\pgfqpoint{1.484809in}{3.194842in}}%
\pgfpathlineto{\pgfqpoint{1.434106in}{3.243639in}}%
\pgfpathlineto{\pgfqpoint{1.402207in}{3.263729in}}%
\pgfpathclose%
\pgfusepath{stroke,fill}%
\end{pgfscope}%
\begin{pgfscope}%
\pgfpathrectangle{\pgfqpoint{0.050000in}{0.050000in}}{\pgfqpoint{4.900000in}{4.900000in}}%
\pgfusepath{clip}%
\pgfsetbuttcap%
\pgfsetroundjoin%
\definecolor{currentfill}{rgb}{0.311880,0.311880,0.311880}%
\pgfsetfillcolor{currentfill}%
\pgfsetfillopacity{0.900000}%
\pgfsetlinewidth{0.501875pt}%
\definecolor{currentstroke}{rgb}{0.200000,0.200000,0.200000}%
\pgfsetstrokecolor{currentstroke}%
\pgfsetdash{}{0pt}%
\pgfpathmoveto{\pgfqpoint{2.159018in}{2.652095in}}%
\pgfpathlineto{\pgfqpoint{2.208023in}{2.604696in}}%
\pgfpathlineto{\pgfqpoint{2.241169in}{2.583574in}}%
\pgfpathlineto{\pgfqpoint{2.192158in}{2.631004in}}%
\pgfpathlineto{\pgfqpoint{2.159018in}{2.652095in}}%
\pgfpathclose%
\pgfusepath{stroke,fill}%
\end{pgfscope}%
\begin{pgfscope}%
\pgfpathrectangle{\pgfqpoint{0.050000in}{0.050000in}}{\pgfqpoint{4.900000in}{4.900000in}}%
\pgfusepath{clip}%
\pgfsetbuttcap%
\pgfsetroundjoin%
\definecolor{currentfill}{rgb}{0.568858,0.568858,0.568858}%
\pgfsetfillcolor{currentfill}%
\pgfsetfillopacity{0.900000}%
\pgfsetlinewidth{0.501875pt}%
\definecolor{currentstroke}{rgb}{0.200000,0.200000,0.200000}%
\pgfsetstrokecolor{currentstroke}%
\pgfsetdash{}{0pt}%
\pgfpathmoveto{\pgfqpoint{2.915937in}{2.040476in}}%
\pgfpathlineto{\pgfqpoint{2.963255in}{1.994526in}}%
\pgfpathlineto{\pgfqpoint{2.997638in}{1.972429in}}%
\pgfpathlineto{\pgfqpoint{2.950320in}{2.018395in}}%
\pgfpathlineto{\pgfqpoint{2.915937in}{2.040476in}}%
\pgfpathclose%
\pgfusepath{stroke,fill}%
\end{pgfscope}%
\begin{pgfscope}%
\pgfpathrectangle{\pgfqpoint{0.050000in}{0.050000in}}{\pgfqpoint{4.900000in}{4.900000in}}%
\pgfusepath{clip}%
\pgfsetbuttcap%
\pgfsetroundjoin%
\definecolor{currentfill}{rgb}{0.190081,0.190081,0.190081}%
\pgfsetfillcolor{currentfill}%
\pgfsetfillopacity{0.900000}%
\pgfsetlinewidth{0.501875pt}%
\definecolor{currentstroke}{rgb}{0.200000,0.200000,0.200000}%
\pgfsetstrokecolor{currentstroke}%
\pgfsetdash{}{0pt}%
\pgfpathmoveto{\pgfqpoint{1.879280in}{2.878880in}}%
\pgfpathlineto{\pgfqpoint{1.928918in}{2.830966in}}%
\pgfpathlineto{\pgfqpoint{1.961614in}{2.810209in}}%
\pgfpathlineto{\pgfqpoint{1.911968in}{2.858156in}}%
\pgfpathlineto{\pgfqpoint{1.879280in}{2.878880in}}%
\pgfpathclose%
\pgfusepath{stroke,fill}%
\end{pgfscope}%
\begin{pgfscope}%
\pgfpathrectangle{\pgfqpoint{0.050000in}{0.050000in}}{\pgfqpoint{4.900000in}{4.900000in}}%
\pgfusepath{clip}%
\pgfsetbuttcap%
\pgfsetroundjoin%
\definecolor{currentfill}{rgb}{0.469819,0.469819,0.469819}%
\pgfsetfillcolor{currentfill}%
\pgfsetfillopacity{0.900000}%
\pgfsetlinewidth{0.501875pt}%
\definecolor{currentstroke}{rgb}{0.200000,0.200000,0.200000}%
\pgfsetstrokecolor{currentstroke}%
\pgfsetdash{}{0pt}%
\pgfpathmoveto{\pgfqpoint{2.636313in}{2.267082in}}%
\pgfpathlineto{\pgfqpoint{2.684263in}{2.220536in}}%
\pgfpathlineto{\pgfqpoint{2.718196in}{2.198784in}}%
\pgfpathlineto{\pgfqpoint{2.670243in}{2.245356in}}%
\pgfpathlineto{\pgfqpoint{2.636313in}{2.267082in}}%
\pgfpathclose%
\pgfusepath{stroke,fill}%
\end{pgfscope}%
\begin{pgfscope}%
\pgfpathrectangle{\pgfqpoint{0.050000in}{0.050000in}}{\pgfqpoint{4.900000in}{4.900000in}}%
\pgfusepath{clip}%
\pgfsetbuttcap%
\pgfsetroundjoin%
\definecolor{currentfill}{rgb}{0.077386,0.077386,0.077386}%
\pgfsetfillcolor{currentfill}%
\pgfsetfillopacity{0.900000}%
\pgfsetlinewidth{0.501875pt}%
\definecolor{currentstroke}{rgb}{0.200000,0.200000,0.200000}%
\pgfsetstrokecolor{currentstroke}%
\pgfsetdash{}{0pt}%
\pgfpathmoveto{\pgfqpoint{1.599442in}{3.105748in}}%
\pgfpathlineto{\pgfqpoint{1.649715in}{3.057318in}}%
\pgfpathlineto{\pgfqpoint{1.681960in}{3.036928in}}%
\pgfpathlineto{\pgfqpoint{1.631677in}{3.085391in}}%
\pgfpathlineto{\pgfqpoint{1.599442in}{3.105748in}}%
\pgfpathclose%
\pgfusepath{stroke,fill}%
\end{pgfscope}%
\begin{pgfscope}%
\pgfpathrectangle{\pgfqpoint{0.050000in}{0.050000in}}{\pgfqpoint{4.900000in}{4.900000in}}%
\pgfusepath{clip}%
\pgfsetbuttcap%
\pgfsetroundjoin%
\definecolor{currentfill}{rgb}{0.379423,0.379423,0.379423}%
\pgfsetfillcolor{currentfill}%
\pgfsetfillopacity{0.900000}%
\pgfsetlinewidth{0.501875pt}%
\definecolor{currentstroke}{rgb}{0.200000,0.200000,0.200000}%
\pgfsetstrokecolor{currentstroke}%
\pgfsetdash{}{0pt}%
\pgfpathmoveto{\pgfqpoint{2.356590in}{2.493847in}}%
\pgfpathlineto{\pgfqpoint{2.405174in}{2.446781in}}%
\pgfpathlineto{\pgfqpoint{2.438657in}{2.425393in}}%
\pgfpathlineto{\pgfqpoint{2.390068in}{2.472488in}}%
\pgfpathlineto{\pgfqpoint{2.356590in}{2.493847in}}%
\pgfpathclose%
\pgfusepath{stroke,fill}%
\end{pgfscope}%
\begin{pgfscope}%
\pgfpathrectangle{\pgfqpoint{0.050000in}{0.050000in}}{\pgfqpoint{4.900000in}{4.900000in}}%
\pgfusepath{clip}%
\pgfsetbuttcap%
\pgfsetroundjoin%
\definecolor{currentfill}{rgb}{0.273126,0.273126,0.273126}%
\pgfsetfillcolor{currentfill}%
\pgfsetfillopacity{0.900000}%
\pgfsetlinewidth{0.501875pt}%
\definecolor{currentstroke}{rgb}{0.200000,0.200000,0.200000}%
\pgfsetstrokecolor{currentstroke}%
\pgfsetdash{}{0pt}%
\pgfpathmoveto{\pgfqpoint{2.076767in}{2.720699in}}%
\pgfpathlineto{\pgfqpoint{2.125986in}{2.673117in}}%
\pgfpathlineto{\pgfqpoint{2.159018in}{2.652095in}}%
\pgfpathlineto{\pgfqpoint{2.109793in}{2.699707in}}%
\pgfpathlineto{\pgfqpoint{2.076767in}{2.720699in}}%
\pgfpathclose%
\pgfusepath{stroke,fill}%
\end{pgfscope}%
\begin{pgfscope}%
\pgfpathrectangle{\pgfqpoint{0.050000in}{0.050000in}}{\pgfqpoint{4.900000in}{4.900000in}}%
\pgfusepath{clip}%
\pgfsetbuttcap%
\pgfsetroundjoin%
\definecolor{currentfill}{rgb}{0.534410,0.534410,0.534410}%
\pgfsetfillcolor{currentfill}%
\pgfsetfillopacity{0.900000}%
\pgfsetlinewidth{0.501875pt}%
\definecolor{currentstroke}{rgb}{0.200000,0.200000,0.200000}%
\pgfsetstrokecolor{currentstroke}%
\pgfsetdash{}{0pt}%
\pgfpathmoveto{\pgfqpoint{2.834138in}{2.108663in}}%
\pgfpathlineto{\pgfqpoint{2.881667in}{2.062486in}}%
\pgfpathlineto{\pgfqpoint{2.915937in}{2.040476in}}%
\pgfpathlineto{\pgfqpoint{2.868407in}{2.086674in}}%
\pgfpathlineto{\pgfqpoint{2.834138in}{2.108663in}}%
\pgfpathclose%
\pgfusepath{stroke,fill}%
\end{pgfscope}%
\begin{pgfscope}%
\pgfpathrectangle{\pgfqpoint{0.050000in}{0.050000in}}{\pgfqpoint{4.900000in}{4.900000in}}%
\pgfusepath{clip}%
\pgfsetbuttcap%
\pgfsetroundjoin%
\definecolor{currentfill}{rgb}{0.145790,0.145790,0.145790}%
\pgfsetfillcolor{currentfill}%
\pgfsetfillopacity{0.900000}%
\pgfsetlinewidth{0.501875pt}%
\definecolor{currentstroke}{rgb}{0.200000,0.200000,0.200000}%
\pgfsetstrokecolor{currentstroke}%
\pgfsetdash{}{0pt}%
\pgfpathmoveto{\pgfqpoint{1.796845in}{2.947634in}}%
\pgfpathlineto{\pgfqpoint{1.846698in}{2.899537in}}%
\pgfpathlineto{\pgfqpoint{1.879280in}{2.878880in}}%
\pgfpathlineto{\pgfqpoint{1.829418in}{2.927009in}}%
\pgfpathlineto{\pgfqpoint{1.796845in}{2.947634in}}%
\pgfpathclose%
\pgfusepath{stroke,fill}%
\end{pgfscope}%
\begin{pgfscope}%
\pgfpathrectangle{\pgfqpoint{0.050000in}{0.050000in}}{\pgfqpoint{4.900000in}{4.900000in}}%
\pgfusepath{clip}%
\pgfsetbuttcap%
\pgfsetroundjoin%
\definecolor{currentfill}{rgb}{0.440323,0.440323,0.440323}%
\pgfsetfillcolor{currentfill}%
\pgfsetfillopacity{0.900000}%
\pgfsetlinewidth{0.501875pt}%
\definecolor{currentstroke}{rgb}{0.200000,0.200000,0.200000}%
\pgfsetstrokecolor{currentstroke}%
\pgfsetdash{}{0pt}%
\pgfpathmoveto{\pgfqpoint{2.554331in}{2.335466in}}%
\pgfpathlineto{\pgfqpoint{2.602493in}{2.288736in}}%
\pgfpathlineto{\pgfqpoint{2.636313in}{2.267082in}}%
\pgfpathlineto{\pgfqpoint{2.588147in}{2.313839in}}%
\pgfpathlineto{\pgfqpoint{2.554331in}{2.335466in}}%
\pgfpathclose%
\pgfusepath{stroke,fill}%
\end{pgfscope}%
\begin{pgfscope}%
\pgfpathrectangle{\pgfqpoint{0.050000in}{0.050000in}}{\pgfqpoint{4.900000in}{4.900000in}}%
\pgfusepath{clip}%
\pgfsetbuttcap%
\pgfsetroundjoin%
\definecolor{currentfill}{rgb}{0.045521,0.045521,0.045521}%
\pgfsetfillcolor{currentfill}%
\pgfsetfillopacity{0.900000}%
\pgfsetlinewidth{0.501875pt}%
\definecolor{currentstroke}{rgb}{0.200000,0.200000,0.200000}%
\pgfsetstrokecolor{currentstroke}%
\pgfsetdash{}{0pt}%
\pgfpathmoveto{\pgfqpoint{1.516824in}{3.174652in}}%
\pgfpathlineto{\pgfqpoint{1.567312in}{3.126039in}}%
\pgfpathlineto{\pgfqpoint{1.599442in}{3.105748in}}%
\pgfpathlineto{\pgfqpoint{1.548943in}{3.154395in}}%
\pgfpathlineto{\pgfqpoint{1.516824in}{3.174652in}}%
\pgfpathclose%
\pgfusepath{stroke,fill}%
\end{pgfscope}%
\begin{pgfscope}%
\pgfpathrectangle{\pgfqpoint{0.050000in}{0.050000in}}{\pgfqpoint{4.900000in}{4.900000in}}%
\pgfusepath{clip}%
\pgfsetbuttcap%
\pgfsetroundjoin%
\definecolor{currentfill}{rgb}{0.346943,0.346943,0.346943}%
\pgfsetfillcolor{currentfill}%
\pgfsetfillopacity{0.900000}%
\pgfsetlinewidth{0.501875pt}%
\definecolor{currentstroke}{rgb}{0.200000,0.200000,0.200000}%
\pgfsetstrokecolor{currentstroke}%
\pgfsetdash{}{0pt}%
\pgfpathmoveto{\pgfqpoint{2.274424in}{2.562384in}}%
\pgfpathlineto{\pgfqpoint{2.323221in}{2.515136in}}%
\pgfpathlineto{\pgfqpoint{2.356590in}{2.493847in}}%
\pgfpathlineto{\pgfqpoint{2.307788in}{2.541124in}}%
\pgfpathlineto{\pgfqpoint{2.274424in}{2.562384in}}%
\pgfpathclose%
\pgfusepath{stroke,fill}%
\end{pgfscope}%
\begin{pgfscope}%
\pgfpathrectangle{\pgfqpoint{0.050000in}{0.050000in}}{\pgfqpoint{4.900000in}{4.900000in}}%
\pgfusepath{clip}%
\pgfsetbuttcap%
\pgfsetroundjoin%
\definecolor{currentfill}{rgb}{0.228835,0.228835,0.228835}%
\pgfsetfillcolor{currentfill}%
\pgfsetfillopacity{0.900000}%
\pgfsetlinewidth{0.501875pt}%
\definecolor{currentstroke}{rgb}{0.200000,0.200000,0.200000}%
\pgfsetstrokecolor{currentstroke}%
\pgfsetdash{}{0pt}%
\pgfpathmoveto{\pgfqpoint{1.994417in}{2.789386in}}%
\pgfpathlineto{\pgfqpoint{2.043849in}{2.741622in}}%
\pgfpathlineto{\pgfqpoint{2.076767in}{2.720699in}}%
\pgfpathlineto{\pgfqpoint{2.027328in}{2.768493in}}%
\pgfpathlineto{\pgfqpoint{1.994417in}{2.789386in}}%
\pgfpathclose%
\pgfusepath{stroke,fill}%
\end{pgfscope}%
\begin{pgfscope}%
\pgfpathrectangle{\pgfqpoint{0.050000in}{0.050000in}}{\pgfqpoint{4.900000in}{4.900000in}}%
\pgfusepath{clip}%
\pgfsetbuttcap%
\pgfsetroundjoin%
\definecolor{currentfill}{rgb}{0.504268,0.504268,0.504268}%
\pgfsetfillcolor{currentfill}%
\pgfsetfillopacity{0.900000}%
\pgfsetlinewidth{0.501875pt}%
\definecolor{currentstroke}{rgb}{0.200000,0.200000,0.200000}%
\pgfsetstrokecolor{currentstroke}%
\pgfsetdash{}{0pt}%
\pgfpathmoveto{\pgfqpoint{2.752240in}{2.176962in}}%
\pgfpathlineto{\pgfqpoint{2.799981in}{2.130581in}}%
\pgfpathlineto{\pgfqpoint{2.834138in}{2.108663in}}%
\pgfpathlineto{\pgfqpoint{2.786396in}{2.155068in}}%
\pgfpathlineto{\pgfqpoint{2.752240in}{2.176962in}}%
\pgfpathclose%
\pgfusepath{stroke,fill}%
\end{pgfscope}%
\begin{pgfscope}%
\pgfpathrectangle{\pgfqpoint{0.050000in}{0.050000in}}{\pgfqpoint{4.900000in}{4.900000in}}%
\pgfusepath{clip}%
\pgfsetbuttcap%
\pgfsetroundjoin%
\definecolor{currentfill}{rgb}{0.113802,0.113802,0.113802}%
\pgfsetfillcolor{currentfill}%
\pgfsetfillopacity{0.900000}%
\pgfsetlinewidth{0.501875pt}%
\definecolor{currentstroke}{rgb}{0.200000,0.200000,0.200000}%
\pgfsetstrokecolor{currentstroke}%
\pgfsetdash{}{0pt}%
\pgfpathmoveto{\pgfqpoint{1.714311in}{3.016471in}}%
\pgfpathlineto{\pgfqpoint{1.764379in}{2.968191in}}%
\pgfpathlineto{\pgfqpoint{1.796845in}{2.947634in}}%
\pgfpathlineto{\pgfqpoint{1.746769in}{2.995947in}}%
\pgfpathlineto{\pgfqpoint{1.714311in}{3.016471in}}%
\pgfpathclose%
\pgfusepath{stroke,fill}%
\end{pgfscope}%
\begin{pgfscope}%
\pgfpathrectangle{\pgfqpoint{0.050000in}{0.050000in}}{\pgfqpoint{4.900000in}{4.900000in}}%
\pgfusepath{clip}%
\pgfsetbuttcap%
\pgfsetroundjoin%
\definecolor{currentfill}{rgb}{0.407843,0.407843,0.407843}%
\pgfsetfillcolor{currentfill}%
\pgfsetfillopacity{0.900000}%
\pgfsetlinewidth{0.501875pt}%
\definecolor{currentstroke}{rgb}{0.200000,0.200000,0.200000}%
\pgfsetstrokecolor{currentstroke}%
\pgfsetdash{}{0pt}%
\pgfpathmoveto{\pgfqpoint{2.472249in}{2.403935in}}%
\pgfpathlineto{\pgfqpoint{2.520624in}{2.357023in}}%
\pgfpathlineto{\pgfqpoint{2.554331in}{2.335466in}}%
\pgfpathlineto{\pgfqpoint{2.505952in}{2.382407in}}%
\pgfpathlineto{\pgfqpoint{2.472249in}{2.403935in}}%
\pgfpathclose%
\pgfusepath{stroke,fill}%
\end{pgfscope}%
\begin{pgfscope}%
\pgfpathrectangle{\pgfqpoint{0.050000in}{0.050000in}}{\pgfqpoint{4.900000in}{4.900000in}}%
\pgfusepath{clip}%
\pgfsetbuttcap%
\pgfsetroundjoin%
\definecolor{currentfill}{rgb}{0.009104,0.009104,0.009104}%
\pgfsetfillcolor{currentfill}%
\pgfsetfillopacity{0.900000}%
\pgfsetlinewidth{0.501875pt}%
\definecolor{currentstroke}{rgb}{0.200000,0.200000,0.200000}%
\pgfsetstrokecolor{currentstroke}%
\pgfsetdash{}{0pt}%
\pgfpathmoveto{\pgfqpoint{1.434106in}{3.243639in}}%
\pgfpathlineto{\pgfqpoint{1.484809in}{3.194842in}}%
\pgfpathlineto{\pgfqpoint{1.516824in}{3.174652in}}%
\pgfpathlineto{\pgfqpoint{1.466109in}{3.223483in}}%
\pgfpathlineto{\pgfqpoint{1.434106in}{3.243639in}}%
\pgfpathclose%
\pgfusepath{stroke,fill}%
\end{pgfscope}%
\begin{pgfscope}%
\pgfpathrectangle{\pgfqpoint{0.050000in}{0.050000in}}{\pgfqpoint{4.900000in}{4.900000in}}%
\pgfusepath{clip}%
\pgfsetbuttcap%
\pgfsetroundjoin%
\definecolor{currentfill}{rgb}{0.317416,0.317416,0.317416}%
\pgfsetfillcolor{currentfill}%
\pgfsetfillopacity{0.900000}%
\pgfsetlinewidth{0.501875pt}%
\definecolor{currentstroke}{rgb}{0.200000,0.200000,0.200000}%
\pgfsetstrokecolor{currentstroke}%
\pgfsetdash{}{0pt}%
\pgfpathmoveto{\pgfqpoint{2.192158in}{2.631004in}}%
\pgfpathlineto{\pgfqpoint{2.241169in}{2.583574in}}%
\pgfpathlineto{\pgfqpoint{2.274424in}{2.562384in}}%
\pgfpathlineto{\pgfqpoint{2.225407in}{2.609843in}}%
\pgfpathlineto{\pgfqpoint{2.192158in}{2.631004in}}%
\pgfpathclose%
\pgfusepath{stroke,fill}%
\end{pgfscope}%
\begin{pgfscope}%
\pgfpathrectangle{\pgfqpoint{0.050000in}{0.050000in}}{\pgfqpoint{4.900000in}{4.900000in}}%
\pgfusepath{clip}%
\pgfsetbuttcap%
\pgfsetroundjoin%
\definecolor{currentfill}{rgb}{0.568858,0.568858,0.568858}%
\pgfsetfillcolor{currentfill}%
\pgfsetfillopacity{0.900000}%
\pgfsetlinewidth{0.501875pt}%
\definecolor{currentstroke}{rgb}{0.200000,0.200000,0.200000}%
\pgfsetstrokecolor{currentstroke}%
\pgfsetdash{}{0pt}%
\pgfpathmoveto{\pgfqpoint{2.950320in}{2.018395in}}%
\pgfpathlineto{\pgfqpoint{2.997638in}{1.972429in}}%
\pgfpathlineto{\pgfqpoint{3.032134in}{1.950261in}}%
\pgfpathlineto{\pgfqpoint{2.984815in}{1.996242in}}%
\pgfpathlineto{\pgfqpoint{2.950320in}{2.018395in}}%
\pgfpathclose%
\pgfusepath{stroke,fill}%
\end{pgfscope}%
\begin{pgfscope}%
\pgfpathrectangle{\pgfqpoint{0.050000in}{0.050000in}}{\pgfqpoint{4.900000in}{4.900000in}}%
\pgfusepath{clip}%
\pgfsetbuttcap%
\pgfsetroundjoin%
\definecolor{currentfill}{rgb}{0.190081,0.190081,0.190081}%
\pgfsetfillcolor{currentfill}%
\pgfsetfillopacity{0.900000}%
\pgfsetlinewidth{0.501875pt}%
\definecolor{currentstroke}{rgb}{0.200000,0.200000,0.200000}%
\pgfsetstrokecolor{currentstroke}%
\pgfsetdash{}{0pt}%
\pgfpathmoveto{\pgfqpoint{1.911968in}{2.858156in}}%
\pgfpathlineto{\pgfqpoint{1.961614in}{2.810209in}}%
\pgfpathlineto{\pgfqpoint{1.994417in}{2.789386in}}%
\pgfpathlineto{\pgfqpoint{1.944763in}{2.837364in}}%
\pgfpathlineto{\pgfqpoint{1.911968in}{2.858156in}}%
\pgfpathclose%
\pgfusepath{stroke,fill}%
\end{pgfscope}%
\begin{pgfscope}%
\pgfpathrectangle{\pgfqpoint{0.050000in}{0.050000in}}{\pgfqpoint{4.900000in}{4.900000in}}%
\pgfusepath{clip}%
\pgfsetbuttcap%
\pgfsetroundjoin%
\definecolor{currentfill}{rgb}{0.469819,0.469819,0.469819}%
\pgfsetfillcolor{currentfill}%
\pgfsetfillopacity{0.900000}%
\pgfsetlinewidth{0.501875pt}%
\definecolor{currentstroke}{rgb}{0.200000,0.200000,0.200000}%
\pgfsetstrokecolor{currentstroke}%
\pgfsetdash{}{0pt}%
\pgfpathmoveto{\pgfqpoint{2.670243in}{2.245356in}}%
\pgfpathlineto{\pgfqpoint{2.718196in}{2.198784in}}%
\pgfpathlineto{\pgfqpoint{2.752240in}{2.176962in}}%
\pgfpathlineto{\pgfqpoint{2.704285in}{2.223560in}}%
\pgfpathlineto{\pgfqpoint{2.670243in}{2.245356in}}%
\pgfpathclose%
\pgfusepath{stroke,fill}%
\end{pgfscope}%
\begin{pgfscope}%
\pgfpathrectangle{\pgfqpoint{0.050000in}{0.050000in}}{\pgfqpoint{4.900000in}{4.900000in}}%
\pgfusepath{clip}%
\pgfsetbuttcap%
\pgfsetroundjoin%
\definecolor{currentfill}{rgb}{0.077386,0.077386,0.077386}%
\pgfsetfillcolor{currentfill}%
\pgfsetfillopacity{0.900000}%
\pgfsetlinewidth{0.501875pt}%
\definecolor{currentstroke}{rgb}{0.200000,0.200000,0.200000}%
\pgfsetstrokecolor{currentstroke}%
\pgfsetdash{}{0pt}%
\pgfpathmoveto{\pgfqpoint{1.631677in}{3.085391in}}%
\pgfpathlineto{\pgfqpoint{1.681960in}{3.036928in}}%
\pgfpathlineto{\pgfqpoint{1.714311in}{3.016471in}}%
\pgfpathlineto{\pgfqpoint{1.664019in}{3.064968in}}%
\pgfpathlineto{\pgfqpoint{1.631677in}{3.085391in}}%
\pgfpathclose%
\pgfusepath{stroke,fill}%
\end{pgfscope}%
\begin{pgfscope}%
\pgfpathrectangle{\pgfqpoint{0.050000in}{0.050000in}}{\pgfqpoint{4.900000in}{4.900000in}}%
\pgfusepath{clip}%
\pgfsetbuttcap%
\pgfsetroundjoin%
\definecolor{currentfill}{rgb}{0.379423,0.379423,0.379423}%
\pgfsetfillcolor{currentfill}%
\pgfsetfillopacity{0.900000}%
\pgfsetlinewidth{0.501875pt}%
\definecolor{currentstroke}{rgb}{0.200000,0.200000,0.200000}%
\pgfsetstrokecolor{currentstroke}%
\pgfsetdash{}{0pt}%
\pgfpathmoveto{\pgfqpoint{2.390068in}{2.472488in}}%
\pgfpathlineto{\pgfqpoint{2.438657in}{2.425393in}}%
\pgfpathlineto{\pgfqpoint{2.472249in}{2.403935in}}%
\pgfpathlineto{\pgfqpoint{2.423656in}{2.451059in}}%
\pgfpathlineto{\pgfqpoint{2.390068in}{2.472488in}}%
\pgfpathclose%
\pgfusepath{stroke,fill}%
\end{pgfscope}%
\begin{pgfscope}%
\pgfpathrectangle{\pgfqpoint{0.050000in}{0.050000in}}{\pgfqpoint{4.900000in}{4.900000in}}%
\pgfusepath{clip}%
\pgfsetbuttcap%
\pgfsetroundjoin%
\definecolor{currentfill}{rgb}{0.273126,0.273126,0.273126}%
\pgfsetfillcolor{currentfill}%
\pgfsetfillopacity{0.900000}%
\pgfsetlinewidth{0.501875pt}%
\definecolor{currentstroke}{rgb}{0.200000,0.200000,0.200000}%
\pgfsetstrokecolor{currentstroke}%
\pgfsetdash{}{0pt}%
\pgfpathmoveto{\pgfqpoint{2.109793in}{2.699707in}}%
\pgfpathlineto{\pgfqpoint{2.159018in}{2.652095in}}%
\pgfpathlineto{\pgfqpoint{2.192158in}{2.631004in}}%
\pgfpathlineto{\pgfqpoint{2.142927in}{2.678646in}}%
\pgfpathlineto{\pgfqpoint{2.109793in}{2.699707in}}%
\pgfpathclose%
\pgfusepath{stroke,fill}%
\end{pgfscope}%
\begin{pgfscope}%
\pgfpathrectangle{\pgfqpoint{0.050000in}{0.050000in}}{\pgfqpoint{4.900000in}{4.900000in}}%
\pgfusepath{clip}%
\pgfsetbuttcap%
\pgfsetroundjoin%
\definecolor{currentfill}{rgb}{0.534410,0.534410,0.534410}%
\pgfsetfillcolor{currentfill}%
\pgfsetfillopacity{0.900000}%
\pgfsetlinewidth{0.501875pt}%
\definecolor{currentstroke}{rgb}{0.200000,0.200000,0.200000}%
\pgfsetstrokecolor{currentstroke}%
\pgfsetdash{}{0pt}%
\pgfpathmoveto{\pgfqpoint{2.868407in}{2.086674in}}%
\pgfpathlineto{\pgfqpoint{2.915937in}{2.040476in}}%
\pgfpathlineto{\pgfqpoint{2.950320in}{2.018395in}}%
\pgfpathlineto{\pgfqpoint{2.902789in}{2.064613in}}%
\pgfpathlineto{\pgfqpoint{2.868407in}{2.086674in}}%
\pgfpathclose%
\pgfusepath{stroke,fill}%
\end{pgfscope}%
\begin{pgfscope}%
\pgfpathrectangle{\pgfqpoint{0.050000in}{0.050000in}}{\pgfqpoint{4.900000in}{4.900000in}}%
\pgfusepath{clip}%
\pgfsetbuttcap%
\pgfsetroundjoin%
\definecolor{currentfill}{rgb}{0.145790,0.145790,0.145790}%
\pgfsetfillcolor{currentfill}%
\pgfsetfillopacity{0.900000}%
\pgfsetlinewidth{0.501875pt}%
\definecolor{currentstroke}{rgb}{0.200000,0.200000,0.200000}%
\pgfsetstrokecolor{currentstroke}%
\pgfsetdash{}{0pt}%
\pgfpathmoveto{\pgfqpoint{1.829418in}{2.927009in}}%
\pgfpathlineto{\pgfqpoint{1.879280in}{2.878880in}}%
\pgfpathlineto{\pgfqpoint{1.911968in}{2.858156in}}%
\pgfpathlineto{\pgfqpoint{1.862098in}{2.906318in}}%
\pgfpathlineto{\pgfqpoint{1.829418in}{2.927009in}}%
\pgfpathclose%
\pgfusepath{stroke,fill}%
\end{pgfscope}%
\begin{pgfscope}%
\pgfpathrectangle{\pgfqpoint{0.050000in}{0.050000in}}{\pgfqpoint{4.900000in}{4.900000in}}%
\pgfusepath{clip}%
\pgfsetbuttcap%
\pgfsetroundjoin%
\definecolor{currentfill}{rgb}{0.440323,0.440323,0.440323}%
\pgfsetfillcolor{currentfill}%
\pgfsetfillopacity{0.900000}%
\pgfsetlinewidth{0.501875pt}%
\definecolor{currentstroke}{rgb}{0.200000,0.200000,0.200000}%
\pgfsetstrokecolor{currentstroke}%
\pgfsetdash{}{0pt}%
\pgfpathmoveto{\pgfqpoint{2.588147in}{2.313839in}}%
\pgfpathlineto{\pgfqpoint{2.636313in}{2.267082in}}%
\pgfpathlineto{\pgfqpoint{2.670243in}{2.245356in}}%
\pgfpathlineto{\pgfqpoint{2.622075in}{2.292141in}}%
\pgfpathlineto{\pgfqpoint{2.588147in}{2.313839in}}%
\pgfpathclose%
\pgfusepath{stroke,fill}%
\end{pgfscope}%
\begin{pgfscope}%
\pgfpathrectangle{\pgfqpoint{0.050000in}{0.050000in}}{\pgfqpoint{4.900000in}{4.900000in}}%
\pgfusepath{clip}%
\pgfsetbuttcap%
\pgfsetroundjoin%
\definecolor{currentfill}{rgb}{0.045521,0.045521,0.045521}%
\pgfsetfillcolor{currentfill}%
\pgfsetfillopacity{0.900000}%
\pgfsetlinewidth{0.501875pt}%
\definecolor{currentstroke}{rgb}{0.200000,0.200000,0.200000}%
\pgfsetstrokecolor{currentstroke}%
\pgfsetdash{}{0pt}%
\pgfpathmoveto{\pgfqpoint{1.548943in}{3.154395in}}%
\pgfpathlineto{\pgfqpoint{1.599442in}{3.105748in}}%
\pgfpathlineto{\pgfqpoint{1.631677in}{3.085391in}}%
\pgfpathlineto{\pgfqpoint{1.581169in}{3.134073in}}%
\pgfpathlineto{\pgfqpoint{1.548943in}{3.154395in}}%
\pgfpathclose%
\pgfusepath{stroke,fill}%
\end{pgfscope}%
\begin{pgfscope}%
\pgfpathrectangle{\pgfqpoint{0.050000in}{0.050000in}}{\pgfqpoint{4.900000in}{4.900000in}}%
\pgfusepath{clip}%
\pgfsetbuttcap%
\pgfsetroundjoin%
\definecolor{currentfill}{rgb}{0.346943,0.346943,0.346943}%
\pgfsetfillcolor{currentfill}%
\pgfsetfillopacity{0.900000}%
\pgfsetlinewidth{0.501875pt}%
\definecolor{currentstroke}{rgb}{0.200000,0.200000,0.200000}%
\pgfsetstrokecolor{currentstroke}%
\pgfsetdash{}{0pt}%
\pgfpathmoveto{\pgfqpoint{2.307788in}{2.541124in}}%
\pgfpathlineto{\pgfqpoint{2.356590in}{2.493847in}}%
\pgfpathlineto{\pgfqpoint{2.390068in}{2.472488in}}%
\pgfpathlineto{\pgfqpoint{2.341261in}{2.519794in}}%
\pgfpathlineto{\pgfqpoint{2.307788in}{2.541124in}}%
\pgfpathclose%
\pgfusepath{stroke,fill}%
\end{pgfscope}%
\begin{pgfscope}%
\pgfpathrectangle{\pgfqpoint{0.050000in}{0.050000in}}{\pgfqpoint{4.900000in}{4.900000in}}%
\pgfusepath{clip}%
\pgfsetbuttcap%
\pgfsetroundjoin%
\definecolor{currentfill}{rgb}{0.234371,0.234371,0.234371}%
\pgfsetfillcolor{currentfill}%
\pgfsetfillopacity{0.900000}%
\pgfsetlinewidth{0.501875pt}%
\definecolor{currentstroke}{rgb}{0.200000,0.200000,0.200000}%
\pgfsetstrokecolor{currentstroke}%
\pgfsetdash{}{0pt}%
\pgfpathmoveto{\pgfqpoint{2.027328in}{2.768493in}}%
\pgfpathlineto{\pgfqpoint{2.076767in}{2.720699in}}%
\pgfpathlineto{\pgfqpoint{2.109793in}{2.699707in}}%
\pgfpathlineto{\pgfqpoint{2.060347in}{2.747533in}}%
\pgfpathlineto{\pgfqpoint{2.027328in}{2.768493in}}%
\pgfpathclose%
\pgfusepath{stroke,fill}%
\end{pgfscope}%
\begin{pgfscope}%
\pgfpathrectangle{\pgfqpoint{0.050000in}{0.050000in}}{\pgfqpoint{4.900000in}{4.900000in}}%
\pgfusepath{clip}%
\pgfsetbuttcap%
\pgfsetroundjoin%
\definecolor{currentfill}{rgb}{0.504268,0.504268,0.504268}%
\pgfsetfillcolor{currentfill}%
\pgfsetfillopacity{0.900000}%
\pgfsetlinewidth{0.501875pt}%
\definecolor{currentstroke}{rgb}{0.200000,0.200000,0.200000}%
\pgfsetstrokecolor{currentstroke}%
\pgfsetdash{}{0pt}%
\pgfpathmoveto{\pgfqpoint{2.786396in}{2.155068in}}%
\pgfpathlineto{\pgfqpoint{2.834138in}{2.108663in}}%
\pgfpathlineto{\pgfqpoint{2.868407in}{2.086674in}}%
\pgfpathlineto{\pgfqpoint{2.820664in}{2.133103in}}%
\pgfpathlineto{\pgfqpoint{2.786396in}{2.155068in}}%
\pgfpathclose%
\pgfusepath{stroke,fill}%
\end{pgfscope}%
\begin{pgfscope}%
\pgfpathrectangle{\pgfqpoint{0.050000in}{0.050000in}}{\pgfqpoint{4.900000in}{4.900000in}}%
\pgfusepath{clip}%
\pgfsetbuttcap%
\pgfsetroundjoin%
\definecolor{currentfill}{rgb}{0.113802,0.113802,0.113802}%
\pgfsetfillcolor{currentfill}%
\pgfsetfillopacity{0.900000}%
\pgfsetlinewidth{0.501875pt}%
\definecolor{currentstroke}{rgb}{0.200000,0.200000,0.200000}%
\pgfsetstrokecolor{currentstroke}%
\pgfsetdash{}{0pt}%
\pgfpathmoveto{\pgfqpoint{1.746769in}{2.995947in}}%
\pgfpathlineto{\pgfqpoint{1.796845in}{2.947634in}}%
\pgfpathlineto{\pgfqpoint{1.829418in}{2.927009in}}%
\pgfpathlineto{\pgfqpoint{1.779333in}{2.975355in}}%
\pgfpathlineto{\pgfqpoint{1.746769in}{2.995947in}}%
\pgfpathclose%
\pgfusepath{stroke,fill}%
\end{pgfscope}%
\begin{pgfscope}%
\pgfpathrectangle{\pgfqpoint{0.050000in}{0.050000in}}{\pgfqpoint{4.900000in}{4.900000in}}%
\pgfusepath{clip}%
\pgfsetbuttcap%
\pgfsetroundjoin%
\definecolor{currentfill}{rgb}{0.407843,0.407843,0.407843}%
\pgfsetfillcolor{currentfill}%
\pgfsetfillopacity{0.900000}%
\pgfsetlinewidth{0.501875pt}%
\definecolor{currentstroke}{rgb}{0.200000,0.200000,0.200000}%
\pgfsetstrokecolor{currentstroke}%
\pgfsetdash{}{0pt}%
\pgfpathmoveto{\pgfqpoint{2.505952in}{2.382407in}}%
\pgfpathlineto{\pgfqpoint{2.554331in}{2.335466in}}%
\pgfpathlineto{\pgfqpoint{2.588147in}{2.313839in}}%
\pgfpathlineto{\pgfqpoint{2.539765in}{2.360808in}}%
\pgfpathlineto{\pgfqpoint{2.505952in}{2.382407in}}%
\pgfpathclose%
\pgfusepath{stroke,fill}%
\end{pgfscope}%
\begin{pgfscope}%
\pgfpathrectangle{\pgfqpoint{0.050000in}{0.050000in}}{\pgfqpoint{4.900000in}{4.900000in}}%
\pgfusepath{clip}%
\pgfsetbuttcap%
\pgfsetroundjoin%
\definecolor{currentfill}{rgb}{0.009104,0.009104,0.009104}%
\pgfsetfillcolor{currentfill}%
\pgfsetfillopacity{0.900000}%
\pgfsetlinewidth{0.501875pt}%
\definecolor{currentstroke}{rgb}{0.200000,0.200000,0.200000}%
\pgfsetstrokecolor{currentstroke}%
\pgfsetdash{}{0pt}%
\pgfpathmoveto{\pgfqpoint{1.466109in}{3.223483in}}%
\pgfpathlineto{\pgfqpoint{1.516824in}{3.174652in}}%
\pgfpathlineto{\pgfqpoint{1.548943in}{3.154395in}}%
\pgfpathlineto{\pgfqpoint{1.498218in}{3.203262in}}%
\pgfpathlineto{\pgfqpoint{1.466109in}{3.223483in}}%
\pgfpathclose%
\pgfusepath{stroke,fill}%
\end{pgfscope}%
\begin{pgfscope}%
\pgfpathrectangle{\pgfqpoint{0.050000in}{0.050000in}}{\pgfqpoint{4.900000in}{4.900000in}}%
\pgfusepath{clip}%
\pgfsetbuttcap%
\pgfsetroundjoin%
\definecolor{currentfill}{rgb}{0.317416,0.317416,0.317416}%
\pgfsetfillcolor{currentfill}%
\pgfsetfillopacity{0.900000}%
\pgfsetlinewidth{0.501875pt}%
\definecolor{currentstroke}{rgb}{0.200000,0.200000,0.200000}%
\pgfsetstrokecolor{currentstroke}%
\pgfsetdash{}{0pt}%
\pgfpathmoveto{\pgfqpoint{2.225407in}{2.609843in}}%
\pgfpathlineto{\pgfqpoint{2.274424in}{2.562384in}}%
\pgfpathlineto{\pgfqpoint{2.307788in}{2.541124in}}%
\pgfpathlineto{\pgfqpoint{2.258766in}{2.588614in}}%
\pgfpathlineto{\pgfqpoint{2.225407in}{2.609843in}}%
\pgfpathclose%
\pgfusepath{stroke,fill}%
\end{pgfscope}%
\begin{pgfscope}%
\pgfpathrectangle{\pgfqpoint{0.050000in}{0.050000in}}{\pgfqpoint{4.900000in}{4.900000in}}%
\pgfusepath{clip}%
\pgfsetbuttcap%
\pgfsetroundjoin%
\definecolor{currentfill}{rgb}{0.190081,0.190081,0.190081}%
\pgfsetfillcolor{currentfill}%
\pgfsetfillopacity{0.900000}%
\pgfsetlinewidth{0.501875pt}%
\definecolor{currentstroke}{rgb}{0.200000,0.200000,0.200000}%
\pgfsetstrokecolor{currentstroke}%
\pgfsetdash{}{0pt}%
\pgfpathmoveto{\pgfqpoint{1.944763in}{2.837364in}}%
\pgfpathlineto{\pgfqpoint{1.994417in}{2.789386in}}%
\pgfpathlineto{\pgfqpoint{2.027328in}{2.768493in}}%
\pgfpathlineto{\pgfqpoint{1.977666in}{2.816503in}}%
\pgfpathlineto{\pgfqpoint{1.944763in}{2.837364in}}%
\pgfpathclose%
\pgfusepath{stroke,fill}%
\end{pgfscope}%
\begin{pgfscope}%
\pgfpathrectangle{\pgfqpoint{0.050000in}{0.050000in}}{\pgfqpoint{4.900000in}{4.900000in}}%
\pgfusepath{clip}%
\pgfsetbuttcap%
\pgfsetroundjoin%
\definecolor{currentfill}{rgb}{0.469819,0.469819,0.469819}%
\pgfsetfillcolor{currentfill}%
\pgfsetfillopacity{0.900000}%
\pgfsetlinewidth{0.501875pt}%
\definecolor{currentstroke}{rgb}{0.200000,0.200000,0.200000}%
\pgfsetstrokecolor{currentstroke}%
\pgfsetdash{}{0pt}%
\pgfpathmoveto{\pgfqpoint{2.704285in}{2.223560in}}%
\pgfpathlineto{\pgfqpoint{2.752240in}{2.176962in}}%
\pgfpathlineto{\pgfqpoint{2.786396in}{2.155068in}}%
\pgfpathlineto{\pgfqpoint{2.738439in}{2.201692in}}%
\pgfpathlineto{\pgfqpoint{2.704285in}{2.223560in}}%
\pgfpathclose%
\pgfusepath{stroke,fill}%
\end{pgfscope}%
\begin{pgfscope}%
\pgfpathrectangle{\pgfqpoint{0.050000in}{0.050000in}}{\pgfqpoint{4.900000in}{4.900000in}}%
\pgfusepath{clip}%
\pgfsetbuttcap%
\pgfsetroundjoin%
\definecolor{currentfill}{rgb}{0.077386,0.077386,0.077386}%
\pgfsetfillcolor{currentfill}%
\pgfsetfillopacity{0.900000}%
\pgfsetlinewidth{0.501875pt}%
\definecolor{currentstroke}{rgb}{0.200000,0.200000,0.200000}%
\pgfsetstrokecolor{currentstroke}%
\pgfsetdash{}{0pt}%
\pgfpathmoveto{\pgfqpoint{1.664019in}{3.064968in}}%
\pgfpathlineto{\pgfqpoint{1.714311in}{3.016471in}}%
\pgfpathlineto{\pgfqpoint{1.746769in}{2.995947in}}%
\pgfpathlineto{\pgfqpoint{1.696467in}{3.044477in}}%
\pgfpathlineto{\pgfqpoint{1.664019in}{3.064968in}}%
\pgfpathclose%
\pgfusepath{stroke,fill}%
\end{pgfscope}%
\begin{pgfscope}%
\pgfpathrectangle{\pgfqpoint{0.050000in}{0.050000in}}{\pgfqpoint{4.900000in}{4.900000in}}%
\pgfusepath{clip}%
\pgfsetbuttcap%
\pgfsetroundjoin%
\definecolor{currentfill}{rgb}{0.379423,0.379423,0.379423}%
\pgfsetfillcolor{currentfill}%
\pgfsetfillopacity{0.900000}%
\pgfsetlinewidth{0.501875pt}%
\definecolor{currentstroke}{rgb}{0.200000,0.200000,0.200000}%
\pgfsetstrokecolor{currentstroke}%
\pgfsetdash{}{0pt}%
\pgfpathmoveto{\pgfqpoint{2.423656in}{2.451059in}}%
\pgfpathlineto{\pgfqpoint{2.472249in}{2.403935in}}%
\pgfpathlineto{\pgfqpoint{2.505952in}{2.382407in}}%
\pgfpathlineto{\pgfqpoint{2.457355in}{2.429560in}}%
\pgfpathlineto{\pgfqpoint{2.423656in}{2.451059in}}%
\pgfpathclose%
\pgfusepath{stroke,fill}%
\end{pgfscope}%
\begin{pgfscope}%
\pgfpathrectangle{\pgfqpoint{0.050000in}{0.050000in}}{\pgfqpoint{4.900000in}{4.900000in}}%
\pgfusepath{clip}%
\pgfsetbuttcap%
\pgfsetroundjoin%
\definecolor{currentfill}{rgb}{0.273126,0.273126,0.273126}%
\pgfsetfillcolor{currentfill}%
\pgfsetfillopacity{0.900000}%
\pgfsetlinewidth{0.501875pt}%
\definecolor{currentstroke}{rgb}{0.200000,0.200000,0.200000}%
\pgfsetstrokecolor{currentstroke}%
\pgfsetdash{}{0pt}%
\pgfpathmoveto{\pgfqpoint{2.142927in}{2.678646in}}%
\pgfpathlineto{\pgfqpoint{2.192158in}{2.631004in}}%
\pgfpathlineto{\pgfqpoint{2.225407in}{2.609843in}}%
\pgfpathlineto{\pgfqpoint{2.176170in}{2.657517in}}%
\pgfpathlineto{\pgfqpoint{2.142927in}{2.678646in}}%
\pgfpathclose%
\pgfusepath{stroke,fill}%
\end{pgfscope}%
\begin{pgfscope}%
\pgfpathrectangle{\pgfqpoint{0.050000in}{0.050000in}}{\pgfqpoint{4.900000in}{4.900000in}}%
\pgfusepath{clip}%
\pgfsetbuttcap%
\pgfsetroundjoin%
\definecolor{currentfill}{rgb}{0.538716,0.538716,0.538716}%
\pgfsetfillcolor{currentfill}%
\pgfsetfillopacity{0.900000}%
\pgfsetlinewidth{0.501875pt}%
\definecolor{currentstroke}{rgb}{0.200000,0.200000,0.200000}%
\pgfsetstrokecolor{currentstroke}%
\pgfsetdash{}{0pt}%
\pgfpathmoveto{\pgfqpoint{2.902789in}{2.064613in}}%
\pgfpathlineto{\pgfqpoint{2.950320in}{2.018395in}}%
\pgfpathlineto{\pgfqpoint{2.984815in}{1.996242in}}%
\pgfpathlineto{\pgfqpoint{2.937284in}{2.042481in}}%
\pgfpathlineto{\pgfqpoint{2.902789in}{2.064613in}}%
\pgfpathclose%
\pgfusepath{stroke,fill}%
\end{pgfscope}%
\begin{pgfscope}%
\pgfpathrectangle{\pgfqpoint{0.050000in}{0.050000in}}{\pgfqpoint{4.900000in}{4.900000in}}%
\pgfusepath{clip}%
\pgfsetbuttcap%
\pgfsetroundjoin%
\definecolor{currentfill}{rgb}{0.145790,0.145790,0.145790}%
\pgfsetfillcolor{currentfill}%
\pgfsetfillopacity{0.900000}%
\pgfsetlinewidth{0.501875pt}%
\definecolor{currentstroke}{rgb}{0.200000,0.200000,0.200000}%
\pgfsetstrokecolor{currentstroke}%
\pgfsetdash{}{0pt}%
\pgfpathmoveto{\pgfqpoint{1.862098in}{2.906318in}}%
\pgfpathlineto{\pgfqpoint{1.911968in}{2.858156in}}%
\pgfpathlineto{\pgfqpoint{1.944763in}{2.837364in}}%
\pgfpathlineto{\pgfqpoint{1.894885in}{2.885558in}}%
\pgfpathlineto{\pgfqpoint{1.862098in}{2.906318in}}%
\pgfpathclose%
\pgfusepath{stroke,fill}%
\end{pgfscope}%
\begin{pgfscope}%
\pgfpathrectangle{\pgfqpoint{0.050000in}{0.050000in}}{\pgfqpoint{4.900000in}{4.900000in}}%
\pgfusepath{clip}%
\pgfsetbuttcap%
\pgfsetroundjoin%
\definecolor{currentfill}{rgb}{0.440323,0.440323,0.440323}%
\pgfsetfillcolor{currentfill}%
\pgfsetfillopacity{0.900000}%
\pgfsetlinewidth{0.501875pt}%
\definecolor{currentstroke}{rgb}{0.200000,0.200000,0.200000}%
\pgfsetstrokecolor{currentstroke}%
\pgfsetdash{}{0pt}%
\pgfpathmoveto{\pgfqpoint{2.622075in}{2.292141in}}%
\pgfpathlineto{\pgfqpoint{2.670243in}{2.245356in}}%
\pgfpathlineto{\pgfqpoint{2.704285in}{2.223560in}}%
\pgfpathlineto{\pgfqpoint{2.656114in}{2.270372in}}%
\pgfpathlineto{\pgfqpoint{2.622075in}{2.292141in}}%
\pgfpathclose%
\pgfusepath{stroke,fill}%
\end{pgfscope}%
\begin{pgfscope}%
\pgfpathrectangle{\pgfqpoint{0.050000in}{0.050000in}}{\pgfqpoint{4.900000in}{4.900000in}}%
\pgfusepath{clip}%
\pgfsetbuttcap%
\pgfsetroundjoin%
\definecolor{currentfill}{rgb}{0.045521,0.045521,0.045521}%
\pgfsetfillcolor{currentfill}%
\pgfsetfillopacity{0.900000}%
\pgfsetlinewidth{0.501875pt}%
\definecolor{currentstroke}{rgb}{0.200000,0.200000,0.200000}%
\pgfsetstrokecolor{currentstroke}%
\pgfsetdash{}{0pt}%
\pgfpathmoveto{\pgfqpoint{1.581169in}{3.134073in}}%
\pgfpathlineto{\pgfqpoint{1.631677in}{3.085391in}}%
\pgfpathlineto{\pgfqpoint{1.664019in}{3.064968in}}%
\pgfpathlineto{\pgfqpoint{1.613500in}{3.113683in}}%
\pgfpathlineto{\pgfqpoint{1.581169in}{3.134073in}}%
\pgfpathclose%
\pgfusepath{stroke,fill}%
\end{pgfscope}%
\begin{pgfscope}%
\pgfpathrectangle{\pgfqpoint{0.050000in}{0.050000in}}{\pgfqpoint{4.900000in}{4.900000in}}%
\pgfusepath{clip}%
\pgfsetbuttcap%
\pgfsetroundjoin%
\definecolor{currentfill}{rgb}{0.346943,0.346943,0.346943}%
\pgfsetfillcolor{currentfill}%
\pgfsetfillopacity{0.900000}%
\pgfsetlinewidth{0.501875pt}%
\definecolor{currentstroke}{rgb}{0.200000,0.200000,0.200000}%
\pgfsetstrokecolor{currentstroke}%
\pgfsetdash{}{0pt}%
\pgfpathmoveto{\pgfqpoint{2.341261in}{2.519794in}}%
\pgfpathlineto{\pgfqpoint{2.390068in}{2.472488in}}%
\pgfpathlineto{\pgfqpoint{2.423656in}{2.451059in}}%
\pgfpathlineto{\pgfqpoint{2.374845in}{2.498395in}}%
\pgfpathlineto{\pgfqpoint{2.341261in}{2.519794in}}%
\pgfpathclose%
\pgfusepath{stroke,fill}%
\end{pgfscope}%
\begin{pgfscope}%
\pgfpathrectangle{\pgfqpoint{0.050000in}{0.050000in}}{\pgfqpoint{4.900000in}{4.900000in}}%
\pgfusepath{clip}%
\pgfsetbuttcap%
\pgfsetroundjoin%
\definecolor{currentfill}{rgb}{0.234371,0.234371,0.234371}%
\pgfsetfillcolor{currentfill}%
\pgfsetfillopacity{0.900000}%
\pgfsetlinewidth{0.501875pt}%
\definecolor{currentstroke}{rgb}{0.200000,0.200000,0.200000}%
\pgfsetstrokecolor{currentstroke}%
\pgfsetdash{}{0pt}%
\pgfpathmoveto{\pgfqpoint{2.060347in}{2.747533in}}%
\pgfpathlineto{\pgfqpoint{2.109793in}{2.699707in}}%
\pgfpathlineto{\pgfqpoint{2.142927in}{2.678646in}}%
\pgfpathlineto{\pgfqpoint{2.093474in}{2.726504in}}%
\pgfpathlineto{\pgfqpoint{2.060347in}{2.747533in}}%
\pgfpathclose%
\pgfusepath{stroke,fill}%
\end{pgfscope}%
\begin{pgfscope}%
\pgfpathrectangle{\pgfqpoint{0.050000in}{0.050000in}}{\pgfqpoint{4.900000in}{4.900000in}}%
\pgfusepath{clip}%
\pgfsetbuttcap%
\pgfsetroundjoin%
\definecolor{currentfill}{rgb}{0.504268,0.504268,0.504268}%
\pgfsetfillcolor{currentfill}%
\pgfsetfillopacity{0.900000}%
\pgfsetlinewidth{0.501875pt}%
\definecolor{currentstroke}{rgb}{0.200000,0.200000,0.200000}%
\pgfsetstrokecolor{currentstroke}%
\pgfsetdash{}{0pt}%
\pgfpathmoveto{\pgfqpoint{2.820664in}{2.133103in}}%
\pgfpathlineto{\pgfqpoint{2.868407in}{2.086674in}}%
\pgfpathlineto{\pgfqpoint{2.902789in}{2.064613in}}%
\pgfpathlineto{\pgfqpoint{2.855044in}{2.111066in}}%
\pgfpathlineto{\pgfqpoint{2.820664in}{2.133103in}}%
\pgfpathclose%
\pgfusepath{stroke,fill}%
\end{pgfscope}%
\begin{pgfscope}%
\pgfpathrectangle{\pgfqpoint{0.050000in}{0.050000in}}{\pgfqpoint{4.900000in}{4.900000in}}%
\pgfusepath{clip}%
\pgfsetbuttcap%
\pgfsetroundjoin%
\definecolor{currentfill}{rgb}{0.113802,0.113802,0.113802}%
\pgfsetfillcolor{currentfill}%
\pgfsetfillopacity{0.900000}%
\pgfsetlinewidth{0.501875pt}%
\definecolor{currentstroke}{rgb}{0.200000,0.200000,0.200000}%
\pgfsetstrokecolor{currentstroke}%
\pgfsetdash{}{0pt}%
\pgfpathmoveto{\pgfqpoint{1.779333in}{2.975355in}}%
\pgfpathlineto{\pgfqpoint{1.829418in}{2.927009in}}%
\pgfpathlineto{\pgfqpoint{1.862098in}{2.906318in}}%
\pgfpathlineto{\pgfqpoint{1.812004in}{2.954696in}}%
\pgfpathlineto{\pgfqpoint{1.779333in}{2.975355in}}%
\pgfpathclose%
\pgfusepath{stroke,fill}%
\end{pgfscope}%
\begin{pgfscope}%
\pgfpathrectangle{\pgfqpoint{0.050000in}{0.050000in}}{\pgfqpoint{4.900000in}{4.900000in}}%
\pgfusepath{clip}%
\pgfsetbuttcap%
\pgfsetroundjoin%
\definecolor{currentfill}{rgb}{0.407843,0.407843,0.407843}%
\pgfsetfillcolor{currentfill}%
\pgfsetfillopacity{0.900000}%
\pgfsetlinewidth{0.501875pt}%
\definecolor{currentstroke}{rgb}{0.200000,0.200000,0.200000}%
\pgfsetstrokecolor{currentstroke}%
\pgfsetdash{}{0pt}%
\pgfpathmoveto{\pgfqpoint{2.539765in}{2.360808in}}%
\pgfpathlineto{\pgfqpoint{2.588147in}{2.313839in}}%
\pgfpathlineto{\pgfqpoint{2.622075in}{2.292141in}}%
\pgfpathlineto{\pgfqpoint{2.573689in}{2.339139in}}%
\pgfpathlineto{\pgfqpoint{2.539765in}{2.360808in}}%
\pgfpathclose%
\pgfusepath{stroke,fill}%
\end{pgfscope}%
\begin{pgfscope}%
\pgfpathrectangle{\pgfqpoint{0.050000in}{0.050000in}}{\pgfqpoint{4.900000in}{4.900000in}}%
\pgfusepath{clip}%
\pgfsetbuttcap%
\pgfsetroundjoin%
\definecolor{currentfill}{rgb}{0.009104,0.009104,0.009104}%
\pgfsetfillcolor{currentfill}%
\pgfsetfillopacity{0.900000}%
\pgfsetlinewidth{0.501875pt}%
\definecolor{currentstroke}{rgb}{0.200000,0.200000,0.200000}%
\pgfsetstrokecolor{currentstroke}%
\pgfsetdash{}{0pt}%
\pgfpathmoveto{\pgfqpoint{1.498218in}{3.203262in}}%
\pgfpathlineto{\pgfqpoint{1.548943in}{3.154395in}}%
\pgfpathlineto{\pgfqpoint{1.581169in}{3.134073in}}%
\pgfpathlineto{\pgfqpoint{1.530432in}{3.182973in}}%
\pgfpathlineto{\pgfqpoint{1.498218in}{3.203262in}}%
\pgfpathclose%
\pgfusepath{stroke,fill}%
\end{pgfscope}%
\begin{pgfscope}%
\pgfpathrectangle{\pgfqpoint{0.050000in}{0.050000in}}{\pgfqpoint{4.900000in}{4.900000in}}%
\pgfusepath{clip}%
\pgfsetbuttcap%
\pgfsetroundjoin%
\definecolor{currentfill}{rgb}{0.317416,0.317416,0.317416}%
\pgfsetfillcolor{currentfill}%
\pgfsetfillopacity{0.900000}%
\pgfsetlinewidth{0.501875pt}%
\definecolor{currentstroke}{rgb}{0.200000,0.200000,0.200000}%
\pgfsetstrokecolor{currentstroke}%
\pgfsetdash{}{0pt}%
\pgfpathmoveto{\pgfqpoint{2.258766in}{2.588614in}}%
\pgfpathlineto{\pgfqpoint{2.307788in}{2.541124in}}%
\pgfpathlineto{\pgfqpoint{2.341261in}{2.519794in}}%
\pgfpathlineto{\pgfqpoint{2.292234in}{2.567314in}}%
\pgfpathlineto{\pgfqpoint{2.258766in}{2.588614in}}%
\pgfpathclose%
\pgfusepath{stroke,fill}%
\end{pgfscope}%
\begin{pgfscope}%
\pgfpathrectangle{\pgfqpoint{0.050000in}{0.050000in}}{\pgfqpoint{4.900000in}{4.900000in}}%
\pgfusepath{clip}%
\pgfsetbuttcap%
\pgfsetroundjoin%
\definecolor{currentfill}{rgb}{0.190081,0.190081,0.190081}%
\pgfsetfillcolor{currentfill}%
\pgfsetfillopacity{0.900000}%
\pgfsetlinewidth{0.501875pt}%
\definecolor{currentstroke}{rgb}{0.200000,0.200000,0.200000}%
\pgfsetstrokecolor{currentstroke}%
\pgfsetdash{}{0pt}%
\pgfpathmoveto{\pgfqpoint{1.977666in}{2.816503in}}%
\pgfpathlineto{\pgfqpoint{2.027328in}{2.768493in}}%
\pgfpathlineto{\pgfqpoint{2.060347in}{2.747533in}}%
\pgfpathlineto{\pgfqpoint{2.010678in}{2.795575in}}%
\pgfpathlineto{\pgfqpoint{1.977666in}{2.816503in}}%
\pgfpathclose%
\pgfusepath{stroke,fill}%
\end{pgfscope}%
\begin{pgfscope}%
\pgfpathrectangle{\pgfqpoint{0.050000in}{0.050000in}}{\pgfqpoint{4.900000in}{4.900000in}}%
\pgfusepath{clip}%
\pgfsetbuttcap%
\pgfsetroundjoin%
\definecolor{currentfill}{rgb}{0.474125,0.474125,0.474125}%
\pgfsetfillcolor{currentfill}%
\pgfsetfillopacity{0.900000}%
\pgfsetlinewidth{0.501875pt}%
\definecolor{currentstroke}{rgb}{0.200000,0.200000,0.200000}%
\pgfsetstrokecolor{currentstroke}%
\pgfsetdash{}{0pt}%
\pgfpathmoveto{\pgfqpoint{2.738439in}{2.201692in}}%
\pgfpathlineto{\pgfqpoint{2.786396in}{2.155068in}}%
\pgfpathlineto{\pgfqpoint{2.820664in}{2.133103in}}%
\pgfpathlineto{\pgfqpoint{2.772705in}{2.179753in}}%
\pgfpathlineto{\pgfqpoint{2.738439in}{2.201692in}}%
\pgfpathclose%
\pgfusepath{stroke,fill}%
\end{pgfscope}%
\begin{pgfscope}%
\pgfpathrectangle{\pgfqpoint{0.050000in}{0.050000in}}{\pgfqpoint{4.900000in}{4.900000in}}%
\pgfusepath{clip}%
\pgfsetbuttcap%
\pgfsetroundjoin%
\definecolor{currentfill}{rgb}{0.077386,0.077386,0.077386}%
\pgfsetfillcolor{currentfill}%
\pgfsetfillopacity{0.900000}%
\pgfsetlinewidth{0.501875pt}%
\definecolor{currentstroke}{rgb}{0.200000,0.200000,0.200000}%
\pgfsetstrokecolor{currentstroke}%
\pgfsetdash{}{0pt}%
\pgfpathmoveto{\pgfqpoint{1.696467in}{3.044477in}}%
\pgfpathlineto{\pgfqpoint{1.746769in}{2.995947in}}%
\pgfpathlineto{\pgfqpoint{1.779333in}{2.975355in}}%
\pgfpathlineto{\pgfqpoint{1.729021in}{3.023919in}}%
\pgfpathlineto{\pgfqpoint{1.696467in}{3.044477in}}%
\pgfpathclose%
\pgfusepath{stroke,fill}%
\end{pgfscope}%
\begin{pgfscope}%
\pgfpathrectangle{\pgfqpoint{0.050000in}{0.050000in}}{\pgfqpoint{4.900000in}{4.900000in}}%
\pgfusepath{clip}%
\pgfsetbuttcap%
\pgfsetroundjoin%
\definecolor{currentfill}{rgb}{0.379423,0.379423,0.379423}%
\pgfsetfillcolor{currentfill}%
\pgfsetfillopacity{0.900000}%
\pgfsetlinewidth{0.501875pt}%
\definecolor{currentstroke}{rgb}{0.200000,0.200000,0.200000}%
\pgfsetstrokecolor{currentstroke}%
\pgfsetdash{}{0pt}%
\pgfpathmoveto{\pgfqpoint{2.457355in}{2.429560in}}%
\pgfpathlineto{\pgfqpoint{2.505952in}{2.382407in}}%
\pgfpathlineto{\pgfqpoint{2.539765in}{2.360808in}}%
\pgfpathlineto{\pgfqpoint{2.491164in}{2.407990in}}%
\pgfpathlineto{\pgfqpoint{2.457355in}{2.429560in}}%
\pgfpathclose%
\pgfusepath{stroke,fill}%
\end{pgfscope}%
\begin{pgfscope}%
\pgfpathrectangle{\pgfqpoint{0.050000in}{0.050000in}}{\pgfqpoint{4.900000in}{4.900000in}}%
\pgfusepath{clip}%
\pgfsetbuttcap%
\pgfsetroundjoin%
\definecolor{currentfill}{rgb}{0.273126,0.273126,0.273126}%
\pgfsetfillcolor{currentfill}%
\pgfsetfillopacity{0.900000}%
\pgfsetlinewidth{0.501875pt}%
\definecolor{currentstroke}{rgb}{0.200000,0.200000,0.200000}%
\pgfsetstrokecolor{currentstroke}%
\pgfsetdash{}{0pt}%
\pgfpathmoveto{\pgfqpoint{2.176170in}{2.657517in}}%
\pgfpathlineto{\pgfqpoint{2.225407in}{2.609843in}}%
\pgfpathlineto{\pgfqpoint{2.258766in}{2.588614in}}%
\pgfpathlineto{\pgfqpoint{2.209523in}{2.636318in}}%
\pgfpathlineto{\pgfqpoint{2.176170in}{2.657517in}}%
\pgfpathclose%
\pgfusepath{stroke,fill}%
\end{pgfscope}%
\begin{pgfscope}%
\pgfpathrectangle{\pgfqpoint{0.050000in}{0.050000in}}{\pgfqpoint{4.900000in}{4.900000in}}%
\pgfusepath{clip}%
\pgfsetbuttcap%
\pgfsetroundjoin%
\definecolor{currentfill}{rgb}{0.151326,0.151326,0.151326}%
\pgfsetfillcolor{currentfill}%
\pgfsetfillopacity{0.900000}%
\pgfsetlinewidth{0.501875pt}%
\definecolor{currentstroke}{rgb}{0.200000,0.200000,0.200000}%
\pgfsetstrokecolor{currentstroke}%
\pgfsetdash{}{0pt}%
\pgfpathmoveto{\pgfqpoint{1.894885in}{2.885558in}}%
\pgfpathlineto{\pgfqpoint{1.944763in}{2.837364in}}%
\pgfpathlineto{\pgfqpoint{1.977666in}{2.816503in}}%
\pgfpathlineto{\pgfqpoint{1.927781in}{2.864730in}}%
\pgfpathlineto{\pgfqpoint{1.894885in}{2.885558in}}%
\pgfpathclose%
\pgfusepath{stroke,fill}%
\end{pgfscope}%
\begin{pgfscope}%
\pgfpathrectangle{\pgfqpoint{0.050000in}{0.050000in}}{\pgfqpoint{4.900000in}{4.900000in}}%
\pgfusepath{clip}%
\pgfsetbuttcap%
\pgfsetroundjoin%
\definecolor{currentfill}{rgb}{0.440323,0.440323,0.440323}%
\pgfsetfillcolor{currentfill}%
\pgfsetfillopacity{0.900000}%
\pgfsetlinewidth{0.501875pt}%
\definecolor{currentstroke}{rgb}{0.200000,0.200000,0.200000}%
\pgfsetstrokecolor{currentstroke}%
\pgfsetdash{}{0pt}%
\pgfpathmoveto{\pgfqpoint{2.656114in}{2.270372in}}%
\pgfpathlineto{\pgfqpoint{2.704285in}{2.223560in}}%
\pgfpathlineto{\pgfqpoint{2.738439in}{2.201692in}}%
\pgfpathlineto{\pgfqpoint{2.690266in}{2.248532in}}%
\pgfpathlineto{\pgfqpoint{2.656114in}{2.270372in}}%
\pgfpathclose%
\pgfusepath{stroke,fill}%
\end{pgfscope}%
\begin{pgfscope}%
\pgfpathrectangle{\pgfqpoint{0.050000in}{0.050000in}}{\pgfqpoint{4.900000in}{4.900000in}}%
\pgfusepath{clip}%
\pgfsetbuttcap%
\pgfsetroundjoin%
\definecolor{currentfill}{rgb}{0.045521,0.045521,0.045521}%
\pgfsetfillcolor{currentfill}%
\pgfsetfillopacity{0.900000}%
\pgfsetlinewidth{0.501875pt}%
\definecolor{currentstroke}{rgb}{0.200000,0.200000,0.200000}%
\pgfsetstrokecolor{currentstroke}%
\pgfsetdash{}{0pt}%
\pgfpathmoveto{\pgfqpoint{1.613500in}{3.113683in}}%
\pgfpathlineto{\pgfqpoint{1.664019in}{3.064968in}}%
\pgfpathlineto{\pgfqpoint{1.696467in}{3.044477in}}%
\pgfpathlineto{\pgfqpoint{1.645937in}{3.093227in}}%
\pgfpathlineto{\pgfqpoint{1.613500in}{3.113683in}}%
\pgfpathclose%
\pgfusepath{stroke,fill}%
\end{pgfscope}%
\begin{pgfscope}%
\pgfpathrectangle{\pgfqpoint{0.050000in}{0.050000in}}{\pgfqpoint{4.900000in}{4.900000in}}%
\pgfusepath{clip}%
\pgfsetbuttcap%
\pgfsetroundjoin%
\definecolor{currentfill}{rgb}{0.346943,0.346943,0.346943}%
\pgfsetfillcolor{currentfill}%
\pgfsetfillopacity{0.900000}%
\pgfsetlinewidth{0.501875pt}%
\definecolor{currentstroke}{rgb}{0.200000,0.200000,0.200000}%
\pgfsetstrokecolor{currentstroke}%
\pgfsetdash{}{0pt}%
\pgfpathmoveto{\pgfqpoint{2.374845in}{2.498395in}}%
\pgfpathlineto{\pgfqpoint{2.423656in}{2.451059in}}%
\pgfpathlineto{\pgfqpoint{2.457355in}{2.429560in}}%
\pgfpathlineto{\pgfqpoint{2.408539in}{2.476925in}}%
\pgfpathlineto{\pgfqpoint{2.374845in}{2.498395in}}%
\pgfpathclose%
\pgfusepath{stroke,fill}%
\end{pgfscope}%
\begin{pgfscope}%
\pgfpathrectangle{\pgfqpoint{0.050000in}{0.050000in}}{\pgfqpoint{4.900000in}{4.900000in}}%
\pgfusepath{clip}%
\pgfsetbuttcap%
\pgfsetroundjoin%
\definecolor{currentfill}{rgb}{0.234371,0.234371,0.234371}%
\pgfsetfillcolor{currentfill}%
\pgfsetfillopacity{0.900000}%
\pgfsetlinewidth{0.501875pt}%
\definecolor{currentstroke}{rgb}{0.200000,0.200000,0.200000}%
\pgfsetstrokecolor{currentstroke}%
\pgfsetdash{}{0pt}%
\pgfpathmoveto{\pgfqpoint{2.093474in}{2.726504in}}%
\pgfpathlineto{\pgfqpoint{2.142927in}{2.678646in}}%
\pgfpathlineto{\pgfqpoint{2.176170in}{2.657517in}}%
\pgfpathlineto{\pgfqpoint{2.126711in}{2.705405in}}%
\pgfpathlineto{\pgfqpoint{2.093474in}{2.726504in}}%
\pgfpathclose%
\pgfusepath{stroke,fill}%
\end{pgfscope}%
\begin{pgfscope}%
\pgfpathrectangle{\pgfqpoint{0.050000in}{0.050000in}}{\pgfqpoint{4.900000in}{4.900000in}}%
\pgfusepath{clip}%
\pgfsetbuttcap%
\pgfsetroundjoin%
\definecolor{currentfill}{rgb}{0.504268,0.504268,0.504268}%
\pgfsetfillcolor{currentfill}%
\pgfsetfillopacity{0.900000}%
\pgfsetlinewidth{0.501875pt}%
\definecolor{currentstroke}{rgb}{0.200000,0.200000,0.200000}%
\pgfsetstrokecolor{currentstroke}%
\pgfsetdash{}{0pt}%
\pgfpathmoveto{\pgfqpoint{2.855044in}{2.111066in}}%
\pgfpathlineto{\pgfqpoint{2.902789in}{2.064613in}}%
\pgfpathlineto{\pgfqpoint{2.937284in}{2.042481in}}%
\pgfpathlineto{\pgfqpoint{2.889538in}{2.088957in}}%
\pgfpathlineto{\pgfqpoint{2.855044in}{2.111066in}}%
\pgfpathclose%
\pgfusepath{stroke,fill}%
\end{pgfscope}%
\begin{pgfscope}%
\pgfpathrectangle{\pgfqpoint{0.050000in}{0.050000in}}{\pgfqpoint{4.900000in}{4.900000in}}%
\pgfusepath{clip}%
\pgfsetbuttcap%
\pgfsetroundjoin%
\definecolor{currentfill}{rgb}{0.113802,0.113802,0.113802}%
\pgfsetfillcolor{currentfill}%
\pgfsetfillopacity{0.900000}%
\pgfsetlinewidth{0.501875pt}%
\definecolor{currentstroke}{rgb}{0.200000,0.200000,0.200000}%
\pgfsetstrokecolor{currentstroke}%
\pgfsetdash{}{0pt}%
\pgfpathmoveto{\pgfqpoint{1.812004in}{2.954696in}}%
\pgfpathlineto{\pgfqpoint{1.862098in}{2.906318in}}%
\pgfpathlineto{\pgfqpoint{1.894885in}{2.885558in}}%
\pgfpathlineto{\pgfqpoint{1.844782in}{2.933970in}}%
\pgfpathlineto{\pgfqpoint{1.812004in}{2.954696in}}%
\pgfpathclose%
\pgfusepath{stroke,fill}%
\end{pgfscope}%
\begin{pgfscope}%
\pgfpathrectangle{\pgfqpoint{0.050000in}{0.050000in}}{\pgfqpoint{4.900000in}{4.900000in}}%
\pgfusepath{clip}%
\pgfsetbuttcap%
\pgfsetroundjoin%
\definecolor{currentfill}{rgb}{0.411903,0.411903,0.411903}%
\pgfsetfillcolor{currentfill}%
\pgfsetfillopacity{0.900000}%
\pgfsetlinewidth{0.501875pt}%
\definecolor{currentstroke}{rgb}{0.200000,0.200000,0.200000}%
\pgfsetstrokecolor{currentstroke}%
\pgfsetdash{}{0pt}%
\pgfpathmoveto{\pgfqpoint{2.573689in}{2.339139in}}%
\pgfpathlineto{\pgfqpoint{2.622075in}{2.292141in}}%
\pgfpathlineto{\pgfqpoint{2.656114in}{2.270372in}}%
\pgfpathlineto{\pgfqpoint{2.607726in}{2.317398in}}%
\pgfpathlineto{\pgfqpoint{2.573689in}{2.339139in}}%
\pgfpathclose%
\pgfusepath{stroke,fill}%
\end{pgfscope}%
\begin{pgfscope}%
\pgfpathrectangle{\pgfqpoint{0.050000in}{0.050000in}}{\pgfqpoint{4.900000in}{4.900000in}}%
\pgfusepath{clip}%
\pgfsetbuttcap%
\pgfsetroundjoin%
\definecolor{currentfill}{rgb}{0.009104,0.009104,0.009104}%
\pgfsetfillcolor{currentfill}%
\pgfsetfillopacity{0.900000}%
\pgfsetlinewidth{0.501875pt}%
\definecolor{currentstroke}{rgb}{0.200000,0.200000,0.200000}%
\pgfsetstrokecolor{currentstroke}%
\pgfsetdash{}{0pt}%
\pgfpathmoveto{\pgfqpoint{1.530432in}{3.182973in}}%
\pgfpathlineto{\pgfqpoint{1.581169in}{3.134073in}}%
\pgfpathlineto{\pgfqpoint{1.613500in}{3.113683in}}%
\pgfpathlineto{\pgfqpoint{1.562753in}{3.162618in}}%
\pgfpathlineto{\pgfqpoint{1.530432in}{3.182973in}}%
\pgfpathclose%
\pgfusepath{stroke,fill}%
\end{pgfscope}%
\begin{pgfscope}%
\pgfpathrectangle{\pgfqpoint{0.050000in}{0.050000in}}{\pgfqpoint{4.900000in}{4.900000in}}%
\pgfusepath{clip}%
\pgfsetbuttcap%
\pgfsetroundjoin%
\definecolor{currentfill}{rgb}{0.317416,0.317416,0.317416}%
\pgfsetfillcolor{currentfill}%
\pgfsetfillopacity{0.900000}%
\pgfsetlinewidth{0.501875pt}%
\definecolor{currentstroke}{rgb}{0.200000,0.200000,0.200000}%
\pgfsetstrokecolor{currentstroke}%
\pgfsetdash{}{0pt}%
\pgfpathmoveto{\pgfqpoint{2.292234in}{2.567314in}}%
\pgfpathlineto{\pgfqpoint{2.341261in}{2.519794in}}%
\pgfpathlineto{\pgfqpoint{2.374845in}{2.498395in}}%
\pgfpathlineto{\pgfqpoint{2.325812in}{2.545945in}}%
\pgfpathlineto{\pgfqpoint{2.292234in}{2.567314in}}%
\pgfpathclose%
\pgfusepath{stroke,fill}%
\end{pgfscope}%
\begin{pgfscope}%
\pgfpathrectangle{\pgfqpoint{0.050000in}{0.050000in}}{\pgfqpoint{4.900000in}{4.900000in}}%
\pgfusepath{clip}%
\pgfsetbuttcap%
\pgfsetroundjoin%
\definecolor{currentfill}{rgb}{0.190081,0.190081,0.190081}%
\pgfsetfillcolor{currentfill}%
\pgfsetfillopacity{0.900000}%
\pgfsetlinewidth{0.501875pt}%
\definecolor{currentstroke}{rgb}{0.200000,0.200000,0.200000}%
\pgfsetstrokecolor{currentstroke}%
\pgfsetdash{}{0pt}%
\pgfpathmoveto{\pgfqpoint{2.010678in}{2.795575in}}%
\pgfpathlineto{\pgfqpoint{2.060347in}{2.747533in}}%
\pgfpathlineto{\pgfqpoint{2.093474in}{2.726504in}}%
\pgfpathlineto{\pgfqpoint{2.043798in}{2.774577in}}%
\pgfpathlineto{\pgfqpoint{2.010678in}{2.795575in}}%
\pgfpathclose%
\pgfusepath{stroke,fill}%
\end{pgfscope}%
\begin{pgfscope}%
\pgfpathrectangle{\pgfqpoint{0.050000in}{0.050000in}}{\pgfqpoint{4.900000in}{4.900000in}}%
\pgfusepath{clip}%
\pgfsetbuttcap%
\pgfsetroundjoin%
\definecolor{currentfill}{rgb}{0.474125,0.474125,0.474125}%
\pgfsetfillcolor{currentfill}%
\pgfsetfillopacity{0.900000}%
\pgfsetlinewidth{0.501875pt}%
\definecolor{currentstroke}{rgb}{0.200000,0.200000,0.200000}%
\pgfsetstrokecolor{currentstroke}%
\pgfsetdash{}{0pt}%
\pgfpathmoveto{\pgfqpoint{2.772705in}{2.179753in}}%
\pgfpathlineto{\pgfqpoint{2.820664in}{2.133103in}}%
\pgfpathlineto{\pgfqpoint{2.855044in}{2.111066in}}%
\pgfpathlineto{\pgfqpoint{2.807084in}{2.157742in}}%
\pgfpathlineto{\pgfqpoint{2.772705in}{2.179753in}}%
\pgfpathclose%
\pgfusepath{stroke,fill}%
\end{pgfscope}%
\begin{pgfscope}%
\pgfpathrectangle{\pgfqpoint{0.050000in}{0.050000in}}{\pgfqpoint{4.900000in}{4.900000in}}%
\pgfusepath{clip}%
\pgfsetbuttcap%
\pgfsetroundjoin%
\definecolor{currentfill}{rgb}{0.081938,0.081938,0.081938}%
\pgfsetfillcolor{currentfill}%
\pgfsetfillopacity{0.900000}%
\pgfsetlinewidth{0.501875pt}%
\definecolor{currentstroke}{rgb}{0.200000,0.200000,0.200000}%
\pgfsetstrokecolor{currentstroke}%
\pgfsetdash{}{0pt}%
\pgfpathmoveto{\pgfqpoint{1.729021in}{3.023919in}}%
\pgfpathlineto{\pgfqpoint{1.779333in}{2.975355in}}%
\pgfpathlineto{\pgfqpoint{1.812004in}{2.954696in}}%
\pgfpathlineto{\pgfqpoint{1.761683in}{3.003294in}}%
\pgfpathlineto{\pgfqpoint{1.729021in}{3.023919in}}%
\pgfpathclose%
\pgfusepath{stroke,fill}%
\end{pgfscope}%
\begin{pgfscope}%
\pgfpathrectangle{\pgfqpoint{0.050000in}{0.050000in}}{\pgfqpoint{4.900000in}{4.900000in}}%
\pgfusepath{clip}%
\pgfsetbuttcap%
\pgfsetroundjoin%
\definecolor{currentfill}{rgb}{0.379423,0.379423,0.379423}%
\pgfsetfillcolor{currentfill}%
\pgfsetfillopacity{0.900000}%
\pgfsetlinewidth{0.501875pt}%
\definecolor{currentstroke}{rgb}{0.200000,0.200000,0.200000}%
\pgfsetstrokecolor{currentstroke}%
\pgfsetdash{}{0pt}%
\pgfpathmoveto{\pgfqpoint{2.491164in}{2.407990in}}%
\pgfpathlineto{\pgfqpoint{2.539765in}{2.360808in}}%
\pgfpathlineto{\pgfqpoint{2.573689in}{2.339139in}}%
\pgfpathlineto{\pgfqpoint{2.525085in}{2.386349in}}%
\pgfpathlineto{\pgfqpoint{2.491164in}{2.407990in}}%
\pgfpathclose%
\pgfusepath{stroke,fill}%
\end{pgfscope}%
\begin{pgfscope}%
\pgfpathrectangle{\pgfqpoint{0.050000in}{0.050000in}}{\pgfqpoint{4.900000in}{4.900000in}}%
\pgfusepath{clip}%
\pgfsetbuttcap%
\pgfsetroundjoin%
\definecolor{currentfill}{rgb}{0.273126,0.273126,0.273126}%
\pgfsetfillcolor{currentfill}%
\pgfsetfillopacity{0.900000}%
\pgfsetlinewidth{0.501875pt}%
\definecolor{currentstroke}{rgb}{0.200000,0.200000,0.200000}%
\pgfsetstrokecolor{currentstroke}%
\pgfsetdash{}{0pt}%
\pgfpathmoveto{\pgfqpoint{2.209523in}{2.636318in}}%
\pgfpathlineto{\pgfqpoint{2.258766in}{2.588614in}}%
\pgfpathlineto{\pgfqpoint{2.292234in}{2.567314in}}%
\pgfpathlineto{\pgfqpoint{2.242985in}{2.615049in}}%
\pgfpathlineto{\pgfqpoint{2.209523in}{2.636318in}}%
\pgfpathclose%
\pgfusepath{stroke,fill}%
\end{pgfscope}%
\begin{pgfscope}%
\pgfpathrectangle{\pgfqpoint{0.050000in}{0.050000in}}{\pgfqpoint{4.900000in}{4.900000in}}%
\pgfusepath{clip}%
\pgfsetbuttcap%
\pgfsetroundjoin%
\definecolor{currentfill}{rgb}{0.151326,0.151326,0.151326}%
\pgfsetfillcolor{currentfill}%
\pgfsetfillopacity{0.900000}%
\pgfsetlinewidth{0.501875pt}%
\definecolor{currentstroke}{rgb}{0.200000,0.200000,0.200000}%
\pgfsetstrokecolor{currentstroke}%
\pgfsetdash{}{0pt}%
\pgfpathmoveto{\pgfqpoint{1.927781in}{2.864730in}}%
\pgfpathlineto{\pgfqpoint{1.977666in}{2.816503in}}%
\pgfpathlineto{\pgfqpoint{2.010678in}{2.795575in}}%
\pgfpathlineto{\pgfqpoint{1.960784in}{2.843833in}}%
\pgfpathlineto{\pgfqpoint{1.927781in}{2.864730in}}%
\pgfpathclose%
\pgfusepath{stroke,fill}%
\end{pgfscope}%
\begin{pgfscope}%
\pgfpathrectangle{\pgfqpoint{0.050000in}{0.050000in}}{\pgfqpoint{4.900000in}{4.900000in}}%
\pgfusepath{clip}%
\pgfsetbuttcap%
\pgfsetroundjoin%
\definecolor{currentfill}{rgb}{0.440323,0.440323,0.440323}%
\pgfsetfillcolor{currentfill}%
\pgfsetfillopacity{0.900000}%
\pgfsetlinewidth{0.501875pt}%
\definecolor{currentstroke}{rgb}{0.200000,0.200000,0.200000}%
\pgfsetstrokecolor{currentstroke}%
\pgfsetdash{}{0pt}%
\pgfpathmoveto{\pgfqpoint{2.690266in}{2.248532in}}%
\pgfpathlineto{\pgfqpoint{2.738439in}{2.201692in}}%
\pgfpathlineto{\pgfqpoint{2.772705in}{2.179753in}}%
\pgfpathlineto{\pgfqpoint{2.724530in}{2.226620in}}%
\pgfpathlineto{\pgfqpoint{2.690266in}{2.248532in}}%
\pgfpathclose%
\pgfusepath{stroke,fill}%
\end{pgfscope}%
\begin{pgfscope}%
\pgfpathrectangle{\pgfqpoint{0.050000in}{0.050000in}}{\pgfqpoint{4.900000in}{4.900000in}}%
\pgfusepath{clip}%
\pgfsetbuttcap%
\pgfsetroundjoin%
\definecolor{currentfill}{rgb}{0.045521,0.045521,0.045521}%
\pgfsetfillcolor{currentfill}%
\pgfsetfillopacity{0.900000}%
\pgfsetlinewidth{0.501875pt}%
\definecolor{currentstroke}{rgb}{0.200000,0.200000,0.200000}%
\pgfsetstrokecolor{currentstroke}%
\pgfsetdash{}{0pt}%
\pgfpathmoveto{\pgfqpoint{1.645937in}{3.093227in}}%
\pgfpathlineto{\pgfqpoint{1.696467in}{3.044477in}}%
\pgfpathlineto{\pgfqpoint{1.729021in}{3.023919in}}%
\pgfpathlineto{\pgfqpoint{1.678482in}{3.072703in}}%
\pgfpathlineto{\pgfqpoint{1.645937in}{3.093227in}}%
\pgfpathclose%
\pgfusepath{stroke,fill}%
\end{pgfscope}%
\begin{pgfscope}%
\pgfpathrectangle{\pgfqpoint{0.050000in}{0.050000in}}{\pgfqpoint{4.900000in}{4.900000in}}%
\pgfusepath{clip}%
\pgfsetbuttcap%
\pgfsetroundjoin%
\definecolor{currentfill}{rgb}{0.346943,0.346943,0.346943}%
\pgfsetfillcolor{currentfill}%
\pgfsetfillopacity{0.900000}%
\pgfsetlinewidth{0.501875pt}%
\definecolor{currentstroke}{rgb}{0.200000,0.200000,0.200000}%
\pgfsetstrokecolor{currentstroke}%
\pgfsetdash{}{0pt}%
\pgfpathmoveto{\pgfqpoint{2.408539in}{2.476925in}}%
\pgfpathlineto{\pgfqpoint{2.457355in}{2.429560in}}%
\pgfpathlineto{\pgfqpoint{2.491164in}{2.407990in}}%
\pgfpathlineto{\pgfqpoint{2.442344in}{2.455385in}}%
\pgfpathlineto{\pgfqpoint{2.408539in}{2.476925in}}%
\pgfpathclose%
\pgfusepath{stroke,fill}%
\end{pgfscope}%
\begin{pgfscope}%
\pgfpathrectangle{\pgfqpoint{0.050000in}{0.050000in}}{\pgfqpoint{4.900000in}{4.900000in}}%
\pgfusepath{clip}%
\pgfsetbuttcap%
\pgfsetroundjoin%
\definecolor{currentfill}{rgb}{0.234371,0.234371,0.234371}%
\pgfsetfillcolor{currentfill}%
\pgfsetfillopacity{0.900000}%
\pgfsetlinewidth{0.501875pt}%
\definecolor{currentstroke}{rgb}{0.200000,0.200000,0.200000}%
\pgfsetstrokecolor{currentstroke}%
\pgfsetdash{}{0pt}%
\pgfpathmoveto{\pgfqpoint{2.126711in}{2.705405in}}%
\pgfpathlineto{\pgfqpoint{2.176170in}{2.657517in}}%
\pgfpathlineto{\pgfqpoint{2.209523in}{2.636318in}}%
\pgfpathlineto{\pgfqpoint{2.160057in}{2.684237in}}%
\pgfpathlineto{\pgfqpoint{2.126711in}{2.705405in}}%
\pgfpathclose%
\pgfusepath{stroke,fill}%
\end{pgfscope}%
\begin{pgfscope}%
\pgfpathrectangle{\pgfqpoint{0.050000in}{0.050000in}}{\pgfqpoint{4.900000in}{4.900000in}}%
\pgfusepath{clip}%
\pgfsetbuttcap%
\pgfsetroundjoin%
\definecolor{currentfill}{rgb}{0.113802,0.113802,0.113802}%
\pgfsetfillcolor{currentfill}%
\pgfsetfillopacity{0.900000}%
\pgfsetlinewidth{0.501875pt}%
\definecolor{currentstroke}{rgb}{0.200000,0.200000,0.200000}%
\pgfsetstrokecolor{currentstroke}%
\pgfsetdash{}{0pt}%
\pgfpathmoveto{\pgfqpoint{1.844782in}{2.933970in}}%
\pgfpathlineto{\pgfqpoint{1.894885in}{2.885558in}}%
\pgfpathlineto{\pgfqpoint{1.927781in}{2.864730in}}%
\pgfpathlineto{\pgfqpoint{1.877669in}{2.913174in}}%
\pgfpathlineto{\pgfqpoint{1.844782in}{2.933970in}}%
\pgfpathclose%
\pgfusepath{stroke,fill}%
\end{pgfscope}%
\begin{pgfscope}%
\pgfpathrectangle{\pgfqpoint{0.050000in}{0.050000in}}{\pgfqpoint{4.900000in}{4.900000in}}%
\pgfusepath{clip}%
\pgfsetbuttcap%
\pgfsetroundjoin%
\definecolor{currentfill}{rgb}{0.411903,0.411903,0.411903}%
\pgfsetfillcolor{currentfill}%
\pgfsetfillopacity{0.900000}%
\pgfsetlinewidth{0.501875pt}%
\definecolor{currentstroke}{rgb}{0.200000,0.200000,0.200000}%
\pgfsetstrokecolor{currentstroke}%
\pgfsetdash{}{0pt}%
\pgfpathmoveto{\pgfqpoint{2.607726in}{2.317398in}}%
\pgfpathlineto{\pgfqpoint{2.656114in}{2.270372in}}%
\pgfpathlineto{\pgfqpoint{2.690266in}{2.248532in}}%
\pgfpathlineto{\pgfqpoint{2.641874in}{2.295585in}}%
\pgfpathlineto{\pgfqpoint{2.607726in}{2.317398in}}%
\pgfpathclose%
\pgfusepath{stroke,fill}%
\end{pgfscope}%
\begin{pgfscope}%
\pgfpathrectangle{\pgfqpoint{0.050000in}{0.050000in}}{\pgfqpoint{4.900000in}{4.900000in}}%
\pgfusepath{clip}%
\pgfsetbuttcap%
\pgfsetroundjoin%
\definecolor{currentfill}{rgb}{0.009104,0.009104,0.009104}%
\pgfsetfillcolor{currentfill}%
\pgfsetfillopacity{0.900000}%
\pgfsetlinewidth{0.501875pt}%
\definecolor{currentstroke}{rgb}{0.200000,0.200000,0.200000}%
\pgfsetstrokecolor{currentstroke}%
\pgfsetdash{}{0pt}%
\pgfpathmoveto{\pgfqpoint{1.562753in}{3.162618in}}%
\pgfpathlineto{\pgfqpoint{1.613500in}{3.113683in}}%
\pgfpathlineto{\pgfqpoint{1.645937in}{3.093227in}}%
\pgfpathlineto{\pgfqpoint{1.595180in}{3.142196in}}%
\pgfpathlineto{\pgfqpoint{1.562753in}{3.162618in}}%
\pgfpathclose%
\pgfusepath{stroke,fill}%
\end{pgfscope}%
\begin{pgfscope}%
\pgfpathrectangle{\pgfqpoint{0.050000in}{0.050000in}}{\pgfqpoint{4.900000in}{4.900000in}}%
\pgfusepath{clip}%
\pgfsetbuttcap%
\pgfsetroundjoin%
\definecolor{currentfill}{rgb}{0.317416,0.317416,0.317416}%
\pgfsetfillcolor{currentfill}%
\pgfsetfillopacity{0.900000}%
\pgfsetlinewidth{0.501875pt}%
\definecolor{currentstroke}{rgb}{0.200000,0.200000,0.200000}%
\pgfsetstrokecolor{currentstroke}%
\pgfsetdash{}{0pt}%
\pgfpathmoveto{\pgfqpoint{2.325812in}{2.545945in}}%
\pgfpathlineto{\pgfqpoint{2.374845in}{2.498395in}}%
\pgfpathlineto{\pgfqpoint{2.408539in}{2.476925in}}%
\pgfpathlineto{\pgfqpoint{2.359502in}{2.524505in}}%
\pgfpathlineto{\pgfqpoint{2.325812in}{2.545945in}}%
\pgfpathclose%
\pgfusepath{stroke,fill}%
\end{pgfscope}%
\begin{pgfscope}%
\pgfpathrectangle{\pgfqpoint{0.050000in}{0.050000in}}{\pgfqpoint{4.900000in}{4.900000in}}%
\pgfusepath{clip}%
\pgfsetbuttcap%
\pgfsetroundjoin%
\definecolor{currentfill}{rgb}{0.190081,0.190081,0.190081}%
\pgfsetfillcolor{currentfill}%
\pgfsetfillopacity{0.900000}%
\pgfsetlinewidth{0.501875pt}%
\definecolor{currentstroke}{rgb}{0.200000,0.200000,0.200000}%
\pgfsetstrokecolor{currentstroke}%
\pgfsetdash{}{0pt}%
\pgfpathmoveto{\pgfqpoint{2.043798in}{2.774577in}}%
\pgfpathlineto{\pgfqpoint{2.093474in}{2.726504in}}%
\pgfpathlineto{\pgfqpoint{2.126711in}{2.705405in}}%
\pgfpathlineto{\pgfqpoint{2.077028in}{2.753510in}}%
\pgfpathlineto{\pgfqpoint{2.043798in}{2.774577in}}%
\pgfpathclose%
\pgfusepath{stroke,fill}%
\end{pgfscope}%
\begin{pgfscope}%
\pgfpathrectangle{\pgfqpoint{0.050000in}{0.050000in}}{\pgfqpoint{4.900000in}{4.900000in}}%
\pgfusepath{clip}%
\pgfsetbuttcap%
\pgfsetroundjoin%
\definecolor{currentfill}{rgb}{0.474125,0.474125,0.474125}%
\pgfsetfillcolor{currentfill}%
\pgfsetfillopacity{0.900000}%
\pgfsetlinewidth{0.501875pt}%
\definecolor{currentstroke}{rgb}{0.200000,0.200000,0.200000}%
\pgfsetstrokecolor{currentstroke}%
\pgfsetdash{}{0pt}%
\pgfpathmoveto{\pgfqpoint{2.807084in}{2.157742in}}%
\pgfpathlineto{\pgfqpoint{2.855044in}{2.111066in}}%
\pgfpathlineto{\pgfqpoint{2.889538in}{2.088957in}}%
\pgfpathlineto{\pgfqpoint{2.841577in}{2.135658in}}%
\pgfpathlineto{\pgfqpoint{2.807084in}{2.157742in}}%
\pgfpathclose%
\pgfusepath{stroke,fill}%
\end{pgfscope}%
\begin{pgfscope}%
\pgfpathrectangle{\pgfqpoint{0.050000in}{0.050000in}}{\pgfqpoint{4.900000in}{4.900000in}}%
\pgfusepath{clip}%
\pgfsetbuttcap%
\pgfsetroundjoin%
\definecolor{currentfill}{rgb}{0.081938,0.081938,0.081938}%
\pgfsetfillcolor{currentfill}%
\pgfsetfillopacity{0.900000}%
\pgfsetlinewidth{0.501875pt}%
\definecolor{currentstroke}{rgb}{0.200000,0.200000,0.200000}%
\pgfsetstrokecolor{currentstroke}%
\pgfsetdash{}{0pt}%
\pgfpathmoveto{\pgfqpoint{1.761683in}{3.003294in}}%
\pgfpathlineto{\pgfqpoint{1.812004in}{2.954696in}}%
\pgfpathlineto{\pgfqpoint{1.844782in}{2.933970in}}%
\pgfpathlineto{\pgfqpoint{1.794453in}{2.982600in}}%
\pgfpathlineto{\pgfqpoint{1.761683in}{3.003294in}}%
\pgfpathclose%
\pgfusepath{stroke,fill}%
\end{pgfscope}%
\begin{pgfscope}%
\pgfpathrectangle{\pgfqpoint{0.050000in}{0.050000in}}{\pgfqpoint{4.900000in}{4.900000in}}%
\pgfusepath{clip}%
\pgfsetbuttcap%
\pgfsetroundjoin%
\definecolor{currentfill}{rgb}{0.379423,0.379423,0.379423}%
\pgfsetfillcolor{currentfill}%
\pgfsetfillopacity{0.900000}%
\pgfsetlinewidth{0.501875pt}%
\definecolor{currentstroke}{rgb}{0.200000,0.200000,0.200000}%
\pgfsetstrokecolor{currentstroke}%
\pgfsetdash{}{0pt}%
\pgfpathmoveto{\pgfqpoint{2.525085in}{2.386349in}}%
\pgfpathlineto{\pgfqpoint{2.573689in}{2.339139in}}%
\pgfpathlineto{\pgfqpoint{2.607726in}{2.317398in}}%
\pgfpathlineto{\pgfqpoint{2.559118in}{2.364637in}}%
\pgfpathlineto{\pgfqpoint{2.525085in}{2.386349in}}%
\pgfpathclose%
\pgfusepath{stroke,fill}%
\end{pgfscope}%
\begin{pgfscope}%
\pgfpathrectangle{\pgfqpoint{0.050000in}{0.050000in}}{\pgfqpoint{4.900000in}{4.900000in}}%
\pgfusepath{clip}%
\pgfsetbuttcap%
\pgfsetroundjoin%
\definecolor{currentfill}{rgb}{0.273126,0.273126,0.273126}%
\pgfsetfillcolor{currentfill}%
\pgfsetfillopacity{0.900000}%
\pgfsetlinewidth{0.501875pt}%
\definecolor{currentstroke}{rgb}{0.200000,0.200000,0.200000}%
\pgfsetstrokecolor{currentstroke}%
\pgfsetdash{}{0pt}%
\pgfpathmoveto{\pgfqpoint{2.242985in}{2.615049in}}%
\pgfpathlineto{\pgfqpoint{2.292234in}{2.567314in}}%
\pgfpathlineto{\pgfqpoint{2.325812in}{2.545945in}}%
\pgfpathlineto{\pgfqpoint{2.276559in}{2.593710in}}%
\pgfpathlineto{\pgfqpoint{2.242985in}{2.615049in}}%
\pgfpathclose%
\pgfusepath{stroke,fill}%
\end{pgfscope}%
\begin{pgfscope}%
\pgfpathrectangle{\pgfqpoint{0.050000in}{0.050000in}}{\pgfqpoint{4.900000in}{4.900000in}}%
\pgfusepath{clip}%
\pgfsetbuttcap%
\pgfsetroundjoin%
\definecolor{currentfill}{rgb}{0.151326,0.151326,0.151326}%
\pgfsetfillcolor{currentfill}%
\pgfsetfillopacity{0.900000}%
\pgfsetlinewidth{0.501875pt}%
\definecolor{currentstroke}{rgb}{0.200000,0.200000,0.200000}%
\pgfsetstrokecolor{currentstroke}%
\pgfsetdash{}{0pt}%
\pgfpathmoveto{\pgfqpoint{1.960784in}{2.843833in}}%
\pgfpathlineto{\pgfqpoint{2.010678in}{2.795575in}}%
\pgfpathlineto{\pgfqpoint{2.043798in}{2.774577in}}%
\pgfpathlineto{\pgfqpoint{1.993897in}{2.822868in}}%
\pgfpathlineto{\pgfqpoint{1.960784in}{2.843833in}}%
\pgfpathclose%
\pgfusepath{stroke,fill}%
\end{pgfscope}%
\begin{pgfscope}%
\pgfpathrectangle{\pgfqpoint{0.050000in}{0.050000in}}{\pgfqpoint{4.900000in}{4.900000in}}%
\pgfusepath{clip}%
\pgfsetbuttcap%
\pgfsetroundjoin%
\definecolor{currentfill}{rgb}{0.440323,0.440323,0.440323}%
\pgfsetfillcolor{currentfill}%
\pgfsetfillopacity{0.900000}%
\pgfsetlinewidth{0.501875pt}%
\definecolor{currentstroke}{rgb}{0.200000,0.200000,0.200000}%
\pgfsetstrokecolor{currentstroke}%
\pgfsetdash{}{0pt}%
\pgfpathmoveto{\pgfqpoint{2.724530in}{2.226620in}}%
\pgfpathlineto{\pgfqpoint{2.772705in}{2.179753in}}%
\pgfpathlineto{\pgfqpoint{2.807084in}{2.157742in}}%
\pgfpathlineto{\pgfqpoint{2.758907in}{2.204635in}}%
\pgfpathlineto{\pgfqpoint{2.724530in}{2.226620in}}%
\pgfpathclose%
\pgfusepath{stroke,fill}%
\end{pgfscope}%
\begin{pgfscope}%
\pgfpathrectangle{\pgfqpoint{0.050000in}{0.050000in}}{\pgfqpoint{4.900000in}{4.900000in}}%
\pgfusepath{clip}%
\pgfsetbuttcap%
\pgfsetroundjoin%
\definecolor{currentfill}{rgb}{0.045521,0.045521,0.045521}%
\pgfsetfillcolor{currentfill}%
\pgfsetfillopacity{0.900000}%
\pgfsetlinewidth{0.501875pt}%
\definecolor{currentstroke}{rgb}{0.200000,0.200000,0.200000}%
\pgfsetstrokecolor{currentstroke}%
\pgfsetdash{}{0pt}%
\pgfpathmoveto{\pgfqpoint{1.678482in}{3.072703in}}%
\pgfpathlineto{\pgfqpoint{1.729021in}{3.023919in}}%
\pgfpathlineto{\pgfqpoint{1.761683in}{3.003294in}}%
\pgfpathlineto{\pgfqpoint{1.711134in}{3.052111in}}%
\pgfpathlineto{\pgfqpoint{1.678482in}{3.072703in}}%
\pgfpathclose%
\pgfusepath{stroke,fill}%
\end{pgfscope}%
\begin{pgfscope}%
\pgfpathrectangle{\pgfqpoint{0.050000in}{0.050000in}}{\pgfqpoint{4.900000in}{4.900000in}}%
\pgfusepath{clip}%
\pgfsetbuttcap%
\pgfsetroundjoin%
\definecolor{currentfill}{rgb}{0.351003,0.351003,0.351003}%
\pgfsetfillcolor{currentfill}%
\pgfsetfillopacity{0.900000}%
\pgfsetlinewidth{0.501875pt}%
\definecolor{currentstroke}{rgb}{0.200000,0.200000,0.200000}%
\pgfsetstrokecolor{currentstroke}%
\pgfsetdash{}{0pt}%
\pgfpathmoveto{\pgfqpoint{2.442344in}{2.455385in}}%
\pgfpathlineto{\pgfqpoint{2.491164in}{2.407990in}}%
\pgfpathlineto{\pgfqpoint{2.525085in}{2.386349in}}%
\pgfpathlineto{\pgfqpoint{2.476261in}{2.433773in}}%
\pgfpathlineto{\pgfqpoint{2.442344in}{2.455385in}}%
\pgfpathclose%
\pgfusepath{stroke,fill}%
\end{pgfscope}%
\begin{pgfscope}%
\pgfpathrectangle{\pgfqpoint{0.050000in}{0.050000in}}{\pgfqpoint{4.900000in}{4.900000in}}%
\pgfusepath{clip}%
\pgfsetbuttcap%
\pgfsetroundjoin%
\definecolor{currentfill}{rgb}{0.234371,0.234371,0.234371}%
\pgfsetfillcolor{currentfill}%
\pgfsetfillopacity{0.900000}%
\pgfsetlinewidth{0.501875pt}%
\definecolor{currentstroke}{rgb}{0.200000,0.200000,0.200000}%
\pgfsetstrokecolor{currentstroke}%
\pgfsetdash{}{0pt}%
\pgfpathmoveto{\pgfqpoint{2.160057in}{2.684237in}}%
\pgfpathlineto{\pgfqpoint{2.209523in}{2.636318in}}%
\pgfpathlineto{\pgfqpoint{2.242985in}{2.615049in}}%
\pgfpathlineto{\pgfqpoint{2.193514in}{2.663000in}}%
\pgfpathlineto{\pgfqpoint{2.160057in}{2.684237in}}%
\pgfpathclose%
\pgfusepath{stroke,fill}%
\end{pgfscope}%
\begin{pgfscope}%
\pgfpathrectangle{\pgfqpoint{0.050000in}{0.050000in}}{\pgfqpoint{4.900000in}{4.900000in}}%
\pgfusepath{clip}%
\pgfsetbuttcap%
\pgfsetroundjoin%
\definecolor{currentfill}{rgb}{0.113802,0.113802,0.113802}%
\pgfsetfillcolor{currentfill}%
\pgfsetfillopacity{0.900000}%
\pgfsetlinewidth{0.501875pt}%
\definecolor{currentstroke}{rgb}{0.200000,0.200000,0.200000}%
\pgfsetstrokecolor{currentstroke}%
\pgfsetdash{}{0pt}%
\pgfpathmoveto{\pgfqpoint{1.877669in}{2.913174in}}%
\pgfpathlineto{\pgfqpoint{1.927781in}{2.864730in}}%
\pgfpathlineto{\pgfqpoint{1.960784in}{2.843833in}}%
\pgfpathlineto{\pgfqpoint{1.910665in}{2.892311in}}%
\pgfpathlineto{\pgfqpoint{1.877669in}{2.913174in}}%
\pgfpathclose%
\pgfusepath{stroke,fill}%
\end{pgfscope}%
\begin{pgfscope}%
\pgfpathrectangle{\pgfqpoint{0.050000in}{0.050000in}}{\pgfqpoint{4.900000in}{4.900000in}}%
\pgfusepath{clip}%
\pgfsetbuttcap%
\pgfsetroundjoin%
\definecolor{currentfill}{rgb}{0.411903,0.411903,0.411903}%
\pgfsetfillcolor{currentfill}%
\pgfsetfillopacity{0.900000}%
\pgfsetlinewidth{0.501875pt}%
\definecolor{currentstroke}{rgb}{0.200000,0.200000,0.200000}%
\pgfsetstrokecolor{currentstroke}%
\pgfsetdash{}{0pt}%
\pgfpathmoveto{\pgfqpoint{2.641874in}{2.295585in}}%
\pgfpathlineto{\pgfqpoint{2.690266in}{2.248532in}}%
\pgfpathlineto{\pgfqpoint{2.724530in}{2.226620in}}%
\pgfpathlineto{\pgfqpoint{2.676136in}{2.273701in}}%
\pgfpathlineto{\pgfqpoint{2.641874in}{2.295585in}}%
\pgfpathclose%
\pgfusepath{stroke,fill}%
\end{pgfscope}%
\begin{pgfscope}%
\pgfpathrectangle{\pgfqpoint{0.050000in}{0.050000in}}{\pgfqpoint{4.900000in}{4.900000in}}%
\pgfusepath{clip}%
\pgfsetbuttcap%
\pgfsetroundjoin%
\definecolor{currentfill}{rgb}{0.013656,0.013656,0.013656}%
\pgfsetfillcolor{currentfill}%
\pgfsetfillopacity{0.900000}%
\pgfsetlinewidth{0.501875pt}%
\definecolor{currentstroke}{rgb}{0.200000,0.200000,0.200000}%
\pgfsetstrokecolor{currentstroke}%
\pgfsetdash{}{0pt}%
\pgfpathmoveto{\pgfqpoint{1.595180in}{3.142196in}}%
\pgfpathlineto{\pgfqpoint{1.645937in}{3.093227in}}%
\pgfpathlineto{\pgfqpoint{1.678482in}{3.072703in}}%
\pgfpathlineto{\pgfqpoint{1.627714in}{3.121707in}}%
\pgfpathlineto{\pgfqpoint{1.595180in}{3.142196in}}%
\pgfpathclose%
\pgfusepath{stroke,fill}%
\end{pgfscope}%
\begin{pgfscope}%
\pgfpathrectangle{\pgfqpoint{0.050000in}{0.050000in}}{\pgfqpoint{4.900000in}{4.900000in}}%
\pgfusepath{clip}%
\pgfsetbuttcap%
\pgfsetroundjoin%
\definecolor{currentfill}{rgb}{0.317416,0.317416,0.317416}%
\pgfsetfillcolor{currentfill}%
\pgfsetfillopacity{0.900000}%
\pgfsetlinewidth{0.501875pt}%
\definecolor{currentstroke}{rgb}{0.200000,0.200000,0.200000}%
\pgfsetstrokecolor{currentstroke}%
\pgfsetdash{}{0pt}%
\pgfpathmoveto{\pgfqpoint{2.359502in}{2.524505in}}%
\pgfpathlineto{\pgfqpoint{2.408539in}{2.476925in}}%
\pgfpathlineto{\pgfqpoint{2.442344in}{2.455385in}}%
\pgfpathlineto{\pgfqpoint{2.393303in}{2.502994in}}%
\pgfpathlineto{\pgfqpoint{2.359502in}{2.524505in}}%
\pgfpathclose%
\pgfusepath{stroke,fill}%
\end{pgfscope}%
\begin{pgfscope}%
\pgfpathrectangle{\pgfqpoint{0.050000in}{0.050000in}}{\pgfqpoint{4.900000in}{4.900000in}}%
\pgfusepath{clip}%
\pgfsetbuttcap%
\pgfsetroundjoin%
\definecolor{currentfill}{rgb}{0.190081,0.190081,0.190081}%
\pgfsetfillcolor{currentfill}%
\pgfsetfillopacity{0.900000}%
\pgfsetlinewidth{0.501875pt}%
\definecolor{currentstroke}{rgb}{0.200000,0.200000,0.200000}%
\pgfsetstrokecolor{currentstroke}%
\pgfsetdash{}{0pt}%
\pgfpathmoveto{\pgfqpoint{2.077028in}{2.753510in}}%
\pgfpathlineto{\pgfqpoint{2.126711in}{2.705405in}}%
\pgfpathlineto{\pgfqpoint{2.160057in}{2.684237in}}%
\pgfpathlineto{\pgfqpoint{2.110368in}{2.732374in}}%
\pgfpathlineto{\pgfqpoint{2.077028in}{2.753510in}}%
\pgfpathclose%
\pgfusepath{stroke,fill}%
\end{pgfscope}%
\begin{pgfscope}%
\pgfpathrectangle{\pgfqpoint{0.050000in}{0.050000in}}{\pgfqpoint{4.900000in}{4.900000in}}%
\pgfusepath{clip}%
\pgfsetbuttcap%
\pgfsetroundjoin%
\definecolor{currentfill}{rgb}{0.081938,0.081938,0.081938}%
\pgfsetfillcolor{currentfill}%
\pgfsetfillopacity{0.900000}%
\pgfsetlinewidth{0.501875pt}%
\definecolor{currentstroke}{rgb}{0.200000,0.200000,0.200000}%
\pgfsetstrokecolor{currentstroke}%
\pgfsetdash{}{0pt}%
\pgfpathmoveto{\pgfqpoint{1.794453in}{2.982600in}}%
\pgfpathlineto{\pgfqpoint{1.844782in}{2.933970in}}%
\pgfpathlineto{\pgfqpoint{1.877669in}{2.913174in}}%
\pgfpathlineto{\pgfqpoint{1.827331in}{2.961839in}}%
\pgfpathlineto{\pgfqpoint{1.794453in}{2.982600in}}%
\pgfpathclose%
\pgfusepath{stroke,fill}%
\end{pgfscope}%
\begin{pgfscope}%
\pgfpathrectangle{\pgfqpoint{0.050000in}{0.050000in}}{\pgfqpoint{4.900000in}{4.900000in}}%
\pgfusepath{clip}%
\pgfsetbuttcap%
\pgfsetroundjoin%
\definecolor{currentfill}{rgb}{0.379423,0.379423,0.379423}%
\pgfsetfillcolor{currentfill}%
\pgfsetfillopacity{0.900000}%
\pgfsetlinewidth{0.501875pt}%
\definecolor{currentstroke}{rgb}{0.200000,0.200000,0.200000}%
\pgfsetstrokecolor{currentstroke}%
\pgfsetdash{}{0pt}%
\pgfpathmoveto{\pgfqpoint{2.559118in}{2.364637in}}%
\pgfpathlineto{\pgfqpoint{2.607726in}{2.317398in}}%
\pgfpathlineto{\pgfqpoint{2.641874in}{2.295585in}}%
\pgfpathlineto{\pgfqpoint{2.593264in}{2.342853in}}%
\pgfpathlineto{\pgfqpoint{2.559118in}{2.364637in}}%
\pgfpathclose%
\pgfusepath{stroke,fill}%
\end{pgfscope}%
\begin{pgfscope}%
\pgfpathrectangle{\pgfqpoint{0.050000in}{0.050000in}}{\pgfqpoint{4.900000in}{4.900000in}}%
\pgfusepath{clip}%
\pgfsetbuttcap%
\pgfsetroundjoin%
\definecolor{currentfill}{rgb}{0.278662,0.278662,0.278662}%
\pgfsetfillcolor{currentfill}%
\pgfsetfillopacity{0.900000}%
\pgfsetlinewidth{0.501875pt}%
\definecolor{currentstroke}{rgb}{0.200000,0.200000,0.200000}%
\pgfsetstrokecolor{currentstroke}%
\pgfsetdash{}{0pt}%
\pgfpathmoveto{\pgfqpoint{2.276559in}{2.593710in}}%
\pgfpathlineto{\pgfqpoint{2.325812in}{2.545945in}}%
\pgfpathlineto{\pgfqpoint{2.359502in}{2.524505in}}%
\pgfpathlineto{\pgfqpoint{2.310243in}{2.572301in}}%
\pgfpathlineto{\pgfqpoint{2.276559in}{2.593710in}}%
\pgfpathclose%
\pgfusepath{stroke,fill}%
\end{pgfscope}%
\begin{pgfscope}%
\pgfpathrectangle{\pgfqpoint{0.050000in}{0.050000in}}{\pgfqpoint{4.900000in}{4.900000in}}%
\pgfusepath{clip}%
\pgfsetbuttcap%
\pgfsetroundjoin%
\definecolor{currentfill}{rgb}{0.151326,0.151326,0.151326}%
\pgfsetfillcolor{currentfill}%
\pgfsetfillopacity{0.900000}%
\pgfsetlinewidth{0.501875pt}%
\definecolor{currentstroke}{rgb}{0.200000,0.200000,0.200000}%
\pgfsetstrokecolor{currentstroke}%
\pgfsetdash{}{0pt}%
\pgfpathmoveto{\pgfqpoint{1.993897in}{2.822868in}}%
\pgfpathlineto{\pgfqpoint{2.043798in}{2.774577in}}%
\pgfpathlineto{\pgfqpoint{2.077028in}{2.753510in}}%
\pgfpathlineto{\pgfqpoint{2.027120in}{2.801833in}}%
\pgfpathlineto{\pgfqpoint{1.993897in}{2.822868in}}%
\pgfpathclose%
\pgfusepath{stroke,fill}%
\end{pgfscope}%
\begin{pgfscope}%
\pgfpathrectangle{\pgfqpoint{0.050000in}{0.050000in}}{\pgfqpoint{4.900000in}{4.900000in}}%
\pgfusepath{clip}%
\pgfsetbuttcap%
\pgfsetroundjoin%
\definecolor{currentfill}{rgb}{0.440323,0.440323,0.440323}%
\pgfsetfillcolor{currentfill}%
\pgfsetfillopacity{0.900000}%
\pgfsetlinewidth{0.501875pt}%
\definecolor{currentstroke}{rgb}{0.200000,0.200000,0.200000}%
\pgfsetstrokecolor{currentstroke}%
\pgfsetdash{}{0pt}%
\pgfpathmoveto{\pgfqpoint{2.758907in}{2.204635in}}%
\pgfpathlineto{\pgfqpoint{2.807084in}{2.157742in}}%
\pgfpathlineto{\pgfqpoint{2.841577in}{2.135658in}}%
\pgfpathlineto{\pgfqpoint{2.793398in}{2.182578in}}%
\pgfpathlineto{\pgfqpoint{2.758907in}{2.204635in}}%
\pgfpathclose%
\pgfusepath{stroke,fill}%
\end{pgfscope}%
\begin{pgfscope}%
\pgfpathrectangle{\pgfqpoint{0.050000in}{0.050000in}}{\pgfqpoint{4.900000in}{4.900000in}}%
\pgfusepath{clip}%
\pgfsetbuttcap%
\pgfsetroundjoin%
\definecolor{currentfill}{rgb}{0.045521,0.045521,0.045521}%
\pgfsetfillcolor{currentfill}%
\pgfsetfillopacity{0.900000}%
\pgfsetlinewidth{0.501875pt}%
\definecolor{currentstroke}{rgb}{0.200000,0.200000,0.200000}%
\pgfsetstrokecolor{currentstroke}%
\pgfsetdash{}{0pt}%
\pgfpathmoveto{\pgfqpoint{1.711134in}{3.052111in}}%
\pgfpathlineto{\pgfqpoint{1.761683in}{3.003294in}}%
\pgfpathlineto{\pgfqpoint{1.794453in}{2.982600in}}%
\pgfpathlineto{\pgfqpoint{1.743895in}{3.031452in}}%
\pgfpathlineto{\pgfqpoint{1.711134in}{3.052111in}}%
\pgfpathclose%
\pgfusepath{stroke,fill}%
\end{pgfscope}%
\begin{pgfscope}%
\pgfpathrectangle{\pgfqpoint{0.050000in}{0.050000in}}{\pgfqpoint{4.900000in}{4.900000in}}%
\pgfusepath{clip}%
\pgfsetbuttcap%
\pgfsetroundjoin%
\definecolor{currentfill}{rgb}{0.351003,0.351003,0.351003}%
\pgfsetfillcolor{currentfill}%
\pgfsetfillopacity{0.900000}%
\pgfsetlinewidth{0.501875pt}%
\definecolor{currentstroke}{rgb}{0.200000,0.200000,0.200000}%
\pgfsetstrokecolor{currentstroke}%
\pgfsetdash{}{0pt}%
\pgfpathmoveto{\pgfqpoint{2.476261in}{2.433773in}}%
\pgfpathlineto{\pgfqpoint{2.525085in}{2.386349in}}%
\pgfpathlineto{\pgfqpoint{2.559118in}{2.364637in}}%
\pgfpathlineto{\pgfqpoint{2.510291in}{2.412090in}}%
\pgfpathlineto{\pgfqpoint{2.476261in}{2.433773in}}%
\pgfpathclose%
\pgfusepath{stroke,fill}%
\end{pgfscope}%
\begin{pgfscope}%
\pgfpathrectangle{\pgfqpoint{0.050000in}{0.050000in}}{\pgfqpoint{4.900000in}{4.900000in}}%
\pgfusepath{clip}%
\pgfsetbuttcap%
\pgfsetroundjoin%
\definecolor{currentfill}{rgb}{0.234371,0.234371,0.234371}%
\pgfsetfillcolor{currentfill}%
\pgfsetfillopacity{0.900000}%
\pgfsetlinewidth{0.501875pt}%
\definecolor{currentstroke}{rgb}{0.200000,0.200000,0.200000}%
\pgfsetstrokecolor{currentstroke}%
\pgfsetdash{}{0pt}%
\pgfpathmoveto{\pgfqpoint{2.193514in}{2.663000in}}%
\pgfpathlineto{\pgfqpoint{2.242985in}{2.615049in}}%
\pgfpathlineto{\pgfqpoint{2.276559in}{2.593710in}}%
\pgfpathlineto{\pgfqpoint{2.227081in}{2.641692in}}%
\pgfpathlineto{\pgfqpoint{2.193514in}{2.663000in}}%
\pgfpathclose%
\pgfusepath{stroke,fill}%
\end{pgfscope}%
\begin{pgfscope}%
\pgfpathrectangle{\pgfqpoint{0.050000in}{0.050000in}}{\pgfqpoint{4.900000in}{4.900000in}}%
\pgfusepath{clip}%
\pgfsetbuttcap%
\pgfsetroundjoin%
\definecolor{currentfill}{rgb}{0.113802,0.113802,0.113802}%
\pgfsetfillcolor{currentfill}%
\pgfsetfillopacity{0.900000}%
\pgfsetlinewidth{0.501875pt}%
\definecolor{currentstroke}{rgb}{0.200000,0.200000,0.200000}%
\pgfsetstrokecolor{currentstroke}%
\pgfsetdash{}{0pt}%
\pgfpathmoveto{\pgfqpoint{1.910665in}{2.892311in}}%
\pgfpathlineto{\pgfqpoint{1.960784in}{2.843833in}}%
\pgfpathlineto{\pgfqpoint{1.993897in}{2.822868in}}%
\pgfpathlineto{\pgfqpoint{1.943770in}{2.871378in}}%
\pgfpathlineto{\pgfqpoint{1.910665in}{2.892311in}}%
\pgfpathclose%
\pgfusepath{stroke,fill}%
\end{pgfscope}%
\begin{pgfscope}%
\pgfpathrectangle{\pgfqpoint{0.050000in}{0.050000in}}{\pgfqpoint{4.900000in}{4.900000in}}%
\pgfusepath{clip}%
\pgfsetbuttcap%
\pgfsetroundjoin%
\definecolor{currentfill}{rgb}{0.411903,0.411903,0.411903}%
\pgfsetfillcolor{currentfill}%
\pgfsetfillopacity{0.900000}%
\pgfsetlinewidth{0.501875pt}%
\definecolor{currentstroke}{rgb}{0.200000,0.200000,0.200000}%
\pgfsetstrokecolor{currentstroke}%
\pgfsetdash{}{0pt}%
\pgfpathmoveto{\pgfqpoint{2.676136in}{2.273701in}}%
\pgfpathlineto{\pgfqpoint{2.724530in}{2.226620in}}%
\pgfpathlineto{\pgfqpoint{2.758907in}{2.204635in}}%
\pgfpathlineto{\pgfqpoint{2.710511in}{2.251744in}}%
\pgfpathlineto{\pgfqpoint{2.676136in}{2.273701in}}%
\pgfpathclose%
\pgfusepath{stroke,fill}%
\end{pgfscope}%
\begin{pgfscope}%
\pgfpathrectangle{\pgfqpoint{0.050000in}{0.050000in}}{\pgfqpoint{4.900000in}{4.900000in}}%
\pgfusepath{clip}%
\pgfsetbuttcap%
\pgfsetroundjoin%
\definecolor{currentfill}{rgb}{0.013656,0.013656,0.013656}%
\pgfsetfillcolor{currentfill}%
\pgfsetfillopacity{0.900000}%
\pgfsetlinewidth{0.501875pt}%
\definecolor{currentstroke}{rgb}{0.200000,0.200000,0.200000}%
\pgfsetstrokecolor{currentstroke}%
\pgfsetdash{}{0pt}%
\pgfpathmoveto{\pgfqpoint{1.627714in}{3.121707in}}%
\pgfpathlineto{\pgfqpoint{1.678482in}{3.072703in}}%
\pgfpathlineto{\pgfqpoint{1.711134in}{3.052111in}}%
\pgfpathlineto{\pgfqpoint{1.660357in}{3.101150in}}%
\pgfpathlineto{\pgfqpoint{1.627714in}{3.121707in}}%
\pgfpathclose%
\pgfusepath{stroke,fill}%
\end{pgfscope}%
\begin{pgfscope}%
\pgfpathrectangle{\pgfqpoint{0.050000in}{0.050000in}}{\pgfqpoint{4.900000in}{4.900000in}}%
\pgfusepath{clip}%
\pgfsetbuttcap%
\pgfsetroundjoin%
\definecolor{currentfill}{rgb}{0.317416,0.317416,0.317416}%
\pgfsetfillcolor{currentfill}%
\pgfsetfillopacity{0.900000}%
\pgfsetlinewidth{0.501875pt}%
\definecolor{currentstroke}{rgb}{0.200000,0.200000,0.200000}%
\pgfsetstrokecolor{currentstroke}%
\pgfsetdash{}{0pt}%
\pgfpathmoveto{\pgfqpoint{2.393303in}{2.502994in}}%
\pgfpathlineto{\pgfqpoint{2.442344in}{2.455385in}}%
\pgfpathlineto{\pgfqpoint{2.476261in}{2.433773in}}%
\pgfpathlineto{\pgfqpoint{2.427215in}{2.481413in}}%
\pgfpathlineto{\pgfqpoint{2.393303in}{2.502994in}}%
\pgfpathclose%
\pgfusepath{stroke,fill}%
\end{pgfscope}%
\begin{pgfscope}%
\pgfpathrectangle{\pgfqpoint{0.050000in}{0.050000in}}{\pgfqpoint{4.900000in}{4.900000in}}%
\pgfusepath{clip}%
\pgfsetbuttcap%
\pgfsetroundjoin%
\definecolor{currentfill}{rgb}{0.190081,0.190081,0.190081}%
\pgfsetfillcolor{currentfill}%
\pgfsetfillopacity{0.900000}%
\pgfsetlinewidth{0.501875pt}%
\definecolor{currentstroke}{rgb}{0.200000,0.200000,0.200000}%
\pgfsetstrokecolor{currentstroke}%
\pgfsetdash{}{0pt}%
\pgfpathmoveto{\pgfqpoint{2.110368in}{2.732374in}}%
\pgfpathlineto{\pgfqpoint{2.160057in}{2.684237in}}%
\pgfpathlineto{\pgfqpoint{2.193514in}{2.663000in}}%
\pgfpathlineto{\pgfqpoint{2.143818in}{2.711168in}}%
\pgfpathlineto{\pgfqpoint{2.110368in}{2.732374in}}%
\pgfpathclose%
\pgfusepath{stroke,fill}%
\end{pgfscope}%
\begin{pgfscope}%
\pgfpathrectangle{\pgfqpoint{0.050000in}{0.050000in}}{\pgfqpoint{4.900000in}{4.900000in}}%
\pgfusepath{clip}%
\pgfsetbuttcap%
\pgfsetroundjoin%
\definecolor{currentfill}{rgb}{0.081938,0.081938,0.081938}%
\pgfsetfillcolor{currentfill}%
\pgfsetfillopacity{0.900000}%
\pgfsetlinewidth{0.501875pt}%
\definecolor{currentstroke}{rgb}{0.200000,0.200000,0.200000}%
\pgfsetstrokecolor{currentstroke}%
\pgfsetdash{}{0pt}%
\pgfpathmoveto{\pgfqpoint{1.827331in}{2.961839in}}%
\pgfpathlineto{\pgfqpoint{1.877669in}{2.913174in}}%
\pgfpathlineto{\pgfqpoint{1.910665in}{2.892311in}}%
\pgfpathlineto{\pgfqpoint{1.860318in}{2.941008in}}%
\pgfpathlineto{\pgfqpoint{1.827331in}{2.961839in}}%
\pgfpathclose%
\pgfusepath{stroke,fill}%
\end{pgfscope}%
\begin{pgfscope}%
\pgfpathrectangle{\pgfqpoint{0.050000in}{0.050000in}}{\pgfqpoint{4.900000in}{4.900000in}}%
\pgfusepath{clip}%
\pgfsetbuttcap%
\pgfsetroundjoin%
\definecolor{currentfill}{rgb}{0.379423,0.379423,0.379423}%
\pgfsetfillcolor{currentfill}%
\pgfsetfillopacity{0.900000}%
\pgfsetlinewidth{0.501875pt}%
\definecolor{currentstroke}{rgb}{0.200000,0.200000,0.200000}%
\pgfsetstrokecolor{currentstroke}%
\pgfsetdash{}{0pt}%
\pgfpathmoveto{\pgfqpoint{2.593264in}{2.342853in}}%
\pgfpathlineto{\pgfqpoint{2.641874in}{2.295585in}}%
\pgfpathlineto{\pgfqpoint{2.676136in}{2.273701in}}%
\pgfpathlineto{\pgfqpoint{2.627523in}{2.320997in}}%
\pgfpathlineto{\pgfqpoint{2.593264in}{2.342853in}}%
\pgfpathclose%
\pgfusepath{stroke,fill}%
\end{pgfscope}%
\begin{pgfscope}%
\pgfpathrectangle{\pgfqpoint{0.050000in}{0.050000in}}{\pgfqpoint{4.900000in}{4.900000in}}%
\pgfusepath{clip}%
\pgfsetbuttcap%
\pgfsetroundjoin%
\definecolor{currentfill}{rgb}{0.278662,0.278662,0.278662}%
\pgfsetfillcolor{currentfill}%
\pgfsetfillopacity{0.900000}%
\pgfsetlinewidth{0.501875pt}%
\definecolor{currentstroke}{rgb}{0.200000,0.200000,0.200000}%
\pgfsetstrokecolor{currentstroke}%
\pgfsetdash{}{0pt}%
\pgfpathmoveto{\pgfqpoint{2.310243in}{2.572301in}}%
\pgfpathlineto{\pgfqpoint{2.359502in}{2.524505in}}%
\pgfpathlineto{\pgfqpoint{2.393303in}{2.502994in}}%
\pgfpathlineto{\pgfqpoint{2.344039in}{2.550820in}}%
\pgfpathlineto{\pgfqpoint{2.310243in}{2.572301in}}%
\pgfpathclose%
\pgfusepath{stroke,fill}%
\end{pgfscope}%
\begin{pgfscope}%
\pgfpathrectangle{\pgfqpoint{0.050000in}{0.050000in}}{\pgfqpoint{4.900000in}{4.900000in}}%
\pgfusepath{clip}%
\pgfsetbuttcap%
\pgfsetroundjoin%
\definecolor{currentfill}{rgb}{0.151326,0.151326,0.151326}%
\pgfsetfillcolor{currentfill}%
\pgfsetfillopacity{0.900000}%
\pgfsetlinewidth{0.501875pt}%
\definecolor{currentstroke}{rgb}{0.200000,0.200000,0.200000}%
\pgfsetstrokecolor{currentstroke}%
\pgfsetdash{}{0pt}%
\pgfpathmoveto{\pgfqpoint{2.027120in}{2.801833in}}%
\pgfpathlineto{\pgfqpoint{2.077028in}{2.753510in}}%
\pgfpathlineto{\pgfqpoint{2.110368in}{2.732374in}}%
\pgfpathlineto{\pgfqpoint{2.060452in}{2.780729in}}%
\pgfpathlineto{\pgfqpoint{2.027120in}{2.801833in}}%
\pgfpathclose%
\pgfusepath{stroke,fill}%
\end{pgfscope}%
\begin{pgfscope}%
\pgfpathrectangle{\pgfqpoint{0.050000in}{0.050000in}}{\pgfqpoint{4.900000in}{4.900000in}}%
\pgfusepath{clip}%
\pgfsetbuttcap%
\pgfsetroundjoin%
\definecolor{currentfill}{rgb}{0.045521,0.045521,0.045521}%
\pgfsetfillcolor{currentfill}%
\pgfsetfillopacity{0.900000}%
\pgfsetlinewidth{0.501875pt}%
\definecolor{currentstroke}{rgb}{0.200000,0.200000,0.200000}%
\pgfsetstrokecolor{currentstroke}%
\pgfsetdash{}{0pt}%
\pgfpathmoveto{\pgfqpoint{1.743895in}{3.031452in}}%
\pgfpathlineto{\pgfqpoint{1.794453in}{2.982600in}}%
\pgfpathlineto{\pgfqpoint{1.827331in}{2.961839in}}%
\pgfpathlineto{\pgfqpoint{1.776764in}{3.010724in}}%
\pgfpathlineto{\pgfqpoint{1.743895in}{3.031452in}}%
\pgfpathclose%
\pgfusepath{stroke,fill}%
\end{pgfscope}%
\begin{pgfscope}%
\pgfpathrectangle{\pgfqpoint{0.050000in}{0.050000in}}{\pgfqpoint{4.900000in}{4.900000in}}%
\pgfusepath{clip}%
\pgfsetbuttcap%
\pgfsetroundjoin%
\definecolor{currentfill}{rgb}{0.351003,0.351003,0.351003}%
\pgfsetfillcolor{currentfill}%
\pgfsetfillopacity{0.900000}%
\pgfsetlinewidth{0.501875pt}%
\definecolor{currentstroke}{rgb}{0.200000,0.200000,0.200000}%
\pgfsetstrokecolor{currentstroke}%
\pgfsetdash{}{0pt}%
\pgfpathmoveto{\pgfqpoint{2.510291in}{2.412090in}}%
\pgfpathlineto{\pgfqpoint{2.559118in}{2.364637in}}%
\pgfpathlineto{\pgfqpoint{2.593264in}{2.342853in}}%
\pgfpathlineto{\pgfqpoint{2.544433in}{2.390336in}}%
\pgfpathlineto{\pgfqpoint{2.510291in}{2.412090in}}%
\pgfpathclose%
\pgfusepath{stroke,fill}%
\end{pgfscope}%
\begin{pgfscope}%
\pgfpathrectangle{\pgfqpoint{0.050000in}{0.050000in}}{\pgfqpoint{4.900000in}{4.900000in}}%
\pgfusepath{clip}%
\pgfsetbuttcap%
\pgfsetroundjoin%
\definecolor{currentfill}{rgb}{0.234371,0.234371,0.234371}%
\pgfsetfillcolor{currentfill}%
\pgfsetfillopacity{0.900000}%
\pgfsetlinewidth{0.501875pt}%
\definecolor{currentstroke}{rgb}{0.200000,0.200000,0.200000}%
\pgfsetstrokecolor{currentstroke}%
\pgfsetdash{}{0pt}%
\pgfpathmoveto{\pgfqpoint{2.227081in}{2.641692in}}%
\pgfpathlineto{\pgfqpoint{2.276559in}{2.593710in}}%
\pgfpathlineto{\pgfqpoint{2.310243in}{2.572301in}}%
\pgfpathlineto{\pgfqpoint{2.260760in}{2.620313in}}%
\pgfpathlineto{\pgfqpoint{2.227081in}{2.641692in}}%
\pgfpathclose%
\pgfusepath{stroke,fill}%
\end{pgfscope}%
\begin{pgfscope}%
\pgfpathrectangle{\pgfqpoint{0.050000in}{0.050000in}}{\pgfqpoint{4.900000in}{4.900000in}}%
\pgfusepath{clip}%
\pgfsetbuttcap%
\pgfsetroundjoin%
\definecolor{currentfill}{rgb}{0.113802,0.113802,0.113802}%
\pgfsetfillcolor{currentfill}%
\pgfsetfillopacity{0.900000}%
\pgfsetlinewidth{0.501875pt}%
\definecolor{currentstroke}{rgb}{0.200000,0.200000,0.200000}%
\pgfsetstrokecolor{currentstroke}%
\pgfsetdash{}{0pt}%
\pgfpathmoveto{\pgfqpoint{1.943770in}{2.871378in}}%
\pgfpathlineto{\pgfqpoint{1.993897in}{2.822868in}}%
\pgfpathlineto{\pgfqpoint{2.027120in}{2.801833in}}%
\pgfpathlineto{\pgfqpoint{1.976985in}{2.850376in}}%
\pgfpathlineto{\pgfqpoint{1.943770in}{2.871378in}}%
\pgfpathclose%
\pgfusepath{stroke,fill}%
\end{pgfscope}%
\begin{pgfscope}%
\pgfpathrectangle{\pgfqpoint{0.050000in}{0.050000in}}{\pgfqpoint{4.900000in}{4.900000in}}%
\pgfusepath{clip}%
\pgfsetbuttcap%
\pgfsetroundjoin%
\definecolor{currentfill}{rgb}{0.411903,0.411903,0.411903}%
\pgfsetfillcolor{currentfill}%
\pgfsetfillopacity{0.900000}%
\pgfsetlinewidth{0.501875pt}%
\definecolor{currentstroke}{rgb}{0.200000,0.200000,0.200000}%
\pgfsetstrokecolor{currentstroke}%
\pgfsetdash{}{0pt}%
\pgfpathmoveto{\pgfqpoint{2.710511in}{2.251744in}}%
\pgfpathlineto{\pgfqpoint{2.758907in}{2.204635in}}%
\pgfpathlineto{\pgfqpoint{2.793398in}{2.182578in}}%
\pgfpathlineto{\pgfqpoint{2.745000in}{2.229715in}}%
\pgfpathlineto{\pgfqpoint{2.710511in}{2.251744in}}%
\pgfpathclose%
\pgfusepath{stroke,fill}%
\end{pgfscope}%
\begin{pgfscope}%
\pgfpathrectangle{\pgfqpoint{0.050000in}{0.050000in}}{\pgfqpoint{4.900000in}{4.900000in}}%
\pgfusepath{clip}%
\pgfsetbuttcap%
\pgfsetroundjoin%
\definecolor{currentfill}{rgb}{0.013656,0.013656,0.013656}%
\pgfsetfillcolor{currentfill}%
\pgfsetfillopacity{0.900000}%
\pgfsetlinewidth{0.501875pt}%
\definecolor{currentstroke}{rgb}{0.200000,0.200000,0.200000}%
\pgfsetstrokecolor{currentstroke}%
\pgfsetdash{}{0pt}%
\pgfpathmoveto{\pgfqpoint{1.660357in}{3.101150in}}%
\pgfpathlineto{\pgfqpoint{1.711134in}{3.052111in}}%
\pgfpathlineto{\pgfqpoint{1.743895in}{3.031452in}}%
\pgfpathlineto{\pgfqpoint{1.693107in}{3.080525in}}%
\pgfpathlineto{\pgfqpoint{1.660357in}{3.101150in}}%
\pgfpathclose%
\pgfusepath{stroke,fill}%
\end{pgfscope}%
\begin{pgfscope}%
\pgfpathrectangle{\pgfqpoint{0.050000in}{0.050000in}}{\pgfqpoint{4.900000in}{4.900000in}}%
\pgfusepath{clip}%
\pgfsetbuttcap%
\pgfsetroundjoin%
\definecolor{currentfill}{rgb}{0.317416,0.317416,0.317416}%
\pgfsetfillcolor{currentfill}%
\pgfsetfillopacity{0.900000}%
\pgfsetlinewidth{0.501875pt}%
\definecolor{currentstroke}{rgb}{0.200000,0.200000,0.200000}%
\pgfsetstrokecolor{currentstroke}%
\pgfsetdash{}{0pt}%
\pgfpathmoveto{\pgfqpoint{2.427215in}{2.481413in}}%
\pgfpathlineto{\pgfqpoint{2.476261in}{2.433773in}}%
\pgfpathlineto{\pgfqpoint{2.510291in}{2.412090in}}%
\pgfpathlineto{\pgfqpoint{2.461241in}{2.459760in}}%
\pgfpathlineto{\pgfqpoint{2.427215in}{2.481413in}}%
\pgfpathclose%
\pgfusepath{stroke,fill}%
\end{pgfscope}%
\begin{pgfscope}%
\pgfpathrectangle{\pgfqpoint{0.050000in}{0.050000in}}{\pgfqpoint{4.900000in}{4.900000in}}%
\pgfusepath{clip}%
\pgfsetbuttcap%
\pgfsetroundjoin%
\definecolor{currentfill}{rgb}{0.195617,0.195617,0.195617}%
\pgfsetfillcolor{currentfill}%
\pgfsetfillopacity{0.900000}%
\pgfsetlinewidth{0.501875pt}%
\definecolor{currentstroke}{rgb}{0.200000,0.200000,0.200000}%
\pgfsetstrokecolor{currentstroke}%
\pgfsetdash{}{0pt}%
\pgfpathmoveto{\pgfqpoint{2.143818in}{2.711168in}}%
\pgfpathlineto{\pgfqpoint{2.193514in}{2.663000in}}%
\pgfpathlineto{\pgfqpoint{2.227081in}{2.641692in}}%
\pgfpathlineto{\pgfqpoint{2.177379in}{2.689891in}}%
\pgfpathlineto{\pgfqpoint{2.143818in}{2.711168in}}%
\pgfpathclose%
\pgfusepath{stroke,fill}%
\end{pgfscope}%
\begin{pgfscope}%
\pgfpathrectangle{\pgfqpoint{0.050000in}{0.050000in}}{\pgfqpoint{4.900000in}{4.900000in}}%
\pgfusepath{clip}%
\pgfsetbuttcap%
\pgfsetroundjoin%
\definecolor{currentfill}{rgb}{0.081938,0.081938,0.081938}%
\pgfsetfillcolor{currentfill}%
\pgfsetfillopacity{0.900000}%
\pgfsetlinewidth{0.501875pt}%
\definecolor{currentstroke}{rgb}{0.200000,0.200000,0.200000}%
\pgfsetstrokecolor{currentstroke}%
\pgfsetdash{}{0pt}%
\pgfpathmoveto{\pgfqpoint{1.860318in}{2.941008in}}%
\pgfpathlineto{\pgfqpoint{1.910665in}{2.892311in}}%
\pgfpathlineto{\pgfqpoint{1.943770in}{2.871378in}}%
\pgfpathlineto{\pgfqpoint{1.893415in}{2.920109in}}%
\pgfpathlineto{\pgfqpoint{1.860318in}{2.941008in}}%
\pgfpathclose%
\pgfusepath{stroke,fill}%
\end{pgfscope}%
\begin{pgfscope}%
\pgfpathrectangle{\pgfqpoint{0.050000in}{0.050000in}}{\pgfqpoint{4.900000in}{4.900000in}}%
\pgfusepath{clip}%
\pgfsetbuttcap%
\pgfsetroundjoin%
\definecolor{currentfill}{rgb}{0.379423,0.379423,0.379423}%
\pgfsetfillcolor{currentfill}%
\pgfsetfillopacity{0.900000}%
\pgfsetlinewidth{0.501875pt}%
\definecolor{currentstroke}{rgb}{0.200000,0.200000,0.200000}%
\pgfsetstrokecolor{currentstroke}%
\pgfsetdash{}{0pt}%
\pgfpathmoveto{\pgfqpoint{2.627523in}{2.320997in}}%
\pgfpathlineto{\pgfqpoint{2.676136in}{2.273701in}}%
\pgfpathlineto{\pgfqpoint{2.710511in}{2.251744in}}%
\pgfpathlineto{\pgfqpoint{2.661895in}{2.299069in}}%
\pgfpathlineto{\pgfqpoint{2.627523in}{2.320997in}}%
\pgfpathclose%
\pgfusepath{stroke,fill}%
\end{pgfscope}%
\begin{pgfscope}%
\pgfpathrectangle{\pgfqpoint{0.050000in}{0.050000in}}{\pgfqpoint{4.900000in}{4.900000in}}%
\pgfusepath{clip}%
\pgfsetbuttcap%
\pgfsetroundjoin%
\definecolor{currentfill}{rgb}{0.278662,0.278662,0.278662}%
\pgfsetfillcolor{currentfill}%
\pgfsetfillopacity{0.900000}%
\pgfsetlinewidth{0.501875pt}%
\definecolor{currentstroke}{rgb}{0.200000,0.200000,0.200000}%
\pgfsetstrokecolor{currentstroke}%
\pgfsetdash{}{0pt}%
\pgfpathmoveto{\pgfqpoint{2.344039in}{2.550820in}}%
\pgfpathlineto{\pgfqpoint{2.393303in}{2.502994in}}%
\pgfpathlineto{\pgfqpoint{2.427215in}{2.481413in}}%
\pgfpathlineto{\pgfqpoint{2.377947in}{2.529269in}}%
\pgfpathlineto{\pgfqpoint{2.344039in}{2.550820in}}%
\pgfpathclose%
\pgfusepath{stroke,fill}%
\end{pgfscope}%
\begin{pgfscope}%
\pgfpathrectangle{\pgfqpoint{0.050000in}{0.050000in}}{\pgfqpoint{4.900000in}{4.900000in}}%
\pgfusepath{clip}%
\pgfsetbuttcap%
\pgfsetroundjoin%
\definecolor{currentfill}{rgb}{0.151326,0.151326,0.151326}%
\pgfsetfillcolor{currentfill}%
\pgfsetfillopacity{0.900000}%
\pgfsetlinewidth{0.501875pt}%
\definecolor{currentstroke}{rgb}{0.200000,0.200000,0.200000}%
\pgfsetstrokecolor{currentstroke}%
\pgfsetdash{}{0pt}%
\pgfpathmoveto{\pgfqpoint{2.060452in}{2.780729in}}%
\pgfpathlineto{\pgfqpoint{2.110368in}{2.732374in}}%
\pgfpathlineto{\pgfqpoint{2.143818in}{2.711168in}}%
\pgfpathlineto{\pgfqpoint{2.093896in}{2.759555in}}%
\pgfpathlineto{\pgfqpoint{2.060452in}{2.780729in}}%
\pgfpathclose%
\pgfusepath{stroke,fill}%
\end{pgfscope}%
\begin{pgfscope}%
\pgfpathrectangle{\pgfqpoint{0.050000in}{0.050000in}}{\pgfqpoint{4.900000in}{4.900000in}}%
\pgfusepath{clip}%
\pgfsetbuttcap%
\pgfsetroundjoin%
\definecolor{currentfill}{rgb}{0.045521,0.045521,0.045521}%
\pgfsetfillcolor{currentfill}%
\pgfsetfillopacity{0.900000}%
\pgfsetlinewidth{0.501875pt}%
\definecolor{currentstroke}{rgb}{0.200000,0.200000,0.200000}%
\pgfsetstrokecolor{currentstroke}%
\pgfsetdash{}{0pt}%
\pgfpathmoveto{\pgfqpoint{1.776764in}{3.010724in}}%
\pgfpathlineto{\pgfqpoint{1.827331in}{2.961839in}}%
\pgfpathlineto{\pgfqpoint{1.860318in}{2.941008in}}%
\pgfpathlineto{\pgfqpoint{1.809742in}{2.989927in}}%
\pgfpathlineto{\pgfqpoint{1.776764in}{3.010724in}}%
\pgfpathclose%
\pgfusepath{stroke,fill}%
\end{pgfscope}%
\begin{pgfscope}%
\pgfpathrectangle{\pgfqpoint{0.050000in}{0.050000in}}{\pgfqpoint{4.900000in}{4.900000in}}%
\pgfusepath{clip}%
\pgfsetbuttcap%
\pgfsetroundjoin%
\definecolor{currentfill}{rgb}{0.351003,0.351003,0.351003}%
\pgfsetfillcolor{currentfill}%
\pgfsetfillopacity{0.900000}%
\pgfsetlinewidth{0.501875pt}%
\definecolor{currentstroke}{rgb}{0.200000,0.200000,0.200000}%
\pgfsetstrokecolor{currentstroke}%
\pgfsetdash{}{0pt}%
\pgfpathmoveto{\pgfqpoint{2.544433in}{2.390336in}}%
\pgfpathlineto{\pgfqpoint{2.593264in}{2.342853in}}%
\pgfpathlineto{\pgfqpoint{2.627523in}{2.320997in}}%
\pgfpathlineto{\pgfqpoint{2.578688in}{2.368509in}}%
\pgfpathlineto{\pgfqpoint{2.544433in}{2.390336in}}%
\pgfpathclose%
\pgfusepath{stroke,fill}%
\end{pgfscope}%
\begin{pgfscope}%
\pgfpathrectangle{\pgfqpoint{0.050000in}{0.050000in}}{\pgfqpoint{4.900000in}{4.900000in}}%
\pgfusepath{clip}%
\pgfsetbuttcap%
\pgfsetroundjoin%
\definecolor{currentfill}{rgb}{0.234371,0.234371,0.234371}%
\pgfsetfillcolor{currentfill}%
\pgfsetfillopacity{0.900000}%
\pgfsetlinewidth{0.501875pt}%
\definecolor{currentstroke}{rgb}{0.200000,0.200000,0.200000}%
\pgfsetstrokecolor{currentstroke}%
\pgfsetdash{}{0pt}%
\pgfpathmoveto{\pgfqpoint{2.260760in}{2.620313in}}%
\pgfpathlineto{\pgfqpoint{2.310243in}{2.572301in}}%
\pgfpathlineto{\pgfqpoint{2.344039in}{2.550820in}}%
\pgfpathlineto{\pgfqpoint{2.294550in}{2.598864in}}%
\pgfpathlineto{\pgfqpoint{2.260760in}{2.620313in}}%
\pgfpathclose%
\pgfusepath{stroke,fill}%
\end{pgfscope}%
\begin{pgfscope}%
\pgfpathrectangle{\pgfqpoint{0.050000in}{0.050000in}}{\pgfqpoint{4.900000in}{4.900000in}}%
\pgfusepath{clip}%
\pgfsetbuttcap%
\pgfsetroundjoin%
\definecolor{currentfill}{rgb}{0.118354,0.118354,0.118354}%
\pgfsetfillcolor{currentfill}%
\pgfsetfillopacity{0.900000}%
\pgfsetlinewidth{0.501875pt}%
\definecolor{currentstroke}{rgb}{0.200000,0.200000,0.200000}%
\pgfsetstrokecolor{currentstroke}%
\pgfsetdash{}{0pt}%
\pgfpathmoveto{\pgfqpoint{1.976985in}{2.850376in}}%
\pgfpathlineto{\pgfqpoint{2.027120in}{2.801833in}}%
\pgfpathlineto{\pgfqpoint{2.060452in}{2.780729in}}%
\pgfpathlineto{\pgfqpoint{2.010310in}{2.829305in}}%
\pgfpathlineto{\pgfqpoint{1.976985in}{2.850376in}}%
\pgfpathclose%
\pgfusepath{stroke,fill}%
\end{pgfscope}%
\begin{pgfscope}%
\pgfpathrectangle{\pgfqpoint{0.050000in}{0.050000in}}{\pgfqpoint{4.900000in}{4.900000in}}%
\pgfusepath{clip}%
\pgfsetbuttcap%
\pgfsetroundjoin%
\definecolor{currentfill}{rgb}{0.013656,0.013656,0.013656}%
\pgfsetfillcolor{currentfill}%
\pgfsetfillopacity{0.900000}%
\pgfsetlinewidth{0.501875pt}%
\definecolor{currentstroke}{rgb}{0.200000,0.200000,0.200000}%
\pgfsetstrokecolor{currentstroke}%
\pgfsetdash{}{0pt}%
\pgfpathmoveto{\pgfqpoint{1.693107in}{3.080525in}}%
\pgfpathlineto{\pgfqpoint{1.743895in}{3.031452in}}%
\pgfpathlineto{\pgfqpoint{1.776764in}{3.010724in}}%
\pgfpathlineto{\pgfqpoint{1.725967in}{3.059831in}}%
\pgfpathlineto{\pgfqpoint{1.693107in}{3.080525in}}%
\pgfpathclose%
\pgfusepath{stroke,fill}%
\end{pgfscope}%
\begin{pgfscope}%
\pgfpathrectangle{\pgfqpoint{0.050000in}{0.050000in}}{\pgfqpoint{4.900000in}{4.900000in}}%
\pgfusepath{clip}%
\pgfsetbuttcap%
\pgfsetroundjoin%
\definecolor{currentfill}{rgb}{0.317416,0.317416,0.317416}%
\pgfsetfillcolor{currentfill}%
\pgfsetfillopacity{0.900000}%
\pgfsetlinewidth{0.501875pt}%
\definecolor{currentstroke}{rgb}{0.200000,0.200000,0.200000}%
\pgfsetstrokecolor{currentstroke}%
\pgfsetdash{}{0pt}%
\pgfpathmoveto{\pgfqpoint{2.461241in}{2.459760in}}%
\pgfpathlineto{\pgfqpoint{2.510291in}{2.412090in}}%
\pgfpathlineto{\pgfqpoint{2.544433in}{2.390336in}}%
\pgfpathlineto{\pgfqpoint{2.495379in}{2.438035in}}%
\pgfpathlineto{\pgfqpoint{2.461241in}{2.459760in}}%
\pgfpathclose%
\pgfusepath{stroke,fill}%
\end{pgfscope}%
\begin{pgfscope}%
\pgfpathrectangle{\pgfqpoint{0.050000in}{0.050000in}}{\pgfqpoint{4.900000in}{4.900000in}}%
\pgfusepath{clip}%
\pgfsetbuttcap%
\pgfsetroundjoin%
\definecolor{currentfill}{rgb}{0.195617,0.195617,0.195617}%
\pgfsetfillcolor{currentfill}%
\pgfsetfillopacity{0.900000}%
\pgfsetlinewidth{0.501875pt}%
\definecolor{currentstroke}{rgb}{0.200000,0.200000,0.200000}%
\pgfsetstrokecolor{currentstroke}%
\pgfsetdash{}{0pt}%
\pgfpathmoveto{\pgfqpoint{2.177379in}{2.689891in}}%
\pgfpathlineto{\pgfqpoint{2.227081in}{2.641692in}}%
\pgfpathlineto{\pgfqpoint{2.260760in}{2.620313in}}%
\pgfpathlineto{\pgfqpoint{2.211052in}{2.668544in}}%
\pgfpathlineto{\pgfqpoint{2.177379in}{2.689891in}}%
\pgfpathclose%
\pgfusepath{stroke,fill}%
\end{pgfscope}%
\begin{pgfscope}%
\pgfpathrectangle{\pgfqpoint{0.050000in}{0.050000in}}{\pgfqpoint{4.900000in}{4.900000in}}%
\pgfusepath{clip}%
\pgfsetbuttcap%
\pgfsetroundjoin%
\definecolor{currentfill}{rgb}{0.081938,0.081938,0.081938}%
\pgfsetfillcolor{currentfill}%
\pgfsetfillopacity{0.900000}%
\pgfsetlinewidth{0.501875pt}%
\definecolor{currentstroke}{rgb}{0.200000,0.200000,0.200000}%
\pgfsetstrokecolor{currentstroke}%
\pgfsetdash{}{0pt}%
\pgfpathmoveto{\pgfqpoint{1.893415in}{2.920109in}}%
\pgfpathlineto{\pgfqpoint{1.943770in}{2.871378in}}%
\pgfpathlineto{\pgfqpoint{1.976985in}{2.850376in}}%
\pgfpathlineto{\pgfqpoint{1.926621in}{2.899140in}}%
\pgfpathlineto{\pgfqpoint{1.893415in}{2.920109in}}%
\pgfpathclose%
\pgfusepath{stroke,fill}%
\end{pgfscope}%
\begin{pgfscope}%
\pgfpathrectangle{\pgfqpoint{0.050000in}{0.050000in}}{\pgfqpoint{4.900000in}{4.900000in}}%
\pgfusepath{clip}%
\pgfsetbuttcap%
\pgfsetroundjoin%
\definecolor{currentfill}{rgb}{0.379423,0.379423,0.379423}%
\pgfsetfillcolor{currentfill}%
\pgfsetfillopacity{0.900000}%
\pgfsetlinewidth{0.501875pt}%
\definecolor{currentstroke}{rgb}{0.200000,0.200000,0.200000}%
\pgfsetstrokecolor{currentstroke}%
\pgfsetdash{}{0pt}%
\pgfpathmoveto{\pgfqpoint{2.661895in}{2.299069in}}%
\pgfpathlineto{\pgfqpoint{2.710511in}{2.251744in}}%
\pgfpathlineto{\pgfqpoint{2.745000in}{2.229715in}}%
\pgfpathlineto{\pgfqpoint{2.696382in}{2.277068in}}%
\pgfpathlineto{\pgfqpoint{2.661895in}{2.299069in}}%
\pgfpathclose%
\pgfusepath{stroke,fill}%
\end{pgfscope}%
\begin{pgfscope}%
\pgfpathrectangle{\pgfqpoint{0.050000in}{0.050000in}}{\pgfqpoint{4.900000in}{4.900000in}}%
\pgfusepath{clip}%
\pgfsetbuttcap%
\pgfsetroundjoin%
\definecolor{currentfill}{rgb}{0.278662,0.278662,0.278662}%
\pgfsetfillcolor{currentfill}%
\pgfsetfillopacity{0.900000}%
\pgfsetlinewidth{0.501875pt}%
\definecolor{currentstroke}{rgb}{0.200000,0.200000,0.200000}%
\pgfsetstrokecolor{currentstroke}%
\pgfsetdash{}{0pt}%
\pgfpathmoveto{\pgfqpoint{2.377947in}{2.529269in}}%
\pgfpathlineto{\pgfqpoint{2.427215in}{2.481413in}}%
\pgfpathlineto{\pgfqpoint{2.461241in}{2.459760in}}%
\pgfpathlineto{\pgfqpoint{2.411968in}{2.507646in}}%
\pgfpathlineto{\pgfqpoint{2.377947in}{2.529269in}}%
\pgfpathclose%
\pgfusepath{stroke,fill}%
\end{pgfscope}%
\begin{pgfscope}%
\pgfpathrectangle{\pgfqpoint{0.050000in}{0.050000in}}{\pgfqpoint{4.900000in}{4.900000in}}%
\pgfusepath{clip}%
\pgfsetbuttcap%
\pgfsetroundjoin%
\definecolor{currentfill}{rgb}{0.151326,0.151326,0.151326}%
\pgfsetfillcolor{currentfill}%
\pgfsetfillopacity{0.900000}%
\pgfsetlinewidth{0.501875pt}%
\definecolor{currentstroke}{rgb}{0.200000,0.200000,0.200000}%
\pgfsetstrokecolor{currentstroke}%
\pgfsetdash{}{0pt}%
\pgfpathmoveto{\pgfqpoint{2.093896in}{2.759555in}}%
\pgfpathlineto{\pgfqpoint{2.143818in}{2.711168in}}%
\pgfpathlineto{\pgfqpoint{2.177379in}{2.689891in}}%
\pgfpathlineto{\pgfqpoint{2.127450in}{2.738311in}}%
\pgfpathlineto{\pgfqpoint{2.093896in}{2.759555in}}%
\pgfpathclose%
\pgfusepath{stroke,fill}%
\end{pgfscope}%
\begin{pgfscope}%
\pgfpathrectangle{\pgfqpoint{0.050000in}{0.050000in}}{\pgfqpoint{4.900000in}{4.900000in}}%
\pgfusepath{clip}%
\pgfsetbuttcap%
\pgfsetroundjoin%
\definecolor{currentfill}{rgb}{0.050073,0.050073,0.050073}%
\pgfsetfillcolor{currentfill}%
\pgfsetfillopacity{0.900000}%
\pgfsetlinewidth{0.501875pt}%
\definecolor{currentstroke}{rgb}{0.200000,0.200000,0.200000}%
\pgfsetstrokecolor{currentstroke}%
\pgfsetdash{}{0pt}%
\pgfpathmoveto{\pgfqpoint{1.809742in}{2.989927in}}%
\pgfpathlineto{\pgfqpoint{1.860318in}{2.941008in}}%
\pgfpathlineto{\pgfqpoint{1.893415in}{2.920109in}}%
\pgfpathlineto{\pgfqpoint{1.842830in}{2.969061in}}%
\pgfpathlineto{\pgfqpoint{1.809742in}{2.989927in}}%
\pgfpathclose%
\pgfusepath{stroke,fill}%
\end{pgfscope}%
\begin{pgfscope}%
\pgfpathrectangle{\pgfqpoint{0.050000in}{0.050000in}}{\pgfqpoint{4.900000in}{4.900000in}}%
\pgfusepath{clip}%
\pgfsetbuttcap%
\pgfsetroundjoin%
\definecolor{currentfill}{rgb}{0.351003,0.351003,0.351003}%
\pgfsetfillcolor{currentfill}%
\pgfsetfillopacity{0.900000}%
\pgfsetlinewidth{0.501875pt}%
\definecolor{currentstroke}{rgb}{0.200000,0.200000,0.200000}%
\pgfsetstrokecolor{currentstroke}%
\pgfsetdash{}{0pt}%
\pgfpathmoveto{\pgfqpoint{2.578688in}{2.368509in}}%
\pgfpathlineto{\pgfqpoint{2.627523in}{2.320997in}}%
\pgfpathlineto{\pgfqpoint{2.661895in}{2.299069in}}%
\pgfpathlineto{\pgfqpoint{2.613058in}{2.346610in}}%
\pgfpathlineto{\pgfqpoint{2.578688in}{2.368509in}}%
\pgfpathclose%
\pgfusepath{stroke,fill}%
\end{pgfscope}%
\begin{pgfscope}%
\pgfpathrectangle{\pgfqpoint{0.050000in}{0.050000in}}{\pgfqpoint{4.900000in}{4.900000in}}%
\pgfusepath{clip}%
\pgfsetbuttcap%
\pgfsetroundjoin%
\definecolor{currentfill}{rgb}{0.234371,0.234371,0.234371}%
\pgfsetfillcolor{currentfill}%
\pgfsetfillopacity{0.900000}%
\pgfsetlinewidth{0.501875pt}%
\definecolor{currentstroke}{rgb}{0.200000,0.200000,0.200000}%
\pgfsetstrokecolor{currentstroke}%
\pgfsetdash{}{0pt}%
\pgfpathmoveto{\pgfqpoint{2.294550in}{2.598864in}}%
\pgfpathlineto{\pgfqpoint{2.344039in}{2.550820in}}%
\pgfpathlineto{\pgfqpoint{2.377947in}{2.529269in}}%
\pgfpathlineto{\pgfqpoint{2.328453in}{2.577343in}}%
\pgfpathlineto{\pgfqpoint{2.294550in}{2.598864in}}%
\pgfpathclose%
\pgfusepath{stroke,fill}%
\end{pgfscope}%
\begin{pgfscope}%
\pgfpathrectangle{\pgfqpoint{0.050000in}{0.050000in}}{\pgfqpoint{4.900000in}{4.900000in}}%
\pgfusepath{clip}%
\pgfsetbuttcap%
\pgfsetroundjoin%
\definecolor{currentfill}{rgb}{0.118354,0.118354,0.118354}%
\pgfsetfillcolor{currentfill}%
\pgfsetfillopacity{0.900000}%
\pgfsetlinewidth{0.501875pt}%
\definecolor{currentstroke}{rgb}{0.200000,0.200000,0.200000}%
\pgfsetstrokecolor{currentstroke}%
\pgfsetdash{}{0pt}%
\pgfpathmoveto{\pgfqpoint{2.010310in}{2.829305in}}%
\pgfpathlineto{\pgfqpoint{2.060452in}{2.780729in}}%
\pgfpathlineto{\pgfqpoint{2.093896in}{2.759555in}}%
\pgfpathlineto{\pgfqpoint{2.043746in}{2.808163in}}%
\pgfpathlineto{\pgfqpoint{2.010310in}{2.829305in}}%
\pgfpathclose%
\pgfusepath{stroke,fill}%
\end{pgfscope}%
\begin{pgfscope}%
\pgfpathrectangle{\pgfqpoint{0.050000in}{0.050000in}}{\pgfqpoint{4.900000in}{4.900000in}}%
\pgfusepath{clip}%
\pgfsetbuttcap%
\pgfsetroundjoin%
\definecolor{currentfill}{rgb}{0.013656,0.013656,0.013656}%
\pgfsetfillcolor{currentfill}%
\pgfsetfillopacity{0.900000}%
\pgfsetlinewidth{0.501875pt}%
\definecolor{currentstroke}{rgb}{0.200000,0.200000,0.200000}%
\pgfsetstrokecolor{currentstroke}%
\pgfsetdash{}{0pt}%
\pgfpathmoveto{\pgfqpoint{1.725967in}{3.059831in}}%
\pgfpathlineto{\pgfqpoint{1.776764in}{3.010724in}}%
\pgfpathlineto{\pgfqpoint{1.809742in}{2.989927in}}%
\pgfpathlineto{\pgfqpoint{1.758935in}{3.039069in}}%
\pgfpathlineto{\pgfqpoint{1.725967in}{3.059831in}}%
\pgfpathclose%
\pgfusepath{stroke,fill}%
\end{pgfscope}%
\begin{pgfscope}%
\pgfpathrectangle{\pgfqpoint{0.050000in}{0.050000in}}{\pgfqpoint{4.900000in}{4.900000in}}%
\pgfusepath{clip}%
\pgfsetbuttcap%
\pgfsetroundjoin%
\definecolor{currentfill}{rgb}{0.317416,0.317416,0.317416}%
\pgfsetfillcolor{currentfill}%
\pgfsetfillopacity{0.900000}%
\pgfsetlinewidth{0.501875pt}%
\definecolor{currentstroke}{rgb}{0.200000,0.200000,0.200000}%
\pgfsetstrokecolor{currentstroke}%
\pgfsetdash{}{0pt}%
\pgfpathmoveto{\pgfqpoint{2.495379in}{2.438035in}}%
\pgfpathlineto{\pgfqpoint{2.544433in}{2.390336in}}%
\pgfpathlineto{\pgfqpoint{2.578688in}{2.368509in}}%
\pgfpathlineto{\pgfqpoint{2.529631in}{2.416237in}}%
\pgfpathlineto{\pgfqpoint{2.495379in}{2.438035in}}%
\pgfpathclose%
\pgfusepath{stroke,fill}%
\end{pgfscope}%
\begin{pgfscope}%
\pgfpathrectangle{\pgfqpoint{0.050000in}{0.050000in}}{\pgfqpoint{4.900000in}{4.900000in}}%
\pgfusepath{clip}%
\pgfsetbuttcap%
\pgfsetroundjoin%
\definecolor{currentfill}{rgb}{0.195617,0.195617,0.195617}%
\pgfsetfillcolor{currentfill}%
\pgfsetfillopacity{0.900000}%
\pgfsetlinewidth{0.501875pt}%
\definecolor{currentstroke}{rgb}{0.200000,0.200000,0.200000}%
\pgfsetstrokecolor{currentstroke}%
\pgfsetdash{}{0pt}%
\pgfpathmoveto{\pgfqpoint{2.211052in}{2.668544in}}%
\pgfpathlineto{\pgfqpoint{2.260760in}{2.620313in}}%
\pgfpathlineto{\pgfqpoint{2.294550in}{2.598864in}}%
\pgfpathlineto{\pgfqpoint{2.244837in}{2.647127in}}%
\pgfpathlineto{\pgfqpoint{2.211052in}{2.668544in}}%
\pgfpathclose%
\pgfusepath{stroke,fill}%
\end{pgfscope}%
\begin{pgfscope}%
\pgfpathrectangle{\pgfqpoint{0.050000in}{0.050000in}}{\pgfqpoint{4.900000in}{4.900000in}}%
\pgfusepath{clip}%
\pgfsetbuttcap%
\pgfsetroundjoin%
\definecolor{currentfill}{rgb}{0.081938,0.081938,0.081938}%
\pgfsetfillcolor{currentfill}%
\pgfsetfillopacity{0.900000}%
\pgfsetlinewidth{0.501875pt}%
\definecolor{currentstroke}{rgb}{0.200000,0.200000,0.200000}%
\pgfsetstrokecolor{currentstroke}%
\pgfsetdash{}{0pt}%
\pgfpathmoveto{\pgfqpoint{1.926621in}{2.899140in}}%
\pgfpathlineto{\pgfqpoint{1.976985in}{2.850376in}}%
\pgfpathlineto{\pgfqpoint{2.010310in}{2.829305in}}%
\pgfpathlineto{\pgfqpoint{1.959939in}{2.878102in}}%
\pgfpathlineto{\pgfqpoint{1.926621in}{2.899140in}}%
\pgfpathclose%
\pgfusepath{stroke,fill}%
\end{pgfscope}%
\begin{pgfscope}%
\pgfpathrectangle{\pgfqpoint{0.050000in}{0.050000in}}{\pgfqpoint{4.900000in}{4.900000in}}%
\pgfusepath{clip}%
\pgfsetbuttcap%
\pgfsetroundjoin%
\definecolor{currentfill}{rgb}{0.278662,0.278662,0.278662}%
\pgfsetfillcolor{currentfill}%
\pgfsetfillopacity{0.900000}%
\pgfsetlinewidth{0.501875pt}%
\definecolor{currentstroke}{rgb}{0.200000,0.200000,0.200000}%
\pgfsetstrokecolor{currentstroke}%
\pgfsetdash{}{0pt}%
\pgfpathmoveto{\pgfqpoint{2.411968in}{2.507646in}}%
\pgfpathlineto{\pgfqpoint{2.461241in}{2.459760in}}%
\pgfpathlineto{\pgfqpoint{2.495379in}{2.438035in}}%
\pgfpathlineto{\pgfqpoint{2.446102in}{2.485951in}}%
\pgfpathlineto{\pgfqpoint{2.411968in}{2.507646in}}%
\pgfpathclose%
\pgfusepath{stroke,fill}%
\end{pgfscope}%
\begin{pgfscope}%
\pgfpathrectangle{\pgfqpoint{0.050000in}{0.050000in}}{\pgfqpoint{4.900000in}{4.900000in}}%
\pgfusepath{clip}%
\pgfsetbuttcap%
\pgfsetroundjoin%
\definecolor{currentfill}{rgb}{0.151326,0.151326,0.151326}%
\pgfsetfillcolor{currentfill}%
\pgfsetfillopacity{0.900000}%
\pgfsetlinewidth{0.501875pt}%
\definecolor{currentstroke}{rgb}{0.200000,0.200000,0.200000}%
\pgfsetstrokecolor{currentstroke}%
\pgfsetdash{}{0pt}%
\pgfpathmoveto{\pgfqpoint{2.127450in}{2.738311in}}%
\pgfpathlineto{\pgfqpoint{2.177379in}{2.689891in}}%
\pgfpathlineto{\pgfqpoint{2.211052in}{2.668544in}}%
\pgfpathlineto{\pgfqpoint{2.161117in}{2.716996in}}%
\pgfpathlineto{\pgfqpoint{2.127450in}{2.738311in}}%
\pgfpathclose%
\pgfusepath{stroke,fill}%
\end{pgfscope}%
\begin{pgfscope}%
\pgfpathrectangle{\pgfqpoint{0.050000in}{0.050000in}}{\pgfqpoint{4.900000in}{4.900000in}}%
\pgfusepath{clip}%
\pgfsetbuttcap%
\pgfsetroundjoin%
\definecolor{currentfill}{rgb}{0.050073,0.050073,0.050073}%
\pgfsetfillcolor{currentfill}%
\pgfsetfillopacity{0.900000}%
\pgfsetlinewidth{0.501875pt}%
\definecolor{currentstroke}{rgb}{0.200000,0.200000,0.200000}%
\pgfsetstrokecolor{currentstroke}%
\pgfsetdash{}{0pt}%
\pgfpathmoveto{\pgfqpoint{1.842830in}{2.969061in}}%
\pgfpathlineto{\pgfqpoint{1.893415in}{2.920109in}}%
\pgfpathlineto{\pgfqpoint{1.926621in}{2.899140in}}%
\pgfpathlineto{\pgfqpoint{1.876028in}{2.948126in}}%
\pgfpathlineto{\pgfqpoint{1.842830in}{2.969061in}}%
\pgfpathclose%
\pgfusepath{stroke,fill}%
\end{pgfscope}%
\begin{pgfscope}%
\pgfpathrectangle{\pgfqpoint{0.050000in}{0.050000in}}{\pgfqpoint{4.900000in}{4.900000in}}%
\pgfusepath{clip}%
\pgfsetbuttcap%
\pgfsetroundjoin%
\definecolor{currentfill}{rgb}{0.351003,0.351003,0.351003}%
\pgfsetfillcolor{currentfill}%
\pgfsetfillopacity{0.900000}%
\pgfsetlinewidth{0.501875pt}%
\definecolor{currentstroke}{rgb}{0.200000,0.200000,0.200000}%
\pgfsetstrokecolor{currentstroke}%
\pgfsetdash{}{0pt}%
\pgfpathmoveto{\pgfqpoint{2.613058in}{2.346610in}}%
\pgfpathlineto{\pgfqpoint{2.661895in}{2.299069in}}%
\pgfpathlineto{\pgfqpoint{2.696382in}{2.277068in}}%
\pgfpathlineto{\pgfqpoint{2.647542in}{2.324638in}}%
\pgfpathlineto{\pgfqpoint{2.613058in}{2.346610in}}%
\pgfpathclose%
\pgfusepath{stroke,fill}%
\end{pgfscope}%
\begin{pgfscope}%
\pgfpathrectangle{\pgfqpoint{0.050000in}{0.050000in}}{\pgfqpoint{4.900000in}{4.900000in}}%
\pgfusepath{clip}%
\pgfsetbuttcap%
\pgfsetroundjoin%
\definecolor{currentfill}{rgb}{0.234371,0.234371,0.234371}%
\pgfsetfillcolor{currentfill}%
\pgfsetfillopacity{0.900000}%
\pgfsetlinewidth{0.501875pt}%
\definecolor{currentstroke}{rgb}{0.200000,0.200000,0.200000}%
\pgfsetstrokecolor{currentstroke}%
\pgfsetdash{}{0pt}%
\pgfpathmoveto{\pgfqpoint{2.328453in}{2.577343in}}%
\pgfpathlineto{\pgfqpoint{2.377947in}{2.529269in}}%
\pgfpathlineto{\pgfqpoint{2.411968in}{2.507646in}}%
\pgfpathlineto{\pgfqpoint{2.362470in}{2.555751in}}%
\pgfpathlineto{\pgfqpoint{2.328453in}{2.577343in}}%
\pgfpathclose%
\pgfusepath{stroke,fill}%
\end{pgfscope}%
\begin{pgfscope}%
\pgfpathrectangle{\pgfqpoint{0.050000in}{0.050000in}}{\pgfqpoint{4.900000in}{4.900000in}}%
\pgfusepath{clip}%
\pgfsetbuttcap%
\pgfsetroundjoin%
\definecolor{currentfill}{rgb}{0.118354,0.118354,0.118354}%
\pgfsetfillcolor{currentfill}%
\pgfsetfillopacity{0.900000}%
\pgfsetlinewidth{0.501875pt}%
\definecolor{currentstroke}{rgb}{0.200000,0.200000,0.200000}%
\pgfsetstrokecolor{currentstroke}%
\pgfsetdash{}{0pt}%
\pgfpathmoveto{\pgfqpoint{2.043746in}{2.808163in}}%
\pgfpathlineto{\pgfqpoint{2.093896in}{2.759555in}}%
\pgfpathlineto{\pgfqpoint{2.127450in}{2.738311in}}%
\pgfpathlineto{\pgfqpoint{2.077294in}{2.786951in}}%
\pgfpathlineto{\pgfqpoint{2.043746in}{2.808163in}}%
\pgfpathclose%
\pgfusepath{stroke,fill}%
\end{pgfscope}%
\begin{pgfscope}%
\pgfpathrectangle{\pgfqpoint{0.050000in}{0.050000in}}{\pgfqpoint{4.900000in}{4.900000in}}%
\pgfusepath{clip}%
\pgfsetbuttcap%
\pgfsetroundjoin%
\definecolor{currentfill}{rgb}{0.013656,0.013656,0.013656}%
\pgfsetfillcolor{currentfill}%
\pgfsetfillopacity{0.900000}%
\pgfsetlinewidth{0.501875pt}%
\definecolor{currentstroke}{rgb}{0.200000,0.200000,0.200000}%
\pgfsetstrokecolor{currentstroke}%
\pgfsetdash{}{0pt}%
\pgfpathmoveto{\pgfqpoint{1.758935in}{3.039069in}}%
\pgfpathlineto{\pgfqpoint{1.809742in}{2.989927in}}%
\pgfpathlineto{\pgfqpoint{1.842830in}{2.969061in}}%
\pgfpathlineto{\pgfqpoint{1.792014in}{3.018238in}}%
\pgfpathlineto{\pgfqpoint{1.758935in}{3.039069in}}%
\pgfpathclose%
\pgfusepath{stroke,fill}%
\end{pgfscope}%
\begin{pgfscope}%
\pgfpathrectangle{\pgfqpoint{0.050000in}{0.050000in}}{\pgfqpoint{4.900000in}{4.900000in}}%
\pgfusepath{clip}%
\pgfsetbuttcap%
\pgfsetroundjoin%
\definecolor{currentfill}{rgb}{0.322584,0.322584,0.322584}%
\pgfsetfillcolor{currentfill}%
\pgfsetfillopacity{0.900000}%
\pgfsetlinewidth{0.501875pt}%
\definecolor{currentstroke}{rgb}{0.200000,0.200000,0.200000}%
\pgfsetstrokecolor{currentstroke}%
\pgfsetdash{}{0pt}%
\pgfpathmoveto{\pgfqpoint{2.529631in}{2.416237in}}%
\pgfpathlineto{\pgfqpoint{2.578688in}{2.368509in}}%
\pgfpathlineto{\pgfqpoint{2.613058in}{2.346610in}}%
\pgfpathlineto{\pgfqpoint{2.563998in}{2.394368in}}%
\pgfpathlineto{\pgfqpoint{2.529631in}{2.416237in}}%
\pgfpathclose%
\pgfusepath{stroke,fill}%
\end{pgfscope}%
\begin{pgfscope}%
\pgfpathrectangle{\pgfqpoint{0.050000in}{0.050000in}}{\pgfqpoint{4.900000in}{4.900000in}}%
\pgfusepath{clip}%
\pgfsetbuttcap%
\pgfsetroundjoin%
\definecolor{currentfill}{rgb}{0.195617,0.195617,0.195617}%
\pgfsetfillcolor{currentfill}%
\pgfsetfillopacity{0.900000}%
\pgfsetlinewidth{0.501875pt}%
\definecolor{currentstroke}{rgb}{0.200000,0.200000,0.200000}%
\pgfsetstrokecolor{currentstroke}%
\pgfsetdash{}{0pt}%
\pgfpathmoveto{\pgfqpoint{2.244837in}{2.647127in}}%
\pgfpathlineto{\pgfqpoint{2.294550in}{2.598864in}}%
\pgfpathlineto{\pgfqpoint{2.328453in}{2.577343in}}%
\pgfpathlineto{\pgfqpoint{2.278734in}{2.625637in}}%
\pgfpathlineto{\pgfqpoint{2.244837in}{2.647127in}}%
\pgfpathclose%
\pgfusepath{stroke,fill}%
\end{pgfscope}%
\begin{pgfscope}%
\pgfpathrectangle{\pgfqpoint{0.050000in}{0.050000in}}{\pgfqpoint{4.900000in}{4.900000in}}%
\pgfusepath{clip}%
\pgfsetbuttcap%
\pgfsetroundjoin%
\definecolor{currentfill}{rgb}{0.081938,0.081938,0.081938}%
\pgfsetfillcolor{currentfill}%
\pgfsetfillopacity{0.900000}%
\pgfsetlinewidth{0.501875pt}%
\definecolor{currentstroke}{rgb}{0.200000,0.200000,0.200000}%
\pgfsetstrokecolor{currentstroke}%
\pgfsetdash{}{0pt}%
\pgfpathmoveto{\pgfqpoint{1.959939in}{2.878102in}}%
\pgfpathlineto{\pgfqpoint{2.010310in}{2.829305in}}%
\pgfpathlineto{\pgfqpoint{2.043746in}{2.808163in}}%
\pgfpathlineto{\pgfqpoint{1.993367in}{2.856993in}}%
\pgfpathlineto{\pgfqpoint{1.959939in}{2.878102in}}%
\pgfpathclose%
\pgfusepath{stroke,fill}%
\end{pgfscope}%
\begin{pgfscope}%
\pgfpathrectangle{\pgfqpoint{0.050000in}{0.050000in}}{\pgfqpoint{4.900000in}{4.900000in}}%
\pgfusepath{clip}%
\pgfsetbuttcap%
\pgfsetroundjoin%
\definecolor{currentfill}{rgb}{0.278662,0.278662,0.278662}%
\pgfsetfillcolor{currentfill}%
\pgfsetfillopacity{0.900000}%
\pgfsetlinewidth{0.501875pt}%
\definecolor{currentstroke}{rgb}{0.200000,0.200000,0.200000}%
\pgfsetstrokecolor{currentstroke}%
\pgfsetdash{}{0pt}%
\pgfpathmoveto{\pgfqpoint{2.446102in}{2.485951in}}%
\pgfpathlineto{\pgfqpoint{2.495379in}{2.438035in}}%
\pgfpathlineto{\pgfqpoint{2.529631in}{2.416237in}}%
\pgfpathlineto{\pgfqpoint{2.480350in}{2.464184in}}%
\pgfpathlineto{\pgfqpoint{2.446102in}{2.485951in}}%
\pgfpathclose%
\pgfusepath{stroke,fill}%
\end{pgfscope}%
\begin{pgfscope}%
\pgfpathrectangle{\pgfqpoint{0.050000in}{0.050000in}}{\pgfqpoint{4.900000in}{4.900000in}}%
\pgfusepath{clip}%
\pgfsetbuttcap%
\pgfsetroundjoin%
\definecolor{currentfill}{rgb}{0.151326,0.151326,0.151326}%
\pgfsetfillcolor{currentfill}%
\pgfsetfillopacity{0.900000}%
\pgfsetlinewidth{0.501875pt}%
\definecolor{currentstroke}{rgb}{0.200000,0.200000,0.200000}%
\pgfsetstrokecolor{currentstroke}%
\pgfsetdash{}{0pt}%
\pgfpathmoveto{\pgfqpoint{2.161117in}{2.716996in}}%
\pgfpathlineto{\pgfqpoint{2.211052in}{2.668544in}}%
\pgfpathlineto{\pgfqpoint{2.244837in}{2.647127in}}%
\pgfpathlineto{\pgfqpoint{2.194895in}{2.695610in}}%
\pgfpathlineto{\pgfqpoint{2.161117in}{2.716996in}}%
\pgfpathclose%
\pgfusepath{stroke,fill}%
\end{pgfscope}%
\begin{pgfscope}%
\pgfpathrectangle{\pgfqpoint{0.050000in}{0.050000in}}{\pgfqpoint{4.900000in}{4.900000in}}%
\pgfusepath{clip}%
\pgfsetbuttcap%
\pgfsetroundjoin%
\definecolor{currentfill}{rgb}{0.050073,0.050073,0.050073}%
\pgfsetfillcolor{currentfill}%
\pgfsetfillopacity{0.900000}%
\pgfsetlinewidth{0.501875pt}%
\definecolor{currentstroke}{rgb}{0.200000,0.200000,0.200000}%
\pgfsetstrokecolor{currentstroke}%
\pgfsetdash{}{0pt}%
\pgfpathmoveto{\pgfqpoint{1.876028in}{2.948126in}}%
\pgfpathlineto{\pgfqpoint{1.926621in}{2.899140in}}%
\pgfpathlineto{\pgfqpoint{1.959939in}{2.878102in}}%
\pgfpathlineto{\pgfqpoint{1.909337in}{2.927122in}}%
\pgfpathlineto{\pgfqpoint{1.876028in}{2.948126in}}%
\pgfpathclose%
\pgfusepath{stroke,fill}%
\end{pgfscope}%
\begin{pgfscope}%
\pgfpathrectangle{\pgfqpoint{0.050000in}{0.050000in}}{\pgfqpoint{4.900000in}{4.900000in}}%
\pgfusepath{clip}%
\pgfsetbuttcap%
\pgfsetroundjoin%
\definecolor{currentfill}{rgb}{0.239908,0.239908,0.239908}%
\pgfsetfillcolor{currentfill}%
\pgfsetfillopacity{0.900000}%
\pgfsetlinewidth{0.501875pt}%
\definecolor{currentstroke}{rgb}{0.200000,0.200000,0.200000}%
\pgfsetstrokecolor{currentstroke}%
\pgfsetdash{}{0pt}%
\pgfpathmoveto{\pgfqpoint{2.362470in}{2.555751in}}%
\pgfpathlineto{\pgfqpoint{2.411968in}{2.507646in}}%
\pgfpathlineto{\pgfqpoint{2.446102in}{2.485951in}}%
\pgfpathlineto{\pgfqpoint{2.396599in}{2.534087in}}%
\pgfpathlineto{\pgfqpoint{2.362470in}{2.555751in}}%
\pgfpathclose%
\pgfusepath{stroke,fill}%
\end{pgfscope}%
\begin{pgfscope}%
\pgfpathrectangle{\pgfqpoint{0.050000in}{0.050000in}}{\pgfqpoint{4.900000in}{4.900000in}}%
\pgfusepath{clip}%
\pgfsetbuttcap%
\pgfsetroundjoin%
\definecolor{currentfill}{rgb}{0.118354,0.118354,0.118354}%
\pgfsetfillcolor{currentfill}%
\pgfsetfillopacity{0.900000}%
\pgfsetlinewidth{0.501875pt}%
\definecolor{currentstroke}{rgb}{0.200000,0.200000,0.200000}%
\pgfsetstrokecolor{currentstroke}%
\pgfsetdash{}{0pt}%
\pgfpathmoveto{\pgfqpoint{2.077294in}{2.786951in}}%
\pgfpathlineto{\pgfqpoint{2.127450in}{2.738311in}}%
\pgfpathlineto{\pgfqpoint{2.161117in}{2.716996in}}%
\pgfpathlineto{\pgfqpoint{2.110953in}{2.765669in}}%
\pgfpathlineto{\pgfqpoint{2.077294in}{2.786951in}}%
\pgfpathclose%
\pgfusepath{stroke,fill}%
\end{pgfscope}%
\begin{pgfscope}%
\pgfpathrectangle{\pgfqpoint{0.050000in}{0.050000in}}{\pgfqpoint{4.900000in}{4.900000in}}%
\pgfusepath{clip}%
\pgfsetbuttcap%
\pgfsetroundjoin%
\definecolor{currentfill}{rgb}{0.013656,0.013656,0.013656}%
\pgfsetfillcolor{currentfill}%
\pgfsetfillopacity{0.900000}%
\pgfsetlinewidth{0.501875pt}%
\definecolor{currentstroke}{rgb}{0.200000,0.200000,0.200000}%
\pgfsetstrokecolor{currentstroke}%
\pgfsetdash{}{0pt}%
\pgfpathmoveto{\pgfqpoint{1.792014in}{3.018238in}}%
\pgfpathlineto{\pgfqpoint{1.842830in}{2.969061in}}%
\pgfpathlineto{\pgfqpoint{1.876028in}{2.948126in}}%
\pgfpathlineto{\pgfqpoint{1.825203in}{2.997337in}}%
\pgfpathlineto{\pgfqpoint{1.792014in}{3.018238in}}%
\pgfpathclose%
\pgfusepath{stroke,fill}%
\end{pgfscope}%
\begin{pgfscope}%
\pgfpathrectangle{\pgfqpoint{0.050000in}{0.050000in}}{\pgfqpoint{4.900000in}{4.900000in}}%
\pgfusepath{clip}%
\pgfsetbuttcap%
\pgfsetroundjoin%
\definecolor{currentfill}{rgb}{0.322584,0.322584,0.322584}%
\pgfsetfillcolor{currentfill}%
\pgfsetfillopacity{0.900000}%
\pgfsetlinewidth{0.501875pt}%
\definecolor{currentstroke}{rgb}{0.200000,0.200000,0.200000}%
\pgfsetstrokecolor{currentstroke}%
\pgfsetdash{}{0pt}%
\pgfpathmoveto{\pgfqpoint{2.563998in}{2.394368in}}%
\pgfpathlineto{\pgfqpoint{2.613058in}{2.346610in}}%
\pgfpathlineto{\pgfqpoint{2.647542in}{2.324638in}}%
\pgfpathlineto{\pgfqpoint{2.598479in}{2.372425in}}%
\pgfpathlineto{\pgfqpoint{2.563998in}{2.394368in}}%
\pgfpathclose%
\pgfusepath{stroke,fill}%
\end{pgfscope}%
\begin{pgfscope}%
\pgfpathrectangle{\pgfqpoint{0.050000in}{0.050000in}}{\pgfqpoint{4.900000in}{4.900000in}}%
\pgfusepath{clip}%
\pgfsetbuttcap%
\pgfsetroundjoin%
\definecolor{currentfill}{rgb}{0.195617,0.195617,0.195617}%
\pgfsetfillcolor{currentfill}%
\pgfsetfillopacity{0.900000}%
\pgfsetlinewidth{0.501875pt}%
\definecolor{currentstroke}{rgb}{0.200000,0.200000,0.200000}%
\pgfsetstrokecolor{currentstroke}%
\pgfsetdash{}{0pt}%
\pgfpathmoveto{\pgfqpoint{2.278734in}{2.625637in}}%
\pgfpathlineto{\pgfqpoint{2.328453in}{2.577343in}}%
\pgfpathlineto{\pgfqpoint{2.362470in}{2.555751in}}%
\pgfpathlineto{\pgfqpoint{2.312745in}{2.604076in}}%
\pgfpathlineto{\pgfqpoint{2.278734in}{2.625637in}}%
\pgfpathclose%
\pgfusepath{stroke,fill}%
\end{pgfscope}%
\begin{pgfscope}%
\pgfpathrectangle{\pgfqpoint{0.050000in}{0.050000in}}{\pgfqpoint{4.900000in}{4.900000in}}%
\pgfusepath{clip}%
\pgfsetbuttcap%
\pgfsetroundjoin%
\definecolor{currentfill}{rgb}{0.081938,0.081938,0.081938}%
\pgfsetfillcolor{currentfill}%
\pgfsetfillopacity{0.900000}%
\pgfsetlinewidth{0.501875pt}%
\definecolor{currentstroke}{rgb}{0.200000,0.200000,0.200000}%
\pgfsetstrokecolor{currentstroke}%
\pgfsetdash{}{0pt}%
\pgfpathmoveto{\pgfqpoint{1.993367in}{2.856993in}}%
\pgfpathlineto{\pgfqpoint{2.043746in}{2.808163in}}%
\pgfpathlineto{\pgfqpoint{2.077294in}{2.786951in}}%
\pgfpathlineto{\pgfqpoint{2.026907in}{2.835814in}}%
\pgfpathlineto{\pgfqpoint{1.993367in}{2.856993in}}%
\pgfpathclose%
\pgfusepath{stroke,fill}%
\end{pgfscope}%
\begin{pgfscope}%
\pgfpathrectangle{\pgfqpoint{0.050000in}{0.050000in}}{\pgfqpoint{4.900000in}{4.900000in}}%
\pgfusepath{clip}%
\pgfsetbuttcap%
\pgfsetroundjoin%
\definecolor{currentfill}{rgb}{0.278662,0.278662,0.278662}%
\pgfsetfillcolor{currentfill}%
\pgfsetfillopacity{0.900000}%
\pgfsetlinewidth{0.501875pt}%
\definecolor{currentstroke}{rgb}{0.200000,0.200000,0.200000}%
\pgfsetstrokecolor{currentstroke}%
\pgfsetdash{}{0pt}%
\pgfpathmoveto{\pgfqpoint{2.480350in}{2.464184in}}%
\pgfpathlineto{\pgfqpoint{2.529631in}{2.416237in}}%
\pgfpathlineto{\pgfqpoint{2.563998in}{2.394368in}}%
\pgfpathlineto{\pgfqpoint{2.514713in}{2.442345in}}%
\pgfpathlineto{\pgfqpoint{2.480350in}{2.464184in}}%
\pgfpathclose%
\pgfusepath{stroke,fill}%
\end{pgfscope}%
\begin{pgfscope}%
\pgfpathrectangle{\pgfqpoint{0.050000in}{0.050000in}}{\pgfqpoint{4.900000in}{4.900000in}}%
\pgfusepath{clip}%
\pgfsetbuttcap%
\pgfsetroundjoin%
\definecolor{currentfill}{rgb}{0.151326,0.151326,0.151326}%
\pgfsetfillcolor{currentfill}%
\pgfsetfillopacity{0.900000}%
\pgfsetlinewidth{0.501875pt}%
\definecolor{currentstroke}{rgb}{0.200000,0.200000,0.200000}%
\pgfsetstrokecolor{currentstroke}%
\pgfsetdash{}{0pt}%
\pgfpathmoveto{\pgfqpoint{2.194895in}{2.695610in}}%
\pgfpathlineto{\pgfqpoint{2.244837in}{2.647127in}}%
\pgfpathlineto{\pgfqpoint{2.278734in}{2.625637in}}%
\pgfpathlineto{\pgfqpoint{2.228787in}{2.674152in}}%
\pgfpathlineto{\pgfqpoint{2.194895in}{2.695610in}}%
\pgfpathclose%
\pgfusepath{stroke,fill}%
\end{pgfscope}%
\begin{pgfscope}%
\pgfpathrectangle{\pgfqpoint{0.050000in}{0.050000in}}{\pgfqpoint{4.900000in}{4.900000in}}%
\pgfusepath{clip}%
\pgfsetbuttcap%
\pgfsetroundjoin%
\definecolor{currentfill}{rgb}{0.050073,0.050073,0.050073}%
\pgfsetfillcolor{currentfill}%
\pgfsetfillopacity{0.900000}%
\pgfsetlinewidth{0.501875pt}%
\definecolor{currentstroke}{rgb}{0.200000,0.200000,0.200000}%
\pgfsetstrokecolor{currentstroke}%
\pgfsetdash{}{0pt}%
\pgfpathmoveto{\pgfqpoint{1.909337in}{2.927122in}}%
\pgfpathlineto{\pgfqpoint{1.959939in}{2.878102in}}%
\pgfpathlineto{\pgfqpoint{1.993367in}{2.856993in}}%
\pgfpathlineto{\pgfqpoint{1.942758in}{2.906047in}}%
\pgfpathlineto{\pgfqpoint{1.909337in}{2.927122in}}%
\pgfpathclose%
\pgfusepath{stroke,fill}%
\end{pgfscope}%
\begin{pgfscope}%
\pgfpathrectangle{\pgfqpoint{0.050000in}{0.050000in}}{\pgfqpoint{4.900000in}{4.900000in}}%
\pgfusepath{clip}%
\pgfsetbuttcap%
\pgfsetroundjoin%
\definecolor{currentfill}{rgb}{0.239908,0.239908,0.239908}%
\pgfsetfillcolor{currentfill}%
\pgfsetfillopacity{0.900000}%
\pgfsetlinewidth{0.501875pt}%
\definecolor{currentstroke}{rgb}{0.200000,0.200000,0.200000}%
\pgfsetstrokecolor{currentstroke}%
\pgfsetdash{}{0pt}%
\pgfpathmoveto{\pgfqpoint{2.396599in}{2.534087in}}%
\pgfpathlineto{\pgfqpoint{2.446102in}{2.485951in}}%
\pgfpathlineto{\pgfqpoint{2.480350in}{2.464184in}}%
\pgfpathlineto{\pgfqpoint{2.430843in}{2.512351in}}%
\pgfpathlineto{\pgfqpoint{2.396599in}{2.534087in}}%
\pgfpathclose%
\pgfusepath{stroke,fill}%
\end{pgfscope}%
\begin{pgfscope}%
\pgfpathrectangle{\pgfqpoint{0.050000in}{0.050000in}}{\pgfqpoint{4.900000in}{4.900000in}}%
\pgfusepath{clip}%
\pgfsetbuttcap%
\pgfsetroundjoin%
\definecolor{currentfill}{rgb}{0.118354,0.118354,0.118354}%
\pgfsetfillcolor{currentfill}%
\pgfsetfillopacity{0.900000}%
\pgfsetlinewidth{0.501875pt}%
\definecolor{currentstroke}{rgb}{0.200000,0.200000,0.200000}%
\pgfsetstrokecolor{currentstroke}%
\pgfsetdash{}{0pt}%
\pgfpathmoveto{\pgfqpoint{2.110953in}{2.765669in}}%
\pgfpathlineto{\pgfqpoint{2.161117in}{2.716996in}}%
\pgfpathlineto{\pgfqpoint{2.194895in}{2.695610in}}%
\pgfpathlineto{\pgfqpoint{2.144725in}{2.744315in}}%
\pgfpathlineto{\pgfqpoint{2.110953in}{2.765669in}}%
\pgfpathclose%
\pgfusepath{stroke,fill}%
\end{pgfscope}%
\begin{pgfscope}%
\pgfpathrectangle{\pgfqpoint{0.050000in}{0.050000in}}{\pgfqpoint{4.900000in}{4.900000in}}%
\pgfusepath{clip}%
\pgfsetbuttcap%
\pgfsetroundjoin%
\definecolor{currentfill}{rgb}{0.013656,0.013656,0.013656}%
\pgfsetfillcolor{currentfill}%
\pgfsetfillopacity{0.900000}%
\pgfsetlinewidth{0.501875pt}%
\definecolor{currentstroke}{rgb}{0.200000,0.200000,0.200000}%
\pgfsetstrokecolor{currentstroke}%
\pgfsetdash{}{0pt}%
\pgfpathmoveto{\pgfqpoint{1.825203in}{2.997337in}}%
\pgfpathlineto{\pgfqpoint{1.876028in}{2.948126in}}%
\pgfpathlineto{\pgfqpoint{1.909337in}{2.927122in}}%
\pgfpathlineto{\pgfqpoint{1.858504in}{2.976366in}}%
\pgfpathlineto{\pgfqpoint{1.825203in}{2.997337in}}%
\pgfpathclose%
\pgfusepath{stroke,fill}%
\end{pgfscope}%
\begin{pgfscope}%
\pgfpathrectangle{\pgfqpoint{0.050000in}{0.050000in}}{\pgfqpoint{4.900000in}{4.900000in}}%
\pgfusepath{clip}%
\pgfsetbuttcap%
\pgfsetroundjoin%
\definecolor{currentfill}{rgb}{0.195617,0.195617,0.195617}%
\pgfsetfillcolor{currentfill}%
\pgfsetfillopacity{0.900000}%
\pgfsetlinewidth{0.501875pt}%
\definecolor{currentstroke}{rgb}{0.200000,0.200000,0.200000}%
\pgfsetstrokecolor{currentstroke}%
\pgfsetdash{}{0pt}%
\pgfpathmoveto{\pgfqpoint{2.312745in}{2.604076in}}%
\pgfpathlineto{\pgfqpoint{2.362470in}{2.555751in}}%
\pgfpathlineto{\pgfqpoint{2.396599in}{2.534087in}}%
\pgfpathlineto{\pgfqpoint{2.346870in}{2.582444in}}%
\pgfpathlineto{\pgfqpoint{2.312745in}{2.604076in}}%
\pgfpathclose%
\pgfusepath{stroke,fill}%
\end{pgfscope}%
\begin{pgfscope}%
\pgfpathrectangle{\pgfqpoint{0.050000in}{0.050000in}}{\pgfqpoint{4.900000in}{4.900000in}}%
\pgfusepath{clip}%
\pgfsetbuttcap%
\pgfsetroundjoin%
\definecolor{currentfill}{rgb}{0.081938,0.081938,0.081938}%
\pgfsetfillcolor{currentfill}%
\pgfsetfillopacity{0.900000}%
\pgfsetlinewidth{0.501875pt}%
\definecolor{currentstroke}{rgb}{0.200000,0.200000,0.200000}%
\pgfsetstrokecolor{currentstroke}%
\pgfsetdash{}{0pt}%
\pgfpathmoveto{\pgfqpoint{2.026907in}{2.835814in}}%
\pgfpathlineto{\pgfqpoint{2.077294in}{2.786951in}}%
\pgfpathlineto{\pgfqpoint{2.110953in}{2.765669in}}%
\pgfpathlineto{\pgfqpoint{2.060560in}{2.814565in}}%
\pgfpathlineto{\pgfqpoint{2.026907in}{2.835814in}}%
\pgfpathclose%
\pgfusepath{stroke,fill}%
\end{pgfscope}%
\begin{pgfscope}%
\pgfpathrectangle{\pgfqpoint{0.050000in}{0.050000in}}{\pgfqpoint{4.900000in}{4.900000in}}%
\pgfusepath{clip}%
\pgfsetbuttcap%
\pgfsetroundjoin%
\definecolor{currentfill}{rgb}{0.278662,0.278662,0.278662}%
\pgfsetfillcolor{currentfill}%
\pgfsetfillopacity{0.900000}%
\pgfsetlinewidth{0.501875pt}%
\definecolor{currentstroke}{rgb}{0.200000,0.200000,0.200000}%
\pgfsetstrokecolor{currentstroke}%
\pgfsetdash{}{0pt}%
\pgfpathmoveto{\pgfqpoint{2.514713in}{2.442345in}}%
\pgfpathlineto{\pgfqpoint{2.563998in}{2.394368in}}%
\pgfpathlineto{\pgfqpoint{2.598479in}{2.372425in}}%
\pgfpathlineto{\pgfqpoint{2.549190in}{2.420432in}}%
\pgfpathlineto{\pgfqpoint{2.514713in}{2.442345in}}%
\pgfpathclose%
\pgfusepath{stroke,fill}%
\end{pgfscope}%
\begin{pgfscope}%
\pgfpathrectangle{\pgfqpoint{0.050000in}{0.050000in}}{\pgfqpoint{4.900000in}{4.900000in}}%
\pgfusepath{clip}%
\pgfsetbuttcap%
\pgfsetroundjoin%
\definecolor{currentfill}{rgb}{0.156863,0.156863,0.156863}%
\pgfsetfillcolor{currentfill}%
\pgfsetfillopacity{0.900000}%
\pgfsetlinewidth{0.501875pt}%
\definecolor{currentstroke}{rgb}{0.200000,0.200000,0.200000}%
\pgfsetstrokecolor{currentstroke}%
\pgfsetdash{}{0pt}%
\pgfpathmoveto{\pgfqpoint{2.228787in}{2.674152in}}%
\pgfpathlineto{\pgfqpoint{2.278734in}{2.625637in}}%
\pgfpathlineto{\pgfqpoint{2.312745in}{2.604076in}}%
\pgfpathlineto{\pgfqpoint{2.262792in}{2.652623in}}%
\pgfpathlineto{\pgfqpoint{2.228787in}{2.674152in}}%
\pgfpathclose%
\pgfusepath{stroke,fill}%
\end{pgfscope}%
\begin{pgfscope}%
\pgfpathrectangle{\pgfqpoint{0.050000in}{0.050000in}}{\pgfqpoint{4.900000in}{4.900000in}}%
\pgfusepath{clip}%
\pgfsetbuttcap%
\pgfsetroundjoin%
\definecolor{currentfill}{rgb}{0.050073,0.050073,0.050073}%
\pgfsetfillcolor{currentfill}%
\pgfsetfillopacity{0.900000}%
\pgfsetlinewidth{0.501875pt}%
\definecolor{currentstroke}{rgb}{0.200000,0.200000,0.200000}%
\pgfsetstrokecolor{currentstroke}%
\pgfsetdash{}{0pt}%
\pgfpathmoveto{\pgfqpoint{1.942758in}{2.906047in}}%
\pgfpathlineto{\pgfqpoint{1.993367in}{2.856993in}}%
\pgfpathlineto{\pgfqpoint{2.026907in}{2.835814in}}%
\pgfpathlineto{\pgfqpoint{1.976290in}{2.884902in}}%
\pgfpathlineto{\pgfqpoint{1.942758in}{2.906047in}}%
\pgfpathclose%
\pgfusepath{stroke,fill}%
\end{pgfscope}%
\begin{pgfscope}%
\pgfpathrectangle{\pgfqpoint{0.050000in}{0.050000in}}{\pgfqpoint{4.900000in}{4.900000in}}%
\pgfusepath{clip}%
\pgfsetbuttcap%
\pgfsetroundjoin%
\definecolor{currentfill}{rgb}{0.239908,0.239908,0.239908}%
\pgfsetfillcolor{currentfill}%
\pgfsetfillopacity{0.900000}%
\pgfsetlinewidth{0.501875pt}%
\definecolor{currentstroke}{rgb}{0.200000,0.200000,0.200000}%
\pgfsetstrokecolor{currentstroke}%
\pgfsetdash{}{0pt}%
\pgfpathmoveto{\pgfqpoint{2.430843in}{2.512351in}}%
\pgfpathlineto{\pgfqpoint{2.480350in}{2.464184in}}%
\pgfpathlineto{\pgfqpoint{2.514713in}{2.442345in}}%
\pgfpathlineto{\pgfqpoint{2.465202in}{2.490542in}}%
\pgfpathlineto{\pgfqpoint{2.430843in}{2.512351in}}%
\pgfpathclose%
\pgfusepath{stroke,fill}%
\end{pgfscope}%
\begin{pgfscope}%
\pgfpathrectangle{\pgfqpoint{0.050000in}{0.050000in}}{\pgfqpoint{4.900000in}{4.900000in}}%
\pgfusepath{clip}%
\pgfsetbuttcap%
\pgfsetroundjoin%
\definecolor{currentfill}{rgb}{0.118354,0.118354,0.118354}%
\pgfsetfillcolor{currentfill}%
\pgfsetfillopacity{0.900000}%
\pgfsetlinewidth{0.501875pt}%
\definecolor{currentstroke}{rgb}{0.200000,0.200000,0.200000}%
\pgfsetstrokecolor{currentstroke}%
\pgfsetdash{}{0pt}%
\pgfpathmoveto{\pgfqpoint{2.144725in}{2.744315in}}%
\pgfpathlineto{\pgfqpoint{2.194895in}{2.695610in}}%
\pgfpathlineto{\pgfqpoint{2.228787in}{2.674152in}}%
\pgfpathlineto{\pgfqpoint{2.178611in}{2.722890in}}%
\pgfpathlineto{\pgfqpoint{2.144725in}{2.744315in}}%
\pgfpathclose%
\pgfusepath{stroke,fill}%
\end{pgfscope}%
\begin{pgfscope}%
\pgfpathrectangle{\pgfqpoint{0.050000in}{0.050000in}}{\pgfqpoint{4.900000in}{4.900000in}}%
\pgfusepath{clip}%
\pgfsetbuttcap%
\pgfsetroundjoin%
\definecolor{currentfill}{rgb}{0.013656,0.013656,0.013656}%
\pgfsetfillcolor{currentfill}%
\pgfsetfillopacity{0.900000}%
\pgfsetlinewidth{0.501875pt}%
\definecolor{currentstroke}{rgb}{0.200000,0.200000,0.200000}%
\pgfsetstrokecolor{currentstroke}%
\pgfsetdash{}{0pt}%
\pgfpathmoveto{\pgfqpoint{1.858504in}{2.976366in}}%
\pgfpathlineto{\pgfqpoint{1.909337in}{2.927122in}}%
\pgfpathlineto{\pgfqpoint{1.942758in}{2.906047in}}%
\pgfpathlineto{\pgfqpoint{1.891916in}{2.955326in}}%
\pgfpathlineto{\pgfqpoint{1.858504in}{2.976366in}}%
\pgfpathclose%
\pgfusepath{stroke,fill}%
\end{pgfscope}%
\begin{pgfscope}%
\pgfpathrectangle{\pgfqpoint{0.050000in}{0.050000in}}{\pgfqpoint{4.900000in}{4.900000in}}%
\pgfusepath{clip}%
\pgfsetbuttcap%
\pgfsetroundjoin%
\definecolor{currentfill}{rgb}{0.195617,0.195617,0.195617}%
\pgfsetfillcolor{currentfill}%
\pgfsetfillopacity{0.900000}%
\pgfsetlinewidth{0.501875pt}%
\definecolor{currentstroke}{rgb}{0.200000,0.200000,0.200000}%
\pgfsetstrokecolor{currentstroke}%
\pgfsetdash{}{0pt}%
\pgfpathmoveto{\pgfqpoint{2.346870in}{2.582444in}}%
\pgfpathlineto{\pgfqpoint{2.396599in}{2.534087in}}%
\pgfpathlineto{\pgfqpoint{2.430843in}{2.512351in}}%
\pgfpathlineto{\pgfqpoint{2.381109in}{2.560738in}}%
\pgfpathlineto{\pgfqpoint{2.346870in}{2.582444in}}%
\pgfpathclose%
\pgfusepath{stroke,fill}%
\end{pgfscope}%
\begin{pgfscope}%
\pgfpathrectangle{\pgfqpoint{0.050000in}{0.050000in}}{\pgfqpoint{4.900000in}{4.900000in}}%
\pgfusepath{clip}%
\pgfsetbuttcap%
\pgfsetroundjoin%
\definecolor{currentfill}{rgb}{0.086490,0.086490,0.086490}%
\pgfsetfillcolor{currentfill}%
\pgfsetfillopacity{0.900000}%
\pgfsetlinewidth{0.501875pt}%
\definecolor{currentstroke}{rgb}{0.200000,0.200000,0.200000}%
\pgfsetstrokecolor{currentstroke}%
\pgfsetdash{}{0pt}%
\pgfpathmoveto{\pgfqpoint{2.060560in}{2.814565in}}%
\pgfpathlineto{\pgfqpoint{2.110953in}{2.765669in}}%
\pgfpathlineto{\pgfqpoint{2.144725in}{2.744315in}}%
\pgfpathlineto{\pgfqpoint{2.094325in}{2.793244in}}%
\pgfpathlineto{\pgfqpoint{2.060560in}{2.814565in}}%
\pgfpathclose%
\pgfusepath{stroke,fill}%
\end{pgfscope}%
\begin{pgfscope}%
\pgfpathrectangle{\pgfqpoint{0.050000in}{0.050000in}}{\pgfqpoint{4.900000in}{4.900000in}}%
\pgfusepath{clip}%
\pgfsetbuttcap%
\pgfsetroundjoin%
\definecolor{currentfill}{rgb}{0.156863,0.156863,0.156863}%
\pgfsetfillcolor{currentfill}%
\pgfsetfillopacity{0.900000}%
\pgfsetlinewidth{0.501875pt}%
\definecolor{currentstroke}{rgb}{0.200000,0.200000,0.200000}%
\pgfsetstrokecolor{currentstroke}%
\pgfsetdash{}{0pt}%
\pgfpathmoveto{\pgfqpoint{2.262792in}{2.652623in}}%
\pgfpathlineto{\pgfqpoint{2.312745in}{2.604076in}}%
\pgfpathlineto{\pgfqpoint{2.346870in}{2.582444in}}%
\pgfpathlineto{\pgfqpoint{2.296912in}{2.631022in}}%
\pgfpathlineto{\pgfqpoint{2.262792in}{2.652623in}}%
\pgfpathclose%
\pgfusepath{stroke,fill}%
\end{pgfscope}%
\begin{pgfscope}%
\pgfpathrectangle{\pgfqpoint{0.050000in}{0.050000in}}{\pgfqpoint{4.900000in}{4.900000in}}%
\pgfusepath{clip}%
\pgfsetbuttcap%
\pgfsetroundjoin%
\definecolor{currentfill}{rgb}{0.050073,0.050073,0.050073}%
\pgfsetfillcolor{currentfill}%
\pgfsetfillopacity{0.900000}%
\pgfsetlinewidth{0.501875pt}%
\definecolor{currentstroke}{rgb}{0.200000,0.200000,0.200000}%
\pgfsetstrokecolor{currentstroke}%
\pgfsetdash{}{0pt}%
\pgfpathmoveto{\pgfqpoint{1.976290in}{2.884902in}}%
\pgfpathlineto{\pgfqpoint{2.026907in}{2.835814in}}%
\pgfpathlineto{\pgfqpoint{2.060560in}{2.814565in}}%
\pgfpathlineto{\pgfqpoint{2.009935in}{2.863686in}}%
\pgfpathlineto{\pgfqpoint{1.976290in}{2.884902in}}%
\pgfpathclose%
\pgfusepath{stroke,fill}%
\end{pgfscope}%
\begin{pgfscope}%
\pgfpathrectangle{\pgfqpoint{0.050000in}{0.050000in}}{\pgfqpoint{4.900000in}{4.900000in}}%
\pgfusepath{clip}%
\pgfsetbuttcap%
\pgfsetroundjoin%
\definecolor{currentfill}{rgb}{0.239908,0.239908,0.239908}%
\pgfsetfillcolor{currentfill}%
\pgfsetfillopacity{0.900000}%
\pgfsetlinewidth{0.501875pt}%
\definecolor{currentstroke}{rgb}{0.200000,0.200000,0.200000}%
\pgfsetstrokecolor{currentstroke}%
\pgfsetdash{}{0pt}%
\pgfpathmoveto{\pgfqpoint{2.465202in}{2.490542in}}%
\pgfpathlineto{\pgfqpoint{2.514713in}{2.442345in}}%
\pgfpathlineto{\pgfqpoint{2.549190in}{2.420432in}}%
\pgfpathlineto{\pgfqpoint{2.499675in}{2.468660in}}%
\pgfpathlineto{\pgfqpoint{2.465202in}{2.490542in}}%
\pgfpathclose%
\pgfusepath{stroke,fill}%
\end{pgfscope}%
\begin{pgfscope}%
\pgfpathrectangle{\pgfqpoint{0.050000in}{0.050000in}}{\pgfqpoint{4.900000in}{4.900000in}}%
\pgfusepath{clip}%
\pgfsetbuttcap%
\pgfsetroundjoin%
\definecolor{currentfill}{rgb}{0.118354,0.118354,0.118354}%
\pgfsetfillcolor{currentfill}%
\pgfsetfillopacity{0.900000}%
\pgfsetlinewidth{0.501875pt}%
\definecolor{currentstroke}{rgb}{0.200000,0.200000,0.200000}%
\pgfsetstrokecolor{currentstroke}%
\pgfsetdash{}{0pt}%
\pgfpathmoveto{\pgfqpoint{2.178611in}{2.722890in}}%
\pgfpathlineto{\pgfqpoint{2.228787in}{2.674152in}}%
\pgfpathlineto{\pgfqpoint{2.262792in}{2.652623in}}%
\pgfpathlineto{\pgfqpoint{2.212610in}{2.701393in}}%
\pgfpathlineto{\pgfqpoint{2.178611in}{2.722890in}}%
\pgfpathclose%
\pgfusepath{stroke,fill}%
\end{pgfscope}%
\begin{pgfscope}%
\pgfpathrectangle{\pgfqpoint{0.050000in}{0.050000in}}{\pgfqpoint{4.900000in}{4.900000in}}%
\pgfusepath{clip}%
\pgfsetbuttcap%
\pgfsetroundjoin%
\definecolor{currentfill}{rgb}{0.013656,0.013656,0.013656}%
\pgfsetfillcolor{currentfill}%
\pgfsetfillopacity{0.900000}%
\pgfsetlinewidth{0.501875pt}%
\definecolor{currentstroke}{rgb}{0.200000,0.200000,0.200000}%
\pgfsetstrokecolor{currentstroke}%
\pgfsetdash{}{0pt}%
\pgfpathmoveto{\pgfqpoint{1.891916in}{2.955326in}}%
\pgfpathlineto{\pgfqpoint{1.942758in}{2.906047in}}%
\pgfpathlineto{\pgfqpoint{1.976290in}{2.884902in}}%
\pgfpathlineto{\pgfqpoint{1.925440in}{2.934215in}}%
\pgfpathlineto{\pgfqpoint{1.891916in}{2.955326in}}%
\pgfpathclose%
\pgfusepath{stroke,fill}%
\end{pgfscope}%
\begin{pgfscope}%
\pgfpathrectangle{\pgfqpoint{0.050000in}{0.050000in}}{\pgfqpoint{4.900000in}{4.900000in}}%
\pgfusepath{clip}%
\pgfsetbuttcap%
\pgfsetroundjoin%
\definecolor{currentfill}{rgb}{0.195617,0.195617,0.195617}%
\pgfsetfillcolor{currentfill}%
\pgfsetfillopacity{0.900000}%
\pgfsetlinewidth{0.501875pt}%
\definecolor{currentstroke}{rgb}{0.200000,0.200000,0.200000}%
\pgfsetstrokecolor{currentstroke}%
\pgfsetdash{}{0pt}%
\pgfpathmoveto{\pgfqpoint{2.381109in}{2.560738in}}%
\pgfpathlineto{\pgfqpoint{2.430843in}{2.512351in}}%
\pgfpathlineto{\pgfqpoint{2.465202in}{2.490542in}}%
\pgfpathlineto{\pgfqpoint{2.415463in}{2.538961in}}%
\pgfpathlineto{\pgfqpoint{2.381109in}{2.560738in}}%
\pgfpathclose%
\pgfusepath{stroke,fill}%
\end{pgfscope}%
\begin{pgfscope}%
\pgfpathrectangle{\pgfqpoint{0.050000in}{0.050000in}}{\pgfqpoint{4.900000in}{4.900000in}}%
\pgfusepath{clip}%
\pgfsetbuttcap%
\pgfsetroundjoin%
\definecolor{currentfill}{rgb}{0.086490,0.086490,0.086490}%
\pgfsetfillcolor{currentfill}%
\pgfsetfillopacity{0.900000}%
\pgfsetlinewidth{0.501875pt}%
\definecolor{currentstroke}{rgb}{0.200000,0.200000,0.200000}%
\pgfsetstrokecolor{currentstroke}%
\pgfsetdash{}{0pt}%
\pgfpathmoveto{\pgfqpoint{2.094325in}{2.793244in}}%
\pgfpathlineto{\pgfqpoint{2.144725in}{2.744315in}}%
\pgfpathlineto{\pgfqpoint{2.178611in}{2.722890in}}%
\pgfpathlineto{\pgfqpoint{2.128204in}{2.771852in}}%
\pgfpathlineto{\pgfqpoint{2.094325in}{2.793244in}}%
\pgfpathclose%
\pgfusepath{stroke,fill}%
\end{pgfscope}%
\begin{pgfscope}%
\pgfpathrectangle{\pgfqpoint{0.050000in}{0.050000in}}{\pgfqpoint{4.900000in}{4.900000in}}%
\pgfusepath{clip}%
\pgfsetbuttcap%
\pgfsetroundjoin%
\definecolor{currentfill}{rgb}{0.156863,0.156863,0.156863}%
\pgfsetfillcolor{currentfill}%
\pgfsetfillopacity{0.900000}%
\pgfsetlinewidth{0.501875pt}%
\definecolor{currentstroke}{rgb}{0.200000,0.200000,0.200000}%
\pgfsetstrokecolor{currentstroke}%
\pgfsetdash{}{0pt}%
\pgfpathmoveto{\pgfqpoint{2.296912in}{2.631022in}}%
\pgfpathlineto{\pgfqpoint{2.346870in}{2.582444in}}%
\pgfpathlineto{\pgfqpoint{2.381109in}{2.560738in}}%
\pgfpathlineto{\pgfqpoint{2.331145in}{2.609349in}}%
\pgfpathlineto{\pgfqpoint{2.296912in}{2.631022in}}%
\pgfpathclose%
\pgfusepath{stroke,fill}%
\end{pgfscope}%
\begin{pgfscope}%
\pgfpathrectangle{\pgfqpoint{0.050000in}{0.050000in}}{\pgfqpoint{4.900000in}{4.900000in}}%
\pgfusepath{clip}%
\pgfsetbuttcap%
\pgfsetroundjoin%
\definecolor{currentfill}{rgb}{0.050073,0.050073,0.050073}%
\pgfsetfillcolor{currentfill}%
\pgfsetfillopacity{0.900000}%
\pgfsetlinewidth{0.501875pt}%
\definecolor{currentstroke}{rgb}{0.200000,0.200000,0.200000}%
\pgfsetstrokecolor{currentstroke}%
\pgfsetdash{}{0pt}%
\pgfpathmoveto{\pgfqpoint{2.009935in}{2.863686in}}%
\pgfpathlineto{\pgfqpoint{2.060560in}{2.814565in}}%
\pgfpathlineto{\pgfqpoint{2.094325in}{2.793244in}}%
\pgfpathlineto{\pgfqpoint{2.043693in}{2.842399in}}%
\pgfpathlineto{\pgfqpoint{2.009935in}{2.863686in}}%
\pgfpathclose%
\pgfusepath{stroke,fill}%
\end{pgfscope}%
\begin{pgfscope}%
\pgfpathrectangle{\pgfqpoint{0.050000in}{0.050000in}}{\pgfqpoint{4.900000in}{4.900000in}}%
\pgfusepath{clip}%
\pgfsetbuttcap%
\pgfsetroundjoin%
\definecolor{currentfill}{rgb}{0.118354,0.118354,0.118354}%
\pgfsetfillcolor{currentfill}%
\pgfsetfillopacity{0.900000}%
\pgfsetlinewidth{0.501875pt}%
\definecolor{currentstroke}{rgb}{0.200000,0.200000,0.200000}%
\pgfsetstrokecolor{currentstroke}%
\pgfsetdash{}{0pt}%
\pgfpathmoveto{\pgfqpoint{2.212610in}{2.701393in}}%
\pgfpathlineto{\pgfqpoint{2.262792in}{2.652623in}}%
\pgfpathlineto{\pgfqpoint{2.296912in}{2.631022in}}%
\pgfpathlineto{\pgfqpoint{2.246723in}{2.679825in}}%
\pgfpathlineto{\pgfqpoint{2.212610in}{2.701393in}}%
\pgfpathclose%
\pgfusepath{stroke,fill}%
\end{pgfscope}%
\begin{pgfscope}%
\pgfpathrectangle{\pgfqpoint{0.050000in}{0.050000in}}{\pgfqpoint{4.900000in}{4.900000in}}%
\pgfusepath{clip}%
\pgfsetbuttcap%
\pgfsetroundjoin%
\definecolor{currentfill}{rgb}{0.018208,0.018208,0.018208}%
\pgfsetfillcolor{currentfill}%
\pgfsetfillopacity{0.900000}%
\pgfsetlinewidth{0.501875pt}%
\definecolor{currentstroke}{rgb}{0.200000,0.200000,0.200000}%
\pgfsetstrokecolor{currentstroke}%
\pgfsetdash{}{0pt}%
\pgfpathmoveto{\pgfqpoint{1.925440in}{2.934215in}}%
\pgfpathlineto{\pgfqpoint{1.976290in}{2.884902in}}%
\pgfpathlineto{\pgfqpoint{2.009935in}{2.863686in}}%
\pgfpathlineto{\pgfqpoint{1.959077in}{2.913033in}}%
\pgfpathlineto{\pgfqpoint{1.925440in}{2.934215in}}%
\pgfpathclose%
\pgfusepath{stroke,fill}%
\end{pgfscope}%
\begin{pgfscope}%
\pgfpathrectangle{\pgfqpoint{0.050000in}{0.050000in}}{\pgfqpoint{4.900000in}{4.900000in}}%
\pgfusepath{clip}%
\pgfsetbuttcap%
\pgfsetroundjoin%
\definecolor{currentfill}{rgb}{0.195617,0.195617,0.195617}%
\pgfsetfillcolor{currentfill}%
\pgfsetfillopacity{0.900000}%
\pgfsetlinewidth{0.501875pt}%
\definecolor{currentstroke}{rgb}{0.200000,0.200000,0.200000}%
\pgfsetstrokecolor{currentstroke}%
\pgfsetdash{}{0pt}%
\pgfpathmoveto{\pgfqpoint{2.415463in}{2.538961in}}%
\pgfpathlineto{\pgfqpoint{2.465202in}{2.490542in}}%
\pgfpathlineto{\pgfqpoint{2.499675in}{2.468660in}}%
\pgfpathlineto{\pgfqpoint{2.449932in}{2.517110in}}%
\pgfpathlineto{\pgfqpoint{2.415463in}{2.538961in}}%
\pgfpathclose%
\pgfusepath{stroke,fill}%
\end{pgfscope}%
\begin{pgfscope}%
\pgfpathrectangle{\pgfqpoint{0.050000in}{0.050000in}}{\pgfqpoint{4.900000in}{4.900000in}}%
\pgfusepath{clip}%
\pgfsetbuttcap%
\pgfsetroundjoin%
\definecolor{currentfill}{rgb}{0.086490,0.086490,0.086490}%
\pgfsetfillcolor{currentfill}%
\pgfsetfillopacity{0.900000}%
\pgfsetlinewidth{0.501875pt}%
\definecolor{currentstroke}{rgb}{0.200000,0.200000,0.200000}%
\pgfsetstrokecolor{currentstroke}%
\pgfsetdash{}{0pt}%
\pgfpathmoveto{\pgfqpoint{2.128204in}{2.771852in}}%
\pgfpathlineto{\pgfqpoint{2.178611in}{2.722890in}}%
\pgfpathlineto{\pgfqpoint{2.212610in}{2.701393in}}%
\pgfpathlineto{\pgfqpoint{2.162196in}{2.750388in}}%
\pgfpathlineto{\pgfqpoint{2.128204in}{2.771852in}}%
\pgfpathclose%
\pgfusepath{stroke,fill}%
\end{pgfscope}%
\begin{pgfscope}%
\pgfpathrectangle{\pgfqpoint{0.050000in}{0.050000in}}{\pgfqpoint{4.900000in}{4.900000in}}%
\pgfusepath{clip}%
\pgfsetbuttcap%
\pgfsetroundjoin%
\definecolor{currentfill}{rgb}{0.156863,0.156863,0.156863}%
\pgfsetfillcolor{currentfill}%
\pgfsetfillopacity{0.900000}%
\pgfsetlinewidth{0.501875pt}%
\definecolor{currentstroke}{rgb}{0.200000,0.200000,0.200000}%
\pgfsetstrokecolor{currentstroke}%
\pgfsetdash{}{0pt}%
\pgfpathmoveto{\pgfqpoint{2.331145in}{2.609349in}}%
\pgfpathlineto{\pgfqpoint{2.381109in}{2.560738in}}%
\pgfpathlineto{\pgfqpoint{2.415463in}{2.538961in}}%
\pgfpathlineto{\pgfqpoint{2.365495in}{2.587602in}}%
\pgfpathlineto{\pgfqpoint{2.331145in}{2.609349in}}%
\pgfpathclose%
\pgfusepath{stroke,fill}%
\end{pgfscope}%
\begin{pgfscope}%
\pgfpathrectangle{\pgfqpoint{0.050000in}{0.050000in}}{\pgfqpoint{4.900000in}{4.900000in}}%
\pgfusepath{clip}%
\pgfsetbuttcap%
\pgfsetroundjoin%
\definecolor{currentfill}{rgb}{0.050073,0.050073,0.050073}%
\pgfsetfillcolor{currentfill}%
\pgfsetfillopacity{0.900000}%
\pgfsetlinewidth{0.501875pt}%
\definecolor{currentstroke}{rgb}{0.200000,0.200000,0.200000}%
\pgfsetstrokecolor{currentstroke}%
\pgfsetdash{}{0pt}%
\pgfpathmoveto{\pgfqpoint{2.043693in}{2.842399in}}%
\pgfpathlineto{\pgfqpoint{2.094325in}{2.793244in}}%
\pgfpathlineto{\pgfqpoint{2.128204in}{2.771852in}}%
\pgfpathlineto{\pgfqpoint{2.077564in}{2.821040in}}%
\pgfpathlineto{\pgfqpoint{2.043693in}{2.842399in}}%
\pgfpathclose%
\pgfusepath{stroke,fill}%
\end{pgfscope}%
\begin{pgfscope}%
\pgfpathrectangle{\pgfqpoint{0.050000in}{0.050000in}}{\pgfqpoint{4.900000in}{4.900000in}}%
\pgfusepath{clip}%
\pgfsetbuttcap%
\pgfsetroundjoin%
\definecolor{currentfill}{rgb}{0.118354,0.118354,0.118354}%
\pgfsetfillcolor{currentfill}%
\pgfsetfillopacity{0.900000}%
\pgfsetlinewidth{0.501875pt}%
\definecolor{currentstroke}{rgb}{0.200000,0.200000,0.200000}%
\pgfsetstrokecolor{currentstroke}%
\pgfsetdash{}{0pt}%
\pgfpathmoveto{\pgfqpoint{2.246723in}{2.679825in}}%
\pgfpathlineto{\pgfqpoint{2.296912in}{2.631022in}}%
\pgfpathlineto{\pgfqpoint{2.331145in}{2.609349in}}%
\pgfpathlineto{\pgfqpoint{2.280952in}{2.658183in}}%
\pgfpathlineto{\pgfqpoint{2.246723in}{2.679825in}}%
\pgfpathclose%
\pgfusepath{stroke,fill}%
\end{pgfscope}%
\begin{pgfscope}%
\pgfpathrectangle{\pgfqpoint{0.050000in}{0.050000in}}{\pgfqpoint{4.900000in}{4.900000in}}%
\pgfusepath{clip}%
\pgfsetbuttcap%
\pgfsetroundjoin%
\definecolor{currentfill}{rgb}{0.018208,0.018208,0.018208}%
\pgfsetfillcolor{currentfill}%
\pgfsetfillopacity{0.900000}%
\pgfsetlinewidth{0.501875pt}%
\definecolor{currentstroke}{rgb}{0.200000,0.200000,0.200000}%
\pgfsetstrokecolor{currentstroke}%
\pgfsetdash{}{0pt}%
\pgfpathmoveto{\pgfqpoint{1.959077in}{2.913033in}}%
\pgfpathlineto{\pgfqpoint{2.009935in}{2.863686in}}%
\pgfpathlineto{\pgfqpoint{2.043693in}{2.842399in}}%
\pgfpathlineto{\pgfqpoint{1.992827in}{2.891780in}}%
\pgfpathlineto{\pgfqpoint{1.959077in}{2.913033in}}%
\pgfpathclose%
\pgfusepath{stroke,fill}%
\end{pgfscope}%
\begin{pgfscope}%
\pgfpathrectangle{\pgfqpoint{0.050000in}{0.050000in}}{\pgfqpoint{4.900000in}{4.900000in}}%
\pgfusepath{clip}%
\pgfsetbuttcap%
\pgfsetroundjoin%
\definecolor{currentfill}{rgb}{0.086490,0.086490,0.086490}%
\pgfsetfillcolor{currentfill}%
\pgfsetfillopacity{0.900000}%
\pgfsetlinewidth{0.501875pt}%
\definecolor{currentstroke}{rgb}{0.200000,0.200000,0.200000}%
\pgfsetstrokecolor{currentstroke}%
\pgfsetdash{}{0pt}%
\pgfpathmoveto{\pgfqpoint{2.162196in}{2.750388in}}%
\pgfpathlineto{\pgfqpoint{2.212610in}{2.701393in}}%
\pgfpathlineto{\pgfqpoint{2.246723in}{2.679825in}}%
\pgfpathlineto{\pgfqpoint{2.196304in}{2.728852in}}%
\pgfpathlineto{\pgfqpoint{2.162196in}{2.750388in}}%
\pgfpathclose%
\pgfusepath{stroke,fill}%
\end{pgfscope}%
\begin{pgfscope}%
\pgfpathrectangle{\pgfqpoint{0.050000in}{0.050000in}}{\pgfqpoint{4.900000in}{4.900000in}}%
\pgfusepath{clip}%
\pgfsetbuttcap%
\pgfsetroundjoin%
\definecolor{currentfill}{rgb}{0.156863,0.156863,0.156863}%
\pgfsetfillcolor{currentfill}%
\pgfsetfillopacity{0.900000}%
\pgfsetlinewidth{0.501875pt}%
\definecolor{currentstroke}{rgb}{0.200000,0.200000,0.200000}%
\pgfsetstrokecolor{currentstroke}%
\pgfsetdash{}{0pt}%
\pgfpathmoveto{\pgfqpoint{2.365495in}{2.587602in}}%
\pgfpathlineto{\pgfqpoint{2.415463in}{2.538961in}}%
\pgfpathlineto{\pgfqpoint{2.449932in}{2.517110in}}%
\pgfpathlineto{\pgfqpoint{2.399960in}{2.565783in}}%
\pgfpathlineto{\pgfqpoint{2.365495in}{2.587602in}}%
\pgfpathclose%
\pgfusepath{stroke,fill}%
\end{pgfscope}%
\begin{pgfscope}%
\pgfpathrectangle{\pgfqpoint{0.050000in}{0.050000in}}{\pgfqpoint{4.900000in}{4.900000in}}%
\pgfusepath{clip}%
\pgfsetbuttcap%
\pgfsetroundjoin%
\definecolor{currentfill}{rgb}{0.050073,0.050073,0.050073}%
\pgfsetfillcolor{currentfill}%
\pgfsetfillopacity{0.900000}%
\pgfsetlinewidth{0.501875pt}%
\definecolor{currentstroke}{rgb}{0.200000,0.200000,0.200000}%
\pgfsetstrokecolor{currentstroke}%
\pgfsetdash{}{0pt}%
\pgfpathmoveto{\pgfqpoint{2.077564in}{2.821040in}}%
\pgfpathlineto{\pgfqpoint{2.128204in}{2.771852in}}%
\pgfpathlineto{\pgfqpoint{2.162196in}{2.750388in}}%
\pgfpathlineto{\pgfqpoint{2.111550in}{2.799609in}}%
\pgfpathlineto{\pgfqpoint{2.077564in}{2.821040in}}%
\pgfpathclose%
\pgfusepath{stroke,fill}%
\end{pgfscope}%
\begin{pgfscope}%
\pgfpathrectangle{\pgfqpoint{0.050000in}{0.050000in}}{\pgfqpoint{4.900000in}{4.900000in}}%
\pgfusepath{clip}%
\pgfsetbuttcap%
\pgfsetroundjoin%
\definecolor{currentfill}{rgb}{0.118354,0.118354,0.118354}%
\pgfsetfillcolor{currentfill}%
\pgfsetfillopacity{0.900000}%
\pgfsetlinewidth{0.501875pt}%
\definecolor{currentstroke}{rgb}{0.200000,0.200000,0.200000}%
\pgfsetstrokecolor{currentstroke}%
\pgfsetdash{}{0pt}%
\pgfpathmoveto{\pgfqpoint{2.280952in}{2.658183in}}%
\pgfpathlineto{\pgfqpoint{2.331145in}{2.609349in}}%
\pgfpathlineto{\pgfqpoint{2.365495in}{2.587602in}}%
\pgfpathlineto{\pgfqpoint{2.315296in}{2.636469in}}%
\pgfpathlineto{\pgfqpoint{2.280952in}{2.658183in}}%
\pgfpathclose%
\pgfusepath{stroke,fill}%
\end{pgfscope}%
\begin{pgfscope}%
\pgfpathrectangle{\pgfqpoint{0.050000in}{0.050000in}}{\pgfqpoint{4.900000in}{4.900000in}}%
\pgfusepath{clip}%
\pgfsetbuttcap%
\pgfsetroundjoin%
\definecolor{currentfill}{rgb}{0.018208,0.018208,0.018208}%
\pgfsetfillcolor{currentfill}%
\pgfsetfillopacity{0.900000}%
\pgfsetlinewidth{0.501875pt}%
\definecolor{currentstroke}{rgb}{0.200000,0.200000,0.200000}%
\pgfsetstrokecolor{currentstroke}%
\pgfsetdash{}{0pt}%
\pgfpathmoveto{\pgfqpoint{1.992827in}{2.891780in}}%
\pgfpathlineto{\pgfqpoint{2.043693in}{2.842399in}}%
\pgfpathlineto{\pgfqpoint{2.077564in}{2.821040in}}%
\pgfpathlineto{\pgfqpoint{2.026691in}{2.870455in}}%
\pgfpathlineto{\pgfqpoint{1.992827in}{2.891780in}}%
\pgfpathclose%
\pgfusepath{stroke,fill}%
\end{pgfscope}%
\begin{pgfscope}%
\pgfpathrectangle{\pgfqpoint{0.050000in}{0.050000in}}{\pgfqpoint{4.900000in}{4.900000in}}%
\pgfusepath{clip}%
\pgfsetbuttcap%
\pgfsetroundjoin%
\definecolor{currentfill}{rgb}{0.086490,0.086490,0.086490}%
\pgfsetfillcolor{currentfill}%
\pgfsetfillopacity{0.900000}%
\pgfsetlinewidth{0.501875pt}%
\definecolor{currentstroke}{rgb}{0.200000,0.200000,0.200000}%
\pgfsetstrokecolor{currentstroke}%
\pgfsetdash{}{0pt}%
\pgfpathmoveto{\pgfqpoint{2.196304in}{2.728852in}}%
\pgfpathlineto{\pgfqpoint{2.246723in}{2.679825in}}%
\pgfpathlineto{\pgfqpoint{2.280952in}{2.658183in}}%
\pgfpathlineto{\pgfqpoint{2.230526in}{2.707244in}}%
\pgfpathlineto{\pgfqpoint{2.196304in}{2.728852in}}%
\pgfpathclose%
\pgfusepath{stroke,fill}%
\end{pgfscope}%
\begin{pgfscope}%
\pgfpathrectangle{\pgfqpoint{0.050000in}{0.050000in}}{\pgfqpoint{4.900000in}{4.900000in}}%
\pgfusepath{clip}%
\pgfsetbuttcap%
\pgfsetroundjoin%
\definecolor{currentfill}{rgb}{0.050073,0.050073,0.050073}%
\pgfsetfillcolor{currentfill}%
\pgfsetfillopacity{0.900000}%
\pgfsetlinewidth{0.501875pt}%
\definecolor{currentstroke}{rgb}{0.200000,0.200000,0.200000}%
\pgfsetstrokecolor{currentstroke}%
\pgfsetdash{}{0pt}%
\pgfpathmoveto{\pgfqpoint{2.111550in}{2.799609in}}%
\pgfpathlineto{\pgfqpoint{2.162196in}{2.750388in}}%
\pgfpathlineto{\pgfqpoint{2.196304in}{2.728852in}}%
\pgfpathlineto{\pgfqpoint{2.145651in}{2.778107in}}%
\pgfpathlineto{\pgfqpoint{2.111550in}{2.799609in}}%
\pgfpathclose%
\pgfusepath{stroke,fill}%
\end{pgfscope}%
\begin{pgfscope}%
\pgfpathrectangle{\pgfqpoint{0.050000in}{0.050000in}}{\pgfqpoint{4.900000in}{4.900000in}}%
\pgfusepath{clip}%
\pgfsetbuttcap%
\pgfsetroundjoin%
\definecolor{currentfill}{rgb}{0.122907,0.122907,0.122907}%
\pgfsetfillcolor{currentfill}%
\pgfsetfillopacity{0.900000}%
\pgfsetlinewidth{0.501875pt}%
\definecolor{currentstroke}{rgb}{0.200000,0.200000,0.200000}%
\pgfsetstrokecolor{currentstroke}%
\pgfsetdash{}{0pt}%
\pgfpathmoveto{\pgfqpoint{2.315296in}{2.636469in}}%
\pgfpathlineto{\pgfqpoint{2.365495in}{2.587602in}}%
\pgfpathlineto{\pgfqpoint{2.399960in}{2.565783in}}%
\pgfpathlineto{\pgfqpoint{2.349755in}{2.614682in}}%
\pgfpathlineto{\pgfqpoint{2.315296in}{2.636469in}}%
\pgfpathclose%
\pgfusepath{stroke,fill}%
\end{pgfscope}%
\begin{pgfscope}%
\pgfpathrectangle{\pgfqpoint{0.050000in}{0.050000in}}{\pgfqpoint{4.900000in}{4.900000in}}%
\pgfusepath{clip}%
\pgfsetbuttcap%
\pgfsetroundjoin%
\definecolor{currentfill}{rgb}{0.018208,0.018208,0.018208}%
\pgfsetfillcolor{currentfill}%
\pgfsetfillopacity{0.900000}%
\pgfsetlinewidth{0.501875pt}%
\definecolor{currentstroke}{rgb}{0.200000,0.200000,0.200000}%
\pgfsetstrokecolor{currentstroke}%
\pgfsetdash{}{0pt}%
\pgfpathmoveto{\pgfqpoint{2.026691in}{2.870455in}}%
\pgfpathlineto{\pgfqpoint{2.077564in}{2.821040in}}%
\pgfpathlineto{\pgfqpoint{2.111550in}{2.799609in}}%
\pgfpathlineto{\pgfqpoint{2.060669in}{2.849058in}}%
\pgfpathlineto{\pgfqpoint{2.026691in}{2.870455in}}%
\pgfpathclose%
\pgfusepath{stroke,fill}%
\end{pgfscope}%
\begin{pgfscope}%
\pgfpathrectangle{\pgfqpoint{0.050000in}{0.050000in}}{\pgfqpoint{4.900000in}{4.900000in}}%
\pgfusepath{clip}%
\pgfsetbuttcap%
\pgfsetroundjoin%
\definecolor{currentfill}{rgb}{0.086490,0.086490,0.086490}%
\pgfsetfillcolor{currentfill}%
\pgfsetfillopacity{0.900000}%
\pgfsetlinewidth{0.501875pt}%
\definecolor{currentstroke}{rgb}{0.200000,0.200000,0.200000}%
\pgfsetstrokecolor{currentstroke}%
\pgfsetdash{}{0pt}%
\pgfpathmoveto{\pgfqpoint{2.230526in}{2.707244in}}%
\pgfpathlineto{\pgfqpoint{2.280952in}{2.658183in}}%
\pgfpathlineto{\pgfqpoint{2.315296in}{2.636469in}}%
\pgfpathlineto{\pgfqpoint{2.264864in}{2.685562in}}%
\pgfpathlineto{\pgfqpoint{2.230526in}{2.707244in}}%
\pgfpathclose%
\pgfusepath{stroke,fill}%
\end{pgfscope}%
\begin{pgfscope}%
\pgfpathrectangle{\pgfqpoint{0.050000in}{0.050000in}}{\pgfqpoint{4.900000in}{4.900000in}}%
\pgfusepath{clip}%
\pgfsetbuttcap%
\pgfsetroundjoin%
\definecolor{currentfill}{rgb}{0.054625,0.054625,0.054625}%
\pgfsetfillcolor{currentfill}%
\pgfsetfillopacity{0.900000}%
\pgfsetlinewidth{0.501875pt}%
\definecolor{currentstroke}{rgb}{0.200000,0.200000,0.200000}%
\pgfsetstrokecolor{currentstroke}%
\pgfsetdash{}{0pt}%
\pgfpathmoveto{\pgfqpoint{2.145651in}{2.778107in}}%
\pgfpathlineto{\pgfqpoint{2.196304in}{2.728852in}}%
\pgfpathlineto{\pgfqpoint{2.230526in}{2.707244in}}%
\pgfpathlineto{\pgfqpoint{2.179867in}{2.756531in}}%
\pgfpathlineto{\pgfqpoint{2.145651in}{2.778107in}}%
\pgfpathclose%
\pgfusepath{stroke,fill}%
\end{pgfscope}%
\begin{pgfscope}%
\pgfpathrectangle{\pgfqpoint{0.050000in}{0.050000in}}{\pgfqpoint{4.900000in}{4.900000in}}%
\pgfusepath{clip}%
\pgfsetbuttcap%
\pgfsetroundjoin%
\definecolor{currentfill}{rgb}{0.018208,0.018208,0.018208}%
\pgfsetfillcolor{currentfill}%
\pgfsetfillopacity{0.900000}%
\pgfsetlinewidth{0.501875pt}%
\definecolor{currentstroke}{rgb}{0.200000,0.200000,0.200000}%
\pgfsetstrokecolor{currentstroke}%
\pgfsetdash{}{0pt}%
\pgfpathmoveto{\pgfqpoint{2.060669in}{2.849058in}}%
\pgfpathlineto{\pgfqpoint{2.111550in}{2.799609in}}%
\pgfpathlineto{\pgfqpoint{2.145651in}{2.778107in}}%
\pgfpathlineto{\pgfqpoint{2.094763in}{2.827589in}}%
\pgfpathlineto{\pgfqpoint{2.060669in}{2.849058in}}%
\pgfpathclose%
\pgfusepath{stroke,fill}%
\end{pgfscope}%
\begin{pgfscope}%
\pgfpathrectangle{\pgfqpoint{0.050000in}{0.050000in}}{\pgfqpoint{4.900000in}{4.900000in}}%
\pgfusepath{clip}%
\pgfsetbuttcap%
\pgfsetroundjoin%
\definecolor{currentfill}{rgb}{0.086490,0.086490,0.086490}%
\pgfsetfillcolor{currentfill}%
\pgfsetfillopacity{0.900000}%
\pgfsetlinewidth{0.501875pt}%
\definecolor{currentstroke}{rgb}{0.200000,0.200000,0.200000}%
\pgfsetstrokecolor{currentstroke}%
\pgfsetdash{}{0pt}%
\pgfpathmoveto{\pgfqpoint{2.264864in}{2.685562in}}%
\pgfpathlineto{\pgfqpoint{2.315296in}{2.636469in}}%
\pgfpathlineto{\pgfqpoint{2.349755in}{2.614682in}}%
\pgfpathlineto{\pgfqpoint{2.299319in}{2.663807in}}%
\pgfpathlineto{\pgfqpoint{2.264864in}{2.685562in}}%
\pgfpathclose%
\pgfusepath{stroke,fill}%
\end{pgfscope}%
\begin{pgfscope}%
\pgfpathrectangle{\pgfqpoint{0.050000in}{0.050000in}}{\pgfqpoint{4.900000in}{4.900000in}}%
\pgfusepath{clip}%
\pgfsetbuttcap%
\pgfsetroundjoin%
\definecolor{currentfill}{rgb}{0.054625,0.054625,0.054625}%
\pgfsetfillcolor{currentfill}%
\pgfsetfillopacity{0.900000}%
\pgfsetlinewidth{0.501875pt}%
\definecolor{currentstroke}{rgb}{0.200000,0.200000,0.200000}%
\pgfsetstrokecolor{currentstroke}%
\pgfsetdash{}{0pt}%
\pgfpathmoveto{\pgfqpoint{2.179867in}{2.756531in}}%
\pgfpathlineto{\pgfqpoint{2.230526in}{2.707244in}}%
\pgfpathlineto{\pgfqpoint{2.264864in}{2.685562in}}%
\pgfpathlineto{\pgfqpoint{2.214199in}{2.734883in}}%
\pgfpathlineto{\pgfqpoint{2.179867in}{2.756531in}}%
\pgfpathclose%
\pgfusepath{stroke,fill}%
\end{pgfscope}%
\begin{pgfscope}%
\pgfpathrectangle{\pgfqpoint{0.050000in}{0.050000in}}{\pgfqpoint{4.900000in}{4.900000in}}%
\pgfusepath{clip}%
\pgfsetbuttcap%
\pgfsetroundjoin%
\definecolor{currentfill}{rgb}{0.018208,0.018208,0.018208}%
\pgfsetfillcolor{currentfill}%
\pgfsetfillopacity{0.900000}%
\pgfsetlinewidth{0.501875pt}%
\definecolor{currentstroke}{rgb}{0.200000,0.200000,0.200000}%
\pgfsetstrokecolor{currentstroke}%
\pgfsetdash{}{0pt}%
\pgfpathmoveto{\pgfqpoint{2.094763in}{2.827589in}}%
\pgfpathlineto{\pgfqpoint{2.145651in}{2.778107in}}%
\pgfpathlineto{\pgfqpoint{2.179867in}{2.756531in}}%
\pgfpathlineto{\pgfqpoint{2.128972in}{2.806048in}}%
\pgfpathlineto{\pgfqpoint{2.094763in}{2.827589in}}%
\pgfpathclose%
\pgfusepath{stroke,fill}%
\end{pgfscope}%
\begin{pgfscope}%
\pgfpathrectangle{\pgfqpoint{0.050000in}{0.050000in}}{\pgfqpoint{4.900000in}{4.900000in}}%
\pgfusepath{clip}%
\pgfsetbuttcap%
\pgfsetroundjoin%
\definecolor{currentfill}{rgb}{0.054625,0.054625,0.054625}%
\pgfsetfillcolor{currentfill}%
\pgfsetfillopacity{0.900000}%
\pgfsetlinewidth{0.501875pt}%
\definecolor{currentstroke}{rgb}{0.200000,0.200000,0.200000}%
\pgfsetstrokecolor{currentstroke}%
\pgfsetdash{}{0pt}%
\pgfpathmoveto{\pgfqpoint{2.214199in}{2.734883in}}%
\pgfpathlineto{\pgfqpoint{2.264864in}{2.685562in}}%
\pgfpathlineto{\pgfqpoint{2.299319in}{2.663807in}}%
\pgfpathlineto{\pgfqpoint{2.248647in}{2.713161in}}%
\pgfpathlineto{\pgfqpoint{2.214199in}{2.734883in}}%
\pgfpathclose%
\pgfusepath{stroke,fill}%
\end{pgfscope}%
\begin{pgfscope}%
\pgfpathrectangle{\pgfqpoint{0.050000in}{0.050000in}}{\pgfqpoint{4.900000in}{4.900000in}}%
\pgfusepath{clip}%
\pgfsetbuttcap%
\pgfsetroundjoin%
\definecolor{currentfill}{rgb}{0.018208,0.018208,0.018208}%
\pgfsetfillcolor{currentfill}%
\pgfsetfillopacity{0.900000}%
\pgfsetlinewidth{0.501875pt}%
\definecolor{currentstroke}{rgb}{0.200000,0.200000,0.200000}%
\pgfsetstrokecolor{currentstroke}%
\pgfsetdash{}{0pt}%
\pgfpathmoveto{\pgfqpoint{2.128972in}{2.806048in}}%
\pgfpathlineto{\pgfqpoint{2.179867in}{2.756531in}}%
\pgfpathlineto{\pgfqpoint{2.214199in}{2.734883in}}%
\pgfpathlineto{\pgfqpoint{2.163297in}{2.784433in}}%
\pgfpathlineto{\pgfqpoint{2.128972in}{2.806048in}}%
\pgfpathclose%
\pgfusepath{stroke,fill}%
\end{pgfscope}%
\begin{pgfscope}%
\pgfpathrectangle{\pgfqpoint{0.050000in}{0.050000in}}{\pgfqpoint{4.900000in}{4.900000in}}%
\pgfusepath{clip}%
\pgfsetbuttcap%
\pgfsetroundjoin%
\definecolor{currentfill}{rgb}{0.018208,0.018208,0.018208}%
\pgfsetfillcolor{currentfill}%
\pgfsetfillopacity{0.900000}%
\pgfsetlinewidth{0.501875pt}%
\definecolor{currentstroke}{rgb}{0.200000,0.200000,0.200000}%
\pgfsetstrokecolor{currentstroke}%
\pgfsetdash{}{0pt}%
\pgfpathmoveto{\pgfqpoint{2.163297in}{2.784433in}}%
\pgfpathlineto{\pgfqpoint{2.214199in}{2.734883in}}%
\pgfpathlineto{\pgfqpoint{2.248647in}{2.713161in}}%
\pgfpathlineto{\pgfqpoint{2.197739in}{2.762745in}}%
\pgfpathlineto{\pgfqpoint{2.163297in}{2.784433in}}%
\pgfpathclose%
\pgfusepath{stroke,fill}%
\end{pgfscope}%
\end{pgfpicture}%
\makeatother%
\endgroup%

	\caption{Graph of the up-and-out put option price}\label{fig:upAndOutPutBarrierSurface}
\end{figure}

\clearpage
\subsubsection{Down-and-out barrier call option}
Let us now consider a down-and-out barrier call option on the underlying asset price \(S_t\) with exercise date \(T\), strike price \(K\), barrier level \(B\), and payoff
\begin{equation}
	C = {(S_T - K)}^+ \mathbbm{1}_{\left\{\min\limits_{0 \le t \le T} S_t > B\right\}} =
	\begin{dcases}
		S_T - K & \text{if } \min\limits_{0 \le t \le T} S_t > B,   \\
		0       & \text{if } \min\limits_{0 \le t \le T} S_t \le B,
	\end{dcases}
\end{equation}
with \(0 \le B \le K\). The down-and-out barrier call option is also called a \textit{Callable Bull Contract}, or a Bull CBBC with no residual value, or a turbo warrant with no rebate, in which \(B\) denotes the call level.

\begin{comment}
\begin{prop}{Down-and-Out Barrier Call Option Pricing}{DownAndOutBarrierCallOptionPricing}
	The price of the down-and-out barrier call option with maturity \(T\), strike price \(K\), and barrier level \(B\) is given by the following formulas, where \(\tau = T-t\). When \(B \le K\), we have
	\begin{align}
		\MoveEqLeft \e^{-(T-t)r} \E^* \left[ {(S_T - K)}^+ \mathbbm{1}_{\left\{ \min\limits_{0 \le t \le T} S_t > B \right\}} \bigg| \mathcal{F}_t \right] \nonumber                                                      \\
		 & = S_t \mathbbm{1}_{\{m_0^t > B\}} \Phi \left( d_1 {\left( \frac{S_t}{K}, T-t \right)} \right) - \e^{-(T-t)r} K \mathbbm{1}_{\{m_0^t > B\}} \Phi \left( d_2 {\left( \frac{S_t}{K}, T-t \right)} \right) \nonumber \\
		 & \quad - B \mathbbm{1}_{\{m_0^t > B\}} {\left( \frac{B}{S_t} \right)}^{2r/\sigma^2} \Phi \left( d_1 {\left( \frac{B^2}{K S_t}, T-t \right)} \right) \nonumber                                                     \\
		 & \quad + \e^{-(T-t)r} K \mathbbm{1}_{\{m_0^t > B\}} {\left( \frac{S_t}{B} \right)}^{1-2r/\sigma^2} \Phi \left( d_2 {\left( \frac{B^2}{K S_t}, T-t \right)} \right) \nonumber                                      \\
		 & = \mathbbm{1}_{\{m_0^t > B\}} \mathrm{Bl}(S_t, K, r, T-t, \sigma) \nonumber                                                                                                                                  \\
		 & \quad - S_t \mathbbm{1}_{\{m_0^t > B\}} {\left( \frac{B}{S_t} \right)}^{1 + 2r/\sigma^2} \Phi \left( d_1 {\left( \frac{B^2}{K S_t}, T-t \right)} \right) \nonumber                                               \\
		 & \quad + K \e^{-r(T-t)} \mathbbm{1}_{\{m_0^t > B\}} {\left( \frac{B}{S_t} \right)}^{-1 + 2r/\sigma^2} \Phi \left( d_2 {\left( \frac{B^2}{K S_t}, T-t \right)} \right) \nonumber                                   \\
		 & = \mathbbm{1}_{\{m_0^t > B\}} \mathrm{Bl}(S_t, K, r, T-t, \sigma) \nonumber                                                                                                                                  \\
		 & \quad - \mathbbm{1}_{\{m_0^t > B\}} {\left( \frac{S_t}{B} \right)}^{1 - 2r/\sigma^2} \mathrm{Bl} \left( \frac{B^2}{S_t}, K, r, T-t, \sigma \right),
	\end{align}
	for \(0 \le t \le T\). When \(B \ge K\), we find
	\begin{align}
		\MoveEqLeft \e^{-(T-t)r} \E^* \left[ {(S_T - K)}^+ \mathbbm{1}_{\left\{ \min\limits_{0 \le t \le T} S_t > B \right\}} \bigg| \mathcal{F}_t \right] \nonumber                                                      \\
		 & = S_t \mathbbm{1}_{\{m_0^t > B\}} \Phi \left( d_1 {\left( \frac{S_t}{B}, T-t \right)} \right) - \e^{-(T-t)r} K \mathbbm{1}_{\{m_0^t > B\}} \Phi \left( d_2 {\left( \frac{S_t}{B}, T-t \right)} \right) \nonumber \\
		 & \quad - B \mathbbm{1}_{\{m_0^t > B\}} {\left( \frac{B}{S_t} \right)}^{2r/\sigma^2} \Phi \left( d_1 {\left( \frac{B}{S_t}, T-t \right)} \right) \nonumber                                                         \\
		 & \quad + \e^{-(T-t)r} K \mathbbm{1}_{\{m_0^t > B\}} {\left( \frac{S_t}{B} \right)}^{1-2r/\sigma^2} \Phi \left( d_2 {\left( \frac{B}{S_t}, T-t \right)} \right),
	\end{align}
	for \(S_t > B\), \(0 \le t \le T\). Here \(m_0^t = \min\limits_{0 \le u \le t} S_u\) and we recall \(d_{1,2}(x, \tau)\) are defined as in \cref{eq:D1Definition,eq:D2Definition} with the spot price term \(S/K\) replaced by the ratio argument \(x\).
\end{prop}
\end{comment}

\begin{prop}{Down-and-Out Barrier Call Option Pricing}{DownAndOutBarrierCall}
	The price of the down-and-out barrier call option with maturity \(T\), strike price \(K\), and barrier level \(B\) is given by the following formulas, where \(\tau = T-t\).

	When \(B \le K\), we have:
	\begin{align}
		\MoveEqLeft C_{DO}(t, S_t) = \mathbbm{1}_{\{m_0^t > B\}} \mathrm{Bl}(S_t, K, r, \tau, \sigma) \nonumber                                                      \\
		 & - B \mathbbm{1}_{\{m_0^t > B\}} {\left( \frac{B}{S_t} \right)}^{1+2r/\sigma^2} \Phi \left( d_1 {\left( \frac{B^2}{K S_t}, \tau \right)} \right) \nonumber     \\
		 & + K \e^{-r\tau} \mathbbm{1}_{\{m_0^t > B\}} {\left( \frac{B}{S_t} \right)}^{-1+2r/\sigma^2} \Phi \left( d_2 {\left( \frac{B^2}{K S_t}, \tau \right)} \right).
	\end{align}

	When \(B \ge K\), we find:
	\begin{align}
		\MoveEqLeft C_{DO}(t, S_t) = S_t \mathbbm{1}_{\{m_0^t > B\}} \Phi \left( d_1 {\left( \frac{S_t}{B}, \tau \right)} \right) - K \e^{-r\tau} \mathbbm{1}_{\{m_0^t > B\}} \Phi \left( d_2 {\left( \frac{S_t}{B}, \tau \right)} \right) \nonumber \\
		 & - B \mathbbm{1}_{\{m_0^t > B\}} {\left( \frac{B}{S_t} \right)}^{2r/\sigma^2} \Phi \left( d_1 {\left( \frac{B}{S_t}, \tau \right)} \right) \nonumber                                                                                       \\
		 & + K \e^{-r\tau} \mathbbm{1}_{\{m_0^t > B\}} {\left( \frac{S_t}{B} \right)}^{1-2r/\sigma^2} \Phi \left( d_2 {\left( \frac{B}{S_t}, \tau \right)} \right).
	\end{align}
	Here \(\mathrm{Bl}\) denotes the standard Black-Scholes call price, and \(d_{1,2}(x, \tau)\) are defined as in \cref{eq:D1Definition,eq:D2Definition} with the spot price term replaced by the ratio argument \(x\). If \(m_0^t \le B\), the price is \(0\).
\end{prop}


\begin{figure}[H]
	\centering
	\input{plotting/downAndOutCallBarrierSurface.pgf}
	\caption{Graph of the down-and-out call option price}\label{fig:downAndOutCallBarrierSurface}
\end{figure}

\clearpage
\subsubsection{Down-and-out barrier put option}
The down-and-out barrier put option with exercise date \(T\), strike price \(K\), and barrier level \(B\) has the payoff:
\begin{equation}
	P = {(K - S_T)}^+ \mathbbm{1}_{\left\{\min_{0 \le t \le T} S_t > B\right\}} =
	\begin{dcases}
		K - S_T & \text{if } \min_{0 \le t \le T} S_t > B,   \\
		0       & \text{if } \min_{0 \le t \le T} S_t \le B.
	\end{dcases}
\end{equation}
For this option to have non-zero value, we typically require \(S_0 > B\). Since it is a put option, the intrinsic value is positive when \(S_T < K\). If \(B \ge K\), the condition \(S_T < K\) combined with the knock-out condition \(\min S_t > B \ge K\) makes the payoff always zero (as \(S_T > B \ge K \implies K - S_T < 0\)). Thus, we only consider the case \(B < K\).

\begin{comment}
\begin{prop}{Down-and-Out Barrier Put Option Pricing}{DownAndOutBarrierPutOptionPricing}
	The price of the down-and-out barrier put option with maturity \(T\), strike price \(K\), and barrier level \(B\) is given by the following formula, where \(\tau = T-t\). When \(B \le K\), we have
	\begin{align}
		\MoveEqLeft \e^{-(T-t)r} \E^* \left[ {(K - S_T)}^+ \mathbbm{1}_{\left\{ \min\limits_{0 \le t \le T} S_t > B \right\}} \bigg| \mathcal{F}_t \right] \nonumber                                                                                            \\
		 & = K \e^{-(T-t)r} \mathbbm{1}_{\{m_0^t > B\}} \left[ \Phi \left( -d_2 {\left( \frac{S_t}{K}, T-t \right)} \right) - \Phi \left( -d_2 {\left( \frac{S_t}{B}, T-t \right)} \right) \right] \nonumber                                                      \\
		 & \quad - S_t \mathbbm{1}_{\{m_0^t > B\}} \left[ \Phi \left( -d_1 {\left( \frac{S_t}{K}, T-t \right)} \right) - \Phi \left( -d_1 {\left( \frac{S_t}{B}, T-t \right)} \right) \right] \nonumber                                                           \\
		 & \quad - K \e^{-(T-t)r} \mathbbm{1}_{\{m_0^t > B\}} {\left( \frac{S_t}{B} \right)}^{1-2r/\sigma^2} \left[ \Phi \left( d_2 {\left( \frac{B^2}{K S_t}, T-t \right)} \right) - \Phi \left( d_2 {\left( \frac{B}{S_t}, T-t \right)} \right) \right] \nonumber \\
		 & \quad + S_t \mathbbm{1}_{\{m_0^t > B\}} {\left( \frac{B}{S_t} \right)}^{1+2r/\sigma^2} \left[ \Phi \left( d_1 {\left( \frac{B^2}{K S_t}, T-t \right)} \right) - \Phi \left( d_1 {\left( \frac{B}{S_t}, T-t \right)} \right) \right] \nonumber            \\
		 & = \mathbbm{1}_{\{m_0^t > B\}} \mathrm{Bl_{put}}(S_t, K, r, T-t, \sigma) \nonumber                                                                                                                                                                  \\
		 & \quad - \mathbbm{1}_{\{m_0^t > B\}} \mathrm{Bl_{put}}(S_t, B, r, T-t, \sigma) \nonumber                                                                                                                                                            \\
		 & \quad + (B - K) \e^{-r(T-t)} \mathbbm{1}_{\{m_0^t > B\}} {\left( \frac{S_t}{B} \right)}^{1-2r/\sigma^2} \Phi \left( d_2 {\left( \frac{B}{S_t}, T-t \right)} \right) \nonumber                                                                          \\
		 & \quad + \mathbbm{1}_{\{m_0^t > B\}} {\left( \frac{S_t}{B} \right)}^{1-2r/\sigma^2} \mathrm{Bl} \left( \frac{B^2}{S_t}, K, r, T-t, \sigma \right) \nonumber                                                                                           \\
		 & \quad - \mathbbm{1}_{\{m_0^t > B\}} {\left( \frac{S_t}{B} \right)}^{1-2r/\sigma^2} \mathrm{Bl} \left( \frac{B^2}{S_t}, B, r, T-t, \sigma \right),
	\end{align}
	for \(S_t > B\), \(0 \le t \le T\). If \(K < B\), the option price is 0.
\end{prop}
\end{comment}
\begin{prop}{Down-and-Out Barrier Put Option Pricing}{DownAndOutBarrierPut}
	The price of the down-and-out barrier put option with maturity \(T\), strike price \(K\), and barrier level \(B\) (where \(B \le K\)) is given by
	\begin{align}
		\MoveEqLeft P_{DO}(t, S_t) = K \e^{-r\tau} \mathbbm{1}_{\{m_0^t > B\}} \left[ \Phi \left( -d_2 {\left( \frac{S_t}{K}, \tau \right)} \right) - \Phi \left( -d_2 {\left( \frac{S_t}{B}, \tau \right)} \right) \right] \nonumber                        \\
		 & - S_t \mathbbm{1}_{\{m_0^t > B\}} \left[ \Phi \left( -d_1 {\left( \frac{S_t}{K}, \tau \right)} \right) - \Phi \left( -d_1 {\left( \frac{S_t}{B}, \tau \right)} \right) \right] \nonumber                                                          \\
		 & - K \e^{-r\tau} \mathbbm{1}_{\{m_0^t > B\}} {\left( \frac{S_t}{B} \right)}^{1-2r/\sigma^2} \left[ \Phi \left( d_2 {\left( \frac{B^2}{K S_t}, \tau \right)} \right) - \Phi \left( d_2 {\left( \frac{B}{S_t}, \tau \right)} \right) \right] \nonumber \\
		 & + S_t \mathbbm{1}_{\{m_0^t > B\}} {\left( \frac{B}{S_t} \right)}^{1+2r/\sigma^2} \left[ \Phi \left( d_1 {\left( \frac{B^2}{K S_t}, \tau \right)} \right) - \Phi \left( d_1 {\left( \frac{B}{S_t}, \tau \right)} \right) \right].
	\end{align}
	Here \(\tau = T-t\), and \(d_{1,2}(x, \tau)\) correspond to the definitions with the spot price term replaced by the ratio argument \(x\).

	If \(B > K\), the price vanishes: \(P_{DO}(t, S_t) = 0\).
\end{prop}

\begin{figure}[H]
	\centering
	%% Creator: Matplotlib, PGF backend
%%
%% To include the figure in your LaTeX document, write
%%   \input{<filename>.pgf}
%%
%% Make sure the required packages are loaded in your preamble
%%   \usepackage{pgf}
%%
%% Also ensure that all the required font packages are loaded; for instance,
%% the lmodern package is sometimes necessary when using math font.
%%   \usepackage{lmodern}
%%
%% Figures using additional raster images can only be included by \input if
%% they are in the same directory as the main LaTeX file. For loading figures
%% from other directories you can use the `import` package
%%   \usepackage{import}
%%
%% and then include the figures with
%%   \import{<path to file>}{<filename>.pgf}
%%
%% Matplotlib used the following preamble
%%   \def\mathdefault#1{#1}
%%   \everymath=\expandafter{\the\everymath\displaystyle}
%%   \IfFileExists{scrextend.sty}{
%%     \usepackage[fontsize=10.000000pt]{scrextend}
%%   }{
%%     \renewcommand{\normalsize}{\fontsize{10.000000}{12.000000}\selectfont}
%%     \normalsize
%%   }
%%   
%%   \makeatletter\@ifpackageloaded{underscore}{}{\usepackage[strings]{underscore}}\makeatother
%%
\begingroup%
\makeatletter%
\begin{pgfpicture}%
\pgfpathrectangle{\pgfpointorigin}{\pgfqpoint{5.000000in}{5.033252in}}%
\pgfusepath{use as bounding box, clip}%
\begin{pgfscope}%
\pgfsetbuttcap%
\pgfsetmiterjoin%
\definecolor{currentfill}{rgb}{1.000000,1.000000,1.000000}%
\pgfsetfillcolor{currentfill}%
\pgfsetlinewidth{0.000000pt}%
\definecolor{currentstroke}{rgb}{1.000000,1.000000,1.000000}%
\pgfsetstrokecolor{currentstroke}%
\pgfsetdash{}{0pt}%
\pgfpathmoveto{\pgfqpoint{0.000000in}{0.000000in}}%
\pgfpathlineto{\pgfqpoint{5.000000in}{0.000000in}}%
\pgfpathlineto{\pgfqpoint{5.000000in}{5.033252in}}%
\pgfpathlineto{\pgfqpoint{0.000000in}{5.033252in}}%
\pgfpathlineto{\pgfqpoint{0.000000in}{0.000000in}}%
\pgfpathclose%
\pgfusepath{fill}%
\end{pgfscope}%
\begin{pgfscope}%
\pgfsetbuttcap%
\pgfsetmiterjoin%
\definecolor{currentfill}{rgb}{1.000000,1.000000,1.000000}%
\pgfsetfillcolor{currentfill}%
\pgfsetlinewidth{0.000000pt}%
\definecolor{currentstroke}{rgb}{0.000000,0.000000,0.000000}%
\pgfsetstrokecolor{currentstroke}%
\pgfsetstrokeopacity{0.000000}%
\pgfsetdash{}{0pt}%
\pgfpathmoveto{\pgfqpoint{0.050000in}{0.083252in}}%
\pgfpathlineto{\pgfqpoint{4.950000in}{0.083252in}}%
\pgfpathlineto{\pgfqpoint{4.950000in}{4.983252in}}%
\pgfpathlineto{\pgfqpoint{0.050000in}{4.983252in}}%
\pgfpathlineto{\pgfqpoint{0.050000in}{0.083252in}}%
\pgfpathclose%
\pgfusepath{fill}%
\end{pgfscope}%
\begin{pgfscope}%
\pgfsetbuttcap%
\pgfsetmiterjoin%
\pgfsetlinewidth{1.003750pt}%
\definecolor{currentstroke}{rgb}{0.950000,0.950000,0.950000}%
\pgfsetstrokecolor{currentstroke}%
\pgfsetstrokeopacity{0.500000}%
\pgfsetdash{}{0pt}%
\pgfpathmoveto{\pgfqpoint{4.712449in}{1.291435in}}%
\pgfpathlineto{\pgfqpoint{3.094309in}{2.647792in}}%
\pgfpathlineto{\pgfqpoint{3.116802in}{4.603902in}}%
\pgfpathlineto{\pgfqpoint{4.812379in}{3.366547in}}%
\pgfusepath{stroke}%
\end{pgfscope}%
\begin{pgfscope}%
\pgfsetbuttcap%
\pgfsetmiterjoin%
\pgfsetlinewidth{1.003750pt}%
\definecolor{currentstroke}{rgb}{0.900000,0.900000,0.900000}%
\pgfsetstrokecolor{currentstroke}%
\pgfsetstrokeopacity{0.500000}%
\pgfsetdash{}{0pt}%
\pgfpathmoveto{\pgfqpoint{0.497771in}{1.893081in}}%
\pgfpathlineto{\pgfqpoint{3.094309in}{2.647792in}}%
\pgfpathlineto{\pgfqpoint{3.116802in}{4.603902in}}%
\pgfpathlineto{\pgfqpoint{0.405110in}{3.916567in}}%
\pgfusepath{stroke}%
\end{pgfscope}%
\begin{pgfscope}%
\pgfsetbuttcap%
\pgfsetmiterjoin%
\pgfsetlinewidth{1.003750pt}%
\definecolor{currentstroke}{rgb}{0.925000,0.925000,0.925000}%
\pgfsetstrokecolor{currentstroke}%
\pgfsetstrokeopacity{0.500000}%
\pgfsetdash{}{0pt}%
\pgfpathmoveto{\pgfqpoint{1.959992in}{0.392483in}}%
\pgfpathlineto{\pgfqpoint{4.712449in}{1.291435in}}%
\pgfpathlineto{\pgfqpoint{3.094309in}{2.647792in}}%
\pgfpathlineto{\pgfqpoint{0.497771in}{1.893081in}}%
\pgfusepath{stroke}%
\end{pgfscope}%
\begin{pgfscope}%
\pgfsetbuttcap%
\pgfsetroundjoin%
\pgfsetlinewidth{0.501875pt}%
\definecolor{currentstroke}{rgb}{0.850000,0.850000,0.850000}%
\pgfsetstrokecolor{currentstroke}%
\pgfsetdash{}{0pt}%
\pgfpathmoveto{\pgfqpoint{1.997811in}{0.404835in}}%
\pgfpathlineto{\pgfqpoint{0.533293in}{1.903406in}}%
\pgfpathlineto{\pgfqpoint{0.442284in}{3.925989in}}%
\pgfusepath{stroke}%
\end{pgfscope}%
\begin{pgfscope}%
\pgfsetbuttcap%
\pgfsetroundjoin%
\pgfsetlinewidth{0.501875pt}%
\definecolor{currentstroke}{rgb}{0.850000,0.850000,0.850000}%
\pgfsetstrokecolor{currentstroke}%
\pgfsetdash{}{0pt}%
\pgfpathmoveto{\pgfqpoint{2.577814in}{0.594264in}}%
\pgfpathlineto{\pgfqpoint{1.078602in}{2.061906in}}%
\pgfpathlineto{\pgfqpoint{1.012695in}{4.070572in}}%
\pgfusepath{stroke}%
\end{pgfscope}%
\begin{pgfscope}%
\pgfsetbuttcap%
\pgfsetroundjoin%
\pgfsetlinewidth{0.501875pt}%
\definecolor{currentstroke}{rgb}{0.850000,0.850000,0.850000}%
\pgfsetstrokecolor{currentstroke}%
\pgfsetdash{}{0pt}%
\pgfpathmoveto{\pgfqpoint{3.145386in}{0.779633in}}%
\pgfpathlineto{\pgfqpoint{1.613206in}{2.217294in}}%
\pgfpathlineto{\pgfqpoint{1.571416in}{4.212191in}}%
\pgfusepath{stroke}%
\end{pgfscope}%
\begin{pgfscope}%
\pgfsetbuttcap%
\pgfsetroundjoin%
\pgfsetlinewidth{0.501875pt}%
\definecolor{currentstroke}{rgb}{0.850000,0.850000,0.850000}%
\pgfsetstrokecolor{currentstroke}%
\pgfsetdash{}{0pt}%
\pgfpathmoveto{\pgfqpoint{3.700922in}{0.961070in}}%
\pgfpathlineto{\pgfqpoint{2.137415in}{2.369661in}}%
\pgfpathlineto{\pgfqpoint{2.118803in}{4.350938in}}%
\pgfusepath{stroke}%
\end{pgfscope}%
\begin{pgfscope}%
\pgfsetbuttcap%
\pgfsetroundjoin%
\pgfsetlinewidth{0.501875pt}%
\definecolor{currentstroke}{rgb}{0.850000,0.850000,0.850000}%
\pgfsetstrokecolor{currentstroke}%
\pgfsetdash{}{0pt}%
\pgfpathmoveto{\pgfqpoint{4.244800in}{1.138701in}}%
\pgfpathlineto{\pgfqpoint{2.651532in}{2.519095in}}%
\pgfpathlineto{\pgfqpoint{2.655198in}{4.486898in}}%
\pgfusepath{stroke}%
\end{pgfscope}%
\begin{pgfscope}%
\pgfsetbuttcap%
\pgfsetroundjoin%
\pgfsetlinewidth{0.501875pt}%
\definecolor{currentstroke}{rgb}{0.850000,0.850000,0.850000}%
\pgfsetstrokecolor{currentstroke}%
\pgfsetdash{}{0pt}%
\pgfpathmoveto{\pgfqpoint{2.144000in}{0.452580in}}%
\pgfpathlineto{\pgfqpoint{0.670642in}{1.943328in}}%
\pgfpathlineto{\pgfqpoint{0.586003in}{3.962418in}}%
\pgfusepath{stroke}%
\end{pgfscope}%
\begin{pgfscope}%
\pgfsetbuttcap%
\pgfsetroundjoin%
\pgfsetlinewidth{0.501875pt}%
\definecolor{currentstroke}{rgb}{0.850000,0.850000,0.850000}%
\pgfsetstrokecolor{currentstroke}%
\pgfsetdash{}{0pt}%
\pgfpathmoveto{\pgfqpoint{2.289392in}{0.500065in}}%
\pgfpathlineto{\pgfqpoint{0.807306in}{1.983051in}}%
\pgfpathlineto{\pgfqpoint{0.728974in}{3.998657in}}%
\pgfusepath{stroke}%
\end{pgfscope}%
\begin{pgfscope}%
\pgfsetbuttcap%
\pgfsetroundjoin%
\pgfsetlinewidth{0.501875pt}%
\definecolor{currentstroke}{rgb}{0.850000,0.850000,0.850000}%
\pgfsetstrokecolor{currentstroke}%
\pgfsetdash{}{0pt}%
\pgfpathmoveto{\pgfqpoint{2.433995in}{0.547292in}}%
\pgfpathlineto{\pgfqpoint{0.943291in}{2.022576in}}%
\pgfpathlineto{\pgfqpoint{0.871203in}{4.034708in}}%
\pgfusepath{stroke}%
\end{pgfscope}%
\begin{pgfscope}%
\pgfsetbuttcap%
\pgfsetroundjoin%
\pgfsetlinewidth{0.501875pt}%
\definecolor{currentstroke}{rgb}{0.850000,0.850000,0.850000}%
\pgfsetstrokecolor{currentstroke}%
\pgfsetdash{}{0pt}%
\pgfpathmoveto{\pgfqpoint{2.720857in}{0.640981in}}%
\pgfpathlineto{\pgfqpoint{1.213244in}{2.101041in}}%
\pgfpathlineto{\pgfqpoint{1.153457in}{4.106251in}}%
\pgfusepath{stroke}%
\end{pgfscope}%
\begin{pgfscope}%
\pgfsetbuttcap%
\pgfsetroundjoin%
\pgfsetlinewidth{0.501875pt}%
\definecolor{currentstroke}{rgb}{0.850000,0.850000,0.850000}%
\pgfsetstrokecolor{currentstroke}%
\pgfsetdash{}{0pt}%
\pgfpathmoveto{\pgfqpoint{2.863129in}{0.687447in}}%
\pgfpathlineto{\pgfqpoint{1.347222in}{2.139983in}}%
\pgfpathlineto{\pgfqpoint{1.293494in}{4.141746in}}%
\pgfusepath{stroke}%
\end{pgfscope}%
\begin{pgfscope}%
\pgfsetbuttcap%
\pgfsetroundjoin%
\pgfsetlinewidth{0.501875pt}%
\definecolor{currentstroke}{rgb}{0.850000,0.850000,0.850000}%
\pgfsetstrokecolor{currentstroke}%
\pgfsetdash{}{0pt}%
\pgfpathmoveto{\pgfqpoint{3.004637in}{0.733664in}}%
\pgfpathlineto{\pgfqpoint{1.480541in}{2.178734in}}%
\pgfpathlineto{\pgfqpoint{1.432812in}{4.177059in}}%
\pgfusepath{stroke}%
\end{pgfscope}%
\begin{pgfscope}%
\pgfsetbuttcap%
\pgfsetroundjoin%
\pgfsetlinewidth{0.501875pt}%
\definecolor{currentstroke}{rgb}{0.850000,0.850000,0.850000}%
\pgfsetstrokecolor{currentstroke}%
\pgfsetdash{}{0pt}%
\pgfpathmoveto{\pgfqpoint{3.285383in}{0.825356in}}%
\pgfpathlineto{\pgfqpoint{1.745221in}{2.255666in}}%
\pgfpathlineto{\pgfqpoint{1.709312in}{4.247144in}}%
\pgfusepath{stroke}%
\end{pgfscope}%
\begin{pgfscope}%
\pgfsetbuttcap%
\pgfsetroundjoin%
\pgfsetlinewidth{0.501875pt}%
\definecolor{currentstroke}{rgb}{0.850000,0.850000,0.850000}%
\pgfsetstrokecolor{currentstroke}%
\pgfsetdash{}{0pt}%
\pgfpathmoveto{\pgfqpoint{3.424634in}{0.870835in}}%
\pgfpathlineto{\pgfqpoint{1.876591in}{2.293850in}}%
\pgfpathlineto{\pgfqpoint{1.846505in}{4.281918in}}%
\pgfusepath{stroke}%
\end{pgfscope}%
\begin{pgfscope}%
\pgfsetbuttcap%
\pgfsetroundjoin%
\pgfsetlinewidth{0.501875pt}%
\definecolor{currentstroke}{rgb}{0.850000,0.850000,0.850000}%
\pgfsetstrokecolor{currentstroke}%
\pgfsetdash{}{0pt}%
\pgfpathmoveto{\pgfqpoint{3.563145in}{0.916073in}}%
\pgfpathlineto{\pgfqpoint{2.007321in}{2.331848in}}%
\pgfpathlineto{\pgfqpoint{1.983000in}{4.316516in}}%
\pgfusepath{stroke}%
\end{pgfscope}%
\begin{pgfscope}%
\pgfsetbuttcap%
\pgfsetroundjoin%
\pgfsetlinewidth{0.501875pt}%
\definecolor{currentstroke}{rgb}{0.850000,0.850000,0.850000}%
\pgfsetstrokecolor{currentstroke}%
\pgfsetdash{}{0pt}%
\pgfpathmoveto{\pgfqpoint{3.837970in}{1.005830in}}%
\pgfpathlineto{\pgfqpoint{2.266879in}{2.407291in}}%
\pgfpathlineto{\pgfqpoint{2.253919in}{4.385186in}}%
\pgfusepath{stroke}%
\end{pgfscope}%
\begin{pgfscope}%
\pgfsetbuttcap%
\pgfsetroundjoin%
\pgfsetlinewidth{0.501875pt}%
\definecolor{currentstroke}{rgb}{0.850000,0.850000,0.850000}%
\pgfsetstrokecolor{currentstroke}%
\pgfsetdash{}{0pt}%
\pgfpathmoveto{\pgfqpoint{3.974295in}{1.050354in}}%
\pgfpathlineto{\pgfqpoint{2.395717in}{2.444739in}}%
\pgfpathlineto{\pgfqpoint{2.388354in}{4.419261in}}%
\pgfusepath{stroke}%
\end{pgfscope}%
\begin{pgfscope}%
\pgfsetbuttcap%
\pgfsetroundjoin%
\pgfsetlinewidth{0.501875pt}%
\definecolor{currentstroke}{rgb}{0.850000,0.850000,0.850000}%
\pgfsetstrokecolor{currentstroke}%
\pgfsetdash{}{0pt}%
\pgfpathmoveto{\pgfqpoint{4.109903in}{1.094644in}}%
\pgfpathlineto{\pgfqpoint{2.523933in}{2.482007in}}%
\pgfpathlineto{\pgfqpoint{2.522111in}{4.453165in}}%
\pgfusepath{stroke}%
\end{pgfscope}%
\begin{pgfscope}%
\pgfsetbuttcap%
\pgfsetroundjoin%
\pgfsetlinewidth{0.501875pt}%
\definecolor{currentstroke}{rgb}{0.850000,0.850000,0.850000}%
\pgfsetstrokecolor{currentstroke}%
\pgfsetdash{}{0pt}%
\pgfpathmoveto{\pgfqpoint{4.378991in}{1.182528in}}%
\pgfpathlineto{\pgfqpoint{2.778518in}{2.556004in}}%
\pgfpathlineto{\pgfqpoint{2.787617in}{4.520463in}}%
\pgfusepath{stroke}%
\end{pgfscope}%
\begin{pgfscope}%
\pgfsetbuttcap%
\pgfsetroundjoin%
\pgfsetlinewidth{0.501875pt}%
\definecolor{currentstroke}{rgb}{0.850000,0.850000,0.850000}%
\pgfsetstrokecolor{currentstroke}%
\pgfsetdash{}{0pt}%
\pgfpathmoveto{\pgfqpoint{4.512482in}{1.226126in}}%
\pgfpathlineto{\pgfqpoint{2.904895in}{2.592737in}}%
\pgfpathlineto{\pgfqpoint{2.919376in}{4.553860in}}%
\pgfusepath{stroke}%
\end{pgfscope}%
\begin{pgfscope}%
\pgfsetbuttcap%
\pgfsetroundjoin%
\pgfsetlinewidth{0.501875pt}%
\definecolor{currentstroke}{rgb}{0.850000,0.850000,0.850000}%
\pgfsetstrokecolor{currentstroke}%
\pgfsetdash{}{0pt}%
\pgfpathmoveto{\pgfqpoint{4.645278in}{1.269497in}}%
\pgfpathlineto{\pgfqpoint{3.030669in}{2.629295in}}%
\pgfpathlineto{\pgfqpoint{3.050477in}{4.587090in}}%
\pgfusepath{stroke}%
\end{pgfscope}%
\begin{pgfscope}%
\pgfsetbuttcap%
\pgfsetroundjoin%
\pgfsetlinewidth{0.501875pt}%
\definecolor{currentstroke}{rgb}{0.850000,0.850000,0.850000}%
\pgfsetstrokecolor{currentstroke}%
\pgfsetdash{}{0pt}%
\pgfpathmoveto{\pgfqpoint{4.575164in}{3.539656in}}%
\pgfpathlineto{\pgfqpoint{4.486806in}{1.480574in}}%
\pgfpathlineto{\pgfqpoint{1.755315in}{0.602532in}}%
\pgfusepath{stroke}%
\end{pgfscope}%
\begin{pgfscope}%
\pgfsetbuttcap%
\pgfsetroundjoin%
\pgfsetlinewidth{0.501875pt}%
\definecolor{currentstroke}{rgb}{0.850000,0.850000,0.850000}%
\pgfsetstrokecolor{currentstroke}%
\pgfsetdash{}{0pt}%
\pgfpathmoveto{\pgfqpoint{4.216811in}{3.801165in}}%
\pgfpathlineto{\pgfqpoint{4.145480in}{1.766680in}}%
\pgfpathlineto{\pgfqpoint{1.446184in}{0.919776in}}%
\pgfusepath{stroke}%
\end{pgfscope}%
\begin{pgfscope}%
\pgfsetbuttcap%
\pgfsetroundjoin%
\pgfsetlinewidth{0.501875pt}%
\definecolor{currentstroke}{rgb}{0.850000,0.850000,0.850000}%
\pgfsetstrokecolor{currentstroke}%
\pgfsetdash{}{0pt}%
\pgfpathmoveto{\pgfqpoint{3.871285in}{4.053315in}}%
\pgfpathlineto{\pgfqpoint{3.815852in}{2.042981in}}%
\pgfpathlineto{\pgfqpoint{1.148194in}{1.225588in}}%
\pgfusepath{stroke}%
\end{pgfscope}%
\begin{pgfscope}%
\pgfsetbuttcap%
\pgfsetroundjoin%
\pgfsetlinewidth{0.501875pt}%
\definecolor{currentstroke}{rgb}{0.850000,0.850000,0.850000}%
\pgfsetstrokecolor{currentstroke}%
\pgfsetdash{}{0pt}%
\pgfpathmoveto{\pgfqpoint{3.537907in}{4.296599in}}%
\pgfpathlineto{\pgfqpoint{3.497330in}{2.309972in}}%
\pgfpathlineto{\pgfqpoint{0.860752in}{1.520574in}}%
\pgfusepath{stroke}%
\end{pgfscope}%
\begin{pgfscope}%
\pgfsetbuttcap%
\pgfsetroundjoin%
\pgfsetlinewidth{0.501875pt}%
\definecolor{currentstroke}{rgb}{0.850000,0.850000,0.850000}%
\pgfsetstrokecolor{currentstroke}%
\pgfsetdash{}{0pt}%
\pgfpathmoveto{\pgfqpoint{3.216049in}{4.531476in}}%
\pgfpathlineto{\pgfqpoint{3.189362in}{2.568117in}}%
\pgfpathlineto{\pgfqpoint{0.583309in}{1.805299in}}%
\pgfusepath{stroke}%
\end{pgfscope}%
\begin{pgfscope}%
\pgfsetbuttcap%
\pgfsetroundjoin%
\pgfsetlinewidth{0.501875pt}%
\definecolor{currentstroke}{rgb}{0.850000,0.850000,0.850000}%
\pgfsetstrokecolor{currentstroke}%
\pgfsetdash{}{0pt}%
\pgfpathmoveto{\pgfqpoint{4.759375in}{3.405227in}}%
\pgfpathlineto{\pgfqpoint{4.662051in}{1.333679in}}%
\pgfpathlineto{\pgfqpoint{1.914256in}{0.439420in}}%
\pgfusepath{stroke}%
\end{pgfscope}%
\begin{pgfscope}%
\pgfsetbuttcap%
\pgfsetroundjoin%
\pgfsetlinewidth{0.501875pt}%
\definecolor{currentstroke}{rgb}{0.850000,0.850000,0.850000}%
\pgfsetstrokecolor{currentstroke}%
\pgfsetdash{}{0pt}%
\pgfpathmoveto{\pgfqpoint{4.666840in}{3.472755in}}%
\pgfpathlineto{\pgfqpoint{4.574038in}{1.407454in}}%
\pgfpathlineto{\pgfqpoint{1.834412in}{0.521359in}}%
\pgfusepath{stroke}%
\end{pgfscope}%
\begin{pgfscope}%
\pgfsetbuttcap%
\pgfsetroundjoin%
\pgfsetlinewidth{0.501875pt}%
\definecolor{currentstroke}{rgb}{0.850000,0.850000,0.850000}%
\pgfsetstrokecolor{currentstroke}%
\pgfsetdash{}{0pt}%
\pgfpathmoveto{\pgfqpoint{4.484334in}{3.605939in}}%
\pgfpathlineto{\pgfqpoint{4.400344in}{1.553048in}}%
\pgfpathlineto{\pgfqpoint{1.676954in}{0.682950in}}%
\pgfusepath{stroke}%
\end{pgfscope}%
\begin{pgfscope}%
\pgfsetbuttcap%
\pgfsetroundjoin%
\pgfsetlinewidth{0.501875pt}%
\definecolor{currentstroke}{rgb}{0.850000,0.850000,0.850000}%
\pgfsetstrokecolor{currentstroke}%
\pgfsetdash{}{0pt}%
\pgfpathmoveto{\pgfqpoint{4.394340in}{3.671613in}}%
\pgfpathlineto{\pgfqpoint{4.314642in}{1.624885in}}%
\pgfpathlineto{\pgfqpoint{1.599319in}{0.762623in}}%
\pgfusepath{stroke}%
\end{pgfscope}%
\begin{pgfscope}%
\pgfsetbuttcap%
\pgfsetroundjoin%
\pgfsetlinewidth{0.501875pt}%
\definecolor{currentstroke}{rgb}{0.850000,0.850000,0.850000}%
\pgfsetstrokecolor{currentstroke}%
\pgfsetdash{}{0pt}%
\pgfpathmoveto{\pgfqpoint{4.305170in}{3.736686in}}%
\pgfpathlineto{\pgfqpoint{4.229691in}{1.696093in}}%
\pgfpathlineto{\pgfqpoint{1.522399in}{0.841562in}}%
\pgfusepath{stroke}%
\end{pgfscope}%
\begin{pgfscope}%
\pgfsetbuttcap%
\pgfsetroundjoin%
\pgfsetlinewidth{0.501875pt}%
\definecolor{currentstroke}{rgb}{0.850000,0.850000,0.850000}%
\pgfsetstrokecolor{currentstroke}%
\pgfsetdash{}{0pt}%
\pgfpathmoveto{\pgfqpoint{4.129255in}{3.865060in}}%
\pgfpathlineto{\pgfqpoint{4.062000in}{1.836654in}}%
\pgfpathlineto{\pgfqpoint{1.370666in}{0.997276in}}%
\pgfusepath{stroke}%
\end{pgfscope}%
\begin{pgfscope}%
\pgfsetbuttcap%
\pgfsetroundjoin%
\pgfsetlinewidth{0.501875pt}%
\definecolor{currentstroke}{rgb}{0.850000,0.850000,0.850000}%
\pgfsetstrokecolor{currentstroke}%
\pgfsetdash{}{0pt}%
\pgfpathmoveto{\pgfqpoint{4.042488in}{3.928378in}}%
\pgfpathlineto{\pgfqpoint{3.979242in}{1.906024in}}%
\pgfpathlineto{\pgfqpoint{1.295835in}{1.074072in}}%
\pgfusepath{stroke}%
\end{pgfscope}%
\begin{pgfscope}%
\pgfsetbuttcap%
\pgfsetroundjoin%
\pgfsetlinewidth{0.501875pt}%
\definecolor{currentstroke}{rgb}{0.850000,0.850000,0.850000}%
\pgfsetstrokecolor{currentstroke}%
\pgfsetdash{}{0pt}%
\pgfpathmoveto{\pgfqpoint{3.956502in}{3.991128in}}%
\pgfpathlineto{\pgfqpoint{3.897195in}{1.974797in}}%
\pgfpathlineto{\pgfqpoint{1.221680in}{1.150173in}}%
\pgfusepath{stroke}%
\end{pgfscope}%
\begin{pgfscope}%
\pgfsetbuttcap%
\pgfsetroundjoin%
\pgfsetlinewidth{0.501875pt}%
\definecolor{currentstroke}{rgb}{0.850000,0.850000,0.850000}%
\pgfsetstrokecolor{currentstroke}%
\pgfsetdash{}{0pt}%
\pgfpathmoveto{\pgfqpoint{3.786827in}{4.114949in}}%
\pgfpathlineto{\pgfqpoint{3.735202in}{2.110583in}}%
\pgfpathlineto{\pgfqpoint{1.075367in}{1.300326in}}%
\pgfusepath{stroke}%
\end{pgfscope}%
\begin{pgfscope}%
\pgfsetbuttcap%
\pgfsetroundjoin%
\pgfsetlinewidth{0.501875pt}%
\definecolor{currentstroke}{rgb}{0.850000,0.850000,0.850000}%
\pgfsetstrokecolor{currentstroke}%
\pgfsetdash{}{0pt}%
\pgfpathmoveto{\pgfqpoint{3.703118in}{4.176036in}}%
\pgfpathlineto{\pgfqpoint{3.655238in}{2.177610in}}%
\pgfpathlineto{\pgfqpoint{1.003190in}{1.374398in}}%
\pgfusepath{stroke}%
\end{pgfscope}%
\begin{pgfscope}%
\pgfsetbuttcap%
\pgfsetroundjoin%
\pgfsetlinewidth{0.501875pt}%
\definecolor{currentstroke}{rgb}{0.850000,0.850000,0.850000}%
\pgfsetstrokecolor{currentstroke}%
\pgfsetdash{}{0pt}%
\pgfpathmoveto{\pgfqpoint{3.620148in}{4.236583in}}%
\pgfpathlineto{\pgfqpoint{3.575950in}{2.244071in}}%
\pgfpathlineto{\pgfqpoint{0.931655in}{1.447811in}}%
\pgfusepath{stroke}%
\end{pgfscope}%
\begin{pgfscope}%
\pgfsetbuttcap%
\pgfsetroundjoin%
\pgfsetlinewidth{0.501875pt}%
\definecolor{currentstroke}{rgb}{0.850000,0.850000,0.850000}%
\pgfsetstrokecolor{currentstroke}%
\pgfsetdash{}{0pt}%
\pgfpathmoveto{\pgfqpoint{3.456386in}{4.356089in}}%
\pgfpathlineto{\pgfqpoint{3.419369in}{2.375320in}}%
\pgfpathlineto{\pgfqpoint{0.790475in}{1.592696in}}%
\pgfusepath{stroke}%
\end{pgfscope}%
\begin{pgfscope}%
\pgfsetbuttcap%
\pgfsetroundjoin%
\pgfsetlinewidth{0.501875pt}%
\definecolor{currentstroke}{rgb}{0.850000,0.850000,0.850000}%
\pgfsetstrokecolor{currentstroke}%
\pgfsetdash{}{0pt}%
\pgfpathmoveto{\pgfqpoint{3.375576in}{4.415061in}}%
\pgfpathlineto{\pgfqpoint{3.342060in}{2.440123in}}%
\pgfpathlineto{\pgfqpoint{0.720814in}{1.664185in}}%
\pgfusepath{stroke}%
\end{pgfscope}%
\begin{pgfscope}%
\pgfsetbuttcap%
\pgfsetroundjoin%
\pgfsetlinewidth{0.501875pt}%
\definecolor{currentstroke}{rgb}{0.850000,0.850000,0.850000}%
\pgfsetstrokecolor{currentstroke}%
\pgfsetdash{}{0pt}%
\pgfpathmoveto{\pgfqpoint{3.295466in}{4.473521in}}%
\pgfpathlineto{\pgfqpoint{3.265393in}{2.504386in}}%
\pgfpathlineto{\pgfqpoint{0.651761in}{1.735050in}}%
\pgfusepath{stroke}%
\end{pgfscope}%
\begin{pgfscope}%
\pgfsetbuttcap%
\pgfsetroundjoin%
\pgfsetlinewidth{0.501875pt}%
\definecolor{currentstroke}{rgb}{0.850000,0.850000,0.850000}%
\pgfsetstrokecolor{currentstroke}%
\pgfsetdash{}{0pt}%
\pgfpathmoveto{\pgfqpoint{3.137315in}{4.588933in}}%
\pgfpathlineto{\pgfqpoint{3.113958in}{2.631322in}}%
\pgfpathlineto{\pgfqpoint{0.515450in}{1.874939in}}%
\pgfusepath{stroke}%
\end{pgfscope}%
\begin{pgfscope}%
\pgfsetbuttcap%
\pgfsetroundjoin%
\pgfsetlinewidth{0.501875pt}%
\definecolor{currentstroke}{rgb}{0.850000,0.850000,0.850000}%
\pgfsetstrokecolor{currentstroke}%
\pgfsetdash{}{0pt}%
\pgfpathmoveto{\pgfqpoint{0.495996in}{1.931849in}}%
\pgfpathlineto{\pgfqpoint{3.094741in}{2.685347in}}%
\pgfpathlineto{\pgfqpoint{4.714360in}{1.331126in}}%
\pgfusepath{stroke}%
\end{pgfscope}%
\begin{pgfscope}%
\pgfsetbuttcap%
\pgfsetroundjoin%
\pgfsetlinewidth{0.501875pt}%
\definecolor{currentstroke}{rgb}{0.850000,0.850000,0.850000}%
\pgfsetstrokecolor{currentstroke}%
\pgfsetdash{}{0pt}%
\pgfpathmoveto{\pgfqpoint{0.479135in}{2.300050in}}%
\pgfpathlineto{\pgfqpoint{3.098840in}{3.041874in}}%
\pgfpathlineto{\pgfqpoint{4.732520in}{1.708225in}}%
\pgfusepath{stroke}%
\end{pgfscope}%
\begin{pgfscope}%
\pgfsetbuttcap%
\pgfsetroundjoin%
\pgfsetlinewidth{0.501875pt}%
\definecolor{currentstroke}{rgb}{0.850000,0.850000,0.850000}%
\pgfsetstrokecolor{currentstroke}%
\pgfsetdash{}{0pt}%
\pgfpathmoveto{\pgfqpoint{0.461997in}{2.674298in}}%
\pgfpathlineto{\pgfqpoint{3.103004in}{3.403976in}}%
\pgfpathlineto{\pgfqpoint{4.750989in}{2.091754in}}%
\pgfusepath{stroke}%
\end{pgfscope}%
\begin{pgfscope}%
\pgfsetbuttcap%
\pgfsetroundjoin%
\pgfsetlinewidth{0.501875pt}%
\definecolor{currentstroke}{rgb}{0.850000,0.850000,0.850000}%
\pgfsetstrokecolor{currentstroke}%
\pgfsetdash{}{0pt}%
\pgfpathmoveto{\pgfqpoint{0.444575in}{3.054743in}}%
\pgfpathlineto{\pgfqpoint{3.107234in}{3.771783in}}%
\pgfpathlineto{\pgfqpoint{4.769776in}{2.481879in}}%
\pgfusepath{stroke}%
\end{pgfscope}%
\begin{pgfscope}%
\pgfsetbuttcap%
\pgfsetroundjoin%
\pgfsetlinewidth{0.501875pt}%
\definecolor{currentstroke}{rgb}{0.850000,0.850000,0.850000}%
\pgfsetstrokecolor{currentstroke}%
\pgfsetdash{}{0pt}%
\pgfpathmoveto{\pgfqpoint{0.426863in}{3.441540in}}%
\pgfpathlineto{\pgfqpoint{3.111530in}{4.145433in}}%
\pgfpathlineto{\pgfqpoint{4.788889in}{2.878772in}}%
\pgfusepath{stroke}%
\end{pgfscope}%
\begin{pgfscope}%
\pgfsetbuttcap%
\pgfsetroundjoin%
\pgfsetlinewidth{0.501875pt}%
\definecolor{currentstroke}{rgb}{0.850000,0.850000,0.850000}%
\pgfsetstrokecolor{currentstroke}%
\pgfsetdash{}{0pt}%
\pgfpathmoveto{\pgfqpoint{0.408852in}{3.834850in}}%
\pgfpathlineto{\pgfqpoint{3.115896in}{4.525066in}}%
\pgfpathlineto{\pgfqpoint{4.808337in}{3.282610in}}%
\pgfusepath{stroke}%
\end{pgfscope}%
\begin{pgfscope}%
\pgfsetrectcap%
\pgfsetroundjoin%
\pgfsetlinewidth{0.501875pt}%
\definecolor{currentstroke}{rgb}{0.000000,0.000000,0.000000}%
\pgfsetstrokecolor{currentstroke}%
\pgfsetdash{}{0pt}%
\pgfpathmoveto{\pgfqpoint{4.712449in}{1.291435in}}%
\pgfpathlineto{\pgfqpoint{1.959992in}{0.392483in}}%
\pgfusepath{stroke}%
\end{pgfscope}%
\begin{pgfscope}%
\pgfsetrectcap%
\pgfsetroundjoin%
\pgfsetlinewidth{0.501875pt}%
\definecolor{currentstroke}{rgb}{0.000000,0.000000,0.000000}%
\pgfsetstrokecolor{currentstroke}%
\pgfsetdash{}{0pt}%
\pgfpathmoveto{\pgfqpoint{1.985007in}{0.417937in}}%
\pgfpathlineto{\pgfqpoint{2.023478in}{0.378571in}}%
\pgfusepath{stroke}%
\end{pgfscope}%
\begin{pgfscope}%
\definecolor{textcolor}{rgb}{0.000000,0.000000,0.000000}%
\pgfsetstrokecolor{textcolor}%
\pgfsetfillcolor{textcolor}%
\pgftext[x=2.100026in,y=0.173457in,,top]{\color{textcolor}{\rmfamily\fontsize{10.000000}{12.000000}\selectfont\catcode`\^=\active\def^{\ifmmode\sp\else\^{}\fi}\catcode`\%=\active\def%{\%}$\mathdefault{0}$}}%
\end{pgfscope}%
\begin{pgfscope}%
\pgfsetrectcap%
\pgfsetroundjoin%
\pgfsetlinewidth{0.501875pt}%
\definecolor{currentstroke}{rgb}{0.000000,0.000000,0.000000}%
\pgfsetstrokecolor{currentstroke}%
\pgfsetdash{}{0pt}%
\pgfpathmoveto{\pgfqpoint{2.564718in}{0.607084in}}%
\pgfpathlineto{\pgfqpoint{2.604065in}{0.568566in}}%
\pgfusepath{stroke}%
\end{pgfscope}%
\begin{pgfscope}%
\definecolor{textcolor}{rgb}{0.000000,0.000000,0.000000}%
\pgfsetstrokecolor{textcolor}%
\pgfsetfillcolor{textcolor}%
\pgftext[x=2.680607in,y=0.365951in,,top]{\color{textcolor}{\rmfamily\fontsize{10.000000}{12.000000}\selectfont\catcode`\^=\active\def^{\ifmmode\sp\else\^{}\fi}\catcode`\%=\active\def%{\%}$\mathdefault{80}$}}%
\end{pgfscope}%
\begin{pgfscope}%
\pgfsetrectcap%
\pgfsetroundjoin%
\pgfsetlinewidth{0.501875pt}%
\definecolor{currentstroke}{rgb}{0.000000,0.000000,0.000000}%
\pgfsetstrokecolor{currentstroke}%
\pgfsetdash{}{0pt}%
\pgfpathmoveto{\pgfqpoint{3.132014in}{0.792180in}}%
\pgfpathlineto{\pgfqpoint{3.172189in}{0.754483in}}%
\pgfusepath{stroke}%
\end{pgfscope}%
\begin{pgfscope}%
\definecolor{textcolor}{rgb}{0.000000,0.000000,0.000000}%
\pgfsetstrokecolor{textcolor}%
\pgfsetfillcolor{textcolor}%
\pgftext[x=3.248707in,y=0.554307in,,top]{\color{textcolor}{\rmfamily\fontsize{10.000000}{12.000000}\selectfont\catcode`\^=\active\def^{\ifmmode\sp\else\^{}\fi}\catcode`\%=\active\def%{\%}$\mathdefault{160}$}}%
\end{pgfscope}%
\begin{pgfscope}%
\pgfsetrectcap%
\pgfsetroundjoin%
\pgfsetlinewidth{0.501875pt}%
\definecolor{currentstroke}{rgb}{0.000000,0.000000,0.000000}%
\pgfsetstrokecolor{currentstroke}%
\pgfsetdash{}{0pt}%
\pgfpathmoveto{\pgfqpoint{3.687288in}{0.973353in}}%
\pgfpathlineto{\pgfqpoint{3.728247in}{0.936452in}}%
\pgfusepath{stroke}%
\end{pgfscope}%
\begin{pgfscope}%
\definecolor{textcolor}{rgb}{0.000000,0.000000,0.000000}%
\pgfsetstrokecolor{textcolor}%
\pgfsetfillcolor{textcolor}%
\pgftext[x=3.804726in,y=0.738657in,,top]{\color{textcolor}{\rmfamily\fontsize{10.000000}{12.000000}\selectfont\catcode`\^=\active\def^{\ifmmode\sp\else\^{}\fi}\catcode`\%=\active\def%{\%}$\mathdefault{240}$}}%
\end{pgfscope}%
\begin{pgfscope}%
\pgfsetrectcap%
\pgfsetroundjoin%
\pgfsetlinewidth{0.501875pt}%
\definecolor{currentstroke}{rgb}{0.000000,0.000000,0.000000}%
\pgfsetstrokecolor{currentstroke}%
\pgfsetdash{}{0pt}%
\pgfpathmoveto{\pgfqpoint{4.230919in}{1.150727in}}%
\pgfpathlineto{\pgfqpoint{4.272621in}{1.114597in}}%
\pgfusepath{stroke}%
\end{pgfscope}%
\begin{pgfscope}%
\definecolor{textcolor}{rgb}{0.000000,0.000000,0.000000}%
\pgfsetstrokecolor{textcolor}%
\pgfsetfillcolor{textcolor}%
\pgftext[x=4.349045in,y=0.919129in,,top]{\color{textcolor}{\rmfamily\fontsize{10.000000}{12.000000}\selectfont\catcode`\^=\active\def^{\ifmmode\sp\else\^{}\fi}\catcode`\%=\active\def%{\%}$\mathdefault{320}$}}%
\end{pgfscope}%
\begin{pgfscope}%
\pgfsetrectcap%
\pgfsetroundjoin%
\pgfsetlinewidth{0.501875pt}%
\definecolor{currentstroke}{rgb}{0.000000,0.000000,0.000000}%
\pgfsetstrokecolor{currentstroke}%
\pgfsetdash{}{0pt}%
\pgfpathmoveto{\pgfqpoint{2.131121in}{0.465611in}}%
\pgfpathlineto{\pgfqpoint{2.169816in}{0.426459in}}%
\pgfusepath{stroke}%
\end{pgfscope}%
\begin{pgfscope}%
\pgfsetrectcap%
\pgfsetroundjoin%
\pgfsetlinewidth{0.501875pt}%
\definecolor{currentstroke}{rgb}{0.000000,0.000000,0.000000}%
\pgfsetstrokecolor{currentstroke}%
\pgfsetdash{}{0pt}%
\pgfpathmoveto{\pgfqpoint{2.276440in}{0.513025in}}%
\pgfpathlineto{\pgfqpoint{2.315355in}{0.474087in}}%
\pgfusepath{stroke}%
\end{pgfscope}%
\begin{pgfscope}%
\pgfsetrectcap%
\pgfsetroundjoin%
\pgfsetlinewidth{0.501875pt}%
\definecolor{currentstroke}{rgb}{0.000000,0.000000,0.000000}%
\pgfsetstrokecolor{currentstroke}%
\pgfsetdash{}{0pt}%
\pgfpathmoveto{\pgfqpoint{2.420970in}{0.560182in}}%
\pgfpathlineto{\pgfqpoint{2.460102in}{0.521455in}}%
\pgfusepath{stroke}%
\end{pgfscope}%
\begin{pgfscope}%
\pgfsetrectcap%
\pgfsetroundjoin%
\pgfsetlinewidth{0.501875pt}%
\definecolor{currentstroke}{rgb}{0.000000,0.000000,0.000000}%
\pgfsetstrokecolor{currentstroke}%
\pgfsetdash{}{0pt}%
\pgfpathmoveto{\pgfqpoint{2.707691in}{0.653732in}}%
\pgfpathlineto{\pgfqpoint{2.747248in}{0.615423in}}%
\pgfusepath{stroke}%
\end{pgfscope}%
\begin{pgfscope}%
\pgfsetrectcap%
\pgfsetroundjoin%
\pgfsetlinewidth{0.501875pt}%
\definecolor{currentstroke}{rgb}{0.000000,0.000000,0.000000}%
\pgfsetstrokecolor{currentstroke}%
\pgfsetdash{}{0pt}%
\pgfpathmoveto{\pgfqpoint{2.849893in}{0.700130in}}%
\pgfpathlineto{\pgfqpoint{2.889659in}{0.662026in}}%
\pgfusepath{stroke}%
\end{pgfscope}%
\begin{pgfscope}%
\pgfsetrectcap%
\pgfsetroundjoin%
\pgfsetlinewidth{0.501875pt}%
\definecolor{currentstroke}{rgb}{0.000000,0.000000,0.000000}%
\pgfsetstrokecolor{currentstroke}%
\pgfsetdash{}{0pt}%
\pgfpathmoveto{\pgfqpoint{2.991332in}{0.746278in}}%
\pgfpathlineto{\pgfqpoint{3.031304in}{0.708379in}}%
\pgfusepath{stroke}%
\end{pgfscope}%
\begin{pgfscope}%
\pgfsetrectcap%
\pgfsetroundjoin%
\pgfsetlinewidth{0.501875pt}%
\definecolor{currentstroke}{rgb}{0.000000,0.000000,0.000000}%
\pgfsetstrokecolor{currentstroke}%
\pgfsetdash{}{0pt}%
\pgfpathmoveto{\pgfqpoint{3.271945in}{0.837836in}}%
\pgfpathlineto{\pgfqpoint{3.312319in}{0.800341in}}%
\pgfusepath{stroke}%
\end{pgfscope}%
\begin{pgfscope}%
\pgfsetrectcap%
\pgfsetroundjoin%
\pgfsetlinewidth{0.501875pt}%
\definecolor{currentstroke}{rgb}{0.000000,0.000000,0.000000}%
\pgfsetstrokecolor{currentstroke}%
\pgfsetdash{}{0pt}%
\pgfpathmoveto{\pgfqpoint{3.411130in}{0.883249in}}%
\pgfpathlineto{\pgfqpoint{3.451702in}{0.845953in}}%
\pgfusepath{stroke}%
\end{pgfscope}%
\begin{pgfscope}%
\pgfsetrectcap%
\pgfsetroundjoin%
\pgfsetlinewidth{0.501875pt}%
\definecolor{currentstroke}{rgb}{0.000000,0.000000,0.000000}%
\pgfsetstrokecolor{currentstroke}%
\pgfsetdash{}{0pt}%
\pgfpathmoveto{\pgfqpoint{3.549576in}{0.928420in}}%
\pgfpathlineto{\pgfqpoint{3.590343in}{0.891323in}}%
\pgfusepath{stroke}%
\end{pgfscope}%
\begin{pgfscope}%
\pgfsetrectcap%
\pgfsetroundjoin%
\pgfsetlinewidth{0.501875pt}%
\definecolor{currentstroke}{rgb}{0.000000,0.000000,0.000000}%
\pgfsetstrokecolor{currentstroke}%
\pgfsetdash{}{0pt}%
\pgfpathmoveto{\pgfqpoint{3.824273in}{1.018048in}}%
\pgfpathlineto{\pgfqpoint{3.865422in}{0.981342in}}%
\pgfusepath{stroke}%
\end{pgfscope}%
\begin{pgfscope}%
\pgfsetrectcap%
\pgfsetroundjoin%
\pgfsetlinewidth{0.501875pt}%
\definecolor{currentstroke}{rgb}{0.000000,0.000000,0.000000}%
\pgfsetstrokecolor{currentstroke}%
\pgfsetdash{}{0pt}%
\pgfpathmoveto{\pgfqpoint{3.960536in}{1.062507in}}%
\pgfpathlineto{\pgfqpoint{4.001872in}{1.025995in}}%
\pgfusepath{stroke}%
\end{pgfscope}%
\begin{pgfscope}%
\pgfsetrectcap%
\pgfsetroundjoin%
\pgfsetlinewidth{0.501875pt}%
\definecolor{currentstroke}{rgb}{0.000000,0.000000,0.000000}%
\pgfsetstrokecolor{currentstroke}%
\pgfsetdash{}{0pt}%
\pgfpathmoveto{\pgfqpoint{4.096083in}{1.106733in}}%
\pgfpathlineto{\pgfqpoint{4.137603in}{1.070413in}}%
\pgfusepath{stroke}%
\end{pgfscope}%
\begin{pgfscope}%
\pgfsetrectcap%
\pgfsetroundjoin%
\pgfsetlinewidth{0.501875pt}%
\definecolor{currentstroke}{rgb}{0.000000,0.000000,0.000000}%
\pgfsetstrokecolor{currentstroke}%
\pgfsetdash{}{0pt}%
\pgfpathmoveto{\pgfqpoint{4.365050in}{1.194491in}}%
\pgfpathlineto{\pgfqpoint{4.406932in}{1.158550in}}%
\pgfusepath{stroke}%
\end{pgfscope}%
\begin{pgfscope}%
\pgfsetrectcap%
\pgfsetroundjoin%
\pgfsetlinewidth{0.501875pt}%
\definecolor{currentstroke}{rgb}{0.000000,0.000000,0.000000}%
\pgfsetstrokecolor{currentstroke}%
\pgfsetdash{}{0pt}%
\pgfpathmoveto{\pgfqpoint{4.498482in}{1.238027in}}%
\pgfpathlineto{\pgfqpoint{4.540541in}{1.202273in}}%
\pgfusepath{stroke}%
\end{pgfscope}%
\begin{pgfscope}%
\pgfsetrectcap%
\pgfsetroundjoin%
\pgfsetlinewidth{0.501875pt}%
\definecolor{currentstroke}{rgb}{0.000000,0.000000,0.000000}%
\pgfsetstrokecolor{currentstroke}%
\pgfsetdash{}{0pt}%
\pgfpathmoveto{\pgfqpoint{4.631220in}{1.281336in}}%
\pgfpathlineto{\pgfqpoint{4.673453in}{1.245768in}}%
\pgfusepath{stroke}%
\end{pgfscope}%
\begin{pgfscope}%
\definecolor{textcolor}{rgb}{0.000000,0.000000,0.000000}%
\pgfsetstrokecolor{textcolor}%
\pgfsetfillcolor{textcolor}%
\pgftext[x=3.616498in,y=0.325541in,,,rotate=18.087038]{\color{textcolor}{\rmfamily\fontsize{9.000000}{10.800000}\selectfont\catcode`\^=\active\def^{\ifmmode\sp\else\^{}\fi}\catcode`\%=\active\def%{\%}Time to Expiry (Days)}}%
\end{pgfscope}%
\begin{pgfscope}%
\pgfsetrectcap%
\pgfsetroundjoin%
\pgfsetlinewidth{0.501875pt}%
\definecolor{currentstroke}{rgb}{0.000000,0.000000,0.000000}%
\pgfsetstrokecolor{currentstroke}%
\pgfsetdash{}{0pt}%
\pgfpathmoveto{\pgfqpoint{0.497771in}{1.893081in}}%
\pgfpathlineto{\pgfqpoint{1.959992in}{0.392483in}}%
\pgfusepath{stroke}%
\end{pgfscope}%
\begin{pgfscope}%
\pgfsetrectcap%
\pgfsetroundjoin%
\pgfsetlinewidth{0.501875pt}%
\definecolor{currentstroke}{rgb}{0.000000,0.000000,0.000000}%
\pgfsetstrokecolor{currentstroke}%
\pgfsetdash{}{0pt}%
\pgfpathmoveto{\pgfqpoint{1.778327in}{0.609929in}}%
\pgfpathlineto{\pgfqpoint{1.709234in}{0.587719in}}%
\pgfusepath{stroke}%
\end{pgfscope}%
\begin{pgfscope}%
\definecolor{textcolor}{rgb}{0.000000,0.000000,0.000000}%
\pgfsetstrokecolor{textcolor}%
\pgfsetfillcolor{textcolor}%
\pgftext[x=1.577490in,y=0.417311in,,top]{\color{textcolor}{\rmfamily\fontsize{10.000000}{12.000000}\selectfont\catcode`\^=\active\def^{\ifmmode\sp\else\^{}\fi}\catcode`\%=\active\def%{\%}$\mathdefault{90}$}}%
\end{pgfscope}%
\begin{pgfscope}%
\pgfsetrectcap%
\pgfsetroundjoin%
\pgfsetlinewidth{0.501875pt}%
\definecolor{currentstroke}{rgb}{0.000000,0.000000,0.000000}%
\pgfsetstrokecolor{currentstroke}%
\pgfsetdash{}{0pt}%
\pgfpathmoveto{\pgfqpoint{1.468903in}{0.926904in}}%
\pgfpathlineto{\pgfqpoint{1.400689in}{0.905502in}}%
\pgfusepath{stroke}%
\end{pgfscope}%
\begin{pgfscope}%
\definecolor{textcolor}{rgb}{0.000000,0.000000,0.000000}%
\pgfsetstrokecolor{textcolor}%
\pgfsetfillcolor{textcolor}%
\pgftext[x=1.271488in,y=0.738095in,,top]{\color{textcolor}{\rmfamily\fontsize{10.000000}{12.000000}\selectfont\catcode`\^=\active\def^{\ifmmode\sp\else\^{}\fi}\catcode`\%=\active\def%{\%}$\mathdefault{105}$}}%
\end{pgfscope}%
\begin{pgfscope}%
\pgfsetrectcap%
\pgfsetroundjoin%
\pgfsetlinewidth{0.501875pt}%
\definecolor{currentstroke}{rgb}{0.000000,0.000000,0.000000}%
\pgfsetstrokecolor{currentstroke}%
\pgfsetdash{}{0pt}%
\pgfpathmoveto{\pgfqpoint{1.170627in}{1.232461in}}%
\pgfpathlineto{\pgfqpoint{1.103273in}{1.211824in}}%
\pgfusepath{stroke}%
\end{pgfscope}%
\begin{pgfscope}%
\definecolor{textcolor}{rgb}{0.000000,0.000000,0.000000}%
\pgfsetstrokecolor{textcolor}%
\pgfsetfillcolor{textcolor}%
\pgftext[x=0.976517in,y=1.047314in,,top]{\color{textcolor}{\rmfamily\fontsize{10.000000}{12.000000}\selectfont\catcode`\^=\active\def^{\ifmmode\sp\else\^{}\fi}\catcode`\%=\active\def%{\%}$\mathdefault{120}$}}%
\end{pgfscope}%
\begin{pgfscope}%
\pgfsetrectcap%
\pgfsetroundjoin%
\pgfsetlinewidth{0.501875pt}%
\definecolor{currentstroke}{rgb}{0.000000,0.000000,0.000000}%
\pgfsetstrokecolor{currentstroke}%
\pgfsetdash{}{0pt}%
\pgfpathmoveto{\pgfqpoint{0.882905in}{1.527206in}}%
\pgfpathlineto{\pgfqpoint{0.816394in}{1.507293in}}%
\pgfusepath{stroke}%
\end{pgfscope}%
\begin{pgfscope}%
\definecolor{textcolor}{rgb}{0.000000,0.000000,0.000000}%
\pgfsetstrokecolor{textcolor}%
\pgfsetfillcolor{textcolor}%
\pgftext[x=0.691994in,y=1.345582in,,top]{\color{textcolor}{\rmfamily\fontsize{10.000000}{12.000000}\selectfont\catcode`\^=\active\def^{\ifmmode\sp\else\^{}\fi}\catcode`\%=\active\def%{\%}$\mathdefault{135}$}}%
\end{pgfscope}%
\begin{pgfscope}%
\pgfsetrectcap%
\pgfsetroundjoin%
\pgfsetlinewidth{0.501875pt}%
\definecolor{currentstroke}{rgb}{0.000000,0.000000,0.000000}%
\pgfsetstrokecolor{currentstroke}%
\pgfsetdash{}{0pt}%
\pgfpathmoveto{\pgfqpoint{0.605187in}{1.811702in}}%
\pgfpathlineto{\pgfqpoint{0.539502in}{1.792476in}}%
\pgfusepath{stroke}%
\end{pgfscope}%
\begin{pgfscope}%
\definecolor{textcolor}{rgb}{0.000000,0.000000,0.000000}%
\pgfsetstrokecolor{textcolor}%
\pgfsetfillcolor{textcolor}%
\pgftext[x=0.417371in,y=1.633471in,,top]{\color{textcolor}{\rmfamily\fontsize{10.000000}{12.000000}\selectfont\catcode`\^=\active\def^{\ifmmode\sp\else\^{}\fi}\catcode`\%=\active\def%{\%}$\mathdefault{150}$}}%
\end{pgfscope}%
\begin{pgfscope}%
\pgfsetrectcap%
\pgfsetroundjoin%
\pgfsetlinewidth{0.501875pt}%
\definecolor{currentstroke}{rgb}{0.000000,0.000000,0.000000}%
\pgfsetstrokecolor{currentstroke}%
\pgfsetdash{}{0pt}%
\pgfpathmoveto{\pgfqpoint{1.937415in}{0.446957in}}%
\pgfpathlineto{\pgfqpoint{1.867876in}{0.424326in}}%
\pgfusepath{stroke}%
\end{pgfscope}%
\begin{pgfscope}%
\pgfsetrectcap%
\pgfsetroundjoin%
\pgfsetlinewidth{0.501875pt}%
\definecolor{currentstroke}{rgb}{0.000000,0.000000,0.000000}%
\pgfsetstrokecolor{currentstroke}%
\pgfsetdash{}{0pt}%
\pgfpathmoveto{\pgfqpoint{1.857498in}{0.528825in}}%
\pgfpathlineto{\pgfqpoint{1.788182in}{0.506406in}}%
\pgfusepath{stroke}%
\end{pgfscope}%
\begin{pgfscope}%
\pgfsetrectcap%
\pgfsetroundjoin%
\pgfsetlinewidth{0.501875pt}%
\definecolor{currentstroke}{rgb}{0.000000,0.000000,0.000000}%
\pgfsetstrokecolor{currentstroke}%
\pgfsetdash{}{0pt}%
\pgfpathmoveto{\pgfqpoint{1.699892in}{0.690278in}}%
\pgfpathlineto{\pgfqpoint{1.631020in}{0.668274in}}%
\pgfusepath{stroke}%
\end{pgfscope}%
\begin{pgfscope}%
\pgfsetrectcap%
\pgfsetroundjoin%
\pgfsetlinewidth{0.501875pt}%
\definecolor{currentstroke}{rgb}{0.000000,0.000000,0.000000}%
\pgfsetstrokecolor{currentstroke}%
\pgfsetdash{}{0pt}%
\pgfpathmoveto{\pgfqpoint{1.622183in}{0.769883in}}%
\pgfpathlineto{\pgfqpoint{1.553532in}{0.748083in}}%
\pgfusepath{stroke}%
\end{pgfscope}%
\begin{pgfscope}%
\pgfsetrectcap%
\pgfsetroundjoin%
\pgfsetlinewidth{0.501875pt}%
\definecolor{currentstroke}{rgb}{0.000000,0.000000,0.000000}%
\pgfsetstrokecolor{currentstroke}%
\pgfsetdash{}{0pt}%
\pgfpathmoveto{\pgfqpoint{1.545190in}{0.848755in}}%
\pgfpathlineto{\pgfqpoint{1.476758in}{0.827156in}}%
\pgfusepath{stroke}%
\end{pgfscope}%
\begin{pgfscope}%
\pgfsetrectcap%
\pgfsetroundjoin%
\pgfsetlinewidth{0.501875pt}%
\definecolor{currentstroke}{rgb}{0.000000,0.000000,0.000000}%
\pgfsetstrokecolor{currentstroke}%
\pgfsetdash{}{0pt}%
\pgfpathmoveto{\pgfqpoint{1.393313in}{1.004340in}}%
\pgfpathlineto{\pgfqpoint{1.325316in}{0.983132in}}%
\pgfusepath{stroke}%
\end{pgfscope}%
\begin{pgfscope}%
\pgfsetrectcap%
\pgfsetroundjoin%
\pgfsetlinewidth{0.501875pt}%
\definecolor{currentstroke}{rgb}{0.000000,0.000000,0.000000}%
\pgfsetstrokecolor{currentstroke}%
\pgfsetdash{}{0pt}%
\pgfpathmoveto{\pgfqpoint{1.318410in}{1.081071in}}%
\pgfpathlineto{\pgfqpoint{1.250628in}{1.060056in}}%
\pgfusepath{stroke}%
\end{pgfscope}%
\begin{pgfscope}%
\pgfsetrectcap%
\pgfsetroundjoin%
\pgfsetlinewidth{0.501875pt}%
\definecolor{currentstroke}{rgb}{0.000000,0.000000,0.000000}%
\pgfsetstrokecolor{currentstroke}%
\pgfsetdash{}{0pt}%
\pgfpathmoveto{\pgfqpoint{1.244184in}{1.157109in}}%
\pgfpathlineto{\pgfqpoint{1.176617in}{1.136284in}}%
\pgfusepath{stroke}%
\end{pgfscope}%
\begin{pgfscope}%
\pgfsetrectcap%
\pgfsetroundjoin%
\pgfsetlinewidth{0.501875pt}%
\definecolor{currentstroke}{rgb}{0.000000,0.000000,0.000000}%
\pgfsetstrokecolor{currentstroke}%
\pgfsetdash{}{0pt}%
\pgfpathmoveto{\pgfqpoint{1.097729in}{1.307138in}}%
\pgfpathlineto{\pgfqpoint{1.030588in}{1.286685in}}%
\pgfusepath{stroke}%
\end{pgfscope}%
\begin{pgfscope}%
\pgfsetrectcap%
\pgfsetroundjoin%
\pgfsetlinewidth{0.501875pt}%
\definecolor{currentstroke}{rgb}{0.000000,0.000000,0.000000}%
\pgfsetstrokecolor{currentstroke}%
\pgfsetdash{}{0pt}%
\pgfpathmoveto{\pgfqpoint{1.025482in}{1.381149in}}%
\pgfpathlineto{\pgfqpoint{0.958552in}{1.360878in}}%
\pgfusepath{stroke}%
\end{pgfscope}%
\begin{pgfscope}%
\pgfsetrectcap%
\pgfsetroundjoin%
\pgfsetlinewidth{0.501875pt}%
\definecolor{currentstroke}{rgb}{0.000000,0.000000,0.000000}%
\pgfsetstrokecolor{currentstroke}%
\pgfsetdash{}{0pt}%
\pgfpathmoveto{\pgfqpoint{0.953876in}{1.454502in}}%
\pgfpathlineto{\pgfqpoint{0.887157in}{1.434411in}}%
\pgfusepath{stroke}%
\end{pgfscope}%
\begin{pgfscope}%
\pgfsetrectcap%
\pgfsetroundjoin%
\pgfsetlinewidth{0.501875pt}%
\definecolor{currentstroke}{rgb}{0.000000,0.000000,0.000000}%
\pgfsetstrokecolor{currentstroke}%
\pgfsetdash{}{0pt}%
\pgfpathmoveto{\pgfqpoint{0.812558in}{1.599270in}}%
\pgfpathlineto{\pgfqpoint{0.746255in}{1.579531in}}%
\pgfusepath{stroke}%
\end{pgfscope}%
\begin{pgfscope}%
\pgfsetrectcap%
\pgfsetroundjoin%
\pgfsetlinewidth{0.501875pt}%
\definecolor{currentstroke}{rgb}{0.000000,0.000000,0.000000}%
\pgfsetstrokecolor{currentstroke}%
\pgfsetdash{}{0pt}%
\pgfpathmoveto{\pgfqpoint{0.742828in}{1.670702in}}%
\pgfpathlineto{\pgfqpoint{0.676732in}{1.651136in}}%
\pgfusepath{stroke}%
\end{pgfscope}%
\begin{pgfscope}%
\pgfsetrectcap%
\pgfsetroundjoin%
\pgfsetlinewidth{0.501875pt}%
\definecolor{currentstroke}{rgb}{0.000000,0.000000,0.000000}%
\pgfsetstrokecolor{currentstroke}%
\pgfsetdash{}{0pt}%
\pgfpathmoveto{\pgfqpoint{0.673707in}{1.741510in}}%
\pgfpathlineto{\pgfqpoint{0.607817in}{1.722115in}}%
\pgfusepath{stroke}%
\end{pgfscope}%
\begin{pgfscope}%
\pgfsetrectcap%
\pgfsetroundjoin%
\pgfsetlinewidth{0.501875pt}%
\definecolor{currentstroke}{rgb}{0.000000,0.000000,0.000000}%
\pgfsetstrokecolor{currentstroke}%
\pgfsetdash{}{0pt}%
\pgfpathmoveto{\pgfqpoint{0.537259in}{1.881288in}}%
\pgfpathlineto{\pgfqpoint{0.471778in}{1.862227in}}%
\pgfusepath{stroke}%
\end{pgfscope}%
\begin{pgfscope}%
\definecolor{textcolor}{rgb}{0.000000,0.000000,0.000000}%
\pgfsetstrokecolor{textcolor}%
\pgfsetfillcolor{textcolor}%
\pgftext[x=0.791685in,y=0.756307in,,,rotate=314.257888]{\color{textcolor}{\rmfamily\fontsize{9.000000}{10.800000}\selectfont\catcode`\^=\active\def^{\ifmmode\sp\else\^{}\fi}\catcode`\%=\active\def%{\%}Underlying ($S_t$)}}%
\end{pgfscope}%
\begin{pgfscope}%
\pgfsetrectcap%
\pgfsetroundjoin%
\pgfsetlinewidth{0.501875pt}%
\definecolor{currentstroke}{rgb}{0.000000,0.000000,0.000000}%
\pgfsetstrokecolor{currentstroke}%
\pgfsetdash{}{0pt}%
\pgfpathmoveto{\pgfqpoint{0.497771in}{1.893081in}}%
\pgfpathlineto{\pgfqpoint{0.405110in}{3.916567in}}%
\pgfusepath{stroke}%
\end{pgfscope}%
\begin{pgfscope}%
\pgfsetrectcap%
\pgfsetroundjoin%
\pgfsetlinewidth{0.501875pt}%
\definecolor{currentstroke}{rgb}{0.000000,0.000000,0.000000}%
\pgfsetstrokecolor{currentstroke}%
\pgfsetdash{}{0pt}%
\pgfpathmoveto{\pgfqpoint{0.517807in}{1.938173in}}%
\pgfpathlineto{\pgfqpoint{0.452321in}{1.919186in}}%
\pgfusepath{stroke}%
\end{pgfscope}%
\begin{pgfscope}%
\definecolor{textcolor}{rgb}{0.000000,0.000000,0.000000}%
\pgfsetstrokecolor{textcolor}%
\pgfsetfillcolor{textcolor}%
\pgftext[x=0.242143in,y=1.968000in,,top]{\color{textcolor}{\rmfamily\fontsize{10.000000}{12.000000}\selectfont\catcode`\^=\active\def^{\ifmmode\sp\else\^{}\fi}\catcode`\%=\active\def%{\%}$\mathdefault{0}$}}%
\end{pgfscope}%
\begin{pgfscope}%
\pgfsetrectcap%
\pgfsetroundjoin%
\pgfsetlinewidth{0.501875pt}%
\definecolor{currentstroke}{rgb}{0.000000,0.000000,0.000000}%
\pgfsetstrokecolor{currentstroke}%
\pgfsetdash{}{0pt}%
\pgfpathmoveto{\pgfqpoint{0.501131in}{2.306279in}}%
\pgfpathlineto{\pgfqpoint{0.435091in}{2.287578in}}%
\pgfusepath{stroke}%
\end{pgfscope}%
\begin{pgfscope}%
\definecolor{textcolor}{rgb}{0.000000,0.000000,0.000000}%
\pgfsetstrokecolor{textcolor}%
\pgfsetfillcolor{textcolor}%
\pgftext[x=0.223257in,y=2.335653in,,top]{\color{textcolor}{\rmfamily\fontsize{10.000000}{12.000000}\selectfont\catcode`\^=\active\def^{\ifmmode\sp\else\^{}\fi}\catcode`\%=\active\def%{\%}$\mathdefault{10}$}}%
\end{pgfscope}%
\begin{pgfscope}%
\pgfsetrectcap%
\pgfsetroundjoin%
\pgfsetlinewidth{0.501875pt}%
\definecolor{currentstroke}{rgb}{0.000000,0.000000,0.000000}%
\pgfsetstrokecolor{currentstroke}%
\pgfsetdash{}{0pt}%
\pgfpathmoveto{\pgfqpoint{0.484180in}{2.680427in}}%
\pgfpathlineto{\pgfqpoint{0.417577in}{2.662025in}}%
\pgfusepath{stroke}%
\end{pgfscope}%
\begin{pgfscope}%
\definecolor{textcolor}{rgb}{0.000000,0.000000,0.000000}%
\pgfsetstrokecolor{textcolor}%
\pgfsetfillcolor{textcolor}%
\pgftext[x=0.204063in,y=2.709331in,,top]{\color{textcolor}{\rmfamily\fontsize{10.000000}{12.000000}\selectfont\catcode`\^=\active\def^{\ifmmode\sp\else\^{}\fi}\catcode`\%=\active\def%{\%}$\mathdefault{20}$}}%
\end{pgfscope}%
\begin{pgfscope}%
\pgfsetrectcap%
\pgfsetroundjoin%
\pgfsetlinewidth{0.501875pt}%
\definecolor{currentstroke}{rgb}{0.000000,0.000000,0.000000}%
\pgfsetstrokecolor{currentstroke}%
\pgfsetdash{}{0pt}%
\pgfpathmoveto{\pgfqpoint{0.466949in}{3.060768in}}%
\pgfpathlineto{\pgfqpoint{0.399773in}{3.042678in}}%
\pgfusepath{stroke}%
\end{pgfscope}%
\begin{pgfscope}%
\definecolor{textcolor}{rgb}{0.000000,0.000000,0.000000}%
\pgfsetstrokecolor{textcolor}%
\pgfsetfillcolor{textcolor}%
\pgftext[x=0.184551in,y=3.089182in,,top]{\color{textcolor}{\rmfamily\fontsize{10.000000}{12.000000}\selectfont\catcode`\^=\active\def^{\ifmmode\sp\else\^{}\fi}\catcode`\%=\active\def%{\%}$\mathdefault{30}$}}%
\end{pgfscope}%
\begin{pgfscope}%
\pgfsetrectcap%
\pgfsetroundjoin%
\pgfsetlinewidth{0.501875pt}%
\definecolor{currentstroke}{rgb}{0.000000,0.000000,0.000000}%
\pgfsetstrokecolor{currentstroke}%
\pgfsetdash{}{0pt}%
\pgfpathmoveto{\pgfqpoint{0.449431in}{3.447457in}}%
\pgfpathlineto{\pgfqpoint{0.381672in}{3.429691in}}%
\pgfusepath{stroke}%
\end{pgfscope}%
\begin{pgfscope}%
\definecolor{textcolor}{rgb}{0.000000,0.000000,0.000000}%
\pgfsetstrokecolor{textcolor}%
\pgfsetfillcolor{textcolor}%
\pgftext[x=0.164714in,y=3.475361in,,top]{\color{textcolor}{\rmfamily\fontsize{10.000000}{12.000000}\selectfont\catcode`\^=\active\def^{\ifmmode\sp\else\^{}\fi}\catcode`\%=\active\def%{\%}$\mathdefault{40}$}}%
\end{pgfscope}%
\begin{pgfscope}%
\pgfsetrectcap%
\pgfsetroundjoin%
\pgfsetlinewidth{0.501875pt}%
\definecolor{currentstroke}{rgb}{0.000000,0.000000,0.000000}%
\pgfsetstrokecolor{currentstroke}%
\pgfsetdash{}{0pt}%
\pgfpathmoveto{\pgfqpoint{0.431617in}{3.840654in}}%
\pgfpathlineto{\pgfqpoint{0.363265in}{3.823226in}}%
\pgfusepath{stroke}%
\end{pgfscope}%
\begin{pgfscope}%
\definecolor{textcolor}{rgb}{0.000000,0.000000,0.000000}%
\pgfsetstrokecolor{textcolor}%
\pgfsetfillcolor{textcolor}%
\pgftext[x=0.144544in,y=3.868027in,,top]{\color{textcolor}{\rmfamily\fontsize{10.000000}{12.000000}\selectfont\catcode`\^=\active\def^{\ifmmode\sp\else\^{}\fi}\catcode`\%=\active\def%{\%}$\mathdefault{50}$}}%
\end{pgfscope}%
\begin{pgfscope}%
\definecolor{textcolor}{rgb}{0.000000,0.000000,0.000000}%
\pgfsetstrokecolor{textcolor}%
\pgfsetfillcolor{textcolor}%
\pgftext[x=-0.198584in,y=2.969882in,,,rotate=272.621908]{\color{textcolor}{\rmfamily\fontsize{9.000000}{10.800000}\selectfont\catcode`\^=\active\def^{\ifmmode\sp\else\^{}\fi}\catcode`\%=\active\def%{\%}Option Price ($P_{DO}$)}}%
\end{pgfscope}%
\begin{pgfscope}%
\pgfpathrectangle{\pgfqpoint{0.050000in}{0.083252in}}{\pgfqpoint{4.900000in}{4.900000in}}%
\pgfusepath{clip}%
\pgfsetbuttcap%
\pgfsetroundjoin%
\definecolor{currentfill}{rgb}{0.270588,0.541176,0.988235}%
\pgfsetfillcolor{currentfill}%
\pgfsetfillopacity{0.950000}%
\pgfsetlinewidth{0.200750pt}%
\definecolor{currentstroke}{rgb}{0.200000,0.200000,0.200000}%
\pgfsetstrokecolor{currentstroke}%
\pgfsetdash{}{0pt}%
\pgfpathmoveto{\pgfqpoint{3.031308in}{4.363944in}}%
\pgfpathlineto{\pgfqpoint{3.004072in}{4.435679in}}%
\pgfpathlineto{\pgfqpoint{3.050456in}{4.448294in}}%
\pgfpathlineto{\pgfqpoint{3.077690in}{4.376627in}}%
\pgfpathlineto{\pgfqpoint{3.031308in}{4.363944in}}%
\pgfpathclose%
\pgfusepath{stroke,fill}%
\end{pgfscope}%
\begin{pgfscope}%
\pgfpathrectangle{\pgfqpoint{0.050000in}{0.083252in}}{\pgfqpoint{4.900000in}{4.900000in}}%
\pgfusepath{clip}%
\pgfsetbuttcap%
\pgfsetroundjoin%
\definecolor{currentfill}{rgb}{0.268128,0.515033,0.953479}%
\pgfsetfillcolor{currentfill}%
\pgfsetfillopacity{0.950000}%
\pgfsetlinewidth{0.200750pt}%
\definecolor{currentstroke}{rgb}{0.200000,0.200000,0.200000}%
\pgfsetstrokecolor{currentstroke}%
\pgfsetdash{}{0pt}%
\pgfpathmoveto{\pgfqpoint{3.085856in}{4.220277in}}%
\pgfpathlineto{\pgfqpoint{3.058569in}{4.292143in}}%
\pgfpathlineto{\pgfqpoint{3.031308in}{4.363944in}}%
\pgfpathlineto{\pgfqpoint{3.077690in}{4.376627in}}%
\pgfpathlineto{\pgfqpoint{3.104949in}{4.304895in}}%
\pgfpathlineto{\pgfqpoint{3.132233in}{4.233097in}}%
\pgfpathlineto{\pgfqpoint{3.085856in}{4.220277in}}%
\pgfpathclose%
\pgfusepath{stroke,fill}%
\end{pgfscope}%
\begin{pgfscope}%
\pgfpathrectangle{\pgfqpoint{0.050000in}{0.083252in}}{\pgfqpoint{4.900000in}{4.900000in}}%
\pgfusepath{clip}%
\pgfsetbuttcap%
\pgfsetroundjoin%
\definecolor{currentfill}{rgb}{0.270588,0.541176,0.988235}%
\pgfsetfillcolor{currentfill}%
\pgfsetfillopacity{0.950000}%
\pgfsetlinewidth{0.200750pt}%
\definecolor{currentstroke}{rgb}{0.200000,0.200000,0.200000}%
\pgfsetstrokecolor{currentstroke}%
\pgfsetdash{}{0pt}%
\pgfpathmoveto{\pgfqpoint{2.938302in}{4.338512in}}%
\pgfpathlineto{\pgfqpoint{2.911061in}{4.410383in}}%
\pgfpathlineto{\pgfqpoint{2.957607in}{4.423042in}}%
\pgfpathlineto{\pgfqpoint{3.004072in}{4.435679in}}%
\pgfpathlineto{\pgfqpoint{3.031308in}{4.363944in}}%
\pgfpathlineto{\pgfqpoint{2.984845in}{4.351239in}}%
\pgfpathlineto{\pgfqpoint{2.938302in}{4.338512in}}%
\pgfpathclose%
\pgfusepath{stroke,fill}%
\end{pgfscope}%
\begin{pgfscope}%
\pgfpathrectangle{\pgfqpoint{0.050000in}{0.083252in}}{\pgfqpoint{4.900000in}{4.900000in}}%
\pgfusepath{clip}%
\pgfsetbuttcap%
\pgfsetroundjoin%
\definecolor{currentfill}{rgb}{0.264929,0.481046,0.908297}%
\pgfsetfillcolor{currentfill}%
\pgfsetfillopacity{0.950000}%
\pgfsetlinewidth{0.200750pt}%
\definecolor{currentstroke}{rgb}{0.200000,0.200000,0.200000}%
\pgfsetstrokecolor{currentstroke}%
\pgfsetdash{}{0pt}%
\pgfpathmoveto{\pgfqpoint{3.140504in}{4.076348in}}%
\pgfpathlineto{\pgfqpoint{3.113167in}{4.148345in}}%
\pgfpathlineto{\pgfqpoint{3.085856in}{4.220277in}}%
\pgfpathlineto{\pgfqpoint{3.132233in}{4.233097in}}%
\pgfpathlineto{\pgfqpoint{3.159542in}{4.161234in}}%
\pgfpathlineto{\pgfqpoint{3.186876in}{4.089306in}}%
\pgfpathlineto{\pgfqpoint{3.140504in}{4.076348in}}%
\pgfpathclose%
\pgfusepath{stroke,fill}%
\end{pgfscope}%
\begin{pgfscope}%
\pgfpathrectangle{\pgfqpoint{0.050000in}{0.083252in}}{\pgfqpoint{4.900000in}{4.900000in}}%
\pgfusepath{clip}%
\pgfsetbuttcap%
\pgfsetroundjoin%
\definecolor{currentfill}{rgb}{0.268128,0.515033,0.953479}%
\pgfsetfillcolor{currentfill}%
\pgfsetfillopacity{0.950000}%
\pgfsetlinewidth{0.200750pt}%
\definecolor{currentstroke}{rgb}{0.200000,0.200000,0.200000}%
\pgfsetstrokecolor{currentstroke}%
\pgfsetdash{}{0pt}%
\pgfpathmoveto{\pgfqpoint{2.992859in}{4.194570in}}%
\pgfpathlineto{\pgfqpoint{2.965568in}{4.266574in}}%
\pgfpathlineto{\pgfqpoint{2.938302in}{4.338512in}}%
\pgfpathlineto{\pgfqpoint{2.984845in}{4.351239in}}%
\pgfpathlineto{\pgfqpoint{3.031308in}{4.363944in}}%
\pgfpathlineto{\pgfqpoint{3.058569in}{4.292143in}}%
\pgfpathlineto{\pgfqpoint{3.085856in}{4.220277in}}%
\pgfpathlineto{\pgfqpoint{3.039398in}{4.207435in}}%
\pgfpathlineto{\pgfqpoint{2.992859in}{4.194570in}}%
\pgfpathclose%
\pgfusepath{stroke,fill}%
\end{pgfscope}%
\begin{pgfscope}%
\pgfpathrectangle{\pgfqpoint{0.050000in}{0.083252in}}{\pgfqpoint{4.900000in}{4.900000in}}%
\pgfusepath{clip}%
\pgfsetbuttcap%
\pgfsetroundjoin%
\definecolor{currentfill}{rgb}{0.261484,0.444444,0.859639}%
\pgfsetfillcolor{currentfill}%
\pgfsetfillopacity{0.950000}%
\pgfsetlinewidth{0.200750pt}%
\definecolor{currentstroke}{rgb}{0.200000,0.200000,0.200000}%
\pgfsetstrokecolor{currentstroke}%
\pgfsetdash{}{0pt}%
\pgfpathmoveto{\pgfqpoint{3.195253in}{3.932160in}}%
\pgfpathlineto{\pgfqpoint{3.167866in}{4.004287in}}%
\pgfpathlineto{\pgfqpoint{3.140504in}{4.076348in}}%
\pgfpathlineto{\pgfqpoint{3.186876in}{4.089306in}}%
\pgfpathlineto{\pgfqpoint{3.214236in}{4.017313in}}%
\pgfpathlineto{\pgfqpoint{3.241621in}{3.945257in}}%
\pgfpathlineto{\pgfqpoint{3.195253in}{3.932160in}}%
\pgfpathclose%
\pgfusepath{stroke,fill}%
\end{pgfscope}%
\begin{pgfscope}%
\pgfpathrectangle{\pgfqpoint{0.050000in}{0.083252in}}{\pgfqpoint{4.900000in}{4.900000in}}%
\pgfusepath{clip}%
\pgfsetbuttcap%
\pgfsetroundjoin%
\definecolor{currentfill}{rgb}{0.264929,0.481046,0.908297}%
\pgfsetfillcolor{currentfill}%
\pgfsetfillopacity{0.950000}%
\pgfsetlinewidth{0.200750pt}%
\definecolor{currentstroke}{rgb}{0.200000,0.200000,0.200000}%
\pgfsetstrokecolor{currentstroke}%
\pgfsetdash{}{0pt}%
\pgfpathmoveto{\pgfqpoint{3.047516in}{4.050365in}}%
\pgfpathlineto{\pgfqpoint{3.020175in}{4.122501in}}%
\pgfpathlineto{\pgfqpoint{2.992859in}{4.194570in}}%
\pgfpathlineto{\pgfqpoint{3.039398in}{4.207435in}}%
\pgfpathlineto{\pgfqpoint{3.085856in}{4.220277in}}%
\pgfpathlineto{\pgfqpoint{3.113167in}{4.148345in}}%
\pgfpathlineto{\pgfqpoint{3.140504in}{4.076348in}}%
\pgfpathlineto{\pgfqpoint{3.094051in}{4.063368in}}%
\pgfpathlineto{\pgfqpoint{3.047516in}{4.050365in}}%
\pgfpathclose%
\pgfusepath{stroke,fill}%
\end{pgfscope}%
\begin{pgfscope}%
\pgfpathrectangle{\pgfqpoint{0.050000in}{0.083252in}}{\pgfqpoint{4.900000in}{4.900000in}}%
\pgfusepath{clip}%
\pgfsetbuttcap%
\pgfsetroundjoin%
\definecolor{currentfill}{rgb}{0.270588,0.541176,0.988235}%
\pgfsetfillcolor{currentfill}%
\pgfsetfillopacity{0.950000}%
\pgfsetlinewidth{0.200750pt}%
\definecolor{currentstroke}{rgb}{0.200000,0.200000,0.200000}%
\pgfsetstrokecolor{currentstroke}%
\pgfsetdash{}{0pt}%
\pgfpathmoveto{\pgfqpoint{2.844971in}{4.312990in}}%
\pgfpathlineto{\pgfqpoint{2.817726in}{4.384999in}}%
\pgfpathlineto{\pgfqpoint{2.864434in}{4.397702in}}%
\pgfpathlineto{\pgfqpoint{2.911061in}{4.410383in}}%
\pgfpathlineto{\pgfqpoint{2.938302in}{4.338512in}}%
\pgfpathlineto{\pgfqpoint{2.891677in}{4.325762in}}%
\pgfpathlineto{\pgfqpoint{2.844971in}{4.312990in}}%
\pgfpathclose%
\pgfusepath{stroke,fill}%
\end{pgfscope}%
\begin{pgfscope}%
\pgfpathrectangle{\pgfqpoint{0.050000in}{0.083252in}}{\pgfqpoint{4.900000in}{4.900000in}}%
\pgfusepath{clip}%
\pgfsetbuttcap%
\pgfsetroundjoin%
\definecolor{currentfill}{rgb}{0.258285,0.410458,0.814456}%
\pgfsetfillcolor{currentfill}%
\pgfsetfillopacity{0.950000}%
\pgfsetlinewidth{0.200750pt}%
\definecolor{currentstroke}{rgb}{0.200000,0.200000,0.200000}%
\pgfsetstrokecolor{currentstroke}%
\pgfsetdash{}{0pt}%
\pgfpathmoveto{\pgfqpoint{3.250103in}{3.787719in}}%
\pgfpathlineto{\pgfqpoint{3.222665in}{3.859971in}}%
\pgfpathlineto{\pgfqpoint{3.195253in}{3.932160in}}%
\pgfpathlineto{\pgfqpoint{3.241621in}{3.945257in}}%
\pgfpathlineto{\pgfqpoint{3.269031in}{3.873138in}}%
\pgfpathlineto{\pgfqpoint{3.296466in}{3.800956in}}%
\pgfpathlineto{\pgfqpoint{3.250103in}{3.787719in}}%
\pgfpathclose%
\pgfusepath{stroke,fill}%
\end{pgfscope}%
\begin{pgfscope}%
\pgfpathrectangle{\pgfqpoint{0.050000in}{0.083252in}}{\pgfqpoint{4.900000in}{4.900000in}}%
\pgfusepath{clip}%
\pgfsetbuttcap%
\pgfsetroundjoin%
\definecolor{currentfill}{rgb}{0.268128,0.515033,0.953479}%
\pgfsetfillcolor{currentfill}%
\pgfsetfillopacity{0.950000}%
\pgfsetlinewidth{0.200750pt}%
\definecolor{currentstroke}{rgb}{0.200000,0.200000,0.200000}%
\pgfsetstrokecolor{currentstroke}%
\pgfsetdash{}{0pt}%
\pgfpathmoveto{\pgfqpoint{2.899537in}{4.168773in}}%
\pgfpathlineto{\pgfqpoint{2.872242in}{4.240915in}}%
\pgfpathlineto{\pgfqpoint{2.844971in}{4.312990in}}%
\pgfpathlineto{\pgfqpoint{2.891677in}{4.325762in}}%
\pgfpathlineto{\pgfqpoint{2.938302in}{4.338512in}}%
\pgfpathlineto{\pgfqpoint{2.965568in}{4.266574in}}%
\pgfpathlineto{\pgfqpoint{2.992859in}{4.194570in}}%
\pgfpathlineto{\pgfqpoint{2.946239in}{4.181683in}}%
\pgfpathlineto{\pgfqpoint{2.899537in}{4.168773in}}%
\pgfpathclose%
\pgfusepath{stroke,fill}%
\end{pgfscope}%
\begin{pgfscope}%
\pgfpathrectangle{\pgfqpoint{0.050000in}{0.083252in}}{\pgfqpoint{4.900000in}{4.900000in}}%
\pgfusepath{clip}%
\pgfsetbuttcap%
\pgfsetroundjoin%
\definecolor{currentfill}{rgb}{0.261484,0.444444,0.859639}%
\pgfsetfillcolor{currentfill}%
\pgfsetfillopacity{0.950000}%
\pgfsetlinewidth{0.200750pt}%
\definecolor{currentstroke}{rgb}{0.200000,0.200000,0.200000}%
\pgfsetstrokecolor{currentstroke}%
\pgfsetdash{}{0pt}%
\pgfpathmoveto{\pgfqpoint{3.102275in}{3.905899in}}%
\pgfpathlineto{\pgfqpoint{3.074883in}{3.978165in}}%
\pgfpathlineto{\pgfqpoint{3.047516in}{4.050365in}}%
\pgfpathlineto{\pgfqpoint{3.094051in}{4.063368in}}%
\pgfpathlineto{\pgfqpoint{3.140504in}{4.076348in}}%
\pgfpathlineto{\pgfqpoint{3.167866in}{4.004287in}}%
\pgfpathlineto{\pgfqpoint{3.195253in}{3.932160in}}%
\pgfpathlineto{\pgfqpoint{3.148805in}{3.919041in}}%
\pgfpathlineto{\pgfqpoint{3.102275in}{3.905899in}}%
\pgfpathclose%
\pgfusepath{stroke,fill}%
\end{pgfscope}%
\begin{pgfscope}%
\pgfpathrectangle{\pgfqpoint{0.050000in}{0.083252in}}{\pgfqpoint{4.900000in}{4.900000in}}%
\pgfusepath{clip}%
\pgfsetbuttcap%
\pgfsetroundjoin%
\definecolor{currentfill}{rgb}{0.254840,0.373856,0.765798}%
\pgfsetfillcolor{currentfill}%
\pgfsetfillopacity{0.950000}%
\pgfsetlinewidth{0.200750pt}%
\definecolor{currentstroke}{rgb}{0.200000,0.200000,0.200000}%
\pgfsetstrokecolor{currentstroke}%
\pgfsetdash{}{0pt}%
\pgfpathmoveto{\pgfqpoint{3.305055in}{3.643035in}}%
\pgfpathlineto{\pgfqpoint{3.277566in}{3.715406in}}%
\pgfpathlineto{\pgfqpoint{3.250103in}{3.787719in}}%
\pgfpathlineto{\pgfqpoint{3.296466in}{3.800956in}}%
\pgfpathlineto{\pgfqpoint{3.323927in}{3.728715in}}%
\pgfpathlineto{\pgfqpoint{3.351413in}{3.656416in}}%
\pgfpathlineto{\pgfqpoint{3.305055in}{3.643035in}}%
\pgfpathclose%
\pgfusepath{stroke,fill}%
\end{pgfscope}%
\begin{pgfscope}%
\pgfpathrectangle{\pgfqpoint{0.050000in}{0.083252in}}{\pgfqpoint{4.900000in}{4.900000in}}%
\pgfusepath{clip}%
\pgfsetbuttcap%
\pgfsetroundjoin%
\definecolor{currentfill}{rgb}{0.264929,0.481046,0.908297}%
\pgfsetfillcolor{currentfill}%
\pgfsetfillopacity{0.950000}%
\pgfsetlinewidth{0.200750pt}%
\definecolor{currentstroke}{rgb}{0.200000,0.200000,0.200000}%
\pgfsetstrokecolor{currentstroke}%
\pgfsetdash{}{0pt}%
\pgfpathmoveto{\pgfqpoint{2.954204in}{4.024292in}}%
\pgfpathlineto{\pgfqpoint{2.926858in}{4.096566in}}%
\pgfpathlineto{\pgfqpoint{2.899537in}{4.168773in}}%
\pgfpathlineto{\pgfqpoint{2.946239in}{4.181683in}}%
\pgfpathlineto{\pgfqpoint{2.992859in}{4.194570in}}%
\pgfpathlineto{\pgfqpoint{3.020175in}{4.122501in}}%
\pgfpathlineto{\pgfqpoint{3.047516in}{4.050365in}}%
\pgfpathlineto{\pgfqpoint{3.000901in}{4.037340in}}%
\pgfpathlineto{\pgfqpoint{2.954204in}{4.024292in}}%
\pgfpathclose%
\pgfusepath{stroke,fill}%
\end{pgfscope}%
\begin{pgfscope}%
\pgfpathrectangle{\pgfqpoint{0.050000in}{0.083252in}}{\pgfqpoint{4.900000in}{4.900000in}}%
\pgfusepath{clip}%
\pgfsetbuttcap%
\pgfsetroundjoin%
\definecolor{currentfill}{rgb}{0.258285,0.410458,0.814456}%
\pgfsetfillcolor{currentfill}%
\pgfsetfillopacity{0.950000}%
\pgfsetlinewidth{0.200750pt}%
\definecolor{currentstroke}{rgb}{0.200000,0.200000,0.200000}%
\pgfsetstrokecolor{currentstroke}%
\pgfsetdash{}{0pt}%
\pgfpathmoveto{\pgfqpoint{3.157135in}{3.761176in}}%
\pgfpathlineto{\pgfqpoint{3.129692in}{3.833569in}}%
\pgfpathlineto{\pgfqpoint{3.102275in}{3.905899in}}%
\pgfpathlineto{\pgfqpoint{3.148805in}{3.919041in}}%
\pgfpathlineto{\pgfqpoint{3.195253in}{3.932160in}}%
\pgfpathlineto{\pgfqpoint{3.222665in}{3.859971in}}%
\pgfpathlineto{\pgfqpoint{3.250103in}{3.787719in}}%
\pgfpathlineto{\pgfqpoint{3.203660in}{3.774459in}}%
\pgfpathlineto{\pgfqpoint{3.157135in}{3.761176in}}%
\pgfpathclose%
\pgfusepath{stroke,fill}%
\end{pgfscope}%
\begin{pgfscope}%
\pgfpathrectangle{\pgfqpoint{0.050000in}{0.083252in}}{\pgfqpoint{4.900000in}{4.900000in}}%
\pgfusepath{clip}%
\pgfsetbuttcap%
\pgfsetroundjoin%
\definecolor{currentfill}{rgb}{0.270588,0.541176,0.988235}%
\pgfsetfillcolor{currentfill}%
\pgfsetfillopacity{0.950000}%
\pgfsetlinewidth{0.200750pt}%
\definecolor{currentstroke}{rgb}{0.200000,0.200000,0.200000}%
\pgfsetstrokecolor{currentstroke}%
\pgfsetdash{}{0pt}%
\pgfpathmoveto{\pgfqpoint{2.751315in}{4.287379in}}%
\pgfpathlineto{\pgfqpoint{2.724065in}{4.359525in}}%
\pgfpathlineto{\pgfqpoint{2.770936in}{4.372273in}}%
\pgfpathlineto{\pgfqpoint{2.817726in}{4.384999in}}%
\pgfpathlineto{\pgfqpoint{2.844971in}{4.312990in}}%
\pgfpathlineto{\pgfqpoint{2.798184in}{4.300196in}}%
\pgfpathlineto{\pgfqpoint{2.751315in}{4.287379in}}%
\pgfpathclose%
\pgfusepath{stroke,fill}%
\end{pgfscope}%
\begin{pgfscope}%
\pgfpathrectangle{\pgfqpoint{0.050000in}{0.083252in}}{\pgfqpoint{4.900000in}{4.900000in}}%
\pgfusepath{clip}%
\pgfsetbuttcap%
\pgfsetroundjoin%
\definecolor{currentfill}{rgb}{0.251642,0.339869,0.720615}%
\pgfsetfillcolor{currentfill}%
\pgfsetfillopacity{0.950000}%
\pgfsetlinewidth{0.200750pt}%
\definecolor{currentstroke}{rgb}{0.200000,0.200000,0.200000}%
\pgfsetstrokecolor{currentstroke}%
\pgfsetdash{}{0pt}%
\pgfpathmoveto{\pgfqpoint{3.360109in}{3.498130in}}%
\pgfpathlineto{\pgfqpoint{3.332569in}{3.570608in}}%
\pgfpathlineto{\pgfqpoint{3.305055in}{3.643035in}}%
\pgfpathlineto{\pgfqpoint{3.351413in}{3.656416in}}%
\pgfpathlineto{\pgfqpoint{3.378925in}{3.584062in}}%
\pgfpathlineto{\pgfqpoint{3.406462in}{3.511659in}}%
\pgfpathlineto{\pgfqpoint{3.360109in}{3.498130in}}%
\pgfpathclose%
\pgfusepath{stroke,fill}%
\end{pgfscope}%
\begin{pgfscope}%
\pgfpathrectangle{\pgfqpoint{0.050000in}{0.083252in}}{\pgfqpoint{4.900000in}{4.900000in}}%
\pgfusepath{clip}%
\pgfsetbuttcap%
\pgfsetroundjoin%
\definecolor{currentfill}{rgb}{0.268128,0.515033,0.953479}%
\pgfsetfillcolor{currentfill}%
\pgfsetfillopacity{0.950000}%
\pgfsetlinewidth{0.200750pt}%
\definecolor{currentstroke}{rgb}{0.200000,0.200000,0.200000}%
\pgfsetstrokecolor{currentstroke}%
\pgfsetdash{}{0pt}%
\pgfpathmoveto{\pgfqpoint{2.805890in}{4.142886in}}%
\pgfpathlineto{\pgfqpoint{2.778589in}{4.215166in}}%
\pgfpathlineto{\pgfqpoint{2.751315in}{4.287379in}}%
\pgfpathlineto{\pgfqpoint{2.798184in}{4.300196in}}%
\pgfpathlineto{\pgfqpoint{2.844971in}{4.312990in}}%
\pgfpathlineto{\pgfqpoint{2.872242in}{4.240915in}}%
\pgfpathlineto{\pgfqpoint{2.899537in}{4.168773in}}%
\pgfpathlineto{\pgfqpoint{2.852754in}{4.155841in}}%
\pgfpathlineto{\pgfqpoint{2.805890in}{4.142886in}}%
\pgfpathclose%
\pgfusepath{stroke,fill}%
\end{pgfscope}%
\begin{pgfscope}%
\pgfpathrectangle{\pgfqpoint{0.050000in}{0.083252in}}{\pgfqpoint{4.900000in}{4.900000in}}%
\pgfusepath{clip}%
\pgfsetbuttcap%
\pgfsetroundjoin%
\definecolor{currentfill}{rgb}{0.261484,0.444444,0.859639}%
\pgfsetfillcolor{currentfill}%
\pgfsetfillopacity{0.950000}%
\pgfsetlinewidth{0.200750pt}%
\definecolor{currentstroke}{rgb}{0.200000,0.200000,0.200000}%
\pgfsetstrokecolor{currentstroke}%
\pgfsetdash{}{0pt}%
\pgfpathmoveto{\pgfqpoint{3.008972in}{3.879547in}}%
\pgfpathlineto{\pgfqpoint{2.981576in}{3.951952in}}%
\pgfpathlineto{\pgfqpoint{2.954204in}{4.024292in}}%
\pgfpathlineto{\pgfqpoint{3.000901in}{4.037340in}}%
\pgfpathlineto{\pgfqpoint{3.047516in}{4.050365in}}%
\pgfpathlineto{\pgfqpoint{3.074883in}{3.978165in}}%
\pgfpathlineto{\pgfqpoint{3.102275in}{3.905899in}}%
\pgfpathlineto{\pgfqpoint{3.055664in}{3.892735in}}%
\pgfpathlineto{\pgfqpoint{3.008972in}{3.879547in}}%
\pgfpathclose%
\pgfusepath{stroke,fill}%
\end{pgfscope}%
\begin{pgfscope}%
\pgfpathrectangle{\pgfqpoint{0.050000in}{0.083252in}}{\pgfqpoint{4.900000in}{4.900000in}}%
\pgfusepath{clip}%
\pgfsetbuttcap%
\pgfsetroundjoin%
\definecolor{currentfill}{rgb}{0.254840,0.373856,0.765798}%
\pgfsetfillcolor{currentfill}%
\pgfsetfillopacity{0.950000}%
\pgfsetlinewidth{0.200750pt}%
\definecolor{currentstroke}{rgb}{0.200000,0.200000,0.200000}%
\pgfsetstrokecolor{currentstroke}%
\pgfsetdash{}{0pt}%
\pgfpathmoveto{\pgfqpoint{3.212097in}{3.616205in}}%
\pgfpathlineto{\pgfqpoint{3.184603in}{3.688721in}}%
\pgfpathlineto{\pgfqpoint{3.157135in}{3.761176in}}%
\pgfpathlineto{\pgfqpoint{3.203660in}{3.774459in}}%
\pgfpathlineto{\pgfqpoint{3.250103in}{3.787719in}}%
\pgfpathlineto{\pgfqpoint{3.277566in}{3.715406in}}%
\pgfpathlineto{\pgfqpoint{3.305055in}{3.643035in}}%
\pgfpathlineto{\pgfqpoint{3.258617in}{3.629631in}}%
\pgfpathlineto{\pgfqpoint{3.212097in}{3.616205in}}%
\pgfpathclose%
\pgfusepath{stroke,fill}%
\end{pgfscope}%
\begin{pgfscope}%
\pgfpathrectangle{\pgfqpoint{0.050000in}{0.083252in}}{\pgfqpoint{4.900000in}{4.900000in}}%
\pgfusepath{clip}%
\pgfsetbuttcap%
\pgfsetroundjoin%
\definecolor{currentfill}{rgb}{0.248443,0.305882,0.675433}%
\pgfsetfillcolor{currentfill}%
\pgfsetfillopacity{0.950000}%
\pgfsetlinewidth{0.200750pt}%
\definecolor{currentstroke}{rgb}{0.200000,0.200000,0.200000}%
\pgfsetstrokecolor{currentstroke}%
\pgfsetdash{}{0pt}%
\pgfpathmoveto{\pgfqpoint{3.415267in}{3.353047in}}%
\pgfpathlineto{\pgfqpoint{3.387675in}{3.425607in}}%
\pgfpathlineto{\pgfqpoint{3.360109in}{3.498130in}}%
\pgfpathlineto{\pgfqpoint{3.406462in}{3.511659in}}%
\pgfpathlineto{\pgfqpoint{3.434026in}{3.439212in}}%
\pgfpathlineto{\pgfqpoint{3.461615in}{3.366730in}}%
\pgfpathlineto{\pgfqpoint{3.415267in}{3.353047in}}%
\pgfpathclose%
\pgfusepath{stroke,fill}%
\end{pgfscope}%
\begin{pgfscope}%
\pgfpathrectangle{\pgfqpoint{0.050000in}{0.083252in}}{\pgfqpoint{4.900000in}{4.900000in}}%
\pgfusepath{clip}%
\pgfsetbuttcap%
\pgfsetroundjoin%
\definecolor{currentfill}{rgb}{0.264683,0.478431,0.904821}%
\pgfsetfillcolor{currentfill}%
\pgfsetfillopacity{0.950000}%
\pgfsetlinewidth{0.200750pt}%
\definecolor{currentstroke}{rgb}{0.200000,0.200000,0.200000}%
\pgfsetstrokecolor{currentstroke}%
\pgfsetdash{}{0pt}%
\pgfpathmoveto{\pgfqpoint{2.860565in}{3.998127in}}%
\pgfpathlineto{\pgfqpoint{2.833215in}{4.070540in}}%
\pgfpathlineto{\pgfqpoint{2.805890in}{4.142886in}}%
\pgfpathlineto{\pgfqpoint{2.852754in}{4.155841in}}%
\pgfpathlineto{\pgfqpoint{2.899537in}{4.168773in}}%
\pgfpathlineto{\pgfqpoint{2.926858in}{4.096566in}}%
\pgfpathlineto{\pgfqpoint{2.954204in}{4.024292in}}%
\pgfpathlineto{\pgfqpoint{2.907426in}{4.011221in}}%
\pgfpathlineto{\pgfqpoint{2.860565in}{3.998127in}}%
\pgfpathclose%
\pgfusepath{stroke,fill}%
\end{pgfscope}%
\begin{pgfscope}%
\pgfpathrectangle{\pgfqpoint{0.050000in}{0.083252in}}{\pgfqpoint{4.900000in}{4.900000in}}%
\pgfusepath{clip}%
\pgfsetbuttcap%
\pgfsetroundjoin%
\definecolor{currentfill}{rgb}{0.258285,0.410458,0.814456}%
\pgfsetfillcolor{currentfill}%
\pgfsetfillopacity{0.950000}%
\pgfsetlinewidth{0.200750pt}%
\definecolor{currentstroke}{rgb}{0.200000,0.200000,0.200000}%
\pgfsetstrokecolor{currentstroke}%
\pgfsetdash{}{0pt}%
\pgfpathmoveto{\pgfqpoint{3.063842in}{3.734542in}}%
\pgfpathlineto{\pgfqpoint{3.036394in}{3.807076in}}%
\pgfpathlineto{\pgfqpoint{3.008972in}{3.879547in}}%
\pgfpathlineto{\pgfqpoint{3.055664in}{3.892735in}}%
\pgfpathlineto{\pgfqpoint{3.102275in}{3.905899in}}%
\pgfpathlineto{\pgfqpoint{3.129692in}{3.833569in}}%
\pgfpathlineto{\pgfqpoint{3.157135in}{3.761176in}}%
\pgfpathlineto{\pgfqpoint{3.110529in}{3.747870in}}%
\pgfpathlineto{\pgfqpoint{3.063842in}{3.734542in}}%
\pgfpathclose%
\pgfusepath{stroke,fill}%
\end{pgfscope}%
\begin{pgfscope}%
\pgfpathrectangle{\pgfqpoint{0.050000in}{0.083252in}}{\pgfqpoint{4.900000in}{4.900000in}}%
\pgfusepath{clip}%
\pgfsetbuttcap%
\pgfsetroundjoin%
\definecolor{currentfill}{rgb}{0.270588,0.541176,0.988235}%
\pgfsetfillcolor{currentfill}%
\pgfsetfillopacity{0.950000}%
\pgfsetlinewidth{0.200750pt}%
\definecolor{currentstroke}{rgb}{0.200000,0.200000,0.200000}%
\pgfsetstrokecolor{currentstroke}%
\pgfsetdash{}{0pt}%
\pgfpathmoveto{\pgfqpoint{2.657330in}{4.261678in}}%
\pgfpathlineto{\pgfqpoint{2.630077in}{4.333962in}}%
\pgfpathlineto{\pgfqpoint{2.677112in}{4.346755in}}%
\pgfpathlineto{\pgfqpoint{2.724065in}{4.359525in}}%
\pgfpathlineto{\pgfqpoint{2.751315in}{4.287379in}}%
\pgfpathlineto{\pgfqpoint{2.704364in}{4.274540in}}%
\pgfpathlineto{\pgfqpoint{2.657330in}{4.261678in}}%
\pgfpathclose%
\pgfusepath{stroke,fill}%
\end{pgfscope}%
\begin{pgfscope}%
\pgfpathrectangle{\pgfqpoint{0.050000in}{0.083252in}}{\pgfqpoint{4.900000in}{4.900000in}}%
\pgfusepath{clip}%
\pgfsetbuttcap%
\pgfsetroundjoin%
\definecolor{currentfill}{rgb}{0.251642,0.339869,0.720615}%
\pgfsetfillcolor{currentfill}%
\pgfsetfillopacity{0.950000}%
\pgfsetlinewidth{0.200750pt}%
\definecolor{currentstroke}{rgb}{0.200000,0.200000,0.200000}%
\pgfsetstrokecolor{currentstroke}%
\pgfsetdash{}{0pt}%
\pgfpathmoveto{\pgfqpoint{3.267161in}{3.471006in}}%
\pgfpathlineto{\pgfqpoint{3.239616in}{3.543632in}}%
\pgfpathlineto{\pgfqpoint{3.212097in}{3.616205in}}%
\pgfpathlineto{\pgfqpoint{3.258617in}{3.629631in}}%
\pgfpathlineto{\pgfqpoint{3.305055in}{3.643035in}}%
\pgfpathlineto{\pgfqpoint{3.332569in}{3.570608in}}%
\pgfpathlineto{\pgfqpoint{3.360109in}{3.498130in}}%
\pgfpathlineto{\pgfqpoint{3.313676in}{3.484579in}}%
\pgfpathlineto{\pgfqpoint{3.267161in}{3.471006in}}%
\pgfpathclose%
\pgfusepath{stroke,fill}%
\end{pgfscope}%
\begin{pgfscope}%
\pgfpathrectangle{\pgfqpoint{0.050000in}{0.083252in}}{\pgfqpoint{4.900000in}{4.900000in}}%
\pgfusepath{clip}%
\pgfsetbuttcap%
\pgfsetroundjoin%
\definecolor{currentfill}{rgb}{0.244998,0.269281,0.626774}%
\pgfsetfillcolor{currentfill}%
\pgfsetfillopacity{0.950000}%
\pgfsetlinewidth{0.200750pt}%
\definecolor{currentstroke}{rgb}{0.200000,0.200000,0.200000}%
\pgfsetstrokecolor{currentstroke}%
\pgfsetdash{}{0pt}%
\pgfpathmoveto{\pgfqpoint{3.470531in}{3.207857in}}%
\pgfpathlineto{\pgfqpoint{3.442886in}{3.280459in}}%
\pgfpathlineto{\pgfqpoint{3.415267in}{3.353047in}}%
\pgfpathlineto{\pgfqpoint{3.461615in}{3.366730in}}%
\pgfpathlineto{\pgfqpoint{3.489231in}{3.294223in}}%
\pgfpathlineto{\pgfqpoint{3.516874in}{3.221705in}}%
\pgfpathlineto{\pgfqpoint{3.470531in}{3.207857in}}%
\pgfpathclose%
\pgfusepath{stroke,fill}%
\end{pgfscope}%
\begin{pgfscope}%
\pgfpathrectangle{\pgfqpoint{0.050000in}{0.083252in}}{\pgfqpoint{4.900000in}{4.900000in}}%
\pgfusepath{clip}%
\pgfsetbuttcap%
\pgfsetroundjoin%
\definecolor{currentfill}{rgb}{0.268128,0.515033,0.953479}%
\pgfsetfillcolor{currentfill}%
\pgfsetfillopacity{0.950000}%
\pgfsetlinewidth{0.200750pt}%
\definecolor{currentstroke}{rgb}{0.200000,0.200000,0.200000}%
\pgfsetstrokecolor{currentstroke}%
\pgfsetdash{}{0pt}%
\pgfpathmoveto{\pgfqpoint{2.711914in}{4.116908in}}%
\pgfpathlineto{\pgfqpoint{2.684610in}{4.189326in}}%
\pgfpathlineto{\pgfqpoint{2.657330in}{4.261678in}}%
\pgfpathlineto{\pgfqpoint{2.704364in}{4.274540in}}%
\pgfpathlineto{\pgfqpoint{2.751315in}{4.287379in}}%
\pgfpathlineto{\pgfqpoint{2.778589in}{4.215166in}}%
\pgfpathlineto{\pgfqpoint{2.805890in}{4.142886in}}%
\pgfpathlineto{\pgfqpoint{2.758943in}{4.129908in}}%
\pgfpathlineto{\pgfqpoint{2.711914in}{4.116908in}}%
\pgfpathclose%
\pgfusepath{stroke,fill}%
\end{pgfscope}%
\begin{pgfscope}%
\pgfpathrectangle{\pgfqpoint{0.050000in}{0.083252in}}{\pgfqpoint{4.900000in}{4.900000in}}%
\pgfusepath{clip}%
\pgfsetbuttcap%
\pgfsetroundjoin%
\definecolor{currentfill}{rgb}{0.261484,0.444444,0.859639}%
\pgfsetfillcolor{currentfill}%
\pgfsetfillopacity{0.950000}%
\pgfsetlinewidth{0.200750pt}%
\definecolor{currentstroke}{rgb}{0.200000,0.200000,0.200000}%
\pgfsetstrokecolor{currentstroke}%
\pgfsetdash{}{0pt}%
\pgfpathmoveto{\pgfqpoint{2.915343in}{3.853103in}}%
\pgfpathlineto{\pgfqpoint{2.887941in}{3.925648in}}%
\pgfpathlineto{\pgfqpoint{2.860565in}{3.998127in}}%
\pgfpathlineto{\pgfqpoint{2.907426in}{4.011221in}}%
\pgfpathlineto{\pgfqpoint{2.954204in}{4.024292in}}%
\pgfpathlineto{\pgfqpoint{2.981576in}{3.951952in}}%
\pgfpathlineto{\pgfqpoint{3.008972in}{3.879547in}}%
\pgfpathlineto{\pgfqpoint{2.962198in}{3.866336in}}%
\pgfpathlineto{\pgfqpoint{2.915343in}{3.853103in}}%
\pgfpathclose%
\pgfusepath{stroke,fill}%
\end{pgfscope}%
\begin{pgfscope}%
\pgfpathrectangle{\pgfqpoint{0.050000in}{0.083252in}}{\pgfqpoint{4.900000in}{4.900000in}}%
\pgfusepath{clip}%
\pgfsetbuttcap%
\pgfsetroundjoin%
\definecolor{currentfill}{rgb}{0.254840,0.373856,0.765798}%
\pgfsetfillcolor{currentfill}%
\pgfsetfillopacity{0.950000}%
\pgfsetlinewidth{0.200750pt}%
\definecolor{currentstroke}{rgb}{0.200000,0.200000,0.200000}%
\pgfsetstrokecolor{currentstroke}%
\pgfsetdash{}{0pt}%
\pgfpathmoveto{\pgfqpoint{3.118813in}{3.589284in}}%
\pgfpathlineto{\pgfqpoint{3.091315in}{3.661944in}}%
\pgfpathlineto{\pgfqpoint{3.063842in}{3.734542in}}%
\pgfpathlineto{\pgfqpoint{3.110529in}{3.747870in}}%
\pgfpathlineto{\pgfqpoint{3.157135in}{3.761176in}}%
\pgfpathlineto{\pgfqpoint{3.184603in}{3.688721in}}%
\pgfpathlineto{\pgfqpoint{3.212097in}{3.616205in}}%
\pgfpathlineto{\pgfqpoint{3.165496in}{3.602756in}}%
\pgfpathlineto{\pgfqpoint{3.118813in}{3.589284in}}%
\pgfpathclose%
\pgfusepath{stroke,fill}%
\end{pgfscope}%
\begin{pgfscope}%
\pgfpathrectangle{\pgfqpoint{0.050000in}{0.083252in}}{\pgfqpoint{4.900000in}{4.900000in}}%
\pgfusepath{clip}%
\pgfsetbuttcap%
\pgfsetroundjoin%
\definecolor{currentfill}{rgb}{0.248197,0.303268,0.671957}%
\pgfsetfillcolor{currentfill}%
\pgfsetfillopacity{0.950000}%
\pgfsetlinewidth{0.200750pt}%
\definecolor{currentstroke}{rgb}{0.200000,0.200000,0.200000}%
\pgfsetstrokecolor{currentstroke}%
\pgfsetdash{}{0pt}%
\pgfpathmoveto{\pgfqpoint{3.322329in}{3.325614in}}%
\pgfpathlineto{\pgfqpoint{3.294732in}{3.398331in}}%
\pgfpathlineto{\pgfqpoint{3.267161in}{3.471006in}}%
\pgfpathlineto{\pgfqpoint{3.313676in}{3.484579in}}%
\pgfpathlineto{\pgfqpoint{3.360109in}{3.498130in}}%
\pgfpathlineto{\pgfqpoint{3.387675in}{3.425607in}}%
\pgfpathlineto{\pgfqpoint{3.415267in}{3.353047in}}%
\pgfpathlineto{\pgfqpoint{3.368839in}{3.339341in}}%
\pgfpathlineto{\pgfqpoint{3.322329in}{3.325614in}}%
\pgfpathclose%
\pgfusepath{stroke,fill}%
\end{pgfscope}%
\begin{pgfscope}%
\pgfpathrectangle{\pgfqpoint{0.050000in}{0.083252in}}{\pgfqpoint{4.900000in}{4.900000in}}%
\pgfusepath{clip}%
\pgfsetbuttcap%
\pgfsetroundjoin%
\definecolor{currentfill}{rgb}{0.241799,0.235294,0.581592}%
\pgfsetfillcolor{currentfill}%
\pgfsetfillopacity{0.950000}%
\pgfsetlinewidth{0.200750pt}%
\definecolor{currentstroke}{rgb}{0.200000,0.200000,0.200000}%
\pgfsetstrokecolor{currentstroke}%
\pgfsetdash{}{0pt}%
\pgfpathmoveto{\pgfqpoint{3.525904in}{3.062687in}}%
\pgfpathlineto{\pgfqpoint{3.498204in}{3.135259in}}%
\pgfpathlineto{\pgfqpoint{3.470531in}{3.207857in}}%
\pgfpathlineto{\pgfqpoint{3.516874in}{3.221705in}}%
\pgfpathlineto{\pgfqpoint{3.544545in}{3.149194in}}%
\pgfpathlineto{\pgfqpoint{3.572243in}{3.076714in}}%
\pgfpathlineto{\pgfqpoint{3.525904in}{3.062687in}}%
\pgfpathclose%
\pgfusepath{stroke,fill}%
\end{pgfscope}%
\begin{pgfscope}%
\pgfpathrectangle{\pgfqpoint{0.050000in}{0.083252in}}{\pgfqpoint{4.900000in}{4.900000in}}%
\pgfusepath{clip}%
\pgfsetbuttcap%
\pgfsetroundjoin%
\definecolor{currentfill}{rgb}{0.264683,0.478431,0.904821}%
\pgfsetfillcolor{currentfill}%
\pgfsetfillopacity{0.950000}%
\pgfsetlinewidth{0.200750pt}%
\definecolor{currentstroke}{rgb}{0.200000,0.200000,0.200000}%
\pgfsetstrokecolor{currentstroke}%
\pgfsetdash{}{0pt}%
\pgfpathmoveto{\pgfqpoint{2.766599in}{3.971870in}}%
\pgfpathlineto{\pgfqpoint{2.739244in}{4.044423in}}%
\pgfpathlineto{\pgfqpoint{2.711914in}{4.116908in}}%
\pgfpathlineto{\pgfqpoint{2.758943in}{4.129908in}}%
\pgfpathlineto{\pgfqpoint{2.805890in}{4.142886in}}%
\pgfpathlineto{\pgfqpoint{2.833215in}{4.070540in}}%
\pgfpathlineto{\pgfqpoint{2.860565in}{3.998127in}}%
\pgfpathlineto{\pgfqpoint{2.813623in}{3.985010in}}%
\pgfpathlineto{\pgfqpoint{2.766599in}{3.971870in}}%
\pgfpathclose%
\pgfusepath{stroke,fill}%
\end{pgfscope}%
\begin{pgfscope}%
\pgfpathrectangle{\pgfqpoint{0.050000in}{0.083252in}}{\pgfqpoint{4.900000in}{4.900000in}}%
\pgfusepath{clip}%
\pgfsetbuttcap%
\pgfsetroundjoin%
\definecolor{currentfill}{rgb}{0.258039,0.407843,0.810980}%
\pgfsetfillcolor{currentfill}%
\pgfsetfillopacity{0.950000}%
\pgfsetlinewidth{0.200750pt}%
\definecolor{currentstroke}{rgb}{0.200000,0.200000,0.200000}%
\pgfsetstrokecolor{currentstroke}%
\pgfsetdash{}{0pt}%
\pgfpathmoveto{\pgfqpoint{2.970222in}{3.707815in}}%
\pgfpathlineto{\pgfqpoint{2.942769in}{3.780492in}}%
\pgfpathlineto{\pgfqpoint{2.915343in}{3.853103in}}%
\pgfpathlineto{\pgfqpoint{2.962198in}{3.866336in}}%
\pgfpathlineto{\pgfqpoint{3.008972in}{3.879547in}}%
\pgfpathlineto{\pgfqpoint{3.036394in}{3.807076in}}%
\pgfpathlineto{\pgfqpoint{3.063842in}{3.734542in}}%
\pgfpathlineto{\pgfqpoint{3.017073in}{3.721190in}}%
\pgfpathlineto{\pgfqpoint{2.970222in}{3.707815in}}%
\pgfpathclose%
\pgfusepath{stroke,fill}%
\end{pgfscope}%
\begin{pgfscope}%
\pgfpathrectangle{\pgfqpoint{0.050000in}{0.083252in}}{\pgfqpoint{4.900000in}{4.900000in}}%
\pgfusepath{clip}%
\pgfsetbuttcap%
\pgfsetroundjoin%
\definecolor{currentfill}{rgb}{0.270588,0.541176,0.988235}%
\pgfsetfillcolor{currentfill}%
\pgfsetfillopacity{0.950000}%
\pgfsetlinewidth{0.200750pt}%
\definecolor{currentstroke}{rgb}{0.200000,0.200000,0.200000}%
\pgfsetstrokecolor{currentstroke}%
\pgfsetdash{}{0pt}%
\pgfpathmoveto{\pgfqpoint{2.563017in}{4.235886in}}%
\pgfpathlineto{\pgfqpoint{2.535759in}{4.308309in}}%
\pgfpathlineto{\pgfqpoint{2.582959in}{4.321147in}}%
\pgfpathlineto{\pgfqpoint{2.630077in}{4.333962in}}%
\pgfpathlineto{\pgfqpoint{2.657330in}{4.261678in}}%
\pgfpathlineto{\pgfqpoint{2.610215in}{4.248793in}}%
\pgfpathlineto{\pgfqpoint{2.563017in}{4.235886in}}%
\pgfpathclose%
\pgfusepath{stroke,fill}%
\end{pgfscope}%
\begin{pgfscope}%
\pgfpathrectangle{\pgfqpoint{0.050000in}{0.083252in}}{\pgfqpoint{4.900000in}{4.900000in}}%
\pgfusepath{clip}%
\pgfsetbuttcap%
\pgfsetroundjoin%
\definecolor{currentfill}{rgb}{0.251642,0.339869,0.720615}%
\pgfsetfillcolor{currentfill}%
\pgfsetfillopacity{0.950000}%
\pgfsetlinewidth{0.200750pt}%
\definecolor{currentstroke}{rgb}{0.200000,0.200000,0.200000}%
\pgfsetstrokecolor{currentstroke}%
\pgfsetdash{}{0pt}%
\pgfpathmoveto{\pgfqpoint{3.173887in}{3.443790in}}%
\pgfpathlineto{\pgfqpoint{3.146337in}{3.516565in}}%
\pgfpathlineto{\pgfqpoint{3.118813in}{3.589284in}}%
\pgfpathlineto{\pgfqpoint{3.165496in}{3.602756in}}%
\pgfpathlineto{\pgfqpoint{3.212097in}{3.616205in}}%
\pgfpathlineto{\pgfqpoint{3.239616in}{3.543632in}}%
\pgfpathlineto{\pgfqpoint{3.267161in}{3.471006in}}%
\pgfpathlineto{\pgfqpoint{3.220565in}{3.457409in}}%
\pgfpathlineto{\pgfqpoint{3.173887in}{3.443790in}}%
\pgfpathclose%
\pgfusepath{stroke,fill}%
\end{pgfscope}%
\begin{pgfscope}%
\pgfpathrectangle{\pgfqpoint{0.050000in}{0.083252in}}{\pgfqpoint{4.900000in}{4.900000in}}%
\pgfusepath{clip}%
\pgfsetbuttcap%
\pgfsetroundjoin%
\definecolor{currentfill}{rgb}{0.244998,0.269281,0.626774}%
\pgfsetfillcolor{currentfill}%
\pgfsetfillopacity{0.950000}%
\pgfsetlinewidth{0.200750pt}%
\definecolor{currentstroke}{rgb}{0.200000,0.200000,0.200000}%
\pgfsetstrokecolor{currentstroke}%
\pgfsetdash{}{0pt}%
\pgfpathmoveto{\pgfqpoint{3.377602in}{3.180096in}}%
\pgfpathlineto{\pgfqpoint{3.349952in}{3.252865in}}%
\pgfpathlineto{\pgfqpoint{3.322329in}{3.325614in}}%
\pgfpathlineto{\pgfqpoint{3.368839in}{3.339341in}}%
\pgfpathlineto{\pgfqpoint{3.415267in}{3.353047in}}%
\pgfpathlineto{\pgfqpoint{3.442886in}{3.280459in}}%
\pgfpathlineto{\pgfqpoint{3.470531in}{3.207857in}}%
\pgfpathlineto{\pgfqpoint{3.424107in}{3.193987in}}%
\pgfpathlineto{\pgfqpoint{3.377602in}{3.180096in}}%
\pgfpathclose%
\pgfusepath{stroke,fill}%
\end{pgfscope}%
\begin{pgfscope}%
\pgfpathrectangle{\pgfqpoint{0.050000in}{0.083252in}}{\pgfqpoint{4.900000in}{4.900000in}}%
\pgfusepath{clip}%
\pgfsetbuttcap%
\pgfsetroundjoin%
\definecolor{currentfill}{rgb}{0.267882,0.512418,0.950004}%
\pgfsetfillcolor{currentfill}%
\pgfsetfillopacity{0.950000}%
\pgfsetlinewidth{0.200750pt}%
\definecolor{currentstroke}{rgb}{0.200000,0.200000,0.200000}%
\pgfsetstrokecolor{currentstroke}%
\pgfsetdash{}{0pt}%
\pgfpathmoveto{\pgfqpoint{2.617609in}{4.090838in}}%
\pgfpathlineto{\pgfqpoint{2.590300in}{4.163396in}}%
\pgfpathlineto{\pgfqpoint{2.563017in}{4.235886in}}%
\pgfpathlineto{\pgfqpoint{2.610215in}{4.248793in}}%
\pgfpathlineto{\pgfqpoint{2.657330in}{4.261678in}}%
\pgfpathlineto{\pgfqpoint{2.684610in}{4.189326in}}%
\pgfpathlineto{\pgfqpoint{2.711914in}{4.116908in}}%
\pgfpathlineto{\pgfqpoint{2.664803in}{4.103884in}}%
\pgfpathlineto{\pgfqpoint{2.617609in}{4.090838in}}%
\pgfpathclose%
\pgfusepath{stroke,fill}%
\end{pgfscope}%
\begin{pgfscope}%
\pgfpathrectangle{\pgfqpoint{0.050000in}{0.083252in}}{\pgfqpoint{4.900000in}{4.900000in}}%
\pgfusepath{clip}%
\pgfsetbuttcap%
\pgfsetroundjoin%
\definecolor{currentfill}{rgb}{0.260100,0.229589,0.557586}%
\pgfsetfillcolor{currentfill}%
\pgfsetfillopacity{0.950000}%
\pgfsetlinewidth{0.200750pt}%
\definecolor{currentstroke}{rgb}{0.200000,0.200000,0.200000}%
\pgfsetstrokecolor{currentstroke}%
\pgfsetdash{}{0pt}%
\pgfpathmoveto{\pgfqpoint{3.581392in}{2.917747in}}%
\pgfpathlineto{\pgfqpoint{3.553633in}{2.990171in}}%
\pgfpathlineto{\pgfqpoint{3.525904in}{3.062687in}}%
\pgfpathlineto{\pgfqpoint{3.572243in}{3.076714in}}%
\pgfpathlineto{\pgfqpoint{3.599971in}{3.004294in}}%
\pgfpathlineto{\pgfqpoint{3.627729in}{2.931972in}}%
\pgfpathlineto{\pgfqpoint{3.581392in}{2.917747in}}%
\pgfpathclose%
\pgfusepath{stroke,fill}%
\end{pgfscope}%
\begin{pgfscope}%
\pgfpathrectangle{\pgfqpoint{0.050000in}{0.083252in}}{\pgfqpoint{4.900000in}{4.900000in}}%
\pgfusepath{clip}%
\pgfsetbuttcap%
\pgfsetroundjoin%
\definecolor{currentfill}{rgb}{0.261484,0.444444,0.859639}%
\pgfsetfillcolor{currentfill}%
\pgfsetfillopacity{0.950000}%
\pgfsetlinewidth{0.200750pt}%
\definecolor{currentstroke}{rgb}{0.200000,0.200000,0.200000}%
\pgfsetstrokecolor{currentstroke}%
\pgfsetdash{}{0pt}%
\pgfpathmoveto{\pgfqpoint{2.821385in}{3.826566in}}%
\pgfpathlineto{\pgfqpoint{2.793979in}{3.899251in}}%
\pgfpathlineto{\pgfqpoint{2.766599in}{3.971870in}}%
\pgfpathlineto{\pgfqpoint{2.813623in}{3.985010in}}%
\pgfpathlineto{\pgfqpoint{2.860565in}{3.998127in}}%
\pgfpathlineto{\pgfqpoint{2.887941in}{3.925648in}}%
\pgfpathlineto{\pgfqpoint{2.915343in}{3.853103in}}%
\pgfpathlineto{\pgfqpoint{2.868405in}{3.839846in}}%
\pgfpathlineto{\pgfqpoint{2.821385in}{3.826566in}}%
\pgfpathclose%
\pgfusepath{stroke,fill}%
\end{pgfscope}%
\begin{pgfscope}%
\pgfpathrectangle{\pgfqpoint{0.050000in}{0.083252in}}{\pgfqpoint{4.900000in}{4.900000in}}%
\pgfusepath{clip}%
\pgfsetbuttcap%
\pgfsetroundjoin%
\definecolor{currentfill}{rgb}{0.254840,0.373856,0.765798}%
\pgfsetfillcolor{currentfill}%
\pgfsetfillopacity{0.950000}%
\pgfsetlinewidth{0.200750pt}%
\definecolor{currentstroke}{rgb}{0.200000,0.200000,0.200000}%
\pgfsetstrokecolor{currentstroke}%
\pgfsetdash{}{0pt}%
\pgfpathmoveto{\pgfqpoint{3.025202in}{3.562271in}}%
\pgfpathlineto{\pgfqpoint{2.997699in}{3.635075in}}%
\pgfpathlineto{\pgfqpoint{2.970222in}{3.707815in}}%
\pgfpathlineto{\pgfqpoint{3.017073in}{3.721190in}}%
\pgfpathlineto{\pgfqpoint{3.063842in}{3.734542in}}%
\pgfpathlineto{\pgfqpoint{3.091315in}{3.661944in}}%
\pgfpathlineto{\pgfqpoint{3.118813in}{3.589284in}}%
\pgfpathlineto{\pgfqpoint{3.072049in}{3.575789in}}%
\pgfpathlineto{\pgfqpoint{3.025202in}{3.562271in}}%
\pgfpathclose%
\pgfusepath{stroke,fill}%
\end{pgfscope}%
\begin{pgfscope}%
\pgfpathrectangle{\pgfqpoint{0.050000in}{0.083252in}}{\pgfqpoint{4.900000in}{4.900000in}}%
\pgfusepath{clip}%
\pgfsetbuttcap%
\pgfsetroundjoin%
\definecolor{currentfill}{rgb}{0.248197,0.303268,0.671957}%
\pgfsetfillcolor{currentfill}%
\pgfsetfillopacity{0.950000}%
\pgfsetlinewidth{0.200750pt}%
\definecolor{currentstroke}{rgb}{0.200000,0.200000,0.200000}%
\pgfsetstrokecolor{currentstroke}%
\pgfsetdash{}{0pt}%
\pgfpathmoveto{\pgfqpoint{3.229064in}{3.298092in}}%
\pgfpathlineto{\pgfqpoint{3.201463in}{3.370964in}}%
\pgfpathlineto{\pgfqpoint{3.173887in}{3.443790in}}%
\pgfpathlineto{\pgfqpoint{3.220565in}{3.457409in}}%
\pgfpathlineto{\pgfqpoint{3.267161in}{3.471006in}}%
\pgfpathlineto{\pgfqpoint{3.294732in}{3.398331in}}%
\pgfpathlineto{\pgfqpoint{3.322329in}{3.325614in}}%
\pgfpathlineto{\pgfqpoint{3.275737in}{3.311864in}}%
\pgfpathlineto{\pgfqpoint{3.229064in}{3.298092in}}%
\pgfpathclose%
\pgfusepath{stroke,fill}%
\end{pgfscope}%
\begin{pgfscope}%
\pgfpathrectangle{\pgfqpoint{0.050000in}{0.083252in}}{\pgfqpoint{4.900000in}{4.900000in}}%
\pgfusepath{clip}%
\pgfsetbuttcap%
\pgfsetroundjoin%
\definecolor{currentfill}{rgb}{0.241799,0.235294,0.581592}%
\pgfsetfillcolor{currentfill}%
\pgfsetfillopacity{0.950000}%
\pgfsetlinewidth{0.200750pt}%
\definecolor{currentstroke}{rgb}{0.200000,0.200000,0.200000}%
\pgfsetstrokecolor{currentstroke}%
\pgfsetdash{}{0pt}%
\pgfpathmoveto{\pgfqpoint{3.432983in}{3.034571in}}%
\pgfpathlineto{\pgfqpoint{3.405278in}{3.107325in}}%
\pgfpathlineto{\pgfqpoint{3.377602in}{3.180096in}}%
\pgfpathlineto{\pgfqpoint{3.424107in}{3.193987in}}%
\pgfpathlineto{\pgfqpoint{3.470531in}{3.207857in}}%
\pgfpathlineto{\pgfqpoint{3.498204in}{3.135259in}}%
\pgfpathlineto{\pgfqpoint{3.525904in}{3.062687in}}%
\pgfpathlineto{\pgfqpoint{3.479484in}{3.048640in}}%
\pgfpathlineto{\pgfqpoint{3.432983in}{3.034571in}}%
\pgfpathclose%
\pgfusepath{stroke,fill}%
\end{pgfscope}%
\begin{pgfscope}%
\pgfpathrectangle{\pgfqpoint{0.050000in}{0.083252in}}{\pgfqpoint{4.900000in}{4.900000in}}%
\pgfusepath{clip}%
\pgfsetbuttcap%
\pgfsetroundjoin%
\definecolor{currentfill}{rgb}{0.264683,0.478431,0.904821}%
\pgfsetfillcolor{currentfill}%
\pgfsetfillopacity{0.950000}%
\pgfsetlinewidth{0.200750pt}%
\definecolor{currentstroke}{rgb}{0.200000,0.200000,0.200000}%
\pgfsetstrokecolor{currentstroke}%
\pgfsetdash{}{0pt}%
\pgfpathmoveto{\pgfqpoint{2.672302in}{3.945521in}}%
\pgfpathlineto{\pgfqpoint{2.644943in}{4.018213in}}%
\pgfpathlineto{\pgfqpoint{2.617609in}{4.090838in}}%
\pgfpathlineto{\pgfqpoint{2.664803in}{4.103884in}}%
\pgfpathlineto{\pgfqpoint{2.711914in}{4.116908in}}%
\pgfpathlineto{\pgfqpoint{2.739244in}{4.044423in}}%
\pgfpathlineto{\pgfqpoint{2.766599in}{3.971870in}}%
\pgfpathlineto{\pgfqpoint{2.719492in}{3.958707in}}%
\pgfpathlineto{\pgfqpoint{2.672302in}{3.945521in}}%
\pgfpathclose%
\pgfusepath{stroke,fill}%
\end{pgfscope}%
\begin{pgfscope}%
\pgfpathrectangle{\pgfqpoint{0.050000in}{0.083252in}}{\pgfqpoint{4.900000in}{4.900000in}}%
\pgfusepath{clip}%
\pgfsetbuttcap%
\pgfsetroundjoin%
\definecolor{currentfill}{rgb}{0.337670,0.310358,0.603968}%
\pgfsetfillcolor{currentfill}%
\pgfsetfillopacity{0.950000}%
\pgfsetlinewidth{0.200750pt}%
\definecolor{currentstroke}{rgb}{0.200000,0.200000,0.200000}%
\pgfsetstrokecolor{currentstroke}%
\pgfsetdash{}{0pt}%
\pgfpathmoveto{\pgfqpoint{3.637006in}{2.773368in}}%
\pgfpathlineto{\pgfqpoint{3.609183in}{2.845461in}}%
\pgfpathlineto{\pgfqpoint{3.581392in}{2.917747in}}%
\pgfpathlineto{\pgfqpoint{3.627729in}{2.931972in}}%
\pgfpathlineto{\pgfqpoint{3.655518in}{2.859794in}}%
\pgfpathlineto{\pgfqpoint{3.683340in}{2.787815in}}%
\pgfpathlineto{\pgfqpoint{3.637006in}{2.773368in}}%
\pgfpathclose%
\pgfusepath{stroke,fill}%
\end{pgfscope}%
\begin{pgfscope}%
\pgfpathrectangle{\pgfqpoint{0.050000in}{0.083252in}}{\pgfqpoint{4.900000in}{4.900000in}}%
\pgfusepath{clip}%
\pgfsetbuttcap%
\pgfsetroundjoin%
\definecolor{currentfill}{rgb}{0.258039,0.407843,0.810980}%
\pgfsetfillcolor{currentfill}%
\pgfsetfillopacity{0.950000}%
\pgfsetlinewidth{0.200750pt}%
\definecolor{currentstroke}{rgb}{0.200000,0.200000,0.200000}%
\pgfsetstrokecolor{currentstroke}%
\pgfsetdash{}{0pt}%
\pgfpathmoveto{\pgfqpoint{2.876273in}{3.680996in}}%
\pgfpathlineto{\pgfqpoint{2.848816in}{3.753814in}}%
\pgfpathlineto{\pgfqpoint{2.821385in}{3.826566in}}%
\pgfpathlineto{\pgfqpoint{2.868405in}{3.839846in}}%
\pgfpathlineto{\pgfqpoint{2.915343in}{3.853103in}}%
\pgfpathlineto{\pgfqpoint{2.942769in}{3.780492in}}%
\pgfpathlineto{\pgfqpoint{2.970222in}{3.707815in}}%
\pgfpathlineto{\pgfqpoint{2.923288in}{3.694417in}}%
\pgfpathlineto{\pgfqpoint{2.876273in}{3.680996in}}%
\pgfpathclose%
\pgfusepath{stroke,fill}%
\end{pgfscope}%
\begin{pgfscope}%
\pgfpathrectangle{\pgfqpoint{0.050000in}{0.083252in}}{\pgfqpoint{4.900000in}{4.900000in}}%
\pgfusepath{clip}%
\pgfsetbuttcap%
\pgfsetroundjoin%
\definecolor{currentfill}{rgb}{0.270342,0.538562,0.984760}%
\pgfsetfillcolor{currentfill}%
\pgfsetfillopacity{0.950000}%
\pgfsetlinewidth{0.200750pt}%
\definecolor{currentstroke}{rgb}{0.200000,0.200000,0.200000}%
\pgfsetstrokecolor{currentstroke}%
\pgfsetdash{}{0pt}%
\pgfpathmoveto{\pgfqpoint{2.468372in}{4.210003in}}%
\pgfpathlineto{\pgfqpoint{2.441111in}{4.282565in}}%
\pgfpathlineto{\pgfqpoint{2.488476in}{4.295449in}}%
\pgfpathlineto{\pgfqpoint{2.535759in}{4.308309in}}%
\pgfpathlineto{\pgfqpoint{2.563017in}{4.235886in}}%
\pgfpathlineto{\pgfqpoint{2.515736in}{4.222956in}}%
\pgfpathlineto{\pgfqpoint{2.468372in}{4.210003in}}%
\pgfpathclose%
\pgfusepath{stroke,fill}%
\end{pgfscope}%
\begin{pgfscope}%
\pgfpathrectangle{\pgfqpoint{0.050000in}{0.083252in}}{\pgfqpoint{4.900000in}{4.900000in}}%
\pgfusepath{clip}%
\pgfsetbuttcap%
\pgfsetroundjoin%
\definecolor{currentfill}{rgb}{0.251396,0.337255,0.717140}%
\pgfsetfillcolor{currentfill}%
\pgfsetfillopacity{0.950000}%
\pgfsetlinewidth{0.200750pt}%
\definecolor{currentstroke}{rgb}{0.200000,0.200000,0.200000}%
\pgfsetstrokecolor{currentstroke}%
\pgfsetdash{}{0pt}%
\pgfpathmoveto{\pgfqpoint{3.080286in}{3.416484in}}%
\pgfpathlineto{\pgfqpoint{3.052731in}{3.489407in}}%
\pgfpathlineto{\pgfqpoint{3.025202in}{3.562271in}}%
\pgfpathlineto{\pgfqpoint{3.072049in}{3.575789in}}%
\pgfpathlineto{\pgfqpoint{3.118813in}{3.589284in}}%
\pgfpathlineto{\pgfqpoint{3.146337in}{3.516565in}}%
\pgfpathlineto{\pgfqpoint{3.173887in}{3.443790in}}%
\pgfpathlineto{\pgfqpoint{3.127128in}{3.430149in}}%
\pgfpathlineto{\pgfqpoint{3.080286in}{3.416484in}}%
\pgfpathclose%
\pgfusepath{stroke,fill}%
\end{pgfscope}%
\begin{pgfscope}%
\pgfpathrectangle{\pgfqpoint{0.050000in}{0.083252in}}{\pgfqpoint{4.900000in}{4.900000in}}%
\pgfusepath{clip}%
\pgfsetbuttcap%
\pgfsetroundjoin%
\definecolor{currentfill}{rgb}{0.244998,0.269281,0.626774}%
\pgfsetfillcolor{currentfill}%
\pgfsetfillopacity{0.950000}%
\pgfsetlinewidth{0.200750pt}%
\definecolor{currentstroke}{rgb}{0.200000,0.200000,0.200000}%
\pgfsetstrokecolor{currentstroke}%
\pgfsetdash{}{0pt}%
\pgfpathmoveto{\pgfqpoint{3.284346in}{3.152249in}}%
\pgfpathlineto{\pgfqpoint{3.256692in}{3.225183in}}%
\pgfpathlineto{\pgfqpoint{3.229064in}{3.298092in}}%
\pgfpathlineto{\pgfqpoint{3.275737in}{3.311864in}}%
\pgfpathlineto{\pgfqpoint{3.322329in}{3.325614in}}%
\pgfpathlineto{\pgfqpoint{3.349952in}{3.252865in}}%
\pgfpathlineto{\pgfqpoint{3.377602in}{3.180096in}}%
\pgfpathlineto{\pgfqpoint{3.331015in}{3.166184in}}%
\pgfpathlineto{\pgfqpoint{3.284346in}{3.152249in}}%
\pgfpathclose%
\pgfusepath{stroke,fill}%
\end{pgfscope}%
\begin{pgfscope}%
\pgfpathrectangle{\pgfqpoint{0.050000in}{0.083252in}}{\pgfqpoint{4.900000in}{4.900000in}}%
\pgfusepath{clip}%
\pgfsetbuttcap%
\pgfsetroundjoin%
\definecolor{currentfill}{rgb}{0.260100,0.229589,0.557586}%
\pgfsetfillcolor{currentfill}%
\pgfsetfillopacity{0.950000}%
\pgfsetlinewidth{0.200750pt}%
\definecolor{currentstroke}{rgb}{0.200000,0.200000,0.200000}%
\pgfsetstrokecolor{currentstroke}%
\pgfsetdash{}{0pt}%
\pgfpathmoveto{\pgfqpoint{3.488477in}{2.889236in}}%
\pgfpathlineto{\pgfqpoint{3.460715in}{2.961863in}}%
\pgfpathlineto{\pgfqpoint{3.432983in}{3.034571in}}%
\pgfpathlineto{\pgfqpoint{3.479484in}{3.048640in}}%
\pgfpathlineto{\pgfqpoint{3.525904in}{3.062687in}}%
\pgfpathlineto{\pgfqpoint{3.553633in}{2.990171in}}%
\pgfpathlineto{\pgfqpoint{3.581392in}{2.917747in}}%
\pgfpathlineto{\pgfqpoint{3.534975in}{2.903502in}}%
\pgfpathlineto{\pgfqpoint{3.488477in}{2.889236in}}%
\pgfpathclose%
\pgfusepath{stroke,fill}%
\end{pgfscope}%
\begin{pgfscope}%
\pgfpathrectangle{\pgfqpoint{0.050000in}{0.083252in}}{\pgfqpoint{4.900000in}{4.900000in}}%
\pgfusepath{clip}%
\pgfsetbuttcap%
\pgfsetroundjoin%
\definecolor{currentfill}{rgb}{0.267882,0.512418,0.950004}%
\pgfsetfillcolor{currentfill}%
\pgfsetfillopacity{0.950000}%
\pgfsetlinewidth{0.200750pt}%
\definecolor{currentstroke}{rgb}{0.200000,0.200000,0.200000}%
\pgfsetstrokecolor{currentstroke}%
\pgfsetdash{}{0pt}%
\pgfpathmoveto{\pgfqpoint{2.522972in}{4.064676in}}%
\pgfpathlineto{\pgfqpoint{2.495660in}{4.137373in}}%
\pgfpathlineto{\pgfqpoint{2.468372in}{4.210003in}}%
\pgfpathlineto{\pgfqpoint{2.515736in}{4.222956in}}%
\pgfpathlineto{\pgfqpoint{2.563017in}{4.235886in}}%
\pgfpathlineto{\pgfqpoint{2.590300in}{4.163396in}}%
\pgfpathlineto{\pgfqpoint{2.617609in}{4.090838in}}%
\pgfpathlineto{\pgfqpoint{2.570332in}{4.077769in}}%
\pgfpathlineto{\pgfqpoint{2.522972in}{4.064676in}}%
\pgfpathclose%
\pgfusepath{stroke,fill}%
\end{pgfscope}%
\begin{pgfscope}%
\pgfpathrectangle{\pgfqpoint{0.050000in}{0.083252in}}{\pgfqpoint{4.900000in}{4.900000in}}%
\pgfusepath{clip}%
\pgfsetbuttcap%
\pgfsetroundjoin%
\definecolor{currentfill}{rgb}{0.261484,0.444444,0.859639}%
\pgfsetfillcolor{currentfill}%
\pgfsetfillopacity{0.950000}%
\pgfsetlinewidth{0.200750pt}%
\definecolor{currentstroke}{rgb}{0.200000,0.200000,0.200000}%
\pgfsetstrokecolor{currentstroke}%
\pgfsetdash{}{0pt}%
\pgfpathmoveto{\pgfqpoint{2.727097in}{3.799936in}}%
\pgfpathlineto{\pgfqpoint{2.699686in}{3.872762in}}%
\pgfpathlineto{\pgfqpoint{2.672302in}{3.945521in}}%
\pgfpathlineto{\pgfqpoint{2.719492in}{3.958707in}}%
\pgfpathlineto{\pgfqpoint{2.766599in}{3.971870in}}%
\pgfpathlineto{\pgfqpoint{2.793979in}{3.899251in}}%
\pgfpathlineto{\pgfqpoint{2.821385in}{3.826566in}}%
\pgfpathlineto{\pgfqpoint{2.774282in}{3.813263in}}%
\pgfpathlineto{\pgfqpoint{2.727097in}{3.799936in}}%
\pgfpathclose%
\pgfusepath{stroke,fill}%
\end{pgfscope}%
\begin{pgfscope}%
\pgfpathrectangle{\pgfqpoint{0.050000in}{0.083252in}}{\pgfqpoint{4.900000in}{4.900000in}}%
\pgfusepath{clip}%
\pgfsetbuttcap%
\pgfsetroundjoin%
\definecolor{currentfill}{rgb}{0.254840,0.373856,0.765798}%
\pgfsetfillcolor{currentfill}%
\pgfsetfillopacity{0.950000}%
\pgfsetlinewidth{0.200750pt}%
\definecolor{currentstroke}{rgb}{0.200000,0.200000,0.200000}%
\pgfsetstrokecolor{currentstroke}%
\pgfsetdash{}{0pt}%
\pgfpathmoveto{\pgfqpoint{2.931263in}{3.535165in}}%
\pgfpathlineto{\pgfqpoint{2.903755in}{3.608113in}}%
\pgfpathlineto{\pgfqpoint{2.876273in}{3.680996in}}%
\pgfpathlineto{\pgfqpoint{2.923288in}{3.694417in}}%
\pgfpathlineto{\pgfqpoint{2.970222in}{3.707815in}}%
\pgfpathlineto{\pgfqpoint{2.997699in}{3.635075in}}%
\pgfpathlineto{\pgfqpoint{3.025202in}{3.562271in}}%
\pgfpathlineto{\pgfqpoint{2.978274in}{3.548730in}}%
\pgfpathlineto{\pgfqpoint{2.931263in}{3.535165in}}%
\pgfpathclose%
\pgfusepath{stroke,fill}%
\end{pgfscope}%
\begin{pgfscope}%
\pgfpathrectangle{\pgfqpoint{0.050000in}{0.083252in}}{\pgfqpoint{4.900000in}{4.900000in}}%
\pgfusepath{clip}%
\pgfsetbuttcap%
\pgfsetroundjoin%
\definecolor{currentfill}{rgb}{0.415240,0.391126,0.650350}%
\pgfsetfillcolor{currentfill}%
\pgfsetfillopacity{0.950000}%
\pgfsetlinewidth{0.200750pt}%
\definecolor{currentstroke}{rgb}{0.200000,0.200000,0.200000}%
\pgfsetstrokecolor{currentstroke}%
\pgfsetdash{}{0pt}%
\pgfpathmoveto{\pgfqpoint{3.692760in}{2.630051in}}%
\pgfpathlineto{\pgfqpoint{3.664864in}{2.701537in}}%
\pgfpathlineto{\pgfqpoint{3.637006in}{2.773368in}}%
\pgfpathlineto{\pgfqpoint{3.683340in}{2.787815in}}%
\pgfpathlineto{\pgfqpoint{3.711199in}{2.716106in}}%
\pgfpathlineto{\pgfqpoint{3.739095in}{2.644749in}}%
\pgfpathlineto{\pgfqpoint{3.692760in}{2.630051in}}%
\pgfpathclose%
\pgfusepath{stroke,fill}%
\end{pgfscope}%
\begin{pgfscope}%
\pgfpathrectangle{\pgfqpoint{0.050000in}{0.083252in}}{\pgfqpoint{4.900000in}{4.900000in}}%
\pgfusepath{clip}%
\pgfsetbuttcap%
\pgfsetroundjoin%
\definecolor{currentfill}{rgb}{0.248197,0.303268,0.671957}%
\pgfsetfillcolor{currentfill}%
\pgfsetfillopacity{0.950000}%
\pgfsetlinewidth{0.200750pt}%
\definecolor{currentstroke}{rgb}{0.200000,0.200000,0.200000}%
\pgfsetstrokecolor{currentstroke}%
\pgfsetdash{}{0pt}%
\pgfpathmoveto{\pgfqpoint{3.135472in}{3.270481in}}%
\pgfpathlineto{\pgfqpoint{3.107866in}{3.343507in}}%
\pgfpathlineto{\pgfqpoint{3.080286in}{3.416484in}}%
\pgfpathlineto{\pgfqpoint{3.127128in}{3.430149in}}%
\pgfpathlineto{\pgfqpoint{3.173887in}{3.443790in}}%
\pgfpathlineto{\pgfqpoint{3.201463in}{3.370964in}}%
\pgfpathlineto{\pgfqpoint{3.229064in}{3.298092in}}%
\pgfpathlineto{\pgfqpoint{3.182309in}{3.284298in}}%
\pgfpathlineto{\pgfqpoint{3.135472in}{3.270481in}}%
\pgfpathclose%
\pgfusepath{stroke,fill}%
\end{pgfscope}%
\begin{pgfscope}%
\pgfpathrectangle{\pgfqpoint{0.050000in}{0.083252in}}{\pgfqpoint{4.900000in}{4.900000in}}%
\pgfusepath{clip}%
\pgfsetbuttcap%
\pgfsetroundjoin%
\definecolor{currentfill}{rgb}{0.241553,0.232680,0.578116}%
\pgfsetfillcolor{currentfill}%
\pgfsetfillopacity{0.950000}%
\pgfsetlinewidth{0.200750pt}%
\definecolor{currentstroke}{rgb}{0.200000,0.200000,0.200000}%
\pgfsetstrokecolor{currentstroke}%
\pgfsetdash{}{0pt}%
\pgfpathmoveto{\pgfqpoint{3.339735in}{3.006371in}}%
\pgfpathlineto{\pgfqpoint{3.312027in}{3.079305in}}%
\pgfpathlineto{\pgfqpoint{3.284346in}{3.152249in}}%
\pgfpathlineto{\pgfqpoint{3.331015in}{3.166184in}}%
\pgfpathlineto{\pgfqpoint{3.377602in}{3.180096in}}%
\pgfpathlineto{\pgfqpoint{3.405278in}{3.107325in}}%
\pgfpathlineto{\pgfqpoint{3.432983in}{3.034571in}}%
\pgfpathlineto{\pgfqpoint{3.386400in}{3.020481in}}%
\pgfpathlineto{\pgfqpoint{3.339735in}{3.006371in}}%
\pgfpathclose%
\pgfusepath{stroke,fill}%
\end{pgfscope}%
\begin{pgfscope}%
\pgfpathrectangle{\pgfqpoint{0.050000in}{0.083252in}}{\pgfqpoint{4.900000in}{4.900000in}}%
\pgfusepath{clip}%
\pgfsetbuttcap%
\pgfsetroundjoin%
\definecolor{currentfill}{rgb}{0.337670,0.310358,0.603968}%
\pgfsetfillcolor{currentfill}%
\pgfsetfillopacity{0.950000}%
\pgfsetlinewidth{0.200750pt}%
\definecolor{currentstroke}{rgb}{0.200000,0.200000,0.200000}%
\pgfsetstrokecolor{currentstroke}%
\pgfsetdash{}{0pt}%
\pgfpathmoveto{\pgfqpoint{3.544094in}{2.744412in}}%
\pgfpathlineto{\pgfqpoint{3.516270in}{2.816734in}}%
\pgfpathlineto{\pgfqpoint{3.488477in}{2.889236in}}%
\pgfpathlineto{\pgfqpoint{3.534975in}{2.903502in}}%
\pgfpathlineto{\pgfqpoint{3.581392in}{2.917747in}}%
\pgfpathlineto{\pgfqpoint{3.609183in}{2.845461in}}%
\pgfpathlineto{\pgfqpoint{3.637006in}{2.773368in}}%
\pgfpathlineto{\pgfqpoint{3.590591in}{2.758900in}}%
\pgfpathlineto{\pgfqpoint{3.544094in}{2.744412in}}%
\pgfpathclose%
\pgfusepath{stroke,fill}%
\end{pgfscope}%
\begin{pgfscope}%
\pgfpathrectangle{\pgfqpoint{0.050000in}{0.083252in}}{\pgfqpoint{4.900000in}{4.900000in}}%
\pgfusepath{clip}%
\pgfsetbuttcap%
\pgfsetroundjoin%
\definecolor{currentfill}{rgb}{0.264683,0.478431,0.904821}%
\pgfsetfillcolor{currentfill}%
\pgfsetfillopacity{0.950000}%
\pgfsetlinewidth{0.200750pt}%
\definecolor{currentstroke}{rgb}{0.200000,0.200000,0.200000}%
\pgfsetstrokecolor{currentstroke}%
\pgfsetdash{}{0pt}%
\pgfpathmoveto{\pgfqpoint{2.577673in}{3.919079in}}%
\pgfpathlineto{\pgfqpoint{2.550310in}{3.991911in}}%
\pgfpathlineto{\pgfqpoint{2.522972in}{4.064676in}}%
\pgfpathlineto{\pgfqpoint{2.570332in}{4.077769in}}%
\pgfpathlineto{\pgfqpoint{2.617609in}{4.090838in}}%
\pgfpathlineto{\pgfqpoint{2.644943in}{4.018213in}}%
\pgfpathlineto{\pgfqpoint{2.672302in}{3.945521in}}%
\pgfpathlineto{\pgfqpoint{2.625029in}{3.932312in}}%
\pgfpathlineto{\pgfqpoint{2.577673in}{3.919079in}}%
\pgfpathclose%
\pgfusepath{stroke,fill}%
\end{pgfscope}%
\begin{pgfscope}%
\pgfpathrectangle{\pgfqpoint{0.050000in}{0.083252in}}{\pgfqpoint{4.900000in}{4.900000in}}%
\pgfusepath{clip}%
\pgfsetbuttcap%
\pgfsetroundjoin%
\definecolor{currentfill}{rgb}{0.258039,0.407843,0.810980}%
\pgfsetfillcolor{currentfill}%
\pgfsetfillopacity{0.950000}%
\pgfsetlinewidth{0.200750pt}%
\definecolor{currentstroke}{rgb}{0.200000,0.200000,0.200000}%
\pgfsetstrokecolor{currentstroke}%
\pgfsetdash{}{0pt}%
\pgfpathmoveto{\pgfqpoint{2.781993in}{3.654083in}}%
\pgfpathlineto{\pgfqpoint{2.754532in}{3.727043in}}%
\pgfpathlineto{\pgfqpoint{2.727097in}{3.799936in}}%
\pgfpathlineto{\pgfqpoint{2.774282in}{3.813263in}}%
\pgfpathlineto{\pgfqpoint{2.821385in}{3.826566in}}%
\pgfpathlineto{\pgfqpoint{2.848816in}{3.753814in}}%
\pgfpathlineto{\pgfqpoint{2.876273in}{3.680996in}}%
\pgfpathlineto{\pgfqpoint{2.829174in}{3.667551in}}%
\pgfpathlineto{\pgfqpoint{2.781993in}{3.654083in}}%
\pgfpathclose%
\pgfusepath{stroke,fill}%
\end{pgfscope}%
\begin{pgfscope}%
\pgfpathrectangle{\pgfqpoint{0.050000in}{0.083252in}}{\pgfqpoint{4.900000in}{4.900000in}}%
\pgfusepath{clip}%
\pgfsetbuttcap%
\pgfsetroundjoin%
\definecolor{currentfill}{rgb}{0.270342,0.538562,0.984760}%
\pgfsetfillcolor{currentfill}%
\pgfsetfillopacity{0.950000}%
\pgfsetlinewidth{0.200750pt}%
\definecolor{currentstroke}{rgb}{0.200000,0.200000,0.200000}%
\pgfsetstrokecolor{currentstroke}%
\pgfsetdash{}{0pt}%
\pgfpathmoveto{\pgfqpoint{2.373395in}{4.184029in}}%
\pgfpathlineto{\pgfqpoint{2.346129in}{4.256731in}}%
\pgfpathlineto{\pgfqpoint{2.393662in}{4.269660in}}%
\pgfpathlineto{\pgfqpoint{2.441111in}{4.282565in}}%
\pgfpathlineto{\pgfqpoint{2.468372in}{4.210003in}}%
\pgfpathlineto{\pgfqpoint{2.420925in}{4.197027in}}%
\pgfpathlineto{\pgfqpoint{2.373395in}{4.184029in}}%
\pgfpathclose%
\pgfusepath{stroke,fill}%
\end{pgfscope}%
\begin{pgfscope}%
\pgfpathrectangle{\pgfqpoint{0.050000in}{0.083252in}}{\pgfqpoint{4.900000in}{4.900000in}}%
\pgfusepath{clip}%
\pgfsetbuttcap%
\pgfsetroundjoin%
\definecolor{currentfill}{rgb}{0.251396,0.337255,0.717140}%
\pgfsetfillcolor{currentfill}%
\pgfsetfillopacity{0.950000}%
\pgfsetlinewidth{0.200750pt}%
\definecolor{currentstroke}{rgb}{0.200000,0.200000,0.200000}%
\pgfsetstrokecolor{currentstroke}%
\pgfsetdash{}{0pt}%
\pgfpathmoveto{\pgfqpoint{2.986355in}{3.389085in}}%
\pgfpathlineto{\pgfqpoint{2.958796in}{3.462155in}}%
\pgfpathlineto{\pgfqpoint{2.931263in}{3.535165in}}%
\pgfpathlineto{\pgfqpoint{2.978274in}{3.548730in}}%
\pgfpathlineto{\pgfqpoint{3.025202in}{3.562271in}}%
\pgfpathlineto{\pgfqpoint{3.052731in}{3.489407in}}%
\pgfpathlineto{\pgfqpoint{3.080286in}{3.416484in}}%
\pgfpathlineto{\pgfqpoint{3.033362in}{3.402796in}}%
\pgfpathlineto{\pgfqpoint{2.986355in}{3.389085in}}%
\pgfpathclose%
\pgfusepath{stroke,fill}%
\end{pgfscope}%
\begin{pgfscope}%
\pgfpathrectangle{\pgfqpoint{0.050000in}{0.083252in}}{\pgfqpoint{4.900000in}{4.900000in}}%
\pgfusepath{clip}%
\pgfsetbuttcap%
\pgfsetroundjoin%
\definecolor{currentfill}{rgb}{0.244998,0.269281,0.626774}%
\pgfsetfillcolor{currentfill}%
\pgfsetfillopacity{0.950000}%
\pgfsetlinewidth{0.200750pt}%
\definecolor{currentstroke}{rgb}{0.200000,0.200000,0.200000}%
\pgfsetstrokecolor{currentstroke}%
\pgfsetdash{}{0pt}%
\pgfpathmoveto{\pgfqpoint{3.190763in}{3.124315in}}%
\pgfpathlineto{\pgfqpoint{3.163105in}{3.197413in}}%
\pgfpathlineto{\pgfqpoint{3.135472in}{3.270481in}}%
\pgfpathlineto{\pgfqpoint{3.182309in}{3.284298in}}%
\pgfpathlineto{\pgfqpoint{3.229064in}{3.298092in}}%
\pgfpathlineto{\pgfqpoint{3.256692in}{3.225183in}}%
\pgfpathlineto{\pgfqpoint{3.284346in}{3.152249in}}%
\pgfpathlineto{\pgfqpoint{3.237596in}{3.138293in}}%
\pgfpathlineto{\pgfqpoint{3.190763in}{3.124315in}}%
\pgfpathclose%
\pgfusepath{stroke,fill}%
\end{pgfscope}%
\begin{pgfscope}%
\pgfpathrectangle{\pgfqpoint{0.050000in}{0.083252in}}{\pgfqpoint{4.900000in}{4.900000in}}%
\pgfusepath{clip}%
\pgfsetbuttcap%
\pgfsetroundjoin%
\definecolor{currentfill}{rgb}{0.492810,0.471895,0.696732}%
\pgfsetfillcolor{currentfill}%
\pgfsetfillopacity{0.950000}%
\pgfsetlinewidth{0.200750pt}%
\definecolor{currentstroke}{rgb}{0.200000,0.200000,0.200000}%
\pgfsetstrokecolor{currentstroke}%
\pgfsetdash{}{0pt}%
\pgfpathmoveto{\pgfqpoint{3.748680in}{2.488519in}}%
\pgfpathlineto{\pgfqpoint{3.720698in}{2.559007in}}%
\pgfpathlineto{\pgfqpoint{3.692760in}{2.630051in}}%
\pgfpathlineto{\pgfqpoint{3.739095in}{2.644749in}}%
\pgfpathlineto{\pgfqpoint{3.767034in}{2.573841in}}%
\pgfpathlineto{\pgfqpoint{3.795019in}{2.503494in}}%
\pgfpathlineto{\pgfqpoint{3.748680in}{2.488519in}}%
\pgfpathclose%
\pgfusepath{stroke,fill}%
\end{pgfscope}%
\begin{pgfscope}%
\pgfpathrectangle{\pgfqpoint{0.050000in}{0.083252in}}{\pgfqpoint{4.900000in}{4.900000in}}%
\pgfusepath{clip}%
\pgfsetbuttcap%
\pgfsetroundjoin%
\definecolor{currentfill}{rgb}{0.260100,0.229589,0.557586}%
\pgfsetfillcolor{currentfill}%
\pgfsetfillopacity{0.950000}%
\pgfsetlinewidth{0.200750pt}%
\definecolor{currentstroke}{rgb}{0.200000,0.200000,0.200000}%
\pgfsetstrokecolor{currentstroke}%
\pgfsetdash{}{0pt}%
\pgfpathmoveto{\pgfqpoint{3.395237in}{2.860642in}}%
\pgfpathlineto{\pgfqpoint{3.367472in}{2.933472in}}%
\pgfpathlineto{\pgfqpoint{3.339735in}{3.006371in}}%
\pgfpathlineto{\pgfqpoint{3.386400in}{3.020481in}}%
\pgfpathlineto{\pgfqpoint{3.432983in}{3.034571in}}%
\pgfpathlineto{\pgfqpoint{3.460715in}{2.961863in}}%
\pgfpathlineto{\pgfqpoint{3.488477in}{2.889236in}}%
\pgfpathlineto{\pgfqpoint{3.441898in}{2.874949in}}%
\pgfpathlineto{\pgfqpoint{3.395237in}{2.860642in}}%
\pgfpathclose%
\pgfusepath{stroke,fill}%
\end{pgfscope}%
\begin{pgfscope}%
\pgfpathrectangle{\pgfqpoint{0.050000in}{0.083252in}}{\pgfqpoint{4.900000in}{4.900000in}}%
\pgfusepath{clip}%
\pgfsetbuttcap%
\pgfsetroundjoin%
\definecolor{currentfill}{rgb}{0.415240,0.391126,0.650350}%
\pgfsetfillcolor{currentfill}%
\pgfsetfillopacity{0.950000}%
\pgfsetlinewidth{0.200750pt}%
\definecolor{currentstroke}{rgb}{0.200000,0.200000,0.200000}%
\pgfsetstrokecolor{currentstroke}%
\pgfsetdash{}{0pt}%
\pgfpathmoveto{\pgfqpoint{3.599848in}{2.600592in}}%
\pgfpathlineto{\pgfqpoint{3.571953in}{2.672337in}}%
\pgfpathlineto{\pgfqpoint{3.544094in}{2.744412in}}%
\pgfpathlineto{\pgfqpoint{3.590591in}{2.758900in}}%
\pgfpathlineto{\pgfqpoint{3.637006in}{2.773368in}}%
\pgfpathlineto{\pgfqpoint{3.664864in}{2.701537in}}%
\pgfpathlineto{\pgfqpoint{3.692760in}{2.630051in}}%
\pgfpathlineto{\pgfqpoint{3.646345in}{2.615332in}}%
\pgfpathlineto{\pgfqpoint{3.599848in}{2.600592in}}%
\pgfpathclose%
\pgfusepath{stroke,fill}%
\end{pgfscope}%
\begin{pgfscope}%
\pgfpathrectangle{\pgfqpoint{0.050000in}{0.083252in}}{\pgfqpoint{4.900000in}{4.900000in}}%
\pgfusepath{clip}%
\pgfsetbuttcap%
\pgfsetroundjoin%
\definecolor{currentfill}{rgb}{0.267882,0.512418,0.950004}%
\pgfsetfillcolor{currentfill}%
\pgfsetfillopacity{0.950000}%
\pgfsetlinewidth{0.200750pt}%
\definecolor{currentstroke}{rgb}{0.200000,0.200000,0.200000}%
\pgfsetstrokecolor{currentstroke}%
\pgfsetdash{}{0pt}%
\pgfpathmoveto{\pgfqpoint{2.428002in}{4.038422in}}%
\pgfpathlineto{\pgfqpoint{2.400686in}{4.111259in}}%
\pgfpathlineto{\pgfqpoint{2.373395in}{4.184029in}}%
\pgfpathlineto{\pgfqpoint{2.420925in}{4.197027in}}%
\pgfpathlineto{\pgfqpoint{2.468372in}{4.210003in}}%
\pgfpathlineto{\pgfqpoint{2.495660in}{4.137373in}}%
\pgfpathlineto{\pgfqpoint{2.522972in}{4.064676in}}%
\pgfpathlineto{\pgfqpoint{2.475529in}{4.051561in}}%
\pgfpathlineto{\pgfqpoint{2.428002in}{4.038422in}}%
\pgfpathclose%
\pgfusepath{stroke,fill}%
\end{pgfscope}%
\begin{pgfscope}%
\pgfpathrectangle{\pgfqpoint{0.050000in}{0.083252in}}{\pgfqpoint{4.900000in}{4.900000in}}%
\pgfusepath{clip}%
\pgfsetbuttcap%
\pgfsetroundjoin%
\definecolor{currentfill}{rgb}{0.261238,0.441830,0.856163}%
\pgfsetfillcolor{currentfill}%
\pgfsetfillopacity{0.950000}%
\pgfsetlinewidth{0.200750pt}%
\definecolor{currentstroke}{rgb}{0.200000,0.200000,0.200000}%
\pgfsetstrokecolor{currentstroke}%
\pgfsetdash{}{0pt}%
\pgfpathmoveto{\pgfqpoint{2.632476in}{3.773212in}}%
\pgfpathlineto{\pgfqpoint{2.605062in}{3.846179in}}%
\pgfpathlineto{\pgfqpoint{2.577673in}{3.919079in}}%
\pgfpathlineto{\pgfqpoint{2.625029in}{3.932312in}}%
\pgfpathlineto{\pgfqpoint{2.672302in}{3.945521in}}%
\pgfpathlineto{\pgfqpoint{2.699686in}{3.872762in}}%
\pgfpathlineto{\pgfqpoint{2.727097in}{3.799936in}}%
\pgfpathlineto{\pgfqpoint{2.679828in}{3.786586in}}%
\pgfpathlineto{\pgfqpoint{2.632476in}{3.773212in}}%
\pgfpathclose%
\pgfusepath{stroke,fill}%
\end{pgfscope}%
\begin{pgfscope}%
\pgfpathrectangle{\pgfqpoint{0.050000in}{0.083252in}}{\pgfqpoint{4.900000in}{4.900000in}}%
\pgfusepath{clip}%
\pgfsetbuttcap%
\pgfsetroundjoin%
\definecolor{currentfill}{rgb}{0.254840,0.373856,0.765798}%
\pgfsetfillcolor{currentfill}%
\pgfsetfillopacity{0.950000}%
\pgfsetlinewidth{0.200750pt}%
\definecolor{currentstroke}{rgb}{0.200000,0.200000,0.200000}%
\pgfsetstrokecolor{currentstroke}%
\pgfsetdash{}{0pt}%
\pgfpathmoveto{\pgfqpoint{2.836992in}{3.507966in}}%
\pgfpathlineto{\pgfqpoint{2.809480in}{3.581057in}}%
\pgfpathlineto{\pgfqpoint{2.781993in}{3.654083in}}%
\pgfpathlineto{\pgfqpoint{2.829174in}{3.667551in}}%
\pgfpathlineto{\pgfqpoint{2.876273in}{3.680996in}}%
\pgfpathlineto{\pgfqpoint{2.903755in}{3.608113in}}%
\pgfpathlineto{\pgfqpoint{2.931263in}{3.535165in}}%
\pgfpathlineto{\pgfqpoint{2.884169in}{3.521577in}}%
\pgfpathlineto{\pgfqpoint{2.836992in}{3.507966in}}%
\pgfpathclose%
\pgfusepath{stroke,fill}%
\end{pgfscope}%
\begin{pgfscope}%
\pgfpathrectangle{\pgfqpoint{0.050000in}{0.083252in}}{\pgfqpoint{4.900000in}{4.900000in}}%
\pgfusepath{clip}%
\pgfsetbuttcap%
\pgfsetroundjoin%
\definecolor{currentfill}{rgb}{0.248197,0.303268,0.671957}%
\pgfsetfillcolor{currentfill}%
\pgfsetfillopacity{0.950000}%
\pgfsetlinewidth{0.200750pt}%
\definecolor{currentstroke}{rgb}{0.200000,0.200000,0.200000}%
\pgfsetstrokecolor{currentstroke}%
\pgfsetdash{}{0pt}%
\pgfpathmoveto{\pgfqpoint{3.041551in}{3.242779in}}%
\pgfpathlineto{\pgfqpoint{3.013940in}{3.315958in}}%
\pgfpathlineto{\pgfqpoint{2.986355in}{3.389085in}}%
\pgfpathlineto{\pgfqpoint{3.033362in}{3.402796in}}%
\pgfpathlineto{\pgfqpoint{3.080286in}{3.416484in}}%
\pgfpathlineto{\pgfqpoint{3.107866in}{3.343507in}}%
\pgfpathlineto{\pgfqpoint{3.135472in}{3.270481in}}%
\pgfpathlineto{\pgfqpoint{3.088553in}{3.256641in}}%
\pgfpathlineto{\pgfqpoint{3.041551in}{3.242779in}}%
\pgfpathclose%
\pgfusepath{stroke,fill}%
\end{pgfscope}%
\begin{pgfscope}%
\pgfpathrectangle{\pgfqpoint{0.050000in}{0.083252in}}{\pgfqpoint{4.900000in}{4.900000in}}%
\pgfusepath{clip}%
\pgfsetbuttcap%
\pgfsetroundjoin%
\definecolor{currentfill}{rgb}{0.241553,0.232680,0.578116}%
\pgfsetfillcolor{currentfill}%
\pgfsetfillopacity{0.950000}%
\pgfsetlinewidth{0.200750pt}%
\definecolor{currentstroke}{rgb}{0.200000,0.200000,0.200000}%
\pgfsetstrokecolor{currentstroke}%
\pgfsetdash{}{0pt}%
\pgfpathmoveto{\pgfqpoint{3.246161in}{2.978086in}}%
\pgfpathlineto{\pgfqpoint{3.218449in}{3.051200in}}%
\pgfpathlineto{\pgfqpoint{3.190763in}{3.124315in}}%
\pgfpathlineto{\pgfqpoint{3.237596in}{3.138293in}}%
\pgfpathlineto{\pgfqpoint{3.284346in}{3.152249in}}%
\pgfpathlineto{\pgfqpoint{3.312027in}{3.079305in}}%
\pgfpathlineto{\pgfqpoint{3.339735in}{3.006371in}}%
\pgfpathlineto{\pgfqpoint{3.292989in}{2.992239in}}%
\pgfpathlineto{\pgfqpoint{3.246161in}{2.978086in}}%
\pgfpathclose%
\pgfusepath{stroke,fill}%
\end{pgfscope}%
\begin{pgfscope}%
\pgfpathrectangle{\pgfqpoint{0.050000in}{0.083252in}}{\pgfqpoint{4.900000in}{4.900000in}}%
\pgfusepath{clip}%
\pgfsetbuttcap%
\pgfsetroundjoin%
\definecolor{currentfill}{rgb}{0.337670,0.310358,0.603968}%
\pgfsetfillcolor{currentfill}%
\pgfsetfillopacity{0.950000}%
\pgfsetlinewidth{0.200750pt}%
\definecolor{currentstroke}{rgb}{0.200000,0.200000,0.200000}%
\pgfsetstrokecolor{currentstroke}%
\pgfsetdash{}{0pt}%
\pgfpathmoveto{\pgfqpoint{3.450858in}{2.715374in}}%
\pgfpathlineto{\pgfqpoint{3.423031in}{2.787925in}}%
\pgfpathlineto{\pgfqpoint{3.395237in}{2.860642in}}%
\pgfpathlineto{\pgfqpoint{3.441898in}{2.874949in}}%
\pgfpathlineto{\pgfqpoint{3.488477in}{2.889236in}}%
\pgfpathlineto{\pgfqpoint{3.516270in}{2.816734in}}%
\pgfpathlineto{\pgfqpoint{3.544094in}{2.744412in}}%
\pgfpathlineto{\pgfqpoint{3.497517in}{2.729903in}}%
\pgfpathlineto{\pgfqpoint{3.450858in}{2.715374in}}%
\pgfpathclose%
\pgfusepath{stroke,fill}%
\end{pgfscope}%
\begin{pgfscope}%
\pgfpathrectangle{\pgfqpoint{0.050000in}{0.083252in}}{\pgfqpoint{4.900000in}{4.900000in}}%
\pgfusepath{clip}%
\pgfsetbuttcap%
\pgfsetroundjoin%
\definecolor{currentfill}{rgb}{0.564414,0.546451,0.739546}%
\pgfsetfillcolor{currentfill}%
\pgfsetfillopacity{0.950000}%
\pgfsetlinewidth{0.200750pt}%
\definecolor{currentstroke}{rgb}{0.200000,0.200000,0.200000}%
\pgfsetstrokecolor{currentstroke}%
\pgfsetdash{}{0pt}%
\pgfpathmoveto{\pgfqpoint{3.804800in}{2.349751in}}%
\pgfpathlineto{\pgfqpoint{3.776713in}{2.418718in}}%
\pgfpathlineto{\pgfqpoint{3.748680in}{2.488519in}}%
\pgfpathlineto{\pgfqpoint{3.795019in}{2.503494in}}%
\pgfpathlineto{\pgfqpoint{3.823054in}{2.433840in}}%
\pgfpathlineto{\pgfqpoint{3.851146in}{2.365024in}}%
\pgfpathlineto{\pgfqpoint{3.804800in}{2.349751in}}%
\pgfpathclose%
\pgfusepath{stroke,fill}%
\end{pgfscope}%
\begin{pgfscope}%
\pgfpathrectangle{\pgfqpoint{0.050000in}{0.083252in}}{\pgfqpoint{4.900000in}{4.900000in}}%
\pgfusepath{clip}%
\pgfsetbuttcap%
\pgfsetroundjoin%
\definecolor{currentfill}{rgb}{0.492810,0.471895,0.696732}%
\pgfsetfillcolor{currentfill}%
\pgfsetfillopacity{0.950000}%
\pgfsetlinewidth{0.200750pt}%
\definecolor{currentstroke}{rgb}{0.200000,0.200000,0.200000}%
\pgfsetstrokecolor{currentstroke}%
\pgfsetdash{}{0pt}%
\pgfpathmoveto{\pgfqpoint{3.655761in}{2.458499in}}%
\pgfpathlineto{\pgfqpoint{3.627783in}{2.529274in}}%
\pgfpathlineto{\pgfqpoint{3.599848in}{2.600592in}}%
\pgfpathlineto{\pgfqpoint{3.646345in}{2.615332in}}%
\pgfpathlineto{\pgfqpoint{3.692760in}{2.630051in}}%
\pgfpathlineto{\pgfqpoint{3.720698in}{2.559007in}}%
\pgfpathlineto{\pgfqpoint{3.748680in}{2.488519in}}%
\pgfpathlineto{\pgfqpoint{3.702261in}{2.473521in}}%
\pgfpathlineto{\pgfqpoint{3.655761in}{2.458499in}}%
\pgfpathclose%
\pgfusepath{stroke,fill}%
\end{pgfscope}%
\begin{pgfscope}%
\pgfpathrectangle{\pgfqpoint{0.050000in}{0.083252in}}{\pgfqpoint{4.900000in}{4.900000in}}%
\pgfusepath{clip}%
\pgfsetbuttcap%
\pgfsetroundjoin%
\definecolor{currentfill}{rgb}{0.264683,0.478431,0.904821}%
\pgfsetfillcolor{currentfill}%
\pgfsetfillopacity{0.950000}%
\pgfsetlinewidth{0.200750pt}%
\definecolor{currentstroke}{rgb}{0.200000,0.200000,0.200000}%
\pgfsetstrokecolor{currentstroke}%
\pgfsetdash{}{0pt}%
\pgfpathmoveto{\pgfqpoint{2.482711in}{3.892544in}}%
\pgfpathlineto{\pgfqpoint{2.455344in}{3.965517in}}%
\pgfpathlineto{\pgfqpoint{2.428002in}{4.038422in}}%
\pgfpathlineto{\pgfqpoint{2.475529in}{4.051561in}}%
\pgfpathlineto{\pgfqpoint{2.522972in}{4.064676in}}%
\pgfpathlineto{\pgfqpoint{2.550310in}{3.991911in}}%
\pgfpathlineto{\pgfqpoint{2.577673in}{3.919079in}}%
\pgfpathlineto{\pgfqpoint{2.530234in}{3.905823in}}%
\pgfpathlineto{\pgfqpoint{2.482711in}{3.892544in}}%
\pgfpathclose%
\pgfusepath{stroke,fill}%
\end{pgfscope}%
\begin{pgfscope}%
\pgfpathrectangle{\pgfqpoint{0.050000in}{0.083252in}}{\pgfqpoint{4.900000in}{4.900000in}}%
\pgfusepath{clip}%
\pgfsetbuttcap%
\pgfsetroundjoin%
\definecolor{currentfill}{rgb}{0.258039,0.407843,0.810980}%
\pgfsetfillcolor{currentfill}%
\pgfsetfillopacity{0.950000}%
\pgfsetlinewidth{0.200750pt}%
\definecolor{currentstroke}{rgb}{0.200000,0.200000,0.200000}%
\pgfsetstrokecolor{currentstroke}%
\pgfsetdash{}{0pt}%
\pgfpathmoveto{\pgfqpoint{2.687382in}{3.627076in}}%
\pgfpathlineto{\pgfqpoint{2.659916in}{3.700177in}}%
\pgfpathlineto{\pgfqpoint{2.632476in}{3.773212in}}%
\pgfpathlineto{\pgfqpoint{2.679828in}{3.786586in}}%
\pgfpathlineto{\pgfqpoint{2.727097in}{3.799936in}}%
\pgfpathlineto{\pgfqpoint{2.754532in}{3.727043in}}%
\pgfpathlineto{\pgfqpoint{2.781993in}{3.654083in}}%
\pgfpathlineto{\pgfqpoint{2.734729in}{3.640591in}}%
\pgfpathlineto{\pgfqpoint{2.687382in}{3.627076in}}%
\pgfpathclose%
\pgfusepath{stroke,fill}%
\end{pgfscope}%
\begin{pgfscope}%
\pgfpathrectangle{\pgfqpoint{0.050000in}{0.083252in}}{\pgfqpoint{4.900000in}{4.900000in}}%
\pgfusepath{clip}%
\pgfsetbuttcap%
\pgfsetroundjoin%
\definecolor{currentfill}{rgb}{0.270342,0.538562,0.984760}%
\pgfsetfillcolor{currentfill}%
\pgfsetfillopacity{0.950000}%
\pgfsetlinewidth{0.200750pt}%
\definecolor{currentstroke}{rgb}{0.200000,0.200000,0.200000}%
\pgfsetstrokecolor{currentstroke}%
\pgfsetdash{}{0pt}%
\pgfpathmoveto{\pgfqpoint{2.278083in}{4.157962in}}%
\pgfpathlineto{\pgfqpoint{2.250814in}{4.230804in}}%
\pgfpathlineto{\pgfqpoint{2.298513in}{4.243779in}}%
\pgfpathlineto{\pgfqpoint{2.346129in}{4.256731in}}%
\pgfpathlineto{\pgfqpoint{2.373395in}{4.184029in}}%
\pgfpathlineto{\pgfqpoint{2.325781in}{4.171007in}}%
\pgfpathlineto{\pgfqpoint{2.278083in}{4.157962in}}%
\pgfpathclose%
\pgfusepath{stroke,fill}%
\end{pgfscope}%
\begin{pgfscope}%
\pgfpathrectangle{\pgfqpoint{0.050000in}{0.083252in}}{\pgfqpoint{4.900000in}{4.900000in}}%
\pgfusepath{clip}%
\pgfsetbuttcap%
\pgfsetroundjoin%
\definecolor{currentfill}{rgb}{0.251396,0.337255,0.717140}%
\pgfsetfillcolor{currentfill}%
\pgfsetfillopacity{0.950000}%
\pgfsetlinewidth{0.200750pt}%
\definecolor{currentstroke}{rgb}{0.200000,0.200000,0.200000}%
\pgfsetstrokecolor{currentstroke}%
\pgfsetdash{}{0pt}%
\pgfpathmoveto{\pgfqpoint{2.892094in}{3.361593in}}%
\pgfpathlineto{\pgfqpoint{2.864530in}{3.434811in}}%
\pgfpathlineto{\pgfqpoint{2.836992in}{3.507966in}}%
\pgfpathlineto{\pgfqpoint{2.884169in}{3.521577in}}%
\pgfpathlineto{\pgfqpoint{2.931263in}{3.535165in}}%
\pgfpathlineto{\pgfqpoint{2.958796in}{3.462155in}}%
\pgfpathlineto{\pgfqpoint{2.986355in}{3.389085in}}%
\pgfpathlineto{\pgfqpoint{2.939266in}{3.375351in}}%
\pgfpathlineto{\pgfqpoint{2.892094in}{3.361593in}}%
\pgfpathclose%
\pgfusepath{stroke,fill}%
\end{pgfscope}%
\begin{pgfscope}%
\pgfpathrectangle{\pgfqpoint{0.050000in}{0.083252in}}{\pgfqpoint{4.900000in}{4.900000in}}%
\pgfusepath{clip}%
\pgfsetbuttcap%
\pgfsetroundjoin%
\definecolor{currentfill}{rgb}{0.244752,0.266667,0.623299}%
\pgfsetfillcolor{currentfill}%
\pgfsetfillopacity{0.950000}%
\pgfsetlinewidth{0.200750pt}%
\definecolor{currentstroke}{rgb}{0.200000,0.200000,0.200000}%
\pgfsetstrokecolor{currentstroke}%
\pgfsetdash{}{0pt}%
\pgfpathmoveto{\pgfqpoint{3.096851in}{3.096292in}}%
\pgfpathlineto{\pgfqpoint{3.069188in}{3.169554in}}%
\pgfpathlineto{\pgfqpoint{3.041551in}{3.242779in}}%
\pgfpathlineto{\pgfqpoint{3.088553in}{3.256641in}}%
\pgfpathlineto{\pgfqpoint{3.135472in}{3.270481in}}%
\pgfpathlineto{\pgfqpoint{3.163105in}{3.197413in}}%
\pgfpathlineto{\pgfqpoint{3.190763in}{3.124315in}}%
\pgfpathlineto{\pgfqpoint{3.143848in}{3.110315in}}%
\pgfpathlineto{\pgfqpoint{3.096851in}{3.096292in}}%
\pgfpathclose%
\pgfusepath{stroke,fill}%
\end{pgfscope}%
\begin{pgfscope}%
\pgfpathrectangle{\pgfqpoint{0.050000in}{0.083252in}}{\pgfqpoint{4.900000in}{4.900000in}}%
\pgfusepath{clip}%
\pgfsetbuttcap%
\pgfsetroundjoin%
\definecolor{currentfill}{rgb}{0.260100,0.229589,0.557586}%
\pgfsetfillcolor{currentfill}%
\pgfsetfillopacity{0.950000}%
\pgfsetlinewidth{0.200750pt}%
\definecolor{currentstroke}{rgb}{0.200000,0.200000,0.200000}%
\pgfsetstrokecolor{currentstroke}%
\pgfsetdash{}{0pt}%
\pgfpathmoveto{\pgfqpoint{3.301669in}{2.831967in}}%
\pgfpathlineto{\pgfqpoint{3.273901in}{2.904998in}}%
\pgfpathlineto{\pgfqpoint{3.246161in}{2.978086in}}%
\pgfpathlineto{\pgfqpoint{3.292989in}{2.992239in}}%
\pgfpathlineto{\pgfqpoint{3.339735in}{3.006371in}}%
\pgfpathlineto{\pgfqpoint{3.367472in}{2.933472in}}%
\pgfpathlineto{\pgfqpoint{3.395237in}{2.860642in}}%
\pgfpathlineto{\pgfqpoint{3.348494in}{2.846315in}}%
\pgfpathlineto{\pgfqpoint{3.301669in}{2.831967in}}%
\pgfpathclose%
\pgfusepath{stroke,fill}%
\end{pgfscope}%
\begin{pgfscope}%
\pgfpathrectangle{\pgfqpoint{0.050000in}{0.083252in}}{\pgfqpoint{4.900000in}{4.900000in}}%
\pgfusepath{clip}%
\pgfsetbuttcap%
\pgfsetroundjoin%
\definecolor{currentfill}{rgb}{0.415240,0.391126,0.650350}%
\pgfsetfillcolor{currentfill}%
\pgfsetfillopacity{0.950000}%
\pgfsetlinewidth{0.200750pt}%
\definecolor{currentstroke}{rgb}{0.200000,0.200000,0.200000}%
\pgfsetstrokecolor{currentstroke}%
\pgfsetdash{}{0pt}%
\pgfpathmoveto{\pgfqpoint{3.506611in}{2.571049in}}%
\pgfpathlineto{\pgfqpoint{3.478717in}{2.643054in}}%
\pgfpathlineto{\pgfqpoint{3.450858in}{2.715374in}}%
\pgfpathlineto{\pgfqpoint{3.497517in}{2.729903in}}%
\pgfpathlineto{\pgfqpoint{3.544094in}{2.744412in}}%
\pgfpathlineto{\pgfqpoint{3.571953in}{2.672337in}}%
\pgfpathlineto{\pgfqpoint{3.599848in}{2.600592in}}%
\pgfpathlineto{\pgfqpoint{3.553270in}{2.585831in}}%
\pgfpathlineto{\pgfqpoint{3.506611in}{2.571049in}}%
\pgfpathclose%
\pgfusepath{stroke,fill}%
\end{pgfscope}%
\begin{pgfscope}%
\pgfpathrectangle{\pgfqpoint{0.050000in}{0.083252in}}{\pgfqpoint{4.900000in}{4.900000in}}%
\pgfusepath{clip}%
\pgfsetbuttcap%
\pgfsetroundjoin%
\definecolor{currentfill}{rgb}{0.267882,0.512418,0.950004}%
\pgfsetfillcolor{currentfill}%
\pgfsetfillopacity{0.950000}%
\pgfsetlinewidth{0.200750pt}%
\definecolor{currentstroke}{rgb}{0.200000,0.200000,0.200000}%
\pgfsetstrokecolor{currentstroke}%
\pgfsetdash{}{0pt}%
\pgfpathmoveto{\pgfqpoint{2.332698in}{4.012074in}}%
\pgfpathlineto{\pgfqpoint{2.305378in}{4.085052in}}%
\pgfpathlineto{\pgfqpoint{2.278083in}{4.157962in}}%
\pgfpathlineto{\pgfqpoint{2.325781in}{4.171007in}}%
\pgfpathlineto{\pgfqpoint{2.373395in}{4.184029in}}%
\pgfpathlineto{\pgfqpoint{2.400686in}{4.111259in}}%
\pgfpathlineto{\pgfqpoint{2.428002in}{4.038422in}}%
\pgfpathlineto{\pgfqpoint{2.380392in}{4.025260in}}%
\pgfpathlineto{\pgfqpoint{2.332698in}{4.012074in}}%
\pgfpathclose%
\pgfusepath{stroke,fill}%
\end{pgfscope}%
\begin{pgfscope}%
\pgfpathrectangle{\pgfqpoint{0.050000in}{0.083252in}}{\pgfqpoint{4.900000in}{4.900000in}}%
\pgfusepath{clip}%
\pgfsetbuttcap%
\pgfsetroundjoin%
\definecolor{currentfill}{rgb}{0.636017,0.621007,0.782361}%
\pgfsetfillcolor{currentfill}%
\pgfsetfillopacity{0.950000}%
\pgfsetlinewidth{0.200750pt}%
\definecolor{currentstroke}{rgb}{0.200000,0.200000,0.200000}%
\pgfsetstrokecolor{currentstroke}%
\pgfsetdash{}{0pt}%
\pgfpathmoveto{\pgfqpoint{3.861169in}{2.214991in}}%
\pgfpathlineto{\pgfqpoint{3.832950in}{2.281782in}}%
\pgfpathlineto{\pgfqpoint{3.804800in}{2.349751in}}%
\pgfpathlineto{\pgfqpoint{3.851146in}{2.365024in}}%
\pgfpathlineto{\pgfqpoint{3.879302in}{2.297208in}}%
\pgfpathlineto{\pgfqpoint{3.907527in}{2.230567in}}%
\pgfpathlineto{\pgfqpoint{3.861169in}{2.214991in}}%
\pgfpathclose%
\pgfusepath{stroke,fill}%
\end{pgfscope}%
\begin{pgfscope}%
\pgfpathrectangle{\pgfqpoint{0.050000in}{0.083252in}}{\pgfqpoint{4.900000in}{4.900000in}}%
\pgfusepath{clip}%
\pgfsetbuttcap%
\pgfsetroundjoin%
\definecolor{currentfill}{rgb}{0.261238,0.441830,0.856163}%
\pgfsetfillcolor{currentfill}%
\pgfsetfillopacity{0.950000}%
\pgfsetlinewidth{0.200750pt}%
\definecolor{currentstroke}{rgb}{0.200000,0.200000,0.200000}%
\pgfsetstrokecolor{currentstroke}%
\pgfsetdash{}{0pt}%
\pgfpathmoveto{\pgfqpoint{2.537523in}{3.746394in}}%
\pgfpathlineto{\pgfqpoint{2.510104in}{3.819503in}}%
\pgfpathlineto{\pgfqpoint{2.482711in}{3.892544in}}%
\pgfpathlineto{\pgfqpoint{2.530234in}{3.905823in}}%
\pgfpathlineto{\pgfqpoint{2.577673in}{3.919079in}}%
\pgfpathlineto{\pgfqpoint{2.605062in}{3.846179in}}%
\pgfpathlineto{\pgfqpoint{2.632476in}{3.773212in}}%
\pgfpathlineto{\pgfqpoint{2.585041in}{3.759815in}}%
\pgfpathlineto{\pgfqpoint{2.537523in}{3.746394in}}%
\pgfpathclose%
\pgfusepath{stroke,fill}%
\end{pgfscope}%
\begin{pgfscope}%
\pgfpathrectangle{\pgfqpoint{0.050000in}{0.083252in}}{\pgfqpoint{4.900000in}{4.900000in}}%
\pgfusepath{clip}%
\pgfsetbuttcap%
\pgfsetroundjoin%
\definecolor{currentfill}{rgb}{0.570381,0.552664,0.743114}%
\pgfsetfillcolor{currentfill}%
\pgfsetfillopacity{0.950000}%
\pgfsetlinewidth{0.200750pt}%
\definecolor{currentstroke}{rgb}{0.200000,0.200000,0.200000}%
\pgfsetstrokecolor{currentstroke}%
\pgfsetdash{}{0pt}%
\pgfpathmoveto{\pgfqpoint{3.711867in}{2.319125in}}%
\pgfpathlineto{\pgfqpoint{3.683788in}{2.388399in}}%
\pgfpathlineto{\pgfqpoint{3.655761in}{2.458499in}}%
\pgfpathlineto{\pgfqpoint{3.702261in}{2.473521in}}%
\pgfpathlineto{\pgfqpoint{3.748680in}{2.488519in}}%
\pgfpathlineto{\pgfqpoint{3.776713in}{2.418718in}}%
\pgfpathlineto{\pgfqpoint{3.804800in}{2.349751in}}%
\pgfpathlineto{\pgfqpoint{3.758374in}{2.334451in}}%
\pgfpathlineto{\pgfqpoint{3.711867in}{2.319125in}}%
\pgfpathclose%
\pgfusepath{stroke,fill}%
\end{pgfscope}%
\begin{pgfscope}%
\pgfpathrectangle{\pgfqpoint{0.050000in}{0.083252in}}{\pgfqpoint{4.900000in}{4.900000in}}%
\pgfusepath{clip}%
\pgfsetbuttcap%
\pgfsetroundjoin%
\definecolor{currentfill}{rgb}{0.254594,0.371242,0.762322}%
\pgfsetfillcolor{currentfill}%
\pgfsetfillopacity{0.950000}%
\pgfsetlinewidth{0.200750pt}%
\definecolor{currentstroke}{rgb}{0.200000,0.200000,0.200000}%
\pgfsetstrokecolor{currentstroke}%
\pgfsetdash{}{0pt}%
\pgfpathmoveto{\pgfqpoint{2.742389in}{3.480672in}}%
\pgfpathlineto{\pgfqpoint{2.714873in}{3.553907in}}%
\pgfpathlineto{\pgfqpoint{2.687382in}{3.627076in}}%
\pgfpathlineto{\pgfqpoint{2.734729in}{3.640591in}}%
\pgfpathlineto{\pgfqpoint{2.781993in}{3.654083in}}%
\pgfpathlineto{\pgfqpoint{2.809480in}{3.581057in}}%
\pgfpathlineto{\pgfqpoint{2.836992in}{3.507966in}}%
\pgfpathlineto{\pgfqpoint{2.789732in}{3.494331in}}%
\pgfpathlineto{\pgfqpoint{2.742389in}{3.480672in}}%
\pgfpathclose%
\pgfusepath{stroke,fill}%
\end{pgfscope}%
\begin{pgfscope}%
\pgfpathrectangle{\pgfqpoint{0.050000in}{0.083252in}}{\pgfqpoint{4.900000in}{4.900000in}}%
\pgfusepath{clip}%
\pgfsetbuttcap%
\pgfsetroundjoin%
\definecolor{currentfill}{rgb}{0.248197,0.303268,0.671957}%
\pgfsetfillcolor{currentfill}%
\pgfsetfillopacity{0.950000}%
\pgfsetlinewidth{0.200750pt}%
\definecolor{currentstroke}{rgb}{0.200000,0.200000,0.200000}%
\pgfsetstrokecolor{currentstroke}%
\pgfsetdash{}{0pt}%
\pgfpathmoveto{\pgfqpoint{2.947299in}{3.214985in}}%
\pgfpathlineto{\pgfqpoint{2.919683in}{3.288317in}}%
\pgfpathlineto{\pgfqpoint{2.892094in}{3.361593in}}%
\pgfpathlineto{\pgfqpoint{2.939266in}{3.375351in}}%
\pgfpathlineto{\pgfqpoint{2.986355in}{3.389085in}}%
\pgfpathlineto{\pgfqpoint{3.013940in}{3.315958in}}%
\pgfpathlineto{\pgfqpoint{3.041551in}{3.242779in}}%
\pgfpathlineto{\pgfqpoint{2.994466in}{3.228893in}}%
\pgfpathlineto{\pgfqpoint{2.947299in}{3.214985in}}%
\pgfpathclose%
\pgfusepath{stroke,fill}%
\end{pgfscope}%
\begin{pgfscope}%
\pgfpathrectangle{\pgfqpoint{0.050000in}{0.083252in}}{\pgfqpoint{4.900000in}{4.900000in}}%
\pgfusepath{clip}%
\pgfsetbuttcap%
\pgfsetroundjoin%
\definecolor{currentfill}{rgb}{0.241553,0.232680,0.578116}%
\pgfsetfillcolor{currentfill}%
\pgfsetfillopacity{0.950000}%
\pgfsetlinewidth{0.200750pt}%
\definecolor{currentstroke}{rgb}{0.200000,0.200000,0.200000}%
\pgfsetstrokecolor{currentstroke}%
\pgfsetdash{}{0pt}%
\pgfpathmoveto{\pgfqpoint{3.152257in}{2.949716in}}%
\pgfpathlineto{\pgfqpoint{3.124541in}{3.023007in}}%
\pgfpathlineto{\pgfqpoint{3.096851in}{3.096292in}}%
\pgfpathlineto{\pgfqpoint{3.143848in}{3.110315in}}%
\pgfpathlineto{\pgfqpoint{3.190763in}{3.124315in}}%
\pgfpathlineto{\pgfqpoint{3.218449in}{3.051200in}}%
\pgfpathlineto{\pgfqpoint{3.246161in}{2.978086in}}%
\pgfpathlineto{\pgfqpoint{3.199250in}{2.963911in}}%
\pgfpathlineto{\pgfqpoint{3.152257in}{2.949716in}}%
\pgfpathclose%
\pgfusepath{stroke,fill}%
\end{pgfscope}%
\begin{pgfscope}%
\pgfpathrectangle{\pgfqpoint{0.050000in}{0.083252in}}{\pgfqpoint{4.900000in}{4.900000in}}%
\pgfusepath{clip}%
\pgfsetbuttcap%
\pgfsetroundjoin%
\definecolor{currentfill}{rgb}{0.343637,0.316571,0.607536}%
\pgfsetfillcolor{currentfill}%
\pgfsetfillopacity{0.950000}%
\pgfsetlinewidth{0.200750pt}%
\definecolor{currentstroke}{rgb}{0.200000,0.200000,0.200000}%
\pgfsetstrokecolor{currentstroke}%
\pgfsetdash{}{0pt}%
\pgfpathmoveto{\pgfqpoint{3.357294in}{2.686255in}}%
\pgfpathlineto{\pgfqpoint{3.329466in}{2.759036in}}%
\pgfpathlineto{\pgfqpoint{3.301669in}{2.831967in}}%
\pgfpathlineto{\pgfqpoint{3.348494in}{2.846315in}}%
\pgfpathlineto{\pgfqpoint{3.395237in}{2.860642in}}%
\pgfpathlineto{\pgfqpoint{3.423031in}{2.787925in}}%
\pgfpathlineto{\pgfqpoint{3.450858in}{2.715374in}}%
\pgfpathlineto{\pgfqpoint{3.404117in}{2.700825in}}%
\pgfpathlineto{\pgfqpoint{3.357294in}{2.686255in}}%
\pgfpathclose%
\pgfusepath{stroke,fill}%
\end{pgfscope}%
\begin{pgfscope}%
\pgfpathrectangle{\pgfqpoint{0.050000in}{0.083252in}}{\pgfqpoint{4.900000in}{4.900000in}}%
\pgfusepath{clip}%
\pgfsetbuttcap%
\pgfsetroundjoin%
\definecolor{currentfill}{rgb}{0.492810,0.471895,0.696732}%
\pgfsetfillcolor{currentfill}%
\pgfsetfillopacity{0.950000}%
\pgfsetlinewidth{0.200750pt}%
\definecolor{currentstroke}{rgb}{0.200000,0.200000,0.200000}%
\pgfsetstrokecolor{currentstroke}%
\pgfsetdash{}{0pt}%
\pgfpathmoveto{\pgfqpoint{3.562519in}{2.428387in}}%
\pgfpathlineto{\pgfqpoint{3.534544in}{2.499454in}}%
\pgfpathlineto{\pgfqpoint{3.506611in}{2.571049in}}%
\pgfpathlineto{\pgfqpoint{3.553270in}{2.585831in}}%
\pgfpathlineto{\pgfqpoint{3.599848in}{2.600592in}}%
\pgfpathlineto{\pgfqpoint{3.627783in}{2.529274in}}%
\pgfpathlineto{\pgfqpoint{3.655761in}{2.458499in}}%
\pgfpathlineto{\pgfqpoint{3.609181in}{2.443454in}}%
\pgfpathlineto{\pgfqpoint{3.562519in}{2.428387in}}%
\pgfpathclose%
\pgfusepath{stroke,fill}%
\end{pgfscope}%
\begin{pgfscope}%
\pgfpathrectangle{\pgfqpoint{0.050000in}{0.083252in}}{\pgfqpoint{4.900000in}{4.900000in}}%
\pgfusepath{clip}%
\pgfsetbuttcap%
\pgfsetroundjoin%
\definecolor{currentfill}{rgb}{0.264437,0.475817,0.901346}%
\pgfsetfillcolor{currentfill}%
\pgfsetfillopacity{0.950000}%
\pgfsetlinewidth{0.200750pt}%
\definecolor{currentstroke}{rgb}{0.200000,0.200000,0.200000}%
\pgfsetstrokecolor{currentstroke}%
\pgfsetdash{}{0pt}%
\pgfpathmoveto{\pgfqpoint{2.387414in}{3.865914in}}%
\pgfpathlineto{\pgfqpoint{2.360043in}{3.939028in}}%
\pgfpathlineto{\pgfqpoint{2.332698in}{4.012074in}}%
\pgfpathlineto{\pgfqpoint{2.380392in}{4.025260in}}%
\pgfpathlineto{\pgfqpoint{2.428002in}{4.038422in}}%
\pgfpathlineto{\pgfqpoint{2.455344in}{3.965517in}}%
\pgfpathlineto{\pgfqpoint{2.482711in}{3.892544in}}%
\pgfpathlineto{\pgfqpoint{2.435105in}{3.879241in}}%
\pgfpathlineto{\pgfqpoint{2.387414in}{3.865914in}}%
\pgfpathclose%
\pgfusepath{stroke,fill}%
\end{pgfscope}%
\begin{pgfscope}%
\pgfpathrectangle{\pgfqpoint{0.050000in}{0.083252in}}{\pgfqpoint{4.900000in}{4.900000in}}%
\pgfusepath{clip}%
\pgfsetbuttcap%
\pgfsetroundjoin%
\definecolor{currentfill}{rgb}{0.258039,0.407843,0.810980}%
\pgfsetfillcolor{currentfill}%
\pgfsetfillopacity{0.950000}%
\pgfsetlinewidth{0.200750pt}%
\definecolor{currentstroke}{rgb}{0.200000,0.200000,0.200000}%
\pgfsetstrokecolor{currentstroke}%
\pgfsetdash{}{0pt}%
\pgfpathmoveto{\pgfqpoint{2.592436in}{3.599973in}}%
\pgfpathlineto{\pgfqpoint{2.564967in}{3.673217in}}%
\pgfpathlineto{\pgfqpoint{2.537523in}{3.746394in}}%
\pgfpathlineto{\pgfqpoint{2.585041in}{3.759815in}}%
\pgfpathlineto{\pgfqpoint{2.632476in}{3.773212in}}%
\pgfpathlineto{\pgfqpoint{2.659916in}{3.700177in}}%
\pgfpathlineto{\pgfqpoint{2.687382in}{3.627076in}}%
\pgfpathlineto{\pgfqpoint{2.639951in}{3.613536in}}%
\pgfpathlineto{\pgfqpoint{2.592436in}{3.599973in}}%
\pgfpathclose%
\pgfusepath{stroke,fill}%
\end{pgfscope}%
\begin{pgfscope}%
\pgfpathrectangle{\pgfqpoint{0.050000in}{0.083252in}}{\pgfqpoint{4.900000in}{4.900000in}}%
\pgfusepath{clip}%
\pgfsetbuttcap%
\pgfsetroundjoin%
\definecolor{currentfill}{rgb}{0.270342,0.538562,0.984760}%
\pgfsetfillcolor{currentfill}%
\pgfsetfillopacity{0.950000}%
\pgfsetlinewidth{0.200750pt}%
\definecolor{currentstroke}{rgb}{0.200000,0.200000,0.200000}%
\pgfsetstrokecolor{currentstroke}%
\pgfsetdash{}{0pt}%
\pgfpathmoveto{\pgfqpoint{2.182435in}{4.131803in}}%
\pgfpathlineto{\pgfqpoint{2.155162in}{4.204786in}}%
\pgfpathlineto{\pgfqpoint{2.203030in}{4.217807in}}%
\pgfpathlineto{\pgfqpoint{2.250814in}{4.230804in}}%
\pgfpathlineto{\pgfqpoint{2.278083in}{4.157962in}}%
\pgfpathlineto{\pgfqpoint{2.230301in}{4.144894in}}%
\pgfpathlineto{\pgfqpoint{2.182435in}{4.131803in}}%
\pgfpathclose%
\pgfusepath{stroke,fill}%
\end{pgfscope}%
\begin{pgfscope}%
\pgfpathrectangle{\pgfqpoint{0.050000in}{0.083252in}}{\pgfqpoint{4.900000in}{4.900000in}}%
\pgfusepath{clip}%
\pgfsetbuttcap%
\pgfsetroundjoin%
\definecolor{currentfill}{rgb}{0.251396,0.337255,0.717140}%
\pgfsetfillcolor{currentfill}%
\pgfsetfillopacity{0.950000}%
\pgfsetlinewidth{0.200750pt}%
\definecolor{currentstroke}{rgb}{0.200000,0.200000,0.200000}%
\pgfsetstrokecolor{currentstroke}%
\pgfsetdash{}{0pt}%
\pgfpathmoveto{\pgfqpoint{2.797500in}{3.334008in}}%
\pgfpathlineto{\pgfqpoint{2.769932in}{3.407372in}}%
\pgfpathlineto{\pgfqpoint{2.742389in}{3.480672in}}%
\pgfpathlineto{\pgfqpoint{2.789732in}{3.494331in}}%
\pgfpathlineto{\pgfqpoint{2.836992in}{3.507966in}}%
\pgfpathlineto{\pgfqpoint{2.864530in}{3.434811in}}%
\pgfpathlineto{\pgfqpoint{2.892094in}{3.361593in}}%
\pgfpathlineto{\pgfqpoint{2.844839in}{3.347812in}}%
\pgfpathlineto{\pgfqpoint{2.797500in}{3.334008in}}%
\pgfpathclose%
\pgfusepath{stroke,fill}%
\end{pgfscope}%
\begin{pgfscope}%
\pgfpathrectangle{\pgfqpoint{0.050000in}{0.083252in}}{\pgfqpoint{4.900000in}{4.900000in}}%
\pgfusepath{clip}%
\pgfsetbuttcap%
\pgfsetroundjoin%
\definecolor{currentfill}{rgb}{0.244752,0.266667,0.623299}%
\pgfsetfillcolor{currentfill}%
\pgfsetfillopacity{0.950000}%
\pgfsetlinewidth{0.200750pt}%
\definecolor{currentstroke}{rgb}{0.200000,0.200000,0.200000}%
\pgfsetstrokecolor{currentstroke}%
\pgfsetdash{}{0pt}%
\pgfpathmoveto{\pgfqpoint{3.002608in}{3.068181in}}%
\pgfpathlineto{\pgfqpoint{2.974940in}{3.141604in}}%
\pgfpathlineto{\pgfqpoint{2.947299in}{3.214985in}}%
\pgfpathlineto{\pgfqpoint{2.994466in}{3.228893in}}%
\pgfpathlineto{\pgfqpoint{3.041551in}{3.242779in}}%
\pgfpathlineto{\pgfqpoint{3.069188in}{3.169554in}}%
\pgfpathlineto{\pgfqpoint{3.096851in}{3.096292in}}%
\pgfpathlineto{\pgfqpoint{3.049771in}{3.082248in}}%
\pgfpathlineto{\pgfqpoint{3.002608in}{3.068181in}}%
\pgfpathclose%
\pgfusepath{stroke,fill}%
\end{pgfscope}%
\begin{pgfscope}%
\pgfpathrectangle{\pgfqpoint{0.050000in}{0.083252in}}{\pgfqpoint{4.900000in}{4.900000in}}%
\pgfusepath{clip}%
\pgfsetbuttcap%
\pgfsetroundjoin%
\definecolor{currentfill}{rgb}{0.707620,0.695563,0.825175}%
\pgfsetfillcolor{currentfill}%
\pgfsetfillopacity{0.950000}%
\pgfsetlinewidth{0.200750pt}%
\definecolor{currentstroke}{rgb}{0.200000,0.200000,0.200000}%
\pgfsetstrokecolor{currentstroke}%
\pgfsetdash{}{0pt}%
\pgfpathmoveto{\pgfqpoint{3.917847in}{2.085704in}}%
\pgfpathlineto{\pgfqpoint{3.889465in}{2.149566in}}%
\pgfpathlineto{\pgfqpoint{3.861169in}{2.214991in}}%
\pgfpathlineto{\pgfqpoint{3.907527in}{2.230567in}}%
\pgfpathlineto{\pgfqpoint{3.935831in}{2.165288in}}%
\pgfpathlineto{\pgfqpoint{3.964223in}{2.101563in}}%
\pgfpathlineto{\pgfqpoint{3.917847in}{2.085704in}}%
\pgfpathclose%
\pgfusepath{stroke,fill}%
\end{pgfscope}%
\begin{pgfscope}%
\pgfpathrectangle{\pgfqpoint{0.050000in}{0.083252in}}{\pgfqpoint{4.900000in}{4.900000in}}%
\pgfusepath{clip}%
\pgfsetbuttcap%
\pgfsetroundjoin%
\definecolor{currentfill}{rgb}{0.641984,0.627220,0.785928}%
\pgfsetfillcolor{currentfill}%
\pgfsetfillopacity{0.950000}%
\pgfsetlinewidth{0.200750pt}%
\definecolor{currentstroke}{rgb}{0.200000,0.200000,0.200000}%
\pgfsetstrokecolor{currentstroke}%
\pgfsetdash{}{0pt}%
\pgfpathmoveto{\pgfqpoint{3.768211in}{2.183746in}}%
\pgfpathlineto{\pgfqpoint{3.740006in}{2.250846in}}%
\pgfpathlineto{\pgfqpoint{3.711867in}{2.319125in}}%
\pgfpathlineto{\pgfqpoint{3.758374in}{2.334451in}}%
\pgfpathlineto{\pgfqpoint{3.804800in}{2.349751in}}%
\pgfpathlineto{\pgfqpoint{3.832950in}{2.281782in}}%
\pgfpathlineto{\pgfqpoint{3.861169in}{2.214991in}}%
\pgfpathlineto{\pgfqpoint{3.814730in}{2.199384in}}%
\pgfpathlineto{\pgfqpoint{3.768211in}{2.183746in}}%
\pgfpathclose%
\pgfusepath{stroke,fill}%
\end{pgfscope}%
\begin{pgfscope}%
\pgfpathrectangle{\pgfqpoint{0.050000in}{0.083252in}}{\pgfqpoint{4.900000in}{4.900000in}}%
\pgfusepath{clip}%
\pgfsetbuttcap%
\pgfsetroundjoin%
\definecolor{currentfill}{rgb}{0.260100,0.229589,0.557586}%
\pgfsetfillcolor{currentfill}%
\pgfsetfillopacity{0.950000}%
\pgfsetlinewidth{0.200750pt}%
\definecolor{currentstroke}{rgb}{0.200000,0.200000,0.200000}%
\pgfsetstrokecolor{currentstroke}%
\pgfsetdash{}{0pt}%
\pgfpathmoveto{\pgfqpoint{3.207772in}{2.803211in}}%
\pgfpathlineto{\pgfqpoint{3.180000in}{2.876440in}}%
\pgfpathlineto{\pgfqpoint{3.152257in}{2.949716in}}%
\pgfpathlineto{\pgfqpoint{3.199250in}{2.963911in}}%
\pgfpathlineto{\pgfqpoint{3.246161in}{2.978086in}}%
\pgfpathlineto{\pgfqpoint{3.273901in}{2.904998in}}%
\pgfpathlineto{\pgfqpoint{3.301669in}{2.831967in}}%
\pgfpathlineto{\pgfqpoint{3.254762in}{2.817599in}}%
\pgfpathlineto{\pgfqpoint{3.207772in}{2.803211in}}%
\pgfpathclose%
\pgfusepath{stroke,fill}%
\end{pgfscope}%
\begin{pgfscope}%
\pgfpathrectangle{\pgfqpoint{0.050000in}{0.083252in}}{\pgfqpoint{4.900000in}{4.900000in}}%
\pgfusepath{clip}%
\pgfsetbuttcap%
\pgfsetroundjoin%
\definecolor{currentfill}{rgb}{0.421207,0.397339,0.653918}%
\pgfsetfillcolor{currentfill}%
\pgfsetfillopacity{0.950000}%
\pgfsetlinewidth{0.200750pt}%
\definecolor{currentstroke}{rgb}{0.200000,0.200000,0.200000}%
\pgfsetstrokecolor{currentstroke}%
\pgfsetdash{}{0pt}%
\pgfpathmoveto{\pgfqpoint{3.413048in}{2.541422in}}%
\pgfpathlineto{\pgfqpoint{3.385154in}{2.613691in}}%
\pgfpathlineto{\pgfqpoint{3.357294in}{2.686255in}}%
\pgfpathlineto{\pgfqpoint{3.404117in}{2.700825in}}%
\pgfpathlineto{\pgfqpoint{3.450858in}{2.715374in}}%
\pgfpathlineto{\pgfqpoint{3.478717in}{2.643054in}}%
\pgfpathlineto{\pgfqpoint{3.506611in}{2.571049in}}%
\pgfpathlineto{\pgfqpoint{3.459871in}{2.556246in}}%
\pgfpathlineto{\pgfqpoint{3.413048in}{2.541422in}}%
\pgfpathclose%
\pgfusepath{stroke,fill}%
\end{pgfscope}%
\begin{pgfscope}%
\pgfpathrectangle{\pgfqpoint{0.050000in}{0.083252in}}{\pgfqpoint{4.900000in}{4.900000in}}%
\pgfusepath{clip}%
\pgfsetbuttcap%
\pgfsetroundjoin%
\definecolor{currentfill}{rgb}{0.267882,0.512418,0.950004}%
\pgfsetfillcolor{currentfill}%
\pgfsetfillopacity{0.950000}%
\pgfsetlinewidth{0.200750pt}%
\definecolor{currentstroke}{rgb}{0.200000,0.200000,0.200000}%
\pgfsetstrokecolor{currentstroke}%
\pgfsetdash{}{0pt}%
\pgfpathmoveto{\pgfqpoint{2.237056in}{3.985633in}}%
\pgfpathlineto{\pgfqpoint{2.209733in}{4.058752in}}%
\pgfpathlineto{\pgfqpoint{2.182435in}{4.131803in}}%
\pgfpathlineto{\pgfqpoint{2.230301in}{4.144894in}}%
\pgfpathlineto{\pgfqpoint{2.278083in}{4.157962in}}%
\pgfpathlineto{\pgfqpoint{2.305378in}{4.085052in}}%
\pgfpathlineto{\pgfqpoint{2.332698in}{4.012074in}}%
\pgfpathlineto{\pgfqpoint{2.284919in}{3.998865in}}%
\pgfpathlineto{\pgfqpoint{2.237056in}{3.985633in}}%
\pgfpathclose%
\pgfusepath{stroke,fill}%
\end{pgfscope}%
\begin{pgfscope}%
\pgfpathrectangle{\pgfqpoint{0.050000in}{0.083252in}}{\pgfqpoint{4.900000in}{4.900000in}}%
\pgfusepath{clip}%
\pgfsetbuttcap%
\pgfsetroundjoin%
\definecolor{currentfill}{rgb}{0.261238,0.441830,0.856163}%
\pgfsetfillcolor{currentfill}%
\pgfsetfillopacity{0.950000}%
\pgfsetlinewidth{0.200750pt}%
\definecolor{currentstroke}{rgb}{0.200000,0.200000,0.200000}%
\pgfsetstrokecolor{currentstroke}%
\pgfsetdash{}{0pt}%
\pgfpathmoveto{\pgfqpoint{2.442233in}{3.719481in}}%
\pgfpathlineto{\pgfqpoint{2.414811in}{3.792732in}}%
\pgfpathlineto{\pgfqpoint{2.387414in}{3.865914in}}%
\pgfpathlineto{\pgfqpoint{2.435105in}{3.879241in}}%
\pgfpathlineto{\pgfqpoint{2.482711in}{3.892544in}}%
\pgfpathlineto{\pgfqpoint{2.510104in}{3.819503in}}%
\pgfpathlineto{\pgfqpoint{2.537523in}{3.746394in}}%
\pgfpathlineto{\pgfqpoint{2.489920in}{3.732949in}}%
\pgfpathlineto{\pgfqpoint{2.442233in}{3.719481in}}%
\pgfpathclose%
\pgfusepath{stroke,fill}%
\end{pgfscope}%
\begin{pgfscope}%
\pgfpathrectangle{\pgfqpoint{0.050000in}{0.083252in}}{\pgfqpoint{4.900000in}{4.900000in}}%
\pgfusepath{clip}%
\pgfsetbuttcap%
\pgfsetroundjoin%
\definecolor{currentfill}{rgb}{0.570381,0.552664,0.743114}%
\pgfsetfillcolor{currentfill}%
\pgfsetfillopacity{0.950000}%
\pgfsetlinewidth{0.200750pt}%
\definecolor{currentstroke}{rgb}{0.200000,0.200000,0.200000}%
\pgfsetstrokecolor{currentstroke}%
\pgfsetdash{}{0pt}%
\pgfpathmoveto{\pgfqpoint{3.618611in}{2.288393in}}%
\pgfpathlineto{\pgfqpoint{3.590539in}{2.357981in}}%
\pgfpathlineto{\pgfqpoint{3.562519in}{2.428387in}}%
\pgfpathlineto{\pgfqpoint{3.609181in}{2.443454in}}%
\pgfpathlineto{\pgfqpoint{3.655761in}{2.458499in}}%
\pgfpathlineto{\pgfqpoint{3.683788in}{2.388399in}}%
\pgfpathlineto{\pgfqpoint{3.711867in}{2.319125in}}%
\pgfpathlineto{\pgfqpoint{3.665279in}{2.303772in}}%
\pgfpathlineto{\pgfqpoint{3.618611in}{2.288393in}}%
\pgfpathclose%
\pgfusepath{stroke,fill}%
\end{pgfscope}%
\begin{pgfscope}%
\pgfpathrectangle{\pgfqpoint{0.050000in}{0.083252in}}{\pgfqpoint{4.900000in}{4.900000in}}%
\pgfusepath{clip}%
\pgfsetbuttcap%
\pgfsetroundjoin%
\definecolor{currentfill}{rgb}{0.254594,0.371242,0.762322}%
\pgfsetfillcolor{currentfill}%
\pgfsetfillopacity{0.950000}%
\pgfsetlinewidth{0.200750pt}%
\definecolor{currentstroke}{rgb}{0.200000,0.200000,0.200000}%
\pgfsetstrokecolor{currentstroke}%
\pgfsetdash{}{0pt}%
\pgfpathmoveto{\pgfqpoint{2.647452in}{3.453283in}}%
\pgfpathlineto{\pgfqpoint{2.619931in}{3.526661in}}%
\pgfpathlineto{\pgfqpoint{2.592436in}{3.599973in}}%
\pgfpathlineto{\pgfqpoint{2.639951in}{3.613536in}}%
\pgfpathlineto{\pgfqpoint{2.687382in}{3.627076in}}%
\pgfpathlineto{\pgfqpoint{2.714873in}{3.553907in}}%
\pgfpathlineto{\pgfqpoint{2.742389in}{3.480672in}}%
\pgfpathlineto{\pgfqpoint{2.694963in}{3.466989in}}%
\pgfpathlineto{\pgfqpoint{2.647452in}{3.453283in}}%
\pgfpathclose%
\pgfusepath{stroke,fill}%
\end{pgfscope}%
\begin{pgfscope}%
\pgfpathrectangle{\pgfqpoint{0.050000in}{0.083252in}}{\pgfqpoint{4.900000in}{4.900000in}}%
\pgfusepath{clip}%
\pgfsetbuttcap%
\pgfsetroundjoin%
\definecolor{currentfill}{rgb}{0.247951,0.300654,0.668481}%
\pgfsetfillcolor{currentfill}%
\pgfsetfillopacity{0.950000}%
\pgfsetlinewidth{0.200750pt}%
\definecolor{currentstroke}{rgb}{0.200000,0.200000,0.200000}%
\pgfsetstrokecolor{currentstroke}%
\pgfsetdash{}{0pt}%
\pgfpathmoveto{\pgfqpoint{2.852714in}{3.187099in}}%
\pgfpathlineto{\pgfqpoint{2.825094in}{3.260583in}}%
\pgfpathlineto{\pgfqpoint{2.797500in}{3.334008in}}%
\pgfpathlineto{\pgfqpoint{2.844839in}{3.347812in}}%
\pgfpathlineto{\pgfqpoint{2.892094in}{3.361593in}}%
\pgfpathlineto{\pgfqpoint{2.919683in}{3.288317in}}%
\pgfpathlineto{\pgfqpoint{2.947299in}{3.214985in}}%
\pgfpathlineto{\pgfqpoint{2.900048in}{3.201054in}}%
\pgfpathlineto{\pgfqpoint{2.852714in}{3.187099in}}%
\pgfpathclose%
\pgfusepath{stroke,fill}%
\end{pgfscope}%
\begin{pgfscope}%
\pgfpathrectangle{\pgfqpoint{0.050000in}{0.083252in}}{\pgfqpoint{4.900000in}{4.900000in}}%
\pgfusepath{clip}%
\pgfsetbuttcap%
\pgfsetroundjoin%
\definecolor{currentfill}{rgb}{0.241553,0.232680,0.578116}%
\pgfsetfillcolor{currentfill}%
\pgfsetfillopacity{0.950000}%
\pgfsetlinewidth{0.200750pt}%
\definecolor{currentstroke}{rgb}{0.200000,0.200000,0.200000}%
\pgfsetstrokecolor{currentstroke}%
\pgfsetdash{}{0pt}%
\pgfpathmoveto{\pgfqpoint{3.058022in}{2.921261in}}%
\pgfpathlineto{\pgfqpoint{3.030301in}{2.994728in}}%
\pgfpathlineto{\pgfqpoint{3.002608in}{3.068181in}}%
\pgfpathlineto{\pgfqpoint{3.049771in}{3.082248in}}%
\pgfpathlineto{\pgfqpoint{3.096851in}{3.096292in}}%
\pgfpathlineto{\pgfqpoint{3.124541in}{3.023007in}}%
\pgfpathlineto{\pgfqpoint{3.152257in}{2.949716in}}%
\pgfpathlineto{\pgfqpoint{3.105181in}{2.935499in}}%
\pgfpathlineto{\pgfqpoint{3.058022in}{2.921261in}}%
\pgfpathclose%
\pgfusepath{stroke,fill}%
\end{pgfscope}%
\begin{pgfscope}%
\pgfpathrectangle{\pgfqpoint{0.050000in}{0.083252in}}{\pgfqpoint{4.900000in}{4.900000in}}%
\pgfusepath{clip}%
\pgfsetbuttcap%
\pgfsetroundjoin%
\definecolor{currentfill}{rgb}{0.343637,0.316571,0.607536}%
\pgfsetfillcolor{currentfill}%
\pgfsetfillopacity{0.950000}%
\pgfsetlinewidth{0.200750pt}%
\definecolor{currentstroke}{rgb}{0.200000,0.200000,0.200000}%
\pgfsetstrokecolor{currentstroke}%
\pgfsetdash{}{0pt}%
\pgfpathmoveto{\pgfqpoint{3.263402in}{2.657057in}}%
\pgfpathlineto{\pgfqpoint{3.235572in}{2.730066in}}%
\pgfpathlineto{\pgfqpoint{3.207772in}{2.803211in}}%
\pgfpathlineto{\pgfqpoint{3.254762in}{2.817599in}}%
\pgfpathlineto{\pgfqpoint{3.301669in}{2.831967in}}%
\pgfpathlineto{\pgfqpoint{3.329466in}{2.759036in}}%
\pgfpathlineto{\pgfqpoint{3.357294in}{2.686255in}}%
\pgfpathlineto{\pgfqpoint{3.310389in}{2.671666in}}%
\pgfpathlineto{\pgfqpoint{3.263402in}{2.657057in}}%
\pgfpathclose%
\pgfusepath{stroke,fill}%
\end{pgfscope}%
\begin{pgfscope}%
\pgfpathrectangle{\pgfqpoint{0.050000in}{0.083252in}}{\pgfqpoint{4.900000in}{4.900000in}}%
\pgfusepath{clip}%
\pgfsetbuttcap%
\pgfsetroundjoin%
\definecolor{currentfill}{rgb}{0.498777,0.478108,0.700300}%
\pgfsetfillcolor{currentfill}%
\pgfsetfillopacity{0.950000}%
\pgfsetlinewidth{0.200750pt}%
\definecolor{currentstroke}{rgb}{0.200000,0.200000,0.200000}%
\pgfsetstrokecolor{currentstroke}%
\pgfsetdash{}{0pt}%
\pgfpathmoveto{\pgfqpoint{3.468951in}{2.398184in}}%
\pgfpathlineto{\pgfqpoint{3.440979in}{2.469548in}}%
\pgfpathlineto{\pgfqpoint{3.413048in}{2.541422in}}%
\pgfpathlineto{\pgfqpoint{3.459871in}{2.556246in}}%
\pgfpathlineto{\pgfqpoint{3.506611in}{2.571049in}}%
\pgfpathlineto{\pgfqpoint{3.534544in}{2.499454in}}%
\pgfpathlineto{\pgfqpoint{3.562519in}{2.428387in}}%
\pgfpathlineto{\pgfqpoint{3.515775in}{2.413297in}}%
\pgfpathlineto{\pgfqpoint{3.468951in}{2.398184in}}%
\pgfpathclose%
\pgfusepath{stroke,fill}%
\end{pgfscope}%
\begin{pgfscope}%
\pgfpathrectangle{\pgfqpoint{0.050000in}{0.083252in}}{\pgfqpoint{4.900000in}{4.900000in}}%
\pgfusepath{clip}%
\pgfsetbuttcap%
\pgfsetroundjoin%
\definecolor{currentfill}{rgb}{0.707620,0.695563,0.825175}%
\pgfsetfillcolor{currentfill}%
\pgfsetfillopacity{0.950000}%
\pgfsetlinewidth{0.200750pt}%
\definecolor{currentstroke}{rgb}{0.200000,0.200000,0.200000}%
\pgfsetstrokecolor{currentstroke}%
\pgfsetdash{}{0pt}%
\pgfpathmoveto{\pgfqpoint{3.824854in}{2.053877in}}%
\pgfpathlineto{\pgfqpoint{3.796491in}{2.118021in}}%
\pgfpathlineto{\pgfqpoint{3.768211in}{2.183746in}}%
\pgfpathlineto{\pgfqpoint{3.814730in}{2.199384in}}%
\pgfpathlineto{\pgfqpoint{3.861169in}{2.214991in}}%
\pgfpathlineto{\pgfqpoint{3.889465in}{2.149566in}}%
\pgfpathlineto{\pgfqpoint{3.917847in}{2.085704in}}%
\pgfpathlineto{\pgfqpoint{3.871391in}{2.069809in}}%
\pgfpathlineto{\pgfqpoint{3.824854in}{2.053877in}}%
\pgfpathclose%
\pgfusepath{stroke,fill}%
\end{pgfscope}%
\begin{pgfscope}%
\pgfpathrectangle{\pgfqpoint{0.050000in}{0.083252in}}{\pgfqpoint{4.900000in}{4.900000in}}%
\pgfusepath{clip}%
\pgfsetbuttcap%
\pgfsetroundjoin%
\definecolor{currentfill}{rgb}{0.264437,0.475817,0.901346}%
\pgfsetfillcolor{currentfill}%
\pgfsetfillopacity{0.950000}%
\pgfsetlinewidth{0.200750pt}%
\definecolor{currentstroke}{rgb}{0.200000,0.200000,0.200000}%
\pgfsetstrokecolor{currentstroke}%
\pgfsetdash{}{0pt}%
\pgfpathmoveto{\pgfqpoint{2.291780in}{3.839190in}}%
\pgfpathlineto{\pgfqpoint{2.264406in}{3.912446in}}%
\pgfpathlineto{\pgfqpoint{2.237056in}{3.985633in}}%
\pgfpathlineto{\pgfqpoint{2.284919in}{3.998865in}}%
\pgfpathlineto{\pgfqpoint{2.332698in}{4.012074in}}%
\pgfpathlineto{\pgfqpoint{2.360043in}{3.939028in}}%
\pgfpathlineto{\pgfqpoint{2.387414in}{3.865914in}}%
\pgfpathlineto{\pgfqpoint{2.339640in}{3.852564in}}%
\pgfpathlineto{\pgfqpoint{2.291780in}{3.839190in}}%
\pgfpathclose%
\pgfusepath{stroke,fill}%
\end{pgfscope}%
\begin{pgfscope}%
\pgfpathrectangle{\pgfqpoint{0.050000in}{0.083252in}}{\pgfqpoint{4.900000in}{4.900000in}}%
\pgfusepath{clip}%
\pgfsetbuttcap%
\pgfsetroundjoin%
\definecolor{currentfill}{rgb}{0.767290,0.757693,0.860854}%
\pgfsetfillcolor{currentfill}%
\pgfsetfillopacity{0.950000}%
\pgfsetlinewidth{0.200750pt}%
\definecolor{currentstroke}{rgb}{0.200000,0.200000,0.200000}%
\pgfsetstrokecolor{currentstroke}%
\pgfsetdash{}{0pt}%
\pgfpathmoveto{\pgfqpoint{3.974907in}{1.963453in}}%
\pgfpathlineto{\pgfqpoint{3.946325in}{2.023602in}}%
\pgfpathlineto{\pgfqpoint{3.917847in}{2.085704in}}%
\pgfpathlineto{\pgfqpoint{3.964223in}{2.101563in}}%
\pgfpathlineto{\pgfqpoint{3.992712in}{2.039586in}}%
\pgfpathlineto{\pgfqpoint{4.021307in}{1.979545in}}%
\pgfpathlineto{\pgfqpoint{3.974907in}{1.963453in}}%
\pgfpathclose%
\pgfusepath{stroke,fill}%
\end{pgfscope}%
\begin{pgfscope}%
\pgfpathrectangle{\pgfqpoint{0.050000in}{0.083252in}}{\pgfqpoint{4.900000in}{4.900000in}}%
\pgfusepath{clip}%
\pgfsetbuttcap%
\pgfsetroundjoin%
\definecolor{currentfill}{rgb}{0.257793,0.405229,0.807505}%
\pgfsetfillcolor{currentfill}%
\pgfsetfillopacity{0.950000}%
\pgfsetlinewidth{0.200750pt}%
\definecolor{currentstroke}{rgb}{0.200000,0.200000,0.200000}%
\pgfsetstrokecolor{currentstroke}%
\pgfsetdash{}{0pt}%
\pgfpathmoveto{\pgfqpoint{2.497155in}{3.572775in}}%
\pgfpathlineto{\pgfqpoint{2.469681in}{3.646162in}}%
\pgfpathlineto{\pgfqpoint{2.442233in}{3.719481in}}%
\pgfpathlineto{\pgfqpoint{2.489920in}{3.732949in}}%
\pgfpathlineto{\pgfqpoint{2.537523in}{3.746394in}}%
\pgfpathlineto{\pgfqpoint{2.564967in}{3.673217in}}%
\pgfpathlineto{\pgfqpoint{2.592436in}{3.599973in}}%
\pgfpathlineto{\pgfqpoint{2.544838in}{3.586386in}}%
\pgfpathlineto{\pgfqpoint{2.497155in}{3.572775in}}%
\pgfpathclose%
\pgfusepath{stroke,fill}%
\end{pgfscope}%
\begin{pgfscope}%
\pgfpathrectangle{\pgfqpoint{0.050000in}{0.083252in}}{\pgfqpoint{4.900000in}{4.900000in}}%
\pgfusepath{clip}%
\pgfsetbuttcap%
\pgfsetroundjoin%
\definecolor{currentfill}{rgb}{0.270342,0.538562,0.984760}%
\pgfsetfillcolor{currentfill}%
\pgfsetfillopacity{0.950000}%
\pgfsetlinewidth{0.200750pt}%
\definecolor{currentstroke}{rgb}{0.200000,0.200000,0.200000}%
\pgfsetstrokecolor{currentstroke}%
\pgfsetdash{}{0pt}%
\pgfpathmoveto{\pgfqpoint{2.086448in}{4.105551in}}%
\pgfpathlineto{\pgfqpoint{2.059172in}{4.178675in}}%
\pgfpathlineto{\pgfqpoint{2.107209in}{4.191742in}}%
\pgfpathlineto{\pgfqpoint{2.155162in}{4.204786in}}%
\pgfpathlineto{\pgfqpoint{2.182435in}{4.131803in}}%
\pgfpathlineto{\pgfqpoint{2.134484in}{4.118689in}}%
\pgfpathlineto{\pgfqpoint{2.086448in}{4.105551in}}%
\pgfpathclose%
\pgfusepath{stroke,fill}%
\end{pgfscope}%
\begin{pgfscope}%
\pgfpathrectangle{\pgfqpoint{0.050000in}{0.083252in}}{\pgfqpoint{4.900000in}{4.900000in}}%
\pgfusepath{clip}%
\pgfsetbuttcap%
\pgfsetroundjoin%
\definecolor{currentfill}{rgb}{0.251396,0.337255,0.717140}%
\pgfsetfillcolor{currentfill}%
\pgfsetfillopacity{0.950000}%
\pgfsetlinewidth{0.200750pt}%
\definecolor{currentstroke}{rgb}{0.200000,0.200000,0.200000}%
\pgfsetstrokecolor{currentstroke}%
\pgfsetdash{}{0pt}%
\pgfpathmoveto{\pgfqpoint{2.702571in}{3.306327in}}%
\pgfpathlineto{\pgfqpoint{2.674999in}{3.379838in}}%
\pgfpathlineto{\pgfqpoint{2.647452in}{3.453283in}}%
\pgfpathlineto{\pgfqpoint{2.694963in}{3.466989in}}%
\pgfpathlineto{\pgfqpoint{2.742389in}{3.480672in}}%
\pgfpathlineto{\pgfqpoint{2.769932in}{3.407372in}}%
\pgfpathlineto{\pgfqpoint{2.797500in}{3.334008in}}%
\pgfpathlineto{\pgfqpoint{2.750078in}{3.320179in}}%
\pgfpathlineto{\pgfqpoint{2.702571in}{3.306327in}}%
\pgfpathclose%
\pgfusepath{stroke,fill}%
\end{pgfscope}%
\begin{pgfscope}%
\pgfpathrectangle{\pgfqpoint{0.050000in}{0.083252in}}{\pgfqpoint{4.900000in}{4.900000in}}%
\pgfusepath{clip}%
\pgfsetbuttcap%
\pgfsetroundjoin%
\definecolor{currentfill}{rgb}{0.244752,0.266667,0.623299}%
\pgfsetfillcolor{currentfill}%
\pgfsetfillopacity{0.950000}%
\pgfsetlinewidth{0.200750pt}%
\definecolor{currentstroke}{rgb}{0.200000,0.200000,0.200000}%
\pgfsetstrokecolor{currentstroke}%
\pgfsetdash{}{0pt}%
\pgfpathmoveto{\pgfqpoint{2.908031in}{3.039980in}}%
\pgfpathlineto{\pgfqpoint{2.880359in}{3.113563in}}%
\pgfpathlineto{\pgfqpoint{2.852714in}{3.187099in}}%
\pgfpathlineto{\pgfqpoint{2.900048in}{3.201054in}}%
\pgfpathlineto{\pgfqpoint{2.947299in}{3.214985in}}%
\pgfpathlineto{\pgfqpoint{2.974940in}{3.141604in}}%
\pgfpathlineto{\pgfqpoint{3.002608in}{3.068181in}}%
\pgfpathlineto{\pgfqpoint{2.955361in}{3.054092in}}%
\pgfpathlineto{\pgfqpoint{2.908031in}{3.039980in}}%
\pgfpathclose%
\pgfusepath{stroke,fill}%
\end{pgfscope}%
\begin{pgfscope}%
\pgfpathrectangle{\pgfqpoint{0.050000in}{0.083252in}}{\pgfqpoint{4.900000in}{4.900000in}}%
\pgfusepath{clip}%
\pgfsetbuttcap%
\pgfsetroundjoin%
\definecolor{currentfill}{rgb}{0.641984,0.627220,0.785928}%
\pgfsetfillcolor{currentfill}%
\pgfsetfillopacity{0.950000}%
\pgfsetlinewidth{0.200750pt}%
\definecolor{currentstroke}{rgb}{0.200000,0.200000,0.200000}%
\pgfsetstrokecolor{currentstroke}%
\pgfsetdash{}{0pt}%
\pgfpathmoveto{\pgfqpoint{3.674932in}{2.152376in}}%
\pgfpathlineto{\pgfqpoint{3.646739in}{2.219794in}}%
\pgfpathlineto{\pgfqpoint{3.618611in}{2.288393in}}%
\pgfpathlineto{\pgfqpoint{3.665279in}{2.303772in}}%
\pgfpathlineto{\pgfqpoint{3.711867in}{2.319125in}}%
\pgfpathlineto{\pgfqpoint{3.740006in}{2.250846in}}%
\pgfpathlineto{\pgfqpoint{3.768211in}{2.183746in}}%
\pgfpathlineto{\pgfqpoint{3.721612in}{2.168076in}}%
\pgfpathlineto{\pgfqpoint{3.674932in}{2.152376in}}%
\pgfpathclose%
\pgfusepath{stroke,fill}%
\end{pgfscope}%
\begin{pgfscope}%
\pgfpathrectangle{\pgfqpoint{0.050000in}{0.083252in}}{\pgfqpoint{4.900000in}{4.900000in}}%
\pgfusepath{clip}%
\pgfsetbuttcap%
\pgfsetroundjoin%
\definecolor{currentfill}{rgb}{0.266067,0.235802,0.561153}%
\pgfsetfillcolor{currentfill}%
\pgfsetfillopacity{0.950000}%
\pgfsetlinewidth{0.200750pt}%
\definecolor{currentstroke}{rgb}{0.200000,0.200000,0.200000}%
\pgfsetstrokecolor{currentstroke}%
\pgfsetdash{}{0pt}%
\pgfpathmoveto{\pgfqpoint{3.113544in}{2.774372in}}%
\pgfpathlineto{\pgfqpoint{3.085769in}{2.847799in}}%
\pgfpathlineto{\pgfqpoint{3.058022in}{2.921261in}}%
\pgfpathlineto{\pgfqpoint{3.105181in}{2.935499in}}%
\pgfpathlineto{\pgfqpoint{3.152257in}{2.949716in}}%
\pgfpathlineto{\pgfqpoint{3.180000in}{2.876440in}}%
\pgfpathlineto{\pgfqpoint{3.207772in}{2.803211in}}%
\pgfpathlineto{\pgfqpoint{3.160700in}{2.788802in}}%
\pgfpathlineto{\pgfqpoint{3.113544in}{2.774372in}}%
\pgfpathclose%
\pgfusepath{stroke,fill}%
\end{pgfscope}%
\begin{pgfscope}%
\pgfpathrectangle{\pgfqpoint{0.050000in}{0.083252in}}{\pgfqpoint{4.900000in}{4.900000in}}%
\pgfusepath{clip}%
\pgfsetbuttcap%
\pgfsetroundjoin%
\definecolor{currentfill}{rgb}{0.421207,0.397339,0.653918}%
\pgfsetfillcolor{currentfill}%
\pgfsetfillopacity{0.950000}%
\pgfsetlinewidth{0.200750pt}%
\definecolor{currentstroke}{rgb}{0.200000,0.200000,0.200000}%
\pgfsetstrokecolor{currentstroke}%
\pgfsetdash{}{0pt}%
\pgfpathmoveto{\pgfqpoint{3.319157in}{2.511715in}}%
\pgfpathlineto{\pgfqpoint{3.291263in}{2.584247in}}%
\pgfpathlineto{\pgfqpoint{3.263402in}{2.657057in}}%
\pgfpathlineto{\pgfqpoint{3.310389in}{2.671666in}}%
\pgfpathlineto{\pgfqpoint{3.357294in}{2.686255in}}%
\pgfpathlineto{\pgfqpoint{3.385154in}{2.613691in}}%
\pgfpathlineto{\pgfqpoint{3.413048in}{2.541422in}}%
\pgfpathlineto{\pgfqpoint{3.366144in}{2.526579in}}%
\pgfpathlineto{\pgfqpoint{3.319157in}{2.511715in}}%
\pgfpathclose%
\pgfusepath{stroke,fill}%
\end{pgfscope}%
\begin{pgfscope}%
\pgfpathrectangle{\pgfqpoint{0.050000in}{0.083252in}}{\pgfqpoint{4.900000in}{4.900000in}}%
\pgfusepath{clip}%
\pgfsetbuttcap%
\pgfsetroundjoin%
\definecolor{currentfill}{rgb}{0.267636,0.509804,0.946528}%
\pgfsetfillcolor{currentfill}%
\pgfsetfillopacity{0.950000}%
\pgfsetlinewidth{0.200750pt}%
\definecolor{currentstroke}{rgb}{0.200000,0.200000,0.200000}%
\pgfsetstrokecolor{currentstroke}%
\pgfsetdash{}{0pt}%
\pgfpathmoveto{\pgfqpoint{2.141077in}{3.959098in}}%
\pgfpathlineto{\pgfqpoint{2.113750in}{4.032359in}}%
\pgfpathlineto{\pgfqpoint{2.086448in}{4.105551in}}%
\pgfpathlineto{\pgfqpoint{2.134484in}{4.118689in}}%
\pgfpathlineto{\pgfqpoint{2.182435in}{4.131803in}}%
\pgfpathlineto{\pgfqpoint{2.209733in}{4.058752in}}%
\pgfpathlineto{\pgfqpoint{2.237056in}{3.985633in}}%
\pgfpathlineto{\pgfqpoint{2.189109in}{3.972377in}}%
\pgfpathlineto{\pgfqpoint{2.141077in}{3.959098in}}%
\pgfpathclose%
\pgfusepath{stroke,fill}%
\end{pgfscope}%
\begin{pgfscope}%
\pgfpathrectangle{\pgfqpoint{0.050000in}{0.083252in}}{\pgfqpoint{4.900000in}{4.900000in}}%
\pgfusepath{clip}%
\pgfsetbuttcap%
\pgfsetroundjoin%
\definecolor{currentfill}{rgb}{0.570381,0.552664,0.743114}%
\pgfsetfillcolor{currentfill}%
\pgfsetfillopacity{0.950000}%
\pgfsetlinewidth{0.200750pt}%
\definecolor{currentstroke}{rgb}{0.200000,0.200000,0.200000}%
\pgfsetstrokecolor{currentstroke}%
\pgfsetdash{}{0pt}%
\pgfpathmoveto{\pgfqpoint{3.525030in}{2.257557in}}%
\pgfpathlineto{\pgfqpoint{3.496966in}{2.327468in}}%
\pgfpathlineto{\pgfqpoint{3.468951in}{2.398184in}}%
\pgfpathlineto{\pgfqpoint{3.515775in}{2.413297in}}%
\pgfpathlineto{\pgfqpoint{3.562519in}{2.428387in}}%
\pgfpathlineto{\pgfqpoint{3.590539in}{2.357981in}}%
\pgfpathlineto{\pgfqpoint{3.618611in}{2.288393in}}%
\pgfpathlineto{\pgfqpoint{3.571861in}{2.272988in}}%
\pgfpathlineto{\pgfqpoint{3.525030in}{2.257557in}}%
\pgfpathclose%
\pgfusepath{stroke,fill}%
\end{pgfscope}%
\begin{pgfscope}%
\pgfpathrectangle{\pgfqpoint{0.050000in}{0.083252in}}{\pgfqpoint{4.900000in}{4.900000in}}%
\pgfusepath{clip}%
\pgfsetbuttcap%
\pgfsetroundjoin%
\definecolor{currentfill}{rgb}{0.261238,0.441830,0.856163}%
\pgfsetfillcolor{currentfill}%
\pgfsetfillopacity{0.950000}%
\pgfsetlinewidth{0.200750pt}%
\definecolor{currentstroke}{rgb}{0.200000,0.200000,0.200000}%
\pgfsetstrokecolor{currentstroke}%
\pgfsetdash{}{0pt}%
\pgfpathmoveto{\pgfqpoint{2.346607in}{3.692472in}}%
\pgfpathlineto{\pgfqpoint{2.319181in}{3.765865in}}%
\pgfpathlineto{\pgfqpoint{2.291780in}{3.839190in}}%
\pgfpathlineto{\pgfqpoint{2.339640in}{3.852564in}}%
\pgfpathlineto{\pgfqpoint{2.387414in}{3.865914in}}%
\pgfpathlineto{\pgfqpoint{2.414811in}{3.792732in}}%
\pgfpathlineto{\pgfqpoint{2.442233in}{3.719481in}}%
\pgfpathlineto{\pgfqpoint{2.394462in}{3.705989in}}%
\pgfpathlineto{\pgfqpoint{2.346607in}{3.692472in}}%
\pgfpathclose%
\pgfusepath{stroke,fill}%
\end{pgfscope}%
\begin{pgfscope}%
\pgfpathrectangle{\pgfqpoint{0.050000in}{0.083252in}}{\pgfqpoint{4.900000in}{4.900000in}}%
\pgfusepath{clip}%
\pgfsetbuttcap%
\pgfsetroundjoin%
\definecolor{currentfill}{rgb}{0.254594,0.371242,0.762322}%
\pgfsetfillcolor{currentfill}%
\pgfsetfillopacity{0.950000}%
\pgfsetlinewidth{0.200750pt}%
\definecolor{currentstroke}{rgb}{0.200000,0.200000,0.200000}%
\pgfsetstrokecolor{currentstroke}%
\pgfsetdash{}{0pt}%
\pgfpathmoveto{\pgfqpoint{2.552179in}{3.425797in}}%
\pgfpathlineto{\pgfqpoint{2.524654in}{3.499320in}}%
\pgfpathlineto{\pgfqpoint{2.497155in}{3.572775in}}%
\pgfpathlineto{\pgfqpoint{2.544838in}{3.586386in}}%
\pgfpathlineto{\pgfqpoint{2.592436in}{3.599973in}}%
\pgfpathlineto{\pgfqpoint{2.619931in}{3.526661in}}%
\pgfpathlineto{\pgfqpoint{2.647452in}{3.453283in}}%
\pgfpathlineto{\pgfqpoint{2.599858in}{3.439552in}}%
\pgfpathlineto{\pgfqpoint{2.552179in}{3.425797in}}%
\pgfpathclose%
\pgfusepath{stroke,fill}%
\end{pgfscope}%
\begin{pgfscope}%
\pgfpathrectangle{\pgfqpoint{0.050000in}{0.083252in}}{\pgfqpoint{4.900000in}{4.900000in}}%
\pgfusepath{clip}%
\pgfsetbuttcap%
\pgfsetroundjoin%
\definecolor{currentfill}{rgb}{0.247951,0.300654,0.668481}%
\pgfsetfillcolor{currentfill}%
\pgfsetfillopacity{0.950000}%
\pgfsetlinewidth{0.200750pt}%
\definecolor{currentstroke}{rgb}{0.200000,0.200000,0.200000}%
\pgfsetstrokecolor{currentstroke}%
\pgfsetdash{}{0pt}%
\pgfpathmoveto{\pgfqpoint{2.757794in}{3.159120in}}%
\pgfpathlineto{\pgfqpoint{2.730170in}{3.232754in}}%
\pgfpathlineto{\pgfqpoint{2.702571in}{3.306327in}}%
\pgfpathlineto{\pgfqpoint{2.750078in}{3.320179in}}%
\pgfpathlineto{\pgfqpoint{2.797500in}{3.334008in}}%
\pgfpathlineto{\pgfqpoint{2.825094in}{3.260583in}}%
\pgfpathlineto{\pgfqpoint{2.852714in}{3.187099in}}%
\pgfpathlineto{\pgfqpoint{2.805296in}{3.173121in}}%
\pgfpathlineto{\pgfqpoint{2.757794in}{3.159120in}}%
\pgfpathclose%
\pgfusepath{stroke,fill}%
\end{pgfscope}%
\begin{pgfscope}%
\pgfpathrectangle{\pgfqpoint{0.050000in}{0.083252in}}{\pgfqpoint{4.900000in}{4.900000in}}%
\pgfusepath{clip}%
\pgfsetbuttcap%
\pgfsetroundjoin%
\definecolor{currentfill}{rgb}{0.241553,0.232680,0.578116}%
\pgfsetfillcolor{currentfill}%
\pgfsetfillopacity{0.950000}%
\pgfsetlinewidth{0.200750pt}%
\definecolor{currentstroke}{rgb}{0.200000,0.200000,0.200000}%
\pgfsetstrokecolor{currentstroke}%
\pgfsetdash{}{0pt}%
\pgfpathmoveto{\pgfqpoint{2.963454in}{2.892719in}}%
\pgfpathlineto{\pgfqpoint{2.935729in}{2.966361in}}%
\pgfpathlineto{\pgfqpoint{2.908031in}{3.039980in}}%
\pgfpathlineto{\pgfqpoint{2.955361in}{3.054092in}}%
\pgfpathlineto{\pgfqpoint{3.002608in}{3.068181in}}%
\pgfpathlineto{\pgfqpoint{3.030301in}{2.994728in}}%
\pgfpathlineto{\pgfqpoint{3.058022in}{2.921261in}}%
\pgfpathlineto{\pgfqpoint{3.010780in}{2.907001in}}%
\pgfpathlineto{\pgfqpoint{2.963454in}{2.892719in}}%
\pgfpathclose%
\pgfusepath{stroke,fill}%
\end{pgfscope}%
\begin{pgfscope}%
\pgfpathrectangle{\pgfqpoint{0.050000in}{0.083252in}}{\pgfqpoint{4.900000in}{4.900000in}}%
\pgfusepath{clip}%
\pgfsetbuttcap%
\pgfsetroundjoin%
\definecolor{currentfill}{rgb}{0.773256,0.763906,0.864421}%
\pgfsetfillcolor{currentfill}%
\pgfsetfillopacity{0.950000}%
\pgfsetlinewidth{0.200750pt}%
\definecolor{currentstroke}{rgb}{0.200000,0.200000,0.200000}%
\pgfsetstrokecolor{currentstroke}%
\pgfsetdash{}{0pt}%
\pgfpathmoveto{\pgfqpoint{3.881866in}{1.931148in}}%
\pgfpathlineto{\pgfqpoint{3.853309in}{1.991519in}}%
\pgfpathlineto{\pgfqpoint{3.824854in}{2.053877in}}%
\pgfpathlineto{\pgfqpoint{3.871391in}{2.069809in}}%
\pgfpathlineto{\pgfqpoint{3.917847in}{2.085704in}}%
\pgfpathlineto{\pgfqpoint{3.946325in}{2.023602in}}%
\pgfpathlineto{\pgfqpoint{3.974907in}{1.963453in}}%
\pgfpathlineto{\pgfqpoint{3.928427in}{1.947321in}}%
\pgfpathlineto{\pgfqpoint{3.881866in}{1.931148in}}%
\pgfpathclose%
\pgfusepath{stroke,fill}%
\end{pgfscope}%
\begin{pgfscope}%
\pgfpathrectangle{\pgfqpoint{0.050000in}{0.083252in}}{\pgfqpoint{4.900000in}{4.900000in}}%
\pgfusepath{clip}%
\pgfsetbuttcap%
\pgfsetroundjoin%
\definecolor{currentfill}{rgb}{0.343637,0.316571,0.607536}%
\pgfsetfillcolor{currentfill}%
\pgfsetfillopacity{0.950000}%
\pgfsetlinewidth{0.200750pt}%
\definecolor{currentstroke}{rgb}{0.200000,0.200000,0.200000}%
\pgfsetstrokecolor{currentstroke}%
\pgfsetdash{}{0pt}%
\pgfpathmoveto{\pgfqpoint{3.169180in}{2.627780in}}%
\pgfpathlineto{\pgfqpoint{3.141347in}{2.701017in}}%
\pgfpathlineto{\pgfqpoint{3.113544in}{2.774372in}}%
\pgfpathlineto{\pgfqpoint{3.160700in}{2.788802in}}%
\pgfpathlineto{\pgfqpoint{3.207772in}{2.803211in}}%
\pgfpathlineto{\pgfqpoint{3.235572in}{2.730066in}}%
\pgfpathlineto{\pgfqpoint{3.263402in}{2.657057in}}%
\pgfpathlineto{\pgfqpoint{3.216332in}{2.642428in}}%
\pgfpathlineto{\pgfqpoint{3.169180in}{2.627780in}}%
\pgfpathclose%
\pgfusepath{stroke,fill}%
\end{pgfscope}%
\begin{pgfscope}%
\pgfpathrectangle{\pgfqpoint{0.050000in}{0.083252in}}{\pgfqpoint{4.900000in}{4.900000in}}%
\pgfusepath{clip}%
\pgfsetbuttcap%
\pgfsetroundjoin%
\definecolor{currentfill}{rgb}{0.713587,0.701776,0.828743}%
\pgfsetfillcolor{currentfill}%
\pgfsetfillopacity{0.950000}%
\pgfsetlinewidth{0.200750pt}%
\definecolor{currentstroke}{rgb}{0.200000,0.200000,0.200000}%
\pgfsetstrokecolor{currentstroke}%
\pgfsetdash{}{0pt}%
\pgfpathmoveto{\pgfqpoint{3.731539in}{2.021904in}}%
\pgfpathlineto{\pgfqpoint{3.703195in}{2.086341in}}%
\pgfpathlineto{\pgfqpoint{3.674932in}{2.152376in}}%
\pgfpathlineto{\pgfqpoint{3.721612in}{2.168076in}}%
\pgfpathlineto{\pgfqpoint{3.768211in}{2.183746in}}%
\pgfpathlineto{\pgfqpoint{3.796491in}{2.118021in}}%
\pgfpathlineto{\pgfqpoint{3.824854in}{2.053877in}}%
\pgfpathlineto{\pgfqpoint{3.778237in}{2.037909in}}%
\pgfpathlineto{\pgfqpoint{3.731539in}{2.021904in}}%
\pgfpathclose%
\pgfusepath{stroke,fill}%
\end{pgfscope}%
\begin{pgfscope}%
\pgfpathrectangle{\pgfqpoint{0.050000in}{0.083252in}}{\pgfqpoint{4.900000in}{4.900000in}}%
\pgfusepath{clip}%
\pgfsetbuttcap%
\pgfsetroundjoin%
\definecolor{currentfill}{rgb}{0.498777,0.478108,0.700300}%
\pgfsetfillcolor{currentfill}%
\pgfsetfillopacity{0.950000}%
\pgfsetlinewidth{0.200750pt}%
\definecolor{currentstroke}{rgb}{0.200000,0.200000,0.200000}%
\pgfsetstrokecolor{currentstroke}%
\pgfsetdash{}{0pt}%
\pgfpathmoveto{\pgfqpoint{3.375055in}{2.367893in}}%
\pgfpathlineto{\pgfqpoint{3.347087in}{2.439558in}}%
\pgfpathlineto{\pgfqpoint{3.319157in}{2.511715in}}%
\pgfpathlineto{\pgfqpoint{3.366144in}{2.526579in}}%
\pgfpathlineto{\pgfqpoint{3.413048in}{2.541422in}}%
\pgfpathlineto{\pgfqpoint{3.440979in}{2.469548in}}%
\pgfpathlineto{\pgfqpoint{3.468951in}{2.398184in}}%
\pgfpathlineto{\pgfqpoint{3.422044in}{2.383050in}}%
\pgfpathlineto{\pgfqpoint{3.375055in}{2.367893in}}%
\pgfpathclose%
\pgfusepath{stroke,fill}%
\end{pgfscope}%
\begin{pgfscope}%
\pgfpathrectangle{\pgfqpoint{0.050000in}{0.083252in}}{\pgfqpoint{4.900000in}{4.900000in}}%
\pgfusepath{clip}%
\pgfsetbuttcap%
\pgfsetroundjoin%
\definecolor{currentfill}{rgb}{0.826959,0.819823,0.896532}%
\pgfsetfillcolor{currentfill}%
\pgfsetfillopacity{0.950000}%
\pgfsetlinewidth{0.200750pt}%
\definecolor{currentstroke}{rgb}{0.200000,0.200000,0.200000}%
\pgfsetstrokecolor{currentstroke}%
\pgfsetdash{}{0pt}%
\pgfpathmoveto{\pgfqpoint{4.032428in}{1.849706in}}%
\pgfpathlineto{\pgfqpoint{4.003605in}{1.905436in}}%
\pgfpathlineto{\pgfqpoint{3.974907in}{1.963453in}}%
\pgfpathlineto{\pgfqpoint{4.021307in}{1.979545in}}%
\pgfpathlineto{\pgfqpoint{4.050020in}{1.921614in}}%
\pgfpathlineto{\pgfqpoint{4.078859in}{1.865948in}}%
\pgfpathlineto{\pgfqpoint{4.032428in}{1.849706in}}%
\pgfpathclose%
\pgfusepath{stroke,fill}%
\end{pgfscope}%
\begin{pgfscope}%
\pgfpathrectangle{\pgfqpoint{0.050000in}{0.083252in}}{\pgfqpoint{4.900000in}{4.900000in}}%
\pgfusepath{clip}%
\pgfsetbuttcap%
\pgfsetroundjoin%
\definecolor{currentfill}{rgb}{0.264437,0.475817,0.901346}%
\pgfsetfillcolor{currentfill}%
\pgfsetfillopacity{0.950000}%
\pgfsetlinewidth{0.200750pt}%
\definecolor{currentstroke}{rgb}{0.200000,0.200000,0.200000}%
\pgfsetstrokecolor{currentstroke}%
\pgfsetdash{}{0pt}%
\pgfpathmoveto{\pgfqpoint{2.195808in}{3.812370in}}%
\pgfpathlineto{\pgfqpoint{2.168429in}{3.885768in}}%
\pgfpathlineto{\pgfqpoint{2.141077in}{3.959098in}}%
\pgfpathlineto{\pgfqpoint{2.189109in}{3.972377in}}%
\pgfpathlineto{\pgfqpoint{2.237056in}{3.985633in}}%
\pgfpathlineto{\pgfqpoint{2.264406in}{3.912446in}}%
\pgfpathlineto{\pgfqpoint{2.291780in}{3.839190in}}%
\pgfpathlineto{\pgfqpoint{2.243837in}{3.825792in}}%
\pgfpathlineto{\pgfqpoint{2.195808in}{3.812370in}}%
\pgfpathclose%
\pgfusepath{stroke,fill}%
\end{pgfscope}%
\begin{pgfscope}%
\pgfpathrectangle{\pgfqpoint{0.050000in}{0.083252in}}{\pgfqpoint{4.900000in}{4.900000in}}%
\pgfusepath{clip}%
\pgfsetbuttcap%
\pgfsetroundjoin%
\definecolor{currentfill}{rgb}{0.257793,0.405229,0.807505}%
\pgfsetfillcolor{currentfill}%
\pgfsetfillopacity{0.950000}%
\pgfsetlinewidth{0.200750pt}%
\definecolor{currentstroke}{rgb}{0.200000,0.200000,0.200000}%
\pgfsetstrokecolor{currentstroke}%
\pgfsetdash{}{0pt}%
\pgfpathmoveto{\pgfqpoint{2.401536in}{3.545481in}}%
\pgfpathlineto{\pgfqpoint{2.374059in}{3.619011in}}%
\pgfpathlineto{\pgfqpoint{2.346607in}{3.692472in}}%
\pgfpathlineto{\pgfqpoint{2.394462in}{3.705989in}}%
\pgfpathlineto{\pgfqpoint{2.442233in}{3.719481in}}%
\pgfpathlineto{\pgfqpoint{2.469681in}{3.646162in}}%
\pgfpathlineto{\pgfqpoint{2.497155in}{3.572775in}}%
\pgfpathlineto{\pgfqpoint{2.449388in}{3.559140in}}%
\pgfpathlineto{\pgfqpoint{2.401536in}{3.545481in}}%
\pgfpathclose%
\pgfusepath{stroke,fill}%
\end{pgfscope}%
\begin{pgfscope}%
\pgfpathrectangle{\pgfqpoint{0.050000in}{0.083252in}}{\pgfqpoint{4.900000in}{4.900000in}}%
\pgfusepath{clip}%
\pgfsetbuttcap%
\pgfsetroundjoin%
\definecolor{currentfill}{rgb}{0.270096,0.535948,0.981284}%
\pgfsetfillcolor{currentfill}%
\pgfsetfillopacity{0.950000}%
\pgfsetlinewidth{0.200750pt}%
\definecolor{currentstroke}{rgb}{0.200000,0.200000,0.200000}%
\pgfsetstrokecolor{currentstroke}%
\pgfsetdash{}{0pt}%
\pgfpathmoveto{\pgfqpoint{1.990122in}{4.079206in}}%
\pgfpathlineto{\pgfqpoint{1.962842in}{4.152472in}}%
\pgfpathlineto{\pgfqpoint{2.011050in}{4.165585in}}%
\pgfpathlineto{\pgfqpoint{2.059172in}{4.178675in}}%
\pgfpathlineto{\pgfqpoint{2.086448in}{4.105551in}}%
\pgfpathlineto{\pgfqpoint{2.038327in}{4.092390in}}%
\pgfpathlineto{\pgfqpoint{1.990122in}{4.079206in}}%
\pgfpathclose%
\pgfusepath{stroke,fill}%
\end{pgfscope}%
\begin{pgfscope}%
\pgfpathrectangle{\pgfqpoint{0.050000in}{0.083252in}}{\pgfqpoint{4.900000in}{4.900000in}}%
\pgfusepath{clip}%
\pgfsetbuttcap%
\pgfsetroundjoin%
\definecolor{currentfill}{rgb}{0.251396,0.337255,0.717140}%
\pgfsetfillcolor{currentfill}%
\pgfsetfillopacity{0.950000}%
\pgfsetlinewidth{0.200750pt}%
\definecolor{currentstroke}{rgb}{0.200000,0.200000,0.200000}%
\pgfsetstrokecolor{currentstroke}%
\pgfsetdash{}{0pt}%
\pgfpathmoveto{\pgfqpoint{2.607307in}{3.278551in}}%
\pgfpathlineto{\pgfqpoint{2.579730in}{3.352207in}}%
\pgfpathlineto{\pgfqpoint{2.552179in}{3.425797in}}%
\pgfpathlineto{\pgfqpoint{2.599858in}{3.439552in}}%
\pgfpathlineto{\pgfqpoint{2.647452in}{3.453283in}}%
\pgfpathlineto{\pgfqpoint{2.674999in}{3.379838in}}%
\pgfpathlineto{\pgfqpoint{2.702571in}{3.306327in}}%
\pgfpathlineto{\pgfqpoint{2.654981in}{3.292451in}}%
\pgfpathlineto{\pgfqpoint{2.607307in}{3.278551in}}%
\pgfpathclose%
\pgfusepath{stroke,fill}%
\end{pgfscope}%
\begin{pgfscope}%
\pgfpathrectangle{\pgfqpoint{0.050000in}{0.083252in}}{\pgfqpoint{4.900000in}{4.900000in}}%
\pgfusepath{clip}%
\pgfsetbuttcap%
\pgfsetroundjoin%
\definecolor{currentfill}{rgb}{0.244752,0.266667,0.623299}%
\pgfsetfillcolor{currentfill}%
\pgfsetfillopacity{0.950000}%
\pgfsetlinewidth{0.200750pt}%
\definecolor{currentstroke}{rgb}{0.200000,0.200000,0.200000}%
\pgfsetstrokecolor{currentstroke}%
\pgfsetdash{}{0pt}%
\pgfpathmoveto{\pgfqpoint{2.813120in}{3.011688in}}%
\pgfpathlineto{\pgfqpoint{2.785444in}{3.085429in}}%
\pgfpathlineto{\pgfqpoint{2.757794in}{3.159120in}}%
\pgfpathlineto{\pgfqpoint{2.805296in}{3.173121in}}%
\pgfpathlineto{\pgfqpoint{2.852714in}{3.187099in}}%
\pgfpathlineto{\pgfqpoint{2.880359in}{3.113563in}}%
\pgfpathlineto{\pgfqpoint{2.908031in}{3.039980in}}%
\pgfpathlineto{\pgfqpoint{2.860618in}{3.025845in}}%
\pgfpathlineto{\pgfqpoint{2.813120in}{3.011688in}}%
\pgfpathclose%
\pgfusepath{stroke,fill}%
\end{pgfscope}%
\begin{pgfscope}%
\pgfpathrectangle{\pgfqpoint{0.050000in}{0.083252in}}{\pgfqpoint{4.900000in}{4.900000in}}%
\pgfusepath{clip}%
\pgfsetbuttcap%
\pgfsetroundjoin%
\definecolor{currentfill}{rgb}{0.647951,0.633433,0.789496}%
\pgfsetfillcolor{currentfill}%
\pgfsetfillopacity{0.950000}%
\pgfsetlinewidth{0.200750pt}%
\definecolor{currentstroke}{rgb}{0.200000,0.200000,0.200000}%
\pgfsetstrokecolor{currentstroke}%
\pgfsetdash{}{0pt}%
\pgfpathmoveto{\pgfqpoint{3.581328in}{2.120881in}}%
\pgfpathlineto{\pgfqpoint{3.553148in}{2.188629in}}%
\pgfpathlineto{\pgfqpoint{3.525030in}{2.257557in}}%
\pgfpathlineto{\pgfqpoint{3.571861in}{2.272988in}}%
\pgfpathlineto{\pgfqpoint{3.618611in}{2.288393in}}%
\pgfpathlineto{\pgfqpoint{3.646739in}{2.219794in}}%
\pgfpathlineto{\pgfqpoint{3.674932in}{2.152376in}}%
\pgfpathlineto{\pgfqpoint{3.628170in}{2.136644in}}%
\pgfpathlineto{\pgfqpoint{3.581328in}{2.120881in}}%
\pgfpathclose%
\pgfusepath{stroke,fill}%
\end{pgfscope}%
\begin{pgfscope}%
\pgfpathrectangle{\pgfqpoint{0.050000in}{0.083252in}}{\pgfqpoint{4.900000in}{4.900000in}}%
\pgfusepath{clip}%
\pgfsetbuttcap%
\pgfsetroundjoin%
\definecolor{currentfill}{rgb}{0.266067,0.235802,0.561153}%
\pgfsetfillcolor{currentfill}%
\pgfsetfillopacity{0.950000}%
\pgfsetlinewidth{0.200750pt}%
\definecolor{currentstroke}{rgb}{0.200000,0.200000,0.200000}%
\pgfsetstrokecolor{currentstroke}%
\pgfsetdash{}{0pt}%
\pgfpathmoveto{\pgfqpoint{3.018984in}{2.745452in}}%
\pgfpathlineto{\pgfqpoint{2.991205in}{2.819075in}}%
\pgfpathlineto{\pgfqpoint{2.963454in}{2.892719in}}%
\pgfpathlineto{\pgfqpoint{3.010780in}{2.907001in}}%
\pgfpathlineto{\pgfqpoint{3.058022in}{2.921261in}}%
\pgfpathlineto{\pgfqpoint{3.085769in}{2.847799in}}%
\pgfpathlineto{\pgfqpoint{3.113544in}{2.774372in}}%
\pgfpathlineto{\pgfqpoint{3.066306in}{2.759922in}}%
\pgfpathlineto{\pgfqpoint{3.018984in}{2.745452in}}%
\pgfpathclose%
\pgfusepath{stroke,fill}%
\end{pgfscope}%
\begin{pgfscope}%
\pgfpathrectangle{\pgfqpoint{0.050000in}{0.083252in}}{\pgfqpoint{4.900000in}{4.900000in}}%
\pgfusepath{clip}%
\pgfsetbuttcap%
\pgfsetroundjoin%
\definecolor{currentfill}{rgb}{0.421207,0.397339,0.653918}%
\pgfsetfillcolor{currentfill}%
\pgfsetfillopacity{0.950000}%
\pgfsetlinewidth{0.200750pt}%
\definecolor{currentstroke}{rgb}{0.200000,0.200000,0.200000}%
\pgfsetstrokecolor{currentstroke}%
\pgfsetdash{}{0pt}%
\pgfpathmoveto{\pgfqpoint{3.224937in}{2.481928in}}%
\pgfpathlineto{\pgfqpoint{3.197042in}{2.554725in}}%
\pgfpathlineto{\pgfqpoint{3.169180in}{2.627780in}}%
\pgfpathlineto{\pgfqpoint{3.216332in}{2.642428in}}%
\pgfpathlineto{\pgfqpoint{3.263402in}{2.657057in}}%
\pgfpathlineto{\pgfqpoint{3.291263in}{2.584247in}}%
\pgfpathlineto{\pgfqpoint{3.319157in}{2.511715in}}%
\pgfpathlineto{\pgfqpoint{3.272088in}{2.496831in}}%
\pgfpathlineto{\pgfqpoint{3.224937in}{2.481928in}}%
\pgfpathclose%
\pgfusepath{stroke,fill}%
\end{pgfscope}%
\begin{pgfscope}%
\pgfpathrectangle{\pgfqpoint{0.050000in}{0.083252in}}{\pgfqpoint{4.900000in}{4.900000in}}%
\pgfusepath{clip}%
\pgfsetbuttcap%
\pgfsetroundjoin%
\definecolor{currentfill}{rgb}{0.267636,0.509804,0.946528}%
\pgfsetfillcolor{currentfill}%
\pgfsetfillopacity{0.950000}%
\pgfsetlinewidth{0.200750pt}%
\definecolor{currentstroke}{rgb}{0.200000,0.200000,0.200000}%
\pgfsetstrokecolor{currentstroke}%
\pgfsetdash{}{0pt}%
\pgfpathmoveto{\pgfqpoint{2.044757in}{3.932468in}}%
\pgfpathlineto{\pgfqpoint{2.017426in}{4.005871in}}%
\pgfpathlineto{\pgfqpoint{1.990122in}{4.079206in}}%
\pgfpathlineto{\pgfqpoint{2.038327in}{4.092390in}}%
\pgfpathlineto{\pgfqpoint{2.086448in}{4.105551in}}%
\pgfpathlineto{\pgfqpoint{2.113750in}{4.032359in}}%
\pgfpathlineto{\pgfqpoint{2.141077in}{3.959098in}}%
\pgfpathlineto{\pgfqpoint{2.092959in}{3.945795in}}%
\pgfpathlineto{\pgfqpoint{2.044757in}{3.932468in}}%
\pgfpathclose%
\pgfusepath{stroke,fill}%
\end{pgfscope}%
\begin{pgfscope}%
\pgfpathrectangle{\pgfqpoint{0.050000in}{0.083252in}}{\pgfqpoint{4.900000in}{4.900000in}}%
\pgfusepath{clip}%
\pgfsetbuttcap%
\pgfsetroundjoin%
\definecolor{currentfill}{rgb}{0.576348,0.558877,0.746682}%
\pgfsetfillcolor{currentfill}%
\pgfsetfillopacity{0.950000}%
\pgfsetlinewidth{0.200750pt}%
\definecolor{currentstroke}{rgb}{0.200000,0.200000,0.200000}%
\pgfsetstrokecolor{currentstroke}%
\pgfsetdash{}{0pt}%
\pgfpathmoveto{\pgfqpoint{3.431122in}{2.226618in}}%
\pgfpathlineto{\pgfqpoint{3.403066in}{2.296859in}}%
\pgfpathlineto{\pgfqpoint{3.375055in}{2.367893in}}%
\pgfpathlineto{\pgfqpoint{3.422044in}{2.383050in}}%
\pgfpathlineto{\pgfqpoint{3.468951in}{2.398184in}}%
\pgfpathlineto{\pgfqpoint{3.496966in}{2.327468in}}%
\pgfpathlineto{\pgfqpoint{3.525030in}{2.257557in}}%
\pgfpathlineto{\pgfqpoint{3.478117in}{2.242100in}}%
\pgfpathlineto{\pgfqpoint{3.431122in}{2.226618in}}%
\pgfpathclose%
\pgfusepath{stroke,fill}%
\end{pgfscope}%
\begin{pgfscope}%
\pgfpathrectangle{\pgfqpoint{0.050000in}{0.083252in}}{\pgfqpoint{4.900000in}{4.900000in}}%
\pgfusepath{clip}%
\pgfsetbuttcap%
\pgfsetroundjoin%
\definecolor{currentfill}{rgb}{0.261238,0.441830,0.856163}%
\pgfsetfillcolor{currentfill}%
\pgfsetfillopacity{0.950000}%
\pgfsetlinewidth{0.200750pt}%
\definecolor{currentstroke}{rgb}{0.200000,0.200000,0.200000}%
\pgfsetstrokecolor{currentstroke}%
\pgfsetdash{}{0pt}%
\pgfpathmoveto{\pgfqpoint{2.250641in}{3.665367in}}%
\pgfpathlineto{\pgfqpoint{2.223212in}{3.738903in}}%
\pgfpathlineto{\pgfqpoint{2.195808in}{3.812370in}}%
\pgfpathlineto{\pgfqpoint{2.243837in}{3.825792in}}%
\pgfpathlineto{\pgfqpoint{2.291780in}{3.839190in}}%
\pgfpathlineto{\pgfqpoint{2.319181in}{3.765865in}}%
\pgfpathlineto{\pgfqpoint{2.346607in}{3.692472in}}%
\pgfpathlineto{\pgfqpoint{2.298667in}{3.678932in}}%
\pgfpathlineto{\pgfqpoint{2.250641in}{3.665367in}}%
\pgfpathclose%
\pgfusepath{stroke,fill}%
\end{pgfscope}%
\begin{pgfscope}%
\pgfpathrectangle{\pgfqpoint{0.050000in}{0.083252in}}{\pgfqpoint{4.900000in}{4.900000in}}%
\pgfusepath{clip}%
\pgfsetbuttcap%
\pgfsetroundjoin%
\definecolor{currentfill}{rgb}{0.254594,0.371242,0.762322}%
\pgfsetfillcolor{currentfill}%
\pgfsetfillopacity{0.950000}%
\pgfsetlinewidth{0.200750pt}%
\definecolor{currentstroke}{rgb}{0.200000,0.200000,0.200000}%
\pgfsetstrokecolor{currentstroke}%
\pgfsetdash{}{0pt}%
\pgfpathmoveto{\pgfqpoint{2.456568in}{3.398215in}}%
\pgfpathlineto{\pgfqpoint{2.429039in}{3.471882in}}%
\pgfpathlineto{\pgfqpoint{2.401536in}{3.545481in}}%
\pgfpathlineto{\pgfqpoint{2.449388in}{3.559140in}}%
\pgfpathlineto{\pgfqpoint{2.497155in}{3.572775in}}%
\pgfpathlineto{\pgfqpoint{2.524654in}{3.499320in}}%
\pgfpathlineto{\pgfqpoint{2.552179in}{3.425797in}}%
\pgfpathlineto{\pgfqpoint{2.504416in}{3.412018in}}%
\pgfpathlineto{\pgfqpoint{2.456568in}{3.398215in}}%
\pgfpathclose%
\pgfusepath{stroke,fill}%
\end{pgfscope}%
\begin{pgfscope}%
\pgfpathrectangle{\pgfqpoint{0.050000in}{0.083252in}}{\pgfqpoint{4.900000in}{4.900000in}}%
\pgfusepath{clip}%
\pgfsetbuttcap%
\pgfsetroundjoin%
\definecolor{currentfill}{rgb}{0.826959,0.819823,0.896532}%
\pgfsetfillcolor{currentfill}%
\pgfsetfillopacity{0.950000}%
\pgfsetlinewidth{0.200750pt}%
\definecolor{currentstroke}{rgb}{0.200000,0.200000,0.200000}%
\pgfsetstrokecolor{currentstroke}%
\pgfsetdash{}{0pt}%
\pgfpathmoveto{\pgfqpoint{3.939326in}{1.817093in}}%
\pgfpathlineto{\pgfqpoint{3.910535in}{1.872952in}}%
\pgfpathlineto{\pgfqpoint{3.881866in}{1.931148in}}%
\pgfpathlineto{\pgfqpoint{3.928427in}{1.947321in}}%
\pgfpathlineto{\pgfqpoint{3.974907in}{1.963453in}}%
\pgfpathlineto{\pgfqpoint{4.003605in}{1.905436in}}%
\pgfpathlineto{\pgfqpoint{4.032428in}{1.849706in}}%
\pgfpathlineto{\pgfqpoint{3.985917in}{1.833421in}}%
\pgfpathlineto{\pgfqpoint{3.939326in}{1.817093in}}%
\pgfpathclose%
\pgfusepath{stroke,fill}%
\end{pgfscope}%
\begin{pgfscope}%
\pgfpathrectangle{\pgfqpoint{0.050000in}{0.083252in}}{\pgfqpoint{4.900000in}{4.900000in}}%
\pgfusepath{clip}%
\pgfsetbuttcap%
\pgfsetroundjoin%
\definecolor{currentfill}{rgb}{0.247951,0.300654,0.668481}%
\pgfsetfillcolor{currentfill}%
\pgfsetfillopacity{0.950000}%
\pgfsetlinewidth{0.200750pt}%
\definecolor{currentstroke}{rgb}{0.200000,0.200000,0.200000}%
\pgfsetstrokecolor{currentstroke}%
\pgfsetdash{}{0pt}%
\pgfpathmoveto{\pgfqpoint{2.662537in}{3.131046in}}%
\pgfpathlineto{\pgfqpoint{2.634909in}{3.204830in}}%
\pgfpathlineto{\pgfqpoint{2.607307in}{3.278551in}}%
\pgfpathlineto{\pgfqpoint{2.654981in}{3.292451in}}%
\pgfpathlineto{\pgfqpoint{2.702571in}{3.306327in}}%
\pgfpathlineto{\pgfqpoint{2.730170in}{3.232754in}}%
\pgfpathlineto{\pgfqpoint{2.757794in}{3.159120in}}%
\pgfpathlineto{\pgfqpoint{2.710208in}{3.145094in}}%
\pgfpathlineto{\pgfqpoint{2.662537in}{3.131046in}}%
\pgfpathclose%
\pgfusepath{stroke,fill}%
\end{pgfscope}%
\begin{pgfscope}%
\pgfpathrectangle{\pgfqpoint{0.050000in}{0.083252in}}{\pgfqpoint{4.900000in}{4.900000in}}%
\pgfusepath{clip}%
\pgfsetbuttcap%
\pgfsetroundjoin%
\definecolor{currentfill}{rgb}{0.241307,0.230065,0.574641}%
\pgfsetfillcolor{currentfill}%
\pgfsetfillopacity{0.950000}%
\pgfsetlinewidth{0.200750pt}%
\definecolor{currentstroke}{rgb}{0.200000,0.200000,0.200000}%
\pgfsetstrokecolor{currentstroke}%
\pgfsetdash{}{0pt}%
\pgfpathmoveto{\pgfqpoint{2.868551in}{2.864091in}}%
\pgfpathlineto{\pgfqpoint{2.840822in}{2.937904in}}%
\pgfpathlineto{\pgfqpoint{2.813120in}{3.011688in}}%
\pgfpathlineto{\pgfqpoint{2.860618in}{3.025845in}}%
\pgfpathlineto{\pgfqpoint{2.908031in}{3.039980in}}%
\pgfpathlineto{\pgfqpoint{2.935729in}{2.966361in}}%
\pgfpathlineto{\pgfqpoint{2.963454in}{2.892719in}}%
\pgfpathlineto{\pgfqpoint{2.916044in}{2.878416in}}%
\pgfpathlineto{\pgfqpoint{2.868551in}{2.864091in}}%
\pgfpathclose%
\pgfusepath{stroke,fill}%
\end{pgfscope}%
\begin{pgfscope}%
\pgfpathrectangle{\pgfqpoint{0.050000in}{0.083252in}}{\pgfqpoint{4.900000in}{4.900000in}}%
\pgfusepath{clip}%
\pgfsetbuttcap%
\pgfsetroundjoin%
\definecolor{currentfill}{rgb}{0.773256,0.763906,0.864421}%
\pgfsetfillcolor{currentfill}%
\pgfsetfillopacity{0.950000}%
\pgfsetlinewidth{0.200750pt}%
\definecolor{currentstroke}{rgb}{0.200000,0.200000,0.200000}%
\pgfsetstrokecolor{currentstroke}%
\pgfsetdash{}{0pt}%
\pgfpathmoveto{\pgfqpoint{3.788503in}{1.898679in}}%
\pgfpathlineto{\pgfqpoint{3.759971in}{1.959280in}}%
\pgfpathlineto{\pgfqpoint{3.731539in}{2.021904in}}%
\pgfpathlineto{\pgfqpoint{3.778237in}{2.037909in}}%
\pgfpathlineto{\pgfqpoint{3.824854in}{2.053877in}}%
\pgfpathlineto{\pgfqpoint{3.853309in}{1.991519in}}%
\pgfpathlineto{\pgfqpoint{3.881866in}{1.931148in}}%
\pgfpathlineto{\pgfqpoint{3.835225in}{1.914934in}}%
\pgfpathlineto{\pgfqpoint{3.788503in}{1.898679in}}%
\pgfpathclose%
\pgfusepath{stroke,fill}%
\end{pgfscope}%
\begin{pgfscope}%
\pgfpathrectangle{\pgfqpoint{0.050000in}{0.083252in}}{\pgfqpoint{4.900000in}{4.900000in}}%
\pgfusepath{clip}%
\pgfsetbuttcap%
\pgfsetroundjoin%
\definecolor{currentfill}{rgb}{0.343637,0.316571,0.607536}%
\pgfsetfillcolor{currentfill}%
\pgfsetfillopacity{0.950000}%
\pgfsetlinewidth{0.200750pt}%
\definecolor{currentstroke}{rgb}{0.200000,0.200000,0.200000}%
\pgfsetstrokecolor{currentstroke}%
\pgfsetdash{}{0pt}%
\pgfpathmoveto{\pgfqpoint{3.074625in}{2.598425in}}%
\pgfpathlineto{\pgfqpoint{3.046790in}{2.671887in}}%
\pgfpathlineto{\pgfqpoint{3.018984in}{2.745452in}}%
\pgfpathlineto{\pgfqpoint{3.066306in}{2.759922in}}%
\pgfpathlineto{\pgfqpoint{3.113544in}{2.774372in}}%
\pgfpathlineto{\pgfqpoint{3.141347in}{2.701017in}}%
\pgfpathlineto{\pgfqpoint{3.169180in}{2.627780in}}%
\pgfpathlineto{\pgfqpoint{3.121944in}{2.613112in}}%
\pgfpathlineto{\pgfqpoint{3.074625in}{2.598425in}}%
\pgfpathclose%
\pgfusepath{stroke,fill}%
\end{pgfscope}%
\begin{pgfscope}%
\pgfpathrectangle{\pgfqpoint{0.050000in}{0.083252in}}{\pgfqpoint{4.900000in}{4.900000in}}%
\pgfusepath{clip}%
\pgfsetbuttcap%
\pgfsetroundjoin%
\definecolor{currentfill}{rgb}{0.713587,0.701776,0.828743}%
\pgfsetfillcolor{currentfill}%
\pgfsetfillopacity{0.950000}%
\pgfsetlinewidth{0.200750pt}%
\definecolor{currentstroke}{rgb}{0.200000,0.200000,0.200000}%
\pgfsetstrokecolor{currentstroke}%
\pgfsetdash{}{0pt}%
\pgfpathmoveto{\pgfqpoint{3.637900in}{1.989784in}}%
\pgfpathlineto{\pgfqpoint{3.609575in}{2.054525in}}%
\pgfpathlineto{\pgfqpoint{3.581328in}{2.120881in}}%
\pgfpathlineto{\pgfqpoint{3.628170in}{2.136644in}}%
\pgfpathlineto{\pgfqpoint{3.674932in}{2.152376in}}%
\pgfpathlineto{\pgfqpoint{3.703195in}{2.086341in}}%
\pgfpathlineto{\pgfqpoint{3.731539in}{2.021904in}}%
\pgfpathlineto{\pgfqpoint{3.684760in}{2.005863in}}%
\pgfpathlineto{\pgfqpoint{3.637900in}{1.989784in}}%
\pgfpathclose%
\pgfusepath{stroke,fill}%
\end{pgfscope}%
\begin{pgfscope}%
\pgfpathrectangle{\pgfqpoint{0.050000in}{0.083252in}}{\pgfqpoint{4.900000in}{4.900000in}}%
\pgfusepath{clip}%
\pgfsetbuttcap%
\pgfsetroundjoin%
\definecolor{currentfill}{rgb}{0.498777,0.478108,0.700300}%
\pgfsetfillcolor{currentfill}%
\pgfsetfillopacity{0.950000}%
\pgfsetlinewidth{0.200750pt}%
\definecolor{currentstroke}{rgb}{0.200000,0.200000,0.200000}%
\pgfsetstrokecolor{currentstroke}%
\pgfsetdash{}{0pt}%
\pgfpathmoveto{\pgfqpoint{3.280831in}{2.337516in}}%
\pgfpathlineto{\pgfqpoint{3.252865in}{2.409486in}}%
\pgfpathlineto{\pgfqpoint{3.224937in}{2.481928in}}%
\pgfpathlineto{\pgfqpoint{3.272088in}{2.496831in}}%
\pgfpathlineto{\pgfqpoint{3.319157in}{2.511715in}}%
\pgfpathlineto{\pgfqpoint{3.347087in}{2.439558in}}%
\pgfpathlineto{\pgfqpoint{3.375055in}{2.367893in}}%
\pgfpathlineto{\pgfqpoint{3.327984in}{2.352715in}}%
\pgfpathlineto{\pgfqpoint{3.280831in}{2.337516in}}%
\pgfpathclose%
\pgfusepath{stroke,fill}%
\end{pgfscope}%
\begin{pgfscope}%
\pgfpathrectangle{\pgfqpoint{0.050000in}{0.083252in}}{\pgfqpoint{4.900000in}{4.900000in}}%
\pgfusepath{clip}%
\pgfsetbuttcap%
\pgfsetroundjoin%
\definecolor{currentfill}{rgb}{0.874694,0.869527,0.925075}%
\pgfsetfillcolor{currentfill}%
\pgfsetfillopacity{0.950000}%
\pgfsetlinewidth{0.200750pt}%
\definecolor{currentstroke}{rgb}{0.200000,0.200000,0.200000}%
\pgfsetstrokecolor{currentstroke}%
\pgfsetdash{}{0pt}%
\pgfpathmoveto{\pgfqpoint{4.090485in}{1.745580in}}%
\pgfpathlineto{\pgfqpoint{4.061385in}{1.796391in}}%
\pgfpathlineto{\pgfqpoint{4.032428in}{1.849706in}}%
\pgfpathlineto{\pgfqpoint{4.078859in}{1.865948in}}%
\pgfpathlineto{\pgfqpoint{4.107833in}{1.812670in}}%
\pgfpathlineto{\pgfqpoint{4.136952in}{1.761866in}}%
\pgfpathlineto{\pgfqpoint{4.090485in}{1.745580in}}%
\pgfpathclose%
\pgfusepath{stroke,fill}%
\end{pgfscope}%
\begin{pgfscope}%
\pgfpathrectangle{\pgfqpoint{0.050000in}{0.083252in}}{\pgfqpoint{4.900000in}{4.900000in}}%
\pgfusepath{clip}%
\pgfsetbuttcap%
\pgfsetroundjoin%
\definecolor{currentfill}{rgb}{0.264437,0.475817,0.901346}%
\pgfsetfillcolor{currentfill}%
\pgfsetfillopacity{0.950000}%
\pgfsetlinewidth{0.200750pt}%
\definecolor{currentstroke}{rgb}{0.200000,0.200000,0.200000}%
\pgfsetstrokecolor{currentstroke}%
\pgfsetdash{}{0pt}%
\pgfpathmoveto{\pgfqpoint{2.099495in}{3.785455in}}%
\pgfpathlineto{\pgfqpoint{2.072113in}{3.858996in}}%
\pgfpathlineto{\pgfqpoint{2.044757in}{3.932468in}}%
\pgfpathlineto{\pgfqpoint{2.092959in}{3.945795in}}%
\pgfpathlineto{\pgfqpoint{2.141077in}{3.959098in}}%
\pgfpathlineto{\pgfqpoint{2.168429in}{3.885768in}}%
\pgfpathlineto{\pgfqpoint{2.195808in}{3.812370in}}%
\pgfpathlineto{\pgfqpoint{2.147694in}{3.798925in}}%
\pgfpathlineto{\pgfqpoint{2.099495in}{3.785455in}}%
\pgfpathclose%
\pgfusepath{stroke,fill}%
\end{pgfscope}%
\begin{pgfscope}%
\pgfpathrectangle{\pgfqpoint{0.050000in}{0.083252in}}{\pgfqpoint{4.900000in}{4.900000in}}%
\pgfusepath{clip}%
\pgfsetbuttcap%
\pgfsetroundjoin%
\definecolor{currentfill}{rgb}{0.257793,0.405229,0.807505}%
\pgfsetfillcolor{currentfill}%
\pgfsetfillopacity{0.950000}%
\pgfsetlinewidth{0.200750pt}%
\definecolor{currentstroke}{rgb}{0.200000,0.200000,0.200000}%
\pgfsetstrokecolor{currentstroke}%
\pgfsetdash{}{0pt}%
\pgfpathmoveto{\pgfqpoint{2.305578in}{3.518089in}}%
\pgfpathlineto{\pgfqpoint{2.278097in}{3.591763in}}%
\pgfpathlineto{\pgfqpoint{2.250641in}{3.665367in}}%
\pgfpathlineto{\pgfqpoint{2.298667in}{3.678932in}}%
\pgfpathlineto{\pgfqpoint{2.346607in}{3.692472in}}%
\pgfpathlineto{\pgfqpoint{2.374059in}{3.619011in}}%
\pgfpathlineto{\pgfqpoint{2.401536in}{3.545481in}}%
\pgfpathlineto{\pgfqpoint{2.353600in}{3.531797in}}%
\pgfpathlineto{\pgfqpoint{2.305578in}{3.518089in}}%
\pgfpathclose%
\pgfusepath{stroke,fill}%
\end{pgfscope}%
\begin{pgfscope}%
\pgfpathrectangle{\pgfqpoint{0.050000in}{0.083252in}}{\pgfqpoint{4.900000in}{4.900000in}}%
\pgfusepath{clip}%
\pgfsetbuttcap%
\pgfsetroundjoin%
\definecolor{currentfill}{rgb}{0.270096,0.535948,0.981284}%
\pgfsetfillcolor{currentfill}%
\pgfsetfillopacity{0.950000}%
\pgfsetlinewidth{0.200750pt}%
\definecolor{currentstroke}{rgb}{0.200000,0.200000,0.200000}%
\pgfsetstrokecolor{currentstroke}%
\pgfsetdash{}{0pt}%
\pgfpathmoveto{\pgfqpoint{1.893453in}{4.052766in}}%
\pgfpathlineto{\pgfqpoint{1.866171in}{4.126174in}}%
\pgfpathlineto{\pgfqpoint{1.914550in}{4.139335in}}%
\pgfpathlineto{\pgfqpoint{1.962842in}{4.152472in}}%
\pgfpathlineto{\pgfqpoint{1.990122in}{4.079206in}}%
\pgfpathlineto{\pgfqpoint{1.941830in}{4.065998in}}%
\pgfpathlineto{\pgfqpoint{1.893453in}{4.052766in}}%
\pgfpathclose%
\pgfusepath{stroke,fill}%
\end{pgfscope}%
\begin{pgfscope}%
\pgfpathrectangle{\pgfqpoint{0.050000in}{0.083252in}}{\pgfqpoint{4.900000in}{4.900000in}}%
\pgfusepath{clip}%
\pgfsetbuttcap%
\pgfsetroundjoin%
\definecolor{currentfill}{rgb}{0.251150,0.334641,0.713664}%
\pgfsetfillcolor{currentfill}%
\pgfsetfillopacity{0.950000}%
\pgfsetlinewidth{0.200750pt}%
\definecolor{currentstroke}{rgb}{0.200000,0.200000,0.200000}%
\pgfsetstrokecolor{currentstroke}%
\pgfsetdash{}{0pt}%
\pgfpathmoveto{\pgfqpoint{2.511704in}{3.250678in}}%
\pgfpathlineto{\pgfqpoint{2.484123in}{3.324480in}}%
\pgfpathlineto{\pgfqpoint{2.456568in}{3.398215in}}%
\pgfpathlineto{\pgfqpoint{2.504416in}{3.412018in}}%
\pgfpathlineto{\pgfqpoint{2.552179in}{3.425797in}}%
\pgfpathlineto{\pgfqpoint{2.579730in}{3.352207in}}%
\pgfpathlineto{\pgfqpoint{2.607307in}{3.278551in}}%
\pgfpathlineto{\pgfqpoint{2.559548in}{3.264627in}}%
\pgfpathlineto{\pgfqpoint{2.511704in}{3.250678in}}%
\pgfpathclose%
\pgfusepath{stroke,fill}%
\end{pgfscope}%
\begin{pgfscope}%
\pgfpathrectangle{\pgfqpoint{0.050000in}{0.083252in}}{\pgfqpoint{4.900000in}{4.900000in}}%
\pgfusepath{clip}%
\pgfsetbuttcap%
\pgfsetroundjoin%
\definecolor{currentfill}{rgb}{0.647951,0.633433,0.789496}%
\pgfsetfillcolor{currentfill}%
\pgfsetfillopacity{0.950000}%
\pgfsetlinewidth{0.200750pt}%
\definecolor{currentstroke}{rgb}{0.200000,0.200000,0.200000}%
\pgfsetstrokecolor{currentstroke}%
\pgfsetdash{}{0pt}%
\pgfpathmoveto{\pgfqpoint{3.487398in}{2.089262in}}%
\pgfpathlineto{\pgfqpoint{3.459231in}{2.157351in}}%
\pgfpathlineto{\pgfqpoint{3.431122in}{2.226618in}}%
\pgfpathlineto{\pgfqpoint{3.478117in}{2.242100in}}%
\pgfpathlineto{\pgfqpoint{3.525030in}{2.257557in}}%
\pgfpathlineto{\pgfqpoint{3.553148in}{2.188629in}}%
\pgfpathlineto{\pgfqpoint{3.581328in}{2.120881in}}%
\pgfpathlineto{\pgfqpoint{3.534404in}{2.105087in}}%
\pgfpathlineto{\pgfqpoint{3.487398in}{2.089262in}}%
\pgfpathclose%
\pgfusepath{stroke,fill}%
\end{pgfscope}%
\begin{pgfscope}%
\pgfpathrectangle{\pgfqpoint{0.050000in}{0.083252in}}{\pgfqpoint{4.900000in}{4.900000in}}%
\pgfusepath{clip}%
\pgfsetbuttcap%
\pgfsetroundjoin%
\definecolor{currentfill}{rgb}{0.244752,0.266667,0.623299}%
\pgfsetfillcolor{currentfill}%
\pgfsetfillopacity{0.950000}%
\pgfsetlinewidth{0.200750pt}%
\definecolor{currentstroke}{rgb}{0.200000,0.200000,0.200000}%
\pgfsetstrokecolor{currentstroke}%
\pgfsetdash{}{0pt}%
\pgfpathmoveto{\pgfqpoint{2.717872in}{2.983304in}}%
\pgfpathlineto{\pgfqpoint{2.690192in}{3.057202in}}%
\pgfpathlineto{\pgfqpoint{2.662537in}{3.131046in}}%
\pgfpathlineto{\pgfqpoint{2.710208in}{3.145094in}}%
\pgfpathlineto{\pgfqpoint{2.757794in}{3.159120in}}%
\pgfpathlineto{\pgfqpoint{2.785444in}{3.085429in}}%
\pgfpathlineto{\pgfqpoint{2.813120in}{3.011688in}}%
\pgfpathlineto{\pgfqpoint{2.765538in}{2.997507in}}%
\pgfpathlineto{\pgfqpoint{2.717872in}{2.983304in}}%
\pgfpathclose%
\pgfusepath{stroke,fill}%
\end{pgfscope}%
\begin{pgfscope}%
\pgfpathrectangle{\pgfqpoint{0.050000in}{0.083252in}}{\pgfqpoint{4.900000in}{4.900000in}}%
\pgfusepath{clip}%
\pgfsetbuttcap%
\pgfsetroundjoin%
\definecolor{currentfill}{rgb}{0.266067,0.235802,0.561153}%
\pgfsetfillcolor{currentfill}%
\pgfsetfillopacity{0.950000}%
\pgfsetlinewidth{0.200750pt}%
\definecolor{currentstroke}{rgb}{0.200000,0.200000,0.200000}%
\pgfsetstrokecolor{currentstroke}%
\pgfsetdash{}{0pt}%
\pgfpathmoveto{\pgfqpoint{2.924088in}{2.716450in}}%
\pgfpathlineto{\pgfqpoint{2.896306in}{2.790265in}}%
\pgfpathlineto{\pgfqpoint{2.868551in}{2.864091in}}%
\pgfpathlineto{\pgfqpoint{2.916044in}{2.878416in}}%
\pgfpathlineto{\pgfqpoint{2.963454in}{2.892719in}}%
\pgfpathlineto{\pgfqpoint{2.991205in}{2.819075in}}%
\pgfpathlineto{\pgfqpoint{3.018984in}{2.745452in}}%
\pgfpathlineto{\pgfqpoint{2.971578in}{2.730962in}}%
\pgfpathlineto{\pgfqpoint{2.924088in}{2.716450in}}%
\pgfpathclose%
\pgfusepath{stroke,fill}%
\end{pgfscope}%
\begin{pgfscope}%
\pgfpathrectangle{\pgfqpoint{0.050000in}{0.083252in}}{\pgfqpoint{4.900000in}{4.900000in}}%
\pgfusepath{clip}%
\pgfsetbuttcap%
\pgfsetroundjoin%
\definecolor{currentfill}{rgb}{0.427174,0.403552,0.657486}%
\pgfsetfillcolor{currentfill}%
\pgfsetfillopacity{0.950000}%
\pgfsetlinewidth{0.200750pt}%
\definecolor{currentstroke}{rgb}{0.200000,0.200000,0.200000}%
\pgfsetstrokecolor{currentstroke}%
\pgfsetdash{}{0pt}%
\pgfpathmoveto{\pgfqpoint{3.130384in}{2.452063in}}%
\pgfpathlineto{\pgfqpoint{3.102489in}{2.525125in}}%
\pgfpathlineto{\pgfqpoint{3.074625in}{2.598425in}}%
\pgfpathlineto{\pgfqpoint{3.121944in}{2.613112in}}%
\pgfpathlineto{\pgfqpoint{3.169180in}{2.627780in}}%
\pgfpathlineto{\pgfqpoint{3.197042in}{2.554725in}}%
\pgfpathlineto{\pgfqpoint{3.224937in}{2.481928in}}%
\pgfpathlineto{\pgfqpoint{3.177702in}{2.467005in}}%
\pgfpathlineto{\pgfqpoint{3.130384in}{2.452063in}}%
\pgfpathclose%
\pgfusepath{stroke,fill}%
\end{pgfscope}%
\begin{pgfscope}%
\pgfpathrectangle{\pgfqpoint{0.050000in}{0.083252in}}{\pgfqpoint{4.900000in}{4.900000in}}%
\pgfusepath{clip}%
\pgfsetbuttcap%
\pgfsetroundjoin%
\definecolor{currentfill}{rgb}{0.267636,0.509804,0.946528}%
\pgfsetfillcolor{currentfill}%
\pgfsetfillopacity{0.950000}%
\pgfsetlinewidth{0.200750pt}%
\definecolor{currentstroke}{rgb}{0.200000,0.200000,0.200000}%
\pgfsetstrokecolor{currentstroke}%
\pgfsetdash{}{0pt}%
\pgfpathmoveto{\pgfqpoint{1.948095in}{3.905743in}}%
\pgfpathlineto{\pgfqpoint{1.920761in}{3.979289in}}%
\pgfpathlineto{\pgfqpoint{1.893453in}{4.052766in}}%
\pgfpathlineto{\pgfqpoint{1.941830in}{4.065998in}}%
\pgfpathlineto{\pgfqpoint{1.990122in}{4.079206in}}%
\pgfpathlineto{\pgfqpoint{2.017426in}{4.005871in}}%
\pgfpathlineto{\pgfqpoint{2.044757in}{3.932468in}}%
\pgfpathlineto{\pgfqpoint{1.996469in}{3.919118in}}%
\pgfpathlineto{\pgfqpoint{1.948095in}{3.905743in}}%
\pgfpathclose%
\pgfusepath{stroke,fill}%
\end{pgfscope}%
\begin{pgfscope}%
\pgfpathrectangle{\pgfqpoint{0.050000in}{0.083252in}}{\pgfqpoint{4.900000in}{4.900000in}}%
\pgfusepath{clip}%
\pgfsetbuttcap%
\pgfsetroundjoin%
\definecolor{currentfill}{rgb}{0.576348,0.558877,0.746682}%
\pgfsetfillcolor{currentfill}%
\pgfsetfillopacity{0.950000}%
\pgfsetlinewidth{0.200750pt}%
\definecolor{currentstroke}{rgb}{0.200000,0.200000,0.200000}%
\pgfsetstrokecolor{currentstroke}%
\pgfsetdash{}{0pt}%
\pgfpathmoveto{\pgfqpoint{3.336887in}{2.195578in}}%
\pgfpathlineto{\pgfqpoint{3.308837in}{2.266159in}}%
\pgfpathlineto{\pgfqpoint{3.280831in}{2.337516in}}%
\pgfpathlineto{\pgfqpoint{3.327984in}{2.352715in}}%
\pgfpathlineto{\pgfqpoint{3.375055in}{2.367893in}}%
\pgfpathlineto{\pgfqpoint{3.403066in}{2.296859in}}%
\pgfpathlineto{\pgfqpoint{3.431122in}{2.226618in}}%
\pgfpathlineto{\pgfqpoint{3.384046in}{2.211110in}}%
\pgfpathlineto{\pgfqpoint{3.336887in}{2.195578in}}%
\pgfpathclose%
\pgfusepath{stroke,fill}%
\end{pgfscope}%
\begin{pgfscope}%
\pgfpathrectangle{\pgfqpoint{0.050000in}{0.083252in}}{\pgfqpoint{4.900000in}{4.900000in}}%
\pgfusepath{clip}%
\pgfsetbuttcap%
\pgfsetroundjoin%
\definecolor{currentfill}{rgb}{0.260992,0.439216,0.852687}%
\pgfsetfillcolor{currentfill}%
\pgfsetfillopacity{0.950000}%
\pgfsetlinewidth{0.200750pt}%
\definecolor{currentstroke}{rgb}{0.200000,0.200000,0.200000}%
\pgfsetstrokecolor{currentstroke}%
\pgfsetdash{}{0pt}%
\pgfpathmoveto{\pgfqpoint{2.154335in}{3.638166in}}%
\pgfpathlineto{\pgfqpoint{2.126902in}{3.711845in}}%
\pgfpathlineto{\pgfqpoint{2.099495in}{3.785455in}}%
\pgfpathlineto{\pgfqpoint{2.147694in}{3.798925in}}%
\pgfpathlineto{\pgfqpoint{2.195808in}{3.812370in}}%
\pgfpathlineto{\pgfqpoint{2.223212in}{3.738903in}}%
\pgfpathlineto{\pgfqpoint{2.250641in}{3.665367in}}%
\pgfpathlineto{\pgfqpoint{2.202531in}{3.651779in}}%
\pgfpathlineto{\pgfqpoint{2.154335in}{3.638166in}}%
\pgfpathclose%
\pgfusepath{stroke,fill}%
\end{pgfscope}%
\begin{pgfscope}%
\pgfpathrectangle{\pgfqpoint{0.050000in}{0.083252in}}{\pgfqpoint{4.900000in}{4.900000in}}%
\pgfusepath{clip}%
\pgfsetbuttcap%
\pgfsetroundjoin%
\definecolor{currentfill}{rgb}{0.254594,0.371242,0.762322}%
\pgfsetfillcolor{currentfill}%
\pgfsetfillopacity{0.950000}%
\pgfsetlinewidth{0.200750pt}%
\definecolor{currentstroke}{rgb}{0.200000,0.200000,0.200000}%
\pgfsetstrokecolor{currentstroke}%
\pgfsetdash{}{0pt}%
\pgfpathmoveto{\pgfqpoint{2.360618in}{3.370535in}}%
\pgfpathlineto{\pgfqpoint{2.333085in}{3.444347in}}%
\pgfpathlineto{\pgfqpoint{2.305578in}{3.518089in}}%
\pgfpathlineto{\pgfqpoint{2.353600in}{3.531797in}}%
\pgfpathlineto{\pgfqpoint{2.401536in}{3.545481in}}%
\pgfpathlineto{\pgfqpoint{2.429039in}{3.471882in}}%
\pgfpathlineto{\pgfqpoint{2.456568in}{3.398215in}}%
\pgfpathlineto{\pgfqpoint{2.408636in}{3.384387in}}%
\pgfpathlineto{\pgfqpoint{2.360618in}{3.370535in}}%
\pgfpathclose%
\pgfusepath{stroke,fill}%
\end{pgfscope}%
\begin{pgfscope}%
\pgfpathrectangle{\pgfqpoint{0.050000in}{0.083252in}}{\pgfqpoint{4.900000in}{4.900000in}}%
\pgfusepath{clip}%
\pgfsetbuttcap%
\pgfsetroundjoin%
\definecolor{currentfill}{rgb}{0.826959,0.819823,0.896532}%
\pgfsetfillcolor{currentfill}%
\pgfsetfillopacity{0.950000}%
\pgfsetlinewidth{0.200750pt}%
\definecolor{currentstroke}{rgb}{0.200000,0.200000,0.200000}%
\pgfsetstrokecolor{currentstroke}%
\pgfsetdash{}{0pt}%
\pgfpathmoveto{\pgfqpoint{3.845900in}{1.784305in}}%
\pgfpathlineto{\pgfqpoint{3.817142in}{1.840297in}}%
\pgfpathlineto{\pgfqpoint{3.788503in}{1.898679in}}%
\pgfpathlineto{\pgfqpoint{3.835225in}{1.914934in}}%
\pgfpathlineto{\pgfqpoint{3.881866in}{1.931148in}}%
\pgfpathlineto{\pgfqpoint{3.910535in}{1.872952in}}%
\pgfpathlineto{\pgfqpoint{3.939326in}{1.817093in}}%
\pgfpathlineto{\pgfqpoint{3.892653in}{1.800721in}}%
\pgfpathlineto{\pgfqpoint{3.845900in}{1.784305in}}%
\pgfpathclose%
\pgfusepath{stroke,fill}%
\end{pgfscope}%
\begin{pgfscope}%
\pgfpathrectangle{\pgfqpoint{0.050000in}{0.083252in}}{\pgfqpoint{4.900000in}{4.900000in}}%
\pgfusepath{clip}%
\pgfsetbuttcap%
\pgfsetroundjoin%
\definecolor{currentfill}{rgb}{0.247951,0.300654,0.668481}%
\pgfsetfillcolor{currentfill}%
\pgfsetfillopacity{0.950000}%
\pgfsetlinewidth{0.200750pt}%
\definecolor{currentstroke}{rgb}{0.200000,0.200000,0.200000}%
\pgfsetstrokecolor{currentstroke}%
\pgfsetdash{}{0pt}%
\pgfpathmoveto{\pgfqpoint{2.566943in}{3.102876in}}%
\pgfpathlineto{\pgfqpoint{2.539310in}{3.176809in}}%
\pgfpathlineto{\pgfqpoint{2.511704in}{3.250678in}}%
\pgfpathlineto{\pgfqpoint{2.559548in}{3.264627in}}%
\pgfpathlineto{\pgfqpoint{2.607307in}{3.278551in}}%
\pgfpathlineto{\pgfqpoint{2.634909in}{3.204830in}}%
\pgfpathlineto{\pgfqpoint{2.662537in}{3.131046in}}%
\pgfpathlineto{\pgfqpoint{2.614783in}{3.116973in}}%
\pgfpathlineto{\pgfqpoint{2.566943in}{3.102876in}}%
\pgfpathclose%
\pgfusepath{stroke,fill}%
\end{pgfscope}%
\begin{pgfscope}%
\pgfpathrectangle{\pgfqpoint{0.050000in}{0.083252in}}{\pgfqpoint{4.900000in}{4.900000in}}%
\pgfusepath{clip}%
\pgfsetbuttcap%
\pgfsetroundjoin%
\definecolor{currentfill}{rgb}{0.874694,0.869527,0.925075}%
\pgfsetfillcolor{currentfill}%
\pgfsetfillopacity{0.950000}%
\pgfsetlinewidth{0.200750pt}%
\definecolor{currentstroke}{rgb}{0.200000,0.200000,0.200000}%
\pgfsetstrokecolor{currentstroke}%
\pgfsetdash{}{0pt}%
\pgfpathmoveto{\pgfqpoint{3.997307in}{1.712879in}}%
\pgfpathlineto{\pgfqpoint{3.968247in}{1.763705in}}%
\pgfpathlineto{\pgfqpoint{3.939326in}{1.817093in}}%
\pgfpathlineto{\pgfqpoint{3.985917in}{1.833421in}}%
\pgfpathlineto{\pgfqpoint{4.032428in}{1.849706in}}%
\pgfpathlineto{\pgfqpoint{4.061385in}{1.796391in}}%
\pgfpathlineto{\pgfqpoint{4.090485in}{1.745580in}}%
\pgfpathlineto{\pgfqpoint{4.043936in}{1.729251in}}%
\pgfpathlineto{\pgfqpoint{3.997307in}{1.712879in}}%
\pgfpathclose%
\pgfusepath{stroke,fill}%
\end{pgfscope}%
\begin{pgfscope}%
\pgfpathrectangle{\pgfqpoint{0.050000in}{0.083252in}}{\pgfqpoint{4.900000in}{4.900000in}}%
\pgfusepath{clip}%
\pgfsetbuttcap%
\pgfsetroundjoin%
\definecolor{currentfill}{rgb}{0.779223,0.770119,0.867989}%
\pgfsetfillcolor{currentfill}%
\pgfsetfillopacity{0.950000}%
\pgfsetlinewidth{0.200750pt}%
\definecolor{currentstroke}{rgb}{0.200000,0.200000,0.200000}%
\pgfsetstrokecolor{currentstroke}%
\pgfsetdash{}{0pt}%
\pgfpathmoveto{\pgfqpoint{3.694815in}{1.866041in}}%
\pgfpathlineto{\pgfqpoint{3.666310in}{1.926882in}}%
\pgfpathlineto{\pgfqpoint{3.637900in}{1.989784in}}%
\pgfpathlineto{\pgfqpoint{3.684760in}{2.005863in}}%
\pgfpathlineto{\pgfqpoint{3.731539in}{2.021904in}}%
\pgfpathlineto{\pgfqpoint{3.759971in}{1.959280in}}%
\pgfpathlineto{\pgfqpoint{3.788503in}{1.898679in}}%
\pgfpathlineto{\pgfqpoint{3.741699in}{1.882381in}}%
\pgfpathlineto{\pgfqpoint{3.694815in}{1.866041in}}%
\pgfpathclose%
\pgfusepath{stroke,fill}%
\end{pgfscope}%
\begin{pgfscope}%
\pgfpathrectangle{\pgfqpoint{0.050000in}{0.083252in}}{\pgfqpoint{4.900000in}{4.900000in}}%
\pgfusepath{clip}%
\pgfsetbuttcap%
\pgfsetroundjoin%
\definecolor{currentfill}{rgb}{0.241307,0.230065,0.574641}%
\pgfsetfillcolor{currentfill}%
\pgfsetfillopacity{0.950000}%
\pgfsetlinewidth{0.200750pt}%
\definecolor{currentstroke}{rgb}{0.200000,0.200000,0.200000}%
\pgfsetstrokecolor{currentstroke}%
\pgfsetdash{}{0pt}%
\pgfpathmoveto{\pgfqpoint{2.773312in}{2.835375in}}%
\pgfpathlineto{\pgfqpoint{2.745579in}{2.909358in}}%
\pgfpathlineto{\pgfqpoint{2.717872in}{2.983304in}}%
\pgfpathlineto{\pgfqpoint{2.765538in}{2.997507in}}%
\pgfpathlineto{\pgfqpoint{2.813120in}{3.011688in}}%
\pgfpathlineto{\pgfqpoint{2.840822in}{2.937904in}}%
\pgfpathlineto{\pgfqpoint{2.868551in}{2.864091in}}%
\pgfpathlineto{\pgfqpoint{2.820973in}{2.849744in}}%
\pgfpathlineto{\pgfqpoint{2.773312in}{2.835375in}}%
\pgfpathclose%
\pgfusepath{stroke,fill}%
\end{pgfscope}%
\begin{pgfscope}%
\pgfpathrectangle{\pgfqpoint{0.050000in}{0.083252in}}{\pgfqpoint{4.900000in}{4.900000in}}%
\pgfusepath{clip}%
\pgfsetbuttcap%
\pgfsetroundjoin%
\definecolor{currentfill}{rgb}{0.343637,0.316571,0.607536}%
\pgfsetfillcolor{currentfill}%
\pgfsetfillopacity{0.950000}%
\pgfsetlinewidth{0.200750pt}%
\definecolor{currentstroke}{rgb}{0.200000,0.200000,0.200000}%
\pgfsetstrokecolor{currentstroke}%
\pgfsetdash{}{0pt}%
\pgfpathmoveto{\pgfqpoint{2.979735in}{2.568993in}}%
\pgfpathlineto{\pgfqpoint{2.951898in}{2.642679in}}%
\pgfpathlineto{\pgfqpoint{2.924088in}{2.716450in}}%
\pgfpathlineto{\pgfqpoint{2.971578in}{2.730962in}}%
\pgfpathlineto{\pgfqpoint{3.018984in}{2.745452in}}%
\pgfpathlineto{\pgfqpoint{3.046790in}{2.671887in}}%
\pgfpathlineto{\pgfqpoint{3.074625in}{2.598425in}}%
\pgfpathlineto{\pgfqpoint{3.027222in}{2.583719in}}%
\pgfpathlineto{\pgfqpoint{2.979735in}{2.568993in}}%
\pgfpathclose%
\pgfusepath{stroke,fill}%
\end{pgfscope}%
\begin{pgfscope}%
\pgfpathrectangle{\pgfqpoint{0.050000in}{0.083252in}}{\pgfqpoint{4.900000in}{4.900000in}}%
\pgfusepath{clip}%
\pgfsetbuttcap%
\pgfsetroundjoin%
\definecolor{currentfill}{rgb}{0.713587,0.701776,0.828743}%
\pgfsetfillcolor{currentfill}%
\pgfsetfillopacity{0.950000}%
\pgfsetlinewidth{0.200750pt}%
\definecolor{currentstroke}{rgb}{0.200000,0.200000,0.200000}%
\pgfsetstrokecolor{currentstroke}%
\pgfsetdash{}{0pt}%
\pgfpathmoveto{\pgfqpoint{3.543936in}{1.957514in}}%
\pgfpathlineto{\pgfqpoint{3.515630in}{2.022573in}}%
\pgfpathlineto{\pgfqpoint{3.487398in}{2.089262in}}%
\pgfpathlineto{\pgfqpoint{3.534404in}{2.105087in}}%
\pgfpathlineto{\pgfqpoint{3.581328in}{2.120881in}}%
\pgfpathlineto{\pgfqpoint{3.609575in}{2.054525in}}%
\pgfpathlineto{\pgfqpoint{3.637900in}{1.989784in}}%
\pgfpathlineto{\pgfqpoint{3.590958in}{1.973668in}}%
\pgfpathlineto{\pgfqpoint{3.543936in}{1.957514in}}%
\pgfpathclose%
\pgfusepath{stroke,fill}%
\end{pgfscope}%
\begin{pgfscope}%
\pgfpathrectangle{\pgfqpoint{0.050000in}{0.083252in}}{\pgfqpoint{4.900000in}{4.900000in}}%
\pgfusepath{clip}%
\pgfsetbuttcap%
\pgfsetroundjoin%
\definecolor{currentfill}{rgb}{0.504744,0.484321,0.703868}%
\pgfsetfillcolor{currentfill}%
\pgfsetfillopacity{0.950000}%
\pgfsetlinewidth{0.200750pt}%
\definecolor{currentstroke}{rgb}{0.200000,0.200000,0.200000}%
\pgfsetstrokecolor{currentstroke}%
\pgfsetdash{}{0pt}%
\pgfpathmoveto{\pgfqpoint{3.186276in}{2.307055in}}%
\pgfpathlineto{\pgfqpoint{3.158312in}{2.379334in}}%
\pgfpathlineto{\pgfqpoint{3.130384in}{2.452063in}}%
\pgfpathlineto{\pgfqpoint{3.177702in}{2.467005in}}%
\pgfpathlineto{\pgfqpoint{3.224937in}{2.481928in}}%
\pgfpathlineto{\pgfqpoint{3.252865in}{2.409486in}}%
\pgfpathlineto{\pgfqpoint{3.280831in}{2.337516in}}%
\pgfpathlineto{\pgfqpoint{3.233595in}{2.322296in}}%
\pgfpathlineto{\pgfqpoint{3.186276in}{2.307055in}}%
\pgfpathclose%
\pgfusepath{stroke,fill}%
\end{pgfscope}%
\begin{pgfscope}%
\pgfpathrectangle{\pgfqpoint{0.050000in}{0.083252in}}{\pgfqpoint{4.900000in}{4.900000in}}%
\pgfusepath{clip}%
\pgfsetbuttcap%
\pgfsetroundjoin%
\definecolor{currentfill}{rgb}{0.264437,0.475817,0.901346}%
\pgfsetfillcolor{currentfill}%
\pgfsetfillopacity{0.950000}%
\pgfsetlinewidth{0.200750pt}%
\definecolor{currentstroke}{rgb}{0.200000,0.200000,0.200000}%
\pgfsetstrokecolor{currentstroke}%
\pgfsetdash{}{0pt}%
\pgfpathmoveto{\pgfqpoint{2.002839in}{3.758444in}}%
\pgfpathlineto{\pgfqpoint{1.975454in}{3.832128in}}%
\pgfpathlineto{\pgfqpoint{1.948095in}{3.905743in}}%
\pgfpathlineto{\pgfqpoint{1.996469in}{3.919118in}}%
\pgfpathlineto{\pgfqpoint{2.044757in}{3.932468in}}%
\pgfpathlineto{\pgfqpoint{2.072113in}{3.858996in}}%
\pgfpathlineto{\pgfqpoint{2.099495in}{3.785455in}}%
\pgfpathlineto{\pgfqpoint{2.051210in}{3.771961in}}%
\pgfpathlineto{\pgfqpoint{2.002839in}{3.758444in}}%
\pgfpathclose%
\pgfusepath{stroke,fill}%
\end{pgfscope}%
\begin{pgfscope}%
\pgfpathrectangle{\pgfqpoint{0.050000in}{0.083252in}}{\pgfqpoint{4.900000in}{4.900000in}}%
\pgfusepath{clip}%
\pgfsetbuttcap%
\pgfsetroundjoin%
\definecolor{currentfill}{rgb}{0.257793,0.405229,0.807505}%
\pgfsetfillcolor{currentfill}%
\pgfsetfillopacity{0.950000}%
\pgfsetlinewidth{0.200750pt}%
\definecolor{currentstroke}{rgb}{0.200000,0.200000,0.200000}%
\pgfsetstrokecolor{currentstroke}%
\pgfsetdash{}{0pt}%
\pgfpathmoveto{\pgfqpoint{2.209279in}{3.490600in}}%
\pgfpathlineto{\pgfqpoint{2.181794in}{3.564418in}}%
\pgfpathlineto{\pgfqpoint{2.154335in}{3.638166in}}%
\pgfpathlineto{\pgfqpoint{2.202531in}{3.651779in}}%
\pgfpathlineto{\pgfqpoint{2.250641in}{3.665367in}}%
\pgfpathlineto{\pgfqpoint{2.278097in}{3.591763in}}%
\pgfpathlineto{\pgfqpoint{2.305578in}{3.518089in}}%
\pgfpathlineto{\pgfqpoint{2.257471in}{3.504357in}}%
\pgfpathlineto{\pgfqpoint{2.209279in}{3.490600in}}%
\pgfpathclose%
\pgfusepath{stroke,fill}%
\end{pgfscope}%
\begin{pgfscope}%
\pgfpathrectangle{\pgfqpoint{0.050000in}{0.083252in}}{\pgfqpoint{4.900000in}{4.900000in}}%
\pgfusepath{clip}%
\pgfsetbuttcap%
\pgfsetroundjoin%
\definecolor{currentfill}{rgb}{0.270096,0.535948,0.981284}%
\pgfsetfillcolor{currentfill}%
\pgfsetfillopacity{0.950000}%
\pgfsetlinewidth{0.200750pt}%
\definecolor{currentstroke}{rgb}{0.200000,0.200000,0.200000}%
\pgfsetstrokecolor{currentstroke}%
\pgfsetdash{}{0pt}%
\pgfpathmoveto{\pgfqpoint{1.796442in}{4.026232in}}%
\pgfpathlineto{\pgfqpoint{1.769157in}{4.099783in}}%
\pgfpathlineto{\pgfqpoint{1.817707in}{4.112990in}}%
\pgfpathlineto{\pgfqpoint{1.866171in}{4.126174in}}%
\pgfpathlineto{\pgfqpoint{1.893453in}{4.052766in}}%
\pgfpathlineto{\pgfqpoint{1.844991in}{4.039511in}}%
\pgfpathlineto{\pgfqpoint{1.796442in}{4.026232in}}%
\pgfpathclose%
\pgfusepath{stroke,fill}%
\end{pgfscope}%
\begin{pgfscope}%
\pgfpathrectangle{\pgfqpoint{0.050000in}{0.083252in}}{\pgfqpoint{4.900000in}{4.900000in}}%
\pgfusepath{clip}%
\pgfsetbuttcap%
\pgfsetroundjoin%
\definecolor{currentfill}{rgb}{0.910496,0.906805,0.946482}%
\pgfsetfillcolor{currentfill}%
\pgfsetfillopacity{0.950000}%
\pgfsetlinewidth{0.200750pt}%
\definecolor{currentstroke}{rgb}{0.200000,0.200000,0.200000}%
\pgfsetstrokecolor{currentstroke}%
\pgfsetdash{}{0pt}%
\pgfpathmoveto{\pgfqpoint{4.149137in}{1.651578in}}%
\pgfpathlineto{\pgfqpoint{4.119733in}{1.697312in}}%
\pgfpathlineto{\pgfqpoint{4.090485in}{1.745580in}}%
\pgfpathlineto{\pgfqpoint{4.136952in}{1.761866in}}%
\pgfpathlineto{\pgfqpoint{4.166221in}{1.713577in}}%
\pgfpathlineto{\pgfqpoint{4.195647in}{1.667794in}}%
\pgfpathlineto{\pgfqpoint{4.149137in}{1.651578in}}%
\pgfpathclose%
\pgfusepath{stroke,fill}%
\end{pgfscope}%
\begin{pgfscope}%
\pgfpathrectangle{\pgfqpoint{0.050000in}{0.083252in}}{\pgfqpoint{4.900000in}{4.900000in}}%
\pgfusepath{clip}%
\pgfsetbuttcap%
\pgfsetroundjoin%
\definecolor{currentfill}{rgb}{0.251150,0.334641,0.713664}%
\pgfsetfillcolor{currentfill}%
\pgfsetfillopacity{0.950000}%
\pgfsetlinewidth{0.200750pt}%
\definecolor{currentstroke}{rgb}{0.200000,0.200000,0.200000}%
\pgfsetstrokecolor{currentstroke}%
\pgfsetdash{}{0pt}%
\pgfpathmoveto{\pgfqpoint{2.415761in}{3.222707in}}%
\pgfpathlineto{\pgfqpoint{2.388176in}{3.296655in}}%
\pgfpathlineto{\pgfqpoint{2.360618in}{3.370535in}}%
\pgfpathlineto{\pgfqpoint{2.408636in}{3.384387in}}%
\pgfpathlineto{\pgfqpoint{2.456568in}{3.398215in}}%
\pgfpathlineto{\pgfqpoint{2.484123in}{3.324480in}}%
\pgfpathlineto{\pgfqpoint{2.511704in}{3.250678in}}%
\pgfpathlineto{\pgfqpoint{2.463775in}{3.236705in}}%
\pgfpathlineto{\pgfqpoint{2.415761in}{3.222707in}}%
\pgfpathclose%
\pgfusepath{stroke,fill}%
\end{pgfscope}%
\begin{pgfscope}%
\pgfpathrectangle{\pgfqpoint{0.050000in}{0.083252in}}{\pgfqpoint{4.900000in}{4.900000in}}%
\pgfusepath{clip}%
\pgfsetbuttcap%
\pgfsetroundjoin%
\definecolor{currentfill}{rgb}{0.647951,0.633433,0.789496}%
\pgfsetfillcolor{currentfill}%
\pgfsetfillopacity{0.950000}%
\pgfsetlinewidth{0.200750pt}%
\definecolor{currentstroke}{rgb}{0.200000,0.200000,0.200000}%
\pgfsetstrokecolor{currentstroke}%
\pgfsetdash{}{0pt}%
\pgfpathmoveto{\pgfqpoint{3.393141in}{2.057519in}}%
\pgfpathlineto{\pgfqpoint{3.364986in}{2.125961in}}%
\pgfpathlineto{\pgfqpoint{3.336887in}{2.195578in}}%
\pgfpathlineto{\pgfqpoint{3.384046in}{2.211110in}}%
\pgfpathlineto{\pgfqpoint{3.431122in}{2.226618in}}%
\pgfpathlineto{\pgfqpoint{3.459231in}{2.157351in}}%
\pgfpathlineto{\pgfqpoint{3.487398in}{2.089262in}}%
\pgfpathlineto{\pgfqpoint{3.440311in}{2.073406in}}%
\pgfpathlineto{\pgfqpoint{3.393141in}{2.057519in}}%
\pgfpathclose%
\pgfusepath{stroke,fill}%
\end{pgfscope}%
\begin{pgfscope}%
\pgfpathrectangle{\pgfqpoint{0.050000in}{0.083252in}}{\pgfqpoint{4.900000in}{4.900000in}}%
\pgfusepath{clip}%
\pgfsetbuttcap%
\pgfsetroundjoin%
\definecolor{currentfill}{rgb}{0.244506,0.264052,0.619823}%
\pgfsetfillcolor{currentfill}%
\pgfsetfillopacity{0.950000}%
\pgfsetlinewidth{0.200750pt}%
\definecolor{currentstroke}{rgb}{0.200000,0.200000,0.200000}%
\pgfsetstrokecolor{currentstroke}%
\pgfsetdash{}{0pt}%
\pgfpathmoveto{\pgfqpoint{2.622286in}{2.954826in}}%
\pgfpathlineto{\pgfqpoint{2.594601in}{3.028880in}}%
\pgfpathlineto{\pgfqpoint{2.566943in}{3.102876in}}%
\pgfpathlineto{\pgfqpoint{2.614783in}{3.116973in}}%
\pgfpathlineto{\pgfqpoint{2.662537in}{3.131046in}}%
\pgfpathlineto{\pgfqpoint{2.690192in}{3.057202in}}%
\pgfpathlineto{\pgfqpoint{2.717872in}{2.983304in}}%
\pgfpathlineto{\pgfqpoint{2.670121in}{2.969077in}}%
\pgfpathlineto{\pgfqpoint{2.622286in}{2.954826in}}%
\pgfpathclose%
\pgfusepath{stroke,fill}%
\end{pgfscope}%
\begin{pgfscope}%
\pgfpathrectangle{\pgfqpoint{0.050000in}{0.083252in}}{\pgfqpoint{4.900000in}{4.900000in}}%
\pgfusepath{clip}%
\pgfsetbuttcap%
\pgfsetroundjoin%
\definecolor{currentfill}{rgb}{0.266067,0.235802,0.561153}%
\pgfsetfillcolor{currentfill}%
\pgfsetfillopacity{0.950000}%
\pgfsetlinewidth{0.200750pt}%
\definecolor{currentstroke}{rgb}{0.200000,0.200000,0.200000}%
\pgfsetstrokecolor{currentstroke}%
\pgfsetdash{}{0pt}%
\pgfpathmoveto{\pgfqpoint{2.828857in}{2.687366in}}%
\pgfpathlineto{\pgfqpoint{2.801071in}{2.761371in}}%
\pgfpathlineto{\pgfqpoint{2.773312in}{2.835375in}}%
\pgfpathlineto{\pgfqpoint{2.820973in}{2.849744in}}%
\pgfpathlineto{\pgfqpoint{2.868551in}{2.864091in}}%
\pgfpathlineto{\pgfqpoint{2.896306in}{2.790265in}}%
\pgfpathlineto{\pgfqpoint{2.924088in}{2.716450in}}%
\pgfpathlineto{\pgfqpoint{2.876515in}{2.701919in}}%
\pgfpathlineto{\pgfqpoint{2.828857in}{2.687366in}}%
\pgfpathclose%
\pgfusepath{stroke,fill}%
\end{pgfscope}%
\begin{pgfscope}%
\pgfpathrectangle{\pgfqpoint{0.050000in}{0.083252in}}{\pgfqpoint{4.900000in}{4.900000in}}%
\pgfusepath{clip}%
\pgfsetbuttcap%
\pgfsetroundjoin%
\definecolor{currentfill}{rgb}{0.427174,0.403552,0.657486}%
\pgfsetfillcolor{currentfill}%
\pgfsetfillopacity{0.950000}%
\pgfsetlinewidth{0.200750pt}%
\definecolor{currentstroke}{rgb}{0.200000,0.200000,0.200000}%
\pgfsetstrokecolor{currentstroke}%
\pgfsetdash{}{0pt}%
\pgfpathmoveto{\pgfqpoint{3.035498in}{2.422123in}}%
\pgfpathlineto{\pgfqpoint{3.007602in}{2.495450in}}%
\pgfpathlineto{\pgfqpoint{2.979735in}{2.568993in}}%
\pgfpathlineto{\pgfqpoint{3.027222in}{2.583719in}}%
\pgfpathlineto{\pgfqpoint{3.074625in}{2.598425in}}%
\pgfpathlineto{\pgfqpoint{3.102489in}{2.525125in}}%
\pgfpathlineto{\pgfqpoint{3.130384in}{2.452063in}}%
\pgfpathlineto{\pgfqpoint{3.082983in}{2.437102in}}%
\pgfpathlineto{\pgfqpoint{3.035498in}{2.422123in}}%
\pgfpathclose%
\pgfusepath{stroke,fill}%
\end{pgfscope}%
\begin{pgfscope}%
\pgfpathrectangle{\pgfqpoint{0.050000in}{0.083252in}}{\pgfqpoint{4.900000in}{4.900000in}}%
\pgfusepath{clip}%
\pgfsetbuttcap%
\pgfsetroundjoin%
\definecolor{currentfill}{rgb}{0.267636,0.509804,0.946528}%
\pgfsetfillcolor{currentfill}%
\pgfsetfillopacity{0.950000}%
\pgfsetlinewidth{0.200750pt}%
\definecolor{currentstroke}{rgb}{0.200000,0.200000,0.200000}%
\pgfsetstrokecolor{currentstroke}%
\pgfsetdash{}{0pt}%
\pgfpathmoveto{\pgfqpoint{1.851089in}{3.878923in}}%
\pgfpathlineto{\pgfqpoint{1.823753in}{3.952612in}}%
\pgfpathlineto{\pgfqpoint{1.796442in}{4.026232in}}%
\pgfpathlineto{\pgfqpoint{1.844991in}{4.039511in}}%
\pgfpathlineto{\pgfqpoint{1.893453in}{4.052766in}}%
\pgfpathlineto{\pgfqpoint{1.920761in}{3.979289in}}%
\pgfpathlineto{\pgfqpoint{1.948095in}{3.905743in}}%
\pgfpathlineto{\pgfqpoint{1.899635in}{3.892345in}}%
\pgfpathlineto{\pgfqpoint{1.851089in}{3.878923in}}%
\pgfpathclose%
\pgfusepath{stroke,fill}%
\end{pgfscope}%
\begin{pgfscope}%
\pgfpathrectangle{\pgfqpoint{0.050000in}{0.083252in}}{\pgfqpoint{4.900000in}{4.900000in}}%
\pgfusepath{clip}%
\pgfsetbuttcap%
\pgfsetroundjoin%
\definecolor{currentfill}{rgb}{0.576348,0.558877,0.746682}%
\pgfsetfillcolor{currentfill}%
\pgfsetfillopacity{0.950000}%
\pgfsetlinewidth{0.200750pt}%
\definecolor{currentstroke}{rgb}{0.200000,0.200000,0.200000}%
\pgfsetstrokecolor{currentstroke}%
\pgfsetdash{}{0pt}%
\pgfpathmoveto{\pgfqpoint{3.242322in}{2.164439in}}%
\pgfpathlineto{\pgfqpoint{3.214278in}{2.235368in}}%
\pgfpathlineto{\pgfqpoint{3.186276in}{2.307055in}}%
\pgfpathlineto{\pgfqpoint{3.233595in}{2.322296in}}%
\pgfpathlineto{\pgfqpoint{3.280831in}{2.337516in}}%
\pgfpathlineto{\pgfqpoint{3.308837in}{2.266159in}}%
\pgfpathlineto{\pgfqpoint{3.336887in}{2.195578in}}%
\pgfpathlineto{\pgfqpoint{3.289646in}{2.180020in}}%
\pgfpathlineto{\pgfqpoint{3.242322in}{2.164439in}}%
\pgfpathclose%
\pgfusepath{stroke,fill}%
\end{pgfscope}%
\begin{pgfscope}%
\pgfpathrectangle{\pgfqpoint{0.050000in}{0.083252in}}{\pgfqpoint{4.900000in}{4.900000in}}%
\pgfusepath{clip}%
\pgfsetbuttcap%
\pgfsetroundjoin%
\definecolor{currentfill}{rgb}{0.260992,0.439216,0.852687}%
\pgfsetfillcolor{currentfill}%
\pgfsetfillopacity{0.950000}%
\pgfsetlinewidth{0.200750pt}%
\definecolor{currentstroke}{rgb}{0.200000,0.200000,0.200000}%
\pgfsetstrokecolor{currentstroke}%
\pgfsetdash{}{0pt}%
\pgfpathmoveto{\pgfqpoint{2.057686in}{3.610867in}}%
\pgfpathlineto{\pgfqpoint{2.030250in}{3.684690in}}%
\pgfpathlineto{\pgfqpoint{2.002839in}{3.758444in}}%
\pgfpathlineto{\pgfqpoint{2.051210in}{3.771961in}}%
\pgfpathlineto{\pgfqpoint{2.099495in}{3.785455in}}%
\pgfpathlineto{\pgfqpoint{2.126902in}{3.711845in}}%
\pgfpathlineto{\pgfqpoint{2.154335in}{3.638166in}}%
\pgfpathlineto{\pgfqpoint{2.106054in}{3.624529in}}%
\pgfpathlineto{\pgfqpoint{2.057686in}{3.610867in}}%
\pgfpathclose%
\pgfusepath{stroke,fill}%
\end{pgfscope}%
\begin{pgfscope}%
\pgfpathrectangle{\pgfqpoint{0.050000in}{0.083252in}}{\pgfqpoint{4.900000in}{4.900000in}}%
\pgfusepath{clip}%
\pgfsetbuttcap%
\pgfsetroundjoin%
\definecolor{currentfill}{rgb}{0.254348,0.368627,0.758847}%
\pgfsetfillcolor{currentfill}%
\pgfsetfillopacity{0.950000}%
\pgfsetlinewidth{0.200750pt}%
\definecolor{currentstroke}{rgb}{0.200000,0.200000,0.200000}%
\pgfsetstrokecolor{currentstroke}%
\pgfsetdash{}{0pt}%
\pgfpathmoveto{\pgfqpoint{2.264326in}{3.342757in}}%
\pgfpathlineto{\pgfqpoint{2.236789in}{3.416713in}}%
\pgfpathlineto{\pgfqpoint{2.209279in}{3.490600in}}%
\pgfpathlineto{\pgfqpoint{2.257471in}{3.504357in}}%
\pgfpathlineto{\pgfqpoint{2.305578in}{3.518089in}}%
\pgfpathlineto{\pgfqpoint{2.333085in}{3.444347in}}%
\pgfpathlineto{\pgfqpoint{2.360618in}{3.370535in}}%
\pgfpathlineto{\pgfqpoint{2.312515in}{3.356658in}}%
\pgfpathlineto{\pgfqpoint{2.264326in}{3.342757in}}%
\pgfpathclose%
\pgfusepath{stroke,fill}%
\end{pgfscope}%
\begin{pgfscope}%
\pgfpathrectangle{\pgfqpoint{0.050000in}{0.083252in}}{\pgfqpoint{4.900000in}{4.900000in}}%
\pgfusepath{clip}%
\pgfsetbuttcap%
\pgfsetroundjoin%
\definecolor{currentfill}{rgb}{0.832926,0.826036,0.900100}%
\pgfsetfillcolor{currentfill}%
\pgfsetfillopacity{0.950000}%
\pgfsetlinewidth{0.200750pt}%
\definecolor{currentstroke}{rgb}{0.200000,0.200000,0.200000}%
\pgfsetstrokecolor{currentstroke}%
\pgfsetdash{}{0pt}%
\pgfpathmoveto{\pgfqpoint{3.752149in}{1.751338in}}%
\pgfpathlineto{\pgfqpoint{3.723425in}{1.807466in}}%
\pgfpathlineto{\pgfqpoint{3.694815in}{1.866041in}}%
\pgfpathlineto{\pgfqpoint{3.741699in}{1.882381in}}%
\pgfpathlineto{\pgfqpoint{3.788503in}{1.898679in}}%
\pgfpathlineto{\pgfqpoint{3.817142in}{1.840297in}}%
\pgfpathlineto{\pgfqpoint{3.845900in}{1.784305in}}%
\pgfpathlineto{\pgfqpoint{3.799065in}{1.767844in}}%
\pgfpathlineto{\pgfqpoint{3.752149in}{1.751338in}}%
\pgfpathclose%
\pgfusepath{stroke,fill}%
\end{pgfscope}%
\begin{pgfscope}%
\pgfpathrectangle{\pgfqpoint{0.050000in}{0.083252in}}{\pgfqpoint{4.900000in}{4.900000in}}%
\pgfusepath{clip}%
\pgfsetbuttcap%
\pgfsetroundjoin%
\definecolor{currentfill}{rgb}{0.874694,0.869527,0.925075}%
\pgfsetfillcolor{currentfill}%
\pgfsetfillopacity{0.950000}%
\pgfsetlinewidth{0.200750pt}%
\definecolor{currentstroke}{rgb}{0.200000,0.200000,0.200000}%
\pgfsetstrokecolor{currentstroke}%
\pgfsetdash{}{0pt}%
\pgfpathmoveto{\pgfqpoint{3.903805in}{1.680005in}}%
\pgfpathlineto{\pgfqpoint{3.874784in}{1.730843in}}%
\pgfpathlineto{\pgfqpoint{3.845900in}{1.784305in}}%
\pgfpathlineto{\pgfqpoint{3.892653in}{1.800721in}}%
\pgfpathlineto{\pgfqpoint{3.939326in}{1.817093in}}%
\pgfpathlineto{\pgfqpoint{3.968247in}{1.763705in}}%
\pgfpathlineto{\pgfqpoint{3.997307in}{1.712879in}}%
\pgfpathlineto{\pgfqpoint{3.950596in}{1.696464in}}%
\pgfpathlineto{\pgfqpoint{3.903805in}{1.680005in}}%
\pgfpathclose%
\pgfusepath{stroke,fill}%
\end{pgfscope}%
\begin{pgfscope}%
\pgfpathrectangle{\pgfqpoint{0.050000in}{0.083252in}}{\pgfqpoint{4.900000in}{4.900000in}}%
\pgfusepath{clip}%
\pgfsetbuttcap%
\pgfsetroundjoin%
\definecolor{currentfill}{rgb}{0.247951,0.300654,0.668481}%
\pgfsetfillcolor{currentfill}%
\pgfsetfillopacity{0.950000}%
\pgfsetlinewidth{0.200750pt}%
\definecolor{currentstroke}{rgb}{0.200000,0.200000,0.200000}%
\pgfsetstrokecolor{currentstroke}%
\pgfsetdash{}{0pt}%
\pgfpathmoveto{\pgfqpoint{2.471008in}{3.074609in}}%
\pgfpathlineto{\pgfqpoint{2.443372in}{3.148691in}}%
\pgfpathlineto{\pgfqpoint{2.415761in}{3.222707in}}%
\pgfpathlineto{\pgfqpoint{2.463775in}{3.236705in}}%
\pgfpathlineto{\pgfqpoint{2.511704in}{3.250678in}}%
\pgfpathlineto{\pgfqpoint{2.539310in}{3.176809in}}%
\pgfpathlineto{\pgfqpoint{2.566943in}{3.102876in}}%
\pgfpathlineto{\pgfqpoint{2.519018in}{3.088754in}}%
\pgfpathlineto{\pgfqpoint{2.471008in}{3.074609in}}%
\pgfpathclose%
\pgfusepath{stroke,fill}%
\end{pgfscope}%
\begin{pgfscope}%
\pgfpathrectangle{\pgfqpoint{0.050000in}{0.083252in}}{\pgfqpoint{4.900000in}{4.900000in}}%
\pgfusepath{clip}%
\pgfsetbuttcap%
\pgfsetroundjoin%
\definecolor{currentfill}{rgb}{0.779223,0.770119,0.867989}%
\pgfsetfillcolor{currentfill}%
\pgfsetfillopacity{0.950000}%
\pgfsetlinewidth{0.200750pt}%
\definecolor{currentstroke}{rgb}{0.200000,0.200000,0.200000}%
\pgfsetstrokecolor{currentstroke}%
\pgfsetdash{}{0pt}%
\pgfpathmoveto{\pgfqpoint{3.600802in}{1.833231in}}%
\pgfpathlineto{\pgfqpoint{3.572323in}{1.894323in}}%
\pgfpathlineto{\pgfqpoint{3.543936in}{1.957514in}}%
\pgfpathlineto{\pgfqpoint{3.590958in}{1.973668in}}%
\pgfpathlineto{\pgfqpoint{3.637900in}{1.989784in}}%
\pgfpathlineto{\pgfqpoint{3.666310in}{1.926882in}}%
\pgfpathlineto{\pgfqpoint{3.694815in}{1.866041in}}%
\pgfpathlineto{\pgfqpoint{3.647849in}{1.849658in}}%
\pgfpathlineto{\pgfqpoint{3.600802in}{1.833231in}}%
\pgfpathclose%
\pgfusepath{stroke,fill}%
\end{pgfscope}%
\begin{pgfscope}%
\pgfpathrectangle{\pgfqpoint{0.050000in}{0.083252in}}{\pgfqpoint{4.900000in}{4.900000in}}%
\pgfusepath{clip}%
\pgfsetbuttcap%
\pgfsetroundjoin%
\definecolor{currentfill}{rgb}{0.241307,0.230065,0.574641}%
\pgfsetfillcolor{currentfill}%
\pgfsetfillopacity{0.950000}%
\pgfsetlinewidth{0.200750pt}%
\definecolor{currentstroke}{rgb}{0.200000,0.200000,0.200000}%
\pgfsetstrokecolor{currentstroke}%
\pgfsetdash{}{0pt}%
\pgfpathmoveto{\pgfqpoint{2.677733in}{2.806571in}}%
\pgfpathlineto{\pgfqpoint{2.649997in}{2.880720in}}%
\pgfpathlineto{\pgfqpoint{2.622286in}{2.954826in}}%
\pgfpathlineto{\pgfqpoint{2.670121in}{2.969077in}}%
\pgfpathlineto{\pgfqpoint{2.717872in}{2.983304in}}%
\pgfpathlineto{\pgfqpoint{2.745579in}{2.909358in}}%
\pgfpathlineto{\pgfqpoint{2.773312in}{2.835375in}}%
\pgfpathlineto{\pgfqpoint{2.725565in}{2.820984in}}%
\pgfpathlineto{\pgfqpoint{2.677733in}{2.806571in}}%
\pgfpathclose%
\pgfusepath{stroke,fill}%
\end{pgfscope}%
\begin{pgfscope}%
\pgfpathrectangle{\pgfqpoint{0.050000in}{0.083252in}}{\pgfqpoint{4.900000in}{4.900000in}}%
\pgfusepath{clip}%
\pgfsetbuttcap%
\pgfsetroundjoin%
\definecolor{currentfill}{rgb}{0.349604,0.322784,0.611103}%
\pgfsetfillcolor{currentfill}%
\pgfsetfillopacity{0.950000}%
\pgfsetlinewidth{0.200750pt}%
\definecolor{currentstroke}{rgb}{0.200000,0.200000,0.200000}%
\pgfsetstrokecolor{currentstroke}%
\pgfsetdash{}{0pt}%
\pgfpathmoveto{\pgfqpoint{2.884510in}{2.539484in}}%
\pgfpathlineto{\pgfqpoint{2.856670in}{2.613391in}}%
\pgfpathlineto{\pgfqpoint{2.828857in}{2.687366in}}%
\pgfpathlineto{\pgfqpoint{2.876515in}{2.701919in}}%
\pgfpathlineto{\pgfqpoint{2.924088in}{2.716450in}}%
\pgfpathlineto{\pgfqpoint{2.951898in}{2.642679in}}%
\pgfpathlineto{\pgfqpoint{2.979735in}{2.568993in}}%
\pgfpathlineto{\pgfqpoint{2.932165in}{2.554248in}}%
\pgfpathlineto{\pgfqpoint{2.884510in}{2.539484in}}%
\pgfpathclose%
\pgfusepath{stroke,fill}%
\end{pgfscope}%
\begin{pgfscope}%
\pgfpathrectangle{\pgfqpoint{0.050000in}{0.083252in}}{\pgfqpoint{4.900000in}{4.900000in}}%
\pgfusepath{clip}%
\pgfsetbuttcap%
\pgfsetroundjoin%
\definecolor{currentfill}{rgb}{0.916463,0.913018,0.950050}%
\pgfsetfillcolor{currentfill}%
\pgfsetfillopacity{0.950000}%
\pgfsetlinewidth{0.200750pt}%
\definecolor{currentstroke}{rgb}{0.200000,0.200000,0.200000}%
\pgfsetstrokecolor{currentstroke}%
\pgfsetdash{}{0pt}%
\pgfpathmoveto{\pgfqpoint{4.055871in}{1.619027in}}%
\pgfpathlineto{\pgfqpoint{4.026513in}{1.664657in}}%
\pgfpathlineto{\pgfqpoint{3.997307in}{1.712879in}}%
\pgfpathlineto{\pgfqpoint{4.043936in}{1.729251in}}%
\pgfpathlineto{\pgfqpoint{4.090485in}{1.745580in}}%
\pgfpathlineto{\pgfqpoint{4.119733in}{1.697312in}}%
\pgfpathlineto{\pgfqpoint{4.149137in}{1.651578in}}%
\pgfpathlineto{\pgfqpoint{4.102545in}{1.635322in}}%
\pgfpathlineto{\pgfqpoint{4.055871in}{1.619027in}}%
\pgfpathclose%
\pgfusepath{stroke,fill}%
\end{pgfscope}%
\begin{pgfscope}%
\pgfpathrectangle{\pgfqpoint{0.050000in}{0.083252in}}{\pgfqpoint{4.900000in}{4.900000in}}%
\pgfusepath{clip}%
\pgfsetbuttcap%
\pgfsetroundjoin%
\definecolor{currentfill}{rgb}{0.719554,0.707989,0.832311}%
\pgfsetfillcolor{currentfill}%
\pgfsetfillopacity{0.950000}%
\pgfsetlinewidth{0.200750pt}%
\definecolor{currentstroke}{rgb}{0.200000,0.200000,0.200000}%
\pgfsetstrokecolor{currentstroke}%
\pgfsetdash{}{0pt}%
\pgfpathmoveto{\pgfqpoint{3.449645in}{1.925093in}}%
\pgfpathlineto{\pgfqpoint{3.421358in}{1.990482in}}%
\pgfpathlineto{\pgfqpoint{3.393141in}{2.057519in}}%
\pgfpathlineto{\pgfqpoint{3.440311in}{2.073406in}}%
\pgfpathlineto{\pgfqpoint{3.487398in}{2.089262in}}%
\pgfpathlineto{\pgfqpoint{3.515630in}{2.022573in}}%
\pgfpathlineto{\pgfqpoint{3.543936in}{1.957514in}}%
\pgfpathlineto{\pgfqpoint{3.496831in}{1.941322in}}%
\pgfpathlineto{\pgfqpoint{3.449645in}{1.925093in}}%
\pgfpathclose%
\pgfusepath{stroke,fill}%
\end{pgfscope}%
\begin{pgfscope}%
\pgfpathrectangle{\pgfqpoint{0.050000in}{0.083252in}}{\pgfqpoint{4.900000in}{4.900000in}}%
\pgfusepath{clip}%
\pgfsetbuttcap%
\pgfsetroundjoin%
\definecolor{currentfill}{rgb}{0.504744,0.484321,0.703868}%
\pgfsetfillcolor{currentfill}%
\pgfsetfillopacity{0.950000}%
\pgfsetlinewidth{0.200750pt}%
\definecolor{currentstroke}{rgb}{0.200000,0.200000,0.200000}%
\pgfsetstrokecolor{currentstroke}%
\pgfsetdash{}{0pt}%
\pgfpathmoveto{\pgfqpoint{3.091388in}{2.276514in}}%
\pgfpathlineto{\pgfqpoint{3.063426in}{2.349105in}}%
\pgfpathlineto{\pgfqpoint{3.035498in}{2.422123in}}%
\pgfpathlineto{\pgfqpoint{3.082983in}{2.437102in}}%
\pgfpathlineto{\pgfqpoint{3.130384in}{2.452063in}}%
\pgfpathlineto{\pgfqpoint{3.158312in}{2.379334in}}%
\pgfpathlineto{\pgfqpoint{3.186276in}{2.307055in}}%
\pgfpathlineto{\pgfqpoint{3.138874in}{2.291794in}}%
\pgfpathlineto{\pgfqpoint{3.091388in}{2.276514in}}%
\pgfpathclose%
\pgfusepath{stroke,fill}%
\end{pgfscope}%
\begin{pgfscope}%
\pgfpathrectangle{\pgfqpoint{0.050000in}{0.083252in}}{\pgfqpoint{4.900000in}{4.900000in}}%
\pgfusepath{clip}%
\pgfsetbuttcap%
\pgfsetroundjoin%
\definecolor{currentfill}{rgb}{0.264191,0.473203,0.897870}%
\pgfsetfillcolor{currentfill}%
\pgfsetfillopacity{0.950000}%
\pgfsetlinewidth{0.200750pt}%
\definecolor{currentstroke}{rgb}{0.200000,0.200000,0.200000}%
\pgfsetstrokecolor{currentstroke}%
\pgfsetdash{}{0pt}%
\pgfpathmoveto{\pgfqpoint{1.905839in}{3.731336in}}%
\pgfpathlineto{\pgfqpoint{1.878451in}{3.805164in}}%
\pgfpathlineto{\pgfqpoint{1.851089in}{3.878923in}}%
\pgfpathlineto{\pgfqpoint{1.899635in}{3.892345in}}%
\pgfpathlineto{\pgfqpoint{1.948095in}{3.905743in}}%
\pgfpathlineto{\pgfqpoint{1.975454in}{3.832128in}}%
\pgfpathlineto{\pgfqpoint{2.002839in}{3.758444in}}%
\pgfpathlineto{\pgfqpoint{1.954382in}{3.744902in}}%
\pgfpathlineto{\pgfqpoint{1.905839in}{3.731336in}}%
\pgfpathclose%
\pgfusepath{stroke,fill}%
\end{pgfscope}%
\begin{pgfscope}%
\pgfpathrectangle{\pgfqpoint{0.050000in}{0.083252in}}{\pgfqpoint{4.900000in}{4.900000in}}%
\pgfusepath{clip}%
\pgfsetbuttcap%
\pgfsetroundjoin%
\definecolor{currentfill}{rgb}{0.257793,0.405229,0.807505}%
\pgfsetfillcolor{currentfill}%
\pgfsetfillopacity{0.950000}%
\pgfsetlinewidth{0.200750pt}%
\definecolor{currentstroke}{rgb}{0.200000,0.200000,0.200000}%
\pgfsetstrokecolor{currentstroke}%
\pgfsetdash{}{0pt}%
\pgfpathmoveto{\pgfqpoint{2.112637in}{3.463013in}}%
\pgfpathlineto{\pgfqpoint{2.085149in}{3.536975in}}%
\pgfpathlineto{\pgfqpoint{2.057686in}{3.610867in}}%
\pgfpathlineto{\pgfqpoint{2.106054in}{3.624529in}}%
\pgfpathlineto{\pgfqpoint{2.154335in}{3.638166in}}%
\pgfpathlineto{\pgfqpoint{2.181794in}{3.564418in}}%
\pgfpathlineto{\pgfqpoint{2.209279in}{3.490600in}}%
\pgfpathlineto{\pgfqpoint{2.161001in}{3.476819in}}%
\pgfpathlineto{\pgfqpoint{2.112637in}{3.463013in}}%
\pgfpathclose%
\pgfusepath{stroke,fill}%
\end{pgfscope}%
\begin{pgfscope}%
\pgfpathrectangle{\pgfqpoint{0.050000in}{0.083252in}}{\pgfqpoint{4.900000in}{4.900000in}}%
\pgfusepath{clip}%
\pgfsetbuttcap%
\pgfsetroundjoin%
\definecolor{currentfill}{rgb}{0.270096,0.535948,0.981284}%
\pgfsetfillcolor{currentfill}%
\pgfsetfillopacity{0.950000}%
\pgfsetlinewidth{0.200750pt}%
\definecolor{currentstroke}{rgb}{0.200000,0.200000,0.200000}%
\pgfsetstrokecolor{currentstroke}%
\pgfsetdash{}{0pt}%
\pgfpathmoveto{\pgfqpoint{1.699085in}{3.999603in}}%
\pgfpathlineto{\pgfqpoint{1.671797in}{4.073297in}}%
\pgfpathlineto{\pgfqpoint{1.720520in}{4.086552in}}%
\pgfpathlineto{\pgfqpoint{1.769157in}{4.099783in}}%
\pgfpathlineto{\pgfqpoint{1.796442in}{4.026232in}}%
\pgfpathlineto{\pgfqpoint{1.747806in}{4.012929in}}%
\pgfpathlineto{\pgfqpoint{1.699085in}{3.999603in}}%
\pgfpathclose%
\pgfusepath{stroke,fill}%
\end{pgfscope}%
\begin{pgfscope}%
\pgfpathrectangle{\pgfqpoint{0.050000in}{0.083252in}}{\pgfqpoint{4.900000in}{4.900000in}}%
\pgfusepath{clip}%
\pgfsetbuttcap%
\pgfsetroundjoin%
\definecolor{currentfill}{rgb}{0.251150,0.334641,0.713664}%
\pgfsetfillcolor{currentfill}%
\pgfsetfillopacity{0.950000}%
\pgfsetlinewidth{0.200750pt}%
\definecolor{currentstroke}{rgb}{0.200000,0.200000,0.200000}%
\pgfsetstrokecolor{currentstroke}%
\pgfsetdash{}{0pt}%
\pgfpathmoveto{\pgfqpoint{2.319477in}{3.194637in}}%
\pgfpathlineto{\pgfqpoint{2.291888in}{3.268732in}}%
\pgfpathlineto{\pgfqpoint{2.264326in}{3.342757in}}%
\pgfpathlineto{\pgfqpoint{2.312515in}{3.356658in}}%
\pgfpathlineto{\pgfqpoint{2.360618in}{3.370535in}}%
\pgfpathlineto{\pgfqpoint{2.388176in}{3.296655in}}%
\pgfpathlineto{\pgfqpoint{2.415761in}{3.222707in}}%
\pgfpathlineto{\pgfqpoint{2.367662in}{3.208685in}}%
\pgfpathlineto{\pgfqpoint{2.319477in}{3.194637in}}%
\pgfpathclose%
\pgfusepath{stroke,fill}%
\end{pgfscope}%
\begin{pgfscope}%
\pgfpathrectangle{\pgfqpoint{0.050000in}{0.083252in}}{\pgfqpoint{4.900000in}{4.900000in}}%
\pgfusepath{clip}%
\pgfsetbuttcap%
\pgfsetroundjoin%
\definecolor{currentfill}{rgb}{0.653918,0.639646,0.793064}%
\pgfsetfillcolor{currentfill}%
\pgfsetfillopacity{0.950000}%
\pgfsetlinewidth{0.200750pt}%
\definecolor{currentstroke}{rgb}{0.200000,0.200000,0.200000}%
\pgfsetstrokecolor{currentstroke}%
\pgfsetdash{}{0pt}%
\pgfpathmoveto{\pgfqpoint{3.298555in}{2.025654in}}%
\pgfpathlineto{\pgfqpoint{3.270412in}{2.094461in}}%
\pgfpathlineto{\pgfqpoint{3.242322in}{2.164439in}}%
\pgfpathlineto{\pgfqpoint{3.289646in}{2.180020in}}%
\pgfpathlineto{\pgfqpoint{3.336887in}{2.195578in}}%
\pgfpathlineto{\pgfqpoint{3.364986in}{2.125961in}}%
\pgfpathlineto{\pgfqpoint{3.393141in}{2.057519in}}%
\pgfpathlineto{\pgfqpoint{3.345889in}{2.041602in}}%
\pgfpathlineto{\pgfqpoint{3.298555in}{2.025654in}}%
\pgfpathclose%
\pgfusepath{stroke,fill}%
\end{pgfscope}%
\begin{pgfscope}%
\pgfpathrectangle{\pgfqpoint{0.050000in}{0.083252in}}{\pgfqpoint{4.900000in}{4.900000in}}%
\pgfusepath{clip}%
\pgfsetbuttcap%
\pgfsetroundjoin%
\definecolor{currentfill}{rgb}{0.244506,0.264052,0.619823}%
\pgfsetfillcolor{currentfill}%
\pgfsetfillopacity{0.950000}%
\pgfsetlinewidth{0.200750pt}%
\definecolor{currentstroke}{rgb}{0.200000,0.200000,0.200000}%
\pgfsetstrokecolor{currentstroke}%
\pgfsetdash{}{0pt}%
\pgfpathmoveto{\pgfqpoint{2.526359in}{2.926254in}}%
\pgfpathlineto{\pgfqpoint{2.498671in}{3.000462in}}%
\pgfpathlineto{\pgfqpoint{2.471008in}{3.074609in}}%
\pgfpathlineto{\pgfqpoint{2.519018in}{3.088754in}}%
\pgfpathlineto{\pgfqpoint{2.566943in}{3.102876in}}%
\pgfpathlineto{\pgfqpoint{2.594601in}{3.028880in}}%
\pgfpathlineto{\pgfqpoint{2.622286in}{2.954826in}}%
\pgfpathlineto{\pgfqpoint{2.574365in}{2.940552in}}%
\pgfpathlineto{\pgfqpoint{2.526359in}{2.926254in}}%
\pgfpathclose%
\pgfusepath{stroke,fill}%
\end{pgfscope}%
\begin{pgfscope}%
\pgfpathrectangle{\pgfqpoint{0.050000in}{0.083252in}}{\pgfqpoint{4.900000in}{4.900000in}}%
\pgfusepath{clip}%
\pgfsetbuttcap%
\pgfsetroundjoin%
\definecolor{currentfill}{rgb}{0.272034,0.242015,0.564721}%
\pgfsetfillcolor{currentfill}%
\pgfsetfillopacity{0.950000}%
\pgfsetlinewidth{0.200750pt}%
\definecolor{currentstroke}{rgb}{0.200000,0.200000,0.200000}%
\pgfsetstrokecolor{currentstroke}%
\pgfsetdash{}{0pt}%
\pgfpathmoveto{\pgfqpoint{2.733286in}{2.658200in}}%
\pgfpathlineto{\pgfqpoint{2.705497in}{2.732391in}}%
\pgfpathlineto{\pgfqpoint{2.677733in}{2.806571in}}%
\pgfpathlineto{\pgfqpoint{2.725565in}{2.820984in}}%
\pgfpathlineto{\pgfqpoint{2.773312in}{2.835375in}}%
\pgfpathlineto{\pgfqpoint{2.801071in}{2.761371in}}%
\pgfpathlineto{\pgfqpoint{2.828857in}{2.687366in}}%
\pgfpathlineto{\pgfqpoint{2.781114in}{2.672793in}}%
\pgfpathlineto{\pgfqpoint{2.733286in}{2.658200in}}%
\pgfpathclose%
\pgfusepath{stroke,fill}%
\end{pgfscope}%
\begin{pgfscope}%
\pgfpathrectangle{\pgfqpoint{0.050000in}{0.083252in}}{\pgfqpoint{4.900000in}{4.900000in}}%
\pgfusepath{clip}%
\pgfsetbuttcap%
\pgfsetroundjoin%
\definecolor{currentfill}{rgb}{0.427174,0.403552,0.657486}%
\pgfsetfillcolor{currentfill}%
\pgfsetfillopacity{0.950000}%
\pgfsetlinewidth{0.200750pt}%
\definecolor{currentstroke}{rgb}{0.200000,0.200000,0.200000}%
\pgfsetstrokecolor{currentstroke}%
\pgfsetdash{}{0pt}%
\pgfpathmoveto{\pgfqpoint{2.940277in}{2.392108in}}%
\pgfpathlineto{\pgfqpoint{2.912379in}{2.465699in}}%
\pgfpathlineto{\pgfqpoint{2.884510in}{2.539484in}}%
\pgfpathlineto{\pgfqpoint{2.932165in}{2.554248in}}%
\pgfpathlineto{\pgfqpoint{2.979735in}{2.568993in}}%
\pgfpathlineto{\pgfqpoint{3.007602in}{2.495450in}}%
\pgfpathlineto{\pgfqpoint{3.035498in}{2.422123in}}%
\pgfpathlineto{\pgfqpoint{2.987929in}{2.407125in}}%
\pgfpathlineto{\pgfqpoint{2.940277in}{2.392108in}}%
\pgfpathclose%
\pgfusepath{stroke,fill}%
\end{pgfscope}%
\begin{pgfscope}%
\pgfpathrectangle{\pgfqpoint{0.050000in}{0.083252in}}{\pgfqpoint{4.900000in}{4.900000in}}%
\pgfusepath{clip}%
\pgfsetbuttcap%
\pgfsetroundjoin%
\definecolor{currentfill}{rgb}{0.946298,0.944083,0.967889}%
\pgfsetfillcolor{currentfill}%
\pgfsetfillopacity{0.950000}%
\pgfsetlinewidth{0.200750pt}%
\definecolor{currentstroke}{rgb}{0.200000,0.200000,0.200000}%
\pgfsetstrokecolor{currentstroke}%
\pgfsetdash{}{0pt}%
\pgfpathmoveto{\pgfqpoint{4.208418in}{1.567380in}}%
\pgfpathlineto{\pgfqpoint{4.178698in}{1.608310in}}%
\pgfpathlineto{\pgfqpoint{4.149137in}{1.651578in}}%
\pgfpathlineto{\pgfqpoint{4.195647in}{1.667794in}}%
\pgfpathlineto{\pgfqpoint{4.225232in}{1.624450in}}%
\pgfpathlineto{\pgfqpoint{4.254977in}{1.583421in}}%
\pgfpathlineto{\pgfqpoint{4.208418in}{1.567380in}}%
\pgfpathclose%
\pgfusepath{stroke,fill}%
\end{pgfscope}%
\begin{pgfscope}%
\pgfpathrectangle{\pgfqpoint{0.050000in}{0.083252in}}{\pgfqpoint{4.900000in}{4.900000in}}%
\pgfusepath{clip}%
\pgfsetbuttcap%
\pgfsetroundjoin%
\definecolor{currentfill}{rgb}{0.267636,0.509804,0.946528}%
\pgfsetfillcolor{currentfill}%
\pgfsetfillopacity{0.950000}%
\pgfsetlinewidth{0.200750pt}%
\definecolor{currentstroke}{rgb}{0.200000,0.200000,0.200000}%
\pgfsetstrokecolor{currentstroke}%
\pgfsetdash{}{0pt}%
\pgfpathmoveto{\pgfqpoint{1.753738in}{3.852006in}}%
\pgfpathlineto{\pgfqpoint{1.726398in}{3.925839in}}%
\pgfpathlineto{\pgfqpoint{1.699085in}{3.999603in}}%
\pgfpathlineto{\pgfqpoint{1.747806in}{4.012929in}}%
\pgfpathlineto{\pgfqpoint{1.796442in}{4.026232in}}%
\pgfpathlineto{\pgfqpoint{1.823753in}{3.952612in}}%
\pgfpathlineto{\pgfqpoint{1.851089in}{3.878923in}}%
\pgfpathlineto{\pgfqpoint{1.802457in}{3.865476in}}%
\pgfpathlineto{\pgfqpoint{1.753738in}{3.852006in}}%
\pgfpathclose%
\pgfusepath{stroke,fill}%
\end{pgfscope}%
\begin{pgfscope}%
\pgfpathrectangle{\pgfqpoint{0.050000in}{0.083252in}}{\pgfqpoint{4.900000in}{4.900000in}}%
\pgfusepath{clip}%
\pgfsetbuttcap%
\pgfsetroundjoin%
\definecolor{currentfill}{rgb}{0.582314,0.565090,0.750250}%
\pgfsetfillcolor{currentfill}%
\pgfsetfillopacity{0.950000}%
\pgfsetlinewidth{0.200750pt}%
\definecolor{currentstroke}{rgb}{0.200000,0.200000,0.200000}%
\pgfsetstrokecolor{currentstroke}%
\pgfsetdash{}{0pt}%
\pgfpathmoveto{\pgfqpoint{3.147425in}{2.133204in}}%
\pgfpathlineto{\pgfqpoint{3.119387in}{2.204490in}}%
\pgfpathlineto{\pgfqpoint{3.091388in}{2.276514in}}%
\pgfpathlineto{\pgfqpoint{3.138874in}{2.291794in}}%
\pgfpathlineto{\pgfqpoint{3.186276in}{2.307055in}}%
\pgfpathlineto{\pgfqpoint{3.214278in}{2.235368in}}%
\pgfpathlineto{\pgfqpoint{3.242322in}{2.164439in}}%
\pgfpathlineto{\pgfqpoint{3.194915in}{2.148833in}}%
\pgfpathlineto{\pgfqpoint{3.147425in}{2.133204in}}%
\pgfpathclose%
\pgfusepath{stroke,fill}%
\end{pgfscope}%
\begin{pgfscope}%
\pgfpathrectangle{\pgfqpoint{0.050000in}{0.083252in}}{\pgfqpoint{4.900000in}{4.900000in}}%
\pgfusepath{clip}%
\pgfsetbuttcap%
\pgfsetroundjoin%
\definecolor{currentfill}{rgb}{0.260992,0.439216,0.852687}%
\pgfsetfillcolor{currentfill}%
\pgfsetfillopacity{0.950000}%
\pgfsetlinewidth{0.200750pt}%
\definecolor{currentstroke}{rgb}{0.200000,0.200000,0.200000}%
\pgfsetstrokecolor{currentstroke}%
\pgfsetdash{}{0pt}%
\pgfpathmoveto{\pgfqpoint{1.960693in}{3.583470in}}%
\pgfpathlineto{\pgfqpoint{1.933253in}{3.657438in}}%
\pgfpathlineto{\pgfqpoint{1.905839in}{3.731336in}}%
\pgfpathlineto{\pgfqpoint{1.954382in}{3.744902in}}%
\pgfpathlineto{\pgfqpoint{2.002839in}{3.758444in}}%
\pgfpathlineto{\pgfqpoint{2.030250in}{3.684690in}}%
\pgfpathlineto{\pgfqpoint{2.057686in}{3.610867in}}%
\pgfpathlineto{\pgfqpoint{2.009233in}{3.597181in}}%
\pgfpathlineto{\pgfqpoint{1.960693in}{3.583470in}}%
\pgfpathclose%
\pgfusepath{stroke,fill}%
\end{pgfscope}%
\begin{pgfscope}%
\pgfpathrectangle{\pgfqpoint{0.050000in}{0.083252in}}{\pgfqpoint{4.900000in}{4.900000in}}%
\pgfusepath{clip}%
\pgfsetbuttcap%
\pgfsetroundjoin%
\definecolor{currentfill}{rgb}{0.832926,0.826036,0.900100}%
\pgfsetfillcolor{currentfill}%
\pgfsetfillopacity{0.950000}%
\pgfsetlinewidth{0.200750pt}%
\definecolor{currentstroke}{rgb}{0.200000,0.200000,0.200000}%
\pgfsetstrokecolor{currentstroke}%
\pgfsetdash{}{0pt}%
\pgfpathmoveto{\pgfqpoint{3.658072in}{1.718185in}}%
\pgfpathlineto{\pgfqpoint{3.629382in}{1.774455in}}%
\pgfpathlineto{\pgfqpoint{3.600802in}{1.833231in}}%
\pgfpathlineto{\pgfqpoint{3.647849in}{1.849658in}}%
\pgfpathlineto{\pgfqpoint{3.694815in}{1.866041in}}%
\pgfpathlineto{\pgfqpoint{3.723425in}{1.807466in}}%
\pgfpathlineto{\pgfqpoint{3.752149in}{1.751338in}}%
\pgfpathlineto{\pgfqpoint{3.705151in}{1.734785in}}%
\pgfpathlineto{\pgfqpoint{3.658072in}{1.718185in}}%
\pgfpathclose%
\pgfusepath{stroke,fill}%
\end{pgfscope}%
\begin{pgfscope}%
\pgfpathrectangle{\pgfqpoint{0.050000in}{0.083252in}}{\pgfqpoint{4.900000in}{4.900000in}}%
\pgfusepath{clip}%
\pgfsetbuttcap%
\pgfsetroundjoin%
\definecolor{currentfill}{rgb}{0.254348,0.368627,0.758847}%
\pgfsetfillcolor{currentfill}%
\pgfsetfillopacity{0.950000}%
\pgfsetlinewidth{0.200750pt}%
\definecolor{currentstroke}{rgb}{0.200000,0.200000,0.200000}%
\pgfsetstrokecolor{currentstroke}%
\pgfsetdash{}{0pt}%
\pgfpathmoveto{\pgfqpoint{2.167691in}{3.314880in}}%
\pgfpathlineto{\pgfqpoint{2.140151in}{3.388981in}}%
\pgfpathlineto{\pgfqpoint{2.112637in}{3.463013in}}%
\pgfpathlineto{\pgfqpoint{2.161001in}{3.476819in}}%
\pgfpathlineto{\pgfqpoint{2.209279in}{3.490600in}}%
\pgfpathlineto{\pgfqpoint{2.236789in}{3.416713in}}%
\pgfpathlineto{\pgfqpoint{2.264326in}{3.342757in}}%
\pgfpathlineto{\pgfqpoint{2.216051in}{3.328831in}}%
\pgfpathlineto{\pgfqpoint{2.167691in}{3.314880in}}%
\pgfpathclose%
\pgfusepath{stroke,fill}%
\end{pgfscope}%
\begin{pgfscope}%
\pgfpathrectangle{\pgfqpoint{0.050000in}{0.083252in}}{\pgfqpoint{4.900000in}{4.900000in}}%
\pgfusepath{clip}%
\pgfsetbuttcap%
\pgfsetroundjoin%
\definecolor{currentfill}{rgb}{0.880661,0.875740,0.928643}%
\pgfsetfillcolor{currentfill}%
\pgfsetfillopacity{0.950000}%
\pgfsetlinewidth{0.200750pt}%
\definecolor{currentstroke}{rgb}{0.200000,0.200000,0.200000}%
\pgfsetstrokecolor{currentstroke}%
\pgfsetdash{}{0pt}%
\pgfpathmoveto{\pgfqpoint{3.809976in}{1.646953in}}%
\pgfpathlineto{\pgfqpoint{3.780996in}{1.697800in}}%
\pgfpathlineto{\pgfqpoint{3.752149in}{1.751338in}}%
\pgfpathlineto{\pgfqpoint{3.799065in}{1.767844in}}%
\pgfpathlineto{\pgfqpoint{3.845900in}{1.784305in}}%
\pgfpathlineto{\pgfqpoint{3.874784in}{1.730843in}}%
\pgfpathlineto{\pgfqpoint{3.903805in}{1.680005in}}%
\pgfpathlineto{\pgfqpoint{3.856931in}{1.663501in}}%
\pgfpathlineto{\pgfqpoint{3.809976in}{1.646953in}}%
\pgfpathclose%
\pgfusepath{stroke,fill}%
\end{pgfscope}%
\begin{pgfscope}%
\pgfpathrectangle{\pgfqpoint{0.050000in}{0.083252in}}{\pgfqpoint{4.900000in}{4.900000in}}%
\pgfusepath{clip}%
\pgfsetbuttcap%
\pgfsetroundjoin%
\definecolor{currentfill}{rgb}{0.785190,0.776332,0.871557}%
\pgfsetfillcolor{currentfill}%
\pgfsetfillopacity{0.950000}%
\pgfsetlinewidth{0.200750pt}%
\definecolor{currentstroke}{rgb}{0.200000,0.200000,0.200000}%
\pgfsetstrokecolor{currentstroke}%
\pgfsetdash{}{0pt}%
\pgfpathmoveto{\pgfqpoint{3.506462in}{1.800245in}}%
\pgfpathlineto{\pgfqpoint{3.478010in}{1.861599in}}%
\pgfpathlineto{\pgfqpoint{3.449645in}{1.925093in}}%
\pgfpathlineto{\pgfqpoint{3.496831in}{1.941322in}}%
\pgfpathlineto{\pgfqpoint{3.543936in}{1.957514in}}%
\pgfpathlineto{\pgfqpoint{3.572323in}{1.894323in}}%
\pgfpathlineto{\pgfqpoint{3.600802in}{1.833231in}}%
\pgfpathlineto{\pgfqpoint{3.553673in}{1.816760in}}%
\pgfpathlineto{\pgfqpoint{3.506462in}{1.800245in}}%
\pgfpathclose%
\pgfusepath{stroke,fill}%
\end{pgfscope}%
\begin{pgfscope}%
\pgfpathrectangle{\pgfqpoint{0.050000in}{0.083252in}}{\pgfqpoint{4.900000in}{4.900000in}}%
\pgfusepath{clip}%
\pgfsetbuttcap%
\pgfsetroundjoin%
\definecolor{currentfill}{rgb}{0.247705,0.298039,0.665006}%
\pgfsetfillcolor{currentfill}%
\pgfsetfillopacity{0.950000}%
\pgfsetlinewidth{0.200750pt}%
\definecolor{currentstroke}{rgb}{0.200000,0.200000,0.200000}%
\pgfsetstrokecolor{currentstroke}%
\pgfsetdash{}{0pt}%
\pgfpathmoveto{\pgfqpoint{2.374731in}{3.046244in}}%
\pgfpathlineto{\pgfqpoint{2.347091in}{3.120475in}}%
\pgfpathlineto{\pgfqpoint{2.319477in}{3.194637in}}%
\pgfpathlineto{\pgfqpoint{2.367662in}{3.208685in}}%
\pgfpathlineto{\pgfqpoint{2.415761in}{3.222707in}}%
\pgfpathlineto{\pgfqpoint{2.443372in}{3.148691in}}%
\pgfpathlineto{\pgfqpoint{2.471008in}{3.074609in}}%
\pgfpathlineto{\pgfqpoint{2.422913in}{3.060439in}}%
\pgfpathlineto{\pgfqpoint{2.374731in}{3.046244in}}%
\pgfpathclose%
\pgfusepath{stroke,fill}%
\end{pgfscope}%
\begin{pgfscope}%
\pgfpathrectangle{\pgfqpoint{0.050000in}{0.083252in}}{\pgfqpoint{4.900000in}{4.900000in}}%
\pgfusepath{clip}%
\pgfsetbuttcap%
\pgfsetroundjoin%
\definecolor{currentfill}{rgb}{0.241307,0.230065,0.574641}%
\pgfsetfillcolor{currentfill}%
\pgfsetfillopacity{0.950000}%
\pgfsetlinewidth{0.200750pt}%
\definecolor{currentstroke}{rgb}{0.200000,0.200000,0.200000}%
\pgfsetstrokecolor{currentstroke}%
\pgfsetdash{}{0pt}%
\pgfpathmoveto{\pgfqpoint{2.581815in}{2.777676in}}%
\pgfpathlineto{\pgfqpoint{2.554074in}{2.851990in}}%
\pgfpathlineto{\pgfqpoint{2.526359in}{2.926254in}}%
\pgfpathlineto{\pgfqpoint{2.574365in}{2.940552in}}%
\pgfpathlineto{\pgfqpoint{2.622286in}{2.954826in}}%
\pgfpathlineto{\pgfqpoint{2.649997in}{2.880720in}}%
\pgfpathlineto{\pgfqpoint{2.677733in}{2.806571in}}%
\pgfpathlineto{\pgfqpoint{2.629817in}{2.792135in}}%
\pgfpathlineto{\pgfqpoint{2.581815in}{2.777676in}}%
\pgfpathclose%
\pgfusepath{stroke,fill}%
\end{pgfscope}%
\begin{pgfscope}%
\pgfpathrectangle{\pgfqpoint{0.050000in}{0.083252in}}{\pgfqpoint{4.900000in}{4.900000in}}%
\pgfusepath{clip}%
\pgfsetbuttcap%
\pgfsetroundjoin%
\definecolor{currentfill}{rgb}{0.349604,0.322784,0.611103}%
\pgfsetfillcolor{currentfill}%
\pgfsetfillopacity{0.950000}%
\pgfsetlinewidth{0.200750pt}%
\definecolor{currentstroke}{rgb}{0.200000,0.200000,0.200000}%
\pgfsetstrokecolor{currentstroke}%
\pgfsetdash{}{0pt}%
\pgfpathmoveto{\pgfqpoint{2.788946in}{2.509899in}}%
\pgfpathlineto{\pgfqpoint{2.761103in}{2.584023in}}%
\pgfpathlineto{\pgfqpoint{2.733286in}{2.658200in}}%
\pgfpathlineto{\pgfqpoint{2.781114in}{2.672793in}}%
\pgfpathlineto{\pgfqpoint{2.828857in}{2.687366in}}%
\pgfpathlineto{\pgfqpoint{2.856670in}{2.613391in}}%
\pgfpathlineto{\pgfqpoint{2.884510in}{2.539484in}}%
\pgfpathlineto{\pgfqpoint{2.836771in}{2.524701in}}%
\pgfpathlineto{\pgfqpoint{2.788946in}{2.509899in}}%
\pgfpathclose%
\pgfusepath{stroke,fill}%
\end{pgfscope}%
\begin{pgfscope}%
\pgfpathrectangle{\pgfqpoint{0.050000in}{0.083252in}}{\pgfqpoint{4.900000in}{4.900000in}}%
\pgfusepath{clip}%
\pgfsetbuttcap%
\pgfsetroundjoin%
\definecolor{currentfill}{rgb}{0.916463,0.913018,0.950050}%
\pgfsetfillcolor{currentfill}%
\pgfsetfillopacity{0.950000}%
\pgfsetlinewidth{0.200750pt}%
\definecolor{currentstroke}{rgb}{0.200000,0.200000,0.200000}%
\pgfsetstrokecolor{currentstroke}%
\pgfsetdash{}{0pt}%
\pgfpathmoveto{\pgfqpoint{3.962278in}{1.586314in}}%
\pgfpathlineto{\pgfqpoint{3.932967in}{1.631834in}}%
\pgfpathlineto{\pgfqpoint{3.903805in}{1.680005in}}%
\pgfpathlineto{\pgfqpoint{3.950596in}{1.696464in}}%
\pgfpathlineto{\pgfqpoint{3.997307in}{1.712879in}}%
\pgfpathlineto{\pgfqpoint{4.026513in}{1.664657in}}%
\pgfpathlineto{\pgfqpoint{4.055871in}{1.619027in}}%
\pgfpathlineto{\pgfqpoint{4.009115in}{1.602691in}}%
\pgfpathlineto{\pgfqpoint{3.962278in}{1.586314in}}%
\pgfpathclose%
\pgfusepath{stroke,fill}%
\end{pgfscope}%
\begin{pgfscope}%
\pgfpathrectangle{\pgfqpoint{0.050000in}{0.083252in}}{\pgfqpoint{4.900000in}{4.900000in}}%
\pgfusepath{clip}%
\pgfsetbuttcap%
\pgfsetroundjoin%
\definecolor{currentfill}{rgb}{0.719554,0.707989,0.832311}%
\pgfsetfillcolor{currentfill}%
\pgfsetfillopacity{0.950000}%
\pgfsetlinewidth{0.200750pt}%
\definecolor{currentstroke}{rgb}{0.200000,0.200000,0.200000}%
\pgfsetstrokecolor{currentstroke}%
\pgfsetdash{}{0pt}%
\pgfpathmoveto{\pgfqpoint{3.355026in}{1.892519in}}%
\pgfpathlineto{\pgfqpoint{3.326757in}{1.958255in}}%
\pgfpathlineto{\pgfqpoint{3.298555in}{2.025654in}}%
\pgfpathlineto{\pgfqpoint{3.345889in}{2.041602in}}%
\pgfpathlineto{\pgfqpoint{3.393141in}{2.057519in}}%
\pgfpathlineto{\pgfqpoint{3.421358in}{1.990482in}}%
\pgfpathlineto{\pgfqpoint{3.449645in}{1.925093in}}%
\pgfpathlineto{\pgfqpoint{3.402376in}{1.908825in}}%
\pgfpathlineto{\pgfqpoint{3.355026in}{1.892519in}}%
\pgfpathclose%
\pgfusepath{stroke,fill}%
\end{pgfscope}%
\begin{pgfscope}%
\pgfpathrectangle{\pgfqpoint{0.050000in}{0.083252in}}{\pgfqpoint{4.900000in}{4.900000in}}%
\pgfusepath{clip}%
\pgfsetbuttcap%
\pgfsetroundjoin%
\definecolor{currentfill}{rgb}{0.504744,0.484321,0.703868}%
\pgfsetfillcolor{currentfill}%
\pgfsetfillopacity{0.950000}%
\pgfsetlinewidth{0.200750pt}%
\definecolor{currentstroke}{rgb}{0.200000,0.200000,0.200000}%
\pgfsetstrokecolor{currentstroke}%
\pgfsetdash{}{0pt}%
\pgfpathmoveto{\pgfqpoint{2.996166in}{2.245894in}}%
\pgfpathlineto{\pgfqpoint{2.968205in}{2.318802in}}%
\pgfpathlineto{\pgfqpoint{2.940277in}{2.392108in}}%
\pgfpathlineto{\pgfqpoint{2.987929in}{2.407125in}}%
\pgfpathlineto{\pgfqpoint{3.035498in}{2.422123in}}%
\pgfpathlineto{\pgfqpoint{3.063426in}{2.349105in}}%
\pgfpathlineto{\pgfqpoint{3.091388in}{2.276514in}}%
\pgfpathlineto{\pgfqpoint{3.043819in}{2.261213in}}%
\pgfpathlineto{\pgfqpoint{2.996166in}{2.245894in}}%
\pgfpathclose%
\pgfusepath{stroke,fill}%
\end{pgfscope}%
\begin{pgfscope}%
\pgfpathrectangle{\pgfqpoint{0.050000in}{0.083252in}}{\pgfqpoint{4.900000in}{4.900000in}}%
\pgfusepath{clip}%
\pgfsetbuttcap%
\pgfsetroundjoin%
\definecolor{currentfill}{rgb}{0.264191,0.473203,0.897870}%
\pgfsetfillcolor{currentfill}%
\pgfsetfillopacity{0.950000}%
\pgfsetlinewidth{0.200750pt}%
\definecolor{currentstroke}{rgb}{0.200000,0.200000,0.200000}%
\pgfsetstrokecolor{currentstroke}%
\pgfsetdash{}{0pt}%
\pgfpathmoveto{\pgfqpoint{1.808494in}{3.704130in}}%
\pgfpathlineto{\pgfqpoint{1.781103in}{3.778103in}}%
\pgfpathlineto{\pgfqpoint{1.753738in}{3.852006in}}%
\pgfpathlineto{\pgfqpoint{1.802457in}{3.865476in}}%
\pgfpathlineto{\pgfqpoint{1.851089in}{3.878923in}}%
\pgfpathlineto{\pgfqpoint{1.878451in}{3.805164in}}%
\pgfpathlineto{\pgfqpoint{1.905839in}{3.731336in}}%
\pgfpathlineto{\pgfqpoint{1.857210in}{3.717745in}}%
\pgfpathlineto{\pgfqpoint{1.808494in}{3.704130in}}%
\pgfpathclose%
\pgfusepath{stroke,fill}%
\end{pgfscope}%
\begin{pgfscope}%
\pgfpathrectangle{\pgfqpoint{0.050000in}{0.083252in}}{\pgfqpoint{4.900000in}{4.900000in}}%
\pgfusepath{clip}%
\pgfsetbuttcap%
\pgfsetroundjoin%
\definecolor{currentfill}{rgb}{0.257547,0.402614,0.804029}%
\pgfsetfillcolor{currentfill}%
\pgfsetfillopacity{0.950000}%
\pgfsetlinewidth{0.200750pt}%
\definecolor{currentstroke}{rgb}{0.200000,0.200000,0.200000}%
\pgfsetstrokecolor{currentstroke}%
\pgfsetdash{}{0pt}%
\pgfpathmoveto{\pgfqpoint{2.015650in}{3.435326in}}%
\pgfpathlineto{\pgfqpoint{1.988158in}{3.509433in}}%
\pgfpathlineto{\pgfqpoint{1.960693in}{3.583470in}}%
\pgfpathlineto{\pgfqpoint{2.009233in}{3.597181in}}%
\pgfpathlineto{\pgfqpoint{2.057686in}{3.610867in}}%
\pgfpathlineto{\pgfqpoint{2.085149in}{3.536975in}}%
\pgfpathlineto{\pgfqpoint{2.112637in}{3.463013in}}%
\pgfpathlineto{\pgfqpoint{2.064186in}{3.449182in}}%
\pgfpathlineto{\pgfqpoint{2.015650in}{3.435326in}}%
\pgfpathclose%
\pgfusepath{stroke,fill}%
\end{pgfscope}%
\begin{pgfscope}%
\pgfpathrectangle{\pgfqpoint{0.050000in}{0.083252in}}{\pgfqpoint{4.900000in}{4.900000in}}%
\pgfusepath{clip}%
\pgfsetbuttcap%
\pgfsetroundjoin%
\definecolor{currentfill}{rgb}{0.270096,0.535948,0.981284}%
\pgfsetfillcolor{currentfill}%
\pgfsetfillopacity{0.950000}%
\pgfsetlinewidth{0.200750pt}%
\definecolor{currentstroke}{rgb}{0.200000,0.200000,0.200000}%
\pgfsetstrokecolor{currentstroke}%
\pgfsetdash{}{0pt}%
\pgfpathmoveto{\pgfqpoint{1.601380in}{3.972878in}}%
\pgfpathlineto{\pgfqpoint{1.574090in}{4.046716in}}%
\pgfpathlineto{\pgfqpoint{1.622987in}{4.060019in}}%
\pgfpathlineto{\pgfqpoint{1.671797in}{4.073297in}}%
\pgfpathlineto{\pgfqpoint{1.699085in}{3.999603in}}%
\pgfpathlineto{\pgfqpoint{1.650276in}{3.986252in}}%
\pgfpathlineto{\pgfqpoint{1.601380in}{3.972878in}}%
\pgfpathclose%
\pgfusepath{stroke,fill}%
\end{pgfscope}%
\begin{pgfscope}%
\pgfpathrectangle{\pgfqpoint{0.050000in}{0.083252in}}{\pgfqpoint{4.900000in}{4.900000in}}%
\pgfusepath{clip}%
\pgfsetbuttcap%
\pgfsetroundjoin%
\definecolor{currentfill}{rgb}{0.653918,0.639646,0.793064}%
\pgfsetfillcolor{currentfill}%
\pgfsetfillopacity{0.950000}%
\pgfsetlinewidth{0.200750pt}%
\definecolor{currentstroke}{rgb}{0.200000,0.200000,0.200000}%
\pgfsetstrokecolor{currentstroke}%
\pgfsetdash{}{0pt}%
\pgfpathmoveto{\pgfqpoint{3.203639in}{1.993666in}}%
\pgfpathlineto{\pgfqpoint{3.175507in}{2.062854in}}%
\pgfpathlineto{\pgfqpoint{3.147425in}{2.133204in}}%
\pgfpathlineto{\pgfqpoint{3.194915in}{2.148833in}}%
\pgfpathlineto{\pgfqpoint{3.242322in}{2.164439in}}%
\pgfpathlineto{\pgfqpoint{3.270412in}{2.094461in}}%
\pgfpathlineto{\pgfqpoint{3.298555in}{2.025654in}}%
\pgfpathlineto{\pgfqpoint{3.251139in}{2.009675in}}%
\pgfpathlineto{\pgfqpoint{3.203639in}{1.993666in}}%
\pgfpathclose%
\pgfusepath{stroke,fill}%
\end{pgfscope}%
\begin{pgfscope}%
\pgfpathrectangle{\pgfqpoint{0.050000in}{0.083252in}}{\pgfqpoint{4.900000in}{4.900000in}}%
\pgfusepath{clip}%
\pgfsetbuttcap%
\pgfsetroundjoin%
\definecolor{currentfill}{rgb}{0.946298,0.944083,0.967889}%
\pgfsetfillcolor{currentfill}%
\pgfsetfillopacity{0.950000}%
\pgfsetlinewidth{0.200750pt}%
\definecolor{currentstroke}{rgb}{0.200000,0.200000,0.200000}%
\pgfsetstrokecolor{currentstroke}%
\pgfsetdash{}{0pt}%
\pgfpathmoveto{\pgfqpoint{4.115051in}{1.535191in}}%
\pgfpathlineto{\pgfqpoint{4.085383in}{1.575916in}}%
\pgfpathlineto{\pgfqpoint{4.055871in}{1.619027in}}%
\pgfpathlineto{\pgfqpoint{4.102545in}{1.635322in}}%
\pgfpathlineto{\pgfqpoint{4.149137in}{1.651578in}}%
\pgfpathlineto{\pgfqpoint{4.178698in}{1.608310in}}%
\pgfpathlineto{\pgfqpoint{4.208418in}{1.567380in}}%
\pgfpathlineto{\pgfqpoint{4.161776in}{1.551304in}}%
\pgfpathlineto{\pgfqpoint{4.115051in}{1.535191in}}%
\pgfpathclose%
\pgfusepath{stroke,fill}%
\end{pgfscope}%
\begin{pgfscope}%
\pgfpathrectangle{\pgfqpoint{0.050000in}{0.083252in}}{\pgfqpoint{4.900000in}{4.900000in}}%
\pgfusepath{clip}%
\pgfsetbuttcap%
\pgfsetroundjoin%
\definecolor{currentfill}{rgb}{0.251150,0.334641,0.713664}%
\pgfsetfillcolor{currentfill}%
\pgfsetfillopacity{0.950000}%
\pgfsetlinewidth{0.200750pt}%
\definecolor{currentstroke}{rgb}{0.200000,0.200000,0.200000}%
\pgfsetstrokecolor{currentstroke}%
\pgfsetdash{}{0pt}%
\pgfpathmoveto{\pgfqpoint{2.222849in}{3.166468in}}%
\pgfpathlineto{\pgfqpoint{2.195257in}{3.240709in}}%
\pgfpathlineto{\pgfqpoint{2.167691in}{3.314880in}}%
\pgfpathlineto{\pgfqpoint{2.216051in}{3.328831in}}%
\pgfpathlineto{\pgfqpoint{2.264326in}{3.342757in}}%
\pgfpathlineto{\pgfqpoint{2.291888in}{3.268732in}}%
\pgfpathlineto{\pgfqpoint{2.319477in}{3.194637in}}%
\pgfpathlineto{\pgfqpoint{2.271206in}{3.180565in}}%
\pgfpathlineto{\pgfqpoint{2.222849in}{3.166468in}}%
\pgfpathclose%
\pgfusepath{stroke,fill}%
\end{pgfscope}%
\begin{pgfscope}%
\pgfpathrectangle{\pgfqpoint{0.050000in}{0.083252in}}{\pgfqpoint{4.900000in}{4.900000in}}%
\pgfusepath{clip}%
\pgfsetbuttcap%
\pgfsetroundjoin%
\definecolor{currentfill}{rgb}{0.244506,0.264052,0.619823}%
\pgfsetfillcolor{currentfill}%
\pgfsetfillopacity{0.950000}%
\pgfsetlinewidth{0.200750pt}%
\definecolor{currentstroke}{rgb}{0.200000,0.200000,0.200000}%
\pgfsetstrokecolor{currentstroke}%
\pgfsetdash{}{0pt}%
\pgfpathmoveto{\pgfqpoint{2.430090in}{2.897586in}}%
\pgfpathlineto{\pgfqpoint{2.402398in}{2.971947in}}%
\pgfpathlineto{\pgfqpoint{2.374731in}{3.046244in}}%
\pgfpathlineto{\pgfqpoint{2.422913in}{3.060439in}}%
\pgfpathlineto{\pgfqpoint{2.471008in}{3.074609in}}%
\pgfpathlineto{\pgfqpoint{2.498671in}{3.000462in}}%
\pgfpathlineto{\pgfqpoint{2.526359in}{2.926254in}}%
\pgfpathlineto{\pgfqpoint{2.478268in}{2.911932in}}%
\pgfpathlineto{\pgfqpoint{2.430090in}{2.897586in}}%
\pgfpathclose%
\pgfusepath{stroke,fill}%
\end{pgfscope}%
\begin{pgfscope}%
\pgfpathrectangle{\pgfqpoint{0.050000in}{0.083252in}}{\pgfqpoint{4.900000in}{4.900000in}}%
\pgfusepath{clip}%
\pgfsetbuttcap%
\pgfsetroundjoin%
\definecolor{currentfill}{rgb}{0.272034,0.242015,0.564721}%
\pgfsetfillcolor{currentfill}%
\pgfsetfillopacity{0.950000}%
\pgfsetlinewidth{0.200750pt}%
\definecolor{currentstroke}{rgb}{0.200000,0.200000,0.200000}%
\pgfsetstrokecolor{currentstroke}%
\pgfsetdash{}{0pt}%
\pgfpathmoveto{\pgfqpoint{2.637376in}{2.628949in}}%
\pgfpathlineto{\pgfqpoint{2.609582in}{2.703324in}}%
\pgfpathlineto{\pgfqpoint{2.581815in}{2.777676in}}%
\pgfpathlineto{\pgfqpoint{2.629817in}{2.792135in}}%
\pgfpathlineto{\pgfqpoint{2.677733in}{2.806571in}}%
\pgfpathlineto{\pgfqpoint{2.705497in}{2.732391in}}%
\pgfpathlineto{\pgfqpoint{2.733286in}{2.658200in}}%
\pgfpathlineto{\pgfqpoint{2.685374in}{2.643585in}}%
\pgfpathlineto{\pgfqpoint{2.637376in}{2.628949in}}%
\pgfpathclose%
\pgfusepath{stroke,fill}%
\end{pgfscope}%
\begin{pgfscope}%
\pgfpathrectangle{\pgfqpoint{0.050000in}{0.083252in}}{\pgfqpoint{4.900000in}{4.900000in}}%
\pgfusepath{clip}%
\pgfsetbuttcap%
\pgfsetroundjoin%
\definecolor{currentfill}{rgb}{0.427174,0.403552,0.657486}%
\pgfsetfillcolor{currentfill}%
\pgfsetfillopacity{0.950000}%
\pgfsetlinewidth{0.200750pt}%
\definecolor{currentstroke}{rgb}{0.200000,0.200000,0.200000}%
\pgfsetstrokecolor{currentstroke}%
\pgfsetdash{}{0pt}%
\pgfpathmoveto{\pgfqpoint{2.844717in}{2.362022in}}%
\pgfpathlineto{\pgfqpoint{2.816818in}{2.435876in}}%
\pgfpathlineto{\pgfqpoint{2.788946in}{2.509899in}}%
\pgfpathlineto{\pgfqpoint{2.836771in}{2.524701in}}%
\pgfpathlineto{\pgfqpoint{2.884510in}{2.539484in}}%
\pgfpathlineto{\pgfqpoint{2.912379in}{2.465699in}}%
\pgfpathlineto{\pgfqpoint{2.940277in}{2.392108in}}%
\pgfpathlineto{\pgfqpoint{2.892539in}{2.377074in}}%
\pgfpathlineto{\pgfqpoint{2.844717in}{2.362022in}}%
\pgfpathclose%
\pgfusepath{stroke,fill}%
\end{pgfscope}%
\begin{pgfscope}%
\pgfpathrectangle{\pgfqpoint{0.050000in}{0.083252in}}{\pgfqpoint{4.900000in}{4.900000in}}%
\pgfusepath{clip}%
\pgfsetbuttcap%
\pgfsetroundjoin%
\definecolor{currentfill}{rgb}{0.582314,0.565090,0.750250}%
\pgfsetfillcolor{currentfill}%
\pgfsetfillopacity{0.950000}%
\pgfsetlinewidth{0.200750pt}%
\definecolor{currentstroke}{rgb}{0.200000,0.200000,0.200000}%
\pgfsetstrokecolor{currentstroke}%
\pgfsetdash{}{0pt}%
\pgfpathmoveto{\pgfqpoint{3.052195in}{2.101877in}}%
\pgfpathlineto{\pgfqpoint{3.024161in}{2.173528in}}%
\pgfpathlineto{\pgfqpoint{2.996166in}{2.245894in}}%
\pgfpathlineto{\pgfqpoint{3.043819in}{2.261213in}}%
\pgfpathlineto{\pgfqpoint{3.091388in}{2.276514in}}%
\pgfpathlineto{\pgfqpoint{3.119387in}{2.204490in}}%
\pgfpathlineto{\pgfqpoint{3.147425in}{2.133204in}}%
\pgfpathlineto{\pgfqpoint{3.099852in}{2.117552in}}%
\pgfpathlineto{\pgfqpoint{3.052195in}{2.101877in}}%
\pgfpathclose%
\pgfusepath{stroke,fill}%
\end{pgfscope}%
\begin{pgfscope}%
\pgfpathrectangle{\pgfqpoint{0.050000in}{0.083252in}}{\pgfqpoint{4.900000in}{4.900000in}}%
\pgfusepath{clip}%
\pgfsetbuttcap%
\pgfsetroundjoin%
\definecolor{currentfill}{rgb}{0.267636,0.509804,0.946528}%
\pgfsetfillcolor{currentfill}%
\pgfsetfillopacity{0.950000}%
\pgfsetlinewidth{0.200750pt}%
\definecolor{currentstroke}{rgb}{0.200000,0.200000,0.200000}%
\pgfsetstrokecolor{currentstroke}%
\pgfsetdash{}{0pt}%
\pgfpathmoveto{\pgfqpoint{1.656039in}{3.824992in}}%
\pgfpathlineto{\pgfqpoint{1.628696in}{3.898970in}}%
\pgfpathlineto{\pgfqpoint{1.601380in}{3.972878in}}%
\pgfpathlineto{\pgfqpoint{1.650276in}{3.986252in}}%
\pgfpathlineto{\pgfqpoint{1.699085in}{3.999603in}}%
\pgfpathlineto{\pgfqpoint{1.726398in}{3.925839in}}%
\pgfpathlineto{\pgfqpoint{1.753738in}{3.852006in}}%
\pgfpathlineto{\pgfqpoint{1.704932in}{3.838511in}}%
\pgfpathlineto{\pgfqpoint{1.656039in}{3.824992in}}%
\pgfpathclose%
\pgfusepath{stroke,fill}%
\end{pgfscope}%
\begin{pgfscope}%
\pgfpathrectangle{\pgfqpoint{0.050000in}{0.083252in}}{\pgfqpoint{4.900000in}{4.900000in}}%
\pgfusepath{clip}%
\pgfsetbuttcap%
\pgfsetroundjoin%
\definecolor{currentfill}{rgb}{0.838893,0.832249,0.903668}%
\pgfsetfillcolor{currentfill}%
\pgfsetfillopacity{0.950000}%
\pgfsetlinewidth{0.200750pt}%
\definecolor{currentstroke}{rgb}{0.200000,0.200000,0.200000}%
\pgfsetstrokecolor{currentstroke}%
\pgfsetdash{}{0pt}%
\pgfpathmoveto{\pgfqpoint{3.563667in}{1.684843in}}%
\pgfpathlineto{\pgfqpoint{3.535012in}{1.741260in}}%
\pgfpathlineto{\pgfqpoint{3.506462in}{1.800245in}}%
\pgfpathlineto{\pgfqpoint{3.553673in}{1.816760in}}%
\pgfpathlineto{\pgfqpoint{3.600802in}{1.833231in}}%
\pgfpathlineto{\pgfqpoint{3.629382in}{1.774455in}}%
\pgfpathlineto{\pgfqpoint{3.658072in}{1.718185in}}%
\pgfpathlineto{\pgfqpoint{3.610911in}{1.701538in}}%
\pgfpathlineto{\pgfqpoint{3.563667in}{1.684843in}}%
\pgfpathclose%
\pgfusepath{stroke,fill}%
\end{pgfscope}%
\begin{pgfscope}%
\pgfpathrectangle{\pgfqpoint{0.050000in}{0.083252in}}{\pgfqpoint{4.900000in}{4.900000in}}%
\pgfusepath{clip}%
\pgfsetbuttcap%
\pgfsetroundjoin%
\definecolor{currentfill}{rgb}{0.260992,0.439216,0.852687}%
\pgfsetfillcolor{currentfill}%
\pgfsetfillopacity{0.950000}%
\pgfsetlinewidth{0.200750pt}%
\definecolor{currentstroke}{rgb}{0.200000,0.200000,0.200000}%
\pgfsetstrokecolor{currentstroke}%
\pgfsetdash{}{0pt}%
\pgfpathmoveto{\pgfqpoint{1.863353in}{3.555975in}}%
\pgfpathlineto{\pgfqpoint{1.835910in}{3.630088in}}%
\pgfpathlineto{\pgfqpoint{1.808494in}{3.704130in}}%
\pgfpathlineto{\pgfqpoint{1.857210in}{3.717745in}}%
\pgfpathlineto{\pgfqpoint{1.905839in}{3.731336in}}%
\pgfpathlineto{\pgfqpoint{1.933253in}{3.657438in}}%
\pgfpathlineto{\pgfqpoint{1.960693in}{3.583470in}}%
\pgfpathlineto{\pgfqpoint{1.912066in}{3.569735in}}%
\pgfpathlineto{\pgfqpoint{1.863353in}{3.555975in}}%
\pgfpathclose%
\pgfusepath{stroke,fill}%
\end{pgfscope}%
\begin{pgfscope}%
\pgfpathrectangle{\pgfqpoint{0.050000in}{0.083252in}}{\pgfqpoint{4.900000in}{4.900000in}}%
\pgfusepath{clip}%
\pgfsetbuttcap%
\pgfsetroundjoin%
\definecolor{currentfill}{rgb}{0.880661,0.875740,0.928643}%
\pgfsetfillcolor{currentfill}%
\pgfsetfillopacity{0.950000}%
\pgfsetlinewidth{0.200750pt}%
\definecolor{currentstroke}{rgb}{0.200000,0.200000,0.200000}%
\pgfsetstrokecolor{currentstroke}%
\pgfsetdash{}{0pt}%
\pgfpathmoveto{\pgfqpoint{3.715819in}{1.613718in}}%
\pgfpathlineto{\pgfqpoint{3.686882in}{1.664572in}}%
\pgfpathlineto{\pgfqpoint{3.658072in}{1.718185in}}%
\pgfpathlineto{\pgfqpoint{3.705151in}{1.734785in}}%
\pgfpathlineto{\pgfqpoint{3.752149in}{1.751338in}}%
\pgfpathlineto{\pgfqpoint{3.780996in}{1.697800in}}%
\pgfpathlineto{\pgfqpoint{3.809976in}{1.646953in}}%
\pgfpathlineto{\pgfqpoint{3.762938in}{1.630359in}}%
\pgfpathlineto{\pgfqpoint{3.715819in}{1.613718in}}%
\pgfpathclose%
\pgfusepath{stroke,fill}%
\end{pgfscope}%
\begin{pgfscope}%
\pgfpathrectangle{\pgfqpoint{0.050000in}{0.083252in}}{\pgfqpoint{4.900000in}{4.900000in}}%
\pgfusepath{clip}%
\pgfsetbuttcap%
\pgfsetroundjoin%
\definecolor{currentfill}{rgb}{0.254348,0.368627,0.758847}%
\pgfsetfillcolor{currentfill}%
\pgfsetfillopacity{0.950000}%
\pgfsetlinewidth{0.200750pt}%
\definecolor{currentstroke}{rgb}{0.200000,0.200000,0.200000}%
\pgfsetstrokecolor{currentstroke}%
\pgfsetdash{}{0pt}%
\pgfpathmoveto{\pgfqpoint{2.070711in}{3.286902in}}%
\pgfpathlineto{\pgfqpoint{2.043167in}{3.361149in}}%
\pgfpathlineto{\pgfqpoint{2.015650in}{3.435326in}}%
\pgfpathlineto{\pgfqpoint{2.064186in}{3.449182in}}%
\pgfpathlineto{\pgfqpoint{2.112637in}{3.463013in}}%
\pgfpathlineto{\pgfqpoint{2.140151in}{3.388981in}}%
\pgfpathlineto{\pgfqpoint{2.167691in}{3.314880in}}%
\pgfpathlineto{\pgfqpoint{2.119244in}{3.300904in}}%
\pgfpathlineto{\pgfqpoint{2.070711in}{3.286902in}}%
\pgfpathclose%
\pgfusepath{stroke,fill}%
\end{pgfscope}%
\begin{pgfscope}%
\pgfpathrectangle{\pgfqpoint{0.050000in}{0.083252in}}{\pgfqpoint{4.900000in}{4.900000in}}%
\pgfusepath{clip}%
\pgfsetbuttcap%
\pgfsetroundjoin%
\definecolor{currentfill}{rgb}{0.785190,0.776332,0.871557}%
\pgfsetfillcolor{currentfill}%
\pgfsetfillopacity{0.950000}%
\pgfsetlinewidth{0.200750pt}%
\definecolor{currentstroke}{rgb}{0.200000,0.200000,0.200000}%
\pgfsetstrokecolor{currentstroke}%
\pgfsetdash{}{0pt}%
\pgfpathmoveto{\pgfqpoint{3.411794in}{1.767078in}}%
\pgfpathlineto{\pgfqpoint{3.383368in}{1.828707in}}%
\pgfpathlineto{\pgfqpoint{3.355026in}{1.892519in}}%
\pgfpathlineto{\pgfqpoint{3.402376in}{1.908825in}}%
\pgfpathlineto{\pgfqpoint{3.449645in}{1.925093in}}%
\pgfpathlineto{\pgfqpoint{3.478010in}{1.861599in}}%
\pgfpathlineto{\pgfqpoint{3.506462in}{1.800245in}}%
\pgfpathlineto{\pgfqpoint{3.459169in}{1.783685in}}%
\pgfpathlineto{\pgfqpoint{3.411794in}{1.767078in}}%
\pgfpathclose%
\pgfusepath{stroke,fill}%
\end{pgfscope}%
\begin{pgfscope}%
\pgfpathrectangle{\pgfqpoint{0.050000in}{0.083252in}}{\pgfqpoint{4.900000in}{4.900000in}}%
\pgfusepath{clip}%
\pgfsetbuttcap%
\pgfsetroundjoin%
\definecolor{currentfill}{rgb}{0.247705,0.298039,0.665006}%
\pgfsetfillcolor{currentfill}%
\pgfsetfillopacity{0.950000}%
\pgfsetlinewidth{0.200750pt}%
\definecolor{currentstroke}{rgb}{0.200000,0.200000,0.200000}%
\pgfsetstrokecolor{currentstroke}%
\pgfsetdash{}{0pt}%
\pgfpathmoveto{\pgfqpoint{2.278111in}{3.017780in}}%
\pgfpathlineto{\pgfqpoint{2.250467in}{3.092158in}}%
\pgfpathlineto{\pgfqpoint{2.222849in}{3.166468in}}%
\pgfpathlineto{\pgfqpoint{2.271206in}{3.180565in}}%
\pgfpathlineto{\pgfqpoint{2.319477in}{3.194637in}}%
\pgfpathlineto{\pgfqpoint{2.347091in}{3.120475in}}%
\pgfpathlineto{\pgfqpoint{2.374731in}{3.046244in}}%
\pgfpathlineto{\pgfqpoint{2.326464in}{3.032024in}}%
\pgfpathlineto{\pgfqpoint{2.278111in}{3.017780in}}%
\pgfpathclose%
\pgfusepath{stroke,fill}%
\end{pgfscope}%
\begin{pgfscope}%
\pgfpathrectangle{\pgfqpoint{0.050000in}{0.083252in}}{\pgfqpoint{4.900000in}{4.900000in}}%
\pgfusepath{clip}%
\pgfsetbuttcap%
\pgfsetroundjoin%
\definecolor{currentfill}{rgb}{0.241307,0.230065,0.574641}%
\pgfsetfillcolor{currentfill}%
\pgfsetfillopacity{0.950000}%
\pgfsetlinewidth{0.200750pt}%
\definecolor{currentstroke}{rgb}{0.200000,0.200000,0.200000}%
\pgfsetstrokecolor{currentstroke}%
\pgfsetdash{}{0pt}%
\pgfpathmoveto{\pgfqpoint{2.485554in}{2.748689in}}%
\pgfpathlineto{\pgfqpoint{2.457809in}{2.823165in}}%
\pgfpathlineto{\pgfqpoint{2.430090in}{2.897586in}}%
\pgfpathlineto{\pgfqpoint{2.478268in}{2.911932in}}%
\pgfpathlineto{\pgfqpoint{2.526359in}{2.926254in}}%
\pgfpathlineto{\pgfqpoint{2.554074in}{2.851990in}}%
\pgfpathlineto{\pgfqpoint{2.581815in}{2.777676in}}%
\pgfpathlineto{\pgfqpoint{2.533727in}{2.763194in}}%
\pgfpathlineto{\pgfqpoint{2.485554in}{2.748689in}}%
\pgfpathclose%
\pgfusepath{stroke,fill}%
\end{pgfscope}%
\begin{pgfscope}%
\pgfpathrectangle{\pgfqpoint{0.050000in}{0.083252in}}{\pgfqpoint{4.900000in}{4.900000in}}%
\pgfusepath{clip}%
\pgfsetbuttcap%
\pgfsetroundjoin%
\definecolor{currentfill}{rgb}{0.349604,0.322784,0.611103}%
\pgfsetfillcolor{currentfill}%
\pgfsetfillopacity{0.950000}%
\pgfsetlinewidth{0.200750pt}%
\definecolor{currentstroke}{rgb}{0.200000,0.200000,0.200000}%
\pgfsetstrokecolor{currentstroke}%
\pgfsetdash{}{0pt}%
\pgfpathmoveto{\pgfqpoint{2.693043in}{2.480238in}}%
\pgfpathlineto{\pgfqpoint{2.665196in}{2.554576in}}%
\pgfpathlineto{\pgfqpoint{2.637376in}{2.628949in}}%
\pgfpathlineto{\pgfqpoint{2.685374in}{2.643585in}}%
\pgfpathlineto{\pgfqpoint{2.733286in}{2.658200in}}%
\pgfpathlineto{\pgfqpoint{2.761103in}{2.584023in}}%
\pgfpathlineto{\pgfqpoint{2.788946in}{2.509899in}}%
\pgfpathlineto{\pgfqpoint{2.741037in}{2.495078in}}%
\pgfpathlineto{\pgfqpoint{2.693043in}{2.480238in}}%
\pgfpathclose%
\pgfusepath{stroke,fill}%
\end{pgfscope}%
\begin{pgfscope}%
\pgfpathrectangle{\pgfqpoint{0.050000in}{0.083252in}}{\pgfqpoint{4.900000in}{4.900000in}}%
\pgfusepath{clip}%
\pgfsetbuttcap%
\pgfsetroundjoin%
\definecolor{currentfill}{rgb}{0.922430,0.919231,0.953618}%
\pgfsetfillcolor{currentfill}%
\pgfsetfillopacity{0.950000}%
\pgfsetlinewidth{0.200750pt}%
\definecolor{currentstroke}{rgb}{0.200000,0.200000,0.200000}%
\pgfsetstrokecolor{currentstroke}%
\pgfsetdash{}{0pt}%
\pgfpathmoveto{\pgfqpoint{3.868356in}{1.553439in}}%
\pgfpathlineto{\pgfqpoint{3.839094in}{1.598838in}}%
\pgfpathlineto{\pgfqpoint{3.809976in}{1.646953in}}%
\pgfpathlineto{\pgfqpoint{3.856931in}{1.663501in}}%
\pgfpathlineto{\pgfqpoint{3.903805in}{1.680005in}}%
\pgfpathlineto{\pgfqpoint{3.932967in}{1.631834in}}%
\pgfpathlineto{\pgfqpoint{3.962278in}{1.586314in}}%
\pgfpathlineto{\pgfqpoint{3.915358in}{1.569897in}}%
\pgfpathlineto{\pgfqpoint{3.868356in}{1.553439in}}%
\pgfpathclose%
\pgfusepath{stroke,fill}%
\end{pgfscope}%
\begin{pgfscope}%
\pgfpathrectangle{\pgfqpoint{0.050000in}{0.083252in}}{\pgfqpoint{4.900000in}{4.900000in}}%
\pgfusepath{clip}%
\pgfsetbuttcap%
\pgfsetroundjoin%
\definecolor{currentfill}{rgb}{0.964198,0.962722,0.978593}%
\pgfsetfillcolor{currentfill}%
\pgfsetfillopacity{0.950000}%
\pgfsetlinewidth{0.200750pt}%
\definecolor{currentstroke}{rgb}{0.200000,0.200000,0.200000}%
\pgfsetstrokecolor{currentstroke}%
\pgfsetdash{}{0pt}%
\pgfpathmoveto{\pgfqpoint{4.268325in}{1.491721in}}%
\pgfpathlineto{\pgfqpoint{4.238295in}{1.528600in}}%
\pgfpathlineto{\pgfqpoint{4.208418in}{1.567380in}}%
\pgfpathlineto{\pgfqpoint{4.254977in}{1.583421in}}%
\pgfpathlineto{\pgfqpoint{4.284881in}{1.544522in}}%
\pgfpathlineto{\pgfqpoint{4.314939in}{1.507506in}}%
\pgfpathlineto{\pgfqpoint{4.268325in}{1.491721in}}%
\pgfpathclose%
\pgfusepath{stroke,fill}%
\end{pgfscope}%
\begin{pgfscope}%
\pgfpathrectangle{\pgfqpoint{0.050000in}{0.083252in}}{\pgfqpoint{4.900000in}{4.900000in}}%
\pgfusepath{clip}%
\pgfsetbuttcap%
\pgfsetroundjoin%
\definecolor{currentfill}{rgb}{0.725521,0.714202,0.835879}%
\pgfsetfillcolor{currentfill}%
\pgfsetfillopacity{0.950000}%
\pgfsetlinewidth{0.200750pt}%
\definecolor{currentstroke}{rgb}{0.200000,0.200000,0.200000}%
\pgfsetstrokecolor{currentstroke}%
\pgfsetdash{}{0pt}%
\pgfpathmoveto{\pgfqpoint{3.260077in}{1.859790in}}%
\pgfpathlineto{\pgfqpoint{3.231827in}{1.925889in}}%
\pgfpathlineto{\pgfqpoint{3.203639in}{1.993666in}}%
\pgfpathlineto{\pgfqpoint{3.251139in}{2.009675in}}%
\pgfpathlineto{\pgfqpoint{3.298555in}{2.025654in}}%
\pgfpathlineto{\pgfqpoint{3.326757in}{1.958255in}}%
\pgfpathlineto{\pgfqpoint{3.355026in}{1.892519in}}%
\pgfpathlineto{\pgfqpoint{3.307593in}{1.876174in}}%
\pgfpathlineto{\pgfqpoint{3.260077in}{1.859790in}}%
\pgfpathclose%
\pgfusepath{stroke,fill}%
\end{pgfscope}%
\begin{pgfscope}%
\pgfpathrectangle{\pgfqpoint{0.050000in}{0.083252in}}{\pgfqpoint{4.900000in}{4.900000in}}%
\pgfusepath{clip}%
\pgfsetbuttcap%
\pgfsetroundjoin%
\definecolor{currentfill}{rgb}{0.504744,0.484321,0.703868}%
\pgfsetfillcolor{currentfill}%
\pgfsetfillopacity{0.950000}%
\pgfsetlinewidth{0.200750pt}%
\definecolor{currentstroke}{rgb}{0.200000,0.200000,0.200000}%
\pgfsetstrokecolor{currentstroke}%
\pgfsetdash{}{0pt}%
\pgfpathmoveto{\pgfqpoint{2.900607in}{2.215200in}}%
\pgfpathlineto{\pgfqpoint{2.872647in}{2.288426in}}%
\pgfpathlineto{\pgfqpoint{2.844717in}{2.362022in}}%
\pgfpathlineto{\pgfqpoint{2.892539in}{2.377074in}}%
\pgfpathlineto{\pgfqpoint{2.940277in}{2.392108in}}%
\pgfpathlineto{\pgfqpoint{2.968205in}{2.318802in}}%
\pgfpathlineto{\pgfqpoint{2.996166in}{2.245894in}}%
\pgfpathlineto{\pgfqpoint{2.948429in}{2.230556in}}%
\pgfpathlineto{\pgfqpoint{2.900607in}{2.215200in}}%
\pgfpathclose%
\pgfusepath{stroke,fill}%
\end{pgfscope}%
\begin{pgfscope}%
\pgfpathrectangle{\pgfqpoint{0.050000in}{0.083252in}}{\pgfqpoint{4.900000in}{4.900000in}}%
\pgfusepath{clip}%
\pgfsetbuttcap%
\pgfsetroundjoin%
\definecolor{currentfill}{rgb}{0.264191,0.473203,0.897870}%
\pgfsetfillcolor{currentfill}%
\pgfsetfillopacity{0.950000}%
\pgfsetlinewidth{0.200750pt}%
\definecolor{currentstroke}{rgb}{0.200000,0.200000,0.200000}%
\pgfsetstrokecolor{currentstroke}%
\pgfsetdash{}{0pt}%
\pgfpathmoveto{\pgfqpoint{1.710800in}{3.676827in}}%
\pgfpathlineto{\pgfqpoint{1.683406in}{3.750945in}}%
\pgfpathlineto{\pgfqpoint{1.656039in}{3.824992in}}%
\pgfpathlineto{\pgfqpoint{1.704932in}{3.838511in}}%
\pgfpathlineto{\pgfqpoint{1.753738in}{3.852006in}}%
\pgfpathlineto{\pgfqpoint{1.781103in}{3.778103in}}%
\pgfpathlineto{\pgfqpoint{1.808494in}{3.704130in}}%
\pgfpathlineto{\pgfqpoint{1.759690in}{3.690491in}}%
\pgfpathlineto{\pgfqpoint{1.710800in}{3.676827in}}%
\pgfpathclose%
\pgfusepath{stroke,fill}%
\end{pgfscope}%
\begin{pgfscope}%
\pgfpathrectangle{\pgfqpoint{0.050000in}{0.083252in}}{\pgfqpoint{4.900000in}{4.900000in}}%
\pgfusepath{clip}%
\pgfsetbuttcap%
\pgfsetroundjoin%
\definecolor{currentfill}{rgb}{0.257547,0.402614,0.804029}%
\pgfsetfillcolor{currentfill}%
\pgfsetfillopacity{0.950000}%
\pgfsetlinewidth{0.200750pt}%
\definecolor{currentstroke}{rgb}{0.200000,0.200000,0.200000}%
\pgfsetstrokecolor{currentstroke}%
\pgfsetdash{}{0pt}%
\pgfpathmoveto{\pgfqpoint{1.918316in}{3.407540in}}%
\pgfpathlineto{\pgfqpoint{1.890822in}{3.481793in}}%
\pgfpathlineto{\pgfqpoint{1.863353in}{3.555975in}}%
\pgfpathlineto{\pgfqpoint{1.912066in}{3.569735in}}%
\pgfpathlineto{\pgfqpoint{1.960693in}{3.583470in}}%
\pgfpathlineto{\pgfqpoint{1.988158in}{3.509433in}}%
\pgfpathlineto{\pgfqpoint{2.015650in}{3.435326in}}%
\pgfpathlineto{\pgfqpoint{1.967027in}{3.421446in}}%
\pgfpathlineto{\pgfqpoint{1.918316in}{3.407540in}}%
\pgfpathclose%
\pgfusepath{stroke,fill}%
\end{pgfscope}%
\begin{pgfscope}%
\pgfpathrectangle{\pgfqpoint{0.050000in}{0.083252in}}{\pgfqpoint{4.900000in}{4.900000in}}%
\pgfusepath{clip}%
\pgfsetbuttcap%
\pgfsetroundjoin%
\definecolor{currentfill}{rgb}{0.269850,0.533333,0.977809}%
\pgfsetfillcolor{currentfill}%
\pgfsetfillopacity{0.950000}%
\pgfsetlinewidth{0.200750pt}%
\definecolor{currentstroke}{rgb}{0.200000,0.200000,0.200000}%
\pgfsetstrokecolor{currentstroke}%
\pgfsetdash{}{0pt}%
\pgfpathmoveto{\pgfqpoint{1.503327in}{3.946057in}}%
\pgfpathlineto{\pgfqpoint{1.476034in}{4.020040in}}%
\pgfpathlineto{\pgfqpoint{1.525106in}{4.033390in}}%
\pgfpathlineto{\pgfqpoint{1.574090in}{4.046716in}}%
\pgfpathlineto{\pgfqpoint{1.601380in}{3.972878in}}%
\pgfpathlineto{\pgfqpoint{1.552397in}{3.959480in}}%
\pgfpathlineto{\pgfqpoint{1.503327in}{3.946057in}}%
\pgfpathclose%
\pgfusepath{stroke,fill}%
\end{pgfscope}%
\begin{pgfscope}%
\pgfpathrectangle{\pgfqpoint{0.050000in}{0.083252in}}{\pgfqpoint{4.900000in}{4.900000in}}%
\pgfusepath{clip}%
\pgfsetbuttcap%
\pgfsetroundjoin%
\definecolor{currentfill}{rgb}{0.659885,0.645859,0.796632}%
\pgfsetfillcolor{currentfill}%
\pgfsetfillopacity{0.950000}%
\pgfsetlinewidth{0.200750pt}%
\definecolor{currentstroke}{rgb}{0.200000,0.200000,0.200000}%
\pgfsetstrokecolor{currentstroke}%
\pgfsetdash{}{0pt}%
\pgfpathmoveto{\pgfqpoint{3.108391in}{1.961559in}}%
\pgfpathlineto{\pgfqpoint{3.080270in}{2.031142in}}%
\pgfpathlineto{\pgfqpoint{3.052195in}{2.101877in}}%
\pgfpathlineto{\pgfqpoint{3.099852in}{2.117552in}}%
\pgfpathlineto{\pgfqpoint{3.147425in}{2.133204in}}%
\pgfpathlineto{\pgfqpoint{3.175507in}{2.062854in}}%
\pgfpathlineto{\pgfqpoint{3.203639in}{1.993666in}}%
\pgfpathlineto{\pgfqpoint{3.156057in}{1.977627in}}%
\pgfpathlineto{\pgfqpoint{3.108391in}{1.961559in}}%
\pgfpathclose%
\pgfusepath{stroke,fill}%
\end{pgfscope}%
\begin{pgfscope}%
\pgfpathrectangle{\pgfqpoint{0.050000in}{0.083252in}}{\pgfqpoint{4.900000in}{4.900000in}}%
\pgfusepath{clip}%
\pgfsetbuttcap%
\pgfsetroundjoin%
\definecolor{currentfill}{rgb}{0.946298,0.944083,0.967889}%
\pgfsetfillcolor{currentfill}%
\pgfsetfillopacity{0.950000}%
\pgfsetlinewidth{0.200750pt}%
\definecolor{currentstroke}{rgb}{0.200000,0.200000,0.200000}%
\pgfsetstrokecolor{currentstroke}%
\pgfsetdash{}{0pt}%
\pgfpathmoveto{\pgfqpoint{4.021355in}{1.502861in}}%
\pgfpathlineto{\pgfqpoint{3.991740in}{1.543370in}}%
\pgfpathlineto{\pgfqpoint{3.962278in}{1.586314in}}%
\pgfpathlineto{\pgfqpoint{4.009115in}{1.602691in}}%
\pgfpathlineto{\pgfqpoint{4.055871in}{1.619027in}}%
\pgfpathlineto{\pgfqpoint{4.085383in}{1.575916in}}%
\pgfpathlineto{\pgfqpoint{4.115051in}{1.535191in}}%
\pgfpathlineto{\pgfqpoint{4.068245in}{1.519044in}}%
\pgfpathlineto{\pgfqpoint{4.021355in}{1.502861in}}%
\pgfpathclose%
\pgfusepath{stroke,fill}%
\end{pgfscope}%
\begin{pgfscope}%
\pgfpathrectangle{\pgfqpoint{0.050000in}{0.083252in}}{\pgfqpoint{4.900000in}{4.900000in}}%
\pgfusepath{clip}%
\pgfsetbuttcap%
\pgfsetroundjoin%
\definecolor{currentfill}{rgb}{0.251150,0.334641,0.713664}%
\pgfsetfillcolor{currentfill}%
\pgfsetfillopacity{0.950000}%
\pgfsetlinewidth{0.200750pt}%
\definecolor{currentstroke}{rgb}{0.200000,0.200000,0.200000}%
\pgfsetstrokecolor{currentstroke}%
\pgfsetdash{}{0pt}%
\pgfpathmoveto{\pgfqpoint{2.125875in}{3.138199in}}%
\pgfpathlineto{\pgfqpoint{2.098280in}{3.212585in}}%
\pgfpathlineto{\pgfqpoint{2.070711in}{3.286902in}}%
\pgfpathlineto{\pgfqpoint{2.119244in}{3.300904in}}%
\pgfpathlineto{\pgfqpoint{2.167691in}{3.314880in}}%
\pgfpathlineto{\pgfqpoint{2.195257in}{3.240709in}}%
\pgfpathlineto{\pgfqpoint{2.222849in}{3.166468in}}%
\pgfpathlineto{\pgfqpoint{2.174405in}{3.152346in}}%
\pgfpathlineto{\pgfqpoint{2.125875in}{3.138199in}}%
\pgfpathclose%
\pgfusepath{stroke,fill}%
\end{pgfscope}%
\begin{pgfscope}%
\pgfpathrectangle{\pgfqpoint{0.050000in}{0.083252in}}{\pgfqpoint{4.900000in}{4.900000in}}%
\pgfusepath{clip}%
\pgfsetbuttcap%
\pgfsetroundjoin%
\definecolor{currentfill}{rgb}{0.244506,0.264052,0.619823}%
\pgfsetfillcolor{currentfill}%
\pgfsetfillopacity{0.950000}%
\pgfsetlinewidth{0.200750pt}%
\definecolor{currentstroke}{rgb}{0.200000,0.200000,0.200000}%
\pgfsetstrokecolor{currentstroke}%
\pgfsetdash{}{0pt}%
\pgfpathmoveto{\pgfqpoint{2.333477in}{2.868821in}}%
\pgfpathlineto{\pgfqpoint{2.305781in}{2.943333in}}%
\pgfpathlineto{\pgfqpoint{2.278111in}{3.017780in}}%
\pgfpathlineto{\pgfqpoint{2.326464in}{3.032024in}}%
\pgfpathlineto{\pgfqpoint{2.374731in}{3.046244in}}%
\pgfpathlineto{\pgfqpoint{2.402398in}{2.971947in}}%
\pgfpathlineto{\pgfqpoint{2.430090in}{2.897586in}}%
\pgfpathlineto{\pgfqpoint{2.381827in}{2.883216in}}%
\pgfpathlineto{\pgfqpoint{2.333477in}{2.868821in}}%
\pgfpathclose%
\pgfusepath{stroke,fill}%
\end{pgfscope}%
\begin{pgfscope}%
\pgfpathrectangle{\pgfqpoint{0.050000in}{0.083252in}}{\pgfqpoint{4.900000in}{4.900000in}}%
\pgfusepath{clip}%
\pgfsetbuttcap%
\pgfsetroundjoin%
\definecolor{currentfill}{rgb}{0.272034,0.242015,0.564721}%
\pgfsetfillcolor{currentfill}%
\pgfsetfillopacity{0.950000}%
\pgfsetlinewidth{0.200750pt}%
\definecolor{currentstroke}{rgb}{0.200000,0.200000,0.200000}%
\pgfsetstrokecolor{currentstroke}%
\pgfsetdash{}{0pt}%
\pgfpathmoveto{\pgfqpoint{2.541123in}{2.599614in}}%
\pgfpathlineto{\pgfqpoint{2.513325in}{2.674168in}}%
\pgfpathlineto{\pgfqpoint{2.485554in}{2.748689in}}%
\pgfpathlineto{\pgfqpoint{2.533727in}{2.763194in}}%
\pgfpathlineto{\pgfqpoint{2.581815in}{2.777676in}}%
\pgfpathlineto{\pgfqpoint{2.609582in}{2.703324in}}%
\pgfpathlineto{\pgfqpoint{2.637376in}{2.628949in}}%
\pgfpathlineto{\pgfqpoint{2.589292in}{2.614292in}}%
\pgfpathlineto{\pgfqpoint{2.541123in}{2.599614in}}%
\pgfpathclose%
\pgfusepath{stroke,fill}%
\end{pgfscope}%
\begin{pgfscope}%
\pgfpathrectangle{\pgfqpoint{0.050000in}{0.083252in}}{\pgfqpoint{4.900000in}{4.900000in}}%
\pgfusepath{clip}%
\pgfsetbuttcap%
\pgfsetroundjoin%
\definecolor{currentfill}{rgb}{0.433141,0.409765,0.661053}%
\pgfsetfillcolor{currentfill}%
\pgfsetfillopacity{0.950000}%
\pgfsetlinewidth{0.200750pt}%
\definecolor{currentstroke}{rgb}{0.200000,0.200000,0.200000}%
\pgfsetstrokecolor{currentstroke}%
\pgfsetdash{}{0pt}%
\pgfpathmoveto{\pgfqpoint{2.748819in}{2.331867in}}%
\pgfpathlineto{\pgfqpoint{2.720917in}{2.405980in}}%
\pgfpathlineto{\pgfqpoint{2.693043in}{2.480238in}}%
\pgfpathlineto{\pgfqpoint{2.741037in}{2.495078in}}%
\pgfpathlineto{\pgfqpoint{2.788946in}{2.509899in}}%
\pgfpathlineto{\pgfqpoint{2.816818in}{2.435876in}}%
\pgfpathlineto{\pgfqpoint{2.844717in}{2.362022in}}%
\pgfpathlineto{\pgfqpoint{2.796811in}{2.346953in}}%
\pgfpathlineto{\pgfqpoint{2.748819in}{2.331867in}}%
\pgfpathclose%
\pgfusepath{stroke,fill}%
\end{pgfscope}%
\begin{pgfscope}%
\pgfpathrectangle{\pgfqpoint{0.050000in}{0.083252in}}{\pgfqpoint{4.900000in}{4.900000in}}%
\pgfusepath{clip}%
\pgfsetbuttcap%
\pgfsetroundjoin%
\definecolor{currentfill}{rgb}{0.582314,0.565090,0.750250}%
\pgfsetfillcolor{currentfill}%
\pgfsetfillopacity{0.950000}%
\pgfsetlinewidth{0.200750pt}%
\definecolor{currentstroke}{rgb}{0.200000,0.200000,0.200000}%
\pgfsetstrokecolor{currentstroke}%
\pgfsetdash{}{0pt}%
\pgfpathmoveto{\pgfqpoint{2.956629in}{2.070460in}}%
\pgfpathlineto{\pgfqpoint{2.928600in}{2.142487in}}%
\pgfpathlineto{\pgfqpoint{2.900607in}{2.215200in}}%
\pgfpathlineto{\pgfqpoint{2.948429in}{2.230556in}}%
\pgfpathlineto{\pgfqpoint{2.996166in}{2.245894in}}%
\pgfpathlineto{\pgfqpoint{3.024161in}{2.173528in}}%
\pgfpathlineto{\pgfqpoint{3.052195in}{2.101877in}}%
\pgfpathlineto{\pgfqpoint{3.004454in}{2.086179in}}%
\pgfpathlineto{\pgfqpoint{2.956629in}{2.070460in}}%
\pgfpathclose%
\pgfusepath{stroke,fill}%
\end{pgfscope}%
\begin{pgfscope}%
\pgfpathrectangle{\pgfqpoint{0.050000in}{0.083252in}}{\pgfqpoint{4.900000in}{4.900000in}}%
\pgfusepath{clip}%
\pgfsetbuttcap%
\pgfsetroundjoin%
\definecolor{currentfill}{rgb}{0.267389,0.507190,0.943053}%
\pgfsetfillcolor{currentfill}%
\pgfsetfillopacity{0.950000}%
\pgfsetlinewidth{0.200750pt}%
\definecolor{currentstroke}{rgb}{0.200000,0.200000,0.200000}%
\pgfsetstrokecolor{currentstroke}%
\pgfsetdash{}{0pt}%
\pgfpathmoveto{\pgfqpoint{1.557990in}{3.797882in}}%
\pgfpathlineto{\pgfqpoint{1.530646in}{3.872004in}}%
\pgfpathlineto{\pgfqpoint{1.503327in}{3.946057in}}%
\pgfpathlineto{\pgfqpoint{1.552397in}{3.959480in}}%
\pgfpathlineto{\pgfqpoint{1.601380in}{3.972878in}}%
\pgfpathlineto{\pgfqpoint{1.628696in}{3.898970in}}%
\pgfpathlineto{\pgfqpoint{1.656039in}{3.824992in}}%
\pgfpathlineto{\pgfqpoint{1.607058in}{3.811449in}}%
\pgfpathlineto{\pgfqpoint{1.557990in}{3.797882in}}%
\pgfpathclose%
\pgfusepath{stroke,fill}%
\end{pgfscope}%
\begin{pgfscope}%
\pgfpathrectangle{\pgfqpoint{0.050000in}{0.083252in}}{\pgfqpoint{4.900000in}{4.900000in}}%
\pgfusepath{clip}%
\pgfsetbuttcap%
\pgfsetroundjoin%
\definecolor{currentfill}{rgb}{0.844860,0.838462,0.907236}%
\pgfsetfillcolor{currentfill}%
\pgfsetfillopacity{0.950000}%
\pgfsetlinewidth{0.200750pt}%
\definecolor{currentstroke}{rgb}{0.200000,0.200000,0.200000}%
\pgfsetstrokecolor{currentstroke}%
\pgfsetdash{}{0pt}%
\pgfpathmoveto{\pgfqpoint{3.468934in}{1.651304in}}%
\pgfpathlineto{\pgfqpoint{3.440313in}{1.707874in}}%
\pgfpathlineto{\pgfqpoint{3.411794in}{1.767078in}}%
\pgfpathlineto{\pgfqpoint{3.459169in}{1.783685in}}%
\pgfpathlineto{\pgfqpoint{3.506462in}{1.800245in}}%
\pgfpathlineto{\pgfqpoint{3.535012in}{1.741260in}}%
\pgfpathlineto{\pgfqpoint{3.563667in}{1.684843in}}%
\pgfpathlineto{\pgfqpoint{3.516342in}{1.668098in}}%
\pgfpathlineto{\pgfqpoint{3.468934in}{1.651304in}}%
\pgfpathclose%
\pgfusepath{stroke,fill}%
\end{pgfscope}%
\begin{pgfscope}%
\pgfpathrectangle{\pgfqpoint{0.050000in}{0.083252in}}{\pgfqpoint{4.900000in}{4.900000in}}%
\pgfusepath{clip}%
\pgfsetbuttcap%
\pgfsetroundjoin%
\definecolor{currentfill}{rgb}{0.970165,0.968935,0.982161}%
\pgfsetfillcolor{currentfill}%
\pgfsetfillopacity{0.950000}%
\pgfsetlinewidth{0.200750pt}%
\definecolor{currentstroke}{rgb}{0.200000,0.200000,0.200000}%
\pgfsetstrokecolor{currentstroke}%
\pgfsetdash{}{0pt}%
\pgfpathmoveto{\pgfqpoint{4.174847in}{1.460057in}}%
\pgfpathlineto{\pgfqpoint{4.144874in}{1.496657in}}%
\pgfpathlineto{\pgfqpoint{4.115051in}{1.535191in}}%
\pgfpathlineto{\pgfqpoint{4.161776in}{1.551304in}}%
\pgfpathlineto{\pgfqpoint{4.208418in}{1.567380in}}%
\pgfpathlineto{\pgfqpoint{4.238295in}{1.528600in}}%
\pgfpathlineto{\pgfqpoint{4.268325in}{1.491721in}}%
\pgfpathlineto{\pgfqpoint{4.221627in}{1.475904in}}%
\pgfpathlineto{\pgfqpoint{4.174847in}{1.460057in}}%
\pgfpathclose%
\pgfusepath{stroke,fill}%
\end{pgfscope}%
\begin{pgfscope}%
\pgfpathrectangle{\pgfqpoint{0.050000in}{0.083252in}}{\pgfqpoint{4.900000in}{4.900000in}}%
\pgfusepath{clip}%
\pgfsetbuttcap%
\pgfsetroundjoin%
\definecolor{currentfill}{rgb}{0.260992,0.439216,0.852687}%
\pgfsetfillcolor{currentfill}%
\pgfsetfillopacity{0.950000}%
\pgfsetlinewidth{0.200750pt}%
\definecolor{currentstroke}{rgb}{0.200000,0.200000,0.200000}%
\pgfsetstrokecolor{currentstroke}%
\pgfsetdash{}{0pt}%
\pgfpathmoveto{\pgfqpoint{1.765665in}{3.528381in}}%
\pgfpathlineto{\pgfqpoint{1.738220in}{3.602639in}}%
\pgfpathlineto{\pgfqpoint{1.710800in}{3.676827in}}%
\pgfpathlineto{\pgfqpoint{1.759690in}{3.690491in}}%
\pgfpathlineto{\pgfqpoint{1.808494in}{3.704130in}}%
\pgfpathlineto{\pgfqpoint{1.835910in}{3.630088in}}%
\pgfpathlineto{\pgfqpoint{1.863353in}{3.555975in}}%
\pgfpathlineto{\pgfqpoint{1.814553in}{3.542191in}}%
\pgfpathlineto{\pgfqpoint{1.765665in}{3.528381in}}%
\pgfpathclose%
\pgfusepath{stroke,fill}%
\end{pgfscope}%
\begin{pgfscope}%
\pgfpathrectangle{\pgfqpoint{0.050000in}{0.083252in}}{\pgfqpoint{4.900000in}{4.900000in}}%
\pgfusepath{clip}%
\pgfsetbuttcap%
\pgfsetroundjoin%
\definecolor{currentfill}{rgb}{0.886628,0.881953,0.932211}%
\pgfsetfillcolor{currentfill}%
\pgfsetfillopacity{0.950000}%
\pgfsetlinewidth{0.200750pt}%
\definecolor{currentstroke}{rgb}{0.200000,0.200000,0.200000}%
\pgfsetstrokecolor{currentstroke}%
\pgfsetdash{}{0pt}%
\pgfpathmoveto{\pgfqpoint{3.621333in}{1.580297in}}%
\pgfpathlineto{\pgfqpoint{3.592438in}{1.631153in}}%
\pgfpathlineto{\pgfqpoint{3.563667in}{1.684843in}}%
\pgfpathlineto{\pgfqpoint{3.610911in}{1.701538in}}%
\pgfpathlineto{\pgfqpoint{3.658072in}{1.718185in}}%
\pgfpathlineto{\pgfqpoint{3.686882in}{1.664572in}}%
\pgfpathlineto{\pgfqpoint{3.715819in}{1.613718in}}%
\pgfpathlineto{\pgfqpoint{3.668617in}{1.597032in}}%
\pgfpathlineto{\pgfqpoint{3.621333in}{1.580297in}}%
\pgfpathclose%
\pgfusepath{stroke,fill}%
\end{pgfscope}%
\begin{pgfscope}%
\pgfpathrectangle{\pgfqpoint{0.050000in}{0.083252in}}{\pgfqpoint{4.900000in}{4.900000in}}%
\pgfusepath{clip}%
\pgfsetbuttcap%
\pgfsetroundjoin%
\definecolor{currentfill}{rgb}{0.791157,0.782545,0.875125}%
\pgfsetfillcolor{currentfill}%
\pgfsetfillopacity{0.950000}%
\pgfsetlinewidth{0.200750pt}%
\definecolor{currentstroke}{rgb}{0.200000,0.200000,0.200000}%
\pgfsetstrokecolor{currentstroke}%
\pgfsetdash{}{0pt}%
\pgfpathmoveto{\pgfqpoint{3.316797in}{1.733726in}}%
\pgfpathlineto{\pgfqpoint{3.288397in}{1.795645in}}%
\pgfpathlineto{\pgfqpoint{3.260077in}{1.859790in}}%
\pgfpathlineto{\pgfqpoint{3.307593in}{1.876174in}}%
\pgfpathlineto{\pgfqpoint{3.355026in}{1.892519in}}%
\pgfpathlineto{\pgfqpoint{3.383368in}{1.828707in}}%
\pgfpathlineto{\pgfqpoint{3.411794in}{1.767078in}}%
\pgfpathlineto{\pgfqpoint{3.364337in}{1.750426in}}%
\pgfpathlineto{\pgfqpoint{3.316797in}{1.733726in}}%
\pgfpathclose%
\pgfusepath{stroke,fill}%
\end{pgfscope}%
\begin{pgfscope}%
\pgfpathrectangle{\pgfqpoint{0.050000in}{0.083252in}}{\pgfqpoint{4.900000in}{4.900000in}}%
\pgfusepath{clip}%
\pgfsetbuttcap%
\pgfsetroundjoin%
\definecolor{currentfill}{rgb}{0.254348,0.368627,0.758847}%
\pgfsetfillcolor{currentfill}%
\pgfsetfillopacity{0.950000}%
\pgfsetlinewidth{0.200750pt}%
\definecolor{currentstroke}{rgb}{0.200000,0.200000,0.200000}%
\pgfsetstrokecolor{currentstroke}%
\pgfsetdash{}{0pt}%
\pgfpathmoveto{\pgfqpoint{1.973384in}{3.258825in}}%
\pgfpathlineto{\pgfqpoint{1.945837in}{3.333218in}}%
\pgfpathlineto{\pgfqpoint{1.918316in}{3.407540in}}%
\pgfpathlineto{\pgfqpoint{1.967027in}{3.421446in}}%
\pgfpathlineto{\pgfqpoint{2.015650in}{3.435326in}}%
\pgfpathlineto{\pgfqpoint{2.043167in}{3.361149in}}%
\pgfpathlineto{\pgfqpoint{2.070711in}{3.286902in}}%
\pgfpathlineto{\pgfqpoint{2.022091in}{3.272876in}}%
\pgfpathlineto{\pgfqpoint{1.973384in}{3.258825in}}%
\pgfpathclose%
\pgfusepath{stroke,fill}%
\end{pgfscope}%
\begin{pgfscope}%
\pgfpathrectangle{\pgfqpoint{0.050000in}{0.083252in}}{\pgfqpoint{4.900000in}{4.900000in}}%
\pgfusepath{clip}%
\pgfsetbuttcap%
\pgfsetroundjoin%
\definecolor{currentfill}{rgb}{0.247705,0.298039,0.665006}%
\pgfsetfillcolor{currentfill}%
\pgfsetfillopacity{0.950000}%
\pgfsetlinewidth{0.200750pt}%
\definecolor{currentstroke}{rgb}{0.200000,0.200000,0.200000}%
\pgfsetstrokecolor{currentstroke}%
\pgfsetdash{}{0pt}%
\pgfpathmoveto{\pgfqpoint{2.181145in}{2.989215in}}%
\pgfpathlineto{\pgfqpoint{2.153497in}{3.063742in}}%
\pgfpathlineto{\pgfqpoint{2.125875in}{3.138199in}}%
\pgfpathlineto{\pgfqpoint{2.174405in}{3.152346in}}%
\pgfpathlineto{\pgfqpoint{2.222849in}{3.166468in}}%
\pgfpathlineto{\pgfqpoint{2.250467in}{3.092158in}}%
\pgfpathlineto{\pgfqpoint{2.278111in}{3.017780in}}%
\pgfpathlineto{\pgfqpoint{2.229671in}{3.003510in}}%
\pgfpathlineto{\pgfqpoint{2.181145in}{2.989215in}}%
\pgfpathclose%
\pgfusepath{stroke,fill}%
\end{pgfscope}%
\begin{pgfscope}%
\pgfpathrectangle{\pgfqpoint{0.050000in}{0.083252in}}{\pgfqpoint{4.900000in}{4.900000in}}%
\pgfusepath{clip}%
\pgfsetbuttcap%
\pgfsetroundjoin%
\definecolor{currentfill}{rgb}{0.241061,0.227451,0.571165}%
\pgfsetfillcolor{currentfill}%
\pgfsetfillopacity{0.950000}%
\pgfsetlinewidth{0.200750pt}%
\definecolor{currentstroke}{rgb}{0.200000,0.200000,0.200000}%
\pgfsetstrokecolor{currentstroke}%
\pgfsetdash{}{0pt}%
\pgfpathmoveto{\pgfqpoint{2.388949in}{2.719609in}}%
\pgfpathlineto{\pgfqpoint{2.361200in}{2.794244in}}%
\pgfpathlineto{\pgfqpoint{2.333477in}{2.868821in}}%
\pgfpathlineto{\pgfqpoint{2.381827in}{2.883216in}}%
\pgfpathlineto{\pgfqpoint{2.430090in}{2.897586in}}%
\pgfpathlineto{\pgfqpoint{2.457809in}{2.823165in}}%
\pgfpathlineto{\pgfqpoint{2.485554in}{2.748689in}}%
\pgfpathlineto{\pgfqpoint{2.437295in}{2.734161in}}%
\pgfpathlineto{\pgfqpoint{2.388949in}{2.719609in}}%
\pgfpathclose%
\pgfusepath{stroke,fill}%
\end{pgfscope}%
\begin{pgfscope}%
\pgfpathrectangle{\pgfqpoint{0.050000in}{0.083252in}}{\pgfqpoint{4.900000in}{4.900000in}}%
\pgfusepath{clip}%
\pgfsetbuttcap%
\pgfsetroundjoin%
\definecolor{currentfill}{rgb}{0.922430,0.919231,0.953618}%
\pgfsetfillcolor{currentfill}%
\pgfsetfillopacity{0.950000}%
\pgfsetlinewidth{0.200750pt}%
\definecolor{currentstroke}{rgb}{0.200000,0.200000,0.200000}%
\pgfsetstrokecolor{currentstroke}%
\pgfsetdash{}{0pt}%
\pgfpathmoveto{\pgfqpoint{3.774105in}{1.520399in}}%
\pgfpathlineto{\pgfqpoint{3.744892in}{1.565668in}}%
\pgfpathlineto{\pgfqpoint{3.715819in}{1.613718in}}%
\pgfpathlineto{\pgfqpoint{3.762938in}{1.630359in}}%
\pgfpathlineto{\pgfqpoint{3.809976in}{1.646953in}}%
\pgfpathlineto{\pgfqpoint{3.839094in}{1.598838in}}%
\pgfpathlineto{\pgfqpoint{3.868356in}{1.553439in}}%
\pgfpathlineto{\pgfqpoint{3.821272in}{1.536940in}}%
\pgfpathlineto{\pgfqpoint{3.774105in}{1.520399in}}%
\pgfpathclose%
\pgfusepath{stroke,fill}%
\end{pgfscope}%
\begin{pgfscope}%
\pgfpathrectangle{\pgfqpoint{0.050000in}{0.083252in}}{\pgfqpoint{4.900000in}{4.900000in}}%
\pgfusepath{clip}%
\pgfsetbuttcap%
\pgfsetroundjoin%
\definecolor{currentfill}{rgb}{0.349604,0.322784,0.611103}%
\pgfsetfillcolor{currentfill}%
\pgfsetfillopacity{0.950000}%
\pgfsetlinewidth{0.200750pt}%
\definecolor{currentstroke}{rgb}{0.200000,0.200000,0.200000}%
\pgfsetstrokecolor{currentstroke}%
\pgfsetdash{}{0pt}%
\pgfpathmoveto{\pgfqpoint{2.596797in}{2.450501in}}%
\pgfpathlineto{\pgfqpoint{2.568946in}{2.525049in}}%
\pgfpathlineto{\pgfqpoint{2.541123in}{2.599614in}}%
\pgfpathlineto{\pgfqpoint{2.589292in}{2.614292in}}%
\pgfpathlineto{\pgfqpoint{2.637376in}{2.628949in}}%
\pgfpathlineto{\pgfqpoint{2.665196in}{2.554576in}}%
\pgfpathlineto{\pgfqpoint{2.693043in}{2.480238in}}%
\pgfpathlineto{\pgfqpoint{2.644963in}{2.465379in}}%
\pgfpathlineto{\pgfqpoint{2.596797in}{2.450501in}}%
\pgfpathclose%
\pgfusepath{stroke,fill}%
\end{pgfscope}%
\begin{pgfscope}%
\pgfpathrectangle{\pgfqpoint{0.050000in}{0.083252in}}{\pgfqpoint{4.900000in}{4.900000in}}%
\pgfusepath{clip}%
\pgfsetbuttcap%
\pgfsetroundjoin%
\definecolor{currentfill}{rgb}{0.725521,0.714202,0.835879}%
\pgfsetfillcolor{currentfill}%
\pgfsetfillopacity{0.950000}%
\pgfsetlinewidth{0.200750pt}%
\definecolor{currentstroke}{rgb}{0.200000,0.200000,0.200000}%
\pgfsetstrokecolor{currentstroke}%
\pgfsetdash{}{0pt}%
\pgfpathmoveto{\pgfqpoint{3.164797in}{1.826905in}}%
\pgfpathlineto{\pgfqpoint{3.136565in}{1.893386in}}%
\pgfpathlineto{\pgfqpoint{3.108391in}{1.961559in}}%
\pgfpathlineto{\pgfqpoint{3.156057in}{1.977627in}}%
\pgfpathlineto{\pgfqpoint{3.203639in}{1.993666in}}%
\pgfpathlineto{\pgfqpoint{3.231827in}{1.925889in}}%
\pgfpathlineto{\pgfqpoint{3.260077in}{1.859790in}}%
\pgfpathlineto{\pgfqpoint{3.212478in}{1.843367in}}%
\pgfpathlineto{\pgfqpoint{3.164797in}{1.826905in}}%
\pgfpathclose%
\pgfusepath{stroke,fill}%
\end{pgfscope}%
\begin{pgfscope}%
\pgfpathrectangle{\pgfqpoint{0.050000in}{0.083252in}}{\pgfqpoint{4.900000in}{4.900000in}}%
\pgfusepath{clip}%
\pgfsetbuttcap%
\pgfsetroundjoin%
\definecolor{currentfill}{rgb}{0.510711,0.490534,0.707436}%
\pgfsetfillcolor{currentfill}%
\pgfsetfillopacity{0.950000}%
\pgfsetlinewidth{0.200750pt}%
\definecolor{currentstroke}{rgb}{0.200000,0.200000,0.200000}%
\pgfsetstrokecolor{currentstroke}%
\pgfsetdash{}{0pt}%
\pgfpathmoveto{\pgfqpoint{2.804710in}{2.184436in}}%
\pgfpathlineto{\pgfqpoint{2.776749in}{2.257982in}}%
\pgfpathlineto{\pgfqpoint{2.748819in}{2.331867in}}%
\pgfpathlineto{\pgfqpoint{2.796811in}{2.346953in}}%
\pgfpathlineto{\pgfqpoint{2.844717in}{2.362022in}}%
\pgfpathlineto{\pgfqpoint{2.872647in}{2.288426in}}%
\pgfpathlineto{\pgfqpoint{2.900607in}{2.215200in}}%
\pgfpathlineto{\pgfqpoint{2.852701in}{2.199827in}}%
\pgfpathlineto{\pgfqpoint{2.804710in}{2.184436in}}%
\pgfpathclose%
\pgfusepath{stroke,fill}%
\end{pgfscope}%
\begin{pgfscope}%
\pgfpathrectangle{\pgfqpoint{0.050000in}{0.083252in}}{\pgfqpoint{4.900000in}{4.900000in}}%
\pgfusepath{clip}%
\pgfsetbuttcap%
\pgfsetroundjoin%
\definecolor{currentfill}{rgb}{0.264191,0.473203,0.897870}%
\pgfsetfillcolor{currentfill}%
\pgfsetfillopacity{0.950000}%
\pgfsetlinewidth{0.200750pt}%
\definecolor{currentstroke}{rgb}{0.200000,0.200000,0.200000}%
\pgfsetstrokecolor{currentstroke}%
\pgfsetdash{}{0pt}%
\pgfpathmoveto{\pgfqpoint{1.612757in}{3.649425in}}%
\pgfpathlineto{\pgfqpoint{1.585361in}{3.723689in}}%
\pgfpathlineto{\pgfqpoint{1.557990in}{3.797882in}}%
\pgfpathlineto{\pgfqpoint{1.607058in}{3.811449in}}%
\pgfpathlineto{\pgfqpoint{1.656039in}{3.824992in}}%
\pgfpathlineto{\pgfqpoint{1.683406in}{3.750945in}}%
\pgfpathlineto{\pgfqpoint{1.710800in}{3.676827in}}%
\pgfpathlineto{\pgfqpoint{1.661822in}{3.663139in}}%
\pgfpathlineto{\pgfqpoint{1.612757in}{3.649425in}}%
\pgfpathclose%
\pgfusepath{stroke,fill}%
\end{pgfscope}%
\begin{pgfscope}%
\pgfpathrectangle{\pgfqpoint{0.050000in}{0.083252in}}{\pgfqpoint{4.900000in}{4.900000in}}%
\pgfusepath{clip}%
\pgfsetbuttcap%
\pgfsetroundjoin%
\definecolor{currentfill}{rgb}{0.257547,0.402614,0.804029}%
\pgfsetfillcolor{currentfill}%
\pgfsetfillopacity{0.950000}%
\pgfsetlinewidth{0.200750pt}%
\definecolor{currentstroke}{rgb}{0.200000,0.200000,0.200000}%
\pgfsetstrokecolor{currentstroke}%
\pgfsetdash{}{0pt}%
\pgfpathmoveto{\pgfqpoint{1.820634in}{3.379655in}}%
\pgfpathlineto{\pgfqpoint{1.793137in}{3.454053in}}%
\pgfpathlineto{\pgfqpoint{1.765665in}{3.528381in}}%
\pgfpathlineto{\pgfqpoint{1.814553in}{3.542191in}}%
\pgfpathlineto{\pgfqpoint{1.863353in}{3.555975in}}%
\pgfpathlineto{\pgfqpoint{1.890822in}{3.481793in}}%
\pgfpathlineto{\pgfqpoint{1.918316in}{3.407540in}}%
\pgfpathlineto{\pgfqpoint{1.869519in}{3.393610in}}%
\pgfpathlineto{\pgfqpoint{1.820634in}{3.379655in}}%
\pgfpathclose%
\pgfusepath{stroke,fill}%
\end{pgfscope}%
\begin{pgfscope}%
\pgfpathrectangle{\pgfqpoint{0.050000in}{0.083252in}}{\pgfqpoint{4.900000in}{4.900000in}}%
\pgfusepath{clip}%
\pgfsetbuttcap%
\pgfsetroundjoin%
\definecolor{currentfill}{rgb}{0.659885,0.645859,0.796632}%
\pgfsetfillcolor{currentfill}%
\pgfsetfillopacity{0.950000}%
\pgfsetlinewidth{0.200750pt}%
\definecolor{currentstroke}{rgb}{0.200000,0.200000,0.200000}%
\pgfsetstrokecolor{currentstroke}%
\pgfsetdash{}{0pt}%
\pgfpathmoveto{\pgfqpoint{3.012809in}{1.929333in}}%
\pgfpathlineto{\pgfqpoint{2.984698in}{1.999329in}}%
\pgfpathlineto{\pgfqpoint{2.956629in}{2.070460in}}%
\pgfpathlineto{\pgfqpoint{3.004454in}{2.086179in}}%
\pgfpathlineto{\pgfqpoint{3.052195in}{2.101877in}}%
\pgfpathlineto{\pgfqpoint{3.080270in}{2.031142in}}%
\pgfpathlineto{\pgfqpoint{3.108391in}{1.961559in}}%
\pgfpathlineto{\pgfqpoint{3.060642in}{1.945460in}}%
\pgfpathlineto{\pgfqpoint{3.012809in}{1.929333in}}%
\pgfpathclose%
\pgfusepath{stroke,fill}%
\end{pgfscope}%
\begin{pgfscope}%
\pgfpathrectangle{\pgfqpoint{0.050000in}{0.083252in}}{\pgfqpoint{4.900000in}{4.900000in}}%
\pgfusepath{clip}%
\pgfsetbuttcap%
\pgfsetroundjoin%
\definecolor{currentfill}{rgb}{0.269850,0.533333,0.977809}%
\pgfsetfillcolor{currentfill}%
\pgfsetfillopacity{0.950000}%
\pgfsetlinewidth{0.200750pt}%
\definecolor{currentstroke}{rgb}{0.200000,0.200000,0.200000}%
\pgfsetstrokecolor{currentstroke}%
\pgfsetdash{}{0pt}%
\pgfpathmoveto{\pgfqpoint{1.404922in}{3.919140in}}%
\pgfpathlineto{\pgfqpoint{1.377627in}{3.993268in}}%
\pgfpathlineto{\pgfqpoint{1.426875in}{4.006666in}}%
\pgfpathlineto{\pgfqpoint{1.476034in}{4.020040in}}%
\pgfpathlineto{\pgfqpoint{1.503327in}{3.946057in}}%
\pgfpathlineto{\pgfqpoint{1.454169in}{3.932611in}}%
\pgfpathlineto{\pgfqpoint{1.404922in}{3.919140in}}%
\pgfpathclose%
\pgfusepath{stroke,fill}%
\end{pgfscope}%
\begin{pgfscope}%
\pgfpathrectangle{\pgfqpoint{0.050000in}{0.083252in}}{\pgfqpoint{4.900000in}{4.900000in}}%
\pgfusepath{clip}%
\pgfsetbuttcap%
\pgfsetroundjoin%
\definecolor{currentfill}{rgb}{0.952265,0.950296,0.971457}%
\pgfsetfillcolor{currentfill}%
\pgfsetfillopacity{0.950000}%
\pgfsetlinewidth{0.200750pt}%
\definecolor{currentstroke}{rgb}{0.200000,0.200000,0.200000}%
\pgfsetstrokecolor{currentstroke}%
\pgfsetdash{}{0pt}%
\pgfpathmoveto{\pgfqpoint{3.927327in}{1.470389in}}%
\pgfpathlineto{\pgfqpoint{3.897767in}{1.510672in}}%
\pgfpathlineto{\pgfqpoint{3.868356in}{1.553439in}}%
\pgfpathlineto{\pgfqpoint{3.915358in}{1.569897in}}%
\pgfpathlineto{\pgfqpoint{3.962278in}{1.586314in}}%
\pgfpathlineto{\pgfqpoint{3.991740in}{1.543370in}}%
\pgfpathlineto{\pgfqpoint{4.021355in}{1.502861in}}%
\pgfpathlineto{\pgfqpoint{3.974382in}{1.486643in}}%
\pgfpathlineto{\pgfqpoint{3.927327in}{1.470389in}}%
\pgfpathclose%
\pgfusepath{stroke,fill}%
\end{pgfscope}%
\begin{pgfscope}%
\pgfpathrectangle{\pgfqpoint{0.050000in}{0.083252in}}{\pgfqpoint{4.900000in}{4.900000in}}%
\pgfusepath{clip}%
\pgfsetbuttcap%
\pgfsetroundjoin%
\definecolor{currentfill}{rgb}{0.250903,0.332026,0.710188}%
\pgfsetfillcolor{currentfill}%
\pgfsetfillopacity{0.950000}%
\pgfsetlinewidth{0.200750pt}%
\definecolor{currentstroke}{rgb}{0.200000,0.200000,0.200000}%
\pgfsetstrokecolor{currentstroke}%
\pgfsetdash{}{0pt}%
\pgfpathmoveto{\pgfqpoint{2.028555in}{3.109828in}}%
\pgfpathlineto{\pgfqpoint{2.000956in}{3.184361in}}%
\pgfpathlineto{\pgfqpoint{1.973384in}{3.258825in}}%
\pgfpathlineto{\pgfqpoint{2.022091in}{3.272876in}}%
\pgfpathlineto{\pgfqpoint{2.070711in}{3.286902in}}%
\pgfpathlineto{\pgfqpoint{2.098280in}{3.212585in}}%
\pgfpathlineto{\pgfqpoint{2.125875in}{3.138199in}}%
\pgfpathlineto{\pgfqpoint{2.077259in}{3.124026in}}%
\pgfpathlineto{\pgfqpoint{2.028555in}{3.109828in}}%
\pgfpathclose%
\pgfusepath{stroke,fill}%
\end{pgfscope}%
\begin{pgfscope}%
\pgfpathrectangle{\pgfqpoint{0.050000in}{0.083252in}}{\pgfqpoint{4.900000in}{4.900000in}}%
\pgfusepath{clip}%
\pgfsetbuttcap%
\pgfsetroundjoin%
\definecolor{currentfill}{rgb}{0.244506,0.264052,0.619823}%
\pgfsetfillcolor{currentfill}%
\pgfsetfillopacity{0.950000}%
\pgfsetlinewidth{0.200750pt}%
\definecolor{currentstroke}{rgb}{0.200000,0.200000,0.200000}%
\pgfsetstrokecolor{currentstroke}%
\pgfsetdash{}{0pt}%
\pgfpathmoveto{\pgfqpoint{2.236519in}{2.839956in}}%
\pgfpathlineto{\pgfqpoint{2.208819in}{2.914620in}}%
\pgfpathlineto{\pgfqpoint{2.181145in}{2.989215in}}%
\pgfpathlineto{\pgfqpoint{2.229671in}{3.003510in}}%
\pgfpathlineto{\pgfqpoint{2.278111in}{3.017780in}}%
\pgfpathlineto{\pgfqpoint{2.305781in}{2.943333in}}%
\pgfpathlineto{\pgfqpoint{2.333477in}{2.868821in}}%
\pgfpathlineto{\pgfqpoint{2.285041in}{2.854401in}}%
\pgfpathlineto{\pgfqpoint{2.236519in}{2.839956in}}%
\pgfpathclose%
\pgfusepath{stroke,fill}%
\end{pgfscope}%
\begin{pgfscope}%
\pgfpathrectangle{\pgfqpoint{0.050000in}{0.083252in}}{\pgfqpoint{4.900000in}{4.900000in}}%
\pgfusepath{clip}%
\pgfsetbuttcap%
\pgfsetroundjoin%
\definecolor{currentfill}{rgb}{0.272034,0.242015,0.564721}%
\pgfsetfillcolor{currentfill}%
\pgfsetfillopacity{0.950000}%
\pgfsetlinewidth{0.200750pt}%
\definecolor{currentstroke}{rgb}{0.200000,0.200000,0.200000}%
\pgfsetstrokecolor{currentstroke}%
\pgfsetdash{}{0pt}%
\pgfpathmoveto{\pgfqpoint{2.444525in}{2.570193in}}%
\pgfpathlineto{\pgfqpoint{2.416724in}{2.644921in}}%
\pgfpathlineto{\pgfqpoint{2.388949in}{2.719609in}}%
\pgfpathlineto{\pgfqpoint{2.437295in}{2.734161in}}%
\pgfpathlineto{\pgfqpoint{2.485554in}{2.748689in}}%
\pgfpathlineto{\pgfqpoint{2.513325in}{2.674168in}}%
\pgfpathlineto{\pgfqpoint{2.541123in}{2.599614in}}%
\pgfpathlineto{\pgfqpoint{2.492867in}{2.584914in}}%
\pgfpathlineto{\pgfqpoint{2.444525in}{2.570193in}}%
\pgfpathclose%
\pgfusepath{stroke,fill}%
\end{pgfscope}%
\begin{pgfscope}%
\pgfpathrectangle{\pgfqpoint{0.050000in}{0.083252in}}{\pgfqpoint{4.900000in}{4.900000in}}%
\pgfusepath{clip}%
\pgfsetbuttcap%
\pgfsetroundjoin%
\definecolor{currentfill}{rgb}{0.982099,0.981361,0.989296}%
\pgfsetfillcolor{currentfill}%
\pgfsetfillopacity{0.950000}%
\pgfsetlinewidth{0.200750pt}%
\definecolor{currentstroke}{rgb}{0.200000,0.200000,0.200000}%
\pgfsetstrokecolor{currentstroke}%
\pgfsetdash{}{0pt}%
\pgfpathmoveto{\pgfqpoint{4.328809in}{1.422365in}}%
\pgfpathlineto{\pgfqpoint{4.298500in}{1.456432in}}%
\pgfpathlineto{\pgfqpoint{4.268325in}{1.491721in}}%
\pgfpathlineto{\pgfqpoint{4.314939in}{1.507506in}}%
\pgfpathlineto{\pgfqpoint{4.345142in}{1.472067in}}%
\pgfpathlineto{\pgfqpoint{4.375481in}{1.437839in}}%
\pgfpathlineto{\pgfqpoint{4.328809in}{1.422365in}}%
\pgfpathclose%
\pgfusepath{stroke,fill}%
\end{pgfscope}%
\begin{pgfscope}%
\pgfpathrectangle{\pgfqpoint{0.050000in}{0.083252in}}{\pgfqpoint{4.900000in}{4.900000in}}%
\pgfusepath{clip}%
\pgfsetbuttcap%
\pgfsetroundjoin%
\definecolor{currentfill}{rgb}{0.433141,0.409765,0.661053}%
\pgfsetfillcolor{currentfill}%
\pgfsetfillopacity{0.950000}%
\pgfsetlinewidth{0.200750pt}%
\definecolor{currentstroke}{rgb}{0.200000,0.200000,0.200000}%
\pgfsetstrokecolor{currentstroke}%
\pgfsetdash{}{0pt}%
\pgfpathmoveto{\pgfqpoint{2.652579in}{2.301643in}}%
\pgfpathlineto{\pgfqpoint{2.624674in}{2.376013in}}%
\pgfpathlineto{\pgfqpoint{2.596797in}{2.450501in}}%
\pgfpathlineto{\pgfqpoint{2.644963in}{2.465379in}}%
\pgfpathlineto{\pgfqpoint{2.693043in}{2.480238in}}%
\pgfpathlineto{\pgfqpoint{2.720917in}{2.405980in}}%
\pgfpathlineto{\pgfqpoint{2.748819in}{2.331867in}}%
\pgfpathlineto{\pgfqpoint{2.700742in}{2.316763in}}%
\pgfpathlineto{\pgfqpoint{2.652579in}{2.301643in}}%
\pgfpathclose%
\pgfusepath{stroke,fill}%
\end{pgfscope}%
\begin{pgfscope}%
\pgfpathrectangle{\pgfqpoint{0.050000in}{0.083252in}}{\pgfqpoint{4.900000in}{4.900000in}}%
\pgfusepath{clip}%
\pgfsetbuttcap%
\pgfsetroundjoin%
\definecolor{currentfill}{rgb}{0.588281,0.571303,0.753818}%
\pgfsetfillcolor{currentfill}%
\pgfsetfillopacity{0.950000}%
\pgfsetlinewidth{0.200750pt}%
\definecolor{currentstroke}{rgb}{0.200000,0.200000,0.200000}%
\pgfsetstrokecolor{currentstroke}%
\pgfsetdash{}{0pt}%
\pgfpathmoveto{\pgfqpoint{2.860727in}{2.038960in}}%
\pgfpathlineto{\pgfqpoint{2.832702in}{2.111370in}}%
\pgfpathlineto{\pgfqpoint{2.804710in}{2.184436in}}%
\pgfpathlineto{\pgfqpoint{2.852701in}{2.199827in}}%
\pgfpathlineto{\pgfqpoint{2.900607in}{2.215200in}}%
\pgfpathlineto{\pgfqpoint{2.928600in}{2.142487in}}%
\pgfpathlineto{\pgfqpoint{2.956629in}{2.070460in}}%
\pgfpathlineto{\pgfqpoint{2.908721in}{2.054720in}}%
\pgfpathlineto{\pgfqpoint{2.860727in}{2.038960in}}%
\pgfpathclose%
\pgfusepath{stroke,fill}%
\end{pgfscope}%
\begin{pgfscope}%
\pgfpathrectangle{\pgfqpoint{0.050000in}{0.083252in}}{\pgfqpoint{4.900000in}{4.900000in}}%
\pgfusepath{clip}%
\pgfsetbuttcap%
\pgfsetroundjoin%
\definecolor{currentfill}{rgb}{0.844860,0.838462,0.907236}%
\pgfsetfillcolor{currentfill}%
\pgfsetfillopacity{0.950000}%
\pgfsetlinewidth{0.200750pt}%
\definecolor{currentstroke}{rgb}{0.200000,0.200000,0.200000}%
\pgfsetstrokecolor{currentstroke}%
\pgfsetdash{}{0pt}%
\pgfpathmoveto{\pgfqpoint{3.373871in}{1.617563in}}%
\pgfpathlineto{\pgfqpoint{3.345285in}{1.674291in}}%
\pgfpathlineto{\pgfqpoint{3.316797in}{1.733726in}}%
\pgfpathlineto{\pgfqpoint{3.364337in}{1.750426in}}%
\pgfpathlineto{\pgfqpoint{3.411794in}{1.767078in}}%
\pgfpathlineto{\pgfqpoint{3.440313in}{1.707874in}}%
\pgfpathlineto{\pgfqpoint{3.468934in}{1.651304in}}%
\pgfpathlineto{\pgfqpoint{3.421444in}{1.634459in}}%
\pgfpathlineto{\pgfqpoint{3.373871in}{1.617563in}}%
\pgfpathclose%
\pgfusepath{stroke,fill}%
\end{pgfscope}%
\begin{pgfscope}%
\pgfpathrectangle{\pgfqpoint{0.050000in}{0.083252in}}{\pgfqpoint{4.900000in}{4.900000in}}%
\pgfusepath{clip}%
\pgfsetbuttcap%
\pgfsetroundjoin%
\definecolor{currentfill}{rgb}{0.267389,0.507190,0.943053}%
\pgfsetfillcolor{currentfill}%
\pgfsetfillopacity{0.950000}%
\pgfsetlinewidth{0.200750pt}%
\definecolor{currentstroke}{rgb}{0.200000,0.200000,0.200000}%
\pgfsetstrokecolor{currentstroke}%
\pgfsetdash{}{0pt}%
\pgfpathmoveto{\pgfqpoint{1.459590in}{3.770673in}}%
\pgfpathlineto{\pgfqpoint{1.432243in}{3.844942in}}%
\pgfpathlineto{\pgfqpoint{1.404922in}{3.919140in}}%
\pgfpathlineto{\pgfqpoint{1.454169in}{3.932611in}}%
\pgfpathlineto{\pgfqpoint{1.503327in}{3.946057in}}%
\pgfpathlineto{\pgfqpoint{1.530646in}{3.872004in}}%
\pgfpathlineto{\pgfqpoint{1.557990in}{3.797882in}}%
\pgfpathlineto{\pgfqpoint{1.508834in}{3.784290in}}%
\pgfpathlineto{\pgfqpoint{1.459590in}{3.770673in}}%
\pgfpathclose%
\pgfusepath{stroke,fill}%
\end{pgfscope}%
\begin{pgfscope}%
\pgfpathrectangle{\pgfqpoint{0.050000in}{0.083252in}}{\pgfqpoint{4.900000in}{4.900000in}}%
\pgfusepath{clip}%
\pgfsetbuttcap%
\pgfsetroundjoin%
\definecolor{currentfill}{rgb}{0.970165,0.968935,0.982161}%
\pgfsetfillcolor{currentfill}%
\pgfsetfillopacity{0.950000}%
\pgfsetlinewidth{0.200750pt}%
\definecolor{currentstroke}{rgb}{0.200000,0.200000,0.200000}%
\pgfsetstrokecolor{currentstroke}%
\pgfsetdash{}{0pt}%
\pgfpathmoveto{\pgfqpoint{4.081035in}{1.428269in}}%
\pgfpathlineto{\pgfqpoint{4.051121in}{1.464582in}}%
\pgfpathlineto{\pgfqpoint{4.021355in}{1.502861in}}%
\pgfpathlineto{\pgfqpoint{4.068245in}{1.519044in}}%
\pgfpathlineto{\pgfqpoint{4.115051in}{1.535191in}}%
\pgfpathlineto{\pgfqpoint{4.144874in}{1.496657in}}%
\pgfpathlineto{\pgfqpoint{4.174847in}{1.460057in}}%
\pgfpathlineto{\pgfqpoint{4.127983in}{1.444178in}}%
\pgfpathlineto{\pgfqpoint{4.081035in}{1.428269in}}%
\pgfpathclose%
\pgfusepath{stroke,fill}%
\end{pgfscope}%
\begin{pgfscope}%
\pgfpathrectangle{\pgfqpoint{0.050000in}{0.083252in}}{\pgfqpoint{4.900000in}{4.900000in}}%
\pgfusepath{clip}%
\pgfsetbuttcap%
\pgfsetroundjoin%
\definecolor{currentfill}{rgb}{0.892595,0.888166,0.935779}%
\pgfsetfillcolor{currentfill}%
\pgfsetfillopacity{0.950000}%
\pgfsetlinewidth{0.200750pt}%
\definecolor{currentstroke}{rgb}{0.200000,0.200000,0.200000}%
\pgfsetstrokecolor{currentstroke}%
\pgfsetdash{}{0pt}%
\pgfpathmoveto{\pgfqpoint{3.526517in}{1.546684in}}%
\pgfpathlineto{\pgfqpoint{3.497666in}{1.597538in}}%
\pgfpathlineto{\pgfqpoint{3.468934in}{1.651304in}}%
\pgfpathlineto{\pgfqpoint{3.516342in}{1.668098in}}%
\pgfpathlineto{\pgfqpoint{3.563667in}{1.684843in}}%
\pgfpathlineto{\pgfqpoint{3.592438in}{1.631153in}}%
\pgfpathlineto{\pgfqpoint{3.621333in}{1.580297in}}%
\pgfpathlineto{\pgfqpoint{3.573967in}{1.563515in}}%
\pgfpathlineto{\pgfqpoint{3.526517in}{1.546684in}}%
\pgfpathclose%
\pgfusepath{stroke,fill}%
\end{pgfscope}%
\begin{pgfscope}%
\pgfpathrectangle{\pgfqpoint{0.050000in}{0.083252in}}{\pgfqpoint{4.900000in}{4.900000in}}%
\pgfusepath{clip}%
\pgfsetbuttcap%
\pgfsetroundjoin%
\definecolor{currentfill}{rgb}{0.260746,0.436601,0.849212}%
\pgfsetfillcolor{currentfill}%
\pgfsetfillopacity{0.950000}%
\pgfsetlinewidth{0.200750pt}%
\definecolor{currentstroke}{rgb}{0.200000,0.200000,0.200000}%
\pgfsetstrokecolor{currentstroke}%
\pgfsetdash{}{0pt}%
\pgfpathmoveto{\pgfqpoint{1.667627in}{3.500688in}}%
\pgfpathlineto{\pgfqpoint{1.640179in}{3.575092in}}%
\pgfpathlineto{\pgfqpoint{1.612757in}{3.649425in}}%
\pgfpathlineto{\pgfqpoint{1.661822in}{3.663139in}}%
\pgfpathlineto{\pgfqpoint{1.710800in}{3.676827in}}%
\pgfpathlineto{\pgfqpoint{1.738220in}{3.602639in}}%
\pgfpathlineto{\pgfqpoint{1.765665in}{3.528381in}}%
\pgfpathlineto{\pgfqpoint{1.716690in}{3.514547in}}%
\pgfpathlineto{\pgfqpoint{1.667627in}{3.500688in}}%
\pgfpathclose%
\pgfusepath{stroke,fill}%
\end{pgfscope}%
\begin{pgfscope}%
\pgfpathrectangle{\pgfqpoint{0.050000in}{0.083252in}}{\pgfqpoint{4.900000in}{4.900000in}}%
\pgfusepath{clip}%
\pgfsetbuttcap%
\pgfsetroundjoin%
\definecolor{currentfill}{rgb}{0.791157,0.782545,0.875125}%
\pgfsetfillcolor{currentfill}%
\pgfsetfillopacity{0.950000}%
\pgfsetlinewidth{0.200750pt}%
\definecolor{currentstroke}{rgb}{0.200000,0.200000,0.200000}%
\pgfsetstrokecolor{currentstroke}%
\pgfsetdash{}{0pt}%
\pgfpathmoveto{\pgfqpoint{3.221469in}{1.700183in}}%
\pgfpathlineto{\pgfqpoint{3.193096in}{1.762407in}}%
\pgfpathlineto{\pgfqpoint{3.164797in}{1.826905in}}%
\pgfpathlineto{\pgfqpoint{3.212478in}{1.843367in}}%
\pgfpathlineto{\pgfqpoint{3.260077in}{1.859790in}}%
\pgfpathlineto{\pgfqpoint{3.288397in}{1.795645in}}%
\pgfpathlineto{\pgfqpoint{3.316797in}{1.733726in}}%
\pgfpathlineto{\pgfqpoint{3.269174in}{1.716979in}}%
\pgfpathlineto{\pgfqpoint{3.221469in}{1.700183in}}%
\pgfpathclose%
\pgfusepath{stroke,fill}%
\end{pgfscope}%
\begin{pgfscope}%
\pgfpathrectangle{\pgfqpoint{0.050000in}{0.083252in}}{\pgfqpoint{4.900000in}{4.900000in}}%
\pgfusepath{clip}%
\pgfsetbuttcap%
\pgfsetroundjoin%
\definecolor{currentfill}{rgb}{0.254348,0.368627,0.758847}%
\pgfsetfillcolor{currentfill}%
\pgfsetfillopacity{0.950000}%
\pgfsetlinewidth{0.200750pt}%
\definecolor{currentstroke}{rgb}{0.200000,0.200000,0.200000}%
\pgfsetstrokecolor{currentstroke}%
\pgfsetdash{}{0pt}%
\pgfpathmoveto{\pgfqpoint{1.875708in}{3.230646in}}%
\pgfpathlineto{\pgfqpoint{1.848158in}{3.305186in}}%
\pgfpathlineto{\pgfqpoint{1.820634in}{3.379655in}}%
\pgfpathlineto{\pgfqpoint{1.869519in}{3.393610in}}%
\pgfpathlineto{\pgfqpoint{1.918316in}{3.407540in}}%
\pgfpathlineto{\pgfqpoint{1.945837in}{3.333218in}}%
\pgfpathlineto{\pgfqpoint{1.973384in}{3.258825in}}%
\pgfpathlineto{\pgfqpoint{1.924589in}{3.244748in}}%
\pgfpathlineto{\pgfqpoint{1.875708in}{3.230646in}}%
\pgfpathclose%
\pgfusepath{stroke,fill}%
\end{pgfscope}%
\begin{pgfscope}%
\pgfpathrectangle{\pgfqpoint{0.050000in}{0.083252in}}{\pgfqpoint{4.900000in}{4.900000in}}%
\pgfusepath{clip}%
\pgfsetbuttcap%
\pgfsetroundjoin%
\definecolor{currentfill}{rgb}{0.247705,0.298039,0.665006}%
\pgfsetfillcolor{currentfill}%
\pgfsetfillopacity{0.950000}%
\pgfsetlinewidth{0.200750pt}%
\definecolor{currentstroke}{rgb}{0.200000,0.200000,0.200000}%
\pgfsetstrokecolor{currentstroke}%
\pgfsetdash{}{0pt}%
\pgfpathmoveto{\pgfqpoint{2.083831in}{2.960549in}}%
\pgfpathlineto{\pgfqpoint{2.056180in}{3.035223in}}%
\pgfpathlineto{\pgfqpoint{2.028555in}{3.109828in}}%
\pgfpathlineto{\pgfqpoint{2.077259in}{3.124026in}}%
\pgfpathlineto{\pgfqpoint{2.125875in}{3.138199in}}%
\pgfpathlineto{\pgfqpoint{2.153497in}{3.063742in}}%
\pgfpathlineto{\pgfqpoint{2.181145in}{2.989215in}}%
\pgfpathlineto{\pgfqpoint{2.132531in}{2.974895in}}%
\pgfpathlineto{\pgfqpoint{2.083831in}{2.960549in}}%
\pgfpathclose%
\pgfusepath{stroke,fill}%
\end{pgfscope}%
\begin{pgfscope}%
\pgfpathrectangle{\pgfqpoint{0.050000in}{0.083252in}}{\pgfqpoint{4.900000in}{4.900000in}}%
\pgfusepath{clip}%
\pgfsetbuttcap%
\pgfsetroundjoin%
\definecolor{currentfill}{rgb}{0.241061,0.227451,0.571165}%
\pgfsetfillcolor{currentfill}%
\pgfsetfillopacity{0.950000}%
\pgfsetlinewidth{0.200750pt}%
\definecolor{currentstroke}{rgb}{0.200000,0.200000,0.200000}%
\pgfsetstrokecolor{currentstroke}%
\pgfsetdash{}{0pt}%
\pgfpathmoveto{\pgfqpoint{2.291997in}{2.690433in}}%
\pgfpathlineto{\pgfqpoint{2.264245in}{2.765226in}}%
\pgfpathlineto{\pgfqpoint{2.236519in}{2.839956in}}%
\pgfpathlineto{\pgfqpoint{2.285041in}{2.854401in}}%
\pgfpathlineto{\pgfqpoint{2.333477in}{2.868821in}}%
\pgfpathlineto{\pgfqpoint{2.361200in}{2.794244in}}%
\pgfpathlineto{\pgfqpoint{2.388949in}{2.719609in}}%
\pgfpathlineto{\pgfqpoint{2.340516in}{2.705033in}}%
\pgfpathlineto{\pgfqpoint{2.291997in}{2.690433in}}%
\pgfpathclose%
\pgfusepath{stroke,fill}%
\end{pgfscope}%
\begin{pgfscope}%
\pgfpathrectangle{\pgfqpoint{0.050000in}{0.083252in}}{\pgfqpoint{4.900000in}{4.900000in}}%
\pgfusepath{clip}%
\pgfsetbuttcap%
\pgfsetroundjoin%
\definecolor{currentfill}{rgb}{0.922430,0.919231,0.953618}%
\pgfsetfillcolor{currentfill}%
\pgfsetfillopacity{0.950000}%
\pgfsetlinewidth{0.200750pt}%
\definecolor{currentstroke}{rgb}{0.200000,0.200000,0.200000}%
\pgfsetstrokecolor{currentstroke}%
\pgfsetdash{}{0pt}%
\pgfpathmoveto{\pgfqpoint{3.679521in}{1.487191in}}%
\pgfpathlineto{\pgfqpoint{3.650359in}{1.532318in}}%
\pgfpathlineto{\pgfqpoint{3.621333in}{1.580297in}}%
\pgfpathlineto{\pgfqpoint{3.668617in}{1.597032in}}%
\pgfpathlineto{\pgfqpoint{3.715819in}{1.613718in}}%
\pgfpathlineto{\pgfqpoint{3.744892in}{1.565668in}}%
\pgfpathlineto{\pgfqpoint{3.774105in}{1.520399in}}%
\pgfpathlineto{\pgfqpoint{3.726855in}{1.503816in}}%
\pgfpathlineto{\pgfqpoint{3.679521in}{1.487191in}}%
\pgfpathclose%
\pgfusepath{stroke,fill}%
\end{pgfscope}%
\begin{pgfscope}%
\pgfpathrectangle{\pgfqpoint{0.050000in}{0.083252in}}{\pgfqpoint{4.900000in}{4.900000in}}%
\pgfusepath{clip}%
\pgfsetbuttcap%
\pgfsetroundjoin%
\definecolor{currentfill}{rgb}{0.731488,0.720415,0.839446}%
\pgfsetfillcolor{currentfill}%
\pgfsetfillopacity{0.950000}%
\pgfsetlinewidth{0.200750pt}%
\definecolor{currentstroke}{rgb}{0.200000,0.200000,0.200000}%
\pgfsetstrokecolor{currentstroke}%
\pgfsetdash{}{0pt}%
\pgfpathmoveto{\pgfqpoint{3.069184in}{1.793862in}}%
\pgfpathlineto{\pgfqpoint{3.040969in}{1.860746in}}%
\pgfpathlineto{\pgfqpoint{3.012809in}{1.929333in}}%
\pgfpathlineto{\pgfqpoint{3.060642in}{1.945460in}}%
\pgfpathlineto{\pgfqpoint{3.108391in}{1.961559in}}%
\pgfpathlineto{\pgfqpoint{3.136565in}{1.893386in}}%
\pgfpathlineto{\pgfqpoint{3.164797in}{1.826905in}}%
\pgfpathlineto{\pgfqpoint{3.117032in}{1.810403in}}%
\pgfpathlineto{\pgfqpoint{3.069184in}{1.793862in}}%
\pgfpathclose%
\pgfusepath{stroke,fill}%
\end{pgfscope}%
\begin{pgfscope}%
\pgfpathrectangle{\pgfqpoint{0.050000in}{0.083252in}}{\pgfqpoint{4.900000in}{4.900000in}}%
\pgfusepath{clip}%
\pgfsetbuttcap%
\pgfsetroundjoin%
\definecolor{currentfill}{rgb}{0.355571,0.328997,0.614671}%
\pgfsetfillcolor{currentfill}%
\pgfsetfillopacity{0.950000}%
\pgfsetlinewidth{0.200750pt}%
\definecolor{currentstroke}{rgb}{0.200000,0.200000,0.200000}%
\pgfsetstrokecolor{currentstroke}%
\pgfsetdash{}{0pt}%
\pgfpathmoveto{\pgfqpoint{2.500207in}{2.420689in}}%
\pgfpathlineto{\pgfqpoint{2.472353in}{2.495440in}}%
\pgfpathlineto{\pgfqpoint{2.444525in}{2.570193in}}%
\pgfpathlineto{\pgfqpoint{2.492867in}{2.584914in}}%
\pgfpathlineto{\pgfqpoint{2.541123in}{2.599614in}}%
\pgfpathlineto{\pgfqpoint{2.568946in}{2.525049in}}%
\pgfpathlineto{\pgfqpoint{2.596797in}{2.450501in}}%
\pgfpathlineto{\pgfqpoint{2.548545in}{2.435604in}}%
\pgfpathlineto{\pgfqpoint{2.500207in}{2.420689in}}%
\pgfpathclose%
\pgfusepath{stroke,fill}%
\end{pgfscope}%
\begin{pgfscope}%
\pgfpathrectangle{\pgfqpoint{0.050000in}{0.083252in}}{\pgfqpoint{4.900000in}{4.900000in}}%
\pgfusepath{clip}%
\pgfsetbuttcap%
\pgfsetroundjoin%
\definecolor{currentfill}{rgb}{0.510711,0.490534,0.707436}%
\pgfsetfillcolor{currentfill}%
\pgfsetfillopacity{0.950000}%
\pgfsetlinewidth{0.200750pt}%
\definecolor{currentstroke}{rgb}{0.200000,0.200000,0.200000}%
\pgfsetstrokecolor{currentstroke}%
\pgfsetdash{}{0pt}%
\pgfpathmoveto{\pgfqpoint{2.708472in}{2.153606in}}%
\pgfpathlineto{\pgfqpoint{2.680511in}{2.227472in}}%
\pgfpathlineto{\pgfqpoint{2.652579in}{2.301643in}}%
\pgfpathlineto{\pgfqpoint{2.700742in}{2.316763in}}%
\pgfpathlineto{\pgfqpoint{2.748819in}{2.331867in}}%
\pgfpathlineto{\pgfqpoint{2.776749in}{2.257982in}}%
\pgfpathlineto{\pgfqpoint{2.804710in}{2.184436in}}%
\pgfpathlineto{\pgfqpoint{2.756634in}{2.169029in}}%
\pgfpathlineto{\pgfqpoint{2.708472in}{2.153606in}}%
\pgfpathclose%
\pgfusepath{stroke,fill}%
\end{pgfscope}%
\begin{pgfscope}%
\pgfpathrectangle{\pgfqpoint{0.050000in}{0.083252in}}{\pgfqpoint{4.900000in}{4.900000in}}%
\pgfusepath{clip}%
\pgfsetbuttcap%
\pgfsetroundjoin%
\definecolor{currentfill}{rgb}{0.264191,0.473203,0.897870}%
\pgfsetfillcolor{currentfill}%
\pgfsetfillopacity{0.950000}%
\pgfsetlinewidth{0.200750pt}%
\definecolor{currentstroke}{rgb}{0.200000,0.200000,0.200000}%
\pgfsetstrokecolor{currentstroke}%
\pgfsetdash{}{0pt}%
\pgfpathmoveto{\pgfqpoint{1.514362in}{3.621925in}}%
\pgfpathlineto{\pgfqpoint{1.486963in}{3.696334in}}%
\pgfpathlineto{\pgfqpoint{1.459590in}{3.770673in}}%
\pgfpathlineto{\pgfqpoint{1.508834in}{3.784290in}}%
\pgfpathlineto{\pgfqpoint{1.557990in}{3.797882in}}%
\pgfpathlineto{\pgfqpoint{1.585361in}{3.723689in}}%
\pgfpathlineto{\pgfqpoint{1.612757in}{3.649425in}}%
\pgfpathlineto{\pgfqpoint{1.563603in}{3.635688in}}%
\pgfpathlineto{\pgfqpoint{1.514362in}{3.621925in}}%
\pgfpathclose%
\pgfusepath{stroke,fill}%
\end{pgfscope}%
\begin{pgfscope}%
\pgfpathrectangle{\pgfqpoint{0.050000in}{0.083252in}}{\pgfqpoint{4.900000in}{4.900000in}}%
\pgfusepath{clip}%
\pgfsetbuttcap%
\pgfsetroundjoin%
\definecolor{currentfill}{rgb}{0.982099,0.981361,0.989296}%
\pgfsetfillcolor{currentfill}%
\pgfsetfillopacity{0.950000}%
\pgfsetlinewidth{0.200750pt}%
\definecolor{currentstroke}{rgb}{0.200000,0.200000,0.200000}%
\pgfsetstrokecolor{currentstroke}%
\pgfsetdash{}{0pt}%
\pgfpathmoveto{\pgfqpoint{4.235212in}{1.391332in}}%
\pgfpathlineto{\pgfqpoint{4.204963in}{1.425073in}}%
\pgfpathlineto{\pgfqpoint{4.174847in}{1.460057in}}%
\pgfpathlineto{\pgfqpoint{4.221627in}{1.475904in}}%
\pgfpathlineto{\pgfqpoint{4.268325in}{1.491721in}}%
\pgfpathlineto{\pgfqpoint{4.298500in}{1.456432in}}%
\pgfpathlineto{\pgfqpoint{4.328809in}{1.422365in}}%
\pgfpathlineto{\pgfqpoint{4.282053in}{1.406863in}}%
\pgfpathlineto{\pgfqpoint{4.235212in}{1.391332in}}%
\pgfpathclose%
\pgfusepath{stroke,fill}%
\end{pgfscope}%
\begin{pgfscope}%
\pgfpathrectangle{\pgfqpoint{0.050000in}{0.083252in}}{\pgfqpoint{4.900000in}{4.900000in}}%
\pgfusepath{clip}%
\pgfsetbuttcap%
\pgfsetroundjoin%
\definecolor{currentfill}{rgb}{0.659885,0.645859,0.796632}%
\pgfsetfillcolor{currentfill}%
\pgfsetfillopacity{0.950000}%
\pgfsetlinewidth{0.200750pt}%
\definecolor{currentstroke}{rgb}{0.200000,0.200000,0.200000}%
\pgfsetstrokecolor{currentstroke}%
\pgfsetdash{}{0pt}%
\pgfpathmoveto{\pgfqpoint{2.916892in}{1.896993in}}%
\pgfpathlineto{\pgfqpoint{2.888789in}{1.967418in}}%
\pgfpathlineto{\pgfqpoint{2.860727in}{2.038960in}}%
\pgfpathlineto{\pgfqpoint{2.908721in}{2.054720in}}%
\pgfpathlineto{\pgfqpoint{2.956629in}{2.070460in}}%
\pgfpathlineto{\pgfqpoint{2.984698in}{1.999329in}}%
\pgfpathlineto{\pgfqpoint{3.012809in}{1.929333in}}%
\pgfpathlineto{\pgfqpoint{2.964893in}{1.913177in}}%
\pgfpathlineto{\pgfqpoint{2.916892in}{1.896993in}}%
\pgfpathclose%
\pgfusepath{stroke,fill}%
\end{pgfscope}%
\begin{pgfscope}%
\pgfpathrectangle{\pgfqpoint{0.050000in}{0.083252in}}{\pgfqpoint{4.900000in}{4.900000in}}%
\pgfusepath{clip}%
\pgfsetbuttcap%
\pgfsetroundjoin%
\definecolor{currentfill}{rgb}{0.257547,0.402614,0.804029}%
\pgfsetfillcolor{currentfill}%
\pgfsetfillopacity{0.950000}%
\pgfsetlinewidth{0.200750pt}%
\definecolor{currentstroke}{rgb}{0.200000,0.200000,0.200000}%
\pgfsetstrokecolor{currentstroke}%
\pgfsetdash{}{0pt}%
\pgfpathmoveto{\pgfqpoint{1.722602in}{3.351668in}}%
\pgfpathlineto{\pgfqpoint{1.695102in}{3.426213in}}%
\pgfpathlineto{\pgfqpoint{1.667627in}{3.500688in}}%
\pgfpathlineto{\pgfqpoint{1.716690in}{3.514547in}}%
\pgfpathlineto{\pgfqpoint{1.765665in}{3.528381in}}%
\pgfpathlineto{\pgfqpoint{1.793137in}{3.454053in}}%
\pgfpathlineto{\pgfqpoint{1.820634in}{3.379655in}}%
\pgfpathlineto{\pgfqpoint{1.771662in}{3.365674in}}%
\pgfpathlineto{\pgfqpoint{1.722602in}{3.351668in}}%
\pgfpathclose%
\pgfusepath{stroke,fill}%
\end{pgfscope}%
\begin{pgfscope}%
\pgfpathrectangle{\pgfqpoint{0.050000in}{0.083252in}}{\pgfqpoint{4.900000in}{4.900000in}}%
\pgfusepath{clip}%
\pgfsetbuttcap%
\pgfsetroundjoin%
\definecolor{currentfill}{rgb}{0.952265,0.950296,0.971457}%
\pgfsetfillcolor{currentfill}%
\pgfsetfillopacity{0.950000}%
\pgfsetlinewidth{0.200750pt}%
\definecolor{currentstroke}{rgb}{0.200000,0.200000,0.200000}%
\pgfsetstrokecolor{currentstroke}%
\pgfsetdash{}{0pt}%
\pgfpathmoveto{\pgfqpoint{3.832965in}{1.437777in}}%
\pgfpathlineto{\pgfqpoint{3.803462in}{1.477821in}}%
\pgfpathlineto{\pgfqpoint{3.774105in}{1.520399in}}%
\pgfpathlineto{\pgfqpoint{3.821272in}{1.536940in}}%
\pgfpathlineto{\pgfqpoint{3.868356in}{1.553439in}}%
\pgfpathlineto{\pgfqpoint{3.897767in}{1.510672in}}%
\pgfpathlineto{\pgfqpoint{3.927327in}{1.470389in}}%
\pgfpathlineto{\pgfqpoint{3.880187in}{1.454100in}}%
\pgfpathlineto{\pgfqpoint{3.832965in}{1.437777in}}%
\pgfpathclose%
\pgfusepath{stroke,fill}%
\end{pgfscope}%
\begin{pgfscope}%
\pgfpathrectangle{\pgfqpoint{0.050000in}{0.083252in}}{\pgfqpoint{4.900000in}{4.900000in}}%
\pgfusepath{clip}%
\pgfsetbuttcap%
\pgfsetroundjoin%
\definecolor{currentfill}{rgb}{0.269850,0.533333,0.977809}%
\pgfsetfillcolor{currentfill}%
\pgfsetfillopacity{0.950000}%
\pgfsetlinewidth{0.200750pt}%
\definecolor{currentstroke}{rgb}{0.200000,0.200000,0.200000}%
\pgfsetstrokecolor{currentstroke}%
\pgfsetdash{}{0pt}%
\pgfpathmoveto{\pgfqpoint{1.306165in}{3.892125in}}%
\pgfpathlineto{\pgfqpoint{1.278868in}{3.966399in}}%
\pgfpathlineto{\pgfqpoint{1.328292in}{3.979845in}}%
\pgfpathlineto{\pgfqpoint{1.377627in}{3.993268in}}%
\pgfpathlineto{\pgfqpoint{1.404922in}{3.919140in}}%
\pgfpathlineto{\pgfqpoint{1.355588in}{3.905645in}}%
\pgfpathlineto{\pgfqpoint{1.306165in}{3.892125in}}%
\pgfpathclose%
\pgfusepath{stroke,fill}%
\end{pgfscope}%
\begin{pgfscope}%
\pgfpathrectangle{\pgfqpoint{0.050000in}{0.083252in}}{\pgfqpoint{4.900000in}{4.900000in}}%
\pgfusepath{clip}%
\pgfsetbuttcap%
\pgfsetroundjoin%
\definecolor{currentfill}{rgb}{0.250903,0.332026,0.710188}%
\pgfsetfillcolor{currentfill}%
\pgfsetfillopacity{0.950000}%
\pgfsetlinewidth{0.200750pt}%
\definecolor{currentstroke}{rgb}{0.200000,0.200000,0.200000}%
\pgfsetstrokecolor{currentstroke}%
\pgfsetdash{}{0pt}%
\pgfpathmoveto{\pgfqpoint{1.930885in}{3.081355in}}%
\pgfpathlineto{\pgfqpoint{1.903283in}{3.156036in}}%
\pgfpathlineto{\pgfqpoint{1.875708in}{3.230646in}}%
\pgfpathlineto{\pgfqpoint{1.924589in}{3.244748in}}%
\pgfpathlineto{\pgfqpoint{1.973384in}{3.258825in}}%
\pgfpathlineto{\pgfqpoint{2.000956in}{3.184361in}}%
\pgfpathlineto{\pgfqpoint{2.028555in}{3.109828in}}%
\pgfpathlineto{\pgfqpoint{1.979764in}{3.095604in}}%
\pgfpathlineto{\pgfqpoint{1.930885in}{3.081355in}}%
\pgfpathclose%
\pgfusepath{stroke,fill}%
\end{pgfscope}%
\begin{pgfscope}%
\pgfpathrectangle{\pgfqpoint{0.050000in}{0.083252in}}{\pgfqpoint{4.900000in}{4.900000in}}%
\pgfusepath{clip}%
\pgfsetbuttcap%
\pgfsetroundjoin%
\definecolor{currentfill}{rgb}{0.244260,0.261438,0.616348}%
\pgfsetfillcolor{currentfill}%
\pgfsetfillopacity{0.950000}%
\pgfsetlinewidth{0.200750pt}%
\definecolor{currentstroke}{rgb}{0.200000,0.200000,0.200000}%
\pgfsetstrokecolor{currentstroke}%
\pgfsetdash{}{0pt}%
\pgfpathmoveto{\pgfqpoint{2.139212in}{2.810991in}}%
\pgfpathlineto{\pgfqpoint{2.111508in}{2.885805in}}%
\pgfpathlineto{\pgfqpoint{2.083831in}{2.960549in}}%
\pgfpathlineto{\pgfqpoint{2.132531in}{2.974895in}}%
\pgfpathlineto{\pgfqpoint{2.181145in}{2.989215in}}%
\pgfpathlineto{\pgfqpoint{2.208819in}{2.914620in}}%
\pgfpathlineto{\pgfqpoint{2.236519in}{2.839956in}}%
\pgfpathlineto{\pgfqpoint{2.187909in}{2.825486in}}%
\pgfpathlineto{\pgfqpoint{2.139212in}{2.810991in}}%
\pgfpathclose%
\pgfusepath{stroke,fill}%
\end{pgfscope}%
\begin{pgfscope}%
\pgfpathrectangle{\pgfqpoint{0.050000in}{0.083252in}}{\pgfqpoint{4.900000in}{4.900000in}}%
\pgfusepath{clip}%
\pgfsetbuttcap%
\pgfsetroundjoin%
\definecolor{currentfill}{rgb}{0.272034,0.242015,0.564721}%
\pgfsetfillcolor{currentfill}%
\pgfsetfillopacity{0.950000}%
\pgfsetlinewidth{0.200750pt}%
\definecolor{currentstroke}{rgb}{0.200000,0.200000,0.200000}%
\pgfsetstrokecolor{currentstroke}%
\pgfsetdash{}{0pt}%
\pgfpathmoveto{\pgfqpoint{2.347581in}{2.540683in}}%
\pgfpathlineto{\pgfqpoint{2.319776in}{2.615583in}}%
\pgfpathlineto{\pgfqpoint{2.291997in}{2.690433in}}%
\pgfpathlineto{\pgfqpoint{2.340516in}{2.705033in}}%
\pgfpathlineto{\pgfqpoint{2.388949in}{2.719609in}}%
\pgfpathlineto{\pgfqpoint{2.416724in}{2.644921in}}%
\pgfpathlineto{\pgfqpoint{2.444525in}{2.570193in}}%
\pgfpathlineto{\pgfqpoint{2.396097in}{2.555449in}}%
\pgfpathlineto{\pgfqpoint{2.347581in}{2.540683in}}%
\pgfpathclose%
\pgfusepath{stroke,fill}%
\end{pgfscope}%
\begin{pgfscope}%
\pgfpathrectangle{\pgfqpoint{0.050000in}{0.083252in}}{\pgfqpoint{4.900000in}{4.900000in}}%
\pgfusepath{clip}%
\pgfsetbuttcap%
\pgfsetroundjoin%
\definecolor{currentfill}{rgb}{0.433141,0.409765,0.661053}%
\pgfsetfillcolor{currentfill}%
\pgfsetfillopacity{0.950000}%
\pgfsetlinewidth{0.200750pt}%
\definecolor{currentstroke}{rgb}{0.200000,0.200000,0.200000}%
\pgfsetstrokecolor{currentstroke}%
\pgfsetdash{}{0pt}%
\pgfpathmoveto{\pgfqpoint{2.555995in}{2.271355in}}%
\pgfpathlineto{\pgfqpoint{2.528087in}{2.345976in}}%
\pgfpathlineto{\pgfqpoint{2.500207in}{2.420689in}}%
\pgfpathlineto{\pgfqpoint{2.548545in}{2.435604in}}%
\pgfpathlineto{\pgfqpoint{2.596797in}{2.450501in}}%
\pgfpathlineto{\pgfqpoint{2.624674in}{2.376013in}}%
\pgfpathlineto{\pgfqpoint{2.652579in}{2.301643in}}%
\pgfpathlineto{\pgfqpoint{2.604330in}{2.286507in}}%
\pgfpathlineto{\pgfqpoint{2.555995in}{2.271355in}}%
\pgfpathclose%
\pgfusepath{stroke,fill}%
\end{pgfscope}%
\begin{pgfscope}%
\pgfpathrectangle{\pgfqpoint{0.050000in}{0.083252in}}{\pgfqpoint{4.900000in}{4.900000in}}%
\pgfusepath{clip}%
\pgfsetbuttcap%
\pgfsetroundjoin%
\definecolor{currentfill}{rgb}{0.588281,0.571303,0.753818}%
\pgfsetfillcolor{currentfill}%
\pgfsetfillopacity{0.950000}%
\pgfsetlinewidth{0.200750pt}%
\definecolor{currentstroke}{rgb}{0.200000,0.200000,0.200000}%
\pgfsetstrokecolor{currentstroke}%
\pgfsetdash{}{0pt}%
\pgfpathmoveto{\pgfqpoint{2.764486in}{2.007381in}}%
\pgfpathlineto{\pgfqpoint{2.736463in}{2.080184in}}%
\pgfpathlineto{\pgfqpoint{2.708472in}{2.153606in}}%
\pgfpathlineto{\pgfqpoint{2.756634in}{2.169029in}}%
\pgfpathlineto{\pgfqpoint{2.804710in}{2.184436in}}%
\pgfpathlineto{\pgfqpoint{2.832702in}{2.111370in}}%
\pgfpathlineto{\pgfqpoint{2.860727in}{2.038960in}}%
\pgfpathlineto{\pgfqpoint{2.812649in}{2.023180in}}%
\pgfpathlineto{\pgfqpoint{2.764486in}{2.007381in}}%
\pgfpathclose%
\pgfusepath{stroke,fill}%
\end{pgfscope}%
\begin{pgfscope}%
\pgfpathrectangle{\pgfqpoint{0.050000in}{0.083252in}}{\pgfqpoint{4.900000in}{4.900000in}}%
\pgfusepath{clip}%
\pgfsetbuttcap%
\pgfsetroundjoin%
\definecolor{currentfill}{rgb}{0.850827,0.844675,0.910804}%
\pgfsetfillcolor{currentfill}%
\pgfsetfillopacity{0.950000}%
\pgfsetlinewidth{0.200750pt}%
\definecolor{currentstroke}{rgb}{0.200000,0.200000,0.200000}%
\pgfsetstrokecolor{currentstroke}%
\pgfsetdash{}{0pt}%
\pgfpathmoveto{\pgfqpoint{3.278476in}{1.583611in}}%
\pgfpathlineto{\pgfqpoint{3.249926in}{1.640505in}}%
\pgfpathlineto{\pgfqpoint{3.221469in}{1.700183in}}%
\pgfpathlineto{\pgfqpoint{3.269174in}{1.716979in}}%
\pgfpathlineto{\pgfqpoint{3.316797in}{1.733726in}}%
\pgfpathlineto{\pgfqpoint{3.345285in}{1.674291in}}%
\pgfpathlineto{\pgfqpoint{3.373871in}{1.617563in}}%
\pgfpathlineto{\pgfqpoint{3.326215in}{1.600614in}}%
\pgfpathlineto{\pgfqpoint{3.278476in}{1.583611in}}%
\pgfpathclose%
\pgfusepath{stroke,fill}%
\end{pgfscope}%
\begin{pgfscope}%
\pgfpathrectangle{\pgfqpoint{0.050000in}{0.083252in}}{\pgfqpoint{4.900000in}{4.900000in}}%
\pgfusepath{clip}%
\pgfsetbuttcap%
\pgfsetroundjoin%
\definecolor{currentfill}{rgb}{0.267389,0.507190,0.943053}%
\pgfsetfillcolor{currentfill}%
\pgfsetfillopacity{0.950000}%
\pgfsetlinewidth{0.200750pt}%
\definecolor{currentstroke}{rgb}{0.200000,0.200000,0.200000}%
\pgfsetstrokecolor{currentstroke}%
\pgfsetdash{}{0pt}%
\pgfpathmoveto{\pgfqpoint{1.360837in}{3.743367in}}%
\pgfpathlineto{\pgfqpoint{1.333488in}{3.817781in}}%
\pgfpathlineto{\pgfqpoint{1.306165in}{3.892125in}}%
\pgfpathlineto{\pgfqpoint{1.355588in}{3.905645in}}%
\pgfpathlineto{\pgfqpoint{1.404922in}{3.919140in}}%
\pgfpathlineto{\pgfqpoint{1.432243in}{3.844942in}}%
\pgfpathlineto{\pgfqpoint{1.459590in}{3.770673in}}%
\pgfpathlineto{\pgfqpoint{1.410258in}{3.757032in}}%
\pgfpathlineto{\pgfqpoint{1.360837in}{3.743367in}}%
\pgfpathclose%
\pgfusepath{stroke,fill}%
\end{pgfscope}%
\begin{pgfscope}%
\pgfpathrectangle{\pgfqpoint{0.050000in}{0.083252in}}{\pgfqpoint{4.900000in}{4.900000in}}%
\pgfusepath{clip}%
\pgfsetbuttcap%
\pgfsetroundjoin%
\definecolor{currentfill}{rgb}{0.892595,0.888166,0.935779}%
\pgfsetfillcolor{currentfill}%
\pgfsetfillopacity{0.950000}%
\pgfsetlinewidth{0.200750pt}%
\definecolor{currentstroke}{rgb}{0.200000,0.200000,0.200000}%
\pgfsetstrokecolor{currentstroke}%
\pgfsetdash{}{0pt}%
\pgfpathmoveto{\pgfqpoint{3.431369in}{1.512872in}}%
\pgfpathlineto{\pgfqpoint{3.402562in}{1.563719in}}%
\pgfpathlineto{\pgfqpoint{3.373871in}{1.617563in}}%
\pgfpathlineto{\pgfqpoint{3.421444in}{1.634459in}}%
\pgfpathlineto{\pgfqpoint{3.468934in}{1.651304in}}%
\pgfpathlineto{\pgfqpoint{3.497666in}{1.597538in}}%
\pgfpathlineto{\pgfqpoint{3.526517in}{1.546684in}}%
\pgfpathlineto{\pgfqpoint{3.478985in}{1.529803in}}%
\pgfpathlineto{\pgfqpoint{3.431369in}{1.512872in}}%
\pgfpathclose%
\pgfusepath{stroke,fill}%
\end{pgfscope}%
\begin{pgfscope}%
\pgfpathrectangle{\pgfqpoint{0.050000in}{0.083252in}}{\pgfqpoint{4.900000in}{4.900000in}}%
\pgfusepath{clip}%
\pgfsetbuttcap%
\pgfsetroundjoin%
\definecolor{currentfill}{rgb}{0.970165,0.968935,0.982161}%
\pgfsetfillcolor{currentfill}%
\pgfsetfillopacity{0.950000}%
\pgfsetlinewidth{0.200750pt}%
\definecolor{currentstroke}{rgb}{0.200000,0.200000,0.200000}%
\pgfsetstrokecolor{currentstroke}%
\pgfsetdash{}{0pt}%
\pgfpathmoveto{\pgfqpoint{3.986887in}{1.396358in}}%
\pgfpathlineto{\pgfqpoint{3.957034in}{1.432376in}}%
\pgfpathlineto{\pgfqpoint{3.927327in}{1.470389in}}%
\pgfpathlineto{\pgfqpoint{3.974382in}{1.486643in}}%
\pgfpathlineto{\pgfqpoint{4.021355in}{1.502861in}}%
\pgfpathlineto{\pgfqpoint{4.051121in}{1.464582in}}%
\pgfpathlineto{\pgfqpoint{4.081035in}{1.428269in}}%
\pgfpathlineto{\pgfqpoint{4.034003in}{1.412329in}}%
\pgfpathlineto{\pgfqpoint{3.986887in}{1.396358in}}%
\pgfpathclose%
\pgfusepath{stroke,fill}%
\end{pgfscope}%
\begin{pgfscope}%
\pgfpathrectangle{\pgfqpoint{0.050000in}{0.083252in}}{\pgfqpoint{4.900000in}{4.900000in}}%
\pgfusepath{clip}%
\pgfsetbuttcap%
\pgfsetroundjoin%
\definecolor{currentfill}{rgb}{0.797124,0.788758,0.878693}%
\pgfsetfillcolor{currentfill}%
\pgfsetfillopacity{0.950000}%
\pgfsetlinewidth{0.200750pt}%
\definecolor{currentstroke}{rgb}{0.200000,0.200000,0.200000}%
\pgfsetstrokecolor{currentstroke}%
\pgfsetdash{}{0pt}%
\pgfpathmoveto{\pgfqpoint{3.125809in}{1.666443in}}%
\pgfpathlineto{\pgfqpoint{3.097462in}{1.728992in}}%
\pgfpathlineto{\pgfqpoint{3.069184in}{1.793862in}}%
\pgfpathlineto{\pgfqpoint{3.117032in}{1.810403in}}%
\pgfpathlineto{\pgfqpoint{3.164797in}{1.826905in}}%
\pgfpathlineto{\pgfqpoint{3.193096in}{1.762407in}}%
\pgfpathlineto{\pgfqpoint{3.221469in}{1.700183in}}%
\pgfpathlineto{\pgfqpoint{3.173680in}{1.683338in}}%
\pgfpathlineto{\pgfqpoint{3.125809in}{1.666443in}}%
\pgfpathclose%
\pgfusepath{stroke,fill}%
\end{pgfscope}%
\begin{pgfscope}%
\pgfpathrectangle{\pgfqpoint{0.050000in}{0.083252in}}{\pgfqpoint{4.900000in}{4.900000in}}%
\pgfusepath{clip}%
\pgfsetbuttcap%
\pgfsetroundjoin%
\definecolor{currentfill}{rgb}{0.260746,0.436601,0.849212}%
\pgfsetfillcolor{currentfill}%
\pgfsetfillopacity{0.950000}%
\pgfsetlinewidth{0.200750pt}%
\definecolor{currentstroke}{rgb}{0.200000,0.200000,0.200000}%
\pgfsetstrokecolor{currentstroke}%
\pgfsetdash{}{0pt}%
\pgfpathmoveto{\pgfqpoint{1.569238in}{3.472894in}}%
\pgfpathlineto{\pgfqpoint{1.541787in}{3.547445in}}%
\pgfpathlineto{\pgfqpoint{1.514362in}{3.621925in}}%
\pgfpathlineto{\pgfqpoint{1.563603in}{3.635688in}}%
\pgfpathlineto{\pgfqpoint{1.612757in}{3.649425in}}%
\pgfpathlineto{\pgfqpoint{1.640179in}{3.575092in}}%
\pgfpathlineto{\pgfqpoint{1.667627in}{3.500688in}}%
\pgfpathlineto{\pgfqpoint{1.618477in}{3.486804in}}%
\pgfpathlineto{\pgfqpoint{1.569238in}{3.472894in}}%
\pgfpathclose%
\pgfusepath{stroke,fill}%
\end{pgfscope}%
\begin{pgfscope}%
\pgfpathrectangle{\pgfqpoint{0.050000in}{0.083252in}}{\pgfqpoint{4.900000in}{4.900000in}}%
\pgfusepath{clip}%
\pgfsetbuttcap%
\pgfsetroundjoin%
\definecolor{currentfill}{rgb}{0.254102,0.366013,0.755371}%
\pgfsetfillcolor{currentfill}%
\pgfsetfillopacity{0.950000}%
\pgfsetlinewidth{0.200750pt}%
\definecolor{currentstroke}{rgb}{0.200000,0.200000,0.200000}%
\pgfsetstrokecolor{currentstroke}%
\pgfsetdash{}{0pt}%
\pgfpathmoveto{\pgfqpoint{1.777681in}{3.202365in}}%
\pgfpathlineto{\pgfqpoint{1.750129in}{3.277052in}}%
\pgfpathlineto{\pgfqpoint{1.722602in}{3.351668in}}%
\pgfpathlineto{\pgfqpoint{1.771662in}{3.365674in}}%
\pgfpathlineto{\pgfqpoint{1.820634in}{3.379655in}}%
\pgfpathlineto{\pgfqpoint{1.848158in}{3.305186in}}%
\pgfpathlineto{\pgfqpoint{1.875708in}{3.230646in}}%
\pgfpathlineto{\pgfqpoint{1.826738in}{3.216518in}}%
\pgfpathlineto{\pgfqpoint{1.777681in}{3.202365in}}%
\pgfpathclose%
\pgfusepath{stroke,fill}%
\end{pgfscope}%
\begin{pgfscope}%
\pgfpathrectangle{\pgfqpoint{0.050000in}{0.083252in}}{\pgfqpoint{4.900000in}{4.900000in}}%
\pgfusepath{clip}%
\pgfsetbuttcap%
\pgfsetroundjoin%
\definecolor{currentfill}{rgb}{0.247705,0.298039,0.665006}%
\pgfsetfillcolor{currentfill}%
\pgfsetfillopacity{0.950000}%
\pgfsetlinewidth{0.200750pt}%
\definecolor{currentstroke}{rgb}{0.200000,0.200000,0.200000}%
\pgfsetstrokecolor{currentstroke}%
\pgfsetdash{}{0pt}%
\pgfpathmoveto{\pgfqpoint{1.986168in}{2.931780in}}%
\pgfpathlineto{\pgfqpoint{1.958514in}{3.006603in}}%
\pgfpathlineto{\pgfqpoint{1.930885in}{3.081355in}}%
\pgfpathlineto{\pgfqpoint{1.979764in}{3.095604in}}%
\pgfpathlineto{\pgfqpoint{2.028555in}{3.109828in}}%
\pgfpathlineto{\pgfqpoint{2.056180in}{3.035223in}}%
\pgfpathlineto{\pgfqpoint{2.083831in}{2.960549in}}%
\pgfpathlineto{\pgfqpoint{2.035043in}{2.946178in}}%
\pgfpathlineto{\pgfqpoint{1.986168in}{2.931780in}}%
\pgfpathclose%
\pgfusepath{stroke,fill}%
\end{pgfscope}%
\begin{pgfscope}%
\pgfpathrectangle{\pgfqpoint{0.050000in}{0.083252in}}{\pgfqpoint{4.900000in}{4.900000in}}%
\pgfusepath{clip}%
\pgfsetbuttcap%
\pgfsetroundjoin%
\definecolor{currentfill}{rgb}{0.928397,0.925444,0.957186}%
\pgfsetfillcolor{currentfill}%
\pgfsetfillopacity{0.950000}%
\pgfsetlinewidth{0.200750pt}%
\definecolor{currentstroke}{rgb}{0.200000,0.200000,0.200000}%
\pgfsetstrokecolor{currentstroke}%
\pgfsetdash{}{0pt}%
\pgfpathmoveto{\pgfqpoint{3.584605in}{1.453813in}}%
\pgfpathlineto{\pgfqpoint{3.555495in}{1.498786in}}%
\pgfpathlineto{\pgfqpoint{3.526517in}{1.546684in}}%
\pgfpathlineto{\pgfqpoint{3.573967in}{1.563515in}}%
\pgfpathlineto{\pgfqpoint{3.621333in}{1.580297in}}%
\pgfpathlineto{\pgfqpoint{3.650359in}{1.532318in}}%
\pgfpathlineto{\pgfqpoint{3.679521in}{1.487191in}}%
\pgfpathlineto{\pgfqpoint{3.632105in}{1.470523in}}%
\pgfpathlineto{\pgfqpoint{3.584605in}{1.453813in}}%
\pgfpathclose%
\pgfusepath{stroke,fill}%
\end{pgfscope}%
\begin{pgfscope}%
\pgfpathrectangle{\pgfqpoint{0.050000in}{0.083252in}}{\pgfqpoint{4.900000in}{4.900000in}}%
\pgfusepath{clip}%
\pgfsetbuttcap%
\pgfsetroundjoin%
\definecolor{currentfill}{rgb}{0.241061,0.227451,0.571165}%
\pgfsetfillcolor{currentfill}%
\pgfsetfillopacity{0.950000}%
\pgfsetlinewidth{0.200750pt}%
\definecolor{currentstroke}{rgb}{0.200000,0.200000,0.200000}%
\pgfsetstrokecolor{currentstroke}%
\pgfsetdash{}{0pt}%
\pgfpathmoveto{\pgfqpoint{2.194698in}{2.661161in}}%
\pgfpathlineto{\pgfqpoint{2.166942in}{2.736109in}}%
\pgfpathlineto{\pgfqpoint{2.139212in}{2.810991in}}%
\pgfpathlineto{\pgfqpoint{2.187909in}{2.825486in}}%
\pgfpathlineto{\pgfqpoint{2.236519in}{2.839956in}}%
\pgfpathlineto{\pgfqpoint{2.264245in}{2.765226in}}%
\pgfpathlineto{\pgfqpoint{2.291997in}{2.690433in}}%
\pgfpathlineto{\pgfqpoint{2.243391in}{2.675809in}}%
\pgfpathlineto{\pgfqpoint{2.194698in}{2.661161in}}%
\pgfpathclose%
\pgfusepath{stroke,fill}%
\end{pgfscope}%
\begin{pgfscope}%
\pgfpathrectangle{\pgfqpoint{0.050000in}{0.083252in}}{\pgfqpoint{4.900000in}{4.900000in}}%
\pgfusepath{clip}%
\pgfsetbuttcap%
\pgfsetroundjoin%
\definecolor{currentfill}{rgb}{0.731488,0.720415,0.839446}%
\pgfsetfillcolor{currentfill}%
\pgfsetfillopacity{0.950000}%
\pgfsetlinewidth{0.200750pt}%
\definecolor{currentstroke}{rgb}{0.200000,0.200000,0.200000}%
\pgfsetstrokecolor{currentstroke}%
\pgfsetdash{}{0pt}%
\pgfpathmoveto{\pgfqpoint{2.973238in}{1.760661in}}%
\pgfpathlineto{\pgfqpoint{2.945039in}{1.827970in}}%
\pgfpathlineto{\pgfqpoint{2.916892in}{1.896993in}}%
\pgfpathlineto{\pgfqpoint{2.964893in}{1.913177in}}%
\pgfpathlineto{\pgfqpoint{3.012809in}{1.929333in}}%
\pgfpathlineto{\pgfqpoint{3.040969in}{1.860746in}}%
\pgfpathlineto{\pgfqpoint{3.069184in}{1.793862in}}%
\pgfpathlineto{\pgfqpoint{3.021253in}{1.777281in}}%
\pgfpathlineto{\pgfqpoint{2.973238in}{1.760661in}}%
\pgfpathclose%
\pgfusepath{stroke,fill}%
\end{pgfscope}%
\begin{pgfscope}%
\pgfpathrectangle{\pgfqpoint{0.050000in}{0.083252in}}{\pgfqpoint{4.900000in}{4.900000in}}%
\pgfusepath{clip}%
\pgfsetbuttcap%
\pgfsetroundjoin%
\definecolor{currentfill}{rgb}{0.355571,0.328997,0.614671}%
\pgfsetfillcolor{currentfill}%
\pgfsetfillopacity{0.950000}%
\pgfsetlinewidth{0.200750pt}%
\definecolor{currentstroke}{rgb}{0.200000,0.200000,0.200000}%
\pgfsetstrokecolor{currentstroke}%
\pgfsetdash{}{0pt}%
\pgfpathmoveto{\pgfqpoint{2.403270in}{2.390799in}}%
\pgfpathlineto{\pgfqpoint{2.375413in}{2.465748in}}%
\pgfpathlineto{\pgfqpoint{2.347581in}{2.540683in}}%
\pgfpathlineto{\pgfqpoint{2.396097in}{2.555449in}}%
\pgfpathlineto{\pgfqpoint{2.444525in}{2.570193in}}%
\pgfpathlineto{\pgfqpoint{2.472353in}{2.495440in}}%
\pgfpathlineto{\pgfqpoint{2.500207in}{2.420689in}}%
\pgfpathlineto{\pgfqpoint{2.451782in}{2.405753in}}%
\pgfpathlineto{\pgfqpoint{2.403270in}{2.390799in}}%
\pgfpathclose%
\pgfusepath{stroke,fill}%
\end{pgfscope}%
\begin{pgfscope}%
\pgfpathrectangle{\pgfqpoint{0.050000in}{0.083252in}}{\pgfqpoint{4.900000in}{4.900000in}}%
\pgfusepath{clip}%
\pgfsetbuttcap%
\pgfsetroundjoin%
\definecolor{currentfill}{rgb}{0.994033,0.993787,0.996432}%
\pgfsetfillcolor{currentfill}%
\pgfsetfillopacity{0.950000}%
\pgfsetlinewidth{0.200750pt}%
\definecolor{currentstroke}{rgb}{0.200000,0.200000,0.200000}%
\pgfsetstrokecolor{currentstroke}%
\pgfsetdash{}{0pt}%
\pgfpathmoveto{\pgfqpoint{4.389771in}{1.356107in}}%
\pgfpathlineto{\pgfqpoint{4.359239in}{1.389088in}}%
\pgfpathlineto{\pgfqpoint{4.328809in}{1.422365in}}%
\pgfpathlineto{\pgfqpoint{4.375481in}{1.437839in}}%
\pgfpathlineto{\pgfqpoint{4.405941in}{1.404391in}}%
\pgfpathlineto{\pgfqpoint{4.436503in}{1.371232in}}%
\pgfpathlineto{\pgfqpoint{4.389771in}{1.356107in}}%
\pgfpathclose%
\pgfusepath{stroke,fill}%
\end{pgfscope}%
\begin{pgfscope}%
\pgfpathrectangle{\pgfqpoint{0.050000in}{0.083252in}}{\pgfqpoint{4.900000in}{4.900000in}}%
\pgfusepath{clip}%
\pgfsetbuttcap%
\pgfsetroundjoin%
\definecolor{currentfill}{rgb}{0.510711,0.490534,0.707436}%
\pgfsetfillcolor{currentfill}%
\pgfsetfillopacity{0.950000}%
\pgfsetlinewidth{0.200750pt}%
\definecolor{currentstroke}{rgb}{0.200000,0.200000,0.200000}%
\pgfsetstrokecolor{currentstroke}%
\pgfsetdash{}{0pt}%
\pgfpathmoveto{\pgfqpoint{2.611892in}{2.122715in}}%
\pgfpathlineto{\pgfqpoint{2.583929in}{2.196900in}}%
\pgfpathlineto{\pgfqpoint{2.555995in}{2.271355in}}%
\pgfpathlineto{\pgfqpoint{2.604330in}{2.286507in}}%
\pgfpathlineto{\pgfqpoint{2.652579in}{2.301643in}}%
\pgfpathlineto{\pgfqpoint{2.680511in}{2.227472in}}%
\pgfpathlineto{\pgfqpoint{2.708472in}{2.153606in}}%
\pgfpathlineto{\pgfqpoint{2.660225in}{2.138168in}}%
\pgfpathlineto{\pgfqpoint{2.611892in}{2.122715in}}%
\pgfpathclose%
\pgfusepath{stroke,fill}%
\end{pgfscope}%
\begin{pgfscope}%
\pgfpathrectangle{\pgfqpoint{0.050000in}{0.083252in}}{\pgfqpoint{4.900000in}{4.900000in}}%
\pgfusepath{clip}%
\pgfsetbuttcap%
\pgfsetroundjoin%
\definecolor{currentfill}{rgb}{0.263945,0.470588,0.894394}%
\pgfsetfillcolor{currentfill}%
\pgfsetfillopacity{0.950000}%
\pgfsetlinewidth{0.200750pt}%
\definecolor{currentstroke}{rgb}{0.200000,0.200000,0.200000}%
\pgfsetstrokecolor{currentstroke}%
\pgfsetdash{}{0pt}%
\pgfpathmoveto{\pgfqpoint{1.415614in}{3.594325in}}%
\pgfpathlineto{\pgfqpoint{1.388213in}{3.668881in}}%
\pgfpathlineto{\pgfqpoint{1.360837in}{3.743367in}}%
\pgfpathlineto{\pgfqpoint{1.410258in}{3.757032in}}%
\pgfpathlineto{\pgfqpoint{1.459590in}{3.770673in}}%
\pgfpathlineto{\pgfqpoint{1.486963in}{3.696334in}}%
\pgfpathlineto{\pgfqpoint{1.514362in}{3.621925in}}%
\pgfpathlineto{\pgfqpoint{1.465032in}{3.608138in}}%
\pgfpathlineto{\pgfqpoint{1.415614in}{3.594325in}}%
\pgfpathclose%
\pgfusepath{stroke,fill}%
\end{pgfscope}%
\begin{pgfscope}%
\pgfpathrectangle{\pgfqpoint{0.050000in}{0.083252in}}{\pgfqpoint{4.900000in}{4.900000in}}%
\pgfusepath{clip}%
\pgfsetbuttcap%
\pgfsetroundjoin%
\definecolor{currentfill}{rgb}{0.665852,0.652072,0.800200}%
\pgfsetfillcolor{currentfill}%
\pgfsetfillopacity{0.950000}%
\pgfsetlinewidth{0.200750pt}%
\definecolor{currentstroke}{rgb}{0.200000,0.200000,0.200000}%
\pgfsetstrokecolor{currentstroke}%
\pgfsetdash{}{0pt}%
\pgfpathmoveto{\pgfqpoint{2.820637in}{1.864542in}}%
\pgfpathlineto{\pgfqpoint{2.792543in}{1.935415in}}%
\pgfpathlineto{\pgfqpoint{2.764486in}{2.007381in}}%
\pgfpathlineto{\pgfqpoint{2.812649in}{2.023180in}}%
\pgfpathlineto{\pgfqpoint{2.860727in}{2.038960in}}%
\pgfpathlineto{\pgfqpoint{2.888789in}{1.967418in}}%
\pgfpathlineto{\pgfqpoint{2.916892in}{1.896993in}}%
\pgfpathlineto{\pgfqpoint{2.868807in}{1.880781in}}%
\pgfpathlineto{\pgfqpoint{2.820637in}{1.864542in}}%
\pgfpathclose%
\pgfusepath{stroke,fill}%
\end{pgfscope}%
\begin{pgfscope}%
\pgfpathrectangle{\pgfqpoint{0.050000in}{0.083252in}}{\pgfqpoint{4.900000in}{4.900000in}}%
\pgfusepath{clip}%
\pgfsetbuttcap%
\pgfsetroundjoin%
\definecolor{currentfill}{rgb}{0.988066,0.987574,0.992864}%
\pgfsetfillcolor{currentfill}%
\pgfsetfillopacity{0.950000}%
\pgfsetlinewidth{0.200750pt}%
\definecolor{currentstroke}{rgb}{0.200000,0.200000,0.200000}%
\pgfsetstrokecolor{currentstroke}%
\pgfsetdash{}{0pt}%
\pgfpathmoveto{\pgfqpoint{4.141276in}{1.360185in}}%
\pgfpathlineto{\pgfqpoint{4.111089in}{1.393596in}}%
\pgfpathlineto{\pgfqpoint{4.081035in}{1.428269in}}%
\pgfpathlineto{\pgfqpoint{4.127983in}{1.444178in}}%
\pgfpathlineto{\pgfqpoint{4.174847in}{1.460057in}}%
\pgfpathlineto{\pgfqpoint{4.204963in}{1.425073in}}%
\pgfpathlineto{\pgfqpoint{4.235212in}{1.391332in}}%
\pgfpathlineto{\pgfqpoint{4.188286in}{1.375772in}}%
\pgfpathlineto{\pgfqpoint{4.141276in}{1.360185in}}%
\pgfpathclose%
\pgfusepath{stroke,fill}%
\end{pgfscope}%
\begin{pgfscope}%
\pgfpathrectangle{\pgfqpoint{0.050000in}{0.083252in}}{\pgfqpoint{4.900000in}{4.900000in}}%
\pgfusepath{clip}%
\pgfsetbuttcap%
\pgfsetroundjoin%
\definecolor{currentfill}{rgb}{0.952265,0.950296,0.971457}%
\pgfsetfillcolor{currentfill}%
\pgfsetfillopacity{0.950000}%
\pgfsetlinewidth{0.200750pt}%
\definecolor{currentstroke}{rgb}{0.200000,0.200000,0.200000}%
\pgfsetstrokecolor{currentstroke}%
\pgfsetdash{}{0pt}%
\pgfpathmoveto{\pgfqpoint{3.738268in}{1.405024in}}%
\pgfpathlineto{\pgfqpoint{3.708824in}{1.444816in}}%
\pgfpathlineto{\pgfqpoint{3.679521in}{1.487191in}}%
\pgfpathlineto{\pgfqpoint{3.726855in}{1.503816in}}%
\pgfpathlineto{\pgfqpoint{3.774105in}{1.520399in}}%
\pgfpathlineto{\pgfqpoint{3.803462in}{1.477821in}}%
\pgfpathlineto{\pgfqpoint{3.832965in}{1.437777in}}%
\pgfpathlineto{\pgfqpoint{3.785658in}{1.421418in}}%
\pgfpathlineto{\pgfqpoint{3.738268in}{1.405024in}}%
\pgfpathclose%
\pgfusepath{stroke,fill}%
\end{pgfscope}%
\begin{pgfscope}%
\pgfpathrectangle{\pgfqpoint{0.050000in}{0.083252in}}{\pgfqpoint{4.900000in}{4.900000in}}%
\pgfusepath{clip}%
\pgfsetbuttcap%
\pgfsetroundjoin%
\definecolor{currentfill}{rgb}{0.257547,0.402614,0.804029}%
\pgfsetfillcolor{currentfill}%
\pgfsetfillopacity{0.950000}%
\pgfsetlinewidth{0.200750pt}%
\definecolor{currentstroke}{rgb}{0.200000,0.200000,0.200000}%
\pgfsetstrokecolor{currentstroke}%
\pgfsetdash{}{0pt}%
\pgfpathmoveto{\pgfqpoint{1.624218in}{3.323581in}}%
\pgfpathlineto{\pgfqpoint{1.596715in}{3.398273in}}%
\pgfpathlineto{\pgfqpoint{1.569238in}{3.472894in}}%
\pgfpathlineto{\pgfqpoint{1.618477in}{3.486804in}}%
\pgfpathlineto{\pgfqpoint{1.667627in}{3.500688in}}%
\pgfpathlineto{\pgfqpoint{1.695102in}{3.426213in}}%
\pgfpathlineto{\pgfqpoint{1.722602in}{3.351668in}}%
\pgfpathlineto{\pgfqpoint{1.673454in}{3.337637in}}%
\pgfpathlineto{\pgfqpoint{1.624218in}{3.323581in}}%
\pgfpathclose%
\pgfusepath{stroke,fill}%
\end{pgfscope}%
\begin{pgfscope}%
\pgfpathrectangle{\pgfqpoint{0.050000in}{0.083252in}}{\pgfqpoint{4.900000in}{4.900000in}}%
\pgfusepath{clip}%
\pgfsetbuttcap%
\pgfsetroundjoin%
\definecolor{currentfill}{rgb}{0.269850,0.533333,0.977809}%
\pgfsetfillcolor{currentfill}%
\pgfsetfillopacity{0.950000}%
\pgfsetlinewidth{0.200750pt}%
\definecolor{currentstroke}{rgb}{0.200000,0.200000,0.200000}%
\pgfsetstrokecolor{currentstroke}%
\pgfsetdash{}{0pt}%
\pgfpathmoveto{\pgfqpoint{1.207053in}{3.865013in}}%
\pgfpathlineto{\pgfqpoint{1.179754in}{3.939433in}}%
\pgfpathlineto{\pgfqpoint{1.229355in}{3.952928in}}%
\pgfpathlineto{\pgfqpoint{1.278868in}{3.966399in}}%
\pgfpathlineto{\pgfqpoint{1.306165in}{3.892125in}}%
\pgfpathlineto{\pgfqpoint{1.256654in}{3.878582in}}%
\pgfpathlineto{\pgfqpoint{1.207053in}{3.865013in}}%
\pgfpathclose%
\pgfusepath{stroke,fill}%
\end{pgfscope}%
\begin{pgfscope}%
\pgfpathrectangle{\pgfqpoint{0.050000in}{0.083252in}}{\pgfqpoint{4.900000in}{4.900000in}}%
\pgfusepath{clip}%
\pgfsetbuttcap%
\pgfsetroundjoin%
\definecolor{currentfill}{rgb}{0.250903,0.332026,0.710188}%
\pgfsetfillcolor{currentfill}%
\pgfsetfillopacity{0.950000}%
\pgfsetlinewidth{0.200750pt}%
\definecolor{currentstroke}{rgb}{0.200000,0.200000,0.200000}%
\pgfsetstrokecolor{currentstroke}%
\pgfsetdash{}{0pt}%
\pgfpathmoveto{\pgfqpoint{1.832865in}{3.052779in}}%
\pgfpathlineto{\pgfqpoint{1.805260in}{3.127608in}}%
\pgfpathlineto{\pgfqpoint{1.777681in}{3.202365in}}%
\pgfpathlineto{\pgfqpoint{1.826738in}{3.216518in}}%
\pgfpathlineto{\pgfqpoint{1.875708in}{3.230646in}}%
\pgfpathlineto{\pgfqpoint{1.903283in}{3.156036in}}%
\pgfpathlineto{\pgfqpoint{1.930885in}{3.081355in}}%
\pgfpathlineto{\pgfqpoint{1.881919in}{3.067080in}}%
\pgfpathlineto{\pgfqpoint{1.832865in}{3.052779in}}%
\pgfpathclose%
\pgfusepath{stroke,fill}%
\end{pgfscope}%
\begin{pgfscope}%
\pgfpathrectangle{\pgfqpoint{0.050000in}{0.083252in}}{\pgfqpoint{4.900000in}{4.900000in}}%
\pgfusepath{clip}%
\pgfsetbuttcap%
\pgfsetroundjoin%
\definecolor{currentfill}{rgb}{0.244260,0.261438,0.616348}%
\pgfsetfillcolor{currentfill}%
\pgfsetfillopacity{0.950000}%
\pgfsetlinewidth{0.200750pt}%
\definecolor{currentstroke}{rgb}{0.200000,0.200000,0.200000}%
\pgfsetstrokecolor{currentstroke}%
\pgfsetdash{}{0pt}%
\pgfpathmoveto{\pgfqpoint{2.041556in}{2.781923in}}%
\pgfpathlineto{\pgfqpoint{2.013849in}{2.856887in}}%
\pgfpathlineto{\pgfqpoint{1.986168in}{2.931780in}}%
\pgfpathlineto{\pgfqpoint{2.035043in}{2.946178in}}%
\pgfpathlineto{\pgfqpoint{2.083831in}{2.960549in}}%
\pgfpathlineto{\pgfqpoint{2.111508in}{2.885805in}}%
\pgfpathlineto{\pgfqpoint{2.139212in}{2.810991in}}%
\pgfpathlineto{\pgfqpoint{2.090428in}{2.796470in}}%
\pgfpathlineto{\pgfqpoint{2.041556in}{2.781923in}}%
\pgfpathclose%
\pgfusepath{stroke,fill}%
\end{pgfscope}%
\begin{pgfscope}%
\pgfpathrectangle{\pgfqpoint{0.050000in}{0.083252in}}{\pgfqpoint{4.900000in}{4.900000in}}%
\pgfusepath{clip}%
\pgfsetbuttcap%
\pgfsetroundjoin%
\definecolor{currentfill}{rgb}{0.278001,0.248228,0.568289}%
\pgfsetfillcolor{currentfill}%
\pgfsetfillopacity{0.950000}%
\pgfsetlinewidth{0.200750pt}%
\definecolor{currentstroke}{rgb}{0.200000,0.200000,0.200000}%
\pgfsetstrokecolor{currentstroke}%
\pgfsetdash{}{0pt}%
\pgfpathmoveto{\pgfqpoint{2.250289in}{2.511084in}}%
\pgfpathlineto{\pgfqpoint{2.222480in}{2.586150in}}%
\pgfpathlineto{\pgfqpoint{2.194698in}{2.661161in}}%
\pgfpathlineto{\pgfqpoint{2.243391in}{2.675809in}}%
\pgfpathlineto{\pgfqpoint{2.291997in}{2.690433in}}%
\pgfpathlineto{\pgfqpoint{2.319776in}{2.615583in}}%
\pgfpathlineto{\pgfqpoint{2.347581in}{2.540683in}}%
\pgfpathlineto{\pgfqpoint{2.298979in}{2.525895in}}%
\pgfpathlineto{\pgfqpoint{2.250289in}{2.511084in}}%
\pgfpathclose%
\pgfusepath{stroke,fill}%
\end{pgfscope}%
\begin{pgfscope}%
\pgfpathrectangle{\pgfqpoint{0.050000in}{0.083252in}}{\pgfqpoint{4.900000in}{4.900000in}}%
\pgfusepath{clip}%
\pgfsetbuttcap%
\pgfsetroundjoin%
\definecolor{currentfill}{rgb}{0.433141,0.409765,0.661053}%
\pgfsetfillcolor{currentfill}%
\pgfsetfillopacity{0.950000}%
\pgfsetlinewidth{0.200750pt}%
\definecolor{currentstroke}{rgb}{0.200000,0.200000,0.200000}%
\pgfsetstrokecolor{currentstroke}%
\pgfsetdash{}{0pt}%
\pgfpathmoveto{\pgfqpoint{2.459065in}{2.241002in}}%
\pgfpathlineto{\pgfqpoint{2.431155in}{2.315868in}}%
\pgfpathlineto{\pgfqpoint{2.403270in}{2.390799in}}%
\pgfpathlineto{\pgfqpoint{2.451782in}{2.405753in}}%
\pgfpathlineto{\pgfqpoint{2.500207in}{2.420689in}}%
\pgfpathlineto{\pgfqpoint{2.528087in}{2.345976in}}%
\pgfpathlineto{\pgfqpoint{2.555995in}{2.271355in}}%
\pgfpathlineto{\pgfqpoint{2.507573in}{2.256186in}}%
\pgfpathlineto{\pgfqpoint{2.459065in}{2.241002in}}%
\pgfpathclose%
\pgfusepath{stroke,fill}%
\end{pgfscope}%
\begin{pgfscope}%
\pgfpathrectangle{\pgfqpoint{0.050000in}{0.083252in}}{\pgfqpoint{4.900000in}{4.900000in}}%
\pgfusepath{clip}%
\pgfsetbuttcap%
\pgfsetroundjoin%
\definecolor{currentfill}{rgb}{0.850827,0.844675,0.910804}%
\pgfsetfillcolor{currentfill}%
\pgfsetfillopacity{0.950000}%
\pgfsetlinewidth{0.200750pt}%
\definecolor{currentstroke}{rgb}{0.200000,0.200000,0.200000}%
\pgfsetstrokecolor{currentstroke}%
\pgfsetdash{}{0pt}%
\pgfpathmoveto{\pgfqpoint{3.182748in}{1.549440in}}%
\pgfpathlineto{\pgfqpoint{3.154235in}{1.606508in}}%
\pgfpathlineto{\pgfqpoint{3.125809in}{1.666443in}}%
\pgfpathlineto{\pgfqpoint{3.173680in}{1.683338in}}%
\pgfpathlineto{\pgfqpoint{3.221469in}{1.700183in}}%
\pgfpathlineto{\pgfqpoint{3.249926in}{1.640505in}}%
\pgfpathlineto{\pgfqpoint{3.278476in}{1.583611in}}%
\pgfpathlineto{\pgfqpoint{3.230653in}{1.566554in}}%
\pgfpathlineto{\pgfqpoint{3.182748in}{1.549440in}}%
\pgfpathclose%
\pgfusepath{stroke,fill}%
\end{pgfscope}%
\begin{pgfscope}%
\pgfpathrectangle{\pgfqpoint{0.050000in}{0.083252in}}{\pgfqpoint{4.900000in}{4.900000in}}%
\pgfusepath{clip}%
\pgfsetbuttcap%
\pgfsetroundjoin%
\definecolor{currentfill}{rgb}{0.588281,0.571303,0.753818}%
\pgfsetfillcolor{currentfill}%
\pgfsetfillopacity{0.950000}%
\pgfsetlinewidth{0.200750pt}%
\definecolor{currentstroke}{rgb}{0.200000,0.200000,0.200000}%
\pgfsetstrokecolor{currentstroke}%
\pgfsetdash{}{0pt}%
\pgfpathmoveto{\pgfqpoint{2.667904in}{1.975731in}}%
\pgfpathlineto{\pgfqpoint{2.639883in}{2.048933in}}%
\pgfpathlineto{\pgfqpoint{2.611892in}{2.122715in}}%
\pgfpathlineto{\pgfqpoint{2.660225in}{2.138168in}}%
\pgfpathlineto{\pgfqpoint{2.708472in}{2.153606in}}%
\pgfpathlineto{\pgfqpoint{2.736463in}{2.080184in}}%
\pgfpathlineto{\pgfqpoint{2.764486in}{2.007381in}}%
\pgfpathlineto{\pgfqpoint{2.716237in}{1.991564in}}%
\pgfpathlineto{\pgfqpoint{2.667904in}{1.975731in}}%
\pgfpathclose%
\pgfusepath{stroke,fill}%
\end{pgfscope}%
\begin{pgfscope}%
\pgfpathrectangle{\pgfqpoint{0.050000in}{0.083252in}}{\pgfqpoint{4.900000in}{4.900000in}}%
\pgfusepath{clip}%
\pgfsetbuttcap%
\pgfsetroundjoin%
\definecolor{currentfill}{rgb}{0.898562,0.894379,0.939346}%
\pgfsetfillcolor{currentfill}%
\pgfsetfillopacity{0.950000}%
\pgfsetlinewidth{0.200750pt}%
\definecolor{currentstroke}{rgb}{0.200000,0.200000,0.200000}%
\pgfsetstrokecolor{currentstroke}%
\pgfsetdash{}{0pt}%
\pgfpathmoveto{\pgfqpoint{3.335888in}{1.478856in}}%
\pgfpathlineto{\pgfqpoint{3.307127in}{1.529688in}}%
\pgfpathlineto{\pgfqpoint{3.278476in}{1.583611in}}%
\pgfpathlineto{\pgfqpoint{3.326215in}{1.600614in}}%
\pgfpathlineto{\pgfqpoint{3.373871in}{1.617563in}}%
\pgfpathlineto{\pgfqpoint{3.402562in}{1.563719in}}%
\pgfpathlineto{\pgfqpoint{3.431369in}{1.512872in}}%
\pgfpathlineto{\pgfqpoint{3.383670in}{1.495890in}}%
\pgfpathlineto{\pgfqpoint{3.335888in}{1.478856in}}%
\pgfpathclose%
\pgfusepath{stroke,fill}%
\end{pgfscope}%
\begin{pgfscope}%
\pgfpathrectangle{\pgfqpoint{0.050000in}{0.083252in}}{\pgfqpoint{4.900000in}{4.900000in}}%
\pgfusepath{clip}%
\pgfsetbuttcap%
\pgfsetroundjoin%
\definecolor{currentfill}{rgb}{0.267389,0.507190,0.943053}%
\pgfsetfillcolor{currentfill}%
\pgfsetfillopacity{0.950000}%
\pgfsetlinewidth{0.200750pt}%
\definecolor{currentstroke}{rgb}{0.200000,0.200000,0.200000}%
\pgfsetstrokecolor{currentstroke}%
\pgfsetdash{}{0pt}%
\pgfpathmoveto{\pgfqpoint{1.261730in}{3.715961in}}%
\pgfpathlineto{\pgfqpoint{1.234378in}{3.790523in}}%
\pgfpathlineto{\pgfqpoint{1.207053in}{3.865013in}}%
\pgfpathlineto{\pgfqpoint{1.256654in}{3.878582in}}%
\pgfpathlineto{\pgfqpoint{1.306165in}{3.892125in}}%
\pgfpathlineto{\pgfqpoint{1.333488in}{3.817781in}}%
\pgfpathlineto{\pgfqpoint{1.360837in}{3.743367in}}%
\pgfpathlineto{\pgfqpoint{1.311328in}{3.729676in}}%
\pgfpathlineto{\pgfqpoint{1.261730in}{3.715961in}}%
\pgfpathclose%
\pgfusepath{stroke,fill}%
\end{pgfscope}%
\begin{pgfscope}%
\pgfpathrectangle{\pgfqpoint{0.050000in}{0.083252in}}{\pgfqpoint{4.900000in}{4.900000in}}%
\pgfusepath{clip}%
\pgfsetbuttcap%
\pgfsetroundjoin%
\definecolor{currentfill}{rgb}{0.797124,0.788758,0.878693}%
\pgfsetfillcolor{currentfill}%
\pgfsetfillopacity{0.950000}%
\pgfsetlinewidth{0.200750pt}%
\definecolor{currentstroke}{rgb}{0.200000,0.200000,0.200000}%
\pgfsetstrokecolor{currentstroke}%
\pgfsetdash{}{0pt}%
\pgfpathmoveto{\pgfqpoint{3.029815in}{1.632500in}}%
\pgfpathlineto{\pgfqpoint{3.001494in}{1.695394in}}%
\pgfpathlineto{\pgfqpoint{2.973238in}{1.760661in}}%
\pgfpathlineto{\pgfqpoint{3.021253in}{1.777281in}}%
\pgfpathlineto{\pgfqpoint{3.069184in}{1.793862in}}%
\pgfpathlineto{\pgfqpoint{3.097462in}{1.728992in}}%
\pgfpathlineto{\pgfqpoint{3.125809in}{1.666443in}}%
\pgfpathlineto{\pgfqpoint{3.077854in}{1.649498in}}%
\pgfpathlineto{\pgfqpoint{3.029815in}{1.632500in}}%
\pgfpathclose%
\pgfusepath{stroke,fill}%
\end{pgfscope}%
\begin{pgfscope}%
\pgfpathrectangle{\pgfqpoint{0.050000in}{0.083252in}}{\pgfqpoint{4.900000in}{4.900000in}}%
\pgfusepath{clip}%
\pgfsetbuttcap%
\pgfsetroundjoin%
\definecolor{currentfill}{rgb}{0.976132,0.975148,0.985729}%
\pgfsetfillcolor{currentfill}%
\pgfsetfillopacity{0.950000}%
\pgfsetlinewidth{0.200750pt}%
\definecolor{currentstroke}{rgb}{0.200000,0.200000,0.200000}%
\pgfsetstrokecolor{currentstroke}%
\pgfsetdash{}{0pt}%
\pgfpathmoveto{\pgfqpoint{3.892402in}{1.364328in}}%
\pgfpathlineto{\pgfqpoint{3.862612in}{1.400041in}}%
\pgfpathlineto{\pgfqpoint{3.832965in}{1.437777in}}%
\pgfpathlineto{\pgfqpoint{3.880187in}{1.454100in}}%
\pgfpathlineto{\pgfqpoint{3.927327in}{1.470389in}}%
\pgfpathlineto{\pgfqpoint{3.957034in}{1.432376in}}%
\pgfpathlineto{\pgfqpoint{3.986887in}{1.396358in}}%
\pgfpathlineto{\pgfqpoint{3.939687in}{1.380358in}}%
\pgfpathlineto{\pgfqpoint{3.892402in}{1.364328in}}%
\pgfpathclose%
\pgfusepath{stroke,fill}%
\end{pgfscope}%
\begin{pgfscope}%
\pgfpathrectangle{\pgfqpoint{0.050000in}{0.083252in}}{\pgfqpoint{4.900000in}{4.900000in}}%
\pgfusepath{clip}%
\pgfsetbuttcap%
\pgfsetroundjoin%
\definecolor{currentfill}{rgb}{0.260746,0.436601,0.849212}%
\pgfsetfillcolor{currentfill}%
\pgfsetfillopacity{0.950000}%
\pgfsetlinewidth{0.200750pt}%
\definecolor{currentstroke}{rgb}{0.200000,0.200000,0.200000}%
\pgfsetstrokecolor{currentstroke}%
\pgfsetdash{}{0pt}%
\pgfpathmoveto{\pgfqpoint{1.470494in}{3.445001in}}%
\pgfpathlineto{\pgfqpoint{1.443041in}{3.519698in}}%
\pgfpathlineto{\pgfqpoint{1.415614in}{3.594325in}}%
\pgfpathlineto{\pgfqpoint{1.465032in}{3.608138in}}%
\pgfpathlineto{\pgfqpoint{1.514362in}{3.621925in}}%
\pgfpathlineto{\pgfqpoint{1.541787in}{3.547445in}}%
\pgfpathlineto{\pgfqpoint{1.569238in}{3.472894in}}%
\pgfpathlineto{\pgfqpoint{1.519910in}{3.458960in}}%
\pgfpathlineto{\pgfqpoint{1.470494in}{3.445001in}}%
\pgfpathclose%
\pgfusepath{stroke,fill}%
\end{pgfscope}%
\begin{pgfscope}%
\pgfpathrectangle{\pgfqpoint{0.050000in}{0.083252in}}{\pgfqpoint{4.900000in}{4.900000in}}%
\pgfusepath{clip}%
\pgfsetbuttcap%
\pgfsetroundjoin%
\definecolor{currentfill}{rgb}{0.254102,0.366013,0.755371}%
\pgfsetfillcolor{currentfill}%
\pgfsetfillopacity{0.950000}%
\pgfsetlinewidth{0.200750pt}%
\definecolor{currentstroke}{rgb}{0.200000,0.200000,0.200000}%
\pgfsetstrokecolor{currentstroke}%
\pgfsetdash{}{0pt}%
\pgfpathmoveto{\pgfqpoint{1.679302in}{3.173983in}}%
\pgfpathlineto{\pgfqpoint{1.651747in}{3.248817in}}%
\pgfpathlineto{\pgfqpoint{1.624218in}{3.323581in}}%
\pgfpathlineto{\pgfqpoint{1.673454in}{3.337637in}}%
\pgfpathlineto{\pgfqpoint{1.722602in}{3.351668in}}%
\pgfpathlineto{\pgfqpoint{1.750129in}{3.277052in}}%
\pgfpathlineto{\pgfqpoint{1.777681in}{3.202365in}}%
\pgfpathlineto{\pgfqpoint{1.728536in}{3.188187in}}%
\pgfpathlineto{\pgfqpoint{1.679302in}{3.173983in}}%
\pgfpathclose%
\pgfusepath{stroke,fill}%
\end{pgfscope}%
\begin{pgfscope}%
\pgfpathrectangle{\pgfqpoint{0.050000in}{0.083252in}}{\pgfqpoint{4.900000in}{4.900000in}}%
\pgfusepath{clip}%
\pgfsetbuttcap%
\pgfsetroundjoin%
\definecolor{currentfill}{rgb}{1.000000,1.000000,1.000000}%
\pgfsetfillcolor{currentfill}%
\pgfsetfillopacity{0.950000}%
\pgfsetlinewidth{0.200750pt}%
\definecolor{currentstroke}{rgb}{0.200000,0.200000,0.200000}%
\pgfsetstrokecolor{currentstroke}%
\pgfsetdash{}{0pt}%
\pgfpathmoveto{\pgfqpoint{4.296052in}{1.325775in}}%
\pgfpathlineto{\pgfqpoint{4.265581in}{1.358398in}}%
\pgfpathlineto{\pgfqpoint{4.235212in}{1.391332in}}%
\pgfpathlineto{\pgfqpoint{4.282053in}{1.406863in}}%
\pgfpathlineto{\pgfqpoint{4.328809in}{1.422365in}}%
\pgfpathlineto{\pgfqpoint{4.359239in}{1.389088in}}%
\pgfpathlineto{\pgfqpoint{4.389771in}{1.356107in}}%
\pgfpathlineto{\pgfqpoint{4.342955in}{1.340955in}}%
\pgfpathlineto{\pgfqpoint{4.296052in}{1.325775in}}%
\pgfpathclose%
\pgfusepath{stroke,fill}%
\end{pgfscope}%
\begin{pgfscope}%
\pgfpathrectangle{\pgfqpoint{0.050000in}{0.083252in}}{\pgfqpoint{4.900000in}{4.900000in}}%
\pgfusepath{clip}%
\pgfsetbuttcap%
\pgfsetroundjoin%
\definecolor{currentfill}{rgb}{0.247459,0.295425,0.661530}%
\pgfsetfillcolor{currentfill}%
\pgfsetfillopacity{0.950000}%
\pgfsetlinewidth{0.200750pt}%
\definecolor{currentstroke}{rgb}{0.200000,0.200000,0.200000}%
\pgfsetstrokecolor{currentstroke}%
\pgfsetdash{}{0pt}%
\pgfpathmoveto{\pgfqpoint{1.888154in}{2.902908in}}%
\pgfpathlineto{\pgfqpoint{1.860496in}{2.977879in}}%
\pgfpathlineto{\pgfqpoint{1.832865in}{3.052779in}}%
\pgfpathlineto{\pgfqpoint{1.881919in}{3.067080in}}%
\pgfpathlineto{\pgfqpoint{1.930885in}{3.081355in}}%
\pgfpathlineto{\pgfqpoint{1.958514in}{3.006603in}}%
\pgfpathlineto{\pgfqpoint{1.986168in}{2.931780in}}%
\pgfpathlineto{\pgfqpoint{1.937205in}{2.917357in}}%
\pgfpathlineto{\pgfqpoint{1.888154in}{2.902908in}}%
\pgfpathclose%
\pgfusepath{stroke,fill}%
\end{pgfscope}%
\begin{pgfscope}%
\pgfpathrectangle{\pgfqpoint{0.050000in}{0.083252in}}{\pgfqpoint{4.900000in}{4.900000in}}%
\pgfusepath{clip}%
\pgfsetbuttcap%
\pgfsetroundjoin%
\definecolor{currentfill}{rgb}{0.928397,0.925444,0.957186}%
\pgfsetfillcolor{currentfill}%
\pgfsetfillopacity{0.950000}%
\pgfsetlinewidth{0.200750pt}%
\definecolor{currentstroke}{rgb}{0.200000,0.200000,0.200000}%
\pgfsetstrokecolor{currentstroke}%
\pgfsetdash{}{0pt}%
\pgfpathmoveto{\pgfqpoint{3.489354in}{1.420262in}}%
\pgfpathlineto{\pgfqpoint{3.460298in}{1.465066in}}%
\pgfpathlineto{\pgfqpoint{3.431369in}{1.512872in}}%
\pgfpathlineto{\pgfqpoint{3.478985in}{1.529803in}}%
\pgfpathlineto{\pgfqpoint{3.526517in}{1.546684in}}%
\pgfpathlineto{\pgfqpoint{3.555495in}{1.498786in}}%
\pgfpathlineto{\pgfqpoint{3.584605in}{1.453813in}}%
\pgfpathlineto{\pgfqpoint{3.537021in}{1.437059in}}%
\pgfpathlineto{\pgfqpoint{3.489354in}{1.420262in}}%
\pgfpathclose%
\pgfusepath{stroke,fill}%
\end{pgfscope}%
\begin{pgfscope}%
\pgfpathrectangle{\pgfqpoint{0.050000in}{0.083252in}}{\pgfqpoint{4.900000in}{4.900000in}}%
\pgfusepath{clip}%
\pgfsetbuttcap%
\pgfsetroundjoin%
\definecolor{currentfill}{rgb}{0.737455,0.726628,0.843014}%
\pgfsetfillcolor{currentfill}%
\pgfsetfillopacity{0.950000}%
\pgfsetlinewidth{0.200750pt}%
\definecolor{currentstroke}{rgb}{0.200000,0.200000,0.200000}%
\pgfsetstrokecolor{currentstroke}%
\pgfsetdash{}{0pt}%
\pgfpathmoveto{\pgfqpoint{2.876956in}{1.727299in}}%
\pgfpathlineto{\pgfqpoint{2.848773in}{1.795060in}}%
\pgfpathlineto{\pgfqpoint{2.820637in}{1.864542in}}%
\pgfpathlineto{\pgfqpoint{2.868807in}{1.880781in}}%
\pgfpathlineto{\pgfqpoint{2.916892in}{1.896993in}}%
\pgfpathlineto{\pgfqpoint{2.945039in}{1.827970in}}%
\pgfpathlineto{\pgfqpoint{2.973238in}{1.760661in}}%
\pgfpathlineto{\pgfqpoint{2.925139in}{1.744000in}}%
\pgfpathlineto{\pgfqpoint{2.876956in}{1.727299in}}%
\pgfpathclose%
\pgfusepath{stroke,fill}%
\end{pgfscope}%
\begin{pgfscope}%
\pgfpathrectangle{\pgfqpoint{0.050000in}{0.083252in}}{\pgfqpoint{4.900000in}{4.900000in}}%
\pgfusepath{clip}%
\pgfsetbuttcap%
\pgfsetroundjoin%
\definecolor{currentfill}{rgb}{0.241061,0.227451,0.571165}%
\pgfsetfillcolor{currentfill}%
\pgfsetfillopacity{0.950000}%
\pgfsetlinewidth{0.200750pt}%
\definecolor{currentstroke}{rgb}{0.200000,0.200000,0.200000}%
\pgfsetstrokecolor{currentstroke}%
\pgfsetdash{}{0pt}%
\pgfpathmoveto{\pgfqpoint{2.097049in}{2.631788in}}%
\pgfpathlineto{\pgfqpoint{2.069289in}{2.706890in}}%
\pgfpathlineto{\pgfqpoint{2.041556in}{2.781923in}}%
\pgfpathlineto{\pgfqpoint{2.090428in}{2.796470in}}%
\pgfpathlineto{\pgfqpoint{2.139212in}{2.810991in}}%
\pgfpathlineto{\pgfqpoint{2.166942in}{2.736109in}}%
\pgfpathlineto{\pgfqpoint{2.194698in}{2.661161in}}%
\pgfpathlineto{\pgfqpoint{2.145917in}{2.646487in}}%
\pgfpathlineto{\pgfqpoint{2.097049in}{2.631788in}}%
\pgfpathclose%
\pgfusepath{stroke,fill}%
\end{pgfscope}%
\begin{pgfscope}%
\pgfpathrectangle{\pgfqpoint{0.050000in}{0.083252in}}{\pgfqpoint{4.900000in}{4.900000in}}%
\pgfusepath{clip}%
\pgfsetbuttcap%
\pgfsetroundjoin%
\definecolor{currentfill}{rgb}{0.355571,0.328997,0.614671}%
\pgfsetfillcolor{currentfill}%
\pgfsetfillopacity{0.950000}%
\pgfsetlinewidth{0.200750pt}%
\definecolor{currentstroke}{rgb}{0.200000,0.200000,0.200000}%
\pgfsetstrokecolor{currentstroke}%
\pgfsetdash{}{0pt}%
\pgfpathmoveto{\pgfqpoint{2.305986in}{2.360831in}}%
\pgfpathlineto{\pgfqpoint{2.278124in}{2.435972in}}%
\pgfpathlineto{\pgfqpoint{2.250289in}{2.511084in}}%
\pgfpathlineto{\pgfqpoint{2.298979in}{2.525895in}}%
\pgfpathlineto{\pgfqpoint{2.347581in}{2.540683in}}%
\pgfpathlineto{\pgfqpoint{2.375413in}{2.465748in}}%
\pgfpathlineto{\pgfqpoint{2.403270in}{2.390799in}}%
\pgfpathlineto{\pgfqpoint{2.354672in}{2.375825in}}%
\pgfpathlineto{\pgfqpoint{2.305986in}{2.360831in}}%
\pgfpathclose%
\pgfusepath{stroke,fill}%
\end{pgfscope}%
\begin{pgfscope}%
\pgfpathrectangle{\pgfqpoint{0.050000in}{0.083252in}}{\pgfqpoint{4.900000in}{4.900000in}}%
\pgfusepath{clip}%
\pgfsetbuttcap%
\pgfsetroundjoin%
\definecolor{currentfill}{rgb}{0.516678,0.496747,0.711003}%
\pgfsetfillcolor{currentfill}%
\pgfsetfillopacity{0.950000}%
\pgfsetlinewidth{0.200750pt}%
\definecolor{currentstroke}{rgb}{0.200000,0.200000,0.200000}%
\pgfsetstrokecolor{currentstroke}%
\pgfsetdash{}{0pt}%
\pgfpathmoveto{\pgfqpoint{2.514967in}{2.091767in}}%
\pgfpathlineto{\pgfqpoint{2.487002in}{2.166270in}}%
\pgfpathlineto{\pgfqpoint{2.459065in}{2.241002in}}%
\pgfpathlineto{\pgfqpoint{2.507573in}{2.256186in}}%
\pgfpathlineto{\pgfqpoint{2.555995in}{2.271355in}}%
\pgfpathlineto{\pgfqpoint{2.583929in}{2.196900in}}%
\pgfpathlineto{\pgfqpoint{2.611892in}{2.122715in}}%
\pgfpathlineto{\pgfqpoint{2.563472in}{2.107248in}}%
\pgfpathlineto{\pgfqpoint{2.514967in}{2.091767in}}%
\pgfpathclose%
\pgfusepath{stroke,fill}%
\end{pgfscope}%
\begin{pgfscope}%
\pgfpathrectangle{\pgfqpoint{0.050000in}{0.083252in}}{\pgfqpoint{4.900000in}{4.900000in}}%
\pgfusepath{clip}%
\pgfsetbuttcap%
\pgfsetroundjoin%
\definecolor{currentfill}{rgb}{0.263945,0.470588,0.894394}%
\pgfsetfillcolor{currentfill}%
\pgfsetfillopacity{0.950000}%
\pgfsetlinewidth{0.200750pt}%
\definecolor{currentstroke}{rgb}{0.200000,0.200000,0.200000}%
\pgfsetstrokecolor{currentstroke}%
\pgfsetdash{}{0pt}%
\pgfpathmoveto{\pgfqpoint{1.316510in}{3.566626in}}%
\pgfpathlineto{\pgfqpoint{1.289107in}{3.641329in}}%
\pgfpathlineto{\pgfqpoint{1.261730in}{3.715961in}}%
\pgfpathlineto{\pgfqpoint{1.311328in}{3.729676in}}%
\pgfpathlineto{\pgfqpoint{1.360837in}{3.743367in}}%
\pgfpathlineto{\pgfqpoint{1.388213in}{3.668881in}}%
\pgfpathlineto{\pgfqpoint{1.415614in}{3.594325in}}%
\pgfpathlineto{\pgfqpoint{1.366106in}{3.580488in}}%
\pgfpathlineto{\pgfqpoint{1.316510in}{3.566626in}}%
\pgfpathclose%
\pgfusepath{stroke,fill}%
\end{pgfscope}%
\begin{pgfscope}%
\pgfpathrectangle{\pgfqpoint{0.050000in}{0.083252in}}{\pgfqpoint{4.900000in}{4.900000in}}%
\pgfusepath{clip}%
\pgfsetbuttcap%
\pgfsetroundjoin%
\definecolor{currentfill}{rgb}{0.665852,0.652072,0.800200}%
\pgfsetfillcolor{currentfill}%
\pgfsetfillopacity{0.950000}%
\pgfsetlinewidth{0.200750pt}%
\definecolor{currentstroke}{rgb}{0.200000,0.200000,0.200000}%
\pgfsetstrokecolor{currentstroke}%
\pgfsetdash{}{0pt}%
\pgfpathmoveto{\pgfqpoint{2.724044in}{1.831986in}}%
\pgfpathlineto{\pgfqpoint{2.695956in}{1.903326in}}%
\pgfpathlineto{\pgfqpoint{2.667904in}{1.975731in}}%
\pgfpathlineto{\pgfqpoint{2.716237in}{1.991564in}}%
\pgfpathlineto{\pgfqpoint{2.764486in}{2.007381in}}%
\pgfpathlineto{\pgfqpoint{2.792543in}{1.935415in}}%
\pgfpathlineto{\pgfqpoint{2.820637in}{1.864542in}}%
\pgfpathlineto{\pgfqpoint{2.772383in}{1.848277in}}%
\pgfpathlineto{\pgfqpoint{2.724044in}{1.831986in}}%
\pgfpathclose%
\pgfusepath{stroke,fill}%
\end{pgfscope}%
\begin{pgfscope}%
\pgfpathrectangle{\pgfqpoint{0.050000in}{0.083252in}}{\pgfqpoint{4.900000in}{4.900000in}}%
\pgfusepath{clip}%
\pgfsetbuttcap%
\pgfsetroundjoin%
\definecolor{currentfill}{rgb}{0.958231,0.956509,0.975025}%
\pgfsetfillcolor{currentfill}%
\pgfsetfillopacity{0.950000}%
\pgfsetlinewidth{0.200750pt}%
\definecolor{currentstroke}{rgb}{0.200000,0.200000,0.200000}%
\pgfsetstrokecolor{currentstroke}%
\pgfsetdash{}{0pt}%
\pgfpathmoveto{\pgfqpoint{3.643235in}{1.372131in}}%
\pgfpathlineto{\pgfqpoint{3.613851in}{1.411655in}}%
\pgfpathlineto{\pgfqpoint{3.584605in}{1.453813in}}%
\pgfpathlineto{\pgfqpoint{3.632105in}{1.470523in}}%
\pgfpathlineto{\pgfqpoint{3.679521in}{1.487191in}}%
\pgfpathlineto{\pgfqpoint{3.708824in}{1.444816in}}%
\pgfpathlineto{\pgfqpoint{3.738268in}{1.405024in}}%
\pgfpathlineto{\pgfqpoint{3.690794in}{1.388595in}}%
\pgfpathlineto{\pgfqpoint{3.643235in}{1.372131in}}%
\pgfpathclose%
\pgfusepath{stroke,fill}%
\end{pgfscope}%
\begin{pgfscope}%
\pgfpathrectangle{\pgfqpoint{0.050000in}{0.083252in}}{\pgfqpoint{4.900000in}{4.900000in}}%
\pgfusepath{clip}%
\pgfsetbuttcap%
\pgfsetroundjoin%
\definecolor{currentfill}{rgb}{0.988066,0.987574,0.992864}%
\pgfsetfillcolor{currentfill}%
\pgfsetfillopacity{0.950000}%
\pgfsetlinewidth{0.200750pt}%
\definecolor{currentstroke}{rgb}{0.200000,0.200000,0.200000}%
\pgfsetstrokecolor{currentstroke}%
\pgfsetdash{}{0pt}%
\pgfpathmoveto{\pgfqpoint{4.046999in}{1.328925in}}%
\pgfpathlineto{\pgfqpoint{4.016878in}{1.362003in}}%
\pgfpathlineto{\pgfqpoint{3.986887in}{1.396358in}}%
\pgfpathlineto{\pgfqpoint{4.034003in}{1.412329in}}%
\pgfpathlineto{\pgfqpoint{4.081035in}{1.428269in}}%
\pgfpathlineto{\pgfqpoint{4.111089in}{1.393596in}}%
\pgfpathlineto{\pgfqpoint{4.141276in}{1.360185in}}%
\pgfpathlineto{\pgfqpoint{4.094180in}{1.344569in}}%
\pgfpathlineto{\pgfqpoint{4.046999in}{1.328925in}}%
\pgfpathclose%
\pgfusepath{stroke,fill}%
\end{pgfscope}%
\begin{pgfscope}%
\pgfpathrectangle{\pgfqpoint{0.050000in}{0.083252in}}{\pgfqpoint{4.900000in}{4.900000in}}%
\pgfusepath{clip}%
\pgfsetbuttcap%
\pgfsetroundjoin%
\definecolor{currentfill}{rgb}{0.257547,0.402614,0.804029}%
\pgfsetfillcolor{currentfill}%
\pgfsetfillopacity{0.950000}%
\pgfsetlinewidth{0.200750pt}%
\definecolor{currentstroke}{rgb}{0.200000,0.200000,0.200000}%
\pgfsetstrokecolor{currentstroke}%
\pgfsetdash{}{0pt}%
\pgfpathmoveto{\pgfqpoint{1.525479in}{3.295391in}}%
\pgfpathlineto{\pgfqpoint{1.497973in}{3.370232in}}%
\pgfpathlineto{\pgfqpoint{1.470494in}{3.445001in}}%
\pgfpathlineto{\pgfqpoint{1.519910in}{3.458960in}}%
\pgfpathlineto{\pgfqpoint{1.569238in}{3.472894in}}%
\pgfpathlineto{\pgfqpoint{1.596715in}{3.398273in}}%
\pgfpathlineto{\pgfqpoint{1.624218in}{3.323581in}}%
\pgfpathlineto{\pgfqpoint{1.574893in}{3.309499in}}%
\pgfpathlineto{\pgfqpoint{1.525479in}{3.295391in}}%
\pgfpathclose%
\pgfusepath{stroke,fill}%
\end{pgfscope}%
\begin{pgfscope}%
\pgfpathrectangle{\pgfqpoint{0.050000in}{0.083252in}}{\pgfqpoint{4.900000in}{4.900000in}}%
\pgfusepath{clip}%
\pgfsetbuttcap%
\pgfsetroundjoin%
\definecolor{currentfill}{rgb}{0.269850,0.533333,0.977809}%
\pgfsetfillcolor{currentfill}%
\pgfsetfillopacity{0.950000}%
\pgfsetlinewidth{0.200750pt}%
\definecolor{currentstroke}{rgb}{0.200000,0.200000,0.200000}%
\pgfsetstrokecolor{currentstroke}%
\pgfsetdash{}{0pt}%
\pgfpathmoveto{\pgfqpoint{1.107585in}{3.837803in}}%
\pgfpathlineto{\pgfqpoint{1.080284in}{3.912369in}}%
\pgfpathlineto{\pgfqpoint{1.130063in}{3.925913in}}%
\pgfpathlineto{\pgfqpoint{1.179754in}{3.939433in}}%
\pgfpathlineto{\pgfqpoint{1.207053in}{3.865013in}}%
\pgfpathlineto{\pgfqpoint{1.157364in}{3.851420in}}%
\pgfpathlineto{\pgfqpoint{1.107585in}{3.837803in}}%
\pgfpathclose%
\pgfusepath{stroke,fill}%
\end{pgfscope}%
\begin{pgfscope}%
\pgfpathrectangle{\pgfqpoint{0.050000in}{0.083252in}}{\pgfqpoint{4.900000in}{4.900000in}}%
\pgfusepath{clip}%
\pgfsetbuttcap%
\pgfsetroundjoin%
\definecolor{currentfill}{rgb}{0.250903,0.332026,0.710188}%
\pgfsetfillcolor{currentfill}%
\pgfsetfillopacity{0.950000}%
\pgfsetlinewidth{0.200750pt}%
\definecolor{currentstroke}{rgb}{0.200000,0.200000,0.200000}%
\pgfsetstrokecolor{currentstroke}%
\pgfsetdash{}{0pt}%
\pgfpathmoveto{\pgfqpoint{1.734492in}{3.024100in}}%
\pgfpathlineto{\pgfqpoint{1.706884in}{3.099077in}}%
\pgfpathlineto{\pgfqpoint{1.679302in}{3.173983in}}%
\pgfpathlineto{\pgfqpoint{1.728536in}{3.188187in}}%
\pgfpathlineto{\pgfqpoint{1.777681in}{3.202365in}}%
\pgfpathlineto{\pgfqpoint{1.805260in}{3.127608in}}%
\pgfpathlineto{\pgfqpoint{1.832865in}{3.052779in}}%
\pgfpathlineto{\pgfqpoint{1.783722in}{3.038452in}}%
\pgfpathlineto{\pgfqpoint{1.734492in}{3.024100in}}%
\pgfpathclose%
\pgfusepath{stroke,fill}%
\end{pgfscope}%
\begin{pgfscope}%
\pgfpathrectangle{\pgfqpoint{0.050000in}{0.083252in}}{\pgfqpoint{4.900000in}{4.900000in}}%
\pgfusepath{clip}%
\pgfsetbuttcap%
\pgfsetroundjoin%
\definecolor{currentfill}{rgb}{0.244260,0.261438,0.616348}%
\pgfsetfillcolor{currentfill}%
\pgfsetfillopacity{0.950000}%
\pgfsetlinewidth{0.200750pt}%
\definecolor{currentstroke}{rgb}{0.200000,0.200000,0.200000}%
\pgfsetstrokecolor{currentstroke}%
\pgfsetdash{}{0pt}%
\pgfpathmoveto{\pgfqpoint{1.943548in}{2.752753in}}%
\pgfpathlineto{\pgfqpoint{1.915837in}{2.827866in}}%
\pgfpathlineto{\pgfqpoint{1.888154in}{2.902908in}}%
\pgfpathlineto{\pgfqpoint{1.937205in}{2.917357in}}%
\pgfpathlineto{\pgfqpoint{1.986168in}{2.931780in}}%
\pgfpathlineto{\pgfqpoint{2.013849in}{2.856887in}}%
\pgfpathlineto{\pgfqpoint{2.041556in}{2.781923in}}%
\pgfpathlineto{\pgfqpoint{1.992596in}{2.767351in}}%
\pgfpathlineto{\pgfqpoint{1.943548in}{2.752753in}}%
\pgfpathclose%
\pgfusepath{stroke,fill}%
\end{pgfscope}%
\begin{pgfscope}%
\pgfpathrectangle{\pgfqpoint{0.050000in}{0.083252in}}{\pgfqpoint{4.900000in}{4.900000in}}%
\pgfusepath{clip}%
\pgfsetbuttcap%
\pgfsetroundjoin%
\definecolor{currentfill}{rgb}{0.278001,0.248228,0.568289}%
\pgfsetfillcolor{currentfill}%
\pgfsetfillopacity{0.950000}%
\pgfsetlinewidth{0.200750pt}%
\definecolor{currentstroke}{rgb}{0.200000,0.200000,0.200000}%
\pgfsetstrokecolor{currentstroke}%
\pgfsetdash{}{0pt}%
\pgfpathmoveto{\pgfqpoint{2.152647in}{2.481391in}}%
\pgfpathlineto{\pgfqpoint{2.124835in}{2.556620in}}%
\pgfpathlineto{\pgfqpoint{2.097049in}{2.631788in}}%
\pgfpathlineto{\pgfqpoint{2.145917in}{2.646487in}}%
\pgfpathlineto{\pgfqpoint{2.194698in}{2.661161in}}%
\pgfpathlineto{\pgfqpoint{2.222480in}{2.586150in}}%
\pgfpathlineto{\pgfqpoint{2.250289in}{2.511084in}}%
\pgfpathlineto{\pgfqpoint{2.201512in}{2.496249in}}%
\pgfpathlineto{\pgfqpoint{2.152647in}{2.481391in}}%
\pgfpathclose%
\pgfusepath{stroke,fill}%
\end{pgfscope}%
\begin{pgfscope}%
\pgfpathrectangle{\pgfqpoint{0.050000in}{0.083252in}}{\pgfqpoint{4.900000in}{4.900000in}}%
\pgfusepath{clip}%
\pgfsetbuttcap%
\pgfsetroundjoin%
\definecolor{currentfill}{rgb}{0.433141,0.409765,0.661053}%
\pgfsetfillcolor{currentfill}%
\pgfsetfillopacity{0.950000}%
\pgfsetlinewidth{0.200750pt}%
\definecolor{currentstroke}{rgb}{0.200000,0.200000,0.200000}%
\pgfsetstrokecolor{currentstroke}%
\pgfsetdash{}{0pt}%
\pgfpathmoveto{\pgfqpoint{2.361788in}{2.210586in}}%
\pgfpathlineto{\pgfqpoint{2.333874in}{2.285689in}}%
\pgfpathlineto{\pgfqpoint{2.305986in}{2.360831in}}%
\pgfpathlineto{\pgfqpoint{2.354672in}{2.375825in}}%
\pgfpathlineto{\pgfqpoint{2.403270in}{2.390799in}}%
\pgfpathlineto{\pgfqpoint{2.431155in}{2.315868in}}%
\pgfpathlineto{\pgfqpoint{2.459065in}{2.241002in}}%
\pgfpathlineto{\pgfqpoint{2.410470in}{2.225802in}}%
\pgfpathlineto{\pgfqpoint{2.361788in}{2.210586in}}%
\pgfpathclose%
\pgfusepath{stroke,fill}%
\end{pgfscope}%
\begin{pgfscope}%
\pgfpathrectangle{\pgfqpoint{0.050000in}{0.083252in}}{\pgfqpoint{4.900000in}{4.900000in}}%
\pgfusepath{clip}%
\pgfsetbuttcap%
\pgfsetroundjoin%
\definecolor{currentfill}{rgb}{0.856794,0.850888,0.914371}%
\pgfsetfillcolor{currentfill}%
\pgfsetfillopacity{0.950000}%
\pgfsetlinewidth{0.200750pt}%
\definecolor{currentstroke}{rgb}{0.200000,0.200000,0.200000}%
\pgfsetstrokecolor{currentstroke}%
\pgfsetdash{}{0pt}%
\pgfpathmoveto{\pgfqpoint{3.086687in}{1.515040in}}%
\pgfpathlineto{\pgfqpoint{3.058210in}{1.572291in}}%
\pgfpathlineto{\pgfqpoint{3.029815in}{1.632500in}}%
\pgfpathlineto{\pgfqpoint{3.077854in}{1.649498in}}%
\pgfpathlineto{\pgfqpoint{3.125809in}{1.666443in}}%
\pgfpathlineto{\pgfqpoint{3.154235in}{1.606508in}}%
\pgfpathlineto{\pgfqpoint{3.182748in}{1.549440in}}%
\pgfpathlineto{\pgfqpoint{3.134759in}{1.532269in}}%
\pgfpathlineto{\pgfqpoint{3.086687in}{1.515040in}}%
\pgfpathclose%
\pgfusepath{stroke,fill}%
\end{pgfscope}%
\begin{pgfscope}%
\pgfpathrectangle{\pgfqpoint{0.050000in}{0.083252in}}{\pgfqpoint{4.900000in}{4.900000in}}%
\pgfusepath{clip}%
\pgfsetbuttcap%
\pgfsetroundjoin%
\definecolor{currentfill}{rgb}{0.594248,0.577516,0.757386}%
\pgfsetfillcolor{currentfill}%
\pgfsetfillopacity{0.950000}%
\pgfsetlinewidth{0.200750pt}%
\definecolor{currentstroke}{rgb}{0.200000,0.200000,0.200000}%
\pgfsetstrokecolor{currentstroke}%
\pgfsetdash{}{0pt}%
\pgfpathmoveto{\pgfqpoint{2.570978in}{1.944016in}}%
\pgfpathlineto{\pgfqpoint{2.542958in}{2.017625in}}%
\pgfpathlineto{\pgfqpoint{2.514967in}{2.091767in}}%
\pgfpathlineto{\pgfqpoint{2.563472in}{2.107248in}}%
\pgfpathlineto{\pgfqpoint{2.611892in}{2.122715in}}%
\pgfpathlineto{\pgfqpoint{2.639883in}{2.048933in}}%
\pgfpathlineto{\pgfqpoint{2.667904in}{1.975731in}}%
\pgfpathlineto{\pgfqpoint{2.619484in}{1.959881in}}%
\pgfpathlineto{\pgfqpoint{2.570978in}{1.944016in}}%
\pgfpathclose%
\pgfusepath{stroke,fill}%
\end{pgfscope}%
\begin{pgfscope}%
\pgfpathrectangle{\pgfqpoint{0.050000in}{0.083252in}}{\pgfqpoint{4.900000in}{4.900000in}}%
\pgfusepath{clip}%
\pgfsetbuttcap%
\pgfsetroundjoin%
\definecolor{currentfill}{rgb}{0.898562,0.894379,0.939346}%
\pgfsetfillcolor{currentfill}%
\pgfsetfillopacity{0.950000}%
\pgfsetlinewidth{0.200750pt}%
\definecolor{currentstroke}{rgb}{0.200000,0.200000,0.200000}%
\pgfsetstrokecolor{currentstroke}%
\pgfsetdash{}{0pt}%
\pgfpathmoveto{\pgfqpoint{3.240072in}{1.444626in}}%
\pgfpathlineto{\pgfqpoint{3.211358in}{1.495438in}}%
\pgfpathlineto{\pgfqpoint{3.182748in}{1.549440in}}%
\pgfpathlineto{\pgfqpoint{3.230653in}{1.566554in}}%
\pgfpathlineto{\pgfqpoint{3.278476in}{1.583611in}}%
\pgfpathlineto{\pgfqpoint{3.307127in}{1.529688in}}%
\pgfpathlineto{\pgfqpoint{3.335888in}{1.478856in}}%
\pgfpathlineto{\pgfqpoint{3.288022in}{1.461768in}}%
\pgfpathlineto{\pgfqpoint{3.240072in}{1.444626in}}%
\pgfpathclose%
\pgfusepath{stroke,fill}%
\end{pgfscope}%
\begin{pgfscope}%
\pgfpathrectangle{\pgfqpoint{0.050000in}{0.083252in}}{\pgfqpoint{4.900000in}{4.900000in}}%
\pgfusepath{clip}%
\pgfsetbuttcap%
\pgfsetroundjoin%
\definecolor{currentfill}{rgb}{0.803091,0.794971,0.882261}%
\pgfsetfillcolor{currentfill}%
\pgfsetfillopacity{0.950000}%
\pgfsetlinewidth{0.200750pt}%
\definecolor{currentstroke}{rgb}{0.200000,0.200000,0.200000}%
\pgfsetstrokecolor{currentstroke}%
\pgfsetdash{}{0pt}%
\pgfpathmoveto{\pgfqpoint{2.933487in}{1.598346in}}%
\pgfpathlineto{\pgfqpoint{2.905191in}{1.661610in}}%
\pgfpathlineto{\pgfqpoint{2.876956in}{1.727299in}}%
\pgfpathlineto{\pgfqpoint{2.925139in}{1.744000in}}%
\pgfpathlineto{\pgfqpoint{2.973238in}{1.760661in}}%
\pgfpathlineto{\pgfqpoint{3.001494in}{1.695394in}}%
\pgfpathlineto{\pgfqpoint{3.029815in}{1.632500in}}%
\pgfpathlineto{\pgfqpoint{2.981693in}{1.615450in}}%
\pgfpathlineto{\pgfqpoint{2.933487in}{1.598346in}}%
\pgfpathclose%
\pgfusepath{stroke,fill}%
\end{pgfscope}%
\begin{pgfscope}%
\pgfpathrectangle{\pgfqpoint{0.050000in}{0.083252in}}{\pgfqpoint{4.900000in}{4.900000in}}%
\pgfusepath{clip}%
\pgfsetbuttcap%
\pgfsetroundjoin%
\definecolor{currentfill}{rgb}{0.267389,0.507190,0.943053}%
\pgfsetfillcolor{currentfill}%
\pgfsetfillopacity{0.950000}%
\pgfsetlinewidth{0.200750pt}%
\definecolor{currentstroke}{rgb}{0.200000,0.200000,0.200000}%
\pgfsetstrokecolor{currentstroke}%
\pgfsetdash{}{0pt}%
\pgfpathmoveto{\pgfqpoint{1.162265in}{3.688457in}}%
\pgfpathlineto{\pgfqpoint{1.134912in}{3.763165in}}%
\pgfpathlineto{\pgfqpoint{1.107585in}{3.837803in}}%
\pgfpathlineto{\pgfqpoint{1.157364in}{3.851420in}}%
\pgfpathlineto{\pgfqpoint{1.207053in}{3.865013in}}%
\pgfpathlineto{\pgfqpoint{1.234378in}{3.790523in}}%
\pgfpathlineto{\pgfqpoint{1.261730in}{3.715961in}}%
\pgfpathlineto{\pgfqpoint{1.212042in}{3.702221in}}%
\pgfpathlineto{\pgfqpoint{1.162265in}{3.688457in}}%
\pgfpathclose%
\pgfusepath{stroke,fill}%
\end{pgfscope}%
\begin{pgfscope}%
\pgfpathrectangle{\pgfqpoint{0.050000in}{0.083252in}}{\pgfqpoint{4.900000in}{4.900000in}}%
\pgfusepath{clip}%
\pgfsetbuttcap%
\pgfsetroundjoin%
\definecolor{currentfill}{rgb}{0.976132,0.975148,0.985729}%
\pgfsetfillcolor{currentfill}%
\pgfsetfillopacity{0.950000}%
\pgfsetlinewidth{0.200750pt}%
\definecolor{currentstroke}{rgb}{0.200000,0.200000,0.200000}%
\pgfsetstrokecolor{currentstroke}%
\pgfsetdash{}{0pt}%
\pgfpathmoveto{\pgfqpoint{3.797577in}{1.332179in}}%
\pgfpathlineto{\pgfqpoint{3.767854in}{1.367577in}}%
\pgfpathlineto{\pgfqpoint{3.738268in}{1.405024in}}%
\pgfpathlineto{\pgfqpoint{3.785658in}{1.421418in}}%
\pgfpathlineto{\pgfqpoint{3.832965in}{1.437777in}}%
\pgfpathlineto{\pgfqpoint{3.862612in}{1.400041in}}%
\pgfpathlineto{\pgfqpoint{3.892402in}{1.364328in}}%
\pgfpathlineto{\pgfqpoint{3.845032in}{1.348268in}}%
\pgfpathlineto{\pgfqpoint{3.797577in}{1.332179in}}%
\pgfpathclose%
\pgfusepath{stroke,fill}%
\end{pgfscope}%
\begin{pgfscope}%
\pgfpathrectangle{\pgfqpoint{0.050000in}{0.083252in}}{\pgfqpoint{4.900000in}{4.900000in}}%
\pgfusepath{clip}%
\pgfsetbuttcap%
\pgfsetroundjoin%
\definecolor{currentfill}{rgb}{0.260746,0.436601,0.849212}%
\pgfsetfillcolor{currentfill}%
\pgfsetfillopacity{0.950000}%
\pgfsetlinewidth{0.200750pt}%
\definecolor{currentstroke}{rgb}{0.200000,0.200000,0.200000}%
\pgfsetstrokecolor{currentstroke}%
\pgfsetdash{}{0pt}%
\pgfpathmoveto{\pgfqpoint{1.371395in}{3.417006in}}%
\pgfpathlineto{\pgfqpoint{1.343939in}{3.491851in}}%
\pgfpathlineto{\pgfqpoint{1.316510in}{3.566626in}}%
\pgfpathlineto{\pgfqpoint{1.366106in}{3.580488in}}%
\pgfpathlineto{\pgfqpoint{1.415614in}{3.594325in}}%
\pgfpathlineto{\pgfqpoint{1.443041in}{3.519698in}}%
\pgfpathlineto{\pgfqpoint{1.470494in}{3.445001in}}%
\pgfpathlineto{\pgfqpoint{1.420989in}{3.431016in}}%
\pgfpathlineto{\pgfqpoint{1.371395in}{3.417006in}}%
\pgfpathclose%
\pgfusepath{stroke,fill}%
\end{pgfscope}%
\begin{pgfscope}%
\pgfpathrectangle{\pgfqpoint{0.050000in}{0.083252in}}{\pgfqpoint{4.900000in}{4.900000in}}%
\pgfusepath{clip}%
\pgfsetbuttcap%
\pgfsetroundjoin%
\definecolor{currentfill}{rgb}{0.254102,0.366013,0.755371}%
\pgfsetfillcolor{currentfill}%
\pgfsetfillopacity{0.950000}%
\pgfsetlinewidth{0.200750pt}%
\definecolor{currentstroke}{rgb}{0.200000,0.200000,0.200000}%
\pgfsetstrokecolor{currentstroke}%
\pgfsetdash{}{0pt}%
\pgfpathmoveto{\pgfqpoint{1.580569in}{3.145497in}}%
\pgfpathlineto{\pgfqpoint{1.553011in}{3.220480in}}%
\pgfpathlineto{\pgfqpoint{1.525479in}{3.295391in}}%
\pgfpathlineto{\pgfqpoint{1.574893in}{3.309499in}}%
\pgfpathlineto{\pgfqpoint{1.624218in}{3.323581in}}%
\pgfpathlineto{\pgfqpoint{1.651747in}{3.248817in}}%
\pgfpathlineto{\pgfqpoint{1.679302in}{3.173983in}}%
\pgfpathlineto{\pgfqpoint{1.629980in}{3.159753in}}%
\pgfpathlineto{\pgfqpoint{1.580569in}{3.145497in}}%
\pgfpathclose%
\pgfusepath{stroke,fill}%
\end{pgfscope}%
\begin{pgfscope}%
\pgfpathrectangle{\pgfqpoint{0.050000in}{0.083252in}}{\pgfqpoint{4.900000in}{4.900000in}}%
\pgfusepath{clip}%
\pgfsetbuttcap%
\pgfsetroundjoin%
\definecolor{currentfill}{rgb}{1.000000,1.000000,1.000000}%
\pgfsetfillcolor{currentfill}%
\pgfsetfillopacity{0.950000}%
\pgfsetlinewidth{0.200750pt}%
\definecolor{currentstroke}{rgb}{0.200000,0.200000,0.200000}%
\pgfsetstrokecolor{currentstroke}%
\pgfsetdash{}{0pt}%
\pgfpathmoveto{\pgfqpoint{4.201989in}{1.295332in}}%
\pgfpathlineto{\pgfqpoint{4.171581in}{1.327596in}}%
\pgfpathlineto{\pgfqpoint{4.141276in}{1.360185in}}%
\pgfpathlineto{\pgfqpoint{4.188286in}{1.375772in}}%
\pgfpathlineto{\pgfqpoint{4.235212in}{1.391332in}}%
\pgfpathlineto{\pgfqpoint{4.265581in}{1.358398in}}%
\pgfpathlineto{\pgfqpoint{4.296052in}{1.325775in}}%
\pgfpathlineto{\pgfqpoint{4.249064in}{1.310567in}}%
\pgfpathlineto{\pgfqpoint{4.201989in}{1.295332in}}%
\pgfpathclose%
\pgfusepath{stroke,fill}%
\end{pgfscope}%
\begin{pgfscope}%
\pgfpathrectangle{\pgfqpoint{0.050000in}{0.083252in}}{\pgfqpoint{4.900000in}{4.900000in}}%
\pgfusepath{clip}%
\pgfsetbuttcap%
\pgfsetroundjoin%
\definecolor{currentfill}{rgb}{0.247459,0.295425,0.661530}%
\pgfsetfillcolor{currentfill}%
\pgfsetfillopacity{0.950000}%
\pgfsetlinewidth{0.200750pt}%
\definecolor{currentstroke}{rgb}{0.200000,0.200000,0.200000}%
\pgfsetstrokecolor{currentstroke}%
\pgfsetdash{}{0pt}%
\pgfpathmoveto{\pgfqpoint{1.789786in}{2.873932in}}%
\pgfpathlineto{\pgfqpoint{1.762126in}{2.949052in}}%
\pgfpathlineto{\pgfqpoint{1.734492in}{3.024100in}}%
\pgfpathlineto{\pgfqpoint{1.783722in}{3.038452in}}%
\pgfpathlineto{\pgfqpoint{1.832865in}{3.052779in}}%
\pgfpathlineto{\pgfqpoint{1.860496in}{2.977879in}}%
\pgfpathlineto{\pgfqpoint{1.888154in}{2.902908in}}%
\pgfpathlineto{\pgfqpoint{1.839014in}{2.888433in}}%
\pgfpathlineto{\pgfqpoint{1.789786in}{2.873932in}}%
\pgfpathclose%
\pgfusepath{stroke,fill}%
\end{pgfscope}%
\begin{pgfscope}%
\pgfpathrectangle{\pgfqpoint{0.050000in}{0.083252in}}{\pgfqpoint{4.900000in}{4.900000in}}%
\pgfusepath{clip}%
\pgfsetbuttcap%
\pgfsetroundjoin%
\definecolor{currentfill}{rgb}{0.934364,0.931657,0.960754}%
\pgfsetfillcolor{currentfill}%
\pgfsetfillopacity{0.950000}%
\pgfsetlinewidth{0.200750pt}%
\definecolor{currentstroke}{rgb}{0.200000,0.200000,0.200000}%
\pgfsetstrokecolor{currentstroke}%
\pgfsetdash{}{0pt}%
\pgfpathmoveto{\pgfqpoint{3.393767in}{1.386536in}}%
\pgfpathlineto{\pgfqpoint{3.364766in}{1.431154in}}%
\pgfpathlineto{\pgfqpoint{3.335888in}{1.478856in}}%
\pgfpathlineto{\pgfqpoint{3.383670in}{1.495890in}}%
\pgfpathlineto{\pgfqpoint{3.431369in}{1.512872in}}%
\pgfpathlineto{\pgfqpoint{3.460298in}{1.465066in}}%
\pgfpathlineto{\pgfqpoint{3.489354in}{1.420262in}}%
\pgfpathlineto{\pgfqpoint{3.441603in}{1.403421in}}%
\pgfpathlineto{\pgfqpoint{3.393767in}{1.386536in}}%
\pgfpathclose%
\pgfusepath{stroke,fill}%
\end{pgfscope}%
\begin{pgfscope}%
\pgfpathrectangle{\pgfqpoint{0.050000in}{0.083252in}}{\pgfqpoint{4.900000in}{4.900000in}}%
\pgfusepath{clip}%
\pgfsetbuttcap%
\pgfsetroundjoin%
\definecolor{currentfill}{rgb}{0.737455,0.726628,0.843014}%
\pgfsetfillcolor{currentfill}%
\pgfsetfillopacity{0.950000}%
\pgfsetlinewidth{0.200750pt}%
\definecolor{currentstroke}{rgb}{0.200000,0.200000,0.200000}%
\pgfsetstrokecolor{currentstroke}%
\pgfsetdash{}{0pt}%
\pgfpathmoveto{\pgfqpoint{2.780337in}{1.693778in}}%
\pgfpathlineto{\pgfqpoint{2.752169in}{1.762020in}}%
\pgfpathlineto{\pgfqpoint{2.724044in}{1.831986in}}%
\pgfpathlineto{\pgfqpoint{2.772383in}{1.848277in}}%
\pgfpathlineto{\pgfqpoint{2.820637in}{1.864542in}}%
\pgfpathlineto{\pgfqpoint{2.848773in}{1.795060in}}%
\pgfpathlineto{\pgfqpoint{2.876956in}{1.727299in}}%
\pgfpathlineto{\pgfqpoint{2.828688in}{1.710559in}}%
\pgfpathlineto{\pgfqpoint{2.780337in}{1.693778in}}%
\pgfpathclose%
\pgfusepath{stroke,fill}%
\end{pgfscope}%
\begin{pgfscope}%
\pgfpathrectangle{\pgfqpoint{0.050000in}{0.083252in}}{\pgfqpoint{4.900000in}{4.900000in}}%
\pgfusepath{clip}%
\pgfsetbuttcap%
\pgfsetroundjoin%
\definecolor{currentfill}{rgb}{0.241061,0.227451,0.571165}%
\pgfsetfillcolor{currentfill}%
\pgfsetfillopacity{0.950000}%
\pgfsetlinewidth{0.200750pt}%
\definecolor{currentstroke}{rgb}{0.200000,0.200000,0.200000}%
\pgfsetstrokecolor{currentstroke}%
\pgfsetdash{}{0pt}%
\pgfpathmoveto{\pgfqpoint{1.999047in}{2.602314in}}%
\pgfpathlineto{\pgfqpoint{1.971284in}{2.677569in}}%
\pgfpathlineto{\pgfqpoint{1.943548in}{2.752753in}}%
\pgfpathlineto{\pgfqpoint{1.992596in}{2.767351in}}%
\pgfpathlineto{\pgfqpoint{2.041556in}{2.781923in}}%
\pgfpathlineto{\pgfqpoint{2.069289in}{2.706890in}}%
\pgfpathlineto{\pgfqpoint{2.097049in}{2.631788in}}%
\pgfpathlineto{\pgfqpoint{2.048092in}{2.617064in}}%
\pgfpathlineto{\pgfqpoint{1.999047in}{2.602314in}}%
\pgfpathclose%
\pgfusepath{stroke,fill}%
\end{pgfscope}%
\begin{pgfscope}%
\pgfpathrectangle{\pgfqpoint{0.050000in}{0.083252in}}{\pgfqpoint{4.900000in}{4.900000in}}%
\pgfusepath{clip}%
\pgfsetbuttcap%
\pgfsetroundjoin%
\definecolor{currentfill}{rgb}{0.355571,0.328997,0.614671}%
\pgfsetfillcolor{currentfill}%
\pgfsetfillopacity{0.950000}%
\pgfsetlinewidth{0.200750pt}%
\definecolor{currentstroke}{rgb}{0.200000,0.200000,0.200000}%
\pgfsetstrokecolor{currentstroke}%
\pgfsetdash{}{0pt}%
\pgfpathmoveto{\pgfqpoint{2.208351in}{2.330783in}}%
\pgfpathlineto{\pgfqpoint{2.180486in}{2.406107in}}%
\pgfpathlineto{\pgfqpoint{2.152647in}{2.481391in}}%
\pgfpathlineto{\pgfqpoint{2.201512in}{2.496249in}}%
\pgfpathlineto{\pgfqpoint{2.250289in}{2.511084in}}%
\pgfpathlineto{\pgfqpoint{2.278124in}{2.435972in}}%
\pgfpathlineto{\pgfqpoint{2.305986in}{2.360831in}}%
\pgfpathlineto{\pgfqpoint{2.257212in}{2.345817in}}%
\pgfpathlineto{\pgfqpoint{2.208351in}{2.330783in}}%
\pgfpathclose%
\pgfusepath{stroke,fill}%
\end{pgfscope}%
\begin{pgfscope}%
\pgfpathrectangle{\pgfqpoint{0.050000in}{0.083252in}}{\pgfqpoint{4.900000in}{4.900000in}}%
\pgfusepath{clip}%
\pgfsetbuttcap%
\pgfsetroundjoin%
\definecolor{currentfill}{rgb}{0.516678,0.496747,0.711003}%
\pgfsetfillcolor{currentfill}%
\pgfsetfillopacity{0.950000}%
\pgfsetlinewidth{0.200750pt}%
\definecolor{currentstroke}{rgb}{0.200000,0.200000,0.200000}%
\pgfsetstrokecolor{currentstroke}%
\pgfsetdash{}{0pt}%
\pgfpathmoveto{\pgfqpoint{2.417694in}{2.060769in}}%
\pgfpathlineto{\pgfqpoint{2.389728in}{2.135583in}}%
\pgfpathlineto{\pgfqpoint{2.361788in}{2.210586in}}%
\pgfpathlineto{\pgfqpoint{2.410470in}{2.225802in}}%
\pgfpathlineto{\pgfqpoint{2.459065in}{2.241002in}}%
\pgfpathlineto{\pgfqpoint{2.487002in}{2.166270in}}%
\pgfpathlineto{\pgfqpoint{2.514967in}{2.091767in}}%
\pgfpathlineto{\pgfqpoint{2.466374in}{2.076274in}}%
\pgfpathlineto{\pgfqpoint{2.417694in}{2.060769in}}%
\pgfpathclose%
\pgfusepath{stroke,fill}%
\end{pgfscope}%
\begin{pgfscope}%
\pgfpathrectangle{\pgfqpoint{0.050000in}{0.083252in}}{\pgfqpoint{4.900000in}{4.900000in}}%
\pgfusepath{clip}%
\pgfsetbuttcap%
\pgfsetroundjoin%
\definecolor{currentfill}{rgb}{0.665852,0.652072,0.800200}%
\pgfsetfillcolor{currentfill}%
\pgfsetfillopacity{0.950000}%
\pgfsetlinewidth{0.200750pt}%
\definecolor{currentstroke}{rgb}{0.200000,0.200000,0.200000}%
\pgfsetstrokecolor{currentstroke}%
\pgfsetdash{}{0pt}%
\pgfpathmoveto{\pgfqpoint{2.627109in}{1.799331in}}%
\pgfpathlineto{\pgfqpoint{2.599028in}{1.871160in}}%
\pgfpathlineto{\pgfqpoint{2.570978in}{1.944016in}}%
\pgfpathlineto{\pgfqpoint{2.619484in}{1.959881in}}%
\pgfpathlineto{\pgfqpoint{2.667904in}{1.975731in}}%
\pgfpathlineto{\pgfqpoint{2.695956in}{1.903326in}}%
\pgfpathlineto{\pgfqpoint{2.724044in}{1.831986in}}%
\pgfpathlineto{\pgfqpoint{2.675619in}{1.815670in}}%
\pgfpathlineto{\pgfqpoint{2.627109in}{1.799331in}}%
\pgfpathclose%
\pgfusepath{stroke,fill}%
\end{pgfscope}%
\begin{pgfscope}%
\pgfpathrectangle{\pgfqpoint{0.050000in}{0.083252in}}{\pgfqpoint{4.900000in}{4.900000in}}%
\pgfusepath{clip}%
\pgfsetbuttcap%
\pgfsetroundjoin%
\definecolor{currentfill}{rgb}{0.263945,0.470588,0.894394}%
\pgfsetfillcolor{currentfill}%
\pgfsetfillopacity{0.950000}%
\pgfsetlinewidth{0.200750pt}%
\definecolor{currentstroke}{rgb}{0.200000,0.200000,0.200000}%
\pgfsetstrokecolor{currentstroke}%
\pgfsetdash{}{0pt}%
\pgfpathmoveto{\pgfqpoint{1.217049in}{3.538826in}}%
\pgfpathlineto{\pgfqpoint{1.189644in}{3.613677in}}%
\pgfpathlineto{\pgfqpoint{1.162265in}{3.688457in}}%
\pgfpathlineto{\pgfqpoint{1.212042in}{3.702221in}}%
\pgfpathlineto{\pgfqpoint{1.261730in}{3.715961in}}%
\pgfpathlineto{\pgfqpoint{1.289107in}{3.641329in}}%
\pgfpathlineto{\pgfqpoint{1.316510in}{3.566626in}}%
\pgfpathlineto{\pgfqpoint{1.266824in}{3.552738in}}%
\pgfpathlineto{\pgfqpoint{1.217049in}{3.538826in}}%
\pgfpathclose%
\pgfusepath{stroke,fill}%
\end{pgfscope}%
\begin{pgfscope}%
\pgfpathrectangle{\pgfqpoint{0.050000in}{0.083252in}}{\pgfqpoint{4.900000in}{4.900000in}}%
\pgfusepath{clip}%
\pgfsetbuttcap%
\pgfsetroundjoin%
\definecolor{currentfill}{rgb}{0.958231,0.956509,0.975025}%
\pgfsetfillcolor{currentfill}%
\pgfsetfillopacity{0.950000}%
\pgfsetlinewidth{0.200750pt}%
\definecolor{currentstroke}{rgb}{0.200000,0.200000,0.200000}%
\pgfsetstrokecolor{currentstroke}%
\pgfsetdash{}{0pt}%
\pgfpathmoveto{\pgfqpoint{3.547863in}{1.339102in}}%
\pgfpathlineto{\pgfqpoint{3.518542in}{1.378340in}}%
\pgfpathlineto{\pgfqpoint{3.489354in}{1.420262in}}%
\pgfpathlineto{\pgfqpoint{3.537021in}{1.437059in}}%
\pgfpathlineto{\pgfqpoint{3.584605in}{1.453813in}}%
\pgfpathlineto{\pgfqpoint{3.613851in}{1.411655in}}%
\pgfpathlineto{\pgfqpoint{3.643235in}{1.372131in}}%
\pgfpathlineto{\pgfqpoint{3.595592in}{1.355634in}}%
\pgfpathlineto{\pgfqpoint{3.547863in}{1.339102in}}%
\pgfpathclose%
\pgfusepath{stroke,fill}%
\end{pgfscope}%
\begin{pgfscope}%
\pgfpathrectangle{\pgfqpoint{0.050000in}{0.083252in}}{\pgfqpoint{4.900000in}{4.900000in}}%
\pgfusepath{clip}%
\pgfsetbuttcap%
\pgfsetroundjoin%
\definecolor{currentfill}{rgb}{0.988066,0.987574,0.992864}%
\pgfsetfillcolor{currentfill}%
\pgfsetfillopacity{0.950000}%
\pgfsetlinewidth{0.200750pt}%
\definecolor{currentstroke}{rgb}{0.200000,0.200000,0.200000}%
\pgfsetstrokecolor{currentstroke}%
\pgfsetdash{}{0pt}%
\pgfpathmoveto{\pgfqpoint{3.952380in}{1.297554in}}%
\pgfpathlineto{\pgfqpoint{3.922327in}{1.330296in}}%
\pgfpathlineto{\pgfqpoint{3.892402in}{1.364328in}}%
\pgfpathlineto{\pgfqpoint{3.939687in}{1.380358in}}%
\pgfpathlineto{\pgfqpoint{3.986887in}{1.396358in}}%
\pgfpathlineto{\pgfqpoint{4.016878in}{1.362003in}}%
\pgfpathlineto{\pgfqpoint{4.046999in}{1.328925in}}%
\pgfpathlineto{\pgfqpoint{3.999732in}{1.313253in}}%
\pgfpathlineto{\pgfqpoint{3.952380in}{1.297554in}}%
\pgfpathclose%
\pgfusepath{stroke,fill}%
\end{pgfscope}%
\begin{pgfscope}%
\pgfpathrectangle{\pgfqpoint{0.050000in}{0.083252in}}{\pgfqpoint{4.900000in}{4.900000in}}%
\pgfusepath{clip}%
\pgfsetbuttcap%
\pgfsetroundjoin%
\definecolor{currentfill}{rgb}{0.257301,0.400000,0.800554}%
\pgfsetfillcolor{currentfill}%
\pgfsetfillopacity{0.950000}%
\pgfsetlinewidth{0.200750pt}%
\definecolor{currentstroke}{rgb}{0.200000,0.200000,0.200000}%
\pgfsetstrokecolor{currentstroke}%
\pgfsetdash{}{0pt}%
\pgfpathmoveto{\pgfqpoint{1.426384in}{3.267100in}}%
\pgfpathlineto{\pgfqpoint{1.398876in}{3.342088in}}%
\pgfpathlineto{\pgfqpoint{1.371395in}{3.417006in}}%
\pgfpathlineto{\pgfqpoint{1.420989in}{3.431016in}}%
\pgfpathlineto{\pgfqpoint{1.470494in}{3.445001in}}%
\pgfpathlineto{\pgfqpoint{1.497973in}{3.370232in}}%
\pgfpathlineto{\pgfqpoint{1.525479in}{3.295391in}}%
\pgfpathlineto{\pgfqpoint{1.475976in}{3.281258in}}%
\pgfpathlineto{\pgfqpoint{1.426384in}{3.267100in}}%
\pgfpathclose%
\pgfusepath{stroke,fill}%
\end{pgfscope}%
\begin{pgfscope}%
\pgfpathrectangle{\pgfqpoint{0.050000in}{0.083252in}}{\pgfqpoint{4.900000in}{4.900000in}}%
\pgfusepath{clip}%
\pgfsetbuttcap%
\pgfsetroundjoin%
\definecolor{currentfill}{rgb}{0.269604,0.530719,0.974333}%
\pgfsetfillcolor{currentfill}%
\pgfsetfillopacity{0.950000}%
\pgfsetlinewidth{0.200750pt}%
\definecolor{currentstroke}{rgb}{0.200000,0.200000,0.200000}%
\pgfsetstrokecolor{currentstroke}%
\pgfsetdash{}{0pt}%
\pgfpathmoveto{\pgfqpoint{1.007757in}{3.810494in}}%
\pgfpathlineto{\pgfqpoint{0.980455in}{3.885208in}}%
\pgfpathlineto{\pgfqpoint{1.030414in}{3.898801in}}%
\pgfpathlineto{\pgfqpoint{1.080284in}{3.912369in}}%
\pgfpathlineto{\pgfqpoint{1.107585in}{3.837803in}}%
\pgfpathlineto{\pgfqpoint{1.057716in}{3.824161in}}%
\pgfpathlineto{\pgfqpoint{1.007757in}{3.810494in}}%
\pgfpathclose%
\pgfusepath{stroke,fill}%
\end{pgfscope}%
\begin{pgfscope}%
\pgfpathrectangle{\pgfqpoint{0.050000in}{0.083252in}}{\pgfqpoint{4.900000in}{4.900000in}}%
\pgfusepath{clip}%
\pgfsetbuttcap%
\pgfsetroundjoin%
\definecolor{currentfill}{rgb}{0.250903,0.332026,0.710188}%
\pgfsetfillcolor{currentfill}%
\pgfsetfillopacity{0.950000}%
\pgfsetlinewidth{0.200750pt}%
\definecolor{currentstroke}{rgb}{0.200000,0.200000,0.200000}%
\pgfsetstrokecolor{currentstroke}%
\pgfsetdash{}{0pt}%
\pgfpathmoveto{\pgfqpoint{1.635763in}{2.995317in}}%
\pgfpathlineto{\pgfqpoint{1.608153in}{3.070443in}}%
\pgfpathlineto{\pgfqpoint{1.580569in}{3.145497in}}%
\pgfpathlineto{\pgfqpoint{1.629980in}{3.159753in}}%
\pgfpathlineto{\pgfqpoint{1.679302in}{3.173983in}}%
\pgfpathlineto{\pgfqpoint{1.706884in}{3.099077in}}%
\pgfpathlineto{\pgfqpoint{1.734492in}{3.024100in}}%
\pgfpathlineto{\pgfqpoint{1.685172in}{3.009721in}}%
\pgfpathlineto{\pgfqpoint{1.635763in}{2.995317in}}%
\pgfpathclose%
\pgfusepath{stroke,fill}%
\end{pgfscope}%
\begin{pgfscope}%
\pgfpathrectangle{\pgfqpoint{0.050000in}{0.083252in}}{\pgfqpoint{4.900000in}{4.900000in}}%
\pgfusepath{clip}%
\pgfsetbuttcap%
\pgfsetroundjoin%
\definecolor{currentfill}{rgb}{0.244260,0.261438,0.616348}%
\pgfsetfillcolor{currentfill}%
\pgfsetfillopacity{0.950000}%
\pgfsetlinewidth{0.200750pt}%
\definecolor{currentstroke}{rgb}{0.200000,0.200000,0.200000}%
\pgfsetstrokecolor{currentstroke}%
\pgfsetdash{}{0pt}%
\pgfpathmoveto{\pgfqpoint{1.845186in}{2.723477in}}%
\pgfpathlineto{\pgfqpoint{1.817473in}{2.798740in}}%
\pgfpathlineto{\pgfqpoint{1.789786in}{2.873932in}}%
\pgfpathlineto{\pgfqpoint{1.839014in}{2.888433in}}%
\pgfpathlineto{\pgfqpoint{1.888154in}{2.902908in}}%
\pgfpathlineto{\pgfqpoint{1.915837in}{2.827866in}}%
\pgfpathlineto{\pgfqpoint{1.943548in}{2.752753in}}%
\pgfpathlineto{\pgfqpoint{1.894411in}{2.738128in}}%
\pgfpathlineto{\pgfqpoint{1.845186in}{2.723477in}}%
\pgfpathclose%
\pgfusepath{stroke,fill}%
\end{pgfscope}%
\begin{pgfscope}%
\pgfpathrectangle{\pgfqpoint{0.050000in}{0.083252in}}{\pgfqpoint{4.900000in}{4.900000in}}%
\pgfusepath{clip}%
\pgfsetbuttcap%
\pgfsetroundjoin%
\definecolor{currentfill}{rgb}{0.278001,0.248228,0.568289}%
\pgfsetfillcolor{currentfill}%
\pgfsetfillopacity{0.950000}%
\pgfsetlinewidth{0.200750pt}%
\definecolor{currentstroke}{rgb}{0.200000,0.200000,0.200000}%
\pgfsetstrokecolor{currentstroke}%
\pgfsetdash{}{0pt}%
\pgfpathmoveto{\pgfqpoint{2.054652in}{2.451602in}}%
\pgfpathlineto{\pgfqpoint{2.026837in}{2.526991in}}%
\pgfpathlineto{\pgfqpoint{1.999047in}{2.602314in}}%
\pgfpathlineto{\pgfqpoint{2.048092in}{2.617064in}}%
\pgfpathlineto{\pgfqpoint{2.097049in}{2.631788in}}%
\pgfpathlineto{\pgfqpoint{2.124835in}{2.556620in}}%
\pgfpathlineto{\pgfqpoint{2.152647in}{2.481391in}}%
\pgfpathlineto{\pgfqpoint{2.103694in}{2.466509in}}%
\pgfpathlineto{\pgfqpoint{2.054652in}{2.451602in}}%
\pgfpathclose%
\pgfusepath{stroke,fill}%
\end{pgfscope}%
\begin{pgfscope}%
\pgfpathrectangle{\pgfqpoint{0.050000in}{0.083252in}}{\pgfqpoint{4.900000in}{4.900000in}}%
\pgfusepath{clip}%
\pgfsetbuttcap%
\pgfsetroundjoin%
\definecolor{currentfill}{rgb}{0.439108,0.415978,0.664621}%
\pgfsetfillcolor{currentfill}%
\pgfsetfillopacity{0.950000}%
\pgfsetlinewidth{0.200750pt}%
\definecolor{currentstroke}{rgb}{0.200000,0.200000,0.200000}%
\pgfsetstrokecolor{currentstroke}%
\pgfsetdash{}{0pt}%
\pgfpathmoveto{\pgfqpoint{2.264160in}{2.180108in}}%
\pgfpathlineto{\pgfqpoint{2.236242in}{2.255438in}}%
\pgfpathlineto{\pgfqpoint{2.208351in}{2.330783in}}%
\pgfpathlineto{\pgfqpoint{2.257212in}{2.345817in}}%
\pgfpathlineto{\pgfqpoint{2.305986in}{2.360831in}}%
\pgfpathlineto{\pgfqpoint{2.333874in}{2.285689in}}%
\pgfpathlineto{\pgfqpoint{2.361788in}{2.210586in}}%
\pgfpathlineto{\pgfqpoint{2.313018in}{2.195355in}}%
\pgfpathlineto{\pgfqpoint{2.264160in}{2.180108in}}%
\pgfpathclose%
\pgfusepath{stroke,fill}%
\end{pgfscope}%
\begin{pgfscope}%
\pgfpathrectangle{\pgfqpoint{0.050000in}{0.083252in}}{\pgfqpoint{4.900000in}{4.900000in}}%
\pgfusepath{clip}%
\pgfsetbuttcap%
\pgfsetroundjoin%
\definecolor{currentfill}{rgb}{0.856794,0.850888,0.914371}%
\pgfsetfillcolor{currentfill}%
\pgfsetfillopacity{0.950000}%
\pgfsetlinewidth{0.200750pt}%
\definecolor{currentstroke}{rgb}{0.200000,0.200000,0.200000}%
\pgfsetstrokecolor{currentstroke}%
\pgfsetdash{}{0pt}%
\pgfpathmoveto{\pgfqpoint{2.990290in}{1.480399in}}%
\pgfpathlineto{\pgfqpoint{2.961850in}{1.537844in}}%
\pgfpathlineto{\pgfqpoint{2.933487in}{1.598346in}}%
\pgfpathlineto{\pgfqpoint{2.981693in}{1.615450in}}%
\pgfpathlineto{\pgfqpoint{3.029815in}{1.632500in}}%
\pgfpathlineto{\pgfqpoint{3.058210in}{1.572291in}}%
\pgfpathlineto{\pgfqpoint{3.086687in}{1.515040in}}%
\pgfpathlineto{\pgfqpoint{3.038530in}{1.497750in}}%
\pgfpathlineto{\pgfqpoint{2.990290in}{1.480399in}}%
\pgfpathclose%
\pgfusepath{stroke,fill}%
\end{pgfscope}%
\begin{pgfscope}%
\pgfpathrectangle{\pgfqpoint{0.050000in}{0.083252in}}{\pgfqpoint{4.900000in}{4.900000in}}%
\pgfusepath{clip}%
\pgfsetbuttcap%
\pgfsetroundjoin%
\definecolor{currentfill}{rgb}{0.594248,0.577516,0.757386}%
\pgfsetfillcolor{currentfill}%
\pgfsetfillopacity{0.950000}%
\pgfsetlinewidth{0.200750pt}%
\definecolor{currentstroke}{rgb}{0.200000,0.200000,0.200000}%
\pgfsetstrokecolor{currentstroke}%
\pgfsetdash{}{0pt}%
\pgfpathmoveto{\pgfqpoint{2.473708in}{1.912247in}}%
\pgfpathlineto{\pgfqpoint{2.445688in}{1.986268in}}%
\pgfpathlineto{\pgfqpoint{2.417694in}{2.060769in}}%
\pgfpathlineto{\pgfqpoint{2.466374in}{2.076274in}}%
\pgfpathlineto{\pgfqpoint{2.514967in}{2.091767in}}%
\pgfpathlineto{\pgfqpoint{2.542958in}{2.017625in}}%
\pgfpathlineto{\pgfqpoint{2.570978in}{1.944016in}}%
\pgfpathlineto{\pgfqpoint{2.522386in}{1.928137in}}%
\pgfpathlineto{\pgfqpoint{2.473708in}{1.912247in}}%
\pgfpathclose%
\pgfusepath{stroke,fill}%
\end{pgfscope}%
\begin{pgfscope}%
\pgfpathrectangle{\pgfqpoint{0.050000in}{0.083252in}}{\pgfqpoint{4.900000in}{4.900000in}}%
\pgfusepath{clip}%
\pgfsetbuttcap%
\pgfsetroundjoin%
\definecolor{currentfill}{rgb}{0.904529,0.900592,0.942914}%
\pgfsetfillcolor{currentfill}%
\pgfsetfillopacity{0.950000}%
\pgfsetlinewidth{0.200750pt}%
\definecolor{currentstroke}{rgb}{0.200000,0.200000,0.200000}%
\pgfsetstrokecolor{currentstroke}%
\pgfsetdash{}{0pt}%
\pgfpathmoveto{\pgfqpoint{3.143921in}{1.410176in}}%
\pgfpathlineto{\pgfqpoint{3.115254in}{1.460958in}}%
\pgfpathlineto{\pgfqpoint{3.086687in}{1.515040in}}%
\pgfpathlineto{\pgfqpoint{3.134759in}{1.532269in}}%
\pgfpathlineto{\pgfqpoint{3.182748in}{1.549440in}}%
\pgfpathlineto{\pgfqpoint{3.211358in}{1.495438in}}%
\pgfpathlineto{\pgfqpoint{3.240072in}{1.444626in}}%
\pgfpathlineto{\pgfqpoint{3.192039in}{1.427429in}}%
\pgfpathlineto{\pgfqpoint{3.143921in}{1.410176in}}%
\pgfpathclose%
\pgfusepath{stroke,fill}%
\end{pgfscope}%
\begin{pgfscope}%
\pgfpathrectangle{\pgfqpoint{0.050000in}{0.083252in}}{\pgfqpoint{4.900000in}{4.900000in}}%
\pgfusepath{clip}%
\pgfsetbuttcap%
\pgfsetroundjoin%
\definecolor{currentfill}{rgb}{0.803091,0.794971,0.882261}%
\pgfsetfillcolor{currentfill}%
\pgfsetfillopacity{0.950000}%
\pgfsetlinewidth{0.200750pt}%
\definecolor{currentstroke}{rgb}{0.200000,0.200000,0.200000}%
\pgfsetstrokecolor{currentstroke}%
\pgfsetdash{}{0pt}%
\pgfpathmoveto{\pgfqpoint{2.836823in}{1.563973in}}%
\pgfpathlineto{\pgfqpoint{2.808552in}{1.627635in}}%
\pgfpathlineto{\pgfqpoint{2.780337in}{1.693778in}}%
\pgfpathlineto{\pgfqpoint{2.828688in}{1.710559in}}%
\pgfpathlineto{\pgfqpoint{2.876956in}{1.727299in}}%
\pgfpathlineto{\pgfqpoint{2.905191in}{1.661610in}}%
\pgfpathlineto{\pgfqpoint{2.933487in}{1.598346in}}%
\pgfpathlineto{\pgfqpoint{2.885197in}{1.581188in}}%
\pgfpathlineto{\pgfqpoint{2.836823in}{1.563973in}}%
\pgfpathclose%
\pgfusepath{stroke,fill}%
\end{pgfscope}%
\begin{pgfscope}%
\pgfpathrectangle{\pgfqpoint{0.050000in}{0.083252in}}{\pgfqpoint{4.900000in}{4.900000in}}%
\pgfusepath{clip}%
\pgfsetbuttcap%
\pgfsetroundjoin%
\definecolor{currentfill}{rgb}{0.267143,0.504575,0.939577}%
\pgfsetfillcolor{currentfill}%
\pgfsetfillopacity{0.950000}%
\pgfsetlinewidth{0.200750pt}%
\definecolor{currentstroke}{rgb}{0.200000,0.200000,0.200000}%
\pgfsetstrokecolor{currentstroke}%
\pgfsetdash{}{0pt}%
\pgfpathmoveto{\pgfqpoint{1.062441in}{3.660852in}}%
\pgfpathlineto{\pgfqpoint{1.035086in}{3.735709in}}%
\pgfpathlineto{\pgfqpoint{1.007757in}{3.810494in}}%
\pgfpathlineto{\pgfqpoint{1.057716in}{3.824161in}}%
\pgfpathlineto{\pgfqpoint{1.107585in}{3.837803in}}%
\pgfpathlineto{\pgfqpoint{1.134912in}{3.763165in}}%
\pgfpathlineto{\pgfqpoint{1.162265in}{3.688457in}}%
\pgfpathlineto{\pgfqpoint{1.112398in}{3.674667in}}%
\pgfpathlineto{\pgfqpoint{1.062441in}{3.660852in}}%
\pgfpathclose%
\pgfusepath{stroke,fill}%
\end{pgfscope}%
\begin{pgfscope}%
\pgfpathrectangle{\pgfqpoint{0.050000in}{0.083252in}}{\pgfqpoint{4.900000in}{4.900000in}}%
\pgfusepath{clip}%
\pgfsetbuttcap%
\pgfsetroundjoin%
\definecolor{currentfill}{rgb}{0.976132,0.975148,0.985729}%
\pgfsetfillcolor{currentfill}%
\pgfsetfillopacity{0.950000}%
\pgfsetlinewidth{0.200750pt}%
\definecolor{currentstroke}{rgb}{0.200000,0.200000,0.200000}%
\pgfsetstrokecolor{currentstroke}%
\pgfsetdash{}{0pt}%
\pgfpathmoveto{\pgfqpoint{3.702412in}{1.299917in}}%
\pgfpathlineto{\pgfqpoint{3.672756in}{1.334989in}}%
\pgfpathlineto{\pgfqpoint{3.643235in}{1.372131in}}%
\pgfpathlineto{\pgfqpoint{3.690794in}{1.388595in}}%
\pgfpathlineto{\pgfqpoint{3.738268in}{1.405024in}}%
\pgfpathlineto{\pgfqpoint{3.767854in}{1.367577in}}%
\pgfpathlineto{\pgfqpoint{3.797577in}{1.332179in}}%
\pgfpathlineto{\pgfqpoint{3.750038in}{1.316062in}}%
\pgfpathlineto{\pgfqpoint{3.702412in}{1.299917in}}%
\pgfpathclose%
\pgfusepath{stroke,fill}%
\end{pgfscope}%
\begin{pgfscope}%
\pgfpathrectangle{\pgfqpoint{0.050000in}{0.083252in}}{\pgfqpoint{4.900000in}{4.900000in}}%
\pgfusepath{clip}%
\pgfsetbuttcap%
\pgfsetroundjoin%
\definecolor{currentfill}{rgb}{0.260746,0.436601,0.849212}%
\pgfsetfillcolor{currentfill}%
\pgfsetfillopacity{0.950000}%
\pgfsetlinewidth{0.200750pt}%
\definecolor{currentstroke}{rgb}{0.200000,0.200000,0.200000}%
\pgfsetstrokecolor{currentstroke}%
\pgfsetdash{}{0pt}%
\pgfpathmoveto{\pgfqpoint{1.271938in}{3.388909in}}%
\pgfpathlineto{\pgfqpoint{1.244480in}{3.463903in}}%
\pgfpathlineto{\pgfqpoint{1.217049in}{3.538826in}}%
\pgfpathlineto{\pgfqpoint{1.266824in}{3.552738in}}%
\pgfpathlineto{\pgfqpoint{1.316510in}{3.566626in}}%
\pgfpathlineto{\pgfqpoint{1.343939in}{3.491851in}}%
\pgfpathlineto{\pgfqpoint{1.371395in}{3.417006in}}%
\pgfpathlineto{\pgfqpoint{1.321711in}{3.402970in}}%
\pgfpathlineto{\pgfqpoint{1.271938in}{3.388909in}}%
\pgfpathclose%
\pgfusepath{stroke,fill}%
\end{pgfscope}%
\begin{pgfscope}%
\pgfpathrectangle{\pgfqpoint{0.050000in}{0.083252in}}{\pgfqpoint{4.900000in}{4.900000in}}%
\pgfusepath{clip}%
\pgfsetbuttcap%
\pgfsetroundjoin%
\definecolor{currentfill}{rgb}{0.254102,0.366013,0.755371}%
\pgfsetfillcolor{currentfill}%
\pgfsetfillopacity{0.950000}%
\pgfsetlinewidth{0.200750pt}%
\definecolor{currentstroke}{rgb}{0.200000,0.200000,0.200000}%
\pgfsetstrokecolor{currentstroke}%
\pgfsetdash{}{0pt}%
\pgfpathmoveto{\pgfqpoint{1.481479in}{3.116908in}}%
\pgfpathlineto{\pgfqpoint{1.453918in}{3.192040in}}%
\pgfpathlineto{\pgfqpoint{1.426384in}{3.267100in}}%
\pgfpathlineto{\pgfqpoint{1.475976in}{3.281258in}}%
\pgfpathlineto{\pgfqpoint{1.525479in}{3.295391in}}%
\pgfpathlineto{\pgfqpoint{1.553011in}{3.220480in}}%
\pgfpathlineto{\pgfqpoint{1.580569in}{3.145497in}}%
\pgfpathlineto{\pgfqpoint{1.531068in}{3.131216in}}%
\pgfpathlineto{\pgfqpoint{1.481479in}{3.116908in}}%
\pgfpathclose%
\pgfusepath{stroke,fill}%
\end{pgfscope}%
\begin{pgfscope}%
\pgfpathrectangle{\pgfqpoint{0.050000in}{0.083252in}}{\pgfqpoint{4.900000in}{4.900000in}}%
\pgfusepath{clip}%
\pgfsetbuttcap%
\pgfsetroundjoin%
\definecolor{currentfill}{rgb}{0.934364,0.931657,0.960754}%
\pgfsetfillcolor{currentfill}%
\pgfsetfillopacity{0.950000}%
\pgfsetlinewidth{0.200750pt}%
\definecolor{currentstroke}{rgb}{0.200000,0.200000,0.200000}%
\pgfsetstrokecolor{currentstroke}%
\pgfsetdash{}{0pt}%
\pgfpathmoveto{\pgfqpoint{3.297843in}{1.352632in}}%
\pgfpathlineto{\pgfqpoint{3.268899in}{1.397043in}}%
\pgfpathlineto{\pgfqpoint{3.240072in}{1.444626in}}%
\pgfpathlineto{\pgfqpoint{3.288022in}{1.461768in}}%
\pgfpathlineto{\pgfqpoint{3.335888in}{1.478856in}}%
\pgfpathlineto{\pgfqpoint{3.364766in}{1.431154in}}%
\pgfpathlineto{\pgfqpoint{3.393767in}{1.386536in}}%
\pgfpathlineto{\pgfqpoint{3.345847in}{1.369607in}}%
\pgfpathlineto{\pgfqpoint{3.297843in}{1.352632in}}%
\pgfpathclose%
\pgfusepath{stroke,fill}%
\end{pgfscope}%
\begin{pgfscope}%
\pgfpathrectangle{\pgfqpoint{0.050000in}{0.083252in}}{\pgfqpoint{4.900000in}{4.900000in}}%
\pgfusepath{clip}%
\pgfsetbuttcap%
\pgfsetroundjoin%
\definecolor{currentfill}{rgb}{1.000000,1.000000,1.000000}%
\pgfsetfillcolor{currentfill}%
\pgfsetfillopacity{0.950000}%
\pgfsetlinewidth{0.200750pt}%
\definecolor{currentstroke}{rgb}{0.200000,0.200000,0.200000}%
\pgfsetstrokecolor{currentstroke}%
\pgfsetdash{}{0pt}%
\pgfpathmoveto{\pgfqpoint{4.107580in}{1.264776in}}%
\pgfpathlineto{\pgfqpoint{4.077238in}{1.296682in}}%
\pgfpathlineto{\pgfqpoint{4.046999in}{1.328925in}}%
\pgfpathlineto{\pgfqpoint{4.094180in}{1.344569in}}%
\pgfpathlineto{\pgfqpoint{4.141276in}{1.360185in}}%
\pgfpathlineto{\pgfqpoint{4.171581in}{1.327596in}}%
\pgfpathlineto{\pgfqpoint{4.201989in}{1.295332in}}%
\pgfpathlineto{\pgfqpoint{4.154828in}{1.280068in}}%
\pgfpathlineto{\pgfqpoint{4.107580in}{1.264776in}}%
\pgfpathclose%
\pgfusepath{stroke,fill}%
\end{pgfscope}%
\begin{pgfscope}%
\pgfpathrectangle{\pgfqpoint{0.050000in}{0.083252in}}{\pgfqpoint{4.900000in}{4.900000in}}%
\pgfusepath{clip}%
\pgfsetbuttcap%
\pgfsetroundjoin%
\definecolor{currentfill}{rgb}{0.743422,0.732841,0.846582}%
\pgfsetfillcolor{currentfill}%
\pgfsetfillopacity{0.950000}%
\pgfsetlinewidth{0.200750pt}%
\definecolor{currentstroke}{rgb}{0.200000,0.200000,0.200000}%
\pgfsetstrokecolor{currentstroke}%
\pgfsetdash{}{0pt}%
\pgfpathmoveto{\pgfqpoint{2.683379in}{1.660098in}}%
\pgfpathlineto{\pgfqpoint{2.655225in}{1.728854in}}%
\pgfpathlineto{\pgfqpoint{2.627109in}{1.799331in}}%
\pgfpathlineto{\pgfqpoint{2.675619in}{1.815670in}}%
\pgfpathlineto{\pgfqpoint{2.724044in}{1.831986in}}%
\pgfpathlineto{\pgfqpoint{2.752169in}{1.762020in}}%
\pgfpathlineto{\pgfqpoint{2.780337in}{1.693778in}}%
\pgfpathlineto{\pgfqpoint{2.731900in}{1.676958in}}%
\pgfpathlineto{\pgfqpoint{2.683379in}{1.660098in}}%
\pgfpathclose%
\pgfusepath{stroke,fill}%
\end{pgfscope}%
\begin{pgfscope}%
\pgfpathrectangle{\pgfqpoint{0.050000in}{0.083252in}}{\pgfqpoint{4.900000in}{4.900000in}}%
\pgfusepath{clip}%
\pgfsetbuttcap%
\pgfsetroundjoin%
\definecolor{currentfill}{rgb}{0.247459,0.295425,0.661530}%
\pgfsetfillcolor{currentfill}%
\pgfsetfillopacity{0.950000}%
\pgfsetlinewidth{0.200750pt}%
\definecolor{currentstroke}{rgb}{0.200000,0.200000,0.200000}%
\pgfsetstrokecolor{currentstroke}%
\pgfsetdash{}{0pt}%
\pgfpathmoveto{\pgfqpoint{1.691064in}{2.844850in}}%
\pgfpathlineto{\pgfqpoint{1.663400in}{2.920119in}}%
\pgfpathlineto{\pgfqpoint{1.635763in}{2.995317in}}%
\pgfpathlineto{\pgfqpoint{1.685172in}{3.009721in}}%
\pgfpathlineto{\pgfqpoint{1.734492in}{3.024100in}}%
\pgfpathlineto{\pgfqpoint{1.762126in}{2.949052in}}%
\pgfpathlineto{\pgfqpoint{1.789786in}{2.873932in}}%
\pgfpathlineto{\pgfqpoint{1.740469in}{2.859404in}}%
\pgfpathlineto{\pgfqpoint{1.691064in}{2.844850in}}%
\pgfpathclose%
\pgfusepath{stroke,fill}%
\end{pgfscope}%
\begin{pgfscope}%
\pgfpathrectangle{\pgfqpoint{0.050000in}{0.083252in}}{\pgfqpoint{4.900000in}{4.900000in}}%
\pgfusepath{clip}%
\pgfsetbuttcap%
\pgfsetroundjoin%
\definecolor{currentfill}{rgb}{0.240815,0.224837,0.567689}%
\pgfsetfillcolor{currentfill}%
\pgfsetfillopacity{0.950000}%
\pgfsetlinewidth{0.200750pt}%
\definecolor{currentstroke}{rgb}{0.200000,0.200000,0.200000}%
\pgfsetstrokecolor{currentstroke}%
\pgfsetdash{}{0pt}%
\pgfpathmoveto{\pgfqpoint{1.900692in}{2.572737in}}%
\pgfpathlineto{\pgfqpoint{1.872926in}{2.648142in}}%
\pgfpathlineto{\pgfqpoint{1.845186in}{2.723477in}}%
\pgfpathlineto{\pgfqpoint{1.894411in}{2.738128in}}%
\pgfpathlineto{\pgfqpoint{1.943548in}{2.752753in}}%
\pgfpathlineto{\pgfqpoint{1.971284in}{2.677569in}}%
\pgfpathlineto{\pgfqpoint{1.999047in}{2.602314in}}%
\pgfpathlineto{\pgfqpoint{1.949914in}{2.587539in}}%
\pgfpathlineto{\pgfqpoint{1.900692in}{2.572737in}}%
\pgfpathclose%
\pgfusepath{stroke,fill}%
\end{pgfscope}%
\begin{pgfscope}%
\pgfpathrectangle{\pgfqpoint{0.050000in}{0.083252in}}{\pgfqpoint{4.900000in}{4.900000in}}%
\pgfusepath{clip}%
\pgfsetbuttcap%
\pgfsetroundjoin%
\definecolor{currentfill}{rgb}{0.355571,0.328997,0.614671}%
\pgfsetfillcolor{currentfill}%
\pgfsetfillopacity{0.950000}%
\pgfsetlinewidth{0.200750pt}%
\definecolor{currentstroke}{rgb}{0.200000,0.200000,0.200000}%
\pgfsetstrokecolor{currentstroke}%
\pgfsetdash{}{0pt}%
\pgfpathmoveto{\pgfqpoint{2.110364in}{2.300651in}}%
\pgfpathlineto{\pgfqpoint{2.082495in}{2.376152in}}%
\pgfpathlineto{\pgfqpoint{2.054652in}{2.451602in}}%
\pgfpathlineto{\pgfqpoint{2.103694in}{2.466509in}}%
\pgfpathlineto{\pgfqpoint{2.152647in}{2.481391in}}%
\pgfpathlineto{\pgfqpoint{2.180486in}{2.406107in}}%
\pgfpathlineto{\pgfqpoint{2.208351in}{2.330783in}}%
\pgfpathlineto{\pgfqpoint{2.159401in}{2.315727in}}%
\pgfpathlineto{\pgfqpoint{2.110364in}{2.300651in}}%
\pgfpathclose%
\pgfusepath{stroke,fill}%
\end{pgfscope}%
\begin{pgfscope}%
\pgfpathrectangle{\pgfqpoint{0.050000in}{0.083252in}}{\pgfqpoint{4.900000in}{4.900000in}}%
\pgfusepath{clip}%
\pgfsetbuttcap%
\pgfsetroundjoin%
\definecolor{currentfill}{rgb}{0.516678,0.496747,0.711003}%
\pgfsetfillcolor{currentfill}%
\pgfsetfillopacity{0.950000}%
\pgfsetlinewidth{0.200750pt}%
\definecolor{currentstroke}{rgb}{0.200000,0.200000,0.200000}%
\pgfsetstrokecolor{currentstroke}%
\pgfsetdash{}{0pt}%
\pgfpathmoveto{\pgfqpoint{2.320073in}{2.029726in}}%
\pgfpathlineto{\pgfqpoint{2.292104in}{2.104844in}}%
\pgfpathlineto{\pgfqpoint{2.264160in}{2.180108in}}%
\pgfpathlineto{\pgfqpoint{2.313018in}{2.195355in}}%
\pgfpathlineto{\pgfqpoint{2.361788in}{2.210586in}}%
\pgfpathlineto{\pgfqpoint{2.389728in}{2.135583in}}%
\pgfpathlineto{\pgfqpoint{2.417694in}{2.060769in}}%
\pgfpathlineto{\pgfqpoint{2.368928in}{2.045253in}}%
\pgfpathlineto{\pgfqpoint{2.320073in}{2.029726in}}%
\pgfpathclose%
\pgfusepath{stroke,fill}%
\end{pgfscope}%
\begin{pgfscope}%
\pgfpathrectangle{\pgfqpoint{0.050000in}{0.083252in}}{\pgfqpoint{4.900000in}{4.900000in}}%
\pgfusepath{clip}%
\pgfsetbuttcap%
\pgfsetroundjoin%
\definecolor{currentfill}{rgb}{0.671819,0.658285,0.803768}%
\pgfsetfillcolor{currentfill}%
\pgfsetfillopacity{0.950000}%
\pgfsetlinewidth{0.200750pt}%
\definecolor{currentstroke}{rgb}{0.200000,0.200000,0.200000}%
\pgfsetstrokecolor{currentstroke}%
\pgfsetdash{}{0pt}%
\pgfpathmoveto{\pgfqpoint{2.529832in}{1.766588in}}%
\pgfpathlineto{\pgfqpoint{2.501755in}{1.838927in}}%
\pgfpathlineto{\pgfqpoint{2.473708in}{1.912247in}}%
\pgfpathlineto{\pgfqpoint{2.522386in}{1.928137in}}%
\pgfpathlineto{\pgfqpoint{2.570978in}{1.944016in}}%
\pgfpathlineto{\pgfqpoint{2.599028in}{1.871160in}}%
\pgfpathlineto{\pgfqpoint{2.627109in}{1.799331in}}%
\pgfpathlineto{\pgfqpoint{2.578514in}{1.782970in}}%
\pgfpathlineto{\pgfqpoint{2.529832in}{1.766588in}}%
\pgfpathclose%
\pgfusepath{stroke,fill}%
\end{pgfscope}%
\begin{pgfscope}%
\pgfpathrectangle{\pgfqpoint{0.050000in}{0.083252in}}{\pgfqpoint{4.900000in}{4.900000in}}%
\pgfusepath{clip}%
\pgfsetbuttcap%
\pgfsetroundjoin%
\definecolor{currentfill}{rgb}{0.964198,0.962722,0.978593}%
\pgfsetfillcolor{currentfill}%
\pgfsetfillopacity{0.950000}%
\pgfsetlinewidth{0.200750pt}%
\definecolor{currentstroke}{rgb}{0.200000,0.200000,0.200000}%
\pgfsetstrokecolor{currentstroke}%
\pgfsetdash{}{0pt}%
\pgfpathmoveto{\pgfqpoint{3.452152in}{1.305937in}}%
\pgfpathlineto{\pgfqpoint{3.422895in}{1.344869in}}%
\pgfpathlineto{\pgfqpoint{3.393767in}{1.386536in}}%
\pgfpathlineto{\pgfqpoint{3.441603in}{1.403421in}}%
\pgfpathlineto{\pgfqpoint{3.489354in}{1.420262in}}%
\pgfpathlineto{\pgfqpoint{3.518542in}{1.378340in}}%
\pgfpathlineto{\pgfqpoint{3.547863in}{1.339102in}}%
\pgfpathlineto{\pgfqpoint{3.500050in}{1.322536in}}%
\pgfpathlineto{\pgfqpoint{3.452152in}{1.305937in}}%
\pgfpathclose%
\pgfusepath{stroke,fill}%
\end{pgfscope}%
\begin{pgfscope}%
\pgfpathrectangle{\pgfqpoint{0.050000in}{0.083252in}}{\pgfqpoint{4.900000in}{4.900000in}}%
\pgfusepath{clip}%
\pgfsetbuttcap%
\pgfsetroundjoin%
\definecolor{currentfill}{rgb}{0.263945,0.470588,0.894394}%
\pgfsetfillcolor{currentfill}%
\pgfsetfillopacity{0.950000}%
\pgfsetlinewidth{0.200750pt}%
\definecolor{currentstroke}{rgb}{0.200000,0.200000,0.200000}%
\pgfsetstrokecolor{currentstroke}%
\pgfsetdash{}{0pt}%
\pgfpathmoveto{\pgfqpoint{1.117229in}{3.510924in}}%
\pgfpathlineto{\pgfqpoint{1.089822in}{3.585924in}}%
\pgfpathlineto{\pgfqpoint{1.062441in}{3.660852in}}%
\pgfpathlineto{\pgfqpoint{1.112398in}{3.674667in}}%
\pgfpathlineto{\pgfqpoint{1.162265in}{3.688457in}}%
\pgfpathlineto{\pgfqpoint{1.189644in}{3.613677in}}%
\pgfpathlineto{\pgfqpoint{1.217049in}{3.538826in}}%
\pgfpathlineto{\pgfqpoint{1.167184in}{3.524888in}}%
\pgfpathlineto{\pgfqpoint{1.117229in}{3.510924in}}%
\pgfpathclose%
\pgfusepath{stroke,fill}%
\end{pgfscope}%
\begin{pgfscope}%
\pgfpathrectangle{\pgfqpoint{0.050000in}{0.083252in}}{\pgfqpoint{4.900000in}{4.900000in}}%
\pgfusepath{clip}%
\pgfsetbuttcap%
\pgfsetroundjoin%
\definecolor{currentfill}{rgb}{0.988066,0.987574,0.992864}%
\pgfsetfillcolor{currentfill}%
\pgfsetfillopacity{0.950000}%
\pgfsetlinewidth{0.200750pt}%
\definecolor{currentstroke}{rgb}{0.200000,0.200000,0.200000}%
\pgfsetstrokecolor{currentstroke}%
\pgfsetdash{}{0pt}%
\pgfpathmoveto{\pgfqpoint{3.857417in}{1.266074in}}%
\pgfpathlineto{\pgfqpoint{3.827434in}{1.298478in}}%
\pgfpathlineto{\pgfqpoint{3.797577in}{1.332179in}}%
\pgfpathlineto{\pgfqpoint{3.845032in}{1.348268in}}%
\pgfpathlineto{\pgfqpoint{3.892402in}{1.364328in}}%
\pgfpathlineto{\pgfqpoint{3.922327in}{1.330296in}}%
\pgfpathlineto{\pgfqpoint{3.952380in}{1.297554in}}%
\pgfpathlineto{\pgfqpoint{3.904942in}{1.281827in}}%
\pgfpathlineto{\pgfqpoint{3.857417in}{1.266074in}}%
\pgfpathclose%
\pgfusepath{stroke,fill}%
\end{pgfscope}%
\begin{pgfscope}%
\pgfpathrectangle{\pgfqpoint{0.050000in}{0.083252in}}{\pgfqpoint{4.900000in}{4.900000in}}%
\pgfusepath{clip}%
\pgfsetbuttcap%
\pgfsetroundjoin%
\definecolor{currentfill}{rgb}{0.257301,0.400000,0.800554}%
\pgfsetfillcolor{currentfill}%
\pgfsetfillopacity{0.950000}%
\pgfsetlinewidth{0.200750pt}%
\definecolor{currentstroke}{rgb}{0.200000,0.200000,0.200000}%
\pgfsetstrokecolor{currentstroke}%
\pgfsetdash{}{0pt}%
\pgfpathmoveto{\pgfqpoint{1.326932in}{3.238706in}}%
\pgfpathlineto{\pgfqpoint{1.299422in}{3.313843in}}%
\pgfpathlineto{\pgfqpoint{1.271938in}{3.388909in}}%
\pgfpathlineto{\pgfqpoint{1.321711in}{3.402970in}}%
\pgfpathlineto{\pgfqpoint{1.371395in}{3.417006in}}%
\pgfpathlineto{\pgfqpoint{1.398876in}{3.342088in}}%
\pgfpathlineto{\pgfqpoint{1.426384in}{3.267100in}}%
\pgfpathlineto{\pgfqpoint{1.376703in}{3.252916in}}%
\pgfpathlineto{\pgfqpoint{1.326932in}{3.238706in}}%
\pgfpathclose%
\pgfusepath{stroke,fill}%
\end{pgfscope}%
\begin{pgfscope}%
\pgfpathrectangle{\pgfqpoint{0.050000in}{0.083252in}}{\pgfqpoint{4.900000in}{4.900000in}}%
\pgfusepath{clip}%
\pgfsetbuttcap%
\pgfsetroundjoin%
\definecolor{currentfill}{rgb}{0.269604,0.530719,0.974333}%
\pgfsetfillcolor{currentfill}%
\pgfsetfillopacity{0.950000}%
\pgfsetlinewidth{0.200750pt}%
\definecolor{currentstroke}{rgb}{0.200000,0.200000,0.200000}%
\pgfsetstrokecolor{currentstroke}%
\pgfsetdash{}{0pt}%
\pgfpathmoveto{\pgfqpoint{0.907570in}{3.783086in}}%
\pgfpathlineto{\pgfqpoint{0.880266in}{3.857948in}}%
\pgfpathlineto{\pgfqpoint{0.930405in}{3.871590in}}%
\pgfpathlineto{\pgfqpoint{0.980455in}{3.885208in}}%
\pgfpathlineto{\pgfqpoint{1.007757in}{3.810494in}}%
\pgfpathlineto{\pgfqpoint{0.957709in}{3.796802in}}%
\pgfpathlineto{\pgfqpoint{0.907570in}{3.783086in}}%
\pgfpathclose%
\pgfusepath{stroke,fill}%
\end{pgfscope}%
\begin{pgfscope}%
\pgfpathrectangle{\pgfqpoint{0.050000in}{0.083252in}}{\pgfqpoint{4.900000in}{4.900000in}}%
\pgfusepath{clip}%
\pgfsetbuttcap%
\pgfsetroundjoin%
\definecolor{currentfill}{rgb}{0.250657,0.329412,0.706713}%
\pgfsetfillcolor{currentfill}%
\pgfsetfillopacity{0.950000}%
\pgfsetlinewidth{0.200750pt}%
\definecolor{currentstroke}{rgb}{0.200000,0.200000,0.200000}%
\pgfsetstrokecolor{currentstroke}%
\pgfsetdash{}{0pt}%
\pgfpathmoveto{\pgfqpoint{1.536679in}{2.966429in}}%
\pgfpathlineto{\pgfqpoint{1.509065in}{3.041705in}}%
\pgfpathlineto{\pgfqpoint{1.481479in}{3.116908in}}%
\pgfpathlineto{\pgfqpoint{1.531068in}{3.131216in}}%
\pgfpathlineto{\pgfqpoint{1.580569in}{3.145497in}}%
\pgfpathlineto{\pgfqpoint{1.608153in}{3.070443in}}%
\pgfpathlineto{\pgfqpoint{1.635763in}{2.995317in}}%
\pgfpathlineto{\pgfqpoint{1.586266in}{2.980886in}}%
\pgfpathlineto{\pgfqpoint{1.536679in}{2.966429in}}%
\pgfpathclose%
\pgfusepath{stroke,fill}%
\end{pgfscope}%
\begin{pgfscope}%
\pgfpathrectangle{\pgfqpoint{0.050000in}{0.083252in}}{\pgfqpoint{4.900000in}{4.900000in}}%
\pgfusepath{clip}%
\pgfsetbuttcap%
\pgfsetroundjoin%
\definecolor{currentfill}{rgb}{0.244260,0.261438,0.616348}%
\pgfsetfillcolor{currentfill}%
\pgfsetfillopacity{0.950000}%
\pgfsetlinewidth{0.200750pt}%
\definecolor{currentstroke}{rgb}{0.200000,0.200000,0.200000}%
\pgfsetstrokecolor{currentstroke}%
\pgfsetdash{}{0pt}%
\pgfpathmoveto{\pgfqpoint{1.746469in}{2.694096in}}%
\pgfpathlineto{\pgfqpoint{1.718753in}{2.769509in}}%
\pgfpathlineto{\pgfqpoint{1.691064in}{2.844850in}}%
\pgfpathlineto{\pgfqpoint{1.740469in}{2.859404in}}%
\pgfpathlineto{\pgfqpoint{1.789786in}{2.873932in}}%
\pgfpathlineto{\pgfqpoint{1.817473in}{2.798740in}}%
\pgfpathlineto{\pgfqpoint{1.845186in}{2.723477in}}%
\pgfpathlineto{\pgfqpoint{1.795872in}{2.708800in}}%
\pgfpathlineto{\pgfqpoint{1.746469in}{2.694096in}}%
\pgfpathclose%
\pgfusepath{stroke,fill}%
\end{pgfscope}%
\begin{pgfscope}%
\pgfpathrectangle{\pgfqpoint{0.050000in}{0.083252in}}{\pgfqpoint{4.900000in}{4.900000in}}%
\pgfusepath{clip}%
\pgfsetbuttcap%
\pgfsetroundjoin%
\definecolor{currentfill}{rgb}{0.278001,0.248228,0.568289}%
\pgfsetfillcolor{currentfill}%
\pgfsetfillopacity{0.950000}%
\pgfsetlinewidth{0.200750pt}%
\definecolor{currentstroke}{rgb}{0.200000,0.200000,0.200000}%
\pgfsetstrokecolor{currentstroke}%
\pgfsetdash{}{0pt}%
\pgfpathmoveto{\pgfqpoint{1.956304in}{2.421715in}}%
\pgfpathlineto{\pgfqpoint{1.928485in}{2.497260in}}%
\pgfpathlineto{\pgfqpoint{1.900692in}{2.572737in}}%
\pgfpathlineto{\pgfqpoint{1.949914in}{2.587539in}}%
\pgfpathlineto{\pgfqpoint{1.999047in}{2.602314in}}%
\pgfpathlineto{\pgfqpoint{2.026837in}{2.526991in}}%
\pgfpathlineto{\pgfqpoint{2.054652in}{2.451602in}}%
\pgfpathlineto{\pgfqpoint{2.005523in}{2.436671in}}%
\pgfpathlineto{\pgfqpoint{1.956304in}{2.421715in}}%
\pgfpathclose%
\pgfusepath{stroke,fill}%
\end{pgfscope}%
\begin{pgfscope}%
\pgfpathrectangle{\pgfqpoint{0.050000in}{0.083252in}}{\pgfqpoint{4.900000in}{4.900000in}}%
\pgfusepath{clip}%
\pgfsetbuttcap%
\pgfsetroundjoin%
\definecolor{currentfill}{rgb}{0.862760,0.857101,0.917939}%
\pgfsetfillcolor{currentfill}%
\pgfsetfillopacity{0.950000}%
\pgfsetlinewidth{0.200750pt}%
\definecolor{currentstroke}{rgb}{0.200000,0.200000,0.200000}%
\pgfsetstrokecolor{currentstroke}%
\pgfsetdash{}{0pt}%
\pgfpathmoveto{\pgfqpoint{2.893558in}{1.445504in}}%
\pgfpathlineto{\pgfqpoint{2.865155in}{1.503156in}}%
\pgfpathlineto{\pgfqpoint{2.836823in}{1.563973in}}%
\pgfpathlineto{\pgfqpoint{2.885197in}{1.581188in}}%
\pgfpathlineto{\pgfqpoint{2.933487in}{1.598346in}}%
\pgfpathlineto{\pgfqpoint{2.961850in}{1.537844in}}%
\pgfpathlineto{\pgfqpoint{2.990290in}{1.480399in}}%
\pgfpathlineto{\pgfqpoint{2.941966in}{1.462984in}}%
\pgfpathlineto{\pgfqpoint{2.893558in}{1.445504in}}%
\pgfpathclose%
\pgfusepath{stroke,fill}%
\end{pgfscope}%
\begin{pgfscope}%
\pgfpathrectangle{\pgfqpoint{0.050000in}{0.083252in}}{\pgfqpoint{4.900000in}{4.900000in}}%
\pgfusepath{clip}%
\pgfsetbuttcap%
\pgfsetroundjoin%
\definecolor{currentfill}{rgb}{0.439108,0.415978,0.664621}%
\pgfsetfillcolor{currentfill}%
\pgfsetfillopacity{0.950000}%
\pgfsetlinewidth{0.200750pt}%
\definecolor{currentstroke}{rgb}{0.200000,0.200000,0.200000}%
\pgfsetstrokecolor{currentstroke}%
\pgfsetdash{}{0pt}%
\pgfpathmoveto{\pgfqpoint{2.166180in}{2.149565in}}%
\pgfpathlineto{\pgfqpoint{2.138259in}{2.225113in}}%
\pgfpathlineto{\pgfqpoint{2.110364in}{2.300651in}}%
\pgfpathlineto{\pgfqpoint{2.159401in}{2.315727in}}%
\pgfpathlineto{\pgfqpoint{2.208351in}{2.330783in}}%
\pgfpathlineto{\pgfqpoint{2.236242in}{2.255438in}}%
\pgfpathlineto{\pgfqpoint{2.264160in}{2.180108in}}%
\pgfpathlineto{\pgfqpoint{2.215214in}{2.164845in}}%
\pgfpathlineto{\pgfqpoint{2.166180in}{2.149565in}}%
\pgfpathclose%
\pgfusepath{stroke,fill}%
\end{pgfscope}%
\begin{pgfscope}%
\pgfpathrectangle{\pgfqpoint{0.050000in}{0.083252in}}{\pgfqpoint{4.900000in}{4.900000in}}%
\pgfusepath{clip}%
\pgfsetbuttcap%
\pgfsetroundjoin%
\definecolor{currentfill}{rgb}{0.904529,0.900592,0.942914}%
\pgfsetfillcolor{currentfill}%
\pgfsetfillopacity{0.950000}%
\pgfsetlinewidth{0.200750pt}%
\definecolor{currentstroke}{rgb}{0.200000,0.200000,0.200000}%
\pgfsetstrokecolor{currentstroke}%
\pgfsetdash{}{0pt}%
\pgfpathmoveto{\pgfqpoint{3.047433in}{1.375493in}}%
\pgfpathlineto{\pgfqpoint{3.018815in}{1.426236in}}%
\pgfpathlineto{\pgfqpoint{2.990290in}{1.480399in}}%
\pgfpathlineto{\pgfqpoint{3.038530in}{1.497750in}}%
\pgfpathlineto{\pgfqpoint{3.086687in}{1.515040in}}%
\pgfpathlineto{\pgfqpoint{3.115254in}{1.460958in}}%
\pgfpathlineto{\pgfqpoint{3.143921in}{1.410176in}}%
\pgfpathlineto{\pgfqpoint{3.095719in}{1.392864in}}%
\pgfpathlineto{\pgfqpoint{3.047433in}{1.375493in}}%
\pgfpathclose%
\pgfusepath{stroke,fill}%
\end{pgfscope}%
\begin{pgfscope}%
\pgfpathrectangle{\pgfqpoint{0.050000in}{0.083252in}}{\pgfqpoint{4.900000in}{4.900000in}}%
\pgfusepath{clip}%
\pgfsetbuttcap%
\pgfsetroundjoin%
\definecolor{currentfill}{rgb}{0.809058,0.801184,0.885829}%
\pgfsetfillcolor{currentfill}%
\pgfsetfillopacity{0.950000}%
\pgfsetlinewidth{0.200750pt}%
\definecolor{currentstroke}{rgb}{0.200000,0.200000,0.200000}%
\pgfsetstrokecolor{currentstroke}%
\pgfsetdash{}{0pt}%
\pgfpathmoveto{\pgfqpoint{2.739822in}{1.529372in}}%
\pgfpathlineto{\pgfqpoint{2.711576in}{1.593466in}}%
\pgfpathlineto{\pgfqpoint{2.683379in}{1.660098in}}%
\pgfpathlineto{\pgfqpoint{2.731900in}{1.676958in}}%
\pgfpathlineto{\pgfqpoint{2.780337in}{1.693778in}}%
\pgfpathlineto{\pgfqpoint{2.808552in}{1.627635in}}%
\pgfpathlineto{\pgfqpoint{2.836823in}{1.563973in}}%
\pgfpathlineto{\pgfqpoint{2.788364in}{1.546702in}}%
\pgfpathlineto{\pgfqpoint{2.739822in}{1.529372in}}%
\pgfpathclose%
\pgfusepath{stroke,fill}%
\end{pgfscope}%
\begin{pgfscope}%
\pgfpathrectangle{\pgfqpoint{0.050000in}{0.083252in}}{\pgfqpoint{4.900000in}{4.900000in}}%
\pgfusepath{clip}%
\pgfsetbuttcap%
\pgfsetroundjoin%
\definecolor{currentfill}{rgb}{0.594248,0.577516,0.757386}%
\pgfsetfillcolor{currentfill}%
\pgfsetfillopacity{0.950000}%
\pgfsetlinewidth{0.200750pt}%
\definecolor{currentstroke}{rgb}{0.200000,0.200000,0.200000}%
\pgfsetstrokecolor{currentstroke}%
\pgfsetdash{}{0pt}%
\pgfpathmoveto{\pgfqpoint{2.376090in}{1.880433in}}%
\pgfpathlineto{\pgfqpoint{2.348069in}{1.954868in}}%
\pgfpathlineto{\pgfqpoint{2.320073in}{2.029726in}}%
\pgfpathlineto{\pgfqpoint{2.368928in}{2.045253in}}%
\pgfpathlineto{\pgfqpoint{2.417694in}{2.060769in}}%
\pgfpathlineto{\pgfqpoint{2.445688in}{1.986268in}}%
\pgfpathlineto{\pgfqpoint{2.473708in}{1.912247in}}%
\pgfpathlineto{\pgfqpoint{2.424942in}{1.896345in}}%
\pgfpathlineto{\pgfqpoint{2.376090in}{1.880433in}}%
\pgfpathclose%
\pgfusepath{stroke,fill}%
\end{pgfscope}%
\begin{pgfscope}%
\pgfpathrectangle{\pgfqpoint{0.050000in}{0.083252in}}{\pgfqpoint{4.900000in}{4.900000in}}%
\pgfusepath{clip}%
\pgfsetbuttcap%
\pgfsetroundjoin%
\definecolor{currentfill}{rgb}{0.976132,0.975148,0.985729}%
\pgfsetfillcolor{currentfill}%
\pgfsetfillopacity{0.950000}%
\pgfsetlinewidth{0.200750pt}%
\definecolor{currentstroke}{rgb}{0.200000,0.200000,0.200000}%
\pgfsetstrokecolor{currentstroke}%
\pgfsetdash{}{0pt}%
\pgfpathmoveto{\pgfqpoint{3.606905in}{1.267544in}}%
\pgfpathlineto{\pgfqpoint{3.577318in}{1.302279in}}%
\pgfpathlineto{\pgfqpoint{3.547863in}{1.339102in}}%
\pgfpathlineto{\pgfqpoint{3.595592in}{1.355634in}}%
\pgfpathlineto{\pgfqpoint{3.643235in}{1.372131in}}%
\pgfpathlineto{\pgfqpoint{3.672756in}{1.334989in}}%
\pgfpathlineto{\pgfqpoint{3.702412in}{1.299917in}}%
\pgfpathlineto{\pgfqpoint{3.654701in}{1.283744in}}%
\pgfpathlineto{\pgfqpoint{3.606905in}{1.267544in}}%
\pgfpathclose%
\pgfusepath{stroke,fill}%
\end{pgfscope}%
\begin{pgfscope}%
\pgfpathrectangle{\pgfqpoint{0.050000in}{0.083252in}}{\pgfqpoint{4.900000in}{4.900000in}}%
\pgfusepath{clip}%
\pgfsetbuttcap%
\pgfsetroundjoin%
\definecolor{currentfill}{rgb}{0.267143,0.504575,0.939577}%
\pgfsetfillcolor{currentfill}%
\pgfsetfillopacity{0.950000}%
\pgfsetlinewidth{0.200750pt}%
\definecolor{currentstroke}{rgb}{0.200000,0.200000,0.200000}%
\pgfsetstrokecolor{currentstroke}%
\pgfsetdash{}{0pt}%
\pgfpathmoveto{\pgfqpoint{0.962256in}{3.633147in}}%
\pgfpathlineto{\pgfqpoint{0.934900in}{3.708152in}}%
\pgfpathlineto{\pgfqpoint{0.907570in}{3.783086in}}%
\pgfpathlineto{\pgfqpoint{0.957709in}{3.796802in}}%
\pgfpathlineto{\pgfqpoint{1.007757in}{3.810494in}}%
\pgfpathlineto{\pgfqpoint{1.035086in}{3.735709in}}%
\pgfpathlineto{\pgfqpoint{1.062441in}{3.660852in}}%
\pgfpathlineto{\pgfqpoint{1.012394in}{3.647012in}}%
\pgfpathlineto{\pgfqpoint{0.962256in}{3.633147in}}%
\pgfpathclose%
\pgfusepath{stroke,fill}%
\end{pgfscope}%
\begin{pgfscope}%
\pgfpathrectangle{\pgfqpoint{0.050000in}{0.083252in}}{\pgfqpoint{4.900000in}{4.900000in}}%
\pgfusepath{clip}%
\pgfsetbuttcap%
\pgfsetroundjoin%
\definecolor{currentfill}{rgb}{0.260500,0.433987,0.845736}%
\pgfsetfillcolor{currentfill}%
\pgfsetfillopacity{0.950000}%
\pgfsetlinewidth{0.200750pt}%
\definecolor{currentstroke}{rgb}{0.200000,0.200000,0.200000}%
\pgfsetstrokecolor{currentstroke}%
\pgfsetdash{}{0pt}%
\pgfpathmoveto{\pgfqpoint{1.172121in}{3.360710in}}%
\pgfpathlineto{\pgfqpoint{1.144662in}{3.435853in}}%
\pgfpathlineto{\pgfqpoint{1.117229in}{3.510924in}}%
\pgfpathlineto{\pgfqpoint{1.167184in}{3.524888in}}%
\pgfpathlineto{\pgfqpoint{1.217049in}{3.538826in}}%
\pgfpathlineto{\pgfqpoint{1.244480in}{3.463903in}}%
\pgfpathlineto{\pgfqpoint{1.271938in}{3.388909in}}%
\pgfpathlineto{\pgfqpoint{1.222075in}{3.374822in}}%
\pgfpathlineto{\pgfqpoint{1.172121in}{3.360710in}}%
\pgfpathclose%
\pgfusepath{stroke,fill}%
\end{pgfscope}%
\begin{pgfscope}%
\pgfpathrectangle{\pgfqpoint{0.050000in}{0.083252in}}{\pgfqpoint{4.900000in}{4.900000in}}%
\pgfusepath{clip}%
\pgfsetbuttcap%
\pgfsetroundjoin%
\definecolor{currentfill}{rgb}{0.940331,0.937870,0.964321}%
\pgfsetfillcolor{currentfill}%
\pgfsetfillopacity{0.950000}%
\pgfsetlinewidth{0.200750pt}%
\definecolor{currentstroke}{rgb}{0.200000,0.200000,0.200000}%
\pgfsetstrokecolor{currentstroke}%
\pgfsetdash{}{0pt}%
\pgfpathmoveto{\pgfqpoint{3.201580in}{1.318547in}}%
\pgfpathlineto{\pgfqpoint{3.172694in}{1.362728in}}%
\pgfpathlineto{\pgfqpoint{3.143921in}{1.410176in}}%
\pgfpathlineto{\pgfqpoint{3.192039in}{1.427429in}}%
\pgfpathlineto{\pgfqpoint{3.240072in}{1.444626in}}%
\pgfpathlineto{\pgfqpoint{3.268899in}{1.397043in}}%
\pgfpathlineto{\pgfqpoint{3.297843in}{1.352632in}}%
\pgfpathlineto{\pgfqpoint{3.249754in}{1.335612in}}%
\pgfpathlineto{\pgfqpoint{3.201580in}{1.318547in}}%
\pgfpathclose%
\pgfusepath{stroke,fill}%
\end{pgfscope}%
\begin{pgfscope}%
\pgfpathrectangle{\pgfqpoint{0.050000in}{0.083252in}}{\pgfqpoint{4.900000in}{4.900000in}}%
\pgfusepath{clip}%
\pgfsetbuttcap%
\pgfsetroundjoin%
\definecolor{currentfill}{rgb}{0.254102,0.366013,0.755371}%
\pgfsetfillcolor{currentfill}%
\pgfsetfillopacity{0.950000}%
\pgfsetlinewidth{0.200750pt}%
\definecolor{currentstroke}{rgb}{0.200000,0.200000,0.200000}%
\pgfsetstrokecolor{currentstroke}%
\pgfsetdash{}{0pt}%
\pgfpathmoveto{\pgfqpoint{1.382031in}{3.088215in}}%
\pgfpathlineto{\pgfqpoint{1.354468in}{3.163496in}}%
\pgfpathlineto{\pgfqpoint{1.326932in}{3.238706in}}%
\pgfpathlineto{\pgfqpoint{1.376703in}{3.252916in}}%
\pgfpathlineto{\pgfqpoint{1.426384in}{3.267100in}}%
\pgfpathlineto{\pgfqpoint{1.453918in}{3.192040in}}%
\pgfpathlineto{\pgfqpoint{1.481479in}{3.116908in}}%
\pgfpathlineto{\pgfqpoint{1.431800in}{3.102575in}}%
\pgfpathlineto{\pgfqpoint{1.382031in}{3.088215in}}%
\pgfpathclose%
\pgfusepath{stroke,fill}%
\end{pgfscope}%
\begin{pgfscope}%
\pgfpathrectangle{\pgfqpoint{0.050000in}{0.083252in}}{\pgfqpoint{4.900000in}{4.900000in}}%
\pgfusepath{clip}%
\pgfsetbuttcap%
\pgfsetroundjoin%
\definecolor{currentfill}{rgb}{0.743422,0.732841,0.846582}%
\pgfsetfillcolor{currentfill}%
\pgfsetfillopacity{0.950000}%
\pgfsetlinewidth{0.200750pt}%
\definecolor{currentstroke}{rgb}{0.200000,0.200000,0.200000}%
\pgfsetstrokecolor{currentstroke}%
\pgfsetdash{}{0pt}%
\pgfpathmoveto{\pgfqpoint{2.586081in}{1.626261in}}%
\pgfpathlineto{\pgfqpoint{2.557940in}{1.695568in}}%
\pgfpathlineto{\pgfqpoint{2.529832in}{1.766588in}}%
\pgfpathlineto{\pgfqpoint{2.578514in}{1.782970in}}%
\pgfpathlineto{\pgfqpoint{2.627109in}{1.799331in}}%
\pgfpathlineto{\pgfqpoint{2.655225in}{1.728854in}}%
\pgfpathlineto{\pgfqpoint{2.683379in}{1.660098in}}%
\pgfpathlineto{\pgfqpoint{2.634773in}{1.643199in}}%
\pgfpathlineto{\pgfqpoint{2.586081in}{1.626261in}}%
\pgfpathclose%
\pgfusepath{stroke,fill}%
\end{pgfscope}%
\begin{pgfscope}%
\pgfpathrectangle{\pgfqpoint{0.050000in}{0.083252in}}{\pgfqpoint{4.900000in}{4.900000in}}%
\pgfusepath{clip}%
\pgfsetbuttcap%
\pgfsetroundjoin%
\definecolor{currentfill}{rgb}{1.000000,1.000000,1.000000}%
\pgfsetfillcolor{currentfill}%
\pgfsetfillopacity{0.950000}%
\pgfsetlinewidth{0.200750pt}%
\definecolor{currentstroke}{rgb}{0.200000,0.200000,0.200000}%
\pgfsetstrokecolor{currentstroke}%
\pgfsetdash{}{0pt}%
\pgfpathmoveto{\pgfqpoint{4.012824in}{1.234109in}}%
\pgfpathlineto{\pgfqpoint{3.982550in}{1.265657in}}%
\pgfpathlineto{\pgfqpoint{3.952380in}{1.297554in}}%
\pgfpathlineto{\pgfqpoint{3.999732in}{1.313253in}}%
\pgfpathlineto{\pgfqpoint{4.046999in}{1.328925in}}%
\pgfpathlineto{\pgfqpoint{4.077238in}{1.296682in}}%
\pgfpathlineto{\pgfqpoint{4.107580in}{1.264776in}}%
\pgfpathlineto{\pgfqpoint{4.060246in}{1.249457in}}%
\pgfpathlineto{\pgfqpoint{4.012824in}{1.234109in}}%
\pgfpathclose%
\pgfusepath{stroke,fill}%
\end{pgfscope}%
\begin{pgfscope}%
\pgfpathrectangle{\pgfqpoint{0.050000in}{0.083252in}}{\pgfqpoint{4.900000in}{4.900000in}}%
\pgfusepath{clip}%
\pgfsetbuttcap%
\pgfsetroundjoin%
\definecolor{currentfill}{rgb}{0.247459,0.295425,0.661530}%
\pgfsetfillcolor{currentfill}%
\pgfsetfillopacity{0.950000}%
\pgfsetlinewidth{0.200750pt}%
\definecolor{currentstroke}{rgb}{0.200000,0.200000,0.200000}%
\pgfsetstrokecolor{currentstroke}%
\pgfsetdash{}{0pt}%
\pgfpathmoveto{\pgfqpoint{1.591984in}{2.815663in}}%
\pgfpathlineto{\pgfqpoint{1.564318in}{2.891082in}}%
\pgfpathlineto{\pgfqpoint{1.536679in}{2.966429in}}%
\pgfpathlineto{\pgfqpoint{1.586266in}{2.980886in}}%
\pgfpathlineto{\pgfqpoint{1.635763in}{2.995317in}}%
\pgfpathlineto{\pgfqpoint{1.663400in}{2.920119in}}%
\pgfpathlineto{\pgfqpoint{1.691064in}{2.844850in}}%
\pgfpathlineto{\pgfqpoint{1.641568in}{2.830270in}}%
\pgfpathlineto{\pgfqpoint{1.591984in}{2.815663in}}%
\pgfpathclose%
\pgfusepath{stroke,fill}%
\end{pgfscope}%
\begin{pgfscope}%
\pgfpathrectangle{\pgfqpoint{0.050000in}{0.083252in}}{\pgfqpoint{4.900000in}{4.900000in}}%
\pgfusepath{clip}%
\pgfsetbuttcap%
\pgfsetroundjoin%
\definecolor{currentfill}{rgb}{0.240815,0.224837,0.567689}%
\pgfsetfillcolor{currentfill}%
\pgfsetfillopacity{0.950000}%
\pgfsetlinewidth{0.200750pt}%
\definecolor{currentstroke}{rgb}{0.200000,0.200000,0.200000}%
\pgfsetstrokecolor{currentstroke}%
\pgfsetdash{}{0pt}%
\pgfpathmoveto{\pgfqpoint{1.801981in}{2.543053in}}%
\pgfpathlineto{\pgfqpoint{1.774212in}{2.618610in}}%
\pgfpathlineto{\pgfqpoint{1.746469in}{2.694096in}}%
\pgfpathlineto{\pgfqpoint{1.795872in}{2.708800in}}%
\pgfpathlineto{\pgfqpoint{1.845186in}{2.723477in}}%
\pgfpathlineto{\pgfqpoint{1.872926in}{2.648142in}}%
\pgfpathlineto{\pgfqpoint{1.900692in}{2.572737in}}%
\pgfpathlineto{\pgfqpoint{1.851381in}{2.557908in}}%
\pgfpathlineto{\pgfqpoint{1.801981in}{2.543053in}}%
\pgfpathclose%
\pgfusepath{stroke,fill}%
\end{pgfscope}%
\begin{pgfscope}%
\pgfpathrectangle{\pgfqpoint{0.050000in}{0.083252in}}{\pgfqpoint{4.900000in}{4.900000in}}%
\pgfusepath{clip}%
\pgfsetbuttcap%
\pgfsetroundjoin%
\definecolor{currentfill}{rgb}{0.361538,0.335210,0.618239}%
\pgfsetfillcolor{currentfill}%
\pgfsetfillopacity{0.950000}%
\pgfsetlinewidth{0.200750pt}%
\definecolor{currentstroke}{rgb}{0.200000,0.200000,0.200000}%
\pgfsetstrokecolor{currentstroke}%
\pgfsetdash{}{0pt}%
\pgfpathmoveto{\pgfqpoint{2.012022in}{2.270431in}}%
\pgfpathlineto{\pgfqpoint{1.984150in}{2.346103in}}%
\pgfpathlineto{\pgfqpoint{1.956304in}{2.421715in}}%
\pgfpathlineto{\pgfqpoint{2.005523in}{2.436671in}}%
\pgfpathlineto{\pgfqpoint{2.054652in}{2.451602in}}%
\pgfpathlineto{\pgfqpoint{2.082495in}{2.376152in}}%
\pgfpathlineto{\pgfqpoint{2.110364in}{2.300651in}}%
\pgfpathlineto{\pgfqpoint{2.061237in}{2.285552in}}%
\pgfpathlineto{\pgfqpoint{2.012022in}{2.270431in}}%
\pgfpathclose%
\pgfusepath{stroke,fill}%
\end{pgfscope}%
\begin{pgfscope}%
\pgfpathrectangle{\pgfqpoint{0.050000in}{0.083252in}}{\pgfqpoint{4.900000in}{4.900000in}}%
\pgfusepath{clip}%
\pgfsetbuttcap%
\pgfsetroundjoin%
\definecolor{currentfill}{rgb}{0.516678,0.496747,0.711003}%
\pgfsetfillcolor{currentfill}%
\pgfsetfillopacity{0.950000}%
\pgfsetlinewidth{0.200750pt}%
\definecolor{currentstroke}{rgb}{0.200000,0.200000,0.200000}%
\pgfsetstrokecolor{currentstroke}%
\pgfsetdash{}{0pt}%
\pgfpathmoveto{\pgfqpoint{2.222100in}{1.998641in}}%
\pgfpathlineto{\pgfqpoint{2.194127in}{2.074052in}}%
\pgfpathlineto{\pgfqpoint{2.166180in}{2.149565in}}%
\pgfpathlineto{\pgfqpoint{2.215214in}{2.164845in}}%
\pgfpathlineto{\pgfqpoint{2.264160in}{2.180108in}}%
\pgfpathlineto{\pgfqpoint{2.292104in}{2.104844in}}%
\pgfpathlineto{\pgfqpoint{2.320073in}{2.029726in}}%
\pgfpathlineto{\pgfqpoint{2.271131in}{2.014188in}}%
\pgfpathlineto{\pgfqpoint{2.222100in}{1.998641in}}%
\pgfpathclose%
\pgfusepath{stroke,fill}%
\end{pgfscope}%
\begin{pgfscope}%
\pgfpathrectangle{\pgfqpoint{0.050000in}{0.083252in}}{\pgfqpoint{4.900000in}{4.900000in}}%
\pgfusepath{clip}%
\pgfsetbuttcap%
\pgfsetroundjoin%
\definecolor{currentfill}{rgb}{0.671819,0.658285,0.803768}%
\pgfsetfillcolor{currentfill}%
\pgfsetfillopacity{0.950000}%
\pgfsetlinewidth{0.200750pt}%
\definecolor{currentstroke}{rgb}{0.200000,0.200000,0.200000}%
\pgfsetstrokecolor{currentstroke}%
\pgfsetdash{}{0pt}%
\pgfpathmoveto{\pgfqpoint{2.432210in}{1.733767in}}%
\pgfpathlineto{\pgfqpoint{2.404137in}{1.806638in}}%
\pgfpathlineto{\pgfqpoint{2.376090in}{1.880433in}}%
\pgfpathlineto{\pgfqpoint{2.424942in}{1.896345in}}%
\pgfpathlineto{\pgfqpoint{2.473708in}{1.912247in}}%
\pgfpathlineto{\pgfqpoint{2.501755in}{1.838927in}}%
\pgfpathlineto{\pgfqpoint{2.529832in}{1.766588in}}%
\pgfpathlineto{\pgfqpoint{2.481065in}{1.750186in}}%
\pgfpathlineto{\pgfqpoint{2.432210in}{1.733767in}}%
\pgfpathclose%
\pgfusepath{stroke,fill}%
\end{pgfscope}%
\begin{pgfscope}%
\pgfpathrectangle{\pgfqpoint{0.050000in}{0.083252in}}{\pgfqpoint{4.900000in}{4.900000in}}%
\pgfusepath{clip}%
\pgfsetbuttcap%
\pgfsetroundjoin%
\definecolor{currentfill}{rgb}{0.964198,0.962722,0.978593}%
\pgfsetfillcolor{currentfill}%
\pgfsetfillopacity{0.950000}%
\pgfsetlinewidth{0.200750pt}%
\definecolor{currentstroke}{rgb}{0.200000,0.200000,0.200000}%
\pgfsetstrokecolor{currentstroke}%
\pgfsetdash{}{0pt}%
\pgfpathmoveto{\pgfqpoint{3.356100in}{1.272639in}}%
\pgfpathlineto{\pgfqpoint{3.326909in}{1.311243in}}%
\pgfpathlineto{\pgfqpoint{3.297843in}{1.352632in}}%
\pgfpathlineto{\pgfqpoint{3.345847in}{1.369607in}}%
\pgfpathlineto{\pgfqpoint{3.393767in}{1.386536in}}%
\pgfpathlineto{\pgfqpoint{3.422895in}{1.344869in}}%
\pgfpathlineto{\pgfqpoint{3.452152in}{1.305937in}}%
\pgfpathlineto{\pgfqpoint{3.404169in}{1.289304in}}%
\pgfpathlineto{\pgfqpoint{3.356100in}{1.272639in}}%
\pgfpathclose%
\pgfusepath{stroke,fill}%
\end{pgfscope}%
\begin{pgfscope}%
\pgfpathrectangle{\pgfqpoint{0.050000in}{0.083252in}}{\pgfqpoint{4.900000in}{4.900000in}}%
\pgfusepath{clip}%
\pgfsetbuttcap%
\pgfsetroundjoin%
\definecolor{currentfill}{rgb}{0.263945,0.470588,0.894394}%
\pgfsetfillcolor{currentfill}%
\pgfsetfillopacity{0.950000}%
\pgfsetlinewidth{0.200750pt}%
\definecolor{currentstroke}{rgb}{0.200000,0.200000,0.200000}%
\pgfsetstrokecolor{currentstroke}%
\pgfsetdash{}{0pt}%
\pgfpathmoveto{\pgfqpoint{1.017047in}{3.482922in}}%
\pgfpathlineto{\pgfqpoint{0.989639in}{3.558070in}}%
\pgfpathlineto{\pgfqpoint{0.962256in}{3.633147in}}%
\pgfpathlineto{\pgfqpoint{1.012394in}{3.647012in}}%
\pgfpathlineto{\pgfqpoint{1.062441in}{3.660852in}}%
\pgfpathlineto{\pgfqpoint{1.089822in}{3.585924in}}%
\pgfpathlineto{\pgfqpoint{1.117229in}{3.510924in}}%
\pgfpathlineto{\pgfqpoint{1.067183in}{3.496936in}}%
\pgfpathlineto{\pgfqpoint{1.017047in}{3.482922in}}%
\pgfpathclose%
\pgfusepath{stroke,fill}%
\end{pgfscope}%
\begin{pgfscope}%
\pgfpathrectangle{\pgfqpoint{0.050000in}{0.083252in}}{\pgfqpoint{4.900000in}{4.900000in}}%
\pgfusepath{clip}%
\pgfsetbuttcap%
\pgfsetroundjoin%
\definecolor{currentfill}{rgb}{0.988066,0.987574,0.992864}%
\pgfsetfillcolor{currentfill}%
\pgfsetfillopacity{0.950000}%
\pgfsetlinewidth{0.200750pt}%
\definecolor{currentstroke}{rgb}{0.200000,0.200000,0.200000}%
\pgfsetstrokecolor{currentstroke}%
\pgfsetdash{}{0pt}%
\pgfpathmoveto{\pgfqpoint{3.762108in}{1.234488in}}%
\pgfpathlineto{\pgfqpoint{3.732199in}{1.266552in}}%
\pgfpathlineto{\pgfqpoint{3.702412in}{1.299917in}}%
\pgfpathlineto{\pgfqpoint{3.750038in}{1.316062in}}%
\pgfpathlineto{\pgfqpoint{3.797577in}{1.332179in}}%
\pgfpathlineto{\pgfqpoint{3.827434in}{1.298478in}}%
\pgfpathlineto{\pgfqpoint{3.857417in}{1.266074in}}%
\pgfpathlineto{\pgfqpoint{3.809806in}{1.250294in}}%
\pgfpathlineto{\pgfqpoint{3.762108in}{1.234488in}}%
\pgfpathclose%
\pgfusepath{stroke,fill}%
\end{pgfscope}%
\begin{pgfscope}%
\pgfpathrectangle{\pgfqpoint{0.050000in}{0.083252in}}{\pgfqpoint{4.900000in}{4.900000in}}%
\pgfusepath{clip}%
\pgfsetbuttcap%
\pgfsetroundjoin%
\definecolor{currentfill}{rgb}{0.257301,0.400000,0.800554}%
\pgfsetfillcolor{currentfill}%
\pgfsetfillopacity{0.950000}%
\pgfsetlinewidth{0.200750pt}%
\definecolor{currentstroke}{rgb}{0.200000,0.200000,0.200000}%
\pgfsetstrokecolor{currentstroke}%
\pgfsetdash{}{0pt}%
\pgfpathmoveto{\pgfqpoint{1.227119in}{3.210208in}}%
\pgfpathlineto{\pgfqpoint{1.199607in}{3.285495in}}%
\pgfpathlineto{\pgfqpoint{1.172121in}{3.360710in}}%
\pgfpathlineto{\pgfqpoint{1.222075in}{3.374822in}}%
\pgfpathlineto{\pgfqpoint{1.271938in}{3.388909in}}%
\pgfpathlineto{\pgfqpoint{1.299422in}{3.313843in}}%
\pgfpathlineto{\pgfqpoint{1.326932in}{3.238706in}}%
\pgfpathlineto{\pgfqpoint{1.277071in}{3.224470in}}%
\pgfpathlineto{\pgfqpoint{1.227119in}{3.210208in}}%
\pgfpathclose%
\pgfusepath{stroke,fill}%
\end{pgfscope}%
\begin{pgfscope}%
\pgfpathrectangle{\pgfqpoint{0.050000in}{0.083252in}}{\pgfqpoint{4.900000in}{4.900000in}}%
\pgfusepath{clip}%
\pgfsetbuttcap%
\pgfsetroundjoin%
\definecolor{currentfill}{rgb}{0.269604,0.530719,0.974333}%
\pgfsetfillcolor{currentfill}%
\pgfsetfillopacity{0.950000}%
\pgfsetlinewidth{0.200750pt}%
\definecolor{currentstroke}{rgb}{0.200000,0.200000,0.200000}%
\pgfsetstrokecolor{currentstroke}%
\pgfsetdash{}{0pt}%
\pgfpathmoveto{\pgfqpoint{0.807019in}{3.755578in}}%
\pgfpathlineto{\pgfqpoint{0.779714in}{3.830588in}}%
\pgfpathlineto{\pgfqpoint{0.830035in}{3.844280in}}%
\pgfpathlineto{\pgfqpoint{0.880266in}{3.857948in}}%
\pgfpathlineto{\pgfqpoint{0.907570in}{3.783086in}}%
\pgfpathlineto{\pgfqpoint{0.857340in}{3.769344in}}%
\pgfpathlineto{\pgfqpoint{0.807019in}{3.755578in}}%
\pgfpathclose%
\pgfusepath{stroke,fill}%
\end{pgfscope}%
\begin{pgfscope}%
\pgfpathrectangle{\pgfqpoint{0.050000in}{0.083252in}}{\pgfqpoint{4.900000in}{4.900000in}}%
\pgfusepath{clip}%
\pgfsetbuttcap%
\pgfsetroundjoin%
\definecolor{currentfill}{rgb}{0.250657,0.329412,0.706713}%
\pgfsetfillcolor{currentfill}%
\pgfsetfillopacity{0.950000}%
\pgfsetlinewidth{0.200750pt}%
\definecolor{currentstroke}{rgb}{0.200000,0.200000,0.200000}%
\pgfsetstrokecolor{currentstroke}%
\pgfsetdash{}{0pt}%
\pgfpathmoveto{\pgfqpoint{1.437235in}{2.937437in}}%
\pgfpathlineto{\pgfqpoint{1.409620in}{3.012862in}}%
\pgfpathlineto{\pgfqpoint{1.382031in}{3.088215in}}%
\pgfpathlineto{\pgfqpoint{1.431800in}{3.102575in}}%
\pgfpathlineto{\pgfqpoint{1.481479in}{3.116908in}}%
\pgfpathlineto{\pgfqpoint{1.509065in}{3.041705in}}%
\pgfpathlineto{\pgfqpoint{1.536679in}{2.966429in}}%
\pgfpathlineto{\pgfqpoint{1.487002in}{2.951946in}}%
\pgfpathlineto{\pgfqpoint{1.437235in}{2.937437in}}%
\pgfpathclose%
\pgfusepath{stroke,fill}%
\end{pgfscope}%
\begin{pgfscope}%
\pgfpathrectangle{\pgfqpoint{0.050000in}{0.083252in}}{\pgfqpoint{4.900000in}{4.900000in}}%
\pgfusepath{clip}%
\pgfsetbuttcap%
\pgfsetroundjoin%
\definecolor{currentfill}{rgb}{0.244014,0.258824,0.612872}%
\pgfsetfillcolor{currentfill}%
\pgfsetfillopacity{0.950000}%
\pgfsetlinewidth{0.200750pt}%
\definecolor{currentstroke}{rgb}{0.200000,0.200000,0.200000}%
\pgfsetstrokecolor{currentstroke}%
\pgfsetdash{}{0pt}%
\pgfpathmoveto{\pgfqpoint{1.647395in}{2.664607in}}%
\pgfpathlineto{\pgfqpoint{1.619676in}{2.740171in}}%
\pgfpathlineto{\pgfqpoint{1.591984in}{2.815663in}}%
\pgfpathlineto{\pgfqpoint{1.641568in}{2.830270in}}%
\pgfpathlineto{\pgfqpoint{1.691064in}{2.844850in}}%
\pgfpathlineto{\pgfqpoint{1.718753in}{2.769509in}}%
\pgfpathlineto{\pgfqpoint{1.746469in}{2.694096in}}%
\pgfpathlineto{\pgfqpoint{1.696977in}{2.679365in}}%
\pgfpathlineto{\pgfqpoint{1.647395in}{2.664607in}}%
\pgfpathclose%
\pgfusepath{stroke,fill}%
\end{pgfscope}%
\begin{pgfscope}%
\pgfpathrectangle{\pgfqpoint{0.050000in}{0.083252in}}{\pgfqpoint{4.900000in}{4.900000in}}%
\pgfusepath{clip}%
\pgfsetbuttcap%
\pgfsetroundjoin%
\definecolor{currentfill}{rgb}{0.278001,0.248228,0.568289}%
\pgfsetfillcolor{currentfill}%
\pgfsetfillopacity{0.950000}%
\pgfsetlinewidth{0.200750pt}%
\definecolor{currentstroke}{rgb}{0.200000,0.200000,0.200000}%
\pgfsetstrokecolor{currentstroke}%
\pgfsetdash{}{0pt}%
\pgfpathmoveto{\pgfqpoint{1.857599in}{2.391724in}}%
\pgfpathlineto{\pgfqpoint{1.829777in}{2.467424in}}%
\pgfpathlineto{\pgfqpoint{1.801981in}{2.543053in}}%
\pgfpathlineto{\pgfqpoint{1.851381in}{2.557908in}}%
\pgfpathlineto{\pgfqpoint{1.900692in}{2.572737in}}%
\pgfpathlineto{\pgfqpoint{1.928485in}{2.497260in}}%
\pgfpathlineto{\pgfqpoint{1.956304in}{2.421715in}}%
\pgfpathlineto{\pgfqpoint{1.906996in}{2.406732in}}%
\pgfpathlineto{\pgfqpoint{1.857599in}{2.391724in}}%
\pgfpathclose%
\pgfusepath{stroke,fill}%
\end{pgfscope}%
\begin{pgfscope}%
\pgfpathrectangle{\pgfqpoint{0.050000in}{0.083252in}}{\pgfqpoint{4.900000in}{4.900000in}}%
\pgfusepath{clip}%
\pgfsetbuttcap%
\pgfsetroundjoin%
\definecolor{currentfill}{rgb}{0.868727,0.863314,0.921507}%
\pgfsetfillcolor{currentfill}%
\pgfsetfillopacity{0.950000}%
\pgfsetlinewidth{0.200750pt}%
\definecolor{currentstroke}{rgb}{0.200000,0.200000,0.200000}%
\pgfsetstrokecolor{currentstroke}%
\pgfsetdash{}{0pt}%
\pgfpathmoveto{\pgfqpoint{2.796488in}{1.410338in}}%
\pgfpathlineto{\pgfqpoint{2.768123in}{1.468214in}}%
\pgfpathlineto{\pgfqpoint{2.739822in}{1.529372in}}%
\pgfpathlineto{\pgfqpoint{2.788364in}{1.546702in}}%
\pgfpathlineto{\pgfqpoint{2.836823in}{1.563973in}}%
\pgfpathlineto{\pgfqpoint{2.865155in}{1.503156in}}%
\pgfpathlineto{\pgfqpoint{2.893558in}{1.445504in}}%
\pgfpathlineto{\pgfqpoint{2.845065in}{1.427956in}}%
\pgfpathlineto{\pgfqpoint{2.796488in}{1.410338in}}%
\pgfpathclose%
\pgfusepath{stroke,fill}%
\end{pgfscope}%
\begin{pgfscope}%
\pgfpathrectangle{\pgfqpoint{0.050000in}{0.083252in}}{\pgfqpoint{4.900000in}{4.900000in}}%
\pgfusepath{clip}%
\pgfsetbuttcap%
\pgfsetroundjoin%
\definecolor{currentfill}{rgb}{0.439108,0.415978,0.664621}%
\pgfsetfillcolor{currentfill}%
\pgfsetfillopacity{0.950000}%
\pgfsetlinewidth{0.200750pt}%
\definecolor{currentstroke}{rgb}{0.200000,0.200000,0.200000}%
\pgfsetstrokecolor{currentstroke}%
\pgfsetdash{}{0pt}%
\pgfpathmoveto{\pgfqpoint{2.067845in}{2.118957in}}%
\pgfpathlineto{\pgfqpoint{2.039921in}{2.194709in}}%
\pgfpathlineto{\pgfqpoint{2.012022in}{2.270431in}}%
\pgfpathlineto{\pgfqpoint{2.061237in}{2.285552in}}%
\pgfpathlineto{\pgfqpoint{2.110364in}{2.300651in}}%
\pgfpathlineto{\pgfqpoint{2.138259in}{2.225113in}}%
\pgfpathlineto{\pgfqpoint{2.166180in}{2.149565in}}%
\pgfpathlineto{\pgfqpoint{2.117057in}{2.134270in}}%
\pgfpathlineto{\pgfqpoint{2.067845in}{2.118957in}}%
\pgfpathclose%
\pgfusepath{stroke,fill}%
\end{pgfscope}%
\begin{pgfscope}%
\pgfpathrectangle{\pgfqpoint{0.050000in}{0.083252in}}{\pgfqpoint{4.900000in}{4.900000in}}%
\pgfusepath{clip}%
\pgfsetbuttcap%
\pgfsetroundjoin%
\definecolor{currentfill}{rgb}{0.910496,0.906805,0.946482}%
\pgfsetfillcolor{currentfill}%
\pgfsetfillopacity{0.950000}%
\pgfsetlinewidth{0.200750pt}%
\definecolor{currentstroke}{rgb}{0.200000,0.200000,0.200000}%
\pgfsetstrokecolor{currentstroke}%
\pgfsetdash{}{0pt}%
\pgfpathmoveto{\pgfqpoint{2.950606in}{1.340567in}}%
\pgfpathlineto{\pgfqpoint{2.922039in}{1.391259in}}%
\pgfpathlineto{\pgfqpoint{2.893558in}{1.445504in}}%
\pgfpathlineto{\pgfqpoint{2.941966in}{1.462984in}}%
\pgfpathlineto{\pgfqpoint{2.990290in}{1.480399in}}%
\pgfpathlineto{\pgfqpoint{3.018815in}{1.426236in}}%
\pgfpathlineto{\pgfqpoint{3.047433in}{1.375493in}}%
\pgfpathlineto{\pgfqpoint{2.999062in}{1.358062in}}%
\pgfpathlineto{\pgfqpoint{2.950606in}{1.340567in}}%
\pgfpathclose%
\pgfusepath{stroke,fill}%
\end{pgfscope}%
\begin{pgfscope}%
\pgfpathrectangle{\pgfqpoint{0.050000in}{0.083252in}}{\pgfqpoint{4.900000in}{4.900000in}}%
\pgfusepath{clip}%
\pgfsetbuttcap%
\pgfsetroundjoin%
\definecolor{currentfill}{rgb}{0.809058,0.801184,0.885829}%
\pgfsetfillcolor{currentfill}%
\pgfsetfillopacity{0.950000}%
\pgfsetlinewidth{0.200750pt}%
\definecolor{currentstroke}{rgb}{0.200000,0.200000,0.200000}%
\pgfsetstrokecolor{currentstroke}%
\pgfsetdash{}{0pt}%
\pgfpathmoveto{\pgfqpoint{2.642483in}{1.494534in}}%
\pgfpathlineto{\pgfqpoint{2.614261in}{1.559099in}}%
\pgfpathlineto{\pgfqpoint{2.586081in}{1.626261in}}%
\pgfpathlineto{\pgfqpoint{2.634773in}{1.643199in}}%
\pgfpathlineto{\pgfqpoint{2.683379in}{1.660098in}}%
\pgfpathlineto{\pgfqpoint{2.711576in}{1.593466in}}%
\pgfpathlineto{\pgfqpoint{2.739822in}{1.529372in}}%
\pgfpathlineto{\pgfqpoint{2.691194in}{1.511983in}}%
\pgfpathlineto{\pgfqpoint{2.642483in}{1.494534in}}%
\pgfpathclose%
\pgfusepath{stroke,fill}%
\end{pgfscope}%
\begin{pgfscope}%
\pgfpathrectangle{\pgfqpoint{0.050000in}{0.083252in}}{\pgfqpoint{4.900000in}{4.900000in}}%
\pgfusepath{clip}%
\pgfsetbuttcap%
\pgfsetroundjoin%
\definecolor{currentfill}{rgb}{0.600215,0.583729,0.760953}%
\pgfsetfillcolor{currentfill}%
\pgfsetfillopacity{0.950000}%
\pgfsetlinewidth{0.200750pt}%
\definecolor{currentstroke}{rgb}{0.200000,0.200000,0.200000}%
\pgfsetstrokecolor{currentstroke}%
\pgfsetdash{}{0pt}%
\pgfpathmoveto{\pgfqpoint{2.278122in}{1.848588in}}%
\pgfpathlineto{\pgfqpoint{2.250099in}{1.923436in}}%
\pgfpathlineto{\pgfqpoint{2.222100in}{1.998641in}}%
\pgfpathlineto{\pgfqpoint{2.271131in}{2.014188in}}%
\pgfpathlineto{\pgfqpoint{2.320073in}{2.029726in}}%
\pgfpathlineto{\pgfqpoint{2.348069in}{1.954868in}}%
\pgfpathlineto{\pgfqpoint{2.376090in}{1.880433in}}%
\pgfpathlineto{\pgfqpoint{2.327150in}{1.864514in}}%
\pgfpathlineto{\pgfqpoint{2.278122in}{1.848588in}}%
\pgfpathclose%
\pgfusepath{stroke,fill}%
\end{pgfscope}%
\begin{pgfscope}%
\pgfpathrectangle{\pgfqpoint{0.050000in}{0.083252in}}{\pgfqpoint{4.900000in}{4.900000in}}%
\pgfusepath{clip}%
\pgfsetbuttcap%
\pgfsetroundjoin%
\definecolor{currentfill}{rgb}{0.982099,0.981361,0.989296}%
\pgfsetfillcolor{currentfill}%
\pgfsetfillopacity{0.950000}%
\pgfsetlinewidth{0.200750pt}%
\definecolor{currentstroke}{rgb}{0.200000,0.200000,0.200000}%
\pgfsetstrokecolor{currentstroke}%
\pgfsetdash{}{0pt}%
\pgfpathmoveto{\pgfqpoint{3.511052in}{1.235065in}}%
\pgfpathlineto{\pgfqpoint{3.481539in}{1.269451in}}%
\pgfpathlineto{\pgfqpoint{3.452152in}{1.305937in}}%
\pgfpathlineto{\pgfqpoint{3.500050in}{1.322536in}}%
\pgfpathlineto{\pgfqpoint{3.547863in}{1.339102in}}%
\pgfpathlineto{\pgfqpoint{3.577318in}{1.302279in}}%
\pgfpathlineto{\pgfqpoint{3.606905in}{1.267544in}}%
\pgfpathlineto{\pgfqpoint{3.559022in}{1.251317in}}%
\pgfpathlineto{\pgfqpoint{3.511052in}{1.235065in}}%
\pgfpathclose%
\pgfusepath{stroke,fill}%
\end{pgfscope}%
\begin{pgfscope}%
\pgfpathrectangle{\pgfqpoint{0.050000in}{0.083252in}}{\pgfqpoint{4.900000in}{4.900000in}}%
\pgfusepath{clip}%
\pgfsetbuttcap%
\pgfsetroundjoin%
\definecolor{currentfill}{rgb}{0.267143,0.504575,0.939577}%
\pgfsetfillcolor{currentfill}%
\pgfsetfillopacity{0.950000}%
\pgfsetlinewidth{0.200750pt}%
\definecolor{currentstroke}{rgb}{0.200000,0.200000,0.200000}%
\pgfsetstrokecolor{currentstroke}%
\pgfsetdash{}{0pt}%
\pgfpathmoveto{\pgfqpoint{0.861709in}{3.605341in}}%
\pgfpathlineto{\pgfqpoint{0.834351in}{3.680496in}}%
\pgfpathlineto{\pgfqpoint{0.807019in}{3.755578in}}%
\pgfpathlineto{\pgfqpoint{0.857340in}{3.769344in}}%
\pgfpathlineto{\pgfqpoint{0.907570in}{3.783086in}}%
\pgfpathlineto{\pgfqpoint{0.934900in}{3.708152in}}%
\pgfpathlineto{\pgfqpoint{0.962256in}{3.633147in}}%
\pgfpathlineto{\pgfqpoint{0.912028in}{3.619257in}}%
\pgfpathlineto{\pgfqpoint{0.861709in}{3.605341in}}%
\pgfpathclose%
\pgfusepath{stroke,fill}%
\end{pgfscope}%
\begin{pgfscope}%
\pgfpathrectangle{\pgfqpoint{0.050000in}{0.083252in}}{\pgfqpoint{4.900000in}{4.900000in}}%
\pgfusepath{clip}%
\pgfsetbuttcap%
\pgfsetroundjoin%
\definecolor{currentfill}{rgb}{0.260500,0.433987,0.845736}%
\pgfsetfillcolor{currentfill}%
\pgfsetfillopacity{0.950000}%
\pgfsetlinewidth{0.200750pt}%
\definecolor{currentstroke}{rgb}{0.200000,0.200000,0.200000}%
\pgfsetstrokecolor{currentstroke}%
\pgfsetdash{}{0pt}%
\pgfpathmoveto{\pgfqpoint{1.071943in}{3.332409in}}%
\pgfpathlineto{\pgfqpoint{1.044482in}{3.407701in}}%
\pgfpathlineto{\pgfqpoint{1.017047in}{3.482922in}}%
\pgfpathlineto{\pgfqpoint{1.067183in}{3.496936in}}%
\pgfpathlineto{\pgfqpoint{1.117229in}{3.510924in}}%
\pgfpathlineto{\pgfqpoint{1.144662in}{3.435853in}}%
\pgfpathlineto{\pgfqpoint{1.172121in}{3.360710in}}%
\pgfpathlineto{\pgfqpoint{1.122078in}{3.346572in}}%
\pgfpathlineto{\pgfqpoint{1.071943in}{3.332409in}}%
\pgfpathclose%
\pgfusepath{stroke,fill}%
\end{pgfscope}%
\begin{pgfscope}%
\pgfpathrectangle{\pgfqpoint{0.050000in}{0.083252in}}{\pgfqpoint{4.900000in}{4.900000in}}%
\pgfusepath{clip}%
\pgfsetbuttcap%
\pgfsetroundjoin%
\definecolor{currentfill}{rgb}{0.946298,0.944083,0.967889}%
\pgfsetfillcolor{currentfill}%
\pgfsetfillopacity{0.950000}%
\pgfsetlinewidth{0.200750pt}%
\definecolor{currentstroke}{rgb}{0.200000,0.200000,0.200000}%
\pgfsetstrokecolor{currentstroke}%
\pgfsetdash{}{0pt}%
\pgfpathmoveto{\pgfqpoint{3.104976in}{1.284276in}}%
\pgfpathlineto{\pgfqpoint{3.076151in}{1.328201in}}%
\pgfpathlineto{\pgfqpoint{3.047433in}{1.375493in}}%
\pgfpathlineto{\pgfqpoint{3.095719in}{1.392864in}}%
\pgfpathlineto{\pgfqpoint{3.143921in}{1.410176in}}%
\pgfpathlineto{\pgfqpoint{3.172694in}{1.362728in}}%
\pgfpathlineto{\pgfqpoint{3.201580in}{1.318547in}}%
\pgfpathlineto{\pgfqpoint{3.153320in}{1.301435in}}%
\pgfpathlineto{\pgfqpoint{3.104976in}{1.284276in}}%
\pgfpathclose%
\pgfusepath{stroke,fill}%
\end{pgfscope}%
\begin{pgfscope}%
\pgfpathrectangle{\pgfqpoint{0.050000in}{0.083252in}}{\pgfqpoint{4.900000in}{4.900000in}}%
\pgfusepath{clip}%
\pgfsetbuttcap%
\pgfsetroundjoin%
\definecolor{currentfill}{rgb}{0.253856,0.363399,0.751895}%
\pgfsetfillcolor{currentfill}%
\pgfsetfillopacity{0.950000}%
\pgfsetlinewidth{0.200750pt}%
\definecolor{currentstroke}{rgb}{0.200000,0.200000,0.200000}%
\pgfsetstrokecolor{currentstroke}%
\pgfsetdash{}{0pt}%
\pgfpathmoveto{\pgfqpoint{1.282222in}{3.059418in}}%
\pgfpathlineto{\pgfqpoint{1.254658in}{3.134849in}}%
\pgfpathlineto{\pgfqpoint{1.227119in}{3.210208in}}%
\pgfpathlineto{\pgfqpoint{1.277071in}{3.224470in}}%
\pgfpathlineto{\pgfqpoint{1.326932in}{3.238706in}}%
\pgfpathlineto{\pgfqpoint{1.354468in}{3.163496in}}%
\pgfpathlineto{\pgfqpoint{1.382031in}{3.088215in}}%
\pgfpathlineto{\pgfqpoint{1.332172in}{3.073830in}}%
\pgfpathlineto{\pgfqpoint{1.282222in}{3.059418in}}%
\pgfpathclose%
\pgfusepath{stroke,fill}%
\end{pgfscope}%
\begin{pgfscope}%
\pgfpathrectangle{\pgfqpoint{0.050000in}{0.083252in}}{\pgfqpoint{4.900000in}{4.900000in}}%
\pgfusepath{clip}%
\pgfsetbuttcap%
\pgfsetroundjoin%
\definecolor{currentfill}{rgb}{0.749389,0.739054,0.850150}%
\pgfsetfillcolor{currentfill}%
\pgfsetfillopacity{0.950000}%
\pgfsetlinewidth{0.200750pt}%
\definecolor{currentstroke}{rgb}{0.200000,0.200000,0.200000}%
\pgfsetstrokecolor{currentstroke}%
\pgfsetdash{}{0pt}%
\pgfpathmoveto{\pgfqpoint{2.488442in}{1.592270in}}%
\pgfpathlineto{\pgfqpoint{2.460312in}{1.662172in}}%
\pgfpathlineto{\pgfqpoint{2.432210in}{1.733767in}}%
\pgfpathlineto{\pgfqpoint{2.481065in}{1.750186in}}%
\pgfpathlineto{\pgfqpoint{2.529832in}{1.766588in}}%
\pgfpathlineto{\pgfqpoint{2.557940in}{1.695568in}}%
\pgfpathlineto{\pgfqpoint{2.586081in}{1.626261in}}%
\pgfpathlineto{\pgfqpoint{2.537305in}{1.609284in}}%
\pgfpathlineto{\pgfqpoint{2.488442in}{1.592270in}}%
\pgfpathclose%
\pgfusepath{stroke,fill}%
\end{pgfscope}%
\begin{pgfscope}%
\pgfpathrectangle{\pgfqpoint{0.050000in}{0.083252in}}{\pgfqpoint{4.900000in}{4.900000in}}%
\pgfusepath{clip}%
\pgfsetbuttcap%
\pgfsetroundjoin%
\definecolor{currentfill}{rgb}{0.247459,0.295425,0.661530}%
\pgfsetfillcolor{currentfill}%
\pgfsetfillopacity{0.950000}%
\pgfsetlinewidth{0.200750pt}%
\definecolor{currentstroke}{rgb}{0.200000,0.200000,0.200000}%
\pgfsetstrokecolor{currentstroke}%
\pgfsetdash{}{0pt}%
\pgfpathmoveto{\pgfqpoint{1.492545in}{2.786369in}}%
\pgfpathlineto{\pgfqpoint{1.464877in}{2.861939in}}%
\pgfpathlineto{\pgfqpoint{1.437235in}{2.937437in}}%
\pgfpathlineto{\pgfqpoint{1.487002in}{2.951946in}}%
\pgfpathlineto{\pgfqpoint{1.536679in}{2.966429in}}%
\pgfpathlineto{\pgfqpoint{1.564318in}{2.891082in}}%
\pgfpathlineto{\pgfqpoint{1.591984in}{2.815663in}}%
\pgfpathlineto{\pgfqpoint{1.542310in}{2.801029in}}%
\pgfpathlineto{\pgfqpoint{1.492545in}{2.786369in}}%
\pgfpathclose%
\pgfusepath{stroke,fill}%
\end{pgfscope}%
\begin{pgfscope}%
\pgfpathrectangle{\pgfqpoint{0.050000in}{0.083252in}}{\pgfqpoint{4.900000in}{4.900000in}}%
\pgfusepath{clip}%
\pgfsetbuttcap%
\pgfsetroundjoin%
\definecolor{currentfill}{rgb}{1.000000,1.000000,1.000000}%
\pgfsetfillcolor{currentfill}%
\pgfsetfillopacity{0.950000}%
\pgfsetlinewidth{0.200750pt}%
\definecolor{currentstroke}{rgb}{0.200000,0.200000,0.200000}%
\pgfsetstrokecolor{currentstroke}%
\pgfsetdash{}{0pt}%
\pgfpathmoveto{\pgfqpoint{3.917718in}{1.203328in}}%
\pgfpathlineto{\pgfqpoint{3.887516in}{1.234522in}}%
\pgfpathlineto{\pgfqpoint{3.857417in}{1.266074in}}%
\pgfpathlineto{\pgfqpoint{3.904942in}{1.281827in}}%
\pgfpathlineto{\pgfqpoint{3.952380in}{1.297554in}}%
\pgfpathlineto{\pgfqpoint{3.982550in}{1.265657in}}%
\pgfpathlineto{\pgfqpoint{4.012824in}{1.234109in}}%
\pgfpathlineto{\pgfqpoint{3.965314in}{1.218733in}}%
\pgfpathlineto{\pgfqpoint{3.917718in}{1.203328in}}%
\pgfpathclose%
\pgfusepath{stroke,fill}%
\end{pgfscope}%
\begin{pgfscope}%
\pgfpathrectangle{\pgfqpoint{0.050000in}{0.083252in}}{\pgfqpoint{4.900000in}{4.900000in}}%
\pgfusepath{clip}%
\pgfsetbuttcap%
\pgfsetroundjoin%
\definecolor{currentfill}{rgb}{0.240815,0.224837,0.567689}%
\pgfsetfillcolor{currentfill}%
\pgfsetfillopacity{0.950000}%
\pgfsetlinewidth{0.200750pt}%
\definecolor{currentstroke}{rgb}{0.200000,0.200000,0.200000}%
\pgfsetstrokecolor{currentstroke}%
\pgfsetdash{}{0pt}%
\pgfpathmoveto{\pgfqpoint{1.702913in}{2.513263in}}%
\pgfpathlineto{\pgfqpoint{1.675141in}{2.588971in}}%
\pgfpathlineto{\pgfqpoint{1.647395in}{2.664607in}}%
\pgfpathlineto{\pgfqpoint{1.696977in}{2.679365in}}%
\pgfpathlineto{\pgfqpoint{1.746469in}{2.694096in}}%
\pgfpathlineto{\pgfqpoint{1.774212in}{2.618610in}}%
\pgfpathlineto{\pgfqpoint{1.801981in}{2.543053in}}%
\pgfpathlineto{\pgfqpoint{1.752492in}{2.528171in}}%
\pgfpathlineto{\pgfqpoint{1.702913in}{2.513263in}}%
\pgfpathclose%
\pgfusepath{stroke,fill}%
\end{pgfscope}%
\begin{pgfscope}%
\pgfpathrectangle{\pgfqpoint{0.050000in}{0.083252in}}{\pgfqpoint{4.900000in}{4.900000in}}%
\pgfusepath{clip}%
\pgfsetbuttcap%
\pgfsetroundjoin%
\definecolor{currentfill}{rgb}{0.361538,0.335210,0.618239}%
\pgfsetfillcolor{currentfill}%
\pgfsetfillopacity{0.950000}%
\pgfsetlinewidth{0.200750pt}%
\definecolor{currentstroke}{rgb}{0.200000,0.200000,0.200000}%
\pgfsetstrokecolor{currentstroke}%
\pgfsetdash{}{0pt}%
\pgfpathmoveto{\pgfqpoint{1.913324in}{2.240118in}}%
\pgfpathlineto{\pgfqpoint{1.885448in}{2.315955in}}%
\pgfpathlineto{\pgfqpoint{1.857599in}{2.391724in}}%
\pgfpathlineto{\pgfqpoint{1.906996in}{2.406732in}}%
\pgfpathlineto{\pgfqpoint{1.956304in}{2.421715in}}%
\pgfpathlineto{\pgfqpoint{1.984150in}{2.346103in}}%
\pgfpathlineto{\pgfqpoint{2.012022in}{2.270431in}}%
\pgfpathlineto{\pgfqpoint{1.962717in}{2.255286in}}%
\pgfpathlineto{\pgfqpoint{1.913324in}{2.240118in}}%
\pgfpathclose%
\pgfusepath{stroke,fill}%
\end{pgfscope}%
\begin{pgfscope}%
\pgfpathrectangle{\pgfqpoint{0.050000in}{0.083252in}}{\pgfqpoint{4.900000in}{4.900000in}}%
\pgfusepath{clip}%
\pgfsetbuttcap%
\pgfsetroundjoin%
\definecolor{currentfill}{rgb}{0.522645,0.502960,0.714571}%
\pgfsetfillcolor{currentfill}%
\pgfsetfillopacity{0.950000}%
\pgfsetlinewidth{0.200750pt}%
\definecolor{currentstroke}{rgb}{0.200000,0.200000,0.200000}%
\pgfsetstrokecolor{currentstroke}%
\pgfsetdash{}{0pt}%
\pgfpathmoveto{\pgfqpoint{2.123773in}{1.967520in}}%
\pgfpathlineto{\pgfqpoint{2.095797in}{2.043209in}}%
\pgfpathlineto{\pgfqpoint{2.067845in}{2.118957in}}%
\pgfpathlineto{\pgfqpoint{2.117057in}{2.134270in}}%
\pgfpathlineto{\pgfqpoint{2.166180in}{2.149565in}}%
\pgfpathlineto{\pgfqpoint{2.194127in}{2.074052in}}%
\pgfpathlineto{\pgfqpoint{2.222100in}{1.998641in}}%
\pgfpathlineto{\pgfqpoint{2.172981in}{1.983085in}}%
\pgfpathlineto{\pgfqpoint{2.123773in}{1.967520in}}%
\pgfpathclose%
\pgfusepath{stroke,fill}%
\end{pgfscope}%
\begin{pgfscope}%
\pgfpathrectangle{\pgfqpoint{0.050000in}{0.083252in}}{\pgfqpoint{4.900000in}{4.900000in}}%
\pgfusepath{clip}%
\pgfsetbuttcap%
\pgfsetroundjoin%
\definecolor{currentfill}{rgb}{0.964198,0.962722,0.978593}%
\pgfsetfillcolor{currentfill}%
\pgfsetfillopacity{0.950000}%
\pgfsetlinewidth{0.200750pt}%
\definecolor{currentstroke}{rgb}{0.200000,0.200000,0.200000}%
\pgfsetstrokecolor{currentstroke}%
\pgfsetdash{}{0pt}%
\pgfpathmoveto{\pgfqpoint{3.259704in}{1.239212in}}%
\pgfpathlineto{\pgfqpoint{3.230582in}{1.277462in}}%
\pgfpathlineto{\pgfqpoint{3.201580in}{1.318547in}}%
\pgfpathlineto{\pgfqpoint{3.249754in}{1.335612in}}%
\pgfpathlineto{\pgfqpoint{3.297843in}{1.352632in}}%
\pgfpathlineto{\pgfqpoint{3.326909in}{1.311243in}}%
\pgfpathlineto{\pgfqpoint{3.356100in}{1.272639in}}%
\pgfpathlineto{\pgfqpoint{3.307945in}{1.255942in}}%
\pgfpathlineto{\pgfqpoint{3.259704in}{1.239212in}}%
\pgfpathclose%
\pgfusepath{stroke,fill}%
\end{pgfscope}%
\begin{pgfscope}%
\pgfpathrectangle{\pgfqpoint{0.050000in}{0.083252in}}{\pgfqpoint{4.900000in}{4.900000in}}%
\pgfusepath{clip}%
\pgfsetbuttcap%
\pgfsetroundjoin%
\definecolor{currentfill}{rgb}{0.677785,0.664498,0.807336}%
\pgfsetfillcolor{currentfill}%
\pgfsetfillopacity{0.950000}%
\pgfsetlinewidth{0.200750pt}%
\definecolor{currentstroke}{rgb}{0.200000,0.200000,0.200000}%
\pgfsetstrokecolor{currentstroke}%
\pgfsetdash{}{0pt}%
\pgfpathmoveto{\pgfqpoint{2.334242in}{1.700885in}}%
\pgfpathlineto{\pgfqpoint{2.306170in}{1.774310in}}%
\pgfpathlineto{\pgfqpoint{2.278122in}{1.848588in}}%
\pgfpathlineto{\pgfqpoint{2.327150in}{1.864514in}}%
\pgfpathlineto{\pgfqpoint{2.376090in}{1.880433in}}%
\pgfpathlineto{\pgfqpoint{2.404137in}{1.806638in}}%
\pgfpathlineto{\pgfqpoint{2.432210in}{1.733767in}}%
\pgfpathlineto{\pgfqpoint{2.383270in}{1.717333in}}%
\pgfpathlineto{\pgfqpoint{2.334242in}{1.700885in}}%
\pgfpathclose%
\pgfusepath{stroke,fill}%
\end{pgfscope}%
\begin{pgfscope}%
\pgfpathrectangle{\pgfqpoint{0.050000in}{0.083252in}}{\pgfqpoint{4.900000in}{4.900000in}}%
\pgfusepath{clip}%
\pgfsetbuttcap%
\pgfsetroundjoin%
\definecolor{currentfill}{rgb}{0.263699,0.467974,0.890919}%
\pgfsetfillcolor{currentfill}%
\pgfsetfillopacity{0.950000}%
\pgfsetlinewidth{0.200750pt}%
\definecolor{currentstroke}{rgb}{0.200000,0.200000,0.200000}%
\pgfsetstrokecolor{currentstroke}%
\pgfsetdash{}{0pt}%
\pgfpathmoveto{\pgfqpoint{0.916503in}{3.454817in}}%
\pgfpathlineto{\pgfqpoint{0.889093in}{3.530115in}}%
\pgfpathlineto{\pgfqpoint{0.861709in}{3.605341in}}%
\pgfpathlineto{\pgfqpoint{0.912028in}{3.619257in}}%
\pgfpathlineto{\pgfqpoint{0.962256in}{3.633147in}}%
\pgfpathlineto{\pgfqpoint{0.989639in}{3.558070in}}%
\pgfpathlineto{\pgfqpoint{1.017047in}{3.482922in}}%
\pgfpathlineto{\pgfqpoint{0.966821in}{3.468882in}}%
\pgfpathlineto{\pgfqpoint{0.916503in}{3.454817in}}%
\pgfpathclose%
\pgfusepath{stroke,fill}%
\end{pgfscope}%
\begin{pgfscope}%
\pgfpathrectangle{\pgfqpoint{0.050000in}{0.083252in}}{\pgfqpoint{4.900000in}{4.900000in}}%
\pgfusepath{clip}%
\pgfsetbuttcap%
\pgfsetroundjoin%
\definecolor{currentfill}{rgb}{0.994033,0.993787,0.996432}%
\pgfsetfillcolor{currentfill}%
\pgfsetfillopacity{0.950000}%
\pgfsetlinewidth{0.200750pt}%
\definecolor{currentstroke}{rgb}{0.200000,0.200000,0.200000}%
\pgfsetstrokecolor{currentstroke}%
\pgfsetdash{}{0pt}%
\pgfpathmoveto{\pgfqpoint{3.666452in}{1.202799in}}%
\pgfpathlineto{\pgfqpoint{3.636618in}{1.234522in}}%
\pgfpathlineto{\pgfqpoint{3.606905in}{1.267544in}}%
\pgfpathlineto{\pgfqpoint{3.654701in}{1.283744in}}%
\pgfpathlineto{\pgfqpoint{3.702412in}{1.299917in}}%
\pgfpathlineto{\pgfqpoint{3.732199in}{1.266552in}}%
\pgfpathlineto{\pgfqpoint{3.762108in}{1.234488in}}%
\pgfpathlineto{\pgfqpoint{3.714324in}{1.218656in}}%
\pgfpathlineto{\pgfqpoint{3.666452in}{1.202799in}}%
\pgfpathclose%
\pgfusepath{stroke,fill}%
\end{pgfscope}%
\begin{pgfscope}%
\pgfpathrectangle{\pgfqpoint{0.050000in}{0.083252in}}{\pgfqpoint{4.900000in}{4.900000in}}%
\pgfusepath{clip}%
\pgfsetbuttcap%
\pgfsetroundjoin%
\definecolor{currentfill}{rgb}{0.257301,0.400000,0.800554}%
\pgfsetfillcolor{currentfill}%
\pgfsetfillopacity{0.950000}%
\pgfsetlinewidth{0.200750pt}%
\definecolor{currentstroke}{rgb}{0.200000,0.200000,0.200000}%
\pgfsetstrokecolor{currentstroke}%
\pgfsetdash{}{0pt}%
\pgfpathmoveto{\pgfqpoint{1.126945in}{3.181607in}}%
\pgfpathlineto{\pgfqpoint{1.099431in}{3.257044in}}%
\pgfpathlineto{\pgfqpoint{1.071943in}{3.332409in}}%
\pgfpathlineto{\pgfqpoint{1.122078in}{3.346572in}}%
\pgfpathlineto{\pgfqpoint{1.172121in}{3.360710in}}%
\pgfpathlineto{\pgfqpoint{1.199607in}{3.285495in}}%
\pgfpathlineto{\pgfqpoint{1.227119in}{3.210208in}}%
\pgfpathlineto{\pgfqpoint{1.177077in}{3.195920in}}%
\pgfpathlineto{\pgfqpoint{1.126945in}{3.181607in}}%
\pgfpathclose%
\pgfusepath{stroke,fill}%
\end{pgfscope}%
\begin{pgfscope}%
\pgfpathrectangle{\pgfqpoint{0.050000in}{0.083252in}}{\pgfqpoint{4.900000in}{4.900000in}}%
\pgfusepath{clip}%
\pgfsetbuttcap%
\pgfsetroundjoin%
\definecolor{currentfill}{rgb}{0.269604,0.530719,0.974333}%
\pgfsetfillcolor{currentfill}%
\pgfsetfillopacity{0.950000}%
\pgfsetlinewidth{0.200750pt}%
\definecolor{currentstroke}{rgb}{0.200000,0.200000,0.200000}%
\pgfsetstrokecolor{currentstroke}%
\pgfsetdash{}{0pt}%
\pgfpathmoveto{\pgfqpoint{0.706105in}{3.727970in}}%
\pgfpathlineto{\pgfqpoint{0.678798in}{3.803129in}}%
\pgfpathlineto{\pgfqpoint{0.729302in}{3.816871in}}%
\pgfpathlineto{\pgfqpoint{0.779714in}{3.830588in}}%
\pgfpathlineto{\pgfqpoint{0.807019in}{3.755578in}}%
\pgfpathlineto{\pgfqpoint{0.756608in}{3.741786in}}%
\pgfpathlineto{\pgfqpoint{0.706105in}{3.727970in}}%
\pgfpathclose%
\pgfusepath{stroke,fill}%
\end{pgfscope}%
\begin{pgfscope}%
\pgfpathrectangle{\pgfqpoint{0.050000in}{0.083252in}}{\pgfqpoint{4.900000in}{4.900000in}}%
\pgfusepath{clip}%
\pgfsetbuttcap%
\pgfsetroundjoin%
\definecolor{currentfill}{rgb}{0.250657,0.329412,0.706713}%
\pgfsetfillcolor{currentfill}%
\pgfsetfillopacity{0.950000}%
\pgfsetlinewidth{0.200750pt}%
\definecolor{currentstroke}{rgb}{0.200000,0.200000,0.200000}%
\pgfsetstrokecolor{currentstroke}%
\pgfsetdash{}{0pt}%
\pgfpathmoveto{\pgfqpoint{1.337431in}{2.908338in}}%
\pgfpathlineto{\pgfqpoint{1.309814in}{2.983914in}}%
\pgfpathlineto{\pgfqpoint{1.282222in}{3.059418in}}%
\pgfpathlineto{\pgfqpoint{1.332172in}{3.073830in}}%
\pgfpathlineto{\pgfqpoint{1.382031in}{3.088215in}}%
\pgfpathlineto{\pgfqpoint{1.409620in}{3.012862in}}%
\pgfpathlineto{\pgfqpoint{1.437235in}{2.937437in}}%
\pgfpathlineto{\pgfqpoint{1.387378in}{2.922901in}}%
\pgfpathlineto{\pgfqpoint{1.337431in}{2.908338in}}%
\pgfpathclose%
\pgfusepath{stroke,fill}%
\end{pgfscope}%
\begin{pgfscope}%
\pgfpathrectangle{\pgfqpoint{0.050000in}{0.083252in}}{\pgfqpoint{4.900000in}{4.900000in}}%
\pgfusepath{clip}%
\pgfsetbuttcap%
\pgfsetroundjoin%
\definecolor{currentfill}{rgb}{0.244014,0.258824,0.612872}%
\pgfsetfillcolor{currentfill}%
\pgfsetfillopacity{0.950000}%
\pgfsetlinewidth{0.200750pt}%
\definecolor{currentstroke}{rgb}{0.200000,0.200000,0.200000}%
\pgfsetstrokecolor{currentstroke}%
\pgfsetdash{}{0pt}%
\pgfpathmoveto{\pgfqpoint{1.547962in}{2.635012in}}%
\pgfpathlineto{\pgfqpoint{1.520240in}{2.710727in}}%
\pgfpathlineto{\pgfqpoint{1.492545in}{2.786369in}}%
\pgfpathlineto{\pgfqpoint{1.542310in}{2.801029in}}%
\pgfpathlineto{\pgfqpoint{1.591984in}{2.815663in}}%
\pgfpathlineto{\pgfqpoint{1.619676in}{2.740171in}}%
\pgfpathlineto{\pgfqpoint{1.647395in}{2.664607in}}%
\pgfpathlineto{\pgfqpoint{1.597724in}{2.649823in}}%
\pgfpathlineto{\pgfqpoint{1.547962in}{2.635012in}}%
\pgfpathclose%
\pgfusepath{stroke,fill}%
\end{pgfscope}%
\begin{pgfscope}%
\pgfpathrectangle{\pgfqpoint{0.050000in}{0.083252in}}{\pgfqpoint{4.900000in}{4.900000in}}%
\pgfusepath{clip}%
\pgfsetbuttcap%
\pgfsetroundjoin%
\definecolor{currentfill}{rgb}{0.868727,0.863314,0.921507}%
\pgfsetfillcolor{currentfill}%
\pgfsetfillopacity{0.950000}%
\pgfsetlinewidth{0.200750pt}%
\definecolor{currentstroke}{rgb}{0.200000,0.200000,0.200000}%
\pgfsetstrokecolor{currentstroke}%
\pgfsetdash{}{0pt}%
\pgfpathmoveto{\pgfqpoint{2.699081in}{1.374884in}}%
\pgfpathlineto{\pgfqpoint{2.670754in}{1.433001in}}%
\pgfpathlineto{\pgfqpoint{2.642483in}{1.494534in}}%
\pgfpathlineto{\pgfqpoint{2.691194in}{1.511983in}}%
\pgfpathlineto{\pgfqpoint{2.739822in}{1.529372in}}%
\pgfpathlineto{\pgfqpoint{2.768123in}{1.468214in}}%
\pgfpathlineto{\pgfqpoint{2.796488in}{1.410338in}}%
\pgfpathlineto{\pgfqpoint{2.747827in}{1.392649in}}%
\pgfpathlineto{\pgfqpoint{2.699081in}{1.374884in}}%
\pgfpathclose%
\pgfusepath{stroke,fill}%
\end{pgfscope}%
\begin{pgfscope}%
\pgfpathrectangle{\pgfqpoint{0.050000in}{0.083252in}}{\pgfqpoint{4.900000in}{4.900000in}}%
\pgfusepath{clip}%
\pgfsetbuttcap%
\pgfsetroundjoin%
\definecolor{currentfill}{rgb}{0.283968,0.254441,0.571857}%
\pgfsetfillcolor{currentfill}%
\pgfsetfillopacity{0.950000}%
\pgfsetlinewidth{0.200750pt}%
\definecolor{currentstroke}{rgb}{0.200000,0.200000,0.200000}%
\pgfsetstrokecolor{currentstroke}%
\pgfsetdash{}{0pt}%
\pgfpathmoveto{\pgfqpoint{1.758537in}{2.361629in}}%
\pgfpathlineto{\pgfqpoint{1.730711in}{2.437482in}}%
\pgfpathlineto{\pgfqpoint{1.702913in}{2.513263in}}%
\pgfpathlineto{\pgfqpoint{1.752492in}{2.528171in}}%
\pgfpathlineto{\pgfqpoint{1.801981in}{2.543053in}}%
\pgfpathlineto{\pgfqpoint{1.829777in}{2.467424in}}%
\pgfpathlineto{\pgfqpoint{1.857599in}{2.391724in}}%
\pgfpathlineto{\pgfqpoint{1.808113in}{2.376690in}}%
\pgfpathlineto{\pgfqpoint{1.758537in}{2.361629in}}%
\pgfpathclose%
\pgfusepath{stroke,fill}%
\end{pgfscope}%
\begin{pgfscope}%
\pgfpathrectangle{\pgfqpoint{0.050000in}{0.083252in}}{\pgfqpoint{4.900000in}{4.900000in}}%
\pgfusepath{clip}%
\pgfsetbuttcap%
\pgfsetroundjoin%
\definecolor{currentfill}{rgb}{0.916463,0.913018,0.950050}%
\pgfsetfillcolor{currentfill}%
\pgfsetfillopacity{0.950000}%
\pgfsetlinewidth{0.200750pt}%
\definecolor{currentstroke}{rgb}{0.200000,0.200000,0.200000}%
\pgfsetstrokecolor{currentstroke}%
\pgfsetdash{}{0pt}%
\pgfpathmoveto{\pgfqpoint{2.853441in}{1.305384in}}%
\pgfpathlineto{\pgfqpoint{2.824925in}{1.356011in}}%
\pgfpathlineto{\pgfqpoint{2.796488in}{1.410338in}}%
\pgfpathlineto{\pgfqpoint{2.845065in}{1.427956in}}%
\pgfpathlineto{\pgfqpoint{2.893558in}{1.445504in}}%
\pgfpathlineto{\pgfqpoint{2.922039in}{1.391259in}}%
\pgfpathlineto{\pgfqpoint{2.950606in}{1.340567in}}%
\pgfpathlineto{\pgfqpoint{2.902066in}{1.323009in}}%
\pgfpathlineto{\pgfqpoint{2.853441in}{1.305384in}}%
\pgfpathclose%
\pgfusepath{stroke,fill}%
\end{pgfscope}%
\begin{pgfscope}%
\pgfpathrectangle{\pgfqpoint{0.050000in}{0.083252in}}{\pgfqpoint{4.900000in}{4.900000in}}%
\pgfusepath{clip}%
\pgfsetbuttcap%
\pgfsetroundjoin%
\definecolor{currentfill}{rgb}{0.439108,0.415978,0.664621}%
\pgfsetfillcolor{currentfill}%
\pgfsetfillopacity{0.950000}%
\pgfsetlinewidth{0.200750pt}%
\definecolor{currentstroke}{rgb}{0.200000,0.200000,0.200000}%
\pgfsetstrokecolor{currentstroke}%
\pgfsetdash{}{0pt}%
\pgfpathmoveto{\pgfqpoint{1.969154in}{2.088279in}}%
\pgfpathlineto{\pgfqpoint{1.941226in}{2.164222in}}%
\pgfpathlineto{\pgfqpoint{1.913324in}{2.240118in}}%
\pgfpathlineto{\pgfqpoint{1.962717in}{2.255286in}}%
\pgfpathlineto{\pgfqpoint{2.012022in}{2.270431in}}%
\pgfpathlineto{\pgfqpoint{2.039921in}{2.194709in}}%
\pgfpathlineto{\pgfqpoint{2.067845in}{2.118957in}}%
\pgfpathlineto{\pgfqpoint{2.018545in}{2.103627in}}%
\pgfpathlineto{\pgfqpoint{1.969154in}{2.088279in}}%
\pgfpathclose%
\pgfusepath{stroke,fill}%
\end{pgfscope}%
\begin{pgfscope}%
\pgfpathrectangle{\pgfqpoint{0.050000in}{0.083252in}}{\pgfqpoint{4.900000in}{4.900000in}}%
\pgfusepath{clip}%
\pgfsetbuttcap%
\pgfsetroundjoin%
\definecolor{currentfill}{rgb}{0.815025,0.807397,0.889396}%
\pgfsetfillcolor{currentfill}%
\pgfsetfillopacity{0.950000}%
\pgfsetlinewidth{0.200750pt}%
\definecolor{currentstroke}{rgb}{0.200000,0.200000,0.200000}%
\pgfsetstrokecolor{currentstroke}%
\pgfsetdash{}{0pt}%
\pgfpathmoveto{\pgfqpoint{2.544804in}{1.459446in}}%
\pgfpathlineto{\pgfqpoint{2.516605in}{1.524530in}}%
\pgfpathlineto{\pgfqpoint{2.488442in}{1.592270in}}%
\pgfpathlineto{\pgfqpoint{2.537305in}{1.609284in}}%
\pgfpathlineto{\pgfqpoint{2.586081in}{1.626261in}}%
\pgfpathlineto{\pgfqpoint{2.614261in}{1.559099in}}%
\pgfpathlineto{\pgfqpoint{2.642483in}{1.494534in}}%
\pgfpathlineto{\pgfqpoint{2.593686in}{1.477022in}}%
\pgfpathlineto{\pgfqpoint{2.544804in}{1.459446in}}%
\pgfpathclose%
\pgfusepath{stroke,fill}%
\end{pgfscope}%
\begin{pgfscope}%
\pgfpathrectangle{\pgfqpoint{0.050000in}{0.083252in}}{\pgfqpoint{4.900000in}{4.900000in}}%
\pgfusepath{clip}%
\pgfsetbuttcap%
\pgfsetroundjoin%
\definecolor{currentfill}{rgb}{0.600215,0.583729,0.760953}%
\pgfsetfillcolor{currentfill}%
\pgfsetfillopacity{0.950000}%
\pgfsetlinewidth{0.200750pt}%
\definecolor{currentstroke}{rgb}{0.200000,0.200000,0.200000}%
\pgfsetstrokecolor{currentstroke}%
\pgfsetdash{}{0pt}%
\pgfpathmoveto{\pgfqpoint{2.179801in}{1.816726in}}%
\pgfpathlineto{\pgfqpoint{2.151775in}{1.891980in}}%
\pgfpathlineto{\pgfqpoint{2.123773in}{1.967520in}}%
\pgfpathlineto{\pgfqpoint{2.172981in}{1.983085in}}%
\pgfpathlineto{\pgfqpoint{2.222100in}{1.998641in}}%
\pgfpathlineto{\pgfqpoint{2.250099in}{1.923436in}}%
\pgfpathlineto{\pgfqpoint{2.278122in}{1.848588in}}%
\pgfpathlineto{\pgfqpoint{2.229006in}{1.832658in}}%
\pgfpathlineto{\pgfqpoint{2.179801in}{1.816726in}}%
\pgfpathclose%
\pgfusepath{stroke,fill}%
\end{pgfscope}%
\begin{pgfscope}%
\pgfpathrectangle{\pgfqpoint{0.050000in}{0.083252in}}{\pgfqpoint{4.900000in}{4.900000in}}%
\pgfusepath{clip}%
\pgfsetbuttcap%
\pgfsetroundjoin%
\definecolor{currentfill}{rgb}{0.982099,0.981361,0.989296}%
\pgfsetfillcolor{currentfill}%
\pgfsetfillopacity{0.950000}%
\pgfsetlinewidth{0.200750pt}%
\definecolor{currentstroke}{rgb}{0.200000,0.200000,0.200000}%
\pgfsetstrokecolor{currentstroke}%
\pgfsetdash{}{0pt}%
\pgfpathmoveto{\pgfqpoint{3.414854in}{1.202488in}}%
\pgfpathlineto{\pgfqpoint{3.385415in}{1.236509in}}%
\pgfpathlineto{\pgfqpoint{3.356100in}{1.272639in}}%
\pgfpathlineto{\pgfqpoint{3.404169in}{1.289304in}}%
\pgfpathlineto{\pgfqpoint{3.452152in}{1.305937in}}%
\pgfpathlineto{\pgfqpoint{3.481539in}{1.269451in}}%
\pgfpathlineto{\pgfqpoint{3.511052in}{1.235065in}}%
\pgfpathlineto{\pgfqpoint{3.462997in}{1.218789in}}%
\pgfpathlineto{\pgfqpoint{3.414854in}{1.202488in}}%
\pgfpathclose%
\pgfusepath{stroke,fill}%
\end{pgfscope}%
\begin{pgfscope}%
\pgfpathrectangle{\pgfqpoint{0.050000in}{0.083252in}}{\pgfqpoint{4.900000in}{4.900000in}}%
\pgfusepath{clip}%
\pgfsetbuttcap%
\pgfsetroundjoin%
\definecolor{currentfill}{rgb}{0.267143,0.504575,0.939577}%
\pgfsetfillcolor{currentfill}%
\pgfsetfillopacity{0.950000}%
\pgfsetlinewidth{0.200750pt}%
\definecolor{currentstroke}{rgb}{0.200000,0.200000,0.200000}%
\pgfsetstrokecolor{currentstroke}%
\pgfsetdash{}{0pt}%
\pgfpathmoveto{\pgfqpoint{0.760796in}{3.577434in}}%
\pgfpathlineto{\pgfqpoint{0.733438in}{3.652738in}}%
\pgfpathlineto{\pgfqpoint{0.706105in}{3.727970in}}%
\pgfpathlineto{\pgfqpoint{0.756608in}{3.741786in}}%
\pgfpathlineto{\pgfqpoint{0.807019in}{3.755578in}}%
\pgfpathlineto{\pgfqpoint{0.834351in}{3.680496in}}%
\pgfpathlineto{\pgfqpoint{0.861709in}{3.605341in}}%
\pgfpathlineto{\pgfqpoint{0.811298in}{3.591401in}}%
\pgfpathlineto{\pgfqpoint{0.760796in}{3.577434in}}%
\pgfpathclose%
\pgfusepath{stroke,fill}%
\end{pgfscope}%
\begin{pgfscope}%
\pgfpathrectangle{\pgfqpoint{0.050000in}{0.083252in}}{\pgfqpoint{4.900000in}{4.900000in}}%
\pgfusepath{clip}%
\pgfsetbuttcap%
\pgfsetroundjoin%
\definecolor{currentfill}{rgb}{0.946298,0.944083,0.967889}%
\pgfsetfillcolor{currentfill}%
\pgfsetfillopacity{0.950000}%
\pgfsetlinewidth{0.200750pt}%
\definecolor{currentstroke}{rgb}{0.200000,0.200000,0.200000}%
\pgfsetstrokecolor{currentstroke}%
\pgfsetdash{}{0pt}%
\pgfpathmoveto{\pgfqpoint{3.008031in}{1.249818in}}%
\pgfpathlineto{\pgfqpoint{2.979268in}{1.293454in}}%
\pgfpathlineto{\pgfqpoint{2.950606in}{1.340567in}}%
\pgfpathlineto{\pgfqpoint{2.999062in}{1.358062in}}%
\pgfpathlineto{\pgfqpoint{3.047433in}{1.375493in}}%
\pgfpathlineto{\pgfqpoint{3.076151in}{1.328201in}}%
\pgfpathlineto{\pgfqpoint{3.104976in}{1.284276in}}%
\pgfpathlineto{\pgfqpoint{3.056546in}{1.267071in}}%
\pgfpathlineto{\pgfqpoint{3.008031in}{1.249818in}}%
\pgfpathclose%
\pgfusepath{stroke,fill}%
\end{pgfscope}%
\begin{pgfscope}%
\pgfpathrectangle{\pgfqpoint{0.050000in}{0.083252in}}{\pgfqpoint{4.900000in}{4.900000in}}%
\pgfusepath{clip}%
\pgfsetbuttcap%
\pgfsetroundjoin%
\definecolor{currentfill}{rgb}{0.260500,0.433987,0.845736}%
\pgfsetfillcolor{currentfill}%
\pgfsetfillopacity{0.950000}%
\pgfsetlinewidth{0.200750pt}%
\definecolor{currentstroke}{rgb}{0.200000,0.200000,0.200000}%
\pgfsetstrokecolor{currentstroke}%
\pgfsetdash{}{0pt}%
\pgfpathmoveto{\pgfqpoint{0.971402in}{3.304004in}}%
\pgfpathlineto{\pgfqpoint{0.943939in}{3.379447in}}%
\pgfpathlineto{\pgfqpoint{0.916503in}{3.454817in}}%
\pgfpathlineto{\pgfqpoint{0.966821in}{3.468882in}}%
\pgfpathlineto{\pgfqpoint{1.017047in}{3.482922in}}%
\pgfpathlineto{\pgfqpoint{1.044482in}{3.407701in}}%
\pgfpathlineto{\pgfqpoint{1.071943in}{3.332409in}}%
\pgfpathlineto{\pgfqpoint{1.021718in}{3.318219in}}%
\pgfpathlineto{\pgfqpoint{0.971402in}{3.304004in}}%
\pgfpathclose%
\pgfusepath{stroke,fill}%
\end{pgfscope}%
\begin{pgfscope}%
\pgfpathrectangle{\pgfqpoint{0.050000in}{0.083252in}}{\pgfqpoint{4.900000in}{4.900000in}}%
\pgfusepath{clip}%
\pgfsetbuttcap%
\pgfsetroundjoin%
\definecolor{currentfill}{rgb}{0.749389,0.739054,0.850150}%
\pgfsetfillcolor{currentfill}%
\pgfsetfillopacity{0.950000}%
\pgfsetlinewidth{0.200750pt}%
\definecolor{currentstroke}{rgb}{0.200000,0.200000,0.200000}%
\pgfsetstrokecolor{currentstroke}%
\pgfsetdash{}{0pt}%
\pgfpathmoveto{\pgfqpoint{2.390460in}{1.558133in}}%
\pgfpathlineto{\pgfqpoint{2.362338in}{1.628678in}}%
\pgfpathlineto{\pgfqpoint{2.334242in}{1.700885in}}%
\pgfpathlineto{\pgfqpoint{2.383270in}{1.717333in}}%
\pgfpathlineto{\pgfqpoint{2.432210in}{1.733767in}}%
\pgfpathlineto{\pgfqpoint{2.460312in}{1.662172in}}%
\pgfpathlineto{\pgfqpoint{2.488442in}{1.592270in}}%
\pgfpathlineto{\pgfqpoint{2.439494in}{1.575219in}}%
\pgfpathlineto{\pgfqpoint{2.390460in}{1.558133in}}%
\pgfpathclose%
\pgfusepath{stroke,fill}%
\end{pgfscope}%
\begin{pgfscope}%
\pgfpathrectangle{\pgfqpoint{0.050000in}{0.083252in}}{\pgfqpoint{4.900000in}{4.900000in}}%
\pgfusepath{clip}%
\pgfsetbuttcap%
\pgfsetroundjoin%
\definecolor{currentfill}{rgb}{0.253856,0.363399,0.751895}%
\pgfsetfillcolor{currentfill}%
\pgfsetfillopacity{0.950000}%
\pgfsetlinewidth{0.200750pt}%
\definecolor{currentstroke}{rgb}{0.200000,0.200000,0.200000}%
\pgfsetstrokecolor{currentstroke}%
\pgfsetdash{}{0pt}%
\pgfpathmoveto{\pgfqpoint{1.182052in}{3.030515in}}%
\pgfpathlineto{\pgfqpoint{1.154485in}{3.106097in}}%
\pgfpathlineto{\pgfqpoint{1.126945in}{3.181607in}}%
\pgfpathlineto{\pgfqpoint{1.177077in}{3.195920in}}%
\pgfpathlineto{\pgfqpoint{1.227119in}{3.210208in}}%
\pgfpathlineto{\pgfqpoint{1.254658in}{3.134849in}}%
\pgfpathlineto{\pgfqpoint{1.282222in}{3.059418in}}%
\pgfpathlineto{\pgfqpoint{1.232182in}{3.044980in}}%
\pgfpathlineto{\pgfqpoint{1.182052in}{3.030515in}}%
\pgfpathclose%
\pgfusepath{stroke,fill}%
\end{pgfscope}%
\begin{pgfscope}%
\pgfpathrectangle{\pgfqpoint{0.050000in}{0.083252in}}{\pgfqpoint{4.900000in}{4.900000in}}%
\pgfusepath{clip}%
\pgfsetbuttcap%
\pgfsetroundjoin%
\definecolor{currentfill}{rgb}{0.247213,0.292810,0.658055}%
\pgfsetfillcolor{currentfill}%
\pgfsetfillopacity{0.950000}%
\pgfsetlinewidth{0.200750pt}%
\definecolor{currentstroke}{rgb}{0.200000,0.200000,0.200000}%
\pgfsetstrokecolor{currentstroke}%
\pgfsetdash{}{0pt}%
\pgfpathmoveto{\pgfqpoint{1.392746in}{2.756969in}}%
\pgfpathlineto{\pgfqpoint{1.365075in}{2.832690in}}%
\pgfpathlineto{\pgfqpoint{1.337431in}{2.908338in}}%
\pgfpathlineto{\pgfqpoint{1.387378in}{2.922901in}}%
\pgfpathlineto{\pgfqpoint{1.437235in}{2.937437in}}%
\pgfpathlineto{\pgfqpoint{1.464877in}{2.861939in}}%
\pgfpathlineto{\pgfqpoint{1.492545in}{2.786369in}}%
\pgfpathlineto{\pgfqpoint{1.442691in}{2.771682in}}%
\pgfpathlineto{\pgfqpoint{1.392746in}{2.756969in}}%
\pgfpathclose%
\pgfusepath{stroke,fill}%
\end{pgfscope}%
\begin{pgfscope}%
\pgfpathrectangle{\pgfqpoint{0.050000in}{0.083252in}}{\pgfqpoint{4.900000in}{4.900000in}}%
\pgfusepath{clip}%
\pgfsetbuttcap%
\pgfsetroundjoin%
\definecolor{currentfill}{rgb}{1.000000,1.000000,1.000000}%
\pgfsetfillcolor{currentfill}%
\pgfsetfillopacity{0.950000}%
\pgfsetlinewidth{0.200750pt}%
\definecolor{currentstroke}{rgb}{0.200000,0.200000,0.200000}%
\pgfsetstrokecolor{currentstroke}%
\pgfsetdash{}{0pt}%
\pgfpathmoveto{\pgfqpoint{3.822260in}{1.172434in}}%
\pgfpathlineto{\pgfqpoint{3.792133in}{1.203279in}}%
\pgfpathlineto{\pgfqpoint{3.762108in}{1.234488in}}%
\pgfpathlineto{\pgfqpoint{3.809806in}{1.250294in}}%
\pgfpathlineto{\pgfqpoint{3.857417in}{1.266074in}}%
\pgfpathlineto{\pgfqpoint{3.887516in}{1.234522in}}%
\pgfpathlineto{\pgfqpoint{3.917718in}{1.203328in}}%
\pgfpathlineto{\pgfqpoint{3.870033in}{1.187895in}}%
\pgfpathlineto{\pgfqpoint{3.822260in}{1.172434in}}%
\pgfpathclose%
\pgfusepath{stroke,fill}%
\end{pgfscope}%
\begin{pgfscope}%
\pgfpathrectangle{\pgfqpoint{0.050000in}{0.083252in}}{\pgfqpoint{4.900000in}{4.900000in}}%
\pgfusepath{clip}%
\pgfsetbuttcap%
\pgfsetroundjoin%
\definecolor{currentfill}{rgb}{0.240815,0.224837,0.567689}%
\pgfsetfillcolor{currentfill}%
\pgfsetfillopacity{0.950000}%
\pgfsetlinewidth{0.200750pt}%
\definecolor{currentstroke}{rgb}{0.200000,0.200000,0.200000}%
\pgfsetstrokecolor{currentstroke}%
\pgfsetdash{}{0pt}%
\pgfpathmoveto{\pgfqpoint{1.603485in}{2.483364in}}%
\pgfpathlineto{\pgfqpoint{1.575710in}{2.559224in}}%
\pgfpathlineto{\pgfqpoint{1.547962in}{2.635012in}}%
\pgfpathlineto{\pgfqpoint{1.597724in}{2.649823in}}%
\pgfpathlineto{\pgfqpoint{1.647395in}{2.664607in}}%
\pgfpathlineto{\pgfqpoint{1.675141in}{2.588971in}}%
\pgfpathlineto{\pgfqpoint{1.702913in}{2.513263in}}%
\pgfpathlineto{\pgfqpoint{1.653244in}{2.498327in}}%
\pgfpathlineto{\pgfqpoint{1.603485in}{2.483364in}}%
\pgfpathclose%
\pgfusepath{stroke,fill}%
\end{pgfscope}%
\begin{pgfscope}%
\pgfpathrectangle{\pgfqpoint{0.050000in}{0.083252in}}{\pgfqpoint{4.900000in}{4.900000in}}%
\pgfusepath{clip}%
\pgfsetbuttcap%
\pgfsetroundjoin%
\definecolor{currentfill}{rgb}{0.361538,0.335210,0.618239}%
\pgfsetfillcolor{currentfill}%
\pgfsetfillopacity{0.950000}%
\pgfsetlinewidth{0.200750pt}%
\definecolor{currentstroke}{rgb}{0.200000,0.200000,0.200000}%
\pgfsetstrokecolor{currentstroke}%
\pgfsetdash{}{0pt}%
\pgfpathmoveto{\pgfqpoint{1.814267in}{2.209708in}}%
\pgfpathlineto{\pgfqpoint{1.786389in}{2.285704in}}%
\pgfpathlineto{\pgfqpoint{1.758537in}{2.361629in}}%
\pgfpathlineto{\pgfqpoint{1.808113in}{2.376690in}}%
\pgfpathlineto{\pgfqpoint{1.857599in}{2.391724in}}%
\pgfpathlineto{\pgfqpoint{1.885448in}{2.315955in}}%
\pgfpathlineto{\pgfqpoint{1.913324in}{2.240118in}}%
\pgfpathlineto{\pgfqpoint{1.863840in}{2.224926in}}%
\pgfpathlineto{\pgfqpoint{1.814267in}{2.209708in}}%
\pgfpathclose%
\pgfusepath{stroke,fill}%
\end{pgfscope}%
\begin{pgfscope}%
\pgfpathrectangle{\pgfqpoint{0.050000in}{0.083252in}}{\pgfqpoint{4.900000in}{4.900000in}}%
\pgfusepath{clip}%
\pgfsetbuttcap%
\pgfsetroundjoin%
\definecolor{currentfill}{rgb}{0.522645,0.502960,0.714571}%
\pgfsetfillcolor{currentfill}%
\pgfsetfillopacity{0.950000}%
\pgfsetlinewidth{0.200750pt}%
\definecolor{currentstroke}{rgb}{0.200000,0.200000,0.200000}%
\pgfsetstrokecolor{currentstroke}%
\pgfsetdash{}{0pt}%
\pgfpathmoveto{\pgfqpoint{2.025090in}{1.936363in}}%
\pgfpathlineto{\pgfqpoint{1.997109in}{2.012312in}}%
\pgfpathlineto{\pgfqpoint{1.969154in}{2.088279in}}%
\pgfpathlineto{\pgfqpoint{2.018545in}{2.103627in}}%
\pgfpathlineto{\pgfqpoint{2.067845in}{2.118957in}}%
\pgfpathlineto{\pgfqpoint{2.095797in}{2.043209in}}%
\pgfpathlineto{\pgfqpoint{2.123773in}{1.967520in}}%
\pgfpathlineto{\pgfqpoint{2.074476in}{1.951946in}}%
\pgfpathlineto{\pgfqpoint{2.025090in}{1.936363in}}%
\pgfpathclose%
\pgfusepath{stroke,fill}%
\end{pgfscope}%
\begin{pgfscope}%
\pgfpathrectangle{\pgfqpoint{0.050000in}{0.083252in}}{\pgfqpoint{4.900000in}{4.900000in}}%
\pgfusepath{clip}%
\pgfsetbuttcap%
\pgfsetroundjoin%
\definecolor{currentfill}{rgb}{0.970165,0.968935,0.982161}%
\pgfsetfillcolor{currentfill}%
\pgfsetfillopacity{0.950000}%
\pgfsetlinewidth{0.200750pt}%
\definecolor{currentstroke}{rgb}{0.200000,0.200000,0.200000}%
\pgfsetstrokecolor{currentstroke}%
\pgfsetdash{}{0pt}%
\pgfpathmoveto{\pgfqpoint{3.162964in}{1.205661in}}%
\pgfpathlineto{\pgfqpoint{3.133913in}{1.243527in}}%
\pgfpathlineto{\pgfqpoint{3.104976in}{1.284276in}}%
\pgfpathlineto{\pgfqpoint{3.153320in}{1.301435in}}%
\pgfpathlineto{\pgfqpoint{3.201580in}{1.318547in}}%
\pgfpathlineto{\pgfqpoint{3.230582in}{1.277462in}}%
\pgfpathlineto{\pgfqpoint{3.259704in}{1.239212in}}%
\pgfpathlineto{\pgfqpoint{3.211377in}{1.222452in}}%
\pgfpathlineto{\pgfqpoint{3.162964in}{1.205661in}}%
\pgfpathclose%
\pgfusepath{stroke,fill}%
\end{pgfscope}%
\begin{pgfscope}%
\pgfpathrectangle{\pgfqpoint{0.050000in}{0.083252in}}{\pgfqpoint{4.900000in}{4.900000in}}%
\pgfusepath{clip}%
\pgfsetbuttcap%
\pgfsetroundjoin%
\definecolor{currentfill}{rgb}{0.677785,0.664498,0.807336}%
\pgfsetfillcolor{currentfill}%
\pgfsetfillopacity{0.950000}%
\pgfsetlinewidth{0.200750pt}%
\definecolor{currentstroke}{rgb}{0.200000,0.200000,0.200000}%
\pgfsetstrokecolor{currentstroke}%
\pgfsetdash{}{0pt}%
\pgfpathmoveto{\pgfqpoint{2.235924in}{1.667961in}}%
\pgfpathlineto{\pgfqpoint{2.207851in}{1.741959in}}%
\pgfpathlineto{\pgfqpoint{2.179801in}{1.816726in}}%
\pgfpathlineto{\pgfqpoint{2.229006in}{1.832658in}}%
\pgfpathlineto{\pgfqpoint{2.278122in}{1.848588in}}%
\pgfpathlineto{\pgfqpoint{2.306170in}{1.774310in}}%
\pgfpathlineto{\pgfqpoint{2.334242in}{1.700885in}}%
\pgfpathlineto{\pgfqpoint{2.285127in}{1.684427in}}%
\pgfpathlineto{\pgfqpoint{2.235924in}{1.667961in}}%
\pgfpathclose%
\pgfusepath{stroke,fill}%
\end{pgfscope}%
\begin{pgfscope}%
\pgfpathrectangle{\pgfqpoint{0.050000in}{0.083252in}}{\pgfqpoint{4.900000in}{4.900000in}}%
\pgfusepath{clip}%
\pgfsetbuttcap%
\pgfsetroundjoin%
\definecolor{currentfill}{rgb}{0.263699,0.467974,0.890919}%
\pgfsetfillcolor{currentfill}%
\pgfsetfillopacity{0.950000}%
\pgfsetlinewidth{0.200750pt}%
\definecolor{currentstroke}{rgb}{0.200000,0.200000,0.200000}%
\pgfsetstrokecolor{currentstroke}%
\pgfsetdash{}{0pt}%
\pgfpathmoveto{\pgfqpoint{0.815593in}{3.426610in}}%
\pgfpathlineto{\pgfqpoint{0.788182in}{3.502058in}}%
\pgfpathlineto{\pgfqpoint{0.760796in}{3.577434in}}%
\pgfpathlineto{\pgfqpoint{0.811298in}{3.591401in}}%
\pgfpathlineto{\pgfqpoint{0.861709in}{3.605341in}}%
\pgfpathlineto{\pgfqpoint{0.889093in}{3.530115in}}%
\pgfpathlineto{\pgfqpoint{0.916503in}{3.454817in}}%
\pgfpathlineto{\pgfqpoint{0.866094in}{3.440726in}}%
\pgfpathlineto{\pgfqpoint{0.815593in}{3.426610in}}%
\pgfpathclose%
\pgfusepath{stroke,fill}%
\end{pgfscope}%
\begin{pgfscope}%
\pgfpathrectangle{\pgfqpoint{0.050000in}{0.083252in}}{\pgfqpoint{4.900000in}{4.900000in}}%
\pgfusepath{clip}%
\pgfsetbuttcap%
\pgfsetroundjoin%
\definecolor{currentfill}{rgb}{0.994033,0.993787,0.996432}%
\pgfsetfillcolor{currentfill}%
\pgfsetfillopacity{0.950000}%
\pgfsetlinewidth{0.200750pt}%
\definecolor{currentstroke}{rgb}{0.200000,0.200000,0.200000}%
\pgfsetstrokecolor{currentstroke}%
\pgfsetdash{}{0pt}%
\pgfpathmoveto{\pgfqpoint{3.570446in}{1.171011in}}%
\pgfpathlineto{\pgfqpoint{3.540690in}{1.202392in}}%
\pgfpathlineto{\pgfqpoint{3.511052in}{1.235065in}}%
\pgfpathlineto{\pgfqpoint{3.559022in}{1.251317in}}%
\pgfpathlineto{\pgfqpoint{3.606905in}{1.267544in}}%
\pgfpathlineto{\pgfqpoint{3.636618in}{1.234522in}}%
\pgfpathlineto{\pgfqpoint{3.666452in}{1.202799in}}%
\pgfpathlineto{\pgfqpoint{3.618493in}{1.186917in}}%
\pgfpathlineto{\pgfqpoint{3.570446in}{1.171011in}}%
\pgfpathclose%
\pgfusepath{stroke,fill}%
\end{pgfscope}%
\begin{pgfscope}%
\pgfpathrectangle{\pgfqpoint{0.050000in}{0.083252in}}{\pgfqpoint{4.900000in}{4.900000in}}%
\pgfusepath{clip}%
\pgfsetbuttcap%
\pgfsetroundjoin%
\definecolor{currentfill}{rgb}{0.257301,0.400000,0.800554}%
\pgfsetfillcolor{currentfill}%
\pgfsetfillopacity{0.950000}%
\pgfsetlinewidth{0.200750pt}%
\definecolor{currentstroke}{rgb}{0.200000,0.200000,0.200000}%
\pgfsetstrokecolor{currentstroke}%
\pgfsetdash{}{0pt}%
\pgfpathmoveto{\pgfqpoint{1.026407in}{3.152901in}}%
\pgfpathlineto{\pgfqpoint{0.998891in}{3.228489in}}%
\pgfpathlineto{\pgfqpoint{0.971402in}{3.304004in}}%
\pgfpathlineto{\pgfqpoint{1.021718in}{3.318219in}}%
\pgfpathlineto{\pgfqpoint{1.071943in}{3.332409in}}%
\pgfpathlineto{\pgfqpoint{1.099431in}{3.257044in}}%
\pgfpathlineto{\pgfqpoint{1.126945in}{3.181607in}}%
\pgfpathlineto{\pgfqpoint{1.076721in}{3.167267in}}%
\pgfpathlineto{\pgfqpoint{1.026407in}{3.152901in}}%
\pgfpathclose%
\pgfusepath{stroke,fill}%
\end{pgfscope}%
\begin{pgfscope}%
\pgfpathrectangle{\pgfqpoint{0.050000in}{0.083252in}}{\pgfqpoint{4.900000in}{4.900000in}}%
\pgfusepath{clip}%
\pgfsetbuttcap%
\pgfsetroundjoin%
\definecolor{currentfill}{rgb}{0.250657,0.329412,0.706713}%
\pgfsetfillcolor{currentfill}%
\pgfsetfillopacity{0.950000}%
\pgfsetlinewidth{0.200750pt}%
\definecolor{currentstroke}{rgb}{0.200000,0.200000,0.200000}%
\pgfsetstrokecolor{currentstroke}%
\pgfsetdash{}{0pt}%
\pgfpathmoveto{\pgfqpoint{1.237265in}{2.879134in}}%
\pgfpathlineto{\pgfqpoint{1.209645in}{2.954861in}}%
\pgfpathlineto{\pgfqpoint{1.182052in}{3.030515in}}%
\pgfpathlineto{\pgfqpoint{1.232182in}{3.044980in}}%
\pgfpathlineto{\pgfqpoint{1.282222in}{3.059418in}}%
\pgfpathlineto{\pgfqpoint{1.309814in}{2.983914in}}%
\pgfpathlineto{\pgfqpoint{1.337431in}{2.908338in}}%
\pgfpathlineto{\pgfqpoint{1.287393in}{2.893749in}}%
\pgfpathlineto{\pgfqpoint{1.237265in}{2.879134in}}%
\pgfpathclose%
\pgfusepath{stroke,fill}%
\end{pgfscope}%
\begin{pgfscope}%
\pgfpathrectangle{\pgfqpoint{0.050000in}{0.083252in}}{\pgfqpoint{4.900000in}{4.900000in}}%
\pgfusepath{clip}%
\pgfsetbuttcap%
\pgfsetroundjoin%
\definecolor{currentfill}{rgb}{0.874694,0.869527,0.925075}%
\pgfsetfillcolor{currentfill}%
\pgfsetfillopacity{0.950000}%
\pgfsetlinewidth{0.200750pt}%
\definecolor{currentstroke}{rgb}{0.200000,0.200000,0.200000}%
\pgfsetstrokecolor{currentstroke}%
\pgfsetdash{}{0pt}%
\pgfpathmoveto{\pgfqpoint{2.601335in}{1.339117in}}%
\pgfpathlineto{\pgfqpoint{2.573045in}{1.397500in}}%
\pgfpathlineto{\pgfqpoint{2.544804in}{1.459446in}}%
\pgfpathlineto{\pgfqpoint{2.593686in}{1.477022in}}%
\pgfpathlineto{\pgfqpoint{2.642483in}{1.494534in}}%
\pgfpathlineto{\pgfqpoint{2.670754in}{1.433001in}}%
\pgfpathlineto{\pgfqpoint{2.699081in}{1.374884in}}%
\pgfpathlineto{\pgfqpoint{2.650250in}{1.357041in}}%
\pgfpathlineto{\pgfqpoint{2.601335in}{1.339117in}}%
\pgfpathclose%
\pgfusepath{stroke,fill}%
\end{pgfscope}%
\begin{pgfscope}%
\pgfpathrectangle{\pgfqpoint{0.050000in}{0.083252in}}{\pgfqpoint{4.900000in}{4.900000in}}%
\pgfusepath{clip}%
\pgfsetbuttcap%
\pgfsetroundjoin%
\definecolor{currentfill}{rgb}{0.244014,0.258824,0.612872}%
\pgfsetfillcolor{currentfill}%
\pgfsetfillopacity{0.950000}%
\pgfsetlinewidth{0.200750pt}%
\definecolor{currentstroke}{rgb}{0.200000,0.200000,0.200000}%
\pgfsetstrokecolor{currentstroke}%
\pgfsetdash{}{0pt}%
\pgfpathmoveto{\pgfqpoint{1.448167in}{2.605308in}}%
\pgfpathlineto{\pgfqpoint{1.420443in}{2.681175in}}%
\pgfpathlineto{\pgfqpoint{1.392746in}{2.756969in}}%
\pgfpathlineto{\pgfqpoint{1.442691in}{2.771682in}}%
\pgfpathlineto{\pgfqpoint{1.492545in}{2.786369in}}%
\pgfpathlineto{\pgfqpoint{1.520240in}{2.710727in}}%
\pgfpathlineto{\pgfqpoint{1.547962in}{2.635012in}}%
\pgfpathlineto{\pgfqpoint{1.498110in}{2.620174in}}%
\pgfpathlineto{\pgfqpoint{1.448167in}{2.605308in}}%
\pgfpathclose%
\pgfusepath{stroke,fill}%
\end{pgfscope}%
\begin{pgfscope}%
\pgfpathrectangle{\pgfqpoint{0.050000in}{0.083252in}}{\pgfqpoint{4.900000in}{4.900000in}}%
\pgfusepath{clip}%
\pgfsetbuttcap%
\pgfsetroundjoin%
\definecolor{currentfill}{rgb}{0.283968,0.254441,0.571857}%
\pgfsetfillcolor{currentfill}%
\pgfsetfillopacity{0.950000}%
\pgfsetlinewidth{0.200750pt}%
\definecolor{currentstroke}{rgb}{0.200000,0.200000,0.200000}%
\pgfsetstrokecolor{currentstroke}%
\pgfsetdash{}{0pt}%
\pgfpathmoveto{\pgfqpoint{1.659114in}{2.331425in}}%
\pgfpathlineto{\pgfqpoint{1.631286in}{2.407431in}}%
\pgfpathlineto{\pgfqpoint{1.603485in}{2.483364in}}%
\pgfpathlineto{\pgfqpoint{1.653244in}{2.498327in}}%
\pgfpathlineto{\pgfqpoint{1.702913in}{2.513263in}}%
\pgfpathlineto{\pgfqpoint{1.730711in}{2.437482in}}%
\pgfpathlineto{\pgfqpoint{1.758537in}{2.361629in}}%
\pgfpathlineto{\pgfqpoint{1.708870in}{2.346540in}}%
\pgfpathlineto{\pgfqpoint{1.659114in}{2.331425in}}%
\pgfpathclose%
\pgfusepath{stroke,fill}%
\end{pgfscope}%
\begin{pgfscope}%
\pgfpathrectangle{\pgfqpoint{0.050000in}{0.083252in}}{\pgfqpoint{4.900000in}{4.900000in}}%
\pgfusepath{clip}%
\pgfsetbuttcap%
\pgfsetroundjoin%
\definecolor{currentfill}{rgb}{0.916463,0.913018,0.950050}%
\pgfsetfillcolor{currentfill}%
\pgfsetfillopacity{0.950000}%
\pgfsetlinewidth{0.200750pt}%
\definecolor{currentstroke}{rgb}{0.200000,0.200000,0.200000}%
\pgfsetstrokecolor{currentstroke}%
\pgfsetdash{}{0pt}%
\pgfpathmoveto{\pgfqpoint{2.755935in}{1.269927in}}%
\pgfpathlineto{\pgfqpoint{2.727472in}{1.320472in}}%
\pgfpathlineto{\pgfqpoint{2.699081in}{1.374884in}}%
\pgfpathlineto{\pgfqpoint{2.747827in}{1.392649in}}%
\pgfpathlineto{\pgfqpoint{2.796488in}{1.410338in}}%
\pgfpathlineto{\pgfqpoint{2.824925in}{1.356011in}}%
\pgfpathlineto{\pgfqpoint{2.853441in}{1.305384in}}%
\pgfpathlineto{\pgfqpoint{2.804731in}{1.287691in}}%
\pgfpathlineto{\pgfqpoint{2.755935in}{1.269927in}}%
\pgfpathclose%
\pgfusepath{stroke,fill}%
\end{pgfscope}%
\begin{pgfscope}%
\pgfpathrectangle{\pgfqpoint{0.050000in}{0.083252in}}{\pgfqpoint{4.900000in}{4.900000in}}%
\pgfusepath{clip}%
\pgfsetbuttcap%
\pgfsetroundjoin%
\definecolor{currentfill}{rgb}{0.820992,0.813610,0.892964}%
\pgfsetfillcolor{currentfill}%
\pgfsetfillopacity{0.950000}%
\pgfsetlinewidth{0.200750pt}%
\definecolor{currentstroke}{rgb}{0.200000,0.200000,0.200000}%
\pgfsetstrokecolor{currentstroke}%
\pgfsetdash{}{0pt}%
\pgfpathmoveto{\pgfqpoint{2.446786in}{1.424095in}}%
\pgfpathlineto{\pgfqpoint{2.418608in}{1.489757in}}%
\pgfpathlineto{\pgfqpoint{2.390460in}{1.558133in}}%
\pgfpathlineto{\pgfqpoint{2.439494in}{1.575219in}}%
\pgfpathlineto{\pgfqpoint{2.488442in}{1.592270in}}%
\pgfpathlineto{\pgfqpoint{2.516605in}{1.524530in}}%
\pgfpathlineto{\pgfqpoint{2.544804in}{1.459446in}}%
\pgfpathlineto{\pgfqpoint{2.495838in}{1.441804in}}%
\pgfpathlineto{\pgfqpoint{2.446786in}{1.424095in}}%
\pgfpathclose%
\pgfusepath{stroke,fill}%
\end{pgfscope}%
\begin{pgfscope}%
\pgfpathrectangle{\pgfqpoint{0.050000in}{0.083252in}}{\pgfqpoint{4.900000in}{4.900000in}}%
\pgfusepath{clip}%
\pgfsetbuttcap%
\pgfsetroundjoin%
\definecolor{currentfill}{rgb}{0.439108,0.415978,0.664621}%
\pgfsetfillcolor{currentfill}%
\pgfsetfillopacity{0.950000}%
\pgfsetlinewidth{0.200750pt}%
\definecolor{currentstroke}{rgb}{0.200000,0.200000,0.200000}%
\pgfsetstrokecolor{currentstroke}%
\pgfsetdash{}{0pt}%
\pgfpathmoveto{\pgfqpoint{1.870105in}{2.057523in}}%
\pgfpathlineto{\pgfqpoint{1.842173in}{2.133646in}}%
\pgfpathlineto{\pgfqpoint{1.814267in}{2.209708in}}%
\pgfpathlineto{\pgfqpoint{1.863840in}{2.224926in}}%
\pgfpathlineto{\pgfqpoint{1.913324in}{2.240118in}}%
\pgfpathlineto{\pgfqpoint{1.941226in}{2.164222in}}%
\pgfpathlineto{\pgfqpoint{1.969154in}{2.088279in}}%
\pgfpathlineto{\pgfqpoint{1.919674in}{2.072911in}}%
\pgfpathlineto{\pgfqpoint{1.870105in}{2.057523in}}%
\pgfpathclose%
\pgfusepath{stroke,fill}%
\end{pgfscope}%
\begin{pgfscope}%
\pgfpathrectangle{\pgfqpoint{0.050000in}{0.083252in}}{\pgfqpoint{4.900000in}{4.900000in}}%
\pgfusepath{clip}%
\pgfsetbuttcap%
\pgfsetroundjoin%
\definecolor{currentfill}{rgb}{0.600215,0.583729,0.760953}%
\pgfsetfillcolor{currentfill}%
\pgfsetfillopacity{0.950000}%
\pgfsetlinewidth{0.200750pt}%
\definecolor{currentstroke}{rgb}{0.200000,0.200000,0.200000}%
\pgfsetstrokecolor{currentstroke}%
\pgfsetdash{}{0pt}%
\pgfpathmoveto{\pgfqpoint{2.081126in}{1.784860in}}%
\pgfpathlineto{\pgfqpoint{2.053096in}{1.860507in}}%
\pgfpathlineto{\pgfqpoint{2.025090in}{1.936363in}}%
\pgfpathlineto{\pgfqpoint{2.074476in}{1.951946in}}%
\pgfpathlineto{\pgfqpoint{2.123773in}{1.967520in}}%
\pgfpathlineto{\pgfqpoint{2.151775in}{1.891980in}}%
\pgfpathlineto{\pgfqpoint{2.179801in}{1.816726in}}%
\pgfpathlineto{\pgfqpoint{2.130508in}{1.800793in}}%
\pgfpathlineto{\pgfqpoint{2.081126in}{1.784860in}}%
\pgfpathclose%
\pgfusepath{stroke,fill}%
\end{pgfscope}%
\begin{pgfscope}%
\pgfpathrectangle{\pgfqpoint{0.050000in}{0.083252in}}{\pgfqpoint{4.900000in}{4.900000in}}%
\pgfusepath{clip}%
\pgfsetbuttcap%
\pgfsetroundjoin%
\definecolor{currentfill}{rgb}{0.982099,0.981361,0.989296}%
\pgfsetfillcolor{currentfill}%
\pgfsetfillopacity{0.950000}%
\pgfsetlinewidth{0.200750pt}%
\definecolor{currentstroke}{rgb}{0.200000,0.200000,0.200000}%
\pgfsetstrokecolor{currentstroke}%
\pgfsetdash{}{0pt}%
\pgfpathmoveto{\pgfqpoint{3.318308in}{1.169817in}}%
\pgfpathlineto{\pgfqpoint{3.288946in}{1.203460in}}%
\pgfpathlineto{\pgfqpoint{3.259704in}{1.239212in}}%
\pgfpathlineto{\pgfqpoint{3.307945in}{1.255942in}}%
\pgfpathlineto{\pgfqpoint{3.356100in}{1.272639in}}%
\pgfpathlineto{\pgfqpoint{3.385415in}{1.236509in}}%
\pgfpathlineto{\pgfqpoint{3.414854in}{1.202488in}}%
\pgfpathlineto{\pgfqpoint{3.366625in}{1.186163in}}%
\pgfpathlineto{\pgfqpoint{3.318308in}{1.169817in}}%
\pgfpathclose%
\pgfusepath{stroke,fill}%
\end{pgfscope}%
\begin{pgfscope}%
\pgfpathrectangle{\pgfqpoint{0.050000in}{0.083252in}}{\pgfqpoint{4.900000in}{4.900000in}}%
\pgfusepath{clip}%
\pgfsetbuttcap%
\pgfsetroundjoin%
\definecolor{currentfill}{rgb}{0.952265,0.950296,0.971457}%
\pgfsetfillcolor{currentfill}%
\pgfsetfillopacity{0.950000}%
\pgfsetlinewidth{0.200750pt}%
\definecolor{currentstroke}{rgb}{0.200000,0.200000,0.200000}%
\pgfsetstrokecolor{currentstroke}%
\pgfsetdash{}{0pt}%
\pgfpathmoveto{\pgfqpoint{2.910743in}{1.215167in}}%
\pgfpathlineto{\pgfqpoint{2.882045in}{1.258477in}}%
\pgfpathlineto{\pgfqpoint{2.853441in}{1.305384in}}%
\pgfpathlineto{\pgfqpoint{2.902066in}{1.323009in}}%
\pgfpathlineto{\pgfqpoint{2.950606in}{1.340567in}}%
\pgfpathlineto{\pgfqpoint{2.979268in}{1.293454in}}%
\pgfpathlineto{\pgfqpoint{3.008031in}{1.249818in}}%
\pgfpathlineto{\pgfqpoint{2.959430in}{1.232517in}}%
\pgfpathlineto{\pgfqpoint{2.910743in}{1.215167in}}%
\pgfpathclose%
\pgfusepath{stroke,fill}%
\end{pgfscope}%
\begin{pgfscope}%
\pgfpathrectangle{\pgfqpoint{0.050000in}{0.083252in}}{\pgfqpoint{4.900000in}{4.900000in}}%
\pgfusepath{clip}%
\pgfsetbuttcap%
\pgfsetroundjoin%
\definecolor{currentfill}{rgb}{0.260500,0.433987,0.845736}%
\pgfsetfillcolor{currentfill}%
\pgfsetfillopacity{0.950000}%
\pgfsetlinewidth{0.200750pt}%
\definecolor{currentstroke}{rgb}{0.200000,0.200000,0.200000}%
\pgfsetstrokecolor{currentstroke}%
\pgfsetdash{}{0pt}%
\pgfpathmoveto{\pgfqpoint{0.870495in}{3.275495in}}%
\pgfpathlineto{\pgfqpoint{0.843031in}{3.351089in}}%
\pgfpathlineto{\pgfqpoint{0.815593in}{3.426610in}}%
\pgfpathlineto{\pgfqpoint{0.866094in}{3.440726in}}%
\pgfpathlineto{\pgfqpoint{0.916503in}{3.454817in}}%
\pgfpathlineto{\pgfqpoint{0.943939in}{3.379447in}}%
\pgfpathlineto{\pgfqpoint{0.971402in}{3.304004in}}%
\pgfpathlineto{\pgfqpoint{0.920994in}{3.289763in}}%
\pgfpathlineto{\pgfqpoint{0.870495in}{3.275495in}}%
\pgfpathclose%
\pgfusepath{stroke,fill}%
\end{pgfscope}%
\begin{pgfscope}%
\pgfpathrectangle{\pgfqpoint{0.050000in}{0.083252in}}{\pgfqpoint{4.900000in}{4.900000in}}%
\pgfusepath{clip}%
\pgfsetbuttcap%
\pgfsetroundjoin%
\definecolor{currentfill}{rgb}{0.755356,0.745267,0.853718}%
\pgfsetfillcolor{currentfill}%
\pgfsetfillopacity{0.950000}%
\pgfsetlinewidth{0.200750pt}%
\definecolor{currentstroke}{rgb}{0.200000,0.200000,0.200000}%
\pgfsetstrokecolor{currentstroke}%
\pgfsetdash{}{0pt}%
\pgfpathmoveto{\pgfqpoint{2.292133in}{1.523860in}}%
\pgfpathlineto{\pgfqpoint{2.264018in}{1.595105in}}%
\pgfpathlineto{\pgfqpoint{2.235924in}{1.667961in}}%
\pgfpathlineto{\pgfqpoint{2.285127in}{1.684427in}}%
\pgfpathlineto{\pgfqpoint{2.334242in}{1.700885in}}%
\pgfpathlineto{\pgfqpoint{2.362338in}{1.628678in}}%
\pgfpathlineto{\pgfqpoint{2.390460in}{1.558133in}}%
\pgfpathlineto{\pgfqpoint{2.341340in}{1.541012in}}%
\pgfpathlineto{\pgfqpoint{2.292133in}{1.523860in}}%
\pgfpathclose%
\pgfusepath{stroke,fill}%
\end{pgfscope}%
\begin{pgfscope}%
\pgfpathrectangle{\pgfqpoint{0.050000in}{0.083252in}}{\pgfqpoint{4.900000in}{4.900000in}}%
\pgfusepath{clip}%
\pgfsetbuttcap%
\pgfsetroundjoin%
\definecolor{currentfill}{rgb}{0.253856,0.363399,0.751895}%
\pgfsetfillcolor{currentfill}%
\pgfsetfillopacity{0.950000}%
\pgfsetlinewidth{0.200750pt}%
\definecolor{currentstroke}{rgb}{0.200000,0.200000,0.200000}%
\pgfsetstrokecolor{currentstroke}%
\pgfsetdash{}{0pt}%
\pgfpathmoveto{\pgfqpoint{1.081517in}{3.001507in}}%
\pgfpathlineto{\pgfqpoint{1.053949in}{3.077240in}}%
\pgfpathlineto{\pgfqpoint{1.026407in}{3.152901in}}%
\pgfpathlineto{\pgfqpoint{1.076721in}{3.167267in}}%
\pgfpathlineto{\pgfqpoint{1.126945in}{3.181607in}}%
\pgfpathlineto{\pgfqpoint{1.154485in}{3.106097in}}%
\pgfpathlineto{\pgfqpoint{1.182052in}{3.030515in}}%
\pgfpathlineto{\pgfqpoint{1.131830in}{3.016024in}}%
\pgfpathlineto{\pgfqpoint{1.081517in}{3.001507in}}%
\pgfpathclose%
\pgfusepath{stroke,fill}%
\end{pgfscope}%
\begin{pgfscope}%
\pgfpathrectangle{\pgfqpoint{0.050000in}{0.083252in}}{\pgfqpoint{4.900000in}{4.900000in}}%
\pgfusepath{clip}%
\pgfsetbuttcap%
\pgfsetroundjoin%
\definecolor{currentfill}{rgb}{0.247213,0.292810,0.658055}%
\pgfsetfillcolor{currentfill}%
\pgfsetfillopacity{0.950000}%
\pgfsetlinewidth{0.200750pt}%
\definecolor{currentstroke}{rgb}{0.200000,0.200000,0.200000}%
\pgfsetstrokecolor{currentstroke}%
\pgfsetdash{}{0pt}%
\pgfpathmoveto{\pgfqpoint{1.292584in}{2.727461in}}%
\pgfpathlineto{\pgfqpoint{1.264911in}{2.803334in}}%
\pgfpathlineto{\pgfqpoint{1.237265in}{2.879134in}}%
\pgfpathlineto{\pgfqpoint{1.287393in}{2.893749in}}%
\pgfpathlineto{\pgfqpoint{1.337431in}{2.908338in}}%
\pgfpathlineto{\pgfqpoint{1.365075in}{2.832690in}}%
\pgfpathlineto{\pgfqpoint{1.392746in}{2.756969in}}%
\pgfpathlineto{\pgfqpoint{1.342710in}{2.742228in}}%
\pgfpathlineto{\pgfqpoint{1.292584in}{2.727461in}}%
\pgfpathclose%
\pgfusepath{stroke,fill}%
\end{pgfscope}%
\begin{pgfscope}%
\pgfpathrectangle{\pgfqpoint{0.050000in}{0.083252in}}{\pgfqpoint{4.900000in}{4.900000in}}%
\pgfusepath{clip}%
\pgfsetbuttcap%
\pgfsetroundjoin%
\definecolor{currentfill}{rgb}{1.000000,1.000000,1.000000}%
\pgfsetfillcolor{currentfill}%
\pgfsetfillopacity{0.950000}%
\pgfsetlinewidth{0.200750pt}%
\definecolor{currentstroke}{rgb}{0.200000,0.200000,0.200000}%
\pgfsetstrokecolor{currentstroke}%
\pgfsetdash{}{0pt}%
\pgfpathmoveto{\pgfqpoint{3.726448in}{1.141427in}}%
\pgfpathlineto{\pgfqpoint{3.696399in}{1.171929in}}%
\pgfpathlineto{\pgfqpoint{3.666452in}{1.202799in}}%
\pgfpathlineto{\pgfqpoint{3.714324in}{1.218656in}}%
\pgfpathlineto{\pgfqpoint{3.762108in}{1.234488in}}%
\pgfpathlineto{\pgfqpoint{3.792133in}{1.203279in}}%
\pgfpathlineto{\pgfqpoint{3.822260in}{1.172434in}}%
\pgfpathlineto{\pgfqpoint{3.774398in}{1.156945in}}%
\pgfpathlineto{\pgfqpoint{3.726448in}{1.141427in}}%
\pgfpathclose%
\pgfusepath{stroke,fill}%
\end{pgfscope}%
\begin{pgfscope}%
\pgfpathrectangle{\pgfqpoint{0.050000in}{0.083252in}}{\pgfqpoint{4.900000in}{4.900000in}}%
\pgfusepath{clip}%
\pgfsetbuttcap%
\pgfsetroundjoin%
\definecolor{currentfill}{rgb}{0.240815,0.224837,0.567689}%
\pgfsetfillcolor{currentfill}%
\pgfsetfillopacity{0.950000}%
\pgfsetlinewidth{0.200750pt}%
\definecolor{currentstroke}{rgb}{0.200000,0.200000,0.200000}%
\pgfsetstrokecolor{currentstroke}%
\pgfsetdash{}{0pt}%
\pgfpathmoveto{\pgfqpoint{1.503695in}{2.453356in}}%
\pgfpathlineto{\pgfqpoint{1.475918in}{2.529369in}}%
\pgfpathlineto{\pgfqpoint{1.448167in}{2.605308in}}%
\pgfpathlineto{\pgfqpoint{1.498110in}{2.620174in}}%
\pgfpathlineto{\pgfqpoint{1.547962in}{2.635012in}}%
\pgfpathlineto{\pgfqpoint{1.575710in}{2.559224in}}%
\pgfpathlineto{\pgfqpoint{1.603485in}{2.483364in}}%
\pgfpathlineto{\pgfqpoint{1.553635in}{2.468374in}}%
\pgfpathlineto{\pgfqpoint{1.503695in}{2.453356in}}%
\pgfpathclose%
\pgfusepath{stroke,fill}%
\end{pgfscope}%
\begin{pgfscope}%
\pgfpathrectangle{\pgfqpoint{0.050000in}{0.083252in}}{\pgfqpoint{4.900000in}{4.900000in}}%
\pgfusepath{clip}%
\pgfsetbuttcap%
\pgfsetroundjoin%
\definecolor{currentfill}{rgb}{0.361538,0.335210,0.618239}%
\pgfsetfillcolor{currentfill}%
\pgfsetfillopacity{0.950000}%
\pgfsetlinewidth{0.200750pt}%
\definecolor{currentstroke}{rgb}{0.200000,0.200000,0.200000}%
\pgfsetstrokecolor{currentstroke}%
\pgfsetdash{}{0pt}%
\pgfpathmoveto{\pgfqpoint{1.714850in}{2.179195in}}%
\pgfpathlineto{\pgfqpoint{1.686969in}{2.255346in}}%
\pgfpathlineto{\pgfqpoint{1.659114in}{2.331425in}}%
\pgfpathlineto{\pgfqpoint{1.708870in}{2.346540in}}%
\pgfpathlineto{\pgfqpoint{1.758537in}{2.361629in}}%
\pgfpathlineto{\pgfqpoint{1.786389in}{2.285704in}}%
\pgfpathlineto{\pgfqpoint{1.814267in}{2.209708in}}%
\pgfpathlineto{\pgfqpoint{1.764604in}{2.194465in}}%
\pgfpathlineto{\pgfqpoint{1.714850in}{2.179195in}}%
\pgfpathclose%
\pgfusepath{stroke,fill}%
\end{pgfscope}%
\begin{pgfscope}%
\pgfpathrectangle{\pgfqpoint{0.050000in}{0.083252in}}{\pgfqpoint{4.900000in}{4.900000in}}%
\pgfusepath{clip}%
\pgfsetbuttcap%
\pgfsetroundjoin%
\definecolor{currentfill}{rgb}{0.522645,0.502960,0.714571}%
\pgfsetfillcolor{currentfill}%
\pgfsetfillopacity{0.950000}%
\pgfsetlinewidth{0.200750pt}%
\definecolor{currentstroke}{rgb}{0.200000,0.200000,0.200000}%
\pgfsetstrokecolor{currentstroke}%
\pgfsetdash{}{0pt}%
\pgfpathmoveto{\pgfqpoint{1.926048in}{1.905170in}}%
\pgfpathlineto{\pgfqpoint{1.898063in}{1.981355in}}%
\pgfpathlineto{\pgfqpoint{1.870105in}{2.057523in}}%
\pgfpathlineto{\pgfqpoint{1.919674in}{2.072911in}}%
\pgfpathlineto{\pgfqpoint{1.969154in}{2.088279in}}%
\pgfpathlineto{\pgfqpoint{1.997109in}{2.012312in}}%
\pgfpathlineto{\pgfqpoint{2.025090in}{1.936363in}}%
\pgfpathlineto{\pgfqpoint{1.975614in}{1.920772in}}%
\pgfpathlineto{\pgfqpoint{1.926048in}{1.905170in}}%
\pgfpathclose%
\pgfusepath{stroke,fill}%
\end{pgfscope}%
\begin{pgfscope}%
\pgfpathrectangle{\pgfqpoint{0.050000in}{0.083252in}}{\pgfqpoint{4.900000in}{4.900000in}}%
\pgfusepath{clip}%
\pgfsetbuttcap%
\pgfsetroundjoin%
\definecolor{currentfill}{rgb}{0.970165,0.968935,0.982161}%
\pgfsetfillcolor{currentfill}%
\pgfsetfillopacity{0.950000}%
\pgfsetlinewidth{0.200750pt}%
\definecolor{currentstroke}{rgb}{0.200000,0.200000,0.200000}%
\pgfsetstrokecolor{currentstroke}%
\pgfsetdash{}{0pt}%
\pgfpathmoveto{\pgfqpoint{3.065877in}{1.171990in}}%
\pgfpathlineto{\pgfqpoint{3.036899in}{1.209438in}}%
\pgfpathlineto{\pgfqpoint{3.008031in}{1.249818in}}%
\pgfpathlineto{\pgfqpoint{3.056546in}{1.267071in}}%
\pgfpathlineto{\pgfqpoint{3.104976in}{1.284276in}}%
\pgfpathlineto{\pgfqpoint{3.133913in}{1.243527in}}%
\pgfpathlineto{\pgfqpoint{3.162964in}{1.205661in}}%
\pgfpathlineto{\pgfqpoint{3.114464in}{1.188840in}}%
\pgfpathlineto{\pgfqpoint{3.065877in}{1.171990in}}%
\pgfpathclose%
\pgfusepath{stroke,fill}%
\end{pgfscope}%
\begin{pgfscope}%
\pgfpathrectangle{\pgfqpoint{0.050000in}{0.083252in}}{\pgfqpoint{4.900000in}{4.900000in}}%
\pgfusepath{clip}%
\pgfsetbuttcap%
\pgfsetroundjoin%
\definecolor{currentfill}{rgb}{0.677785,0.664498,0.807336}%
\pgfsetfillcolor{currentfill}%
\pgfsetfillopacity{0.950000}%
\pgfsetlinewidth{0.200750pt}%
\definecolor{currentstroke}{rgb}{0.200000,0.200000,0.200000}%
\pgfsetstrokecolor{currentstroke}%
\pgfsetdash{}{0pt}%
\pgfpathmoveto{\pgfqpoint{2.137254in}{1.635021in}}%
\pgfpathlineto{\pgfqpoint{2.109179in}{1.709608in}}%
\pgfpathlineto{\pgfqpoint{2.081126in}{1.784860in}}%
\pgfpathlineto{\pgfqpoint{2.130508in}{1.800793in}}%
\pgfpathlineto{\pgfqpoint{2.179801in}{1.816726in}}%
\pgfpathlineto{\pgfqpoint{2.207851in}{1.741959in}}%
\pgfpathlineto{\pgfqpoint{2.235924in}{1.667961in}}%
\pgfpathlineto{\pgfqpoint{2.186633in}{1.651491in}}%
\pgfpathlineto{\pgfqpoint{2.137254in}{1.635021in}}%
\pgfpathclose%
\pgfusepath{stroke,fill}%
\end{pgfscope}%
\begin{pgfscope}%
\pgfpathrectangle{\pgfqpoint{0.050000in}{0.083252in}}{\pgfqpoint{4.900000in}{4.900000in}}%
\pgfusepath{clip}%
\pgfsetbuttcap%
\pgfsetroundjoin%
\definecolor{currentfill}{rgb}{0.994033,0.993787,0.996432}%
\pgfsetfillcolor{currentfill}%
\pgfsetfillopacity{0.950000}%
\pgfsetlinewidth{0.200750pt}%
\definecolor{currentstroke}{rgb}{0.200000,0.200000,0.200000}%
\pgfsetstrokecolor{currentstroke}%
\pgfsetdash{}{0pt}%
\pgfpathmoveto{\pgfqpoint{3.474088in}{1.139128in}}%
\pgfpathlineto{\pgfqpoint{3.444413in}{1.170170in}}%
\pgfpathlineto{\pgfqpoint{3.414854in}{1.202488in}}%
\pgfpathlineto{\pgfqpoint{3.462997in}{1.218789in}}%
\pgfpathlineto{\pgfqpoint{3.511052in}{1.235065in}}%
\pgfpathlineto{\pgfqpoint{3.540690in}{1.202392in}}%
\pgfpathlineto{\pgfqpoint{3.570446in}{1.171011in}}%
\pgfpathlineto{\pgfqpoint{3.522311in}{1.155081in}}%
\pgfpathlineto{\pgfqpoint{3.474088in}{1.139128in}}%
\pgfpathclose%
\pgfusepath{stroke,fill}%
\end{pgfscope}%
\begin{pgfscope}%
\pgfpathrectangle{\pgfqpoint{0.050000in}{0.083252in}}{\pgfqpoint{4.900000in}{4.900000in}}%
\pgfusepath{clip}%
\pgfsetbuttcap%
\pgfsetroundjoin%
\definecolor{currentfill}{rgb}{0.257055,0.397386,0.797078}%
\pgfsetfillcolor{currentfill}%
\pgfsetfillopacity{0.950000}%
\pgfsetlinewidth{0.200750pt}%
\definecolor{currentstroke}{rgb}{0.200000,0.200000,0.200000}%
\pgfsetstrokecolor{currentstroke}%
\pgfsetdash{}{0pt}%
\pgfpathmoveto{\pgfqpoint{0.925503in}{3.124090in}}%
\pgfpathlineto{\pgfqpoint{0.897986in}{3.199829in}}%
\pgfpathlineto{\pgfqpoint{0.870495in}{3.275495in}}%
\pgfpathlineto{\pgfqpoint{0.920994in}{3.289763in}}%
\pgfpathlineto{\pgfqpoint{0.971402in}{3.304004in}}%
\pgfpathlineto{\pgfqpoint{0.998891in}{3.228489in}}%
\pgfpathlineto{\pgfqpoint{1.026407in}{3.152901in}}%
\pgfpathlineto{\pgfqpoint{0.976001in}{3.138508in}}%
\pgfpathlineto{\pgfqpoint{0.925503in}{3.124090in}}%
\pgfpathclose%
\pgfusepath{stroke,fill}%
\end{pgfscope}%
\begin{pgfscope}%
\pgfpathrectangle{\pgfqpoint{0.050000in}{0.083252in}}{\pgfqpoint{4.900000in}{4.900000in}}%
\pgfusepath{clip}%
\pgfsetbuttcap%
\pgfsetroundjoin%
\definecolor{currentfill}{rgb}{0.250657,0.329412,0.706713}%
\pgfsetfillcolor{currentfill}%
\pgfsetfillopacity{0.950000}%
\pgfsetlinewidth{0.200750pt}%
\definecolor{currentstroke}{rgb}{0.200000,0.200000,0.200000}%
\pgfsetstrokecolor{currentstroke}%
\pgfsetdash{}{0pt}%
\pgfpathmoveto{\pgfqpoint{1.136734in}{2.849822in}}%
\pgfpathlineto{\pgfqpoint{1.109112in}{2.925701in}}%
\pgfpathlineto{\pgfqpoint{1.081517in}{3.001507in}}%
\pgfpathlineto{\pgfqpoint{1.131830in}{3.016024in}}%
\pgfpathlineto{\pgfqpoint{1.182052in}{3.030515in}}%
\pgfpathlineto{\pgfqpoint{1.209645in}{2.954861in}}%
\pgfpathlineto{\pgfqpoint{1.237265in}{2.879134in}}%
\pgfpathlineto{\pgfqpoint{1.187045in}{2.864491in}}%
\pgfpathlineto{\pgfqpoint{1.136734in}{2.849822in}}%
\pgfpathclose%
\pgfusepath{stroke,fill}%
\end{pgfscope}%
\begin{pgfscope}%
\pgfpathrectangle{\pgfqpoint{0.050000in}{0.083252in}}{\pgfqpoint{4.900000in}{4.900000in}}%
\pgfusepath{clip}%
\pgfsetbuttcap%
\pgfsetroundjoin%
\definecolor{currentfill}{rgb}{0.880661,0.875740,0.928643}%
\pgfsetfillcolor{currentfill}%
\pgfsetfillopacity{0.950000}%
\pgfsetlinewidth{0.200750pt}%
\definecolor{currentstroke}{rgb}{0.200000,0.200000,0.200000}%
\pgfsetstrokecolor{currentstroke}%
\pgfsetdash{}{0pt}%
\pgfpathmoveto{\pgfqpoint{2.503250in}{1.303010in}}%
\pgfpathlineto{\pgfqpoint{2.474998in}{1.361687in}}%
\pgfpathlineto{\pgfqpoint{2.446786in}{1.424095in}}%
\pgfpathlineto{\pgfqpoint{2.495838in}{1.441804in}}%
\pgfpathlineto{\pgfqpoint{2.544804in}{1.459446in}}%
\pgfpathlineto{\pgfqpoint{2.573045in}{1.397500in}}%
\pgfpathlineto{\pgfqpoint{2.601335in}{1.339117in}}%
\pgfpathlineto{\pgfqpoint{2.552335in}{1.321108in}}%
\pgfpathlineto{\pgfqpoint{2.503250in}{1.303010in}}%
\pgfpathclose%
\pgfusepath{stroke,fill}%
\end{pgfscope}%
\begin{pgfscope}%
\pgfpathrectangle{\pgfqpoint{0.050000in}{0.083252in}}{\pgfqpoint{4.900000in}{4.900000in}}%
\pgfusepath{clip}%
\pgfsetbuttcap%
\pgfsetroundjoin%
\definecolor{currentfill}{rgb}{0.244014,0.258824,0.612872}%
\pgfsetfillcolor{currentfill}%
\pgfsetfillopacity{0.950000}%
\pgfsetlinewidth{0.200750pt}%
\definecolor{currentstroke}{rgb}{0.200000,0.200000,0.200000}%
\pgfsetstrokecolor{currentstroke}%
\pgfsetdash{}{0pt}%
\pgfpathmoveto{\pgfqpoint{1.348009in}{2.575496in}}%
\pgfpathlineto{\pgfqpoint{1.320283in}{2.651515in}}%
\pgfpathlineto{\pgfqpoint{1.292584in}{2.727461in}}%
\pgfpathlineto{\pgfqpoint{1.342710in}{2.742228in}}%
\pgfpathlineto{\pgfqpoint{1.392746in}{2.756969in}}%
\pgfpathlineto{\pgfqpoint{1.420443in}{2.681175in}}%
\pgfpathlineto{\pgfqpoint{1.448167in}{2.605308in}}%
\pgfpathlineto{\pgfqpoint{1.398134in}{2.590416in}}%
\pgfpathlineto{\pgfqpoint{1.348009in}{2.575496in}}%
\pgfpathclose%
\pgfusepath{stroke,fill}%
\end{pgfscope}%
\begin{pgfscope}%
\pgfpathrectangle{\pgfqpoint{0.050000in}{0.083252in}}{\pgfqpoint{4.900000in}{4.900000in}}%
\pgfusepath{clip}%
\pgfsetbuttcap%
\pgfsetroundjoin%
\definecolor{currentfill}{rgb}{0.922430,0.919231,0.953618}%
\pgfsetfillcolor{currentfill}%
\pgfsetfillopacity{0.950000}%
\pgfsetlinewidth{0.200750pt}%
\definecolor{currentstroke}{rgb}{0.200000,0.200000,0.200000}%
\pgfsetstrokecolor{currentstroke}%
\pgfsetdash{}{0pt}%
\pgfpathmoveto{\pgfqpoint{2.658089in}{1.234176in}}%
\pgfpathlineto{\pgfqpoint{2.629680in}{1.284618in}}%
\pgfpathlineto{\pgfqpoint{2.601335in}{1.339117in}}%
\pgfpathlineto{\pgfqpoint{2.650250in}{1.357041in}}%
\pgfpathlineto{\pgfqpoint{2.699081in}{1.374884in}}%
\pgfpathlineto{\pgfqpoint{2.727472in}{1.320472in}}%
\pgfpathlineto{\pgfqpoint{2.755935in}{1.269927in}}%
\pgfpathlineto{\pgfqpoint{2.707055in}{1.252090in}}%
\pgfpathlineto{\pgfqpoint{2.658089in}{1.234176in}}%
\pgfpathclose%
\pgfusepath{stroke,fill}%
\end{pgfscope}%
\begin{pgfscope}%
\pgfpathrectangle{\pgfqpoint{0.050000in}{0.083252in}}{\pgfqpoint{4.900000in}{4.900000in}}%
\pgfusepath{clip}%
\pgfsetbuttcap%
\pgfsetroundjoin%
\definecolor{currentfill}{rgb}{0.283968,0.254441,0.571857}%
\pgfsetfillcolor{currentfill}%
\pgfsetfillopacity{0.950000}%
\pgfsetlinewidth{0.200750pt}%
\definecolor{currentstroke}{rgb}{0.200000,0.200000,0.200000}%
\pgfsetstrokecolor{currentstroke}%
\pgfsetdash{}{0pt}%
\pgfpathmoveto{\pgfqpoint{1.559329in}{2.301111in}}%
\pgfpathlineto{\pgfqpoint{1.531499in}{2.377270in}}%
\pgfpathlineto{\pgfqpoint{1.503695in}{2.453356in}}%
\pgfpathlineto{\pgfqpoint{1.553635in}{2.468374in}}%
\pgfpathlineto{\pgfqpoint{1.603485in}{2.483364in}}%
\pgfpathlineto{\pgfqpoint{1.631286in}{2.407431in}}%
\pgfpathlineto{\pgfqpoint{1.659114in}{2.331425in}}%
\pgfpathlineto{\pgfqpoint{1.609267in}{2.316282in}}%
\pgfpathlineto{\pgfqpoint{1.559329in}{2.301111in}}%
\pgfpathclose%
\pgfusepath{stroke,fill}%
\end{pgfscope}%
\begin{pgfscope}%
\pgfpathrectangle{\pgfqpoint{0.050000in}{0.083252in}}{\pgfqpoint{4.900000in}{4.900000in}}%
\pgfusepath{clip}%
\pgfsetbuttcap%
\pgfsetroundjoin%
\definecolor{currentfill}{rgb}{0.820992,0.813610,0.892964}%
\pgfsetfillcolor{currentfill}%
\pgfsetfillopacity{0.950000}%
\pgfsetlinewidth{0.200750pt}%
\definecolor{currentstroke}{rgb}{0.200000,0.200000,0.200000}%
\pgfsetstrokecolor{currentstroke}%
\pgfsetdash{}{0pt}%
\pgfpathmoveto{\pgfqpoint{2.348428in}{1.388469in}}%
\pgfpathlineto{\pgfqpoint{2.320269in}{1.454780in}}%
\pgfpathlineto{\pgfqpoint{2.292133in}{1.523860in}}%
\pgfpathlineto{\pgfqpoint{2.341340in}{1.541012in}}%
\pgfpathlineto{\pgfqpoint{2.390460in}{1.558133in}}%
\pgfpathlineto{\pgfqpoint{2.418608in}{1.489757in}}%
\pgfpathlineto{\pgfqpoint{2.446786in}{1.424095in}}%
\pgfpathlineto{\pgfqpoint{2.397650in}{1.406318in}}%
\pgfpathlineto{\pgfqpoint{2.348428in}{1.388469in}}%
\pgfpathclose%
\pgfusepath{stroke,fill}%
\end{pgfscope}%
\begin{pgfscope}%
\pgfpathrectangle{\pgfqpoint{0.050000in}{0.083252in}}{\pgfqpoint{4.900000in}{4.900000in}}%
\pgfusepath{clip}%
\pgfsetbuttcap%
\pgfsetroundjoin%
\definecolor{currentfill}{rgb}{0.445075,0.422191,0.668189}%
\pgfsetfillcolor{currentfill}%
\pgfsetfillopacity{0.950000}%
\pgfsetlinewidth{0.200750pt}%
\definecolor{currentstroke}{rgb}{0.200000,0.200000,0.200000}%
\pgfsetstrokecolor{currentstroke}%
\pgfsetdash{}{0pt}%
\pgfpathmoveto{\pgfqpoint{1.770694in}{2.026681in}}%
\pgfpathlineto{\pgfqpoint{1.742759in}{2.102972in}}%
\pgfpathlineto{\pgfqpoint{1.714850in}{2.179195in}}%
\pgfpathlineto{\pgfqpoint{1.764604in}{2.194465in}}%
\pgfpathlineto{\pgfqpoint{1.814267in}{2.209708in}}%
\pgfpathlineto{\pgfqpoint{1.842173in}{2.133646in}}%
\pgfpathlineto{\pgfqpoint{1.870105in}{2.057523in}}%
\pgfpathlineto{\pgfqpoint{1.820445in}{2.042114in}}%
\pgfpathlineto{\pgfqpoint{1.770694in}{2.026681in}}%
\pgfpathclose%
\pgfusepath{stroke,fill}%
\end{pgfscope}%
\begin{pgfscope}%
\pgfpathrectangle{\pgfqpoint{0.050000in}{0.083252in}}{\pgfqpoint{4.900000in}{4.900000in}}%
\pgfusepath{clip}%
\pgfsetbuttcap%
\pgfsetroundjoin%
\definecolor{currentfill}{rgb}{0.600215,0.583729,0.760953}%
\pgfsetfillcolor{currentfill}%
\pgfsetfillopacity{0.950000}%
\pgfsetlinewidth{0.200750pt}%
\definecolor{currentstroke}{rgb}{0.200000,0.200000,0.200000}%
\pgfsetstrokecolor{currentstroke}%
\pgfsetdash{}{0pt}%
\pgfpathmoveto{\pgfqpoint{1.982093in}{1.753006in}}%
\pgfpathlineto{\pgfqpoint{1.954058in}{1.829022in}}%
\pgfpathlineto{\pgfqpoint{1.926048in}{1.905170in}}%
\pgfpathlineto{\pgfqpoint{1.975614in}{1.920772in}}%
\pgfpathlineto{\pgfqpoint{2.025090in}{1.936363in}}%
\pgfpathlineto{\pgfqpoint{2.053096in}{1.860507in}}%
\pgfpathlineto{\pgfqpoint{2.081126in}{1.784860in}}%
\pgfpathlineto{\pgfqpoint{2.031654in}{1.768931in}}%
\pgfpathlineto{\pgfqpoint{1.982093in}{1.753006in}}%
\pgfpathclose%
\pgfusepath{stroke,fill}%
\end{pgfscope}%
\begin{pgfscope}%
\pgfpathrectangle{\pgfqpoint{0.050000in}{0.083252in}}{\pgfqpoint{4.900000in}{4.900000in}}%
\pgfusepath{clip}%
\pgfsetbuttcap%
\pgfsetroundjoin%
\definecolor{currentfill}{rgb}{0.988066,0.987574,0.992864}%
\pgfsetfillcolor{currentfill}%
\pgfsetfillopacity{0.950000}%
\pgfsetlinewidth{0.200750pt}%
\definecolor{currentstroke}{rgb}{0.200000,0.200000,0.200000}%
\pgfsetstrokecolor{currentstroke}%
\pgfsetdash{}{0pt}%
\pgfpathmoveto{\pgfqpoint{3.221412in}{1.137062in}}%
\pgfpathlineto{\pgfqpoint{3.192130in}{1.170311in}}%
\pgfpathlineto{\pgfqpoint{3.162964in}{1.205661in}}%
\pgfpathlineto{\pgfqpoint{3.211377in}{1.222452in}}%
\pgfpathlineto{\pgfqpoint{3.259704in}{1.239212in}}%
\pgfpathlineto{\pgfqpoint{3.288946in}{1.203460in}}%
\pgfpathlineto{\pgfqpoint{3.318308in}{1.169817in}}%
\pgfpathlineto{\pgfqpoint{3.269904in}{1.153450in}}%
\pgfpathlineto{\pgfqpoint{3.221412in}{1.137062in}}%
\pgfpathclose%
\pgfusepath{stroke,fill}%
\end{pgfscope}%
\begin{pgfscope}%
\pgfpathrectangle{\pgfqpoint{0.050000in}{0.083252in}}{\pgfqpoint{4.900000in}{4.900000in}}%
\pgfusepath{clip}%
\pgfsetbuttcap%
\pgfsetroundjoin%
\definecolor{currentfill}{rgb}{0.952265,0.950296,0.971457}%
\pgfsetfillcolor{currentfill}%
\pgfsetfillopacity{0.950000}%
\pgfsetlinewidth{0.200750pt}%
\definecolor{currentstroke}{rgb}{0.200000,0.200000,0.200000}%
\pgfsetstrokecolor{currentstroke}%
\pgfsetdash{}{0pt}%
\pgfpathmoveto{\pgfqpoint{2.813110in}{1.180321in}}%
\pgfpathlineto{\pgfqpoint{2.784479in}{1.223260in}}%
\pgfpathlineto{\pgfqpoint{2.755935in}{1.269927in}}%
\pgfpathlineto{\pgfqpoint{2.804731in}{1.287691in}}%
\pgfpathlineto{\pgfqpoint{2.853441in}{1.305384in}}%
\pgfpathlineto{\pgfqpoint{2.882045in}{1.258477in}}%
\pgfpathlineto{\pgfqpoint{2.910743in}{1.215167in}}%
\pgfpathlineto{\pgfqpoint{2.861970in}{1.197769in}}%
\pgfpathlineto{\pgfqpoint{2.813110in}{1.180321in}}%
\pgfpathclose%
\pgfusepath{stroke,fill}%
\end{pgfscope}%
\begin{pgfscope}%
\pgfpathrectangle{\pgfqpoint{0.050000in}{0.083252in}}{\pgfqpoint{4.900000in}{4.900000in}}%
\pgfusepath{clip}%
\pgfsetbuttcap%
\pgfsetroundjoin%
\definecolor{currentfill}{rgb}{0.755356,0.745267,0.853718}%
\pgfsetfillcolor{currentfill}%
\pgfsetfillopacity{0.950000}%
\pgfsetlinewidth{0.200750pt}%
\definecolor{currentstroke}{rgb}{0.200000,0.200000,0.200000}%
\pgfsetstrokecolor{currentstroke}%
\pgfsetdash{}{0pt}%
\pgfpathmoveto{\pgfqpoint{2.193459in}{1.489471in}}%
\pgfpathlineto{\pgfqpoint{2.165348in}{1.561478in}}%
\pgfpathlineto{\pgfqpoint{2.137254in}{1.635021in}}%
\pgfpathlineto{\pgfqpoint{2.186633in}{1.651491in}}%
\pgfpathlineto{\pgfqpoint{2.235924in}{1.667961in}}%
\pgfpathlineto{\pgfqpoint{2.264018in}{1.595105in}}%
\pgfpathlineto{\pgfqpoint{2.292133in}{1.523860in}}%
\pgfpathlineto{\pgfqpoint{2.242839in}{1.506679in}}%
\pgfpathlineto{\pgfqpoint{2.193459in}{1.489471in}}%
\pgfpathclose%
\pgfusepath{stroke,fill}%
\end{pgfscope}%
\begin{pgfscope}%
\pgfpathrectangle{\pgfqpoint{0.050000in}{0.083252in}}{\pgfqpoint{4.900000in}{4.900000in}}%
\pgfusepath{clip}%
\pgfsetbuttcap%
\pgfsetroundjoin%
\definecolor{currentfill}{rgb}{0.253856,0.363399,0.751895}%
\pgfsetfillcolor{currentfill}%
\pgfsetfillopacity{0.950000}%
\pgfsetlinewidth{0.200750pt}%
\definecolor{currentstroke}{rgb}{0.200000,0.200000,0.200000}%
\pgfsetstrokecolor{currentstroke}%
\pgfsetdash{}{0pt}%
\pgfpathmoveto{\pgfqpoint{0.980616in}{2.972393in}}%
\pgfpathlineto{\pgfqpoint{0.953046in}{3.048278in}}%
\pgfpathlineto{\pgfqpoint{0.925503in}{3.124090in}}%
\pgfpathlineto{\pgfqpoint{0.976001in}{3.138508in}}%
\pgfpathlineto{\pgfqpoint{1.026407in}{3.152901in}}%
\pgfpathlineto{\pgfqpoint{1.053949in}{3.077240in}}%
\pgfpathlineto{\pgfqpoint{1.081517in}{3.001507in}}%
\pgfpathlineto{\pgfqpoint{1.031113in}{2.986963in}}%
\pgfpathlineto{\pgfqpoint{0.980616in}{2.972393in}}%
\pgfpathclose%
\pgfusepath{stroke,fill}%
\end{pgfscope}%
\begin{pgfscope}%
\pgfpathrectangle{\pgfqpoint{0.050000in}{0.083252in}}{\pgfqpoint{4.900000in}{4.900000in}}%
\pgfusepath{clip}%
\pgfsetbuttcap%
\pgfsetroundjoin%
\definecolor{currentfill}{rgb}{0.247213,0.292810,0.658055}%
\pgfsetfillcolor{currentfill}%
\pgfsetfillopacity{0.950000}%
\pgfsetlinewidth{0.200750pt}%
\definecolor{currentstroke}{rgb}{0.200000,0.200000,0.200000}%
\pgfsetstrokecolor{currentstroke}%
\pgfsetdash{}{0pt}%
\pgfpathmoveto{\pgfqpoint{1.192057in}{2.697845in}}%
\pgfpathlineto{\pgfqpoint{1.164382in}{2.773870in}}%
\pgfpathlineto{\pgfqpoint{1.136734in}{2.849822in}}%
\pgfpathlineto{\pgfqpoint{1.187045in}{2.864491in}}%
\pgfpathlineto{\pgfqpoint{1.237265in}{2.879134in}}%
\pgfpathlineto{\pgfqpoint{1.264911in}{2.803334in}}%
\pgfpathlineto{\pgfqpoint{1.292584in}{2.727461in}}%
\pgfpathlineto{\pgfqpoint{1.242366in}{2.712666in}}%
\pgfpathlineto{\pgfqpoint{1.192057in}{2.697845in}}%
\pgfpathclose%
\pgfusepath{stroke,fill}%
\end{pgfscope}%
\begin{pgfscope}%
\pgfpathrectangle{\pgfqpoint{0.050000in}{0.083252in}}{\pgfqpoint{4.900000in}{4.900000in}}%
\pgfusepath{clip}%
\pgfsetbuttcap%
\pgfsetroundjoin%
\definecolor{currentfill}{rgb}{1.000000,1.000000,1.000000}%
\pgfsetfillcolor{currentfill}%
\pgfsetfillopacity{0.950000}%
\pgfsetlinewidth{0.200750pt}%
\definecolor{currentstroke}{rgb}{0.200000,0.200000,0.200000}%
\pgfsetstrokecolor{currentstroke}%
\pgfsetdash{}{0pt}%
\pgfpathmoveto{\pgfqpoint{3.630282in}{1.110305in}}%
\pgfpathlineto{\pgfqpoint{3.600313in}{1.140474in}}%
\pgfpathlineto{\pgfqpoint{3.570446in}{1.171011in}}%
\pgfpathlineto{\pgfqpoint{3.618493in}{1.186917in}}%
\pgfpathlineto{\pgfqpoint{3.666452in}{1.202799in}}%
\pgfpathlineto{\pgfqpoint{3.696399in}{1.171929in}}%
\pgfpathlineto{\pgfqpoint{3.726448in}{1.141427in}}%
\pgfpathlineto{\pgfqpoint{3.678410in}{1.125880in}}%
\pgfpathlineto{\pgfqpoint{3.630282in}{1.110305in}}%
\pgfpathclose%
\pgfusepath{stroke,fill}%
\end{pgfscope}%
\begin{pgfscope}%
\pgfpathrectangle{\pgfqpoint{0.050000in}{0.083252in}}{\pgfqpoint{4.900000in}{4.900000in}}%
\pgfusepath{clip}%
\pgfsetbuttcap%
\pgfsetroundjoin%
\definecolor{currentfill}{rgb}{0.240569,0.222222,0.564214}%
\pgfsetfillcolor{currentfill}%
\pgfsetfillopacity{0.950000}%
\pgfsetlinewidth{0.200750pt}%
\definecolor{currentstroke}{rgb}{0.200000,0.200000,0.200000}%
\pgfsetstrokecolor{currentstroke}%
\pgfsetdash{}{0pt}%
\pgfpathmoveto{\pgfqpoint{1.403542in}{2.423238in}}%
\pgfpathlineto{\pgfqpoint{1.375762in}{2.499404in}}%
\pgfpathlineto{\pgfqpoint{1.348009in}{2.575496in}}%
\pgfpathlineto{\pgfqpoint{1.398134in}{2.590416in}}%
\pgfpathlineto{\pgfqpoint{1.448167in}{2.605308in}}%
\pgfpathlineto{\pgfqpoint{1.475918in}{2.529369in}}%
\pgfpathlineto{\pgfqpoint{1.503695in}{2.453356in}}%
\pgfpathlineto{\pgfqpoint{1.453664in}{2.438311in}}%
\pgfpathlineto{\pgfqpoint{1.403542in}{2.423238in}}%
\pgfpathclose%
\pgfusepath{stroke,fill}%
\end{pgfscope}%
\begin{pgfscope}%
\pgfpathrectangle{\pgfqpoint{0.050000in}{0.083252in}}{\pgfqpoint{4.900000in}{4.900000in}}%
\pgfusepath{clip}%
\pgfsetbuttcap%
\pgfsetroundjoin%
\definecolor{currentfill}{rgb}{0.361538,0.335210,0.618239}%
\pgfsetfillcolor{currentfill}%
\pgfsetfillopacity{0.950000}%
\pgfsetlinewidth{0.200750pt}%
\definecolor{currentstroke}{rgb}{0.200000,0.200000,0.200000}%
\pgfsetstrokecolor{currentstroke}%
\pgfsetdash{}{0pt}%
\pgfpathmoveto{\pgfqpoint{1.615071in}{2.148573in}}%
\pgfpathlineto{\pgfqpoint{1.587187in}{2.224879in}}%
\pgfpathlineto{\pgfqpoint{1.559329in}{2.301111in}}%
\pgfpathlineto{\pgfqpoint{1.609267in}{2.316282in}}%
\pgfpathlineto{\pgfqpoint{1.659114in}{2.331425in}}%
\pgfpathlineto{\pgfqpoint{1.686969in}{2.255346in}}%
\pgfpathlineto{\pgfqpoint{1.714850in}{2.179195in}}%
\pgfpathlineto{\pgfqpoint{1.665006in}{2.163898in}}%
\pgfpathlineto{\pgfqpoint{1.615071in}{2.148573in}}%
\pgfpathclose%
\pgfusepath{stroke,fill}%
\end{pgfscope}%
\begin{pgfscope}%
\pgfpathrectangle{\pgfqpoint{0.050000in}{0.083252in}}{\pgfqpoint{4.900000in}{4.900000in}}%
\pgfusepath{clip}%
\pgfsetbuttcap%
\pgfsetroundjoin%
\definecolor{currentfill}{rgb}{0.522645,0.502960,0.714571}%
\pgfsetfillcolor{currentfill}%
\pgfsetfillopacity{0.950000}%
\pgfsetlinewidth{0.200750pt}%
\definecolor{currentstroke}{rgb}{0.200000,0.200000,0.200000}%
\pgfsetstrokecolor{currentstroke}%
\pgfsetdash{}{0pt}%
\pgfpathmoveto{\pgfqpoint{1.826644in}{1.873933in}}%
\pgfpathlineto{\pgfqpoint{1.798656in}{1.950329in}}%
\pgfpathlineto{\pgfqpoint{1.770694in}{2.026681in}}%
\pgfpathlineto{\pgfqpoint{1.820445in}{2.042114in}}%
\pgfpathlineto{\pgfqpoint{1.870105in}{2.057523in}}%
\pgfpathlineto{\pgfqpoint{1.898063in}{1.981355in}}%
\pgfpathlineto{\pgfqpoint{1.926048in}{1.905170in}}%
\pgfpathlineto{\pgfqpoint{1.876391in}{1.889557in}}%
\pgfpathlineto{\pgfqpoint{1.826644in}{1.873933in}}%
\pgfpathclose%
\pgfusepath{stroke,fill}%
\end{pgfscope}%
\begin{pgfscope}%
\pgfpathrectangle{\pgfqpoint{0.050000in}{0.083252in}}{\pgfqpoint{4.900000in}{4.900000in}}%
\pgfusepath{clip}%
\pgfsetbuttcap%
\pgfsetroundjoin%
\definecolor{currentfill}{rgb}{0.976132,0.975148,0.985729}%
\pgfsetfillcolor{currentfill}%
\pgfsetfillopacity{0.950000}%
\pgfsetlinewidth{0.200750pt}%
\definecolor{currentstroke}{rgb}{0.200000,0.200000,0.200000}%
\pgfsetstrokecolor{currentstroke}%
\pgfsetdash{}{0pt}%
\pgfpathmoveto{\pgfqpoint{2.968442in}{1.138209in}}%
\pgfpathlineto{\pgfqpoint{2.939541in}{1.175200in}}%
\pgfpathlineto{\pgfqpoint{2.910743in}{1.215167in}}%
\pgfpathlineto{\pgfqpoint{2.959430in}{1.232517in}}%
\pgfpathlineto{\pgfqpoint{3.008031in}{1.249818in}}%
\pgfpathlineto{\pgfqpoint{3.036899in}{1.209438in}}%
\pgfpathlineto{\pgfqpoint{3.065877in}{1.171990in}}%
\pgfpathlineto{\pgfqpoint{3.017203in}{1.155113in}}%
\pgfpathlineto{\pgfqpoint{2.968442in}{1.138209in}}%
\pgfpathclose%
\pgfusepath{stroke,fill}%
\end{pgfscope}%
\begin{pgfscope}%
\pgfpathrectangle{\pgfqpoint{0.050000in}{0.083252in}}{\pgfqpoint{4.900000in}{4.900000in}}%
\pgfusepath{clip}%
\pgfsetbuttcap%
\pgfsetroundjoin%
\definecolor{currentfill}{rgb}{0.683752,0.670711,0.810903}%
\pgfsetfillcolor{currentfill}%
\pgfsetfillopacity{0.950000}%
\pgfsetlinewidth{0.200750pt}%
\definecolor{currentstroke}{rgb}{0.200000,0.200000,0.200000}%
\pgfsetstrokecolor{currentstroke}%
\pgfsetdash{}{0pt}%
\pgfpathmoveto{\pgfqpoint{2.038229in}{1.602100in}}%
\pgfpathlineto{\pgfqpoint{2.010150in}{1.677282in}}%
\pgfpathlineto{\pgfqpoint{1.982093in}{1.753006in}}%
\pgfpathlineto{\pgfqpoint{2.031654in}{1.768931in}}%
\pgfpathlineto{\pgfqpoint{2.081126in}{1.784860in}}%
\pgfpathlineto{\pgfqpoint{2.109179in}{1.709608in}}%
\pgfpathlineto{\pgfqpoint{2.137254in}{1.635021in}}%
\pgfpathlineto{\pgfqpoint{2.087786in}{1.618556in}}%
\pgfpathlineto{\pgfqpoint{2.038229in}{1.602100in}}%
\pgfpathclose%
\pgfusepath{stroke,fill}%
\end{pgfscope}%
\begin{pgfscope}%
\pgfpathrectangle{\pgfqpoint{0.050000in}{0.083252in}}{\pgfqpoint{4.900000in}{4.900000in}}%
\pgfusepath{clip}%
\pgfsetbuttcap%
\pgfsetroundjoin%
\definecolor{currentfill}{rgb}{0.994033,0.993787,0.996432}%
\pgfsetfillcolor{currentfill}%
\pgfsetfillopacity{0.950000}%
\pgfsetlinewidth{0.200750pt}%
\definecolor{currentstroke}{rgb}{0.200000,0.200000,0.200000}%
\pgfsetstrokecolor{currentstroke}%
\pgfsetdash{}{0pt}%
\pgfpathmoveto{\pgfqpoint{3.377377in}{1.107156in}}%
\pgfpathlineto{\pgfqpoint{3.347786in}{1.137861in}}%
\pgfpathlineto{\pgfqpoint{3.318308in}{1.169817in}}%
\pgfpathlineto{\pgfqpoint{3.366625in}{1.186163in}}%
\pgfpathlineto{\pgfqpoint{3.414854in}{1.202488in}}%
\pgfpathlineto{\pgfqpoint{3.444413in}{1.170170in}}%
\pgfpathlineto{\pgfqpoint{3.474088in}{1.139128in}}%
\pgfpathlineto{\pgfqpoint{3.425777in}{1.123153in}}%
\pgfpathlineto{\pgfqpoint{3.377377in}{1.107156in}}%
\pgfpathclose%
\pgfusepath{stroke,fill}%
\end{pgfscope}%
\begin{pgfscope}%
\pgfpathrectangle{\pgfqpoint{0.050000in}{0.083252in}}{\pgfqpoint{4.900000in}{4.900000in}}%
\pgfusepath{clip}%
\pgfsetbuttcap%
\pgfsetroundjoin%
\definecolor{currentfill}{rgb}{0.880661,0.875740,0.928643}%
\pgfsetfillcolor{currentfill}%
\pgfsetfillopacity{0.950000}%
\pgfsetlinewidth{0.200750pt}%
\definecolor{currentstroke}{rgb}{0.200000,0.200000,0.200000}%
\pgfsetstrokecolor{currentstroke}%
\pgfsetdash{}{0pt}%
\pgfpathmoveto{\pgfqpoint{2.404826in}{1.266526in}}%
\pgfpathlineto{\pgfqpoint{2.376612in}{1.325537in}}%
\pgfpathlineto{\pgfqpoint{2.348428in}{1.388469in}}%
\pgfpathlineto{\pgfqpoint{2.397650in}{1.406318in}}%
\pgfpathlineto{\pgfqpoint{2.446786in}{1.424095in}}%
\pgfpathlineto{\pgfqpoint{2.474998in}{1.361687in}}%
\pgfpathlineto{\pgfqpoint{2.503250in}{1.303010in}}%
\pgfpathlineto{\pgfqpoint{2.454080in}{1.284818in}}%
\pgfpathlineto{\pgfqpoint{2.404826in}{1.266526in}}%
\pgfpathclose%
\pgfusepath{stroke,fill}%
\end{pgfscope}%
\begin{pgfscope}%
\pgfpathrectangle{\pgfqpoint{0.050000in}{0.083252in}}{\pgfqpoint{4.900000in}{4.900000in}}%
\pgfusepath{clip}%
\pgfsetbuttcap%
\pgfsetroundjoin%
\definecolor{currentfill}{rgb}{0.250411,0.326797,0.703237}%
\pgfsetfillcolor{currentfill}%
\pgfsetfillopacity{0.950000}%
\pgfsetlinewidth{0.200750pt}%
\definecolor{currentstroke}{rgb}{0.200000,0.200000,0.200000}%
\pgfsetstrokecolor{currentstroke}%
\pgfsetdash{}{0pt}%
\pgfpathmoveto{\pgfqpoint{1.035836in}{2.820403in}}%
\pgfpathlineto{\pgfqpoint{1.008213in}{2.896435in}}%
\pgfpathlineto{\pgfqpoint{0.980616in}{2.972393in}}%
\pgfpathlineto{\pgfqpoint{1.031113in}{2.986963in}}%
\pgfpathlineto{\pgfqpoint{1.081517in}{3.001507in}}%
\pgfpathlineto{\pgfqpoint{1.109112in}{2.925701in}}%
\pgfpathlineto{\pgfqpoint{1.136734in}{2.849822in}}%
\pgfpathlineto{\pgfqpoint{1.086331in}{2.835126in}}%
\pgfpathlineto{\pgfqpoint{1.035836in}{2.820403in}}%
\pgfpathclose%
\pgfusepath{stroke,fill}%
\end{pgfscope}%
\begin{pgfscope}%
\pgfpathrectangle{\pgfqpoint{0.050000in}{0.083252in}}{\pgfqpoint{4.900000in}{4.900000in}}%
\pgfusepath{clip}%
\pgfsetbuttcap%
\pgfsetroundjoin%
\definecolor{currentfill}{rgb}{0.928397,0.925444,0.957186}%
\pgfsetfillcolor{currentfill}%
\pgfsetfillopacity{0.950000}%
\pgfsetlinewidth{0.200750pt}%
\definecolor{currentstroke}{rgb}{0.200000,0.200000,0.200000}%
\pgfsetstrokecolor{currentstroke}%
\pgfsetdash{}{0pt}%
\pgfpathmoveto{\pgfqpoint{2.559901in}{1.198107in}}%
\pgfpathlineto{\pgfqpoint{2.531548in}{1.248421in}}%
\pgfpathlineto{\pgfqpoint{2.503250in}{1.303010in}}%
\pgfpathlineto{\pgfqpoint{2.552335in}{1.321108in}}%
\pgfpathlineto{\pgfqpoint{2.601335in}{1.339117in}}%
\pgfpathlineto{\pgfqpoint{2.629680in}{1.284618in}}%
\pgfpathlineto{\pgfqpoint{2.658089in}{1.234176in}}%
\pgfpathlineto{\pgfqpoint{2.609038in}{1.216183in}}%
\pgfpathlineto{\pgfqpoint{2.559901in}{1.198107in}}%
\pgfpathclose%
\pgfusepath{stroke,fill}%
\end{pgfscope}%
\begin{pgfscope}%
\pgfpathrectangle{\pgfqpoint{0.050000in}{0.083252in}}{\pgfqpoint{4.900000in}{4.900000in}}%
\pgfusepath{clip}%
\pgfsetbuttcap%
\pgfsetroundjoin%
\definecolor{currentfill}{rgb}{0.244014,0.258824,0.612872}%
\pgfsetfillcolor{currentfill}%
\pgfsetfillopacity{0.950000}%
\pgfsetlinewidth{0.200750pt}%
\definecolor{currentstroke}{rgb}{0.200000,0.200000,0.200000}%
\pgfsetstrokecolor{currentstroke}%
\pgfsetdash{}{0pt}%
\pgfpathmoveto{\pgfqpoint{1.247486in}{2.545574in}}%
\pgfpathlineto{\pgfqpoint{1.219758in}{2.621746in}}%
\pgfpathlineto{\pgfqpoint{1.192057in}{2.697845in}}%
\pgfpathlineto{\pgfqpoint{1.242366in}{2.712666in}}%
\pgfpathlineto{\pgfqpoint{1.292584in}{2.727461in}}%
\pgfpathlineto{\pgfqpoint{1.320283in}{2.651515in}}%
\pgfpathlineto{\pgfqpoint{1.348009in}{2.575496in}}%
\pgfpathlineto{\pgfqpoint{1.297794in}{2.560549in}}%
\pgfpathlineto{\pgfqpoint{1.247486in}{2.545574in}}%
\pgfpathclose%
\pgfusepath{stroke,fill}%
\end{pgfscope}%
\begin{pgfscope}%
\pgfpathrectangle{\pgfqpoint{0.050000in}{0.083252in}}{\pgfqpoint{4.900000in}{4.900000in}}%
\pgfusepath{clip}%
\pgfsetbuttcap%
\pgfsetroundjoin%
\definecolor{currentfill}{rgb}{0.826959,0.819823,0.896532}%
\pgfsetfillcolor{currentfill}%
\pgfsetfillopacity{0.950000}%
\pgfsetlinewidth{0.200750pt}%
\definecolor{currentstroke}{rgb}{0.200000,0.200000,0.200000}%
\pgfsetstrokecolor{currentstroke}%
\pgfsetdash{}{0pt}%
\pgfpathmoveto{\pgfqpoint{2.249727in}{1.352553in}}%
\pgfpathlineto{\pgfqpoint{2.221586in}{1.419603in}}%
\pgfpathlineto{\pgfqpoint{2.193459in}{1.489471in}}%
\pgfpathlineto{\pgfqpoint{2.242839in}{1.506679in}}%
\pgfpathlineto{\pgfqpoint{2.292133in}{1.523860in}}%
\pgfpathlineto{\pgfqpoint{2.320269in}{1.454780in}}%
\pgfpathlineto{\pgfqpoint{2.348428in}{1.388469in}}%
\pgfpathlineto{\pgfqpoint{2.299120in}{1.370548in}}%
\pgfpathlineto{\pgfqpoint{2.249727in}{1.352553in}}%
\pgfpathclose%
\pgfusepath{stroke,fill}%
\end{pgfscope}%
\begin{pgfscope}%
\pgfpathrectangle{\pgfqpoint{0.050000in}{0.083252in}}{\pgfqpoint{4.900000in}{4.900000in}}%
\pgfusepath{clip}%
\pgfsetbuttcap%
\pgfsetroundjoin%
\definecolor{currentfill}{rgb}{0.283968,0.254441,0.571857}%
\pgfsetfillcolor{currentfill}%
\pgfsetfillopacity{0.950000}%
\pgfsetlinewidth{0.200750pt}%
\definecolor{currentstroke}{rgb}{0.200000,0.200000,0.200000}%
\pgfsetstrokecolor{currentstroke}%
\pgfsetdash{}{0pt}%
\pgfpathmoveto{\pgfqpoint{1.459181in}{2.270687in}}%
\pgfpathlineto{\pgfqpoint{1.431348in}{2.346999in}}%
\pgfpathlineto{\pgfqpoint{1.403542in}{2.423238in}}%
\pgfpathlineto{\pgfqpoint{1.453664in}{2.438311in}}%
\pgfpathlineto{\pgfqpoint{1.503695in}{2.453356in}}%
\pgfpathlineto{\pgfqpoint{1.531499in}{2.377270in}}%
\pgfpathlineto{\pgfqpoint{1.559329in}{2.301111in}}%
\pgfpathlineto{\pgfqpoint{1.509301in}{2.285913in}}%
\pgfpathlineto{\pgfqpoint{1.459181in}{2.270687in}}%
\pgfpathclose%
\pgfusepath{stroke,fill}%
\end{pgfscope}%
\begin{pgfscope}%
\pgfpathrectangle{\pgfqpoint{0.050000in}{0.083252in}}{\pgfqpoint{4.900000in}{4.900000in}}%
\pgfusepath{clip}%
\pgfsetbuttcap%
\pgfsetroundjoin%
\definecolor{currentfill}{rgb}{0.445075,0.422191,0.668189}%
\pgfsetfillcolor{currentfill}%
\pgfsetfillopacity{0.950000}%
\pgfsetlinewidth{0.200750pt}%
\definecolor{currentstroke}{rgb}{0.200000,0.200000,0.200000}%
\pgfsetstrokecolor{currentstroke}%
\pgfsetdash{}{0pt}%
\pgfpathmoveto{\pgfqpoint{1.670921in}{1.995744in}}%
\pgfpathlineto{\pgfqpoint{1.642983in}{2.072195in}}%
\pgfpathlineto{\pgfqpoint{1.615071in}{2.148573in}}%
\pgfpathlineto{\pgfqpoint{1.665006in}{2.163898in}}%
\pgfpathlineto{\pgfqpoint{1.714850in}{2.179195in}}%
\pgfpathlineto{\pgfqpoint{1.742759in}{2.102972in}}%
\pgfpathlineto{\pgfqpoint{1.770694in}{2.026681in}}%
\pgfpathlineto{\pgfqpoint{1.720853in}{2.011225in}}%
\pgfpathlineto{\pgfqpoint{1.670921in}{1.995744in}}%
\pgfpathclose%
\pgfusepath{stroke,fill}%
\end{pgfscope}%
\begin{pgfscope}%
\pgfpathrectangle{\pgfqpoint{0.050000in}{0.083252in}}{\pgfqpoint{4.900000in}{4.900000in}}%
\pgfusepath{clip}%
\pgfsetbuttcap%
\pgfsetroundjoin%
\definecolor{currentfill}{rgb}{0.606182,0.589942,0.764521}%
\pgfsetfillcolor{currentfill}%
\pgfsetfillopacity{0.950000}%
\pgfsetlinewidth{0.200750pt}%
\definecolor{currentstroke}{rgb}{0.200000,0.200000,0.200000}%
\pgfsetstrokecolor{currentstroke}%
\pgfsetdash{}{0pt}%
\pgfpathmoveto{\pgfqpoint{1.882698in}{1.721174in}}%
\pgfpathlineto{\pgfqpoint{1.854659in}{1.797525in}}%
\pgfpathlineto{\pgfqpoint{1.826644in}{1.873933in}}%
\pgfpathlineto{\pgfqpoint{1.876391in}{1.889557in}}%
\pgfpathlineto{\pgfqpoint{1.926048in}{1.905170in}}%
\pgfpathlineto{\pgfqpoint{1.954058in}{1.829022in}}%
\pgfpathlineto{\pgfqpoint{1.982093in}{1.753006in}}%
\pgfpathlineto{\pgfqpoint{1.932441in}{1.737087in}}%
\pgfpathlineto{\pgfqpoint{1.882698in}{1.721174in}}%
\pgfpathclose%
\pgfusepath{stroke,fill}%
\end{pgfscope}%
\begin{pgfscope}%
\pgfpathrectangle{\pgfqpoint{0.050000in}{0.083252in}}{\pgfqpoint{4.900000in}{4.900000in}}%
\pgfusepath{clip}%
\pgfsetbuttcap%
\pgfsetroundjoin%
\definecolor{currentfill}{rgb}{0.958231,0.956509,0.975025}%
\pgfsetfillcolor{currentfill}%
\pgfsetfillopacity{0.950000}%
\pgfsetlinewidth{0.200750pt}%
\definecolor{currentstroke}{rgb}{0.200000,0.200000,0.200000}%
\pgfsetstrokecolor{currentstroke}%
\pgfsetdash{}{0pt}%
\pgfpathmoveto{\pgfqpoint{2.715132in}{1.145277in}}%
\pgfpathlineto{\pgfqpoint{2.686571in}{1.187790in}}%
\pgfpathlineto{\pgfqpoint{2.658089in}{1.234176in}}%
\pgfpathlineto{\pgfqpoint{2.707055in}{1.252090in}}%
\pgfpathlineto{\pgfqpoint{2.755935in}{1.269927in}}%
\pgfpathlineto{\pgfqpoint{2.784479in}{1.223260in}}%
\pgfpathlineto{\pgfqpoint{2.813110in}{1.180321in}}%
\pgfpathlineto{\pgfqpoint{2.764164in}{1.162824in}}%
\pgfpathlineto{\pgfqpoint{2.715132in}{1.145277in}}%
\pgfpathclose%
\pgfusepath{stroke,fill}%
\end{pgfscope}%
\begin{pgfscope}%
\pgfpathrectangle{\pgfqpoint{0.050000in}{0.083252in}}{\pgfqpoint{4.900000in}{4.900000in}}%
\pgfusepath{clip}%
\pgfsetbuttcap%
\pgfsetroundjoin%
\definecolor{currentfill}{rgb}{0.988066,0.987574,0.992864}%
\pgfsetfillcolor{currentfill}%
\pgfsetfillopacity{0.950000}%
\pgfsetlinewidth{0.200750pt}%
\definecolor{currentstroke}{rgb}{0.200000,0.200000,0.200000}%
\pgfsetstrokecolor{currentstroke}%
\pgfsetdash{}{0pt}%
\pgfpathmoveto{\pgfqpoint{3.124163in}{1.104234in}}%
\pgfpathlineto{\pgfqpoint{3.094965in}{1.137071in}}%
\pgfpathlineto{\pgfqpoint{3.065877in}{1.171990in}}%
\pgfpathlineto{\pgfqpoint{3.114464in}{1.188840in}}%
\pgfpathlineto{\pgfqpoint{3.162964in}{1.205661in}}%
\pgfpathlineto{\pgfqpoint{3.192130in}{1.170311in}}%
\pgfpathlineto{\pgfqpoint{3.221412in}{1.137062in}}%
\pgfpathlineto{\pgfqpoint{3.172832in}{1.120657in}}%
\pgfpathlineto{\pgfqpoint{3.124163in}{1.104234in}}%
\pgfpathclose%
\pgfusepath{stroke,fill}%
\end{pgfscope}%
\begin{pgfscope}%
\pgfpathrectangle{\pgfqpoint{0.050000in}{0.083252in}}{\pgfqpoint{4.900000in}{4.900000in}}%
\pgfusepath{clip}%
\pgfsetbuttcap%
\pgfsetroundjoin%
\definecolor{currentfill}{rgb}{0.755356,0.745267,0.853718}%
\pgfsetfillcolor{currentfill}%
\pgfsetfillopacity{0.950000}%
\pgfsetlinewidth{0.200750pt}%
\definecolor{currentstroke}{rgb}{0.200000,0.200000,0.200000}%
\pgfsetstrokecolor{currentstroke}%
\pgfsetdash{}{0pt}%
\pgfpathmoveto{\pgfqpoint{2.094436in}{1.454994in}}%
\pgfpathlineto{\pgfqpoint{2.066325in}{1.527833in}}%
\pgfpathlineto{\pgfqpoint{2.038229in}{1.602100in}}%
\pgfpathlineto{\pgfqpoint{2.087786in}{1.618556in}}%
\pgfpathlineto{\pgfqpoint{2.137254in}{1.635021in}}%
\pgfpathlineto{\pgfqpoint{2.165348in}{1.561478in}}%
\pgfpathlineto{\pgfqpoint{2.193459in}{1.489471in}}%
\pgfpathlineto{\pgfqpoint{2.143991in}{1.472241in}}%
\pgfpathlineto{\pgfqpoint{2.094436in}{1.454994in}}%
\pgfpathclose%
\pgfusepath{stroke,fill}%
\end{pgfscope}%
\begin{pgfscope}%
\pgfpathrectangle{\pgfqpoint{0.050000in}{0.083252in}}{\pgfqpoint{4.900000in}{4.900000in}}%
\pgfusepath{clip}%
\pgfsetbuttcap%
\pgfsetroundjoin%
\definecolor{currentfill}{rgb}{0.247213,0.292810,0.658055}%
\pgfsetfillcolor{currentfill}%
\pgfsetfillopacity{0.950000}%
\pgfsetlinewidth{0.200750pt}%
\definecolor{currentstroke}{rgb}{0.200000,0.200000,0.200000}%
\pgfsetstrokecolor{currentstroke}%
\pgfsetdash{}{0pt}%
\pgfpathmoveto{\pgfqpoint{1.091163in}{2.668120in}}%
\pgfpathlineto{\pgfqpoint{1.063486in}{2.744298in}}%
\pgfpathlineto{\pgfqpoint{1.035836in}{2.820403in}}%
\pgfpathlineto{\pgfqpoint{1.086331in}{2.835126in}}%
\pgfpathlineto{\pgfqpoint{1.136734in}{2.849822in}}%
\pgfpathlineto{\pgfqpoint{1.164382in}{2.773870in}}%
\pgfpathlineto{\pgfqpoint{1.192057in}{2.697845in}}%
\pgfpathlineto{\pgfqpoint{1.141656in}{2.682996in}}%
\pgfpathlineto{\pgfqpoint{1.091163in}{2.668120in}}%
\pgfpathclose%
\pgfusepath{stroke,fill}%
\end{pgfscope}%
\begin{pgfscope}%
\pgfpathrectangle{\pgfqpoint{0.050000in}{0.083252in}}{\pgfqpoint{4.900000in}{4.900000in}}%
\pgfusepath{clip}%
\pgfsetbuttcap%
\pgfsetroundjoin%
\definecolor{currentfill}{rgb}{1.000000,1.000000,1.000000}%
\pgfsetfillcolor{currentfill}%
\pgfsetfillopacity{0.950000}%
\pgfsetlinewidth{0.200750pt}%
\definecolor{currentstroke}{rgb}{0.200000,0.200000,0.200000}%
\pgfsetstrokecolor{currentstroke}%
\pgfsetdash{}{0pt}%
\pgfpathmoveto{\pgfqpoint{3.533757in}{1.079068in}}%
\pgfpathlineto{\pgfqpoint{3.503872in}{1.108917in}}%
\pgfpathlineto{\pgfqpoint{3.474088in}{1.139128in}}%
\pgfpathlineto{\pgfqpoint{3.522311in}{1.155081in}}%
\pgfpathlineto{\pgfqpoint{3.570446in}{1.171011in}}%
\pgfpathlineto{\pgfqpoint{3.600313in}{1.140474in}}%
\pgfpathlineto{\pgfqpoint{3.630282in}{1.110305in}}%
\pgfpathlineto{\pgfqpoint{3.582064in}{1.094701in}}%
\pgfpathlineto{\pgfqpoint{3.533757in}{1.079068in}}%
\pgfpathclose%
\pgfusepath{stroke,fill}%
\end{pgfscope}%
\begin{pgfscope}%
\pgfpathrectangle{\pgfqpoint{0.050000in}{0.083252in}}{\pgfqpoint{4.900000in}{4.900000in}}%
\pgfusepath{clip}%
\pgfsetbuttcap%
\pgfsetroundjoin%
\definecolor{currentfill}{rgb}{0.240569,0.222222,0.564214}%
\pgfsetfillcolor{currentfill}%
\pgfsetfillopacity{0.950000}%
\pgfsetlinewidth{0.200750pt}%
\definecolor{currentstroke}{rgb}{0.200000,0.200000,0.200000}%
\pgfsetstrokecolor{currentstroke}%
\pgfsetdash{}{0pt}%
\pgfpathmoveto{\pgfqpoint{1.303023in}{2.393010in}}%
\pgfpathlineto{\pgfqpoint{1.275241in}{2.469329in}}%
\pgfpathlineto{\pgfqpoint{1.247486in}{2.545574in}}%
\pgfpathlineto{\pgfqpoint{1.297794in}{2.560549in}}%
\pgfpathlineto{\pgfqpoint{1.348009in}{2.575496in}}%
\pgfpathlineto{\pgfqpoint{1.375762in}{2.499404in}}%
\pgfpathlineto{\pgfqpoint{1.403542in}{2.423238in}}%
\pgfpathlineto{\pgfqpoint{1.353328in}{2.408138in}}%
\pgfpathlineto{\pgfqpoint{1.303023in}{2.393010in}}%
\pgfpathclose%
\pgfusepath{stroke,fill}%
\end{pgfscope}%
\begin{pgfscope}%
\pgfpathrectangle{\pgfqpoint{0.050000in}{0.083252in}}{\pgfqpoint{4.900000in}{4.900000in}}%
\pgfusepath{clip}%
\pgfsetbuttcap%
\pgfsetroundjoin%
\definecolor{currentfill}{rgb}{0.367505,0.341423,0.621807}%
\pgfsetfillcolor{currentfill}%
\pgfsetfillopacity{0.950000}%
\pgfsetlinewidth{0.200750pt}%
\definecolor{currentstroke}{rgb}{0.200000,0.200000,0.200000}%
\pgfsetstrokecolor{currentstroke}%
\pgfsetdash{}{0pt}%
\pgfpathmoveto{\pgfqpoint{1.514928in}{2.117841in}}%
\pgfpathlineto{\pgfqpoint{1.487041in}{2.194301in}}%
\pgfpathlineto{\pgfqpoint{1.459181in}{2.270687in}}%
\pgfpathlineto{\pgfqpoint{1.509301in}{2.285913in}}%
\pgfpathlineto{\pgfqpoint{1.559329in}{2.301111in}}%
\pgfpathlineto{\pgfqpoint{1.587187in}{2.224879in}}%
\pgfpathlineto{\pgfqpoint{1.615071in}{2.148573in}}%
\pgfpathlineto{\pgfqpoint{1.565045in}{2.133221in}}%
\pgfpathlineto{\pgfqpoint{1.514928in}{2.117841in}}%
\pgfpathclose%
\pgfusepath{stroke,fill}%
\end{pgfscope}%
\begin{pgfscope}%
\pgfpathrectangle{\pgfqpoint{0.050000in}{0.083252in}}{\pgfqpoint{4.900000in}{4.900000in}}%
\pgfusepath{clip}%
\pgfsetbuttcap%
\pgfsetroundjoin%
\definecolor{currentfill}{rgb}{0.522645,0.502960,0.714571}%
\pgfsetfillcolor{currentfill}%
\pgfsetfillopacity{0.950000}%
\pgfsetlinewidth{0.200750pt}%
\definecolor{currentstroke}{rgb}{0.200000,0.200000,0.200000}%
\pgfsetstrokecolor{currentstroke}%
\pgfsetdash{}{0pt}%
\pgfpathmoveto{\pgfqpoint{1.726877in}{1.842638in}}%
\pgfpathlineto{\pgfqpoint{1.698886in}{1.919222in}}%
\pgfpathlineto{\pgfqpoint{1.670921in}{1.995744in}}%
\pgfpathlineto{\pgfqpoint{1.720853in}{2.011225in}}%
\pgfpathlineto{\pgfqpoint{1.770694in}{2.026681in}}%
\pgfpathlineto{\pgfqpoint{1.798656in}{1.950329in}}%
\pgfpathlineto{\pgfqpoint{1.826644in}{1.873933in}}%
\pgfpathlineto{\pgfqpoint{1.776806in}{1.858293in}}%
\pgfpathlineto{\pgfqpoint{1.726877in}{1.842638in}}%
\pgfpathclose%
\pgfusepath{stroke,fill}%
\end{pgfscope}%
\begin{pgfscope}%
\pgfpathrectangle{\pgfqpoint{0.050000in}{0.083252in}}{\pgfqpoint{4.900000in}{4.900000in}}%
\pgfusepath{clip}%
\pgfsetbuttcap%
\pgfsetroundjoin%
\definecolor{currentfill}{rgb}{0.976132,0.975148,0.985729}%
\pgfsetfillcolor{currentfill}%
\pgfsetfillopacity{0.950000}%
\pgfsetlinewidth{0.200750pt}%
\definecolor{currentstroke}{rgb}{0.200000,0.200000,0.200000}%
\pgfsetstrokecolor{currentstroke}%
\pgfsetdash{}{0pt}%
\pgfpathmoveto{\pgfqpoint{2.870656in}{1.104326in}}%
\pgfpathlineto{\pgfqpoint{2.841834in}{1.140815in}}%
\pgfpathlineto{\pgfqpoint{2.813110in}{1.180321in}}%
\pgfpathlineto{\pgfqpoint{2.861970in}{1.197769in}}%
\pgfpathlineto{\pgfqpoint{2.910743in}{1.215167in}}%
\pgfpathlineto{\pgfqpoint{2.939541in}{1.175200in}}%
\pgfpathlineto{\pgfqpoint{2.968442in}{1.138209in}}%
\pgfpathlineto{\pgfqpoint{2.919593in}{1.121279in}}%
\pgfpathlineto{\pgfqpoint{2.870656in}{1.104326in}}%
\pgfpathclose%
\pgfusepath{stroke,fill}%
\end{pgfscope}%
\begin{pgfscope}%
\pgfpathrectangle{\pgfqpoint{0.050000in}{0.083252in}}{\pgfqpoint{4.900000in}{4.900000in}}%
\pgfusepath{clip}%
\pgfsetbuttcap%
\pgfsetroundjoin%
\definecolor{currentfill}{rgb}{0.683752,0.670711,0.810903}%
\pgfsetfillcolor{currentfill}%
\pgfsetfillopacity{0.950000}%
\pgfsetlinewidth{0.200750pt}%
\definecolor{currentstroke}{rgb}{0.200000,0.200000,0.200000}%
\pgfsetstrokecolor{currentstroke}%
\pgfsetdash{}{0pt}%
\pgfpathmoveto{\pgfqpoint{1.938846in}{1.569239in}}%
\pgfpathlineto{\pgfqpoint{1.910762in}{1.645006in}}%
\pgfpathlineto{\pgfqpoint{1.882698in}{1.721174in}}%
\pgfpathlineto{\pgfqpoint{1.932441in}{1.737087in}}%
\pgfpathlineto{\pgfqpoint{1.982093in}{1.753006in}}%
\pgfpathlineto{\pgfqpoint{2.010150in}{1.677282in}}%
\pgfpathlineto{\pgfqpoint{2.038229in}{1.602100in}}%
\pgfpathlineto{\pgfqpoint{1.988582in}{1.585659in}}%
\pgfpathlineto{\pgfqpoint{1.938846in}{1.569239in}}%
\pgfpathclose%
\pgfusepath{stroke,fill}%
\end{pgfscope}%
\begin{pgfscope}%
\pgfpathrectangle{\pgfqpoint{0.050000in}{0.083252in}}{\pgfqpoint{4.900000in}{4.900000in}}%
\pgfusepath{clip}%
\pgfsetbuttcap%
\pgfsetroundjoin%
\definecolor{currentfill}{rgb}{0.886628,0.881953,0.932211}%
\pgfsetfillcolor{currentfill}%
\pgfsetfillopacity{0.950000}%
\pgfsetlinewidth{0.200750pt}%
\definecolor{currentstroke}{rgb}{0.200000,0.200000,0.200000}%
\pgfsetstrokecolor{currentstroke}%
\pgfsetdash{}{0pt}%
\pgfpathmoveto{\pgfqpoint{2.306063in}{1.229620in}}%
\pgfpathlineto{\pgfqpoint{2.277885in}{1.289016in}}%
\pgfpathlineto{\pgfqpoint{2.249727in}{1.352553in}}%
\pgfpathlineto{\pgfqpoint{2.299120in}{1.370548in}}%
\pgfpathlineto{\pgfqpoint{2.348428in}{1.388469in}}%
\pgfpathlineto{\pgfqpoint{2.376612in}{1.325537in}}%
\pgfpathlineto{\pgfqpoint{2.404826in}{1.266526in}}%
\pgfpathlineto{\pgfqpoint{2.355487in}{1.248130in}}%
\pgfpathlineto{\pgfqpoint{2.306063in}{1.229620in}}%
\pgfpathclose%
\pgfusepath{stroke,fill}%
\end{pgfscope}%
\begin{pgfscope}%
\pgfpathrectangle{\pgfqpoint{0.050000in}{0.083252in}}{\pgfqpoint{4.900000in}{4.900000in}}%
\pgfusepath{clip}%
\pgfsetbuttcap%
\pgfsetroundjoin%
\definecolor{currentfill}{rgb}{0.994033,0.993787,0.996432}%
\pgfsetfillcolor{currentfill}%
\pgfsetfillopacity{0.950000}%
\pgfsetlinewidth{0.200750pt}%
\definecolor{currentstroke}{rgb}{0.200000,0.200000,0.200000}%
\pgfsetstrokecolor{currentstroke}%
\pgfsetdash{}{0pt}%
\pgfpathmoveto{\pgfqpoint{3.280310in}{1.075101in}}%
\pgfpathlineto{\pgfqpoint{3.250806in}{1.105473in}}%
\pgfpathlineto{\pgfqpoint{3.221412in}{1.137062in}}%
\pgfpathlineto{\pgfqpoint{3.269904in}{1.153450in}}%
\pgfpathlineto{\pgfqpoint{3.318308in}{1.169817in}}%
\pgfpathlineto{\pgfqpoint{3.347786in}{1.137861in}}%
\pgfpathlineto{\pgfqpoint{3.377377in}{1.107156in}}%
\pgfpathlineto{\pgfqpoint{3.328888in}{1.091138in}}%
\pgfpathlineto{\pgfqpoint{3.280310in}{1.075101in}}%
\pgfpathclose%
\pgfusepath{stroke,fill}%
\end{pgfscope}%
\begin{pgfscope}%
\pgfpathrectangle{\pgfqpoint{0.050000in}{0.083252in}}{\pgfqpoint{4.900000in}{4.900000in}}%
\pgfusepath{clip}%
\pgfsetbuttcap%
\pgfsetroundjoin%
\definecolor{currentfill}{rgb}{0.928397,0.925444,0.957186}%
\pgfsetfillcolor{currentfill}%
\pgfsetfillopacity{0.950000}%
\pgfsetlinewidth{0.200750pt}%
\definecolor{currentstroke}{rgb}{0.200000,0.200000,0.200000}%
\pgfsetstrokecolor{currentstroke}%
\pgfsetdash{}{0pt}%
\pgfpathmoveto{\pgfqpoint{2.461372in}{1.161690in}}%
\pgfpathlineto{\pgfqpoint{2.433076in}{1.211842in}}%
\pgfpathlineto{\pgfqpoint{2.404826in}{1.266526in}}%
\pgfpathlineto{\pgfqpoint{2.454080in}{1.284818in}}%
\pgfpathlineto{\pgfqpoint{2.503250in}{1.303010in}}%
\pgfpathlineto{\pgfqpoint{2.531548in}{1.248421in}}%
\pgfpathlineto{\pgfqpoint{2.559901in}{1.198107in}}%
\pgfpathlineto{\pgfqpoint{2.510679in}{1.179944in}}%
\pgfpathlineto{\pgfqpoint{2.461372in}{1.161690in}}%
\pgfpathclose%
\pgfusepath{stroke,fill}%
\end{pgfscope}%
\begin{pgfscope}%
\pgfpathrectangle{\pgfqpoint{0.050000in}{0.083252in}}{\pgfqpoint{4.900000in}{4.900000in}}%
\pgfusepath{clip}%
\pgfsetbuttcap%
\pgfsetroundjoin%
\definecolor{currentfill}{rgb}{0.826959,0.819823,0.896532}%
\pgfsetfillcolor{currentfill}%
\pgfsetfillopacity{0.950000}%
\pgfsetlinewidth{0.200750pt}%
\definecolor{currentstroke}{rgb}{0.200000,0.200000,0.200000}%
\pgfsetstrokecolor{currentstroke}%
\pgfsetdash{}{0pt}%
\pgfpathmoveto{\pgfqpoint{2.150685in}{1.316330in}}%
\pgfpathlineto{\pgfqpoint{2.122557in}{1.384237in}}%
\pgfpathlineto{\pgfqpoint{2.094436in}{1.454994in}}%
\pgfpathlineto{\pgfqpoint{2.143991in}{1.472241in}}%
\pgfpathlineto{\pgfqpoint{2.193459in}{1.489471in}}%
\pgfpathlineto{\pgfqpoint{2.221586in}{1.419603in}}%
\pgfpathlineto{\pgfqpoint{2.249727in}{1.352553in}}%
\pgfpathlineto{\pgfqpoint{2.200249in}{1.334480in}}%
\pgfpathlineto{\pgfqpoint{2.150685in}{1.316330in}}%
\pgfpathclose%
\pgfusepath{stroke,fill}%
\end{pgfscope}%
\begin{pgfscope}%
\pgfpathrectangle{\pgfqpoint{0.050000in}{0.083252in}}{\pgfqpoint{4.900000in}{4.900000in}}%
\pgfusepath{clip}%
\pgfsetbuttcap%
\pgfsetroundjoin%
\definecolor{currentfill}{rgb}{0.243768,0.256209,0.609396}%
\pgfsetfillcolor{currentfill}%
\pgfsetfillopacity{0.950000}%
\pgfsetlinewidth{0.200750pt}%
\definecolor{currentstroke}{rgb}{0.200000,0.200000,0.200000}%
\pgfsetstrokecolor{currentstroke}%
\pgfsetdash{}{0pt}%
\pgfpathmoveto{\pgfqpoint{1.146596in}{2.515543in}}%
\pgfpathlineto{\pgfqpoint{1.118866in}{2.591868in}}%
\pgfpathlineto{\pgfqpoint{1.091163in}{2.668120in}}%
\pgfpathlineto{\pgfqpoint{1.141656in}{2.682996in}}%
\pgfpathlineto{\pgfqpoint{1.192057in}{2.697845in}}%
\pgfpathlineto{\pgfqpoint{1.219758in}{2.621746in}}%
\pgfpathlineto{\pgfqpoint{1.247486in}{2.545574in}}%
\pgfpathlineto{\pgfqpoint{1.197087in}{2.530572in}}%
\pgfpathlineto{\pgfqpoint{1.146596in}{2.515543in}}%
\pgfpathclose%
\pgfusepath{stroke,fill}%
\end{pgfscope}%
\begin{pgfscope}%
\pgfpathrectangle{\pgfqpoint{0.050000in}{0.083252in}}{\pgfqpoint{4.900000in}{4.900000in}}%
\pgfusepath{clip}%
\pgfsetbuttcap%
\pgfsetroundjoin%
\definecolor{currentfill}{rgb}{0.283968,0.254441,0.571857}%
\pgfsetfillcolor{currentfill}%
\pgfsetfillopacity{0.950000}%
\pgfsetlinewidth{0.200750pt}%
\definecolor{currentstroke}{rgb}{0.200000,0.200000,0.200000}%
\pgfsetstrokecolor{currentstroke}%
\pgfsetdash{}{0pt}%
\pgfpathmoveto{\pgfqpoint{1.358667in}{2.240151in}}%
\pgfpathlineto{\pgfqpoint{1.330832in}{2.316617in}}%
\pgfpathlineto{\pgfqpoint{1.303023in}{2.393010in}}%
\pgfpathlineto{\pgfqpoint{1.353328in}{2.408138in}}%
\pgfpathlineto{\pgfqpoint{1.403542in}{2.423238in}}%
\pgfpathlineto{\pgfqpoint{1.431348in}{2.346999in}}%
\pgfpathlineto{\pgfqpoint{1.459181in}{2.270687in}}%
\pgfpathlineto{\pgfqpoint{1.408970in}{2.255433in}}%
\pgfpathlineto{\pgfqpoint{1.358667in}{2.240151in}}%
\pgfpathclose%
\pgfusepath{stroke,fill}%
\end{pgfscope}%
\begin{pgfscope}%
\pgfpathrectangle{\pgfqpoint{0.050000in}{0.083252in}}{\pgfqpoint{4.900000in}{4.900000in}}%
\pgfusepath{clip}%
\pgfsetbuttcap%
\pgfsetroundjoin%
\definecolor{currentfill}{rgb}{0.445075,0.422191,0.668189}%
\pgfsetfillcolor{currentfill}%
\pgfsetfillopacity{0.950000}%
\pgfsetlinewidth{0.200750pt}%
\definecolor{currentstroke}{rgb}{0.200000,0.200000,0.200000}%
\pgfsetstrokecolor{currentstroke}%
\pgfsetdash{}{0pt}%
\pgfpathmoveto{\pgfqpoint{1.570783in}{1.964700in}}%
\pgfpathlineto{\pgfqpoint{1.542842in}{2.041307in}}%
\pgfpathlineto{\pgfqpoint{1.514928in}{2.117841in}}%
\pgfpathlineto{\pgfqpoint{1.565045in}{2.133221in}}%
\pgfpathlineto{\pgfqpoint{1.615071in}{2.148573in}}%
\pgfpathlineto{\pgfqpoint{1.642983in}{2.072195in}}%
\pgfpathlineto{\pgfqpoint{1.670921in}{1.995744in}}%
\pgfpathlineto{\pgfqpoint{1.620897in}{1.980235in}}%
\pgfpathlineto{\pgfqpoint{1.570783in}{1.964700in}}%
\pgfpathclose%
\pgfusepath{stroke,fill}%
\end{pgfscope}%
\begin{pgfscope}%
\pgfpathrectangle{\pgfqpoint{0.050000in}{0.083252in}}{\pgfqpoint{4.900000in}{4.900000in}}%
\pgfusepath{clip}%
\pgfsetbuttcap%
\pgfsetroundjoin%
\definecolor{currentfill}{rgb}{0.606182,0.589942,0.764521}%
\pgfsetfillcolor{currentfill}%
\pgfsetfillopacity{0.950000}%
\pgfsetlinewidth{0.200750pt}%
\definecolor{currentstroke}{rgb}{0.200000,0.200000,0.200000}%
\pgfsetstrokecolor{currentstroke}%
\pgfsetdash{}{0pt}%
\pgfpathmoveto{\pgfqpoint{1.782940in}{1.689368in}}%
\pgfpathlineto{\pgfqpoint{1.754896in}{1.766007in}}%
\pgfpathlineto{\pgfqpoint{1.726877in}{1.842638in}}%
\pgfpathlineto{\pgfqpoint{1.776806in}{1.858293in}}%
\pgfpathlineto{\pgfqpoint{1.826644in}{1.873933in}}%
\pgfpathlineto{\pgfqpoint{1.854659in}{1.797525in}}%
\pgfpathlineto{\pgfqpoint{1.882698in}{1.721174in}}%
\pgfpathlineto{\pgfqpoint{1.832865in}{1.705268in}}%
\pgfpathlineto{\pgfqpoint{1.782940in}{1.689368in}}%
\pgfpathclose%
\pgfusepath{stroke,fill}%
\end{pgfscope}%
\begin{pgfscope}%
\pgfpathrectangle{\pgfqpoint{0.050000in}{0.083252in}}{\pgfqpoint{4.900000in}{4.900000in}}%
\pgfusepath{clip}%
\pgfsetbuttcap%
\pgfsetroundjoin%
\definecolor{currentfill}{rgb}{0.958231,0.956509,0.975025}%
\pgfsetfillcolor{currentfill}%
\pgfsetfillopacity{0.950000}%
\pgfsetlinewidth{0.200750pt}%
\definecolor{currentstroke}{rgb}{0.200000,0.200000,0.200000}%
\pgfsetstrokecolor{currentstroke}%
\pgfsetdash{}{0pt}%
\pgfpathmoveto{\pgfqpoint{2.616806in}{1.110031in}}%
\pgfpathlineto{\pgfqpoint{2.588318in}{1.152052in}}%
\pgfpathlineto{\pgfqpoint{2.559901in}{1.198107in}}%
\pgfpathlineto{\pgfqpoint{2.609038in}{1.216183in}}%
\pgfpathlineto{\pgfqpoint{2.658089in}{1.234176in}}%
\pgfpathlineto{\pgfqpoint{2.686571in}{1.187790in}}%
\pgfpathlineto{\pgfqpoint{2.715132in}{1.145277in}}%
\pgfpathlineto{\pgfqpoint{2.666013in}{1.127679in}}%
\pgfpathlineto{\pgfqpoint{2.616806in}{1.110031in}}%
\pgfpathclose%
\pgfusepath{stroke,fill}%
\end{pgfscope}%
\begin{pgfscope}%
\pgfpathrectangle{\pgfqpoint{0.050000in}{0.083252in}}{\pgfqpoint{4.900000in}{4.900000in}}%
\pgfusepath{clip}%
\pgfsetbuttcap%
\pgfsetroundjoin%
\definecolor{currentfill}{rgb}{0.988066,0.987574,0.992864}%
\pgfsetfillcolor{currentfill}%
\pgfsetfillopacity{0.950000}%
\pgfsetlinewidth{0.200750pt}%
\definecolor{currentstroke}{rgb}{0.200000,0.200000,0.200000}%
\pgfsetstrokecolor{currentstroke}%
\pgfsetdash{}{0pt}%
\pgfpathmoveto{\pgfqpoint{3.026561in}{1.071343in}}%
\pgfpathlineto{\pgfqpoint{2.997448in}{1.103750in}}%
\pgfpathlineto{\pgfqpoint{2.968442in}{1.138209in}}%
\pgfpathlineto{\pgfqpoint{3.017203in}{1.155113in}}%
\pgfpathlineto{\pgfqpoint{3.065877in}{1.171990in}}%
\pgfpathlineto{\pgfqpoint{3.094965in}{1.137071in}}%
\pgfpathlineto{\pgfqpoint{3.124163in}{1.104234in}}%
\pgfpathlineto{\pgfqpoint{3.075407in}{1.087795in}}%
\pgfpathlineto{\pgfqpoint{3.026561in}{1.071343in}}%
\pgfpathclose%
\pgfusepath{stroke,fill}%
\end{pgfscope}%
\begin{pgfscope}%
\pgfpathrectangle{\pgfqpoint{0.050000in}{0.083252in}}{\pgfqpoint{4.900000in}{4.900000in}}%
\pgfusepath{clip}%
\pgfsetbuttcap%
\pgfsetroundjoin%
\definecolor{currentfill}{rgb}{0.761323,0.751480,0.857286}%
\pgfsetfillcolor{currentfill}%
\pgfsetfillopacity{0.950000}%
\pgfsetlinewidth{0.200750pt}%
\definecolor{currentstroke}{rgb}{0.200000,0.200000,0.200000}%
\pgfsetstrokecolor{currentstroke}%
\pgfsetdash{}{0pt}%
\pgfpathmoveto{\pgfqpoint{1.995061in}{1.420475in}}%
\pgfpathlineto{\pgfqpoint{1.966948in}{1.494222in}}%
\pgfpathlineto{\pgfqpoint{1.938846in}{1.569239in}}%
\pgfpathlineto{\pgfqpoint{1.988582in}{1.585659in}}%
\pgfpathlineto{\pgfqpoint{2.038229in}{1.602100in}}%
\pgfpathlineto{\pgfqpoint{2.066325in}{1.527833in}}%
\pgfpathlineto{\pgfqpoint{2.094436in}{1.454994in}}%
\pgfpathlineto{\pgfqpoint{2.044793in}{1.437736in}}%
\pgfpathlineto{\pgfqpoint{1.995061in}{1.420475in}}%
\pgfpathclose%
\pgfusepath{stroke,fill}%
\end{pgfscope}%
\begin{pgfscope}%
\pgfpathrectangle{\pgfqpoint{0.050000in}{0.083252in}}{\pgfqpoint{4.900000in}{4.900000in}}%
\pgfusepath{clip}%
\pgfsetbuttcap%
\pgfsetroundjoin%
\definecolor{currentfill}{rgb}{0.240569,0.222222,0.564214}%
\pgfsetfillcolor{currentfill}%
\pgfsetfillopacity{0.950000}%
\pgfsetlinewidth{0.200750pt}%
\definecolor{currentstroke}{rgb}{0.200000,0.200000,0.200000}%
\pgfsetstrokecolor{currentstroke}%
\pgfsetdash{}{0pt}%
\pgfpathmoveto{\pgfqpoint{1.202137in}{2.362671in}}%
\pgfpathlineto{\pgfqpoint{1.174353in}{2.439144in}}%
\pgfpathlineto{\pgfqpoint{1.146596in}{2.515543in}}%
\pgfpathlineto{\pgfqpoint{1.197087in}{2.530572in}}%
\pgfpathlineto{\pgfqpoint{1.247486in}{2.545574in}}%
\pgfpathlineto{\pgfqpoint{1.275241in}{2.469329in}}%
\pgfpathlineto{\pgfqpoint{1.303023in}{2.393010in}}%
\pgfpathlineto{\pgfqpoint{1.252626in}{2.377854in}}%
\pgfpathlineto{\pgfqpoint{1.202137in}{2.362671in}}%
\pgfpathclose%
\pgfusepath{stroke,fill}%
\end{pgfscope}%
\begin{pgfscope}%
\pgfpathrectangle{\pgfqpoint{0.050000in}{0.083252in}}{\pgfqpoint{4.900000in}{4.900000in}}%
\pgfusepath{clip}%
\pgfsetbuttcap%
\pgfsetroundjoin%
\definecolor{currentfill}{rgb}{1.000000,1.000000,1.000000}%
\pgfsetfillcolor{currentfill}%
\pgfsetfillopacity{0.950000}%
\pgfsetlinewidth{0.200750pt}%
\definecolor{currentstroke}{rgb}{0.200000,0.200000,0.200000}%
\pgfsetstrokecolor{currentstroke}%
\pgfsetdash{}{0pt}%
\pgfpathmoveto{\pgfqpoint{3.436874in}{1.047717in}}%
\pgfpathlineto{\pgfqpoint{3.407075in}{1.077259in}}%
\pgfpathlineto{\pgfqpoint{3.377377in}{1.107156in}}%
\pgfpathlineto{\pgfqpoint{3.425777in}{1.123153in}}%
\pgfpathlineto{\pgfqpoint{3.474088in}{1.139128in}}%
\pgfpathlineto{\pgfqpoint{3.503872in}{1.108917in}}%
\pgfpathlineto{\pgfqpoint{3.533757in}{1.079068in}}%
\pgfpathlineto{\pgfqpoint{3.485361in}{1.063407in}}%
\pgfpathlineto{\pgfqpoint{3.436874in}{1.047717in}}%
\pgfpathclose%
\pgfusepath{stroke,fill}%
\end{pgfscope}%
\begin{pgfscope}%
\pgfpathrectangle{\pgfqpoint{0.050000in}{0.083252in}}{\pgfqpoint{4.900000in}{4.900000in}}%
\pgfusepath{clip}%
\pgfsetbuttcap%
\pgfsetroundjoin%
\definecolor{currentfill}{rgb}{0.367505,0.341423,0.621807}%
\pgfsetfillcolor{currentfill}%
\pgfsetfillopacity{0.950000}%
\pgfsetlinewidth{0.200750pt}%
\definecolor{currentstroke}{rgb}{0.200000,0.200000,0.200000}%
\pgfsetstrokecolor{currentstroke}%
\pgfsetdash{}{0pt}%
\pgfpathmoveto{\pgfqpoint{1.414419in}{2.086996in}}%
\pgfpathlineto{\pgfqpoint{1.386529in}{2.163610in}}%
\pgfpathlineto{\pgfqpoint{1.358667in}{2.240151in}}%
\pgfpathlineto{\pgfqpoint{1.408970in}{2.255433in}}%
\pgfpathlineto{\pgfqpoint{1.459181in}{2.270687in}}%
\pgfpathlineto{\pgfqpoint{1.487041in}{2.194301in}}%
\pgfpathlineto{\pgfqpoint{1.514928in}{2.117841in}}%
\pgfpathlineto{\pgfqpoint{1.464719in}{2.102433in}}%
\pgfpathlineto{\pgfqpoint{1.414419in}{2.086996in}}%
\pgfpathclose%
\pgfusepath{stroke,fill}%
\end{pgfscope}%
\begin{pgfscope}%
\pgfpathrectangle{\pgfqpoint{0.050000in}{0.083252in}}{\pgfqpoint{4.900000in}{4.900000in}}%
\pgfusepath{clip}%
\pgfsetbuttcap%
\pgfsetroundjoin%
\definecolor{currentfill}{rgb}{0.528612,0.509173,0.718139}%
\pgfsetfillcolor{currentfill}%
\pgfsetfillopacity{0.950000}%
\pgfsetlinewidth{0.200750pt}%
\definecolor{currentstroke}{rgb}{0.200000,0.200000,0.200000}%
\pgfsetstrokecolor{currentstroke}%
\pgfsetdash{}{0pt}%
\pgfpathmoveto{\pgfqpoint{1.626745in}{1.811266in}}%
\pgfpathlineto{\pgfqpoint{1.598751in}{1.888019in}}%
\pgfpathlineto{\pgfqpoint{1.570783in}{1.964700in}}%
\pgfpathlineto{\pgfqpoint{1.620897in}{1.980235in}}%
\pgfpathlineto{\pgfqpoint{1.670921in}{1.995744in}}%
\pgfpathlineto{\pgfqpoint{1.698886in}{1.919222in}}%
\pgfpathlineto{\pgfqpoint{1.726877in}{1.842638in}}%
\pgfpathlineto{\pgfqpoint{1.676857in}{1.826963in}}%
\pgfpathlineto{\pgfqpoint{1.626745in}{1.811266in}}%
\pgfpathclose%
\pgfusepath{stroke,fill}%
\end{pgfscope}%
\begin{pgfscope}%
\pgfpathrectangle{\pgfqpoint{0.050000in}{0.083252in}}{\pgfqpoint{4.900000in}{4.900000in}}%
\pgfusepath{clip}%
\pgfsetbuttcap%
\pgfsetroundjoin%
\definecolor{currentfill}{rgb}{0.982099,0.981361,0.989296}%
\pgfsetfillcolor{currentfill}%
\pgfsetfillopacity{0.950000}%
\pgfsetlinewidth{0.200750pt}%
\definecolor{currentstroke}{rgb}{0.200000,0.200000,0.200000}%
\pgfsetstrokecolor{currentstroke}%
\pgfsetdash{}{0pt}%
\pgfpathmoveto{\pgfqpoint{2.772519in}{1.070354in}}%
\pgfpathlineto{\pgfqpoint{2.743780in}{1.106289in}}%
\pgfpathlineto{\pgfqpoint{2.715132in}{1.145277in}}%
\pgfpathlineto{\pgfqpoint{2.764164in}{1.162824in}}%
\pgfpathlineto{\pgfqpoint{2.813110in}{1.180321in}}%
\pgfpathlineto{\pgfqpoint{2.841834in}{1.140815in}}%
\pgfpathlineto{\pgfqpoint{2.870656in}{1.104326in}}%
\pgfpathlineto{\pgfqpoint{2.821632in}{1.087350in}}%
\pgfpathlineto{\pgfqpoint{2.772519in}{1.070354in}}%
\pgfpathclose%
\pgfusepath{stroke,fill}%
\end{pgfscope}%
\begin{pgfscope}%
\pgfpathrectangle{\pgfqpoint{0.050000in}{0.083252in}}{\pgfqpoint{4.900000in}{4.900000in}}%
\pgfusepath{clip}%
\pgfsetbuttcap%
\pgfsetroundjoin%
\definecolor{currentfill}{rgb}{0.683752,0.670711,0.810903}%
\pgfsetfillcolor{currentfill}%
\pgfsetfillopacity{0.950000}%
\pgfsetlinewidth{0.200750pt}%
\definecolor{currentstroke}{rgb}{0.200000,0.200000,0.200000}%
\pgfsetstrokecolor{currentstroke}%
\pgfsetdash{}{0pt}%
\pgfpathmoveto{\pgfqpoint{1.839101in}{1.536487in}}%
\pgfpathlineto{\pgfqpoint{1.811009in}{1.612806in}}%
\pgfpathlineto{\pgfqpoint{1.782940in}{1.689368in}}%
\pgfpathlineto{\pgfqpoint{1.832865in}{1.705268in}}%
\pgfpathlineto{\pgfqpoint{1.882698in}{1.721174in}}%
\pgfpathlineto{\pgfqpoint{1.910762in}{1.645006in}}%
\pgfpathlineto{\pgfqpoint{1.938846in}{1.569239in}}%
\pgfpathlineto{\pgfqpoint{1.889019in}{1.552846in}}%
\pgfpathlineto{\pgfqpoint{1.839101in}{1.536487in}}%
\pgfpathclose%
\pgfusepath{stroke,fill}%
\end{pgfscope}%
\begin{pgfscope}%
\pgfpathrectangle{\pgfqpoint{0.050000in}{0.083252in}}{\pgfqpoint{4.900000in}{4.900000in}}%
\pgfusepath{clip}%
\pgfsetbuttcap%
\pgfsetroundjoin%
\definecolor{currentfill}{rgb}{0.892595,0.888166,0.935779}%
\pgfsetfillcolor{currentfill}%
\pgfsetfillopacity{0.950000}%
\pgfsetlinewidth{0.200750pt}%
\definecolor{currentstroke}{rgb}{0.200000,0.200000,0.200000}%
\pgfsetstrokecolor{currentstroke}%
\pgfsetdash{}{0pt}%
\pgfpathmoveto{\pgfqpoint{2.206961in}{1.192230in}}%
\pgfpathlineto{\pgfqpoint{2.178820in}{1.252080in}}%
\pgfpathlineto{\pgfqpoint{2.150685in}{1.316330in}}%
\pgfpathlineto{\pgfqpoint{2.200249in}{1.334480in}}%
\pgfpathlineto{\pgfqpoint{2.249727in}{1.352553in}}%
\pgfpathlineto{\pgfqpoint{2.277885in}{1.289016in}}%
\pgfpathlineto{\pgfqpoint{2.306063in}{1.229620in}}%
\pgfpathlineto{\pgfqpoint{2.256554in}{1.210990in}}%
\pgfpathlineto{\pgfqpoint{2.206961in}{1.192230in}}%
\pgfpathclose%
\pgfusepath{stroke,fill}%
\end{pgfscope}%
\begin{pgfscope}%
\pgfpathrectangle{\pgfqpoint{0.050000in}{0.083252in}}{\pgfqpoint{4.900000in}{4.900000in}}%
\pgfusepath{clip}%
\pgfsetbuttcap%
\pgfsetroundjoin%
\definecolor{currentfill}{rgb}{0.934364,0.931657,0.960754}%
\pgfsetfillcolor{currentfill}%
\pgfsetfillopacity{0.950000}%
\pgfsetlinewidth{0.200750pt}%
\definecolor{currentstroke}{rgb}{0.200000,0.200000,0.200000}%
\pgfsetstrokecolor{currentstroke}%
\pgfsetdash{}{0pt}%
\pgfpathmoveto{\pgfqpoint{2.362500in}{1.124886in}}%
\pgfpathlineto{\pgfqpoint{2.334265in}{1.174835in}}%
\pgfpathlineto{\pgfqpoint{2.306063in}{1.229620in}}%
\pgfpathlineto{\pgfqpoint{2.355487in}{1.248130in}}%
\pgfpathlineto{\pgfqpoint{2.404826in}{1.266526in}}%
\pgfpathlineto{\pgfqpoint{2.433076in}{1.211842in}}%
\pgfpathlineto{\pgfqpoint{2.461372in}{1.161690in}}%
\pgfpathlineto{\pgfqpoint{2.411978in}{1.143339in}}%
\pgfpathlineto{\pgfqpoint{2.362500in}{1.124886in}}%
\pgfpathclose%
\pgfusepath{stroke,fill}%
\end{pgfscope}%
\begin{pgfscope}%
\pgfpathrectangle{\pgfqpoint{0.050000in}{0.083252in}}{\pgfqpoint{4.900000in}{4.900000in}}%
\pgfusepath{clip}%
\pgfsetbuttcap%
\pgfsetroundjoin%
\definecolor{currentfill}{rgb}{0.994033,0.993787,0.996432}%
\pgfsetfillcolor{currentfill}%
\pgfsetfillopacity{0.950000}%
\pgfsetlinewidth{0.200750pt}%
\definecolor{currentstroke}{rgb}{0.200000,0.200000,0.200000}%
\pgfsetstrokecolor{currentstroke}%
\pgfsetdash{}{0pt}%
\pgfpathmoveto{\pgfqpoint{3.182885in}{1.042969in}}%
\pgfpathlineto{\pgfqpoint{3.153471in}{1.073017in}}%
\pgfpathlineto{\pgfqpoint{3.124163in}{1.104234in}}%
\pgfpathlineto{\pgfqpoint{3.172832in}{1.120657in}}%
\pgfpathlineto{\pgfqpoint{3.221412in}{1.137062in}}%
\pgfpathlineto{\pgfqpoint{3.250806in}{1.105473in}}%
\pgfpathlineto{\pgfqpoint{3.280310in}{1.075101in}}%
\pgfpathlineto{\pgfqpoint{3.231642in}{1.059044in}}%
\pgfpathlineto{\pgfqpoint{3.182885in}{1.042969in}}%
\pgfpathclose%
\pgfusepath{stroke,fill}%
\end{pgfscope}%
\begin{pgfscope}%
\pgfpathrectangle{\pgfqpoint{0.050000in}{0.083252in}}{\pgfqpoint{4.900000in}{4.900000in}}%
\pgfusepath{clip}%
\pgfsetbuttcap%
\pgfsetroundjoin%
\definecolor{currentfill}{rgb}{0.832926,0.826036,0.900100}%
\pgfsetfillcolor{currentfill}%
\pgfsetfillopacity{0.950000}%
\pgfsetlinewidth{0.200750pt}%
\definecolor{currentstroke}{rgb}{0.200000,0.200000,0.200000}%
\pgfsetstrokecolor{currentstroke}%
\pgfsetdash{}{0pt}%
\pgfpathmoveto{\pgfqpoint{2.051301in}{1.279788in}}%
\pgfpathlineto{\pgfqpoint{2.023181in}{1.348707in}}%
\pgfpathlineto{\pgfqpoint{1.995061in}{1.420475in}}%
\pgfpathlineto{\pgfqpoint{2.044793in}{1.437736in}}%
\pgfpathlineto{\pgfqpoint{2.094436in}{1.454994in}}%
\pgfpathlineto{\pgfqpoint{2.122557in}{1.384237in}}%
\pgfpathlineto{\pgfqpoint{2.150685in}{1.316330in}}%
\pgfpathlineto{\pgfqpoint{2.101036in}{1.298100in}}%
\pgfpathlineto{\pgfqpoint{2.051301in}{1.279788in}}%
\pgfpathclose%
\pgfusepath{stroke,fill}%
\end{pgfscope}%
\begin{pgfscope}%
\pgfpathrectangle{\pgfqpoint{0.050000in}{0.083252in}}{\pgfqpoint{4.900000in}{4.900000in}}%
\pgfusepath{clip}%
\pgfsetbuttcap%
\pgfsetroundjoin%
\definecolor{currentfill}{rgb}{0.289935,0.260654,0.575425}%
\pgfsetfillcolor{currentfill}%
\pgfsetfillopacity{0.950000}%
\pgfsetlinewidth{0.200750pt}%
\definecolor{currentstroke}{rgb}{0.200000,0.200000,0.200000}%
\pgfsetstrokecolor{currentstroke}%
\pgfsetdash{}{0pt}%
\pgfpathmoveto{\pgfqpoint{1.257785in}{2.209502in}}%
\pgfpathlineto{\pgfqpoint{1.229947in}{2.286123in}}%
\pgfpathlineto{\pgfqpoint{1.202137in}{2.362671in}}%
\pgfpathlineto{\pgfqpoint{1.252626in}{2.377854in}}%
\pgfpathlineto{\pgfqpoint{1.303023in}{2.393010in}}%
\pgfpathlineto{\pgfqpoint{1.330832in}{2.316617in}}%
\pgfpathlineto{\pgfqpoint{1.358667in}{2.240151in}}%
\pgfpathlineto{\pgfqpoint{1.308272in}{2.224841in}}%
\pgfpathlineto{\pgfqpoint{1.257785in}{2.209502in}}%
\pgfpathclose%
\pgfusepath{stroke,fill}%
\end{pgfscope}%
\begin{pgfscope}%
\pgfpathrectangle{\pgfqpoint{0.050000in}{0.083252in}}{\pgfqpoint{4.900000in}{4.900000in}}%
\pgfusepath{clip}%
\pgfsetbuttcap%
\pgfsetroundjoin%
\definecolor{currentfill}{rgb}{0.445075,0.422191,0.668189}%
\pgfsetfillcolor{currentfill}%
\pgfsetfillopacity{0.950000}%
\pgfsetlinewidth{0.200750pt}%
\definecolor{currentstroke}{rgb}{0.200000,0.200000,0.200000}%
\pgfsetstrokecolor{currentstroke}%
\pgfsetdash{}{0pt}%
\pgfpathmoveto{\pgfqpoint{1.470278in}{1.933545in}}%
\pgfpathlineto{\pgfqpoint{1.442335in}{2.010307in}}%
\pgfpathlineto{\pgfqpoint{1.414419in}{2.086996in}}%
\pgfpathlineto{\pgfqpoint{1.464719in}{2.102433in}}%
\pgfpathlineto{\pgfqpoint{1.514928in}{2.117841in}}%
\pgfpathlineto{\pgfqpoint{1.542842in}{2.041307in}}%
\pgfpathlineto{\pgfqpoint{1.570783in}{1.964700in}}%
\pgfpathlineto{\pgfqpoint{1.520576in}{1.949137in}}%
\pgfpathlineto{\pgfqpoint{1.470278in}{1.933545in}}%
\pgfpathclose%
\pgfusepath{stroke,fill}%
\end{pgfscope}%
\begin{pgfscope}%
\pgfpathrectangle{\pgfqpoint{0.050000in}{0.083252in}}{\pgfqpoint{4.900000in}{4.900000in}}%
\pgfusepath{clip}%
\pgfsetbuttcap%
\pgfsetroundjoin%
\definecolor{currentfill}{rgb}{0.964198,0.962722,0.978593}%
\pgfsetfillcolor{currentfill}%
\pgfsetfillopacity{0.950000}%
\pgfsetlinewidth{0.200750pt}%
\definecolor{currentstroke}{rgb}{0.200000,0.200000,0.200000}%
\pgfsetstrokecolor{currentstroke}%
\pgfsetdash{}{0pt}%
\pgfpathmoveto{\pgfqpoint{2.518132in}{1.074584in}}%
\pgfpathlineto{\pgfqpoint{2.489720in}{1.116031in}}%
\pgfpathlineto{\pgfqpoint{2.461372in}{1.161690in}}%
\pgfpathlineto{\pgfqpoint{2.510679in}{1.179944in}}%
\pgfpathlineto{\pgfqpoint{2.559901in}{1.198107in}}%
\pgfpathlineto{\pgfqpoint{2.588318in}{1.152052in}}%
\pgfpathlineto{\pgfqpoint{2.616806in}{1.110031in}}%
\pgfpathlineto{\pgfqpoint{2.567513in}{1.092333in}}%
\pgfpathlineto{\pgfqpoint{2.518132in}{1.074584in}}%
\pgfpathclose%
\pgfusepath{stroke,fill}%
\end{pgfscope}%
\begin{pgfscope}%
\pgfpathrectangle{\pgfqpoint{0.050000in}{0.083252in}}{\pgfqpoint{4.900000in}{4.900000in}}%
\pgfusepath{clip}%
\pgfsetbuttcap%
\pgfsetroundjoin%
\definecolor{currentfill}{rgb}{0.606182,0.589942,0.764521}%
\pgfsetfillcolor{currentfill}%
\pgfsetfillopacity{0.950000}%
\pgfsetlinewidth{0.200750pt}%
\definecolor{currentstroke}{rgb}{0.200000,0.200000,0.200000}%
\pgfsetstrokecolor{currentstroke}%
\pgfsetdash{}{0pt}%
\pgfpathmoveto{\pgfqpoint{1.682816in}{1.657574in}}%
\pgfpathlineto{\pgfqpoint{1.654767in}{1.734446in}}%
\pgfpathlineto{\pgfqpoint{1.626745in}{1.811266in}}%
\pgfpathlineto{\pgfqpoint{1.676857in}{1.826963in}}%
\pgfpathlineto{\pgfqpoint{1.726877in}{1.842638in}}%
\pgfpathlineto{\pgfqpoint{1.754896in}{1.766007in}}%
\pgfpathlineto{\pgfqpoint{1.782940in}{1.689368in}}%
\pgfpathlineto{\pgfqpoint{1.732924in}{1.673471in}}%
\pgfpathlineto{\pgfqpoint{1.682816in}{1.657574in}}%
\pgfpathclose%
\pgfusepath{stroke,fill}%
\end{pgfscope}%
\begin{pgfscope}%
\pgfpathrectangle{\pgfqpoint{0.050000in}{0.083252in}}{\pgfqpoint{4.900000in}{4.900000in}}%
\pgfusepath{clip}%
\pgfsetbuttcap%
\pgfsetroundjoin%
\definecolor{currentfill}{rgb}{0.994033,0.993787,0.996432}%
\pgfsetfillcolor{currentfill}%
\pgfsetfillopacity{0.950000}%
\pgfsetlinewidth{0.200750pt}%
\definecolor{currentstroke}{rgb}{0.200000,0.200000,0.200000}%
\pgfsetstrokecolor{currentstroke}%
\pgfsetdash{}{0pt}%
\pgfpathmoveto{\pgfqpoint{2.928602in}{1.038403in}}%
\pgfpathlineto{\pgfqpoint{2.899579in}{1.070362in}}%
\pgfpathlineto{\pgfqpoint{2.870656in}{1.104326in}}%
\pgfpathlineto{\pgfqpoint{2.919593in}{1.121279in}}%
\pgfpathlineto{\pgfqpoint{2.968442in}{1.138209in}}%
\pgfpathlineto{\pgfqpoint{2.997448in}{1.103750in}}%
\pgfpathlineto{\pgfqpoint{3.026561in}{1.071343in}}%
\pgfpathlineto{\pgfqpoint{2.977626in}{1.054878in}}%
\pgfpathlineto{\pgfqpoint{2.928602in}{1.038403in}}%
\pgfpathclose%
\pgfusepath{stroke,fill}%
\end{pgfscope}%
\begin{pgfscope}%
\pgfpathrectangle{\pgfqpoint{0.050000in}{0.083252in}}{\pgfqpoint{4.900000in}{4.900000in}}%
\pgfusepath{clip}%
\pgfsetbuttcap%
\pgfsetroundjoin%
\definecolor{currentfill}{rgb}{0.761323,0.751480,0.857286}%
\pgfsetfillcolor{currentfill}%
\pgfsetfillopacity{0.950000}%
\pgfsetlinewidth{0.200750pt}%
\definecolor{currentstroke}{rgb}{0.200000,0.200000,0.200000}%
\pgfsetstrokecolor{currentstroke}%
\pgfsetdash{}{0pt}%
\pgfpathmoveto{\pgfqpoint{1.895332in}{1.385992in}}%
\pgfpathlineto{\pgfqpoint{1.867210in}{1.460716in}}%
\pgfpathlineto{\pgfqpoint{1.839101in}{1.536487in}}%
\pgfpathlineto{\pgfqpoint{1.889019in}{1.552846in}}%
\pgfpathlineto{\pgfqpoint{1.938846in}{1.569239in}}%
\pgfpathlineto{\pgfqpoint{1.966948in}{1.494222in}}%
\pgfpathlineto{\pgfqpoint{1.995061in}{1.420475in}}%
\pgfpathlineto{\pgfqpoint{1.945241in}{1.403223in}}%
\pgfpathlineto{\pgfqpoint{1.895332in}{1.385992in}}%
\pgfpathclose%
\pgfusepath{stroke,fill}%
\end{pgfscope}%
\begin{pgfscope}%
\pgfpathrectangle{\pgfqpoint{0.050000in}{0.083252in}}{\pgfqpoint{4.900000in}{4.900000in}}%
\pgfusepath{clip}%
\pgfsetbuttcap%
\pgfsetroundjoin%
\definecolor{currentfill}{rgb}{1.000000,1.000000,1.000000}%
\pgfsetfillcolor{currentfill}%
\pgfsetfillopacity{0.950000}%
\pgfsetlinewidth{0.200750pt}%
\definecolor{currentstroke}{rgb}{0.200000,0.200000,0.200000}%
\pgfsetstrokecolor{currentstroke}%
\pgfsetdash{}{0pt}%
\pgfpathmoveto{\pgfqpoint{3.339628in}{1.016250in}}%
\pgfpathlineto{\pgfqpoint{3.309919in}{1.045505in}}%
\pgfpathlineto{\pgfqpoint{3.280310in}{1.075101in}}%
\pgfpathlineto{\pgfqpoint{3.328888in}{1.091138in}}%
\pgfpathlineto{\pgfqpoint{3.377377in}{1.107156in}}%
\pgfpathlineto{\pgfqpoint{3.407075in}{1.077259in}}%
\pgfpathlineto{\pgfqpoint{3.436874in}{1.047717in}}%
\pgfpathlineto{\pgfqpoint{3.388296in}{1.031998in}}%
\pgfpathlineto{\pgfqpoint{3.339628in}{1.016250in}}%
\pgfpathclose%
\pgfusepath{stroke,fill}%
\end{pgfscope}%
\begin{pgfscope}%
\pgfpathrectangle{\pgfqpoint{0.050000in}{0.083252in}}{\pgfqpoint{4.900000in}{4.900000in}}%
\pgfusepath{clip}%
\pgfsetbuttcap%
\pgfsetroundjoin%
\definecolor{currentfill}{rgb}{0.367505,0.341423,0.621807}%
\pgfsetfillcolor{currentfill}%
\pgfsetfillopacity{0.950000}%
\pgfsetlinewidth{0.200750pt}%
\definecolor{currentstroke}{rgb}{0.200000,0.200000,0.200000}%
\pgfsetstrokecolor{currentstroke}%
\pgfsetdash{}{0pt}%
\pgfpathmoveto{\pgfqpoint{1.313541in}{2.056037in}}%
\pgfpathlineto{\pgfqpoint{1.285650in}{2.132807in}}%
\pgfpathlineto{\pgfqpoint{1.257785in}{2.209502in}}%
\pgfpathlineto{\pgfqpoint{1.308272in}{2.224841in}}%
\pgfpathlineto{\pgfqpoint{1.358667in}{2.240151in}}%
\pgfpathlineto{\pgfqpoint{1.386529in}{2.163610in}}%
\pgfpathlineto{\pgfqpoint{1.414419in}{2.086996in}}%
\pgfpathlineto{\pgfqpoint{1.364026in}{2.071531in}}%
\pgfpathlineto{\pgfqpoint{1.313541in}{2.056037in}}%
\pgfpathclose%
\pgfusepath{stroke,fill}%
\end{pgfscope}%
\begin{pgfscope}%
\pgfpathrectangle{\pgfqpoint{0.050000in}{0.083252in}}{\pgfqpoint{4.900000in}{4.900000in}}%
\pgfusepath{clip}%
\pgfsetbuttcap%
\pgfsetroundjoin%
\definecolor{currentfill}{rgb}{0.982099,0.981361,0.989296}%
\pgfsetfillcolor{currentfill}%
\pgfsetfillopacity{0.950000}%
\pgfsetlinewidth{0.200750pt}%
\definecolor{currentstroke}{rgb}{0.200000,0.200000,0.200000}%
\pgfsetstrokecolor{currentstroke}%
\pgfsetdash{}{0pt}%
\pgfpathmoveto{\pgfqpoint{2.674027in}{1.036310in}}%
\pgfpathlineto{\pgfqpoint{2.645374in}{1.071632in}}%
\pgfpathlineto{\pgfqpoint{2.616806in}{1.110031in}}%
\pgfpathlineto{\pgfqpoint{2.666013in}{1.127679in}}%
\pgfpathlineto{\pgfqpoint{2.715132in}{1.145277in}}%
\pgfpathlineto{\pgfqpoint{2.743780in}{1.106289in}}%
\pgfpathlineto{\pgfqpoint{2.772519in}{1.070354in}}%
\pgfpathlineto{\pgfqpoint{2.723317in}{1.053339in}}%
\pgfpathlineto{\pgfqpoint{2.674027in}{1.036310in}}%
\pgfpathclose%
\pgfusepath{stroke,fill}%
\end{pgfscope}%
\begin{pgfscope}%
\pgfpathrectangle{\pgfqpoint{0.050000in}{0.083252in}}{\pgfqpoint{4.900000in}{4.900000in}}%
\pgfusepath{clip}%
\pgfsetbuttcap%
\pgfsetroundjoin%
\definecolor{currentfill}{rgb}{0.528612,0.509173,0.718139}%
\pgfsetfillcolor{currentfill}%
\pgfsetfillopacity{0.950000}%
\pgfsetlinewidth{0.200750pt}%
\definecolor{currentstroke}{rgb}{0.200000,0.200000,0.200000}%
\pgfsetstrokecolor{currentstroke}%
\pgfsetdash{}{0pt}%
\pgfpathmoveto{\pgfqpoint{1.526246in}{1.779796in}}%
\pgfpathlineto{\pgfqpoint{1.498249in}{1.856707in}}%
\pgfpathlineto{\pgfqpoint{1.470278in}{1.933545in}}%
\pgfpathlineto{\pgfqpoint{1.520576in}{1.949137in}}%
\pgfpathlineto{\pgfqpoint{1.570783in}{1.964700in}}%
\pgfpathlineto{\pgfqpoint{1.598751in}{1.888019in}}%
\pgfpathlineto{\pgfqpoint{1.626745in}{1.811266in}}%
\pgfpathlineto{\pgfqpoint{1.576542in}{1.795544in}}%
\pgfpathlineto{\pgfqpoint{1.526246in}{1.779796in}}%
\pgfpathclose%
\pgfusepath{stroke,fill}%
\end{pgfscope}%
\begin{pgfscope}%
\pgfpathrectangle{\pgfqpoint{0.050000in}{0.083252in}}{\pgfqpoint{4.900000in}{4.900000in}}%
\pgfusepath{clip}%
\pgfsetbuttcap%
\pgfsetroundjoin%
\definecolor{currentfill}{rgb}{0.683752,0.670711,0.810903}%
\pgfsetfillcolor{currentfill}%
\pgfsetfillopacity{0.950000}%
\pgfsetlinewidth{0.200750pt}%
\definecolor{currentstroke}{rgb}{0.200000,0.200000,0.200000}%
\pgfsetstrokecolor{currentstroke}%
\pgfsetdash{}{0pt}%
\pgfpathmoveto{\pgfqpoint{1.738989in}{1.503895in}}%
\pgfpathlineto{\pgfqpoint{1.710890in}{1.580690in}}%
\pgfpathlineto{\pgfqpoint{1.682816in}{1.657574in}}%
\pgfpathlineto{\pgfqpoint{1.732924in}{1.673471in}}%
\pgfpathlineto{\pgfqpoint{1.782940in}{1.689368in}}%
\pgfpathlineto{\pgfqpoint{1.811009in}{1.612806in}}%
\pgfpathlineto{\pgfqpoint{1.839101in}{1.536487in}}%
\pgfpathlineto{\pgfqpoint{1.789091in}{1.520169in}}%
\pgfpathlineto{\pgfqpoint{1.738989in}{1.503895in}}%
\pgfpathclose%
\pgfusepath{stroke,fill}%
\end{pgfscope}%
\begin{pgfscope}%
\pgfpathrectangle{\pgfqpoint{0.050000in}{0.083252in}}{\pgfqpoint{4.900000in}{4.900000in}}%
\pgfusepath{clip}%
\pgfsetbuttcap%
\pgfsetroundjoin%
\definecolor{currentfill}{rgb}{0.898562,0.894379,0.939346}%
\pgfsetfillcolor{currentfill}%
\pgfsetfillopacity{0.950000}%
\pgfsetlinewidth{0.200750pt}%
\definecolor{currentstroke}{rgb}{0.200000,0.200000,0.200000}%
\pgfsetstrokecolor{currentstroke}%
\pgfsetdash{}{0pt}%
\pgfpathmoveto{\pgfqpoint{2.107523in}{1.154271in}}%
\pgfpathlineto{\pgfqpoint{2.079415in}{1.214676in}}%
\pgfpathlineto{\pgfqpoint{2.051301in}{1.279788in}}%
\pgfpathlineto{\pgfqpoint{2.101036in}{1.298100in}}%
\pgfpathlineto{\pgfqpoint{2.150685in}{1.316330in}}%
\pgfpathlineto{\pgfqpoint{2.178820in}{1.252080in}}%
\pgfpathlineto{\pgfqpoint{2.206961in}{1.192230in}}%
\pgfpathlineto{\pgfqpoint{2.157284in}{1.173329in}}%
\pgfpathlineto{\pgfqpoint{2.107523in}{1.154271in}}%
\pgfpathclose%
\pgfusepath{stroke,fill}%
\end{pgfscope}%
\begin{pgfscope}%
\pgfpathrectangle{\pgfqpoint{0.050000in}{0.083252in}}{\pgfqpoint{4.900000in}{4.900000in}}%
\pgfusepath{clip}%
\pgfsetbuttcap%
\pgfsetroundjoin%
\definecolor{currentfill}{rgb}{0.940331,0.937870,0.964321}%
\pgfsetfillcolor{currentfill}%
\pgfsetfillopacity{0.950000}%
\pgfsetlinewidth{0.200750pt}%
\definecolor{currentstroke}{rgb}{0.200000,0.200000,0.200000}%
\pgfsetstrokecolor{currentstroke}%
\pgfsetdash{}{0pt}%
\pgfpathmoveto{\pgfqpoint{2.263286in}{1.087646in}}%
\pgfpathlineto{\pgfqpoint{2.235114in}{1.137335in}}%
\pgfpathlineto{\pgfqpoint{2.206961in}{1.192230in}}%
\pgfpathlineto{\pgfqpoint{2.256554in}{1.210990in}}%
\pgfpathlineto{\pgfqpoint{2.306063in}{1.229620in}}%
\pgfpathlineto{\pgfqpoint{2.334265in}{1.174835in}}%
\pgfpathlineto{\pgfqpoint{2.362500in}{1.124886in}}%
\pgfpathlineto{\pgfqpoint{2.312936in}{1.106324in}}%
\pgfpathlineto{\pgfqpoint{2.263286in}{1.087646in}}%
\pgfpathclose%
\pgfusepath{stroke,fill}%
\end{pgfscope}%
\begin{pgfscope}%
\pgfpathrectangle{\pgfqpoint{0.050000in}{0.083252in}}{\pgfqpoint{4.900000in}{4.900000in}}%
\pgfusepath{clip}%
\pgfsetbuttcap%
\pgfsetroundjoin%
\definecolor{currentfill}{rgb}{1.000000,1.000000,1.000000}%
\pgfsetfillcolor{currentfill}%
\pgfsetfillopacity{0.950000}%
\pgfsetlinewidth{0.200750pt}%
\definecolor{currentstroke}{rgb}{0.200000,0.200000,0.200000}%
\pgfsetstrokecolor{currentstroke}%
\pgfsetdash{}{0pt}%
\pgfpathmoveto{\pgfqpoint{3.085101in}{1.010769in}}%
\pgfpathlineto{\pgfqpoint{3.055779in}{1.040502in}}%
\pgfpathlineto{\pgfqpoint{3.026561in}{1.071343in}}%
\pgfpathlineto{\pgfqpoint{3.075407in}{1.087795in}}%
\pgfpathlineto{\pgfqpoint{3.124163in}{1.104234in}}%
\pgfpathlineto{\pgfqpoint{3.153471in}{1.073017in}}%
\pgfpathlineto{\pgfqpoint{3.182885in}{1.042969in}}%
\pgfpathlineto{\pgfqpoint{3.134038in}{1.026877in}}%
\pgfpathlineto{\pgfqpoint{3.085101in}{1.010769in}}%
\pgfpathclose%
\pgfusepath{stroke,fill}%
\end{pgfscope}%
\begin{pgfscope}%
\pgfpathrectangle{\pgfqpoint{0.050000in}{0.083252in}}{\pgfqpoint{4.900000in}{4.900000in}}%
\pgfusepath{clip}%
\pgfsetbuttcap%
\pgfsetroundjoin%
\definecolor{currentfill}{rgb}{0.838893,0.832249,0.903668}%
\pgfsetfillcolor{currentfill}%
\pgfsetfillopacity{0.950000}%
\pgfsetlinewidth{0.200750pt}%
\definecolor{currentstroke}{rgb}{0.200000,0.200000,0.200000}%
\pgfsetstrokecolor{currentstroke}%
\pgfsetdash{}{0pt}%
\pgfpathmoveto{\pgfqpoint{1.951574in}{1.242925in}}%
\pgfpathlineto{\pgfqpoint{1.923456in}{1.313064in}}%
\pgfpathlineto{\pgfqpoint{1.895332in}{1.385992in}}%
\pgfpathlineto{\pgfqpoint{1.945241in}{1.403223in}}%
\pgfpathlineto{\pgfqpoint{1.995061in}{1.420475in}}%
\pgfpathlineto{\pgfqpoint{2.023181in}{1.348707in}}%
\pgfpathlineto{\pgfqpoint{2.051301in}{1.279788in}}%
\pgfpathlineto{\pgfqpoint{2.001480in}{1.261397in}}%
\pgfpathlineto{\pgfqpoint{1.951574in}{1.242925in}}%
\pgfpathclose%
\pgfusepath{stroke,fill}%
\end{pgfscope}%
\begin{pgfscope}%
\pgfpathrectangle{\pgfqpoint{0.050000in}{0.083252in}}{\pgfqpoint{4.900000in}{4.900000in}}%
\pgfusepath{clip}%
\pgfsetbuttcap%
\pgfsetroundjoin%
\definecolor{currentfill}{rgb}{0.445075,0.422191,0.668189}%
\pgfsetfillcolor{currentfill}%
\pgfsetfillopacity{0.950000}%
\pgfsetlinewidth{0.200750pt}%
\definecolor{currentstroke}{rgb}{0.200000,0.200000,0.200000}%
\pgfsetstrokecolor{currentstroke}%
\pgfsetdash{}{0pt}%
\pgfpathmoveto{\pgfqpoint{1.369405in}{1.902274in}}%
\pgfpathlineto{\pgfqpoint{1.341460in}{1.979193in}}%
\pgfpathlineto{\pgfqpoint{1.313541in}{2.056037in}}%
\pgfpathlineto{\pgfqpoint{1.364026in}{2.071531in}}%
\pgfpathlineto{\pgfqpoint{1.414419in}{2.086996in}}%
\pgfpathlineto{\pgfqpoint{1.442335in}{2.010307in}}%
\pgfpathlineto{\pgfqpoint{1.470278in}{1.933545in}}%
\pgfpathlineto{\pgfqpoint{1.419888in}{1.917924in}}%
\pgfpathlineto{\pgfqpoint{1.369405in}{1.902274in}}%
\pgfpathclose%
\pgfusepath{stroke,fill}%
\end{pgfscope}%
\begin{pgfscope}%
\pgfpathrectangle{\pgfqpoint{0.050000in}{0.083252in}}{\pgfqpoint{4.900000in}{4.900000in}}%
\pgfusepath{clip}%
\pgfsetbuttcap%
\pgfsetroundjoin%
\definecolor{currentfill}{rgb}{0.970165,0.968935,0.982161}%
\pgfsetfillcolor{currentfill}%
\pgfsetfillopacity{0.950000}%
\pgfsetlinewidth{0.200750pt}%
\definecolor{currentstroke}{rgb}{0.200000,0.200000,0.200000}%
\pgfsetstrokecolor{currentstroke}%
\pgfsetdash{}{0pt}%
\pgfpathmoveto{\pgfqpoint{2.419108in}{1.038938in}}%
\pgfpathlineto{\pgfqpoint{2.390777in}{1.079706in}}%
\pgfpathlineto{\pgfqpoint{2.362500in}{1.124886in}}%
\pgfpathlineto{\pgfqpoint{2.411978in}{1.143339in}}%
\pgfpathlineto{\pgfqpoint{2.461372in}{1.161690in}}%
\pgfpathlineto{\pgfqpoint{2.489720in}{1.116031in}}%
\pgfpathlineto{\pgfqpoint{2.518132in}{1.074584in}}%
\pgfpathlineto{\pgfqpoint{2.468664in}{1.056785in}}%
\pgfpathlineto{\pgfqpoint{2.419108in}{1.038938in}}%
\pgfpathclose%
\pgfusepath{stroke,fill}%
\end{pgfscope}%
\begin{pgfscope}%
\pgfpathrectangle{\pgfqpoint{0.050000in}{0.083252in}}{\pgfqpoint{4.900000in}{4.900000in}}%
\pgfusepath{clip}%
\pgfsetbuttcap%
\pgfsetroundjoin%
\definecolor{currentfill}{rgb}{0.606182,0.589942,0.764521}%
\pgfsetfillcolor{currentfill}%
\pgfsetfillopacity{0.950000}%
\pgfsetlinewidth{0.200750pt}%
\definecolor{currentstroke}{rgb}{0.200000,0.200000,0.200000}%
\pgfsetstrokecolor{currentstroke}%
\pgfsetdash{}{0pt}%
\pgfpathmoveto{\pgfqpoint{1.582322in}{1.625753in}}%
\pgfpathlineto{\pgfqpoint{1.554271in}{1.702810in}}%
\pgfpathlineto{\pgfqpoint{1.526246in}{1.779796in}}%
\pgfpathlineto{\pgfqpoint{1.576542in}{1.795544in}}%
\pgfpathlineto{\pgfqpoint{1.626745in}{1.811266in}}%
\pgfpathlineto{\pgfqpoint{1.654767in}{1.734446in}}%
\pgfpathlineto{\pgfqpoint{1.682816in}{1.657574in}}%
\pgfpathlineto{\pgfqpoint{1.632615in}{1.641670in}}%
\pgfpathlineto{\pgfqpoint{1.582322in}{1.625753in}}%
\pgfpathclose%
\pgfusepath{stroke,fill}%
\end{pgfscope}%
\begin{pgfscope}%
\pgfpathrectangle{\pgfqpoint{0.050000in}{0.083252in}}{\pgfqpoint{4.900000in}{4.900000in}}%
\pgfusepath{clip}%
\pgfsetbuttcap%
\pgfsetroundjoin%
\definecolor{currentfill}{rgb}{0.994033,0.993787,0.996432}%
\pgfsetfillcolor{currentfill}%
\pgfsetfillopacity{0.950000}%
\pgfsetlinewidth{0.200750pt}%
\definecolor{currentstroke}{rgb}{0.200000,0.200000,0.200000}%
\pgfsetstrokecolor{currentstroke}%
\pgfsetdash{}{0pt}%
\pgfpathmoveto{\pgfqpoint{2.830285in}{1.005431in}}%
\pgfpathlineto{\pgfqpoint{2.801353in}{1.036923in}}%
\pgfpathlineto{\pgfqpoint{2.772519in}{1.070354in}}%
\pgfpathlineto{\pgfqpoint{2.821632in}{1.087350in}}%
\pgfpathlineto{\pgfqpoint{2.870656in}{1.104326in}}%
\pgfpathlineto{\pgfqpoint{2.899579in}{1.070362in}}%
\pgfpathlineto{\pgfqpoint{2.928602in}{1.038403in}}%
\pgfpathlineto{\pgfqpoint{2.879489in}{1.021920in}}%
\pgfpathlineto{\pgfqpoint{2.830285in}{1.005431in}}%
\pgfpathclose%
\pgfusepath{stroke,fill}%
\end{pgfscope}%
\begin{pgfscope}%
\pgfpathrectangle{\pgfqpoint{0.050000in}{0.083252in}}{\pgfqpoint{4.900000in}{4.900000in}}%
\pgfusepath{clip}%
\pgfsetbuttcap%
\pgfsetroundjoin%
\definecolor{currentfill}{rgb}{0.767290,0.757693,0.860854}%
\pgfsetfillcolor{currentfill}%
\pgfsetfillopacity{0.950000}%
\pgfsetlinewidth{0.200750pt}%
\definecolor{currentstroke}{rgb}{0.200000,0.200000,0.200000}%
\pgfsetstrokecolor{currentstroke}%
\pgfsetdash{}{0pt}%
\pgfpathmoveto{\pgfqpoint{1.795242in}{1.351668in}}%
\pgfpathlineto{\pgfqpoint{1.767108in}{1.427411in}}%
\pgfpathlineto{\pgfqpoint{1.738989in}{1.503895in}}%
\pgfpathlineto{\pgfqpoint{1.789091in}{1.520169in}}%
\pgfpathlineto{\pgfqpoint{1.839101in}{1.536487in}}%
\pgfpathlineto{\pgfqpoint{1.867210in}{1.460716in}}%
\pgfpathlineto{\pgfqpoint{1.895332in}{1.385992in}}%
\pgfpathlineto{\pgfqpoint{1.845332in}{1.368800in}}%
\pgfpathlineto{\pgfqpoint{1.795242in}{1.351668in}}%
\pgfpathclose%
\pgfusepath{stroke,fill}%
\end{pgfscope}%
\begin{pgfscope}%
\pgfpathrectangle{\pgfqpoint{0.050000in}{0.083252in}}{\pgfqpoint{4.900000in}{4.900000in}}%
\pgfusepath{clip}%
\pgfsetbuttcap%
\pgfsetroundjoin%
\definecolor{currentfill}{rgb}{1.000000,1.000000,1.000000}%
\pgfsetfillcolor{currentfill}%
\pgfsetfillopacity{0.950000}%
\pgfsetlinewidth{0.200750pt}%
\definecolor{currentstroke}{rgb}{0.200000,0.200000,0.200000}%
\pgfsetstrokecolor{currentstroke}%
\pgfsetdash{}{0pt}%
\pgfpathmoveto{\pgfqpoint{3.242020in}{0.984667in}}%
\pgfpathlineto{\pgfqpoint{3.212403in}{1.013659in}}%
\pgfpathlineto{\pgfqpoint{3.182885in}{1.042969in}}%
\pgfpathlineto{\pgfqpoint{3.231642in}{1.059044in}}%
\pgfpathlineto{\pgfqpoint{3.280310in}{1.075101in}}%
\pgfpathlineto{\pgfqpoint{3.309919in}{1.045505in}}%
\pgfpathlineto{\pgfqpoint{3.339628in}{1.016250in}}%
\pgfpathlineto{\pgfqpoint{3.290870in}{1.000473in}}%
\pgfpathlineto{\pgfqpoint{3.242020in}{0.984667in}}%
\pgfpathclose%
\pgfusepath{stroke,fill}%
\end{pgfscope}%
\begin{pgfscope}%
\pgfpathrectangle{\pgfqpoint{0.050000in}{0.083252in}}{\pgfqpoint{4.900000in}{4.900000in}}%
\pgfusepath{clip}%
\pgfsetbuttcap%
\pgfsetroundjoin%
\definecolor{currentfill}{rgb}{0.988066,0.987574,0.992864}%
\pgfsetfillcolor{currentfill}%
\pgfsetfillopacity{0.950000}%
\pgfsetlinewidth{0.200750pt}%
\definecolor{currentstroke}{rgb}{0.200000,0.200000,0.200000}%
\pgfsetstrokecolor{currentstroke}%
\pgfsetdash{}{0pt}%
\pgfpathmoveto{\pgfqpoint{2.575179in}{1.002216in}}%
\pgfpathlineto{\pgfqpoint{2.546616in}{1.036856in}}%
\pgfpathlineto{\pgfqpoint{2.518132in}{1.074584in}}%
\pgfpathlineto{\pgfqpoint{2.567513in}{1.092333in}}%
\pgfpathlineto{\pgfqpoint{2.616806in}{1.110031in}}%
\pgfpathlineto{\pgfqpoint{2.645374in}{1.071632in}}%
\pgfpathlineto{\pgfqpoint{2.674027in}{1.036310in}}%
\pgfpathlineto{\pgfqpoint{2.624648in}{1.019267in}}%
\pgfpathlineto{\pgfqpoint{2.575179in}{1.002216in}}%
\pgfpathclose%
\pgfusepath{stroke,fill}%
\end{pgfscope}%
\begin{pgfscope}%
\pgfpathrectangle{\pgfqpoint{0.050000in}{0.083252in}}{\pgfqpoint{4.900000in}{4.900000in}}%
\pgfusepath{clip}%
\pgfsetbuttcap%
\pgfsetroundjoin%
\definecolor{currentfill}{rgb}{0.528612,0.509173,0.718139}%
\pgfsetfillcolor{currentfill}%
\pgfsetfillopacity{0.950000}%
\pgfsetlinewidth{0.200750pt}%
\definecolor{currentstroke}{rgb}{0.200000,0.200000,0.200000}%
\pgfsetstrokecolor{currentstroke}%
\pgfsetdash{}{0pt}%
\pgfpathmoveto{\pgfqpoint{1.425378in}{1.748213in}}%
\pgfpathlineto{\pgfqpoint{1.397378in}{1.825281in}}%
\pgfpathlineto{\pgfqpoint{1.369405in}{1.902274in}}%
\pgfpathlineto{\pgfqpoint{1.419888in}{1.917924in}}%
\pgfpathlineto{\pgfqpoint{1.470278in}{1.933545in}}%
\pgfpathlineto{\pgfqpoint{1.498249in}{1.856707in}}%
\pgfpathlineto{\pgfqpoint{1.526246in}{1.779796in}}%
\pgfpathlineto{\pgfqpoint{1.475858in}{1.764019in}}%
\pgfpathlineto{\pgfqpoint{1.425378in}{1.748213in}}%
\pgfpathclose%
\pgfusepath{stroke,fill}%
\end{pgfscope}%
\begin{pgfscope}%
\pgfpathrectangle{\pgfqpoint{0.050000in}{0.083252in}}{\pgfqpoint{4.900000in}{4.900000in}}%
\pgfusepath{clip}%
\pgfsetbuttcap%
\pgfsetroundjoin%
\definecolor{currentfill}{rgb}{0.904529,0.900592,0.942914}%
\pgfsetfillcolor{currentfill}%
\pgfsetfillopacity{0.950000}%
\pgfsetlinewidth{0.200750pt}%
\definecolor{currentstroke}{rgb}{0.200000,0.200000,0.200000}%
\pgfsetstrokecolor{currentstroke}%
\pgfsetdash{}{0pt}%
\pgfpathmoveto{\pgfqpoint{2.007750in}{1.115619in}}%
\pgfpathlineto{\pgfqpoint{1.979674in}{1.176732in}}%
\pgfpathlineto{\pgfqpoint{1.951574in}{1.242925in}}%
\pgfpathlineto{\pgfqpoint{2.001480in}{1.261397in}}%
\pgfpathlineto{\pgfqpoint{2.051301in}{1.279788in}}%
\pgfpathlineto{\pgfqpoint{2.079415in}{1.214676in}}%
\pgfpathlineto{\pgfqpoint{2.107523in}{1.154271in}}%
\pgfpathlineto{\pgfqpoint{2.057678in}{1.135042in}}%
\pgfpathlineto{\pgfqpoint{2.007750in}{1.115619in}}%
\pgfpathclose%
\pgfusepath{stroke,fill}%
\end{pgfscope}%
\begin{pgfscope}%
\pgfpathrectangle{\pgfqpoint{0.050000in}{0.083252in}}{\pgfqpoint{4.900000in}{4.900000in}}%
\pgfusepath{clip}%
\pgfsetbuttcap%
\pgfsetroundjoin%
\definecolor{currentfill}{rgb}{0.689719,0.676924,0.814471}%
\pgfsetfillcolor{currentfill}%
\pgfsetfillopacity{0.950000}%
\pgfsetlinewidth{0.200750pt}%
\definecolor{currentstroke}{rgb}{0.200000,0.200000,0.200000}%
\pgfsetstrokecolor{currentstroke}%
\pgfsetdash{}{0pt}%
\pgfpathmoveto{\pgfqpoint{1.638506in}{1.471485in}}%
\pgfpathlineto{\pgfqpoint{1.610401in}{1.548634in}}%
\pgfpathlineto{\pgfqpoint{1.582322in}{1.625753in}}%
\pgfpathlineto{\pgfqpoint{1.632615in}{1.641670in}}%
\pgfpathlineto{\pgfqpoint{1.682816in}{1.657574in}}%
\pgfpathlineto{\pgfqpoint{1.710890in}{1.580690in}}%
\pgfpathlineto{\pgfqpoint{1.738989in}{1.503895in}}%
\pgfpathlineto{\pgfqpoint{1.688794in}{1.487669in}}%
\pgfpathlineto{\pgfqpoint{1.638506in}{1.471485in}}%
\pgfpathclose%
\pgfusepath{stroke,fill}%
\end{pgfscope}%
\begin{pgfscope}%
\pgfpathrectangle{\pgfqpoint{0.050000in}{0.083252in}}{\pgfqpoint{4.900000in}{4.900000in}}%
\pgfusepath{clip}%
\pgfsetbuttcap%
\pgfsetroundjoin%
\definecolor{currentfill}{rgb}{0.946298,0.944083,0.967889}%
\pgfsetfillcolor{currentfill}%
\pgfsetfillopacity{0.950000}%
\pgfsetlinewidth{0.200750pt}%
\definecolor{currentstroke}{rgb}{0.200000,0.200000,0.200000}%
\pgfsetstrokecolor{currentstroke}%
\pgfsetdash{}{0pt}%
\pgfpathmoveto{\pgfqpoint{2.163731in}{1.049903in}}%
\pgfpathlineto{\pgfqpoint{2.135626in}{1.099249in}}%
\pgfpathlineto{\pgfqpoint{2.107523in}{1.154271in}}%
\pgfpathlineto{\pgfqpoint{2.157284in}{1.173329in}}%
\pgfpathlineto{\pgfqpoint{2.206961in}{1.192230in}}%
\pgfpathlineto{\pgfqpoint{2.235114in}{1.137335in}}%
\pgfpathlineto{\pgfqpoint{2.263286in}{1.087646in}}%
\pgfpathlineto{\pgfqpoint{2.213551in}{1.068842in}}%
\pgfpathlineto{\pgfqpoint{2.163731in}{1.049903in}}%
\pgfpathclose%
\pgfusepath{stroke,fill}%
\end{pgfscope}%
\begin{pgfscope}%
\pgfpathrectangle{\pgfqpoint{0.050000in}{0.083252in}}{\pgfqpoint{4.900000in}{4.900000in}}%
\pgfusepath{clip}%
\pgfsetbuttcap%
\pgfsetroundjoin%
\definecolor{currentfill}{rgb}{0.838893,0.832249,0.903668}%
\pgfsetfillcolor{currentfill}%
\pgfsetfillopacity{0.950000}%
\pgfsetlinewidth{0.200750pt}%
\definecolor{currentstroke}{rgb}{0.200000,0.200000,0.200000}%
\pgfsetstrokecolor{currentstroke}%
\pgfsetdash{}{0pt}%
\pgfpathmoveto{\pgfqpoint{1.851502in}{1.205766in}}%
\pgfpathlineto{\pgfqpoint{1.823378in}{1.277413in}}%
\pgfpathlineto{\pgfqpoint{1.795242in}{1.351668in}}%
\pgfpathlineto{\pgfqpoint{1.845332in}{1.368800in}}%
\pgfpathlineto{\pgfqpoint{1.895332in}{1.385992in}}%
\pgfpathlineto{\pgfqpoint{1.923456in}{1.313064in}}%
\pgfpathlineto{\pgfqpoint{1.951574in}{1.242925in}}%
\pgfpathlineto{\pgfqpoint{1.901581in}{1.224379in}}%
\pgfpathlineto{\pgfqpoint{1.851502in}{1.205766in}}%
\pgfpathclose%
\pgfusepath{stroke,fill}%
\end{pgfscope}%
\begin{pgfscope}%
\pgfpathrectangle{\pgfqpoint{0.050000in}{0.083252in}}{\pgfqpoint{4.900000in}{4.900000in}}%
\pgfusepath{clip}%
\pgfsetbuttcap%
\pgfsetroundjoin%
\definecolor{currentfill}{rgb}{1.000000,1.000000,1.000000}%
\pgfsetfillcolor{currentfill}%
\pgfsetfillopacity{0.950000}%
\pgfsetlinewidth{0.200750pt}%
\definecolor{currentstroke}{rgb}{0.200000,0.200000,0.200000}%
\pgfsetstrokecolor{currentstroke}%
\pgfsetdash{}{0pt}%
\pgfpathmoveto{\pgfqpoint{2.986954in}{0.978511in}}%
\pgfpathlineto{\pgfqpoint{2.957728in}{1.007941in}}%
\pgfpathlineto{\pgfqpoint{2.928602in}{1.038403in}}%
\pgfpathlineto{\pgfqpoint{2.977626in}{1.054878in}}%
\pgfpathlineto{\pgfqpoint{3.026561in}{1.071343in}}%
\pgfpathlineto{\pgfqpoint{3.055779in}{1.040502in}}%
\pgfpathlineto{\pgfqpoint{3.085101in}{1.010769in}}%
\pgfpathlineto{\pgfqpoint{3.036073in}{0.994647in}}%
\pgfpathlineto{\pgfqpoint{2.986954in}{0.978511in}}%
\pgfpathclose%
\pgfusepath{stroke,fill}%
\end{pgfscope}%
\begin{pgfscope}%
\pgfpathrectangle{\pgfqpoint{0.050000in}{0.083252in}}{\pgfqpoint{4.900000in}{4.900000in}}%
\pgfusepath{clip}%
\pgfsetbuttcap%
\pgfsetroundjoin%
\definecolor{currentfill}{rgb}{0.970165,0.968935,0.982161}%
\pgfsetfillcolor{currentfill}%
\pgfsetfillopacity{0.950000}%
\pgfsetlinewidth{0.200750pt}%
\definecolor{currentstroke}{rgb}{0.200000,0.200000,0.200000}%
\pgfsetstrokecolor{currentstroke}%
\pgfsetdash{}{0pt}%
\pgfpathmoveto{\pgfqpoint{2.319733in}{1.003102in}}%
\pgfpathlineto{\pgfqpoint{2.291488in}{1.043057in}}%
\pgfpathlineto{\pgfqpoint{2.263286in}{1.087646in}}%
\pgfpathlineto{\pgfqpoint{2.312936in}{1.106324in}}%
\pgfpathlineto{\pgfqpoint{2.362500in}{1.124886in}}%
\pgfpathlineto{\pgfqpoint{2.390777in}{1.079706in}}%
\pgfpathlineto{\pgfqpoint{2.419108in}{1.038938in}}%
\pgfpathlineto{\pgfqpoint{2.369465in}{1.021043in}}%
\pgfpathlineto{\pgfqpoint{2.319733in}{1.003102in}}%
\pgfpathclose%
\pgfusepath{stroke,fill}%
\end{pgfscope}%
\begin{pgfscope}%
\pgfpathrectangle{\pgfqpoint{0.050000in}{0.083252in}}{\pgfqpoint{4.900000in}{4.900000in}}%
\pgfusepath{clip}%
\pgfsetbuttcap%
\pgfsetroundjoin%
\definecolor{currentfill}{rgb}{0.606182,0.589942,0.764521}%
\pgfsetfillcolor{currentfill}%
\pgfsetfillopacity{0.950000}%
\pgfsetlinewidth{0.200750pt}%
\definecolor{currentstroke}{rgb}{0.200000,0.200000,0.200000}%
\pgfsetstrokecolor{currentstroke}%
\pgfsetdash{}{0pt}%
\pgfpathmoveto{\pgfqpoint{1.481459in}{1.593853in}}%
\pgfpathlineto{\pgfqpoint{1.453405in}{1.671070in}}%
\pgfpathlineto{\pgfqpoint{1.425378in}{1.748213in}}%
\pgfpathlineto{\pgfqpoint{1.475858in}{1.764019in}}%
\pgfpathlineto{\pgfqpoint{1.526246in}{1.779796in}}%
\pgfpathlineto{\pgfqpoint{1.554271in}{1.702810in}}%
\pgfpathlineto{\pgfqpoint{1.582322in}{1.625753in}}%
\pgfpathlineto{\pgfqpoint{1.531937in}{1.609816in}}%
\pgfpathlineto{\pgfqpoint{1.481459in}{1.593853in}}%
\pgfpathclose%
\pgfusepath{stroke,fill}%
\end{pgfscope}%
\begin{pgfscope}%
\pgfpathrectangle{\pgfqpoint{0.050000in}{0.083252in}}{\pgfqpoint{4.900000in}{4.900000in}}%
\pgfusepath{clip}%
\pgfsetbuttcap%
\pgfsetroundjoin%
\definecolor{currentfill}{rgb}{0.994033,0.993787,0.996432}%
\pgfsetfillcolor{currentfill}%
\pgfsetfillopacity{0.950000}%
\pgfsetlinewidth{0.200750pt}%
\definecolor{currentstroke}{rgb}{0.200000,0.200000,0.200000}%
\pgfsetstrokecolor{currentstroke}%
\pgfsetdash{}{0pt}%
\pgfpathmoveto{\pgfqpoint{2.731607in}{0.972445in}}%
\pgfpathlineto{\pgfqpoint{2.702770in}{1.003452in}}%
\pgfpathlineto{\pgfqpoint{2.674027in}{1.036310in}}%
\pgfpathlineto{\pgfqpoint{2.723317in}{1.053339in}}%
\pgfpathlineto{\pgfqpoint{2.772519in}{1.070354in}}%
\pgfpathlineto{\pgfqpoint{2.801353in}{1.036923in}}%
\pgfpathlineto{\pgfqpoint{2.830285in}{1.005431in}}%
\pgfpathlineto{\pgfqpoint{2.780991in}{0.988939in}}%
\pgfpathlineto{\pgfqpoint{2.731607in}{0.972445in}}%
\pgfpathclose%
\pgfusepath{stroke,fill}%
\end{pgfscope}%
\begin{pgfscope}%
\pgfpathrectangle{\pgfqpoint{0.050000in}{0.083252in}}{\pgfqpoint{4.900000in}{4.900000in}}%
\pgfusepath{clip}%
\pgfsetbuttcap%
\pgfsetroundjoin%
\definecolor{currentfill}{rgb}{0.767290,0.757693,0.860854}%
\pgfsetfillcolor{currentfill}%
\pgfsetfillopacity{0.950000}%
\pgfsetlinewidth{0.200750pt}%
\definecolor{currentstroke}{rgb}{0.200000,0.200000,0.200000}%
\pgfsetstrokecolor{currentstroke}%
\pgfsetdash{}{0pt}%
\pgfpathmoveto{\pgfqpoint{1.694784in}{1.317702in}}%
\pgfpathlineto{\pgfqpoint{1.666635in}{1.394415in}}%
\pgfpathlineto{\pgfqpoint{1.638506in}{1.471485in}}%
\pgfpathlineto{\pgfqpoint{1.688794in}{1.487669in}}%
\pgfpathlineto{\pgfqpoint{1.738989in}{1.503895in}}%
\pgfpathlineto{\pgfqpoint{1.767108in}{1.427411in}}%
\pgfpathlineto{\pgfqpoint{1.795242in}{1.351668in}}%
\pgfpathlineto{\pgfqpoint{1.745059in}{1.334625in}}%
\pgfpathlineto{\pgfqpoint{1.694784in}{1.317702in}}%
\pgfpathclose%
\pgfusepath{stroke,fill}%
\end{pgfscope}%
\begin{pgfscope}%
\pgfpathrectangle{\pgfqpoint{0.050000in}{0.083252in}}{\pgfqpoint{4.900000in}{4.900000in}}%
\pgfusepath{clip}%
\pgfsetbuttcap%
\pgfsetroundjoin%
\definecolor{currentfill}{rgb}{1.000000,1.000000,1.000000}%
\pgfsetfillcolor{currentfill}%
\pgfsetfillopacity{0.950000}%
\pgfsetlinewidth{0.200750pt}%
\definecolor{currentstroke}{rgb}{0.200000,0.200000,0.200000}%
\pgfsetstrokecolor{currentstroke}%
\pgfsetdash{}{0pt}%
\pgfpathmoveto{\pgfqpoint{3.144045in}{0.952968in}}%
\pgfpathlineto{\pgfqpoint{3.114524in}{0.981723in}}%
\pgfpathlineto{\pgfqpoint{3.085101in}{1.010769in}}%
\pgfpathlineto{\pgfqpoint{3.134038in}{1.026877in}}%
\pgfpathlineto{\pgfqpoint{3.182885in}{1.042969in}}%
\pgfpathlineto{\pgfqpoint{3.212403in}{1.013659in}}%
\pgfpathlineto{\pgfqpoint{3.242020in}{0.984667in}}%
\pgfpathlineto{\pgfqpoint{3.193078in}{0.968832in}}%
\pgfpathlineto{\pgfqpoint{3.144045in}{0.952968in}}%
\pgfpathclose%
\pgfusepath{stroke,fill}%
\end{pgfscope}%
\begin{pgfscope}%
\pgfpathrectangle{\pgfqpoint{0.050000in}{0.083252in}}{\pgfqpoint{4.900000in}{4.900000in}}%
\pgfusepath{clip}%
\pgfsetbuttcap%
\pgfsetroundjoin%
\definecolor{currentfill}{rgb}{0.988066,0.987574,0.992864}%
\pgfsetfillcolor{currentfill}%
\pgfsetfillopacity{0.950000}%
\pgfsetlinewidth{0.200750pt}%
\definecolor{currentstroke}{rgb}{0.200000,0.200000,0.200000}%
\pgfsetstrokecolor{currentstroke}%
\pgfsetdash{}{0pt}%
\pgfpathmoveto{\pgfqpoint{2.475972in}{0.968104in}}%
\pgfpathlineto{\pgfqpoint{2.447503in}{1.001984in}}%
\pgfpathlineto{\pgfqpoint{2.419108in}{1.038938in}}%
\pgfpathlineto{\pgfqpoint{2.468664in}{1.056785in}}%
\pgfpathlineto{\pgfqpoint{2.518132in}{1.074584in}}%
\pgfpathlineto{\pgfqpoint{2.546616in}{1.036856in}}%
\pgfpathlineto{\pgfqpoint{2.575179in}{1.002216in}}%
\pgfpathlineto{\pgfqpoint{2.525620in}{0.985160in}}%
\pgfpathlineto{\pgfqpoint{2.475972in}{0.968104in}}%
\pgfpathclose%
\pgfusepath{stroke,fill}%
\end{pgfscope}%
\begin{pgfscope}%
\pgfpathrectangle{\pgfqpoint{0.050000in}{0.083252in}}{\pgfqpoint{4.900000in}{4.900000in}}%
\pgfusepath{clip}%
\pgfsetbuttcap%
\pgfsetroundjoin%
\definecolor{currentfill}{rgb}{0.904529,0.900592,0.942914}%
\pgfsetfillcolor{currentfill}%
\pgfsetfillopacity{0.950000}%
\pgfsetlinewidth{0.200750pt}%
\definecolor{currentstroke}{rgb}{0.200000,0.200000,0.200000}%
\pgfsetstrokecolor{currentstroke}%
\pgfsetdash{}{0pt}%
\pgfpathmoveto{\pgfqpoint{1.907649in}{1.076078in}}%
\pgfpathlineto{\pgfqpoint{1.879598in}{1.138155in}}%
\pgfpathlineto{\pgfqpoint{1.851502in}{1.205766in}}%
\pgfpathlineto{\pgfqpoint{1.901581in}{1.224379in}}%
\pgfpathlineto{\pgfqpoint{1.951574in}{1.242925in}}%
\pgfpathlineto{\pgfqpoint{1.979674in}{1.176732in}}%
\pgfpathlineto{\pgfqpoint{2.007750in}{1.115619in}}%
\pgfpathlineto{\pgfqpoint{1.957740in}{1.095976in}}%
\pgfpathlineto{\pgfqpoint{1.907649in}{1.076078in}}%
\pgfpathclose%
\pgfusepath{stroke,fill}%
\end{pgfscope}%
\begin{pgfscope}%
\pgfpathrectangle{\pgfqpoint{0.050000in}{0.083252in}}{\pgfqpoint{4.900000in}{4.900000in}}%
\pgfusepath{clip}%
\pgfsetbuttcap%
\pgfsetroundjoin%
\definecolor{currentfill}{rgb}{0.952265,0.950296,0.971457}%
\pgfsetfillcolor{currentfill}%
\pgfsetfillopacity{0.950000}%
\pgfsetlinewidth{0.200750pt}%
\definecolor{currentstroke}{rgb}{0.200000,0.200000,0.200000}%
\pgfsetstrokecolor{currentstroke}%
\pgfsetdash{}{0pt}%
\pgfpathmoveto{\pgfqpoint{2.063838in}{1.011561in}}%
\pgfpathlineto{\pgfqpoint{2.035803in}{1.060445in}}%
\pgfpathlineto{\pgfqpoint{2.007750in}{1.115619in}}%
\pgfpathlineto{\pgfqpoint{2.057678in}{1.135042in}}%
\pgfpathlineto{\pgfqpoint{2.107523in}{1.154271in}}%
\pgfpathlineto{\pgfqpoint{2.135626in}{1.099249in}}%
\pgfpathlineto{\pgfqpoint{2.163731in}{1.049903in}}%
\pgfpathlineto{\pgfqpoint{2.113827in}{1.030814in}}%
\pgfpathlineto{\pgfqpoint{2.063838in}{1.011561in}}%
\pgfpathclose%
\pgfusepath{stroke,fill}%
\end{pgfscope}%
\begin{pgfscope}%
\pgfpathrectangle{\pgfqpoint{0.050000in}{0.083252in}}{\pgfqpoint{4.900000in}{4.900000in}}%
\pgfusepath{clip}%
\pgfsetbuttcap%
\pgfsetroundjoin%
\definecolor{currentfill}{rgb}{0.689719,0.676924,0.814471}%
\pgfsetfillcolor{currentfill}%
\pgfsetfillopacity{0.950000}%
\pgfsetlinewidth{0.200750pt}%
\definecolor{currentstroke}{rgb}{0.200000,0.200000,0.200000}%
\pgfsetstrokecolor{currentstroke}%
\pgfsetdash{}{0pt}%
\pgfpathmoveto{\pgfqpoint{1.537650in}{1.439194in}}%
\pgfpathlineto{\pgfqpoint{1.509541in}{1.516560in}}%
\pgfpathlineto{\pgfqpoint{1.481459in}{1.593853in}}%
\pgfpathlineto{\pgfqpoint{1.531937in}{1.609816in}}%
\pgfpathlineto{\pgfqpoint{1.582322in}{1.625753in}}%
\pgfpathlineto{\pgfqpoint{1.610401in}{1.548634in}}%
\pgfpathlineto{\pgfqpoint{1.638506in}{1.471485in}}%
\pgfpathlineto{\pgfqpoint{1.588124in}{1.455334in}}%
\pgfpathlineto{\pgfqpoint{1.537650in}{1.439194in}}%
\pgfpathclose%
\pgfusepath{stroke,fill}%
\end{pgfscope}%
\begin{pgfscope}%
\pgfpathrectangle{\pgfqpoint{0.050000in}{0.083252in}}{\pgfqpoint{4.900000in}{4.900000in}}%
\pgfusepath{clip}%
\pgfsetbuttcap%
\pgfsetroundjoin%
\definecolor{currentfill}{rgb}{0.844860,0.838462,0.907236}%
\pgfsetfillcolor{currentfill}%
\pgfsetfillopacity{0.950000}%
\pgfsetlinewidth{0.200750pt}%
\definecolor{currentstroke}{rgb}{0.200000,0.200000,0.200000}%
\pgfsetstrokecolor{currentstroke}%
\pgfsetdash{}{0pt}%
\pgfpathmoveto{\pgfqpoint{1.751085in}{1.168424in}}%
\pgfpathlineto{\pgfqpoint{1.722940in}{1.241979in}}%
\pgfpathlineto{\pgfqpoint{1.694784in}{1.317702in}}%
\pgfpathlineto{\pgfqpoint{1.745059in}{1.334625in}}%
\pgfpathlineto{\pgfqpoint{1.795242in}{1.351668in}}%
\pgfpathlineto{\pgfqpoint{1.823378in}{1.277413in}}%
\pgfpathlineto{\pgfqpoint{1.851502in}{1.205766in}}%
\pgfpathlineto{\pgfqpoint{1.801337in}{1.187103in}}%
\pgfpathlineto{\pgfqpoint{1.751085in}{1.168424in}}%
\pgfpathclose%
\pgfusepath{stroke,fill}%
\end{pgfscope}%
\begin{pgfscope}%
\pgfpathrectangle{\pgfqpoint{0.050000in}{0.083252in}}{\pgfqpoint{4.900000in}{4.900000in}}%
\pgfusepath{clip}%
\pgfsetbuttcap%
\pgfsetroundjoin%
\definecolor{currentfill}{rgb}{1.000000,1.000000,1.000000}%
\pgfsetfillcolor{currentfill}%
\pgfsetfillopacity{0.950000}%
\pgfsetlinewidth{0.200750pt}%
\definecolor{currentstroke}{rgb}{0.200000,0.200000,0.200000}%
\pgfsetstrokecolor{currentstroke}%
\pgfsetdash{}{0pt}%
\pgfpathmoveto{\pgfqpoint{2.888443in}{0.946201in}}%
\pgfpathlineto{\pgfqpoint{2.859315in}{0.975347in}}%
\pgfpathlineto{\pgfqpoint{2.830285in}{1.005431in}}%
\pgfpathlineto{\pgfqpoint{2.879489in}{1.021920in}}%
\pgfpathlineto{\pgfqpoint{2.928602in}{1.038403in}}%
\pgfpathlineto{\pgfqpoint{2.957728in}{1.007941in}}%
\pgfpathlineto{\pgfqpoint{2.986954in}{0.978511in}}%
\pgfpathlineto{\pgfqpoint{2.937745in}{0.962362in}}%
\pgfpathlineto{\pgfqpoint{2.888443in}{0.946201in}}%
\pgfpathclose%
\pgfusepath{stroke,fill}%
\end{pgfscope}%
\begin{pgfscope}%
\pgfpathrectangle{\pgfqpoint{0.050000in}{0.083252in}}{\pgfqpoint{4.900000in}{4.900000in}}%
\pgfusepath{clip}%
\pgfsetbuttcap%
\pgfsetroundjoin%
\definecolor{currentfill}{rgb}{0.976132,0.975148,0.985729}%
\pgfsetfillcolor{currentfill}%
\pgfsetfillopacity{0.950000}%
\pgfsetlinewidth{0.200750pt}%
\definecolor{currentstroke}{rgb}{0.200000,0.200000,0.200000}%
\pgfsetstrokecolor{currentstroke}%
\pgfsetdash{}{0pt}%
\pgfpathmoveto{\pgfqpoint{2.220005in}{0.967099in}}%
\pgfpathlineto{\pgfqpoint{2.191853in}{1.006058in}}%
\pgfpathlineto{\pgfqpoint{2.163731in}{1.049903in}}%
\pgfpathlineto{\pgfqpoint{2.213551in}{1.068842in}}%
\pgfpathlineto{\pgfqpoint{2.263286in}{1.087646in}}%
\pgfpathlineto{\pgfqpoint{2.291488in}{1.043057in}}%
\pgfpathlineto{\pgfqpoint{2.319733in}{1.003102in}}%
\pgfpathlineto{\pgfqpoint{2.269914in}{0.985119in}}%
\pgfpathlineto{\pgfqpoint{2.220005in}{0.967099in}}%
\pgfpathclose%
\pgfusepath{stroke,fill}%
\end{pgfscope}%
\begin{pgfscope}%
\pgfpathrectangle{\pgfqpoint{0.050000in}{0.083252in}}{\pgfqpoint{4.900000in}{4.900000in}}%
\pgfusepath{clip}%
\pgfsetbuttcap%
\pgfsetroundjoin%
\definecolor{currentfill}{rgb}{0.994033,0.993787,0.996432}%
\pgfsetfillcolor{currentfill}%
\pgfsetfillopacity{0.950000}%
\pgfsetlinewidth{0.200750pt}%
\definecolor{currentstroke}{rgb}{0.200000,0.200000,0.200000}%
\pgfsetstrokecolor{currentstroke}%
\pgfsetdash{}{0pt}%
\pgfpathmoveto{\pgfqpoint{2.632565in}{0.939467in}}%
\pgfpathlineto{\pgfqpoint{2.603827in}{0.969974in}}%
\pgfpathlineto{\pgfqpoint{2.575179in}{1.002216in}}%
\pgfpathlineto{\pgfqpoint{2.624648in}{1.019267in}}%
\pgfpathlineto{\pgfqpoint{2.674027in}{1.036310in}}%
\pgfpathlineto{\pgfqpoint{2.702770in}{1.003452in}}%
\pgfpathlineto{\pgfqpoint{2.731607in}{0.972445in}}%
\pgfpathlineto{\pgfqpoint{2.682131in}{0.955954in}}%
\pgfpathlineto{\pgfqpoint{2.632565in}{0.939467in}}%
\pgfpathclose%
\pgfusepath{stroke,fill}%
\end{pgfscope}%
\begin{pgfscope}%
\pgfpathrectangle{\pgfqpoint{0.050000in}{0.083252in}}{\pgfqpoint{4.900000in}{4.900000in}}%
\pgfusepath{clip}%
\pgfsetbuttcap%
\pgfsetroundjoin%
\definecolor{currentfill}{rgb}{0.767290,0.757693,0.860854}%
\pgfsetfillcolor{currentfill}%
\pgfsetfillopacity{0.950000}%
\pgfsetlinewidth{0.200750pt}%
\definecolor{currentstroke}{rgb}{0.200000,0.200000,0.200000}%
\pgfsetstrokecolor{currentstroke}%
\pgfsetdash{}{0pt}%
\pgfpathmoveto{\pgfqpoint{1.593947in}{1.284352in}}%
\pgfpathlineto{\pgfqpoint{1.565785in}{1.361768in}}%
\pgfpathlineto{\pgfqpoint{1.537650in}{1.439194in}}%
\pgfpathlineto{\pgfqpoint{1.588124in}{1.455334in}}%
\pgfpathlineto{\pgfqpoint{1.638506in}{1.471485in}}%
\pgfpathlineto{\pgfqpoint{1.666635in}{1.394415in}}%
\pgfpathlineto{\pgfqpoint{1.694784in}{1.317702in}}%
\pgfpathlineto{\pgfqpoint{1.644413in}{1.300936in}}%
\pgfpathlineto{\pgfqpoint{1.593947in}{1.284352in}}%
\pgfpathclose%
\pgfusepath{stroke,fill}%
\end{pgfscope}%
\begin{pgfscope}%
\pgfpathrectangle{\pgfqpoint{0.050000in}{0.083252in}}{\pgfqpoint{4.900000in}{4.900000in}}%
\pgfusepath{clip}%
\pgfsetbuttcap%
\pgfsetroundjoin%
\definecolor{currentfill}{rgb}{0.910496,0.906805,0.946482}%
\pgfsetfillcolor{currentfill}%
\pgfsetfillopacity{0.950000}%
\pgfsetlinewidth{0.200750pt}%
\definecolor{currentstroke}{rgb}{0.200000,0.200000,0.200000}%
\pgfsetstrokecolor{currentstroke}%
\pgfsetdash{}{0pt}%
\pgfpathmoveto{\pgfqpoint{1.807227in}{1.035310in}}%
\pgfpathlineto{\pgfqpoint{1.779190in}{1.098827in}}%
\pgfpathlineto{\pgfqpoint{1.751085in}{1.168424in}}%
\pgfpathlineto{\pgfqpoint{1.801337in}{1.187103in}}%
\pgfpathlineto{\pgfqpoint{1.851502in}{1.205766in}}%
\pgfpathlineto{\pgfqpoint{1.879598in}{1.138155in}}%
\pgfpathlineto{\pgfqpoint{1.907649in}{1.076078in}}%
\pgfpathlineto{\pgfqpoint{1.857477in}{1.055878in}}%
\pgfpathlineto{\pgfqpoint{1.807227in}{1.035310in}}%
\pgfpathclose%
\pgfusepath{stroke,fill}%
\end{pgfscope}%
\begin{pgfscope}%
\pgfpathrectangle{\pgfqpoint{0.050000in}{0.083252in}}{\pgfqpoint{4.900000in}{4.900000in}}%
\pgfusepath{clip}%
\pgfsetbuttcap%
\pgfsetroundjoin%
\definecolor{currentfill}{rgb}{0.988066,0.987574,0.992864}%
\pgfsetfillcolor{currentfill}%
\pgfsetfillopacity{0.950000}%
\pgfsetlinewidth{0.200750pt}%
\definecolor{currentstroke}{rgb}{0.200000,0.200000,0.200000}%
\pgfsetstrokecolor{currentstroke}%
\pgfsetdash{}{0pt}%
\pgfpathmoveto{\pgfqpoint{2.376402in}{0.934015in}}%
\pgfpathlineto{\pgfqpoint{2.348034in}{0.967047in}}%
\pgfpathlineto{\pgfqpoint{2.319733in}{1.003102in}}%
\pgfpathlineto{\pgfqpoint{2.369465in}{1.021043in}}%
\pgfpathlineto{\pgfqpoint{2.419108in}{1.038938in}}%
\pgfpathlineto{\pgfqpoint{2.447503in}{1.001984in}}%
\pgfpathlineto{\pgfqpoint{2.475972in}{0.968104in}}%
\pgfpathlineto{\pgfqpoint{2.426232in}{0.951053in}}%
\pgfpathlineto{\pgfqpoint{2.376402in}{0.934015in}}%
\pgfpathclose%
\pgfusepath{stroke,fill}%
\end{pgfscope}%
\begin{pgfscope}%
\pgfpathrectangle{\pgfqpoint{0.050000in}{0.083252in}}{\pgfqpoint{4.900000in}{4.900000in}}%
\pgfusepath{clip}%
\pgfsetbuttcap%
\pgfsetroundjoin%
\definecolor{currentfill}{rgb}{1.000000,1.000000,1.000000}%
\pgfsetfillcolor{currentfill}%
\pgfsetfillopacity{0.950000}%
\pgfsetlinewidth{0.200750pt}%
\definecolor{currentstroke}{rgb}{0.200000,0.200000,0.200000}%
\pgfsetstrokecolor{currentstroke}%
\pgfsetdash{}{0pt}%
\pgfpathmoveto{\pgfqpoint{3.045703in}{0.921153in}}%
\pgfpathlineto{\pgfqpoint{3.016280in}{0.949703in}}%
\pgfpathlineto{\pgfqpoint{2.986954in}{0.978511in}}%
\pgfpathlineto{\pgfqpoint{3.036073in}{0.994647in}}%
\pgfpathlineto{\pgfqpoint{3.085101in}{1.010769in}}%
\pgfpathlineto{\pgfqpoint{3.114524in}{0.981723in}}%
\pgfpathlineto{\pgfqpoint{3.144045in}{0.952968in}}%
\pgfpathlineto{\pgfqpoint{3.094920in}{0.937075in}}%
\pgfpathlineto{\pgfqpoint{3.045703in}{0.921153in}}%
\pgfpathclose%
\pgfusepath{stroke,fill}%
\end{pgfscope}%
\begin{pgfscope}%
\pgfpathrectangle{\pgfqpoint{0.050000in}{0.083252in}}{\pgfqpoint{4.900000in}{4.900000in}}%
\pgfusepath{clip}%
\pgfsetbuttcap%
\pgfsetroundjoin%
\definecolor{currentfill}{rgb}{0.958231,0.956509,0.975025}%
\pgfsetfillcolor{currentfill}%
\pgfsetfillopacity{0.950000}%
\pgfsetlinewidth{0.200750pt}%
\definecolor{currentstroke}{rgb}{0.200000,0.200000,0.200000}%
\pgfsetstrokecolor{currentstroke}%
\pgfsetdash{}{0pt}%
\pgfpathmoveto{\pgfqpoint{1.963608in}{0.972483in}}%
\pgfpathlineto{\pgfqpoint{1.935650in}{1.020706in}}%
\pgfpathlineto{\pgfqpoint{1.907649in}{1.076078in}}%
\pgfpathlineto{\pgfqpoint{1.957740in}{1.095976in}}%
\pgfpathlineto{\pgfqpoint{2.007750in}{1.115619in}}%
\pgfpathlineto{\pgfqpoint{2.035803in}{1.060445in}}%
\pgfpathlineto{\pgfqpoint{2.063838in}{1.011561in}}%
\pgfpathlineto{\pgfqpoint{2.013765in}{0.992125in}}%
\pgfpathlineto{\pgfqpoint{1.963608in}{0.972483in}}%
\pgfpathclose%
\pgfusepath{stroke,fill}%
\end{pgfscope}%
\begin{pgfscope}%
\pgfpathrectangle{\pgfqpoint{0.050000in}{0.083252in}}{\pgfqpoint{4.900000in}{4.900000in}}%
\pgfusepath{clip}%
\pgfsetbuttcap%
\pgfsetroundjoin%
\definecolor{currentfill}{rgb}{0.850827,0.844675,0.910804}%
\pgfsetfillcolor{currentfill}%
\pgfsetfillopacity{0.950000}%
\pgfsetlinewidth{0.200750pt}%
\definecolor{currentstroke}{rgb}{0.200000,0.200000,0.200000}%
\pgfsetstrokecolor{currentstroke}%
\pgfsetdash{}{0pt}%
\pgfpathmoveto{\pgfqpoint{1.650310in}{1.131323in}}%
\pgfpathlineto{\pgfqpoint{1.622128in}{1.207235in}}%
\pgfpathlineto{\pgfqpoint{1.593947in}{1.284352in}}%
\pgfpathlineto{\pgfqpoint{1.644413in}{1.300936in}}%
\pgfpathlineto{\pgfqpoint{1.694784in}{1.317702in}}%
\pgfpathlineto{\pgfqpoint{1.722940in}{1.241979in}}%
\pgfpathlineto{\pgfqpoint{1.751085in}{1.168424in}}%
\pgfpathlineto{\pgfqpoint{1.700744in}{1.149790in}}%
\pgfpathlineto{\pgfqpoint{1.650310in}{1.131323in}}%
\pgfpathclose%
\pgfusepath{stroke,fill}%
\end{pgfscope}%
\begin{pgfscope}%
\pgfpathrectangle{\pgfqpoint{0.050000in}{0.083252in}}{\pgfqpoint{4.900000in}{4.900000in}}%
\pgfusepath{clip}%
\pgfsetbuttcap%
\pgfsetroundjoin%
\definecolor{currentfill}{rgb}{1.000000,1.000000,1.000000}%
\pgfsetfillcolor{currentfill}%
\pgfsetfillopacity{0.950000}%
\pgfsetlinewidth{0.200750pt}%
\definecolor{currentstroke}{rgb}{0.200000,0.200000,0.200000}%
\pgfsetstrokecolor{currentstroke}%
\pgfsetdash{}{0pt}%
\pgfpathmoveto{\pgfqpoint{2.789566in}{0.913851in}}%
\pgfpathlineto{\pgfqpoint{2.760538in}{0.942734in}}%
\pgfpathlineto{\pgfqpoint{2.731607in}{0.972445in}}%
\pgfpathlineto{\pgfqpoint{2.780991in}{0.988939in}}%
\pgfpathlineto{\pgfqpoint{2.830285in}{1.005431in}}%
\pgfpathlineto{\pgfqpoint{2.859315in}{0.975347in}}%
\pgfpathlineto{\pgfqpoint{2.888443in}{0.946201in}}%
\pgfpathlineto{\pgfqpoint{2.839051in}{0.930031in}}%
\pgfpathlineto{\pgfqpoint{2.789566in}{0.913851in}}%
\pgfpathclose%
\pgfusepath{stroke,fill}%
\end{pgfscope}%
\begin{pgfscope}%
\pgfpathrectangle{\pgfqpoint{0.050000in}{0.083252in}}{\pgfqpoint{4.900000in}{4.900000in}}%
\pgfusepath{clip}%
\pgfsetbuttcap%
\pgfsetroundjoin%
\definecolor{currentfill}{rgb}{0.982099,0.981361,0.989296}%
\pgfsetfillcolor{currentfill}%
\pgfsetfillopacity{0.950000}%
\pgfsetlinewidth{0.200750pt}%
\definecolor{currentstroke}{rgb}{0.200000,0.200000,0.200000}%
\pgfsetstrokecolor{currentstroke}%
\pgfsetdash{}{0pt}%
\pgfpathmoveto{\pgfqpoint{2.119922in}{0.930977in}}%
\pgfpathlineto{\pgfqpoint{2.091871in}{0.968688in}}%
\pgfpathlineto{\pgfqpoint{2.063838in}{1.011561in}}%
\pgfpathlineto{\pgfqpoint{2.113827in}{1.030814in}}%
\pgfpathlineto{\pgfqpoint{2.163731in}{1.049903in}}%
\pgfpathlineto{\pgfqpoint{2.191853in}{1.006058in}}%
\pgfpathlineto{\pgfqpoint{2.220005in}{0.967099in}}%
\pgfpathlineto{\pgfqpoint{2.170008in}{0.949048in}}%
\pgfpathlineto{\pgfqpoint{2.119922in}{0.930977in}}%
\pgfpathclose%
\pgfusepath{stroke,fill}%
\end{pgfscope}%
\begin{pgfscope}%
\pgfpathrectangle{\pgfqpoint{0.050000in}{0.083252in}}{\pgfqpoint{4.900000in}{4.900000in}}%
\pgfusepath{clip}%
\pgfsetbuttcap%
\pgfsetroundjoin%
\definecolor{currentfill}{rgb}{1.000000,1.000000,1.000000}%
\pgfsetfillcolor{currentfill}%
\pgfsetfillopacity{0.950000}%
\pgfsetlinewidth{0.200750pt}%
\definecolor{currentstroke}{rgb}{0.200000,0.200000,0.200000}%
\pgfsetstrokecolor{currentstroke}%
\pgfsetdash{}{0pt}%
\pgfpathmoveto{\pgfqpoint{2.533156in}{0.906520in}}%
\pgfpathlineto{\pgfqpoint{2.504521in}{0.936521in}}%
\pgfpathlineto{\pgfqpoint{2.475972in}{0.968104in}}%
\pgfpathlineto{\pgfqpoint{2.525620in}{0.985160in}}%
\pgfpathlineto{\pgfqpoint{2.575179in}{1.002216in}}%
\pgfpathlineto{\pgfqpoint{2.603827in}{0.969974in}}%
\pgfpathlineto{\pgfqpoint{2.632565in}{0.939467in}}%
\pgfpathlineto{\pgfqpoint{2.582906in}{0.922988in}}%
\pgfpathlineto{\pgfqpoint{2.533156in}{0.906520in}}%
\pgfpathclose%
\pgfusepath{stroke,fill}%
\end{pgfscope}%
\begin{pgfscope}%
\pgfpathrectangle{\pgfqpoint{0.050000in}{0.083252in}}{\pgfqpoint{4.900000in}{4.900000in}}%
\pgfusepath{clip}%
\pgfsetbuttcap%
\pgfsetroundjoin%
\definecolor{currentfill}{rgb}{0.922430,0.919231,0.953618}%
\pgfsetfillcolor{currentfill}%
\pgfsetfillopacity{0.950000}%
\pgfsetlinewidth{0.200750pt}%
\definecolor{currentstroke}{rgb}{0.200000,0.200000,0.200000}%
\pgfsetstrokecolor{currentstroke}%
\pgfsetdash{}{0pt}%
\pgfpathmoveto{\pgfqpoint{1.706505in}{0.992620in}}%
\pgfpathlineto{\pgfqpoint{1.678456in}{1.058685in}}%
\pgfpathlineto{\pgfqpoint{1.650310in}{1.131323in}}%
\pgfpathlineto{\pgfqpoint{1.700744in}{1.149790in}}%
\pgfpathlineto{\pgfqpoint{1.751085in}{1.168424in}}%
\pgfpathlineto{\pgfqpoint{1.779190in}{1.098827in}}%
\pgfpathlineto{\pgfqpoint{1.807227in}{1.035310in}}%
\pgfpathlineto{\pgfqpoint{1.756902in}{1.014276in}}%
\pgfpathlineto{\pgfqpoint{1.706505in}{0.992620in}}%
\pgfpathclose%
\pgfusepath{stroke,fill}%
\end{pgfscope}%
\begin{pgfscope}%
\pgfpathrectangle{\pgfqpoint{0.050000in}{0.083252in}}{\pgfqpoint{4.900000in}{4.900000in}}%
\pgfusepath{clip}%
\pgfsetbuttcap%
\pgfsetroundjoin%
\definecolor{currentfill}{rgb}{0.964198,0.962722,0.978593}%
\pgfsetfillcolor{currentfill}%
\pgfsetfillopacity{0.950000}%
\pgfsetlinewidth{0.200750pt}%
\definecolor{currentstroke}{rgb}{0.200000,0.200000,0.200000}%
\pgfsetstrokecolor{currentstroke}%
\pgfsetdash{}{0pt}%
\pgfpathmoveto{\pgfqpoint{1.863049in}{0.932454in}}%
\pgfpathlineto{\pgfqpoint{1.835179in}{0.979644in}}%
\pgfpathlineto{\pgfqpoint{1.807227in}{1.035310in}}%
\pgfpathlineto{\pgfqpoint{1.857477in}{1.055878in}}%
\pgfpathlineto{\pgfqpoint{1.907649in}{1.076078in}}%
\pgfpathlineto{\pgfqpoint{1.935650in}{1.020706in}}%
\pgfpathlineto{\pgfqpoint{1.963608in}{0.972483in}}%
\pgfpathlineto{\pgfqpoint{1.913369in}{0.952605in}}%
\pgfpathlineto{\pgfqpoint{1.863049in}{0.932454in}}%
\pgfpathclose%
\pgfusepath{stroke,fill}%
\end{pgfscope}%
\begin{pgfscope}%
\pgfpathrectangle{\pgfqpoint{0.050000in}{0.083252in}}{\pgfqpoint{4.900000in}{4.900000in}}%
\pgfusepath{clip}%
\pgfsetbuttcap%
\pgfsetroundjoin%
\definecolor{currentfill}{rgb}{0.994033,0.993787,0.996432}%
\pgfsetfillcolor{currentfill}%
\pgfsetfillopacity{0.950000}%
\pgfsetlinewidth{0.200750pt}%
\definecolor{currentstroke}{rgb}{0.200000,0.200000,0.200000}%
\pgfsetstrokecolor{currentstroke}%
\pgfsetdash{}{0pt}%
\pgfpathmoveto{\pgfqpoint{2.276468in}{0.900009in}}%
\pgfpathlineto{\pgfqpoint{2.248206in}{0.932099in}}%
\pgfpathlineto{\pgfqpoint{2.220005in}{0.967099in}}%
\pgfpathlineto{\pgfqpoint{2.269914in}{0.985119in}}%
\pgfpathlineto{\pgfqpoint{2.319733in}{1.003102in}}%
\pgfpathlineto{\pgfqpoint{2.348034in}{0.967047in}}%
\pgfpathlineto{\pgfqpoint{2.376402in}{0.934015in}}%
\pgfpathlineto{\pgfqpoint{2.326481in}{0.916997in}}%
\pgfpathlineto{\pgfqpoint{2.276468in}{0.900009in}}%
\pgfpathclose%
\pgfusepath{stroke,fill}%
\end{pgfscope}%
\begin{pgfscope}%
\pgfpathrectangle{\pgfqpoint{0.050000in}{0.083252in}}{\pgfqpoint{4.900000in}{4.900000in}}%
\pgfusepath{clip}%
\pgfsetbuttcap%
\pgfsetroundjoin%
\definecolor{currentfill}{rgb}{1.000000,1.000000,1.000000}%
\pgfsetfillcolor{currentfill}%
\pgfsetfillopacity{0.950000}%
\pgfsetlinewidth{0.200750pt}%
\definecolor{currentstroke}{rgb}{0.200000,0.200000,0.200000}%
\pgfsetstrokecolor{currentstroke}%
\pgfsetdash{}{0pt}%
\pgfpathmoveto{\pgfqpoint{2.946992in}{0.889221in}}%
\pgfpathlineto{\pgfqpoint{2.917669in}{0.917603in}}%
\pgfpathlineto{\pgfqpoint{2.888443in}{0.946201in}}%
\pgfpathlineto{\pgfqpoint{2.937745in}{0.962362in}}%
\pgfpathlineto{\pgfqpoint{2.986954in}{0.978511in}}%
\pgfpathlineto{\pgfqpoint{3.016280in}{0.949703in}}%
\pgfpathlineto{\pgfqpoint{3.045703in}{0.921153in}}%
\pgfpathlineto{\pgfqpoint{2.996394in}{0.905201in}}%
\pgfpathlineto{\pgfqpoint{2.946992in}{0.889221in}}%
\pgfpathclose%
\pgfusepath{stroke,fill}%
\end{pgfscope}%
\begin{pgfscope}%
\pgfpathrectangle{\pgfqpoint{0.050000in}{0.083252in}}{\pgfqpoint{4.900000in}{4.900000in}}%
\pgfusepath{clip}%
\pgfsetbuttcap%
\pgfsetroundjoin%
\definecolor{currentfill}{rgb}{1.000000,1.000000,1.000000}%
\pgfsetfillcolor{currentfill}%
\pgfsetfillopacity{0.950000}%
\pgfsetlinewidth{0.200750pt}%
\definecolor{currentstroke}{rgb}{0.200000,0.200000,0.200000}%
\pgfsetstrokecolor{currentstroke}%
\pgfsetdash{}{0pt}%
\pgfpathmoveto{\pgfqpoint{2.690318in}{0.881469in}}%
\pgfpathlineto{\pgfqpoint{2.661395in}{0.910118in}}%
\pgfpathlineto{\pgfqpoint{2.632565in}{0.939467in}}%
\pgfpathlineto{\pgfqpoint{2.682131in}{0.955954in}}%
\pgfpathlineto{\pgfqpoint{2.731607in}{0.972445in}}%
\pgfpathlineto{\pgfqpoint{2.760538in}{0.942734in}}%
\pgfpathlineto{\pgfqpoint{2.789566in}{0.913851in}}%
\pgfpathlineto{\pgfqpoint{2.739988in}{0.897664in}}%
\pgfpathlineto{\pgfqpoint{2.690318in}{0.881469in}}%
\pgfpathclose%
\pgfusepath{stroke,fill}%
\end{pgfscope}%
\begin{pgfscope}%
\pgfpathrectangle{\pgfqpoint{0.050000in}{0.083252in}}{\pgfqpoint{4.900000in}{4.900000in}}%
\pgfusepath{clip}%
\pgfsetbuttcap%
\pgfsetroundjoin%
\definecolor{currentfill}{rgb}{0.988066,0.987574,0.992864}%
\pgfsetfillcolor{currentfill}%
\pgfsetfillopacity{0.950000}%
\pgfsetlinewidth{0.200750pt}%
\definecolor{currentstroke}{rgb}{0.200000,0.200000,0.200000}%
\pgfsetstrokecolor{currentstroke}%
\pgfsetdash{}{0pt}%
\pgfpathmoveto{\pgfqpoint{2.019481in}{0.894837in}}%
\pgfpathlineto{\pgfqpoint{1.991542in}{0.930941in}}%
\pgfpathlineto{\pgfqpoint{1.963608in}{0.972483in}}%
\pgfpathlineto{\pgfqpoint{2.013765in}{0.992125in}}%
\pgfpathlineto{\pgfqpoint{2.063838in}{1.011561in}}%
\pgfpathlineto{\pgfqpoint{2.091871in}{0.968688in}}%
\pgfpathlineto{\pgfqpoint{2.119922in}{0.930977in}}%
\pgfpathlineto{\pgfqpoint{2.069747in}{0.912899in}}%
\pgfpathlineto{\pgfqpoint{2.019481in}{0.894837in}}%
\pgfpathclose%
\pgfusepath{stroke,fill}%
\end{pgfscope}%
\begin{pgfscope}%
\pgfpathrectangle{\pgfqpoint{0.050000in}{0.083252in}}{\pgfqpoint{4.900000in}{4.900000in}}%
\pgfusepath{clip}%
\pgfsetbuttcap%
\pgfsetroundjoin%
\definecolor{currentfill}{rgb}{1.000000,1.000000,1.000000}%
\pgfsetfillcolor{currentfill}%
\pgfsetfillopacity{0.950000}%
\pgfsetlinewidth{0.200750pt}%
\definecolor{currentstroke}{rgb}{0.200000,0.200000,0.200000}%
\pgfsetstrokecolor{currentstroke}%
\pgfsetdash{}{0pt}%
\pgfpathmoveto{\pgfqpoint{2.433377in}{0.873630in}}%
\pgfpathlineto{\pgfqpoint{2.404848in}{0.903130in}}%
\pgfpathlineto{\pgfqpoint{2.376402in}{0.934015in}}%
\pgfpathlineto{\pgfqpoint{2.426232in}{0.951053in}}%
\pgfpathlineto{\pgfqpoint{2.475972in}{0.968104in}}%
\pgfpathlineto{\pgfqpoint{2.504521in}{0.936521in}}%
\pgfpathlineto{\pgfqpoint{2.533156in}{0.906520in}}%
\pgfpathlineto{\pgfqpoint{2.483313in}{0.890066in}}%
\pgfpathlineto{\pgfqpoint{2.433377in}{0.873630in}}%
\pgfpathclose%
\pgfusepath{stroke,fill}%
\end{pgfscope}%
\begin{pgfscope}%
\pgfpathrectangle{\pgfqpoint{0.050000in}{0.083252in}}{\pgfqpoint{4.900000in}{4.900000in}}%
\pgfusepath{clip}%
\pgfsetbuttcap%
\pgfsetroundjoin%
\definecolor{currentfill}{rgb}{0.970165,0.968935,0.982161}%
\pgfsetfillcolor{currentfill}%
\pgfsetfillopacity{0.950000}%
\pgfsetlinewidth{0.200750pt}%
\definecolor{currentstroke}{rgb}{0.200000,0.200000,0.200000}%
\pgfsetstrokecolor{currentstroke}%
\pgfsetdash{}{0pt}%
\pgfpathmoveto{\pgfqpoint{1.762169in}{0.891126in}}%
\pgfpathlineto{\pgfqpoint{1.734411in}{0.936400in}}%
\pgfpathlineto{\pgfqpoint{1.706505in}{0.992620in}}%
\pgfpathlineto{\pgfqpoint{1.756902in}{1.014276in}}%
\pgfpathlineto{\pgfqpoint{1.807227in}{1.035310in}}%
\pgfpathlineto{\pgfqpoint{1.835179in}{0.979644in}}%
\pgfpathlineto{\pgfqpoint{1.863049in}{0.932454in}}%
\pgfpathlineto{\pgfqpoint{1.812648in}{0.911979in}}%
\pgfpathlineto{\pgfqpoint{1.762169in}{0.891126in}}%
\pgfpathclose%
\pgfusepath{stroke,fill}%
\end{pgfscope}%
\begin{pgfscope}%
\pgfpathrectangle{\pgfqpoint{0.050000in}{0.083252in}}{\pgfqpoint{4.900000in}{4.900000in}}%
\pgfusepath{clip}%
\pgfsetbuttcap%
\pgfsetroundjoin%
\definecolor{currentfill}{rgb}{0.994033,0.993787,0.996432}%
\pgfsetfillcolor{currentfill}%
\pgfsetfillopacity{0.950000}%
\pgfsetlinewidth{0.200750pt}%
\definecolor{currentstroke}{rgb}{0.200000,0.200000,0.200000}%
\pgfsetstrokecolor{currentstroke}%
\pgfsetdash{}{0pt}%
\pgfpathmoveto{\pgfqpoint{2.176165in}{0.866167in}}%
\pgfpathlineto{\pgfqpoint{2.148014in}{0.897224in}}%
\pgfpathlineto{\pgfqpoint{2.119922in}{0.930977in}}%
\pgfpathlineto{\pgfqpoint{2.170008in}{0.949048in}}%
\pgfpathlineto{\pgfqpoint{2.220005in}{0.967099in}}%
\pgfpathlineto{\pgfqpoint{2.248206in}{0.932099in}}%
\pgfpathlineto{\pgfqpoint{2.276468in}{0.900009in}}%
\pgfpathlineto{\pgfqpoint{2.226363in}{0.883062in}}%
\pgfpathlineto{\pgfqpoint{2.176165in}{0.866167in}}%
\pgfpathclose%
\pgfusepath{stroke,fill}%
\end{pgfscope}%
\begin{pgfscope}%
\pgfpathrectangle{\pgfqpoint{0.050000in}{0.083252in}}{\pgfqpoint{4.900000in}{4.900000in}}%
\pgfusepath{clip}%
\pgfsetbuttcap%
\pgfsetroundjoin%
\definecolor{currentfill}{rgb}{1.000000,1.000000,1.000000}%
\pgfsetfillcolor{currentfill}%
\pgfsetfillopacity{0.950000}%
\pgfsetlinewidth{0.200750pt}%
\definecolor{currentstroke}{rgb}{0.200000,0.200000,0.200000}%
\pgfsetstrokecolor{currentstroke}%
\pgfsetdash{}{0pt}%
\pgfpathmoveto{\pgfqpoint{2.847908in}{0.857172in}}%
\pgfpathlineto{\pgfqpoint{2.818689in}{0.885428in}}%
\pgfpathlineto{\pgfqpoint{2.789566in}{0.913851in}}%
\pgfpathlineto{\pgfqpoint{2.839051in}{0.930031in}}%
\pgfpathlineto{\pgfqpoint{2.888443in}{0.946201in}}%
\pgfpathlineto{\pgfqpoint{2.917669in}{0.917603in}}%
\pgfpathlineto{\pgfqpoint{2.946992in}{0.889221in}}%
\pgfpathlineto{\pgfqpoint{2.897496in}{0.873211in}}%
\pgfpathlineto{\pgfqpoint{2.847908in}{0.857172in}}%
\pgfpathclose%
\pgfusepath{stroke,fill}%
\end{pgfscope}%
\begin{pgfscope}%
\pgfpathrectangle{\pgfqpoint{0.050000in}{0.083252in}}{\pgfqpoint{4.900000in}{4.900000in}}%
\pgfusepath{clip}%
\pgfsetbuttcap%
\pgfsetroundjoin%
\definecolor{currentfill}{rgb}{0.988066,0.987574,0.992864}%
\pgfsetfillcolor{currentfill}%
\pgfsetfillopacity{0.950000}%
\pgfsetlinewidth{0.200750pt}%
\definecolor{currentstroke}{rgb}{0.200000,0.200000,0.200000}%
\pgfsetstrokecolor{currentstroke}%
\pgfsetdash{}{0pt}%
\pgfpathmoveto{\pgfqpoint{1.918676in}{0.858912in}}%
\pgfpathlineto{\pgfqpoint{1.890866in}{0.892869in}}%
\pgfpathlineto{\pgfqpoint{1.863049in}{0.932454in}}%
\pgfpathlineto{\pgfqpoint{1.913369in}{0.952605in}}%
\pgfpathlineto{\pgfqpoint{1.963608in}{0.972483in}}%
\pgfpathlineto{\pgfqpoint{1.991542in}{0.930941in}}%
\pgfpathlineto{\pgfqpoint{2.019481in}{0.894837in}}%
\pgfpathlineto{\pgfqpoint{1.969125in}{0.876825in}}%
\pgfpathlineto{\pgfqpoint{1.918676in}{0.858912in}}%
\pgfpathclose%
\pgfusepath{stroke,fill}%
\end{pgfscope}%
\begin{pgfscope}%
\pgfpathrectangle{\pgfqpoint{0.050000in}{0.083252in}}{\pgfqpoint{4.900000in}{4.900000in}}%
\pgfusepath{clip}%
\pgfsetbuttcap%
\pgfsetroundjoin%
\definecolor{currentfill}{rgb}{1.000000,1.000000,1.000000}%
\pgfsetfillcolor{currentfill}%
\pgfsetfillopacity{0.950000}%
\pgfsetlinewidth{0.200750pt}%
\definecolor{currentstroke}{rgb}{0.200000,0.200000,0.200000}%
\pgfsetstrokecolor{currentstroke}%
\pgfsetdash{}{0pt}%
\pgfpathmoveto{\pgfqpoint{2.590699in}{0.849059in}}%
\pgfpathlineto{\pgfqpoint{2.561881in}{0.877513in}}%
\pgfpathlineto{\pgfqpoint{2.533156in}{0.906520in}}%
\pgfpathlineto{\pgfqpoint{2.582906in}{0.922988in}}%
\pgfpathlineto{\pgfqpoint{2.632565in}{0.939467in}}%
\pgfpathlineto{\pgfqpoint{2.661395in}{0.910118in}}%
\pgfpathlineto{\pgfqpoint{2.690318in}{0.881469in}}%
\pgfpathlineto{\pgfqpoint{2.640556in}{0.865267in}}%
\pgfpathlineto{\pgfqpoint{2.590699in}{0.849059in}}%
\pgfpathclose%
\pgfusepath{stroke,fill}%
\end{pgfscope}%
\begin{pgfscope}%
\pgfpathrectangle{\pgfqpoint{0.050000in}{0.083252in}}{\pgfqpoint{4.900000in}{4.900000in}}%
\pgfusepath{clip}%
\pgfsetbuttcap%
\pgfsetroundjoin%
\definecolor{currentfill}{rgb}{1.000000,1.000000,1.000000}%
\pgfsetfillcolor{currentfill}%
\pgfsetfillopacity{0.950000}%
\pgfsetlinewidth{0.200750pt}%
\definecolor{currentstroke}{rgb}{0.200000,0.200000,0.200000}%
\pgfsetstrokecolor{currentstroke}%
\pgfsetdash{}{0pt}%
\pgfpathmoveto{\pgfqpoint{2.333225in}{0.840819in}}%
\pgfpathlineto{\pgfqpoint{2.304805in}{0.869846in}}%
\pgfpathlineto{\pgfqpoint{2.276468in}{0.900009in}}%
\pgfpathlineto{\pgfqpoint{2.326481in}{0.916997in}}%
\pgfpathlineto{\pgfqpoint{2.376402in}{0.934015in}}%
\pgfpathlineto{\pgfqpoint{2.404848in}{0.903130in}}%
\pgfpathlineto{\pgfqpoint{2.433377in}{0.873630in}}%
\pgfpathlineto{\pgfqpoint{2.383348in}{0.857213in}}%
\pgfpathlineto{\pgfqpoint{2.333225in}{0.840819in}}%
\pgfpathclose%
\pgfusepath{stroke,fill}%
\end{pgfscope}%
\begin{pgfscope}%
\pgfpathrectangle{\pgfqpoint{0.050000in}{0.083252in}}{\pgfqpoint{4.900000in}{4.900000in}}%
\pgfusepath{clip}%
\pgfsetbuttcap%
\pgfsetroundjoin%
\definecolor{currentfill}{rgb}{1.000000,1.000000,1.000000}%
\pgfsetfillcolor{currentfill}%
\pgfsetfillopacity{0.950000}%
\pgfsetlinewidth{0.200750pt}%
\definecolor{currentstroke}{rgb}{0.200000,0.200000,0.200000}%
\pgfsetstrokecolor{currentstroke}%
\pgfsetdash{}{0pt}%
\pgfpathmoveto{\pgfqpoint{2.075486in}{0.832599in}}%
\pgfpathlineto{\pgfqpoint{2.047454in}{0.862563in}}%
\pgfpathlineto{\pgfqpoint{2.019481in}{0.894837in}}%
\pgfpathlineto{\pgfqpoint{2.069747in}{0.912899in}}%
\pgfpathlineto{\pgfqpoint{2.119922in}{0.930977in}}%
\pgfpathlineto{\pgfqpoint{2.148014in}{0.897224in}}%
\pgfpathlineto{\pgfqpoint{2.176165in}{0.866167in}}%
\pgfpathlineto{\pgfqpoint{2.125873in}{0.849341in}}%
\pgfpathlineto{\pgfqpoint{2.075486in}{0.832599in}}%
\pgfpathclose%
\pgfusepath{stroke,fill}%
\end{pgfscope}%
\begin{pgfscope}%
\pgfpathrectangle{\pgfqpoint{0.050000in}{0.083252in}}{\pgfqpoint{4.900000in}{4.900000in}}%
\pgfusepath{clip}%
\pgfsetbuttcap%
\pgfsetroundjoin%
\definecolor{currentfill}{rgb}{1.000000,1.000000,1.000000}%
\pgfsetfillcolor{currentfill}%
\pgfsetfillopacity{0.950000}%
\pgfsetlinewidth{0.200750pt}%
\definecolor{currentstroke}{rgb}{0.200000,0.200000,0.200000}%
\pgfsetstrokecolor{currentstroke}%
\pgfsetdash{}{0pt}%
\pgfpathmoveto{\pgfqpoint{2.748451in}{0.825005in}}%
\pgfpathlineto{\pgfqpoint{2.719337in}{0.853180in}}%
\pgfpathlineto{\pgfqpoint{2.690318in}{0.881469in}}%
\pgfpathlineto{\pgfqpoint{2.739988in}{0.897664in}}%
\pgfpathlineto{\pgfqpoint{2.789566in}{0.913851in}}%
\pgfpathlineto{\pgfqpoint{2.818689in}{0.885428in}}%
\pgfpathlineto{\pgfqpoint{2.847908in}{0.857172in}}%
\pgfpathlineto{\pgfqpoint{2.798226in}{0.841103in}}%
\pgfpathlineto{\pgfqpoint{2.748451in}{0.825005in}}%
\pgfpathclose%
\pgfusepath{stroke,fill}%
\end{pgfscope}%
\begin{pgfscope}%
\pgfpathrectangle{\pgfqpoint{0.050000in}{0.083252in}}{\pgfqpoint{4.900000in}{4.900000in}}%
\pgfusepath{clip}%
\pgfsetbuttcap%
\pgfsetroundjoin%
\definecolor{currentfill}{rgb}{0.994033,0.993787,0.996432}%
\pgfsetfillcolor{currentfill}%
\pgfsetfillopacity{0.950000}%
\pgfsetlinewidth{0.200750pt}%
\definecolor{currentstroke}{rgb}{0.200000,0.200000,0.200000}%
\pgfsetstrokecolor{currentstroke}%
\pgfsetdash{}{0pt}%
\pgfpathmoveto{\pgfqpoint{1.817492in}{0.823748in}}%
\pgfpathlineto{\pgfqpoint{1.789837in}{0.854792in}}%
\pgfpathlineto{\pgfqpoint{1.762169in}{0.891126in}}%
\pgfpathlineto{\pgfqpoint{1.812648in}{0.911979in}}%
\pgfpathlineto{\pgfqpoint{1.863049in}{0.932454in}}%
\pgfpathlineto{\pgfqpoint{1.890866in}{0.892869in}}%
\pgfpathlineto{\pgfqpoint{1.918676in}{0.858912in}}%
\pgfpathlineto{\pgfqpoint{1.868133in}{0.841179in}}%
\pgfpathlineto{\pgfqpoint{1.817492in}{0.823748in}}%
\pgfpathclose%
\pgfusepath{stroke,fill}%
\end{pgfscope}%
\begin{pgfscope}%
\pgfpathrectangle{\pgfqpoint{0.050000in}{0.083252in}}{\pgfqpoint{4.900000in}{4.900000in}}%
\pgfusepath{clip}%
\pgfsetbuttcap%
\pgfsetroundjoin%
\definecolor{currentfill}{rgb}{1.000000,1.000000,1.000000}%
\pgfsetfillcolor{currentfill}%
\pgfsetfillopacity{0.950000}%
\pgfsetlinewidth{0.200750pt}%
\definecolor{currentstroke}{rgb}{0.200000,0.200000,0.200000}%
\pgfsetstrokecolor{currentstroke}%
\pgfsetdash{}{0pt}%
\pgfpathmoveto{\pgfqpoint{2.490705in}{0.816624in}}%
\pgfpathlineto{\pgfqpoint{2.461995in}{0.844928in}}%
\pgfpathlineto{\pgfqpoint{2.433377in}{0.873630in}}%
\pgfpathlineto{\pgfqpoint{2.483313in}{0.890066in}}%
\pgfpathlineto{\pgfqpoint{2.533156in}{0.906520in}}%
\pgfpathlineto{\pgfqpoint{2.561881in}{0.877513in}}%
\pgfpathlineto{\pgfqpoint{2.590699in}{0.849059in}}%
\pgfpathlineto{\pgfqpoint{2.540749in}{0.832845in}}%
\pgfpathlineto{\pgfqpoint{2.490705in}{0.816624in}}%
\pgfpathclose%
\pgfusepath{stroke,fill}%
\end{pgfscope}%
\begin{pgfscope}%
\pgfpathrectangle{\pgfqpoint{0.050000in}{0.083252in}}{\pgfqpoint{4.900000in}{4.900000in}}%
\pgfusepath{clip}%
\pgfsetbuttcap%
\pgfsetroundjoin%
\definecolor{currentfill}{rgb}{1.000000,1.000000,1.000000}%
\pgfsetfillcolor{currentfill}%
\pgfsetfillopacity{0.950000}%
\pgfsetlinewidth{0.200750pt}%
\definecolor{currentstroke}{rgb}{0.200000,0.200000,0.200000}%
\pgfsetstrokecolor{currentstroke}%
\pgfsetdash{}{0pt}%
\pgfpathmoveto{\pgfqpoint{2.232696in}{0.808104in}}%
\pgfpathlineto{\pgfqpoint{2.204389in}{0.836722in}}%
\pgfpathlineto{\pgfqpoint{2.176165in}{0.866167in}}%
\pgfpathlineto{\pgfqpoint{2.226363in}{0.883062in}}%
\pgfpathlineto{\pgfqpoint{2.276468in}{0.900009in}}%
\pgfpathlineto{\pgfqpoint{2.304805in}{0.869846in}}%
\pgfpathlineto{\pgfqpoint{2.333225in}{0.840819in}}%
\pgfpathlineto{\pgfqpoint{2.283008in}{0.824449in}}%
\pgfpathlineto{\pgfqpoint{2.232696in}{0.808104in}}%
\pgfpathclose%
\pgfusepath{stroke,fill}%
\end{pgfscope}%
\begin{pgfscope}%
\pgfpathrectangle{\pgfqpoint{0.050000in}{0.083252in}}{\pgfqpoint{4.900000in}{4.900000in}}%
\pgfusepath{clip}%
\pgfsetbuttcap%
\pgfsetroundjoin%
\definecolor{currentfill}{rgb}{1.000000,1.000000,1.000000}%
\pgfsetfillcolor{currentfill}%
\pgfsetfillopacity{0.950000}%
\pgfsetlinewidth{0.200750pt}%
\definecolor{currentstroke}{rgb}{0.200000,0.200000,0.200000}%
\pgfsetstrokecolor{currentstroke}%
\pgfsetdash{}{0pt}%
\pgfpathmoveto{\pgfqpoint{1.974427in}{0.799427in}}%
\pgfpathlineto{\pgfqpoint{1.946519in}{0.828350in}}%
\pgfpathlineto{\pgfqpoint{1.918676in}{0.858912in}}%
\pgfpathlineto{\pgfqpoint{1.969125in}{0.876825in}}%
\pgfpathlineto{\pgfqpoint{2.019481in}{0.894837in}}%
\pgfpathlineto{\pgfqpoint{2.047454in}{0.862563in}}%
\pgfpathlineto{\pgfqpoint{2.075486in}{0.832599in}}%
\pgfpathlineto{\pgfqpoint{2.025005in}{0.815956in}}%
\pgfpathlineto{\pgfqpoint{1.974427in}{0.799427in}}%
\pgfpathclose%
\pgfusepath{stroke,fill}%
\end{pgfscope}%
\begin{pgfscope}%
\pgfpathrectangle{\pgfqpoint{0.050000in}{0.083252in}}{\pgfqpoint{4.900000in}{4.900000in}}%
\pgfusepath{clip}%
\pgfsetbuttcap%
\pgfsetroundjoin%
\definecolor{currentfill}{rgb}{1.000000,1.000000,1.000000}%
\pgfsetfillcolor{currentfill}%
\pgfsetfillopacity{0.950000}%
\pgfsetlinewidth{0.200750pt}%
\definecolor{currentstroke}{rgb}{0.200000,0.200000,0.200000}%
\pgfsetstrokecolor{currentstroke}%
\pgfsetdash{}{0pt}%
\pgfpathmoveto{\pgfqpoint{2.648617in}{0.792720in}}%
\pgfpathlineto{\pgfqpoint{2.619611in}{0.820861in}}%
\pgfpathlineto{\pgfqpoint{2.590699in}{0.849059in}}%
\pgfpathlineto{\pgfqpoint{2.640556in}{0.865267in}}%
\pgfpathlineto{\pgfqpoint{2.690318in}{0.881469in}}%
\pgfpathlineto{\pgfqpoint{2.719337in}{0.853180in}}%
\pgfpathlineto{\pgfqpoint{2.748451in}{0.825005in}}%
\pgfpathlineto{\pgfqpoint{2.698581in}{0.808878in}}%
\pgfpathlineto{\pgfqpoint{2.648617in}{0.792720in}}%
\pgfpathclose%
\pgfusepath{stroke,fill}%
\end{pgfscope}%
\begin{pgfscope}%
\pgfpathrectangle{\pgfqpoint{0.050000in}{0.083252in}}{\pgfqpoint{4.900000in}{4.900000in}}%
\pgfusepath{clip}%
\pgfsetbuttcap%
\pgfsetroundjoin%
\definecolor{currentfill}{rgb}{1.000000,1.000000,1.000000}%
\pgfsetfillcolor{currentfill}%
\pgfsetfillopacity{0.950000}%
\pgfsetlinewidth{0.200750pt}%
\definecolor{currentstroke}{rgb}{0.200000,0.200000,0.200000}%
\pgfsetstrokecolor{currentstroke}%
\pgfsetdash{}{0pt}%
\pgfpathmoveto{\pgfqpoint{2.390333in}{0.784158in}}%
\pgfpathlineto{\pgfqpoint{2.361733in}{0.812370in}}%
\pgfpathlineto{\pgfqpoint{2.333225in}{0.840819in}}%
\pgfpathlineto{\pgfqpoint{2.383348in}{0.857213in}}%
\pgfpathlineto{\pgfqpoint{2.433377in}{0.873630in}}%
\pgfpathlineto{\pgfqpoint{2.461995in}{0.844928in}}%
\pgfpathlineto{\pgfqpoint{2.490705in}{0.816624in}}%
\pgfpathlineto{\pgfqpoint{2.440566in}{0.800396in}}%
\pgfpathlineto{\pgfqpoint{2.390333in}{0.784158in}}%
\pgfpathclose%
\pgfusepath{stroke,fill}%
\end{pgfscope}%
\begin{pgfscope}%
\pgfpathrectangle{\pgfqpoint{0.050000in}{0.083252in}}{\pgfqpoint{4.900000in}{4.900000in}}%
\pgfusepath{clip}%
\pgfsetbuttcap%
\pgfsetroundjoin%
\definecolor{currentfill}{rgb}{1.000000,1.000000,1.000000}%
\pgfsetfillcolor{currentfill}%
\pgfsetfillopacity{0.950000}%
\pgfsetlinewidth{0.200750pt}%
\definecolor{currentstroke}{rgb}{0.200000,0.200000,0.200000}%
\pgfsetstrokecolor{currentstroke}%
\pgfsetdash{}{0pt}%
\pgfpathmoveto{\pgfqpoint{2.131787in}{0.775483in}}%
\pgfpathlineto{\pgfqpoint{2.103594in}{0.803804in}}%
\pgfpathlineto{\pgfqpoint{2.075486in}{0.832599in}}%
\pgfpathlineto{\pgfqpoint{2.125873in}{0.849341in}}%
\pgfpathlineto{\pgfqpoint{2.176165in}{0.866167in}}%
\pgfpathlineto{\pgfqpoint{2.204389in}{0.836722in}}%
\pgfpathlineto{\pgfqpoint{2.232696in}{0.808104in}}%
\pgfpathlineto{\pgfqpoint{2.182290in}{0.791783in}}%
\pgfpathlineto{\pgfqpoint{2.131787in}{0.775483in}}%
\pgfpathclose%
\pgfusepath{stroke,fill}%
\end{pgfscope}%
\begin{pgfscope}%
\pgfpathrectangle{\pgfqpoint{0.050000in}{0.083252in}}{\pgfqpoint{4.900000in}{4.900000in}}%
\pgfusepath{clip}%
\pgfsetbuttcap%
\pgfsetroundjoin%
\definecolor{currentfill}{rgb}{1.000000,1.000000,1.000000}%
\pgfsetfillcolor{currentfill}%
\pgfsetfillopacity{0.950000}%
\pgfsetlinewidth{0.200750pt}%
\definecolor{currentstroke}{rgb}{0.200000,0.200000,0.200000}%
\pgfsetstrokecolor{currentstroke}%
\pgfsetdash{}{0pt}%
\pgfpathmoveto{\pgfqpoint{1.872980in}{0.766703in}}%
\pgfpathlineto{\pgfqpoint{1.845197in}{0.794920in}}%
\pgfpathlineto{\pgfqpoint{1.817492in}{0.823748in}}%
\pgfpathlineto{\pgfqpoint{1.868133in}{0.841179in}}%
\pgfpathlineto{\pgfqpoint{1.918676in}{0.858912in}}%
\pgfpathlineto{\pgfqpoint{1.946519in}{0.828350in}}%
\pgfpathlineto{\pgfqpoint{1.974427in}{0.799427in}}%
\pgfpathlineto{\pgfqpoint{1.923752in}{0.783015in}}%
\pgfpathlineto{\pgfqpoint{1.872980in}{0.766703in}}%
\pgfpathclose%
\pgfusepath{stroke,fill}%
\end{pgfscope}%
\begin{pgfscope}%
\pgfpathrectangle{\pgfqpoint{0.050000in}{0.083252in}}{\pgfqpoint{4.900000in}{4.900000in}}%
\pgfusepath{clip}%
\pgfsetbuttcap%
\pgfsetroundjoin%
\definecolor{currentfill}{rgb}{1.000000,1.000000,1.000000}%
\pgfsetfillcolor{currentfill}%
\pgfsetfillopacity{0.950000}%
\pgfsetlinewidth{0.200750pt}%
\definecolor{currentstroke}{rgb}{0.200000,0.200000,0.200000}%
\pgfsetstrokecolor{currentstroke}%
\pgfsetdash{}{0pt}%
\pgfpathmoveto{\pgfqpoint{2.548405in}{0.760317in}}%
\pgfpathlineto{\pgfqpoint{2.519508in}{0.788468in}}%
\pgfpathlineto{\pgfqpoint{2.490705in}{0.816624in}}%
\pgfpathlineto{\pgfqpoint{2.540749in}{0.832845in}}%
\pgfpathlineto{\pgfqpoint{2.590699in}{0.849059in}}%
\pgfpathlineto{\pgfqpoint{2.619611in}{0.820861in}}%
\pgfpathlineto{\pgfqpoint{2.648617in}{0.792720in}}%
\pgfpathlineto{\pgfqpoint{2.598558in}{0.776534in}}%
\pgfpathlineto{\pgfqpoint{2.548405in}{0.760317in}}%
\pgfpathclose%
\pgfusepath{stroke,fill}%
\end{pgfscope}%
\begin{pgfscope}%
\pgfpathrectangle{\pgfqpoint{0.050000in}{0.083252in}}{\pgfqpoint{4.900000in}{4.900000in}}%
\pgfusepath{clip}%
\pgfsetbuttcap%
\pgfsetroundjoin%
\definecolor{currentfill}{rgb}{1.000000,1.000000,1.000000}%
\pgfsetfillcolor{currentfill}%
\pgfsetfillopacity{0.950000}%
\pgfsetlinewidth{0.200750pt}%
\definecolor{currentstroke}{rgb}{0.200000,0.200000,0.200000}%
\pgfsetstrokecolor{currentstroke}%
\pgfsetdash{}{0pt}%
\pgfpathmoveto{\pgfqpoint{2.289581in}{0.751646in}}%
\pgfpathlineto{\pgfqpoint{2.261093in}{0.779831in}}%
\pgfpathlineto{\pgfqpoint{2.232696in}{0.808104in}}%
\pgfpathlineto{\pgfqpoint{2.283008in}{0.824449in}}%
\pgfpathlineto{\pgfqpoint{2.333225in}{0.840819in}}%
\pgfpathlineto{\pgfqpoint{2.361733in}{0.812370in}}%
\pgfpathlineto{\pgfqpoint{2.390333in}{0.784158in}}%
\pgfpathlineto{\pgfqpoint{2.340005in}{0.767909in}}%
\pgfpathlineto{\pgfqpoint{2.289581in}{0.751646in}}%
\pgfpathclose%
\pgfusepath{stroke,fill}%
\end{pgfscope}%
\begin{pgfscope}%
\pgfpathrectangle{\pgfqpoint{0.050000in}{0.083252in}}{\pgfqpoint{4.900000in}{4.900000in}}%
\pgfusepath{clip}%
\pgfsetbuttcap%
\pgfsetroundjoin%
\definecolor{currentfill}{rgb}{1.000000,1.000000,1.000000}%
\pgfsetfillcolor{currentfill}%
\pgfsetfillopacity{0.950000}%
\pgfsetlinewidth{0.200750pt}%
\definecolor{currentstroke}{rgb}{0.200000,0.200000,0.200000}%
\pgfsetstrokecolor{currentstroke}%
\pgfsetdash{}{0pt}%
\pgfpathmoveto{\pgfqpoint{2.030494in}{0.742915in}}%
\pgfpathlineto{\pgfqpoint{2.002416in}{0.771106in}}%
\pgfpathlineto{\pgfqpoint{1.974427in}{0.799427in}}%
\pgfpathlineto{\pgfqpoint{2.025005in}{0.815956in}}%
\pgfpathlineto{\pgfqpoint{2.075486in}{0.832599in}}%
\pgfpathlineto{\pgfqpoint{2.103594in}{0.803804in}}%
\pgfpathlineto{\pgfqpoint{2.131787in}{0.775483in}}%
\pgfpathlineto{\pgfqpoint{2.081189in}{0.759197in}}%
\pgfpathlineto{\pgfqpoint{2.030494in}{0.742915in}}%
\pgfpathclose%
\pgfusepath{stroke,fill}%
\end{pgfscope}%
\begin{pgfscope}%
\pgfpathrectangle{\pgfqpoint{0.050000in}{0.083252in}}{\pgfqpoint{4.900000in}{4.900000in}}%
\pgfusepath{clip}%
\pgfsetbuttcap%
\pgfsetroundjoin%
\definecolor{currentfill}{rgb}{1.000000,1.000000,1.000000}%
\pgfsetfillcolor{currentfill}%
\pgfsetfillopacity{0.950000}%
\pgfsetlinewidth{0.200750pt}%
\definecolor{currentstroke}{rgb}{0.200000,0.200000,0.200000}%
\pgfsetstrokecolor{currentstroke}%
\pgfsetdash{}{0pt}%
\pgfpathmoveto{\pgfqpoint{2.447813in}{0.727793in}}%
\pgfpathlineto{\pgfqpoint{2.419026in}{0.755997in}}%
\pgfpathlineto{\pgfqpoint{2.390333in}{0.784158in}}%
\pgfpathlineto{\pgfqpoint{2.440566in}{0.800396in}}%
\pgfpathlineto{\pgfqpoint{2.490705in}{0.816624in}}%
\pgfpathlineto{\pgfqpoint{2.519508in}{0.788468in}}%
\pgfpathlineto{\pgfqpoint{2.548405in}{0.760317in}}%
\pgfpathlineto{\pgfqpoint{2.498157in}{0.744070in}}%
\pgfpathlineto{\pgfqpoint{2.447813in}{0.727793in}}%
\pgfpathclose%
\pgfusepath{stroke,fill}%
\end{pgfscope}%
\begin{pgfscope}%
\pgfpathrectangle{\pgfqpoint{0.050000in}{0.083252in}}{\pgfqpoint{4.900000in}{4.900000in}}%
\pgfusepath{clip}%
\pgfsetbuttcap%
\pgfsetroundjoin%
\definecolor{currentfill}{rgb}{1.000000,1.000000,1.000000}%
\pgfsetfillcolor{currentfill}%
\pgfsetfillopacity{0.950000}%
\pgfsetlinewidth{0.200750pt}%
\definecolor{currentstroke}{rgb}{0.200000,0.200000,0.200000}%
\pgfsetstrokecolor{currentstroke}%
\pgfsetdash{}{0pt}%
\pgfpathmoveto{\pgfqpoint{2.188446in}{0.719063in}}%
\pgfpathlineto{\pgfqpoint{2.160070in}{0.747285in}}%
\pgfpathlineto{\pgfqpoint{2.131787in}{0.775483in}}%
\pgfpathlineto{\pgfqpoint{2.182290in}{0.791783in}}%
\pgfpathlineto{\pgfqpoint{2.232696in}{0.808104in}}%
\pgfpathlineto{\pgfqpoint{2.261093in}{0.779831in}}%
\pgfpathlineto{\pgfqpoint{2.289581in}{0.751646in}}%
\pgfpathlineto{\pgfqpoint{2.239062in}{0.735365in}}%
\pgfpathlineto{\pgfqpoint{2.188446in}{0.719063in}}%
\pgfpathclose%
\pgfusepath{stroke,fill}%
\end{pgfscope}%
\begin{pgfscope}%
\pgfpathrectangle{\pgfqpoint{0.050000in}{0.083252in}}{\pgfqpoint{4.900000in}{4.900000in}}%
\pgfusepath{clip}%
\pgfsetbuttcap%
\pgfsetroundjoin%
\definecolor{currentfill}{rgb}{1.000000,1.000000,1.000000}%
\pgfsetfillcolor{currentfill}%
\pgfsetfillopacity{0.950000}%
\pgfsetlinewidth{0.200750pt}%
\definecolor{currentstroke}{rgb}{0.200000,0.200000,0.200000}%
\pgfsetstrokecolor{currentstroke}%
\pgfsetdash{}{0pt}%
\pgfpathmoveto{\pgfqpoint{1.928815in}{0.710313in}}%
\pgfpathlineto{\pgfqpoint{1.900852in}{0.738543in}}%
\pgfpathlineto{\pgfqpoint{1.872980in}{0.766703in}}%
\pgfpathlineto{\pgfqpoint{1.923752in}{0.783015in}}%
\pgfpathlineto{\pgfqpoint{1.974427in}{0.799427in}}%
\pgfpathlineto{\pgfqpoint{2.002416in}{0.771106in}}%
\pgfpathlineto{\pgfqpoint{2.030494in}{0.742915in}}%
\pgfpathlineto{\pgfqpoint{1.979703in}{0.726625in}}%
\pgfpathlineto{\pgfqpoint{1.928815in}{0.710313in}}%
\pgfpathclose%
\pgfusepath{stroke,fill}%
\end{pgfscope}%
\begin{pgfscope}%
\pgfpathrectangle{\pgfqpoint{0.050000in}{0.083252in}}{\pgfqpoint{4.900000in}{4.900000in}}%
\pgfusepath{clip}%
\pgfsetbuttcap%
\pgfsetroundjoin%
\definecolor{currentfill}{rgb}{1.000000,1.000000,1.000000}%
\pgfsetfillcolor{currentfill}%
\pgfsetfillopacity{0.950000}%
\pgfsetlinewidth{0.200750pt}%
\definecolor{currentstroke}{rgb}{0.200000,0.200000,0.200000}%
\pgfsetstrokecolor{currentstroke}%
\pgfsetdash{}{0pt}%
\pgfpathmoveto{\pgfqpoint{2.346838in}{0.695148in}}%
\pgfpathlineto{\pgfqpoint{2.318162in}{0.723434in}}%
\pgfpathlineto{\pgfqpoint{2.289581in}{0.751646in}}%
\pgfpathlineto{\pgfqpoint{2.340005in}{0.767909in}}%
\pgfpathlineto{\pgfqpoint{2.390333in}{0.784158in}}%
\pgfpathlineto{\pgfqpoint{2.419026in}{0.755997in}}%
\pgfpathlineto{\pgfqpoint{2.447813in}{0.727793in}}%
\pgfpathlineto{\pgfqpoint{2.397373in}{0.711486in}}%
\pgfpathlineto{\pgfqpoint{2.346838in}{0.695148in}}%
\pgfpathclose%
\pgfusepath{stroke,fill}%
\end{pgfscope}%
\begin{pgfscope}%
\pgfpathrectangle{\pgfqpoint{0.050000in}{0.083252in}}{\pgfqpoint{4.900000in}{4.900000in}}%
\pgfusepath{clip}%
\pgfsetbuttcap%
\pgfsetroundjoin%
\definecolor{currentfill}{rgb}{1.000000,1.000000,1.000000}%
\pgfsetfillcolor{currentfill}%
\pgfsetfillopacity{0.950000}%
\pgfsetlinewidth{0.200750pt}%
\definecolor{currentstroke}{rgb}{0.200000,0.200000,0.200000}%
\pgfsetstrokecolor{currentstroke}%
\pgfsetdash{}{0pt}%
\pgfpathmoveto{\pgfqpoint{2.086926in}{0.686380in}}%
\pgfpathlineto{\pgfqpoint{2.058664in}{0.714688in}}%
\pgfpathlineto{\pgfqpoint{2.030494in}{0.742915in}}%
\pgfpathlineto{\pgfqpoint{2.081189in}{0.759197in}}%
\pgfpathlineto{\pgfqpoint{2.131787in}{0.775483in}}%
\pgfpathlineto{\pgfqpoint{2.160070in}{0.747285in}}%
\pgfpathlineto{\pgfqpoint{2.188446in}{0.719063in}}%
\pgfpathlineto{\pgfqpoint{2.137734in}{0.702736in}}%
\pgfpathlineto{\pgfqpoint{2.086926in}{0.686380in}}%
\pgfpathclose%
\pgfusepath{stroke,fill}%
\end{pgfscope}%
\begin{pgfscope}%
\pgfpathrectangle{\pgfqpoint{0.050000in}{0.083252in}}{\pgfqpoint{4.900000in}{4.900000in}}%
\pgfusepath{clip}%
\pgfsetbuttcap%
\pgfsetroundjoin%
\definecolor{currentfill}{rgb}{1.000000,1.000000,1.000000}%
\pgfsetfillcolor{currentfill}%
\pgfsetfillopacity{0.950000}%
\pgfsetlinewidth{0.200750pt}%
\definecolor{currentstroke}{rgb}{0.200000,0.200000,0.200000}%
\pgfsetstrokecolor{currentstroke}%
\pgfsetdash{}{0pt}%
\pgfpathmoveto{\pgfqpoint{2.245478in}{0.662380in}}%
\pgfpathlineto{\pgfqpoint{2.216915in}{0.690766in}}%
\pgfpathlineto{\pgfqpoint{2.188446in}{0.719063in}}%
\pgfpathlineto{\pgfqpoint{2.239062in}{0.735365in}}%
\pgfpathlineto{\pgfqpoint{2.289581in}{0.751646in}}%
\pgfpathlineto{\pgfqpoint{2.318162in}{0.723434in}}%
\pgfpathlineto{\pgfqpoint{2.346838in}{0.695148in}}%
\pgfpathlineto{\pgfqpoint{2.296206in}{0.678780in}}%
\pgfpathlineto{\pgfqpoint{2.245478in}{0.662380in}}%
\pgfpathclose%
\pgfusepath{stroke,fill}%
\end{pgfscope}%
\begin{pgfscope}%
\pgfpathrectangle{\pgfqpoint{0.050000in}{0.083252in}}{\pgfqpoint{4.900000in}{4.900000in}}%
\pgfusepath{clip}%
\pgfsetbuttcap%
\pgfsetroundjoin%
\definecolor{currentfill}{rgb}{1.000000,1.000000,1.000000}%
\pgfsetfillcolor{currentfill}%
\pgfsetfillopacity{0.950000}%
\pgfsetlinewidth{0.200750pt}%
\definecolor{currentstroke}{rgb}{0.200000,0.200000,0.200000}%
\pgfsetstrokecolor{currentstroke}%
\pgfsetdash{}{0pt}%
\pgfpathmoveto{\pgfqpoint{1.985019in}{0.653578in}}%
\pgfpathlineto{\pgfqpoint{1.956871in}{0.681993in}}%
\pgfpathlineto{\pgfqpoint{1.928815in}{0.710313in}}%
\pgfpathlineto{\pgfqpoint{1.979703in}{0.726625in}}%
\pgfpathlineto{\pgfqpoint{2.030494in}{0.742915in}}%
\pgfpathlineto{\pgfqpoint{2.058664in}{0.714688in}}%
\pgfpathlineto{\pgfqpoint{2.086926in}{0.686380in}}%
\pgfpathlineto{\pgfqpoint{2.036021in}{0.669995in}}%
\pgfpathlineto{\pgfqpoint{1.985019in}{0.653578in}}%
\pgfpathclose%
\pgfusepath{stroke,fill}%
\end{pgfscope}%
\begin{pgfscope}%
\pgfpathrectangle{\pgfqpoint{0.050000in}{0.083252in}}{\pgfqpoint{4.900000in}{4.900000in}}%
\pgfusepath{clip}%
\pgfsetbuttcap%
\pgfsetroundjoin%
\definecolor{currentfill}{rgb}{1.000000,1.000000,1.000000}%
\pgfsetfillcolor{currentfill}%
\pgfsetfillopacity{0.950000}%
\pgfsetlinewidth{0.200750pt}%
\definecolor{currentstroke}{rgb}{0.200000,0.200000,0.200000}%
\pgfsetstrokecolor{currentstroke}%
\pgfsetdash{}{0pt}%
\pgfpathmoveto{\pgfqpoint{2.143731in}{0.629487in}}%
\pgfpathlineto{\pgfqpoint{2.115282in}{0.657980in}}%
\pgfpathlineto{\pgfqpoint{2.086926in}{0.686380in}}%
\pgfpathlineto{\pgfqpoint{2.137734in}{0.702736in}}%
\pgfpathlineto{\pgfqpoint{2.188446in}{0.719063in}}%
\pgfpathlineto{\pgfqpoint{2.216915in}{0.690766in}}%
\pgfpathlineto{\pgfqpoint{2.245478in}{0.662380in}}%
\pgfpathlineto{\pgfqpoint{2.194653in}{0.645949in}}%
\pgfpathlineto{\pgfqpoint{2.143731in}{0.629487in}}%
\pgfpathclose%
\pgfusepath{stroke,fill}%
\end{pgfscope}%
\begin{pgfscope}%
\pgfpathrectangle{\pgfqpoint{0.050000in}{0.083252in}}{\pgfqpoint{4.900000in}{4.900000in}}%
\pgfusepath{clip}%
\pgfsetbuttcap%
\pgfsetroundjoin%
\definecolor{currentfill}{rgb}{1.000000,1.000000,1.000000}%
\pgfsetfillcolor{currentfill}%
\pgfsetfillopacity{0.950000}%
\pgfsetlinewidth{0.200750pt}%
\definecolor{currentstroke}{rgb}{0.200000,0.200000,0.200000}%
\pgfsetstrokecolor{currentstroke}%
\pgfsetdash{}{0pt}%
\pgfpathmoveto{\pgfqpoint{2.041595in}{0.596467in}}%
\pgfpathlineto{\pgfqpoint{2.013260in}{0.625070in}}%
\pgfpathlineto{\pgfqpoint{1.985019in}{0.653578in}}%
\pgfpathlineto{\pgfqpoint{2.036021in}{0.669995in}}%
\pgfpathlineto{\pgfqpoint{2.086926in}{0.686380in}}%
\pgfpathlineto{\pgfqpoint{2.115282in}{0.657980in}}%
\pgfpathlineto{\pgfqpoint{2.143731in}{0.629487in}}%
\pgfpathlineto{\pgfqpoint{2.092712in}{0.612993in}}%
\pgfpathlineto{\pgfqpoint{2.041595in}{0.596467in}}%
\pgfpathclose%
\pgfusepath{stroke,fill}%
\end{pgfscope}%
\end{pgfpicture}%
\makeatother%
\endgroup%

	\caption{Graph of the down-and-out put option price}\label{fig:downAndOutPutBarrierSurface}
\end{figure}

\clearpage
\subsection{Knock-in Barrier}
